\paragraph{纤维代数、碰撞统计与零温 primitive 谱的统一热力学接口}
在 bin-fold 的纤维群胚代数、折叠同余商半环、碰撞统计、规范群表示谱以及零温四态极限之间，存在一条可由现有对象直接闭合的接口：有限尺度上的迹态极值由纤维多重度决定，折叠同余商的局部奇异性由 Fibonacci 模数的平方因子完全控制，二次碰撞矩正好等于相应半单代数的归一化维数，规范群的不可约表示按纤维分拆完全分离，而零温 primitive 轨道与 fold-gauge 常数的 Bernoulli 层级又分别把 Perron 谱、\(\,2\pi\) 的上侧逼近以及偶 \(\zeta\)-值的递阶恢复固定在同一条热力学账本上。

\begin{theorem}[纤维群胚代数上的最大熵迹态唯一由均匀微态推前给出]\label{thm:xi-foldbin-groupoid-maxent-tracial-state}
\leanverified{paper\_xi\_foldbin\_groupoid\_maxent\_tracial\_state}
令
\[
A_m^{\mathrm{bin}}:=\CC[\mathcal R_m^{\mathrm{bin}}]
\cong
\bigoplus_{w\in X_m} M_{d_m^{\mathrm{bin}}(w)}(\CC),
\qquad
\mu_m(w):=\frac{d_m^{\mathrm{bin}}(w)}{2^m}.
\]
对任意 \(\pi\in\Delta(X_m)\)，定义
\[
\tau_\pi\!\left(\bigoplus_{w\in X_m}A_w\right)
:=
\sum_{w\in X_m}\pi(w)\,\operatorname{tr}_{d_m^{\mathrm{bin}}(w)}(A_w),
\]
其中 \(\operatorname{tr}_d\) 为 \(M_d(\CC)\) 上的归一化矩阵迹。则 \(A_m^{\mathrm{bin}}\) 的全部迹态恰为 \(\{\tau_\pi:\pi\in\Delta(X_m)\}\)，并且其 von Neumann 熵满足
\[
S(\tau_\pi)
=
H(\pi)+\sum_{w\in X_m}\pi(w)\log d_m^{\mathrm{bin}}(w),
\qquad
H(\pi):=-\sum_{w\in X_m}\pi(w)\log\pi(w).
\]
此外，
\[
m\log 2-S(\tau_\pi)=D_{\mathrm{KL}}(\pi\|\mu_m).
\]
特别地，
\[
S(\tau_{\mu_m})=m\log 2,
\]
且 \(\tau_{\mu_m}\) 是 \(A_m^{\mathrm{bin}}\) 上唯一的最大熵迹态。
\end{theorem}
\begin{proof}
由命题 \ref{prop:fold-bin-groupoid-algebra-wedderburn}，
\[
A_m^{\mathrm{bin}}
\cong
\bigoplus_{w\in X_m} M_{d_m^{\mathrm{bin}}(w)}(\CC).
\]
每个矩阵块上的迹态唯一等于归一化矩阵迹，故有限直和代数上的任意迹态都唯一写成上述凸组合；其系数恰构成一个概率向量 \(\pi\in\Delta(X_m)\)。
\par
对总代数上的普通迹
\(
\Tr:=\sum_{w\in X_m}\Tr_{d_m^{\mathrm{bin}}(w)}
\)，
状态 \(\tau_\pi\) 的密度矩阵为
\[
\rho_\pi
=
\bigoplus_{w\in X_m}
\frac{\pi(w)}{d_m^{\mathrm{bin}}(w)}I_{d_m^{\mathrm{bin}}(w)}.
\]
于是
\[
\begin{aligned}
S(\tau_\pi)
&=-\Tr(\rho_\pi\log\rho_\pi)\\
&=
-\sum_{w\in X_m}d_m^{\mathrm{bin}}(w)\,
\frac{\pi(w)}{d_m^{\mathrm{bin}}(w)}
\log\!\left(\frac{\pi(w)}{d_m^{\mathrm{bin}}(w)}\right)\\
&=
H(\pi)+\sum_{w\in X_m}\pi(w)\log d_m^{\mathrm{bin}}(w).
\end{aligned}
\]
另一方面，
\[
\begin{aligned}
D_{\mathrm{KL}}(\pi\|\mu_m)
&=
\sum_{w\in X_m}\pi(w)
\log\frac{\pi(w)}{d_m^{\mathrm{bin}}(w)/2^m}\\
&=
-H(\pi)-\sum_{w\in X_m}\pi(w)\log d_m^{\mathrm{bin}}(w)+m\log 2\\
&=
m\log 2-S(\tau_\pi).
\end{aligned}
\]
取 \(\pi=\mu_m\) 即得 \(S(\tau_{\mu_m})=m\log 2\)。又由 \(D_{\mathrm{KL}}(\pi\|\mu_m)\ge 0\) 且等号当且仅当 \(\pi=\mu_m\)，得到唯一最大熵性。证毕。
\end{proof}

\begin{theorem}[碰撞概率等于纤维群胚代数的归一化维数]\label{thm:xi-foldbin-collision-probability-normalized-groupoid-dimension}
\leanverified{paper\_xi\_foldbin\_collision\_probability\_normalized\_groupoid\_dimension}
沿用定理 \ref{thm:xi-foldbin-groupoid-maxent-tracial-state} 的记号。若 \(U,V\) 为 \(\Omega_m=\{0,1\}^m\) 上独立均匀微态，则
\[
\PP\!\bigl(\Fold^{\mathrm{bin}}_m(U)=\Fold^{\mathrm{bin}}_m(V)\bigr)
=
\sum_{w\in X_m}\mu_m(w)^2
=
\frac{\dim_{\CC}A_m^{\mathrm{bin}}}{4^m}.
\]
等价地，
\[
\dim_{\CC}A_m^{\mathrm{bin}}
=
\sum_{w\in X_m}\bigl(d_m^{\mathrm{bin}}(w)\bigr)^2.
\]
\end{theorem}
\begin{proof}
仍由命题 \ref{prop:fold-bin-groupoid-algebra-wedderburn}，
\[
\dim_{\CC}A_m^{\mathrm{bin}}
=
\sum_{w\in X_m}
\dim_{\CC}M_{d_m^{\mathrm{bin}}(w)}(\CC)
=
\sum_{w\in X_m}\bigl(d_m^{\mathrm{bin}}(w)\bigr)^2.
\]
另一方面，独立均匀抽样给出
\[
\PP\!\bigl(\Fold^{\mathrm{bin}}_m(U)=\Fold^{\mathrm{bin}}_m(V)\bigr)
=
\sum_{w\in X_m}
\left(\frac{d_m^{\mathrm{bin}}(w)}{2^m}\right)^2
=
\frac{1}{4^m}
\sum_{w\in X_m}\bigl(d_m^{\mathrm{bin}}(w)\bigr)^2.
\]
合并两式即得结论。证毕。
\end{proof}

\begin{theorem}[bin-fold 规范群的表示分配律与 \texorpdfstring{$m=6$}{m=6} 的精确谱]\label{thm:xi-foldbin-gauge-representation-distribution-law-m6-spectrum}
\leanverified{paper\_xi\_foldbin\_gauge\_representation\_distribution\_law\_m6\_spectrum}
令
\[
G_m^{\mathrm{bin}}:=\Aut(\Fold^{\mathrm{bin}}_m)
\cong
\prod_{w\in X_m}\mathfrak S_{d_m^{\mathrm{bin}}(w)}.
\]
则 \(\widehat{G_m^{\mathrm{bin}}}\) 中的不可约复表示与分拆族
\[
\{\lambda_w\vdash d_m^{\mathrm{bin}}(w)\}_{w\in X_m}
\]
一一对应，因而
\[
|\widehat{G_m^{\mathrm{bin}}}|
=
\prod_{w\in X_m}p\!\bigl(d_m^{\mathrm{bin}}(w)\bigr),
\]
其中 \(p(n)\) 为分拆函数。并且表示 \(\zeta\)-函数具有完全乘积分解
\[
\zeta_{G_m^{\mathrm{bin}}}^{\mathrm{irr}}(s)
:=
\sum_{\rho\in\widehat{G_m^{\mathrm{bin}}}}(\dim\rho)^{-s}
=
\prod_{w\in X_m}
\left(
\sum_{\lambda\vdash d_m^{\mathrm{bin}}(w)}(\dim\lambda)^{-s}
\right).
\]
若
\[
N_m:=\#\{w\in X_m:\ d_m^{\mathrm{bin}}(w)\ge 2\},
\]
则
\[
\bigl(G_m^{\mathrm{bin}}\bigr)^{\mathrm{ab}}\cong (\ZZ/2\ZZ)^{N_m},
\qquad
\bigl|\Hom(G_m^{\mathrm{bin}},\{\pm 1\})\bigr|=2^{N_m},
\]
从而 \(G_m^{\mathrm{bin}}\) 恰有 \(2^{N_m}-1\) 个指数为 \(2\) 的正规子群。
特别地，在 \(m=6\) 时，
\[
G_6^{\mathrm{bin}}
\cong
S_2^{\,8}\times S_3^{\,4}\times S_4^{\,9},
\qquad
\bigl(G_6^{\mathrm{bin}}\bigr)^{\mathrm{ab}}\cong (\ZZ/2\ZZ)^{21},
\]
并且
\[
|\widehat{G_6^{\mathrm{bin}}}|
=
p(2)^8p(3)^4p(4)^9
=
2^8\cdot 3^4\cdot 5^9,
\qquad
\bigl|\Hom(G_6^{\mathrm{bin}},\{\pm 1\})\bigr|=2^{21}.
\]
\end{theorem}
\begin{proof}
命题 \ref{prop:fold-bin-gauge-decomposition} 给出纤维直积分解。有限直积群的不可约表示恰为各因子不可约表示的外张量积，而 \(\mathfrak S_d\) 的不可约复表示由分拆 \(\lambda\vdash d\) 分类，故得到对 \(\widehat{G_m^{\mathrm{bin}}}\) 的参数化以及
\[
|\widehat{G_m^{\mathrm{bin}}}|
=
\prod_{w\in X_m}p\!\bigl(d_m^{\mathrm{bin}}(w)\bigr).
\]
维数在外张量积下相乘，遂有
\[
\zeta_{G_m^{\mathrm{bin}}}^{\mathrm{irr}}(s)
=
\prod_{w\in X_m}
\left(
\sum_{\lambda\vdash d_m^{\mathrm{bin}}(w)}(\dim\lambda)^{-s}
\right).
\]
对一维表示，只需注意：\(\mathfrak S_1\) 只有平凡表示，而 \(\mathfrak S_d\)（\(d\ge 2\)）恰有平凡与符号两种一维表示。
因此
\[
\Hom(G_m^{\mathrm{bin}},\{\pm 1\})
\cong
\prod_{w:\,d_m^{\mathrm{bin}}(w)\ge 2}\ZZ/2\ZZ,
\]
其基数为 \(2^{N_m}\)，并与 Abel 化同构。
每个非平凡同态的核都为指数 \(2\) 的正规子群，反之任一指数 \(2\) 正规子群都来自某个非平凡 \(\{\pm1\}\)-角色，故数目为 \(2^{N_m}-1\)。
最后在 \(m=6\) 时，代入推论 \ref{cor:fold-bin-gauge-m6} 的分解即可。证毕。
\end{proof}

\begin{theorem}[折叠同余商半环的奇异环判别与幂零层大小]\label{thm:xi-fold-congruence-semiring-singularity-nilpotent-profile}
\leanverified{paper\_xi\_fold\_congruence\_semiring\_singularity\_nilpotent\_profile}
设
\[
\mathcal R_m:=\Omega_m/\!\equiv_m.
\]
并将 Fibonacci 模数分解写为
\[
F_{m+2}=\prod_{i=1}^{s}p_i^{e_i}.
\]
则在定理 \ref{thm:fold-congruence-mod-semirings} 的半环同构下，\(\mathcal R_m\) 作为交换环满足：
\begin{enumerate}
  \item 存在规范分解
  \[
  \mathcal R_m\cong \prod_{i=1}^{s}\ZZ/(p_i^{e_i}\ZZ).
  \]
  \item 幂等元个数恰为
  \[
  2^{s}=2^{\omega(F_{m+2})}.
  \]
  \item 幂零元个数恰为
  \[
  \frac{F_{m+2}}{\operatorname{rad}(F_{m+2})}
  =
  \prod_{i=1}^{s}p_i^{e_i-1}.
  \]
  \item \(\mathcal R_m\) 为约化环当且仅当 \(F_{m+2}\) 平方自由；为局部环当且仅当 \(F_{m+2}\) 为素数幂。
\end{enumerate}
因此，折叠同余商的非约化奇异性完全由 Fibonacci 模数的平方因子控制。
\end{theorem}
\begin{proof}
由定理 \ref{thm:fold-congruence-mod-semirings}，
\[
\mathcal R_m\cong \ZZ/(F_{m+2}\ZZ).
\]
对整数模的 Chinese remainder 分解给出
\[
\ZZ/(F_{m+2}\ZZ)\cong \prod_{i=1}^{s}\ZZ/(p_i^{e_i}\ZZ),
\]
从而得到第一条。

在每个局部因子 \(\ZZ/(p_i^{e_i}\ZZ)\) 中，幂等元只能取 \(0\) 或 \(1\)。因此全体幂等元可由各坐标独立选择 \(0/1\) 得到，数量恰为 \(2^s\)。这也与推论 \ref{cor:fold-mod-semirings-idempotents-crt-boolean} 一致。

另一方面，在 \(\ZZ/(p_i^{e_i}\ZZ)\) 中，幂零元恰为所有 \(p_i\) 的倍数，其个数为 \(p_i^{e_i-1}\)。故乘积分解下的幂零元个数为
\[
\prod_{i=1}^{s}p_i^{e_i-1}
=
\frac{\prod_i p_i^{e_i}}{\prod_i p_i}
=
\frac{F_{m+2}}{\operatorname{rad}(F_{m+2})}.
\]

有限交换环约化当且仅当不存在非零幂零元，故等价于每个 \(e_i=1\)，亦即 \(F_{m+2}\) 平方自由。又有限交换环为局部环当且仅当其 Chinese remainder 分解只有一个因子，故等价于 \(s=1\)，亦即 \(F_{m+2}\) 为素数幂。证毕。
\end{proof}

\begin{theorem}[Fold 商半环 Jacobson 根基的 Loewy 长度与分层商公式]\label{thm:xi-fold-congruence-semiring-jacobson-loewy-profile}
沿用定理 \ref{thm:fold-congruence-mod-semirings} 的记号，并写
\[
R_m:=\mathcal R_m\cong \ZZ/(F_{m+2}\ZZ),
\qquad
F_{m+2}=\prod_{i=1}^{s}p_i^{e_i}.
\]
记 \(J_m:=J(R_m)\) 为 Jacobson 根基。则有：
\begin{enumerate}
  \item 根基的显式描述为
  \[
  \boxed{\;
  J_m=\left(\prod_{i=1}^{s}p_i\right)R_m.\;}
  \]
  \item 若记
  \[
  \ell_m:=\min\{t\ge 1:\ J_m^t=0\},
  \]
  则
  \[
  \boxed{\;
  \ell_m=\max_{1\le i\le s}e_i.\;}
  \]
  \item 对每个整数 \(0\le t<\ell_m\)，有
  \[
  \boxed{\;
  J_m^t/J_m^{t+1}\cong \prod_{i:\,e_i>t}\ZZ/p_i\ZZ,\qquad
  \bigl|J_m^t/J_m^{t+1}\bigr|=\prod_{i:\,e_i>t}p_i,\;}
  \]
  其中约定 \(J_m^0:=R_m\)。
\end{enumerate}
因此，Fold 商半环的 Loewy 深度恰由 Fibonacci 模数中最大的素数幂指数决定。
\end{theorem}
\begin{proof}
由定理 \ref{thm:fold-congruence-mod-semirings} 与 Chinese remainder 分解，
\[
R_m\cong \prod_{i=1}^{s}\ZZ/(p_i^{e_i}\ZZ).
\]
每个局部因子 \(\ZZ/(p_i^{e_i}\ZZ)\) 的 Jacobson 根基都是主理想
\[
(p_i)/(p_i^{e_i}),
\]
故整体根基为
\[
J_m\cong \prod_{i=1}^{s}(p_i)/(p_i^{e_i}),
\]
这恰等价于
\[
J_m=\left(\prod_{i=1}^{s}p_i\right)R_m.
\]

对任意 \(t\ge 0\)，于是有
\[
J_m^t\cong \prod_{i=1}^{s}(p_i^t)/(p_i^{e_i}).
\]
其中第 \(i\) 个因子当且仅当 \(t\ge e_i\) 时为零。因此 \(J_m^t=0\) 当且仅当对所有 \(i\) 都有 \(t\ge e_i\)，故满足 \(J_m^t=0\) 的最小整数恰为
\[
\ell_m=\max_{1\le i\le s}e_i.
\]

最后，对 \(0\le t<\ell_m\)，由上式得到
\[
J_m^t/J_m^{t+1}
\cong
\prod_{i=1}^{s}
\frac{(p_i^t)/(p_i^{e_i})}{(p_i^{t+1})/(p_i^{e_i})}.
\]
若 \(e_i\le t\)，则该因子为零；若 \(e_i>t\)，则
\[
\frac{(p_i^t)/(p_i^{e_i})}{(p_i^{t+1})/(p_i^{e_i})}
\cong
\ZZ/p_i\ZZ.
\]
故
\[
J_m^t/J_m^{t+1}\cong \prod_{i:\,e_i>t}\ZZ/p_i\ZZ,
\]
并取基数即得
\[
\bigl|J_m^t/J_m^{t+1}\bigr|=\prod_{i:\,e_i>t}p_i.
\]
证毕。
\end{proof}

\paragraph{window-\texorpdfstring{$6$}{6} 残基半环与奇偶角色空间的计数分离}
在 \(m=6\) 的具体窗口上，折叠同余商的有限残基算术与 fold-gauge 的奇偶角色空间可以直接并列比较，从而得到一条算术分层与一条严格的有限状态分离结论。

\begin{theorem}[window-\texorpdfstring{$6$}{6} 残基半环的单位--零因子分层与互补 CRT 幂等元]\label{thm:xi-window6-residue-semiring-unit-zerodivisor-idempotent-trichotomy}
记
\[
\mathcal R_6:=\Omega_6/\!\equiv_6.
\]
并令
\[
\mathrm{ZD}_6^\times
:=
\{x\in \mathcal R_6:\ x\neq 0,\ \exists\,y\neq 0\ \text{使}\ xy=0\}.
\]
则
\[
\mathcal R_6\cong \ZZ/(F_8\ZZ)=\ZZ/(21\ZZ),
\]
\[
\mathcal R_6
=
\mathcal R_6^\times\sqcup \mathrm{ZD}_6^\times\sqcup \{\overline 0\},
\qquad
|\mathcal R_6^\times|=12,\quad
|\mathrm{ZD}_6^\times|=8.
\]
并且全部幂等元恰为
\[
\operatorname{Idem}(\mathcal R_6)=\{\overline 0,\overline 1,e_{(3)},e_{(7)}\},
\qquad
e_{(3)}=\overline 7,\quad e_{(7)}=\overline{15},
\]
并满足
\[
e_{(3)}+e_{(7)}=\overline 1,
\qquad
e_{(3)}e_{(7)}=\overline 0.
\]
\leanverified{zmod21\_zerodiv\_count}
\leanverified{zmod21\_trichotomy}
\end{theorem}
\begin{proof}
由定理 \ref{thm:xi-fold-congruence-semiring-singularity-nilpotent-profile}，
\[
\mathcal R_6\cong \ZZ/(F_8\ZZ)=\ZZ/(21\ZZ).
\]
故单位元个数为
\[
|\mathcal R_6^\times|=\varphi(21)=12.
\]

若 \(\overline a\in \ZZ/(21\ZZ)\) 是非零非单位元，则 \(\gcd(a,21)\in\{3,7\}\)。记 \(d:=\gcd(a,21)\)，则
\[
\overline a\cdot \overline{21/d}=\overline 0,
\qquad
\overline{21/d}\neq \overline 0.
\]
因此每个非零非单位元都是零因子。反过来，零因子当然不可能是单位，于是非零零因子恰为全部非零非单位元，因而
\[
|\mathrm{ZD}_6^\times|=21-12-1=8.
\]

为求幂等元的显式代表，只需利用中国剩余同构
\[
\ZZ/(21\ZZ)\cong \ZZ/(3\ZZ)\times \ZZ/(7\ZZ).
\]
幂等元在每个局部分量上只能取 \(0\) 或 \(1\)，故四个 CRT 坐标为
\[
(0,0),\qquad (1,1),\qquad (1,0),\qquad (0,1),
\]
它们回拉到 \(\ZZ/(21\ZZ)\) 中恰为
\[
\overline 0,\qquad \overline 1,\qquad \overline 7,\qquad \overline{15}.
\]
最后，
\[
\overline 7+\overline{15}=\overline{22}=\overline 1,
\qquad
\overline 7\cdot \overline{15}=\overline{105}=\overline 0.
\]
证毕。
\end{proof}

\begin{corollary}[window-\texorpdfstring{$6$}{6} 残基半环无法忠实承载全部奇偶角色]\label{cor:xi-window6-residue-semiring-parity-character-gap}
\leanverified{paper\_xi\_window6\_residue\_semiring\_parity\_character\_gap}
沿用定理 \ref{thm:xi-window6-residue-semiring-unit-zerodivisor-idempotent-trichotomy} 的记号。另一方面，
\[
\bigl(G_6^{\mathrm{bin}}\bigr)^{\mathrm{ab}}\cong (\ZZ/2\ZZ)^{21},
\qquad
\bigl|\Hom(G_6^{\mathrm{bin}},\{\pm 1\})\bigr|=2^{21}.
\]
因此，任何仅通过残基半环 \(\mathcal R_6\) 取值的读出都不可能分离全部奇偶角色；特别地，不存在单射
\[
\Hom(G_6^{\mathrm{bin}},\{\pm 1\})\hookrightarrow \mathcal R_6.
\]
\end{corollary}
\begin{proof}
由定理 \ref{thm:xi-foldbin-gauge-representation-distribution-law-m6-spectrum}，
\[
\bigl(G_6^{\mathrm{bin}}\bigr)^{\mathrm{ab}}\cong (\ZZ/2\ZZ)^{21},
\qquad
\bigl|\Hom(G_6^{\mathrm{bin}},\{\pm 1\})\bigr|=2^{21}.
\]
而由定理 \ref{thm:xi-window6-residue-semiring-unit-zerodivisor-idempotent-trichotomy}，
\[
|\mathcal R_6|=21<2^{21}.
\]
因此，从奇偶角色空间到 \(\mathcal R_6\) 的任何映射都不可能是单射。于是，任何仅通过 \(\mathcal R_6\) 取值的残基读出都无法忠实区分全部奇偶角色。证毕。
\end{proof}

\begin{theorem}[零温四态极限的 Artin--Mazur \texorpdfstring{$\zeta$}{zeta} 函数、迹递推与 primitive 轨道精确主项]\label{thm:xi-zero-temperature-fourstate-artin-mazur-zeta-primitive-asymptotic}
\leanverified{paper\_xi\_zero\_temperature\_fourstate\_artin\_mazur\_zeta\_primitive\_asymptotic}
令 \(\Sigma_\infty\) 为由
\[
B_\infty=
\begin{pmatrix}
0&1&1&0\\
0&0&1&0\\
0&0&3&1\\
1&0&0&3
\end{pmatrix}
\]
定义的四态一步 SFT，并记其 Perron 根为 \(\rho:=\rho(B_\infty)\)。再令
\[
a_n:=\#\operatorname{Fix}(\sigma^n|_{\Sigma_\infty})=\Tr(B_\infty^n),
\qquad
\pi_n:=\frac{1}{n}\sum_{d\mid n}\mu(d)\,a_{n/d},
\]
其中 \(\mu\) 为 Möbius 函数。则
\[
\zeta_\infty(z)
:=
\exp\!\left(\sum_{n\ge 1}\frac{a_n}{n}z^n\right)
=
\frac{1}{\det(I-zB_\infty)}
=
\frac{1}{1-6z+9z^2-z^3-z^4}.
\]
此外，序列 \((a_n)_{n\ge 1}\) 满足精确递推
\[
a_{n+4}=6a_{n+3}-9a_{n+2}+a_{n+1}+a_n
\qquad(n\ge 1),
\]
其初值为
\[
a_1=6,\qquad a_2=18,\qquad a_3=57,\qquad a_4=190.
\]
若将 \(B_\infty\) 的特征值按模长排序为
\[
\rho=\lambda_1,
\qquad
\eta:=|\lambda_2|\ge |\lambda_3|\ge |\lambda_4|,
\qquad
\eta<\rho,
\]
则 primitive 轨道数满足
\[
\pi_n=\frac{\rho^n}{n}+O\!\left(\frac{\eta^n}{n}\right).
\]
并且
\[
h_{\mathrm{top}}(\Sigma_\infty)=\log\rho.
\]
等价地，在推论 \ref{cor:real-input-40-ground-entropy} 所采用的原始系统步长上，零温地基每步熵密度为 \(\tfrac12\log\rho\)。
\end{theorem}
\begin{proof}
零温四态极限段已给出
\[
\det(I-zB_\infty)=1-6z+9z^2-z^3-z^4.
\]
对任意一步 SFT，Artin--Mazur \(\zeta\)-函数满足
\[
\zeta(z)=\det(I-zB)^{-1},
\qquad
\#\operatorname{Fix}(\sigma^n)=\Tr(B^n),
\]
故第一式立即成立。

\(B_\infty\) 的特征多项式为
\[
\chi(\lambda)=\lambda^4-6\lambda^3+9\lambda^2-\lambda-1.
\]
由 Cayley--Hamilton 定理，
\[
B_\infty^{n+4}=6B_\infty^{n+3}-9B_\infty^{n+2}+B_\infty^{n+1}+B_\infty^{n}
\qquad(n\ge 1),
\]
两侧取迹即得递推。初值由直接矩阵乘法给出
\[
\Tr(B_\infty)=6,\qquad
\Tr(B_\infty^2)=18,\qquad
\Tr(B_\infty^3)=57,\qquad
\Tr(B_\infty^4)=190.
\]

再看 primitive 轨道主项。对 \(\chi\) 直接求根可知其四个根两两不同，且按模长依次为
\[
\rho\approx 3.596154676009,\qquad
\eta\approx 2.182644323596,\qquad
|\lambda_3|\approx 0.484278410709,\qquad
|\lambda_4|\approx 0.263077410313.
\]
特别地，
\[
\eta>\sqrt{\rho}.
\]
因此由谱分解，
\[
a_n=\lambda_1^n+\lambda_2^n+\lambda_3^n+\lambda_4^n
=
\rho^n+O(\eta^n).
\]
再由 Möbius 反演公式
\[
n\pi_n=\sum_{d\mid n}\mu(d)\,a_{n/d}
=
a_n+\sum_{\substack{d\mid n\\ d\ge 2}}\mu(d)\,a_{n/d},
\]
其中主项 \(a_n=\rho^n+O(\eta^n)\)。对任意真约数 \(d\ge 2\)，有 \(n/d\le n/2\)，故
\[
a_{n/d}=O(\rho^{n/d})=O(\rho^{n/2}).
\]
又因 \(\eta>\sqrt{\rho}\)，故 \(\rho^{n/2}=O(\eta^n)\)；再结合约数函数 \(\tau(n)=n^{o(1)}\) 的次指数增长，可得
\[
\sum_{\substack{d\mid n\\ d\ge 2}}\mu(d)\,a_{n/d}=O(\eta^n).
\]
于是
\[
n\pi_n=\rho^n+O(\eta^n),
\]
即
\[
\pi_n=\frac{\rho^n}{n}+O\!\left(\frac{\eta^n}{n}\right).
\]

最后，边移位 \(\Sigma_\infty\) 的拓扑熵等于 \(\log\rho(B_\infty)\)；与推论 \ref{cor:real-input-40-ground-entropy} 比较，即得原始系统步长上的熵密度为 \(\tfrac12\log\rho\)。证毕。
\end{proof}

\begin{theorem}[零温四态地基 \texorpdfstring{SFT}{SFT} 的 Perron--Parry 数据完全闭式]\label{thm:xi-zero-temperature-fourstate-perron-parry-closed-form}
\leanverified{paper\_xi\_zero\_temperature\_fourstate\_perron\_parry\_closed\_form}
沿用定理 \ref{thm:xi-zero-temperature-fourstate-artin-mazur-zeta-primitive-asymptotic} 与
\ref{thm:xi-terminal-real-input-positive-entropy-freezing-face} 的记号。记
\[
\widehat F(w)=1-6w+9w^2-w^3-w^4=\det(I-wB_\infty),
\qquad
\rho:=\rho(B_\infty).
\]
则成立：
\begin{enumerate}
  \item \(\rho\) 是四次多项式
  \[
  x^4-6x^3+9x^2-x-1
  \]
  的最大实根；该多项式在 \(\QQ\) 上不可约。
  \item \(B_\infty\) 的一组正右 Perron 特征向量与正左 Perron 特征向量可分别取为
  \[
  v=
  \begin{pmatrix}
  (\rho-3)^2\\
  \rho^{-1}\\
  1\\
  \rho-3
  \end{pmatrix},
  \qquad
  \ell=
  \begin{pmatrix}
  1\\
  \rho^{-1}\\
  \rho(\rho-3)\\
  \rho
  \end{pmatrix},
  \]
  并满足
  \[
  B_\infty v=\rho v,
  \qquad
  B_\infty^{\mathsf T}\ell=\rho\ell.
  \]
  \item 由 Parry 公式
  \[
  P_{ij}=\frac{(B_\infty)_{ij}v_j}{\rho v_i},
  \qquad
  \pi_i\propto \ell_i v_i
  \]
  所定义的极大熵平衡态可显式写成
  \[
  P=
  \begin{pmatrix}
  0 & \dfrac{1}{1+\rho} & \dfrac{\rho}{1+\rho} & 0\\[2mm]
  0 & 0 & 1 & 0\\[2mm]
  0 & 0 & \dfrac{3}{\rho} & \dfrac{\rho-3}{\rho}\\[2mm]
  \dfrac{\rho-3}{\rho} & 0 & 0 & \dfrac{3}{\rho}
  \end{pmatrix},
  \]
  \[
  \pi=
  \frac{1}{3\rho^2+\rho-6}
  \Bigl((\rho-3)(\rho+1),\ \rho-3,\ \rho(\rho+1),\ \rho(\rho+1)\Bigr).
  \]
\end{enumerate}
\end{theorem}
\begin{proof}
由定理 \ref{thm:xi-zero-temperature-fourstate-artin-mazur-zeta-primitive-asymptotic}，
\(B_\infty\) 的特征多项式正是
\[
\chi_{B_\infty}(x)=x^4-6x^3+9x^2-x-1.
\]
先证其在 \(\QQ\) 上不可约。若可约，则由 Gauss 引理可写成
\[
x^4-6x^3+9x^2-x-1=(x^2+ax+b)(x^2+cx+d),
\qquad a,b,c,d\in\ZZ.
\]
比较系数得到
\[
a+c=-6,\qquad ac+b+d=9,\qquad ad+bc=-1,\qquad bd=-1.
\]
由 \(bd=-1\)，只能有 \((b,d)=(1,-1)\) 或 \((-1,1)\)。

若 \((b,d)=(1,-1)\)，则
\[
ac=9,\qquad c-a=-1,\qquad a+c=-6,
\]
从而
\[
(a,c)=\left(-\frac52,-\frac72\right),
\]
与 \(ac=9\) 矛盾。若 \((b,d)=(-1,1)\)，则
\[
ac=9,\qquad a-c=-1,\qquad a+c=-6,
\]
同样推出
\[
(a,c)=\left(-\frac72,-\frac52\right),
\]
仍与 \(ac=9\) 矛盾。故 \(\chi_{B_\infty}\) 在 \(\QQ\) 上不可约。又
\[
\chi_{B_\infty}(3)=-4<0,
\qquad
\chi_{B_\infty}(4)=11>0,
\]
故 \(\rho\in(3,4)\) 为其最大实根。

现取 \(v=(a,b,c,d)^{\mathsf T}\) 满足 \(B_\infty v=\rho v\)。逐项展开得
\[
b+c=\rho a,\qquad c=\rho b,\qquad 3c+d=\rho c,\qquad a+3d=\rho d.
\]
令 \(c=1\)，则
\[
b=\rho^{-1},\qquad d=\rho-3,\qquad a=(\rho-3)^2.
\]
此时第一式化为
\[
\rho(\rho-3)^2=1+\rho^{-1},
\]
亦即
\[
\rho^2(\rho-3)^2=\rho+1.
\]
而
\[
\rho^2(\rho-3)^2-(\rho+1)
=
\rho^4-6\rho^3+9\rho^2-\rho-1
=
0,
\]
故 \(v\) 确为右特征向量。又因 \(\rho>3\)，\(v\) 的各分量全部严格为正。

对左特征向量，令 \(\ell=(\alpha,\beta,\gamma,\delta)^{\mathsf T}\) 满足
\[
B_\infty^{\mathsf T}\ell=\rho\ell.
\]
逐项展开得
\[
\delta=\rho\alpha,\qquad \alpha=\rho\beta,\qquad \alpha+\beta+3\gamma=\rho\gamma,\qquad \gamma+3\delta=\rho\delta.
\]
令 \(\alpha=1\)，则
\[
\beta=\rho^{-1},\qquad \delta=\rho,\qquad \gamma=\rho(\rho-3),
\]
而第三式同样等价于
\[
\rho^4-6\rho^3+9\rho^2-\rho-1=0.
\]
故 \(\ell\) 确为左特征向量，且其各分量亦严格为正。

再由 Parry 公式计算转移矩阵。对第一行，
\[
P_{12}
=
\frac{1\cdot \rho^{-1}}{\rho(\rho-3)^2}
=
\frac{1}{\rho^2(\rho-3)^2}
=
\frac{1}{1+\rho},
\]
\[
P_{13}
=
\frac{1\cdot 1}{\rho(\rho-3)^2}
=
\frac{\rho}{\rho^2(\rho-3)^2}
=
\frac{\rho}{1+\rho}.
\]
其余非零项直接得到
\[
P_{23}=1,\qquad
P_{33}=\frac{3}{\rho},\qquad
P_{34}=\frac{\rho-3}{\rho},\qquad
P_{41}=\frac{\rho-3}{\rho},\qquad
P_{44}=\frac{3}{\rho},
\]
从而得到所示矩阵 \(P\)。

对平稳分布，
\[
\pi_i\propto \ell_i v_i
=
\bigl((\rho-3)^2,\ \rho^{-2},\ \rho(\rho-3),\ \rho(\rho-3)\bigr).
\]
将该向量乘以公因子 \(\rho^2(\rho-3)\)，并利用
\[
\rho^2(\rho-3)^2=\rho+1,
\]
便得到等价权重
\[
\bigl((\rho-3)(\rho+1),\ \rho-3,\ \rho(\rho+1),\ \rho(\rho+1)\bigr).
\]
其分量和为
\[
(\rho-3)(\rho+1)+(\rho-3)+2\rho(\rho+1)=3\rho^2+\rho-6,
\]
故归一化后恰得所示 \(\pi\)。证毕。
\end{proof}

\begin{theorem}[零温 Parry 平衡态的严格不可逆性与无穷熵产生]\label{thm:xi-zero-temperature-parry-irreversible-infinite-entropy-production}
沿用定理 \ref{thm:xi-zero-temperature-fourstate-perron-parry-closed-form} 的记号。则对应的 Parry 马尔可夫链 \((P,\pi)\) 严格不可逆。更强地，按扩展值约定
\[
a\log\frac{a}{0}=+\infty
\qquad(a>0)
\]
定义的稳态熵产生率
\[
\mathcal E(P,\pi)
:=
\sum_{i,j}\pi_iP_{ij}\log\frac{\pi_iP_{ij}}{\pi_jP_{ji}}
\]
满足
\[
\mathcal E(P,\pi)=+\infty.
\]
\end{theorem}
\begin{proof}
由定理 \ref{thm:xi-zero-temperature-fourstate-perron-parry-closed-form}，
\[
P_{13}=\frac{\rho}{1+\rho}>0,
\qquad
P_{31}=0.
\]
又因 \(\rho>3\)，平稳分布的第一分量满足
\[
\pi_1=\frac{(\rho-3)(\rho+1)}{3\rho^2+\rho-6}>0.
\]
从而
\[
\pi_1P_{13}>0,
\qquad
\pi_3P_{31}=0.
\]
若该链可逆，则必须满足详细平衡
\[
\pi_iP_{ij}=\pi_jP_{ji}
\qquad(\forall i,j),
\]
这与 \((i,j)=(1,3)\) 时的上式矛盾，故链严格不可逆。

再看熵产生率中的 \((i,j)=(1,3)\) 项。由上式可得
\[
\pi_1P_{13}\log\frac{\pi_1P_{13}}{\pi_3P_{31}}
=
\pi_1P_{13}\log\frac{\pi_1P_{13}}{0}
=
+\infty.
\]
因此整和必发散到 \(+\infty\)。证毕。
\end{proof}

\begin{theorem}[零温四态地基的有限矩阵行列式表示极小性]\label{thm:xi-zero-temperature-fourstate-determinant-state-minimality}
\leanverified{paper\_xi\_zero\_temperature\_fourstate\_determinant\_state\_minimality}
沿用定理 \ref{thm:xi-terminal-real-input-positive-entropy-freezing-face} 的记号。不存在任何维数 \(k\le 3\) 的矩阵
\[
A\in M_k(\ZZ_{\ge 0})
\]
使
\[
\det(I-wA)=\widehat F(w)=1-6w+9w^2-w^3-w^4.
\]
因而在“有限矩阵行列式表示”的意义下，零温四态地基的最小状态数恰为 \(4\)。
\end{theorem}
\begin{proof}
更强地，对任意 \(A\in M_k(\ZZ)\) 都有
\[
\det(I-wA)
\]
是关于 \(w\) 的次数至多为 \(k\) 的多项式，其最高次项为
\[
(-1)^k\det(A)\,w^k.
\]
而
\[
\widehat F(w)=1-6w+9w^2-w^3-w^4
\]
的次数恰为 \(4\)。因此若 \(k\le 3\)，则不可能有
\[
\det(I-wA)=\widehat F(w).
\]
另一方面，定理 \ref{thm:xi-terminal-real-input-positive-entropy-freezing-face} 已给出
\[
\widehat F(w)=\det(I-wB_\infty)
\]
且 \(B_\infty\in M_4(\ZZ_{\ge 0})\)。故最小矩阵维数恰为 \(4\)。证毕。
\end{proof}

\begin{theorem}[fold-gauge 常数序列的非 \texorpdfstring{$C$}{C}-有限性与前两阶 \texorpdfstring{$\zeta(2r)$}{zeta(2r)} 重整化极限]\label{thm:xi-foldbin-gauge-constant-not-cfinite-first-two-zeta-limits}
令
\[
u_m:=\log C_m-\log(2\pi),
\qquad
q:=\frac{\varphi}{2}.
\]
则序列 \((u_m)_{m\ge 1}\) 不是 \(C\)-有限序列。
并且有两阶显式展开
\[
u_m
=
\frac{1}{3\varphi}\,q^m
-\frac{\sqrt5}{180}\,q^{3m}
+O(q^{5m}).
\]
因此
\[
\lim_{m\to\infty}q^{-m}u_m=\frac{1}{3\varphi},
\qquad
\lim_{m\to\infty}q^{-3m}\left(u_m-\frac{1}{3\varphi}q^m\right)
=
-\frac{\sqrt5}{180}.
\]
从而前两阶重整化极限已足以恢复
\[
B_2=\frac16,\qquad B_4=-\frac1{30},
\]
并进而恢复
\[
\zeta(2)=\frac{\pi^2}{6},
\qquad
\zeta(4)=\frac{\pi^4}{90}.
\]
\end{theorem}
\begin{proof}
非 \(C\)-有限性正是定理 \ref{thm:fold-bin-gauge-constant-not-c-finite} 的内容。
另一方面，定理 \ref{thm:fold-bin-gauge-constant-stirling-bernoulli-hierarchy} 在 \(R=2\) 时给出
\[
u_m
=
A_1q^m+A_2q^{3m}+O(q^{5m}),
\]
其中
\[
A_1=
\frac{B_2}{1\cdot 1}\bigl(\varphi^{-1}+\varphi^{-1}\bigr)
=
\frac16\cdot\frac{2}{\varphi}
=
\frac{1}{3\varphi},
\]
\[
A_2=
\frac{B_4}{2\cdot 3}\bigl(\varphi^{-1}+\varphi\bigr)
=
-\frac{1}{30}\cdot\frac{\varphi+\varphi^{-1}}{6}
=
-\frac{\sqrt5}{180}.
\]
于是得到所示两阶展开与两条重整化极限。
最后用标准恒等式
\[
B_{2r}=(-1)^{r-1}\frac{2(2r)!}{(2\pi)^{2r}}\zeta(2r)
\]
即可从 \(B_2,B_4\) 恢复 \(\zeta(2),\zeta(4)\)。
证毕。
\end{proof}

\begin{theorem}[fold-gauge 常数对 \texorpdfstring{$2\pi$}{2pi} 的最终上侧逼近与首项刚性]\label{thm:xi-foldbin-gauge-constant-one-sided-approach-2pi}
沿用定理 \ref{thm:xi-foldbin-gauge-constant-not-cfinite-first-two-zeta-limits} 的记号，令
\[
q:=\frac{\varphi}{2}\in(0,1).
\]
则当 \(m\to\infty\) 时，
\[
\log C_m
=
\log(2\pi)+\frac{1}{3\varphi}q^m+O(q^{3m}),
\]
从而
\[
C_m
=
2\pi\left(1+\frac{1}{3\varphi}q^m+O(q^{2m})\right).
\]
特别地，存在 \(m_0\) 使得对一切 \(m\ge m_0\) 都有
\[
C_m>2\pi.
\]
因此 \(C_m\) 对 \(2\pi\) 的收敛最终始终来自上侧，而非振荡式逼近。
\end{theorem}
\begin{proof}
由定理 \ref{thm:xi-foldbin-gauge-constant-not-cfinite-first-two-zeta-limits}，
\[
\log C_m-\log(2\pi)
=
\frac{1}{3\varphi}q^m-\frac{\sqrt5}{180}q^{3m}+O(q^{5m})
=
\frac{1}{3\varphi}q^m+O(q^{3m}).
\]
由于首项系数 \(1/(3\varphi)\) 为正，故对充分大的 \(m\) 有 \(\log C_m>\log(2\pi)\)，即 \(C_m>2\pi\)。

再令
\[
x_m:=\frac{1}{3\varphi}q^m+O(q^{3m})=O(q^m).
\]
则
\[
C_m=2\pi e^{x_m}
=
2\pi\bigl(1+x_m+O(x_m^2)\bigr)
=
2\pi\left(1+\frac{1}{3\varphi}q^m+O(q^{2m})\right),
\]
因为 \(x_m^2=O(q^{2m})\)。证毕。
\end{proof}

\begin{theorem}[fold-gauge 序列对 Bernoulli 数与偶整数 \texorpdfstring{$\zeta$}{zeta} 值的递阶逆谱公式]\label{thm:xi-foldbin-gauge-constant-recursive-bernoulli-zeta-recovery}
沿用定理 \ref{thm:fold-bin-gauge-constant-stirling-bernoulli-hierarchy} 的记号，令
\[
E_m:=\log C_m-\log(2\pi),
\qquad
q:=\frac{\varphi}{2}.
\]
递归定义抽取序列
\[
R_m^{(1)}:=q^{-m}E_m,
\qquad
a_1^{\mathrm{ext}}:=\lim_{m\to\infty}R_m^{(1)},
\]
并对每个 \(r\ge 2\) 定义
\[
R_m^{(r)}
:=
q^{-(2r-1)m}
\left(
E_m-\sum_{j=1}^{r-1}a_j^{\mathrm{ext}}\,q^{(2j-1)m}
\right),
\qquad
a_r^{\mathrm{ext}}:=\lim_{m\to\infty}R_m^{(r)}.
\]
则上述各极限都存在，并满足
\[
a_r^{\mathrm{ext}}
=
\frac{B_{2r}}{r(2r-1)}\bigl(\varphi^{-1}+\varphi^{2r-3}\bigr)
\qquad(r\ge 1).
\]
因此对每个 \(r\ge 1\)，有显式恢复公式
\[
B_{2r}
=
\frac{r(2r-1)}{\varphi^{-1}+\varphi^{2r-3}}
\lim_{m\to\infty}R_m^{(r)},
\]
以及
\[
\zeta(2r)
=
(-1)^{r-1}\frac{(2\pi)^{2r}}{2(2r)!}\,
\frac{r(2r-1)}{\varphi^{-1}+\varphi^{2r-3}}
\lim_{m\to\infty}R_m^{(r)}.
\]
\end{theorem}
\begin{proof}
由定理 \ref{thm:fold-bin-gauge-constant-stirling-bernoulli-hierarchy}，对任意固定 \(r\ge 1\) 有展开
\[
E_m
=
\sum_{j=1}^{r}\alpha_j q^{(2j-1)m}+O(q^{(2r+1)m}),
\]
其中
\[
\alpha_j
:=
\frac{B_{2j}}{j(2j-1)}\bigl(\varphi^{-1}+\varphi^{2j-3}\bigr).
\]
当 \(r=1\) 时，
\[
R_m^{(1)}=q^{-m}E_m=\alpha_1+O(q^{2m}),
\]
故极限存在且 \(a_1^{\mathrm{ext}}=\alpha_1\)。

设对某个 \(r\ge 2\)，已知 \(a_j^{\mathrm{ext}}=\alpha_j\) 对 \(1\le j\le r-1\) 成立。则
\[
E_m-\sum_{j=1}^{r-1}a_j^{\mathrm{ext}}\,q^{(2j-1)m}
=
\alpha_r q^{(2r-1)m}+O(q^{(2r+1)m}),
\]
从而
\[
R_m^{(r)}
=
\alpha_r+O(q^{2m}).
\]
故极限存在且 \(a_r^{\mathrm{ext}}=\alpha_r\)。归纳完成。

最后对上式解出 \(B_{2r}\)，再与标准恒等式
\[
B_{2r}=(-1)^{r-1}\frac{2(2r)!}{(2\pi)^{2r}}\zeta(2r)
\]
联立，即得所示偶整数 \(\zeta\)-值恢复公式。证毕。
\end{proof}

% (User input window)

\paragraph{尾熵变分、容量平台与有限素数局部化的对偶序结构}
前两条把纤维尾计数在指数尺度上的几何压缩为压力谱与连续容量谱的统一变分公式，
后两条则把有限素数局部化系统的离散偏序与其 Pontryagin 对偶端的满射偏序写成同一条
\(\operatorname{Hom}\)--商映射链。

\begin{theorem}[尾熵函数同时生成压力谱与连续容量谱]\label{thm:xi-tail-entropy-pressure-capacity-variational}
\leanverified{paper\_xi\_tail\_entropy\_pressure\_capacity\_variational}
沿用命题 \ref{prop:op-algebra-capacity-kinks-equal-spectrum-levels} 与定理
\ref{thm:pom-oracle-capacity-stieltjes-inversion-mellin} 的记号。
于是对每个 \(m\) 都有
\[
\mathcal C_m^{\mathrm{cont}}(T)=\int_0^T N_m(t)\,\dd t,
\qquad
S_q(m)=q\int_0^\infty t^{q-1}N_m(t)\,\dd t
\qquad(q>0),
\]
并且 \(\mathcal C_m^{\mathrm{cont}}\) 与 \(N_m\) 彼此决定。
记
\[
M_m:=\max_{x\in X_m}d_m(x),
\qquad
\alpha_\ast:=\lim_{m\to\infty}\frac{1}{m}\log M_m.
\]
进一步假设：对每个 \(\gamma\in[0,\alpha_\ast]\)，极限
\[
s(\gamma):=\lim_{m\to\infty}\frac{1}{m}\log N_m(e^{\gamma m})
\]
存在，且
\(
\frac1m\log N_m(e^{\gamma m})
\)
在 \([0,\alpha_\ast]\) 的每个紧子区间上一致收敛到 \(s(\gamma)\)。
定义
\[
P(q):=\lim_{m\to\infty}\frac{1}{m}\log S_q(m)
\qquad(q>0),
\]
\[
\Gamma(\beta):=\lim_{m\to\infty}\frac{1}{m}\log \mathcal C_m^{\mathrm{cont}}(e^{\beta m})
\qquad(\beta\ge 0),
\]
则上述极限存在，且满足精确变分公式
\[
\boxed{\;
P(q)=\sup_{0\le \gamma\le \alpha_\ast}\bigl(q\gamma+s(\gamma)\bigr)\qquad(q>0)\; }
\]
以及
\[
\boxed{\;
\Gamma(\beta)=\sup_{0\le \gamma\le \min\{\beta,\alpha_\ast\}}
\bigl(\gamma+s(\gamma)\bigr)\qquad(\beta\ge 0)\; }.
\]
特别地，压力谱与容量谱并非两套独立对象，而是同一尾熵函数 \(s(\gamma)\) 的
Laplace--Abel 投影。
\end{theorem}
\begin{proof}
由定理 \ref{thm:pom-oracle-capacity-stieltjes-inversion-mellin}，
\[
S_q(m)=q\int_0^\infty t^{q-1}N_m(t)\,\dd t.
\]
由于 \(N_m(t)=0\) 当 \(t>M_m\) 时成立，作代换 \(t=e^{\gamma m}\) 得
\[
\begin{aligned}
S_q(m)
&=
qm\int_0^{\frac1m\log M_m}
\exp\!\left(
m\left(
q\gamma+\frac{1}{m}\log N_m(e^{\gamma m})
\right)
\right)\,\dd\gamma.
\end{aligned}
\]
又因为
\(
\frac1m\log M_m\to\alpha_\ast
\)
且
\(
\frac1m\log N_m(e^{\gamma m})\to s(\gamma)
\)
在任意紧子区间上一致收敛，
故可对紧区间上的指数积分应用标准 Laplace 原理，得到
\[
\lim_{m\to\infty}\frac{1}{m}\log S_q(m)
=
\sup_{0\le \gamma\le \alpha_\ast}\bigl(q\gamma+s(\gamma)\bigr).
\]
这就给出 \(P(q)\) 的公式。
\par
同理，由命题 \ref{prop:op-algebra-capacity-kinks-equal-spectrum-levels}，
\[
\mathcal C_m^{\mathrm{cont}}(e^{\beta m})
=
\int_0^{e^{\beta m}}N_m(t)\,\dd t.
\]
再次代换 \(t=e^{\gamma m}\)，并注意当 \(\gamma>\frac1m\log M_m\) 时
\(N_m(e^{\gamma m})=0\)，可得
\[
\mathcal C_m^{\mathrm{cont}}(e^{\beta m})
=
m\int_0^{\min\{\beta,\frac1m\log M_m\}}
\exp\!\left(
m\left(
\gamma+\frac{1}{m}\log N_m(e^{\gamma m})
\right)
\right)\,\dd\gamma.
\]
再次应用 Laplace 原理，即得
\[
\lim_{m\to\infty}\frac{1}{m}\log \mathcal C_m^{\mathrm{cont}}(e^{\beta m})
=
\sup_{0\le \gamma\le \min\{\beta,\alpha_\ast\}}
\bigl(\gamma+s(\gamma)\bigr).
\]
证毕。
\end{proof}

\begin{theorem}[高预算连续容量的线性--平台律]\label{thm:xi-high-budget-continuous-capacity-linear-plateau}
\leanverified{paper\_xi\_high\_budget\_continuous\_capacity\_linear\_plateau}
对每个 \(m\)，令
\[
M_m:=\max_{x\in X_m}d_m(x),
\qquad
\mathcal M_m:=\{x\in X_m:\ d_m(x)=M_m\},
\qquad
K_m:=|\mathcal M_m|,
\]
并定义
\[
M_m^{(2)}:=\max\{d_m(x):\ d_m(x)<M_m\},
\]
若次大纤维不存在则取 \(M_m^{(2)}:=1\)。
进一步设极限
\[
\alpha_\ast:=\lim_{m\to\infty}\frac{1}{m}\log M_m,
\qquad
g_\ast:=\lim_{m\to\infty}\frac{1}{m}\log K_m
\]
存在，并记
\[
\alpha_\ast^{(2)}:=\limsup_{m\to\infty}\frac{1}{m}\log M_m^{(2)}.
\]
并假设
\[
\alpha_\ast^{(2)}<\alpha_\ast.
\]
定义连续容量指数
\[
\Gamma(\beta):=\lim_{m\to\infty}\frac{1}{m}\log \mathcal C_m^{\mathrm{cont}}(e^{\beta m})
\qquad(\beta\ge 0),
\]
以及
\[
P(1):=\lim_{m\to\infty}\frac{1}{m}\log S_1(m).
\]
再设
\[
\beta_c:=\log\varphi+\alpha_\ast^{(2)}-g_\ast,
\qquad
a_c:=\frac{\log\varphi-g_\ast}{\alpha_\ast-\alpha_\ast^{(2)}}.
\]
则成立：
\begin{enumerate}
  \item 若 \(\beta_c<\beta<\alpha_\ast\)，则
  \[
  \boxed{\;
  \Gamma(\beta)=\beta+g_\ast\; }.
  \]
  \item 若 \(\beta>\alpha_\ast\)，则
  \[
  \boxed{\;
  \Gamma(\beta)=P(1)=\log 2\; }.
  \]
  \item 若进一步 \(a_c<1\)（等价地 \(\beta_c<\alpha_\ast\)），则
  \[
  P(1)=\alpha_\ast+g_\ast,
  \]
  从而平台值可重写为 \(\alpha_\ast+g_\ast\)，并且
  \[
  \boxed{\;
  \Gamma(\beta)=
  \begin{cases}
  \beta+g_\ast,& \beta_c<\beta<\alpha_\ast,\\[2mm]
  \alpha_\ast+g_\ast,& \beta\ge \alpha_\ast.
  \end{cases}\;}
  \]
\end{enumerate}
\end{theorem}
\begin{proof}
由定义
\[
\mathcal C_m^{\mathrm{cont}}(T)
=
\sum_{x\in X_m}\min\bigl(d_m(x),T\bigr).
\]
先取 \(\beta<\alpha_\ast\) 并令 \(T=e^{\beta m}\)。
由 \(M_m=e^{\alpha_\ast m+o(m)}\) 可知，对充分大的 \(m\)，最大纤维满足 \(M_m>T\)，因此
\[
\mathcal C_m^{\mathrm{cont}}(T)\ge K_mT.
\]
另一方面，
\[
\mathcal C_m^{\mathrm{cont}}(T)
\le
K_mT+\sum_{x\notin\mathcal M_m}d_m(x)
\le
K_mT+|X_m|\,M_m^{(2)}.
\]
由 \(|X_m|=F_{m+2}\) 可知
\(
\frac1m\log|X_m|\to\log\varphi
\)，
故第一项的指数为 \(g_\ast+\beta\)，第二项的指数至多为
\(\log\varphi+\alpha_\ast^{(2)}\)。
当 \(\beta>\beta_c=\log\varphi+\alpha_\ast^{(2)}-g_\ast\) 时，前者严格支配后者，于是
\[
\Gamma(\beta)=g_\ast+\beta.
\]
这证明了第 \(1\) 条。
\par
若 \(\beta>\alpha_\ast\)，则 \(e^{\beta m}>M_m\) 终将对所有纤维成立，从而
\[
\mathcal C_m^{\mathrm{cont}}(e^{\beta m})
=
\sum_{x\in X_m}d_m(x)
=
S_1(m).
\]
而 \(\sum_x d_m(x)=2^m\) 是折叠纤维基数的恒等式，故
\[
P(1)=\lim_{m\to\infty}\frac1m\log S_1(m)=\log 2,
\]
并得到第 \(2\) 条。
\par
最后设 \(a_c<1\)，即
\[
g_\ast+\alpha_\ast>\log\varphi+\alpha_\ast^{(2)}.
\]
注意
\[
S_1(m)
=
\sum_{x\in\mathcal M_m}M_m+\sum_{x\notin\mathcal M_m}d_m(x)
=
K_mM_m+\sum_{x\notin\mathcal M_m}d_m(x),
\]
于是
\[
K_mM_m\le S_1(m)\le K_mM_m+|X_m|\,M_m^{(2)}.
\]
第一项的指数为 \(g_\ast+\alpha_\ast\)，第二项的指数至多为
\(\log\varphi+\alpha_\ast^{(2)}\)，故在 \(a_c<1\) 下有
\[
P(1)=\lim_{m\to\infty}\frac1m\log S_1(m)=g_\ast+\alpha_\ast.
\]
结合第 \(2\) 条，这一平台值同时等于 \(\log 2\)。
再把该恒等式代回前两步的结论，便得到所示分段律。
\par
对边界点 \(\beta=\alpha_\ast\)，只需注意
\[
K_m\min\{M_m,e^{\alpha_\ast m}\}
\le
\mathcal C_m^{\mathrm{cont}}(e^{\alpha_\ast m})
\le
K_m e^{\alpha_\ast m}+|X_m|\,M_m^{(2)},
\]
其中
\(
\min\{M_m,e^{\alpha_\ast m}\}=e^{\alpha_\ast m+o(m)}
\)。
在 \(a_c<1\) 的支配关系下，两边皆具有指数 \(g_\ast+\alpha_\ast\)，故平台公式延伸到
\(\beta=\alpha_\ast\)。
证毕。
\end{proof}

\begin{theorem}[有限素数局部化数系之间的 \texorpdfstring{$\operatorname{Hom}$}{Hom} 分类]\label{thm:xi-localized-integers-hom-classification}
对有限素数集 \(S,T\subset\mathbb P\)，记
\[
G_S:=\ZZ[S^{-1}],
\qquad
G_T:=\ZZ[T^{-1}],
\]
并视它们为加法群。
由定理 \ref{thm:cdim-star-z1s-dual-extension}，
\[
0\to \ZZ\to G_S\to \bigoplus_{p\in S}\ZZ(p^\infty)\to 0
\]
以及
\[
0\to \prod_{p\in S}\ZZ_p\to \widehat{G_S}\to \TT\to 0
\]
皆为短正合列。
则有自然同构
\[
\boxed{\;
\operatorname{Hom}_{\mathbf{Ab}}(G_S,G_T)\cong
\begin{cases}
G_T,& S\subseteq T,\\[1mm]
0,& S\not\subseteq T,
\end{cases}\;}
\]
并且当 \(S\subseteq T\) 时，每个同态都唯一写成
\[
\varphi_u:G_S\to G_T,
\qquad
\varphi_u(x):=ux,
\qquad
u\in G_T.
\]
\(\operatorname{Hom}(G_S,G_T)\cong G_T\) 可由 \(\varphi_u\mapsto u=\varphi_u(1)\) 给出。
此外有：
\begin{enumerate}
  \item \(\varphi_u\neq 0\) 当且仅当 \(u\neq 0\)；若 \(u\neq 0\)，则 \(\varphi_u\) 必为单射；
  \item \(\varphi_u\) 为满射当且仅当 \(S=T\) 且 \(u\in G_T^\times\) 是局部化环 \(\ZZ[T^{-1}]\) 的单位。
\end{enumerate}
\end{theorem}
\leanverified{paper\_xi\_localized\_integers\_hom\_classification}
\begin{proof}
先设 \(S\subseteq T\)。
任取 \(u\in G_T\)，定义
\[
\varphi_u:G_S\to G_T,
\qquad
\varphi_u(x):=ux.
\]
由于 \(x\in G_S\) 的分母素数只来自 \(S\subseteq T\)，而 \(u\in G_T\) 的分母素数也只来自
\(T\)，故乘积 \(ux\) 仍属于 \(G_T\)，从而 \(\varphi_u\) 为良定义的加法同态。
\par
另一方面，任取 \(\varphi\in\operatorname{Hom}(G_S,G_T)\)，记 \(u:=\varphi(1)\)。
对任意
\[
x=\frac{a}{\prod_{p\in S}p^{k_p}}\in G_S
\]
都有
\[
\left(\prod_{p\in S}p^{k_p}\right)\varphi(x)=\varphi(a)=au.
\]
另一方面，
\[
\left(\prod_{p\in S}p^{k_p}\right)(ux)=au.
\]
故
\[
\left(\prod_{p\in S}p^{k_p}\right)\bigl(\varphi(x)-ux\bigr)=0.
\]
由于 \(G_T\subset\QQ\) 无挠，遂有
\[
\varphi(x)=ux.
\]
故每个同态都唯一等于某个 \(\varphi_u\)，因而
\(
\operatorname{Hom}(G_S,G_T)\cong G_T
\)。
\par
再设 \(S\not\subseteq T\)，取 \(p\in S\setminus T\)。
任意 \(\varphi:G_S\to G_T\) 记 \(r:=\varphi(1)\)。
由于在 \(G_S\) 中对每个 \(n\ge 1\) 都有 \(1=p^n\cdot p^{-n}\)，于是
\[
r=\varphi(1)=p^n\varphi(p^{-n}),
\]
故 \(r\in p^nG_T\) 对一切 \(n\ge 1\) 都成立。
现证明
\[
\bigcap_{n\ge 1}p^nG_T=\{0\}
\qquad (p\notin T).
\]
若 \(x\in G_T\setminus\{0\}\)，可写成既约分式
\[
x=\frac{a}{\prod_{q\in T}q^{m_q}}
\]
其中 \(a\in\ZZ\setminus\{0\}\)，且分母与 \(p\) 互素。
若 \(x\in p^nG_T\)，则存在 \(y_n\in G_T\) 使 \(x=p^n y_n\)；
把两边写成既约分式便知 \(p^n\mid a\)。
这不可能对所有 \(n\) 同时成立，故只能有 \(x=0\)。
因此 \(r=0\)。
\par
最后，对任意 \(x=a/b\in G_S\) 仍有
\[
b\,\varphi(x)=\varphi(a)=a\,r=0.
\]
由于 \(G_T\) 无挠，得 \(\varphi(x)=0\)。
故 \(\varphi=0\)，从而
\(
\operatorname{Hom}(G_S,G_T)=0
\)。
\par
现设 \(S\subseteq T\) 且 \(u\in G_T\)。
显然 \(\varphi_u=0\) 当且仅当 \(u=0\)。
若 \(u\neq 0\) 且 \(\varphi_u(x)=0\)，则 \(ux=0\)。
由于 \(G_S\subset\QQ\) 且 \(\QQ\) 无挠，乘以非零有理数不可能消去非零元素，故 \(x=0\)。
因此 \(\varphi_u\) 必为单射。
\par
最后判别满射性。
若 \(S=T\)，则 \(\varphi_u\) 满射当且仅当 \(uG_S=G_S\)，这等价于主理想 \((u)\) 等于整环
\(\ZZ[S^{-1}]\) 本身，亦即 \(u\in G_S^\times\)。
反之若 \(S\subsetneq T\)，取 \(q\in T\setminus S\)。
对任意 \(x\in G_S\)，由于 \(x\) 的分母中不含 \(q\)，故 \(ux\) 在 \(q\)-方向上的分母指数至多由
\(u\) 的 \(q\)-分母指数给出一个固定上界。
然而 \(G_T\) 中包含全部 \(q^{-N}\)（\(N\ge 1\)），其 \(q\)-分母指数无上界，故
\(uG_S\neq G_T\)。
于是 \(\varphi_u\) 不可能满射。
综上即得全部结论。
\end{proof}

\begin{corollary}[有限素数局部化系统的序结构由 \texorpdfstring{$\operatorname{Hom}$}{Hom} 与对偶商映射完全刻画]\label{cor:xi-localized-integers-order-dual-quotient}
\leanverified{paper\_xi\_localized\_integers\_order\_dual\_quotient}
沿用定理 \ref{thm:xi-localized-integers-hom-classification} 的记号。
对有限素数集 \(S,T\subset\mathbb P\)，下列命题等价：
\begin{enumerate}
  \item \(S\subseteq T\)；
  \item 存在非零群同态 \(G_S\to G_T\)；
  \item 存在单射群同态 \(G_S\hookrightarrow G_T\)；
  \item 存在连续满同态 \(\widehat{G_T}\twoheadrightarrow \widehat{G_S}\)；
  \item 存在连续满同态
  \[
  \prod_{p\in T}\ZZ_p\twoheadrightarrow \prod_{p\in S}\ZZ_p.
  \]
\end{enumerate}
特别地，
\[
G_S\cong G_T
\quad\Longleftrightarrow\quad
\widehat{G_S}\cong \widehat{G_T}
\quad\Longleftrightarrow\quad
S=T.
\]
\end{corollary}
\begin{proof}
\((1)\Rightarrow(3)\) 由自然包含 \(G_S\subseteq G_T\) 立即成立；
\((3)\Rightarrow(2)\) 显然；
\((2)\Rightarrow(1)\) 则由定理
\ref{thm:xi-localized-integers-hom-classification} 中
\(\operatorname{Hom}(G_S,G_T)\) 的分类直接给出。
\par
\((1)\Rightarrow(4)\)：由包含 \(G_S\hookrightarrow G_T\) 对偶化，
得到连续满同态
\(
\widehat{G_T}\twoheadrightarrow\widehat{G_S}
\)。
\((4)\Rightarrow(3)\)：Pontryagin 对偶把紧交换群的连续满射送到离散交换群的单射，
故对偶化后得到 \(G_S\hookrightarrow G_T\)。
\par
\((1)\Rightarrow(5)\) 由坐标投影
\[
\prod_{p\in T}\ZZ_p\longrightarrow \prod_{p\in S}\ZZ_p
\]
直接给出。
\par
\((5)\Rightarrow(1)\)：由定理 \ref{thm:cdim-star-prime-coordinate-orthogonality}，
任意连续同态
\(
\prod_{p\in T}\ZZ_p\to\prod_{p\in S}\ZZ_p
\)
都按素数坐标分解，且互异素数坐标之间不存在非零连续态射。
若某个 \(p\in S\setminus T\) 存在，则目标中的 \(p\)-坐标 \(\ZZ_p\) 不可能由源端产生非零像，
从而不可能得到满射，矛盾。
故必有 \(S\subseteq T\)。
\par
于是五条命题全部等价。
若 \(G_S\cong G_T\)，则同时有 \(S\subseteq T\) 与 \(T\subseteq S\)，故 \(S=T\)；
反向显然。
对偶系统的同构判别完全同理。
证毕。
\end{proof}

\begin{theorem}[有限素数局部化数系的张量乘法律]\label{thm:xi-localized-integers-tensor-union-law}
沿用定理 \ref{thm:xi-localized-integers-hom-classification} 的记号。则有自然同构
\[
\boxed{\;
G_S\otimes_{\ZZ}G_T\cong G_{S\cup T},
\qquad
\operatorname{Tor}_1^{\ZZ}(G_S,G_T)=0.\;}
\]
更一般地，对任意有限素数集族 \(S_1,\dots,S_r\subset\mathbb P\)，有
\[
\boxed{\;
G_{S_1}\otimes_{\ZZ}\cdots\otimes_{\ZZ}G_{S_r}
\cong
G_{S_1\cup\cdots\cup S_r}.\;}
\]
因此，有限素数集 \(S\mapsto G_S\) 在对象层不仅携带由包含关系控制的偏序态射结构，而且其并集运算恰对应于局部化数系的张量乘法。
\leanverified{paper\_xi\_localized\_integers\_tensor\_union\_law}
\end{theorem}
\begin{proof}
记
\[
M_S:=\left\{\prod_{p\in S}p^{k_p}:k_p\in\NN\right\},
\qquad
M_T:=\left\{\prod_{p\in T}p^{\ell_p}:\ell_p\in\NN\right\}.
\]
则 \(G_S=M_S^{-1}\ZZ\)、\(G_T=M_T^{-1}\ZZ\)。
局部化的标准张量公式给出
\[
M_S^{-1}\ZZ\otimes_{\ZZ}M_T^{-1}\ZZ
\cong
(M_SM_T)^{-1}\ZZ.
\]
而 \(M_SM_T\) 正是由 \(S\cup T\) 中素数生成的乘法集合，故
\[
G_S\otimes_{\ZZ}G_T\cong \ZZ[(S\cup T)^{-1}]=G_{S\cup T}.
\]
对 \(r\) 个因子的情形只需迭代应用同一公式即可。
\par
另一方面，局部化函子 \(M_S^{-1}(-)\) 对 \(\ZZ\)-模是精确的，故 \(G_S\) 为平坦
\(\ZZ\)-模。于是对任意 \(\ZZ\)-模 \(N\) 都有
\[
\operatorname{Tor}_1^{\ZZ}(G_S,N)=0.
\]
取 \(N=G_T\) 即得所示消失结论。证毕。
\end{proof}

\paragraph{有限素数支撑乘法单子与局部化数系的结构分界}
在有限素数支撑情形，乘法对象的结构层复杂度与局部化加法群的对偶分解可同时压缩为一组完全显式的秩公式、阻碍判据与内部对称分类。

\begin{theorem}[有限素数支撑乘法单子的结构层同态半圆维公式]\label{thm:xi-finite-prime-support-multiplicative-half-circle-dimension}
设 \(S\subset\mathbb P\) 为非空有限素数集，记
\[
\NN_S^\times:=\left\{\prod_{p\in S}p^{a_p}: a_p\in \ZZ_{\ge 0}\right\}\subset (\NN_{\ge 1},\times),
\qquad
r:=|S|.
\]
在附录 \(H\) 的结构层同态半圆维与实现层半圆维口径下，
\[
\cdim^{\mathrm{hom}}_{+}(\NN_S^\times)=\frac{r}{2},
\qquad
\cdim^{\mathrm{impl}}(\NN_S^\times)=\frac12.
\]
特别地，
\[
\cdim^{\mathrm{hom}}_{+}(\NN_S^\times)-\cdim^{\mathrm{impl}}(\NN_S^\times)
=\frac{r-1}{2},
\]
故当 \(r\ge 2\) 时，结构层与实现层之间存在严格分离。
\end{theorem}
\begin{proof}
定义估值向量映射
\[
v_S:\NN_S^\times\longrightarrow \ZZ_{\ge 0}^{S},
\qquad
n\longmapsto \bigl(v_p(n)\bigr)_{p\in S}.
\]
由唯一分解定理，\(v_S\) 为交换幺半群同构，因而
\[
\NN_S^\times\cong (\ZZ_{\ge 0}^{S},+)\cong (\ZZ_{\ge 0}^{r},+).
\]
于是
\(
\cdim^{\mathrm{hom}}_{+}(\NN_S^\times)\le r/2
\)。

反之，若存在单射幺半群同态
\[
\NN_S^\times\hookrightarrow (\NN^k,+),
\]
则由命题 \ref{prop:killo-grothendieck-completion-preserves-injection}，Grothendieck 完成后得到群单射
\[
\ZZ^r\cong G(\NN_S^\times)\hookrightarrow G(\NN^k)=\ZZ^k.
\]
自由阿贝尔群之间存在单射 \(\ZZ^r\hookrightarrow \ZZ^k\) 当且仅当 \(r\le k\)，故必有 \(k\ge r\)。因此
\(
\cdim^{\mathrm{hom}}_{+}(\NN_S^\times)\ge r/2
\)，从而第一式成立。

对实现层半圆维，只需注意 \(r\ge 1\) 时 \(\NN_S^\times\) 为可数无限集，故存在集合单射
\(
\NN_S^\times\hookrightarrow \NN
\)，于是
\(
\cdim^{\mathrm{impl}}(\NN_S^\times)\le 1/2
\)；另一方面 \(\cdim^{\mathrm{impl}}(X)=0\) 只可能发生在有限集上，故必有
\(
\cdim^{\mathrm{impl}}(\NN_S^\times)=1/2
\)。
两式相减即得差值公式。证毕。
\end{proof}

\begin{corollary}[有限素数截断乘法账本的线性寄存器律与全局不可有限化]\label{cor:xi-finite-prime-support-linear-register-law}
\leanverified{paper\_xi\_finite\_prime\_support\_linear\_register\_law}
记 \(p_1<p_2<\cdots\) 为递增素数序列，并对 \(r\ge 1\) 定义
\[
S_r:=\{p_1,\dots,p_r\},
\qquad
M_r:=\NN_{S_r}^{\times}.
\]
则：
\begin{enumerate}
  \item 任意单射幺半群同态
  \[
  \Phi_r:M_r\hookrightarrow (\NN^{k_r},+)
  \]
  都必满足
  \[
  k_r\ge r.
  \]
  因而不存在任何满足 \(k_r=o(r)\) 的精确同态化方案。
  \item 直系极限满足
  \[
  \varinjlim_{r}M_r=\NN_{\times},
  \]
  并且
  \[
  \sup_{r\ge 1}\cdim^{\mathrm{hom}}_{+}(M_r)
  =
  +\infty
  =
  \cdim^{\mathrm{hom}}_{+}(\NN_{\times}),
  \qquad
  \cdim^{\mathrm{impl}}(\NN_{\times})=\frac12.
  \]
\end{enumerate}
\end{corollary}
\begin{proof}
由定理 \ref{thm:xi-finite-prime-support-multiplicative-half-circle-dimension}，
\[
\cdim^{\mathrm{hom}}_{+}(M_r)=\frac{r}{2}.
\]
按定义，这等价于任意单射幺半群同态 \(M_r\hookrightarrow (\NN^{k_r},+)\) 都必须满足 \(k_r\ge r\)。因此 \(k_r=o(r)\) 不可能成立。

另一方面，每个正整数仅含有限个素因子，故
\[
\NN_{\times}=\bigcup_{r\ge 1}M_r
\]
且该并与自然包含映射一致地给出直系极限 \(\varinjlim_r M_r=\NN_{\times}\)。
又由上式，
\[
\sup_{r\ge 1}\cdim^{\mathrm{hom}}_{+}(M_r)
=
\sup_{r\ge 1}\frac{r}{2}
=+\infty.
\]
再结合命题 \ref{prop:killo-implementation-vs-structural-half-circle-dimension}，
\[
\cdim^{\mathrm{hom}}_{+}(\NN_{\times})=+\infty,
\qquad
\cdim^{\mathrm{impl}}(\NN_{\times})=\frac12.
\]
证毕。
\end{proof}

\begin{corollary}[有限素数支撑乘法半群与局部化数系的三层圆维分裂]\label{cor:xi-finite-prime-support-threeway-cdim-split}
设 \(S\subset\mathbb P\) 为非空有限素数集，记
\[
M_S:=\NN_S^\times
=
\left\{\prod_{p\in S}p^{a_p}:a_p\in\ZZ_{\ge 0}\right\},
\qquad
G_S:=\ZZ[S^{-1}],
\qquad
r:=|S|.
\]
则有精确公式
\[
\cdim^{\mathrm{impl}}(M_S)=\frac12,
\qquad
\cdim^{\mathrm{hom}}_{+}(M_S)=\frac{r}{2},
\qquad
\cdim_{\mathrm{fr}}(G_S)=1.
\]
\leanverified{paper\_xi\_finite\_prime\_support\_threeway\_cdim\_split}
特别地，当 \(r\ge 3\) 时，
\[
\cdim^{\mathrm{impl}}(M_S)
<
\cdim_{\mathrm{fr}}(G_S)
<
\cdim^{\mathrm{hom}}_{+}(M_S),
\]
并且结构层缺口满足
\[
\cdim^{\mathrm{hom}}_{+}(M_S)-\cdim_{\mathrm{fr}}(G_S)=\frac{r-2}{2}.
\]
\end{corollary}
\begin{proof}
前两式正是定理 \ref{thm:xi-finite-prime-support-multiplicative-half-circle-dimension} 的内容。第三式由定理 \ref{thm:killo-localization-circle-dimension-prime-ledger-splitting} 给出，因为
\[
G_S\otimes_{\ZZ}\QQ\cong \QQ
\]
具有有理秩 \(1\)。当 \(r\ge 3\) 时，
\[
\frac12<1<\frac{r}{2}
\]
立即推出所示严格不等式链；最后一式只是代入化简。证毕。
\end{proof}

\begin{theorem}[秩一无扭加法账本无法忠实承载两素数以上的乘法算术]\label{thm:xi-finite-prime-support-no-rankone-additive-host}
\leanverified{paper\_xi\_finite\_prime\_support\_no\_rankone\_additive\_host}
设 \(S\subset\mathbb P\) 为有限素数集，记 \(r:=|S|\ge 2\) 与 \(M_S:=\NN_S^\times\)。则对任意秩 \(1\) 的无扭加法群 \(A\)，都不存在单射幺半群同态
\[
M_S\hookrightarrow A.
\]
特别地，对任意有限素数集 \(T\subset\mathbb P\)，都不存在单射幺半群同态
\[
M_S\hookrightarrow G_T=\ZZ[T^{-1}].
\]
\end{theorem}
\begin{proof}
反设存在单射幺半群同态
\[
f:M_S\hookrightarrow A.
\]
由定理 \ref{thm:xi-finite-prime-support-multiplicative-half-circle-dimension}，\(M_S\cong \ZZ_{\ge 0}^{r}\)。
再由命题 \ref{prop:killo-grothendieck-completion-preserves-injection}，Grothendieck 完成后得到群单射
\[
\widetilde f:\ZZ^r\hookrightarrow G(A),
\]
其中 \(G(A)\) 由 \(A\) 的加法半群完成而来。由于 \(A\) 本身已是无扭加法群，故 \(G(A)=A\)。于是
\[
\ZZ^r\hookrightarrow A.
\]
然而 \(A\) 为秩 \(1\) 无扭群，故
\[
\dim_{\QQ}(A\otimes_{\ZZ}\QQ)=1,
\]
不可能含有秩 \(r\ge 2\) 的自由阿贝尔子群，矛盾。

当 \(A=G_T=\ZZ[T^{-1}]\) 时，由定理 \ref{thm:killo-localization-circle-dimension-prime-ledger-splitting} 或定理 \ref{thm:conclusion-cdim-finite-rank-extension} 可知其有理秩正为 \(1\)，故亦属上述情形。证毕。
\end{proof}

\begin{theorem}[单整数首坐标无法参与有限素数支撑乘法单子的精确加性账本]\label{thm:xi-single-integer-coordinate-no-exact-additive-ledger}
设 \(S\subset\mathbb P\) 为有限素数集且 \(|S|=r\ge 2\)。令
\[
e:\NN_S^\times\to \ZZ
\]
为任意单射。则不存在任何交换群 \(L\) 与映射
\[
s:\NN_S^\times\to L
\]
使得
\[
\Phi:\NN_S^\times\to \ZZ\oplus L,
\qquad
\Phi(n):=(e(n),s(n))
\]
同时满足：
\begin{enumerate}
  \item \(\Phi\) 为单射；
  \item 对一切 \(a,b\in \NN_S^\times\) 都有
  \[
  \Phi(ab)=\Phi(a)+\Phi(b).
  \]
\end{enumerate}
换言之，当独立素数方向至少为两个时，任何单整数首坐标都不可能成为精确乘法--加法账本的一部分。
\leanverified{paper\_xi\_single\_integer\_coordinate\_no\_exact\_additive\_ledger}
\end{theorem}
\begin{proof}
若上述 \(\Phi\) 存在，则取第一坐标即得
\[
e(ab)=e(a)+e(b)
\qquad
(\forall a,b\in \NN_S^\times),
\]
因而 \(e\) 本身已是交换幺半群同态。经由定理
\ref{thm:xi-finite-prime-support-multiplicative-half-circle-dimension} 证明中的估值同构，\(e\) 对应于一个加法幺半群同态
\[
\ell:\ZZ_{\ge 0}^{r}\to \ZZ.
\]
该同态唯一延拓为群同态
\[
\widetilde\ell:\ZZ^r\to \ZZ.
\]
由于 \(r\ge 2\)，\(\widetilde\ell\) 不可能为单射，故存在
\(
u\in \ker(\widetilde\ell)\setminus\{0\}
\)。
将 \(u\) 写成 Jordan 分解
\[
u=u^+-u^-,
\qquad
u^\pm\in \ZZ_{\ge 0}^{r},
\qquad
u^+\neq u^-.
\]
于是
\[
\ell(u^+)-\ell(u^-)=\widetilde\ell(u)=0,
\]
即 \(\ell(u^+)=\ell(u^-)\)。这说明 \(\ell\) 在 \(\ZZ_{\ge 0}^{r}\) 上不是单射，从而 \(e\) 也不是单射，矛盾。故所述 \((L,s)\) 不存在。
证毕。
\end{proof}

\begin{theorem}[局部化数系的自同态环与自同构群完全显式]\label{thm:xi-localized-integers-endomorphism-automorphism-explicit}
设 \(S\subset\mathbb P\) 为有限素数集，并记
\[
G_S:=\ZZ[S^{-1}]
=
\left\{
a\prod_{p\in S}p^{-k_p}
:
a\in\ZZ,\ k_p\in \ZZ_{\ge 0}
\right\}
\]
为其加法群。则每个群自同态 \(f:G_S\to G_S\) 都唯一地由某个 \(u\in \ZZ[S^{-1}]\) 给出，且
\[
f(x)=ux
\qquad
(\forall x\in G_S).
\]
因此存在自然环同构
\[
\operatorname{End}_{\ZZ}(G_S)\cong \ZZ[S^{-1}],
\]
并且
\[
\operatorname{Aut}_{\ZZ}(G_S)
\cong
\ZZ[S^{-1}]^\times
=
\left\{\pm\prod_{p\in S}p^{n_p}: n_p\in\ZZ\right\}
\cong
\{\pm1\}\times \ZZ^{r}.
\]
\end{theorem}
\begin{proof}
任取群自同态 \(f:G_S\to G_S\)，记
\(
u:=f(1)\in G_S
\)。
对任意
\[
x=\frac{a}{b}\in G_S,
\qquad
a\in\ZZ,
\qquad
b=\prod_{p\in S}p^{k_p}\in \NN,
\]
在加法群 \(G_S\) 中有关系 \(bx=a\)。对其施加 \(f\)，得到
\[
b\,f(x)=f(a)=a\,f(1)=au.
\]
由于 \(G_S\subset \QQ\) 为无挠群，上式在 \(\QQ\) 中唯一解为
\[
f(x)=\frac{a}{b}u=ux.
\]
故每个自同态都由乘以某个 \(u\in \ZZ[S^{-1}]\) 唯一给出；反之，任意 \(u\in \ZZ[S^{-1}]\) 显然定义一个自同态，故
\[
\operatorname{End}_{\ZZ}(G_S)\cong \ZZ[S^{-1}]
\]
作为环同构成立。

进一步，\(f_u(x)=ux\) 为自同构当且仅当 \(u\) 在局部化环 \(\ZZ[S^{-1}]\) 中可逆，即
\[
u\in \ZZ[S^{-1}]^\times
=
\left\{\pm\prod_{p\in S}p^{n_p}: n_p\in\ZZ\right\}.
\]
该单位群同构于 \(\{\pm1\}\times \ZZ^r\)，从而得到所述自同构群公式。
证毕。
\end{proof}
\leanverified{paper\_xi\_localized\_integers\_endomorphism\_automorphism\_explicit}

\begin{theorem}[局部化数系之间的全部加法同态由局部化乘法完全刻画]\label{thm:xi-localized-integers-cross-hom-classification}
设 \(S,T\subset\mathbb P\) 为有限素数集，并记
\[
G_S:=\ZZ[S^{-1}],
\qquad
G_T:=\ZZ[T^{-1}]
\]
为其加法群。则成立严格分类
\[
\operatorname{Hom}_{\ZZ}(G_S,G_T)\cong
\begin{cases}
0,& S\not\subseteq T,\\[1mm]
G_T,& S\subseteq T,
\end{cases}
\]
其中当 \(S\subseteq T\) 时，同构由
\[
x\in G_T\longmapsto m_x,
\qquad
m_x(y):=xy
\]
给出。
\end{theorem}
\begin{proof}
先设 \(S\subseteq T\)。任取 \(x\in G_T\)，对任意 \(y\in G_S\) 可写成
\[
y=\frac{a}{D},
\qquad
a\in\ZZ,
\qquad
D=\prod_{p\in S}p^{k_p}.
\]
由于 \(S\subseteq T\)，\(xy\) 的分母仍只含 \(T\)-素数，故 \(xy\in G_T\)。于是
\[
m_x:G_S\to G_T,\qquad y\mapsto xy
\]
为良定义加法同态。映射 \(x\mapsto m_x\) 显然保持加法，且由 \(m_x(1)=x\) 可知其单射。

现取任意 \(\phi\in \operatorname{Hom}_{\ZZ}(G_S,G_T)\)，记
\[
x:=\phi(1)\in G_T.
\]
对任意 \(y=a/D\in G_S\)，在加法群 \(G_S\) 中有关系 \(Dy=a\)。施加 \(\phi\) 得
\[
D\,\phi(y)=\phi(a)=ax.
\]
由于 \(G_T\subset \QQ\) 无挠，乘以非零整数 \(D\) 的映射单射，故
\[
\phi(y)=\frac{ax}{D}=xy=m_x(y).
\]
于是 \(\phi=m_x\)，从而 \(\operatorname{Hom}_{\ZZ}(G_S,G_T)\cong G_T\)。

再设 \(S\not\subseteq T\)。取某个素数
\[
p\in S\setminus T.
\]
任取 \(\phi\in \operatorname{Hom}_{\ZZ}(G_S,G_T)\)，并记 \(x:=\phi(1)\)。由于 \(p^{-k}\in G_S\) 对每个 \(k\ge 1\) 都成立，故
\[
p^k\,\phi(p^{-k})=\phi(1)=x,
\]
从而 \(x\in p^kG_T\) 对一切 \(k\ge 1\) 都成立。若 \(x\neq 0\)，写其既约形式为
\[
x=\frac{a}{b},
\qquad
a\in\ZZ\setminus\{0\},
\qquad
b\in\NN,
\]
且 \(b\) 的全部素因子都属于 \(T\)。由于 \(p\notin T\)，\(b\) 不被 \(p\) 整除。若 \(x\in p^kG_T\)，则存在 \(y_k=c_k/d_k\in G_T\) 使 \(x=p^k y_k\)，其中 \(d_k\) 的全部素因子仍属于 \(T\)。交叉相乘得到
\[
ad_k=p^kbc_k.
\]
因 \(p\nmid b\) 且 \(p\nmid d_k\)，故 \(p^k\mid a\)。此事对所有 \(k\) 同时成立，只能推出 \(a=0\)，矛盾。故
\[
\bigcap_{k\ge 1}p^kG_T=\{0\},
\qquad\text{从而}\qquad
x=0.
\]
最后任取 \(y=a/D\in G_S\)，仍由 \(Dy=a\) 得
\[
D\,\phi(y)=\phi(a)=ax=0.
\]
再由 \(G_T\) 无挠，得到 \(\phi(y)=0\)。故 \(\phi=0\)，于是
\[
\operatorname{Hom}_{\ZZ}(G_S,G_T)=0.
\]
证毕。
\end{proof}

\begin{corollary}[局部化数系及其对偶 solenoid 的全部连续可逆自同态恰由局部化单位给出]\label{cor:xi-localized-integers-dual-automorphism-units}
沿用上述记号。取 \(T=S\) 时，有自然环同构
\[
\operatorname{End}_{\ZZ}(G_S)\cong G_S,
\qquad
\phi\longmapsto \phi(1),
\]
并且
\[
\operatorname{Aut}_{\ZZ}(G_S)\cong G_S^\times
=
\left\{\pm\prod_{p\in S}p^{n_p}: n_p\in\ZZ,\ \text{仅有限多个 }n_p\neq 0\right\}.
\]
若记文稿中的 Pontryagin 对偶短正合列为
\[
0\longrightarrow \prod_{p\in S}\ZZ_p\longrightarrow \widehat{G_S}\longrightarrow \TT\longrightarrow 0,
\]
再由 Pontryagin 对偶，连续端同态环满足
\[
\operatorname{End}_{\mathrm{cts}}(\widehat{G_S})
\cong
\operatorname{End}_{\ZZ}(G_S)^{\mathrm{op}}
\cong
G_S,
\]
而连续自同构群满足
\[
\operatorname{Aut}_{\mathrm{cts}}(\widehat{G_S})\cong G_S^\times.
\]
因此，\(\widehat{G_S}\) 的全部连续可逆自同态都只对应于 \(S\)-局部化有理数的可逆乘法，不会额外产生新的内部自由度。
\leanverified{paper\_xi\_localized\_integers\_dual\_automorphism\_units}
\end{corollary}
\begin{proof}
取 \(T=S\) 时，定理 \ref{thm:xi-localized-integers-cross-hom-classification} 立刻给出
\[
\operatorname{End}_{\ZZ}(G_S)\cong G_S.
\]
其单位群正是定理 \ref{thm:xi-localized-integers-endomorphism-automorphism-explicit} 中所刻画的
\[
G_S^\times
=
\left\{\pm\prod_{p\in S}p^{n_p}: n_p\in\ZZ,\ \text{仅有限多个 }n_p\neq 0\right\},
\]
故 \(\operatorname{Aut}_{\ZZ}(G_S)\cong G_S^\times\)。

Pontryagin 对偶在局部紧交换群范畴中给出精确反变等价，因此
\[
\operatorname{End}_{\mathrm{cts}}(\widehat{G_S})
\cong
\operatorname{End}_{\ZZ}(G_S)^{\mathrm{op}}.
\]
又因 \(G_S\) 为交换环，反环与自身一致，故
\[
\operatorname{End}_{\mathrm{cts}}(\widehat{G_S})\cong G_S.
\]
单位在对偶等价下恰对应自同构，遂有
\[
\operatorname{Aut}_{\mathrm{cts}}(\widehat{G_S})\cong G_S^\times.
\]
证毕。
\end{proof}

\begin{corollary}[有限素数支撑同时决定局部化数系的嵌入偏序与内部对称]\label{cor:xi-localized-integers-order-endomorphism-rigidity}
\leanverified{paper\_xi\_localized\_integers\_order\_endomorphism\_rigidity}
对任意有限素数集 \(S,T\subset\mathbb P\)，有
\[
S\subseteq T
\iff
\exists\, \text{群单射 }G_S\hookrightarrow G_T
\iff
\exists\, \text{连续满同态 }\widehat{G_T}\twoheadrightarrow \widehat{G_S}.
\]
特别地，
\[
G_S\cong G_T
\iff
\widehat{G_S}\cong \widehat{G_T}
\iff
S=T.
\]
同时，
\[
\operatorname{End}_{\ZZ}(G_S)\cong \ZZ[S^{-1}],
\qquad
\operatorname{Aut}_{\ZZ}(G_S)\cong \ZZ[S^{-1}]^\times.
\]
因此，有限素数支撑 \(S\) 同时唯一决定了 \(G_S\) 的嵌入偏序、对偶满商偏序与全部内部对称。
\end{corollary}
\begin{proof}
偏序刻画正是推论 \ref{cor:xi-localized-integers-order-dual-quotient} 的内容，而端同态环与自同构群公式由定理 \ref{thm:xi-localized-integers-endomorphism-automorphism-explicit} 给出。合并即得。
\end{proof}

\begin{corollary}[无界独立素数方向迫使结构层同态半圆维发散]\label{cor:xi-unbounded-prime-directions-force-infinite-half-circle-dimension}
\leanverified{paper\_xi\_unbounded\_prime\_directions\_force\_infinite\_half\_circle\_dimension}
设 \(M\) 为交换可消去幺半群。若对任意 \(N\ge 1\)，都存在有限素数集 \(S_N\subset\mathbb P\) 满足 \(|S_N|=N\)，并存在单射幺半群同态
\[
\NN_{S_N}^\times\hookrightarrow M,
\]
则
\[
\cdim^{\mathrm{hom}}_{+}(M)=+\infty.
\]
\end{corollary}
\begin{proof}
反设 \(\cdim^{\mathrm{hom}}_{+}(M)<\infty\)。则存在某个有限 \(k\) 与单射幺半群同态
\[
M\hookrightarrow (\NN^k,+).
\]
与
\(
\NN_{S_N}^\times\hookrightarrow M
\)
复合，得到对每个 \(N\) 都存在单射幺半群同态
\[
\NN_{S_N}^\times\hookrightarrow (\NN^k,+).
\]
但由定理 \ref{thm:xi-finite-prime-support-multiplicative-half-circle-dimension}，
\[
\cdim^{\mathrm{hom}}_{+}(\NN_{S_N}^\times)=\frac{N}{2},
\]
故必有 \(N\le k\)。这与 \(N\) 可任意大矛盾。故
\(
\cdim^{\mathrm{hom}}_{+}(M)=+\infty
\)。
证毕。
\end{proof}
\paragraph{有限素数支撑乘法单子的结构层与实现层分离}
在有限素数支撑情形，结构层半圆维随独立素数方向线性增长，而实现层半圆维保持可数对象的极小值；因此两层只在单素数支撑时重合。

\begin{theorem}[有限素数支撑乘法单子的结构层与实现层精确分离]\label{thm:xi-finite-prime-support-halfdimension-gap-formula}
\leanverified{paper\_xi\_finite\_prime\_support\_halfdimension\_gap\_formula}
设 \(S\subset\mathbb P\) 为有限非空素数集，记
\[
M_S:=\NN_S^\times
=
\left\{\prod_{p\in S}p^{a_p}: a_p\in\ZZ_{\ge 0}\right\},
\qquad
r:=|S|.
\]
则
\[
\cdim^{\mathrm{hom}}_{+}(M_S)=\frac{r}{2},
\qquad
\cdim^{\mathrm{impl}}(M_S)=\frac12,
\qquad
\cdim^{\mathrm{hom}}_{+}(M_S)-\cdim^{\mathrm{impl}}(M_S)=\frac{r-1}{2}.
\]
特别地，
\[
\cdim^{\mathrm{hom}}_{+}(M_S)=\cdim^{\mathrm{impl}}(M_S)
\iff
r=1,
\]
而当 \(r\ge 2\) 时两层严格分离。
\end{theorem}
\begin{proof}
前两式正是定理 \ref{thm:xi-finite-prime-support-multiplicative-half-circle-dimension} 的内容。相减即得缺口公式。又
\[
\frac{r}{2}=\frac12
\iff
r=1,
\]
故所述等价与严格分离判据随即成立。证毕。
\end{proof}
\begin{theorem}[有限素数局部化群的可除素数恢复与拓扑分类刚性]\label{thm:xi-localized-integers-divisibility-topological-rigidity}
\leanverified{paper\_xi\_localized\_integers\_divisibility\_topological\_rigidity}
沿用定理 \ref{thm:xi-localized-integers-hom-classification} 的记号。对每个有限素数集
\[
S\subset\mathbb P,
\qquad
G_S:=\ZZ[S^{-1}],
\]
定义其可除素数集
\[
\operatorname{Div}(G_S):=\{p\in\mathbb P:\ pG_S=G_S\}.
\]
则有严格恒等式
\[
\boxed{\;
\operatorname{Div}(G_S)=S.\;}
\]
因此
\[
\boxed{\;
G_S\cong G_T
\quad\Longleftrightarrow\quad
\widehat{G_S}\cong \widehat{G_T}
\quad\Longleftrightarrow\quad
S=T.\;}
\]
\end{theorem}
\begin{proof}
任取素数 \(p\)。若 \(p\in S\)，则 \(1/p\in G_S\)。对任意 \(x\in G_S\)，有
\[
x=p\left(\frac{x}{p}\right),
\qquad
\frac{x}{p}\in G_S,
\]
故 \(pG_S=G_S\)，从而 \(p\in \operatorname{Div}(G_S)\)。

反之，若 \(pG_S=G_S\)，则存在某个 \(y\in G_S\) 使
\[
py=1.
\]
由于 \(G_S\subset\QQ\) 为无挠加法群，该方程在 \(\QQ\) 中的解唯一，故
\[
y=\frac1p\in G_S.
\]
这等价于 \(p\in S\)。于是 \(\operatorname{Div}(G_S)=S\)。

现在若 \(G_S\cong G_T\)，则可除素数集作为纯群论不变量保持不变，故
\[
S=\operatorname{Div}(G_S)=\operatorname{Div}(G_T)=T.
\]
反向显然。

再由 Pontryagin 对偶在离散交换群与紧交换群之间给出反变范畴等价，
\[
\widehat{G_S}\cong \widehat{G_T}
\quad\Longleftrightarrow\quad
G_S\cong G_T,
\]
故同样等价于 \(S=T\)。证毕。
\end{proof}

\leanverified{paper\_xi\_localized\_solenoid\_terminal\_object\_poset}
\begin{theorem}[有限局部化 solenoid 的反序满射格与 universal solenoid 的终对象性质]\label{thm:xi-localized-solenoid-terminal-object-poset}
对每个有限素数集 \(S\subset\mathbb P\)，记
\[
\Sigma_S:=\widehat{G_S}.
\]
则对任意有限素数集 \(S,T\subset\mathbb P\)，存在连续满同态
\[
\Sigma_S\twoheadrightarrow \Sigma_T
\]
当且仅当
\[
\boxed{\;T\subseteq S.\;}
\]
并且当 \(T\subseteq S\) 时，有规范短正合列
\[
\boxed{\;
0\longrightarrow \prod_{p\in S\setminus T}\ZZ_p
\longrightarrow \Sigma_S
\longrightarrow \Sigma_T
\longrightarrow 0.\;}
\]

进一步，记 universal solenoid
\[
\Sigma_{\mathbb P}:=\widehat{\QQ}.
\]
则对每个有限 \(S\subset\mathbb P\) 都存在规范满同态
\[
\Sigma_{\mathbb P}\twoheadrightarrow \Sigma_S
\]
且其核满足
\[
\boxed{\;
\ker\!\bigl(\Sigma_{\mathbb P}\twoheadrightarrow \Sigma_S\bigr)
\cong
\prod_{p\notin S}\ZZ_p.\;}
\]
因此 \(\Sigma_{\mathbb P}\) 在“有限局部化 solenoid 与连续满同态”所成范畴中是一个终对象；每个 \(\Sigma_S\) 都可由 \(\Sigma_{\mathbb P}\) 沿补集素数轴作规范商得到。
\end{theorem}
\begin{proof}
由推论 \ref{cor:xi-localized-integers-order-dual-quotient}，连续满同态
\[
\Sigma_S\twoheadrightarrow \Sigma_T
\]
存在当且仅当离散对偶端存在单射
\[
G_T\hookrightarrow G_S,
\]
而这又等价于 \(T\subseteq S\)。因此第一条判别成立。

当 \(T\subseteq S\) 时，定理 \ref{thm:xi-cdim-localization-quotient-solenoid-surjections} 已给出规范短正合列
\[
0\longrightarrow \prod_{p\in S\setminus T}\ZZ_p
\longrightarrow \Sigma_S
\longrightarrow \Sigma_T
\longrightarrow 0,
\]
这便得到核的精确描述。

最后固定有限 \(S\subset\mathbb P\)。由
\[
\QQ/\ZZ\cong \bigoplus_{p\in\mathbb P}C_{p^\infty},
\qquad
G_S/\ZZ\cong \bigoplus_{p\in S}C_{p^\infty},
\]
可得
\[
\QQ/G_S
\cong
(\QQ/\ZZ)/(G_S/\ZZ)
\cong
\bigoplus_{p\notin S}C_{p^\infty}.
\]
对该短正合列
\[
0\longrightarrow G_S\longrightarrow \QQ\longrightarrow \QQ/G_S\longrightarrow 0
\]
作 Pontryagin 对偶，便得到
\[
0\longrightarrow \prod_{p\notin S}\ZZ_p
\longrightarrow \widehat{\QQ}
\longrightarrow \widehat{G_S}
\longrightarrow 0.
\]
亦即
\[
0\longrightarrow \prod_{p\notin S}\ZZ_p
\longrightarrow \Sigma_{\mathbb P}
\longrightarrow \Sigma_S
\longrightarrow 0.
\]
故 \(\Sigma_{\mathbb P}\) 向每个有限局部化 solenoid 都有规范满射，且核正由补集素数方向给出。终对象性质随之成立。证毕。
\end{proof}
% (User input window)

\paragraph{整数点谱率、原子 Witt 同余分层与解析--RH 分离}
附录 \ref{app:pom-projective-pressure-operator} 已把整数 \(q\) 的 moment-kernel
\(\widetilde A_q\) 提升到实参数射影转移算子 \(\mathcal L_q\)；
附录 \ref{app:real-input-40-zeta-u} 与
\ref{app:real-input-40-finite-rh} 则把 primitive 的原子 Witt 分量、
解析半径、Dirichlet 扭曲与平方根误差窗口全部写成可核验的代数对象。
下列条目记录这些接口在当前文稿中的几条直接结构后果。

\begin{theorem}[整数 moment-kernel 谱半径满足次乘法律]\label{thm:xi-projective-moment-radius-submultiplicative}
沿用定义 \ref{def:pom-projective-transfer-operator}、
命题 \ref{prop:pom-projective-operator-invariant-homogeneous}
与定理 \ref{thm:pom-projective-operator-degenerates-to-moment-kernel} 的记号。
对每个整数 \(q\ge 2\)，记
\[
r_q:=\rho(\widetilde A_q)=\rho(\mathcal L_q|_{\mathcal H_q}).
\]
则对任意整数 \(a,b\ge 2\) 有
\[
\boxed{\ r_{a+b}\le r_a\,r_b\ }.
\]
\end{theorem}
\leanverified{paper\_xi\_projective\_moment\_radius\_submultiplicative}
\begin{proof}
对每个输出符号 \(c\in A_{\mathrm{out}}\)，定义线性算子
\[
\mathcal T_c^{(q)}:\mathcal H_q\longrightarrow \mathcal H_q,
\qquad
(\mathcal T_c^{(q)}p)(w):=p(B_cw).
\]
由命题 \ref{prop:pom-projective-operator-invariant-homogeneous}，
\[
\mathcal L_q|_{\mathcal H_q}=\sum_{c\in A_{\mathrm{out}}}\mathcal T_c^{(q)}.
\]
以下在各 \(\mathcal H_q\) 的单项式基下，不再区分这些算子与其矩阵表示。
\par
现在考虑乘法映射
\[
m_{a,b}:\mathcal H_a\otimes \mathcal H_b\longrightarrow \mathcal H_{a+b},
\qquad
m_{a,b}(f\otimes g):=fg.
\]
它是满射：任意次数 \(a+b\) 的单项式都可分解为一项次数 \(a\) 与一项次数 \(b\)
单项式之积。
在 \(\mathcal H_a\otimes\mathcal H_b\) 上定义
\[
\mathcal S_{a,b}:=\sum_{c\in A_{\mathrm{out}}}\mathcal T_c^{(a)}\otimes \mathcal T_c^{(b)}.
\]
则对纯张量 \(f\otimes g\) 有
\[
\begin{aligned}
m_{a,b}\!\bigl(\mathcal S_{a,b}(f\otimes g)\bigr)(w)
&=
\sum_{c\in A_{\mathrm{out}}} f(B_cw)\,g(B_cw)\\
&=
\bigl(\mathcal L_{a+b}(fg)\bigr)(w)
=
\bigl((\mathcal L_{a+b}|_{\mathcal H_{a+b}})\circ m_{a,b}\bigr)(f\otimes g)(w).
\end{aligned}
\]
故
\[
m_{a,b}\circ \mathcal S_{a,b}
=
(\mathcal L_{a+b}|_{\mathcal H_{a+b}})\circ m_{a,b}.
\]
于是 \(\mathcal H_{a+b}\) 是 \(\mathcal S_{a,b}\) 的一个商表示，从而
\[
\rho(\mathcal L_{a+b}|_{\mathcal H_{a+b}})\le \rho(\mathcal S_{a,b}).
\]
\par
固定 \(\varepsilon>0\)。
在 \(\mathcal H_q\) 的单项式基下记全 \(1\) 向量为 \(\mathbf 1_q\)，并定义
\[
x_q^{(\varepsilon)}
:=
\bigl((r_q+\varepsilon)I-(\mathcal L_q|_{\mathcal H_q})\bigr)^{-1}\mathbf 1_q
\qquad(q=a,b).
\]
由于 \(r_q+\varepsilon>\rho(\mathcal L_q|_{\mathcal H_q})\)，Neumann 级数给出
\[
x_q^{(\varepsilon)}
=
\sum_{n\ge 0}(r_q+\varepsilon)^{-n-1}
\bigl(\mathcal L_q|_{\mathcal H_q}\bigr)^n\mathbf 1_q,
\]
故 \(x_q^{(\varepsilon)}\) 的全部坐标严格为正，并满足
\[
\mathcal L_q|_{\mathcal H_q}\,x_q^{(\varepsilon)}
=(r_q+\varepsilon)x_q^{(\varepsilon)}-\mathbf 1_q
\le
(r_q+\varepsilon)x_q^{(\varepsilon)}.
\]
于是
\[
\begin{aligned}
\mathcal S_{a,b}\bigl(x_a^{(\varepsilon)}\otimes x_b^{(\varepsilon)}\bigr)
&=
\sum_c \mathcal T_c^{(a)}x_a^{(\varepsilon)}
\otimes
\mathcal T_c^{(b)}x_b^{(\varepsilon)}\\
&\le
\left(\sum_c \mathcal T_c^{(a)}x_a^{(\varepsilon)}\right)
\otimes
\left(\sum_c \mathcal T_c^{(b)}x_b^{(\varepsilon)}\right)\\
&=
\bigl(\mathcal L_a|_{\mathcal H_a}x_a^{(\varepsilon)}\bigr)
\otimes
\bigl(\mathcal L_b|_{\mathcal H_b}x_b^{(\varepsilon)}\bigr)\\
&\le
(r_a+\varepsilon)(r_b+\varepsilon)\,
x_a^{(\varepsilon)}\otimes x_b^{(\varepsilon)}.
\end{aligned}
\]
由于 \(x_a^{(\varepsilon)}\otimes x_b^{(\varepsilon)}\) 在坐标上严格为正，Collatz--Wielandt
原理给出
\[
\rho(\mathcal S_{a,b})\le (r_a+\varepsilon)(r_b+\varepsilon).
\]
合并上述两式并令 \(\varepsilon\downarrow 0\)，便得
\[
r_{a+b}
=
\rho(\mathcal L_{a+b}|_{\mathcal H_{a+b}})
\le
r_a r_b.
\]
再由定理 \ref{thm:pom-projective-operator-degenerates-to-moment-kernel}，
\(
\rho(\mathcal L_q|_{\mathcal H_q})=r_q
\)
对 \(q\ge 2\) 成立，故结论成立。证毕。
\end{proof}

\begin{corollary}[高阶谱斜率的 Fekete 极限存在]\label{cor:xi-projective-moment-fekete-slope}
\leanverified{paper\_xi\_projective\_moment\_fekete\_slope}
沿用定理 \ref{thm:xi-projective-moment-radius-submultiplicative} 的记号，定义
\[
\tau_\ast:=\lim_{q\to\infty}\frac{1}{q}\log r_q.
\]
则该极限存在，并且
\[
\boxed{\ \tau_\ast=\inf_{q\ge 2}\frac{1}{q}\log r_q\ }.
\]
\end{corollary}
\begin{proof}
记
\[
a_q:=\log r_q
\qquad(q\ge 2),
\]
并令
\[
L:=\inf_{q\ge 2}\frac{a_q}{q}.
\]
由定理 \ref{thm:xi-projective-moment-radius-submultiplicative}，
\[
a_{m+n}\le a_m+a_n
\qquad(m,n\ge 2).
\]
因此显然有
\[
\liminf_{q\to\infty}\frac{a_q}{q}\ge L.
\]
\par
任取 \(\varepsilon>0\)，可取某个 \(q_0\ge 2\) 使
\[
\frac{a_{q_0}}{q_0}<L+\varepsilon.
\]
对任意充分大的 \(q\)，取
\[
k:=\left\lfloor\frac{q-2}{q_0}\right\rfloor,
\qquad
s:=q-kq_0.
\]
则
\[
2\le s\le q_0+1.
\]
反复应用上式的次可加性，得到
\[
a_q\le k\,a_{q_0}+C_{q_0},
\qquad
C_{q_0}:=\max_{2\le s\le q_0+1}a_s.
\]
两边除以 \(q\) 得
\[
\frac{a_q}{q}
\le
\frac{kq_0}{q}\cdot \frac{a_{q_0}}{q_0}+\frac{C_{q_0}}{q}
\le
\frac{kq_0}{q}(L+\varepsilon)+\frac{C_{q_0}}{q}.
\]
令 \(q\to\infty\)，注意 \(kq_0/q\to 1\)，便得到
\[
\limsup_{q\to\infty}\frac{a_q}{q}\le L+\varepsilon.
\]
再令 \(\varepsilon\downarrow 0\)，遂有
\[
\limsup_{q\to\infty}\frac{a_q}{q}\le L.
\]
与前述 \(\liminf\) 界合并，即得极限存在且等于 \(L\)。证毕。
\end{proof}

\leanverified{paper\_xi\_atomic\_witt\_mod\_class\_localization}
\begin{theorem}[原子 Witt 因子对长度同余类的贡献严格局域于 \texorpdfstring{$2\bmod m$}{2 mod m}]\label{thm:xi-atomic-witt-mod-class-localization}
沿用命题 \ref{prop:atomic-2-orbit-split-logM} 的记号。
记
\[
P(z,u)=\sum_{n\ge 1}p_n(u)\,z^n,
\qquad
P_{\mathrm{core}}(z,u)=\sum_{n\ge 1}p_n^{\mathrm{core}}(u)\,z^n,
\]
并满足精确剥离
\[
P(z,u)=u z^2+P_{\mathrm{core}}(z,u).
\]
固定模数 \(m\ge 2\)。
对每个剩余类 \(a\in \ZZ/m\ZZ\)，定义 primitive 长度同余分量
\[
P^{[a]}(z,u)
:=
\sum_{\substack{n\ge 1\\ n\equiv a\;(\mathrm{mod}\,m)}}p_n(u)\,z^n,
\]
\[
P_{\mathrm{core}}^{[a]}(z,u)
:=
\sum_{\substack{n\ge 1\\ n\equiv a\;(\mathrm{mod}\,m)}}
p_n^{\mathrm{core}}(u)\,z^n.
\]
则有严格恒等式
\[
\boxed{\;
P^{[a]}(z,u)=\delta_{a,[2]}\,u z^2+P_{\mathrm{core}}^{[a]}(z,u)\; }.
\]
等价地，若以根单位滤波定义
\[
P^{[a]}(z,u)
=
\frac1m\sum_{j=0}^{m-1}\omega_m^{-aj}P(\omega_m^j z,u),
\qquad
\omega_m:=e^{2\pi i/m},
\]
则原子 Witt 因子只在剩余类 \(a\equiv 2\pmod m\) 中留下显式项 \(u z^2\)，
对其余全部长度同余类零泄漏。
\end{theorem}
\begin{proof}
由
\[
P(z,u)=u z^2+P_{\mathrm{core}}(z,u)
\]
直接按系数比较可得
\[
p_n(u)=
\begin{cases}
u+p_2^{\mathrm{core}}(u),& n=2,\\[2mm]
p_n^{\mathrm{core}}(u),& n\neq 2.
\end{cases}
\]
故对任意剩余类 \(a\in\ZZ/m\ZZ\)，由定义立即得到
\[
\begin{aligned}
P^{[a]}(z,u)
&=
\sum_{\substack{n\ge 1\\ n\equiv a\;(\mathrm{mod}\,m)}}p_n(u)\,z^n\\
&=
\delta_{a,[2]}\,u z^2
+\sum_{\substack{n\ge 1\\ n\equiv a\;(\mathrm{mod}\,m)}}
p_n^{\mathrm{core}}(u)\,z^n\\
&=
\delta_{a,[2]}\,u z^2+P_{\mathrm{core}}^{[a]}(z,u).
\end{aligned}
\]
\par
对根单位滤波形式，只需把分解代入即可：
\[
\frac1m\sum_{j=0}^{m-1}\omega_m^{-aj}u(\omega_m^j z)^2
=
u z^2\cdot
\frac1m\sum_{j=0}^{m-1}\omega_m^{(2-a)j}
=
\delta_{a,[2]}\,u z^2.
\]
于是 Fourier 投影与按系数分组两种表述完全一致。证毕。
\end{proof}

\begin{theorem}[解析但非 RH 的实轴相确实存在]\label{thm:xi-analytic-but-not-rh-phase-realinput40}
\leanverified{paper\_xi\_analytic\_but\_not\_rh\_phase\_realinput40}
沿用命题 \ref{prop:real-input-40-collision-pressure}、
推论 \ref{cor:real-input-40-collision-branch-radius}
与命题 \ref{prop:real-input-40-collision-rhk-window} 的记号。
记
\[
P_2(\theta):=\log\rho(e^\theta),
\qquad
R_\theta:=2.76230289346087\ldots,
\qquad
\theta_R:=\log u_R.
\]
则
\[
\boxed{\ 0<\theta_R<R_\theta\ }.
\]
特别地，开区间
\[
(\theta_R,R_\theta)\subset \RR_{>0}
\]
非空，并且对任意 \(\theta\in(\theta_R,R_\theta)\) 有：
\begin{enumerate}
  \item \(P_2(\theta)\) 仍位于 \(\theta=0\) 处 Perron 解析芽的收敛盘内部；
  \item 相应参数 \(u=e^\theta\) 已满足 \(u>u_R\)，故 kernel RH/Ramanujan 不等式失效。
\end{enumerate}
因此在真实输入 \(40\) 状态核的纯碰撞位势上，存在一个真正的
“解析但非 RH”实轴相。
\end{theorem}
\begin{proof}
由推论 \ref{cor:real-input-40-collision-branch-radius}，
\[
R_\theta\approx 2.76230289346087.
\]
由命题 \ref{prop:real-input-40-collision-rhk-window}，
\[
u_R\approx 3.59262181344,
\qquad
\theta_R=\log u_R\approx 1.27888.
\]
数值比较立即给出
\[
0<\theta_R<R_\theta.
\]
\par
若 \(\theta\in(\theta_R,R_\theta)\)，则首先 \(|\theta|<R_\theta\)，故
\(P_2(\theta)\) 仍处于 \(\theta=0\) 的解析收敛盘内。
另一方面，\(\theta>\theta_R\) 等价于 \(u=e^\theta>u_R\)，
而命题 \ref{prop:real-input-40-collision-rhk-window} 已证明
\[
\textsf{RH}(u)\Longleftrightarrow 0<u\le u_R.
\]
故当 \(u>u_R\) 时 kernel RH/Ramanujan 条件已失效。
两者合并即得结论。证毕。
\end{proof}

\begin{theorem}[典型密度附近的 Cram\'er germ 具有二次数域刚性，且右尾严格更昂贵]\label{thm:xi-real-input-cramer-germ-quadratic-field-rigidity}
沿用附录 \ref{app:real-input-40-zeta-u} 的输出位势记号，并定义归一化大偏差速率函数
\[
I(\alpha):=\sup_{\theta\in\RR}\bigl(\theta\alpha-(P(\theta)-P(0))\bigr).
\]
设
\[
\alpha_0:=P'(0),\qquad
\kappa_2:=P''(0),\qquad
\kappa_3:=P^{(3)}(0),\qquad
\kappa_4:=P^{(4)}(0).
\]
则表 \ref{tab:real_input_40_output_potential_cumulants_closed} 给出
\[
\alpha_0=\frac{23}{20}-\frac{7\sqrt5}{20},
\qquad
\kappa_2=\frac{1791}{200}-\frac{3983\sqrt5}{1000}>0,
\]
\[
\kappa_3=\frac{1416369}{5000}-\frac{126691\sqrt5}{1000}<0,
\qquad
\kappa_4=\frac{1354428303}{100000}-\frac{605718497\sqrt5}{100000}>0,
\]
并且这些常数全部属于 \(\QQ(\sqrt5)\)。
于是当 \(\delta\to 0\) 时，Legendre 对偶在 \(\alpha=\alpha_0\) 附近具有展开
\[
\boxed{\;
I(\alpha_0+\delta)=
\frac{\delta^2}{2\kappa_2}
-\frac{\kappa_3}{6\kappa_2^3}\delta^3
+\frac{3\kappa_3^2-\kappa_2\kappa_4}{24\kappa_2^5}\delta^4
+O(\delta^5).\;}
\]
特别地，前三个非平凡系数都属于 \(\QQ(\sqrt5)\)。
此外，由于 \(\kappa_3<0\)，有奇阶不对称律
\[
\boxed{\;
I(\alpha_0+\delta)-I(\alpha_0-\delta)
=
-\frac{\kappa_3}{3\kappa_2^3}\delta^3
+O(\delta^5)
>0
\qquad(\delta>0\ \text{充分小}).\;}
\]
从而对充分小的 \(\delta>0\) 还成立
\[
I(\alpha_0+\delta)>\frac{\delta^2}{2\kappa_2}>I(\alpha_0-\delta).
\]
故在典型密度附近，右尾偏离不仅严格比同幅度的左尾偏离更昂贵，而且已严格高于 Gaussian 二次主项。
\end{theorem}
\begin{proof}
把 \(P(\theta)-P(0)\) 在 \(\theta=0\) 处展开为
\[
P(\theta)-P(0)
=
\alpha_0\theta+\frac{\kappa_2}{2}\theta^2+\frac{\kappa_3}{6}\theta^3+\frac{\kappa_4}{24}\theta^4+O(\theta^5).
\]
Legendre 极值条件为
\[
\alpha_0+\delta=P'(\theta)
=
\alpha_0+\kappa_2\theta+\frac{\kappa_3}{2}\theta^2+\frac{\kappa_4}{6}\theta^3+O(\theta^4).
\]
由 \(\kappa_2>0\)，解析反函数定理给出 \(\theta=\theta(\delta)\) 在 \(\delta=0\) 邻域内解析存在。
将其逐项反解并代回
\[
I(\alpha_0+\delta)=\theta(\delta)(\alpha_0+\delta)-\bigl(P(\theta(\delta))-P(0)\bigr),
\]
即得到所示展开。
各系数由 \(\kappa_2,\kappa_3,\kappa_4\) 经过有理运算组成，因此仍属于 \(\QQ(\sqrt5)\)。

最后把 \(\delta\) 与 \(-\delta\) 的展开相减，偶次项相消，得到
\[
I(\alpha_0+\delta)-I(\alpha_0-\delta)
=
-\frac{\kappa_3}{3\kappa_2^3}\delta^3+O(\delta^5).
\]
由于 \(\kappa_3<0\)，主项为正，故对充分小的 \(\delta>0\) 有右尾代价严格大于左尾代价。再将 \(\pm\delta\) 代回前一展开式，并注意 \(-\kappa_3/(6\kappa_2^3)>0\)，便得与 Gaussian 二次主项的严格夹逼。证毕。
\end{proof}

\begin{theorem}[非平凡 Dirichlet 模态的谱半径强制同余类残差下界]\label{thm:xi-dirichlet-mode-forces-congruence-residual}
\leanverified{paper\_xi\_dirichlet\_mode\_forces\_congruence\_residual}
固定模数 \(m\ge 2\)，并对每个剩余类 \(a\in\ZZ/m\ZZ\) 定义 periodic 层的同余计数
\[
A_n(a;m):=\#\{\gamma:\ \gamma\text{ 为长度 }n\text{ 的闭路径},\ E(\gamma)\equiv a\pmod m\},
\qquad
A_n:=\sum_{a\in\ZZ/m\ZZ}A_n(a;m).
\]
记 \(\omega_m:=e^{2\pi i/m}\)，并对每个 \(1\le j\le m-1\) 定义
\[
\rho_{m,j}:=\rho\!\bigl(M(\omega_m^j)\bigr),
\qquad
\lambda:=\rho(M(1))=\varphi^2.
\]
则对每个非平凡模态 \(1\le j\le m-1\)，均有
\[
\boxed{\;
\limsup_{n\to\infty}
\Bigl(
\max_{a\in\ZZ/m\ZZ}\Bigl|A_n(a;m)-\frac{A_n}{m}\Bigr|
\Bigr)^{1/n}
\ge
\rho_{m,j}.\;}
\]
因而
\[
\boxed{\;
\limsup_{n\to\infty}
\Bigl(
\max_{a\in\ZZ/m\ZZ}\Bigl|A_n(a;m)-\frac{A_n}{m}\Bigr|
\Bigr)^{1/n}
\ge
\max_{1\le j\le m-1}\rho_{m,j}.\;}
\]
特别地，若该模数满足
\[
\theta_m:=\max_{1\le j\le m-1}\frac{\log\rho_{m,j}}{\log\lambda}>\frac12,
\]
则存在某个剩余类 \(a_\ast\in\ZZ/m\ZZ\) 使
\[
\limsup_{n\to\infty}
\frac{1}{n\log\lambda}
\log\Bigl|A_n(a_\ast;m)-\frac{A_n}{m}\Bigr|
\ge
\theta_m
>
\frac12.
\]
因此该模数下的同余类分布不可能满足任何 RH 形的平方根指数误差界。
\end{theorem}
\begin{proof}
对每个 \(j\neq 0\)，离散 Fourier 反演给出
\[
\widehat A_n(j)
:=
\sum_{a\in\ZZ/m\ZZ}A_n(a;m)\omega_m^{ja}
=
a_n(\omega_m^j)
=
\Tr\!\bigl(M(\omega_m^j)^n\bigr),
\]
其中最后一步正是命题 \ref{prop:real-input-40-output-dirichlet-nonrh} 中 periodic 层的识别。

设 \(M(\omega_m^j)\) 的特征值为 \(\lambda_1,\dots,\lambda_r\)，并设
\[
|\lambda_1|=\cdots=|\lambda_s|=\rho_{m,j}>|\lambda_{s+1}|\ge \cdots \ge |\lambda_r|.
\]
令 \(\xi_i:=\lambda_i/\rho_{m,j}\in\TT\)（\(1\le i\le s\)）。
则
\[
\widehat A_n(j)
=
\sum_{i=1}^{r}\lambda_i^n
=
\rho_{m,j}^n\Bigl(\sum_{i=1}^{s}\xi_i^n\Bigr)+O((\rho_{m,j}-\varepsilon)^n)
\]
对某个 \(\varepsilon>0\) 成立。
由于序列 \((\xi_1^n,\dots,\xi_s^n)\) 在紧群 \(\TT^s\) 中的闭包包含单位元，存在子序列 \(n_r\to\infty\) 使
\[
\xi_i^{n_r}\to 1
\qquad(1\le i\le s).
\]
故沿该子序列有
\[
\widehat A_{n_r}(j)=s\,\rho_{m,j}^{n_r}+o(\rho_{m,j}^{n_r}),
\]
从而
\[
\limsup_{n\to\infty}|\widehat A_n(j)|^{1/n}=\rho_{m,j}.
\]

现在定义偏差
\[
b_n(a;m):=A_n(a;m)-\frac{A_n}{m}.
\]
则对 \(j\neq 0\) 有
\[
\widehat b_n(j)
:=
\sum_{a\in\ZZ/m\ZZ}b_n(a;m)\omega_m^{ja}
=
\widehat A_n(j),
\]
因为 \(\sum_a \omega_m^{ja}=0\)。
于是
\[
|\widehat A_n(j)|
=
\Bigl|\sum_{a\in\ZZ/m\ZZ}b_n(a;m)\omega_m^{ja}\Bigr|
\le
\sum_{a\in\ZZ/m\ZZ}|b_n(a;m)|
\le
m\max_{a\in\ZZ/m\ZZ}|b_n(a;m)|.
\]
取 \(n\to\infty\) 的上极限并开 \(n\) 次方，即得
\[
\limsup_{n\to\infty}
\Bigl(
\max_{a\in\ZZ/m\ZZ}\Bigl|A_n(a;m)-\frac{A_n}{m}\Bigr|
\Bigr)^{1/n}
\ge
\rho_{m,j}.
\]
再对 \(j\neq 0\) 取最大即得第二式。

若 \(\theta_m>1/2\)，则 \(\max_{j\neq 0}\rho_{m,j}=\lambda^{\theta_m}\)。
又由于剩余类只有有限个，恒有
\[
\limsup_{n\to\infty}\max_{a} c_{n,a}\le \max_{a}\limsup_{n\to\infty} c_{n,a}
\]
对任意非负数组 \(c_{n,a}\) 成立。
应用于
\[
c_{n,a}:=\Bigl|A_n(a;m)-\frac{A_n}{m}\Bigr|^{1/n},
\]
可知存在某个 \(a_\ast\) 使相应单类偏差的上极限至少为 \(\lambda^{\theta_m}\)。
对数除以 \(n\log\lambda\) 即得最后一式。证毕。
\end{proof}

\begin{theorem}[Dirichlet 扭曲 Fourier 级数的收敛半径恰由谱半径决定，故平方根解析窗必然破裂]\label{thm:xi-output-potential-dirichlet-fourier-radius-barrier}
\leanverified{paper\_xi\_output\_potential\_dirichlet\_fourier\_radius\_barrier}
沿用命题 \ref{prop:real-input-40-output-dirichlet-nonrh} 的记号。固定整数 \(m\ge 2\) 与 \(1\le j\le m-1\)，记
\[
M_j:=M(\omega_m^j),
\qquad
\rho_{m,j}:=\rho(M_j),
\qquad
a_n(\omega_m^j):=\Tr(M_j^n).
\]
定义普通生成函数
\[
A_{m,j}(z):=\sum_{n\ge 1}a_n(\omega_m^j)\,z^n.
\]
则 \(A_{m,j}(z)\) 为有理函数，且其收敛半径恰为
\[
\boxed{\;R_{m,j}=\frac{1}{\rho_{m,j}}.\;}
\]
特别地，若某个模数 \(m\) 满足
\[
\theta_m:=\max_{1\le j\le m-1}\frac{\log\rho_{m,j}}{\log\lambda}>\frac12,
\qquad
\lambda:=\rho(M(1))=\varphi^2,
\]
则存在 \(j\) 使
\[
R_{m,j}<\lambda^{-1/2}.
\]
等价地，归一化级数
\[
A_{m,j}(\lambda^{-1/2}z)
\]
在单位圆内部已出现极点。因而不存在统一的平方根型解析窗能够同时覆盖全部非平凡 Dirichlet 扭曲模态。
\end{theorem}
\begin{proof}
在代数闭包上把 \(M_j\) 的特征值按代数重数记为
\[
\eta_1,\dots,\eta_{40}.
\]
则对每个 \(n\ge 1\) 都有
\[
a_n(\omega_m^j)=\Tr(M_j^n)=\sum_{\ell=1}^{40}\eta_\ell^n.
\]
故在 \(|z|<1/\rho_{m,j}\) 内可逐项求和，得到
\[
\begin{aligned}
A_{m,j}(z)
&=
\sum_{n\ge 1}\sum_{\ell=1}^{40}\eta_\ell^n z^n\\
&=
\sum_{\ell=1}^{40}\sum_{n\ge 1}(\eta_\ell z)^n\\
&=
\sum_{\ell=1}^{40}\frac{\eta_\ell z}{1-\eta_\ell z}.
\end{aligned}
\]
因此 \(A_{m,j}\) 为有理函数。其全部非零极点正位于
\[
z=\eta_\ell^{-1}
\qquad(\eta_\ell\neq 0).
\]
若某个非零特征值 \(\eta\) 的代数重数为 \(r\ge 1\)，则对应分式项合并为
\[
\frac{r\,\eta z}{1-\eta z},
\]
其留数非零，故该极点不会被消去。于是 \(A_{m,j}\) 的最近极点到原点的距离恰为
\[
\min_{\eta_\ell\neq 0}\frac{1}{|\eta_\ell|}
=
\frac{1}{\max_\ell|\eta_\ell|}
=
\frac{1}{\rho_{m,j}},
\]
这就证明了
\[
R_{m,j}=\frac{1}{\rho_{m,j}}.
\]

现若 \(\theta_m>1/2\)，则由定义存在某个 \(j\) 使
\[
\rho_{m,j}=\lambda^{\theta_{m,j}}
>
\lambda^{1/2}.
\]
于是
\[
R_{m,j}=\rho_{m,j}^{-1}<\lambda^{-1/2}.
\]
把变量按 \(\lambda^{-1/2}\) 归一化后，级数 \(A_{m,j}(\lambda^{-1/2}z)\) 的收敛半径变为
\[
\lambda^{1/2}R_{m,j}=\frac{\lambda^{1/2}}{\rho_{m,j}}<1.
\]
由于该函数为有理函数，收敛半径严格小于 \(1\) 等价于单位圆内部已有极点。故任何试图以 \(\lambda^{n/2}\) 作为统一尺度的解析窗，在该 Dirichlet 模态上都必然失败。证毕。
\end{proof}

\paragraph{射影压力的凸包原则与 residue Fourier 指数的精确闭包}
附录 \ref{app:pom-projective-pressure-operator} 的实参数射影转移算子 \(\mathcal L_q\) 与命题 \ref{prop:real-input-40-output-dirichlet-nonrh} 的 cyclotomic 扭曲谱之间，还可进一步抽取出下列彼此衔接的结构结论。前三条先把整数阶压力样本压缩为有限谱证书，并给出连续压力对离散 moment-kernel 数据的无条件凸控制；中间部分把命题 \ref{prop:real-input-40-output-dirichlet-nonrh} 与定理 \ref{thm:xi-dirichlet-mode-forces-congruence-residual} 中的上界链与下界链闭合为精确指数定理，并将平方根障碍直接提升为模分布局部极限定理的失效；最后一条则把前述凸性整理为离散压力数据的规范上包络表述。

\begin{theorem}[射影压力算子族的路径级 Hölder 不等式与对数凸压力]\label{thm:xi-projective-pressure-path-holder-convexity}
\leanverified{paper\_xi\_projective\_pressure\_path\_holder\_convexity}
沿用定义 \ref{def:pom-projective-transfer-operator} 的记号，把 \(\mathcal L_q\) 视为 \(C(\Delta)\) 上的正算子。对 \(q\ge 1\) 与 \(n\ge 1\)，定义
\[
\Lambda_n(q):=\log \|\mathcal L_q^n\mathbf 1\|_\infty,
\qquad
\Lambda(q):=\log \rho(\mathcal L_q).
\]
则：
\begin{enumerate}
  \item 对每个固定 \(n\ge 1\)，函数 \(q\mapsto \Lambda_n(q)\) 在 \([1,\infty)\) 上凸；
  \item 因而 \(q\mapsto \Lambda(q)\) 在 \([1,\infty)\) 上凸；
  \item 对每个整数 \(q\ge 2\)，有
  \[
  \Lambda(q)\ge \log r_q.
  \]
\end{enumerate}
\end{theorem}
\begin{proof}
对任意输出词
\(
\beta=(b_1,\dots,b_n)\in A_{\mathrm{out}}^n
\)
与 \(w\in\Delta\)，递归定义
\[
w_0:=w,
\qquad
w_t:=\Phi_{b_t}(w_{t-1})\quad(1\le t\le n),
\]
并记路径权
\[
\Gamma_\beta(w):=\prod_{t=1}^n g_{b_t}(w_{t-1})\ge 0.
\]
由定义反复展开 \(\mathcal L_q\) 得
\[
(\mathcal L_q^n\mathbf 1)(w)=\sum_{\beta\in A_{\mathrm{out}}^n}\Gamma_\beta(w)^q.
\]
现取 \(q_0,q_1\ge 1\)、\(\theta\in[0,1]\)，并记
\(
q_\theta:=(1-\theta)q_0+\theta q_1
\)。
对固定 \(w\)，上式右端是有限个非负指数项之和，故由 Hölder 不等式（端点 \(\theta=0,1\) 为平凡情形）有
\[
\sum_{\beta}\Gamma_\beta(w)^{q_\theta}
\le
\left(\sum_{\beta}\Gamma_\beta(w)^{q_1}\right)^\theta
\left(\sum_{\beta}\Gamma_\beta(w)^{q_0}\right)^{1-\theta}.
\]
取上确界便得
\[
\|\mathcal L_{q_\theta}^n\mathbf 1\|_\infty
\le
\|\mathcal L_{q_1}^n\mathbf 1\|_\infty^\theta
\|\mathcal L_{q_0}^n\mathbf 1\|_\infty^{1-\theta},
\]
从而 \(\Lambda_n\) 凸。

另一方面，\(\mathcal L_q\) 为正算子，故对任意 \(f\in C(\Delta)\) 满足
\[
|\mathcal L_q^n f|
\le
\|f\|_\infty\,\mathcal L_q^n\mathbf 1,
\]
于是
\[
\|\mathcal L_q^n\|_{\infty\to\infty}
=
\|\mathcal L_q^n\mathbf 1\|_\infty.
\]
由 Gelfand 谱半径公式，
\[
\Lambda(q)
=
\log \rho(\mathcal L_q)
=
\lim_{n\to\infty}\frac{1}{n}\log \|\mathcal L_q^n\|_{\infty\to\infty}
=
\lim_{n\to\infty}\frac{1}{n}\Lambda_n(q).
\]
凸函数列 \(n^{-1}\Lambda_n\) 的逐点极限仍为凸函数，故 \(\Lambda\) 在 \([1,\infty)\) 上凸。

最后，由定理 \ref{thm:pom-projective-operator-degenerates-to-moment-kernel}，
\[
\rho(\mathcal L_q|_{\mathcal H_q})=\rho(\widetilde A_q)=r_q
\qquad(q\in\mathbb Z_{\ge 2}),
\]
而限制算子的谱半径不超过全空间算子谱半径，故
\[
\log r_q
=
\log\rho(\mathcal L_q|_{\mathcal H_q})
\le
\log\rho(\mathcal L_q)
=
\Lambda(q).
\]
证毕。
\end{proof}

\begin{theorem}[整数阶压力样本由有限谱证书支配并且皆为代数整数的对数]\label{thm:xi-projective-pressure-integer-spectral-certificate}
沿用定理 \ref{thm:pom-projective-operator-degenerates-to-moment-kernel} 的记号。对每个整数 \(q\ge 2\)，记
\[
\widetilde A_q\in M_{d_q}(\ZZ_{\ge 0}),
\qquad
r_q:=\rho(\widetilde A_q),
\qquad
\Lambda_{\mathrm{int}}(q):=\log r_q.
\]
则成立：
\begin{enumerate}
  \item \(r_q\) 是有限整矩阵 \(\widetilde A_q\) 的 Perron 根，故 \(r_q\) 为代数数；
  \item 若记
  \[
  \chi_q(\lambda):=\det(\lambda I-\widetilde A_q),
  \]
  则 \(r_q\) 是 \(\chi_q\) 的最大正根，从而整数样本 \(\Lambda_{\mathrm{int}}(q)\) 由有限谱证书 \((\chi_q,r_q)\) 完全支配；
  \item 由于 \(\widetilde A_q\) 具有整系数条目，\(r_q\) 实为代数整数。
\end{enumerate}
特别地，整数点压力样本并非仅以极限定义隐式存在，而是落在一族有限维、整系数、可审计的代数谱对象上。若把连续归一化压力写作
\[
\Lambda(t)=\log\lambda(t+1)-\log|A_{\mathrm{out}}|,
\]
则在整数点 \(t=q-1\) 处有
\[
\Lambda(q-1)=\Lambda_{\mathrm{int}}(q)-\log|A_{\mathrm{out}}|.
\]
\end{theorem}
\begin{proof}
由定理 \ref{thm:pom-projective-operator-degenerates-to-moment-kernel}，
\[
\bigl[\mathcal L_q|_{\mathcal H_q}\bigr]_{(w^\nu)}=\widetilde A_q,
\qquad
\rho(\mathcal L_q|_{\mathcal H_q})=\rho(\widetilde A_q)=r_q.
\]
因此 \(r_q\) 是有限矩阵 \(\widetilde A_q\) 的特征值，立即推出 \(r_q\) 为代数数。

再记
\[
\chi_q(\lambda)=\det(\lambda I-\widetilde A_q).
\]
按定义，\(r_q\) 是 \(\chi_q\) 的某个实正根。又由于 \(\widetilde A_q\) 为非负矩阵，Perron--Frobenius 理论表明其谱半径正好是最大正实特征值，故 \(r_q\) 也是 \(\chi_q\) 的最大正根。于是整数点压力样本
\[
\Lambda_{\mathrm{int}}(q)=\log r_q
\]
便可由有限谱证书 \((\chi_q,r_q)\) 直接恢复；连续归一化压力在整数点的取值只需再减去常数 \(\log|A_{\mathrm{out}}|\)。

最后，由注 \ref{rem:pom-moment-symmetric-power-substitution-operator} 与定理 \ref{thm:pom-projective-operator-degenerates-to-moment-kernel} 的构造，\(\widetilde A_q\) 来自整系数线性替换在单项式基上的矩阵表示，故 \(\widetilde A_q\) 具有整数条目。于是其特征多项式 \(\chi_q\) 为首一整系数多项式，从而全部特征值皆为代数整数，特别 \(r_q\) 为代数整数。证毕。
\end{proof}
\leanverified{paper\_xi\_projective\_pressure\_integer\_spectral\_certificate}
\begin{theorem}[高阶碰撞相的齐次多项式 Perron 极值原理]\label{thm:xi-projective-moment-homogeneous-polynomial-perron-extremizer}
沿用命题 \ref{prop:pom-projective-operator-invariant-homogeneous} 与定理 \ref{thm:pom-projective-operator-degenerates-to-moment-kernel} 的记号。对每个整数 \(q\ge 2\)，存在一个非零齐次多项式
\[
p_q\in \mathcal H_q
\]
且其单项式系数全为非负，使得
\[
(\mathcal L_q|_{\mathcal H_q})p_q=r_q\,p_q,
\qquad
r_q=\rho(\widetilde A_q).
\]
换言之，每一个整数阶碰撞谱的 Perron 极值都可在有限维齐次多项式扇区中实现，而无需引入超出多项式范畴的测试函数族。
\end{theorem}
\begin{proof}
由定理 \ref{thm:pom-projective-operator-degenerates-to-moment-kernel}，\(\mathcal L_q\) 在 \(\mathcal H_q\) 上的限制在单项式基下恰由非负矩阵 \(\widetilde A_q\) 表示，并且
\[
\rho(\mathcal L_q|_{\mathcal H_q})=\rho(\widetilde A_q)=r_q.
\]
由 Perron--Frobenius 定理，\(\widetilde A_q\) 存在一个非零非负特征向量 \(v_q\)，满足
\[
\widetilde A_qv_q=r_qv_q.
\]
把 \(v_q\) 重新解释为 \(\mathcal H_q\) 的单项式基坐标，即得到一个齐次多项式 \(p_q\in \mathcal H_q\)，其系数全为非负，且满足
\[
(\mathcal L_q|_{\mathcal H_q})p_q=r_q\,p_q.
\]
证毕。
\end{proof}

\begin{theorem}[模 \texorpdfstring{$m$}{m} 输出计数的精确 Fourier--Parseval 反演公式]\label{thm:xi-dirichlet-residue-fourier-parseval-exact}
固定整数 \(m\ge 2\)，并记 \(\omega_m:=e^{2\pi i/m}\)。对每个 \(r\in\ZZ/m\ZZ\) 与 \(n\ge 1\)，定义 periodic 与 primitive 的 residue counts
\[
A_{m,r}(n):=\#\{\gamma:\ |\gamma|=n,\ E(\gamma)\equiv r\pmod m\},
\]
\[
P_{m,r}(n):=\#\{\mathfrak p:\ |\mathfrak p|=n,\ E(\mathfrak p)\equiv r\pmod m\}.
\]
则成立精确 Fourier 反演公式
\[
\boxed{\;
A_{m,r}(n)=\frac1m\sum_{j=0}^{m-1}a_n(\omega_m^j)\,\omega_m^{-jr},\;}
\]
\[
\boxed{\;
P_{m,r}(n)=\frac1m\sum_{j=0}^{m-1}p_n(\omega_m^j)\,\omega_m^{-jr}.\;}
\]
并且有 Parseval 恒等式
\[
\boxed{\;
\sum_{r\in\ZZ/m\ZZ}
\left|A_{m,r}(n)-\frac{a_n(1)}{m}\right|^2
=
\frac1m\sum_{j=1}^{m-1}|a_n(\omega_m^j)|^2,\;}
\]
\[
\boxed{\;
\sum_{r\in\ZZ/m\ZZ}
\left|P_{m,r}(n)-\frac{p_n(1)}{m}\right|^2
=
\frac1m\sum_{j=1}^{m-1}|p_n(\omega_m^j)|^2.\;}
\]
\end{theorem}
\begin{proof}
由命题 \ref{prop:real-input-40-output-dirichlet-nonrh}，
\[
a_n(\omega_m^j)
=
\sum_{\gamma:\ |\gamma|=n}\omega_m^{jE(\gamma)}
=
\sum_{r\in\ZZ/m\ZZ}A_{m,r}(n)\,\omega_m^{jr}.
\]
同理，primitive 层的权和满足
\[
p_n(\omega_m^j)
=
\sum_{\mathfrak p:\ |\mathfrak p|=n}\omega_m^{jE(\mathfrak p)}
=
\sum_{r\in\ZZ/m\ZZ}P_{m,r}(n)\,\omega_m^{jr},
\]
这正是同一有限循环群上的 Fourier 变换。对 \(\ZZ/m\ZZ\) 应用离散 Fourier 反演，立得前两式。

再定义中心化序列
\[
\widetilde A_{m,r}(n):=A_{m,r}(n)-\frac{a_n(1)}{m},
\qquad
\widetilde P_{m,r}(n):=P_{m,r}(n)-\frac{p_n(1)}{m}.
\]
它们在 \(j=0\) 频率处的 Fourier 系数都为零，而在 \(j\neq 0\) 处分别恰为 \(a_n(\omega_m^j)\) 与 \(p_n(\omega_m^j)\)。于是有限循环群上的 Parseval 恒等式给出
\[
\sum_r |\widetilde A_{m,r}(n)|^2
=
\frac1m\sum_{j=1}^{m-1}|a_n(\omega_m^j)|^2,
\qquad
\sum_r |\widetilde P_{m,r}(n)|^2
=
\frac1m\sum_{j=1}^{m-1}|p_n(\omega_m^j)|^2.
\]
证毕。
\end{proof}

\begin{theorem}[twisted trace 与 residue bias 的最坏模态具有相同精确指数]\label{thm:xi-dirichlet-residue-bias-exponent-exact}
固定整数 \(m\ge 2\) 与 \(1\le j\le m-1\)。记
\[
a_n^{(j)}:=a_n(\omega_m^j)=\Tr\!\bigl(M(\omega_m^j)^n\bigr),
\qquad
\rho_m:=\max_{1\le j\le m-1}\rho_{m,j}.
\]
则有
\[
\boxed{\;
\limsup_{n\to\infty}|a_n^{(j)}|^{1/n}=\rho_{m,j}.\;}
\]
进一步，若定义 centered periodic residue deviation
\[
\Delta_m(n):=
\max_{r\in\ZZ/m\ZZ}
\left|A_{m,r}(n)-\frac{a_n(1)}{m}\right|,
\]
则其最坏指数满足精确恒等式
\[
\boxed{\;
\limsup_{n\to\infty}\Delta_m(n)^{1/n}=\rho_m.\;}
\]
\end{theorem}
\begin{proof}
对固定 \(j\)，设 \(M(\omega_m^j)\) 的特征值为 \(\lambda_{j,1},\dots,\lambda_{j,d}\)。则
\[
a_n^{(j)}=\sum_{\ell=1}^d \lambda_{j,\ell}^n.
\]
上界
\[
\limsup_{n\to\infty}|a_n^{(j)}|^{1/n}\le \rho(M(\omega_m^j))=\rho_{m,j}
\]
显然成立。对反向下界，取所有满足 \(|\lambda_{j,\ell}|=\rho_{m,j}\) 的主模长特征值，写成
\[
\lambda_{j,\ell}=\rho_{m,j}e^{i\theta_\ell},
\qquad
\ell=1,\dots,s.
\]
则
\[
\rho_{m,j}^{-n}a_n^{(j)}=\sum_{\ell=1}^{s}e^{in\theta_\ell}+o(1).
\]
由 Kronecker 逼近定理，存在整数序列 \(n_k\to\infty\) 使
\[
e^{in_k\theta_\ell}\to 1
\qquad(1\le \ell\le s).
\]
于是
\[
\rho_{m,j}^{-n_k}a_{n_k}^{(j)}\to s\neq 0,
\]
故 \(\limsup |a_n^{(j)}|^{1/n}\ge \rho_{m,j}\)。合并上下界即得第一式。

再证 \(\Delta_m(n)\) 的精确指数。由定理 \ref{thm:xi-dirichlet-residue-fourier-parseval-exact} 的 Fourier 反演，
\[
a_n(\omega_m^j)
=
\sum_{r\in\ZZ/m\ZZ}
\left(A_{m,r}(n)-\frac{a_n(1)}{m}\right)\omega_m^{jr}.
\]
因此对每个 \(1\le j\le m-1\)，有
\[
|a_n(\omega_m^j)|\le m\,\Delta_m(n),
\]
从而
\[
\limsup_{n\to\infty}\Delta_m(n)^{1/n}\ge \rho_{m,j}.
\]
对 \(j\) 取最大值得
\[
\limsup_{n\to\infty}\Delta_m(n)^{1/n}\ge \rho_m.
\]

另一方面，再由定理 \ref{thm:xi-dirichlet-residue-fourier-parseval-exact}，
\[
\left|A_{m,r}(n)-\frac{a_n(1)}{m}\right|
\le
\frac1m\sum_{j=1}^{m-1}|a_n(\omega_m^j)|.
\]
结合命题 \ref{prop:real-input-40-output-dirichlet-nonrh} 的上界
\[
|a_n(\omega_m^j)|\le C_{m,j}\rho_{m,j}^n,
\]
便得
\[
\Delta_m(n)\le \frac1m\sum_{j=1}^{m-1}C_{m,j}\rho_{m,j}^n.
\]
对有限和取 \(\limsup^{1/n}\) 即有
\[
\limsup_{n\to\infty}\Delta_m(n)^{1/n}\le \rho_m.
\]
上下界合并即得第二式。证毕。
\end{proof}

\begin{corollary}[平方根 RH 型误差在某些模数上以严格指数方式失败]\label{cor:xi-dirichlet-residue-rh-failure-strict-exponential}
沿用命题 \ref{prop:real-input-40-output-dirichlet-nonrh} 的记号。若某个模数 \(m\ge 2\) 满足
\[
\theta_m:=\max_{1\le j\le m-1}\frac{\log\rho_{m,j}}{\log\lambda}>\frac12,
\qquad
\lambda=\rho(M(1))=\varphi^2,
\]
则
\[
\limsup_{n\to\infty}\frac{\log \Delta_m(n)}{n\log\lambda}
=
\theta_m
>
\frac12.
\]
等价地，存在常数 \(\eta>0\) 与无穷多个整数 \(n\)，使
\[
\Delta_m(n)\ge \lambda^{(\frac12+\eta)n}.
\]
因此，对应模数 \(m\) 的 periodic residue counts 不可能满足任何统一的 RH 型平方根误差上界
\[
\Delta_m(n)\ll \lambda^{n/2}.
\]
\end{corollary}
\begin{proof}
由定理 \ref{thm:xi-dirichlet-residue-bias-exponent-exact}，
\[
\limsup_{n\to\infty}\Delta_m(n)^{1/n}
=
\rho_m
=
\lambda^{\theta_m}.
\]
取对数并除以 \(\log\lambda\) 即得
\[
\limsup_{n\to\infty}\frac{\log \Delta_m(n)}{n\log\lambda}
=
\theta_m.
\]
若 \(\theta_m>1/2\)，取 \(\eta:=(\theta_m-1/2)/2>0\)。由 \(\limsup\) 的定义，存在无穷多个 \(n\) 使
\[
\Delta_m(n)\ge \lambda^{(\frac12+\eta)n}.
\]
这与任何统一的 \(\Delta_m(n)\ll \lambda^{n/2}\) 上界矛盾。证毕。
\end{proof}

\begin{theorem}[模 \texorpdfstring{$m$}{m} 局部极限定理的平方根障碍]\label{thm:xi-dirichlet-local-limit-sqrt-barrier}
沿用定理 \ref{thm:xi-dirichlet-residue-fourier-parseval-exact} 与命题 \ref{prop:real-input-40-output-dirichlet-nonrh} 的记号。若某个模数 \(m\ge 2\) 满足
\[
\theta_m:=\max_{1\le j\le m-1}\frac{\log\rho_{m,j}}{\log\lambda}>\frac12,
\qquad
\lambda=\rho(M(1))=\varphi^2,
\]
则不存在常数 \(C<\infty\) 使得
\[
\max_{r\in\ZZ/m\ZZ}\left|A_{m,r}(n)-\frac{a_n(1)}{m}\right|
\le C\,\lambda^{n/2}
\qquad (n\ge 1)
\]
恒成立。同样地，也不存在常数 \(C_{\mathrm{prim}}<\infty\) 使
\[
\max_{r\in\ZZ/m\ZZ}\left|P_{m,r}(n)-\frac{p_n(1)}{m}\right|
\le C_{\mathrm{prim}}\,\lambda^{n/2}
\qquad (n\ge 1)
\]
对所有 \(n\ge 1\) 同时成立。换言之，一旦 \(\theta_m>\frac12\)，则该模数上的 periodic 与 primitive residue counts 都不可能满足 RH 型平方根局部极限定理。
\end{theorem}
\begin{proof}
若存在常数 \(C<\infty\) 使 periodic 层上界成立，则对每个 \(1\le j\le m-1\)，由定理 \ref{thm:xi-dirichlet-residue-fourier-parseval-exact} 的逆 Fourier 变换有
\[
a_n(\omega_m^j)
=
\sum_{r\in\ZZ/m\ZZ}
\left(A_{m,r}(n)-\frac{a_n(1)}{m}\right)\omega_m^{jr},
\]
从而
\[
|a_n(\omega_m^j)|
\le
\sum_{r\in\ZZ/m\ZZ}\left|A_{m,r}(n)-\frac{a_n(1)}{m}\right|
\le
mC\,\lambda^{n/2}.
\]
对 \(n\to\infty\) 取 \(\limsup^{1/n}\)，再用定理 \ref{thm:xi-dirichlet-residue-bias-exponent-exact} 的精确指数公式
\[
\limsup_{n\to\infty}|a_n(\omega_m^j)|^{1/n}=\rho_{m,j},
\]
便得 \(\rho_{m,j}\le \sqrt{\lambda}\) 对一切 \(j\neq 0\) 成立，从而 \(\theta_m\le \frac12\)，与假设矛盾。故 periodic 层的统一平方根上界不存在。

primitive 层完全同理。若存在 \(C_{\mathrm{prim}}<\infty\) 使
\[
\max_{r\in\ZZ/m\ZZ}\left|P_{m,r}(n)-\frac{p_n(1)}{m}\right|
\le C_{\mathrm{prim}}\,\lambda^{n/2},
\]
则由同一定理的 primitive Fourier 反演公式
\[
p_n(\omega_m^j)
=
\sum_{r\in\ZZ/m\ZZ}
\left(P_{m,r}(n)-\frac{p_n(1)}{m}\right)\omega_m^{jr}
\]
得到
\[
|p_n(\omega_m^j)|\le mC_{\mathrm{prim}}\,\lambda^{n/2}
\qquad (1\le j\le m-1).
\]
而定理 \ref{thm:xi-dirichlet-residue-sqrtlaw-twist-equivalence} 已表明：对固定模数 \(m\)，primitive 全剩余类平方根律与全部非平凡扭曲谱半径界 \(\rho_{m,j}\le \sqrt{\lambda}\) 严格等价。于是仍推出 \(\theta_m\le \frac12\)，与假设矛盾。证毕。
\end{proof}

\begin{theorem}[模 \texorpdfstring{$m$}{m} 全剩余类平方根律与全部扭曲谱半径界的等价性]\label{thm:xi-dirichlet-residue-sqrtlaw-twist-equivalence}
固定整数 \(m\ge 2\)，沿用定理 \ref{thm:xi-dirichlet-residue-fourier-parseval-exact} 与
\ref{thm:xi-dirichlet-residue-bias-exponent-exact} 的记号。则下列三条等价：
\begin{enumerate}
  \item 对所有 \(1\le j\le m-1\)，都有
  \[
  \rho_{m,j}\le \sqrt{\lambda}.
  \]
  \item 对每个剩余类 \(r\in\ZZ/m\ZZ\)，都有 periodic 平方根律
  \[
  A_{m,r}(n)=\frac{a_n(1)}{m}+O\!\bigl(\lambda^{n/2+o(n)}\bigr).
  \]
  \item 对每个剩余类 \(r\in\ZZ/m\ZZ\)，都有 primitive 平方根律
  \[
  P_{m,r}(n)=\frac{p_n(1)}{m}+O\!\bigl(\lambda^{n/2+o(n)}\bigr).
  \]
\end{enumerate}
换言之，固定模数上的全剩余类 RH 型平方根律与全部非平凡 twisted spectral radii 落在 \(\sqrt{\lambda}\) 之下严格等价。
\end{theorem}
\begin{proof}
先证 (1)\(\Rightarrow\)(2)。由定理 \ref{thm:xi-dirichlet-residue-fourier-parseval-exact} 的 Fourier 反演，
\[
A_{m,r}(n)-\frac{a_n(1)}{m}
=
\frac1m\sum_{j=1}^{m-1}a_n(\omega_m^j)\omega_m^{-jr}.
\]
再由命题 \ref{prop:real-input-40-output-dirichlet-nonrh}，
\[
|a_n(\omega_m^j)|\le C_{m,j}\rho_{m,j}^n\le C_{m,j}\lambda^{n/2}.
\]
对有限个 \(j\) 求和即得
\[
A_{m,r}(n)-\frac{a_n(1)}{m}=O(\lambda^{n/2}),
\]
从而 (2) 成立。

再证 (2)\(\Rightarrow\)(1)。由逆 Fourier 变换，
\[
a_n(\omega_m^j)
=
\sum_{r\in\ZZ/m\ZZ}\omega_m^{jr}
\left(A_{m,r}(n)-\frac{a_n(1)}{m}\right).
\]
若 (2) 成立，则右端为 \(O(\lambda^{n/2+o(n)})\)。结合定理 \ref{thm:xi-dirichlet-residue-bias-exponent-exact} 的精确指数公式
\[
\limsup_{n\to\infty}|a_n(\omega_m^j)|^{1/n}=\rho_{m,j},
\]
便得 \(\rho_{m,j}\le \sqrt{\lambda}\) 对所有 \(j\neq 0\) 成立，即 (1)。

接着证 (1)\(\Rightarrow\)(3)。仍由定理 \ref{thm:xi-dirichlet-residue-fourier-parseval-exact}，
\[
P_{m,r}(n)-\frac{p_n(1)}{m}
=
\frac1m\sum_{j=1}^{m-1}p_n(\omega_m^j)\omega_m^{-jr}.
\]
而 Witt/M\"obius 反演给出
\[
p_n(\omega_m^j)
=
\frac1n\sum_{d\mid n}\mu(d)\,a_{n/d}(\omega_m^{jd}).
\]
在假设 (1) 下，对每个 \(d\mid n\) 与每个 \(j\neq 0\)，都有
\[
a_{n/d}(\omega_m^{jd})=
\begin{cases}
O(\lambda^{n/2}),& d=1,\\[2pt]
O(\lambda^{n/d})=O(\lambda^{n/2}),& d\ge 2,
\end{cases}
\]
故
\[
p_n(\omega_m^j)=O\!\bigl(\tau(n)\lambda^{n/2}\bigr)=O\!\bigl(\lambda^{n/2+o(n)}\bigr),
\]
其中用到约数函数满足 \(\tau(n)=e^{o(n)}\)。代回 Fourier 反演即得 (3)。

最后证 (3)\(\Rightarrow\)(1)。若 (1) 不成立，则存在某个 \(j\neq 0\) 使 \(\rho_{m,j}>\sqrt{\lambda}\)，从而
\[
\theta_m:=\max_{1\le \ell\le m-1}\frac{\log\rho_{m,\ell}}{\log\lambda}>\frac12.
\]
于是由定理 \ref{thm:xi-primitive-congruence-residual-nonrh}，存在某个剩余类 \(r_\ast\) 使
\[
\limsup_{n\to\infty}
\frac{1}{n\log\lambda}
\log\Bigl|P_{m,r_\ast}(n)-\frac{p_n(1)}{m}\Bigr|
\ge \theta_m>\frac12,
\]
这与 (3) 矛盾。故 (3) 亦推出 (1)。证毕。
\end{proof}

\begin{theorem}[连续射影压力是离散 moment-kernel 数据的规范凸上包络]\label{thm:xi-projective-pressure-canonical-convex-majorant}
定义连续射影压力
\[
\Lambda(q):=\log \rho(\mathcal L_q)
\qquad(q\ge 1),
\]
并在整数点 \(q\in\mathbb Z_{\ge 2}\) 定义离散压力数据
\[
\lambda_q:=\log r_q.
\]
则：
\begin{enumerate}
  \item \(\Lambda\) 在 \([1,\infty)\) 上凸；
  \item 对每个整数 \(q\ge 2\)，有
  \[
  \Lambda(q)\ge \lambda_q;
  \]
  \item 若 \(q_0<q_1<\cdots<q_s\) 为任意整数、\(\vartheta_i\ge 0\)、\(\sum_{i=0}^s\vartheta_i=1\)，并记
  \[
  \bar q:=\sum_{i=0}^s \vartheta_i q_i,
  \]
  则有 Jensen 型不等式
  \[
  \Lambda(\bar q)\le \sum_{i=0}^s \vartheta_i \Lambda(q_i).
  \]
  若 \(\bar q\in\mathbb Z_{\ge 2}\)，则进一步有链式估计
  \[
  \lambda_{\bar q}\le \Lambda(\bar q)\le \sum_{i=0}^s \vartheta_i \Lambda(q_i).
  \]
\end{enumerate}
因此，\(\Lambda\) 为离散数据 \(\{\lambda_q\}_{q\ge 2}\) 提供了一条不依赖外加拟合选择的规范凸上包络；而在定理 \ref{thm:pom-real-q-pressure-analytic-interpolation} 的谱隙制度下，上述整数点不等式进一步升级为精确插值 \(\Lambda(q)=\lambda_q\)。
\end{theorem}
\begin{proof}
前两条即定理 \ref{thm:xi-projective-pressure-path-holder-convexity} 的直接内容。第三条只是凸函数的 Jensen 不等式；若 \(\bar q\) 恰为整数，再与 \(\lambda_{\bar q}\le \Lambda(\bar q)\) 合并，即得所示链式估计。最后一句则由定理 \ref{thm:pom-real-q-pressure-analytic-interpolation} 直接给出。证毕。
\end{proof}

\begin{corollary}[解析插值制度下整数 moment 谱的离散对数凸性]\label{cor:xi-projective-pressure-integer-logconvex-under-interpolation}
\leanverified{paper\_xi\_projective\_pressure\_integer\_logconvex\_under\_interpolation}
在定理 \ref{thm:pom-real-q-pressure-analytic-interpolation} 的设定下，记
\[
I:=J\cap\ZZ_{\ge 2}.
\]
则序列
\[
q\longmapsto \log r_q
\]
在 \(I\) 上离散凸。特别地，若 \(q-1,q,q+1\in I\)，则
\[
\boxed{\;
r_q^2\le r_{q-1}\,r_{q+1}.\;}
\]
并且离散斜率
\[
q\longmapsto \log r_{q+1}-\log r_q
\]
在其定义域上单调不减。
\end{corollary}
\begin{proof}
由定理 \ref{thm:pom-real-q-pressure-analytic-interpolation}，对每个 \(q\in I\) 都有
\[
\log r_q=\tau(q).
\]
而 \(\tau\) 在 \(J\) 上为凸函数，故其限制到整数子集 \(I\) 后仍给出一个离散凸序列。对任意满足 \(q-1,q,q+1\in I\) 的整数 \(q\)，由中点凸性
\[
2\tau(q)\le \tau(q-1)+\tau(q+1)
\]
得到
\[
2\log r_q\le \log r_{q-1}+\log r_{q+1},
\]
指数化即得盒装不等式。离散凸序列的相邻差分单调不减，遂得第二条。证毕。
\end{proof}

\paragraph{凸压力与 Fourier 指数的闭合接口}
上述结果表明：附录 E 提供的并不只是“把整数 \(q\) 推广为实参数”的形式构造，而是一条无条件凸压力曲线；更进一步，输出分桶矩阵族已经完备决定整数矩谱塔，而连续归一化压力在整点格上的全部取值也被该离散谱塔唯一锚定。附录 D.527 提供的也不只是 cyclotomic 扭曲的上界证书，而是经由精确 Fourier--Parseval 反演所闭合的真实 residue bias 指数。于是，连续 \(q\)-压力与模 \(m\) Dirichlet 扭曲不再是两条并列技术线，而构成了同一解析控制理论的两侧边界：一侧以输出桶化刚性与凸性支配离散 moment 数据，另一侧以 twisted spectral radius 精确锁定 residue 偏差的指数尺度，并把平方根门槛的失效直接转写为模分布局部极限定理的不可实现性。
\begin{theorem}[整数阶射影压力给出全部输出 cylinder \texorpdfstring{$q$}{q}-矩的精确有限维实现]\label{thm:xi-projective-cylinder-qmoment-finite-realization}
沿用定义 \ref{def:pom-projective-transfer-operator}、定义 \ref{def:pom-homogeneous-polynomials-simplex}、命题 \ref{prop:pom-projective-operator-invariant-homogeneous} 与定理 \ref{thm:pom-projective-operator-degenerates-to-moment-kernel} 的记号。对任意输出词
\[
\mathbf b=b_1\cdots b_n\in A_{\mathrm{out}}^n
\]
记
\[
B_{\mathbf b}:=B_{b_n}\cdots B_{b_1}.
\]
则成立：
\begin{enumerate}
  \item 对任意整数 \(q\ge 1\)、任意 \(p\in\mathcal H_q\)、任意 \(n\ge 1\) 与任意 \(w\in\Delta\)，都有
  \[
  \boxed{\;
  (\mathcal L_q^n p)(w)=\sum_{|\mathbf b|=n}p(B_{\mathbf b}w).\;}
  \]
  \item 若定义
  \[
  p_q(w):=\langle \mathbf 1,w\rangle^q\in\mathcal H_q,
  \]
  则 \(p_q|_{\Delta}\equiv 1\)，从而
  \[
  \boxed{\;
  (\mathcal L_q^n\mathbf 1)(w)
  =
  \sum_{|\mathbf b|=n}\langle \mathbf 1,B_{\mathbf b}w\rangle^q.\;}
  \]
  \item 对每个整数 \(q\ge 2\)，序列
  \[
  M_n^{(q)}(w):=(\mathcal L_q^n\mathbf 1)(w)
  \]
  是有限矩阵幂 \(\widetilde A_q^n\) 的一个矩阵系数；因而对每个固定 \(w\in\Delta\)，它满足阶数至多为
  \[
  \dim\mathcal H_q=\binom{q+|Q|-1}{|Q|-1}
  \]
  的线性递推，而其所处的 Perron 增长尺度由
  \[
  r_q=\rho(\widetilde A_q)
  \]
  精确控制。
\end{enumerate}
换言之，整数截面上的实参数射影压力并不只在谱半径层面退化到 moment-kernel；全部输出 cylinder 的 \(q\)-阶计数矩本身，就已被同一有限维线性系统严格承载。
\end{theorem}
\begin{proof}
先证第一条。\(n=1\) 时，这正是命题 \ref{prop:pom-projective-operator-invariant-homogeneous} 给出的公式
\[
(\mathcal L_q p)(w)=\sum_{b\in A_{\mathrm{out}}}p(B_bw).
\]
若结论对某个 \(n\ge 1\) 已成立，则
\[
(\mathcal L_q^{n+1}p)(w)
=
\sum_{b\in A_{\mathrm{out}}}g_b(w)^q\,(\mathcal L_q^n p)(\Phi_b(w)).
\]
代入归纳假设得到
\[
(\mathcal L_q^{n+1}p)(w)
=
\sum_{b\in A_{\mathrm{out}}}
\sum_{|\mathbf c|=n}
g_b(w)^q\,p(B_{\mathbf c}\Phi_b(w)).
\]
由于 \(p\) 为 \(q\) 次齐次多项式且
\[
\Phi_b(w)=\frac{B_bw}{g_b(w)},
\]
于是
\[
g_b(w)^q\,p(B_{\mathbf c}\Phi_b(w))
=
p(B_{\mathbf c}B_bw).
\]
对所有 \(b\) 与 \(\mathbf c\) 求和，便得
\[
(\mathcal L_q^{n+1}p)(w)=\sum_{|\mathbf d|=n+1}p(B_{\mathbf d}w),
\]
归纳完成。

第二条取
\[
p_q(w)=\langle \mathbf 1,w\rangle^q.
\]
由于 \(w\in\Delta\) 时 \(\langle \mathbf 1,w\rangle=1\)，故 \(p_q|_{\Delta}\equiv 1\)。再由第一条，
\[
(\mathcal L_q^n\mathbf 1)(w)
=
(\mathcal L_q^n p_q)(w)
=
\sum_{|\mathbf b|=n}p_q(B_{\mathbf b}w)
=
\sum_{|\mathbf b|=n}\langle \mathbf 1,B_{\mathbf b}w\rangle^q.
\]

第三条中，对 \(q\ge 2\)，定理 \ref{thm:pom-projective-operator-degenerates-to-moment-kernel} 表明 \(\mathcal L_q|_{\mathcal H_q}\) 在任一单项式基下的矩阵表示恰为 \(\widetilde A_q\)。因此 \((\mathcal L_q^n p_q)(w)\) 是 \(\widetilde A_q^n\) 的一个矩阵系数。Cayley--Hamilton 定理遂给出阶数至多为 \(\dim\mathcal H_q\) 的线性递推，而该有限维系统的 Perron 主尺度正是
\[
r_q=\rho(\widetilde A_q).
\]
证毕。
\end{proof}

\begin{theorem}[输出分桶矩阵对整数矩谱塔的完备决定性]\label{thm:xi-projective-output-bucket-determines-integer-moment-tower}
沿用定义 \ref{def:pom-projective-transfer-operator}、
命题 \ref{prop:pom-projective-operator-invariant-homogeneous}
与定理 \ref{thm:pom-projective-operator-degenerates-to-moment-kernel} 的记号。设
\(\mathcal T,\mathcal T'\)
为两套确定性在线核，具有同一状态集 \(Q\) 与同一输出字母表
\(A_{\mathrm{out}}\)。若对每个输出符号 \(b\in A_{\mathrm{out}}\) 都有
\[
B_b(\mathcal T)=B_b(\mathcal T')\in \Mat_{Q\times Q}(\NN),
\]
并分别记两套在线核在整数阶 \(q\) 上的对称幂商 moment-kernel 与 Perron 根为
\[
\widetilde A_q(\mathcal T),\ \widetilde A_q(\mathcal T'),
\qquad
r_q(\mathcal T),\ r_q(\mathcal T'),
\]
则对每个整数 \(q\ge 2\) 都成立
\[
\widetilde A_q(\mathcal T)=\widetilde A_q(\mathcal T'),
\]
从而
\[
r_q(\mathcal T)=r_q(\mathcal T'),
\qquad
\log r_q(\mathcal T)=\log r_q(\mathcal T').
\]
换言之，整数矩层上的对称幂商 moment-kernel 塔及其 Perron 率只依赖输出分桶矩阵族
\(\{B_b\}_{b\in A_{\mathrm{out}}}\)，而与输入字母表的内部标号及同输出边的进一步细分无关。
\end{theorem}
\begin{proof}
由定义 \ref{def:pom-projective-transfer-operator}，对每个 \(b\in A_{\mathrm{out}}\)，
\[
g_b(w)=\langle \mathbf 1,B_b w\rangle,
\qquad
\Phi_b(w)=\frac{B_b w}{\langle \mathbf 1,B_b w\rangle}
\quad\bigl(g_b(w)>0\bigr),
\]
故 \(g_b\) 与 \(\Phi_b\) 仅由矩阵 \(B_b\) 决定。于是当
\(B_b(\mathcal T)=B_b(\mathcal T')\)
对所有 \(b\) 成立时，两套在线核诱导出的射影算子
\(\mathcal L_q\) 与 \(\mathcal L_q'\)
在每个整数 \(q\ge 2\) 的齐次子空间 \(\mathcal H_q\) 上满足同一作用公式
\[
(\mathcal L_q p)(w)=\sum_{b\in A_{\mathrm{out}}}p(B_b w)
=
(\mathcal L_q' p)(w)
\qquad (p\in \mathcal H_q).
\]
因此
\[
\mathcal L_q|_{\mathcal H_q}=\mathcal L_q'|_{\mathcal H_q}.
\]
再由定理 \ref{thm:pom-projective-operator-degenerates-to-moment-kernel}，上述公共限制在单项式基下的矩阵表示恰为
\(\widetilde A_q\)，故
\[
\widetilde A_q(\mathcal T)=\widetilde A_q(\mathcal T').
\]
谱半径与其对数相等式随之立即得到。证毕。
\end{proof}

\begin{theorem}[离散矩谱的整点嵌入唯一锚定连续射影压力]\label{thm:xi-projective-moment-tower-anchors-normalized-pressure}
在定理 \ref{thm:pom-real-q-pressure-analytic-interpolation} 的设定下，记
\[
\Lambda(t):=\log \lambda(t+1)-\log |A_{\mathrm{out}}|
\qquad (t+1\in J),
\]
其中 \(\lambda(\cdot)\) 为射影算子族 \(\mathcal L_q\) 的 Perron 主特征值。则对每个整数
\(q\ge 2\) 且 \(q\in J\)，都有
\[
\Lambda(q-1)=\log r_q-\log |A_{\mathrm{out}}|.
\]
因此离散矩谱
\[
\{\log r_q-\log |A_{\mathrm{out}}|\}_{q\in J\cap \ZZ_{\ge 2}}
\]
恰为连续归一化射影压力在整点格 \(\ZZ_{\ge 1}\) 上的取值。特别地，一旦射影算子族
\(\mathcal L_{t+1}\)
及其谱隙制度固定，连续压力曲线在整数层便不再拥有独立自由度，而被 moment-kernel 塔完全锁定。
\end{theorem}
\begin{proof}
由定理 \ref{thm:pom-real-q-pressure-analytic-interpolation}，对每个整数
\(q\ge 2\) 且 \(q\in J\)，有
\[
\lambda(q)=r_q.
\]
于是把 \(t=q-1\) 代入
\(
\Lambda(t)=\log \lambda(t+1)-\log |A_{\mathrm{out}}|
\)
便得到
\[
\Lambda(q-1)=\log \lambda(q)-\log |A_{\mathrm{out}}|
=
\log r_q-\log |A_{\mathrm{out}}|.
\]
这正说明离散矩谱给出了连续归一化射影压力的全部整数锚点。最后一句只是把这一恒等式解释为：在同一算子族与同一热力学制度下，整数层取值已被完全预定。证毕。
\end{proof}
为便于把前述凸压力链与长度同余类的 \(\zeta\)-接口对齐，下面把射影压力、periodic residue 生成函数以及 residue zeta 的 cyclotomic--Witt 分解整理为一组可直接调用的标准形式。

\begin{theorem}[射影压力算子谱半径的单纯形凸性证书]\label{thm:xi-time-part50dc-projective-pressure-perron-logconvex}
\leanverified{paper\_xi\_time\_part50dc\_projective\_pressure\_perron\_logconvex}
沿用定义 \ref{def:pom-projective-transfer-operator} 的记号。在单纯形
\[
\Delta:=\{w\in\mathbb R_{\ge 0}^{Q}:\ \langle \mathbf 1,w\rangle=1\}
\]
上，对每个输出符号 \(b\in A_{\mathrm{out}}\) 记
\[
g_b(w):=\langle \mathbf 1,B_bw\rangle,
\qquad
\Phi_b(w):=\frac{B_bw}{\langle \mathbf 1,B_bw\rangle}
\quad \bigl(g_b(w)>0\bigr),
\]
并对任意实参数 \(q\ge 1\) 定义射影转移算子
\[
(\mathcal L_qf)(w):=
\sum_{b\in A_{\mathrm{out}}}g_b(w)^q\,f(\Phi_b(w)),
\]
其中 \(g_b(w)=0\) 的项按 \(0\) 计。再记
\[
\lambda(q):=\rho_{\mathrm{proj}}(q):=\rho(\mathcal L_q).
\]
则函数
\[
q\longmapsto \log \lambda(q)
\]
在 \([1,\infty)\) 上凸。更精确地，对每个 \(n\ge 1\)，函数
\[
q\longmapsto \log \|\mathcal L_q^n\|_{\infty\to\infty}
=
\log \|\mathcal L_q^n\mathbf 1\|_\infty
\]
在 \([1,\infty)\) 上凸；从而对归一化连续压力
\[
\Lambda(t):=\log \lambda(t+1)-\log |A_{\mathrm{out}}|
\qquad (t\ge 0),
\]
也有
\[
t\longmapsto \Lambda(t)
\]
在 \([0,\infty)\) 上凸。等价地，对任意 \(q_0,q_1\ge 1\) 与 \(t\in[0,1]\)，都有
\[
\rho_{\mathrm{proj}}\bigl((1-t)q_0+tq_1\bigr)
\le
\rho_{\mathrm{proj}}(q_0)^{1-t}\rho_{\mathrm{proj}}(q_1)^t.
\]
\end{theorem}
\begin{proof}
先证算子范数可由正锥元 \(\mathbf 1\) 实现。对任意 \(f\in C(\Delta)\) 与任意 \(w\in\Delta\)，有
\[
|(\mathcal L_qf)(w)|
\le
\sum_{b\in A_{\mathrm{out}}}g_b(w)^q\,\|f\|_\infty
=
(\mathcal L_q\mathbf 1)(w)\,\|f\|_\infty.
\]
取上确界可得
\[
\|\mathcal L_qf\|_\infty\le \|\mathcal L_q\mathbf 1\|_\infty\,\|f\|_\infty,
\]
故
\[
\|\mathcal L_q\|_{\infty\to\infty}\le \|\mathcal L_q\mathbf 1\|_\infty.
\]
反向不等式由取 \(f=\mathbf 1\) 立即得到，因此
\[
\|\mathcal L_q\|_{\infty\to\infty}=\|\mathcal L_q\mathbf 1\|_\infty.
\]
对 \(\mathcal L_q^n\) 重复同一论证，便有
\[
\|\mathcal L_q^n\|_{\infty\to\infty}=\|\mathcal L_q^n\mathbf 1\|_\infty
\qquad (n\ge 1).
\]

现固定 \(n\ge 1\)。对任意输出路径
\[
\beta=(b_1,\dots,b_n)\in A_{\mathrm{out}}^n
\]
与任意 \(w\in\Delta\)，递归定义
\[
w_0:=w,
\qquad
w_i:=\Phi_{b_i}(w_{i-1})
\quad (1\le i\le n),
\]
并记路径权
\[
G_\beta(w):=\prod_{i=1}^{n}g_{b_i}(w_{i-1})\ge 0.
\]
反复展开 \(\mathcal L_q\) 得
\[
(\mathcal L_q^n\mathbf 1)(w)=\sum_{\beta\in A_{\mathrm{out}}^n}G_\beta(w)^q.
\]
若对某个 \(w\) 有全部 \(G_\beta(w)=0\)，则上式恒为 \(0\)，该点对上确界无贡献；否则在指标集
\[
I(w):=\{\beta\in A_{\mathrm{out}}^n:\ G_\beta(w)>0\}
\]
上有
\[
\log (\mathcal L_q^n\mathbf 1)(w)
=
\log\!\left(\sum_{\beta\in I(w)}e^{q\log G_\beta(w)}\right),
\]
这是有限个仿射函数 \(q\mapsto q\log G_\beta(w)\) 的 \(\log\)-sum-exp，因而关于 \(q\) 凸。对 \(w\) 取上确界仍保持凸性，于是
\[
q\longmapsto
\log\|\mathcal L_q^n\mathbf 1\|_\infty
=
\log\|\mathcal L_q^n\|_{\infty\to\infty}
\]
在 \([1,\infty)\) 上凸。

最后，由 Gelfand 公式，
\[
\log \lambda(q)
=
\log \rho(\mathcal L_q)
=
\lim_{n\to\infty}\frac{1}{n}\log\|\mathcal L_q^n\|_{\infty\to\infty}.
\]
右端是一列凸函数的点态极限，因此仍为凸函数。这就证明了
\[
q\longmapsto \log \lambda(q)
\]
在 \([1,\infty)\) 上凸。对
\[
\Lambda(t)=\log \lambda(t+1)-\log |A_{\mathrm{out}}|
\]
仅作平移与常数平移，故 \(t\mapsto \Lambda(t)\) 在 \([0,\infty)\) 上同样凸。最后将对数凸性指数化，即得所示谱半径不等式。证毕。
\end{proof}

\begin{corollary}[整数阶 moment-kernel 谱半径嵌入一条规范的凸包络曲线]\label{cor:xi-time-part50dc-moment-kernel-convex-envelope}
\leanverified{paper\_xi\_time\_part50dc\_moment\_kernel\_convex\_envelope}
对每个整数 \(q\ge 2\)，记
\[
r_q:=\rho(\widetilde A_q).
\]
则有
\[
\log r_q\le \log \rho_{\mathrm{proj}}(q).
\]
因此，离散压力样本
\[
q\longmapsto \log r_q
\]
内嵌于连续凸曲线
\[
q\longmapsto \log \rho_{\mathrm{proj}}(q)
\]
之下。更精确地，若 \(q_0<q_1<\cdots<q_s\) 为整数、\(\vartheta_i\ge 0\)、\(\sum_i\vartheta_i=1\)，并记
\[
\bar q:=\sum_{i=0}^s\vartheta_i q_i\in\mathbb Z_{\ge 2},
\]
则有链式估计
\[
\log r_{\bar q}
\le
\log \rho_{\mathrm{proj}}(\bar q)
\le
\sum_{i=0}^s \vartheta_i\log \rho_{\mathrm{proj}}(q_i).
\]
\end{corollary}
\begin{proof}
定理 \ref{thm:xi-projective-pressure-canonical-convex-majorant} 的第二、三条正是上述结论在
\[
\Lambda(q)=\log\rho_{\mathrm{proj}}(q)
\]
记号下的改写。证毕。
\end{proof}

\begin{theorem}[长度同余类 periodic residue 生成函数的有理性与 core 八次控制]\label{thm:xi-time-part50dc-periodic-residue-rational-core-octic}
固定整数 \(m\ge 2\)，记 \(\omega_m:=e^{2\pi i/m}\)，并把定理 \ref{thm:xi-dirichlet-residue-fourier-parseval-exact} 的 periodic residue counts 改写为
\[
A_n^{(r;m)}:=A_{m,r}(n)
\qquad (r\in\mathbb Z/m\mathbb Z,\ n\ge 1).
\]
定义普通生成函数
\[
G_{r,m}(z):=\sum_{n\ge 1}A_n^{(r;m)}z^n.
\]
则有 Fourier--resolvent 表示
\[
G_{r,m}(z)
=
\frac1m\sum_{j=0}^{m-1}\omega_m^{-jr}\,
\Tr\!\left(\frac{zM(\omega_m^j)}{I-zM(\omega_m^j)}\right).
\]
进一步，沿用附录 \ref{app:real-input-40-zeta-u} 的因式分解
\[
\Delta(z,u)=\det(I-zM(u))=(1-uz^2)F(z,u),
\]
并定义 core 生成函数
\[
G_{r,m}^{\mathrm{core}}(z)
:=
-\frac1m\sum_{j=0}^{m-1}\omega_m^{-jr}\,
z\partial_z\log F(z,\omega_m^j).
\]
则有精确分解
\[
G_{r,m}(z)
=
2\sum_{\substack{k\ge 1\\ k\equiv r\ (\mathrm{mod}\ m)}}z^{2k}
+G_{r,m}^{\mathrm{core}}(z).
\]
因此：
\begin{enumerate}
  \item \(G_{r,m}(z)\) 是有理函数，其分母整除
  \[
  (1-z^{2m})\prod_{j=0}^{m-1}F(z,\omega_m^j);
  \]
  \item \(G_{r,m}^{\mathrm{core}}(z)\) 的分母整除
  \[
  \prod_{j=0}^{m-1}F(z,\omega_m^j);
  \]
  \item 系数序列 \(\bigl(A_n^{(r;m)}\bigr)_{n\ge 1}\) 满足长度至多 \(10m\) 的线性递推，而扣除显式原子项之后的序列满足长度至多 \(8m\) 的线性递推。
\end{enumerate}
\end{theorem}
\begin{proof}
由定理 \ref{thm:xi-dirichlet-residue-fourier-parseval-exact}，
\[
A_n^{(r;m)}
=
\frac1m\sum_{j=0}^{m-1}\omega_m^{-jr}a_n(\omega_m^j).
\]
对 \(n\ge 1\) 乘以 \(z^n\) 并求和，得到
\[
G_{r,m}(z)
=
\frac1m\sum_{j=0}^{m-1}\omega_m^{-jr}
\sum_{n\ge 1}a_n(\omega_m^j)z^n.
\]
另一方面，由矩阵几何级数，
\[
\sum_{n\ge 1}a_n(u)z^n
=
\sum_{n\ge 1}\Tr\!\bigl((zM(u))^n\bigr)
=
\Tr\!\left(\frac{zM(u)}{I-zM(u)}\right),
\]
从而得到第一条显示公式。

再利用
\[
\Tr\!\left(\frac{zM(u)}{I-zM(u)}\right)
=
-z\partial_z\log\det(I-zM(u))
=
-z\partial_z\log\Delta(z,u),
\]
并代入
\[
\Delta(z,u)=(1-uz^2)F(z,u),
\]
便有
\[
\Tr\!\left(\frac{zM(u)}{I-zM(u)}\right)
=
\frac{2uz^2}{1-uz^2}
-z\partial_z\log F(z,u).
\]
将 \(u=\omega_m^j\) 代回 Fourier 平均，可得
\[
\frac1m\sum_{j=0}^{m-1}\omega_m^{-jr}\frac{2\omega_m^j z^2}{1-\omega_m^j z^2}
=
\frac1m\sum_{j=0}^{m-1}\omega_m^{-jr}
\sum_{k\ge 1}2\omega_m^{jk}z^{2k}
=
2\sum_{k\ge 1}z^{2k}\cdot
\frac1m\sum_{j=0}^{m-1}\omega_m^{j(k-r)}.
\]
单位根正交性给出
\[
\frac1m\sum_{j=0}^{m-1}\omega_m^{j(k-r)}
=
\mathbf 1_{\{k\equiv r\ (\mathrm{mod}\ m)\}},
\]
故
\[
G_{r,m}(z)
=
2\sum_{\substack{k\ge 1\\ k\equiv r\ (\mathrm{mod}\ m)}}z^{2k}
+G_{r,m}^{\mathrm{core}}(z).
\]

原子项显然是分母整除 \(1-z^{2m}\) 的有理函数，而 \(G_{r,m}^{\mathrm{core}}(z)\) 是有限个
\[
\frac{\text{多项式}}{F(z,\omega_m^j)}
\]
之和，因此其分母整除 \(\prod_{j=0}^{m-1}F(z,\omega_m^j)\)。由注
\ref{rem:real-input-40-zero-activation}，对 \(j=0\) 时 \(F(z,1)\) 的 \(z\)-次数为 \(7\)，而对 \(j\neq 0\) 时次数至多为 \(8\)；因此
\[
\deg_z\prod_{j=0}^{m-1}F(z,\omega_m^j)\le 8m,
\]
而加上原子分母 \(1-z^{2m}\) 后，完整分母次数至多为 \(10m\)。由有理生成函数与常系数线性递推的标准对应，便得到最后一条。证毕。
\end{proof}

\begin{corollary}[single atom 在 periodic residue 生成函数层形成显式 cyclotomic 壳]\label{cor:xi-time-part50dc-atomic-periodic-residue-cyclotomic-shell}
\leanverified{paper\_xi\_time\_part50dc\_atomic\_periodic\_residue\_cyclotomic\_shell}
沿用定理 \ref{thm:xi-time-part50dc-periodic-residue-rational-core-octic} 的记号，并定义原子 generating function
\[
G_{r,m}^{\mathrm{cycle}}(z):=G_{r,m}(z)-G_{r,m}^{\mathrm{core}}(z).
\]
令 \(\ell_r\in\{2,4,\dots,2m\}\) 为满足
\[
\ell_r\equiv 2r\pmod{2m}
\]
的唯一正代表。则有显式闭式
\[
\boxed{\;
G_{r,m}^{\mathrm{cycle}}(z)=\frac{2z^{\ell_r}}{1-z^{2m}}.\;}
\]
特别地：
\begin{enumerate}
  \item \(G_{r,m}^{\mathrm{cycle}}(z)\) 的全部极点都是 simple cyclotomic poles，且恰位于
  \[
  \mu_{2m}:=\{\zeta\in\CC:\ \zeta^{2m}=1\};
  \]
  \item 原子 generating function \(G_{r,m}^{\mathrm{cycle}}(z)\) 的系数序列 \(\bigl([z^n]G_{r,m}^{\mathrm{cycle}}(z)\bigr)_{n\ge 1}\) 的最小周期恰为 \(2m\)。
\end{enumerate}
\end{corollary}
\begin{proof}
由定理 \ref{thm:xi-time-part50dc-periodic-residue-rational-core-octic}，
\[
G_{r,m}^{\mathrm{cycle}}(z)
=
2\sum_{\substack{k\ge 1\\ k\equiv r\ (\mathrm{mod}\ m)}}z^{2k}.
\]
当 \(r\neq 0\) 时，满足 \(k\equiv r\pmod m\) 的正整数可唯一写成 \(k=r+qm\)（\(q\ge 0\)），于是
\[
G_{r,m}^{\mathrm{cycle}}(z)
=
2\sum_{q\ge 0}z^{2r+2qm}
=
\frac{2z^{2r}}{1-z^{2m}}.
\]
当 \(r=0\) 时，最小正解为 \(k=m\)，故
\[
G_{0,m}^{\mathrm{cycle}}(z)
=
2\sum_{q\ge 0}z^{2m+2qm}
=
\frac{2z^{2m}}{1-z^{2m}}.
\]
两种情形统一写成
\[
G_{r,m}^{\mathrm{cycle}}(z)=\frac{2z^{\ell_r}}{1-z^{2m}}.
\]
因此其全部极点恰为 \(1-z^{2m}\) 的零点，并且均为单极点。另一方面，系数序列恰是在模 \(2m\) 的唯一同余类 \(n\equiv \ell_r\pmod{2m}\) 上取值 \(2\)、其余位置取值 \(0\)，故其最小周期就是 \(2m\)。证毕。
\end{proof}

\begin{corollary}[full periodic residue 生成函数的奇点来源分裂为 core 谱与 cyclotomic 壳]\label{cor:xi-time-part50dc-full-periodic-residue-core-cyclotomic-splitting}
\leanverified{paper\_xi\_time\_part50dc\_full\_periodic\_residue\_core\_cyclotomic\_splitting}
沿用上述记号。对每个固定的剩余类 \(r\in\ZZ/m\ZZ\)，full periodic residue generating function 满足精确分解
\[
\boxed{\;
G_{r,m}(z)
=
G_{r,m}^{\mathrm{core}}(z)+\frac{2z^{\ell_r}}{1-z^{2m}}.\;}
\]
因此，single atom 所能引入的新奇点只可能落在 \(\mu_{2m}\) 上，并且都为 simple cyclotomic poles；全部模长不等于 \(1\) 的奇点全部来自 \(G_{r,m}^{\mathrm{core}}(z)\)。特别地，一切仅由奇点模长决定的指数型量，例如 periodic residue counts 的 \(\limsup^{1/n}\) 与相应的平方根门槛失效，均完全由 core 谱控制。
\end{corollary}
\begin{proof}
第一式只是定理 \ref{thm:xi-time-part50dc-periodic-residue-rational-core-octic} 与推论 \ref{cor:xi-time-part50dc-atomic-periodic-residue-cyclotomic-shell} 的直接合并。又
\[
\frac{2z^{\ell_r}}{1-z^{2m}}
\]
的全部极点都位于单位圆上的 \(\mu_{2m}\)，故 atom 部分不可能贡献任何模长不等于 \(1\) 的奇点。于是全部此类奇点都必须来自 \(G_{r,m}^{\mathrm{core}}(z)\)；若在个别单位根处发生可去抵消，也只会删除 \(\mu_{2m}\) 中的点，而不会在其外产生新的奇点。证毕。
\end{proof}

\begin{theorem}[长度同余类 residue zeta 的 cyclotomic--Witt 分解与原子梳齿骨架]\label{thm:xi-time-part50dc-residue-zeta-cyclotomic-witt}
固定整数 \(m\ge 2\) 与剩余类 \(r\in\mathbb Z/m\mathbb Z\)。定义 residue zeta 函数（作为 \(z=0\) 附近常数项为 \(1\) 的形式幂级数）
\[
\mathcal Z_{r,m}(z)
:=
\exp\!\Bigl(\sum_{n\ge 1}\frac{A_n^{(r;m)}}{n}z^n\Bigr).
\]
则有精确分解
\[
\mathcal Z_{r,m}(z)
=
\prod_{j=0}^{m-1}
\Delta(z,\omega_m^j)^{-\omega_m^{-jr}/m},
\]
其中幂指数按 \(z=0\) 附近常数项为 \(1\) 的分支解释。进一步，
\[
\mathcal Z_{r,m}(z)=\mathcal A_{r,m}(z)\,\mathcal R_{r,m}(z),
\]
其中
\[
\mathcal A_{r,m}(z)
:=
\prod_{j=0}^{m-1}(1-\omega_m^j z^2)^{-\omega_m^{-jr}/m},
\qquad
\mathcal R_{r,m}(z)
:=
\prod_{j=0}^{m-1}F(z,\omega_m^j)^{-\omega_m^{-jr}/m}.
\]
并且原子因子满足显式梳齿展开
\[
\log \mathcal A_{r,m}(z)
=
\sum_{\substack{n\ge 1\\ n\equiv r\ (\mathrm{mod}\ m)}}
\frac{z^{2n}}{n}.
\]
特别地，
\[
\mathcal A_{0,m}(z)=(1-z^{2m})^{-1/m}.
\]
此外，对全部剩余类求积有
\[
\prod_{r\in\mathbb Z/m\mathbb Z}\mathcal Z_{r,m}(z)
=
\Delta(z,1)^{-1}
=
\zeta(z,1).
\]
换言之，\(\mathcal A_{r,m}\) 完全记录二步原子在长度同余类上的梳齿骨架，而 \(\mathcal R_{r,m}\) 只承载由 core 因子 \(F(z,\omega_m^j)\) 所决定的剩余谱数据。
\end{theorem}
\begin{proof}
由定理 \ref{thm:xi-dirichlet-residue-fourier-parseval-exact} 的 Fourier 反演，
\[
\sum_{n\ge 1}\frac{A_n^{(r;m)}}{n}z^n
=
\frac1m\sum_{j=0}^{m-1}\omega_m^{-jr}
\sum_{n\ge 1}\frac{a_n(\omega_m^j)}{n}z^n.
\]
而由
\[
\zeta(z,u)=\Delta(z,u)^{-1}
=
\exp\!\Bigl(\sum_{n\ge 1}\frac{a_n(u)}{n}z^n\Bigr),
\]
可得
\[
\sum_{n\ge 1}\frac{a_n(u)}{n}z^n=-\log \Delta(z,u).
\]
于是
\[
\sum_{n\ge 1}\frac{A_n^{(r;m)}}{n}z^n
=
-\frac1m\sum_{j=0}^{m-1}\omega_m^{-jr}\log\Delta(z,\omega_m^j).
\]
指数化便得到第一条 cyclotomic 分解。

再把
\[
\Delta(z,u)=(1-uz^2)F(z,u)
\]
代入，即得
\[
\mathcal Z_{r,m}(z)=\mathcal A_{r,m}(z)\,\mathcal R_{r,m}(z).
\]
对原子因子取对数，有
\[
\log\mathcal A_{r,m}(z)
=
-\frac1m\sum_{j=0}^{m-1}\omega_m^{-jr}\log(1-\omega_m^j z^2).
\]
展开
\[
-\log(1-\omega_m^j z^2)
=
\sum_{n\ge 1}\frac{\omega_m^{jn}}{n}z^{2n},
\]
得到
\[
\log\mathcal A_{r,m}(z)
=
\sum_{n\ge 1}\frac{z^{2n}}{n}\cdot
\frac1m\sum_{j=0}^{m-1}\omega_m^{j(n-r)}.
\]
内层平均等于 \(1\) 当且仅当 \(n\equiv r\pmod m\)，否则为 \(0\)，故
\[
\log \mathcal A_{r,m}(z)
=
\sum_{\substack{n\ge 1\\ n\equiv r\ (\mathrm{mod}\ m)}}
\frac{z^{2n}}{n}.
\]
当 \(r=0\) 时，上式化为
\[
\log\mathcal A_{0,m}(z)
=
\sum_{k\ge 1}\frac{z^{2mk}}{mk}
=
-\frac1m\log(1-z^{2m}),
\]
从而
\[
\mathcal A_{0,m}(z)=(1-z^{2m})^{-1/m}.
\]

最后，对全部 \(r\in\mathbb Z/m\mathbb Z\) 求积并利用单位根正交性
\[
\sum_{r\in\mathbb Z/m\mathbb Z}\omega_m^{-jr}
=
m\,\delta_{j0},
\]
便得到
\[
\prod_{r\in\mathbb Z/m\mathbb Z}\mathcal Z_{r,m}(z)
=
\Delta(z,1)^{-1}
=
\zeta(z,1).
\]
证毕。
\end{proof}

\begin{conjecture}[Dirichlet 挠点上的 \texorpdfstring{$\mathrm{RH}(u)$}{RH(u)} 失败应在单位圆上稠密]\label{conj:xi-dirichlet-torsion-rhu-failure-dense}
设
\[
\mathcal E_{\mathrm{tor}}
:=
\{u\in \mathrm{Tor}(\TT)\setminus\{1\}:\ \text{对应扭曲上的 RH 型平方根律失败}\}.
\]
则应有如下更强现象：除有限多个例外外，全部非平凡挠点都属于 \(\mathcal E_{\mathrm{tor}}\)。特别地，
\[
\overline{\mathcal E_{\mathrm{tor}}}=\TT.
\]
\end{conjecture}
\begin{remark}
该猜想与本文已有三条证据完全同向：其一，定理 \ref{thm:xi-output-potential-large-odd-primes-bad} 已给出充分大奇素数模数全部为坏模数；其二，推论 \ref{cor:xi-output-potential-bad-moduli-density-one} 已证明坏模数集合在 \(\NN\) 中具有自然密度 \(1\)；其三，结论 \ref{con:xi-time-part9l-dirichlet-twist-divisibility-monotonicity} 表明平方根门槛失效会沿整除链向上传播。因而，现有结果已经把扭曲平方根律的失败推进到一个自然密度为 \(1\) 的模数族上；本猜想只是把这种“模数典型性”进一步提升为“挠点典型性”，并断言其在单位圆上形成稠密失败集。
\end{remark}

\begin{corollary}[同一 40 状态核上的 Gaussian bulk 与 arithmetic obstruction 构成严格双速率分裂]\label{cor:xi-gaussian-but-nonrh-congruence-separation}
\leanverified{paper\_xi\_gaussian\_but\_nonrh\_congruence\_separation}
在真实输入 \(40\) 状态核上，不存在单一的平方根谱律同时统治实参数压力方向与全部 cyclotomic 扭曲方向。更精确地说：
\begin{enumerate}
  \item \textbf{实方向。} 在
  \[
  \alpha_0=P'(0)=\frac{23}{20}-\frac{7\sqrt5}{20}
  \]
  的邻域内，归一化速率函数 \(I(\alpha)\) 具有定理 \ref{thm:xi-real-input-cramer-germ-quadratic-field-rigidity} 的 Cram\'er germ；其二次主项给出 Gaussian bulk，而奇次修正由 \(\kappa_3<0\) 精确控制，故右尾严格比左尾更昂贵。
  \item \textbf{单位根方向。} 存在某个模数 \(m\ge 2\) 与某个剩余类 \(a_\ast\in\ZZ/m\ZZ\) 使
  \[
  \limsup_{n\to\infty}
  \frac{1}{n\log\lambda}
  \log\Bigl|A_n(a_\ast;m)-\frac{A_n}{m}\Bigr|
  >
  \frac12.
  \]
\end{enumerate}
因此，该核呈现出严格的实轴--单位圆双速率分裂：中心极限定理与 Edgeworth 喷射在实方向成立，而平方根型算术等分布在某些 Dirichlet 方向必然破裂。换言之，核版 RH 窗口只约束正实轴压力方向，并不向同模长的 cyclotomic 扭曲截面径向传播。
\end{corollary}
\begin{proof}
第 \(1\) 条由定理 \ref{thm:xi-real-input-cramer-germ-quadratic-field-rigidity} 直接给出。

对第 \(2\) 条，命题 \ref{prop:real-input-40-output-dirichlet-nonrh} 已证明存在某个模数 \(m\) 满足
\[
\theta_m=\max_{1\le j\le m-1}\frac{\log\rho_{m,j}}{\log\lambda}>\frac12.
\]
再由推论 \ref{cor:xi-dirichlet-residue-rh-failure-strict-exponential}，对应模数的 periodic residue bias 以精确指数 \(\theta_m\) 违反平方根门槛，从而严格超过 \(1/2\)。
于是实参数压力方向与 cyclotomic 扭曲方向不可能受同一平方根谱律统一控制。证毕。
\end{proof}

\paragraph{输出位势的端点能垒、局部 Legendre 曲率与激活支刚性}
在输出位势的冻结端点、局部累积量闭式、Dirichlet 扭曲门槛与激活分支的解析展开已经建立之后，还可进一步抽取出下列若干条彼此闭合的解析后果。

\begin{theorem}[饱和端点的大偏差缺口等于地基熵指数]\label{thm:xi-output-potential-endpoint-gap-ground-state}
沿用命题 \ref{prop:real-input-40-pressure-freezing}、推论 \ref{cor:real-input-40-ground-entropy} 与推论 \ref{cor:real-input-40-rate-endpoints} 的记号。记
\[
P(\theta)=\log\lambda(e^\theta),
\qquad
h_{\mathrm{ground}}:=\log c>0,
\]
并令 \(I(\alpha)\) 为输出位势的归一化大偏差速率函数。则端点差额满足严格恒等式
\[
\boxed{\;
I(0)-I\!\left(\frac12\right)=h_{\mathrm{ground}}.
\;}
\]
\end{theorem}
\begin{proof}
由推论 \ref{cor:real-input-40-rate-endpoints}，
\[
I(0)=\log\lambda(1)=\log(\varphi^2),
\qquad
I\!\left(\frac12\right)=\log\!\left(\frac{\varphi^2}{c}\right).
\]
两式相减得到
\[
I(0)-I\!\left(\frac12\right)
=
\log(\varphi^2)-\log\!\left(\frac{\varphi^2}{c}\right)
=
\log c.
\]
再由推论 \ref{cor:real-input-40-ground-entropy}，\(\log c=h_{\mathrm{ground}}\)。证毕。
\end{proof}
\leanverified{paper_xi_output_potential_endpoint_gap_ground_state}

\begin{theorem}[局部 Legendre 几何在 \texorpdfstring{$\mathbb Q(\sqrt5)$}{Q(sqrt5)} 上闭合]\label{thm:xi-output-potential-local-legendre-quadratic-field}
沿用命题 \ref{prop:real-input-40-output-cumulants-closed} 的记号。令
\[
\alpha_0:=P'(0),
\qquad
I(\alpha):=\sup_{\theta\in\RR}\Bigl\{\theta\alpha-\bigl(P(\theta)-P(0)\bigr)\Bigr\}.
\]
则 \(I\) 在 \(\alpha_0\) 的某邻域内解析，并且存在系数 \(c_n\in \mathbb Q(\sqrt5)\) 使
\[
\boxed{\;
I(\alpha)=\sum_{n\ge 2}c_n(\alpha-\alpha_0)^n
\;}
\]
在该邻域内成立。
\end{theorem}
\begin{proof}
由命题 \ref{prop:real-input-40-output-cumulants-closed}，\(P\) 在 \(\theta=0\) 处的全部 Taylor 系数属于 \(\mathbb Q(\sqrt5)\)，并且
\[
P''(0)>0
\]
（见表 \ref{tab:real_input_40_output_potential_cumulants_closed}），故 \(P'\) 在 \(\theta=0\) 邻域内由解析反函数定理可逆。于是存在唯一解析函数
\[
\Theta(\alpha)
\]
定义在 \(\alpha_0\) 的某邻域内，使得
\[
P'(\Theta(\alpha))=\alpha,
\qquad
\Theta(\alpha_0)=0.
\]
由于 \(P'(\theta)\) 的 Taylor 系数全在 \(\mathbb Q(\sqrt5)\) 中，且一次项 \(P''(0)\neq 0\)，由 Lagrange 反演公式可知 \(\Theta(\alpha)\) 的 Taylor 系数仍属于 \(\mathbb Q(\sqrt5)\)。

再由 Legendre 对偶的驻点表示，
\[
I(\alpha)=\alpha\,\Theta(\alpha)-\bigl(P(\Theta(\alpha))-P(0)\bigr).
\]
右端由 \(\mathbb Q(\sqrt5)\) 上的解析幂级数复合、乘法与加法构成，因此其 Taylor 系数仍全部属于 \(\mathbb Q(\sqrt5)\)。最后，由归一化定义，
\[
I(\alpha_0)=0,
\qquad
I'(\alpha_0)=\Theta(\alpha_0)=0,
\]
故展开从二次项开始。证毕。
\end{proof}

\begin{theorem}[输出位势大偏差率函数的局部非对称精确展开]\label{thm:xi-output-potential-local-legendre-skew-steep}
沿用定理 \ref{thm:xi-output-potential-local-legendre-quadratic-field} 的记号，并记
\[
\mu:=\alpha_0=P'(0)=\frac{23}{20}-\frac{7\sqrt5}{20},
\qquad
\sigma^2:=P''(0)=\frac{1791}{200}-\frac{3983}{1000}\sqrt5>0,
\]
\[
\kappa_3:=P^{(3)}(0)=\frac{1416369}{5000}-\frac{126691}{1000}\sqrt5<0,
\qquad
\kappa_4:=P^{(4)}(0)=\frac{1354428303}{100000}-\frac{605718497}{100000}\sqrt5>0.
\]
则当 \(\delta\to 0\) 时，局部速率函数满足严格展开
\[
\boxed{\;
I(\mu+\delta)
=
\frac{\delta^2}{2\sigma^2}
-\frac{\kappa_3}{6\sigma^6}\delta^3
+\frac{3\kappa_3^2-\sigma^2\kappa_4}{24\sigma^{10}}\delta^4
+O(\delta^5).\;}
\]
特别地，前三个非平凡系数都属于 \(\QQ(\sqrt5)\)，并且
\[
\boxed{\;
-\frac{\kappa_3}{6\sigma^6}\approx 22.8681939098>0,\qquad
\frac{3\kappa_3^2-\sigma^2\kappa_4}{24\sigma^{10}}\approx 72.7346495328>0.\;}
\]
数值上有
\begin{equation}
\label{eq:real-input-40-output-potential-local-rate-series}
I(\mu+\delta)=10.2582524\,\delta^2+22.8681939\,\delta^3+72.7346495\,\delta^4+O(\delta^5).
\end{equation}
等价地，
\[
I''(\mu)=\frac{1}{\sigma^2}>0,\qquad
I^{(3)}(\mu)= -\frac{\kappa_3}{\sigma^6}>0,\qquad
I^{(4)}(\mu)=\frac{3\kappa_3^2-\sigma^2\kappa_4}{\sigma^{10}}>0.
\]
因此，输出位势的局部大偏差既不是关于 \(\mu\) 的对称 Gaussian 扰动，也不是关于 \(\mu\) 的偶函数；其中心邻域内右侧偏离的指数代价严格高于左侧偏离，并伴随正的四阶局部陡峭性。
\end{theorem}
\begin{proof}
由关系
\[
P'(\Theta(\alpha))=\alpha
\]
令 \(\alpha=\mu+\delta\)。由
\[
\Theta(\mu)=0
\]
与 \(P''(0)=\sigma^2>0\)，函数 \(P'\) 在 \(0\) 的邻域内解析可逆。把
\[
P'(\theta)
=
\mu+\sigma^2\theta+\frac{\kappa_3}{2}\theta^2+\frac{\kappa_4}{6}\theta^3+O(\theta^4)
\]
按 \(\delta\) 反解，得到
\[
\Theta(\mu+\delta)
=
\frac{\delta}{\sigma^2}
-\frac{\kappa_3}{2\sigma^6}\delta^2
+\frac{3\kappa_3^2-\sigma^2\kappa_4}{6\sigma^{10}}\delta^3
+O(\delta^4).
\]
再由 Legendre 对偶的驻点表示，
\[
I(\mu+\delta)
=
(\mu+\delta)\Theta(\mu+\delta)-\bigl(P(\Theta(\mu+\delta))-P(0)\bigr).
\]
再把
\[
P(\theta)
=
P(0)+\mu\theta+\frac{\sigma^2}{2}\theta^2+\frac{\kappa_3}{6}\theta^3+\frac{\kappa_4}{24}\theta^4+O(\theta^5)
\]
代入并整理各阶 \(\delta\) 项，便得到所示四阶展开。

三次项系数的符号由 \(\kappa_3<0\) 直接给出。对四次项，代入闭式可化简为
\[
3\kappa_3^2-\sigma^2\kappa_4
=
\frac{5989334657578-2678511881499\sqrt5}{25000000}
>0,
\]
故四次项系数亦严格为正，从而
\[
I^{(4)}(\mu)=\frac{3\kappa_3^2-\sigma^2\kappa_4}{\sigma^{10}}>0.
\]
生成式
\eqref{eq:real-input-40-output-potential-local-rate-series}
中的数值系数与上文两条数值近似都由上述显式公式直接代入得到。

最后，若 \(I\) 在 \(\mu\) 邻域内关于 \(\mu\) 反射对称，则其所有奇阶导数在 \(\mu\) 处都必须为零；这与
\(
I^{(3)}(\mu)>0
\)
矛盾。故 \(I\) 不是偶函数。由于原始累积量满足 \(\kappa_3<0\)，这表明局部 Cram\'er germ 在均值点右侧具有更高的指数代价；同时 \(I^{(4)}(\mu)>0\) 给出正的四阶局部陡峭性。证毕。
\end{proof}

\begin{corollary}[输出位势局部偏差的严格方向不等式]\label{cor:xi-output-potential-local-legendre-one-sided-asymmetry}
沿用定理 \ref{thm:xi-output-potential-local-legendre-skew-steep} 的记号。对充分小的 \(\delta>0\)，有
\[
\boxed{\;
I(\mu+\delta)-I(\mu-\delta)
=
-\frac{\kappa_3}{3\sigma^6}\delta^3+O(\delta^5)
>0.\;}
\]
从而
\[
\boxed{\;
I(\mu+\delta)>\frac{\delta^2}{2\sigma^2}>I(\mu-\delta)
\qquad(0<\delta\ll1).\;}
\]
因此在真实输入 \(40\) 状态核的输出位势下，同等幅度的右偏差比高斯主项与左偏差都具有更高的指数代价；换言之，局部中尺度偏差在速率函数层面呈现严格的右侧代价增强。
\end{corollary}
\begin{proof}
由定理 \ref{thm:xi-output-potential-local-legendre-skew-steep} 的 Taylor 展开，偶次项在
\[
I(\mu+\delta)-I(\mu-\delta)
\]
中完全相消，主导项即为
\[
2\left(-\frac{\kappa_3}{6\sigma^6}\right)\delta^3
=
-\frac{\kappa_3}{3\sigma^6}\delta^3.
\]
由于 \(\kappa_3<0\)，该主项严格为正，故对充分小的 \(\delta>0\) 差值必为正。
\par
同理，由同一定理的四阶展开，
\[
I(\mu+\delta)-\frac{\delta^2}{2\sigma^2}
=
-\frac{\kappa_3}{6\sigma^6}\delta^3+O(\delta^4),
\]
\[
I(\mu-\delta)-\frac{\delta^2}{2\sigma^2}
=
\frac{\kappa_3}{6\sigma^6}\delta^3+O(\delta^4).
\]
由于
\(
-\kappa_3/(6\sigma^6)>0
\)，故当 \(\delta>0\) 充分小时，第一式右端严格为正而第二式右端严格为负，于是得到所示双边不等式。证毕。
\end{proof}

\paragraph{输出计数半质量中心的常数右移}

\begin{corollary}[输出计数的一阶 Edgeworth 连续中位点具有正的常数右移]\label{cor:xi-real-input-40-output-edgeworth-continuous-median-right-shift}
沿用定理 \ref{thm:xi-output-potential-local-legendre-skew-steep} 与
\ref{thm:xi-real-input-40-output-edgeworth-quasipowers} 的记号。令 \(S_n\) 为长度 \(n\) 上输出 \(1\) 的总数，并记
\[
\sigma:=\sqrt{\sigma^2}>0.
\]
再把 \(\mathfrak m_n^{\mathrm E}\) 定义为下列一阶 Edgeworth 半质量方程在 \(n\mu\) 邻域内的唯一实解：
\[
\Phi\!\left(\frac{\mathfrak m_n^{\mathrm E}-n\mu}{\sigma\sqrt n}\right)
+\frac{\kappa_3}{6\sigma^3\sqrt n}
\left(
1-\left(\frac{\mathfrak m_n^{\mathrm E}-n\mu}{\sigma\sqrt n}\right)^2
\right)
\phi\!\left(\frac{\mathfrak m_n^{\mathrm E}-n\mu}{\sigma\sqrt n}\right)
=
\frac12.
\]
则有
\[
\boxed{\;
\mathfrak m_n^{\mathrm E}
=
n\mu-\frac{\kappa_3}{6\sigma^2}+O(n^{-1/2})
=
n\mu+c_*+O(n^{-1/2}).\;}
\]
其中
\[
\boxed{\;
c_*:=-\frac{\kappa_3}{6\sigma^2}
=
\frac{-34164035+15595899\sqrt5}{13058700}
\approx 0.0543281746438>0.\;}
\]
因此，一阶 Edgeworth 连续中位点稳定位于均值右侧。
\end{corollary}
\begin{proof}
由定理 \ref{thm:xi-real-input-40-output-edgeworth-quasipowers}，标准化变量
\[
Y_n:=\frac{S_n-n\mu}{\sigma\sqrt n}
\]
在中心窗口内满足
\[
\mathbb P(Y_n\le x)
=
\Phi(x)+\frac{\kappa_3}{6\sigma^3\sqrt n}(1-x^2)\phi(x)+O(n^{-1}).
\]
令
\[
x_n:=\frac{\mathfrak m_n^{\mathrm E}-n\mu}{\sigma\sqrt n}.
\]
由于 \(\Phi(0)=1/2\) 且 \(\phi(0)>0\)，隐函数定理保证上面的半质量方程在 \(0\) 邻域内有唯一解，并且 \(x_n\to 0\)。把 \(x=x_n\) 代入，并用
\[
\Phi(x_n)-\frac12=\phi(0)x_n+O(x_n^3),
\qquad
(1-x_n^2)\phi(x_n)=\phi(0)+O(x_n^2),
\]
得到
\[
0
=
\phi(0)x_n+\frac{\kappa_3}{6\sigma^3\sqrt n}\phi(0)+O(n^{-1}),
\]
从而
\[
x_n=-\frac{\kappa_3}{6\sigma^3\sqrt n}+O(n^{-1}).
\]
乘回 \(\sigma\sqrt n\) 即得
\[
\mathfrak m_n^{\mathrm E}
=
n\mu-\frac{\kappa_3}{6\sigma^2}+O(n^{-1/2}).
\]
再把 \(\sigma^2\) 与 \(\kappa_3\) 的显式闭式代入并化简，得到
\[
-\frac{\kappa_3}{6\sigma^2}
=
\frac{-34164035+15595899\sqrt5}{13058700}.
\]
由于 \(\kappa_3<0\)，故 \(c_*>0\)。证毕。
\end{proof}

\paragraph{输出位势局部正规形与中心化 residue 的单原子截断律}
在输出位势的局部累积量展开、模 \(m\) 扭曲迹的 Fourier--Parseval 反演以及单原子 Euler 因子的 periodic/core 剥离都已建立之后，还可把这两条链进一步压缩为同一组中心化 residue 结论：局部大偏差在均值点具有四阶 Legendre--Birkhoff 正规形，residue 偏差的平方平均能量具有精确扭曲指数，而长度 \(2\) 原子因子在中心化 residue 空间中只产生一维坐标扰动，并把所有小于 \(1\) 的 core 指数精确截断到 \(1\)。

\begin{theorem}[输出位势速率函数的局部 Legendre--Birkhoff 正规形]\label{thm:xi-output-potential-local-legendre-birkhoff-normal-form}
沿用定理 \ref{thm:xi-output-potential-local-legendre-skew-steep} 的记号，并记
\[
I_{\mathrm{out}}(\alpha):=I(\alpha),
\qquad
\alpha_0:=\mu=P'(0).
\]
则当 \(\delta\to 0\) 时，
\[
\boxed{\;
I_{\mathrm{out}}(\alpha_0+\delta)
=
\frac{\delta^2}{2\sigma^2}
-\frac{\kappa_3}{6\sigma^6}\delta^3
+\frac{3\kappa_3^2-\sigma^2\kappa_4}{24\sigma^{10}}\delta^4
+O(\delta^5).\;}
\]
并且
\[
\boxed{\;
I_{\mathrm{out}}(\alpha_0+\delta)-I_{\mathrm{out}}(\alpha_0-\delta)
=
-\frac{\kappa_3}{3\sigma^6}\delta^3+O(\delta^5).\;}
\]
特别地，由 \(\kappa_3<0\) 得对充分小的 \(\delta>0\) 恒有
\[
I_{\mathrm{out}}(\alpha_0+\delta)>I_{\mathrm{out}}(\alpha_0-\delta),
\]
故均值点附近的局部大偏差呈现严格的右侧代价增强。
\end{theorem}
\begin{proof}
定理 \ref{thm:xi-output-potential-local-legendre-skew-steep} 已给出
\[
I(\mu+\delta)
=
\frac{\delta^2}{2\sigma^2}
-\frac{\kappa_3}{6\sigma^6}\delta^3
+\frac{3\kappa_3^2-\sigma^2\kappa_4}{24\sigma^{10}}\delta^4
+O(\delta^5).
\]
把其中 \(\mu\) 改写为 \(\alpha_0\)，便得到第一式。第二式正是推论 \ref{cor:xi-output-potential-local-legendre-one-sided-asymmetry} 的内容。最后由 \(\kappa_3<0\) 立得对充分小的 \(\delta>0\) 差值严格为正。证毕。
\end{proof}

\begin{theorem}[模 \texorpdfstring{$m$}{m} residue 偏差平方平均能量的精确扭曲指数]\label{thm:xi-output-residue-l2-energy-exact-exponent}
\leanverified{paper\_xi\_output\_residue\_l2\_energy\_exact\_exponent}
固定整数 \(m\ge 2\)，并记 \(\omega_m:=e^{2\pi i/m}\)。定义中心化 residue 计数与其平方平均能量
\[
\widetilde A_{m,r}(n):=A_{m,r}(n)-\frac{a_n(1)}{m},
\qquad
\mathcal V_m(n):=\sum_{r\in\ZZ/m\ZZ}\abs{\widetilde A_{m,r}(n)}^2,
\]
并记
\[
\rho_{m,j}:=\rho\!\bigl(M(\omega_m^j)\bigr),
\qquad
\rho_m:=\max_{1\le j\le m-1}\rho_{m,j}.
\]
则有
\[
\boxed{\;
\mathcal V_m(n)=\frac1m\sum_{j=1}^{m-1}\abs{a_n(\omega_m^j)}^2.\;}
\]
并且
\[
\boxed{\;
\limsup_{n\to\infty}\mathcal V_m(n)^{1/n}=\rho_m^2.\;}
\]
\end{theorem}
\begin{proof}
第一式正是定理 \ref{thm:xi-dirichlet-residue-fourier-parseval-exact} 在 periodic residue counts 上的 Parseval 恒等式。再由命题 \ref{prop:real-input-40-output-dirichlet-nonrh} 的扭曲迹上界，
\[
\abs{a_n(\omega_m^j)}\le C_{m,j}\rho_{m,j}^n
\qquad(1\le j\le m-1),
\]
得到
\[
\mathcal V_m(n)
\le
\frac1m\sum_{j=1}^{m-1}C_{m,j}^2\rho_{m,j}^{2n}
\le
C\,\rho_m^{2n},
\]
故
\(
\limsup_{n\to\infty}\mathcal V_m(n)^{1/n}\le \rho_m^2
\)。

另一方面，取 \(j_\ast\) 使 \(\rho_{m,j_\ast}=\rho_m\)。由定理 \ref{thm:xi-dirichlet-residue-bias-exponent-exact} 的 twisted trace 精确指数公式，
\[
\limsup_{n\to\infty}\abs{a_n(\omega_m^{j_\ast})}^{1/n}=\rho_m.
\]
再由 Parseval 公式，
\[
\mathcal V_m(n)\ge \frac1m\abs{a_n(\omega_m^{j_\ast})}^2,
\]
故
\(
\limsup_{n\to\infty}\mathcal V_m(n)^{1/n}\ge \rho_m^2
\)。
上下界合并即得结论。证毕。
\end{proof}

\begin{theorem}[单原子 Euler 因子在中心化 residue 超平面中的秩一坐标尖峰]\label{thm:xi-single-atom-centered-residue-rankone-spike}
\leanverified{paper\_xi\_single\_atom\_centered\_residue\_rankone\_spike}
沿用单原子剥离
\[
\Delta(z,u)=(1-uz^2)\,F(z,u)
\]
的记号。固定 \(m\ge 2\)，并记 \(\omega_m:=e^{2\pi i/m}\)。对 \(n\ge 1\) 与 \(r\in \ZZ/m\ZZ\)，定义
\[
\widetilde A_{m,r}^{\mathrm{full}}(n)
:=
A_{m,r}^{\mathrm{full}}(n)-\frac{a_n^{\mathrm{full}}(1)}{m},
\qquad
\widetilde A_{m,r}^{\mathrm{core}}(n)
:=
A_{m,r}^{\mathrm{core}}(n)-\frac{a_n^{\mathrm{core}}(1)}{m},
\]
以及中心化差
\[
\Delta_{m,r}(n)
:=
\widetilde A_{m,r}^{\mathrm{full}}(n)-\widetilde A_{m,r}^{\mathrm{core}}(n).
\]
则当 \(n\) 为奇数时，
\[
\Delta_{m,r}(n)=0
\qquad(\forall r\in\ZZ/m\ZZ),
\]
而当 \(n=2\ell\) 为偶数时，
\[
\boxed{\;
\Delta_{m,r}(2\ell)
=
2\,\mathbf 1_{\{r\equiv \ell\;(\mathrm{mod}\,m)\}}-\frac{2}{m}.\;}
\]
若把
\(
\Delta_m(2\ell):=(\Delta_{m,r}(2\ell))_{r\in\ZZ/m\ZZ}
\)
视为 \(\CC^{\ZZ/m\ZZ}\) 中的向量，并记 \(r_0\equiv \ell\pmod m\)，则
\[
\Delta_m(2\ell)=2\left(e_{r_0}-\frac{1}{m}\mathbf 1_m\right),
\qquad
\sum_{r\in\ZZ/m\ZZ}\Delta_{m,r}(2\ell)=0,
\]
其中 \(e_{r_0}\) 与 \(\mathbf 1_m\) 分别表示标准基向量与全一向量。进一步，其 Fourier 变换
\[
\widehat\Delta_{m,j}(2\ell)
:=
\sum_{r\in\ZZ/m\ZZ}\Delta_{m,r}(2\ell)\,\omega_m^{jr}
\]
满足
\[
\boxed{\;
\widehat\Delta_{m,0}(2\ell)=0,\qquad
\widehat\Delta_{m,j}(2\ell)=2\,\omega_m^{j\ell}\quad(1\le j\le m-1).\;}
\]
并且
\[
\boxed{\;
\|\Delta_m(2\ell)\|_{\ell^1}=4\Bigl(1-\frac1m\Bigr),\qquad
\|\Delta_m(2\ell)\|_{\ell^2}^2=4\Bigl(1-\frac1m\Bigr),\qquad
\|\Delta_m(2\ell)\|_{\ell^\infty}=2\Bigl(1-\frac1m\Bigr).\;}
\]
\end{theorem}
\begin{proof}
由
\[
\log\zeta(z,u)-\log\zeta_{\mathrm{core}}(z,u)
=
-\log(1-uz^2)
=
\sum_{k\ge 1}\frac{u^k}{k}z^{2k},
\]
可知单原子在 periodic/ghost 层上的附加项满足
\[
a_n^{\mathrm{cycle}}(u)
:=
a_n^{\mathrm{full}}(u)-a_n^{\mathrm{core}}(u)
=
2u^{n/2}\mathbf 1_{\{2\mid n\}}.
\]
因此当 \(n\) 为奇数时没有 residue 修正。对偶长度 \(n=2\ell\)，由 Fourier 反演，
\[
\begin{aligned}
A_{m,r}^{\mathrm{cycle}}(2\ell)
&:=
A_{m,r}^{\mathrm{full}}(2\ell)-A_{m,r}^{\mathrm{core}}(2\ell)\\
&=
\frac1m\sum_{j=0}^{m-1}a_{2\ell}^{\mathrm{cycle}}(\omega_m^j)\,\omega_m^{-jr}\\
&=
\frac{2}{m}\sum_{j=0}^{m-1}\omega_m^{j(\ell-r)}
=
2\,\mathbf 1_{\{r\equiv \ell\;(\mathrm{mod}\,m)\}}.
\end{aligned}
\]
另一方面，
\[
a_{2\ell}^{\mathrm{cycle}}(1)=2.
\]
故
\[
\Delta_{m,r}(2\ell)
=
A_{m,r}^{\mathrm{cycle}}(2\ell)-\frac{a_{2\ell}^{\mathrm{cycle}}(1)}{m}
=
2\,\mathbf 1_{\{r\equiv \ell\;(\mathrm{mod}\,m)\}}-\frac{2}{m},
\]
从而得到中心化尖峰的显式公式。

再对 Fourier 变换直接计算：
\[
\widehat\Delta_{m,j}(2\ell)
=
2\omega_m^{j\ell}-\frac{2}{m}\sum_{r\in\ZZ/m\ZZ}\omega_m^{jr}.
\]
当 \(j=0\) 时两项相消；当 \(1\le j\le m-1\) 时第二项由单位根正交性为零，于是得到
\(
\widehat\Delta_{m,j}(2\ell)=2\omega_m^{j\ell}
\)。

最后，\(\Delta_m(2\ell)\) 的一个坐标等于 \(2-2/m\)，其余 \(m-1\) 个坐标都等于 \(-2/m\)。因此
\[
\|\Delta_m(2\ell)\|_{\ell^1}
=
\left(2-\frac{2}{m}\right)+(m-1)\frac{2}{m}
=
4\left(1-\frac{1}{m}\right),
\]
\[
\|\Delta_m(2\ell)\|_{\ell^2}^2
=
\left(2-\frac{2}{m}\right)^2+(m-1)\left(\frac{2}{m}\right)^2
=
4\left(1-\frac{1}{m}\right),
\]
\[
\|\Delta_m(2\ell)\|_{\ell^\infty}=2-\frac{2}{m}.
\]
证毕。
\end{proof}

\begin{corollary}[单原子 Euler 因子对中心化 residue \texorpdfstring{$L^2$}{L2} 指数的精确截断律]\label{cor:xi-single-atom-centered-residue-l2-exponent-truncation}
固定 \(m\ge 2\)。定义 full/core 的中心化 residue 平方平均能量
\[
\mathcal V_m^{\mathrm{core}}(n):=\sum_{r\in\ZZ/m\ZZ}\abs{\widetilde A_{m,r}^{\mathrm{core}}(n)}^2,
\qquad
\mathcal V_m^{\mathrm{full}}(n):=\sum_{r\in\ZZ/m\ZZ}\abs{\widetilde A_{m,r}^{\mathrm{full}}(n)}^2,
\]
并对 \(1\le j\le m-1\) 定义
\[
\rho_{m,j}^{\mathrm{core}}
:=
\rho\!\bigl(M_{\mathrm{core}}(\omega_m^j)\bigr)
=
\frac{1}{\min\{|z|:\ F(z,\omega_m^j)=0\}},
\qquad
\rho_{m,j}^{\mathrm{full}}
:=
\rho\!\bigl(M(\omega_m^j)\bigr)
=
\frac{1}{\min\{|z|:\ \Delta(z,\omega_m^j)=0\}}.
\]
再记
\[
\rho_m^{\mathrm{core}}:=\max_{1\le j\le m-1}\rho_{m,j}^{\mathrm{core}},
\qquad
\rho_m^{\mathrm{full}}:=\max_{1\le j\le m-1}\rho_{m,j}^{\mathrm{full}}.
\]
则有
\[
\boxed{\;
\limsup_{n\to\infty}\bigl(\mathcal V_m^{\mathrm{core}}(n)\bigr)^{1/(2n)}
=
\rho_m^{\mathrm{core}}.\;}
\]
以及
\[
\boxed{\;
\limsup_{n\to\infty}\bigl(\mathcal V_m^{\mathrm{full}}(n)\bigr)^{1/(2n)}
=
\rho_m^{\mathrm{full}}
=
\max\{1,\rho_m^{\mathrm{core}}\}.\;}
\]
因此，当 \(\rho_m^{\mathrm{core}}>1\) 时，单原子因子不改变中心化 residue 的 \(L^2\) 指数；当 \(\rho_m^{\mathrm{core}}<1\) 时，它把 full 系统的该指数精确截断到 \(1\)。
\end{corollary}
\begin{proof}
把定理 \ref{thm:xi-output-residue-l2-energy-exact-exponent} 的证明逐字应用于 core 系统，即把 \(A_{m,r}\)、\(a_n\)、\(M\) 分别替换为 \(A_{m,r}^{\mathrm{core}}\)、\(a_n^{\mathrm{core}}\)、\(M_{\mathrm{core}}\)，便得到
\[
\limsup_{n\to\infty}\bigl(\mathcal V_m^{\mathrm{core}}(n)\bigr)^{1/n}
=
\bigl(\rho_m^{\mathrm{core}}\bigr)^2.
\]
取平方根即得第一式。

另一方面，由分解
\[
\Delta(z,\omega_m^j)=(1-\omega_m^j z^2)\,F(z,\omega_m^j)
\]
可知原子因子 \(1-\omega_m^j z^2\) 的两个零点都位于单位圆上，即都满足 \(|z|=1\)。因此
\[
\rho_{m,j}^{\mathrm{full}}
=
\max\{1,\rho_{m,j}^{\mathrm{core}}\}
\qquad(1\le j\le m-1).
\]
对 \(j\) 取最大值得
\[
\rho_m^{\mathrm{full}}=\max\{1,\rho_m^{\mathrm{core}}\}.
\]
再把定理 \ref{thm:xi-output-residue-l2-energy-exact-exponent} 应用于 full 系统，便有
\[
\limsup_{n\to\infty}\bigl(\mathcal V_m^{\mathrm{full}}(n)\bigr)^{1/n}
=
\bigl(\rho_m^{\mathrm{full}}\bigr)^2.
\]
再取平方根即可得到第二式。最后一句只是对
\(
\rho_m^{\mathrm{full}}=\max\{1,\rho_m^{\mathrm{core}}\}
\)
的直接解释。证毕。
\end{proof}

\begin{theorem}[零温端点的平方根临界律与普适比值 \texorpdfstring{$2$}{2}]\label{thm:xi-output-potential-zero-temp-square-root-critical-law}
沿用命题 \ref{prop:real-input-40-pressure-freezing} 的记号。设
\[
P(\theta)=\log\lambda(e^\theta)
=\frac{\theta}{2}+\log c+\frac{d}{c}e^{-\theta/2}
+\Bigl(\frac{e}{c}-\frac{d^2}{2c^2}\Bigr)e^{-\theta}
+O(e^{-3\theta/2})
\qquad (\theta\to+\infty),
\]
其中 \(c,d,e\) 为同命题中的 Puiseux 系数。则有
\[
\boxed{\;
\frac12-P'(\theta)=\frac{d}{2c}e^{-\theta/2}+O(e^{-\theta}),
\qquad
P''(\theta)=\frac{d}{4c}e^{-\theta/2}+O(e^{-\theta}).\;}
\]
因此
\[
\boxed{\;
\lim_{\theta\to+\infty}\frac{\frac12-P'(\theta)}{P''(\theta)}=2.\;}
\]
等价地，若记 \(u=e^\theta\)，则
\[
\frac12-P'(\theta)=\frac{d}{2c}u^{-1/2}+O(u^{-1}),
\qquad
P''(\theta)=\frac{d}{4c}u^{-1/2}+O(u^{-1}).
\]
\end{theorem}
\begin{proof}
对命题 \ref{prop:real-input-40-pressure-freezing} 中的 Puiseux 展开逐项求导，得到
\[
P'(\theta)=\frac12-\frac{d}{2c}e^{-\theta/2}
-\Bigl(\frac{e}{c}-\frac{d^2}{2c^2}\Bigr)e^{-\theta}
+O(e^{-3\theta/2}),
\]
从而
\[
\frac12-P'(\theta)=\frac{d}{2c}e^{-\theta/2}+O(e^{-\theta}).
\]
再求一次导数即得
\[
P''(\theta)=\frac{d}{4c}e^{-\theta/2}+O(e^{-\theta}).
\]
由于命题 \ref{prop:real-input-40-pressure-freezing} 已给出 \(c,d>0\)，上述两式主项同为 \(e^{-\theta/2}\) 量级，故
\[
\frac{\frac12-P'(\theta)}{P''(\theta)}
=
\frac{d/(2c)+O(e^{-\theta/2})}{d/(4c)+O(e^{-\theta/2})}
\longrightarrow
\frac{d/(2c)}{d/(4c)}
=2.
\]
将 \(e^{-\theta/2}=u^{-1/2}\) 代回，便得到 \(u\)-变量下的等价表达。证毕。
\end{proof}
\begin{theorem}[激活分支的三阶接触刚性与唯一小根性]\label{thm:xi-output-potential-activated-branch-rigidity}
沿用命题 \ref{prop:real-input-40-activated-branch-linear} 的记号，令
\[
G(\lambda,u)=\lambda^8F(\lambda^{-1},u).
\]
则在 \(u=1\) 的某邻域内，激活分支 \(\lambda_{\mathrm{act}}(u)\) 满足：
\begin{enumerate}
  \item \textbf{（三阶接触刚性）}
  \[
  \boxed{\;
  \lambda_{\mathrm{act}}(u)-(u-1)=3(u-1)^3+O\!\bigl((u-1)^4\bigr).
  \;}
  \]
  换言之，\(\lambda_{\mathrm{act}}\) 对线性分支 \(u\mapsto u-1\) 恰有二阶切触，而首个非线性偏离出现在三次项。
  \item \textbf{（唯一小根性）} 存在解析函数 \(\widetilde G(\lambda,u)\)，使
  \[
  G(\lambda,u)=\bigl(\lambda-\lambda_{\mathrm{act}}(u)\bigr)\widetilde G(\lambda,u),
  \qquad
  \widetilde G(0,1)\neq 0.
  \]
  因而在 \(u=1\) 的某个复邻域内，\(\lambda_{\mathrm{act}}(u)\) 是唯一趋于 \(0\) 的根分支，而其余七条根分支都与 \(0\) 保持正距离。
\end{enumerate}
\end{theorem}
\begin{proof}
由注 \ref{rem:real-input-40-activated-branch-series}，
\[
\lambda_{\mathrm{act}}(u)
:=(u-1)+3(u-1)^3-(u-1)^4+18(u-1)^5-22(u-1)^6
+O\!\bigl((u-1)^7\bigr).
\]
因此二次项系数严格为零，而三次项系数等于 \(3\neq 0\)，从而立刻得到第一条。

再由命题 \ref{prop:real-input-40-activated-branch-linear}，\(\lambda=0\) 是 \(u=1\) 时的简单根，且
\[
\lambda_{\mathrm{act}}(1)=0,
\qquad
\partial_\lambda G(0,1)=1\neq 0.
\]
由隐函数定理，存在唯一解析分支 \(\lambda_{\mathrm{act}}(u)\) 满足
\[
G(\lambda_{\mathrm{act}}(u),u)=0
\]
并经过 \((\lambda,u)=(0,1)\)。

对解析函数 \(G(\lambda,u)\) 在 \(\lambda=\lambda_{\mathrm{act}}(u)\) 处应用 Weierstrass 除法，可写成
\[
G(\lambda,u)=\bigl(\lambda-\lambda_{\mathrm{act}}(u)\bigr)\widetilde G(\lambda,u),
\]
其中 \(\widetilde G\) 在 \((0,1)\) 邻域解析。把 \(\lambda=\lambda_{\mathrm{act}}(u)\) 处求导代入上式可得
\[
\partial_\lambda G(\lambda_{\mathrm{act}}(u),u)=\widetilde G(\lambda_{\mathrm{act}}(u),u).
\]
因此在 \((0,1)\) 处，
\[
\widetilde G(0,1)=\partial_\lambda G(0,1)=1\neq 0.
\]
由连续性可知，存在 \(\delta,r>0\)，使得当 \(|u-1|<\delta\) 且 \(|\lambda|<r\) 时，
\[
\widetilde G(\lambda,u)\neq 0.
\]
于是 \(G(\lambda,u)=0\) 在该邻域内只能由
\(
\lambda=\lambda_{\mathrm{act}}(u)
\)
给出。这说明 \(\lambda_{\mathrm{act}}(u)\) 是唯一小根，而其余七条根分支全部远离 \(0\)。证毕。
\end{proof}

\paragraph{单位根扭曲的固定模态展开与四阶修正增益}

\begin{theorem}[固定单位根模态误差指数的四阶展开]\label{thm:xi-output-potential-dirichlet-fixed-j-fourth-order}
沿用命题 \ref{prop:real-input-40-output-cumulants-closed} 与命题 \ref{prop:real-input-40-output-dirichlet-nonrh} 的记号。令
\[
\omega_m:=e^{2\pi i/m},
\qquad
\rho_{m,j}:=\rho\!\bigl(M(\omega_m^j)\bigr),
\qquad
\theta_{m,j}:=\frac{\log \rho_{m,j}}{\log \lambda},
\qquad
\lambda=\lambda(1)=\varphi^2,
\]
并记
\[
\kappa_2:=P''(0)>0,
\qquad
\kappa_4:=P^{(4)}(0)>0.
\]
则对每个固定整数 \(j\ge 1\)，当 \(m\to\infty\) 且 \(m>j\) 时有
\[
\log\frac{\rho_{m,j}}{\lambda}
:=
-\frac{\kappa_2}{2}\Bigl(\frac{2\pi j}{m}\Bigr)^2
+\frac{\kappa_4}{24}\Bigl(\frac{2\pi j}{m}\Bigr)^4
+O_j(m^{-6}),
\]
从而
\[
\boxed{\;
\theta_{m,j}
:=
1-\frac{2\pi^2\kappa_2}{\log\lambda}\frac{j^2}{m^2}
+\frac{2\pi^4\kappa_4}{3\log\lambda}\frac{j^4}{m^4}
+O_j(m^{-6}).\;}
\]
特别地，
\[
\boxed{\;
\theta_{m,1}
:=
1-\frac{2\pi^2\kappa_2}{\log\lambda}\,m^{-2}
+\frac{2\pi^4\kappa_4}{3\log\lambda}\,m^{-4}
+O(m^{-6}).\;}
\]
\end{theorem}
\begin{proof}
由命题 \ref{prop:real-input-40-output-cumulants-closed}，压力
\[
P(\theta)=\log \lambda(e^\theta)
\]
在 \(\theta=0\) 邻域解析，且各阶导数均为实数。另一方面，命题 \ref{prop:real-input-40-output-dirichlet-nonrh} 给出的单位根扭曲满足
\[
\log\frac{\rho_{m,j}}{\lambda}
:=
\Re\!\left(P\!\left(\frac{2\pi i j}{m}\right)-P(0)\right).
\]
把 \(P(\theta)-P(0)\) 在 \(\theta=0\) 处展开为 Taylor 级数，并取
\(
\theta=2\pi i j/m
\)，得到
\[
P(\theta)-P(0)
:=
P'(0)\theta+\frac{\kappa_2}{2}\theta^2+\frac{P^{(3)}(0)}{6}\theta^3
+\frac{\kappa_4}{24}\theta^4+O_j(m^{-5}).
\]
由于奇次项在纯虚方向上贡献纯虚部，而首个未写出的偶次实部出现在六阶，故
\[
\Re\!\left(P\!\left(\frac{2\pi i j}{m}\right)-P(0)\right)
:=
-\frac{\kappa_2}{2}\Bigl(\frac{2\pi j}{m}\Bigr)^2
+\frac{\kappa_4}{24}\Bigl(\frac{2\pi j}{m}\Bigr)^4
+O_j(m^{-6}).
\]
这就给出第一式。再除以 \(\log \lambda\) 即得 \(\theta_{m,j}\) 的展开。\(j=1\) 情形只是其特例。证毕。
\end{proof}

\begin{corollary}[有界单位根窗口内的最坏模态最终由 \texorpdfstring{$j=1$}{j=1} 实现]\label{cor:xi-output-potential-dirichlet-bounded-window-j1}
表 \ref{tab:real_input_40_output_potential_dirichlet_twists} 在已审计模数上记录了
\[
j_{\mathrm{worst}}=1.
\]
更进一步，对任意固定整数 \(J\ge 2\)，存在 \(m_0(J)\) 使得对所有 \(m\ge m_0(J)\) 与一切 \(2\le j\le J\)，均有
\[
\theta_{m,1}>\theta_{m,j}.
\]
因此，在任何有界 Dirichlet 扭曲窗口内，\(j=1\) 最终严格支配全部其余固定模态。
\end{corollary}
\begin{proof}
令
\[
A:=\frac{2\pi^2\kappa_2}{\log\lambda}>0,
\qquad
B:=\frac{2\pi^4\kappa_4}{3\log\lambda}>0.
\]
由定理 \ref{thm:xi-output-potential-dirichlet-fixed-j-fourth-order}，对每个固定 \(J\ge 2\)，可把 \(2\le j\le J\) 上的误差统一写成 \(O_J(m^{-6})\)，从而
\[
\theta_{m,1}-\theta_{m,j}
:=
A\frac{j^2-1}{m^2}
-B\frac{j^4-1}{m^4}
+O_J(m^{-6}).
\]
右端主项为正且阶为 \(m^{-2}\)；其余项至多为 \(m^{-4}\)。因此对每个固定 \(j\in\{2,\dots,J\}\)，当 \(m\) 足够大时差值严格为正。对有限集合 \(\{2,\dots,J\}\) 取所需阈值的最大者，即得结论。证毕。
\end{proof}

\noindent 上述推论只使用了 \(j\) 有界的局部 Taylor 窗口。再结合单位圆上谱半径函数在 \(t=0\) 邻域的严格下降与远离 \(1\) 时的统一谱隙，还可把表 \ref{tab:real_input_40_output_potential_dirichlet_twists} 中的经验模式提升为充分大模数下的全局严格结论。
\begin{theorem}[单位圆扭曲谱半径的严格降阶与最坏模态的首模坍缩]\label{thm:xi-output-potential-dirichlet-global-worst-mode-j1}
\leanverified{paper\_xi\_output\_potential\_dirichlet\_global\_worst\_mode\_j1}
沿用命题 \ref{prop:real-input-40-output-dirichlet-nonrh}、推论 \ref{cor:real-input-40-output-dirichlet-cumulant-bridge} 与定理 \ref{thm:xi-atomic-witt-cycle-primitive-exact-peeling} 的记号。定义
\[
r(t):=\frac{\rho(M(e^{it}))}{\lambda},
\qquad
\lambda:=\rho(M(1))=\varphi^2,
\qquad
t\in[-\pi,\pi].
\]
则成立：
\begin{enumerate}
  \item 对每个 \(t\in(0,\pi]\)，都有
  \[
  r(t)<1.
  \]
  等价地，对每个 \(u\in \mathbb S^1\setminus\{1\}\)，都有
  \[
  \rho(M(u))<\lambda.
  \]
  \item 存在 \(m_0\) 使得对一切 \(m\ge m_0\) 都有
  \[
  \rho_m=\rho_{m,1}=\rho_{m,m-1},
  \qquad
  \theta_m=\theta_{m,1}=\theta_{m,m-1},
  \]
  并且对所有 \(2\le j\le m-2\) 严格满足
  \[
  \rho_{m,j}<\rho_{m,1}.
  \]
\end{enumerate}
换言之，全部非平凡单位根扭曲中的最慢衰减通道在充分大模数下必由基本频率 \(j=1\) 及其共轭 \(j=m-1\) 唯一实现。
\end{theorem}
\begin{proof}
由推论 \ref{cor:real-input-40-output-dirichlet-cumulant-bridge} 的证明，主特征值在 \(u=1\) 处为单根，故存在 \(\eta_0>0\) 使得对 \(|t|\le \eta_0\) 成立
\[
\log r(t)=\Re\bigl(P(it)-P(0)\bigr)
=-\frac{\kappa_2}{2}t^2+\frac{\kappa_4}{24}t^4+O(t^6).
\]
由于 \(\kappa_2>0\)，可缩小 \(\eta_0\) 为某个 \(\eta\in(0,\eta_0]\)，使 \(r(t)\) 在 \([0,\eta]\) 上随 \(t\) 严格下降。

另一方面，\(|M(e^{it})|=M(1)\) 且 \(M(1)\) 为非负不可约矩阵，故由 Wielandt 不等式（见 \cite{Wielandt1950UnzerlegbareNichtnegative}）有 \(r(t)\le 1\)。若某个 \(t\in(0,\pi]\) 取等，则 Wielandt 等号刻画迫使该单位模扭曲与 \(M(1)\) 只差一个全局相位与单位模对角共轭，也即出现 coboundary 可规约情形。命题 \ref{prop:real-input-40-pressure-freezing} 给出 \(\lambda(0)=1\)，故输出 \(0\) 子图含有闭路；而定理 \ref{thm:xi-atomic-witt-cycle-primitive-exact-peeling} 给出总输出 \(1\) 的 primitive 二步类。比较这两类闭路的总相位，可知上述可规约情形只能发生在 \(e^{it}=1\)。因此对全部 \(t\neq 0\) 都有
\[
r(t)<1.
\]
由连续性与紧性，存在 \(\delta_\eta>0\) 使得
\[
r(t)\le 1-\delta_\eta
\qquad (\eta\le |t|\le \pi).
\]

再由
\[
r(2\pi/m)
=\exp\!\left(
-\frac{\kappa_2}{2}\Bigl(\frac{2\pi}{m}\Bigr)^2
+O(m^{-4})
\right)
\longrightarrow 1,
\]
可取 \(m_0\) 使对一切 \(m\ge m_0\) 都有 \(2\pi/m<\eta\) 且
\[
r(2\pi/m)>1-\delta_\eta.
\]
于是凡满足 \(2\pi j/m\in[\eta,\pi]\) 的模态都不可能达到最坏值；而对满足 \(2\pi j/m<\eta\) 的其余模态，由 \(r(t)\) 在 \([0,\eta]\) 上严格下降知最大值只能在最小正角 \(2\pi/m\) 处取得，即 \(j=1\)。复共轭给出 \(\rho_{m,m-1}=\rho_{m,1}\)，其余 \(2\le j\le m-2\) 严格更小。证毕。
\end{proof}

\begin{theorem}[首模 Dirichlet 扭曲的二次主导衰减与指数趋一]\label{thm:xi-output-potential-dirichlet-index-to-one-all-moduli}
沿用定理 \ref{thm:xi-output-potential-dirichlet-global-worst-mode-j1} 与推论 \ref{cor:real-input-40-output-dirichlet-cumulant-bridge} 的记号。对每个 \(m\ge 2\)，定义
\[
\theta_{m,j}:=\frac{\log\rho(M(\omega_m^j))}{\log\lambda},
\qquad
\theta_m:=\max_{1\le j\le m-1}\theta_{m,j}.
\]
并记
\[
\rho_{m,1}:=\rho(M(\omega_m)).
\]
则存在 \(m_0\)，使得对一切 \(m\ge m_0\) 都有
\[
\theta_m=\theta_{m,1}=\theta_{m,m-1},
\]
并且
\[
\log\frac{\rho_{m,1}}{\lambda}
=
-\frac{\kappa_2}{2}\Bigl(\frac{2\pi}{m}\Bigr)^2
+\frac{\kappa_4}{24}\Bigl(\frac{2\pi}{m}\Bigr)^4
+O(m^{-6}),
\qquad
\kappa_2=P''(0)>0,
\]
并且
\[
\theta_m\longrightarrow 1
\qquad(m\to\infty).
\]
特别地，存在 \(m_1\)，使得对所有 \(m\ge m_1\) 都有
\[
\theta_m>\frac12.
\]
因此，对充分大的模数 \(m\)，periodic 与 primitive residue counts 都不可能满足统一的 RH 型平方根局部极限定理。
\end{theorem}
\begin{proof}
由定理 \ref{thm:xi-output-potential-dirichlet-global-worst-mode-j1}，存在 \(m_0\)，使得对一切 \(m\ge m_0\) 都有
\[
\theta_m=\theta_{m,1}=\theta_{m,m-1}.
\]
另一方面，由推论 \ref{cor:real-input-40-output-dirichlet-cumulant-bridge}，
\[
\log\frac{\rho_{m,1}}{\lambda}
=
-\frac{\kappa_2}{2}\Bigl(\frac{2\pi}{m}\Bigr)^2
+\frac{\kappa_4}{24}\Bigl(\frac{2\pi}{m}\Bigr)^4
+O(m^{-6})
\longrightarrow 0.
\]
故
\[
\theta_{m,1}
=
1+\frac{1}{\log\lambda}\log\frac{\rho_{m,1}}{\lambda}
\longrightarrow 1.
\]
与前述最终等式合并，即得
\[
\theta_m\longrightarrow 1.
\]
从而对充分大的 \(m\) 必有 \(\theta_m>\frac12\)。再结合推论
\ref{cor:xi-dirichlet-residue-rh-failure-strict-exponential} 与定理
\ref{thm:xi-dirichlet-local-limit-sqrt-barrier}，便知充分大模数上的 periodic 与 primitive residue counts 都不可能满足统一的 RH 型平方根局部极限定理。证毕。
\end{proof}

\begin{corollary}[\texorpdfstring{$\theta_m$}{theta_m} 的四阶显式展开把失效强度直接可计算化]\label{cor:xi-output-potential-dirichlet-index-expansion-all-moduli}
沿用定理 \ref{thm:xi-output-potential-dirichlet-index-to-one-all-moduli} 的记号。则对充分大的 \(m\) 都有
\[
\theta_m
=
1-\frac{\kappa_2}{2\log\lambda}\Bigl(\frac{2\pi}{m}\Bigr)^2
+\frac{\kappa_4}{24\log\lambda}\Bigl(\frac{2\pi}{m}\Bigr)^4
+O(m^{-6}).
\]
特别地，
\[
\theta_m=1-\Theta(m^{-2}),
\]
故最坏 Dirichlet 扭曲指数并非仅仅越过平方根门槛，而是以 \(m^{-2}\) 量级逼近 \(1\)。

若再代入命题 \ref{prop:real-input-40-output-cumulants-closed} 中的闭式常数
\[
\kappa_2=\frac{6\sqrt5-5}{125},
\qquad
\kappa_4=\frac{7+24\sqrt5}{3125},
\]
则进一步得到
\[
\theta_m
=
1-\frac{6\sqrt5-5}{250\log\lambda}\Bigl(\frac{2\pi}{m}\Bigr)^2
+\frac{7+24\sqrt5}{75000\log\lambda}\Bigl(\frac{2\pi}{m}\Bigr)^4
+O(m^{-6}).
\]
\end{corollary}
\begin{proof}
由定理 \ref{thm:xi-output-potential-dirichlet-index-to-one-all-moduli}，对充分大的 \(m\) 有
\[
\theta_m=\theta_{m,1}.
\]
再由推论 \ref{cor:real-input-40-output-dirichlet-cumulant-bridge}，
\[
\log\frac{\rho_{m,1}}{\lambda}
=
-\frac{\kappa_2}{2}\Bigl(\frac{2\pi}{m}\Bigr)^2
+\frac{\kappa_4}{24}\Bigl(\frac{2\pi}{m}\Bigr)^4
+O(m^{-6}),
\]
故
\[
\theta_m
=
\theta_{m,1}
=
1+\frac{1}{\log\lambda}\log\frac{\rho_{m,1}}{\lambda},
\]
代入即得第一条展开。由于 \(\kappa_2>0\)，主导修正项严格为负，故
\[
\theta_m=1-\Theta(m^{-2}).
\]
最后再把命题 \ref{prop:real-input-40-output-cumulants-closed} 给出的 \(\kappa_2,\kappa_4\) 闭式代入即可得到第二条公式。证毕。
\end{proof}

\begin{conjecture}[Dirichlet 扭曲谱的全局单峰性与首模恒最坏]\label{conj:xi-output-potential-dirichlet-all-m-j1}
表 \ref{tab:real_input_40_output_potential_dirichlet_twists} 中全部已审计模数都满足
\[
j_{\mathrm{worst}}=1,
\qquad
\rho_m=\rho_{m,1}=\rho_{m,m-1}.
\]
结合定理 \ref{thm:xi-output-potential-dirichlet-global-worst-mode-j1}，自然提出更强的全局单峰猜想：函数
\[
\theta\longmapsto \rho(M(e^{i\theta}))
\qquad (\theta\in[0,\pi])
\]
严格单调递减。因而对每个整数 \(m\ge 2\)，恒有
\[
\rho_m=\rho_{m,1}=\rho_{m,m-1},
\qquad
\theta_m=\theta_{m,1}.
\]
若该猜想成立，则全部非平凡 Dirichlet 误差指数都可由单一显式序列 \(\{\theta_{m,1}\}_{m\ge 2}\) 统一控制，而无需再对 \(1\le j\le m-1\) 逐一取最坏值。
\end{conjecture}
\begin{theorem}[充分大奇素数模数的平方根门槛必然失效]\label{thm:xi-output-potential-large-odd-primes-bad}
记
\[
c_0:=\frac{2\pi^2\kappa_2}{\log\lambda}
:=\frac{(6\sqrt5-5)\,4\pi^2}{250\log(\varphi^2)}
\approx 1.3809571881649567.
\]
再定义
\[
\theta_m:=\max_{1\le j\le m-1}\theta_{m,j},
\qquad
\mathcal B:=\{m\ge 2:\ \theta_m>\tfrac12\}.
\]
则对每个奇素数 \(p\) 都有
\[
\theta_{p,1}
:=
1-\frac{c_0}{p^2}+O(p^{-4})
\qquad(p\to\infty).
\]
特别地，
\[
\theta_{p,1}\longrightarrow 1
\qquad(p\to\infty),
\]
从而存在 \(p_0\) 使得一切奇素数 \(p>p_0\) 都满足
\[
\theta_p\ge \theta_{p,1}>\frac12.
\]
等价地，所有充分大的奇素数模数都属于坏模数集合 \(\mathcal B\)。
\end{theorem}
\begin{proof}
由定理 \ref{thm:xi-output-potential-dirichlet-fixed-j-fourth-order} 的 \(j=1\) 情形，
\[
\theta_{m,1}
:=
1-\frac{2\pi^2\kappa_2}{\log\lambda}\,m^{-2}
+\frac{2\pi^4\kappa_4}{3\log\lambda}\,m^{-4}
+O(m^{-6}).
\]
将
\(
c_0=2\pi^2\kappa_2/\log\lambda
\)
代入并把四阶及更高阶项并入误差，即得
\[
\theta_{m,1}=1-\frac{c_0}{m^2}+O(m^{-4}).
\]
取 \(m=p\) 为奇素数便得到所示展开。由于右端趋于 \(1\)，故 \(\theta_{p,1}\to 1\)。于是对充分大的奇素数 \(p\) 有 \(\theta_{p,1}>1/2\)，再由
\[
\theta_p=\max_{1\le j\le p-1}\theta_{p,j}\ge \theta_{p,1}
\]
即知 \(p\in\mathcal B\)。证毕。
\end{proof}

\begin{corollary}[坏模数集合具有自然密度 \texorpdfstring{$1$}{1}]\label{cor:xi-output-potential-bad-moduli-density-one}
沿用定理 \ref{thm:xi-output-potential-large-odd-primes-bad} 的记号，则
\[
\#\bigl(\mathcal B\cap [1,x]\bigr)=x-o(x)
\qquad(x\to\infty),
\]
从而 \(\mathcal B\) 的自然密度存在且等于 \(1\)。
\end{corollary}
\begin{proof}
由定理 \ref{thm:xi-output-potential-large-odd-primes-bad}，存在 \(p_0\) 使所有奇素数 \(p>p_0\) 都属于 \(\mathcal B\)。再由结论
\ref{con:xi-time-part9l-dirichlet-twist-divisibility-monotonicity}，若 \(m_0\in\mathcal B\)，则其任一倍数仍属于 \(\mathcal B\)。因此任何含有某个奇素因子 \(p>p_0\) 的整数都属于 \(\mathcal B\)。

于是补集 \(\NN_{\ge 2}\setminus\mathcal B\) 只能由所有素因子都落在有限集合
\[
\mathcal P_0:=\{2\}\cup\{p\in\mathbb P:\ p\le p_0\}
\]
中的整数构成。对任意 \(x\ge 2\)，这类整数个数满足粗上界
\[
\#\{n\le x:\ \text{全部素因子都属于 }\mathcal P_0\}
\le
\prod_{p\in\mathcal P_0}
\left(1+\left\lfloor\frac{\log x}{\log p}\right\rfloor\right)
:=
O\!\bigl((\log x)^{|\mathcal P_0|}\bigr)
:=
o(x).
\]
因此
\[
\#\bigl(\mathcal B\cap [1,x]\bigr)=x-o(x),
\]
故 \(\mathcal B\) 的自然密度为 \(1\)。证毕。
\end{proof}

\begin{theorem}[四阶修正相对二阶近似的二次量级增益]\label{thm:xi-output-potential-dirichlet-fourth-order-gain}
沿用定理 \ref{thm:xi-output-potential-dirichlet-fixed-j-fourth-order} 的记号，并记
\[
A:=2\pi^2\kappa_2,
\qquad
B:=\frac{2\pi^4\kappa_4}{3}.
\]
定义 \(j=1\) 通道上的二阶与四阶近似
\[
\Bigl(\frac{\rho_{m,1}}{\lambda}\Bigr)^{(2)}
:=
\exp(-A m^{-2}),
\qquad
\Bigl(\frac{\rho_{m,1}}{\lambda}\Bigr)^{(4)}
:=
\exp(-A m^{-2}+B m^{-4}),
\]
再定义对应误差
\[
\widetilde\varepsilon_m^{(2)}
:=
\left\lvert
\frac{\rho_{m,1}}{\lambda}
-\Bigl(\frac{\rho_{m,1}}{\lambda}\Bigr)^{(2)}
\right\rvert,
\qquad
\widetilde\varepsilon_m^{(4)}
:=
\left\lvert
\frac{\rho_{m,1}}{\lambda}
-\Bigl(\frac{\rho_{m,1}}{\lambda}\Bigr)^{(4)}
\right\rvert.
\]
则当 \(m\to\infty\) 时有
\[
\boxed{\;
\widetilde\varepsilon_m^{(2)}=B\,m^{-4}+O(m^{-6}),
\qquad
\widetilde\varepsilon_m^{(4)}=O(m^{-6}).\;}
\]
特别地，
\[
\boxed{\;
\frac{\widetilde\varepsilon_m^{(2)}}{\widetilde\varepsilon_m^{(4)}}=\Omega(m^2).\;}
\]
亦即，四阶修正相对于二阶高斯近似的误差压缩至少具有二次幂量级增益。
\end{theorem}
\begin{proof}
由定理 \ref{thm:xi-output-potential-dirichlet-fixed-j-fourth-order} 在 \(j=1\) 时的展开，
\[
\log\frac{\rho_{m,1}}{\lambda}
:=
-A m^{-2}+B m^{-4}+O(m^{-6}).
\]
于是
\[
\frac{\rho_{m,1}}{\lambda}
:=
\exp\!\bigl(-A m^{-2}+B m^{-4}+O(m^{-6})\bigr).
\]
与二阶近似比较，得到
\[
\widetilde\varepsilon_m^{(2)}
:=
e^{-A m^{-2}}
\left|e^{B m^{-4}+O(m^{-6})}-1\right|
:=
B\,m^{-4}+O(m^{-6}),
\]
其中最后一步使用了 \(B>0\)。

与四阶近似比较，则有
\[
\widetilde\varepsilon_m^{(4)}
:=
e^{-A m^{-2}+B m^{-4}}
\left|e^{O(m^{-6})}-1\right|
:=
O(m^{-6}).
\]
因此存在常数 \(c,C>0\) 使得对充分大的 \(m\),
\[
\widetilde\varepsilon_m^{(2)}\ge c\,m^{-4},
\qquad
\widetilde\varepsilon_m^{(4)}\le C\,m^{-6},
\]
从而
\[
\frac{\widetilde\varepsilon_m^{(2)}}{\widetilde\varepsilon_m^{(4)}}
\ge
\frac{c}{C}\,m^2.
\]
这正是所述 \(\Omega(m^2)\) 增益。证毕。
\end{proof}
\noindent 附录注 \ref{rem:real-input-40-dirichlet-cumulant-error-order} 进一步给出 \(m=5,10,20\) 的审计值：在这些模数上 \(j_{\mathrm{worst}}=1\)，故上面的 \(j=1\) 通道误差与最坏模态误差一致，并对应约 \(5.3\times 10^2\)、\(1.3\times 10^3\)、\(3.7\times 10^3\) 的实测改进因子。
\begin{theorem}[饱和端点的次谱近简并与 \texorpdfstring{$\sqrt u$}{sqrt(u)} 级弛豫时间]\label{thm:xi-output-potential-endpoint-near-degeneracy-sqrtu-relaxation}
沿用命题 \ref{prop:real-input-40-output-mirror-branch} 的记号。设 \(\lambda_1(u)\) 为 Perron 主特征值，\(\Lambda_2(u)\) 为最大非主模特征值，并定义
\[
\gamma(u):=1-\frac{\Lambda_2(u)}{\lambda_1(u)},
\qquad
\tau_{\mathrm{rel}}(u):=\gamma(u)^{-1}.
\]
则当 \(u\to+\infty\) 时，
\[
\Lambda_2(u)=|\lambda_{\mathrm{mir}}(u)|,
\qquad
\gamma(u)=\frac{2d}{c}\,u^{-1/2}+O(u^{-1}),
\]
其中
\[
c\approx 1.8963529935137691,
\qquad
d\approx 0.8079433173235080.
\]
特别地，
\[
\boxed{\;
\tau_{\mathrm{rel}}(u)=\frac{c}{2d}\,u^{1/2}+O(1).\;}
\]
亦即饱和端点的次谱与主谱在相对意义下以 \(u^{-1/2}\) 的速度碰撞，而相应弛豫时间恰按 \(\sqrt u\) 量级发散。
\end{theorem}
\begin{proof}
命题 \ref{prop:real-input-40-output-mirror-branch} 已给出
\[
\Lambda_2(u)=|\lambda_{\mathrm{mir}}(u)|,
\qquad
1-\frac{\Lambda_2(u)}{\lambda_1(u)}
=
\frac{2d}{c}\,u^{-1/2}+O(u^{-1}).
\]
由于 \(2d/c>0\)，对右端取倒数便得到
\[
\tau_{\mathrm{rel}}(u)
=
\left(\frac{2d}{c}\,u^{-1/2}+O(u^{-1})\right)^{-1}
=
\frac{c}{2d}\,u^{1/2}+O(1).
\]
证毕。
\end{proof}

\begin{corollary}[饱和端点邻域不存在任何统一次谱幂指数节约]\label{cor:xi-output-potential-endpoint-no-uniform-subunit-spectral-exponent}
沿用定理 \ref{thm:xi-output-potential-endpoint-near-degeneracy-sqrtu-relaxation} 的记号。对任意固定 \(\beta<1\)，存在 \(u_0(\beta)\)，使得当 \(u\ge u_0(\beta)\) 时恒有
\[
\Lambda_2(u)>\lambda_1(u)^\beta.
\]
因此，不存在任何与 \(u\) 无关的指数 \(\beta<1\) 能在饱和端点邻域给出统一的次谱上界
\[
\Lambda_2(u)\le \lambda_1(u)^\beta.
\]
\end{corollary}
\begin{proof}
由定理 \ref{thm:xi-output-potential-endpoint-near-degeneracy-sqrtu-relaxation} 与命题 \ref{prop:real-input-40-pressure-freezing}，
\[
\log\frac{\Lambda_2(u)}{\lambda_1(u)^\beta}
=
(1-\beta)\log\lambda_1(u)+\log\frac{\Lambda_2(u)}{\lambda_1(u)},
\]
其中
\[
(1-\beta)\log\lambda_1(u)
=
\frac{1-\beta}{2}\log u+O(1),
\]
而
\[
\log\frac{\Lambda_2(u)}{\lambda_1(u)}
=
\log\!\Bigl(1-\frac{2d}{c}\,u^{-1/2}+O(u^{-1})\Bigr)
=
-\frac{2d}{c}\,u^{-1/2}+O(u^{-1}).
\]
第一项随 \(u\to+\infty\) 趋于 \(+\infty\)，第二项趋于 \(0\)。故
\[
\log\frac{\Lambda_2(u)}{\lambda_1(u)^\beta}\to +\infty,
\]
从而对充分大的 \(u\) 必有 \(\Lambda_2(u)>\lambda_1(u)^\beta\)。证毕。
\end{proof}
\paragraph{镜像 Perron 对的奇偶抑制与介观二周期干涉}
在饱和端点的镜像 Perron 分支、次谱近简并与 \(\sqrt u\) 级弛豫时间已经建立之后，还可把主支与镜像支的联合贡献进一步压缩为下列两条奇偶接口。

\begin{theorem}[零温极限中主--镜像双分支幂和的奇偶抑制律]\label{thm:xi-time-part50cb-mirror-branch-evenodd-power-sum-suppression}
沿用命题 \ref{prop:real-input-40-output-mirror-branch} 的记号。设
\[
\lambda_1(u)=c\sqrt u+d+\frac{e}{\sqrt u}+\frac{f}{u}+O(u^{-3/2}),
\]
\[
\lambda_{\mathrm{mir}}(u)=-c\sqrt u+d-\frac{e}{\sqrt u}+\frac{f}{u}+O(u^{-3/2}),
\qquad
c>1,
\]
并定义双分支幂和
\[
S_n(u):=\lambda_1(u)^n+\lambda_{\mathrm{mir}}(u)^n.
\]
则对每个固定整数 \(r\ge 0\)，都有
\[
\boxed{\;
S_{2r}(u)=2c^{2r}u^r+O(u^{r-1/2}),\;}
\]
\[
\boxed{\;
S_{2r+1}(u)=2(2r+1)d\,c^{2r}u^r+O(u^{r-1/2}).\;}
\]
因此，按 Perron 标度归一化后，
\[
\boxed{\;
\frac{S_{2r+1}(u)}{\lambda_1(u)^{2r+1}}=O(u^{-1/2}),
\qquad
\frac{S_{2r}(u)}{\lambda_1(u)^{2r}}\longrightarrow 2.\;}
\]
换言之，镜像 Perron 分支并不只是在模值上逼近主支；它在奇阶幂和上精确造成半个幂次的系统性损失。
\end{theorem}
\begin{proof}
记
\[
a:=c\sqrt u,
\qquad
b:=d+\frac{e}{\sqrt u}+\frac{f}{u}+O(u^{-3/2}),
\qquad
b':=d-\frac{e}{\sqrt u}+\frac{f}{u}+O(u^{-3/2}),
\]
则
\[
\lambda_1(u)=a+b,
\qquad
\lambda_{\mathrm{mir}}(u)=-a+b'.
\]

对偶数 \(2r\)，二项展开给出
\[
(a+b)^{2r}=a^{2r}+2r a^{2r-1}b+O(a^{2r-2}),
\]
\[
(-a+b')^{2r}=a^{2r}-2r a^{2r-1}b'+O(a^{2r-2}),
\]
故
\[
S_{2r}(u)=2a^{2r}+2r a^{2r-1}(b-b')+O(a^{2r-2}).
\]
而
\[
b-b'=\frac{2e}{\sqrt u}+O(u^{-3/2})=O(a^{-1}),
\]
于是
\[
S_{2r}(u)=2a^{2r}+O(a^{2r-1})=2c^{2r}u^r+O(u^{r-1/2}).
\]

对奇数 \(2r+1\)，最高阶项严格抵消：
\[
(a+b)^{2r+1}=a^{2r+1}+(2r+1)a^{2r}b+O(a^{2r-1}),
\]
\[
(-a+b')^{2r+1}=-a^{2r+1}+(2r+1)a^{2r}b'+O(a^{2r-1}),
\]
从而
\[
S_{2r+1}(u)=(2r+1)a^{2r}(b+b')+O(a^{2r-1}).
\]
再由
\[
b+b'=2d+O(u^{-1}),
\]
得到
\[
S_{2r+1}(u)=2(2r+1)d\,c^{2r}u^r+O(u^{r-1/2}).
\]

最后，由 \(\lambda_1(u)=c\sqrt u+O(1)\)，可知
\[
\lambda_1(u)^{2r}\sim c^{2r}u^r,
\qquad
\lambda_1(u)^{2r+1}\sim c^{2r+1}u^{r+1/2},
\]
代回即得归一化后的两条结论。证毕。
\end{proof}

\begin{corollary}[介观长度上的严格二周期干涉律]\label{cor:xi-time-part50cb-mesoscopic-period2-interference}
沿用定理 \ref{thm:xi-time-part50cb-mirror-branch-evenodd-power-sum-suppression} 的记号。设整数值函数 \(n(u)\ge 1\) 满足
\[
n(u)=o(\sqrt u)
\qquad (u\to+\infty).
\]
则有
\[
\left|
\frac{\lambda_{\mathrm{mir}}(u)^{n(u)}}{\lambda_1(u)^{n(u)}}-(-1)^{n(u)}
\right|
\longrightarrow 0,
\]
以及
\[
\left|
\frac{S_{n(u)}(u)}{\lambda_1(u)^{n(u)}}-\bigl(1+(-1)^{n(u)}\bigr)
\right|
\longrightarrow 0.
\]
因此，在介观时间窗 \(n=o(\sqrt u)\) 上，主支与镜像支对总谱贡献呈现严格的二周期干涉：偶长度归一化极限为 \(2\)，奇长度归一化极限为 \(0\)。
\end{corollary}
\begin{proof}
由命题 \ref{prop:real-input-40-output-mirror-branch} 的相对谱隙展开，
\[
1-\frac{|\lambda_{\mathrm{mir}}(u)|}{\lambda_1(u)}
=
\frac{2d}{c}u^{-1/2}+O(u^{-1}),
\]
故
\[
\frac{\lambda_{\mathrm{mir}}(u)}{\lambda_1(u)}
=
-\left(1-\frac{2d}{c}u^{-1/2}+O(u^{-1})\right).
\]
于是
\[
\frac{\lambda_{\mathrm{mir}}(u)^{n(u)}}{\lambda_1(u)^{n(u)}}
=
(-1)^{n(u)}
\left(1-\frac{2d}{c}u^{-1/2}+O(u^{-1})\right)^{n(u)}.
\]
由于 \(n(u)=o(\sqrt u)\)，括号内第二因子的对数满足
\[
n(u)\,O(u^{-1/2})=o(1),
\]
故其趋于 \(1\)。这就得到第一条。

第二条则由恒等式
\[
\frac{S_{n(u)}(u)}{\lambda_1(u)^{n(u)}}
=
1+\frac{\lambda_{\mathrm{mir}}(u)^{n(u)}}{\lambda_1(u)^{n(u)}}
\]
立即推出。证毕。
\end{proof}

\paragraph{输出位势的闭合解析链}
上述结果表明：饱和端点的能垒缺口、零温端点的平方根临界律、镜像 Perron 支诱导的次谱近简并与 \(\sqrt u\) 级弛豫时间、任意 \(\beta<1\) 统一次谱幂界在端点邻域的失效、局部 Legendre 几何的二次数域封闭、均值附近的非对称陡峭性与 Edgeworth 半质量中心的正常数右移、唯一小根的激活支结构，以及固定单位根模态的四阶展开、有界窗口内 \(j=1\) 的最终支配和四阶修正的二次量级收益，实际上都由同一条输出位势代数曲线统一控制。端点零温极限、中心局部 Legendre 对偶几何、中心极限定域、慢模近简并与单位根方向的 Dirichlet 扭曲并非彼此孤立的机制，而是同一解析对象在实轴、局部对偶坐标域与 cyclotomic 边界上的不同投影。

更具体地说，局部累积量闭式
\[
P^{(k)}(0)\in\mathbb Q(\sqrt5)
\]
经由定理 \ref{thm:xi-output-potential-dirichlet-fixed-j-fourth-order} 传递为 \(\theta_{m,j}\) 的四阶展开，再经推论 \ref{cor:xi-output-potential-dirichlet-bounded-window-j1} 固定出有界窗口内的最坏模态，并最终在定理 \ref{thm:xi-output-potential-dirichlet-fourth-order-gain} 中表现为四阶修正相对二阶近似的二次量级收益。
\begin{theorem}[奇素数模扭曲的 RH 障碍指数趋于 \texorpdfstring{$1$}{1}，因而充分大奇素数全部为坏模数]\label{thm:xi-output-potential-odd-prime-rh-obstruction-index-to-one}
\leanverified{paper\_xi\_output\_potential\_odd\_prime\_rh\_obstruction\_index\_to\_one}
沿用定理 \ref{thm:xi-output-potential-dirichlet-fixed-j-fourth-order} 的记号，并记
\[
\theta_p:=\max_{1\le j\le p-1}\theta_{p,j}
\qquad(p\ge 3\ \text{为奇素数}).
\]
则当奇素数 \(p\to\infty\) 时，有显式四阶展开
\[
\boxed{\;
\theta_{p,1}
=
1-\frac{6\sqrt5-5}{250\log\lambda}\Bigl(\frac{2\pi}{p}\Bigr)^2
+\frac{7+24\sqrt5}{75000\log\lambda}\Bigl(\frac{2\pi}{p}\Bigr)^4
+O(p^{-6}).\;}
\]
从而
\[
\boxed{\;
\theta_p\ge \theta_{p,1}\longrightarrow 1
\qquad(p\to\infty,\ p\ \text{为奇素数}).\;}
\]
特别地，存在 \(p_0\) 使得一切奇素数 \(p\ge p_0\) 都满足
\[
\theta_p>\frac12.
\]
换言之，除有限多个例外外，所有奇素数模数都属于平方根门槛失效的坏模数层。
\end{theorem}
\begin{proof}
由定理 \ref{thm:xi-output-potential-dirichlet-fixed-j-fourth-order} 的 \(j=1\) 情形，
\[
\theta_{m,1}
=
1-\frac{2\pi^2\kappa_2}{\log\lambda}\,m^{-2}
+\frac{2\pi^4\kappa_4}{3\log\lambda}\,m^{-4}
+O(m^{-6}).
\]
再由命题 \ref{prop:real-input-40-output-cumulants-closed} 中的显式常数
\[
\kappa_2=\frac{6\sqrt5-5}{125},
\qquad
\kappa_4=\frac{7+24\sqrt5}{3125},
\]
代入并取 \(m=p\) 为奇素数，即得
\[
\theta_{p,1}
=
1-\frac{6\sqrt5-5}{250\log\lambda}\Bigl(\frac{2\pi}{p}\Bigr)^2
+\frac{7+24\sqrt5}{75000\log\lambda}\Bigl(\frac{2\pi}{p}\Bigr)^4
+O(p^{-6}).
\]
因此 \(\theta_{p,1}\to 1\)。另一方面，由定义
\[
\theta_p=\max_{1\le j\le p-1}\theta_{p,j}\ge \theta_{p,1},
\]
遂得 \(\theta_p\to 1\)。于是对充分大的奇素数 \(p\)，必有 \(\theta_p>\tfrac12\)。证毕。
\end{proof}

\begin{corollary}[奇素数坏模数在素数集中的自然密度为 \texorpdfstring{$1$}{1}]\label{cor:xi-output-potential-odd-prime-bad-moduli-prime-density-one}
\leanverified{paper\_xi\_output\_potential\_odd\_prime\_bad\_moduli\_prime\_density\_one}
定义奇素数坏模数集
\[
\mathcal P_{\mathrm{bad}}
:=
\{p\ge 3:\ p\ \text{为素数且}\ \theta_p>\tfrac12\}.
\]
则有
\[
\boxed{\;
\lim_{X\to\infty}
\frac{\#\{p\le X:\ p\in\mathcal P_{\mathrm{bad}}\}}{\pi(X)}
=1.\;}
\]
\end{corollary}
\begin{proof}
由定理 \ref{thm:xi-output-potential-odd-prime-rh-obstruction-index-to-one}，存在 \(p_0\) 使得所有奇素数 \(p\ge p_0\) 都属于 \(\mathcal P_{\mathrm{bad}}\)。因此其补集至多包含有限多个奇素数，再加上素数 \(2\)。故
\[
\#\{p\le X:\ p\notin\mathcal P_{\mathrm{bad}}\}=O(1),
\]
从而
\[
\frac{\#\{p\le X:\ p\in\mathcal P_{\mathrm{bad}}\}}{\pi(X)}
=
1-\frac{O(1)}{\pi(X)}
\longrightarrow 1.
\]
证毕。
\end{proof}

\paragraph{零温 witness 复杂度、Ramanujan 谱缝与平衡均值间隙}
在输出位势的端点大偏差、地基熵、局部 Legendre 几何与碰撞位势的 RH/Ramanujan 临界窗已经建立之后，还可进一步把零温极端层的可恢复性、可采样性与谱窗口失效压缩为下列若干条结果。统一记
\[
\mathcal G_n:=
\{\gamma:\ |\gamma|=n,\ \text{其输出 }1\text{ 的总数为 }k_{\max}(n)\}
\]
为极端饱和层。

\begin{theorem}[极端饱和层的完备外置阈值恰为 \texorpdfstring{$\log_2 c$}{log2 c} 比特/步]\label{thm:xi-real-input-ground-layer-oracle-threshold}
\leanverified{paper\_xi\_real\_input\_ground\_layer\_oracle\_threshold}
沿用推论 \ref{cor:real-input-40-ground-entropy} 的记号，并令 \(\Gamma_n\) 在 \(\mathcal G_n\) 上均匀分布。对任意只允许 \(B_n\) 比特外置寄存器的编码--解码机制，记恢复输出为 \(\widehat\Gamma_n\)。则有
\[
\boxed{\;
\mathbb P(\widehat\Gamma_n=\Gamma_n)
\le
\frac{2^{B_n}}{A_{n,k_{\max}(n)}}.
\;}
\]
从而：
\begin{enumerate}
  \item 若存在 \(\varepsilon>0\) 使
  \[
  B_n\le (\log_2 c-\varepsilon)n
  \]
  对充分大的 \(n\) 成立，则
  \[
  \mathbb P(\widehat\Gamma_n=\Gamma_n)
  \le
  \exp\!\bigl(-\varepsilon(\log 2)\,n+o(n)\bigr).
  \]
  \item 若希望存在固定 \(\eta\in(0,1]\) 使
  \[
  \mathbb P(\widehat\Gamma_n=\Gamma_n)\ge \eta
  \]
  对充分大的 \(n\) 成立，则必有
  \[
  B_n\ge n\log_2 c+\log_2\eta+o(n).
  \]
\end{enumerate}
换言之，极端饱和层的最小完备外置率恰为 \(\log_2 c\) 比特/步。
\end{theorem}
\begin{proof}
由定义，
\[
|\mathcal G_n|=A_{n,k_{\max}(n)}.
\]
设外置寄存器的取值集合为 \(\mathcal R_n\)，则
\(
|\mathcal R_n|\le 2^{B_n}
\)。
对任意确定性编码--解码机制，每个寄存器值 \(r\in\mathcal R_n\) 至多只能对应一个被正确恢复的候选路径；因此在 \(\mathcal G_n\) 中至多只有 \(2^{B_n}\) 条路径可能被正确译回。由于 \(\Gamma_n\) 在 \(\mathcal G_n\) 上均匀分布，得到
\[
\mathbb P(\widehat\Gamma_n=\Gamma_n)\le \frac{2^{B_n}}{|\mathcal G_n|}
=\frac{2^{B_n}}{A_{n,k_{\max}(n)}}.
\]
随机机制只是确定性机制的凸组合，故同一上界仍然成立。

再由推论 \ref{cor:real-input-40-ground-entropy}，
\[
A_{n,k_{\max}(n)}=\exp\bigl((\log c)n+o(n)\bigr).
\]
将其代回上式，便得到
\[
\mathbb P(\widehat\Gamma_n=\Gamma_n)
\le
\exp\bigl(B_n\log 2-(\log c)n+o(n)\bigr).
\]
若 \(B_n\le (\log_2 c-\varepsilon)n\)，则
\[
B_n\log 2-(\log c)n\le -\varepsilon(\log 2)\,n,
\]
从而得到第一条。
反之，若成功概率对充分大的 \(n\) 满足
\(
\mathbb P(\widehat\Gamma_n=\Gamma_n)\ge \eta
\)，则必有
\[
\eta
\le
\frac{2^{B_n}}{A_{n,k_{\max}(n)}},
\]
即
\[
B_n\ge \log_2 A_{n,k_{\max}(n)}+\log_2\eta
=n\log_2 c+\log_2\eta+o(n).
\]
证毕。
\end{proof}

\begin{theorem}[极端饱和层在无偏路径总体中的稀有指数等于 \texorpdfstring{$\log(\varphi^2/c)$}{log(phi^2/c)}]\label{thm:xi-real-input-ground-layer-unbiased-rarity}
\leanverified{paper\_xi\_real\_input\_ground\_layer\_unbiased\_rarity}
沿用命题 \ref{prop:real-input-40-pressure-freezing} 与推论 \ref{cor:real-input-40-ground-entropy} 的记号，记
\[
Z_n(1):=\mathbf{1}^{\mathsf T}M(1)^n\mathbf{1},
\qquad
p_n:=\frac{A_{n,k_{\max}(n)}}{Z_n(1)}.
\]
则有精确指数律
\[
\boxed{\;
\lim_{n\to\infty}\frac{1}{n}\log p_n
=
\log c-\log\lambda(1)
=
-\delta_\ast,
\;}
\]
其中
\[
\lambda(1)=\frac{3+\sqrt5}{2}=\varphi^2,
\qquad
\delta_\ast:=\log\frac{\varphi^2}{c}\approx 0.32249108560171.
\]
特别地，
\[
p_n=\exp(-\delta_\ast n+o(n)).
\]
若在全部长度 \(n\) 的可行路径上做独立无偏抽样，直到首次命中 \(\mathcal G_n\) 的等待时间记为 \(T_n\)，则
\[
\mathbb E[T_n]=\frac{1}{p_n}=\exp(\delta_\ast n+o(n)).
\]
\end{theorem}
\begin{proof}
由命题 \ref{prop:real-input-40-pressure-freezing}，
\[
\lim_{n\to\infty}\frac{1}{n}\log Z_n(1)=P(0)=\log\lambda(1)=\log(\varphi^2).
\]
另一方面，推论 \ref{cor:real-input-40-ground-entropy} 给出
\[
\lim_{n\to\infty}\frac{1}{n}\log A_{n,k_{\max}(n)}=\log c.
\]
两式相减即得
\[
\lim_{n\to\infty}\frac{1}{n}\log p_n
=
\log c-\log(\varphi^2)
=
-\delta_\ast.
\]
这与推论 \ref{cor:real-input-40-rate-endpoints} 的端点代价
\(
I(1/2)=\log(\varphi^2/c)
\)
完全一致。
若独立重复无偏抽样，则 \(T_n\) 服从参数为 \(p_n\) 的几何分布，因而
\(
\mathbb E[T_n]=1/p_n
\)，从而得到所示指数律。证毕。
\end{proof}

\begin{theorem}[RH-sharp 临界参数严格位于零温四次锚点之前]\label{thm:xi-real-input-rh-window-zero-temp-perron-gap}
沿用命题 \ref{prop:real-input-40-collision-rhk-window} 与定理 \ref{thm:xi-zero-temperature-fourstate-artin-mazur-zeta-primitive-asymptotic} 的记号。设 \(u_R>1\) 为唯一满足
\[
u_R^5+4u_R^4+3u_R^3-96u_R^2-42u_R-14=0
\]
的正实根，并记
\[
\mathcal Q(u):=u^5+4u^4+3u^3-96u^2-42u-14,
\]
\[
\mathcal R(y):=y^4-6y^3+9y^2-y-1,
\qquad
y_\infty:=\rho(B_\infty)=c^2.
\]
则有严格有理夹逼
\[
\boxed{\;
\frac{359}{100}<u_R<\frac{3593}{1000}<y_\infty<\frac{18}{5}.\;}
\]
特别地，
\[
\boxed{\;
u_R<y_\infty=\rho(B_\infty)=c^2.
\;}
\]
因此在严格非空区间
\[
u\in\bigl(u_R,\rho(B_\infty)\bigr)
\]
上，核版 RH/Ramanujan 条件已经失效，但零温 quartic 地基尺度尚未达到其自身 Perron 阈值。
\end{theorem}
\begin{proof}
由命题 \ref{prop:real-input-40-collision-rhk-window}，\(\mathcal Q\) 在 \((0,\infty)\) 上恰有一个正根 \(u_R\)。直接代入得到
\[
\mathcal Q\!\left(\frac{359}{100}\right)
=
-\frac{25100699801}{10^{10}}<0,
\qquad
\mathcal Q\!\left(\frac{3593}{1000}\right)
=
\frac{362872814205193}{10^{15}}>0.
\]
由介值定理与唯一正根性，立得
\[
\frac{359}{100}<u_R<\frac{3593}{1000}.
\]

另一方面，定理 \ref{thm:xi-zero-temperature-fourstate-artin-mazur-zeta-primitive-asymptotic} 给出 \(y_\infty\) 是四次方程
\[
\mathcal R(y)=y^4-6y^3+9y^2-y-1=0
\]
的最大正根。
\[
\mathcal R\!\left(\frac{3593}{1000}\right)
=
-\frac{53334838799}{10^{12}}<0,
\qquad
\mathcal R\!\left(\frac{18}{5}\right)
=
\frac{41}{625}>0.
\]
因此同样由介值定理可得
\[
\frac{3593}{1000}<y_\infty<\frac{18}{5}.
\]
\[
u_R<y_\infty=\rho(B_\infty)=c^2.
\]
\[
u_R<\frac{3593}{1000}<y_\infty=\rho(B_\infty)=c^2.
\]
于是所述严格非空区间
\[
u\in\bigl(u_R,\rho(B_\infty)\bigr)
\]
立刻成立。证毕。
\end{proof}

\begin{theorem}[RH-sharp 阈与零温冻结阈的代数分裂]\label{thm:xi-real-input-rh-threshold-zero-temp-algebraic-splitting}
\leanverified{paper\_xi\_real\_input\_rh\_threshold\_zero\_temp\_algebraic\_splitting}
沿用定理 \ref{thm:xi-real-input-rh-window-zero-temp-perron-gap} 的记号。则
\[
\gcd_{\ZZ[X]}\bigl(\mathcal Q(X),\mathcal R(X)\bigr)=1.
\]
等价地，
\[
\operatorname{Res}\bigl(\mathcal Q,\mathcal R\bigr)\neq 0.
\]
因此 \(u_R\) 与 \(y_\infty=\rho(B_\infty)\) 不可能由同一个代数根同时满足这两条定义方程；特别地，
\[
u_R\neq y_\infty.
\]
结合定理 \ref{thm:xi-real-input-rh-window-zero-temp-perron-gap}，可知 RH-sharp 谱阈与零温冻结阈不仅在实轴上严格分离，而且在代数上亦不可合并。
\end{theorem}
\begin{proof}
对整系数多项式 \(\mathcal Q\) 与 \(\mathcal R\) 施行 Euclid 算法可得
\[
\gcd_{\ZZ[X]}\bigl(\mathcal Q(X),\mathcal R(X)\bigr)=1,
\]
等价地，
\[
\operatorname{Res}\bigl(\mathcal Q,\mathcal R\bigr)=-720892\neq 0.
\]
故 \(\mathcal Q\) 与 \(\mathcal R\) 没有公共代数根。

另一方面，定理 \ref{thm:xi-real-input-rh-window-zero-temp-perron-gap} 已给出
\[
\mathcal Q(u_R)=0,
\qquad
\mathcal R(y_\infty)=0,
\qquad
u_R<y_\infty.
\]
若 \(u_R=y_\infty\)，则该代数数将同时是 \(\mathcal Q\) 与 \(\mathcal R\) 的公共根，矛盾。故
\[
u_R\neq y_\infty.
\]
于是两条阈值边界在定义方程层面已经代数分裂，而定理 \ref{thm:xi-real-input-rh-window-zero-temp-perron-gap} 又进一步给出其实轴顺序为
\(
u_R<y_\infty
\)。
证毕。
\end{proof}

\paragraph{输出位势速率函数的端点对数尖点与不对称刚性}
在定理 \ref{thm:xi-real-input-rh-window-zero-temp-perron-gap} 已把 RH-sharp 五次阈值与零温四次锚点严格分离、定理 \ref{thm:xi-output-potential-zero-temp-square-root-critical-law} 已锁定右端零温尾部的 \(e^{-\theta/2}\) 量级之后，归一化速率函数本身在两端还呈现出更细且彼此不对称的对数边界几何。

\begin{theorem}[输出位势速率函数在右端点具有显式对数尖点]\label{thm:xi-output-potential-rate-right-endpoint-log-cusp}
沿用命题 \ref{prop:real-input-40-pressure-freezing}、推论 \ref{cor:real-input-40-rate-endpoints} 与结论 \ref{con:xi-terminal-real-input-output-microcanonical-logcusp-second-order} 的记号。记
\[
I(\alpha):=\sup_{\theta\in\RR}\bigl(\theta\alpha-(P(\theta)-P(0))\bigr),
\qquad
A:=\frac{d}{2c}>0,
\]
并令
\[
\varepsilon:=\frac12-\alpha.
\]
则当 \(\alpha\uparrow \frac12\)（即 \(\varepsilon\downarrow 0\)）时，有
\[
\boxed{\;
I(\alpha)
=
I\!\left(\frac12\right)
+2\varepsilon\log\varepsilon
-2(1+\log A)\varepsilon
+O(\varepsilon^2).\;}
\]
等价地，
\[
\boxed{\;
I\!\left(\frac12\right)-I(\alpha)
=
2\varepsilon|\log\varepsilon|
+2(1+\log A)\varepsilon
+O(\varepsilon^2).\;}
\]
并且导数满足精确发散律
\[
\boxed{\;
I'(\alpha)
=
-2\log\!\Bigl(\frac12-\alpha\Bigr)+2\log A+O\!\Bigl(\frac12-\alpha\Bigr),
\qquad
\alpha\uparrow \frac12.\;}
\]
特别地，\(I\) 在右端点具有垂直切线，而该端点不是解析边界点，而是对数型尖点。
\end{theorem}
\begin{proof}
记
\[
s(\alpha):=\inf_{\theta\in\RR}\bigl(P(\theta)-\alpha\theta\bigr).
\]
由 Legendre--Fenchel 对偶的归一化关系，
\[
I(\alpha)=P(0)-s(\alpha).
\]
另一方面，结论 \ref{con:xi-terminal-real-input-output-microcanonical-logcusp-second-order} 已给出：若记
\[
a:=\frac{d}{c}=2A,
\qquad
\delta:=\frac12-\alpha,
\]
则
\[
s\!\left(\frac12-\delta\right)
=
\log c-2\delta\log\delta+2\left(1+\log\frac{a}{2}\right)\delta+O(\delta^2).
\]
由于 \(a/2=A\)，并且推论 \ref{cor:real-input-40-rate-endpoints} 给出
\[
I\!\left(\frac12\right)=P(0)-\log c,
\]
故直接相减便得
\[
I(\alpha)
=
I\!\left(\frac12\right)
+2\varepsilon\log\varepsilon
-2(1+\log A)\varepsilon
+O(\varepsilon^2).
\]
第二个盒装式只是将 \(\log\varepsilon=-|\log\varepsilon|\) 代回后的重写。

再证导数公式。对每个 \(\alpha\in(0,\frac12)\)，由严格凸性可取唯一极值点 \(\theta(\alpha)\) 满足
\[
P'(\theta(\alpha))=\alpha,
\]
并有
\[
I'(\alpha)=\theta(\alpha).
\]
由定理 \ref{thm:xi-output-potential-zero-temp-square-root-critical-law}，
\[
\frac12-P'(\theta)=\frac{d}{2c}e^{-\theta/2}+O(e^{-\theta})
=
A e^{-\theta/2}\bigl(1+O(e^{-\theta/2})\bigr).
\]
于是当 \(\alpha\uparrow\frac12\) 时，
\[
\varepsilon
=
A e^{-\theta(\alpha)/2}\bigl(1+O(e^{-\theta(\alpha)/2})\bigr),
\]
从而
\[
e^{-\theta(\alpha)/2}=\frac{\varepsilon}{A}+O(\varepsilon^2).
\]
两边取对数得到
\[
\theta(\alpha)
=
-2\log\varepsilon+2\log A+O(\varepsilon),
\]
亦即
\[
I'(\alpha)
=
-2\log\!\Bigl(\frac12-\alpha\Bigr)+2\log A+O\!\Bigl(\frac12-\alpha\Bigr).
\]
因此 \(I'(\alpha)\to+\infty\)，右端点具有垂直切线；而 \(\varepsilon\log\varepsilon\) 项则表明该端点不可能解析。证毕。
\end{proof}

\begin{theorem}[输出位势右端点对偶参数的二阶反演与速率函数边界正规形]\label{thm:xi-output-potential-right-endpoint-second-order-inversion-normal-form}
\leanverified{paper\_xi\_output\_potential\_right\_endpoint\_second\_order\_inversion\_normal\_form}
沿用命题 \ref{prop:real-input-40-pressure-freezing}、推论 \ref{cor:real-input-40-rate-endpoints} 与结论 \ref{con:xi-terminal-real-input-output-microcanonical-logcusp-second-order} 的记号。令 \(\theta(\alpha)\) 为满足
\[
P'(\theta(\alpha))=\alpha
\]
的唯一对偶参数，并记
\[
\delta:=\frac12-\alpha\downarrow 0.
\]
则有二阶反演公式
\[
\boxed{\;
\theta\!\left(\frac12-\delta\right)
=
2\log\frac{d}{2c\,\delta}
+
\left(\frac{8ce}{d^2}-4\right)\delta
+O(\delta^2).\;}
\]
进一步，归一化速率函数满足
\[
\boxed{\;
I\!\left(\frac12-\delta\right)
=
\log\frac{\lambda(1)}{c}
-2\delta\log\frac{d}{2c\,\delta}
-2\delta
+\left(2-\frac{4ce}{d^2}\right)\delta^2
+O(\delta^3).\;}
\]
因此右端点的对数尖点不仅具有主系数 \(2\)，其二阶曲率修正也由 Puiseux 常数 \(c,d,e\) 完全显式决定。
\end{theorem}
\begin{proof}
记
\[
a:=\frac{d}{c},
\qquad
b:=\frac{e}{c}-\frac{d^2}{2c^2},
\qquad
x:=e^{-\theta/2}.
\]
由命题 \ref{prop:real-input-40-pressure-freezing}，
\[
P(\theta)
=
\frac{\theta}{2}+\log c+a x+b x^2+O(x^3),
\]
从而
\[
P'(\theta)
=
\frac12-\frac{a}{2}x-bx^2+O(x^3).
\]
于是极值条件 \(P'(\theta)=\frac12-\delta\) 等价于
\[
\delta=\frac{a}{2}x+bx^2+O(x^3).
\]
对该关系作幂级数反演可得
\[
x=\frac{2\delta}{a}-\frac{8b}{a^3}\delta^2+O(\delta^3).
\]
再取对数：
\[
\log x
=
\log\frac{2\delta}{a}-\frac{4b}{a^2}\delta+O(\delta^2),
\]
故
\[
\theta\!\left(\frac12-\delta\right)
=
-2\log x
=
2\log\frac{a}{2\delta}+\frac{8b}{a^2}\delta+O(\delta^2).
\]
将 \(a,b\) 代回，便得到
\[
\theta\!\left(\frac12-\delta\right)
=
2\log\frac{d}{2c\,\delta}
+
\left(\frac{8ce}{d^2}-4\right)\delta
+O(\delta^2).
\]

另一方面，结论 \ref{con:xi-terminal-real-input-output-microcanonical-logcusp-second-order} 已给出
\[
s\!\left(\frac12-\delta\right)
=
\log c
-2\delta\log\delta
+2\left(1+\log\frac{a}{2}\right)\delta
+\frac{4b}{a^2}\delta^2
+O(\delta^3).
\]
再由
\[
I(\alpha)=P(0)-s(\alpha),
\qquad
P(0)=\log\lambda(1),
\]
可得
\[
I\!\left(\frac12-\delta\right)
=
\log\frac{\lambda(1)}{c}
+2\delta\log\delta
-2\left(1+\log\frac{a}{2}\right)\delta
-\frac{4b}{a^2}\delta^2
+O(\delta^3).
\]
将其改写为
\[
I\!\left(\frac12-\delta\right)
=
\log\frac{\lambda(1)}{c}
-2\delta\log\frac{a}{2\delta}
-2\delta
-\frac{4b}{a^2}\delta^2
+O(\delta^3),
\]
再代回 \(a=d/c\) 与 \(b=e/c-d^2/(2c^2)\)，便得到
\[
I\!\left(\frac12-\delta\right)
=
\log\frac{\lambda(1)}{c}
-2\delta\log\frac{d}{2c\,\delta}
-2\delta
+\left(2-\frac{4ce}{d^2}\right)\delta^2
+O(\delta^3).
\]
证毕。
\end{proof}

\begin{theorem}[输出位势速率函数在左端点具有 Cram\'er--Lambert 型尖点]\label{thm:xi-output-potential-rate-left-endpoint-cramer-lambert}
沿用命题 \ref{prop:real-input-40-pressure-freezing} 与推论 \ref{cor:real-input-40-rate-endpoints} 的记号。则当 \(\alpha\downarrow 0\) 时，有
\[
\boxed{\;
I(\alpha)
=
I(0)+\alpha\log\alpha-(1+\log 4)\alpha+O(\alpha^2).\;}
\]
等价地，
\[
\boxed{\;
I(0)-I(\alpha)
=
\alpha|\log\alpha|+(1+\log 4)\alpha+O(\alpha^2).\;}
\]
并且
\[
\boxed{\;
I'(\alpha)=\log\alpha-\log4+O(\alpha)\to-\infty,\;}
\]
\[
\boxed{\;
I''(\alpha)=\frac1\alpha+O(1).\;}
\]
因此左端点同样不是解析端点，而是以主系数 \(1\) 控制的对数尖点。
\end{theorem}
\begin{proof}
由命题 \ref{prop:real-input-40-pressure-freezing} 在 \(u\to 0\) 时的展开，
\[
\lambda(u)=1+4u-5u^2-115u^3+2317u^4+O(u^5),
\]
故
\[
\log\lambda(u)=4u+O(u^2).
\]
另一方面，若记 \(\theta=\log u\)，则
\[
\alpha(u)=P'(\theta)=u\frac{\lambda'(u)}{\lambda(u)}=4u+O(u^2).
\]
因此可反解为
\[
u=\frac{\alpha}{4}+O(\alpha^2),
\qquad
\log u=\log\alpha-\log4+O(\alpha).
\]

再由 Legendre 对偶在内部点上的参数表示，
\[
I(\alpha(u))=\alpha(u)\log u-\log\lambda(u)+P(0).
\]
推论 \ref{cor:real-input-40-rate-endpoints} 给出
\[
I(0)=P(0)=\log\lambda(1)=\log(\varphi^2).
\]
并且由
\[
\log\lambda(u)=4u+O(u^2)=\alpha+O(\alpha^2)
\]
可得
\[
\begin{aligned}
I(\alpha)
&=
I(0)+\alpha\log u-\log\lambda(u)\\
&=
I(0)+\alpha(\log\alpha-\log4+O(\alpha))-\alpha+O(\alpha^2)\\
&=
I(0)+\alpha\log\alpha-(1+\log4)\alpha+O(\alpha^2).
\end{aligned}
\]
这便得到第一条与其等价重写。

对导数，仍记内部极值点为 \(\theta(\alpha)\)。则
\[
I'(\alpha)=\theta(\alpha)=\log u=\log\alpha-\log4+O(\alpha),
\]
从而 \(I'(\alpha)\to-\infty\)。再对上式求导，便得到
\[
I''(\alpha)=\frac1\alpha+O(1).
\]
因此左端点同样具有对数型非解析尖点。证毕。
\end{proof}

\begin{corollary}[输出位势速率函数两端尖点强度的刚性不对称]\label{cor:xi-output-potential-endpoint-log-cusp-rigidity-pair}
沿用定理 \ref{thm:xi-output-potential-rate-right-endpoint-log-cusp} 与定理 \ref{thm:xi-output-potential-rate-left-endpoint-cramer-lambert} 的记号。则有极限比值
\[
\boxed{\;
\lim_{\alpha\downarrow0}\frac{I(0)-I(\alpha)}{\alpha|\log\alpha|}=1,
\qquad
\lim_{\alpha\uparrow1/2}\frac{I(1/2)-I(\alpha)}{(1/2-\alpha)|\log(1/2-\alpha)|}=2.\;}
\]
因此，在本文固定的频率坐标 \(\alpha\) 中，两端对数尖点强度形成有序对
\[
\boxed{\;(1,2).\;}
\]
特别地，不存在端点交换局部对称
\[
I(0)-I(\alpha)=I\!\left(\frac12\right)-I\!\left(\frac12-\alpha\right)+o(\alpha|\log\alpha|)
\qquad
(\alpha\downarrow0),
\]
因而更不存在以
\[
\alpha\longmapsto \frac12-\alpha
\]
为主导阶模型的仿射端点反演对称。
\end{corollary}
\begin{proof}
第一条极限直接由定理 \ref{thm:xi-output-potential-rate-left-endpoint-cramer-lambert} 的主导项
\[
I(0)-I(\alpha)=\alpha|\log\alpha|+O(\alpha)
\]
得到。第二条极限则由定理 \ref{thm:xi-output-potential-rate-right-endpoint-log-cusp} 的主导项
\[
I\!\left(\frac12\right)-I(\alpha)
=
2\Bigl(\frac12-\alpha\Bigr)\Bigl|\log\!\Bigl(\frac12-\alpha\Bigr)\Bigr|
+O\!\left(\frac12-\alpha\right)
\]
立得。

若存在所述端点交换局部对称，则两侧在 \(\alpha|\log\alpha|\) 量级上的主导系数必须相同；但左端系数为 \(1\)，右端系数为 \(2\)，矛盾。故该局部对称不可能成立，而仿射反演 \(\alpha\mapsto\frac12-\alpha\) 只是其最直接的特例。证毕。
\end{proof}

\noindent 这几条结果进一步表明：输出位势的端点几何并非由单一对称冻结机制主宰。右端点同时受到零温四次 Perron 地基与 \(e^{-\theta/2}\) 尾部的控制，因此不仅产生主系数为 \(2\) 的对数尖点，而且其对偶反演与二阶曲率修正也被 \(c,d,e\) 完全显式锁定；左端点则由 \(u\to0\) 的 Cram\'er 型稀薄展开主导，仅留下主系数为 \(1\) 的对数尖点。于是，RH-sharp 五次阈值先于四次锚点出现的谱事实，与速率函数两端不可对称化的边界几何，便在同一条输出位势终端链上汇合起来。

\begin{theorem}[最大密度层与平衡均值之间存在显式代数间隙而仍保留正熵]\label{thm:xi-real-input-ground-layer-mean-gap-positive-entropy}
\leanverified{paper\_xi\_real\_input\_ground\_layer\_mean\_gap\_positive\_entropy}
沿用推论 \ref{cor:real-input-40-alpha-max}、推论 \ref{cor:real-input-40-ground-entropy} 与表 \ref{tab:real_input_40_output_potential_cumulants_closed} 的记号，记
\[
\alpha_{\max}:=\frac12,
\qquad
\alpha_0:=P'(0)=\frac{23}{20}-\frac{7\sqrt5}{20},
\]
\[
\Delta_\mu:=\alpha_{\max}-\alpha_0
:=
\frac{7\sqrt5-13}{20}
\approx 0.132623792124926.
\]
则对每个 \(n\) 及每条 \(\gamma\in\mathcal G_n\)，都有
\[
\frac{k_{\max}(n)}{n}
:=
\frac12+O\!\left(\frac{1}{n}\right)
:=
\alpha_0+\Delta_\mu+O\!\left(\frac{1}{n}\right).
\]
并且极端饱和层的总数满足
\[
|\mathcal G_n|
:=
\exp\bigl((\log c)n+o(n)\bigr).
\]
因此在距离平衡均值 \(\Delta_\mu n+O(1)\) 的线性偏差层上，仍承载指数多的可行路径。
\end{theorem}
\begin{proof}
由推论 \ref{cor:real-input-40-ground-entropy}，
\[
k_{\max}(n)=\lfloor n/2\rfloor,
\qquad
|\mathcal G_n|=A_{n,k_{\max}(n)}
:=
\exp\bigl((\log c)n+o(n)\bigr).
\]
因此
\[
\frac{k_{\max}(n)}{n}=\frac12+O\!\left(\frac{1}{n}\right).
\]
另一方面，表 \ref{tab:real_input_40_output_potential_cumulants_closed} 给出
\[
P'(0)=\frac{23}{20}-\frac{7\sqrt5}{20},
\]
故
\[
\frac12-P'(0)=\frac{7\sqrt5-13}{20}=\Delta_\mu>0.
\]
将两式合并，即得
\[
\frac{k_{\max}(n)}{n}
:=
P'(0)+\Delta_\mu+O\!\left(\frac{1}{n}\right),
\]
并且同一偏差层的路径总数仍具有指数增长率 \(\log c\)。证毕。
\end{proof}
\paragraph{bin-fold 饱和窗口、零温微正则尖点与近地基壳层上界}
在 bin-fold 最大纤维的精确饱和判据、输出位势右端点的对数尖点以及地基层的正熵增长率都已固定之后，还可进一步把离散最大纤维、零温 Puiseux 展开与近端壳层计数压缩为同一条显式链。其共同特征在于：饱和窗口只由单个 Fibonacci 差值决定，而同一差值又经 Legendre--Fenchel 变换传递为右端点的对数尖点与近地基壳层的指数上界。

\begin{theorem}[最大纤维饱和的差窗闭式与最大尾和公式]\label{thm:xi-foldbin-maxfiber-saturation-gap-window-closed}
沿用定义 \ref{def:K-of-m}、命题 \ref{prop:fold-bin-digit-dp} 与推论 \ref{cor:fold-bin-saturation} 的记号。对任意 \(m\ge 1\)，令
\[
K:=K(m),
\qquad
L:=K-m,
\qquad
d_{\max}^{\mathrm{bin}}(m):=d_m^{\mathrm{bin}}(0^m).
\]
再定义
\[
S_{\max}(m):=
\max\Bigl\{
\sum_{k=m+1}^{K}c_kF_{k+1}:\ c_k\in\{0,1\},\ c_kc_{k+1}=0
\Bigr\},
\]
并记
\[
\varepsilon_m:=
\begin{cases}
1,& L\equiv 0\pmod 2,\\
0,& L\equiv 1\pmod 2.
\end{cases}
\]
则 \(S_{\max}(m)\) 具有统一闭式
\[
\boxed{\;
S_{\max}(m)
=
\sum_{j=0}^{\lfloor (L-1)/2\rfloor}F_{K+1-2j}
=
F_{K+2}-F_{m+1+\varepsilon_m},
\;}
\]
其中当 \(L=0\) 时上式中的有限和按空和解释为 \(0\)。特别地，
\[
\boxed{\;
d_{\max}^{\mathrm{bin}}(m)=F_{L+2}
\iff
S_{\max}(m)\le 2^m-1
\iff
0\le F_{K+2}-2^m\le F_{m+1+\varepsilon_m}-1.
\;}
\]
\end{theorem}
\begin{proof}
由推论 \ref{cor:fold-bin-saturation}，最大纤维由 \(0^m\) 取得，并且
\[
d_{\max}^{\mathrm{bin}}(m)=F_{L+2}
\iff
S_{\max}(m)\le 2^m-1.
\]
因此只需把 \(S_{\max}(m)\) 写成统一闭式。

当 \(L=0\) 时尾部为空，故 \(S_{\max}(m)=0\)，而此时 \(K=m\) 且
\[
F_{K+2}-F_{m+1+\varepsilon_m}=F_{m+2}-F_{m+2}=0.
\]
以下设 \(L\ge 1\)。

由于权重 \(F_{k+1}\) 随 \(k\) 严格递增，若某个 admissible 尾部在位置 \(K\) 未取 \(1\)，则分两种情形：若 \(c_{K-1}=0\)，直接把 \(c_K\) 改为 \(1\) 即使尾和严格增大；若 \(c_{K-1}=1\)，则把 \(c_{K-1}\) 改为 \(0\) 并把 \(c_K\) 改为 \(1\)，尾和仍然严格增大。故任一极大尾部都必满足 \(c_K=1\)。去掉第 \(K\) 位后，同一论证可递归施加到区间 \(\{m+1,\dots,K-2\}\)，故极大模式只能是
\[
c_K=c_{K-2}=c_{K-4}=\cdots=1,
\]
其余各位皆为 \(0\)。因此
\[
S_{\max}(m)=\sum_{j=0}^{\lfloor (L-1)/2\rfloor}F_{K+1-2j}.
\]

再用恒等式
\[
\sum_{j=0}^{r}F_{n-2j}=F_{n+1}-F_{n-2r-1},
\]
取 \(n=K+1\)、\(r=\lfloor (L-1)/2\rfloor\)，得到
\[
S_{\max}(m)
=
F_{K+2}-F_{K-2\lfloor (L-1)/2\rfloor}.
\]
若 \(L=2s\) 为偶数，则 \(\lfloor (L-1)/2\rfloor=s-1\)，减项为 \(F_{m+2}\)；若 \(L=2s+1\) 为奇数，则减项为 \(F_{m+1}\)。这正好统一为
\[
S_{\max}(m)=F_{K+2}-F_{m+1+\varepsilon_m}.
\]

将该闭式代回饱和判据，便得到
\[
F_{K+2}-F_{m+1+\varepsilon_m}\le 2^m-1
\iff
F_{K+2}-2^m\le F_{m+1+\varepsilon_m}-1.
\]
而由 \(F_{K+1}\le 2^m-1<F_{K+2}\) 可知 \(2^m\le F_{K+2}\)，故左侧差值自动非负。证毕。
\end{proof}

\begin{theorem}[饱和点的指数级 Diophantine 逼近与有限性]\label{thm:xi-foldbin-saturation-diophantine-finiteness}
定义饱和点集合
\[
\mathcal S
:=
\Bigl\{
m\ge 1:\ d_{\max}^{\mathrm{bin}}(m)=F_{K(m)-m+2}
\Bigr\}.
\]
则存在常数 \(C>0\)，使得对每个 \(m\in\mathcal S\) 都有
\[
\boxed{\;
\Bigl|(K(m)+2)\log\varphi-m\log 2-\log\sqrt5\Bigr|
\le
C\Bigl(\frac{\varphi}{2}\Bigr)^m.
\;}
\]
因此 \(\mathcal S\) 只能是有限集；特别地，它具有零自然密度。
\end{theorem}
\begin{proof}
若 \(m\in\mathcal S\)，则由定理 \ref{thm:xi-foldbin-maxfiber-saturation-gap-window-closed}，
\[
0\le F_{K+2}-2^m\le F_{m+1+\varepsilon_m}-1\le F_{m+2}.
\]
令
\[
\widehat\varphi:=-\varphi^{-1},
\qquad
\Delta_m:=(K+2)\log\varphi-m\log 2-\log\sqrt5.
\]
把 Binet 公式
\[
F_n=\frac{\varphi^n-\widehat\varphi^n}{\sqrt5}
\]
代入，得到
\[
\left|
\frac{\varphi^{K+2}}{\sqrt5}-2^m
\right|
\le
F_{m+2}+\frac{|\widehat\varphi|^{K+2}}{\sqrt5}.
\]
由 \(F_{K+1}\le 2^m-1<F_{K+2}\) 可知 \(K=(\log 2/\log\varphi)m+O(1)\)，故
\[
F_{m+2}+\frac{|\widehat\varphi|^{K+2}}{\sqrt5}
=
O(\varphi^m).
\]
两边除以 \(2^m\) 即得
\[
|e^{\Delta_m}-1|
\le
C_1\Bigl(\frac{\varphi}{2}\Bigr)^m.
\]
由于 \(\varphi/2<1\)，当 \(m\) 足够大时右侧小于 \(1/2\)，于是
\[
|\Delta_m|
\le
2|e^{\Delta_m}-1|
\le
2C_1\Bigl(\frac{\varphi}{2}\Bigr)^m.
\]
把有限多个小 \(m\) 吸收到常数中，即得所示不等式。

另一方面，\(\Delta_m\) 是关于代数数 \(2\) 与 \(\varphi\) 的非零对数线性形式；由 Baker--W\"ustholz/Matveev 型下界，它在指数 \(m\) 上只能达到至多多项式级的小量，不可能对无穷多个 \(m\) 满足上述指数级估计。故 \(\mathcal S\) 必为有限集。这正是注 \ref{rem:fold-bin-finite} 所记录的机制。有限集当然具有零自然密度。证毕。
\end{proof}

\begin{theorem}[零温端点微正则包络的极小点与显式尖点闭式]\label{thm:xi-output-microcanonical-logcusp-legendre-closed}
沿用命题 \ref{prop:real-input-40-pressure-freezing} 的记号，定义
\[
h(\alpha):=\inf_{\theta\ge 0}\bigl(P(\theta)-\alpha\theta\bigr),
\qquad
0\le \alpha\le \frac12.
\]
再记
\[
a:=\frac{d}{c}>0,
\qquad
b:=\frac{e}{c}-\frac{d^2}{2c^2}.
\]
则存在 \(\delta_0>0\)，使得当
\[
\delta:=\frac12-\alpha\in(0,\delta_0)
\]
时，\(h(\alpha)\) 在 \([0,\infty)\) 上的极小点唯一，记为 \(\theta(\alpha)\)。并且
\[
\boxed{\;
\theta(\alpha)
=
-2\log\frac{2\delta}{a}
+\frac{8b}{a^2}\delta
+O(\delta^2),
\qquad
\delta\downarrow 0,
\;}
\]
以及
\[
\boxed{\;
h\!\left(\frac12-\delta\right)
=
\log c
-2\delta\log\delta
+2\left(1+\log\frac{a}{2}\right)\delta
+\frac{4b}{a^2}\delta^2
+O(\delta^3).
\;}
\]
\end{theorem}
\begin{proof}
由推论 \ref{cor:real-input-40-alpha-max} 与命题 \ref{prop:real-input-40-output-cumulants-closed}，
\[
P'(0)<\frac12.
\]
因此当 \(\delta>0\) 充分小时，极小点不可能停在边界 \(\theta=0\)，而必须满足驻点方程
\[
P'(\theta)=\frac12-\delta.
\]
再由命题 \ref{prop:real-input-40-pressure-freezing}，
\[
P(\theta)
=
\frac{\theta}{2}+\log c+a e^{-\theta/2}+b e^{-\theta}+O(e^{-3\theta/2}),
\]
从而
\[
P'(\theta)
=
\frac12-\frac{a}{2}e^{-\theta/2}-b e^{-\theta}+O(e^{-3\theta/2}).
\]
令
\[
u:=e^{-\theta/2},
\]
则驻点方程改写为
\[
\delta=\frac{a}{2}u+b u^2+O(u^3).
\]
对该幂级数作反演，得到
\[
u=\frac{2\delta}{a}-\frac{8b}{a^3}\delta^2+O(\delta^3).
\]
取对数并乘以 \(-2\)，便有
\[
\theta(\alpha)
=
-2\log u
=
-2\log\frac{2\delta}{a}
+\frac{8b}{a^2}\delta
+O(\delta^2).
\]
由于同一邻域内
\[
P''(\theta)=\frac{a}{4}e^{-\theta/2}+O(e^{-\theta})>0,
\]
故该驻点唯一，并且就是 \(h(\alpha)\) 的唯一极小点。

第二个展开正是结论 \ref{con:xi-terminal-real-input-output-microcanonical-logcusp-second-order} 的内容；将上式得到的 \(u\)-反演代回
\[
h(\alpha)=P(\theta(\alpha))-\alpha\theta(\alpha)
\]
亦可直接复得
\[
h\!\left(\frac12-\delta\right)
=
\log c
-2\delta\log\delta
+2\left(1+\log\frac{a}{2}\right)\delta
+\frac{4b}{a^2}\delta^2
+O(\delta^3).
\]
证毕。
\end{proof}

\begin{corollary}[零温端点微正则包络的斜率与曲率同时发散]\label{cor:xi-output-microcanonical-endpoint-slope-curvature-diverge}
沿用定理 \ref{thm:xi-output-microcanonical-logcusp-legendre-closed} 的记号。当 \(\alpha\uparrow 1/2\) 时，有
\[
\boxed{\;
h'(\alpha)
=
2\log\!\Bigl(\frac{2(1/2-\alpha)}{a}\Bigr)
+O\!\bigl(1/2-\alpha\bigr)
\longrightarrow -\infty,
\;}
\]
以及
\[
\boxed{\;
h''(\alpha)
=
-\frac{2}{1/2-\alpha}+O(1)
\longrightarrow -\infty.
\;}
\]
因此 \(h\) 在右端点不可能具有有限斜率的 \(C^1\) 延拓，更不存在有限曲率的二阶延拓。
\end{corollary}
\begin{proof}
令 \(\delta:=1/2-\alpha\)。由定理 \ref{thm:xi-output-microcanonical-logcusp-legendre-closed}，
\[
h\!\left(\frac12-\delta\right)
=
\log c
-2\delta\log\delta
+2\left(1+\log\frac{a}{2}\right)\delta
+\frac{4b}{a^2}\delta^2
+O(\delta^3).
\]
对 \(\delta\) 求导得
\[
\frac{d}{d\delta}
h\!\left(\frac12-\delta\right)
=
-2\log\delta
+2\log\frac{a}{2}
+\frac{8b}{a^2}\delta
+O(\delta^2).
\]
而 \(h'(\alpha)=-(d/d\delta)h(1/2-\delta)\)，故
\[
h'(\alpha)
=
2\log\frac{2\delta}{a}+O(\delta),
\]
这就给出第一式。再对上式求导，得到
\[
h''(\alpha)
=
-\frac{2}{\delta}+O(1)
=
-\frac{2}{1/2-\alpha}+O(1),
\]
因而二者都趋于 \(-\infty\)。证毕。
\end{proof}

\begin{theorem}[近地基壳层计数的统一 Legendre 上界]\label{thm:xi-near-ground-shell-counting-legendre-upper-bound}
沿用命题 \ref{prop:real-input-40-pressure-freezing} 与定理 \ref{thm:xi-real-input-40-output-edgeworth-quasipowers} 的记号。记
\[
Z_n(u)=\sum_{k=0}^{n}A_{n,k}u^k,
\qquad
k_n(\alpha):=\lfloor \alpha n\rfloor
\qquad
\bigl(0\le \alpha\le 1/2\bigr).
\]
则对任意固定的 \(\alpha\in[0,1/2]\)，都有
\[
\boxed{\;
\limsup_{n\to\infty}\frac{1}{n}\log A_{n,k_n(\alpha)}
\le
h(\alpha),
\;}
\]
其中 \(h\) 为定理 \ref{thm:xi-output-microcanonical-logcusp-legendre-closed} 中的微正则包络。特别地，当 \(\delta\downarrow 0\) 且 \(\alpha=1/2-\delta\) 时，
\[
\boxed{\;
\limsup_{n\to\infty}\frac{1}{n}\log A_{n,\lfloor(1/2-\delta)n\rfloor}
\le
\log c
-2\delta\log\delta
+2\left(1+\log\frac{a}{2}\right)\delta
+\frac{4b}{a^2}\delta^2
+O(\delta^3).
\;}
\]
\end{theorem}
\begin{proof}
任取 \(\theta\ge 0\)。由
\[
Z_n(e^\theta)=\sum_{k=0}^{n}A_{n,k}e^{\theta k},
\]
可得
\[
A_{n,k_n(\alpha)}e^{\theta k_n(\alpha)}
\le
Z_n(e^\theta).
\]
两边取对数并除以 \(n\)，注意到 \(k_n(\alpha)=\alpha n+O(1)\)，得到
\[
\frac{1}{n}\log A_{n,k_n(\alpha)}
\le
\frac{1}{n}\log Z_n(e^\theta)-\alpha\theta+O\!\left(\frac{1}{n}\right).
\]
对 \(n\to\infty\) 取上极限，并利用 Perron--Frobenius 压力极限
\[
\lim_{n\to\infty}\frac{1}{n}\log Z_n(e^\theta)=P(\theta),
\]
便有
\[
\limsup_{n\to\infty}\frac{1}{n}\log A_{n,k_n(\alpha)}
\le
P(\theta)-\alpha\theta.
\]
再对全部 \(\theta\ge 0\) 取下确界，便得到
\[
\limsup_{n\to\infty}\frac{1}{n}\log A_{n,k_n(\alpha)}
\le
\inf_{\theta\ge 0}\bigl(P(\theta)-\alpha\theta\bigr)
=
h(\alpha).
\]
当 \(\alpha=1/2-\delta\) 且 \(\delta\downarrow 0\) 时，只需把定理
\ref{thm:xi-output-microcanonical-logcusp-legendre-closed} 的端点展开代入即可。证毕。
\end{proof}

\noindent
综上，\(\log c\) 同时控制极端饱和层的最小完备外置率、无偏采样等待代价与微正则端点值；\(\Delta_\mu\) 继续刻画最大密度层与平衡均值之间的线性代数距离，而不消灭其正的指数熵；\(u_R\) 与 \(c^2=\rho(B_\infty)\) 之间的谱缝则说明 RH/Ramanujan 失效严格先于零温四次 Perron 尺度。进一步，bin-fold 最大纤维的离散异常饱和被压缩为单个 Fibonacci 差窗与指数级 Diophantine 逼近，而输出位势的零温 Puiseux 展开又把这一离散窗口传递为微正则端点的对数尖点与近地基壳层的统一上界。因而本段的零温结构可概括为
\[
\text{最大纤维饱和判据}
\Longrightarrow
\text{指数级 Diophantine 逼近}
\Longrightarrow
\text{零温端点 Puiseux 几何}
\Longrightarrow
\text{微正则对数尖点}
\Longrightarrow
\text{近地基壳层指数上界}.
\]
其中，从极端饱和层到其线性邻壳，主导几何始终由同一 Puiseux--Legendre 机制锁定。

% (User input window)

\paragraph{离散预算容量指数、基态驱动界与 window-\texorpdfstring{$6$}{6} 异常分离}
前一组结果把离散比特预算下的预言机容量压缩为尾谱变分公式与由最大纤维数据驱动的粗相图；
后一组结果则把 window-\(6\) 的奇偶荷格、\(B_3/C_3\) 根字典、中心结构与精确可逆阈值写成同一组
可核验的离散分级结论。

\begin{theorem}[预言机容量指数的尾谱变分界]\label{thm:xi-oracle-capacity-exponent-tail-variational-bounds}
\leanverified{paper\_xi\_oracle\_capacity\_exponent\_tail\_variational\_bounds}
设折叠族满足
\[
\Fold_m:\Omega_m\to X_m,
\qquad
|X_m|=F_{m+2},
\]
并记纤维多重度
\[
d_m(x):=\bigl|\Fold_m^{-1}(x)\bigr|
\qquad(x\in X_m).
\]
对整数预算 \(B\ge 0\)，定义预言机容量
\[
C_m(B):=\sum_{x\in X_m}\min\!\bigl(d_m(x),2^B\bigr).
\]
对 \(\lambda\ge 0\)，定义指数尾计数
\[
N_m(\lambda):=\#\{x\in X_m:\ d_m(x)\ge e^{\lambda m}\},
\]
并记
\[
s^{-}(\lambda):=\liminf_{m\to\infty}\frac1m\log N_m(\lambda),
\qquad
s^{+}(\lambda):=\limsup_{m\to\infty}\frac1m\log N_m(\lambda).
\]
对任意 \(\beta\ge 0\)，令
\[
B_m:=\left\lfloor \frac{\beta m}{\log 2}\right\rfloor,
\]
并定义容量指数
\[
\underline\Sigma(\beta):=\liminf_{m\to\infty}\frac1m\log C_m(B_m),
\qquad
\overline\Sigma(\beta):=\limsup_{m\to\infty}\frac1m\log C_m(B_m).
\]
则有严格变分界
\[
\boxed{\;
\max\!\left\{\log\varphi,\ \sup_{0\le \lambda<\beta}\bigl(\lambda+s^{-}(\lambda)\bigr)\right\}
\le
\underline\Sigma(\beta)
\le
\overline\Sigma(\beta)
\le
\max\!\left\{\log\varphi,\ \sup_{0\le \lambda\le \beta}\bigl(\lambda+s^{+}(\lambda)\bigr)\right\}. \;}
\]
\end{theorem}
\begin{proof}
当 \(\beta=0\) 时，\(B_m=0\) 且
\[
C_m(B_m)=C_m(0)=\sum_{x\in X_m}1=|X_m|.
\]
于是
\[
\underline\Sigma(0)=\overline\Sigma(0)
=
\lim_{m\to\infty}\frac1m\log|X_m|
=
\log\varphi,
\]
结论成立。以下设 \(\beta>0\)。
\par
由
\(
B_m=\lfloor \beta m/\log 2\rfloor
\)
得
\[
2^{B_m}=e^{\beta m+O(1)}.
\]
\par
\textbf{下界。}
固定 \(0\le \lambda<\beta\)。
当 \(m\) 充分大时，
\(
2^{B_m}\ge e^{\lambda m}
\)。
故若 \(x\in X_m\) 满足 \(d_m(x)\ge e^{\lambda m}\)，则
\[
\min\bigl(d_m(x),2^{B_m}\bigr)\ge e^{\lambda m}.
\]
于是
\[
C_m(B_m)\ge e^{\lambda m}N_m(\lambda).
\]
取 \(\liminf\) 得
\[
\underline\Sigma(\beta)\ge \lambda+s^{-}(\lambda).
\]
再对全部 \(0\le \lambda<\beta\) 取上确界，得到
\[
\underline\Sigma(\beta)\ge \sup_{0\le \lambda<\beta}\bigl(\lambda+s^{-}(\lambda)\bigr).
\]
另一方面，对每个 \(x\in X_m\) 都有
\(
\min(d_m(x),2^{B_m})\ge 1
\)，故
\[
C_m(B_m)\ge |X_m|.
\]
于是
\[
\underline\Sigma(\beta)\ge
\lim_{m\to\infty}\frac1m\log|X_m|
=\log\varphi.
\]
二者合并即得左端不等式。
\par
\textbf{上界。}
固定 \(\varepsilon\in(0,\beta)\)，并令
\[
J:=\left\lfloor\frac{\beta}{\varepsilon}\right\rfloor.
\]
由于 \(J\varepsilon\le \beta<(J+1)\varepsilon\)，对任意 \(x\in X_m\) 有三种可能：
\begin{enumerate}
  \item 若 \(d_m(x)<e^{\varepsilon m}\)，则
  \(
  \min(d_m(x),2^{B_m})\le e^{\varepsilon m}
  \)；
  \item 若对某个 \(1\le j\le J-1\) 有
  \(
  e^{j\varepsilon m}\le d_m(x)<e^{(j+1)\varepsilon m}
  \)，则
  \(
  \min(d_m(x),2^{B_m})\le e^{(j+1)\varepsilon m}
  \)；
  \item 若 \(d_m(x)\ge e^{J\varepsilon m}\)，则
  \(
  \min(d_m(x),2^{B_m})\le 2^{B_m}=e^{\beta m+O(1)}
  \)。
\end{enumerate}
因此
\[
\begin{aligned}
C_m(B_m)
&\le
e^{\varepsilon m}|X_m|
+\sum_{j=1}^{J-1} e^{(j+1)\varepsilon m}N_m(j\varepsilon)
+e^{\beta m+O(1)}N_m(J\varepsilon).
\end{aligned}
\]
取 \(\limsup\) 并用有限和的指数主导原则，得到
\[
\overline\Sigma(\beta)
\le
\max\!\left\{
\log\varphi+\varepsilon,\
\max_{1\le j\le J-1}\bigl((j+1)\varepsilon+s^{+}(j\varepsilon)\bigr),\
\beta+s^{+}(J\varepsilon)
\right\}.
\]
若 \(J=1\)，则中间最大值项省略。
对中间各项，
\[
(j+1)\varepsilon+s^{+}(j\varepsilon)
=
\varepsilon+\bigl(j\varepsilon+s^{+}(j\varepsilon)\bigr)
\le
\varepsilon+\sup_{0\le \lambda\le \beta}\bigl(\lambda+s^{+}(\lambda)\bigr).
\]
而对最后一项，由 \(0\le \beta-J\varepsilon<\varepsilon\) 得
\[
\beta+s^{+}(J\varepsilon)
\le
\varepsilon+\bigl(J\varepsilon+s^{+}(J\varepsilon)\bigr)
\le
\varepsilon+\sup_{0\le \lambda\le \beta}\bigl(\lambda+s^{+}(\lambda)\bigr).
\]
故
\[
\overline\Sigma(\beta)
\le
\max\!\left\{
\log\varphi+\varepsilon,\
\varepsilon+\sup_{0\le \lambda\le \beta}\bigl(\lambda+s^{+}(\lambda)\bigr)
\right\}.
\]
令 \(\varepsilon\downarrow 0\)，即得
\[
\overline\Sigma(\beta)
\le
\max\!\left\{
\log\varphi,\
\sup_{0\le \lambda\le \beta}\bigl(\lambda+s^{+}(\lambda)\bigr)
\right\}.
\]
证毕。
\end{proof}

\begin{corollary}[若尾谱存在，则预言机容量指数具有精确闭式]\label{cor:xi-oracle-capacity-exponent-closed-form-tail-spectrum}
沿用定理 \ref{thm:xi-oracle-capacity-exponent-tail-variational-bounds} 的记号。
进一步假设对每个 \(\lambda\ge 0\)，极限
\[
s(\lambda):=\lim_{m\to\infty}\frac1m\log N_m(\lambda)
\]
存在，且 \(s\) 上半连续。
则极限
\[
\Sigma(\beta):=\lim_{m\to\infty}\frac1m\log C_m(B_m)
\]
存在，并满足
\[
\boxed{\;
\Sigma(\beta)=
\max\!\left\{\log\varphi,\ \sup_{0\le \lambda\le \beta}\bigl(\lambda+s(\lambda)\bigr)\right\}. \;}
\]
特别地，若定义完美恢复的临界比特率
\[
\beta_c:=\inf\{\beta\ge 0:\ \Sigma(\beta)=\log 2\},
\]
则
\[
\boxed{\;
\beta_c=
\inf\!\left\{
\beta\ge 0:\ 
\sup_{0\le \lambda\le \beta}\bigl(\lambda+s(\lambda)\bigr)=\log 2
\right\}. \;}
\]
\end{corollary}
\begin{proof}
由定理 \ref{thm:xi-oracle-capacity-exponent-tail-variational-bounds}，
\[
\max\!\left\{\log\varphi,\ \sup_{0\le \lambda<\beta}\bigl(\lambda+s(\lambda)\bigr)\right\}
\le
\underline\Sigma(\beta)
\le
\overline\Sigma(\beta)
\le
\max\!\left\{\log\varphi,\ \sup_{0\le \lambda\le \beta}\bigl(\lambda+s(\lambda)\bigr)\right\}.
\]
当 \(\beta=0\) 时，由 \(N_m(0)=|X_m|=F_{m+2}\) 得
\[
s(0)=\lim_{m\to\infty}\frac1m\log N_m(0)=\log\varphi,
\]
故上、下界同为 \(\log\varphi\)，从而
\[
\Sigma(0)=\log\varphi.
\]
以下设 \(\beta>0\)。
又由于每个 \(N_m(\lambda)\) 关于 \(\lambda\) 单调不增，故其指数极限 \(s(\lambda)\) 亦单调不增。
于是对任意 \(\beta>0\) 与任意 \(\lambda<\beta\) 充分接近 \(\beta\)，都有
\[
\lambda+s(\lambda)\ge \lambda+s(\beta).
\]
令 \(\lambda\uparrow \beta\) 得
\[
\sup_{0\le \lambda<\beta}\bigl(\lambda+s(\lambda)\bigr)\ge \beta+s(\beta),
\]
从而
\[
\sup_{0\le \lambda<\beta}\bigl(\lambda+s(\lambda)\bigr)
=
\sup_{0\le \lambda\le \beta}\bigl(\lambda+s(\lambda)\bigr).
\]
故上下界一致，极限 \(\Sigma(\beta)\) 存在并等于所示闭式。
\par
此外，由平凡上界
\[
C_m(B_m)\le \sum_{x\in X_m}d_m(x)=2^m
\]
可知
\[
\Sigma(\beta)\le \log 2.
\]
又 \(\log\varphi<\log 2\)，因此
\[
\Sigma(\beta)=\log 2
\quad\Longleftrightarrow\quad
\sup_{0\le \lambda\le \beta}\bigl(\lambda+s(\lambda)\bigr)=\log 2,
\]
临界比特率公式随即由定义得到。
最后，\(s\) 的上半连续性保证上述上确界在紧区间 \([0,\beta]\) 上可由某个极值点达到。证毕。
\end{proof}

\paragraph{基态驱动的离散容量粗界与有限窗极端分离}
在尾谱极限完全建立之前，仍可仅凭最大纤维层的两项指数 \((\alpha_*,g_*)\) 给出离散预言机容量的一条双侧闭式界；而在 \(m=6\) 的有限窗上，该粗界退化为完全可枚举的精确阈值。

\begin{theorem}[线性比特率下离散预言机容量的基态驱动双侧界]\label{thm:xi-oracle-capacity-exponent-groundstate-two-sided-bounds}
沿用定理 \ref{thm:xi-oracle-capacity-exponent-tail-variational-bounds} 的折叠记号。
记
\[
M_m:=\max_{x\in X_m}d_m(x),\qquad
K_m:=\#\{x\in X_m:\ d_m(x)=M_m\},
\]
\[
A_m^{\max}:=\{x\in X_m:\ d_m(x)=M_m\},
\qquad
\alpha_*:=\lim_{m\to\infty}\frac{1}{m}\log M_m,\qquad
g_*:=\lim_{m\to\infty}\frac{1}{m}\log K_m.
\]
对 \(\beta\ge 0\) 令
\[
B_m:=\lfloor \beta m\rfloor,
\]
并定义离散容量指数
\[
\Gamma_{\mathrm{disc}}(\beta):=\limsup_{m\to\infty}\frac1m\log C_m(B_m),
\qquad
\mathrm{Succ}_m(B):=\frac{C_m(B)}{2^m}.
\]
则有普适双侧界
\[
\boxed{\;
g_*+\min\{\alpha_*,\beta\log 2\}
\le
\Gamma_{\mathrm{disc}}(\beta)
\le
\min\{\log 2,\ \log\varphi+\beta\log 2\}.\;}
\]
此外，若
\[
\beta>\frac{\alpha_*}{\log 2},
\]
则对充分大的 \(m\) 有 \(2^{B_m}\ge M_m\)，从而
\[
C_m(B_m)=2^m,\qquad \mathrm{Succ}_m(B_m)=1.
\]
\end{theorem}
\begin{proof}
只取最大纤维集合 \(A_m^{\max}\) 即得
\[
C_m(B_m)\ge K_m\min(M_m,2^{B_m}).
\]
对两端取对数并除以 \(m\)，再取上极限，得到
\[
\Gamma_{\mathrm{disc}}(\beta)\ge g_*+\min\{\alpha_*,\beta\log 2\}.
\]

另一方面，对每个 \(x\in X_m\) 都有
\[
\min(d_m(x),2^{B_m})\le 2^{B_m},
\]
故
\[
C_m(B_m)\le |X_m|\,2^{B_m}.
\]
由 \(|X_m|=F_{m+2}\) 得
\[
\Gamma_{\mathrm{disc}}(\beta)\le \log\varphi+\beta\log 2.
\]
再由
\[
\min(d_m(x),2^{B_m})\le d_m(x)
\]
与
\[
\sum_{x\in X_m}d_m(x)=2^m
\]
得到
\[
C_m(B_m)\le 2^m,
\]
故 \(\Gamma_{\mathrm{disc}}(\beta)\le \log 2\)。
两式合并即得上界。

最后，若 \(\beta>\alpha_*/\log 2\)，则
\[
2^{B_m}=\exp(\beta m\log 2+O(1))
\]
而
\[
M_m=\exp(\alpha_*m+o(m)),
\]
故 \(2^{B_m}\ge M_m\) 对充分大 \(m\) 成立。
于是
\[
C_m(B_m)=\sum_{x\in X_m}d_m(x)=2^m,
\]
从而 \(\mathrm{Succ}_m(B_m)=1\)。
证毕。
\end{proof}

\begin{proposition}[window-\texorpdfstring{$6$}{6} 中精确可逆阈值与完备异常签名维数的极端分离]\label{prop:xi-window6-two-bit-recovery-vs-21bit-signature}
由表 \ref{tab:fold_bin_gauge_m6_invariants}，window-\(6\) 的纤维直方图为
\[
2:8,\qquad 3:4,\qquad 4:9,
\]
且由定理 \ref{thm:xi-foldbin-gauge-representation-distribution-law-m6-spectrum} 可知
\[
\Aut(\Fold_6^{\mathrm{bin}})^{\mathrm{ab}}\cong (\ZZ/2\ZZ)^{21}.
\]
定义最小精确反演比特预算
\[
B_{\mathrm{inv}}^{\min}(6):=\min\{B\in\ZZ_{\ge 0}:\ \mathrm{Succ}_6(B)=1\}.
\]
则
\[
C_6(0)=21,\qquad
C_6(1)=42,\qquad
C_6(2)=64=2^6,
\]
从而
\[
\mathrm{Succ}_6(0)=\frac{21}{64},\qquad
\mathrm{Succ}_6(1)=\frac{21}{32},\qquad
\mathrm{Succ}_6(2)=1.
\]
特别地，
\[
\boxed{\;
B_{\mathrm{inv}}^{\min}(6)=2,\qquad
\Aut(\Fold_6^{\mathrm{bin}})^{\mathrm{ab}}\cong (\ZZ/2\ZZ)^{21}.\;}
\]
等价地，window-\(6\) 的精确微态恢复仅需 \(2\) 比特侧信息，而最小完备二元异常签名需要 \(21\) 个独立坐标。
\end{proposition}
\begin{proof}
由纤维直方图与预言机容量的闭式
\[
C_6(B)=\sum_{x\in X_6}\min(d_6(x),2^B)
=8\min(2,2^B)+4\min(3,2^B)+9\min(4,2^B),
\]
直接得到
\[
C_6(0)=8+4+9=21,
\qquad
C_6(1)=8\cdot 2+4\cdot 2+9\cdot 2=42,
\]
\[
C_6(2)=8\cdot 2+4\cdot 3+9\cdot 4=64=2^6.
\]
故相应成功率如上所示。

又由于 \(\max_{x\in X_6}d_6(x)=4\)，当 \(B=1\) 时仍有纤维无法被完全分离，亦即
\[
\mathrm{Succ}_6(1)=\frac{42}{64}<1,
\]
而当 \(B=2\) 时已有 \(2^B=4\ge \max_x d_6(x)\)，从而每条纤维均可精确译码。
因此 \(B_{\mathrm{inv}}^{\min}(6)=2\)。

另一方面，定理 \ref{thm:xi-foldbin-gauge-representation-distribution-law-m6-spectrum} 给出
\[
\Aut(\Fold_6^{\mathrm{bin}})^{\mathrm{ab}}\cong (\ZZ/2\ZZ)^{21},
\]
故最小完备二元异常签名空间的维数恰为 \(21\)。
证毕。
\end{proof}
\leanverified{paper\_xi\_window6\_two\_bit\_recovery\_vs\_21bit\_signature}

\begin{theorem}[window-\texorpdfstring{$6$}{6} 的恢复资源与异常资源形成严格的二维分叉]\label{thm:xi-window6-oracle-anomaly-coordinate-bifurcation}
继续沿用命题 \ref{prop:xi-window6-two-bit-recovery-vs-21bit-signature} 的记号。称预算对
\[
(B_{\mathrm{oracle}},B_{\mathrm{anom}})\in \ZZ_{\ge 0}^2
\]
为完备的，若同时满足下列两条语义约束：
\begin{enumerate}
  \item 在 oracle 语义下，解码不得跨纤维，且完美恢复必须由各纤维上的侧信息单射实现；
  \item 在 fold-gauge 语义下，异常签名仍由纤维独立置换的直积结构
  \[
  \Aut(\Fold_6^{\mathrm{bin}})
  \cong
  \prod_{w\in X_6}S_{d_6(w)}
  \]
  所承载，故完备二元异常签名的坐标数等于其可交换化秩。
\end{enumerate}
则在乘积偏序
\(
(a,b)\preceq(a',b')\iff a\le a',\ b\le b'
\)
下，window-\(6\) 的唯一坐标极小完备预算对为
\[
\boxed{\;
\bigl(B_{\mathrm{oracle}}^{\min},B_{\mathrm{anom}}^{\min}\bigr)
=
(2,21).\;}
\]
等价地：
\begin{enumerate}
  \item 完美 oracle 反演需要且只需要 \(2\) 比特侧信息；
  \item 完备二元异常签名需要且只需要 \(21\) 个独立坐标；
  \item 任意试图在坐标意义上压到 \(B_{\mathrm{oracle}}<2\) 或 \(B_{\mathrm{anom}}<21\) 的方案，必然分别破坏上述 oracle 语义或 fold-gauge 语义。
\end{enumerate}
\end{theorem}
\begin{proof}
由命题 \ref{prop:xi-window6-two-bit-recovery-vs-21bit-signature}，
\[
B_{\mathrm{oracle}}^{\min}=B_{\mathrm{inv}}^{\min}(6)=2,
\]
并且
\[
\Aut(\Fold_6^{\mathrm{bin}})^{\mathrm{ab}}\cong (\ZZ/2\ZZ)^{21}.
\]
因此，在保持纤维内单射恢复机制不变的前提下，凡 \(B_{\mathrm{oracle}}<2\) 都不可能实现全微态恢复；这给出了 oracle 坐标的最小值。

另一方面，可交换化群
\(
\Aut(\Fold_6^{\mathrm{bin}})^{\mathrm{ab}}
\)
作为完备二元异常签名的普适承载对象，恰需要 \(21\) 个 \(\FF_2\)-坐标才能忠实表示。因而任何 \(B_{\mathrm{anom}}<21\) 的方案若仍欲完备承载全部异常信息，就必然引入跨纤维关系，从而破坏纤维独立置换的直积语义。

故 \((2,21)\) 同时可达，且在乘积偏序下不能再向任一坐标方向下降，因而它是唯一的坐标极小完备预算对。证毕。
\end{proof}

\begin{theorem}[window-\texorpdfstring{$6$}{6} 奇偶荷格在 \texorpdfstring{$B_3$}{B3} 字典下的三分级]\label{thm:xi-window6-parity-charge-trigrading-b3}
记
\[
G_6^{\mathrm{bin}}:=\Aut(\Fold_6^{\mathrm{bin}})=\mathcal G_6^{\mathrm{bin}},
\]
并沿用 \(X_6=X_6^{\mathrm{cyc}}\sqcup X_6^{\mathrm{bdry}}\) 的刚性分裂。
固定表 \ref{tab:fold6_b3c3_root_dictionary_b3} 的 \(B_3\) 根字典
\[
\iota_B:X_6^{\mathrm{cyc}}\xrightarrow{\sim}\Phi(B_3).
\]
定义
\[
X_{6,\mathrm{short}}
:=
\iota_B^{-1}\bigl(\{\pm e_1,\pm e_2,\pm e_3\}\bigr),
\]
\[
X_{6,\mathrm{long}}
:=
\iota_B^{-1}\bigl(\{\pm e_i\pm e_j:\ 1\le i<j\le 3\}\bigr).
\]
则
\[
X_6=X_6^{\mathrm{bdry}}\sqcup X_{6,\mathrm{short}}\sqcup X_{6,\mathrm{long}},
\qquad
|X_6^{\mathrm{bdry}}|=3,\quad |X_{6,\mathrm{short}}|=6,\quad |X_{6,\mathrm{long}}|=12,
\]
并且经由符号映射给出的坐标识别
\[
\bigl(G_6^{\mathrm{bin}}\bigr)^{\mathrm{ab}}\cong (\ZZ/2\ZZ)^{X_6}
\]
下，存在规范直和分解
\[
\boxed{\;
\bigl(G_6^{\mathrm{bin}}\bigr)^{\mathrm{ab}}
\cong
(\ZZ/2\ZZ)^{X_6^{\mathrm{bdry}}}
\oplus
(\ZZ/2\ZZ)^{X_{6,\mathrm{short}}}
\oplus
(\ZZ/2\ZZ)^{X_{6,\mathrm{long}}}
\cong
(\ZZ/2\ZZ)^3_{\mathrm{bdry}}
\oplus
(\ZZ/2\ZZ)^6_{\mathrm{short}}
\oplus
(\ZZ/2\ZZ)^{12}_{\mathrm{long}}.\;}
\]
若改用表 \ref{tab:fold6_b3c3_root_dictionary_c3} 的 \(C_3\) 规范，则后两项仅交换“长/短”的名称，而各直和因子的维数保持不变。
\end{theorem}
\begin{proof}
由推论 \ref{cor:fold-bin-m6-fiber-hist}，window-\(6\) 的全部纤维大小只取 \(2,3,4\)，故每个
\(w\in X_6\) 都贡献一个非平凡的 \(\ZZ/2\ZZ\) 符号坐标。
再由定理 \ref{thm:xi-foldbin-gauge-representation-distribution-law-m6-spectrum}，
\[
\bigl(G_6^{\mathrm{bin}}\bigr)^{\mathrm{ab}}\cong (\ZZ/2\ZZ)^{21},
\]
并且该阿贝尔化可经由纤维符号映射识别为 \((\ZZ/2\ZZ)^{X_6}\) 的坐标格。
\par
另一方面，表 \ref{tab:fold6_b3c3_root_dictionary_b3} 给出一个确定的 \(B_3\) 根字典。
其中
\[
\{\pm e_1,\pm e_2,\pm e_3\}
\]
恰对应六个轴根，故其原像集合 \(X_{6,\mathrm{short}}\) 具有 \(6\) 个元素；
而
\[
\{\pm e_i\pm e_j:\ 1\le i<j\le 3\}
\]
恰为 \(B_3\) 的 \(12\) 个长根，故其原像集合 \(X_{6,\mathrm{long}}\) 具有 \(12\) 个元素。
再与 boundary 集合 \(X_6^{\mathrm{bdry}}\)（大小 \(3\)）合并，即得
\[
X_6=X_6^{\mathrm{bdry}}\sqcup X_{6,\mathrm{short}}\sqcup X_{6,\mathrm{long}}.
\]
对坐标格按此分块，便得到所示直和分解。
\par
若改用 \(C_3\) 字典，则表 \ref{tab:fold6_b3c3_root_dictionary_c3} 仅把 \(B_3\) 规范中的长根与短根名称交换，而根集总数与 boundary 三维子格不变，故维数结论保持不变。证毕。
\end{proof}

\begin{theorem}[window-\texorpdfstring{$6$}{6} 的纤维多重度场在 \texorpdfstring{$B_3$}{B3} 根系上严格破坏 Weyl 轨道均匀性]\label{thm:xi-window6-fiber-multiplicity-weyl-breaking-b3}
沿用定理 \ref{thm:xi-window6-parity-charge-trigrading-b3} 的记号。
将纤维多重度函数
\[
d_6:X_6\to \{2,3,4\}
\]
经 \(B_3\) 字典传送到根系上，定义
\[
\widetilde d_6:\Phi(B_3)\to \{2,3,4\},
\qquad
\widetilde d_6(\alpha):=d_6\bigl(\iota_B^{-1}(\alpha)\bigr).
\]
则 \(\widetilde d_6\) 在 \(W(B_3)\) 的两个根轨道上都不是常数。更精确地：
\begin{enumerate}
  \item 在短根轨道上，数值 \(2\) 与 \(4\) 同时出现；
  \item 在长根轨道上，数值 \(2,3,4\) 三者全部出现。
\end{enumerate}
因此 \(\widetilde d_6\) 既不是 \(W(B_3)\)-不变函数，也不是仅依赖根长的函数。
\end{theorem}
\begin{proof}
由表 \ref{tab:fold6_b3c3_root_dictionary_b3}，
\[
\iota_B(\texttt{100000})=(1,0,0),
\qquad
\iota_B(\texttt{000001})=(-1,0,0).
\]
这两个根同属 \(B_3\) 的短根 Weyl 轨道。
而由表 \ref{tab:fold_bin_gauge_m6_fibers}，
\[
d_6(\texttt{100000})=4,
\qquad
d_6(\texttt{000001})=2.
\]
故 \(\widetilde d_6\) 在短根轨道上不是常数。
\par
再由表 \ref{tab:fold6_b3c3_root_dictionary_b3}，
\[
\iota_B(\texttt{101000})=(1,-1,0),\qquad
\iota_B(\texttt{100010})=(1,0,1),\qquad
\iota_B(\texttt{010001})=(0,1,1).
\]
这三者都是 \(B_3\) 的长根，且同属唯一的长根 Weyl 轨道。
由表 \ref{tab:fold_bin_gauge_m6_fibers}，
\[
d_6(\texttt{101000})=4,\qquad
d_6(\texttt{100010})=3,\qquad
d_6(\texttt{010001})=2.
\]
故在长根轨道上，\(\widetilde d_6\) 取到全部三种数值 \(2,3,4\)。
\par
综上，\(\widetilde d_6\) 不可能在任一 Weyl 根轨道上常值，特别地不可能只依赖“根长”这一粗分层。证毕。
\end{proof}

\begin{theorem}[window-\texorpdfstring{$6$}{6} 规范群中心的精确结构与 boundary 中心的余维]\label{thm:xi-window6-center-exact-splitting-bdry-cyclic}
\leanverified{paper\_xi\_window6\_center\_exact\_splitting\_bdry\_cyclic}
沿用定理 \ref{thm:xi-foldbin-gauge-representation-distribution-law-m6-spectrum} 的记号。
令
\[
d_6(w):=\bigl|\Fold_6^{\mathrm{bin}}{}^{-1}(w)\bigr|
\qquad(w\in X_6),
\]
\[
X_6^{(2)}:=\{w\in X_6:\ d_6(w)=2\},
\qquad
X_{6,\mathrm{cyc},2}:=X_6^{(2)}\setminus X_6^{\mathrm{bdry}}.
\]
则
\[
X_6^{(2)}
:=
\{\texttt{000001},\texttt{000101},\texttt{001001},\texttt{010001},\texttt{010101},
\texttt{100001},\texttt{100101},\texttt{101001}\},
\]
\[
X_{6,\mathrm{cyc},2}
:=
\{\texttt{000001},\texttt{000101},\texttt{001001},\texttt{010001},\texttt{010101}\}.
\]
并且规范群中心满足精确分解
\[
\boxed{\;
Z\!\bigl(G_6^{\mathrm{bin}}\bigr)
\cong
(\ZZ/2\ZZ)^{X_6^{(2)}}
\cong
(\ZZ/2\ZZ)^{X_6^{\mathrm{bdry}}}
\oplus
(\ZZ/2\ZZ)^{X_{6,\mathrm{cyc},2}}
\cong
(\ZZ/2\ZZ)^3_{\mathrm{bdry}}
\oplus
(\ZZ/2\ZZ)^5_{\mathrm{cyc},\,d=2}
\cong
(\ZZ/2\ZZ)^8.\;}
\]
特别地，boundary 中心子群 \((\ZZ/2\ZZ)^3_{\mathrm{bdry}}\) 是整个中心的一个规范直和因子，
其在 \(Z(G_6^{\mathrm{bin}})\) 中的余维恰为 \(5\)。
\end{theorem}
\begin{proof}
由定理 \ref{thm:xi-foldbin-gauge-representation-distribution-law-m6-spectrum}，
\[
G_6^{\mathrm{bin}}\cong \prod_{w\in X_6}\mathfrak S_{d_6(w)}.
\]
对有限直积群，中心满足
\[
Z\!\left(\prod_{w\in X_6}\mathfrak S_{d_6(w)}\right)
\cong
\prod_{w\in X_6} Z\!\left(\mathfrak S_{d_6(w)}\right).
\]
而
\[
Z(\mathfrak S_d)\cong
\begin{cases}
\ZZ/2\ZZ,& d=2,\\
0,& d\ge 3,
\end{cases}
\]
因为 \(\mathfrak S_2\) 本身阿贝尔，而 \(\mathfrak S_3,\mathfrak S_4\) 的中心平凡。
故
\[
Z(G_6^{\mathrm{bin}})\cong (\ZZ/2\ZZ)^{X_6^{(2)}}.
\]
\par
再由表 \ref{tab:fold_bin_gauge_m6_fibers} 可直接读出恰有八条 \(d_6=2\) 纤维，
其词列表正是上式所列的 \(X_6^{(2)}\)。
其中后三条
\[
X_6^{\mathrm{bdry}}=\{\texttt{100001},\texttt{100101},\texttt{101001}\}
\]
正是 boundary 扇区；前五条构成 cyclic 的 \(d=2\) 子集 \(X_{6,\mathrm{cyc},2}\)。
于是
\[
X_6^{(2)}=X_6^{\mathrm{bdry}}\sqcup X_{6,\mathrm{cyc},2},
\]
从而中心按坐标分裂为
\[
(\ZZ/2\ZZ)^{X_6^{(2)}}
\cong
(\ZZ/2\ZZ)^{X_6^{\mathrm{bdry}}}
\oplus
(\ZZ/2\ZZ)^{X_{6,\mathrm{cyc},2}}
\cong
(\ZZ/2\ZZ)^3\oplus (\ZZ/2\ZZ)^5.
\]
这同时给出总中心维数 \(8\) 以及 boundary 中心子群的余维 \(5\)。
\par
最后，推论 \ref{cor:fold-bin-boundary-center-z2} 已把
\((\ZZ/2\ZZ)^{X_6^{\mathrm{bdry}}}\)
识别为 boundary sheet-flip 所生成的规范中心子格；故这里得到的是对该结果的完整中心精化。证毕。
\end{proof}
\begin{theorem}[window-\texorpdfstring{$6$}{6} 的 cyclic sector 在 \texorpdfstring{$C_3$}{C3} 规范下出现严格的 \texorpdfstring{$9:4:5$}{9:4:5} 纤维三分]\label{thm:xi-window6-c3-cyclic-fiber-trisection}
\leanverified{paper\_xi\_window6\_c3\_cyclic\_fiber\_trisection}
沿用表 \ref{tab:fold6_b3c3_root_dictionary_c3} 的 \(C_3\) 根字典
\[
\iota_C:X_6^{\mathrm{cyc}}\xrightarrow{\sim}\Phi(C_3)
\]
与 window-\(6\) 的纤维多重度函数 \(d_6^{\mathrm{bin}}\)。
定义
\[
\mathcal R_j:=\{r\in X_6^{\mathrm{cyc}}:\ d_6^{\mathrm{bin}}(r)=j\}
\qquad(j=2,3,4).
\]
则在 \(C_3\) 规范下有精确描述
\[
\boxed{\;
\iota_C(\mathcal R_4)
=
\{2e_1,\pm 2e_2,\pm 2e_3,\pm e_1\pm e_2\},\;}
\]
\[
\boxed{\;
\iota_C(\mathcal R_3)
=
\{\pm e_1\pm e_3\},\;}
\]
\[
\boxed{\;
\iota_C(\mathcal R_2)
=
\{-2e_1,\pm e_2\pm e_3\}.\;}
\]
特别地，
\[
\boxed{\;
|\mathcal R_4|=9,\qquad |\mathcal R_3|=4,\qquad |\mathcal R_2|=5.\;}
\]
\end{theorem}
\begin{proof}
由表 \ref{tab:fold6_b3c3_root_dictionary_c3}，\(X_6^{\mathrm{cyc}}\) 的 \(18\) 个词与 \(C_3\) 根系逐项对应。另一方面，由 window-\(6\) 的纤维判别与直方图
\[
d_6^{\mathrm{bin}}(w)\in\{2,3,4\},
\qquad
(2:8,3:4,4:9),
\]
可直接按词逐项核对。

首先，末位 \(w_6=1\) 的五个 cyclic 词为
\[
\texttt{000001},\ \texttt{010001},\ \texttt{001001},\ \texttt{000101},\ \texttt{010101},
\]
其在表 \ref{tab:fold6_b3c3_root_dictionary_c3} 下分别对应
\[
-2e_1,\ e_2+e_3,\ e_2-e_3,\ -e_2+e_3,\ -e_2-e_3.
\]
由 \(m=6\) 的精确纤维判别，这五个词恰全都满足 \(d_6^{\mathrm{bin}}=2\)，故
\[
\iota_C(\mathcal R_2)=\{-2e_1,\pm e_2\pm e_3\}.
\]

其次，满足 \(w_6=0\) 且落在 \(d_6^{\mathrm{bin}}=3\) 层的四个 cyclic 词为
\[
\texttt{100010},\ \texttt{010010},\ \texttt{001010},\ \texttt{101010},
\]
其对应向量分别为
\[
e_1+e_3,\ e_1-e_3,\ -e_1+e_3,\ -e_1-e_3.
\]
故
\[
\iota_C(\mathcal R_3)=\{\pm e_1\pm e_3\}.
\]

最后，其余九个 cyclic 词自动落在 \(d_6^{\mathrm{bin}}=4\) 层。由表 \ref{tab:fold6_b3c3_root_dictionary_c3} 逐项读出，它们恰对应
\[
2e_1,\ \pm 2e_2,\ \pm 2e_3,\ \pm e_1\pm e_2.
\]
于是得到
\[
\iota_C(\mathcal R_4)=\{2e_1,\pm 2e_2,\pm 2e_3,\pm e_1\pm e_2\},
\]
故上式恰含 \(9\) 个向量。三分计数遂得。证毕。
\end{proof}

\begin{theorem}[window-\texorpdfstring{$6$}{6} 中心 \texorpdfstring{$2$}{2} 群在 \texorpdfstring{$C_3$}{C3} 根切片下的规范 \texorpdfstring{$5\oplus 3$}{5+3} 分裂]\label{thm:xi-window6-c3-center-five-plus-three-splitting}
\leanverified{paper\_xi\_window6\_c3\_center\_five\_plus\_three\_splitting}
沿用定理 \ref{thm:xi-window6-center-exact-splitting-bdry-cyclic} 与定理 \ref{thm:xi-window6-c3-cyclic-fiber-trisection} 的记号。记
\[
Z_{\mathrm{bdry}}:=(\ZZ/2\ZZ)^{X_6^{\mathrm{bdry}}},
\qquad
Z_{\mathrm{cyc}}:=(\ZZ/2\ZZ)^{\mathcal R_2}.
\]
则
\[
\boxed{\;
Z\!\bigl(G_6^{\mathrm{bin}}\bigr)
\cong
Z_{\mathrm{cyc}}\oplus Z_{\mathrm{bdry}}
\cong
(\ZZ/2\ZZ)^5\oplus(\ZZ/2\ZZ)^3.\;}
\]
其中 \(Z_{\mathrm{bdry}}\) 是 boundary parity 的规范三维子格，而 \(Z_{\mathrm{cyc}}\) 由 \(C_3\) 根切片
\[
\{-2e_1,\pm e_2\pm e_3\}
\]
上的五条最小纤维生成。
\end{theorem}
\begin{proof}
由定理 \ref{thm:xi-window6-center-exact-splitting-bdry-cyclic}，
\[
Z\!\bigl(G_6^{\mathrm{bin}}\bigr)
\cong
(\ZZ/2\ZZ)^{X_6^{(2)}}
\cong
(\ZZ/2\ZZ)^{X_6^{\mathrm{bdry}}}
\oplus
(\ZZ/2\ZZ)^{X_{6,\mathrm{cyc},2}}.
\]
而该定理已给出
\[
X_{6,\mathrm{cyc},2}
=
\{\texttt{000001},\texttt{000101},\texttt{001001},\texttt{010001},\texttt{010101}\}.
\]
再由定理 \ref{thm:xi-window6-c3-cyclic-fiber-trisection}，上述五个 cyclic 最小纤维在 \(C_3\) 规范下恰对应
\[
\{-2e_1,\pm e_2\pm e_3\}.
\]
故
\[
Z_{\mathrm{cyc}}\cong (\ZZ/2\ZZ)^5,
\qquad
Z_{\mathrm{bdry}}\cong (\ZZ/2\ZZ)^3.
\]
二者支撑于不交的纤维坐标，因此中心即按这两个坐标块规范分裂为
\[
Z\!\bigl(G_6^{\mathrm{bin}}\bigr)\cong Z_{\mathrm{cyc}}\oplus Z_{\mathrm{bdry}}
\cong
(\ZZ/2\ZZ)^5\oplus(\ZZ/2\ZZ)^3.
\]
证毕。
\end{proof}

\begin{theorem}[window-\texorpdfstring{$6$}{6} 的 boundary parity 模在 \texorpdfstring{$C_3$}{C3} 作用下的唯一分裂]\label{thm:xi-window6-boundary-parity-c3-module-splitting}
沿用定理 \ref{thm:xi-window6-center-exact-splitting-bdry-cyclic}、定理 \ref{thm:xi-time-part53da-boundary-z6-torsor} 与定理 \ref{thm:xi-time-part65e-window6-geometric-diagonal-z2-parity} 的记号。记
\[
X_6^{\mathrm{bdry}}
=
\{\texttt{100001},\texttt{100101},\texttt{101001}\},
\qquad
P_{\partial}:=(\FF_2)^{X_6^{\mathrm{bdry}}}\cong \FF_2^3.
\]
令 \(\sigma\) 为由三条 boundary 轴循环置换所诱导的 \(C_3\)-生成元作用，则作为 \(\FF_2[C_3]\)-模有规范分裂
\[
\boxed{\;
P_{\partial}\cong \mathbf 1\oplus V_2,\;}
\]
其中
\[
\mathbf 1=\langle(1,1,1)\rangle
\]
是唯一平凡子模，而 \(V_2\) 是唯一二维不可约 \(\FF_2[C_3]\)-模。等价地，
\[
\boxed{\;
P_{\partial}/\langle(1,1,1)\rangle\cong \FF_4\;}
\]
作为加法群成立，且 \(C_3\) 的生成元作用对应于乘上 \(\FF_4^\times\) 中的一个三次单位根。
\end{theorem}
\begin{proof}
由定理 \ref{thm:xi-window6-center-exact-splitting-bdry-cyclic}，
\[
P_{\partial}\cong (\FF_2)^{X_6^{\mathrm{bdry}}}
\cong
(\ZZ/2\ZZ)^3
\]
正是 boundary parity 的规范三维向量空间。另一方面，定理 \ref{thm:xi-time-part53da-boundary-z6-torsor} 的证明已给出：三条 boundary 轴存在一个循环置换的 \(C_3\)-作用。于是可取标准基 \(e_1,e_2,e_3\) 使
\[
\sigma(e_1)=e_2,\qquad \sigma(e_2)=e_3,\qquad \sigma(e_3)=e_1.
\]
因此 \(P_{\partial}\) 作为 \(\FF_2[C_3]\)-模同构于三循环的置换模，也即自由秩 \(1\) 的正则模
\[
P_{\partial}\cong \FF_2[C_3].
\]

在 \(\FF_2\) 上，\(\sigma\) 的特征多项式为
\[
t^3-1=(t+1)(t^2+t+1),
\]
且
\[
t^2+t+1
\]
在 \(\FF_2[t]\) 中不可约。故 \(P_{\partial}\) 的 Jordan--H\"older 组成因子只能是一个一维平凡因子和一个二维不可约因子。

向量
\[
(1,1,1)=e_1+e_2+e_3
\]
显然被 \(\sigma\) 固定，并且由定理 \ref{thm:xi-time-part65e-window6-geometric-diagonal-z2-parity}，它恰对应几何可实现的对角 parity 子群
\[
\langle(1,1,1)\rangle\cong \ZZ/2\ZZ.
\]
因此
\[
\mathbf 1:=\langle(1,1,1)\rangle
\]
是一个非零平凡子模。由于平凡因子重数为 \(1\)，它必是唯一的非零平凡子模。商模
\[
P_{\partial}/\mathbf 1
\]
于是只能是唯一的二维不可约模 \(V_2\)。

最后，
\[
\FF_2[t]/(t^2+t+1)\cong \FF_4,
\]
故该二维不可约模可识别为 \(\FF_4\) 的加法群，而 \(\sigma\) 的作用则对应于乘以 \(\FF_4^\times\) 中阶为 \(3\) 的元素。证毕。
\end{proof}

\begin{theorem}[window-\texorpdfstring{$6$}{6} boundary parity 的 \texorpdfstring{$C_3$}{C3}-等变投影仅有四种]\label{thm:xi-window6-boundary-parity-c3-equivariant-idempotents}
\leanverified{paper\_xi\_window6\_boundary\_parity\_c3\_equivariant\_idempotents}
沿用定理 \ref{thm:xi-window6-boundary-parity-c3-module-splitting} 的记号，则
\[
\operatorname{End}_{C_3}(P_{\partial})
\cong
\FF_2[C_3]
\cong
\FF_2[t]/(t^3-1)
\cong
\FF_2\times \FF_4.
\]
从而其幂等元恰有四个，也即 \(C_3\)-等变投影恰有四种：
\[
0,\qquad
\pi_{\mathrm{diag}}:P_{\partial}\twoheadrightarrow \langle(1,1,1)\rangle,\qquad
\pi_{V_2}:P_{\partial}\twoheadrightarrow V_2,\qquad
\mathrm{id}_{P_{\partial}}.
\]
特别地，不存在第五种非平凡 \(C_3\)-等变线性读出。
\end{theorem}
\begin{proof}
由定理 \ref{thm:xi-window6-boundary-parity-c3-module-splitting} 的证明，\(P_{\partial}\) 是 \(\FF_2[C_3]\) 上的自由秩 \(1\) 模，故
\[
P_{\partial}\cong \FF_2[C_3]
\]
作为左 \(\FF_2[C_3]\)-模成立。因此
\[
\operatorname{End}_{C_3}(P_{\partial})
\cong
\operatorname{End}_{\FF_2[C_3]}(\FF_2[C_3])
\cong
\FF_2[C_3]^{\mathrm{op}}
\cong
\FF_2[C_3],
\]
其中最后一步用到群代数交换。

再由分解
\[
t^3-1=(t+1)(t^2+t+1)
\]
及中国剩余定理，
\[
\FF_2[C_3]
\cong
\FF_2[t]/(t^3-1)
\cong
\FF_2[t]/(t+1)\times \FF_2[t]/(t^2+t+1)
\cong
\FF_2\times \FF_4.
\]
有限域直积环的幂等元恰为各分量上取 \(0\) 或 \(1\) 的四种组合，因此这里的幂等元正好只有
\[
(0,0),\qquad (1,0),\qquad (0,1),\qquad (1,1).
\]
把它们经上述同构拉回到 \(\operatorname{End}_{C_3}(P_{\partial})\)，便分别得到零投影、到平凡对角因子 \(\mathbf 1=\langle(1,1,1)\rangle\) 的投影、到二维不可约因子 \(V_2\) 的投影，以及恒等投影。故全部 \(C_3\)-等变投影恰有四种，不存在第五种非平凡等变线性读出。证毕。
\end{proof}
% (User input window)

\paragraph{奇偶谱盲点、boundary 中心扇区与 Lee--Yang 分支覆叠}
下列条目把折叠多重度的奇偶特征、window-\(6\) 的 boundary 中心分解、
规范体积的 Stirling 指数律，以及 Lee--Yang 振幅域与 RH 临界半径之间的若干
直接后果并列写成统一的定理化表述。

\begin{theorem}[折叠多重度的严格奇偶平衡]\label{thm:xi-fold-multiplicity-strict-parity-balance}
沿用命题 \ref{prop:fold-multiplicity-group-algebra} 与定理 \ref{thm:fold-spectrum-factorization}
的记号。记
\[
M:=F_{m+2},
\qquad
D_m(x):=\sum_{r=0}^{M-1}d_m(r)x^r\in \CC[x]/(x^M-1),
\]
\[
\widehat d_m(k):=\sum_{r=0}^{M-1} d_m(r)e^{-2\pi i kr/M}.
\]
若 \(m\equiv 1\pmod 3\)，则 \(M\) 为偶数，并且
\[
\boxed{\;
\widehat d_m(M/2)=0,
\qquad
D_m(-1)=0.\;}
\]
等价地，若定义
\[
E_m:=\sum_{\substack{0\le r<M\\ r\ \mathrm{even}}}d_m(r),
\qquad
O_m:=\sum_{\substack{0\le r<M\\ r\ \mathrm{odd}}}d_m(r),
\]
则有严格平衡
\[
\boxed{\;
E_m=O_m=2^{m-1}.\;}
\]
\end{theorem}
\begin{proof}
Fibonacci 数列模 \(2\) 的奇偶性以 \(3\) 为周期：
\(
F_n
\)
为偶数当且仅当
\(
n\equiv 0\pmod 3
\)。
因此当 \(m\equiv 1\pmod 3\) 时，
\(
m+2\equiv 0\pmod 3
\)，故
\[
M=F_{m+2}
\]
为偶数，从而
\(
M/2\in \ZZ
\)。
\par
由定理 \ref{thm:fold-spectrum-factorization}，
\[
\widehat d_m(M/2)
=
\prod_{j=1}^{m}\Bigl(1+e^{-\pi i F_{j+1}}\Bigr).
\]
取 \(j=1\) 即出现因子
\[
1+e^{-\pi i F_2}=1+e^{-\pi i}=0,
\]
故
\(
\widehat d_m(M/2)=0
\)。
\par
另一方面，由 \(M\) 为偶数可知
\(
(-1)^M-1=0
\)，因而取值映射 \(x\mapsto -1\) 在商环
\(
\CC[x]/(x^M-1)
\)
上良定。
再由命题 \ref{prop:fold-multiplicity-group-algebra}，
\[
D_m(-1)
=
\prod_{j=1}^{m}\Bigl(1+(-1)^{F_{j+1}}\Bigr)=0.
\]
事实上，
\[
D_m(-1)=\sum_{r=0}^{M-1}d_m(r)(-1)^r=\widehat d_m(M/2),
\]
故两种表述完全一致。
\par
最后，
\[
E_m-O_m=\sum_{r=0}^{M-1}d_m(r)(-1)^r=\widehat d_m(M/2)=0,
\]
而
\[
E_m+O_m=\sum_{r=0}^{M-1}d_m(r)=2^m,
\]
因为每个 \(\omega\in\Omega_m=\{0,1\}^m\) 恰贡献一个模 \(M\) 的余类。
联立即得
\(
E_m=O_m=2^{m-1}
\)。
证毕。
\leanverified{paper\_xi\_fold\_multiplicity\_strict\_parity\_balance}
\end{proof}

\begin{proposition}[window-\texorpdfstring{$6$}{6} boundary 奇偶荷的八扇区中心分解]\label{prop:xi-window6-boundary-eight-sector-central-decomposition}
沿用推论 \ref{cor:fold-bin-boundary-center-z2} 的记号。令
\[
G_6^{\mathrm{bin}}:=\mathcal G_6^{\mathrm{bin}}=\Aut(\Fold_6^{\mathrm{bin}}),
\qquad
Z_{\partial}:=\langle \tau_w:\ w\in X_6^{\mathrm{bdry}}\rangle
\cong (\ZZ/2\ZZ)^3
\hookrightarrow Z(G_6^{\mathrm{bin}}),
\]
并记其角色群
\[
\widehat Z_{\partial}:=\Hom(Z_{\partial},\{\pm1\}),
\qquad
|\widehat Z_{\partial}|=8.
\]
对每个 \(\chi\in \widehat Z_{\partial}\)，定义
\[
e_\chi:=\frac1{8}\sum_{z\in Z_{\partial}}\chi(z)\,z\in \CC[G_6^{\mathrm{bin}}].
\]
则每个 \(e_\chi\) 都是中心幂等元，并满足
\[
e_\chi e_{\chi'}=\delta_{\chi,\chi'}\,e_\chi,
\qquad
\sum_{\chi\in\widehat Z_{\partial}} e_\chi=1.
\]
因此复群代数存在规范直和分解
\[
\boxed{\;
\CC[G_6^{\mathrm{bin}}]
=
\bigoplus_{\chi\in\widehat Z_{\partial}}
\CC[G_6^{\mathrm{bin}}]\,e_\chi.\;}
\]
进一步，对任意有限维表示 \(V\)，令
\[
V_\chi:=e_\chi V.
\]
则
\[
\boxed{\;
V=\bigoplus_{\chi\in\widehat Z_{\partial}}V_\chi,\;}
\]
并且对任意 \(z\in Z_{\partial}\)，其在 \(V_\chi\) 上的作用为标量 \(\chi(z)\)。
特别地，三条 boundary sheet-flip 在每个 \(V_\chi\) 上同时对角化，
其特征值三元组恰由 \(\chi\) 唯一确定。
\end{proposition}
\begin{proof}
由于 \(Z_{\partial}\subset Z(G_6^{\mathrm{bin}})\)，子代数
\(
\CC[Z_{\partial}]
\)
自然嵌入
\(
Z(\CC[G_6^{\mathrm{bin}}])
\)。
又因 \(Z_{\partial}\cong(\ZZ/2\ZZ)^3\) 为有限 Abel 群，角色正交关系给出
\[
\frac1{8}\sum_{z\in Z_{\partial}}\chi(z)\chi'(z)
=
\delta_{\chi,\chi'}
\qquad
(\chi,\chi'\in \widehat Z_{\partial}),
\]
其中用到了 \(z^{-1}=z\) 与 \(\chi(z)\in\{\pm1\}\)。
因此标准 Fourier 幂等元满足
\[
e_\chi e_{\chi'}=\delta_{\chi,\chi'}e_\chi,
\qquad
\sum_{\chi}e_\chi=1.
\]
并且每个 \(e_\chi\) 由中心元线性组合而成，故自身亦为中心元。
\par
于是
\[
\CC[G_6^{\mathrm{bin}}]
=
\CC[G_6^{\mathrm{bin}}]\Bigl(\sum_{\chi}e_\chi\Bigr)
=
\sum_{\chi}\CC[G_6^{\mathrm{bin}}]e_\chi.
\]
由于这些幂等元两两正交，上式实为直和分解。
\par
对任意有限维表示 \(V\)，由
\(
\sum_\chi e_\chi=1
\)
得
\[
V=\sum_\chi e_\chi V,
\]
而由正交性
\(
e_\chi e_{\chi'}=0
\)
（\(\chi\neq\chi'\)）知该和为直和。
记 \(V_\chi:=e_\chi V\)，即得所示分解。
\par
最后，对任意 \(z\in Z_{\partial}\) 有
\[
ze_\chi
=
\frac1{8}\sum_{y\in Z_{\partial}}\chi(y)\,zy
=
\frac1{8}\sum_{y'\in Z_{\partial}}\chi(z^{-1}y')\,y'
=
\chi(z)e_\chi,
\]
其中第二个等号取 \(y'=zy\)，第三个等号用到 \(z^{-1}=z\) 与角色乘法性。
因此 \(z\) 在 \(V_\chi=e_\chi V\) 上恰以标量 \(\chi(z)\) 作用。
对三条生成元 \(\tau_w\) 同时应用此式，即得八个 boundary 奇偶荷扇区的中心分解。证毕。
\leanverified{paper\_xi\_window6\_boundary\_eight\_sector\_central\_decomposition}
\end{proof}

\begin{theorem}[window-\texorpdfstring{$6$}{6} boundary 奇偶荷的最小线性可见维数]\label{thm:xi-window6-boundary-minimal-linear-visible-dimension}
沿用命题 \ref{prop:xi-window6-boundary-eight-sector-central-decomposition} 与
定理 \ref{thm:xi-window6-center-exact-splitting-bdry-cyclic} 的记号，记
\[
Z_{\mathrm{bdry}}:=Z_{\partial}\cong (\ZZ/2\ZZ)^3
\hookrightarrow Z(G_6^{\mathrm{bin}}).
\]
设 \(k\) 为满足 \(\mathrm{char}(k)\neq 2\) 的域，且
\[
\rho:G_6^{\mathrm{bin}}\longrightarrow \GL_n(k)
\]
为一个在 \(Z_{\mathrm{bdry}}\) 上忠实的线性表示。
则必有
\[
n\ge 3.
\]
等价地，在奇特征或特征 \(0\) 的线性审计口径下，任何维数小于 \(3\) 的表示都不可能同时分辨三条 boundary sheet-flip。
\end{theorem}
\begin{proof}
由于 \(Z_{\mathrm{bdry}}\cong (\ZZ/2\ZZ)^3\) 由三个彼此对易的 involution 生成，且 \(\mathrm{char}(k)\neq 2\)，对任意 \(z\in Z_{\mathrm{bdry}}\) 而言，\(\rho(z)\) 的最小多项式整除
\[
t^2-1=(t-1)(t+1),
\]
故 \(\rho(z)\) 可对角化，且其特征值仅可能为 \(\pm 1\)。
又因这些算子两两对易，它们可在 \(k^n\) 上同时对角化。

于是 \(k^n\) 可分解为若干一维公共特征子空间之直和；等价地，存在角色
\[
\chi_1,\dots,\chi_n\in \Hom(Z_{\mathrm{bdry}},\{\pm 1\})
\]
使得每个 \(z\in Z_{\mathrm{bdry}}\) 在第 \(j\) 个特征方向上以标量 \(\chi_j(z)\) 作用。
若 \(\rho|_{Z_{\mathrm{bdry}}}\) 忠实，则
\[
\bigcap_{j=1}^{n}\ker(\chi_j)=\{1\}.
\]
这等价于 \(\chi_1,\dots,\chi_n\) 在对偶空间
\[
\Hom(Z_{\mathrm{bdry}},\{\pm 1\})
\cong (\FF_2)^3
\]
中张成整个三维空间。
因此至少需要三个线性无关的角色，遂得 \(n\ge 3\)。
\leanverified{paper\_xi\_window6\_boundary\_minimal\_linear\_visible\_dimension}
\end{proof}

\begin{theorem}[规范体积的 Boltzmann--Stirling 指数公式]\label{thm:xi-gauge-volume-boltzmann-stirling-exponential}
沿用命题 \ref{prop:fold-bin-gauge-volume} 与定理 \ref{thm:fold-bin-gauge-volume-stirling}
的记号。记
\[
\mathcal G_m^{\mathrm{bin}}:=\Aut(\Fold_m^{\mathrm{bin}}),
\qquad
g_m:=2^{-m}\log|\mathcal G_m^{\mathrm{bin}}|,
\]
\[
\kappa_m:=2^{-m}\sum_{w\in X_m}d_m^{\mathrm{bin}}(w)\log d_m^{\mathrm{bin}}(w),
\qquad
\pi_m(w):=\frac{d_m^{\mathrm{bin}}(w)}{2^m}.
\]
则成立指数等价
\[
\boxed{\;
|\mathcal G_m^{\mathrm{bin}}|
=
\exp\!\bigl(-2^m+o(2^m)\bigr)\,
\prod_{w\in X_m}\Bigl(d_m^{\mathrm{bin}}(w)\Bigr)^{d_m^{\mathrm{bin}}(w)}
\qquad(m\to\infty).\;}
\]
等价地，
\[
\boxed{\;
\lim_{m\to\infty}
\frac{1}{2^m}\log
\frac{\prod_{w\in X_m}(d_m^{\mathrm{bin}}(w))^{d_m^{\mathrm{bin}}(w)}}{|\mathcal G_m^{\mathrm{bin}}|}
=1.\;}
\]
并且在以 \(\log_2\) 为单位时，
\[
\boxed{\;
\lim_{m\to\infty}
\left(
\mathbb E_{\pi_m}\!\bigl[\log_2 d_m^{\mathrm{bin}}(X)\bigr]
-\frac{1}{2^m\log 2}\log|\mathcal G_m^{\mathrm{bin}}|
\right)
=
\log_2 e.\;}
\]
\end{theorem}
\begin{proof}
由定理 \ref{thm:fold-bin-gauge-volume-stirling}，
\[
g_m=\kappa_m-1+O\!\left(\frac{|X_m|\,m}{2^m}\right).
\]
两边乘以 \(2^m\) 得
\[
\log|\mathcal G_m^{\mathrm{bin}}|
=
\sum_{w\in X_m}d_m^{\mathrm{bin}}(w)\log d_m^{\mathrm{bin}}(w)
-2^m
+O\!\bigl(|X_m|\,m\bigr).
\]
又因
\[
\sum_{w\in X_m}d_m^{\mathrm{bin}}(w)\log d_m^{\mathrm{bin}}(w)
=
\log \prod_{w\in X_m}\Bigl(d_m^{\mathrm{bin}}(w)\Bigr)^{d_m^{\mathrm{bin}}(w)},
\]
且
\[
|X_m|=F_{m+2}=\Theta(\varphi^m)=o(2^m),
\]
故
\[
\log|\mathcal G_m^{\mathrm{bin}}|
=
\log \prod_{w\in X_m}\Bigl(d_m^{\mathrm{bin}}(w)\Bigr)^{d_m^{\mathrm{bin}}(w)}
-2^m+o(2^m).
\]
指数化即得第一式。
\par
将上式移项后除以 \(2^m\)，得到
\[
\frac{1}{2^m}\log
\frac{\prod_{w\in X_m}(d_m^{\mathrm{bin}}(w))^{d_m^{\mathrm{bin}}(w)}}{|\mathcal G_m^{\mathrm{bin}}|}
=1+o(1),
\]
故第二式成立。
\par
最后，由
\[
\mathbb E_{\pi_m}\!\bigl[\log_2 d_m^{\mathrm{bin}}(X)\bigr]
=
\frac{\kappa_m}{\log 2}
\]
与
\[
\frac{1}{2^m\log 2}\log|\mathcal G_m^{\mathrm{bin}}|
=
\frac{g_m}{\log 2},
\]
可得所求差值等于
\[
\frac{\kappa_m-g_m}{\log 2}.
\]
再由
\(
\kappa_m-g_m\to 1
\)
得到极限
\(
1/\log 2=\log_2 e
\)。
证毕。
\end{proof}

\paragraph{规范体积密度、纤维奇偶荷与可解性阈值}
在定理 \ref{thm:xi-foldbin-collision-probability-normalized-groupoid-dimension} 的碰撞公式与定理 \ref{thm:xi-gauge-volume-boltzmann-stirling-exponential} 的指数律之间，还存在一条由 escort 相对熵控制的精确桥接；同时，定义 \ref{def:fold-bin-gauge-sign} 的纤维奇偶荷可把 Abel 化、中心、window-\(6\) 的全自同构结构以及低阶上同调统一为纯纤维判别。

\begin{theorem}[规范体积密度的二次 escort--KL 精确分解]\label{thm:xi-foldbin-gauge-escort-kl-decomposition}
\leanverified{paper\_xi\_foldbin\_gauge\_escort\_kl\_decomposition}
沿用命题 \ref{prop:fold-bin-chi2-col} 与命题 \ref{prop:fold-bin-gauge-volume} 的记号。记
\[
p_m^{\mathrm{bin}}(w):=\frac{d_m^{\mathrm{bin}}(w)}{2^m},
\qquad
\widetilde p_m(w):=\frac{\bigl(p_m^{\mathrm{bin}}(w)\bigr)^2}{\mathsf{Col}^{\mathrm{bin}}_m}
\qquad (w\in X_m).
\]
则有严格恒等式
\[
\boxed{\;
\kappa_m
=
\log\!\bigl(2^m\mathsf{Col}^{\mathrm{bin}}_m\bigr)
-D_{\mathrm{KL}}\!\bigl(p_m^{\mathrm{bin}}\|\widetilde p_m\bigr).
\;}
\]
等价地，
\[
\boxed{\;
\kappa_m
=
\log\!\Bigl(\frac{2^m}{|X_m|}\Bigr)
+\log\!\Bigl(1+\chi^2\!\bigl(p_m^{\mathrm{bin}}\|u_m\bigr)\Bigr)
-D_{\mathrm{KL}}\!\bigl(p_m^{\mathrm{bin}}\|\widetilde p_m\bigr),
\;}
\]
其中 \(u_m(w)=1/|X_m|\)。
特别地，
\[
\log\!\Bigl(\frac{2^m}{|X_m|}\Bigr)
\le
\kappa_m
\le
\log\!\bigl(2^m\mathsf{Col}^{\mathrm{bin}}_m\bigr),
\]
且两端取等都当且仅当 \(p_m^{\mathrm{bin}}=u_m\)。
\end{theorem}
\begin{proof}
为简记起见，写
\[
d_w:=d_m^{\mathrm{bin}}(w),
\qquad
p_w:=p_m^{\mathrm{bin}}(w),
\qquad
H(p):=-\sum_{w\in X_m}p_w\log p_w.
\]
由于 \(X_m\) 由实际出现的稳定类型组成，故对每个 \(w\in X_m\) 都有 \(d_w\ge 1\)，从而 \(p_w>0\)。
由定义
\[
\kappa_m
=
2^{-m}\sum_{w\in X_m}d_w\log d_w
=
\sum_{w\in X_m}p_w\log(2^m p_w)
=
m\log 2-H(p).
\]
另一方面，
\[
\begin{aligned}
D_{\mathrm{KL}}(p\|\widetilde p)
&=
\sum_{w\in X_m}p_w\log\frac{p_w}{p_w^2/\mathsf{Col}^{\mathrm{bin}}_m}\\
&=
\log\!\bigl(\mathsf{Col}^{\mathrm{bin}}_m\bigr)-\sum_{w\in X_m}p_w\log p_w\\
&=
\log\!\bigl(\mathsf{Col}^{\mathrm{bin}}_m\bigr)+H(p).
\end{aligned}
\]
消去 \(H(p)\) 即得
\[
\kappa_m
=
\log\!\bigl(2^m\mathsf{Col}^{\mathrm{bin}}_m\bigr)
-D_{\mathrm{KL}}(p\|\widetilde p).
\]
再由命题 \ref{prop:fold-bin-chi2-col} 的恒等式
\[
\chi^2\!\bigl(p_m^{\mathrm{bin}}\|u_m\bigr)
=
|X_m|\,\mathsf{Col}^{\mathrm{bin}}_m-1,
\]
便得到第二种写法。

下界由 Shannon 熵上界 \(H(p)\le \log|X_m|\) 立即推出，即
\[
\kappa_m=m\log 2-H(p)\ge m\log 2-\log|X_m|.
\]
且等号成立当且仅当 \(p=u_m\)。
上界则由 \(D_{\mathrm{KL}}(p\|\widetilde p)\ge 0\) 得到；若上界取等，则 \(p=\widetilde p\)，即
\[
p_w=\frac{p_w^2}{\mathsf{Col}^{\mathrm{bin}}_m}
\qquad (w\in X_m).
\]
由于所有 \(p_w>0\)，故 \(p_w=\mathsf{Col}^{\mathrm{bin}}_m\) 对所有 \(w\) 同时成立，于是 \(p\) 在 \(X_m\) 上均匀，即 \(p=u_m\)。
反之若 \(p=u_m\)，则显然同时达到上下界。证毕。
\end{proof}

\begin{corollary}[规范体积密度的碰撞表达与常数主差额]\label{cor:xi-foldbin-gauge-collision-rigid-shift}
\leanverified{paper\_xi\_foldbin\_gauge\_collision\_rigid\_shift}
沿用定理 \ref{thm:xi-foldbin-gauge-escort-kl-decomposition} 的记号，则
\[
\boxed{\;
g_m
=
\log\!\bigl(2^m\mathsf{Col}^{\mathrm{bin}}_m\bigr)
-1
-D_{\mathrm{KL}}\!\bigl(p_m^{\mathrm{bin}}\|\widetilde p_m\bigr)
+O\!\Bigl(m\bigl(\tfrac{\varphi}{2}\bigr)^m\Bigr).
\;}
\]
等价地，
\[
\boxed{\;
g_m
=
\log\!\Bigl(\frac{2^m}{|X_m|}\Bigr)
+\log\!\Bigl(1+\chi^2\!\bigl(p_m^{\mathrm{bin}}\|u_m\bigr)\Bigr)
-1
-D_{\mathrm{KL}}\!\bigl(p_m^{\mathrm{bin}}\|\widetilde p_m\bigr)
+O\!\Bigl(m\bigl(\tfrac{\varphi}{2}\bigr)^m\Bigr).
\;}
\]
因此 \(g_m\) 与碰撞项 \(\log(2^m\mathsf{Col}^{\mathrm{bin}}_m)\) 的主差额由常数 \(1\) 固定，余项仅为指数小修正。
\end{corollary}
\begin{proof}
由定理 \ref{thm:fold-bin-gauge-volume-stirling}，
\[
g_m=\kappa_m-1+O\!\left(\frac{|X_m|\,m}{2^m}\right)
=
\kappa_m-1+O\!\Bigl(m\bigl(\tfrac{\varphi}{2}\bigr)^m\Bigr),
\]
其中用到了 \(|X_m|=F_{m+2}=\Theta(\varphi^m)\)。
将定理 \ref{thm:xi-foldbin-gauge-escort-kl-decomposition} 的两种表示分别代入即得。证毕。
\end{proof}

\begin{theorem}[纤维奇偶荷的分裂正合列与最小完备秩]\label{thm:xi-foldbin-parity-charge-split-exact-minrank}
\leanverified{paper\_xi\_foldbin\_parity\_charge\_split\_exact\_minrank}
沿用命题 \ref{prop:fold-bin-gauge-decomposition} 与定义 \ref{def:fold-bin-gauge-sign} 的记号。记
\[
G_m^{\mathrm{bin}}:=\mathcal{G}^{\mathrm{bin}}_m=\Aut(\Fold_m^{\mathrm{bin}}),
\qquad
d_w:=d_m^{\mathrm{bin}}(w),
\qquad
N_m:=\#\{w\in X_m:\ d_w\ge 2\}.
\]
将定义 \ref{def:fold-bin-gauge-sign} 中的
\(
\mathrm{sgn}_m:G_m^{\mathrm{bin}}\to (\ZZ/2\ZZ)^{X_m}
\)
与坐标投影
\[
(\ZZ/2\ZZ)^{X_m}
\longrightarrow
(\ZZ/2\ZZ)^{\{w\in X_m:\ d_w\ge 2\}}
\cong
(\ZZ/2\ZZ)^{N_m}
\]
复合后仍记为 \(\mathrm{sgn}_m\)。
约定 \(\mathfrak A_1=\mathfrak A_2=1\)。
则存在分裂正合列
\[
\boxed{\;
1\longrightarrow \prod_{w\in X_m}\mathfrak A_{d_w}
\longrightarrow G_m^{\mathrm{bin}}
\xrightarrow{\ \mathrm{sgn}_m\ }
(\ZZ/2\ZZ)^{N_m}
\longrightarrow 1.
\;}
\]
从而：
\begin{enumerate}
  \item \(\bigl(G_m^{\mathrm{bin}}\bigr)^{\mathrm{ab}}\cong (\ZZ/2\ZZ)^{N_m}\)；
  \item 若群同态 \(\Phi:G_m^{\mathrm{bin}}\to (\ZZ/2\ZZ)^r\) 满足存在线性映射
  \[
  L:(\ZZ/2\ZZ)^r\to (\ZZ/2\ZZ)^{N_m}
  \]
  使得 \(\mathrm{sgn}_m=L\circ\Phi\)，则必有 \(r\ge N_m\)；
  \item 上述下界由 \(\mathrm{sgn}_m\) 自身达到，因此 \(N_m\) 正是纤维奇偶荷的最小完备秩。
\end{enumerate}
特别地，在 \(m=6\) 时，由推论 \ref{cor:fold-bin-m6-fiber-hist} 的直方图 \(2:8,3:4,4:9\) 可得 \(N_6=21\)，于是
\[
\bigl(G_6^{\mathrm{bin}}\bigr)^{\mathrm{ab}}\cong (\ZZ/2\ZZ)^{21}.
\]
\end{theorem}
\begin{proof}
由命题 \ref{prop:fold-bin-gauge-decomposition}，
\[
G_m^{\mathrm{bin}}
\cong
\prod_{w\in X_m}\mathfrak S_{d_w}.
\]
于是只需逐纤维考察对称群。

当 \(d=1\) 时 \(\mathfrak S_1\) 平凡。
当 \(d=2\) 时 \(\mathfrak S_2\cong \ZZ/2\ZZ\)，其符号映射即为同构，核为平凡群，也即 \(\mathfrak A_2=1\)。
当 \(d\ge 3\) 时，有标准分裂正合列
\[
1\longrightarrow \mathfrak A_d
\longrightarrow \mathfrak S_d
\xrightarrow{\ \mathrm{sgn}\ }
\ZZ/2\ZZ
\longrightarrow 1,
\]
任取一换位即可给出截面。
对全部 \(w\in X_m\) 取直积，便得到所述分裂正合列。

第 (1) 条随即成立，因为右端已是所有因子 Abel 化的直积。
对第 (2) 条，条件 \(\mathrm{sgn}_m=L\circ\Phi\) 表明 \((\ZZ/2\ZZ)^{N_m}\) 是 \(\Phi(G_m^{\mathrm{bin}})\subseteq (\ZZ/2\ZZ)^r\) 的一个线性商。
由于 \(\mathrm{sgn}_m\) 满射，故
\[
\operatorname{rank}_{\FF_2}\Phi(G_m^{\mathrm{bin}})
\ge N_m,
\]
而 \(\Phi(G_m^{\mathrm{bin}})\) 是 \((\ZZ/2\ZZ)^r\) 的子空间，故其秩至多为 \(r\)，于是 \(r\ge N_m\)。
第 (3) 条由 \(\mathrm{sgn}_m\) 本身即为到 \((\ZZ/2\ZZ)^{N_m}\) 的满射立即得到。

最后，在 \(m=6\) 时，表 \ref{tab:fold_bin_degeneracy_spectrum_m6_m12} 的直方图给出全部 \(21\) 条纤维都满足 \(d_w\ge 2\)，故 \(N_6=21\)。证毕。
\end{proof}

\begin{theorem}[中心与可解性的纤维判别]\label{thm:xi-foldbin-center-solvable-fiber-criterion}
\leanverified{paper\_xi\_foldbin\_center\_solvable\_fiber\_criterion}
沿用命题 \ref{prop:fold-bin-gauge-decomposition} 的记号。记
\[
G_m^{\mathrm{bin}}:=\mathcal{G}^{\mathrm{bin}}_m,
\qquad
d_w:=d_m^{\mathrm{bin}}(w),
\qquad
B_m:=\#\{w\in X_m:\ d_w=2\},
\qquad
d_{\max}(m):=\max_{w\in X_m}d_w.
\]
则中心满足
\[
\boxed{\;
Z\!\bigl(G_m^{\mathrm{bin}}\bigr)\cong (\ZZ/2\ZZ)^{B_m}.
\;}
\]
此外，\(G_m^{\mathrm{bin}}\) 可解当且仅当 \(d_{\max}(m)\le 4\)。
若以 \(\ell_{\mathrm{der}}(G)\) 记群 \(G\) 的导出长度，则在可解情形下
\[
\ell_{\mathrm{der}}\!\bigl(G_m^{\mathrm{bin}}\bigr)
=
\begin{cases}
0,& d_{\max}(m)=1,\\
1,& d_{\max}(m)=2,\\
2,& d_{\max}(m)=3,\\
3,& d_{\max}(m)=4.
\end{cases}
\]
特别地，表 \ref{tab:fold_bin_degeneracy_spectrum_m6_m12} 的 \(m=6\) 行给出 \(B_6=8\) 与 \(d_{\max}(6)=4\)，故
\[
Z\!\bigl(G_6^{\mathrm{bin}}\bigr)\cong (\ZZ/2\ZZ)^8,
\qquad
\ell_{\mathrm{der}}\!\bigl(G_6^{\mathrm{bin}}\bigr)=3.
\]
并且定理 \ref{thm:xi-window6-center-exact-splitting-bdry-cyclic} 中的 boundary 中心子群
\[
(\ZZ/2\ZZ)^3_{\mathrm{bdry}}\hookrightarrow Z\!\bigl(G_6^{\mathrm{bin}}\bigr)
\]
正是该 \(8\) 维中心的规范直和因子。
另一方面，表 \ref{tab:fold_bin_degeneracy_spectrum_m6_m12} 给出 \(d_{\max}(10)=9\)，而定理 \ref{thm:fold-bin-min-degeneracy-fib-audited} 给出 \(d_{\min}(12)=8\)，故 \(G_{10}^{\mathrm{bin}}\) 与 \(G_{12}^{\mathrm{bin}}\) 都不可解。
\end{theorem}
\begin{proof}
由命题 \ref{prop:fold-bin-gauge-decomposition}，
\[
G_m^{\mathrm{bin}}
\cong
\prod_{w\in X_m}\mathfrak S_{d_w}.
\]
对有限直积群有
\[
Z\!\left(\prod_{w\in X_m}\mathfrak S_{d_w}\right)
\cong
\prod_{w\in X_m}Z(\mathfrak S_{d_w}).
\]
而
\[
Z(\mathfrak S_d)\cong
\begin{cases}
\ZZ/2\ZZ,& d=2,\\
1,& d\neq 2,
\end{cases}
\]
故立得 \(Z(G_m^{\mathrm{bin}})\cong (\ZZ/2\ZZ)^{B_m}\)。

有限直积群可解当且仅当每个因子都可解，而 \(\mathfrak S_d\) 可解当且仅当 \(d\le 4\)。
因此
\[
G_m^{\mathrm{bin}}\ \text{可解}
\iff
d_{\max}(m)\le 4.
\]
在可解情形下，导出长度等于各因子导出长度的最大值。
再用
\[
\ell_{\mathrm{der}}(\mathfrak S_1)=0,\qquad
\ell_{\mathrm{der}}(\mathfrak S_2)=1,\qquad
\ell_{\mathrm{der}}(\mathfrak S_3)=2,\qquad
\ell_{\mathrm{der}}(\mathfrak S_4)=3,
\]
便得到所示分段公式。

\(m=6\) 的结论由表 \ref{tab:fold_bin_degeneracy_spectrum_m6_m12} 的直方图 \(2:8,3:4,4:9\) 直接给出，而 boundary 中心子群为规范直和因子正是定理 \ref{thm:xi-window6-center-exact-splitting-bdry-cyclic} 的内容。
对 \(m=10\)，表中存在 \(d=9\) 纤维，故含不可解因子 \(\mathfrak S_9\)。
对 \(m=12\)，由定理 \ref{thm:fold-bin-min-degeneracy-fib-audited} 的 \(d_{\min}(12)=8\) 可知每条纤维都满足 \(d_w\ge 8\)，从而同样不可解。证毕。
\end{proof}

\begin{theorem}[window-\texorpdfstring{$6$}{6} 规范群全自同构群的完全结构]\label{thm:xi-window6-full-automorphism-group-structure}
\leanverified{paper\_xi\_window6\_full\_automorphism\_group\_structure}
设
\[
G:=G_6^{\mathrm{bin}}
\cong
(\ZZ/2\ZZ)^8\times S_3^{\,4}\times S_4^{\,9}.
\]
则存在自然半直积同构
\[
\boxed{\;
\Aut(G)
\cong
\Hom\!\bigl((\ZZ/2\ZZ)^{13},(\ZZ/2\ZZ)^8\bigr)
\rtimes
\Bigl(
\GL_8(\FF_2)\times (S_3\wr S_4)\times (S_4\wr S_9)
\Bigr).\;}
\]
特别地，
\[
\Hom\!\bigl((\ZZ/2\ZZ)^{13},(\ZZ/2\ZZ)^8\bigr)\cong (\ZZ/2\ZZ)^{104},
\]
从而
\[
\boxed{\;
\Aut(G)
\cong
(\ZZ/2\ZZ)^{104}
\rtimes
\Bigl(
\GL_8(\FF_2)\times (S_3\wr S_4)\times (S_4\wr S_9)
\Bigr).\;}
\]
\end{theorem}
\begin{proof}
记
\[
A:=(\ZZ/2\ZZ)^8,
\qquad
H:=S_3^{\,4}\times S_4^{\,9},
\]
则 \(G\cong A\times H\)。由定理 \ref{thm:xi-foldbin-center-solvable-fiber-criterion}，
\[
Z(G)\cong (\ZZ/2\ZZ)^8\cong A.
\]
又 \(Z(S_3)=Z(S_4)=1\)，故
\[
Z(H)=1.
\]
因此 \(A=Z(G)\) 为特征子群；任意 \(\varphi\in\Aut(G)\) 都保持 \(A\)，并诱导
\[
\alpha\in\Aut(A),
\qquad
\beta\in\Aut(H).
\]

对任意 \(h\in H\)，写
\[
\varphi(0,h)=(\gamma(h),\beta(h))
\]
其中 \(\gamma(h)\in A\)。由于 \(A\) 中心化整个 \(G\)，有
\[
\varphi(0,h_1h_2)=\varphi(0,h_1)\varphi(0,h_2)
=
\bigl(\gamma(h_1)+\gamma(h_2),\,\beta(h_1)\beta(h_2)\bigr),
\]
故 \(\gamma:H\to A\) 为群同态。反之，任取
\[
\alpha\in\Aut(A),
\qquad
\beta\in\Aut(H),
\qquad
\gamma\in\Hom(H,A),
\]
映射
\[
(a,h)\longmapsto (\alpha(a)+\gamma(h),\,\beta(h))
\]
定义 \(G\) 的自同构。于是
\[
\Aut(G)\cong \Hom(H,A)\rtimes\bigl(\Aut(A)\times\Aut(H)\bigr),
\]
其中半直积作用由
\[
(\alpha,\beta)\cdot \gamma=\alpha\circ \gamma\circ \beta^{-1}
\]
给出。

由 \(A\cong(\ZZ/2\ZZ)^8\) 立得
\[
\Aut(A)\cong \GL_8(\FF_2).
\]
另一方面，
\[
H^{\mathrm{ab}}
\cong
(S_3^{\mathrm{ab}})^4\times (S_4^{\mathrm{ab}})^9
\cong
(\ZZ/2\ZZ)^{4+9}
=(\ZZ/2\ZZ)^{13},
\]
故
\[
\Hom(H,A)
\cong
\Hom(H^{\mathrm{ab}},A)
\cong
\Hom\!\bigl((\ZZ/2\ZZ)^{13},(\ZZ/2\ZZ)^8\bigr)
\cong
(\ZZ/2\ZZ)^{104}.
\]

最后，\(S_3\) 与 \(S_4\) 都是 complete 群，且彼此不同构，因而 \(H\) 的自同构不能在两类因子之间混合；并且对 complete 群 \(K\) 的 \(n\) 重直积有标准识别
\[
\Aut(K^n)\cong K\wr S_n.
\]
于是
\[
\Aut(H)
\cong
\Aut(S_3^4)\times \Aut(S_4^9)
\cong
(S_3\wr S_4)\times (S_4\wr S_9).
\]
综上即得结论。证毕。
\end{proof}

\begin{theorem}[window-\texorpdfstring{$6$}{6} 规范群 Schur 乘子的 \texorpdfstring{$2$}{2}-初等结构与秩]\label{thm:xi-window6-schur-multiplier-219}
\leanverified{paper\_xi\_window6\_schur\_multiplier\_219}
沿用定理 \ref{thm:xi-window6-full-automorphism-group-structure} 的记号。则
\[
\boxed{\;
M(G)\cong (\ZZ/2\ZZ)^{219}.\;}
\]
等价地，
\[
\boxed{\;
H^2(G,\CC^\times)\cong (\ZZ/2\ZZ)^{219}.\;}
\]
因此 window-\(6\) 规范群的复射影表示障碍完全由 \(219\) 维 \(\FF_2\)-向量空间参数化。
\end{theorem}
\begin{proof}
由直积群的 K\"unneth--Schur 公式，
\[
M(G_1\times G_2)
\cong
M(G_1)\oplus M(G_2)\oplus \bigl(G_1^{\mathrm{ab}}\otimes G_2^{\mathrm{ab}}\bigr).
\]
把
\[
G\cong (\ZZ/2\ZZ)^8\times S_3^{\,4}\times S_4^{\,9}
\]
逐步分解即可。

先看单因子。已知
\[
M(\ZZ/2\ZZ)=0,
\qquad
M(S_3)=0,
\qquad
M(S_4)\cong \ZZ/2\ZZ,
\]
并且
\[
(\ZZ/2\ZZ)^{\mathrm{ab}}\cong \ZZ/2\ZZ,
\qquad
S_3^{\mathrm{ab}}\cong \ZZ/2\ZZ,
\qquad
S_4^{\mathrm{ab}}\cong \ZZ/2\ZZ.
\]
因此各部分贡献分别为：
\[
M\bigl((\ZZ/2\ZZ)^8\bigr)\cong (\ZZ/2\ZZ)^{\binom82}=(\ZZ/2\ZZ)^{28},
\]
\[
M(S_3^4)\cong (\ZZ/2\ZZ)^{\binom42}=(\ZZ/2\ZZ)^6,
\]
\[
M(S_4^9)\cong (\ZZ/2\ZZ)^{9+\binom92}=(\ZZ/2\ZZ)^{45}.
\]
再加上三类交叉张量项
\[
(\ZZ/2\ZZ)^8\otimes (S_3^4)^{\mathrm{ab}}\cong (\ZZ/2\ZZ)^{8\cdot 4}=(\ZZ/2\ZZ)^{32},
\]
\[
(\ZZ/2\ZZ)^8\otimes (S_4^9)^{\mathrm{ab}}\cong (\ZZ/2\ZZ)^{8\cdot 9}=(\ZZ/2\ZZ)^{72},
\]
\[
(S_3^4)^{\mathrm{ab}}\otimes (S_4^9)^{\mathrm{ab}}\cong (\ZZ/2\ZZ)^{4\cdot 9}=(\ZZ/2\ZZ)^{36}.
\]
总秩于是为
\[
28+6+45+32+72+36=219.
\]
故
\[
M(G)\cong (\ZZ/2\ZZ)^{219}.
\]

对有限群 \(G\)，Schur 乘子满足
\[
M(G)\cong H_2(G;\ZZ),
\]
而 \(\CC^\times\) 可除，因此泛系数定理给出
\[
H^2(G,\CC^\times)\cong \Hom(M(G),\CC^\times).
\]
由于 \(M(G)\) 已是 \(2\)-初等阿贝尔群，其 Pontryagin 对偶与自身同构，故
\[
H^2(G,\CC^\times)\cong (\ZZ/2\ZZ)^{219}.
\]
证毕。
\end{proof}

\begin{theorem}[window-\texorpdfstring{$6$}{6} boundary 二层种子的规范置换表示分解]\label{thm:xi-window6-boundary-two-sheet-permutation-decomposition}
\leanverified{paper\_xi\_window6\_boundary\_two\_sheet\_permutation\_decomposition}
沿用定理 \ref{thm:xi-window6-center-exact-splitting-bdry-cyclic} 与命题
\ref{prop:xi-window6-sheet-parity-anchor-m6} 的记号。记
\[
X_6^{\mathrm{bdry}}
=
\{\texttt{100001},\texttt{100101},\texttt{101001}\},
\]
\[
\widetilde X_6^{\mathrm{bdry}}
:=
\bigsqcup_{w\in X_6^{\mathrm{bdry}}}
\bigl(\Fold_6^{\mathrm{bin}}\bigr)^{-1}(w).
\]
则 \(\widetilde X_6^{\mathrm{bdry}}\) 是由三条两点纤维组成的六点集；每条纤维上的唯一非平凡交换诱导一个全局对合
\[
\zeta:\widetilde X_6^{\mathrm{bdry}}\to \widetilde X_6^{\mathrm{bdry}},
\]
而基底三点集 \(X_6^{\mathrm{bdry}}\) 的置换群给出自然作用
\[
S_3\curvearrowright \widetilde X_6^{\mathrm{bdry}}.
\]
于是复置换表示满足规范分解
\[
\boxed{\;
\CC[\widetilde X_6^{\mathrm{bdry}}]
\cong
\mathbf 1_{+}\oplus \mathrm{Std}_{+}\oplus \mathbf 1_{-}\oplus \mathrm{Std}_{-}.\;}
\]
其中下标 \(\pm\) 表示 \(\zeta\) 的 \(\pm 1\)-本征空间，而 \(\mathrm{Std}\) 为 \(S_3\) 的二维标准表示。特别地，
\[
\boxed{\;
\CC[\widetilde X_6^{\mathrm{bdry}}]
\ \text{中不出现}\ 
\mathrm{sgn}_{+}\ \text{或}\ \mathrm{sgn}_{-}. \;}
\]
\end{theorem}
\begin{proof}
由命题 \ref{prop:xi-window6-sheet-parity-anchor-m6}，
\[
\Delta_6^{\mathrm{bdry}}=\{0,34\},
\]
故每条 boundary 纤维恰有两点，从而
\[
|\widetilde X_6^{\mathrm{bdry}}|=2\,|X_6^{\mathrm{bdry}}|=6.
\]
对 \(X_6^{\mathrm{bdry}}\) 的三个元素取编号 \(w_1,w_2,w_3\)，并把相应纤维写成
\[
\widetilde w_i=\{w_i^+,w_i^-\}\qquad (1\le i\le 3).
\]
记 \(e_i^\pm\) 为 \(\CC[\widetilde X_6^{\mathrm{bdry}}]\) 的对应基向量，并定义
\[
u_i:=e_i^++e_i^-,
\qquad
v_i:=e_i^+-e_i^-.
\]
由构造可知
\[
\zeta(u_i)=u_i,
\qquad
\zeta(v_i)=-v_i,
\]
故有直和分解
\[
\CC[\widetilde X_6^{\mathrm{bdry}}]
=
U_+\oplus U_-,
\qquad
U_+:=\Span\{u_1,u_2,u_3\},
\quad
U_-:=\Span\{v_1,v_2,v_3\}.
\]
另一方面，\(S_3\) 只在三条基底纤维之间作置换，因此在 \(U_+\) 与 \(U_-\) 上都实现同一个三点置换表示。于是
\[
U_\pm\cong \mathbf 1\oplus \mathrm{Std}.
\]
把 \(\zeta\)-本征标签附回，即得
\[
U_+\cong \mathbf 1_+\oplus \mathrm{Std}_+,
\qquad
U_-\cong \mathbf 1_-\oplus \mathrm{Std}_-.
\]
两式合并便得到所示分解。

最后，三点置换表示的特征标分解只含平凡表示与标准表示，不含符号表示；因此在 \(U_+\) 与 \(U_-\) 中都不会出现 \(\mathrm{sgn}\) 分量，从而全表示中亦无 \(\mathrm{sgn}_{\pm}\)。证毕。
\end{proof}

\begin{corollary}[boundary 二层种子的对称交换子代数恰为四维]\label{cor:xi-window6-boundary-two-sheet-commutant-fourdimensional}
\leanverified{paper\_xi\_window6\_boundary\_two\_sheet\_commutant\_fourdimensional}
沿用定理 \ref{thm:xi-window6-boundary-two-sheet-permutation-decomposition} 的记号，则
\[
\boxed{\;
\End_{S_3\times \langle\zeta\rangle}
\bigl(\CC[\widetilde X_6^{\mathrm{bdry}}]\bigr)
\cong \CC^4.\;}
\]
等价地，任一同时尊重三条 boundary 纤维无标号置换对称与 sheet parity 的复线性可观测量，恰有四个独立标量自由度，分别对应
\[
\mathbf 1_{+},\qquad
\mathrm{Std}_{+},\qquad
\mathbf 1_{-},\qquad
\mathrm{Std}_{-}.
\]
\end{corollary}
\begin{proof}
由定理 \ref{thm:xi-window6-boundary-two-sheet-permutation-decomposition}，
\[
\CC[\widetilde X_6^{\mathrm{bdry}}]
\cong
\mathbf 1_{+}\oplus \mathrm{Std}_{+}\oplus \mathbf 1_{-}\oplus \mathrm{Std}_{-}.
\]
这四个 \(S_3\times \langle\zeta\rangle\)-不可约分量两两不等价：\(\mathbf 1\) 与 \(\mathrm{Std}\) 作为 \(S_3\)-表示本已不等价，而 \(\zeta\) 的本征值不同也排除了跨奇偶空间的等价。于是由 Schur 引理，
\[
\End_{S_3\times \langle\zeta\rangle}
\bigl(\CC[\widetilde X_6^{\mathrm{bdry}}]\bigr)
\]
恰由在四个不可约块上分别作标量乘法的算子组成，从而同构于 \(\CC^4\)。证毕。
\end{proof}

\begin{theorem}[二层与三层无标号 boundary uplift 的统一外加信息下界]\label{thm:xi-boundary-uplift-mixed-radix-register-lower-bound}
\leanverified{paper\_xi\_boundary\_uplift\_mixed\_radix\_register\_lower\_bound}
设有 \(b_2\) 条两点纤维与 \(b_3\) 条三点纤维，并要求在不预先固定纤维内部标号的前提下，同时记录下列独立选择：
\begin{enumerate}
  \item 对每条两点纤维指定一个 sheet parity；
  \item 对每条三点纤维指定一个非平凡对合。
\end{enumerate}
若某外加寄存器的状态数为 \(R\)，并足以无歧义地区分全部上述全局选择，则必有
\[
\boxed{\;
R\ge 2^{b_2}3^{b_3}.\;}
\]
等价地，若记所需二进制位数为
\[
r:=\lceil \log_2 R\rceil,
\]
则
\[
\boxed{\;
r\ge b_2+b_3\log_2 3.\;}
\]

进一步，将此一般结论应用于 boundary uplift，可得
\[
\boxed{\;
r_{\min}(6)\ge 3,\qquad
r_{\min}(7)\ge 5\log_2 3,\qquad
r_{\min}(8)\ge 8\log_2 3.\;}
\]
若还要求把 window-\(6\) 的三位 sheet parity 作为独立因子完整继承到 \(m=7,8\)，则更有
\[
\boxed{\;
r_{\mathrm{ext}}(7)\ge 3+5\log_2 3>10,\qquad
r_{\mathrm{ext}}(8)\ge 3+8\log_2 3>15.\;}
\]
从而
\[
\boxed{\;
r_{\mathrm{ext}}(7)\ge 11,\qquad
r_{\mathrm{ext}}(8)\ge 16.\;}
\]
\end{theorem}
\begin{proof}
两点纤维的无标号二元集上恰有两种 parity 选择；三点纤维上的非平凡对合等价于指定唯一不动点，故恰有三种可能。由于各纤维上的选择彼此独立，全部全局选择总数为
\[
2^{b_2}3^{b_3}.
\]
若寄存器状态数为 \(R\)，则至多只能区分 \(R\) 个全局选择；因而必要条件是
\[
R\ge 2^{b_2}3^{b_3}.
\]
取二进制对数便得
\[
\log_2 R\ge b_2+b_3\log_2 3,
\]
再取上整即得位数下界。

现应用到 boundary uplift。由定理 \ref{thm:xi-family-number-window-quantization-minimal-window}，
\[
|X_6^{\mathrm{bdry}}|=3,\qquad
|X_7^{\mathrm{bdry}}|=5,\qquad
|X_8^{\mathrm{bdry}}|=8.
\]
又由命题 \ref{prop:xi-window6-sheet-parity-anchor-m6}，
\[
\Delta_6^{\mathrm{bdry}}=\{0,34\},
\qquad
\Delta_7^{\mathrm{bdry}}=\{0,55,89\},
\qquad
\Delta_8^{\mathrm{bdry}}=\{0,89,144\}.
\]
因此 \(m=6\) 的每条 boundary 纤维都是二层，而 \(m=7,8\) 的每条 boundary 纤维都是三层，于是
\[
(b_2,b_3)=
\begin{cases}
(3,0),& m=6,\\
(0,5),& m=7,\\
(0,8),& m=8.
\end{cases}
\]
代入一般下界便得到
\[
r_{\min}(6)\ge 3,\qquad
r_{\min}(7)\ge 5\log_2 3,\qquad
r_{\min}(8)\ge 8\log_2 3.
\]

若还要求把 \(m=6\) 的三位 sheet parity 作为独立因子保留，则总状态数至少再乘以 \(2^3\)，故所需位数下界在 \(m=7,8\) 时分别增加 \(3\)。于是
\[
r_{\mathrm{ext}}(7)\ge 3+5\log_2 3,\qquad
r_{\mathrm{ext}}(8)\ge 3+8\log_2 3.
\]
最后由于
\[
3+5\log_2 3>10,
\qquad
3+8\log_2 3>15,
\]
取上整即得
\[
r_{\mathrm{ext}}(7)\ge 11,
\qquad
r_{\mathrm{ext}}(8)\ge 16.
\]
证毕。
\end{proof}

\begin{theorem}[boundary 中心三维子群的非特征性与自同构轨道]\label{thm:xi-window6-boundary-center-nonintrinsic-orbit}
设
\[
Z_{\partial}\cong (\ZZ/2\ZZ)^3_{\mathrm{bdry}}
\hookrightarrow
Z(G)\cong (\ZZ/2\ZZ)^8
\]
为定理 \ref{thm:xi-window6-center-exact-splitting-bdry-cyclic} 中的 boundary 中心子群。则：
\begin{enumerate}
  \item \(Z_{\partial}\) 不是 \(G\) 的特征子群；
  \item \(\Aut(G)\) 在 \(Z(G)\) 的全部三维 \(\FF_2\)-子空间上作用传递；
  \item \(Z_{\partial}\) 的自同构轨道大小恰为
  \[
  \boxed{\;
  \genfrac{[}{]}{0pt}{}{8}{3}_{2}
  =
  \frac{(2^8-1)(2^7-1)(2^6-1)}{(2^3-1)(2^2-1)(2-1)}
  =97155.\;}
  \]
\end{enumerate}
因此，从抽象群 \(G\) 本身无法恢复 boundary parity 的几何来源；它只是 \(8\) 维中心 \(\FF_2\)-向量空间中的某个三维子空间。
\end{theorem}
\begin{proof}
由定理 \ref{thm:xi-window6-full-automorphism-group-structure}，对任意
\[
\alpha\in \GL_8(\FF_2)\cong \Aut\bigl(Z(G)\bigr)
\]
都可定义自同构
\[
(a,h)\longmapsto (\alpha(a),h)
\qquad (a\in Z(G),\ h\in H).
\]
因此限制映射
\[
\Aut(G)\longrightarrow \Aut\bigl(Z(G)\bigr)\cong \GL_8(\FF_2)
\]
是满射。于是 \(\Aut(G)\) 在 \(Z(G)\cong \FF_2^8\) 的三维子空间集合上诱导出完整的 \(\GL_8(\FF_2)\)-作用。

而 \(\GL_8(\FF_2)\) 在 Grassmannian
\[
\mathrm{Gr}(3,8;\FF_2)
\]
上作用传递，故任意两个三维子空间都在 \(\Aut(G)\) 下共轭。特别地，\(Z_{\partial}\) 不是特征子群。

轨道大小即三维子空间总数，也就是 Gaussian 二项式系数
\[
\genfrac{[}{]}{0pt}{}{8}{3}_{2}
=
\frac{(2^8-1)(2^7-1)(2^6-1)}{(2^3-1)(2^2-1)(2-1)}
=97155.
\]
证毕。
\end{proof}

\begin{theorem}[boundary parity 的 \texorpdfstring{$2$}{2}-群相位扭转可全局延拓]\label{thm:xi-window6-boundary-parity-cocycle-extension}
\leanverified{paper\_xi\_window6\_boundary\_parity\_cocycle\_extension}
设
\[
W:=Z_{\partial}\cong (\ZZ/2\ZZ)^3\subset Z(G)\cong (\ZZ/2\ZZ)^8.
\]
则限制映射
\[
\mathrm{res}_W:\ H^2(G,\CC^\times)\longrightarrow H^2(W,\CC^\times)
\]
为满射。更精确地，
\[
H^2(W,\CC^\times)\cong (\ZZ/2\ZZ)^3,
\]
并且
\[
\ker(\mathrm{res}_W)\cong (\ZZ/2\ZZ)^{216}.
\]
因此，boundary parity 子群上的每一个中心 \(2\)-cocycle 都可以延拓为 \(G\) 上的全局射影相位。
\end{theorem}
\begin{proof}
记
\[
A:=Z(G)\cong (\ZZ/2\ZZ)^8.
\]
由于 \(W\) 是 \(A\) 的三维 \(\FF_2\)-子空间，可取补空间 \(W'\) 使得
\[
A=W\oplus W'.
\]
对任意 elementary abelian \(2\)-群 \(E\)，都有标准识别
\[
H^2(E,\CC^\times)\cong \Lambda^2(E^\vee).
\]
因此
\[
H^2(W,\CC^\times)\cong \Lambda^2(W^\vee)\cong (\ZZ/2\ZZ)^{\binom32}=(\ZZ/2\ZZ)^3.
\]

现取任意
\[
\eta\in H^2(W,\CC^\times).
\]
在上述识别下，\(\eta\) 对应于 \(W\) 上某个交替双线性形式。将其在分解
\[
A=W\oplus W'
\]
上按零延拓到 \(A\)，便得到某个
\[
\widetilde\eta\in H^2(A,\CC^\times)
\]
使得
\[
\widetilde\eta|_W=\eta.
\]

再考虑直积投影
\[
\pi_A:G=A\times H\longrightarrow A.
\]
由上同调的反变性，
\[
\pi_A^*:H^2(A,\CC^\times)\longrightarrow H^2(G,\CC^\times)
\]
给出提升类 \(\pi_A^*(\widetilde\eta)\)。其在 \(W\subset A\subset G\) 上的限制正好仍为 \(\eta\)。因此 \(\mathrm{res}_W\) 满射。

最后，由定理 \ref{thm:xi-window6-schur-multiplier-219}，
\[
\dim_{\FF_2}H^2(G,\CC^\times)=219,
\]
而
\[
\dim_{\FF_2}H^2(W,\CC^\times)=3.
\]
由于 \(\mathrm{res}_W\) 已证为满射，故
\[
\dim_{\FF_2}\ker(\mathrm{res}_W)=219-3=216.
\]
于是
\[
\ker(\mathrm{res}_W)\cong (\ZZ/2\ZZ)^{216}.
\]
证毕。
\end{proof}

\paragraph{window-\texorpdfstring{$6$}{6} 的根--权交会刚性、orbit--Levi 骨架与二层 boundary 种子的 \texorpdfstring{$\ZZ_6$}{Z6} 主齐性}
在四步 uplift 同构、boundary-\(10\) 的根--Cartan 四块分裂、window-\(6\) 的 \(C_6\) 轨道骨架，以及二层 boundary 种子的两层覆盖已经固定之后，weight-layer、orbit--Levi、sheet parity 与等变编码四条离散接口可进一步压缩为下列四条结构性结论。其共同特征是：\(9+6+3+3\) 的 weight-layer 分块与 \(6+6+6+3\) 的根--Cartan 分块之间只有唯一的五原子共同细化；唯一三元轨道被锁定于 \(A_2^{(-)}\)；二层 boundary 种子形成正则 \(\ZZ_6\)-torsor；而任何保持全部内生对称的外置编码都必须真实含有阶 \(6\) 元。

\begin{theorem}[window-\texorpdfstring{$6$}{6} 的 weight-layer 分块与 \texorpdfstring{$B_3/C_3$}{B3/C3} 根分块的唯一五原子共同细化]\label{thm:xi-time-part53da-root-weight-common-refinement-five-atoms}
\[
X_6=X_6^{\mathrm{cyc}}\sqcup X_6^{\mathrm{bdry}},
\qquad
|X_6|=21,\quad |X_6^{\mathrm{cyc}}|=18,\quad |X_6^{\mathrm{bdry}}|=3,
\]
且
\[
X_6^{\mathrm{bdry}}
=
\{\texttt{100001},\ \texttt{100101},\ \texttt{101001}\}.
\]
记表 \ref{tab:window6_patisalam_9633_partition} 给出的 weight-layer 刚性分块为
\[
X_6^{\mathrm{cyc}}
=
X_6^{(2)}\sqcup X_6^{(1)}\sqcup X_6^{(L)},
\qquad
|X_6^{(2)}|=9,\quad |X_6^{(1)}|=6,\quad |X_6^{(L)}|=3,
\]
并以 \(X_6^{(R)}:=X_6^{\mathrm{bdry}}\) 补成
\[
X_6=X_6^{(2)}\sqcup X_6^{(1)}\sqcup X_6^{(L)}\sqcup X_6^{(R)}.
\]
另一方面，沿用定理 \ref{thm:boundary-shift4-uplift-isomorphism} 的 uplift 同构
\[
\iota_6:X_6\xrightarrow{\sim}X_{10}^{\mathrm{bdry}},
\qquad
\iota_6(v)=10\,v\,01,
\]
\[
X_6
=
X_6^{\mathrm{short}}
\sqcup
X_6^{A_2(+)}
\sqcup
X_6^{A_2(-)}
\sqcup
X_6^{\mathrm{Cartan}},
\]
其中 \(X_6^{\mathrm{short}},X_6^{A_2(+)},X_6^{A_2(-)},X_6^{\mathrm{Cartan}}\) 分别是定理 \ref{thm:boundary-m10-b3c3-rigid-four-block-split} 的四块分裂在 \(\iota_6\) 下的原像，故
\[
21=6+6+6+3.
\]
则这两种分块的共同细化恰有且仅有五个非空原子，其交会矩阵唯一为
\[
\begin{array}{c|cccc}
& \mathrm{short} & A_2^{(+)} & A_2^{(-)} & \mathrm{Cartan}\\
\hline
X_6^{(2)} & 0 & 3 & 6 & 0\\
X_6^{(1)} & 6 & 0 & 0 & 0\\
X_6^{(L)} & 0 & 3 & 0 & 0\\
X_6^{(R)} & 0 & 0 & 0 & 3
\end{array}
\]
从而共同细化的五个原子大小严格为
\[
6\ \oplus\ 3\ \oplus\ 6\ \oplus\ 3\ \oplus\ 3.
\]
\end{theorem}
\begin{proof}
表 \ref{tab:window6_patisalam_9633_partition} 与表 \ref{tab:boundary_uplift_m10_dictionary} 允许直接逐词核对。

首先，六个 weight-\(1\) 词
\[
\texttt{100000},\ \texttt{010000},\ \texttt{001000},\ \texttt{000100},\ \texttt{000010},\ \texttt{000001}
\]
恰对应六条 short roots，因此
\[
X_6^{(1)}=X_6^{\mathrm{short}}.
\]

其次，
\[
X_6^{(L)}=\{\texttt{000000},\ \texttt{010101},\ \texttt{101010}\}
\]
在根字典中全部落入 \(A_2^{(+)}\)，于是
\[
X_6^{(L)}\subset X_6^{A_2(+)}
\quad\text{且}\quad
|X_6^{(L)}\cap X_6^{A_2(+)}|=3.
\]

再者，九个 weight-\(2\) 词中
\[
\texttt{010100},\ \texttt{100010},\ \texttt{010001}
\]
落入 \(A_2^{(+)}\)，其余六个
\[
\texttt{101000},\ \texttt{100100},\ \texttt{010010},\ \texttt{001010},\ \texttt{001001},\ \texttt{000101}
\]
落入 \(A_2^{(-)}\)。因此
\[
|X_6^{(2)}\cap X_6^{A_2(+)}|=3,
\qquad
|X_6^{(2)}\cap X_6^{A_2(-)}|=6.
\]

最后，boundary 三词
\[
\texttt{100001},\ \texttt{100101},\ \texttt{101001}
\]
在表 \ref{tab:boundary_uplift_m10_dictionary} 中全部对应 Cartan 方向，因此
\[
X_6^{(R)}=X_6^{\mathrm{Cartan}}
\quad\text{且}\quad
|X_6^{(R)}\cap X_6^{\mathrm{Cartan}}|=3.
\]

四列总数已经分别达到 \(6,6,6,3\)，故除此五个非空交块外不再可能有新的非空交会。于是交会矩阵被唯一锁定，共同细化恰由五个非空原子组成。证毕。
\end{proof}

\begin{theorem}[orbit--Levi 骨架与根--Cartan 分裂的交会矩阵完全唯一]\label{thm:xi-time-part53da-orbit-levi-root-cartan-intersection}
沿用表 \ref{tab:window6_c6_orbit_decomposition} 的 \(C_6\) 轨道分解，并记
\[
X_6^{\mathrm{cyc},3}:=\mathcal{O}(\texttt{001001})
=
\{\texttt{001001},\texttt{010010},\texttt{100100}\},
\qquad
X_6^{\mathrm{cyc},15}:=X_6^{\mathrm{cyc}}\setminus X_6^{\mathrm{cyc},3}.
\]
则在
\[
X_6=X_6^{\mathrm{cyc},15}\sqcup X_6^{\mathrm{cyc},3}\sqcup X_6^{\mathrm{bdry}}
\qquad (21=15\oplus 3\oplus 3)
\]
与定理 \ref{thm:xi-time-part53da-root-weight-common-refinement-five-atoms} 的根--Cartan 分块之间，交会矩阵唯一为
\[
\begin{array}{c|cccc}
& \mathrm{short} & A_2^{(+)} & A_2^{(-)} & \mathrm{Cartan}\\
\hline
X_6^{\mathrm{cyc},15} & 6 & 6 & 3 & 0\\
X_6^{\mathrm{cyc},3}  & 0 & 0 & 3 & 0\\
X_6^{\mathrm{bdry}}   & 0 & 0 & 0 & 3
\end{array}
\]
并且同时在 orbit--Levi 骨架与根--Cartan 分块上常值的函数空间维数恰为 \(4\)。
\end{theorem}
\begin{proof}
表 \ref{tab:window6_c6_orbit_decomposition} 表明，唯一的三元轨道正是
\[
X_6^{\mathrm{cyc},3}
=
\{\texttt{001001},\texttt{010010},\texttt{100100}\}.
\]
再由表 \ref{tab:boundary_uplift_m10_dictionary}，这三个词分别对应
\[
(0,1,-1),\qquad (1,0,-1),\qquad (-1,1,0),
\]
因而全部落在 \(A_2^{(-)}\)。故
\[
X_6^{\mathrm{cyc},3}\subset X_6^{A_2(-)},
\qquad
|X_6^{\mathrm{cyc},3}\cap X_6^{A_2(-)}|=3.
\]

又 \(X_6^{\mathrm{bdry}}\) 的三个词在同一字典中全部对应 \(\mathrm{Cartan}\) 方向，因此
\[
|X_6^{\mathrm{bdry}}\cap X_6^{\mathrm{Cartan}}|=3.
\]

由定理 \ref{thm:xi-time-part53da-root-weight-common-refinement-five-atoms}，整个 cyclic 扇区与根--Cartan 字典的交会总数为
\[
6\ \mathrm{short}\ +\ 6\ A_2^{(+)}\ +\ 6\ A_2^{(-)}.
\]
从中扣除 \(X_6^{\mathrm{cyc},3}\) 已经占去的 \(3\) 个 \(A_2^{(-)}\) 方向，剩余的 \(X_6^{\mathrm{cyc},15}\) 便必然恰好分布为
\[
6\ \mathrm{short}\ +\ 6\ A_2^{(+)}\ +\ 3\ A_2^{(-)}.
\]
于是第一行也被唯一锁定，得到所示交会矩阵。

最后，一个函数同时在 orbit--Levi 三块与根--Cartan 四块上常值，当且仅当它在全部非空交块上常值。由于这里只有四个非空交块，这样的函数空间便恰为四维。特别地，唯一三元轨道被根字典强制置于 \(A_2^{(-)}\)，而不能落入 \(A_2^{(+)}\)。证毕。
\end{proof}

\begin{theorem}[window-\texorpdfstring{$6$}{6} 的二层 boundary 种子是正则 \texorpdfstring{$\ZZ_6$}{Z6}-torsor]\label{thm:xi-time-part53da-boundary-z6-torsor}
设
\[
X_6^{\mathrm{bdry}}
=
\{\texttt{100001},\ \texttt{100101},\ \texttt{101001}\},
\]
并沿用表 \ref{tab:fold6_boundary_sheet_pairs} 的 dyadic uplift 词典：每个 boundary 词在 \(\Fold_6^{\mathrm{bin}}\) 下恰有两个原像，且同一纤维内两点的差值恒为 \(34=F_9\)。记
\[
\widetilde X_6^{\mathrm{bdry}}
:=
\bigsqcup_{w\in X_6^{\mathrm{bdry}}}
\bigl(\Fold_6^{\mathrm{bin}}\bigr)^{-1}(w),
\qquad
|\widetilde X_6^{\mathrm{bdry}}|=6.
\]
则 \(\widetilde X_6^{\mathrm{bdry}}\) 上存在一个自由且传递的 \(\ZZ_6\)-作用。特别地，
\[
\CC[\widetilde X_6^{\mathrm{bdry}}]\cong \CC[\ZZ_6].
\]
作为复置换模，\(\CC[\widetilde X_6^{\mathrm{bdry}}]\) 恰为 \(\ZZ_6\) 的正则表示。
\end{theorem}
\begin{proof}
表 \ref{tab:fold6_boundary_sheet_pairs} 表明，\(\widetilde X_6^{\mathrm{bdry}}\) 由三条两点纤维组成，因此确有
\[
|\widetilde X_6^{\mathrm{bdry}}|=6.
\]
记
\[
\zeta:\widetilde X_6^{\mathrm{bdry}}\to \widetilde X_6^{\mathrm{bdry}}
\]
为每条两点纤维上的对角 sheet-flip。另一方面，三条 boundary 轴的循环置换给出一个 \(C_3\)-作用于基底三点集 \(X_6^{\mathrm{bdry}}\)，并自然提升为 \(\widetilde X_6^{\mathrm{bdry}}\) 上的置换作用。

由于 \(\zeta\) 只翻转 sheet 而不改变 boundary 轴，\(C_3\) 只循环 boundary 轴而不改变 sheet，二者可交换，故得到直积作用
\[
\langle\zeta\rangle\times C_3\curvearrowright \widetilde X_6^{\mathrm{bdry}}.
\]
再由
\[
\langle\zeta\rangle\times C_3
\cong
\ZZ_2\times \ZZ_3
\cong
\ZZ_6
\]
便可把该作用视为 \(\ZZ_6\)-作用。

该作用是传递的：给定任意两点 \(x,y\in \widetilde X_6^{\mathrm{bdry}}\)，先用 \(C_3\) 把 \(x\) 所在纤维送到 \(y\) 所在纤维；若此时 sheet 标号不同，再施加一次 \(\zeta\) 即可把像点送到 \(y\)。

另一方面，
\[
\bigl|\langle\zeta\rangle\times C_3\bigr|=2\cdot 3=6=|\widetilde X_6^{\mathrm{bdry}}|.
\]
有限群在同基数集合上的传递作用自动自由，故 \(\widetilde X_6^{\mathrm{bdry}}\) 是一个 \(\ZZ_6\)-torsor。最后，任一有限群在自身 torsor 上的置换表示都与其正则表示同构，因此
\[
\CC[\widetilde X_6^{\mathrm{bdry}}]\cong \CC[\ZZ_6].
\]
证毕。
\end{proof}

\begin{corollary}[保持全部内生对称的 boundary 编码必须含有阶六元]\label{cor:xi-time-part53da-equivariant-boundary-coding-order-six}
沿用定理 \ref{thm:xi-time-part53da-boundary-z6-torsor} 的 \(\ZZ_6\)-作用。设 \(A\) 为有限阿贝尔群，则下述两件事等价：
\begin{enumerate}
  \item 存在单射 \(\Phi:\widetilde X_6^{\mathrm{bdry}}\to A\) 与群同态 \(\rho:\ZZ_6\to A\)，使得
  \[
  \Phi(g\cdot x)=\rho(g)+\Phi(x)
  \qquad
  (g\in \ZZ_6,\ x\in \widetilde X_6^{\mathrm{bdry}});
  \]
  \item \(A\) 含有一个循环子群同构于 \(\ZZ_6\)，等价地，\(A\) 含有一个阶 \(6\) 元。
\end{enumerate}
特别地：
\begin{enumerate}
  \item 对任意 \(r\ge 1\)，纯二进寄存器 \(A\cong (\ZZ/2\ZZ)^r\) 不可能承载这样的等变编码；
  \item 对任意 \(r\ge 1\)，纯三进寄存器 \(A\cong (\ZZ/3\ZZ)^r\) 也不可能承载这样的等变编码；
  \item 在全部有限阿贝尔目标中，最小可行者恰为 \(\ZZ_6\cong \ZZ_2\times \ZZ_3\)。
\end{enumerate}
\end{corollary}
\begin{proof}
把定理 \ref{thm:xi-time-part53da-boundary-z6-torsor} 中的作用群识别为 \(\ZZ_6\)。若存在这样的等变单射 \(\Phi\)，则必存在群同态
\[
\rho:\ZZ_6\to A
\]
使得
\[
\Phi(g\cdot x)=\rho(g)+\Phi(x)
\qquad
(g\in \ZZ_6,\ x\in \widetilde X_6^{\mathrm{bdry}}).
\]
固定任意 \(x_0\in \widetilde X_6^{\mathrm{bdry}}\)。由于 \(\widetilde X_6^{\mathrm{bdry}}\) 是 \(\ZZ_6\)-torsor，其轨道恰有 \(6\) 个点，故单射 \(\Phi\) 迫使
\[
\{\Phi(g\cdot x_0):\ g\in \ZZ_6\}
\]
也有 \(6\) 个不同元素。由等变性，
\[
\Phi(g\cdot x_0)=\rho(g)+\Phi(x_0),
\]
故该像轨道正是子群 \(\rho(\ZZ_6)\) 的一个平移副本；其大小只能是 \(|\rho(\ZZ_6)|\)。于是
\[
|\rho(\ZZ_6)|=6,
\]
从而 \(\rho(1)\in A\) 的阶恰为 \(6\)。

反之，若 \(A\) 含有一个阶 \(6\) 元 \(a\)，则取同态
\[
\rho:\ZZ_6\to A,\qquad \rho(1)=a.
\]
固定基点 \(x_0\in \widetilde X_6^{\mathrm{bdry}}\)，并对任意 \(g\in \ZZ_6\) 定义
\[
\Phi(g\cdot x_0):=\rho(g).
\]
由于 \(\widetilde X_6^{\mathrm{bdry}}\) 上的 \(\ZZ_6\)-作用自由且传递，该定义良好并且自动等变；又因 \(a\) 的阶为 \(6\)，\(\rho\) 为单射，故 \(\Phi\) 亦为单射。

最后，\((\ZZ/2\ZZ)^r\) 的每个元素阶都整除 \(2\)，\((\ZZ/3\ZZ)^r\) 的每个元素阶都整除 \(3\)，因此都不含阶 \(6\) 元；而 \(\ZZ_6\cong\ZZ_2\times\ZZ_3\) 显然含有阶 \(6\) 元，且阶数最小。证毕。
\end{proof}

\noindent
由此，window-\(6\) 的 weight-layer 分块、orbit--Levi 骨架、根--Cartan 字典与二层 boundary 覆盖并非彼此松散兼容的四组离散证书，而是由唯一交会矩阵与正则 \(\ZZ_6\)-torsor 共同锁定的一组刚性接口。于是，定理 \ref{thm:xi-boundary-uplift-mixed-radix-register-lower-bound} 给出的 mixed-radix 状态数下界，在这一最小窗口上可进一步强化为严格的群论判据：任何保留全部 boundary 内生对称的目标寄存器都必须真实含有阶 \(6\) 元。

\paragraph{bin-fold 临界尾谱、连续容量双折线与复温零点晶格}
在附录 \ref{app:fold-multiplicity} 已经给出 bin-fold 的两态一致渐近、容量--尾函数等价、Mellin--Stieltjes 双对偶以及 Rényi 速率塌缩之后，还可把临界阈值附近的尾谱几何与复温度解析结构进一步压缩为下列四条结果。

\begin{theorem}[bin-fold 的临界尾谱呈现三阶梯刚性]\label{thm:xi-foldbin-critical-tail-three-step-rigidity}
\leanverified{paper\_xi\_foldbin\_critical\_tail\_three\_step\_rigidity}
记
\[
\alpha_{\mathrm{bin}}:=\log\frac{2}{\varphi},
\qquad
T_m(s):=\exp(\alpha_{\mathrm{bin}}m+s)=\Bigl(\frac{2}{\varphi}\Bigr)^m e^s,
\qquad s\in\RR,
\]
并定义 bin-fold 尾计数
\[
N_m^{\mathrm{bin}}(t):=\card{\left\{w\in X_m:\ d_m^{\mathrm{bin}}(w)\ge t\right\}}.
\]
则对任意固定 \(s\in\RR\)，当 \(m\to\infty\) 时有如下严格三阶梯律：若 \(m\) 充分大，则
\[
\boxed{
N_m^{\mathrm{bin}}(T_m(s))
=
\begin{cases}
F_{m+2}, & s<-\log\varphi,\\[2mm]
F_{m+1}, & -\log\varphi<s<0,\\[2mm]
0, & s>0.
\end{cases}}
\]
从而归一化尾比例满足
\[
\boxed{
\frac{N_m^{\mathrm{bin}}(T_m(s))}{|X_m|}
\longrightarrow
\begin{cases}
1, & s<-\log\varphi,\\[2mm]
\varphi^{-1}, & -\log\varphi<s<0,\\[2mm]
0, & s>0.
\end{cases}}
\]
\end{theorem}
\begin{proof}
记
\[
c_m:=\Bigl(\frac{2}{\varphi}\Bigr)^m,
\qquad
\varepsilon_m:=\Bigl(\frac{\varphi}{2}\Bigr)^m.
\]
由定理 \ref{thm:fold-bin-two-state-asymptotic}，存在常数 \(C>0\) 使对充分大 \(m\) 与任意 \(w\in X_m\) 都有
\[
d_m^{\mathrm{bin}}(w)=c_m\varphi^{-w_m}\bigl(1+\delta_{m,w}\bigr),
\qquad
|\delta_{m,w}|\le C\varepsilon_m.
\]
由于 \(s\) 固定且不位于两个临界值 \(-\log\varphi\) 与 \(0\)，三种情形分别讨论。

若 \(s<-\log\varphi\)，则 \(e^s<\varphi^{-1}\)。取 \(\eta>0\) 使
\[
e^s<\varphi^{-1}(1-2\eta).
\]
当 \(m\) 充分大时，\(C\varepsilon_m\le \eta\)，于是
\[
d_m^{\mathrm{bin}}(w)\ge c_m\varphi^{-w_m}(1-\eta)\ge c_m\varphi^{-1}(1-\eta)>c_m e^s=T_m(s)
\]
对全部 \(w\in X_m\) 成立。因此 \(N_m^{\mathrm{bin}}(T_m(s))=|X_m|=F_{m+2}\)。

若 \(-\log\varphi<s<0\)，则 \(\varphi^{-1}<e^s<1\)。取 \(\eta>0\) 使
\[
\varphi^{-1}(1+2\eta)<e^s<1-2\eta.
\]
当 \(m\) 充分大时，\(C\varepsilon_m\le \eta/2\)。于是对末位 \(w_m=0\) 的词有
\[
d_m^{\mathrm{bin}}(w)\ge c_m(1-\eta/2)>c_m e^s=T_m(s),
\]
而对末位 \(w_m=1\) 的词有
\[
d_m^{\mathrm{bin}}(w)\le c_m\varphi^{-1}(1+\eta/2)<c_m e^s=T_m(s).
\]
故此时恰有末位为 \(0\) 的稳定词落在阈值之上。由 Fibonacci 递推下的标准计数，
\[
\card{\{w\in X_m:w_m=0\}}=F_{m+1},
\qquad
\card{\{w\in X_m:w_m=1\}}=F_m,
\]
遂得 \(N_m^{\mathrm{bin}}(T_m(s))=F_{m+1}\)。

若 \(s>0\)，则 \(e^s>1\)。取 \(\eta>0\) 使 \(1+\eta<e^s\)。当 \(m\) 充分大时，\(C\varepsilon_m\le \eta/2\)，故
\[
d_m^{\mathrm{bin}}(w)\le c_m(1+\eta/2)<c_m e^s=T_m(s)
\]
对全部 \(w\in X_m\) 成立，从而 \(N_m^{\mathrm{bin}}(T_m(s))=0\)。

最后，由 \(|X_m|=F_{m+2}\) 与 \(F_{m+1}/F_{m+2}\to \varphi^{-1}\) 立得比例极限。证毕。
\end{proof}

\begin{theorem}[临界连续容量的精确双折线极限]\label{thm:xi-foldbin-critical-capacity-broken-line}
\leanverified{paper\_xi\_foldbin\_critical\_capacity\_broken\_line}
沿用命题 \ref{prop:fold-capacity-tail-equivalence} 的连续容量记号，并在 bin-fold 口径下定义
\[
\mathcal C_m^{\mathrm{bin,cont}}(T):=\sum_{w\in X_m}\min\!\bigl(d_m^{\mathrm{bin}}(w),T\bigr).
\]
则对任意固定 \(s\in\RR\)，沿临界尺度 \(T_m(s)=\bigl(\tfrac{2}{\varphi}\bigr)^m e^s\) 有极限
\[
\boxed{
\lim_{m\to\infty}\frac{\mathcal C_m^{\mathrm{bin,cont}}(T_m(s))}{2^m}
=
\mathcal C_{\mathrm{bin}}(s):=
\begin{cases}
\dfrac{\varphi^2}{\sqrt5}e^s, & s<-\log\varphi,\\[3mm]
\dfrac{\varphi}{\sqrt5}e^s+\dfrac{1}{\varphi\sqrt5}, & -\log\varphi<s<0,\\[3mm]
1, & s>0.
\end{cases}}
\]
并且 \(\mathcal C_{\mathrm{bin}}\) 在 \(s=-\log\varphi\) 与 \(s=0\) 处连续，但其导数在这两点均发生跃迁。因此临界窗口在极限中形成一个具有两处折点的双折线轮廓。
\end{theorem}
\begin{proof}
记
\[
c_m:=\Bigl(\frac{2}{\varphi}\Bigr)^m,
\qquad
\varepsilon_m:=\Bigl(\frac{\varphi}{2}\Bigr)^m.
\]
由定理 \ref{thm:xi-foldbin-critical-tail-three-step-rigidity} 的分类与定理 \ref{thm:fold-bin-two-state-asymptotic} 的一致展开，可按三区间逐段计算。

若 \(s<-\log\varphi\)，则对充分大 \(m\) 有 \(T_m(s)\) 低于全部纤维多重度，故
\[
\mathcal C_m^{\mathrm{bin,cont}}(T_m(s))
=
T_m(s)\,|X_m|
=
c_m e^s F_{m+2}.
\]
除以 \(2^m\) 并用 \(F_{m+2}\sim \varphi^{m+2}/\sqrt5\)，得到
\[
\frac{\mathcal C_m^{\mathrm{bin,cont}}(T_m(s))}{2^m}
=
\frac{F_{m+2}}{\varphi^m}e^s
\longrightarrow
\frac{\varphi^2}{\sqrt5}e^s.
\]

若 \(-\log\varphi<s<0\)，则对充分大 \(m\) 恰有末位 \(w_m=0\) 的纤维处在线性段，而末位 \(w_m=1\) 的纤维已经饱和。因此
\[
\mathcal C_m^{\mathrm{bin,cont}}(T_m(s))
=
F_{m+1}T_m(s)+\sum_{w_m=1} d_m^{\mathrm{bin}}(w).
\]
再由定理 \ref{thm:fold-bin-two-state-asymptotic} 的一致展开，
\[
\sum_{w_m=1} d_m^{\mathrm{bin}}(w)
=
F_m\,c_m\varphi^{-1}\bigl(1+O(\varepsilon_m)\bigr).
\]
故
\[
\mathcal C_m^{\mathrm{bin,cont}}(T_m(s))
=
F_{m+1}c_m e^s+F_m c_m\varphi^{-1}\bigl(1+O(\varepsilon_m)\bigr).
\]
两端除以 \(2^m\) 即得
\[
\frac{\mathcal C_m^{\mathrm{bin,cont}}(T_m(s))}{2^m}
=
\frac{F_{m+1}}{\varphi^m}e^s+\frac{F_m}{\varphi^{m+1}}+o(1)
\longrightarrow
\frac{\varphi}{\sqrt5}e^s+\frac{1}{\varphi\sqrt5}.
\]

若 \(s>0\)，则对充分大 \(m\) 有 \(T_m(s)\) 高于全部纤维多重度，故
\[
\mathcal C_m^{\mathrm{bin,cont}}(T_m(s))
=
\sum_{w\in X_m} d_m^{\mathrm{bin}}(w)
=
2^m,
\]
从而极限为 \(1\)。

连续性只需在两个连接点代入检查：
\[
\frac{\varphi^2}{\sqrt5}e^{-\log\varphi}
=
\frac{\varphi}{\sqrt5}
=
\frac{\varphi}{\sqrt5}\varphi^{-1}+\frac{1}{\varphi\sqrt5},
\]
\[
\frac{\varphi}{\sqrt5}e^0+\frac{1}{\varphi\sqrt5}
=
\frac{\varphi+\varphi^{-1}}{\sqrt5}
=
1.
\]
另一方面，
\[
\mathcal C_{\mathrm{bin}}'(s)
=
\begin{cases}
\dfrac{\varphi^2}{\sqrt5}e^s, & s<-\log\varphi,\\[2mm]
\dfrac{\varphi}{\sqrt5}e^s, & -\log\varphi<s<0,\\[2mm]
0, & s>0,
\end{cases}
\]
故在 \(s=-\log\varphi\) 与 \(s=0\) 处导数均不连续。证毕。
\end{proof}

\begin{theorem}[bin-fold 幂和配分函数的复温度零点在 \texorpdfstring{$\Re q=-1$}{Re q=-1} 上形成算术晶格]\label{thm:xi-foldbin-complex-temperature-zero-lattice}
\leanverified{paper\_xi\_foldbin\_complex\_temperature\_zero\_lattice}
对每个固定 \(m\)，把
\[
q\longmapsto S_q^{\mathrm{bin}}(m):=\sum_{w\in X_m}\bigl(d_m^{\mathrm{bin}}(w)\bigr)^q
\]
视为复变量 \(q\) 的整函数。定义近似零点
\[
q_{m,k}^{(0)}
:=
-\frac{\log(F_{m+1}/F_m)}{\log\varphi}
+\frac{(2k+1)\pi i}{\log\varphi},
\qquad k\in\ZZ.
\]
则对每个固定整数 \(k\)，存在唯一真零点 \(q_{m,k}\) 满足
\[
\boxed{\;
q_{m,k}=q_{m,k}^{(0)}+O\!\left(\Bigl(\frac{\varphi}{2}\Bigr)^m\right)
\qquad (m\to\infty).\;}
\]
特别地，
\[
\boxed{\;
\Re q_{m,k}\to -1,
\qquad
\Im q_{m,k}=\frac{(2k+1)\pi}{\log\varphi}+o(1).\;}
\]
因此，bin-fold 配分函数的复温度零点在极限中沿垂直直线
\[
\boxed{\ \Re q=-1\ }
\]
形成等间距的 Lee--Yang 型算术晶格。
\end{theorem}
\begin{proof}
记
\[
c_m:=\Bigl(\frac{2}{\varphi}\Bigr)^m,
\qquad
\varepsilon_m:=\Bigl(\frac{\varphi}{2}\Bigr)^m.
\]
固定 \(k\in\ZZ\)，并取围绕极限点
\[
q_k^\ast:=-1+\frac{(2k+1)\pi i}{\log\varphi}
\]
的闭圆盘
\[
K_k:=\{q\in\CC:\ |q-q_k^\ast|\le 1\}.
\]
由定理 \ref{thm:fold-bin-two-state-asymptotic}，存在常数 \(C>0\) 与实数 \(\delta_{m,w}\) 使
\[
d_m^{\mathrm{bin}}(w)=c_m\varphi^{-w_m}(1+\delta_{m,w}),
\qquad
|\delta_{m,w}|\le C\varepsilon_m
\]
对充分大 \(m\) 与全部 \(w\in X_m\) 成立。由于 \(K_k\) 有界且 \(|\delta_{m,w}|<1/2\) 对充分大 \(m\) 成立，复指数函数的解析展开给出一致估计
\[
\bigl(1+\delta_{m,w}\bigr)^q=1+O_k(\varepsilon_m)
\qquad (q\in K_k).
\]
于是
\[
\bigl(d_m^{\mathrm{bin}}(w)\bigr)^q
=
c_m^q\varphi^{-q w_m}\bigl(1+O_k(\varepsilon_m)\bigr)
\qquad (q\in K_k),
\]
并且误差在 \(w\in X_m\) 上一致。按末位拆分求和，得到
\[
S_q^{\mathrm{bin}}(m)
=
c_m^q\bigl(B_m(q)+E_m(q)\bigr),
\]
其中
\[
B_m(q):=F_{m+1}+\varphi^{-q}F_m,
\qquad
\sup_{q\in K_k}|E_m(q)|\le C_k F_{m+1}\varepsilon_m
\]
对某个与 \(k\) 有关但与 \(m\) 无关的常数 \(C_k>0\) 成立。

由于 \(c_m^q\neq 0\)，故 \(S_q^{\mathrm{bin}}(m)\) 的零点与 \(B_m(q)+E_m(q)\) 的零点一致。近似零点由 \(B_m(q)=0\) 给出，即
\[
\varphi^{-q}=-\frac{F_{m+1}}{F_m},
\]
故其解恰为 \(q_{m,\ell}^{(0)}\)（\(\ell\in\ZZ\)）。在 \(q=q_{m,k}^{(0)}\) 处，
\[
B_m'(q)=-(\log\varphi)\varphi^{-q}F_m,
\qquad
B_m'\!\bigl(q_{m,k}^{(0)}\bigr)=(\log\varphi)F_{m+1}\neq 0,
\]
因此 \(q_{m,k}^{(0)}\) 是 \(B_m\) 的单根。

取圆周
\[
\Gamma_{m,k}:=\{q:\ |q-q_{m,k}^{(0)}|=R\varepsilon_m\},
\]
其中 \(R>0\) 为待定常数。由于 \(B_m\) 在 \(q_{m,k}^{(0)}\) 处为单根，Taylor 展开给出
\[
B_m(q)=B_m'\!\bigl(q_{m,k}^{(0)}\bigr)(q-q_{m,k}^{(0)})+O_k\!\bigl(F_{m+1}|q-q_{m,k}^{(0)}|^2\bigr).
\]
因此对充分大 \(m\) 与 \(q\in \Gamma_{m,k}\)，有
\[
|B_m(q)|\ge c_k R F_{m+1}\varepsilon_m
\]
其中 \(c_k>0\) 与 \(k\) 有关但与 \(m\) 无关。另一方面，
\[
|E_m(q)|\le C_k F_{m+1}\varepsilon_m.
\]
取 \(R>2C_k/c_k\)，即可保证在 \(\Gamma_{m,k}\) 上 \(|E_m(q)|<|B_m(q)|\)。由 Rouch\'e 定理，\(B_m+E_m\) 与 \(B_m\) 在 \(\Gamma_{m,k}\) 围成的圆盘内具有同样数目的零点。由于 \(B_m\) 在该圆盘内恰有一枚零点 \(q_{m,k}^{(0)}\)，故 \(S_q^{\mathrm{bin}}(m)\) 在其中亦恰有一枚零点 \(q_{m,k}\)，并满足
\[
q_{m,k}-q_{m,k}^{(0)}=O(\varepsilon_m).
\]

最后，由 \(F_{m+1}/F_m\to \varphi\) 得
\[
-\frac{\log(F_{m+1}/F_m)}{\log\varphi}\to -1,
\]
从而
\[
\Re q_{m,k}\to -1,
\qquad
\Im q_{m,k}=\frac{(2k+1)\pi}{\log\varphi}+o(1).
\]
证毕。
\end{proof}

\begin{corollary}[正温半平面无奇性：冻结不会进入 \texorpdfstring{$\Re q>-1$}{Re q>-1}]\label{cor:xi-foldbin-positive-halfplane-no-freezing}
\leanverified{paper\_xi\_foldbin\_positive\_halfplane\_no\_freezing}
对任意紧集 \(K\Subset\{\Re q>-1\}\)，当 \(m\) 充分大时，\(S_q^{\mathrm{bin}}(m)\) 在 \(K\) 上无零点；因此存在与正实轴一致的全纯对数支 \(\log_K S_q^{\mathrm{bin}}(m)\)。并且
\[
\frac{1}{m}\log_K S_q^{\mathrm{bin}}(m)
\longrightarrow
\tau_{\mathrm{bin}}(q):=
q\log\frac{2}{\varphi}+\log\varphi
\]
在 \(K\) 上一致成立。
特别地，\(\tau_{\mathrm{bin}}\) 在整个半平面 \(\Re q>-1\) 上全纯；因而对全部实参数 \(q>0\)，bin-fold 都不存在冻结相变。
\end{corollary}
\begin{proof}
记
\[
\varepsilon_m:=\Bigl(\frac{\varphi}{2}\Bigr)^m.
\]
固定紧集 \(K\Subset\{\Re q>-1\}\)。由定理 \ref{thm:xi-foldbin-complex-temperature-zero-lattice}，\(S_q^{\mathrm{bin}}(m)\) 的全部零点都位于 \(\Re q=-1+O(\varepsilon_m)\) 的垂直晶格附近。由于 \(K\) 与直线 \(\Re q=-1\) 具有正距离，故对充分大 \(m\) 而言，\(S_q^{\mathrm{bin}}(m)\) 在 \(K\) 上无零点，从而存在全纯对数支 \(\log_K S_q^{\mathrm{bin}}(m)\)。

再由上一条定理证明中的复参数估计，
\[
S_q^{\mathrm{bin}}(m)
=
\Bigl(\frac{2}{\varphi}\Bigr)^{qm}
\bigl(F_{m+1}+\varphi^{-q}F_m\bigr)\bigl(1+O_K(\varepsilon_m)\bigr)
\]
在 \(K\) 上一致成立。另一方面，由 Binet 公式，
\[
F_{m+1}+\varphi^{-q}F_m
=
\frac{\varphi^m}{\sqrt5}\Bigl(\varphi+\varphi^{-q}+O(\varphi^{-2m})\Bigr)
\]
在 \(K\) 上一致成立。由于 \(K\Subset\{\Re q>-1\}\)，存在 \(\delta>0\) 使 \(\Re q\ge -1+\delta\) 对全部 \(q\in K\) 成立，故
\[
|\varphi^{-q}|\le \varphi^{1-\delta}<\varphi,
\]
于是 \(\varphi+\varphi^{-q}\) 在 \(K\) 上一致远离零点。因而
\[
\frac{1}{m}\log_K\!\bigl(F_{m+1}+\varphi^{-q}F_m\bigr)
=
\log\varphi+O_K(m^{-1}),
\]
并且
\[
\frac{1}{m}\log_K\!\bigl(1+O_K(\varepsilon_m)\bigr)=o(1)
\]
在 \(K\) 上一致成立。

综合得到
\[
\frac{1}{m}\log_K S_q^{\mathrm{bin}}(m)
=
q\log\frac{2}{\varphi}
+\log\varphi+o(1)
\]
在 \(K\) 上一致成立。右端显然在 \(\Re q>-1\) 上全纯，于是 \(\tau_{\mathrm{bin}}(q)=q\log(2/\varphi)+\log\varphi\) 即为所求极限函数。对实 \(q>0\) 自然无任何奇性；故所谓 Rényi 速率塌缩并不是冻结后的线性支，而是半平面 \(\Re q>-1\) 上先验存在的复解析退化。证毕。
\end{proof}

\begin{remark}[临界尾谱与复温零点的数值锚点]\label{rem:xi-foldbin-critical-tail-zero-lattice-audit}
脚本 \texttt{scripts/exp\_fold\_bin\_critical\_tail\_zero\_lattice\_audit.py} 对阈值 \(s=-0.75,-0.25,0.20\) 的尾谱与容量值、以及 \(k=0\) 分支的复温零点偏移给出如下核验片段：
\begin{equation}\label{eq_fold_bin_critical_tail_zero_lattice_audit}
\begin{aligned}
&s=-0.75:\ \frac{N_{24}^{\mathrm{bin}}(T_{24}(-0.75))}{|X_{24}|}=1\to 1,\quad 2^{-24}\mathcal C_{24}^{\mathrm{bin,cont}}(T_{24}(-0.75))=0.553056\to 0.553056\\
&s=-0.25:\ \frac{N_{24}^{\mathrm{bin}}(T_{24}(-0.25))}{|X_{24}|}=0.618034\to 0.618034,\quad 2^{-24}\mathcal C_{24}^{\mathrm{bin,cont}}(T_{24}(-0.25))=0.83992\to 0.839939\\
&s=0.2:\ \frac{N_{24}^{\mathrm{bin}}(T_{24}(0.2))}{|X_{24}|}=0\to 0,\quad 2^{-24}\mathcal C_{24}^{\mathrm{bin,cont}}(T_{24}(0.2))=1\to 1\\
&|q_{20,0}-q_{20,0}^{(0)}|=0.071626,\ |q_{24,0}-q_{24,0}^{(0)}|=0.001213\\
&q_{24,0}^{(0)}=-1+6.528503\,\mathrm{i},\qquad q_{24,0}=-0.999567+6.52737\,\mathrm{i},\qquad |S_{q_{24,0}}^{\mathrm{bin}}(24)|=1.00099\times 10^{-87}
\end{aligned}
\end{equation}
\end{remark}

\paragraph{bin-fold 临界窗的复温读法}
上述四条结果表明：在临界尺度 \(T_m(s)\) 附近，bin-fold 尾谱并不形成连续展宽的肩部，而是被末位两态律压缩为严格三阶梯；相应连续容量则呈现具有两处折点的双折线极限。与此同时，幂和配分函数的复温零点并不进入正温半平面，而是沿 \(\Re q=-1\) 的算术晶格聚集。因此，正文中全部 Rényi 速率塌缩到 \(\log\varphi\) 的事实，应被理解为一条复解析刚性，而非冻结之后才出现的线性支。

\paragraph{bin-fold 欠分辨曲率、fold-gauge 的 even-\texorpdfstring{$\zeta$}{zeta} 塔与素数寄存器极小理想}
在 bin-fold 的两态一致渐近、欠分辨尾谱/容量极限、fold-gauge 常数的 Stirling--Bernoulli 层级以及素数寄存器估值算术化已经建立之后，还可进一步压缩出下列五条结构性后果。

\begin{theorem}[bin-fold 欠分辨极限由二原子曲率测度完全决定]\label{thm:xi-foldbin-underresolution-two-atomic-curvature-measure}
\leanverified{paper\_xi\_foldbin\_underresolution\_two\_atomic\_curvature\_measure}
沿用命题 \ref{prop:fold-bin-degeneracy-tail-capacity-kinks} 的记号，定义
\[
\widehat N_m(s):=\frac{1}{|X_m|}\,N_m^{\mathrm{bin}}\!\left(s\left(\frac{2}{\varphi}\right)^m\right),
\qquad
\widehat C_m(s):=2^{-m}\,\mathcal C_m^{\mathrm{bin,cont}}\!\left(s\left(\frac{2}{\varphi}\right)^m\right),
\qquad s>0.
\]
则当 \(s\notin\{\varphi^{-1},1\}\) 时，
\[
\widehat N_m(s)\longrightarrow \widehat N_\infty(s):=
\begin{cases}
1,& 0<s<\varphi^{-1},\\[2pt]
\varphi^{-1},& \varphi^{-1}<s<1,\\[2pt]
0,& s>1,
\end{cases}
\]
并且 \(\widehat C_m\) 在每个紧区间上一致收敛到
\[
\widehat C_\infty(s):=\frac{\varphi}{\sqrt5}\min\{1,s\}+\frac{1}{\sqrt5}\min\{\varphi^{-1},s\}.
\]
记 \(\nu_\infty:=-D^2\widehat C_\infty\) 为 \(\widehat C_\infty\) 的分布二阶导数所对应的曲率测度，则
\[
\boxed{\;
\nu_\infty=
\frac{1}{\sqrt5}\,\delta_{\varphi^{-1}}+\frac{\varphi}{\sqrt5}\,\delta_{1}.
\;}
\]
并且对每个 \(s>0\) 有
\[
\boxed{\;
\widehat C_\infty(s)=\int_0^s \nu_\infty([t,\infty))\,\dd t,
\;}
\]
而对 \(s\notin\{\varphi^{-1},1\}\) 有
\[
\boxed{\;
\widehat N_\infty(s)=\frac{\sqrt5}{\varphi^2}\,\nu_\infty([s,\infty)).
\;}
\]
因此，bin-fold 的欠分辨极限并非仅表现为两处折点的分段线性图像，而严格等价于一条支撑在 \(\{\varphi^{-1},1\}\) 上的二原子曲率测度；折点位置与原子质量共同构成其完备证书。
\end{theorem}
\begin{proof}
命题 \ref{prop:fold-bin-degeneracy-tail-capacity-kinks} 已给出 \(\widehat C_m\to \widehat C_\infty\) 的紧集上一致收敛，以及 \(\widehat N_m\to \widehat N_\infty\) 的点态极限。对 \(\widehat C_\infty\) 直接分段求导可得
\[
\widehat C_\infty'(s)=
\begin{cases}
\dfrac{\varphi+1}{\sqrt5}=\dfrac{\varphi^2}{\sqrt5},& 0<s<\varphi^{-1},\\[3mm]
\dfrac{\varphi}{\sqrt5},& \varphi^{-1}<s<1,\\[3mm]
0,& s>1.
\end{cases}
\]
故在 \(s=\varphi^{-1}\) 处的斜率跃降为
\[
\frac{\varphi^2}{\sqrt5}-\frac{\varphi}{\sqrt5}=\frac{1}{\sqrt5},
\]
而在 \(s=1\) 处的斜率跃降为
\[
\frac{\varphi}{\sqrt5}-0=\frac{\varphi}{\sqrt5}.
\]
因此分布意义下
\[
-D^2\widehat C_\infty=\frac{1}{\sqrt5}\,\delta_{\varphi^{-1}}+\frac{\varphi}{\sqrt5}\,\delta_1.
\]
于是其尾函数为
\[
\nu_\infty([s,\infty))=
\begin{cases}
\dfrac{\varphi^2}{\sqrt5},& 0<s<\varphi^{-1},\\[3mm]
\dfrac{\varphi}{\sqrt5},& \varphi^{-1}<s<1,\\[3mm]
0,& s>1.
\end{cases}
\]
对该尾函数从 \(0\) 积分到 \(s\) 即恢复 \(\widehat C_\infty\)。另一方面，把上式与 \(\widehat N_\infty\) 的分段表达比较，立得
\[
\widehat N_\infty(s)=\frac{\sqrt5}{\varphi^2}\,\nu_\infty([s,\infty))
\qquad
\bigl(s\notin\{\varphi^{-1},1\}\bigr).
\]
证毕。
\end{proof}

\begin{theorem}[fold-gauge 规范化常数序列唯一恢复全部偶 \texorpdfstring{$\zeta$}{zeta} 值]\label{thm:xi-foldbin-gauge-constant-even-zeta-recovery}
\leanverified{paper\_xi\_foldbin\_gauge\_constant\_even\_zeta\_recovery}
沿用推论 \ref{cor:fold-bin-recover-2pi} 的 \(C_m\) 记号，并令
\[
q:=\frac{\varphi}{2}\in(0,1).
\]
则对每个 \(r\ge 1\)，递归定义残差极限
\[
\mathcal R_1:=\lim_{m\to\infty}q^{-m}\bigl(\log C_m-\log(2\pi)\bigr),
\]
\[
\mathcal R_r:=
\lim_{m\to\infty}
q^{-(2r-1)m}
\left[
\log C_m-\log(2\pi)-\sum_{j=1}^{r-1}\mathcal R_j\,q^{(2j-1)m}
\right]
\qquad (r\ge 2).
\]
则每个极限都存在，并满足
\[
\boxed{\;
\mathcal R_r=
\frac{B_{2r}}{r(2r-1)}\bigl(\varphi^{-1}+\varphi^{2r-3}\bigr).
\;}
\]
因此
\[
\boxed{\;
B_{2r}
=
\frac{r(2r-1)}{\varphi^{-1}+\varphi^{2r-3}}\,\mathcal R_r,
\;}
\]
并由标准恒等式
\[
B_{2r}=(-1)^{r-1}\frac{2(2r)!}{(2\pi)^{2r}}\zeta(2r)
\]
得到
\[
\boxed{\;
\zeta(2r)
=
(-1)^{r-1}\frac{(2\pi)^{2r}}{2(2r)!}\,
\frac{r(2r-1)}{\varphi^{-1}+\varphi^{2r-3}}\,
\mathcal R_r.
\;}
\]
故单一序列 \(\{\log C_m\}_{m\ge 1}\) 已构成完整的 even-\(\zeta\) 塔编码器：它不仅内生恢复 \(2\pi\)，而且逐阶恢复全部 Bernoulli 数与全部偶值 \(\zeta(2r)\)。
\end{theorem}
\begin{proof}
由定理 \ref{thm:fold-bin-gauge-constant-stirling-bernoulli-hierarchy}，对任意固定 \(R\ge 1\) 有
\[
\log C_m
=
\log(2\pi)+
\sum_{r=1}^{R}
\frac{B_{2r}}{r(2r-1)}
\bigl(\varphi^{-1}+\varphi^{2r-3}\bigr)
q^{(2r-1)m}
\;+\;
O\!\bigl(q^{(2R+1)m}\bigr).
\]
因此对每个固定 \(r\)，移去前 \(r-1\) 层后得到
\[
\log C_m-\log(2\pi)-\sum_{j=1}^{r-1}\mathcal R_j q^{(2j-1)m}
=
\frac{B_{2r}}{r(2r-1)}\bigl(\varphi^{-1}+\varphi^{2r-3}\bigr)q^{(2r-1)m}
+O\!\bigl(q^{(2r+1)m}\bigr),
\]
其中对 \(j<r\) 的 \(\mathcal R_j\) 已由同一公式递归恢复。两边同乘 \(q^{-(2r-1)m}\) 并令 \(m\to\infty\)，即得
\[
\mathcal R_r=
\frac{B_{2r}}{r(2r-1)}\bigl(\varphi^{-1}+\varphi^{2r-3}\bigr).
\]
于是 \(B_{2r}\) 的显式反演随即成立；再代入 Bernoulli 与 \(\zeta(2r)\) 的标准恒等式，即得到偶值 \(\zeta\) 的恢复公式。证毕。
\end{proof}

\begin{theorem}[素数寄存器算术半群的幂等元精确计数]\label{thm:xi-prime-register-idempotent-exact-count}
\leanverified{paper\_xi\_prime\_register\_idempotent\_exact\_count}
沿用定理 \ref{thm:killo-finite-semigroup-prime-valuation-arithmetization} 的记号。记
\[
\operatorname{Idem}(S_n):=\{N\in S_n:\ N\star N=N\},
\]
\[
\operatorname{Idem}_k(S_n):=
\{N\in \operatorname{Idem}(S_n):\ |\operatorname{Im}(f_N)|=k\}
\qquad (1\le k\le n).
\]
则有精确计数公式
\[
\boxed{\;
\bigl|\operatorname{Idem}_k(S_n)\bigr|=\binom{n}{k}k^{n-k}
\qquad (1\le k\le n),
\;}
\]
从而
\[
\boxed{\;
\bigl|\operatorname{Idem}(S_n)\bigr|
=
\sum_{k=1}^{n}\binom{n}{k}k^{n-k}.
\;}
\]
换言之，素数寄存器算术半群中的幂等元个数与 \([n]\) 上幂等自映射的经典计数严格一致，并被完全算术化为估值方程
\[
v_{q_{v_{q_i}(N)}}(N)=v_{q_i}(N)\qquad (i\in[n]).
\]
\end{theorem}
\begin{proof}
由定理 \ref{thm:killo-finite-semigroup-prime-valuation-arithmetization} 与推论 \ref{cor:killo-projection-idempotent-prime-register}，\(\operatorname{Idem}(S_n)\) 与集合 \([n]\) 上的全部幂等自映射一一对应。

固定 \(k\in\{1,\dots,n\}\)。若 \(f:[n]\to[n]\) 幂等且 \(|\operatorname{Im}(f)|=k\)，记
\[
A:=\operatorname{Im}(f),\qquad |A|=k.
\]
由幂等性 \(f(f(i))=f(i)\) 可知，对每个 \(a\in A\) 都有 \(f(a)=a\)。因此 \(f\) 在像集 \(A\) 上必为恒等，而对每个 \(x\in[n]\setminus A\) 只能任意映到 \(A\) 中的某一点。

反之，任取一个 \(k\)-元子集 \(A\subseteq[n]\)，令 \(f|_A=\mathrm{id}_A\)，并把每个 \(x\in[n]\setminus A\) 任意送入 \(A\)，所得映射自动满足 \(f^2=f\) 且 \(|\operatorname{Im}(f)|=k\)。

故先选像集 \(A\) 有 \(\binom{n}{k}\) 种，再给其余 \(n-k\) 个点各指定一个像到 \(A\) 中，有 \(k^{n-k}\) 种。于是
\[
\bigl|\operatorname{Idem}_k(S_n)\bigr|=\binom{n}{k}k^{n-k}.
\]
对 \(k\) 求和即得总数。证毕。
\end{proof}

\begin{theorem}[素数寄存器半群的唯一极小双边理想与左零带结构]\label{thm:xi-prime-register-minimal-ideal-left-zero-band}
\leanverified{paper\_xi\_prime\_register\_minimal\_ideal\_left\_zero\_band}
沿用定理 \ref{thm:killo-finite-semigroup-prime-valuation-arithmetization} 的 \((S_n,\star)\) 记号，并设
\[
Q_n:=\prod_{i=0}^{n-1}q_i,
\qquad
\mathcal K_n:=\{Q_n^c:\ c\in[n]\}\subset S_n.
\]
则：
\begin{enumerate}
  \item \(N\in \mathcal K_n\) 当且仅当 \(f_N:[n]\to[n]\) 为常值映射；
  \item \(\mathcal K_n\) 恰是全部秩 \(1\) 元素组成的集合，即
  \[
  \mathcal K_n=\{N\in S_n:\ |\operatorname{Im}(f_N)|=1\};
  \]
  \item 对任意 \(a,b\in[n]\)，有
  \[
  \boxed{\;
  Q_n^a\star Q_n^b=Q_n^a,
  \;}
  \]
  故 \(\mathcal K_n\) 作为子半群同构于一个 \(n\) 元左零带；
  \item \(\mathcal K_n\) 是 \((S_n,\star)\) 的唯一极小双边理想：任意非空双边理想 \(I\subseteq S_n\) 都满足 \(\mathcal K_n\subseteq I\)。
\end{enumerate}
\end{theorem}
\begin{proof}
若 \(N=Q_n^c\)，则对每个 \(i\in[n]\) 都有
\[
v_{q_i}(N)=c,
\]
故 \(f_N\) 为常值 \(c\)。反之，若 \(f_N\equiv c\)，则全部估值都等于 \(c\)，于是
\[
N=\prod_{i=0}^{n-1}q_i^c=Q_n^c.
\]
这就证明了第 (1) 条，而第 (2) 条只是“常值映射恰有秩 \(1\)”的直接转写。

对第 (3) 条，记 \(c_a,c_b\) 分别为常值 \(a,b\) 的映射。由定理 \ref{thm:killo-finite-semigroup-prime-valuation-arithmetization},
\[
f_{Q_n^a\star Q_n^b}=f_{Q_n^a}\circ f_{Q_n^b}=c_a\circ c_b=c_a,
\]
故 \(Q_n^a\star Q_n^b=Q_n^a\)。

最后证明第 (4) 条。任取非空双边理想 \(I\subseteq S_n\)，取 \(N\in I\)。对任意 \(a,b\in[n]\)，由于 \(I\) 为双边理想，有
\[
Q_n^a\star N\star Q_n^b\in I.
\]
而在变换半群侧，
\[
f_{Q_n^a\star N\star Q_n^b}
=
f_{Q_n^a}\circ f_N\circ f_{Q_n^b}
=
c_a\circ f_N\circ c_b
=
c_a.
\]
故
\[
Q_n^a\star N\star Q_n^b=Q_n^a\in I.
\]
由于 \(a\) 任意，\(\mathcal K_n\subseteq I\)。因此 \(\mathcal K_n\) 包含于每个非空双边理想之中，遂为唯一极小双边理想。证毕。
\end{proof}

\begin{theorem}[bin-fold 的尺度化多重度最终成为末位的完备读取统计量]\label{thm:xi-foldbin-lastbit-complete-readout-statistic}
\leanverified{paper\_xi\_foldbin\_lastbit\_complete\_readout\_statistic}
沿用定理 \ref{thm:fold-bin-two-state-asymptotic} 与推论 \ref{cor:fold-bin-two-point-limit-law} 的记号。对每个 \(w\in X_m\) 定义尺度化多重度
\[
Z_m^{\mathrm{bin}}(w):=\left(\frac{\varphi}{2}\right)^m d_m^{\mathrm{bin}}(w),
\]
并置阈值
\[
\tau:=\frac{1+\varphi^{-1}}{2}.
\]
则存在 \(m_0\) 使得对一切 \(m\ge m_0\) 与一切 \(w\in X_m\)，都有精确判别
\[
\boxed{\;
w_m=0 \iff Z_m^{\mathrm{bin}}(w)>\tau,
\qquad
w_m=1 \iff Z_m^{\mathrm{bin}}(w)<\tau.
\;}
\]
换言之，对所有充分大的 \(m\)，末位比特 \(w_m\) 是尺度化多重度 \(Z_m^{\mathrm{bin}}\) 的确定性函数。

若记 \(W_m\sim \mu_m\) 为推论 \ref{cor:fold-bin-two-point-limit-law} 的 bin-fold 输出，并记 \(H_{\mu_m}\)、\(I_{\mu_m}\) 分别为 \(\mu_m\) 下的 Shannon 条件熵与互信息，则对充分大的 \(m\) 有
\[
H_{\mu_m}\!\bigl(W_m(m)\mid Z_m^{\mathrm{bin}}(W_m)\bigr)=0,
\qquad
I_{\mu_m}\!\bigl(W_m(m);Z_m^{\mathrm{bin}}(W_m)\bigr)=H_{\mu_m}\!\bigl(W_m(m)\bigr).
\]
并且
\[
\mu_m\bigl(W_m(m)=1\bigr)\longrightarrow \frac{1}{\varphi\sqrt5},
\]
\[
I_{\mu_m}\!\bigl(W_m(m);Z_m^{\mathrm{bin}}(W_m)\bigr)\longrightarrow
h_2\!\left(\frac{1}{\varphi\sqrt5}\right),
\qquad
h_2(p):=-p\log p-(1-p)\log(1-p).
\]
因此，尺度化多重度并不只是末位的渐近充分统计量，而在充分大尺度上已经成为末位的精确读取统计量。
\end{theorem}
\begin{proof}
由定理 \ref{thm:fold-bin-two-state-asymptotic}，存在常数 \(C>0\) 使对充分大 \(m\) 与任意 \(w\in X_m\) 成立
\[
\bigl|Z_m^{\mathrm{bin}}(w)-1\bigr|
\le
C\left(\frac{\varphi}{2}\right)^m
\qquad (w_m=0),
\]
\[
\bigl|Z_m^{\mathrm{bin}}(w)-\varphi^{-1}\bigr|
\le
C\left(\frac{\varphi}{2}\right)^m
\qquad (w_m=1).
\]
由于两个中心之间的间隔为
\[
1-\varphi^{-1}=\varphi^{-2}>0,
\]
可取 \(m_0\) 使得当 \(m\ge m_0\) 时
\[
C\left(\frac{\varphi}{2}\right)^m<\frac{1-\varphi^{-1}}{2}.
\]
于是当 \(m\ge m_0\) 时，若 \(w_m=0\)，则
\[
Z_m^{\mathrm{bin}}(w)>\frac{1+\varphi^{-1}}{2}=\tau,
\]
而若 \(w_m=1\)，则
\[
Z_m^{\mathrm{bin}}(w)<\frac{1+\varphi^{-1}}{2}=\tau.
\]
故阈值判别式成立，从而 \(W_m(m)\) 是 \(Z_m^{\mathrm{bin}}(W_m)\) 的确定性函数。于是条件熵为零，并有
\[
I_{\mu_m}\!\bigl(W_m(m);Z_m^{\mathrm{bin}}(W_m)\bigr)
=
H_{\mu_m}\!\bigl(W_m(m)\bigr).
\]

另一方面，由推论 \ref{cor:fold-bin-two-point-limit-law},
\[
\mu_m\bigl(W_m(m)=1\bigr)\longrightarrow \frac{1}{\varphi\sqrt5}.
\]
故二值 Shannon 熵满足
\[
H_{\mu_m}\!\bigl(W_m(m)\bigr)
=
h_2\!\bigl(\mu_m(W_m(m)=1)\bigr)
\longrightarrow
h_2\!\left(\frac{1}{\varphi\sqrt5}\right).
\]
于是互信息极限随之成立。证毕。
\end{proof}

\paragraph{欠分辨曲率、even-\texorpdfstring{$\zeta$}{zeta} 编码与估值算术化的闭链}
上述五条结果把此前分散的三条线索压缩为同一条解析链：在线性尺度上，bin-fold 的欠分辨极限等价于一条支撑于 \(\{\varphi^{-1},1\}\) 的二原子曲率测度；在规范体积常数序列上，\(\{\log C_m\}_{m\ge 1}\) 逐阶内生恢复全部 Bernoulli 系数与 even-\(\zeta\) 塔；在素数寄存器估值侧，变换半群的幂等谱与极小双边理想则获得纯算术实现。回到 bin-fold 输出本身，尺度化多重度最终把末位比特完全读出，从而使“二原子极限”不只停留在分布层，而上升为逐点可判别的读取结构。
\paragraph{Bernoulli jet 的 Hankel 非退化、冻结相强逆指数与 Fibonacci 共振整函数}
在 fold-gauge 常数 \(\log C_m\) 的 Stirling--Bernoulli 层级、固定冻结相的预言机容量阈值以及共振梯整函数 \(\mathcal F_\varphi\) 的闭包已经建立之后，还可进一步抽出下列五条后果。

\begin{theorem}[\texorpdfstring{$\log C_m$}{log Cm} 的 Hankel 行列式具有严格的 Vandermonde 主项]\label{thm:xi-logcm-hankel-vandermonde-mainterm}
\leanverified{paper\_xi\_logcm\_hankel\_vandermonde\_mainterm}
令
\[
a_m:=\log C_m-\log(2\pi),
\qquad
\lambda_r:=\left(\frac{\varphi}{2}\right)^{2r-1}\in(0,1)
\qquad (r\ge 1),
\]
\[
c_r:=\frac{B_{2r}}{r(2r-1)}\bigl(\varphi^{-1}+\varphi^{2r-3}\bigr).
\]
则由定理 \ref{thm:fold-bin-gauge-constant-stirling-bernoulli-hierarchy}，对每个固定 \(R\ge 1\) 有
\[
a_m=\sum_{r=1}^{R}c_r\lambda_r^m+O(\lambda_{R+1}^m),
\qquad
0<\lambda_{R+1}<\lambda_R<\cdots<\lambda_1<1.
\]
对固定 \(R\) 定义 Hankel 块
\[
H_R(m):=\bigl[a_{m+i+j}\bigr]_{0\le i,j\le R-1},
\qquad
\Delta_R(m):=\det H_R(m).
\]
则当 \(m\to\infty\) 时，
\[
\boxed{\;
\Delta_R(m)=
\Bigl(\prod_{r=1}^{R}c_r\Bigr)
\Bigl(\prod_{r=1}^{R}\lambda_r\Bigr)^m
\prod_{1\le i<j\le R}(\lambda_j-\lambda_i)^2
\bigl(1+o(1)\bigr).
\;}
\]
特别地，\(\Delta_R(m)\neq 0\) 对充分大的 \(m\) 必成立；因此 \(\{\log C_m\}\) 的前 \(R\) 阶 Bernoulli jet 在有限 Hankel 阶上最终完全非退化。
\end{theorem}
\begin{proof}
记 Vandermonde 矩阵
\[
V_R:=\bigl(\lambda_j^{\,i}\bigr)_{0\le i\le R-1,\ 1\le j\le R},
\]
以及对角矩阵
\[
D_R(m):=\operatorname{diag}\bigl(c_1\lambda_1^m,\dots,c_R\lambda_R^m\bigr).
\]
由 \(a_{m+i+j}\) 的展开，对每个固定 \(i,j\in\{0,\dots,R-1\}\) 都有
\[
a_{m+i+j}=\sum_{r=1}^{R}c_r\lambda_r^{m+i+j}+O(\lambda_{R+1}^{m+i+j})
=\sum_{r=1}^{R}c_r\lambda_r^{m+i+j}+O(\lambda_{R+1}^{m}),
\]
其中因为 \(i+j\) 有界，误差项可统一吸收为 \(O(\lambda_{R+1}^m)\)。
故
\[
H_R(m)=V_R D_R(m) V_R^{\mathsf T}+E_R(m),
\]
且 \(E_R(m)\) 的全部条目均为 \(O(\lambda_{R+1}^m)\)。

主项矩阵
\[
M_R(m):=V_R D_R(m) V_R^{\mathsf T}
\]
的行列式满足
\[
\det M_R(m)=\det(V_R)^2\det D_R(m)
=\Bigl(\prod_{r=1}^{R}c_r\lambda_r^m\Bigr)\det(V_R)^2.
\]
又因 \(\lambda_1,\dots,\lambda_R\) 两两不同，
\[
\det(V_R)=\prod_{1\le i<j\le R}(\lambda_j-\lambda_i)\neq 0.
\]
此外，偶 Bernoulli 数 \(B_{2r}\) 对一切 \(r\ge 1\) 都非零，且
\(
\varphi^{-1}+\varphi^{2r-3}>0
\)，故 \(c_r\neq 0\)。
因此 \(M_R(m)\) 可逆，并且
\[
M_R(m)^{-1}=V_R^{-\mathsf T}D_R(m)^{-1}V_R^{-1}.
\]
由于 \(V_R\) 与 \(V_R^{-1}\) 与 \(m\) 无关，而 \(D_R(m)^{-1}\) 的最大对角元尺度为 \(O(\lambda_R^{-m})\)，故
\[
\|M_R(m)^{-1}\|=O(\lambda_R^{-m}).
\]
另一方面，\(\|E_R(m)\|=O(\lambda_{R+1}^m)\)，从而
\[
\|M_R(m)^{-1}E_R(m)\|=O\!\left(\left(\frac{\lambda_{R+1}}{\lambda_R}\right)^m\right)=o(1).
\]
于是
\[
\det H_R(m)=\det M_R(m)\cdot
\det\!\bigl(I+M_R(m)^{-1}E_R(m)\bigr)
=\det M_R(m)\bigl(1+o(1)\bigr),
\]
这就给出了所示主项公式。由于主项常数非零，\(\Delta_R(m)\neq 0\) 对充分大的 \(m\) 必成立。证毕。
\end{proof}

\begin{corollary}[固定阶 Bernoulli 系数可由有限连续窗口指数精度恢复]\label{cor:xi-logcm-finite-window-bernoulli-recovery}
\leanverified{paper\_xi\_logcm\_finite\_window\_bernoulli\_recovery}
沿用定理 \ref{thm:xi-logcm-hankel-vandermonde-mainterm} 的记号。对固定 \(R\ge 1\)，令
\[
\mathbf a^{(m)}:=
\bigl(a_m,a_{m+1},\dots,a_{m+R-1}\bigr)^{\mathsf T},
\qquad
\mathbf c:=(c_1,\dots,c_R)^{\mathsf T},
\]
\[
V_R:=\bigl(\lambda_j^{\,i}\bigr)_{0\le i\le R-1,\ 1\le j\le R},
\qquad
D_m:=\operatorname{diag}(\lambda_1^m,\dots,\lambda_R^m).
\]
定义有限窗口恢复向量
\[
\widehat{\mathbf c}^{(m)}:=D_m^{-1}V_R^{-1}\mathbf a^{(m)}.
\]
则对每个 \(1\le r\le R\) 都有
\[
\boxed{\;
\widehat c_r^{(m)}=c_r+O\!\left(\left(\frac{\lambda_{R+1}}{\lambda_r}\right)^m\right).
\;}
\]
因而前 \(R\) 个 Bernoulli 系数与偶值 \(\zeta\)-jet 都可由仅 \(R\) 个连续样本
\[
\log C_m,\ \log C_{m+1},\ \dots,\ \log C_{m+R-1}
\]
以指数精度恢复。更具体地，
\[
B_{2r}=
\frac{r(2r-1)}{\varphi^{-1}+\varphi^{2r-3}}\,c_r,
\qquad
\zeta(2r)=(-1)^{r-1}\frac{(2\pi)^{2r}}{2(2r)!}\,B_{2r}.
\]
\end{corollary}
\begin{proof}
由定理 \ref{thm:fold-bin-gauge-constant-stirling-bernoulli-hierarchy}，
\[
a_{m+i}=\sum_{r=1}^{R}c_r\lambda_r^{m+i}+O(\lambda_{R+1}^{m+i})
\qquad (0\le i\le R-1).
\]
故可写成向量形式
\[
\mathbf a^{(m)}=V_R D_m \mathbf c+\mathbf e^{(m)},
\]
其中 \(\mathbf e^{(m)}\) 的每个分量都是 \(O(\lambda_{R+1}^{m})\)。
左乘 \(D_m^{-1}V_R^{-1}\) 得
\[
\widehat{\mathbf c}^{(m)}
=
\mathbf c+D_m^{-1}V_R^{-1}\mathbf e^{(m)}.
\]
因为 \(V_R^{-1}\) 与 \(m\) 无关，第 \(r\) 个分量的误差至多被 \(\lambda_r^{-m}\) 放大，故
\[
\widehat c_r^{(m)}-c_r=
O\!\left(\frac{\lambda_{R+1}^m}{\lambda_r^m}\right)
=
O\!\left(\left(\frac{\lambda_{R+1}}{\lambda_r}\right)^m\right).
\]
从而有限窗口对 \((c_1,\dots,c_R)\) 的恢复具有指数精度。
最后把
\[
c_r=\frac{B_{2r}}{r(2r-1)}\bigl(\varphi^{-1}+\varphi^{2r-3}\bigr)
\]
与
\[
B_{2r}=(-1)^{r-1}\frac{2(2r)!}{(2\pi)^{2r}}\zeta(2r)
\]
代回即可。证毕。
\end{proof}

\begin{corollary}[fold-gauge 常数第二 Bernoulli 系数的显式闭式与一步比值极限]\label{cor:xi-logcm-second-bernoulli-coefficient-ratio}
设
\[
E_m:=\log C_m-\log(2\pi).
\]
则有精确展开
\[
E_m
=
\frac{1}{3\varphi}\left(\frac{\varphi}{2}\right)^m
-\frac{\sqrt5}{180}\left(\frac{\varphi}{2}\right)^{3m}
+O\!\left(\left(\frac{\varphi}{2}\right)^{5m}\right).
\]
从而
\[
\lim_{m\to\infty}
\left(\frac{2}{\varphi}\right)^{3m}
\left(
E_m-\frac{1}{3\varphi}\left(\frac{\varphi}{2}\right)^m
\right)
=
-\frac{\sqrt5}{180},
\]
并且
\[
\lim_{m\to\infty}\frac{E_{m+1}}{E_m}=\frac{\varphi}{2}.
\]
\end{corollary}
\begin{proof}
由定理 \ref{thm:fold-bin-gauge-constant-stirling-bernoulli-hierarchy} 在 \(R=2\) 时的展开，
\[
E_m=c_1\left(\frac{\varphi}{2}\right)^m+c_2\left(\frac{\varphi}{2}\right)^{3m}
+O\!\left(\left(\frac{\varphi}{2}\right)^{5m}\right),
\]
其中
\[
c_r=\frac{B_{2r}}{r(2r-1)}\bigl(\varphi^{-1}+\varphi^{2r-3}\bigr).
\]
对 \(r=1\)，由 \(B_2=\frac16\) 得
\[
c_1=\frac{B_2}{1\cdot 1}\bigl(\varphi^{-1}+\varphi^{-1}\bigr)
=
\frac16\cdot 2\varphi^{-1}
=
\frac{1}{3\varphi}.
\]
对 \(r=2\)，由 \(B_4=-\frac1{30}\) 得
\[
c_2=\frac{B_4}{2\cdot 3}\bigl(\varphi^{-1}+\varphi\bigr)
=
-\frac1{180}\bigl(\varphi+\varphi^{-1}\bigr).
\]
再用恒等式
\[
\varphi+\varphi^{-1}=\sqrt5,
\]
便得到
\[
c_2=-\frac{\sqrt5}{180}.
\]
代回即得所示二阶展开。于是
\[
\left(\frac{2}{\varphi}\right)^{3m}
\left(
E_m-\frac{1}{3\varphi}\left(\frac{\varphi}{2}\right)^m
\right)
=
-\frac{\sqrt5}{180}
+O\!\left(\left(\frac{\varphi}{2}\right)^{2m}\right),
\]
令 \(m\to\infty\) 即得第一条极限。

最后，由首项系数 \(\frac{1}{3\varphi}\neq 0\)，有
\[
E_m=\frac{1}{3\varphi}\left(\frac{\varphi}{2}\right)^m\bigl(1+o(1)\bigr),
\]
故
\[
\frac{E_{m+1}}{E_m}
=
\frac{\frac{1}{3\varphi}(\varphi/2)^{m+1}(1+o(1))}
{\frac{1}{3\varphi}(\varphi/2)^m(1+o(1))}
\longrightarrow
\frac{\varphi}{2}.
\]
证毕。
\end{proof}
\begin{theorem}[固定冻结相预言机反演的强逆指数]\label{thm:xi-fixed-freezing-oracle-strong-converse-exponent}
\leanverified{paper\_xi\_fixed\_freezing\_oracle\_strong\_converse\_exponent}
沿用结论 \ref{con:xi-terminal-fixed-freezing-escort-oracle-sharp-threshold} 的记号，并记
\[
P_{m,a}^{\mathrm{opt}}(B):=\mathrm{Succ}^{\mathrm{esc}}_{m,a}(B).
\]
设 \(a>a_{\mathrm c}\)，并令
\[
\Delta(a):=a(\alpha_*-\alpha_*^{(2)})-(\log\varphi-g_*)>0.
\]
若预算满足
\[
B_m=\beta m+o(m),
\]
则有如下精确强逆指数：
\begin{enumerate}
  \item 若 \(\beta<\alpha_*/\log 2\)，则
  \[
  \boxed{\;
  \limsup_{m\to\infty}\frac{1}{m}\log P_{m,a}^{\mathrm{opt}}(B_m)
  \le
  -\min\bigl\{\Delta(a),\ \alpha_*-\beta\log 2\bigr\}.
  \;}
  \]
  \item 若 \(\beta>\alpha_*/\log 2\)，则
  \[
  \boxed{\;
  1-P_{m,a}^{\mathrm{opt}}(B_m)
  \le
  \exp\!\bigl(-m\Delta(a)+o(m)\bigr).
  \;}
  \]
\end{enumerate}
因此临界斜率 \(\alpha_*/\log 2\) 不仅给出可/不可恢复阈值，而且给出成功率与失败率的精确指数分界。
\end{theorem}
\begin{proof}
记
\[
A_m^{\max}:=\{x\in X_m:\ d_m(x)=M_m\}.
\]
由结论 \ref{con:xi-terminal-fixed-freezing-escort-oracle-sharp-threshold} 的证明与定理 \ref{thm:fold-pressure-exact-freezing-gap-threshold} 的冻结集中估计，对固定 \(a>a_{\mathrm c}\) 有
\[
\pi_{m,a}(A_m^{\max})
\ge
1-\exp\!\bigl(-m\Delta(a)+o(m)\bigr).
\]

先证第 (1) 条。任取一个允许 \(B_m\) 比特外部寄存器的编码--解码机制。对任意 \(x\in A_m^{\max}\)，纤维 \(\Fold_m^{-1}(x)\) 的大小恰为 \(M_m\)。
在该纤维内，至多只有 \(2^{B_m}\) 个微态可以被彼此区分，故条件成功概率不超过
\[
\min\!\left(1,\frac{2^{B_m}}{M_m}\right).
\]
于是整体成功率满足
\[
P_{m,a}^{\mathrm{opt}}(B_m)
\le
\pi_{m,a}(A_m^{\max})\min\!\left(1,\frac{2^{B_m}}{M_m}\right)
+\pi_{m,a}(X_m\setminus A_m^{\max}).
\]
当 \(\beta<\alpha_*/\log 2\) 时，
\[
\frac{2^{B_m}}{M_m}
\le
\exp\!\bigl(-m(\alpha_*-\beta\log 2)+o(m)\bigr).
\]
同时
\[
\pi_{m,a}(X_m\setminus A_m^{\max})
\le
\exp\!\bigl(-m\Delta(a)+o(m)\bigr).
\]
故
\[
P_{m,a}^{\mathrm{opt}}(B_m)
\le
\exp\!\bigl(-m(\alpha_*-\beta\log 2)+o(m)\bigr)
+\exp\!\bigl(-m\Delta(a)+o(m)\bigr).
\]
对数除以 \(m\) 后取上极限，即得
\[
\limsup_{m\to\infty}\frac{1}{m}\log P_{m,a}^{\mathrm{opt}}(B_m)
\le
-\min\bigl\{\Delta(a),\alpha_*-\beta\log 2\bigr\}.
\]

再证第 (2) 条。若 \(\beta>\alpha_*/\log 2\)，则可取 \(\varepsilon>0\) 使
\[
\beta\log 2\ge \alpha_*+\varepsilon.
\]
由 \(B_m=\beta m+o(m)\) 与 \(M_m=\exp(\alpha_*m+o(m))\)，对充分大的 \(m\) 有
\[
2^{B_m}\ge M_m.
\]
于是可以在每条最大纤维上施加完全一一编码，并在其外任意定义译码机制。该机制在 \(A_m^{\max}\) 上成功率为 \(1\)，故
\[
P_{m,a}^{\mathrm{opt}}(B_m)\ge \pi_{m,a}(A_m^{\max}).
\]
从而
\[
1-P_{m,a}^{\mathrm{opt}}(B_m)
\le
1-\pi_{m,a}(A_m^{\max})
\le
\exp\!\bigl(-m\Delta(a)+o(m)\bigr).
\]
证毕。
\end{proof}

\begin{theorem}[Fibonacci 共振整函数的零点计数满足自相似递推与精确线性主项]\label{thm:xi-fold-resonance-zero-count-selfsimilar-linear}
\leanverified{paper\_xi\_fold\_resonance\_zero\_count\_selfsimilar\_linear}
沿用定理 \ref{thm:fold-resonance-entire-lp} 的记号，记 Fibonacci 共振整函数为
\[
\mathcal F_\varphi(z)=\prod_{n=2}^{\infty}\cos(\pi z\varphi^{-n}).
\]
定义正半轴零点计数函数
\[
N_\varphi^+(R):=\#\bigl(\mathcal Z(\mathcal F_\varphi)\cap(0,R]\bigr),
\qquad R>0.
\]
约定当上限小于下限时求和为空和 \(0\)。则对每个 \(R>0\) 都有严格公式
\[
\boxed{\;
N_\varphi^+(R)=
\sum_{n=2}^{\lfloor \log_\varphi(2R)\rfloor}
\left\lfloor R\varphi^{-n}+\frac12\right\rfloor.
\;}
\]
从而得到自相似递推
\[
\boxed{\;
N_\varphi^+(R)=N_\varphi^+\!\left(\frac{R}{\varphi}\right)
+\left\lfloor R\varphi^{-2}+\frac12\right\rfloor.
\;}
\]
进一步，若记
\[
N_\varphi(R):=\#\bigl(\mathcal Z(\mathcal F_\varphi)\cap[-R,R]\bigr),
\]
则当 \(R\to\infty\) 时，
\[
\boxed{\;
N_\varphi^+(R)=R+O(\log R),
\qquad
N_\varphi(R)=2R+O(\log R).
\;}
\]
\end{theorem}
\begin{proof}
由定理 \ref{thm:fold-resonance-entire-lp}，
\[
\mathcal Z(\mathcal F_\varphi)=
\Bigl\{(k+\tfrac12)\varphi^n:\ k\in\ZZ,\ n\ge 2\Bigr\},
\]
故正零点恰为
\[
\Bigl\{(k+\tfrac12)\varphi^n:\ k\in\ZZ_{\ge 0},\ n\ge 2\Bigr\}.
\]
对固定 \(n\ge 2\)，满足
\[
(k+\tfrac12)\varphi^n\le R
\]
的非负整数 \(k\) 的个数正是
\[
\left\lfloor R\varphi^{-n}+\frac12\right\rfloor.
\]
只有当 \(R\varphi^{-n}\ge \frac12\) 时这一项才可能非零，即 \(n\le \lfloor\log_\varphi(2R)\rfloor\)。因此得到精确求和公式。

设 \(L(R):=\lfloor\log_\varphi(2R)\rfloor\)。则
\[
L(R/\varphi)=L(R)-1.
\]
于是
\[
N_\varphi^+\!\left(\frac{R}{\varphi}\right)
=
\sum_{n=2}^{L(R)-1}
\left\lfloor R\varphi^{-(n+1)}+\frac12\right\rfloor
=
\sum_{n=3}^{L(R)}
\left\lfloor R\varphi^{-n}+\frac12\right\rfloor.
\]
将其与 \(N_\varphi^+(R)\) 的求和式相减，即得
\[
N_\varphi^+(R)-N_\varphi^+\!\left(\frac{R}{\varphi}\right)
=
\left\lfloor R\varphi^{-2}+\frac12\right\rfloor.
\]

最后，由
\[
\left\lfloor R\varphi^{-n}+\frac12\right\rfloor
=
R\varphi^{-n}+O(1)
\]
并对 \(n=2,\dots,L(R)\) 求和，得到
\[
N_\varphi^+(R)
=
R\sum_{n=2}^{L(R)}\varphi^{-n}+O(L(R))
=
R+O(\log R),
\]
其中使用 \(\sum_{n=2}^{\infty}\varphi^{-n}=1\) 与 \(L(R)=O(\log R)\)。
由于零点关于原点对称且 \(0\) 不是零点，
\[
N_\varphi(R)=2N_\varphi^+(R)=2R+O(\log R).
\]
证毕。
\end{proof}

\begin{corollary}[Fibonacci 共振整函数零点计数的前向自相似更新公式]\label{cor:xi-fold-resonance-zero-count-forward-update}
沿用定理 \ref{thm:xi-fold-resonance-zero-count-selfsimilar-linear} 的记号。则对任意
\[
R\ge 0
\]
都有
\[
N_\varphi^+(\varphi R)
=
N_\varphi^+(R)
+
\left\lfloor \frac{R}{\varphi}+\frac12\right\rfloor.
\]
若再记
\[
N_\varphi(R):=\#\bigl(\mathcal Z(\mathcal F_\varphi)\cap[-R,R]\bigr),
\]
则同样有
\[
N_\varphi(\varphi R)
=
N_\varphi(R)
+
2\left\lfloor \frac{R}{\varphi}+\frac12\right\rfloor.
\]
\end{corollary}
\begin{proof}
定理 \ref{thm:xi-fold-resonance-zero-count-selfsimilar-linear} 已给出
\[
N_\varphi^+(T)=N_\varphi^+\!\left(\frac{T}{\varphi}\right)+\left\lfloor T\varphi^{-2}+\frac12\right\rfloor
\qquad (T>0).
\]
当 \(R=0\) 时两式显然成立。以下设 \(R>0\)，并对上式取
\[
T=\varphi R
\]
便得到
\[
N_\varphi^+(\varphi R)
=
N_\varphi^+(R)
+
\left\lfloor \frac{R}{\varphi}+\frac12\right\rfloor,
\]
这就证明了第一式。又由于零点关于原点对称且 \(0\notin\mathcal Z(\mathcal F_\varphi)\)，有
\[
N_\varphi(R)=2N_\varphi^+(R),
\]
从而
\[
N_\varphi(\varphi R)
=
2N_\varphi^+(\varphi R)
=
2N_\varphi^+(R)
+
2\left\lfloor \frac{R}{\varphi}+\frac12\right\rfloor
=
N_\varphi(R)
+
2\left\lfloor \frac{R}{\varphi}+\frac12\right\rfloor.
\]
证毕。
\end{proof}
\begin{corollary}[Fibonacci 共振整函数零点的精确线性密度与指数收敛阶]\label{cor:xi-fold-resonance-zero-density-exponent-one}
\leanverified{paper\_xi\_fold\_resonance\_zero\_density\_exponent\_one}
沿用定理 \ref{thm:xi-fold-resonance-zero-count-selfsimilar-linear} 的记号，则零点线性密度存在且满足
\[
\boxed{\;
\lim_{R\to\infty}\frac{N_\varphi(R)}{2R}=1.\;}
\]
特别地，\(\mathcal Z(\mathcal F_\varphi)\) 的指数收敛阶恰为 \(1\)。
\end{corollary}
\begin{proof}
由定理 \ref{thm:xi-fold-resonance-zero-count-selfsimilar-linear}，
\[
N_\varphi(R)=2R+O(\log R),
\]
两边同除以 \(2R\) 并令 \(R\to\infty\)，立即得到
\[
\frac{N_\varphi(R)}{2R}\to 1.
\]

再证指数收敛阶。记 \(\rho\) 为零点集的指数收敛阶。若 \(\sigma>1\)，则由上式存在常数 \(C>0\) 使
\[
N_\varphi(2^{j+1})-N_\varphi(2^j)\le C\,2^j
\qquad(j\gg 1).
\]
故
\[
\sum_{\substack{z\in\mathcal Z(\mathcal F_\varphi)\\ 2^j<|z|\le 2^{j+1}}}|z|^{-\sigma}
\le
C\,2^j\cdot 2^{-j\sigma}
=
C\,2^{-(\sigma-1)j},
\]
从而 \(\sum_{z\in\mathcal Z(\mathcal F_\varphi)\setminus\{0\}}|z|^{-\sigma}\) 收敛。

反之，若 \(0<\sigma<1\)，则同样由 \(N_\varphi(R)=2R+O(\log R)\) 可知存在常数 \(c>0\) 使对充分大的 \(j\) 有
\[
N_\varphi(2^{j+1})-N_\varphi(2^j)\ge c\,2^j.
\]
于是
\[
\sum_{\substack{z\in\mathcal Z(\mathcal F_\varphi)\\ 2^j<|z|\le 2^{j+1}}}|z|^{-\sigma}
\ge
c\,2^j\cdot 2^{-(j+1)\sigma}
\asymp
2^{(1-\sigma)j},
\]
其级数发散。故指数收敛阶正是 \(1\)。证毕。
\end{proof}

\begin{remark}[正弦型强化的障碍]\label{rem:xi-fold-resonance-sine-type-obstruction}
附录中的定理 \ref{thm:fold-resonance-cartwright-indicator} 已经给出
\[
h(\theta)=\pi|\sin\theta|,
\]
因此 \(\mathcal F_\varphi\) 的指示函数并非猜想，而是现有论证链中的已证结论。
另一方面，定理 \ref{thm:fold-resonance-lucas-near-resonance-gap} 与命题 \ref{prop:fold-resonance-min-gap-global-lower} 共同表明
\[
\delta(R)\asymp R^{-1},
\]
故零点并不满足经典正弦型理论所要求的统一分离。
因此，任何进一步的 Paley--Wiener 或采样型强化，都必须在非一致分离的 Cartwright 框架中表述，而不能直接归约为标准 sine-type 整函数。
\end{remark}

\begin{theorem}[Fibonacci 共振整函数对数导数的严格自相似 Riccati 型展开]\label{thm:xi-fold-resonance-logderivative-selfsimilar}
\leanverified{paper\_xi\_fold\_resonance\_logderivative\_selfsimilar}
沿用定理 \ref{thm:fold-resonance-entire-lp} 的记号，并在
\(\CC\setminus\mathcal Z(\mathcal F_\varphi)\) 上定义
\[
\Psi_\varphi(z):=\frac{\mathcal F_\varphi'(z)}{\mathcal F_\varphi(z)}.
\]
则成立严格恒等式
\[
\boxed{\;
\Psi_\varphi(z)=
-\frac{\pi}{\varphi^2}\tan\!\left(\frac{\pi z}{\varphi^2}\right)
+\frac{1}{\varphi}\,\Psi_\varphi\!\left(\frac{z}{\varphi}\right).
\;}
\]
并且对每个远离零点集的紧集，上式可迭代为局部一致收敛的级数表示
\[
\boxed{\;
\Psi_\varphi(z)=
-\pi\sum_{n=2}^{\infty}\varphi^{-n}\tan(\pi z\varphi^{-n}),
\qquad
z\notin \mathcal Z(\mathcal F_\varphi).
\;}
\]
特别地，对每个实数 \(u\in\RR\setminus\mathcal Z(\mathcal F_\varphi)\)，由于
\[
C_\varphi(u)=|\mathcal F_\varphi(u)|,
\]
故有
\[
\boxed{\;
\frac{\dd}{\dd u}\log C_\varphi(u)
=
-\pi\sum_{n=2}^{\infty}\varphi^{-n}\tan(\pi u\varphi^{-n}).
\;}
\]
\end{theorem}
\begin{proof}
由定理 \ref{thm:fold-resonance-entire-lp} 的函数方程
\[
\mathcal F_\varphi(z)=
\cos\!\left(\frac{\pi z}{\varphi^2}\right)
\mathcal F_\varphi\!\left(\frac{z}{\varphi}\right),
\]
在 \(z\notin \mathcal Z(\mathcal F_\varphi)\) 处取对数导数，得到
\[
\Psi_\varphi(z)=
-\frac{\pi}{\varphi^2}\tan\!\left(\frac{\pi z}{\varphi^2}\right)
+\frac{1}{\varphi}\,\Psi_\varphi\!\left(\frac{z}{\varphi}\right).
\]
这就证明了第一式。

迭代 \(N\) 次后得到
\[
\Psi_\varphi(z)=
-\pi\sum_{n=2}^{N+1}\varphi^{-n}\tan(\pi z\varphi^{-n})
+\varphi^{-N}\Psi_\varphi\!\left(\frac{z}{\varphi^{N}}\right).
\]
又因 \(\mathcal F_\varphi\) 为偶函数且 \(\mathcal F_\varphi(0)=1\)，故
\[
\mathcal F_\varphi'(0)=0,
\qquad
\Psi_\varphi(0)=0.
\]
于是当 \(\zeta\to 0\) 时，
\[
\Psi_\varphi(\zeta)=O(\zeta).
\]
因此余项满足
\[
\varphi^{-N}\Psi_\varphi\!\left(\frac{z}{\varphi^{N}}\right)
=
O\!\left(\varphi^{-N}\frac{|z|}{\varphi^{N}}\right)
=
O(\varphi^{-2N})\to 0.
\]
令 \(N\to\infty\) 即得级数表示。

最后，当 \(u\in\RR\setminus \mathcal Z(\mathcal F_\varphi)\) 时，\(\mathcal F_\varphi(u)\) 为非零实数，故
\[
\frac{\dd}{\dd u}\log|\mathcal F_\varphi(u)|
=
\frac{\mathcal F_\varphi'(u)}{\mathcal F_\varphi(u)}
=
\Psi_\varphi(u).
\]
再用
\[
C_\varphi(u)=|\mathcal F_\varphi(u)|
\]
即可得到最后一式。证毕。
\end{proof}

\paragraph{指数衰减 jet、冻结相强逆与指数缩放实零梯的对照闭链}
上述五条结果把两条此前分散的新增主线压缩到同一层级上：一方面，\(\log C_m\) 给出一条指数衰减的 Bernoulli/even-\(\zeta\) jet，其前 \(R\) 阶在有限 Hankel 阶上最终非退化，并可由有限连续窗口指数精度恢复；另一方面，\(\mathcal F_\varphi\) 给出一条指数缩放的实零点梯，其总体零点计数却呈现线性主项，并且局部斜率完全由自相似 tangent 级数控制。夹在二者之间的冻结相预言机反演，则把“可否恢复”精化为严格的强逆指数，从而把阈值现象提升为速率现象。

\begin{theorem}[Lee--Yang 振幅常数 \texorpdfstring{$B$}{B} 的六次性与真正二次覆叠]\label{thm:xi-leyang-amplitude-B-true-quadratic-cover}
沿用定理 \ref{thm:killo-leyang-kappa-square-cubic-field} 与
\ref{thm:killo-leyang-B-square-cubic-and-rational-intertwine} 的记号。
记
\[
u:=\kappa_\star^2,
\qquad
b:=B^2,
\qquad
K:=\QQ(b)=\QQ(u)=\QQ(\lambda_\star).
\]
则
\[
\boxed{\;
[\QQ(B):\QQ]=6,
\qquad
[\QQ(B):\QQ(B^2)]=2,
\qquad
B\notin \QQ(u).\;}
\]
换言之，振幅常数 \(B\) 本身不属于由分支参数 \(u\) 所生成的三次数域，
而是在该三次数域之上的一个真正二次覆叠点。
\end{theorem}
\begin{proof}
由定理 \ref{thm:killo-leyang-B-square-cubic-and-rational-intertwine}，
\[
162b^3+1593b^2+1998b+1184=0
\]
在 \(\QQ\) 上不可约，且
\[
[K:\QQ]=[\QQ(b):\QQ]=3.
\]
把该多项式首一化为
\[
x^3+\frac{1593}{162}x^2+\frac{1998}{162}x+\frac{1184}{162},
\]
于是
\[
N_{K/\QQ}(b)
:=
-\frac{1184}{162}
:=
-\frac{592}{81}<0.
\]
若 \(b\) 在 \(K\) 中为平方，即存在 \(\beta\in K\) 使
\(
b=\beta^2
\)，则
\[
N_{K/\QQ}(b)=N_{K/\QQ}(\beta)^2
\]
必须是 \(\QQ\) 中的平方，特别应为非负有理数，这与上式矛盾。
故 \(b\) 在 \(K\) 中不是平方。
\par
因此多项式
\[
X^2-b\in K[X]
\]
不可约，从而
\[
[K(B):K]=2.
\]
又因 \(B^2=b\in \QQ(B)\)，故
\(
K=\QQ(b)\subseteq \QQ(B)
\)，于是 \(K(B)=\QQ(B)\)。
由塔式公式得到
\[
[\QQ(B):\QQ]=[\QQ(B):K]\,[K:\QQ]=2\cdot 3=6.
\]
同时
\[
[\QQ(B):\QQ(B^2)]=[\QQ(B):\QQ(b)]=[\QQ(B):K]=2.
\]
特别地，若 \(B\in \QQ(u)=K\)，则 \(b=B^2\) 将是 \(K\) 中平方，这与前述矛盾相违。
故 \(B\notin \QQ(u)\)。
证毕。
\end{proof}

\begin{corollary}[分支立方域与振幅平方立方域的共同二次解消子域与判别式平方比]\label{cor:xi-leyang-u-b-square-common-quadratic-resolvent}
沿用定理 \ref{thm:killo-leyang-kappa-square-cubic-field}、
\ref{thm:killo-leyang-B-square-cubic-and-rational-intertwine} 与
\ref{thm:xi-leyang-amplitude-B-true-quadratic-cover} 的记号。
记
\[
\Delta_u:=\Disc\!\bigl(324u^3-540u^2+333u-74\bigr),
\qquad
\Delta_b:=\Disc\!\bigl(162b^3+1593b^2+1998b+1184\bigr).
\]
则
\[
\frac{\Delta_b}{\Delta_u}
:=
\left(\frac{1199}{4}\right)^2.
\]
因而 \(\Delta_u,\Delta_b\) 属于同一平方类，其平方自由部分同为
\(-111\)。特别地，这两条不可约 \(S_3\)-三次模型虽然书写不同，却具有同一个唯一二次解消子域
\[
\boxed{\;
\QQ(\sqrt{-111}).\;}
\]
等价地，它们的 Galois 闭包共享同一二次子域。
\end{corollary}
\begin{proof}
由定理 \ref{thm:killo-leyang-kappa-square-cubic-field} 与
\ref{thm:killo-leyang-B-square-cubic-and-rational-intertwine}，
\[
\Delta_u=-2^{6}\cdot 3^{9}\cdot 37,
\qquad
\Delta_b=-2^{2}\cdot 3^{9}\cdot 11^{2}\cdot 37\cdot 109^{2}.
\]
直接相除得
\[
\frac{\Delta_b}{\Delta_u}
:=
\frac{-2^{2}\cdot 3^{9}\cdot 11^{2}\cdot 37\cdot 109^{2}}
{-2^{6}\cdot 3^{9}\cdot 37}
:=
\frac{11^{2}109^{2}}{2^{4}}
:=
\left(\frac{1199}{4}\right)^2.
\]
故 \(\Delta_u\) 与 \(\Delta_b\) 属于同一平方类。又因
\[
\Delta_u=-(2^3 3^4)^2\cdot 111,
\]
其平方自由部分为 \(-111\)，故 \(\Delta_b\) 的平方自由部分亦同为 \(-111\)。
而对任一不可约 \(S_3\)-三次多项式，其 Galois 闭包中的唯一二次子域恰为
\[
\QQ(\sqrt{\Disc}).
\]
故上述两条三次模型的唯一二次解消子域都等于
\[
\QQ(\sqrt{-111}).
\]
证毕。
\end{proof}
\paragraph{Lee--Yang 三次闭包的二次角色筛选}
在共同二次解消子域确定之后，Chebotarev 意义下的完全分裂首先受到单一二次角色的刚性约束。

\begin{corollary}[Lee--Yang 全分裂素数的二次角色筛选]\label{cor:xi-leyang-split-primes-quadratic-character-filter}
记
\[
K_{\mathrm{LY}}:=\QQ(u)=\QQ(b),
\]
并令 \(L_{\mathrm{LY}}/\QQ\) 为其 Galois 闭包，\(\chi_{-111}\) 为二次域
\(\QQ(\sqrt{-111})/\QQ\) 所对应的二次 Dirichlet 角色。则对任意素数
\[
p\notin\{3,37\},
\]
若 \(p\) 在 \(L_{\mathrm{LY}}\) 中完全分裂，则
\[
\chi_{-111}(p)=\left(\frac{-111}{p}\right)=1.
\]
特别地，除去分歧素数 \(3,37\) 后，Lee--Yang 完全分裂素数必位于该唯一二次角色的核中。
\leanverified{paper\_xi\_leyang\_split\_primes\_quadratic\_character\_filter}
\end{corollary}
\begin{proof}
由推论 \ref{cor:xi-leyang-u-b-square-common-quadratic-resolvent}，
\(L_{\mathrm{LY}}\) 的唯一二次子域为
\[
\QQ(\sqrt{-111}).
\]
若素数 \(p\) 在 \(L_{\mathrm{LY}}\) 中完全分裂，则其 Frobenius 元在
\(\Gal(L_{\mathrm{LY}}/\QQ)\) 中为恒等元；限制到任意中间域后仍为恒等，特别是在二次子域
\(\QQ(\sqrt{-111})\) 上平凡。

对二次扩张而言，未分歧素数上的 Frobenius 平凡当且仅当对应二次角色取值为 \(1\)。由于 \(p\notin\{3,37\}\)，该素数在
\(\QQ(\sqrt{-111})\) 中不分歧，故
\[
\chi_{-111}(p)=\left(\frac{-111}{p}\right)=1.
\]
证毕。
\end{proof}

\begin{theorem}[核版 RH 临界点先于最近复分歧，且边界为纯谱相变]\label{thm:xi-rh-critical-before-nearest-complex-branch}
沿用推论 \ref{cor:real-input-40-collision-branch-radius} 与
命题 \ref{prop:real-input-40-collision-rhk-window} 的记号。
令
\[
P_2(\theta):=\log \rho(e^\theta).
\]
则：
\begin{enumerate}
  \item \(P_2(\theta)\) 的代数分歧点恰为
  \[
  \theta=\log u_j+2\pi i k,\qquad k\in\ZZ,
  \]
  其中 \(u_j\) 遍历三次判别式方程
  \[
  \Delta(u)=4u^3+25u^2+88u+8=0
  \]
  的三个根；
  \item 在 \(\theta=0\) 处的 Taylor 收敛半径为
  \[
  R_\theta=\min_j\min_{k\in\ZZ}\bigl|\log u_j+2\pi i k\bigr|
  \approx 2.76230289346087;
  \]
  \item 核版 RH/Ramanujan 条件的唯一临界参数满足
  \[
  u_R^5+4u_R^4+3u_R^3-96u_R^2-42u_R-14=0,
  \qquad
  \theta_R:=\log u_R\approx 1.27888,
  \]
  且
  \[
  \textsf{RH}(u)\Longleftrightarrow 0<u\le u_R.
  \]
\end{enumerate}
特别地，
\[
\boxed{\;
\theta_R<R_\theta.\;}
\]
并且在临界点 \(u_R\) 处有
\[
\Lambda(u_R)=\lambda_1(u_R)^{1/2},
\]
其中 \(\Lambda(u)\) 为命题 \ref{prop:real-input-40-collision-rhk-window} 的次谱半径。因而存在非空实区间
\[
(\theta_R,R_\theta)\subset \RR_{>0}
\]
使得在该区间上 \(P_2(\theta)\) 仍解析，而 kernel RH 条件已经失效。
换言之，此处 RH 型失效并非由解析边界触发，而是由解析域内部的纯谱交叉触发。
\end{theorem}
\begin{proof}
前三条分别是推论 \ref{cor:real-input-40-collision-branch-radius} 与
命题 \ref{prop:real-input-40-collision-rhk-window} 的直接重述。
故只需比较显式数值：
\[
\theta_R\approx 1.27888
\qquad\text{而}\qquad
R_\theta\approx 2.76230289346087.
\]
于是立即得到
\[
\theta_R<R_\theta.
\]
若 \(\theta\in(\theta_R,R_\theta)\)，则一方面 \(|\theta|<R_\theta\)，故
\(P_2(\theta)\) 仍位于 \(\theta=0\) 的解析收敛盘内；
另一方面 \(u=e^\theta>u_R\)，由
\(
\textsf{RH}(u)\Longleftrightarrow 0<u\le u_R
\)
知 kernel RH 条件已失效。
又由命题 \ref{prop:real-input-40-collision-rhk-window} 的 RH-sharp 等号，边界点满足
\(\Lambda(u_R)=\lambda_1(u_R)^{1/2}\)。故所述“解析仍成立而 RH 已失效”的区间确实非空，并且临界边界由次谱半径与平方根 Perron 尺度的等号刻画，而不是由分支奇点或极点触发。
证毕。
\end{proof}
\paragraph{临界五次方程的伽罗瓦障碍与半平面根分布}
在定理 \ref{thm:xi-rh-critical-before-nearest-complex-branch} 与注 \ref{rem:collision-cubic-rhsharp-quintic} 的阈值多项式
\[
P(u):=u^5+4u^4+3u^3-96u^2-42u-14
\]
之上，还可进一步把唯一正实临界根 \(u_R\) 的代数复杂度与其余四根的稳定半平面位置完全锁定。

\begin{proposition}[临界五次方程的伽罗瓦障碍与 \texorpdfstring{$u_R$}{uR} 的根式不可解性]\label{prop:xi-collision-threshold-quintic-unsolvability}
\leanverified{paper\_xi\_collision\_threshold\_quintic\_unsolvability}
设 \(u_R>1\) 为定理 \ref{thm:xi-rh-critical-before-nearest-complex-branch} 中满足 \(P(u_R)=0\) 的唯一正实根。
则 \(P\) 在 \(\QQ\) 上不可约，且其伽罗瓦群满足
\[
\Gal(P/\QQ)\cong S_5.
\]
因此 \(u_R\) 不能由有理数经有限次四则运算与开方、开立方等根式运算表达。
\end{proposition}
\begin{proof}
先看不可约性。把 \(P\) 模 \(11\) 约化，直接检验可知其在 \(\FF_{11}[u]\) 中既无一次因子也无二次因子。
由于 \(\deg P=5\)，故 \(P\bmod 11\) 必为不可约五次，进而由 Gauss 引理知 \(P\) 在 \(\QQ\) 上不可约。

再看伽罗瓦群。模 \(11\) 不可约说明
\(
G:=\Gal(P/\QQ)\subset S_5
\)
是传递子群且含有一个 \(5\)-循环。
另一方面，模 \(31\) 约化时 \(P\) 分解为一个不可约二次因子与一个不可约三次因子之积。
由 Dedekind 分解型判据，\(G\) 含有循环型为 \((2)(3)\) 的元素；其平方给出一个 \(3\)-循环。

最后计算判别式：
\[
\mathrm{Disc}(P)=2^6\cdot 3^2\cdot 7\cdot 743\cdot 1693^2,
\]
它不是 \(\QQ\) 中的平方，故 \(G\nsubseteq A_5\)。
而 \(S_5\) 的含 \(5\)-循环的传递子群只有
\[
C_5,\qquad D_5,\qquad F_{20},\qquad A_5,\qquad S_5.
\]
其中 \(C_5,D_5,F_{20}\) 均不含 \(3\)-循环，\(A_5\) 又已由判别式非平方排除，
故唯一可能为 \(G=S_5\)。

由 Abel--Ruffini 定理，\(P\) 的根不可能由根式统一表达；特别地，其唯一正实根 \(u_R\) 亦不可由根式表达。证毕。
\end{proof}

\begin{proposition}[临界五次方程的唯一右半平面根与左半平面余根]\label{prop:xi-collision-threshold-quintic-halfplane-root-distribution}
\leanverified{paper\_xi\_collision\_threshold\_quintic\_halfplane\_root\_distribution}
对同一多项式
\[
P(u)=u^5+4u^4+3u^3-96u^2-42u-14,
\]
开右半平面 \(\{\Re(u)>0\}\) 内恰有一个零点；该零点正是 \(u_R\)。
其余四个零点全部满足
\[
\Re(u)<0.
\]
\end{proposition}
\begin{proof}
对 \(P\) 作 Routh 表，其首列为
\[
1,\qquad
4,\qquad
27,\qquad
-\frac{2438}{27},\qquad
-\frac{104069}{2438},\qquad
-14.
\]
该首列无零项且仅发生一次变号。
因此由 Routh--Hurwitz 判据，\(P\) 在开右半平面内恰有一个零点，并且不存在纯虚根。

另一方面，\(P\) 的系数符号序列为
\[
(+,+,+,-,-,-),
\]
仅有一次变号。
由 Descartes 符号法则，\(P\) 恰有一个正实根。
结合定理 \ref{thm:xi-rh-critical-before-nearest-complex-branch} 中对临界参数的唯一性，这个正实根正是 \(u_R\)。

由于 \(P\) 的系数为实数，非实根必成共轭对出现。
既然右半平面内只有一个零点且不存在纯虚根，除 \(u_R\) 之外的其余四根便只能全部落在开左半平面。证毕。
\end{proof}

\paragraph{Fibonacci 标度缺口、共振梯认证、原子 Witt 奇偶手术与黄金分支传输振荡}
下列条目把折叠碰撞缺口、整数共振梯的可认证逼近、真实输入 \(40\) 状态核的原子 Witt 手术、连分分母 tick 上的单箱传输置中、黄金分支的正向缺口，以及 window-\(6\) 根字典的二次等分与线性障碍统一压缩为可独立引用的终端结论。

\begin{theorem}[Fibonacci 标度碰撞缺口的显式正下界]\label{thm:xi-fold-fibonacci-collision-gap-positive-floor}
\leanverified{paper\_xi\_fold\_fibonacci\_collision\_gap\_positive\_floor}
沿用附录 \ref{app:fold-multiplicity} 中 Fold 多重度的记号。记
\[
g_{\mathrm{col}}^2(m):=\mathsf{Col}_m-\frac{1}{F_{m+2}},
\qquad
\overline d_m:=\frac{2^m}{F_{m+2}},
\qquad
M_m:=\max_r d_m(r).
\]
设临界共振常数 \(C_\varphi>0\) 如定理 \ref{thm:fold-critical-resonance-constant}。则有显式下界
\[
\boxed{\;
\liminf_{m\to\infty}F_{m+2}\,g_{\mathrm{col}}^2(m)
\ge
2C_\varphi^2
\approx 0.1007267225141179.\;}
\]
并且峰值纤维满足
\[
\boxed{\;
\liminf_{m\to\infty}\frac{M_m}{\overline d_m}
\ge
1+C_\varphi
\approx 1.2244178274047294.\;}
\]
因此折叠均匀性缺口在 Fibonacci 尺度上不可能比 \(F_{m+2}^{-1}\) 更快消失，而峰值纤维相对均值的偏置亦保持严格正常数。
\end{theorem}
\begin{proof}
第一条即推论 \ref{cor:fold-critical-resonance-gap} 的直接内容；第二条即推论 \ref{cor:fold-max-fiber-constant-bias}。两式合并便得到 Fibonacci 尺度上的 \(L^2\)--\(L^\infty\) 双重刚性。证毕。
\end{proof}

\begin{theorem}[整数共振梯常数的对数复杂度可认证性与最终持久性]\label{thm:xi-fold-resonance-ladder-logarithmic-auditability}
\leanverified{paper\_xi\_fold\_resonance\_ladder\_logarithmic\_auditability}
沿用定义 \ref{def:fold-critical-resonance-ladder} 的记号。对任意非零整数 \(u\in\mathbb Z\setminus\{0\}\)，定义
\[
C_\varphi(u):=\prod_{n=2}^{\infty}\bigl|\cos(\pi u\varphi^{-n})\bigr|,
\qquad
C_{\varphi,N}(u):=\prod_{n=2}^{N}\bigl|\cos(\pi u\varphi^{-n})\bigr|.
\]
并令
\[
k_m(u)\equiv uF_m\pmod{F_{m+2}},
\qquad 0\le k_m(u)<F_{m+2}.
\]
则对任意 \(\varepsilon\in(0,1)\)，只要
\[
N\ge
\max\!\left\{
\left\lceil \log_\varphi(2\pi |u|)-1\right\rceil,
\left\lceil \frac12\log_\varphi
\frac{\pi^2u^2}{(1-\varphi^{-2})\varepsilon}-1\right\rceil
\right\},
\]
便有显式绝对误差界
\[
\boxed{\;0\le C_{\varphi,N}(u)-C_\varphi(u)\le \varepsilon.\;}
\]
从而 \(C_\varphi(u)\) 可由
\[
O\bigl(\log |u|+\log \varepsilon^{-1}\bigr)
\]
个因子完成可审计逼近。并且存在 \(m_0(u)\) 使得对所有 \(m\ge m_0(u)\),
\[
a_m\!\bigl(k_m(u)\bigr)\ge \frac12\,C_\varphi(u)>0.
\]
\end{theorem}
\begin{proof}
记
\[
\eta_N(u):=
\frac{\pi^2u^2\,\varphi^{-2(N+1)}}{1-\varphi^{-2}}.
\]
由命题 \ref{prop:fold-resonance-truncation-error}，
\[
C_{\varphi,N}(u)e^{-\eta_N(u)}\le C_\varphi(u)\le C_{\varphi,N}(u).
\]
故
\[
0\le C_{\varphi,N}(u)-C_\varphi(u)
\le C_{\varphi,N}(u)\bigl(1-e^{-\eta_N(u)}\bigr)
\le \eta_N(u),
\]
其中最后一步使用 \(C_{\varphi,N}(u)\le 1\) 与 \(1-e^{-x}\le x\)。给定的 \(N\) 同时保证
\[
\pi |u|\varphi^{-(N+1)}\le \frac12
\qquad\text{与}\qquad
\eta_N(u)\le \varepsilon,
\]
于是得到绝对误差界。由于所需 \(N\) 由两条对数级上界控制，截断长度满足
\(
N=O(\log|u|+\log\varepsilon^{-1})
\)。
\par
再由定理 \ref{thm:fold-critical-resonance-constant-integer-ladder}，
\[
a_m\!\bigl(k_m(u)\bigr)\longrightarrow C_\varphi(u)>0.
\]
故存在 \(m_0(u)\) 使得对所有 \(m\ge m_0(u)\) 都有
\[
a_m\!\bigl(k_m(u)\bigr)\ge \frac12\,C_\varphi(u).
\]
证毕。
\end{proof}

\begin{theorem}[原子 Witt 因子的偶奇手术刚性]\label{thm:xi-real-input-40-atomic-witt-even-odd-surgery}
对 \(u>0\)，沿用命题 \ref{prop:atomic-2-orbit-split-logM} 的记号，记
\[
\zeta(z,u)=\frac{1}{\Delta(z,u)},
\qquad
\Delta(z,u)=(1-uz^2)\,F(z,u),
\]
\[
P(z,u)=\sum_{n\ge 1}p_n(u)z^n,
\qquad
P_{\mathrm{core}}(z,u)=\sum_{n\ge 1}p_n^{\mathrm{core}}(u)z^n.
\]
任选一个满足
\[
\det(I-zM_{\mathrm{core}}(u))=F(z,u)
\]
的有限维矩阵实现 \(M_{\mathrm{core}}(u)\)。则 primitive 坐标满足精确偶奇手术公式
\[
\boxed{\;
p_n(u)=p_n^{\mathrm{core}}(u)+u\,\mathbf 1_{\{n=2\}}.\;}
\]
原子 Euler 因子的 ghost 坐标满足
\[
a_n^{\mathrm{cycle}}(u)=
\begin{cases}
2u^{n/2},& 2\mid n,\\
0,& 2\nmid n,
\end{cases}
\]
并且对任意 \(n\ge 1\) 有迹差公式
\[
\boxed{\;
\Tr(M(u)^n)-\Tr(M_{\mathrm{core}}(u)^n)
=
(\sqrt u)^n+(-\sqrt u)^n
=
\begin{cases}
2u^{n/2},& 2\mid n,\\
0,& 2\nmid n.
\end{cases}\;}
\]
特别地，所有奇长度 primitive 数据与奇长度迹数据都完全不受该原子 Euler 因子影响；全部修正严格局域在长度 \(2\) 的 primitive 坐标与偶长度的 ghost/trace 层。
\end{theorem}
\begin{proof}
命题 \ref{prop:atomic-2-orbit-split-logM} 给出精确剥离
\[
P(z,u)=u z^2+P_{\mathrm{core}}(z,u),
\]
逐项比较系数即得
\[
p_n(u)=p_n^{\mathrm{core}}(u)+u\,\mathbf 1_{\{n=2\}}.
\]
更一般地，这也是命题 \ref{prop:primitive-orbit-surgery} 在 \(m=1\)、\(E=1\)、\(\ell=2\) 下的特例。
\par
原子 ghost 坐标的显式公式来自注 \ref{rem:real-input-40-single-euler}。最后，推论 \ref{cor:real-input-40-factor-sqrtu-eigs} 说明因子 \((1-uz^2)\) 在谱上对应两条显式特征值 \(\pm\sqrt u\)，故其对 \(n\)-步迹的贡献正是
\[
(\sqrt u)^n+(-\sqrt u)^n,
\]
从而随即得到偶奇分离。证毕。
\end{proof}

\begin{theorem}[连分分母处 \texorpdfstring{$W_1$}{W1} 传输的单箱最优正向置中]\label{thm:xi-kronecker-w1-single-box-optimal-centering}
\leanverified{paper\_xi\_kronecker\_w1\_single\_box\_optimal\_centering}
沿用结论 \ref{con:xi-terminal-kronecker-w1-denominator-one-sided-sensitivity} 的记号。
设 \(p/q\) 为既约有理数且 \(q\ge 2\)，记
\[
\delta:=\alpha-\frac pq.
\]
若满足无跨箱条件
\[
|\delta|<\frac{1}{q(q-1)},
\]
则定理 \ref{thm:kronecker-w1-denominator-closed-form} 的量化耦合计算在该窗口内仍逐字成立，从而
\[
\mathsf{T}_q(\alpha)=
\begin{cases}
\dfrac{1}{2q}+\dfrac{q-1}{2}\left(\dfrac pq-\alpha\right),& \delta<0,\\[2mm]
\dfrac{1}{2q}-\dfrac{q-1}{2}\left(\alpha-\dfrac pq\right)
+\dfrac{q(q-1)(2q-1)}{6}\left(\alpha-\dfrac pq\right)^2,& \delta>0.
\end{cases}
\]
定义
\[
\delta_\ast(q):=\frac{3}{2q(2q-1)}.
\]
则 \(\mathsf{T}_q\) 在整个单箱窗口上的唯一全局极小值恰出现于好侧点
\[
\delta=\delta_\ast(q)>0.
\]
相应最小值为
\[
\boxed{\;
\mathsf{T}_q^{\min}
:=
\mathsf{T}_q\!\left(\frac pq+\delta_\ast(q)\right)
=
\frac1{2q}-\frac{3(q-1)}{8q(2q-1)}.\;}
\]
特别地，
\[
\boxed{\;\mathsf{T}_q^{\min}<\frac1{2q}=\mathsf{T}_q\!\left(\frac pq\right).\;}
\]
因此最优传输置中既不发生在坏侧，也不发生在严格对齐点 \(\alpha=p/q\)，而必须以一个正向箱内偏移实现。
\end{theorem}
\begin{proof}
无跨箱条件给出
\[
(q-1)|\delta|<\frac1q,
\]
故附录定理 \ref{thm:kronecker-w1-denominator-closed-form} 证明中的排序与分位积分计算保持不变，从而同一闭式在该窗口内成立。
当 \(\delta<0\) 时，\(\mathsf{T}_q(\alpha)\) 关于 \(\delta\) 为严格线性函数
\[
\mathsf{T}_q(\alpha)=\frac1{2q}-\frac{q-1}{2}\delta,
\]
故坏侧不存在内点极小值，且其下确界为 \(\delta\uparrow 0\) 时的 \(1/(2q)\)。

当 \(\delta>0\) 时，
\[
\mathsf{T}_q(\alpha)
=
\frac1{2q}-\frac{q-1}{2}\delta
+\frac{q(q-1)(2q-1)}{6}\delta^2,
\]
故
\[
\frac{d^2}{d\delta^2}\mathsf{T}_q(\alpha)=\frac{q(q-1)(2q-1)}3>0,
\]
从而该侧严格凸。唯一临界点由
\[
-\frac{q-1}{2}+\frac{q(q-1)(2q-1)}{3}\delta=0
\]
给出，即
\[
\delta=\delta_\ast(q)=\frac{3}{2q(2q-1)}.
\]
又
\[
\delta_\ast(q)<\frac1{q(q-1)}
\]
等价于 \(3(q-1)<2(2q-1)\)，对一切 \(q\ge 2\) 成立，故该临界点确在窗口内部，并且为好侧唯一极小点。代回得
\[
\mathsf{T}_q^{\min}
=
\frac1{2q}
-\frac{q-1}{2}\cdot\frac{3}{2q(2q-1)}
+\frac{q(q-1)(2q-1)}{6}\cdot\frac{9}{4q^2(2q-1)^2}
=
\frac1{2q}-\frac{3(q-1)}{8q(2q-1)}.
\]
由于减去的第二项严格为正，故 \(\mathsf{T}_q^{\min}<1/(2q)\)。于是全局最优只能发生在好侧正向偏移处。证毕。
\end{proof}

\begin{theorem}[有理阈值处 \texorpdfstring{$W_1$}{W1} 传输曲线的 \texorpdfstring{$C^1$}{C1} 非 \texorpdfstring{$C^2$}{C2} 尖曲率]\label{thm:xi-kronecker-w1-rational-threshold-c1-nonc2-cusp}
沿用定理 \ref{thm:xi-kronecker-w1-single-box-optimal-centering} 的记号。在同一单箱窗口
\[
\left|\alpha-\frac pq\right|<\frac{1}{q(q-1)}
\]
内，函数 \(\alpha\mapsto \mathsf{T}_q(\alpha)\) 在 \(\alpha=p/q\) 处可微，但不属于 \(C^2\)。更准确地，
\[
\partial_-\mathsf{T}_q\!\left(\frac pq\right)
=
\partial_+\mathsf{T}_q\!\left(\frac pq\right)
=
-\frac{q-1}{2},
\]
而左右二阶导数满足
\[
\partial_-^2\mathsf{T}_q\!\left(\frac pq\right)=0,
\qquad
\partial_+^2\mathsf{T}_q\!\left(\frac pq\right)=\frac{q(q-1)(2q-1)}{3}.
\]
因此 \(\alpha=p/q\) 是一条真正的曲率跳跃阈值：一阶几何在两侧无缝拼接，而二阶几何在有理分母处发生不可消去的断裂。
\end{theorem}
\begin{proof}
由定理 \ref{thm:xi-kronecker-w1-single-box-optimal-centering}，当
\(\alpha<p/q\) 时，
\[
\mathsf{T}_q(\alpha)=\frac{1}{2q}+\frac{q-1}{2}\left(\frac pq-\alpha\right),
\]
故左侧导数与二阶导数分别为
\[
\partial_-\mathsf{T}_q\!\left(\frac pq\right)=-\frac{q-1}{2},
\qquad
\partial_-^2\mathsf{T}_q\!\left(\frac pq\right)=0.
\]
当 \(\alpha>p/q\) 时，
\[
\mathsf{T}_q(\alpha)
=
\frac{1}{2q}
-\frac{q-1}{2}\left(\alpha-\frac pq\right)
+\frac{q(q-1)(2q-1)}{6}\left(\alpha-\frac pq\right)^2,
\]
于是
\[
\partial_+\mathsf{T}_q\!\left(\frac pq\right)=-\frac{q-1}{2},
\qquad
\partial_+^2\mathsf{T}_q\!\left(\frac pq\right)=\frac{q(q-1)(2q-1)}{3}.
\]
左右一阶导数相等，故 \(\mathsf{T}_q\) 在 \(\alpha=p/q\) 处属于 \(C^1\)；但左右二阶导数不同，故该点不属于 \(C^2\)。证毕。
\end{proof}


\begin{theorem}[黄金分支 Wasserstein 审计的真正二相极限]\label{thm:xi-golden-w1-true-two-phase-limit}
\leanverified{paper\_xi\_golden\_w1\_true\_two\_phase\_limit}
沿用附录 \ref{app:kronecker-w1-transport} 的记号。设
\[
\alpha=\varphi^{-1},
\qquad
\mu_N:=\frac1N\sum_{t=0}^{N-1}\delta_{\{t\alpha\}},
\qquad
\mathsf{T}_N(\alpha):=W_1(\mu_N,\lambda),
\]
其中 \(\lambda\) 为 \([0,1)\) 上的 Lebesgue 测度。则对任意 Lipschitz 函数 \(f:[0,1)\to\mathbb R\) 有审计不等式
\[
\left|
\frac1N\sum_{t=0}^{N-1}f(\{t\alpha\})-\int_0^1 f(x)\,dx
\right|
\le
\mathrm{Lip}(f)\,\mathsf{T}_N(\alpha).
\]
并且在 Fibonacci 长度 \(N=F_n\) 上，存在两个真正不同的极限常数：
\[
\boxed{\;
\lim_{n\to\infty,\ n\ \mathrm{even}}
F_n\,\mathsf{T}_{F_n}(\varphi^{-1})
=
\frac12+\frac{1}{2\sqrt5}\;},
\]
\[
\boxed{\;
\lim_{n\to\infty,\ n\ \mathrm{odd}}
F_n\,\mathsf{T}_{F_n}(\varphi^{-1})
=
\frac12-\frac{1}{2\sqrt5}+\frac1{15}\;}.
\]
等价地，若记
\[
\delta_n:=\varphi^{-1}-\frac{F_{n-1}}{F_n}
=(-1)^{n-1}\frac{\varphi^{-n}}{F_n},
\]
则 \(\delta_n<0\) 与 \(\delta_n>0\) 分别对应上述两个极限常数。因而
\[
\limsup_{n\to\infty}F_n\,\mathsf{T}_{F_n}(\varphi^{-1})
-
\liminf_{n\to\infty}F_n\,\mathsf{T}_{F_n}(\varphi^{-1})
=
\frac1{\sqrt5}-\frac1{15}>0,
\]
特别地，不存在常数 \(C\) 使
\[
\mathsf{T}_{F_n}(\varphi^{-1})=\frac{C}{F_n}+o(F_n^{-1}).
\]
\end{theorem}
\begin{proof}
Lipschitz 审计不等式即定理 \ref{thm:kronecker-w1-lipschitz-readout}。两条 Fibonacci 子序列的极限常数是推论 \ref{cor:kronecker-w1-fibonacci-limits} 的直接内容；而引理 \ref{lem:golden-convergent-error-closed} 给出
\[
\delta_n=(-1)^{n-1}\frac{\varphi^{-n}}{F_n},
\]
故偶数 \(n\) 与奇数 \(n\) 的符号分相恰与两条交错子序列对应。若存在单一常数 \(C\) 使
\[
\mathsf{T}_{F_n}(\varphi^{-1})=\frac{C}{F_n}+o(F_n^{-1}),
\]
则 \(F_n\mathsf{T}_{F_n}(\varphi^{-1})\) 全序列应收敛到 \(C\)，这与两条子序列极限不同矛盾。证毕。
\end{proof}

\leanverified{paper\_xi\_golden\_w1\_star\_certificate\_parity\_splitting}
\begin{theorem}[黄金分支上 \texorpdfstring{$W_1$}{W1} 证书与星差证书的奇偶分裂常数]\label{thm:xi-golden-w1-star-certificate-parity-splitting}
沿用附录 \ref{app:kronecker-w1-transport} 的记号。令
\[
D_n^\ast:=D_{F_n}^\ast(P_{F_n}(\varphi^{-1})),
\qquad
\mathsf{T}_n:=\mathsf{T}_{F_n}(\varphi^{-1}).
\]
则当 \(n\) 为偶数时，有精确恒等式
\[
\boxed{\;
\frac{\mathsf{T}_n}{D_n^\ast}=\frac12.\;}
\]
当 \(n\) 为奇数时，则有极限常数
\[
\boxed{\;
\lim_{\substack{n\to\infty\\ n\ \mathrm{odd}}}
\frac{\mathsf{T}_n}{D_n^\ast}
=
\frac12-\frac{1}{2\sqrt5}+\frac1{15}.\;}
\]
等价地，在奇数子序列上成立精确缺口
\[
\boxed{\;
\lim_{\substack{n\to\infty\\ n\ \mathrm{odd}}}
F_n\bigl(D_n^\ast-2\mathsf{T}_n\bigr)
=
\frac1{\sqrt5}-\frac2{15}>0.\;}
\]
因此，Fibonacci 分母尺度上的两类误差证书出现严格奇偶分相：坏侧恰取半因子，而好侧则收敛到更小的显式常数。
\end{theorem}
\begin{proof}
由结论 \ref{con:xi-terminal-kronecker-w1-denominator-one-sided-sensitivity} 与引理 \ref{lem:golden-convergent-error-closed}，
\[
\delta_n:=\varphi^{-1}-\frac{F_{n-1}}{F_n}
=(-1)^{n-1}\frac{\varphi^{-n}}{F_n},
\]
故 \(\delta_n<0\) 当且仅当 \(n\) 为偶数，而 \(\delta_n>0\) 当且仅当 \(n\) 为奇数。

当 \(n\) 为偶数时，黄金分支处于坏侧，因而由结论 \ref{con:xi-terminal-kronecker-w1-denominator-one-sided-sensitivity} 直接得到
\[
\mathsf{T}_n=\frac12\,D_n^\ast.
\]

当 \(n\) 为奇数时，黄金分支处于好侧，于是同一结论给出
\[
D_n^\ast=\frac1{F_n}.
\]
另一方面，推论 \ref{cor:kronecker-w1-fibonacci-limits} 给出
\[
F_n\mathsf{T}_n
\longrightarrow
\frac12-\frac{1}{2\sqrt5}+\frac1{15}
\qquad
\bigl(n\to\infty,\ n\ \mathrm{odd}\bigr).
\]
将 \(D_n^\ast=F_n^{-1}\) 代入，便得到
\[
\frac{\mathsf{T}_n}{D_n^\ast}
=
F_n\mathsf{T}_n
\longrightarrow
\frac12-\frac{1}{2\sqrt5}+\frac1{15}.
\]
最后，
\[
F_n\bigl(D_n^\ast-2\mathsf{T}_n\bigr)
=
1-2F_n\mathsf{T}_n
\longrightarrow
1-2\left(\frac12-\frac{1}{2\sqrt5}+\frac1{15}\right)
=
\frac1{\sqrt5}-\frac2{15},
\]
即得所示缺口公式。证毕。
\end{proof}

\begin{corollary}[黄金分支在好侧未达单箱最优置中且保有显式正缺口]\label{cor:xi-golden-w1-good-side-gap-to-single-box-optimum}\leanverified{paper\_xi\_golden\_w1\_good\_side\_gap\_to\_single\_box\_optimum}
沿用定理 \ref{thm:xi-kronecker-w1-single-box-optimal-centering} 与
定理 \ref{thm:xi-golden-w1-true-two-phase-limit} 的记号。取
\[
\alpha=\varphi^{-1},
\qquad
q=F_n,
\qquad
p=F_{n-1},
\qquad
\delta_n:=\alpha-\frac{F_{n-1}}{F_n}
=
\frac{(-1)^{n-1}}{\sqrt5\,F_n^2}\Bigl(1-(-1)^n\varphi^{-2n}\Bigr).
\]
记 \(\mathsf{T}_{F_n}^{\min}\) 为定理 \ref{thm:xi-kronecker-w1-single-box-optimal-centering} 在 \(q=F_n\) 处给出的理论最小值。则在好侧子序列 \(n\) 为奇数时：
\begin{enumerate}
  \item 存在 \(n_0\) 使得对一切奇数 \(n\ge n_0\)，都有
  \[
  \delta_n<\delta_\ast(F_n).
  \]
  因而黄金分支的真实箱内偏移始终落在好侧抛物线的下降支上，尚未达到单箱最优置中点。
  \item 存在显式正极限缺口
  \[
  \boxed{\;
  \lim_{\substack{n\to\infty\\ n\ \mathrm{odd}}}
  F_n\Bigl(\mathsf{T}_{F_n}(\varphi^{-1})-\mathsf{T}_{F_n}^{\min}\Bigr)
  =
  \frac{61}{240}-\frac{1}{2\sqrt5}
  >0.\;}
  \]
\end{enumerate}
\end{corollary}
\begin{proof}
当 \(n\) 为奇数时，由引理 \ref{lem:golden-convergent-error-closed}，
\[
F_n^2\delta_n
=
\frac{1}{\sqrt5}\bigl(1+\varphi^{-2n}\bigr)
\longrightarrow \frac1{\sqrt5}.
\]
另一方面，
\[
\delta_\ast(F_n)=\frac{3}{2F_n(2F_n-1)},
\qquad
F_n^2\delta_\ast(F_n)=\frac{3F_n}{2(2F_n-1)}
\longrightarrow \frac34.
\]
由于 \(1/\sqrt5<3/4\)，故最终必有 \(\delta_n<\delta_\ast(F_n)\)，得到第一条。

再由定理 \ref{thm:xi-kronecker-w1-single-box-optimal-centering},
\[
F_n\,\mathsf{T}_{F_n}^{\min}
=
\frac12-\frac{3(F_n-1)}{8(2F_n-1)}
\longrightarrow
\frac12-\frac{3}{16}
=
\frac{5}{16}.
\]
而定理 \ref{thm:xi-golden-w1-true-two-phase-limit} 给出
\[
F_n\,\mathsf{T}_{F_n}(\varphi^{-1})
\longrightarrow
\frac12-\frac{1}{2\sqrt5}+\frac1{15}
\qquad
(n\to\infty,\ n\ \mathrm{odd}).
\]
两式相减便得
\[
\lim_{\substack{n\to\infty\\ n\ \mathrm{odd}}}
F_n\Bigl(\mathsf{T}_{F_n}(\varphi^{-1})-\mathsf{T}_{F_n}^{\min}\Bigr)
=
\left(\frac12-\frac{1}{2\sqrt5}+\frac1{15}\right)-\frac{5}{16}
=
\frac{61}{240}-\frac{1}{2\sqrt5}.
\]
该常数严格为正，故得到第二条。证毕。
\end{proof}


\begin{corollary}[黄金分支 Fibonacci 审计常数的双相算术平均恰为 \texorpdfstring{$8/15$}{8/15}]\label{cor:xi-golden-w1-two-phase-arithmetic-mean-eight-fifteenths}
\leanverified{paper\_xi\_golden\_w1\_two\_phase\_arithmetic\_mean\_eight\_fifteenths}
沿用定理 \ref{thm:xi-golden-w1-true-two-phase-limit} 的记号。则黄金分支在 Fibonacci 分母长度上的两相极限常数之算术平均满足
\[
\boxed{\;
\frac12\lim_{\substack{n\to\infty\\ n\ \mathrm{even}}}
F_n\,\mathsf{T}_{F_n}(\varphi^{-1})
+
\frac12\lim_{\substack{n\to\infty\\ n\ \mathrm{odd}}}
F_n\,\mathsf{T}_{F_n}(\varphi^{-1})
=
\frac{8}{15}.\;}
\]
等价地，
\[
\boxed{\;
\frac12\left(\limsup_{n\to\infty}F_n\,\mathsf{T}_{F_n}(\varphi^{-1})
+
\liminf_{n\to\infty}F_n\,\mathsf{T}_{F_n}(\varphi^{-1})\right)
=
\frac{8}{15}.\;}
\]
\end{corollary}
\begin{proof}
由定理 \ref{thm:xi-golden-w1-true-two-phase-limit}，
\[
\lim_{\substack{n\to\infty\\ n\ \mathrm{even}}}
F_n\,\mathsf{T}_{F_n}(\varphi^{-1})
=
\frac12+\frac{1}{2\sqrt5},
\]
\[
\lim_{\substack{n\to\infty\\ n\ \mathrm{odd}}}
F_n\,\mathsf{T}_{F_n}(\varphi^{-1})
=
\frac12-\frac{1}{2\sqrt5}+\frac1{15}.
\]
两式取算术平均，立即得到
\[
\frac12\left(\frac12+\frac{1}{2\sqrt5}\right)
+
\frac12\left(\frac12-\frac{1}{2\sqrt5}+\frac1{15}\right)
=
\frac{8}{15}.
\]
又因该定理已给出两子列分别对应于全序列的 \(\limsup\) 与 \(\liminf\)，故第二个盒装公式只是同一恒等式的重述。证毕。
\end{proof}

\begin{corollary}[黄金分支 Fibonacci 传输常数的双相振幅与中心值]\label{cor:xi-golden-w1-two-phase-amplitude-center}
\leanverified{paper\_xi\_golden\_w1\_two\_phase\_amplitude\_center}
沿用定理 \ref{thm:xi-golden-w1-true-two-phase-limit} 的记号。则数列
\[
F_n\,\mathsf{T}_{F_n}(\varphi^{-1})
\]
不收敛，并且恰有两个聚点
\[
\frac12+\frac{1}{2\sqrt5},
\qquad
\frac12-\frac{1}{2\sqrt5}+\frac1{15}.
\]
特别地，
\[
\boxed{\;
\limsup_{n\to\infty}F_n\,\mathsf{T}_{F_n}(\varphi^{-1})
=
\frac12+\frac{1}{2\sqrt5},\;}
\]
\[
\boxed{\;
\liminf_{n\to\infty}F_n\,\mathsf{T}_{F_n}(\varphi^{-1})
=
\frac12-\frac{1}{2\sqrt5}+\frac1{15}.\;}
\]
若记双相振幅
\[
\Omega_\varphi
:=
\limsup_{n\to\infty}F_n\,\mathsf{T}_{F_n}(\varphi^{-1})
-
\liminf_{n\to\infty}F_n\,\mathsf{T}_{F_n}(\varphi^{-1}),
\]
则
\[
\boxed{\;
\Omega_\varphi
=
\frac1{\sqrt5}-\frac1{15}.\;}
\]
同时，两聚点的算术中心满足
\[
\boxed{\;
\frac12\left(\limsup_{n\to\infty}F_n\,\mathsf{T}_{F_n}(\varphi^{-1})
+
\liminf_{n\to\infty}F_n\,\mathsf{T}_{F_n}(\varphi^{-1})\right)
=
\frac{8}{15}.\;}
\]
\end{corollary}
\begin{proof}
由定理 \ref{thm:xi-golden-w1-true-two-phase-limit}，偶指标与奇指标 Fibonacci 子序列分别收敛到
\[
\frac12+\frac{1}{2\sqrt5},
\qquad
\frac12-\frac{1}{2\sqrt5}+\frac1{15},
\]
且二者不同，故全序列不收敛而恰有这两个聚点；相应的 \(\limsup\) 与 \(\liminf\) 公式随即成立。

两式相减得到
\[
\Omega_\varphi
=
\left(\frac12+\frac{1}{2\sqrt5}\right)
-
\left(\frac12-\frac{1}{2\sqrt5}+\frac1{15}\right)
=
\frac1{\sqrt5}-\frac1{15}.
\]
而两式的算术平均恰为推论
\ref{cor:xi-golden-w1-two-phase-arithmetic-mean-eight-fifteenths} 的内容，因此中心值为 \(8/15\)。证毕。
\end{proof}

\begin{theorem}[黄金分支在金属平均族中唯一最大化单位符号预算下的分母熵率]\label{thm:xi-golden-metallic-symbol-budget-entropy-optimality}
\leanverified{paper\_xi\_golden\_metallic\_symbol\_budget\_entropy\_optimality}
沿用定义 \ref{def:const-type} 与定理 \ref{thm:metallic-gap} 的记号。对每个整数 \(A\ge 1\)，令
\[
\alpha_A=[0;A,A,A,\dots],
\]
并记其第 \(n\) 个收敛分母为 \(q_n(A)\)。定义单位符号预算下的分母熵率
\[
\eta(A):=
\lim_{n\to\infty}\frac{1}{An}\log q_n(A).
\]
则极限存在，且有显式公式
\[
\boxed{\;
\eta(A)=\frac{\log\lambda_A}{A}=\kappa(A)^{-1}.\;}
\]
因此函数 \(A\mapsto \eta(A)\) 在 \(A\ge 1\) 上严格递减，并在黄金分支 \(A=1\) 处取得唯一最大值
\[
\boxed{\;
\eta_{\max}=\eta(1)=\log\varphi.\;}
\]
换言之，在常数型金属平均族中，黄金分支以最少的单位符号预算实现最大的分母指数增长率；任意 \(A\ge 2\) 都严格逊于 \(A=1\)。
\end{theorem}
\begin{proof}
常数型连分数的分母满足标准递推
\[
q_{n+1}(A)=Aq_n(A)+q_{n-1}(A),
\]
其特征多项式为
\[
X^2-AX-1.
\]
由定义 \ref{def:const-type}，该多项式的 Perron 根正是
\[
\lambda_A=\frac{A+\sqrt{A^2+4}}{2}.
\]
因此存在常数 \(c_A\neq 0\) 使
\[
q_n(A)=c_A\lambda_A^n+O(\lambda_A^{-n}),
\]
从而
\[
\lim_{n\to\infty}\frac{1}{n}\log q_n(A)=\log\lambda_A.
\]
两边再除以 \(A\)，便得到
\[
\eta(A)=\frac{\log\lambda_A}{A}.
\]
而定理 \ref{thm:metallic-gap} 已知
\[
\kappa(A):=\frac{A}{\log\lambda_A}
\]
在 \(A\ge 1\) 上严格递增，故
\[
\eta(A)=\kappa(A)^{-1}
\]
必严格递减，并在最小允许值 \(A=1\) 处取得唯一最大值。最后，由
\[
\lambda_1=\varphi
\]
即得
\[
\eta(1)=\log\varphi.
\]
证毕。
\end{proof}

\begin{corollary}[金属平均族的逐位指数效率由黄金端点唯一极大化]\label{cor:xi-golden-metallic-perdigit-exponential-efficiency}
\leanverified{paper\_xi\_golden\_metallic\_perdigit\_exponential\_efficiency}
沿用定义 \ref{def:const-type} 与定理 \ref{thm:metallic-gap} 的记号，并把
\[
h(A):=\frac{\log\lambda_A}{A}
\qquad (A\in[1,\infty))
\]
视为实变量 \(A\) 上的解析函数。则：
\begin{enumerate}
  \item \(h(A)\) 在 \([1,\infty)\) 上严格递减；
  \item 等价地，\(\lambda_A^{1/A}\) 在 \([1,\infty)\) 上严格递减；
  \item 对每个 \(A>1\) 都有严格不等式
  \[
  \lambda_A<\varphi^A,
  \]
  特别地，对所有整数 \(A\ge 2\) 都成立。
\end{enumerate}
因此在全部常数型金属平均口径中，黄金端点 \(A=1\) 唯一最大化“每个连分位数所产生的指数增长效率”。
\end{corollary}
\begin{proof}
由定理 \ref{thm:metallic-gap}，
\[
\kappa(A):=\frac{A}{\log\lambda_A}
\]
在 \([1,\infty)\) 上严格递增且恒为正。故其倒数
\[
h(A)=\kappa(A)^{-1}=\frac{\log\lambda_A}{A}
\]
严格递减，得到第一条。

对正函数 \(h(A)\) 逐点取指数，便得
\[
\lambda_A^{1/A}=e^{h(A)}
\]
在 \([1,\infty)\) 上亦严格递减，得到第二条。

最后，由 \(\lambda_1=\varphi\) 与第二条，对任意 \(A>1\) 都有
\[
\lambda_A^{1/A}<\lambda_1=\varphi.
\]
两边取 \(A\) 次幂即得
\[
\lambda_A<\varphi^A.
\]
证毕。
\end{proof}


\begin{theorem}[输出位势的一阶 Edgeworth 左偏厚尾非高斯性]\label{thm:xi-real-input-40-output-edgeworth-nongaussian}
沿用命题 \ref{prop:real-input-40-output-cumulants-closed} 与注 \ref{rem:real-input-40-output-cumulants-clt} 的记号。令
\[
u=e^\theta,
\qquad
P(\theta)=\log \lambda(e^\theta).
\]
则 \(P^{(k)}(0)\in \QQ(\sqrt5)\)，且前四阶累积量密度满足
\[
P''(0)=\frac{1791}{200}-\frac{3983\sqrt5}{1000}
\approx 0.0487412456183376392>0,
\]
\[
P^{(3)}(0)=\frac{1416369}{5000}-\frac{126691\sqrt5}{1000}
\approx -0.0158881374258564278<0,
\]
\[
P^{(4)}(0)=\frac{1354428303}{100000}-\frac{605718497\sqrt5}{100000}
\approx 0.00568478997567274949>0.
\]
定义标准化第三、第四累积量
\[
\gamma_1:=\frac{P^{(3)}(0)}{(P''(0))^{3/2}},
\qquad
\gamma_2:=\frac{P^{(4)}(0)}{(P''(0))^2}.
\]
则
\[
\boxed{\;
\gamma_1\approx -1.476481644016533,
\qquad
\gamma_2\approx 2.3928814165461985.\;}
\]
因此输出计数统计在一阶 Edgeworth 修正层面严格左偏且严格厚尾；任何对称或近高斯的首项修正都与该符号指纹不相容。
\end{theorem}
\begin{proof}
命题 \ref{prop:real-input-40-output-cumulants-closed} 与表 \ref{tab:real_input_40_output_potential_cumulants_closed} 给出 \(P''(0)\)、\(P^{(3)}(0)\)、\(P^{(4)}(0)\) 的闭式。将其直接代入 \(\gamma_1,\gamma_2\) 的定义即可得到所示数值。由于
\[
P''(0)>0,\qquad P^{(3)}(0)<0,\qquad P^{(4)}(0)>0,
\]
而注 \ref{rem:real-input-40-output-cumulants-clt} 已说明这些量正是 CLT/Edgeworth 首项修正的可审计输入常数，故结论成立。证毕。
\end{proof}

\begin{theorem}[输出计数的准幂极限定律与显式一阶 Edgeworth 系数]\label{thm:xi-real-input-40-output-edgeworth-quasipowers}
沿用推论 \ref{cor:real-input-40-ground-entropy} 与命题 \ref{prop:real-input-40-output-cumulants-closed} 的记号。令
\[
Z_n(u)=\mathbf{1}^{\mathsf T}M(u)^n\mathbf{1}
=\sum_{k=0}^n A_{n,k}u^k,
\qquad
\mathbb P(X_n=k)=\frac{A_{n,k}}{Z_n(1)}.
\]
再记
\[
\mu:=P'(0),
\qquad
\kappa_r:=P^{(r)}(0),
\qquad
\widetilde X_n:=\frac{X_n-n\mu}{\sqrt{n\kappa_2}}.
\]
则 \(\widetilde X_n\Longrightarrow \mathcal N(0,1)\)。并且在中心窗口内一致成立
\[
\mathbb P(\widetilde X_n\le x)
=
\Phi(x)
+\frac{\kappa_3}{6\kappa_2^{3/2}\sqrt n}(1-x^2)\phi(x)
+O(n^{-1}).
\]
特别地，
\[
\frac{\kappa_3}{6\kappa_2^{3/2}}
\approx -0.2460802740027555.
\]
此外，标准化偏度与超额峰度满足
\[
\mathrm{Skew}(X_n)
=
\frac{\kappa_3}{\kappa_2^{3/2}}\frac{1}{\sqrt n}
+O(n^{-1})
=
-\frac{1.4764816440165331}{\sqrt n}+O(n^{-1}),
\]
\[
\mathrm{ExKurt}(X_n)
=
\frac{\kappa_4}{\kappa_2^2}\frac{1}{n}
+O(n^{-3/2})
=
\frac{2.3928814165461985}{n}+O(n^{-3/2}).
\]
因此第一 Edgeworth 修正的符号与尺度都被单一代数常数 \(\kappa_3/(6\kappa_2^{3/2})\) 完全锁定；而离散分位点的精细描述仍需再附加独立的 lattice 校正。
\end{theorem}
\begin{proof}
由主特征值在 \(u=1\) 处的单性与解析扰动，存在 \(\theta=0\) 的复邻域、解析且不为零的函数 \(C(\theta)\) 与常数 \(0<\beta<1\)，使得
\[
Z_n(e^\theta)
=
C(\theta)e^{nP(\theta)}\bigl(1+O(\beta^n)\bigr)
\]
在该邻域内一致成立。于是
\[
\log \mathbb E[e^{\theta X_n}]
=
\log Z_n(e^\theta)-\log Z_n(1)
=
n\bigl(P(\theta)-P(0)\bigr)
+\log\frac{C(\theta)}{C(0)}
+O(\beta^n).
\]
这正是附录 \ref{app:pressure-analytic} 中有限状态解析加权矩阵的标准准幂情形，故中心极限定理与一阶 Edgeworth 展开在中心窗口内成立；其首项系数恰为
\(
\kappa_3/(6\kappa_2^{3/2})
\)。

另一方面，由上式在 \(\theta=0\) 处逐阶求导可得
\[
\operatorname{Var}(X_n)=n\kappa_2+O(1),
\qquad
\kappa_3(X_n)=n\kappa_3+O(1),
\qquad
\kappa_4(X_n)=n\kappa_4+O(1),
\]
从而
\[
\mathrm{Skew}(X_n)
=
\frac{\kappa_3(X_n)}{\operatorname{Var}(X_n)^{3/2}}
=
\frac{\kappa_3}{\kappa_2^{3/2}}\frac{1}{\sqrt n}
+O(n^{-1}),
\]
\[
\mathrm{ExKurt}(X_n)
=
\frac{\kappa_4(X_n)}{\operatorname{Var}(X_n)^2}
=
\frac{\kappa_4}{\kappa_2^2}\frac{1}{n}
+O(n^{-3/2}).
\]
最后，把命题 \ref{prop:real-input-40-output-cumulants-closed} 的闭式代入，并用定理 \ref{thm:xi-real-input-40-output-edgeworth-nongaussian} 中已给出的数值化简，即得各个常数。证毕。
\end{proof}

\begin{theorem}[window-\texorpdfstring{$6$}{6} 的 \texorpdfstring{$C_3$}{C3} 根字典满足二次矩各向同性与二次能量等分]\label{thm:xi-window6-c3-quadratic-energy-equipartition}
沿用定理 \ref{thm:xi-window6-center-exact-splitting-bdry-cyclic} 的 window-\(6\) 记号，并令
\[
\iota_C:X_6^{\mathrm{cyc}}\to \mathbb R^3
\]
为表 \ref{tab:fold6_b3c3_root_dictionary_c3} 给出的 \(C_3\) 根字典。记
\[
L:=\{w\in X_6^{\mathrm{cyc}}:\ \mathrm{wt}(w)=1\},
\qquad
S:=X_6^{\mathrm{cyc}}\setminus L.
\]
则 \(L\) 恰对应六个长根
\[
\{\pm(2,0,0),\ \pm(0,2,0),\ \pm(0,0,2)\},
\]
而 \(S\) 恰对应十二个短根
\[
\{(\pm1,\pm1,0),\ (\pm1,0,\pm1),\ (0,\pm1,\pm1)\}.
\]
并且成立精确二次框架恒等式
\[
\boxed{\;
\sum_{w\in L}\iota_C(w)\iota_C(w)^{\mathsf T}=8I_3,\qquad
\sum_{w\in S}\iota_C(w)\iota_C(w)^{\mathsf T}=8I_3.\;}
\]
从而
\[
\boxed{\;
\sum_{w\in X_6^{\mathrm{cyc}}}\iota_C(w)\iota_C(w)^{\mathsf T}=16I_3.\;}
\]
等价地，
\[
\boxed{\;
\frac1{18}\sum_{w\in X_6^{\mathrm{cyc}}}\iota_C(w)\iota_C(w)^{\mathsf T}
=
\frac89\,I_3.\;}
\]
并且对任意 \(v\in\mathbb R^3\) 有
\[
\boxed{\;
\frac1{18}\sum_{w\in X_6^{\mathrm{cyc}}}\langle \iota_C(w),v\rangle^2
=
\frac89\,\|v\|_2^2.\;}
\]
等价地，对任意实对称矩阵 \(A\in M_3(\mathbb R)\)，若记 \(Q_A(v):=v^{\mathsf T}Av\)，则
\[
\boxed{\;
\sum_{w\in L}Q_A(\iota_C(w))
=
\sum_{w\in S}Q_A(\iota_C(w))
=
8\,\operatorname{tr}(A).\;}
\]
因此尽管两层分别具有 \(6\) 与 \(12\) 个向量、半径分别为 \(2\) 与 \(\sqrt2\)，它们对全部二次可观测量的总贡献完全相等；特别地，cyclic 层在该根字典下具有严格的二阶各向同性。
\end{theorem}
\begin{proof}
由表 \ref{tab:fold6_b3c3_root_dictionary_c3}，集合 \(L\) 正是
\[
\{\pm 2e_1,\ \pm 2e_2,\ \pm 2e_3\}.
\]
于是
\[
\sum_{w\in L}\iota_C(w)\iota_C(w)^{\mathsf T}
=
2\cdot(2e_1)(2e_1)^{\mathsf T}
+2\cdot(2e_2)(2e_2)^{\mathsf T}
+2\cdot(2e_3)(2e_3)^{\mathsf T}
=
8I_3.
\]

对集合 \(S\)，同一张表给出三组短根：
\[
\{(\pm1,\pm1,0)\},
\qquad
\{(\pm1,0,\pm1)\},
\qquad
\{(0,\pm1,\pm1)\}.
\]
对第一组有
\[
\sum_{(\varepsilon_1,\varepsilon_2)\in\{\pm1\}^2}
(\varepsilon_1 e_1+\varepsilon_2 e_2)
(\varepsilon_1 e_1+\varepsilon_2 e_2)^{\mathsf T}
=
4\bigl(e_1e_1^{\mathsf T}+e_2e_2^{\mathsf T}\bigr),
\]
因为所有混合项都按符号对称相消。其余两组同理，分别给出
\[
4\bigl(e_1e_1^{\mathsf T}+e_3e_3^{\mathsf T}\bigr),
\qquad
4\bigl(e_2e_2^{\mathsf T}+e_3e_3^{\mathsf T}\bigr).
\]
三式相加即得
\[
\sum_{w\in S}\iota_C(w)\iota_C(w)^{\mathsf T}=8I_3.
\]
总和公式随即成立。两边再除以 \(18\)，便得到归一化二次矩公式。

对任意 \(v\in\mathbb R^3\)，由
\[
\langle \iota_C(w),v\rangle^2
=
v^{\mathsf T}\bigl(\iota_C(w)\iota_C(w)^{\mathsf T}\bigr)v
\]
与总和公式
\[
\sum_{w\in X_6^{\mathrm{cyc}}}\iota_C(w)\iota_C(w)^{\mathsf T}=16I_3
\]
可得
\[
\sum_{w\in X_6^{\mathrm{cyc}}}\langle \iota_C(w),v\rangle^2
=
v^{\mathsf T}(16I_3)v
=
16\|v\|_2^2.
\]
再除以 \(18\) 即得
\[
\frac1{18}\sum_{w\in X_6^{\mathrm{cyc}}}\langle \iota_C(w),v\rangle^2
=
\frac89\,\|v\|_2^2.
\]

最后，对任意实对称矩阵 \(A\)，利用
\[
Q_A(v)=\operatorname{tr}\!\bigl(A\,vv^{\mathsf T}\bigr)
\]
与迹的线性性，可得
\[
\sum_{w\in L}Q_A(\iota_C(w))
=
\operatorname{tr}\!\left(
A\sum_{w\in L}\iota_C(w)\iota_C(w)^{\mathsf T}
\right)
=
\operatorname{tr}(A\cdot 8I_3)
=
8\,\operatorname{tr}(A),
\]
而 \(S\) 的情形完全相同。证毕。
\end{proof}

\begin{corollary}[window-\texorpdfstring{$6$}{6} 的 \texorpdfstring{$C_3$}{C3} 根字典不存在非零一次不变量]\label{cor:xi-window6-c3-no-nonzero-linear-invariant}
沿用定理 \ref{thm:xi-window6-c3-quadratic-energy-equipartition} 的记号，则
\[
\sum_{w\in X_6^{\mathrm{cyc}}}\iota_C(w)=0.
\]
因此对任意线性泛函 \(\ell:\mathbb R^3\to\mathbb R\)，都有
\[
\sum_{w\in X_6^{\mathrm{cyc}}}\ell\bigl(\iota_C(w)\bigr)=0.
\]
\end{corollary}
\begin{proof}
由表 \ref{tab:fold6_b3c3_root_dictionary_c3}，\(X_6^{\mathrm{cyc}}\) 中的每个向量都与其相反向量成对出现：长根部分为
\[
\{\pm(2,0,0),\ \pm(0,2,0),\ \pm(0,0,2)\},
\]
短根部分为
\[
\{(\pm1,\pm1,0),\ (\pm1,0,\pm1),\ (0,\pm1,\pm1)\}.
\]
故逐对相加即得
\[
\sum_{w\in X_6^{\mathrm{cyc}}}\iota_C(w)=0.
\]
对任意线性泛函 \(\ell\)，再由线性性立得
\[
\sum_{w\in X_6^{\mathrm{cyc}}}\ell\bigl(\iota_C(w)\bigr)
=
\ell\!\left(\sum_{w\in X_6^{\mathrm{cyc}}}\iota_C(w)\right)
=
0.
\]
证毕。
\end{proof}
\leanverified{paper\_xi\_window6\_c3\_no\_nonzero\_linear\_invariant}


\begin{theorem}[window-\texorpdfstring{$6$}{6} 的 \texorpdfstring{$B_3/C_3$}{B3/C3} 根系字典不存在全局线性统一器]\label{thm:xi-window6-b3c3-no-global-linear-unifier}
\leanverified{paper\_xi\_window6\_b3c3\_no\_global\_linear\_unifier}
沿用定理 \ref{thm:xi-window6-parity-charge-trigrading-b3} 的分解记号。记
\[
|X_6|=21,\qquad |X_6^{\mathrm{cyc}}|=18,\qquad |X_6^{\mathrm{bdry}}|=3,
\]
\[
X_6^{\mathrm{bdry}}
=
\{\texttt{100001},\texttt{100101},\texttt{101001}\}.
\]
令
\[
\iota_B:X_6^{\mathrm{cyc}}\to \mathbb R^3,
\qquad
\iota_C:X_6^{\mathrm{cyc}}\to \mathbb R^3
\]
分别为表 \ref{tab:fold6_b3c3_root_dictionary_b3} 与表 \ref{tab:fold6_b3c3_root_dictionary_c3} 给出的 \(B_3\) 与 \(C_3\) 根字典。则不存在线性映射
\[
T\in \mathrm{End}(\mathbb R^3)
\]
使得
\[
T\bigl(\iota_B(w)\bigr)=\iota_C(w)
\qquad
\forall\,w\in X_6^{\mathrm{cyc}}.
\]
\end{theorem}
\begin{proof}
由表 \ref{tab:fold6_b3c3_root_dictionary_b3} 与表 \ref{tab:fold6_b3c3_root_dictionary_c3} 的 weight-\(1\) 词条，
\[
\iota_B(\texttt{100000})=(1,0,0),\quad
\iota_B(\texttt{010000})=(0,1,0),\quad
\iota_B(\texttt{001000})=(0,0,1),
\]
\[
\iota_C(\texttt{100000})=(2,0,0),\quad
\iota_C(\texttt{010000})=(0,2,0),\quad
\iota_C(\texttt{001000})=(0,0,2).
\]
若这样的线性映射 \(T\) 存在，则必有
\[
T(e_1)=2e_1,\qquad T(e_2)=2e_2,\qquad T(e_3)=2e_3,
\]
其中 \(e_1,e_2,e_3\) 为标准基。于是
\[
T(e_1-e_2)=2e_1-2e_2.
\]
然而同两张表对词 \(\texttt{101000}\) 的记录都给出
\[
\iota_B(\texttt{101000})=(1,-1,0),
\qquad
\iota_C(\texttt{101000})=(1,-1,0).
\]
故若 \(T(\iota_B(\texttt{101000}))=\iota_C(\texttt{101000})\)，则必须有
\[
T(e_1-e_2)=e_1-e_2,
\]
与前式矛盾。证毕。
\end{proof}

\paragraph{Weyl 不变量像、双框架矩与黄金分支证书振荡}
在 window-\(6\) 的 \(B_3/C_3\) 根字典、\(C_3\) 二次矩各向同性、\(B_3/C_3\) 线性统一障碍以及黄金分支的 Kronecker--Wasserstein 奇偶证书已经建立之后，还可进一步把 Weyl 不变量、方向矩、有限框架与 Fibonacci 归一化误差证书压缩为同一组闭合后果。

\begin{theorem}[window-\texorpdfstring{$6$}{6} 循环根字典上的 Weyl 不变量像恰由平方范数生成]\label{thm:xi-window6-b3c3-weyl-invariant-image-quadratic}
沿用定理 \ref{thm:xi-window6-b3c3-no-global-linear-unifier} 的记号，并记
\[
R_B:=\iota_B(X_6^{\mathrm{cyc}}),
\qquad
R_C:=\iota_C(X_6^{\mathrm{cyc}}).
\]
分别写
\[
R_{B,\mathrm{sh}}:=\{\pm e_1,\pm e_2,\pm e_3\},
\qquad
R_{B,\mathrm{lg}}:=\{\pm e_i\pm e_j:\ 1\le i<j\le 3\},
\]
\[
R_{C,\mathrm{sh}}:=\{\pm e_i\pm e_j:\ 1\le i<j\le 3\},
\qquad
R_{C,\mathrm{lg}}:=\{\pm 2e_1,\pm 2e_2,\pm 2e_3\}.
\]
令
\[
W_B:=W(B_3)\cong (\ZZ/2\ZZ)^3\rtimes S_3,
\qquad
W_C:=W(C_3)\cong (\ZZ/2\ZZ)^3\rtimes S_3.
\]
则：
\begin{enumerate}
  \item \(W_B\) 在 \(R_B\) 上恰有两条轨道 \(R_{B,\mathrm{sh}}\) 与 \(R_{B,\mathrm{lg}}\)，其大小分别为 \(6\) 与 \(12\)；\(W_C\) 在 \(R_C\) 上同样恰有两条轨道 \(R_{C,\mathrm{sh}}\) 与 \(R_{C,\mathrm{lg}}\)，其大小亦分别为 \(12\) 与 \(6\)；
  \item 限制映射
  \[
  \operatorname{res}_B:\mathbb R[x_1,x_2,x_3]^{W_B}\to \mathbb R^{R_B},
  \qquad
  \operatorname{res}_C:\mathbb R[x_1,x_2,x_3]^{W_C}\to \mathbb R^{R_C}
  \]
  的像都恰为二维，并且都由常数函数 \(1\) 与平方范数
  \[
  Q(x):=|x|^2
  \]
  的限制生成；
  \item 在 \(B_3\) 规范下，短根与长根的指示函数满足
  \[
  \boxed{\;
  \mathbf 1_{R_{B,\mathrm{sh}}}(x)=2-|x|^2,
  \qquad
  \mathbf 1_{R_{B,\mathrm{lg}}}(x)=|x|^2-1
  \qquad(x\in R_B).\;}
  \]
  \item 在 \(C_3\) 规范下，相应地有
  \[
  \boxed{\;
  \mathbf 1_{R_{C,\mathrm{sh}}}(x)=\frac{4-|x|^2}{2},
  \qquad
  \mathbf 1_{R_{C,\mathrm{lg}}}(x)=\frac{|x|^2-2}{2}
  \qquad(x\in R_C).\;}
  \]
\end{enumerate}
\end{theorem}
\begin{proof}
由表 \ref{tab:fold6_b3c3_root_dictionary_b3}，\(R_B\) 正是
\[
\{\pm e_1,\pm e_2,\pm e_3\}\cup \{\pm e_i\pm e_j:\ 1\le i<j\le 3\},
\]
而由表 \ref{tab:fold6_b3c3_root_dictionary_c3}，\(R_C\) 正是
\[
\{\pm 2e_1,\pm 2e_2,\pm 2e_3\}\cup \{\pm e_i\pm e_j:\ 1\le i<j\le 3\}.
\]
群 \(W_B\) 与 \(W_C\) 都由坐标置换与独立符号翻转生成，因此分别在各自的轴向根集合与非轴向根集合上作用传递，于是两侧都恰有两条轨道，大小由集合基数直接读出。

任一 Weyl 不变多项式在有限根集上的限制都必须在每条 Weyl 轨道上取常值，因此 \(\operatorname{im}(\operatorname{res}_B)\) 与 \(\operatorname{im}(\operatorname{res}_C)\) 的维数都至多为 \(2\)。另一方面，在 \(R_B\) 上
\[
Q(x)=|x|^2=
\begin{cases}
1,& x\in R_{B,\mathrm{sh}},\\
2,& x\in R_{B,\mathrm{lg}},
\end{cases}
\]
而在 \(R_C\) 上
\[
Q(x)=|x|^2=
\begin{cases}
2,& x\in R_{C,\mathrm{sh}},\\
4,& x\in R_{C,\mathrm{lg}}.
\end{cases}
\]
故 \(1\) 与 \(Q\) 在两条轨道上分别取不同值，从而在线性上独立，并生成全部轨道常值函数空间。于是两侧像空间都恰为二维。

最后，只需对两条轨道上的取值解线性插值方程即可。在 \(R_B\) 上，
\[
a+b|x|^2=
\begin{cases}
1,& |x|^2=1,\\
0,& |x|^2=2
\end{cases}
\]
解得 \(a=2,b=-1\)，从而
\[
\mathbf 1_{R_{B,\mathrm{sh}}}(x)=2-|x|^2.
\]
互补轨道的指示函数即为
\[
\mathbf 1_{R_{B,\mathrm{lg}}}(x)=1-\mathbf 1_{R_{B,\mathrm{sh}}}(x)=|x|^2-1.
\]
在 \(R_C\) 上同理，对
\[
a+b|x|^2=
\begin{cases}
1,& |x|^2=2,\\
0,& |x|^2=4
\end{cases}
\]
求解得到 \(a=2,b=-\frac12\)，故
\[
\mathbf 1_{R_{C,\mathrm{sh}}}(x)=\frac{4-|x|^2}{2},
\qquad
\mathbf 1_{R_{C,\mathrm{lg}}}(x)=1-\mathbf 1_{R_{C,\mathrm{sh}}}(x)=\frac{|x|^2-2}{2}.
\]
证毕。
\end{proof}
\leanverified{paper\_xi\_window6\_b3c3\_weyl\_invariant\_image\_quadratic}

\begin{theorem}[window-\texorpdfstring{$6$}{6} 的 \texorpdfstring{$B_3/C_3$}{B3/C3} 字典形成双紧框架且不存在单一 Euclid 相似统一]\label{thm:xi-window6-b3c3-tight-frame-fourth-moment-nonsimilarity}
沿用定理 \ref{thm:xi-window6-b3c3-weyl-invariant-image-quadratic} 的记号。则二阶矩算子满足
\[
\boxed{\;
\sum_{r\in R_B}rr^{\mathsf T}=10I_3,
\qquad
\sum_{r\in R_C}rr^{\mathsf T}=16I_3.\;}
\]
因而 \(R_B\) 与 \(R_C\) 分别构成 \(\mathbb R^3\) 上框架常数为 \(10\) 与 \(16\) 的有限紧框架。

进一步，对任意 \(v=(v_1,v_2,v_3)\in\mathbb R^3\)，四阶方向矩满足
\[
\boxed{\;
\sum_{r\in R_B}\langle r,v\rangle^4
=
10\sum_{i=1}^3 v_i^4
+24\sum_{1\le i<j\le 3}v_i^2v_j^2,\;}
\]
\[
\boxed{\;
\sum_{r\in R_C}\langle r,v\rangle^4
=
40\sum_{i=1}^3 v_i^4
+24\sum_{1\le i<j\le 3}v_i^2v_j^2.\;}
\]
最后，不存在 \(s>0\) 与 \(O\in O(3)\) 使
\[
R_C=s\,O(R_B).
\]
换言之，\(B_3/C_3\) 两套字典之间的长短根互换并非单一 Euclid 相似变换，而必然是按 Weyl 轨道分段的重标定。
\leanverified{paper\_xi\_window6\_b3c3\_tight\_frame\_fourth\_moment\_nonsimilarity}
\end{theorem}
\begin{proof}
对 \(R_B\)，六个短根 \(\pm e_i\) 的贡献为
\[
\sum_{i=1}^3\bigl(e_ie_i^{\mathsf T}+(-e_i)(-e_i)^{\mathsf T}\bigr)=2I_3.
\]
而对每一对指标 \(1\le i<j\le 3\)，四个长根 \(\pm e_i\pm e_j\) 的和满足
\[
\sum_{\varepsilon_i,\varepsilon_j\in\{\pm1\}}
(\varepsilon_i e_i+\varepsilon_j e_j)(\varepsilon_i e_i+\varepsilon_j e_j)^{\mathsf T}
=
4(e_ie_i^{\mathsf T}+e_je_j^{\mathsf T}),
\]
因为交叉项按符号对称抵消。对三对 \((i,j)\) 求和便得到总贡献 \(8I_3\)，从而
\[
\sum_{r\in R_B}rr^{\mathsf T}=2I_3+8I_3=10I_3.
\]

对 \(R_C\)，六个长根 \(\pm2e_i\) 的贡献为
\[
\sum_{i=1}^3\bigl((2e_i)(2e_i)^{\mathsf T}+(-2e_i)(-2e_i)^{\mathsf T}\bigr)=8I_3,
\]
而十二个短根 \(\pm e_i\pm e_j\) 的贡献与上面完全相同，亦为 \(8I_3\)。故
\[
\sum_{r\in R_C}rr^{\mathsf T}=16I_3.
\]
于是对任意 \(v\in\mathbb R^3\)，都有
\[
\sum_{r\in R_B}\langle r,v\rangle^2=10\|v\|_2^2,
\qquad
\sum_{r\in R_C}\langle r,v\rangle^2=16\|v\|_2^2,
\]
即二者分别为框架常数 \(10\) 与 \(16\) 的有限紧框架。

再看四阶矩。对 \(R_B\) 的六个短根，
\[
\sum_{r\in R_{B,\mathrm{sh}}}\langle r,v\rangle^4=2\sum_{i=1}^3 v_i^4.
\]
而对任一固定对 \((i,j)\)，有恒等式
\[
\sum_{\varepsilon_i,\varepsilon_j\in\{\pm1\}}
(\varepsilon_i v_i+\varepsilon_j v_j)^4
=
4(v_i^4+6v_i^2v_j^2+v_j^4).
\]
对三对 \((i,j)\) 求和便得到长根部分贡献
\[
8\sum_{i=1}^3 v_i^4+24\sum_{1\le i<j\le 3}v_i^2v_j^2.
\]
两者相加即得
\[
\sum_{r\in R_B}\langle r,v\rangle^4
=
10\sum_{i=1}^3 v_i^4+24\sum_{1\le i<j\le 3}v_i^2v_j^2.
\]
对 \(R_C\)，长根 \(\pm 2e_i\) 部分给出
\[
\sum_{r\in R_{C,\mathrm{lg}}}\langle r,v\rangle^4=32\sum_{i=1}^3 v_i^4,
\]
而十二个短根仍给出
\[
8\sum_{i=1}^3 v_i^4+24\sum_{1\le i<j\le 3}v_i^2v_j^2,
\]
故
\[
\sum_{r\in R_C}\langle r,v\rangle^4
=
40\sum_{i=1}^3 v_i^4+24\sum_{1\le i<j\le 3}v_i^2v_j^2.
\]

最后，若存在 \(s>0\) 与 \(O\in O(3)\) 使 \(R_C=s\,O(R_B)\)，则由二阶矩公式必有
\[
16I_3
=
\sum_{r\in R_C}rr^{\mathsf T}
=
s^2 O\left(\sum_{r\in R_B}rr^{\mathsf T}\right)O^{\mathsf T}
=
10s^2I_3,
\]
故
\[
s^2=\frac85,
\qquad
s^4=\frac{64}{25}.
\]
对任意单位向量 \(v\)，记 \(x_i:=v_i^2\)，则 \(\sum_i x_i=1\)。上面的四阶矩公式给出
\[
\sum_{r\in R_B}\langle r,v\rangle^4
=
10\sum_i x_i^2+24\sum_{i<j}x_ix_j
=
12-2\sum_i x_i^2,
\]
故其取值范围为
\[
\left[10,\frac{34}{3}\right],
\]
因为 \(\sum_i x_i^2\in[1/3,1]\)。若乘上 \(s^4=64/25\)，则可能范围变为
\[
\left[\frac{128}{5},\frac{2176}{75}\right].
\]
而对 \(R_C\) 同理，
\[
\sum_{r\in R_C}\langle r,v\rangle^4
=
40\sum_i x_i^2+24\sum_{i<j}x_ix_j
=
12+28\sum_i x_i^2,
\]
其取值范围为
\[
\left[\frac{64}{3},40\right].
\]
两者范围不同，因此不可能满足
\[
\sum_{r\in R_C}\langle r,v\rangle^4
=
s^4\sum_{r\in R_B}\langle r,O^{\mathsf T}v\rangle^4
\qquad(\forall\,\|v\|_2=1).
\]
故这样的 Euclid 相似变换不存在。证毕。
\end{proof}

\leanverified{paper\_xi\_golden\_fibonacci\_star\_w1\_normalized\_parity}
\begin{corollary}[黄金分支在 Fibonacci 分母上的星差与传输证书具有严格奇偶分裂]\label{cor:xi-golden-fibonacci-star-w1-normalized-parity}
沿用定理 \ref{thm:xi-golden-w1-star-certificate-parity-splitting} 的记号。记
\[
D_n^\ast:=D_{F_n}^\ast(P_{F_n}(\varphi^{-1})),
\qquad
\mathsf{T}_n:=\mathsf{T}_{F_n}(\varphi^{-1}).
\]
则 Fibonacci 归一化的星差满足
\[
\boxed{\;
F_nD_n^\ast=
\begin{cases}
1+\dfrac{1}{\sqrt5}+o(1),& n\ \mathrm{even},\\[2mm]
1,& n\ \mathrm{odd},
\end{cases}}
\]
其中奇数支是精确恒等式。

进一步，
\[
\boxed{\;
\frac{\mathsf{T}_n}{D_n^\ast}\to
\begin{cases}
\dfrac12,& n\ \mathrm{even},\\[2mm]
\dfrac12-\dfrac{1}{2\sqrt5}+\dfrac1{15},& n\ \mathrm{odd}.
\end{cases}}
\]
特别地，好侧奇数子序列上保留下来的额外箱内几何常数正是
\[
\frac12-\frac{1}{2\sqrt5}+\frac1{15}.
\]
\end{corollary}
\begin{proof}
由定理 \ref{thm:xi-golden-w1-star-certificate-parity-splitting}，当 \(n\) 为偶数时有精确恒等式
\[
\mathsf{T}_n=\frac12 D_n^\ast.
\]
再由定理 \ref{thm:xi-golden-w1-true-two-phase-limit} 的偶数支，
\[
F_n\mathsf{T}_n\to \frac12+\frac{1}{2\sqrt5},
\]
故
\[
F_nD_n^\ast=2F_n\mathsf{T}_n\to 1+\frac1{\sqrt5}.
\]

当 \(n\) 为奇数时，由结论 \ref{con:xi-terminal-kronecker-w1-denominator-one-sided-sensitivity}，黄金分支处于好侧，因而
\[
D_n^\ast=\frac1{F_n},
\]
从而 \(F_nD_n^\ast=1\) 为精确恒等式。再由定理
\ref{thm:xi-golden-w1-star-certificate-parity-splitting} 的奇数支，
\[
\frac{\mathsf{T}_n}{D_n^\ast}
=
F_n\mathsf{T}_n
\to
\frac12-\frac{1}{2\sqrt5}+\frac1{15}.
\]
证毕。
\end{proof}

\begin{corollary}[黄金分支 Fibonacci 归一化证书的振幅与 Ces\`aro 常数]\label{cor:xi-golden-fibonacci-star-w1-oscillation-cesaro}
\leanverified{paper\_xi\_golden\_fibonacci\_star\_w1\_oscillation\_cesaro}
沿用推论 \ref{cor:xi-golden-fibonacci-star-w1-normalized-parity} 的记号。则 Fibonacci 归一化传输差异与星差的振幅分别为
\[
\boxed{\;
\limsup_{n\to\infty}F_n\mathsf{T}_n
-
\liminf_{n\to\infty}F_n\mathsf{T}_n
=
\frac1{\sqrt5}-\frac1{15},\;}
\]
\[
\boxed{\;
\limsup_{n\to\infty}F_nD_n^\ast
-
\liminf_{n\to\infty}F_nD_n^\ast
=
\frac1{\sqrt5}.\;}
\]
并且 Ces\`aro 平均满足
\[
\boxed{\;
\lim_{N\to\infty}\frac1N\sum_{n=1}^N F_n\mathsf{T}_n
=
\frac{8}{15},\;}
\]
\[
\boxed{\;
\lim_{N\to\infty}\frac1N\sum_{n=1}^N F_nD_n^\ast
=
1+\frac{1}{2\sqrt5}.\;}
\]
\end{corollary}
\begin{proof}
传输差异的振幅公式与 Ces\`aro 平均公式分别就是推论
\ref{cor:xi-golden-w1-two-phase-amplitude-center} 与
\ref{cor:xi-golden-w1-two-phase-arithmetic-mean-eight-fifteenths} 的内容。

对星差，由推论 \ref{cor:xi-golden-fibonacci-star-w1-normalized-parity}，
\[
\limsup_{n\to\infty}F_nD_n^\ast=1+\frac1{\sqrt5},
\qquad
\liminf_{n\to\infty}F_nD_n^\ast=1,
\]
相减即得振幅
\[
\frac1{\sqrt5}.
\]

同样，由偶数与奇数索引在自然数中各占密度 \(1/2\)，且对应两条有界子序列分别收敛到
\[
1+\frac1{\sqrt5},
\qquad
1,
\]
故其 Ces\`aro 平均收敛到二者的算术平均
\[
\frac12\left(1+\frac1{\sqrt5}\right)+\frac12\cdot 1
=
1+\frac{1}{2\sqrt5}.
\]
证毕。
\end{proof}

\begin{theorem}[window-\texorpdfstring{$6$}{6} 循环根字典上的 Weyl 平均概率测度完全由单一二次矩恢复]\label{thm:xi-window6-b3c3-weyl-average-measure-second-moment}
沿用定理 \ref{thm:xi-window6-b3c3-weyl-invariant-image-quadratic} 的记号。令
\[
\sigma_{B,\mathrm{sh}},\ \sigma_{B,\mathrm{lg}}
\]
分别为 \(R_{B,\mathrm{sh}}\)、\(R_{B,\mathrm{lg}}\) 上的均匀概率测度，令
\[
\sigma_{C,\mathrm{sh}},\ \sigma_{C,\mathrm{lg}}
\]
分别为 \(R_{C,\mathrm{sh}}\)、\(R_{C,\mathrm{lg}}\) 上的均匀概率测度。对任意支撑在 \(R_B\) 上的概率测度 \(\nu_B\)，定义其 Weyl 平均为
\[
\overline{\nu}_B:=\frac{1}{|W_B|}\sum_{g\in W_B}g_\ast\nu_B,
\]
并记
\[
m_2^{(B)}(\nu_B):=\int_{R_B}|x|^2\,d\nu_B(x).
\]
则有显式恢复公式
\[
\boxed{\;
\overline{\nu}_B
=
\bigl(2-m_2^{(B)}(\nu_B)\bigr)\sigma_{B,\mathrm{sh}}
+
\bigl(m_2^{(B)}(\nu_B)-1\bigr)\sigma_{B,\mathrm{lg}}.\;}
\]

同样，对任意支撑在 \(R_C\) 上的概率测度 \(\nu_C\)，定义
\[
\overline{\nu}_C:=\frac{1}{|W_C|}\sum_{g\in W_C}g_\ast\nu_C,
\qquad
m_2^{(C)}(\nu_C):=\int_{R_C}|x|^2\,d\nu_C(x),
\]
则有
\[
\boxed{\;
\overline{\nu}_C
=
\frac{4-m_2^{(C)}(\nu_C)}{2}\,\sigma_{C,\mathrm{sh}}
+
\frac{m_2^{(C)}(\nu_C)-2}{2}\,\sigma_{C,\mathrm{lg}}.\;}
\]
因此，\(R_B\) 与 \(R_C\) 上全部 Weyl 不变概率测度构成的单纯形都退化为一条闭线段，其唯一自由参数就是二次矩。
\leanverified{paper\_xi\_window6\_b3c3\_weyl\_average\_measure\_second\_moment}
\end{theorem}
\begin{proof}
对 \(R_B\)，任意 Weyl 平均测度 \(\overline{\nu}_B\) 都在两条 Weyl 轨道上取常值，故必可唯一写成
\[
\overline{\nu}_B=a\,\sigma_{B,\mathrm{sh}}+(1-a)\,\sigma_{B,\mathrm{lg}}
\qquad (0\le a\le 1).
\]
由于 \(|x|^2\) 是 \(W_B\)-不变函数，故
\[
m_2^{(B)}(\nu_B)=m_2^{(B)}(\overline{\nu}_B).
\]
再利用 \(R_{B,\mathrm{sh}}\) 上 \(|x|^2=1\)、\(R_{B,\mathrm{lg}}\) 上 \(|x|^2=2\)，得到
\[
m_2^{(B)}(\nu_B)
=
a\cdot 1+(1-a)\cdot 2
=
2-a.
\]
解得
\[
a=2-m_2^{(B)}(\nu_B),
\qquad
1-a=m_2^{(B)}(\nu_B)-1,
\]
从而得到第一条公式。

对 \(R_C\) 同理，任意 Weyl 平均都唯一写成
\[
\overline{\nu}_C=b\,\sigma_{C,\mathrm{sh}}+(1-b)\,\sigma_{C,\mathrm{lg}},
\]
并由 \(R_{C,\mathrm{sh}}\) 上 \(|x|^2=2\)、\(R_{C,\mathrm{lg}}\) 上 \(|x|^2=4\) 得
\[
m_2^{(C)}(\nu_C)=m_2^{(C)}(\overline{\nu}_C)=2b+4(1-b)=4-2b.
\]
故
\[
b=\frac{4-m_2^{(C)}(\nu_C)}{2},
\qquad
1-b=\frac{m_2^{(C)}(\nu_C)-2}{2},
\]
即得第二条公式。

最后，Weyl 不变概率测度恰是两轨道均匀测度的凸组合，因此其参数空间正是一条闭线段，而上式表明二次矩已经给出这一线段上的唯一坐标。证毕。
\end{proof}

\noindent
由此，window-\(6\) 的 \(B_3/C_3\) 根字典不再只是一个离散配对表，而在函数层、矩层与概率层都表现出同一条低维刚性：Weyl 不变量的限制像由单一二次函数生成，二组根字典分别形成显式可审计的紧框架而彼此又不能由单一 Euclid 相似统一，黄金分支的 Fibonacci 证书则在星差与传输之间呈现严格奇偶振荡，而对称概率信息最终又坍缩到单一二次矩这一唯一自由度。


\paragraph{window-\texorpdfstring{$6$}{6} 轨道级对偶、坐标四重态分层与黄金分支的最优 Lipschitz 常数}
在定理 \ref{thm:xi-window6-cyclic-kernel-bc3-root-datum} 的 \(BC_3\) 规范根数据、定理 \ref{thm:xi-window6-fiber-multiplicity-weyl-breaking-b3} 的 Weyl 破缺、定理 \ref{thm:xi-golden-w1-true-two-phase-limit} 与推论 \ref{cor:xi-golden-fibonacci-star-w1-oscillation-cesaro} 的黄金分支 Fibonacci 证书已经固定之后，还可进一步把轨道级对偶、坐标异常与 Lipschitz 推前极限压缩为下列一组结果。

\begin{theorem}[window-\texorpdfstring{$6$}{6} 的 \texorpdfstring{$X_6^{\mathrm{cyc}}$}{X6cyc} 实现 \texorpdfstring{$B_3^\vee=C_3$}{B3v=C3} 的轨道级 Langlands 对偶]\label{thm:xi-window6-b3c3-orbitwise-langlands-duality}
沿用定理 \ref{thm:xi-window6-cyclic-kernel-bc3-root-datum} 的记号。记
\[
R_B:=\iota_B(X_6^{\mathrm{cyc}}),
\qquad
R_C:=\iota_C(X_6^{\mathrm{cyc}}),
\]
\[
X_{6,\mathrm{ax}}
:=
\{w\in X_6^{\mathrm{cyc}}:\ \mathrm{wt}(w)=1\},
\qquad
X_{6,\mathrm{off}}
:=
X_6^{\mathrm{cyc}}\setminus X_{6,\mathrm{ax}}.
\]
则有
\[
\iota_B(X_{6,\mathrm{ax}})
=
\{\pm e_1,\pm e_2,\pm e_3\},
\qquad
\iota_B(X_{6,\mathrm{off}})
=
\{\pm e_i\pm e_j:\ 1\le i<j\le 3\},
\]
\[
\iota_C(X_{6,\mathrm{ax}})
=
\{\pm 2e_1,\pm 2e_2,\pm 2e_3\},
\qquad
\iota_C(X_{6,\mathrm{off}})
=
\{\pm e_i\pm e_j:\ 1\le i<j\le 3\}.
\]
因此映射
\[
\mathcal D:R_B\to R_C,
\qquad
\mathcal D(\pm e_i)=\pm 2e_i,
\qquad
\mathcal D(\pm e_i\pm e_j)=\pm e_i\pm e_j
\]
满足
\[
\mathcal D\circ \iota_B=\iota_C
\qquad\text{于 }X_6^{\mathrm{cyc}}\text{ 上成立},
\]
亦即
\[
\mathcal D=\iota_C\circ \iota_B^{-1}.
\]
此外，\(\mathcal D\) 对公共 Weyl 群
\[
W:=W(B_3)\cong W(C_3)\cong (\ZZ/2\ZZ)^3\rtimes S_3
\]
是等变的；它把 \(B_3\) 的短根轨道与 \(C_3\) 的长根轨道识别，并把 \(B_3\) 的长根轨道与 \(C_3\) 的短根轨道识别。换言之，\(X_6^{\mathrm{cyc}}\) 精确实现了 \(B_3^\vee=C_3\) 的轨道级 Langlands 对偶。
\end{theorem}
\begin{proof}
定理 \ref{thm:xi-window6-cyclic-kernel-bc3-root-datum} 已经给出两张字典上的 \(6+12\) 轨道分裂，故上述四个集合恒等式立即成立。

由此可知 \(\mathcal D=\iota_C\circ \iota_B^{-1}\) 在 \(R_B\) 上被唯一确定；其显式公式正是
\[
\mathcal D(\pm e_i)=\pm 2e_i,
\qquad
\mathcal D(\pm e_i\pm e_j)=\pm e_i\pm e_j.
\]
再取任意 \(g\in W\)。由于 \(g\) 由坐标置换与独立符号翻转生成，故对轴根有
\[
\mathcal D(g(\pm e_i))
=
g(\pm 2e_i)
=
g\bigl(\mathcal D(\pm e_i)\bigr),
\]
对非轴根有
\[
\mathcal D(g(\pm e_i\pm e_j))
=
g(\pm e_i\pm e_j)
=
g\bigl(\mathcal D(\pm e_i\pm e_j)\bigr).
\]
于是 \(\mathcal D\) 与 \(W\)-作用可交换，即为 Weyl 等变双射。最后，由 \(R_B\) 中的六个短根被送到 \(R_C\) 中的六个长根，而十二个长根被送到 \(R_C\) 中的十二个短根，便得到所述轨道级对偶。证毕。
\end{proof}

\begin{corollary}[window-\texorpdfstring{$6$}{6} 的 Hamming 壳精确实现 \texorpdfstring{$B_3/C_3$}{B3/C3} 的长短根交换]\label{cor:xi-window6-hamming-shell-b3c3-length-swap}
沿用定理 \ref{thm:xi-window6-b3c3-orbitwise-langlands-duality} 的记号。则
\[
X_{6,\mathrm{ax}}
=
\{w\in X_6^{\mathrm{cyc}}:\ \mathrm{wt}(w)=1\}
\]
恰为 \(X_6^{\mathrm{cyc}}\) 的六点 Hamming 壳，并满足：
\begin{enumerate}
  \item 在 \(B_3\) 字典 \(\iota_B\) 下，\(X_{6,\mathrm{ax}}\) 恰对应六条短根
  \[
  \{\pm e_1,\pm e_2,\pm e_3\},
  \]
  而其补集 \(X_6^{\mathrm{cyc}}\setminus X_{6,\mathrm{ax}}\) 恰对应十二条长根
  \[
  \{\pm e_i\pm e_j:\ 1\le i<j\le 3\}.
  \]
  \item 在 \(C_3\) 字典 \(\iota_C\) 下，同一 Hamming 壳 \(X_{6,\mathrm{ax}}\) 恰对应六条长根
  \[
  \{\pm 2e_1,\pm 2e_2,\pm 2e_3\},
  \]
  而其补集恰对应十二条短根
  \[
  \{\pm e_i\pm e_j:\ 1\le i<j\le 3\}.
  \]
\end{enumerate}
因此，window-\(6\) 的 weight-\(1\) Hamming 壳在 \(B_3\) 与 \(C_3\) 两套规范之间精确充当长短根交换器。
\end{corollary}
\begin{proof}
由定理 \ref{thm:xi-window6-cyclic-kernel-bc3-root-datum}，\(|X_{6,\mathrm{ax}}|=6\) 且其补集大小为 \(12\)；同一定理同时给出
\[
\iota_B(X_{6,\mathrm{ax}})
=
\{\pm e_1,\pm e_2,\pm e_3\},
\qquad
\iota_C(X_{6,\mathrm{ax}})
=
\{\pm 2e_1,\pm 2e_2,\pm 2e_3\}.
\]
而定理 \ref{thm:xi-window6-b3c3-orbitwise-langlands-duality} 已表明两套字典下的补集分别给出 \(B_3\) 的十二条长根与 \(C_3\) 的十二条短根。合并即得结论。证毕。
\end{proof}

\begin{theorem}[window-\texorpdfstring{$6$}{6} 的 \texorpdfstring{$d_6^{\mathrm{bin}}$}{d6bin} 在 \texorpdfstring{$B_3$}{B3} 字典下呈现精确的坐标四重态分层]\label{thm:xi-window6-dbin-coordinate-quadruple-stratification}
沿用定理 \ref{thm:xi-window6-parity-charge-trigrading-b3} 的 \(B_3\) 根字典，定义
\[
\widetilde d_6^{\mathrm{bin}}
:=
d_6^{\mathrm{bin}}\circ \iota_B^{-1}:R_B\to\{2,3,4\}.
\]
再记
\[
L_x:=\{(0,\pm 1,\pm 1)\},
\qquad
L_y:=\{(\pm 1,0,\pm 1)\},
\qquad
L_z:=\{(\pm 1,\pm 1,0)\}.
\]
则有精确分层公式
\[
\widetilde d_6^{\mathrm{bin}}(r)=
\begin{cases}
2,& r\in L_x\cup \{-e_1\},\\[2mm]
3,& r\in L_y,\\[2mm]
4,& r\in L_z\cup \{e_1,\pm e_2,\pm e_3\}.
\end{cases}
\]
特别地：
\begin{enumerate}
  \item 在 \(12\) 条长根上，\(\widetilde d_6^{\mathrm{bin}}\) 不是 Weyl 不变量，但恰按
  \[
  L_x\mapsto 2,
  \qquad
  L_y\mapsto 3,
  \qquad
  L_z\mapsto 4
  \]
  分裂为三个坐标四重态；
  \item 在 \(6\) 条短根上，只有 \(-e_1\) 取到最小值 \(2\)，而
  \[
  e_1,\ \pm e_2,\ \pm e_3
  \]
  全部取值 \(4\)；
  \item 在循环扇区内部已有精确直方图
  \[
  \#\{w\in X_6^{\mathrm{cyc}}:\ d_6^{\mathrm{bin}}(w)=2\}=5,\qquad
  \#\{w\in X_6^{\mathrm{cyc}}:\ d_6^{\mathrm{bin}}(w)=3\}=4,\qquad
  \#\{w\in X_6^{\mathrm{cyc}}:\ d_6^{\mathrm{bin}}(w)=4\}=9,
  \]
  再并入三条 boundary 状态后，恢复全文 \(X_6\) 的直方图
  \[
  2:8,\qquad 3:4,\qquad 4:9.
  \]
\end{enumerate}
\end{theorem}
\begin{proof}
由表 \ref{tab:fold6_b3c3_root_dictionary_b3} 可直接读出
\[
\iota_B(\texttt{000001})=-e_1,
\qquad
\iota_B(\texttt{010001}),\iota_B(\texttt{001001}),
\iota_B(\texttt{000101}),\iota_B(\texttt{010101})\in L_x.
\]
而表 \ref{tab:fold_bin_gauge_m6_fibers} 给出这五个 cyclic 词的纤维大小全部等于 \(2\)。故
\[
\widetilde d_6^{\mathrm{bin}}=2
\qquad\text{于 }L_x\cup\{-e_1\}\text{ 上成立}.
\]

同样，由表 \ref{tab:fold6_b3c3_root_dictionary_b3}，
\[
\iota_B(\texttt{100010}),\ \iota_B(\texttt{010010}),\ \iota_B(\texttt{001010}),\ \iota_B(\texttt{101010})
\]
恰为 \(L_y\) 的四个点；而由表 \ref{tab:fold_bin_gauge_m6_fibers}，这四个词的纤维大小全部等于 \(3\)。故
\[
\widetilde d_6^{\mathrm{bin}}=3
\qquad\text{于 }L_y\text{ 上成立}.
\]

剩余的九个 cyclic 词
\[
\texttt{000000},\ \texttt{100000},\ \texttt{010000},\ \texttt{001000},\ \texttt{000100},
\ \texttt{000010},\ \texttt{101000},\ \texttt{100100},\ \texttt{010100}
\]
在表 \ref{tab:fold6_b3c3_root_dictionary_b3} 下恰对应
\[
L_z\cup \{e_1,\pm e_2,\pm e_3\},
\]
并且由表 \ref{tab:fold_bin_gauge_m6_fibers} 可见其纤维大小全部等于 \(4\)。这便得到整体分层公式。

三条附加断言只是对上述公式逐点计数。最后，boundary 集
\[
X_6^{\mathrm{bdry}}
=
\{\texttt{100001},\texttt{100101},\texttt{101001}\}
\]
的三条词在表 \ref{tab:fold_bin_gauge_m6_fibers} 中全部满足 \(d_6^{\mathrm{bin}}=2\)，故全局直方图由 cyclic sector 的 \(5,4,9\) 增补为 \(8,4,9\)。证毕。
\end{proof}

\begin{theorem}[黄金分支在 Fibonacci 分母上的 \texorpdfstring{$W_1$}{W1} 与星差证书满足严格奇偶双相律]\label{thm:xi-golden-fibonacci-w1-star-strict-parity-law}
沿用定理 \ref{thm:xi-golden-w1-true-two-phase-limit} 与推论 \ref{cor:xi-golden-fibonacci-star-w1-normalized-parity} 的记号。令
\[
\mathsf T_n:=\mathsf T_{F_n}(\varphi^{-1}),
\qquad
D_n^\ast:=D_{F_n}^\ast(P_{F_n}(\varphi^{-1})),
\]
\[
\delta_n:=\varphi^{-1}-\frac{F_{n-1}}{F_n}
=
\frac{(-1)^{n-1}}{\sqrt5\,F_n^2}\Bigl(1-(-1)^n\varphi^{-2n}\Bigr),
\]
并记
\[
A:=\frac12+\frac{1}{2\sqrt5},
\qquad
B:=\frac12-\frac{1}{2\sqrt5}+\frac1{15}.
\]
则：
\begin{enumerate}
  \item \(\delta_n<0\) 当且仅当 \(n\) 为偶数，\(\delta_n>0\) 当且仅当 \(n\) 为奇数；
  \item 归一化传输差异满足
  \[
  \limsup_{n\to\infty}F_n\mathsf T_n=A,
  \qquad
  \liminf_{n\to\infty}F_n\mathsf T_n=B;
  \]
  \item 归一化星差满足
  \[
  \limsup_{n\to\infty}F_nD_n^\ast=1+\frac1{\sqrt5},
  \qquad
  \liminf_{n\to\infty}F_nD_n^\ast=1;
  \]
  \item 两者的振幅分别为
  \[
  \limsup_{n\to\infty}F_n\mathsf T_n-\liminf_{n\to\infty}F_n\mathsf T_n
  =
  \frac1{\sqrt5}-\frac1{15},
  \]
  \[
  \limsup_{n\to\infty}F_nD_n^\ast-\liminf_{n\to\infty}F_nD_n^\ast
  =
  \frac1{\sqrt5};
  \]
  \item Ces\`aro 平均满足
  \[
  \lim_{N\to\infty}\frac1N\sum_{n=1}^N F_n\mathsf T_n
  =
  \frac{8}{15},
  \qquad
  \lim_{N\to\infty}\frac1N\sum_{n=1}^N F_nD_n^\ast
  =
  1+\frac{1}{2\sqrt5}.
  \]
\end{enumerate}
\end{theorem}
\begin{proof}
引理 \ref{lem:golden-convergent-error-closed} 已经给出 \(\delta_n\) 的闭式表达，因而其符号满足
\[
\operatorname{sign}(\delta_n)=(-1)^{n-1}.
\]
于是偶数 \(n\) 对应坏侧，奇数 \(n\) 对应好侧。

再由定理 \ref{thm:xi-golden-w1-true-two-phase-limit}，偶数与奇数两条子序列上的极限常数分别为 \(A\) 与 \(B\)，故第二条成立。由推论 \ref{cor:xi-golden-fibonacci-star-w1-normalized-parity}，偶数子序列上
\[
F_nD_n^\ast\to 1+\frac1{\sqrt5},
\]
而奇数子序列上恒有
\[
F_nD_n^\ast=1,
\]
从而得到第三条。

第四条与第五条正是推论 \ref{cor:xi-golden-fibonacci-star-w1-oscillation-cesaro} 的内容。证毕。
\end{proof}

\begin{theorem}[黄金分支 Fibonacci 分母刻度上的 \texorpdfstring{$W_1$}{W1} Lipschitz 推前常数具有最优奇偶上界]\label{thm:xi-golden-fibonacci-lipschitz-pushforward-sharp}
沿用定理 \ref{thm:xi-golden-w1-true-two-phase-limit} 与定理 \ref{thm:xi-golden-fibonacci-w1-star-strict-parity-law} 的记号。对任意 Lipschitz 映射
\[
\Phi:[0,1)\to[0,1)
\]
定义
\[
W_n(\Phi):=
W_1\bigl(\Phi_\#\mu_{F_n},\Phi_\#\lambda\bigr).
\]
则对任意 \(L\)-Lipschitz 的 \(\Phi\) 有
\[
\limsup_{\substack{n\to\infty\\ n\ \mathrm{even}}}
F_nW_n(\Phi)
\le
L\left(\frac12+\frac{1}{2\sqrt5}\right),
\]
\[
\limsup_{\substack{n\to\infty\\ n\ \mathrm{odd}}}
F_nW_n(\Phi)
\le
L\left(\frac12-\frac{1}{2\sqrt5}+\frac1{15}\right).
\]
而且这两个系数在统一不等式
\[
\limsup_{\substack{n\to\infty\\ n\ \mathrm{even}}}
F_nW_n(\Phi)\le \mathrm{Lip}(\Phi)\,A,
\qquad
\limsup_{\substack{n\to\infty\\ n\ \mathrm{odd}}}
F_nW_n(\Phi)\le \mathrm{Lip}(\Phi)\,B
\]
中都是最优的：当 \(\Phi=\mathrm{id}_{[0,1)}\) 时，两条上界同时取等。

特别地，由于
\[
B=\frac12-\frac{1}{2\sqrt5}+\frac1{15}<1,
\]
故对任意 \(1\)-Lipschitz 读出都有
\[
\limsup_{\substack{n\to\infty\\ n\ \mathrm{odd}}}
F_nW_n(\Phi)
\le
\frac12-\frac{1}{2\sqrt5}+\frac1{15}
<
1.
\]
而同一奇数子序列上的星差基线恒满足
\[
F_nD_n^\ast=1.
\]
\end{theorem}
\begin{proof}
由命题 \ref{prop:w1-lipschitz-pushforward}，
\[
W_n(\Phi)\le \mathrm{Lip}(\Phi)\,\mathsf T_n.
\]
若 \(\Phi\) 为 \(L\)-Lipschitz，则
\[
W_n(\Phi)\le L\,\mathsf T_n.
\]
两边乘以 \(F_n\)，并分别在偶数与奇数子序列上取 \(\limsup\)，再代入定理 \ref{thm:xi-golden-fibonacci-w1-star-strict-parity-law} 的第二条，即得两条奇偶上界。

若取 \(\Phi=\mathrm{id}_{[0,1)}\)，则
\[
W_n(\Phi)=W_1(\mu_{F_n},\lambda)=\mathsf T_n,
\qquad
\mathrm{Lip}(\Phi)=1.
\]
因此偶数与奇数子序列上的上界常数分别恰达到 \(A\) 与 \(B\)，故这两个系数不能再减小，亦即为最优常数。

最后，由定理 \ref{thm:xi-golden-fibonacci-w1-star-strict-parity-law}，奇数 \(n\) 上恒有
\[
F_nD_n^\ast=1.
\]
再结合 \(B<1\)，便得到最后一条严格不等式。证毕。
\end{proof}

\noindent
由此，window-\(6\) 的 \(B_3/C_3\) 字典被进一步提升为一条精确的轨道级对偶与坐标四重态分层，而黄金分支的 Fibonacci 证书则进一步推出全部 Lipschitz 推前上的最优二相常数。于是根系字典、纤维异常与传输读出三者在同一有限审计链上获得了统一闭合。


\begin{corollary}[六倍窗大零块与非消失碰撞缺口的共存]\label{cor:xi-window6-zero-block-and-collision-gap-coexistence}
\leanverified{paper\_xi\_window6\_zero\_block\_and\_collision\_gap\_coexistence}
沿用定理 \ref{thm:xi-fold-fibonacci-collision-gap-positive-floor} 的记号。若
\[
m+2\equiv 0\pmod 6,
\]
则由推论 \ref{cor:fold-zero-half-index-multiple6} 有
\[
\boxed{\;
|\mathcal Z_m|\ge F_{(m+2)/2},
\qquad
S_{F_{(m+2)/2}}(F_{m+2})\subseteq \mathcal Z_m.\;}
\]
同时
\[
\boxed{\;
\liminf_{\substack{m\to\infty\\ m+2\equiv 0\ (\mathrm{mod}\ 6)}}
F_{m+2}\,g_{\mathrm{col}}^2(m)
\ge
2C_\varphi^2
\approx 0.1007267225141179>0.\;}
\]
因此在最有利于产生大零点块的六倍窗上，也仍存在严格正的 Fibonacci 级碰撞缺口；零点增殖本身并不足以迫使折叠趋于均匀。
\end{corollary}
\begin{proof}
第一条是推论 \ref{cor:fold-zero-half-index-multiple6} 的直接内容。第二条只需把定理 \ref{thm:xi-fold-fibonacci-collision-gap-positive-floor} 的下极限限制到子序列 \(m+2\equiv 0\pmod 6\)；任一子序列的下极限都不小于全序列下极限，故结论成立。证毕。
\end{proof}

\begin{theorem}[六倍窗完整余类零块与非平凡谱峰抑制]\label{thm:xi-window6-full-coset-zero-block-peak-suppression}
沿用推论 \ref{cor:fold-zero-half-index-multiple6}、注 \ref{rem:fold-zero-lb-m58} 与定义 \ref{def:fold-normalized-amplitude} 的记号。若
\[
m+2\equiv 0\pmod 6,
\]
并记
\[
M:=F_{m+2},
\qquad
d:=\frac{m+2}{2},
\qquad
\mathcal Z_m:=\{k\in \ZZ/(M\ZZ):\ \widehat d_m(k)=0\},
\]
\[
\Sigma_m:=\{k\in \ZZ/(M\ZZ):\ \widehat d_m(k)\neq 0\},
\qquad
S_m:=\max_{k\ne 0}|\widehat d_m(k)|.
\]
则有完整余类消失
\[
S_{F_d}(M)\subseteq \mathcal Z_m,
\qquad
\Sigma_m\subseteq \bigl(\ZZ/(M\ZZ)\bigr)\setminus S_{F_d}(M),
\]
从而
\[
\boxed{\;
|\Sigma_m|
\le
M-F_d
=
F_{m+2}-F_{(m+2)/2}.\;}
\]
并且非平凡 Fourier 峰值满足
\[
\boxed{\;
S_m
\le
2^m\sqrt{M-1-|\mathcal Z_m|}
\le
2^m\sqrt{M-1-F_d}
=
2^m\sqrt{F_{m+2}-1-F_{(m+2)/2}}.\;}
\]
特别地，当 \(m=58\) 时，
\[
\boxed{\;
|\Sigma_{58}|
\le
F_{60}-839415,
\qquad
S_{58}
\le
2^{58}\sqrt{F_{60}-1-839415}.\;}
\]
\end{theorem}
\begin{proof}
由推论 \ref{cor:fold-zero-half-index-multiple6}，
\[
S_{F_d}(M)\subseteq \mathcal Z_m,
\qquad
|S_{F_d}(M)|=F_d.
\]
因此全部 Fourier 支撑都落在其补集内，从而
\[
\Sigma_m\subseteq \bigl(\ZZ/(M\ZZ)\bigr)\setminus S_{F_d}(M),
\]
并立即得到
\[
|\Sigma_m|\le M-F_d.
\]

再记
\[
a_m(k):=\frac{|\widehat d_m(k)|}{2^m}.
\]
由定义 \ref{def:fold-normalized-amplitude}，
\[
0\le a_m(k)\le 1,
\qquad
a_m(k)=0\iff k\in\mathcal Z_m.
\]
于是
\[
\frac{S_m}{2^m}=\max_{k\ne 0}a_m(k)
\le
\left(\sum_{k\ne 0}a_m(k)^2\right)^{1/2}
\le
\sqrt{\#\{k\ne 0:\ a_m(k)\ne 0\}}
=
\sqrt{M-1-|\mathcal Z_m|}.
\]
又由推论 \ref{cor:fold-zero-half-index-multiple6}，\(|\mathcal Z_m|\ge F_d\)，故
\[
S_m\le 2^m\sqrt{M-1-|\mathcal Z_m|}
\le
2^m\sqrt{M-1-F_d}.
\]
最后在 \(m=58\) 时，由注 \ref{rem:fold-zero-lb-m58} 的精确计数
\[
|\mathcal Z_{58}|=839415
\]
代回即可得到两条数值化上界。证毕。
\end{proof}

\paragraph{连续容量导数、有限 primitive 手术与 window-\texorpdfstring{$6$}{6} 异常分裂}
下列条目把连续容量曲线在低阈值处的右斜率、有限个 Euler 原子的 primitive 手术可加性、真实输入 \(40\) 状态核零温 Abel 常数的核心分离，以及 window-\(6\) 的中心/非中心异常直和分裂压缩为一组可独立调用的终端结论。

\begin{theorem}[连续容量曲线前两处右导数完全刻画 fold-gauge 的阿贝尔化与中心]\label{thm:xi-foldbin-continuous-capacity-right-derivative-abelian-center}
\leanverified{paper\_xi\_foldbin\_continuous\_capacity\_right\_derivative\_abelian\_center}
沿用命题 \ref{prop:op-algebra-capacity-kinks-equal-spectrum-levels} 与定理 \ref{thm:xi-foldbin-gauge-representation-distribution-law-m6-spectrum} 的记号。对固定 \(m\)，令
\[
d_m^{\mathrm{bin}}(w):=\bigl|\Fold_m^{\mathrm{bin}}{}^{-1}(w)\bigr|,
\qquad
N_m(t):=\#\{w\in X_m:\ d_m^{\mathrm{bin}}(w)\ge t\},
\]
\[
\mathcal C_m^{\mathrm{cont}}(T):=\sum_{w\in X_m}\min\bigl(d_m^{\mathrm{bin}}(w),T\bigr),
\qquad
G_m^{\mathrm{bin}}:=\Aut(\Fold_m^{\mathrm{bin}}).
\]
并记
\[
R_m:=\dim_{\FF_2}\bigl(G_m^{\mathrm{bin}}\bigr)^{\mathrm{ab}},
\qquad
Z_m:=\dim_{\FF_2} Z\!\bigl(G_m^{\mathrm{bin}}\bigr).
\]
则 \(\mathcal C_m^{\mathrm{cont}}\) 为分段线性凹函数，且对每个整数 \(k\ge 0\) 都有
\[
\partial_T^{+}\mathcal C_m^{\mathrm{cont}}(k)=N_m(k^+)=N_m(k+1).
\]
特别地，
\[
\boxed{\;
R_m
=
N_m(2)
=
\partial_T^{+}\mathcal C_m^{\mathrm{cont}}(1).\;}
\]
并且
\[
\boxed{\;
Z_m
=
N_m(2)-N_m(3)
=
\partial_T^{+}\mathcal C_m^{\mathrm{cont}}(1)-\partial_T^{+}\mathcal C_m^{\mathrm{cont}}(2).\;}
\]
若经由自然交换化映射把 \(Z(G_m^{\mathrm{bin}})\) 视为 \(\bigl(G_m^{\mathrm{bin}}\bigr)^{\mathrm{ab}}\) 的子空间，则
\[
\boxed{\;
\dim_{\FF_2}
\left(
\bigl(G_m^{\mathrm{bin}}\bigr)^{\mathrm{ab}}/Z\!\bigl(G_m^{\mathrm{bin}}\bigr)
\right)
=
N_m(3)
=
\partial_T^{+}\mathcal C_m^{\mathrm{cont}}(2).\;}
\]
在 \(m=6\) 时，由表 \ref{tab:fold_bin_gauge_m6_fibers} 给出的纤维直方图
\[
d=2:8,\qquad d=3:4,\qquad d=4:9,
\]
可得
\[
R_6=21,
\qquad
Z_6=8,
\qquad
\dim_{\FF_2}
\left(
\bigl(G_6^{\mathrm{bin}}\bigr)^{\mathrm{ab}}/Z\!\bigl(G_6^{\mathrm{bin}}\bigr)
\right)
=13.
\]
\end{theorem}
\begin{proof}
命题 \ref{prop:op-algebra-capacity-kinks-equal-spectrum-levels} 已给出 \(\mathcal C_m^{\mathrm{cont}}\) 的分段线性凹性。对单项
\[
f_d(T):=\min(d,T)
\qquad(d\in \ZZ_{\ge 1}),
\]
其在整数点 \(k\) 的右导数满足
\[
\partial_T^{+}f_d(k)=
\begin{cases}
1,& d\ge k+1,\\
0,& d\le k.
\end{cases}
\]
对全部纤维求和便得到
\[
\partial_T^{+}\mathcal C_m^{\mathrm{cont}}(k)
=
\sum_{w\in X_m}\partial_T^{+}\min\bigl(d_m^{\mathrm{bin}}(w),T\bigr)\big|_{T=k}
=
\#\{w:\ d_m^{\mathrm{bin}}(w)\ge k+1\}
=
N_m(k+1).
\]
由于各 \(d_m^{\mathrm{bin}}(w)\) 为整数，上式亦即 \(N_m(k^+)\)。
\par
另一方面，定理 \ref{thm:xi-foldbin-gauge-representation-distribution-law-m6-spectrum} 给出
\[
\bigl(G_m^{\mathrm{bin}}\bigr)^{\mathrm{ab}}
\cong
\prod_{w:\,d_m^{\mathrm{bin}}(w)\ge 2}\ZZ/2\ZZ,
\]
故
\[
R_m=\#\{w:\ d_m^{\mathrm{bin}}(w)\ge 2\}=N_m(2)=\partial_T^{+}\mathcal C_m^{\mathrm{cont}}(1).
\]
又因为
\[
G_m^{\mathrm{bin}}
\cong
\prod_{w\in X_m}\mathfrak S_{d_m^{\mathrm{bin}}(w)},
\]
而有限对称群的中心满足
\[
Z(\mathfrak S_d)\cong
\begin{cases}
\ZZ/2\ZZ,& d=2,\\
\{e\},& d\neq 2,
\end{cases}
\]
所以
\[
Z\!\bigl(G_m^{\mathrm{bin}}\bigr)
\cong
\prod_{w:\,d_m^{\mathrm{bin}}(w)=2}\ZZ/2\ZZ,
\]
从而
\[
Z_m=\#\{w:\ d_m^{\mathrm{bin}}(w)=2\}=N_m(2)-N_m(3).
\]
在交换化中，\(d=2\) 纤维恰给出中心坐标，而 \(d\ge 3\) 纤维给出其余全部符号坐标，因此商空间的 \(\FF_2\)-维数正为 \(N_m(3)\)。最后，\(m=6\) 的数值直接由表 \ref{tab:fold_bin_gauge_m6_fibers} 的直方图读出。证毕。
\end{proof}

\begin{theorem}[有限个可剥离 Euler 原子的 primitive 手术严格可加]\label{thm:xi-finite-euler-primitive-surgery-strict-additivity}
\leanverified{paper\_xi\_finite\_euler\_primitive\_surgery\_strict\_additivity}
沿用 \S\ref{app:witt-weighted} 中加权 Witt 变换的记号。设某个加权有限矩阵族/有限型子移位的 \(\zeta\)-函数在形式幂级数环中分解为
\[
\zeta(z,u)=
\left(
\prod_{j=1}^{J}(1-u^{E_j}z^{\ell_j})^{-m_j}
\right)
\zeta_{\mathrm{core}}(z,u),
\qquad
m_j\in \ZZ_{\ge 1},\ \ell_j\ge 1,\ E_j\ge 0,
\]
并假设每个因子 \(1-u^{E_j}z^{\ell_j}\) 在主极点 \(z=z_\star(u)\) 处解析且不为零，因此主增长率 \(\lambda(u)=z_\star(u)^{-1}\) 完全由 \(\zeta_{\mathrm{core}}\) 的极点决定。令
\[
P(z,u):=\sum_{n\ge 1}p_n(u)z^n
=
\sum_{r\ge 1}\frac{\mu(r)}{r}\log \zeta(z^r,u^r)
\]
为 primitive 生成函数，并记 \(P_{\mathrm{core}}(z,u)\) 为 \(\zeta_{\mathrm{core}}(z,u)\) 的 primitive 生成函数。则有严格分解
\[
\boxed{\;
P(z,u)
=
P_{\mathrm{core}}(z,u)
+\sum_{j=1}^{J}m_j\,u^{E_j}z^{\ell_j}.\;}
\]
进一步，若在 \(z\to z_\star(u)\) 时有展开
\[
P(z,u)=-\log(1-\lambda(u)z)+\log\mathfrak M(u)+o(1),
\]
\[
P_{\mathrm{core}}(z,u)
=
-\log(1-\lambda(u)z)+\log\mathfrak M_{\mathrm{core}}(u)+o(1),
\]
则必有精确常数修正
\[
\boxed{\;
\log\mathfrak M(u)
=
\log\mathfrak M_{\mathrm{core}}(u)
+\sum_{j=1}^{J}m_j\,u^{E_j}z_\star(u)^{\ell_j}.\;}
\]
\end{theorem}
\begin{proof}
这是命题 \ref{prop:primitive-orbit-surgery} 的有限和版本。把分解式取对数得
\[
\log\zeta(z,u)
=
-\sum_{j=1}^{J}m_j\log(1-u^{E_j}z^{\ell_j})
+\log\zeta_{\mathrm{core}}(z,u).
\]
将其代入 Möbius--Witt 变换，得到
\[
\begin{aligned}
P(z,u)
&=
\sum_{r\ge 1}\frac{\mu(r)}{r}\log\zeta(z^r,u^r)\\
&=
P_{\mathrm{core}}(z,u)
-\sum_{j=1}^{J}m_j\sum_{r\ge 1}\frac{\mu(r)}{r}
\log\bigl(1-u^{rE_j}z^{r\ell_j}\bigr).
\end{aligned}
\]
对任意形式变量 \(X\) 的经典 Witt 恒等式
\[
-\sum_{r\ge 1}\frac{\mu(r)}{r}\log(1-X^r)=X
\]
逐项应用于 \(X=u^{E_j}z^{\ell_j}\)，便得到
\[
P(z,u)=P_{\mathrm{core}}(z,u)+\sum_{j=1}^{J}m_j\,u^{E_j}z^{\ell_j}.
\]
再令 \(z\to z_\star(u)\)，并比较两边展开中的常数项，即得
\[
\log\mathfrak M(u)-\log\mathfrak M_{\mathrm{core}}(u)
=
\sum_{j=1}^{J}m_j\,u^{E_j}z_\star(u)^{\ell_j}.
\]
证毕。
\end{proof}

\begin{theorem}[原子 Witt 因子的主极点热力学不变性与 Abel 有限部的刚性偏移]\label{thm:xi-real-input40-atomic-factor-thermodynamic-invariance}
\leanverified{paper\_xi\_real\_input40\_atomic\_factor\_thermodynamic\_invariance}
沿用附录 \ref{app:real-input-40-zeta-u}、定理
\ref{thm:xi-finite-euler-primitive-surgery-strict-additivity} 与命题
\ref{prop:real-input-40-output-cumulants-closed} 的记号。对真实输入 \(40\) 状态核，记
\[
\zeta(z,u)=\Delta(z,u)^{-1},
\qquad
\Delta(z,u)=(1-uz^2)\,F(z,u),
\qquad
\zeta_{\mathrm{core}}(z,u):=F(z,u)^{-1}.
\]
在输出位势的主 Perron 分支上，令
\[
z_\star(u):=\lambda(u)^{-1}
\]
为由
\[
F(z_\star(u),u)=0
\]
确定的主极点，并定义 primitive 生成函数
\[
\mathcal P(z,u):=\sum_{n\ge 1}p_n(u)z^n,
\qquad
\mathcal P_{\mathrm{core}}(z,u):=\sum_{n\ge 1}p_n^{\mathrm{core}}(u)z^n.
\]
则有精确剥离
\[
\boxed{\;
\mathcal P(z,u)=uz^2+\mathcal P_{\mathrm{core}}(z,u).\;}
\]
若进一步在 \(z\to z_\star(u)\) 时有 Abel 有限部展开
\[
\mathcal P_{\mathrm{core}}(z,u)
=
-\log\!\bigl(1-\lambda(u)z\bigr)
+\log\mathfrak M_{\mathrm{core}}(u)+o(1),
\]
则总体 primitive 生成函数满足
\[
\mathcal P(z,u)
=
-\log\!\bigl(1-\lambda(u)z\bigr)
+\log\mathfrak M(u)+o(1),
\]
且有限部常数发生精确刚性偏移
\[
\boxed{\;
\log\mathfrak M(u)
=
\log\mathfrak M_{\mathrm{core}}(u)+u\,z_\star(u)^2
=
\log\mathfrak M_{\mathrm{core}}(u)+\frac{u}{\lambda(u)^2}.\;}
\]
特别地，由
\[
P(\theta):=\log\lambda(e^\theta)
\]
定义的压力函数、其全部累积量 \(P^{(k)}(0)\) 以及与之对应的
Legendre--Fenchel 大偏差率函数，都完全由 core 方程 \(F(z,u)=0\) 决定，而不受原子因子 \((1-uz^2)\) 的影响。
\end{theorem}
\begin{proof}
定理 \ref{thm:xi-atomic-witt-cycle-primitive-exact-peeling} 已给出
\[
\mathcal P(z,u)=uz^2+\mathcal P_{\mathrm{core}}(z,u).
\]
由于 \(z_\star(u)\) 是由 \(F(z,u)=0\) 所定义的主极点，而 \(uz^2\) 在 \(z=z_\star(u)\) 处解析，故当
\(z\to z_\star(u)\) 时有
\[
uz^2=u\,z_\star(u)^2+o(1).
\]
把这一式与 \(\mathcal P_{\mathrm{core}}(z,u)\) 的 Abel 有限部展开相加，即得
\[
\mathcal P(z,u)
=
-\log\!\bigl(1-\lambda(u)z\bigr)
+\Bigl(\log\mathfrak M_{\mathrm{core}}(u)+u\,z_\star(u)^2\Bigr)+o(1),
\]
从而
\[
\log\mathfrak M(u)=\log\mathfrak M_{\mathrm{core}}(u)+u\,z_\star(u)^2
=
\log\mathfrak M_{\mathrm{core}}(u)+\frac{u}{\lambda(u)^2}.
\]
另一方面，压力函数 \(P(\theta)=\log\lambda(e^\theta)\) 只依赖主极点倒数 \(\lambda(u)\)，而 \(\lambda(u)\) 由 \(F(z,u)=0\) 的主根定义；因此一切仅由 \(P\) 导出的热力学量都与原子因子无关。证毕。
\end{proof}

\begin{corollary}[原子有限部偏移的微分塔与平衡点初值]\label{cor:xi-real-input40-atomic-abel-shift-differential-tower}
\leanverified{paper\_xi\_real\_input40\_atomic\_abel\_shift\_differential\_tower}
沿用定理 \ref{thm:xi-real-input40-atomic-factor-thermodynamic-invariance} 的记号，并定义
\[
\beta(\theta):=e^\theta z_\star(e^\theta)^2
=
e^\theta\lambda(e^\theta)^{-2}
=
e^\theta e^{-2P(\theta)}.
\]
则有精确恒等式
\[
\boxed{\;
\log\mathfrak M(e^\theta)-\log\mathfrak M_{\mathrm{core}}(e^\theta)=\beta(\theta).\;}
\]
并且 \(\beta\) 满足微分闭式
\[
\boxed{\;
\beta'(\theta)=\bigl(1-2P'(\theta)\bigr)\beta(\theta),\;}
\]
\[
\boxed{\;
\beta''(\theta)=\Bigl(\bigl(1-2P'(\theta)\bigr)^2-2P''(\theta)\Bigr)\beta(\theta).\;}
\]
在平衡点 \(\theta=0\) 处，由表
\ref{tab:real_input_40_output_potential_cumulants_closed} 中的闭式
\[
P'(0)=\frac{23}{20}-\frac{7\sqrt5}{20},
\qquad
P''(0)=\frac{1791}{200}-\frac{3983\sqrt5}{1000},
\qquad
\lambda(1)=\varphi^2
\]
可得
\[
\boxed{\;
\beta(0)=\varphi^{-4}=\frac{7-3\sqrt5}{2},\;}
\]
\[
\boxed{\;
\beta'(0)=\frac{-49+22\sqrt5}{5}>0,\;}
\]
\[
\boxed{\;
\beta''(0)=\frac{-47145+21083\sqrt5}{500}<0.\;}
\]
\end{corollary}
\begin{proof}
第一条恒等式是定理
\ref{thm:xi-real-input40-atomic-factor-thermodynamic-invariance} 在 \(u=e^\theta\) 下的直接改写。

由
\[
\beta(\theta)=e^\theta e^{-2P(\theta)}
\]
取对数微分，得到
\[
\frac{\beta'(\theta)}{\beta(\theta)}=1-2P'(\theta),
\]
从而
\[
\beta'(\theta)=\bigl(1-2P'(\theta)\bigr)\beta(\theta).
\]
再对上式求导，即得
\[
\beta''(\theta)
=
\Bigl(\bigl(1-2P'(\theta)\bigr)^2-2P''(\theta)\Bigr)\beta(\theta).
\]

在 \(\theta=0\) 处，由 \(\lambda(1)=\varphi^2\) 可知
\[
\beta(0)=\lambda(1)^{-2}=\varphi^{-4}=\frac{7-3\sqrt5}{2}.
\]
再将表 \ref{tab:real_input_40_output_potential_cumulants_closed} 中的
\(P'(0)\) 与 \(P''(0)\) 代入前两条微分公式，即得
\[
\beta'(0)=\bigl(1-2P'(0)\bigr)\beta(0)=\frac{-49+22\sqrt5}{5},
\]
\[
\beta''(0)
=
\Bigl(\bigl(1-2P'(0)\bigr)^2-2P''(0)\Bigr)\beta(0)
=
\frac{-47145+21083\sqrt5}{500}.
\]
由其精确值的数值符号即可得到
\(
\beta'(0)>0
\)
与
\(
\beta''(0)<0
\)。
证毕。
\end{proof}


\begin{corollary}[真实输入 \texorpdfstring{$40$}{40} 状态核零温极限下核心 Abel 常数的闭式分离]\label{cor:xi-real-input40-zero-temp-core-abel-constant-splitting}
\leanverified{paper\_xi\_real\_input40\_zero\_temp\_core\_abel\_constant\_splitting}
沿用附录 \ref{app:real-input-40-zeta-u} 的记号。对真实输入 \(40\) 状态核，
\[
\Delta(z,u)=(1-uz^2)\,F(z,u),
\qquad
\zeta_{\mathrm{core}}(z,u):=F(z,u)^{-1}.
\]
令 \(\log\mathfrak M(u)\) 与 \(\log\mathfrak M_{\mathrm{core}}(u)\) 分别表示总体与核心部分的 Abel 有限部常数，并定义零温极限
\[
\log\mathfrak M_\infty:=\lim_{u\to\infty}\log\mathfrak M(u),
\qquad
\log\mathfrak M_{\infty,\mathrm{core}}
:=
\lim_{u\to\infty}\log\mathfrak M_{\mathrm{core}}(u).
\]
再令 \(B_\infty\) 为零温四态地基矩阵，\(\rho(B_\infty)\) 为其 Perron 根，并记
\[
b:=\rho(B_\infty)^{-1}.
\]
则有闭式分离
\[
\boxed{\;
\log\mathfrak M_\infty
=
\log\mathfrak M_{\infty,\mathrm{core}}+b,
\qquad
b=\rho(B_\infty)^{-1}.\;}
\]
等价地，
\[
\boxed{\;
\log\mathfrak M_{\infty,\mathrm{core}}
=
\log\mathfrak M_\infty-b
\approx 0.004215453121.\;}
\]
因此零温极限中由原子 Euler 因子 \((1-uz^2)^{-1}\) 所贡献的 Abel 常数位移，不是自由参数，而恰由地基系统 \(B_\infty\) 的 Perron 逆根刚性锁定。
\end{corollary}
\begin{proof}
定理 \ref{thm:xi-finite-euler-primitive-surgery-strict-additivity} 在 \(J=1\)、\(m_1=1\)、\(E_1=1\)、\(\ell_1=2\) 的特例给出
\[
\log\mathfrak M(u)
=
\log\mathfrak M_{\mathrm{core}}(u)+u\,z_\star(u)^2.
\]
另一方面，命题 \ref{prop:real-input-40-pressure-freezing} 给出
\[
\lambda(u)=c\sqrt u\,(1+o(1))
\qquad(u\to\infty),
\]
因此
\[
u\,z_\star(u)^2
=
u\,\lambda(u)^{-2}
\longrightarrow
\frac{1}{c^2}.
\]
由命题 \ref{prop:real-input-40-logM-infty} 与零温四态极限的地基证书知
\[
\frac{1}{c^2}=b=\rho(B_\infty)^{-1},
\qquad
\log\mathfrak M_\infty=\lim_{u\to\infty}\log\mathfrak M(u)
\]
确实存在。于是 \(\log\mathfrak M_{\mathrm{core}}(u)=\log\mathfrak M(u)-u z_\star(u)^2\) 亦收敛，并满足
\[
\log\mathfrak M_{\infty,\mathrm{core}}
=
\log\mathfrak M_\infty-b.
\]
将命题 \ref{prop:real-input-40-logM-infty} 中的数值
\(
\log\mathfrak M_\infty\approx 0.282290255262
\)
与
\(
b\approx 0.278074802141
\)
代入，便得到
\(
\log\mathfrak M_{\infty,\mathrm{core}}\approx 0.004215453121
\)。
证毕。
\end{proof}

\begin{theorem}[真实输入 \texorpdfstring{$40$}{40} 状态核的零温 Abel 常数跃迁由四次代数数唯一锁定]\label{thm:xi-real-input40-zero-temp-abel-jump-quartic-locked}
\leanverified{paper\_xi\_real\_input40\_zero\_temp\_abel\_jump\_quartic\_locked}
沿用推论 \ref{cor:xi-real-input40-zero-temp-core-abel-constant-splitting} 与命题 \ref{prop:real-input-40-logM-infty} 的记号。令
\[
b:=\lim_{u\to\infty}u\,z_\star(u)^2.
\]
则 \(b\) 是四次方程
\[
1-6b+9b^2-b^3-b^4=0
\]
的最大正根，且
\[
b\approx 0.27807480214113.
\]
进一步，真实输入 \(40\) 状态核满足精确极限位移
\[
\boxed{\;
\lim_{u\to\infty}
\bigl(\log\mathfrak M(u)-\log\mathfrak M_{\mathrm{core}}(u)\bigr)
=
b.\;}
\]
从而
\[
\boxed{\;
\lim_{u\to\infty}
\frac{\mathfrak M(u)}{\mathfrak M_{\mathrm{core}}(u)}
=
e^{b}
\approx 1.320584976098.\;}
\]
\end{theorem}
\begin{proof}
命题 \ref{prop:real-input-40-logM-infty} 已证明 \(b\) 正是方程
\[
1-6b+9b^2-b^3-b^4=0
\]
的最大正根，并给出其数值近似。

另一方面，推论 \ref{cor:xi-real-input40-zero-temp-core-abel-constant-splitting} 给出
\[
\log\mathfrak M_\infty
=
\log\mathfrak M_{\infty,\mathrm{core}}+b.
\]
因此
\[
\lim_{u\to\infty}
\bigl(\log\mathfrak M(u)-\log\mathfrak M_{\mathrm{core}}(u)\bigr)
=
b.
\]
再由
\[
\frac{\mathfrak M(u)}{\mathfrak M_{\mathrm{core}}(u)}
=
\exp\!\bigl(\log\mathfrak M(u)-\log\mathfrak M_{\mathrm{core}}(u)\bigr),
\]
对上式取极限，即得
\[
\lim_{u\to\infty}
\frac{\mathfrak M(u)}{\mathfrak M_{\mathrm{core}}(u)}
=
e^b
\approx 1.320584976098.
\]
证毕。
\end{proof}

\begin{proposition}[零温 prefactor 常数 \texorpdfstring{$C_\infty$}{C∞} 的四次整系数方程]\label{prop:xi-real-input40-cinf-quartic-elimination}
\leanverified{paper\_xi\_real\_input40\_cinf\_quartic\_elimination}
沿用定理 \ref{thm:xi-real-input40-zero-temp-abel-jump-quartic-locked} 的记号。定义
\[
C_\infty:=
\frac{1}{2b(1-b)(6-18b+3b^2+4b^3)}.
\]
则 \(C_\infty\) 满足四次整系数方程
\[
\boxed{\;
88864\,C_\infty^4
-22216\,C_\infty^3
-275864\,C_\infty^2
-3646\,C_\infty
+1=0.\;}
\]
换言之，零温 prefactor 常数本身就是一个四次代数数，而不必再经由中间变量 \(b\) 间接描述。
\end{proposition}
\begin{proof}
由定理 \ref{thm:xi-real-input40-zero-temp-abel-jump-quartic-locked}，
\[
b^4+b^3-9b^2+6b-1=0.
\]
另一方面，\(C_\infty\) 的定义等价于
\[
2C_\infty\,b(1-b)(6-18b+3b^2+4b^3)-1=0.
\]
于是 \((b,C_\infty)\) 满足一组二元代数方程。对上述两式关于变量 \(b\) 取 Sylvester 结式，可得到仅依赖于 \(C_\infty\) 的整系数多项式；直接化简该结式恰为
\[
88864\,C_\infty^4
-22216\,C_\infty^3
-275864\,C_\infty^2
-3646\,C_\infty
+1.
\]
因此 \(C_\infty\) 必为该四次多项式的根。相应消元核验可由脚本
\texttt{scripts/exp\_xi\_real\_input\_40\_c\_infty\_quartic\_audit.py}
逐项复现。证毕。
\end{proof}

\noindent 在上述四次代数闭式之外，命题 \ref{prop:real-input-40-logM-infty} 的 Möbius--Euler 表达式还允许直接量化 Abel 有限部常数的截断误差。
\begin{theorem}[Abel 有限部常数的显式截断误差界]\label{thm:xi-real-input40-logM-infty-truncation-bound}
沿用命题 \ref{prop:real-input-40-logM-infty} 的记号。定义部分和
\[
\log\mathfrak M_\infty^{(N)}
:=
\log C_\infty
-\sum_{2\le m\le N}\frac{\mu(m)}{m}
\log\Bigl[(1-b^m)\bigl(1-6b^m+9b^{2m}-b^{3m}-b^{4m}\bigr)\Bigr].
\]
则对每个 \(N\ge 1\) 都有显式尾界
\[
\boxed{\;
\bigl|\log\mathfrak M_\infty-\log\mathfrak M_\infty^{(N)}\bigr|
\le
\frac{K_b}{N+1}\frac{b^{N+1}}{1-b},
\qquad
K_b:=\frac{7}{1-7b^2}.\;}
\]
特别地，上述 Möbius--Euler 级数绝对收敛，并且以几何速度收敛到 \(\log\mathfrak M_\infty\)。
\leanverified{paper\_xi\_real\_input40\_logM\_infty\_truncation\_bound}
\end{theorem}
\begin{proof}
设
\[
f(t):=1-6t+9t^2-t^3-t^4.
\]
由直接计算，
\[
f(1/4)=\frac{11}{256}>0,
\qquad
f(1/3)=-\frac{4}{81}<0.
\]
由于 \(b\) 是 \(f\) 的最大正根，遂有
\[
b\in(1/4,1/3),
\]
从而对一切 \(m\ge 2\) 都满足
\[
b^m\le b^2<\frac19.
\]

再定义
\[
g(t):=(1-t)(1-6t+9t^2-t^3-t^4)=1-7t+15t^2-10t^3+t^5.
\]
对 \(t\in[0,1]\)，有
\[
15t^2-10t^3+t^5=t^2(15-10t+t^3)\ge 0,
\]
故
\[
g(t)\ge 1-7t.
\]
取 \(t=b^m\)（\(m\ge 2\)），便得到
\[
g(b^m)\ge 1-7b^m\ge 1-7b^2>0.
\]
另一方面，由 \(0<b<1\) 可知 \(0<g(b^m)\le 1\)，于是
\[
\bigl|\log g(b^m)\bigr|
=-\log g(b^m)
\le -\log(1-7b^m).
\]
再利用 \(0\le x<1\) 时的基本不等式
\[
-\log(1-x)\le \frac{x}{1-x},
\]
便得
\[
\bigl|\log g(b^m)\bigr|
\le \frac{7b^m}{1-7b^m}
\le \frac{7}{1-7b^2}\,b^m
=K_b\,b^m.
\]
因此
\[
\bigl|\log\mathfrak M_\infty-\log\mathfrak M_\infty^{(N)}\bigr|
\le
\sum_{m>N}\frac{|\mu(m)|}{m}\,\bigl|\log g(b^m)\bigr|
\le
K_b\sum_{m>N}\frac{b^m}{m}.
\]
对 \(m>N\) 有 \(1/m\le 1/(N+1)\)，故
\[
\sum_{m>N}\frac{b^m}{m}
\le
\frac{1}{N+1}\sum_{m>N}b^m
=
\frac{1}{N+1}\frac{b^{N+1}}{1-b}.
\]
合并即得所示显式尾界。由该上界可见级数绝对收敛，且尾项按 \(b^{N+1}\) 几何衰减。证毕。
\end{proof}

在上述零温 Abel 跃迁已被四次代数数完全锁定之后，输出位势零温尾区与 core 系统之间还可进一步抽取出两条彼此闭合的热力学结论。

\begin{theorem}[输出位势零温尾部的 core 热力学同一化]\label{thm:xi-real-input40-core-thermodynamic-tail-equivalence}
\leanverified{paper\_xi\_real\_input40\_core\_thermodynamic\_tail\_equivalence}
沿用推论 \ref{cor:xi-real-input40-zero-temp-core-abel-constant-splitting} 与命题
\ref{prop:real-input-40-pressure-freezing} 的记号。记
\[
\lambda_{\mathrm{core}}(u)
\]
为 \(\zeta_{\mathrm{core}}(z,u)=F(z,u)^{-1}\) 的主增长率，并定义对应热力学压力
\[
\mathcal P_{\mathrm{core}}(\theta):=\log \lambda_{\mathrm{core}}(e^\theta).
\]
则存在 \(u_0>1\)（等价地 \(\theta_0:=\log u_0>0\)），使得对所有 \(\theta\ge \theta_0\) 都有
\[
\boxed{\;
\lambda(e^\theta)=\lambda_{\mathrm{core}}(e^\theta),
\qquad
P(\theta)=\mathcal P_{\mathrm{core}}(\theta).\;}
\]
此外，在 \(\theta=0\) 亦有
\[
P(0)=\mathcal P_{\mathrm{core}}(0).
\]
若定义右尾阈值
\[
\alpha_{\mathrm{tail}}:=P'(\theta_0),
\]
并记其右尾 Legendre 分支
\[
I_{\mathrm{core}}^{\mathrm{tail}}(\alpha):=
\sup_{\theta\ge \theta_0}
\Bigl\{\theta\alpha-\bigl(\mathcal P_{\mathrm{core}}(\theta)-\mathcal P_{\mathrm{core}}(0)\bigr)\Bigr\},
\]
则在右端零温支上成立
\[
\boxed{\;
I(\alpha)=I_{\mathrm{core}}^{\mathrm{tail}}(\alpha)
\qquad
\bigl(\alpha\in[\alpha_{\mathrm{tail}},\tfrac12]\bigr).\;}
\]
特别地，若定义
\[
\alpha_{\max}^{\mathrm{core}}
:=
\lim_{\theta\to+\infty}\mathcal P_{\mathrm{core}}'(\theta),
\qquad
h_{\mathrm{ground}}^{\mathrm{core}}
:=
\lim_{\theta\to+\infty}\left(\mathcal P_{\mathrm{core}}(\theta)-\frac{\theta}{2}\right),
\]
则有
\[
\boxed{\;
\alpha_{\max}^{\mathrm{core}}=\frac12,
\qquad
h_{\mathrm{ground}}^{\mathrm{core}}=\log c>0.\;}
\]
\end{theorem}
\begin{proof}
由命题 \ref{prop:real-input-40-pressure-freezing}，
\[
\lambda(u)=c\sqrt u\,(1+o(1))
\qquad(u\to\infty),
\]
故主极点 \(z_\star(u)=\lambda(u)^{-1}\) 满足
\[
u\,z_\star(u)^2\longrightarrow c^{-2}.
\]
再由推论 \ref{cor:xi-real-input40-zero-temp-core-abel-constant-splitting}，
\[
c^{-2}=b<1.
\]
因此存在 \(u_0>1\) 使得对所有 \(u\ge u_0\) 都有
\[
u\,z_\star(u)^2<1.
\]
于是原子因子 \(1-uz^2\) 在主极点 \(z=z_\star(u)\) 处解析且不为零。将定理
\ref{thm:xi-finite-euler-primitive-surgery-strict-additivity} 应用于分解
\[
\Delta(z,u)=(1-uz^2)\,F(z,u)
\]
（此处 \(m=1,E=1,\ell=2\)），可知主增长率完全由 \(\zeta_{\mathrm{core}}(z,u)=F(z,u)^{-1}\) 决定，从而
\[
\lambda(u)=\lambda_{\mathrm{core}}(u)
\qquad(u\ge u_0).
\]
写成 \(\theta=\log u\) 即得
\[
P(\theta)=\mathcal P_{\mathrm{core}}(\theta)
\qquad(\theta\ge\theta_0).
\]

另一方面，在 \(u=1\) 时
\[
z_\star(1)=\lambda(1)^{-1}=\varphi^{-2},
\qquad
z_\star(1)^2=\varphi^{-4}<1,
\]
故同一定理在 \(u=1\) 处亦可应用，得到
\[
\lambda(1)=\lambda_{\mathrm{core}}(1),
\qquad
P(0)=\mathcal P_{\mathrm{core}}(0).
\]

再取任意 \(\alpha\in[\alpha_{\mathrm{tail}},\frac12)\)。由于 \(P\) 在实轴上凸且
\[
\lim_{\theta\to+\infty}P'(\theta)=\frac12,
\]
其 Legendre--Fenchel 对偶在该区间的极值必由某个 \(\theta\ge \theta_0\) 取得。于是
\[
I(\alpha)
=
\sup_{\theta\ge \theta_0}
\Bigl\{\theta\alpha-\bigl(P(\theta)-P(0)\bigr)\Bigr\}.
\]
在 \([\theta_0,\infty)\) 上有 \(P(\theta)=\mathcal P_{\mathrm{core}}(\theta)\)，且 \(P(0)=\mathcal P_{\mathrm{core}}(0)\)，故
\[
I(\alpha)
=
\sup_{\theta\ge \theta_0}
\Bigl\{\theta\alpha-\bigl(\mathcal P_{\mathrm{core}}(\theta)-\mathcal P_{\mathrm{core}}(0)\bigr)\Bigr\}
=
I_{\mathrm{core}}^{\mathrm{tail}}(\alpha).
\]
端点 \(\alpha=\frac12\) 由连续延拓得到同样结论。

最后，由尾部压力同一化，
\[
\mathcal P_{\mathrm{core}}(\theta)=P(\theta)
\qquad(\theta\gg 1),
\]
故其右端斜率与冻结常数与完整模型一致。再由定理
\ref{thm:xi-terminal-real-input-positive-entropy-freezing-face}，
\[
\alpha_{\max}=\frac12,
\qquad
h_{\mathrm{ground}}=\log c>0,
\]
遂得
\[
\alpha_{\max}^{\mathrm{core}}=\frac12,
\qquad
h_{\mathrm{ground}}^{\mathrm{core}}=\log c.
\]
证毕。
\end{proof}

\begin{proposition}[零温 Abel 有限部的原子占比]\label{prop:xi-real-input40-zero-temp-abel-atomic-share}
沿用推论 \ref{cor:xi-real-input40-zero-temp-core-abel-constant-splitting} 与定理
\ref{thm:xi-real-input40-zero-temp-abel-jump-quartic-locked} 的记号。则
\[
\boxed{\;
\log\mathfrak M_{\infty,\mathrm{core}}
=
\log\mathfrak M_\infty-b
\approx 0.00421545312087,\;}
\]
\[
\boxed{\;
\mathfrak M_{\infty,\mathrm{core}}
=
\exp(\log\mathfrak M_\infty-b)
\approx 1.00422435064134.\;}
\]
并且原子 Euler 因子 \((1-uz^2)^{-1}\) 在零温 Abel 有限部中的加性占比满足
\[
\boxed{\;
\frac{b}{\log\mathfrak M_\infty}
\approx 0.9850669548725387,\qquad
\frac{\log\mathfrak M_{\infty,\mathrm{core}}}{\log\mathfrak M_\infty}
\approx 0.0149330451274613.\;}
\]
换言之，零温有限部的 \(98.506695487\ldots\%\) 由单一长度 \(2\) 的原子 Euler 因子承担，而 core 剩余部分仅占 \(1.493304512\ldots\%\)。
\end{proposition}
\begin{proof}
由推论 \ref{cor:xi-real-input40-zero-temp-core-abel-constant-splitting}，
\[
\log\mathfrak M_{\infty,\mathrm{core}}=\log\mathfrak M_\infty-b.
\]
将
\[
\log\mathfrak M_\infty\approx 0.282290255262,
\qquad
b\approx 0.27807480214113
\]
代入，即得
\[
\log\mathfrak M_{\infty,\mathrm{core}}\approx 0.00421545312087.
\]
对其取指数便得到
\[
\mathfrak M_{\infty,\mathrm{core}}
=
\exp(\log\mathfrak M_\infty-b)
\approx 1.00422435064134.
\]
最后，两项加性占比正是
\[
\frac{b}{\log\mathfrak M_\infty},
\qquad
\frac{\log\mathfrak M_\infty-b}{\log\mathfrak M_\infty},
\]
直接代入上述数值即得所示比例。证毕。
\end{proof}

\begin{theorem}[window-\texorpdfstring{$6$}{6} 的中心异常与非中心异常存在规范直和分裂]\label{thm:xi-window6-central-noncentral-anomaly-splitting}
沿用定理 \ref{thm:xi-foldbin-gauge-representation-distribution-law-m6-spectrum}、
定理 \ref{thm:xi-window6-center-exact-splitting-bdry-cyclic}
与命题 \ref{prop:xi-window6-boundary-eight-sector-central-decomposition} 的记号。令
\[
X_6^{(2)}:=\{w\in X_6:\ d_6(w)=2\},
\qquad
X_6^{(\ge 3)}:=\{w\in X_6:\ d_6(w)\ge 3\}.
\]
在符号坐标识别
\[
\bigl(G_6^{\mathrm{bin}}\bigr)^{\mathrm{ab}}
\cong
(\ZZ/2\ZZ)^{X_6}
\]
下，定义
\[
Q_6:= (\ZZ/2\ZZ)^{X_6^{(\ge 3)}}.
\]
则存在规范直和分裂
\[
\boxed{\;
\bigl(G_6^{\mathrm{bin}}\bigr)^{\mathrm{ab}}
\cong
Z\!\bigl(G_6^{\mathrm{bin}}\bigr)\oplus Q_6
\cong
(\ZZ/2\ZZ)^8\oplus (\ZZ/2\ZZ)^{13}.\;}
\]
其中 \(Z(G_6^{\mathrm{bin}})\) 通过自然交换化映射嵌入 \(\bigl(G_6^{\mathrm{bin}}\bigr)^{\mathrm{ab}}\)，而 \(Q_6\) 恰由全部 \(d_6=3,4\) 纤维所贡献的非中心异常坐标组成。
\par
进一步，boundary 中心子群
\[
Z_{\partial}
=
\langle \tau_w:\ w\in X_6^{\mathrm{bdry}}\rangle
\cong
(\ZZ/2\ZZ)^3
\]
是 \(Z(G_6^{\mathrm{bin}})\) 的一个规范直和因子，从而
\[
\boxed{\;
Z\!\bigl(G_6^{\mathrm{bin}}\bigr)/Z_{\partial}
\cong
(\ZZ/2\ZZ)^5,
\qquad
\bigl(G_6^{\mathrm{bin}}\bigr)^{\mathrm{ab}}/Z_{\partial}
\cong
(\ZZ/2\ZZ)^{18}.\;}
\]
\end{theorem}
\leanverified{paper\_xi\_window6\_central\_noncentral\_anomaly\_splitting}
\begin{proof}
定理 \ref{thm:xi-foldbin-gauge-representation-distribution-law-m6-spectrum} 给出
\[
\bigl(G_6^{\mathrm{bin}}\bigr)^{\mathrm{ab}}
\cong
(\ZZ/2\ZZ)^{X_6}.
\]
由表 \ref{tab:fold_bin_gauge_m6_fibers} 可知 \(m=6\) 时纤维直方图为
\[
d=2:8,\qquad d=3:4,\qquad d=4:9.
\]
因此
\[
|X_6^{(2)}|=8,
\qquad
|X_6^{(\ge 3)}|=4+9=13,
\]
且坐标格按集合分拆 \(X_6=X_6^{(2)}\sqcup X_6^{(\ge 3)}\) 立刻得到
\[
(\ZZ/2\ZZ)^{X_6}
\cong
(\ZZ/2\ZZ)^{X_6^{(2)}}
\oplus
(\ZZ/2\ZZ)^{X_6^{(\ge 3)}}.
\]
另一方面，定理 \ref{thm:xi-window6-center-exact-splitting-bdry-cyclic} 已证明
\[
Z\!\bigl(G_6^{\mathrm{bin}}\bigr)
\cong
(\ZZ/2\ZZ)^{X_6^{(2)}}
\cong
(\ZZ/2\ZZ)^8.
\]
将此与上式合并，便得到
\[
\bigl(G_6^{\mathrm{bin}}\bigr)^{\mathrm{ab}}
\cong
Z\!\bigl(G_6^{\mathrm{bin}}\bigr)\oplus Q_6
\cong
(\ZZ/2\ZZ)^8\oplus(\ZZ/2\ZZ)^{13},
\]
其中 \(Q_6=(\ZZ/2\ZZ)^{X_6^{(\ge 3)}}\) 正是全部非中心异常坐标的和。
\par
再由命题 \ref{prop:xi-window6-boundary-eight-sector-central-decomposition} 与定理 \ref{thm:xi-window6-center-exact-splitting-bdry-cyclic}，
\[
Z_{\partial}\cong (\ZZ/2\ZZ)^3
\]
是 \(Z(G_6^{\mathrm{bin}})\) 的规范直和因子，且其余维数为 \(5\)。故
\[
Z\!\bigl(G_6^{\mathrm{bin}}\bigr)/Z_{\partial}\cong (\ZZ/2\ZZ)^5.
\]
最后将 \(Z_{\partial}\) 在上述直和分裂中商去，得到
\[
\bigl(G_6^{\mathrm{bin}}\bigr)^{\mathrm{ab}}/Z_{\partial}
\cong
\left(Z\!\bigl(G_6^{\mathrm{bin}}\bigr)/Z_{\partial}\right)\oplus Q_6
\cong
(\ZZ/2\ZZ)^5\oplus(\ZZ/2\ZZ)^{13}
\cong
(\ZZ/2\ZZ)^{18}.
\]
证毕。
\end{proof}

\paragraph{有限矩刚性、dyadic Mellin 读出与冻结相仿射谱}
下列条目把有限分辨率上的纤维多重度谱矩刚性、dyadic 预算下容量曲线对幂和的积分读出、冻结相压力谱的仿射锁定，以及 window-\(6\) 锚点的完整有限尺度闭式统一压缩为一组可直接调用的终端结论。

\begin{theorem}[纤维多重度谱的有限矩刚性：前 \texorpdfstring{$2s_m$}{2s\_m} 个幂和已决定全部直方图]\label{thm:xi-fiber-spectrum-finite-moment-rigidity}
\leanverified{paper\_xi\_fiber\_spectrum\_finite\_moment\_rigidity}
固定分辨率 \(m\)。把幂和记号延拓到整数 \(q\ge 0\)：
\[
S_q(m):=\sum_{x\in X_m}d_m(x)^q,
\qquad
S_0(m)=|X_m|.
\]
设纤维多重度的不同取值为
\[
1\le \delta_{1,m}<\delta_{2,m}<\cdots<\delta_{s_m,m}=M_m,
\]
并记其重数
\[
n_{j,m}:=\#\{x\in X_m:\ d_m(x)=\delta_{j,m}\}
\qquad(1\le j\le s_m).
\]
则对所有整数 \(q\ge 0\) 都有严格分解
\[
\boxed{\;
S_q(m)=\sum_{j=1}^{s_m}n_{j,m}\,\delta_{j,m}^{\,q}.\;}
\]
相应普通生成函数
\[
R_m(z):=\sum_{q\ge 0}S_q(m)\,z^q
\]
为有理函数，且满足
\[
\boxed{\;
R_m(z)=\sum_{j=1}^{s_m}\frac{n_{j,m}}{1-\delta_{j,m}z}
=
\frac{A_m(z)}{\prod_{j=1}^{s_m}(1-\delta_{j,m}z)},\qquad
\deg A_m<s_m.\;}
\]
若定义特征多项式
\[
\chi_m(X):=\prod_{j=1}^{s_m}(X-\delta_{j,m})
=
X^{s_m}-c_{1,m}X^{s_m-1}+\cdots+(-1)^{s_m}c_{s_m,m},
\]
则序列 \(\bigl(S_q(m)\bigr)_{q\ge 0}\) 满足阶数恰为 \(s_m\) 的线性递推
\[
\boxed{\;
S_{q+s_m}(m)
=
c_{1,m}S_{q+s_m-1}(m)-\cdots+(-1)^{s_m+1}c_{s_m,m}S_q(m)
\qquad(q\ge 0).\;}
\]
进一步，截断向量
\[
S_0(m),S_1(m),\dots,S_{2s_m-1}(m)
\]
已足以唯一恢复全部 \(\{(\delta_{j,m},n_{j,m})\}_{j=1}^{s_m}\)，从而唯一恢复尾函数
\[
N_m(t)=\#\{x\in X_m:\ d_m(x)\ge t\}
\]
与连续容量曲线
\[
\mathcal C_m^{\mathrm{cont}}(T)=\sum_{x\in X_m}\min\bigl(d_m(x),T\bigr).
\]
\end{theorem}
\begin{proof}
按纤维多重度值分层即可得到
\[
S_q(m)=\sum_{x\in X_m}d_m(x)^q=\sum_{j=1}^{s_m}n_{j,m}\delta_{j,m}^{\,q},
\]
这就给出第一式。
\par
于是对形式幂级数有
\[
\begin{aligned}
R_m(z)
&=
\sum_{q\ge 0}\sum_{j=1}^{s_m}n_{j,m}\delta_{j,m}^{\,q}z^q\\
&=
\sum_{j=1}^{s_m}n_{j,m}\sum_{q\ge 0}(\delta_{j,m}z)^q
=
\sum_{j=1}^{s_m}\frac{n_{j,m}}{1-\delta_{j,m}z}.
\end{aligned}
\]
把各项通分便得到所示有理表达。由于各极点 \(z=\delta_{j,m}^{-1}\) 两两不同且对应留数 \(n_{j,m}\delta_{j,m}^{-1}\) 非零，分母不存在约消，故最小分母次数恰为 \(s_m\)。因此 \(\bigl(S_q(m)\bigr)_{q\ge 0}\) 的最小线性递推阶正是 \(s_m\)，其特征多项式即为 \(\chi_m\)，从而得到上式递推。
\par
再看唯一恢复。令 Hankel 主块
\[
H_m:=\bigl(S_{r+s}(m)\bigr)_{0\le r,s\le s_m-1}.
\]
由第一式可分解为
\[
H_m=V_m\,\mathrm{diag}(n_{1,m},\dots,n_{s_m,m})\,V_m^{\mathsf T},
\]
其中
\[
V_m:=\bigl(\delta_{j,m}^{\,r}\bigr)_{0\le r\le s_m-1,\ 1\le j\le s_m}
\]
为 Vandermonde 矩阵。由于 \(\delta_{j,m}\) 两两不同且 \(n_{j,m}>0\)，故 \(V_m\) 可逆，从而 \(H_m\) 非奇异。于是递推系数 \((c_{1,m},\dots,c_{s_m,m})\) 由线性系统
\[
\bigl(S_{r+s}(m)\bigr)_{0\le r,s\le s_m-1}\,
\begin{pmatrix}
(-1)^{s_m+1}c_{s_m,m}\\
\vdots\\
-c_{2,m}\\
c_{1,m}
\end{pmatrix}
=
\begin{pmatrix}
S_{s_m}(m)\\
\vdots\\
S_{2s_m-1}(m)
\end{pmatrix}
\]
唯一确定；右侧正好只用到前 \(2s_m\) 个幂和。求得 \(\chi_m\) 后，其根集便是 \(\{\delta_{j,m}\}_{j=1}^{s_m}\)。再把这些值代回前 \(s_m\) 个矩方程
\[
S_q(m)=\sum_{j=1}^{s_m}n_{j,m}\delta_{j,m}^{\,q}
\qquad(0\le q\le s_m-1),
\]
即可由 Vandermonde 线性系统唯一求出 \((n_{1,m},\dots,n_{s_m,m})\)。
\par
最后，
\[
N_m(t)=\sum_{\delta_{j,m}\ge t}n_{j,m},
\qquad
\mathcal C_m^{\mathrm{cont}}(T)=\sum_{j=1}^{s_m}n_{j,m}\min(\delta_{j,m},T),
\]
故尾函数与容量曲线也随之唯一恢复。证毕。
\end{proof}

\begin{theorem}[容量亏损的 Mellin--Stieltjes 精确反演与 dyadic 读出]\label{thm:xi-dyadic-oracle-capacity-mellin-recovery}
沿用定义 \ref{def:fold-capacity-tail}、命题 \ref{prop:fold-capacity-tail-equivalence} 与命题 \ref{prop:fold-moment-mellin-stieltjes} 的记号。对固定 \(m\) 与任意实数 \(q>1\)，有精确恒等式
\[
\boxed{\;
S_q(m)
=
q(q-1)\int_{0}^{\infty}T^{q-2}\Bigl(2^m-\mathcal C_m^{\mathrm{cont}}(T)\Bigr)\,\dd T.\;}
\]
进一步，把离散预言机容量沿 dyadic 尺度连续化为
\[
C_m(B):=\mathcal C_m^{\mathrm{cont}}(2^B)
\qquad(B\in \RR),
\]
则对 Lebesgue--a.e.\ \(B\in\RR\) 都有
\[
\boxed{\;
\frac{\dd}{\dd B}C_m(B)
=
(\log 2)\,2^B\,N_m(2^B).\;}
\]
从而
\[
\boxed{\;
S_q(m)
=
q\log 2\int_{-\infty}^{\infty}2^{qB}N_m(2^B)\,\dd B
\;}
\]
以及
\[
\boxed{\;
S_q(m)
=
q(q-1)\log 2
\int_{-\infty}^{\infty}
2^{(q-1)B}\Bigl(2^m-C_m(B)\Bigr)\,\dd B.\;}
\]
因此一旦 dyadic 预算曲线 \(C_m(B)\) 已知，全部正阶幂和 \(\{S_q(m)\}_{q>1}\) 都可由显式 Mellin 积分直接重构；换言之，\(S_q(m)\) 不仅是尾函数 \(N_m\) 的 Mellin 变换，同时也是容量亏损 \(2^m-\mathcal C_m^{\mathrm{cont}}\) 的 Mellin--Stieltjes 精确反演。
\end{theorem}
\begin{proof}
由命题 \ref{prop:fold-capacity-tail-equivalence}，
\[
\mathcal C_m^{\mathrm{cont}}(T)=\int_{0}^{T}N_m(t)\,\dd t.
\]
又因
\[
\sum_{x\in X_m}d_m(x)=2^m,
\]
故
\[
2^m-\mathcal C_m^{\mathrm{cont}}(T)
=
\int_{T}^{\infty}N_m(t)\,\dd t.
\]
于是对 \(q>1\) 可交换积分次序：
\[
\begin{aligned}
q(q-1)\int_{0}^{\infty}T^{q-2}\Bigl(2^m-\mathcal C_m^{\mathrm{cont}}(T)\Bigr)\,\dd T
&=
q(q-1)\int_{0}^{\infty}T^{q-2}\left(\int_{T}^{\infty}N_m(t)\,\dd t\right)\dd T\\
&=
q(q-1)\int_{0}^{\infty}N_m(t)\left(\int_{0}^{t}T^{q-2}\,\dd T\right)\dd t\\
&=
q\int_{0}^{\infty}t^{q-1}N_m(t)\,\dd t\\
&=
S_q(m),
\end{aligned}
\]
最后一步正是命题 \ref{prop:fold-moment-mellin-stieltjes}。
\par
对 dyadic 变量，命题 \ref{prop:fold-capacity-tail-equivalence} 还给出 \(\mathcal C_m^{\mathrm{cont}}\) 在几乎处处可微，且
\[
\frac{\dd}{\dd T}\mathcal C_m^{\mathrm{cont}}(T)=N_m(T)
\qquad\text{a.e.}
\]
因此链式法则推出
\[
\frac{\dd}{\dd B}C_m(B)
=
\frac{\dd}{\dd B}\mathcal C_m^{\mathrm{cont}}(2^B)
=
(\log 2)\,2^B\,N_m(2^B)
\]
于几乎处处成立。
\par
再在
\[
S_q(m)=q\int_{0}^{\infty}t^{q-1}N_m(t)\,\dd t
\]
中作变量代换 \(t=2^B\)，利用 \(\dd t=(\log 2)2^B\,\dd B\)，得到
\[
S_q(m)=q\log 2\int_{-\infty}^{\infty}2^{qB}N_m(2^B)\,\dd B.
\]
同理，在
\[
S_q(m)=q(q-1)\int_{0}^{\infty}T^{q-2}\Bigl(2^m-\mathcal C_m^{\mathrm{cont}}(T)\Bigr)\,\dd T
\]
中令 \(T=2^B\)，便得到最后一式。证毕。
\end{proof}

\begin{theorem}[容量亏损指数的 Mellin--Laplace 变分与冻结暴露平面]\label{thm:xi-capacity-defect-exponent-mellin-laplace-face}
沿用定理 \ref{thm:xi-dyadic-oracle-capacity-mellin-recovery} 的记号，并设
\[
\alpha_*:=\lim_{m\to\infty}\frac{1}{m}\log M_m,
\qquad
g_*:=\lim_{m\to\infty}\frac{1}{m}\log K_m
\]
存在。对每个 \(\beta\in[0,\alpha_*]\)，定义容量亏损指数近似
\[
\Lambda_m(\beta):=
\frac{1}{m}\log\Bigl(2^m-\mathcal C_m^{\mathrm{cont}}(e^{\beta m})\Bigr).
\]
进一步假设极限
\[
\Lambda(\beta):=\lim_{m\to\infty}\Lambda_m(\beta)
\]
对每个 \(\beta\in[0,\alpha_*]\) 存在，且 \(\Lambda_m\to\Lambda\) 在 \([0,\alpha_*]\) 的每个紧子区间上一致收敛。则对任意实数 \(a>1\)，极限
\[
P(a):=\lim_{m\to\infty}\frac{1}{m}\log S_a(m)
\]
存在，并满足精确变分公式
\[
\boxed{\;
P(a)=\sup_{0\le \beta\le \alpha_*}\Bigl((a-1)\beta+\Lambda(\beta)\Bigr).\;}
\]
若进一步设
\[
\alpha_*^{(2)}:=\limsup_{m\to\infty}\frac{1}{m}\log M_m^{(2)}<\alpha_*,
\]
则对每个
\[
\alpha_*^{(2)}<\beta<\alpha_*
\]
都成立
\[
\boxed{\;
\Lambda(\beta)=\alpha_*+g_*.\;}
\]
因此容量亏损指数在严格谱隙区间上形成常值暴露平面；并且当
\[
a>a_{\mathrm c}:=\frac{\log\varphi-g_*}{\alpha_*-\alpha_*^{(2)}}
\]
时，冻结线性支
\[
P(a)=a\alpha_*+g_*
\]
正由该常值平面的右端在 Laplace 极限中暴露出来。
\end{theorem}
\begin{proof}
由定理 \ref{thm:xi-dyadic-oracle-capacity-mellin-recovery}，对任意 \(a>1\) 都有
\[
S_a(m)
=
a(a-1)\int_0^\infty T^{a-2}\Bigl(2^m-\mathcal C_m^{\mathrm{cont}}(T)\Bigr)\,\dd T.
\]
又因
\[
\mathcal C_m^{\mathrm{cont}}(T)=2^m
\qquad (T\ge M_m),
\]
故积分可截断到 \([0,M_m]\)。令 \(T=e^{\beta m}\)，则 \(\dd T=me^{\beta m}\dd\beta\)，从而
\[
S_a(m)
=
a(a-1)m
\int_0^{\frac1m\log M_m}
\exp\!\Bigl(
m\bigl((a-1)\beta+\Lambda_m(\beta)\bigr)
\Bigr)\,\dd\beta.
\]
由于
\[
\frac1m\log M_m\to \alpha_*,
\]
由于 \(\Lambda_m\to\Lambda\) 在 \([0,\alpha_*]\) 上可由紧性统一控制，故可直接对上式的指数积分应用 Laplace 原理，得到
\[
\lim_{m\to\infty}\frac1m\log S_a(m)
=
\sup_{0\le \beta\le \alpha_*}
\Bigl((a-1)\beta+\Lambda(\beta)\Bigr).
\]
这就得到了所示变分公式。

再设 \(\alpha_*^{(2)}<\beta<\alpha_*\)。由定义，
\[
\frac1m\log M_m^{(2)}\le \alpha_*^{(2)}+o(1),
\qquad
\frac1m\log M_m=\alpha_*+o(1),
\]
故对充分大的 \(m\) 必有
\[
M_m^{(2)}<e^{\beta m}<M_m.
\]
于是
\[
2^m-\mathcal C_m^{\mathrm{cont}}(e^{\beta m})
=
\sum_{x\in X_m}\bigl(d_m(x)-e^{\beta m}\bigr)_+
=
K_m\bigl(M_m-e^{\beta m}\bigr),
\]
因为只有最大纤维对上式仍有贡献。又由 \(\beta<\alpha_*\) 得
\[
\frac{e^{\beta m}}{M_m}
=
\exp\!\bigl(-(\alpha_*-\beta)m+o(m)\bigr)\to 0,
\]
故
\[
M_m-e^{\beta m}=M_m(1+o(1)).
\]
从而
\[
\Lambda(\beta)
=
\lim_{m\to\infty}\frac1m\log\bigl(K_mM_m(1+o(1))\bigr)
=
\alpha_*+g_*.
\]

最后，在上述常值区间上
\[
(a-1)\beta+\Lambda(\beta)
=
(a-1)\beta+\alpha_*+g_*,
\]
其值随 \(\beta\) 严格递增。因而该表达式在
\(
\beta\uparrow \alpha_*
\)
时趋于
\[
a\alpha_*+g_*.
\]
这与定理 \ref{thm:fold-pressure-freezing-threshold} 在 \(a>a_{\mathrm c}\) 时给出的冻结公式一致，故所述线性支正由该常值暴露平面的右端在 Laplace 极限中暴露出来。证毕。
\end{proof}


\begin{theorem}[尾指数的普适压力包络与超最大纤维的指数排除]\label{thm:xi-tail-exponent-pressure-envelope-no-supermax}
沿用附录 \ref{subsubsec:fold-fiber-groundstate-freezing-oracle-capacity} 的记号。对 \(\tau\ge 0\)，定义尾指数上界函数
\[
\overline\beta(\tau):=
\limsup_{m\to\infty}\frac{1}{m}\log N_m(e^{\tau m}),
\]
并约定若 \(N_m(e^{\tau m})=0\) 则 \(\log 0:=-\infty\)。则对每个整数 \(a\ge 1\) 恒有
\[
\boxed{\;
\overline\beta(\tau)\le P_a-a\tau.\;}
\]
从而得到普适包络
\[
\boxed{\;
\overline\beta(\tau)\le \inf_{a\ge 1}\bigl(P_a-a\tau\bigr).\;}
\]
进一步，由定理 \ref{thm:fold-pressure-groundstate-bounds} 的全局上界
\[
P_a\le a\alpha_*+\log\varphi
\qquad(a\ge 1),
\]
可推出：
\[
\boxed{\;
\tau>\alpha_*
\quad\Longrightarrow\quad
\overline\beta(\tau)=-\infty.\;}
\]
亦即：当 \(\tau>\alpha_*\) 时，指数尺度上不可能存在大小达到 \(e^{\tau m}\) 的纤维。
\par
若再假设严格基态间隙 \(\alpha_*^{(2)}<\alpha_*\)，并令
\[
a_{\mathrm c}:=\frac{\log\varphi-g_*}{\alpha_*-\alpha_*^{(2)}},
\qquad
a_0:=\min\{a\in\ZZ_{\ge 1}:\ a>a_{\mathrm c}\},
\]
则对所有 \(\tau\in[\alpha_*^{(2)},\alpha_*]\) 都有
\[
\boxed{\;
\overline\beta(\tau)\le g_*+a_0(\alpha_*-\tau).\;}
\]
\end{theorem}
\begin{proof}
若 \(d_m(x)\ge e^{\tau m}\)，则对任意整数 \(a\ge 1\) 有
\[
d_m(x)^a\ge e^{a\tau m}.
\]
于是
\[
S_a(m)=\sum_{x\in X_m}d_m(x)^a
\ge
\sum_{d_m(x)\ge e^{\tau m}}d_m(x)^a
\ge
e^{a\tau m}N_m(e^{\tau m}).
\]
两边取对数、除以 \(m\)、再取 \(\limsup\) 即得
\[
\overline\beta(\tau)\le P_a-a\tau.
\]
对 \(a\ge 1\) 取下确界，便得到普适包络。
\par
现在取 \(\tau>\alpha_*\)。由定理 \ref{thm:fold-pressure-groundstate-bounds}，
\[
\overline\beta(\tau)\le a(\alpha_*-\tau)+\log\varphi
\qquad(a\ge 1).
\]
由于 \(\alpha_*-\tau<0\)，令 \(a\to\infty\) 即得右端趋于 \(-\infty\)，故
\[
\overline\beta(\tau)=-\infty.
\]
因 \(N_m(e^{\tau m})\) 为非负整数，这也等价于：对充分大的 \(m\) 必有 \(N_m(e^{\tau m})=0\)。
\par
最后，在严格基态间隙下，定理 \ref{thm:fold-pressure-freezing-threshold} 给出
\[
P_{a_0}=a_0\alpha_*+g_*.
\]
把这代回一般上界便得到
\[
\overline\beta(\tau)\le P_{a_0}-a_0\tau=g_*+a_0(\alpha_*-\tau),
\]
结论成立。证毕。
\end{proof}

\begin{theorem}[冻结相的仿射刚性：任意两个冻结点唯一恢复 \texorpdfstring{$\alpha_*$}{alpha\_*} 与 \texorpdfstring{$g_*$}{g\_*}]\label{thm:xi-freezing-affine-rigidity-two-point-recovery}
在严格基态间隙 \(\alpha_*^{(2)}<\alpha_*\) 下，沿用定理 \ref{thm:fold-pressure-freezing-threshold} 的阈值
\[
a_{\mathrm c}:=\frac{\log\varphi-g_*}{\alpha_*-\alpha_*^{(2)}}.
\]
若 \(a,b\in \ZZ_{\ge 1}\) 且 \(a_{\mathrm c}<a<b\)，则冻结相内的压力值满足
\[
\boxed{\;
\alpha_*=\frac{P_b-P_a}{\,b-a\,},
\qquad
g_*=\frac{bP_a-aP_b}{\,b-a\,}.\;}
\]
特别地，对任意整数 \(a\ge 1\) 且 \(a>a_{\mathrm c}\)，有相邻差分公式
\[
\boxed{\;
\alpha_*=P_{a+1}-P_a,
\qquad
g_*=(a+1)P_a-aP_{a+1}.\;}
\]
并且对任意整数 \(a\ge 2\) 且 \(a-1>a_{\mathrm c}\)，冻结相内的全部二阶离散曲率严格消失：
\[
\boxed{\;
P_{a+1}-2P_a+P_{a-1}=0.\;}
\]
\end{theorem}
\begin{proof}
由定理 \ref{thm:fold-pressure-freezing-threshold}，对所有整数 \(q>a_{\mathrm c}\) 都有
\[
P_q=q\alpha_*+g_*.
\]
把 \(q=a,b\) 代入即得线性方程组
\[
P_a=a\alpha_*+g_*,
\qquad
P_b=b\alpha_*+g_*.
\]
联立消元便得到
\[
\alpha_*=\frac{P_b-P_a}{b-a},
\qquad
g_*=\frac{bP_a-aP_b}{b-a}.
\]
令 \(b=a+1\) 即得相邻差分公式。
\par
若 \(a-1>a_{\mathrm c}\)，则 \(a-1,a,a+1\) 三点都落在冻结相内，于是
\[
P_{a-1}=(a-1)\alpha_*+g_*,
\quad
P_a=a\alpha_*+g_*,
\quad
P_{a+1}=(a+1)\alpha_*+g_*.
\]
代入即得
\[
P_{a+1}-2P_a+P_{a-1}=0.
\]
证毕。
\end{proof}

\begin{proposition}[window-\texorpdfstring{$6$}{6} 的完整有限尺度谱学：幂和、尾函数、容量曲线与最小递推]\label{prop:xi-window6-complete-finite-scale-spectrum-closed-form}
取 \(m=6\)。由推论 \ref{cor:fold-bin-m6-fiber-hist}，其纤维多重度直方图为
\[
2:8,\qquad 3:4,\qquad 4:9,
\qquad |X_6|=21.
\]
因此 window-\(6\) 的纤维谱可完全显式化如下。
\par
对所有整数 \(q\ge 0\)，幂和满足
\[
\boxed{\;
S_q(6)=8\cdot 2^q+4\cdot 3^q+9\cdot 4^q.\;}
\]
相应普通生成函数为
\[
\boxed{\;
R_6(z):=\sum_{q\ge 0}S_q(6)z^q
=
\frac{21-125z+182z^2}{(1-2z)(1-3z)(1-4z)}.\;}
\]
因此 \(\bigl(S_q(6)\bigr)_{q\ge 0}\) 满足最小三阶递推
\[
\boxed{\;
S_{q+3}(6)=9S_{q+2}(6)-26S_{q+1}(6)+24S_q(6)
\qquad(q\ge 0),\;}
\]
其初值可取
\[
S_0(6)=21,\qquad S_1(6)=64,\qquad S_2(6)=212.
\]
尾函数为
\[
\boxed{\;
N_6(t)=
\begin{cases}
21,& 0\le t\le 2,\\
13,& 2<t\le 3,\\
9,& 3<t\le 4,\\
0,& t>4,
\end{cases}\;}
\]
连续容量曲线为
\[
\boxed{\;
\mathcal C_6^{\mathrm{cont}}(T)=
\begin{cases}
21T,& 0\le T\le 2,\\
13T+16,& 2\le T\le 3,\\
9T+28,& 3\le T\le 4,\\
64,& T\ge 4,
\end{cases}\;}
\]
而对 \(B\ge 0\) 的 dyadic 预算参数
\[
C_6(B):=\mathcal C_6^{\mathrm{cont}}(2^B)
\]
则有
\[
\boxed{\;
C_6(B)=
\begin{cases}
21\cdot 2^B,& 0\le B\le 1,\\
13\cdot 2^B+16,& 1\le B\le \log_2 3,\\
9\cdot 2^B+28,& \log_2 3\le B\le 2,\\
64,& B\ge 2.
\end{cases}\;}
\]
特别地，由 \(s_6=3\) 可知，前六个幂和 \(S_0(6),\dots,S_5(6)\) 已足以唯一重构整个 window-\(6\) 的纤维直方图、尾函数与容量曲线。
\end{proposition}
\begin{proof}
由推论 \ref{cor:fold-bin-m6-fiber-hist}，window-\(6\) 的不同纤维大小恰为 \(2,3,4\)，重数分别为 \(8,4,9\)。于是对所有 \(q\ge 0\)，直接有
\[
S_q(6)=8\cdot 2^q+4\cdot 3^q+9\cdot 4^q.
\]
这就给出了幂和闭式。
\par
普通生成函数因此为
\[
R_6(z)=\frac{8}{1-2z}+\frac{4}{1-3z}+\frac{9}{1-4z},
\]
通分后得到
\[
R_6(z)=\frac{21-125z+182z^2}{(1-2z)(1-3z)(1-4z)}.
\]
又因为
\[
(1-2z)(1-3z)(1-4z)=1-9z+26z^2-24z^3,
\]
按生成函数与递推的对应关系即可得到
\[
S_{q+3}(6)=9S_{q+2}(6)-26S_{q+1}(6)+24S_q(6).
\]
初值则由幂和闭式直接算得
\[
S_0(6)=8+4+9=21,\qquad
S_1(6)=16+12+36=64,\qquad
S_2(6)=32+36+144=212.
\]
\par
尾函数由定义
\[
N_6(t)=\#\{x\in X_6:\ d_6(x)\ge t\}
\]
直接从直方图读出：当 \(0\le t\le 2\) 时全部 \(21\) 条纤维都被计入；当 \(2<t\le 3\) 时仅剩 \(d=3,4\) 两层，共 \(4+9=13\) 条；当 \(3<t\le 4\) 时仅剩 \(9\) 条；当 \(t>4\) 时为空。
\par
再由命题 \ref{prop:fold-capacity-tail-equivalence}，
\[
\mathcal C_6^{\mathrm{cont}}(T)=\int_0^T N_6(t)\,\dd t.
\]
逐段积分便得到
\[
\mathcal C_6^{\mathrm{cont}}(T)=21T
\qquad(0\le T\le 2),
\]
\[
\mathcal C_6^{\mathrm{cont}}(T)=42+13(T-2)=13T+16
\qquad(2\le T\le 3),
\]
\[
\mathcal C_6^{\mathrm{cont}}(T)=55+9(T-3)=9T+28
\qquad(3\le T\le 4),
\]
而对 \(T\ge 4\) 则达到总和
\[
\mathcal C_6^{\mathrm{cont}}(T)=\sum_{x\in X_6}d_6(x)=2^6=64.
\]
最后令 \(T=2^B\) 即得 dyadic 预算下的分段闭式。由定理 \ref{thm:xi-fiber-spectrum-finite-moment-rigidity}，因 \(s_6=3\)，前六个幂和已足以唯一恢复整个有限尺度谱。证毕。
\end{proof}

\begin{theorem}[window-\texorpdfstring{$6$}{6} 的 oracle 几何完全可解]\label{thm:xi-window6-oracle-geometry-completely-solvable}
沿用命题 \ref{prop:xi-window6-complete-finite-scale-spectrum-closed-form} 的记号。则 window-\(6\) 的尾函数、连续容量曲线与全阶矩谱全部具有闭式
\[
N_6(t)=
\begin{cases}
21,& 0\le t\le 2,\\
13,& 2<t\le 3,\\
9,& 3<t\le 4,\\
0,& t>4,
\end{cases}
\]
\[
\mathcal C_6^{\mathrm{cont}}(T)=
\begin{cases}
21T,& 0\le T\le 2,\\
13T+16,& 2\le T\le 3,\\
9T+28,& 3\le T\le 4,\\
64,& T\ge 4,
\end{cases}
\]
\[
S_q(6)=8\cdot 2^q+4\cdot 3^q+9\cdot 4^q
\qquad(q>0).
\]
进一步，若写离散 dyadic 预言机容量与最优成功率为
\[
C_6(B):=\mathcal C_6^{\mathrm{cont}}(2^B),
\qquad
\mathrm{Succ}_6^{\mathrm{bin}}(B):=\frac{C_6(B)}{64},
\]
则有
\[
\boxed{\;
C_6(0)=21,\qquad
C_6(1)=42,\qquad
C_6(B)=64\ \ (B\ge 2),\;}
\]
\[
\boxed{\;
\mathrm{Succ}_6^{\mathrm{bin}}(0)=\frac{21}{64},\qquad
\mathrm{Succ}_6^{\mathrm{bin}}(1)=\frac{21}{32},\qquad
\mathrm{Succ}_6^{\mathrm{bin}}(B)=1\ \ (B\ge 2).\;}
\]
\end{theorem}
\begin{proof}
前三条闭式正是命题 \ref{prop:xi-window6-complete-finite-scale-spectrum-closed-form} 的内容。再由
\[
C_6(B)=\mathcal C_6^{\mathrm{cont}}(2^B),
\]
代入 \(B=0,1\) 分别得到
\[
C_6(0)=\mathcal C_6^{\mathrm{cont}}(1)=21,
\qquad
C_6(1)=\mathcal C_6^{\mathrm{cont}}(2)=42.
\]
而当 \(B\ge 2\) 时 \(2^B\ge 4\)，故由上式的末段可知
\[
C_6(B)=64.
\]
两边除以 \(2^6=64\) 即得成功率公式。证毕。
\end{proof}

\paragraph{预言机容量的 layer-cake 对偶与纤维尾谱完备性}
在 dyadic 预算下的预言机容量已经写成纤维多重度的截断和之后，还可把预算变量延拓为连续阈值，并将容量曲线、尾谱与 multiplicity 直方图之间的互译关系压缩为一条完全无损的 layer-cake 对偶。

\begin{theorem}[预言机容量与纤维尾谱的 layer-cake 对偶及完备性]\label{thm:xi-oracle-capacity-layercake-tail-completeness}
\leanverified{paper\_xi\_oracle\_capacity\_layercake\_tail\_completeness}
沿用定义 \ref{def:fold-capacity-tail} 与定理 \ref{thm:fold-oracle-capacity-closed-form} 的记号。对固定分辨率 \(m\)，继续记
\[
d_m(x)\qquad(x\in X_m)
\]
为折叠纤维多重度，并定义连续阈值容量
\[
\mathcal C_m^{\mathrm{cont}}(T)
:=
\sum_{x\in X_m}\min\bigl(d_m(x),T\bigr)
\qquad(T\ge 0).
\]
则 dyadic 预言机容量正是其 dyadic 采样：
\[
C_m(B)=\mathcal C_m^{\mathrm{cont}}(2^B)
\qquad(B\in\ZZ_{\ge 0}).
\]
若再记尾计数函数
\[
N_m(t):=\#\{x\in X_m:\ d_m(x)\ge t\}
\qquad(t\ge 0),
\]
则对一切 \(T\ge 0\) 都有精确 layer-cake 恒等式
\[
\mathcal C_m^{\mathrm{cont}}(T)=\int_0^T N_m(t)\,\dd t.
\]
并且对每个整数 \(k\ge 1\)，成立差分恢复公式
\[
N_m(k)=\mathcal C_m^{\mathrm{cont}}(k)-\mathcal C_m^{\mathrm{cont}}(k-1).
\]
若定义 multiplicity 直方图
\[
\mathsf M_m(k):=\#\{x\in X_m:\ d_m(x)=k\}
\qquad(k\ge 1),
\]
则
\[
\mathsf M_m(k)=N_m(k)-N_m(k+1).
\]
因此 \(\mathcal C_m^{\mathrm{cont}}\) 与 \(N_m\) 互相唯一决定，而 \(\mathsf M_m\) 可由二阶离散差分无损恢复。
\end{theorem}
\begin{proof}
由定理 \ref{thm:fold-oracle-capacity-closed-form}，
\[
C_m(B)=\sum_{x\in X_m}\min\bigl(d_m(x),2^B\bigr),
\]
故直接代入 \(T=2^B\) 即得
\[
C_m(B)=\mathcal C_m^{\mathrm{cont}}(2^B).
\]

对任意整数 \(d\ge 0\) 与任意 \(T\ge 0\)，都有单点恒等式
\[
\min(d,T)=\int_0^T \mathbf 1_{\{t\le d\}}\,\dd t.
\]
令 \(d=d_m(x)\) 并在 \(x\in X_m\) 上求和，交换有限求和与积分，便得到
\[
\begin{aligned}
\mathcal C_m^{\mathrm{cont}}(T)
&=
\sum_{x\in X_m}\int_0^T \mathbf 1_{\{t\le d_m(x)\}}\,\dd t\\
&=
\int_0^T \sum_{x\in X_m}\mathbf 1_{\{d_m(x)\ge t\}}\,\dd t\\
&=
\int_0^T N_m(t)\,\dd t.
\end{aligned}
\]
这正是容量曲线与尾谱之间的 layer-cake 对偶。

再取整数 \(k\ge 1\)。由上式与逐项差分，
\[
\begin{aligned}
\mathcal C_m^{\mathrm{cont}}(k)-\mathcal C_m^{\mathrm{cont}}(k-1)
&=
\sum_{x\in X_m}
\Bigl(\min\bigl(d_m(x),k\bigr)-\min\bigl(d_m(x),k-1\bigr)\Bigr)\\
&=
\sum_{x\in X_m}\mathbf 1_{\{d_m(x)\ge k\}}\\
&=
N_m(k).
\end{aligned}
\]
最后，
\[
\mathbf 1_{\{d_m(x)\ge k\}}-\mathbf 1_{\{d_m(x)\ge k+1\}}
=
\mathbf 1_{\{d_m(x)=k\}},
\]
故
\[
\mathsf M_m(k)
=
\sum_{x\in X_m}\mathbf 1_{\{d_m(x)=k\}}
=
N_m(k)-N_m(k+1).
\]
于是 \(\mathcal C_m^{\mathrm{cont}}\)、\(N_m\) 与 \(\mathsf M_m\) 的相互恢复全部完成。证毕。
\end{proof}

\begin{theorem}[全整数预算采样下的预言机成功率闭式与多重度直方图可逆性]\label{thm:xi-oracle-capacity-full-budget-discrete-inversion}
\leanverified{paper\_xi\_oracle\_capacity\_full\_budget\_discrete\_inversion}
沿用定理 \ref{thm:fold-oracle-capacity-closed-form} 与定理 \ref{thm:xi-oracle-capacity-layercake-tail-completeness} 的记号。对每个整数
\[
s\in\{0,1,\dots,2^m\},
\]
定义字母表大小至多为 \(s\) 的预言机最优成功率
\[
\mathrm{Succ}_m(s):=
\sup\Bigl\{\PP(\widehat\omega=\omega):\ |\mathcal A|\le s\Bigr\},
\]
其中 \(\mathcal A\) 表示侧信息字母表。则有精确闭式
\[
\boxed{\;
\mathrm{Succ}_m(s)
=
2^{-m}\sum_{x\in X_m}\min\bigl(d_m(x),s\bigr).\;}
\]
若再记整数预算容量
\[
C_m(s):=\sum_{x\in X_m}\min\bigl(d_m(x),s\bigr),
\qquad
\Delta_m(s):=C_m(s)-C_m(s-1)\quad(1\le s\le 2^m),
\]
并约定
\[
\Delta_m(2^m+1):=0,
\]
则成立显式反演公式
\[
\boxed{\;
\Delta_m(s)=\#\{x\in X_m:\ d_m(x)\ge s\},\;}
\]
\[
\boxed{\;
\#\{x\in X_m:\ d_m(x)=s\}
=
\Delta_m(s)-\Delta_m(s+1).\;}
\]
因此，整条整数预算曲线 \(s\mapsto \mathrm{Succ}_m(s)\) 与 \(s\mapsto C_m(s)\) 都唯一确定纤维多重度多重集 \(\{d_m(x):x\in X_m\}\)。
\end{theorem}
\begin{proof}
固定一条纤维 \(x\in X_m\)，并记其大小为 \(d=d_m(x)\)。若侧信息字母表大小至多为 \(s\)，则在该纤维内部至多只能给 \(s\) 个微观态分配彼此可区分的标签；故这条纤维对正确恢复的最优贡献恰为
\[
\min(d,s).
\]
对全部纤维求和并除以总微观态数 \(2^m\)，便得到
\[
\mathrm{Succ}_m(s)=2^{-m}\sum_{x\in X_m}\min(d_m(x),s).
\]

再看反演。由定义，
\[
C_m(s)=\sum_{x\in X_m}\min(d_m(x),s).
\]
当 \(s\) 增加 \(1\) 时，每个满足 \(d_m(x)\ge s\) 的项恰增加 \(1\)，其余项保持不变，因此
\[
\Delta_m(s)
=
C_m(s)-C_m(s-1)
=
\sum_{x\in X_m}\mathbf 1_{\{d_m(x)\ge s\}}
=
\#\{x\in X_m:\ d_m(x)\ge s\}.
\]
再作一次相邻差分，便得到
\[
\#\{x\in X_m:\ d_m(x)=s\}
=
\#\{d_m\ge s\}-\#\{d_m\ge s+1\}
=
\Delta_m(s)-\Delta_m(s+1).
\]
故整条整数预算曲线唯一恢复纤维多重度直方图。证毕。
\end{proof}

\begin{theorem}[精确反演阈值等于峰值纤维指数]\label{thm:xi-exact-inversion-threshold-peak-index}
沿用定理 \ref{thm:xi-oracle-capacity-full-budget-discrete-inversion} 的记号，并定义
\[
s_{\mathrm{inv}}(m)
:=
\min\{s\ge 1:\ \mathrm{Succ}_m(s)=1\},
\qquad
\mathrm{Succ}^{\mathrm{bin}}_m(B):=\mathrm{Succ}_m(2^B),
\]
\[
B_{\mathrm{inv}}(m)
:=
\min\{B\ge 0:\ \mathrm{Succ}^{\mathrm{bin}}_m(B)=1\}.
\]
则有闭式
\[
\boxed{\;
s_{\mathrm{inv}}(m)=\max_{x\in X_m}d_m(x),\qquad
B_{\mathrm{inv}}(m)=\Bigl\lceil\log_2\max_{x\in X_m}d_m(x)\Bigr\rceil.\;}
\]
在交换情形下，若 \(E_m\) 为折叠 \(\Fold_m\) 诱导的纤维平均条件期望，则由 Watatani 指数元的闭式（定理 \ref{thm:op-algebra-fold-watatani-index-equals-multiplicity-field}）进一步得到
\[
\boxed{\;
B_{\mathrm{inv}}(m)
=
\Bigl\lceil
\log_2 \operatorname*{ess\,sup}_{x}\Ind(E_m)(x)
\Bigr\rceil.\;}
\]
更一般地，若每步可读字母表大小为 \(|\mathfrak A_0|\)，则最小精确反演时间预算满足
\[
\boxed{\;
T_{\mathrm{inv}}(m)
=
\Bigl\lceil
\log_{|\mathfrak A_0|}
\operatorname*{ess\,sup}_{x}\Ind(E_m)(x)
\Bigr\rceil.\;}
\]
\end{theorem}
\begin{proof}
由定理 \ref{thm:xi-oracle-capacity-full-budget-discrete-inversion}，
\[
\mathrm{Succ}_m(s)=1
\quad\Longleftrightarrow\quad
\sum_{x\in X_m}\min\bigl(d_m(x),s\bigr)=\sum_{x\in X_m}d_m(x).
\]
右端成立当且仅当对每个 \(x\in X_m\) 都有 \(\min(d_m(x),s)=d_m(x)\)，亦即
\[
s\ge d_m(x)\qquad(\forall x\in X_m).
\]
故
\[
s_{\mathrm{inv}}(m)=\max_{x\in X_m}d_m(x),
\]
从而 dyadic 预算下
\[
B_{\mathrm{inv}}(m)
=
\min\{B\ge 0:\ 2^B\ge \max_x d_m(x)\}
=
\Bigl\lceil\log_2\max_x d_m(x)\Bigr\rceil.
\]
交换情形下，定理 \ref{thm:op-algebra-fold-watatani-index-equals-multiplicity-field} 给出
\(
\Ind(E_m)(x)=d_m(x)
\)
（a.e.），故第二式成立。
若每步可读字母表大小为 \(|\mathfrak A_0|\)，则 \(T\) 步至多提供 \(|\mathfrak A_0|^T\) 个可区分标签；将上式中的 \(s\) 替换为 \(|\mathfrak A_0|^T\) 即得
\[
T_{\mathrm{inv}}(m)
=
\min\{T\ge 0:\ |\mathfrak A_0|^T\ge \operatorname*{ess\,sup}_x \Ind(E_m)(x)\},
\]
亦即所述闭式。证毕。
\end{proof}

\begin{corollary}[平均缺陷预算与峰值指数预算的分层]\label{cor:xi-average-defect-vs-peak-index-budget}
固定分辨率 \(m\) 与宏观边缘 \(\pi\in\Delta(X_m)\)。
若记录轴 \(r\) 使扩展读出 \((\Fold_m,r)\) 为单射，并令 \(A\sim\pi^{\mathrm e}\)、\(X:=\Fold_m(A)\sim\pi\)，则由定理 \ref{thm:address-defect-law} 有平均预算下界
\[
H(r(A)\mid X)\ \ge\ \kappa_m(\pi)
=
\mathbb E_{\pi}\bigl[\log d_m(X)\bigr].
\]
若进一步要求在全部纤维上一致地实现精确反演，则由定理 \ref{thm:xi-exact-inversion-threshold-peak-index}，任意 \(|\mathfrak A_0|\)-进实现都必须满足
\[
T\ \ge\
\log_{|\mathfrak A_0|}
\operatorname*{ess\,sup}_{x}\Ind(E_m)(x).
\]
因此，平均型辅助资源由 \(\kappa_m(\pi)\) 控制，而一致精确可逆性由峰值指数
\(
\operatorname*{ess\,sup}_{x}\Ind(E_m)(x)
\)
控制；二者一般不能塌缩为同一标尺。
更具体地，
\[
\kappa_m(\pi)
\le
\log \operatorname*{ess\,sup}_{x}\Ind(E_m)(x),
\]
且只有当 \(d_m(x)\) 在 \(\pi\)-几乎处处的支撑上本质常值时，上述两个量才可能取同一数值。
特别地，window-\(6\) 的纤维谱 \(\{2,3,4\}\) 排除了把平均预算与峰值预算压成单一数的可能性。
\end{corollary}
\begin{proof}
第一条不等式是定理 \ref{thm:address-defect-law} 第 (2) 条的直接重述。
第二条不等式由定理 \ref{thm:xi-exact-inversion-threshold-peak-index} 直接给出。
最后，由
\[
\log d_m(X)\le \log \operatorname*{ess\,sup}_{x}\Ind(E_m)(x)
\qquad\text{\(\pi\)-a.e.}
\]
取 \(\pi\)-期望即得
\(
\kappa_m(\pi)\le \log \operatorname*{ess\,sup}_{x}\Ind(E_m)(x)
\)。
只有当 \(\log d_m(X)\) 在 \(\pi\)-几乎处处取等时，期望才可取等；这等价于 \(d_m(x)\) 在相关支撑上本质常值。证毕。
\end{proof}

\paragraph{连续容量的完备统计量、分布导数与 window-\texorpdfstring{$6$}{6} 异常初始斜率}
在前述有限矩刚性、Mellin--Stieltjes 反演与 window-\(6\) 的有限尺度闭式之后，还可把连续容量曲线、尾计数函数、纤维直方图、纤维经验测度与全矩谱收束为同一份完全可逆的数据。

\begin{theorem}[连续容量曲线是纤维直方图与全矩谱的完备统计量]\label{thm:xi-capacity-tail-histogram-moment-complete-statistic}
固定分辨率 \(m\)。记
\[
h_m(d):=\#\{x\in X_m:\ d_m(x)=d\}
\qquad(d\in \ZZ_{\ge 1}),
\]
\[
N_m(t):=\#\{x\in X_m:\ d_m(x)\ge t\},
\qquad
\mathcal C_m^{\mathrm{cont}}(T):=\sum_{x\in X_m}\min\bigl(d_m(x),T\bigr).
\]
则 \(\mathcal C_m^{\mathrm{cont}}\)、\(N_m\) 与 \(h_m\) 两两唯一决定。更具体地，
\[
\mathcal C_m^{\mathrm{cont}}(T)
=
\int_0^T N_m(t)\,\dd t
=
\sum_{d\ge 1}h_m(d)\min(d,T),
\]
并且
\[
\partial_T^+\mathcal C_m^{\mathrm{cont}}(t)=N_m(t)
\qquad\text{对 Lebesgue--a.e.\ }t>0\text{ 成立},
\]
从而
\[
h_m(d)=N_m(d)-N_m(d+1)
\qquad(d\ge 1),
\]
\[
N_m(t)=\sum_{d\ge \lceil t\rceil}h_m(d)
\qquad(t>0).
\]
此外，对任意 \(q>0\) 都有 Mellin--Stieltjes 公式
\[
S_q(m):=\sum_{x\in X_m}d_m(x)^q
=
q\int_0^\infty t^{q-1}N_m(t)\,\dd t.
\]
因此在固定 \(m\) 上，下列四类对象完全等价：
\[
\mathcal C_m^{\mathrm{cont}},
\qquad
N_m,
\qquad
h_m,
\qquad
\{S_q(m)\}_{q>0}.
\]
\end{theorem}
\begin{proof}
由命题 \ref{prop:fold-capacity-tail-equivalence}，
\[
\mathcal C_m^{\mathrm{cont}}(T)=\int_0^T N_m(t)\,\dd t.
\]
又因全部纤维大小 \(d_m(x)\) 都是正整数，\(N_m\) 是整数值右连续阶梯函数，其跳点恰为出现过的纤维尺寸。故
\[
\partial_T^+\mathcal C_m^{\mathrm{cont}}(t)=N_m(t)
\qquad\text{对 a.e.\ }t>0\text{ 成立},
\]
并且由右连续性，\(N_m\) 由该 a.e.\ 导数唯一延拓回整个函数。

若按纤维大小分层，则
\[
N_m(t)
=
\sum_{d\ge 1}\#\{x\in X_m:\ d_m(x)=d,\ d\ge t\}
=
\sum_{d\ge \lceil t\rceil}h_m(d),
\]
于是
\[
h_m(d)=N_m(d)-N_m(d+1).
\]
反向地，把同一分层代回容量定义即得
\[
\mathcal C_m^{\mathrm{cont}}(T)=\sum_{d\ge 1}h_m(d)\min(d,T).
\]
故 \(\mathcal C_m^{\mathrm{cont}}\)、\(N_m\) 与 \(h_m\) 彼此唯一决定。

再由定理 \ref{thm:xi-time-part57b-capacity-moment-exact-mellin-stieltjes}，
\[
S_q(m)=q\int_0^\infty t^{q-1}N_m(t)\,\dd t
\qquad(q>0),
\]
故任一对象 \(\mathcal C_m^{\mathrm{cont}}\)、\(N_m\)、\(h_m\) 都唯一决定全部正阶矩谱。

最后说明矩谱反向决定直方图。设不同纤维大小为
\[
1\le \delta_{1,m}<\cdots<\delta_{s_m,m},
\]
其重数为 \(n_{j,m}\)。由于每个 \(\delta_{j,m}\ge 1\)，有
\[
\lim_{q\downarrow 0}S_q(m)
=
\sum_{j=1}^{s_m}n_{j,m}
=
|X_m|
=
S_0(m).
\]
因此从全谱 \(\{S_q(m)\}_{q>0}\) 可恢复 \(S_0(m)\)，而整数矩
\[
S_1(m),\dots,S_{2s_m-1}(m)
\]
本就包含在该全谱中。于是定理 \ref{thm:xi-fiber-spectrum-finite-moment-rigidity} 适用，并唯一恢复全部 \(\{(\delta_{j,m},n_{j,m})\}_{j=1}^{s_m}\)，从而唯一恢复 \(h_m\)。因此四类对象完全等价。证毕。
\end{proof}

\begin{theorem}[连续容量曲线等于推前分布与均匀层测度的截断重叠泛函]\label{thm:xi-capacity-overlap-functional-l1-identity}
\leanverified{paper\_xi\_capacity\_overlap\_functional\_l1\_identity}
沿用定理 \ref{thm:xi-capacity-tail-histogram-moment-complete-statistic} 的记号，并定义
\[
\mu_m(x):=\frac{d_m(x)}{2^m},
\qquad
u_m(x):=\frac{1}{|X_m|}
\qquad (x\in X_m),
\]
\[
\tau_m(T):=\frac{T|X_m|}{2^m}
\qquad (T\ge 0),
\]
以及
\[
\|f\|_{\ell^1(X_m)}:=\sum_{x\in X_m}|f(x)|.
\]
则对任意 \(T\ge 0\) 都有严格恒等式
\[
\boxed{\;
\frac{\mathcal C_m^{\mathrm{cont}}(T)}{2^m}
=
\sum_{x\in X_m}\min\bigl(\mu_m(x),\tau_m(T)u_m(x)\bigr)
=
\frac{1+\tau_m(T)-\|\mu_m-\tau_m(T)u_m\|_{\ell^1(X_m)}}{2}.\;}
\]
特别地，若对 bin-fold 写
\[
C_m^{\mathrm{bin}}(B):=\mathcal C_m^{\mathrm{cont}}(2^B),
\qquad
\mathrm{Succ}_m^{\mathrm{bin}}(B):=\frac{C_m^{\mathrm{bin}}(B)}{2^m},
\]
则对任意整数 \(B\ge 0\) 都有
\[
\boxed{\;
\mathrm{Succ}_m^{\mathrm{bin}}(B)
=
\frac{1+\tau_m(2^B)-\|\mu_m-\tau_m(2^B)u_m\|_{\ell^1(X_m)}}{2}.\;}
\]
\end{theorem}
\begin{proof}
由连续容量曲线的定义，
\[
\frac{\mathcal C_m^{\mathrm{cont}}(T)}{2^m}
=
\sum_{x\in X_m}\min\!\left(\frac{d_m(x)}{2^m},\frac{T}{2^m}\right)
=
\sum_{x\in X_m}\min\bigl(\mu_m(x),\tau_m(T)u_m(x)\bigr),
\]
因为对每个 \(x\in X_m\) 都有
\[
\tau_m(T)u_m(x)=\frac{T|X_m|}{2^m}\cdot\frac{1}{|X_m|}=\frac{T}{2^m}.
\]
再对每个 \(x\) 应用标量恒等式
\[
\min(a,b)=\frac{a+b-|a-b|}{2}
\qquad(a,b\ge 0),
\]
并利用
\[
\sum_{x\in X_m}\mu_m(x)=1,
\qquad
\sum_{x\in X_m}u_m(x)=1,
\]
便得到
\[
\sum_{x\in X_m}\min\bigl(\mu_m(x),\tau_m(T)u_m(x)\bigr)
=
\frac{1+\tau_m(T)-\|\mu_m-\tau_m(T)u_m\|_{\ell^1(X_m)}}{2}.
\]
这就证明了第一条恒等式。取 \(T=2^B\) 并用
\(
\mathrm{Succ}_m^{\mathrm{bin}}(B)=2^{-m}\mathcal C_m^{\mathrm{cont}}(2^B)
\)
即可得到 dyadic 预算下的成功率公式。证毕。
\end{proof}

\begin{corollary}[有限支持纤维谱的 Hankel--Vandermonde 显式重构公式]\label{cor:xi-fiber-spectrum-hankel-vandermonde-explicit-recovery}
\leanverified{paper\_xi\_fiber\_spectrum\_hankel\_vandermonde\_explicit\_recovery}
固定分辨率 \(m\)。设纤维多重度只取 \(R\) 个不同值
\[
1\le a_1<\cdots<a_R,
\]
并记其重数为
\[
n_j:=\#\{x\in X_m:\ d_m(x)=a_j\}
\qquad(1\le j\le R).
\]
则对一切整数 \(q\ge 0\) 都有
\[
S_q(m)=\sum_{j=1}^{R}n_j a_j^q.
\]
若记消去多项式
\[
P_m(X):=\prod_{j=1}^{R}(X-a_j)
=
X^R+b_{R-1}X^{R-1}+\cdots+b_0,
\]
则系数向量 \((b_0,\dots,b_{R-1})\) 是 Hankel 线性系统
\[
\sum_{\ell=0}^{R-1}b_\ell\,S_{\ell+i}(m)=-S_{R+i}(m)
\qquad(0\le i\le R-1)
\]
的唯一解。求得根 \(\{a_j\}_{j=1}^{R}\) 后，重数向量 \((n_1,\dots,n_R)\) 再由 Vandermonde 系统
\[
S_q(m)=\sum_{j=1}^{R}n_j a_j^q
\qquad(0\le q\le R-1)
\]
唯一恢复。

因此，仅凭
\[
S_0(m),S_1(m),\dots,S_{2R-1}(m)
\]
便已足以唯一决定完整纤维直方图。
\end{corollary}
\begin{proof}
这正是定理 \ref{thm:xi-fiber-spectrum-finite-moment-rigidity} 的坐标化重写；其中
\[
R=s_m,\qquad a_j=\delta_{j,m},\qquad n_j=n_{j,m}.
\]
为便于直接调用，这里把相应线性系统写成显式形式。

由 \(P_m(a_j)=0\) 对每个 \(1\le j\le R\) 成立，故对任意 \(i\ge 0\) 有
\[
0=\sum_{j=1}^{R}n_j a_j^i P_m(a_j)
=
\sum_{j=1}^{R}n_j a_j^i
\Bigl(a_j^R+\sum_{\ell=0}^{R-1}b_\ell a_j^\ell\Bigr).
\]
按 \(S_q(m)=\sum_j n_j a_j^q\) 整理即得
\[
S_{R+i}(m)+\sum_{\ell=0}^{R-1}b_\ell S_{\ell+i}(m)=0
\qquad(i\ge 0),
\]
特别地，对 \(0\le i\le R-1\) 就得到所示 Hankel 线性系统。

现证其解唯一。记
\[
H_m:=\bigl(S_{r+s}(m)\bigr)_{0\le r,s\le R-1}.
\]
对任意向量 \(u=(u_0,\dots,u_{R-1})^{\mathsf T}\)，有
\[
u^{\mathsf T}H_m u
=
\sum_{j=1}^{R}n_j\Bigl(\sum_{\ell=0}^{R-1}u_\ell a_j^\ell\Bigr)^2.
\]
若 \(u^{\mathsf T}H_m u=0\)，则右端各项皆为零，于是多项式
\[
Q(X):=\sum_{\ell=0}^{R-1}u_\ell X^\ell
\]
在 \(R\) 个互异点 \(a_1,\dots,a_R\) 上全为零。由于 \(\deg Q\le R-1\)，只能 \(Q\equiv 0\)，故 \(u=0\)。因此 \(H_m\) 正定，特别可逆，所以上述 Hankel 系统有唯一解。

得到 \(P_m\) 后，其根集即为全部不同纤维大小 \(\{a_j\}_{j=1}^{R}\)。再把这些根代回
\[
S_q(m)=\sum_{j=1}^{R}n_j a_j^q
\qquad(0\le q\le R-1),
\]
便得到一个 Vandermonde 线性系统。由于 \(a_j\) 两两不同，其系数矩阵可逆，从而 \((n_1,\dots,n_R)\) 唯一确定。证毕。
\end{proof}

\begin{theorem}[连续容量曲线的分布二阶导数恰为纤维多重度经验测度]\label{thm:xi-capacity-second-distribution-derivative-fiber-empirical-measure}
沿用定理 \ref{thm:xi-capacity-tail-histogram-moment-complete-statistic} 的记号，并定义纤维多重度经验测度
\[
\nu_m:=\sum_{x\in X_m}\delta_{d_m(x)}
=
\sum_{d\ge 1}h_m(d)\,\delta_d.
\]
则在 \((0,\infty)\) 上按分布成立精确恒等式
\[
\boxed{\;
-D_T^2\mathcal C_m^{\mathrm{cont}}=\nu_m.\;}
\]
等价地，\(\mathcal C_m^{\mathrm{cont}}\) 的折点支撑恰为全部出现过的纤维大小，而对任意整数 \(d\ge 1\) 都有
\[
\boxed{\;
h_m(d)
=
\nu_m(\{d\})
=
N_m(d)-N_m(d+0)
=
\partial_T^{-}\mathcal C_m^{\mathrm{cont}}(d)-\partial_T^{+}\mathcal C_m^{\mathrm{cont}}(d).\;}
\]
因此，连续容量曲线不仅决定尾函数与全矩谱，而且已经完整且无损地决定了纤维多重度多重集本身。
\end{theorem}
\begin{proof}
在 \((0,\infty)\) 上，
\[
N_m(T)=\sum_{x\in X_m}\mathbf 1_{\{T\le d_m(x)\}}
\]
是右连续阶梯函数。对任意固定 \(d>0\)，分布导数满足
\[
D_T\mathbf 1_{\{T\le d\}}=-\delta_d.
\]
故在分布意义下
\[
D_TN_m
=
\sum_{x\in X_m}D_T\mathbf 1_{\{T\le d_m(x)\}}
=
-\sum_{x\in X_m}\delta_{d_m(x)}
=
-\nu_m.
\]

另一方面，由定理 \ref{thm:xi-capacity-tail-histogram-moment-complete-statistic}，
\[
\mathcal C_m^{\mathrm{cont}}(T)=\int_0^T N_m(t)\,\dd t,
\]
从而
\[
D_T\mathcal C_m^{\mathrm{cont}}=N_m,
\qquad
D_T^2\mathcal C_m^{\mathrm{cont}}=D_TN_m=-\nu_m
\]
按分布成立，即得第一式。

又因 \(N_m\) 在 \(T=d\) 处的向下跳跃恰等于
\[
N_m(d)-N_m(d+0)=\#\{x\in X_m:\ d_m(x)=d\}=h_m(d),
\]
而 \(\mathcal C_m^{\mathrm{cont}}\) 在该点的左右导数分别为
\[
\partial_T^{-}\mathcal C_m^{\mathrm{cont}}(d)=N_m(d),
\qquad
\partial_T^{+}\mathcal C_m^{\mathrm{cont}}(d)=N_m(d+0),
\]
遂得第二式。证毕。
\end{proof}

\begin{theorem}[矩谱对尾函数与连续容量曲线具有严格的 Mellin 反演公式]\label{thm:xi-capacity-moment-mellin-inversion-tail-capacity}
沿用定理 \ref{thm:xi-capacity-tail-histogram-moment-complete-statistic} 的记号。对复数 \(s\in\CC\) 定义
\[
S_s(m):=\sum_{x\in X_m}d_m(x)^s
=
\sum_{d\ge 1}h_m(d)\,d^s,
\]
其中 \(d^s:=e^{s\log d}\) 取实对数支。则当 \(\Re s>0\) 时有
\[
\boxed{\;
\frac{S_s(m)}{s}
=
\int_0^\infty T^{s-1}N_m(T)\,\dd T.\;}
\]
因此，对任意 \(c>0\) 与任意不是纤维大小跳点的 \(T>0\)，有 Mellin 反演公式
\[
\boxed{\;
N_m(T)
=
\frac{1}{2\pi i}
\int_{c-i\infty}^{c+i\infty}
\frac{S_s(m)}{s}\,T^{-s}\,\dd s.\;}
\]
进一步，对任意 \(c>1\) 与任意 \(T>0\)，有
\[
\boxed{\;
2^m-\mathcal C_m^{\mathrm{cont}}(T)
=
\frac{1}{2\pi i}
\int_{c-i\infty}^{c+i\infty}
\frac{S_s(m)}{s(s-1)}\,T^{1-s}\,\dd s.\;}
\]
等价地，
\[
\boxed{\;
\mathcal C_m^{\mathrm{cont}}(T)
=
2^m+
\frac{1}{2\pi i}
\int_{c-i\infty}^{c+i\infty}
\frac{S_s(m)}{s(1-s)}\,T^{1-s}\,\dd s.\;}
\]
因此，全体矩谱 \(\{S_s(m)\}_{\Re s>0}\) 不仅与尾函数等价，而且可通过显式解析核无损反演出连续容量曲线。
\end{theorem}
\begin{proof}
由定理 \ref{thm:xi-capacity-tail-histogram-moment-complete-statistic} 的 Mellin--Stieltjes 公式，对 \(\Re s>0\) 有
\[
S_s(m)=s\int_0^\infty T^{s-1}N_m(T)\,\dd T.
\]
这就得到第一式。由于 \(N_m\) 有界、右连续且紧支撑，标准 Mellin 反演定理适用，从而在 \(N_m\) 的连续点 \(T\) 上，
\[
N_m(T)
=
\frac{1}{2\pi i}
\int_{c-i\infty}^{c+i\infty}
\frac{S_s(m)}{s}\,T^{-s}\,\dd s
\qquad(c>0).
\]

另一方面，由
\[
\sum_{x\in X_m}d_m(x)=2^m
\]
与定理 \ref{thm:xi-capacity-tail-histogram-moment-complete-statistic} 中的
\[
\mathcal C_m^{\mathrm{cont}}(T)=\int_0^T N_m(t)\,\dd t
\]
可得
\[
2^m-\mathcal C_m^{\mathrm{cont}}(T)=\int_T^\infty N_m(t)\,\dd t.
\]
将上式中的 \(N_m(t)\) 代入其 Mellin 反演公式，并在 \(c>1\) 时交换积分次序，得到
\[
\begin{aligned}
2^m-\mathcal C_m^{\mathrm{cont}}(T)
&=
\frac{1}{2\pi i}
\int_{c-i\infty}^{c+i\infty}
\frac{S_s(m)}{s}
\left(\int_T^\infty t^{-s}\,\dd t\right)\dd s\\
&=
\frac{1}{2\pi i}
\int_{c-i\infty}^{c+i\infty}
\frac{S_s(m)}{s(s-1)}\,T^{1-s}\,\dd s,
\end{aligned}
\]
因为当 \(\Re s>1\) 时，
\[
\int_T^\infty t^{-s}\,\dd t=\frac{T^{1-s}}{s-1}.
\]
这便得到第三式；最后一式只需利用
\[
\frac{1}{s(1-s)}=-\frac{1}{s(s-1)}
\]
改写即可。证毕。
\end{proof}

\begin{theorem}[连续容量曲线与全阶矩谱之间的 Mellin--Barnes 完备互逆]\label{thm:xi-capacity-moment-mellin-barnes-complete-inversion}
沿用定理 \ref{thm:xi-capacity-moment-mellin-inversion-tail-capacity} 的记号。由于 \(X_m\) 有限，
\[
S_s(m):=\sum_{x\in X_m}d_m(x)^s
\]
在整个复平面上全纯。对任意满足 \(0<\Re s<1\) 的复数 \(s\)，有
\[
\boxed{\;
S_s(m)
=
s(1-s)\int_0^\infty T^{s-2}\mathcal C_m^{\mathrm{cont}}(T)\,\dd T.\;}
\]
反之，对任意 \(c\in(0,1)\) 与任意 \(T>0\)，都有 Mellin--Barnes 反演公式
\[
\boxed{\;
\mathcal C_m^{\mathrm{cont}}(T)
=
\frac{1}{2\pi i}
\int_{c-i\infty}^{c+i\infty}
\frac{S_s(m)}{s(1-s)}\,T^{1-s}\,\dd s.\;}
\]
因此，固定尺度上的连续容量曲线与全阶矩谱在 Mellin--Barnes 变换下完备可逆。
\end{theorem}
\begin{proof}
先说明积分收敛性。由
\[
\mathcal C_m^{\mathrm{cont}}(T)=\sum_{x\in X_m}\min\bigl(d_m(x),T\bigr),
\]
当 \(T\downarrow 0\) 时有
\(
\mathcal C_m^{\mathrm{cont}}(T)=|X_m|\,T
\)，而当 \(T\ge \max_x d_m(x)\) 时有
\(
\mathcal C_m^{\mathrm{cont}}(T)=2^m
\)。
因此
\(
T^{s-2}\mathcal C_m^{\mathrm{cont}}(T)
\)
在 \(0\) 附近如 \(T^{s-1}\)，在 \(\infty\) 附近如 \(T^{s-2}\)，故当 \(0<\Re s<1\) 时绝对可积。

再由定理 \ref{thm:xi-capacity-tail-histogram-moment-complete-statistic}，
\[
\mathcal C_m^{\mathrm{cont}}(T)=\int_0^T N_m(t)\,\dd t.
\]
于是对 \(0<\Re s<1\)，可用 Fubini 定理交换积分次序，得到
\[
\begin{aligned}
\int_0^\infty T^{s-2}\mathcal C_m^{\mathrm{cont}}(T)\,\dd T
&=
\int_0^\infty T^{s-2}\int_0^T N_m(t)\,\dd t\,\dd T\\
&=
\int_0^\infty N_m(t)\left(\int_t^\infty T^{s-2}\,\dd T\right)\dd t\\
&=
\frac{1}{1-s}\int_0^\infty t^{s-1}N_m(t)\,\dd t.
\end{aligned}
\]
而定理 \ref{thm:xi-capacity-moment-mellin-inversion-tail-capacity} 已给出
\[
\frac{S_s(m)}{s}
=
\int_0^\infty t^{s-1}N_m(t)\,\dd t
\qquad(\Re s>0),
\]
故立得
\[
\int_0^\infty T^{s-2}\mathcal C_m^{\mathrm{cont}}(T)\,\dd T
=
\frac{S_s(m)}{s(1-s)},
\]
也即第一条公式。

最后，对函数
\(
T\mapsto \mathcal C_m^{\mathrm{cont}}(T)
\)
应用标准 Mellin 反演。由于它连续、分段线性、在 \(0\) 附近至多线性增长且在无穷远处有界，其 Mellin 变换
\[
\int_0^\infty T^{s-2}\mathcal C_m^{\mathrm{cont}}(T)\,\dd T
\]
恰在条带 \(0<\Re s<1\) 上全纯并满足反演条件。因此对任意 \(c\in(0,1)\) 与任意 \(T>0\)，
\[
\mathcal C_m^{\mathrm{cont}}(T)
=
\frac{1}{2\pi i}
\int_{c-i\infty}^{c+i\infty}
\left(\int_0^\infty u^{s-2}\mathcal C_m^{\mathrm{cont}}(u)\,\dd u\right)
T^{1-s}\,\dd s.
\]
再把第一条公式代回，即得所示 Mellin--Barnes 反演式。证毕。
\end{proof}

\begin{theorem}[window-\texorpdfstring{$6$}{6} 异常签名维数等于容量曲线在 \texorpdfstring{$T=1$}{T=1} 的初始右导数]\label{thm:xi-window6-anomaly-rank-initial-capacity-right-derivative}
沿用定理 \ref{thm:xi-foldbin-audited-even-center-collapse-abel-full-rank}、
定理 \ref{thm:xi-capacity-tail-histogram-moment-complete-statistic}
与命题 \ref{prop:xi-window6-complete-finite-scale-spectrum-closed-form} 的记号，则
\[
\dim_{\FF_2}\bigl(G_6^{\mathrm{bin}}\bigr)^{\mathrm{ab}}
=
N_6(1^+)
=
\partial_T^+\mathcal C_6^{\mathrm{cont}}(1)
=
21.
\]
\end{theorem}
\begin{proof}
由定理 \ref{thm:xi-foldbin-audited-even-center-collapse-abel-full-rank}，
\[
\bigl(G_6^{\mathrm{bin}}\bigr)^{\mathrm{ab}}\cong (\ZZ/2\ZZ)^{|X_6|},
\]
且 \(|X_6|=21\)，故
\[
\dim_{\FF_2}\bigl(G_6^{\mathrm{bin}}\bigr)^{\mathrm{ab}}=21.
\]

另一方面，命题 \ref{prop:xi-window6-complete-finite-scale-spectrum-closed-form} 已给出
\[
N_6(t)=21
\qquad(0\le t\le 2).
\]
因此
\[
N_6(1^+)=21.
\]
再由定理 \ref{thm:xi-capacity-tail-histogram-moment-complete-statistic}，
\[
\partial_T^+\mathcal C_6^{\mathrm{cont}}(1)=N_6(1^+)=21.
\]
合并即得所述恒等式。证毕。
\end{proof}

\paragraph{window-\texorpdfstring{$6$}{6} 上精确反演预算与异常审计维数的分离}
在 window-\(6\) 的容量阈值与 gauge 结构都已闭式化之后，还可把“恢复微观态”与“审计异常”两种资源明确区分为两条不同的量纲。

\begin{proposition}[window-\texorpdfstring{$6$}{6} 上精确可逆预算与异常审计维数的正交分离]\label{prop:xi-window6-exact-inversion-vs-anomaly-audit}
\leanverified{paper\_xi\_window6\_exact\_inversion\_vs\_anomaly\_audit}
在 rigid window-\(6\) 模型中，精确反演所需的最小 dyadic 预算为
\[
\boxed{\;
B_{\mathrm{inv}}^{(6)}=2.\;}
\]
另一方面，若异常审计签名要求同时满足：
\begin{enumerate}
  \item 保留各折叠纤维因子的独立置换语义；
  \item 以二元寄存器记录每个纤维因子的奇偶异常；
  \item 不在不同纤维因子之间引入额外关系，
\end{enumerate}
则其最小完备二元签名维数恰为
\[
\boxed{\;
B_{\mathrm{anom}}^{(6)}=21.\;}
\]
更具体地，
\[
\bigl(\mathcal G^{\mathrm{bin}}_6\bigr)^{\mathrm{ab}}\cong(\ZZ_2)^{21},
\]
故任何保持上述直积语义的二元异常审计器都至少需要 \(21\) 个独立比特，而 \(21\) 个纤维符号坐标又确实足以实现完备审计。
因此，window-\(6\) 中“精确恢复微观态”与“完整审计异常”由两种不同不变量控制：前者只看最大纤维，后者则看纤维因子数目的 Abel 化维数；二者是正交而不可互换的资源。
\end{proposition}
\begin{proof}
由前述 window-\(6\) 容量闭式，
\[
C_6(0)=21,\qquad C_6(1)=42,\qquad C_6(B)=64\quad(B\ge 2),
\]
故两比特预算已足以达到 \(\mathrm{Succ}^{\mathrm{bin}}_6(B)=1\)，即 \(B_{\mathrm{inv}}^{(6)}=2\)；这也与定理 \ref{thm:xi-exact-inversion-threshold-peak-index} 在纤维谱 \(\{2,3,4\}\) 上给出的
\(
\lceil \log_2 4\rceil=2
\)
一致。

另一方面，推论 \ref{cor:fold-bin-gauge-m6} 给出
\[
\mathcal G^{\mathrm{bin}}_6\cong S_2^{\,8}\times S_3^{\,4}\times S_4^{\,9},
\qquad
\bigl(\mathcal G^{\mathrm{bin}}_6\bigr)^{\mathrm{ab}}\cong(\ZZ_2)^{21}.
\]
若异常审计器以二元寄存器记录各纤维因子的奇偶异常，并保持上述直积语义，则其输出必因子化为某个
\(
\bigl(\mathcal G^{\mathrm{bin}}_6\bigr)^{\mathrm{ab}}\to(\ZZ_2)^b
\)
的群同态。
要使该审计在 Abel 化层面完备，必须在 \((\ZZ_2)^{21}\) 上忠实，故必有 \(b\ge 21\)；反过来，\(21\) 个纤维符号坐标本身就给出一个完备的二元审计器。
因此 \(B_{\mathrm{anom}}^{(6)}=21\)。

最后，\(B_{\mathrm{inv}}^{(6)}\) 由最大纤维大小决定，而 \(B_{\mathrm{anom}}^{(6)}\) 由 Abel 化秩决定，二者所依赖的不变量不同，故不能压缩为同一预算轴。证毕。
\end{proof}

\begin{corollary}[window-\texorpdfstring{$6$}{6} 上局域精确反演与异常完备审计之间的 \texorpdfstring{$19$}{19} bit 通胀刚性]\label{cor:xi-window6-nineteen-bit-inflation-rigidity}
沿用命题 \ref{prop:xi-window6-exact-inversion-vs-anomaly-audit} 的记号，则有精确恒等式
\[
\boxed{\;
B_{\mathrm{anom}}^{(6)}-B_{\mathrm{inv}}^{(6)}=19,\qquad
\frac{2^{B_{\mathrm{anom}}^{(6)}}}{2^{B_{\mathrm{inv}}^{(6)}}}=2^{19}.\;}
\]
因此，任何既要保持 window-\(6\) 局域精确反演、又要保持全部纤维独立异常签名的二元证书机制，都必须承受不可消去的 \(19\)-bit 通胀。
\end{corollary}
\begin{proof}
命题 \ref{prop:xi-window6-exact-inversion-vs-anomaly-audit} 已给出
\[
B_{\mathrm{inv}}^{(6)}=2,
\qquad
B_{\mathrm{anom}}^{(6)}=21.
\]
两式相减即得第一条恒等式；再取指数便得第二条恒等式。最后一句只是把该数值差额读作证书维数的不可压缩差距。证毕。
\end{proof}
\leanverified{paper\_xi\_window6\_nineteen\_bit\_inflation\_rigidity}

\paragraph{bin-fold 终位扇区的序理想分层、临界容量相图与黄金共振的不可达常数隙}
在尾部 digit-DP 计数、连续容量曲线的 layer-cake 对偶、两态一致渐近以及 Fibonacci 共振频率上的 Fourier 刚性都已建立之后，bin-fold 的纤维谱还可进一步压缩为一条由“终位扇区内的序理想几何”所驱动的主链。其一，固定末位后的全部纤维层级不再是任意散乱的，而是严格沿 \(V_m\)-次序形成初始段；其二，dyadic 预算下的最优反演成功率在临界尺度上出现精确双折点相图；其三，共振常数 \(C_\varphi\) 又把这种纤维非均匀性提升为总变差与碰撞缺口上的常数级障碍。

\begin{theorem}[终位扇区内的序理想分层]\label{thm:xi-time-part55ba-terminal-sector-order-ideal}
\leanverified{paper\_xi\_time\_part55ba\_terminal\_sector\_order\_ideal}
固定 \(m\ge 1\)，并沿用命题~\ref{prop:fold-bin-digit-dp} 与命题~\ref{prop:xi-foldbin-fiber-wallcrossing-representation} 的记号。对 \(\tau\in\{0,1\}\) 定义终位扇区
\[
X_m^{(\tau)}:=\{w\in X_m:\ w_m=\tau\}.
\]
则存在一个有限集合 \(\mathcal F_{m,\tau}\) 及权函数
\[
W_{m,\tau}:\mathcal F_{m,\tau}\to \NN,
\]
使得对任意 \(w\in X_m^{(\tau)}\) 都有
\[
d_m^{\mathrm{bin}}(w)
=
\#\Bigl\{c\in\mathcal F_{m,\tau}:\
W_{m,\tau}(c)\le 2^m-1-V_m(w)\Bigr\}.
\]
若记
\[
d_{m,\tau}(V):=d_m^{\mathrm{bin}}(w)
\qquad
\bigl(w\in X_m^{(\tau)},\ V_m(w)=V\bigr),
\]
则该定义良定。
因此，在每个固定终位扇区内：
\begin{enumerate}
  \item \(d_m^{\mathrm{bin}}(w)\) 仅依赖于 \(V_m(w)\)；
  \item 映射 \(V\mapsto d_{m,\tau}(V)\) 单调不增；
  \item 对任意整数 \(r\ge 1\)，上水平集
  \[
  \{w\in X_m^{(\tau)}:\ d_m^{\mathrm{bin}}(w)\ge r\}
  \]
  在 \(V_m\)-次序下总是初始段。
\end{enumerate}
换言之，全部纤维层级在每个终位扇区内都呈现严格的序理想分层。
\end{theorem}
\begin{proof}
由命题~\ref{prop:fold-bin-digit-dp}，固定 \(w\) 后，其全部原像可与满足相邻约束、边界约束与阈值约束的尾部位串一一对应；而命题~\ref{prop:xi-foldbin-fiber-wallcrossing-representation} 已将该计数重写为“尾部权重不超过 \(\mathrm{slack}(w)\)”的可行解个数，其中
\[
\mathrm{slack}(w)=2^m-1-V_m(w).
\]
当末位 \(\tau=w_m\) 固定时，允许的尾部状态空间与边界约束一并固定，因此可把全部可行尾部配置编码成某个有限集合 \(\mathcal F_{m,\tau}\)，并令 \(W_{m,\tau}(c)\) 表示相应尾部的权重和。
于是
\[
d_m^{\mathrm{bin}}(w)
=
\#\{c\in\mathcal F_{m,\tau}:W_{m,\tau}(c)\le \mathrm{slack}(w)\}.
\]
这就得到主公式。

由于右端仅通过阈值 \(\mathrm{slack}(w)=2^m-1-V_m(w)\) 依赖于 \(w\)，故 \(d_m^{\mathrm{bin}}(w)\) 在固定扇区内只依赖于 \(V_m(w)\)。
又因阈值随 \(V_m(w)\) 增大而单调减小，故 \(V\mapsto d_{m,\tau}(V)\) 单调不增。
最后，对任意 \(r\ge 1\)，集合
\[
\{V:\ d_{m,\tau}(V)\ge r\}
\]
必为某个形如 \((-\infty,V_\ast]\cap V_m(X_m^{(\tau)})\) 的初始段，于是对应的词集合也在 \(V_m\)-次序下形成初始段。证毕。
\end{proof}

\begin{theorem}[临界 bit-budget 的精确双折点相图]\label{thm:xi-time-part55ba-critical-oracle-phase-diagram}
沿用定理~\ref{thm:fold-oracle-capacity-closed-form} 的记号，记
\[
\mathrm{Succ}_m(B):=2^{-m}\sum_{w\in X_m}\min\bigl(d_m^{\mathrm{bin}}(w),2^B\bigr).
\]
设 \(s>0\)，并取任意整数序列 \((B_m)_{m\ge 1}\) 满足
\[
2^{B_m}=s\left(\frac{2}{\varphi}\right)^m(1+o(1)).
\]
则
\[
\mathrm{Succ}_m(B_m)\longrightarrow \mathcal S_\varphi(s),
\]
其中极限相图为
\[
\mathcal S_\varphi(s)
=
\frac{\varphi}{\sqrt5}\min\{1,s\}
+
\frac{1}{\sqrt5}\min\{\varphi^{-1},s\}.
\]
等价地，
\[
\mathcal S_\varphi(s)=
\begin{cases}
\dfrac{\varphi^2}{\sqrt5}\,s, & 0\le s\le \varphi^{-1},\\[1ex]
\dfrac{\varphi}{\sqrt5}\,s+\dfrac{1}{\varphi\sqrt5}, & \varphi^{-1}\le s\le 1,\\[1ex]
1, & s\ge 1.
\end{cases}
\]
因此，bin-fold 预言机容量在临界尺度上恰出现两个不可消去折点 \(s=\varphi^{-1}\) 与 \(s=1\)。
\end{theorem}
\begin{proof}
由定理~\ref{thm:fold-oracle-capacity-closed-form}，
\[
\mathrm{Succ}_m(B_m)=2^{-m}\sum_{w\in X_m}\min\bigl(d_m^{\mathrm{bin}}(w),2^{B_m}\bigr).
\]
把 \(X_m\) 按末位分解为
\[
A_0:=X_m^{(0)},
\qquad
A_1:=X_m^{(1)}.
\]
由定理~\ref{thm:fold-bin-two-state-asymptotic}，对 \(w\in A_b\) 一致有
\[
d_m^{\mathrm{bin}}(w)=\varphi^{-b}\left(\frac{2}{\varphi}\right)^m+O(1).
\]
结合
\[
2^{B_m}=s\left(\frac{2}{\varphi}\right)^m(1+o(1)),
\]
可知对每个固定 \(b\in\{0,1\}\)，
\[
\min\bigl(d_m^{\mathrm{bin}}(w),2^{B_m}\bigr)
=
\left(\frac{2}{\varphi}\right)^m
\min\{\varphi^{-b},s\}
+O(1)
\]
在 \(A_b\) 上一致成立。
于是
\[
\sum_{w\in A_b}\min\bigl(d_m^{\mathrm{bin}}(w),2^{B_m}\bigr)
=
|A_b|\left(\frac{2}{\varphi}\right)^m\min\{\varphi^{-b},s\}+O(|A_b|).
\]
又
\[
|A_0|=F_{m+1},
\qquad
|A_1|=F_m,
\]
故
\[
\mathrm{Succ}_m(B_m)
=
\varphi^{-m}\Bigl(F_{m+1}\min\{1,s\}+F_m\min\{\varphi^{-1},s\}\Bigr)
+O\!\left(\frac{|X_m|}{2^m}\right).
\]
最后用 Binet 渐近
\[
\frac{F_{m+1}}{\varphi^m}\to \frac{\varphi}{\sqrt5},
\qquad
\frac{F_m}{\varphi^m}\to \frac{1}{\sqrt5},
\qquad
\frac{|X_m|}{2^m}=O\!\left(\left(\frac{\varphi}{2}\right)^m\right)\to 0,
\]
即得所示极限相图。证毕。
\end{proof}

\begin{corollary}[强逆定理与常数比特精确量化律]\label{cor:xi-time-part55ba-strong-converse-quantized-budget}
设 \(B_m=\alpha m+o(m)\)。则：
\begin{enumerate}
  \item 若
  \[
  \alpha<\log_2\frac{2}{\varphi},
  \]
  则
  \[
  \mathrm{Succ}_m(B_m)\to 0;
  \]
  \item 若
  \[
  \alpha>\log_2\frac{2}{\varphi},
  \]
  则
  \[
  \mathrm{Succ}_m(B_m)\to 1.
  \]
\end{enumerate}
此外，对任意目标成功率 \(\eta\in(0,1)\)，使 \(\mathrm{Succ}_m(B_m)\to\eta\) 的最小比特预算满足
\[
B_m
=
m\log_2\frac{2}{\varphi}
+\log_2 s_\eta
+o(1),
\]
其中
\[
s_\eta=
\begin{cases}
\dfrac{\sqrt5}{\varphi^2}\eta, & 0<\eta\le \dfrac{\varphi}{\sqrt5},\\[1.2ex]
\dfrac{\sqrt5\,\eta-\varphi^{-1}}{\varphi}, & \dfrac{\varphi}{\sqrt5}\le \eta<1.
\end{cases}
\]
因此，临界过渡完全发生在常数比特尺度，而不存在中心极限定理型的 \(\sqrt m\)-窗口。
\end{corollary}
\begin{proof}
记
\[
s_m:=2^{B_m-m\log_2(2/\varphi)}.
\]
若 \(\alpha<\log_2(2/\varphi)\)，则 \(s_m\to 0\)；若 \(\alpha>\log_2(2/\varphi)\)，则 \(s_m\to\infty\)。由定理~\ref{thm:xi-time-part55ba-critical-oracle-phase-diagram} 中 \(\mathcal S_\varphi(s)\) 的三段闭式，分别得到
\[
\mathcal S_\varphi(0)=0,
\qquad
\lim_{s\to\infty}\mathcal S_\varphi(s)=1,
\]
从而前两条成立。

对第三条，只需对严格递增函数 \(\mathcal S_\varphi\) 在 \((0,1)\) 上作显式反演。
当 \(0<\eta\le \varphi/\sqrt5\) 时，
\[
\eta=\frac{\varphi^2}{\sqrt5}s
\]
故
\[
s_\eta=\frac{\sqrt5}{\varphi^2}\eta.
\]
当 \(\varphi/\sqrt5\le \eta<1\) 时，
\[
\eta=\frac{\varphi}{\sqrt5}s+\frac{1}{\varphi\sqrt5},
\]
故
\[
s_\eta=\frac{\sqrt5\,\eta-\varphi^{-1}}{\varphi}.
\]
再代回
\[
2^{B_m}\sim s_\eta\left(\frac{2}{\varphi}\right)^m
\]
即可得到所示预算公式。最后一句只是把这一预算公式改写为“临界修正仅为常数项”的表述。证毕。
\end{proof}

\begin{theorem}[黄金共振导致的不可达常数隙]\label{thm:xi-time-part55ba-golden-resonance-nonuniform-gap}
设 \(\mu_m\) 为 bin-fold 的推前分布，
\[
\mu_m(r):=\frac{d_m(r)}{2^m}
\qquad
\bigl(r\in \ZZ/F_{m+2}\ZZ\bigr),
\]
其中 \(d_m(r)\) 表示二进制区间折叠在余类 \(r\) 上的推前 multiplicity。记 \(u_m\) 为同一循环群上的 Haar 均匀分布。
定义
\[
a_m(k):=\frac{|\widehat d_m(k)|}{2^m},
\qquad
\mathsf{Col}_m:=\sum_{r}\mu_m(r)^2.
\]
则有严格下界
\[
\liminf_{m\to\infty}\|\mu_m-u_m\|_{\mathrm{TV}}
\ge \frac{C_\varphi}{2},
\]
\[
\liminf_{m\to\infty}
F_{m+2}\left(\mathsf{Col}_m-\frac1{F_{m+2}}\right)
\ge 2C_\varphi^2.
\]
因此，黄金共振在 bin-fold 推前分布上留下了不可消去的常数级指纹：它不仅阻止二阶碰撞收敛到均匀基线，也在总变差意义上强制保持一个严格正的极限缺口。
\end{theorem}
\begin{proof}
由定理~\ref{thm:fold-critical-resonance-constant}，
\[
a_m(F_m)\longrightarrow C_\varphi,
\qquad
a_m(F_{m+1})\longrightarrow C_\varphi.
\]
另一方面，对任意 \(k\neq 0\)，分布差的 Fourier 系数满足
\[
\widehat{\mu_m-u_m}(k)=\frac{\widehat d_m(k)}{2^m},
\]
故
\[
2\|\mu_m-u_m\|_{\mathrm{TV}}
=
\sum_r|\mu_m(r)-u_m(r)|
\ge
\left|\sum_r(\mu_m(r)-u_m(r))e^{-2\pi i kr/F_{m+2}}\right|
=
a_m(k).
\]
取 \(k=F_m\) 并令 \(m\to\infty\)，即得
\[
\liminf_{m\to\infty}\|\mu_m-u_m\|_{\mathrm{TV}}
\ge \frac{C_\varphi}{2}.
\]

再由 Parseval 表示
\[
F_{m+2}\,\mathsf{Col}_m=\sum_{k\in \ZZ/(F_{m+2}\ZZ)}a_m(k)^2,
\]
并注意到零频 \(a_m(0)=1\)，可得
\[
F_{m+2}\left(\mathsf{Col}_m-\frac1{F_{m+2}}\right)
=
\sum_{k\neq 0}a_m(k)^2
\ge
a_m(F_m)^2+a_m(F_{m+1})^2.
\]
令 \(m\to\infty\) 即得
\[
\liminf_{m\to\infty}
F_{m+2}\left(\mathsf{Col}_m-\frac1{F_{m+2}}\right)
\ge 2C_\varphi^2.
\]
证毕。
\end{proof}

\paragraph{局部化相位采样、solenoid 刚性与分层 prime-register 的极小性}
以下条目把有限素数局部化数系的相位采样闭式、基本扩张的不分裂性、solenoid 端同态环以及有限指数子群/有限商/profinite 完成/连通 Lie 商的刚性分类、dyadic 读出的奇素失明、与有限生成离散交换群的严格不可混同性，以及分层 prime-register 外置在最坏情形下的锐最优性，压缩为同一条可直接调用的终端链条。

\begin{theorem}[局部化数系相位采样函数的闭式与素数剥离律]\label{thm:xi-localized-phase-sampling-closed-form}
\leanverified{paper\_xi\_localized\_phase\_sampling\_closed\_form}
设 \(S\subset\mathbb P\) 为有限素数集，并记
\[
G_S:=\ZZ[S^{-1}]
\]
为加法群。定义相位采样函数
\[
\Sigma^{\mathrm{ph}}_{G_S}(N):=\card{\operatorname{Hom}(G_S,\ZZ/N\ZZ)}
\qquad(N\ge 1).
\]
则对一切 \(N\ge 1\) 都有闭式
\[
\boxed{\;
\Sigma^{\mathrm{ph}}_{G_S}(N)
=
\prod_{\substack{q^{v_q(N)}\parallel N\\ q\notin S}}q^{v_q(N)}
=
\frac{N}{\prod_{p\in S}p^{v_p(N)}}.\;}
\]
等价地，\(\Sigma^{\mathrm{ph}}_{G_S}(N)\) 恰由 \(N\) 去掉全部 \(S\)-主部分得到。特别地，
\[
p\in S \iff \Sigma^{\mathrm{ph}}_{G_S}(p)=1,
\qquad
p\notin S \iff \Sigma^{\mathrm{ph}}_{G_S}(p)=p.
\]
因此 \(\Sigma^{\mathrm{ph}}_{G_S}\) 唯一恢复素数集 \(S\)，并因定理 \ref{thm:killo-localization-circle-dimension-prime-ledger-splitting} 进一步唯一恢复 Pontryagin 对偶短正合列
\[
0\to \prod_{p\in S}\ZZ_p\to \widehat{G_S}\to \TT\to 0
\]
中的 profinite 核；与此同时 \(\cdim_{\mathrm{fr}}(G_S)=1\) 与 \(S\) 无关。
\end{theorem}
\begin{proof}
定理 \ref{thm:xi-cdim-localization-zeta-AS} 已证明：若把同一加法群记作
\(
A_S=\ZZ[S^{-1}]
\)，则
\[
\card{\operatorname{Hom}(A_S,\ZZ/N\ZZ)}
=
\prod_{\substack{q^{v_q(N)}\parallel N\\ q\notin S}}q^{v_q(N)}
=
\frac{N}{\prod_{p\in S}p^{v_p(N)}}
\qquad(N\ge 1).
\]
将 \(A_S\) 记为 \(G_S\) 即得所述闭式。
\par
对任意素数 \(p\)，代入 \(N=p\) 得
\[
\Sigma^{\mathrm{ph}}_{G_S}(p)=
\begin{cases}
1,& p\in S,\\[1mm]
p,& p\notin S.
\end{cases}
\]
故 \(S\) 由 \(\Sigma^{\mathrm{ph}}_{G_S}\) 唯一恢复。
再由定理 \ref{thm:killo-localization-circle-dimension-prime-ledger-splitting}，核
\(\prod_{p\in S}\ZZ_p\) 的非平凡 \(p\)-主因子与 \(p\in S\) 等价，而 \(\cdim_{\mathrm{fr}}(G_S)=1\) 始终成立。证毕。
\end{proof}

\begin{corollary}[dyadic 相位谱对奇素支撑的完全失明]\label{cor:xi-localized-dyadic-phase-blindness}
\leanverified{paper\_xi\_localized\_dyadic\_phase\_blindness}
沿用定理 \ref{thm:xi-localized-phase-sampling-closed-form} 的记号，定义 dyadic 相位谱
\[
\Sigma^{(2)}_{G_S}(a):=\Sigma^{\mathrm{ph}}_{G_S}(2^a)
\qquad(a\ge 1).
\]
则有严格二分
\[
\boxed{\;
\Sigma^{(2)}_{G_S}(a)=
\begin{cases}
2^a,& 2\notin S,\\[1mm]
1,& 2\in S.
\end{cases}\;}
\]
因此若有限素数集 \(S,T\subset\mathbb P\) 满足
\[
2\in S \iff 2\in T,
\]
则
\[
\Sigma^{(2)}_{G_S}(a)=\Sigma^{(2)}_{G_T}(a)
\qquad(a\ge 1).
\]
换言之，dyadic 相位谱只能探测 \(2\) 是否进入素数账本，而对全部奇素支撑严格失明。
\end{corollary}
\begin{proof}
把 \(N=2^a\) 代入定理 \ref{thm:xi-localized-phase-sampling-closed-form} 即得
\[
\Sigma^{\mathrm{ph}}_{G_S}(2^a)
=
\frac{2^a}{\prod_{p\in S}p^{v_p(2^a)}}
=
\begin{cases}
2^a,& 2\notin S,\\[1mm]
1,& 2\in S.
\end{cases}
\]
而对任意奇素 \(p\) 恒有 \(v_p(2^a)=0\)，故其是否属于 \(S\) 不会改变数值。证毕。
\end{proof}

\begin{theorem}[局部化数系对偶短正合列的分裂判据]\label{thm:xi-localized-basic-extension-strictly-nonsplit}
\leanverified{paper\_xi\_localized\_basic\_extension\_strictly\_nonsplit}
设 \(S\subset\mathbb P\) 为有限素数集，并记
\[
G_S:=\ZZ[S^{-1}].
\]
再记
\[
\Sigma_S:=\widehat{G_S}.
\]
前述结果给出互为对偶的短正合列
\[
0\longrightarrow \ZZ\longrightarrow G_S\longrightarrow \bigoplus_{p\in S}\ZZ(p^\infty)\longrightarrow 0
\]
\[
0\longrightarrow \prod_{p\in S}\ZZ_p\longrightarrow \widehat{G_S}\longrightarrow \TT\longrightarrow 0.
\]
则下列两条等价：
\begin{enumerate}
  \item \(S=\varnothing\)；
  \item 上述两条短正合列在各自范畴中都分裂。
\end{enumerate}
等价地，
\[
S\neq\varnothing
\quad\Longrightarrow\quad
\Sigma_S\not\cong \TT\times \prod_{p\in S}\ZZ_p.
\]
特别地，当 \(S\neq\varnothing\) 时，不存在群同态截面
\[
\bigoplus_{p\in S}\ZZ(p^\infty)\longrightarrow G_S,
\]
也不存在连续群同态截面
\[
\TT\longrightarrow \widehat{G_S}.
\]
\end{theorem}
\begin{proof}
若 \(S=\varnothing\)，则 \(G_S=\ZZ\) 且 \(\Sigma_S=\widehat{\ZZ}\cong \TT\)，两条短正合列都退化为平凡扩张，故显然分裂。

现设 \(S\neq\varnothing\)。若第一条短正合列分裂，则
\[
G_S\cong \ZZ\oplus \bigoplus_{p\in S}\ZZ(p^\infty).
\]
右端含有非零挠子群，而
\[
G_S=\ZZ[S^{-1}]\subset \QQ
\]
作为有理数加法群的子群无挠，矛盾。故离散侧不分裂。

Pontryagin 对偶在局部紧交换群范畴中给出精确反变等价。若第二条短正合列分裂，则再对偶化便得到第一条离散短正合列分裂，与上一步矛盾。故紧侧亦不分裂。证毕。
\end{proof}

\begin{theorem}[局部化数系端同态环的完全分类与 solenoid 的直积不可分解性]\label{thm:xi-localized-endomorphism-ring-solenoid-indecomposable}
沿用定理 \ref{thm:xi-localized-basic-extension-strictly-nonsplit} 的记号。则：
\begin{enumerate}
  \item 任意群同态 \(\phi:G_S\to G_S\) 都唯一地写成
  \[
  \phi(x)=qx
  \qquad (x\in G_S)
  \]
  的形式，其中 \(q\in G_S\)。
  \item 因而
  \[
  \operatorname{End}(G_S)\cong G_S
  \]
  作为交换环同构成立，其中加法为点态相加，乘法为复合。
  \item \(\operatorname{End}(G_S)\) 的幂等元只有 \(0,1\)。等价地，
  \[
  \operatorname{End}_{\mathrm{cts}}(\widehat{G_S})
  \cong
  \operatorname{End}(G_S)^{\mathrm{op}}
  \cong
  G_S
  \]
  亦只含平凡幂等元；特别地，\(\widehat{G_S}\) 不存在非平凡连续幂等自同态，因而不能分解为两个非平凡紧交换群的直积。
\end{enumerate}
\end{theorem}
\begin{proof}
取任意 \(\phi\in \operatorname{End}(G_S)\)，记
\[
q:=\phi(1)\in G_S.
\]
对任意 \(x\in G_S\)，可写成
\[
x=\frac{a}{D},
\qquad
a\in \ZZ,
\qquad
D=\prod_{p\in S}p^{k_p}
\]
的形式。由 \(Dx=a\) 可得
\[
D\,\phi(x)=\phi(a)=aq.
\]
由于 \(G_S\subset \QQ\) 无挠，乘以 \(D\) 的映射单射，故
\[
\phi(x)=\frac{aq}{D}=qx.
\]
这证明了第 \(1\) 条；唯一性由在 \(x=1\) 处取值立即得到。反过来，任取 \(q\in G_S\)，映射 \(x\mapsto qx\) 因 \(G_S\) 为 \(\QQ\) 的子环而确为 \(G_S\) 的自同态。于是
\[
q\longmapsto m_q,\qquad m_q(x):=qx
\]
给出 \(\operatorname{End}(G_S)\) 与 \(G_S\) 之间的双射。

该双射保持加法，且满足
\[
m_q\circ m_r=m_{qr},
\]
故第 \(2\) 条成立。

由于 \(G_S=\ZZ[S^{-1}]\subset \QQ\) 是整环的子环，幂等方程
\[
e^2=e
\]
只可能有解 \(e=0\) 或 \(e=1\)。故 \(\operatorname{End}(G_S)\) 仅有平凡幂等元。Pontryagin 对偶把连续端同态环识别为离散端端同态环的反环；又因 \(G_S\) 交换，反环与自身一致，故 \(\operatorname{End}_{\mathrm{cts}}(\widehat{G_S})\) 同样只含平凡幂等元。

若 \(\widehat{G_S}\cong K_1\times K_2\) 为两个非平凡紧交换群的直积，则合成
\[
\widehat{G_S}\xrightarrow{\ \mathrm{pr}_1\ }K_1
\hookrightarrow
K_1\times K_2\cong \widehat{G_S}
\]
给出一个既非零也非恒等的连续幂等自同态，矛盾。故 \(\widehat{G_S}\) 不能作非平凡直积分解。证毕。
\end{proof}

\begin{theorem}[有限商与有限相位层的完全分类]\label{thm:xi-localized-finite-quotient-phase-layer-classification}
\leanverified{paper\_xi\_localized\_finite\_quotient\_phase\_layer\_classification}
沿用定理 \ref{thm:xi-localized-basic-extension-strictly-nonsplit} 的记号。对任意正整数 \(n\)，定义
\[
n_{S^\perp}:=\frac{n}{\prod_{p\in S}p^{v_p(n)}},
\]
即从 \(n\) 中剥离全部 \(S\)-主因子后得到的 \(S\)-自由部分。则：
\begin{enumerate}
  \item 商群满足
  \[
  G_S/nG_S\cong \ZZ/n_{S^\perp}\ZZ.
  \]
  \item 因而 \(G_S\) 的全部有限商恰为循环群 \(\ZZ/d\ZZ\)，其中 \(d\) 的全部素因子都不属于 \(S\)。
  \item 对偶地，\(\widehat{G_S}\) 的全部有限闭子群也恰为这些循环群。
\end{enumerate}
\end{theorem}
\begin{proof}
把 \(n\) 分解为
\[
n=n_S\,n_{S^\perp},
\qquad
n_S:=\prod_{p\in S}p^{v_p(n)}.
\]
对每个 \(p\in S\)，乘法映射 \(x\mapsto px\) 在 \(G_S=\ZZ[S^{-1}]\) 上是自同构，因为 \(x\mapsto x/p\) 仍取值于 \(G_S\)。故
\[
nG_S=n_{S^\perp}G_S.
\]
因此只需证明当 \(n\) 与 \(\prod_{p\in S}p\) 互素时，
\[
G_S/nG_S\cong \ZZ/n\ZZ.
\]

考虑自然映射
\[
\rho:\ZZ\longrightarrow G_S/nG_S.
\]
任取 \(x=a/b\in G_S\)，其中 \(a\in\ZZ\)，且 \(b\) 的全部素因子属于 \(S\)。由于 \(\gcd(b,n)=1\)，可取 \(c,k\in\ZZ\) 使
\[
bc=1+nk.
\]
于是
\[
x-ac
=
\frac{a}{b}-ac
=
\frac{a(1-bc)}{b}
=
-n\,\frac{ak}{b}\in nG_S,
\]
故 \(x\equiv ac\pmod{nG_S}\)。这说明 \(\rho\) 满射。

再看核。若 \(m\in\ker\rho\)，则
\[
m=n\frac{a}{b}
\]
对某个 \(\frac{a}{b}\in G_S\) 成立，且 \(\gcd(b,n)=1\)。于是
\[
bm=na.
\]
由于 \(\gcd(b,n)=1\)，必有 \(b\mid a\)，从而 \(a/b\in \ZZ\)，进而 \(m\in n\ZZ\)。故
\[
\ker(\rho)=n\ZZ,
\]
于是
\[
G_S/nG_S\cong \ZZ/n\ZZ.
\]
把 \(n\) 替换回 \(n_{S^\perp}\) 即得第 \(1\) 条。

若 \(F\) 为 \(G_S\) 的任一有限商，令 \(n\) 为 \(F\) 的指数，则满射 \(G_S\twoheadrightarrow F\) 经由 \(G_S/nG_S\) 因子化。由第 \(1\) 条，\(G_S/nG_S\) 为循环群，故 \(F\) 亦为循环群，且其阶的全部素因子都不属于 \(S\)。反之，若 \(d\) 的全部素因子都不属于 \(S\)，则 \(d=d_{S^\perp}\)，从而
\[
G_S/dG_S\cong \ZZ/d\ZZ.
\]
第 \(2\) 条得证。

第 \(3\) 条由 Pontryagin 对偶把有限离散商与有限闭子群一一对应，并且有限循环群的对偶仍为同阶循环群，立即得到。证毕。
\end{proof}

\paragraph{solenoid 覆盖度谱、乘法周期律与有限闭子群的循环实现}
在局部化数系的端同态环、有限商与有限闭子群已经完成分类之后，还可把整数乘法在 \(\widehat{G_S}\) 上进一步压缩为一组完全显式的覆盖与动力学结论。其核心现象是：所有有限层数据都只由 \(S\)-自由部分
\[
n_{S^\perp}:=\frac{n}{\prod_{p\in S}p^{v_p(n)}}
\]
控制。

\begin{theorem}[整数乘法在 \texorpdfstring{$\widehat{G_S}$}{G-hat-S} 上的核度公式]\label{thm:xi-localized-solenoid-multiplication-kernel-degree}
设 \(S\subset\mathbb P\) 为有限非空素数集，记
\[
G_S:=\ZZ[S^{-1}],
\qquad
\Sigma_S:=\widehat{G_S}.
\]
对任意正整数 \(n\)，考虑整数乘法自同态
\[
[n]:\Sigma_S\longrightarrow \Sigma_S,
\qquad
x\longmapsto nx.
\]
则 \([n]\) 必为满射，其核为有限循环群，并且
\[
\boxed{\;
\ker[n]\cong \ZZ/n_{S^\perp}\ZZ,
\qquad
\deg([n])=n_{S^\perp}.\;}
\]
\end{theorem}
\begin{proof}
令
\[
m_n:G_S\to G_S,
\qquad
x\mapsto nx
\]
为 Pontryagin 对偶端映射。由于
\[
G_S=\ZZ[S^{-1}]\subset \QQ
\]
无挠，\(m_n\) 为单射。于是对偶映射 \([n]\) 为满射。

再看短正合列
\[
0\longrightarrow G_S\xrightarrow{\,m_n\,}G_S\longrightarrow G_S/nG_S\longrightarrow 0.
\]
Pontryagin 对偶的精确性给出
\[
0\longrightarrow \widehat{G_S/nG_S}\longrightarrow \Sigma_S\xrightarrow{\,[n]\,}\Sigma_S\longrightarrow 0,
\]
故
\[
\ker[n]\cong \widehat{G_S/nG_S}.
\]
而定理 \ref{thm:xi-localized-finite-quotient-phase-layer-classification} 已证明
\[
G_S/nG_S\cong \ZZ/n_{S^\perp}\ZZ.
\]
有限循环群的 Pontryagin 对偶仍为同阶循环群，因此
\[
\ker[n]\cong \ZZ/n_{S^\perp}\ZZ.
\]

最后求覆盖度。记 \(F:=\ker[n]\)。由于 \([n]\) 为满射群同态，诱导同构
\[
\Sigma_S/F\cong \Sigma_S.
\]
有限群 \(F\) 通过平移自由作用于 \(\Sigma_S\)。取单位元的一个开邻域 \(U\) 使
\[
(U-U)\cap F=\{0\},
\]
则不同 \(F\)-平移 \(U+\eta\) 两两不交。故商映射
\[
q:\Sigma_S\to \Sigma_S/F
\]
是度数为 \(|F|\) 的有限正则覆盖，而 \([n]\) 与该商映射只差一个群同构。于是
\[
\deg([n])=|F|=n_{S^\perp}.
\]
证毕。
\end{proof}

\begin{theorem}[局部化标量乘法在 \texorpdfstring{$\widehat{G_S}$}{G-hat-S} 上的核度公式]\label{thm:xi-localized-solenoid-localized-scalar-kernel-degree}
设 \(S\subset\mathbb P\) 为有限素数集，记
\[
G_S:=\ZZ[S^{-1}],
\qquad
\Sigma_S:=\widehat{G_S}.
\]
任取非零标量
\[
q\in \ZZ[S^{-1}],
\qquad
q=\varepsilon\,a\prod_{p\in S}p^{-k_p},
\qquad
\varepsilon\in\{\pm1\},
\quad
a\in\NN,
\]
并定义其 \(S\)-自由部分
\[
a_{S^\perp}:=\frac{a}{\prod_{p\in S}p^{v_p(a)}}.
\]
令
\[
[q]:\Sigma_S\longrightarrow \Sigma_S
\]
表示与离散侧“乘以 \(q\)”对偶的连续自同态。则 \([q]\) 为满射，其核为有限循环群，并且
\[
\boxed{\;
\ker[q]\cong \ZZ/a_{S^\perp}\ZZ,
\qquad
\deg([q])=a_{S^\perp}.\;}
\]
特别地，
\[
[q]\ \text{是同构}
\iff
a_{S^\perp}=1
\iff
q\in \ZZ[S^{-1}]^\times.
\]
\end{theorem}
\begin{proof}
令
\[
m_q:G_S\to G_S,
\qquad
x\mapsto qx
\]
为离散侧端映射。由于 \(G_S\subset \QQ\) 无挠且 \(q\neq 0\)，故 \(m_q\) 为单射，从而其对偶 \([q]\) 为满射。

把 \(a\) 分解为
\[
a=a_S\,a_{S^\perp},
\qquad
a_S:=\prod_{p\in S}p^{v_p(a)}.
\]
则
\[
u:=\varepsilon\,a_S\prod_{p\in S}p^{-k_p}
\]
是局部化环 \(\ZZ[S^{-1}]\) 的单位，且
\[
q=u\,a_{S^\perp}.
\]
于是
\[
qG_S=a_{S^\perp}G_S.
\]
对短正合列
\[
0\longrightarrow G_S\xrightarrow{\,m_q\,}G_S\longrightarrow G_S/qG_S\longrightarrow 0
\]
作 Pontryagin 对偶，得到
\[
0\longrightarrow \widehat{G_S/qG_S}\longrightarrow \Sigma_S\xrightarrow{\,[q]\,}\Sigma_S\longrightarrow 0,
\]
故
\[
\ker[q]\cong \widehat{G_S/qG_S}\cong \widehat{G_S/a_{S^\perp}G_S}.
\]
而定理 \ref{thm:xi-localized-finite-quotient-phase-layer-classification} 已给出
\[
G_S/a_{S^\perp}G_S\cong \ZZ/a_{S^\perp}\ZZ.
\]
有限循环群的 Pontryagin 对偶仍为同阶循环群，遂得
\[
\ker[q]\cong \ZZ/a_{S^\perp}\ZZ.
\]

最后，\([q]\) 是紧交换群上的满同态且核有限，故如定理 \ref{thm:xi-localized-solenoid-multiplication-kernel-degree} 的证明所示，它是度数等于核大小的有限正则覆盖。因此
\[
\deg([q])=\#\ker[q]=a_{S^\perp}.
\]
核平凡当且仅当 \(a_{S^\perp}=1\)，这又等价于 \(q\) 没有任何逃离 \(S\) 的素因子，即 \(q\in \ZZ[S^{-1}]^\times\)。证毕。
\end{proof}

\begin{corollary}[覆盖度谱完全由被剥离后的素数支撑决定]\label{cor:xi-localized-solenoid-cover-degree-spectrum}
沿用定理 \ref{thm:xi-localized-solenoid-multiplication-kernel-degree} 的记号。则 \(\Sigma_S\) 上全部整数乘法覆盖的度数集合恰为
\[
\boxed{\;
\operatorname{Deg}(\Sigma_S)
:=
\{\deg([n]): n\ge 1\}
=
\{d\in\NN:\ \text{\(d\) 的全部素因子都不属于 }S\}.\;}
\]
并且对任意 \(n\ge 1\)，以下三条等价：
\[
[n]\ \text{是同构}
\iff
n_{S^\perp}=1
\iff
\text{\(n\) 的全部素因子都属于 }S.
\]
\end{corollary}
\begin{proof}
由定理 \ref{thm:xi-localized-solenoid-multiplication-kernel-degree}，
\[
\deg([n])=n_{S^\perp}.
\]
故任意整数乘法所产生的覆盖度都自动剥离了全部 \(S\)-主因子。反之，若 \(d\) 的全部素因子都不属于 \(S\)，则
\[
d_{S^\perp}=d,
\]
取 \(n=d\) 即得 \(\deg([d])=d\)。这就给出度数谱的精确描述。

又因为 \([n]\) 已知为满射，所以它为同构当且仅当核平凡；而定理
\ref{thm:xi-localized-solenoid-multiplication-kernel-degree}
给出
\[
|\ker[n]|=n_{S^\perp}.
\]
故 \([n]\) 为同构当且仅当 \(n_{S^\perp}=1\)，也即 \(n\) 的全部素因子都落在 \(S\) 中。证毕。
\end{proof}

\begin{theorem}[整数乘法动力学的周期点计数与 primitive 轨道公式]\label{thm:xi-localized-solenoid-periodic-point-formula}
固定整数 \(a\ge 2\)，考虑连续群自映射
\[
T_a:\Sigma_S\longrightarrow \Sigma_S,
\qquad
x\longmapsto ax.
\]
则对任意 \(m\ge 1\)，其 \(m\)-步不动点群满足
\[
\boxed{\;
\mathrm{Fix}(T_a^m)\cong \ZZ/(a^m-1)_{S^\perp}\ZZ,
\qquad
\#\mathrm{Fix}(T_a^m)=(a^m-1)_{S^\perp}.\;}
\]
于是 exact period 为 \(m\) 的 primitive 轨道数满足
\[
\boxed{\;
\mathcal O_m(T_a)
=
\frac1m\sum_{d\mid m}\mu(d)\,(a^{m/d}-1)_{S^\perp}.\;}
\]
\end{theorem}
\begin{proof}
由定义，
\[
T_a^m=[a^m].
\]
故
\[
x\in \mathrm{Fix}(T_a^m)
\iff
a^m x=x
\iff
(a^m-1)x=0
\iff
x\in \ker[a^m-1].
\]
因此
\[
\mathrm{Fix}(T_a^m)=\ker[a^m-1].
\]
将定理 \ref{thm:xi-localized-solenoid-multiplication-kernel-degree} 应用于整数
\[
n=a^m-1
\]
即得
\[
\ker[a^m-1]\cong \ZZ/(a^m-1)_{S^\perp}\ZZ,
\]
并且
\[
\#\mathrm{Fix}(T_a^m)=|\ker[a^m-1]|=(a^m-1)_{S^\perp}.
\]

对第二条，设 \(\mathcal O_d(T_a)\) 为 exact period 为 \(d\) 的 primitive 轨道数。则标准周期分解给出
\[
\#\mathrm{Fix}(T_a^m)=\sum_{d\mid m} d\,\mathcal O_d(T_a).
\]
对该等式施行 Möbius 反演，得到
\[
\mathcal O_m(T_a)
=
\frac1m\sum_{d\mid m}\mu(d)\,\#\mathrm{Fix}(T_a^{m/d})
=
\frac1m\sum_{d\mid m}\mu(d)\,(a^{m/d}-1)_{S^\perp}.
\]
证毕。
\end{proof}

\begin{corollary}[单质数核实验即可完全恢复 prime-register]\label{cor:xi-localized-solenoid-prime-kernel-tomography}
\leanverified{paper\_xi\_localized\_solenoid\_prime\_kernel\_tomography}
对任意素数 \(q\)，有
\[
\#\ker[q]=
\begin{cases}
1,& q\in S,\\[1mm]
q,& q\notin S.
\end{cases}
\]
等价地，
\[
\boxed{\;
q\in S
\iff
\#\ker[q]=1.\;}
\]
因此，\(\Sigma_S\) 的一阶覆盖核谱
\[
q\longmapsto \#\ker[q]
\qquad (q\ \text{为素数})
\]
已经足以唯一恢复素数账本 \(S\)。
\end{corollary}
\begin{proof}
把定理 \ref{thm:xi-localized-solenoid-multiplication-kernel-degree} 中的 \(n\) 取为素数 \(q\)。若 \(q\in S\)，则
\[
q_{S^\perp}=1,
\]
故 \(\ker[q]\) 平凡；若 \(q\notin S\)，则
\[
q_{S^\perp}=q,
\]
故
\[
\ker[q]\cong \ZZ/q\ZZ.
\]
于是所述判别式立即成立。证毕。
\end{proof}

\begin{theorem}[\texorpdfstring{$n$}{n}-挠子群与全部有限闭子群由整数乘法核唯一实现]\label{thm:xi-localized-solenoid-finite-subgroups-are-kernels}
\leanverified{paper\_xi\_localized\_solenoid\_finite\_subgroups\_are\_kernels}
对任意 \(n\ge 1\)，\(\Sigma_S=\widehat{G_S}\) 的 \(n\)-挠子群满足
\[
\Sigma_S[n]=\ker[n]\cong \ZZ/n_{S^\perp}\ZZ.
\]
更一般地，\(\Sigma_S\) 的每一个有限闭子群 \(F\le \Sigma_S\) 都必为循环群；更精确地，存在唯一满足“其全部素因子都不属于 \(S\)”的正整数 \(d\)，使得
\[
\boxed{\;
F=\ker[d]\cong \ZZ/d\ZZ.\;}
\]
反之，对每个这样的 \(d\)，\(\ker[d]\) 都给出一个有限闭子群。特别地，\(\Sigma_S\) 中不存在任何非循环有限闭子群。
\end{theorem}
\begin{proof}
按定义，
\[
\Sigma_S[n]=\{x\in \Sigma_S:\ nx=0\}=\ker[n].
\]
故第一条由定理 \ref{thm:xi-localized-solenoid-multiplication-kernel-degree} 立即得到。

任取有限闭子群 \(F\le \Sigma_S\)。由定理
\ref{thm:xi-localized-finite-quotient-phase-layer-classification}
第 \(3\) 条，\(F\) 同构于某个循环群 \(\ZZ/d\ZZ\)，其中 \(d\) 的全部素因子都不属于 \(S\)。因此 \(F\) 为阶 \(d\) 的循环群，尤其其每个元素都被 \(d\) 杀掉，故
\[
F\subseteq \ker[d].
\]
另一方面，由定理 \ref{thm:xi-localized-solenoid-multiplication-kernel-degree}，
\[
\ker[d]\cong \ZZ/d_{S^\perp}\ZZ.
\]
由于 \(d\) 本身已是 \(S\)-自由整数，故 \(d_{S^\perp}=d\)。于是
\[
|\ker[d]|=d=|F|.
\]
结合 \(F\subseteq \ker[d]\)，可得
\[
F=\ker[d].
\]
这就证明了存在性。

唯一性也随之成立：若
\[
\ker[d_1]=\ker[d_2],
\]
则两边阶数相同。再由定理 \ref{thm:xi-localized-solenoid-multiplication-kernel-degree}，
\[
d_1=|\ker[d_1]|=|\ker[d_2]|=d_2.
\]

反过来，若 \(d\) 的全部素因子都不属于 \(S\)，则 \(d_{S^\perp}=d\)，故定理
\ref{thm:xi-localized-solenoid-multiplication-kernel-degree}
给出
\[
\ker[d]\cong \ZZ/d\ZZ,
\]
从而它确为一个有限闭子群。

最后，由所有有限闭子群都等于某个 \(\ker[d]\) 且 \(\ker[d]\) 总为循环群，立即推出 \(\Sigma_S\) 中不存在任何非循环有限闭子群。证毕。
\end{proof}

若将整数乘法进一步推广为由任意 \(x\in G_S\setminus\{0\}\) 所诱导的连续自同态，则同样的 \(S\)-自由剥离律仍然保持，只是控制有限核与周期点计数的整数由 \(n\) 变为 \(x\) 的 \(S\)-外分子以及差分 \(a^n-b^n\) 的 \(S\)-外部分。

\begin{theorem}[连续自同态的核与余核由 \texorpdfstring{$S$}{S}-外分子唯一决定]\label{thm:xi-localized-solenoid-endomorphism-kernel-cokernel-sfree-numerator}
设
\[
x=\frac{a}{b}\in G_S=\ZZ[S^{-1}]\setminus\{0\},
\qquad
a\in\ZZ\setminus\{0\},
\qquad
b\in\langle S\rangle,
\qquad
\gcd(a,b)=1.
\]
定义
\[
v:=|a|_{S^\perp}
=
\frac{|a|}{\prod_{p\in S}p^{v_p(a)}}.
\]
记
\[
m_x:G_S\to G_S,
\qquad
y\mapsto xy,
\]
并令
\[
\alpha_x:=\widehat{m_x}:\Sigma_S\to \Sigma_S
\]
为其 Pontryagin 对偶连续自同态。则：
\begin{enumerate}
  \item
  \[
  \operatorname{coker}(m_x)\cong \ZZ/v\ZZ;
  \]
  \item
  \[
  \ker(\alpha_x)\cong \ZZ/v\ZZ;
  \]
  \item \(\alpha_x\) 为同构当且仅当 \(v=1\)，亦即 \(x\in G_S^\times\)。
\end{enumerate}
\end{theorem}
\begin{proof}
写
\[
|a|=uv,
\]
其中 \(u\) 的全部素因子都属于 \(S\)，而 \(v\) 与 \(\prod_{p\in S}p\) 互素。再记 \(\varepsilon:=\operatorname{sgn}(a)\in\{\pm1\}\)。于是
\[
x=\left(\varepsilon\frac{u}{b}\right)v.
\]
由于 \(u/b\in \ZZ[S^{-1}]^\times\) 且 \(\varepsilon=\pm1\) 亦为单位，故
\[
\varepsilon\frac{u}{b}\in G_S^\times.
\]
从而 \(m_x\) 与 \(m_v\) 只差一个群自同构，故
\[
\operatorname{coker}(m_x)\cong \operatorname{coker}(m_v)\cong G_S/vG_S.
\]
又因为 \(v\) 与 \(\prod_{p\in S}p\) 互素，定理 \ref{thm:xi-localized-finite-quotient-phase-layer-classification} 给出
\[
G_S/vG_S\cong \ZZ/v\ZZ,
\]
从而得到第 \(1\) 条。

另一方面，\(x\neq 0\) 且 \(G_S\subset \QQ\) 无挠，故 \(m_x\) 为单射。对短正合列
\[
0\longrightarrow G_S\xrightarrow{\,m_x\,}G_S\longrightarrow \operatorname{coker}(m_x)\longrightarrow 0
\]
施行 Pontryagin 对偶，得到
\[
0\longrightarrow \widehat{\operatorname{coker}(m_x)}\longrightarrow \Sigma_S\xrightarrow{\,\alpha_x\,}\Sigma_S\longrightarrow 0.
\]
因此
\[
\ker(\alpha_x)\cong \widehat{\operatorname{coker}(m_x)}.
\]
而有限循环群的 Pontryagin 对偶仍为同阶循环群，故第 \(2\) 条成立。

最后，由上式可知 \(\alpha_x\) 总为满射；因此 \(\alpha_x\) 为同构当且仅当 \(\ker(\alpha_x)=0\)，等价于 \(v=1\)。而 \(v=1\) 恰当且仅当 \(a\) 的全部素因子都属于 \(S\)，也即 \(x\) 在局部化环 \(G_S\) 中可逆。证毕。
\end{proof}

\paragraph{一般连续自同态的有限覆盖度与覆盖半群的算术谱}
在一般连续自同态的核与余核已经由 \(S\)-外分子完全刻画之后，还可把其拓扑非可逆性进一步压缩为覆盖度层面的两条闭式结论。

\leanverified{paper\_xi\_localized\_solenoid\_nonzero\_endomorphism\_finite\_cover\_degree}
\begin{corollary}[非零连续自同态必为有限覆盖，且覆盖度等于 \texorpdfstring{$S$}{S}-外分子]\label{cor:xi-localized-solenoid-nonzero-endomorphism-finite-cover-degree}
沿用定理 \ref{thm:xi-localized-solenoid-endomorphism-kernel-cokernel-sfree-numerator} 的记号。则 \(\alpha_x\) 是满射有限覆盖，并且
\[
\boxed{\;
\ker(\alpha_x)\cong \ZZ/v\ZZ,
\qquad
\deg(\alpha_x)=v.\;}
\]
特别地，\(\alpha_x\) 为同构当且仅当 \(v=1\)，亦即 \(x\in G_S^\times\)。
\end{corollary}
\begin{proof}
定理 \ref{thm:xi-localized-solenoid-endomorphism-kernel-cokernel-sfree-numerator} 已给出
\[
\ker(\alpha_x)\cong \ZZ/v\ZZ
\]
以及 \(\alpha_x\) 为同构当且仅当 \(v=1\)。
另一方面，由 \(x\neq 0\) 与 \(G_S\subset \QQ\) 无挠，离散端乘法映射
\[
m_x:G_S\to G_S
\]
是单射，故其对偶 \(\alpha_x\) 为满射。

再证有限覆盖性。由于 \(\ker(\alpha_x)\) 有限，可取单位元的开邻域 \(U\subset \Sigma_S\) 满足
\[
(U-U)\cap \ker(\alpha_x)=\{0\}.
\]
于是对不同的 \(\eta,\eta'\in \ker(\alpha_x)\)，平移集 \(U+\eta\) 与 \(U+\eta'\) 两两不交。又因 \(\alpha_x\) 为群同态，
\[
\alpha_x(U+\eta)=\alpha_x(U)
\qquad
(\eta\in \ker(\alpha_x)),
\]
并且若 \(u_1,u_2\in U\) 满足 \(\alpha_x(u_1)=\alpha_x(u_2)\)，则
\[
u_1-u_2\in (U-U)\cap \ker(\alpha_x)=\{0\},
\]
故 \(\alpha_x|_U\) 单射。由于连续满射群同态是开映射，\(\alpha_x(U)\) 为开集，且
\[
\alpha_x|_U:U\longrightarrow \alpha_x(U)
\]
为同胚。再利用平移共轭，可知每个 \(\alpha_x|_{U+\eta}\) 都是到 \(\alpha_x(U)\) 的同胚，故在该开邻域上给出恰好 \(|\ker(\alpha_x)|\) 个局部分支。于是 \(\alpha_x\) 是度数为
\[
\card{\ker(\alpha_x)}=v
\]
的有限覆盖。证毕。
\end{proof}

\begin{corollary}[非零连续自同态覆盖度半群的严格算术分类]\label{cor:xi-localized-solenoid-nonzero-endomorphism-degree-semigroup}
沿用定理 \ref{thm:xi-localized-solenoid-endomorphism-kernel-cokernel-sfree-numerator} 的记号。定义
\[
\deg:\operatorname{End}_{\mathrm{cts}}(\Sigma_S)\setminus\{0\}\to \NN_{\ge 1},
\qquad
\deg(\alpha):=\card{\ker(\alpha)}.
\]
则其像恰为
\[
\boxed{\;
\operatorname{Im}(\deg)
=
\left\{
n\in \NN_{\ge 1}:\ \gcd\!\left(n,\prod_{p\in S}p\right)=1
\right\}.\;}
\]
并且对任意非零连续自同态 \(\alpha_x,\alpha_y\) 都有
\[
\boxed{\;
\deg(\alpha_x\circ \alpha_y)=\deg(\alpha_x)\deg(\alpha_y).\;}
\]
因此，\(\Sigma_S\) 的全部非平凡连续覆盖半群与去掉 \(S\)-素数后的正整数乘法半群完全同构。
\end{corollary}
\begin{proof}
由定理 \ref{thm:xi-localized-endomorphism-ring-solenoid-indecomposable}，
\[
\operatorname{End}_{\mathrm{cts}}(\Sigma_S)\cong G_S,
\]
故每个非零连续自同态都唯一写成某个 \(\alpha_x\)。

任取
\[
x=\frac{a}{b}\in G_S\setminus\{0\},
\qquad
\gcd(a,b)=1,
\qquad
b\in \langle S\rangle.
\]
把
\[
|a|=uv
\]
写成 \(u\) 的全部素因子都属于 \(S\)、而 \(v\) 与 \(\prod_{p\in S}p\) 互素的分解。由推论 \ref{cor:xi-localized-solenoid-nonzero-endomorphism-finite-cover-degree}，
\[
\deg(\alpha_x)=v.
\]
因此每个覆盖度都与 \(S\) 互素。

反之，若 \(n\in\NN_{\ge1}\) 与 \(\prod_{p\in S}p\) 互素，则 \(n\in G_S\)，并且它的 \(S\)-外分子就是 \(n\) 本身。故
\[
\deg(\alpha_n)=n.
\]
从而上述集合中的每个正整数都可实现，这就给出 \(\operatorname{Im}(\deg)\) 的精确刻画。

最后证明乘法性。设
\[
x=\varepsilon_x\frac{u_xv_x}{b_x},
\qquad
y=\varepsilon_y\frac{u_yv_y}{b_y},
\]
其中 \(\varepsilon_x,\varepsilon_y\in\{\pm1\}\)，\(u_x,u_y,b_x,b_y\) 的全部素因子都属于 \(S\)，而 \(v_x,v_y\) 都与 \(\prod_{p\in S}p\) 互素。则
\[
xy=\varepsilon_x\varepsilon_y\frac{u_xu_y}{b_xb_y}\,v_xv_y.
\]
由于 \(\frac{u_xu_y}{b_xb_y}\in G_S^\times\)，故 \(xy\) 的 \(S\)-外分子恰为 \(v_xv_y\)。再由交换性，
\[
\alpha_x\circ \alpha_y=\alpha_{xy},
\]
于是
\[
\deg(\alpha_x\circ \alpha_y)
=
\deg(\alpha_{xy})
=
v_xv_y
=
\deg(\alpha_x)\deg(\alpha_y).
\]
证毕。
\end{proof}


\begin{theorem}[一般连续自同态的周期点计数与 Artin--Mazur \texorpdfstring{$\zeta$}{zeta} 的算术闭式]\label{thm:xi-localized-solenoid-general-endomorphism-periodic-zeta}
\leanverified{paper\_xi\_localized\_solenoid\_general\_endomorphism\_periodic\_zeta}
沿用定理 \ref{thm:xi-localized-solenoid-endomorphism-kernel-cokernel-sfree-numerator} 的记号，并额外假设
\[
x\notin\{\pm 1\}.
\]
对任意非零整数 \(m\)，记
\[
|m|_{S^\perp}:=\frac{|m|}{\prod_{p\in S}p^{v_p(m)}}.
\]
则对每个 \(n\ge 1\)，不动点群满足
\[
\boxed{\;
\mathrm{Fix}(\alpha_x^n)
\cong
\ZZ/|a^n-b^n|_{S^\perp}\ZZ,
\qquad
\#\mathrm{Fix}(\alpha_x^n)=|a^n-b^n|_{S^\perp}.\;}
\]
因此其 Artin--Mazur \(\zeta\)-函数作为形式幂级数满足
\[
\boxed{\;
\zeta_{\alpha_x}(z)
:=
\exp\!\left(
\sum_{n\ge 1}\frac{\#\mathrm{Fix}(\alpha_x^n)}{n}z^n
\right)
=
\exp\!\left(
\sum_{n\ge 1}\frac{|a^n-b^n|_{S^\perp}}{n}z^n
\right).\;}
\]
\end{theorem}
\begin{proof}
局部化环 \(G_S=\ZZ[S^{-1}]\subset \QQ\) 的有限阶单位只有 \(\pm1\)。故 \(x\notin\{\pm1\}\) 蕴含 \(x^n\neq 1\) 对全部 \(n\ge 1\) 成立，从而下面出现的差分同态都对应非零有理数乘法。

由定义，
\[
\alpha_x^n=\widehat{m_x}^n=\widehat{m_{x^n}}.
\]
因此
\[
\alpha_x^n-\mathrm{id}_{\Sigma_S}
=
\widehat{m_{x^n}-\mathrm{id}_{G_S}}
=
\widehat{m_{x^n-1}}.
\]
于是
\[
\mathrm{Fix}(\alpha_x^n)
=
\ker(\alpha_x^n-\mathrm{id}_{\Sigma_S})
=
\ker(\widehat{m_{x^n-1}}).
\]

再写
\[
x^n-1=\frac{a^n-b^n}{b^n}.
\]
由于 \(b^n\in\langle S\rangle\)，它在 \(G_S\) 中为单位，故 \(x^n-1\) 的 \(S\)-外分子恰为 \(|a^n-b^n|_{S^\perp}\)。将定理 \ref{thm:xi-localized-solenoid-endomorphism-kernel-cokernel-sfree-numerator} 应用于非零元素 \(x^n-1\)，得到
\[
\ker(\widehat{m_{x^n-1}})
\cong
\ZZ/|a^n-b^n|_{S^\perp}\ZZ.
\]
从而所示不动点计数公式成立。

最后把
\[
\#\mathrm{Fix}(\alpha_x^n)=|a^n-b^n|_{S^\perp}
\]
代入 Artin--Mazur \(\zeta\)-函数的定义即可得到所示形式幂级数闭式。证毕。
\end{proof}

\noindent
上述结果把 prime-register 机制推进为一条完整的覆盖--动力学--群论链：\(\Sigma_S=\widehat{G_S}\) 的整数乘法覆盖度、\(n\)-挠子群、全部有限闭子群、整数乘法的周期点计数与 primitive 轨道数，以及任意连续自同态的核、余核与 Artin--Mazur \(\zeta\)-函数，都由 \(S\)-自由剥离操作严格控制。由此，\(\Sigma_S\) 不再只是“带有 \(p\)-进核的 solenoid”这一拓扑模型，而成为一个其有限层对称性、非可逆性与周期动力学都可被 prime-register 完全算术化的对象。


\begin{theorem}[局部化数系有限指数子群与 profinite 完成的完全分类]\label{thm:xi-localized-finite-index-profinite-completion-classification}
\leanverified{paper\_xi\_localized\_finite\_index\_profinite\_completion\_classification}
沿用定理 \ref{thm:xi-localized-basic-extension-strictly-nonsplit} 的记号。对任意子群 \(H\le G_S\)，以下两条等价：
\begin{enumerate}
  \item \(H\) 在 \(G_S\) 中具有有限指数；
  \item 存在唯一正整数 \(m\) 使得对全部 \(p\in S\) 都有 \(p\nmid m\)，并且
  \[
  H=mG_S.
  \]
\end{enumerate}
在此情形下必有
\[
G_S/H\cong \ZZ/m\ZZ.
\]
进一步，若以
\[
\widehat{G_S}^{\,\mathrm{pf}}
\]
表示 \(G_S\) 的 profinite 完成，则
\[
\widehat{G_S}^{\,\mathrm{pf}}
\cong
\varprojlim_{\substack{m\ge 1\\ p\nmid m\ \forall p\in S}}\ZZ/m\ZZ
\cong
\prod_{\ell\notin S}\ZZ_\ell.
\]
\end{theorem}
\begin{proof}
若 \(H=mG_S\) 且对全部 \(p\in S\) 都有 \(p\nmid m\)，则定理 \ref{thm:xi-localized-finite-quotient-phase-layer-classification} 第 \(1\) 条给出
\[
G_S/H=G_S/mG_S\cong \ZZ/m\ZZ,
\]
故 \(H\) 具有有限指数。

反过来，设 \(H\le G_S\) 具有有限指数，并记
\[
Q:=G_S/H.
\]
由定理 \ref{thm:xi-localized-finite-quotient-phase-layer-classification} 第 \(2\) 条，存在唯一正整数 \(m\) 使得对全部 \(p\in S\) 都有 \(p\nmid m\)，且
\[
Q\cong \ZZ/m\ZZ.
\]
由于 \(Q\) 的指数为 \(m\)，任意同态 \(G_S\to Q\) 都杀掉 \(mG_S\)。因此自然商映射 \(G_S\twoheadrightarrow Q\) 经由
\[
G_S/mG_S
\]
因子化。再由定理 \ref{thm:xi-localized-finite-quotient-phase-layer-classification} 第 \(1\) 条，
\[
G_S/mG_S\cong \ZZ/m\ZZ
\]
也恰有 \(m\) 个元素。于是该因子化给出一个从 \(G_S/mG_S\) 到 \(Q\) 的满射，而两者基数同为 \(m\)，故该满射必为同构。其核只能为零，从而
\[
H=mG_S.
\]
唯一性随即由商群阶数的唯一性得到。

由前半部分，\(\{mG_S:\ p\nmid m\ \forall p\in S\}\) 与 \(G_S\) 的全部有限指数子群完全一致，故这些子群构成定义 profinite 完成的共尾系统。于是
\[
\widehat{G_S}^{\,\mathrm{pf}}
\cong
\varprojlim_{\substack{m\ge 1\\ p\nmid m\ \forall p\in S}}G_S/mG_S
\cong
\varprojlim_{\substack{m\ge 1\\ p\nmid m\ \forall p\in S}}\ZZ/m\ZZ.
\]
最后，由中国剩余定理与 profinite 整数的标准分解，
\[
\varprojlim_{\substack{m\ge 1\\ p\nmid m\ \forall p\in S}}\ZZ/m\ZZ
\cong
\prod_{\ell\notin S}\ZZ_\ell.
\]
证毕。
\end{proof}

上述分类表明，Pontryagin 对偶短正合列所显露的 \(S\)-adic 核与 profinite 完成所显露的 \(S^c\)-adic 有限商方向彼此互补；因而 \(G_S\) 的连续圆因子与离散有限商账本在两侧同时被锁定。


\paragraph{局部化 solenoid 的周期指数、熵闭式、子群 \texorpdfstring{$\zeta$}{zeta} 与原始素因子强迫}
在 \(\Sigma_S=\widehat{G_S}\) 的连续自同态、周期点计数以及有限指数子群已经分别得到完全闭式之后，还可把增长指数、拓扑熵、子群 \(\zeta\) 函数与新素数强迫进一步压缩为下列几条纯算术后果。其共同特征是：\(S\)-内素数只改变有限误差、Euler 因子或有限例外，而指数级复杂度完全由分子--分母的最大尺度与 \(S\)-外账本决定。

\begin{theorem}[solenoid 标量自同态的周期点指数由分子--分母最大尺度精确控制]\label{thm:xi-localized-solenoid-periodic-growth-max-scale}
设 \(S\subset\mathbb P\) 为有限素数集，并取
\[
x=\frac{a}{b}\in G_S=\ZZ[S^{-1}]\setminus\{\pm1\},
\qquad
\gcd(a,b)=1,
\qquad
b\in\langle S\rangle.
\]
记
\[
\alpha_x:=\widehat{m_x}:\Sigma_S\to\Sigma_S,
\qquad
F_n(x):=\#\mathrm{Fix}(\alpha_x^n).
\]
则当 \(n\to\infty\) 时有
\[
\boxed{\;
\log F_n(x)
=
n\log \max\{|a|,|b|\}+O(\log n).\;}
\]
从而
\[
\boxed{\;
\lim_{n\to\infty}\frac{1}{n}\log F_n(x)
=
\log \max\{|a|,|b|\}.\;}
\]
\leanverified{paper\_xi\_localized\_solenoid\_periodic\_growth\_max\_scale}
\end{theorem}
\begin{proof}
由定理 \ref{thm:xi-localized-solenoid-general-endomorphism-periodic-zeta}，
\[
F_n(x)=|a^n-b^n|_{S^\perp}.
\]
因此
\[
\log F_n(x)
=
\log|a^n-b^n|-\sum_{p\in S}v_p(a^n-b^n)\log p.
\]

先估计 \(S\)-主部分。固定 \(p\in S\)。
若 \(p\mid ab\)，由于 \(\gcd(a,b)=1\)，\(p\) 只能整除 \(a,b\) 之一，从而
\[
a^n-b^n\equiv \pm b^n \not\equiv 0 \pmod p
\quad\text{或}\quad
a^n-b^n\equiv a^n \not\equiv 0 \pmod p,
\]
故
\[
v_p(a^n-b^n)=0.
\]

以下设 \(p\nmid ab\)。令 \(r_p\) 为 \(ab^{-1}\) 在 \((\ZZ/p\ZZ)^\times\) 中的乘法阶。
若 \(p\mid a^n-b^n\)，则 \(r_p\mid n\)；又因 \(r_p\mid p-1\)，故 \(p\nmid r_p\)。
把 \(n\) 写成
\[
n=r_p\,p^t s,
\qquad
t=v_p(n),
\qquad
p\nmid s.
\]
再令
\[
u_p:=a^{r_p}b^{-r_p}\in \ZZ_p^\times.
\]
由于 \(a^{r_p}\equiv b^{r_p}\pmod p\)，有 \(u_p\equiv 1\pmod p\)。
并且
\[
a^n-b^n
=
b^n\bigl(u_p^{p^t s}-1\bigr),
\]
而 \(b\in \ZZ_p^\times\)，故
\[
v_p(a^n-b^n)=v_p\bigl(u_p^{p^t s}-1\bigr).
\]
又因 \(u_p^{p^t}\equiv 1\pmod p\) 且 \(p\nmid s\)，几何级数
\[
1+u_p^{p^t}+\cdots+\bigl(u_p^{p^t}\bigr)^{s-1}
\equiv s \not\equiv 0 \pmod p
\]
是 \(p\)-进单位，于是
\[
v_p\bigl(u_p^{p^t s}-1\bigr)=v_p\bigl(u_p^{p^t}-1\bigr).
\]

若 \(p\) 为奇素数，写 \(u_p^{p^j}=1+\pi_j\) 且 \(v_p(\pi_j)\ge 1\)。则二项式展开给出
\[
(1+\pi_j)^p-1
=
p\pi_j+\sum_{k=2}^{p-1}\binom{p}{k}\pi_j^k+\pi_j^p,
\]
其中第一项赋值为 \(v_p(\pi_j)+1\)，其余各项赋值严格更大，故
\[
v_p\bigl(u_p^{p^{j+1}}-1\bigr)
=
v_p\bigl(u_p^{p^j}-1\bigr)+1.
\]
迭代可得
\[
v_p\bigl(u_p^{p^t}-1\bigr)=v_p(u_p-1)+t.
\]

若 \(p=2\)，则 \(u_p\) 为奇 \(2\)-进单位，从而 \(u_p^2\equiv 1\pmod 8\)。因此对每个 \(j\ge 1\)，有
\[
v_2\bigl(u_p^{2^{j+1}}-1\bigr)
=
v_2\bigl(u_p^{2^j}-1\bigr)+v_2\bigl(u_p^{2^j}+1\bigr)
=
v_2\bigl(u_p^{2^j}-1\bigr)+1,
\]
因为 \(u_p^{2^j}+1\equiv 2\pmod 8\)。于是存在只依赖于 \(a,b,p\) 的常数 \(C_p\) 使
\[
v_2\bigl(u_p^{2^t}-1\bigr)\le C_p+t.
\]

综上，对每个 \(p\in S\) 都存在常数 \(C_p'\) 使
\[
v_p(a^n-b^n)\le C_p'+v_p(n).
\]
由于 \(S\) 有限，
\[
\sum_{p\in S}v_p(a^n-b^n)\log p
\le
\sum_{p\in S}C_p'\log p+\sum_{p\in S}v_p(n)\log p
=
O(1)+O(\log n).
\]
故
\[
\log F_n(x)=\log|a^n-b^n|+O(\log n).
\]

最后设
\[
M:=\max\{|a|,|b|\},
\qquad
m:=\min\{|a|,|b|\}.
\]
由 \(x\notin\{\pm1\}\) 与 \(\gcd(a,b)=1\) 可知 \(m<M\)，记 \(\theta:=m/M\in(0,1)\)。
于是
\[
|a^n-b^n|
=
M^n\Bigl|1-\varepsilon\,\theta^n\Bigr|
\]
对某个 \(\varepsilon\in\{\pm1\}\) 成立，从而
\[
\log|a^n-b^n|
=
n\log M+O(1).
\]
代回即得
\[
\log F_n(x)=n\log M+O(\log n),
\]
进而极限公式成立。证毕。
\end{proof}

\begin{corollary}[周期点计数级数与 logarithmic 级数的收敛半径同由最大尺度决定]\label{cor:xi-localized-solenoid-periodic-series-radius-max-scale}
沿用定理 \ref{thm:xi-localized-solenoid-periodic-growth-max-scale} 的记号，并记
\[
M:=\max\{|a|,|b|\}.
\]
定义两类周期点级数
\[
\mathcal P_x(z):=\sum_{n\ge 1}F_n(x)\,z^n,
\qquad
\mathcal L_x(z):=\sum_{n\ge 1}\frac{F_n(x)}{n}\,z^n.
\]
则二者的收敛半径都恰为
\[
\boxed{\;
R_x=\frac{1}{M}.\;}
\]
\end{corollary}
\begin{proof}
由定理 \ref{thm:xi-localized-solenoid-periodic-growth-max-scale}，
\[
\log F_n(x)=n\log M+O(\log n).
\]
因此
\[
\limsup_{n\to\infty}F_n(x)^{1/n}=M.
\]
又因为
\[
\left(\frac{F_n(x)}{n}\right)^{1/n}
=
F_n(x)^{1/n}\,n^{-1/n},
\]
而 \(n^{-1/n}\to 1\)，故
\[
\limsup_{n\to\infty}\left(\frac{F_n(x)}{n}\right)^{1/n}=M.
\]
对 \(\mathcal P_x\) 与 \(\mathcal L_x\) 分别应用 Cauchy--Hadamard 公式，即得二者的收敛半径都等于 \(M^{-1}\)。证毕。
\end{proof}

\begin{corollary}[单位情形的拓扑熵仅由有理尺度 \texorpdfstring{$\max\{|a|,|b|\}$}{max\{|a|,|b|\}} 决定]\label{cor:xi-localized-solenoid-unit-entropy-max-scale}
在定理 \ref{thm:xi-localized-solenoid-periodic-growth-max-scale} 的设定下，若 \(x\in G_S^\times\)，则 \(\alpha_x\) 为连续自同构，并且
\[
\boxed{\;
h_{\mathrm{top}}(\alpha_x)=\log\max\{|a|,|b|\}.\;}
\]
\end{corollary}
\begin{proof}
由定理 \ref{thm:xi-localized-solenoid-endomorphism-kernel-cokernel-sfree-numerator} 第 \(3\) 条，\(x\in G_S^\times\) 当且仅当 \(\alpha_x\) 为同构。对 solenoid 自同构，Lind--Ward 的 \(p\)-进熵分解给出
\[
h_{\mathrm{top}}(\alpha_x)
=
\sum_{v\le \infty}\max\{\log|x|_v,0\},
\]
其中 \(v\) 遍历 \(\QQ\) 的 Archimedean 与非 Archimedean 赋值 \cite{LindWard1988SolenoidsPadicEntropy}。

若 \(|a|\ge |b|\)，则
\[
\max\{\log|x|_\infty,0\}=\log\frac{|a|}{|b|}.
\]
又因 \(\gcd(a,b)=1\)，对任意素数 \(p\) 有
\[
\max\{\log|x|_p,0\}=
\begin{cases}
v_p(b)\log p,& p\mid b,\\
0,& p\nmid b.
\end{cases}
\]
故
\[
h_{\mathrm{top}}(\alpha_x)
=
\log\frac{|a|}{|b|}
+\sum_{p\mid b}v_p(b)\log p
=
\log|a|.
\]

若 \(|a|<|b|\)，则 Archimedean 贡献为 \(0\)，并且只有 \(p\mid b\) 时才有正的 \(p\)-进贡献，因此
\[
h_{\mathrm{top}}(\alpha_x)
=
\sum_{p\mid b}v_p(b)\log p
=
\log|b|.
\]
两种情形合并即得
\[
h_{\mathrm{top}}(\alpha_x)=\log\max\{|a|,|b|\}.
\]
证毕。
\end{proof}

\begin{theorem}[局部化数系有限指数子群的 \texorpdfstring{$\zeta$}{zeta} 函数闭式]\label{thm:xi-localized-subgroup-zeta-euler-product}
\leanverified{paper\_xi\_localized\_subgroup\_zeta\_euler\_product}
设
\[
P_S:=\prod_{p\in S}p.
\]
则 \(G_S=\ZZ[S^{-1}]\) 的有限指数子群恰为
\[
H=mG_S,
\qquad
m\ge 1,
\qquad
\gcd(m,P_S)=1,
\]
并且该表示唯一。在此基础上，其子群 \(\zeta\) 函数
\[
\zeta_{G_S}^{\le}(s)
:=
\sum_{H\le_f G_S}[G_S:H]^{-s}
\]
在 \(\Re(s)>1\) 上满足精确 Euler 乘积
\[
\boxed{\;
\zeta_{G_S}^{\le}(s)
=
\sum_{\substack{m\ge 1\\ (m,P_S)=1}}m^{-s}
=
\zeta(s)\prod_{p\in S}(1-p^{-s}).\;}
\]
\end{theorem}
\begin{proof}
有限指数子群的完全分类已由定理 \ref{thm:xi-localized-finite-index-profinite-completion-classification} 给出：\(H\le G_S\) 具有有限指数，当且仅当
\[
H=mG_S
\]
且 \(\gcd(m,P_S)=1\)，并且该 \(m\) 唯一。又由同一定理，
\[
G_S/mG_S\cong \ZZ/m\ZZ,
\]
故
\[
[G_S:mG_S]=m.
\]
因此
\[
\zeta_{G_S}^{\le}(s)
=
\sum_{\substack{m\ge 1\\ (m,P_S)=1}}m^{-s}.
\]
在 \(\Re(s)>1\) 上按素数分解 Euler 乘积，
\[
\sum_{\substack{m\ge 1\\ (m,P_S)=1}}m^{-s}
=
\prod_{p\notin S}(1-p^{-s})^{-1}
=
\zeta(s)\prod_{p\in S}(1-p^{-s}).
\]
证毕。
\end{proof}

\begin{corollary}[有限指数子群计数的主项、显式误差与极点留数]\label{cor:xi-localized-subgroup-counting-main-term}
沿用定理 \ref{thm:xi-localized-subgroup-zeta-euler-product} 的记号。定义
\[
N_S(X)
:=
\#\{H\le_f G_S:\ [G_S:H]\le X\}.
\]
则对任意 \(X\ge 1\)，有精确公式
\[
N_S(X)
=
\#\{1\le m\le X:\ \gcd(m,P_S)=1\}
=
\sum_{d\mid P_S}\mu(d)\left\lfloor\frac{X}{d}\right\rfloor.
\]
从而
\[
\boxed{\;
N_S(X)
=
\frac{\varphi(P_S)}{P_S}\,X+O(2^{|S|}),
\qquad
X\to\infty.\;}
\]
并且
\[
\boxed{\;
\Res_{s=1}\zeta_{G_S}^{\le}(s)
=
\frac{\varphi(P_S)}{P_S}.\;}
\]
\end{corollary}
\begin{proof}
由定理 \ref{thm:xi-localized-subgroup-zeta-euler-product}，
\[
N_S(X)=\#\{1\le m\le X:\ \gcd(m,P_S)=1\}.
\]
对互素条件施行 Möbius 反演，
\[
\mathbf 1_{\gcd(m,P_S)=1}
=
\sum_{d\mid \gcd(m,P_S)}\mu(d)
=
\sum_{\substack{d\mid P_S\\ d\mid m}}\mu(d).
\]
对 \(m\le X\) 求和并交换顺序，得到
\[
N_S(X)
=
\sum_{d\mid P_S}\mu(d)\sum_{\substack{m\le X\\ d\mid m}}1
=
\sum_{d\mid P_S}\mu(d)\left\lfloor\frac{X}{d}\right\rfloor.
\]
再用
\[
\left\lfloor\frac{X}{d}\right\rfloor=\frac{X}{d}+O(1)
\]
得到
\[
N_S(X)
=
X\sum_{d\mid P_S}\frac{\mu(d)}{d}
+O\!\left(\sum_{d\mid P_S}1\right).
\]
由于 \(P_S\) 为平方自由数，
\[
\sum_{d\mid P_S}1=2^{|S|},
\qquad
\sum_{d\mid P_S}\frac{\mu(d)}{d}
=
\prod_{p\in S}\left(1-\frac1p\right)
=
\frac{\varphi(P_S)}{P_S},
\]
故主项展开成立。

留数公式由定理 \ref{thm:xi-localized-subgroup-zeta-euler-product} 的 Euler 乘积与 \(\zeta(s)\) 在 \(s=1\) 的单极点立即得到。证毕。
\end{proof}

\begin{theorem}[周期点账本的 \texorpdfstring{$S$}{S}-外原始素因子强迫]\label{thm:xi-localized-solenoid-periodic-zsigmondy-forcing}
在定理 \ref{thm:xi-localized-solenoid-periodic-growth-max-scale} 的设定下，存在有限集合
\[
E=E(a,b,S)\subset \NN
\]
使得对每个 \(n\notin E\)，都存在素数
\[
q_n\notin S
\]
满足
\[
q_n\mid F_n(x),
\qquad
q_n\nmid F_k(x)\quad(1\le k<n).
\]
换言之，周期点计数序列 \(\{F_n(x)\}_{n\ge 1}\) 在 \(S\)-外素数方向上具有原始素因子强迫性。
\end{theorem}
\begin{proof}
记
\[
U_n:=\frac{a^n-b^n}{a-b}\qquad(n\ge 1).
\]
由于 \(a/b\notin\{\pm1\}\)，序列 \((U_n)\) 是非退化 Lucas 型序列。Bilu--Hanrot--Voutier 关于 Lucas 与 Lehmer 序列 primitive divisor 的定理表明：存在有限集合 \(E_0=E_0(a,b)\subset\NN\)，使得对每个 \(n\notin E_0\)，\(U_n\) 都存在 primitive prime divisor \(q_n\) \cite{BiluHanrotVoutier2001PrimitiveDivisorsLucas}。

对这样的 \(n\)，由 primitive 性可知
\[
q_n\mid U_n,
\qquad
q_n\nmid (a-b)U_1U_2\cdots U_{n-1}.
\]
因此
\[
q_n\mid a^n-b^n,
\qquad
q_n\nmid a^k-b^k\quad(1\le k<n),
\]
即 \(q_n\) 也是 \(a^n-b^n\) 的 primitive prime divisor。

现证明当 \(n\) 充分大时 \(q_n\notin S\)。若反设 \(q_n\in S\)，则上式自动给出 \(q_n\nmid ab\)。于是 \(ab^{-1}\) 在 \((\ZZ/q_n\ZZ)^\times\) 中有定义，并且 primitive 性推出其乘法阶恰为 \(n\)。故
\[
n\mid (q_n-1).
\]
记
\[
M_S:=
\begin{cases}
0,& S=\varnothing,\\
\max_{q\in S}(q-1),& S\neq\varnothing.
\end{cases}
\]
则当 \(n>M_S\) 时，上述整除关系不可能成立。于是对全部
\[
n\notin E_1:=\{1,\dots,M_S\},
\]
primitive prime divisor 必不属于 \(S\)。

最后令
\[
E:=E_0\cup E_1.
\]
对任意 \(n\notin E\)，已有 \(q_n\notin S\) 且 \(q_n\mid a^n-b^n\)。再由定理 \ref{thm:xi-localized-solenoid-general-endomorphism-periodic-zeta}，
\[
F_n(x)=|a^n-b^n|_{S^\perp},
\]
所以
\[
q_n\mid F_n(x).
\]
另一方面，若对某个 \(k<n\) 有 \(q_n\mid F_k(x)\)，由于 \(q_n\notin S\)，便必有 \(q_n\mid a^k-b^k\)，这与 \(q_n\) 的 primitive 性矛盾。故
\[
q_n\nmid F_k(x)\qquad(1\le k<n).
\]
证毕。
\end{proof}

\noindent
这组结果把 \(G_S\) 的有限指数层、\(\Sigma_S\) 的周期层以及单位自同态的熵层压缩到同一条可计算账本上：指数主项由 \(\max\{|a|,|b|\}\) 给出，有限指数复杂度由删去 \(S\)-Euler 因子的整数计数控制，而足够高的周期层则必然在 \(S\)-外方向引入新的原始素因子。


\paragraph{局部化有限指数子群格、补素完成与全素数账本的闭合}
在有限指数子群的完全分类、相应子群 \(\zeta\) 函数的 Euler 乘积以及局部化对偶短正合列的严格不分裂都已建立之后，同一 rank-one 局部化数系的有限指数算术还可进一步压缩为下列三条 lattice/profinite 后果。

\begin{theorem}[局部化数系有限指数子群格的指标保持分类]\label{thm:xi-time-part56e-localized-finite-index-lattice-classification}
设 \(S\subset \mathbb P\) 为有限素数集，并记
\[
P_S:=\prod_{p\in S}p.
\]
则映射
\[
n\longmapsto nG_S
\qquad
\bigl(n\ge 1,\ \gcd(n,P_S)=1\bigr)
\]
在满足 \(\gcd(n,P_S)=1\) 的正整数与 \(G_S\) 的全部有限指数子群之间给出双射，并且保持指标：
\[
[G_S:nG_S]=n.
\]
等价地，\(\operatorname{Sub}_f(G_S)\) 与 \(S\)-自由正整数按整除关系形成的偏序格反向同构。
\end{theorem}
\begin{proof}
定理 \ref{thm:xi-localized-finite-index-profinite-completion-classification} 已证明：\(H\le G_S\) 具有有限指数，当且仅当存在唯一正整数 \(n\) 使
\[
H=nG_S,
\qquad
\gcd(n,P_S)=1,
\]
并且
\[
G_S/H\cong \ZZ/n\ZZ.
\]
因此
\[
[G_S:H]=n=[G_S:nG_S].
\]
双射与指标保持性遂同时成立。最后，由
\[
nG_S\subseteq mG_S
\iff
m\mid n
\]
可知该对应把子群包含关系转化为整数整除关系的反向，故得到所述反向同构。证毕。
\end{proof}

\leanverified{paper\_xi\_time\_part56e\_localized\_finite\_index\_lattice\_gcd\_lcm\_mobius}
\begin{corollary}[局部化有限指数子群格的 \texorpdfstring{$\gcd/\mathrm{lcm}$}{gcd/lcm} 算术与区间 Möbius 闭式]\label{cor:xi-time-part56e-localized-finite-index-lattice-gcd-lcm-mobius}
沿用定理 \ref{thm:xi-time-part56e-localized-finite-index-lattice-classification} 的记号。对任意满足
\[
\gcd(m,P_S)=\gcd(n,P_S)=1
\]
的正整数，记
\[
H_m:=mG_S,
\qquad
H_n:=nG_S.
\]
则在有限指数子群格 \(\operatorname{Sub}_f(G_S)\) 中有
\[
H_m\wedge H_n
=
H_m\cap H_n
=
\operatorname{lcm}(m,n)\,G_S,
\]
\[
H_m\vee H_n
=
H_m+H_n
=
\gcd(m,n)\,G_S.
\]
进一步，若 \(m\mid n\) 且 \(\gcd(n,P_S)=1\)，则区间 Möbius 函数满足
\[
\mu_{\operatorname{Sub}_f(G_S)}(H_n,H_m)=\mu(n/m),
\]
其中右端为经典数论 Möbius 函数。
\end{corollary}
\begin{proof}
由定理 \ref{thm:xi-time-part56e-localized-finite-index-lattice-classification}，
\(\operatorname{Sub}_f(G_S)\) 与 \(S\)-自由正整数的整除格反向同构。故其 meet 与 join 分别对应于整数侧的 \(\operatorname{lcm}\) 与 \(\gcd\)，从而得到
\[
H_m\cap H_n=\operatorname{lcm}(m,n)\,G_S,
\qquad
H_m+H_n=\gcd(m,n)\,G_S.
\]
若愿意直接验证第二式，只需记 \(d:=\gcd(m,n)\)。由 Bézout 恒等式可取 \(u,v\in\ZZ\) 使
\[
um+vn=d,
\]
故 \(dG_S\subseteq mG_S+nG_S\)；反向包含由 \(d\mid m,n\) 立即成立。

再设 \(m\mid n\)。则
\[
H_n\subseteq H_m,
\]
且区间 \([H_n,H_m]\) 与正整数除子格
\[
\{\,d:\ d\mid n/m\,\}
\]
反向同构。有限除子格的 Möbius 函数正是经典 Möbius 函数 \(\mu\)，故
\[
\mu_{\operatorname{Sub}_f(G_S)}(H_n,H_m)=\mu(n/m).
\]
证毕。
\end{proof}

\begin{theorem}[局部化数系的子群 \texorpdfstring{$\zeta$}{zeta}、profinite 完成与补素核的全素数分裂]\label{thm:xi-time-part56e-localized-subgroup-zeta-profinite-complement-splitting}
设 \(S\subset\mathbb P\) 为有限素数集，并记
\[
\pi_S:\widehat{G_S}\twoheadrightarrow \TT,
\qquad
K_S:=\ker(\pi_S).
\]
则同时有
\[
\zeta_{G_S}^{\le}(s)
=
\zeta(s)\prod_{p\in S}(1-p^{-s}),
\]
\[
\widehat{G_S}^{\,\mathrm{pf}}
\cong
\prod_{\ell\notin S}\ZZ_\ell,
\qquad
K_S
\cong
\prod_{p\in S}\ZZ_p,
\]
从而得到规范拓扑群同构
\[
\widehat{G_S}^{\,\mathrm{pf}}\times K_S
\cong
\prod_{p\in\mathbb P}\ZZ_p
=
\widehat{\ZZ}^{\,\mathrm{pf}}.
\]
等价地，\(G_S\) 的有限商方向与其 Pontryagin 对偶上的全不连通核方向恰按 \(S^c\)-素数与 \(S\)-素数互补分裂，并在合并后恢复完整的全素数 profinite 账本。
\end{theorem}
\begin{proof}
第一式正是定理 \ref{thm:xi-localized-subgroup-zeta-euler-product}。
第二式中的 profinite 完成分类由定理
\ref{thm:xi-localized-finite-index-profinite-completion-classification}
给出：
\[
\widehat{G_S}^{\,\mathrm{pf}}
\cong
\prod_{\ell\notin S}\ZZ_\ell.
\]
另一方面，定理 \ref{thm:xi-localized-basic-extension-strictly-nonsplit}
给出 Pontryagin 对偶短正合列
\[
0\longrightarrow \prod_{p\in S}\ZZ_p
\longrightarrow \widehat{G_S}
\xrightarrow{\ \pi_S\ }
\TT
\longrightarrow 0,
\]
故
\[
K_S\cong \prod_{p\in S}\ZZ_p.
\]
将两式相乘即可得到
\[
\widehat{G_S}^{\,\mathrm{pf}}\times K_S
\cong
\left(\prod_{\ell\notin S}\ZZ_\ell\right)\times
\left(\prod_{p\in S}\ZZ_p\right)
\cong
\prod_{p\in\mathbb P}\ZZ_p.
\]
证毕。
\end{proof}

\paragraph{有限素数支撑与紧核族之间的 Galois 型偏序恢复}
在有限指数子群格、profinite 完成以及补素核的全素数分裂都已闭合之后，有限素数支撑 \(S\subset\mathbb P\) 与紧全不连通核
\[
K_S\cong \prod_{p\in S}\ZZ_p
\]
之间还存在一条更强的偏序对应：不仅每个有限素数集可由核本身恢复，而且包含关系也可由核间的可商性精确读取。

\begin{theorem}[有限素数支撑的 Galois 型对应]\label{thm:xi-time-part56ea-prime-support-kernel-galois-correspondence}
\leanverified{paper\_xi\_time\_part56ea\_prime\_support\_kernel\_galois\_correspondence}
对任意有限素数集 \(S,T\subset\mathbb P\)，记
\[
G_S:=\ZZ[S^{-1}],
\qquad
K_S:=\ker(\widehat{G_S}\to \TT).
\]
则下列命题严格等价：
\begin{enumerate}
  \item \(S\subseteq T\)；
  \item 存在连续开满射
  \[
  \pi_{T,S}:K_T\to K_S;
  \]
  \item 存在短正合列
  \[
  0\longrightarrow \prod_{p\in T\setminus S}\ZZ_p
  \longrightarrow K_T
  \xrightarrow{\ \pi_{T,S}\ }
  K_S
  \longrightarrow 0.
  \]
\end{enumerate}
特别地，有限素数集按包含形成的偏序与紧核族 \(\{K_S\}\) 按连续开商映射形成的偏序完全一致。
\end{theorem}
\begin{proof}
由定理~\ref{thm:xi-time-part56e-localized-subgroup-zeta-profinite-complement-splitting}，
\[
K_S\cong \prod_{p\in S}\ZZ_p,
\qquad
K_T\cong \prod_{p\in T}\ZZ_p.
\]
若 \(S\subseteq T\)，则坐标投影
\[
\prod_{p\in T}\ZZ_p\longrightarrow \prod_{p\in S}\ZZ_p
\]
给出连续开满射，其核正是
\[
\prod_{p\in T\setminus S}\ZZ_p.
\]
故 \((1)\Rightarrow(2)\Rightarrow(3)\)。

反过来，设存在连续开满射
\[
\pi_{T,S}:K_T\to K_S.
\]
对任意素数 \(p\)，取模 \(p\) 商得到诱导满射
\[
K_T/pK_T\longrightarrow K_S/pK_S.
\]
而由
\[
K_U\cong \prod_{\ell\in U}\ZZ_\ell
\]
可直接计算
\[
K_U/pK_U\cong
\begin{cases}
\FF_p, & p\in U,\\
0, & p\notin U.
\end{cases}
\]
因此，若 \(p\in S\)，则 \(K_S/pK_S\neq 0\)，从而满射迫使 \(K_T/pK_T\neq 0\)，故 \(p\in T\)。
这就证明了 \(S\subseteq T\)，即 \((2)\Rightarrow(1)\)。
至于 \((3)\Rightarrow(2)\) 则为正合列定义的直接内容。证毕。
\end{proof}


\paragraph{局部化数系的有限商、torsion 谱与 Dirichlet 账本}
在有限商分类、solenoid 的整数乘法核度公式与有限指数子群 Euler 乘积已经确立之后，局部化数系的 quotient 层与 torsion 层还可进一步压缩为同一组 Dirichlet 账本。

\begin{theorem}[局部化数系有限商的指数账本与循环实现判据]\label{thm:xi-localized-quotient-index-ledger}
沿用前述记号。对任意正整数 \(n\)，记
\[
n_{S^\perp}:=\frac{n}{\prod_{p\in S}p^{v_p(n)}}.
\]
则有
\[
G_S/nG_S\cong \ZZ/n_{S^\perp}\ZZ,
\qquad
[G_S:nG_S]=n_{S^\perp}.
\]
特别地，有限循环群 \(\ZZ/m\ZZ\) 出现在 \(G_S\) 的有限商中，当且仅当 \(m\) 的全部素因子都不属于 \(S\)。
\end{theorem}
\begin{proof}
定理 \ref{thm:xi-localized-finite-quotient-phase-layer-classification} 已给出
\[
G_S/nG_S\cong \ZZ/n_{S^\perp}\ZZ.
\]
两边取基数即得指数公式。最后一句正是同一定理对有限商分类的重述。证毕。
\end{proof}

\begin{theorem}[Pontryagin 对偶的 \texorpdfstring{$n$}{n}-挠谱与 exact-order 计数]\label{thm:xi-localized-torsion-exact-order-ledger}
沿用上述记号，并记 \(\Sigma_S:=\widehat{G_S}\)。对任意正整数 \(n\)，有
\[
\Sigma_S[n]\cong \ZZ/n_{S^\perp}\ZZ\cong \mu_{n_{S^\perp}},
\qquad
|\Sigma_S[n]|=n_{S^\perp}.
\]
其中 \(\mu_m:=\{z\in\TT:\ z^m=1\}\)。更进一步，恰为 \(n\) 阶的元素个数满足
\[
\#\{x\in\Sigma_S:\ \operatorname{ord}(x)=n\}
=
\begin{cases}
\varphi(n),& \text{\(n\) 的全部素因子都不属于 }S,\\[1mm]
0,& \text{否则}.
\end{cases}
\]
\end{theorem}
\begin{proof}
由定理 \ref{thm:xi-localized-solenoid-finite-subgroups-are-kernels}，
\[
\Sigma_S[n]=\ker[n]\cong \ZZ/n_{S^\perp}\ZZ.
\]
有限循环群与 \(\mu_{n_{S^\perp}}\) 同构，故第一式成立，取基数即得第二式。
又因 \(\Sigma_S[n]\) 为阶 \(n_{S^\perp}\) 的循环群，其中恰为 \(n\) 阶的元素个数在 \(n\mid n_{S^\perp}\) 时为 \(\varphi(n)\)，否则为 \(0\)。而 \(n\mid n_{S^\perp}\) 当且仅当 \(n\) 的全部素因子都不属于 \(S\)。证毕。
\end{proof}

\begin{theorem}[quotient-zeta 与 torsion-zeta 的 Euler 乘积闭式]\label{thm:xi-localized-quotient-torsion-zeta-euler-product}
沿用上述记号，定义 Dirichlet 级数
\[
Z_S^{\mathrm{quot}}(s):=\sum_{n\ge 1}\frac{[G_S:nG_S]}{n^s},
\qquad
Z_S^{\mathrm{tor}}(s):=\sum_{n\ge 1}\frac{|\Sigma_S[n]|}{n^s}.
\]
则二者在 \(\Re(s)>2\) 上绝对收敛，并满足
\[
Z_S^{\mathrm{quot}}(s)=Z_S^{\mathrm{tor}}(s)
=
\prod_{p\in S}\frac{1}{1-p^{-s}}
\prod_{p\notin S}\frac{1}{1-p^{1-s}}
=
\zeta(s-1)\prod_{p\in S}\frac{1-p^{1-s}}{1-p^{-s}}.
\]
若再定义 exact-order Dirichlet 级数
\[
E_S(s):=\sum_{n\ge 1}\frac{\#\{x\in\Sigma_S:\ \operatorname{ord}(x)=n\}}{n^s},
\]
则在 \(\Re(s)>2\) 上有
\[
E_S(s)
=
\frac{\zeta(s-1)}{\zeta(s)}
\prod_{p\in S}\frac{1-p^{1-s}}{1-p^{-s}}.
\]
\leanverified{paper\_xi\_torsion\_exact\_order\_ledger}
\end{theorem}
\begin{proof}
由定理 \ref{thm:xi-localized-quotient-index-ledger} 与定理
\ref{thm:xi-localized-torsion-exact-order-ledger}，
\[
[G_S:nG_S]=n_{S^\perp}=|\Sigma_S[n]|.
\]
故两条 Dirichlet 级数逐项相等。
又因
\[
0\le n_{S^\perp}\le n,
\]
故当 \(\sigma:=\Re(s)>2\) 时，
\[
\sum_{n\ge 1}\left|\frac{[G_S:nG_S]}{n^s}\right|
=
\sum_{n\ge 1}\frac{n_{S^\perp}}{n^\sigma}
\le
\sum_{n\ge 1}n^{1-\sigma}<\infty,
\]
从而 \(Z_S^{\mathrm{quot}}\) 与 \(Z_S^{\mathrm{tor}}\) 都绝对收敛。

函数 \(n\mapsto n_{S^\perp}\) 为乘法函数，且
\[
p^e\longmapsto
\begin{cases}
1,& p\in S,\\[1mm]
p^e,& p\notin S.
\end{cases}
\]
因此 Euler 因子分别为
\[
\sum_{e\ge 0}\frac{1}{p^{es}}=\frac{1}{1-p^{-s}}
\qquad (p\in S),
\]
\[
\sum_{e\ge 0}\frac{p^e}{p^{es}}=\frac{1}{1-p^{1-s}}
\qquad (p\notin S),
\]
从而得到所示 Euler 乘积；再把 \(\prod_{p\notin S}(1-p^{1-s})^{-1}\) 写成 \(\zeta(s-1)\) 去掉 \(S\)-内 Euler 因子后的形式，即得第二个闭式。

对 \(E_S(s)\)，由定理 \ref{thm:xi-localized-torsion-exact-order-ledger}，其系数函数为
\[
e_S(n):=\#\{x\in\Sigma_S:\ \operatorname{ord}(x)=n\}
=
\begin{cases}
\varphi(n),& \text{\(n\) 的全部素因子都不属于 }S,\\[1mm]
0,& \text{否则}.
\end{cases}
\]
又因 \(0\le e_S(n)\le \varphi(n)\le n\)，故 \(E_S(s)\) 在 \(\Re(s)>2\) 上亦绝对收敛。
故 \(e_S\) 为乘法函数，且对应 Euler 因子为
\[
1
\qquad (p\in S),
\]
以及
\[
1+\sum_{e\ge 1}\frac{\varphi(p^e)}{p^{es}}
=
1+\sum_{e\ge 1}\frac{p^e-p^{e-1}}{p^{es}}
=
\frac{1-p^{-s}}{1-p^{1-s}}
\qquad (p\notin S).
\]
将全部素数相乘，并把 \(p\notin S\) 的乘积重写为
\[
\prod_{p}\frac{1-p^{-s}}{1-p^{1-s}}
\prod_{p\in S}\frac{1-p^{1-s}}{1-p^{-s}}
=
\frac{\zeta(s-1)}{\zeta(s)}
\prod_{p\in S}\frac{1-p^{1-s}}{1-p^{-s}},
\]
便得所示公式。证毕。
\end{proof}

\begin{corollary}[有限商指数函数在素数点上的完全恢复]\label{cor:xi-localized-quotient-index-prime-recovery}
定义
\[
I_S(n):=[G_S:nG_S].
\]
则对每个 \(n\ge 1\) 都有
\[
I_S(n)=n_{S^\perp}.
\]
特别地，对任意素数 \(p\) 有
\[
p\in S \iff I_S(p)=1,
\qquad
p\notin S \iff I_S(p)=p.
\]
因此，仅由素数点取值模式 \(p\mapsto I_S(p)\) 即可完全且唯一恢复 \(S\)。
\end{corollary}
\begin{proof}
第一式是定理 \ref{thm:xi-localized-quotient-index-ledger} 的指数公式。对 \(n=p\) 代入即得第二式，故结论成立。证毕。
\end{proof}

\begin{theorem}[连通 Lie 商的圆周刚性]\label{thm:xi-localized-connected-lie-quotient-circle-rigidity}
沿用定理 \ref{thm:xi-localized-basic-extension-strictly-nonsplit} 的记号。设 \(K\) 是 \(\widehat{G_S}\) 的连通紧 Lie 商。则必有
\[
K\cong \TT^r
\qquad\text{且}\qquad
r\le 1.
\]
特别地，任何非平凡连通紧 Lie 商都同构于 \(\TT\)。
\end{theorem}
\begin{proof}
连通紧 Lie 交换群必同构于某个环面 \(\TT^r\)。设
\[
\pi:\widehat{G_S}\twoheadrightarrow K
\]
为满射，则对偶化得到单射
\[
\widehat K\hookrightarrow G_S.
\]
而
\[
\widehat K\cong \ZZ^r.
\]
由于 \(G_S=\ZZ[S^{-1}]\subset \QQ\)，有
\[
G_S\otimes_{\ZZ}\QQ\cong \QQ,
\]
故 \(G_S\) 的有理秩为 \(1\)。于是其任意自由阿贝尔子群的秩至多为 \(1\)，从而 \(r\le 1\)。若 \(K\) 非平凡，则 \(r\neq 0\)，故必有 \(r=1\)，即
\[
K\cong \TT.
\]
证毕。
\end{proof}

\begin{theorem}[非平凡局部化数系与有限生成离散交换群的相位采样不可混同]\label{thm:xi-localized-phase-sampling-not-finitely-generated}
\leanverified{paper\_xi\_localized\_phase\_sampling\_not\_finitely\_generated}
设 \(S\subset\mathbb P\) 为非空有限集。则不存在有限生成离散交换群 \(H\) 使得
\[
\card{\operatorname{Hom}(H,\ZZ/N\ZZ)}
=
\Sigma^{\mathrm{ph}}_{G_S}(N)
\qquad(\forall\,N\ge 1).
\]
\end{theorem}
\begin{proof}
反设存在
\(
H\cong \ZZ^{r}\oplus F
\)
满足上述恒等式。对任意素数 \(p\) 与 \(a\ge 1\)，由注 \ref{rem:xi-cdim-hom-quotient-consistency}，
\[
\card{\operatorname{Hom}(H,\ZZ/p^a\ZZ)}=\card{H/p^aH}.
\]
再由命题 \ref{prop:xi-cdim-local-inversion-delta}，定义
\[
c_p(a):=
v_p\!\bigl(\card{H/p^aH}\bigr)-ar
=
\sum_{j=1}^{t_p}\min(\lambda_{p,j},a),
\]
其中 \(F_p\cong\bigoplus_{j=1}^{t_p}\ZZ/p^{\lambda_{p,j}}\ZZ\) 为 \(F\) 的 \(p\)-primary 分量。因而 \(c_p(a)\) 对 \(a\) 有界，从而
\[
\lim_{a\to\infty}\frac1a
v_p\!\bigl(\card{\operatorname{Hom}(H,\ZZ/p^a\ZZ)}\bigr)
=
r
\qquad\text{对每个素数 }p\text{ 都成立}.
\]
另一方面，由定理 \ref{thm:xi-localized-phase-sampling-closed-form}，
\[
\Sigma^{\mathrm{ph}}_{G_S}(p^a)=
\begin{cases}
1,& p\in S,\\[1mm]
p^a,& p\notin S,
\end{cases}
\]
故
\[
\frac1a v_p\!\bigl(\Sigma^{\mathrm{ph}}_{G_S}(p^a)\bigr)
=
\begin{cases}
0,& p\in S,\\[1mm]
1,& p\notin S.
\end{cases}
\]
由于 \(S\) 非空且有限，可同时取到 \(p\in S\) 与 \(q\notin S\)，从而上述极限在素数间取两个不同值 \(0\) 与 \(1\)。这与有限生成离散交换群在全部素数上共享同一秩斜率 \(r\) 矛盾。故不存在这样的 \(H\)。证毕。
\end{proof}

\paragraph{有限值连续 prime-audit 的圆柱因子化与指数级歧义}
局部化数系的 Pontryagin 对偶已把 prime-register 核显式识别为有限个 \(p\)-进坐标的乘积；在此基础上，任何有限值连续审计都必定降到有限精度的圆柱商。

\begin{theorem}[有限值连续 prime-audit 必因子化到有限圆柱商]\label{thm:xi-localized-finite-prime-audit-cylinder-factorization}
设 \(S\subset\mathbb P\) 为有限素数集，并记
\[
\mathfrak K_S:=\prod_{p\in S}\ZZ_p.
\]
由定理 \ref{thm:killo-localization-circle-dimension-prime-ledger-splitting}，\(\mathfrak K_S\) 正是 \(\widehat{G_S}\) 的 prime-register 核。令 \(A\) 为有限离散集合，\(F:\mathfrak K_S\to A\) 为连续映射。则存在子集 \(T\subset S\) 与整数族 \((n_p)_{p\in T}\)（其中 \(n_p\ge 1\)），使得 \(F\) 因子化为
\[
\mathfrak K_S
\longrightarrow
\prod_{p\in T}\ZZ/p^{n_p}\ZZ
\xrightarrow{\ \overline F\ }
A.
\]
换言之，任何有限输出的连续 prime-audit 都只依赖有限精度的 \(p\)-进截断；在有限输出与连续性同时成立时，不存在 genuinely 全精度的 prime-audit。
\end{theorem}
\begin{proof}
\(\mathfrak K_S\) 是紧全不连通 Hausdorff 群，而 \(A\) 为有限离散空间。故对每个 \(a\in A\)，纤维
\[
F^{-1}(a)
\]
都是紧开集。

取任意 \(a\in A\) 与任意 \(x=(x_p)_{p\in S}\in F^{-1}(a)\)。由乘积拓扑的基可知，存在整数族
\[
n_{x,p}\in \ZZ_{\ge 0}\qquad(p\in S)
\]
使得圆柱邻域
\[
U_x:=\prod_{p\in S}\bigl(x_p+p^{n_{x,p}}\ZZ_p\bigr)
\]
满足
\[
x\in U_x\subseteq F^{-1}(a).
\]
由于 \(F^{-1}(a)\) 紧，可从 \(\{U_x\}_{x\in F^{-1}(a)}\) 中抽取有限子覆盖
\[
F^{-1}(a)\subseteq U_{x^{(a,1)}}\cup\cdots\cup U_{x^{(a,r_a)}}.
\]
对全部 \(a\in A\) 与全部所选圆柱，把每个素数坐标上的指数取最大值：
\[
n_p:=\max_{a\in A}\max_{1\le j\le r_a} n_{x^{(a,j)},p}
\qquad(p\in S).
\]
再记
\[
T:=\{p\in S:\ n_p\ge 1\}.
\]
现设 \(y,z\in\mathfrak K_S\) 满足
\[
y\equiv z \pmod{p^{n_p}\ZZ_p}
\qquad(\forall\,p\in S).
\]
若 \(F(y)=a\)，则 \(y\in U_{x^{(a,j)}}\) 对某个 \(j\) 成立。由于
\[
n_p\ge n_{x^{(a,j)},p}
\qquad(\forall\,p\in S),
\]
并且 \(z\) 与 \(y\) 在每个坐标上满足更强的同余条件，必有
\[
z\in U_{x^{(a,j)}}\subseteq F^{-1}(a).
\]
故 \(F(z)=a=F(y)\)。这说明 \(F\) 在同一有限商
\[
\prod_{p\in T}\ZZ/p^{n_p}\ZZ
\]
的每条纤维上常值，从而存在唯一映射 \(\overline F\) 使上述因子化成立。证毕。
\end{proof}

\begin{corollary}[有限 prime-audit 的不可完备性具有任意指数级歧义]\label{cor:xi-localized-finite-prime-audit-arbitrary-exponential-ambiguity}
\leanverified{paper\_xi\_localized\_finite\_prime\_audit\_arbitrary\_exponential\_ambiguity}
固定一个经定理 \ref{thm:xi-localized-finite-prime-audit-cylinder-factorization} 因子化到有限圆柱商
\[
Q_T:=\prod_{p\in T}\ZZ/p^{n_p}\ZZ
\]
的 finite prime-audit，其中 \(T\subset\mathbb P\) 有限。则对任意 \(N\ge 1\)，都存在至少 \(2^N\) 个两两非同构的局部化群
\[
G_{S_E}:=\ZZ[S_E^{-1}],
\qquad
S_E:=T\cup\{q_i:\ i\in E\},
\qquad
E\subseteq\{1,\dots,N\},
\]
其中 \(q_1,\dots,q_N\in\mathbb P\setminus T\) 互异，并且这些局部化群都诱导同一 audited finite quotient \(Q_T\)。因而在给定 audit 下，它们不可区分。
\end{corollary}
\begin{proof}
任选互异素数
\[
q_1,\dots,q_N\in\mathbb P\setminus T,
\]
并对每个 \(E\subseteq\{1,\dots,N\}\) 定义
\[
S_E:=T\cup\{q_i:\ i\in E\},
\qquad
G_{S_E}:=\ZZ[S_E^{-1}].
\]
由定理 \ref{thm:killo-localization-circle-dimension-prime-ledger-splitting}，
\[
\widehat{G_{S_E}}
\]
的 prime-register 核恰为
\[
\mathfrak K_{S_E}=\prod_{p\in S_E}\ZZ_p
=
\left(\prod_{p\in T}\ZZ_p\right)\times
\left(\prod_{q_i\in E}\ZZ_{q_i}\right).
\]
而审计因子商 \(Q_T\) 只读取 \(T\) 中素数的有限截断，因此对全部 \(E\) 都有同一投影
\[
\mathfrak K_{S_E}\longrightarrow Q_T,
\]
额外加入的素数坐标 \(q_i\) 全部被遗忘。故这 \(2^N\) 个局部化群诱导完全相同的 audited finite quotient \(Q_T\)，从而在该 audit 下不可区分。

另一方面，若 \(E\neq E'\)，则 \(S_E\neq S_{E'}\)。由定理 \ref{thm:xi-localized-integers-padic-kernel-rigidity}，
\[
G_{S_E}\cong G_{S_{E'}}
\iff
S_E=S_{E'},
\]
故这些 \(G_{S_E}\) 两两非同构。由于子集总数恰为 \(2^N\)，结论成立。证毕。
\end{proof}


\begin{theorem}[有限 prime-audit 的指数级不完备性]\label{thm:xi-localized-finite-prime-audit-exponential-incompleteness}
设 \(T\subset\mathbb P\) 为有限素数集。对任意 \(N\ge 1\)，存在至少 \(2^N\) 个两两非同构的有限素数局部化群
\[
G_E:=\ZZ[S_E^{-1}],
\qquad
S_E:=T\cup\{q_i:\ i\in E\},
\qquad
E\subseteq\{1,\dots,N\},
\]
其中 \(q_1,\dots,q_N\in\mathbb P\setminus T\) 互异，并且满足
\[
\cdim_{\mathrm{fr}}(G_E)=1,
\qquad
K_E^{(T)}:=\prod_{p\in S_E\cap T}\ZZ_p=\prod_{p\in T}\ZZ_p.
\]
因此，任何仅审计有限秩圆维与有限素数窗口 \(T\) 上可见 \(p\)-进核的制度，至少会把 \(2^N\) 个互不同构类压缩为同一个审计像。
\end{theorem}
\begin{proof}
任选互异素数
\[
q_1,\dots,q_N\in\mathbb P\setminus T.
\]
对每个子集 \(E\subseteq\{1,\dots,N\}\)，定义
\[
S_E:=T\cup\{q_i:\ i\in E\},
\qquad
G_E:=\ZZ[S_E^{-1}].
\]
由定理 \ref{thm:xi-localized-phase-sampling-closed-form}，对一切有限素数集 \(S\) 都有
\[
\cdim_{\mathrm{fr}}(\ZZ[S^{-1}])=1,
\]
故全部 \(G_E\) 具有相同的有限秩圆维。

另一方面，由构造
\[
S_E\cap T=T,
\]
于是
\[
K_E^{(T)}=\prod_{p\in S_E\cap T}\ZZ_p=\prod_{p\in T}\ZZ_p
\]
与 \(E\) 无关。故在有限素数窗口 \(T\) 上可见的 \(p\)-进核对这 \(2^N\) 个对象完全一致。

若 \(E\neq E'\)，则 \(S_E\neq S_{E'}\)。由定理 \ref{thm:xi-localized-integers-padic-kernel-rigidity}，
\[
G_E\cong G_{E'}
\iff
S_E=S_{E'},
\]
故这些 \(G_E\) 两两非同构。由于子集总数恰为 \(2^N\)，结论成立。证毕。
\end{proof}

\begin{theorem}[局部化对偶的同构刚性]\label{thm:xi-time-part56-localized-pontryagin-isomorphism-rigidity}
对任意有限素数集 \(S,T\subset \mathbb P\)，记
\[
G_S:=\ZZ[S^{-1}],
\qquad
\Sigma_S:=\widehat{G_S},
\]
则
\[
\boxed{\;
\Sigma_S\cong \Sigma_T
\quad\Longleftrightarrow\quad
S=T.\;}
\]
并且每个 \(\Sigma_S\) 都满足短正合列
\[
\boxed{\;
0\longrightarrow \prod_{p\in S}\ZZ_p
\longrightarrow \Sigma_S
\longrightarrow \TT
\longrightarrow 0,\;}
\]
同时
\[
\boxed{\;
\cdim_{\mathrm{fr}}(G_S)=1.\;}
\]
因此，所有有限素数局部化在连通圆维与有限秩扩展维数层上都塌缩到同一个值，而全部算术差异只能由全不连通的 \(p\)-进核恢复。
\end{theorem}
\begin{proof}
由定理 \ref{thm:xi-localized-phase-sampling-closed-form}，每个 \(G_S\) 的相位采样函数
\[
\Sigma^{\mathrm{ph}}_{G_S}(N)=\card{\operatorname{Hom}(G_S,\ZZ/N\ZZ)}
\]
唯一恢复素数集 \(S\)，并且同一定理已给出
\[
0\longrightarrow \prod_{p\in S}\ZZ_p
\longrightarrow \Sigma_S
\longrightarrow \TT
\longrightarrow 0,
\qquad
\cdim_{\mathrm{fr}}(G_S)=1.
\]
若 \(\Sigma_S\cong \Sigma_T\)，则由 Pontryagin 对偶的反变等价得到
\[
G_S\cong G_T.
\]
于是对每个 \(N\ge 1\) 都有
\[
\card{\operatorname{Hom}(G_S,\ZZ/N\ZZ)}
=
\card{\operatorname{Hom}(G_T,\ZZ/N\ZZ)},
\]
即 \(\Sigma^{\mathrm{ph}}_{G_S}=\Sigma^{\mathrm{ph}}_{G_T}\)。由相位采样闭式的唯一恢复性，必有 \(S=T\)。反向 implication 显然。证毕。
\end{proof}

\begin{corollary}[单圆维下的指数级局部化多样性]\label{cor:xi-time-part56-single-circle-dimension-exponential-diversity}
对每个整数 \(n\ge 1\)，都存在
\[
2^n
\]
个两两不同构的紧交换群，它们都具有相同的连通商 \(\TT\)，并且其 Pontryagin 对偶离散群都满足
\[
\cdim_{\mathrm{fr}}=1.
\]
\end{corollary}
\begin{proof}
取互异素数
\[
p_1,\dots,p_n,
\]
并对每个子集 \(E\subseteq\{1,\dots,n\}\) 定义
\[
S_E:=\{p_i:\ i\in E\},
\qquad
\Sigma_E:=\widehat{\ZZ[S_E^{-1}]}.
\]
由于子集总数为 \(2^n\)，只需证明这些 \(\Sigma_E\) 两两不同构。若 \(E\neq E'\)，则 \(S_E\neq S_{E'}\)，由定理 \ref{thm:xi-time-part56-localized-pontryagin-isomorphism-rigidity} 立得
\[
\Sigma_E\not\cong \Sigma_{E'}.
\]
另一方面，同一定理又给出每个 \(\Sigma_E\) 都满足短正合列
\[
0\longrightarrow \prod_{p\in S_E}\ZZ_p
\longrightarrow \Sigma_E
\longrightarrow \TT
\longrightarrow 0,
\]
故其连通商恒为 \(\TT\)，并且
\[
\cdim_{\mathrm{fr}}\!\bigl(\ZZ[S_E^{-1}]\bigr)=1.
\]
结论成立。证毕。
\end{proof}

\begin{theorem}[分层 prime-register 外置的最坏情形下界与自然扩张闭合]\label{thm:xi-layered-prime-register-sharp-natural-extension}
\leanverified{paper\_xi\_layered\_prime\_register\_sharp\_natural\_extension}
设 \((X_j)_{j=0}^{k}\) 为可数集合列，\(\pi_j:X_j\twoheadrightarrow X_{j+1}\) 为有限纤维满射，并假设对每个 \(j=0,\dots,k-1\) 都存在 \(y_j\in X_{j+1}\) 使
\[
\#\pi_j^{-1}(y_j)\ge 2.
\]
考虑任何分层外置
\[
\Psi(x_0)=(x_k,H(x_0)),
\qquad
H(x_0)=\prod_{j=0}^{k-1}h_j(x_0),
\]
其中每个 \(h_j(x_0)\) 的全部素因子均被限制在预先指定的 prime slice \(\mathcal P_j\)。若 \(\Psi\) 为单射，则最坏情形下每一层都必须有效占用对应的 prime slice；等价地，对每个 \(j\) 都存在微态 \(x_0^{(j)}\) 使
\[
h_j\bigl(x_0^{(j)}\bigr)\neq 1.
\]
另一方面，定理 \ref{thm:killo-layered-prime-slices-finite-depth} 给出恰在每层写入一枚素数的单射构造。故在“每层至少一处分支”的最坏情形下，分层 prime-register 外置所需的最小 prime-slice 数目恰等于层数 \(k\)。

特别地，若 \(T:X\twoheadrightarrow X\) 为非单射有限纤维满射，则任意深度 \(k\) 的前史消歧都至少需要 \(k\) 个有效 prime slices；而定理 \ref{thm:killo-natural-extension-branch-register} 的分支指标寄存器实现把整个无穷前史歧义共轭为右移插入系统，从而把上述有限深度最优外置在逆极限上精确封闭。
\end{theorem}
\begin{proof}
推论 \ref{cor:killo-layered-necessity-minimality} 断言：在每层都存在分支点的前提下，任何单射分层外置都迫使每个 prime slice \(\mathcal P_j\) 在某个微态上非平凡使用。故最坏情形下至少需要 \(k\) 个彼此独立的有效 prime slices。
\par
反向上，定理 \ref{thm:killo-layered-prime-slices-finite-depth} 已构造出单射映射
\[
\Phi(x_0)=(x_k,G(x_0)),
\qquad
G(x_0)=\prod_{j=0}^{k-1}q_j(x_0),
\qquad
q_j(x_0)\in\mathcal P_j,
\]
并且每层恰写入一枚素数。故上述下界可达，因而最小 prime-slice 数目正是 \(k\)。
\par
若 \(T:X\twoheadrightarrow X\) 为非单射有限纤维满射，则存在 \(y\in X\) 使 \(\#T^{-1}(y)\ge 2\)。对任意固定深度 \(k\)，取 \(X_j:=X\) 且 \(\pi_j:=T\)（\(0\le j\le k-1\)），便满足前述最坏情形假设，故至少需要 \(k\) 个有效 prime slices。
最后，定理 \ref{thm:killo-natural-extension-branch-register} 把自然扩张
\(
X^{\leftarrow}
\)
与右移插入系统
\[
(x,(a_1,a_2,\dots))\longmapsto (T(x),(\beta(x),a_1,a_2,\dots))
\]
共轭，从而说明有限深度上的最优分层外置在逆极限上恰由规范分支指标寄存器闭合。证毕。
\end{proof}

\paragraph{\texorpdfstring{$S$}{S}-solenoid 的泛终性质与 prime-register 轴的函子恢复}
在局部化数系对偶 \(\Sigma_S=\widehat{G_S}\) 的刚性、短正合列
\[
0\longrightarrow \prod_{p\in S}\ZZ_p\longrightarrow \Sigma_S\longrightarrow \TT\longrightarrow 0
\]
以及 \(\Sigma_S\cong \Sigma_T\iff S=T\) 的判据已经确定之后，还可把“prime-register 轴并非实现偶然，而是结构强制不变量”压缩为一条范畴论表述：在固定 \(S\)-局部化纪律的实现范畴内，\(\Sigma_S\) 是唯一的终对象，而其账本核的 \(p\)-主分量恰把素数集 \(S\) 函子地恢复出来；对偶地，离散侧的 \(G_S=\ZZ[S^{-1}]\) 则是同一局部化纪律下的初对象。

\begin{theorem}[\texorpdfstring{$S$}{S}-solenoid 的终对象泛性质与 prime-register 轴的函子恢复]\label{thm:xi-time-part56-solenoid-terminal-prime-register-recovery}
固定有限素数集 \(S\subset\mathbb P\)，并记
\[
G_S:=\ZZ[S^{-1}],
\qquad
\Sigma_S:=\widehat{G_S}.
\]
记 \(\pi_S:\Sigma_S\twoheadrightarrow\TT\) 为标准投影，并定义范畴 \(\mathcal C_S\) 如下：
\begin{enumerate}
  \item 对象为一对 \((K,\pi)\)，其中 \(K\) 为紧交换群，\(\pi:K\twoheadrightarrow\TT\) 为连续满射；
  \item 若记 \(A:=\widehat K\) 且 \(j:\ZZ\hookrightarrow A\) 为 \(\pi\) 的对偶嵌入，则要求对每个 \(p\in S\)，乘法映射
  \[
  [p]_A:A\longrightarrow A
  \]
  为双射；
  \item 态射 \(f:(K,\pi)\to(K',\pi')\) 为满足
  \[
  \pi' \circ f=\pi
  \]
  的连续群同态。
\end{enumerate}
则成立：
\begin{enumerate}
  \item \((\Sigma_S,\pi_S)\) 在 \(\mathcal C_S\) 中为终对象；等价地，对任意 \((K,\pi)\in\mathcal C_S\)，都存在唯一连续群同态
  \[
  f:K\longrightarrow \Sigma_S
  \]
  使
  \[
  \pi_S\circ f=\pi.
  \]
  \item 其账本核
  \[
  K_S^{\mathrm{reg}}:=\ker(\pi_S)
  \]
  满足
  \[
  K_S^{\mathrm{reg}}\cong \prod_{p\in S}\ZZ_p;
  \]
  \item 因而素数集 \(S\) 可由终对象的账本核函子地恢复为
  \[
  \boxed{\;
  S=\{\,p\in\mathbb P:\ (K_S^{\mathrm{reg}})_p\neq 0\,\},\;}
  \]
  其中 \((K_S^{\mathrm{reg}})_p\) 表示其 Sylow pro-\(p\) 分量。
\end{enumerate}
\end{theorem}

\begin{proof}
第 \(1\) 条正是局部化的泛性质在紧侧的对偶表述，而其离散对偶表述正是 \(G_S=\ZZ[S^{-1}]\) 在含 \(\ZZ\) 的 \(S\)-局部离散群范畴中具有初对象性质。
确切地说，定理 \ref{thm:cdim-s-solenoid-terminal-object} 已证明：对任意对象 \((K,\pi)\in\mathcal C_S\)，存在唯一连续群同态
\[
f:K\longrightarrow \Sigma_S
\]
使
\[
\pi_S\circ f=\pi.
\]
故 \((\Sigma_S,\pi_S)\) 为终对象。

第 \(2\) 条由定理 \ref{thm:xi-localized-phase-sampling-closed-form} 与定理 \ref{thm:xi-time-part56-localized-pontryagin-isomorphism-rigidity} 已给出的标准短正合列直接得到：
\[
0\longrightarrow \prod_{p\in S}\ZZ_p
\longrightarrow \Sigma_S
\xrightarrow{\ \pi_S\ }
\TT
\longrightarrow 0.
\]
因此
\[
\ker(\pi_S)\cong \prod_{p\in S}\ZZ_p.
\]

对第 \(3\) 条，\(\prod_{p\in S}\ZZ_p\) 的 Sylow pro-\(p\) 分量在且仅在 \(p\in S\) 时非平凡；
而不同素数坐标彼此正交，故核的 \(p\)-主分量结构唯一决定 \(S\)。
这正说明 prime-register 轴并非额外实现选择，而是由终对象的核函子强制记录的结构不变量。证毕。
\end{proof}

\paragraph{冻结 escort 谱、全微态位率阈值与三重算术独立的乘法密度}
以下条目把固定冻结相中的 escort Rényi 熵谱与速率压力差公式、bin-fold 本征推前的绝对 Rényi 常数项、dyadic 容量临界窗与二相自由能变分、全微态精确反演所需比特率的有限层闭式，以及分枝三次场、零温四次场与 RH 临界五次场之间的算术独立性，压缩为一组可直接调用的终端结论。

\begin{theorem}[escort Rényi 速率的压力差公式、双阈值塌缩与 \texorpdfstring{$q=0$}{q=0} 跃迁]\label{thm:xi-fixed-freezing-escort-renyi-spectrum-collapse}
\leanverified{paper\_xi\_fixed\_freezing\_escort\_renyi\_spectrum\_collapse}
沿用结论 \ref{con:xi-terminal-fixed-freezing-escort-oracle-sharp-threshold} 的记号。记
\[
a_{\mathrm c}:=\frac{\log\varphi-g_*}{\alpha_*-\alpha_*^{(2)}}<\infty,
\]
并记
\[
S_t(m):=\sum_{x\in X_m}d_m(x)^t
\qquad (t>0).
\]
对任意固定 \(a>0\) 定义 escort 分布
\[
\pi_{m,a}(x):=\frac{d_m(x)^a}{S_a(m)},
\qquad
A_m^{\max}:=\{x\in X_m:\ d_m(x)=M_m\},
\qquad
K_m:=|A_m^{\max}|.
\]
若极限
\[
P_t:=\lim_{m\to\infty}\frac1m\log S_t(m)
\]
存在，则称其为 \(t\)-阶压力。
记 \(H_q\) 为以自然对数定义的 Rényi 熵族，其中
\[
H_q(\mu):=\frac{1}{1-q}\log\sum_x \mu(x)^q\quad(q\in(0,\infty)\setminus\{1\}),
\]
\[
H_1(\mu):=-\sum_x \mu(x)\log\mu(x),
\qquad
H_\infty(\mu):=-\log\max_x \mu(x),
\qquad
H_0(\mu):=\log|\supp \mu|.
\]
则成立：
\begin{enumerate}
  \item 对任意 \(q\in(0,\infty)\setminus\{1\}\)，若 \(P_a\) 与 \(P_{aq}\) 存在，则
  \[
  \boxed{\;
  \lim_{m\to\infty}\frac1m H_q(\pi_{m,a})
  =
  \frac{P_{aq}-qP_a}{1-q}.\;}
  \]
  \item 若进一步落在双阈值区
  \[
  a>a_{\mathrm c},
  \qquad
  aq>a_{\mathrm c},
  \]
  则对每个 \(q\in(0,\infty)\setminus\{1\}\) 都有
  \[
  \lim_{m\to\infty}\frac1m H_q(\pi_{m,a})=g_*.
  \]
  当 \(q>1\) 时，第二个阈值由 \(aq>a>a_{\mathrm c}\) 自动成立。
  \item 对 \(q\in\{1,\infty\}\)，若 \(a>a_{\mathrm c}\)，则同样有
  \[
  \lim_{m\to\infty}\frac1m H_q(\pi_{m,a})=g_*.
  \]
  \item 若 \(a>a_{\mathrm c}\) 且记
  \[
  U_m^{\max}:=\mathrm{Unif}(A_m^{\max}),
  \]
  则总变差距离满足精确恒等式
  \[
  \boxed{\;
  \|\pi_{m,a}-U_m^{\max}\|_{\mathrm{TV}}
  =
  1-\pi_{m,a}(A_m^{\max})
  \le
  \exp\!\bigl(-m\Delta(a)+o(m)\bigr),\;}
  \]
  其中
  \[
  \Delta(a):=a(\alpha_*-\alpha_*^{(2)})-(\log\varphi-g_*)>0.
  \]
  \item Hartley 熵率不塌缩，而满足
  \[
  \boxed{\;
  \lim_{m\to\infty}\frac1m H_0(\pi_{m,a})
  =
  \lim_{m\to\infty}\frac1m\log|X_m|
  =
  \log\varphi.\;}
  \]
\end{enumerate}
因此，正阶 Rényi 速率先由压力差公式统一表示；一旦 \((a,aq)\) 同时越过冻结阈值，全部正阶 Rényi 熵率便共同塌缩到 \(g_*\)；而当 \(\log\varphi>g_*\) 时，Hartley 熵率仍停留在 \(\log\varphi\)，从而在 \(q=0\) 处发生真实跃迁。
\end{theorem}
\begin{proof}
先取 \(q\in(0,\infty)\setminus\{1\}\)。由定义
\[
\sum_{x\in X_m}\pi_{m,a}(x)^q
=
\frac{S_{aq}(m)}{S_a(m)^q},
\]
故
\[
H_q(\pi_{m,a})
=
\frac{1}{1-q}\Bigl(\log S_{aq}(m)-q\log S_a(m)\Bigr).
\]
若 \(P_a\) 与 \(P_{aq}\) 存在，两边同除以 \(m\) 并令 \(m\to\infty\)，便得到第一条。

若进一步 \(a>a_{\mathrm c}\) 且 \(aq>a_{\mathrm c}\)，则定理 \ref{thm:fold-pressure-freezing-threshold} 同时适用于指数 \(a\) 与 \(aq\)，从而
\[
\frac1m\log S_a(m)\to a\alpha_*+g_*,
\qquad
\frac1m\log S_{aq}(m)\to aq\alpha_*+g_*.
\]
代入即得
\[
\frac1m H_q(\pi_{m,a})
\to
\frac{(aq\alpha_*+g_*)-q(a\alpha_*+g_*)}{1-q}
=
g_*.
\]
这便证明了第二条。
\par
对 \(q=\infty\)，推论 \ref{cor:fold-escort-minentropy-rate} 已给出
\[
\lim_{m\to\infty}\frac1m H_\infty(\pi_{m,a})=h_\infty(a)=g_*.
\]
这给出第三条中 \(q=\infty\) 的情形。
\par
再证 Shannon 情形 \(q=1\)。记
\[
p_m:=\pi_{m,a}(A_m^{\max}).
\]
由于 \(d_m(x)=M_m\) 在 \(A_m^{\max}\) 上常值，\(\pi_{m,a}\) 在该集合上的条件分布本身就是均匀分布 \(U_m^{\max}\)。因此
\[
\|\pi_{m,a}-U_m^{\max}\|_{\mathrm{TV}}
=
\frac12\sum_{x\in A_m^{\max}}\left|\frac{p_m}{K_m}-\frac1{K_m}\right|
+\frac12\sum_{x\notin A_m^{\max}}\pi_{m,a}(x)
=
1-p_m.
\]
再由定理 \ref{thm:fold-escort-groundstate-concentration}，
\[
1-p_m\le \exp\!\bigl(-m\Delta(a)+o(m)\bigr),
\]
即得第四条中的总变差公式。
\par
若 \(p_m<1\)，记 \(\eta_m\) 为 \(\pi_{m,a}\) 在补集 \(X_m\setminus A_m^{\max}\) 上的条件分布；若 \(p_m=1\)，任取 \(\eta_m\) 即可。于是
\[
\pi_{m,a}=p_m\,U_m^{\max}+(1-p_m)\eta_m
\]
且两项支撑互不相交，从而 Shannon 熵满足精确分解
\[
H(\pi_{m,a})
=
h_{\mathrm{bin}}(p_m)+p_m\log K_m+(1-p_m)H(\eta_m),
\]
其中
\[
h_{\mathrm{bin}}(t):=-t\log t-(1-t)\log(1-t).
\]
又因为
\[
0\le H(\eta_m)\le \log|X_m|=\log F_{m+2}=m\log\varphi+O(1),
\]
并且 \(1-p_m\) 指数衰减，故
\[
h_{\mathrm{bin}}(p_m)=o(m),
\qquad
(1-p_m)H(\eta_m)=o(m),
\qquad
p_m\log K_m=\log K_m+o(m).
\]
于是
\[
H(\pi_{m,a})=\log K_m+o(m),
\]
从而
\[
\lim_{m\to\infty}\frac1m H(\pi_{m,a})=g_*.
\]
这便完成了第三条中 \(q=1\) 的情形。
\par
最后，由定义对一切 \(x\in X_m\) 都有 \(d_m(x)\ge 1\)，故 \(\pi_{m,a}(x)>0\) 在整个 \(X_m\) 上成立。因此
\[
H_0(\pi_{m,a})=\log|\supp(\pi_{m,a})|=\log|X_m|=\log F_{m+2},
\]
从而
\[
\lim_{m\to\infty}\frac1m H_0(\pi_{m,a})=\log\varphi.
\]
这正是第五条。
证毕。
\end{proof}

\begin{corollary}[冻结相中 Rényi 阶阈值的显式闭式]\label{cor:xi-fixed-freezing-renyi-critical-order-threshold}
沿用定理 \ref{thm:xi-fixed-freezing-escort-renyi-spectrum-collapse} 的记号。对固定
\[
a>a_{\mathrm c},
\]
定义
\[
r_{\mathrm c}(a):=\frac{a_{\mathrm c}}{a}\in(0,1).
\]
则对任意
\[
r\in[r_{\mathrm c}(a),\infty]
\]
都有
\[
\lim_{m\to\infty}\frac1m H_r(\pi_{m,a})=g_*.
\]
另一方面，
\[
\lim_{m\to\infty}\frac1m H_0(\pi_{m,a})=\log\varphi.
\]
因此，冻结相内部的 Rényi 谱在阶参数上具有显式阈值 \(r_{\mathrm c}(a)\)：阈值以上的全部信息型熵率统一塌缩到基态简并指数 \(g_*\)，而支撑熵率仍停留在环境熵率 \(\log\varphi\)。
\end{corollary}
\begin{proof}
对 \(r>r_{\mathrm c}(a)\) 且 \(r\notin\{1,\infty\}\)，由于
\[
ar>a_{\mathrm c},
\]
定理 \ref{thm:xi-fixed-freezing-escort-renyi-spectrum-collapse} 的第二条直接给出
\[
\lim_{m\to\infty}\frac1m H_r(\pi_{m,a})=g_*.
\]
对 \(r\in\{1,\infty\}\)，同一定理第三条给出同样结论；对 \(r=0\)，第五条给出
\[
\lim_{m\to\infty}\frac1m H_0(\pi_{m,a})=\log\varphi.
\]
因此只需处理临界点
\[
r=r_{\mathrm c}(a)=\frac{a_{\mathrm c}}{a}\in(0,1).
\]

记
\[
\varepsilon_m:=1-\pi_{m,a}(A_m^{\max}).
\]
由定理 \ref{thm:xi-fixed-freezing-escort-renyi-spectrum-collapse} 的第四条，
\[
\varepsilon_m\le \exp\!\bigl(-m\Delta(a)+o(m)\bigr),
\qquad
\Delta(a):=a(\alpha_*-\alpha_*^{(2)})-(\log\varphi-g_*)>0.
\]
并且 \(\pi_{m,a}\) 在 \(A_m^{\max}\) 上均匀，故
\[
\sum_{x\in A_m^{\max}}\pi_{m,a}(x)^r
=
K_m\left(\frac{1-\varepsilon_m}{K_m}\right)^r
=
K_m^{1-r}(1-\varepsilon_m)^r.
\]
另一方面，由 \(0<r<1\) 时的凹性不等式
\[
\sum_{i=1}^{N}b_i^r\le N^{1-r}\Bigl(\sum_{i=1}^{N}b_i\Bigr)^r
\qquad (b_i\ge 0),
\]
可得
\[
\sum_{x\notin A_m^{\max}}\pi_{m,a}(x)^r
\le
|X_m|^{1-r}\varepsilon_m^r.
\]
由于 \(|X_m|=\exp(m\log\varphi+o(m))\)，上式右端满足
\[
|X_m|^{1-r}\varepsilon_m^r
\le
\exp\!\bigl(m[(1-r)\log\varphi-r\Delta(a)]+o(m)\bigr).
\]
当 \(r=r_{\mathrm c}(a)\) 时，由
\[
ra(\alpha_*-\alpha_*^{(2)})=a_{\mathrm c}(\alpha_*-\alpha_*^{(2)})=\log\varphi-g_*
\]
可得
\[
r\Delta(a)=r\bigl[a(\alpha_*-\alpha_*^{(2)})-(\log\varphi-g_*)\bigr]
=(1-r)(\log\varphi-g_*),
\]
故
\[
(1-r)\log\varphi-r\Delta(a)=(1-r)g_*.
\]
同时
\[
K_m^{1-r}(1-\varepsilon_m)^r
=
\exp\!\bigl((1-r)g_*\,m+o(m)\bigr).
\]
于是
\[
\exp\!\bigl((1-r)g_*\,m+o(m)\bigr)
\le
\sum_x\pi_{m,a}(x)^r
\le
2\exp\!\bigl((1-r)g_*\,m+o(m)\bigr),
\]
从而
\[
\sum_x\pi_{m,a}(x)^r
=
\exp\!\bigl((1-r)g_*\,m+o(m)\bigr).
\]
两边取对数并除以 \(1-r\)，便得到
\[
H_r(\pi_{m,a})=g_*\,m+o(m),
\]
即
\[
\lim_{m\to\infty}\frac1m H_r(\pi_{m,a})=g_*.
\]
证毕。
\end{proof}

\begin{corollary}[冻结 escort 态到基态均匀分布的精确总变差公式]\label{cor:xi-fixed-freezing-escort-groundstate-tv-identity}
\leanverified{paper\_xi\_fixed\_freezing\_escort\_groundstate\_tv\_identity}
沿用定理 \ref{thm:xi-fixed-freezing-escort-renyi-spectrum-collapse} 的记号。对任意整数
\[
a>a_{\mathrm c},
\]
记
\[
U_m^{\max}:=\mathrm{Unif}(A_m^{\max}).
\]
则有精确恒等式
\[
\|\pi_{m,a}-U_m^{\max}\|_{\mathrm{TV}}
=
1-\pi_{m,a}(A_m^{\max}),
\]
并且
\[
\|\pi_{m,a}-U_m^{\max}\|_{\mathrm{TV}}
\le
\exp\!\bigl(-m\Delta(a)+o(m)\bigr).
\]
\end{corollary}
\begin{proof}
这正是定理 \ref{thm:xi-fixed-freezing-escort-renyi-spectrum-collapse} 第四条在整数
\(
a>a_{\mathrm c}
\)
上的直接重述。证毕。
\end{proof}

\begin{theorem}[冻结 escort 态对一切有界可观测量的指数统计塌缩]\label{thm:xi-fixed-freezing-escort-bounded-observable-collapse}
沿用推论 \ref{cor:xi-fixed-freezing-escort-groundstate-tv-identity} 的记号。对任意固定
\[
a>a_{\mathrm c}
\]
与任意有界可观测量 \(f_m:X_m\to\CC\)，都有
\[
\boxed{\;
\left|
\EE_{\pi_{m,a}}f_m-\EE_{U_m^{\max}}f_m
\right|
\le
2\|f_m\|_\infty
\exp\!\bigl(-m\Delta(a)+o(m)\bigr).\;}
\]
若进一步 \(f_m\) 在 \(A_m^{\max}\) 上为常数 \(c_m\)，则有精确恒等式
\[
\boxed{\;
\EE_{\pi_{m,a}}f_m-\EE_{U_m^{\max}}f_m
=
\sum_{x\notin A_m^{\max}}\pi_{m,a}(x)\bigl(f_m(x)-c_m\bigr).\;}
\]
因此在冻结区内，一切有界可观测量都以指数速度塌缩到基态均匀分布的统计值；而对基态层上常值的可观测量，全部误差恰由基态外质量单独承担。
\leanverified{paper\_xi\_fixed\_freezing\_escort\_bounded\_observable\_collapse}
\end{theorem}
\begin{proof}
由总变差距离与有界测试函数的标准对偶不等式，
\[
\left|
\EE_{\pi_{m,a}}f_m-\EE_{U_m^{\max}}f_m
\right|
\le
2\|f_m\|_\infty
\|\pi_{m,a}-U_m^{\max}\|_{\mathrm{TV}}.
\]
再由推论 \ref{cor:xi-fixed-freezing-escort-groundstate-tv-identity}，
\[
\|\pi_{m,a}-U_m^{\max}\|_{\mathrm{TV}}
\le
\exp\!\bigl(-m\Delta(a)+o(m)\bigr),
\]
便得到第一条估计。

若 \(f_m|_{A_m^{\max}}\equiv c_m\)，则
\[
\EE_{U_m^{\max}}f_m=c_m.
\]
另一方面，由于 \(\pi_{m,a}\) 在 \(A_m^{\max}\) 上本身均匀，故
\[
\sum_{x\in A_m^{\max}}\pi_{m,a}(x)f_m(x)
=
\pi_{m,a}(A_m^{\max})\,c_m.
\]
于是
\[
\begin{aligned}
\EE_{\pi_{m,a}}f_m-\EE_{U_m^{\max}}f_m
&=
\sum_{x\in A_m^{\max}}\pi_{m,a}(x)f_m(x)
+
\sum_{x\notin A_m^{\max}}\pi_{m,a}(x)f_m(x)-c_m\\
&=
\bigl(\pi_{m,a}(A_m^{\max})-1\bigr)c_m
+
\sum_{x\notin A_m^{\max}}\pi_{m,a}(x)f_m(x)\\
&=
\sum_{x\notin A_m^{\max}}\pi_{m,a}(x)\bigl(f_m(x)-c_m\bigr).
\end{aligned}
\]
证毕。
\end{proof}

\leanverified{paper\_xi\_fixed\_freezing\_escort\_pairwise\_tv\_coalescence}
\begin{corollary}[冻结相内部不同 escort 态的指数全变差趋同]\label{cor:xi-fixed-freezing-escort-pairwise-tv-coalescence}
沿用定理 \ref{thm:xi-fixed-freezing-escort-renyi-spectrum-collapse} 的记号。对任意固定
\[
a,b>a_{\mathrm c},
\]
都有
\[
\|\pi_{m,a}-\pi_{m,b}\|_{\mathrm{TV}}
\le
\exp\!\bigl(-m\Delta(a)+o(m)\bigr)
+
\exp\!\bigl(-m\Delta(b)+o(m)\bigr).
\]
特别地，
\[
\|\pi_{m,a}-\pi_{m,b}\|_{\mathrm{TV}}
\le
\exp\!\bigl(-m\min\{\Delta(a),\Delta(b)\}+o(m)\bigr).
\]
\end{corollary}
\begin{proof}
记
\[
U_m^{\max}:=\mathrm{Unif}(A_m^{\max}).
\]
由三角不等式，
\[
\|\pi_{m,a}-\pi_{m,b}\|_{\mathrm{TV}}
\le
\|\pi_{m,a}-U_m^{\max}\|_{\mathrm{TV}}
+
\|U_m^{\max}-\pi_{m,b}\|_{\mathrm{TV}}.
\]
再由定理 \ref{thm:xi-fixed-freezing-escort-renyi-spectrum-collapse} 的第四条，
\[
\|\pi_{m,a}-U_m^{\max}\|_{\mathrm{TV}}
\le
\exp\!\bigl(-m\Delta(a)+o(m)\bigr),
\]
\[
\|U_m^{\max}-\pi_{m,b}\|_{\mathrm{TV}}
\le
\exp\!\bigl(-m\Delta(b)+o(m)\bigr).
\]
相加即得第一条。第二条由
\[
u_m+v_m
\le
2\max\{u_m,v_m\}
\]
应用于
\[
u_m:=\exp\!\bigl(-m\Delta(a)+o(m)\bigr),
\qquad
v_m:=\exp\!\bigl(-m\Delta(b)+o(m)\bigr)
\]
后直接得到。证毕。
\end{proof}

\begin{theorem}[全微态精确反演的有限层比特阈值与线性速率]\label{thm:xi-full-microstate-exact-inversion-bitrate-threshold}
\leanverified{paper\_xi\_full\_microstate\_exact\_inversion\_bitrate\_threshold}
对整数预算 \(B\ge 0\)，定义 dyadic 预言机容量
\[
C_m(B):=\sum_{x\in X_m}\min\bigl(d_m(x),2^B\bigr),
\]
并记最大纤维规模
\[
M_m:=\max_{x\in X_m}d_m(x).
\]
定义全微态精确反演所需的最小比特数
\[
B_m^{\mathrm{full}}
:=
\min\{B\in\ZZ_{\ge 0}: C_m(B)=2^m\}.
\]
则对每个 \(m\) 都有严格恒等式
\[
\boxed{\;
B_m^{\mathrm{full}}=\left\lceil \log_2 M_m\right\rceil.\;}
\]
特别地，若
\[
\alpha_*=\lim_{m\to\infty}\frac1m\log M_m
\]
存在，则
\[
\boxed{\;
\lim_{m\to\infty}\frac{B_m^{\mathrm{full}}}{m}
=
\frac{\alpha_*}{\log 2}.\;}
\]
更进一步，对任意 \(\varepsilon>0\) 与任意整数序列 \(B_m\ge 0\)，成立：
\begin{enumerate}
  \item 若
  \[
  B_m\ge \left(\frac{\alpha_*}{\log2}+\varepsilon\right)m
  \]
  对充分大 \(m\) 成立，则最终
  \[
  C_m(B_m)=2^m.
  \]
  \item 若
  \[
  B_m\le \left(\frac{\alpha_*}{\log2}-\varepsilon\right)m
  \]
  对充分大 \(m\) 成立，则
  \[
  \boxed{\;
  2^m-C_m(B_m)\ \ge\ K_m\bigl(M_m-2^{B_m}\bigr),\;}
  \]
  其中 \(K_m:=|A_m^{\max}|\)。若再假设
  \[
  g_*=\lim_{m\to\infty}\frac1m\log K_m
  \]
  存在，则
  \[
  \boxed{\;
  \liminf_{m\to\infty}\frac1m\log\bigl(2^m-C_m(B_m)\bigr)
  \ge
  \alpha_*+g_*.\;}
  \]
\end{enumerate}
因此 \(\alpha_*/\log 2\) 不是近似阈值，而是全微态精确反演的精确信息速率阈值。
\end{theorem}
\begin{proof}
由定理 \ref{thm:fold-oracle-capacity-closed-form}，
\[
C_m(B)=\sum_{x\in X_m}\min\bigl(d_m(x),2^B\bigr).
\]
注意到
\[
\sum_{x\in X_m}d_m(x)=2^m.
\]
因此 \(C_m(B)=2^m\) 当且仅当每一项都未被截断，即
\[
\min(d_m(x),2^B)=d_m(x)\qquad(\forall x\in X_m),
\]
等价于
\[
2^B\ge d_m(x)\qquad(\forall x\in X_m),
\]
也即
\[
2^B\ge M_m.
\]
满足此条件的最小非负整数 \(B\) 正是 \(\lceil \log_2 M_m\rceil\)，故得到第一式。
\par
若 \(\alpha_*=\lim m^{-1}\log M_m\) 存在，则
\[
\frac1m\log_2 M_m=\frac{1}{m\log2}\log M_m\to \frac{\alpha_*}{\log2}.
\]
而上取整仅引入 \(O(1)\) 误差，故
\[
\frac{B_m^{\mathrm{full}}}{m}\to \frac{\alpha_*}{\log2}.
\]
\par
若
\[
B_m\ge \left(\frac{\alpha_*}{\log2}+\varepsilon\right)m,
\]
则
\[
2^{B_m}\ge \exp\!\bigl((\alpha_*+\varepsilon\log2)m\bigr).
\]
由 \(M_m=\exp(\alpha_*m+o(m))\)，对充分大 \(m\) 必有 \(2^{B_m}\ge M_m\)，于是 \(C_m(B_m)=2^m\)。
\par
反之，若
\[
B_m\le \left(\frac{\alpha_*}{\log2}-\varepsilon\right)m,
\]
则
\[
2^m-C_m(B_m)
=
\sum_{x\in X_m}\bigl(d_m(x)-2^{B_m}\bigr)_+,
\]
其中 \((t)_+:=\max\{t,0\}\)。对每个 \(x\in A_m^{\max}\) 都有 \(d_m(x)=M_m\)，故
\[
2^m-C_m(B_m)\ge \sum_{x\in A_m^{\max}}\bigl(M_m-2^{B_m}\bigr)
=
K_m\bigl(M_m-2^{B_m}\bigr).
\]
这便得到所示有限层下界。
\par
若 \(g_*\) 存在，则 \(K_m=\exp(g_*m+o(m))\)。同时由预算假设
\[
2^{B_m}\le \exp\!\bigl((\alpha_*-\varepsilon\log2)m\bigr),
\]
故
\[
M_m-2^{B_m}
=
\exp(\alpha_*m+o(m))
\left(1-\exp\!\bigl(-\varepsilon(\log2)m+o(m)\bigr)\right)
=
\exp(\alpha_*m+o(m)).
\]
于是
\[
K_m\bigl(M_m-2^{B_m}\bigr)=\exp\!\bigl((\alpha_*+g_*)m+o(m)\bigr),
\]
从而
\[
\liminf_{m\to\infty}\frac1m\log\bigl(2^m-C_m(B_m)\bigr)\ge \alpha_*+g_*.
\]
证毕。
\end{proof}

\paragraph{bin-fold 绝对熵常数、容量临界窗、阈值预算与上界刚性}
在前述冻结 escort 熵率、全微态位率阈值之外，bin-fold 本征推前自身的绝对 Rényi 常数项、dyadic 预算下的临界成功率窗口、阈值预算的常数跃迁、最大纤维组合上界的指数精确性，以及其零温修正的二相变分结构，也可由既有两态渐近与 escort 末位律完全闭合。

\begin{theorem}[bin-fold 本征推前的绝对 Rényi 熵常数项闭式]\label{thm:xi-foldbin-absolute-renyi-entropy-constant-closed}
\leanverified{paper\_xi\_foldbin\_absolute\_renyi\_entropy\_constant\_closed}
令
\[
\mu_m(w):=\frac{d_m^{\mathrm{bin}}(w)}{2^m},
\qquad
w\in X_m.
\]
则对任意固定 \(q>0\)、\(q\neq 1\)，有
\[
H_q(\mu_m)
=
m\log\varphi
+
\frac{1}{1-q}\log\!\Bigl(\frac{\varphi+\varphi^{-q}}{\sqrt5}\Bigr)
+
O\!\Bigl(\Bigl(\frac{\varphi}{2}\Bigr)^m\Bigr).
\]
相应的 Shannon 熵满足
\[
H(\mu_m)
=
m\log\varphi
+
\frac{\log\varphi}{1+\varphi^2}
+
O\!\Bigl(\Bigl(\frac{\varphi}{2}\Bigr)^m\Bigr).
\]
\end{theorem}
\begin{proof}
在定理 \ref{thm:xi-fold-escort-renyi-shannon-constants} 中取 escort 参数为 \(1\)，并把其中的 Rényi 阶改记为 \(q\)，即可得到
\[
H_q(\mu_m)
=
m\log\varphi
+
\frac{1}{1-q}\left[
\log\!\Bigl(\frac{\varphi+\varphi^{-q}}{\sqrt5}\Bigr)
-q\log\!\Bigl(\frac{\varphi+\varphi^{-1}}{\sqrt5}\Bigr)
\right]
+
O\!\Bigl(\Bigl(\frac{\varphi}{2}\Bigr)^m\Bigr).
\]
而 \(\varphi+\varphi^{-1}=\sqrt5\)，故第二项化简为所示闭式。Shannon 情形正是推论 \ref{cor:xi-foldbin-shannon-gauge-constant-mirror} 的内容；再用恒等式 \(1+\varphi^2=\varphi\sqrt5\) 即得此处写法。证毕。
\end{proof}

\begin{theorem}[bin-fold dyadic 预言机容量的临界窗口极限]\label{thm:xi-foldbin-dyadic-capacity-critical-window-limit}
\leanverified{paper\_xi\_foldbin\_dyadic\_capacity\_critical\_window\_limit}
对整数预算 \(B\ge 0\)，定义 dyadic 预言机容量与最优成功率
\[
C_m^{\mathrm{bin}}(B):=\sum_{w\in X_m}\min\!\bigl(d_m^{\mathrm{bin}}(w),2^B\bigr),
\qquad
\mathrm{Succ}_m^{\mathrm{bin}}(B):=\frac{C_m^{\mathrm{bin}}(B)}{2^m}=P_m^\star(2^B).
\]
记
\[
\lambda:=\frac{2}{\varphi}.
\]
若整数序列 \((B_m)_{m\ge 1}\) 满足
\[
\frac{2^{B_m}}{\lambda^m}\longrightarrow \zeta\in(0,\infty),
\]
则极限
\[
\lim_{m\to\infty}\mathrm{Succ}_m^{\mathrm{bin}}(B_m)
\]
存在，且等于
\[
\mathcal S_{\mathrm{bin}}(\zeta)
:=
\begin{cases}
\dfrac{\varphi^2}{\sqrt5}\,\zeta, & 0<\zeta\le \varphi^{-1},\\[3mm]
\dfrac{\varphi}{\sqrt5}\,\zeta+\dfrac{1}{\varphi\sqrt5}, & \varphi^{-1}\le \zeta\le 1,\\[3mm]
1, & \zeta\ge 1.
\end{cases}
\]
\end{theorem}
\begin{proof}
记
\[
A_b^{(m)}:=\{w\in X_m:\ w_m=b\},
\qquad
b\in\{0,1\}.
\]
由定理 \ref{thm:fold-bin-two-state-asymptotic}，存在常数 \(C>0\) 使得对充分大 \(m\) 与任意 \(w\in X_m\) 都有
\[
d_m^{\mathrm{bin}}(w)
=
\lambda^m\varphi^{-w_m}\bigl(1+\varepsilon_{m,w}\bigr),
\qquad
|\varepsilon_{m,w}|\le C\Bigl(\frac{\varphi}{2}\Bigr)^m.
\]
又由稳定词计数，
\[
|A_0^{(m)}|=F_{m+1},
\qquad
|A_1^{(m)}|=F_m.
\]

若 \(0<\zeta<\varphi^{-1}\)，则对充分大 \(m\) 有
\[
2^{B_m}<\lambda^m\varphi^{-1}\left(1-C\Bigl(\frac{\varphi}{2}\Bigr)^m\right),
\]
故全部项都被阈值截断，从而
\[
C_m^{\mathrm{bin}}(B_m)=2^{B_m}|X_m|=2^{B_m}F_{m+2}.
\]
于是
\[
\mathrm{Succ}_m^{\mathrm{bin}}(B_m)
=
\frac{2^{B_m}}{\lambda^m}\cdot \frac{F_{m+2}}{\varphi^m}
\longrightarrow
\zeta\cdot \frac{\varphi^2}{\sqrt5}.
\]

若 \(\varphi^{-1}<\zeta<1\)，则对充分大 \(m\) 有
\[
\lambda^m\varphi^{-1}\left(1+C\Bigl(\frac{\varphi}{2}\Bigr)^m\right)
<
2^{B_m}
<
\lambda^m\left(1-C\Bigl(\frac{\varphi}{2}\Bigr)^m\right).
\]
因此末位 \(0\) 扇区仍被截断，而末位 \(1\) 扇区已全部通过，故
\[
C_m^{\mathrm{bin}}(B_m)
=
F_{m+1}2^{B_m}+\sum_{w\in A_1^{(m)}}d_m^{\mathrm{bin}}(w).
\]
再由上一致展开，
\[
\sum_{w\in A_1^{(m)}}d_m^{\mathrm{bin}}(w)
=
F_m\lambda^m\varphi^{-1}\left(1+O\!\Bigl(\Bigl(\frac{\varphi}{2}\Bigr)^m\Bigr)\right).
\]
两端除以 \(2^m\) 得
\[
\mathrm{Succ}_m^{\mathrm{bin}}(B_m)
=
\frac{2^{B_m}}{\lambda^m}\cdot \frac{F_{m+1}}{\varphi^m}
+
\frac{F_m}{\varphi^{m+1}}
+
o(1)
\longrightarrow
\zeta\frac{\varphi}{\sqrt5}+\frac{1}{\varphi\sqrt5}.
\]

若 \(\zeta>1\)，则对充分大 \(m\) 有
\[
2^{B_m}>\lambda^m\left(1+C\Bigl(\frac{\varphi}{2}\Bigr)^m\right),
\]
故所有纤维都完全通过，从而
\[
C_m^{\mathrm{bin}}(B_m)=\sum_{w\in X_m}d_m^{\mathrm{bin}}(w)=2^m,
\]
即
\[
\mathrm{Succ}_m^{\mathrm{bin}}(B_m)\to 1.
\]
边界点 \(\zeta=\varphi^{-1},1\) 由上述分段公式的连续性直接补上。证毕。
\end{proof}

\begin{corollary}[critical-scale 预言机反演成功曲线的双折点与最小资源反演闭式]\label{cor:xi-foldbin-critical-scale-success-curve-inversion}
\leanverified{paper\_xi\_foldbin\_critical\_scale\_success\_curve\_inversion}
沿用定理 \ref{thm:xi-foldbin-dyadic-capacity-critical-window-limit} 的记号。对 \(s\ge 0\) 定义连续临界尺度
\[
T_m(s):=s\Bigl(\frac{2}{\varphi}\Bigr)^m,
\]
以及连续化成功曲线
\[
\mathcal S_m(s):=
2^{-m}\,\mathcal C_m^{\mathrm{bin,cont}}\!\bigl(T_m(s)\bigr).
\]
则 \(\mathcal S_m\) 在每个紧集 \(K\Subset(0,\infty)\) 上一致收敛到分段线性极限
\[
\boxed{\;
\mathcal S_\infty(s)=
\begin{cases}
\dfrac{\varphi^2}{\sqrt5}\,s, & 0\le s\le \varphi^{-1},\\[6pt]
\dfrac{\varphi}{\sqrt5}\,s+\dfrac{1}{\varphi\sqrt5}, & \varphi^{-1}\le s\le 1,\\[8pt]
1, & s\ge 1.
\end{cases}}
\]
从而极限成功曲线恰在
\[
s=\varphi^{-1},
\qquad
s=1
\]
处出现且仅出现两处刚性折点。

若进一步定义其最小资源广义逆
\[
\mathcal S_\infty^{-1}(\alpha):=\inf\{s\ge 0:\ \mathcal S_\infty(s)\ge \alpha\}
\qquad (0\le \alpha\le 1),
\]
则有显式公式
\[
\boxed{\;
\mathcal S_\infty^{-1}(\alpha)=
\begin{cases}
\dfrac{\sqrt5}{\varphi^2}\,\alpha, & 0\le \alpha\le \dfrac{\varphi}{\sqrt5},\\[8pt]
\dfrac{\sqrt5\,\alpha-\varphi^{-1}}{\varphi}, & \dfrac{\varphi}{\sqrt5}\le \alpha\le 1.
\end{cases}}
\]
\end{corollary}
\begin{proof}
由定理 \ref{thm:xi-foldbin-underresolution-two-atomic-curvature-measure}，
\[
\mathcal S_m(s)=\widehat C_m(s)
\]
在每个紧集 \(K\Subset(0,\infty)\) 上一致收敛到
\[
\widehat C_\infty(s)=\frac{\varphi}{\sqrt5}\min\{1,s\}+\frac{1}{\sqrt5}\min\{\varphi^{-1},s\}.
\]
再注意 \(\mathcal S_m(0)=0=\mathcal S_\infty(0)\)。将右端按 \(s\in[0,\varphi^{-1}]\)、\(s\in[\varphi^{-1},1]\) 与 \(s\ge 1\) 分段展开，即得所示三段式。因此折点恰位于两个最小值函数发生切换的
\[
s=\varphi^{-1},\qquad s=1.
\]

又因 \(\mathcal S_\infty\) 在 \([0,1]\) 上分段线性且严格递增，而在 \([1,\infty)\) 上恒等于 \(1\)，故其最小资源广义逆由两段线性方程直接求解而得。对第一段，
\[
\alpha=\frac{\varphi^2}{\sqrt5}\,s
\]
给出
\[
s=\frac{\sqrt5}{\varphi^2}\,\alpha.
\]
对第二段，
\[
\alpha=\frac{\varphi}{\sqrt5}\,s+\frac{1}{\varphi\sqrt5}
\]
给出
\[
s=\frac{\sqrt5\,\alpha-\varphi^{-1}}{\varphi}.
\]
证毕。
\end{proof}

\begin{corollary}[critical-scale 资源边际收益在首折点处恰按 \texorpdfstring{$\varphi$}{phi} 倍坍缩]\label{cor:xi-foldbin-critical-scale-marginal-gain-phi-drop}
\leanverified{paper\_xi\_foldbin\_critical\_scale\_marginal\_gain\_phi\_drop}
沿用推论 \ref{cor:xi-foldbin-critical-scale-success-curve-inversion} 的记号。则极限成功曲线的边际收益满足
\[
\mathcal S_\infty'(s)=
\begin{cases}
\dfrac{\varphi^2}{\sqrt5}, & 0<s<\varphi^{-1},\\[6pt]
\dfrac{\varphi}{\sqrt5}, & \varphi^{-1}<s<1,\\[6pt]
0, & s>1.
\end{cases}
\]
特别地，在首折点 \(s=\varphi^{-1}\) 处，左右导数之比满足
\[
\boxed{\;
\frac{\mathcal S_\infty'(\varphi^{-1}-0)}{\mathcal S_\infty'(\varphi^{-1}+0)}=\varphi.\;}
\]
亦即，一旦末位 \(1\) 层率先饱和，继续增加临界尺度资源所获得的边际成功收益恰好下降一个黄金比例因子。
\end{corollary}
\begin{proof}
对推论 \ref{cor:xi-foldbin-critical-scale-success-curve-inversion} 中的分段线性表达式逐段求导即可。证毕。
\end{proof}

\begin{theorem}[bin-fold 零温自由能修正的二相变分公式]\label{thm:xi-foldbin-two-phase-free-energy-variational}
\leanverified{paper\_xi\_foldbin\_two\_phase\_free\_energy\_variational}
定义
\[
\Psi(q):=\log\!\Bigl(\frac{\varphi+\varphi^{-q}}{\sqrt5}\Bigr),
\qquad
q\ge 0.
\]
则有精确变分公式
\[
\Psi(q)
=
\log\!\Bigl(\frac{\varphi}{\sqrt5}\Bigr)
+
\sup_{\theta\in[0,1]}
\Bigl\{
h(\theta)-\theta(q+1)\log\varphi
\Bigr\},
\]
其中
\[
h(\theta):=-\theta\log\theta-(1-\theta)\log(1-\theta)
\]
为二元 Shannon 熵。该上确界由唯一点
\[
\theta(q)=\frac{1}{1+\varphi^{q+1}}
\]
取得，并满足
\[
\Psi'(q)=-\theta(q)\log\varphi.
\]
若进一步沿用定理 \ref{thm:xi-fold-escort-renyi-shannon-constants} 的 escort 记号 \(\mu_{m,q}\)，并定义
\[
\kappa_{m,q}^{\mathrm{esc}}
:=
\mathbb E_{\mu_{m,q}}\!\bigl[\log d_m^{\mathrm{bin}}(W_{m,q})\bigr],
\qquad
W_{m,q}\sim \mu_{m,q},
\]
则对每个固定 \(q\ge 0\) 都有
\[
\kappa_{m,q}^{\mathrm{esc}}
=
m\log\!\Bigl(\frac{2}{\varphi}\Bigr)
+
\Psi'(q)
+
O\!\Bigl(\Bigl(\frac{\varphi}{2}\Bigr)^m\Bigr).
\]
\end{theorem}
\begin{proof}
设
\[
F_q(\theta):=h(\theta)-\theta(q+1)\log\varphi.
\]
则
\[
F_q'(\theta)=\log\frac{1-\theta}{\theta}-(q+1)\log\varphi,
\qquad
F_q''(\theta)=-\frac{1}{\theta(1-\theta)}<0.
\]
因此 \(F_q\) 在 \((0,1)\) 上严格凹，且唯一驻点满足
\[
\frac{1-\theta}{\theta}=\varphi^{q+1},
\]
亦即
\[
\theta=\theta(q)=\frac{1}{1+\varphi^{q+1}}.
\]
这就给出了唯一极大点。再用标准二元自由能恒等式
\[
\sup_{\theta\in[0,1]}\{h(\theta)-\theta L\}=\log(1+e^{-L}),
\qquad
L\in\RR,
\]
并在此处代入 \(L=(q+1)\log\varphi\)，得到
\[
\sup_{\theta\in[0,1]}F_q(\theta)
=
\log\bigl(1+\varphi^{-(q+1)}\bigr).
\]
因此
\[
\log\!\Bigl(\frac{\varphi}{\sqrt5}\Bigr)+\sup_{\theta\in[0,1]}F_q(\theta)
=
\log\!\Bigl(\frac{\varphi+\varphi^{-q}}{\sqrt5}\Bigr)
=
\Psi(q).
\]

直接对 \(\Psi\) 求导可得
\[
\Psi'(q)
=
-\frac{\varphi^{-q}\log\varphi}{\varphi+\varphi^{-q}}
=
-\frac{\log\varphi}{1+\varphi^{q+1}}
=
-\theta(q)\log\varphi.
\]

最后，由定理 \ref{thm:fold-bin-two-state-asymptotic} 的一致对数展开，
\[
\log d_m^{\mathrm{bin}}(w)
=
m\log\!\Bigl(\frac{2}{\varphi}\Bigr)
-w_m\log\varphi
+
O\!\Bigl(\Bigl(\frac{\varphi}{2}\Bigr)^m\Bigr)
\]
在 \(w\in X_m\) 上一致成立。另一方面，命题 \ref{prop:fold-bin-escort-lastbit} 给出
\[
\mu_{m,q}(w_m=1)=\theta(q)+O\!\Bigl(\Bigl(\frac{\varphi}{2}\Bigr)^m\Bigr).
\]
对上式在 \(\mu_{m,q}\) 下取期望，得到
\[
\kappa_{m,q}^{\mathrm{esc}}
=
m\log\!\Bigl(\frac{2}{\varphi}\Bigr)
-\theta(q)\log\varphi
+
O\!\Bigl(\Bigl(\frac{\varphi}{2}\Bigr)^m\Bigr).
\]
再代入 \(\Psi'(q)=-\theta(q)\log\varphi\) 即得所述展开。证毕。
\end{proof}

\begin{theorem}[bin-fold 隐藏位预算、可见熵与均匀基线 KL 缺口共享同一常数核]\label{thm:xi-foldbin-hidden-budget-visible-entropy-kl-one-constant}
\leanverified{paper\_xi\_foldbin\_hidden\_budget\_visible\_entropy\_kl\_one\_constant}
沿用定理 \ref{thm:xi-foldbin-absolute-renyi-entropy-constant-closed} 与定理
\ref{thm:xi-foldbin-two-phase-free-energy-variational} 的记号。令
\[
\kappa_m:=\EE_{\mu_m}\!\bigl[\log d_m^{\mathrm{bin}}(W_m)\bigr],
\qquad
W_m\sim\mu_m,
\]
并记稳定层均匀分布为
\[
u_m(w):=\frac{1}{|X_m|}
\qquad (w\in X_m).
\]
则有三条共享同一隐藏位常数核的精确展开：
\[
\boxed{\;
\kappa_m
=
m\log\!\Bigl(\frac{2}{\varphi}\Bigr)
-\frac{\log\varphi}{\varphi\sqrt5}
+O\!\Bigl(\Bigl(\frac{\varphi}{2}\Bigr)^m\Bigr),\;}
\]
\[
\boxed{\;
H(\mu_m)
=
m\log\varphi
+\frac{\log\varphi}{\varphi\sqrt5}
+O\!\Bigl(\Bigl(\frac{\varphi}{2}\Bigr)^m\Bigr),\;}
\]
\[
\boxed{\;
\log|X_m|-H(\mu_m)
=
2\log\varphi-\frac12\log5-\frac{\log\varphi}{\varphi\sqrt5}
+O\!\Bigl(\Bigl(\frac{\varphi}{2}\Bigr)^m\Bigr).\;}
\]
因此
\[
\boxed{\;
D_{\mathrm{KL}}(\mu_m\Vert u_m)
=
2\log\varphi-\frac12\log5-\frac{\log\varphi}{\varphi\sqrt5}
+O\!\Bigl(\Bigl(\frac{\varphi}{2}\Bigr)^m\Bigr),\;}
\]
从而极限
\[
\boxed{\;
\lim_{m\to\infty}D_{\mathrm{KL}}(\mu_m\Vert u_m)
=
2\log\varphi-\frac12\log5-\frac{\log\varphi}{\varphi\sqrt5}>0.\;}
\]
存在且为严格正数。换言之，隐藏位预算、可见 Shannon 熵与相对于稳定层均匀基线的 KL 缺口并非三条彼此独立的数值链，而共享同一个刚性常数核
\[
\frac{\log\varphi}{\varphi\sqrt5}.
\]
\end{theorem}
\begin{proof}
在定理 \ref{thm:xi-foldbin-two-phase-free-energy-variational} 中取 \(q=1\)，并注意
\[
\kappa_m=\kappa_{m,1}^{\mathrm{esc}},
\qquad
\theta(1)=\frac{1}{1+\varphi^2}=\frac{1}{\varphi\sqrt5},
\]
便得到
\[
\kappa_m
=
m\log\!\Bigl(\frac{2}{\varphi}\Bigr)
-\frac{\log\varphi}{\varphi\sqrt5}
+O\!\Bigl(\Bigl(\frac{\varphi}{2}\Bigr)^m\Bigr).
\]
这证明了第一式。

第二式由 Shannon 熵恒等式
\[
H(\mu_m)=m\log2-\kappa_m
\]
直接推出。

再由稳定词计数
\[
|X_m|=F_{m+2}
=
\frac{\varphi^{m+2}}{\sqrt5}+O(\varphi^{-m}),
\]
可得
\[
\log|X_m|
=
(m+2)\log\varphi-\frac12\log5+O(\varphi^{-2m}).
\]
将其减去第二式，便得到第三式；而
\[
D_{\mathrm{KL}}(\mu_m\Vert u_m)=\log|X_m|-H(\mu_m)
\]
是均匀基线下的标准熵差公式，因此第四式与其极限立即成立。严格正性则可由
\[
\frac{1}{\varphi\sqrt5}\neq \varphi^{-2}
\]
以及推论 \ref{cor:xi-fold-uniform-baseline-classical-fdiv-bernoulli-collapse} 在
\[
f(u)=u\log u
\]
处的 Bernoulli 闭式看出。证毕。
\end{proof}

\begin{theorem}[bin-fold 容错阈值预算的 \texorpdfstring{$\log_2\varphi$}{log2 phi} 常数跃迁与指数斜率锁定]\label{thm:xi-foldbin-threshold-budget-logphi-jump}
\leanverified{paper\_xi\_foldbin\_threshold\_budget\_logphi\_jump}
沿用推论 \ref{cor:fold-bin-quantile-threshold-constant} 的记号，定义规范阈值预算
\[
P_m^{\mathrm{thr}}(\varepsilon):=2B^{\mathrm{bin}}_m(\varepsilon),
\qquad
\varepsilon\in(0,1),
\]
并记临界误差阈值
\[
\varepsilon_{\mathrm c}:=\frac{\varphi}{\sqrt5}.
\]
则有：
\begin{enumerate}
  \item 对任意固定
  \[
  0<\varepsilon_-<\varepsilon_{\mathrm c}<\varepsilon_+<1,
  \]
  都有
  \[
  \frac{P_m^{\mathrm{thr}}(\varepsilon_+)}{P_m^{\mathrm{thr}}(\varepsilon_-)}
  \longrightarrow
  \varphi^{-1},
  \]
  从而二进制对数预算差满足
  \[
  \log_2 P_m^{\mathrm{thr}}(\varepsilon_-)
  -
  \log_2 P_m^{\mathrm{thr}}(\varepsilon_+)
  \longrightarrow
  \log_2\varphi.
  \]
  \item 对任意固定
  \[
  \varepsilon\in(0,\varepsilon_{\mathrm c})\cup(\varepsilon_{\mathrm c},1),
  \]
  都有
  \[
  \lim_{m\to\infty}\frac1m\log P_m^{\mathrm{thr}}(\varepsilon)
  =
  \log\frac{2}{\varphi}
  =
  \lim_{m\to\infty}\frac1m\log d_{\max}^{\mathrm{bin}}(m).
  \]
\end{enumerate}
因此，阈值 \(\varepsilon_{\mathrm c}\) 只引入 \(\log_2\varphi\) 量级的常数预算跃迁，而不改变由最大纤维指数锁定的指数主斜率。
\end{theorem}
\begin{proof}
由推论 \ref{cor:fold-bin-quantile-threshold-constant}，
\[
\Bigl(\frac{\varphi}{2}\Bigr)^m B^{\mathrm{bin}}_m(\varepsilon)\longrightarrow
\begin{cases}
1,& 0<\varepsilon<\varepsilon_{\mathrm c},\\
\varphi^{-1},& \varepsilon_{\mathrm c}<\varepsilon<1.
\end{cases}
\]
乘以 \(2\) 即得
\[
P_m^{\mathrm{thr}}(\varepsilon_-)
=
2\Bigl(\frac{2}{\varphi}\Bigr)^m(1+o(1)),
\qquad
P_m^{\mathrm{thr}}(\varepsilon_+)
=
2\varphi^{-1}\Bigl(\frac{2}{\varphi}\Bigr)^m(1+o(1)),
\]
从而
\[
\frac{P_m^{\mathrm{thr}}(\varepsilon_+)}{P_m^{\mathrm{thr}}(\varepsilon_-)}
\longrightarrow
\varphi^{-1}.
\]
再取二进制对数即可得到
\[
\log_2 P_m^{\mathrm{thr}}(\varepsilon_-)
-
\log_2 P_m^{\mathrm{thr}}(\varepsilon_+)
\longrightarrow
\log_2\varphi.
\]
这就证明了第一条。

同样由上式，对任意固定
\[
\varepsilon\in(0,\varepsilon_{\mathrm c})\cup(\varepsilon_{\mathrm c},1)
\]
都有
\[
\lim_{m\to\infty}\frac1m\log P_m^{\mathrm{thr}}(\varepsilon)
=
\log\frac{2}{\varphi}
\]
另一方面，定理 \ref{thm:fold-bin-max-fiber-exponent} 给出
\[
\lim_{m\to\infty}\frac1m\log d_{\max}^{\mathrm{bin}}(m)
=
\log\frac{2}{\varphi}.
\]
两式合并即得所示指数斜率锁定。证毕。
\end{proof}

\leanverified{paper\_xi\_foldbin\_maxfiber\_upper\_bound\_exponential\_sharp\_nonsaturation}
\begin{theorem}[bin-fold 最大纤维组合上界的指数精确性与最终严格失饱和]\label{thm:xi-foldbin-maxfiber-upper-bound-exponential-sharp-nonsaturation}
沿用推论 \ref{cor:fold-bin-saturation} 的记号，定义缺口
\[
\delta_m:=F_{K(m)-m+2}-d_{\max}^{\mathrm{bin}}(m).
\]
则有：
\begin{enumerate}
  \item 存在 \(m_0\)，使得对全部 \(m\ge m_0\) 都有
  \[
  \delta_m>0.
  \]
  \item 并且
  \[
  \lim_{m\to\infty}\frac1m\log F_{K(m)-m+2}
  =
  \lim_{m\to\infty}\frac1m\log d_{\max}^{\mathrm{bin}}(m)
  =
  \log\frac{2}{\varphi},
  \]
  从而
  \[
  \lim_{m\to\infty}\frac1m\log\frac{F_{K(m)-m+2}}{d_{\max}^{\mathrm{bin}}(m)}=0.
  \]
\end{enumerate}
换言之，组合上界在指数尺度上是精确的，但在算术精确尺度上最终严格失饱和。
\end{theorem}
\begin{proof}
由注 \ref{rem:fold-bin-finite}，上界饱和
\[
d_{\max}^{\mathrm{bin}}(m)=F_{K(m)-m+2}
\]
只能发生有限次。故存在 \(m_0\)，使得对全部 \(m\ge m_0\) 都有
\[
\delta_m=F_{K(m)-m+2}-d_{\max}^{\mathrm{bin}}(m)>0.
\]
这就证明了第一条。

对第二条，由定义 \eqref{eq:K-of-m}
\[
F_{K(m)+1}\le 2^m-1<F_{K(m)+2}.
\]
结合 Binet 渐近
\[
\log F_n=n\log\varphi+O(1)
\]
可得
\[
K(m)=m\log_{\varphi}2+O(1).
\]
于是
\[
\log F_{K(m)-m+2}
=
\bigl(K(m)-m+2\bigr)\log\varphi+O(1)
=
m(\log2-\log\varphi)+O(1),
\]
从而
\[
\lim_{m\to\infty}\frac1m\log F_{K(m)-m+2}
=
\log\frac{2}{\varphi}.
\]
另一方面，由定理 \ref{thm:fold-bin-max-fiber-exponent}，
\[
\lim_{m\to\infty}\frac1m\log d_{\max}^{\mathrm{bin}}(m)
=
\log\frac{2}{\varphi}.
\]
两式相减即得
\[
\lim_{m\to\infty}\frac1m\log\frac{F_{K(m)-m+2}}{d_{\max}^{\mathrm{bin}}(m)}=0.
\]
证毕。
\end{proof}

\paragraph{escort Rényi 熵常数的 \texorpdfstring{$\theta(q)$}{theta(q)}-闭式化}
在 Rényi 熵率塌缩与 escort 常数谱已经确立之后，温度参数 \(q\) 的全部 \(O(1)\) 熵信息还可继续压缩为末位 Bernoulli 质量
\(
\theta(q)=1/(1+\varphi^{q+1})
\)。

\begin{theorem}[escort Rényi 熵的速率塌缩与常数项全解]\label{thm:xi-fold-escort-renyi-constant-theta-closed}
沿用定理 \ref{thm:xi-fold-escort-renyi-shannon-constants} 的记号。记
\[
\theta(q):=\frac{1}{1+\varphi^{q+1}}.
\]
则对任意固定 \(q\ge 0\) 与任意 \(r>0\)、\(r\neq 1\)，有
\[
H_r(\mu_{m,q})
=
m\log\varphi
-\frac12\log 5
+
\frac{1}{1-r}
\log\!\Bigl(
\varphi^{1-r}(1-\theta(q))^r+\theta(q)^r
\Bigr)
+
O\!\Bigl(\Bigl(\frac{\varphi}{2}\Bigr)^m\Bigr).
\]
等价地，
\[
H_r(\mu_{m,q})
=
m\log\varphi
-\frac12\log 5
+
\frac{1}{1-r}
\Bigl[
\log\!\bigl(1+\varphi^{1+qr}\bigr)
-r\log\!\bigl(1+\varphi^{q+1}\bigr)
\Bigr]
+
O\!\Bigl(\Bigl(\frac{\varphi}{2}\Bigr)^m\Bigr).
\]
相应的 Shannon 熵满足
\[
H(\mu_{m,q})
=
m\log\varphi
-\frac12\log 5
+
h_2(\theta(q))
+
(1-\theta(q))\log\varphi
+
O\!\Bigl(\Bigl(\frac{\varphi}{2}\Bigr)^m\Bigr),
\]
其中
\[
h_2(p):=-p\log p-(1-p)\log(1-p).
\]
因此 escort 温度 \(q\) 的全部 \(O(1)\) 熵信息都被压缩为单变量 \(\theta(q)\)。
\end{theorem}
\begin{proof}
由定理 \ref{thm:xi-fold-escort-renyi-shannon-constants}，
\[
H_r(\mu_{m,q})
=
m\log\varphi
+
\frac{1}{1-r}\left[
\log\!\Bigl(\frac{\varphi+\varphi^{-qr}}{\sqrt5}\Bigr)
-r\log\!\Bigl(\frac{\varphi+\varphi^{-q}}{\sqrt5}\Bigr)
\right]
+
O\!\Bigl(\Bigl(\frac{\varphi}{2}\Bigr)^m\Bigr).
\]
再用
\[
\varphi+\varphi^{-a}=\varphi^{-a}(1+\varphi^{a+1}),
\]
便得到第二个 Rényi 公式。

另一方面，由
\[
1-\theta(q)=\frac{\varphi^{q+1}}{1+\varphi^{q+1}}
\]
可直接化简出
\[
\varphi^{1-r}(1-\theta(q))^r+\theta(q)^r
=
\frac{1+\varphi^{1+qr}}{(1+\varphi^{q+1})^r},
\]
故第一个 Rényi 公式与第二个完全等价。

对 Shannon 熵，定理 \ref{thm:xi-fold-escort-renyi-shannon-constants} 同样给出
\[
H(\mu_{m,q})
=
m\log\varphi
+
\log\!\Bigl(\frac{\varphi+\varphi^{-q}}{\sqrt5}\Bigr)
+
\frac{q\log\varphi}{1+\varphi^{q+1}}
+
O\!\Bigl(\Bigl(\frac{\varphi}{2}\Bigr)^m\Bigr).
\]
再注意到
\[
h_2(\theta(q))
=
\log\!\bigl(1+\varphi^{q+1}\bigr)
-(q+1)(1-\theta(q))\log\varphi,
\]
于是
\[
-\frac12\log 5
+
h_2(\theta(q))
+
(1-\theta(q))\log\varphi
=
\log\!\Bigl(\frac{\varphi+\varphi^{-q}}{\sqrt5}\Bigr)
+
\frac{q\log\varphi}{1+\varphi^{q+1}}.
\]
这就给出了 Shannon 常数项的所述写法。证毕。
\end{proof}

\paragraph{冻结锥平台、极大纤维骨架与纤维指数自平均}
在固定冻结相的 Rényi 熵率塌缩、基态质量指数集中以及精确反演位率阈值已经建立之后，还可把冻结锥上的平台律、极大纤维骨架投影、辨路径预算相变与纤维指数自平均进一步压缩为下列四条结果。

\begin{theorem}[冻结锥上的 escort Rényi 平台]\label{thm:xi-fixed-freezing-cone-escort-renyi-plateau}
沿用定理 \ref{thm:xi-fixed-freezing-escort-renyi-spectrum-collapse} 的记号。对 \(a>0\) 定义 escort 分布
\[
\pi_{m,a}(x):=\frac{d_m(x)^a}{S_a(m)},
\qquad
S_a(m):=\sum_{x\in X_m}d_m(x)^a,
\]
并对 \(q>0\)、\(q\neq 1\) 定义 escort Rényi 熵
\[
H_q^{\mathrm{esc}}(m;a):=
\frac{1}{1-q}\log\sum_{x\in X_m}\pi_{m,a}(x)^q.
\]
若
\[
a>a_{\mathrm c},
\qquad
aq>a_{\mathrm c},
\]
则极限
\[
h_q^{\mathrm{esc}}(a):=
\lim_{m\to\infty}\frac1m H_q^{\mathrm{esc}}(m;a)
\]
存在，并且满足
\[
h_q^{\mathrm{esc}}(a)=g_*.
\]
\end{theorem}
\begin{proof}
由定义
\[
\sum_{x\in X_m}\pi_{m,a}(x)^q
=
\frac{S_{aq}(m)}{S_a(m)^q},
\]
故
\[
\frac1m H_q^{\mathrm{esc}}(m;a)
=
\frac{1}{1-q}
\left(
\frac1m\log S_{aq}(m)-q\frac1m\log S_a(m)
\right).
\]
当 \(a>a_{\mathrm c}\) 且 \(aq>a_{\mathrm c}\) 时，定理 \ref{thm:fold-pressure-freezing-threshold} 给出
\[
\lim_{m\to\infty}\frac1m\log S_t(m)=t\alpha_*+g_*
\qquad
(t>a_{\mathrm c}).
\]
将 \(t=a\) 与 \(t=aq\) 代回，便得到
\[
h_q^{\mathrm{esc}}(a)
=
\frac{(aq\alpha_*+g_*)-q(a\alpha_*+g_*)}{1-q}
=
g_*.
\]
证毕。
\end{proof}

\begin{theorem}[冻结相的极大纤维骨架投影]\label{thm:xi-fixed-freezing-maxfiber-skeleton-projection}
\leanverified{paper\_xi\_fixed\_freezing\_maxfiber\_skeleton\_projection}
记
\[
A_m^{\max}:=\{x\in X_m:\ d_m(x)=M_m\},
\qquad
K_m:=|A_m^{\max}|,
\]
并令 \(\nu_m\) 为 \(A_m^{\max}\) 上的均匀分布，即
\[
\nu_m(x):=
\begin{cases}
K_m^{-1}, & x\in A_m^{\max},\\
0, & x\notin A_m^{\max}.
\end{cases}
\]
若 \(a>a_{\mathrm c}\)，则对任意满足
\[
\|f_m\|_{\infty}\le 1
\qquad
\bigl(f_m:X_m\to\CC\bigr)
\]
的测试函数，都有
\[
\left|
\sum_{x\in X_m}\pi_{m,a}(x)f_m(x)
-
\sum_{x\in X_m}\nu_m(x)f_m(x)
\right|
\le
2\exp\!\bigl(-m\Delta(a)+o(m)\bigr),
\]
其中
\[
\Delta(a):=a(\alpha_*-\alpha_*^{(2)})-(\log\varphi-g_*)>0.
\]
\end{theorem}
\begin{proof}
由定理 \ref{thm:xi-fixed-freezing-escort-renyi-spectrum-collapse} 的第四条，
\[
\|\pi_{m,a}-\nu_m\|_{\mathrm{TV}}
=
1-\pi_{m,a}(A_m^{\max})
\le
\exp\!\bigl(-m\Delta(a)+o(m)\bigr).
\]
再由总变差距离的对偶表述，对任意 \(\|f_m\|_\infty\le 1\) 都有
\[
\left|
\sum_{x\in X_m}\pi_{m,a}(x)f_m(x)
-
\sum_{x\in X_m}\nu_m(x)f_m(x)
\right|
\le
2\|\pi_{m,a}-\nu_m\|_{\mathrm{TV}}.
\]
将前式代回即得结论。证毕。
\end{proof}

\begin{theorem}[冻结相中的辨路径位率相变]\label{thm:xi-fixed-freezing-path-budget-phase-transition}
设 \(a>a_{\mathrm c}\)。对整数预算 \(B\ge 0\)，定义 escort-可解码质量
\[
\Gamma_{m,a}(B):=
\pi_{m,a}\bigl(\{x\in X_m:\ d_m(x)\le 2^B\}\bigr).
\]
令
\[
B_m:=\lfloor \rho m\rfloor.
\]
并记
\[
\Delta(a):=a(\alpha_*-\alpha_*^{(2)})-(\log\varphi-g_*)>0.
\]
则存在严格相变：
\begin{enumerate}
  \item 若
  \[
  \rho<\frac{\alpha_*}{\log 2},
  \]
  则
  \[
  \Gamma_{m,a}(B_m)\le \exp\!\bigl(-m\Delta(a)+o(m)\bigr)\xrightarrow[m\to\infty]{}0.
  \]
  \item 若
  \[
  \rho>\frac{\alpha_*}{\log 2},
  \]
  则对充分大的 \(m\) 有
  \[
  \Gamma_{m,a}(B_m)=1.
  \]
\end{enumerate}
\end{theorem}
\begin{proof}
若 \(\rho<\alpha_*/\log 2\)，则由
\[
B_m=\rho m+O(1),
\qquad
M_m=\exp(\alpha_*m+o(m))
\]
可得对充分大的 \(m\)，
\[
2^{B_m}=\exp\!\bigl((\rho\log 2)m+o(m)\bigr)<M_m.
\]
因此每个 \(x\in A_m^{\max}\) 都满足 \(d_m(x)=M_m>2^{B_m}\)，从而
\[
\Gamma_{m,a}(B_m)\le \pi_{m,a}(X_m\setminus A_m^{\max}).
\]
再由定理 \ref{thm:xi-fixed-freezing-escort-renyi-spectrum-collapse} 的第四条，
\[
\pi_{m,a}(X_m\setminus A_m^{\max})
=
1-\pi_{m,a}(A_m^{\max})
\le
\exp\!\bigl(-m\Delta(a)+o(m)\bigr),
\]
即得第一条。

若 \(\rho>\alpha_*/\log 2\)，则对充分大的 \(m\) 有
\[
2^{B_m}>M_m.
\]
由于对一切 \(x\in X_m\) 都有 \(d_m(x)\le M_m\)，故
\[
\{x\in X_m:\ d_m(x)\le 2^{B_m}\}=X_m,
\]
于是
\[
\Gamma_{m,a}(B_m)=1.
\]
证毕。
\end{proof}

\begin{theorem}[冻结相中的纤维指数自平均]\label{thm:xi-fixed-freezing-fiber-exponent-selfaveraging}
设 \(a>a_{\mathrm c}\)，并定义
\[
Y_m(x):=\frac{1}{m}\log d_m(x),
\qquad
x\in X_m.
\]
对任意满足
\[
a+t>a_{\mathrm c}
\]
的实数 \(t\)，极限
\[
\Lambda_a(t):=
\lim_{m\to\infty}
\frac1m
\log \EE_{\pi_{m,a}}\!\left[e^{mtY_m}\right]
\]
存在，且
\[
\Lambda_a(t)=\alpha_* t.
\]
进一步，对任意有限 \(p>0\)，都有
\[
\lim_{m\to\infty}
\EE_{\pi_{m,a}}
\left[
\left|Y_m-\alpha_*\right|^p
\right]
=0.
\]
\end{theorem}
\begin{proof}
先计算指数母函数。由定义
\[
\EE_{\pi_{m,a}}\!\left[e^{mtY_m}\right]
=
\sum_{x\in X_m}\frac{d_m(x)^a}{S_a(m)}\,d_m(x)^t
=
\frac{S_{a+t}(m)}{S_a(m)}.
\]
因此
\[
\frac1m
\log \EE_{\pi_{m,a}}\!\left[e^{mtY_m}\right]
=
\frac1m\log S_{a+t}(m)-\frac1m\log S_a(m).
\]
由于 \(a>a_{\mathrm c}\) 且 \(a+t>a_{\mathrm c}\)，定理 \ref{thm:fold-pressure-freezing-threshold} 给出
\[
\lim_{m\to\infty}\frac1m\log S_s(m)=s\alpha_*+g_*
\qquad
(s>a_{\mathrm c}),
\]
故
\[
\Lambda_a(t)
=
(a+t)\alpha_*+g_*-(a\alpha_*+g_*)
=
\alpha_* t.
\]

再证 \(L^p\) 收敛。由 \(1\le d_m(x)\le 2^m\)，对全部 \(x\in X_m\) 都有
\[
0\le Y_m(x)\le \log 2.
\]
定义有界测试函数
\[
f_m(x):=\frac{|Y_m(x)-\alpha_*|^p}{(\log 2+|\alpha_*|)^p},
\]
则 \(\|f_m\|_\infty\le 1\)。另一方面，在 \(A_m^{\max}\) 上 \(Y_m\) 常值，故
\[
\sum_{x\in X_m}\nu_m(x)f_m(x)
=
\left|
\frac1m\log M_m-\alpha_*
\right|^p
(\log 2+|\alpha_*|)^{-p}.
\]
由定理 \ref{thm:xi-fixed-freezing-maxfiber-skeleton-projection}，
\[
\left|
\EE_{\pi_{m,a}}[f_m]-\sum_{x\in X_m}\nu_m(x)f_m(x)
\right|
\le
2\exp\!\bigl(-m\Delta(a)+o(m)\bigr).
\]
再乘回 \((\log 2+|\alpha_*|)^p\)，便得到
\[
\EE_{\pi_{m,a}}\!\left[|Y_m-\alpha_*|^p\right]
\le
\left|
\frac1m\log M_m-\alpha_*
\right|^p
+2(\log 2+|\alpha_*|)^p
\exp\!\bigl(-m\Delta(a)+o(m)\bigr).
\]
由 \(\frac1m\log M_m\to\alpha_*\)，右端趋于 \(0\)，故结论成立。证毕。
\end{proof}

\paragraph{冻结尾部的凸延拓刚性、Legendre 硬壁与 Abel 有限部余项}
在定理 \ref{thm:xi-time-part57b-pressure-spectrum-discrete-convexity-monotone-excess} 已把整数压力谱锚定为离散凸序列、定理 \ref{thm:xi-freezing-affine-rigidity-two-point-recovery} 已把冻结相整数尾部压缩为两点即可恢复的仿射律之后，还可继续把连续凸延拓、Legendre 对偶边界与 Möbius--Euler 截断余项进一步收束为下列四条后果。

\begin{theorem}[冻结阈值之后任意凸压力延拓的仿射刚性]\label{thm:xi-time-part57f-convex-extension-affine-tail}
在严格基态间隙
\[
\alpha_*^{(2)}<\alpha_*
\]
下，记
\[
a_{\mathrm c}:=\frac{\log\varphi-g_*}{\alpha_*-\alpha_*^{(2)}},
\qquad
a_0:=\lfloor a_{\mathrm c}\rfloor+1.
\]
设
\[
\Lambda:[a_0,\infty)\to\RR
\]
为任意凸函数，并满足
\[
\Lambda(n)=P_n
\qquad
\bigl(n\in\ZZ,\ n\ge a_0\bigr).
\]
则必有
\[
\boxed{\;
\Lambda(t)=\alpha_* t+g_*
\qquad
(t\ge a_0).\;}
\]
特别地，任何保持凸性并在冻结区整数锚点与离散压力序列对齐的连续压力外推，在阈值之后都不再允许残留曲率；因此，定理 \ref{thm:xi-projective-pressure-canonical-convex-majorant} 所给出的连续压力接口一旦在这些整数锚点上实现精确对齐，其尾部也只能是同一条仿射直线。
\end{theorem}
\begin{proof}
由定理 \ref{thm:fold-pressure-freezing-threshold}，对每个整数 \(n\ge a_0\) 都有
\[
P_n=n\alpha_*+g_*,
\]
故
\[
G(t):=\Lambda(t)-(\alpha_* t+g_*)
\]
是定义在 \([a_0,\infty)\) 上的凸函数，并满足
\[
G(n)=0
\qquad
\bigl(n\in\ZZ,\ n\ge a_0\bigr).
\]

先看首个区间 \([a_0,a_0+1]\)。端点 \(t=a_0,a_0+1\) 时结论显然，故只需取
\[
t\in(a_0,a_0+1).
\]
由凸函数割线斜率单调性，对
\[
t<a_0+1<a_0+2
\]
应用
\[
\frac{G(a_0+1)-G(t)}{a_0+1-t}
\le
\frac{G(a_0+2)-G(a_0+1)}{1},
\]
便得
\[
\frac{-G(t)}{a_0+1-t}\le 0,
\]
从而 \(G(t)\ge 0\)。另一方面，由
\[
a_0\le t\le a_0+1
\]
与凸性，
\[
G(t)\le
\frac{a_0+1-t}{1}G(a_0)+\frac{t-a_0}{1}G(a_0+1)=0,
\]
故 \(G(t)=0\)。于是
\[
G\equiv 0
\qquad\text{于 }[a_0,a_0+1].
\]

再取任意整数 \(n\ge a_0+1\)。端点 \(t=n,n+1\) 时同样平凡，故只需考虑
\[
t\in(n,n+1).
\]
同样由割线斜率单调性，对
\[
n-1<n\le t
\]
有
\[
\frac{G(n)-G(n-1)}{1}
\le
\frac{G(t)-G(n)}{t-n},
\]
即
\[
0\le \frac{G(t)}{t-n},
\]
故 \(G(t)\ge 0\)。而由 \(G(n)=G(n+1)=0\) 与凸性，
\[
G(t)\le
\frac{n+1-t}{1}G(n)+\frac{t-n}{1}G(n+1)=0.
\]
因此 \(G(t)=0\)。这表明 \(G\equiv 0\) 于每个区间 \([n,n+1]\) 上成立。并合全部相邻整数区间，即得
\[
G(t)\equiv 0
\qquad
(t\ge a_0),
\]
亦即
\[
\Lambda(t)=\alpha_* t+g_*.
\]
证毕。
\end{proof}

\begin{theorem}[冻结尾部的 Legendre 硬壁与边界代价]\label{thm:xi-time-part57f-legendre-hard-wall-boundary-cost}
沿用定理 \ref{thm:xi-time-part57f-convex-extension-affine-tail} 的记号。在 \(X_m\) 上取均匀分布 \(u_m\)，并定义随机变量
\[
Y_m(x):=\frac{1}{m}\log d_m(x)
\qquad
(x\in X_m).
\]
则对每个整数 \(a\ge 1\)，有
\[
\frac1m\log \EE_{u_m}\!\left[e^{amY_m}\right]
=
\frac1m\log\frac{S_a(m)}{|X_m|}
\longrightarrow
P_a-\log\varphi.
\]
再设
\[
I(y):=
\sup_{t\ge a_0}
\Bigl\{
ty-\bigl(\Lambda(t)-\log\varphi\bigr)
\Bigr\}
\]
为任一此类凸延拓 \(\Lambda\) 所对应的 Legendre--Fenchel 对偶。则成立：
\begin{enumerate}
  \item 对任意 \(y>\alpha_*\)，有
  \[
  \boxed{\;I(y)=+\infty.\;}
  \]
  \item 在边界点 \(y=\alpha_*\) 处，恰有
  \[
  \boxed{\;I(\alpha_*)=\log\varphi-g_*.\;}
  \]
\end{enumerate}
换言之，\(\alpha_*\) 是冻结尾部在 Legendre 对偶中的硬上壁，而抵达该边界的指数代价恰由环境熵率与基态简并指数之差给出。
\end{theorem}
\begin{proof}
由定义
\[
\EE_{u_m}\!\left[e^{amY_m}\right]
=
\frac{1}{|X_m|}
\sum_{x\in X_m} d_m(x)^a
=
\frac{S_a(m)}{|X_m|}.
\]
又因 \(|X_m|=F_{m+2}\) 且
\[
\lim_{m\to\infty}\frac1m\log F_{m+2}=\log\varphi,
\]
故
\[
\frac1m\log \EE_{u_m}\!\left[e^{amY_m}\right]
=
\frac1m\log S_a(m)-\frac1m\log|X_m|
\longrightarrow
P_a-\log\varphi.
\]

另一方面，由定理 \ref{thm:xi-time-part57f-convex-extension-affine-tail}，
\[
\Lambda(t)=\alpha_* t+g_*
\qquad
(t\ge a_0).
\]
于是对任意 \(t\ge a_0\)，都有
\[
ty-\bigl(\Lambda(t)-\log\varphi\bigr)
=
t(y-\alpha_*)+\log\varphi-g_*.
\]
若 \(y>\alpha_*\)，右端随 \(t\to+\infty\) 线性发散到 \(+\infty\)，从而
\[
I(y)=+\infty.
\]
若 \(y=\alpha_*\)，则上式对所有 \(t\ge a_0\) 都恒等于
\[
\log\varphi-g_*,
\]
故其上确界亦为该常数，即
\[
I(\alpha_*)=\log\varphi-g_*.
\]
证毕。
\end{proof}

\begin{theorem}[Abel 有限部常数的内禀 \texorpdfstring{$\eta$}{eta}-截断误差界]\label{thm:xi-time-part57f-abel-finite-part-eta-tail}
沿用命题 \ref{prop:real-input-40-logM-infty} 的记号。记
\[
H(x):=(1-x)(1-6x+9x^2-x^3-x^4)
=
1-7x+15x^2-10x^3+x^5.
\]
并定义
\[
\eta:=\min_{0\le x\le b^2}H(x)>0.
\]
对每个整数 \(N\ge 2\)，记
\[
\log\mathfrak M_\infty^{(N)}
:=
\log C_\infty
-\sum_{m=2}^{N}\frac{\mu(m)}{m}\log H(b^m).
\]
则截断误差满足显式上界
\[
\boxed{\;
\bigl|\log\mathfrak M_\infty-\log\mathfrak M_\infty^{(N)}\bigr|
\le
\frac{7}{\eta(1-b)}\cdot \frac{b^{N+1}}{N+1}.\;}
\]
因此，Abel 有限部常数的 Möbius--Euler 展开不只几何收敛，而且其余项可由区间 \([0,b^2]\) 上的内禀最小值 \(\eta\) 直接审计。
\end{theorem}
\begin{proof}
由于 \(b\) 是多项式
\[
1-6x+9x^2-x^3-x^4
\]
的最大正根，故对 \(0\le x\le b^2<b\) 有
\[
1-6x+9x^2-x^3-x^4>0
\qquad\text{且}\qquad
1-x>0.
\]
因此
\[
H(x)>0
\qquad
(0\le x\le b^2),
\]
从而 \(\eta>0\)。

又由命题 \ref{prop:real-input-40-logM-infty} 的同一根定位，可知
\[
\frac14<b<\frac13,
\]
于是
\[
b^m\le b^2<\frac19
\qquad
(m\ge 2).
\]
对 \(x\in[0,b^2]\) 直接计算得
\[
1-H(x)=7x-15x^2+10x^3-x^5\le 7x.
\]
另一方面，由 \(x\le 1/9\) 可知
\[
1-H(x)
=
x\bigl(7-15x+10x^2-x^4\bigr)
\ge
x\Bigl(7-\frac{15}{9}\Bigr)\ge 0,
\]
故
\[
0<H(x)\le 1
\qquad
(0\le x\le b^2).
\]
因此对每个 \(m\ge 2\)，都可应用基本不等式
\[
-\log t\le \frac{1-t}{t}
\qquad
(0<t\le 1),
\]
得到
\[
\bigl|\log H(b^m)\bigr|
=
-\log H(b^m)
\le
\frac{1-H(b^m)}{H(b^m)}
\le
\frac{7b^m}{\eta}.
\]
于是
\[
\bigl|\log\mathfrak M_\infty-\log\mathfrak M_\infty^{(N)}\bigr|
\le
\sum_{m>N}\frac{|\mu(m)|}{m}\,\bigl|\log H(b^m)\bigr|
\le
\frac{7}{\eta}\sum_{m>N}\frac{b^m}{m}.
\]
而对 \(m>N\) 有 \(1/m\le 1/(N+1)\)，故
\[
\sum_{m>N}\frac{b^m}{m}
\le
\frac{1}{N+1}\sum_{m>N}b^m
=
\frac{1}{N+1}\cdot \frac{b^{N+1}}{1-b}.
\]
代回即得所示盒装上界。证毕。
\end{proof}

\begin{proposition}[Möbius--Euler 余项的首项由下一平方自由层主导]\label{prop:xi-time-part57f-mobius-euler-next-squarefree-layer}
沿用定理 \ref{thm:xi-time-part57f-abel-finite-part-eta-tail} 的记号。定义截断余项
\[
R_N:=
\log\mathfrak M_\infty-\log\mathfrak M_\infty^{(N)}
=
-\sum_{m>N}\frac{\mu(m)}{m}\log H(b^m).
\]
则当 \(N\to\infty\) 时，有统一渐近展开
\[
\boxed{\;
R_N
=
7\sum_{m>N}\frac{\mu(m)}{m}b^m
+
O\!\left(\frac{b^{2N+2}}{N+1}\right).\;}
\]
特别地，若 \(N+1\) 为平方自由数，则
\[
\boxed{\;
R_N
=
\frac{7\,\mu(N+1)}{N+1}b^{N+1}
+
O\!\left(\frac{b^{N+2}}{N+2}\right).\;}
\]
因此，余项的首个非零主项由下一层平方自由指数唯一决定，其符号直接由 \(\mu(N+1)\) 控制。
\end{proposition}
\begin{proof}
由
\[
H(x)=1-7x+O(x^2)
\qquad
(x\to 0),
\]
可得
\[
\log H(x)=-7x+O(x^2)
\qquad
(x\to 0).
\]
把 \(x=b^m\) 代入并带回余项展开，得到
\[
\begin{aligned}
R_N
&=
-\sum_{m>N}\frac{\mu(m)}{m}\Bigl(-7b^m+O(b^{2m})\Bigr)\\
&=
7\sum_{m>N}\frac{\mu(m)}{m}b^m
+
O\!\left(\sum_{m>N}\frac{b^{2m}}{m}\right).
\end{aligned}
\]
又因
\[
\sum_{m>N}\frac{b^{2m}}{m}
\le
\frac{1}{N+1}\sum_{m>N}b^{2m}
=
O\!\left(\frac{b^{2N+2}}{N+1}\right),
\]
便得第一条渐近式。

若 \(N+1\) 为平方自由数，则 \(\mu(N+1)\neq 0\)，从而
\[
\sum_{m>N}\frac{\mu(m)}{m}b^m
=
\frac{\mu(N+1)}{N+1}b^{N+1}
+
O\!\left(\sum_{m\ge N+2}\frac{b^m}{m}\right).
\]
而
\[
\sum_{m\ge N+2}\frac{b^m}{m}
\le
\frac{1}{N+2}\sum_{m\ge N+2}b^m
=
O\!\left(\frac{b^{N+2}}{N+2}\right).
\]
将其代回第一条式子，即得
\[
R_N
=
\frac{7\,\mu(N+1)}{N+1}b^{N+1}
+
O\!\left(\frac{b^{N+2}}{N+2}\right).
\]
证毕。
\end{proof}

\noindent 结合定理 \ref{thm:xi-time-part57b-pressure-spectrum-discrete-convexity-monotone-excess} 的离散凸性、定理 \ref{thm:xi-time-part57f-convex-extension-affine-tail} 的凸延拓刚性、定理 \ref{thm:xi-time-part57f-legendre-hard-wall-boundary-cost} 的 Legendre 边界代价，以及定理 \ref{thm:xi-time-part57f-abel-finite-part-eta-tail} 与命题 \ref{prop:xi-time-part57f-mobius-euler-next-squarefree-layer} 对 Möbius--Euler 余项的显式控制可见，冻结相在压力对偶与 Abel 有限部两侧都呈现出同一种硬壁结构：一侧把可达纤维指数锁定在 \(y\le \alpha_*\) 并在边界处留下精确代价 \(\log\varphi-g_*\)，另一侧则把有限部余项压缩为由下一平方自由层主导的首项与可闭式审计的尾界。这为进一步把谱硬壁与有限部硬余项统一到同一 Legendre--Witt 框架提供了直接接口。

\paragraph{冻结预算相图、尾计数平台与 escort 预言机竞争律}
在冻结锥上的极大纤维骨架、连续容量曲线的 Mellin--Stieltjes 对偶与全微态位率阈值已经建立之后，还可把顶端平台刚性、容量缺口闭式与 escort 预算竞争进一步压缩为下列四条结果。它们把最大纤维指数 \(\alpha_*\)、最大层复杂度 \(g_*\)、次大纤维指数 \(\alpha_*^{(2)}\) 与预算速率 \(\beta\) 统一组织为一张可审计的冻结预算相图。

\begin{theorem}[尾计数函数的顶端平台刚性]\label{thm:xi-time-part57ea-tailcount-top-platform-rigidity}
\leanverified{paper\_xi\_time\_part57ea\_tailcount\_top\_platform\_rigidity}
沿用定理 \ref{thm:xi-time-part57b-escort-maxfiber-layer-concentration} 的记号。设
\[
\alpha_*^{(2)}<\beta<\alpha_*,
\qquad
B_m:=\left\lfloor \frac{\beta}{\log 2}\,m\right\rfloor,
\qquad
T_m:=2^{B_m}.
\]
并记连续阈值尾计数函数
\[
N_m(T):=\#\{x\in X_m:\ d_m(x)\ge T\}.
\]
再记
\[
A_m^{\max}:=\{x\in X_m:\ d_m(x)=M_m\}.
\]
则对充分大的 \(m\)，有严格恒等式
\[
\boxed{\;
N_m(T_m)=K_m.\;}
\]
换言之，只要指数阈值 \(T_m\) 最终落在 \(M_m^{(2)}\) 与 \(M_m\) 之间，则超过该阈值的纤维恰好就是全部最大纤维。
\end{theorem}
\begin{proof}
由 \(\beta<\alpha_*\) 及
\[
\frac1m\log M_m\to\alpha_*,
\]
可得对充分大的 \(m\)，
\[
\frac1m\log T_m=\frac{B_m\log 2}{m}
=
\beta+O(m^{-1})
<
\alpha_*+o(1)
=
\frac1m\log M_m,
\]
从而
\[
T_m<M_m.
\]
另一方面，由 \(\beta>\alpha_*^{(2)}\) 及
\[
\limsup_{m\to\infty}\frac1m\log M_m^{(2)}=\alpha_*^{(2)},
\]
可得对充分大的 \(m\)，
\[
\frac1m\log M_m^{(2)}<\beta+O(m^{-1})=\frac1m\log T_m,
\]
从而
\[
M_m^{(2)}<T_m.
\]
于是对 \(x\in A_m^{\max}\) 有 \(d_m(x)=M_m>T_m\)，而对 \(x\notin A_m^{\max}\) 有 \(d_m(x)\le M_m^{(2)}<T_m\)。因此
\[
N_m(T_m)=\#A_m^{\max}=K_m.
\]
证毕。
\end{proof}

\begin{theorem}[预言机容量缺口的顶端精确公式]\label{thm:xi-time-part57ea-oracle-capacity-defect-top-formula}
在定理 \ref{thm:xi-time-part57ea-tailcount-top-platform-rigidity} 的假设下，对充分大的 \(m\) 有
\[
\boxed{\;
C_m(B_m)=2^m-K_m(M_m-T_m),\;}
\]
其中
\[
C_m(B):=\sum_{x\in X_m}\min(d_m(x),2^B),
\qquad
\mathrm{Succ}_m(B):=\frac{C_m(B)}{2^m}.
\]
从而
\[
\boxed{\;
1-\mathrm{Succ}_m(B_m)
=
\frac{K_m(M_m-T_m)}{2^m}.\;}
\]
并且
\[
\boxed{\;
\lim_{m\to\infty}\frac1m
\log\bigl(1-\mathrm{Succ}_m(B_m)\bigr)
=
\alpha_*+g_*-\log 2.\;}
\]
\end{theorem}
\begin{proof}
由定理 \ref{thm:xi-time-part57ea-tailcount-top-platform-rigidity}，对充分大的 \(m\) 有
\[
M_m^{(2)}<T_m<M_m.
\]
于是对 \(x\notin A_m^{\max}\) 有 \(\min(d_m(x),T_m)=d_m(x)\)，对 \(x\in A_m^{\max}\) 有 \(\min(d_m(x),T_m)=T_m\)。故
\[
C_m(B_m)
=
\sum_{x\notin A_m^{\max}}d_m(x)+\sum_{x\in A_m^{\max}}T_m.
\]
再用总质量恒等式 \(\sum_x d_m(x)=2^m\)，得到
\[
C_m(B_m)=2^m-K_mM_m+K_mT_m
=
2^m-K_m(M_m-T_m).
\]
两边除以 \(2^m\) 即得成功率缺口公式。

最后，由 \(\beta<\alpha_*\) 可知
\[
\frac{T_m}{M_m}
=
\exp\!\bigl(-(\alpha_*-\beta)m+o(m)\bigr)\to 0,
\]
故
\[
M_m-T_m=M_m(1+o(1)).
\]
再结合
\[
M_m=\exp(\alpha_*m+o(m)),
\qquad
K_m=\exp(g_*m+o(m)),
\]
得到
\[
1-\mathrm{Succ}_m(B_m)
=
\exp\!\bigl((\alpha_*+g_*-\log 2)m+o(m)\bigr),
\]
从而结论成立。证毕。
\end{proof}

\begin{theorem}[冻结相 escort 预言机成功率的精确分解与竞争上界]\label{thm:xi-time-part57ea-frozen-escort-oracle-success-decomposition}
设 \(a>a_{\mathrm c}\)，并沿用定理 \ref{thm:xi-fixed-freezing-escort-renyi-spectrum-collapse} 的 escort 分布
\[
\pi_{m,a}(x):=\frac{d_m(x)^a}{S_a(m)}.
\]
定义 escort 预算 \(B\) 下的预言机成功率
\[
\mathrm{Succ}_{m,a}^{\mathrm{esc}}(B)
:=
\sum_{x\in X_m}\pi_{m,a}(x)\min\!\left(1,\frac{2^B}{d_m(x)}\right).
\]
仍取
\[
\alpha_*^{(2)}<\beta<\alpha_*,
\qquad
B_m=\left\lfloor \frac{\beta}{\log 2}\,m\right\rfloor,
\qquad
T_m=2^{B_m}.
\]
则对充分大的 \(m\)，有严格分解
\[
\boxed{\;
\mathrm{Succ}_{m,a}^{\mathrm{esc}}(B_m)
=
\frac{T_m}{M_m}\,\pi_{m,a}(A_m^{\max})
+
\bigl(1-\pi_{m,a}(A_m^{\max})\bigr).\;}
\]
从而
\[
\boxed{\;
\mathrm{Succ}_{m,a}^{\mathrm{esc}}(B_m)
=
\frac{T_m}{M_m}
+
O\!\bigl(e^{-m\Delta(a)+o(m)}\bigr),\;}
\]
其中
\[
\Delta(a):=a(\alpha_*-\alpha_*^{(2)})-(\log\varphi-g_*)>0.
\]
特别地，
\[
\boxed{\;
\limsup_{m\to\infty}\frac1m
\log \mathrm{Succ}_{m,a}^{\mathrm{esc}}(B_m)
\le
-\min\{\alpha_*-\beta,\Delta(a)\}.\;}
\]
并且总有下界
\[
\boxed{\;
\mathrm{Succ}_{m,a}^{\mathrm{esc}}(B_m)
\ge
\exp\!\bigl(-(\alpha_*-\beta)m+o(m)\bigr).\;}
\]
\end{theorem}
\begin{proof}
由定理 \ref{thm:xi-time-part57ea-tailcount-top-platform-rigidity}，对充分大的 \(m\) 只有最大纤维满足 \(d_m(x)>T_m\)，而对 \(x\notin A_m^{\max}\) 都有 \(d_m(x)\le T_m\)。因此
\[
\min\!\left(1,\frac{T_m}{d_m(x)}\right)
=
\begin{cases}
T_m/M_m, & x\in A_m^{\max},\\
1, & x\notin A_m^{\max}.
\end{cases}
\]
代回定义即得
\[
\mathrm{Succ}_{m,a}^{\mathrm{esc}}(B_m)
=
\frac{T_m}{M_m}\,\pi_{m,a}(A_m^{\max})
+
\bigl(1-\pi_{m,a}(A_m^{\max})\bigr).
\]

再由定理 \ref{thm:xi-fixed-freezing-escort-renyi-spectrum-collapse} 第四条，
\[
1-\pi_{m,a}(A_m^{\max})
\le
\exp\!\bigl(-m\Delta(a)+o(m)\bigr).
\]
于是
\[
\mathrm{Succ}_{m,a}^{\mathrm{esc}}(B_m)
=
\frac{T_m}{M_m}
+
\left(1-\frac{T_m}{M_m}\right)
\bigl(1-\pi_{m,a}(A_m^{\max})\bigr),
\]
从而
\[
\mathrm{Succ}_{m,a}^{\mathrm{esc}}(B_m)
=
\frac{T_m}{M_m}
+
O\!\bigl(e^{-m\Delta(a)+o(m)}\bigr).
\]
又因为
\[
\frac{T_m}{M_m}
=
\exp\!\bigl(-(\alpha_*-\beta)m+o(m)\bigr),
\]
故上式立即给出
\[
\limsup_{m\to\infty}\frac1m
\log \mathrm{Succ}_{m,a}^{\mathrm{esc}}(B_m)
\le
-\min\{\alpha_*-\beta,\Delta(a)\}.
\]

另一方面，由 \(\pi_{m,a}(A_m^{\max})\to 1\) 可知对充分大的 \(m\) 有 \(\pi_{m,a}(A_m^{\max})\ge \tfrac12\)。因此
\[
\mathrm{Succ}_{m,a}^{\mathrm{esc}}(B_m)
\ge
\frac{T_m}{M_m}\,\pi_{m,a}(A_m^{\max})
\ge
\frac12\,\frac{T_m}{M_m},
\]
即得到所示下界。证毕。
\end{proof}

\begin{theorem}[冻结预算相图中的 oracle--escort 临界线]\label{thm:xi-time-part57ea-oracle-escort-critical-line}
在定理 \ref{thm:xi-time-part57ea-frozen-escort-oracle-success-decomposition} 的设定下，定义
\[
\beta_{\mathrm c}(a)
:=
\alpha_*-\Delta(a)
=
\alpha_*^{(2)}+\frac{\log\varphi-g_*}{a}.
\]
则成立：
\begin{enumerate}
  \item 若
  \[
  \beta_{\mathrm c}(a)<\beta<\alpha_*,
  \]
  则 oracle 截断项严格主导，并且
  \[
  \boxed{\;
  \mathrm{Succ}_{m,a}^{\mathrm{esc}}(B_m)
  =
  \exp\!\bigl(-(\alpha_*-\beta)m+o(m)\bigr).\;}
  \]
  \item 若
  \[
  \alpha_*^{(2)}<\beta\le \beta_{\mathrm c}(a),
  \]
  则有统一上界
  \[
  \boxed{\;
  \mathrm{Succ}_{m,a}^{\mathrm{esc}}(B_m)
  \le
  \exp\!\bigl(-\Delta(a)m+o(m)\bigr).\;}
  \]
  在临界点 \(\beta=\beta_{\mathrm c}(a)\) 处，oracle 项与 escort 泄漏上界恰处于同一指数尺度。
  \item 若冻结相的 min-entropy 率存在，则由定理 \ref{thm:xi-fixed-freezing-escort-renyi-spectrum-collapse} 的第三条，
  \[
  h_\infty(a)=g_*,
  \]
  因而临界线可重写为
  \[
  \boxed{\;
  \beta_{\mathrm c}(a)
  =
  \alpha_*^{(2)}+\frac{\log\varphi-h_\infty(a)}{a}.\;}
  \]
\end{enumerate}
\end{theorem}
\begin{proof}
由定义
\[
\alpha_*-\beta_{\mathrm c}(a)=\Delta(a).
\]
若 \(\beta>\beta_{\mathrm c}(a)\)，则
\[
\alpha_*-\beta<\Delta(a).
\]
于是定理 \ref{thm:xi-time-part57ea-frozen-escort-oracle-success-decomposition} 中
\[
\frac{T_m}{M_m}
=
\exp\!\bigl(-(\alpha_*-\beta)m+o(m)\bigr)
\]
比误差项
\(
O(e^{-m\Delta(a)+o(m)})
\)
衰减更慢，从而成为主导项，故得到第一条。

若 \(\beta\le \beta_{\mathrm c}(a)\)，则
\[
\alpha_*-\beta\ge \Delta(a),
\]
故
\[
\frac{T_m}{M_m}
=
O\!\bigl(e^{-m\Delta(a)+o(m)}\bigr).
\]
再与定理 \ref{thm:xi-time-part57ea-frozen-escort-oracle-success-decomposition} 中的误差控制合并，即得第二条上界；当 \(\beta=\beta_{\mathrm c}(a)\) 时，两项指数相同。

第三条只需把 \(g_*\) 用
\[
h_\infty(a)=g_*
\]
替换即可。证毕。
\end{proof}

\paragraph{尾计数指数、冻结阈值与 oracle 前沿的变分闭包}
在定理 \ref{thm:xi-time-part57b-capacity-moment-exact-mellin-stieltjes} 已把幂和矩写成尾计数与容量曲线的 Mellin--Stieltjes 积分、定理 \ref{thm:xi-time-part57ea-tailcount-top-platform-rigidity}--\ref{thm:xi-time-part57ea-oracle-escort-critical-line} 已把冻结预算相图压缩为最大纤维平台律之后，还可进一步把尾计数指数函数、压力谱、均匀采样下的对数多重度大偏差、尺寸偏置倾斜以及连续容量指数统一组织为一套纯变分闭包。该闭包把前文的充分条件提升为必要且充分的边界公式，并把预算前沿、倾斜前沿与冻结前沿一并置于同一尾指数函数 \(\Psi\) 的支撑之上。

\begin{theorem}[尾计数指数的 Laplace--Mellin 变分原理]\label{thm:xi-time-part57eb-tail-laplace-mellin-variational}
设存在上半连续函数
\[
\Psi:[0,\alpha_*]\to[-\infty,\log\varphi]
\]
满足：对任意紧区间 \(I\subset[0,\alpha_*]\)，当 \(m\to\infty\) 时都有一致指数渐近
\[
N_m(e^{\alpha m})
=
\exp\!\bigl(m\Psi(\alpha)+o(m)\bigr)
\qquad
(\alpha\in I),
\]
其中
\[
N_m(T):=\#\{x\in X_m:\ d_m(x)\ge T\}.
\]
则对每个实数 \(a>0\)，压力满足
\[
\boxed{\;
P_a=\sup_{0\le \alpha\le \alpha_*}\bigl(a\alpha+\Psi(\alpha)\bigr).\;}
\]
\end{theorem}
\begin{proof}
由定理 \ref{thm:xi-time-part57b-capacity-moment-exact-mellin-stieltjes}，
\[
S_a(m)=a\int_0^\infty T^{a-1}N_m(T)\,\dd T.
\]
又因 \(d_m(x)\le M_m\) 对所有 \(x\in X_m\) 成立，故
\[
S_a(m)=a\int_0^{M_m}T^{a-1}N_m(T)\,\dd T.
\]
令 \(T=e^{\alpha m}\)，则 \(\dd T=me^{\alpha m}\dd\alpha\)，从而
\[
S_a(m)
=
am\int_{-\infty}^{\frac1m\log M_m}
\exp(am\alpha)\,N_m(e^{\alpha m})\,\dd\alpha.
\]

当 \(\alpha<0\) 时，\(e^{\alpha m}<1\)，而 \(d_m(x)\in\ZZ_{\ge 1}\)，故
\[
N_m(e^{\alpha m})=|X_m|
=
\exp\!\bigl((\log\varphi)m+o(m)\bigr).
\]
于是该区间上的指数至多为 \(a\alpha+\log\varphi\le \log\varphi\)，不会优于 \(\alpha=0\) 处的候选值。另一方面，
\[
\frac1m\log M_m\to \alpha_*.
\]
因此只需在 \([0,\alpha_*]\) 上分析积分主项。

由假设，对每个紧区间 \(I\subset[0,\alpha_*]\) 都有
\[
N_m(e^{\alpha m})=\exp\!\bigl(m\Psi(\alpha)+o(m)\bigr)
\qquad(\alpha\in I)
\]
且 \(o(m)\) 对 \(\alpha\in I\) 一致。将 \([0,\alpha_*]\) 以有限个小区间覆盖，并对每段分别应用 Laplace 原理，再利用 \(\Psi\) 的上半连续性与有限覆盖取极大值，得到
\[
S_a(m)
=
\exp\!\Bigl(
m\sup_{0\le \alpha\le \alpha_*}(a\alpha+\Psi(\alpha))+o(m)
\Bigr).
\]
两边取 \(\frac1m\log\) 并令 \(m\to\infty\)，即得
\[
P_a=\sup_{0\le \alpha\le \alpha_*}(a\alpha+\Psi(\alpha)).
\]
证毕。
\end{proof}

在尾计数指数与压力谱的变分对应已经建立之后，还可把稳定层均匀采样下的对数多重度、其 \(q=1\) 倾斜以及 oracle 容量缺口的指数上界整理为下列三条接口。

\begin{theorem}[均匀类型采样下纤维对数多重度的压力谱大偏差原理]\label{thm:xi-time-part57eb-uniform-logmultiplicity-ldp}
设
\[
S_q(m):=\sum_{x\in X_m}d_m(x)^q,
\qquad
P_q:=\lim_{m\to\infty}\frac1m\log S_q(m).
\]
设存在包含 \(0\) 的开区间 \(I\subset\RR\)，使得对每个 \(q\in I\) 上述极限存在，并记
\[
\psi(q):=P_q-\log\varphi.
\]
若 \(\psi\) 在 \(I\) 上可微，则对稳定类型层上的均匀分布
\[
u_m(x):=\frac{1}{|X_m|}
\qquad (x\in X_m)
\]
所诱导的随机变量
\[
Y_m:=\frac1m\log d_m(X),
\qquad
X\sim u_m,
\]
在速度 \(m\) 下满足以
\[
I_u(y):=\sup_{q\in I}\bigl(qy-\psi(q)\bigr)
\]
为速率函数的压力谱大偏差原理。
\end{theorem}
\begin{proof}
由定义，
\[
\mathbb E_{u_m}\!\bigl[e^{qmY_m}\bigr]
=
\frac1{|X_m|}\sum_{x\in X_m}d_m(x)^q
=
\frac{S_q(m)}{|X_m|}.
\]
因此对每个 \(q\in I\)，都有
\[
\frac1m\log \mathbb E_{u_m}\!\bigl[e^{qmY_m}\bigr]
=
\frac1m\log S_q(m)-\frac1m\log|X_m|.
\]
又因
\[
|X_m|=F_{m+2}
\]
且
\[
\frac1m\log|X_m|\longrightarrow \log\varphi,
\]
故
\[
\lim_{m\to\infty}\frac1m\log \mathbb E_{u_m}\!\bigl[e^{qmY_m}\bigr]
=
P_q-\log\varphi
=
\psi(q).
\]
另一方面，由 \(1\le d_m(x)\le 2^m\) 可知
\[
0\le Y_m\le \log 2,
\]
从而 \((Y_m)\) 自动指数紧。于是按 Gärtner--Ellis 定理，在 \(\psi\) 于 \(I\) 上可微的假设下，即得到所述速度 \(m\) 的大偏差原理，其速率函数正是 Legendre--Fenchel 变换 \(I_u\)。证毕。
\end{proof}

\begin{theorem}[\texorpdfstring{$q=1$}{q=1} 指数倾斜把均匀层速率函数作刚性线性平移]\label{thm:xi-time-part57eb-sizebias-tilt-rate-shift}
沿用定理 \ref{thm:xi-time-part57eb-uniform-logmultiplicity-ldp} 的记号，并进一步定义
\[
\mu_m(x):=\frac{d_m(x)}{2^m}
\qquad (x\in X_m).
\]
则对任意可测子集 \(A\subseteq X_m\)，有精确恒等式
\[
\mu_m(A)
=
\frac{|X_m|}{2^m}\,
\mathbb E_{u_m}\!\bigl[e^{mY_m}\mathbf 1_A\bigr]
=
\frac{\mathbb E_{u_m}\!\bigl[e^{mY_m}\mathbf 1_A\bigr]}
{\mathbb E_{u_m}\!\bigl[e^{mY_m}\bigr]}.
\]
特别地，
\[
\frac1m\log \mathbb E_{u_m}\!\bigl[e^{mY_m}\bigr]
\longrightarrow
\log 2-\log\varphi.
\]
若定理 \ref{thm:xi-time-part57eb-uniform-logmultiplicity-ldp} 的假设成立且 \(1\in I\)，则 \((Y_m)\) 在测度 \(\mu_m\) 下同样满足速度 \(m\) 的大偏差原理，且其速率函数为
\[
I_\mu(y)=I_u(y)-y+\bigl(\log 2-\log\varphi\bigr).
\]
\end{theorem}
\begin{proof}
由纤维分割恒等式
\[
\sum_{x\in X_m}d_m(x)=2^m
\]
可知 \(\mu_m\) 为概率测度。再由 \(e^{mY_m}=d_m(X)\)，对任意 \(A\subseteq X_m\) 都有
\[
\mu_m(A)
=
\frac1{2^m}\sum_{x\in A}d_m(x)
=
\frac{|X_m|}{2^m}\,
\mathbb E_{u_m}\!\bigl[e^{mY_m}\mathbf 1_A\bigr].
\]
另一方面，
\[
\mathbb E_{u_m}\!\bigl[e^{mY_m}\bigr]
=
\frac1{|X_m|}\sum_{x\in X_m}d_m(x)
=
\frac{2^m}{|X_m|},
\]
故第一条恒等式化为第二条归一化写法。

再由
\[
S_1(m)=\sum_{x\in X_m}d_m(x)=2^m
\]
可得
\[
P_1=\log 2,
\qquad
\psi(1)=P_1-\log\varphi=\log 2-\log\varphi.
\]
于是
\[
\frac1m\log \mathbb E_{u_m}\!\bigl[e^{mY_m}\bigr]
=
\log 2-\frac1m\log|X_m|
\longrightarrow
\log 2-\log\varphi.
\]
最后，由上式可见 \(\mu_m\) 正是 \(u_m\) 以势函数 \(y\mapsto y\) 作出的标准指数倾斜；故按大偏差论中的倾斜原理，其速率函数由 \(I_u\) 刚性平移为
\[
I_\mu(y)=I_u(y)-y+\psi(1)
=
I_u(y)-y+\bigl(\log 2-\log\varphi\bigr).
\]
证毕。
\end{proof}

\begin{theorem}[oracle 容量缺口的 Rényi--压力谱指数上界]\label{thm:xi-time-part57eb-oracle-gap-renyi-pressure-upper-bound}
\leanverified{paper\_xi\_time\_part57eb\_oracle\_gap\_renyi\_pressure\_upper\_bound}
对整数预算 \(B\ge 0\)，定义
\[
C_m(B):=\sum_{x\in X_m}\min\bigl(d_m(x),2^B\bigr),
\qquad
\mathrm{Succ}_m(B):=\frac{C_m(B)}{2^m},
\]
以及容量缺口
\[
\mathcal G_m(B):=1-\mathrm{Succ}_m(B)
=
\frac1{2^m}\sum_{x\in X_m}\bigl(d_m(x)-2^B\bigr)_+.
\]
则对任意 \(q>1\)，都有
\[
\mathcal G_m(B)\le \frac{S_q(m)}{2^m\,2^{(q-1)B}}.
\]
若进一步取
\[
B_m:=\lfloor \beta m\rfloor
\qquad(\beta\ge 0),
\]
并假设 \(P_q\) 存在，则
\[
\limsup_{m\to\infty}\frac1m\log \mathcal G_m(B_m)
\le
P_q-\log 2-(q-1)\beta\log 2.
\]
因此定义
\[
E(\beta):=
\sup_{q>1}\Bigl((q-1)\beta\log 2-\bigl(P_q-\log 2\bigr)\Bigr),
\]
便有统一上界
\[
\limsup_{m\to\infty}\frac1m\log \mathcal G_m(B_m)\le -E(\beta).
\]
\end{theorem}
\begin{proof}
由定义，
\[
\mathcal G_m(B)
=
\frac1{2^m}\sum_{d_m(x)>2^B}\bigl(d_m(x)-2^B\bigr)
\le
\frac1{2^m}\sum_{d_m(x)>2^B}d_m(x).
\]
对任意 \(q>1\)，若 \(d_m(x)>2^B\)，则
\[
d_m(x)\le \frac{d_m(x)^q}{2^{(q-1)B}}.
\]
代回上式即得
\[
\mathcal G_m(B)
\le
\frac1{2^m}\sum_{x\in X_m}\frac{d_m(x)^q}{2^{(q-1)B}}
=
\frac{S_q(m)}{2^m\,2^{(q-1)B}}.
\]
现在取 \(B=B_m=\lfloor \beta m\rfloor\)。由于
\[
2^{(q-1)B_m}
=
\exp\!\bigl((q-1)\beta(\log 2)m+O(1)\bigr),
\]
对上式取对数、除以 \(m\) 并令 \(m\to\infty\)，便得到
\[
\limsup_{m\to\infty}\frac1m\log \mathcal G_m(B_m)
\le
P_q-\log 2-(q-1)\beta\log 2.
\]
最后对全部 \(q>1\) 取上确界，即得所示 \(E(\beta)\) 上界。证毕。
\end{proof}


\begin{theorem}[连续容量曲线的变分公式与成功率指数]\label{thm:xi-time-part57eb-continuous-capacity-success-variational}
定义连续容量曲线
\[
\mathcal C_m^{\mathrm{cont}}(T):=\sum_{x\in X_m}\min(d_m(x),T)
\qquad(T\ge 0),
\]
以及预算指数
\[
Q(\beta):=\limsup_{m\to\infty}\frac1m\log \mathcal C_m^{\mathrm{cont}}(e^{\beta m})
\qquad(\beta\ge 0).
\]
在定理 \ref{thm:xi-time-part57eb-tail-laplace-mellin-variational} 的同一假设下，对每个 \(\beta\ge 0\) 都有
\[
\boxed{\;
Q(\beta)=\sup_{0\le \alpha\le \min(\beta,\alpha_*)}\bigl(\alpha+\Psi(\alpha)\bigr).\;}
\]
进一步，若取离散预算
\[
B_m:=\left\lfloor \frac{\beta m}{\log 2}\right\rfloor,
\]
并记
\[
\mathrm{Succ}_m(B):=\frac{C_m(B)}{2^m},
\qquad
C_m(B):=\sum_{x\in X_m}\min(d_m(x),2^B),
\]
则有
\[
\boxed{\;
\limsup_{m\to\infty}\frac1m\log \mathrm{Succ}_m(B_m)
=
Q(\beta)-\log 2.\;}
\]
\end{theorem}
\begin{proof}
由定理 \ref{thm:xi-capacity-tail-histogram-moment-complete-statistic}，
\[
\mathcal C_m^{\mathrm{cont}}(T)=\int_0^T N_m(t)\,\dd t.
\]
取 \(T=e^{\beta m}\)，再作变量代换 \(t=e^{\alpha m}\)，得到
\[
\mathcal C_m^{\mathrm{cont}}(e^{\beta m})
=
m\int_{-\infty}^{\beta}\exp(\alpha m)\,N_m(e^{\alpha m})\,\dd\alpha.
\]

当 \(\alpha<0\) 时，仍有 \(N_m(e^{\alpha m})=|X_m|=\exp((\log\varphi)m+o(m))\)，故对应指数不超过 \(\log\varphi+\alpha\le \log\varphi\)，不会优于 \(\alpha=0\) 的候选值。又由于 \(N_m(t)=0\) 当 \(t>M_m\)，并且 \(\frac1m\log M_m\to\alpha_*\)，故只需在 \([0,\min(\beta,\alpha_*)]\) 上取主项。由定理 \ref{thm:xi-time-part57eb-tail-laplace-mellin-variational} 的同一一致指数假设，对该积分再施行 Laplace 原理即得
\[
Q(\beta)=\sup_{0\le \alpha\le \min(\beta,\alpha_*)}\bigl(\alpha+\Psi(\alpha)\bigr).
\]

另一方面，由命题 \ref{prop:fold-capacity-tail-equivalence}，
\[
C_m(B)=\mathcal C_m^{\mathrm{cont}}(2^B)
\qquad(B\in\ZZ_{\ge 0}).
\]
对 \(B=B_m\) 有
\[
2^{B_m}=e^{\beta m+O(1)},
\]
故
\[
\limsup_{m\to\infty}\frac1m\log C_m(B_m)=Q(\beta).
\]
再由 \(\mathrm{Succ}_m(B_m)=C_m(B_m)/2^m\)，得到
\[
\limsup_{m\to\infty}\frac1m\log \mathrm{Succ}_m(B_m)=Q(\beta)-\log 2.
\]
证毕。
\end{proof}

\begin{theorem}[冻结阈值的精确变分刻画]\label{thm:xi-time-part57eb-freezing-threshold-exact-variational}
设在区间 \((\alpha_*^{(2)},\alpha_*)\) 上存在尾计数平台
\[
\Psi(\alpha)=g_*
\qquad
(\alpha_*^{(2)}<\alpha<\alpha_*).
\]
定义
\[
a_{\mathrm{fr}}
:=
\sup_{0\le \alpha<\alpha_*}
\frac{\Psi(\alpha)-g_*}{\alpha_*-\alpha}.
\]
则对任意实数 \(a>0\)，有等价判别
\[
\boxed{\;
P_a=a\alpha_*+g_*
\iff
a\ge a_{\mathrm{fr}}.\;}
\]
并且
\[
\boxed{\;
a_{\mathrm{fr}}
\le
\frac{\log\varphi-g_*}{\alpha_*-\alpha_*^{(2)}}.\;}
\]
因此前文的阈值
\[
a_{\mathrm c}:=\frac{\log\varphi-g_*}{\alpha_*-\alpha_*^{(2)}}
\]
不仅是充分条件，而且是精确冻结阈值的统一上界。
\end{theorem}
\begin{proof}
由定理 \ref{thm:xi-time-part57eb-tail-laplace-mellin-variational}，
\[
P_a=\sup_{0\le \alpha\le \alpha_*}\bigl(a\alpha+\Psi(\alpha)\bigr).
\]
由平台假设，对每个 \(\alpha\in(\alpha_*^{(2)},\alpha_*)\) 都有
\[
a\alpha+\Psi(\alpha)=a\alpha+g_*.
\]
令 \(\alpha\uparrow\alpha_*\)，得到
\[
\sup_{\alpha_*^{(2)}<\alpha<\alpha_*}\bigl(a\alpha+\Psi(\alpha)\bigr)
=
a\alpha_*+g_*.
\]
因此
\[
P_a=a\alpha_*+g_*
\]
成立，当且仅当对所有 \(0\le \alpha<\alpha_*\) 都有
\[
a\alpha+\Psi(\alpha)\le a\alpha_*+g_*,
\]
等价于
\[
a\ge \frac{\Psi(\alpha)-g_*}{\alpha_*-\alpha}.
\]
对全部 \(0\le \alpha<\alpha_*\) 取上确界，即得
\[
P_a=a\alpha_*+g_*
\iff
a\ge a_{\mathrm{fr}}.
\]

再由平台假设，当 \(\alpha\in(\alpha_*^{(2)},\alpha_*)\) 时，
\[
\Psi(\alpha)=g_*,
\]
故该区间对上确界没有正贡献。于是
\[
a_{\mathrm{fr}}
=
\sup_{0\le \alpha\le \alpha_*^{(2)}}
\frac{\Psi(\alpha)-g_*}{\alpha_*-\alpha}.
\]
另一方面，对所有 \(\alpha\in[0,\alpha_*]\) 都有 \(\Psi(\alpha)\le \log\varphi\)，从而
\[
a_{\mathrm{fr}}
\le
\sup_{0\le \alpha\le \alpha_*^{(2)}}
\frac{\log\varphi-g_*}{\alpha_*-\alpha}
=
\frac{\log\varphi-g_*}{\alpha_*-\alpha_*^{(2)}}.
\]
证毕。
\end{proof}

\begin{theorem}[oracle--escort 双前沿的精确相图]\label{thm:xi-time-part57eb-oracle-escort-double-frontier}
保持定理 \ref{thm:xi-time-part57eb-freezing-threshold-exact-variational} 的尾计数平台假设。定义连续容量指数
\[
Q(\beta)=\sup_{0\le \alpha\le \min(\beta,\alpha_*)}\bigl(\alpha+\Psi(\alpha)\bigr),
\]
以及预算饱和阈值
\[
\beta_{\mathrm{sat}}
:=
\sup_{0\le \alpha\le \alpha_*^{(2)}}\bigl(\alpha+\Psi(\alpha)-g_*\bigr).
\]
则对每个 \(\beta\in(\alpha_*^{(2)},\alpha_*)\)，成立如下分类：
\begin{enumerate}
  \item 若 \(\beta\ge \beta_{\mathrm{sat}}\)，则
  \[
  \boxed{\;
  Q(\beta)=\beta+g_*.\;}
  \]
  因而对
  \[
  B_m=\left\lfloor \frac{\beta m}{\log 2}\right\rfloor
  \]
  有
  \[
  \limsup_{m\to\infty}\frac1m\log \mathrm{Succ}_m(B_m)
  =
  \beta+g_*-\log 2.
  \]
  \item 若 \(\beta<\beta_{\mathrm{sat}}\)，则
  \[
  \boxed{\;
  Q(\beta)=
  \sup_{0\le \alpha\le \alpha_*^{(2)}}\bigl(\alpha+\Psi(\alpha)\bigr).\;}
  \]
  换言之，容量指数仍由次大尺度区间决定，尚未进入最大纤维主导相。
  \item 总有显式上界
  \[
  \boxed{\;
  \beta_{\mathrm{sat}}
  \le
  \alpha_*^{(2)}+\log\varphi-g_*.\;}
  \]
  因此只要
  \[
  \beta>\alpha_*^{(2)}+\log\varphi-g_*,
  \]
  就必进入线性饱和相 \(Q(\beta)=\beta+g_*\)。
\end{enumerate}
\end{theorem}
\begin{proof}
对 \(\beta\in(\alpha_*^{(2)},\alpha_*)\)，由平台性质，
\[
\Psi(\alpha)=g_*
\qquad
(\alpha_*^{(2)}<\alpha\le \beta).
\]
故在该区间上
\[
\alpha+\Psi(\alpha)=\alpha+g_*
\]
严格随 \(\alpha\) 递增，其最大值为 \(\beta+g_*\)。于是
\[
Q(\beta)
=
\max\Bigl\{
\beta+g_*,
\sup_{0\le \alpha\le \alpha_*^{(2)}}(\alpha+\Psi(\alpha))
\Bigr\}.
\]

若 \(\beta\ge \beta_{\mathrm{sat}}\)，则按定义有
\[
\beta+g_*
\ge
\sup_{0\le \alpha\le \alpha_*^{(2)}}(\alpha+\Psi(\alpha)),
\]
从而
\[
Q(\beta)=\beta+g_*.
\]
再由定理 \ref{thm:xi-time-part57eb-continuous-capacity-success-variational}，得到
\[
\limsup_{m\to\infty}\frac1m\log \mathrm{Succ}_m(B_m)
=
Q(\beta)-\log 2
=
\beta+g_*-\log 2.
\]

若 \(\beta<\beta_{\mathrm{sat}}\)，则
\[
\beta+g_*
<
\sup_{0\le \alpha\le \alpha_*^{(2)}}(\alpha+\Psi(\alpha)),
\]
故
\[
Q(\beta)=
\sup_{0\le \alpha\le \alpha_*^{(2)}}(\alpha+\Psi(\alpha)).
\]
这就给出第二种相。

最后，由 \(\Psi(\alpha)\le \log\varphi\)，有
\[
\beta_{\mathrm{sat}}
=
\sup_{0\le \alpha\le \alpha_*^{(2)}}(\alpha+\Psi(\alpha)-g_*)
\le
\sup_{0\le \alpha\le \alpha_*^{(2)}}(\alpha+\log\varphi-g_*)
=
\alpha_*^{(2)}+\log\varphi-g_*.
\]
故当
\(
\beta>\alpha_*^{(2)}+\log\varphi-g_*
\)
时，必有 \(\beta>\beta_{\mathrm{sat}}\)，从而落入线性饱和相。证毕。
\end{proof}

再进一步，均匀层大偏差与尺寸偏置倾斜应共同决定 oracle 容量缺口的宏观相图。

\begin{conjecture}[oracle 容量缺口指数的双 Legendre 最小包络]\label{conj:xi-time-part57eb-oracle-gap-min-envelope-phase-diagram}
沿用定理 \ref{thm:xi-time-part57eb-uniform-logmultiplicity-ldp} 与定理 \ref{thm:xi-time-part57eb-sizebias-tilt-rate-shift} 的记号。取
\[
B_m:=\lfloor \beta m\rfloor,
\qquad
\tau:=\beta\log 2,
\]
并定义上尾速率函数
\[
J_u(\tau):=\inf_{y>\tau}I_u(y),
\qquad
J_\mu(\tau):=\inf_{y>\tau}I_\mu(y).
\]
则容量缺口
\[
\mathcal G_m(B_m)=1-\mathrm{Succ}_m(B_m)
\]
的指数阶应满足
\[
\lim_{m\to\infty}-\frac1m\log \mathcal G_m(B_m)
=
\min\Bigl\{
J_\mu(\tau),\,
J_u(\tau)-\bigl(\log\varphi-\log 2+\tau\bigr)
\Bigr\}.
\]
特别地，两项相等的曲线应给出 oracle 容量缺口宏观相图的一阶相变边界。

事实上，若记
\[
\tau_m:=\frac{B_m\log 2}{m},
\]
则有精确恒等分解
\[
\mathcal G_m(B_m)
=
\mu_m(Y_m>\tau_m)
-\exp\!\Bigl(m\Bigl(\tau_m+\frac1m\log|X_m|-\log 2\Bigr)\Bigr)\,
u_m(Y_m>\tau_m).
\]
该猜想断言：在大偏差尺度上，上式两项不存在额外的指数级相消，因此缺口指数恰由两条尾率中衰减较慢者决定。
\end{conjecture}




\begin{theorem}[分枝三次场、零温四次场与 RH 临界五次场的算术独立]\label{thm:xi-branch-rh-zero-temp-arithmetic-independence}
\leanverified{paper\_xi\_branch\_rh\_zero\_temp\_arithmetic\_independence}
设
\[
\Delta(u):=4u^3+25u^2+88u+8,
\]
\[
R_{\mathrm{gr}}(y):=y^4-6y^3+9y^2-y-1,
\]
\[
Q_{\mathrm{RH}}(u):=u^5+4u^4+3u^3-96u^2-42u-14.
\]
分别记其分裂域为 \(K_\Delta\)、\(K_{\mathrm{gr}}\)、\(K_{\mathrm{RH}}\)。则
\[
\boxed{\;
\Gal(K_\Delta/\QQ)\cong S_3,\qquad
\Gal(K_{\mathrm{gr}}/\QQ)\cong S_4,\qquad
\Gal(K_{\mathrm{RH}}/\QQ)\cong S_5.\;}
\]
并且三者在 \(\QQ\) 上算术独立，满足
\[
\boxed{\;
\Gal(K_\Delta K_{\mathrm{gr}}K_{\mathrm{RH}}/\QQ)
\cong
S_3\times S_4\times S_5.\;}
\]
等价地，分枝几何、零温地基与 RH 临界阈值并非同一伽罗瓦源头的不同投影，而对应三个彼此正交的对称场。
\end{theorem}
\begin{proof}
先看三次多项式 \(\Delta(u)\)。由有理根判别法，\(\Delta\) 在 \(\QQ\) 上不可约。其判别式为
\[
\Disc(\Delta)=-2^5\cdot 5^3\cdot 11^3,
\]
非平方，故其分裂域伽罗瓦群只能为 \(S_3\) 而非 \(A_3\)。于是
\[
\Gal(K_\Delta/\QQ)\cong S_3,
\qquad
K_\Delta^{A_3}=\QQ(\sqrt{-110}).
\]
\par
再看四次多项式 \(R_{\mathrm{gr}}(y)\)。模 \(3\) 约化后 \(R_{\mathrm{gr}}\) 仍不可约，因此分裂域伽罗瓦群在 \(S_4\) 中传递并含一个 \(4\)-循环。其判别式为
\[
\Disc(R_{\mathrm{gr}})=2777,
\]
非平方，故群不包含于 \(A_4\)。另一方面，模 \(11\) 的分解型为 \([1,3]\)，故 Frobenius 元给出一个 \(3\)-循环。传递子群 \(G\le S_4\) 若同时含 \(4\)-循环与 \(3\)-循环，则只能是 \(S_4\)。故
\[
\Gal(K_{\mathrm{gr}}/\QQ)\cong S_4,
\qquad
K_{\mathrm{gr}}^{A_4}=\QQ(\sqrt{2777}).
\]
\par
对五次多项式 \(Q_{\mathrm{RH}}(u)\)，模 \(5\) 约化后仍不可约，因此分裂域伽罗瓦群在 \(S_5\) 中传递并含一个 \(5\)-循环。其判别式为
\[
\Disc(Q_{\mathrm{RH}})
=
2^6\cdot 3^2\cdot 7\cdot 743\cdot 1693^2,
\]
非平方，故群不包含于 \(A_5\)。又模 \(17\) 的分解型为 \([1,1,3]\)，故 Frobenius 元给出一个 \(3\)-循环。按五次传递子群分类，含 \(5\)-循环与 \(3\)-循环的传递子群只能为 \(A_5\) 或 \(S_5\)；判别式非平方排除 \(A_5\)，从而
\[
\Gal(K_{\mathrm{RH}}/\QQ)\cong S_5,
\qquad
K_{\mathrm{RH}}^{A_5}=\QQ(\sqrt{5201}).
\]
\par
下面证明独立性。先看 \(K_\Delta\) 与 \(K_{\mathrm{gr}}\)。二者交域是 \(\QQ\) 上的 Galois 扩张，其 Galois 群必须同时是 \(S_3\) 与 \(S_4\) 的商群，因而只能为 \(1\)、\(C_2\) 或 \(S_3\)。若为 \(C_2\)，则二者的唯一二次子域必须相同，但
\[
\QQ(\sqrt{-110})\neq \QQ(\sqrt{2777}),
\]
矛盾。若为 \(S_3\)，则 \(K_\Delta\subseteq K_{\mathrm{gr}}\)，从而 \(K_\Delta\) 必等于 \(K_{\mathrm{gr}}^{V_4}\) 这一 \(S_3\)-商域；然而其唯一二次子域应为
\[
K_{\mathrm{gr}}^{A_4}=\QQ(\sqrt{2777}),
\]
仍与 \(\QQ(\sqrt{-110})\) 矛盾。故
\[
K_\Delta\cap K_{\mathrm{gr}}=\QQ
\]
并得到
\[
\Gal(K_\Delta K_{\mathrm{gr}}/\QQ)\cong S_3\times S_4.
\]
\par
再看 \(K_{\mathrm{RH}}\) 与前两者。首先，\(S_5\) 的非平凡真正规商只有 \(C_2\)，故
\[
K_\Delta\cap K_{\mathrm{RH}}=\QQ,
\qquad
K_{\mathrm{gr}}\cap K_{\mathrm{RH}}=\QQ
\]
只需排除二次子域相同的可能，而
\[
\QQ(\sqrt{5201})
\neq
\QQ(\sqrt{-110}),
\qquad
\QQ(\sqrt{5201})
\neq
\QQ(\sqrt{2777}).
\]
接着考虑与合成域 \(K_\Delta K_{\mathrm{gr}}\) 的交。由于
\[
\bigl(S_3\times S_4\bigr)^{\mathrm{ab}}\cong C_2\times C_2,
\]
其全部二次子域恰由三条非零特征给出，即
\[
\QQ(\sqrt{-110}),\qquad
\QQ(\sqrt{2777}),\qquad
\QQ(\sqrt{-110\cdot 2777}).
\]
这些二次域都不同于 \(\QQ(\sqrt{5201})\)。而 \(S_5\) 侧任何非平凡 Galois 交域只能是其唯一二次子域 \(\QQ(\sqrt{5201})\)，故
\[
K_{\mathrm{RH}}\cap K_\Delta K_{\mathrm{gr}}=\QQ.
\]
于是
\[
\Gal(K_\Delta K_{\mathrm{gr}}K_{\mathrm{RH}}/\QQ)
\cong
\Gal(K_\Delta/\QQ)\times\Gal(K_{\mathrm{gr}}/\QQ)\times\Gal(K_{\mathrm{RH}}/\QQ)
\cong
S_3\times S_4\times S_5.
\]
可复现核验脚本：\texttt{scripts/exp\_xi\_branch\_rh\_zero\_temp\_galois\_chebotarev\_audit.py}。证毕。
\end{proof}

\begin{corollary}[三重 Chebotarev 独立律与显式联合密度]\label{cor:xi-branch-rh-zero-temp-chebotarev-product-densities}
\leanverified{paper\_xi\_branch\_rh\_zero\_temp\_chebotarev\_product\_densities}
沿用定理 \ref{thm:xi-branch-rh-zero-temp-arithmetic-independence} 的记号。对一切不分歧于
\[
\Disc(\Delta)\Disc(R_{\mathrm{gr}})\Disc(Q_{\mathrm{RH}})
\]
的素数 \(p\)，三条多项式
\[
\Delta(u),\qquad R_{\mathrm{gr}}(y),\qquad Q_{\mathrm{RH}}(u)
\]
在 \(\FF_p\) 上的分解型彼此独立；更精确地，对任意共轭类
\[
C_3\subset S_3,\qquad C_4\subset S_4,\qquad C_5\subset S_5,
\]
满足
\[
\delta\Bigl(\mathrm{Frob}_p\in C_3\times C_4\times C_5\Bigr)
=
\frac{|C_3|}{|S_3|}\cdot \frac{|C_4|}{|S_4|}\cdot \frac{|C_5|}{|S_5|}.
\]
特别地：
\begin{enumerate}
  \item 三者同时不可约的素数密度为
  \[
  \boxed{\;
  \frac{2}{6}\cdot\frac{6}{24}\cdot\frac{24}{120}
  =
  \frac{1}{60}.\;}
  \]
  \item 三者同时完全分裂的素数密度为
  \[
  \boxed{\;
  \frac{1}{6}\cdot\frac{1}{24}\cdot\frac{1}{120}
  =
  \frac{1}{17280}.\;}
  \]
  \item \(\Delta\) 不可约、\(R_{\mathrm{gr}}\) 具有 \((3+1)\) 分解、且 \(Q_{\mathrm{RH}}\) 不可约的联合密度为
  \[
  \boxed{\;
  \frac{2}{6}\cdot\frac{8}{24}\cdot\frac{24}{120}
  =
  \frac{1}{45},\;}
  \]
  因为 \(S_4\) 中 \((3+1)\) 型共轭类大小为 \(8\)。
\end{enumerate}
\end{corollary}
\begin{proof}
定理 \ref{thm:xi-branch-rh-zero-temp-arithmetic-independence} 已给出
\[
\Gal(K_\Delta K_{\mathrm{gr}}K_{\mathrm{RH}}/\QQ)\cong S_3\times S_4\times S_5.
\]
Chebotarev 密度定理说明：对不分歧素数，Frobenius 共轭类在该直积群中按计数测度等分布。故任意联合分解型事件的密度都等于三个因子中对应共轭类比例的乘积。
\par
另一方面，多项式在 \(\FF_p\) 上的分解型与 Frobenius 元的循环型一致：三次不可约对应 \(S_3\) 中的 \(3\)-循环，四次不可约对应 \(S_4\) 中的 \(4\)-循环，五次不可约对应 \(S_5\) 中的 \(5\)-循环，而四次的 \((3+1)\) 分解对应 \(S_4\) 中大小为 \(8\) 的 \((3+1)\) 共轭类。将这些类大小代入上式，即得到三条显式密度。证毕。
\end{proof}

\paragraph{原子长度二残差、根单位滤波与 Lee--Yang 符号坍缩}
下列条目把定理 \ref{thm:xi-real-input-40-atomic-witt-even-odd-surgery}、
\ref{thm:xi-atomic-witt-mod-class-localization}、定理
\ref{thm:xi-dirichlet-mode-forces-congruence-residual}、推论
\ref{cor:fold-gauge-anomaly-P10-leyang-chebotarev-product-law} 与结论
\ref{con:xi-terminal-zm-leyang-cubic-modp-splitting-density} 重新压缩为同一条
“原子剥离--同余注入--core 门槛--算术筛选”链；并且该链与
\S\ref{subsubsec:zeta-online-kernel-length-modq-chebotarev} 的 root-of-unity 长度滤波接口完全对齐。

\begin{theorem}[原子 Witt 因子的周期--primitive 精确剥离]\label{thm:xi-atomic-witt-cycle-primitive-exact-peeling}
\leanverified{paper\_xi\_atomic\_witt\_cycle\_primitive\_exact\_peeling}
沿用附录 \ref{app:real-input-40-zeta-u} 的记号。设
\[
\zeta(z,u)=\Delta(z,u)^{-1},
\qquad
\Delta(z,u)=(1-uz^2)\,F(z,u),
\qquad
\zeta_{\mathrm{core}}(z,u):=F(z,u)^{-1}.
\]
定义周期对数系数
\[
\log\zeta(z,u)=\sum_{n\ge 1}\frac{a_n(u)}{n}z^n,
\qquad
\log\zeta_{\mathrm{core}}(z,u)=\sum_{n\ge 1}\frac{a_n^{\mathrm{core}}(u)}{n}z^n,
\]
并定义 primitive 生成函数
\[
P(z,u)=\sum_{n\ge 1}p_n(u)z^n,
\qquad
P_{\mathrm{core}}(z,u)=\sum_{n\ge 1}p_n^{\mathrm{core}}(u)z^n.
\]
则对一切 \(n\ge 1\) 都有严格恒等式
\[
\boxed{\;
a_n(u)=a_n^{\mathrm{core}}(u)+2u^{n/2}\mathbf 1_{\{2\mid n\}}\;}
\]
以及
\[
\boxed{\;
p_n(u)=p_n^{\mathrm{core}}(u)+u\,\mathbf 1_{\{n=2\}}.\;}
\]
因此，原子因子在 primitive 层只占据长度 \(2\) 的单一坐标，而在周期/ghost 层则恰好沿全部偶长度产生几何级数修正。
\end{theorem}
\begin{proof}
由
\[
\zeta(z,u)=(1-uz^2)^{-1}\zeta_{\mathrm{core}}(z,u)
\]
取对数得
\[
\log\zeta(z,u)
=
-\log(1-uz^2)+\log\zeta_{\mathrm{core}}(z,u)
=
\sum_{k\ge 1}\frac{u^k}{k}z^{2k}
+\log\zeta_{\mathrm{core}}(z,u).
\]
比较 \(z^n/n\) 的系数，便得到
\[
a_n(u)-a_n^{\mathrm{core}}(u)=
\begin{cases}
2u^{n/2},&2\mid n,\\
0,&2\nmid n.
\end{cases}
\]
另一方面，命题 \ref{prop:atomic-2-orbit-split-logM} 已给出精确 primitive 剥离
\[
P(z,u)=uz^2+P_{\mathrm{core}}(z,u),
\]
逐项比较系数即得
\[
p_n(u)=p_n^{\mathrm{core}}(u)+u\,\mathbf 1_{\{n=2\}}.
\]
证毕。
\end{proof}

\begin{corollary}[无权核相对 core 恰多出唯一一个 primitive 二步类]\label{cor:xi-atomic-witt-unweighted-unique-two-step-class}
沿用定理 \ref{thm:xi-atomic-witt-cycle-primitive-exact-peeling} 的记号，并取 \(u=1\)。则
\[
p_n(1)=p_n^{\mathrm{core}}(1)+\mathbf 1_{\{n=2\}},
\qquad
a_n(1)=a_n^{\mathrm{core}}(1)+2\,\mathbf 1_{\{2\mid n\}}.
\]
因此完整系统相对于 core 恰多出唯一一个 primitive 长度 \(2\) 轨道类；其在 periodic 层上的全部额外贡献恰为偶长度上的二步原子反复：
\[
a_{2m}(1)-a_{2m}^{\mathrm{core}}(1)=2,
\qquad
a_{2m+1}(1)-a_{2m+1}^{\mathrm{core}}(1)=0
\qquad (m\ge 1).
\]
\end{corollary}
\begin{proof}
由定理 \ref{thm:xi-atomic-witt-cycle-primitive-exact-peeling} 直接在 \(u=1\) 代入即可。证毕。
\leanverified{paper\_xi\_atomic\_witt\_unweighted\_unique\_two\_step\_class}
\end{proof}

把同一原子剥离沿能量权变量 \(u\) 施以根单位 Fourier 反演，还可得到模 \(m\) 能量 residue 侧的完全局域化修正。

\begin{theorem}[原子 Witt 因子的能量 residue 单类注入与 \texorpdfstring{$L^2$}{L2} 原子律]\label{thm:xi-atomic-witt-energy-residue-l2-law}
固定整数 \(m\ge 2\)，记 \(\omega_m:=e^{2\pi i/m}\)。对 \(r\in \ZZ/m\ZZ\) 与 \(n\ge 1\)，定义 periodic 与 primitive 的能量 residue 计数
\[
A_{m,r}(n):=\#\{\gamma:\ |\gamma|=n,\ E(\gamma)\equiv r\pmod m\},
\]
\[
P_{m,r}(n):=\#\{\mathfrak p:\ |\mathfrak p|=n,\ E(\mathfrak p)\equiv r\pmod m\}.
\]
并以离散 Fourier 反演定义 core residue 坐标
\[
A_{m,r}^{\mathrm{core}}(n)
:=
\frac1m\sum_{j=0}^{m-1}a_n^{\mathrm{core}}(\omega_m^j)\,\omega_m^{-jr},
\]
\[
P_{m,r}^{\mathrm{core}}(n)
:=
\frac1m\sum_{j=0}^{m-1}p_n^{\mathrm{core}}(\omega_m^j)\,\omega_m^{-jr}.
\]
则对一切 \(n\ge 1\) 与 \(r\in \ZZ/m\ZZ\) 都有严格恒等式
\[
\boxed{\;
A_{m,r}(n)
=
A_{m,r}^{\mathrm{core}}(n)
+2\,\mathbf 1_{\{2\mid n\}}\mathbf 1_{\{r\equiv n/2\;(\mathrm{mod}\,m)\}}.\;}
\]
\[
\boxed{\;
P_{m,r}(n)
=
P_{m,r}^{\mathrm{core}}(n)
+\mathbf 1_{\{n=2\}}\mathbf 1_{\{r\equiv 1\;(\mathrm{mod}\,m)\}}.\;}
\]
并且 twisted 系数差满足精确平方平均律
\[
\boxed{\;
\frac1m\sum_{j=0}^{m-1}
\bigl|a_n(\omega_m^j)-a_n^{\mathrm{core}}(\omega_m^j)\bigr|^2
=
4\,\mathbf 1_{\{2\mid n\}}.\;}
\]
\[
\boxed{\;
\frac1m\sum_{j=0}^{m-1}
\bigl|p_n(\omega_m^j)-p_n^{\mathrm{core}}(\omega_m^j)\bigr|^2
=
\mathbf 1_{\{n=2\}}.\;}
\]
因此，长度 \(2\) 的原子因子在能量 residue 侧并不产生扩散型修正，而是对每个偶长度 \(n\) 只向单一 residue 类 \(r\equiv n/2\pmod m\) 注入恰好两个 periodic 原子；primitive 层则仅在 \((n,r)=(2,1)\) 处留下单点质量。
\end{theorem}
\begin{proof}
由定理 \ref{thm:xi-atomic-witt-cycle-primitive-exact-peeling}，对每个 \(0\le j\le m-1\) 都有
\[
a_n(\omega_m^j)-a_n^{\mathrm{core}}(\omega_m^j)
=
2\,\omega_m^{j n/2}\mathbf 1_{\{2\mid n\}},
\]
\[
p_n(\omega_m^j)-p_n^{\mathrm{core}}(\omega_m^j)
=
\omega_m^j\,\mathbf 1_{\{n=2\}}.
\]
因此
\[
\begin{aligned}
A_{m,r}(n)-A_{m,r}^{\mathrm{core}}(n)
&=
\frac1m\sum_{j=0}^{m-1}
\bigl(a_n(\omega_m^j)-a_n^{\mathrm{core}}(\omega_m^j)\bigr)\omega_m^{-jr} \\
&=
\frac{2\,\mathbf 1_{\{2\mid n\}}}{m}
\sum_{j=0}^{m-1}\omega_m^{j(n/2-r)}.
\end{aligned}
\]
单位根正交性给出
\[
\sum_{j=0}^{m-1}\omega_m^{j(n/2-r)}
=
\begin{cases}
m,& r\equiv n/2\pmod m,\\
0,& r\not\equiv n/2\pmod m,
\end{cases}
\]
从而得到 periodic residue 修正公式。

primitive 层同理：
\[
\begin{aligned}
P_{m,r}(n)-P_{m,r}^{\mathrm{core}}(n)
&=
\frac1m\sum_{j=0}^{m-1}
\bigl(p_n(\omega_m^j)-p_n^{\mathrm{core}}(\omega_m^j)\bigr)\omega_m^{-jr} \\
&=
\frac{\mathbf 1_{\{n=2\}}}{m}
\sum_{j=0}^{m-1}\omega_m^{j(1-r)},
\end{aligned}
\]
故仅当 \(n=2\) 且 \(r\equiv 1\pmod m\) 时取值为 \(1\)。

最后，由 \(|\omega_m^{j n/2}|=1\) 与 \(|\omega_m^j|=1\)，直接得到
\[
\frac1m\sum_{j=0}^{m-1}
\bigl|a_n(\omega_m^j)-a_n^{\mathrm{core}}(\omega_m^j)\bigr|^2
=
\frac1m\sum_{j=0}^{m-1}4\,\mathbf 1_{\{2\mid n\}}
=
4\,\mathbf 1_{\{2\mid n\}},
\]
\[
\frac1m\sum_{j=0}^{m-1}
\bigl|p_n(\omega_m^j)-p_n^{\mathrm{core}}(\omega_m^j)\bigr|^2
=
\frac1m\sum_{j=0}^{m-1}\mathbf 1_{\{n=2\}}
=
\mathbf 1_{\{n=2\}}.
\]
证毕。
\end{proof}

\begin{theorem}[长度同余类滤波中的单类注入律]\label{thm:xi-root-unity-filter-single-class-injection}
固定模数 \(q\ge 2\)，记 \(\omega_q:=e^{2\pi i/q}\)，并取剩余类 \(a\in\ZZ/q\ZZ\)。
定义 primitive 长度滤波
\[
P_{a,q}(z,u):=\frac1q\sum_{j=0}^{q-1}\omega_q^{-aj}P(\omega_q^j z,u),
\qquad
P_{a,q}^{\mathrm{core}}(z,u):=\frac1q\sum_{j=0}^{q-1}\omega_q^{-aj}P_{\mathrm{core}}(\omega_q^j z,u).
\]
则有严格恒等式
\[
\boxed{\;
P_{a,q}(z,u)=P_{a,q}^{\mathrm{core}}(z,u)+u z^2\,\mathbf 1_{\{a\equiv 2\;(\mathrm{mod}\,q)\}}.\;}
\]
进一步，定义周期对数滤波
\[
A_{a,q}(z,u):=\frac1q\sum_{j=0}^{q-1}\omega_q^{-aj}\log\zeta(\omega_q^j z,u),
\]
\[
A_{a,q}^{\mathrm{core}}(z,u):=\frac1q\sum_{j=0}^{q-1}\omega_q^{-aj}\log\zeta_{\mathrm{core}}(\omega_q^j z,u).
\]
则有闭式修正
\[
\boxed{\;
A_{a,q}(z,u)
=
A_{a,q}^{\mathrm{core}}(z,u)
+\sum_{\substack{k\ge 1\\ 2k\equiv a\;(\mathrm{mod}\,q)}}
\frac{u^k z^{2k}}{k}.\;}
\]
因此，原子长度 \(2\) 因子对 root-of-unity 长度筛的作用不是扩散到全部同余类，而是严格局域于单一剩余类 \(2\bmod q\)。
\end{theorem}
\begin{proof}
由定理 \ref{thm:xi-atomic-witt-cycle-primitive-exact-peeling}，
\[
P(z,u)=P_{\mathrm{core}}(z,u)+u z^2.
\]
对其施加 root-of-unity 投影，得到
\[
\begin{aligned}
P_{a,q}(z,u)-P_{a,q}^{\mathrm{core}}(z,u)
&=
\frac1q\sum_{j=0}^{q-1}\omega_q^{-aj}u(\omega_q^j z)^2\\
&=
u z^2\cdot \frac1q\sum_{j=0}^{q-1}\omega_q^{j(2-a)}\\
&=
u z^2\,\mathbf 1_{\{a\equiv 2\;(\mathrm{mod}\,q)\}},
\end{aligned}
\]
从而得到第一式。这正是定理 \ref{thm:xi-atomic-witt-mod-class-localization} 的 Fourier 版本。

再由
\[
\log\zeta(z,u)
=
\log\zeta_{\mathrm{core}}(z,u)
-\log(1-uz^2)
=
\log\zeta_{\mathrm{core}}(z,u)
+\sum_{k\ge 1}\frac{u^k z^{2k}}{k},
\]
同样施加 root-of-unity 投影，便有
\[
\begin{aligned}
A_{a,q}(z,u)-A_{a,q}^{\mathrm{core}}(z,u)
&=
\frac1q\sum_{j=0}^{q-1}\omega_q^{-aj}
\sum_{k\ge 1}\frac{u^k(\omega_q^j z)^{2k}}{k}\\
&=
\sum_{k\ge 1}\frac{u^k z^{2k}}{k}
\left(
\frac1q\sum_{j=0}^{q-1}\omega_q^{j(2k-a)}
\right)\\
&=
\sum_{\substack{k\ge 1\\ 2k\equiv a\;(\mathrm{mod}\,q)}}\frac{u^k z^{2k}}{k}.
\end{aligned}
\]
证毕。
\end{proof}

\begin{theorem}[平方根门槛失败的 core 不可约性]\label{thm:xi-square-root-threshold-core-irreducibility}
沿用命题 \ref{prop:real-input-40-output-dirichlet-nonrh} 的记号。对任意 \(m\ge 2\) 与 \(1\le j\le m-1\)，令
\[
\omega_m:=e^{2\pi i/m},
\qquad
\rho_{m,j}:=\frac{1}{\min\{|z|:\ \Delta(z,\omega_m^j)=0\}}.
\]
若 \(\rho_{m,j}>1\)，则必有
\[
\boxed{\;
\rho_{m,j}
=
\frac{1}{\min\{|z|:\ F(z,\omega_m^j)=0\}}.\;}
\]
特别地，任一满足
\[
\theta_{m,j}:=\frac{\log\rho_{m,j}}{\log\lambda}>\frac12
\]
的非平凡扭曲模态，都完全由 core 因子 \(F(z,\omega_m^j)\) 决定；长度 \(2\) 原子因子 \((1-uz^2)\) 不可能制造这类超平方根门槛现象。
\end{theorem}
\begin{proof}
对 \(|u|=1\) 且 \(u=\omega_m^j\)，原子因子满足
\[
1-uz^2=0
\iff
z^2=u^{-1},
\]
故其全部零点都位于单位圆上，即 \(|z|=1\)。
另一方面，
\[
\Delta(z,\omega_m^j)=(1-\omega_m^j z^2)F(z,\omega_m^j).
\]
若 \(\rho_{m,j}>1\)，则按定义
\[
\min\{|z|:\ \Delta(z,\omega_m^j)=0\}<1.
\]
由于原子因子不可能在 \(|z|<1\) 上产生零点，故达到该最小模的零点只能来自 \(F(z,\omega_m^j)\)。
于是
\[
\min\{|z|:\ \Delta(z,\omega_m^j)=0\}
=
\min\{|z|:\ F(z,\omega_m^j)=0\},
\]
对两边取倒数即得结论。证毕。
\end{proof}

\begin{theorem}[Lee--Yang 三次域符号条件下的根存在事件坍缩]\label{thm:xi-leyang-sign-conditional-root-collapse}
令
\[
g(y):=256y^3+411y^2+165y+32,
\]
并记 \(P_{10}(u)\) 为推论 \ref{cor:fold-gauge-anomaly-P10-modp-splitting-densities} 中的十次分歧多项式。
对除去有限坏素后的素数 \(p\)，定义事件
\[
A:=\{\text{\(P_{10}\bmod p\) 含一次因子}\},
\]
\[
B_{\mathrm{split}}:=\{\text{\(g\bmod p\) 全分裂}\},
\qquad
B_{\mathrm{root}}:=\{\text{\(g\bmod p\) 至少有一根}\}.
\]
并定义 Lee--Yang 二次特征
\[
\chi_{\mathrm{LY}}(p):=\mathrm{sgn}\!\bigl(\sigma_3(p)\bigr)
=
\left(\frac{-111}{p}\right)\in\{\pm1\},
\]
其中 \(\sigma_3(p)\in S_3\) 为 \(g\) 的分裂域 Frobenius 元。
则条件密度满足
\[
\boxed{\;
\delta\!\bigl(B_{\mathrm{root}}\mid \chi_{\mathrm{LY}}=+1\bigr)=\frac13,
\qquad
\delta\!\bigl(B_{\mathrm{root}}\mid \chi_{\mathrm{LY}}=-1\bigr)=1,\;}
\]
\[
\boxed{\;
\delta\!\bigl(B_{\mathrm{split}}\mid \chi_{\mathrm{LY}}=+1\bigr)=\frac13,
\qquad
\delta\!\bigl(B_{\mathrm{split}}\mid \chi_{\mathrm{LY}}=-1\bigr)=0.\;}
\]
进一步，由 \(A\) 只依赖 \(S_{10}\) 分量，而
\(\chi_{\mathrm{LY}},B_{\mathrm{split}},B_{\mathrm{root}}\) 只依赖 \(S_3\) 分量，得到
\[
\boxed{\;
\delta\!\bigl(A\cap B_{\mathrm{root}}\mid \chi_{\mathrm{LY}}=+1\bigr)
=
\frac{28319}{134400},
\qquad
\delta\!\bigl(A\cap B_{\mathrm{root}}\mid \chi_{\mathrm{LY}}=-1\bigr)
=
\frac{28319}{44800},\;}
\]
\[
\boxed{\;
\delta\!\bigl(A\cap B_{\mathrm{split}}\mid \chi_{\mathrm{LY}}=+1\bigr)
=
\frac{28319}{134400},
\qquad
\delta\!\bigl(A\cap B_{\mathrm{split}}\mid \chi_{\mathrm{LY}}=-1\bigr)
=
0.\;}
\]
\end{theorem}
\begin{proof}
由结论 \ref{con:xi-terminal-zm-leyang-cubic-modp-splitting-density}，\(g\bmod p\) 的三种分解型与 \(S_3\) 的三类共轭类对应如下：
恒等类对应全分裂 \((1)(1)(1)\)，换位类对应 \((1)(2)\)，三循环类对应不可约 \((3)\)。
因此
\[
B_{\mathrm{split}} \Longleftrightarrow \sigma_3(p)\in \{e\},
\qquad
B_{\mathrm{root}} \Longleftrightarrow \sigma_3(p)\in \{e\}\cup \{\text{换位类}\}.
\]
另一方面，
\[
\chi_{\mathrm{LY}}(p)=+1
\Longleftrightarrow
\sigma_3(p)\in \{e\}\cup\{\text{三循环类}\},
\]
\[
\chi_{\mathrm{LY}}(p)=-1
\Longleftrightarrow
\sigma_3(p)\in \{\text{换位类}\}.
\]
在 \(S_3\) 中，这三类大小分别为 \(1,3,2\)。
故在 \(\chi_{\mathrm{LY}}=+1\) 的条件下，只剩下大小为 \(1\) 与 \(2\) 的两类，其中只有恒等类产生根，于是
\[
\delta\!\bigl(B_{\mathrm{root}}\mid \chi_{\mathrm{LY}}=+1\bigr)
=
\delta\!\bigl(B_{\mathrm{split}}\mid \chi_{\mathrm{LY}}=+1\bigr)
=
\frac{1}{1+2}
=
\frac13.
\]
在 \(\chi_{\mathrm{LY}}=-1\) 的条件下，只剩下换位类，故
\[
\delta\!\bigl(B_{\mathrm{root}}\mid \chi_{\mathrm{LY}}=-1\bigr)=1,
\qquad
\delta\!\bigl(B_{\mathrm{split}}\mid \chi_{\mathrm{LY}}=-1\bigr)=0.
\]

再由命题 \ref{prop:fold-gauge-anomaly-P10-leyang-linear-disjointness} 与推论
\ref{cor:fold-gauge-anomaly-P10-leyang-chebotarev-product-law}，
\(P_{10}\) 的分解型事件与 \(g\) 的 \(S_3\)-分解型事件严格独立；
并由推论 \ref{cor:fold-gauge-anomaly-P10-modp-splitting-densities}，
\[
\delta(A)=\frac{28319}{44800}.
\]
故
\[
\delta\!\bigl(A\cap B_{\mathrm{root}}\mid \chi_{\mathrm{LY}}=\pm1\bigr)
=
\delta(A)\,
\delta\!\bigl(B_{\mathrm{root}}\mid \chi_{\mathrm{LY}}=\pm1\bigr),
\]
\[
\delta\!\bigl(A\cap B_{\mathrm{split}}\mid \chi_{\mathrm{LY}}=\pm1\bigr)
=
\delta(A)\,
\delta\!\bigl(B_{\mathrm{split}}\mid \chi_{\mathrm{LY}}=\pm1\bigr).
\]
代入前述条件密度即得四条显式公式。证毕。
\end{proof}

\begin{corollary}[原子--core--算术三分律]\label{cor:xi-atomic-core-arithmetic-trichotomy}
在真实输入 \(40\) 状态核的一参数输出位势链与 \(P_{10}\times g\) 的算术筛选链上，可将可分辨结构规范分解为
\[
\boxed{\;
\text{原子层}\ \oplus\ \text{core 谱层}\ \oplus\ \text{算术筛选层}\;}
\]
更精确地：
\begin{enumerate}
  \item \textbf{原子层：}
        由定理 \ref{thm:xi-atomic-witt-cycle-primitive-exact-peeling}、
        \ref{thm:xi-atomic-witt-energy-residue-l2-law} 与
        \ref{thm:xi-root-unity-filter-single-class-injection}，
        原子因子 \((1-uz^2)^{-1}\) 在 primitive 层只占据长度 \(2\) 的单一坐标，
        在能量 residue 侧只向 \(r\equiv n/2\pmod m\) 注入单类 periodic 修正，
        并且在长度同余筛中只向剩余类 \(2\bmod q\) 注入局域残差；
  \item \textbf{core 谱层：}
        由定理 \ref{thm:xi-square-root-threshold-core-irreducibility}，
        一切满足 \(\theta_{m,j}>1/2\) 的非平凡扭曲模态都由 core 因子 \(F(z,\omega_m^j)\) 产生，
        因而平方根门槛的失效是纯粹的 core 谱现象；
  \item \textbf{算术筛选层：}
        由定理 \ref{thm:xi-root-unity-filter-single-class-injection} 与
        \ref{thm:xi-leyang-sign-conditional-root-collapse}，
        长度同余类的 root-of-unity 残差与
        \(\chi_{\mathrm{LY}}(p)=\left(\frac{-111}{p}\right)\) 的 Chebotarev 条件化共同构成一条显式筛选链，
        其模素端表现为“\(\chi_{\mathrm{LY}}=-1\) 时必有根、\(\chi_{\mathrm{LY}}=+1\) 时仅以密度 \(1/3\) 留根”的坍缩律。
\end{enumerate}
因此，长度 \(2\) 原子只负责能量 residue 与长度同余类上的局域残差；超越平方根门槛的全部指数复杂性由 core 谱承担；而把这些谱残差投影到算术根存在事件上的，则是由 \(\left(\frac{-111}{p}\right)\) 控制的 Chebotarev--Legendre 筛选。
\end{corollary}

\paragraph{冻结质量反演、容量矩的 Mellin--Stieltjes 精确式与 core 扭曲门槛}
沿着固定冻结相中的 escort/pressure 链、连续容量曲线的 Mellin--Stieltjes 对偶，
以及真实输入 \(40\) 状态核的原子--core 剥离链，还可进一步抽出下列七条互相咬合的接口。
前五条把冻结质量、压力斜率、离散凸结构、整数阈值刚性与容量矩写成可逆账本；后两条则说明非平凡 twisted 平方根门槛完全由 core 因子承担。

\begin{theorem}[冻结区中 escort 质量对最大纤维层的有限层指数集中]\label{thm:xi-time-part57b-escort-maxfiber-layer-concentration}
沿用定理 \ref{thm:xi-terminal-fixed-freezing-tv-collapse} 与定理
\ref{thm:fold-pressure-freezing-threshold} 的记号。对固定实数 \(a>0\)，定义 escort 测度
\[
\nu_{m,a}(x):=\frac{d_m(x)^a}{S_a(m)},
\qquad
S_a(m):=\sum_{x\in X_m}d_m(x)^a.
\]
再记
\[
M_m:=\max_{x\in X_m}d_m(x),\qquad
K_m:=\#\{x\in X_m:\ d_m(x)=M_m\},
\]
\[
M_m^{(2)}:=\max\{d_m(x):\ d_m(x)<M_m\},
\qquad
\mathcal M_m:=\{x\in X_m:\ d_m(x)=M_m\},
\]
其中若次大纤维不存在则约定 \(M_m^{(2)}:=1\)。则对每个 \(m\) 都有有限层界
\[
\boxed{\;
1-\nu_{m,a}(\mathcal M_m)
\le
\frac{|X_m|-K_m}{K_m}\left(\frac{M_m^{(2)}}{M_m}\right)^a
\le
\frac{|X_m|}{K_m}\left(\frac{M_m^{(2)}}{M_m}\right)^a.\;}
\]
若进一步满足
\[
a>a_{\mathrm c}:=\frac{\log\varphi-g_*}{\alpha_*-\alpha_*^{(2)}},
\]
则 escort 测度在 \(\mathcal M_m\) 上指数集中，并且
\[
\limsup_{m\to\infty}\frac1m
\log\bigl(1-\nu_{m,a}(\mathcal M_m)\bigr)
\le
\log\varphi+a\alpha_*^{(2)}-g_*-a\alpha_*<0.
\]
\end{theorem}
\begin{proof}
写
\[
S_a(m)=K_mM_m^a+R_a(m),
\qquad
R_a(m):=\sum_{d_m(x)<M_m}d_m(x)^a.
\]
于是
\[
1-\nu_{m,a}(\mathcal M_m)=\frac{R_a(m)}{S_a(m)}
\le
\frac{R_a(m)}{K_mM_m^a}.
\]
对补集 \(X_m\setminus\mathcal M_m\) 上的每个 \(x\)，都有 \(d_m(x)\le M_m^{(2)}\)，故
\[
R_a(m)\le (|X_m|-K_m)(M_m^{(2)})^a
\le |X_m|(M_m^{(2)})^a,
\]
从而得到两条有限层不等式。

若 \(a>a_{\mathrm c}\)，则由
\[
\frac1m\log|X_m|\to\log\varphi,\qquad
\frac1m\log K_m\to g_*,
\qquad
\frac1m\log M_m\to\alpha_*,
\]
\[
\limsup_{m\to\infty}\frac1m\log M_m^{(2)}=\alpha_*^{(2)},
\]
可得
\[
\limsup_{m\to\infty}\frac1m
\log\bigl(1-\nu_{m,a}(\mathcal M_m)\bigr)
\le
\log\varphi+a\alpha_*^{(2)}-g_*-a\alpha_*.
\]
由 \(a>a_{\mathrm c}\) 的定义，右端严格为负，故得到指数集中。证毕。
\end{proof}
\leanverified{paper\_xi\_time\_part57b\_escort\_maxfiber\_layer\_concentration}

\begin{theorem}[压力谱对 \texorpdfstring{$\alpha_*$}{alpha*} 的高阶斜率恢复与冻结截距恢复]\label{thm:xi-time-part57b-pressure-spectrum-slope-intercept-recovery}
\leanverified{paper\_xi\_time\_part57b\_pressure\_spectrum\_slope\_intercept\_recovery}
设压力函数
\[
P_a:=\lim_{m\to\infty}\frac1m\log S_a(m)
\qquad(a\in\ZZ_{\ge 1})
\]
存在，并满足定理 \ref{thm:fold-pressure-groundstate-bounds} 的双边估计
\[
a\alpha_*+g_*\le P_a\le a\alpha_*+\log\varphi.
\]
则：
\begin{enumerate}
  \item 斜率参数 \(\alpha_*\) 可由压力谱唯一恢复：
  \[
  \boxed{\;
  \alpha_*=\lim_{a\to\infty}\frac{P_a}{a}.\;}
  \]
  \item 若进一步存在某个整数 \(a_0\) 使冻结等式
  \[
  P_a=a\alpha_*+g_*
  \qquad(\forall a\ge a_0)
  \]
  成立，则对每个 \(a\ge a_0\) 都有
  \[
  \boxed{\;
  g_*=P_a-a\alpha_*.\;}
  \]
\end{enumerate}
换言之，高阶压力谱首先唯一锁定最大纤维指数 \(\alpha_*\)；而一旦进入冻结区，截距 \(g_*\) 亦随之由同一压力谱唯一锁定。
\end{theorem}
\begin{proof}
由
\[
a\alpha_*+g_*\le P_a\le a\alpha_*+\log\varphi
\]
两边同时除以 \(a\)，得到
\[
\alpha_*+\frac{g_*}{a}\le \frac{P_a}{a}\le \alpha_*+\frac{\log\varphi}{a}.
\]
令 \(a\to\infty\)，左右两端都收敛到 \(\alpha_*\)，故
\[
\lim_{a\to\infty}\frac{P_a}{a}=\alpha_*.
\]
这证明了第一条。

若对所有 \(a\ge a_0\) 都有冻结等式 \(P_a=a\alpha_*+g_*\)，则直接移项即得
\[
g_*=P_a-a\alpha_*.
\]
证毕。
\end{proof}

\begin{theorem}[纤维压力谱的离散凸性与基态超额的单调下降]\label{thm:xi-time-part57b-pressure-spectrum-discrete-convexity-monotone-excess}
沿用定理 \ref{thm:xi-time-part57b-pressure-spectrum-slope-intercept-recovery} 的记号。对每个整数 \(a\ge 1\)，记
\[
S_a(m):=\sum_{x\in X_m}d_m(x)^a,
\qquad
P_a:=\lim_{m\to\infty}\frac1m\log S_a(m),
\]
并定义基态超额
\[
E_a:=P_a-a\alpha_*.
\]
则成立：
\begin{enumerate}
  \item 对每个整数 \(a\ge 1\)，压力谱满足离散凸性
  \[
  \boxed{\;2P_{a+1}\le P_a+P_{a+2}.\;}
  \]
  \item 序列 \((E_a)_{a\ge 1}\) 单调不增，即
  \[
  \boxed{\;E_{a+1}\le E_a\qquad(a\ge 1).\;}
  \]
  \item 对每个整数 \(a\ge 1\)，都有
  \[
  \boxed{\;g_*\le E_a\le \log\varphi.\;}
  \]
\end{enumerate}
\end{theorem}
\begin{proof}
对固定 \(m\)，记
\[
\mathbf d^{(m)}:=\bigl(d_m(x)\bigr)_{x\in X_m}.
\]
由幂和的对数凸性（等价于 Hölder 不等式），对每个整数 \(a\ge 1\) 有
\[
S_{a+1}(m)^2\le S_a(m)\,S_{a+2}(m).
\]
两边取对数、除以 \(m\)，再令 \(m\to\infty\)，便得到
\[
2P_{a+1}\le P_a+P_{a+2}.
\]

另一方面，由 \(d_m(x)\le M_m\) 对一切 \(x\in X_m\) 成立，
\[
S_{a+1}(m)=\sum_{x\in X_m}d_m(x)^{a+1}
\le
M_m\sum_{x\in X_m}d_m(x)^a
=
M_m\,S_a(m).
\]
故
\[
\frac1m\log S_{a+1}(m)-\frac{a+1}{m}\log M_m
\le
\frac1m\log S_a(m)-\frac{a}{m}\log M_m.
\]
令 \(m\to\infty\)，利用
\[
\lim_{m\to\infty}\frac1m\log M_m=\alpha_*,
\]
便得到
\[
E_{a+1}\le E_a.
\]

最后，由定理 \ref{thm:fold-pressure-groundstate-bounds}，
\[
a\alpha_*+g_*\le P_a\le a\alpha_*+\log\varphi.
\]
两边同时减去 \(a\alpha_*\)，即得
\[
g_*\le E_a\le \log\varphi.
\]
证毕。
\end{proof}

\begin{corollary}[冻结相整数起始点的单阈值刚性]\label{cor:xi-time-part57b-freezing-first-contact-threshold}
\leanverified{paper\_xi\_time\_part57b\_freezing\_first\_contact\_threshold}
若严格基态间隙 \(\alpha_*^{(2)}<\alpha_*\) 成立，并记
\[
a_{\mathrm c}:=\frac{\log\varphi-g_*}{\alpha_*-\alpha_*^{(2)}}<\infty,
\qquad
a_{\mathrm{fr}}:=\min\{a\in\ZZ_{\ge 1}:\ E_a=g_*\},
\]
则 \(a_{\mathrm{fr}}\) 存在，且满足
\[
\boxed{\;a_{\mathrm{fr}}\le \lfloor a_{\mathrm c}\rfloor+1.\;}
\]
并且对每个整数 \(a\ge a_{\mathrm{fr}}\)，都有
\[
\boxed{\;E_a=g_*.\;}
\]
换言之，冻结相的整数起始点是“第一次触底后永不抬升”的单阈值事件。
\end{corollary}
\begin{proof}
由定理 \ref{thm:fold-pressure-freezing-threshold}，对每个实数 \(a>a_{\mathrm c}\) 都有
\[
P_a=a\alpha_*+g_*,
\]
故对整数
\[
a_0:=\lfloor a_{\mathrm c}\rfloor+1
\]
必有 \(a_0>a_{\mathrm c}\)，从而
\[
E_{a_0}=P_{a_0}-a_0\alpha_*=g_*.
\]
因此集合 \(\{a\in\ZZ_{\ge 1}:E_a=g_*\}\) 非空，\(a_{\mathrm{fr}}\) 良定义，且
\[
a_{\mathrm{fr}}\le a_0=\lfloor a_{\mathrm c}\rfloor+1.
\]

另一方面，由定理 \ref{thm:xi-time-part57b-pressure-spectrum-discrete-convexity-monotone-excess}，
\[
E_{a+1}\le E_a
\qquad\text{且}\qquad
E_a\ge g_*.
\]
于是，一旦某个整数 \(a_1\) 满足 \(E_{a_1}=g_*\)，便对所有 \(a\ge a_1\) 都只能有
\[
g_*\le E_a\le E_{a_1}=g_*,
\]
故
\[
E_a=g_*
\qquad(a\ge a_1).
\]
取 \(a_1=a_{\mathrm{fr}}\) 即得结论。证毕。
\end{proof}

\begin{corollary}[冻结相整数矩阶尾部的精确仿射半群律]\label{cor:xi-time-part57b-pressure-affine-semigroup-law}
\leanverified{paper\_xi\_time\_part57b\_pressure\_affine\_semigroup\_law}
沿用推论 \ref{cor:xi-time-part57b-freezing-first-contact-threshold} 的记号。若整数
\[
a,b>a_{\mathrm c},
\]
则有严格恒等式
\[
\boxed{\;P_a=a\alpha_*+g_*.\;}
\]
并且
\[
\boxed{\;P_{a+1}-P_a=\alpha_*.\;}
\]
以及
\[
\boxed{\;P_{a+b}=P_a+P_b-g_*.\;}
\]
换言之，一旦进入冻结相，整数矩阶上的压力谱尾部满足精确的仿射半群律。
\end{corollary}
\begin{proof}
由定理 \ref{thm:fold-pressure-freezing-threshold}，对任意实数 \(t>a_{\mathrm c}\) 都有
\[
P_t=t\alpha_*+g_*.
\]
特别地，对任意整数 \(a,b>a_{\mathrm c}\)，立得
\[
P_a=a\alpha_*+g_*,
\qquad
P_b=b\alpha_*+g_*,
\qquad
P_{a+b}=(a+b)\alpha_*+g_*.
\]
于是
\[
P_{a+1}-P_a
=
\bigl((a+1)\alpha_*+g_*\bigr)-\bigl(a\alpha_*+g_*\bigr)
=
\alpha_*,
\]
并且
\[
P_{a+b}
=
(a+b)\alpha_*+g_*
=
\bigl(a\alpha_*+g_*\bigr)+\bigl(b\alpha_*+g_*\bigr)-g_*
=
P_a+P_b-g_*.
\]
证毕。
\end{proof}

\begin{theorem}[连续容量曲线与全部幂和矩之间的精确 Mellin--Stieltjes 反演]\label{thm:xi-time-part57b-capacity-moment-exact-mellin-stieltjes}
沿用定义 \ref{def:fold-capacity-tail}、命题 \ref{prop:fold-capacity-tail-equivalence} 与命题
\ref{prop:fold-moment-mellin-stieltjes} 的记号。对固定分辨率 \(m\)，记
\[
N_m(T):=\#\{x\in X_m:\ d_m(x)\ge T\},
\qquad
\mathcal C_m^{\mathrm{cont}}(T):=\sum_{x\in X_m}\min(d_m(x),T).
\]
则对任意实数 \(a>0\) 都有精确公式
\[
\boxed{\;
S_a(m):=\sum_{x\in X_m}d_m(x)^a
=
a\int_0^\infty T^{a-1}N_m(T)\,\dd T.\;}
\]
并且对任意 \(a>1\) 还有
\[
\boxed{\;
S_a(m)
=
a(a-1)\int_0^\infty
T^{a-2}\bigl(2^m-\mathcal C_m^{\mathrm{cont}}(T)\bigr)\,\dd T.\;}
\]
因此，连续容量曲线 \(\mathcal C_m^{\mathrm{cont}}\) 不仅控制可逆容量门槛，而且已经携带了全部实阶幂和矩 \(\{S_a(m)\}_{a>1}\) 的信息。
\end{theorem}
\begin{proof}
对任意整数 \(d\ge 1\) 与任意 \(a>0\)，都有单点恒等式
\[
d^a=a\int_0^d T^{a-1}\,\dd T.
\]
对 \(d=d_m(x)\) 在 \(x\in X_m\) 上求和，并交换有限和与积分，便得到
\[
S_a(m)=a\int_0^\infty T^{a-1}\#\{x\in X_m:\ d_m(x)\ge T\}\,\dd T
=
a\int_0^\infty T^{a-1}N_m(T)\,\dd T.
\]
这正是命题 \ref{prop:fold-moment-mellin-stieltjes} 的实阶写法。

再对 \(a>1\)，有另一单点恒等式
\[
d^a=a(a-1)\int_0^\infty T^{a-2}(d-T)_+\,\dd T.
\]
对 \(x\in X_m\) 求和得
\[
S_a(m)=
a(a-1)\int_0^\infty
T^{a-2}\sum_{x\in X_m}(d_m(x)-T)_+\,\dd T.
\]
而
\[
\sum_{x\in X_m}(d_m(x)-T)_+
=
\sum_{x\in X_m}\bigl(d_m(x)-\min(d_m(x),T)\bigr)
\]
\[
=
\sum_{x\in X_m}d_m(x)-\mathcal C_m^{\mathrm{cont}}(T)
=
2^m-\mathcal C_m^{\mathrm{cont}}(T),
\]
其中最后一步利用
\(
\sum_{x\in X_m}d_m(x)=2^m
\)。
代回即得第二式。证毕。
\end{proof}

\begin{theorem}[原子 Euler 因子不改变任何非平凡 twisted 指数]\label{thm:xi-time-part57b-atomic-factor-zero-twisted-exponent}
沿用定理 \ref{thm:xi-atomic-witt-cycle-primitive-exact-peeling} 的记号。对任意 \(m\ge 2\) 与 \(1\le j\le m-1\)，令
\[
u=\omega_m^j,\qquad
\rho_{m,j}^{\mathrm{full}}
:=
\frac{1}{\min\{|z|:\ \Delta(z,u)=0\}},
\qquad
\rho_{m,j}^{\mathrm{core}}
:=
\frac{1}{\min\{|z|:\ F(z,u)=0\}}.
\]
若记 \(\lambda:=\rho(M(1))=\varphi^2\)，并定义
\[
\theta_{m,j}^{\mathrm{full}}
:=
\frac{\log \rho_{m,j}^{\mathrm{full}}}{\log\lambda},
\qquad
\theta_{m,j}^{\mathrm{core}}
:=
\frac{\log \rho_{m,j}^{\mathrm{core}}}{\log\lambda},
\]
则成立精确关系
\[
\boxed{\;
\rho_{m,j}^{\mathrm{full}}=\max\{1,\rho_{m,j}^{\mathrm{core}}\},\;}
\]
\[
\boxed{\;
\theta_{m,j}^{\mathrm{full}}=\max\{0,\theta_{m,j}^{\mathrm{core}}\}.\;}
\]
特别地，原子 Euler 因子 \((1-uz^2)^{-1}\) 至多把 twisted 指数抬到 \(\theta=0\)；任何真正的正指数 twisted 增长都只能来自 core 因子 \(F(z,u)\)。
\end{theorem}
\begin{proof}
由分解
\[
\Delta(z,u)=(1-uz^2)\,F(z,u),
\]
full 分母的零点由两部分组成。原子因子 \(1-uz^2\) 的零点满足
\[
z^2=u^{-1},
\]
而 \(|u|=1\)，故其全部零点都位于单位圆上，即其贡献的最小模恰为 \(1\)。
另一方面，core 因子 \(F(z,u)\) 的最小零点模长按定义为
\[
\frac{1}{\rho_{m,j}^{\mathrm{core}}}.
\]
因此 full 分母全部零点中的最小模满足
\[
\min\{|z|:\ \Delta(z,u)=0\}
=
\min\left\{1,\frac{1}{\rho_{m,j}^{\mathrm{core}}}\right\}.
\]
取倒数即得
\[
\rho_{m,j}^{\mathrm{full}}=\max\{1,\rho_{m,j}^{\mathrm{core}}\}.
\]
再除以 \(\log\lambda\) 取对数，便得到
\[
\theta_{m,j}^{\mathrm{full}}=\max\{0,\theta_{m,j}^{\mathrm{core}}\}.
\]
证毕。
\end{proof}

\begin{corollary}[twisted 平方根门槛在 core 与 full 之间完全等价]\label{cor:xi-time-part57b-core-full-square-root-barrier-equivalence}
沿用定理 \ref{thm:xi-time-part57b-atomic-factor-zero-twisted-exponent} 的记号，并定义
\[
\theta_m^{\mathrm{full}}
:=
\max_{1\le j\le m-1}\theta_{m,j}^{\mathrm{full}},
\qquad
\theta_m^{\mathrm{core}}
:=
\max_{1\le j\le m-1}\theta_{m,j}^{\mathrm{core}}.
\]
则有等价关系
\[
\boxed{\;
\theta_m^{\mathrm{full}}>\frac12
\quad\Longleftrightarrow\quad
\theta_m^{\mathrm{core}}>\frac12.\;}
\]
因此，文稿中 Dirichlet/primitive 扭曲误差不具备 RH 形平方根指数的现象，在剥离原子因子 \((1-uz^2)^{-1}\) 后完全保留；相反，将该门槛失效归因于原子 Euler 因子的解释并不成立。
\end{corollary}
\begin{proof}
由定理 \ref{thm:xi-time-part57b-atomic-factor-zero-twisted-exponent}，
\[
\theta_{m,j}^{\mathrm{full}}=\max\{0,\theta_{m,j}^{\mathrm{core}}\}
\qquad(1\le j\le m-1).
\]
由于阈值 \(1/2\) 严格大于 \(0\)，故对每个 \(j\) 都有
\[
\theta_{m,j}^{\mathrm{full}}>\frac12
\quad\Longleftrightarrow\quad
\theta_{m,j}^{\mathrm{core}}>\frac12.
\]
对 \(j\) 取最大值即得结论。证毕。
\end{proof}

\begin{remark}[有限层审计接口]\label{rem:xi-time-part57b-freezing-capacity-core-twist-audit}
脚本 \texttt{scripts/exp\_xi\_freezing\_capacity\_core\_twist\_audit.py}
对三类有限层数据作逐项核验：其一是最大纤维层上的 escort 质量与有限层集中上界，
其二是尾函数/连续容量曲线与实阶矩的 Mellin--Stieltjes 两种积分恒等式，
其三是真实输入 \(40\) 状态核在若干模数下的 full/core twisted 半径与指数一致性。
\end{remark}

\paragraph{冻结微正则分解与原子 residue 的单维扰动}
在冻结 escort 态向最大纤维层的指数集中、以及单原子 Euler 因子在 periodic/core residue 侧的精确剥离都已建立之后，还可把这两条链进一步压缩为一组完全闭式的终端结论：冻结相中的 escort 态相对于基态微正则分布具有精确二分解，而原子 residue 修正在有限 Fourier 侧则形成单维等幅相位向量。

\begin{conclusion}[冻结相 escort 态相对于基态微正则分布的精确二分解与 \texorpdfstring{$f$}{f}-散度闭式]\label{con:xi-time-part57ba-freezing-escort-groundstate-fdiv-closed}
沿用定理 \ref{thm:xi-fixed-freezing-escort-renyi-spectrum-collapse} 与推论 \ref{cor:xi-fixed-freezing-escort-groundstate-tv-identity} 的记号，固定 \(a>a_{\mathrm c}\)。记
\[
\sigma_m:=\mathrm{Unif}(A_m^{\max})
=\frac{1}{K_m}\ind_{A_m^{\max}}(\cdot),
\qquad
p_{m,a}:=\pi_{m,a}(A_m^{\max}).
\]
则存在概率测度 \(\eta_{m,a}\)；当 \(p_{m,a}<1\) 时其支撑于 \(X_m\setminus A_m^{\max}\)，并且
\[
\pi_{m,a}=p_{m,a}\sigma_m+(1-p_{m,a})\eta_{m,a},
\qquad
\pi_{m,a}(\,\cdot\mid A_m^{\max})=\sigma_m.
\]
进一步，对任意满足 \(f(1)=0\) 且 \(f(0)<\infty\) 的 Csisz\'ar--Morimoto 散度
\[
D_f(\sigma_m\|\pi_{m,a})
:=
\sum_{x\in X_m}\pi_{m,a}(x)\,
f\!\left(\frac{\sigma_m(x)}{\pi_{m,a}(x)}\right),
\]
都有精确闭式
\[
\boxed{\;
D_f(\sigma_m\|\pi_{m,a})
=
p_{m,a}\,f\!\left(\frac{1}{p_{m,a}}\right)
+(1-p_{m,a})f(0).\;}
\]
特别地，
\[
\boxed{\;
\|\pi_{m,a}-\sigma_m\|_{\mathrm{TV}}=1-p_{m,a},\;}
\]
\[
\boxed{\;
D_{\mathrm{KL}}(\sigma_m\|\pi_{m,a})=-\log p_{m,a},\;}
\]
\[
\boxed{\;
\chi^2(\sigma_m\|\pi_{m,a})=\frac{1}{p_{m,a}}-1.\;}
\]
并且对任意固定的 \(f\)，上述散度满足
\[
D_f(\sigma_m\|\pi_{m,a})
=
O\!\bigl(\exp(-m\Delta(a)+o(m))\bigr).
\]
\end{conclusion}
\begin{proof}
由于 \(d_m(x)=M_m\) 在 \(A_m^{\max}\) 上恒定，对任意 \(x\in A_m^{\max}\) 都有
\[
\pi_{m,a}(x)=\frac{M_m^a}{S_a(m)}=\frac{p_{m,a}}{K_m}.
\]
因此 \(\pi_{m,a}\) 在 \(A_m^{\max}\) 上的条件分布恰为 \(\sigma_m\)。若 \(p_{m,a}<1\)，取 \(\eta_{m,a}\) 为 \(\pi_{m,a}\) 在补集 \(X_m\setminus A_m^{\max}\) 上的条件分布；若 \(p_{m,a}=1\)，由于第二项系数为零，任取 \(X_m\) 上概率测度 \(\eta_{m,a}\) 即可。于是得到所述精确分解。

再计算 \(f\)-散度。对 \(x\in A_m^{\max}\)，有
\[
\frac{\sigma_m(x)}{\pi_{m,a}(x)}
=
\frac{1/K_m}{p_{m,a}/K_m}
=
\frac{1}{p_{m,a}},
\]
而对 \(x\notin A_m^{\max}\)，有 \(\sigma_m(x)=0\)。故
\[
\begin{aligned}
D_f(\sigma_m\|\pi_{m,a})
&=
\sum_{x\in A_m^{\max}}\pi_{m,a}(x)\,
f\!\left(\frac{1}{p_{m,a}}\right)
+
\sum_{x\notin A_m^{\max}}\pi_{m,a}(x)\,f(0)\\
&=
p_{m,a}\,f\!\left(\frac{1}{p_{m,a}}\right)
+(1-p_{m,a})f(0).
\end{aligned}
\]
这给出一般闭式。

对 \(f(t)=t\log t\)，得到
\[
D_{\mathrm{KL}}(\sigma_m\|\pi_{m,a})
=
p_{m,a}\cdot \frac{1}{p_{m,a}}\log\frac{1}{p_{m,a}}
=
-\log p_{m,a}.
\]
对 \(f(t)=(t-1)^2\)，得到
\[
\chi^2(\sigma_m\|\pi_{m,a})
=
p_{m,a}\left(\frac{1}{p_{m,a}}-1\right)^2+(1-p_{m,a})
=
\frac{1}{p_{m,a}}-1.
\]
总变差公式则由支撑分解直接给出：
\[
\|\pi_{m,a}-\sigma_m\|_{\mathrm{TV}}
=
\frac12\sum_{x\in A_m^{\max}}\frac{1-p_{m,a}}{K_m}
+\frac12\sum_{x\notin A_m^{\max}}\pi_{m,a}(x)
=
1-p_{m,a}.
\]

最后，由定理 \ref{thm:xi-fixed-freezing-escort-renyi-spectrum-collapse} 的第四条，
\[
1-p_{m,a}\le \exp\!\bigl(-m\Delta(a)+o(m)\bigr).
\]
当 \(m\) 充分大时有 \(1/p_{m,a}\in[1,2]\)。由于固定的凸函数 \(f\) 在紧区间 \([1,2]\) 上 Lipschitz，故
\[
f\!\left(\frac{1}{p_{m,a}}\right)=O(1-p_{m,a}).
\]
代回一般闭式即得
\[
D_f(\sigma_m\|\pi_{m,a})
=
O(1-p_{m,a})
=
O\!\bigl(\exp(-m\Delta(a)+o(m))\bigr).
\]
证毕。
\end{proof}

\begin{conclusion}[冻结相 \texorpdfstring{$\min$}{min}-entropy 的精确有限尺度修正]\label{con:xi-time-part57ba-freezing-minentropy-finite-size}
沿用结论 \ref{con:xi-time-part57ba-freezing-escort-groundstate-fdiv-closed} 的记号。则 escort 分布的 \(\min\)-entropy 满足精确恒等式
\[
\boxed{\;
H_\infty(\pi_{m,a})
:=
-\log\max_{x\in X_m}\pi_{m,a}(x)
=
\log K_m-\log p_{m,a}.\;}
\]
从而
\[
\boxed{\;
H_\infty(\pi_{m,a})
=
\log K_m
+O\!\bigl(\exp(-m\Delta(a)+o(m))\bigr).\;}
\]
特别地，若 \(g_*=\lim_{m\to\infty}m^{-1}\log K_m\)，则
\[
\boxed{\;
H_\infty(\pi_{m,a})-g_*m
=
\bigl(\log K_m-g_*m\bigr)
+O\!\bigl(\exp(-m\Delta(a)+o(m))\bigr).\;}
\]
\end{conclusion}
\begin{proof}
由上一结论，对 \(x\in A_m^{\max}\) 有
\[
\pi_{m,a}(x)=\frac{p_{m,a}}{K_m},
\]
而对 \(x\notin A_m^{\max}\) 有 \(\pi_{m,a}(x)\le p_{m,a}/K_m\)，因为这些点的权重对应严格小于 \(M_m\) 的纤维。故最大概率恰在 \(A_m^{\max}\) 上达到，并满足
\[
\max_{x\in X_m}\pi_{m,a}(x)=\frac{p_{m,a}}{K_m}.
\]
取负对数即得第一式。

再由
\[
1-p_{m,a}\le \exp\!\bigl(-m\Delta(a)+o(m)\bigr)
\]
与标准估计
\[
-\log p_{m,a}=O(1-p_{m,a}),
\]
便得到
\[
H_\infty(\pi_{m,a})
=
\log K_m+O\!\bigl(\exp(-m\Delta(a)+o(m))\bigr).
\]
最后一式只是把主项 \(g_*m\) 从 \(\log K_m\) 中分离出来。证毕。
\end{proof}

\begin{conclusion}[原子 residue 修正的 delta--flat 极值结构]\label{con:xi-time-part57ba-atomic-residue-delta-flat-extremizer}
沿用定理 \ref{thm:xi-atomic-witt-cycle-primitive-exact-peeling} 与定理 \ref{thm:xi-atomic-witt-energy-residue-l2-law} 的记号。固定偶整数 \(n\) 与模数 \(m\ge 2\)，定义原子 residue 修正
\[
A_{m,r}^{\mathrm{cycle}}(n)
:=
A_{m,r}^{\mathrm{full}}(n)-A_{m,r}^{\mathrm{core}}(n),
\]
以及其 twisted Fourier 向量
\[
a_n^{\mathrm{cycle}}(\omega_m^j)
:=
a_n(\omega_m^j)-a_n^{\mathrm{core}}(\omega_m^j).
\]
则有显式公式
\[
\boxed{\;
A_{m,r}^{\mathrm{cycle}}(n)
=
2\,\ind_{\{r\equiv n/2\;(\mathrm{mod}\,m)\}},\;}
\]
\[
\boxed{\;
a_n^{\mathrm{cycle}}(\omega_m^j)=2\,\omega_m^{j n/2}\qquad(j\in\ZZ/m\ZZ).\;}
\]
因此：
\begin{enumerate}
  \item residue 侧支撑恰为单点，即
  \[
  \boxed{\;\#\supp A_{m,\bullet}^{\mathrm{cycle}}(n)=1.\;}
  \]
  \item twisted Fourier 侧全部模长相同，即
  \[
  \boxed{\;\abs{a_n^{\mathrm{cycle}}(\omega_m^j)}=2\qquad(\forall j\in\ZZ/m\ZZ).\;}
  \]
  \item Parseval 能量满足
  \[
  \boxed{\;
  \sum_{r\in\ZZ/m\ZZ}\abs{A_{m,r}^{\mathrm{cycle}}(n)}^2=4,\qquad
  \frac1m\sum_{j=0}^{m-1}\abs{a_n^{\mathrm{cycle}}(\omega_m^j)}^2=4.\;}
  \]
  \item 对任意 \(1\le p\le \infty\)，都有
  \[
  \boxed{\;\|A_{m,\bullet}^{\mathrm{cycle}}(n)\|_{\ell^p}=2,\;}
  \]
  以及
  \[
  \boxed{\;
  \|a_{m,\bullet}^{\mathrm{cycle}}(n)\|_{\ell^p}
  =
  2m^{1/p}\quad(1\le p<\infty),\qquad
  \|a_{m,\bullet}^{\mathrm{cycle}}(n)\|_{\ell^\infty}=2,\;}
  \]
  其中
  \(
  a_{m,\bullet}^{\mathrm{cycle}}(n):=
  (a_n^{\mathrm{cycle}}(\omega_m^j))_{j\in\ZZ/m\ZZ}.
  \)
\end{enumerate}
换言之，原子 residue 修正在 residue 侧呈现单点支撑，而在 twisted Fourier 侧呈现等幅相位向量；二者由有限循环群上的精确 Fourier 对偶所连接。
\end{conclusion}
\begin{proof}
上述两条显式公式正是定理 \ref{thm:xi-atomic-witt-energy-residue-l2-law} 与定理 \ref{thm:xi-atomic-witt-cycle-primitive-exact-peeling} 的直接重述。由 \(A_{m,r}^{\mathrm{cycle}}(n)\) 只在唯一的剩余类 \(r\equiv n/2\pmod m\) 上取值 \(2\)，立得单点支撑与
\[
\|A_{m,\bullet}^{\mathrm{cycle}}(n)\|_{\ell^p}=2
\qquad(1\le p\le \infty).
\]
另一方面，对每个 \(j\) 都有 \(\abs{\omega_m^{j n/2}}=1\)，故
\[
\abs{a_n^{\mathrm{cycle}}(\omega_m^j)}=2
\qquad(\forall j\in\ZZ/m\ZZ),
\]
从而
\[
\frac1m\sum_{j=0}^{m-1}\abs{a_n^{\mathrm{cycle}}(\omega_m^j)}^2
=
\frac1m\sum_{j=0}^{m-1}4
=
4
\]
以及
\[
\|a_{m,\bullet}^{\mathrm{cycle}}(n)\|_{\ell^p}
=
\left(\sum_{j=0}^{m-1}2^p\right)^{1/p}
=
2m^{1/p}
\qquad(1\le p<\infty),
\]
\[
\|a_{m,\bullet}^{\mathrm{cycle}}(n)\|_{\ell^\infty}=2.
\]
最后，residue 侧平方和只有唯一非零坐标贡献，故
\[
\sum_{r\in\ZZ/m\ZZ}\abs{A_{m,r}^{\mathrm{cycle}}(n)}^2=4.
\]
证毕。
\end{proof}

\begin{conclusion}[full/core residue 差的单维秩一扰动与统一投影界]\label{con:xi-time-part57ba-core-full-residue-rankone-perturbation}
沿用结论 \ref{con:xi-time-part57ba-atomic-residue-delta-flat-extremizer} 的记号，并定义差向量
\[
\Delta A_{m,r}(n)
:=
A_{m,r}^{\mathrm{full}}(n)-A_{m,r}^{\mathrm{core}}(n)
=
A_{m,r}^{\mathrm{cycle}}(n).
\]
若记 \(r_0\equiv n/2\pmod m\)，则 residue 侧有
\[
\boxed{\;
\Delta A_{m,\bullet}(n)=2e_{r_0},\;}
\]
其中 \(e_{r_0}\) 为 \(\CC^{\ZZ/m\ZZ}\) 的标准基向量；而 twisted Fourier 侧则有
\[
\boxed{\;
\bigl(a_n(\omega_m^j)-a_n^{\mathrm{core}}(\omega_m^j)\bigr)_{j\in\ZZ/m\ZZ}
=
2(\omega_m^{j n/2})_{j\in\ZZ/m\ZZ}.\;}
\]
因此 full/core residue 差始终局域在一维坐标子空间中。特别地，对任意满足
\[
Q=Q^*=Q^2,
\qquad
Q\mathbf 1=0
\]
的 Hermitian 投影 \(Q\)，都有统一界
\[
\boxed{\;
\|Q\Delta A_{m,\bullet}(n)\|_2\le 2.\;}
\]
故任意由 residue 向量经过固定正交投影所得到的去均值统计量，其 full/core 差都至多表现为一个单维扰动，并且该扰动的 \(\ell^2\) 幅值与模数 \(m\) 无关。
\end{conclusion}
\begin{proof}
由上一结论，
\[
\Delta A_{m,r}(n)=2\,\ind_{\{r=r_0\}},
\]
故 \(\Delta A_{m,\bullet}(n)=2e_{r_0}\)，这说明 residue 差确实落在一维坐标子空间中。其 twisted Fourier 图像同样由上一结论的显式公式直接给出：
\[
a_n(\omega_m^j)-a_n^{\mathrm{core}}(\omega_m^j)=2\omega_m^{j n/2}.
\]

再由 \(Q=Q^*=Q^2\) 可知 \(Q\) 是 \(\ell^2\) 上的正交投影，从而是压缩算子：
\[
\|Qv\|_2\le \|v\|_2
\qquad(\forall v\in\CC^{\ZZ/m\ZZ}).
\]
取 \(v=\Delta A_{m,\bullet}(n)\)，并利用
\[
\|\Delta A_{m,\bullet}(n)\|_2=\|2e_{r_0}\|_2=2,
\]
即可得到
\[
\|Q\Delta A_{m,\bullet}(n)\|_2\le 2.
\]
证毕。
\end{proof}

\paragraph{冻结 escort 的基态等权塌缩与实阶 Rényi 熵率刚性}
在冻结相的总变差恒等式、基态微正则二分解与实阶 Rényi 压力差公式都已建立之后，还可把基态层上的等权塌缩与全部信息型熵率的刚性收束为同一条直接接口。

\begin{corollary}[冻结 escort 的基态等权塌缩与实阶 Rényi 熵率刚性]\label{cor:xi-time-part57bb-freezing-groundstate-equidistribution-renyi-rigidity}
\leanverified{paper\_xi\_time\_part57bb\_freezing\_groundstate\_equidistribution\_renyi\_rigidity}
沿用定理 \ref{thm:xi-fixed-freezing-escort-renyi-spectrum-collapse}、推论 \ref{cor:xi-fixed-freezing-escort-groundstate-tv-identity} 与定理 \ref{thm:xi-fixed-freezing-escort-bounded-observable-collapse} 的记号。设严格基态间隙
\[
\alpha_*^{(2)}<\alpha_*,
\]
并固定
\[
a>a_{\mathrm c}.
\]
记
\[
u_m^{\max}:=\mathrm{Unif}(A_m^{\max}),
\qquad
\Delta(a):=a(\alpha_*-\alpha_*^{(2)})-(\log\varphi-g_*)>0.
\]
则成立：
\begin{enumerate}
  \item 全变差满足精确恒等式
  \[
  \|\pi_{m,a}-u_m^{\max}\|_{\mathrm{TV}}
  =
  1-\pi_{m,a}(A_m^{\max})
  \le
  \exp\!\bigl(-m\Delta(a)+o(m)\bigr).
  \]
  因而对任意有界函数 \(f_m:X_m\to\mathbb R\)，都有
  \[
  \left|
  \int f_m\,\dd\pi_{m,a}
  -
  \frac1{K_m}\sum_{x\in A_m^{\max}}f_m(x)
  \right|
  \le
  2\|f_m\|_\infty
  \exp\!\bigl(-m\Delta(a)+o(m)\bigr).
  \]
  \item 对任意实数 \(s\in(1,\infty)\)，都有
  \[
  \lim_{m\to\infty}\frac1m H_s(\pi_{m,a})=g_*.
  \]
  若进一步取 \(s=\infty\)，则同样有
  \[
  \lim_{m\to\infty}\frac1m H_\infty(\pi_{m,a})=g_*.
  \]
\end{enumerate}
换言之，冻结区内的 escort 态在分布层面指数塌缩到最大纤维层上的等权微正则态，而在信息层面则把全部实阶 Rényi 熵率统一压缩到同一基态复杂度 \(g_*\)。
\end{corollary}
\begin{proof}
第一条中的总变差恒等式与指数界正是推论 \ref{cor:xi-fixed-freezing-escort-groundstate-tv-identity} 的内容；对有界测试函数的估计则由定理 \ref{thm:xi-fixed-freezing-escort-bounded-observable-collapse} 立即得到。

对第二条，任取 \(s\in(1,\infty)\)。由于 \(a>a_{\mathrm c}\) 且 \(s>1\)，自动有
\[
as>a>a_{\mathrm c}.
\]
于是定理 \ref{thm:xi-fixed-freezing-escort-renyi-spectrum-collapse} 的第二条给出
\[
\lim_{m\to\infty}\frac1m H_s(\pi_{m,a})=g_*.
\]
而 \(s=\infty\) 的情形正是同一定理第三条的内容。证毕。
\end{proof}

\begin{corollary}[冻结 escort 在最大纤维子集上的微正则计数公式]\label{cor:xi-time-part57bb-freezing-groundstate-subset-equidistribution}
沿用推论 \ref{cor:xi-time-part57bb-freezing-groundstate-equidistribution-renyi-rigidity} 的记号。对任意子集
\[
B_m\subseteq A_m^{\max},
\]
都有精确公式
\[
\pi_{m,a}(B_m)
=
\pi_{m,a}(A_m^{\max})\frac{|B_m|}{K_m}.
\]
从而
\[
\left|
\pi_{m,a}(B_m)-\frac{|B_m|}{K_m}
\right|
=
\bigl(1-\pi_{m,a}(A_m^{\max})\bigr)\frac{|B_m|}{K_m}
\le
\exp\!\bigl(-m\Delta(a)+o(m)\bigr).
\]
因此，冻结区中最大纤维内的任意子事件都以指数速度微正则化。
\end{corollary}
\begin{proof}
由推论 \ref{cor:xi-time-part57bb-freezing-groundstate-equidistribution-renyi-rigidity} 的第一条，\(\pi_{m,a}\) 在 \(A_m^{\max}\) 上的条件分布恰为均匀分布 \(u_m^{\max}\)。故对任意 \(x\in A_m^{\max}\) 都有
\[
\pi_{m,a}(x)=\frac{\pi_{m,a}(A_m^{\max})}{K_m}.
\]
于是
\[
\pi_{m,a}(B_m)
=
\sum_{x\in B_m}\pi_{m,a}(x)
=
\pi_{m,a}(A_m^{\max})\frac{|B_m|}{K_m}.
\]
两边减去 \(|B_m|/K_m\)，便得到
\[
\left|
\pi_{m,a}(B_m)-\frac{|B_m|}{K_m}
\right|
=
\bigl(1-\pi_{m,a}(A_m^{\max})\bigr)\frac{|B_m|}{K_m}.
\]
最后再用同一推论中的指数估计
\[
1-\pi_{m,a}(A_m^{\max})\le \exp\!\bigl(-m\Delta(a)+o(m)\bigr)
\]
即得结论。证毕。
\leanverified{paper\_xi\_time\_part57bb\_freezing\_groundstate\_subset\_equidistribution}
\end{proof}

\paragraph{冻结相的均匀简并化、局部精确反演率与双温度分离}
在冻结 escort 对最大纤维层的指数集中、整数矩阶尾部的精确仿射律以及全微态精确反演阈值已经建立之后，还可把冻结端面的配分函数、局部精确编码率与全局位预算阈值进一步压缩为同一条接口。

\begin{theorem}[冻结相的均匀简并化]\label{thm:xi-time-part57bc-freezing-uniform-degeneracy}
\leanverified{paper\_xi\_time\_part57bc\_freezing\_uniform\_degeneracy}
沿用定理 \ref{thm:xi-fixed-freezing-escort-renyi-spectrum-collapse} 与推论 \ref{cor:xi-time-part57bb-freezing-groundstate-equidistribution-renyi-rigidity} 的记号。设
\[
\alpha_*^{(2)}<\alpha_*,
\qquad
a>a_{\mathrm c}:=\frac{\log\varphi-g_*}{\alpha_*-\alpha_*^{(2)}},
\]
并记
\[
A_m^{\max}:=\{x\in X_m:\ d_m(x)=M_m\},
\qquad
K_m:=|A_m^{\max}|,
\qquad
u_m^{\max}:=\mathrm{Unif}(A_m^{\max}),
\]
\[
\pi_{m,a}(x):=\frac{d_m(x)^a}{S_a(m)},
\qquad
\Delta(a):=a(\alpha_*-\alpha_*^{(2)})-(\log\varphi-g_*)>0.
\]
则有
\[
S_a(m)=K_mM_m^a\Bigl(1+O\!\bigl(e^{-m\Delta(a)+o(m)}\bigr)\Bigr),
\]
并且
\[
\|\pi_{m,a}-u_m^{\max}\|_{\mathrm{TV}}
=
O\!\bigl(e^{-m\Delta(a)+o(m)}\bigr).
\]
换言之，冻结区内的 \(a\)-阶配分函数已经彻底由最大纤维层支配，而 escort 态则以同一指数尺度塌缩到该层上的均匀分布。
\end{theorem}
\begin{proof}
记
\[
R_a(m):=\sum_{d_m(x)<M_m}d_m(x)^a.
\]
则
\[
S_a(m)=K_mM_m^a+R_a(m).
\]
另一方面，
\[
\pi_{m,a}(A_m^{\max})
=
\frac{K_mM_m^a}{S_a(m)}
=
\frac{1}{1+\varepsilon_m},
\qquad
\varepsilon_m:=\frac{R_a(m)}{K_mM_m^a}.
\]
由定理 \ref{thm:xi-time-part57b-escort-maxfiber-layer-concentration} 以及推论 \ref{cor:xi-time-part57bb-freezing-groundstate-equidistribution-renyi-rigidity}，
\[
1-\pi_{m,a}(A_m^{\max})
=
\frac{\varepsilon_m}{1+\varepsilon_m}
\le
\exp\!\bigl(-m\Delta(a)+o(m)\bigr).
\]
于是对充分大的 \(m\)，右端小于 \(1/2\)，从而
\[
\varepsilon_m
\le
2\exp\!\bigl(-m\Delta(a)+o(m)\bigr)
=
O\!\bigl(e^{-m\Delta(a)+o(m)}\bigr).
\]
代回
\[
S_a(m)=K_mM_m^a(1+\varepsilon_m)
\]
便得第一式。

第二式则由同一推论的总变差恒等式
\[
\|\pi_{m,a}-u_m^{\max}\|_{\mathrm{TV}}
=
1-\pi_{m,a}(A_m^{\max})
\]
与上式立得。证毕。
\end{proof}

\begin{theorem}[冻结支撑上的精确局部反演率]\label{thm:xi-time-part57bc-frozen-support-exact-local-inversion-rate}
沿用定理 \ref{thm:xi-full-microstate-exact-inversion-bitrate-threshold} 的记号，并定义最大纤维支撑
\[
\Omega_m^{\max}
:=
\bigcup_{x\in A_m^{\max}}\Fold_m^{-1}(x).
\]
记 \(B_m^{\max}\) 为下述最小比特数：给定读出 \(x\in A_m^{\max}\) 后，在 \(\Omega_m^{\max}\) 上实现精确微态反演所需的最小侧信息预算。则对每个 \(m\) 都有严格恒等式
\[
B_m^{\max}
=
\left\lceil \log_2 M_m\right\rceil.
\]
特别地，
\[
\lim_{m\to\infty}\frac{B_m^{\max}}{m}
=
\frac{\alpha_*}{\log 2}.
\]
因此，在冻结支撑上，局部精确可逆率恰由最大纤维指数 \(\alpha_*\) 单独锁定。
\end{theorem}
\begin{proof}
固定 \(x\in A_m^{\max}\)。由定义，
\[
\bigl|\Fold_m^{-1}(x)\bigr|=M_m.
\]
若只允许 \(B\) 比特侧信息，则在该纤维上至多能区分 \(2^B\) 个微态。要对整条最大纤维精确反演，必要条件为
\[
2^B\ge M_m,
\]
故
\[
B_m^{\max}\ge \left\lceil \log_2 M_m\right\rceil.
\]

反之，对每个 \(x\in A_m^{\max}\) 任选双射
\[
\iota_x:\Fold_m^{-1}(x)\xrightarrow{\sim}\{0,1,\dots,M_m-1\}.
\]
把 \(\iota_x(\omega)\) 编成长度 \(\lceil\log_2M_m\rceil\) 的二进制字并作为侧信息，即可在固定 \(x\) 的条件下精确恢复任意
\(\omega\in\Fold_m^{-1}(x)\)。
故下界可达，从而
\[
B_m^{\max}
=
\left\lceil \log_2 M_m\right\rceil.
\]
最后再由定理 \ref{thm:xi-full-microstate-exact-inversion-bitrate-threshold} 中
\[
\lim_{m\to\infty}\frac{1}{m}\log M_m=\alpha_*
\]
即可得到
\[
\lim_{m\to\infty}\frac{B_m^{\max}}{m}
=
\frac{\alpha_*}{\log 2}.
\]
证毕。
\end{proof}

\begin{corollary}[两种信息温度的严格分离]\label{cor:xi-time-part57bc-two-information-temperature-separation}
定义
\[
\rho_{\mathrm{global}}:=\log_2\frac{2}{\varphi},
\qquad
\rho_{\mathrm{local}}:=\frac{\alpha_*}{\log 2}.
\]
则：
\begin{enumerate}
  \item 在均匀微态输入的全局反演问题中，若一列比特预算 \(B_m\) 满足
  \[
  \mathrm{Succ}_m^{\mathrm{bin}}(B_m)\ge 1-\varepsilon
  \qquad
  (0<\varepsilon<1\ \text{固定}),
  \]
  则必有
  \[
  \liminf_{m\to\infty}\frac{B_m}{m}\ge \rho_{\mathrm{global}}.
  \]
  \item 在冻结相渐近全部质量所支撑的最大纤维族 \(\Omega_m^{\max}\) 上，精确局部反演的最小比特率恰为
  \[
  \lim_{m\to\infty}\frac{B_m^{\max}}{m}
  =
  \rho_{\mathrm{local}}.
  \]
\end{enumerate}
因此，该体系同时携带一个外层像空间亏损率 \(\rho_{\mathrm{global}}\) 与一个内层纤维定位率 \(\rho_{\mathrm{local}}\)；二者源于完全不同的机制，不能混同。
\end{corollary}
\begin{proof}
第一条正是定理 \ref{thm:xi-time-part68-entropy-deficit-triple-identification} 的后半部分。第二条则是定理 \ref{thm:xi-time-part57bc-frozen-support-exact-local-inversion-rate} 的直接重述。证毕。
\end{proof}

\begin{corollary}[冻结斜率即编码斜率]\label{cor:xi-time-part57bc-freezing-slope-coding-rate}
记
\[
\beta_{\max}
:=
\lim_{m\to\infty}\frac{B_m^{\max}}{m}
=
\frac{\alpha_*}{\log 2}.
\]
则对任意整数 \(a>a_{\mathrm c}\)，都有
\[
P_a
=
g_*+a(\log 2)\beta_{\max}.
\]
特别地，冻结相中的高阶自由能直线，其斜率恰等于最大纤维精确反演率的自然对数版本。
\end{corollary}
\begin{proof}
由定理 \ref{thm:xi-time-part57bc-frozen-support-exact-local-inversion-rate}，
\[
(\log 2)\beta_{\max}=\alpha_*.
\]
再由推论 \ref{cor:xi-time-part57b-pressure-affine-semigroup-law}，
\[
P_a=a\alpha_*+g_*
\qquad
(a>a_{\mathrm c},\ a\in\ZZ),
\]
代入即得
\[
P_a=g_*+a(\log 2)\beta_{\max}.
\]
证毕。
\end{proof}

\paragraph{碰撞势的局部对偶、素数模阈值与 dyadic 三谱迹}
碰撞位势 \(P_2(\theta)=\log\rho(e^\theta)\) 已在附录中被压缩为一条三次代数压力链：它一方面给出 \(\theta=0\) 处全部低阶累积量的 \(\sqrt5\)-闭式，另一方面在单位根扭曲与 \(u=-1\) 奇偶通道上保留足够低阶的谱数据，使局部 Legendre 对偶、介观模数阈值与 dyadic 周期迹振荡可以继续闭合为一组不再依赖外加参数化的终端结论。

\begin{theorem}[碰撞速率函数在平衡点的四阶 Cramér--Edgeworth 局部几何]\label{thm:xi-real-input-40-collision-rate-cramer-edgeworth}
沿用命题 \ref{prop:real-input-40-collision-pressure}、推论 \ref{cor:real-input-40-collision-cumulants}、\ref{cor:real-input-40-collision-cumulants-higher} 与注 \ref{rem:real-input-40-collision-ldp-param} 的记号。按注 \ref{rem:real-input-40-collision-ldp-param} 的归一化，定义
\[
I_2(\alpha):=\sup_{\theta\in\RR}\bigl(\alpha\theta-(P_2(\theta)-P_2(0))\bigr),
\qquad
\alpha_*:=P_2'(0)=\frac{3-\sqrt5}{10}.
\]
再记
\[
\kappa_2:=P_2''(0)=\frac{6\sqrt5-5}{125},\qquad
\kappa_3:=P_2^{(3)}(0)=\frac{9+10\sqrt5}{625},\qquad
\kappa_4:=P_2^{(4)}(0)=\frac{7+24\sqrt5}{3125}.
\]
则在 \(\alpha_*\) 的某个邻域内，若写
\[
\delta:=\alpha-\alpha_*,
\]
便有解析展开
\[
\boxed{\;
I_2(\alpha_*+\delta)
=
\frac{\delta^2}{2\kappa_2}
-\frac{\kappa_3}{6\kappa_2^3}\,\delta^3
+\frac{3\kappa_3^2-\kappa_2\kappa_4}{24\kappa_2^5}\,\delta^4
+O(\delta^5).\;}
\]
并且前三个非平凡局部系数可完全显式写成
\[
\boxed{\;
\frac{1}{2\kappa_2}
=
\frac{125}{62}+\frac{75\sqrt5}{31}
\approx 7.4259709133059428\;}
\]
\[
\boxed{\;
-\frac{\kappa_3}{6\kappa_2^3}
=
-\frac{849375}{59582}-\frac{525250\sqrt5}{89373}
\approx -27.3970573347852767\;}
\]
\[
\boxed{\;
\frac{3\kappa_3^2-\kappa_2\kappa_4}{24\kappa_2^5}
=
\frac{32390871875}{343549812}
+\frac{4803058125\sqrt5}{114516604}
\approx 188.068114201582058.\;}
\]
特别地，由于 \(\kappa_3>0\)，局部奇次项严格不对称：
\[
\boxed{\;
I_2(\alpha_*+\delta)-I_2(\alpha_*-\delta)
=
-\frac{\kappa_3}{3\kappa_2^3}\,\delta^3+O(\delta^5)
<0
\qquad(\delta>0\ \text{充分小}).\;}
\]
因此在均值点附近，右侧偏差严格比左侧偏差“更便宜”。
\end{theorem}
\begin{proof}
由推论 \ref{cor:real-input-40-collision-cumulants} 与推论 \ref{cor:real-input-40-collision-cumulants-higher}，
\[
P_2(\theta)-P_2(0)
=
\alpha_*\theta+\frac{\kappa_2}{2}\theta^2+\frac{\kappa_3}{6}\theta^3+\frac{\kappa_4}{24}\theta^4+O(\theta^5),
\]
故
\[
\alpha(\theta):=P_2'(\theta)
=
\alpha_*+\kappa_2\theta+\frac{\kappa_3}{2}\theta^2+\frac{\kappa_4}{6}\theta^3+O(\theta^4).
\]
由 \(\kappa_2>0\) 可知 \(\theta=0\) 处存在局部反函数 \(\theta=\theta(\delta)\)。逐次比较幂次，得到
\[
\theta(\delta)
=
\frac{\delta}{\kappa_2}
-\frac{\kappa_3}{2\kappa_2^3}\,\delta^2
+\frac{3\kappa_3^2-\kappa_2\kappa_4}{6\kappa_2^5}\,\delta^3
+O(\delta^4).
\]
另一方面，归一化 Legendre 对偶满足 \(I_2(\alpha_*)=0\) 且
\[
I_2'(\alpha)=\theta(\alpha)
\]
在 \(\alpha_*\) 邻域内成立，因此
\[
I_2(\alpha_*+\delta)=\int_0^\delta \theta(s)\,\dd s.
\]
逐项积分便得盒装展开式。再把 \(\kappa_2,\kappa_3,\kappa_4\) 的闭式代入并化简，即得三条显式系数公式。最后以 \(\delta\) 与 \(-\delta\) 相减并保留奇次项，即得第二条盒装式；由于 \(\kappa_3=(9+10\sqrt5)/625>0\)，右侧严格小于左侧。脚本 \texttt{scripts/exp\_xi\_real\_input\_40\_collision\_ldp\_median\_twist\_audit.py} 对反函数系数与速率函数系数给出符号核验。证毕。
\end{proof}

\begin{corollary}[碰撞计数的 Edgeworth 连续中位点具有显式常数位移]\label{cor:xi-real-input-40-collision-edgeworth-continuous-median-shift}
沿用定理 \ref{thm:xi-real-input-40-collision-rate-cramer-edgeworth} 的记号。令 \(S_n\) 为长度 \(n\) 上碰撞观测 \(\varphi_2\) 的总和，并把 \(\mathfrak m_n\) 定义为下列一阶 Edgeworth 半质量方程在 \(n\alpha_*\) 附近的唯一实解：
\[
\Phi\!\left(\frac{\mathfrak m_n-n\alpha_*}{\sqrt{n\kappa_2}}\right)
+\frac{\kappa_3}{6\kappa_2^{3/2}\sqrt n}
\left(
1-\left(\frac{\mathfrak m_n-n\alpha_*}{\sqrt{n\kappa_2}}\right)^2
\right)
\phi\!\left(\frac{\mathfrak m_n-n\alpha_*}{\sqrt{n\kappa_2}}\right)
=
\frac12,
\]
其中 \(\Phi,\phi\) 分别为标准正态分布函数与密度。则
\[
\boxed{\;
\mathfrak m_n
=
n\alpha_*-\frac{\kappa_3}{6\kappa_2}+o(1)
=
n\frac{3-\sqrt5}{10}
-\left(\frac{23}{310}+\frac{52\sqrt5}{2325}\right)+o(1).\;}
\]
特别地，连续中心位置始终位于 \(n\alpha_*\) 的左侧一个非零常数量级处。另一方面，由于 \(\alpha_*=(3-\sqrt5)/10\notin\QQ\)，任何整数中位数序列都不可能满足
\[
m_n=n\alpha_*+C+o(1)
\]
对某个固定常数 \(C\) 成立；因此上式刻画的是 Edgeworth 中心位置本身，而离散中位数的精确描述必须再附加独立的 lattice 修正。
\end{corollary}
\begin{proof}
对于有限状态解析加权矩阵的一维加性泛函，附录 \ref{app:pressure-analytic} 的标准解析扰动法逐字适用；结合推论 \ref{cor:real-input-40-collision-cumulants} 与推论 \ref{cor:real-input-40-collision-cumulants-higher}，碰撞总数的标准化变量
\[
Z_n:=\frac{S_n-n\alpha_*}{\sqrt{n\kappa_2}}
\]
满足一阶 Edgeworth 近似
\[
\PP(Z_n\le x)
=
\Phi(x)
+\frac{\kappa_3}{6\kappa_2^{3/2}\sqrt n}(1-x^2)\phi(x)
+O(n^{-1})
\]
在中心窗口内一致成立。令
\[
x_n:=\frac{\mathfrak m_n-n\alpha_*}{\sqrt{n\kappa_2}}.
\]
由中心极限定理知 \(x_n\to0\)。将 \(x=x_n\) 代入上式，并用
\[
\Phi(x_n)-\frac12=\phi(0)x_n+O(x_n^3),
\qquad
(1-x_n^2)\phi(x_n)=\phi(0)+O(x_n^2),
\]
得到
\[
0=\phi(0)x_n+\frac{\kappa_3}{6\kappa_2^{3/2}\sqrt n}\phi(0)+O(n^{-1}),
\]
从而
\[
x_n=-\frac{\kappa_3}{6\kappa_2^{3/2}\sqrt n}+O(n^{-1}).
\]
乘回 \(\sqrt{n\kappa_2}\) 即得
\[
\mathfrak m_n-n\alpha_*=-\frac{\kappa_3}{6\kappa_2}+o(1).
\]
最后，将 \(\kappa_2,\kappa_3\) 的闭式直接代入并化简，便得到
\[
-\frac{\kappa_3}{6\kappa_2}
=
-\frac{23}{310}-\frac{52\sqrt5}{2325}.
\]
又因 \(\alpha_*\notin\QQ\)，序列 \(n\alpha_*\) 的小数部分不可能收敛，故整数值中位数不可能收敛到某个固定常数平移的形式；若要对其作精确刻画，必须额外引入 lattice 修正。证毕。
\end{proof}

\begin{theorem}[模素数扭曲的精确介观窗口：\texorpdfstring{$p^2$}{p2} 尺度阈值律]\label{thm:xi-real-input-40-collision-prime-twist-p2-threshold}
沿用命题 \ref{prop:real-input-40-collision-pressure} 与推论 \ref{cor:real-input-40-collision-twist-fourth} 的记号。对每个奇素数 \(p\)，记
\[
\rho_p:=\rho\!\bigl(W(e^{2\pi i/p})\bigr),
\qquad
\lambda:=\rho(W(1))=\varphi^2,
\]
并定义常数
\[
c_2:=\frac{2\pi^2}{125}(6\sqrt5-5),
\qquad
c_4:=\frac{2\pi^4}{9375}(7+24\sqrt5).
\]
则当 \(p\to\infty\) 时有精确渐近
\[
\boxed{\;
\log\frac{\rho_p}{\lambda}
=
-\frac{c_2}{p^2}+\frac{c_4}{p^4}+O(p^{-6}).\;}
\]
因此，对任意沿奇素数指标取的正整数列 \(\{n_p\}\)：
\begin{enumerate}
  \item 若 \(n_p/p^2\to 0\)，则
  \[
  \left(\frac{\rho_p}{\lambda}\right)^{n_p}\to 1;
  \]
  \item 若 \(n_p/p^2\to \tau\in(0,\infty)\)，则
  \[
  \left(\frac{\rho_p}{\lambda}\right)^{n_p}\to e^{-c_2\tau};
  \]
  \item 若 \(n_p/p^2\to\infty\)，则
  \[
  \left(\frac{\rho_p}{\lambda}\right)^{n_p}\to 0.
  \]
\end{enumerate}
换言之，模 \(p\) 的非平凡扭曲通道在长度尺度 \(n\asymp p^2\) 上发生唯一的介观过渡：低于该尺度时几乎不可见；高于该尺度时则被指数压制。
\end{theorem}
\begin{proof}
推论 \ref{cor:real-input-40-collision-twist-fourth} 在 \(t_p=2\pi/p\) 处给出
\[
\log\frac{\rho_p}{\lambda}
=
-\frac{P_2''(0)}{2}\left(\frac{2\pi}{p}\right)^2
+\frac{P_2^{(4)}(0)}{24}\left(\frac{2\pi}{p}\right)^4
+O(p^{-6}),
\]
将推论 \ref{cor:real-input-40-collision-cumulants} 与推论 \ref{cor:real-input-40-collision-cumulants-higher} 中的闭式代入，即得第一条盒装式。等价地，
\[
\log\frac{\rho_p}{\lambda}
=
-\frac{c_2}{p^2}\left(1-\frac{c_4}{c_2p^2}+O(p^{-4})\right)
=
-\frac{c_2+o(1)}{p^2}.
\]
于是
\[
n_p\log\frac{\rho_p}{\lambda}
=
-\bigl(c_2+o(1)\bigr)\frac{n_p}{p^2},
\]
三种极限情形立刻分别给出 \(0\)、\(-c_2\tau\) 与 \(-\infty\)。指数化即得所述结论。
\par
由于周期迹层的模 \(p\) 扭曲通道正是 \(\Tr(W(e^{2\pi i/p})^n)\)，而 primitive 层只是在其上再施加 Möbius/Witt 投影，上述阈值便同时给出两层通道的唯一介观尺度。脚本 \texttt{scripts/exp\_xi\_real\_input\_40\_collision\_ldp\_median\_twist\_audit.py} 对若干奇素数样本的 \(p^{-6}\) 余项进行数值核验。证毕。
\end{proof}

\begin{theorem}[模 \texorpdfstring{$2$}{2} 扭曲周期迹的精确三谱分解]\label{thm:xi-real-input-40-collision-mod2-trace-three-spectrum}
沿用命题 \ref{prop:real-input-40-collision-pressure} 与注 \ref{rem:real-input-40-collision-parity} 的记号。令
\[
F=\begin{pmatrix}1&1\\ 1&0\end{pmatrix},
\qquad
A:=F\otimes F,
\qquad
W(-1):=D(-1)A,
\]
并定义
\[
a_n(-1):=\Tr\!\bigl(W(-1)^n\bigr)\qquad(n\ge 1).
\]
则 \(\det(xI-W(-1))\) 精确分解为
\[
(x+1)\bigl(x^3-2x^2-1\bigr).
\]
记 \(\rho_-\) 为三次方程
\[
x^3-2x^2-1=0
\]
的最大实根，并记其余两根为 \(\xi,\bar\xi\)。则成立：
\begin{enumerate}
  \item
  \[
  |\xi|=|\bar\xi|=\rho_-^{-1/2},
  \qquad
  \xi+\bar\xi=2-\rho_-;
  \]
  \item 存在唯一角 \(\vartheta\in(0,\pi)\) 使
  \[
  \xi=\rho_-^{-1/2}e^{i\vartheta},
  \qquad
  \cos\vartheta=\frac{(2-\rho_-)\sqrt{\rho_-}}{2};
  \]
  \item 对一切 \(n\ge 1\)，有完全闭式
  \[
  \boxed{\;
  a_n(-1)
  =
  \rho_-^n+(-1)^n+2\rho_-^{-n/2}\cos(n\vartheta).\;}
  \]
\end{enumerate}
特别地，
\[
a_n(-1)=\rho_-^n+O(1),
\qquad
\frac1n\log|a_n(-1)|\to \log\rho_-.
\]
因此模 \(2\) 奇偶扭曲下不存在额外隐藏模态：全部振荡都由单一实根 \(\rho_-\) 与单一角频率 \(\vartheta\) 完全控制。
\end{theorem}
\begin{proof}
命题 \ref{prop:real-input-40-collision-pressure} 的特征多项式公式在 \(u=-1\) 处退化为
\[
\det(xI-W(-1))=(x+1)(x^3-2x^2-1).
\]
因此 \(W(-1)\) 的四个特征值恰为
\[
-1,\qquad \rho_-,\qquad \xi,\qquad \bar\xi.
\]
对 monic 三次多项式 \(x^3-2x^2+0x-1\) 应用 Vi\`ete 关系，得到
\[
\rho_-+\xi+\bar\xi=2,
\qquad
\rho_-\xi\bar\xi=1.
\]
由于 \(\xi,\bar\xi\) 为共轭复根，第二式给出
\[
|\xi|^2=\xi\bar\xi=\rho_-^{-1},
\]
从而 \(|\xi|=|\bar\xi|=\rho_-^{-1/2}\)；第一式则给出
\[
2\Re\xi=2-\rho_-.
\]
于是存在唯一 \(\vartheta\in(0,\pi)\) 使
\[
\xi=\rho_-^{-1/2}e^{i\vartheta},
\qquad
\cos\vartheta=\frac{\Re\xi}{|\xi|}
=
\frac{(2-\rho_-)\sqrt{\rho_-}}{2}.
\]
最后，迹等于全部特征值的 \(n\) 次幂之和，故
\[
a_n(-1)
=
(-1)^n+\rho_-^n+\xi^n+\bar\xi^n
=
\rho_-^n+(-1)^n+2\rho_-^{-n/2}\cos(n\vartheta).
\]
将右式除以 \(\rho_-^n\) 得
\[
\rho_-^{-n}a_n(-1)
=
1+(-1)^n\rho_-^{-n}+2\rho_-^{-3n/2}\cos(n\vartheta),
\]
故 \(\rho_-^{-n}a_n(-1)\to1\)，于是既有 \(a_n(-1)=\rho_-^n+O(1)\)，又有
\[
\frac1n\log|a_n(-1)|\to\log\rho_-.
\]
脚本 \texttt{scripts/exp\_xi\_real\_input\_40\_collision\_ldp\_median\_twist\_audit.py} 对特征多项式分解以及前若干 \(n\) 的迹样本进行逐项核验。证毕。
\end{proof}

\begin{theorem}[模 \texorpdfstring{$2$}{2} 奇偶扭曲对应的负频率代数支]\label{thm:xi-real-input-40-collision-mod2-negative-frequency-branch}
沿用定理 \ref{thm:xi-real-input-40-collision-mod2-trace-three-spectrum} 的记号，并令
\[
\alpha_-:=
\frac{(\rho_-^2-2\rho_--1)(\rho_--1)}
{2\rho_-^3-5\rho_-^2+4\rho_-+1}.
\]
则 \(\alpha_-\) 是三次代数数，并满足不可约方程
\[
\boxed{\;
59\alpha^3-59\alpha^2+15\alpha+2=0.\;}
\]
该三次多项式仅有一个实根，且
\[
\alpha_-\approx -0.0947099154877496.
\]
特别地，
\[
\alpha_-\notin \left[0,\frac12\right].
\]
因此，模 \(2\) 奇偶扭曲不落在定理 \ref{thm:xi-time-part9p-positive-branch-global-coordinate} 的实正扭曲分支上，而对应于同一代数参数化中的唯一负频率实支。
\end{theorem}
\begin{proof}
由定理 \ref{thm:xi-real-input-40-collision-mod2-trace-three-spectrum}，
\(\rho_-\) 满足
\[
\rho_-^3-2\rho_-^2-1=0.
\]
另一方面，定理 \ref{thm:xi-time-part9p-collision-frequency-algebraic-ldp} 已给出参数化公式
\[
\alpha(\rho)=
\frac{(\rho^2-2\rho-1)(\rho-1)}
{2\rho^3-5\rho^2+4\rho+1}.
\]
因此 \(\alpha_-=\alpha(\rho_-)\) 满足二元代数系统
\[
\rho^3-2\rho^2-1=0,
\]
\[
\alpha(2\rho^3-5\rho^2+4\rho+1)-(\rho^2-2\rho-1)(\rho-1)=0.
\]
对 \(\rho\) 取 Sylvester 结果式可得
\[
2\bigl(59\alpha^3-59\alpha^2+15\alpha+2\bigr)=0,
\]
故 \(\alpha_-\) 必满足所示三次方程。由有理根判别法，其可能有理根只能取
\[
\pm 1,\ \pm 2,\ \pm \frac{1}{59},\ \pm \frac{2}{59},
\]
直接代入均不为零，因此该三次在 \(\QQ\) 上不可约。

记
\[
f(\alpha):=59\alpha^3-59\alpha^2+15\alpha+2.
\]
其判别式为
\[
\Delta_f=-625931<0,
\]
故 \(f\) 恰有一个实根。又
\[
f(0)=2>0,
\qquad
f\!\left(-\frac{1}{10}\right)=-\frac{149}{1000}<0,
\]
故该唯一实根位于 \((-\tfrac{1}{10},0)\)。由于 \(\alpha_-\) 本身就是该三次的实根，只能等于这一唯一实根，从而
\[
\alpha_-\approx -0.0947099154877496.
\]
最后，由定理 \ref{thm:xi-time-part9p-positive-branch-global-coordinate}，实正扭曲支与 \(\alpha\in(0,\tfrac12)\) 严格一一对应；既然 \(\alpha_-<0\)，该模 \(2\) 分支便不可能属于实正扭曲支。上述结式消元、残差与数值根可由脚本 \texttt{scripts/exp\_xi\_real\_input\_40\_collision\_ldp\_median\_twist\_audit.py} 逐项核验。证毕。
\end{proof}

\begin{theorem}[奇素数单位根扭曲的近主模态刚性]\label{thm:xi-real-input-40-collision-prime-near-perron-rigidity}
沿用定理 \ref{thm:xi-real-input-40-collision-prime-twist-p2-threshold} 的记号。对奇素数 \(p\) 定义
\[
\Theta_p:=\frac{\log \rho_p}{\log\lambda}.
\]
则当 \(p\to\infty\) 时有四阶渐近
\[
\boxed{\;
\Theta_p
=
1-\frac{6\sqrt5-5}{250\log\lambda}\Bigl(\frac{2\pi}{p}\Bigr)^2
+\frac{7+24\sqrt5}{75000\log\lambda}\Bigl(\frac{2\pi}{p}\Bigr)^4
+O(p^{-6}).\;}
\]
特别地，
\[
\Theta_p
=
1-\frac{1.3809571881649568\ldots}{p^2}
+\frac{1.3098894548658713\ldots}{p^4}
+O(p^{-6}),
\]
从而 \(\Theta_p\to 1\)。
\end{theorem}
\begin{proof}
由定理 \ref{thm:xi-real-input-40-collision-prime-twist-p2-threshold}，
\[
\log\frac{\rho_p}{\lambda}
=-\frac{c_2}{p^2}+\frac{c_4}{p^4}+O(p^{-6}),
\]
其中
\[
c_2=\frac{2\pi^2}{125}(6\sqrt5-5),
\qquad
c_4=\frac{2\pi^4}{9375}(7+24\sqrt5).
\]
因此
\[
\log\rho_p
=\log\lambda-\frac{c_2}{p^2}+\frac{c_4}{p^4}+O(p^{-6}).
\]
两边除以 \(\log\lambda\) 即得
\[
\Theta_p
=
1-\frac{c_2}{(\log\lambda)\,p^2}
+\frac{c_4}{(\log\lambda)\,p^4}
+O(p^{-6}),
\]
这正是盒装式。将 \(\lambda=\varphi^2\) 代入并化简，便得到后面的数值系数。最后由主项显然可知 \(\Theta_p\to 1\)。证毕。
\end{proof}

\begin{remark}[平方根门槛失效的强化]\label{rem:xi-real-input-40-collision-prime-near-perron-meaning}
上一定理强于单纯的 \(\Theta_p>\frac12\)。它表明奇素数模上的非平凡单位根扭曲并非只是在若干模数上越过平方根门槛，而是以显式 \(p^{-2}\) 速率贴近未扭曲 Perron 指数 \(1\)；换言之，这族扭曲通道在大素数极限中几乎保留了全部主模态指数。
\end{remark}

\begin{corollary}[素数模上的平方根门槛呈余有限失效]\label{cor:xi-real-input-40-collision-prime-cofinite-sqrt-failure}
沿用定理 \ref{thm:xi-real-input-40-collision-prime-near-perron-rigidity} 与
\ref{thm:xi-real-input-40-collision-mod2-trace-three-spectrum} 的记号。则存在 \(p_0\)，使得对所有奇素数 \(p>p_0\) 都有
\[
\Theta_p>\frac12.
\]
另一方面，模 \(2\) 扭曲指数
\[
\theta_2:=\frac{\log\rho_-}{\log\lambda}
\]
亦满足
\[
\theta_2>\frac12.
\]
因此，在素数模集合中，满足 RH 形平方根门槛的 Dirichlet 扭曲误差至多只剩有限个例外；等价地，平方根门槛在素数模方向呈余有限失效。
\end{corollary}
\begin{proof}
由定理 \ref{thm:xi-real-input-40-collision-prime-near-perron-rigidity}，\(\Theta_p\to 1\)，故对充分大的奇素数 \(p\) 必有 \(\Theta_p>\frac12\)。

对模 \(2\) 的情形，由定理 \ref{thm:xi-real-input-40-collision-mod2-trace-three-spectrum}，\(\rho_-\) 是三次方程
\[
x^3-2x^2-1=0
\]
的最大实根。又该多项式在 \(x=2\) 处取值为 \(-1\)，在 \(x=3\) 处取值为 \(8\)，故
\[
\rho_->2>\varphi=\lambda^{1/2}.
\]
于是
\[
\theta_2=\frac{\log\rho_-}{\log\lambda}>\frac12.
\]
综上，素数模上的平方根门槛失效只可能保留有限个例外。证毕。
\end{proof}

\begin{theorem}[模 \texorpdfstring{$2$}{2} 已是平方根门槛失效的最小证人]\label{thm:xi-real-input-40-collision-minimal-modulus-two-witness}
沿用定理 \ref{thm:xi-real-input-40-collision-mod2-trace-three-spectrum} 的记号。对每个 \(m\ge 2\)，定义
\[
\omega_m:=e^{2\pi i/m},
\qquad
\rho_{m,j}:=\rho\!\bigl(W(\omega_m^j)\bigr),
\qquad
\theta_{m,j}:=\frac{\log \rho_{m,j}}{\log\lambda},
\]
\[
\theta_m:=\max_{1\le j\le m-1}\theta_{m,j},
\qquad
\lambda=\varphi^2.
\]
再记
\[
\rho_-:=\rho_{2,1}=\rho(W(-1)).
\]
则 \(\rho_-\) 是三次方程 \(x^3-2x^2-1=0\) 的最大实根，并且
\[
\sqrt{\lambda}<\rho_-<1+\sqrt2<\lambda,
\]
从而
\[
\boxed{\;
\frac12
<
\theta_2=\theta_{2,1}
<
\frac{\log(1+\sqrt2)}{\log\lambda}
<
1.\;}
\]
特别地，\(\theta_2>\frac12\)；而由于 \(m=1\) 不存在非平凡单位根扭曲，\(m=2\) 便是平方根门槛失效的最小证人。
\end{theorem}
\begin{proof}
当 \(m=2\) 时，唯一非平凡单位根为
\[
\omega_2=-1,
\]
故
\[
\theta_2=\theta_{2,1}=\frac{\log\rho_-}{\log\lambda},
\]
其中 \(\rho_-\) 是定理 \ref{thm:xi-real-input-40-collision-mod2-trace-three-spectrum} 中三次方程
\[
f(x):=x^3-2x^2-1=0
\]
的最大实根。

现直接在两个代数点 \(x=\varphi\) 与 \(x=1+\sqrt2\) 处检验。由
\[
\varphi^2=\varphi+1,
\qquad
\varphi^3=2\varphi+1,
\]
得到
\[
f(\varphi)
=
\varphi^3-2\varphi^2-1
=
(2\varphi+1)-2(\varphi+1)-1
=
-2<0.
\]
另一方面，
\[
(1+\sqrt2)^2=3+2\sqrt2,
\qquad
(1+\sqrt2)^3=7+5\sqrt2,
\]
故
\[
f(1+\sqrt2)
=(7+5\sqrt2)-2(3+2\sqrt2)-1
=\sqrt2>0.
\]
再由
\[
f'(x)=3x^2-4x=x(3x-4)>0
\qquad \left(x>\frac43\right),
\]
可知 \(f\) 在 \((\frac43,\infty)\) 上严格递增。由于
\[
\varphi=\frac{1+\sqrt5}{2}>\frac43,
\]
故由
\[
f(\varphi)<0,
\qquad
f(1+\sqrt2)>0,
\]
得到
\[
\varphi<\rho_-<1+\sqrt2.
\]
又因
\[
\sqrt{\lambda}=\varphi,
\]
故
\[
\sqrt{\lambda}<\rho_-<1+\sqrt2.
\]
两边取对数并除以 \(\log\lambda>0\)，即得
\[
 \frac12
 =
 \frac{\log\sqrt{\lambda}}{\log\lambda}
 <
 \theta_2
 =
 \frac{\log\rho_-}{\log\lambda}
<
\frac{\log(1+\sqrt2)}{\log\lambda}.
\]
最后只需说明右端严格小于 \(1\)。由
\[
\lambda-(1+\sqrt2)=\frac{1+\sqrt5-2\sqrt2}{2},
\]
而
\[
(1+\sqrt5)^2-8=2(\sqrt5-1)>0,
\]
故 \(1+\sqrt5>2\sqrt2\)，从而
\[
\lambda-(1+\sqrt2)>0.
\]
于是
\[
\frac{\log(1+\sqrt2)}{\log\lambda}<1.
\]
再注意到 \(m=2\) 是存在非平凡单位根扭曲的最小模数，结论成立。证毕。
\end{proof}
\leanverified{paper\_xi\_real\_input\_40\_collision\_minimal\_modulus\_two\_witness}

\begin{theorem}[最坏 Dirichlet 扭曲指数趋于 \texorpdfstring{$1$}{1} 并排斥任何统一的次主 Perron 指数]\label{thm:xi-real-input-40-collision-worst-twist-exponent-tends-to-one}
沿用定理 \ref{thm:xi-real-input-40-collision-minimal-modulus-two-witness} 的记号，并定义
\[
\rho_m:=\max_{1\le j\le m-1}\rho_{m,j}.
\]
则当 \(m\to\infty\) 时有显式下界
\[
\boxed{\;
\theta_m
\ge
1-
\frac{(6\sqrt5-5)(2\pi)^2}{250\log(\varphi^2)}\,\frac1{m^2}
+
\frac{(7+24\sqrt5)(2\pi)^4}{75000\log(\varphi^2)}\,\frac1{m^4}
+
O(m^{-6}).\;}
\]
特别地，
\[
\theta_m\longrightarrow 1
\qquad (m\to\infty).
\]
因此不存在任何统一的次主 Perron 指数 \(\sigma<1\) 使得对全部 \(m\ge 2\) 都有
\[
\rho_m\le \lambda^\sigma.
\]
等价地，对任意 \(\sigma<1\)，都有充分大的全部模数 \(m\) 满足
\[
\theta_m>\sigma.
\]
\end{theorem}
\begin{proof}
取最小非平凡角
\[
t_m:=\frac{2\pi}{m}.
\]
则
\[
\omega_m=e^{it_m}
\]
是 \(m\)-次单位根，因此
\[
\rho_m\ge \rho_{m,1}=\rho\!\bigl(W(\omega_m)\bigr).
\]
由推论 \ref{cor:real-input-40-collision-twist-fourth}，
\[
\log\frac{\rho(W(e^{it}))}{\lambda}
=
-\frac{6\sqrt5-5}{250}\,t^2
+\frac{7+24\sqrt5}{75000}\,t^4
+O(t^6)
\qquad (t\to 0).
\]
代入 \(t=t_m=2\pi/m\)，得到
\[
\log\frac{\rho_{m,1}}{\lambda}
=
-\frac{(6\sqrt5-5)(2\pi)^2}{250}\,\frac1{m^2}
+\frac{(7+24\sqrt5)(2\pi)^4}{75000}\,\frac1{m^4}
+O(m^{-6}).
\]
两边除以 \(\log\lambda\)，并利用 \(\theta_m\ge \theta_{m,1}\)，即得所示下界。

另一方面，对任意单位根 \(\omega_m^j\)，都有
\[
|W(\omega_m^j)|=W(1)
\]
逐项按模取值成立。故由谱半径比较，
\[
\rho_{m,j}\le \rho(W(1))=\lambda,
\]
从而
\[
\theta_m\le 1.
\]
于是由下界与上界夹逼可知
\[
\theta_m\to 1.
\]

若存在某个 \(\sigma<1\) 使得 \(\rho_m\le \lambda^\sigma\) 对所有 \(m\ge 2\) 都成立，则必有
\[
\theta_m=\frac{\log\rho_m}{\log\lambda}\le \sigma<1
\qquad(\forall m\ge 2),
\]
这与 \(\theta_m\to 1\) 矛盾。故不存在这样的统一 \(\sigma\)。证毕。
\end{proof}

\paragraph{实输入四十状态核的同余平方根障碍与局部大偏差不对称}
在 Dirichlet 扭曲谱半径、periodic/primitive 同余残差的 Fourier--Witt 控制，以及输出位势的 Cram\'er germ 已经建立之后，还可把它们进一步压缩为两条直接服务于投稿主叙事的综合结论。

\begin{theorem}[实输入四十状态核的长度同余分布不存在统一平方根障碍律]\label{thm:xi-real-input40-congruence-square-root-barrier-synthesis}
\leanverified{paper\_xi\_real\_input40\_congruence\_square\_root\_barrier\_synthesis}
对每个模数 \(m\ge 2\) 与剩余类 \(r\in \ZZ/m\ZZ\)，定义 periodic 与 primitive 层的同余计数
\[
a_n^{(r;m)}:=A_n(r;m),
\qquad
p_n^{(r;m)}:=P_n(r;m),
\]
其中
\[
A_n:=\sum_{r\in\ZZ/m\ZZ}a_n^{(r;m)}=a_n(1),
\qquad
P_n:=\sum_{r\in\ZZ/m\ZZ}p_n^{(r;m)}=p_n(1).
\]
记
\[
\lambda:=\rho(M(1))=\varphi^2.
\]
则存在某个模数 \(m_0\ge 2\)，使得以下两组断言不可能同时成立：对某个固定 \(\varepsilon\ge 0\)，对全部剩余类 \(r\in\ZZ/m_0\ZZ\) 都有
\[
a_n^{(r;m_0)}
=
\frac{A_n}{m_0}
+O\!\bigl(\lambda^{n/2+\varepsilon}\bigr),
\]
以及
\[
p_n^{(r;m_0)}
=
\frac{P_n}{m_0}
+O\!\bigl(\lambda^{n/2+\varepsilon}\bigr).
\]
等价地，该核的长度同余类计数不存在统一的 RH 型平方根误差律：至少一个模数上的 periodic 与 primitive 两层同余分布都会突破 \(\lambda^{n/2}\) 屏障。
\end{theorem}
\begin{proof}
由命题 \ref{prop:real-input-40-output-dirichlet-nonrh}，存在某个模数 \(m_0\ge 2\) 满足
\[
\theta_{m_0}:=
\max_{1\le j\le m_0-1}\frac{\log \rho_{m_0,j}}{\log\lambda}
>
\frac12.
\]
对该模数，定理 \ref{thm:xi-dirichlet-local-limit-sqrt-barrier} 已经断言：不存在常数
\(
C<\infty
\)
使 periodic 层满足统一平方根局部极限定理，也不存在常数
\(
C_{\mathrm{prim}}<\infty
\)
使 primitive 层满足相应平方根局部极限定理。

而
\(
O(\lambda^{n/2+\varepsilon})
\)
与
\(
O(\lambda^{n/2})
\)
只相差固定乘子 \(\lambda^\varepsilon\)，故若上面两组 \(O(\lambda^{n/2+\varepsilon})\) 估计对全部剩余类同时成立，则必分别推出 periodic 与 primitive 层上的统一平方根上界，和定理 \ref{thm:xi-dirichlet-local-limit-sqrt-barrier} 矛盾。证毕。
\end{proof}

\begin{theorem}[实输入四十状态核的局部大偏差率函数严格偏向下尾]\label{thm:xi-real-input40-local-rate-strict-downward-asymmetry}
令
\[
P(\theta):=\log\lambda(e^\theta),
\qquad
\alpha_0:=P'(0),
\qquad
\sigma^2:=P''(0),
\]
\[
\tau_3:=P^{(3)}(0),
\qquad
\tau_4:=P^{(4)}(0).
\]
则表 \ref{tab:real_input_40_output_potential_cumulants_closed} 给出
\[
\sigma^2>0,
\qquad
\tau_3<0,
\qquad
\tau_4>0,
\]
且 \(\alpha_0,\sigma^2,\tau_3,\tau_4\in\QQ(\sqrt5)\)。

定义归一化 Legendre--Fenchel 对偶
\[
I(\alpha):=\sup_{\theta\in\RR}\bigl(\theta\alpha-(P(\theta)-P(0))\bigr),
\]
则当 \(\delta\to 0\) 时有
\[
\boxed{\;
I(\alpha_0+\delta)
=
\frac{\delta^2}{2\sigma^2}
-\frac{\tau_3}{6\sigma^6}\delta^3
+\frac{3\tau_3^2-\sigma^2\tau_4}{24\sigma^{10}}\delta^4
+O(\delta^5).\;}
\]
特别地，
\[
\boxed{\;
I(\alpha_0+\delta)-I(\alpha_0-\delta)
=
-\frac{\tau_3}{3\sigma^6}\delta^3
+O(\delta^5).\;}
\]
由于 \(\tau_3<0\)，故存在 \(\delta_0>0\)，使得对全部 \(0<\delta<\delta_0\) 都有
\[
I(\alpha_0+\delta)>I(\alpha_0-\delta).
\]
因此，该核的局部大偏差几何严格偏向下尾：把输出 \(1\) 的密度向均值以上推高，比向下作同幅度偏移代价更高。
\end{theorem}
\begin{proof}
定理 \ref{thm:xi-real-input-cramer-germ-quadratic-field-rigidity} 已经给出完全相同的 Legendre 对偶展开，只是使用
\[
\kappa_2:=P''(0),\qquad \kappa_3:=P^{(3)}(0),\qquad \kappa_4:=P^{(4)}(0)
\]
这一记号。将其中
\(
\kappa_2,\kappa_3,\kappa_4
\)
分别改写为
\(
\sigma^2,\tau_3,\tau_4
\)
便得到所示公式。

再把 \(\delta\) 与 \(-\delta\) 的展开相减，偶次项相消，即得第二条盒装式。由于 \(\tau_3<0\)，其主项系数严格为正，因此对充分小的正 \(\delta\) 必有
\[
I(\alpha_0+\delta)-I(\alpha_0-\delta)>0.
\]
证毕。
\end{proof}

\begin{corollary}[局部速率函数的微分刚性与右偏判据]\label{cor:xi-real-input40-local-rate-derivative-rigidity}
沿用定理 \ref{thm:xi-real-input40-local-rate-strict-downward-asymmetry} 的记号，并记
\[
\bar\alpha:=\alpha_0.
\]
则有微分恒等式
\[
\boxed{\;
I''(\bar\alpha)=\frac{1}{P''(0)},
\qquad
I^{(3)}(\bar\alpha)=-\frac{P^{(3)}(0)}{(P''(0))^3}.\;}
\]
对当前核，
\[
P''(0)=\frac{1791}{200}-\frac{3983}{1000}\sqrt5>0,
\qquad
P^{(3)}(0)=\frac{1416369}{5000}-\frac{126691}{1000}\sqrt5<0,
\]
故
\[
I^{(3)}(\bar\alpha)>0.
\]
从而当 \(h\to 0\) 时，
\[
\boxed{\;
I(\bar\alpha+h)-I(\bar\alpha-h)
=
\frac{I^{(3)}(\bar\alpha)}{3}h^3+O(h^5).\;}
\]
因此，右侧中等偏差的指数代价严格高于左侧同幅偏差；这一局部不对称由累积量闭式强制决定，而非数值偶然。
\end{corollary}
\begin{proof}
由 Legendre 驻点关系，存在解析函数 \(\Theta(\alpha)\) 使
\[
P'(\Theta(\alpha))=\alpha,
\qquad
I'(\alpha)=\Theta(\alpha).
\]
对 \(\alpha\) 求导得
\[
I''(\alpha)=\Theta'(\alpha)=\frac{1}{P''(\Theta(\alpha))}.
\]
在 \(\alpha=\bar\alpha\) 处有 \(\Theta(\bar\alpha)=0\)，于是
\[
I''(\bar\alpha)=\frac{1}{P''(0)}.
\]
再对上式求导，
\[
I^{(3)}(\alpha)
=
-\frac{P^{(3)}(\Theta(\alpha))}{(P''(\Theta(\alpha)))^2}\,\Theta'(\alpha)
=
-\frac{P^{(3)}(\Theta(\alpha))}{(P''(\Theta(\alpha)))^3},
\]
从而
\[
I^{(3)}(\bar\alpha)=-\frac{P^{(3)}(0)}{(P''(0))^3}.
\]
代入上面的闭式即得 \(I^{(3)}(\bar\alpha)>0\)。

最后，对 \(I\) 在 \(\bar\alpha\) 处作 Taylor 展开，并将 \(h\) 与 \(-h\) 的展开相减，偶次项完全相消，便得到
\[
I(\bar\alpha+h)-I(\bar\alpha-h)
=
\frac{I^{(3)}(\bar\alpha)}{3}h^3+O(h^5).
\]
由 \(I^{(3)}(\bar\alpha)>0\) 立得所述右偏判据。证毕。
\end{proof}


\paragraph{Stirling--Bernoulli 恢复、半维零块、自然扩张账本与共振整函数计数}
折叠规范常数的 Bernoulli 层级、谱零点集合的约数稀疏上界、自然扩张的 prime-register 外置实现以及共振整函数的 Cartwright 结构，在现有命题链下还可继续压缩为下列若干条直接后果。

\begin{theorem}[折叠规范常数列逐阶内生恢复全部 Bernoulli 数]\label{thm:xi-foldbin-gauge-constant-recovers-all-bernoulli}
沿用推论 \ref{cor:fold-bin-recover-2pi} 的规范常数列 \((C_m)_{m\ge 1}\)。对每个整数 \(r\ge 1\)，记
\[
\lambda_r:=\Bigl(\frac{\varphi}{2}\Bigr)^{2r-1},
\qquad
\gamma_r:=\frac{\varphi^{-1}+\varphi^{2r-3}}{r(2r-1)}.
\]
则对每个固定整数 \(R\ge 1\)，极限
\[
\boxed{\;
B_{2R}
=
\gamma_R^{-1}
\lim_{m\to\infty}
\lambda_R^{-m}
\left(
\log C_m-\log(2\pi)-\sum_{r=1}^{R-1}\gamma_r B_{2r}\lambda_r^m
\right)\;}
\]
存在；其中当 \(R=1\) 时，和式按空和解释。
因此序列 \((\log C_m)_{m\ge 1}\) 通过逐阶消元唯一决定全部偶 Bernoulli 数。
\end{theorem}
\begin{proof}
由定理 \ref{thm:fold-bin-gauge-constant-stirling-bernoulli-hierarchy}，对任意固定 \(R\ge 1\) 有渐近展开
\[
\log C_m
=
\log(2\pi)+
\sum_{r=1}^{R}\gamma_r B_{2r}\lambda_r^m
+O(\lambda_{R+1}^m).
\]
又因为
\[
\frac{\lambda_{R+1}}{\lambda_R}
=
\Bigl(\frac{\varphi}{2}\Bigr)^2<1,
\qquad
\gamma_R>0,
\]
故在已知 \(B_2,\dots,B_{2R-2}\) 时，减去前 \(R-1\) 阶项并乘以 \(\lambda_R^{-m}\)，得到
\[
\lambda_R^{-m}
\left(
\log C_m-\log(2\pi)-\sum_{r=1}^{R-1}\gamma_r B_{2r}\lambda_r^m
\right)
=
\gamma_R B_{2R}+O\!\left(\frac{\lambda_{R+1}^m}{\lambda_R^m}\right).
\]
令 \(m\to\infty\) 即得极限存在且等于 \(\gamma_R B_{2R}\)，整理便得到所示公式。证毕。
\end{proof}

\begin{corollary}[单一规范常数列完整编码 \texorpdfstring{\(\pi\)}{pi} 与全部偶值 \texorpdfstring{\(\zeta(2R)\)}{zeta(2R)}]\label{cor:xi-foldbin-gauge-constant-recovers-pi-even-zeta}
沿用定理 \ref{thm:xi-foldbin-gauge-constant-recovers-all-bernoulli} 的记号，则
\[
\boxed{\;
\pi=\frac12\exp\!\Bigl(\lim_{m\to\infty}\log C_m\Bigr).\;}
\]
并且对每个整数 \(R\ge 1\)，有
\[
\boxed{\;
\zeta(2R)
=
(-1)^{R-1}\frac{(2\pi)^{2R}}{2(2R)!}\,B_{2R},\;}
\]
其中 \(B_{2R}\) 由定理 \ref{thm:xi-foldbin-gauge-constant-recovers-all-bernoulli} 的极限公式唯一恢复。
因此，规范常数列 \((C_m)\) 单独即编码了 \(\pi\) 与全部偶值 \(\zeta\)-谱。
\end{corollary}
\begin{proof}
由推论 \ref{cor:fold-bin-recover-2pi}，
\(
\log C_m\to\log(2\pi)
\)，
故第一式立即成立。
再由定理 \ref{thm:xi-foldbin-gauge-constant-recovers-all-bernoulli}，全部 \(B_{2R}\) 可由 \((C_m)\) 逐阶恢复；而定理 \ref{thm:fold-bin-gauge-constant-stirling-bernoulli-hierarchy} 已给出标准恒等式
\[
B_{2R}=(-1)^{R-1}\frac{2(2R)!}{(2\pi)^{2R}}\zeta(2R).
\]
移项即得第二式。证毕。
\end{proof}

\begin{theorem}[fold-gauge 常数对 \texorpdfstring{$2\pi$}{2pi} 的一步 Richardson 型加速]\label{thm:xi-foldbin-gauge-constant-one-step-acceleration-2pi}
沿用推论 \ref{cor:fold-bin-recover-2pi} 的规范常数列 \((C_m)_{m\ge 1}\)，并记
\[
q:=\frac{\varphi}{2}\in(0,1),
\qquad
C_m^{\mathrm{acc}}
:=
C_m\exp\!\left(-\frac{1}{3\varphi}q^m\right).
\]
则有
\[
\boxed{\;
\log C_m^{\mathrm{acc}}-\log(2\pi)
=
-\frac{\sqrt5}{180}\,q^{3m}+O(q^{5m}).\;}
\]
从而
\[
\boxed{\;
C_m^{\mathrm{acc}}
=
2\pi\left(1-\frac{\sqrt5}{180}\,q^{3m}+O(q^{5m})\right).\;}
\]
特别地，去除首个显式 Bernoulli 模态后，\(2\pi\) 逼近误差阶立刻由 \(q^m\) 提升到 \(q^{3m}\)。
\end{theorem}
\begin{proof}
由定理 \ref{thm:xi-foldbin-gauge-constant-not-cfinite-first-two-zeta-limits}，
\[
\log C_m-\log(2\pi)
=
\frac{1}{3\varphi}\,q^m-\frac{\sqrt5}{180}\,q^{3m}+O(q^{5m}).
\]
故按定义直接得到
\[
\log C_m^{\mathrm{acc}}-\log(2\pi)
=
-\frac{\sqrt5}{180}\,q^{3m}+O(q^{5m}).
\]
令
\[
x_m:=-\frac{\sqrt5}{180}\,q^{3m}+O(q^{5m})=O(q^{3m}).
\]
则
\[
C_m^{\mathrm{acc}}
=
2\pi e^{x_m}
=
2\pi\bigl(1+x_m+O(x_m^2)\bigr).
\]
由于 \(x_m^2=O(q^{6m})=O(q^{5m})\)，遂得
\[
C_m^{\mathrm{acc}}
=
2\pi\left(1-\frac{\sqrt5}{180}\,q^{3m}+O(q^{5m})\right).
\]
证毕。
\end{proof}

\begin{theorem}[fold-gauge 常数的任意阶 Bernoulli 消模器与 \texorpdfstring{$(2R+1)$}{(2R+1)} 级加速]\label{thm:xi-foldbin-gauge-constant-arbitrary-order-bernoulli-annihilator}
沿用定理 \ref{thm:xi-foldbin-gauge-constant-recovers-all-bernoulli} 的记号。对任意固定整数 \(R\ge 1\)，定义唯一的次数至多为 \(R\) 的插值多项式
\[
\boxed{\;
p_R(z):=\prod_{r=1}^{R}\frac{z-\lambda_r}{1-\lambda_r}
=\sum_{j=0}^{R}\beta_j^{(R)}z^j.\;}
\]
则
\[
p_R(1)=1,
\qquad
p_R(\lambda_r)=0\quad(1\le r\le R).
\]
再定义加速组合
\[
\mathcal A_m^{(R)}:=\sum_{j=0}^{R}\beta_j^{(R)}\log C_{m+j}.
\]
则有
\[
\boxed{\;
\mathcal A_m^{(R)}
=
\log(2\pi)+O(\lambda_{R+1}^m)
=
\log(2\pi)+O\!\left(\left(\frac{\varphi}{2}\right)^{(2R+1)m}\right).\;}
\]
更精确地，
\[
\boxed{\;
\mathcal A_m^{(R)}
=
\log(2\pi)
+\gamma_{R+1}B_{2R+2}\,p_R(\lambda_{R+1})\,\lambda_{R+1}^m
+O(\lambda_{R+2}^m).\;}
\]
因此，对任意预先指定的阶数 \(R\)，上述线性组合都能同时消去前 \(R\) 个 Bernoulli 模态，并把逼近 \(\log(2\pi)\) 的主误差阶精确提升到 \((\varphi/2)^{(2R+1)m}\)。
\end{theorem}
\begin{proof}
多项式 \(p_R\) 的定义立即给出
\[
p_R(1)=1,
\qquad
p_R(\lambda_r)=0\quad(1\le r\le R),
\]
而唯一性来自 \(R+1\) 个插值条件对次数至多 \(R\) 的多项式的唯一确定。

由定理 \ref{thm:fold-bin-gauge-constant-stirling-bernoulli-hierarchy}，对固定 \(R\) 有渐近展开
\[
\log C_{m+j}
=
\log(2\pi)+\sum_{r=1}^{R+1}\gamma_r B_{2r}\lambda_r^{m+j}+O(\lambda_{R+2}^{m+j})
\]
对全部 \(0\le j\le R\) 同时成立。由于 \(j\) 仅在有限集合 \(\{0,\dots,R\}\) 中变化，余项可统一写成 \(O(\lambda_{R+2}^m)\)。
将该展开代入 \(\mathcal A_m^{(R)}\)，得到
\[
\mathcal A_m^{(R)}
=
\log(2\pi)\sum_{j=0}^{R}\beta_j^{(R)}
+\sum_{r=1}^{R+1}\gamma_r B_{2r}\lambda_r^m
\sum_{j=0}^{R}\beta_j^{(R)}\lambda_r^j
+O(\lambda_{R+2}^m).
\]
而
\[
\sum_{j=0}^{R}\beta_j^{(R)}=p_R(1)=1,
\qquad
\sum_{j=0}^{R}\beta_j^{(R)}\lambda_r^j=p_R(\lambda_r).
\]
故前 \(R\) 个模态全部被精确消去，仅剩
\[
\mathcal A_m^{(R)}
=
\log(2\pi)+\gamma_{R+1}B_{2R+2}p_R(\lambda_{R+1})\lambda_{R+1}^m+O(\lambda_{R+2}^m),
\]
即得更精确的盒装式。由于
\[
\lambda_{R+1}=\left(\frac{\varphi}{2}\right)^{2R+1},
\]
主误差阶便是 \((\varphi/2)^{(2R+1)m}\)。证毕。
\end{proof}

\begin{theorem}[谱零点集合的半维指数律]\label{thm:xi-fold-zero-set-half-dimensional-exponent-law}
\leanverified{paper\_xi\_fold\_zero\_set\_half\_dimensional\_exponent\_law}
沿用定理 \ref{thm:fold-zero-density-sparse} 与推论 \ref{cor:fold-zero-half-index-multiple6} 的记号，记
\[
\mathcal Z_m:=\{k\in \ZZ/(F_{m+2}\ZZ):\ \widehat d_m(k)=0\}.
\]
则沿同余子序列
\[
m+2\equiv 0\pmod 6
\]
有极限
\[
\boxed{\;
\lim_{\substack{m\to\infty\\ m+2\equiv 0\ (6)}}
\frac{1}{m}\log|\mathcal Z_m|
=
\frac12\log\varphi,\;}
\]
并且相对密度满足
\[
\boxed{\;
\lim_{\substack{m\to\infty\\ m+2\equiv 0\ (6)}}
\frac{1}{m}\log\frac{|\mathcal Z_m|}{F_{m+2}}
=
-\frac12\log\varphi.\;}
\]
换言之，\(|\mathcal Z_m|=\varphi^{m/2+o(m)}\)，而 \(|\mathcal Z_m|/F_{m+2}=\varphi^{-m/2+o(m)}\)。
\end{theorem}
\begin{proof}
若 \(m+2\equiv 0\pmod 6\)，记
\[
d:=\frac{m+2}{2}.
\]
由推论 \ref{cor:fold-zero-half-index-multiple6}，
\[
|\mathcal Z_m|\ge F_d.
\]
另一方面，由定理 \ref{thm:fold-zero-density-sparse}，
\[
\frac{|\mathcal Z_m|}{F_{m+2}}
\le
\tau(m+2)\frac{F_{\lfloor (m+2)/2\rfloor}}{F_{m+2}}
=
\tau(m+2)\frac{F_d}{F_{m+2}},
\]
故
\[
|\mathcal Z_m|\le \tau(m+2)\,F_d.
\]
由于约数函数满足 \(\tau(n)=e^{o(n)}\)，而 Fibonacci 数满足 \(F_n=\varphi^{n+o(n)}\)，遂得
\[
F_d\le |\mathcal Z_m|\le e^{o(m)}F_d=\varphi^{m/2+o(m)}.
\]
两边取 \((1/m)\log\) 即得第一条极限。
再由
\[
F_{m+2}=\varphi^{m+o(m)},
\]
从而
\[
\frac{|\mathcal Z_m|}{F_{m+2}}
=
\varphi^{-m/2+o(m)},
\]
第二条极限随即成立。证毕。
\end{proof}
\noindent 该定理表明：六倍窗上由半阶余类块
\(
S_{F_{(m+2)/2}}(F_{m+2})
\)
触发的大零点簇并非仅提供一个粗下界现象，而是把零点集合的绝对规模沿
\(
m+2\equiv 0\pmod 6
\)
这一子序列精确锁定在
\(
\exp((\log\varphi/2)m)
\)
的指数尺度上；相应相对密度则以同一半维速率衰减。

\begin{theorem}[约数驱动零余类并集的精确二元包含--排斥公式]\label{thm:xi-fold-zero-divisor-two-class-inclusion-exclusion}
\leanverified{paper_xi_fold_zero_divisor_two_class_inclusion_exclusion}
沿用推论 \ref{cor:fold-zero-divisor-lb} 与引理 \ref{lem:fold-zero-coset-intersection} 的记号，记
\[
M:=F_{m+2}.
\]
设 \(d,e\) 为 \(m+2\) 的真约数，并满足
\[
F_d\mid \frac{M}{2},
\qquad
F_e\mid \frac{M}{2}.
\]
则
\[
S_{F_d}(M)\subseteq \mathcal Z_m,
\qquad
S_{F_e}(M)\subseteq \mathcal Z_m,
\]
并且
\[
\left|S_{F_d}(M)\cup S_{F_e}(M)\right|
=
F_d+F_e-\varepsilon(d,e)F_{\gcd(d,e)},
\]
其中
\[
\varepsilon(d,e):=
\begin{cases}
1,& 2F_{\gcd(d,e)}\mid(F_e-F_d),\\
0,& \text{否则}.
\end{cases}
\]
特别地，
\[
|\mathcal Z_m|
\ge
F_d+F_e-\varepsilon(d,e)F_{\gcd(d,e)}.
\]
\end{theorem}
\begin{proof}
由推论 \ref{cor:fold-zero-divisor-lb}，
\[
S_{F_d}(M)\subseteq \mathcal Z_m,
\qquad
S_{F_e}(M)\subseteq \mathcal Z_m,
\qquad
|S_{F_d}(M)|=F_d,
\qquad
|S_{F_e}(M)|=F_e.
\]
另一方面，引理 \ref{lem:fold-zero-coset-intersection} 给出
\[
S_g(M)\cap S_h(M)\neq\varnothing
\iff
2\gcd(g,h)\mid(h-g),
\]
且一旦交集非空，就有
\[
|S_g(M)\cap S_h(M)|=\gcd(g,h).
\]
取 \(g=F_d\)、\(h=F_e\)，并使用 Fibonacci 最大公因子恒等式
\[
\gcd(F_d,F_e)=F_{\gcd(d,e)},
\]
便得到
\[
\left|S_{F_d}(M)\cap S_{F_e}(M)\right|
=
\varepsilon(d,e)F_{\gcd(d,e)}.
\]
于是由包含--排斥公式可得
\[
\left|S_{F_d}(M)\cup S_{F_e}(M)\right|
=
F_d+F_e-\varepsilon(d,e)F_{\gcd(d,e)}.
\]
再由该并集包含于 \(\mathcal Z_m\)，立得最后一个下界。证毕。
\end{proof}

上述二元包含--排斥公式还可进一步压缩为一条由 \(\gcd\) 与 \(2\)-进层级共同锁定的交结构定理，并由此得到整个零点集合的分层交半格描述。

\begin{theorem}[零块交集的 \texorpdfstring{$\gcd$}{gcd} 刚性与 \texorpdfstring{$2$}{2}-进分层判别]\label{thm:xi-fold-zero-coset-v2-intersection-rigidity}
\leanverified{paper\_xi\_fold\_zero\_coset\_v2\_intersection\_rigidity}
沿用定理 \ref{thm:fold-zero-coset-rigidity} 的记号。设 \(M\) 为偶整数，且
\[
g,h\mid \frac{M}{2}.
\]
则有精确交公式
\[
\boxed{\;
S_g(M)\cap S_h(M)=
\begin{cases}
S_{\gcd(g,h)}(M),& v_2(g)=v_2(h),\\[2mm]
\varnothing,& v_2(g)\neq v_2(h).
\end{cases}\;}
\]
\end{theorem}
\begin{proof}
由定理 \ref{thm:fold-zero-coset-rigidity}，
\[
S_g(M)=\left\{k\in \ZZ/(M\ZZ):\ k\equiv \frac{M}{2g}\pmod{M/g}\right\},
\]
等价地，
\[
S_g(M)=\left\{\frac{M}{2g}(1+2t)\ \ (\bmod M):\ t=0,1,\dots,g-1\right\}.
\]
由于这些代表元都落在区间 \([0,M)\) 内，任意 \(k\in S_g(M)\) 都满足
\[
v_2(k)=v_2(M)-v_2(g)-1.
\]
因此若 \(v_2(g)\neq v_2(h)\)，则 \(S_g(M)\) 与 \(S_h(M)\) 中元素的 \(2\)-进赋值恒不相同，故二者不可能相交。

现设 \(v_2(g)=v_2(h)\)。令
\[
d:=\gcd(g,h),\qquad a:=\frac{M}{g},\qquad b:=\frac{M}{h},\qquad L:=\operatorname{lcm}(a,b)=\frac{M}{d}.
\]
若再写
\[
g=d\alpha,\qquad h=d\beta,\qquad \gcd(\alpha,\beta)=1,
\]
则
\[
v_2(g)=v_2(h)
\iff
\alpha,\beta\ \text{皆为奇数}
\iff
2d\mid (h-g),
\]
这正与原始交非空判据一致。
于是
\[
S_g(M)=\left\{k:\ k\equiv \frac{a}{2}\pmod a\right\},
\qquad
S_h(M)=\left\{k:\ k\equiv \frac{b}{2}\pmod b\right\}.
\]
又因为 \(v_2(g)=v_2(h)\)，故
\[
\frac{g}{d},\qquad \frac{h}{d}
\]
都是奇数，从而
\[
\frac{L}{a}=\frac{g}{d},
\qquad
\frac{L}{b}=\frac{h}{d}
\]
亦皆为奇数。于是
\[
\frac{L}{2}-\frac{a}{2}
=
a\cdot \frac{L/a-1}{2}\in a\ZZ,
\qquad
\frac{L}{2}-\frac{b}{2}
=
b\cdot \frac{L/b-1}{2}\in b\ZZ,
\]
即
\[
\frac{L}{2}\equiv \frac{a}{2}\pmod a,
\qquad
\frac{L}{2}\equiv \frac{b}{2}\pmod b.
\]
故两组同余相容。由中国剩余定理，其公共解恰为单一余类
\[
k\equiv \frac{L}{2}\pmod L.
\]
而 \(L=M/d\)，故该余类正是
\[
S_d(M)=S_{\gcd(g,h)}(M).
\]
证毕。
\end{proof}

\begin{corollary}[Fibonacci 权重零点集合的 \texorpdfstring{$2$}{2}-进分层交半格分解]\label{cor:xi-fold-zero-set-v2-semilattice-decomposition}
沿用推论 \ref{cor:fold-zero-coset-union} 的记号，令
\[
M:=F_{m+2},
\qquad
\mathcal Z_m
=
\bigcup_{\substack{1\le j\le m\\ g_j\mid M/2}}S_{g_j}(M),
\qquad
g_j:=\gcd(F_{j+1},M)=F_{\gcd(j+1,m+2)}.
\]
再定义
\[
\mathcal G_m:=\left\{g_j:\ 1\le j\le m,\ g_j\mid \frac{M}{2}\right\},
\qquad
\mathcal G_{m,r}:=\{g\in\mathcal G_m:\ v_2(g)=r\},
\]
\[
\mathcal Z_{m,r}:=\bigcup_{g\in\mathcal G_{m,r}}S_g(M).
\]
则有：
\begin{enumerate}
  \item 各层 \(\mathcal Z_{m,r}\) 两两不交，且
  \[
  \boxed{\;
  \mathcal Z_m=\bigsqcup_r \mathcal Z_{m,r}.\;}
  \]
  \item 对每个固定的 \(r\)，集合族 \(\{S_g(M): g\in \mathcal G_{m,r}\}\) 对有限交封闭，并满足
  \[
  \boxed{\;
  S_g(M)\cap S_h(M)=S_{\gcd(g,h)}(M)
  \qquad (g,h\in \mathcal G_{m,r}).\;}
  \]
  \item 若 \(g,h\in \mathcal G_{m,r}\) 且 \(g\mid h\)（等价地，\(h/g\) 为奇数），则
  \[
  \boxed{\;
  S_g(M)\subseteq S_h(M).\;}
  \]
\end{enumerate}
因此，每一层 \(\mathcal Z_{m,r}\) 都是一族按 \(\gcd\) 封闭的有限交半格并；特别地，该层中的整除链只能沿奇倍方向嵌套，故只需取整除关系下的极大元作并即可生成。
\leanverified{paper\_xi\_fold\_zero\_set\_v2\_semilattice\_decomposition}
\end{corollary}
\begin{proof}
第一条与第二条直接由定理 \ref{thm:xi-fold-zero-coset-v2-intersection-rigidity} 给出：不同 \(2\)-进层的余类两两不交，而同层中任意两条余类的交恰为对应的 \(\gcd\)-余类。

第三条中，若 \(g,h\in\mathcal G_{m,r}\) 且 \(g\mid h\)，则
\[
\gcd(g,h)=g.
\]
于是由第二条
\[
S_g(M)
=
S_g(M)\cap S_h(M)
\subseteq
S_h(M).
\]
故同层中较小的除数对应较小的余类，而整个层只需取整除极大元作并即可生成。证毕。
\end{proof}

\leanverified{paper\_xi\_fold\_zero\_dyadic\_tower\_disjoint\_fibonacci\_cosets}
\begin{theorem}[六倍窗零块的 dyadic 塔给出一列两两不交的 Fibonacci 余类块]\label{thm:xi-fold-zero-dyadic-tower-disjoint-fibonacci-cosets}
沿用
\[
M:=F_{m+2},
\qquad
\mathcal Z_m:=\{k\in \ZZ/(M\ZZ):\ \widehat d_m(k)=0\}
\]
的记号。设
\[
m+2=3\cdot 2^r\cdot s,
\qquad
r\ge 1,
\qquad
s\in \NN.
\]
对每个 \(j=1,\dots,r\)，定义
\[
d_j:=\frac{m+2}{2^j},
\qquad
g_j:=F_{d_j}.
\]
则对所有 \(1\le j\le r\)，都有
\[
S_{g_j}(M)\subseteq \mathcal Z_m,
\]
并且这些余类块两两不交：
\[
S_{g_i}(M)\cap S_{g_j}(M)=\varnothing
\qquad(i\neq j).
\]
因此
\[
\boxed{\;
|\mathcal Z_m|
\ge
\sum_{j=1}^{r}F_{(m+2)/2^j}.\;}
\]
\end{theorem}
\begin{proof}
先证包含关系。对每个 \(j\)，由于 \(3\mid d_j\)，Lucas 数列模 \(2\) 的周期性给出
\[
2\mid L_{d_j}.
\]
更一般地，对所有 \(t=0,\dots,j-1\)，都有
\[
3\mid 2^t d_j,
\]
故
\[
2\mid L_{2^t d_j}.
\]
另一方面，反复使用倍角恒等式
\[
F_{2n}=F_nL_n
\]
可得
\[
F_{m+2}
=
F_{2^j d_j}
=
F_{d_j}\prod_{t=0}^{j-1}L_{2^t d_j}.
\]
由于上式乘积中的每个 Lucas 因子都为偶数，故
\[
F_{d_j}\mid \frac{F_{m+2}}{2}=\frac{M}{2}.
\]
于是由推论 \ref{cor:fold-zero-divisor-lb} 立即得到
\[
S_{g_j}(M)=S_{F_{d_j}}(M)\subseteq \mathcal Z_m,
\qquad
|S_{g_j}(M)|=F_{d_j}.
\]

再证两两不交。设 \(i<j\)。则
\[
d_i=2^{j-i}d_j,
\qquad
g_i=F_{2^{j-i}d_j},
\qquad
g_j=F_{d_j}.
\]
由 Fibonacci 最大公因子恒等式，
\[
\gcd(g_i,g_j)
=
\gcd(F_{d_i},F_{d_j})
=
F_{\gcd(d_i,d_j)}
=
F_{d_j}
=
g_j.
\]
若 \(S_{g_i}(M)\cap S_{g_j}(M)\neq\varnothing\)，则由引理 \ref{lem:fold-zero-coset-intersection} 必有
\[
2g_j\mid (g_i-g_j).
\]
但由前述倍角分解，
\[
g_i-g_j
=
F_{2^{j-i}d_j}-F_{d_j}
=
F_{d_j}\left(\prod_{t=0}^{j-i-1}L_{2^t d_j}-1\right).
\]
而每个 \(L_{2^t d_j}\) 都是偶数，所以括号中是“偶数减一”，必为奇数。故
\[
2F_{d_j}\nmid g_i-g_j,
\]
即
\[
2g_j\nmid (g_i-g_j),
\]
矛盾。于是
\[
S_{g_i}(M)\cap S_{g_j}(M)=\varnothing
\qquad(i\neq j).
\]

最后，由这些两两不交的余类块都包含于 \(\mathcal Z_m\)，得到
\[
|\mathcal Z_m|
\ge
\sum_{j=1}^{r}|S_{g_j}(M)|
=
\sum_{j=1}^{r}F_{d_j}
=
\sum_{j=1}^{r}F_{(m+2)/2^j}.
\]
证毕。
\end{proof}

\begin{corollary}[六倍窗零点集合半维指数律的尖锐性可由 dyadic 零块塔重导]\label{cor:xi-fold-zero-dyadic-tower-sharp-halfdim-exponent}
沿同余子序列
\[
m+2\equiv 0\pmod 6,
\]
有
\[
\lim_{\substack{m\to\infty\\ m+2\equiv 0\ (6)}}
\frac{1}{m}\log|\mathcal Z_m|
=
\frac12\log\varphi.
\]
等价地，
\[
|\mathcal Z_m|=\varphi^{m/2+o(m)}.
\]
因此定理 \ref{thm:xi-fold-zero-set-half-dimensional-exponent-law} 中的半维指数率上界在六倍窗上是尖锐的。
\end{corollary}
\begin{proof}
当 \(m+2\equiv 0\pmod 6\) 时，可写
\[
m+2=3\cdot 2^r\cdot s
\]
且 \(r\ge 1\)。由定理 \ref{thm:xi-fold-zero-dyadic-tower-disjoint-fibonacci-cosets}，至少有
\[
|\mathcal Z_m|
\ge
F_{(m+2)/2}.
\]
另一方面，由定理 \ref{thm:fold-zero-density-sparse}，
\[
|\mathcal Z_m|
\le
\tau(m+2)\,F_{\lfloor (m+2)/2\rfloor}.
\]
再结合 Fibonacci 渐近
\[
F_n=\varphi^{n+o(n)}
\]
以及约数函数的次指数增长
\[
\tau(n)=e^{o(n)},
\]
便得
\[
\frac{1}{m}\log F_{(m+2)/2}
\longrightarrow
\frac12\log\varphi,
\]
且
\[
\frac{1}{m}\log\!\Bigl(\tau(m+2)F_{\lfloor (m+2)/2\rfloor}\Bigr)
\longrightarrow
\frac12\log\varphi.
\]
夹逼定理即得所示极限。证毕。
\end{proof}

\begin{corollary}[最小三轴窗零点集合的三余类完全分解]\label{cor:xi-fold-zero-m58-three-coset-exhaustion}
\leanverified{paper\_xi\_fold\_zero\_m58\_three\_coset\_exhaustion}
当 \(m=58\) 时，有
\[
\mathcal Z_{58}
=
S_{F_{15}}(F_{60})
\sqcup
S_{F_{20}}(F_{60})
\sqcup
S_{F_{30}}(F_{60}),
\]
从而
\[
|\mathcal Z_{58}|=F_{15}+F_{20}+F_{30}=610+6765+832040=839415.
\]
换言之，最小三轴窗上的全部谱零点恰被三条由约数格触发的 Fibonacci 余类无重叠地完全耗尽。
\end{corollary}
\begin{proof}
记 \(M:=F_{60}\)。由
\[
\frac{F_{60}}{2}
=
774004377960
=
610\cdot 1268859636
=
6765\cdot 114413064,
\]
可见
\[
F_{15}\mid \frac{F_{60}}{2},
\qquad
F_{20}\mid \frac{F_{60}}{2}.
\]
另一方面，由推论 \ref{cor:fold-zero-half-index-multiple6}，
\[
F_{30}\mid \frac{F_{60}}{2}
\qquad\text{且}\qquad
S_{F_{30}}(F_{60})\subseteq \mathcal Z_{58}.
\]
故由推论 \ref{cor:fold-zero-divisor-lb}，
\[
S_{F_{15}}(F_{60})\subseteq \mathcal Z_{58},
\qquad
S_{F_{20}}(F_{60})\subseteq \mathcal Z_{58},
\qquad
S_{F_{30}}(F_{60})\subseteq \mathcal Z_{58}.
\]
再由定理 \ref{thm:xi-fold-zero-divisor-two-class-inclusion-exclusion} 的交叠判据，
\[
2F_{\gcd(15,20)}=2F_5=10\nmid(F_{20}-F_{15})=6155,
\]
\[
2F_{\gcd(15,30)}=2F_{15}=1220\nmid(F_{30}-F_{15})=831430,
\]
\[
2F_{\gcd(20,30)}=2F_{10}=110\nmid(F_{30}-F_{20})=825275.
\]
因此三条余类两两不交，从而
\[
\left|
S_{F_{15}}(F_{60})
\sqcup
S_{F_{20}}(F_{60})
\sqcup
S_{F_{30}}(F_{60})
\right|
=
F_{15}+F_{20}+F_{30}
=
839415.
\]
而注 \ref{rem:fold-zero-lb-m58} 已给出
\[
|\mathcal Z_{58}|=839415.
\]
故上述三余类并集必与整个 \(\mathcal Z_{58}\) 相等。证毕。
\end{proof}

\begin{theorem}[无界逆分支深度强制无限层 prime-register 外置]\label{thm:xi-natural-extension-infinite-prime-register-necessity}
\leanverified{paper\_xi\_natural\_extension\_infinite\_prime\_register\_necessity}
设 \(T:X\twoheadrightarrow X\) 为有限纤维满射，并假设其逆向分支深度无界，即对每个整数 \(N\ge 1\)，都存在长度 \(N\) 的前史片段
\[
x_{-N}\xrightarrow{\,T\,}x_{-N+1}\xrightarrow{\,T\,}\cdots\xrightarrow{\,T\,}x_{-1}\xrightarrow{\,T\,}x_0
\]
满足
\[
\#T^{-1}(x_{-j+1})\ge 2
\qquad(1\le j\le N).
\]
则在定理 \ref{thm:xi-layered-prime-register-sharp-natural-extension} 的分层 prime-register 外置语义下，不存在任何仅使用有限个 prime slices 的精确单射实现能够把自然扩张
\[
X^{\leftarrow}
:=
\{(x_0,x_{-1},x_{-2},\dots):T(x_{-n})=x_{-n+1}\}
\]
共轭到某个有限层账本系统上。
等价地，对全自然扩张的精确可逆外置，所需有效 prime-register 层数必为无穷。

另一方面，定理 \ref{thm:killo-natural-extension-branch-register} 给出共轭
\[
X^{\leftarrow}\cong \widehat X\subseteq X\times\NN^\NN,
\]
故可数无穷层账本足以完成精确编码。
\end{theorem}
\begin{proof}
反设存在某个有限整数 \(k\ge 1\)，使自然扩张在同一分层 prime-register 语义下可由仅含 \(k\) 个 prime slices 的账本精确单射外置。
由于该外置与自然扩张共轭，其码字在像空间上具有精确逆映射；再与前史截断
\[
\pi_N:X^{\leftarrow}\to X^{N+1},
\qquad
\pi_N(x_0,x_{-1},x_{-2},\dots):=(x_0,x_{-1},\dots,x_{-N})
\]
复合，便知同一有限层账本可对任意固定深度 \(N\) 的前史片段作精确恢复。

取 \(N:=k+1\)。
由无界逆分支深度假设，存在一条长度 \(N\) 的前史片段，其前 \(N\) 层每层都至少出现一次二歧分支。
于是上述有限层账本在该前史族上给出一个深度 \(N\) 的精确分层外置，而所用 prime slices 仍只有原来的 \(k\) 个。
这与定理 \ref{thm:xi-layered-prime-register-sharp-natural-extension} 矛盾：在最坏情形下，深度 \(N\) 的精确分层外置至少需要 \(N\) 个有效 prime slices。

故有限层 prime-register 外置不可能精确实现整个自然扩张。
另一方面，定理 \ref{thm:killo-natural-extension-branch-register} 已把 \(X^{\leftarrow}\) 共轭到 \(X\times\NN^\NN\) 上的右移插入系统，说明可数无穷层账本足以给出精确实现。证毕。
\end{proof}

\begin{theorem}[共振整函数零点计数的线性主项]\label{thm:xi-fold-resonance-entire-zero-counting-linear}
沿用定理 \ref{thm:fold-resonance-entire-lp} 的记号，定义
\[
\mathcal F_\varphi(z):=\prod_{n=2}^{\infty}\cos(\pi z\varphi^{-n}).
\]
则 \(\mathcal F_\varphi\) 属于 Laguerre--P\'olya 类，且全部零点均为实单零点：
\[
\mathcal Z(\mathcal F_\varphi)=\Bigl\{(k+\tfrac12)\varphi^n:\ k\in\ZZ,\ n\ge 2\Bigr\}.
\]
若记正零点计数函数
\[
N_+(R):=\#\{x\in(0,R]:\ \mathcal F_\varphi(x)=0\},
\]
则当 \(R\to\infty\) 时有
\[
\boxed{\;
N_+(R)=R+O(\log R).\;}
\]
相应地，双侧零点计数满足
\[
\boxed{\;
\#\bigl(\mathcal Z(\mathcal F_\varphi)\cap[-R,R]\bigr)=2R+O(\log R).\;}
\]
\end{theorem}
\begin{proof}
由定理 \ref{thm:fold-resonance-entire-lp}，正零点恰为
\[
\Bigl\{(k+\tfrac12)\varphi^n:\ k\in\ZZ_{\ge 0},\ n\ge 2\Bigr\}.
\]
对每个固定的 \(n\ge 2\)，落在 \((0,R]\) 内的该层零点数为
\[
N_n(R)
:=
\#\Bigl\{k\in\ZZ_{\ge 0}:\ (k+\tfrac12)\varphi^n\le R\Bigr\}
=
\left\lfloor \frac{R}{\varphi^n}+\frac12\right\rfloor
=
\frac{R}{\varphi^n}+O(1).
\]
仅当 \(\varphi^n\le 2R\) 时该项才可能非零，故
\[
N_+(R)
=
\sum_{n=2}^{\lfloor \log_\varphi(2R)\rfloor}N_n(R)
=
R\sum_{n=2}^{\lfloor \log_\varphi(2R)\rfloor}\varphi^{-n}
+O(\log R).
\]
又因为
\[
\sum_{n=2}^{\infty}\varphi^{-n}=1,
\]
且有限截断尾和为 \(O(R^{-1})\)，遂得
\[
N_+(R)=R+O(\log R).
\]
最后，由零点集关于原点对称，立得
\[
\#\bigl(\mathcal Z(\mathcal F_\varphi)\cap[-R,R]\bigr)=2R+O(\log R).
\]
证毕。
\end{proof}

\paragraph{fold-gauge 常数对 \texorpdfstring{$2\pi$}{2pi} 的最终单侧逼近与差分完全单调性}
在 \(C_m\) 的 Stirling--Bernoulli 层级、首个显式消模器以及任意阶 Bernoulli 消模加速已经建立之后，原始误差序列本身还可进一步压缩为一条具有固定差分符号的尾部刚性链。

\begin{theorem}[fold-gauge 常数最终从上方单调逼近 \texorpdfstring{$2\pi$}{2pi}，并具有任意固定阶的最终完全单调性]\label{thm:xi-foldbin-gauge-constant-eventual-complete-monotonicity}
沿用推论 \ref{cor:fold-bin-recover-2pi} 的规范常数列 \((C_m)_{m\ge 1}\)，并定义
\[
E_m:=\log C_m-\log(2\pi),
\qquad
(\Delta u)_m:=u_{m+1}-u_m.
\]
则成立：
\begin{enumerate}
  \item 存在 \(m_0\)，使得对全部 \(m\ge m_0\) 都有
  \[
  C_m>2\pi,
  \qquad
  C_{m+1}<C_m.
  \]
  \item 对每个固定整数 \(r\ge 0\)，存在 \(m_r\)，使得对全部 \(m\ge m_r\) 都有
  \[
  (-1)^r(\Delta^r E)_m>0.
  \]
  换言之，误差序列 \((E_m)_{m\ge 1}\) 在充分大 \(m\) 上具有任意固定阶的最终完全单调性。
  \item 若记
  \[
  q:=\frac{\varphi}{2}\in(0,1),
  \]
  则有显式主导展开
  \[
  \boxed{\;
  \log C_m
  =
  \log(2\pi)
  +\frac{1}{3\varphi}q^m
  -\frac{\sqrt5}{180}q^{3m}
  +O(q^{5m}).\;}
  \]
  从而
  \[
  \boxed{\;
  \lim_{m\to\infty}
  \Bigl(\frac{2}{\varphi}\Bigr)^m\bigl(C_m-2\pi\bigr)
  =
  \frac{2\pi}{3\varphi}.\;}
  \]
\end{enumerate}
\end{theorem}
\begin{proof}
由定理 \ref{thm:fold-bin-gauge-constant-stirling-bernoulli-hierarchy} 在 \(R=2\) 时的展开，或等价地由定理 \ref{thm:xi-foldbin-gauge-constant-not-cfinite-first-two-zeta-limits}，有
\[
E_m
=
\frac{1}{3\varphi}q^m
-\frac{\sqrt5}{180}q^{3m}
+O(q^{5m}).
\]
这就得到第 \(3\) 条中的主导展开。

现固定 \(r\ge 0\)。对任意 \(\alpha\in(0,1)\)，几何序列满足
\[
\Delta^r(\alpha^m)=(\alpha-1)^r\alpha^m.
\]
因此
\[
\Delta^r E_m
=
\frac{1}{3\varphi}(q-1)^r q^m
-\frac{\sqrt5}{180}(q^3-1)^r q^{3m}
+O(q^{5m}).
\]
把右端提取出主项，可写为
\[
\Delta^r E_m
=
\frac{1}{3\varphi}(q-1)^r q^m
\left(
1+O(q^{2m})
\right).
\]
由于 \(q\in(0,1)\) 且 \(1/(3\varphi)>0\)，主项符号恰为 \((q-1)^r\) 的符号，即 \((-1)^r\)。故存在 \(m_r\)，使得对全部 \(m\ge m_r\) 都有
\[
(-1)^r(\Delta^r E)_m>0.
\]
这便得到第 \(2\) 条。

取 \(r=0\) 得 \(E_m>0\) 对充分大 \(m\) 成立，于是
\[
\log C_m>\log(2\pi)
\qquad\Longrightarrow\qquad
C_m>2\pi.
\]
取 \(r=1\) 得
\[
(\Delta E)_m<0
\]
对充分大 \(m\) 成立，即 \(E_{m+1}<E_m\)。由于指数函数严格递增，遂有
\[
C_{m+1}<C_m
\]
对充分大 \(m\) 也成立，从而得到第 \(1\) 条。

最后，由 \(E_m=O(q^m)\) 与指数展开
\[
e^{E_m}=1+E_m+O(E_m^2)
\]
可得
\[
C_m-2\pi
=
2\pi\bigl(E_m+O(E_m^2)\bigr)
=
\frac{2\pi}{3\varphi}q^m+O(q^{2m}).
\]
两边同乘 \(q^{-m}=(2/\varphi)^m\)，即得
\[
\lim_{m\to\infty}
\Bigl(\frac{2}{\varphi}\Bigr)^m\bigl(C_m-2\pi\bigr)
=
\frac{2\pi}{3\varphi}.
\]
证毕。
\end{proof}

\paragraph{fold-gauge 规范常数的显式三模渐近、归一化缺陷与超收敛上包络}
在最终单侧逼近、任意固定阶差分尾符号与有限阶 Bernoulli 模态消去都已建立之后，规范常数列还可继续压缩为一组完全显式的有效渐近公式。它们把 Stirling--Bernoulli 层级的前三个非平凡模态、首项消元后的具体上侧包络以及相邻比值的首个二次修正统一到同一条闭式链中。

\begin{theorem}[\texorpdfstring{$\log C_m$}{log Cm} 的前三个非平凡 Bernoulli 模态完全显式化]\label{thm:xi-time-part59ac-foldgauge-three-bernoulli-modes}
记
\[
q_0:=\frac{\varphi}{2},
\qquad
L_m:=\log C_m-\log(2\pi).
\]
则有精确三项展开
\[
\boxed{\;
L_m
=
\frac{1}{3\varphi}\,q_0^m
-\frac{\sqrt5}{180}\,q_0^{3m}
+\frac{\varphi}{210}\,q_0^{5m}
+O(q_0^{7m}).\;}
\]
从而
\[
\boxed{\;
C_m
=
2\pi\exp\!\left(
\frac{1}{3\varphi}\,q_0^m
-\frac{\sqrt5}{180}\,q_0^{3m}
+\frac{\varphi}{210}\,q_0^{5m}
+O(q_0^{7m})
\right).\;}
\]
进一步展开指数，得到
\[
\boxed{\;
C_m
=
2\pi\left(
1+\frac{1}{3\varphi}\,q_0^m
+\frac{1}{18\varphi^2}\,q_0^{2m}
+O(q_0^{3m})
\right).\;}
\]
\end{theorem}
\begin{proof}
对 \(R=3\) 应用定理 \ref{thm:fold-bin-gauge-constant-stirling-bernoulli-hierarchy}，得到
\[
L_m=\sum_{r=1}^{3}a_r q_0^{(2r-1)m}+O(q_0^{7m}),
\qquad
a_r=
\frac{B_{2r}}{r(2r-1)}
\bigl(\varphi^{-1}+\varphi^{2r-3}\bigr).
\]
对 \(r=1,2,3\) 分别计算系数。由
\[
B_2=\frac16,\qquad B_4=-\frac1{30},\qquad B_6=\frac1{42},
\]
得
\[
a_1=\frac16\bigl(\varphi^{-1}+\varphi^{-1}\bigr)=\frac{1}{3\varphi},
\]
\[
a_2=
\frac{-1/30}{2\cdot 3}\bigl(\varphi^{-1}+\varphi\bigr)
=-\frac{1}{180}\bigl(\varphi+\varphi^{-1}\bigr)
=-\frac{\sqrt5}{180},
\]
其中用了恒等式 \(\varphi+\varphi^{-1}=\sqrt5\)。又
\[
a_3=
\frac{1/42}{3\cdot 5}\bigl(\varphi^{-1}+\varphi^3\bigr).
\]
由
\[
\varphi^3=2\varphi+1,
\qquad
\varphi^{-1}=\varphi-1,
\]
可得
\[
\varphi^{-1}+\varphi^3=3\varphi,
\]
因此
\[
a_3=\frac{3\varphi}{630}=\frac{\varphi}{210}.
\]
代回即得 \(L_m\) 的三项展开。

由 \(C_m=2\pi e^{L_m}\)，立得第二式。再由 \(L_m=O(q_0^m)\) 可作指数展开
\[
e^{L_m}=1+L_m+\frac12L_m^2+O(L_m^3)
=1+\frac{1}{3\varphi}\,q_0^m+\frac{1}{18\varphi^2}\,q_0^{2m}+O(q_0^{3m}),
\]
从而得到第三式。证毕。
\end{proof}

\begin{theorem}[fold-gauge 常数对 \texorpdfstring{$2\pi$}{2pi} 的归一化缺陷最终严格递增，并从下方趋向主系数]\label{thm:xi-time-part59ac-foldgauge-normalized-defect-monotone}
沿用定理 \ref{thm:xi-time-part59ac-foldgauge-three-bernoulli-modes} 的记号，并定义
\[
\widetilde D_m:=q_0^{-m}L_m.
\]
则存在 \(m_0\)，使得对全部 \(m\ge m_0\) 都有
\[
C_m>2\pi,
\qquad
C_{m+1}<C_m,
\qquad
\widetilde D_m<\frac{1}{3\varphi},
\qquad
\widetilde D_{m+1}>\widetilde D_m.
\]
并且
\[
\lim_{m\to\infty}\widetilde D_m=\frac{1}{3\varphi}.
\]
\end{theorem}
\begin{proof}
由定理 \ref{thm:xi-time-part59ac-foldgauge-three-bernoulli-modes}，
\[
L_m
=
\frac{1}{3\varphi}\,q_0^m
-\frac{\sqrt5}{180}\,q_0^{3m}
+O(q_0^{5m}).
\]
首项系数为正，故对充分大的 \(m\) 有 \(L_m>0\)，即
\[
\log C_m>\log(2\pi),
\]
从而 \(C_m>2\pi\)。

再作差分，得到
\[
L_{m+1}-L_m
=
\frac{1}{3\varphi}(q_0-1)\,q_0^m
-\frac{\sqrt5}{180}(q_0^3-1)\,q_0^{3m}
+O(q_0^{5m}).
\]
由于 \(0<q_0<1\)，第一项为负且支配其余项，故对充分大的 \(m\) 有
\[
L_{m+1}<L_m.
\]
指数函数严格递增，遂得 \(C_{m+1}<C_m\)。

对归一化缺陷，直接由上式除以 \(q_0^m\) 得
\[
\widetilde D_m
=
\frac{1}{3\varphi}
-\frac{\sqrt5}{180}\,q_0^{2m}
+O(q_0^{4m}),
\]
从而对充分大的 \(m\) 有
\[
\widetilde D_m<\frac{1}{3\varphi},
\qquad
\lim_{m\to\infty}\widetilde D_m=\frac{1}{3\varphi}.
\]
再作差可得
\[
\widetilde D_{m+1}-\widetilde D_m
=
\frac{\sqrt5}{180}(1-q_0^2)\,q_0^{2m}
+O(q_0^{4m}),
\]
其主项严格为正，故对充分大的 \(m\) 有
\[
\widetilde D_{m+1}>\widetilde D_m.
\]
证毕。
\end{proof}

\begin{theorem}[首项消元后的立方级超收敛上侧包络]\label{thm:xi-time-part59ac-foldgauge-first-mode-cancellation-envelope}
沿用定理 \ref{thm:xi-time-part59ac-foldgauge-three-bernoulli-modes} 的记号，并定义
\[
U_m:=\frac{\log C_{m+1}-q_0\log C_m}{1-q_0}.
\]
则有精确渐近
\[
\boxed{\;
U_m
=
\log(2\pi)
+\frac{\sqrt5}{180}(1+q_0)\,q_0^{3m+1}
+O(q_0^{5m}).\;}
\]
特别地，存在 \(m_1\)，使得对全部 \(m\ge m_1\) 都有
\[
\log(2\pi)<U_m<\log C_m.
\]
因此，\(U_m\) 给出 \(\log(2\pi)\) 的一个严格优于 \(\log C_m\) 的上侧近似，其误差从 \(O(q_0^m)\) 提升到 \(O(q_0^{3m})\)。
\end{theorem}
\begin{proof}
由定义，
\[
U_m-\log(2\pi)
=
\frac{L_{m+1}-q_0L_m}{1-q_0}.
\]
再由定理 \ref{thm:xi-time-part59ac-foldgauge-three-bernoulli-modes}，
\[
L_m
=
a_1q_0^m+a_2q_0^{3m}+O(q_0^{5m}),
\qquad
a_1=\frac{1}{3\varphi},
\qquad
a_2=-\frac{\sqrt5}{180}.
\]
于是
\[
L_{m+1}-q_0L_m
=
a_2q_0^{3m+1}(q_0^2-1)+O(q_0^{5m}),
\]
即首个 Bernoulli 模态被完全消去。两边除以 \(1-q_0\)，并利用
\[
\frac{q_0^2-1}{1-q_0}=-(1+q_0),
\]
得到
\[
U_m-\log(2\pi)
=
-a_2(1+q_0)\,q_0^{3m+1}+O(q_0^{5m})
=
\frac{\sqrt5}{180}(1+q_0)\,q_0^{3m+1}+O(q_0^{5m}).
\]
故对充分大的 \(m\) 有 \(U_m>\log(2\pi)\)。

另一方面，
\[
\log C_m-\log(2\pi)=L_m=\frac{1}{3\varphi}\,q_0^m+O(q_0^{3m}),
\]
其主阶为 \(q_0^m\)，而 \(U_m-\log(2\pi)=O(q_0^{3m})\)。于是对充分大的 \(m\) 有
\[
U_m<\log C_m.
\]
证毕。
\end{proof}

\begin{theorem}[连续比值极限的二次修正与普适有理常数 \texorpdfstring{$1/48$}{1/48}]\label{thm:xi-time-part59ac-foldgauge-ratio-one-forty-eighth}
沿用定理 \ref{thm:xi-time-part59ac-foldgauge-three-bernoulli-modes} 的记号，则有精确渐近
\[
\boxed{\;
\frac{L_{m+1}}{q_0L_m}
=
1+\frac{1}{48}\,q_0^{2m}+O(q_0^{4m}).\;}
\]
因此
\[
\lim_{m\to\infty}\frac{L_{m+1}}{L_m}=q_0=\frac{\varphi}{2},
\]
而且该比值从上方趋向 \(\varphi/2\)；其首个二次修正系数恰为普适有理常数 \(1/48\)。
\end{theorem}
\begin{proof}
由定理 \ref{thm:xi-time-part59ac-foldgauge-three-bernoulli-modes}，
\[
L_m
=
a_1q_0^m+a_2q_0^{3m}+O(q_0^{5m})
=
a_1q_0^m\left(1+\frac{a_2}{a_1}q_0^{2m}+O(q_0^{4m})\right),
\]
其中
\[
a_1=\frac{1}{3\varphi},
\qquad
a_2=-\frac{\sqrt5}{180}.
\]
故
\[
\frac{L_{m+1}}{q_0L_m}
=
\frac{1+\frac{a_2}{a_1}q_0^{2m+2}+O(q_0^{4m})}
{1+\frac{a_2}{a_1}q_0^{2m}+O(q_0^{4m})}
=
1+\frac{a_2}{a_1}(q_0^2-1)\,q_0^{2m}+O(q_0^{4m}).
\]
现化简系数。由
\[
\frac{a_2}{a_1}
=
-\frac{\sqrt5}{180}\cdot 3\varphi
=
-\frac{\sqrt5\,\varphi}{60},
\]
以及
\[
q_0^2-1=\frac{\varphi^2}{4}-1=\frac{\varphi-3}{4},
\]
可得
\[
\frac{a_2}{a_1}(q_0^2-1)
=
\frac{\sqrt5\,\varphi(3-\varphi)}{240}.
\]
再由
\[
\varphi=\frac{1+\sqrt5}{2},
\qquad
3-\varphi=\frac{5-\sqrt5}{2},
\]
有
\[
\sqrt5\,\varphi(3-\varphi)
=
\sqrt5\cdot\frac{(1+\sqrt5)(5-\sqrt5)}{4}
=
5.
\]
因此
\[
\frac{a_2}{a_1}(q_0^2-1)=\frac{5}{240}=\frac{1}{48}.
\]
代回即得所示展开。由于修正项系数为正，故该比值对充分大的 \(m\) 从上方趋向 \(1\)；等价地，
\[
\frac{L_{m+1}}{L_m}
=
q_0\left(1+\frac{1}{48}\,q_0^{2m}+O(q_0^{4m})\right)
\]
从上方趋向 \(q_0=\varphi/2\)。证毕。
\end{proof}

在 Bernoulli 层级与最终单调逼近之外，fold-gauge 常数序列的极限本身还保留了一个不能被代数闭包替代的超越指纹。

\begin{proposition}[圆常数的折叠恢复]\label{prop:xi-time-part59-foldgauge-circle-constant-recovery}
\leanverified{paper\_xi\_time\_part59\_foldgauge\_circle\_constant\_recovery}
沿用推论 \ref{cor:fold-bin-recover-2pi} 的记号，定义
\[
C_m
:=
\exp\!\Biggl(
\frac{2^{m+1}}{|X_m|}\bigl(g_m-\kappa_m+1\bigr)
-\frac{1}{|X_m|}\sum_{w\in X_m}\log d_m^{\mathrm{bin}}(w)
\Biggr).
\]
则
\[
\lim_{m\to\infty}C_m=2\pi.
\]
因此，这一完全由离散折叠数据构成的常数列在极限上强制恢复出圆常数；尤其，任何要求把该极限常数代数化的替代描述都不可能与原序列的极限严格一致。
\end{proposition}
\begin{proof}
推论 \ref{cor:fold-bin-recover-2pi} 已给出上述极限。另一方面，\(\pi\) 的超越性是经典事实，故 \(2\pi\) 亦为超越数。于是，若某种替代描述宣称与 \((C_m)\) 具有同一极限而该极限仍为代数常数，便与
\[
\lim_{m\to\infty}C_m=2\pi
\]
矛盾。证毕。
\end{proof}


\noindent 上述结果表明：一方面，fold-gauge 常数的 Stirling--Bernoulli 层级不仅允许逐阶恢复全部偶 Bernoulli 数、构造任意预先指定阶数的显式消模器，而且其前三个非平凡 Bernoulli 模态、归一化缺陷的最终单调趋近、首项消元后的立方级上侧包络以及相邻比值的普适 \(\frac{1}{48}\) 二次修正也都可完全显式写出；原始误差序列 \(\log C_m-\log(2\pi)\) 本身则在充分大 \(m\) 上呈现任意固定阶的最终完全单调性，其极限更由纯离散折叠数据内生恢复出超越常数 \(2\pi\)。另一方面，与定理 \ref{thm:xi-output-potential-dirichlet-fixed-j-fourth-order}、推论 \ref{cor:xi-output-potential-dirichlet-bounded-window-j1} 及定理 \ref{thm:xi-output-potential-dirichlet-fourth-order-gain} 所建立的单位根扭曲四阶展开链合并后，prime-shadow 侧的解析结构同样形成闭合：六倍窗零点集合的半维指数率被精确固定，而局部累积量则把最坏 Dirichlet 扭曲的局部模态、窗口极值机制与高阶近似收益统一到同一输出位势接口之下。于是相关结论不再只是审计意义上的现象陈列，而具备了可独立组织为主定理链的解析骨架。

\paragraph{推前熵账本、Rényi 散度冻结、临界共振底噪与 window-\texorpdfstring{$6$}{6} 编码冗余}
折叠推前分布的 Shannon/Rényi 散度账本、规范体积密度的 Stirling 差常数、临界共振常数强制的二阶底噪，以及 window-\(6\) 的最小完备奇偶荷签名之间，还存在下列可由既有结论直接闭合的关系。

\begin{theorem}[折叠推前熵与规范体积密度的单位 nat 差额]\label{thm:xi-foldbin-entropy-gauge-one-nat-gap}
沿用定理 \ref{thm:xi-gauge-volume-boltzmann-stirling-exponential} 与
\ref{thm:xi-foldbin-gauge-escort-kl-decomposition} 的记号。记
\[
p_m^{\mathrm{bin}}(w):=\frac{d_m^{\mathrm{bin}}(w)}{2^m},
\qquad
u_m(w):=\frac{1}{|X_m|}
\qquad (w\in X_m),
\]
并以自然对数定义 Shannon 熵
\[
H(p):=-\sum_{w\in X_m}p(w)\log p(w).
\]
则有严格恒等式
\[
\boxed{\;
H\!\bigl(p_m^{\mathrm{bin}}\bigr)=m\log 2-\kappa_m,\;}
\]
\[
\boxed{\;
D_{\mathrm{KL}}\!\bigl(p_m^{\mathrm{bin}}\|u_m\bigr)
=
\kappa_m+\log|X_m|-m\log 2.\;}
\]
从而
\[
\boxed{\;
H\!\bigl(p_m^{\mathrm{bin}}\bigr)+g_m
=
m\log 2-1+O\!\Bigl(m\bigl(\tfrac{\varphi}{2}\bigr)^m\Bigr),\;}
\]
\[
\boxed{\;
D_{\mathrm{KL}}\!\bigl(p_m^{\mathrm{bin}}\|u_m\bigr)
=
g_m+1+\log F_{m+2}-m\log 2
+O\!\Bigl(m\bigl(\tfrac{\varphi}{2}\bigr)^m\Bigr).\;}
\]
特别地，折叠推前熵与规范体积密度之和，相对满 \(m\)-比特熵 \(m\log 2\) 仅差一个普适常数 \(1\)。
\end{theorem}
\begin{proof}
记 \(d_w:=d_m^{\mathrm{bin}}(w)\) 与 \(p_w:=p_m^{\mathrm{bin}}(w)\)。由 \(\sum_{w\in X_m}d_w=2^m\) 得
\[
\begin{aligned}
H\!\bigl(p_m^{\mathrm{bin}}\bigr)
&=-\sum_{w\in X_m}\frac{d_w}{2^m}\log\frac{d_w}{2^m}\\
&=-2^{-m}\sum_{w\in X_m}d_w\log d_w
+2^{-m}\sum_{w\in X_m}d_w\,m\log 2\\
&=m\log 2-\kappa_m,
\end{aligned}
\]
即得第一式。

再由 \(u_m(w)=1/|X_m|\) 可得
\[
\begin{aligned}
D_{\mathrm{KL}}\!\bigl(p_m^{\mathrm{bin}}\|u_m\bigr)
&=\sum_{w\in X_m}p_w\log\frac{p_w}{u_m(w)}\\
&=-H\!\bigl(p_m^{\mathrm{bin}}\bigr)+\log|X_m|\\
&=\kappa_m+\log|X_m|-m\log 2.
\end{aligned}
\]
最后，由定理 \ref{thm:xi-gauge-volume-boltzmann-stirling-exponential}，
\[
g_m=\kappa_m-1+O\!\Bigl(m\bigl(\tfrac{\varphi}{2}\bigr)^m\Bigr),
\]
并且 \(|X_m|=F_{m+2}\)。将其代入前两式即得后两式。证毕。
\end{proof}

\begin{theorem}[相对均匀基线的 Rényi 散度公式与冻结相仿射律]\label{thm:xi-foldbin-renyi-divergence-freezing-affine}
沿用定理 \ref{thm:xi-foldbin-entropy-gauge-one-nat-gap} 的记号。对 \(q>0\)、\(q\neq 1\)，定义 Rényi 散度
\[
D_q\!\bigl(p_m^{\mathrm{bin}}\|u_m\bigr)
:=
\frac{1}{q-1}\log\sum_{w\in X_m}
\bigl(p_m^{\mathrm{bin}}(w)\bigr)^q\,u_m(w)^{1-q}.
\]
再定义
\[
S_q(m):=\sum_{w\in X_m}\bigl(d_m^{\mathrm{bin}}(w)\bigr)^q.
\]
则对任意 \(q>0\)、\(q\neq 1\)，都有精确公式
\[
\boxed{\;
D_q\!\bigl(p_m^{\mathrm{bin}}\|u_m\bigr)
=
\log|X_m|+\frac{1}{q-1}\log\frac{S_q(m)}{2^{mq}}.\;}
\]

进一步，对整数 \(q\ge 2\)，若压力极限
\[
P_q:=\lim_{m\to\infty}\frac{1}{m}\log S_q(m)
\]
存在，则
\[
\boxed{\;
\lim_{m\to\infty}\frac{1}{m}
D_q\!\bigl(p_m^{\mathrm{bin}}\|u_m\bigr)
=
\log\varphi+\frac{P_q-q\log 2}{q-1}.\;}
\]

若再假设严格基态间隙 \(\alpha_*^{(2)}<\alpha_*\)，并记
\[
a_{\mathrm c}:=\frac{\log\varphi-g_*}{\alpha_*-\alpha_*^{(2)}},
\]
则对每个整数 \(q>a_{\mathrm c}\)，散度密度
\[
\delta(q):=
\lim_{m\to\infty}\frac{1}{m}
D_q\!\bigl(p_m^{\mathrm{bin}}\|u_m\bigr)
\]
满足显式仿射公式
\[
\boxed{\;
\delta(q)
=
\log\varphi+\frac{q(\alpha_*-\log 2)+g_*}{q-1}.\;}
\]
特别地，
\[
\boxed{\;
\lim_{q\to\infty}\delta(q)=\alpha_*-\log\!\Bigl(\frac{2}{\varphi}\Bigr).\;}
\]
\end{theorem}
\begin{proof}
由 \(u_m(w)=1/|X_m|\)，直接计算得
\[
\sum_{w\in X_m}
\bigl(p_m^{\mathrm{bin}}(w)\bigr)^q\,u_m(w)^{1-q}
=
|X_m|^{q-1}\sum_{w\in X_m}\bigl(p_m^{\mathrm{bin}}(w)\bigr)^q.
\]
而
\[
\sum_{w\in X_m}\bigl(p_m^{\mathrm{bin}}(w)\bigr)^q
=
\sum_{w\in X_m}\left(\frac{d_m^{\mathrm{bin}}(w)}{2^m}\right)^q
=
\frac{S_q(m)}{2^{mq}}.
\]
代回定义即得第一式。

若 \(q\ge 2\) 为整数且 \(P_q\) 存在，则两边除以 \(m\) 并令 \(m\to\infty\)，再用
\[
\lim_{m\to\infty}\frac{1}{m}\log|X_m|
=
\lim_{m\to\infty}\frac{1}{m}\log F_{m+2}
=
\log\varphi,
\]
便得到第二式。

最后，在严格基态间隙下，由定理 \ref{thm:fold-pressure-freezing-threshold}，对每个整数 \(q>a_{\mathrm c}\) 都有
\[
P_q=q\alpha_*+g_*.
\]
将其代入第二式即可得到 \(\delta(q)\) 的仿射公式。再令 \(q\to\infty\)，便得
\[
\delta(q)\to \log\varphi+\alpha_*-\log 2
=
\alpha_*-\log\!\Bigl(\frac{2}{\varphi}\Bigr).
\]
证毕。
\end{proof}

\begin{theorem}[临界共振常数强制 Rényi-\texorpdfstring{$2$}{2} 散度的严格正底噪]\label{thm:xi-foldbin-critical-resonance-renyi2-floor}
沿用定义 \ref{def:fold-normalized-amplitude} 与定理 \ref{thm:fold-collision-spectrum} 的记号，记
\[
\mathsf{Col}^{\mathrm{bin}}_m
:=
\sum_{w\in X_m}\bigl(p_m^{\mathrm{bin}}(w)\bigr)^2
=
\frac{1}{F_{m+2}}\sum_{k}a_m(k)^2.
\]
则 Rényi-\(2\) 散度满足严格恒等式
\[
\boxed{\;
D_2\!\bigl(p_m^{\mathrm{bin}}\|u_m\bigr)
=
\log\!\bigl(F_{m+2}\,\mathsf{Col}^{\mathrm{bin}}_m\bigr).\;}
\]
并且对每个 \(m\) 都有点态下界
\[
\boxed{\;
D_2\!\bigl(p_m^{\mathrm{bin}}\|u_m\bigr)
\ge
\log\!\bigl(1+a_m(F_m)^2+a_m(F_{m+1})^2\bigr).\;}
\]
因此，若 \(C_\varphi>0\) 为定理 \ref{thm:fold-critical-resonance-constant} 的临界共振常数，则
\[
\boxed{\;
\liminf_{m\to\infty}
D_2\!\bigl(p_m^{\mathrm{bin}}\|u_m\bigr)
\ge
\log(1+2C_\varphi^2)
\approx 0.09597061849768761.\;}
\]

等价地，若定义 Rényi-\(2\) 有效支撑
\[
N_{2,m}^{\mathrm{eff}}
:=
\exp\!\Bigl(H_2\!\bigl(p_m^{\mathrm{bin}}\bigr)\Bigr)
=
\frac{1}{\mathsf{Col}^{\mathrm{bin}}_m},
\qquad
H_2(p):=-\log\sum_w p(w)^2,
\]
则
\[
\boxed{\;
N_{2,m}^{\mathrm{eff}}
\le
\frac{F_{m+2}}{1+a_m(F_m)^2+a_m(F_{m+1})^2},\;}
\]
并且
\[
\boxed{\;
\limsup_{m\to\infty}\frac{N_{2,m}^{\mathrm{eff}}}{F_{m+2}}
\le
\frac{1}{1+2C_\varphi^2}
\approx 0.908490708498425.\;}
\]
\end{theorem}
\begin{proof}
定理 \ref{thm:xi-foldbin-renyi-divergence-freezing-affine} 在 \(q=2\) 时给出
\[
D_2\!\bigl(p_m^{\mathrm{bin}}\|u_m\bigr)
=
\log|X_m|+\log\frac{S_2(m)}{2^{2m}}.
\]
又由
\[
S_2(m)=\sum_{w\in X_m}\bigl(d_m^{\mathrm{bin}}(w)\bigr)^2
=
2^{2m}\,\mathsf{Col}^{\mathrm{bin}}_m,
\qquad
|X_m|=F_{m+2},
\]
立得第一条恒等式。

另一方面，由 Parseval 恒等式
\[
F_{m+2}\,\mathsf{Col}^{\mathrm{bin}}_m=\sum_k a_m(k)^2,
\]
并注意 \(a_m(0)=1\)，故
\[
F_{m+2}\,\mathsf{Col}^{\mathrm{bin}}_m
\ge
1+a_m(F_m)^2+a_m(F_{m+1})^2.
\]
取对数即得点态下界。

再由定理 \ref{thm:fold-critical-resonance-constant}，
\[
a_m(F_m)\to C_\varphi,
\qquad
a_m(F_{m+1})\to C_\varphi,
\]
令 \(m\to\infty\) 即得 Rényi-\(2\) 散度的正底噪下界。

最后，由 \(H_2(p_m^{\mathrm{bin}})=-\log\mathsf{Col}^{\mathrm{bin}}_m\) 可知
\[
N_{2,m}^{\mathrm{eff}}=\frac{1}{\mathsf{Col}^{\mathrm{bin}}_m}.
\]
将上式与
\(
F_{m+2}\mathsf{Col}^{\mathrm{bin}}_m
\ge
1+a_m(F_m)^2+a_m(F_{m+1})^2
\)
重排即可得到有效支撑上界；再令 \(m\to\infty\) 便得最后一式。证毕。
\end{proof}

\begin{theorem}[Rényi--2 投影熵的常数缺口]\label{thm:xi-foldbin-renyi2-projection-entropy-constant-gap}
沿用定理 \ref{thm:xi-foldbin-critical-resonance-renyi2-floor} 的记号，定义二阶投影熵
\[
H_2^{\mathrm{proj}}(m):=-\log \mathsf{Col}^{\mathrm{bin}}_m,
\]
并记
\[
\delta_2:=\log(1+2C_\varphi^2)>0.
\]
则
\[
H_2^{\mathrm{proj}}(m)
\le
\log F_{m+2}-\delta_2+o(1).
\]
数值上，
\[
\delta_2\approx 0.09597061849768761\ \text{nats}
\approx 0.138456335393854\ \text{bits}.
\]
换言之，相对于 \(|X_m|=F_{m+2}\) 的理想均匀基线，bin-fold 投影在 Rényi--\(2\) 意义下始终保留一个不可消除的常数缺口。
\end{theorem}
\begin{proof}
由定理 \ref{thm:xi-foldbin-critical-resonance-renyi2-floor}，
\[
D_2\!\bigl(p_m^{\mathrm{bin}}\|u_m\bigr)
=
\log\!\bigl(F_{m+2}\mathsf{Col}^{\mathrm{bin}}_m\bigr).
\]
于是
\[
H_2^{\mathrm{proj}}(m)
=-\log \mathsf{Col}^{\mathrm{bin}}_m
=
\log F_{m+2}
-D_2\!\bigl(p_m^{\mathrm{bin}}\|u_m\bigr).
\]
同一定理又给出
\[
\liminf_{m\to\infty}D_2\!\bigl(p_m^{\mathrm{bin}}\|u_m\bigr)\ge \delta_2,
\]
故
\[
H_2^{\mathrm{proj}}(m)\le \log F_{m+2}-\delta_2+o(1).
\]
比特值只需再除以 \(\log 2\) 即得。证毕。
\end{proof}

\begin{theorem}[临界共振常数强制 bin-fold 推前分布在 Rényi-\texorpdfstring{$2$}{2} 几何中永久远离均匀层]\label{thm:xi-foldbin-critical-resonance-chi2-renyi2-separation}
沿用定理 \ref{thm:xi-foldbin-critical-resonance-renyi2-floor} 的记号。则有精确恒等式
\[
\boxed{\;
\chi^2\!\bigl(p_m^{\mathrm{bin}}\|u_m\bigr)
=
F_{m+2}\left(\mathsf{Col}^{\mathrm{bin}}_m-\frac{1}{F_{m+2}}\right).\;}
\]
进一步，
\[
\boxed{\;
\liminf_{m\to\infty}\chi^2\!\bigl(p_m^{\mathrm{bin}}\|u_m\bigr)
\ge
2C_\varphi^2
\approx 0.1007267225141179.\;}
\]
等价地，
\[
\boxed{\;
\liminf_{m\to\infty}D_2\!\bigl(p_m^{\mathrm{bin}}\|u_m\bigr)
\ge
\log(1+2C_\varphi^2)
\approx 0.09597061849768761.\;}
\]
因此，bin-fold 推前分布在 \(\chi^2\) 与 Rényi-\(2\) 意义下都不可能渐近均匀化。
\end{theorem}
\begin{proof}
第一条恒等式正是命题 \ref{prop:fold-bin-chi2-col} 的内容。另一方面，由定理 \ref{thm:xi-foldbin-critical-resonance-renyi2-floor} 的点态下界，
\[
F_{m+2}\mathsf{Col}^{\mathrm{bin}}_m
\ge
1+a_m(F_m)^2+a_m(F_{m+1})^2,
\]
故
\[
\chi^2\!\bigl(p_m^{\mathrm{bin}}\|u_m\bigr)
=
F_{m+2}\mathsf{Col}^{\mathrm{bin}}_m-1
\ge
a_m(F_m)^2+a_m(F_{m+1})^2.
\]
再由定理 \ref{thm:fold-critical-resonance-constant}，
\[
a_m(F_m)\to C_\varphi,
\qquad
a_m(F_{m+1})\to C_\varphi,
\]
便得到
\[
\liminf_{m\to\infty}\chi^2\!\bigl(p_m^{\mathrm{bin}}\|u_m\bigr)\ge 2C_\varphi^2.
\]
最后，由
\[
D_2\!\bigl(p_m^{\mathrm{bin}}\|u_m\bigr)
=
\log\!\bigl(F_{m+2}\mathsf{Col}^{\mathrm{bin}}_m\bigr)
=
\log\!\Bigl(1+\chi^2\!\bigl(p_m^{\mathrm{bin}}\|u_m\bigr)\Bigr),
\]
取下极限即得 Rényi-\(2\) 版本。证毕。
\end{proof}

\leanverified{paper\_xi\_window6\_renyi\_divergence\_parity\_charge\_redundancy}
\begin{proposition}[window-\texorpdfstring{$6$}{6} 的 Rényi 散度闭式与最小完备奇偶荷签名的熵冗余]\label{prop:xi-window6-renyi-divergence-parity-charge-redundancy}
取 \(m=6\)。由推论 \ref{cor:fold-bin-m6-fiber-hist}，
\[
\#\{w\in X_6:\ d_6^{\mathrm{bin}}(w)=2\}=8,\qquad
\#\{w\in X_6:\ d_6^{\mathrm{bin}}(w)=3\}=4,\qquad
\#\{w\in X_6:\ d_6^{\mathrm{bin}}(w)=4\}=9,
\]
且 \(|X_6|=21\)。因此
\[
p_6^{\mathrm{bin}}
=
\underbrace{\tfrac{1}{32},\dots,\tfrac{1}{32}}_{8\text{ 次}},
\underbrace{\tfrac{3}{64},\dots,\tfrac{3}{64}}_{4\text{ 次}},
\underbrace{\tfrac{1}{16},\dots,\tfrac{1}{16}}_{9\text{ 次}}.
\]
并记 \(u_6(w)\equiv 1/21\) 为 \(X_6\) 上的均匀分布。
对任意 \(q>0\)、\(q\neq 1\)，Rényi 散度曲线具有闭式
\[
\boxed{\;
D_q\!\bigl(p_6^{\mathrm{bin}}\|u_6\bigr)
=
\log 21+\frac{1}{q-1}
\log\frac{8\cdot 2^q+4\cdot 3^q+9\cdot 4^q}{64^q}.\;}
\]

特别地，Shannon 熵与 KL 散度为
\[
\boxed{\;
H\!\bigl(p_6^{\mathrm{bin}}\bigr)
=
\frac{37}{8}\log 2-\frac{3}{16}\log 3
\approx 2.9998159059644762,\;}
\]
\[
\boxed{\;
D_{\mathrm{KL}}\!\bigl(p_6^{\mathrm{bin}}\|u_6\bigr)
=
\log 21-\frac{37}{8}\log 2+\frac{3}{16}\log 3
\approx 0.04470653175894679.\;}
\]

若把 Shannon 熵改以 bit 为单位写成
\[
H_{\mathrm{bit}}\!\bigl(p_6^{\mathrm{bin}}\bigr)
:=
\frac{H(p_6^{\mathrm{bin}})}{\log 2},
\]
则
\[
\boxed{\;
H_{\mathrm{bit}}\!\bigl(p_6^{\mathrm{bin}}\bigr)
=
\frac{37}{8}-\frac{3}{16}\log_2 3
\approx 4.327819531114783.\;}
\]

另一方面，由定理 \ref{thm:xi-foldbin-parity-charge-split-exact-minrank}，window-\(6\) 的最小完备奇偶荷签名秩恰为 \(21\)。因此其相对于输出 Shannon 熵的不可压缩冗余为
\[
\boxed{\;
\operatorname{Red}_6
:=
21-H_{\mathrm{bit}}\!\bigl(p_6^{\mathrm{bin}}\bigr)
=
\frac{131}{8}+\frac{3}{16}\log_2 3
\approx 16.672180468885216\ \mathrm{bits}.\;}
\]
\end{proposition}
\begin{proof}
由 window-\(6\) 直方图，直接有
\[
S_q(6)=8\cdot 2^q+4\cdot 3^q+9\cdot 4^q
\qquad(q>0).
\]
再由定理 \ref{thm:xi-foldbin-renyi-divergence-freezing-affine} 的精确公式，
\[
D_q\!\bigl(p_6^{\mathrm{bin}}\|u_6\bigr)
=
\log|X_6|+\frac{1}{q-1}\log\frac{S_q(6)}{2^{6q}},
\]
而 \(|X_6|=21\)，故得第一条闭式。

Shannon 熵直接按三类纤维质量求和：
\[
\begin{aligned}
H\!\bigl(p_6^{\mathrm{bin}}\bigr)
&=
8\cdot \frac{1}{32}\log 32
+4\cdot \frac{3}{64}\log \frac{64}{3}
+9\cdot \frac{1}{16}\log 16\\
&=
\frac{37}{8}\log 2-\frac{3}{16}\log 3.
\end{aligned}
\]
于是由定理 \ref{thm:xi-foldbin-entropy-gauge-one-nat-gap} 的
\[
D_{\mathrm{KL}}\!\bigl(p_6^{\mathrm{bin}}\|u_6\bigr)
=
\log|X_6|-H\!\bigl(p_6^{\mathrm{bin}}\bigr)
\]
即可得到 KL 散度的闭式。再除以 \(\log 2\) 即得 \(H_{\mathrm{bit}}(p_6^{\mathrm{bin}})\)。

最后，由定理 \ref{thm:xi-foldbin-parity-charge-split-exact-minrank}，最小完备奇偶荷签名秩为 \(21\)，故
\[
\operatorname{Red}_6
=
21-\frac{H(p_6^{\mathrm{bin}})}{\log 2}
=
21-\left(\frac{37}{8}-\frac{3}{16}\log_2 3\right),
\]
整理即得所示公式。证毕。
\end{proof}

\paragraph{审计偶数窗上的中心坍缩、饱和判据与熵--体积--碰撞恒等链}
本段把附录中已建立的 bin-fold 退化谱、规范群与 Stirling 展开结果，回收为六条可直接调用的终端结论：它们分别刻画审计偶数窗上的中心坍缩与 Abel 满秩分离、最大纤维组合上界的精确饱和判据、任意纤维多重度的普适 Fibonacci 上界及其末位强化、\(\kappa_m\)--KL--碰撞概率之间的恒等链、规范群体积的三因子乘法分解，以及最小退化扇区单独强制的显式碰撞下界。

\begin{theorem}[审计偶数窗上中心坍缩与 Abel 满秩的严格分离]\label{thm:xi-foldbin-audited-even-center-collapse-abel-full-rank}
令
\[
G_m^{\mathrm{bin}}
:=
\Aut(\Fold_m^{\mathrm{bin}})
\cong
\prod_{w\in X_m}S_{d_m^{\mathrm{bin}}(w)}.
\]
若
\(
m\in\{6,8,10,12\}
\)，则
\[
\boxed{\;
\bigl(G_m^{\mathrm{bin}}\bigr)^{\mathrm{ab}}
\cong
(\ZZ/2\ZZ)^{|X_m|}.\;}
\]
更强地，
\[
\boxed{\;
Z\!\bigl(G_6^{\mathrm{bin}}\bigr)\cong (\ZZ/2\ZZ)^8,\qquad
Z\!\bigl(G_m^{\mathrm{bin}}\bigr)=1\ \ (m=8,10,12).\;}
\]
并且在 \(m=6\) 时，boundary sheet-parity 规范子群
\[
Z_{\partial}\cong (\ZZ/2\ZZ)^3\hookrightarrow Z\!\bigl(G_6^{\mathrm{bin}}\bigr)
\]
只是整个中心中的一个规范三维坐标子格；事实上
\[
\boxed{\;
Z\!\bigl(G_6^{\mathrm{bin}}\bigr)/Z_{\partial}\cong (\ZZ/2\ZZ)^5.\;}
\]
\leanverified{paper\_xi\_foldbin\_audited\_even\_center\_collapse\_abel\_full\_rank}
\end{theorem}
\begin{proof}
由定理 \ref{thm:xi-audited-even-abel-center-fibonacci-splitting}，
\[
\bigl(G_m^{\mathrm{bin}}\bigr)^{\mathrm{ab}}
\cong
(\ZZ/2\ZZ)^{F_m}\oplus(\ZZ/2\ZZ)^{D_{2m-2}}.
\]
而
\(
F_m+D_{2m-2}=F_m+F_{m+1}=F_{m+2}=|X_m|
\)，故得到
\[
\bigl(G_m^{\mathrm{bin}}\bigr)^{\mathrm{ab}}\cong (\ZZ/2\ZZ)^{|X_m|}.
\]

另一方面，定理 \ref{thm:xi-foldbin-center-solvable-fiber-criterion} 给出
\[
Z\!\bigl(G_m^{\mathrm{bin}}\bigr)\cong (\ZZ/2\ZZ)^{B_m},
\qquad
B_m:=\#\{w\in X_m:\ d_m^{\mathrm{bin}}(w)=2\}.
\]
由表 \ref{tab:fold_bin_degeneracy_spectrum_m6_m12} 的直方图，
\[
m=6:\ 2:8,\,3:4,\,4:9,
\]
故 \(B_6=8\)；而
\[
m=8:\ 3:21,\,5:11,\,6:23,\qquad
m=10:\ 5:55,\,8:52,\,9:37,\qquad
m=12:\ 8:144,\,12:85,\,13:148,
\]
均无 \(d=2\) 纤维，因此 \(B_8=B_{10}=B_{12}=0\)。从而
\[
Z\!\bigl(G_6^{\mathrm{bin}}\bigr)\cong (\ZZ/2\ZZ)^8,
\qquad
Z\!\bigl(G_m^{\mathrm{bin}}\bigr)=1\ \ (m=8,10,12).
\]

最后，由定理 \ref{thm:xi-window6-central-noncentral-anomaly-splitting}，
\(
Z_{\partial}\cong (\ZZ/2\ZZ)^3
\)
是
\(
Z(G_6^{\mathrm{bin}})
\)
的规范直和因子，并且
\[
Z\!\bigl(G_6^{\mathrm{bin}}\bigr)/Z_{\partial}\cong (\ZZ/2\ZZ)^5.
\]
证毕。
\end{proof}

\begin{theorem}[bin-fold 最大纤维的精确饱和判据与最终严格次 Fibonacci 律]\label{thm:xi-foldbin-maxfiber-saturation-final-subfibonacci}
\leanverified{paper\_xi\_foldbin\_maxfiber\_saturation\_final\_subfibonacci}
令 \(K=K(m)\) 为满足
\[
F_{K+1}\le 2^m-1<F_{K+2}
\]
的唯一整数，并置 \(L:=K-m\)。
定义合法尾部集合
\[
\mathcal T_m
:=
\Bigl\{
(c_{m+1},\dots,c_K)\in\{0,1\}^{L}:\ c_kc_{k+1}=0
\Bigr\},
\]
以及最大尾和值
\[
S_{\max}(m)
:=
\max\Bigl\{
\sum_{k=m+1}^{K}c_kF_{k+1}:\ (c_{m+1},\dots,c_K)\in\mathcal T_m
\Bigr\}.
\]
则最大纤维满足
\[
\boxed{\;
d_{\max}^{\mathrm{bin}}(m)
:=
\max_{w\in X_m}d_m^{\mathrm{bin}}(w)
=
d_m^{\mathrm{bin}}(0^m)
=
\#\Bigl\{
c\in\mathcal T_m:\ \sum_{k=m+1}^{K}c_kF_{k+1}\le 2^m-1
\Bigr\}.\;}
\]
并且，记组合上界 \(F_{L+2}=|\mathcal T_m|\)，则
\[
\boxed{\;
d_{\max}^{\mathrm{bin}}(m)=F_{L+2}
\iff
S_{\max}(m)\le 2^m-1.\;}
\]
进一步，存在 \(m_0\) 使得对所有 \(m\ge m_0\) 都有
\[
\boxed{\;
d_{\max}^{\mathrm{bin}}(m)\le F_{L+2}-1.\;}
\]
\end{theorem}
\begin{proof}
由命题 \ref{prop:fold-bin-digit-dp}，对任意 \(w\in X_m\)，
\(
d_m^{\mathrm{bin}}(w)
\)
恰等于满足相邻约束、边界约束与阈值约束的尾部位串数目。
取 \(w=0^m\) 时，\(V_m(w)=0\) 且 \(w_m=0\)，因而边界约束消失，得到
\[
d_m^{\mathrm{bin}}(0^m)
=
\#\Bigl\{
c\in\mathcal T_m:\ \sum_{k=m+1}^{K}c_kF_{k+1}\le 2^m-1
\Bigr\}.
\]
再由推论 \ref{cor:fold-bin-saturation}，最大纤维确由 \(0^m\) 取得，故第一条公式成立。

同一推论又给出
\[
d_{\max}^{\mathrm{bin}}(m)\le F_{L+2}=|\mathcal T_m|,
\]
且该上界饱和当且仅当阈值约束对全部合法尾部都冗余，也即
\[
2^m-1\ge S_{\max}(m).
\]
这就证明了第二条等价判据。

最后，由注 \ref{rem:fold-bin-finite}，上述组合上界的饱和只能发生有限次。
因此存在 \(m_0\) 使得对一切 \(m\ge m_0\)，恒有
\(
d_{\max}^{\mathrm{bin}}(m)<F_{L+2}
\)；
由于两边都是整数，便得到
\[
d_{\max}^{\mathrm{bin}}(m)\le F_{L+2}-1.
\]
证毕。
\end{proof}

\begin{conjecture}[bin-fold 最大纤维组合上界的精确饱和集是有限孤立集]\label{conj:xi-foldbin-maxfiber-saturation-sporadic-set}
沿用定理 \ref{thm:xi-foldbin-maxfiber-saturation-final-subfibonacci} 的记号。生成表
\ref{tab:fold_binary_saturation_points} 对扫描窗口 \(m\le 400\) 给出的审计结果表明，组合上界
\[
d_{\max}^{\mathrm{bin}}(m)\le F_{L+2}
\]
的饱和点恰为
\[
\{1,2,3,4,5,7,12\}.
\]
据此提出猜想
\[
\boxed{\;
d_{\max}^{\mathrm{bin}}(m)=F_{L+2}
\iff
m\in\{1,2,3,4,5,7,12\}.\;}
\]
等价地，
\[
\boxed{\;
2^m-1\ge S_{\max}(m)
\iff
m\in\{1,2,3,4,5,7,12\}.\;}
\]
若此猜想成立，则定理 \ref{thm:xi-foldbin-maxfiber-saturation-final-subfibonacci} 中“饱和只能发生有限次”的结论可进一步精化为一个完全显式的有限异常集判定。
\end{conjecture}

\begin{theorem}[bin-fold 纤维多重度的普适 Fibonacci 上界与末位强化]\label{thm:xi-foldbin-universal-fibonacci-upper-bound-lastbit-improvement}
\leanverified{paper\_xi\_foldbin\_universal\_fibonacci\_upper\_bound\_lastbit\_improvement}
沿用定理 \ref{thm:xi-foldbin-maxfiber-saturation-final-subfibonacci} 的记号，记
\[
L:=K(m)-m.
\]
则对任意 \(w\in X_m\)，都有
\[
\boxed{\;
d_m^{\mathrm{bin}}(w)\le F_{L+2}=F_{K(m)-m+2}.\;}
\]
若进一步满足 \(w_m=1\)，则更强地有
\[
\boxed{\;
d_m^{\mathrm{bin}}(w)\le F_{L+1}=F_{K(m)-m+1}.\;}
\]
\end{theorem}
\begin{proof}
由命题 \ref{prop:fold-bin-digit-dp}，对任意 \(w\in X_m\)，
\(
d_m^{\mathrm{bin}}(w)
\)
恰等于满足
\[
c_kc_{k+1}=0,\qquad
w_mc_{m+1}=0,\qquad
V_m(w)+\sum_{k=m+1}^{K}c_kF_{k+1}\le 2^m-1
\]
的尾部位串 \((c_{m+1},\dots,c_K)\) 的数目。

先忽略阈值约束。长度 \(L\) 的二进制串中不含相邻 \(1\) 的总数是标准 Fibonacci 计数 \(F_{L+2}\)，故
\[
d_m^{\mathrm{bin}}(w)\le F_{L+2}.
\]

若 \(w_m=1\)，则边界约束强制 \(c_{m+1}=0\)。当 \(L=0\) 时尾部为空串，结论平凡成立；当 \(L\ge 1\) 时，剩余自由尾部是长度 \(L-1\) 的无相邻 \(1\) 串，其数目为
\[
F_{(L-1)+2}=F_{L+1}.
\]
再加入阈值约束只会进一步减少可行尾部，故
\[
d_m^{\mathrm{bin}}(w)\le F_{L+1}.
\]
证毕。
\end{proof}

\begin{theorem}[\texorpdfstring{$\kappa_m$}{kappa_m}、KL 散度与碰撞概率的精确恒等链]\label{thm:xi-foldbin-kappa-kl-collision-identity-chain}
\leanverified{paper\_xi\_foldbin\_kappa\_kl\_collision\_identity\_chain}
令
\[
p_m^{\mathrm{bin}}(w):=\frac{d_m^{\mathrm{bin}}(w)}{2^m},
\qquad
u_m(w):=\frac{1}{|X_m|}
\qquad (w\in X_m).
\]
则有严格恒等式
\[
\boxed{\;
D_{\mathrm{KL}}\!\bigl(p_m^{\mathrm{bin}}\|u_m\bigr)
=
\kappa_m-\log\!\Bigl(\frac{2^m}{|X_m|}\Bigr).\;}
\]
因而
\[
\boxed{\;
\log\!\Bigl(\frac{2^m}{|X_m|}\Bigr)
\le
\kappa_m
\le
\log\!\bigl(2^m\mathsf{Col}^{\mathrm{bin}}_m\bigr).\;}
\]
又由
\[
\chi^2\!\bigl(p_m^{\mathrm{bin}}\|u_m\bigr)
=
|X_m|\,\mathsf{Col}^{\mathrm{bin}}_m-1,
\]
可把上界改写为
\[
\kappa_m
\le
\log\!\Bigl(\frac{2^m}{|X_m|}\Bigr)
\;+\;
\log\!\Bigl(1+\chi^2\!\bigl(p_m^{\mathrm{bin}}\|u_m\bigr)\Bigr).
\]
并且两端取等都当且仅当 \(p_m^{\mathrm{bin}}=u_m\)。
\end{theorem}
\begin{proof}
第一条恒等式正是定理 \ref{thm:xi-foldbin-entropy-gauge-one-nat-gap} 的
\[
D_{\mathrm{KL}}\!\bigl(p_m^{\mathrm{bin}}\|u_m\bigr)
=
\kappa_m+\log|X_m|-m\log 2
\]
的改写。
由于 \(D_{\mathrm{KL}}\ge 0\)，立得
\[
\kappa_m\ge \log\!\Bigl(\frac{2^m}{|X_m|}\Bigr),
\]
且等号成立当且仅当 \(p_m^{\mathrm{bin}}=u_m\)。

另一方面，由定理 \ref{thm:xi-foldbin-gauge-escort-kl-decomposition}，
\[
\kappa_m
=
\log\!\bigl(2^m\mathsf{Col}^{\mathrm{bin}}_m\bigr)
-D_{\mathrm{KL}}\!\bigl(p_m^{\mathrm{bin}}\|\widetilde p_m\bigr),
\]
其中 \(\widetilde p_m(w)=(p_m^{\mathrm{bin}}(w))^2/\mathsf{Col}^{\mathrm{bin}}_m\)。
于是
\[
\kappa_m\le \log\!\bigl(2^m\mathsf{Col}^{\mathrm{bin}}_m\bigr),
\]
且等号成立当且仅当 \(p_m^{\mathrm{bin}}=\widetilde p_m\)，即所有 \(p_m^{\mathrm{bin}}(w)\) 在 \(X_m\) 上相等，从而 \(p_m^{\mathrm{bin}}=u_m\)。

最后，由命题 \ref{prop:fold-bin-chi2-col}，
\[
1+\chi^2\!\bigl(p_m^{\mathrm{bin}}\|u_m\bigr)
=
|X_m|\,\mathsf{Col}^{\mathrm{bin}}_m,
\]
故
\[
\log\!\bigl(2^m\mathsf{Col}^{\mathrm{bin}}_m\bigr)
=
\log\!\Bigl(\frac{2^m}{|X_m|}\Bigr)
\;+\;
\log\!\Bigl(1+\chi^2\!\bigl(p_m^{\mathrm{bin}}\|u_m\bigr)\Bigr),
\]
从而得到所述恒等链。证毕。
\end{proof}

\begin{theorem}[规范群体积的三因子精确乘法分解]\label{thm:xi-foldbin-gauge-volume-three-factor-product}
\leanverified{paper\_xi\_foldbin\_gauge\_volume\_three\_factor\_product}
沿用
\(
G_m^{\mathrm{bin}}:=\Aut(\Fold_m^{\mathrm{bin}})
\)
的记号。
则存在实数 \(r_m\) 满足
\[
0\le r_m\le \frac{|X_m|}{12},
\]
使得
\[
\boxed{\;
|G_m^{\mathrm{bin}}|
=
\exp\!\bigl(2^m(\kappa_m-1)\bigr)\,
(2\pi)^{|X_m|/2}\,
\Bigl(\prod_{w\in X_m}d_m^{\mathrm{bin}}(w)\Bigr)^{1/2}
e^{r_m}.\;}
\]
因此，规范群体积在主熵项之外，精确分裂为 Gaussian 半密度因子与纤维谱几何均值因子，而剩余乘法误差严格受
\(
\exp(|X_m|/12)
\)
控制。
\end{theorem}
\begin{proof}
由定理 \ref{thm:fold-bin-gauge-volume-stirling-second-order}，
\[
g_m
=
\kappa_m-1+\frac{1}{2^{m+1}}
\Bigl(
|X_m|\log(2\pi)+\sum_{w\in X_m}\log d_m^{\mathrm{bin}}(w)
\Bigr)+R_m,
\]
其中
\[
0\le R_m\le \frac{|X_m|}{12\cdot 2^m}.
\]
两边同乘 \(2^m\)，并记
\(
r_m:=2^mR_m
\)，则
\[
\log|G_m^{\mathrm{bin}}|
=
2^m(\kappa_m-1)
+\frac12\Bigl(
|X_m|\log(2\pi)+\sum_{w\in X_m}\log d_m^{\mathrm{bin}}(w)
\Bigr)
+r_m,
\]
且 \(0\le r_m\le |X_m|/12\)。
对该恒等式指数化即得所示乘法分解。证毕。
\end{proof}

\begin{theorem}[审计偶数窗上最小退化扇区强制的显式碰撞下界]\label{thm:xi-foldbin-audited-even-minsector-collision-explicit-floor}
\leanverified{paper\_xi\_foldbin\_audited\_even\_minsector\_collision\_explicit\_floor}
若
\(
m\in\{6,8,10,12\}
\)，则
\[
\boxed{\;
\mathsf{Col}^{\mathrm{bin}}_m
\ge
\frac{1}{2^{2m}}
\left(
F_mF_{m/2}^{\,2}
+
\frac{\bigl(2^m-F_mF_{m/2}\bigr)^2}{F_{m+1}}
\right).\;}
\]
等价地，
\[
\boxed{\;
\chi^2\!\bigl(p_m^{\mathrm{bin}}\|u_m\bigr)
\ge
\frac{F_{m+2}}{2^{2m}}
\left(
F_mF_{m/2}^{\,2}
+
\frac{\bigl(2^m-F_mF_{m/2}\bigr)^2}{F_{m+1}}
\right)-1.\;}
\]
并且等号成立当且仅当所有非最小退化纤维具有同一尺寸。
\end{theorem}
\begin{proof}
记
\[
n:=|X_m|=F_{m+2},
\qquad
r:=\bigl|\mathcal S_{m,d_{\min}(m)}\bigr|=F_m,
\qquad
a:=d_{\min}(m)=F_{m/2},
\]
其中第一、第二、第三个等式分别来自稳定类型计数律、定理 \ref{thm:fold-bin-min-degeneracy-fib-audited} 与同一定理。
再由推论 \ref{cor:fold-bin-minsector-complement-maxfiber-audited}，
\[
n-r=F_{m+1}.
\]

把全部纤维大小写为
\[
\underbrace{a,\dots,a}_{r\text{ 次}},\ e_1,\dots,e_{n-r},
\]
则
\[
\sum_{i=1}^{n-r}e_i=2^m-ra=2^m-F_mF_{m/2}.
\]
于是
\[
\sum_{w\in X_m}\bigl(d_m^{\mathrm{bin}}(w)\bigr)^2
=
ra^2+\sum_{i=1}^{n-r}e_i^2.
\]
对后者应用 Cauchy--Schwarz 不等式，
\[
\sum_{i=1}^{n-r}e_i^2
\ge
\frac{\bigl(\sum_{i=1}^{n-r}e_i\bigr)^2}{n-r}
=
\frac{\bigl(2^m-F_mF_{m/2}\bigr)^2}{F_{m+1}}.
\]
因此
\[
\sum_{w\in X_m}\bigl(d_m^{\mathrm{bin}}(w)\bigr)^2
\ge
F_mF_{m/2}^{\,2}
+
\frac{\bigl(2^m-F_mF_{m/2}\bigr)^2}{F_{m+1}}.
\]
再由
\(
\mathsf{Col}^{\mathrm{bin}}_m
=
2^{-2m}\sum_w(d_m^{\mathrm{bin}}(w))^2
\)
得到碰撞下界。
最后由命题 \ref{prop:fold-bin-chi2-col} 的恒等式
\(
\chi^2(p_m^{\mathrm{bin}}\|u_m)=|X_m|\mathsf{Col}^{\mathrm{bin}}_m-1
\)
即得第二条公式。

Cauchy--Schwarz 取等当且仅当
\(
e_1=\cdots=e_{n-r}
\)，也就是全部非最小退化纤维等大。证毕。
\end{proof}

\begin{theorem}[已审计偶数窗存在完全线性的 Fibonacci 低预算相]\label{thm:xi-foldbin-audited-even-fibonacci-linear-low-budget-phase}
\leanverified{paper\_xi\_foldbin\_audited\_even\_fibonacci\_linear\_low\_budget\_phase}
若
\[
m\in\{6,8,10,12\}
\]
且整数预算 \(B\ge 0\) 满足
\[
2^B\le F_{m/2},
\]
则每一条 bin-fold 纤维都超过预算阈值，因而
\[
\boxed{\;
C_m^{\mathrm{bin}}(B)=2^B\,F_{m+2},\qquad
\mathrm{Succ}_m^{\mathrm{bin}}(B)=2^{B-m}F_{m+2}.\;}
\]
换言之，在已审计偶数窗上存在一段完全线性的低预算相；其斜率与截距仅由
\(
F_{m/2}
\)
与
\(
F_{m+2}
\)
两条 Fibonacci 尺度共同控制，而与更高退化层的细节无关。
\end{theorem}
\begin{proof}
由定理 \ref{thm:fold-bin-min-degeneracy-fib-audited}，当 \(m\in\{6,8,10,12\}\) 时，对每个 \(w\in X_m\) 都有
\[
d_m^{\mathrm{bin}}(w)\ge d_{\min}(m)=F_{m/2}.
\]
若 \(2^B\le F_{m/2}\)，则对所有 \(w\in X_m\) 皆有
\[
\min\bigl(d_m^{\mathrm{bin}}(w),2^B\bigr)=2^B.
\]
因此
\[
C_m^{\mathrm{bin}}(B)
=
\sum_{w\in X_m}\min\bigl(d_m^{\mathrm{bin}}(w),2^B\bigr)
=
\sum_{w\in X_m}2^B
=
2^B|X_m|.
\]
再由 \(|X_m|=F_{m+2}\)，便得到
\[
C_m^{\mathrm{bin}}(B)=2^B F_{m+2}.
\]
最后两边除以 \(2^m\)，即得
\[
\mathrm{Succ}_m^{\mathrm{bin}}(B)=2^{B-m}F_{m+2}.
\]
证毕。
\leanverified{paper\_xi\_time\_part69\_audited\_even\_pure\_pigeonhole\_linear\_phase}
\end{proof}

\paragraph{审计偶数窗的奇偶荷截面、可解阴影与纯奇偶荷阶段}
在审计偶数窗上的中心坍缩、Abel 满秩、规范体积与低预算相已经确定之后，还可把
\(
G_m^{\mathrm{bin}}
\)
的可解阴影层进一步压缩为下列几条直接结论。它们分别给出阿贝尔化的一个规范奇偶荷截面、最大可解商的显式相图、中心坍缩与可解非交换坍缩的不同步、纯奇偶荷阶段的半单可解残核、任意可解阴影相对完整规范复杂度的指数级稀薄性，以及
\(
\kappa_m
\)
对规范体积的双侧压缩。

\begin{theorem}[奇偶荷截面的中心交公式]\label{thm:xi-foldbin-even-window-parity-section-center-intersection}
对每个审计偶数窗
\[
m\in\{6,8,10,12\},
\]
沿用定理 \ref{thm:xi-foldbin-audited-even-center-collapse-abel-full-rank} 的记号，并在每个因子
\(
\mathfrak S_{d_m^{\mathrm{bin}}(w)}
\)
中任选一条换位
\(
\tau_w
\)。
定义
\[
E_m:=\bigl\langle \tau_w:\ w\in X_m\bigr\rangle\subset G_m^{\mathrm{bin}}.
\]
则有
\[
E_m\cong (\ZZ/2\ZZ)^{|X_m|},
\]
并且复合映射
\[
E_m\hookrightarrow G_m^{\mathrm{bin}}\twoheadrightarrow
\bigl(G_m^{\mathrm{bin}}\bigr)^{\mathrm{ab}}
\]
是同构。此外，
\[
E_m\cap Z\!\bigl(G_m^{\mathrm{bin}}\bigr)\cong (\ZZ/2\ZZ)^{B_m},
\qquad
B_m:=\#\{w\in X_m:\ d_m^{\mathrm{bin}}(w)=2\}.
\]
特别地，
\[
B_6=8,
\qquad
B_8=B_{10}=B_{12}=0.
\]
\end{theorem}
\begin{proof}
不同的 \(w\) 对应于直积分解
\[
G_m^{\mathrm{bin}}
\cong
\prod_{w\in X_m}\mathfrak S_{d_m^{\mathrm{bin}}(w)}
\]
中的不同坐标因子，因此所选换位两两可交换，且每个 \(\tau_w\) 都满足 \(\tau_w^2=1\)。于是
\[
E_m\cong (\ZZ/2\ZZ)^{|X_m|}.
\]

又由定理 \ref{thm:fold-bin-min-degeneracy-fib-audited}，审计偶数窗上每条 bin-fold 纤维都满足
\(
d_m^{\mathrm{bin}}(w)\ge d_{\min}(m)=F_{m/2}\ge 2
\)。
因此每个因子的符号映射都给出满同态
\[
\operatorname{sgn}_w:\mathfrak S_{d_m^{\mathrm{bin}}(w)}\twoheadrightarrow \ZZ/2\ZZ.
\]
把这些符号映射相乘，得到
\[
\prod_{w\in X_m}\operatorname{sgn}_w:
G_m^{\mathrm{bin}}
\twoheadrightarrow
(\ZZ/2\ZZ)^{|X_m|}.
\]
而每个 \(\tau_w\) 在第 \(w\) 个坐标上给出唯一的奇元，在其余坐标上为单位元，所以该映射限制到 \(E_m\) 后恰是标准坐标同构。再结合定理 \ref{thm:xi-foldbin-audited-even-center-collapse-abel-full-rank} 的
\[
\bigl(G_m^{\mathrm{bin}}\bigr)^{\mathrm{ab}}
\cong
(\ZZ/2\ZZ)^{|X_m|},
\]
便知复合映射
\(
E_m\hookrightarrow G_m^{\mathrm{bin}}\twoheadrightarrow (G_m^{\mathrm{bin}})^{\mathrm{ab}}
\)
为同构。

最后，由定理 \ref{thm:xi-foldbin-audited-even-center-collapse-abel-full-rank}，
\[
Z\!\bigl(G_m^{\mathrm{bin}}\bigr)\cong (\ZZ/2\ZZ)^{B_m}.
\]
直积群的中心等于各因子中心的直积，而
\[
Z(\mathfrak S_d)=
\begin{cases}
\mathfrak S_2,& d=2,\\
1,& d\ge 3.
\end{cases}
\]
故 \(\tau_w\) 落在中心中当且仅当 \(d_m^{\mathrm{bin}}(w)=2\)。于是
\[
E_m\cap Z\!\bigl(G_m^{\mathrm{bin}}\bigr)
=
\bigl\langle \tau_w:\ d_m^{\mathrm{bin}}(w)=2\bigr\rangle
\cong
(\ZZ/2\ZZ)^{B_m}.
\]
最后再由表 \ref{tab:fold_bin_degeneracy_spectrum_m6_m12} 的审计直方图读出
\(
B_6=8
\)
与
\(
B_8=B_{10}=B_{12}=0
\)。
证毕。
\end{proof}
\leanverified{paper\_xi\_foldbin\_even\_window\_parity\_section\_center\_intersection}

\begin{theorem}[审计偶数窗最大可解商的显式相图]\label{thm:xi-foldbin-audited-even-maximal-solvable-quotient}
对
\[
m\in\{6,8,10,12\},
\]
记
\(
Q_m^{\mathrm{sol}}
\)
为
\(
G_m^{\mathrm{bin}}
\)
的最大可解商。定义
\[
q(d):=
\begin{cases}
\mathfrak S_d,& d\le 4,\\
\ZZ/2\ZZ,& d\ge 5.
\end{cases}
\]
则有自然同构
\[
Q_m^{\mathrm{sol}}
\cong
\prod_{w\in X_m}q\!\bigl(d_m^{\mathrm{bin}}(w)\bigr).
\]
特别地，
\[
Q_6^{\mathrm{sol}}
\cong
\mathfrak S_2^{\,8}\times \mathfrak S_3^{\,4}\times \mathfrak S_4^{\,9},
\]
\[
Q_8^{\mathrm{sol}}
\cong
\mathfrak S_3^{\,21}\times (\ZZ/2\ZZ)^{34},
\]
\[
Q_{10}^{\mathrm{sol}}
\cong
(\ZZ/2\ZZ)^{144},
\qquad
Q_{12}^{\mathrm{sol}}
\cong
(\ZZ/2\ZZ)^{377}.
\]
\end{theorem}
\begin{proof}
对称群的标准结构表明：
\[
\mathfrak S_d\ \text{可解}
\iff
d\le 4,
\]
而当 \(d\ge 5\) 时，
\(
A_d
\)
是非阿贝尔单群，因此
\[
\mathfrak S_d/A_d\cong \ZZ/2\ZZ
\]
是 \(\mathfrak S_d\) 的最大可解商。故上式中定义的 \(q(d)\) 正是因子 \(\mathfrak S_d\) 的最大可解商。

现在令
\[
\phi:
G_m^{\mathrm{bin}}
\cong
\prod_{w\in X_m}\mathfrak S_{d_m^{\mathrm{bin}}(w)}
\longrightarrow H
\]
为任意到可解群 \(H\) 的同态。则其在每个坐标因子上的限制
\[
\phi_w:\mathfrak S_{d_m^{\mathrm{bin}}(w)}\to H
\]
具有可解像，因此都必经由 \(q(d_m^{\mathrm{bin}}(w))\)。由直积的泛性质，\(\phi\) 必经由
\[
\prod_{w\in X_m}q\!\bigl(d_m^{\mathrm{bin}}(w)\bigr).
\]
这就证明了该直积正是
\(
G_m^{\mathrm{bin}}
\)
的最大可解商。

最后只需把表 \ref{tab:fold_bin_degeneracy_spectrum_m6_m12} 的审计直方图代入。对
\(
m=6
\)
有
\[
2:8,\qquad 3:4,\qquad 4:9,
\]
故
\[
Q_6^{\mathrm{sol}}
\cong
\mathfrak S_2^{\,8}\times \mathfrak S_3^{\,4}\times \mathfrak S_4^{\,9}.
\]
对
\(
m=8
\)
有
\[
3:21,\qquad 5:11,\qquad 6:23,
\]
故
\[
Q_8^{\mathrm{sol}}
\cong
\mathfrak S_3^{\,21}\times (\ZZ/2\ZZ)^{11+23}
=
\mathfrak S_3^{\,21}\times (\ZZ/2\ZZ)^{34}.
\]
对
\(
m=10
\)
有
\[
5:55,\qquad 8:52,\qquad 9:37,
\]
全部因子都贡献符号商，所以
\[
Q_{10}^{\mathrm{sol}}\cong (\ZZ/2\ZZ)^{55+52+37}=(\ZZ/2\ZZ)^{144}.
\]
同理，对
\(
m=12
\)
的直方图
\[
8:144,\qquad 12:85,\qquad 13:148
\]
可得
\[
Q_{12}^{\mathrm{sol}}\cong (\ZZ/2\ZZ)^{144+85+148}=(\ZZ/2\ZZ)^{377}.
\]
证毕。
\end{proof}
\leanverified{paper\_xi\_foldbin\_audited\_even\_maximal\_solvable\_quotient}

\begin{corollary}[中心坍缩与可解非交换坍缩不同步]\label{cor:xi-foldbin-center-solvable-collapse-two-stage}
在审计偶数窗上，中心坍缩与可解非交换坍缩是两件不同层级的事件。
更具体地，
\[
Z\!\bigl(G_m^{\mathrm{bin}}\bigr)=1
\qquad (m=8,10,12),
\]
但
\[
Q_8^{\mathrm{sol}}
\cong
\mathfrak S_3^{\,21}\times (\ZZ/2\ZZ)^{34}
\]
仍然保留非阿贝尔可解因子，而
\[
Q_{10}^{\mathrm{sol}}
\cong
(\ZZ/2\ZZ)^{144},
\qquad
Q_{12}^{\mathrm{sol}}
\cong
(\ZZ/2\ZZ)^{377}
\]
才首次退化为纯奇偶荷格。
因此，在已审计偶数窗中，“中心已死而可解非交换尚存”的现象恰发生在 \(m=8\)。
\end{corollary}
\begin{proof}
由定理 \ref{thm:xi-foldbin-audited-even-center-collapse-abel-full-rank}，
\[
Z\!\bigl(G_m^{\mathrm{bin}}\bigr)=1
\qquad (m=8,10,12).
\]
另一方面，定理 \ref{thm:xi-foldbin-audited-even-maximal-solvable-quotient} 给出
\[
Q_8^{\mathrm{sol}}
\cong
\mathfrak S_3^{\,21}\times (\ZZ/2\ZZ)^{34},
\]
故在 \(m=8\) 时最大可解商仍含非阿贝尔可解因子；而在
\(
m=10,12
\)
时，最大可解商都已坍缩为纯二元奇偶荷格。证毕。
\end{proof}
\leanverified{paper\_xi\_foldbin\_center\_solvable\_collapse\_two\_stage}

\leanverified{paper\_xi\_foldbin\_pure\_parity\_solvable\_residual\_kernel}
\begin{theorem}[纯奇偶荷阶段的半单可解残核]\label{thm:xi-foldbin-pure-parity-solvable-residual-kernel}
设分辨率 \(m\) 满足
\[
d_m^{\mathrm{bin}}(w)\ge 5
\qquad (\forall\, w\in X_m).
\]
定义正规子群
\[
N_m:=\prod_{w\in X_m}A_{d_m^{\mathrm{bin}}(w)}
\triangleleft
G_m^{\mathrm{bin}}.
\]
则有短正合列
\[
1\longrightarrow N_m
\longrightarrow G_m^{\mathrm{bin}}
\longrightarrow (\ZZ/2\ZZ)^{|X_m|}
\longrightarrow 1.
\]
并且：
\begin{enumerate}
\item \(N_m\) 是 \(G_m^{\mathrm{bin}}\) 的可解残核，即所有到可解群同态之核的交；
\item \(N_m\) 是 \(|X_m|\) 个非阿贝尔单群的直积；
\item 任意到可解群的同态都必经由奇偶荷格
\(
(\ZZ/2\ZZ)^{|X_m|}
\)。
\end{enumerate}
特别地，表 \ref{tab:fold_bin_degeneracy_spectrum_m6_m12} 表明该结论适用于
\[
m=10,\qquad m=12.
\]
\end{theorem}
\begin{proof}
在假设
\(
d_m^{\mathrm{bin}}(w)\ge 5
\)
下，每个交错群
\(
A_{d_m^{\mathrm{bin}}(w)}
\)
都是非阿贝尔单群，并满足
\[
\mathfrak S_{d_m^{\mathrm{bin}}(w)}/A_{d_m^{\mathrm{bin}}(w)}
\cong
\ZZ/2\ZZ.
\]
因此
\[
G_m^{\mathrm{bin}}/N_m
\cong
\prod_{w\in X_m}
\mathfrak S_{d_m^{\mathrm{bin}}(w)}/A_{d_m^{\mathrm{bin}}(w)}
\cong
(\ZZ/2\ZZ)^{|X_m|},
\]
从而得到所示短正合列。

再令
\(
\phi:G_m^{\mathrm{bin}}\to H
\)
为任意到可解群 \(H\) 的同态。其在每个坐标因子上的限制
\[
\phi_w:\mathfrak S_{d_m^{\mathrm{bin}}(w)}\to H
\]
具有可解像，而
\(
A_{d_m^{\mathrm{bin}}(w)}
\)
是 \(\mathfrak S_{d_m^{\mathrm{bin}}(w)}\) 的唯一非平凡非阿贝尔单正规子群，因此必落在 \(\ker\phi_w\) 中。对全部 \(w\) 同时成立，于是
\[
N_m\subseteq \ker\phi.
\]
这说明任意到可解群的同态都经由
\(
G_m^{\mathrm{bin}}/N_m\cong (\ZZ/2\ZZ)^{|X_m|}
\)，故 \(N_m\) 正是可解残核。

第二条由 \(N_m\) 的定义与每个 \(A_{d_m^{\mathrm{bin}}(w)}\) 的非阿贝尔单性立即得出；第三条只是第一条在商层的重述。最后，由表 \ref{tab:fold_bin_degeneracy_spectrum_m6_m12} 可见 \(m=10,12\) 时所有纤维大小都至少为 \(5\)，故这两个审计偶数窗均处于纯奇偶荷阶段。证毕。
\end{proof}

\begin{theorem}[任意可解阴影相对于规范复杂度的指数稀薄律]\label{thm:xi-foldbin-solvable-shadow-exponential-thinning}
设分辨率 \(m\) 满足
\[
d_m^{\mathrm{bin}}(w)\ge 5
\qquad (\forall\, w\in X_m),
\]
并令 \(Q\) 为 \(G_m^{\mathrm{bin}}\) 的任意非平凡可解商。则有
\[
\log |Q|
\le
|X_m|\log 2
=
F_{m+2}\log 2.
\]
另一方面，
\[
\log\bigl|G_m^{\mathrm{bin}}\bigr|
=
2^m g_m
=
2^m
\left(
m\log\!\frac{2}{\varphi}
-1-\frac{\log\varphi}{1+\varphi^2}
+O\!\left(m\left(\frac{\varphi}{2}\right)^m\right)
\right).
\]
因此
\[
\frac{\log|G_m^{\mathrm{bin}}|}{\log|Q|}
\ge
\frac{2^m}{F_{m+2}\log 2}
\left(
m\log\!\frac{2}{\varphi}
-1-\frac{\log\varphi}{1+\varphi^2}
+O\!\left(m\left(\frac{\varphi}{2}\right)^m\right)
\right).
\]
沿任意一条满足上述假设的共尾子列，进一步有
\[
\frac{\log|G_m^{\mathrm{bin}}|}{\log|Q|}
\ge
\frac{\sqrt5}{\varphi^2\log 2}
\left(
m\log\!\frac{2}{\varphi}
-1-\frac{\log\varphi}{1+\varphi^2}
+o(m)
\right)
\left(\frac{2}{\varphi}\right)^m.
\]
\end{theorem}
\leanverified{paper\_xi\_foldbin\_solvable\_shadow\_exponential\_thinning}
\begin{proof}
由定理 \ref{thm:xi-foldbin-pure-parity-solvable-residual-kernel}，任意可解商都必经由奇偶荷格
\[
(\ZZ/2\ZZ)^{|X_m|}.
\]
因此
\[
|Q|\le 2^{|X_m|},
\]
也就是
\[
\log |Q|\le |X_m|\log 2.
\]
再由稳定类型计数律 \(|X_m|=F_{m+2}\)，得到第一条公式。

另一方面，定理 \ref{thm:xi-foldbin-gauge-density-exact-thermodynamic-constant} 给出
\[
g_m
=
m\log\!\Bigl(\frac{2}{\varphi}\Bigr)
-1-\frac{\log\varphi}{1+\varphi^2}
+O\!\Bigl(m\Bigl(\frac{\varphi}{2}\Bigr)^m\Bigr).
\]
乘以 \(2^m\) 便得到
\[
\log|G_m^{\mathrm{bin}}|=2^m g_m
\]
的所示展开。把这两个估计合并，便得到第二条不等式。

最后，沿任意满足假设的共尾子列，Binet 公式给出
\[
F_{m+2}
=
\frac{\varphi^{m+2}}{\sqrt5}+o(\varphi^m),
\]
因此
\[
\frac{2^m}{F_{m+2}\log 2}
=
\frac{\sqrt5}{\varphi^2\log 2}
\left(\frac{2}{\varphi}\right)^m
\bigl(1+o(1)\bigr).
\]
与前式相乘后，主项来自
\(
m\log(2/\varphi)
\)，而全部误差可统一吸收到括号中的 \(o(m)\) 中，故得到最后一式。证毕。
\end{proof}

因此，一旦进入纯奇偶荷阶段，任何可解证书所能承载的全部信息量都至多线性于
\(
F_{m+2}
\)，而完整规范复杂度的主尺度却是
\(
2^m\log(2/\varphi)
\)；
可解阴影相对于完整 gauge 体积因此只是指数级稀薄的投影。

\leanverified{paper\_xi\_foldbin\_gauge\_volume\_kappa\_two\_sided\_compression}
\begin{theorem}[规范体积对 \texorpdfstring{$\kappa_m$}{kappa_m} 的双侧压缩]\label{thm:xi-foldbin-gauge-volume-kappa-two-sided-compression}
对任意分辨率 \(m\)，记
\[
d_{\min}(m):=\min_{w\in X_m}d_m^{\mathrm{bin}}(w).
\]
则成立严格双侧界
\[
2^m(\kappa_m-1)
+\frac{|X_m|}{2}\log\!\bigl(2\pi d_{\min}(m)\bigr)
\le
\log\bigl|G_m^{\mathrm{bin}}\bigr|
\]
以及
\[
\log\bigl|G_m^{\mathrm{bin}}\bigr|
\le
2^m(\kappa_m-1)
+\frac{|X_m|}{2}\log\!\Bigl(2\pi\frac{2^m}{|X_m|}\Bigr)
+\frac{|X_m|}{12}.
\]
特别地，
\[
\log\bigl|G_m^{\mathrm{bin}}\bigr|
=
2^m(\kappa_m-1)+O(|X_m|\,m).
\]
\end{theorem}
\begin{proof}
由定理 \ref{thm:xi-foldbin-gauge-volume-three-factor-product}，存在实数 \(r_m\) 满足
\[
0\le r_m\le \frac{|X_m|}{12},
\]
并且
\[
\log\bigl|G_m^{\mathrm{bin}}\bigr|
=
2^m(\kappa_m-1)
+\frac{|X_m|}{2}\log(2\pi)
+\frac12\sum_{w\in X_m}\log d_m^{\mathrm{bin}}(w)
+r_m.
\]
下界由
\[
\sum_{w\in X_m}\log d_m^{\mathrm{bin}}(w)
\ge
|X_m|\log d_{\min}(m)
\]
直接推出，即
\[
\log\bigl|G_m^{\mathrm{bin}}\bigr|
\ge
2^m(\kappa_m-1)
+\frac{|X_m|}{2}\log\!\bigl(2\pi d_{\min}(m)\bigr).
\]

上界则由 \(\log\) 的凹性与质量守恒
\[
\sum_{w\in X_m}d_m^{\mathrm{bin}}(w)=2^m
\]
给出：
\[
\frac1{|X_m|}
\sum_{w\in X_m}\log d_m^{\mathrm{bin}}(w)
\le
\log\!\left(
\frac1{|X_m|}
\sum_{w\in X_m}d_m^{\mathrm{bin}}(w)
\right)
=
\log\!\left(\frac{2^m}{|X_m|}\right).
\]
代回并使用
\(
r_m\le |X_m|/12
\)
便得
\[
\log\bigl|G_m^{\mathrm{bin}}\bigr|
\le
2^m(\kappa_m-1)
+\frac{|X_m|}{2}\log\!\Bigl(2\pi\frac{2^m}{|X_m|}\Bigr)
+\frac{|X_m|}{12}.
\]

最后，由
\(
|X_m|=F_{m+2}=\Theta(\varphi^m)
\)
且
\(
\log(2^m/|X_m|)=O(m)
\)，
可知两端修正项都只有 \(O(|X_m|m)\) 的量级，从而
\[
\log\bigl|G_m^{\mathrm{bin}}\bigr|
=
2^m(\kappa_m-1)+O(|X_m|\,m).
\]
证毕。
\end{proof}

\begin{conjecture}[可解阴影的最终纯奇偶荷化]\label{conj:xi-foldbin-solvable-shadow-eventual-pure-parity}
存在偶数阈值 \(m_0\)，使得对一切充分大的偶数 \(m\)，最大可解商都满足
\[
Q_m^{\mathrm{sol}}
\cong
(\ZZ/2\ZZ)^{|X_m|}.
\]
等价地说，从某个偶数窗口开始，最大可解商中不再保留任何
\(
\mathfrak S_3
\)
或
\(
\mathfrak S_4
\)
型的非阿贝尔可解残余，全部可解信息完全坍缩为纤维奇偶荷格。

当前已审计数据进一步指向更强的版本：\(m=8\) 是最后一个“中心已死而可解非交换尚存”的审计偶数窗。
\end{conjecture}

综上，bin-fold 规范群中的阿贝尔可见性、中心可见性与可解可见性并不属于同一层级：偶数审计窗先失去中心，再失去可解非交换阴影，而一旦进入纯奇偶荷阶段，全部可解信息便被压缩为
\(
(\ZZ/2\ZZ)^{|X_m|}
\)
这一指数级稀薄的外显投影。



\paragraph{规范群熵账本、有限群代数障碍与偶 \texorpdfstring{$\zeta$}{zeta} 反演}
在规范群体积的二阶 Stirling 展开、审计偶数窗的最小退化律、约数触发零余类并的碰撞削减机制以及 fold-gauge 常数的 Stirling--Bernoulli 层级已经建立之后，还可进一步压缩出下列五条直接结论。

\begin{theorem}[规范群熵的精确账本恒等式]\label{thm:xi-time-part60aa-gauge-entropy-ledger-exact-identity}
沿用定理 \ref{thm:fold-bin-gauge-volume-stirling-second-order} 的记号，定义
\[
p_m^{\mathrm{bin}}(w):=\frac{d_m^{\mathrm{bin}}(w)}{2^m}
\qquad (w\in X_m),
\]
并记其 Shannon 熵为
\[
H\!\bigl(p_m^{\mathrm{bin}}\bigr):=
-\sum_{w\in X_m}p_m^{\mathrm{bin}}(w)\log p_m^{\mathrm{bin}}(w).
\]
则规范群体积密度 \(g_m\) 满足严格恒等式
\[
\boxed{\;
g_m
=
m\log 2
-H\!\bigl(p_m^{\mathrm{bin}}\bigr)
-1
+\frac{|X_m|\log(2\pi)+\sum_{w\in X_m}\log d_m^{\mathrm{bin}}(w)}{2^{m+1}}
+R_m,\;}
\]
其中余项满足
\[
\boxed{\;
0\le R_m\le \frac{|X_m|}{12\cdot 2^m}.\;}
\]
特别地，
\[
\boxed{\;
g_m
=
m\log 2
-H\!\bigl(p_m^{\mathrm{bin}}\bigr)
-1
+O\!\Bigl(\frac{|X_m|\,m}{2^m}\Bigr),\;}
\]
并且
\[
\boxed{\;
\kappa_m=m\log 2-H\!\bigl(p_m^{\mathrm{bin}}\bigr).\;}
\]
因此，\(\kappa_m\) 恰是 \(g_m\) 的零阶主项。
\end{theorem}
\begin{proof}
由定理 \ref{thm:fold-bin-gauge-volume-stirling-second-order}，
\[
g_m
=
\kappa_m-1+\frac{|X_m|\log(2\pi)+\sum_{w\in X_m}\log d_m^{\mathrm{bin}}(w)}{2^{m+1}}
+R_m,
\]
且
\[
0\le R_m\le \frac{|X_m|}{12\cdot 2^m}.
\]

另一方面，由
\[
p_m^{\mathrm{bin}}(w)=\frac{d_m^{\mathrm{bin}}(w)}{2^m},
\qquad
\sum_{w\in X_m}d_m^{\mathrm{bin}}(w)=2^m,
\]
可直接计算
\[
\begin{aligned}
H\!\bigl(p_m^{\mathrm{bin}}\bigr)
&=
-\sum_{w\in X_m}\frac{d_m^{\mathrm{bin}}(w)}{2^m}
\log\frac{d_m^{\mathrm{bin}}(w)}{2^m}\\
&=
-\frac{1}{2^m}\sum_{w\in X_m}d_m^{\mathrm{bin}}(w)\log d_m^{\mathrm{bin}}(w)
+\frac{m\log 2}{2^m}\sum_{w\in X_m}d_m^{\mathrm{bin}}(w)\\
&=
-\kappa_m+m\log 2.
\end{aligned}
\]
即
\[
\kappa_m=m\log 2-H\!\bigl(p_m^{\mathrm{bin}}\bigr).
\]
把该式代回 \(g_m\) 的 Stirling 展开，即得第一条盒装恒等式。

最后，由 \(1\le d_m^{\mathrm{bin}}(w)\le 2^m\) 可知
\[
0\le \sum_{w\in X_m}\log d_m^{\mathrm{bin}}(w)\le |X_m|\,m\log 2,
\]
再结合余项界，便得到
\[
g_m
=
m\log 2-H\!\bigl(p_m^{\mathrm{bin}}\bigr)-1
+O\!\Bigl(\frac{|X_m|\,m}{2^m}\Bigr).
\]
证毕。
\end{proof}

\begin{theorem}[连续容量曲线的对数 Stieltjes 表示与规范体积密度的普适常数差]\label{thm:xi-time-part60aa-log-stieltjes-capacity-gauge-density}
沿用定理 \ref{thm:xi-time-part60aa-gauge-entropy-ledger-exact-identity} 的记号，并定义连续容量曲线
\[
\mathcal C_m^{\mathrm{cont}}(T):=\sum_{w\in X_m}\min\bigl(d_m^{\mathrm{bin}}(w),T\bigr)
\qquad(T\ge 0).
\]
则有精确 Stieltjes 表示
\[
\boxed{\;
\kappa_m
=
1+\frac{1}{2^m}\int_0^\infty \log T\,\dd \mathcal C_m^{\mathrm{cont}}(T).\;}
\]
进一步，规范体积密度满足精确恒等式
\[
\boxed{\;
g_m
=
\frac{1}{2^m}\int_0^\infty \log T\,\dd \mathcal C_m^{\mathrm{cont}}(T)
+\frac{|X_m|\log(2\pi)+\sum_{w\in X_m}\log d_m^{\mathrm{bin}}(w)}{2^{m+1}}
+R_m,\;}
\]
其中
\[
0\le R_m\le \frac{|X_m|}{12\cdot 2^m}.
\]
因此
\[
\boxed{\;
g_m
=
\frac{1}{2^m}\int_0^\infty \log T\,\dd \mathcal C_m^{\mathrm{cont}}(T)
+O\!\Bigl(m\Bigl(\frac{\varphi}{2}\Bigr)^m\Bigr),\;}
\]
并且
\[
\boxed{\;
\kappa_m-g_m
=
1+O\!\Bigl(m\Bigl(\frac{\varphi}{2}\Bigr)^m\Bigr).\;}
\]
换言之，\(\kappa_m\) 恰是连续容量曲线的对数 Stieltjes 积分再加上一枚普适常数 \(1\)，而 \(g_m\) 与该积分之间只差指数小校正。
\end{theorem}
\begin{proof}
对每个固定 \(w\in X_m\)，函数
\[
T\longmapsto \min\bigl(d_m^{\mathrm{bin}}(w),T\bigr)
\]
是有界变差函数，其 Stieltjes 测度正是区间 \([0,d_m^{\mathrm{bin}}(w)]\) 上的 Lebesgue 测度。故
\[
\int_0^\infty \log T\,\dd\min\bigl(d_m^{\mathrm{bin}}(w),T\bigr)
=
\int_0^{d_m^{\mathrm{bin}}(w)}\log T\,\dd T
=
d_m^{\mathrm{bin}}(w)\log d_m^{\mathrm{bin}}(w)-d_m^{\mathrm{bin}}(w).
\]
对全部 \(w\in X_m\) 求和，得到
\[
\int_0^\infty \log T\,\dd\mathcal C_m^{\mathrm{cont}}(T)
=
\sum_{w\in X_m}
\Bigl(
d_m^{\mathrm{bin}}(w)\log d_m^{\mathrm{bin}}(w)-d_m^{\mathrm{bin}}(w)
\Bigr).
\]
由于
\[
\sum_{w\in X_m}d_m^{\mathrm{bin}}(w)=2^m,
\]
而
\[
\kappa_m=\frac{1}{2^m}\sum_{w\in X_m}d_m^{\mathrm{bin}}(w)\log d_m^{\mathrm{bin}}(w),
\]
遂有
\[
\frac{1}{2^m}\int_0^\infty \log T\,\dd\mathcal C_m^{\mathrm{cont}}(T)=\kappa_m-1,
\]
即第一条盒装公式。

再将
\[
\kappa_m-1=\frac{1}{2^m}\int_0^\infty \log T\,\dd\mathcal C_m^{\mathrm{cont}}(T)
\]
代回定理 \ref{thm:xi-time-part60aa-gauge-entropy-ledger-exact-identity} 的精确账本恒等式，便得到 \(g_m\) 的精确表达。

最后，由 \(|X_m|=F_{m+2}\) 且
\[
0\le \sum_{w\in X_m}\log d_m^{\mathrm{bin}}(w)\le |X_m|\,m\log 2,
\]
可把显示校正项与 \(R_m\) 一并吸收到
\[
O\!\Bigl(\frac{|X_m|\,m}{2^m}\Bigr)
=
O\!\Bigl(m\Bigl(\frac{\varphi}{2}\Bigr)^m\Bigr)
\]
中，从而得到两条渐近式。证毕。
\end{proof}

\begin{theorem}[审计偶数窗上的纤维群胚代数不可能是有限群代数]\label{thm:xi-time-part60aa-audited-even-groupoid-no-finite-group-algebra}
对审计偶数窗
\[
m\in\{6,8,10,12\},
\]
纤维群胚代数
\[
\CC[\mathcal R_m^{\mathrm{bin}}]
\cong
\bigoplus_{w\in X_m}M_{d_m^{\mathrm{bin}}(w)}(\CC)
\]
不可能与任何有限群 \(H\) 的复群代数 \(\CC[H]\) 同构。更强地，不存在任何非零代数同态
\[
\chi:\CC[\mathcal R_m^{\mathrm{bin}}]\to\CC.
\]
\end{theorem}
\begin{proof}
由命题 \ref{prop:fold-bin-groupoid-algebra-wedderburn}，
\[
\CC[\mathcal R_m^{\mathrm{bin}}]
\cong
\bigoplus_{w\in X_m}M_{d_m^{\mathrm{bin}}(w)}(\CC).
\]
而定理 \ref{thm:fold-bin-min-degeneracy-fib-audited} 在审计偶数窗上给出
\[
d_m^{\mathrm{bin}}(w)\ge d_{\min}(m)=F_{m/2}\ge 2
\qquad(\forall\,w\in X_m).
\]
因此上述 Wedderburn 分解中不存在任何 \(M_1(\CC)\) 块。

现证不存在非零代数同态 \(\chi:\CC[\mathcal R_m^{\mathrm{bin}}]\to\CC\)。若此类 \(\chi\) 存在，则其在某个直和因子
\[
M_{d_m^{\mathrm{bin}}(w)}(\CC)
\]
上的限制必非零。由于矩阵代数 \(M_n(\CC)\) 为简单代数，非零代数同态的核只能为零，故该限制必为单射。但当 \(n\ge 2\) 时，\(M_n(\CC)\) 非交换，而其像落在交换代数 \(\CC\) 中，矛盾。故不存在非零代数同态。

若反设存在有限群 \(H\) 使
\[
\CC[\mathcal R_m^{\mathrm{bin}}]\cong \CC[H],
\]
则群代数上的增广映射
\[
\varepsilon:\CC[H]\to\CC,
\qquad
\varepsilon\!\left(\sum_{h\in H}a_h h\right):=\sum_{h\in H}a_h
\]
给出一个非零代数同态。经由上述同构传回 \(\CC[\mathcal R_m^{\mathrm{bin}}]\)，便得到矛盾。故
\[
\CC[\mathcal R_m^{\mathrm{bin}}]\not\cong \CC[H]
\]
对一切有限群 \(H\) 都成立。证毕。
\end{proof}

\begin{theorem}[约数触发零余类并给出的规范群熵显式上界]\label{thm:xi-time-part60aa-divisor-zero-class-gauge-entropy-upper-bound}
记
\[
M:=F_{m+2}.
\]
若存在真约数 \(d\mid (m+2)\)、\(d<m+2\)，满足
\[
F_d\mid \frac{M}{2},
\]
则规范群体积密度 \(g_m\) 满足严格上界
\[
\boxed{\;
g_m
\le
m\log 2
-1
+\log\!\Bigl(1-\frac{F_d}{M}\Bigr)
+\frac{|X_m|\log(2\pi)+\sum_{w\in X_m}\log d_m^{\mathrm{bin}}(w)}{2^{m+1}}
+\frac{|X_m|}{12\cdot 2^m}.\;}
\]
特别地，
\[
\boxed{\;
g_m
\le
m\log 2
-1
+\log\!\Bigl(1-\frac{F_d}{M}\Bigr)
+O\!\Bigl(\frac{|X_m|\,m}{2^m}\Bigr).\;}
\]
\end{theorem}
\begin{proof}
由定理 \ref{thm:xi-time-part60aa-gauge-entropy-ledger-exact-identity}，
\[
g_m
=
m\log 2
-H\!\bigl(p_m^{\mathrm{bin}}\bigr)
-1
+\frac{|X_m|\log(2\pi)+\sum_{w\in X_m}\log d_m^{\mathrm{bin}}(w)}{2^{m+1}}
+R_m,
\]
其中
\[
0\le R_m\le \frac{|X_m|}{12\cdot 2^m}.
\]
另一方面，Shannon 熵支配 Rényi--\(2\) 熵，故
\[
H\!\bigl(p_m^{\mathrm{bin}}\bigr)\ge H_2\!\bigl(p_m^{\mathrm{bin}}\bigr)
=-\log \mathsf{Col}^{\mathrm{bin}}_m.
\]
于是
\[
g_m
\le
m\log 2
+\log \mathsf{Col}^{\mathrm{bin}}_m
-1
+\frac{|X_m|\log(2\pi)+\sum_{w\in X_m}\log d_m^{\mathrm{bin}}(w)}{2^{m+1}}
+\frac{|X_m|}{12\cdot 2^m}.
\]

再由推论 \ref{cor:fold-zero-divisor-lb}，
\[
|\mathcal Z_m|\ge F_d,
\]
而定理 \ref{thm:fold-collision-zero-reduction} 给出
\[
\mathsf{Col}^{\mathrm{bin}}_m
\le
1-\frac{|\mathcal Z_m|}{M}
\le
1-\frac{F_d}{M}.
\]
代回上式即得第一条盒装不等式。

最后，与定理 \ref{thm:xi-time-part60aa-gauge-entropy-ledger-exact-identity} 证明末尾相同，由
\[
0\le \sum_{w\in X_m}\log d_m^{\mathrm{bin}}(w)\le |X_m|\,m\log 2
\]
可将显示校正项统一吸收到
\(
O(|X_m|m/2^m)
\)
中，从而得到第二条公式。证毕。
\end{proof}

\begin{theorem}[从规范群常数逐阶显式恢复全部偶值 \texorpdfstring{$\zeta(2r)$}{zeta(2r)}]\label{thm:xi-time-part60aa-gauge-constant-explicit-even-zeta-recovery}
沿用定理 \ref{thm:fold-bin-gauge-constant-stirling-bernoulli-hierarchy} 的规范群常数列 \((C_m)_{m\ge 1}\)，并对每个整数 \(r\ge 1\) 定义
\[
A_r:=
\frac{B_{2r}}{r(2r-1)}
\bigl(\varphi^{-1}+\varphi^{2r-3}\bigr).
\]
再定义递阶剥离后的极限算子
\[
L_r:=
\lim_{m\to\infty}
\Bigl(\frac{2}{\varphi}\Bigr)^{(2r-1)m}
\left[
\log C_m-\log(2\pi)
-\sum_{s=1}^{r-1}A_s\Bigl(\frac{\varphi}{2}\Bigr)^{(2s-1)m}
\right].
\]
则极限存在且满足
\[
\boxed{\;L_r=A_r.\;}
\]
因此对每个 \(r\ge 1\)，全部偶值 \(\zeta(2r)\) 都可由显式公式恢复：
\[
\boxed{\;
\zeta(2r)
=
(-1)^{r-1}
\frac{(2\pi)^{2r}}{2(2r)!}\,
\frac{r(2r-1)}{\varphi^{-1}+\varphi^{2r-3}}\,
L_r.\;}
\]
\end{theorem}
\begin{proof}
由定理 \ref{thm:fold-bin-gauge-constant-stirling-bernoulli-hierarchy}，对任意固定 \(r\ge 1\) 有展开
\[
\log C_m
=
\log(2\pi)
+\sum_{s=1}^{r}A_s\Bigl(\frac{\varphi}{2}\Bigr)^{(2s-1)m}
+O\!\Bigl(\Bigl(\frac{\varphi}{2}\Bigr)^{(2r+1)m}\Bigr).
\]
将前 \(r-1\) 阶项移到左边，得到
\[
\log C_m-\log(2\pi)
-\sum_{s=1}^{r-1}A_s\Bigl(\frac{\varphi}{2}\Bigr)^{(2s-1)m}
=
A_r\Bigl(\frac{\varphi}{2}\Bigr)^{(2r-1)m}
+O\!\Bigl(\Bigl(\frac{\varphi}{2}\Bigr)^{(2r+1)m}\Bigr).
\]
两边同乘
\[
\Bigl(\frac{2}{\varphi}\Bigr)^{(2r-1)m},
\]
便得到
\[
\Bigl(\frac{2}{\varphi}\Bigr)^{(2r-1)m}
\left[
\log C_m-\log(2\pi)
-\sum_{s=1}^{r-1}A_s\Bigl(\frac{\varphi}{2}\Bigr)^{(2s-1)m}
\right]
=
A_r
+O\!\Bigl(\Bigl(\frac{\varphi}{2}\Bigr)^{2m}\Bigr).
\]
由于 \(\varphi/2<1\)，右端误差趋于 \(0\)，故极限存在且等于 \(A_r\)，即 \(L_r=A_r\)。

最后，由标准 Bernoulli--zeta 恒等式
\[
B_{2r}
=
(-1)^{r-1}\frac{2(2r)!}{(2\pi)^{2r}}\zeta(2r),
\]
以及
\[
A_r=
\frac{\varphi^{-1}+\varphi^{2r-3}}{r(2r-1)}\,B_{2r}
=L_r,
\]
解出 \(\zeta(2r)\) 即得所示恢复公式。证毕。
\end{proof}

\begin{theorem}[Stirling--Bernoulli 级数截断余项的最终交错变号]\label{thm:xi-time-part60aa-stirling-bernoulli-tail-alternation}
沿用定理 \ref{thm:fold-bin-gauge-constant-stirling-bernoulli-hierarchy} 与定理 \ref{thm:xi-time-part60aa-gauge-constant-explicit-even-zeta-recovery} 的记号。记
\[
q_r:=\Bigl(\frac{\varphi}{2}\Bigr)^{2r-1},
\qquad
a_r:=A_r=
\frac{B_{2r}}{r(2r-1)}\bigl(\varphi^{-1}+\varphi^{2r-3}\bigr).
\]
对任意 \(R\ge 1\)，定义截断余项
\[
\mathcal R_R(m):=
\log C_m-\log(2\pi)-\sum_{r=1}^{R-1}a_r q_r^m.
\]
则有
\[
\mathcal R_R(m)=a_R q_R^m+O(q_{R+1}^m),
\qquad
\frac{q_{R+1}}{q_R}=\Bigl(\frac{\varphi}{2}\Bigr)^2<1.
\]
从而存在 \(m_0(R)\)，使得对一切 \(m\ge m_0(R)\)，都有
\[
\operatorname{sgn}\mathcal R_R(m)=\operatorname{sgn}(a_R)=\operatorname{sgn}(B_{2R})=(-1)^{R-1}.
\]
特别地，截断和
\[
S_{R-1}(m):=\log(2\pi)+\sum_{r=1}^{R-1}a_r q_r^m
\]
对 \(\log C_m\) 构成最终交错的上、下界：当 \(R\) 为奇数时 \(S_{R-1}(m)<\log C_m\)，当 \(R\) 为偶数时 \(S_{R-1}(m)>\log C_m\)（均对充分大的 \(m\) 成立）。由于指数函数在实轴上严格单调，相应的
\[
\exp(S_{R-1}(m))
\]
也对 \(C_m\) 构成同样交错的最终上、下界。
\end{theorem}
\begin{proof}
由定理 \ref{thm:fold-bin-gauge-constant-stirling-bernoulli-hierarchy}，对固定的 \(R\ge 1\) 有渐近展开
\[
\log C_m
=
\log(2\pi)+\sum_{r=1}^{R}a_r q_r^m+O(q_{R+1}^m).
\]
将前 \(R-1\) 项移到左边，便得到
\[
\mathcal R_R(m)=a_R q_R^m+O(q_{R+1}^m).
\]
又因为
\[
\frac{q_{R+1}}{q_R}=\Bigl(\frac{\varphi}{2}\Bigr)^2<1,
\]
所以 \(a_R q_R^m\) 终将支配余项，故存在 \(m_0(R)\) 使得对一切 \(m\ge m_0(R)\)，
\[
\operatorname{sgn}\mathcal R_R(m)=\operatorname{sgn}(a_R).
\]

再注意到
\[
\varphi^{-1}+\varphi^{2R-3}>0,
\]
故 \(a_R\) 的符号完全由 \(B_{2R}\) 决定。Bernoulli 数满足经典符号律
\[
\operatorname{sgn}(B_{2R})=(-1)^{R-1},
\]
遂得
\[
\operatorname{sgn}\mathcal R_R(m)=(-1)^{R-1}.
\]
于是当 \(R\) 为奇数时余项为正，从而
\[
S_{R-1}(m)<\log C_m;
\]
当 \(R\) 为偶数时余项为负，从而
\[
S_{R-1}(m)>\log C_m.
\]
最后，由指数函数在实轴上严格递增，同样的上下界关系自动传递到 \(C_m=\exp(\log C_m)\)。证毕。
\end{proof}

\noindent
上述五条结论把规范群与 fold-gauge 常数的两条层级链收束为同一份显式账本：一方面，\(g_m\) 精确等于原始 \(m\)-比特预算减去折叠 Shannon 熵，并附带完全显示的 Stirling 校正；另一方面，\(\log C_m\) 的 Bernoulli 模态不仅可逐阶识别，而且既可经由显式极限算子直接恢复全部偶值 \(\zeta(2r)\)，也可在任意固定截断阶上形成最终交错的上下包络。与审计偶数窗上的最小退化扇区及其零余类机制合并后，群熵损失、群代数实现障碍与偶 \(\zeta\)-塔反演遂被压缩进同一条闭合链。

\paragraph{容量层析、Wedderburn 信息完备性与最小退化块账本}
在连续容量曲线的完备统计性、规范群熵账本、纤维群胚代数的 Wedderburn 直和，以及审计偶数窗的 Fibonacci 最小退化律都已建立之后，还可把这些对象层接口进一步压缩为下列几条直接结论。

\begin{theorem}[连续容量族对实参数冻结面的层析完备性]\label{thm:xi-time-part60aaa-capacity-family-freezing-face-tomography}
设对每个 \(m\ge 1\) 给定连续容量曲线
\[
\mathcal C_m^{\mathrm{cont}}(T):=\sum_{x\in X_m}\min\bigl(d_m(x),T\bigr)
\qquad (T\ge 0).
\]
若对每个 \(q>0\) 压力极限
\[
P_q:=\lim_{m\to\infty}\frac1m\log S_q(m),
\qquad
S_q(m):=\sum_{x\in X_m}d_m(x)^q,
\]
都存在，则整族 \(\{\mathcal C_m^{\mathrm{cont}}\}_{m\ge 1}\) 唯一决定全部实参数压力谱 \(\{P_q\}_{q>0}\)。在严格基态间隙
\[
\alpha_*^{(2)}<\alpha_*
\]
下，它因而唯一决定冻结仿射面
\[
q\longmapsto q\alpha_*+g_*
\]
以及其实参数观测起点
\[
a_{\mathrm{fr}}^{\mathbb R}
:=
\inf\bigl\{
q>0:\ P_r=r\alpha_*+g_*
\ \text{对一切 }r\ge q\text{成立}
\bigr\}.
\]
\end{theorem}
\begin{proof}
由定理 \ref{thm:xi-capacity-tail-histogram-moment-complete-statistic}，对每个固定的 \(m\ge 1\) 与每个 \(q>0\)，连续容量曲线 \(\mathcal C_m^{\mathrm{cont}}\) 唯一决定幂和矩
\[
S_q(m)=\sum_{x\in X_m}d_m(x)^q.
\]
因此整族 \(\{\mathcal C_m^{\mathrm{cont}}\}_{m\ge 1}\) 唯一决定每个 \(q>0\) 的数列 \(m\mapsto S_q(m)\)，从而在极限存在时唯一决定
\[
P_q=\lim_{m\to\infty}\frac1m\log S_q(m).
\]

若进一步满足严格基态间隙，则定理 \ref{thm:fold-pressure-freezing-threshold} 给出：对每个实数 \(q>a_{\mathrm c}\)，都有
\[
P_q=q\alpha_*+g_*.
\]
故由已恢复的整条压力谱可唯一读出其高温区外的线性冻结支，并唯一恢复上述实参数起点 \(a_{\mathrm{fr}}^{\mathbb R}\)。证毕。
\end{proof}

\begin{theorem}[固定分辨率上的 Wedderburn--信息完备接口]\label{thm:xi-time-part60aaa-wedderburn-information-completeness}
\leanverified{paper\_xi\_time\_part60aaa\_wedderburn\_information\_completeness}
固定分辨率 \(m\ge 1\)。定义纤维直方图
\[
n_d(m):=\#\{x\in X_m:\ d_m(x)=d\}
\qquad(d\ge 1).
\]
则下列四类数据两两等价，并且彼此唯一决定：
\[
\textup{(i)}\ \{n_d(m)\}_{d\ge 1},
\qquad
\textup{(ii)}\ \mathcal C_m^{\mathrm{cont}},
\qquad
\textup{(iii)}\ \{S_q(m)\}_{q>0},
\qquad
\textup{(iv)}\ \mathbb C[\mathcal R_m^{\mathrm{bin}}]\ \text{的半单代数同构型}.
\]
此外，这四类等价数据还共同决定 fold-gauge 群的直积分解
\[
G_m^{\mathrm{bin}}:=\Aut(\Fold_m^{\mathrm{bin}})
\cong
\prod_{d\ge 1}\mathfrak S_d^{\,n_d(m)}.
\]
\end{theorem}
\begin{proof}
由定理 \ref{thm:xi-capacity-tail-histogram-moment-complete-statistic}，固定 \(m\) 时，连续容量曲线 \(\mathcal C_m^{\mathrm{cont}}\)、纤维直方图 \(\{n_d(m)\}_{d\ge 1}\) 与全矩谱 \(\{S_q(m)\}_{q>0}\) 彼此唯一决定。

另一方面，由命题 \ref{prop:fold-bin-groupoid-algebra-wedderburn}，
\[
\mathbb C[\mathcal R_m^{\mathrm{bin}}]
\cong
\bigoplus_{x\in X_m}M_{d_m(x)}(\mathbb C)
\cong
\bigoplus_{d\ge 1}M_d(\mathbb C)^{\oplus n_d(m)}.
\]
有限维半单复代数的矩阵块大小及其重数由 Artin--Wedderburn 定理唯一决定，因此半单代数同构型与直方图 \(\{n_d(m)\}_{d\ge 1}\) 等价。

最后，由命题 \ref{prop:fold-bin-gauge-decomposition}，
\[
G_m^{\mathrm{bin}}
\cong
\prod_{x\in X_m}\mathfrak S_{d_m(x)}
\cong
\prod_{d\ge 1}\mathfrak S_d^{\,n_d(m)}.
\]
故前述四类等价数据又共同决定 gauge 群的直积分解。证毕。
\end{proof}

\begin{proposition}[审计偶数窗上的最小退化块构成显式 Fibonacci 半单理想]\label{prop:xi-time-part60aaa-audited-even-minideal-fibonacci-semisimple}
\leanverified{paper_xi_time_part60aaa_audited_even_minideal_fibonacci_semisimple}
对审计偶数窗
\[
m\in\{6,8,10,12\},
\]
定义最小退化扇区
\[
\mathcal S_{m,*}:=\{x\in X_m:\ d_m^{\mathrm{bin}}(x)=d_{\min}(m)\},
\]
以及相应块直和
\[
A_{m,\min}
:=
\bigoplus_{x\in\mathcal S_{m,*}}M_{d_{\min}(m)}(\mathbb C)
\subset
\mathbb C[\mathcal R_m^{\mathrm{bin}}].
\]
则 \(A_{m,\min}\) 是一个半单理想，并且存在显式同构
\[
A_{m,\min}\cong M_{F_{m/2}}(\mathbb C)^{\oplus F_m}.
\]
特别地，
\[
\dim_{\mathbb C}A_{m,\min}=F_m\,F_{m/2}^{\,2},
\qquad
Z(A_{m,\min})\cong \mathbb C^{F_m}.
\]
\end{proposition}
\begin{proof}
由定理 \ref{thm:fold-bin-min-degeneracy-fib-audited}，在审计偶数窗上有
\[
d_{\min}(m)=F_{m/2},
\qquad
|\mathcal S_{m,*}|=F_m.
\]
另一方面，命题 \ref{prop:fold-bin-groupoid-algebra-wedderburn} 给出
\[
\mathbb C[\mathcal R_m^{\mathrm{bin}}]
\cong
\bigoplus_{x\in X_m}M_{d_m^{\mathrm{bin}}(x)}(\mathbb C).
\]
因此 \(A_{m,\min}\) 正是其中由全部最小退化块组成的直和，故自动是一个半单理想，并满足
\[
A_{m,\min}
\cong
M_{d_{\min}(m)}(\mathbb C)^{\oplus |\mathcal S_{m,*}|}
\cong
M_{F_{m/2}}(\mathbb C)^{\oplus F_m}.
\]
维数公式
\[
\dim_{\mathbb C}M_n(\mathbb C)=n^2
\]
立即给出
\[
\dim_{\mathbb C}A_{m,\min}=F_m\,F_{m/2}^{\,2},
\]
而矩阵代数直和的中心同构于对应块数目的复数直和，故
\[
Z(A_{m,\min})\cong \mathbb C^{F_m}.
\]
证毕。
\end{proof}

\begin{corollary}[exact gauge 账本相对 naive 缺陷账本的一纳特主阶回扣]\label{cor:xi-time-part60aaa-one-nat-unnormalized-ledger-rebate}
定义 naive 缺陷账本与 exact gauge 账本
\[
A_m:=\sum_{x\in X_m}d_m^{\mathrm{bin}}(x)\log d_m^{\mathrm{bin}}(x),
\qquad
L_m:=\log|G_m^{\mathrm{bin}}|.
\]
则有
\[
L_m
=
A_m-2^m+O(m\varphi^m).
\]
等价地，
\[
\frac{L_m}{2^m}
=
\kappa_m-1+O\!\Bigl(m\Bigl(\frac{\varphi}{2}\Bigr)^m\Bigr).
\]
\end{corollary}
\begin{proof}
由定理 \ref{thm:xi-time-part60aa-log-stieltjes-capacity-gauge-density}，
\[
\kappa_m-g_m
=
1+O\!\Bigl(m\Bigl(\frac{\varphi}{2}\Bigr)^m\Bigr),
\]
其中
\[
\kappa_m=\frac1{2^m}\sum_{x\in X_m}d_m^{\mathrm{bin}}(x)\log d_m^{\mathrm{bin}}(x)=\frac{A_m}{2^m},
\qquad
g_m=\frac1{2^m}\log|G_m^{\mathrm{bin}}|=\frac{L_m}{2^m}.
\]
代回即得第二式。再将两边同乘 \(2^m\)，并用
\[
2^m\cdot O\!\Bigl(m\Bigl(\frac{\varphi}{2}\Bigr)^m\Bigr)=O(m\varphi^m),
\]
便得到第一式。证毕。
\end{proof}

\noindent
与冻结相中的基态等权塌缩相合，上述结果表明：离散多重度谱中的层纹与波动在连续容量口径下并未消失，而是分别迁移为冻结面的线性支、群胚代数的半单块结构以及 exact gauge 账本相对于 naive 缺陷账本的一纳特主阶回扣。因而，连续平滑并不是离散结构的抹除，而是离散结构向基态等权、容量层析与 Wedderburn 分块三类函子像的重编码。

\paragraph{隐藏位压力微分、Bernoulli 递阶反演与谱零块的二项式因子化}
在推前熵账本、Rényi 散度冻结、规范群体积的 Stirling--Bernoulli 层级，以及约数驱动的谱零块均已建立之后，还可把隐藏位统计、碰撞指纹、Bernoulli--\(\zeta\) 反演与群环多项式的代数刚性进一步压缩为下列一组闭合结论。

\begin{theorem}[压力谱在 \texorpdfstring{$a=1$}{a=1} 的一、二阶导数分别给出隐藏对数多重度的密度与涨落]\label{thm:xi-time-part60ab-hidden-logmultiplicity-pressure-derivatives}
\leanverified{paper\_xi\_time\_part60ab\_hidden\_logmultiplicity\_pressure\_derivatives}
沿用 H.3.1 的纤维多重度记号。对每个 \(m\) 记
\[
\mu_m(x):=\frac{d_m(x)}{2^m},
\qquad
S_a(m):=\sum_{x\in X_m} d_m(x)^a,
\qquad
P_m(a):=\frac1m\log S_a(m).
\]
设开区间 \(I\subset (0,\infty)\) 满足 \(1\in I\)，并假设每个 \(P_m\) 在 \(I\) 上属于 \(C^2\)，且存在 \(P\in C^2(I)\) 使
\[
P_m\longrightarrow P
\qquad\text{于 }C^2_{\mathrm{loc}}(I).
\]
若 \(X\sim \mu_m\)，定义隐藏对数多重度随机变量
\[
\mathscr H_m:=\log d_m(X).
\]
则对每个 \(m\) 都有精确恒等式
\[
\frac1m\EE_{\mu_m}[\mathscr H_m]=P_m'(1),
\qquad
\frac1m\Var_{\mu_m}(\mathscr H_m)=P_m''(1),
\]
从而
\[
\boxed{\;
\lim_{m\to\infty}\frac1m\EE_{\mu_m}[\mathscr H_m]=P'(1),\;}
\]
\[
\boxed{\;
\lim_{m\to\infty}\frac1m\Var_{\mu_m}(\mathscr H_m)=P''(1).\;}
\]
\end{theorem}
\begin{proof}
定义
\[
\Phi_m(a):=\frac1m\log\sum_{x\in X_m}\mu_m(x)^a.
\]
由
\[
\mu_m(x)=\frac{d_m(x)}{2^m}
\]
可得
\[
\Phi_m(a)
=
\frac1m\log\sum_x d_m(x)^a-a\log 2
=
P_m(a)-a\log 2.
\]

另一方面，
\[
\Phi_m'(a)
=
\frac1m
\frac{\sum_x \mu_m(x)^a\log \mu_m(x)}{\sum_x\mu_m(x)^a}.
\]
在 \(a=1\) 处利用 \(\sum_x\mu_m(x)=1\)，得到
\[
\Phi_m'(1)
=
\frac1m\sum_x\mu_m(x)\log\mu_m(x)
=
\frac1m\sum_x\mu_m(x)\log d_m(x)-\log 2
=
\frac1m\EE_{\mu_m}[\mathscr H_m]-\log 2.
\]
而由 \(\Phi_m(a)=P_m(a)-a\log 2\) 又有
\[
\Phi_m'(1)=P_m'(1)-\log 2.
\]
两式相减即得
\[
\frac1m\EE_{\mu_m}[\mathscr H_m]=P_m'(1).
\]

再看二阶导数。标准指数族计算给出
\[
\Phi_m''(1)
=
\frac1m
\left(
\sum_x\mu_m(x)(\log\mu_m(x))^2
\;-\;
\Bigl(\sum_x\mu_m(x)\log\mu_m(x)\Bigr)^2
\right)
=
\frac1m\Var_{\mu_m}(\log\mu_m(X)).
\]
又因
\[
\log\mu_m(X)=\mathscr H_m-m\log 2,
\]
加常数不改变方差，故
\[
\Var_{\mu_m}(\log\mu_m(X))=\Var_{\mu_m}(\mathscr H_m).
\]
另一方面，
\[
\Phi_m''(1)=P_m''(1).
\]
遂得
\[
\frac1m\Var_{\mu_m}(\mathscr H_m)=P_m''(1).
\]

最后，由 \(C^2_{\mathrm{loc}}(I)\) 收敛，必有
\[
P_m'(1)\to P'(1),
\qquad
P_m''(1)\to P''(1).
\]
代入前述有限尺度恒等式即得极限公式。证毕。
\end{proof}

\leanverified{paper\_xi\_time\_part60ab\_hidden\_gap\_kl\_renyi2\_collision}
\begin{theorem}[隐藏位缺口的 KL 识别与 Rényi--\texorpdfstring{$2$}{2} 碰撞证书]\label{thm:xi-time-part60ab-hidden-gap-kl-renyi2-collision}
沿用表 \ref{tab:fold_multiplicity_histogram} 的隐藏位预算记号。记
\[
u_m(x):=\frac1{|X_m|}
\qquad (x\in X_m),
\]
并定义隐藏位、可见位与容量缺口
\[
h_m:=\EE_{\mu_m}\bigl[\log_2 d_m(X)\bigr],
\qquad
I_m:=m-h_m,
\qquad
\mathrm{gap}_m:=\log_2|X_m|-I_m.
\]
再定义二阶碰撞指纹
\[
C_{\mathrm{col}}(m):=
|X_m|\sum_{x\in X_m}\mu_m(x)^2.
\]
则有严格恒等式
\[
\boxed{\;
\mathrm{gap}_m
=
\frac{1}{\log 2}\,
D_{\mathrm{KL}}(\mu_m\|u_m),\;}
\]
以及精确上界
\[
\boxed{\;
\mathrm{gap}_m
\le
\log_2 C_{\mathrm{col}}(m).\;}
\]
并且下列三条等价：
\begin{enumerate}
  \item \(\mathrm{gap}_m=\log_2 C_{\mathrm{col}}(m)\)；
  \item \(\mu_m=u_m\)；
  \item \(d_m(x)\) 在 \(X_m\) 上常值。
\end{enumerate}
\end{theorem}
\begin{proof}
由定义，
\[
\mathrm{gap}_m
=
\log_2|X_m|-m+\sum_{x\in X_m}\mu_m(x)\log_2 d_m(x).
\]
再用
\[
\mu_m(x)=\frac{d_m(x)}{2^m},
\qquad
u_m(x)=\frac1{|X_m|},
\]
得到
\[
\mathrm{gap}_m
=
\sum_x \mu_m(x)\log_2\frac{d_m(x)|X_m|}{2^m}
=
\sum_x \mu_m(x)\log_2\frac{\mu_m(x)}{u_m(x)}
=
\frac{1}{\log 2}\,D_{\mathrm{KL}}(\mu_m\|u_m).
\]
这便给出第一条恒等式。

另一方面，Rényi 散度对阶数单调，故
\[
D_{\mathrm{KL}}(\mu_m\|u_m)\le D_2(\mu_m\|u_m).
\]
而
\[
D_2(\mu_m\|u_m)
=
\log\sum_x \frac{\mu_m(x)^2}{u_m(x)}
=
\log\!\Bigl(|X_m|\sum_x \mu_m(x)^2\Bigr)
=
\log C_{\mathrm{col}}(m).
\]
两边除以 \(\log 2\) 即得
\[
\mathrm{gap}_m\le \log_2 C_{\mathrm{col}}(m).
\]

最后，\(D_{\mathrm{KL}}(\mu_m\|u_m)=D_2(\mu_m\|u_m)\) 当且仅当对 \(\mu_m\)-几乎处处的 \(x\)，比值 \(\mu_m(x)/u_m(x)\) 为常数。由于 \(u_m\) 在 \(X_m\) 上严格正、且二者都是概率分布，这等价于 \(\mu_m=u_m\)。而
\[
\mu_m=u_m
\iff
\frac{d_m(x)}{2^m}=\frac1{|X_m|}
\ \ (\forall x\in X_m)
\iff
d_m(x)\ \text{在 }X_m\text{ 上常值}.
\]
证毕。
\end{proof}

\begin{theorem}[fold-gauge 常数列的逐阶极限消元反演给出全部偶 Bernoulli 数与偶值 \texorpdfstring{$\zeta(2r)$}{zeta(2r)}]\label{thm:xi-time-part60ab-gauge-constant-recursive-bernoulli-stripping}
\leanverified{paper\_xi\_time\_part60ab\_gauge\_constant\_recursive\_bernoulli\_stripping}
沿用定理 \ref{thm:fold-bin-gauge-constant-stirling-bernoulli-hierarchy} 的记号，记
\[
q:=\frac{\varphi}{2}\in(0,1),
\qquad
E_m:=\log C_m-\log(2\pi).
\]
并对每个整数 \(r\ge 1\) 定义
\[
a_r:=
\frac{B_{2r}}{r(2r-1)}
\bigl(\varphi^{-1}+\varphi^{2r-3}\bigr).
\]
再递归定义残差
\[
R_m^{(1)}:=E_m,
\qquad
R_m^{(r)}:=R_m^{(r-1)}-a_{r-1}q^{(2r-3)m}
\qquad (r\ge 2).
\]
则对每个 \(r\ge 1\) 都有极限
\[
\boxed{\;
\lim_{m\to\infty}q^{-(2r-1)m}R_m^{(r)}=a_r.\;}
\]
从而
\[
\boxed{\;
B_{2r}
=
\frac{r(2r-1)}{\varphi^{-1}+\varphi^{2r-3}}
\lim_{m\to\infty}
\Bigl(\frac{2}{\varphi}\Bigr)^{(2r-1)m}R_m^{(r)}.\;}
\]
\[
\boxed{\;
\zeta(2r)
=
(-1)^{r-1}\frac{(2\pi)^{2r}}{2(2r)!}\,
\frac{r(2r-1)}{\varphi^{-1}+\varphi^{2r-3}}
\lim_{m\to\infty}
\Bigl(\frac{2}{\varphi}\Bigr)^{(2r-1)m}R_m^{(r)}.\;}
\]
因此，单一标量序列 \((\log C_m)_{m\ge 1}\) 已足以逐阶恢复全部偶 Bernoulli 数与全部偶值 \(\zeta(2r)\)。
\end{theorem}
\begin{proof}
由定理 \ref{thm:fold-bin-gauge-constant-stirling-bernoulli-hierarchy}，对任意固定 \(R\ge 1\) 有
\[
E_m
=
\sum_{s=1}^{R}a_s q^{(2s-1)m}
+O\!\bigl(q^{(2R+1)m}\bigr).
\]
对上式取 \(R=r\)，并逐次减去前 \(r-1\) 个主项，得到
\[
R_m^{(r)}
=
a_r q^{(2r-1)m}+O\!\bigl(q^{(2r+1)m}\bigr).
\]
两边同乘 \(q^{-(2r-1)m}\) 得
\[
q^{-(2r-1)m}
R_m^{(r)}
=
a_r+O(q^{2m}),
\]
令 \(m\to\infty\) 即得极限公式。
再由
\[
a_r=
\frac{B_{2r}}{r(2r-1)}
\bigl(\varphi^{-1}+\varphi^{2r-3}\bigr)
\]
即可反解出 \(B_{2r}\)。
再由标准 Bernoulli--\(\zeta\) 恒等式
\[
B_{2r}
=
(-1)^{r-1}\frac{2(2r)!}{(2\pi)^{2r}}\zeta(2r),
\]
可立即反解出 \(\zeta(2r)\) 的显式恢复公式。证毕。
\end{proof}

\begin{corollary}[fold-gauge 常数的唯一主模与比例极限]\label{cor:xi-time-part60ab-gauge-constant-leading-mode-ratio}
\leanverified{paper\_xi\_time\_part60ab\_gauge\_constant\_leading\_mode\_ratio}
沿用定理 \ref{thm:xi-time-part60ab-gauge-constant-recursive-bernoulli-stripping} 的记号，则
\[
\boxed{\;
\log C_m-\log(2\pi)
=
a_1 q^m+O(q^{3m}),
\qquad
q=\frac{\varphi}{2}.\;}
\]
从而
\[
\boxed{\;
\lim_{m\to\infty}
\frac{\log C_{m+1}-\log(2\pi)}
{\log C_m-\log(2\pi)}
=
\frac{\varphi}{2},\;}
\]
以及
\[
\boxed{\;
\lim_{m\to\infty}
\Bigl(\frac{2}{\varphi}\Bigr)^m
\bigl(\log C_m-\log(2\pi)\bigr)
=
a_1.\;}
\]
\end{corollary}
\begin{proof}
由上一条定理在 \(r=1\) 时的系数公式，
\[
a_1=
\frac{B_2}{1\cdot 1}
\bigl(\varphi^{-1}+\varphi^{-1}\bigr)
=
\frac16\cdot 2\varphi^{-1}
=
\frac{1}{3\varphi}.
\]
故
\[
\log C_m-\log(2\pi)
=
a_1 q^m+O(q^{3m}).
\]
于是
\[
\frac{\log C_{m+1}-\log(2\pi)}
{\log C_m-\log(2\pi)}
=
\frac{\frac{1}{3\varphi}q^{m+1}+O(q^{3m+3})}
{\frac{1}{3\varphi}q^m+O(q^{3m})}
\longrightarrow q=\frac{\varphi}{2},
\]
并且
\[
\Bigl(\frac{2}{\varphi}\Bigr)^m\bigl(\log C_m-\log(2\pi)\bigr)
=
\frac{1}{3\varphi}+O(q^{2m})
\longrightarrow
\frac{1}{3\varphi}.
\]
证毕。
\end{proof}

在上述逐阶极限消元公式之后，\(\log C_m\) 的 Bernoulli--\(\zeta\) 谱链还可进一步压缩为以下三条结构后果。

\begin{theorem}[fold-gauge 常数的 Richardson--Neville 任意阶消模器与 \texorpdfstring{$(2R+1)$}{(2R+1)} 级加速]\label{thm:xi-time-part60ab-gauge-constant-richardson-neville-accelerator}
沿用定理 \ref{thm:xi-time-part60ab-gauge-constant-recursive-bernoulli-stripping} 的记号，并定义
\[
\lambda_r:=q^{2r-1}
\qquad (r\ge 1).
\]
对任意固定整数 \(R\ge 1\)，令
\[
Q_R(z):=\prod_{r=1}^{R}\frac{z-\lambda_r}{1-\lambda_r}
=
\sum_{j=0}^{R}\gamma_{R,j}z^j,
\]
并定义线性加速器
\[
\mathcal A_R(m):=\sum_{j=0}^{R}\gamma_{R,j}\log C_{m+j}.
\]
则有
\[
\boxed{\;
\mathcal A_R(m)=\log(2\pi)+O\!\bigl(q^{(2R+1)m}\bigr).\;}
\]
换言之，\(\mathcal A_R\) 精确消去 \(\log C_m\) 渐近展开中的前 \(R\) 个 Bernoulli 模态，并把逼近 \(\log(2\pi)\) 的误差阶提升到 \(q^{(2R+1)m}\)。
\end{theorem}
\begin{proof}
由定义立即有
\[
Q_R(1)=1,
\qquad
Q_R(\lambda_r)=0
\qquad (1\le r\le R).
\]
另一方面，由定理 \ref{thm:fold-bin-gauge-constant-stirling-bernoulli-hierarchy}，对固定 \(R\) 有
\[
\log C_{m+j}
=
\log(2\pi)+\sum_{r=1}^{R}a_r\lambda_r^{m+j}
+O\!\bigl(\lambda_{R+1}^{m+j}\bigr)
\qquad (0\le j\le R).
\]
由于 \(j\) 仅在有限集合 \(\{0,\dots,R\}\) 中变化，上式余项可统一写为 \(O(\lambda_{R+1}^m)\)。将其代入 \(\mathcal A_R(m)\)，得到
\[
\mathcal A_R(m)
=
\log(2\pi)\sum_{j=0}^{R}\gamma_{R,j}
+\sum_{r=1}^{R}a_r\lambda_r^m\sum_{j=0}^{R}\gamma_{R,j}\lambda_r^j
+O\!\bigl(\lambda_{R+1}^m\bigr).
\]
再由
\[
\sum_{j=0}^{R}\gamma_{R,j}=Q_R(1)=1,
\qquad
\sum_{j=0}^{R}\gamma_{R,j}\lambda_r^j=Q_R(\lambda_r)=0
\qquad (1\le r\le R),
\]
主项中的前 \(R\) 个 Bernoulli 模态全部被精确消去，从而
\[
\mathcal A_R(m)
=
\log(2\pi)+O\!\bigl(\lambda_{R+1}^m\bigr)
=
\log(2\pi)+O\!\bigl(q^{(2R+1)m}\bigr).
\]
证毕。
\end{proof}

\begin{theorem}[fold-gauge 误差序列的非 \texorpdfstring{$C$}{C}-有限性与普通生成函数的非有理性]\label{thm:xi-time-part60ab-gauge-constant-non-cfinite-ogf-obstruction}
\leanverified{paper\_xi\_time\_part60ab\_gauge\_constant\_non\_cfinite\_ogf\_obstruction}
令
\[
u_m:=\log C_m-\log(2\pi).
\]
则序列 \((u_m)_{m\ge 1}\) 不是任何常系数线性递推序列，即不是 \(C\)-有限序列。因而其普通生成函数
\[
U(z):=\sum_{m\ge 1}u_m z^m
\]
不可能是有理函数。
\end{theorem}
\begin{proof}
反设 \((u_m)_{m\ge 1}\) 为 \(C\)-有限序列。则存在有限个复数 \(\mu_1,\dots,\mu_s\) 与多项式 \(P_1,\dots,P_s\)，使得对充分大的 \(m\) 有
\[
u_m=\sum_{j=1}^{s}P_j(m)\mu_j^m.
\]
因此 \(u_m\) 只可能由有限多个指数底数 \(\mu_j\) 生成。

另一方面，由定理 \ref{thm:xi-time-part60ab-gauge-constant-recursive-bernoulli-stripping}，对每个 \(R\ge 1\) 都有
\[
u_m-\sum_{r=1}^{R-1}a_r q^{(2r-1)m}
=
a_R q^{(2R-1)m}+O\!\bigl(q^{(2R+1)m}\bigr).
\]
又由于
\[
a_R=
\frac{B_{2R}}{R(2R-1)}
\bigl(\varphi^{-1}+\varphi^{2R-3}\bigr)
\]
且 \(B_{2R}\neq 0\)，故 \(a_R\neq 0\) 对每个 \(R\ge 1\) 都成立。于是对每个 \(R\)，尺度 \(q^{(2R-1)m}\) 都真实出现在 \(u_m\) 的渐近链中。另一方面，
\[
u_m-\sum_{r=1}^{R-1}a_r q^{(2r-1)m}
\]
仍是由有限个指数底数
\[
\{\mu_1,\dots,\mu_s,q,q^3,\dots,q^{2R-3}\}
\]
生成的多项式--指数和。由其渐近主项为非零的 \(a_R q^{(2R-1)m}\)，可知底数 \(q^{2R-1}\) 必属于上式有限集合。由于
\[
q^{2R-1}\notin \{q,q^3,\dots,q^{2R-3}\},
\]
遂只能有
\[
q^{2R-1}\in \{\mu_1,\dots,\mu_s\}.
\]
这与 \(\{\mu_1,\dots,\mu_s\}\) 有限而
\[
q,q^3,q^5,\dots
\]
两两不同矛盾。故 \((u_m)\) 不是 \(C\)-有限序列。

最后，普通生成函数有理当且仅当其系数序列 \(C\)-有限，这是经典等价命题，故 \(U(z)\) 不可能为有理函数。证毕。
\end{proof}

\begin{corollary}[fold-gauge Bernoulli 塔的无限模态障碍]\label{cor:xi-time-part60ab-gauge-constant-infinite-modal-obstruction}
令
\[
E_m:=\log C_m-\log(2\pi).
\]
则不存在有限组互异复数 \(\lambda_1,\dots,\lambda_s\) 与多项式 \(P_1,\dots,P_s\)，使得对充分大的 \(m\) 恒有
\[
E_m=\sum_{j=1}^{s}P_j(m)\lambda_j^m.
\]
等价地，H.111 的 Stirling--Bernoulli 层级展开对应一条真正的无限模态谱链，而非任何有限维 companion 机制所能封装的有限指数叠加。
\end{corollary}
\begin{proof}
这正是定理 \ref{thm:xi-time-part60ab-gauge-constant-non-cfinite-ogf-obstruction} 的直接重述，因为常系数线性递推序列与有限个多项式--指数项的和完全等价。证毕。
\end{proof}

\begin{theorem}[高阶尾部精度强制前导 Bernoulli--\texorpdfstring{$\zeta$}{zeta} 模态一致]\label{thm:xi-time-part60ab-gauge-constant-tail-rigidity}
沿用定理 \ref{thm:xi-time-part60ab-gauge-constant-recursive-bernoulli-stripping} 的记号。设另一实序列 \((D_m)_{m\ge 1}\) 满足：对某个固定整数 \(R\ge 1\)，存在实系数 \(b_1,\dots,b_R\) 使
\[
D_m-\log(2\pi)
=
\sum_{r=1}^{R}b_r q^{(2r-1)m}
+O\!\bigl(q^{(2R+1)m}\bigr),
\]
并且同时有
\[
D_m-\log C_m
=
O\!\bigl(q^{(2R+1)m}\bigr).
\]
则必有
\[
b_r=a_r
\qquad (1\le r\le R).
\]
因而前 \(R\) 个偶 Bernoulli 数 \(B_{2r}\) 与偶值 \(\zeta(2r)\) 也被同一尾部精度唯一锁定。
\end{theorem}
\begin{proof}
由定理 \ref{thm:fold-bin-gauge-constant-stirling-bernoulli-hierarchy}，
\[
\log C_m-\log(2\pi)
=
\sum_{r=1}^{R}a_r q^{(2r-1)m}
+O\!\bigl(q^{(2R+1)m}\bigr).
\]
将其与 \(D_m\) 的展开相减，得到
\[
D_m-\log C_m
=
\sum_{r=1}^{R}(b_r-a_r)q^{(2r-1)m}
+O\!\bigl(q^{(2R+1)m}\bigr).
\]
若存在最小的 \(r_0\le R\) 使 \(b_{r_0}\neq a_{r_0}\)，则上式可改写为
\[
D_m-\log C_m
=
(b_{r_0}-a_{r_0})q^{(2r_0-1)m}
+O\!\bigl(q^{(2r_0+1)m}\bigr).
\]
两边同乘 \(q^{-(2r_0-1)m}\)，并令 \(m\to\infty\)，便得到
\[
q^{-(2r_0-1)m}\bigl(D_m-\log C_m\bigr)\to b_{r_0}-a_{r_0}\neq 0,
\]
这与假设
\[
D_m-\log C_m
=
O\!\bigl(q^{(2R+1)m}\bigr)
\]
矛盾，因为后者特别蕴含
\(
D_m-\log C_m=O(q^{(2r_0+1)m})
\)。
故 \(b_r=a_r\) 对全部 \(1\le r\le R\) 成立。

最后，由定理 \ref{thm:xi-time-part60ab-gauge-constant-recursive-bernoulli-stripping} 中 \(a_r\) 与 \(B_{2r}\)、\(\zeta(2r)\) 的显式对应关系，前 \(R\) 个 Bernoulli--\(\zeta\) 模态也随之被唯一锁定。证毕。
\end{proof}


\paragraph{Stirling--Bernoulli 模态的移位投影与偶值 \texorpdfstring{$\zeta$}{zeta} 显式恢复}
附录中的 Bernoulli 层级、递归抽取与有限阶移位消元已经给出三条彼此兼容的读出接口。把它们压回正文主线，可得到如下三条直接后果：每个 Stirling--Bernoulli 模态都可由一个显式有限阶移位投影子直接抽出；偶值 \(\zeta(2r)\) 因而成为单一规范常数序列上的直接可审计读出；而无限多个互异衰减模的同时存在，则排除了任何有限 companion 型封装。

\begin{theorem}[fold-gauge Bernoulli 模态的有限阶移位投影公式]\label{thm:xi-time-part60ab2-logcm-shift-projector}
沿用定理 \ref{thm:fold-bin-gauge-constant-stirling-bernoulli-hierarchy} 的记号，记
\[
f_m:=\log C_m-\log(2\pi),
\qquad
q_r:=\left(\frac{\varphi}{2}\right)^{2r-1},
\qquad
a_r:=\frac{B_{2r}}{r(2r-1)}\bigl(\varphi^{-1}+\varphi^{2r-3}\bigr).
\]
并令移位算子 \(E\) 作用于序列 \(f=(f_m)_{m\ge 1}\) 为
\[
(Ef)_m:=f_{m+1}.
\]
对每个 \(r\ge 1\) 定义
\[
\mathcal P_{r-1}(E):=\prod_{j=1}^{r-1}(E-q_j),
\]
其中 \(r=1\) 时取空积 \(\mathcal P_0(E)=1\)。则极限
\[
\lim_{m\to\infty}
\frac{q_r^{-m}\bigl(\mathcal P_{r-1}(E)f\bigr)_m}
{\prod_{j=1}^{r-1}(q_r-q_j)}
\]
存在，并且恰等于 \(a_r\)。特别地，每个模态系数 \(a_r\) 都可由 \((\log C_m)_{m\ge 1}\) 的一个有限阶移位投影极限直接读出。
\end{theorem}
\begin{proof}
定理 \ref{thm:derived-fold-gauge-shift-annihilator-even-zeta} 已给出
\[
\lim_{m\to\infty}q_r^{-m}\bigl(\mathcal P_{r-1}(E)f\bigr)_m
=
a_r\prod_{j=1}^{r-1}(q_r-q_j).
\]
又由于
\[
q_1>q_2>\cdots>0,
\]
故对每个 \(j<r\) 都有 \(q_r-q_j\neq 0\)。两边同除以该非零乘积，即得所示公式。证毕。
\end{proof}

\begin{corollary}[有限阶移位投影给出的偶值 \texorpdfstring{$\zeta(2r)$}{zeta(2r)} 显式恢复公式]\label{cor:xi-time-part60ab2-logcm-shift-even-zeta}
\leanverified{paper\_xi\_time\_part60ab2\_logcm\_shift\_even\_zeta}
沿用定理 \ref{thm:xi-time-part60ab2-logcm-shift-projector} 的记号。对每个 \(r\ge 1\)，都有
\[
\zeta(2r)
=
(-1)^{r-1}\frac{(2\pi)^{2r}}{2(2r)!}\,
\frac{r(2r-1)}{\varphi^{-1}+\varphi^{2r-3}}\,
\frac{
\lim\limits_{m\to\infty}
q_r^{-m}\bigl(\mathcal P_{r-1}(E)f\bigr)_m
}
{\prod_{j=1}^{r-1}(q_r-q_j)}.
\]
因此，偶值 \(\zeta\)-谱并不只作为 \((\log C_m)\) 的渐近尾项存在，而是可由一族显式有限阶移位投影子逐模态直接恢复。
\end{corollary}
\begin{proof}
由定理 \ref{thm:xi-time-part60ab2-logcm-shift-projector}，
\[
a_r=
\frac{
\lim\limits_{m\to\infty}
q_r^{-m}\bigl(\mathcal P_{r-1}(E)f\bigr)_m
}
{\prod_{j=1}^{r-1}(q_r-q_j)}.
\]
再代入
\[
a_r=\frac{B_{2r}}{r(2r-1)}\bigl(\varphi^{-1}+\varphi^{2r-3}\bigr),
\qquad
B_{2r}=(-1)^{r-1}\frac{2(2r)!}{(2\pi)^{2r}}\zeta(2r),
\]
整理即得。证毕。
\end{proof}

\begin{corollary}[归一化移位投影子的单步 Bernoulli 恢复]\label{cor:xi-time-part60ab2-logcm-normalized-projector-bernoulli}
\leanverified{paper\_xi\_time\_part60ab2\_logcm\_normalized\_projector\_bernoulli}
沿用定理 \ref{thm:xi-time-part60ab2-logcm-shift-projector} 的记号。对每个 \(r\ge 1\)，定义归一化移位投影子
\[
\Pi_r(E):=
\prod_{j=1}^{r-1}\frac{E-q_j}{q_r-q_j},
\]
其中 \(r=1\) 时取空积 \(\Pi_1(E)=1\)。
则
\[
\lim_{m\to\infty}q_r^{-m}\bigl(\Pi_r(E)f\bigr)_m=a_r,
\]
从而
\[
\lim_{m\to\infty}q_r^{-m}\bigl(\Pi_r(E)f\bigr)_m
=
\frac{B_{2r}}{r(2r-1)}\bigl(\varphi^{-1}+\varphi^{2r-3}\bigr),
\]
以及
\[
B_{2r}
=
\frac{r(2r-1)}{\varphi^{-1}+\varphi^{2r-3}}
\lim_{m\to\infty}q_r^{-m}\bigl(\Pi_r(E)f\bigr)_m.
\]
因此，偶 Bernoulli 数可由一族封闭形式的归一化移位投影子直接逐阶读出，而无需求助于递归剥离此前各级模态。
\end{corollary}
\begin{proof}
由定义，
\[
\Pi_r(E)
=
\frac{\mathcal P_{r-1}(E)}{\prod_{j=1}^{r-1}(q_r-q_j)}.
\]
故定理 \ref{thm:xi-time-part60ab2-logcm-shift-projector} 立即给出
\[
\lim_{m\to\infty}q_r^{-m}\bigl(\Pi_r(E)f\bigr)_m=a_r.
\]
再代入
\[
a_r=\frac{B_{2r}}{r(2r-1)}\bigl(\varphi^{-1}+\varphi^{2r-3}\bigr)
\]
并整理，即得所示 Bernoulli 恢复公式。证毕。
\end{proof}

\begin{theorem}[有限阶移位投影强制 fold-gauge 谱塔无限]\label{thm:xi-time-part60ab2-logcm-shift-infinite-modal}
序列
\[
f_m=\log C_m-\log(2\pi)
\]
不是任何有限阶常系数线性递推序列；等价地，其普通生成函数
\[
F(z):=\sum_{m\ge 1}f_m z^m
\]
不是有理函数。
\end{theorem}
\leanverified{paper\_xi\_time\_part60ab2\_logcm\_shift\_infinite\_modal}
\begin{proof}
由定理 \ref{thm:xi-time-part60ab2-logcm-shift-projector}，对每个 \(r\ge 1\) 都有
\[
\lim_{m\to\infty}
\frac{q_r^{-m}\bigl(\mathcal P_{r-1}(E)f\bigr)_m}
{\prod_{j=1}^{r-1}(q_r-q_j)}
=a_r.
\]
而
\[
a_r=
\frac{B_{2r}}{r(2r-1)}\bigl(\varphi^{-1}+\varphi^{2r-3}\bigr)\neq 0,
\]
因为偶 Bernoulli 数 \(B_{2r}\) 从不消失，且 \(\varphi^{-1}+\varphi^{2r-3}>0\)。因此，\(f_m\) 的渐近谱包含无限多个互异底数
\[
q_r=\left(\frac{\varphi}{2}\right)^{2r-1},
\qquad
r=1,2,\dots.
\]
另一方面，任何有限阶常系数线性递推序列的渐近谱都只能由有限多个指数底数构成，故不可能容纳上述无限模态塔。于是 \(f_m\) 非 \(C\)-有限。

普通生成函数有理当且仅当其系数序列 \(C\)-有限，这是经典等价命题，遂得 \(F(z)\) 非有理。证毕。
\end{proof}

\noindent
上述三条结果把 \((\log C_m)\) 的 Bernoulli 层级从“逐阶剥离尾项”的恢复程序，提升为一族显式有限阶移位滤波器的直接读出接口。因而，偶值 \(\zeta\)-谱并非只是在渐近展开中被动出现，而可由单一规范常数序列经过确定性的有限阶投影逐模态抽出；同时，无限模态障碍又说明这种读出结构本质上超出任何有限状态线性递推的封装能力。


\begin{theorem}[约数触发的谱零块强制群环多项式含有显式二项式因子]\label{thm:xi-time-part60ab-zero-block-forces-binomial-factor}
沿用命题 \ref{prop:fold-multiplicity-group-algebra} 的记号。记
\[
M:=F_{m+2},
\qquad
D_m(x):=\sum_{r=0}^{M-1}d_m(r)x^r\in\ZZ[x],
\qquad
\deg D_m<M.
\]
若 \(d\mid (m+2)\) 是真约数且满足
\[
F_d\mid \frac{M}{2},
\]
则有整系数因子结论
\[
\boxed{\;
x^{F_d}+1\mid D_m(x)\qquad\text{于 }\ZZ[x].\;}
\]
特别地，当
\[
m+2\equiv 0\pmod 6
\]
并取 \(d=(m+2)/2\) 时，
\[
\boxed{\;
x^{F_{(m+2)/2}}+1\mid D_m(x).\;}
\]
\end{theorem}
\begin{proof}
令
\[
t:=\frac{M}{2F_d}\in\NN.
\]
由推论 \ref{cor:fold-zero-divisor-lb}，
\[
S_{F_d}(M)\subseteq \mathcal Z_m
:=
\{k\in\ZZ/(M\ZZ):\ \widehat d_m(k)=0\}.
\]
对每个 \(\ell=0,1,\dots,F_d-1\)，取
\[
k_\ell:=(2\ell+1)t.
\]
则
\[
k_\ell F_d=(2\ell+1)\frac{M}{2}\equiv \frac{M}{2}\pmod M,
\]
故 \(k_\ell\in S_{F_d}(M)\subseteq \mathcal Z_m\)。于是由离散 Fourier 变换的定义，
\[
D_m(\omega_M^{-k_\ell})=\widehat d_m(k_\ell)=0,
\qquad
\omega_M:=e^{2\pi i/M}.
\]

另一方面，
\[
(\omega_M^{-k_\ell})^{F_d}
=
\exp\!\left(-2\pi i\,\frac{k_\ell F_d}{M}\right)
=
\exp\!\left(-(2\ell+1)\pi i\right)
=
-1.
\]
故每个 \(\omega_M^{-k_\ell}\) 都是 \(x^{F_d}+1\) 的根。并且
\[
\omega_M^{-k_\ell}
=
\exp\!\left(-\frac{(2\ell+1)\pi i}{F_d}\right)
\qquad (0\le \ell\le F_d-1)
\]
恰给出 \(x^{F_d}+1\) 的全部 \(F_d\) 个互异复根。因此 \(x^{F_d}+1\) 的全部根都是 \(D_m(x)\) 的根，从而
\[
x^{F_d}+1\mid D_m(x)
\qquad\text{于 }\CC[x].
\]
由于二者都属于 \(\ZZ[x]\)，Gauss 引理进一步给出
\[
x^{F_d}+1\mid D_m(x)
\qquad\text{于 }\ZZ[x].
\]

当 \(m+2\equiv 0\pmod 6\) 时，推论 \ref{cor:fold-zero-half-index-multiple6} 恰给出
\[
F_{(m+2)/2}\mid \frac{F_{m+2}}{2},
\]
故特殊情形立即成立。证毕。
\end{proof}

\noindent
上述链条把表 \ref{tab:fold_multiplicity_histogram} 中的隐藏位、容量缺口与碰撞指纹严格并入统一的信息几何结构；同时又把 \(\log C_m\) 的渐近模态提升为一条逐阶可反演、可任意阶消模并伴随非 \(C\)-有限障碍与尾部刚性的 Bernoulli--\(\zeta\) 谱链，并把约数驱动的谱零块最终刚化为群环多项式上的显式整系数因子。

\paragraph{bin-fold 群论三重分解、中心态熵同一、规范体积常数项与零余类算术放大}
在纤维置换分解、群胚代数的 Wedderburn 直和、输出 Rényi 熵的常数项闭式以及零余类的约数触发机制已经建立之后，还可把同一套退化直方图进一步压缩为下列六条对象层结论。

\begin{theorem}[bin-fold 规范群的中心--交换化--导群三重分解]\label{thm:xi-foldbin-gauge-group-center-abel-derived-triple-decomposition}
沿用定义 \ref{def:fold-bin-gauge-group} 与定义 \ref{def:fold-bin-degeneracy-sectors} 的记号，记
\[
G_m^{\mathrm{bin}}:=\Aut(\Fold_m^{\mathrm{bin}}),
\qquad
\mathcal S_{m,d}:=\{w\in X_m:\ d_m^{\mathrm{bin}}(w)=d\}.
\]
则有自然同构
\[
Z\!\bigl(G_m^{\mathrm{bin}}\bigr)\cong (\ZZ/2\ZZ)^{|\mathcal S_{m,2}|},
\]
\[
\bigl(G_m^{\mathrm{bin}}\bigr)^{\mathrm{ab}}
\cong
(\ZZ/2\ZZ)^{\sum_{d\ge 2}|\mathcal S_{m,d}|},
\]
\[
[G_m^{\mathrm{bin}},G_m^{\mathrm{bin}}]
\cong
\prod_{d\ge 3}\mathfrak A_d^{\,|\mathcal S_{m,d}|}.
\]
特别地，若 \(d_{\min}(m)\ge 2\)，则
\[
\bigl(G_m^{\mathrm{bin}}\bigr)^{\mathrm{ab}}
\cong
(\ZZ/2\ZZ)^{|X_m|}.
\]
对审计偶数窗
\(
m\in\{6,8,10,12\}
\)，进一步有
\[
\bigl(G_m^{\mathrm{bin}}\bigr)^{\mathrm{ab}}
\cong
(\ZZ/2\ZZ)^{F_{m+2}}.
\]
\end{theorem}
\begin{proof}
由命题 \ref{prop:fold-bin-gauge-decomposition}，
\[
G_m^{\mathrm{bin}}
\cong
\prod_{w\in X_m}\mathfrak S_{d_m^{\mathrm{bin}}(w)}.
\]
有限直积对中心、交换化与导群分别可按因子逐项计算，而对称群满足
\[
Z(\mathfrak S_d)\cong
\begin{cases}
\ZZ/2\ZZ,& d=2,\\
1,& d\neq 2,
\end{cases}
\]
\[
\mathfrak S_d^{\mathrm{ab}}\cong
\begin{cases}
1,& d=1,\\
\ZZ/2\ZZ,& d\ge 2,
\end{cases}
\qquad
[\mathfrak S_d,\mathfrak S_d]\cong
\begin{cases}
1,& d\le 2,\\
\mathfrak A_d,& d\ge 3.
\end{cases}
\]
把这些公式逐纤维代入，便得到
\[
Z\!\bigl(G_m^{\mathrm{bin}}\bigr)\cong (\ZZ/2\ZZ)^{|\mathcal S_{m,2}|},
\]
\[
\bigl(G_m^{\mathrm{bin}}\bigr)^{\mathrm{ab}}
\cong
(\ZZ/2\ZZ)^{\sum_{d\ge 2}|\mathcal S_{m,d}|},
\]
\[
[G_m^{\mathrm{bin}},G_m^{\mathrm{bin}}]
\cong
\prod_{d\ge 3}\mathfrak A_d^{\,|\mathcal S_{m,d}|}.
\]

若 \(d_{\min}(m)\ge 2\)，则每条纤维都贡献一个 \(\ZZ/2\ZZ\) 因子，故
\[
\sum_{d\ge 2}|\mathcal S_{m,d}|=|X_m|,
\]
从而
\[
\bigl(G_m^{\mathrm{bin}}\bigr)^{\mathrm{ab}}\cong (\ZZ/2\ZZ)^{|X_m|}.
\]
对 \(m\in\{6,8,10,12\}\)，定理 \ref{thm:fold-bin-min-degeneracy-fib-audited} 给出
\[
d_{\min}(m)=F_{m/2}\ge 2,
\]
再由稳定类型计数律 \(|X_m|=F_{m+2}\)，即得最后一式。证毕。
\end{proof}

\begin{theorem}[群胚代数中心态的 Rényi--Shannon 熵与输出熵的严格同一性]\label{thm:xi-foldbin-groupoid-central-renyi-shannon-identity}
沿用命题 \ref{prop:fold-bin-groupoid-algebra-wedderburn} 与定理 \ref{thm:xi-foldbin-fkdet-logvolume-kappa} 的记号，记
\[
A_m:=\CC[\mathcal R_m^{\mathrm{bin}}]
\cong
\bigoplus_{w\in X_m}M_{d_m^{\mathrm{bin}}(w)}(\CC),
\qquad
Z(A_m)=\bigoplus_{w\in X_m}\CC z_w,
\]
其中 \(z_w\) 为第 \(w\)-块的最小中心幂等元，并在 \(A_m\) 上赋予概率迹
\[
\tau_m\!\left(\bigoplus_{w\in X_m}a_w\right):=
2^{-m}\sum_{w\in X_m}\operatorname{Tr}(a_w).
\]
则
\[
\tau_m(z_w)=\frac{d_m^{\mathrm{bin}}(w)}{2^m}=:\mu_m(w),
\]
故中心极小幂等元的谱权分布恰与 bin-fold 输出分布 \(\mu_m\) 一致。

对 \(q>0\)、\(q\neq 1\)，定义中心 Rényi 熵
\[
H_q^{\mathrm{cen}}(m):=
\frac{1}{1-q}\log\sum_{w\in X_m}\tau_m(z_w)^q,
\]
并把 \(q=1\) 情形按连续延拓记为
\[
H_1^{\mathrm{cen}}(m):=
-\sum_{w\in X_m}\tau_m(z_w)\log\tau_m(z_w).
\]
则对一切 \(q>0\) 都有
\[
H_q^{\mathrm{cen}}(m)=H_q(\mu_m).
\]
特别地，对 \(q>0\)、\(q\neq 1\)，有
\[
H_q^{\mathrm{cen}}(m)
=
m\log\varphi
+\frac{1}{1-q}\log\!\Bigl(\frac{\varphi+\varphi^{-q}}{\sqrt5}\Bigr)
+O\!\Bigl(\Bigl(\frac{\varphi}{2}\Bigr)^m\Bigr),
\]
并且
\[
H_1^{\mathrm{cen}}(m)
=
m\log\varphi
+\frac{\log\varphi}{1+\varphi^2}
+O\!\Bigl(\Bigl(\frac{\varphi}{2}\Bigr)^m\Bigr)
=
\log|X_m|-C_{\mathrm{KL}}+o(1),
\]
其中
\[
C_{\mathrm{KL}}
=
2\log\varphi-\frac12\log 5-\frac{\log\varphi}{\varphi\sqrt5}.
\]
从而
\[
\frac{1}{m}H_q^{\mathrm{cen}}(m)\longrightarrow \log\varphi
\qquad (q>0).
\]
\end{theorem}
\begin{proof}
在 Wedderburn 分解下，\(z_w\) 正是第 \(w\)-块上的单位元与其余各块上的零元，故
\[
\tau_m(z_w)
=
2^{-m}\operatorname{Tr}(I_{d_m^{\mathrm{bin}}(w)})
=
2^{-m}d_m^{\mathrm{bin}}(w)
=
\mu_m(w).
\]
因此中心原子 \(\{z_w\}_{w\in X_m}\) 的谱权分布恰为 \(\mu_m\)，从而
\[
H_q^{\mathrm{cen}}(m)=H_q(\mu_m)
\]
对一切 \(q>0\) 都成立。

当 \(q\neq 1\) 时，把该恒等式代入定理 \ref{thm:xi-foldbin-absolute-renyi-entropy-constant-closed} 即得 Rényi 常数项闭式；\(q=1\) 的精确常数项同样由该定理给出。再由命题 \ref{prop:fold-bin-kl-constant} 的
\[
H(\mu_m)=\log|X_m|-C_{\mathrm{KL}}+o(1)
\]
便得到 Shannon 形式的第二种写法。最后，
\[
\frac1m H_q^{\mathrm{cen}}(m)\to \log\varphi
\]
可由上述显式展开直接读出，亦可视为定理 \ref{thm:fold-bin-renyi-rate-collapse} 的中心态版本。证毕。
\end{proof}

\begin{theorem}[规范体积密度的精确热力学常数项]\label{thm:xi-foldbin-gauge-density-exact-thermodynamic-constant}
沿用定理 \ref{thm:xi-foldbin-groupoid-central-renyi-shannon-identity} 的记号。则有
\[
g_m
=
m\log\!\Bigl(\frac{2}{\varphi}\Bigr)
-1
-\frac{\log\varphi}{1+\varphi^2}
+O\!\Bigl(m\Bigl(\frac{\varphi}{2}\Bigr)^m\Bigr).
\]
从而
\[
\lim_{m\to\infty}
\left(
g_m-m\log\!\Bigl(\frac{2}{\varphi}\Bigr)
\right)
=
-1-\frac{\log\varphi}{1+\varphi^2}.
\]
\end{theorem}
\begin{proof}
由定理 \ref{thm:xi-foldbin-groupoid-central-renyi-shannon-identity} 在 \(q=1\) 的结论，
\[
H(\mu_m)=H_1^{\mathrm{cen}}(m)
=
m\log\varphi+\frac{\log\varphi}{1+\varphi^2}
+O\!\Bigl(\Bigl(\frac{\varphi}{2}\Bigr)^m\Bigr).
\]
另一方面，定理 \ref{thm:xi-foldbin-gauge-entropy-one-nat-law} 给出
\[
g_m
=
m\log 2-H(\mu_m)-1
+O\!\Bigl(m\Bigl(\frac{\varphi}{2}\Bigr)^m\Bigr).
\]
将前式代入即得
\[
g_m
=
m\log\!\Bigl(\frac{2}{\varphi}\Bigr)
-1
-\frac{\log\varphi}{1+\varphi^2}
+O\!\Bigl(m\Bigl(\frac{\varphi}{2}\Bigr)^m\Bigr),
\]
并立刻推出极限常数。证毕。
\end{proof}
\leanverified{paper\_xi\_foldbin\_gauge\_density\_exact\_thermodynamic\_constant}

\begin{theorem}[无限算术子序列上的零余类放大律]\label{thm:xi-fold-zero-arithmetic-subsequence-amplification}
令
\[
M:=F_{m+2},
\qquad
\mathcal Z_m:=\{k\in\ZZ/(M\ZZ):\ \widehat d_m(k)=0\},
\qquad
\bar d_m:=\frac{2^m}{F_{m+2}}.
\]
并沿用定理 \ref{thm:fold-zero-uncertainty} 的记号
\[
p_m(r):=\frac{d_m(r)}{2^m},
\qquad
u_m(r):=\frac1M.
\]
若
\[
m\equiv 4\pmod 6,
\]
则
\[
|\mathcal Z_m|\ge F_{(m+2)/2}.
\]
因而
\[
\mathsf{Col}_m
\le
1-\frac{F_{(m+2)/2}}{F_{m+2}},
\qquad
D_{\mathrm{KL}}(p_m\|u_m)
\le
\log\!\bigl(F_{m+2}-F_{(m+2)/2}\bigr),
\]
\[
M_m-\bar d_m
\le
2^m
\sqrt{
\frac{F_{m+2}-1-F_{(m+2)/2}}{F_{m+2}}
}.
\]
此外，
\[
\frac{F_{(m+2)/2}}{F_{m+2}}=\Theta(\varphi^{-m/2}),
\]
故零余类所强制的碰撞削减项在该无限算术子序列上具有 \(\Theta(\varphi^{-m/2})\) 的确定性尺度。
\leanverified{paper\_xi\_fold\_zero\_arithmetic\_subsequence\_amplification}
\end{theorem}
\begin{proof}
若 \(m\equiv 4\pmod 6\)，则 \(m+2\equiv 0\pmod 6\)。由推论 \ref{cor:fold-zero-half-index-multiple6} 立即得到
\[
|\mathcal Z_m|\ge F_{(m+2)/2}.
\]
把该下界分别代入定理 \ref{thm:fold-collision-zero-reduction} 与定理 \ref{thm:fold-zero-uncertainty} 的第 \(2\)、\(3\) 条，便得到所示碰撞上界、KL 上界以及最大纤维偏离上界。

最后，由 Binet 公式
\[
F_n=\frac{\varphi^n-(-\varphi^{-1})^n}{\sqrt5}
\]
可知
\[
\frac{F_{(m+2)/2}}{F_{m+2}}
=
\Theta\!\bigl(\varphi^{-(m+2)/2}\bigr)
=
\Theta(\varphi^{-m/2}),
\]
故最后一条结论成立。证毕。
\end{proof}

\begin{theorem}[六倍窗零块以平方根尺度注入稳定层，并对碰撞与相对熵给出同阶改进]\label{thm:xi-time-part60ab-window6-zero-block-halfscale-information-improvement}
沿用定理 \ref{thm:xi-fold-zero-multiple6-information-gap-improvement} 的记号。若
\[
m+2\equiv 0\pmod 6,
\]
并记
\[
d:=\frac{m+2}{2},
\qquad
|X_m|=F_{m+2},
\]
则有
\[
\frac{|\mathcal Z_m|}{|X_m|}
\ge
\frac{F_d}{F_{m+2}}
=
\varphi^{-d}\bigl(1+O(\varphi^{-2d})\bigr).
\]
特别地，
\[
|\mathcal Z_m|\asymp |X_m|^{1/2}.
\]
进一步，
\[
1-\mathsf{Col}_m\ge \frac{F_d}{F_{m+2}}
=
\varphi^{-d}\bigl(1+o(1)\bigr),
\]
并且
\[
\log|X_m|-D_{\mathrm{KL}}(p_m\|u_m)
\ge
\log\frac{F_{m+2}}{F_{m+2}-F_d}
=
\varphi^{-d}+O(\varphi^{-2d}).
\]
\end{theorem}
\begin{proof}
由定理 \ref{thm:xi-fold-zero-multiple6-information-gap-improvement}，
\[
|\mathcal Z_m|\ge F_d,
\qquad
\mathsf{Col}_m\le 1-\frac{F_d}{F_{m+2}},
\qquad
D_{\mathrm{KL}}(p_m\|u_m)\le \log(F_{m+2}-F_d).
\]
两边分别除以
\[
|X_m|=F_{m+2},
\]
便得
\[
\frac{|\mathcal Z_m|}{|X_m|}\ge \frac{F_d}{F_{m+2}},
\qquad
1-\mathsf{Col}_m\ge \frac{F_d}{F_{m+2}},
\]
以及
\[
\log|X_m|-D_{\mathrm{KL}}(p_m\|u_m)
\ge
\log\frac{F_{m+2}}{F_{m+2}-F_d}.
\]

再由 Binet 公式，
\[
F_n=\frac{\varphi^n-(-\varphi^{-1})^n}{\sqrt5},
\]
可知
\[
\frac{F_d}{F_{m+2}}
=
\varphi^{-d}\bigl(1+O(\varphi^{-2d})\bigr).
\]
由于
\[
|X_m|=F_{m+2}\asymp \varphi^{m+2}=\varphi^{2d},
\]
于是
\[
F_d\asymp \varphi^d\asymp |X_m|^{1/2},
\]
从而得到
\[
|\mathcal Z_m|\asymp |X_m|^{1/2}.
\]

最后，令
\[
x:=\frac{F_d}{F_{m+2}},
\]
则 \(x\to 0\)，并且
\[
\log\frac{F_{m+2}}{F_{m+2}-F_d}
=
-\log(1-x)
=
x+O(x^2)
=
\varphi^{-d}+O(\varphi^{-2d}).
\]
证毕。
\end{proof}

\begin{theorem}[dyadic 零余类塔触发的碰撞、KL、峰值与最大纤维偏离同步强化]\label{thm:xi-fold-zero-dyadic-tower-synchronous-information-gap-improvement}\label{thm:xi-fold-zero-multiple6-information-gap-improvement}
\leanverified{paper\_xi\_fold\_zero\_dyadic\_tower\_synchronous\_information\_gap\_improvement}
沿用定理 \ref{thm:xi-fold-zero-arithmetic-subsequence-amplification} 与定理 \ref{thm:fold-zero-uncertainty} 的记号。若
\[
m+2=3\cdot 2^r\cdot s,
\qquad
r\ge 1,
\qquad
s\ \text{为奇数},
\]
并定义
\[
\Xi_m:=\sum_{j=1}^{r}F_{(m+2)/2^j},
\]
则
\[
|\mathcal Z_m|\ge \Xi_m.
\]
因而
\[
\mathsf{Col}_m
\le
1-\frac{\Xi_m}{F_{m+2}},
\qquad
D_{\mathrm{KL}}(p_m\|u_m)
\le
\log\!\bigl(F_{m+2}-\Xi_m\bigr),
\]
\[
S_m\le 2^m\sqrt{F_{m+2}-1-\Xi_m},
\]
\[
M_m-\bar d_m
\le
2^m\sqrt{\frac{F_{m+2}-1-\Xi_m}{F_{m+2}}}.
\]
特别地，当 \(m=58\) 时，
\[
\Xi_{58}=F_{30}+F_{15}=832650,
\]
故上述四条界已经严格强于仅取半阶零块 \(F_{30}=832040\) 时得到的对应上界；而若再与推论 \ref{cor:xi-fold-zero-m58-three-coset-exhaustion} 的精确三余类分解合并，则相同四条界还可进一步提升到由 \(|\mathcal Z_{58}|=839415\) 所决定的更强数值。
\end{theorem}
\begin{proof}
由定理 \ref{thm:xi-fold-zero-dyadic-tower-disjoint-fibonacci-cosets}，
\[
|\mathcal Z_m|\ge \Xi_m.
\]
将该下界分别代入定理 \ref{thm:fold-collision-zero-reduction} 与定理
\ref{thm:fold-zero-uncertainty} 的第 \(2\)、\(3\)、\(4\) 条，即得所示碰撞上界、KL 上界、Fourier 峰值上界以及最大纤维偏离上界。

当 \(m=58\) 时，
\[
\Xi_{58}=F_{30}+F_{15}=832650
\]
由直接代入即得。最后一句只是把推论
\ref{cor:xi-fold-zero-m58-three-coset-exhaustion} 的精确值
\[
|\mathcal Z_{58}|=839415
\]
再次代入同四条不等式。证毕。
\end{proof}

\begin{proposition}[规范体积与奇偶异常的 Schur 横截分离]\label{prop:xi-time-part60abb-gauge-volume-parity-transverse-separation}
\leanverified{paper\_xi\_time\_part60abb\_gauge\_volume\_parity\_transverse\_separation}
固定分辨率 \(m\)，并设
\[
d=(d_x)_{x\in X_m},
\qquad
d'=(d_x')_{x\in X_m}
\]
是两组正整数退化谱，满足
\[
\sum_{x\in X_m}d_x=\sum_{x\in X_m}d_x'=2^m.
\]
定义
\[
G(d):=\prod_{x\in X_m}\mathfrak S_{d_x},
\qquad
H_{\mathrm{gauge}}(d):=\sum_{x\in X_m}\log(d_x!),
\qquad
N_2(d):=\dim_{\FF_2}G(d)^{\mathrm{ab}}.
\]
设两组谱具有相同的非平凡支撑
\[
\Sigma:=\{x\in X_m:\ d_x\ge 2\}=\{x\in X_m:\ d_x'\ge 2\},
\]
并且 \(d'|_\Sigma\) 在共同支撑上严格 majorize \(d|_\Sigma\)。则有
\[
N_2(d')=N_2(d)=|\Sigma|,
\qquad
H_{\mathrm{gauge}}(d')>H_{\mathrm{gauge}}(d).
\]
因此，模二奇偶异常只记录“哪些纤维是非平凡的”，而规范体积则会严格响应这些非平凡纤维内部的凸性重分配；二者在 Schur 方向上横截而不可互代。
\end{proposition}
\begin{proof}
由对称群的交换化公式
\[
\mathfrak S_n^{\mathrm{ab}}
\cong
\begin{cases}
1,& n=1,\\
\ZZ/2\ZZ,& n\ge 2,
\end{cases}
\]
立得
\[
G(d)^{\mathrm{ab}}
\cong
(\ZZ/2\ZZ)^{|\Sigma|},
\qquad
G(d')^{\mathrm{ab}}
\cong
(\ZZ/2\ZZ)^{|\Sigma|}.
\]
故
\[
N_2(d)=N_2(d')=|\Sigma|.
\]

再看规范体积。由于共同支撑外的坐标都等于 \(1\)，它们对
\[
H_{\mathrm{gauge}}(\cdot)=\sum_x\log(d_x!)
\]
没有任何差异性贡献。于是只需在 \(\Sigma\) 上比较
\[
\sum_{x\in\Sigma}\log(d_x!)
\qquad\text{与}\qquad
\sum_{x\in\Sigma}\log(d_x'!).
\]
令
\[
f(t):=\log\Gamma(t+1)
\qquad (t\ge 1).
\]
则
\[
f''(t)=\psi_1(t+1)>0,
\]
故 \(f\) 在 \([1,\infty)\) 上严格凸。由 Karamata 不等式，既然 \(d'|_\Sigma\) 严格 majorize \(d|_\Sigma\)，便有
\[
\sum_{x\in\Sigma}f(d_x')>\sum_{x\in\Sigma}f(d_x).
\]
而 \(f(n)=\log(n!)\) 对整数 \(n\ge 1\) 成立，故
\[
H_{\mathrm{gauge}}(d')>H_{\mathrm{gauge}}(d).
\]
证毕。
\end{proof}


\noindent
上述各条结论把 bin-fold 的退化直方图同时提升到群论骨架、群胚中心态、规范热力学与 Schur 凸性横截四个层面，并把 fold 主线上的零余类约数触发机制收束为一条完整的算术子序列规律：退化谱不仅决定 \(G_m^{\mathrm{bin}}\) 的中心、交换化与导群，也决定中心极小幂等元的谱权与规范体积密度的零温偏移；另一方面，\(m\equiv 4\pmod 6\) 不仅给出一条统一的零余类放大序列，而且其 dyadic 零块塔的总量 \(\Xi_m\) 还会线性传递到碰撞、KL、Fourier 峰值与最大纤维偏离的四重上界，而对同一退化支撑上的 Schur 重分配，模二奇偶异常仅记录非平凡支撑本身，规范体积却会严格响应凸性集中度，从而进一步把“群论账本”与“体积账本”的横截不可比性固定下来。

\paragraph{稳定层均匀纤维谱、精确尺寸偏置与中心纤维尺寸算子的两迹谱测度}
在稳定层均匀基线下的缩放退化度二原子极限、微态推前分布的两点极限以及 bin-fold 纤维群胚代数的 Wedderburn 直和都已建立之后，还可把这三条链进一步收束为一个统一的中心谱对象：均匀输出纤维谱给出基线两原子测度，微态推前恰是这一谱的精确尺寸偏置，而同一中心纤维尺寸算子在两条自然迹下的谱测度正好分别实现这两类极限。

\begin{theorem}[bin-fold 在稳定层均匀基线下的二原子谱与固定阶矩闭式]\label{thm:xi-time-part60ab4-uniform-fiber-spectrum-twoatom-moment-closure}
\leanverified{paper\_xi\_time\_part60ab4\_uniform\_fiber\_spectrum\_twoatom\_moment\_closure}
记
\[
u_m(w):=\frac{1}{|X_m|}
\qquad (w\in X_m),
\qquad
U_m\sim u_m,
\]
\[
Y_m:=\left(\frac{\varphi}{2}\right)^m d_m^{\mathrm{bin}}(U_m).
\]
则 \(Y_m\) 依分布收敛到两点随机变量
\[
Y=
\begin{cases}
1, & \text{以概率 } \varphi^{-1},\\[3pt]
\varphi^{-1}, & \text{以概率 } \varphi^{-2}.
\end{cases}
\]
并且对任意固定实数 \(q\ge 0\)，都有
\[
\lim_{m\to\infty}\mathbb E[Y_m^q]
=
\frac{1}{\varphi}+\varphi^{-q-2}.
\]
特别地，若记
\[
S_q^{\mathrm{bin}}(m):=\sum_{w\in X_m}\bigl(d_m^{\mathrm{bin}}(w)\bigr)^q,
\]
则有
\[
S_q^{\mathrm{bin}}(m)
\sim
\frac{\varphi+\varphi^{-q}}{\sqrt5}
\left(\frac{2^q}{\varphi^{q-1}}\right)^m.
\]
\end{theorem}
\begin{proof}
第一条正是定理 \ref{thm:xi-time-part70a-uniform-scaled-degeneracy-quantitative} 的极限形式。对任意固定 \(q\ge 0\)，推论 \ref{cor:xi-time-part70a-scaled-moments-leading-constant} 给出
\[
\lim_{m\to\infty}\mathbb E[Y_m^q]
=
\lim_{m\to\infty}
\frac{1}{|X_m|}
\sum_{w\in X_m}
\left(
\left(\frac{\varphi}{2}\right)^m d_m^{\mathrm{bin}}(w)
\right)^q
=
\frac{1}{\varphi}+\varphi^{-q-2}.
\]
再由
\[
\mathbb E[Y_m^q]
=
\left(\frac{\varphi}{2}\right)^{mq}\frac{S_q^{\mathrm{bin}}(m)}{|X_m|}
\]
以及
\[
|X_m|=F_{m+2}\sim \frac{\varphi^{m+2}}{\sqrt5},
\]
即可整理得到
\[
S_q^{\mathrm{bin}}(m)
\sim
\frac{\varphi+\varphi^{-q}}{\sqrt5}
\left(\frac{2^q}{\varphi^{q-1}}\right)^m.
\]
证毕。
\end{proof}

\begin{theorem}[微态推前谱是均匀基线纤维谱的精确尺寸偏置变换]\label{thm:xi-time-part60ab4-exact-sizebias-pushforward-law}
\leanverified{paper\_xi\_time\_part60ab4\_exact\_sizebias\_pushforward\_law}
记
\[
\mu_m(w):=\frac{d_m^{\mathrm{bin}}(w)}{2^m}
\qquad (w\in X_m),
\qquad
W_m\sim \mu_m,
\]
\[
Z_m:=\left(\frac{\varphi}{2}\right)^m d_m^{\mathrm{bin}}(W_m).
\]
则对任意有界函数 \(f:\RR_{\ge 0}\to\CC\)，都有精确恒等式
\[
\mathbb E[f(Z_m)]
=
\frac{\mathbb E[Y_m f(Y_m)]}{\mathbb E[Y_m]}.
\]
等价地，对任意有界函数 \(g:\RR_{\ge 0}\to\CC\)，都有
\[
\mathbb E_{\mu_m}\!\bigl[g(d_m^{\mathrm{bin}}(W_m))\bigr]
=
\frac{
\mathbb E_{u_m}\!\bigl[d_m^{\mathrm{bin}}(U_m)\,g(d_m^{\mathrm{bin}}(U_m))\bigr]
}{
\mathbb E_{u_m}\!\bigl[d_m^{\mathrm{bin}}(U_m)\bigr]
}.
\]
从而 \(Z_m\Longrightarrow Z\)，其中
\[
Z=
\begin{cases}
1, & \text{以概率 } \dfrac{\varphi}{\sqrt5},\\[6pt]
\varphi^{-1}, & \text{以概率 } \dfrac{1}{\varphi\sqrt5},
\end{cases}
\]
并且 \(Z\) 正好是 \(Y\) 的尺寸偏置变换：
\[
\mathbb P(Z\in A)
=
\frac{\mathbb E\!\bigl[Y\,\mathbf 1_{\{Y\in A\}}\bigr]}{\mathbb E[Y]}
\qquad
(A\subseteq \RR_{\ge 0}\ \text{Borel}).
\]
特别地，
\[
\mathbb E[Y]
=
\frac{1}{\varphi}+\varphi^{-3}
=
\frac{\sqrt5}{\varphi^2},
\]
且
\[
\mathbb P(Z=1)=\frac{\varphi}{\sqrt5},
\qquad
\mathbb P(Z=\varphi^{-1})=\frac{1}{\varphi\sqrt5}.
\]
\end{theorem}
\begin{proof}
由定义，
\[
\mathbb E[f(Z_m)]
=
\sum_{w\in X_m}\frac{d_m^{\mathrm{bin}}(w)}{2^m}
f\!\left(\left(\frac{\varphi}{2}\right)^m d_m^{\mathrm{bin}}(w)\right).
\]
另一方面，
\[
\mathbb E[Y_m f(Y_m)]
=
\frac{1}{|X_m|}
\sum_{w\in X_m}
\left(\frac{\varphi}{2}\right)^m
d_m^{\mathrm{bin}}(w)
f\!\left(\left(\frac{\varphi}{2}\right)^m d_m^{\mathrm{bin}}(w)\right),
\]
而纤维分割恒等式
\[
\sum_{w\in X_m}d_m^{\mathrm{bin}}(w)=2^m
\]
给出
\[
\mathbb E[Y_m]
=
\left(\frac{\varphi}{2}\right)^m\frac{2^m}{|X_m|}.
\]
因此
\[
\frac{\mathbb E[Y_m f(Y_m)]}{\mathbb E[Y_m]}
=
\sum_{w\in X_m}\frac{d_m^{\mathrm{bin}}(w)}{2^m}
f\!\left(\left(\frac{\varphi}{2}\right)^m d_m^{\mathrm{bin}}(w)\right)
=
\mathbb E[f(Z_m)],
\]
从而第一条精确恒等式成立。把 \(g(t):=f((\varphi/2)^m t)\) 代回，即得未缩放形式。

再由定理 \ref{thm:xi-time-part70a-uniform-scaled-degeneracy-quantitative} 的一致近似可知，当 \(m\) 充分大时，\(Y_m\) 几乎处处落在固定紧区间
\[
\left[\frac{\varphi^{-1}}{2},2\right]
\]
之内。于是对任意有界连续函数 \(f\)，函数 \(y\mapsto y f(y)\) 在该区间上仍为有界连续函数。结合定理 \ref{thm:xi-time-part60ab4-uniform-fiber-spectrum-twoatom-moment-closure}，得到
\[
\mathbb E[f(Z_m)]
=
\frac{\mathbb E[Y_m f(Y_m)]}{\mathbb E[Y_m]}
\longrightarrow
\frac{\mathbb E[Y f(Y)]}{\mathbb E[Y]}.
\]
故 \(Z_m\) 的极限分布就是 \(Y\) 的尺寸偏置变换。

最后，由
\[
\mathbb E[Y]
=
1\cdot \varphi^{-1}
+\varphi^{-1}\cdot \varphi^{-2}
=
\frac{\sqrt5}{\varphi^2}
\]
以及两点尺寸偏置公式
\[
\mathbb P(Z=1)
=
\frac{1\cdot \varphi^{-1}}{\mathbb E[Y]}
=
\frac{\varphi}{\sqrt5},
\qquad
\mathbb P(Z=\varphi^{-1})
=
\frac{\varphi^{-1}\cdot \varphi^{-2}}{\mathbb E[Y]}
=
\frac{1}{\varphi\sqrt5},
\]
便得到所示闭式。证毕。
\end{proof}

\begin{theorem}[中心纤维尺寸算子在两条自然迹下的谱测度识别]\label{thm:xi-time-part60ab4-central-size-operator-two-trace-spectral-identification}
沿用定义 \ref{def:fold-bin-fiber-groupoid}、命题 \ref{prop:fold-bin-groupoid-algebra-wedderburn} 与定理 \ref{thm:xi-foldbin-groupoid-central-renyi-shannon-identity} 的记号，记
\[
A_m:=\CC[\mathcal R_m^{\mathrm{bin}}]
\cong
\bigoplus_{w\in X_m}M_{d_m^{\mathrm{bin}}(w)}(\CC).
\]
定义中心纤维尺寸算子
\[
\Lambda_m:=
\bigoplus_{w\in X_m}
d_m^{\mathrm{bin}}(w)\,I_{d_m^{\mathrm{bin}}(w)}
\in Z(A_m),
\]
以及其尺度化版本
\[
\widetilde\Lambda_m:=
\left(\frac{\varphi}{2}\right)^m\Lambda_m.
\]
再定义两条迹
\[
\tau_m^{\mathrm{cen}}\!\left(\bigoplus_{w\in X_m}a_w\right)
:=
\frac{1}{|X_m|}
\sum_{w\in X_m}\operatorname{tr}_{d_m^{\mathrm{bin}}(w)}(a_w),
\]
\[
\tau_m^{\mathrm{reg}}\!\left(\bigoplus_{w\in X_m}a_w\right)
:=
\frac{1}{2^m}
\sum_{w\in X_m}\operatorname{Tr}(a_w),
\]
其中 \(\operatorname{tr}_d\) 为归一化矩阵迹；\(\tau_m^{\mathrm{reg}}\) 即定理 \ref{thm:xi-foldbin-groupoid-central-renyi-shannon-identity} 中的概率迹 \(\tau_m\)。则在最小中心幂等元 \(z_w\) 上有
\[
\tau_m^{\mathrm{cen}}(z_w)=u_m(w)=\frac{1}{|X_m|},
\qquad
\tau_m^{\mathrm{reg}}(z_w)=\mu_m(w)=\frac{d_m^{\mathrm{bin}}(w)}{2^m}.
\]
并且对任意有界函数 \(f:\RR_{\ge 0}\to\CC\)，都有精确恒等式
\[
\tau_m^{\mathrm{reg}}\!\bigl(f(\Lambda_m)\bigr)
=
\frac{
\tau_m^{\mathrm{cen}}\!\bigl(\Lambda_m f(\Lambda_m)\bigr)
}{
\tau_m^{\mathrm{cen}}(\Lambda_m)
}.
\]
此外，
\[
\tau_m^{\mathrm{cen}}\!\bigl(f(\widetilde\Lambda_m)\bigr)
=
\mathbb E[f(Y_m)],
\qquad
\tau_m^{\mathrm{reg}}\!\bigl(f(\widetilde\Lambda_m)\bigr)
=
\mathbb E[f(Z_m)].
\]
因此 \(\widetilde\Lambda_m\) 在 \(\tau_m^{\mathrm{cen}}\) 下的谱测度收敛到 \(Y\)，而在 \(\tau_m^{\mathrm{reg}}\) 下的谱测度收敛到 \(Z\)。
\end{theorem}
\begin{proof}
在第 \(w\)-块上，\(\Lambda_m\) 只是标量
\[
d_m^{\mathrm{bin}}(w)\,I_{d_m^{\mathrm{bin}}(w)}.
\]
因此对任意有界函数 \(f\)，都有
\[
f(\Lambda_m)
=
\bigoplus_{w\in X_m}
f\!\bigl(d_m^{\mathrm{bin}}(w)\bigr)\,I_{d_m^{\mathrm{bin}}(w)}.
\]
由此立即得到
\[
\tau_m^{\mathrm{cen}}(z_w)
=
\frac{1}{|X_m|},
\qquad
\tau_m^{\mathrm{reg}}(z_w)
=
\frac{d_m^{\mathrm{bin}}(w)}{2^m}.
\]
同理，
\[
\tau_m^{\mathrm{cen}}\!\bigl(f(\Lambda_m)\bigr)
=
\frac{1}{|X_m|}
\sum_{w\in X_m}
f\!\bigl(d_m^{\mathrm{bin}}(w)\bigr),
\]
\[
\tau_m^{\mathrm{cen}}\!\bigl(\Lambda_m f(\Lambda_m)\bigr)
=
\frac{1}{|X_m|}
\sum_{w\in X_m}
d_m^{\mathrm{bin}}(w)\,
f\!\bigl(d_m^{\mathrm{bin}}(w)\bigr),
\]
\[
\tau_m^{\mathrm{cen}}(\Lambda_m)
=
\frac{1}{|X_m|}
\sum_{w\in X_m}d_m^{\mathrm{bin}}(w)
=
\frac{2^m}{|X_m|},
\]
\[
\tau_m^{\mathrm{reg}}\!\bigl(f(\Lambda_m)\bigr)
=
\frac{1}{2^m}
\sum_{w\in X_m}
d_m^{\mathrm{bin}}(w)\,
f\!\bigl(d_m^{\mathrm{bin}}(w)\bigr).
\]
比较后便得到
\[
\tau_m^{\mathrm{reg}}\!\bigl(f(\Lambda_m)\bigr)
=
\frac{
\tau_m^{\mathrm{cen}}\!\bigl(\Lambda_m f(\Lambda_m)\bigr)
}{
\tau_m^{\mathrm{cen}}(\Lambda_m)
}.
\]

再由 \(\widetilde\Lambda_m=((\varphi/2)^m)\Lambda_m\)，可知
\[
\tau_m^{\mathrm{cen}}\!\bigl(f(\widetilde\Lambda_m)\bigr)
=
\frac{1}{|X_m|}
\sum_{w\in X_m}
f\!\left(\left(\frac{\varphi}{2}\right)^m d_m^{\mathrm{bin}}(w)\right)
=
\mathbb E[f(Y_m)],
\]
\[
\tau_m^{\mathrm{reg}}\!\bigl(f(\widetilde\Lambda_m)\bigr)
=
\frac{1}{2^m}
\sum_{w\in X_m}
d_m^{\mathrm{bin}}(w)\,
f\!\left(\left(\frac{\varphi}{2}\right)^m d_m^{\mathrm{bin}}(w)\right)
=
\mathbb E[f(Z_m)].
\]
最后，分别应用定理 \ref{thm:xi-time-part60ab4-uniform-fiber-spectrum-twoatom-moment-closure} 与定理 \ref{thm:xi-time-part60ab4-exact-sizebias-pushforward-law}，便得到两条谱测度收敛。证毕。
\end{proof}
\leanverified{paper\_xi\_time\_part60ab4\_central\_size\_operator\_two\_trace\_spectral\_identification}

\begin{theorem}[中心纤维尺寸算子的全阶矩与群胚代数维数主常数]\label{thm:xi-time-part60ab4-central-size-schatten-moment-dimension-constant}
\leanverified{paper\_xi\_time\_part60ab4\_central\_size\_schatten\_moment\_dimension\_constant}
在定理 \ref{thm:xi-time-part60ab4-central-size-operator-two-trace-spectral-identification} 的设定下，对任意固定实数 \(q\ge 0\)，都有
\[
\lim_{m\to\infty}
\tau_m^{\mathrm{cen}}\!\bigl(\widetilde\Lambda_m^q\bigr)
=
\frac{1}{\varphi}+\varphi^{-q-2},
\]
\[
\lim_{m\to\infty}
\tau_m^{\mathrm{reg}}\!\bigl(\widetilde\Lambda_m^q\bigr)
=
\frac{\varphi+\varphi^{-q-1}}{\sqrt5}.
\]
特别地，
\[
\lim_{m\to\infty}\tau_m^{\mathrm{cen}}(\widetilde\Lambda_m)
=
\frac{\sqrt5}{\varphi^2},
\qquad
\lim_{m\to\infty}\tau_m^{\mathrm{reg}}(\widetilde\Lambda_m)
=
\frac{2}{\sqrt5}.
\]
并且群胚代数维数满足精确主项
\[
\dim_{\CC}A_m
\sim
\frac{2}{\sqrt5}
\left(\frac{4}{\varphi}\right)^m.
\]
\end{theorem}
\begin{proof}
由定理 \ref{thm:xi-time-part60ab4-central-size-operator-two-trace-spectral-identification}，
\[
\tau_m^{\mathrm{cen}}\!\bigl(\widetilde\Lambda_m^q\bigr)
=
\mathbb E[Y_m^q],
\qquad
\tau_m^{\mathrm{reg}}\!\bigl(\widetilde\Lambda_m^q\bigr)
=
\mathbb E[Z_m^q].
\]
第一条极限于是直接由定理 \ref{thm:xi-time-part60ab4-uniform-fiber-spectrum-twoatom-moment-closure} 给出。

再由定理 \ref{thm:xi-time-part60ab4-exact-sizebias-pushforward-law} 的精确尺寸偏置恒等式，
\[
\mathbb E[Z_m^q]
=
\frac{\mathbb E[Y_m^{q+1}]}{\mathbb E[Y_m]}.
\]
令 \(m\to\infty\)，并代入
\[
\mathbb E[Y^{q+1}]
=
\frac{1}{\varphi}+\varphi^{-q-3},
\qquad
\mathbb E[Y]
=
\frac{\sqrt5}{\varphi^2},
\]
便得到
\[
\lim_{m\to\infty}\tau_m^{\mathrm{reg}}\!\bigl(\widetilde\Lambda_m^q\bigr)
=
\frac{\varphi+\varphi^{-q-1}}{\sqrt5}.
\]
取 \(q=1\) 即得
\[
\lim_{m\to\infty}\tau_m^{\mathrm{cen}}(\widetilde\Lambda_m)
=
\frac{\sqrt5}{\varphi^2},
\qquad
\lim_{m\to\infty}\tau_m^{\mathrm{reg}}(\widetilde\Lambda_m)
=
\frac{2}{\sqrt5}.
\]

另一方面，在 Wedderburn 分解下有
\[
\tau_m^{\mathrm{reg}}(\Lambda_m)
=
\frac{1}{2^m}
\sum_{w\in X_m}
\operatorname{Tr}\!\bigl(
d_m^{\mathrm{bin}}(w)\,I_{d_m^{\mathrm{bin}}(w)}
\bigr)
=
\frac{1}{2^m}
\sum_{w\in X_m}\bigl(d_m^{\mathrm{bin}}(w)\bigr)^2
=
\frac{\dim_{\CC}A_m}{2^m}.
\]
由于
\[
\widetilde\Lambda_m
=
\left(\frac{\varphi}{2}\right)^m\Lambda_m,
\]
故
\[
\tau_m^{\mathrm{reg}}(\Lambda_m)
=
\left(\frac{2}{\varphi}\right)^m
\tau_m^{\mathrm{reg}}(\widetilde\Lambda_m).
\]
于是
\[
\dim_{\CC}A_m
=
2^m\tau_m^{\mathrm{reg}}(\Lambda_m)
=
\left(\frac{4}{\varphi}\right)^m
\tau_m^{\mathrm{reg}}(\widetilde\Lambda_m),
\]
再代入
\[
\tau_m^{\mathrm{reg}}(\widetilde\Lambda_m)\longrightarrow \frac{2}{\sqrt5},
\]
便得到
\[
\dim_{\CC}A_m
\sim
\frac{2}{\sqrt5}
\left(\frac{4}{\varphi}\right)^m.
\]
证毕。
\end{proof}

因此，bin-fold 的稳定层均匀谱 \(Y\)、其尺寸偏置像 \(Z\) 与纤维群胚代数中的中心纤维尺寸算子 \(\widetilde\Lambda_m\) 之间并不存在三条彼此独立的极限律：前二者正是后一对象在中心均匀迹与概率迹下的两类谱测度，而 \(\dim_{\CC}A_m\) 的主常数 \(2/\sqrt5\) 则恰好等于该概率谱的一阶矩。
\paragraph{稳定层自反中心、Fourier 相位锁定与 boundary 中心的超选择分裂}
在 bin-fold 稳定层的群环生成函数、Fourier 乘积分解、规范群熵账本以及 window-\(6\) 的 boundary 中心直和结构都已建立之后，中心投影型读出与保纤维相位型读出还可继续压缩为同一条代数骨架：稳定层多重度谱不仅具有严格的 Fibonacci 自反中心，而且其 Fourier 相位被该中心完全锁定；与此同时，规范群 bits 的全局极值由严格离散凸性唯一决定，而 boundary sheet-flip 在中心层上则强制出现八个不可相混的超选择扇区。

\begin{theorem}[稳定层多重度谱的自反中心]\label{thm:xi-time-part60ab2-stable-spectrum-selfreciprocal-center}
\leanverified{paper\_xi\_time\_part60ab2\_stable\_spectrum\_selfreciprocal\_center}
记
\[
M:=F_{m+2},
\qquad
D_m(x):=\sum_{r=0}^{M-1}d_m(r)x^r\in \CC[x]/(x^M-1),
\]
并定义
\[
\Sigma_m:=\sum_{j=1}^{m}F_{j+1}=F_{m+3}-2\equiv F_{m+1}-2\pmod M.
\]
则在 \(\CC[x]/(x^M-1)\) 中成立严格自反恒等式
\[
D_m(x)\equiv x^{\Sigma_m}D_m(x^{-1})\pmod{x^M-1}.
\]
等价地，
\[
d_m(r)=d_m(\Sigma_m-r)
\qquad (\forall r\in \ZZ/(M\ZZ)).
\]
\end{theorem}
\begin{proof}
由命题 \ref{prop:fold-multiplicity-group-algebra}，
\[
D_m(x)\equiv \prod_{j=1}^{m}\bigl(1+x^{F_{j+1}}\bigr)\pmod{x^M-1}.
\]
于是
\[
\begin{aligned}
x^{\Sigma_m}D_m(x^{-1})
&\equiv
x^{\sum_{j=1}^{m}F_{j+1}}
\prod_{j=1}^{m}\bigl(1+x^{-F_{j+1}}\bigr)\\
&=
\prod_{j=1}^{m}\bigl(x^{F_{j+1}}+1\bigr)
\equiv
D_m(x)
\pmod{x^M-1}.
\end{aligned}
\]
比较模 \(x^M-1\) 下的系数便得
\[
d_m(r)=d_m(\Sigma_m-r).
\]
证毕。
\end{proof}

\begin{theorem}[Fourier 相位的刚性锁定]\label{thm:xi-time-part60ab2-stable-spectrum-fourier-phase-locking}
\leanverified{paper\_xi\_time\_part60ab2\_stable\_spectrum\_fourier\_phase\_locking}
沿用定理 \ref{thm:xi-time-part60ab2-stable-spectrum-selfreciprocal-center} 的记号。对每个
\[
k\in \ZZ/(M\ZZ)
\]
定义 Fourier 模
\[
\widehat d_m(k):=\sum_{r=0}^{M-1}d_m(r)e^{-2\pi i kr/M}.
\]
则对所有 \(k\) 都有
\[
\widehat d_m(k)=e^{-2\pi i k\Sigma_m/M}\,\overline{\widehat d_m(k)}.
\]
从而
\[
e^{\pi i k\Sigma_m/M}\widehat d_m(k)\in \RR.
\]
亦即每一个 Fourier 模的相位都被 \(\Sigma_m\) 强制锁定到一条固定直径上，而不再保留一般的复平面旋转自由度。
\end{theorem}
\begin{proof}
由定理 \ref{thm:xi-time-part60ab2-stable-spectrum-selfreciprocal-center}，
\[
d_m(r)=d_m(\Sigma_m-r).
\]
故
\[
\begin{aligned}
\widehat d_m(k)
&=
\sum_{r=0}^{M-1}d_m(\Sigma_m-r)e^{-2\pi i kr/M}\\
&=
e^{-2\pi i k\Sigma_m/M}
\sum_{r=0}^{M-1}d_m(r)e^{2\pi i kr/M}\\
&=
e^{-2\pi i k\Sigma_m/M}\,\overline{\widehat d_m(k)},
\end{aligned}
\]
其中最后一步用到 \(d_m(r)\in \RR\)。于是
\[
e^{\pi i k\Sigma_m/M}\widehat d_m(k)
=
\overline{e^{\pi i k\Sigma_m/M}\widehat d_m(k)},
\]
故该量为实数。证毕。
\end{proof}

\begin{theorem}[精确暗模的算术判别与模 \texorpdfstring{$3$}{3} 出现律]\label{thm:xi-time-part60ab2-exact-dark-mode-arithmetic-criterion}
\leanverified{paper\_xi\_time\_part60ab2\_exact\_dark\_mode\_arithmetic\_criterion}
沿用前述记号。则对任意 \(k\in \ZZ/(M\ZZ)\)，有
\[
\widehat d_m(k)=0
\]
当且仅当 \(M\) 为偶数且存在
\[
j\in\{1,\dots,m\}
\]
使
\[
kF_{j+1}\equiv \frac{M}{2}\pmod M.
\]
因此：
\begin{enumerate}
  \item 若 \(M=F_{m+2}\) 为奇数，则不存在任何精确暗模；
  \item 若 \(M\) 为偶数，则
  \[
  \widehat d_m(M/2)=0;
  \]
  \item 由 Fibonacci 奇偶周期为 \(3\)，有
  \[
  F_{m+2}\ \text{为偶数}
  \iff
  m\equiv 1\pmod 3,
  \]
  从而
  \[
  m\equiv 1\pmod 3
  \iff
  \widehat d_m(M/2)=0.
  \]
\end{enumerate}
\end{theorem}
\begin{proof}
由定理 \ref{thm:fold-spectrum-factorization}，
\[
\widehat d_m(k)=\prod_{j=1}^{m}\left(1+e^{-2\pi i kF_{j+1}/M}\right).
\]
故 \(\widehat d_m(k)=0\) 当且仅当某一因子为零，即
\[
1+e^{-2\pi i kF_{j+1}/M}=0
\]
对某个 \(j\) 成立。该式等价于
\[
e^{-2\pi i kF_{j+1}/M}=-1.
\]
而 \(-1\) 作为 \(M\)-次单位根出现当且仅当 \(M\) 为偶数；在此情形下，上式又等价于
\[
kF_{j+1}\equiv \frac{M}{2}\pmod M.
\]
这就证明了首条判别。

若 \(M\) 为偶数，则取 \(j=1\) 且 \(F_2=1\)，立得
\[
\frac{M}{2}F_2\equiv \frac{M}{2}\pmod M,
\]
故 \(\widehat d_m(M/2)=0\)。

最后，由 Fibonacci 数列模 \(2\) 的周期为 \(3\)，\(F_n\) 为偶数当且仅当 \(3\mid n\)。代入 \(n=m+2\) 即得
\[
F_{m+2}\ \text{为偶数}
\iff
3\mid (m+2)
\iff
m\equiv 1\pmod 3.
\]
证毕。
\end{proof}

\begin{theorem}[规范群 bits 泛函的全局极小值与极大值]\label{thm:xi-time-part60ab2-gauge-bits-global-extrema}
\leanverified{paper\_xi\_time\_part60ab2\_gauge\_bits\_global\_extrema}
令
\[
N:=2^m,
\qquad
L:=|X_m|=F_{m+2}.
\]
在所有满足
\[
(d_1,\dots,d_L)\in \NN^{L},
\qquad
\sum_{i=1}^{L}d_i=N
\]
的正整数向量上，定义
\[
\mathcal H(d_1,\dots,d_L):=\sum_{i=1}^{L}\log_2(d_i!).
\]
若记
\[
N=qL+r,
\qquad
0\le r<L,
\]
则 \(\mathcal H\) 的全局极小值与极大值分别为
\[
\mathcal H_{\min}
=(L-r)\log_2(q!)+r\log_2((q+1)!),
\]
\[
\mathcal H_{\max}
=
\log_2((N-L+1)!).
\]
并且：
\begin{enumerate}
  \item 极小分布在置换意义下唯一，恰为
  \[
  (d_i)\sim
  (\underbrace{q,\dots,q}_{L-r},
  \underbrace{q+1,\dots,q+1}_{r});
  \]
  \item 极大分布在置换意义下唯一，恰为
  \[
  (d_i)\sim (N-L+1,1,\dots,1).
  \]
\end{enumerate}
特别地，对 bin-fold 的实际纤维多重度谱 \(\bigl(d_m^{\mathrm{bin}}(w)\bigr)_{w\in X_m}\)，其规范群 bits
\[
H_m^{\mathrm{gauge}}=\sum_{w\in X_m}\log_2\bigl(d_m^{\mathrm{bin}}(w)!\bigr)
\]
必落在上述两端显式包络之间。
\end{theorem}
\begin{proof}
令
\[
f(n):=\log_2(n!)
\qquad (n\in\NN).
\]
则
\[
f(n+1)-f(n)=\log_2(n+1)
\]
严格递增，因此 \(f\) 在正整数上严格离散凸。

先证极小值。若某一向量中存在 \(a\ge b+2\)，则
\[
f(a)+f(b)-f(a-1)-f(b+1)
=
\log_2 a-\log_2(b+1)>0.
\]
故把一单位质量从较大坐标移向较小坐标会严格减小 \(\mathcal H\)。反复执行这一调平操作，直至所有坐标之差至多为 \(1\)，便得到唯一可能的极小构型，即全部坐标只取
\[
q,\qquad q+1
\]
两值，其中恰有 \(r\) 个取 \(q+1\)。代入即得 \(\mathcal H_{\min}\)。

再证极大值。若存在两个坐标 \(a\ge b\ge 2\)，则
\[
f(a+1)+f(b-1)-f(a)-f(b)
=
\log_2(a+1)-\log_2 b>0.
\]
故把一单位质量从较小的非平凡坐标移向较大的坐标会严格增大 \(\mathcal H\)。反复执行该集中操作，直到至多只有一个坐标大于 \(1\)，便得到唯一可能的极大构型
\[
(N-L+1,1,\dots,1),
\]
其函数值正是
\[
\log_2((N-L+1)!).
\]
证毕。
\end{proof}

\leanverified{paper\_xi\_time\_part60ab2\_window6\_boundary\_eight\_superselection}
\begin{theorem}[window-\texorpdfstring{$6$}{6} boundary 中心的八重超选择分解]\label{thm:xi-time-part60ab2-window6-boundary-eight-superselection}
取 \(m=6\)，并记
\[
X_6^{\mathrm{bdry}}=\{100001,100101,101001\}.
\]
设
\[
C_{\mathrm{bdry}}
:=
\langle \tau_w:\ w\in X_6^{\mathrm{bdry}}\rangle
=
Z_{\partial}
\cong
(\ZZ/2\ZZ)^3
\subset Z(G_6^{\mathrm{bin}}).
\]
对每个 \(\chi\in \widehat{C_{\mathrm{bdry}}}\)，定义
\[
e_\chi:=2^{-3}\sum_{g\in C_{\mathrm{bdry}}}\chi(g)\,g
\in \CC[C_{\mathrm{bdry}}].
\]
则
\[
\CC[C_{\mathrm{bdry}}]\cong \CC^{8}
\]
具有八个两两正交的极小中心幂等元 \(e_\chi\)。进一步，对任意在 \(\CC[C_{\mathrm{bdry}}]\) 上忠实的 \(\CC[G_6^{\mathrm{bin}}]\)-模 \(V\)，都有规范直和分解
\[
V=\bigoplus_{\chi\in \widehat{C_{\mathrm{bdry}}}}V_\chi,
\qquad
V_\chi:=e_\chi V,
\]
并且所有 \(V_\chi\) 都非零。换言之，boundary 中心强制出现八个非空的超选择扇区。
\end{theorem}
\begin{proof}
命题 \ref{prop:xi-window6-boundary-eight-sector-central-decomposition} 已给出 \(\{e_\chi\}_{\chi\in\widehat{C_{\mathrm{bdry}}}}\) 的两两正交性、极小中心幂等元性质以及
\[
\sum_{\chi\in\widehat{C_{\mathrm{bdry}}}}e_\chi=1.
\]
因此对任意 \(\CC[G_6^{\mathrm{bin}}]\)-模 \(V\)，都有
\[
V=\left(\sum_{\chi}e_\chi\right)V=\sum_{\chi}e_\chi V,
\]
而由 \(e_\chi e_{\chi'}=0\)（\(\chi\neq\chi'\)）知此和实为直和，故
\[
V=\bigoplus_{\chi\in\widehat{C_{\mathrm{bdry}}}}V_\chi.
\]

现在只需证明忠实性蕴含每个分量非零。若对某个 \(\chi\) 有
\[
V_\chi=e_\chi V=0,
\]
则 \(e_\chi\) 在 \(V\) 上作用为零。由于 \(e_\chi\neq 0\) 且 \(e_\chi\in \CC[C_{\mathrm{bdry}}]\)，这说明 \(\CC[C_{\mathrm{bdry}}]\to \End_{\CC}(V)\) 的限制映射具有非零核，与其忠实性矛盾。故所有 \(V_\chi\) 都非零。证毕。
\end{proof}

\leanverified{paper\_xi\_time\_part60ab2\_window6\_central\_bulk\_parity\_splitting\_ratio}
\begin{corollary}[window-\texorpdfstring{$6$}{6} parity-charge 的中心--体相分裂比例]\label{cor:xi-time-part60ab2-window6-central-bulk-parity-splitting-ratio}
在交换化口径下有规范直和分解
\[
\bigl(G_6^{\mathrm{bin}}\bigr)^{\mathrm{ab}}
\cong
(\ZZ/2\ZZ)^3\oplus (\ZZ/2\ZZ)^{18},
\]
其中第一因子恰由
\[
100001,\qquad 100101,\qquad 101001
\]
三条 boundary parity 所生成的中心像给出，而第二因子由非中心 bulk parity 组成。特别地，window-\(6\) 的 parity-charge 渠道中，中心稳健者所占比例恰为
\[
\frac{3}{21}=\frac17,
\]
其余
\[
\frac{18}{21}=\frac67
\]
都属于非中心通道。
\end{corollary}
\begin{proof}
定理 \ref{thm:xi-window6-boundary-parity-canonical-direct-summand} 已给出规范直和分解
\[
\bigl(G_6^{\mathrm{bin}}\bigr)^{\mathrm{ab}}
\cong
C_{\partial}\oplus (\ZZ/2\ZZ)^{18}
\cong
(\ZZ/2\ZZ)^3\oplus(\ZZ/2\ZZ)^{18},
\]
其中
\[
C_{\partial}=(\ZZ/2\ZZ)^{X_6^{\mathrm{bdry}}}
\]
正是 boundary 三条中心 sheet-flip 在交换化中的像。于是三维因子给出全部中心稳健 parity-channel，余下十八维因子给出非中心 bulk parity-channel。比例公式随之立即成立。证毕。
\end{proof}

\noindent
由此，稳定层的代数读出与 window-\(6\) 的中心结构便被压缩到同一组可核验对象之中：一方面，多重度剖面的自反中心把可见条纹的对称轴固定为显式的 Fibonacci 余类，Fourier 相位则被该中心锁定到一条一维直径；另一方面，规范群 bits 的极值几何由严格离散凸性完全支配，而 boundary sheet-flip 在中心层上又把保纤维相位型读出严格分解为八个彼此不可相混的超选择扇区。

\paragraph{window-\texorpdfstring{$6$}{6} 自由商障碍、boundary 中心残余与有限加法账本的不可能性}
在 window-\(6\) 的纤维退化多值性、boundary 中心直和结构以及有限秩加法账本障碍已经确立之后，还可把“全局自由对称解释失败”与“局域中心残余仍然存在”压缩为如下两条直接后果。

\begin{theorem}[自由对称解释与有限加法账本的联合不可能性]\label{thm:xi-time-part60ab3-free-symmetry-finite-ledger-impossible}
不存在任何三元组 \((H,L,\Phi)\) 同时满足下列条件：
\begin{enumerate}
  \item \(H\) 为有限群，并且对每个 \(m\ge 1\)，都存在一个 \(H\) 在
  \[
  \Omega_m:=\{0,1\}^m
  \]
  上的自由作用，使其轨道划分恰等于 \(\Fold_m^{\mathrm{bin}}\) 的纤维划分；
  \item \(L\) 为有限生成交换群；
  \item \(\Phi:(\NN_{>0},\times)\to(\ZZ\oplus L,+)\) 为单射幺半群同态。
\end{enumerate}
\end{theorem}
\begin{proof}
若第 \(1\) 条成立，则取 \(m=6\) 即得到 \(\Fold_6^{\mathrm{bin}}\) 的纤维划分可由某个有限群的自由轨道划分实现。然而推论 \ref{cor:foldbin6-degeneracy-multivalued-nonfree} 已证明这在 window-\(6\) 上不可能，因为退化度直方图
\[
2:8,\qquad 3:4,\qquad 4:9
\]
具有多值轨道大小，不能来自任何自由有限群作用。

另一方面，第 \(2\) 条与第 \(3\) 条的组合又被定理 \ref{thm:killo-godel-compression-not-finite-rank-homologizable} 直接排除：不存在有限生成交换群 \(L\) 使
\[
\Phi:(\NN_{>0},\times)\to(\ZZ\oplus L,+)
\]
成为单射幺半群同态。故上述三元组不可能存在。证毕。
\end{proof}

\begin{proposition}[window-\texorpdfstring{$6$}{6} 的 boundary 中心残余]\label{prop:xi-time-part60ab3-window6-boundary-center-residual}
\leanverified{paper\_xi\_time\_part60ab3\_window6\_boundary\_center\_residual}
记
\[
X_6^{\mathrm{bdry}}=\{100001,100101,101001\}.
\]
则这三条 boundary 词都对应二点纤维 \(d_6^{\mathrm{bin}}(w)=2\)。对每个 \(w\in X_6^{\mathrm{bdry}}\)，记 \(\tau_w\) 为该二点纤维上的唯一非平凡对合。则
\[
C_{\mathrm{bdry}}
:=
\langle \tau_w:\ w\in X_6^{\mathrm{bdry}}\rangle
\cong
(\ZZ/2\ZZ)^3
\hookrightarrow
Z(G_6^{\mathrm{bin}}).
\]
并且在交换化映射下，
\[
C_{\mathrm{bdry}}\longrightarrow \bigl(G_6^{\mathrm{bin}}\bigr)^{\mathrm{ab}}
\]
的像构成一个规范三维坐标子格，亦即 boundary parity 的中心残余在 abelianization 中仍保持独立可辨。
\end{proposition}
\begin{proof}
推论 \ref{cor:fold-bin-boundary-center-z2} 已把三条 boundary 二点纤维上的 sheet-flip 识别为规范中心嵌入
\[
(\ZZ/2\ZZ)^3\hookrightarrow Z(G_6^{\mathrm{bin}}),
\]
其中三条生成元正对应于 \(X_6^{\mathrm{bdry}}\) 中的三个 \(w\)。这就给出第一条结论。

再由定理 \ref{thm:xi-window6-boundary-parity-canonical-direct-summand}，
\[
\bigl(G_6^{\mathrm{bin}}\bigr)^{\mathrm{ab}}
\cong
(\ZZ/2\ZZ)^3\oplus (\ZZ/2\ZZ)^{18},
\]
其中第一因子恰由 boundary 三条中心 sheet-flip 的像生成。故 \(C_{\mathrm{bdry}}\) 在交换化中的像彼此独立，并形成一个规范三维坐标子格。证毕。
\end{proof}

\noindent
由此，window-\(6\) 的多对一折叠既不能被解释为自由有限对称的纯商，也不能被有限秩加法账本完全线性化；但这并不意味着群论痕迹全部消失。恰相反，boundary 三条二点纤维仍在中心层留下一个可规范识别的 \((\ZZ/2\ZZ)^3\) 残余，并在 parity-charge 的交换化通道中保持三维独立坐标。

\paragraph{bin-fold 群胚代数的 Morita 刚性、规范群剖面与最大纤维饱和有限性}
在 bin-fold 的纤维群胚 Wedderburn 分解、碰撞--Parseval 维数量、纤维置换分解以及最大纤维组合上界的阈值判据已经建立之后，还可把同一套退化谱进一步压缩为下列四条对象层结论。

\begin{theorem}[bin-fold 纤维群胚代数的 Morita--\(K\) 理论、中心计数与精确维数下界]\label{thm:xi-foldbin-groupoid-morita-k-center-dimension-lb}
沿用定义 \ref{def:fold-bin-fiber-groupoid}、命题 \ref{prop:fold-bin-groupoid-algebra-wedderburn} 以及定理 \ref{thm:fold-collision-spectrum} 的 Fourier 记号，记
\[
A_m:=\CC[\mathcal R_m^{\mathrm{bin}}]
\cong
\bigoplus_{w\in X_m}M_{d_m^{\mathrm{bin}}(w)}(\CC).
\]
则 \(A_m\) 作为有限维复代数与 \(\CC^{F_{m+2}}\) Morita 等价；在自然的 \(\ast\)-结构下将其视为有限维 \(C^\ast\)-代数，则
\[
K_0(A_m)\cong \ZZ^{F_{m+2}},
\qquad
K_1(A_m)=0.
\]
并且
\[
Z(A_m)\cong \CC^{F_{m+2}},
\qquad
\dim_{\CC}Z(A_m)
=
\#\operatorname{Irr}(A_m)
=
\#\operatorname{Prim}(A_m)
=
F_{m+2}.
\]
此外，
\[
\dim_{\CC}A_m
=
\sum_{w\in X_m}\bigl(d_m^{\mathrm{bin}}(w)\bigr)^2
=
\frac{1}{F_{m+2}}
\sum_{k\in \ZZ/(F_{m+2}\ZZ)}
\bigl|\widehat d_m(k)\bigr|^2,
\]
并满足精确下界
\[
\dim_{\CC}A_m\ge \frac{2^{2m}}{F_{m+2}}.
\]
等号成立当且仅当
\[
d_m^{\mathrm{bin}}(w)\equiv \frac{2^m}{F_{m+2}}
\qquad (w\in X_m),
\]
亦即全部纤维完全平衡。
\end{theorem}
\begin{proof}
由命题 \ref{prop:fold-bin-groupoid-algebra-wedderburn}，
\[
A_m\cong \bigoplus_{w\in X_m}M_{d_m^{\mathrm{bin}}(w)}(\CC).
\]
每个矩阵块 \(M_d(\CC)\) 与 \(\CC\) Morita 等价，故其有限直和与
\[
\bigoplus_{w\in X_m}\CC\cong \CC^{|X_m|}=\CC^{F_{m+2}}
\]
Morita 等价。
有限维 \(C^\ast\)-代数的 \(K\)-理论在 Morita 等价下不变，而
\[
K_0(M_d(\CC))\cong \ZZ,
\qquad
K_1(M_d(\CC))=0.
\]
对有限直和逐块相加，便得到
\[
K_0(A_m)\cong \ZZ^{F_{m+2}},
\qquad
K_1(A_m)=0.
\]

中心由块对角形式直接给出
\[
Z(A_m)\cong \bigoplus_{w\in X_m}\CC\cong \CC^{F_{m+2}}.
\]
因此中心维数、简单模个数与原始理想个数都恰等于块数 \(F_{m+2}\)。

另一方面，
\[
\dim_{\CC}A_m
=
\sum_{w\in X_m}\dim_{\CC}M_{d_m^{\mathrm{bin}}(w)}(\CC)
=
\sum_{w\in X_m}\bigl(d_m^{\mathrm{bin}}(w)\bigr)^2.
\]
再由定理 \ref{thm:fold-collision-spectrum} 的 Parseval 恒等式与
\(
|X_m|=F_{m+2}
\)
可得
\[
\sum_{w\in X_m}\bigl(d_m^{\mathrm{bin}}(w)\bigr)^2
=
\frac{1}{F_{m+2}}
\sum_{k\in \ZZ/(F_{m+2}\ZZ)}
\bigl|\widehat d_m(k)\bigr|^2.
\]

最后，由 Cauchy--Schwarz，
\[
\sum_{w\in X_m}\bigl(d_m^{\mathrm{bin}}(w)\bigr)^2
\ge
\frac{\bigl(\sum_w d_m^{\mathrm{bin}}(w)\bigr)^2}{|X_m|}
=
\frac{2^{2m}}{F_{m+2}},
\]
因为
\(
\sum_w d_m^{\mathrm{bin}}(w)=2^m
\)。
取等当且仅当全部 \(d_m^{\mathrm{bin}}(w)\) 相等，亦即纤维完全平衡。证毕。
\end{proof}

\begin{theorem}[bin-fold 规范群的群阶、交换化与中心的完整闭式]\label{thm:xi-foldbin-gauge-group-order-abel-center-closed}
\leanverified{paper\_xi\_foldbin\_gauge\_group\_order\_abel\_center\_closed}
令
\[
G_m^{\mathrm{bin}}:=\Aut(\Fold_m^{\mathrm{bin}}).
\]
则
\[
|G_m^{\mathrm{bin}}|
=
\prod_{w\in X_m}d_m^{\mathrm{bin}}(w)!.
\]
并且若记
\[
N_m:=\#\{w\in X_m:\ d_m^{\mathrm{bin}}(w)\ge 2\},
\qquad
M_m^{(2)}:=\#\{w\in X_m:\ d_m^{\mathrm{bin}}(w)=2\},
\]
则有自然同构
\[
\bigl(G_m^{\mathrm{bin}}\bigr)^{\mathrm{ab}}
\cong
(\ZZ/2\ZZ)^{N_m},
\qquad
Z\!\bigl(G_m^{\mathrm{bin}}\bigr)
\cong
(\ZZ/2\ZZ)^{M_m^{(2)}}.
\]
特别地，在 \(m=6\) 时，由直方图 \(2:8,\ 3:4,\ 4:9\) 得
\[
\bigl(G_6^{\mathrm{bin}}\bigr)^{\mathrm{ab}}\cong (\ZZ/2\ZZ)^{21},
\qquad
Z\!\bigl(G_6^{\mathrm{bin}}\bigr)\cong (\ZZ/2\ZZ)^8.
\]
\end{theorem}
\begin{proof}
由命题 \ref{prop:fold-bin-gauge-decomposition}，
\[
G_m^{\mathrm{bin}}
\cong
\prod_{w\in X_m}\mathfrak S_{d_m^{\mathrm{bin}}(w)}.
\]
故群阶立刻为
\[
|G_m^{\mathrm{bin}}|
=
\prod_{w\in X_m}\bigl|\mathfrak S_{d_m^{\mathrm{bin}}(w)}\bigr|
=
\prod_{w\in X_m}d_m^{\mathrm{bin}}(w)!.
\]

又因
\[
\mathfrak S_d^{\mathrm{ab}}\cong
\begin{cases}
1,& d\le 1,\\
\ZZ/2\ZZ,& d\ge 2,
\end{cases}
\qquad
Z(\mathfrak S_d)\cong
\begin{cases}
\ZZ/2\ZZ,& d=2,\\
1,& d\neq 2,
\end{cases}
\]
有限直积对交换化与中心逐因子计算，便得到
\[
\bigl(G_m^{\mathrm{bin}}\bigr)^{\mathrm{ab}}
\cong
(\ZZ/2\ZZ)^{N_m},
\qquad
Z\!\bigl(G_m^{\mathrm{bin}}\bigr)
\cong
(\ZZ/2\ZZ)^{M_m^{(2)}}.
\]

最后，当 \(m=6\) 时，由推论 \ref{cor:fold-bin-m6-fiber-hist} 的直方图 \(2:8,\ 3:4,\ 4:9\) 可知
\[
N_6=8+4+9=21,
\qquad
M_6^{(2)}=8,
\]
从而得到所述两式。证毕。
\end{proof}

\begin{theorem}[bin-fold 最大纤维组合上界的逐纤维精确饱和判据]\label{thm:xi-foldbin-fiberwise-saturation-sector-criterion}
\leanverified{paper\_xi\_foldbin\_fiberwise\_saturation\_sector\_criterion}
沿用定义 \ref{def:K-of-m} 与命题 \ref{prop:fold-bin-fiber-tail} 的记号。
设 \(m\ge 1\)，取 \(K=K(m)\) 并记 \(L:=K-m\)。
定义
\[
S_{\max}(m):=
\max\Bigl\{
\sum_{k=m+1}^{K}c_kF_{k+1}:
\ (c_{m+1},\dots,c_K)\in\{0,1\}^{L},\ c_kc_{k+1}=0
\Bigr\}.
\]
则对任意 \(w\in X_m\) 有上界
\[
d_m^{\mathrm{bin}}(w)\le
\begin{cases}
F_{L+2},& w_m=0,\\
F_{L+1},& w_m=1.
\end{cases}
\]
并且达到组合上界 \(F_{L+2}\) 当且仅当
\[
w_m=0
\quad\text{且}\quad
V_m(w)\le 2^m-1-S_{\max}(m).
\]
等价地，
\[
\{w\in X_m:\ d_m^{\mathrm{bin}}(w)=F_{L+2}\}
=
\{w\in X_m:\ w_m=0,\ V_m(w)\le 2^m-1-S_{\max}(m)\}.
\]
\end{theorem}
\begin{proof}
由命题 \ref{prop:fold-bin-fiber-tail}，纤维
\(
(\Fold_m^{\mathrm{bin}})^{-1}(w)
\)
与满足
\[
\sum_{k=m+1}^{K}c_kF_{k+1}\le 2^m-1-V_m(w),
\qquad
c_kc_{k+1}=0,
\qquad
w_mc_{m+1}=0
\]
的尾部位串一一对应。

先忽略阈值约束，只保留相邻约束 \(c_kc_{k+1}=0\)。
长度 \(L\) 的无相邻 \(1\) 二进制串总数是标准 Fibonacci 计数 \(F_{L+2}\)。
当 \(L=0\) 时，尾部只剩空串，上述结论平凡成立。
以下设 \(L\ge 1\)。
若 \(w_m=1\)，则还必须满足 \(c_{m+1}=0\)，故可选尾部退化为长度 \(L-1\) 的无相邻 \(1\) 串，总数为 \(F_{L+1}\)。
于是得到上述逐纤维上界。

若 \(w_m=0\) 且
\[
V_m(w)\le 2^m-1-S_{\max}(m),
\]
则每一条满足相邻约束的尾部串都自动满足阈值约束，因而全部 \(F_{L+2}\) 条合法尾部都被保留，故
\[
d_m^{\mathrm{bin}}(w)=F_{L+2}.
\]

反之，若 \(w_m=1\)，则上界已降为 \(F_{L+1}<F_{L+2}\)，故不可能饱和。
若 \(w_m=0\) 但
\[
V_m(w)>2^m-1-S_{\max}(m),
\]
则任取实现 \(S_{\max}(m)\) 的尾部串，便有
\[
V_m(w)+\sum_{k=m+1}^{K}c_kF_{k+1}>2^m-1,
\]
从而该尾部被阈值约束排除，所以可行尾部数严格小于 \(F_{L+2}\)。
综上，达到 \(F_{L+2}\) 的充要条件正是所述阈值判据。证毕。
\end{proof}

\begin{corollary}[最大纤维饱和的精确等价条件与仅有限次发生]\label{cor:xi-foldbin-maxfiber-saturation-finite-occurrence}
沿用定理 \ref{thm:xi-foldbin-fiberwise-saturation-sector-criterion} 的记号，并记
\[
d_{\max}^{\mathrm{bin}}(m):=
\max_{w\in X_m}d_m^{\mathrm{bin}}(w).
\]
则有精确等价
\[
d_{\max}^{\mathrm{bin}}(m)=F_{L+2}
\iff
S_{\max}(m)\le 2^m-1.
\]
因此，最大纤维达到组合上界 \(F_{L+2}\) 只会对有限多个 \(m\) 发生。
并且在可复核扫描窗口 \(m\le 400\) 内，
\[
\{\,m\le 400:\ d_{\max}^{\mathrm{bin}}(m)=F_{L+2}\,\}
=
\{1,2,3,4,5,7,12\},
\]
如表 \ref{tab:fold_binary_saturation_points} 所示。
\end{corollary}
\begin{proof}
由推论 \ref{cor:fold-bin-saturation}，最大纤维由 \(0^m\) 取得，而 \(0^m\) 满足
\[
(0^m)_m=0,
\qquad
V_m(0^m)=0.
\]
因此定理 \ref{thm:xi-foldbin-fiberwise-saturation-sector-criterion} 作用于 \(w=0^m\) 即给出
\[
d_{\max}^{\mathrm{bin}}(m)=F_{L+2}
\iff
S_{\max}(m)\le 2^m-1.
\]

再由注 \ref{rem:fold-bin-finite}，上述饱和条件只能发生有限次，故最大纤维达到组合上界也只能有限次发生。
最后，在 \(m\le 400\) 的审计窗口内，表 \ref{tab:fold_binary_saturation_points} 已逐项列出全部饱和点，故得到所示有限集合。证毕。
\end{proof}

\paragraph{bin-fold 群胚代数的忠实表示尺度、PI 阈值与六倍窗半尺度缺陷}
在 bin-fold 纤维群胚代数的 Morita--\(K\) 理论、中心计数、规范群交换化以及六倍窗零余类放大律都已建立之后，还可继续从同一份 Wedderburn 退化直方图抽取出下列五条对象层后果。它们分别刻画最小忠实表示尺度、Morita 类型对块大小的遗忘、PI 阈值的最大纤维判别、六倍窗信息缺口的黄金半尺度锁定，以及 window-\(6\) 上三种彼此不重合的刚性尺度。

\begin{theorem}[bin-fold 纤维群胚代数的最小忠实表示维数等于全部微态数]\label{thm:xi-time-part60aca-minimal-faithful-dimension-microstate-count}
\leanverified{paper\_xi\_time\_part60aca\_minimal\_faithful\_dimension\_microstate\_count}
沿用定义 \ref{def:fold-bin-fiber-groupoid} 与命题 \ref{prop:fold-bin-groupoid-algebra-wedderburn} 的记号，记
\[
A_m:=\CC[\mathcal R_m^{\mathrm{bin}}]
\cong
\bigoplus_{w\in X_m}M_{d_w}(\CC),
\qquad
d_w:=d_m^{\mathrm{bin}}(w).
\]
定义
\[
\mu_{\mathrm{faith}}(A_m)
:=
\min\Bigl\{
\dim_{\CC}H:\ 
\exists\ \rho:A_m\hookrightarrow \operatorname{End}_{\CC}(H)\ \text{为忠实幺代数表示}
\Bigr\}.
\]
则有严格恒等式
\[
\mu_{\mathrm{faith}}(A_m)
=
\sum_{w\in X_m} d_w
=
2^m.
\]
等价地，\(A_m\) 的最小忠实 Hilbert 载体维数恰等于全部微态总数。
\end{theorem}
\begin{proof}
对每个 \(w\in X_m\)，记 \(e_w\in A_m\) 为对应于 \(M_{d_w}(\CC)\) 这一 Wedderburn 块的中心极小幂等元。它们满足
\[
e_we_{w'}=\delta_{w,w'}e_w,
\qquad
\sum_{w\in X_m}e_w=1.
\]
任取忠实幺表示
\[
\rho:A_m\hookrightarrow \operatorname{End}_{\CC}(H).
\]
由于 \(\rho\) 单射，必有 \(\rho(e_w)\neq 0\) 对一切 \(w\) 成立。令
\[
H_w:=\rho(e_w)H.
\]
由幂等元两两正交且和为 \(1\)，得到直和分解
\[
H=\bigoplus_{w\in X_m}H_w.
\]

限制到第 \(w\) 块，得到一个非零表示
\[
M_{d_w}(\CC)\cong e_wA_m\longrightarrow \operatorname{End}_{\CC}(H_w).
\]
由于 \(M_{d_w}(\CC)\) 是简单代数，非零表示的核只能是零，故该表示自动忠实。另一方面，有限维 \(M_{d_w}(\CC)\)-模必分解为标准模 \(\CC^{d_w}\) 的若干重直和，因此
\[
\dim_{\CC}H_w\ge d_w.
\]
求和即得
\[
\dim_{\CC}H=\sum_{w\in X_m}\dim_{\CC}H_w\ge \sum_{w\in X_m}d_w.
\]

反向地，对每个 \(w\) 取标准表示
\[
M_{d_w}(\CC)\curvearrowright \CC^{d_w},
\]
再作外直和
\[
\bigoplus_{w\in X_m}\CC^{d_w},
\]
便得到 \(A_m\) 的忠实表示，其维数恰为 \(\sum_w d_w\)。因此
\[
\mu_{\mathrm{faith}}(A_m)=\sum_{w\in X_m}d_w.
\]
最后，\(\Omega_m=\{0,1\}^m\) 被各条纤维
\[
\Fold_m^{\mathrm{bin}}{}^{-1}(w)
\qquad (w\in X_m)
\]
不交分割，故
\[
\sum_{w\in X_m}d_w=|\Omega_m|=2^m.
\]
证毕。
\end{proof}

\leanverified{paper\_xi\_time\_part60aca\_morita\_type\_forgets\_block\_sizes}
\begin{theorem}[bin-fold 纤维群胚代数的 Morita 类型完全坍缩到 Fibonacci 中心秩]\label{thm:xi-time-part60aca-morita-type-forgets-block-sizes}
沿用定理 \ref{thm:xi-time-part60aca-minimal-faithful-dimension-microstate-count} 的记号。则 \(A_m\) 与
\[
\CC^{F_{m+2}}
\]
Morita 等价。进一步，在自然的 \(\ast\)-结构下将 \(A_m\) 视为有限维 \(C^\ast\)-代数，则其有序 \(K_0\) 群满足
\[
K_0(A_m)\cong \ZZ^{F_{m+2}},
\qquad
K_0(A_m)^+\cong \NN^{F_{m+2}}.
\]
因此，bin-fold 纤维群胚代数的 Morita 类型完全忘记了块大小 \((d_w)\)，只保留块数 \(|X_m|=F_{m+2}\)。
\end{theorem}
\begin{proof}
Morita 等价
\[
A_m\ \text{与}\ \CC^{F_{m+2}}
\]
以及群同构
\[
K_0(A_m)\cong \ZZ^{F_{m+2}}
\]
已由定理 \ref{thm:xi-foldbin-groupoid-morita-k-center-dimension-lb} 给出。现只需说明正锥。

由 Wedderburn 分解，
\[
A_m\cong \bigoplus_{w\in X_m}M_{d_w}(\CC)
\]
作为有限维 \(C^\ast\)-代数成立。于是
\[
K_0(A_m)\cong \bigoplus_{w\in X_m}K_0\bigl(M_{d_w}(\CC)\bigr)
\cong \ZZ^{|X_m|},
\]
其中每个直和因子由相应矩阵块中的投影类生成。对矩阵代数 \(M_d(\CC)\)，投影类恰按其秩由 \(\NN\) 参数化，因此正锥逐块相加后得到
\[
K_0(A_m)^+\cong \NN^{|X_m|}.
\]
再由 \(|X_m|=F_{m+2}\)，遂得
\[
K_0(A_m)^+\cong \NN^{F_{m+2}}.
\]
证毕。
\end{proof}

\begin{theorem}[bin-fold 纤维群胚代数的 PI 阈值由最大纤维完全决定]\label{thm:xi-time-part60aca-pideg-maxfiber-threshold}
沿用定理 \ref{thm:xi-time-part60aca-minimal-faithful-dimension-microstate-count} 的记号，并定义
\[
D_m^{\mathrm{bin}}
:=
\max_{w\in X_m}d_m^{\mathrm{bin}}(w)
=
d_{\max}^{\mathrm{bin}}(m).
\]
对 \(r\ge 1\)，记 \(2r\) 次标准多项式为
\[
s_{2r}(x_1,\dots,x_{2r})
:=
\sum_{\sigma\in \mathfrak S_{2r}}
\operatorname{sgn}(\sigma)\,
x_{\sigma(1)}\cdots x_{\sigma(2r)}.
\]
则
\[
\operatorname{PIdeg}(A_m)=D_m^{\mathrm{bin}}.
\]
等价地，
\[
A_m\models s_{2D_m^{\mathrm{bin}}},
\qquad
A_m\not\models s_{2D_m^{\mathrm{bin}}-2}.
\]
特别地，当 \(m=6\) 时，由直方图 \(2:8,\ 3:4,\ 4:9\) 得
\[
\operatorname{PIdeg}(A_6)=4,
\qquad
A_6\models s_8,
\qquad
A_6\not\models s_6.
\]
\end{theorem}
\begin{proof}
由 Wedderburn 分解，
\[
A_m\cong \bigoplus_{w\in X_m}M_{d_w}(\CC).
\]
对有限维半单复代数
\[
A\cong \bigoplus_{j=1}^{r}M_{n_j}(\CC),
\]
一个多项式恒等式在 \(A\) 上成立，当且仅当它在每个矩阵块 \(M_{n_j}(\CC)\) 上都成立。由 Amitsur--Levitzki 定理，
\[
M_{n_j}(\CC)\models s_{2n_j},
\qquad
M_{n_j}(\CC)\not\models s_{2n_j-2}.
\]
因此 \(A\) 的 PI 度数恰为 \(\max_j n_j\)。将
\[
n_j=d_w
\]
代入即得
\[
\operatorname{PIdeg}(A_m)=\max_{w\in X_m}d_w=D_m^{\mathrm{bin}}.
\]
等价的标准恒等式表述随即成立。

当 \(m=6\) 时，推论 \ref{cor:fold-bin-m6-fiber-hist} 给出不同块尺寸恰为 \(2,3,4\)，故 \(D_6^{\mathrm{bin}}=4\)。于是
\[
\operatorname{PIdeg}(A_6)=4,
\qquad
A_6\models s_8,
\qquad
A_6\not\models s_6.
\]
证毕。
\end{proof}

\begin{theorem}[六倍窗共振把碰撞与 KL 缺口锁定在黄金半尺度]\label{thm:xi-time-part60aca-sixfold-halfscale-defect}
沿用定理 \ref{thm:xi-fold-zero-arithmetic-subsequence-amplification} 的记号，特别记
\[
p_m(r):=\frac{d_m(r)}{2^m},
\qquad
u_m(r):=\frac{1}{F_{m+2}}.
\]
若
\[
m+2\equiv 0\pmod 6,
\qquad
d:=\frac{m+2}{2},
\]
则 bin-fold 的碰撞缺口与 KL 缺口满足
\[
1-\mathsf{Col}_m
\ge
\frac{F_d}{F_{m+2}}
=
\frac{1}{L_d}
=
\varphi^{-d}\bigl(1+O(\varphi^{-2d})\bigr),
\]
以及
\[
\log F_{m+2}-D_{\mathrm{KL}}(p_m\|u_m)
\ge
-\log\!\left(1-\frac{F_d}{F_{m+2}}\right)
=
\varphi^{-d}\bigl(1+O(\varphi^{-d})\bigr).
\]
等价地，
\[
1-\mathsf{Col}_m
\ge
\varphi^{-(m+2)/2}\bigl(1+O(\varphi^{-m})\bigr),
\]
\[
\log F_{m+2}-D_{\mathrm{KL}}(p_m\|u_m)
\ge
\varphi^{-(m+2)/2}\bigl(1+O(\varphi^{-m/2})\bigr).
\]
因此，在六倍窗子序列上，bin-fold 对均匀分布的接近程度不可能优于一个黄金半尺度 \(\varphi^{-m/2}\)。
\end{theorem}
\begin{proof}
由结论 \ref{con:xi-fold-window6-lucas-barrier-collision-gap}，
\[
1-\mathsf{Col}_m
\ge
\frac{F_d}{F_{2d}}
=
\frac{1}{L_d},
\]
其中 \(L_d\) 为第 \(d\) 个 Lucas 数。由 Binet 公式，
\[
L_d=\varphi^d+(-\varphi)^{-d}
=
\varphi^d\bigl(1+O(\varphi^{-2d})\bigr),
\]
从而
\[
\frac{1}{L_d}
=
\varphi^{-d}\bigl(1+O(\varphi^{-2d})\bigr).
\]
又因 \(2d=m+2\)，故第一条渐近式可改写为
\[
1-\mathsf{Col}_m
\ge
\varphi^{-(m+2)/2}\bigl(1+O(\varphi^{-m})\bigr).
\]

另一方面，条件 \(m+2\equiv 0\pmod 6\) 等价于 \(m\equiv 4\pmod 6\)。于是定理 \ref{thm:xi-fold-zero-arithmetic-subsequence-amplification} 给出
\[
D_{\mathrm{KL}}(p_m\|u_m)
\le
\log\!\bigl(F_{m+2}-F_d\bigr).
\]
因此
\[
\log F_{m+2}-D_{\mathrm{KL}}(p_m\|u_m)
\ge
\log F_{m+2}-\log\!\bigl(F_{m+2}-F_d\bigr)
=
-\log\!\left(1-\frac{F_d}{F_{m+2}}\right).
\]
再由
\[
\frac{F_d}{F_{m+2}}
=
\frac{F_d}{F_{2d}}
=
\frac{1}{L_d}
=
\varphi^{-d}\bigl(1+O(\varphi^{-2d})\bigr),
\]
以及
\[
-\log(1-x)=x+O(x^2)
\qquad(x\to 0),
\]
便得到
\[
\log F_{m+2}-D_{\mathrm{KL}}(p_m\|u_m)
\ge
\varphi^{-d}\bigl(1+O(\varphi^{-d})\bigr).
\]
把 \(d=(m+2)/2\) 代回即得所示第二条渐近式。证毕。
\end{proof}

\begin{theorem}[window-\texorpdfstring{$6$}{6} 上 Morita 秩、PI 阈值与忠实表示维数的三尺度分裂]\label{thm:xi-time-part60aca-window6-triple-rigidity-scales}
取 \(m=6\)。则
\[
A_6
\cong
M_2(\CC)^{\oplus 8}\oplus M_3(\CC)^{\oplus 4}\oplus M_4(\CC)^{\oplus 9}.
\]
从而同一非交换对象上同时出现三条互不重合的刚性尺度：
\[
\#\operatorname{Irr}(A_6)
=
\operatorname{rank}_{\ZZ}K_0(A_6)
=
21,
\qquad
\operatorname{PIdeg}(A_6)=4,
\qquad
\mu_{\mathrm{faith}}(A_6)=64.
\]
此外，
\[
A_6\models s_8,
\qquad
A_6\not\models s_6,
\qquad
\bigl(G_6^{\mathrm{bin}}\bigr)^{\mathrm{ab}}\cong (\ZZ/2\ZZ)^{21}.
\]
因此，window-\(6\) 上的 Morita--\(K\) 理论尺度、PI 尺度与忠实表示尺度并不塌缩为同一数值。
\end{theorem}
\begin{proof}
由推论 \ref{cor:fold-bin-m6-fiber-hist}，
\[
\#\{w\in X_6:\ d_6^{\mathrm{bin}}(w)=2\}=8,
\qquad
\#\{w\in X_6:\ d_6^{\mathrm{bin}}(w)=3\}=4,
\qquad
\#\{w\in X_6:\ d_6^{\mathrm{bin}}(w)=4\}=9.
\]
代入 Wedderburn 分解便得
\[
A_6
\cong
M_2(\CC)^{\oplus 8}\oplus M_3(\CC)^{\oplus 4}\oplus M_4(\CC)^{\oplus 9}.
\]

由定理 \ref{thm:xi-time-part60aca-morita-type-forgets-block-sizes}，
\[
\#\operatorname{Irr}(A_6)
=
\operatorname{rank}_{\ZZ}K_0(A_6)
=
|X_6|
=
21.
\]
由定理 \ref{thm:xi-time-part60aca-pideg-maxfiber-threshold}，
\[
\operatorname{PIdeg}(A_6)=4,
\qquad
A_6\models s_8,
\qquad
A_6\not\models s_6.
\]
再由定理 \ref{thm:xi-time-part60aca-minimal-faithful-dimension-microstate-count}，
\[
\mu_{\mathrm{faith}}(A_6)=2^6=64.
\]
等价地，直接由直方图求和也可写成
\[
\mu_{\mathrm{faith}}(A_6)
=
8\cdot 2+4\cdot 3+9\cdot 4
=
64.
\]
最后，定理 \ref{thm:xi-foldbin-gauge-group-order-abel-center-closed} 给出
\[
\bigl(G_6^{\mathrm{bin}}\bigr)^{\mathrm{ab}}\cong (\ZZ/2\ZZ)^{21},
\]
故其 \(\FF_2\)-维数与第一条尺度一致，而与后两条尺度不同。证毕。
\end{proof}

\paragraph{bin-fold 群胚代数的有序 \texorpdfstring{$K_0$}{K0}、稳定刚性与自同构阈值}
在 bin-fold 纤维群胚代数的 Wedderburn 信息完备性、Morita--\(K\) 理论、PI 阈值与 window-\(6\) 的异常资源分叉已经建立之后，还可把同一退化谱直方图继续压缩为如下六条对象层结论：带单位元类的有序 \(K_0\) 完全分类、稳定同构刚性、PI 度数的黄金下界、自同构群的半直积分解、中心块数与异常比特数的严格对齐，以及首次非交换阈值的精确定位。

\begin{theorem}[bin-fold 纤维群胚代数的有序 \texorpdfstring{$K_0$}{K0} 完全分类]\label{thm:xi-time-part60acb-binfold-ordered-k0-complete-classification}
沿用定义 \ref{def:fold-bin-fiber-groupoid} 与命题 \ref{prop:fold-bin-groupoid-algebra-wedderburn} 的记号，记
\[
A_m:=\CC[\mathcal R_m^{\mathrm{bin}}],
\qquad
\mathcal R_m^{\mathrm{bin}}
=
\bigsqcup_{w\in X_m}(F_w\times F_w),
\qquad
d_w:=|F_w|=d_m^{\mathrm{bin}}(w).
\]
则在由 Wedderburn 块给出的自然坐标下，有
\[
\boxed{\;
\bigl(K_0(A_m),K_0(A_m)^+,[1_{A_m}]\bigr)
\cong
\bigl(\ZZ^{X_m},\NN^{X_m},(d_w)_{w\in X_m}\bigr).\;}
\]
特别地，带 order unit 的有序 \(K_0\) 唯一恢复纤维多重度多重集
\[
\{d_w:w\in X_m\},
\]
反过来，该多重度多重集也唯一决定 \(A_m\) 的复代数同构类。
\end{theorem}
\begin{proof}
由命题 \ref{prop:fold-bin-groupoid-algebra-wedderburn}，
\[
A_m\cong \bigoplus_{w\in X_m}M_{d_w}(\CC).
\]
因此
\[
K_0(A_m)\cong \bigoplus_{w\in X_m}K_0\!\bigl(M_{d_w}(\CC)\bigr).
\]
对每个矩阵块 \(M_{d_w}(\CC)\)，以秩一投影类作为 \(K_0\)-生成元，则
\[
K_0\!\bigl(M_{d_w}(\CC)\bigr)\cong \ZZ,
\qquad
K_0\!\bigl(M_{d_w}(\CC)\bigr)^+\cong \NN,
\qquad
[1_{M_{d_w}(\CC)}]=d_w.
\]
逐块直和便得到
\[
K_0(A_m)\cong \ZZ^{X_m},
\qquad
K_0(A_m)^+\cong \NN^{X_m},
\qquad
[1_{A_m}]=(d_w)_{w\in X_m}.
\]

由于 \(\NN^{X_m}\) 的极端射线恰对应于各个 Wedderburn 块，故在不固定块顺序时，order unit 的坐标仍唯一恢复多重度多重集 \(\{d_w\}_{w\in X_m}\) 仅差一个置换。反向方向由定理 \ref{thm:xi-time-part60aaa-wedderburn-information-completeness} 立即给出：纤维直方图，等价地多重度多重集，唯一决定 \(A_m\) 的半单代数同构型。证毕。
\end{proof}

\begin{corollary}[Morita 压缩与稳定刚性]\label{cor:xi-time-part60acb-binfold-morita-stable-rigidity}
设 \(\mathcal K\) 为可分无限维 Hilbert 空间上的紧算子代数。则对每个 \(m\ge 1\)，有
\[
\boxed{\;
A_m\ \text{与}\ \CC^{F_{m+2}}\ \text{Morita 等价}.\;}
\]
并且对任意 \(m,m'\ge 1\)，成立等价
\[
\boxed{\;
A_m\otimes\mathcal K\cong A_{m'}\otimes\mathcal K
\iff
F_{m+2}=F_{m'+2}
\iff
m=m'.\;}
\]
\end{corollary}
\begin{proof}
第一条正是定理 \ref{thm:xi-time-part60aca-morita-type-forgets-block-sizes} 的内容。

另一方面，由 Wedderburn 分解，
\[
A_m\otimes\mathcal K
\cong
\bigoplus_{w\in X_m}\bigl(M_{d_w}(\CC)\otimes\mathcal K\bigr)
\cong
\bigoplus_{w\in X_m}\mathcal K
\cong
\mathcal K^{\oplus F_{m+2}},
\]
因为 \(M_d(\CC)\otimes\mathcal K\cong \mathcal K\) 对每个 \(d\ge 1\) 都成立。
故若
\[
A_m\otimes\mathcal K\cong A_{m'}\otimes\mathcal K,
\]
则由稳定性
\[
K_0(A_m\otimes\mathcal K)\cong K_0(A_m),
\qquad
K_0(A_{m'}\otimes\mathcal K)\cong K_0(A_{m'}),
\]
再结合定理 \ref{thm:xi-time-part60acb-binfold-ordered-k0-complete-classification}，得到
\[
\ZZ^{F_{m+2}}\cong K_0(A_m)\cong K_0(A_{m'})\cong \ZZ^{F_{m'+2}},
\]
从而
\[
F_{m+2}=F_{m'+2}.
\]
反之，若两者相等，则两边都同构于同一个
\(
\mathcal K^{\oplus F_{m+2}}
\)。

最后，由 Fibonacci 数列严格递增，\(F_{m+2}=F_{m'+2}\) 等价于 \(m=m'\)。证毕。
\end{proof}

\begin{theorem}[PI 度数的纤维极值律与黄金下界]\label{thm:xi-time-part60acb-binfold-pideg-golden-lower-bound}
沿用记号
\[
A_m\cong \bigoplus_{w\in X_m}M_{d_w}(\CC),
\qquad
M_m:=\max_{w\in X_m}d_w,
\qquad
\overline d_m:=\frac{2^m}{F_{m+2}}.
\]
则
\[
\boxed{\;
\operatorname{PIdeg}(A_m)=M_m.\;}
\]
等价地，\(A_m\) 的最小标准多项式恒等式阶数恰为
\[
\boxed{\;
2M_m.\;}
\]
进一步，若 \(C_\varphi>0\) 为定理 \ref{thm:fold-critical-resonance-constant} 的黄金共振常数，则
\[
\boxed{\;
\liminf_{m\to\infty}
\frac{\operatorname{PIdeg}(A_m)}{\overline d_m}
\ge
1+C_\varphi.\;}
\]
\end{theorem}
\begin{proof}
由定理 \ref{thm:xi-time-part60aca-pideg-maxfiber-threshold}，已知
\[
\operatorname{PIdeg}(A_m)=M_m,
\qquad
A_m\models s_{2M_m},
\qquad
A_m\not\models s_{2M_m-2},
\]
故前两条断言成立。

再由定理 \ref{thm:xi-time-part9m-maxfiber-uniformity-gap-chain}，
\[
\liminf_{m\to\infty}\frac{M_m}{\overline d_m}\ge 1+C_\varphi.
\]
把 \(\operatorname{PIdeg}(A_m)=M_m\) 代入即得
\[
\liminf_{m\to\infty}
\frac{\operatorname{PIdeg}(A_m)}{\overline d_m}
\ge
1+C_\varphi.
\]
证毕。
\end{proof}

\begin{theorem}[纤维群胚代数自同构群的半直积分解]\label{thm:xi-time-part60acb-binfold-algebra-automorphism-semidirect}
对每个 \(d\ge 1\) 记
\[
n_d(m):=\#\{w\in X_m:\ d_w=d\}.
\]
则
\[
\boxed{\;
\operatorname{Aut}_{\CC\text{-alg}}(A_m)
\cong
\left(\prod_{\substack{d\ge 1\\ n_d(m)\neq 0}}\operatorname{PGL}_d(\CC)^{\,n_d(m)}\right)
\rtimes
\left(\prod_{\substack{d\ge 1\\ n_d(m)\neq 0}}\mathfrak S_{n_d(m)}\right),\;}
\]
其中右端半直积作用由各 \(\mathfrak S_{n_d(m)}\) 对同尺寸 \(\operatorname{PGL}_d(\CC)\)-因子的置换给出。

特别地，在 \(m=6\) 时，由直方图
\[
2:8,\qquad 3:4,\qquad 4:9,
\]
有
\[
\boxed{\;
\operatorname{Aut}_{\CC\text{-alg}}(A_6)
\cong
\bigl(\operatorname{PGL}_2(\CC)^8\times \operatorname{PGL}_3(\CC)^4\times \operatorname{PGL}_4(\CC)^9\bigr)
\rtimes
\bigl(\mathfrak S_8\times \mathfrak S_4\times \mathfrak S_9\bigr).\;}
\]
并且
\[
\boxed{\;
\pi_0\!\bigl(\operatorname{Aut}_{\CC\text{-alg}}(A_6)\bigr)
\cong
\mathfrak S_8\times \mathfrak S_4\times \mathfrak S_9.\;}
\]
\end{theorem}
\begin{proof}
由定理 \ref{thm:xi-time-part60aaa-wedderburn-information-completeness}，
\[
A_m
\cong
\bigoplus_{\substack{d\ge 1\\ n_d(m)\neq 0}}M_d(\CC)^{\oplus n_d(m)}.
\]
由于 \(M_d(\CC)\not\cong M_e(\CC)\)（\(d\neq e\)）作为简单复代数彼此不同构，任一代数自同构都不能在不同尺寸的块之间混合；故
\[
\operatorname{Aut}_{\CC\text{-alg}}(A_m)
\cong
\prod_{\substack{d\ge 1\\ n_d(m)\neq 0}}
\operatorname{Aut}_{\CC\text{-alg}}\!\bigl(M_d(\CC)^{\oplus n_d(m)}\bigr).
\]

现固定一个 \(d\)，记 \(n:=n_d(m)\) 且
\[
B_d:=M_d(\CC)^{\oplus n}.
\]
任一 \(\CC\)-代数自同构 \(\Phi\) 必将 \(B_d\) 的极小中心幂等元两两置换，从而给出一个排列
\[
\sigma\in \mathfrak S_n.
\]
对每个分量，\(\Phi\) 在 \(M_d(\CC)\) 上诱导一个复代数自同构；由 Skolem--Noether 定理，每个此类自同构都由某个
\[
g_i\in \operatorname{PGL}_d(\CC)
\]
给出。于是
\[
\operatorname{Aut}_{\CC\text{-alg}}(B_d)
\cong
\operatorname{PGL}_d(\CC)^n\rtimes \mathfrak S_n,
\]
其中 \(\mathfrak S_n\) 通过置换 \(n\) 个等尺寸分量作用。
对全部 \(d\) 同时取直积，即得所示总分解。

当 \(m=6\) 时，推论 \ref{cor:fold-bin-m6-fiber-hist} 给出
\[
n_2(6)=8,\qquad n_3(6)=4,\qquad n_4(6)=9,
\]
代入即得显式公式。

最后，\(\operatorname{PGL}_d(\CC)\) 在复拓扑下连通，而 \(\mathfrak S_n\) 离散，故
\[
\pi_0\!\bigl(\operatorname{PGL}_d(\CC)^n\rtimes \mathfrak S_n\bigr)\cong \mathfrak S_n.
\]
对 \(d=2,3,4\) 取直积，便得到
\[
\pi_0\!\bigl(\operatorname{Aut}_{\CC\text{-alg}}(A_6)\bigr)
\cong
\mathfrak S_8\times \mathfrak S_4\times \mathfrak S_9.
\]
证毕。
\end{proof}

\begin{theorem}[window-\texorpdfstring{$6$}{6} 的中心块数与异常比特数严格对齐]\label{thm:xi-time-part60acb-window6-center-anomaly-bit-alignment}
记
\[
A_6:=\CC[\mathcal R_6^{\mathrm{bin}}],
\qquad
G_6^{\mathrm{bin}}:=\Aut(\Fold_6^{\mathrm{bin}}).
\]
则
\[
\boxed{\;
\dim_{\CC}Z(A_6)=21,
\qquad
\dim_{\FF_2}\bigl((G_6^{\mathrm{bin}})^{\mathrm{ab}}\bigr)=21.\;}
\]
更具体地：
\begin{enumerate}
  \item 任何保持全部 simple block 逐块可分的交换块账本，至少需要 \(21\) 个独立复坐标；该下界由
  \[
  Z(A_6)\cong \CC^{21}
  \]
  精确达到。
  \item 在纤维独立置换语义下，任何完备承载全部奇偶异常信息的二元异常账本，至少需要 \(21\) 个独立 \(\FF_2\)-坐标；该下界由
  \[
  (G_6^{\mathrm{bin}})^{\mathrm{ab}}\cong (\ZZ/2\ZZ)^{21}
  \]
  精确达到。
\end{enumerate}
因此，window-\(6\) 上最小块分离交换坐标数与最小异常比特数严格相等，二者都恰为 \(21\)。
\end{theorem}
\begin{proof}
由定理 \ref{thm:xi-foldbin-groupoid-morita-k-center-dimension-lb}，
\[
Z(A_6)\cong \CC^{F_8}=\CC^{21},
\]
故
\[
\dim_{\CC}Z(A_6)=21.
\]
记 \(e_w\in Z(A_6)\) 为与各个 Wedderburn 块对应的 \(21\) 个中心原始幂等元。它们两两正交且非零，并满足
\[
\sum_{w\in X_6}e_w=1.
\]
若一个交换块账本仍要把全部 simple block 逐块分开，则这些 \(e_w\) 在该账本中的像必须仍然两两可区分；特别地，它们必须给出 \(21\) 个两两正交的非零幂等元。任何有限维交换半单复代数中，两两正交非零幂等元的个数不超过其复维数，故此类账本的交换坐标数必至少为 \(21\)。而中心 \(Z(A_6)\) 本身恰以 \(21\) 个坐标达到该下界。

另一方面，由定理 \ref{thm:xi-foldbin-gauge-group-order-abel-center-closed}，
\[
\bigl(G_6^{\mathrm{bin}}\bigr)^{\mathrm{ab}}\cong (\ZZ/2\ZZ)^{21},
\]
故
\[
\dim_{\FF_2}\bigl((G_6^{\mathrm{bin}})^{\mathrm{ab}}\bigr)=21.
\]
再由定理 \ref{thm:xi-window6-oracle-anomaly-coordinate-bifurcation}，在保持纤维独立置换语义不变的前提下，完备异常信息的最小二元承载对象恰需 \(21\) 个独立 \(\FF_2\)-坐标，任何更小的方案都必然破坏该逐纤维直积语义。

故块分离交换坐标数与异常比特数在 window-\(6\) 上严格对齐，并且都以 \(21\) 为最小值。证毕。
\end{proof}

\begin{proposition}[bin-fold 群胚代数的首次非交换阈值]\label{prop:xi-time-part60acb-binfold-first-noncommutative-threshold}
\leanverified{paper\_xi\_time\_part60acb\_binfold\_first\_noncommutative\_threshold}
对每个 \(m\ge 1\)，记
\[
A_m\cong \bigoplus_{w\in X_m}M_{d_w}(\CC),
\qquad
|X_m|=F_{m+2}.
\]
则：
\begin{enumerate}
  \item 对所有 \(m\ge 1\)，\(A_m\) 都不是单代数；
  \item \(A_m\) 可交换当且仅当 \(m=1\)；
  \item 因而 bin-fold 群胚代数的首次非交换阈值恰发生在
  \[
  \boxed{\;m=2.\;}
  \]
\end{enumerate}
\end{proposition}
\begin{proof}
由 \(F_{m+2}\ge 2\)（\(m\ge 1\)）可知
\[
|X_m|=F_{m+2}\ge 2.
\]
因此 \(A_m\) 的 Wedderburn 分解中至少含有两个简单块，从而 \(A_m\) 对所有 \(m\ge 1\) 都不是单代数。

再看交换性。有限维半单复代数
\[
\bigoplus_{w\in X_m}M_{d_w}(\CC)
\]
可交换，当且仅当每个块都为 \(1\times 1\)，亦即
\[
d_w=1
\qquad (\forall\,w\in X_m).
\]
这等价于
\[
\sum_{w\in X_m}d_w=|X_m|.
\]
而左端恒等于 \(2^m\)，右端恒等于 \(F_{m+2}\)，故交换性等价于
\[
2^m=F_{m+2}.
\]
当 \(m=1\) 时，
\[
2^1=2=F_3,
\]
于是 \(A_1\cong \CC\oplus \CC\) 确为交换。
当 \(m\ge 2\) 时，Fibonacci 数满足
\[
F_{m+2}<2^m,
\]
例如可由归纳得到：\(F_4=3<4=2^2\)，且若
\[
F_{n+2}<2^n,\qquad F_{n+1}<2^{n-1},
\]
则
\[
F_{n+3}=F_{n+2}+F_{n+1}<2^n+2^{n-1}<2^{n+1}.
\]
从而平均块大小
\[
\frac{1}{|X_m|}\sum_{w\in X_m}d_w=\frac{2^m}{F_{m+2}}>1.
\]
故必存在某个 \(w\) 使 \(d_w\ge 2\)，于是 \(A_m\) 非交换。

综上，\(A_m\) 可交换当且仅当 \(m=1\)，故首次非交换阈值恰为 \(m=2\)。证毕。
\end{proof}


\noindent
由此，bin-fold 纤维群胚代数上的对象层刚性已进一步闭合：Morita--\(K\) 理论只保留 Fibonacci 块数，而带单位元类的有序 \(K_0\) 又把全部块大小逐坐标恢复；稳定范畴中幸存的不变量仍只有块数，并因此对分辨率给出严格刚性；PI 理论则精确读取最大纤维，并继承黄金共振常数所强制的统一下界；同时，代数自同构群按等尺寸矩阵块分裂为内自同构 Lie 群与离散置换群的半直积，window-\(6\) 上的中心块数与异常比特数又同时锁定为 \(21\)；最后，非交换性并非渐近出现，而是在 \(m=2\) 处发生首次精确阈值跳变。因而，退化谱直方图不再只是静态组合摘要，而是同时支配有序 \(K\) 理论、稳定类型、PI 阈值、自同构群与非交换起始尺度的统一刚性对象。


\paragraph{容量直方图、Morita 坍缩、Bernoulli 反演、维数超载与自由商判据的统一闭合}
在连续容量曲线的完备统计性、bin-fold 纤维群胚代数的 Wedderburn 分解、fold-gauge 常数的 Stirling--Bernoulli 层级、两态一致渐近以及 window-\(6\) 的非自由群商障碍都已建立之后，还可把这些链条进一步压缩为下列五条彼此衔接的终端结论。

\begin{theorem}[连续容量曲线同时完备决定纤维直方图、规范群与群胚代数]\label{thm:xi-time-part60acb-capacity-histogram-gauge-groupoid-complete}
\leanverified{paper\_xi\_time\_part60acb\_capacity\_histogram\_gauge\_groupoid\_complete}
固定窗口 \(m\ge 1\)。记
\[
h_m(d):=\#\{w\in X_m:\ d_m^{\mathrm{bin}}(w)=d\}
\qquad(d\ge 1),
\]
\[
N_m(t):=\#\{w\in X_m:\ d_m^{\mathrm{bin}}(w)\ge t\},
\qquad
\mathcal C_m^{\mathrm{cont}}(T):=\sum_{w\in X_m}\min\bigl(d_m^{\mathrm{bin}}(w),T\bigr).
\]
则对每个整数 \(d\ge 1\) 都有
\[
h_m(d)=N_m(d)-N_m(d+1).
\]
因此，\(\mathcal C_m^{\mathrm{cont}}\) 唯一决定全部纤维直方图 \(\{h_m(d)\}_{d\ge 1}\)。进一步，
\[
G_m^{\mathrm{bin}}
\cong
\prod_{d\ge 1}\mathfrak S_d^{\,h_m(d)},
\qquad
\CC[\mathcal R_m^{\mathrm{bin}}]
\cong
\bigoplus_{d\ge 1}M_d(\CC)^{\oplus h_m(d)}.
\]
换言之，在固定窗口上，连续容量曲线同时唯一决定 bin-fold 的纤维多重度直方图、fold-gauge 群的同构型以及纤维群胚代数的 Wedderburn 分解。
\end{theorem}
\begin{proof}
定理 \ref{thm:xi-capacity-tail-histogram-moment-complete-statistic} 已给出
\[
\mathcal C_m^{\mathrm{cont}}(T)=\int_0^T N_m(t)\,\dd t,
\qquad
\partial_T^+\mathcal C_m^{\mathrm{cont}}(t)=N_m(t)\ \text{a.e.}
\]
故 \(N_m\) 由 \(\mathcal C_m^{\mathrm{cont}}\) 唯一恢复。由于 \(N_m\) 是右连续整数值阶梯函数，对每个整数 \(d\ge 1\) 有
\[
N_m(d)-N_m(d+1)
=
\#\{w:\ d_m^{\mathrm{bin}}(w)\ge d\}
-\#\{w:\ d_m^{\mathrm{bin}}(w)\ge d+1\}
=
\#\{w:\ d_m^{\mathrm{bin}}(w)=d\}
=
h_m(d).
\]
这便恢复了全部直方图。

再由命题 \ref{prop:fold-bin-gauge-decomposition}，
\[
G_m^{\mathrm{bin}}
\cong
\prod_{w\in X_m}\mathfrak S_{d_m^{\mathrm{bin}}(w)},
\]
把相同纤维大小的因子归并，即得
\[
G_m^{\mathrm{bin}}
\cong
\prod_{d\ge 1}\mathfrak S_d^{\,h_m(d)}.
\]
同理，由命题 \ref{prop:fold-bin-groupoid-algebra-wedderburn}，
\[
\CC[\mathcal R_m^{\mathrm{bin}}]
\cong
\bigoplus_{w\in X_m}M_{d_m^{\mathrm{bin}}(w)}(\CC),
\]
按相同矩阵尺寸归并便得到
\[
\CC[\mathcal R_m^{\mathrm{bin}}]
\cong
\bigoplus_{d\ge 1}M_d(\CC)^{\oplus h_m(d)}.
\]
结论成立。
\end{proof}

\begin{theorem}[bin-fold 纤维群胚代数的 Morita--\(K\) 理论类型仅由 Fibonacci 块数决定]\label{thm:xi-time-part60acb-groupoid-morita-k-stable-classification}
\leanverified{paper\_xi\_time\_part60acb\_groupoid\_morita\_k\_stable\_classification}
记
\[
A_m:=\CC[\mathcal R_m^{\mathrm{bin}}].
\]
则
\[
A_m\ \text{Morita 等价于}\ \CC^{F_{m+2}},
\qquad
A_m\otimes\mathcal K\cong \mathcal K^{\oplus F_{m+2}},
\]
其中 \(\mathcal K\) 表示可分无限维 Hilbert 空间上的紧算子代数。并且
\[
K_0(A_m)\cong \ZZ^{F_{m+2}},
\qquad
K_1(A_m)=0.
\]
若 \(m\neq n\)，则 \(A_m\) 与 \(A_n\) 不 Morita 等价。换言之，bin-fold 纤维群胚代数的 Morita 类仅由块数 \(|X_m|=F_{m+2}\) 决定，而不再记录各矩阵块的尺寸分布。
\end{theorem}
\begin{proof}
由定理 \ref{thm:xi-foldbin-groupoid-morita-k-center-dimension-lb}，
\[
A_m\cong \bigoplus_{w\in X_m}M_{d_m^{\mathrm{bin}}(w)}(\CC),
\qquad
K_0(A_m)\cong \ZZ^{F_{m+2}},
\qquad
K_1(A_m)=0.
\]
由于每个矩阵代数 \(M_d(\CC)\) 与 \(\CC\) Morita 等价，故有限直和 \(A_m\) 与
\[
\bigoplus_{w\in X_m}\CC
\cong
\CC^{F_{m+2}}
\]
Morita 等价。

再由稳定化恒等式 \(M_d(\CC)\otimes\mathcal K\cong\mathcal K\)，得到
\[
A_m\otimes\mathcal K
\cong
\bigoplus_{w\in X_m}\bigl(M_{d_m^{\mathrm{bin}}(w)}(\CC)\otimes\mathcal K\bigr)
\cong
\mathcal K^{\oplus F_{m+2}}.
\]
最后，定理 \ref{thm:xi-foldbin-groupoid-morita-k-center-dimension-lb} 同时给出
\[
Z(A_m)\cong \CC^{F_{m+2}}.
\]
Morita 等价保持中心的同构型，而 Fibonacci 数列严格递增，故 \(m\neq n\) 时 \(F_{m+2}\neq F_{n+2}\)，从而 \(Z(A_m)\not\cong Z(A_n)\)。于是 \(A_m\) 与 \(A_n\) 不可能 Morita 等价。
\end{proof}

\begin{theorem}[fold-gauge 常数列逐阶恢复全部偶 Bernoulli 数与偶值 \texorpdfstring{$\zeta(2r)$}{zeta(2r)}]\label{thm:xi-time-part60acb-gauge-constant-recovers-even-zeta}
\leanverified{paper\_xi\_time\_part60acb\_gauge\_constant\_recovers\_even\_zeta}
沿用定理 \ref{thm:fold-bin-gauge-constant-stirling-bernoulli-hierarchy} 的记号。定义
\[
\beta_r:=\left(\frac{\varphi}{2}\right)^{2r-1},
\qquad
\gamma_r:=\frac{\varphi^{-1}+\varphi^{2r-3}}{r(2r-1)}.
\]
则对每个整数 \(r\ge 1\)，都有显式极限公式
\[
B_{2r}
=
\frac{1}{\gamma_r}
\lim_{m\to\infty}
\beta_r^{-m}
\left(
\log C_m-\log(2\pi)-\sum_{j=1}^{r-1}\gamma_j B_{2j}\beta_j^m
\right).
\]
从而
\[
\zeta(2r)
=
(-1)^{r-1}\frac{(2\pi)^{2r}}{2(2r)!}\,B_{2r}.
\]
因此，单一序列 \((\log C_m)_{m\ge 1}\) 已足以逐阶唯一恢复全部偶 Bernoulli 数与全部偶值 \(\zeta(2r)\)。
\end{theorem}
\begin{proof}
由定理 \ref{thm:fold-bin-gauge-constant-stirling-bernoulli-hierarchy}，对任意固定 \(R\ge 1\) 有展开
\[
\log C_m
=
\log(2\pi)+
\sum_{r=1}^{R}
\gamma_r B_{2r}\beta_r^m
+
O\!\bigl(\beta_{R+1}^m\bigr).
\]
于是对固定的 \(r\ge 1\)，把前 \(r-1\) 个主项移到左端，便得到
\[
\log C_m-\log(2\pi)-\sum_{j=1}^{r-1}\gamma_j B_{2j}\beta_j^m
=
\gamma_r B_{2r}\beta_r^m+O(\beta_{r+1}^m).
\]
注意到
\[
0<\beta_{r+1}<\beta_r<1,
\qquad
\frac{\beta_{r+1}}{\beta_r}
=
\left(\frac{\varphi}{2}\right)^2<1.
\]
两边同乘 \(\beta_r^{-m}\) 后令 \(m\to\infty\)，便得到所示极限公式。

再由 Euler--Bernoulli 恒等式
\[
B_{2r}
=
(-1)^{r-1}\frac{2(2r)!}{(2\pi)^{2r}}\zeta(2r),
\]
即可反解出 \(\zeta(2r)\) 的表达式。结论成立。
\end{proof}

\begin{theorem}[纤维群胚代数的维数主项与 Morita 不可见的指数超载]\label{thm:xi-time-part60acb-groupoid-dimension-overload}
\leanverified{paper\_xi\_time\_part60acb\_groupoid\_dimension\_overload}
仍记
\[
A_m=\CC[\mathcal R_m^{\mathrm{bin}}]
\cong
\bigoplus_{w\in X_m}M_{d_m^{\mathrm{bin}}(w)}(\CC).
\]
则
\[
\dim_{\CC}A_m=\sum_{w\in X_m}\bigl(d_m^{\mathrm{bin}}(w)\bigr)^2.
\]
并且当 \(m\to\infty\) 时，
\[
\dim_{\CC}A_m
=
\frac{\varphi+\varphi^{-2}}{\sqrt5}
\left(\frac{4}{\varphi}\right)^m
+O(2^m).
\]
进一步，由
\[
\rank K_0(A_m)=F_{m+2}\sim \frac{\varphi^{m+2}}{\sqrt5}
\]
可得
\[
\frac{\dim_{\CC}A_m}{\rank K_0(A_m)}
=
\bigl(\varphi^{-1}+\varphi^{-4}\bigr)
\left(\frac{4}{\varphi^2}\right)^m
+o\!\left(\left(\frac{4}{\varphi^2}\right)^m\right).
\]
特别地，
\[
\frac{\dim_{\CC}A_m}{\rank K_0(A_m)}\longrightarrow +\infty,
\qquad
\frac{4}{\varphi^2}>1.
\]
因此，Morita--\(K\) 理论层面只记录 \(F_{m+2}\) 个简单块，而完整群胚代数的线性维数则以底数 \(4/\varphi\) 指数增长；二者之间存在一个 Morita 不可见的指数超载。
\end{theorem}
\begin{proof}
维数公式由 Wedderburn 分解直接给出：
\[
\dim_{\CC}A_m
=
\sum_{w\in X_m}\dim_{\CC}M_{d_m^{\mathrm{bin}}(w)}(\CC)
=
\sum_{w\in X_m}\bigl(d_m^{\mathrm{bin}}(w)\bigr)^2.
\]
另一方面，定理 \ref{thm:fold-bin-two-state-asymptotic} 给出一致两态渐近
\[
d_m^{\mathrm{bin}}(w)=\varphi^{-w_m}\left(\frac{2}{\varphi}\right)^m+O(1),
\]
且
\[
\#\{w\in X_m:\ w_m=0\}=F_{m+1},
\qquad
\#\{w\in X_m:\ w_m=1\}=F_m.
\]
于是
\[
\bigl(d_m^{\mathrm{bin}}(w)\bigr)^2
=
\varphi^{-2w_m}\left(\frac{2}{\varphi}\right)^{2m}
+O\!\left(\left(\frac{2}{\varphi}\right)^m\right)+O(1),
\]
从而
\[
\dim_{\CC}A_m
=
\left(\frac{2}{\varphi}\right)^{2m}
\bigl(F_{m+1}+\varphi^{-2}F_m\bigr)
+O\!\left(F_{m+2}\left(\frac{2}{\varphi}\right)^m\right)
+O(F_{m+2}).
\]
再由 Binet 公式，
\[
F_{m+1}+\varphi^{-2}F_m
=
\frac{\varphi+\varphi^{-2}}{\sqrt5}\,\varphi^m+O(\varphi^{-m}),
\]
故主项化为
\[
\left(\frac{2}{\varphi}\right)^{2m}
\cdot
\frac{\varphi+\varphi^{-2}}{\sqrt5}\,\varphi^m
=
\frac{\varphi+\varphi^{-2}}{\sqrt5}\left(\frac{4}{\varphi}\right)^m.
\]
又因
\[
F_{m+2}\left(\frac{2}{\varphi}\right)^m=O(2^m),
\qquad
F_{m+2}=O(\varphi^m)=o(2^m),
\]
便得到
\[
\dim_{\CC}A_m
=
\frac{\varphi+\varphi^{-2}}{\sqrt5}\left(\frac{4}{\varphi}\right)^m+O(2^m).
\]

最后用定理 \ref{thm:xi-foldbin-groupoid-morita-k-center-dimension-lb} 中的
\(
K_0(A_m)\cong \ZZ^{F_{m+2}}
\)
与
\(
F_{m+2}\sim \varphi^{m+2}/\sqrt5
\)
相除，即得
\[
\frac{\dim_{\CC}A_m}{\rank K_0(A_m)}
\sim
\frac{\varphi+\varphi^{-2}}{\varphi^2}
\left(\frac{4}{\varphi^2}\right)^m
=
\bigl(\varphi^{-1}+\varphi^{-4}\bigr)
\left(\frac{4}{\varphi^2}\right)^m.
\]
结论成立。
\end{proof}

\begin{theorem}[常纤维当且仅当自由有限群作用可实现]\label{thm:xi-time-part60acb-constant-fiber-free-group-action-criterion}
\leanverified{paper\_xi\_time\_part60acb\_constant\_fiber\_free\_group\_action\_criterion}
设 \(f:Y\twoheadrightarrow X\) 为有限集之间的满射，并记纤维直方图
\[
h(d):=\#\{x\in X:\ |f^{-1}(x)|=d\}
\qquad(d\ge 1).
\]
则下列条件等价：
\begin{enumerate}
  \item 等价关系 \(y\sim y'\iff f(y)=f(y')\) 是某个有限群在 \(Y\) 上自由作用的轨道划分；
  \item 存在整数 \(D\ge 1\) 使得 \(h(d)=0\) 对一切 \(d\neq D\) 成立，即全部纤维大小恒等于同一个常数 \(D\)。
\end{enumerate}
若进一步要求该划分来自有限阿贝尔群上的线性陪集划分，则同样必须满足上述常纤维条件。

应用于 window-\(6\) 的 bin-fold
\[
F_6:=\Fold_6^{\mathrm{bin}}:\Omega_6\to X_6,
\]
其纤维直方图
\[
2:8,\qquad 3:4,\qquad 4:9
\]
同时出现三种纤维大小，故 \(F_6\) 的纤维划分既不可能由任何有限群的自由作用实现，也不可能来自任何有限阿贝尔群的线性陪集划分。
\end{theorem}
\begin{proof}
若该等价关系来自某个有限群 \(H\) 在 \(Y\) 上的自由作用，则每个轨道都与 \(H\) 自身等势，故轨道大小恒等于 \(|H|\)。而轨道正是 \(f\) 的纤维，因此所有纤维大小都相等，得到 \(1\Rightarrow 2\)。

反之，若全部纤维大小都等于同一个常数 \(D\)，则取任一阶为 \(D\) 的有限群 \(H\)（例如循环群 \(\ZZ/D\ZZ\)）。对每个 \(x\in X\)，任选一个双射
\[
f^{-1}(x)\cong H,
\]
把 \(H\) 在自身上的左正则作用搬运到纤维 \(f^{-1}(x)\) 上，并在不同纤维上彼此独立地施行。于是得到 \(H\) 在 \(Y\) 上的自由作用，其轨道恰为 \(f\) 的各条纤维，从而 \(2\Rightarrow 1\)。

若该划分来自有限阿贝尔群上的线性陪集划分，则每个等价类都是某个固定子群陪集的原像，因而大小恒等于该子群的阶，仍必为常纤维。

最后，对 window-\(6\) 而言，纤维直方图同时出现 \(2,3,4\) 三种大小，故不满足常纤维条件。于是自由有限群作用实现与阿贝尔线性陪集实现均被排除。
\end{proof}

\noindent
由此，连续容量曲线、Morita 类型、Bernoulli--\(\zeta\) 谱链、完整群胚代数维数以及自由群商可实现性这五条看似分离的对象层结论，最终都被压缩到同一套纤维多重度数据之上：容量曲线完整恢复直方图，直方图决定规范群与 Wedderburn 分块，Morita 理论进一步只保留 Fibonacci 块数，维数主项则重新显露被 Morita 理论遗忘的指数级块尺寸负载，而自由群商解释是否可能则恰由直方图是否退化为单值常纤维来判定。

\paragraph{审计偶数窗的 Fibonacci 邻项分割与 window-\(6\) 最大纤维的算术阈值}
在审计偶数窗上的最小退化 Fibonacci 律、最小扇区补集与 Fold 极大纤维的跨模型锁定，以及 window-\(6\) 的单跳阶梯分类已经确定之后，还可把偶数窗的集合分割与最小窗口上的极大 bin-fold 纤维进一步压缩为下列两条直接结论。

\begin{theorem}[审计偶数窗中的 Fibonacci 邻项分割律]\label{thm:xi-time-part60ad-audited-even-fibonacci-neighbor-partition}
对审计窗口
\[
m\in\{6,8,10,12\},
\]
沿用定义 \ref{def:fold-bin-degeneracy-sectors} 的记号，记
\[
\mathcal S_{m,d_{\min}(m)}\subset X_m
\]
为 bin-fold 的最小退化扇区，并记 Fold 在层级 \(n\) 的最大纤维基数为
\[
D_n:=\max_{x\in X_n}d_n(x).
\]
则成立精确分割
\[
|X_m|
=
\bigl|\mathcal S_{m,d_{\min}(m)}\bigr|+D_{2m-2}.
\]
更精确地，三项恰由三枚相邻 Fibonacci 数给出：
\[
\bigl|\mathcal S_{m,d_{\min}(m)}\bigr|=F_m,
\qquad
D_{2m-2}=F_{m+1},
\qquad
|X_m|=F_{m+2}.
\]
因此在这些审计窗口中，最小退化扇区与其补集并非任意分裂，而是被一对相邻 Fibonacci 数严格锁定。
\end{theorem}
\begin{proof}
由定理 \ref{thm:fold-bin-min-degeneracy-fib-audited}，
\[
d_{\min}(m)=F_{m/2},
\qquad
\bigl|\mathcal S_{m,d_{\min}(m)}\bigr|=F_m=|X_{m-2}|.
\]
另一方面，推论 \ref{cor:fold-bin-minsector-complement-maxfiber-audited} 已给出
\[
\bigl|X_m\setminus \mathcal S_{m,d_{\min}(m)}\bigr|=D_{2m-2}.
\]
由稳定类型计数律 \(|X_n|=F_{n+2}\)，遂得
\[
D_{2m-2}
=
|X_m|-|X_{m-2}|
=
F_{m+2}-F_m
=
F_{m+1}.
\]
代回便得到
\[
|X_m|
=
\bigl|\mathcal S_{m,d_{\min}(m)}\bigr|+D_{2m-2}
=
F_m+F_{m+1}
=
F_{m+2}.
\]
证毕。
\end{proof}

\begin{theorem}[window-\(6\) 最大 bin-fold 纤维的算术阈值刻画]\label{thm:xi-time-part60ad-window6-max-binfold-arithmetic-threshold}
\leanverified{paper\_xi\_time\_part60ad\_window6\_max\_binfold\_arithmetic\_threshold}
在 \(m=6\) 时，令 \(w\in X_6\) 且记
\[
V:=V_6(w).
\]
则二进制区间折叠 \(\mathrm{Fold}^{\mathrm{bin}}_6\) 的最大纤维满足
\[
d_6^{\mathrm{bin}}(w)=4
\quad\Longleftrightarrow\quad
w_6=0\ \text{且}\ V\le 8.
\]
并且一旦上述条件成立，就有精确原像公式
\[
\bigl(\mathrm{Fold}^{\mathrm{bin}}_6\bigr)^{-1}(w)
=
\{V,\ V+F_8,\ V+F_9,\ V+F_{10}\}
=
\{V,\ V+21,\ V+34,\ V+55\}.
\]
此外，
\[
\#\{\,w\in X_6:\ d_6^{\mathrm{bin}}(w)=4\,\}=9.
\]
因此，window-\(6\) 中的最大退化并不在全部稳定类型上弥散出现，而是被阈值 \(V_6(w)\le 8\) 与末位条件 \(w_6=0\) 的算术截面完全选出。
\end{theorem}
\begin{proof}
由定理 \ref{thm:xi-foldbin-lastbit-tail-spectrum-single-jump-staircase}，在 \(m=6\) 时有精确三档公式
\[
d_6^{\mathrm{bin}}(w)=
\begin{cases}
2,& w_6=1,\\[4pt]
4,& w_6=0,\ V_6(w)\le 8,\\[4pt]
3,& w_6=0,\ V_6(w)>8.
\end{cases}
\]
故
\[
d_6^{\mathrm{bin}}(w)=4
\quad\Longleftrightarrow\quad
w_6=0\ \text{且}\ V_6(w)\le 8.
\]

再由表 \ref{tab:fold6_bin_uplift_choice_collapse} 的第一行，
\[
\Delta(w):=\{\,n-V_6(w):\ n\in(\mathrm{Fold}^{\mathrm{bin}}_6)^{-1}(w)\,\}
=
\{0,21,34,55\}
\]
恰当且仅当 \(w_6=0,\ V_6(w)\le 8\)。因此在上述条件下，
\[
\bigl(\mathrm{Fold}^{\mathrm{bin}}_6\bigr)^{-1}(w)
=
V+\Delta(w)
=
\{V,\ V+21,\ V+34,\ V+55\},
\]
即所示原像公式。

同一表的第一行同时列出了全部满足该条件的稳定词，共 \(9\) 个，故
\[
\#\{\,w\in X_6:\ d_6^{\mathrm{bin}}(w)=4\,\}=9.
\]
证毕。
\end{proof}

\noindent
由此，审计偶数窗中的最小退化扇区已不仅仅给出一个局部 Fibonacci 计数，它与更高层 Fold 极大纤维一起把 \(X_m\) 精确切分为两块相邻 Fibonacci 尺度；而在最小窗口 \(m=6\) 上，最大 bin-fold 纤维又被末位与 Zeckendorf 阈值的联合算术判据唯一选出。于是，全局审计分割与局部最大退化阈值在同一条离散 Fibonacci 链上闭合。

\paragraph{Jensen 缺陷反演、排斥半径幂律与共动缺陷剖面的 Fourier--Prony 线性化}
以第 \ref{subsec:dephysicalized-horizon-timefiber-nohair} 小节中的
Jensen 缺陷--排斥半径链、第 \ref{subsubsec:xi-horizon-spectral-measure-tomography}
小节中的共动 Cayley 层析，以及
Toeplitz/Hankel/Prony 的既有恢复定理为基础，还可把“径向零点账本”与“共动缺陷剖面”进一步压缩为下列五条专门化结论。

\begin{theorem}[Jensen 缺陷的分布反演与零点半径测度的完全恢复]\label{thm:xi-time-part60a-jensen-defect-radial-measure-inversion}
\leanverified{paper\_xi\_time\_part60a\_jensen\_defect\_radial\_measure\_inversion}
沿用定理 \ref{thm:dephys-defect-repulsion-radius} 与命题
\ref{prop:dephys-jensen-defect-log-derivative} 的记号。设
\(F_{\zeta}\) 的盘内零点按重数记为 \(\{a\}\subset\mathbb D\)，并定义
\((-\infty,0)\) 上的离散正测度
\[
\mu_{\zeta}:=
\sum_{a:\,F_{\zeta}(a)=0}\mathrm{mult}(a)\,\delta_{\log|a|}.
\]
再令
\[
u:=\log\rho<0,
\qquad
\mathcal D_{\zeta}(u):=D_{\zeta}(e^u).
\]
则对每个 \(u<0\) 都有严格积分表示
\begin{equation}\label{eq:xi-time-part60a-jensen-defect-measure-potential}
\boxed{\;
\mathcal D_{\zeta}(u)=\int_{(-\infty,u)}(u-s)\,\dd\mu_{\zeta}(s).\;}
\end{equation}
从而：
\begin{enumerate}
  \item \(\mathcal D_{\zeta}\) 在 \((-\infty,0)\) 上为凸、分段仿射函数；
  \item 几乎处处有
  \[
  \mathcal D_{\zeta}'(u)=N_{F_{\zeta}}(e^u);
  \]
  \item 在分布意义下成立
  \[
  \boxed{\;
  \mathcal D_{\zeta}''=\mu_{\zeta}.\;}
  \]
\end{enumerate}
换言之，Jensen 缺陷的分布二阶导数恰等于零点半径多重集所定义的离散测度。
\end{theorem}
\begin{proof}
由 Jensen 求和式与命题 \ref{prop:dephys-jensen-defect-log-derivative} 的记号，
\[
D_{\zeta}(\rho)=\sum_{|a|<\rho}\log\frac{\rho}{|a|}.
\]
代入 \(\rho=e^u\) 得
\[
\mathcal D_{\zeta}(u)
=
\sum_{\log|a|<u}\bigl(u-\log|a|\bigr),
\]
而右端正是 \eqref{eq:xi-time-part60a-jensen-defect-measure-potential}
中的离散势函数表示。每一项 \((u-s)_+\) 关于 \(u\) 为凸且分段仿射，故其局部有限和
\(\mathcal D_{\zeta}\) 亦如此。

在可导点，
\[
\frac{\dd}{\dd u}(u-\log|a|)_+=\mathbf 1_{\{u>\log|a|\}},
\]
于是
\[
\mathcal D_{\zeta}'(u)
=
\sum_a \mathbf 1_{\{u>\log|a|\}}
=
N_{F_{\zeta}}(e^u),
\]
这与命题 \ref{prop:dephys-jensen-defect-log-derivative} 完全一致。
再对该分段常值函数取分布导数，即在每个跳点 \(\log|a|\) 处产生权重
\(\mathrm{mult}(a)\) 的 Dirac 质量，从而
\[
\mathcal D_{\zeta}''=
\sum_a \mathrm{mult}(a)\,\delta_{\log|a|}
=\mu_{\zeta}.
\]
证毕。
\end{proof}

\begin{theorem}[排斥半径的分段幂律正规形与整数斜率层级]\label{thm:xi-time-part60a-repulsion-radius-piecewise-power-law}
\leanverified{paper\_xi\_time\_part60a\_repulsion\_radius\_piecewise\_power\_law}
沿用定理 \ref{thm:dephys-defect-repulsion-radius} 的排斥半径记号
\[
r_{\ast}(\rho):=\rho\,e^{-D_{\zeta}(\rho)}\qquad(0<\rho<1).
\]
设 \(F_{\zeta}\) 的盘内零点模长的不同取值按严格递增排列为
\[
0<\varrho_1<\varrho_2<\cdots<\varrho_J<1,
\]
并记模长恰为 \(\varrho_j\) 的零点总重数为 \(\mu_j\)，再定义累计重数
\[
M_j:=\mu_1+\cdots+\mu_j,
\qquad
M_0:=0.
\]
约定 \(\varrho_0:=0\)、\(\varrho_{J+1}:=1\)。则对每个
\(j=0,1,\dots,J\) 与每个 \(\rho\in(\varrho_j,\varrho_{j+1})\)，都有严格闭式
\begin{equation}\label{eq:xi-time-part60a-repulsion-piecewise-power-law}
\boxed{\;
r_{\ast}(\rho)=
\rho^{\,1-M_j}
\prod_{i=1}^{j}\varrho_i^{\,\mu_i}.\;}
\end{equation}
因此：
\begin{enumerate}
  \item 在每个零点模长间隙 \((\varrho_j,\varrho_{j+1})\) 上，\(r_{\ast}\) 是幂函数；
  \item 其对数斜率恰为整数 \(1-M_j\)；
  \item 尤其有
  \[
  r_{\ast}(\rho)=\rho\qquad(0<\rho<\varrho_1),
  \]
  若 \(\mu_1=1\)，则
  \[
  r_{\ast}(\rho)\equiv \varrho_1
  \qquad(\varrho_1<\rho<\varrho_2),
  \]
  而一旦 \(M_j\ge 2\)，则 \(r_{\ast}\) 在
  \((\varrho_j,\varrho_{j+1})\) 上严格递减。
\end{enumerate}
\end{theorem}
\begin{proof}
对 \(\rho\in(\varrho_j,\varrho_{j+1})\)，恰有前 \(j\) 个模长层进入 Jensen 求和，故
\[
D_{\zeta}(\rho)
=
\sum_{i=1}^{j}\mu_i\log\frac{\rho}{\varrho_i}
=
M_j\log\rho-\sum_{i=1}^{j}\mu_i\log\varrho_i.
\]
代回 \(r_{\ast}(\rho)=\rho e^{-D_{\zeta}(\rho)}\) 即得
\eqref{eq:xi-time-part60a-repulsion-piecewise-power-law}。
对数求导得到
\[
\frac{\dd}{\dd(\log\rho)}\log r_{\ast}(\rho)=1-M_j,
\]
这与推论 \ref{cor:dephys-repulsion-radius-zero-count-driving} 在该区间上
\(N_{F_{\zeta}}(\rho)=M_j\) 的表达一致。由 \(M_j\) 的整数取值，三种单调性结论立即推出。证毕。
\end{proof}

\begin{theorem}[有限 Jensen 证书的绝对盲区]\label{thm:xi-time-part60a-finite-jensen-certificate-blindspot}
\leanverified{paper\_xi\_time\_part60a\_finite\_jensen\_certificate\_blindspot}
固定任意有限半径采样集
\[
0<\rho_1<\rho_2<\cdots<\rho_M<1,
\]
并对单位圆盘内解析对象 \(F\) 定义 Jensen 缺陷样本向量
\[
\mathbf D_{\boldsymbol\rho}(F):=
\bigl(D_F(\rho_1),\dots,D_F(\rho_M)\bigr)\in\RR^M.
\]
则存在两个有限 Blaschke 乘积 \(B_0,B_1\) 满足：
\begin{enumerate}
  \item \(B_0\) 在 \(\mathbb D\) 内无零点；
  \item \(B_1\) 在 \(\mathbb D\) 内恰有一个零点 \(a\)，并且 \(|a|>\rho_M\)；
  \item 但二者具有完全相同的有限 Jensen 证书：
  \[
  \boxed{\;
  \mathbf D_{\boldsymbol\rho}(B_0)=\mathbf D_{\boldsymbol\rho}(B_1).\;}
  \]
\end{enumerate}
因此，任何仅依赖有限半径样本
\(\{D_{\zeta}(\rho_j)\}_{j=1}^{M}\) 且 \(\rho_M<1\) 的判定协议，都不可能决定盘内是否存在零点；从而不能作为 RH 的决定性有限证书。
\end{theorem}
\begin{proof}
取
\[
B_0(w)\equiv 1.
\]
再任取满足 \(\rho_M<|a|<1\) 的点 \(a\in\mathbb D\)，并令
\[
B_1(w):=\frac{w-a}{1-\overline a\,w}.
\]
这是一个恰有单零点 \(a\) 的有限 Blaschke 因子。对任意
\(\rho<|a|\)，Jensen 缺陷只对 \(|z|<\rho\) 内的零点求和，因此
\[
D_{B_1}(\rho)=0.
\]
另一方面，\(B_0\) 在盘内无零点，故
\[
D_{B_0}(\rho)=0.
\]
由于 \(\rho_j<|a|\) 对全部 \(j=1,\dots,M\) 都成立，遂有
\[
D_{B_0}(\rho_j)=D_{B_1}(\rho_j)=0
\qquad (1\le j\le M),
\]
即
\(
\mathbf D_{\boldsymbol\rho}(B_0)=\mathbf D_{\boldsymbol\rho}(B_1)
\)。
但二者的盘内零点性质完全不同，故任何仅依赖该有限样本向量的协议都无法判定“盘内是否有零点”。对圆盘完成化对象 \(F_{\zeta}\) 而言，这正排除了仅用有限 Jensen 样本决定 RH 的可能性。证毕。
\end{proof}

\begin{theorem}[共动缺陷剖面的 Fourier--Prony 线性化]\label{thm:xi-time-part60a-comoving-defect-fourier-prony-linearization}
沿用 \S\ref{subsubsec:xi-horizon-spectral-measure-tomography} 的有限点缺陷模型。
设离切片零点为
\[
\rho_j=\beta_j+\mathrm{i}\gamma_j,
\qquad
\delta_j:=\frac12-\beta_j>0,
\qquad
1\le j\le N,
\]
并记对应的共动 Cayley 像为
\[
w_j(x_0):=\mathcal C\bigl((\gamma_j-x_0)+\mathrm{i}\delta_j\bigr)\in\mathbb D.
\]
定义缺陷剖面
\[
h(x_0):=\sum_{j=1}^{N}m_j\bigl(1-|w_j(x_0)|^2\bigr),
\qquad
m_j\in\NN.
\]
则其 Fourier 变换
\[
\widehat h(\xi):=\int_{\RR}h(x_0)e^{-\mathrm{i}\xi x_0}\,\dd x_0
\]
满足完全显式分解
\begin{equation}\label{eq:xi-time-part60a-comoving-defect-fourier-prony}
\boxed{\;
\widehat h(\xi)
=
4\pi\sum_{j=1}^{N}
\frac{m_j\delta_j}{1+\delta_j}\,
e^{-(1+\delta_j)|\xi|}\,
e^{-\mathrm{i}\gamma_j\xi},
\qquad \xi\in\RR.\;}
\end{equation}
特别地，取单位采样步长并定义
\[
s_n:=\widehat h(n)
\qquad (n\in\NN_0),
\]
再记
\[
\lambda_j:=e^{-(1+\delta_j)-\mathrm{i}\gamma_j},
\qquad
c_j:=4\pi\,\frac{m_j\delta_j}{1+\delta_j},
\]
则有有限指数和表示
\begin{equation}\label{eq:xi-time-part60a-comoving-defect-samples-prony}
\boxed{\;
s_n=\sum_{j=1}^{N}c_j\,\lambda_j^{\,n}.\;}
\end{equation}
因此：
\begin{enumerate}
  \item \((s_n)_{n\ge 0}\) 满足阶数至多 \(N\) 的常系数线性递推；
  \item 其无穷 Hankel 矩阵 \(H:=(s_{i+j})_{i,j\ge 0}\) 的秩不超过互异 \(\lambda_j\) 的个数；
  \item 若 \(\lambda_1,\dots,\lambda_N\) 两两不同，则 \(\operatorname{rank}(H)=N\)。
\end{enumerate}
\end{theorem}
\begin{proof}
由 \eqref{eq:xi-comoving-scan-lorentz-peak}，
\[
h(x_0)=
\sum_{j=1}^{N}
\frac{4m_j\delta_j}{(\gamma_j-x_0)^2+(1+\delta_j)^2}.
\]
对每一项取平移变量 \(t=x_0-\gamma_j\)，并应用标准积分公式
\[
\int_{\RR}\frac{e^{-\mathrm{i}\xi t}}{t^2+a^2}\,\dd t
=
\frac{\pi}{a}e^{-a|\xi|}
\qquad(a>0),
\]
即得 \eqref{eq:xi-time-part60a-comoving-defect-fourier-prony}。
再取 \(\xi=n\in\NN_0\)，便得到
\eqref{eq:xi-time-part60a-comoving-defect-samples-prony}。

有限指数和自动满足以
\[
Q(z):=\prod_{j=1}^{N}(z-\lambda_j)
\]
为特征多项式的常系数线性递推。另一方面，
\[
H=(s_{i+j})_{i,j\ge 0}
=
V\,\mathrm{diag}(c_1,\dots,c_N)\,V^{\mathsf T},
\qquad
V_{ij}:=\lambda_j^{\,i},
\]
故 \(\operatorname{rank}(H)\) 不超过互异 \(\lambda_j\) 的个数；若
\(\lambda_j\) 两两不同，则相应 Vandermonde 矩阵列满秩，从而
\(\operatorname{rank}(H)=N\)。证毕。
\end{proof}

\begin{corollary}[有限缺陷共动谱的最小完备采样阈值]\label{cor:xi-time-part60a-comoving-defect-minimal-sampling-sharp}
\leanverified{paper\_xi\_time\_part60a\_comoving\_defect\_minimal\_sampling\_sharp}
在定理 \ref{thm:xi-time-part60a-comoving-defect-fourier-prony-linearization}
的设定下，进一步假设 \(\lambda_1,\dots,\lambda_N\) 两两不同。则：
\begin{enumerate}
  \item 任意长度小于 \(2N\) 的初始样本段
  \[
  s_0,\dots,s_{2N-2}
  \]
  在一般位置上不足以唯一恢复全部 \((\lambda_j,c_j)_{j=1}^{N}\)；
  \item 长度恰为 \(2N\) 的初始样本段
  \[
  s_0,\dots,s_{2N-1}
  \]
  在一般位置下已足以唯一恢复 \((\lambda_j,c_j)_{j=1}^{N}\)（忽略排列）；
  \item 因而在给定高度先验窗口长度为 \(2\pi\) 的条件下，可唯一恢复
  \((\gamma_j,\delta_j,m_j)_{j=1}^{N}\)。
\end{enumerate}
\end{corollary}
\begin{proof}
将定理 \ref{thm:xi-time-part60a-comoving-defect-fourier-prony-linearization}
中的指数和表示
\(
s_n=\sum_{j=1}^{N}c_j\lambda_j^n
\)
专门化到推论 \ref{cor:xi-prony-decision-vs-reconstruction-one-sample-gap}，
取其中的节点 \(a_j=\lambda_j\)、权重 \(\omega_j=c_j\)，便得到前两条结论。

对第三条，由
\[
|\lambda_j|=e^{-(1+\delta_j)},
\qquad
\arg(\lambda_j)\equiv -\gamma_j\pmod{2\pi},
\]
可恢复
\[
\delta_j=-\log|\lambda_j|-1,
\qquad
\gamma_j\equiv-\arg(\lambda_j)\pmod{2\pi}.
\]
若事先知道 \(\gamma_j\) 落在某个固定长度为 \(2\pi\) 的区间内，则该同余类具有唯一代表。
最后由
\[
m_j=
\frac{c_j}{4\pi}\cdot\frac{1+\delta_j}{\delta_j}
\]
恢复重数 \(m_j\)。证毕。
\end{proof}

\begin{remark}[有限原子审计接口]\label{rem:xi-time-part60a-jensen-comoving-prony-audit}
脚本 \texttt{scripts/exp\_xi\_jensen\_comoving\_prony\_audit.py}
给出一个有限 Blaschke/有限缺陷样本上的逐项核验：
它同时检查 Jensen 缺陷的离散势表示、排斥半径的分段幂律、
有限半径采样的绝对盲区，以及单位步长共动剖面的 Fourier--Prony
恢复。
\end{remark}

\paragraph{Fourier--Joukowsky 干涉曲线、奇异环接触计数与 Fibonacci 采样谱}
在共动缺陷剖面的 Fourier--Prony 线性化、碰撞谱的单位圆采样表达以及 Joukowsky 门的既有解析接口都已建立之后，还可把若干看似分离的平面参数曲线统一理解为“单位圆旋转轨道经 Joukowsky 门读出”的有限干涉对象，并把这一几何接口直接接到 Fibonacci 生成多项式 \(D_m\) 的碰撞谱链上。

\begin{theorem}[Fourier--Joukowsky 生成原理：玫瑰线与 Lissajous 曲线作为概率特征曲线的单位圆外推]\label{thm:xi-time-part60ae-fourier-joukowsky-generation-principle}
设
\[
\mu=\sum_{j=1}^{N}p_j\delta_{\lambda_j}
\]
为有限支撑概率分布，其中 \(p_j\ge 0\)、\(\sum_j p_j=1\)。定义其特征函数
\[
\widehat\mu(t):=\int_{\RR}e^{itx}\,\dd\mu(x)=\sum_{j=1}^{N}p_j e^{i\lambda_j t}.
\]
则轨道
\[
t\longmapsto \widehat\mu(t)\in\CC
\]
恰是若干单位圆匀速旋转的凸组合；反之，任意形如
\[
\sum_{j=1}^{N}p_j e^{i\lambda_j t}
\]
的三角多项式，只要系数非负且和为 \(1\)，都可实现为某个有限支撑概率分布的特征函数轨道。

进一步，定义归一化 Joukowsky 门
\[
J:S^1\to [-1,1],
\qquad
J(z):=\frac{z+z^{-1}}{2}=\Re z.
\]
则有：
\begin{enumerate}
  \item 对任意有理斜率 \(n/d>0\)，玫瑰线的复表示
  \[
  z(\theta):=\cos\!\Bigl(\frac{n}{d}\theta\Bigr)e^{i\theta}
  \]
  满足严格恒等式
  \[
  z(\theta)=\frac12\Bigl(e^{i(1+n/d)\theta}+e^{i(1-n/d)\theta}\Bigr),
  \]
  因而正是二点分布
  \[
  \frac12\bigl(\delta_{1+n/d}+\delta_{1-n/d}\bigr)
  \]
  的特征函数轨道；
  \item 对任意整数 \(a,b\ge 1\) 与相位 \(\delta\in\RR\)，Lissajous 曲线
  \[
  (x(t),y(t))=\bigl(\cos(at+\delta),\cos(bt)\bigr)
  \]
  可写为
  \[
  \gamma_{a,b,\delta}(t):=\bigl(e^{i(at+\delta)},e^{ibt}\bigr)\in S^1\times S^1
  \]
  经坐标门 \((J,J)\) 的像：
  \[
  (x(t),y(t))
  =
  \bigl(J(\gamma_{a,b,\delta}^{(1)}(t)),J(\gamma_{a,b,\delta}^{(2)}(t))\bigr).
  \]
\end{enumerate}
\end{theorem}
\begin{proof}
第一条只是特征函数定义的直接展开；反向上取离散随机变量 \(X\) 满足
\[
\PP(X=\lambda_j)=p_j,
\]
便得到
\[
\EE[e^{itX}]=\sum_{j=1}^{N}p_j e^{i\lambda_j t}.
\]

对玫瑰线，只需用 Euler 公式
\[
\cos\alpha=\frac{e^{i\alpha}+e^{-i\alpha}}{2}
\]
展开：
\[
\cos\!\Bigl(\frac{n}{d}\theta\Bigr)e^{i\theta}
=
\frac12\Bigl(e^{i(1+n/d)\theta}+e^{i(1-n/d)\theta}\Bigr).
\]

对 Lissajous 曲线，由
\[
J(e^{i\vartheta})=\frac{e^{i\vartheta}+e^{-i\vartheta}}{2}=\cos\vartheta
\]
可知
\[
J(\gamma_{a,b,\delta}^{(1)}(t))=\cos(at+\delta),
\qquad
J(\gamma_{a,b,\delta}^{(2)}(t))=\cos(bt).
\]
于是所示分解成立。证毕。
\end{proof}

\begin{proposition}[奇异环作为 Joukowsky 门的分歧环及其接触次数的完全计数]\label{prop:xi-time-part60ae-lissajous-singular-ring-contact-count}
\leanverified{paper\_xi\_time\_part60ae\_lissajous\_singular\_ring\_contact\_count}
沿用定理 \ref{thm:xi-time-part60ae-fourier-joukowsky-generation-principle} 的归一化 Joukowsky 门 \(J\)。则乘积映射
\[
J\times J:S^1\times S^1\to [-1,1]^2
\]
的 Jacobian 退化集合为
\[
\Sigma:=
\bigl(\{\pm 1\}\times S^1\bigr)\cup \bigl(S^1\times \{\pm 1\}\bigr),
\]
其像恰为正方形边界
\[
\partial([-1,1]^2).
\]

对闭测地线
\[
\gamma_{a,b,\delta}(t)=\bigl(e^{i(at+\delta)},e^{ibt}\bigr)
\qquad (t\in[0,2\pi))
\]
定义奇异环接触集合
\[
\mathcal T_{a,b,\delta}:=
\{t\in[0,2\pi):\gamma_{a,b,\delta}(t)\in\Sigma\}.
\]
再定义角点重合项
\[
\mathcal C_{a,b,\delta}
:=
\Bigl\{
t:\ e^{i(at+\delta)}\in\{\pm 1\},\ e^{ibt}\in\{\pm 1\}
\Bigr\}.
\]
则有精确计数公式
\[
\boxed{\;
|\mathcal T_{a,b,\delta}|=2a+2b-|\mathcal C_{a,b,\delta}|.\;}
\]
并且若记 \(g:=\gcd(a,b)\)，则
\[
\mathcal C_{a,b,\delta}\neq\varnothing
\iff
\delta\in \frac{\pi}{g}\ZZ,
\]
且在此情形下
\[
\boxed{\;
|\mathcal C_{a,b,\delta}|=2g.\;}
\]
\end{proposition}
\begin{proof}
由
\[
J(e^{i\vartheta})=\cos\vartheta
\]
可知 \(\mathrm dJ\) 退化当且仅当 \(\sin\vartheta=0\)，亦即 \(e^{i\vartheta}\in\{\pm 1\}\)。故 \(J\times J\) 的分歧环正是
\[
\Sigma=
\bigl(\{\pm 1\}\times S^1\bigr)\cup \bigl(S^1\times \{\pm 1\}\bigr),
\]
而其像显然是 \([-1,1]^2\) 的边界。

对接触计数，条件 \(\gamma_{a,b,\delta}(t)\in\Sigma\) 等价于
\[
at+\delta\in\pi\ZZ
\qquad\text{或}\qquad
bt\in\pi\ZZ.
\]
在区间 \([0,2\pi)\) 内，前者恰有 \(2a\) 个解，后者恰有 \(2b\) 个解，故由容斥得
\[
|\mathcal T_{a,b,\delta}|=2a+2b-|\mathcal C_{a,b,\delta}|.
\]

现求 \(|\mathcal C_{a,b,\delta}|\)。它等价于同余系统
\[
at\equiv -\delta\pmod{\pi},
\qquad
bt\equiv 0\pmod{\pi}.
\]
令
\[
a=ga',
\qquad
b=gb',
\qquad
\gcd(a',b')=1.
\]
第二式给出
\[
t=\frac{\pi k}{b}
\qquad (k\in\ZZ).
\]
代回第一式，得到
\[
\frac{a}{b}\pi k\equiv -\delta\pmod{\pi}
\iff
\frac{a'}{b'}k\equiv -\frac{\delta}{\pi}\pmod{1}.
\]
由于 \(a'\) 与 \(b'\) 互素，该同余有解当且仅当
\[
\frac{\delta}{\pi}\in \frac{1}{g}\ZZ,
\]
即
\[
\delta\in\frac{\pi}{g}\ZZ.
\]
一旦可解，\(k\) 在模 \(2b'=2b/g\) 的剩余类中恰有 \(2\) 个解，而每个剩余类在 \([0,2\pi)\) 上对应 \(g\) 个不同的 \(t\)。故总数为
\[
|\mathcal C_{a,b,\delta}|=2g.
\]
证毕。
\end{proof}

\begin{theorem}[Fibonacci 子集和的单位圆干涉曲线与离散碰撞谱的模采样]\label{thm:xi-time-part60ae-fibonacci-interference-curve-sampling}
沿用定理 \ref{thm:fold-collision-spectrum} 与定义 \ref{def:fold-normalized-amplitude} 的记号，记
\[
M:=F_{m+2},
\qquad
D_m(x)\equiv \prod_{j=1}^{m}\bigl(1+x^{F_{j+1}}\bigr)\pmod{x^M-1},
\]
\[
\widehat d_m(k)=D_m(e^{-2\pi i k/M}),
\qquad
a_m(k):=\frac{|\widehat d_m(k)|}{2^m}.
\]
定义连续单位圆干涉曲线
\[
\Phi_m(t):=2^{-m}D_m(e^{it})
=
2^{-m}\prod_{j=1}^{m}\bigl(1+e^{iF_{j+1}t}\bigr),
\qquad
t\in\RR.
\]
则成立：
\begin{enumerate}
  \item 离散碰撞谱正是 \(\Phi_m\) 在单位圆等分网格上的模采样：
  \[
  \boxed{\;
  a_m(k)=\bigl|\Phi_m(2\pi k/M)\bigr|
  \qquad (k\in \ZZ/(M\ZZ)).\;}
  \]
  \item 若定义零点集合
  \[
  \mathcal S_m:=\{e^{it}\in S^1:\ \Phi_m(t)=0\},
  \]
  则
  \[
  \boxed{\;
  \mathcal S_m=
  \bigcup_{j=1}^{m}
  \left\{
  e^{i(2\ell+1)\pi/F_{j+1}}:\ 0\le \ell\le F_{j+1}-1
  \right\}.\;}
  \]
  \item 并集
  \[
  \bigcup_{m\ge 1}\mathcal S_m
  \]
  在 \(S^1\) 中稠密。
\end{enumerate}
\end{theorem}
\begin{proof}
由
\[
\widehat d_m(k)=D_m(e^{-2\pi i k/M})
\]
与
\[
a_m(k)=\frac{|\widehat d_m(k)|}{2^m}
\]
可知
\[
a_m(k)=2^{-m}\bigl|D_m(e^{-2\pi i k/M})\bigr|
=
\bigl|\Phi_m(-2\pi k/M)\bigr|.
\]
又因为取模后对 \(k\mapsto -k\) 不变，故亦可写成
\[
a_m(k)=\bigl|\Phi_m(2\pi k/M)\bigr|.
\]

对零点集合，由乘积表达式，
\[
\Phi_m(t)=0
\iff
\exists\,j\in\{1,\dots,m\}\ \text{使}\ 1+e^{iF_{j+1}t}=0.
\]
这等价于
\[
F_{j+1}t\equiv \pi \pmod{2\pi},
\]
亦即
\[
t\equiv \frac{(2\ell+1)\pi}{F_{j+1}}
\pmod{2\pi}
\qquad (\ell\in\ZZ).
\]
投影到单位圆便得到 \(\mathcal S_m\) 的显式表达。

最后证稠密性。任取单位圆上的非空弧段 \(I\subset S^1\)，其角长度为某个 \(\eta>0\)。取 \(n\) 充分大，使
\[
\frac{2\pi}{F_{n+1}}<\eta.
\]
则步长为 \(2\pi/F_{n+1}\) 的奇半格点
\[
\frac{(2\ell+1)\pi}{F_{n+1}}
\qquad (0\le \ell\le F_{n+1}-1)
\]
在圆周上的相邻间距不超过 \(\eta\)，故必有一个点落入 \(I\)。该点属于 \(\mathcal S_n\)，从而
\[
I\cap \bigcup_{m\ge 1}\mathcal S_m\neq \varnothing.
\]
故该并集在 \(S^1\) 中稠密。证毕。
\end{proof}

\begin{theorem}[准 Lissajous 轨道接近角点奇异环的 Diophantine 定量律]\label{thm:xi-time-part60ae-diophantine-singular-ring-approximation}
设 \(\alpha>0\) 为无理数，并考虑环面直线流
\[
\gamma_{\alpha}(t):=(e^{it},e^{i\alpha t})\in S^1\times S^1.
\]
记角点奇异环为
\[
\Sigma_{\square}:=\{\pm 1\}\times \{\pm 1\}.
\]
设 \(p_n/q_n\) 为 \(\alpha\) 的连分数主收敛。定义时刻
\[
t_n:=\pi q_n.
\]
则第一坐标严格满足
\[
e^{it_n}=(-1)^{q_n}\in\{\pm 1\},
\]
而第二坐标具有近角点表示
\[
e^{i\alpha t_n}=(-1)^{p_n}e^{i\pi(\alpha q_n-p_n)},
\qquad
|\alpha q_n-p_n|<\frac{1}{q_{n+1}}.
\]
从而
\[
\boxed{\;
\bigl|\cos(\alpha t_n)-(-1)^{p_n}\bigr|
=
O\!\Bigl(\frac{1}{q_{n+1}^2}\Bigr).\;}
\]
换言之，\(\gamma_{\alpha}\) 在时刻 \(t_n\) 以二次量级贴近角点奇异环 \(\Sigma_{\square}\)，而贴近速率完全由下一收敛分母 \(q_{n+1}\) 控制。

特别地，当 \(\alpha=\varphi\) 时，
\[
q_n=F_n,
\]
故该贴近序列按 Fibonacci 尺度分层；其误差层级与定理 \ref{thm:fold-critical-resonance-constant} 中黄金比共振乘积
\[
C_{\varphi}(u)=\prod_{n\ge 2}\Bigl|\cos\bigl(\pi u\varphi^{-n}\bigr)\Bigr|
\]
所依赖的小角衰减阶具有同源的 Fibonacci 逼近结构。
\end{theorem}
\begin{proof}
连分数理论给出
\[
\left|\alpha-\frac{p_n}{q_n}\right|
<
\frac{1}{q_nq_{n+1}},
\]
故
\[
|\alpha q_n-p_n|<\frac{1}{q_{n+1}}.
\]
取
\[
t_n=\pi q_n,
\]
则
\[
e^{it_n}=e^{i\pi q_n}=(-1)^{q_n}\in\{\pm 1\}.
\]
另一方面，
\[
\alpha t_n=\pi\alpha q_n=\pi p_n+\pi(\alpha q_n-p_n),
\]
于是
\[
e^{i\alpha t_n}=(-1)^{p_n}e^{i\pi(\alpha q_n-p_n)}.
\]
从而
\[
\cos(\alpha t_n)
=
\cos\bigl(\pi p_n+\pi(\alpha q_n-p_n)\bigr)
=
(-1)^{p_n}\cos\bigl(\pi(\alpha q_n-p_n)\bigr).
\]
令
\[
\varepsilon_n:=\alpha q_n-p_n.
\]
由 Taylor 展开
\[
\cos(\pi\varepsilon_n)=1-\frac{\pi^2}{2}\varepsilon_n^2+O(\varepsilon_n^4)
\]
可得
\[
\bigl|\cos(\alpha t_n)-(-1)^{p_n}\bigr|
=
O(\varepsilon_n^2)
=
O\!\Bigl(\frac{1}{q_{n+1}^2}\Bigr).
\]
当 \(\alpha=\varphi\) 时，主收敛分母即为 Fibonacci 数 \(F_n\)，故误差分层按 Fibonacci 尺度组织。最后一条仅指出：\(C_{\varphi}(u)\) 的每一因子都由角尺度 \(\varphi^{-n}\) 控制，而 \(\varphi\) 的 Diophantine 最优逼近同样由 Fibonacci 分母主导，故二者共享同一误差层级来源。证毕。
\end{proof}

\begin{conjecture}[奇异环稠密化诱导碰撞谱异常质量]\label{conj:xi-time-part60ae-singular-ring-densification-collision-anomaly}
沿用定理 \ref{thm:xi-time-part60ae-fibonacci-interference-curve-sampling} 的记号，定义模值经验测度
\[
\nu_m:=
\frac{1}{F_{m+2}}
\sum_{k=0}^{F_{m+2}-1}
\delta_{\left|\Phi_m(2\pi k/F_{m+2})\right|}.
\]
预期存在一个非原子极限测度 \(\nu\) 使
\[
\nu_m\Rightarrow \nu
\qquad (m\to\infty),
\]
并且其近零质量满足以下统一原则：对小参数 \(\varepsilon\downarrow 0\)，
\[
\nu([0,\varepsilon])
\]
的主导衰减速率由奇异环并集
\[
\bigcup_{m\ge 1}\mathcal S_m
\]
在 \(S^1\) 上的加密速率所支配，而该主导项应可用黄金比共振常数族 \(C_{\varphi}(u)\) 的小角级联近似来刻画。

换言之，Fibonacci 干涉曲线的奇异环稠密化应当是离散碰撞谱异常质量的统一几何来源：单位圆上的层级稠密零环一旦通过等分采样被读出，便必然在 \(\{a_m(k)\}\) 的经验分布中留下可审计的近零异常质量。
\end{conjecture}

\noindent
由此，有限支撑特征函数、玫瑰线、Lissajous 曲线与 Fibonacci 子集和多项式 \(D_m\) 不再是彼此分离的对象：它们都可统一理解为单位圆旋转经 Joukowsky 门读出的有限干涉曲线。对于折叠谱而言，\(D_m\) 给出的并非单一相位花纹，而是一族随 Fibonacci 尺度递归加密的半转根环；其离散采样模值正是碰撞谱 \(a_m(k)\)，而黄金比的 Diophantine 最优逼近则进一步解释了这些干涉曲线在奇异环附近所呈现的层级接近律。

\paragraph{Gödel--Zeckendorf 地址极限的黄金自仿射闭合、自然尺度维数与查询刚性}
在地址前缀等距、Zeckendorf 有限截面递推与归一化复杂度界之间，还可抽出一组更紧的终端闭合：其一把“禁止相邻 \(1\)”直接压缩为一维 \(2^L\)-进地址线上的严格二分自仿射方程；其二把该极限集的自然尺度覆盖数、Hausdorff 维数与 Minkowski 内容写成 Fibonacci 常数的显式闭式；其三把折叠与稳定运算的可计算性提升为可精确审计的位查询刚性，并进一步把稳定算术层锁定为与自然数半环 PTIME 双向可解释的同一表示；最终它再与真实输入 \(40\)-状态核的一阶 Edgeworth 指纹汇合为同一条可复核链。

\begin{theorem}[Gödel--Zeckendorf 地址极限的严格二分自仿射方程]\label{thm:xi-zg-address-binary-self-affine}
取整数 \(L\ge 1\)，令
\[
B:=2^L,
\qquad
v\in T_2(E_{\min})\setminus 2T_2(E_{\min}),
\qquad
L_v:=\ZZ_2v.
\]
定义 admissible 序列空间
\[
\Sigma_{\mathrm{ZG}}
:=
\bigl\{
(z_t)_{t\ge 1}\in\{0,1\}^{\NN}:\ z_tz_{t+1}=0\ \ (t\ge 1)
\bigr\},
\]
以及地址极限集
\[
K_{B,v}
:=
\left\{
\sum_{t\ge 1}z_tB^{t-1}v:\ (z_t)_{t\ge 1}\in\Sigma_{\mathrm{ZG}}
\right\}
\subseteq L_v.
\]
则上述级数在 \(L_v\) 中收敛，\(\Sigma_{\mathrm{ZG}}\to K_{B,v}\) 的地址映射为单射，并且成立严格不交并分解
\[
\boxed{\;
K_{B,v}
=
BK_{B,v}\ \sqcup\ \bigl(v+B^2K_{B,v}\bigr).\;}
\]
\end{theorem}
\begin{proof}
由于 \(B=2^L\) 在 \(2\)-进绝对值下满足 \(|B|_2<1\)，故对任意 \((z_t)\in\Sigma_{\mathrm{ZG}}\)，级数
\[
\sum_{t\ge 1}z_tB^{t-1}v
\]
在完备群 \(L_v\) 中收敛。

再把 \(\mathcal E=\{0,1\}\)、\(\mathrm{code}(0)=0\)、\(\mathrm{code}(1)=1\) 代入定理
\ref{thm:xi-time-part9g-holographic-prefix-isometry-on-line}，即可知不同 admissible 无限序列给出不同地址，因而地址映射在 \(\Sigma_{\mathrm{ZG}}\) 上为单射。

现证自仿射分解。任取
\[
x=\sum_{t\ge 1}z_tB^{t-1}v\in K_{B,v},
\qquad
(z_t)\in\Sigma_{\mathrm{ZG}}.
\]
若 \(z_1=0\)，则尾序列 \((z_{t+1})_{t\ge 1}\) 仍属 \(\Sigma_{\mathrm{ZG}}\)，并且
\[
x
=
\sum_{t\ge 2}z_tB^{t-1}v
=
B\sum_{u\ge 1}z_{u+1}B^{u-1}v
\in BK_{B,v}.
\]
若 \(z_1=1\)，则 admissibility 强制 \(z_2=0\)，故 \((z_{t+2})_{t\ge 1}\in\Sigma_{\mathrm{ZG}}\)，并且
\[
x
=
v+\sum_{t\ge 3}z_tB^{t-1}v
=
v+B^2\sum_{u\ge 1}z_{u+2}B^{u-1}v
\in v+B^2K_{B,v}.
\]
因此
\[
K_{B,v}\subseteq BK_{B,v}\cup\bigl(v+B^2K_{B,v}\bigr).
\]

反向包含直接成立：若 \(y=\sum_{t\ge 1}z_tB^{t-1}v\in K_{B,v}\)，则
\[
By=\sum_{t\ge 2}z_{t-1}B^{t-1}v\in K_{B,v},
\]
而
\[
v+B^2y
=
v+\sum_{t\ge 3}z_{t-2}B^{t-1}v
\in K_{B,v},
\]
因为两条扩展序列仍满足禁止相邻 \(1\) 约束。

最后证不交。若
\[
Bx=v+B^2y
\qquad
(x,y\in K_{B,v}),
\]
则左端对应一条首位为 \(0\) 的 admissible 序列，右端对应一条首两位为 \(10\) 的 admissible 序列。由定理
\ref{thm:xi-time-part9g-holographic-prefix-isometry-on-line} 的地址单射性，二者不可能相等，故并为不交并。证毕。
\end{proof}

\begin{theorem}[Gödel--Zeckendorf 地址极限的精确 Minkowski 维数与自然尺度内容]\label{thm:xi-zg-address-minkowski-dimension-content}\leanverified{paper\_xi\_zg\_address\_minkowski\_dimension\_content}
沿用定理 \ref{thm:xi-zg-address-binary-self-affine} 的记号，并在 \(L_v\) 上取定理 \ref{thm:xi-time-part9g-holographic-prefix-isometry-on-line} 的 \(B\)-进超度量。对 \(n\ge 0\)，令 \(N_{B,v}(n)\) 表示覆盖 \(K_{B,v}\) 所需的最少 \(B^{-n}\)-球个数；等价地，\(N_{B,v}(n)\) 为商 \(L_v/B^nL_v\) 中与 \(K_{B,v}\) 相交的余类数。则
\[
\boxed{\;
N_{B,v}(n)=F_{n+2}\qquad (n\ge 0).\;}
\]
从而
\[
\boxed{\;
\dim_{\mathrm H}(K_{B,v})
=
\dim_{\mathrm M}(K_{B,v})
=
\frac{\log\varphi}{\log B}
=
\frac{\log\varphi}{L\log 2}.\;}
\]
并且沿自然 \(B\)-进尺度存在精确内容极限
\[
\boxed{\;
\lim_{n\to\infty}
N_{B,v}(n)\,B^{-n\dim_{\mathrm M}(K_{B,v})}
=
\frac{\varphi^2}{\sqrt5}.\;}
\]
此外，\(K_{B,v}\) 在 \(L_v\) 的 Haar 概率测度下为零测集。
\end{theorem}
\begin{proof}
对每个 \(n\ge 0\)，记
\[
\mathcal A_n
:=
\bigl\{
(z_1,\dots,z_n)\in\{0,1\}^n:\ z_tz_{t+1}=0\ \ (1\le t<n)
\bigr\},
\]
并约定 \(\mathcal A_0=\{\varnothing\}\)。对 \(a=(z_1,\dots,z_n)\in\mathcal A_n\)，定义对应 cylinder
\[
[a]
:=
\bigl\{
(w_t)_{t\ge 1}\in\Sigma_{\mathrm{ZG}}:\ w_1=z_1,\dots,w_n=z_n
\bigr\},
\]
以及前缀地址
\[
X_n(a):=\sum_{t=1}^{n}z_tB^{t-1}v.
\]
由定理 \ref{thm:xi-time-part9g-holographic-prefix-isometry-on-line}，
\[
X([a])=\bigl(X_n(a)+B^nL_v\bigr)\cap K_{B,v},
\]
且不同 \(a\in\mathcal A_n\) 给出不同的 \(B^nL_v\)-余类。于是 \(K_{B,v}\) 在尺度 \(B^{-n}\) 上恰分解为 \(|\mathcal A_n|\) 个互不相交 cylinder，故
\[
N_{B,v}(n)=|\mathcal A_n|.
\]

另一方面，定理 \ref{thm:xi-time-part9g-zg-primorial-fibonacci-affine-recursion} 已证明
\[
|\mathcal A_n|=F_{n+2},
\]
故第一条公式成立。

记
\[
D:=\frac{\log\varphi}{\log B}.
\]
则由 Binet 公式
\[
F_{n+2}\sim \frac{\varphi^{n+2}}{\sqrt5}
\]
得到
\[
\frac{\log N_{B,v}(n)}{n\log B}
=
\frac{\log F_{n+2}}{n\log B}
\longrightarrow
\frac{\log\varphi}{\log B}
=
D.
\]
因此 Minkowski 维数为 \(D\)。

再证 Hausdorff 维数。若 \(s>D\)，则第 \(n\) 层 cylinder 覆盖给出
\[
\mathcal H^s_{B^{-n}}(K_{B,v})
\le
N_{B,v}(n)\,B^{-ns}
=
F_{n+2}B^{-ns}
\to 0,
\]
故 \(\mathcal H^s(K_{B,v})=0\)。

反之，若 \(s<D\)，则任一直径不超过 \(B^{-n}\) 的集合都包含于某个 \(B^{-n}\)-球中，而任一此类球至多与一个第 \(n\) 层 cylinder 相交；否则它将同时命中两个不同的 \(B^nL_v\)-余类，与定理
\ref{thm:xi-time-part9g-holographic-prefix-isometry-on-line} 给出的严格余类分离矛盾。故任意直径不超过 \(B^{-n}\) 的覆盖都至少需要 \(N_{B,v}(n)\) 个集合，从而
\[
\mathcal H^s_{B^{-n}}(K_{B,v})
\ge
N_{B,v}(n)\,B^{-ns}
=
F_{n+2}B^{-ns}
\to \infty.
\]
因此 \(\mathcal H^s(K_{B,v})=\infty\)。综上 \(\dim_{\mathrm H}(K_{B,v})=D=\dim_{\mathrm M}(K_{B,v})\)。

再由 \(B^D=\varphi\)，有
\[
N_{B,v}(n)\,B^{-nD}
=
F_{n+2}\varphi^{-n}
\longrightarrow
\frac{\varphi^2}{\sqrt5},
\]
即得自然尺度内容极限。

最后，记 \(\lambda_v\) 为 \(L_v\) 的 Haar 概率测度。每个 \(B^nL_v\)-余类的测度均为 \(B^{-n}\)，而 \(K_{B,v}\) 只与其中 \(N_{B,v}(n)=F_{n+2}\) 个余类相交，故
\[
\lambda_v(K_{B,v})
\le
F_{n+2}B^{-n}
\sim
\frac{\varphi^2}{\sqrt5}\Bigl(\frac{\varphi}{B}\Bigr)^n.
\]
由于 \(B=2^L\ge 2>\varphi\)，令 \(n\to\infty\) 即得 \(\lambda_v(K_{B,v})=0\)。证毕。
\end{proof}

\begin{theorem}[折叠算子的确定性位查询复杂度恰为分辨率]\label{thm:xi-fold-query-complexity-exact-m}\leanverified{paper\_xi\_fold\_query\_complexity\_exact\_m}
对 \(m\ge 1\)，定义 \(Q_{\det}(\Fold_m)\) 为：在最坏情形下，为精确输出 \(\Fold_m(\omega)\)，确定性算法必须读取的输入位数的最小值。则
\[
\boxed{\;
Q_{\det}(\Fold_m)=m.\;}
\]
\end{theorem}
\begin{proof}
上界由命题 \ref{prop:fold-complexity} 直接给出：顺序读取全部 \(m\) 位并执行 Zeckendorf 规范化即可，故
\[
Q_{\det}(\Fold_m)\le m.
\]

现证下界。任取一个精确计算 \(\Fold_m\) 的确定性算法 \(\mathcal A\)，并考察其在输入
\[
\mathbf0:=0^m
\]
上的运行。若 \(\mathcal A\) 在该输入上未查询某一位置 \(j\in\{1,\dots,m\}\)，则令
\[
e_j:=00\cdots010\cdots0
\]
为仅第 \(j\) 位为 \(1\) 的字。由于 \(\mathbf0\) 与 \(e_j\) 在 \(\mathcal A\) 已查询的全部坐标上取值相同，算法在两输入上将经历完全相同的查询转录并输出同一结果。

然而 \(e_j\in X_m\)，由命题 \ref{prop:fold-basic} 的幂等性，
\[
\Fold_m(\mathbf0)=\mathbf0,
\qquad
\Fold_m(e_j)=e_j,
\]
二者显然不同，矛盾。故 \(\mathcal A\) 在输入 \(\mathbf0\) 上必须读取全部 \(m\) 位，从而
\[
Q_{\det}(\Fold_m)\ge m.
\]
合并上下界即得结论。证毕。
\end{proof}

\begin{corollary}[稳定加法的全输入查询饱和与稳定乘法单位切片的满查询刚性]\label{cor:xi-stable-add-query-saturation-mul-unit-slice}
\leanverified{paper\_xi\_stable\_add\_query\_saturation\_mul\_unit\_slice}
沿用定义 \ref{def:stable-add} 与 \ref{def:stable-mul} 的记号。对支撑包含于 \(\{1,\dots,m\}\) 的合法输入，记 \(Q_{\det}^{\oplus}(m)\) 为稳定加法 \(c\oplus d\) 的最坏情形总位查询复杂度。再令
\[
\mathbf1:=10\cdots0\in X_\infty,
\qquad
\mathrm{Val}(\mathbf1)=1,
\]
并定义乘法的单位切片
\[
L_m^{\otimes}(c):=c\otimes \mathbf1,
\qquad
R_m^{\otimes}(d):=\mathbf1\otimes d.
\]
则有
\[
\boxed{\;
Q_{\det}^{\oplus}(m)=2m,\qquad
Q_{\det}(L_m^{\otimes})=Q_{\det}(R_m^{\otimes})=m.\;}
\]
\end{corollary}
\begin{proof}
先看稳定加法。上界由命题 \ref{prop:stable-arith-complexity} 直接给出：读取两输入全部 \(2m\) 位后执行值层加法与 Zeckendorf 规范化即可，故
\[
Q_{\det}^{\oplus}(m)\le 2m.
\]

现证下界。任取一个精确计算 \(c\oplus d\) 的确定性算法 \(\mathcal A\)，并考察其在输入
\[
(\mathbf0,\mathbf0)
\]
上的运行，其中 \(\mathbf0:=0^\infty\)，并令 \(e_j\) 表示第 \(j\) 位为 \(1\)、其余位为 \(0\) 的合法无限字。若 \(\mathcal A\) 在该运行中未查询第一输入的某一位置 \(j\le m\)，则与输入
\[
(e_j,\mathbf0)
\]
相比，所有已查询坐标取值仍完全相同，故 \(\mathcal A\) 将给出同一输出；但
\[
e_j\oplus \mathbf0=e_j,
\qquad
\mathbf0\oplus \mathbf0=\mathbf0,
\]
矛盾。故第一输入的前 \(m\) 位必须全部被查询。

若 \(\mathcal A\) 未查询第二输入的某一位置 \(j\le m\)，同理比较
\[
(\mathbf0,\mathbf0)
\qquad\text{与}\qquad
(\mathbf0,e_j),
\]
即可得到矛盾，因为
\[
\mathbf0\oplus e_j=e_j
\neq
\mathbf0.
\]
故在输入 \((\mathbf0,\mathbf0)\) 上，两侧前 \(m\) 位都必须全部读取，从而
\[
Q_{\det}^{\oplus}(m)\ge 2m.
\]
于是 \(Q_{\det}^{\oplus}(m)=2m\)。

再看乘法单位切片。由 \(\mathrm{Val}(\mathbf1)=1\) 与定义 \ref{def:stable-mul}，
\[
L_m^{\otimes}(c)=c\otimes\mathbf1=c,
\qquad
R_m^{\otimes}(d)=\mathbf1\otimes d=d.
\]
因此两条切片在约束输入类上就是恒等映射。其查询下界可用与定理 \ref{thm:xi-fold-query-complexity-exact-m} 同样的零字--单点字对抗得到：若在输入 \(\mathbf0\) 上遗漏某个位置 \(j\le m\)，则无法区分 \(\mathbf0\) 与 \(e_j\)。故
\[
Q_{\det}(L_m^{\otimes})\ge m,
\qquad
Q_{\det}(R_m^{\otimes})\ge m.
\]
显然读取前 \(m\) 位即可输出该合法字本身，故反向不等式也成立，从而
\[
Q_{\det}(L_m^{\otimes})=Q_{\det}(R_m^{\otimes})=m.
\]
证毕。
\end{proof}

\begin{theorem}[稳定算术层与自然数半环在 PTIME 内双向可解释]\label{thm:xi-stable-semiring-ptime-biinterpretability}
\leanverified{paper\_xi\_stable\_semiring\_ptime\_biinterpretability}
沿用定义 \ref{def:stable-add}、\ref{def:stable-mul} 与命题
\ref{prop:stable-arith-complexity} 的记号。则值映射
\[
\mathrm{Val}:X_\infty\longrightarrow \NN
\]
及其逆映射 \(Z:=\mathrm{Val}^{-1}\) 给出交换半环同构
\[
\boxed{\;
\mathrm{Val}(c\oplus d)=\mathrm{Val}(c)+\mathrm{Val}(d),\qquad
\mathrm{Val}(c\otimes d)=\mathrm{Val}(c)\,\mathrm{Val}(d).\;}
\]
更进一步，在有效分辨率 \(m\) 的输入口径下，两侧运算可在多项式时间内双向模拟：
\begin{enumerate}
  \item 从 \(c,d\in X_\infty\) 计算 \(c\oplus d\) 与 \(c\otimes d\)，在比特口径下分别需要
  \[
  O(m^2),\qquad O(M(m)+m^2)
  \]
  位运算，其中 \(M(m)\) 为 \(m\)-比特整数乘法成本；
  \item 从 \(m\)-比特整数 \(N\) 计算 Zeckendorf 规范形 \(Z(N)\)，以及从 \(c\in X_\infty\) 计算 \(\mathrm{Val}(c)\)，都可在 \(O(m^2)\) 位运算内完成。
\end{enumerate}
因此
\[
\boxed{\;
(X_\infty,\oplus,\otimes)\ \text{与}\ (\NN,+,\times)
\ \text{在 semiring 意义下同构，并在 PTIME 内双向可解释。}\;}
\]
\end{theorem}
\begin{proof}
交换半环同构本身是定义性的：由定理 \ref{thm:stable-add-value-consistency} 与定义 \ref{def:stable-mul}，
\[
\mathrm{Val}(c\oplus d)=\mathrm{Val}(c)+\mathrm{Val}(d),
\qquad
c\otimes d=\mathrm{Val}^{-1}\!\bigl(\mathrm{Val}(c)\mathrm{Val}(d)\bigr),
\]
而 \(\mathrm{Val}\) 由 Zeckendorf 唯一性给出双射，因此半环同构正是定理 \ref{thm:mul-definitional} 的内容。

算法复杂度则给出真正的计算内容。对第一条，命题 \ref{prop:stable-arith-complexity} 已经证明：在有效分辨率 \(m\) 的输入上，稳定加法与稳定乘法可按“值层运算 \(+\) Zeckendorf 规范化”实现，并分别满足
\[
O(m^2),\qquad O(M(m)+m^2)
\]
的比特复杂度。

对第二条，若给定 \(m\)-比特整数 \(N\)，则其 Zeckendorf 规范形 \(Z(N)\) 可由附录复杂度界中的贪心规范化实现；命题 \ref{prop:fold-complexity} 对 Fibonacci 权重的分析表明，这一步需要 \(O(m^2)\) 位运算。反过来，若给定支撑包含于 \(\{1,\dots,m\}\) 的 \(c\in X_\infty\)，则
\[
\mathrm{Val}(c)=\sum_{k=1}^{m}c_kF_{k+1}
\]
只是 \(O(m)\) 项、每项位长 \(O(m)\) 的加权求和，因此同样可在 \(O(m^2)\) 位运算内完成。

于是，自然数加乘可先经 \(Z\) 送入稳定地址层，再用 \(\oplus,\otimes\) 模拟并用 \(\mathrm{Val}\) 读回；稳定地址层运算也可先经 \(\mathrm{Val}\) 送回整数值层完成。两侧都具有多项式时间实现，故得到所述 PTIME 双向可解释性。证毕。
\end{proof}

\paragraph{真实输入 \texorpdfstring{$40$}{40}-状态核的一阶 Edgeworth 指纹闭合}
与上述地址几何与查询刚性并行，定理 \ref{thm:xi-real-input-40-output-edgeworth-nongaussian} 已给出真实输入 \(40\)-状态核在输出位势方向上的显式符号指纹
\[
\gamma_1\approx -1.476481644016533,
\qquad
\gamma_2\approx 2.3928814165461985.
\]
因此输出计数的一阶 Edgeworth 修正严格左偏且严格厚尾。于是，Gödel--Zeckendorf 地址极限的黄金维数律、折叠与稳定接口的查询刚性，以及真实输入核的非高斯统计指纹，便在同一审计口径下封闭为一条统一的可复核链。

\paragraph{末位充分化、均匀基线判别与零点几何对碰撞缺口的指数失明}
由末位两层塌缩、局部信息密度的一阶离散修正、均匀基线下的两原子似然比极限以及谱零点的半维稀疏律，可进一步抽出下列八条彼此衔接的终端后果：它们分别给出末位统计量的一样本渐近充分性、亚临界样本规模下完整输出实验与末位 Bernoulli 实验的 Le Cam 等价、隐藏对数多重度涨落的二点极限律、其 Laplace--累积闭式、相对稳定层均匀基线的一样本最优判别常数、全部经典散度的 Bernoulli 塌缩、零点数量与碰撞缺口主尺度之间的严格分离，以及零点相对密度在近零邻域无法控制归一化碰撞尺度。

\begin{theorem}[末位统计量的渐近充分性]\label{thm:xi-fold-lastbit-asymptotic-sufficiency}
\leanverified{paper\_xi\_fold\_lastbit\_asymptotic\_sufficiency}
设
\[
X_m:=\{w\in\{0,1\}^m:\ w_iw_{i+1}=0\ \ (1\le i<m)\},
\qquad
\mu_m(w):=\frac{d_m^{\mathrm{bin}}(w)}{2^m},
\qquad
\ell_m(w):=w_m\in\{0,1\},
\]
并将稳定层分成两块
\[
A_i:=\{w\in X_m:\ \ell_m(w)=i\}
\qquad(i=0,1).
\]
定义 Markov 核
\[
K_m(i,\cdot):=\mathrm{Unif}(A_i),
\]
以及由末位统计量重构的块均匀近似
\[
\bar\mu_m:=K_m\circ (\ell_m)_\#\mu_m.
\]
则有
\[
\boxed{\;
D_{\mathrm{TV}}(\mu_m,\bar\mu_m)
=
O\!\Bigl(\Bigl(\frac{\varphi}{2}\Bigr)^m\Bigr).\;}
\]
因此 \(\ell_m\) 在 Le Cam 意义下对 \(\mu_m\) 是渐近充分统计量。

更具体地，对任意有界函数 \(f_m:X_m\to\RR\)，若满足
\[
\sum_{w\in A_i}f_m(w)=0
\qquad(i=0,1),
\]
则
\[
\boxed{\;
\bigl|\mathbb E_{\mu_m}[f_m]\bigr|
\le
2\|f_m\|_\infty\,
O\!\Bigl(\Bigl(\frac{\varphi}{2}\Bigr)^m\Bigr).\;}
\]
\leanverified{paper\_xi\_fold\_lastbit\_asymptotic\_sufficiency}
\end{theorem}
\begin{proof}
由命题 \ref{prop:fold-bin-lastbit-sufficient-tv}，存在绝对常数 \(C>0\) 使
\[
D_{\mathrm{TV}}(\mu_m,\bar\mu_m)\le C\Bigl(\frac{\varphi}{2}\Bigr)^m.
\]
这正是第一式。又因为 \(\bar\mu_m=K_m\circ(\ell_m)_\#\mu_m\)，已构造出一条把末位统计量提升回稳定层的重构核，其误差在全变差意义下指数消失，故 \(\ell_m\) 对 \(\mu_m\) 在 Le Cam 意义下渐近充分。

对第二条，只需注意 \(\bar\mu_m\) 在每个 \(A_i\) 上为均匀分布，故
\[
\mathbb E_{\bar\mu_m}[f_m]
=
\sum_{i=0}^{1}\bar\mu_m(A_i)\,\frac{1}{|A_i|}\sum_{w\in A_i}f_m(w)=0.
\]
于是
\[
\bigl|\mathbb E_{\mu_m}[f_m]\bigr|
=
\bigl|\mathbb E_{\mu_m}[f_m]-\mathbb E_{\bar\mu_m}[f_m]\bigr|
\le
2\|f_m\|_\infty D_{\mathrm{TV}}(\mu_m,\bar\mu_m).
\]
代入第一式即得所示估计。证毕。
\end{proof}

\begin{theorem}[亚临界样本规模下完整 bin-fold 输出实验与末位 Bernoulli 实验的 Le Cam 等价]\label{thm:xi-fold-lastbit-subcritical-sample-lecam-equivalence}
\leanverified{paper\_xi\_fold\_lastbit\_subcritical\_sample\_lecam\_equivalence}
沿用定理 \ref{thm:xi-fold-lastbit-asymptotic-sufficiency} 的记号，并记
\[
p_m:=\mu_m(A_1)=\bar\mu_m(A_1).
\]
对任意 \(n\ge 1\)，定义单点统计实验
\[
\mathcal E_m^{(n)}:=\{\mu_m^{\otimes n}\}
\qquad\text{定义在 }X_m^n\text{ 上},
\]
\[
\mathcal B_m^{(n)}:=\{\mathrm{Bernoulli}(p_m)^{\otimes n}\}
\qquad\text{定义在 }\{0,1\}^n\text{ 上}.
\]
并记
\[
\Delta(\mathcal E,\mathcal F):=\max\{\delta(\mathcal E,\mathcal F),\delta(\mathcal F,\mathcal E)\}
\]
为 Le Cam 距离。再定义确定性读出核
\[
L_m^{(n)}(w_1,\dots,w_n):=\bigl(\ell_m(w_1),\dots,\ell_m(w_n)\bigr),
\]
以及重构核
\[
K_m^{(n)}(b_1,\dots,b_n,\cdot):=
K_m(b_1,\cdot)\otimes\cdots\otimes K_m(b_n,\cdot).
\]
则有精确恒等式
\[
L_{m\#}^{(n)}\mu_m^{\otimes n}
=
\mathrm{Bernoulli}(p_m)^{\otimes n},
\qquad
K_{m\#}^{(n)}\mathrm{Bernoulli}(p_m)^{\otimes n}
=
\bar\mu_m^{\otimes n},
\]
并且存在绝对常数 \(C>0\) 使
\[
\boxed{\;
D_{\mathrm{TV}}\!\Bigl(
\mu_m^{\otimes n},
K_{m\#}^{(n)}\mathrm{Bernoulli}(p_m)^{\otimes n}
\Bigr)
\le
nC\Bigl(\frac{\varphi}{2}\Bigr)^m.\;}
\]
因此 Le Cam 距离满足
\[
\boxed{\;
\Delta\!\bigl(\mathcal E_m^{(n)},\mathcal B_m^{(n)}\bigr)
\le
nC\Bigl(\frac{\varphi}{2}\Bigr)^m.\;}
\]
特别地，若 \(n_m=o((2/\varphi)^m)\)，则
\[
\boxed{\;
\Delta\!\bigl(\mathcal E_m^{(n_m)},\mathcal B_m^{(n_m)}\bigr)\longrightarrow 0.\;}
\]
\end{theorem}
\begin{proof}
由定义，\(L_m^{(n)}\) 只是逐坐标读取末位，因此
\[
L_{m\#}^{(n)}\mu_m^{\otimes n}
=
\bigl((\ell_m)_\#\mu_m\bigr)^{\otimes n}
=
\mathrm{Bernoulli}(p_m)^{\otimes n}.
\]
另一方面，\(K_m(i,\cdot)=\mathrm{Unif}(A_i)\) 且 \(\bar\mu_m=K_m\circ(\ell_m)_\#\mu_m\)，故
\[
K_{m\#}^{(n)}\mathrm{Bernoulli}(p_m)^{\otimes n}
=
\bar\mu_m^{\otimes n}.
\]

再由乘积测度的望远镜估计，对任意概率分布 \(\mu,\nu\) 都有
\[
D_{\mathrm{TV}}(\mu^{\otimes n},\nu^{\otimes n})
\le
n\,D_{\mathrm{TV}}(\mu,\nu).
\]
将 \((\mu,\nu)=(\mu_m,\bar\mu_m)\) 代入，并使用定理 \ref{thm:xi-fold-lastbit-asymptotic-sufficiency} 中
\[
D_{\mathrm{TV}}(\mu_m,\bar\mu_m)\le C\Bigl(\frac{\varphi}{2}\Bigr)^m,
\]
即得所示全变差界。

按 deficiency 的定义，从 \(\mathcal E_m^{(n)}\) 到 \(\mathcal B_m^{(n)}\) 的缺陷为 \(0\)，因为 \(L_m^{(n)}\) 已实现精确推送；从 \(\mathcal B_m^{(n)}\) 到 \(\mathcal E_m^{(n)}\) 的缺陷至多为上述全变差误差。再由
定理中的 Le Cam 距离定义，便得到所示距离界。若 \(n_m=o((2/\varphi)^m)\)，则右端趋于 \(0\)，故两实验渐近等价。证毕。
\end{proof}

\begin{corollary}[隐藏对数多重度涨落收敛为末位支配的二点律]\label{cor:xi-fold-hidden-logmultiplicity-branchbit-two-point-law}
\leanverified{paper\_xi\_fold\_hidden\_logmultiplicity\_branchbit\_two\_point\_law}
沿用 \(\mu_m\) 的记号，并令 \(W_m\sim\mu_m\)。定义
\[
Y_m:=\log d_m^{\mathrm{bin}}(W_m)-m\log\!\Bigl(\frac{2}{\varphi}\Bigr).
\]
则有分布收敛
\[
\boxed{\;
Y_m\Longrightarrow
\frac{\varphi}{\sqrt5}\,\delta_0
+\frac{1}{\varphi\sqrt5}\,\delta_{-\log\varphi}.\;}
\]
等价地，若 \(Y_\infty\) 为满足
\[
\mathbb P(Y_\infty=0)=\frac{\varphi}{\sqrt5},
\qquad
\mathbb P(Y_\infty=-\log\varphi)=\frac{1}{\varphi\sqrt5}
\]
的二点随机变量，则 \(Y_m\Rightarrow Y_\infty\)。特别地，
\[
\boxed{\;
\lim_{m\to\infty}\operatorname{Var}(Y_m)
=
\frac{1}{\varphi\sqrt5}
\left(1-\frac{1}{\varphi\sqrt5}\right)(\log\varphi)^2.\;}
\]
因此，隐藏对数多重度的首个非平凡涨落并不趋于非退化 Gaussian 律，而直接塌缩为由末位比特控制的 Bernoulli 二点极限。
\end{corollary}
\begin{proof}
由推论 \ref{cor:xi-fold-local-information-density-first-order-discrete-law}，若记
\[
\gamma_m:=\frac{1}{m}\log d_m^{\mathrm{bin}}(W_m),
\]
则有
\[
m\left(\gamma_m-\log\!\frac{2}{\varphi}\right)\Longrightarrow
\frac{\varphi}{\sqrt5}\,\delta_0
+\frac{1}{\varphi\sqrt5}\,\delta_{-\log\varphi}.
\]
而左端恰为 \(Y_m\)。方差极限亦即该推论中的 \(m^2\operatorname{Var}(\gamma_m)\) 极限，因为
\[
\operatorname{Var}(Y_m)=m^2\operatorname{Var}(\gamma_m).
\]
证毕。
\end{proof}

\begin{theorem}[隐藏对数多重度涨落的二原子 Laplace--累积母函数闭式]\label{thm:xi-fold-hidden-logmultiplicity-laplace-cumulant-closed}
沿用推论 \ref{cor:xi-fold-hidden-logmultiplicity-branchbit-two-point-law} 的记号。对 \(t\in\RR\) 定义
\[
\Lambda_m(t):=\log \EE_{\mu_m}\!\bigl[e^{tY_m}\bigr].
\]
再记
\[
\theta_\star:=\frac{1}{1+\varphi^2}=\frac{1}{\varphi\sqrt5}.
\]
则对任意紧区间 \(I\Subset\RR\)，都有一致渐近公式
\[
\boxed{\;
\sup_{t\in I}
\left|
\Lambda_m(t)
-\log\!\Bigl(\frac{\varphi+\varphi^{-t-1}}{\sqrt5}\Bigr)
\right|
=
O_I\!\Bigl(\Bigl(\frac{\varphi}{2}\Bigr)^m\Bigr).\;}
\]
特别地，
\[
\boxed{\;
\Lambda(t):=\lim_{m\to\infty}\Lambda_m(t)
=
\log\!\Bigl(\frac{\varphi+\varphi^{-t-1}}{\sqrt5}\Bigr).\;}
\]
因此 \(Y_m\) 的一切累积量都收敛到对应二原子极限律 \(Y_\infty\) 的累积量；若记
\[
\kappa_r^\infty:=\lim_{m\to\infty}\kappa_r(Y_m)=\Lambda^{(r)}(0),
\]
则前四阶为
\[
\boxed{\;
\kappa_1^\infty=-\frac{\log\varphi}{\varphi\sqrt5},\qquad
\kappa_2^\infty=\frac{(\log\varphi)^2}{5},\qquad
\kappa_3^\infty=-\frac{(\log\varphi)^3}{5\sqrt5},\qquad
\kappa_4^\infty=-\frac{(\log\varphi)^4}{25}.\;}
\]
换言之，隐藏对数多重度涨落并不落入任何非退化 Gaussian 中心极限定域，而是严格保留一个由末位比特控制的二原子极限核心。
\end{theorem}
\begin{proof}
记
\[
a:=\log\varphi,
\qquad
B_m:=W_m(m)\in\{0,1\}.
\]
由定理 \ref{thm:fold-bin-two-state-asymptotic} 的一致对数展开，
\[
Y_m=-a\,B_m+R_m,
\qquad
\sup_{w\in X_m}|R_m(w)|
=
O\!\Bigl(\Bigl(\frac{\varphi}{2}\Bigr)^m\Bigr).
\]
因此对任意紧区间 \(I\Subset\RR\) 与任意 \(t\in I\)，都有一致展开
\[
e^{tY_m}
=
e^{-atB_m}e^{tR_m}
=
e^{-atB_m}
\left(
1+O_I\!\Bigl(\Bigl(\frac{\varphi}{2}\Bigr)^m\Bigr)
\right).
\]
对 \(\mu_m\) 取期望，并利用命题 \ref{prop:fold-bin-escort-lastbit} 在 \(q=1\) 处的末位质量公式
\[
\mu_m(B_m=1)=\theta_\star+O\!\Bigl(\Bigl(\frac{\varphi}{2}\Bigr)^m\Bigr),
\qquad
\mu_m(B_m=0)=1-\theta_\star+O\!\Bigl(\Bigl(\frac{\varphi}{2}\Bigr)^m\Bigr),
\]
得到
\[
\EE_{\mu_m}[e^{tY_m}]
=
(1-\theta_\star)+\theta_\star e^{-at}
+O_I\!\Bigl(\Bigl(\frac{\varphi}{2}\Bigr)^m\Bigr).
\]
而
\[
(1-\theta_\star)+\theta_\star e^{-at}
=
\frac{\varphi}{\sqrt5}+\frac{1}{\varphi\sqrt5}\varphi^{-t}
=
\frac{\varphi+\varphi^{-t-1}}{\sqrt5}.
\]
由于右端在每个紧区间 \(I\) 上严格正且有正下界，取对数即得
\[
\Lambda_m(t)
=
\log\!\Bigl(\frac{\varphi+\varphi^{-t-1}}{\sqrt5}\Bigr)
+O_I\!\Bigl(\Bigl(\frac{\varphi}{2}\Bigr)^m\Bigr),
\]
从而得到 \(\Lambda\) 的闭式。

另一方面，由上式可知 \(Y_m\) 一致有界，故分布收敛
\[
Y_m\Rightarrow Y_\infty
\]
自动推出各阶矩收敛；累积量作为矩的通用多项式函数，也随之逐阶收敛到 \(Y_\infty\) 的累积量。由于
\[
Y_\infty=-a\,B_\star,
\qquad
\PP(B_\star=1)=\theta_\star,
\qquad
\PP(B_\star=0)=1-\theta_\star,
\]
Bernoulli 两点律的标准累积量公式给出
\[
\kappa_1^\infty=-a\theta_\star,\qquad
\kappa_2^\infty=a^2\theta_\star(1-\theta_\star),
\]
\[
\kappa_3^\infty=-a^3\theta_\star(1-\theta_\star)(1-2\theta_\star),
\]
\[
\kappa_4^\infty=a^4\theta_\star(1-\theta_\star)\bigl(1-6\theta_\star(1-\theta_\star)\bigr).
\]
再用恒等式
\[
\theta_\star(1-\theta_\star)=\frac15,
\qquad
1-2\theta_\star=\frac{1}{\sqrt5},
\]
即可化简得到所列四个闭式。证毕。
\end{proof}

\begin{theorem}[纤维信息与规范体积密度的全显式两尺度渐近]\label{thm:xi-fold-hidden-logmultiplicity-kappa-gauge-two-scale-explicit}
沿用定理 \ref{thm:fold-bin-two-state-asymptotic} 与命题 \ref{prop:fold-bin-escort-lastbit} 在 \(q=1\) 处的记号，并记
\[
\kappa_m
:=
\EE_{\mu_m}\!\bigl[\log d_m^{\mathrm{bin}}(W_m)\bigr].
\]
则有显式渐近公式
\[
\kappa_m
=
m\log\!\Bigl(\frac{2}{\varphi}\Bigr)
-\frac{\log\varphi}{1+\varphi^2}
+O\!\Bigl(\Bigl(\frac{\varphi}{2}\Bigr)^m\Bigr).
\]
若进一步记 \(g_m\) 为规范体积密度，则
\[
g_m
=
m\log\!\Bigl(\frac{2}{\varphi}\Bigr)
-\Bigl(1+\frac{\log\varphi}{1+\varphi^2}\Bigr)
+\frac{(\varphi/2)^m}{2\sqrt5}
\Bigl(
\varphi^2 m\log\!\Bigl(\frac{2}{\varphi}\Bigr)
+\varphi^2\log(2\pi)-\log\varphi
\Bigr)
+O\!\Bigl(\Bigl(\frac{\varphi}{2}\Bigr)^m\Bigr).
\]
\end{theorem}
\begin{proof}
由定理 \ref{thm:fold-bin-two-state-asymptotic} 的一致对数展开，
\[
\log d_m^{\mathrm{bin}}(w)
=
m\log\!\Bigl(\frac{2}{\varphi}\Bigr)
-w_m\log\varphi
+O\!\Bigl(\Bigl(\frac{\varphi}{2}\Bigr)^m\Bigr)
\]
在 \(X_m\) 上一致成立。对 \(\mu_m\) 取期望得
\[
\kappa_m
=
m\log\!\Bigl(\frac{2}{\varphi}\Bigr)
-\log\varphi\cdot \mu_m\!\bigl(W_m(m)=1\bigr)
+O\!\Bigl(\Bigl(\frac{\varphi}{2}\Bigr)^m\Bigr).
\]
再由命题 \ref{prop:fold-bin-escort-lastbit} 在 \(q=1\) 处的末位极限，
\[
\mu_m\!\bigl(W_m(m)=1\bigr)
=
\frac{1}{1+\varphi^2}
+O\!\Bigl(\Bigl(\frac{\varphi}{2}\Bigr)^m\Bigr),
\]
代回即得 \(\kappa_m\) 的闭式渐近。

另一方面，由命题 \ref{prop:fold-bin-gauge-density-two-scale}，
\[
g_m
=
\kappa_m-1
+\frac{(\varphi/2)^m}{2\sqrt5}
\Bigl(
\varphi^2 m\log\!\Bigl(\frac{2}{\varphi}\Bigr)
+\varphi^2\log(2\pi)-\log\varphi
\Bigr)
+O\!\Bigl(\Bigl(\frac{\varphi}{2}\Bigr)^m\Bigr).
\]
将上式中的 \(\kappa_m\) 代入，便得到所示 \(g_m\) 的全显式两尺度展开。证毕。
\end{proof}


\begin{theorem}[与稳定层均匀分布的一样本最优判别常数]\label{thm:xi-fold-uniform-baseline-testing-constant}
\leanverified{paper\_xi\_fold\_uniform\_baseline\_testing\_constant}
沿用 \(\mu_m(w)=d_m^{\mathrm{bin}}(w)/2^m\) 的记号。令
\[
u_m(w):=\frac{1}{|X_m|}
\qquad (w\in X_m)
\]
为稳定层均匀分布，并记似然比
\[
R_m(w):=\frac{\mu_m(w)}{u_m(w)}.
\]
则在 \(u_m\) 下，\(R_m\) 依分布收敛到二点随机变量
\[
R=
\begin{cases}
\dfrac{\varphi^2}{\sqrt5},& \text{以概率 }\dfrac1\varphi,\\[8pt]
\dfrac{\varphi}{\sqrt5},& \text{以概率 }\dfrac1{\varphi^2}.
\end{cases}
\]
从而有精确极限
\[
\boxed{\;
\lim_{m\to\infty}D_{\mathrm{TV}}(\mu_m,u_m)
=
1-\frac{2}{\sqrt5}.\;}
\]
特别地，在等先验的一样本假设检验
\[
H_0:u_m,
\qquad
H_1:\mu_m
\]
中，最优 Bayes 错误概率满足
\[
\boxed{\;
\lim_{m\to\infty}P_e^\ast(\mu_m,u_m)
=
\frac1{\sqrt5}.\;}
\]
\end{theorem}
\begin{proof}
由定理 \ref{thm:conclusion-foldbin-likelihood-ratio-two-atom-transfer}，
\[
R_m(W_m)\Rightarrow R,
\qquad
W_m\sim u_m,
\]
并且对充分大的 \(m\)，随机变量 \(R_m(W_m)\) 落在某个固定闭区间 \(I\subset(0,\infty)\) 中。故
\[
D_{\mathrm{TV}}(\mu_m,u_m)
=
\frac12\,\mathbb E_{u_m}|R_m-1|
\longrightarrow
\frac12\,\mathbb E|R-1|.
\]
直接计算得到
\[
\frac12\,\mathbb E|R-1|
=
\frac12\left[
\frac1\varphi\Bigl(\frac{\varphi^2}{\sqrt5}-1\Bigr)
+
\frac1{\varphi^2}\Bigl(1-\frac{\varphi}{\sqrt5}\Bigr)
\right].
\]
这与定理 \ref{thm:xi-foldbin-pushforward-uniform-two-atom-ledger} 给出的全变差极限常数
\[
\frac12\left(\frac1{\sqrt5}-\varphi^{-3}\right)
\]
一致；再用恒等式 \(\varphi^{-3}=\sqrt5-2\)，即得
\[
\lim_{m\to\infty}D_{\mathrm{TV}}(\mu_m,u_m)=1-\frac{2}{\sqrt5}.
\]

最后，由 Helstrom 公式
\[
P_e^\ast(\mu_m,u_m)=\frac12\bigl(1-D_{\mathrm{TV}}(\mu_m,u_m)\bigr),
\]
令 \(m\to\infty\) 即得
\[
\lim_{m\to\infty}P_e^\ast(\mu_m,u_m)=\frac1{\sqrt5}.
\]
证毕。
\end{proof}

\begin{corollary}[相对于稳定层均匀基线的全部经典散度塌缩为单一 Bernoulli 几何]\label{cor:xi-fold-uniform-baseline-classical-fdiv-bernoulli-collapse}
\leanverified{paper\_xi\_fold\_uniform\_baseline\_classical\_fdiv\_bernoulli\_collapse}
沿用定理 \ref{thm:xi-fold-uniform-baseline-testing-constant} 的记号，并记
\[
q_*:=\varphi^{-2},
\qquad
p_*:=\frac{1}{\varphi\sqrt5}.
\]
则对任意凸函数 \(f:(0,\infty)\to\RR\) 满足 \(f(1)=0\)，有
\[
\boxed{\;
D_f(\mu_m\Vert u_m)
=
D_f\!\bigl(\mathrm{Bernoulli}(p_*)\Vert \mathrm{Bernoulli}(q_*)\bigr)
+O_f\!\Bigl(\Bigl(\frac{\varphi}{2}\Bigr)^m\Bigr).\;}
\]
特别地，
\[
\boxed{\;
\lim_{m\to\infty}D_{\mathrm{TV}}(\mu_m,u_m)
=
\left|p_*-q_*\right|
=
1-\frac{2}{\sqrt5},\;}
\]
\[
\boxed{\;
\lim_{m\to\infty}D_{\mathrm{KL}}(\mu_m\Vert u_m)
=
D_{\mathrm{KL}}\!\bigl(\mathrm{Bernoulli}(p_*)\Vert \mathrm{Bernoulli}(q_*)\bigr),\;}
\]
\[
\boxed{\;
\lim_{m\to\infty}\chi^2(\mu_m\Vert u_m)
=
\chi^2\!\bigl(\mathrm{Bernoulli}(p_*)\Vert \mathrm{Bernoulli}(q_*)\bigr).\;}
\]
换言之，\(\mu_m\) 相对于稳定层均匀基线 \(u_m\) 的全部经典 Csisz\'ar 几何，在极限中都严格塌缩为同一个二点 Bernoulli 对。
\end{corollary}
\begin{proof}
由定理 \ref{thm:xi-foldbin-pushforward-uniform-two-atom-ledger}，对任意满足条件的 \(f\) 都有
\[
D_f(\mu_m\Vert u_m)
=
\frac1\varphi\,f\!\left(\frac{\varphi^2}{\sqrt5}\right)
+\frac1{\varphi^2}\,f\!\left(\frac{\varphi}{\sqrt5}\right)
+O_f\!\Bigl(\Bigl(\frac{\varphi}{2}\Bigr)^m\Bigr).
\]
另一方面，
\[
1-q_*=\frac1\varphi,
\qquad
1-p_*=\frac{\varphi}{\sqrt5},
\]
故
\[
\frac{1-p_*}{1-q_*}
=
\frac{\varphi/\sqrt5}{1/\varphi}
=
\frac{\varphi^2}{\sqrt5},
\qquad
\frac{p_*}{q_*}
=
\frac{1/(\varphi\sqrt5)}{\varphi^{-2}}
=
\frac{\varphi}{\sqrt5}.
\]
这正说明上式等于
\[
D_f\!\bigl(\mathrm{Bernoulli}(p_*)\Vert \mathrm{Bernoulli}(q_*)\bigr)
+O_f\!\Bigl(\Bigl(\frac{\varphi}{2}\Bigr)^m\Bigr).
\]
对全变差、KL 与 \(\chi^2\) 只需分别取
\[
f_{\mathrm{TV}}(t)=\frac12|t-1|,
\qquad
f_{\mathrm{KL}}(t)=t\log t,
\qquad
f_{\chi^2}(t)=(t-1)^2.
\]
其中全变差的极限常数也可写作
\[
|p_*-q_*|
=
\varphi^{-2}-\frac{1}{\varphi\sqrt5}
=
1-\frac{2}{\sqrt5},
\]
与定理 \ref{thm:xi-fold-uniform-baseline-testing-constant} 一致。证毕。
\end{proof}

\begin{theorem}[零点几何对碰撞缺口的指数失明]\label{thm:xi-fold-zero-geometry-collision-gap-blindness}
\leanverified{paper\_xi\_fold\_zero\_geometry\_collision\_gap\_blindness}
沿用结论 \ref{con:xi-fold-collision-spectral-trace-parseval} 的记号，令
\[
M:=F_{m+2},
\qquad
\mathcal Z_m:=\{k\in \ZZ/(M\ZZ):\ \widehat d_m(k)=0\}.
\]
其中 \(C_\varphi>0\) 为定理 \ref{thm:fold-critical-resonance-constant} 的黄金共振常数。则碰撞缺口与零点密度满足严格分离：
\[
\boxed{\;
M\Bigl(\mathsf{Col}_m-\frac1M\Bigr)
\ge
2C_\varphi^2+o(1),
\qquad
\frac{|\mathcal Z_m|}{M}
\le
\tau(m+2)\,O(\varphi^{-m/2}).\;}
\]
从而
\[
\boxed{\;
\frac{\mathsf{Col}_m-\frac1M}{|\mathcal Z_m|/M^2}
=
\frac{M\bigl(\mathsf{Col}_m-\frac1M\bigr)}{|\mathcal Z_m|/M}
\longrightarrow +\infty.\;}
\]
更显式地，
\[
\boxed{\;
\frac{\mathsf{Col}_m-\frac1M}{|\mathcal Z_m|/M^2}
\gtrsim
\frac{2C_\varphi^2}{\tau(m+2)}\,\varphi^{m/2}.\;}
\]
\end{theorem}
\begin{proof}
由结论 \ref{con:xi-fold-chi2-constant-residual-energy}，
\[
M\Bigl(\mathsf{Col}_m-\frac1M\Bigr)\ge 2C_\varphi^2+o(1).
\]
另一方面，由定理 \ref{thm:fold-zero-density-sparse}，
\[
\frac{|\mathcal Z_m|}{M}
\le
\tau(m+2)\,\frac{F_{\lfloor (m+2)/2\rfloor}}{F_{m+2}}
=
\tau(m+2)\,O(\varphi^{-m/2}).
\]
将两式相除，即得
\[
\frac{M\bigl(\mathsf{Col}_m-\frac1M\bigr)}{|\mathcal Z_m|/M}
\to +\infty.
\]
而
\[
\frac{M\bigl(\mathsf{Col}_m-\frac1M\bigr)}{|\mathcal Z_m|/M}
=
\frac{\mathsf{Col}_m-\frac1M}{|\mathcal Z_m|/M^2},
\]
故两种写法完全等价。

最后，把上界
\(
|\mathcal Z_m|/M\le \tau(m+2)\,O(\varphi^{-m/2})
\)
写成倒数下界，再与
\(
M(\mathsf{Col}_m-1/M)\ge 2C_\varphi^2+o(1)
\)
合并，即得所示 \(\gtrsim\) 估计。证毕。
\end{proof}

\begin{theorem}[零点计数对碰撞缺口不存在消失模量]\label{cor:xi-fold-zero-density-no-bounded-control-on-normalized-collision}\leanverified{paper\_xi\_fold\_zero\_density\_no\_bounded\_control\_on\_normalized\_collision}
沿用定理 \ref{thm:xi-fold-zero-geometry-collision-gap-blindness} 的记号，并定义
\[
M_m:=F_{m+2},
\qquad
\delta_m:=\frac{|\mathcal Z_m|}{M_m},
\qquad
\Gamma_m:=M_m\!\left(\mathsf{Col}_m-\frac{1}{M_m}\right).
\]
则沿子列
\[
m\equiv 1\pmod 3
\]
有
\[
\boxed{\;
\delta_m\longrightarrow 0,
\qquad
\liminf_{\substack{m\to\infty\\ m\equiv 1\!\!\!\!\pmod 3}}
\Gamma_m
\ge
\frac{2C_\varphi^2}{\varphi}>0.\;}
\]
从而不存在任何函数 \(\eta:[0,1]\to[0,\infty)\) 满足
\[
\eta(t)\longrightarrow 0
\qquad (t\to 0),
\]
并且对所有充分大的 \(m\equiv 1\pmod 3\) 都有
\[
\Gamma_m\le \eta(\delta_m).
\]
换言之，即使谱零点的相对密度趋于 \(0\)，它也不足以连续控制折叠碰撞的归一化缺口；非零频率振幅层仍是不可删除的数据层。
\end{theorem}
\begin{proof}
当 \(m\equiv 1\pmod 3\) 时，\(m+2\equiv 0\pmod 3\)，故 Fibonacci 数 \(M_m=F_{m+2}\) 为偶数。于是由定理 \ref{thm:fold-zero-density-sparse}，
\[
\delta_m=\frac{|\mathcal Z_m|}{M_m}
\le
\tau(m+2)\,O(\varphi^{-m/2})
\longrightarrow 0
\qquad
\bigl(m\to\infty,\ m\equiv 1\!\!\!\!\pmod 3\bigr).
\]

另一方面，定理 \ref{thm:xi-fold-zero-geometry-collision-gap-blindness} 已给出更强的全序列下界
\[
\Gamma_m
=
M_m\!\left(\mathsf{Col}_m-\frac1{M_m}\right)
\ge
2C_\varphi^2+o(1).
\]
因此沿任意子列，特别是沿 \(m\equiv 1\pmod 3\)，都有
\[
\liminf_{\substack{m\to\infty\\ m\equiv 1\!\!\!\!\pmod 3}}\Gamma_m
\ge
2C_\varphi^2
>
\frac{2C_\varphi^2}{\varphi},
\]
从而所示盒装下界成立。

若存在函数 \(\eta\) 满足 \(\eta(t)\to 0\) 且对所有充分大的 \(m\equiv 1\pmod 3\) 都有
\[
\Gamma_m\le \eta(\delta_m),
\]
则由 \(\delta_m\to 0\) 可得
\[
\eta(\delta_m)\longrightarrow 0
\qquad
\bigl(m\to\infty,\ m\equiv 1\!\!\!\!\pmod 3\bigr),
\]
从而 \(\Gamma_m\to 0\) 沿该子列成立，这与上式的严格正下极限矛盾。故这样的消失模量不存在。证毕。
\end{proof}

\paragraph{末位统计充分性、黄金判别门槛与 fold-gauge 一次项的统一压缩}
在末位充分化、均匀基线二原子极限、规范群熵账本以及 fold-gauge 常数的 Bernoulli--Stirling 层级均已建立之后，还可把这些输入进一步压缩为一条直接服务于判别几何与规范体积审计的闭式链条。

\begin{theorem}[bin-fold 判别几何在末位上的统计充分性塌缩]\label{thm:xi-fold-lastbit-statistical-sufficiency-collapse}
\leanverified{paper\_xi\_fold\_lastbit\_statistical\_sufficiency\_collapse}
设
\[
\mu_m(w):=\frac{d_m^{\mathrm{bin}}(w)}{2^m},
\qquad
u_m(w):=\frac{1}{|X_m|}
\qquad (w\in X_m),
\]
并按末位分层
\[
A_i:=\{w\in X_m:\ w_m=i\}
\qquad (i=0,1).
\]
定义粗化测度
\[
\bar\mu_m
:=
\mu_m(A_0)\,u_m(\cdot\mid A_0)
\;+\;
\mu_m(A_1)\,u_m(\cdot\mid A_1).
\]
则存在常数 \(C>0\) 使
\[
\boxed{\;
D_{\mathrm{TV}}(\mu_m,\bar\mu_m)
\le
C\Bigl(\frac{\varphi}{2}\Bigr)^m.\;}
\]
特别地，统计实验
\[
(X_m,\mu_m,u_m)
\]
与其二元末位压缩
\[
\bigl(\{0,1\},(\ell_m)_\#\mu_m,(\ell_m)_\#u_m\bigr),
\qquad
\ell_m(w):=w_m,
\]
在全变差意义下指数等价；因而末位统计量 \(\ell_m\) 对 bin-fold 非均匀性是渐近充分的。
\end{theorem}
\begin{proof}
记
\[
K_m(i,\cdot):=\mathrm{Unif}(A_i)
\qquad (i=0,1).
\]
则
\[
K_m\circ (\ell_m)_\#\mu_m
=
\mu_m(A_0)\,u_m(\cdot\mid A_0)
+\mu_m(A_1)\,u_m(\cdot\mid A_1)
=
\bar\mu_m.
\]
而定理 \ref{thm:xi-fold-lastbit-asymptotic-sufficiency} 已给出
\[
D_{\mathrm{TV}}\!\bigl(\mu_m,K_m\circ(\ell_m)_\#\mu_m\bigr)
=
O\!\Bigl(\Bigl(\frac{\varphi}{2}\Bigr)^m\Bigr).
\]
将右端常数记为 \(C\) 即得所示估计；指数等价与渐近充分性亦随之成立。证毕。
\end{proof}

\begin{corollary}[相对稳定层均匀基线的黄金判别门槛]\label{cor:xi-fold-uniform-baseline-golden-threshold}
沿用定理 \ref{thm:xi-fold-lastbit-statistical-sufficiency-collapse} 的记号，则
\[
u_m(A_0)=\frac{F_{m+1}}{F_{m+2}}\longrightarrow \varphi^{-1},
\qquad
u_m(A_1)=\frac{F_m}{F_{m+2}}\longrightarrow \varphi^{-2},
\]
并且
\[
\mu_m(A_0)\longrightarrow \frac{\varphi}{\sqrt5},
\qquad
\mu_m(A_1)\longrightarrow \frac{1}{\varphi\sqrt5}.
\]
从而
\[
\boxed{\;
\lim_{m\to\infty}D_{\mathrm{TV}}(\mu_m,u_m)
=
1-\frac{2}{\sqrt5}.\;}
\]
若记等先验下判别 \(\mu_m\) 与 \(u_m\) 的最优 Bayes 错误概率为
\[
\mathsf P_{\mathrm{err}}^\ast(m)
:=
\frac{1-D_{\mathrm{TV}}(\mu_m,u_m)}{2},
\]
则
\[
\boxed{\;
\lim_{m\to\infty}\mathsf P_{\mathrm{err}}^\ast(m)
=
\frac{1}{\sqrt5}.\;}
\]
\end{corollary}
\begin{proof}
由定理 \ref{thm:xi-fold-lastbit-statistical-sufficiency-collapse}，
\[
D_{\mathrm{TV}}(\mu_m,u_m)
=
D_{\mathrm{TV}}(\bar\mu_m,u_m)
+O\!\Bigl(\Bigl(\frac{\varphi}{2}\Bigr)^m\Bigr).
\]
又因为 \(\bar\mu_m\) 与 \(u_m\) 在分层 \(A_0\sqcup A_1\) 上均为块常值分布，故
\[
D_{\mathrm{TV}}(\bar\mu_m,u_m)
=
\bigl|\mu_m(A_0)-u_m(A_0)\bigr|.
\]
其中 \(u_m(A_0)=|A_0|/|X_m|=F_{m+1}/F_{m+2}\) 与
\(
u_m(A_1)=F_m/F_{m+2}
\)
由 Fibonacci 计数立得；而
\(
\mu_m(A_0),\mu_m(A_1)
\)
的极限可由定理 \ref{thm:xi-fold-uniform-baseline-testing-constant}，或等价地由
定理 \ref{thm:xi-foldbin-pushforward-uniform-two-atom-ledger} 的二原子极限直接得到。于是
\[
\lim_{m\to\infty}D_{\mathrm{TV}}(\mu_m,u_m)
=
\left|\frac{\varphi}{\sqrt5}-\frac1\varphi\right|
=
1-\frac{2}{\sqrt5}.
\]
最后把该极限代入 Bayes 公式即可得到
\(
\lim_{m\to\infty}\mathsf P_{\mathrm{err}}^\ast(m)=1/\sqrt5
\)。
证毕。
\end{proof}

\leanverified{paper\_xi\_fold\_uniform\_baseline\_golden\_threshold}

\begin{theorem}[纤维信息项与规范体积密度的一次项显式化]\label{thm:xi-fold-kappa-gauge-first-order-constants}
\leanverified{paper\_xi\_fold\_kappa\_gauge\_first\_order\_constants}
令
\[
\kappa_m:=\sum_{w\in X_m}\mu_m(w)\log d_m^{\mathrm{bin}}(w),
\qquad
g_m:=2^{-m}\log |G_m^{\mathrm{bin}}|.
\]
则当 \(m\to\infty\) 时有
\[
\boxed{\;
\kappa_m
=
m\log\!\Bigl(\frac{2}{\varphi}\Bigr)
-\frac{\log\varphi}{\varphi\sqrt5}
+O\!\Bigl(\Bigl(\frac{\varphi}{2}\Bigr)^m\Bigr),\;}
\]
以及
\[
\boxed{\;
g_m
=
m\log\!\Bigl(\frac{2}{\varphi}\Bigr)
-1
-\frac{\log\varphi}{\varphi\sqrt5}
+\frac{1}{2\sqrt5}\Bigl(\frac{\varphi}{2}\Bigr)^m
\Bigl(
\varphi^2 m\log\!\Bigl(\frac{2}{\varphi}\Bigr)
+\varphi^2\log(2\pi)-\log\varphi
\Bigr)
+O\!\Bigl(\Bigl(\frac{\varphi}{2}\Bigr)^m\Bigr).\;}
\]
因此 \(\kappa_m\) 与 \(g_m\) 具有相同的一次主斜率
\(
\log(2/\varphi)
\)，
且其常数项之差严格等于 \(1\)。
\end{theorem}
\begin{proof}
第一式正是推论 \ref{cor:xi-foldbin-shannon-gauge-constant-mirror} 中
\(\kappa_m\) 的显式展开。

再由定理 \ref{thm:xi-foldbin-gauge-entropy-one-nat-law}，
\[
g_m
=
\kappa_m-1
+\frac{1}{2\sqrt5}\Bigl(\frac{\varphi}{2}\Bigr)^m
\Bigl(
\varphi^2 m\log\!\Bigl(\frac{2}{\varphi}\Bigr)
+\varphi^2\log(2\pi)-\log\varphi
\Bigr)
+O\!\Bigl(\Bigl(\frac{\varphi}{2}\Bigr)^m\Bigr).
\]
把 \(\kappa_m\) 的展开代入，即得第二式。证毕。
\end{proof}

\begin{theorem}[fold-gauge 的 \texorpdfstring{$2\pi$}{2pi} 估计序列具有最终单侧性与黄金公比锁定]\label{thm:xi-fold-gauge-2pi-estimator-geometric-lock}
\leanverified{paper\_xi\_fold\_gauge\_2pi\_estimator\_geometric\_lock}
沿用推论 \ref{cor:fold-bin-recover-2pi} 的规范常数列 \((C_m)_{m\ge 1}\)，并记
\[
E_m:=\log C_m-\log(2\pi).
\]
则存在 \(m_0\)，使得对全部 \(m\ge m_0\) 都有
\[
C_m>2\pi,
\qquad
C_{m+1}<C_m.
\]
更具体地，
\[
\boxed{\;
E_m
=
\frac{1}{3\varphi}\Bigl(\frac{\varphi}{2}\Bigr)^m
+O\!\Bigl(\Bigl(\frac{\varphi}{2}\Bigr)^{3m}\Bigr),\;}
\]
从而
\[
\boxed{\;
\frac{E_{m+1}}{E_m}
\longrightarrow
\frac{\varphi}{2}.\;}
\]
换言之，\(C_m\) 不仅收敛到 \(2\pi\)，而且其对数误差尾部形成公比为 \(\varphi/2\) 的黄金几何级数。
\end{theorem}
\begin{proof}
由定理 \ref{thm:xi-foldbin-gauge-constant-eventual-complete-monotonicity}，
对充分大的 \(m\) 已有
\[
C_m>2\pi,
\qquad
C_{m+1}<C_m.
\]
同一定理还给出更强展开
\[
E_m
=
\frac{1}{3\varphi}\Bigl(\frac{\varphi}{2}\Bigr)^m
-\frac{\sqrt5}{180}\Bigl(\frac{\varphi}{2}\Bigr)^{3m}
+O\!\Bigl(\Bigl(\frac{\varphi}{2}\Bigr)^{5m}\Bigr),
\]
故特别可写成所示的首项加三次余项形式。于是
\[
\frac{E_{m+1}}{E_m}
=
\frac{\frac{1}{3\varphi}(\varphi/2)^{m+1}+O((\varphi/2)^{3m+3})}
{\frac{1}{3\varphi}(\varphi/2)^m+O((\varphi/2)^{3m})}
\longrightarrow
\frac{\varphi}{2}.
\]
证毕。
\end{proof}

\paragraph{window-\texorpdfstring{$6$}{6} 群商实现的稳定子分层与 Lee--Yang 振幅常数的双二次阴影}
在 window-\(6\) 的纤维退化直方图、自由作用与线性核障碍已经锁定之后，任何试图把纤维分解解释为有限群轨道分解的方案都必须携带多层稳定子；另一方面，Lee--Yang 振幅常数 \(B\) 的真正六次性并不止于 \(\QQ(B)/\QQ\) 的次数陈述，其分裂域还显现出一条与三次分辨子域不同的额外二次阴影。

\begin{theorem}[window-\texorpdfstring{$6$}{6} 的任何群商实现都强制三层稳定子分层]\label{thm:xi-window6-group-quotient-isotropy-strata}
设 \(\Omega_6:=\{0,1\}^6\)，并令
\[
F_6:=\Fold^{\mathrm{bin}}_6:\Omega_6\to X_6.
\]
若存在某个有限群 \(G\) 在 \(\Omega_6\) 上作用，使其轨道分解恰等于 \(F_6\) 的纤维分解，则
\[
\boxed{\;12\mid |G|.\;}
\]
并且轨道稳定子必分成三种不同的阶层：对每个 \(\omega\in\Omega_6\)，有
\[
|\mathrm{Orb}_G(\omega)|\in\{2,3,4\}
\quad\Longleftrightarrow\quad
|\mathrm{Stab}_G(\omega)|\in\left\{\frac{|G|}{2},\frac{|G|}{3},\frac{|G|}{4}\right\}.
\]
更具体地，大小分别为 \(2,3,4\) 的轨道个数恰为
\[
8,\qquad 4,\qquad 9.
\]
因此任何这样的群商实现都不可能是自由作用，而只能是带有三层稳定子分层的离散 orbifold 型商。
\end{theorem}
\begin{proof}
由推论 \ref{cor:foldbin6-degeneracy-multivalued-nonfree}，\(F_6\) 的纤维划分既不可能来自自由有限群作用，也不可能来自阿贝尔线性陪集划分。另一方面，由推论 \ref{cor:fold-bin-m6-fiber-hist}，
\[
|\{w\in X_6:\ |F_6^{-1}(w)|=2\}|=8,\qquad
|\{w\in X_6:\ |F_6^{-1}(w)|=3\}|=4,\qquad
|\{w\in X_6:\ |F_6^{-1}(w)|=4\}|=9.
\]
若群 \(G\) 的轨道分解恰与纤维分解一致，则轨道大小只能取 \(2,3,4\)，且对应轨道数正是上述 \(8,4,9\)。

对任意 \(\omega\in\Omega_6\)，轨道--稳定子定理给出
\[
|\mathrm{Orb}_G(\omega)|=\frac{|G|}{|\mathrm{Stab}_G(\omega)|}.
\]
因此 \(2,3,4\) 都必须整除 \(|G|\)，遂有
\[
\mathrm{lcm}(2,3,4)=12\mid |G|.
\]
并且当轨道大小分别为 \(2,3,4\) 时，稳定子阶只能分别为
\[
\frac{|G|}{2},\qquad \frac{|G|}{3},\qquad \frac{|G|}{4}.
\]
这三者彼此不同，故点稳定子至少分成三层；又因为全部轨道只出现这三种大小，所以稳定子阶也恰只出现这三层。证毕。
\end{proof}

\begin{theorem}[Lee--Yang 振幅常数 \texorpdfstring{$B$}{B} 的六次性与双二次阴影]\label{thm:xi-leyang-amplitude-B-sextic-biquadratic-shadow}
\leanverified{paper\_xi\_leyang\_amplitude\_b\_sextic\_biquadratic\_shadow}
设
\[
f_B(X):=162X^6+1593X^4+1998X^2+1184.
\]
则
\[
\boxed{\;[\QQ(B):\QQ]=6.\;}
\]
并且 \(f_B\) 在 \(\QQ\) 上不可约，正是 \(B\) 的六次最小多项式；其判别式满足
\[
\boxed{\;
\Disc(f_B)=
-2^{16}\cdot 3^{22}\cdot 11^4\cdot 37^3\cdot 109^4,\;}
\]
故平方自由部分为
\[
\boxed{\;-37.\;}
\]
若 \(N\) 为 \(f_B\) 的分裂域，则
\[
\QQ(\sqrt{-37})\subset N.
\]
另一方面，\(N\) 同时包含 \(B^2\) 的三次最小多项式的分裂域，因而也包含其唯一二次分辨子域
\[
\QQ(\sqrt{-111}).
\]
于是
\[
\boxed{\;
\QQ(\sqrt{-37},\sqrt{-111})\subset N,
\qquad
12\mid [N:\QQ].\;}
\]
\end{theorem}
\begin{proof}
由定理 \ref{thm:xi-leyang-amplitude-B-true-quadratic-cover}，
\[
[\QQ(B):\QQ]=6.
\]
又因 \(B^2\) 满足三次方程
\[
g(T):=162T^3+1593T^2+1998T+1184,
\]
故 \(g(B^2)=0\)，亦即 \(f_B(B)=g(B^2)=0\)。由于 \(f_B\) 的次数也是 \(6\)，它必为 \(B\) 的最小多项式，从而在 \(\QQ\) 上不可约。

设 \(r_1,r_2,r_3\) 为 \(g\) 的根，则
\[
f_B(X)=162\prod_{j=1}^{3}(X^2-r_j),
\]
其六个根为 \(\pm \sqrt{r_1},\pm \sqrt{r_2},\pm \sqrt{r_3}\)。于是
\[
\Disc(f_B)
=
162^{10}\Bigl(\prod_{j=1}^{3}4r_j\Bigr)
\prod_{1\le i<j\le 3}(r_i-r_j)^4.
\]
另一方面，
\[
\Disc(g)=162^4\prod_{1\le i<j\le 3}(r_i-r_j)^2,
\qquad
\prod_{j=1}^{3}r_j=-\frac{1184}{162}.
\]
故可化为
\[
\Disc(f_B)
=
-64\cdot 162\cdot 1184\cdot \Disc(g)^2.
\]
再由推论 \ref{cor:xi-leyang-u-b-square-common-quadratic-resolvent}，
\[
\Disc(g)=
-2^2\cdot 3^9\cdot 11^2\cdot 37\cdot 109^2.
\]
代入即得
\[
\Disc(f_B)=
-2^{16}\cdot 3^{22}\cdot 11^4\cdot 37^3\cdot 109^4.
\]
因此其平方自由部分为 \(-37\)。分裂域 \(N\) 因而包含
\[
\QQ(\sqrt{\Disc(f_B)})=\QQ(\sqrt{-37}).
\]

再者，\(N\) 含有 \(f_B\) 的全部根 \(\pm\sqrt{r_j}\)，平方后即得到 \(g\) 的全部根 \(r_j\)，故 \(g\) 的分裂域 \(L_b\) 满足
\[
L_b\subset N.
\]
由推论 \ref{cor:xi-leyang-u-b-square-common-quadratic-resolvent}，\(L_b\) 的唯一二次子域是
\[
\QQ(\sqrt{-111}).
\]
于是 \(N\) 同时包含 \(\QQ(\sqrt{-37})\) 与 \(\QQ(\sqrt{-111})\)。由于这两个虚二次域不同，遂得
\[
\QQ(\sqrt{-37},\sqrt{-111})\subset N.
\]
又因 \([L_b:\QQ]=6\)，且 \(\QQ(\sqrt{-37})\not\subset L_b\)（因为 \(L_b\) 的唯一二次子域为 \(\QQ(\sqrt{-111})\)），故
\[
[L_b(\sqrt{-37}):\QQ]=12.
\]
最后注意 \(L_b(\sqrt{-37})\subset N\)，于是
\[
12\mid [N:\QQ].
\]
证毕。
\end{proof}

\paragraph{截断 Jensen 剖面、solenoid 终对象与有限可见证书的双重不完备性}
在 Jensen 缺陷的分布反演与有限半径绝对盲区、有限素数局部化群与其对偶 solenoid 的 \(\mathrm{Hom}\)-范畴分类，以及 prime-register 轴的无限外置语义已经建立之后，还可把“半径方向的失明”与“素数方向的失明”压缩为下列五条闭合结论。

\begin{theorem}[截断 Jensen 剖面的完全性与完全失明的精确分界]\label{thm:xi-time-part60e-truncated-jensen-profile-completeness-blindness}
设 \(F\) 为单位圆盘 \(\mathbb D\) 上非常值全纯函数，零点按重数记为 \(\{a\}\subset\mathbb D\)，并定义 Jensen 缺陷
\[
D_F(\rho):=\sum_{|a|<\rho}\log\frac{\rho}{|a|}
\qquad (0<\rho<1).
\]
固定 \(0<\rho_\star<1\)，并记截断剖面
\[
\mathcal J_{\rho_\star}(F):=D_F|_{(0,\rho_\star]}.
\]
再定义径向零测度
\[
\mu_F^{\mathrm{rad}}
:=
\sum_a \operatorname{mult}(a)\,\delta_{\log|a|}
\qquad\text{于 }(-\infty,0).
\]
则成立：
\begin{enumerate}
  \item \(\mathcal J_{\rho_\star}(F)\) 完全决定，且仅决定，
  \[
  \mu_F^{\mathrm{rad}}\big|_{(-\infty,\log\rho_\star)}.
  \]
  \item 对任意两个盘内全纯函数 \(F,G\)，有等价关系
  \[
  \mathcal J_{\rho_\star}(F)=\mathcal J_{\rho_\star}(G)
  \iff
  \mu_F^{\mathrm{rad}}\big|_{(-\infty,\log\rho_\star)}
  =
  \mu_G^{\mathrm{rad}}\big|_{(-\infty,\log\rho_\star)}.
  \]
  \item 因而，对零点多重集在闭环带
  \[
  \{z\in\mathbb D:\ \rho_\star\le |z|<1\}
  \]
  上所作的任意增删，都不会改变整个截断剖面 \(\mathcal J_{\rho_\star}\)。
\end{enumerate}
\end{theorem}
\begin{proof}
置
\[
u:=\log\rho<0,
\qquad
u_\star:=\log\rho_\star,
\qquad
\mathcal D_F(u):=D_F(e^u).
\]
由定理 \ref{thm:xi-time-part60a-jensen-defect-radial-measure-inversion}，
\[
\mathcal D_F(u)=\int_{(-\infty,u)}(u-s)\,\dd\mu_F^{\mathrm{rad}}(s),
\qquad
\mathcal D_F''=\mu_F^{\mathrm{rad}}
\]
在分布意义下于 \((-\infty,0)\) 成立。

因此，一旦 \(\mathcal D_F\) 已知于区间 \((-\infty,u_\star]\)，便可在开区间 \((-\infty,u_\star)\) 上取其分布二阶导数，从而唯一恢复
\[
\mu_F^{\mathrm{rad}}\big|_{(-\infty,u_\star)}.
\]
反之，若该限制测度已知，则对每个 \(u\le u_\star\)，势函数表示
\[
\mathcal D_F(u)=\int_{(-\infty,u)}(u-s)\,\dd\mu_F^{\mathrm{rad}}(s)
\]
只依赖于 \((-\infty,u_\star)\) 上的测度，因为积分核中始终要求 \(s<u\le u_\star\)。故第 \(1\) 条成立。

第 \(2\) 条是第 \(1\) 条的直接重述：若两条截断剖面相同，则其在 \((-\infty,u_\star)\) 上的分布二阶导数相同；反向由同一势函数公式立即得到。

对第 \(3\) 条，只需注意：若某个零点满足 \(|a|\ge \rho_\star\)，则对任意 \(\rho\le \rho_\star\) 都有 \(|a|<\rho\) 失效，故该零点在
\[
D_F(\rho)=\sum_{|a|<\rho}\log\frac{\rho}{|a|}
\]
中从不出现。因此它对整个 \((0,\rho_\star]\) 上的截断剖面完全不可见。证毕。
\end{proof}

\begin{theorem}[solenoid 终对象之间连续态射的完全分类]\label{thm:xi-time-part60e-solenoid-terminal-morphism-classification}
设 \(S,T\subset\mathbb P\) 为有限素数集，并记
\[
G_S:=\ZZ[S^{-1}],
\qquad
G_T:=\ZZ[T^{-1}],
\qquad
\Sigma_S:=\widehat{G_S},
\qquad
\Sigma_T:=\widehat{G_T}.
\]
则存在自然同构
\[
\operatorname{Hom}_{\mathrm{cts}}(\Sigma_T,\Sigma_S)
\cong
\operatorname{Hom}_{\ZZ}(G_S,G_T)
\cong
\begin{cases}
0,& S\not\subseteq T,\\[1mm]
G_T,& S\subseteq T.
\end{cases}
\]
更具体地：
\begin{enumerate}
  \item 存在非零连续群同态 \(\Sigma_T\to\Sigma_S\) 当且仅当 \(S\subseteq T\)；
  \item 在 \(S\subseteq T\) 时，每个连续同态都唯一对应于某个 \(x\in G_T\)，其对偶映射为
  \[
  m_x:G_S\to G_T,
  \qquad
  y\mapsto xy;
  \]
  \item 对应连续同态为满射当且仅当 \(x\neq 0\)；
  \item 对应连续同态为同构当且仅当 \(S=T\) 且 \(x\in G_T^\times\) 为 \(T\)-局部化单位。
\end{enumerate}
\end{theorem}
\begin{proof}
Pontryagin 对偶把紧交换群范畴与离散交换群范畴反变等价，因此
\[
\operatorname{Hom}_{\mathrm{cts}}(\Sigma_T,\Sigma_S)
\cong
\operatorname{Hom}_{\ZZ}(G_S,G_T).
\]
再由定理 \ref{thm:xi-cdim-localization-hom-category-classification}，后者恰在 \(S\subseteq T\) 时同构于 \(G_T\)，而在 \(S\not\subseteq T\) 时为零群。这给出前两条。

当 \(S\subseteq T\) 时，\(x\in G_T\) 对应的离散端同态正是
\[
m_x:G_S\to G_T,\qquad y\mapsto xy.
\]
由于 \(G_T\subset\QQ\) 无挠，\(m_x\) 在 \(x\neq 0\) 时为单射；对偶化后得到连续同态 \(\widehat{m_x}:\Sigma_T\to\Sigma_S\) 为满射。若 \(x=0\)，则对应映射显然为零映射，非满。故第 \(3\) 条成立。

最后，连续同态 \(\widehat{m_x}\) 为同构当且仅当 \(m_x\) 为离散群同构。由定理 \ref{thm:xi-cdim-localization-hom-category-classification}，这首先迫使 \(S=T\)；而当 \(S=T\) 时，\(m_x\) 为同构当且仅当 \(x\) 在局部化环 \(\ZZ[T^{-1}]\) 中可逆，即 \(x\in G_T^\times\)。证毕。
\end{proof}

\begin{corollary}[无穷 prime-register solenoid 的有限可见因子化]\label{cor:xi-time-part60e-infinite-solenoid-finite-visible-factorization}
设 \(S\subset\mathbb P\) 为可数无穷素数集，并记
\[
G_S:=\ZZ[S^{-1}],
\qquad
\Sigma_S:=\widehat{G_S}.
\]
则对任意有限子集 \(S_0\subset S\)，包含映射 \(G_{S_0}\hookrightarrow G_S\) 对偶化得到自然满射
\[
\pi_{S\to S_0}:\Sigma_S\twoheadrightarrow \Sigma_{S_0}.
\]
并且成立：
\begin{enumerate}
  \item 任意连续特征
  \[
  \chi:\Sigma_S\to \TT
  \]
  都经过某个有限子集 \(S_0\subset S\) 因子化，即存在
  \[
  \chi_0:\Sigma_{S_0}\to \TT
  \]
  使 \(\chi=\chi_0\circ\pi_{S\to S_0}\)。
  \item 更一般地，若 \(K\) 为紧 Abel Lie 群，则任意连续群同态
  \[
  f:\Sigma_S\to K
  \]
  都经过某个有限子集 \(S_0\subset S\) 因子化。
  \item 对任意有限子集 \(S_0\subset S\)，存在自然短正合列
  \[
  0\longrightarrow \prod_{p\in S\setminus S_0}\ZZ_p
  \longrightarrow \Sigma_S
  \xrightarrow{\ \pi_{S\to S_0}\ }
  \Sigma_{S_0}
  \longrightarrow 0.
  \]
  因而，任何有限族连续可见同态都不可能穷尽无穷 prime-register 核。
\end{enumerate}
\end{corollary}
\begin{proof}
由定义，
\[
G_S=\ZZ[S^{-1}]=\bigcup_{\substack{S_0\subset S\\ |S_0|<\infty}}G_{S_0}.
\]
因此每个 \(x\in G_S\) 的分母素数支撑都只落在某个有限子集 \(S_0\subset S\) 内。

对第 \(1\) 条，Pontryagin 对偶给出自然同构
\[
\operatorname{Hom}_{\mathrm{cts}}(\Sigma_S,\TT)\cong G_S.
\]
任取连续特征 \(\chi\)，令 \(x\in G_S\) 为其对偶元素。由上一步存在有限 \(S_0\subset S\) 使 \(x\in G_{S_0}\)。再对偶化有限情形的定理 \ref{thm:xi-time-part60e-solenoid-terminal-morphism-classification}，便得到 \(\chi\) 经由 \(\pi_{S\to S_0}\) 因子化。

对第 \(2\) 条，设 \(f:\Sigma_S\to K\) 为连续同态。由于 \(K\) 为紧 Abel Lie 群，其对偶群 \(\widehat K\) 为有限生成离散阿贝尔群。于是 \(f^\vee:\widehat K\to G_S\) 的像由有限多个元素生成，而每个生成元都落在某个有限 \(G_{S_j}\) 中；取这些 \(S_j\) 的并得有限集合 \(S_0\subset S\)，便有
\[
\operatorname{im}(f^\vee)\subseteq G_{S_0}.
\]
故 \(f^\vee\) 经由包含 \(G_{S_0}\hookrightarrow G_S\) 因子化。再对偶化即得
\[
f=f_0\circ \pi_{S\to S_0}
\]
对某个连续同态 \(f_0:\Sigma_{S_0}\to K\) 成立。

对第 \(3\) 条，只需计算离散端商群。由
\[
G_S/\ZZ\cong \bigoplus_{p\in S}C_{p^\infty},
\qquad
G_{S_0}/\ZZ\cong \bigoplus_{p\in S_0}C_{p^\infty},
\]
可得
\[
G_S/G_{S_0}
\cong
\bigoplus_{p\in S\setminus S_0}C_{p^\infty}.
\]
再取 Pontryagin 对偶，并用
\[
\widehat{C_{p^\infty}}\cong \ZZ_p,
\qquad
\widehat{\bigoplus_{p\in S\setminus S_0}C_{p^\infty}}
\cong
\prod_{p\in S\setminus S_0}\ZZ_p,
\]
便得到所示短正合列。

若给定有限族连续可见同态 \(f_j:\Sigma_S\to K_j\)（其中每个 \(K_j\) 都是紧 Abel Lie 群），则由第 \(2\) 条可取有限 \(S_j\subset S\) 使每个 \(f_j\) 都经由 \(\pi_{S\to S_j}\) 因子化。令 \(S_0:=\bigcup_j S_j\)，则全部 \(f_j\) 共同经由 \(\pi_{S\to S_0}\) 因子化。由于 \(S\setminus S_0\) 仍为无穷集，核
\[
\ker(\pi_{S\to S_0})\cong \prod_{p\in S\setminus S_0}\ZZ_p
\]
为无限紧群，因此这些有限可见同态不可能把 \(\Sigma_S\) 分离到有限歧义。证毕。
\end{proof}

\begin{theorem}[有限半径证书与有限素数证书的统一不可完备性]\label{thm:xi-time-part60e-finite-radius-finite-prime-certificate-incompleteness}
固定 \(0<\rho_\star<1\)，并取可数无穷素数集 \(S\subset\mathbb P\) 与有限子集 \(S_0\subset S\)。考虑对象对
\[
(F,\xi),
\]
其中 \(F\) 为单位圆盘上的有限 Blaschke 乘积，\(\xi\in \Sigma_S\)。定义联合有限证书
\[
\mathcal C_{\rho_\star,S_0}(F,\xi)
:=
\Bigl(
\mathcal J_{\rho_\star}(F),\ \pi_{S\to S_0}(\xi)
\Bigr).
\]
则每个非空纤维
\[
\mathcal C_{\rho_\star,S_0}^{-1}(\eta,u)
\]
至少是可数无限集。更精确地，对任意固定基对象 \((F_0,\xi_0)\)，存在一列对象
\[
\{(F_n,\xi_n)\}_{n\ge 1}
\]
使得：
\begin{enumerate}
  \item \(\mathcal C_{\rho_\star,S_0}(F_n,\xi_n)\) 全部相同；
  \item \(F_n\) 两两不同；
  \item \(\xi_n\) 两两不同。
\end{enumerate}
\end{theorem}
\begin{proof}
固定任意基对象 \((F_0,\xi_0)\)。

先处理 Jensen 侧。取一列两两不同的实数
\[
a_n\in(\rho_\star,1)
\qquad(n\ge 1),
\]
并定义有限 Blaschke 因子
\[
B_{a_n}(z):=\frac{z-a_n}{1-a_n z},
\qquad
F_n:=F_0\cdot B_{a_n}.
\]
由于每个新零点 \(a_n\) 都满足 \(|a_n|\ge \rho_\star\)，定理 \ref{thm:xi-time-part60e-truncated-jensen-profile-completeness-blindness} 第 \(3\) 条给出
\[
\mathcal J_{\rho_\star}(F_n)=\mathcal J_{\rho_\star}(F_0)
\qquad(n\ge 1).
\]
而 \(\{a_n\}\) 两两不同，故 \(F_n\) 两两不同。

再处理 arithmetic 侧。由推论 \ref{cor:xi-time-part60e-infinite-solenoid-finite-visible-factorization}，
\[
\ker(\pi_{S\to S_0})
\cong
\prod_{p\in S\setminus S_0}\ZZ_p
\]
为无限紧阿贝尔群。任选一列两两不同的元素
\[
\eta_n\in \ker(\pi_{S\to S_0})
\qquad(n\ge 1),
\]
并置
\[
\xi_n:=\xi_0+\eta_n.
\]
则
\[
\pi_{S\to S_0}(\xi_n)=\pi_{S\to S_0}(\xi_0)
\qquad(n\ge 1),
\]
且 \(\xi_n\) 两两不同。

于是对全部 \(n\ge 1\)，都有
\[
\mathcal C_{\rho_\star,S_0}(F_n,\xi_n)
=
\mathcal C_{\rho_\star,S_0}(F_0,\xi_0).
\]
这就构造出同一纤维中的可数无限列，故每个非空纤维至少可数无限。证毕。
\end{proof}

\begin{theorem}[有限可见统计塌缩与无限算术核强制的统一二分律]\label{thm:xi-time-part60e-finite-visible-collapse-infinite-padic-ledger-dichotomy}
固定可数无穷素数集 \(S\subset\mathbb P\)，并定义径向--算术对象类
\[
\mathfrak O_S^{\mathrm{rad}}
:=
\Bigl\{
(\mu,\xi):
\mu \text{ 为某个有限 Blaschke 乘积的径向零测度， }
\xi\in\Sigma_S
\Bigr\}.
\]
对任意这类 \(\mu\) 与任意 \(0<\rho_\star<1\)，定义其截断 Jensen 势函数
\[
\mathcal J_{\rho_\star}(\mu)(\rho)
:=
\int_{(-\infty,\log\rho)}(\log\rho-s)\,\dd\mu(s)
\qquad (0<\rho\le \rho_\star).
\]
则成立：
\begin{enumerate}
  \item 任意经过某个有限证书
  \[
  (\mu,\xi)\longmapsto
  \bigl(\mathcal J_{\rho_\star}(\mu),\ \pi_{S\to S_0}(\xi)\bigr),
  \qquad
  0<\rho_\star<1,\ \ S_0\subset S\ \text{有限},
  \]
  的可见不变量，都不是 \(\mathfrak O_S^{\mathrm{rad}}\) 上的完备不变量；事实上，在有限 Blaschke 子类上，每个非空证书纤维都至少可数无限。
  \item 若要得到径向--算术层面的完备可逆证书，则在解析侧必须保留全连续 Jensen 剖面
  \[
  \rho\longmapsto
  \int_{(-\infty,\log\rho)}(\log\rho-s)\,\dd\mu(s)
  \qquad(0<\rho<1),
  \]
  在算术侧必须保留全部有限素数投影
  \[
  \{\pi_{S\to S_0}(\xi)\}_{S_0\subset S,\ |S_0|<\infty}.
  \]
  等价地，算术侧必须保留整个 \(\xi\in\Sigma_S\)。
  \item 因而，有限可见统计在径向半径方向与素数方向上同时塌缩；而完备可逆账本在算术侧被迫是无限 \(p\)-进的。
\end{enumerate}
\end{theorem}
\begin{proof}
第 \(1\) 条由定理 \ref{thm:xi-time-part60e-finite-radius-finite-prime-certificate-incompleteness} 直接给出；该定理甚至已经在有限 Blaschke 子类上给出显式的可数无限纤维。

对第 \(2\) 条，取 \(\mu=\mu_F^{\mathrm{rad}}\) 为某个有限 Blaschke 乘积 \(F\) 的径向零测度。由定理 \ref{thm:xi-time-part60e-truncated-jensen-profile-completeness-blindness} 可知：任取 \(\rho_\star<1\)，截断剖面只能恢复
\[
\mu\big|_{(-\infty,\log\rho_\star)};
\]
只有让 \(\rho_\star\uparrow 1\)，即保留整个 \(0<\rho<1\) 上的 Jensen 势函数，才能恢复全部径向零测度 \(\mu\)。

算术侧则由
\[
G_S=\bigcup_{\substack{S_0\subset S\\ |S_0|<\infty}}G_{S_0}
=
\varinjlim_{\substack{S_0\subset S\\ |S_0|<\infty}}G_{S_0}
\]
与 Pontryagin 对偶的反变等价得到
\[
\Sigma_S
\cong
\varprojlim_{\substack{S_0\subset S\\ |S_0|<\infty}}\Sigma_{S_0}.
\]
因此，\(\xi\in\Sigma_S\) 与其全部有限投影族
\[
\{\pi_{S\to S_0}(\xi)\}_{S_0\subset S,\ |S_0|<\infty}
\]
完全等价；任何有限截断都会丢失推论 \ref{cor:xi-time-part60e-infinite-solenoid-finite-visible-factorization} 中的无限核
\(
\prod_{p\in S\setminus S_0}\ZZ_p
\)。

第 \(3\) 条只是前两条的合并表述：在解析侧，有限半径截断留下不可见外环；在算术侧，有限素数截断留下不可见无限 \(p\)-进核。故有限可见统计必然塌缩，而完备可逆账本在算术侧必然保持无限 \(p\)-进结构。证毕。
\end{proof}

\paragraph{有限视界 classical RH 审计、谱替身与共尾 prime-ledger 的分离障碍}
在有限半径/有限素数证书的不完备性、无一致 realization 维数上界时的谱替身分岔、纯相位制度下的 \(\mathsf{NULL}\) 类型安全，以及扭子塔首次偏离深度的 Busy--Beaver 下界都已建立之后，还可把有限可见审计的障碍机制压缩为如下七条终端结论。它们分别刻画有限视界完备模数的不可计算增长、有限可见语言中的谱替身不可分辨、可靠性/完备性/可计算性的三难并立不可能、障碍来源的精确定位、纯相位制度下离切片信息的二歧、共尾 prime-ledger 的有限视界不可重构，以及相应的视界--账本分离原则。

\begin{theorem}[有限视界完备模数的 Busy--Beaver 障碍]\label{thm:xi-time-part60eb-finite-horizon-complete-modulus-busybeaver}
\leanverified{paper\_xi\_time\_part60eb\_finite\_horizon\_complete\_modulus\_busybeaver}
设 \(\mathcal V\) 为一种 classical RH 审计制度。对每个输入长度上界 \(\ell\)，\(\mathcal V\) 只允许访问受同一资源参数控制的
\[
\text{有限 prime-stage } \Sigma_S,\qquad
\text{有限 Toeplitz 主子式阶数},\qquad
\text{有限 dyadic 深度},\qquad
\text{有限 torsion 层级}.
\]
若存在函数 \(\Phi:\NN\to\NN\) 使得：对所有长度不超过 \(\ell\) 的输入实例，审计制度 \(\mathcal V\) 仅访问资源不超过 \(\Phi(\ell)\) 的上述四类有限视界，便总能给出正确判定，则称 \(\Phi\) 为 \(\mathcal V\) 的完备模数。

若 \(\mathcal V\) 在命题 \ref{prop:xi-time-part9g-deviation-depth-computability-separation} 的 torsion-tower 编码所给出的停机型偏离实例族上同时可靠且完备，则 \(\Phi\) 不可计算。更强地，存在常数 \(c>0\) 使得对充分大的 \(\ell\)，有
\[
\Phi(\ell)\ \ge\ L\,\mathrm{BB}(\ell-c)+1.
\]
\end{theorem}
\begin{proof}
沿用命题 \ref{prop:xi-time-part9g-first-deviation-depth} 与命题
\ref{prop:xi-time-part9g-deviation-depth-computability-separation}
的记号。对每个编码实例 \((M,x)\)，其有限层 torsion 投影
\[
Q_{M,x}^{(n)}
\]
都是全可计算的，并且
\[
\delta(M,x)<\infty
\quad\Longleftrightarrow\quad
M(x)\ \text{停机}.
\]

现取一个停机实例 \((M,x)\)，并记其首次偏离深度为 \(\delta(M,x)\)。由命题
\ref{prop:xi-time-part9g-first-deviation-depth}，
在全部层级 \(n<\delta(M,x)\) 上，编码对象与非停机基线对象的 torsion 投影完全一致。由于 prime-stage、有限 Toeplitz 主子式、有限 dyadic 深度与有限 torsion 层级都被同一资源上界控制，只要资源预算严格小于 \(\delta(M,x)\)，审计制度就仍停留在“首次偏离尚未出现”的有限视界之内，因而无法把该停机实例与对应的非停机基线区分开来。

若 \(\mathcal V\) 在这族实例上既可靠又完备，则对每个长度不超过 \(\ell\) 的停机实例都必须在资源 \(\Phi(\ell)\) 以内检出首次偏离；否则它就会把某个停机实例与非停机基线混淆。故必有
\[
\delta(M,x)\le \Phi(|(M,x)|)
\]
对全部停机实例成立。于是 \(\Phi\) 支配了一条统一验证深度函数。命题
\ref{prop:xi-time-part9g-deviation-depth-computability-separation}
第 \(3\) 条立即给出
\[
\Phi(\ell)\ \ge\ L\,\mathrm{BB}(\ell-c)+1
\]
对充分大的 \(\ell\) 成立，从而 \(\Phi\) 不可计算。证毕。
\end{proof}

\begin{theorem}[有限可见语言中的谱替身不可分辨]\label{thm:xi-time-part60eb-finite-visible-language-spectral-surrogates}
设 \(\mathcal L_{\mathrm{vis}}\) 为由下列可见对象生成的对象语言：
\[
\text{(i) disk-visible 相位输出},\qquad
\text{(ii) 有限多个有理频率角色},\qquad
\text{(iii) 有限阶 Toeplitz/Gram 主子式},\qquad
\text{(iv) 有限 dyadic 证书对象}.
\]
再设 admissible completion 族 \(\mathfrak X\) 在文稿既定的 Toeplitz--Carath\'eodory 审计语义下满足：其最小 realization 维数不存在一致上界。

则对任意有限子语言 \(\mathcal L_0\subset \mathcal L_{\mathrm{vis}}\)，存在两个 completion 对象
\[
X^+,X^-\in\mathfrak X
\]
满足
\[
X^+\equiv_{\mathcal L_0}X^-,
\]
其中 \(\equiv_{\mathcal L_0}\) 表示二者在 \(\mathcal L_0\) 的全部可见项与关系解释上完全一致；
但 \(X^+\) 满足全局 Toeplitz--Carath\'eodory 正性，而 \(X^-\) 不满足。等价地，在本文的 classical RH 语义中，真值不能由任意固定有限可见语言完全分离。
\end{theorem}
\leanverified{paper\_xi\_time\_part60eb\_finite\_visible\_language\_spectral\_surrogates}
\begin{proof}
先把 \(\mathcal L_0\) 所访问的数据压缩为一个有限审计切片。由推论 \ref{cor:xi-time-part60e-infinite-solenoid-finite-visible-factorization}，任意有限族连续可见同态都必经由某个有限 prime-stage 因子化。另一方面，命题 \ref{prop:xi-time-part62d-universal-solenoid-phase-completed-scan} 第 \(3\) 条说明：任意有限有理频率扫描都只依赖某个有限层因子。于是，disk-visible 相位输出与有限多个有理频率角色可共同压缩到某个有限 prime-stage。

有限阶 Toeplitz/Gram 主子式与有限 dyadic 证书对象只涉及有限多个可见坐标以及有限次代数运算，因此在把上述有限 prime-stage 进一步放大后，也都经由同一有限视界因子化。故 \(\mathcal L_0\) 所能读取的全部信息可等价压缩为一个固定有限审计规则 \(\mathcal A_{\mathcal L_0}\)。

再由定理 \ref{thm:xi-spectral-surrogate-finite-audit-undecidable}，在不存在一致 realization 维数上界时，任意固定有限审计都存在两组谱替身对象，它们在该审计访问的一切有限数据上完全一致，而全局 Toeplitz--PSD 真值相反。又由定理 \ref{thm:xi-slice-horizon-carath-toeplitz}，全局 Toeplitz--Carath\'eodory 正性与相应的 RH 真值等价。于是存在 \(X^+,X^-\in\mathfrak X\) 使
\[
X^+\equiv_{\mathcal L_0}X^-,
\]
但其全局真值相反。证毕。
\end{proof}

\begin{corollary}[无一致 realization 维数界时有限视界 RH 审计的三难并立不可能]\label{cor:xi-time-part60eb-finite-horizon-rh-audit-trilemma}
\leanverified{paper\_xi\_time\_part60eb\_finite\_horizon\_rh\_audit\_trilemma}
在无一致 realization 维数上界的前提下，任何把 classical RH 限制在 \(\mathcal L_{\mathrm{vis}}\) 中进行的有限视界审计制度，都不可能同时满足：
\[
\text{可靠性},\qquad
\text{完备性},\qquad
\text{可计算性}.
\]
\end{corollary}
\begin{proof}
由定理 \ref{thm:xi-time-part60eb-finite-visible-language-spectral-surrogates}，任意固定有限可见语言都存在不可分辨而全局真值相反的谱替身，故若坚持停留在有限视界内，则完备性失效。

若试图通过引入一个统一的完备模数来恢复完备性，则在 torsion-tower 编码的停机型偏离实例族上，定理 \ref{thm:xi-time-part60eb-finite-horizon-complete-modulus-busybeaver} 强制该模数具有 Busy--Beaver 级不可计算增长。故可计算性失效。

因此，在无一致 realization 维数界时，上述三者不能并存。证毕。
\end{proof}

\begin{proposition}[决定性障碍并非“处于无限维背景”本身]\label{prop:xi-time-part60eb-obstruction-not-ambient-infinity}
\leanverified{paper\_xi\_time\_part60eb\_obstruction\_not\_ambient\_infinity}
在本文框架内，classical RH 的 no-go 机制并不等价于“对象处于无限维背景”这一事实本身。真正不可压缩的障碍是：
\[
\text{(a) realization 维数无一致上界},\qquad
\text{(b) 可见层对 prime-ledger 核失真},\qquad
\text{(c) 因而全阶 Toeplitz 量词无法统一消去为有限证书}.
\]
更精确地，若一族候选 completion 对象具有一致 realization 维数上界 \(D\)，则存在统一截断阶
\[
N_0\le 2D
\]
使全局 Toeplitz--Carath\'eodory 正性等价于有限个主子式条件。故“处于无限维背景中”本身并不推出制度性不可判定。
\end{proposition}
\begin{proof}
若最小 realization 维数有一致上界 \(D\)，则定理 \ref{thm:xi-toeplitz-psd-coherence-horizon-threshold} 与定理 \ref{thm:xi-oracle-collapse-toeplitz-psd-finite-truncation} 同时给出：存在统一阈值 \(N_0\le 2D\)，使
\[
\forall N\ge 0,\ \mathbf T_N\succeq 0
\quad\Longleftrightarrow\quad
\forall N\le N_0,\ \mathbf T_N\succeq 0.
\]
于是全局正性层级在该有界类上可被压缩为有限证书对象。

反之，若一致维数上界失效，则定理 \ref{thm:xi-spectral-surrogate-finite-audit-undecidable} 给出有限审计不可判定的谱替身；而定理 \ref{thm:xi-time-part67-continuous-phase-external-arithmetic-kernel} 与推论 \ref{cor:xi-time-part60e-infinite-solenoid-finite-visible-factorization} 表明，连通可见层并不忠实保留外置 prime-ledger 核。故真正导致 no-go 的，是“无一致维数界”与“核失真”这两条结构性机制的合成，而非背景空间的无限维性本身。证毕。
\end{proof}

\begin{theorem}[纯相位制度下离切片信息的二歧结构]\label{thm:xi-time-part60eb-pure-phase-offslice-bifurcation}
\leanverified{paper\_xi\_time\_part60eb\_pure\_phase\_offslice\_bifurcation}
在纯相位可见协议中，任何关于离切片信息的断言都只能以如下两种方式之一出现：
\[
\text{(i) 显式外置模式轴 }\rho=\log|u|,\qquad
\text{(ii) 作为 }\mathsf{NULL}\text{ 或失败见证被外置。}
\]
不存在第三种纯相位内部表述，可在不增加模式轴的前提下稳定承载离切片信息。并且，若企图以统一可计算的有限模式轴预算来恢复全部离切片信息，则该预算必受定理 \ref{thm:xi-time-part60eb-finite-horizon-complete-modulus-busybeaver} 的 Busy--Beaver 障碍支配。
\end{theorem}
\begin{proof}
命题 \ref{prop:xi-offcritical-falsifiable-restatement} 已经给出严格的接口二分：离切片断言在体系内只能以“模式轴外延”或“\(\mathsf{NULL}\) 失败见证”两种方式出现。定理 \ref{thm:xi-offset-null-type-safety} 与命题 \ref{prop:xi-null-three-way} 进一步说明：若拒绝外置相应的正交记录轴，则该自由度在类型层不可被内部吸收，只能稳定退化为 \(\mathsf{NULL}\)。故不存在第三种纯相位内部表述。

若反设存在一条统一可计算的有限模式轴预算，能够对某族输入完全恢复离切片信息，则它与其余有限视界资源合并后便给出一条可计算的完备模数。把此模数施用于 torsion-tower 编码的停机型偏离实例族，即与定理 \ref{thm:xi-time-part60eb-finite-horizon-complete-modulus-busybeaver} 矛盾。故任何这类统一预算都必受 Busy--Beaver 级障碍支配。证毕。
\end{proof}

\begin{theorem}[共尾 prime-ledger 的有限视界不可重构性]\label{thm:xi-time-part60eb-cofinal-prime-ledger-finite-horizon-nonreconstruction}
\leanverified{paper\_xi\_time\_part60eb\_cofinal\_prime\_ledger\_finite\_horizon\_nonreconstruction}
设 \(\mathcal O\) 为任意有限族可见观测量，其成员皆由 disk-visible 输出、有限有理频率角色、有限阶 Toeplitz/Gram 主子式与有限 dyadic 证书对象编译而成。则存在一个有限素数集 \(S_0\) 使得
\[
\forall O\in\mathcal O,\qquad
O\ \text{都经由}\ \Sigma_{S_0}\ \text{因子化}.
\]

进一步，对任意严格共尾扩张链
\[
S_0\subsetneq S_1\subsetneq S_2\subsetneq \cdots,
\qquad
\bigcup_{n\ge 0}S_n=\mathbb P,
\]
相应 solenoid \(\Sigma_{S_n}\) 在全部 \(\mathcal O\)-观测上完全不可区分；然而它们的 prime-ledger 核
\[
K_{S_n}\cong \prod_{p\in S_n}\ZZ_p
\]
两两不可同构，并且每个 \(S_n\) 都可由 \(K_{S_n}\) 唯一恢复。
\end{theorem}
\begin{proof}
由推论 \ref{cor:xi-time-part60e-infinite-solenoid-finite-visible-factorization}，任意有限族连续可见同态都共同经由某个有限 prime-stage \(\Sigma_{S_0}\) 因子化。命题 \ref{prop:xi-time-part62d-universal-solenoid-phase-completed-scan} 第 \(3\) 条又说明：任意有限有理频率扫描本身就是某个有限层因子观测。于是 disk-visible 输出与有限有理频率角色都经由某个有限阶段因子化。

有限阶 Toeplitz/Gram 主子式以及有限 dyadic 证书对象只依赖于有限多个可见坐标和有限次代数运算；把所需有限阶段一并并入 \(S_0\) 后，它们也都经由 \(\Sigma_{S_0}\) 因子化。故对任意 \(O\in\mathcal O\)，都有
\[
O=O_0\circ \pi_{S_n\to S_0}
\qquad(n\ge 0)
\]
对某个 \(O_0\) 成立。这证明了 \(\Sigma_{S_n}\) 在全部 \(\mathcal O\)-观测上不可区分。

另一方面，若 \(S_n\neq S_m\)，则存在素数 \(p\) 使 \(p\in S_n\triangle S_m\)。对任意有限素数集 \(S\)，有
\[
K_S\cong \prod_{q\in S}\ZZ_q,
\]
故
\[
p\in S
\quad\Longleftrightarrow\quad
K_S/pK_S\neq 0
\quad\Longleftrightarrow\quad
\operatorname{Hom}_{\mathrm{cts}}(K_S,\ZZ/p\ZZ)\neq 0.
\]
因此 \(K_{S_n}\) 与 \(K_{S_m}\) 的 \(p\)-主因子结构不同，从而二者不同构；并且 \(S\) 可由判别全部素数 \(p\) 的上述等价条件唯一恢复。证毕。
\end{proof}

\begin{corollary}[视界--账本分离原则]\label{cor:xi-time-part60eb-horizon-ledger-separation-principle}
\leanverified{paper\_xi\_time\_part60eb\_horizon\_ledger\_separation\_principle}
若 classical RH 的真值最终依赖于共尾 prime-ledger 的整体几何，而审计协议仅驻留于有限视界可见层，则该真值的完全恢复必然失败。更准确地说，失败机制并非“样本不足”，而是可见函子对内部 prime-ledger 的非忠实压缩。
\end{corollary}
\begin{proof}
定理 \ref{thm:xi-time-part60eb-cofinal-prime-ledger-finite-horizon-nonreconstruction} 表明：任意有限视界可见层都只能看到某个有限 prime-stage，而不同的共尾 prime-ledger 可以在全部有限可见观测上保持不可区分。若 classical RH 的真值确实依赖于这种共尾账本几何，则任何忠实恢复过程都必须在有限视界内区分这些对象；这与定理 \ref{thm:xi-time-part60eb-cofinal-prime-ledger-finite-horizon-nonreconstruction} 相矛盾。

若进一步无一致 realization 维数上界，则定理 \ref{thm:xi-time-part60eb-finite-visible-language-spectral-surrogates} 把这种函子性失真强化为一对有限可见完全同像而全局 Toeplitz--Carath\'eodory 真值相反的谱替身对象。故所述恢复失败是结构性的，而非样本数量意义下的偶然不足。证毕。
\end{proof}

\noindent
综上，本文中的有限可见障碍机制可被严格压缩为一条共同结构：一方面，只要可靠完备性试图覆盖 torsion-tower 编码的停机型偏离实例，就必然遭遇 Busy--Beaver 级不可计算完备模数；另一方面，只要一致 realization 维数上界缺失，有限可见语言便不可避免地出现谱替身，而纯相位制度对离切片信息的处理又只能在“模式轴外延”与“\(\mathsf{NULL}\) 失败见证”之间二歧。于是真正不可压缩的障碍并不在于抽象的无限维背景本身，而在于无界 realization 维数、prime-ledger 核的可见失真以及有限视界预算的 Busy--Beaver 增长三者的共同作用。

\paragraph{坏模数的倍数密度与 window-\texorpdfstring{$6$}{6} 长根扇区的隐藏退化}
Dirichlet 扭曲坏模数沿整除链的传播，与 window-\(6\) 的 \(B_3\) 长根扇区在纤维多重度层上的残余对称，可进一步压缩为下列一组结构性结论。

\begin{corollary}[整除传播单独强迫坏模数集具有正下密度]\label{cor:xi-dirichlet-bad-moduli-positive-lower-density-by-divisibility}
\leanverified{paper\_xi\_dirichlet\_bad\_moduli\_positive\_lower\_density\_by\_divisibility}
沿用命题 \ref{prop:real-input-40-output-dirichlet-nonrh} 与结论 \ref{con:xi-time-part9l-dirichlet-twist-divisibility-monotonicity} 的记号，定义
\[
\mathcal B:=\{m\ge 2:\ \theta_m>\tfrac12\}.
\]
则 \(\mathcal B\neq\varnothing\)，并且若记
\[
m_0:=\min \mathcal B,
\]
则
\[
\underline d(\mathcal B)
:=
\liminf_{X\to\infty}\frac{\#(\mathcal B\cap [1,X])}{X}
\ge
\frac{1}{m_0}
>
0.
\]
换言之，即使不调用大奇素数模数的渐近展开，仅由“坏模数沿倍数链传播”这一整除刚性，也已足以推出坏模数集的正下密度。
\end{corollary}
\begin{proof}
由命题 \ref{prop:real-input-40-output-dirichlet-nonrh}，存在某个模数 \(m\ge 2\) 满足 \(\theta_m>\tfrac12\)，故 \(\mathcal B\neq\varnothing\)，从而最小元 \(m_0\) 存在。
再由结论 \ref{con:xi-time-part9l-dirichlet-twist-divisibility-monotonicity}，若 \(m_0\in\mathcal B\)，则其任意倍数 \(rm_0\)（\(r\ge 1\)）仍属于 \(\mathcal B\)。因此
\[
m_0\NN_{\ge 1}\subseteq \mathcal B.
\]
于是对任意 \(X\ge 1\)，有
\[
\#(\mathcal B\cap [1,X])\ge \left\lfloor \frac{X}{m_0}\right\rfloor.
\]
两边除以 \(X\) 并取 \(\liminf\)，即得
\[
\underline d(\mathcal B)\ge \frac{1}{m_0}.
\]
证毕。
\end{proof}

\begin{theorem}[window-\texorpdfstring{$6$}{6} dyadic 微态质量的 \texorpdfstring{$16+36+12$}{16+36+12} 精确三块分解]\label{thm:xi-window6-dyadic-microstate-exact-three-block-mass-split}
\leanverified{paper\_xi\_window6\_dyadic\_microstate\_exact\_three\_block\_mass\_split}
沿用
\[
\mathrm{Fold}^{\mathrm{bin}}_6:\{0,\dots,63\}\to X_6
\]
与 \(V_6\) 的记号。则
\[
2^6
=
\sum_{\substack{w\in X_6\\ w_6=1}} d_6^{\mathrm{bin}}(w)
+
\sum_{\substack{w\in X_6\\ w_6=0,\ V_6(w)\le 8}} d_6^{\mathrm{bin}}(w)
+
\sum_{\substack{w\in X_6\\ w_6=0,\ V_6(w)>8}} d_6^{\mathrm{bin}}(w)
=
16+36+12.
\]
等价地，全部 \(64\) 个 dyadic 微态在 \(\mathrm{Fold}^{\mathrm{bin}}_6\) 下被精确分解为三块总质量分别为 \(16\)、\(36\)、\(12\) 的互斥簇。进一步，在均匀微态基线下有
\[
\mathbb P(w_6=1)=\frac14,\qquad
\mathbb P(w_6=0,\ V_6(w)\le 8)=\frac{9}{16},\qquad
\mathbb P(w_6=0,\ V_6(w)>8)=\frac{3}{16}.
\]
\end{theorem}
\begin{proof}
由推论 \ref{cor:fold-bin-m6-fiber-hist}，
\[
\#\{w\in X_6:\ d_6^{\mathrm{bin}}(w)=2\}=8,\qquad
\#\{w\in X_6:\ d_6^{\mathrm{bin}}(w)=3\}=4,\qquad
\#\{w\in X_6:\ d_6^{\mathrm{bin}}(w)=4\}=9.
\]
另一方面，表 \ref{tab:fold6_bin_uplift_choice_collapse} 给出精确对应
\[
w_6=1 \Longleftrightarrow d_6^{\mathrm{bin}}(w)=2,
\]
\[
w_6=0,\ V_6(w)\le 8 \Longleftrightarrow d_6^{\mathrm{bin}}(w)=4,
\]
\[
w_6=0,\ V_6(w)>8 \Longleftrightarrow d_6^{\mathrm{bin}}(w)=3.
\]
因此三类纤维的总微态质量分别为
\[
8\cdot 2=16,\qquad 9\cdot 4=36,\qquad 4\cdot 3=12.
\]
又因为不同 \(w\) 的纤维两两不交且其并恰为 \(\{0,\dots,63\}\)，故
\[
64=16+36+12.
\]
最后两边除以 \(64\) 即得所示概率。证毕。
\end{proof}

\begin{theorem}[window-\texorpdfstring{$6$}{6} 的 \texorpdfstring{$A_2(+)$}{A2(+)} 与 \texorpdfstring{$A_2(-)$}{A2(-)} 长根扇区具有完全相同的纤维多重度谱]\label{thm:xi-window6-a2pm-identical-fiber-spectrum}
\leanverified{paper\_xi\_window6\_a2pm\_identical\_fiber\_spectrum}
沿用定理 \ref{thm:boundary-m10-b3c3-rigid-four-block-split} 的 \(B_3\) 根字典与分块记号。定义
\[
\Phi_{A_2(+)}
:=
\{\pm(e_1+e_2),\ \pm(e_1+e_3),\ \pm(e_2+e_3)\},
\]
\[
\Phi_{A_2(-)}
:=
\{\pm(e_1-e_2),\ \pm(e_1-e_3),\ \pm(e_2-e_3)\},
\]
\[
X_6^{A_2(+)}:=\iota_B^{-1}\bigl(\Phi_{A_2(+)}\bigr),
\qquad
X_6^{A_2(-)}:=\iota_B^{-1}\bigl(\Phi_{A_2(-)}\bigr).
\]
则两块的纤维多重度多重集完全一致，且都等于
\[
\{4,4,3,3,2,2\}.
\]
等价地，若定义生成多项式
\[
P_{A_2(\pm)}(t):=\sum_{w\in X_6^{A_2(\pm)}} t^{d_6^{\mathrm{bin}}(w)},
\]
则
\[
P_{A_2(+)}(t)=P_{A_2(-)}(t)=2t^4+2t^3+2t^2.
\]
\end{theorem}
\begin{proof}
由表 \ref{tab:boundary_uplift_m10_dictionary} 可直接读出
\[
X_6^{A_2(+)}
=
\{\texttt{000000},\texttt{010100},\texttt{100010},
\texttt{101010},\texttt{010001},\texttt{010101}\},
\]
\[
X_6^{A_2(-)}
=
\{\texttt{101000},\texttt{100100},\texttt{010010},
\texttt{001010},\texttt{001001},\texttt{000101}\}.
\]
另一方面，推论 \ref{cor:fold-bin-m6-fiber-hist} 与表 \ref{tab:fold6_bin_uplift_choice_collapse} 给出三档精确分类：
\[
w_6=0,\ V_6(w)\le 8 \Longleftrightarrow d_6^{\mathrm{bin}}(w)=4,
\]
\[
w_6=0,\ V_6(w)>8 \Longleftrightarrow d_6^{\mathrm{bin}}(w)=3,
\]
\[
w_6=1 \Longleftrightarrow d_6^{\mathrm{bin}}(w)=2.
\]
其中在 \(X_6^{A_2(+)}\) 中，
\[
\texttt{000000},\texttt{010100}\in\{w_6=0,\ V_6(w)\le 8\},
\]
\[
\texttt{100010},\texttt{101010}\in\{w_6=0,\ V_6(w)>8\},
\]
\[
\texttt{010001},\texttt{010101}\in\{w_6=1\}.
\]
故其多重度多重集为
\[
\{4,4,3,3,2,2\}.
\]
同理，在 \(X_6^{A_2(-)}\) 中，
\[
\texttt{101000},\texttt{100100}\in\{w_6=0,\ V_6(w)\le 8\},
\]
\[
\texttt{010010},\texttt{001010}\in\{w_6=0,\ V_6(w)>8\},
\]
\[
\texttt{001001},\texttt{000101}\in\{w_6=1\},
\]
故其多重度多重集同样为
\[
\{4,4,3,3,2,2\}.
\]
于是两块的生成多项式均为
\[
2t^4+2t^3+2t^2.
\]
证毕。
\end{proof}

\begin{corollary}[任何仅依赖纤维多重度的统计量都无法区分 \texorpdfstring{$A_2(+)$}{A2(+)} 与 \texorpdfstring{$A_2(-)$}{A2(-)}]\label{cor:xi-window6-a2pm-multiplicity-blind-statistics}
沿用定理 \ref{thm:xi-window6-a2pm-identical-fiber-spectrum} 的记号。对任意函数
\[
\Phi:\{2,3,4\}\to \RR,
\]
都有
\[
\sum_{w\in X_6^{A_2(+)}}\Phi\!\bigl(d_6^{\mathrm{bin}}(w)\bigr)
=
\sum_{w\in X_6^{A_2(-)}}\Phi\!\bigl(d_6^{\mathrm{bin}}(w)\bigr).
\]
特别地，对任意 \(q\in\RR\)，有
\[
\sum_{w\in X_6^{A_2(+)}}\bigl(d_6^{\mathrm{bin}}(w)\bigr)^q
=
\sum_{w\in X_6^{A_2(-)}}\bigl(d_6^{\mathrm{bin}}(w)\bigr)^q
=
2(2^q+3^q+4^q).
\]
\end{corollary}
\begin{proof}
由定理 \ref{thm:xi-window6-a2pm-identical-fiber-spectrum}，两块的纤维多重度多重集都等于
\[
\{4,4,3,3,2,2\}.
\]
因此对任意只依赖多重度数值的标量函数 \(\Phi\)，两侧求和都仅由同一个多重集决定，故必相等。取 \(\Phi(d)=d^q\) 即得第二式。证毕。
\end{proof}

\begin{theorem}[纤维大小场在 \texorpdfstring{$W(B_3)$}{W(B3)}-不变子空间中的最优两轨道压缩]\label{thm:xi-window6-weyl-invariant-best-two-orbit-compression}
沿用推论 \ref{cor:fold6-weyl-two-orbit-compression} 的分块记号，记
\[
d_6^{\mathrm{cyc}}:=d_6^{\mathrm{bin}}\big|_{X_6^{\mathrm{cyc}}}:X_6^{\mathrm{cyc}}\to\{2,3,4\}.
\]
在实向量空间 \(\RR^{X_6^{\mathrm{cyc}}}\) 上取均匀内积
\[
\langle f,g\rangle
:=
\frac{1}{18}\sum_{w\in X_6^{\mathrm{cyc}}}f(w)g(w),
\]
并定义 \(W(B_3)\)-不变子空间
\[
\mathcal W
:=
\{f:X_6^{\mathrm{cyc}}\to\RR:\ f\ \text{在}\ X_6^{\mathrm{short}}\ \text{与}\ X_6^{\mathrm{long}}\ \text{上分别常值}\}.
\]
则成立：
\begin{enumerate}
\item \(d_6^{\mathrm{cyc}}\) 在 \(\mathcal W\) 上的唯一正交投影 \(\Pi_{\mathcal W}d_6^{\mathrm{cyc}}\) 为
\[
(\Pi_{\mathcal W}d_6^{\mathrm{cyc}})(w)
=
\begin{cases}
\frac{11}{3},& w\in X_6^{\mathrm{short}},\\[4pt]
3,& w\in X_6^{\mathrm{long}}.
\end{cases}
\]
\item 最小均方误差恰为
\[
\inf_{f\in\mathcal W}
\frac{1}{18}\sum_{w\in X_6^{\mathrm{cyc}}}\bigl(d_6^{\mathrm{cyc}}(w)-f(w)\bigr)^2
=
\frac{17}{27}.
\]
\item 一致范数下的最佳 \(W(B_3)\)-不变逼近误差恰为
\[
\inf_{f\in\mathcal W}\|d_6^{\mathrm{cyc}}-f\|_{\infty}=1,
\]
且唯一极小元为常函数
\[
f\equiv 3.
\]
\end{enumerate}
\leanverified{paper\_xi\_window6\_weyl\_invariant\_best\_two\_orbit\_compression}
\end{theorem}
\begin{proof}
由表 \ref{tab:fold6_b3c3_root_dictionary_b3} 与表 \ref{tab:fold_bin_gauge_m6_fibers}，六个短根原像的多重度多重集为
\[
\{4,4,4,4,4,2\}.
\]
另一方面，由定理 \ref{thm:xi-window6-a2pm-identical-fiber-spectrum}，
\[
X_6^{\mathrm{long}}=X_6^{A_2(+)}\sqcup X_6^{A_2(-)}
\]
且两块各自的多重度多重集都为 \(\{4,4,3,3,2,2\}\)。故十二个长根原像的多重度多重集为
\[
\{4,4,4,4,3,3,3,3,2,2,2,2\}.
\]

由推论 \ref{cor:fold6-weyl-two-orbit-compression}，\(\mathcal W\) 正是“短根块常值、长根块常值”的二维子空间。因此其正交投影由两条轨道上的平均值给出。短根平均值为
\[
\frac{5\cdot 4+2}{6}=\frac{11}{3},
\]
长根平均值为
\[
\frac{4\cdot 4+4\cdot 3+4\cdot 2}{12}=3.
\]
这就得到第一条。

对第二条，短根块的最小平方误差贡献为
\[
\frac{1}{18}\left(
5\left(4-\frac{11}{3}\right)^2+\left(2-\frac{11}{3}\right)^2
\right)
=
\frac{1}{18}\left(5\cdot \frac{1}{9}+\frac{25}{9}\right)
=
\frac{5}{27},
\]
长根块的最小平方误差贡献为
\[
\frac{1}{18}\left(
4(4-3)^2+4(3-3)^2+4(2-3)^2
\right)
=
\frac{8}{18}
=
\frac{4}{9}.
\]
两者相加即得
\[
\frac{5}{27}+\frac{4}{9}=\frac{17}{27}.
\]

最后考察一致范数。任取 \(f\in\mathcal W\)，则存在 \(a,b\in\RR\) 使
\[
f|_{X_6^{\mathrm{short}}}\equiv a,
\qquad
f|_{X_6^{\mathrm{long}}}\equiv b.
\]
于是
\[
\sup_{w\in X_6^{\mathrm{short}}}|d_6^{\mathrm{cyc}}(w)-f(w)|
=
\max\{|2-a|,|4-a|\},
\]
其最小值为 \(1\)，且唯一在 \(a=3\) 时取得；
\[
\sup_{w\in X_6^{\mathrm{long}}}|d_6^{\mathrm{cyc}}(w)-f(w)|
=
\max\{|2-b|,|3-b|,|4-b|\},
\]
其最小值同样为 \(1\)，且唯一在 \(b=3\) 时取得。故
\[
\|d_6^{\mathrm{cyc}}-f\|_\infty
=
\max\left\{
\max\{|2-a|,|4-a|\},
\max\{|2-b|,|3-b|,|4-b|\}
\right\}
\ge 1,
\]
并且等号当且仅当 \(a=b=3\) 时成立。于是最优一致逼近误差为 \(1\)，唯一极小元为常函数 \(f\equiv 3\)。证毕。
\end{proof}

\noindent
这些结论说明：Dirichlet 扭曲侧的平方根门槛失效一旦出现，便会沿倍数链产生具有正下密度的坏模数族；而在 window-\(6\) 的 \(B_3\) 离散根字典中，虽然纤维大小场严格打破了 Weyl 轨道均匀性，长根层内部却仍保留一条精确的 \(A_2(+)/A_2(-)\) 多重度退化。于是“Weyl 破缺”与“残余二块等谱”并非矛盾，而是在不同粗细尺度上同时成立的两级结构。

\begin{theorem}[window-\texorpdfstring{$6$}{6} 的 \texorpdfstring{$B_3$}{B3} 根云上同类等值二次型的完全分类]\label{thm:xi-window6-b3-rootcloud-quadratic-isotropy-forcing}
沿用定理 \ref{thm:boundary-m10-b3c3-rigid-four-block-split} 的 \(B_3\) 根字典与定理 \ref{thm:xi-window6-a2pm-identical-fiber-spectrum} 的记号。记
\[
\Phi_{\mathrm{short}}
:=
\{\pm e_1,\pm e_2,\pm e_3\},
\]
\[
\Phi_{A_2(+)}
:=
\{\pm(e_1+e_2),\pm(e_1+e_3),\pm(e_2+e_3)\},
\]
\[
\Phi_{A_2(-)}
:=
\{\pm(e_1-e_2),\pm(e_1-e_3),\pm(e_2-e_3)\}.
\]
设
\[
Q(x)=x^{\mathsf T}Ax
\qquad
\bigl(A=(a_{ij})\in \mathrm{Sym}_3(\RR)\bigr)
\]
为 \(\RR^3\) 上实对称二次型。若存在实数 \(q_s,q_+,q_-\) 使得
\[
Q(\alpha)=q_s\quad (\alpha\in \Phi_{\mathrm{short}}),
\]
\[
Q(\beta)=q_+\quad (\beta\in \Phi_{A_2(+)}),
\]
\[
Q(\gamma)=q_-\quad (\gamma\in \Phi_{A_2(-)}),
\]
则 \(A\) 必且仅必具有形式
\[
\boxed{\;
A=
\begin{pmatrix}
a&b&b\\
b&a&b\\
b&b&a
\end{pmatrix},
\qquad
q_s=a,\quad
q_+=2a+2b,\quad
q_-=2a-2b.\;}
\]
特别地，若进一步满足
\[
q_+=q_-,
\]
则必有
\[
\boxed{\;
A=aI_3.\;}
\]
换言之，一旦 \(A_2(+)\) 与 \(A_2(-)\) 两个长根块在二次能量上不可区分，整个 \(B_3\) 根云的二次几何便被强制为欧氏各向同性。
\end{theorem}
\begin{proof}
由
\[
\Phi_{\mathrm{short}}=\{\pm e_1,\pm e_2,\pm e_3\},
\]
条件 \(Q(\pm e_i)=q_s\) 对 \(i=1,2,3\) 立刻给出
\[
a_{11}=a_{22}=a_{33}=:a=q_s.
\]

再看 \(A_2(+)\) 块。因为
\[
Q(e_1+e_2)=a_{11}+a_{22}+2a_{12}=2a+2a_{12},
\]
\[
Q(e_1+e_3)=2a+2a_{13},
\qquad
Q(e_2+e_3)=2a+2a_{23},
\]
而三者都等于同一个常数 \(q_+\)，故
\[
a_{12}=a_{13}=a_{23}=:b.
\]
于是
\[
A=
\begin{pmatrix}
a&b&b\\
b&a&b\\
b&b&a
\end{pmatrix}.
\]

对这种矩阵直接计算可得
\[
Q(e_i)=a\qquad (1\le i\le 3),
\]
\[
Q(e_i+e_j)=2a+2b\qquad (1\le i<j\le 3),
\]
\[
Q(e_i-e_j)=2a-2b\qquad (1\le i<j\le 3).
\]
因此
\[
q_s=a,\qquad q_+=2a+2b,\qquad q_-=2a-2b,
\]
从而前述形式既是必要条件，也是充分条件。

最后，若 \(q_+=q_-\)，则
\[
2a+2b=2a-2b,
\]
故 \(b=0\)，于是
\[
A=aI_3.
\]
这正是所述各向同性结论。证毕。
\end{proof}

\paragraph{黄金地址柱密度、对数级模糊质量与有限证书的双重强失明}
在 Gödel--Zeckendorf 地址极限的严格前缀等距、素数码集合 \(\mathrm{ZG}\) 的自然密度闭式，以及有限半径/有限素数证书的不完备性已经分别建立之后，还可把两条链进一步压缩为一组更锋利的结论：任意有限地址柱集都承载严格正的算术密度，而零前缀地址球中的质量衰减仅为对数级，从而把有限半径解析证书与有限 prime shadow 的联合失明定量化为一条显式下界。

\begin{theorem}[任意 admissible 前缀柱集的自然密度矩阵公式与严格正性]\label{thm:xi-time-part60f-zg-cylinder-density-matrix-positive}
\leanverified{paper\_xi\_time\_part60f\_zg\_cylinder\_density\_matrix\_positive}
设
\[
\mathrm{ZG}
:=
\Bigl\{
\prod_{k\ge 1}p_k^{z_k}:\ 
z_k\in\{0,1\},\ z_kz_{k+1}=0,\ z_k=0\ \text{对充分大 }k
\Bigr\}.
\]
对任意长度 \(n\) 的二元词
\[
\mathbf z=(z_1,\dots,z_n)\in\{0,1\}^n
\]
定义算术柱集
\[
\mathrm{ZG}[\mathbf z]
:=
\Bigl\{
m\in \mathrm{ZG}:\ v_{p_1}(m)=z_1,\dots,v_{p_n}(m)=z_n
\Bigr\}.
\]
再定义一步矩阵
\[
M_k^{(0)}
:=
\begin{pmatrix}
\dfrac{p_k}{p_k+1} & \dfrac{p_k}{p_k+1}\\[2mm]
0&0
\end{pmatrix},
\qquad
M_k^{(1)}
:=
\begin{pmatrix}
0&0\\[2mm]
\dfrac{1}{p_k+1}&0
\end{pmatrix},
\]
\[
M_k:=M_k^{(0)}+M_k^{(1)}
=
\begin{pmatrix}
\dfrac{p_k}{p_k+1} & \dfrac{p_k}{p_k+1}\\[2mm]
\dfrac{1}{p_k+1} & 0
\end{pmatrix}.
\]
记
\[
\mathbf e_0:=\begin{pmatrix}1\\0\end{pmatrix},
\qquad
\mathbf 1:=\begin{pmatrix}1\\1\end{pmatrix}.
\]
则成立：
\begin{enumerate}
  \item 若 \(\mathbf z\) 不 admissible，即存在 \(i<n\) 使 \(z_iz_{i+1}=1\)，则
  \[
  \mathrm{dens}\bigl(\mathrm{ZG}[\mathbf z]\bigr)=0.
  \]
  \item 若 \(\mathbf z\) admissible，则自然密度存在，且满足精确矩阵公式
  \[
  \mathrm{dens}\bigl(\mathrm{ZG}[\mathbf z]\bigr)
  =
  \frac{1}{\zeta(2)}
  \lim_{N\to\infty}
  \mathbf 1^\top
  M_NM_{N-1}\cdots M_{n+1}
  M_n^{(z_n)}\cdots M_1^{(z_1)}\mathbf e_0.
  \]
  \item 对每个 admissible \(\mathbf z\)，上述极限严格为正。
\end{enumerate}
\end{theorem}
\begin{proof}
对 \(N\ge n\)，定义截断柱集
\[
\mathrm{ZG}^{(N)}[\mathbf z]
:=
\Bigl\{
m\in\NN:\ 
v_{p_i}(m)=z_i\ (1\le i\le n),\
v_{p_k}(m)\in\{0,1\}\ (n<k\le N),
\]
\[
v_{p_k}(m)v_{p_{k+1}}(m)=0\ (1\le k<N)
\Bigr\}.
\]
其成员资格仅依赖于模
\[
\prod_{k=1}^{N}p_k^2
\]
的剩余类，故自然密度存在。

若 \(\mathbf z\) 不 admissible，则任一长度 \(N\) 的扩展都不 admissible，从而
\[
\mathrm{ZG}^{(N)}[\mathbf z]=\varnothing
\qquad(N\ge n),
\]
故 \(\mathrm{dens}(\mathrm{ZG}[\mathbf z])=0\)。
这也可直接从
\[
M_{i+1}^{(1)}M_i^{(1)}=0
\]
看出。

以下设 \(\mathbf z\) admissible。
对任意 admissible 扩展
\[
\mathbf x=(x_1,\dots,x_N)\in\{0,1\}^N
\]
且 \(x_i=z_i\)（\(1\le i\le n\)），恰有
\[
\mathrm{dens}\bigl\{m:\ v_{p_k}(m)=x_k\ (1\le k\le N)\bigr\}
=
\prod_{x_k=0}\left(1-\frac1{p_k}\right)
\prod_{x_k=1}\left(\frac1{p_k}-\frac1{p_k^2}\right).
\]
将两类因子重写为
\[
1-\frac1{p_k}
=
\left(1-\frac1{p_k^2}\right)\frac{p_k}{p_k+1},
\qquad
\frac1{p_k}-\frac1{p_k^2}
=
\left(1-\frac1{p_k^2}\right)\frac1{p_k+1},
\]
便得到
\[
\mathrm{dens}\bigl\{m:\ v_{p_k}(m)=x_k\ (1\le k\le N)\bigr\}
=
\left(\prod_{k=1}^{N}\left(1-\frac1{p_k^2}\right)\right)
\prod_{x_k=0}\frac{p_k}{p_k+1}
\prod_{x_k=1}\frac1{p_k+1}.
\]
对所有 admissible 扩展求和，即得
\[
\mathrm{dens}\bigl(\mathrm{ZG}^{(N)}[\mathbf z]\bigr)
=
\left(\prod_{k=1}^{N}\left(1-\frac1{p_k^2}\right)\right)
U_N(\mathbf z),
\]
其中
\[
U_N(\mathbf z)
:=
\mathbf 1^\top
M_N\cdots M_{n+1}
M_n^{(z_n)}\cdots M_1^{(z_1)}\mathbf e_0.
\]
这正是所述矩阵量。

另一方面，
\[
\mathrm{ZG}[\mathbf z]=\bigcap_{N\ge n}\mathrm{ZG}^{(N)}[\mathbf z].
\]
若 \(m\in \mathrm{ZG}^{(N)}[\mathbf z]\setminus \mathrm{ZG}[\mathbf z]\)，则尾部必发生违约事件之一：
\[
\exists\,k>N,\ p_k^2\mid m
\qquad\text{或}\qquad
\exists\,k\ge N,\ p_kp_{k+1}\mid m.
\]
故
\[
\mathrm{ZG}^{(N)}[\mathbf z]\setminus \mathrm{ZG}[\mathbf z]
\subseteq
\Bigl(\bigcup_{k>N}\{p_k^2\mid m\}\Bigr)
\cup
\Bigl(\bigcup_{k\ge N}\{p_kp_{k+1}\mid m\}\Bigr).
\]
由并集上界，
\[
0\le
\mathrm{dens}\bigl(\mathrm{ZG}^{(N)}[\mathbf z]\bigr)
-\mathrm{dens}\bigl(\mathrm{ZG}[\mathbf z]\bigr)
\le
\sum_{k>N}\frac1{p_k^2}
+
\sum_{k\ge N}\frac1{p_kp_{k+1}}.
\]
令 \(N\to\infty\)，右端趋于 \(0\)，故自然密度存在并满足
\[
\mathrm{dens}\bigl(\mathrm{ZG}[\mathbf z]\bigr)
=
\lim_{N\to\infty}
\left(\prod_{k=1}^{N}\left(1-\frac1{p_k^2}\right)\right)U_N(\mathbf z).
\]
再由
\[
\prod_{k=1}^{N}\left(1-\frac1{p_k^2}\right)\longrightarrow \frac1{\zeta(2)},
\]
便得到第二条。

最后证明严格正性。
由 admissibility，矩阵积
\[
M_n^{(z_n)}\cdots M_1^{(z_1)}\mathbf e_0
\]
必为
\[
c_{\mathbf z}\,\mathbf e_{z_n},
\qquad
c_{\mathbf z}:=
\prod_{i:\ z_i=0}\frac{p_i}{p_i+1}
\prod_{i:\ z_i=1}\frac1{p_i+1}
>0,
\]
其中 \(\mathbf e_1:=(0,1)^\top\)。

若 \(z_n=0\)，则
\[
U_N(\mathbf z)=
c_{\mathbf z}\,\mathbf 1^\top M_N\cdots M_{n+1}\mathbf e_0.
\]
此时右端正是将定理 \ref{con:xi-zg-density-closed-form-tail-bound} 的证明逐字应用到移位素数列
\[
(p_{n+1},p_{n+2},\dots)
\]
所得的尾部归一化权重，其极限严格为正。

若 \(z_n=1\)，则
\[
M_{n+1}\mathbf e_1
=
\begin{pmatrix}
\dfrac{p_{n+1}}{p_{n+1}+1}\\[2mm]
0
\end{pmatrix}
=
\frac{p_{n+1}}{p_{n+1}+1}\mathbf e_0,
\]
故
\[
U_N(\mathbf z)=
c_{\mathbf z}\frac{p_{n+1}}{p_{n+1}+1}\,
\mathbf 1^\top M_N\cdots M_{n+2}\mathbf e_0,
\]
同样化归到移位尾部的正极限情形。
因此对每个 admissible \(\mathbf z\)，极限严格为正。证毕。
\end{proof}

\begin{theorem}[零前缀柱集的对数级自然密度律]\label{thm:xi-time-part60f-zg-zero-prefix-log-density}
记
\[
\mathbf 0^{(n)}:=(0,\dots,0)\in\{0,1\}^n,
\qquad
\mathrm{ZG}[\mathbf 0^{(n)}]
=
\{m\in\mathrm{ZG}:\ p_1,\dots,p_n\nmid m\}.
\]
则当 \(n\to\infty\) 时有精确渐近
\[
\mathrm{dens}\bigl(\mathrm{ZG}[\mathbf 0^{(n)}]\bigr)
\sim
\frac{e^{-\gamma}}{\log p_n},
\]
其中 \(\gamma\) 为 Euler--Mascheroni 常数。
\end{theorem}
\begin{proof}
定义尾部约束集合
\[
\mathrm{Tail}_n
:=
\Bigl\{
m\in\NN:\ 
\forall k>n,\ v_{p_k}(m)\in\{0,1\},\
v_{p_k}(m)v_{p_{k+1}}(m)=0
\Bigr\}.
\]
则
\[
\mathrm{ZG}[\mathbf 0^{(n)}]
=
\Bigl(\bigcap_{i=1}^{n}\{p_i\nmid m\}\Bigr)\cap \mathrm{Tail}_n.
\]
对任意固定 \(N>n\)，有限截断版本的密度按前 \(n\) 个素数与尾部条件相乘，令 \(N\to\infty\) 即得
\[
\mathrm{dens}\bigl(\mathrm{ZG}[\mathbf 0^{(n)}]\bigr)
=
\left(\prod_{i=1}^{n}\left(1-\frac1{p_i}\right)\right)\tau_n,
\]
其中
\[
\tau_n:=\mathrm{dens}(\mathrm{Tail}_n).
\]

现证 \(\tau_n\to 1\)。
若 \(m\notin \mathrm{Tail}_n\)，则尾部必发生违约事件之一：
\[
\exists\,k>n,\ p_k^2\mid m
\qquad\text{或}\qquad
\exists\,k\ge n+1,\ p_kp_{k+1}\mid m.
\]
故
\[
1-\tau_n
\le
\sum_{k>n}\frac1{p_k^2}
+
\sum_{k\ge n+1}\frac1{p_kp_{k+1}}
=o(1).
\]
因此
\[
\tau_n=1+o(1).
\]
于是
\[
\mathrm{dens}\bigl(\mathrm{ZG}[\mathbf 0^{(n)}]\bigr)
=
\left(\prod_{i=1}^{n}\left(1-\frac1{p_i}\right)\right)(1+o(1)).
\]
最后由 Mertens 乘积公式
\[
\prod_{p\le x}\left(1-\frac1p\right)\sim \frac{e^{-\gamma}}{\log x}
\]
并取 \(x=p_n\)，得到
\[
\mathrm{dens}\bigl(\mathrm{ZG}[\mathbf 0^{(n)}]\bigr)
\sim
\frac{e^{-\gamma}}{\log p_n}.
\]
证毕。
\end{proof}

\begin{theorem}[自然密度诱导的黄金柱测度、全支撑与有限支撑点的零局部维数]\label{thm:xi-time-part60f-zg-cylinder-measure-local-dimension-zero}
记
\[
\Sigma_{\mathrm{ZG}}
:=
\{(z_k)_{k\ge 1}\in\{0,1\}^{\NN}:\ z_kz_{k+1}=0\ \forall k\},
\]
并对任意 admissible 前缀
\[
\mathbf z=(z_1,\dots,z_n)
\]
定义 cylinder
\[
[\mathbf z]
:=
\{x\in\Sigma_{\mathrm{ZG}}:\ x_1=z_1,\dots,x_n=z_n\}.
\]
再定义柱权
\[
\nu_{\mathrm{ZG}}([\mathbf z])
:=
\frac{\mathrm{dens}\bigl(\mathrm{ZG}[\mathbf z]\bigr)}{\delta_{\mathrm{ZG}}},
\qquad
\delta_{\mathrm{ZG}}:=\mathrm{dens}(\mathrm{ZG}).
\]
则成立：
\begin{enumerate}
  \item \(\nu_{\mathrm{ZG}}\) 唯一延拓为 \(\Sigma_{\mathrm{ZG}}\) 上的 Borel 概率测度；
  \item \(\nu_{\mathrm{ZG}}\) 具有全支撑：每个非空 cylinder 都有正质量；
  \item 对任意 \(B=2^L\) 与 \(v\in T_2(E_{\min})\setminus 2T_2(E_{\min})\)，地址映射
  \[
  \Xi_B:\Sigma_{\mathrm{ZG}}\to K_{B,v},
  \qquad
  \Xi_B((z_t)_{t\ge 1})=\sum_{t\ge 1}z_tB^{t-1}v
  \]
  的推前
  \[
  \mu_{B}:=(\Xi_B)_\#\nu_{\mathrm{ZG}}
  \]
  是 \(K_{B,v}\) 上的全支撑 Borel 概率测度；
  \item 若 \(e=(e_k)_{k\ge 1}\in\Sigma_{\mathrm{ZG}}\) 有限支撑，即仅有限多个 \(e_k\) 非零，则存在常数 \(C(e)\in(0,\infty)\) 使得
  \[
  \nu_{\mathrm{ZG}}([e_1,\dots,e_n])
  \sim
  \frac{C(e)}{\log p_n}
  \qquad(n\to\infty),
  \]
  并且在点 \(x_e:=\Xi_B(e)\) 处，推前测度的局部维数沿自然半径 \(B^{-n}\) 满足
  \[
  \lim_{n\to\infty}
  \frac{\log \mu_B\bigl(\mathrm B_{K_{B,v}}(x_e,B^{-n})\bigr)}
       {\log B^{-n}}
  =0.
  \]
\end{enumerate}
\end{theorem}
\begin{proof}
先证第 \(1\) 条。
对任意 admissible 前缀 \(\mathbf z\)，有不交并分解
\[
[\mathbf z]
=
\bigsqcup_{\substack{a\in\{0,1\}\\ \mathbf z a\ \mathrm{admissible}}}
[\mathbf z a].
\]
算术侧同样满足
\[
\mathrm{ZG}[\mathbf z]
=
\bigsqcup_{\substack{a\in\{0,1\}\\ \mathbf z a\ \mathrm{admissible}}}
\mathrm{ZG}[\mathbf z a],
\]
而各项自然密度由定理
\ref{thm:xi-time-part60f-zg-cylinder-density-matrix-positive}
存在。
故柱权满足一致的有限可加兼容关系；又
\[
\nu_{\mathrm{ZG}}(\Sigma_{\mathrm{ZG}})
=
\frac{\mathrm{dens}(\mathrm{ZG})}{\delta_{\mathrm{ZG}}}
=1.
\]
由 cylinder 代数上的 Carath\'eodory 扩张定理，\(\nu_{\mathrm{ZG}}\) 唯一延拓为 \(\Sigma_{\mathrm{ZG}}\) 上的 Borel 概率测度。

第 \(2\) 条由定理 \ref{thm:xi-time-part60f-zg-cylinder-density-matrix-positive} 的严格正性立即得到：每个非空 cylinder 对应 admissible 前缀，因而其质量严格为正。

对第 \(3\) 条，定理 \ref{thm:xi-zg-address-binary-self-affine} 已证明 \(\Xi_B\) 在 \(\Sigma_{\mathrm{ZG}}\) 上单射，其像为 \(K_{B,v}\)；连续性显然，而 \(\Sigma_{\mathrm{ZG}}\) 紧、\(K_{B,v}\subset \ZZ_2v\) 为 Hausdorff 空间，故 \(\Xi_B\) 是到其像的同胚。于是推前 \(\mu_B\) 是 \(K_{B,v}\) 上的 Borel 概率测度。又 \(\Xi_B\) 将非空 cylinder 送到 \(K_{B,v}\) 的非空相对开集，结合第 \(2\) 条即知 \(\mu_B\) 具有全支撑。

现证第 \(4\) 条。
设 \(e\) 的支撑包含于 \(\{1,\dots,m_0\}\)。
对 \(n\ge m_0\)，前缀
\[
(e_1,\dots,e_n)
\]
恰等于固定有限非零前缀后接 \(n-m_0\) 个零。故
\[
\mathrm{ZG}[e_1,\dots,e_n]
\]
可写为满足前 \(m_0\) 位固定、其后到第 \(n\) 位全为零、而尾部 \(>n\) 继续满足 \(\mathrm{ZG}\) 约束的整数集合。于是存在尾部密度 \(\tau_n=1+o(1)\)（与定理 \ref{thm:xi-time-part60f-zg-zero-prefix-log-density} 同样的并集估计）使
\[
\mathrm{dens}\bigl(\mathrm{ZG}[e_1,\dots,e_n]\bigr)
=
\left(\prod_{\substack{1\le k\le m_0\\ e_k=1}}
\left(\frac1{p_k}-\frac1{p_k^2}\right)\right)
\left(\prod_{\substack{1\le k\le n\\ e_k=0}}
\left(1-\frac1{p_k}\right)\right)\tau_n.
\]
把有限前缀部分整理为
\[
\prod_{\substack{1\le k\le m_0\\ e_k=1}}
\left(\frac1{p_k}-\frac1{p_k^2}\right)
\prod_{\substack{1\le k\le m_0\\ e_k=0}}
\left(1-\frac1{p_k}\right)
=
\left(\prod_{k=1}^{m_0}\left(1-\frac1{p_k}\right)\right)
\prod_{\substack{1\le k\le m_0\\ e_k=1}}\frac1{p_k}.
\]
故
\[
\mathrm{dens}\bigl(\mathrm{ZG}[e_1,\dots,e_n]\bigr)
=
\left(\prod_{k=1}^{n}\left(1-\frac1{p_k}\right)\right)
\left(\prod_{\substack{1\le k\le m_0\\ e_k=1}}\frac1{p_k}\right)
(1+o(1)).
\]
由 Mertens 公式，存在
\[
C(e):=
\frac{e^{-\gamma}}{\delta_{\mathrm{ZG}}}
\prod_{\substack{1\le k\le m_0\\ e_k=1}}\frac1{p_k}
\in(0,\infty)
\]
使
\[
\nu_{\mathrm{ZG}}([e_1,\dots,e_n])
=
\frac{\mathrm{dens}(\mathrm{ZG}[e_1,\dots,e_n])}{\delta_{\mathrm{ZG}}}
\sim
\frac{C(e)}{\log p_n}.
\]

最后，由定理 \ref{thm:xi-time-part9g-holographic-prefix-isometry-on-line} 的柱--球对应，
\[
\Xi_B([e_1,\dots,e_n])
=
\left(\sum_{t=1}^{n}e_tB^{t-1}v+B^n\ZZ_2v\right)\cap K_{B,v}
=
\mathrm B_{K_{B,v}}(x_e,B^{-n}),
\]
其中 \(x_e=\Xi_B(e)\)。
因此
\[
\mu_B\bigl(\mathrm B_{K_{B,v}}(x_e,B^{-n})\bigr)
=
\nu_{\mathrm{ZG}}([e_1,\dots,e_n])
\sim
\frac{C(e)}{\log p_n}.
\]
于是
\[
\frac{\log \mu_B\bigl(\mathrm B_{K_{B,v}}(x_e,B^{-n})\bigr)}{\log B^{-n}}
=
\frac{\log C(e)-\log\log p_n+o(1)}{-n\log B}.
\]
由素数定理 \(p_n\sim n\log n\) 可知 \(\log\log p_n=o(n)\)，故上式趋于 \(0\)。证毕。
\end{proof}

\begin{corollary}[有限地址分辨率的不可消解模糊质量]\label{cor:xi-time-part60f-address-resolution-logarithmic-ambiguity-mass}
沿用定理 \ref{thm:xi-time-part60f-zg-cylinder-measure-local-dimension-zero} 的记号，对每个 \(n\ge 1\) 定义
\[
U_n
:=
\{x\in K_{B,v}:\ x\equiv 0\pmod{B^n\ZZ_2v}\}.
\]
则
\[
U_n=\Xi_B([0^n]),
\]
并且
\[
\mu_B(U_n)
=
\nu_{\mathrm{ZG}}([0^n])
\sim
\frac{e^{-\gamma}}{\delta_{\mathrm{ZG}}\log p_n}.
\]
因此，任何只读取前 \(n\) 个 \(2^L\)-进地址位的有限分辨率观测，其不可消解模糊质量至多只能降到 \((\log p_n)^{-1}\) 量级，而不可能指数消失。
\end{corollary}
\begin{proof}
由定理 \ref{thm:xi-time-part9g-holographic-prefix-isometry-on-line} 的柱--球对应，
\[
\Xi_B([0^n])=(0+B^n\ZZ_2v)\cap K_{B,v}=U_n.
\]
于是
\[
\mu_B(U_n)=\nu_{\mathrm{ZG}}([0^n])
=
\frac{\mathrm{dens}(\mathrm{ZG}[0^n])}{\delta_{\mathrm{ZG}}}.
\]
再由定理 \ref{thm:xi-time-part60f-zg-zero-prefix-log-density} 即得
\[
\mu_B(U_n)\sim \frac{e^{-\gamma}}{\delta_{\mathrm{ZG}}\log p_n}.
\]
证毕。
\end{proof}

\begin{theorem}[有限半径解析证书与有限 prime shadow 的统一强失明]\label{thm:xi-time-part60f-finite-radius-primeshadow-unified-strong-blindness}
固定 \(0<\rho_\star<1\) 与 \(n\ge 1\)。
考虑对象对
\[
(F,N),
\]
其中 \(F\) 为单位圆盘上的有限 Blaschke 乘积，\(N\in \mathrm{ZG}\)。
定义联合有限证书
\[
\mathcal C_{n,\rho_\star}(F,N)
:=
\Bigl(
\mathcal J_{\rho_\star}(F),\ 
(v_{p_1}(N),\dots,v_{p_n}(N))
\Bigr),
\]
其中 \(\mathcal J_{\rho_\star}(F)\) 为定理
\ref{thm:xi-time-part60e-truncated-jensen-profile-completeness-blindness}
中的截断 Jensen 剖面。
则存在一列两两不同的对象对
\[
\{(F_j,N_j)\}_{j\ge 1}
\]
使得它们具有完全相同的联合有限证书，并满足：
\begin{enumerate}
  \item \(F_j\) 两两不同，且每一步新增零点都位于环带
  \[
  \rho_\star<|z|<1;
  \]
  \item \(N_j\in \mathrm{ZG}[0^n]\) 两两不同；
  \item 该共同算术纤维满足
  \[
  \mathrm{dens}\bigl(\mathrm{ZG}[0^n]\bigr)
  \sim
  \frac{e^{-\gamma}}{\log p_n}.
  \]
\end{enumerate}
\end{theorem}
\begin{proof}
先处理解析侧。
取一列两两不同的实数
\[
a_j\in(\rho_\star,1)
\qquad(j\ge 1),
\]
并定义 Blaschke 因子
\[
B_{a_j}(z):=\frac{z-a_j}{1-a_j z},
\qquad
F_j:=F_1\cdot \prod_{r=2}^{j}B_{a_r}(z),
\]
其中 \(F_1\) 固定任意给定的有限 Blaschke 乘积。
由于所有新增零点都满足 \(|a_j|>\rho_\star\)，定理
\ref{thm:xi-time-part60e-truncated-jensen-profile-completeness-blindness}
给出
\[
\mathcal J_{\rho_\star}(F_j)=\mathcal J_{\rho_\star}(F_1)
\qquad(j\ge 1),
\]
而 \(\{a_j\}\) 两两不同，故 \(F_j\) 两两不同。

再处理算术侧。
由定理 \ref{thm:xi-time-part60f-zg-zero-prefix-log-density}，
\[
\mathrm{dens}\bigl(\mathrm{ZG}[0^n]\bigr)>0.
\]
故 \(\mathrm{ZG}[0^n]\) 必为无限集。任选一列两两不同的元素
\[
N_j\in \mathrm{ZG}[0^n]
\qquad(j\ge 1).
\]
则对每个 \(j\)，都有
\[
v_{p_1}(N_j)=\cdots=v_{p_n}(N_j)=0,
\]
从而它们的有限 prime shadow 完全相同。

综上，对所有 \(j\ge 1\)，都有
\[
\mathcal C_{n,\rho_\star}(F_j,N_j)
=
\Bigl(\mathcal J_{\rho_\star}(F_1),(0,\dots,0)\Bigr),
\]
即得到一列联合有限证书完全相同而对象两两不同的序列。
最后，第 \(3\) 条正是定理 \ref{thm:xi-time-part60f-zg-zero-prefix-log-density} 的内容。证毕。
\end{proof}

\noindent
以上结论把地址侧与算术侧压缩为同一条可直接调用的主链：有限前缀不仅在 \(2^L\)-进地址空间中切出一个精确 cylinder，同时也在 \(\mathrm{ZG}\) 内切出一个严格正密度的算术纤维；当该前缀取为全零时，这一纤维的质量只按 \((\log p_n)^{-1}\) 衰减，而非按几何半径的幂级衰减。于是，有限半径解析证书与有限 prime shadow 即便同时给出，也仍然保留可数无限的对象纤维，并且这条失明不是纯存在性陈述，而带有明确的对数量级下界。


\paragraph{算术自然密度测度、可数 prime-register solenoid 与有限维可见性的统一塌缩}
在黄金地址柱测度的对数级质量律、无穷 prime-register solenoid 的有限可见因子化，以及联合有限证书不完备性已经建立之后，还可进一步抽取出下列五条结构后果。

\begin{theorem}[算术自然密度诱导的黄金柱测度不是任何连续势的 Gibbs 态]\label{thm:xi-time-part60f2-zg-density-measure-non-gibbs}
\leanverified{paper\_xi\_time\_part60f2\_zg\_density\_measure\_non\_gibbs}
沿用定理 \ref{thm:xi-time-part60f-zg-cylinder-measure-local-dimension-zero} 的记号，记
\[
\sigma:\Sigma_{\mathrm{ZG}}\to \Sigma_{\mathrm{ZG}},
\qquad
\sigma(z_1,z_2,z_3,\dots)=(z_2,z_3,\dots).
\]
则自然密度诱导的柱测度
\[
\nu_{\mathrm{ZG}}
\]
不是 \((\Sigma_{\mathrm{ZG}},\sigma)\) 上任何连续势
\[
\phi\in C(\Sigma_{\mathrm{ZG}})
\]
的 Gibbs 测度。
\end{theorem}
\begin{proof}
对每个 \(n\ge 1\)，记
\[
[0^n]:=[(0,\dots,0)]\subset \Sigma_{\mathrm{ZG}}.
\]
由定理 \ref{thm:xi-time-part60f-zg-zero-prefix-log-density} 与定理 \ref{thm:xi-time-part60f-zg-cylinder-measure-local-dimension-zero}，
\[
\nu_{\mathrm{ZG}}([0^n])
=
\frac{\mathrm{dens}(\mathrm{ZG}[0^n])}{\delta_{\mathrm{ZG}}}
\sim
\frac{e^{-\gamma}}{\delta_{\mathrm{ZG}}\log p_n}.
\]
特别地，
\[
\nu_{\mathrm{ZG}}([0^n])\to 0,
\]
并且该衰减仅为次指数级。

反设 \(\nu_{\mathrm{ZG}}\) 是某个连续势 \(\phi\in C(\Sigma_{\mathrm{ZG}})\) 的 Gibbs 测度。则按 Gibbs 定义，存在常数 \(C\ge 1\) 使对一切 \(n\ge 1\) 都有
\[
C^{-1}
\le
\nu_{\mathrm{ZG}}([0^n])
\exp\!\bigl(nP(\phi)-S_n\phi(0^\infty)\bigr)
\le
C,
\]
其中 \(0^\infty=(0,0,\dots)\in\Sigma_{\mathrm{ZG}}\) 为全零固定点。
由于
\[
S_n\phi(0^\infty)=n\phi(0^\infty),
\]
便得到
\[
\nu_{\mathrm{ZG}}([0^n])\asymp \exp\!\bigl(-n(P(\phi)-\phi(0^\infty))\bigr).
\]
若 \(P(\phi)<\phi(0^\infty)\)，则右端随 \(n\) 指数增长，与 \(\nu_{\mathrm{ZG}}([0^n])\le 1\) 矛盾。
若 \(P(\phi)=\phi(0^\infty)\)，则 \(\nu_{\mathrm{ZG}}([0^n])\) 有正下界，这与 \(\nu_{\mathrm{ZG}}([0^n])\to 0\) 矛盾。
若 \(P(\phi)>\phi(0^\infty)\)，则 \(\nu_{\mathrm{ZG}}([0^n])\) 必指数衰减；而
\[
\nu_{\mathrm{ZG}}([0^n])\sim \frac{e^{-\gamma}}{\delta_{\mathrm{ZG}}\log p_n}
\]
仅按 \((\log p_n)^{-1}\) 衰减。由素数定理 \(p_n\sim n\log n\) 可知
\[
\log\log p_n=o(n),
\]
故这不可能与任意指数律等价。三种情形皆矛盾，因此 \(\nu_{\mathrm{ZG}}\) 不可能是任何连续势的 Gibbs 测度。证毕。
\end{proof}

\begin{theorem}[可数 prime-register solenoid 之间连续态射的完全分类]\label{thm:xi-time-part60f2-countable-solenoid-morphism-classification}
\leanverified{paper\_xi\_time\_part60f2\_countable\_solenoid\_morphism\_classification}
设 \(S,T\subset\mathbb P\) 为可数素数集，并记
\[
G_S:=\ZZ[S^{-1}],
\qquad
G_T:=\ZZ[T^{-1}],
\qquad
\Sigma_S:=\widehat{G_S},
\qquad
\Sigma_T:=\widehat{G_T}.
\]
则存在自然同构
\[
\operatorname{Hom}_{\ZZ}(G_S,G_T)\cong
\begin{cases}
0,& S\not\subseteq T,\\[1mm]
G_T,& S\subseteq T,
\end{cases}
\]
从而对偶化得到
\[
\operatorname{Hom}_{\mathrm{cts}}(\Sigma_T,\Sigma_S)\cong
\begin{cases}
0,& S\not\subseteq T,\\[1mm]
G_T,& S\subseteq T.
\end{cases}
\]
\end{theorem}
\begin{proof}
先设 \(S\subseteq T\)。
任取 \(x\in G_T\)，定义
\[
m_x:G_S\to G_T,
\qquad
y\mapsto xy.
\]
由于 \(G_S,G_T\subset\QQ\)，且 \(y\in G_S\) 的分母素数仅来自 \(S\subseteq T\)，故 \(xy\in G_T\)，从而 \(m_x\) 为良定义的加法同态。

反过来，任取 \(\phi\in \operatorname{Hom}_{\ZZ}(G_S,G_T)\)，并记
\[
x:=\phi(1)\in G_T.
\]
对任意 \(y=a/b\in G_S\)（其中 \(a\in\ZZ\)，\(b\) 的全部素因子属于 \(S\)），有
\[
b\,\phi(y)=\phi(by)=\phi(a)=a\,\phi(1)=ax.
\]
由于 \(G_T\subset\QQ\) 无挠，故
\[
\phi(y)=\frac{a}{b}x=xy=m_x(y).
\]
因此 \(\phi=m_x\)，从而
\[
\operatorname{Hom}_{\ZZ}(G_S,G_T)\cong G_T.
\]

再设 \(S\not\subseteq T\)，取 \(p\in S\setminus T\)。
任取 \(\phi\in\operatorname{Hom}_{\ZZ}(G_S,G_T)\)，并记 \(x:=\phi(1)\)。
由于 \(p^{-n}\in G_S\) 对每个 \(n\ge 1\) 都成立，故
\[
x=\phi(1)=p^n\phi(p^{-n})\in p^nG_T
\qquad(n\ge 1).
\]
于是
\[
x\in \bigcap_{n\ge 1}p^nG_T.
\]
若 \(u\in G_T\setminus\{0\}\)，则可写成
\[
u=\frac{a}{\prod_{q\in T}q^{m_q}}
\]
其中 \(a\in\ZZ\setminus\{0\}\)，且分母与 \(p\) 互素。
若 \(u\in p^nG_T\)，则把两边写成既约分式可知 \(p^n\mid a\)。
这不可能对所有 \(n\) 同时成立，故
\[
\bigcap_{n\ge 1}p^nG_T=\{0\}.
\]
于是 \(x=0\)。
再对任意 \(y=a/b\in G_S\) 有
\[
b\,\phi(y)=\phi(a)=a\,\phi(1)=0,
\]
而 \(G_T\) 无挠，故 \(\phi(y)=0\)。
因此
\[
\operatorname{Hom}_{\ZZ}(G_S,G_T)=0.
\]

最后，Pontryagin 对偶给出反变等价
\[
\operatorname{Hom}_{\mathrm{cts}}(\Sigma_T,\Sigma_S)\cong \operatorname{Hom}_{\ZZ}(G_S,G_T),
\]
故连续态射的分类随即成立。证毕。
\end{proof}

\begin{theorem}[无穷 prime-register solenoid 的有限维连续表示只看见有限个素数轴]\label{thm:xi-time-part60f2-finite-dimensional-representation-finite-prime-factorization}
\leanverified{paper\_xi\_time\_part60f2\_finite\_dimensional\_representation\_finite\_prime\_factorization}
设 \(S\subset\mathbb P\) 为可数无穷素数集，并记
\[
\Sigma_S:=\widehat{\ZZ[S^{-1}]}.
\]
任取连续酉表示
\[
\rho:\Sigma_S\to U(d).
\]
则存在有限子集
\[
S_0\Subset S
\]
与连续酉表示
\[
\rho_0:\Sigma_{S_0}\to U(d)
\]
使得
\[
\rho=\rho_0\circ \pi_{S\to S_0}.
\]
因此余有限尾核
\[
K_{S\setminus S_0}:=\prod_{p\in S\setminus S_0}\ZZ_p
\]
满足
\[
K_{S\setminus S_0}\subseteq \ker\rho.
\]
此外，\(\rho(\Sigma_S)\) 的 Lie 秩至多为 \(1\)；更强地，
\[
\rho(\Sigma_S)
\]
只能是平凡群或一维圆群 \(\TT\)。
\end{theorem}
\begin{proof}
由于 \(\Sigma_S\) 为紧 Abel 群，\(\rho(\Sigma_S)\) 是一族彼此交换的酉矩阵，故可同时酉对角化。
于是存在 \(U\in U(d)\) 与连续特征
\[
\chi_1,\dots,\chi_d:\Sigma_S\to \TT
\]
使得对每个 \(\xi\in\Sigma_S\) 都有
\[
U\rho(\xi)U^{-1}
=
\operatorname{diag}\bigl(\chi_1(\xi),\dots,\chi_d(\xi)\bigr).
\]
由 Pontryagin 对偶，
\[
\operatorname{Hom}_{\mathrm{cts}}(\Sigma_S,\TT)\cong G_S:=\ZZ[S^{-1}],
\]
故每个 \(\chi_j\) 对应于某个 \(x_j\in G_S\)。
每个 \(x_j\) 的分母只涉及有限多个素数。
令 \(S_0\subset S\) 为所有 \(x_j\) 的分母中出现之素数的并集，则 \(S_0\) 有限，且 \(x_j\in G_{S_0}\) 对所有 \(j\) 都成立。
于是每个 \(\chi_j\) 都经由
\[
\pi_{S\to S_0}:\Sigma_S\twoheadrightarrow \Sigma_{S_0}
\]
因子化，从而整个 \(\rho\) 也经由 \(\pi_{S\to S_0}\) 因子化。

由推论 \ref{cor:xi-time-part60e-infinite-solenoid-finite-visible-factorization} 的短正合列，
\[
0\to \prod_{p\in S\setminus S_0}\ZZ_p\to \Sigma_S\xrightarrow{\ \pi_{S\to S_0}\ }\Sigma_{S_0}\to 0,
\]
立得
\[
K_{S\setminus S_0}\subseteq \ker\rho.
\]

最后分析像的 Lie 秩。
在对角化后，\(\rho\) 可视为映射
\[
\Sigma_S\to \TT^d.
\]
其对偶离散映射为
\[
\rho^\vee:\ZZ^d\to G_S,
\qquad
e_j\mapsto x_j.
\]
因此
\[
\widehat{\rho(\Sigma_S)}\cong \operatorname{im}(\rho^\vee)=:H
=
\langle x_1,\dots,x_d\rangle\subset G_S\subset\QQ.
\]
而 \(\QQ\) 的任意有限生成子群都循环：取 \(x_1,\dots,x_d\) 的一个公共分母 \(m\)，则
\[
H\subseteq \frac1m\ZZ,
\]
故 \(H\) 是 \((1/m)\ZZ\) 的有限生成子群，从而必为零群或无限循环群。
因此其 Pontryagin 对偶
\[
\rho(\Sigma_S)\cong \widehat H
\]
只能是平凡群或 \(\TT\)。
特别地，\(\rho(\Sigma_S)\) 的 Lie 秩至多为 \(1\)。证毕。
\end{proof}

\begin{corollary}[无穷 prime-register 终对象不存在忠实有限维 Lie 模型]\label{cor:xi-time-part60f2-no-faithful-finite-dimensional-lie-model}
\leanverified{paper\_xi\_time\_part60f2\_no\_faithful\_finite\_dimensional\_lie\_model}
在定理 \ref{thm:xi-time-part60f2-finite-dimensional-representation-finite-prime-factorization} 的设定下，\(\Sigma_S\) 不存在忠实的有限维连续酉表示。
更一般地，不存在到任何有限维紧 Lie 群的忠实连续商实现。
\end{corollary}
\begin{proof}
若 \(\rho:\Sigma_S\to U(d)\) 忠实，则由定理 \ref{thm:xi-time-part60f2-finite-dimensional-representation-finite-prime-factorization}，存在有限 \(S_0\Subset S\) 使
\[
K_{S\setminus S_0}=\prod_{p\in S\setminus S_0}\ZZ_p
\subseteq \ker\rho.
\]
由于 \(S\setminus S_0\) 非空，\(K_{S\setminus S_0}\) 为非平凡闭子群，从而 \(\ker\rho\neq \{0\}\)，矛盾。

若进一步存在一个到有限维紧 Lie 群的忠实连续商实现
\[
q:\Sigma_S\twoheadrightarrow G,
\qquad
\ker q=\{0\},
\]
则由 Peter--Weyl 定理，\(G\) 存在某个忠实有限维酉表示
\[
\iota:G\hookrightarrow U(d).
\]
于是
\[
\iota\circ q:\Sigma_S\to U(d)
\]
给出忠实有限维连续酉表示，与上一步矛盾。故所述 Lie 模型亦不存在。证毕。
\end{proof}

\begin{theorem}[有限半径 Jensen 证书与有限维连续表示的统一塌缩]\label{thm:xi-time-part60f2-finite-radius-finite-dimensional-unified-collapse}
\leanverified{paper\_xi\_time\_part60f2\_finite\_radius\_finite\_dimensional\_unified\_collapse}
设 \(S\subset\mathbb P\) 为可数无穷素数集，固定
\[
0<\rho_\star<1,
\]
并取任意有限维连续酉表示
\[
\varrho:\Sigma_S\to U(d).
\]
则存在一列对象对
\[
\{(F_n,\xi_n)\}_{n\ge 1},
\]
其中每个 \(F_n\) 都是单位圆盘上的有限 Blaschke 乘积，且 \(\xi_n\in \Sigma_S\)，满足
\begin{enumerate}
  \item \(\mathcal J_{\rho_\star}(F_n)\) 全部相同；
  \item \(\varrho(\xi_n)\) 全部相同；
  \item \(F_n\) 两两不同；
  \item \(\xi_n\) 两两不同。
\end{enumerate}
因此，任何把有限半径 Jensen 剖面与有限维连续算术观测并合的有限可见证书都不可能完备。
\end{theorem}
\begin{proof}
由定理 \ref{thm:xi-time-part60f2-finite-dimensional-representation-finite-prime-factorization}，存在有限子集 \(S_0\Subset S\) 与连续酉表示
\[
\varrho_0:\Sigma_{S_0}\to U(d)
\]
使
\[
\varrho=\varrho_0\circ \pi_{S\to S_0}.
\]
再由定理 \ref{thm:xi-time-part60e-finite-radius-finite-prime-certificate-incompleteness}，存在一列对象对
\[
\{(F_n,\xi_n)\}_{n\ge 1}
\]
使得
\[
\mathcal J_{\rho_\star}(F_n)
=
\mathcal J_{\rho_\star}(F_1),
\qquad
\pi_{S\to S_0}(\xi_n)
=
\pi_{S\to S_0}(\xi_1),
\]
并且 \(F_n\) 两两不同、\(\xi_n\) 两两不同。
于是对所有 \(n\ge 1\)，
\[
\varrho(\xi_n)
=
\varrho_0\!\bigl(\pi_{S\to S_0}(\xi_n)\bigr)
=
\varrho_0\!\bigl(\pi_{S\to S_0}(\xi_1)\bigr)
=
\varrho(\xi_1).
\]
这就得到所需序列。证毕。
\end{proof}

\noindent
上述五条结果把 H.161--H.162 的地址柱密度链与局部化数系的 prime-register 对偶链压缩为同一条有限可见性判别律：算术自然密度诱导的地址测度不是任何连续势的 Gibbs 平衡态；而在 prime-register 终对象一旦进入可数无穷层级之后，所有有限维连续观测都只能通过有限个素数轴，并必然湮灭一条余有限 \(p\)-进尾核。因而，有限半径解析证书与有限维连续表示即使联合，也仍然在解析外环与算术尾核两个方向上同时失明。

\paragraph{局部化有限商、周期动力学、容量原子恢复与缺陷重整化}
在 prime-register 的有限商分类、局部化 solenoid 的周期谱、bin-fold 临界容量的原子 defect law 以及 fold-gauge 常数列的前两阶 Bernoulli 模态已经建立之后，还可进一步抽取出下列六条彼此闭合的结构后果。

\begin{theorem}[局部化数系有限商的精确外部化公式]\label{thm:xi-localized-quotient-exact-sfree-externalization}
\leanverified{paper\_xi\_localized\_quotient\_exact\_sfree\_externalization}
设 \(S\subset\mathbb P\) 为有限素数集，并记
\[
G_S:=\ZZ[S^{-1}].
\]
对任意 \(n\ge 1\)，令
\[
n_S:=\prod_{p\in S}p^{v_p(n)},
\qquad
n_{S^\perp}:=\frac{n}{n_S}.
\]
则有自然同构
\[
\boxed{\;
G_S/nG_S\cong \ZZ/n_{S^\perp}\ZZ.\;}
\]
特别地，\(G_S\) 的每一个有限商都只由 \(n\) 的 \(S\)-自由部分决定。
\end{theorem}
\begin{proof}
对每个 \(p\in S\)，元素 \(p^{-1}\in G_S\) 存在，故乘法映射
\[
x\longmapsto px
\]
在 \(G_S\) 上可逆。于是对
\[
n=n_S\,n_{S^\perp}
\]
有
\[
nG_S=n_Sn_{S^\perp}G_S=n_{S^\perp}G_S.
\]
因此只需证明当
\[
\gcd\!\Bigl(n,\prod_{p\in S}p\Bigr)=1
\]
时，
\[
G_S/nG_S\cong \ZZ/n\ZZ.
\]

任取 \(x\in G_S\)，可写成
\[
x=\frac{a}{s},
\qquad
a\in\ZZ,
\]
其中 \(s\) 的全部素因子都属于 \(S\)。由于 \(\gcd(s,n)=1\)，\(s\) 在 \(\ZZ/n\ZZ\) 中可逆。定义
\[
\pi_n:G_S\longrightarrow \ZZ/n\ZZ,
\qquad
\pi_n\!\left(\frac{a}{s}\right):=a\,s^{-1}\pmod n.
\]
若 \(a/s=a'/s'\) 于 \(G_S\) 中相等，则 \(as'=a's\)，从而
\[
a\,s^{-1}\equiv a'\,s'^{-1}\pmod n,
\]
故 \(\pi_n\) 良定义。又因整数子群 \(\ZZ\subset G_S\) 已映满 \(\ZZ/n\ZZ\)，故 \(\pi_n\) 为满射。

再看其核。若
\[
\pi_n\!\left(\frac{a}{s}\right)=0,
\]
则 \(n\mid a\)。写 \(a=nb\)，便有
\[
\frac{a}{s}=n\cdot \frac{b}{s}\in nG_S.
\]
故
\[
\ker\pi_n\subset nG_S.
\]
反向包含显然成立，于是
\[
\ker\pi_n=nG_S.
\]
由第一同构定理得到
\[
G_S/nG_S\cong \ZZ/n\ZZ.
\]
最后把 \(n\) 替换为 \(n_{S^\perp}\) 即得结论。证毕。
\end{proof}

\begin{theorem}[局部化 solenoid 自覆盖的固定点公式与 Artin--Mazur \texorpdfstring{$\zeta$}{zeta} 函数]\label{thm:xi-localized-solenoid-fixedpoints-artin-mazur-series}
沿用定理 \ref{thm:xi-localized-quotient-exact-sfree-externalization} 的记号，并记
\[
\Sigma_S:=\widehat{G_S}.
\]
设整数 \(a\ge 2\) 满足
\[
\gcd\!\Bigl(a,\prod_{p\in S}p\Bigr)=1,
\]
令
\[
f_a:\Sigma_S\longrightarrow \Sigma_S
\]
为离散端“乘以 \(a\)”的 Pontryagin 对偶。则对每个 \(k\ge 1\)，有
\[
\boxed{\;
\mathrm{Fix}(f_a^k)\cong \widehat{G_S/(a^k-1)G_S},
\qquad
\#\mathrm{Fix}(f_a^k)=(a^k-1)_{S^\perp}.\;}
\]
并且其 Artin--Mazur \(\zeta\)-函数满足
\[
\boxed{\;
\zeta_{f_a}(z)
:=
\exp\!\left(
\sum_{k\ge 1}\frac{\#\mathrm{Fix}(f_a^k)}{k}z^k
\right)
=
\exp\!\left(
\sum_{k\ge 1}\frac{(a^k-1)_{S^\perp}}{k}z^k
\right).\;}
\]
特别地，当 \(S=\varnothing\) 时，
\[
\boxed{\;
\zeta_{f_a}(z)=\frac{1-z}{1-az}.\;}
\]
\end{theorem}
\begin{proof}
由对偶定义，对任意 \(x\in \Sigma_S\) 与 \(g\in G_S\)，有
\[
(f_a^k x)(g)=x(a^k g).
\]
因此
\[
x\in \mathrm{Fix}(f_a^k)
\iff
x(a^k g)=x(g)\quad(\forall g\in G_S),
\]
亦即
\[
x\bigl((a^k-1)g\bigr)=1\quad(\forall g\in G_S).
\]
这等价于 \(x\) 因子化经过商群
\[
G_S/(a^k-1)G_S.
\]
故有自然同构
\[
\mathrm{Fix}(f_a^k)\cong \widehat{G_S/(a^k-1)G_S}.
\]

再由定理 \ref{thm:xi-localized-quotient-exact-sfree-externalization}，
\[
G_S/(a^k-1)G_S\cong \ZZ/(a^k-1)_{S^\perp}\ZZ.
\]
有限循环群的 Pontryagin 对偶仍为同阶循环群，于是
\[
\#\mathrm{Fix}(f_a^k)=(a^k-1)_{S^\perp}.
\]
把该计数代入 Artin--Mazur \(\zeta\)-函数的定义，即得所示显式级数。

当 \(S=\varnothing\) 时，\((a^k-1)_{S^\perp}=a^k-1\)，故
\[
\log\zeta_{f_a}(z)
=
\sum_{k\ge 1}\frac{a^k-1}{k}z^k
=
-\log(1-az)+\log(1-z).
\]
指数化即得
\[
\zeta_{f_a}(z)=\frac{1-z}{1-az}.
\]
证毕。
\end{proof}

\begin{theorem}[bin-fold 极限容量曲线的原子测度唯一恢复]\label{thm:xi-binfold-capacity-atomic-measure-unique-recovery}
把定理 \ref{thm:xi-binfold-critical-capacity-defect-atomic-law} 的极限容量曲线写为
\[
\widehat C_\infty(s)
:=
\frac{\varphi}{\sqrt5}\min\{1,s\}
+
\frac{1}{\sqrt5}\min\{\varphi^{-1},s\}
\qquad (s>0).
\]
则存在唯一有限正测度
\[
\boxed{\;
\eta_\infty
:=
\frac{1}{\sqrt5}\delta_{\varphi^{-1}}
+
\frac{\varphi}{\sqrt5}\delta_{1}\;}
\]
使得对一切 \(s>0\) 都有
\[
\boxed{\;
\widehat C_\infty(s)=\int_{(0,\infty)}\min\{s,y\}\,\eta_\infty(\dd y).\;}
\]
并且在分布意义下，
\[
\boxed{\;
\frac{\dd}{\dd s}\widehat C_\infty(s)=\eta_\infty([s,\infty)),
\qquad
-\frac{\dd^2}{\dd s^2}\widehat C_\infty=\eta_\infty.\;}
\]
\end{theorem}
\begin{proof}
取
\[
\eta_\infty
:=
\frac{1}{\sqrt5}\delta_{\varphi^{-1}}
+
\frac{\varphi}{\sqrt5}\delta_1.
\]
则直接计算得到
\[
\int_{(0,\infty)}\min\{s,y\}\,\eta_\infty(\dd y)
=
\frac{1}{\sqrt5}\min\{s,\varphi^{-1}\}
+
\frac{\varphi}{\sqrt5}\min\{s,1\}
=
\widehat C_\infty(s).
\]
故积分表示成立。

再对各分段求导，得到
\[
\widehat C_\infty'(s)=
\begin{cases}
\dfrac{\varphi+1}{\sqrt5},& 0<s<\varphi^{-1},\\[2mm]
\dfrac{\varphi}{\sqrt5},& \varphi^{-1}<s<1,\\[2mm]
0,& s>1.
\end{cases}
\]
另一方面，
\[
\eta_\infty([s,\infty))=
\begin{cases}
\dfrac{1}{\sqrt5}+\dfrac{\varphi}{\sqrt5},& 0<s<\varphi^{-1},\\[2mm]
\dfrac{\varphi}{\sqrt5},& \varphi^{-1}<s<1,\\[2mm]
0,& s>1,
\end{cases}
\]
两者一致，故第一条分布恒等式成立。

再对上式取分布导数。由于 \(\widehat C_\infty'\) 只在
\[
s=\varphi^{-1},\qquad s=1
\]
发生跃迁，且两处跳跃大小分别为
\[
\frac{1}{\sqrt5},
\qquad
\frac{\varphi}{\sqrt5},
\]
故
\[
-\frac{\dd^2}{\dd s^2}\widehat C_\infty
=
\frac{1}{\sqrt5}\delta_{\varphi^{-1}}
+
\frac{\varphi}{\sqrt5}\delta_1
=
\eta_\infty.
\]

唯一性亦随之而来：若另一有限正测度 \(\eta\) 满足同一积分表示，则对函数
\[
s\longmapsto \int_{(0,\infty)}\min\{s,y\}\,\eta(\dd y)
\]
作二阶分布导数便得到
\[
\eta=-\widehat C_\infty''=\eta_\infty.
\]
故 \(\eta_\infty\) 唯一。证毕。
\end{proof}

\begin{theorem}[bin-fold 容量极限是稳定层均匀二原子律的尺寸偏置平均]\label{thm:xi-binfold-capacity-size-bias-uniform-law}
沿用定理 \ref{thm:xi-foldbin-uniform-layer-degeneracy-two-atom-law} 的记号，记 \(Y_\infty^{\mathrm{unif}}\) 为该定理中的极限随机变量，即
\[
\PP\!\bigl(Y_\infty^{\mathrm{unif}}=1\bigr)=\varphi^{-1},
\qquad
\PP\!\bigl(Y_\infty^{\mathrm{unif}}=\varphi^{-1}\bigr)=\varphi^{-2}.
\]
则前一定理中的原子测度满足
\[
\boxed{\;
\eta_\infty
=
\frac{\varphi^2}{\sqrt5}\,
\mathcal L\!\bigl(Y_\infty^{\mathrm{unif}}\bigr).\;}
\]
等价地，对每个 \(s>0\) 都有
\[
\boxed{\;
\widehat C_\infty(s)
=
\frac{\varphi^2}{\sqrt5}\,
\EE\!\bigl[\min\{Y_\infty^{\mathrm{unif}},s\}\bigr].\;}
\]

进一步，记 \(Y_\infty^\mu\) 为推论 \ref{cor:fold-bin-two-point-limit-law} 的 bin-fold 推前极限随机变量。则 \(Y_\infty^\mu\) 正是 \(Y_\infty^{\mathrm{unif}}\) 的尺寸偏置变换，即
\[
\boxed{\;
\PP\!\bigl(Y_\infty^\mu=y\bigr)
=
\frac{
y\,\PP\!\bigl(Y_\infty^{\mathrm{unif}}=y\bigr)
}{
\EE\!\bigl[Y_\infty^{\mathrm{unif}}\bigr]
}
\qquad
\bigl(y\in\{1,\varphi^{-1}\}\bigr).\;}
\]
特别地，
\[
\PP\!\bigl(Y_\infty^\mu=1\bigr)=\frac{\varphi}{\sqrt5},
\qquad
\PP\!\bigl(Y_\infty^\mu=\varphi^{-1}\bigr)=\frac{1}{\varphi\sqrt5}.
\]
\end{theorem}
\begin{proof}
由定理 \ref{thm:xi-foldbin-uniform-layer-degeneracy-two-atom-law}，
\[
\mathcal L\!\bigl(Y_\infty^{\mathrm{unif}}\bigr)
=
\frac{1}{\varphi}\delta_1+\frac{1}{\varphi^2}\delta_{\varphi^{-1}}.
\]
两边同乘 \(\varphi^2/\sqrt5\)，便得到
\[
\frac{\varphi^2}{\sqrt5}\,
\mathcal L\!\bigl(Y_\infty^{\mathrm{unif}}\bigr)
=
\frac{\varphi}{\sqrt5}\delta_1+\frac{1}{\sqrt5}\delta_{\varphi^{-1}}
=
\eta_\infty,
\]
其中最后一步正是定理 \ref{thm:xi-binfold-capacity-atomic-measure-unique-recovery} 中的原子测度公式。于是
\[
\widehat C_\infty(s)
=
\int_{(0,\infty)}\min\{s,y\}\,\eta_\infty(\dd y)
=
\frac{\varphi^2}{\sqrt5}
\EE\!\bigl[\min\{Y_\infty^{\mathrm{unif}},s\}\bigr].
\]

再由定理 \ref{thm:xi-foldbin-uniform-layer-degeneracy-two-atom-law} 在 \(a=1\) 时的矩公式，
\[
\EE\!\bigl[Y_\infty^{\mathrm{unif}}\bigr]
=
\frac{1}{\varphi}+\varphi^{-3}
=
\frac{\sqrt5}{\varphi^2},
\]
其中最后一步由 \(\varphi^2=\varphi+1\) 直接化简可得。故尺寸偏置概率为
\[
\frac{1\cdot \varphi^{-1}}{\sqrt5/\varphi^2}=\frac{\varphi}{\sqrt5},
\qquad
\frac{\varphi^{-1}\cdot \varphi^{-2}}{\sqrt5/\varphi^2}
=
\frac{1}{\varphi\sqrt5}.
\]
这与推论 \ref{cor:fold-bin-two-point-limit-law} 的两点极限概率完全一致，故 \(Y_\infty^\mu\) 正是 \(Y_\infty^{\mathrm{unif}}\) 的尺寸偏置变换。证毕。
\end{proof}

\begin{theorem}[bin-fold 容量极限是统一层矩谱的 Hardy--Littlewood 变换并满足 Mellin--Stieltjes 完全对偶]\label{thm:xi-binfold-capacity-hl-mellin-duality-uniform-law}
沿用定理 \ref{thm:xi-binfold-capacity-size-bias-uniform-law} 与定理 \ref{thm:xi-foldbin-uniform-layer-two-point-exponential-family} 的记号。记
\[
Y_\infty^{\mathrm{unif}}
\sim
\frac{1}{\varphi}\delta_1+\frac{1}{\varphi^2}\delta_{\varphi^{-1}},
\qquad
Z(q):=\EE\!\left[\bigl(Y_\infty^{\mathrm{unif}}\bigr)^q\right]
=
\frac{1}{\varphi}+\varphi^{-q-2}.
\]
则：
\begin{enumerate}
  \item 对每个 \(s>0\)，bin-fold 极限容量曲线满足 Hardy--Littlewood 变换公式
  \[
  \boxed{\;
  \widehat C_\infty(s)
  =
  \frac{\varphi^2}{\sqrt5}\,
  \EE\!\bigl[\min\{Y_\infty^{\mathrm{unif}},s\}\bigr].\;}
  \]
  \item 对任意 \(0<q<1\)，有精确 Mellin--Stieltjes 对偶
  \[
  \boxed{\;
  Z(q)
  =
  \frac{\sqrt5}{\varphi^2}\,
  q(1-q)
  \int_0^\infty s^{q-2}\widehat C_\infty(s)\,\dd s.\;}
  \]
\end{enumerate}
\end{theorem}
\begin{proof}
第一条正是定理 \ref{thm:xi-binfold-capacity-size-bias-uniform-law} 的第二个盒装公式。

现证第二条。先记
\[
F(s):=\EE\!\bigl[\min\{Y_\infty^{\mathrm{unif}},s\}\bigr]
\qquad (s>0).
\]
对任意非负随机变量 \(Y\) 与任意 \(0<q<1\)，都有通用恒等式
\[
\EE[Y^q]
=
q(1-q)\int_0^\infty s^{q-2}\EE[\min\{Y,s\}]\,\dd s.
\]
事实上，由 layer-cake 公式，
\[
\EE[Y^q]
=
q\int_0^\infty s^{q-1}\PP(Y\ge s)\,\dd s,
\]
而
\[
\EE[\min\{Y,s\}]
=
\int_0^s \PP(Y\ge t)\,\dd t
\]
对 \(s\) 的导数正是 \(\PP(Y\ge s)\)。对右端作一次分部积分，即得上式。

将其应用于 \(Y=Y_\infty^{\mathrm{unif}}\)，便得到
\[
Z(q)
=
q(1-q)\int_0^\infty s^{q-2}F(s)\,\dd s.
\]
再由第一条，
\[
F(s)=\frac{\sqrt5}{\varphi^2}\,\widehat C_\infty(s),
\]
故
\[
Z(q)
=
\frac{\sqrt5}{\varphi^2}\,
q(1-q)\int_0^\infty s^{q-2}\widehat C_\infty(s)\,\dd s.
\]
证毕。
\end{proof}

\begin{theorem}[fold-gauge 缺陷序列的三次重整化正规形]\label{thm:xi-fold-gauge-defect-cubic-renormalization-normal-form}
沿用定理 \ref{thm:xi-foldbin-gauge-constant-not-cfinite-first-two-zeta-limits} 的记号，令
\[
E_m:=\log C_m-\log(2\pi),
\qquad
q:=\frac{\varphi}{2}.
\]
则当 \(m\to\infty\) 时，
\[
\boxed{\;
E_{m+1}
=
qE_m+\frac{3(2+\sqrt5)}{32}\,E_m^3+O(E_m^5).\;}
\]
等价地，
\[
\boxed{\;
\lim_{m\to\infty}
\frac{E_{m+1}-qE_m}{E_m^3}
=
\frac{3(2+\sqrt5)}{32}.\;}
\]
\end{theorem}
\begin{proof}
由定理 \ref{thm:xi-foldbin-gauge-constant-not-cfinite-first-two-zeta-limits}，
\[
E_m
=
\frac{1}{3\varphi}q^m-\frac{\sqrt5}{180}q^{3m}+O(q^{5m}).
\]
记
\[
a:=\frac{1}{3\varphi},
\qquad
b:=-\frac{\sqrt5}{180}.
\]
则
\[
E_m=aq^m+bq^{3m}+O(q^{5m}).
\]
因此
\[
E_{m+1}-qE_m
=
b(q^3-q)q^{3m}+O(q^{5m}),
\]
而
\[
E_m^3
=
a^3 q^{3m}+O(q^{5m}).
\]
于是
\[
\frac{E_{m+1}-qE_m}{E_m^3}
=
\frac{b(q^3-q)}{a^3}+O(q^{2m}).
\]
直接代入
\[
a=\frac{1}{3\varphi},
\qquad
b=-\frac{\sqrt5}{180},
\qquad
q=\frac{\varphi}{2}
\]
化简可得
\[
\frac{b(q^3-q)}{a^3}
=
\frac{3(2+\sqrt5)}{32}.
\]
故所示极限成立。

最后，由 \(a>0\) 可知
\[
E_m\sim aq^m,
\]
从而 \(q^{5m}=O(E_m^5)\)。把上式改写即得
\[
E_{m+1}
=
qE_m+\frac{3(2+\sqrt5)}{32}\,E_m^3+O(E_m^5).
\]
证毕。
\end{proof}

\paragraph{布尔投影格、Fibonacci 模环分解与有限秩账本障碍}
链上内算子的固定点算术、有限分辨率可见算术的 Fibonacci 模环结构以及有限秩加法账本的同态不可能性，在现有命题链下还可继续压缩为下列七条直接后果。前四条把链上规范投影完全布尔化，写成有限 Euler 乘积，并进一步刻画全部平方自由 \(\gcd/\mathrm{lcm}\) 算术化的唯一正规形及其 prime-register 码长的最优阶；其后三条把可见算术的幂等结构推进为分解刚性与跨分辨率障碍，并以 prime-power 分辨率收束为单链相图。

\begin{theorem}[链上内算子格的布尔刚性与旗标闭式]\label{thm:xi-chain-interior-boolean-flag-closed-form}
\leanverified{paper\_xi\_chain\_interior\_boolean\_flag\_closed\_form}
沿用定理 \ref{thm:killo-chain-interior-godel-gcd-lcm} 的记号，记
\[
\mathcal I_n:=\mathrm{Int}(C_n),
\qquad
\operatorname{rk}(I):=|\mathrm{Fix}(I)|-1.
\]
则：
\begin{enumerate}
  \item 有序格 \((\mathcal I_n,\preceq)\) 与布尔格 \(B_{n-1}\) 规范同构；
  \item
  \[
  |\mathcal I_n|=2^{n-1};
  \]
  \item 秩生成多项式满足
  \[
  \sum_{I\in\mathcal I_n}q^{\operatorname{rk}(I)}=(1+q)^{n-1};
  \]
  \item 若定义
  \[
  Z_{\mathcal I_n}(k)
  :=
  \#\{I_1\preceq\cdots\preceq I_{k-1}:\ I_j\in\mathcal I_n\}
  \qquad(k\ge 1),
  \]
  则
  \[
  Z_{\mathcal I_n}(k)=k^{n-1};
  \]
  \item 对任意 \(I\preceq J\)，区间 \([I,J]\) 仍为布尔格，且
  \[
  \mu_{\mathcal I_n}(I,J)=(-1)^{|\mathrm{Fix}(J)\setminus \mathrm{Fix}(I)|}.
  \]
\end{enumerate}
\end{theorem}
\begin{proof}
由定理 \ref{thm:killo-chain-interior-godel-gcd-lcm}，映射
\[
I\longmapsto \mathrm{Fix}(I)
\]
给出 \(\mathcal I_n\) 与
\[
\mathcal F_n:=\{F\subseteq \{0,1,\dots,n-1\}:0\in F\}
\]
之间的格同构，且 meet/join 分别对应集合交与并。去掉强制出现的元素 \(0\) 之后，\(\mathcal F_n\) 与 \(\mathcal P(\{1,\dots,n-1\})\) 规范同构，故 \((\mathcal I_n,\preceq)\) 即为秩 \(n-1\) 的布尔格。这给出第 \(1\) 条与第 \(2\) 条。

在上述同构下，秩函数 \(\operatorname{rk}(I)\) 恰对应于子集大小，故
\[
\sum_{I\in\mathcal I_n}q^{\operatorname{rk}(I)}
=
\sum_{S\subseteq\{1,\dots,n-1\}}q^{|S|}
=(1+q)^{n-1},
\]
从而第 \(3\) 条成立。

再看 zeta 多项式。一个多重链
\[
I_1\preceq\cdots\preceq I_{k-1}
\]
对应于子集链
\[
S_1\subseteq \cdots\subseteq S_{k-1}\subseteq \{1,\dots,n-1\}.
\]
对每个 \(i\in\{1,\dots,n-1\}\)，只需记录它是在第几步首次进入该链；若始终不进入，则记为第 \(0\) 类。故每个 \(i\) 恰有 \(k\) 种独立选择，从而
\[
Z_{\mathcal I_n}(k)=k^{n-1}.
\]

最后，对任意 \(I\preceq J\)，对应子集满足
\[
\mathrm{Fix}(I)\subseteq \mathrm{Fix}(J).
\]
于是区间 \([I,J]\) 同构于差集 \(\mathrm{Fix}(J)\setminus \mathrm{Fix}(I)\) 上的布尔格，故 Möbius 函数等于相对秩的符号：
\[
\mu_{\mathcal I_n}(I,J)=(-1)^{|\mathrm{Fix}(J)\setminus \mathrm{Fix}(I)|}.
\]
证毕。
\end{proof}

\begin{theorem}[平方自由 Gödel 像的 Euler 因子化与 Möbius 扭曲闭式]\label{thm:xi-chain-interior-godel-euler-mobius-factorization}
\leanverified{paper\_xi\_chain\_interior\_godel\_euler\_mobius\_factorization}
沿用定理 \ref{thm:killo-chain-interior-godel-gcd-lcm} 与定理 \ref{thm:xi-chain-interior-boolean-flag-closed-form} 的记号，并固定互异素数 \(q_0,\dots,q_{n-1}\)。则：
\begin{enumerate}
  \item 平方自由像集满足
  \[
  \kappa(\mathcal I_n)
  =
  \left\{
  q_0\prod_{i\in S}q_i:\ S\subseteq\{1,\dots,n-1\}
  \right\};
  \]
  \item 对任意 \(s\in\CC\)，有限 Dirichlet 生成函数满足
  \[
  \sum_{I\in\mathcal I_n}\kappa(I)^{-s}
  =
  q_0^{-s}\prod_{i=1}^{n-1}(1+q_i^{-s});
  \]
  \item 若 \(\widehat 0\) 表示最小内算子，则 Möbius 扭曲生成函数满足
  \[
  \sum_{I\in\mathcal I_n}\mu_{\mathcal I_n}(\widehat 0,I)\,\kappa(I)^{-s}
  =
  q_0^{-s}\prod_{i=1}^{n-1}(1-q_i^{-s}).
  \]
\end{enumerate}
\end{theorem}
\begin{proof}
由定理 \ref{thm:killo-chain-interior-godel-gcd-lcm}，任意 \(I\in\mathcal I_n\) 唯一对应于一个固定点集
\[
\mathrm{Fix}(I)=\{0\}\cup S,
\qquad
S\subseteq\{1,\dots,n-1\},
\]
并且
\[
\kappa(I)=q_0\prod_{i\in S}q_i.
\]
这立刻给出第 \(1\) 条。

于是
\[
\sum_{I\in\mathcal I_n}\kappa(I)^{-s}
=
\sum_{S\subseteq\{1,\dots,n-1\}}
\left(q_0\prod_{i\in S}q_i\right)^{-s}
=
q_0^{-s}\prod_{i=1}^{n-1}(1+q_i^{-s}),
\]
得到第 \(2\) 条。

再由定理 \ref{thm:xi-chain-interior-boolean-flag-closed-form}，对满足 \(\mathrm{Fix}(I)=\{0\}\cup S\) 的内算子有
\[
\mu_{\mathcal I_n}(\widehat 0,I)=(-1)^{|S|}.
\]
故
\[
\sum_{I\in\mathcal I_n}\mu_{\mathcal I_n}(\widehat 0,I)\,\kappa(I)^{-s}
=
\sum_{S\subseteq\{1,\dots,n-1\}}
(-1)^{|S|}
\left(q_0\prod_{i\in S}q_i\right)^{-s}
=
q_0^{-s}\prod_{i=1}^{n-1}(1-q_i^{-s}),
\]
即得第 \(3\) 条。证毕。
\end{proof}

\begin{corollary}[链上内算子格的 Dirichlet--Möbius 封闭与特征多项式]\label{cor:xi-chain-interior-dirichlet-mobius-characteristic-closure}
\leanverified{paper\_xi\_chain\_interior\_dirichlet\_mobius\_characteristic\_closure}
沿用定理 \ref{thm:xi-chain-interior-boolean-flag-closed-form} 与定理 \ref{thm:xi-chain-interior-godel-euler-mobius-factorization} 的记号。记
\[
\Xi_n(s):=\sum_{I\in\mathcal I_n}\kappa(I)^{-s},
\qquad
\mathfrak M_n(s):=\sum_{I\in\mathcal I_n}\mu_{\mathcal I_n}(\widehat 0,I)\,\kappa(I)^{-s}.
\]
则有
\[
\Xi_n(s)=q_0^{-s}\prod_{i=1}^{n-1}(1+q_i^{-s}),
\qquad
\mathfrak M_n(s)=q_0^{-s}\prod_{i=1}^{n-1}(1-q_i^{-s}).
\]
更一般地，对任意 \(I\preceq J\) 都有
\[
\mu_{\mathcal I_n}(I,J)
=
(-1)^{|\mathrm{Fix}(J)\setminus \mathrm{Fix}(I)|}
=
\mu_{\mathrm{arith}}\!\left(\frac{\kappa(J)}{\kappa(I)}\right),
\]
其中 \(\mu_{\mathrm{arith}}\) 为经典算术 Möbius 函数。特别地，若定义特征多项式
\[
\chi_{\mathcal I_n}(x)
:=
\sum_{I\in\mathcal I_n}\mu_{\mathcal I_n}(\widehat 0,I)\,
x^{\,n-1-\operatorname{rk}(I)},
\]
则
\[
\chi_{\mathcal I_n}(x)=(x-1)^{n-1}.
\]
\end{corollary}
\begin{proof}
前两式正是定理 \ref{thm:xi-chain-interior-godel-euler-mobius-factorization} 的第 \(2\)、\(3\) 条。
对任意 \(I\preceq J\)，定理 \ref{thm:xi-chain-interior-boolean-flag-closed-form} 已给出
\[
\mu_{\mathcal I_n}(I,J)=(-1)^{|\mathrm{Fix}(J)\setminus \mathrm{Fix}(I)|}.
\]
另一方面，由定理 \ref{thm:killo-chain-interior-godel-gcd-lcm}，
\[
\frac{\kappa(J)}{\kappa(I)}
=
\prod_{i\in \mathrm{Fix}(J)\setminus \mathrm{Fix}(I)}q_i
\]
为平方自由整数，故其算术 Möbius 函数恰等于相同符号，得到第三式。

最后，由 \(\widehat 0\) 对应的固定点集仅为 \(\{0\}\)，对每个满足 \(\mathrm{Fix}(I)=\{0\}\cup S\) 的内算子有
\[
\mu_{\mathcal I_n}(\widehat 0,I)=(-1)^{|S|},
\qquad
\operatorname{rk}(I)=|S|.
\]
于是
\[
\chi_{\mathcal I_n}(x)
=
\sum_{S\subseteq\{1,\dots,n-1\}}(-1)^{|S|}x^{n-1-|S|}
=
(x-1)^{n-1},
\]
其中最后一步是二项式展开。证毕。
\end{proof}

\begin{theorem}[链上规范投影的平方自由 \texorpdfstring{$\gcd/\mathrm{lcm}$}{gcd/lcm} 算术化具有唯一正规形]\label{thm:xi-chain-interior-squarefree-gcd-lcm-normalform}
沿用定理 \ref{thm:killo-chain-interior-godel-gcd-lcm} 的记号，记
\[
\mathcal I_n:=\mathrm{Int}(C_n).
\]
设映射
\[
\eta:\mathcal I_n\to \NN_{\ge 1}
\]
满足：
\begin{enumerate}
  \item 对每个 \(I\in\mathcal I_n\)，\(\eta(I)\) 都是平方自由整数；
  \item \(\eta\) 是单射；
  \item 对任意 \(I,J\in\mathcal I_n\)，都有
  \[
  \eta(I\wedge J)=\gcd(\eta(I),\eta(J)),
  \qquad
  \eta(I\vee J)=\mathrm{lcm}(\eta(I),\eta(J)).
  \]
\end{enumerate}
则存在唯一确定的平方自由整数 \(b\) 与两两互素的平方自由整数
\[
r_1,\dots,r_{n-1}>1
\]
使得对每个满足 \(0\in F\subseteq\{0,\dots,n-1\}\) 的集合 \(F\)，都有
\[
\boxed{\;
\eta(I_F)=b\prod_{i\in F\setminus\{0\}}r_i,
\;}
\]
其中空积按 \(1\) 解释。特别地，任何把链上规范投影算术化为平方自由整数整除格的 \(\gcd/\mathrm{lcm}\)-语义，都只可能与定理 \ref{thm:killo-chain-interior-godel-gcd-lcm} 中的规范编码相差一个公共底因子与一组两两互素坐标标签的重命名。
\end{theorem}
\begin{proof}
记
\[
\widehat 0:=I_{\{0\}},
\qquad
A_i:=I_{\{0,i\}}
\qquad (1\le i\le n-1).
\]
由定理 \ref{thm:killo-chain-interior-godel-gcd-lcm}，\(\mathcal I_n\) 与包含 \(0\) 的子集布尔格同构，因此 \(A_1,\dots,A_{n-1}\) 正是 \(\mathcal I_n\) 的全部原子，并且对 \(i\neq j\) 有
\[
A_i\wedge A_j=\widehat 0.
\]

令
\[
b:=\eta(\widehat 0).
\]
由假设，\(b\) 是平方自由整数。又因为
\[
\gcd(\eta(A_i),\eta(A_j))
=
\eta(A_i\wedge A_j)
=
\eta(\widehat 0)
=
b
\qquad(i\neq j),
\]
可知 \(b\mid \eta(A_i)\) 对每个 \(i\) 都成立。于是可写
\[
\eta(A_i)=b\,r_i
\qquad (1\le i\le n-1).
\]
由于 \(\eta(A_i)\) 与 \(b\) 都是平方自由整数，故 \(r_i\) 也是平方自由整数，且自动满足
\[
\gcd(b,r_i)=1.
\]
若某个 \(r_i=1\)，则
\[
\eta(A_i)=b=\eta(\widehat 0),
\]
与 \(\eta\) 的单射性矛盾，因此 \(r_i>1\)。

再由
\[
\gcd(\eta(A_i),\eta(A_j))
=
\gcd(br_i,br_j)
=
b\,\gcd(r_i,r_j)
=
b
\]
得到
\[
\gcd(r_i,r_j)=1
\qquad(i\neq j),
\]
故 \(r_1,\dots,r_{n-1}\) 两两互素。

现在任取 \(F\subseteq\{0,\dots,n-1\}\) 且 \(0\in F\)。由同一定理的布尔格描述，
\[
I_F=\bigvee_{i\in F\setminus\{0\}}A_i,
\]
其中当 \(F=\{0\}\) 时右端为空 join，其值即为 \(\widehat 0\)。于是由 \(\eta\) 对 join 的保持性可得
\[
\eta(I_F)
=
\mathrm{lcm}\bigl(\eta(A_i):i\in F\setminus\{0\}\bigr).
\]
而 \(b,r_1,\dots,r_{n-1}\) 全部平方自由，且 \(b\) 与各 \(r_i\) 互素、各 \(r_i\) 彼此互素，故上式化为
\[
\eta(I_F)
=
\mathrm{lcm}\bigl(br_i:i\in F\setminus\{0\}\bigr)
=
b\prod_{i\in F\setminus\{0\}}r_i.
\]
这就证明了所述正规形。

最后，\(b=\eta(\widehat 0)\) 且 \(r_i=\eta(A_i)/b\) 由 \(\eta\) 唯一决定，故该正规形唯一。证毕。
\end{proof}

\begin{corollary}[链上 canonical arithmeticization 在 prime-register 码长意义下达到最优阶]\label{cor:xi-chain-interior-canonical-arithmeticization-optimal-order}
沿用定理 \ref{thm:xi-chain-interior-squarefree-gcd-lcm-normalform} 的记号。记 \(p_j\) 为第 \(j\) 个素数，并把 \(\{0,1\}^{n-1}\) 与 \(\mathcal I_n\) 通过对应
\[
(\varepsilon_1,\dots,\varepsilon_{n-1})
\longleftrightarrow
I_{\{0\}\cup\{i:\varepsilon_i=1\}}
\]
识别。定义规范编码
\[
E_n^{\mathrm{can}}(\varepsilon_1,\dots,\varepsilon_{n-1})
:=
p_1\prod_{i=1}^{n-1}p_{i+1}^{\varepsilon_i}.
\]
则有
\[
\boxed{\;
\max_{\varepsilon\in\{0,1\}^{n-1}}
\log E_n^{\mathrm{can}}(\varepsilon)
=
\sum_{j=1}^{n}\log p_j
=
\vartheta(p_n)
=
(1+o(1))\,n\log n.\;}
\]
另一方面，对任意满足定理 \ref{thm:xi-chain-interior-squarefree-gcd-lcm-normalform} 三条假设的算术化
\[
\eta:\mathcal I_n\to \NN_{\ge 1},
\]
都有
\[
\boxed{\;
\max_{I\in\mathcal I_n}\log \eta(I)=\Omega(n\log n).\;}
\]
因此，在保持 \(\wedge/\vee\) 与 \(\gcd/\mathrm{lcm}\) 对应的平方自由 prime-register 算术化之中，规范编码的最坏码长阶恰为
\[
\Theta(n\log n),
\]
从而达到最优阶。并且，这一下界与定理 \ref{thm:xi-godel-prime-bit-complexity-lower-bound} 中对一般 prime-valuation 可寻址 Gödel 整数化给出的 \(T\log T\) 障碍完全一致。
\end{corollary}
\begin{proof}
对规范编码，只需取全 \(1\) 向量，便有
\[
\max_{\varepsilon\in\{0,1\}^{n-1}}
\log E_n^{\mathrm{can}}(\varepsilon)
=
\sum_{j=1}^{n}\log p_j
=
\vartheta(p_n).
\]
由素数定理的标准推论
\[
\vartheta(p_n)\sim p_n\sim n\log n,
\]
即得所示上界渐近式。

再证下界。任取满足三条假设的 \(\eta\)。由定理 \ref{thm:xi-chain-interior-squarefree-gcd-lcm-normalform}，存在平方自由整数 \(b\) 与两两互素的平方自由整数 \(r_1,\dots,r_{n-1}>1\)，使
\[
\eta(I_F)=b\prod_{i\in F\setminus\{0\}}r_i.
\]
对每个 \(i\)，任选一个素因子 \(s_i\mid r_i\)。由于 \(r_i\) 两两互素，\(s_1,\dots,s_{n-1}\) 彼此不同。于是
\[
\max_{I\in\mathcal I_n}\eta(I)
\ge
\eta(I_{\{0,1,\dots,n-1\}})
=
b\prod_{i=1}^{n-1}r_i
\ge
\prod_{i=1}^{n-1}s_i.
\]
而任意 \(n-1\) 个互异素数的乘积都不小于前 \(n-1\) 个素数的乘积，故
\[
\max_{I\in\mathcal I_n}\eta(I)
\ge
\prod_{j=1}^{n-1}p_j.
\]
取对数即得
\[
\max_{I\in\mathcal I_n}\log \eta(I)
\ge
\sum_{j=1}^{n-1}\log p_j
=
\vartheta(p_{n-1})
=
\Omega(n\log n).
\]
于是规范编码给出的 \(O(n\log n)\) 上界与任何同类算术化都必须满足的 \(\Omega(n\log n)\) 下界同阶，故其最坏码长达到最优阶。最后一条只是在更一般的 prime-valuation 可寻址 Gödel 整数化语义下援引定理 \ref{thm:xi-godel-prime-bit-complexity-lower-bound} 的同型下界。证毕。
\end{proof}

\begin{theorem}[分辨率商半环的分解刚性由 \texorpdfstring{$\omega(F_{m+2})$}{omega(Fm+2)} 完全控制]\label{thm:xi-visible-arithmetic-omega-controlled-splitting}
沿用定理 \ref{thm:fold-congruence-mod-semirings} 的记号，记
\[
R_m^{\mathrm{vis}}
:=
\Omega_m/\!\equiv_m
\cong
\ZZ/(N_m\ZZ),
\qquad
N_m:=F_{m+2},
\qquad
r_m:=\omega(N_m).
\]
则：
\begin{enumerate}
  \item \(R_m^{\mathrm{vis}}\) 的幂等元总数恰为
  \[
  2^{r_m};
  \]
  \item \(R_m^{\mathrm{vis}}\) 存在非平凡直积分解，当且仅当 \(r_m\ge 2\)；
  \item \(R_m^{\mathrm{vis}}\) 直接不可分解，当且仅当 \(N_m\) 为素数幂；
  \item 若 \(N_m=p\) 为素数，则
  \[
  R_m^{\mathrm{vis}}\cong \FF_p,
  \]
  从而其幂等元仅有 \(0,1\)。
\end{enumerate}
\end{theorem}
\begin{proof}
将 \(N_m\) 分解为
\[
N_m=\prod_{j=1}^{r_m}p_j^{a_j},
\]
其中 \(p_j\) 两两不同。由中国剩余定理，
\[
R_m^{\mathrm{vis}}
\cong
\ZZ/(N_m\ZZ)
\cong
\prod_{j=1}^{r_m}\ZZ/(p_j^{a_j}\ZZ).
\]
第 \(1\) 条正是定理 \ref{thm:xi-visible-arithmetic-unit-idempotent-count} 的幂等元计数公式。

对交换环 \(A\)，非平凡直积分解
\[
A\simeq A_1\times A_2
\qquad
(A_1,A_2\neq 0)
\]
等价于存在非平凡幂等元：若分解存在，则 \(e=(1,0)\) 满足 \(e^2=e\) 且 \(e\neq 0,1\)；反之，若 \(e^2=e\) 且 \(e\neq 0,1\)，则
\[
A\cong eA\times (1-e)A.
\]
因此第 \(2\) 条等价于“存在非平凡幂等元”。由第 \(1\) 条，这恰当且仅当 \(2^{r_m}>2\)，亦即 \(r_m\ge 2\)。

第 \(3\) 条随即成立，因为
\[
r_m=1
\quad\Longleftrightarrow\quad
N_m \text{ 为素数幂}.
\]
最后，若 \(N_m=p\) 为素数，则
\[
R_m^{\mathrm{vis}}\cong \ZZ/(p\ZZ)=\FF_p,
\]
故其仅有平凡幂等元 \(0,1\)。这也与推论 \ref{cor:field-phase-fib-prime} 一致。证毕。
\end{proof}

\leanverified{paper\_xi\_visible\_arithmetic\_no\_compatible\_finite\_rank\_ledger}
\begin{theorem}[跨分辨率相容的有限秩加法账本不能同态化全部可见乘法]\label{thm:xi-visible-arithmetic-no-compatible-finite-rank-ledger}
设 \(L\) 为有限生成交换群。假设对每个 \(m\ge 1\) 都存在映射
\[
\Phi_m:\ZZ/(F_{m+2}\ZZ)\longrightarrow \ZZ\oplus L
\]
满足：
\begin{enumerate}
  \item \(\Phi_m\) 在乘法幺半群上单射；
  \item 对任意 \(\overline a,\overline b\in \ZZ/(F_{m+2}\ZZ)\)，都有
  \[
  \Phi_m(\overline a\,\overline b)=\Phi_m(\overline a)+\Phi_m(\overline b);
  \]
  \item 对每个整数 \(n\ge 1\)，存在 \(m_0(n)\) 使得当 \(m\ge m_0(n)\) 时，\(\Phi_m(\overline n)\) 与 \(m\) 无关，其中 \(\overline n\) 表示 \(n\) 在 \(\ZZ/(F_{m+2}\ZZ)\) 中的剩余类。
\end{enumerate}
则这样的映射族不存在。
\end{theorem}
\begin{proof}
由第 \(3\) 条，可对每个 \(n\ge 1\) 定义稳定值
\[
\Phi(n):=\Phi_m(\overline n)
\qquad
(m\ge m_0(n)).
\]
这给出映射
\[
\Phi:(\NN_{\ge 1},\times)\longrightarrow (\ZZ\oplus L,+).
\]

先证其单射。若 \(a\neq b\)，则由 \(F_{m+2}\to\infty\)，可取 \(m\) 足够大，使
\[
m\ge m_0(a),\qquad m\ge m_0(b),\qquad F_{m+2}>\max\{a,b\}.
\]
于是 \(\overline a\neq \overline b\) 于 \(\ZZ/(F_{m+2}\ZZ)\) 中成立。由第 \(1\) 条，
\[
\Phi(a)=\Phi_m(\overline a)\neq \Phi_m(\overline b)=\Phi(b),
\]
故 \(\Phi\) 为单射。

再证其乘法--加法同态性。对任意 \(a,b\ge 1\)，取
\[
m\ge \max\{m_0(a),m_0(b),m_0(ab)\}.
\]
则由第 \(2\) 条，
\[
\Phi(ab)
=
\Phi_m(\overline{ab})
=
\Phi_m(\overline a\,\overline b)
=
\Phi_m(\overline a)+\Phi_m(\overline b)
=
\Phi(a)+\Phi(b).
\]
因此 \(\Phi:(\NN_{\ge 1},\times)\to (\ZZ\oplus L,+)\) 是单射幺半群同态。

另一方面，由 \(L\) 有限生成，可知交换群 \(\ZZ\oplus L\) 作为加法幺半群亦为有限生成且可消去。于是定理 \ref{thm:killo-prime-freedom-non-finitizability} 排除任何从正整数乘法幺半群到 \((\ZZ\oplus L,+)\) 的单射幺半群同态。
这与上面构造出的 \(\Phi\) 矛盾。故所述映射族不存在。证毕。
\end{proof}

\leanverified{paper\_xi\_visible\_arithmetic\_primepower\_quotient\_chain}
\begin{corollary}[素数幂分辨率上的商格单链化]\label{cor:xi-visible-arithmetic-primepower-quotient-chain}
沿用定理 \ref{thm:xi-visible-arithmetic-omega-controlled-splitting} 的记号。若
\[
N_m=F_{m+2}=p^a
\]
为某个素数幂，则 \(R_m^{\mathrm{vis}}\cong \ZZ/(p^a\ZZ)\) 的全部双边理想恰为
\[
(0)=(p^a)\subset (p^{a-1})\subset \cdots \subset (p)\subset (1),
\]
并且由定理 \ref{thm:xi-fold-congruence-quotient-divisor-lattice}，\(R_m^{\mathrm{vis}}\) 的全部半环商构成长度 \(a+1\) 的线性链。特别地，此时不存在彼此横向独立的算术粗化方向。
\end{corollary}
\begin{proof}
由定理 \ref{thm:xi-visible-arithmetic-omega-controlled-splitting}，当 \(N_m=p^a\) 时，
\[
R_m^{\mathrm{vis}}\cong \ZZ/(p^a\ZZ).
\]
而 \(p^a\) 的全部正约数正好是
\[
1,p,\dots,p^a,
\]
它们在整除偏序下形成一条线性链。定理 \ref{thm:xi-fold-congruence-quotient-divisor-lattice} 已将 \(R_m^{\mathrm{vis}}\) 的商格与 \(N_m\) 的约数格规范等价起来，故全部半环商构成长度 \(a+1\) 的线性链。

等价地，\(\ZZ/(p^a\ZZ)\) 的全部理想就是
\[
(p^j)/(p^a)
\qquad
(0\le j\le a),
\]
从而得到所示理想链。证毕。
\end{proof}

\paragraph{素数寄存器幂等锥、有限素窗口半圆维与二原子信息几何的闭合描述}
在素数寄存器估值算术化、链上规范投影的布尔化、有限素数支撑乘法单子的半圆维分离以及 bin-fold 相对均匀基线的普适 \(f\)-散度塌缩已经建立之后，还可进一步抽取出下列七条直接后果。

\leanverified{paper\_xi\_prime\_register\_idempotent\_cone\_rank\_dirichlet}
\begin{theorem}[素数寄存器幂等锥的秩分层、生成多项式与 Dirichlet 闭式]\label{thm:xi-prime-register-idempotent-cone-rank-dirichlet}
沿用定理 \ref{thm:killo-finite-semigroup-prime-valuation-arithmetization} 与定理 \ref{thm:xi-prime-register-idempotent-exact-count} 的记号，并记
\[
S_n:=\left\{\prod_{i=0}^{n-1}q_i^{a_i}: a_i\in\{0,\dots,n-1\}\right\},
\]
\[
E_n:=\operatorname{Idem}(S_n),
\qquad
E_{n,k}:=\operatorname{Idem}_k(S_n),
\qquad
\operatorname{rk}(N):=|\operatorname{Im}(f_N)|.
\]
则对每个 \(1\le k\le n\) 都有
\[
|E_{n,k}|=\binom{n}{k}k^{n-k},
\qquad
|E_n|=\sum_{k=1}^{n}\binom{n}{k}k^{n-k}.
\]
其秩生成多项式满足
\[
\mathcal E_n(y):=\sum_{N\in E_n}y^{\operatorname{rk}(N)}
=
\sum_{k=1}^{n}\binom{n}{k}k^{n-k}y^k.
\]
此外，有限 Dirichlet 生成式具有显式闭式
\[
\mathfrak E_n(s):=\sum_{N\in E_n}N^{-s}
=
\sum_{\varnothing\neq F\subseteq\{0,\dots,n-1\}}
\left(\prod_{j\in F}q_j^{-sj}\right)
\prod_{i\notin F}\left(\sum_{j\in F}q_i^{-sj}\right).
\]
\end{theorem}
\begin{proof}
由推论 \ref{cor:killo-projection-idempotent-prime-register}，\(N\in E_n\) 当且仅当 \(f_N\) 为幂等映射，且此时
\[
\operatorname{Im}(f_N)=\operatorname{Fix}(f_N).
\]
因此，定理 \ref{thm:xi-prime-register-idempotent-exact-count} 已直接给出
\[
|E_{n,k}|=\binom{n}{k}k^{n-k},
\qquad
|E_n|=\sum_{k=1}^{n}\binom{n}{k}k^{n-k}.
\]
将秩 \(k\) 的计数按 \(y^k\) 加权求和，即得 \(\mathcal E_n(y)\) 的闭式。

再求 Dirichlet 生成式。固定非空子集 \(F\subseteq\{0,\dots,n-1\}\)。由
\[
\operatorname{Im}(f_N)=\operatorname{Fix}(f_N)=F
\]
可知，此类幂等元恰由下列独立条件刻画：
\[
v_{q_j}(N)=j\qquad (j\in F),
\qquad
v_{q_i}(N)\in F\qquad (i\notin F).
\]
因此，对 \(j\in F\) 必有贡献因子 \(q_j^{-sj}\)；而对 \(i\notin F\)，其指数可在 \(F\) 中任取，故贡献
\[
\sum_{j\in F}q_i^{-sj}.
\]
各坐标彼此独立，相乘后即得到固定像集 \(F\) 的全部贡献
\[
\left(\prod_{j\in F}q_j^{-sj}\right)
\prod_{i\notin F}\left(\sum_{j\in F}q_i^{-sj}\right).
\]
再对全部非空 \(F\) 求和，即得所示 Dirichlet 闭式。证毕。
\end{proof}

\begin{theorem}[单调幂等骨架的 Boolean 交格与指数稀疏性]\label{thm:xi-monotone-idempotent-skeleton-gcd-lcm-sparsity}
沿用定理 \ref{thm:killo-chain-interior-godel-gcd-lcm} 的记号，记
\[
\mathcal I_n:=\mathrm{Int}(C_n),
\qquad
\mathcal I_{n,k}:=\{I\in\mathcal I_n:\ |\mathrm{Fix}(I)|=k\}
\qquad(1\le k\le n).
\]
则对任意满足 \(0\in F\cap G\subseteq\{0,\dots,n-1\}\) 的子集 \(F,G\)，有
\[
I_F\wedge I_G=I_{F\cap G},
\qquad
I_F\vee I_G=I_{F\cup G},
\]
从而
\[
\kappa(I_F\wedge I_G)=\gcd(\kappa(I_F),\kappa(I_G)),
\qquad
\kappa(I_F\vee I_G)=\mathrm{lcm}(\kappa(I_F),\kappa(I_G)).
\]
特别地，\(\mathcal I_n\) 在 meet 运算下构成交换幂等半群，并经 \(\kappa\) 与平方自由整数上的 \(\gcd\)-半群同构。

此外，
\[
|\mathcal I_{n,k}|=\binom{n-1}{k-1},
\]
并且其在全部秩 \(k\) 幂等元中所占比例满足
\[
\frac{|\mathcal I_{n,k}|}{|E_{n,k}|}
=
\frac{\binom{n-1}{k-1}}{\binom{n}{k}k^{n-k}}
=
\frac{k}{n}\,k^{-(n-k)}.
\]
因而只要 \(1\le k\le n-1\)，该比例便随 \(n-k\) 指数衰减。
\leanverified{paper\_xi\_monotone\_idempotent\_skeleton\_gcd\_lcm\_sparsity}
\end{theorem}
\begin{proof}
由定理 \ref{thm:killo-chain-interior-godel-gcd-lcm}，\(\mathcal I_n\) 与
\[
\mathcal F_n:=\{F\subseteq[n]:0\in F\}
\]
通过 \(F\mapsto I_F\) 一一对应，且格运算严格对应于集合交并。因而
\[
I_F\wedge I_G=I_{F\cap G},
\qquad
I_F\vee I_G=I_{F\cup G}.
\]
再由同一定理的算术化公式，
\[
\kappa(I_F\wedge I_G)=\gcd(\kappa(I_F),\kappa(I_G)),
\qquad
\kappa(I_F\vee I_G)=\mathrm{lcm}(\kappa(I_F),\kappa(I_G)),
\]
即得第一部分。特别地，meet 的交换律与幂等律使 \(\mathcal I_n\) 成为交换幂等半群。

对计数部分，只需从 \(\{1,\dots,n-1\}\) 中选出 \(k-1\) 个元素并与强制出现的 \(0\) 合并为固定点集 \(F\)，故
\[
|\mathcal I_{n,k}|=\binom{n-1}{k-1}.
\]
再与定理 \ref{thm:xi-prime-register-idempotent-cone-rank-dirichlet} 中
\[
|E_{n,k}|=\binom{n}{k}k^{n-k}
\]
相除，便得到所示比例公式。最后一条由 \(k\ge 1\) 时 \(k^{-(n-k)}\) 的指数衰减立即成立。证毕。
\end{proof}

\begin{corollary}[规范投影骨架在全部素数寄存器幂等元中的超指数稀薄性]\label{cor:xi-monotone-idempotent-skeleton-superexponential-sparsity}
沿用定理 \ref{thm:xi-chain-interior-boolean-flag-closed-form} 与定理 \ref{thm:xi-prime-register-idempotent-cone-rank-dirichlet} 的记号。对每个 \(n\ge 2\)，记
\[
c_n:=|\mathcal I_n|=2^{n-1},
\qquad
e_n:=|E_n|=\sum_{k=1}^{n}\binom{n}{k}k^{n-k}.
\]
则
\[
\frac{c_n}{e_n}
\le
\frac{n+1}{2\,\lfloor n/2\rfloor^{\,n-\lfloor n/2\rfloor}},
\]
并且
\[
\frac{c_n}{e_n}
=
\exp\!\Bigl(-\frac12 n\log n+O(n)\Bigr)
\longrightarrow 0
\qquad (n\to\infty).
\]
因而链上规范投影在全部素数寄存器幂等元中仅占超指数稀薄的比例。
\end{corollary}
\begin{proof}
记
\[
m:=\lfloor n/2\rfloor.
\]
由定理 \ref{thm:xi-prime-register-idempotent-cone-rank-dirichlet}，
\[
e_n=\sum_{k=1}^{n}\binom{n}{k}k^{n-k}
\ge
\binom{n}{m}m^{n-m}.
\]
另一方面，由
\[
2^n=\sum_{k=0}^{n}\binom{n}{k}
\]
可知中心二项式系数满足
\[
\binom{n}{m}\ge \frac{2^n}{n+1}.
\]
再结合
\[
c_n=2^{n-1},
\]
得到
\[
\frac{c_n}{e_n}
\le
\frac{2^{n-1}}{(2^n/(n+1))\,m^{n-m}}
=
\frac{n+1}{2\,m^{n-m}},
\]
即得第一条不等式。

由于 \(m=\lfloor n/2\rfloor\sim n/2\)，于是
\[
\log\frac{c_n}{e_n}
\le
\log(n+1)-\log 2-(n-m)\log m
=
-\frac12 n\log n+O(n).
\]
指数化后即得
\[
\frac{c_n}{e_n}
=
\exp\!\Bigl(-\frac12 n\log n+O(n)\Bigr)\to 0.
\]
证毕。
\end{proof}

\begin{theorem}[有限素窗口同态半圆维的严格线性差距]\label{thm:xi-finite-prime-window-halfcircle-linear-gap}
取互异素数 \(p_1,\dots,p_r\)，并记
\[
M_r:=\langle p_1,\dots,p_r\rangle
=
\left\{p_1^{a_1}\cdots p_r^{a_r}: a_1,\dots,a_r\in\ZZ_{\ge 0}\right\}.
\]
则
\[
\cdim^{\mathrm{hom}}_{+}(M_r)=\frac{r}{2},
\qquad
\cdim^{\mathrm{impl}}(M_r)=\frac12.
\]
特别地，
\[
\cdim^{\mathrm{hom}}_{+}(M_r)-\cdim^{\mathrm{impl}}(M_r)=\frac{r-1}{2},
\]
故有限素窗口上的结构层与实现层半圆维差距已经呈严格线性增长。
\leanverified{paper\_xi\_finite\_prime\_window\_halfcircle\_linear\_gap}
\end{theorem}
\begin{proof}
令
\[
S_r:=\{p_1,\dots,p_r\}\subset\mathbb P.
\]
则
\[
M_r=\NN_{S_r}^\times.
\]
所述等式正是定理 \ref{thm:xi-finite-prime-support-multiplicative-half-circle-dimension} 在 \(S=S_r\) 时的直接转写。证毕。
\end{proof}

\begin{theorem}[单一 \texorpdfstring{$\chi^2$}{chi2} 极限常数完备编码全部 Csisz\'ar 几何]\label{thm:xi-chi2-complete-csiszar-geometry}
沿用命题 \ref{prop:fold-bin-f-divergence-collapse} 的记号，并记
\[
f_2(t):=(t-1)^2,
\qquad
\mathcal C_2:=D_{f_2}^\infty=\lim_{m\to\infty}\chi^2(\mu_m\Vert u_m).
\]
则
\[
\varphi=\bigl(5(\mathcal C_2+1)-1\bigr)^{1/3}.
\]
若进一步置
\[
\xi:=\bigl(5(\mathcal C_2+1)-1\bigr)^{1/3},
\]
则对任意凸函数 \(f:(0,\infty)\to\RR\) 满足 \(f(1)=0\)，其 bin-fold 极限 \(f\)-散度都满足
\[
D_f^\infty
=
\frac{1}{\xi}\,f\!\left(\frac{\xi^2}{\sqrt5}\right)
+\frac{1}{\xi^2}\,f\!\left(\frac{\xi}{\sqrt5}\right).
\]
因此，bin-fold 相对均匀基线的全部渐近 Csisz\'ar 信息几何并非无限维数据，而由单一标量 \(\mathcal C_2\) 完整决定。
\end{theorem}
\begin{proof}
由命题 \ref{prop:fold-bin-f-divergence-collapse} 取 \(f=f_2\)，得到
\[
\mathcal C_2
=
\frac{1}{\varphi}\left(\frac{\varphi^2}{\sqrt5}\right)^2
+\frac{1}{\varphi^2}\left(\frac{\varphi}{\sqrt5}\right)^2
-1
=
\frac{\varphi^3+1}{5}-1.
\]
整理即得
\[
\varphi^3=5(\mathcal C_2+1)-1,
\]
从而
\[
\varphi=\bigl(5(\mathcal C_2+1)-1\bigr)^{1/3}.
\]
将
\[
\xi=\bigl(5(\mathcal C_2+1)-1\bigr)^{1/3}=\varphi
\]
代回命题 \ref{prop:fold-bin-f-divergence-collapse} 的两点闭式，便得到任意 \(f\) 的极限表达式。故全部 \(D_f^\infty\) 都由 \(\mathcal C_2\) 唯一决定。证毕。
\end{proof}

\begin{theorem}[渐近 \texorpdfstring{$f$}{f}-散度泛函的二原子刚性]\label{thm:xi-asymptotic-f-divergence-functional-two-atom-rigidity}
记
\[
\nu_\varphi
:=
\frac{1}{\varphi}\,\delta_{\varphi^2/\sqrt5}
+\frac{1}{\varphi^2}\,\delta_{\varphi/\sqrt5}.
\]
设 \(K\subset(0,\infty)\) 为包含
\(
\{\varphi/\sqrt5,\varphi^2/\sqrt5\}
\)
的紧区间，\(\nu\) 为 \(K\) 上的 Borel 概率测度。若对一切 \(\alpha>1\) 都有
\[
\int_K \bigl(t^\alpha-1-\alpha(t-1)\bigr)\,\dd\nu(t)
=
D_{f_\alpha}^\infty
=
\int_K \bigl(t^\alpha-1-\alpha(t-1)\bigr)\,\dd\nu_\varphi(t),
\]
其中
\[
f_\alpha(t):=t^\alpha-1-\alpha(t-1),
\]
则必有
\[
\nu=\nu_\varphi.
\]
换言之，bin-fold 的渐近 \(f\)-散度泛函在测度论层面唯一对应于该二原子模型。
\end{theorem}
\begin{proof}
设
\[
\sigma:=\nu-\nu_\varphi.
\]
对 \(z\in\CC\) 定义
\[
G(z):=\int_K \bigl(t^z-1-z(t-1)\bigr)\,\dd\sigma(t),
\qquad
t^z:=e^{z\log t}.
\]
由于 \(K\subset(0,\infty)\) 为紧集，\(\log t\) 在 \(K\) 上有界，故 \(G\) 是整函数。由假设，对每个实数 \(\alpha>1\) 都有
\[
G(\alpha)=0.
\]
零点集含有聚点，因此由恒等定理可知
\[
G\equiv 0 \qquad \text{于 } \CC.
\]
于是对一切 \(z\in\CC\) 都有
\[
0=G''(z)=\int_K (\log t)^2 t^z\,\dd\sigma(t).
\]
定义有限符号测度
\[
\dd\lambda(t):=(\log t)^2 t^2\,\dd\sigma(t).
\]
则对任意整数 \(n\ge 0\)，
\[
\int_K t^n\,\dd\lambda(t)=G''(n+2)=0.
\]
因此 \(\lambda\) 对 \(K\) 上所有多项式的积分皆为零。由 Stone--Weierstrass 定理，多项式在 \(C(K)\) 中稠密，故
\[
\lambda=0.
\]
而对 \(t\neq 1\) 有
\[
(\log t)^2 t^2>0,
\]
故 \(\sigma\) 只能支撑在 \(\{1\}\) 上。另一方面，\(\nu\) 与 \(\nu_\varphi\) 都是概率测度，故
\[
\sigma(K)=0.
\]
因此 \(\sigma\) 不可能在 \(\{1\}\) 上留下任何非零质量，只能有
\[
\sigma=0.
\]
于是 \(\nu=\nu_\varphi\)。证毕。
\end{proof}

\begin{proposition}[bin-fold 渐近 \texorpdfstring{$f$}{f}-散度塌缩为唯一二点实验]\label{prop:xi-foldbin-f-divergence-two-point-experiment}
沿用命题 \ref{prop:fold-bin-f-divergence-collapse} 的记号。定义 \(\{0,1\}\) 上的两点实验
\[
Q_\infty:=\left(\frac1\varphi,\frac1{\varphi^2}\right),
\qquad
P_\infty:=\left(\frac{\varphi}{\sqrt5},\frac{1}{\varphi\sqrt5}\right).
\]
则对任意凸函数 \(f:(0,\infty)\to\RR\) 满足 \(f(1)=0\)，都有
\[
\boxed{\;
\lim_{m\to\infty}D_f(\mu_m\Vert u_m)
=
D_f(P_\infty\Vert Q_\infty),\;}
\]
并且
\[
\boxed{\;
D_f(\mu_m\Vert u_m)
=
D_f(P_\infty\Vert Q_\infty)
+O_f\bigl((\varphi/2)^m\bigr).\;}
\]
其中
\[
D_f(P_\infty\Vert Q_\infty)
:=
\frac1\varphi\,f\!\left(\frac{P_{\infty,0}}{Q_{\infty,0}}\right)
+\frac1{\varphi^2}\,f\!\left(\frac{P_{\infty,1}}{Q_{\infty,1}}\right).
\]
换言之，\((\mu_m,u_m)\) 的全部渐近 Csisz\'ar 几何完全塌缩到单一二点实验 \((P_\infty,Q_\infty)\)。
\end{proposition}
\begin{proof}
直接计算得到
\[
\frac{P_{\infty,0}}{Q_{\infty,0}}
=
\frac{\varphi/\sqrt5}{1/\varphi}
=
\frac{\varphi^2}{\sqrt5},
\qquad
\frac{P_{\infty,1}}{Q_{\infty,1}}
=
\frac{1/(\varphi\sqrt5)}{1/\varphi^2}
=
\frac{\varphi}{\sqrt5}.
\]
因此
\[
D_f(P_\infty\Vert Q_\infty)
=
\frac1\varphi\,f\!\left(\frac{\varphi^2}{\sqrt5}\right)
+\frac1{\varphi^2}\,f\!\left(\frac{\varphi}{\sqrt5}\right).
\]
而命题 \ref{prop:fold-bin-f-divergence-collapse} 已给出
\[
D_f(\mu_m\Vert u_m)
=
\frac1\varphi\,f\!\left(\frac{\varphi^2}{\sqrt5}\right)
+\frac1{\varphi^2}\,f\!\left(\frac{\varphi}{\sqrt5}\right)
+O_f\bigl((\varphi/2)^m\bigr).
\]
将两式比较，即得极限恒等式与误差阶。证毕。
\end{proof}

\paragraph{素数寄存器骨架与二原子信息几何的双侧闭合}
上述七条结果表明，素数寄存器侧的幂等结构已经可由秩分层、Dirichlet 级数、布尔骨架的秩层指数稀疏性以及总量级的超指数稀薄性完全控制；与此同时，bin-fold 相对均匀基线的渐近信息几何则最终塌缩为由单一 \(\chi^2\) 常数编码的唯一二原子测度。两侧共同给出一条算术外置与信息塌缩相互对照的闭合描述。

\paragraph{链型幂等元的自指合成与可比链上的 \texorpdfstring{$\gcd$}{gcd} 阴影}
在素数寄存器幂等锥的秩分层、链上规范投影的布尔骨架与超指数稀薄性已经明确之后，还可进一步把链型幂等元在自指合成下的闭合规律压缩为下列两条后果。

\begin{theorem}[链型幂等元的自指合成饱和律与可比链上的 \texorpdfstring{$\gcd$}{gcd} 坍缩]\label{thm:xi-chain-idempotent-star-saturation-comparable-gcd}
\leanverified{paper\_xi\_chain\_idempotent\_star\_saturation\_comparable\_gcd}
沿用定理 \ref{thm:killo-finite-semigroup-prime-valuation-arithmetization}、
\ref{thm:killo-chain-interior-godel-gcd-lcm} 与
\ref{thm:killo-prime-register-chain-idempotent-boolean-meet} 的记号。对每个满足
\[
0\in F\subseteq [n]
\]
的子集，记
\[
N_F:=\Phi^{-1}(I_F)\in E_n^{\mathrm{ch}},
\qquad
\kappa(N_F):=\kappa(I_F)=\prod_{i\in F}q_i.
\]
并定义
\[
F\triangleright G:=I_F(G)=\{I_F(g):g\in G\}.
\]
则对任意 \(0\in F,G\subseteq[n]\)，都有
\[
\boxed{\;
N_F\star N_G=N_{F\triangleright G},
\qquad
\kappa(N_F\star N_G)=\prod_{j\in F\triangleright G}q_j.\;}
\]
特别地，若 \(F\) 与 \(G\) 在包含关系下可比，则
\[
\boxed{\;
F\triangleright G=F\cap G,\qquad
N_F\star N_G=N_G\star N_F=N_{F\cap G},\qquad
\kappa(N_F\star N_G)=\gcd(\kappa(N_F),\kappa(N_G)).\;}
\]
因此，在任意包含链所生成的链型幂等子族上，自指合成 \(\star\) 本身便塌缩为交换幂等的 \(\gcd\)-带。
\end{theorem}
\begin{proof}
由定理 \ref{thm:killo-finite-semigroup-prime-valuation-arithmetization}，
\[
f_{N_F\star N_G}=f_{N_F}\circ f_{N_G}=I_F\circ I_G.
\]
因此只需确定 \(I_F\circ I_G\) 的固定点集。记
\[
H:=F\triangleright G=I_F(G).
\]
因 \(0\in G\) 且 \(I_F(0)=0\)，故 \(0\in H\)。

任取 \(i\in[n]\)，置
\[
h:=(I_F\circ I_G)(i)=I_F(I_G(i)).
\]
由于 \(I_G(i)\in G\) 且 \(h\le I_G(i)\le i\)，可知
\[
h\in H\cap\{0,\dots,i\}.
\]
反过来，任取
\[
x\in H\cap\{0,\dots,i\},
\]
则存在 \(g_x\in G\) 使
\[
x=I_F(g_x).
\]
由于 \(x\in F\) 且 \(x\le i\)，由 \(I_G\) 的单调性有
\[
I_G(x)\le I_G(i).
\]
于是
\[
x=I_F(x)=I_F(I_G(x))\le I_F(I_G(i))=h.
\]
故 \(h\) 正是 \(H\cap\{0,\dots,i\}\) 的最大元，也即
\[
(I_F\circ I_G)(i)=I_H(i).
\]
因此
\[
I_F\circ I_G=I_{F\triangleright G}.
\]
再送回 \(S_n\) 侧即得
\[
N_F\star N_G=N_{F\triangleright G}.
\]
由 \(\kappa\) 的定义，
\[
\kappa(N_F\star N_G)=\kappa(I_{F\triangleright G})=\prod_{j\in F\triangleright G}q_j.
\]

若 \(F\subseteq G\)，则对任意 \(g\in G\) 有 \(I_F(g)\in F\)，而对每个 \(f\in F\) 取 \(g=f\in G\) 又得 \(f\in I_F(G)\)，故
\[
F\triangleright G=F=F\cap G.
\]
若 \(G\subseteq F\)，则每个 \(g\in G\) 都满足 \(I_F(g)=g\)，故
\[
F\triangleright G=G=F\cap G.
\]
由此可比情形下
\[
N_F\star N_G=N_{F\cap G}.
\]
交换性来自 \(F\cap G=G\cap F\)，而
\[
\kappa(N_{F\cap G})=\prod_{j\in F\cap G}q_j=\gcd(\kappa(N_F),\kappa(N_G))
\]
则由定理 \ref{thm:killo-chain-interior-godel-gcd-lcm} 的平方自由算术化立得。证毕。
\end{proof}

\leanverified{paper\_xi\_chain\_idempotent\_comparable\_finite\_gcd\_collapse}
\begin{corollary}[可比链式投影复合的一步交换化与像集闭式]\label{cor:xi-chain-idempotent-comparable-finite-gcd-collapse}
沿用定理 \ref{thm:xi-chain-idempotent-star-saturation-comparable-gcd} 的记号。设
\[
0\in F_1,\dots,F_r\subseteq[n]
\qquad (r\ge 2)
\]
两两在包含关系下可比。则有
\[
\boxed{\;
N_{F_1}\star\cdots\star N_{F_r}
=
N_{\bigcap_{j=1}^{r}F_j},\;}
\]
从而
\[
\boxed{\;
\kappa(N_{F_1}\star\cdots\star N_{F_r})
=
\gcd\bigl(\kappa(N_{F_1}),\dots,\kappa(N_{F_r})\bigr).\;}
\]
特别地，该复合与排列顺序无关，并且
\[
\boxed{\;
\operatorname{Im}\!\bigl(f_{N_{F_1}\star\cdots\star N_{F_r}}\bigr)
=
\operatorname{Fix}\!\bigl(f_{N_{F_1}\star\cdots\star N_{F_r}}\bigr)
=
\bigcap_{j=1}^{r}F_j.\;}
\]
\end{corollary}
\begin{proof}
由于 \(F_1,\dots,F_r\) 两两可比，可按包含关系重排为
\[
F_{\sigma(1)}\supseteq F_{\sigma(2)}\supseteq \cdots \supseteq F_{\sigma(r)}.
\]
由定理 \ref{thm:xi-chain-idempotent-star-saturation-comparable-gcd} 的可比情形，
\[
N_{F_{\sigma(1)}}\star N_{F_{\sigma(2)}}=N_{F_{\sigma(2)}}.
\]
继续迭代便得
\[
N_{F_{\sigma(1)}}\star\cdots\star N_{F_{\sigma(r)}}=N_{F_{\sigma(r)}}
=
N_{\bigcap_{j=1}^{r}F_j}.
\]
由于右端仅依赖于最小元 \(\bigcap_j F_j\)，故复合结果与排列顺序无关。

再由定理 \ref{thm:xi-chain-idempotent-star-saturation-comparable-gcd}，
\[
\kappa(N_{\bigcap_j F_j})
=
\prod_{i\in \cap_j F_j}q_i
=
\gcd\bigl(\kappa(N_{F_1}),\dots,\kappa(N_{F_r})\bigr),
\]
得第二式。

最后，由 \(f_{N_H}=I_H\) 以及推论 \ref{cor:killo-projection-idempotent-prime-register}，
\[
\operatorname{Im}(f_{N_H})=\operatorname{Fix}(f_{N_H})=H
\]
对任意 \(H\subseteq[n]\) 且 \(0\in H\) 成立。取
\[
H=\bigcap_{j=1}^{r}F_j
\]
即得像集与不动点集的闭式。证毕。
\end{proof}

\noindent 结合定理 \ref{thm:xi-prime-register-minimal-ideal-left-zero-band} 的显式最小理想与推论 \ref{cor:xi-monotone-idempotent-skeleton-superexponential-sparsity} 的超指数稀薄性可见，素数寄存器幂等宇宙内部同时存在两类彼此对照的刚性对象：一类是由 \(\{Q_n^c\}_{c\in[n]}\) 构成的左零吸收核，另一类是由可比链型幂等元所诱导的交换 \(\gcd\)-带阴影；后者虽然在包含链上完全交换化，但在全部幂等元中仍只占超指数稀薄的比例。

\paragraph{素数寄存器有限变换半群的自同构刚性、幂等分层与 Green 秩几何}
在素数寄存器的有限变换半群算术化、幂等锥的秩分层、链上内算子格的布尔化以及秩 \(1\) 左零层已经建立之后，还可把这条有限半群链进一步压缩为四条严格后果：全部半群自同构都由同一个坐标置换共轭给出，幂等元按固定指数集获得完全分层，链上内算子格在 Möbius 函数、特征多项式与极大链层面保持纯 Boolean 结构，而全部双边理想与 Green 的 \(\mathcal J\)-类则仅由秩决定。

\begin{theorem}[素数寄存器变换半群的自同构群完全刚性]\label{thm:xi-prime-register-semigroup-automorphism-rigidity}
沿用定理 \ref{thm:killo-finite-semigroup-prime-valuation-arithmetization} 的记号，并记
\[
\mathfrak S_n:=\operatorname{Sym}([n]).
\]
则半群自同构群满足
\[
\boxed{\;
\operatorname{Aut}(S_n,\star)\cong \mathfrak S_n.\;}
\]
更精确地，对每个置换 \(\sigma\in \mathfrak S_n\) 定义
\[
\alpha_\sigma\!\left(\prod_{i=0}^{n-1}q_i^{a_i}\right)
:=
\prod_{i=0}^{n-1}q_i^{\,\sigma(a_{\sigma^{-1}(i)})}.
\]
则
\[
\boxed{\;
\sigma\longmapsto \alpha_\sigma
\;}
\]
给出群同构
\[
\mathfrak S_n\xrightarrow{\sim}\operatorname{Aut}(S_n,\star).
\]
\leanverified{paper\_xi\_prime\_register\_semigroup\_automorphism\_rigidity}
\end{theorem}
\begin{proof}
由定理 \ref{thm:killo-finite-semigroup-prime-valuation-arithmetization}，
\[
\Phi:(S_n,\star)\xrightarrow{\sim}T_n:=([n]^{[n]},\circ),
\qquad
\Phi(N)=f_N.
\]
故只需求 \(T_n\) 的自同构群。

对每个 \(a\in[n]\)，记常值映射
\[
c_a:[n]\to[n],
\qquad
c_a(i)\equiv a.
\]
先刻画这些常值映射的纯半群论性质。若 \(x=c_a\)，则对任意 \(g\in T_n\) 都有
\[
x\circ g=c_a=x.
\]
反之，若某个 \(x\in T_n\) 满足
\[
x\circ g=x
\qquad
(\forall g\in T_n),
\]
则对任意 \(b\in[n]\) 取 \(g=c_b\)，便得
\[
x=x\circ c_b=c_{x(b)},
\]
故 \(x\) 必为常值映射。于是，\(\{c_a:a\in[n]\}\) 恰是 \(T_n\) 中全部左零元的集合。任意半群自同构都保持这一纯乘法性质，因此对每个
\[
\beta\in \operatorname{Aut}(T_n)
\]
必存在唯一置换 \(\sigma\in\mathfrak S_n\) 使
\[
\beta(c_a)=c_{\sigma(a)}
\qquad
(\forall a\in[n]).
\]

任取 \(f\in T_n\)。由于
\[
f\circ c_a=c_{f(a)},
\]
可得
\[
\beta(f)\circ c_{\sigma(a)}
=
\beta(f)\circ \beta(c_a)
=
\beta(f\circ c_a)
=
\beta(c_{f(a)})
=
c_{\sigma(f(a))}.
\]
另一方面，对任意 \(h\in T_n\) 与 \(b\in[n]\) 都有
\[
h\circ c_b=c_{h(b)}.
\]
故上式等价于
\[
\beta(f)(\sigma(a))=\sigma(f(a))
\qquad
(\forall a\in[n]),
\]
亦即
\[
\beta(f)=\sigma\circ f\circ \sigma^{-1}.
\]
因此 \(T_n\) 的每个自同构都由某个唯一的 \(\sigma\in\mathfrak S_n\) 给出，且复合律对应于置换乘法，从而
\[
\operatorname{Aut}(T_n)\cong \mathfrak S_n.
\]

最后把该描述经 \(\Phi\) 共轭回 \(S_n\)。若
\[
N=\prod_{i=0}^{n-1}q_i^{a_i},
\qquad
f_N(i)=a_i,
\]
则
\[
(\sigma\circ f_N\circ \sigma^{-1})(i)
=
\sigma(a_{\sigma^{-1}(i)}).
\]
再由 \(\Phi^{-1}\) 的定义，恰对应于
\[
\Phi^{-1}(\sigma\circ f_N\circ \sigma^{-1})
=
\prod_{i=0}^{n-1}q_i^{\,\sigma(a_{\sigma^{-1}(i)})}
=
\alpha_\sigma(N).
\]
故 \(\sigma\mapsto \alpha_\sigma\) 给出所述群同构。证毕。
\end{proof}

\begin{theorem}[素数寄存器幂等元的固定指数分层与完全计数]\label{thm:xi-prime-register-idempotent-fixed-index-stratification}
沿用定理 \ref{thm:killo-finite-semigroup-prime-valuation-arithmetization} 与推论 \ref{cor:killo-projection-idempotent-prime-register} 的记号，并记
\[
E_n:=\{N\in S_n:\ N\star N=N\}.
\]
对
\[
N=\prod_{i=0}^{n-1}q_i^{a_i}\in S_n
\]
定义固定指数集
\[
F_N:=\{j\in[n]:a_j=j\}.
\]
则下列命题等价：
\begin{enumerate}
  \item \(N\in E_n\)；
  \item 对每个 \(i\in[n]\) 都有
  \[
  a_i\in F_N;
  \]
  \item 存在唯一非空子集 \(F\subseteq[n]\) 使
  \[
  a_j=j\qquad (j\in F),
  \qquad
  a_i\in F\qquad (i\notin F),
  \]
  且此时 \(F=F_N=\operatorname{Im}(f_N)=\operatorname{Fix}(f_N)\)。
\end{enumerate}
进一步，若 \(F\subseteq[n]\) 非空且 \(|F|=k\)，则
\[
\boxed{\;
\#\{N\in E_n:\ F_N=F\}=k^{\,n-k}.\;}
\]
从而对每个 \(1\le k\le n\)，若记
\[
E_{n,k}:=\{N\in E_n:\ |F_N|=k\},
\]
则有
\[
\boxed{\;
|E_{n,k}|=\binom{n}{k}k^{\,n-k},
\qquad
|E_n|=\sum_{k=1}^{n}\binom{n}{k}k^{\,n-k}.\;}
\]
\leanverified{paper\_xi\_prime\_register\_idempotent\_fixed\_index\_stratification}
\end{theorem}
\begin{proof}
由推论 \ref{cor:killo-projection-idempotent-prime-register}，
\[
N\in E_n
\iff
f_N^2=f_N,
\]
并且一旦成立，就有
\[
\operatorname{Im}(f_N)=\operatorname{Fix}(f_N).
\]
由于
\[
f_N(i)=a_i,
\]
故
\[
\operatorname{Fix}(f_N)=\{j\in[n]:f_N(j)=j\}=\{j\in[n]:a_j=j\}=F_N.
\]
于是
\[
N\in E_n
\iff
a_i=f_N(i)\in \operatorname{Fix}(f_N)=F_N
\qquad
(\forall i\in[n]),
\]
这就给出了第 \(1\) 条与第 \(2\) 条的等价性。

再把 \(F:=F_N\) 展开。对 \(j\in F\)，定义已强制
\[
a_j=j.
\]
而对 \(i\notin F\)，第 \(2\) 条恰断言
\[
a_i\in F.
\]
反之，若某个非空子集 \(F\subseteq[n]\) 满足
\[
a_j=j\quad (j\in F),
\qquad
a_i\in F\quad (i\notin F),
\]
则 \(f_N\) 的像集恰为 \(F\)，并且 \(F\) 上每点都被固定，因此
\[
\operatorname{Im}(f_N)=F=\operatorname{Fix}(f_N),
\]
从而 \(f_N\) 幂等，亦即 \(N\in E_n\)。这就得到第 \(3\) 条。

固定非空子集 \(F\subseteq[n]\)，设 \(|F|=k\)。若要求 \(F_N=F\)，则对 \(j\in F\) 的坐标 \(a_j\) 被唯一强制为 \(j\)；而对其余 \(n-k\) 个坐标，每个都可独立取 \(F\) 中任一元素，故总数恰为
\[
k^{\,n-k}.
\]
再对所有大小为 \(k\) 的子集 \(F\) 求和，即得
\[
|E_{n,k}|=\binom{n}{k}k^{\,n-k}.
\]
最后对 \(k\) 求和得到
\[
|E_n|=\sum_{k=1}^{n}\binom{n}{k}k^{\,n-k}.
\]
证毕。
\end{proof}

\leanverified{paper\_xi\_chain\_interior\_boolean\_mobius\_characteristic\_maxchains}
\begin{theorem}[链上内算子格的 Boolean 秩、Möbius 函数、特征多项式与极大链闭式]\label{thm:xi-chain-interior-boolean-mobius-characteristic-maxchains}
沿用定理 \ref{thm:killo-chain-interior-godel-gcd-lcm}、定理 \ref{thm:xi-chain-interior-boolean-flag-closed-form} 与推论 \ref{cor:xi-chain-interior-dirichlet-mobius-characteristic-closure} 的记号。则
\[
(\mathrm{Int}(C_n),\preceq)\cong \mathcal P(\{1,\dots,n-1\})
\]
作为分次格同构；特别地，它是一个秩为 \(n-1\) 的 Boolean 格。并且对任意 \(I\preceq J\) 都有
\[
\boxed{\;
\mu_{\mathrm{Int}(C_n)}(I,J)
=
(-1)^{|\mathrm{Fix}(J)\setminus \mathrm{Fix}(I)|}
=
\mu_{\mathrm{arith}}\!\left(\frac{\kappa(J)}{\kappa(I)}\right).\;}
\]
其特征多项式满足
\[
\boxed{\;
\chi_{\mathrm{Int}(C_n)}(t)=(t-1)^{n-1}.\;}
\]
此外，极大链条数恰为
\[
\boxed{\;
\#\{\text{maximal chains in }\mathrm{Int}(C_n)\}=(n-1)!.\;}
\]
等价地，经 \(\kappa\) 算术化后，全部含 \(q_0\) 的平方自由因子构成一个 Boolean 格，而其 Möbius 函数与特征多项式完全由新增素因子集合控制。
\end{theorem}
\begin{proof}
由定理 \ref{thm:xi-chain-interior-boolean-flag-closed-form}，
\[
(\mathrm{Int}(C_n),\preceq)
\]
与包含 \(0\) 的子集格同构。去掉强制出现的元素 \(0\) 之后，即得
\[
(\mathrm{Int}(C_n),\preceq)\cong \mathcal P(\{1,\dots,n-1\}),
\]
故其确为秩 \(n-1\) 的 Boolean 格。

同一定理已给出
\[
\mu_{\mathrm{Int}(C_n)}(I,J)
=
(-1)^{|\mathrm{Fix}(J)\setminus \mathrm{Fix}(I)|},
\]
而推论 \ref{cor:xi-chain-interior-dirichlet-mobius-characteristic-closure} 进一步给出
\[
\mu_{\mathrm{Int}(C_n)}(I,J)
=
\mu_{\mathrm{arith}}\!\left(\frac{\kappa(J)}{\kappa(I)}\right)
\]
以及
\[
\chi_{\mathrm{Int}(C_n)}(t)=(t-1)^{n-1}.
\]
故前两条盒装公式已经成立。

最后计算极大链。于上述幂集同构下，一条极大链恰对应于
\[
\varnothing
\subset
\{i_1\}
\subset
\{i_1,i_2\}
\subset\cdots\subset
\{i_1,\dots,i_{n-1}\},
\]
其中
\[
(i_1,\dots,i_{n-1})
\]
是 \(\{1,\dots,n-1\}\) 的一个排列。反之，每个排列都唯一给出一条极大链。因此极大链总数正好是
\[
(n-1)!.
\]
证毕。
\end{proof}

\begin{theorem}[素数寄存器半群的 Green 秩分层与左零极小理想]\label{thm:xi-prime-register-green-rank-stratification}
\leanverified{paper\_xi\_prime\_register\_green\_rank\_stratification}
沿用定理 \ref{thm:killo-finite-semigroup-prime-valuation-arithmetization} 的记号，并对
\[
N\in S_n
\]
定义
\[
\operatorname{rk}(N):=|\operatorname{Im}(f_N)|.
\]
则成立下列结论：
\begin{enumerate}
  \item 对任意 \(N\in S_n\)，双边理想生成闭包满足
  \[
  \boxed{\;
  S_n\star N\star S_n
  =
  \{M\in S_n:\ \operatorname{rk}(M)\le \operatorname{rk}(N)\}.\;}
  \]
  因而 Green 的 \(\mathcal J\)-类恰由秩分层：
  \[
  \boxed{\;
  N\mathrel{\mathcal J}M
  \iff
  \operatorname{rk}(N)=\operatorname{rk}(M).\;}
  \]
  \item 对每个 \(1\le k\le n\)，秩恰为 \(k\) 的元素个数为
  \[
  \boxed{\;
  \#\{N\in S_n:\ \operatorname{rk}(N)=k\}
  =
  \binom{n}{k}\,k!\,S(n,k),\;}
  \]
  其中 \(S(n,k)\) 表示第二类 Stirling 数。
  \item 记
  \[
  Q_n:=\prod_{i=0}^{n-1}q_i,
  \qquad
  \mathcal K_n:=\{Q_n^c:\ c\in[n]\}.
  \]
  则
  \[
  \boxed{\;
  \mathcal K_n
  =
  \{N\in S_n:\ \operatorname{rk}(N)=1\},\qquad |\mathcal K_n|=n,\;}
  \]
  并且对任意 \(a,b\in[n]\) 都有
  \[
  \boxed{\;
  Q_n^a\star Q_n^b=Q_n^a.\;}
  \]
  特别地，\(\mathcal K_n\) 是 \((S_n,\star)\) 的唯一极小双边理想，且作为子半群同构于一个 \(n\) 元左零带。
\end{enumerate}
\end{theorem}
\begin{proof}
仍记
\[
T_n:=[n]^{[n]}.
\]
由定理 \ref{thm:killo-finite-semigroup-prime-valuation-arithmetization}，
\[
\Phi:(S_n,\star)\xrightarrow{\sim}(T_n,\circ)
\]
为半群同构，且
\[
\operatorname{rk}(N)=|\operatorname{Im}(f_N)|
\]
正是 \(f_N\) 在 \(T_n\) 中的像集秩。

先证第 \(1\) 条。对任意 \(u,f,v\in T_n\)，有
\[
\operatorname{Im}(f\circ v)\subseteq \operatorname{Im}(f),
\qquad
\operatorname{Im}(u\circ f)\subseteq u(\operatorname{Im}(f)),
\]
故
\[
\operatorname{rk}(u\circ f\circ v)\le \operatorname{rk}(f).
\]
因此
\[
T_n f T_n\subseteq \{g\in T_n:\operatorname{rk}(g)\le \operatorname{rk}(f)\}.
\]

反过来，设 \(g\in T_n\) 满足
\[
\ell:=\operatorname{rk}(g)\le \operatorname{rk}(f)=:k.
\]
取子集
\[
A\subseteq \operatorname{Im}(f),
\qquad
|A|=\ell,
\]
并记
\[
B:=\operatorname{Im}(g).
\]
选定一个双射
\[
\theta:A\xrightarrow{\sim}B.
\]
对每个 \(a\in A\)，任选一点 \(x_a\in[n]\) 使
\[
f(x_a)=a.
\]
定义映射
\[
v:[n]\to[n],
\qquad
v(t):=x_{\theta^{-1}(g(t))},
\]
以及
\[
u:[n]\to[n]
\]
满足
\[
u(a)=\theta(a)\qquad (a\in A),
\]
在 \([n]\setminus A\) 上任意取值。则对每个 \(t\in[n]\)，有
\[
f(v(t))
=
f(x_{\theta^{-1}(g(t))})
=
\theta^{-1}(g(t)),
\]
故
\[
u(f(v(t)))=\theta(\theta^{-1}(g(t)))=g(t).
\]
于是
\[
g=u\circ f\circ v\in T_n f T_n.
\]
因此
\[
T_n f T_n=\{g\in T_n:\operatorname{rk}(g)\le \operatorname{rk}(f)\}.
\]
再经 \(\Phi^{-1}\) 传回 \(S_n\)，便得到所述双边理想公式。由 Green 的 \(\mathcal J\)-关系定义，
\[
N\mathrel{\mathcal J}M
\iff
S_n\star N\star S_n=S_n\star M\star S_n,
\]
故等价于
\[
\operatorname{rk}(N)=\operatorname{rk}(M).
\]

第 \(2\) 条是满射计数的直接转写。构造一个秩为 \(k\) 的映射
\[
f:[n]\to[n]
\]
时，先选像集
\[
\operatorname{Im}(f)\subseteq[n],
\qquad
|\operatorname{Im}(f)|=k,
\]
有 \(\binom{n}{k}\) 种；再在该 \(k\) 元像集上选一个满射，其数量为经典公式
\[
k!\,S(n,k).
\]
故秩为 \(k\) 的映射总数为
\[
\binom{n}{k}k!S(n,k).
\]
经 \(\Phi\) 与 \(S_n\) 一一对应，即得第 \(2\) 条。

第 \(3\) 条对应于 \(k=1\) 层。秩 \(1\) 的映射恰是常值映射
\[
c_a(i)\equiv a
\qquad
(a\in[n]).
\]
而在 \(S_n\) 中，\(c_a\) 的原像正是
\[
Q_n^a=\prod_{i=0}^{n-1}q_i^a.
\]
因此
\[
\mathcal K_n=\{N\in S_n:\operatorname{rk}(N)=1\},
\qquad
|\mathcal K_n|=n.
\]
并且
\[
f_{Q_n^a\star Q_n^b}
=
f_{Q_n^a}\circ f_{Q_n^b}
=
c_a\circ c_b
=
c_a
=
f_{Q_n^a},
\]
故
\[
Q_n^a\star Q_n^b=Q_n^a.
\]
由第 \(1\) 条，秩 \(1\) 层正是最小的非空双边理想层，因此 \(\mathcal K_n\) 是唯一极小双边理想；又因其乘法满足左零律，故它同构于一个 \(n\) 元左零带。证毕。
\end{proof}

\paragraph{链上内算子的关联代数反演与随机 Euler 统计}
在链上内算子格的 Boolean 刚性、平方自由 Gödel 像的 \(\gcd/\mathrm{lcm}\) 算术化以及 Möbius 函数闭式已经固定之后，还可把这条有限格链进一步推进到关联代数与概率统计两侧的完全显式闭式。

\begin{theorem}[链上内算子 Gödel 整除格的 Möbius 反演]\label{thm:xi-chain-interior-incidence-algebra-mobius-inversion}
沿用定理 \ref{thm:killo-chain-interior-godel-gcd-lcm} 与定理 \ref{thm:xi-chain-interior-boolean-mobius-characteristic-maxchains} 的记号，记
\[
\mathcal I_n:=\mathrm{Int}(C_n),
\qquad
Q_n:=\prod_{i=0}^{n-1}q_i,
\]
\[
\mathcal D_n
:=
\{d\mid Q_n:\ d\ \text{平方自由且}\ q_0\mid d\}.
\]
则映射
\[
\kappa:\mathcal I_n\longrightarrow \mathcal D_n
\]
在偏序 \(\preceq\) 与整除 \(\mid\) 之间给出格同构，并且对任意 \(I\preceq J\) 都有
\[
\mu_{\mathcal I_n}(I,J)
=
(-1)^{\omega(\kappa(J)/\kappa(I))}
=
\mu_{\mathrm{arith}}\!\left(\frac{\kappa(J)}{\kappa(I)}\right).
\]
等价地，对任意交换群 \(A\) 上的函数 \(f,g:\mathcal D_n\to A\)，关系
\[
g(d)=\sum_{\substack{e\in \mathcal D_n\\ e\mid d}}f(e)
\qquad
(d\in \mathcal D_n)
\]
成立当且仅当
\[
f(d)
=
\sum_{\substack{e\in \mathcal D_n\\ e\mid d}}
\mu_{\mathrm{arith}}\!\left(\frac{d}{e}\right)\,g(e)
=
\sum_{\substack{e\in \mathcal D_n\\ e\mid d}}
(-1)^{\omega(d/e)}\,g(e).
\]
\end{theorem}
\begin{proof}
由定理 \ref{thm:killo-chain-interior-godel-gcd-lcm}，\(\kappa\) 把 \(\mathcal I_n\) 规范识别为全部含 \(q_0\) 的平方自由因子所成之整除格，并把 meet/join 精确送到 \(\gcd/\mathrm{lcm}\)。因此 \(\kappa\) 的确给出 \((\mathcal I_n,\preceq)\) 与 \((\mathcal D_n,\mid)\) 的格同构。

再由定理 \ref{thm:xi-chain-interior-boolean-mobius-characteristic-maxchains}，
\[
\mu_{\mathcal I_n}(I,J)
=
(-1)^{|\mathrm{Fix}(J)\setminus \mathrm{Fix}(I)|}.
\]
而 \(\kappa(J)/\kappa(I)\) 正好是由新增固定点对应素数所组成的平方自由整数，故
\[
|\mathrm{Fix}(J)\setminus \mathrm{Fix}(I)|
=
\omega(\kappa(J)/\kappa(I)),
\]
从而得到 Möbius 函数的算术闭式。

最后，有限偏序集上的一般 Möbius 反演公式适用于 \(\mathcal I_n\)；经由格同构 \(\kappa\) 转写到 \(\mathcal D_n\) 上，即得所示整除偏序反演公式。证毕。
\end{proof}

\begin{theorem}[随机内算子平方自由像的 Euler 乘积与对数方差闭式]\label{thm:xi-chain-interior-random-godel-euler-logvariance}
\leanverified{paper\_xi\_chain\_interior\_random\_godel\_euler\_logvariance}
沿用定理 \ref{thm:killo-chain-interior-godel-gcd-lcm} 的记号，并设随机变量
\[
I\sim \mathrm{Unif}\bigl(\mathrm{Int}(C_n)\bigr).
\]
则存在独立同分布随机变量
\[
\varepsilon_1,\dots,\varepsilon_{n-1}
\stackrel{\mathrm{i.i.d.}}{\sim}
\mathrm{Bernoulli}\!\left(\frac12\right)
\]
使得
\[
\kappa(I)=q_0\prod_{j=1}^{n-1}q_j^{\varepsilon_j}.
\]
因而对任意 \(s\in\CC\) 都有
\[
\EE\!\bigl[\kappa(I)^s\bigr]
=
q_0^s\prod_{j=1}^{n-1}\frac{1+q_j^s}{2}.
\]
特别地，
\[
\EE[\log \kappa(I)]
=
\log q_0+\frac12\sum_{j=1}^{n-1}\log q_j,
\qquad
\mathrm{Var}(\log \kappa(I))
=
\frac14\sum_{j=1}^{n-1}(\log q_j)^2.
\]
\end{theorem}
\begin{proof}
由定理 \ref{thm:xi-chain-interior-boolean-flag-closed-form}，
\[
\mathcal I_n\cong \mathcal P(\{1,\dots,n-1\})
\]
作为 Boolean 格同构；而在均匀分布下，这正对应于每个指标 \(j\in\{1,\dots,n-1\}\) 是否属于固定点集的独立等概率选择。于是存在独立 Bernoulli\((1/2)\) 变量 \(\varepsilon_j\) 使
\[
\mathrm{Fix}(I)=\{0\}\cup \{j:\varepsilon_j=1\}.
\]
再由定理 \ref{thm:killo-chain-interior-godel-gcd-lcm} 的定义，
\[
\kappa(I)=q_0\prod_{j=1}^{n-1}q_j^{\varepsilon_j},
\]
得到第一式。

对任意 \(s\in\CC\)，由独立性可得
\[
\EE\!\bigl[\kappa(I)^s\bigr]
=
q_0^s\prod_{j=1}^{n-1}\EE\!\bigl[q_j^{s\varepsilon_j}\bigr]
=
q_0^s\prod_{j=1}^{n-1}\frac{1+q_j^s}{2},
\]
即得 Euler 乘积闭式。

最后，
\[
\log\kappa(I)=\log q_0+\sum_{j=1}^{n-1}\varepsilon_j\log q_j.
\]
由
\[
\EE[\varepsilon_j]=\frac12,
\qquad
\mathrm{Var}(\varepsilon_j)=\frac14,
\]
以及独立性消去协方差项，便得到
\[
\EE[\log \kappa(I)]
=
\log q_0+\frac12\sum_{j=1}^{n-1}\log q_j,
\]
\[
\mathrm{Var}(\log \kappa(I))
=
\frac14\sum_{j=1}^{n-1}(\log q_j)^2.
\]
证毕。
\end{proof}

\paragraph{素数寄存器幂等元的唯一标准形与链型规范截面的固定点稀薄率}
在素数寄存器幂等元的固定指数分层、链上内算子格的布尔化以及链型幂等元的自指合成规律已经确定之后，还可把固定点纤维的内部结构进一步压缩为如下两条完全显式的后果。

\begin{theorem}[素数寄存器幂等元的唯一标准形与固定点纤维计数]\label{thm:xi-prime-register-idempotent-unique-normal-form}
\leanverified{paper\_xi\_prime\_register\_idempotent\_unique\_normal\_form}
沿用定理 \ref{thm:killo-finite-semigroup-prime-valuation-arithmetization}、推论 \ref{cor:killo-projection-idempotent-prime-register} 与定理 \ref{thm:xi-prime-register-idempotent-fixed-index-stratification} 的记号。则 \(N\in S_n\) 满足 \(N\star N=N\) 当且仅当存在唯一非空子集 \(F\subseteq[n]\) 与唯一映射
\[
\phi:[n]\setminus F\to F
\]
使得
\[
N=
\left(\prod_{j\in F}q_j^{\,j}\right)
\left(\prod_{i\notin F}q_i^{\,\phi(i)}\right).
\]
此时
\[
F=\operatorname{Fix}(f_N)=\operatorname{Im}(f_N).
\]
特别地，若 \(|F|=k\)，则满足 \(\operatorname{Fix}(f_N)=F\) 的幂等元恰有
\[
k^{\,n-k}
\]
个；因此固定点数为 \(k\) 的幂等元总数为
\[
\binom{n}{k}k^{\,n-k},
\]
而全部幂等元总数为
\[
\sum_{k=1}^{n}\binom{n}{k}k^{\,n-k}.
\]
\end{theorem}
\begin{proof}
把
\[
N=\prod_{i=0}^{n-1}q_i^{a_i}
\]
写成估值坐标形式。由定理 \ref{thm:xi-prime-register-idempotent-fixed-index-stratification}，\(N\star N=N\) 当且仅当存在唯一非空子集 \(F\subseteq[n]\) 使
\[
a_j=j\qquad (j\in F),
\qquad
a_i\in F\qquad (i\notin F),
\]
并且此时
\[
F=\operatorname{Fix}(f_N)=\operatorname{Im}(f_N).
\]
于是对每个 \(i\notin F\) 定义
\[
\phi(i):=a_i\in F,
\]
便得到所示分解式
\[
N=
\left(\prod_{j\in F}q_j^{\,j}\right)
\left(\prod_{i\notin F}q_i^{\,\phi(i)}\right).
\]
由于素因子分解唯一，故给定 \(N\) 后，\(F\) 与 \(\phi\) 都被唯一确定。

反过来，若给定非空 \(F\subseteq[n]\) 与映射 \(\phi:[n]\setminus F\to F\)，并按上式定义 \(N\)，则
\[
f_N(j)=j\qquad (j\in F),
\qquad
f_N(i)=\phi(i)\in F\qquad (i\notin F).
\]
因此对任意 \(i\in[n]\) 都有
\[
f_N(f_N(i))=f_N(i),
\]
故 \(f_N\) 幂等，等价地 \(N\star N=N\)。

若固定 \(\operatorname{Fix}(f_N)=F\) 且 \(|F|=k\)，则 \(j\in F\) 的指数被唯一强制，而对每个 \(i\notin F\) 都可独立选择 \(\phi(i)\in F\)，故恰有
\[
k^{\,n-k}
\]
个幂等元。再对所有大小为 \(k\) 的子集 \(F\) 求和，即得
\[
\binom{n}{k}k^{\,n-k}
\]
与总数公式。证毕。
\end{proof}

\begin{theorem}[链上规范截面在固定零点幂等纤维中的唯一单调--收缩实现与超指数稀薄率]\label{thm:xi-prime-register-canonical-fixedzero-section-sparsity}
\leanverified{paper\_xi\_prime\_register\_canonical\_fixedzero\_section\_sparsity}
沿用定理 \ref{thm:killo-chain-interior-godel-gcd-lcm}、定理 \ref{thm:xi-prime-register-idempotent-unique-normal-form} 与定理 \ref{thm:xi-chain-idempotent-star-saturation-comparable-gcd} 的记号。对任意满足 \(0\in F\subseteq[n]\) 的子集，记
\[
N_F^{\mathrm{can}}:=\Phi^{-1}(I_F)=N_F.
\]
则：
\begin{enumerate}
  \item 在全部满足 \(\operatorname{Fix}(f_N)=F\) 的幂等元中，恰有一个对应变换 \(f_N\) 同时单调且收缩；该幂等元正是 \(N_F^{\mathrm{can}}\)；
  \item 因而
  \[
  \mathrm{Int}(C_n)\cong \{N_F^{\mathrm{can}}:0\in F\subseteq[n]\}
  \]
  给出 prime-register 幂等层中的一个规范截面；
  \item 该截面的基数为
  \[
  |\mathrm{Int}(C_n)|=2^{n-1};
  \]
  \item 若记
  \[
  E_n^{(0)}
  :=
  \#\{N\in S_n:\ N\star N=N,\ f_N(0)=0\},
  \]
  则
  \[
  E_n^{(0)}
  =
  \sum_{k=1}^{n}\binom{n-1}{k-1}k^{\,n-k},
  \]
  并且相对密度满足
  \[
  \frac{|\mathrm{Int}(C_n)|}{E_n^{(0)}}
  \le
  \exp\!\Bigl(-\frac12 n\log n+O(n)\Bigr)
  \qquad (n\to\infty).
  \]
\end{enumerate}
\end{theorem}
\begin{proof}
对任意 \(0\in F\subseteq[n]\)，定理 \ref{thm:killo-chain-interior-godel-gcd-lcm} 已给出唯一的内算子
\[
I_F(i)=\max\{j\in F:j\le i\},
\]
其固定点集正为 \(F\)。经 \(\Phi^{-1}\) 拉回，得到幂等元 \(N_F^{\mathrm{can}}\)，且对应变换 \(f_{N_F^{\mathrm{can}}}=I_F\) 显然同时单调且收缩。

反之，设 \(N\in S_n\) 满足 \(N\star N=N\)、\(\operatorname{Fix}(f_N)=F\)，并且 \(f_N\) 单调且收缩。由推论 \ref{cor:killo-projection-idempotent-prime-register}，
\[
\operatorname{Im}(f_N)=\operatorname{Fix}(f_N)=F.
\]
任取 \(i\in[n]\)，记
\[
j:=\max(F\cap\{0,\dots,i\}),
\]
其中该最大元存在，因为 \(0\in F\)。由于 \(f_N(i)\in F\) 且收缩性给出 \(f_N(i)\le i\)，故
\[
f_N(i)\in F\cap\{0,\dots,i\},
\]
从而
\[
f_N(i)\le j.
\]
另一方面，\(j\in F\) 意味着 \(f_N(j)=j\)，再由单调性与 \(j\le i\) 得
\[
j=f_N(j)\le f_N(i).
\]
故 \(f_N(i)=j=I_F(i)\)。由于 \(i\) 任意，遂有
\[
f_N=I_F,
\]
从而 \(N=N_F^{\mathrm{can}}\)。这证明了第 \(1\) 条，也得到第 \(2\) 条。

第 \(3\) 条即定理 \ref{thm:xi-chain-interior-boolean-flag-closed-form} 的基数公式
\[
|\mathrm{Int}(C_n)|=2^{n-1}.
\]

再来计算 \(E_n^{(0)}\)。由定理 \ref{thm:xi-prime-register-idempotent-unique-normal-form}，\(N\star N=N\) 且 \(f_N(0)=0\) 等价于其固定点集 \(F=\operatorname{Fix}(f_N)\) 满足 \(0\in F\)。当 \(|F|=k\) 时，包含 \(0\) 的此类子集共有
\[
\binom{n-1}{k-1}
\]
个，而每个固定点集对应的幂等元数目为
\[
k^{\,n-k}.
\]
故
\[
E_n^{(0)}=\sum_{k=1}^{n}\binom{n-1}{k-1}k^{\,n-k}.
\]

最后估计相对密度。取
\[
m:=\left\lfloor \frac n2\right\rfloor.
\]
则
\[
E_n^{(0)}
\ge
\binom{n-1}{m-1}m^{\,n-m}.
\]
由 Stirling 公式，
\[
\binom{n-1}{m-1}=\exp(O(n)),
\qquad
m^{\,n-m}
=
\exp\!\Bigl(\frac12 n\log n+O(n)\Bigr).
\]
于是
\[
E_n^{(0)}
\ge
\exp\!\Bigl(\frac12 n\log n+O(n)\Bigr).
\]
另一方面，
\[
|\mathrm{Int}(C_n)|=2^{n-1}=\exp(O(n)).
\]
二式相除即得
\[
\frac{|\mathrm{Int}(C_n)|}{E_n^{(0)}}
\le
\exp\!\Bigl(-\frac12 n\log n+O(n)\Bigr).
\]
证毕。
\end{proof}
\paragraph{素数加权二维读出的白化圆对称、最小层析与 Mahalanobis 阈值链}
在定理 \ref{thm:killo-prime-weighted-elliptic-fingerprint-clt} 已把素数加权二维读出识别为椭圆型高斯极限指纹，而推论 \ref{cor:killo-prime-weighted-angular-law-sigma-recovery} 又已把其角向边际压缩为对 \(\Sigma_\infty\) 的唯一恢复之后，还可进一步抽取出如下五条彼此闭合的后果。它们分别把椭圆极限白化为严格圆对称高斯，把固定扇区概率压缩为纯角长度，把两条线性读出的去耦合判据写成 \(\Sigma_\infty\)-正交条件，把二维椭圆指纹压缩为三探针最小层析，并把全部 Mahalanobis 型二值判定统一为同一个 \(\chi_2^2\) 半径变量所支配的嵌套 Bernoulli 链。

\begin{theorem}[白化后极限严格圆对称，Mahalanobis 半径平方满足 \texorpdfstring{$\chi_2^2$}{chi-square-2} 极限律]\label{thm:xi-time-part61aca-primeweighted-whitened-radial-chi2}
\leanverified{paper\_xi\_time\_part61aca\_primeweighted\_whitened\_radial\_chi2}
沿用定理 \ref{thm:killo-prime-weighted-elliptic-fingerprint-clt} 的设定，记
\[
Y_T:=\frac{1}{\sigma}\Sigma_\infty^{-1/2}
\frac{S_T}{\sqrt{\operatorname{tr}(V_T)}}\in\RR^2.
\]
则有
\[
Y_T\Longrightarrow \mathcal N(0,I_2).
\]
特别地，Mahalanobis 半径平方
\[
R_T^2
:=
\frac{1}{\sigma^2\operatorname{tr}(V_T)}
S_T^\top \Sigma_\infty^{-1}S_T
=
\|Y_T\|^2
\]
满足
\[
R_T^2\Longrightarrow \chi_2^2.
\]
因此对任意 \(q\ge 0\)，都有
\[
\lim_{T\to\infty}\PP(R_T^2\le q)=1-e^{-q/2},
\qquad
\lim_{T\to\infty}\PP(R_T^2>q)=e^{-q/2}.
\]
\end{theorem}
\begin{proof}
由定理 \ref{thm:killo-prime-weighted-elliptic-fingerprint-clt}，
\[
\frac{S_T}{\sqrt{\operatorname{tr}(V_T)}}
\Longrightarrow
\mathcal N(0,\sigma^2\Sigma_\infty),
\]
且 \(\Sigma_\infty\) 正定。对极限施加线性变换
\[
x\longmapsto \sigma^{-1}\Sigma_\infty^{-1/2}x
\]
并应用连续映射定理，立得
\[
Y_T\Longrightarrow \mathcal N(0,I_2).
\]
再由映射 \(y\mapsto \|y\|^2\) 的连续性，
\[
R_T^2=\|Y_T\|^2\Longrightarrow \|Z\|^2,
\qquad
Z\sim \mathcal N(0,I_2).
\]
而二维标准高斯范数平方即 \(\chi_2^2\) 分布，其分布函数为
\[
\PP(\chi_2^2\le q)=1-e^{-q/2}.
\]
结论成立。证毕。
\end{proof}

\begin{corollary}[白化角变量渐近均匀，固定扇区概率仅由夹角决定]\label{cor:xi-time-part61aca-primeweighted-whitened-sector-uniformity}
\leanverified{paper\_xi\_time\_part61aca\_primeweighted\_whitened\_sector\_uniformity}
沿用定理 \ref{thm:xi-time-part61aca-primeweighted-whitened-radial-chi2} 的记号。对任意 Borel 扇区
\[
\mathcal W(\alpha,\beta)
:=
\{(r\cos\theta,r\sin\theta):r\ge 0,\ \theta\in[\alpha,\beta]\},
\qquad
0\le \beta-\alpha\le 2\pi,
\]
若其边界对极限分布为零测，则有
\[
\PP\bigl(Y_T\in \mathcal W(\alpha,\beta)\bigr)
\longrightarrow
\frac{\beta-\alpha}{2\pi}.
\]
等价地，白化坐标中的极限方向变量在 \([0,2\pi)\) 上服从均匀分布。
\end{corollary}
\begin{proof}
由定理 \ref{thm:xi-time-part61aca-primeweighted-whitened-radial-chi2}，
\[
Y_T\Longrightarrow Z\sim \mathcal N(0,I_2).
\]
二维标准高斯分布严格旋转不变，故其极坐标中的角变量在 \([0,2\pi)\) 上均匀分布，并与半径独立。因此
\[
\PP\bigl(Z\in\mathcal W(\alpha,\beta)\bigr)=\frac{\beta-\alpha}{2\pi}.
\]
再由 Portmanteau 定理即得结论。证毕。
\end{proof}

\begin{theorem}[任意两条线性读出的渐近独立性当且仅当 \texorpdfstring{$\Sigma_\infty$}{Sigma-infty}-正交]\label{thm:xi-time-part61aca-primeweighted-sigma-orthogonal-independence}
\leanverified{paper\_xi\_time\_part61aca\_primeweighted\_sigma\_orthogonal\_independence}
沿用定理 \ref{thm:killo-prime-weighted-elliptic-fingerprint-clt} 的设定。对任意非零向量 \(a,b\in\RR^2\)，定义标量读出
\[
X_T(a):=\frac{a^\top S_T}{\sqrt{\operatorname{tr}(V_T)}},
\qquad
X_T(b):=\frac{b^\top S_T}{\sqrt{\operatorname{tr}(V_T)}}.
\]
则
\[
\bigl(X_T(a),X_T(b)\bigr)
\Longrightarrow
\mathcal N(0,\Gamma_{a,b}),
\]
其中
\[
\Gamma_{a,b}
:=
\sigma^2
\begin{pmatrix}
a^\top\Sigma_\infty a & a^\top\Sigma_\infty b\\
a^\top\Sigma_\infty b & b^\top\Sigma_\infty b
\end{pmatrix}.
\]
特别地，下列两命题等价：
\[
a^\top\Sigma_\infty b=0
\qquad\Longleftrightarrow\qquad
X_T(a)\ \text{与}\ X_T(b)\ \text{渐近独立}.
\]
若进一步满足 \(a^\top\Sigma_\infty b=0\)，则对任意 \(\varepsilon_1,\varepsilon_2\in\{\pm1\}\)，都有
\[
\PP\bigl(\operatorname{sgn}X_T(a)=\varepsilon_1,\ \operatorname{sgn}X_T(b)=\varepsilon_2\bigr)
\longrightarrow
\frac14.
\]
\end{theorem}
\begin{proof}
记
\[
A:=
\begin{pmatrix}
a^\top\\
b^\top
\end{pmatrix}.
\]
则
\[
\begin{pmatrix}
X_T(a)\\
X_T(b)
\end{pmatrix}
=
A\frac{S_T}{\sqrt{\operatorname{tr}(V_T)}}.
\]
由定理 \ref{thm:killo-prime-weighted-elliptic-fingerprint-clt} 与连续映射定理，
\[
A\frac{S_T}{\sqrt{\operatorname{tr}(V_T)}}
\Longrightarrow
\mathcal N(0,\sigma^2A\Sigma_\infty A^\top),
\]
而
\[
\sigma^2A\Sigma_\infty A^\top=\Gamma_{a,b}.
\]
这就得到所示联合高斯极限。

二维中心高斯向量独立当且仅当协方差为零，因此渐近独立当且仅当
\[
a^\top\Sigma_\infty b=0.
\]
在该条件下，极限分布是两个独立的一维中心高斯的乘积，故四个符号象限的极限概率均为 \(1/4\)。由于各象限边界由两条直线组成，对连续高斯分布均为零测，Portmanteau 定理遂给出有限 \(T\) 的极限结论。证毕。
\end{proof}

\begin{theorem}[三探针方向方差完全恢复椭圆极限指纹，且两探针永远不足]\label{thm:xi-time-part61aca-primeweighted-three-probe-minimal-tomography}
\leanverified{paper\_xi\_time\_part61aca\_primeweighted\_three\_probe\_minimal\_tomography}
沿用定理 \ref{thm:killo-prime-weighted-elliptic-fingerprint-clt} 的设定，并记
\[
M_\infty:=\sigma^2\Sigma_\infty=
\begin{pmatrix}
M_{11}&M_{12}\\
M_{12}&M_{22}
\end{pmatrix}.
\]
令 \(e_1=(1,0)^\top\)、\(e_2=(0,1)^\top\)，并定义三个方向上的归一化方差极限
\[
v_1:=\lim_{T\to\infty}\frac{\Var(e_1^\top S_T)}{\operatorname{tr}(V_T)},
\qquad
v_2:=\lim_{T\to\infty}\frac{\Var(e_2^\top S_T)}{\operatorname{tr}(V_T)},
\]
\[
v_3:=\lim_{T\to\infty}\frac{\Var((e_1+e_2)^\top S_T)}{\operatorname{tr}(V_T)}.
\]
则有显式恢复公式
\[
M_{11}=v_1,
\qquad
M_{22}=v_2,
\qquad
M_{12}=\frac{v_3-v_1-v_2}{2}.
\]
因此，仅凭三个标量探针 \(e_1,e_2,e_1+e_2\) 的极限方差，即可完全恢复二维椭圆极限指纹 \(M_\infty=\sigma^2\Sigma_\infty\)。

进一步，三探针在此意义下是最小的：任意只含两个方向的方差数据都不足以唯一确定 \(M_\infty\)。
\end{theorem}
\begin{proof}
由于 \(\xi_t\) 独立同分布且 \(\EE\xi_t=0\)、\(\EE\xi_t^2=\sigma^2\)，对任意固定 \(c\in\RR^2\) 都有精确恒等式
\[
\Var(c^\top S_T)=\sigma^2\,c^\top V_T c.
\]
两边同除以 \(\operatorname{tr}(V_T)\)，再令 \(T\to\infty\)，由
\[
\frac{V_T}{\operatorname{tr}(V_T)}\to \Sigma_\infty
\]
得到
\[
\lim_{T\to\infty}\frac{\Var(c^\top S_T)}{\operatorname{tr}(V_T)}
=
\sigma^2 c^\top\Sigma_\infty c
=
c^\top M_\infty c.
\]
分别取 \(c=e_1,e_2,e_1+e_2\)，便得
\[
v_1=M_{11},\qquad v_2=M_{22},
\]
\[
v_3=M_{11}+2M_{12}+M_{22}.
\]
消元即得
\[
M_{12}=\frac{v_3-v_1-v_2}{2}.
\]
故三探针足以唯一恢复 \(M_\infty\)。

再证两探针不足。设只给定两个方向 \(a,b\in\RR^2\) 的方差数据
\[
a^\top M_\infty a,\qquad b^\top M_\infty b.
\]
若 \(a,b\) 共线，则两条数据本质上只提供一个方向上的标量约束，显然不足以确定二维对称矩阵。

若 \(a,b\) 不共线，则对称矩阵空间 \(\mathrm{Sym}_2(\RR)\) 维数为 \(3\)，而线性映射
\[
\mathcal L:\mathrm{Sym}_2(\RR)\to \RR^2,
\qquad
H\longmapsto (a^\top Ha,\ b^\top Hb)
\]
的秩至多为 \(2\)，故其核中存在非零矩阵 \(H\neq 0\)。于是
\[
a^\top (M_\infty+tH)a=a^\top M_\infty a,
\qquad
b^\top (M_\infty+tH)b=b^\top M_\infty b
\]
对一切 \(t\in\RR\) 都成立。又因正定锥是开集，而 \(M_\infty\) 正定，故当 \(|t|\) 充分小时，\(M_\infty+tH\) 仍正定。于是存在不止一个正定矩阵与同一两探针方差数据相容，故两探针永远不足。证毕。
\end{proof}

\begin{theorem}[Mahalanobis 半径诱导的全部二值阈值在极限中形成单参数嵌套 Bernoulli 链]\label{thm:xi-time-part61aca-primeweighted-mahalanobis-threshold-bernoulli-chain}
\leanverified{paper\_xi\_time\_part61aca\_primeweighted\_mahalanobis\_threshold\_bernoulli\_chain}
沿用定理 \ref{thm:xi-time-part61aca-primeweighted-whitened-radial-chi2} 的记号。对任意 \(q\ge 0\)，定义阈值指标
\[
J_T(q):=\mathbf 1_{\{R_T^2>q\}}.
\]
则对每个固定 \(q\)，有
\[
J_T(q)\Longrightarrow \mathrm{Bernoulli}(e^{-q/2}).
\]
更强地，对任意有限阈值链
\[
0\le q_1<q_2<\cdots<q_r,
\]
向量
\[
\bigl(J_T(q_1),\dots,J_T(q_r)\bigr)
\]
在分布上收敛到嵌套 Bernoulli 链
\[
\bigl(\mathbf 1_{\{E>q_1\}},\dots,\mathbf 1_{\{E>q_r\}}\bigr),
\qquad
E\sim \chi_2^2.
\]
因此对任意 \(1\le i\le j\le r\)，都有
\[
\PP\bigl(J_T(q_i)=1,\ J_T(q_j)=1\bigr)\longrightarrow e^{-q_j/2},
\]
\[
\Cov\bigl(J_T(q_i),J_T(q_j)\bigr)\longrightarrow
e^{-q_j/2}\bigl(1-e^{-q_i/2}\bigr).
\]
\end{theorem}
\begin{proof}
由定理 \ref{thm:xi-time-part61aca-primeweighted-whitened-radial-chi2}，
\[
R_T^2\Longrightarrow E\sim \chi_2^2.
\]
而 \(\chi_2^2\) 的尾分布满足
\[
\PP(E>q)=e^{-q/2}.
\]
由于集合 \((q,\infty)\) 的边界点 \(\{q\}\) 对连续分布 \(E\) 为零测，Portmanteau 定理给出
\[
J_T(q)=\mathbf 1_{\{R_T^2>q\}}
\Longrightarrow
\mathbf 1_{\{E>q\}}
\sim \mathrm{Bernoulli}(e^{-q/2}).
\]

对有限阈值链，映射
\[
x\longmapsto
\bigl(\mathbf 1_{\{x>q_1\}},\dots,\mathbf 1_{\{x>q_r\}}\bigr)
\]
仅在有限集合 \(\{q_1,\dots,q_r\}\) 处不连续，而 \(\chi_2^2\) 分布对该有限集赋零质量，因此再次由连续映射定理得到
\[
\bigl(J_T(q_1),\dots,J_T(q_r)\bigr)
\Longrightarrow
\bigl(\mathbf 1_{\{E>q_1\}},\dots,\mathbf 1_{\{E>q_r\}}\bigr).
\]
再由 \(q_i\le q_j\) 时事件包含关系
\[
\{E>q_j\}\subseteq \{E>q_i\}
\]
可得
\[
\PP(E>q_i,\ E>q_j)=\PP(E>q_j)=e^{-q_j/2},
\]
以及
\[
\Cov(\mathbf 1_{\{E>q_i\}},\mathbf 1_{\{E>q_j\}})
=
e^{-q_j/2}-e^{-q_i/2}e^{-q_j/2}
=
e^{-q_j/2}\bigl(1-e^{-q_i/2}\bigr).
\]
结论成立。证毕。
\end{proof}

\noindent
由此，素数加权二维读出的椭圆高斯极限在正文中获得一条完全闭合的白化--层析--判别链：白化后极限严格圆对称，全部固定方向扇区只留下均匀角流；任意两条线性探针的去耦合恰由 \(\Sigma_\infty\)-正交控制；二维椭圆指纹 \(\sigma^2\Sigma_\infty\) 又可由三条一维方差探针以最小方式完全恢复；而一切 Mahalanobis 型二值判定最终都退化为同一个 \(\chi_2^2\) 半径变量的嵌套阈值过程。于是，素数加权读出的二维统计结构便不再是开放的高维对象，而被压缩为一个可白化、可最小层析、且可由单参数半径完全判别的显式闭式体系。

\paragraph{链上平方自由 Gödel 编码的平均复杂度最优律}
在定理 \ref{thm:xi-chain-interior-godel-euler-mobius-factorization} 已把链上规范投影的平方自由像写成 Euler--Dirichlet 与 Möbius 扭曲闭式，而推论 \ref{cor:xi-chain-interior-canonical-arithmeticization-optimal-order} 又已把其最坏码长锁定在 \(\Theta(n\log n)\) 之后，还可进一步把同一平方自由编码的平均对数复杂度完全闭式化。由此，链上素数寄存器外置的制度成本便不再只是极端输入所强制，而在均匀分布下同样保持精确的 \(n\log n\) 级下界。

\begin{theorem}[链上归一化平方自由 Gödel 编码的平均对数复杂度达到精确极小值]\label{thm:xi-chain-interior-godel-average-budget-optimality}
\leanverified{paper\_xi\_chain\_interior\_godel\_average\_budget\_optimality}
沿用定理 \ref{thm:killo-chain-interior-godel-gcd-lcm} 与定理 \ref{thm:xi-chain-interior-boolean-flag-closed-form} 的记号，记
\[
\mathcal I_n:=\mathrm{Int}(C_n),
\qquad
C_n=\{0<1<\cdots<n-1\}.
\]
固定互异素数 \(q_0,\dots,q_{n-1}\)，并定义归一化平方自由 Gödel 编码
\[
\kappa'(I_F):=\frac{\kappa(I_F)}{q_0}
=
\prod_{i\in F\setminus\{0\}}q_i
\qquad
(0\in F\subseteq\{0,\dots,n-1\}).
\]
把 \(\mathcal I_n\) 与二元字集 \(\{0,1\}^{n-1}\) 通过对应
\[
I\longleftrightarrow
\bigl(\mathbf 1_{\{1\in \mathrm{Fix}(I)\}},\dots,\mathbf 1_{\{n-1\in \mathrm{Fix}(I)\}}\bigr)
\]
识别。设
\[
E:\mathcal I_n\hookrightarrow \NN_{\ge 1}
\]
为任意可寻址 Gödel 编码，即在上述识别下，存在互异素数 \(p_1,\dots,p_{n-1}\) 与解码函数
\[
\delta_t:\NN\to\{0,1\}
\qquad(1\le t\le n-1)
\]
使得对每个 \(I\in\mathcal I_n\) 都有
\[
\mathbf 1_{\{t\in\mathrm{Fix}(I)\}}
=
\delta_t\!\bigl(v_{p_t}(E(I))\bigr)
\qquad(1\le t\le n-1).
\]
则平均对数复杂度满足精确下界
\[
\frac{1}{|\mathcal I_n|}
\sum_{I\in\mathcal I_n}\log E(I)
\ge
\frac12\sum_{t=1}^{n-1}\log p_t.
\]
特别地，对适当重排后的解码素数有
\[
\frac{1}{|\mathcal I_n|}
\sum_{I\in\mathcal I_n}\log E(I)
\ge
\frac12\log(n!)
=
\Omega(n\log n).
\]

另一方面，规范编码 \(\kappa'\) 精确达到相应极小值：
\[
\frac{1}{|\mathcal I_n|}
\sum_{I\in\mathcal I_n}\log \kappa'(I)
=
\frac12\sum_{t=1}^{n-1}\log q_t.
\]
因而在以 \(q_1,\dots,q_{n-1}\) 为解码素数的全部可寻址编码之中，\(\kappa'\) 的平均对数复杂度达到精确极小值。并且未归一化编码 \(\kappa\) 仅多出一个固定常数项：
\[
\frac{1}{|\mathcal I_n|}
\sum_{I\in\mathcal I_n}\log \kappa(I)
=
\log q_0+\frac12\sum_{t=1}^{n-1}\log q_t.
\]
\end{theorem}
\begin{proof}
由定理 \ref{thm:xi-chain-interior-boolean-flag-closed-form}，\(\mathcal I_n\) 与 \(\{0,1\}^{n-1}\) 规范同构，且
\[
|\mathcal I_n|=2^{n-1}.
\]
因此，题设中的 \(E\) 正是字母表 \(\Sigma=\{0,1\}\)、长度 \(T=n-1\) 的可寻址 Gödel 编码。于是定理 \ref{thm:xi-addressable-godel-product-lower-bound-classification} 在 \(q=2\) 时立即给出
\[
\frac{1}{2^{n-1}}
\sum_{I\in\mathcal I_n}\log E(I)
\ge
\frac12\sum_{t=1}^{n-1}\log p_t.
\]

将 \(p_1,\dots,p_{n-1}\) 递增重排后，不等式左端不变，而对每个 \(t\) 都有
\[
p_t\ge t+1.
\]
故
\[
\sum_{t=1}^{n-1}\log p_t
\ge
\sum_{t=1}^{n-1}\log(t+1)
=
\log(n!),
\]
从而得到
\[
\frac{1}{|\mathcal I_n|}
\sum_{I\in\mathcal I_n}\log E(I)
\ge
\frac12\log(n!)
=
\Omega(n\log n).
\]

再看 \(\kappa'\)。对每个 \(1\le i\le n-1\)，在 \(\mathcal I_n\) 的均匀分布下，事件
\[
i\in\mathrm{Fix}(I)
\]
恰有一半元素满足，因为 \(\mathcal I_n\cong \mathcal P(\{1,\dots,n-1\})\) 是布尔格。故
\[
\frac{1}{|\mathcal I_n|}
\sum_{I\in\mathcal I_n}\log \kappa'(I)
=
\frac{1}{|\mathcal I_n|}
\sum_{I\in\mathcal I_n}\sum_{i\in\mathrm{Fix}(I)\setminus\{0\}}\log q_i
=
\frac12\sum_{i=1}^{n-1}\log q_i.
\]
这与上面的普适下界在解码素数取为 \(q_1,\dots,q_{n-1}\) 时完全相等，故 \(\kappa'\) 达到平均复杂度的精确极小值。

最后，由 \(\kappa(I)=q_0\kappa'(I)\)，对两边取对数并平均即可得到
\[
\frac{1}{|\mathcal I_n|}
\sum_{I\in\mathcal I_n}\log \kappa(I)
=
\log q_0+\frac12\sum_{t=1}^{n-1}\log q_t.
\]
证毕。
\end{proof}

\noindent
于是，链上规范投影格的平方自由 Gödel 化在平均口径下也完全闭合：Dirichlet--Möbius 闭式给出其乘法像的精确算术边界，而可寻址 Gödel 理论则说明，无论最坏情形还是均匀平均，制度成本都同样被 \(n\log n\) 级增长所刚性锁定；规范编码 \(\kappa'\) 因而不仅在实现上自然，而且在平均与最坏两种复杂度口径下都达到精确极小。

在链上内算子的随机平方自由像与平均复杂度最优律都已固定之后，归一化规范编码 \(\kappa'\) 的平均位预算还可进一步收束为一条同时包含独立性、方差与最优性的综合闭式。

\begin{theorem}[链上规范平方自由 Gödel 编码的平均位预算精确最优性]\label{thm:xi-time-part61acc-chain-godel-average-variance-budget-optimality}
\leanverified{paper\_xi\_time\_part61acc\_chain\_godel\_average\_variance\_budget\_optimality}
沿用定理 \ref{thm:killo-chain-interior-godel-gcd-lcm} 与定理 \ref{thm:xi-chain-interior-boolean-flag-closed-form} 的记号，记
\[
\mathcal I_n:=\mathrm{Int}(C_n),
\qquad
C_n=\{0<1<\cdots<n-1\},
\]
并固定互异素数 \(q_0,\dots,q_{n-1}\)。对每个 \(I\in\mathcal I_n\)，定义
\[
\kappa'(I):=\frac{\kappa(I)}{q_0}
=
\prod_{i\in \mathrm{Fix}(I)\setminus\{0\}}q_i.
\]
若 \(I\) 在 \(\mathcal I_n\) 上服从均匀分布，则成立：
\begin{enumerate}
  \item 指示变量
        \[
        \mathbf 1_{\{i\in \mathrm{Fix}(I)\}}
        \qquad (1\le i\le n-1)
        \]
        相互独立，且都服从 \(\mathrm{Bernoulli}(1/2)\)；
  \item 因而
        \[
        \EE[\log \kappa'(I)]
        =
        \frac12\sum_{i=1}^{n-1}\log q_i,
        \qquad
        \mathrm{Var}(\log \kappa'(I))
        =
        \frac14\sum_{i=1}^{n-1}(\log q_i)^2;
        \]
  \item 更强地，若
        \[
        E:\mathcal I_n\hookrightarrow \NN_{\ge 1}
        \]
        是任意使用同一组解码素数 \(q_1,\dots,q_{n-1}\) 的可寻址 Gödel 编码，即对每个 \(1\le i\le n-1\) 都存在函数
        \[
        \delta_i:\NN\to\{0,1\}
        \]
        使
        \[
        \mathbf 1_{\{i\in \mathrm{Fix}(I)\}}
        =
        \delta_i\!\bigl(v_{q_i}(E(I))\bigr)
        \qquad (I\in\mathcal I_n),
        \]
        则
        \[
        \frac{1}{|\mathcal I_n|}
        \sum_{I\in\mathcal I_n}\log E(I)
        \ge
        \frac12\sum_{i=1}^{n-1}\log q_i,
        \]
        且 \(\kappa'\) 取等。
\end{enumerate}
\end{theorem}
\begin{proof}
由定理 \ref{thm:xi-chain-interior-random-godel-euler-logvariance}，在 \(\mathcal I_n\) 的均匀分布下，存在独立同分布随机变量
\[
\varepsilon_1,\dots,\varepsilon_{n-1}
\stackrel{\mathrm{i.i.d.}}{\sim}
\mathrm{Bernoulli}\!\left(\frac12\right)
\]
使得
\[
\kappa(I)=q_0\prod_{i=1}^{n-1}q_i^{\varepsilon_i},
\qquad
\varepsilon_i=\mathbf 1_{\{i\in \mathrm{Fix}(I)\}}.
\]
这立即给出第 \(1\) 条。

再由
\[
\kappa'(I)=\frac{\kappa(I)}{q_0}
=
\prod_{i=1}^{n-1}q_i^{\varepsilon_i},
\]
得到
\[
\log \kappa'(I)=\sum_{i=1}^{n-1}\varepsilon_i\log q_i.
\]
由于
\[
\EE[\varepsilon_i]=\frac12,
\qquad
\mathrm{Var}(\varepsilon_i)=\frac14,
\]
并且各 \(\varepsilon_i\) 相互独立，故
\[
\EE[\log \kappa'(I)]
=
\sum_{i=1}^{n-1}\EE[\varepsilon_i]\log q_i
=
\frac12\sum_{i=1}^{n-1}\log q_i,
\]
\[
\mathrm{Var}(\log \kappa'(I))
=
\sum_{i=1}^{n-1}\mathrm{Var}(\varepsilon_i)(\log q_i)^2
=
\frac14\sum_{i=1}^{n-1}(\log q_i)^2.
\]
第 \(2\) 条成立。

最后，由定理 \ref{thm:xi-chain-interior-godel-average-budget-optimality}，任意可寻址 Gödel 编码的平均对数复杂度都至少等于对应解码素数对数和的一半。将该定理应用于此处固定的解码素数族 \(q_1,\dots,q_{n-1}\)，便得
\[
\frac{1}{|\mathcal I_n|}
\sum_{I\in\mathcal I_n}\log E(I)
\ge
\frac12\sum_{i=1}^{n-1}\log q_i.
\]
而同一定理也已证明规范编码 \(\kappa'\) 在同一组解码素数下恰达到该下界。故第 \(3\) 条成立。证毕。
\end{proof}

\noindent
因此，链上规范平方自由 Gödel 编码的平均制度成本不仅在均匀源下具有完全显式的均值与方差，而且在固定解码素数族的全部可寻址编码中仍保持逐坐标意义下的精确最优性。


在分解刚性与素数幂单链相图之外，有限分辨率可见算术的幂零层与约化核也可由同一 Fibonacci 模环结构完全恢复。

\begin{theorem}[有限分辨率可见算术的半单核与幂零厚度由 Fibonacci 模数的素因子结构完全决定]\label{thm:xi-visible-arithmetic-nilradical-semisimple-thickness}
沿用定理 \ref{thm:fold-congruence-mod-semirings} 与定理 \ref{thm:xi-visible-arithmetic-omega-controlled-splitting} 的记号，记
\[
R_m^{\mathrm{vis}}
:=
\Omega_m/\!\equiv_m
\cong
\ZZ/(N_m\ZZ),
\qquad
N_m:=F_{m+2},
\qquad
\operatorname{rad}(N_m):=\prod_{p\mid N_m}p.
\]
定义理想
\[
\mathfrak N_m:=\operatorname{rad}\bigl(R_m^{\mathrm{vis}}\bigr)
=
\bigl(\operatorname{rad}(N_m)\bigr)/(N_m)
\subset R_m^{\mathrm{vis}}.
\]
则：
\begin{enumerate}
  \item 幂零根基恰为 \(\mathfrak N_m\)，并且
  \[
  |\mathfrak N_m|
  =
  \frac{N_m}{\operatorname{rad}(N_m)}.
  \]
  \item 约化商满足
  \[
  R_m^{\mathrm{vis}}/\mathfrak N_m
  \cong
  \ZZ/\operatorname{rad}(N_m)\ZZ
  \cong
  \prod_{p\mid N_m}\FF_p.
  \]
  \item \(\mathfrak N_m\) 的最小幂零指数为
  \[
  \nu_m
  :=
  \min\{k\ge 1:\ \mathfrak N_m^k=0\}
  =
  \max_{p^\alpha\parallel N_m}\alpha.
  \]
  \item \(R_m^{\mathrm{vis}}\) 为局部半环，当且仅当 \(N_m\) 为素数幂；\(R_m^{\mathrm{vis}}\) 为约化半环，当且仅当 \(N_m\) 平方自由。
\end{enumerate}
\end{theorem}
\begin{proof}
将 \(N_m\) 分解为
\[
N_m=\prod_{j=1}^{r}p_j^{\alpha_j}.
\]
由中国剩余定理与定理 \ref{thm:fold-congruence-mod-semirings}，
\[
R_m^{\mathrm{vis}}
\cong
\ZZ/(N_m\ZZ)
\cong
\prod_{j=1}^{r}\ZZ/(p_j^{\alpha_j}\ZZ).
\]

对每个局部因子 \(\ZZ/(p_j^{\alpha_j}\ZZ)\)，其幂零元恰为理想 \((p_j)/(p_j^{\alpha_j})\)；故全局幂零根基正是这些局部幂零理想的直积，也即
\[
\mathfrak N_m
=
\bigl(\operatorname{rad}(N_m)\bigr)/(N_m).
\]
其基数因而为
\[
|\mathfrak N_m|
=
\frac{N_m}{\operatorname{rad}(N_m)}.
\]
这证明了第 \(1\) 条。

将每个局部因子模去其 maximal ideal \((p_j)\) 后，便得到剩余域 \(\FF_{p_j}\)。于是
\[
R_m^{\mathrm{vis}}/\mathfrak N_m
\cong
\prod_{j=1}^{r}\FF_{p_j}
\cong
\ZZ/\operatorname{rad}(N_m)\ZZ,
\]
从而第 \(2\) 条成立。

再看幂零指数。对任意 \(k\ge 1\)，条件 \(\mathfrak N_m^k=0\) 等价于
\[
\operatorname{rad}(N_m)^k\in (N_m),
\]
亦即对每个 \(j\) 都有 \(k\ge \alpha_j\)。因此最小此类 \(k\) 正是
\[
\max_j\alpha_j
=
\max_{p^\alpha\parallel N_m}\alpha,
\]
故第 \(3\) 条成立。

最后，\(R_m^{\mathrm{vis}}\) 为局部半环，当且仅当 CRT 分解中只有一个局部因子，亦即 \(r=1\)，等价于 \(N_m\) 为素数幂；而 \(R_m^{\mathrm{vis}}\) 为约化半环，当且仅当幂零根基为零，等价于所有 \(\alpha_j=1\)，亦即 \(N_m\) 平方自由。第 \(4\) 条成立。证毕。
\end{proof}

\paragraph{分辨率算术几何：合同格、幂等元、幂零根与四相分类}
上一小节已经把固定分辨率 \(m\) 的折叠同余商
\[
R_m:=\Omega_m/\!\equiv_m
\]
识别为有限剩余类半环，并与 Fibonacci 模数层上的剩余类半环规范同构：
\[
R_m\cong \ZZ/(N_m\ZZ),
\qquad
N_m:=F_{m+2}.
\]
记
\[
\bar\pi_m:R_m\xrightarrow{\ \sim\ }\ZZ/(N_m\ZZ)
\]
为这一规范同构。于是，固定分辨率对象的代数性质不再依赖 \(\Omega_m\) 的具体组合呈现，而完全由 \(N_m\) 的素因子型控制。本段把这一事实组织为一层分辨率算术几何：全部粗粒化合同、投影型端同态、幂零方向、局部/分裂/约化相以及 window-\(6\) 的算术读出，均由单个 Fibonacci 模数的约数格与 Chinese remainder 结构统一决定。

\begin{theorem}[分辨率商半环的合同格与 Fibonacci 约数格]\label{thm:xi-fold-congruence-quotient-divisor-lattice}
\leanverified{paper\_xi\_fold\_congruence\_quotient\_divisor\_lattice}
对每个约数 \(d\mid N_m\)，令
\[
\rho_d:\ZZ/(N_m\ZZ)\longrightarrow \ZZ/(d\ZZ),
\qquad
\overline n\longmapsto \overline n \pmod d,
\]
并在 \(R_m\) 上定义关系
\[
x\,\Theta_d\,y
\quad\Longleftrightarrow\quad
\rho_d\!\bigl(\bar\pi_m(x)\bigr)=\rho_d\!\bigl(\bar\pi_m(y)\bigr).
\]
则 \(\Theta_d\) 是 \(R_m\) 的半环合同，且
\[
R_m/\Theta_d\cong \ZZ/(d\ZZ).
\]
此外，\(R_m\) 的任意半环合同都唯一地写成某个 \(\Theta_d\)；并且在偏序上有
\[
\Theta_{d_1}\le \Theta_{d_2}
\iff
d_2\mid d_1.
\]
换言之，
\[
\operatorname{Con}(R_m)
\cong
\operatorname{Idl}\bigl(\ZZ/(N_m\ZZ)\bigr)
\cong
\operatorname{Div}(N_m)^{\mathrm{op}}.
\]
\end{theorem}
\begin{proof}
由于 \(d\mid N_m\)，映射 \(\rho_d\) 良定义且保持加法与乘法；于是
\[
\rho_{m,d}:=\rho_d\circ \bar\pi_m:R_m\longrightarrow \ZZ/(d\ZZ)
\]
是满半环同态，而 \(\Theta_d\) 正是其核合同，从而 \(R_m/\Theta_d\cong \ZZ/(d\ZZ)\)。

现设 \(\Theta\) 为 \(R_m\) 上任意半环合同。经 \(\bar\pi_m\) 输运后，得到 \(\ZZ/(N_m\ZZ)\) 上的半环合同 \(\approx\)。由于 \(\ZZ/(N_m\ZZ)\) 是交换环，其任意半环商自动带有加法逆元，故 \(\approx\) 实际上就是环合同，并由某个理想唯一决定。另一方面，
\(\ZZ/(N_m\ZZ)\) 的全部理想恰为
\[
(d)/(N_m)
\qquad
(d\mid N_m),
\]
故 \(\Theta\) 必唯一对应某个 \(\Theta_d\)。

最后，\(\Theta_{d_1}\le \Theta_{d_2}\) 当且仅当对应理想满足
\[
(d_1)/(N_m)\subseteq (d_2)/(N_m),
\]
而这又等价于 \(d_2\mid d_1\)。于是合同格与约数格确为反同构。证毕。
\end{proof}

\begin{theorem}[分辨率商半环的非幺端同态与幺自同构刚性]\label{thm:xi-fold-congruence-unital-automorphism-rigidity}
\leanverified{paper\_xi\_fold\_congruence\_unital\_automorphism\_rigidity}
记
\[
\operatorname{End}_0(R_m)
:=
\{\phi:R_m\to R_m:\ \phi\ \text{保持 }0,+,\cdot\},
\]
\[
\operatorname{Idem}(R_m)
:=
\{e\in R_m:\ e^2=e\}.
\]
则映射
\[
\operatorname{End}_0(R_m)\longrightarrow \operatorname{Idem}(R_m),
\qquad
\phi\longmapsto \phi(1)
\]
为双射，其逆由
\[
e\longmapsto \phi_e,
\qquad
\phi_e(x):=ex
\]
给出。若
\[
N_m=\prod_{j=1}^{r}p_j^{a_j},
\qquad r=\omega(N_m),
\]
则
\[
|\operatorname{End}_0(R_m)|
=
|\operatorname{Idem}(R_m)|
=
2^r.
\]
特别地，
\[
\Aut_{\mathrm{unital\; semiring}}(R_m)=\{\mathrm{id}_{R_m}\}.
\]
\end{theorem}
\begin{proof}
经 \(\bar\pi_m\) 识别后，只需在 \(A:=\ZZ/(N_m\ZZ)\) 上证明。任取 \(\phi\in \operatorname{End}_0(A)\)，令
\[
e:=\phi(\overline 1).
\]
由乘法保持性，
\[
e^2=\phi(\overline 1)\phi(\overline 1)=\phi(\overline 1\cdot \overline 1)=\phi(\overline 1)=e,
\]
故 \(e\) 为幂等元。对任意 \(\overline k\in A\)，有
\[
\phi(\overline k)
=
\phi(k\cdot \overline 1)
=
k\,\phi(\overline 1)
=
\overline k\,e,
\]
因此 \(\phi=\phi_e\)，从而 \(\phi\mapsto \phi(\overline 1)\) 为单射。

反之，若 \(e^2=e\)，则
\[
\phi_e(\overline x+\overline y)=e(\overline x+\overline y)=e\overline x+e\overline y,
\]
\[
\phi_e(\overline x\,\overline y)=e\,\overline x\,\overline y
=
e^2\,\overline x\,\overline y
=
(e\overline x)(e\overline y)
=
\phi_e(\overline x)\phi_e(\overline y),
\]
故 \(\phi_e\) 为半环端同态。这就证明了所述双射。

再由 Chinese remainder 分解
\[
A\cong \prod_{j=1}^{r}\ZZ/(p_j^{a_j}\ZZ),
\]
每个局部因子中的幂等元只能是 \(0\) 或 \(1\)，故全体幂等元数恰为 \(2^r\)。

若 \(\phi\) 还是保持单位元的半环自同构，则 \(\phi(\overline 1)=\overline 1\)，于是对全部 \(\overline x\in A\) 都有
\[
\phi(\overline x)=\overline x\cdot \overline 1=\overline x.
\]
故 \(\phi=\mathrm{id}_A\)，传回 \(R_m\) 即得结论。证毕。
\end{proof}

\begin{theorem}[分辨率商半环的幂零根、约化核与平方自由判别]\label{thm:xi-fold-congruence-nilradical-reduced-quotient}
\leanverified{paper\_xi\_fold\_congruence\_nilradical\_reduced\_quotient}
设
\[
\operatorname{rad}(N_m):=\prod_{p\mid N_m}p,
\qquad
\mathcal N_m:=\sqrt{(0)}\subset R_m.
\]
则
\[
\mathcal N_m
=
\bigl(\operatorname{rad}(N_m)\bigr)/(N_m),
\qquad
R_m/\mathcal N_m
\cong
\ZZ/\bigl(\operatorname{rad}(N_m)\ZZ\bigr)
\cong
\prod_{p\mid N_m}\FF_p.
\]
并且 \(R_m\) 约化当且仅当 \(N_m\) 平方自由；等价地，当且仅当 \(F_{m+2}\) 不含平方因子。
\end{theorem}
\begin{proof}
这正是定理 \ref{thm:xi-visible-arithmetic-nilradical-semisimple-thickness} 在记号 \(R_m=\Omega_m/\!\equiv_m\) 下的直接改写。该定理已给出幂零根基的显式理想描述、约化商的 CRT 闭式以及“约化当且仅当平方自由”的等价判别。证毕。
\end{proof}

\begin{theorem}[分辨率算术几何的四相分类]\label{thm:xi-fold-congruence-crt-primitive-idempotents}
设
\[
N_m=\prod_{j=1}^{r}p_j^{a_j}
\]
为 \(N_m\) 的素因子分解，其中 \(p_j\) 两两不同。则
\[
R_m\cong \prod_{j=1}^{r}\ZZ/(p_j^{a_j}\ZZ),
\]
并且恰落在下列四种互斥且穷尽的算术几何相之一：
\begin{enumerate}
  \item \textbf{域相}：\(N_m=p\) 为素数，此时
  \[
  R_m\cong \FF_p.
  \]
  \item \textbf{局部非约化相}：\(N_m=p^a\) 且 \(a\ge 2\)，此时 \(R_m\) 为局部对象，但含非零幂零元。
  \item \textbf{分裂约化相}：\(r\ge 2\) 且 \(a_1=\cdots=a_r=1\)，此时
  \[
  R_m\cong \prod_{j=1}^{r}\FF_{p_j}.
  \]
  \item \textbf{分裂非约化相}：\(r\ge 2\) 且至少有一个 \(a_j\ge 2\)，此时 \(R_m\) 既发生分裂，又在至少一个局部分量中携带幂零方向。
\end{enumerate}
此外，\(R_m\) 的原始中心幂等元恰有 \(r\) 个，而全部幂等元总数恰为 \(2^r\)。
\end{theorem}
\begin{proof}
Chinese remainder theorem 给出
\[
R_m
\cong
\ZZ/(N_m\ZZ)
\cong
\prod_{j=1}^{r}\ZZ/(p_j^{a_j}\ZZ).
\]
若 \(r=1\) 且 \(a_1=1\)，则 \(R_m\cong \FF_{p_1}\)，故为域；若 \(r=1\) 且 \(a_1\ge 2\)，则 \(R_m\cong \ZZ/(p_1^{a_1}\ZZ)\) 为局部环，并含非零幂零元 \(p_1\)。这给出前两相。

若 \(r\ge 2\) 且全部 \(a_j=1\)，则所有局部因子都约化，故
\[
R_m\cong \prod_{j=1}^{r}\FF_{p_j}
\]
约化且非局部；若 \(r\ge 2\) 且某个 \(a_j\ge 2\)，则整体仍非局部，但至少一个因子含非零幂零元，从而 \(R_m\) 非约化。这给出后三相，且四种情形显然互斥并且穷尽。

最后，由定理 \ref{thm:xi-fold-congruence-unital-automorphism-rigidity}，全部幂等元数等于 \(2^r\)。而原始中心幂等元正是各坐标投影
\[
e_j:=(0,\dots,0,1,0,\dots,0)
\qquad
(1\le j\le r),
\]
故恰有 \(r\) 个。证毕。
\end{proof}
\leanverified{paper\_xi\_fold\_congruence\_crt\_primitive\_idempotents}

\begin{proposition}[window-\texorpdfstring{$6$}{6} 的算术分裂与组合刚性应严格区分]\label{prop:xi-window6-resolution-arithmetic-split-separation}
当 \(m=6\) 时，
\[
N_6=F_8=21=3\cdot 7,
\qquad
R_6\cong \ZZ/(21\ZZ)\cong \FF_3\times \FF_7.
\]
因此 \(R_6\) 恰有 \(4\) 个幂等元，并且只对应 \(2\) 个非平凡互补中央幂等元，也即 \(2\) 个非平凡直因子。另一方面，window-\(6\) 的组合层刚性分裂为
\[
X_6=X_6^{\mathrm{cyc}}\sqcup X_6^{\mathrm{bdry}},
\qquad
|X_6|=21,\quad |X_6^{\mathrm{cyc}}|=18,\quad |X_6^{\mathrm{bdry}}|=3,
\]
而二进区间折叠的纤维大小分布则为
\[
d_6^{\mathrm{bin}}(w)\in\{2,3,4\}.
\]
因而，window-\(6\) 上的 \(4\) 个幂等元只参数化投影型非幺端同态的数目，并不对应于“四块 primitive 不变块”的组合分解。
\end{proposition}
\begin{proof}
由定理 \ref{thm:xi-fold-congruence-crt-primitive-idempotents}，\(R_6\) 处于分裂约化相，并且
\[
R_6\cong \FF_3\times \FF_7.
\]
于是由定理 \ref{thm:xi-fold-congruence-unital-automorphism-rigidity} 得
\[
|\operatorname{Idem}(R_6)|=2^{\omega(21)}=4.
\]
同时，由积环的坐标分解可知，其原始中心幂等元仅为 \((1,0)\) 与 \((0,1)\) 两个，故恰有两个非平凡互补直因子；更细的单位--零因子与幂等元显式代表则由定理 \ref{thm:xi-window6-residue-semiring-unit-zerodivisor-idempotent-trichotomy} 给出。

另一方面，window-\(6\) 的组合刚性分裂
\[
X_6=X_6^{\mathrm{cyc}}\sqcup X_6^{\mathrm{bdry}},
\qquad
|X_6^{\mathrm{cyc}}|=18,\quad |X_6^{\mathrm{bdry}}|=3,
\]
已由定理 \ref{thm:xi-window6-cyclic-kernel-bc3-root-datum} 记录，而推论 \ref{cor:fold-bin-m6-fiber-hist} 表明其纤维直方图只取三档 \(2,3,4\)。因此，算术分裂、纤维大小分布与组合刚性分裂分别编码不同层级的信息，不可彼此混同。证毕。
\end{proof}
\leanverified{paper\_xi\_window6\_resolution\_arithmetic\_split\_separation}

\noindent 综上，固定分辨率层上的代数几何型完全由 Fibonacci 模数 \(N_m=F_{m+2}\) 的约数结构与指数向量决定；下一小节则把这种有限层刚性提升为跨分辨率兼容系统上的 Fibonacci-adic 逆极限环。

\paragraph{局部化整数族的统一有限化障碍与连续时钟/素数账本裂解}
在定理 \ref{thm:xi-localized-integers-padic-kernel-rigidity}、
推论 \ref{cor:xi-localized-integers-prime-axis-increment-law} 与
定理 \ref{thm:xi-localized-integers-connected-rational-blindness} 已经把
\(\ZZ[S^{-1}]\) 的连通圆维部分与全不连通 \(p\)-进核部分严格分离之后，
还可进一步抽出两条结构性后果：其一是该对象族不存在统一的有限生成忠实素数账本，
其二是连续可见时钟冻结为一维，而离散审计通道数随素数支撑无界增长。

\begin{theorem}[素数账本的统一有限忠实化不可能性]\label{thm:xi-localized-integers-no-uniform-finitely-generated-ledger}
\leanverified{paper\_xi\_localized\_integers\_no\_uniform\_finitely\_generated\_ledger}
不存在固定有限生成交换群 \(L\) 与一族映射
\[
\Phi_S:G_S=\ZZ[S^{-1}]_{>0}\longrightarrow L
\qquad
(\,|S|<\infty\,)
\]
同时满足下列条件：
\begin{enumerate}
  \item \(\Phi_S(xy)=\Phi_S(x)+\Phi_S(y)\) 对任意 \(x,y\in G_S\) 成立；
  \item 记录在 prime-ledger 意义下忠实，即随着 \(S\) 增长所新增的素数支撑轴可由记录统一且无碰撞地恢复。
\end{enumerate}
换言之，有限素数局部化对象族不存在统一的有限生成可加账本，既能保持乘法对象的可组合性，又能对素数支撑给出忠实审计。
\end{theorem}
\begin{proof}
由推论 \ref{cor:xi-localized-integers-prime-axis-increment-law}，
每加入一个新素数 \(p\)，对偶端就恰增加一条新的 \(p\)-进轴
\[
\ZZ_p,
\]
而连通商 \(\TT\) 与有限秩圆维 \(\cdim_{\mathrm{fr}}=1\) 完全不变。
因此，若记录要忠实地恢复素数支撑，
它就必须能够区分任意大的有限素数集所对应的两两独立的 prime-axis 增量。

另一方面，固定有限生成交换群 \(L\) 只具有有限自由秩与有限 torsion 支撑，
不可能承载无界增长的彼此独立的素数账本方向。
故不存在满足上述两条要求的统一协议族 \((\Phi_S)_S\)。

从另一侧看，若把 \(\Phi_S\) 限制到正整数乘法幺半群，
则得到一类从乘法对象到有限生成可消去加法账本的忠实压缩；
这与定理 \ref{thm:killo-prime-freedom-non-finitizability} 所排除的有限生成 ledgerization 障碍一致。
故统一有限忠实化不可能。证毕。
\end{proof}

\begin{corollary}[一维连续时钟与无界离散账本道的裂解]\label{cor:xi-localized-integers-oneclock-unbounded-prime-ledgers}
对每个固定有限素数集 \(S\subset\mathbb P\)，都有
\[
\cdim_{\mathrm{fr}}(G_S)=1.
\]
但任何忠实的外置审计若要精确保留 \(S\)，至少需要 \(|S|\) 条两两独立的全不连通素数账本道；
等价地，在对象族
\[
\{\,\ZZ[S^{-1}]:\ |S|<\infty\,\}
\]
上，连续可见维数冻结为 \(1\)，而离散素数账本通道数无上界。
\end{corollary}
\begin{proof}
定理 \ref{thm:xi-localized-integers-padic-kernel-rigidity} 已给出
\[
(\widehat{G_S})^0\cong \TT,
\qquad
\cdim_{\mathrm{fr}}(G_S)=1,
\qquad
K_S=\prod_{p\in S}\ZZ_p.
\]
因此连续可见部分始终只有单一圆时钟。

另一方面，\(K_S\) 的每个因子 \(\ZZ_p\) 都对应一个不可与其它素数混淆的 \(p\)-主方向；
由定理 \ref{thm:xi-localized-integers-padic-kernel-rigidity}，
\[
p\in S
\iff
\text{\(K_S\) 的 \(p\)-主因子非平凡}.
\]
故若要忠实地保留 \(S\)，至少必须为每个 \(p\in S\) 保留一条独立的离散账本道。
这给出下界 \(|S|\)；当 \(S\) 在所有有限素数集上变化时，该下界无界。
因此“一维连续时钟/无界离散账本道”的裂解严格成立。证毕。
\end{proof}

\paragraph{dyadic Fibonacci 完成层的局部分解、出生过滤与有限扇区障碍}
有限分辨率可见算术已经把每一层折叠同余商识别为单个 Fibonacci 模环，而定理 \ref{thm:xi-time-consistent-inverse-limit-readout} 与定理 \ref{thm:xi-rmin-is-zhat} 也已给出跨层兼容系统的逆极限接口。若把这两条链限制到 dyadic 子塔 \(m=2^k-2\)，则完成化对象表现出一种与通常 \(p\)-adic 增厚显著不同的局部结构：固定素数方向上的厚度在其首次出现后立即稳定，而全部无穷性仅由新素数方向的持续出生提供。

\begin{theorem}[dyadic Fibonacci 完成层的稳定局部分解]\label{thm:xi-dyadic-fibonacci-completion-stable-local-splitting}
沿用定理 \ref{thm:fold-congruence-mod-semirings} 的记号，并定义
\[
R_k:=\Omega_{2^k-2}/\!\equiv_{2^k-2}
\cong
\ZZ/(F_{2^k}\ZZ)
\qquad(k\ge 2),
\]
\[
\mathfrak R_{\mathrm{dy}}
:=
\varprojlim_{k\ge 2}R_k,
\qquad
\mathcal P_{\mathrm{dy}}
:=
\{p\in\mathbb P:\ \exists k\ge 2,\ p\mid F_{2^k}\}.
\]
对每个 \(p\in\mathcal P_{\mathrm{dy}}\)，记
\[
k_p:=\min\{k\ge 2:\ p\mid F_{2^k}\},
\qquad
e_p:=v_p(F_{2^{k_p}}).
\]
则存在与过渡映射相容的规范环同构
\[
R_k
\cong
\prod_{\substack{p\in\mathcal P_{\mathrm{dy}}\\ k_p\le k}}
\ZZ/(p^{e_p}\ZZ)
\qquad(k\ge 2),
\]
并且存在规范拓扑环同构
\[
\mathfrak R_{\mathrm{dy}}
\cong
\prod_{p\in\mathcal P_{\mathrm{dy}}}\ZZ/(p^{e_p}\ZZ).
\]
特别地，\(\mathfrak R_{\mathrm{dy}}\) 的无穷性并非来自任何固定素数方向上的无界 \(p\)-进增厚，而完全来自无穷多个新生素数方向的横向累积。
\end{theorem}
\begin{proof}
由定理 \ref{thm:fold-congruence-mod-semirings}，
\[
R_k\cong \ZZ/(F_{2^k}\ZZ).
\]
再由 Fibonacci 数的恒等式
\[
\gcd(F_a,F_b)=F_{\gcd(a,b)},
\]
可知 \(a\mid b\) 时有 \(F_a\mid F_b\)。因此 \((R_k)_{k\ge 2}\) 通过自然约化映射构成逆系统。

先注意 \(2,5\notin\mathcal P_{\mathrm{dy}}\)：\(F_n\) 为偶数当且仅当 \(3\mid n\)，而 \(5\mid F_n\) 当且仅当 \(5\mid n\)；但 \(2^k\) 既不被 \(3\) 整除，也不被 \(5\) 整除。故对任意 \(p\in\mathcal P_{\mathrm{dy}}\)，必有 \(p\neq 2,5\) 且 \(p\) 为奇素数。

记 \(\rho(p)\) 为 Fibonacci 数列模 \(p\) 的出现阶。由 Wall 的标准结果，
\[
p\mid F_n \iff \rho(p)\mid n,
\]
并且对 \(p\neq 2,5\) 有 \(p\)-进赋值公式\cite{Wall1960FibonacciModm,RowlandMedina2009PadicValuationFibonacci}
\[
v_p(F_n)=
\begin{cases}
v_p(n)+v_p(F_{\rho(p)}), & \rho(p)\mid n,\\[1mm]
0, & \rho(p)\nmid n.
\end{cases}
\]
由于 \(p\mid F_{2^{k_p}}\)，可知 \(\rho(p)\mid 2^{k_p}\)，故 \(\rho(p)=2^j\) 对某个 \(j\le k_p\)。若 \(j<k_p\)，则
\[
p\mid F_{\rho(p)}=F_{2^j},
\]
与 \(k_p\) 的极小性矛盾。因此
\[
\rho(p)=2^{k_p}.
\]
于是对任意 \(k\ge k_p\)，上述赋值公式给出
\[
v_p(F_{2^k})
=
v_p(2^{k-k_p})+v_p(F_{2^{k_p}})
=
e_p,
\]
因为 \(p\) 为奇素数，故 \(v_p(2^{k-k_p})=0\)。这说明每个素数方向在首次出现后其指数立即稳定。

对固定 \(k\)，中国剩余定理给出
\[
R_k
\cong
\prod_{p\mid F_{2^k}}\ZZ/(p^{v_p(F_{2^k})}\ZZ).
\]
而上一段已表明：若 \(p\mid F_{2^k}\)，则 \(k_p\le k\) 且 \(v_p(F_{2^k})=e_p\)；反之，若 \(k_p\le k\)，则 \(p\mid F_{2^k}\)。故
\[
R_k
\cong
\prod_{\substack{p\in\mathcal P_{\mathrm{dy}}\\ k_p\le k}}
\ZZ/(p^{e_p}\ZZ).
\]
在此识别下，过渡映射 \(R_{k+1}\to R_k\) 仅删除满足 \(k_p=k+1\) 的新生坐标，而对既有坐标保持恒等。于是映射
\[
\prod_{p\in\mathcal P_{\mathrm{dy}}}\ZZ/(p^{e_p}\ZZ)
\longrightarrow
\varprojlim_{k\ge 2}R_k,
\qquad
(a_p)_p\longmapsto
\bigl((a_p)_{k_p\le k}\bigr)_{k\ge 2}
\]
是一个拓扑环同构，证毕。
\end{proof}

\begin{corollary}[Cantor 型幂等骨架、非半局部性与非 Noether 性]\label{cor:xi-dyadic-fibonacci-completion-cantor-idempotent-skeleton}
在定理 \ref{thm:xi-dyadic-fibonacci-completion-stable-local-splitting} 的设定下，
\[
\operatorname{Idem}(\mathfrak R_{\mathrm{dy}})
\cong
\{0,1\}^{\mathcal P_{\mathrm{dy}}}.
\]
并且 \(\mathcal P_{\mathrm{dy}}\) 是可数无穷集，因此 \(\operatorname{Idem}(\mathfrak R_{\mathrm{dy}})\) 在乘积拓扑下同胚于经典 Cantor 集。特别地：
\begin{enumerate}
  \item \(\mathfrak R_{\mathrm{dy}}\) 具有 \(2^{\aleph_0}\) 个幂等元；
  \item \(\mathfrak R_{\mathrm{dy}}\) 不是半局部环；
  \item \(\mathfrak R_{\mathrm{dy}}\) 不是 Noether 环。
\end{enumerate}
\end{corollary}
\begin{proof}
记
\[
A_p:=\ZZ/(p^{e_p}\ZZ)
\qquad(p\in\mathcal P_{\mathrm{dy}}).
\]
由定理 \ref{thm:xi-dyadic-fibonacci-completion-stable-local-splitting}，
\[
\mathfrak R_{\mathrm{dy}}\cong \prod_{p\in\mathcal P_{\mathrm{dy}}}A_p.
\]
每个 \(A_p\) 都是有限局部环，故其幂等元只能是 \(0\) 或 \(1\)。因此
\[
\operatorname{Idem}(\mathfrak R_{\mathrm{dy}})
\cong
\prod_{p\in\mathcal P_{\mathrm{dy}}}\operatorname{Idem}(A_p)
\cong
\{0,1\}^{\mathcal P_{\mathrm{dy}}}.
\]

集合 \(\mathcal P_{\mathrm{dy}}\) 可数，因为它是素数集的子集。它又是无限集：当 \(k=3\) 时，\(7\mid F_8\) 且 \(7\nmid F_4\)，故已有一个新生奇素数；而对每个 \(k\ge 4\)，由推论 \ref{cor:xi-leyang-fold-infinite-mean-gamma-platform} 取 \(n=2^k\)，存在素数 \(p\) 满足 \(\rho(p)=2^k\)，于是 \(p\mid F_{2^k}\) 且 \(p\nmid F_{2^j}\)（\(j<k\)），从而 \(p\in\mathcal P_{\mathrm{dy}}\) 且 \(k_p=k\)。故 \(\mathcal P_{\mathrm{dy}}\) 为可数无穷集。任选枚举 \(\mathcal P_{\mathrm{dy}}=\{p_1,p_2,\dots\}\)，便得到
\[
\{0,1\}^{\mathcal P_{\mathrm{dy}}}\cong \{0,1\}^{\mathbb N},
\]
即 Cantor 集的标准模型。这给出枚举中的第一条。

对每个 \(p\in\mathcal P_{\mathrm{dy}}\)，理想
\[
\mathfrak m_p
:=
pA_p\times \prod_{q\in\mathcal P_{\mathrm{dy}}\setminus\{p\}}A_q
\]
的商同构于 \(\FF_p\)，故 \(\mathfrak m_p\) 是极大理想。不同 \(p\) 给出不同的极大理想，因此 \(\mathfrak R_{\mathrm{dy}}\) 至少有可数无穷多个极大理想，不能是半局部环。

最后取互异素数枚举 \(\mathcal P_{\mathrm{dy}}=\{p_1,p_2,\dots\}\)，并定义理想
\[
I_n
:=
\left(\prod_{j\le n}A_{p_j}\right)\times
\left(\prod_{j>n}0\right).
\]
则
\[
I_1\subsetneq I_2\subsetneq I_3\subsetneq \cdots,
\]
因为第 \(n+1\) 个原始坐标幂等元属于 \(I_{n+1}\setminus I_n\)。故 \(\mathfrak R_{\mathrm{dy}}\) 不是 Noether 环。证毕。
\end{proof}

\begin{theorem}[Jacobson 骨架与素数出生账本的半单恢复]\label{thm:xi-dyadic-fibonacci-completion-jacobson-semisimple-recovery}
在定理 \ref{thm:xi-dyadic-fibonacci-completion-stable-local-splitting} 的设定下，设
\[
A_p:=\ZZ/(p^{e_p}\ZZ)
\qquad(p\in\mathcal P_{\mathrm{dy}}).
\]
则
\[
J(\mathfrak R_{\mathrm{dy}})
=
\prod_{p\in\mathcal P_{\mathrm{dy}}}pA_p,
\qquad
\mathfrak R_{\mathrm{dy}}/J(\mathfrak R_{\mathrm{dy}})
\cong
\prod_{p\in\mathcal P_{\mathrm{dy}}}\FF_p.
\]
因而 \(\mathcal P_{\mathrm{dy}}\) 可由半单商本身完全恢复：
\[
p\in\mathcal P_{\mathrm{dy}}
\iff
\FF_p\ \text{作为直因子出现于}\ 
\mathfrak R_{\mathrm{dy}}/J(\mathfrak R_{\mathrm{dy}}).
\]
\end{theorem}
\begin{proof}
由定理 \ref{thm:xi-dyadic-fibonacci-completion-stable-local-splitting}，
\[
\mathfrak R_{\mathrm{dy}}\cong \prod_{p\in\mathcal P_{\mathrm{dy}}}A_p.
\]
每个 \(A_p\) 都是 Artin 局部环，其 Jacobson 根基为 \(pA_p\)，残余域为 \(\FF_p\)。

对任意环族 \((B_i)\)，有
\[
J\!\left(\prod_i B_i\right)=\prod_i J(B_i),
\]
因为 \((b_i)_i\) 在 \(\prod_i B_i\) 中为单位当且仅当每个 \(b_i\) 在 \(B_i\) 中都是单位。故
\[
J(\mathfrak R_{\mathrm{dy}})
=
\prod_{p\in\mathcal P_{\mathrm{dy}}}J(A_p)
=
\prod_{p\in\mathcal P_{\mathrm{dy}}}pA_p.
\]
再按坐标取商得到
\[
\mathfrak R_{\mathrm{dy}}/J(\mathfrak R_{\mathrm{dy}})
\cong
\prod_{p\in\mathcal P_{\mathrm{dy}}}A_p/J(A_p)
\cong
\prod_{p\in\mathcal P_{\mathrm{dy}}}\FF_p.
\]
由于不同素数对应的有限域 \(\FF_p\) 两两不同构，故 \(\FF_p\) 出现为直因子当且仅当 \(p\in\mathcal P_{\mathrm{dy}}\)。证毕。
\end{proof}

与定理 \ref{thm:xi-localized-phase-sampling-closed-form} 的局部化数系情形相比，这里的恢复机制更为刚性：对 \(\ZZ[S^{-1}]\) 而言，素数支撑 \(S\) 需经 Pontryagin 对偶短正合列中的 profinite 核 \(\prod_{p\in S}\ZZ_p\) 才能恢复；而在 dyadic Fibonacci 完成层中，半单商 \(\mathfrak R_{\mathrm{dy}}/J(\mathfrak R_{\mathrm{dy}})\) 已足以恢复全部素数出生支撑 \(\mathcal P_{\mathrm{dy}}\)。

\begin{theorem}[出生过滤与非 \texorpdfstring{$p$}{p}-adic 累厚]\label{thm:xi-dyadic-fibonacci-completion-birth-filtration}
\leanverified{paper\_xi\_dyadic\_fibonacci\_completion\_birth\_filtration}
在定理 \ref{thm:xi-dyadic-fibonacci-completion-stable-local-splitting} 的记号下，定义第 \(k\) 层的出生素数集
\[
\mathcal P_{\mathrm{dy}}^{\mathrm{birth}}(k)
:=
\{p\in\mathcal P_{\mathrm{dy}}:\ k_p=k\}.
\]
则对每个 \(n\ge 2\)，有规范环同构
\[
R_n
\cong
\prod_{2\le k\le n}\ \prod_{p\in\mathcal P_{\mathrm{dy}}^{\mathrm{birth}}(k)}
\ZZ/(p^{e_p}\ZZ),
\]
\[
\mathfrak R_{\mathrm{dy}}
\cong
\prod_{k\ge 2}\ \prod_{p\in\mathcal P_{\mathrm{dy}}^{\mathrm{birth}}(k)}
\ZZ/(p^{e_p}\ZZ).
\]
在此识别下，过渡映射 \(R_{n+1}\to R_n\) 精确等于删除全部满足 \(k_p=n+1\) 的新生坐标。并且对每个 \(k\ge 3\)，都有
\[
\mathcal P_{\mathrm{dy}}^{\mathrm{birth}}(k)\neq \varnothing.
\]
\end{theorem}
\begin{proof}
前两式只是把定理 \ref{thm:xi-dyadic-fibonacci-completion-stable-local-splitting} 中的局部分解按首次出现层级 \(k_p\) 重新分组。由于每个既有素数坐标的指数自出生之后即固定为 \(e_p\)，故过渡映射 \(R_{n+1}\to R_n\) 不会改变任何既有局部因子的指数，只会删除尚未在第 \(n\) 层出生的坐标；这正是所述的坐标删除。

对非空性，\(k=3\) 时有
\[
F_8=21=3\cdot 7,
\]
而 \(3\mid F_4\)、\(7\nmid F_4\)，故 \(7\in\mathcal P_{\mathrm{dy}}^{\mathrm{birth}}(3)\)。若 \(k\ge 4\)，则由推论 \ref{cor:xi-leyang-fold-infinite-mean-gamma-platform} 取 \(n=2^k\)，存在素数 \(p\) 使 \(\rho(p)=2^k\)。于是
\[
p\mid F_{2^k},
\qquad
p\nmid F_{2^j}\quad(j<k),
\]
故 \(p\in\mathcal P_{\mathrm{dy}}^{\mathrm{birth}}(k)\)。证毕。
\end{proof}

因此，dyadic Fibonacci 商塔的完成化并不是沿固定素数方向的 \(p\)-adic 累厚，而是一条按分辨率不断接入新生局部块的出生过滤。

\begin{theorem}[有限局部扇区上的忠实压缩不可能性]\label{thm:xi-dyadic-fibonacci-completion-no-semilocal-embedding}
不存在将 \(\mathfrak R_{\mathrm{dy}}\) 单射嵌入任何交换半局部环的环同态。特别地，不存在单射环同态
\[
\mathfrak R_{\mathrm{dy}}
\hookrightarrow
\prod_{i=1}^{r}B_i,
\]
其中每个 \(B_i\) 都是交换局部环。
\end{theorem}
\begin{proof}
由推论 \ref{cor:xi-dyadic-fibonacci-completion-cantor-idempotent-skeleton}，
\[
\left|\operatorname{Idem}(\mathfrak R_{\mathrm{dy}})\right|=2^{\aleph_0}.
\]
反设 \(f:\mathfrak R_{\mathrm{dy}}\hookrightarrow B\) 为某个交换半局部环 \(B\) 上的单射环同态。设 \(B\) 共有 \(r\) 个极大理想，则
\[
B/J(B)\cong \prod_{i=1}^{r}\kappa_i
\]
为有限个域的直积。交换环上幂等元模去 Jacobson 根基后可唯一提升，因此
\[
\operatorname{Idem}(B)\cong \operatorname{Idem}(B/J(B))
\cong
\{0,1\}^{r},
\]
从而 \(|\operatorname{Idem}(B)|=2^r<\infty\)。

另一方面，单射环同态保持不同元素不同，故不同幂等元在 \(f\) 下仍保持不同。这将迫使 \(B\) 至少含有 \(2^{\aleph_0}\) 个幂等元，与上式矛盾。故此类单射不存在。最后，有限个交换局部环的直积本身即为交换半局部环，故特别情形随即成立。证毕。
\end{proof}

\begin{conjecture}[dyadic Fibonacci 完成层的约化性]\label{conj:xi-dyadic-fibonacci-completion-reducedness}
在定理 \ref{thm:xi-dyadic-fibonacci-completion-stable-local-splitting} 的记号下，对每个 \(p\in\mathcal P_{\mathrm{dy}}\) 都有
\[
e_p=1.
\]
\end{conjecture}

\begin{proposition}[dyadic Fibonacci 完成层约化性的等价表述]\label{prop:xi-dyadic-fibonacci-completion-reduced-equivalences}
在定理 \ref{thm:xi-dyadic-fibonacci-completion-stable-local-splitting} 的记号下，下列四条等价：
\begin{enumerate}
  \item 对每个 \(p\in\mathcal P_{\mathrm{dy}}\)，都有 \(e_p=1\)；
  \item \(\mathfrak R_{\mathrm{dy}}\) 为约化环；
  \item \(J(\mathfrak R_{\mathrm{dy}})=0\)；
  \item
  \[
  \mathfrak R_{\mathrm{dy}}
  \cong
  \prod_{p\in\mathcal P_{\mathrm{dy}}}\FF_p.
  \]
\end{enumerate}
\end{proposition}
\begin{proof}
由定理 \ref{thm:xi-dyadic-fibonacci-completion-stable-local-splitting}，
\[
\mathfrak R_{\mathrm{dy}}
\cong
\prod_{p\in\mathcal P_{\mathrm{dy}}}\ZZ/(p^{e_p}\ZZ).
\]
对单个局部因子 \(A_p=\ZZ/(p^{e_p}\ZZ)\) 而言，以下条件等价：
\begin{enumerate}
  \item \(e_p=1\)；
  \item \(A_p\) 为域，亦即 \(A_p\cong \FF_p\)；
  \item \(A_p\) 为约化环；
  \item \(J(A_p)=pA_p=0\)。
\end{enumerate}
逐坐标应用这一判别便知：全部 \(e_p\) 同时等于 \(1\)，当且仅当整个直积环是约化的；也当且仅当其 Jacobson 根基
\[
J(\mathfrak R_{\mathrm{dy}})
=
\prod_{p\in\mathcal P_{\mathrm{dy}}}p\,\ZZ/(p^{e_p}\ZZ)
\]
为零；又这与直积退化为
\[
\prod_{p\in\mathcal P_{\mathrm{dy}}}\FF_p
\]
完全等价。证毕。
\end{proof}

该猜想把全部非约化障碍压缩到了 dyadic 出现阶支撑 \(\mathcal P_{\mathrm{dy}}\) 的局部指数上：若猜想成立，则 \(\mathfrak R_{\mathrm{dy}}\) 将退化为一个由可数无穷个残余域直积而成的纯残余对象，同时保留推论 \ref{cor:xi-dyadic-fibonacci-completion-cantor-idempotent-skeleton} 所揭示的 Cantor 型幂等骨架。


\paragraph{Smith 谱分类、模核 \texorpdfstring{$\zeta$}{zeta} 主因子、整数椭圆原子分解与有限时域/单寄存器账本预算}
前述有限秩账本障碍之后，Smith 信息亏损谱的局部分类、模核计数的 Dirichlet 全局化、整数轴对齐椭圆的原子结构，以及有限时域与单寄存器口径下自然扩张的极小二进制外置预算还可继续压缩为下列结构性结论。

\begin{theorem}[\texorpdfstring{$p$}{p}-进信息亏损谱的完全分类]\label{thm:xi-smith-padic-loss-spectrum-classification}
设 \(f:\ZZ_{\ge 0}\to \ZZ_{\ge 0}\) 满足 \(f(0)=0\)，并对 \(k\ge 1\) 定义前向差分
\[
\Delta(k):=f(k)-f(k-1).
\]
则下列三条等价：
\begin{enumerate}
  \item 存在某个满行秩整数矩阵 \(A\in M_{m\times n}(\ZZ)\) 与素数 \(p\)，使得 \(f=f_p\)，其中 \(f_p\) 为定理 \ref{thm:killo-smith-loss-spectrum} 的 \(\,p\)-进信息亏损谱；
  \item 序列 \((\Delta(k))_{k\ge 1}\) 是非负、非增、有限支撑的整数列；
  \item \(f\) 是整数值、单调不减、离散凹并最终常数的函数。
\end{enumerate}
在这些等价条件成立时，约定对充分大的 \(k\) 有 \(\Delta(k)=0\)，并定义
\[
m_e:=\Delta(e)-\Delta(e+1)\qquad(e\ge 1).
\]
则 \(m_e\in\ZZ_{\ge 0}\)，且只有有限多个 \(m_e\) 非零，并且
\[
f(k)=\sum_{e\ge 1}m_e\min(k,e)\qquad(k\ge 0).
\]
该表示唯一。特别地，\(f\) 所对应的正 \(\,p\)-进 Smith 指数多重集由“指数 \(e\) 重复 \(m_e\) 次”唯一确定。
\leanverified{paper\_xi\_smith\_padic\_loss\_spectrum\_classification}
\end{theorem}
\begin{proof}
若 \(f=f_p\)，则定理 \ref{thm:killo-smith-loss-spectrum} 给出
\[
f(k)=\sum_{i=1}^{m}\min(k,e_i(p)),
\qquad
\Delta(k)=\#\{i:e_i(p)\ge k\}.
\]
对每个固定 \(e\)，函数 \(k\mapsto \min(k,e)\) 单调不减、离散凹且最终常数；其有限和 \(f\) 亦然。并且 \(\Delta(k)\) 显然是非负、非增且有限支撑的整数列，故第 \(1\) 条推出第 \(2\)、\(3\) 条。

第 \(2\) 条与第 \(3\) 条等价只需注意：整数值函数的离散凹性恰等价于前向差分单调不增；单调不减恰等价于前向差分非负；最终常数恰等价于前向差分最终为零。

现证第 \(2\) 条推出第 \(1\) 条。由 \(\Delta\) 非增且有限支撑，
\[
m_e=\Delta(e)-\Delta(e+1)\ge 0,
\]
并且只有有限多个 \(m_e\) 非零。对任意 \(k\ge 1\)，望远镜求和给出
\[
\sum_{e\ge k}m_e
=
\sum_{e\ge k}\bigl(\Delta(e)-\Delta(e+1)\bigr)
=
\Delta(k).
\]
于是
\[
f(k)
=
\sum_{j=1}^{k}\Delta(j)
=
\sum_{j=1}^{k}\sum_{e\ge j}m_e
=
\sum_{e\ge 1}m_e\sum_{j=1}^{k}\mathbf 1_{\{e\ge j\}}
=
\sum_{e\ge 1}m_e\min(k,e).
\]
若所有 \(m_e\) 都为零，则 \(f\equiv 0\)，取 \(A=(1)\) 即有 \(f=f_p\equiv 0\)。否则记
\[
r:=\sum_{e\ge 1}m_e\ge 1,
\]
并把每个整数 \(e\) 按重数 \(m_e\) 排成序列 \(e_1,\dots,e_r\)。固定任意素数 \(p\)，取对角矩阵
\[
A=\operatorname{diag}(p^{e_1},\dots,p^{e_r})\in M_r(\ZZ).
\]
由定理 \ref{thm:killo-smith-loss-spectrum}，
\[
f_p(k)=\sum_{\nu=1}^{r}\min(k,e_\nu)=\sum_{e\ge 1}m_e\min(k,e)=f(k),
\]
故第 \(1\) 条成立。

最后，唯一性由
\[
\Delta(k)=\sum_{e\ge k}m_e
\]
立即推出
\[
m_e=\Delta(e)-\Delta(e+1).
\]
因此正 \(\,p\)-进 Smith 指数多重集也随之唯一确定。证毕。
\end{proof}

\begin{corollary}[Smith 指数的二阶离散曲率恢复律]\label{cor:xi-smith-second-discrete-curvature-recovery}
\leanverified{paper\_xi\_smith\_second\_discrete\_curvature\_recovery}
沿用定理 \ref{thm:killo-smith-loss-spectrum} 的记号，并定义
\[
f_p(0):=0,
\qquad
\Delta_p(k):=f_p(k)-f_p(k-1)\qquad(k\ge 1).
\]
则对每个 \(k\ge 1\) 都有
\[
\#\{i:e_i(p)=k\}
=
\Delta_p(k)-\Delta_p(k+1)
=
2f_p(k)-f_p(k-1)-f_p(k+1).
\]
因此下列三个局部不变量均可由 \(f_p\) 无损恢复：
\[
r_p:=\#\{i:e_i(p)\ge 1\}=f_p(1),
\]
\[
\tau_p:=\max\{k:\Delta_p(k)>0\}=\max_i e_i(p),
\]
\[
\sigma_p:=\lim_{k\to\infty}f_p(k)=\sum_{i=1}^{m}e_i(p)=v_p\!\Bigl(\prod_{i=1}^{m}d_i\Bigr).
\]
若进一步 \(A\) 为方阵满秩，则对一切 \(k\ge \tau_p\) 都有
\[
\#\ker(A\bmod p^k)=p^{\sigma_p},
\]
即模 \(p^k\) 的核大小在阈值 \(\tau_p\) 之后精确稳定为 \(|\det A|\) 的 \(\,p\)-进部分。
\end{corollary}
\begin{proof}
定理 \ref{thm:killo-smith-loss-spectrum} 给出
\[
\Delta_p(k)=\#\{i:e_i(p)\ge k\}.
\]
于是
\[
\Delta_p(k)-\Delta_p(k+1)
=
\#\{i:e_i(p)\ge k\}-\#\{i:e_i(p)\ge k+1\}
=
\#\{i:e_i(p)=k\}.
\]
再由
\[
\Delta_p(k)-\Delta_p(k+1)
=
\bigl(f_p(k)-f_p(k-1)\bigr)-\bigl(f_p(k+1)-f_p(k)\bigr),
\]
即得二阶离散曲率公式。

三个局部不变量随即成立。第一式取 \(k=1\) 即得；第二式来自 \(\Delta_p(k)>0\iff \exists\,i,\ e_i(p)\ge k\)；第三式因为 \(f_p(k)=\sum_i\min(k,e_i(p))\) 最终稳定于 \(\sum_i e_i(p)\)。

若 \(A\) 为方阵满秩，则 \(m=n\)。对 \(k\ge \tau_p\)，一切 \(i\) 都满足 \(\min(k,e_i(p))=e_i(p)\)，故
\[
f_p(k)=\sum_{i=1}^{n}e_i(p)=\sigma_p.
\]
由定理 \ref{thm:killo-smith-loss-spectrum} 中
\[
K_p(k)=\#\ker(A\bmod p^k)=p^{\,k(n-m)+f_p(k)},
\]
便得
\[
\#\ker(A\bmod p^k)=p^{\sigma_p}.
\]
又因 \(\prod_{i=1}^{n}d_i=|\det A|\)，该稳定值正是 \(|\det A|\) 的 \(\,p\)-进部分。证毕。
\end{proof}

\begin{theorem}[Smith 不变量的有限 \texorpdfstring{$p$}{p}-窗完全恢复与顶层行列式指数]\label{thm:xi-smith-finite-pwindow-complete-recovery-determinant-index}
设 \(A\in M_{m\times n}(\ZZ)\) 满足 \(\operatorname{rank}_{\QQ}(A)=m\)，其 Smith 不变量为
\[
d_1\mid d_2\mid\cdots\mid d_m.
\]
固定素数 \(p\)。对 \(k\ge 1\) 定义
\[
K_p(k):=\#\ker(A\bmod p^k),
\qquad
f_p(k):=\log_p K_p(k)-k(n-m),
\]
并记
\[
e_i(p):=v_p(d_i),
\qquad
E_p:=\max_{1\le i\le m}e_i(p).
\]
则成立：
\begin{enumerate}
  \item 有限窗口 \(\bigl(K_p(k)\bigr)_{1\le k\le E_p}\) 已充要决定 \(p\)-初级 Smith 多重集 \(\{e_i(p)\}_{i=1}^{m}\)；对一切 \(k>E_p\) 的数据均为冗余。
  \item 极限
  \[
  \nu_p(A):=\lim_{k\to\infty}f_p(k)
  \]
  存在，且满足
  \[
  \nu_p(A)
  =
  \sum_{i=1}^{m}e_i(p)
  =
  v_p\!\Bigl(\prod_{i=1}^{m}d_i\Bigr).
  \]
  \item 若 \(A\) 为方阵且 \(\det A\neq 0\)，则
  \[
  \nu_p(A)=v_p(\det A).
  \]
\end{enumerate}
\end{theorem}
\begin{proof}
由定理 \ref{thm:killo-smith-loss-spectrum}，
\[
f_p(k)=\sum_{i=1}^{m}\min\bigl(k,e_i(p)\bigr)
\qquad(k\ge 1),
\]
并且
\[
\Delta_p(k):=f_p(k)-f_p(k-1)=\#\{i:\ e_i(p)\ge k\}.
\]
若 \(k>E_p\)，则对所有 \(i\) 都有 \(e_i(p)<k\)，故 \(\Delta_p(k)=0\)。因此 \(f_p\) 在 \(k\ge E_p\) 后恒定，全部 \(k>E_p\) 的数据均不再携带新的信息。另一方面，由
\[
K_p(k)=p^{\,k(n-m)+f_p(k)}
\]
可由窗口 \(\bigl(K_p(k)\bigr)_{1\le k\le E_p}\) 唯一恢复 \(\bigl(f_p(k)\bigr)_{1\le k\le E_p}\)，进而恢复
\[
\Delta_p(k)=f_p(k)-f_p(k-1)
\qquad(1\le k\le E_p),
\qquad
\Delta_p(k)=0
\quad(k>E_p).
\]
于是对每个 \(\ell\ge 1\)，都有
\[
\#\{i:\ e_i(p)=\ell\}
=
\Delta_p(\ell)-\Delta_p(\ell+1),
\]
故整个多重集 \(\{e_i(p)\}_{i=1}^{m}\) 被唯一恢复。

反之，若仅知某个长度 \(L<E_p\) 的初始窗口，则无法唯一确定该多重集。事实上，取某个满足 \(e_j(p)=E_p\) 的指标 \(j\)，把它替换为 \(e_j'(p):=E_p+1\)，其余赋值保持不变。则对所有 \(k\le L\) 都有
\[
\min(k,E_p)=\min(k,E_p+1)=k,
\]
故修改前后的 \(f_p(k)\) 与 \(K_p(k)\) 完全一致，但对应的 \(p\)-初级 Smith 多重集不同。于是长度短于 \(E_p\) 的窗口不可能给出无损恢复，第 \(1\) 条成立。

第 \(2\) 条由平台稳定性立即得到：当 \(k\ge E_p\) 时，
\[
f_p(k)=\sum_{i=1}^{m}e_i(p),
\]
故极限存在并等于该和值。又因 \(\prod_i d_i\) 的 \(p\)-赋值正是 \(\sum_i e_i(p)\)，故
\[
\nu_p(A)=v_p\!\Bigl(\prod_{i=1}^{m}d_i\Bigr).
\]

若 \(A\) 为方阵且 \(\det A\neq 0\)，则 Smith 乘积满足
\[
\prod_{i=1}^{m}d_i=|\det A|,
\]
故
\[
\nu_p(A)=v_p(|\det A|)=v_p(\det A).
\]
第 \(3\) 条成立。证毕。
\end{proof}
\leanverified{paper\_xi\_smith\_finite\_pwindow\_complete\_recovery\_determinant\_index}


\begin{theorem}[Smith 缺损谱的共轭分拆恢复式与完全实现]\label{thm:xi-smith-loss-spectrum-conjugate-partition-realization}
\leanverified{paper\_xi\_smith\_loss\_spectrum\_conjugate\_partition\_realization}
在定理 \ref{thm:xi-smith-padic-loss-spectrum-classification} 的设定下，进一步固定行秩为 \(m\)，并设
\[
\Delta(k):=f(k)-f(k-1)\in\{0,1,\dots,m\}
\qquad(k\ge 1)
\]
为非增且最终为零的整数列。对 \(1\le i\le m\)，定义
\[
e_i:=\#\{k\ge 1:\ \Delta(k)\ge i\}.
\]
则
\[
e_1\ge e_2\ge \cdots\ge e_m\ge 0,
\]
并且对一切 \(k\ge 0\) 都有
\[
\boxed{\;
f(k)=\sum_{i=1}^{m}\min(k,e_i).\;}
\]
进一步，对每个 \(r\ge 1\) 都有显式恢复公式
\[
\boxed{\;
\#\{i:\ e_i=r\}
=
\Delta(r)-\Delta(r+1)
=
2f(r)-f(r-1)-f(r+1).\;}
\]
而零指数的重数满足
\[
\boxed{\;
\#\{i:\ e_i=0\}=m-\Delta(1).\;}
\]
反过来，任意满足上述条件的 \(f\) 都可由某个满行秩整数矩阵在某个素数 \(p\) 上的 Smith 缺损谱实现。因而，\(p\)-进核计数序列与长度恰为 \(m\) 的 Smith 指数分拆之间存在严格的共轭分拆双向等价。
\end{theorem}
\begin{proof}
由定理 \ref{thm:xi-smith-padic-loss-spectrum-classification}，此类 \(f\) 唯一对应于某个正指数多重集 \(\{e>0\}\) 的重数函数
\[
m_e=\Delta(e)-\Delta(e+1).
\]
现把零指数一并补足到长度 \(m\)。由定义，序列 \((e_i)_{i=1}^{m}\) 显然非增且非负；并且对任意 \(j\ge 1\)，有
\[
\#\{i:\ e_i\ge j\}=\Delta(j).
\]
因此对任意 \(k\ge 0\)，应用共轭分拆恒等式得到
\[
\sum_{i=1}^{m}\min(k,e_i)
=
\sum_{j=1}^{k}\#\{i:\ e_i\ge j\}
=
\sum_{j=1}^{k}\Delta(j)
=
f(k).
\]
这就得到第一条公式。

对 \(r\ge 1\)，指数恰等于 \(r\) 的重数为
\[
\#\{i:\ e_i=r\}
=
\#\{i:\ e_i\ge r\}-\#\{i:\ e_i\ge r+1\}
=
\Delta(r)-\Delta(r+1).
\]
再利用 \(\Delta(r)=f(r)-f(r-1)\) 即得
\[
\Delta(r)-\Delta(r+1)=2f(r)-f(r-1)-f(r+1).
\]
零指数的重数则由总长度减去正指数个数给出：
\[
\#\{i:\ e_i=0\}=m-\#\{i:\ e_i\ge 1\}=m-\Delta(1).
\]

最后给出实现。任选素数 \(p\)，取对角矩阵
\[
A=\operatorname{diag}(p^{e_1},\dots,p^{e_m})\in M_m(\ZZ).
\]
该矩阵满秩，且其 \(p\)-进 Smith 指数恰为 \((e_1,\dots,e_m)\)。于是由定理 \ref{thm:killo-smith-loss-spectrum}，
\[
f_p(k)=\sum_{i=1}^{m}\min(k,e_i)=f(k).
\]
故所求实现成立，且由上式与各恢复公式知其唯一性亦成立。证毕。
\end{proof}

\begin{corollary}[Smith 体积由 \texorpdfstring{$p$}{p}-进亏损谱的平台值精确恢复]\label{cor:xi-smith-volume-recovered-from-padic-platform}
设 \(A\in M_{m\times n}(\ZZ)\) 满足 \(\operatorname{rank}_{\QQ}(A)=m\)，其 Smith 不变量为
\[
d_1\mid d_2\mid\cdots\mid d_m.
\]
定义 Smith 体积
\[
\mathfrak D(A):=\prod_{i=1}^{m}d_i.
\]
则对每个素数 \(p\) 都有
\[
\boxed{\;
v_p(\mathfrak D(A))
=
\lim_{k\to\infty}f_p(k).\;}
\]
从而
\[
\boxed{\;
\mathfrak D(A)=\prod_{p}p^{\,\lim_{k\to\infty}f_p(k)}.\;}
\]
\end{corollary}
\begin{proof}
由推论 \ref{cor:xi-smith-second-discrete-curvature-recovery}，
\[
\lim_{k\to\infty}f_p(k)=\sum_{i=1}^{m}e_i(p)
=
v_p\!\Bigl(\prod_{i=1}^{m}d_i\Bigr)
=
v_p(\mathfrak D(A)).
\]
由于仅有有限多个素数 \(p\) 对 \(\mathfrak D(A)\) 具有正赋值，把各素数上的赋值重新拼合便得到第二式。证毕。
\end{proof}


\begin{theorem}[模核 Dirichlet 生成函数的 \texorpdfstring{$\zeta$}{zeta} 主因子与有限 Euler 修正]\label{thm:xi-kernel-dirichlet-zeta-finite-euler-correction}
设 \(A\in M_{m\times n}(\ZZ)\) 满足 \(\operatorname{rank}_{\QQ}(A)=m\)，其 Smith 不变量为
\[
d_1\mid d_2\mid\cdots\mid d_m.
\]
对每个 \(b\ge 1\)，记
\[
A_b:(\ZZ/b\ZZ)^n\longrightarrow(\ZZ/b\ZZ)^m
\]
为模 \(b\) 的诱导映射，并定义归一化模核计数
\[
\kappa_A(b):=\frac{\#\ker(A_b)}{b^{\,n-m}}
=
\prod_{i=1}^{m}\gcd(d_i,b).
\]
再定义 Dirichlet 生成函数
\[
Z_A(s):=\sum_{b\ge 1}\frac{\kappa_A(b)}{b^s},
\qquad
\Re s>1.
\]
令
\[
\Delta_A:=\prod_{i=1}^{m}d_i,
\]
并对每个素数 \(p\mid \Delta_A\) 记
\[
e_i(p):=v_p(d_i),\qquad
E_p:=\max_i e_i(p),\qquad
\sigma_p:=\sum_{i=1}^{m}e_i(p),
\]
\[
f_p(k):=\sum_{i=1}^{m}\min(k,e_i(p)).
\]
则成立精确 Euler 分解
\[
Z_A(s)=\zeta(s)\prod_{p\mid \Delta_A}P_{A,p}(p^{-s}),
\]
其中
\[
P_{A,p}(T)
:=
(1-T)\sum_{k=0}^{E_p-1}p^{f_p(k)}T^k
+p^{\sigma_p}T^{E_p}
\in\ZZ[T].
\]
特别地，\(Z_A(s)\) 在复平面上亚纯延拓，仅在 \(s=1\) 具有一个单极点，并且
\[
\operatorname*{Res}_{s=1}Z_A(s)
=
\prod_{p\mid \Delta_A}P_{A,p}(p^{-1}).
\]
\end{theorem}
\begin{proof}
推论 \ref{cor:killo-kernel-size-general-modulus} 已给出
\[
\kappa_A(b)=\prod_{i=1}^{m}\gcd(d_i,b).
\]
若 \((b_1,b_2)=1\)，则
\[
\gcd(d_i,b_1b_2)=\gcd(d_i,b_1)\gcd(d_i,b_2),
\]
从而 \(\kappa_A\) 为乘法函数。故在 \(\Re s>1\) 上，
\[
Z_A(s)=\prod_{p}Z_{A,p}(s),
\qquad
Z_{A,p}(s):=\sum_{k\ge 0}\frac{\kappa_A(p^k)}{p^{ks}}.
\]

再由同一推论与定理 \ref{thm:killo-smith-loss-spectrum}，
\[
\kappa_A(p^k)=\prod_{i=1}^{m}p^{\min(k,e_i(p))}=p^{f_p(k)}.
\]
若 \(k\ge E_p\)，则 \(\min(k,e_i(p))=e_i(p)\) 对一切 \(i\) 都成立，故 \(f_p(k)=\sigma_p\)。记 \(T:=p^{-s}\)，则
\[
Z_{A,p}(s)
=
\sum_{k=0}^{E_p-1}p^{f_p(k)}T^k+\sum_{k\ge E_p}p^{\sigma_p}T^k
=
\sum_{k=0}^{E_p-1}p^{f_p(k)}T^k+\frac{p^{\sigma_p}T^{E_p}}{1-T}.
\]
整理可得
\[
Z_{A,p}(s)=\frac{P_{A,p}(T)}{1-T}.
\]

当 \(p\nmid \Delta_A\) 时，一切 \(e_i(p)=0\)，于是 \(\kappa_A(p^k)=1\) 对所有 \(k\) 都成立，故
\[
Z_{A,p}(s)=\sum_{k\ge 0}p^{-ks}=\frac{1}{1-p^{-s}}.
\]
将这些素数上的因子合并，便得到
\[
Z_A(s)
=
\prod_{p\nmid \Delta_A}\frac{1}{1-p^{-s}}
\prod_{p\mid \Delta_A}\frac{P_{A,p}(p^{-s})}{1-p^{-s}}
=
\zeta(s)\prod_{p\mid \Delta_A}P_{A,p}(p^{-s}).
\]
右端是 \(\zeta(s)\) 乘以有限个整函数因子，因此仅在 \(s=1\) 保留 \(\zeta(s)\) 的单极点。又 \(\operatorname*{Res}_{s=1}\zeta(s)=1\)，故
\[
\operatorname*{Res}_{s=1}Z_A(s)
=
\prod_{p\mid \Delta_A}P_{A,p}(p^{-1}).
\]
证毕。
\end{proof}

\begin{theorem}[整数轴对齐椭圆半群的自由交换原子分解]\label{thm:xi-integer-ellipse-free-commutative-monoid}
沿用定理 \ref{thm:killo-ellipse-diagonal-prime-register-equivalence} 的记号，记
\[
\mathcal A:=\{E_{p,1},E_{1,p}:p\in\mathcal P\}\subset \mathsf E_{\NN}.
\]
则：
\begin{enumerate}
  \item \((\mathsf E_{\NN},\odot)\) 是以 \(\mathcal A\) 为原子集的自由交换幺半群；
  \item 若约定 \(M^{\odot r}\) 表示 \(r\) 次 \(\odot\)-乘积，并取 \(M^{\odot 0}=E_{1,1}\)，则每个 \(E_{a,b}\in\mathsf E_{\NN}\) 都具有唯一原子分解
  \[
  E_{a,b}
  =
  \left(\bigodot_{p\in\mathcal P}E_{p,1}^{\odot v_p(a)}\right)
  \odot
  \left(\bigodot_{p\in\mathcal P}E_{1,p}^{\odot v_p(b)}\right),
  \]
  其中仅有限多个因子不同于单位元 \(E_{1,1}\)；
  \item 其幺半群自同构群满足
  \[
  \operatorname{Aut}_{\mathrm{Mon}}(\mathsf E_{\NN},\odot)
  \cong
  \mathrm{Sym}(\mathcal A)
  \cong
  \mathrm{Sym}(\mathcal P\times\{x,y\}).
  \]
\end{enumerate}
\end{theorem}
\begin{proof}
定理 \ref{thm:killo-ellipse-diagonal-prime-register-equivalence} 已给出幺半群同构
\[
(\NN^{(\mathcal P)}\times \NN^{(\mathcal P)},+)
\cong
(\mathsf E_{\NN},\odot),\qquad
(r_1,r_2)\longmapsto E_{N(r_1),N(r_2)}.
\]
设 \(\delta_p\in\NN^{(\mathcal P)}\) 为在素数 \(p\) 处取值 \(1\)、其余处取值 \(0\) 的基向量。则
\[
(\delta_p,0)\longmapsto E_{p,1},
\qquad
(0,\delta_p)\longmapsto E_{1,p}.
\]
而 \((\NN^{(\mathcal P)}\times \NN^{(\mathcal P)},+)\) 正是以这些基向量为自由基的交换幺半群，因此其原子（不可约元）恰为
\[
\{(\delta_p,0),(0,\delta_p):p\in\mathcal P\}.
\]
经同构传回，便得到第 \(1\) 条。

任取
\[
a=\prod_{p\in\mathcal P}p^{v_p(a)},
\qquad
b=\prod_{p\in\mathcal P}p^{v_p(b)},
\]
则在自由基坐标下，
\[
(r_1,r_2)=\sum_{p}v_p(a)(\delta_p,0)+\sum_{p}v_p(b)(0,\delta_p).
\]
把该式送回 \((\mathsf E_{\NN},\odot)\) 便得到所示原子分解。由于自由交换幺半群中的基坐标唯一，故该分解唯一，第 \(2\) 条成立。

最后，幺半群自同构保持不可约元，因此必将原子集 \(\mathcal A\) 置换到自身。反之，任一 \(\mathcal A\) 上的置换都可在自由交换幺半群上按乘法唯一延拓为自同构。于是
\[
\operatorname{Aut}_{\mathrm{Mon}}(\mathsf E_{\NN},\odot)
\cong
\mathrm{Sym}(\mathcal A).
\]
再由
\[
\mathcal A=\{E_{p,1},E_{1,p}:p\in\mathcal P\}
\cong
\mathcal P\times\{x,y\},
\]
即得最后一个同构。证毕。
\end{proof}

\begin{corollary}[保持面积的整数轴对齐椭圆幺半群自同构仅由各素数轴交换组成]\label{cor:xi-integer-ellipse-area-preserving-automorphism-product}
沿用定理 \ref{thm:xi-integer-ellipse-free-commutative-monoid} 的记号，并定义面积函数
\[
\operatorname{Area}(E_{a,b}):=\pi ab.
\]
记
\[
\operatorname{Aut}_{\mathrm{Mon},\operatorname{Area}}(\mathsf E_{\NN},\odot)
:=
\left\{
\Phi\in \operatorname{Aut}_{\mathrm{Mon}}(\mathsf E_{\NN},\odot):
\operatorname{Area}(\Phi(E))=\operatorname{Area}(E)\ \forall E\in\mathsf E_{\NN}
\right\}.
\]
则有规范同构
\[
\operatorname{Aut}_{\mathrm{Mon},\operatorname{Area}}(\mathsf E_{\NN},\odot)
\cong
\prod_{p\in\mathcal P}(\ZZ/2\ZZ),
\]
其中第 \(p\) 个因子由两原子
\[
E_{p,1}\longleftrightarrow E_{1,p}
\]
的交换生成。
\end{corollary}
\begin{proof}
由定理 \ref{thm:xi-integer-ellipse-free-commutative-monoid}，任意幺半群自同构都唯一对应于原子集
\[
\mathcal A=\{E_{p,1},E_{1,p}:p\in\mathcal P\}
\]
上的一个置换。若 \(\Phi\) 还保持面积逐点不变，则对每个素数 \(p\) 有
\[
\operatorname{Area}(E_{p,1})=\operatorname{Area}(E_{1,p})=\pi p.
\]
而当 \(q\neq p\) 时，
\[
\operatorname{Area}(E_{q,1})=\operatorname{Area}(E_{1,q})=\pi q\neq \pi p.
\]
故 \(\Phi(E_{p,1})\) 只能取 \(E_{p,1}\) 或 \(E_{1,p}\)，同理 \(\Phi(E_{1,p})\) 亦只能取该二元组中的另一项。于是 \(\Phi\) 必逐素数保持二点块
\[
\{E_{p,1},E_{1,p}\}
\]
并在每个块上仅有“恒等”与“交换”两种可能。

反之，对每个 \(p\) 独立选择是否交换 \(E_{p,1}\) 与 \(E_{1,p}\)，便得到原子集 \(\mathcal A\) 上的一个置换；由定理 \ref{thm:xi-integer-ellipse-free-commutative-monoid}，该置换唯一延拓为 \((\mathsf E_{\NN},\odot)\) 的幺半群自同构。由于这种延拓保持每个素数对 \(\{E_{p,1},E_{1,p}\}\) 的面积 \(\pi p\)，故它逐点保持全部椭圆的面积。
更具体地，若
\[
E_{a,b}
=
\left(\bigodot_{p}E_{p,1}^{\odot \alpha_p}\right)
\odot
\left(\bigodot_{p}E_{1,p}^{\odot \beta_p}\right),
\]
则任意这类自同构只会把每个固定 \(p\) 上的 \(\alpha_p,\beta_p\) 交换或保持不变，因此
\[
v_p(ab)=\alpha_p+\beta_p
\]
逐素数保持不变，从而
\[
\operatorname{Area}(E_{a,b})=\pi ab
\]
逐点保持不变。

因此，保持面积的自同构群恰为各素数上独立交换所成的直积
\[
\prod_{p\in\mathcal P}(\ZZ/2\ZZ).
\]
证毕。
\end{proof}

\begin{theorem}[有限时域自然扩张的最小二进制外置预算]\label{thm:xi-finite-time-natural-extension-active-branch-budget}
沿用定理 \ref{thm:xi-layered-prime-register-sharp-natural-extension} 的分层外置语义，但把第 \(j\) 层的 prime-slice 通道替换为有限字母表 \(A_j\)。设
\[
X_0\xrightarrow{\pi_0}X_1\xrightarrow{\pi_1}\cdots\xrightarrow{\pi_{N-1}}X_N
\]
为有限时域链，并假设每层纤维满足
\[
\#\pi_j^{-1}(y)\le 2
\qquad
(\forall\,y\in X_{j+1},\ 0\le j\le N-1).
\]
定义活跃分支层集合
\[
B_N:=\{\,j\in\{0,\dots,N-1\}:\exists\,y\in X_{j+1}\text{ 使 }\#\pi_j^{-1}(y)=2\,\}.
\]
考虑任意分层外置编码
\[
\Psi(x_0)=\bigl(x_N,h_0(x_0),\dots,h_{N-1}(x_0)\bigr)\in X_N\times\prod_{j=0}^{N-1}A_j
\]
并假设 \(\Psi\) 为单射。则对每个 \(j\in B_N\) 都必有
\[
|A_j|\ge 2.
\]
因此总二进制预算满足
\[
\sum_{j=0}^{N-1}\left\lceil \log_2|A_j|\right\rceil\ge |B_N|.
\]
反之，存在单射分层外置使得
\[
|A_j|=
\begin{cases}
2,& j\in B_N,\\
1,& j\notin B_N.
\end{cases}
\]
故最小总二进制外置预算恰等于 \(|B_N|\)。
\end{theorem}
\begin{proof}
对任意活跃层 \(j\in B_N\)，存在 \(y\in X_{j+1}\) 使 \(\#\pi_j^{-1}(y)=2\)。推论 \ref{cor:killo-layered-necessity-minimality} 的论证对该单层可逐字局部化：若第 \(j\) 层通道退化为单值，则可以构造两条具有相同终端宏态、仅在第 \(j\) 层分支选择不同的历史，使得分层外置无法保持单射。故每个活跃层都必须配备非平凡字母表，亦即 \(|A_j|\ge 2\)。对所有 \(j\in B_N\) 求和便得
\[
\sum_{j=0}^{N-1}\left\lceil \log_2|A_j|\right\rceil\ge |B_N|.
\]

反向上，对每个 \(j\in B_N\) 与每个 \(y\in X_{j+1}\)，由于 \(\#\pi_j^{-1}(y)\le 2\)，可取单射
\[
\iota_{j,y}:\pi_j^{-1}(y)\hookrightarrow \{0,1\}.
\]
若 \(j\notin B_N\)，则每条纤维都是单点集，可唯一地取
\[
\iota_{j,y}:\pi_j^{-1}(y)\to \{0\}.
\]
对任意 \(x_0\in X_0\)，递推定义
\[
x_{j+1}:=\pi_j(x_j)\qquad(0\le j\le N-1),
\]
并设
\[
h_j(x_0):=\iota_{j,x_{j+1}}(x_j)\in A_j,
\qquad
A_j:=
\begin{cases}
\{0,1\},& j\in B_N,\\
\{0\},& j\notin B_N.
\end{cases}
\]
于是
\[
\Psi(x_0)=\bigl(x_N,h_0(x_0),\dots,h_{N-1}(x_0)\bigr).
\]
给定 \(\Psi(x_0)\)，可由 \(x_N\) 与 \(h_{N-1}\) 先恢复唯一的 \(x_{N-1}\)，再由 \(x_{N-1}\) 与 \(h_{N-2}\) 恢复唯一的 \(x_{N-2}\)，如此逐层逆推，最终恢复唯一的 \(x_0\)。故 \(\Psi\) 单射。

该构造在每个活跃层恰使用一位二进制字母，在每个非活跃层使用平凡单字母，因此总预算正好等于 \(|B_N|\)。证毕。
\end{proof}

\begin{proposition}[自然扩张单寄存器实现的最小字母表与二进长度]\label{prop:xi-natural-extension-single-register-minimal-alphabet}
\leanverified{paper\_xi\_natural\_extension\_single\_register\_minimal\_alphabet}
设 \(T:X\twoheadrightarrow X\) 为有限纤维满射，并记
\[
M:=\sup_{y\in X}\#T^{-1}(y)<\infty.
\]
考虑定理 \ref{thm:killo-natural-extension-branch-register} 中把自然扩张编码为右移插入系统的单寄存器模型：取有限字母表 \(\mathcal A\)，并对每个 \(y\in X\) 选取单射
\[
\iota_y:T^{-1}(y)\hookrightarrow \mathcal A,
\qquad
\beta(x):=\iota_{T(x)}(x).
\]
则下列两条等价：
\begin{enumerate}
  \item 对全部 \(y\in X\)，上述纤维编码可同时取到；
  \item \(|\mathcal A|\ge M\)。
\end{enumerate}
因此，若记单个历史符号所需的最小二进制寄存器长度为
\[
L_{\min}(T),
\]
则有精确公式
\[
\boxed{\;
L_{\min}(T)=\left\lceil \log_2 M \right\rceil.\;}
\]
并且该下界可达：取
\[
\mathcal A=\{0,1,\dots,M-1\}
\]
并在每条纤维上任选单射 \(\iota_y\)，即可得到与定理 \ref{thm:killo-natural-extension-branch-register} 同型的右移插入双射模型而不丢失任何前史信息。
\end{proposition}
\begin{proof}
若存在某个 \(y\in X\) 使 \(\#T^{-1}(y)=M\)，而 \(|\mathcal A|<M\)，则不存在单射
\[
T^{-1}(y)\hookrightarrow \mathcal A.
\]
于是无法仅凭单一历史符号 \(\beta(x)\in\mathcal A\) 无歧义地区分该纤维内的全部前史分支。故必要条件为
\[
|\mathcal A|\ge M.
\]

反过来，若 \(|\mathcal A|\ge M\)，则对每个 \(y\in X\) 都可选取单射
\[
\iota_y:T^{-1}(y)\hookrightarrow \mathcal A.
\]
由此定义
\[
\beta(x):=\iota_{T(x)}(x)
\]
后，完全沿定理 \ref{thm:killo-natural-extension-branch-register} 的递推构造，即得编码
\[
C:X^{\leftarrow}\longrightarrow X\times \mathcal A^{\NN},
\qquad
C\bigl((x_0,x_{-1},x_{-2},\dots)\bigr)
=
\bigl(x_0,(\beta(x_{-1}),\beta(x_{-2}),\dots)\bigr),
\]
其为到像上的双射，并把自然扩张左移共轭为右移插入变换
\[
\widehat T(x,(a_1,a_2,\dots))=(T(x),(\beta(x),a_1,a_2,\dots)).
\]
故充分性也成立。

于是最小可行字母表大小恰为 \(M\)。把 \(M\) 个字母固定编码为二进制串所需的最小位长正是
\[
\left\lceil \log_2 M \right\rceil,
\]
故
\[
L_{\min}(T)=\left\lceil \log_2 M \right\rceil.
\]
证毕。
\end{proof}

\begin{theorem}[有界纤维满射的最小字母自然扩张实现]\label{thm:xi-natural-extension-minimal-d-alphabet-realization}
\leanverified{paper\_xi\_natural\_extension\_minimal\_d\_alphabet\_realization}
设 \(T:X\twoheadrightarrow X\) 为可数集合上的有限纤维满射，并记
\[
D:=\sup_{y\in X}\#T^{-1}(y)<\infty,
\qquad
[D]:=\{1,\dots,D\}.
\]
对每个 \(y\in X\) 任选单射
\[
\iota_y:T^{-1}(y)\hookrightarrow [D],
\qquad
\beta(x):=\iota_{T(x)}(x).
\]
定义编码
\[
C_D:X^{\leftarrow}\longrightarrow X\times [D]^{\NN},
\]
\[
C_D\bigl((x_0,x_{-1},x_{-2},\dots)\bigr)
:=
\bigl(x_0,(\beta(x_{-1}),\beta(x_{-2}),\dots)\bigr),
\]
并记
\[
\widehat X_D:=C_D(X^{\leftarrow})\subset X\times [D]^{\NN}.
\]
则 \(C_D\) 把自然扩张 \((X^{\leftarrow},\sigma)\) 共轭到 \(\widehat X_D\) 上的双射
\[
\widehat T_D\bigl(x,(a_1,a_2,\dots)\bigr)
:=
\bigl(Tx,(\beta(x),a_1,a_2,\dots)\bigr).
\]
若进一步存在某个 \(y_0\in X\) 满足
\[
\#T^{-1}(y_0)=D,
\]
则任何此类右移插入实现都不可能把字母表缩小到大小 \(D-1\) 的集合。因而 \(D\) 正是自然扩张单寄存器实现的最小可行字母数。
\end{theorem}
\begin{proof}
由定理 \ref{thm:killo-natural-extension-branch-register}，若允许目标字母表为 \(\NN\)，则上述编码把自然扩张 \((X^{\leftarrow},\sigma)\) 共轭为右移插入系统。由于每条纤维的大小都不超过 \(D\)，全部单射 \(\iota_y\) 可统一取值于 \([D]\)，故像空间自动落在 \(X\times [D]^{\NN}\) 中，从而得到所示子空间 \(\widehat X_D\) 与双射 \(\widehat T_D\)。

最小性则正是命题 \ref{prop:xi-natural-extension-single-register-minimal-alphabet} 的内容：若某条纤维大小等于 \(D\)，则任何单寄存器实现的字母表大小都至少为 \(D\)；反之取字母表 \([D]\) 即可实现。因此最小可行字母数恰为 \(D\)。证毕。
\end{proof}

\paragraph{局部化数系的逐素数自然扩张}
把定理 \ref{thm:xi-natural-extension-minimal-d-alphabet-realization} 专用于局部化数系上的素数邻接覆盖时，自然扩张对象与新增的 prime-register 方向都可写成完全闭式。

\begin{theorem}[局部化数系在新素数邻接下的自然扩张恰为扩大的 solenoid]\label{thm:xi-time-part62ac-localized-prime-adjoining-natural-extension}
设 \(S\subset\mathbb P\) 为有限素数集，并记
\[
\Sigma_S:=\widehat{G_S},
\qquad
G_S=\ZZ[S^{-1}],
\qquad
0\longrightarrow \prod_{q\in S}\ZZ_q\longrightarrow \Sigma_S\xrightarrow{\ \pi_S\ }\TT\longrightarrow 0.
\]
取素数 \(p\notin S\)，并令
\[
m_p:G_S\to G_S,
\qquad
x\longmapsto px,
\]
其 Pontryagin 对偶记为
\[
\Phi_{S,p}:=\widehat{m_p}:\Sigma_S\twoheadrightarrow \Sigma_S.
\]
则成立：
\begin{enumerate}
  \item \(\Phi_{S,p}\) 的每条纤维恰含 \(p\) 个点；
  \item 其自然扩张
  \[
  \Sigma_S^{\leftarrow,\Phi_{S,p}}
  :=
  \varprojlim\bigl(\Sigma_S\xleftarrow{\Phi_{S,p}}\Sigma_S\xleftarrow{\Phi_{S,p}}\cdots\bigr)
  \]
  与 \(\Sigma_{S\cup\{p\}}=\widehat{G_{S\cup\{p\}}}\) 规范同构；
  \item 由定理 \ref{thm:xi-natural-extension-minimal-d-alphabet-realization}，该自然扩张的单寄存器右移插入实现可精确取 \(p\)-元字母表 \([p]\)，且不能进一步缩小。
\end{enumerate}
\end{theorem}
\begin{proof}
由于 \(m_p\) 为离散群单射，其对偶 \(\Phi_{S,p}\) 为紧群满射。又由短正合列
\[
0\longrightarrow G_S\xrightarrow{m_p}G_S\longrightarrow G_S/pG_S\longrightarrow 0
\]
对偶化得到
\[
0\longrightarrow \widehat{G_S/pG_S}\longrightarrow \Sigma_S\xrightarrow{\Phi_{S,p}}\Sigma_S\longrightarrow 0,
\]
故
\[
\ker(\Phi_{S,p})\cong \widehat{G_S/pG_S}.
\]
因为 \(p\notin S\)，\(G_S=\ZZ[S^{-1}]\) 中一切分母在 \(\ZZ/p\ZZ\) 中都可逆，所以模 \(p\) 约化给出同构
\[
G_S/pG_S\cong \ZZ/p\ZZ.
\]
于是
\[
\ker(\Phi_{S,p})\cong \ZZ/p\ZZ,
\]
从而每条纤维都恰是一个大小为 \(p\) 的陪集，第一条成立。

对第二条，Pontryagin 对偶把逆极限变为直极限，因此
\[
\bigl(\Sigma_S^{\leftarrow,\Phi_{S,p}}\bigr)^\vee
\cong
\varinjlim\bigl(G_S\xrightarrow{m_p}G_S\xrightarrow{m_p}\cdots\bigr).
\]
由于 \(m_p\) 为乘以 \(p\)，该直极限可与并集
\[
\bigcup_{k\ge 0}p^{-k}G_S
\]
规范识别，而后者恰为
\[
\ZZ[(S\cup\{p\})^{-1}]=G_{S\cup\{p\}}.
\]
再对偶化即得
\[
\Sigma_S^{\leftarrow,\Phi_{S,p}}\cong \widehat{G_{S\cup\{p\}}}=\Sigma_{S\cup\{p\}}.
\]

第三条由前两条与定理 \ref{thm:xi-natural-extension-minimal-d-alphabet-realization} 立即推出：此时纤维大小统一等于 \(p\)，故自然扩张恰可由 \(p\)-元分支指标寄存器实现，而最小字母数也正是 \(p\)。证毕。
\end{proof}

\begin{corollary}[逐素数邻接的顺序无关性与缺失素方向的精确核]\label{cor:xi-time-part62ac-prime-adjoining-order-independent-kernel-product}
设 \(S\subset\mathbb P\) 为有限素数集，\(T\subset\mathbb P\setminus S\) 为有限子集。则无论按何种顺序逐个对 \(T\) 中素数执行定理 \ref{thm:xi-time-part62ac-localized-prime-adjoining-natural-extension} 的邻接扩张，所得最终自然扩张对象都规范同构于
\[
\Sigma_{S\cup T}=\widehat{G_{S\cup T}}.
\]
并且存在规范短正合列
\[
0\longrightarrow \prod_{p\in T}\ZZ_p
\longrightarrow
\Sigma_{S\cup T}
\longrightarrow
\Sigma_S
\longrightarrow 0.
\]
\end{corollary}
\begin{proof}
任取 \(T=\{p_1,\dots,p_r\}\) 的一个排列。反复应用定理 \ref{thm:xi-time-part62ac-localized-prime-adjoining-natural-extension} 的第二条可得，逐步邻接后的对象依次为
\[
\Sigma_{S\cup\{p_1\}},
\quad
\Sigma_{S\cup\{p_1,p_2\}},
\quad
\dots,
\quad
\Sigma_{S\cup T}.
\]
故最终对象只取决于并集 \(S\cup T\)，与邻接顺序无关。

至于短正合列，这正是推论 \ref{cor:xi-localized-integers-prime-axis-increment-law} 在集合差 \(T=(S\cup T)\setminus S\) 上的直接应用。证毕。
\end{proof}


\paragraph{Smith 扭结的 Dirichlet 主项、torsion 解析指纹与整数椭圆的 Grothendieck 排斥}
在模核 Dirichlet 生成函数的 \(\zeta\)-主因子、\(\,p\)-进 Smith 指数的局部恢复以及整数轴对齐椭圆幺半群的自由原子分解已经建立之后，还可进一步把平均阶、torsion 余核的一维解析刻画与有限生成可消去目标上的排斥律压缩为下列四条后果。

\begin{theorem}[规范化模核 Dirichlet 级数可写成正系数有限 Euler 修正与 \texorpdfstring{$\zeta$}{zeta} 的精确乘积]\label{thm:xi-smith-normalized-kernel-positive-finite-euler-correction}
\leanverified{paper\_xi\_smith\_normalized\_kernel\_positive\_finite\_euler\_correction}
设 \(A\in M_{m\times n}(\ZZ)\) 满足 \(\operatorname{rank}_{\QQ}(A)=m\)，其 Smith 不变量为
\[
d_1\mid d_2\mid \cdots \mid d_m.
\]
定义
\[
h_A(b):=\frac{\#\ker(A_b)}{b^{\,n-m}}
=
\prod_{i=1}^{m}\gcd(d_i,b),
\]
\[
\mathcal H_A(s):=\sum_{b\ge 1}\frac{h_A(b)}{b^s},
\qquad \Re s>1,
\]
并记
\[
\Delta_A:=\prod_{i=1}^{m}d_i.
\]
对每个素数 \(p\)，定义
\[
e_i(p):=v_p(d_i),\qquad
e_p:=\max_i e_i(p),
\]
\[
s_p(k):=\sum_{i=1}^{m}\min(k,e_i(p))
\qquad (k\ge 0).
\]
则有精确 Euler 分解
\[
\boxed{\;
\mathcal H_A(s)
=
\zeta(s)\prod_{p\mid \Delta_A}Q_{A,p}(p^{-s}),
\;}
\]
其中每个局部修正项都是正系数多项式
\[
\boxed{\;
Q_{A,p}(X)
=
1+\sum_{k=1}^{e_p}
\Bigl(p^{s_p(k)}-p^{s_p(k-1)}\Bigr)X^k
\in 1+X\ZZ_{>0}[X].\;}
\]
特别地，\(\mathcal H_A(s)\) 可亚纯延拓到整个复平面，且唯一极点是 \(s=1\) 处的单极点。
\end{theorem}
\begin{proof}
由推论 \ref{cor:killo-kernel-size-general-modulus}，
\[
h_A(b)=\prod_{i=1}^{m}\gcd(d_i,b),
\]
故 \(h_A\) 为乘法函数，从而
\[
\mathcal H_A(s)=\prod_p L_{A,p}(s),
\qquad
L_{A,p}(s):=\sum_{k\ge 0}\frac{h_A(p^k)}{p^{ks}}.
\]
再由定理 \ref{thm:killo-smith-loss-spectrum}，
\[
h_A(p^k)=p^{s_p(k)}.
\]
若 \(p\nmid \Delta_A\)，则全部 \(e_i(p)=0\)，于是 \(s_p(k)=0\) 对一切 \(k\) 都成立，从而
\[
L_{A,p}(s)=\sum_{k\ge 0}p^{-ks}=\frac{1}{1-p^{-s}}.
\]

以下设 \(p\mid \Delta_A\)。令
\[
a_k:=p^{s_p(k)}.
\]
当 \(k\ge e_p\) 时，\(\min(k,e_i(p))=e_i(p)\) 对全部 \(i\) 都成立，故 \((a_k)\) 在 \(k\ge e_p\) 后稳定。于是
\[
(1-X)\sum_{k\ge 0}a_kX^k
=
a_0+\sum_{k\ge 1}(a_k-a_{k-1})X^k
\]
在 \(k>e_p\) 时系数全部消失，因此
\[
(1-X)\sum_{k\ge 0}a_kX^k
=
1+\sum_{k=1}^{e_p}(a_k-a_{k-1})X^k
=
Q_{A,p}(X).
\]
把 \(X=p^{-s}\) 代回，便得
\[
L_{A,p}(s)
=
\sum_{k\ge 0}p^{s_p(k)-ks}
=
\frac{Q_{A,p}(p^{-s})}{1-p^{-s}}.
\]
将全部素数上的局部因子相乘，得到
\[
\mathcal H_A(s)
=
\prod_{p\nmid \Delta_A}\frac{1}{1-p^{-s}}
\prod_{p\mid \Delta_A}\frac{Q_{A,p}(p^{-s})}{1-p^{-s}}
=
\zeta(s)\prod_{p\mid \Delta_A}Q_{A,p}(p^{-s}).
\]

最后证明系数为正整数。对 \(1\le k\le e_p\)，有
\[
s_p(k)-s_p(k-1)=\#\{\,i:e_i(p)\ge k\,\}\ge 1,
\]
因为 \(k\le e_p\) 时至少有一个 \(i\) 满足 \(e_i(p)\ge k\)。故
\[
p^{s_p(k)}-p^{s_p(k-1)}\in \ZZ_{>0}.
\]
有限乘积 \(\prod_{p\mid \Delta_A}Q_{A,p}(p^{-s})\) 是整函数，因此 \(\mathcal H_A(s)\) 的唯一极点只能来自 \(\zeta(s)\) 在 \(s=1\) 的单极点。证毕。
\end{proof}

\begin{corollary}[规范化模核计数的平均阶主项与 Smith 扭结检测常数]\label{cor:xi-smith-normalized-kernel-average-order}
\leanverified{paper\_xi\_smith\_normalized\_kernel\_average\_order}
沿用定理 \ref{thm:xi-smith-normalized-kernel-positive-finite-euler-correction} 的记号，记
\[
G_A(s):=\prod_{p\mid \Delta_A}Q_{A,p}(p^{-s}),
\qquad
R_A:=G_A(1)=\prod_{p\mid \Delta_A}Q_{A,p}(p^{-1}).
\]
则有：
\begin{enumerate}
  \item \(R_A>0\)，并且
  \[
  R_A=1
  \iff
  \Delta_A=1
  \iff
  d_1=\cdots=d_m=1.
  \]
  \item 规范化模核计数满足显式平均阶
  \[
  \boxed{\;
  \sum_{b\le x}\frac{\#\ker(A_b)}{b^{\,n-m}}
  =
  \sum_{b\le x}h_A(b)
  \sim
  R_A\,x
  \qquad (x\to\infty).\;}
  \]
\end{enumerate}
\end{corollary}
\begin{proof}
由定理 \ref{thm:xi-smith-normalized-kernel-positive-finite-euler-correction}，每个 \(Q_{A,p}(X)\) 具有正整数系数，因此对 \(X=1/p\in(0,1)\) 有
\[
Q_{A,p}(1/p)>0.
\]
故 \(R_A>0\)。

若 \(\Delta_A=1\)，则空乘积给出 \(R_A=1\)。反之若 \(\Delta_A>1\)，则存在某个 \(p\mid \Delta_A\)。此时 \(e_p\ge 1\)，而
\[
Q_{A,p}(1/p)
=
1+\sum_{k=1}^{e_p}\Bigl(p^{s_p(k)}-p^{s_p(k-1)}\Bigr)p^{-k}
>1.
\]
其余因子都为正，因此 \(R_A>1\)。故
\[
R_A=1\iff \Delta_A=1.
\]
又 \(\Delta_A=\prod_i d_i\)，且每个 \(d_i\ge 1\)，故 \(\Delta_A=1\) 等价于全部 \(d_i=1\)。

再由定理 \ref{thm:xi-smith-normalized-kernel-positive-finite-euler-correction}，
\[
\mathcal H_A(s)=\zeta(s)G_A(s),
\]
其中 \(G_A(s)\) 为整函数，因而 \(\mathcal H_A(s)\) 在 \(\Re s\ge 1\) 上唯一奇点是 \(s=1\) 的单极点，留数为
\[
\operatorname*{Res}_{s=1}\mathcal H_A(s)=G_A(1)=R_A.
\]
另一方面，\(h_A(b)\ge 0\) 对全部 \(b\) 成立。由 Wiener--Ikehara 定理，
\[
\sum_{b\le x}h_A(b)\sim R_Ax.
\]
代回 \(h_A(b)=\#\ker(A_b)/b^{\,n-m}\) 即得所示平均阶。证毕。
\end{proof}

\begin{theorem}[规范化模核 Dirichlet 级数完整编码 Smith 扭结余核]\label{thm:xi-smith-dirichlet-complete-torsion-invariant}
\leanverified{paper\_xi\_smith\_dirichlet\_complete\_torsion\_invariant}
设 \(A\in M_{m_A\times n_A}(\ZZ)\)、\(B\in M_{m_B\times n_B}(\ZZ)\) 都满足满行秩条件。则下列命题等价：
\begin{enumerate}
  \item \(\mathcal H_A(s)=\mathcal H_B(s)\) 作为 Dirichlet 级数恒等；
  \item 对每个素数 \(p\) 与每个 \(k\ge 1\)，都有
  \[
  \frac{\#\ker(A_{p^k})}{p^{k(n_A-m_A)}}
  =
  \frac{\#\ker(B_{p^k})}{p^{k(n_B-m_B)}}.
  \]
  等价地，\(h_A(p^k)=h_B(p^k)\)；
  \item \(A\) 与 \(B\) 的全部非平凡 Smith 不变量（即大于 \(1\) 的不变量因子）多重集完全一致；
  \item
  \[
  \operatorname{coker}(A)_{\mathrm{tors}}
  \cong
  \operatorname{coker}(B)_{\mathrm{tors}}.
  \]
\end{enumerate}
\end{theorem}
\begin{proof}
第 \(1\) 条推出第 \(2\) 条是显然的，因为 Dirichlet 级数相等意味着全部系数相等，特别在 \(b=p^k\) 时系数相等。

现证第 \(2\) 条推出第 \(3\) 条。对每个素数 \(p\)，定义
\[
f_{A,p}(k):=v_p\!\bigl(h_A(p^k)\bigr),
\qquad
f_{B,p}(k):=v_p\!\bigl(h_B(p^k)\bigr).
\]
由第 \(2\) 条，\(f_{A,p}(k)=f_{B,p}(k)\) 对一切 \(p,k\) 都成立。另一方面，
\[
f_{A,p}(k)=\sum_i \min(k,e_i^A(p)),
\qquad
f_{B,p}(k)=\sum_i \min(k,e_i^B(p)).
\]
由推论 \ref{cor:xi-smith-second-discrete-curvature-recovery}，
\[
\#\{\,i:e_i^A(p)=k\,\}
=
2f_{A,p}(k)-f_{A,p}(k-1)-f_{A,p}(k+1),
\]
而右端完全由 \(f_{A,p}\) 决定；同理适用于 \(B\)。因此对每个 \(p\)，正 \(p\)-进指数的多重集
\[
\{\,e_i^A(p)>0\,\}
\]
与
\[
\{\,e_i^B(p)>0\,\}
\]
完全一致。

对每个 \(p\)，把这些正指数按非降序排列，并在左端补零到同一长度 \(r\)，得到
\[
0\le e_1(p)\le \cdots \le e_r(p).
\]
则 torsion 部分的 Smith 不变量正是
\[
t_j:=\prod_{p}p^{e_j(p)}
\qquad (1\le j\le r),
\]
其中省去等于 \(1\) 的项即可得到全部非平凡不变量。由于 \(A\) 与 \(B\) 对每个 \(p\) 的指数列一致，故它们的全部非平凡 Smith 不变量多重集也一致，第 \(3\) 条成立。

第 \(3\) 条与第 \(4\) 条等价，是 Smith 正规形与有限生成交换群结构定理的直接内容。

最后证第 \(4\) 条推出第 \(1\) 条。若 torsion 余核同构，则 \(A\) 与 \(B\) 的全部非平凡 Smith 不变量一致。记这些不变量分别为
\[
t_1,\dots,t_r.
\]
则
\[
h_A(b)=\prod_{j=1}^{r}\gcd(t_j,b)=h_B(b)
\qquad (\forall\, b\ge 1),
\]
因为平凡不变量 \(1\) 对乘积没有贡献。故全部 Dirichlet 系数相同，从而
\[
\mathcal H_A(s)=\mathcal H_B(s).
\]
证毕。
\end{proof}

\begin{theorem}[整数轴对齐椭圆幺半群的 Grothendieck 完成具有两份可数秩自由度，并排斥一切有限生成可消去忠实线性化]\label{thm:xi-integer-ellipse-grothendieck-double-countable-rank}
\leanverified{paper\_xi\_integer\_ellipse\_grothendieck\_double\_countable\_rank}
沿用定理 \ref{thm:killo-ellipse-diagonal-prime-register-equivalence} 与定理 \ref{thm:xi-integer-ellipse-free-commutative-monoid} 的记号。则：
\begin{enumerate}
  \item \((\mathsf E_{\NN},\odot)\) 的 Grothendieck 群完成满足
  \[
  \boxed{\;
  G(\mathsf E_{\NN})
  \cong
  \left(\bigoplus_{p\in\mathcal P}\ZZ\right)
  \oplus
  \left(\bigoplus_{p\in\mathcal P}\ZZ\right).\;}
  \]
  \item 不存在任何单射幺半群同态
  \[
  \mathsf E_{\NN}\hookrightarrow M,
  \]
  其中 \(M\) 为有限生成可消去交换幺半群。
  \item 特别地，对任意有限 \(k\)，都不存在忠实线性化
  \[
  \mathsf E_{\NN}\hookrightarrow (\NN^k,+).
  \]
\end{enumerate}
\end{theorem}
\begin{proof}
由定理 \ref{thm:killo-ellipse-diagonal-prime-register-equivalence}，
\[
(\mathsf E_{\NN},\odot)
\cong
\bigl(\NN^{(\mathcal P)}\times \NN^{(\mathcal P)},+\bigr),
\]
其中 \(\NN^{(\mathcal P)}\) 表示以素数集 \(\mathcal P\) 为基的自由交换幺半群。对 Grothendieck 完成取直积，得到
\[
G(\mathsf E_{\NN})
\cong
G\!\bigl(\NN^{(\mathcal P)}\bigr)\oplus G\!\bigl(\NN^{(\mathcal P)}\bigr)
\cong
\left(\bigoplus_{p\in\mathcal P}\ZZ\right)\oplus
\left(\bigoplus_{p\in\mathcal P}\ZZ\right),
\]
这就得到第 \(1\) 条。

为证第 \(2\) 条，考虑子幺半群
\[
\mathsf E_x:=\{\,E_{a,1}:a\in\NN_{>0}\,\}\subset \mathsf E_{\NN}.
\]
由同一定理，
\[
(\mathsf E_x,\odot)\cong \NN_{\times}:=(\NN_{\ge 1},\times).
\]
若存在单射幺半群同态
\[
\Phi:\mathsf E_{\NN}\hookrightarrow M
\]
且 \(M\) 有限生成、可消去、交换，则其限制
\[
\Phi|_{\mathsf E_x}:\mathsf E_x\hookrightarrow M
\]
仍为单射幺半群同态。经由 \(\mathsf E_x\cong \NN_{\times}\)，便得到
\[
\NN_{\times}\hookrightarrow M,
\]
这与定理 \ref{thm:killo-prime-freedom-non-finitizability} 矛盾。故此类 \(\Phi\) 不存在。

第 \(3\) 条只是把 \(M\) 取为有限生成可消去交换幺半群 \((\NN^k,+)\) 的特例。证毕。
\end{proof}

\noindent
这组结果把 Smith 模核谱与整数轴对齐椭圆的素数寄存器结构收束为同一条解析--代数链：一方面，规范化模核计数的 Dirichlet 级数严格分解为 \(\zeta\) 与有限 Euler 修正的乘积，其平均阶主项与 torsion 余核均被显式常数完全控制；另一方面，整数轴对齐椭圆幺半群虽具有完全自由的原子分解，却仍保留两份可数秩 Grothendieck 自由度，从而排斥一切有限生成可消去目标上的忠实压缩。


\paragraph{ZG 的临界极点、Lee--Yang dyadic 分支逆极限与局部/branch 数域分离}
在 Zeckendorf--Gödel 素数码的密度--Dirichlet 链、Lee--Yang 分歧塔的 \(2\)-进 torsor 几何，以及真实输入 \(40\) 状态核的局部累积量闭式与 branch 三次数域已经建立之后，还可进一步抽取出下列五条互相闭合的终端结论。前两条把 \(s=1\) 的主极点精确锁定为自然密度；中间两条将 dyadic 分歧逆极限压缩为五片 \(2\)-进二维纤维并完全计算其周期点 \(\zeta\)-函数；最后一条则把局部累积量的二次数域与 branch 参数的三次数域严格分离。

\begin{theorem}[Zeckendorf--Gödel Dirichlet 级数在 \texorpdfstring{$s=1$}{s=1} 的简单极点与自然密度留数]\label{thm:xi-time-part62a-zg-simple-pole-density-residue}
\leanverified{paper\_xi\_time\_part62a\_zg\_simple\_pole\_density\_residue}
沿用定理 \ref{thm:xi-zg-zeta-factorization}、推论 \ref{cor:xi-zg-dirichlet-residue} 与定理 \ref{thm:xi-zg-abel-residue-log-density} 的记号。记
\[
D_{\mathrm{ZG}}(s):=\sum_{n\in\mathrm{ZG}}n^{-s}
\qquad (\Re(s)>1).
\]
则有
\[
\boxed{\;
\lim_{\sigma\downarrow 1}(\sigma-1)D_{\mathrm{ZG}}(\sigma)=\delta_{\mathrm{ZG}}.
\;}
\]
等价地，\(D_{\mathrm{ZG}}(s)\) 在 \(s=1\) 处具有留数 \(\delta_{\mathrm{ZG}}\) 的简单极点。
\end{theorem}
\begin{proof}
推论 \ref{cor:xi-zg-dirichlet-residue} 已给出 \(D_{\mathrm{ZG}}(s)\) 在 \(s=1\) 处具有简单极点，且其留数等于 \(\mathrm{ZG}\) 的 Dirichlet 密度。另一方面，定理 \ref{thm:xi-zg-abel-residue-log-density} 已把同一留数识别为自然密度 \(\delta_{\mathrm{ZG}}\)。故所示极限与简单极点结论同时成立。证毕。
\end{proof}

\begin{corollary}[Zeckendorf--Gödel Dirichlet 级数的 \texorpdfstring{$\zeta$}{zeta} 归一化密度极限]\label{cor:xi-time-part62a-zg-zeta-normalized-density-limit}
\leanverified{paper\_xi\_time\_part62a\_zg\_zeta\_normalized\_density\_limit}
沿用定理 \ref{thm:xi-time-part62a-zg-simple-pole-density-residue} 的记号，则
\[
\boxed{\;
\lim_{\sigma\downarrow 1}\frac{D_{\mathrm{ZG}}(\sigma)}{\zeta(\sigma)}
=
\delta_{\mathrm{ZG}}.
\;}
\]
\end{corollary}
\begin{proof}
由定理 \ref{thm:xi-time-part62a-zg-simple-pole-density-residue}，
\[
D_{\mathrm{ZG}}(\sigma)=\frac{\delta_{\mathrm{ZG}}}{\sigma-1}+O(1)
\qquad (\sigma\downarrow 1).
\]
而 Riemann \(\zeta\) 函数满足
\[
\zeta(\sigma)=\frac{1}{\sigma-1}+O(1)
\qquad (\sigma\downarrow 1).
\]
两式相除即得结论。再由定理 \ref{thm:xi-zg-zeta-factorization} 的二维 hard-core transfer 表示，上述极限常数同时也是 prime-register 矩阵 Euler 动力学所生成的临界常数。证毕。
\end{proof}

\begin{theorem}[Lee--Yang dyadic 分支逆极限的五分量 \texorpdfstring{$\ZZ_2^2$}{Z2^2} 识别]\label{thm:xi-time-part62a-leyang-branch-inverse-limit-five-z2x2}
\leanverified{paper\_xi\_time\_part62a\_leyang\_branch\_inverse\_limit\_five\_z2x2}
沿用推论 \ref{cor:xi-terminal-zm-leyang-finite-branch-regular-4ary-address} 的记号，记
\[
\mathcal R_\infty:=\Pi_\infty^{\mathrm{fin}}
=
\varprojlim_n \Pi_n^{\mathrm{fin}}
=
\bigsqcup_{i=0}^{4}\Pi_{i,\infty}.
\]
则在每个分量选择一条相容提升，并在 \(T_2(E_{\min})\) 上选取一组 \(\ZZ_2\)-基之后，有同构
\[
\boxed{\;
\mathcal R_\infty
\cong
\bigsqcup_{i=0}^{4}T_2(E_{\min})
\cong
\bigsqcup_{i=0}^{4}\ZZ_2^2.
\;}
\]
换言之，Lee--Yang 分支塔的 dyadic 逆极限恰为五片彼此不交的二维 \(2\)-进纤维。
\end{theorem}
\begin{proof}
推论 \ref{cor:xi-terminal-zm-leyang-finite-branch-regular-4ary-address} 的第三条已经给出
\[
\mathcal R_\infty
\cong
\bigsqcup_{i=0}^{4}\bigl(\mathbf P_i+T_2(E_{\min})\bigr)
\cong
\bigsqcup_{i=0}^{4}T_2(E_{\min}),
\]
其中 \(\mathbf P_i\) 为各分支上的相容提升。又 \(T_2(E_{\min})\) 是秩 \(2\) 的自由 \(\ZZ_2\)-模；一旦选定 \(\ZZ_2\)-基，便得到
\[
T_2(E_{\min})\cong \ZZ_2^2.
\]
合并即得所示同构。证毕。
\end{proof}

\begin{corollary}[五分支 dyadic 自映射的周期点常数化与 Artin--Mazur \texorpdfstring{$\zeta$}{zeta} 闭式]\label{cor:xi-time-part62a-leyang-branch-artin-mazur-zeta}
\leanverified{paper\_xi\_time\_part62a\_leyang\_branch\_artin\_mazur\_zeta}
在定理 \ref{thm:xi-time-part62a-leyang-branch-inverse-limit-five-z2x2} 的识别下，于第 \(i\) 个分量上定义
\[
F(\mathbf P_i+x):=\mathbf P_i+2x
\qquad
\bigl(x\in T_2(E_{\min}),\ 0\le i\le 4\bigr).
\]
则对任意 \(n\ge 1\)，\(F^n\) 的不动点数恒为
\[
\#\operatorname{Fix}(F^n)=5,
\]
从而 Artin--Mazur \(\zeta\)-函数满足闭式
\[
\boxed{\;
\zeta_F(z):=
\exp\!\Bigl(
\sum_{n\ge 1}\frac{\#\operatorname{Fix}(F^n)}{n}z^n
\Bigr)
=(1-z)^{-5}.
\;}
\]
\end{corollary}
\begin{proof}
在任一分量上，\(F^n\) 等于 \(x\mapsto 2^n x\)。于是其不动点方程为
\[
2^n x=x
\quad\Longleftrightarrow\quad
(2^n-1)x=0.
\]
由于 \(2^n-1\) 为奇数，故它在 \(\ZZ_2\) 中是单位，从而在
\(
T_2(E_{\min})\cong \ZZ_2^2
\)
上只有零解 \(x=0\)。因此每个分量恰有一个 \(F^n\)-不动点，总计
\[
\#\operatorname{Fix}(F^n)=5.
\]
代回 Artin--Mazur 定义即得
\[
\zeta_F(z)
=
\exp\!\Bigl(\sum_{n\ge 1}\frac{5}{n}z^n\Bigr)
=(1-z)^{-5}.
\]
证毕。
\end{proof}

\begin{theorem}[局部累积量二次数域与 branch 三次数域的线性无交]\label{thm:xi-time-part62a-local-cumulant-branch-cubic-linear-disjointness}
\leanverified{paper\_xi\_time\_part62a\_local\_cumulant\_branch\_cubic\_linear\_disjointness}
沿用命题 \ref{prop:real-input-40-output-cumulants-closed}、定理 \ref{thm:xi-terminal-zm-kappa-square-cubic-field-s3}、\ref{thm:xi-terminal-zm-stokes-leyang-shared-artin-representation} 与推论 \ref{cor:xi-terminal-zm-stokes-leyang-common-quadratic-resolvent} 的记号。记
\[
u:=\kappa_*^2,
\qquad
b:=B^2,
\qquad
K_{\mathrm{loc}}:=\QQ(\sqrt5),
\qquad
K_{\mathrm{br}}:=\QQ(u)=\QQ(b)=\QQ(\lambda_*).
\]
则有
\[
\boxed{\;
K_{\mathrm{loc}}\cap K_{\mathrm{br}}=\QQ,
\qquad
[K_{\mathrm{loc}}K_{\mathrm{br}}:\QQ]=6.
\;}
\]
特别地，任一非有理局部累积量
\(
P^{(k)}(0)\in K_{\mathrm{loc}}\setminus\QQ
\)
都不属于 \(K_{\mathrm{br}}\)；反之，只要 \(R(T)\in\QQ(T)\) 在 \(u\)、\(b\) 或 \(\lambda_*\) 处有定义且其取值不在 \(\QQ\) 中，则该值亦不属于 \(K_{\mathrm{loc}}\)。
\end{theorem}
\begin{proof}
命题 \ref{prop:real-input-40-output-cumulants-closed} 给出全部局部累积量
\[
P^{(k)}(0)\in K_{\mathrm{loc}}=\QQ(\sqrt5).
\]
另一方面，定理 \ref{thm:xi-terminal-zm-kappa-square-cubic-field-s3} 表明 \([\,\QQ(u):\QQ\,]=3\)，而定理
\ref{thm:xi-terminal-zm-stokes-leyang-shared-artin-representation} 给出
\[
K_{\mathrm{br}}=\QQ(u)=\QQ(b)=\QQ(\lambda_*).
\]
故 \([K_{\mathrm{br}}:\QQ]=3\)。

再由推论 \ref{cor:xi-terminal-zm-stokes-leyang-common-quadratic-resolvent}，\(K_{\mathrm{br}}\) 的伽罗瓦闭包唯一二次子域为
\[
\QQ(\sqrt{-111}),
\]
显然不同于 \(K_{\mathrm{loc}}=\QQ(\sqrt5)\)。尤其
\[
[K_{\mathrm{loc}}\cap K_{\mathrm{br}}:\QQ]
\mid
[K_{\mathrm{loc}}:\QQ]=2,
\qquad
[K_{\mathrm{loc}}\cap K_{\mathrm{br}}:\QQ]
\mid
[K_{\mathrm{br}}:\QQ]=3,
\]
故
\[
K_{\mathrm{loc}}\cap K_{\mathrm{br}}=\QQ.
\]
于是由塔式公式
\[
[K_{\mathrm{loc}}K_{\mathrm{br}}:\QQ]
=
\frac{[K_{\mathrm{loc}}:\QQ]\,[K_{\mathrm{br}}:\QQ]}
{[K_{\mathrm{loc}}\cap K_{\mathrm{br}}:\QQ]}
=2\cdot 3=6.
\]

若某个非有理局部累积量 \(P^{(k)}(0)\) 属于 \(K_{\mathrm{br}}\)，则它同时属于
\(
K_{\mathrm{loc}}\cap K_{\mathrm{br}}=\QQ
\)，矛盾。对 \(R(u)\)、\(R(b)\) 与 \(R(\lambda_*)\) 的情形同理。证毕。
\end{proof}

\begin{remark}[有限层审计接口]\label{rem:xi-time-part62a-zg-branch-field-audit}
脚本 \texttt{scripts/exp\_xi\_zg\_branch\_tower\_field\_separation\_audit.py}
同时给出三类可复核输出：其一是 \(D_{\mathrm{ZG}}(s)\) 在 \(s\downarrow 1\) 时的极点缩放值与 \(\zeta\)-归一化值如何逼近同一密度常数；其二是五分支 dyadic 塔在有限 \(2\)-幂截断上的固定点计数恒等于每分支 \(1\) 个、总计 \(5\) 个；其三是 Lee--Yang branch 三次在 \(\QQ(\sqrt5)\) 上仍不可约，且 \(\sqrt5+u\) 的最小多项式次数恰为 \(6\)。
\end{remark}

\paragraph{单向量全息像集、Lee--Yang dyadic 层化与矩阵 Euler 刚性}
在规范 Gödel 半直积的 \(2^L\)-进仿射实现、Lee--Yang 分歧塔的 dyadic 地址化以及 Zeckendorf--Gödel Dirichlet 链已经建立之后，还可进一步抽出下列五条彼此闭合的结构结论。第一条给出单向量全息像集的精确自相似、Haar 零测与维数公式；第二、三条把 Lee--Yang 分支逆极限与单向量模型之间的秩障碍压缩为纯 \(\ZZ_2\)-模论判据；第四条把 Zeckendorf--Gödel Dirichlet 级数的临界指数与矩阵 Euler 乘积的不可退化性同时锁定；第五条则把同一矩阵 Euler 动力学进一步压缩为 prime-indexed continuant 递推。

\begin{conclusion}[单向量 \texorpdfstring{$2^L$}{2^L}-进全息像集的精确自相似、Haar 零测与维数公式]\label{con:xi-time-part62b-rankone-hologram-selfsimilar-measure-dimension}
沿用文稿关于 \(E_{\min}\) 与 \(T_2(E_{\min})\) 的记号。设 \(\mathcal E\) 为有限事件集，\(\mathrm{code}:\mathcal E\hookrightarrow \NN\) 为单射，并记
\[
M:=\max_{e\in\mathcal E}\mathrm{code}(e),
\qquad
B:=2^L>M,
\qquad
T_v:=\ZZ_2v\subset T_2(E_{\min}),
\]
其中 \(v\in T_2(E_{\min})\setminus\{0\}\)。对 \(\mathbf e=(e_t)_{t\ge 1}\in\mathcal E^{\NN}\) 定义
\[
X(\mathbf e):=\sum_{t\ge 1}\mathrm{code}(e_t)\,B^{t-1}v,
\qquad
D:=\mathrm{code}(\mathcal E)\subset\{0,\dots,B-1\},
\]
以及像集
\[
\mathcal X_D:=X(\mathcal E^{\NN})\subset T_v.
\]
则成立：
\begin{enumerate}
  \item 地址级数在 \(T_v\) 中收敛，映射 \(\mathbf e\mapsto X(\mathbf e)\) 为单射，并且
  \[
  \mathcal X_D=\bigsqcup_{d\in D}\bigl(dv+B\mathcal X_D\bigr).
  \]
  更精确地，每一块 \(dv+B\mathcal X_D\) 都等于 \(\mathcal X_D\) 与某个 \(B T_v\)-陪集的交，因此该分解在 \(\mathcal X_D\) 的相对拓扑下为互不相交的 clopen 分解。
  \item 若 \(\mu_v\) 表示 \(T_v\cong \ZZ_2\) 上的标准 Haar 概率测度，则
  \[
  \mu_v(\mathcal X_D)=
  \begin{cases}
  1,& |D|=B,\\
  0,& |D|<B.
  \end{cases}
  \]
  并且当 \(|D|<B\) 时，\(\mathcal X_D\) 在 \(T_v\) 中内点为空。
  \item 在 \(T_v\) 的标准 \(2\)-进超度量下，
  \[
  \dim_{\mathrm H}(\mathcal X_D)
  =
  \dim_{\mathrm M}(\mathcal X_D)
  =
  \frac{\log |D|}{\log B}
  =
  \frac{\log_2|\mathcal E|}{L}.
  \]
\end{enumerate}
\end{conclusion}
\begin{proof}
由于 \(T_v=\ZZ_2v\) 是秩 \(1\) 的自由 \(\ZZ_2\)-模，标量乘法给出同构
\[
\iota_v:\ZZ_2\longrightarrow T_v,
\qquad
a\longmapsto av.
\]
又因 \(|B|_2<1\) 且 \(\mathrm{code}(e_t)\in\{0,\dots,M\}\) 有界，定义 \(\mathcal X_D\) 的级数在完备群 \(T_v\) 中逐项收敛。

现证单射。若 \(X(\mathbf e)=X(\mathbf e')\)，则
\[
\sum_{t\ge 1}\bigl(\mathrm{code}(e_t)-\mathrm{code}(e_t')\bigr)B^{t-1}v=0.
\]
经由 \(\iota_v^{-1}\) 可化为
\[
\sum_{t\ge 1}c_tB^{t-1}=0
\qquad\text{于 }\ZZ_2,
\]
其中 \(c_t:=\mathrm{code}(e_t)-\mathrm{code}(e_t')\in\{-(B-1),\dots,B-1\}\)。若并非所有 \(c_t\) 都为零，取最小的 \(n\) 使 \(c_n\neq 0\)，则
\[
\sum_{t\ge 1}c_tB^{t-1}=B^{n-1}(c_n+Bu)
\]
对某个 \(u\in\ZZ_2\) 成立。由于 \(0<|c_n|<B\)，标量 \(c_n\) 不属于 \(B\ZZ_2\)，故 \(c_n+Bu\notin B\ZZ_2\)，与上式为零矛盾。故一切 \(c_t=0\)，从而 \(\mathbf e=\mathbf e'\)。

对任意 \(\mathbf e=(e_t)_{t\ge 1}\in\mathcal E^{\NN}\)，把首位拆开可得
\[
X(\mathbf e)
=
\mathrm{code}(e_1)v
+B\sum_{u\ge 1}\mathrm{code}(e_{u+1})\,B^{u-1}v.
\]
故
\[
\mathcal X_D\subseteq \bigcup_{d\in D}\bigl(dv+B\mathcal X_D\bigr).
\]
反向包含同样直接：若 \(x=\sum_{u\ge 1}d_uB^{u-1}v\in \mathcal X_D\)，则对任意 \(d\in D\) 都有
\[
dv+Bx
=
\sum_{t\ge 1}\widetilde d_t B^{t-1}v\in \mathcal X_D,
\]
其中 \(\widetilde d_1=d\)、\(\widetilde d_{t+1}=d_t\)。于是并式成立。

进一步，若 \(x\in dv+B\mathcal X_D\)，则显然 \(x\in (dv+B T_v)\cap \mathcal X_D\)。反过来，若
\[
x\in (dv+B T_v)\cap \mathcal X_D,
\]
则可写成
\[
x=\sum_{t\ge 1}d_tB^{t-1}v
\qquad (d_t\in D),
\]
且又有 \(x-dv\in B T_v\)。因此
\[
\sum_{t\ge 1}(d_t-\delta_{1t}d)B^{t-1}v\in B T_v.
\]
经由 \(\iota_v^{-1}\) 传回 \(\ZZ_2\)，便知
\[
(d_1-d)+B u\in B\ZZ_2
\]
对某个 \(u\in\ZZ_2\) 成立。由于 \(d_1,d\in\{0,\dots,B-1\}\)，只可能有 \(d_1=d\)。于是
\[
x=dv+B\sum_{u\ge 1}d_{u+1}B^{u-1}v\in dv+B\mathcal X_D.
\]
故
\[
dv+B\mathcal X_D=(dv+B T_v)\cap \mathcal X_D.
\]
因此每一块都是 \(\mathcal X_D\) 与某个 clopen 陪集 \(dv+B T_v\) 的交，故在 \(\mathcal X_D\) 的相对拓扑下为 clopen。若 \(d\neq d'\)，则上式表明两块的首位数字分别为 \(d\) 与 \(d'\)，由刚才的唯一性它们必互不相交。

现证 Haar 测度。由于乘法 \(x\mapsto Bx\) 在 \(T_v\cong\ZZ_2\) 上的像是指数 \(B\) 的开子群 \(B T_v\)，故对任意可测集 \(A\subseteq T_v\) 都有
\[
\mu_v(BA)=B^{-1}\mu_v(A).
\]
再利用平移不变性与上述不交并分解，
\[
\mu_v(\mathcal X_D)
=
\sum_{d\in D}\mu_v(dv+B\mathcal X_D)
=
|D|\,B^{-1}\mu_v(\mathcal X_D).
\]
若 \(|D|<B\)，则只能有 \(\mu_v(\mathcal X_D)=0\)。

若 \(|D|=B\)，则 \(D\subseteq\{0,\dots,B-1\}\) 必强制 \(D=\{0,\dots,B-1\}\)。此时 \(T_v\cong\ZZ_2\) 的每个元素都具有唯一的 \(B\)-进展开
\[
x=\sum_{t\ge 1}d_t B^{t-1}v,
\qquad
d_t\in D,
\]
故 \(x\in \mathcal X_D\)，从而 \(\mathcal X_D=T_v\) 且 \(\mu_v(\mathcal X_D)=1\)。

若 \(|D|<B\) 而 \(\mathcal X_D\) 含有非空内点，则它必包含某个 \(B^nT_v\)-陪集，因而 Haar 测度至少为 \(B^{-n}>0\)，与 \(\mu_v(\mathcal X_D)=0\) 矛盾。因此内点为空。

最后证明维数。对每个 \(n\ge 1\)，长度 \(n\) 的前缀集合正好是 \(D^n\)，共有 \(|D|^n\) 个。每个前缀
\[
(d_1,\dots,d_n)\in D^n
\]
对应的 cylinder 正好是
\[
\left(\sum_{t=1}^{n}d_tB^{t-1}v+B^nT_v\right)\cap \mathcal X_D.
\]
由上面的唯一性论证，不同前缀给出不同的 \(B^nT_v\)-陪集；而在 \(T_v\cong \ZZ_2\) 的标准 \(2\)-进超度量下，每个这样的陪集直径恰为 \(B^{-n}\)。因此
\[
N_n(\mathcal X_D)=|D|^n,
\]
其中 \(N_n(\mathcal X_D)\) 表示尺度 \(B^{-n}\) 下所需的最少球数。于是
\[
\dim_{\mathrm M}(\mathcal X_D)
=
\lim_{n\to\infty}\frac{\log |D|^n}{-\log B^{-n}}
=
\frac{\log |D|}{\log B}.
\]

若 \(s>\frac{\log |D|}{\log B}\)，则由长度 \(n\) cylinder 覆盖得到
\[
\mathcal H^s_{B^{-n}}(\mathcal X_D)
\le
|D|^nB^{-ns}
\longrightarrow 0,
\]
故 \(\mathcal H^s(\mathcal X_D)=0\)。

反过来，令 \(\nu\) 为 \(D^\NN\) 上等权 Bernoulli 测度，并通过地址映射推前到 \(\mathcal X_D\)。任意半径 \(B^{-n}\) 的球至多与一个长度 \(n\) cylinder 相交，因为不同长度 \(n\) 前缀对应不同的 \(B^nT_v\)-陪集。故对任意此类球 \(U\) 都有
\[
\nu(U)\le |D|^{-n}=B^{-n\log_B|D|}.
\]
由质量分布原理，
\[
\dim_{\mathrm H}(\mathcal X_D)\ge \frac{\log |D|}{\log B}.
\]
综上
\[
\dim_{\mathrm H}(\mathcal X_D)=\dim_{\mathrm M}(\mathcal X_D)=\frac{\log |D|}{\log B}.
\]
又因 \(|D|=|\mathcal E|\) 且 \(B=2^L\)，最后一式可重写为 \(\frac{\log_2|\mathcal E|}{L}\)。证毕。
\end{proof}

\begin{conclusion}[Lee--Yang 分支逆极限恰为五份 \texorpdfstring{$\ZZ_2^2$}{Z2^2} 地址片]\label{con:xi-time-part62b-leyang-branch-inverse-limit-five-z2x2}
沿用推论 \ref{cor:xi-terminal-zm-leyang-finite-branch-regular-4ary-address} 与定理 \ref{thm:xi-time-part62a-leyang-branch-inverse-limit-five-z2x2} 的记号。记
\[
\Pi_\infty^{\mathrm{fin}}
:=
\varprojlim_n \Pi_n^{\mathrm{fin}}
=
\bigsqcup_{i=0}^{4}\Pi_{i,\infty}.
\]
则在每个分支上选取一条相容提升 \(\mathbf P_i\in\Pi_{i,\infty}\)，并在 \(T_2(E_{\min})\) 上选取一组 \(\ZZ_2\)-基之后，有规范同构
\[
\Pi_\infty^{\mathrm{fin}}
\cong
\bigsqcup_{i=0}^{4}\bigl(\mathbf P_i+T_2(E_{\min})\bigr)
\cong
\bigsqcup_{i=0}^{4}T_2(E_{\min})
\cong
\bigsqcup_{i=0}^{4}\ZZ_2^2.
\]
因此，Lee--Yang 分支塔的 dyadic 逆极限不是抽象的 Cantor 型地址集，而是五片彼此不交的 \(2\)-进二维地址片。
\end{conclusion}
\begin{proof}
推论 \ref{cor:xi-terminal-zm-leyang-finite-branch-regular-4ary-address} 的第三条已经给出
\[
\Pi_\infty^{\mathrm{fin}}
\cong
\bigsqcup_{i=0}^{4}\bigl(\mathbf P_i+T_2(E_{\min})\bigr)
\cong
\bigsqcup_{i=0}^{4}T_2(E_{\min}),
\]
其中 \(\mathbf P_i\) 是各分支上的相容提升。又 \(T_2(E_{\min})\) 是秩 \(2\) 的自由 \(\ZZ_2\)-模，一旦选定基便有
\[
T_2(E_{\min})\cong \ZZ_2^2.
\]
合并即得所示同构；这与定理 \ref{thm:xi-time-part62a-leyang-branch-inverse-limit-five-z2x2} 的识别完全一致。证毕。
\end{proof}

\begin{conclusion}[单向量 Gödel--Tate 全息模型对 Lee--Yang 全分支逆极限的结构性障碍]\label{con:xi-time-part62b-rankone-hologram-leyang-rank-obstruction}
沿用结论 \ref{con:xi-time-part62b-rankone-hologram-selfsimilar-measure-dimension} 的记号，并令
\[
\mathcal R_\infty:=\Pi_\infty^{\mathrm{fin}}.
\]
则成立：
\begin{enumerate}
  \item 任意连续群同态
  \[
  \Phi:T_v\longrightarrow T_2(E_{\min})\cong \ZZ_2^2
  \]
  的像都包含在某个秩至多为 \(1\) 的闭 \(\ZZ_2\)-子模中；特别地，\(\Phi(T_v)\) 不可能是 \(\ZZ_2^2\) 的开子群。
  \item 因而，沿单一方向 \(v\) 的 affine \(B\)-进 Gödel--Tate 全息模型不可能在 \(\ZZ_2\)-模意义下实现任一 Lee--Yang 分支分量 \(\mathbf P_i+T_2(E_{\min})\)，更不可能实现五分支逆极限 \(\mathcal R_\infty\)。任何真正覆盖 \(\mathcal R_\infty\) 的 affine 地址模型，其数字集在 \(T_2(E_{\min})\) 中至少必须生成秩 \(2\) 的 \(\ZZ_2\)-子模。
\end{enumerate}
\end{conclusion}
\begin{proof}
在 \(T_v\cong \ZZ_2\) 的识别下，任意连续群同态 \(\Phi:T_v\to T_2(E_{\min})\cong \ZZ_2^2\) 都由单个向量 \(w:=\Phi(v)\) 唯一决定，并满足
\[
\Phi(T_v)=\ZZ_2 w.
\]
故其像始终是一个循环 \(\ZZ_2\)-子模，秩至多为 \(1\)。而 \(\ZZ_2^2\) 的任意开子群都具有有限指数，因而仍是秩 \(2\) 的自由 \(\ZZ_2\)-模。于是不存在具有开像的连续群同态 \(T_v\to \ZZ_2^2\)。

另一方面，结论 \ref{con:xi-time-part62b-rankone-hologram-selfsimilar-measure-dimension} 已表明单向量全息像集始终落在 rank-one 子模 \(T_v\) 内；而结论 \ref{con:xi-time-part62b-leyang-branch-inverse-limit-five-z2x2} 则把每个 Lee--Yang 分支分量识别为某个仿射平移 \(\mathbf P_i+T_2(E_{\min})\)，其线性部分是 rank-two 模 \(T_2(E_{\min})\cong \ZZ_2^2\)。若存在某个沿单向量的 affine 地址实现把 \(\mathcal X_D\) 或 \(T_v\) 在模意义下识别为某个分支分量，则平移归零后将导出一个连续 \(\ZZ_2\)-模同构
\[
T_v\longrightarrow T_2(E_{\min}),
\]
这与第一条矛盾。故任何 affine 地址实现若要覆盖 \(\mathcal R_\infty\)，其数字方向至少需要两条 \(\ZZ_2\)-线性独立的生成方向。证毕。
\end{proof}

\begin{conclusion}[Zeckendorf--Gödel Dirichlet 级数的临界指数与矩阵 Euler 乘积刚性]\label{con:xi-time-part62b-zg-critical-exponent-matrix-euler-rigidity}
沿用定理 \ref{thm:xi-zg-zeta-factorization}、定理 \ref{thm:xi-zg-abel-residue-log-density} 与定理 \ref{thm:xi-time-part62a-zg-simple-pole-density-residue} 的记号，记
\[
D_{\mathrm{ZG}}(s):=\sum_{n\in \mathrm{ZG}}n^{-s},
\qquad
\Re(s)>1.
\]
则成立：
\begin{enumerate}
  \item \(D_{\mathrm{ZG}}(s)\) 的收敛 abscissa 满足
  \[
  \sigma_c(D_{\mathrm{ZG}})=1.
  \]
  \item 不存在幂级数族
  \[
  f_i(x)\in 1+x\RR[[x]]
  \qquad (i\ge 1)
  \]
  使得
  \[
  D_{\mathrm{ZG}}(s)=\prod_{i\ge 1}f_i(p_i^{-s})
  \]
  在 Dirichlet 级数意义下恒等成立。换言之，Zeckendorf--Gödel 约束所导出的 Euler 结构本质上是矩阵型的，不能退化为彼此独立的标量局部因子乘积。
\end{enumerate}
\end{conclusion}
\begin{proof}
由定理 \ref{thm:xi-zg-abel-residue-log-density} 与定理 \ref{thm:xi-time-part62a-zg-simple-pole-density-residue}，当 \(\Re(s)>1\) 时级数绝对收敛，因此 \(\sigma_c(D_{\mathrm{ZG}})\le 1\)。另一方面，每个单素数 \(p_i\) 都属于 \(\mathrm{ZG}\)，因为只激活一个素数位不会违反“相邻素因子不可同时出现”的约束。故对 \(\sigma:=\Re(s)\le 1\) 有
\[
D_{\mathrm{ZG}}(\sigma)\ge \sum_{i\ge 1}p_i^{-\sigma}.
\]
右端在 \(\sigma\le 1\) 时发散，尤其 \(\sigma=1\) 时即为素数调和级数。于是 \(\sigma_c(D_{\mathrm{ZG}})\ge 1\)，从而
\[
\sigma_c(D_{\mathrm{ZG}})=1.
\]

现证第二条。反设存在标量局部因子分解
\[
D_{\mathrm{ZG}}(s)=\prod_{i\ge 1}f_i(p_i^{-s}),
\qquad
f_i(x)=1+c_{i,1}x+c_{i,2}x^2+\cdots .
\]
由于 \(\mathrm{ZG}\) 中不含任何平方因子 \(p_i^2\)，Dirichlet 级数中 \(p_i^{-2s}\) 的系数必须为 \(0\)。而该系数恰等于 \(f_i(x)\) 的 \(x^2\) 项系数 \(c_{i,2}\)，故 \(c_{i,2}=0\)。同理，对所有 \(k\ge 2\) 都有 \(c_{i,k}=0\)。因此每个局部因子只能是
\[
f_i(x)=1+c_{i,1}x.
\]

再由于每个单素数 \(p_i\in \mathrm{ZG}\)，Dirichlet 级数中 \(p_i^{-s}\) 的系数为 \(1\)，故 \(c_{i,1}=1\)。于是
\[
f_i(x)=1+x
\qquad (i\ge 1).
\]
但如此一来，乘积
\[
\prod_{i\ge 1}(1+p_i^{-s})
\]
中的 \(p_i^{-s}p_{i+1}^{-s}\) 项系数等于 \(1\)。这意味着相邻素数积 \(p_ip_{i+1}\) 会以非零系数出现。然而按 \(\mathrm{ZG}\) 的定义，相邻素数 \(p_i,p_{i+1}\) 不能同时出现，故 \(p_ip_{i+1}\notin \mathrm{ZG}\)，对应系数必须为 \(0\)。矛盾。

因此不存在这样的标量 Euler 因子分解。故定理 \ref{thm:xi-zg-zeta-factorization} 中的矩阵 Euler 乘积并非表述上的便利，而是相邻耦合约束所强制的结构刚性。证毕。
\end{proof}

\begin{theorem}[Zeckendorf--Gödel 矩阵 Euler 积的 continuant 递推压缩]\label{thm:xi-time-part62b-zg-matrix-euler-continuant}
\leanverified{paper\_xi\_time\_part62b\_zg\_matrix\_euler\_continuant}
沿用定理 \ref{thm:killo-zg-dirichlet-matrix-euler} 证明中的记号。对 \(n\ge 0\) 记
\[
D_n(s):=A_n+B_n.
\]
则对 \(\Re(s)>1\) 有
\[
D_0(s)=1,
\qquad
D_1(s)=1+p_1^{-s},
\]
并且对一切 \(n\ge 2\)，满足非自治二阶递推
\[
D_n(s)=D_{n-1}(s)+p_n^{-s}D_{n-2}(s).
\]
因而
\[
D_{\mathrm{ZG}}(s)=\lim_{n\to\infty}D_n(s).
\]
\end{theorem}
\begin{proof}
由定理 \ref{thm:killo-zg-dirichlet-matrix-euler} 的证明，
\[
A_n=A_{n-1}+B_{n-1},
\qquad
B_n=p_n^{-s}A_{n-1},
\qquad
(A_0,B_0)=(1,0).
\]
因此
\[
A_n=A_{n-1}+B_{n-1}=D_{n-1}(s).
\]
代回 \(D_n(s)=A_n+B_n\) 得
\[
D_n(s)
=
D_{n-1}(s)+p_n^{-s}A_{n-1}
=
D_{n-1}(s)+p_n^{-s}D_{n-2}(s),
\]
这就得到所示递推。初值
\[
D_0(s)=A_0+B_0=1,
\qquad
D_1(s)=A_1+B_1=1+p_1^{-s}
\]
同样直接由定义给出。

最后，定理 \ref{thm:killo-zg-dirichlet-matrix-euler} 已证明
\[
D_{\mathrm{ZG}}(s)=\lim_{n\to\infty}(A_n+B_n),
\]
故极限正是 \(\lim_{n\to\infty}D_n(s)\)。证毕。
\end{proof}

\noindent
上述五条结论把单向量 \(2^L\)-进全息像集的自相似维数律、Lee--Yang 分支逆极限的秩二地址几何、单向量 affine 实现的模论障碍以及 Zeckendorf--Gödel Dirichlet 链的矩阵 Euler 刚性压缩为同一条闭合链：单向量模型可以精确生成一条 \(2\)-进地址线及其分形维数，却不能在模意义下覆盖 Lee--Yang 的二维 dyadic 分支层化；与此同时，Zeckendorf--Gödel 约束在解析侧所产生的 Euler 结构既保持了不可压缩的相邻耦合，又允许在 continuant 递推层面被完全标量化追踪。

\paragraph{dyadic multiplicity 的复杂度屏障与周期化 Fourier 的绝对连续性排斥}
在局部化数系的连通 Lie 商刚性与 Zeckendorf--Gödel Dirichlet 级数的 \(s=1\) 密度--留数定理已经确立之后，还可沿 dyadic multiplicity 与 Fibonacci 共振两条链进一步抽取出下列若干条后果。

\begin{theorem}[固定元数 dyadic multiplicity 的 Presburger 谓词类属于 \texorpdfstring{$\mathbf P$}{P}]\label{thm:xi-fixed-arity-dyadic-multiplicity-presburger-in-p}
\leanverified{paper\_xi\_fixed\_arity\_dyadic\_multiplicity\_presburger\_in\_p}
固定整数 \(r\ge 1\)，并固定一个 Presburger 公式
\[
\Phi=\Phi\bigl(m;V_1,\epsilon_1,d_1;\dots;V_r,\epsilon_r,d_r\bigr).
\]
定义语言
\[
L_\Phi
:=
\left\{
\bigl(m,w^{(1)},\dots,w^{(r)}\bigr):
w^{(i)}\in X_m,\ 
\Phi\!\bigl(m;V_m(w^{(1)}),w_m^{(1)},d_m^{\mathrm{bin}}(w^{(1)});\dots;V_m(w^{(r)}),w_m^{(r)},d_m^{\mathrm{bin}}(w^{(r)})\bigr)
\right\}.
\]
则
\[
L_\Phi\in \mathbf P.
\]
特别地，若存在 many-one 归约
\[
\mathrm{SAT}\le_m^p L_\Phi,
\]
则必有
\[
\mathbf P=\mathbf{NP}.
\]
\end{theorem}
\begin{proof}
设输入为
\[
\bigl(m,w^{(1)},\dots,w^{(r)}\bigr).
\]
由于 \(r\) 固定，输入总长度与 \(m\) 同阶。先逐个 \(i=1,\dots,r\) 扫描字 \(w^{(i)}\)，在线性时间内检验其是否属于 \(X_m\)，并同时计算
\[
V_m\bigl(w^{(i)}\bigr),
\qquad
\epsilon_i:=w_m^{(i)}.
\]
再由命题 \ref{prop:fold-bin-digit-dp}（亦即结论 \ref{con:xi-terminal-foldbin-digit-dp-linear-time} 的算法核），可在 \(O(m)\) 时间内精确计算
\[
d_i:=d_m^{\mathrm{bin}}\bigl(w^{(i)}\bigr).
\]
因此全部参数
\[
m,\ V_i,\ \epsilon_i,\ d_i
\qquad(1\le i\le r)
\]
都可在总时间 \(O(m)\) 内求得，并且其二进制长度均为 \(O(m)\)。

由于 \(\Phi\) 是固定的 Presburger 公式，可对其预先执行一次量词消去，得到一个仅依赖于 \(\Phi\) 的等价量词自由公式；该公式是有限个线性等式、不等式与同余条件的布尔组合。于是，对上述 \(O(m)\)-比特整数代入求值可在多项式时间内完成。故 \(L_\Phi\) 可由多项式时间算法判定，即
\[
L_\Phi\in\mathbf P.
\]

最后，若 \(\mathrm{SAT}\le_m^p L_\Phi\)，则由 \(L_\Phi\in \mathbf P\) 立得 \(\mathrm{SAT}\in \mathbf P\)，从而
\[
\mathbf P=\mathbf{NP}.
\]
证毕。
\end{proof}

\begin{theorem}[黄金参数二点 Bernoulli 卷积的周期化 Fourier 系数不属于任何有限 \texorpdfstring{$\ell^r$}{lr}，因而不具有任何 \texorpdfstring{$L^p$}{Lp} 密度]\label{thm:xi-bernoulli-two-periodization-no-lp-density}
令
\[
q:=\varphi^{-1},
\]
并令 \(\mu_q^{(2)}\) 如定理 \ref{thm:fold-pisot-bernoulli-convolution-representation}。记
\[
\nu:=\bigl(\RR\to \mathbb T_2=\RR/(2\ZZ)\bigr)_\#\mu_q^{(2)}
\]
为其 \(2\)-周期化测度。则对任意
\[
1\le r<\infty,
\]
\[
\widehat\nu\notin \ell^r(\ZZ).
\]
特别地，对任意
\[
1<p\le\infty,
\]
\(\nu\) 都不可能具有 \(L^p(\mathbb T_2)\) 密度。
\end{theorem}
\begin{proof}
由定理 \ref{thm:xi-pisot-bernoulli-fibonacci-direction-nonrajchman}，对任意非零整数初值对
\[
(u_0,u_1)\in\ZZ^2\setminus\{(0,0)\},
\]
若 \((u_n)_{n\ge 0}\) 满足 Fibonacci 递推
\[
u_{n+2}=u_{n+1}+u_n,
\]
则存在常数 \(c(u_0,u_1)>0\) 使
\[
\bigl|\widehat\nu(u_n)\bigr|
=
\bigl|\widehat{\mu_q^{(2)}}(\pi u_n)\bigr|
\longrightarrow c(u_0,u_1),
\qquad
|u_n|\to\infty.
\]
固定任意一组此类非零初值对，并记 \(c:=c(u_0,u_1)\)。则存在 \(N_0\) 使得对一切 \(n\ge N_0\)，都有
\[
\bigl|\widehat\nu(u_n)\bigr|\ge \frac{c}{2}.
\]
因此对任意 \(1\le r<\infty\)，
\[
\sum_{u\in\ZZ}\bigl|\widehat\nu(u)\bigr|^r
\ge
\sum_{n\ge N_0}\bigl|\widehat\nu(u_n)\bigr|^r
\ge
\sum_{n\ge N_0}\left(\frac{c}{2}\right)^r
=+\infty.
\]
故 \(\widehat\nu\notin \ell^r(\ZZ)\) 对所有有限 \(r\) 都成立。

再反设 \(\nu\) 具有某个 \(L^p(\mathbb T_2)\) 密度。若 \(1<p\le 2\)，则 Hausdorff--Young 不等式给出
\[
\widehat\nu\in \ell^{p'}(\ZZ),
\qquad
\frac1p+\frac1{p'}=1,
\]
与上式矛盾。若 \(2<p\le\infty\)，则由于 \(\mathbb T_2\) 的 Haar 测度有限，有
\[
L^p(\mathbb T_2)\subseteq L^2(\mathbb T_2),
\]
从而 Plancherel 定理给出 \(\widehat\nu\in \ell^2(\ZZ)\)，同样矛盾。故 \(\nu\) 不可能具有任何 \(L^p(\mathbb T_2)\) 密度。证毕。
\end{proof}

\begin{theorem}[黄金参数二点 Bernoulli 卷积的周期化测度在全部 \texorpdfstring{$L^r$}{Lr} 标度上均无密度]\label{thm:xi-bernoulli-two-periodization-no-any-lr-density}
沿用定理 \ref{thm:xi-bernoulli-two-periodization-no-lp-density} 的记号。则对任意
\[
1\le r\le \infty,
\]
\(\nu\) 都不可能具有 \(L^r(\mathbb T_2)\) 密度。
\end{theorem}
\begin{proof}
当 \(1<r\le \infty\) 时，这正是定理 \ref{thm:xi-bernoulli-two-periodization-no-lp-density} 的结论。

只剩 \(r=1\)。若 \(\nu\) 具有 \(L^1(\mathbb T_2)\) 密度，则由 Riemann--Lebesgue 引理必有
\[
\widehat\nu(u)\longrightarrow 0
\qquad (|u|\to\infty).
\]
然而定理 \ref{thm:xi-pisot-bernoulli-fibonacci-direction-nonrajchman} 已给出沿任意非平凡 Fibonacci 递推逃逸序列 \((u_n)\) 都有
\[
|u_n|\to\infty,
\qquad
\bigl|\widehat\nu(u_n)\bigr|\longrightarrow c(u_0,u_1)>0.
\]
这与 Riemann--Lebesgue 引理矛盾。故 \(L^1\) 密度亦不存在。综上，\(\nu\) 在全部 \(L^r\) 标度上都无密度。证毕。
\end{proof}

\paragraph{Gödel--Tate 地址动力系统、Ahlfors 正则与前缀证书预算}
在单向量 \(2^L\)-进全息像集的精确自相似、Lee--Yang 分支逆极限的秩障碍以及 Zeckendorf--Gödel 矩阵 Euler 刚性已经建立之后，还可把地址像本身的动力系统、维数测度的精确锐性、最小进位阈值、规范前缀刚性与有限前缀的外显证书预算进一步压缩为下列诸条结论。

\begin{conclusion}[Gödel--Tate 地址像作为严格 \texorpdfstring{$B$}{B}-进全移位系统的显式 \texorpdfstring{$\zeta$}{zeta} 函数与熵]\label{con:xi-time-part62c-godel-tate-address-fullshift-zeta-entropy}
沿用结论 \ref{con:xi-time-part62b-rankone-hologram-selfsimilar-measure-dimension} 的记号，并进一步假设
\[
v\in T_2(E_{\min})\setminus 2T_2(E_{\min}).
\]
记
\[
N:=|\mathcal E|=|D|,
\qquad
X:\mathcal E^{\NN}\longrightarrow \mathcal X_D\subset T_v
\]
为该结论中的地址映射，\(\sigma:\mathcal E^{\NN}\to \mathcal E^{\NN}\) 为单侧全移位
\[
\sigma(e_1,e_2,e_3,\dots)=(e_2,e_3,\dots).
\]
则：
\begin{enumerate}
  \item 每个 \(x\in\mathcal X_D\) 都唯一写成
  \[
  x=dv+By
  \qquad (d\in D,\ y\in\mathcal X_D).
  \]
  因而映射
  \[
  S_B:\mathcal X_D\to \mathcal X_D,
  \qquad
  S_B(dv+By):=y
  \]
  良定义且连续。等价地，当
  \[
  x\in dv+B\mathcal X_D
  \qquad (d\in D)
  \]
  时，有分支公式
  \[
  S_B(x)=\frac{x-dv}{B}.
  \]
  \item \(X\) 把 \((\mathcal E^{\NN},\sigma)\) 与 \((\mathcal X_D,S_B)\) 拓扑共轭：
  \[
  X\circ \sigma=S_B\circ X.
  \]
  特别地，\((\mathcal X_D,S_B)\) 是一个严格的 \(N\)-符号单侧全移位系统。
  \item 对每个 \(n\ge 1\)，有
  \[
  \#\operatorname{Fix}(S_B^n)=N^n.
  \]
  因而其 Artin--Mazur \(\zeta\) 函数与拓扑熵满足
  \[
  \zeta_{S_B}(z)
  :=
  \exp\!\left(\sum_{n\ge 1}\frac{\#\operatorname{Fix}(S_B^n)}{n}z^n\right)
  =
  \frac{1}{1-Nz},
  \qquad
  h_{\mathrm{top}}(S_B)=\log N.
  \]
\end{enumerate}
\end{conclusion}
\begin{proof}
由结论 \ref{con:xi-time-part62b-rankone-hologram-selfsimilar-measure-dimension} 的第 \(1\) 条，
\[
\mathcal X_D=\bigsqcup_{d\in D}(dv+B\mathcal X_D),
\]
且该并严格不交，因此每个 \(x\in\mathcal X_D\) 都唯一写成 \(x=dv+By\) 的形式，其中 \(d\in D\)、\(y\in\mathcal X_D\)。故 \(S_B\) 良定义。

再由定理 \ref{thm:xi-time-part9g-holographic-prefix-isometry-on-line}，映射 \(X\) 把事件流的前缀超度量等距地送到 \(\mathcal X_D\) 上的 \(B\)-进超度量，因而是从紧空间 \(\mathcal E^{\NN}\) 到 Hausdorff 空间 \(\mathcal X_D\) 的同胚。对任意 \(\mathbf e=(e_t)_{t\ge 1}\in\mathcal E^{\NN}\)，按定义
\[
X(\mathbf e)=\mathrm{code}(e_1)v+B\,X(\sigma\mathbf e),
\]
于是
\[
S_B\bigl(X(\mathbf e)\bigr)=X(\sigma\mathbf e),
\]
即
\[
X\circ \sigma=S_B\circ X.
\]
因此 \(S_B=X\circ \sigma\circ X^{-1}\) 连续，且 \((\mathcal X_D,S_B)\) 与 \(N\)-符号单侧全移位共轭。

对每个 \(n\ge 1\)，\(\sigma^n\) 的不动点恰为长度 \(n\) 纯周期块的无限重复，故
\[
\#\operatorname{Fix}(\sigma^n)=N^n.
\]
由共轭性立得
\[
\#\operatorname{Fix}(S_B^n)=N^n.
\]
于是
\[
\sum_{n\ge 1}\frac{\#\operatorname{Fix}(S_B^n)}{n}z^n
=
\sum_{n\ge 1}\frac{N^n}{n}z^n
=
-\log(1-Nz),
\]
指数化得
\[
\zeta_{S_B}(z)=\frac{1}{1-Nz}.
\]
又单侧全移位的拓扑熵为 \(\log N\)，故由共轭性
\[
h_{\mathrm{top}}(S_B)=\log N.
\]
证毕。
\end{proof}

\begin{corollary}[Gödel--Tate 地址动力系统的 primitive 周期轨道数闭式]\label{cor:xi-time-part62c-godel-tate-address-primitive-periodic-count}
沿用结论 \ref{con:xi-time-part62c-godel-tate-address-fullshift-zeta-entropy} 的记号。对每个 \(n\ge 1\)，记
\[
\operatorname{Prim}_{S_B}(n)
\]
为 \((\mathcal X_D,S_B)\) 中最小周期恰为 \(n\) 的 primitive 周期轨道条数。则有闭式
\[
\operatorname{Prim}_{S_B}(n)
=
\frac{1}{n}\sum_{d\mid n}\mu(d)\,N^{n/d}.
\]
等价地，最小周期恰为 \(n\) 的周期点总数为
\[
\sum_{d\mid n}\mu(d)\,N^{n/d}.
\]
\leanverified{paper\_xi\_time\_part62c\_godel\_tate\_address\_primitive\_periodic\_count}
\end{corollary}
\begin{proof}
对任意离散动力系统，长度 \(n\) 的全部不动点按其最小周期分层，满足恒等式
\[
\#\operatorname{Fix}(S_B^n)
=
\sum_{d\mid n} d\,\operatorname{Prim}_{S_B}(d).
\]
由结论 \ref{con:xi-time-part62c-godel-tate-address-fullshift-zeta-entropy}，
\[
\#\operatorname{Fix}(S_B^n)=N^n.
\]
故
\[
N^n=\sum_{d\mid n} d\,\operatorname{Prim}_{S_B}(d).
\]
对该整除卷积施行 Möbius 反演，即得
\[
n\,\operatorname{Prim}_{S_B}(n)=\sum_{d\mid n}\mu(d)\,N^{n/d},
\]
亦即所示公式。最后一式只是把每条 primitive \(n\)-周期轨道含有恰好 \(n\) 个周期点这一事实重写。证毕。
\end{proof}

\begin{conclusion}[Gödel--Tate 地址集的精确 Hausdorff/Minkowski 维数与 Ahlfors 正则性]\label{con:xi-time-part62c-godel-tate-address-ahlfors-regularity}
沿用结论 \ref{con:xi-time-part62c-godel-tate-address-fullshift-zeta-entropy} 的记号，并在 \(\mathcal X_D\) 上赋予定理 \ref{thm:xi-time-part9g-holographic-prefix-isometry-on-line} 中的 \(B\)-进超度量
\[
d_B(x,y)=
\begin{cases}
0,& x=y,\\[2pt]
|B|_2^{\,\nu_B(x-y)},& x\neq y,
\end{cases}
\qquad
\nu_B(z):=\max\{n\ge 0:\ z\in B^nT_v\}.
\]
记 \(\nu_D\) 为均匀 Bernoulli 测度在地址映射下的推前：
\[
\nu_D
:=
X_*\!\left(
\left(\frac1N\sum_{e\in \mathcal E}\delta_e\right)^{\!\otimes\NN}
\right).
\]
则：
\begin{enumerate}
  \item \(\mathcal X_D\) 的 Hausdorff 维数与 Minkowski 维数一致，并满足
  \[
  \dim_{\mathrm H}(\mathcal X_D)
  =
  \dim_{\mathrm M}(\mathcal X_D)
  =
  \frac{\log N}{\log B}.
  \]
  \item 令 \(s:=\frac{\log N}{\log B}\)。则 \(\nu_D\) 是 \((\mathcal X_D,d_B)\) 上的 \(s\)-Ahlfors 正则测度：存在常数 \(c_1,c_2>0\)，使得对所有 \(x\in\mathcal X_D\) 与所有充分小的 \(r>0\)，若记 \(\mathrm B_{\mathcal X_D}(x,r)\) 为 \(\mathcal X_D\) 中半径 \(r\) 的闭球，则
  \[
  c_1r^s\le \nu_D\bigl(\mathrm B_{\mathcal X_D}(x,r)\bigr)\le c_2r^s.
  \]
  事实上，可取
  \[
  c_1=N^{-1},
  \qquad
  c_2=N.
  \]
  \item 若 \(N<B\)，则 \(\mathcal X_D\) 在地址线 \(T_v\cong\ZZ_2\) 中 Haar 测度为 \(0\)。
\end{enumerate}
\end{conclusion}
\begin{proof}
第 \(1\) 条已经由结论 \ref{con:xi-time-part62b-rankone-hologram-selfsimilar-measure-dimension} 的第 \(3\) 条给出。故只需证明第 \(2\) 条。

对任意长度 \(n\) 前缀
\[
a=(e_1,\dots,e_n)\in\mathcal E^n,
\]
记其 cylinder 为
\[
[a]:=\{\mathbf e\in\mathcal E^{\NN}:\ \mathrm{pref}_n(\mathbf e)=a\}.
\]
由定理 \ref{thm:xi-time-part9g-holographic-prefix-isometry-on-line}，
\[
X_n(a):=\sum_{t=1}^{n}\mathrm{code}(e_t)\,B^{t-1}v,
\qquad
X([a])=\bigl(X_n(a)+B^nT_v\bigr)\cap \mathcal X_D,
\]
而这正是 \(\mathcal X_D\) 中一个半径 \(B^{-n}\) 的闭球。另一方面，均匀 Bernoulli 测度给出
\[
\nu_D\bigl(X([a])\bigr)=N^{-n}=(B^{-n})^s.
\]

现取任意 \(x\in\mathcal X_D\) 与充分小的 \(r>0\)。存在唯一整数 \(n\ge 0\) 使得
\[
B^{-(n+1)}<r\le B^{-n}.
\]
由于超度量的球只在离散半径 \(B^{-m}\) 处发生变化，遂有包含关系
\[
\mathrm B_{\mathcal X_D}(x,B^{-(n+1)})
\subseteq
\mathrm B_{\mathcal X_D}(x,r)
\subseteq
\mathrm B_{\mathcal X_D}(x,B^{-n}),
\]
并且左右两端都是某个长度分别为 \(n+1\) 与 \(n\) 的 cylinder 像。因此
\[
N^{-(n+1)}
\le
\nu_D\bigl(\mathrm B_{\mathcal X_D}(x,r)\bigr)
\le
N^{-n}.
\]
又因 \(s=\log_B N\)，有
\[
N^{-(n+1)}=B^{-(n+1)s},
\qquad
N^{-n}=B^{-ns}.
\]
同时由 \(B^{-(n+1)}<r\le B^{-n}\) 得
\[
B^{-(n+1)s}<r^s\le B^{-ns}.
\]
于是
\[
N^{-1}r^s\le \nu_D\bigl(\mathrm B_{\mathcal X_D}(x,r)\bigr)\le Nr^s.
\]
这就证明了 \(s\)-Ahlfors 正则性。

第 \(3\) 条亦可由结论 \ref{con:xi-time-part62b-rankone-hologram-selfsimilar-measure-dimension} 的第 \(2\) 条直接得到；等价地，当 \(N<B\) 时有 \(s<1\)，故 \(\mathcal X_D\) 作为 \(T_v\cong\ZZ_2\) 中的真自相似子集具有零 Haar 测度。证毕。
\end{proof}

\begin{conclusion}[faithful Gödel--Tate 编码的最小 bit 宽度与地址线饱和阈值]\label{con:xi-time-part62c-godel-tate-minimal-bitwidth-threshold}
取整数 \(M\ge 0\) 与 \(B=2^L\ge 2\)，并固定原始向量
\[
v\in T_2(E_{\min})\setminus 2T_2(E_{\min}).
\]
定义 digit 集与地址映射
\[
\mathcal A_M:=\{0,1,\dots,M\},
\qquad
\Xi_{B,M}:\mathcal A_M^{\NN}\longrightarrow T_v,
\]
\[
\Xi_{B,M}(a_1,a_2,\dots):=\sum_{t\ge 1}a_tB^{t-1}v.
\]
则：
\begin{enumerate}
  \item \(\Xi_{B,M}\) 为单射当且仅当
  \[
  B>M.
  \]
  因而 faithful Gödel--Tate 实现所需的最小 bit 宽度精确为
  \[
  L_{\min}=\left\lceil \log_2(M+1)\right\rceil.
  \]
  \item 当 \(M=B-1\) 时，地址像恰好充满整条 \(B\)-进地址线：
  \[
  \Xi_{B,B-1}\bigl(\mathcal A_{B-1}^{\NN}\bigr)=T_v=\ZZ_2v.
  \]
  \item 当 \(M<B-1\) 时，\(\Xi_{B,M}(\mathcal A_M^{\NN})\) 是 \(T_v\) 的真 Cantor 型子集，并且在 \(T_v\) 中 Haar 测度为 \(0\)。
\end{enumerate}
\end{conclusion}
\begin{proof}
先证第 \(1\) 条。若 \(B>M\)，则 \(\mathcal A_M\subset\{0,\dots,B-1\}\)。若两列 digit 在首个分歧位置 \(n\) 处不同，则其像之差可写成
\[
B^{n-1}(c_n+Bu)v,
\]
其中 \(u\in\ZZ_2\) 且
\[
c_n\in\{-(B-1),\dots,B-1\}\setminus\{0\}.
\]
由于 \(0<|c_n|<B\)，有 \(c_n\notin B\ZZ_2\)，从而 \(c_n+Bu\notin B\ZZ_2\)。再由 \(T_v=\ZZ_2v\) 对 \(\ZZ_2\) 无挠，得上式不可能为 \(0\)。故 \(\Xi_{B,M}\) 单射。

反之，若 \(B\le M\)，则两列
\[
(B,0,0,\dots),
\qquad
(0,1,0,0,\dots)
\]
都属于 \(\mathcal A_M^{\NN}\)，且
\[
\Xi_{B,M}(B,0,0,\dots)=Bv=\Xi_{B,M}(0,1,0,0,\dots).
\]
故 \(\Xi_{B,M}\) 不单射。这就证明了
\[
\Xi_{B,M}\ \text{单射}
\quad\Longleftrightarrow\quad
B>M.
\]
由于 \(B=2^L\)，最小的 \(L\) 恰是满足 \(2^L>M\) 的最小整数，即
\[
L_{\min}=\left\lceil\log_2(M+1)\right\rceil.
\]

现证第 \(2\) 条。当 \(M=B-1\) 时，digit 集恰为
\[
D=\{0,\dots,B-1\},
\]
故由结论 \ref{con:xi-time-part62b-rankone-hologram-selfsimilar-measure-dimension} 的第 \(2\) 条，地址像正好等于整条地址线 \(T_v\)。

最后设 \(M<B-1\)。此时 \(N:=M+1<B\)。由结论 \ref{con:xi-time-part62c-godel-tate-address-ahlfors-regularity}，
\[
\dim_{\mathrm H}\bigl(\Xi_{B,M}(\mathcal A_M^{\NN})\bigr)
=
\frac{\log(M+1)}{\log B}<1.
\]
因此该像不可能等于一维地址线 \(T_v\cong\ZZ_2\)，并且其 Haar 测度为 \(0\)。证毕。
\end{proof}

\begin{conclusion}[faithful Gödel--Tate 规范编码的精确 \texorpdfstring{$B$}{B}-进前缀刚性]\label{con:xi-time-part62c-godel-tate-prefix-rigidity}
沿用定理 \ref{thm:killo-godel-semidir-affine-2adic} 的记号。记
\[
H:=\mathcal R\rtimes_\Sigma \NN,
\qquad
H_{\mathrm{can}}
:=
\{(r,k)\in H:\ \mathrm{supp}(r)\subseteq\{p_1,\dots,p_k\},\ 0\le r(p_t)\le M\},
\]
并在地址线 \(T_v=\ZZ_2v\subset T_2(E_{\min})\) 上定义
\[
\operatorname{ev}_B(r):=\sum_{t\ge 1}r(p_t)B^{t-1}v.
\]
对不同的 \(r,s\in H_{\mathrm{can}}\)，定义首个素数分歧位置
\[
\iota(r,s):=\min\{t\ge 1:\ r(p_t)\neq s(p_t)\}.
\]
又对非零元素 \(z\in T_v\) 定义 \(B\)-进赋值
\[
\nu_B(z):=\max\{n\ge 0:\ z\in B^nT_v\}.
\]
则有精确恒等式
\[
\boxed{\;
\nu_B\bigl(\operatorname{ev}_B(r)-\operatorname{ev}_B(s)\bigr)
=
\iota(r,s)-1.\;}
\]
等价地，若定义
\[
d_B(r,s):=
\begin{cases}
0,& r=s,\\[2pt]
|B|_2^{\,\nu_B(\operatorname{ev}_B(r)-\operatorname{ev}_B(s))},& r\neq s,
\end{cases}
\]
则
\[
\boxed{\;
d_B(r,s)=|B|_2^{\,\iota(r,s)-1}
\qquad (r\neq s).\;}
\]
因此，规范编码
\[
r\longmapsto \operatorname{ev}_B(r)
\]
把 prime-register 的前缀超度量等距嵌入到地址线 \(T_v\) 的 \(B\)-进超度量中。
\end{conclusion}
\begin{proof}
由结论 \ref{con:xi-time-part62c-godel-tate-minimal-bitwidth-threshold}，faithful 规范编码等价于 \(B>M\)。取不同的 \(r,s\in H_{\mathrm{can}}\)，并记
\[
t_0:=\iota(r,s).
\]
则
\[
\operatorname{ev}_B(r)-\operatorname{ev}_B(s)
=
\sum_{t\ge 1}(r(p_t)-s(p_t))B^{t-1}v
\]
\[
=
B^{t_0-1}
\left(
r(p_{t_0})-s(p_{t_0})
+
\sum_{t>t_0}(r(p_t)-s(p_t))B^{t-t_0}
\right)v.
\]
记括号中的 \(\ZZ_2\)-系数为 \(u_{r,s}\)。由于
\[
0\le r(p_t),s(p_t)\le M<B,
\]
首项
\[
c_0:=r(p_{t_0})-s(p_{t_0})
\]
满足
\[
-(B-1)\le c_0\le B-1,
\qquad
c_0\neq 0.
\]
故 \(c_0\notin B\ZZ_2\)。而其余各项都属于 \(B\ZZ_2\)，于是
\[
u_{r,s}\equiv c_0\pmod{B\ZZ_2},
\]
从而 \(u_{r,s}\notin B\ZZ_2\)。这就说明
\[
u_{r,s}v\in T_v\setminus BT_v.
\]
因此
\[
\operatorname{ev}_B(r)-\operatorname{ev}_B(s)
\in
B^{t_0-1}T_v\setminus B^{t_0}T_v,
\]
亦即
\[
\nu_B\bigl(\operatorname{ev}_B(r)-\operatorname{ev}_B(s)\bigr)=t_0-1.
\]
第二个盒装公式只是按定义改写。证毕。
\end{proof}

\begin{conclusion}[有限前缀 cylinder 的最小 oracle 预算]\label{con:xi-time-part62c-godel-tate-prefix-oracle-budget}
沿用结论 \ref{con:xi-time-part62c-godel-tate-address-fullshift-zeta-entropy} 的记号。对每个 \(n\ge 1\) 与每个前缀
\[
a=(e_1,\dots,e_n)\in\mathcal E^n
\]
定义其截断地址
\[
X_n(a):=\sum_{t=1}^{n}\mathrm{code}(e_t)\,B^{t-1}v\in T_v.
\]
则：
\begin{enumerate}
  \item 映射
  \[
  a\longmapsto X_n(a)\bmod B^nT_v
  \]
  从 \(\mathcal E^n\) 到 \(T_v/B^nT_v\) 为单射；其像恰由
  \[
  N^n
  \]
  个两两不同的剩余类组成。
  \item 因而，若要把长度 \(n\) 的 cylinder 指数外显为一个有限二进制 oracle 证书，则任何无误方案至少需要
  \[
  \left\lceil n\log_2N\right\rceil
  \]
  个 bit。
  \item 上述下界可达：存在注入
  \[
  \mathrm{enc}_n:\mathcal E^n\hookrightarrow \{0,1\}^{\lceil n\log_2N\rceil},
  \]
  因而长度 \(n\) 前缀的外显 oracle 预算精确等于
  \[
  \left\lceil n\log_2N\right\rceil .
  \]
\end{enumerate}
\end{conclusion}
\begin{proof}
由定理 \ref{thm:xi-time-part9g-holographic-prefix-isometry-on-line}，
\[
X([a])=\bigl(X_n(a)+B^nT_v\bigr)\cap \mathcal X_D.
\]
若 \(a\neq b\in\mathcal E^n\)，则长度 \(n\) 的两个 cylinders 不相交，故对应的 \(B^nT_v\)-剩余类不同。于是
\[
a\longmapsto X_n(a)\bmod B^nT_v
\]
为单射。由于 \(|\mathcal E^n|=N^n\)，其像恰含 \(N^n\) 个不同剩余类，从而第 \(1\) 条成立。

若一个二进制 oracle 证书长度为 \(m\)，则其可能输出至多为 \(2^m\) 种。要想无歧义地区分 \(N^n\) 个不同长度 \(n\) 前缀，必须有
\[
2^m\ge N^n,
\]
即
\[
m\ge \left\lceil n\log_2N\right\rceil.
\]
这证明了第 \(2\) 条。

反过来，有限集 \(\mathcal E^n\) 的基数正是 \(N^n\)，故总可选取一个注入
\[
\mathrm{enc}_n:\mathcal E^n\hookrightarrow \{0,1\}^{\lceil n\log_2N\rceil}.
\]
因此长度 \(n\) 前缀确可由恰好 \(\lceil n\log_2N\rceil\) 个 bit 的外显 oracle 证书无误标识。证毕。
\end{proof}

\paragraph{Gödel--Tate 仿射作用的闭不变核与轨道闭包}
在 faithful \(2\)-进仿射实现、地址线前缀刚性与有限 oracle 预算已经建立之后，还可把整个 Tate 模块上的闭不变加法结构与轨道闭包几何继续压缩为如下两条刚性结论。它们表明：尽管作用发生在二维 \(2\)-进模块 \(T_2(E_{\min})\) 上，真正的极小核仍唯一地坍缩到由地址向量 \(v\) 所张成的一条 \(2\)-进直线。

\begin{theorem}[Gödel--Tate 仿射作用的闭不变加法子群完全分类]\label{thm:xi-godel-tate-affine-closed-subgroup-classification}
沿用定理 \ref{thm:xi-time-part9g-affine-semigroup-unique-fixed-point} 的 \(T_2(E_{\min})\cong \ZZ_2^2\) 识别。固定
\[
B:=2^L,
\qquad
\mathcal R:=\NN^{(\mathbb P)},
\qquad
H:=\mathcal R\rtimes_\Sigma \NN,
\]
并取原始向量
\[
v\in T:=T_2(E_{\min})\setminus 2T.
\]
定义
\[
L_v:=\ZZ_2v\subset T,
\qquad
\operatorname{ev}_B(r):=\sum_{t\ge 1}r(p_t)B^{t-1}v,
\]
\[
\Psi(r,k)(x):=B^k x+\operatorname{ev}_B(r).
\]
再令
\[
\pi:T\longrightarrow T/L_v\cong \ZZ_2
\]
为自然商映射，并对每个 \(n\in\NN\) 定义
\[
M_n:=\pi^{-1}(2^n\ZZ_2),
\qquad
M_\infty:=L_v.
\]
则 \(\Psi(H)\)-不变的闭加法子群恰为
\[
\boxed{\;
M_0\supsetneq M_1\supsetneq M_2\supsetneq \cdots \supsetneq M_\infty=L_v,\;}
\]
除此之外再无其他闭不变加法子群。特别地，
\[
\boxed{\;\{0\}\ \text{不是}\ \Psi(H)\text{-不变的}.\;}
\]
\end{theorem}
\begin{proof}
先证
\[
\overline{\operatorname{ev}_B(\mathcal R)}=L_v.
\]
对任意 \(r\in\mathcal R\)，\(\operatorname{ev}_B(r)\) 显然属于 \(L_v\)，故
\[
\operatorname{ev}_B(\mathcal R)\subseteq L_v.
\]
反之，对每个 \(n\in\NN\)，取 \(r_n\in\mathcal R\) 满足
\[
r_n(p_1)=n,
\qquad
r_n(p)=0\ \ (p\neq p_1).
\]
则
\[
\operatorname{ev}_B(r_n)=nv.
\]
由于 \(\NN\) 在 \(\ZZ_2\) 中稠密，\(\{nv:n\in\NN\}\) 在 \(L_v=\ZZ_2v\) 中稠密，故上式成立。

现设 \(M\subset T\) 为非空闭加法子群，且对所有 \((r,k)\in H\) 都有
\[
\Psi(r,k)(M)\subseteq M.
\]
由于 \(0\in M\)，取 \(k=0\) 得
\[
\operatorname{ev}_B(r)=\Psi(r,0)(0)\in M
\qquad (r\in\mathcal R).
\]
由闭性与上一步稠密性，
\[
L_v\subseteq M.
\]

因此 \(\bar M:=\pi(M)\) 是 \(\ZZ_2\) 的闭加法子群。又对任意 \(x\in T\)、
\((r,k)\in H\)，有
\[
\pi(\Psi(r,k)(x))
=
\pi(B^k x+\operatorname{ev}_B(r))
=
B^k\pi(x),
\]
因为 \(\operatorname{ev}_B(r)\in L_v\)。故 \(\bar M\) 对乘法
\[
y\longmapsto B^k y
\]
不变。

而 \(\ZZ_2\) 的闭加法子群恰为
\[
2^n\ZZ_2\qquad (n=0,1,2,\dots,\infty),
\]
其中 \(2^\infty\ZZ_2:=0\)。这些子群都被 \(B^k=2^{Lk}\) 保持，故
\[
\bar M=2^n\ZZ_2
\]
对某个 \(n\in\{0,1,2,\dots,\infty\}\) 成立，从而
\[
M=\pi^{-1}(\bar M)=\pi^{-1}(2^n\ZZ_2).
\]
这就得到所有可能的闭不变加法子群恰为
\[
M_n\qquad (n=0,1,2,\dots,\infty).
\]

反过来，每个 \(M_n\) 都显然为闭加法子群；并且 \(L_v\subseteq M_n\)，故平移部分保持
\(M_n\)；商上又有
\[
\pi(\Psi(r,k)(M_n))=B^k\pi(M_n)\subseteq 2^n\ZZ_2,
\]
故 \(M_n\) 确为 \(\Psi(H)\)-不变。

最后，\(\{0\}\) 不是不变子群，因为对前述 \(r_1\) 有
\[
\Psi(r_1,0)(0)=v\neq 0.
\]
证毕。
\end{proof}

\begin{theorem}[Gödel--Tate 仿射作用的轨道闭包与唯一极小集]\label{thm:xi-godel-tate-affine-orbit-closure-minimal-set}
沿用定理 \ref{thm:xi-godel-tate-affine-closed-subgroup-classification} 的记号。则对任意 \(x\in T\)，其
\(\Psi(H)\)-轨道闭包满足精确公式
\[
\boxed{\;
\overline{\Psi(H)\cdot x}
=
L_v\ \cup\ \bigcup_{n\ge 0}\bigl(B^n x+L_v\bigr).\;}
\]
从而：
\begin{enumerate}
  \item 若 \(x\in L_v\)，则
  \[
  \overline{\Psi(H)\cdot x}=L_v;
  \]
  \item \(L_v\) 是 \(\Psi(H)\) 作用下唯一的极小非空闭不变集；
  \item 任意连续 \(\Psi(H)\)-不变函数
  \[
  f:T\to\CC
  \]
  必为常数。
\end{enumerate}
\end{theorem}
\begin{proof}
对每个固定 \(n\ge 0\)，有
\[
\Psi(\mathcal R,n)(x):=\{\Psi(r,n)(x):r\in\mathcal R\}
=
B^n x+\operatorname{ev}_B(\mathcal R).
\]
由定理 \ref{thm:xi-godel-tate-affine-closed-subgroup-classification} 证明开头的稠密性，
\[
\overline{\Psi(\mathcal R,n)(x)}=B^n x+L_v.
\]
因此
\[
L_v\cup\bigcup_{n\ge 0}(B^n x+L_v)\subseteq \overline{\Psi(H)\cdot x}.
\]

另一方面，考虑商映射 \(\pi:T\to T/L_v\cong\ZZ_2\)。由
\[
\pi(\Psi(r,n)(x))=B^n\pi(x)
\]
可知 \(\Psi(H)\cdot x\) 在商上的像恰为
\[
\{B^n\pi(x):n\ge 0\}.
\]
由于 \(|B|_2<1\)，有
\[
B^n\pi(x)\longrightarrow 0
\qquad (n\to\infty)
\]
于 \(\ZZ_2\) 中成立，故集合
\[
\{0\}\cup\{B^n\pi(x):n\ge 0\}
\]
是闭的。其原像正是
\[
L_v\cup\bigcup_{n\ge 0}(B^n x+L_v),
\]
从而右端本身已闭，并包含整个轨道；于是得到所示闭包公式。

若 \(x\in L_v\)，则对每个 \(n\ge 0\) 都有 \(B^n x+L_v=L_v\)，故
\[
\overline{\Psi(H)\cdot x}=L_v.
\]
这证明了第 \(1\) 条。

再设 \(Y\subset T\) 为任意非空闭不变集，取 \(x\in Y\)。由轨道闭包公式，
\[
L_v\subseteq \overline{\Psi(H)\cdot x}\subseteq Y.
\]
因此每个非空闭不变集都包含 \(L_v\)。而第 \(1\) 条表明 \(L_v\) 自身就是非空闭不变集，故它是唯一极小集。这证明了第 \(2\) 条。

最后设 \(f:T\to\CC\) 连续且 \(\Psi(H)\)-不变。由不变性，\(f\) 在每条轨道上常数；由连续性，\(f\) 在每条轨道闭包上仍常数。对任意 \(y\in L_v\)，第 \(1\) 条给出
\[
\overline{\Psi(H)\cdot y}=L_v,
\]
故 \(f\) 在 \(L_v\) 上为常数。再对任意 \(x\in T\)，其轨道闭包包含 \(L_v\)，于是 \(f(x)\) 必等于 \(L_v\) 上的该常值。故 \(f\) 在整个 \(T\) 上常数。证毕。
\end{proof}
\paragraph{事件词的自由仿射单子与前缀仿射像的精确 Haar 体积}
在忠实的 Gödel--Tate 规范编码的前缀刚性与有限证书预算已经确定之后，事件词的连接本身还可在环境 Tate 模块上被直接重写为一个自由仿射子单子；与此同时，每个有限前缀所切出的环境仿射像也具有完全显式的 Haar 体积律。这样，符号连接、仿射作用与 \(2^L\)-进体积便在同一层面被统一为可直接计算的有限证书。

\begin{theorem}[事件词在 \texorpdfstring{$T_2(E_{\min})$}{T2(Emin)} 上生成自由仿射子单子]\label{thm:xi-time-part62cc-eventword-free-affine-submonoid}
沿用定理 \ref{thm:killo-godel-semidir-affine-2adic} 的记号。设 \(E\) 为有限事件集，
\[
\mathrm{code}:E\hookrightarrow \NN,
\qquad
M:=\max_{e\in E}\mathrm{code}(e),
\qquad
B:=2^L>M,
\]
并记
\[
T:=T_2(E_{\min})\cong \ZZ_2^2,
\qquad
v\in T\setminus 2T.
\]
对任意有限事件词
\[
a=a_1\cdots a_n\in E^\ast
\]
定义前缀向量与仿射映射
\[
x_a:=\sum_{t=1}^{n}\mathrm{code}(a_t)B^{t-1}v,
\qquad
F_a(x):=x_a+B^n x
\qquad(x\in T).
\]
并约定空词 \(\varnothing\) 满足 \(x_{\varnothing}=0\)、\(F_{\varnothing}=\mathrm{id}_T\)。再对无限事件流
\[
\mathbf e=(e_t)_{t\ge 1}\in E^{\NN}
\]
定义全息点
\[
X(\mathbf e):=\sum_{t\ge 1}\mathrm{code}(e_t)B^{t-1}v.
\]
则有：
\begin{enumerate}
  \item 对任意 \(a,b\in E^\ast\)，都有
  \[
  F_a\circ F_b=F_{ab},
  \]
  其中 \(ab\) 表示词连接；
  \item 对任意 \(a\in E^\ast\) 与 \(\mathbf e\in E^{\NN}\)，都有
  \[
  X(a\mathbf e)=F_a\bigl(X(\mathbf e)\bigr);
  \]
  \item 映射
  \[
  E^\ast\hookrightarrow \mathrm{Aff}(T),
  \qquad
  a\longmapsto F_a
  \]
  为单射。
\end{enumerate}
因此，\(\{F_a:a\in E^\ast\}\) 在 \(\mathrm{Aff}(T)\) 中构成一个与事件词自由单子 \(E^\ast\) 规范同构的自由仿射子单子。
\end{theorem}
\begin{proof}
设 \(a\in E^m\)、\(b\in E^n\)。由定义，
\[
x_{ab}
=
\sum_{t=1}^{m}\mathrm{code}(a_t)B^{t-1}v
+
\sum_{s=1}^{n}\mathrm{code}(b_s)B^{m+s-1}v
=
x_a+B^m x_b.
\]
故对任意 \(x\in T\) 都有
\[
F_a(F_b(x))
=
x_a+B^m(x_b+B^n x)
=
(x_a+B^m x_b)+B^{m+n}x
=
F_{ab}(x),
\]
从而得到第 \(1\) 条。

第 \(2\) 条同样直接计算即可：
\[
X(a\mathbf e)
=
\sum_{t=1}^{m}\mathrm{code}(a_t)B^{t-1}v
+
\sum_{s\ge 1}\mathrm{code}(e_s)B^{m+s-1}v
=
x_a+B^m X(\mathbf e)
=
F_a(X(\mathbf e)).
\]

最后证明单射。若 \(F_a=F_b\)，设 \(|a|=m\)、\(|b|=n\)。先取 \(x=0\)，得
\[
x_a=x_b.
\]
再对任意 \(x\in T\) 比较
\[
F_a(x)-F_a(0)=B^m x,
\qquad
F_b(x)-F_b(0)=B^n x,
\]
遂有
\[
B^m x=B^n x
\qquad(\forall x\in T).
\]
由于 \(T\cong \ZZ_2^2\) 无挠且 \(B>1\)，故必有 \(m=n\)。现设公共长度为 \(n\)。若 \(a\neq b\)，记 \(j\) 为首个满足 \(a_j\neq b_j\) 的位置，则
\[
x_a-x_b
=
B^{j-1}
\left(
\mathrm{code}(a_j)-\mathrm{code}(b_j)
+
\sum_{t>j}\bigl(\mathrm{code}(a_t)-\mathrm{code}(b_t)\bigr)B^{t-j}
\right)v.
\]
其中首项
\[
c_j:=\mathrm{code}(a_j)-\mathrm{code}(b_j)
\]
满足
\[
-(B-1)\le c_j\le B-1,
\qquad
c_j\neq 0,
\]
故 \(c_j\notin B\ZZ_2\)。而其余各项都落在 \(B\ZZ_2\) 中，因此括号内的 \(2\)-进系数不属于 \(B\ZZ_2\)。由 \(v\notin 2T\) 可知 \(v\notin BT\)，于是
\[
x_a-x_b\in B^{j-1}T\setminus B^jT,
\]
特别地 \(x_a-x_b\neq 0\)，与 \(x_a=x_b\) 矛盾。故只能有 \(a=b\)。单射性得证。证毕。
\end{proof}
\leanverified{paper\_xi\_time\_part62cc\_eventword\_free\_affine\_submonoid}

\begin{corollary}[前缀仿射像的精确 Haar 体积律]\label{cor:xi-time-part62cc-prefix-affine-haar-volume}
\leanverified{paper\_xi\_time\_part62cc\_prefix\_affine\_haar\_volume}
沿用定理 \ref{thm:xi-time-part62cc-eventword-free-affine-submonoid} 的记号。对任意长度 \(n\) 的事件词
\[
a\in E^n
\]
定义其前缀仿射像
\[
U_a:=F_a(T)=x_a+B^nT\subset T.
\]
若 \(m_T\) 为 \(T\cong \ZZ_2^2\) 上的归一化 Haar 概率测度，则：
\begin{enumerate}
  \item 每个 \(U_a\) 都是子群 \(B^nT\) 的一个陪集，因而亦即半径 \(B^{-n}\) 的闭球，并满足
  \[
  m_T(U_a)=B^{-2n};
  \]
  \item 对固定的 \(n\)，若 \(a\neq b\in E^n\)，则
  \[
  U_a\cap U_b=\varnothing;
  \]
  \item 因而长度 \(n\) 前缀的总体环境阴影满足
  \[
  m_T\!\left(\bigcup_{a\in E^n}U_a\right)=|E|^n B^{-2n}.
  \]
\end{enumerate}
\end{corollary}
\begin{proof}
由定义
\[
U_a=x_a+B^nT
\]
显然是 \(B^nT\) 的陪集。又 \(T\cong \ZZ_2^2\) 为秩 \(2\) 的自由 \(\ZZ_2\)-模，故
\[
T/B^nT\cong (\ZZ/2^{Ln}\ZZ)^2,
\]
从而
\[
|T:B^nT|=2^{2Ln}=B^{2n}.
\]
于是每个陪集的 Haar 质量都等于指数倒数 \(B^{-2n}\)，得到第 \(1\) 条。

再证第 \(2\) 条。若 \(U_a\cap U_b\neq\varnothing\)，则同为 \(B^nT\) 的陪集，必有
\[
U_a=U_b,
\qquad
x_a-x_b\in B^nT.
\]
若 \(a\neq b\)，记 \(j\le n\) 为其首个分歧位置。与上一定理证明中的同一分解完全相同，可得
\[
x_a-x_b\in B^{j-1}T\setminus B^jT.
\]
由于 \(j\le n\)，这尤其推出
\[
x_a-x_b\notin B^nT,
\]
与上式矛盾。故只能有 \(a=b\)，于是不同长度 \(n\) 词对应的 \(U_a\) 两两不交。

第 \(3\) 条于是由 \(|E|^n\) 个两两不交陪集的体积相加立即得到：
\[
m_T\!\left(\bigcup_{a\in E^n}U_a\right)
=
\sum_{a\in E^n}m_T(U_a)
=
|E|^n B^{-2n}.
\]
证毕。
\end{proof}

\noindent
因此，事件词的时间前缀在 Tate 模块中并非只是被动记录，而是以自由仿射单子的方式主动生成其环境几何；而每个长度 \(n\) 的前缀又恰好切出一个 Haar 质量为 \(B^{-2n}\) 的仿射球。于是，词连接的组合增长与 \(2^L\)-进环境体积之间便获得了完全显式的乘法对应。


\noindent
上述诸条结论把 H.159--H.160 所给出的忠实仿射实现进一步压缩为同一条严格的地址动力学链：单向量 \(B\)-进地址像不仅是一个精确自相似的非阿基米德 Cantor 系统，而且与全移位严格共轭，从而同时拥有显式的 Artin--Mazur \(\zeta\) 函数、拓扑熵与 primitive 周期轨道计数闭式；与此同时，Ahlfors 正则维数律与最小 bit-width 阈值又把该系统的几何尺度与编码门槛完全锁定。进一步，规范编码在 \(H_{\mathrm{can}}\) 上并非仅仅忠实，而是把首个素数分歧位置严格翻译为地址线上的 \(B\)-进赋值；在有限深度上，长度 \(n\) 的 cylinder 族又恰好给出 \(N^n\) 个彼此分离的剩余类，并在整个环境 Tate 模块中切出 Haar 质量为 \(B^{-2n}\) 的显式仿射球，从而把前缀信息的外显证书预算与环境体积同时锁定在精确计数律之下；而当作用提升为事件词的自由仿射子单子并进一步延伸到整个 Tate 模块 \(T_2(E_{\min})\) 时，全部闭不变加法子群与全部轨道闭包又最终坍缩到由 \(L_v=\ZZ_2v\) 生成的一条唯一极小地址核上。

\paragraph{\texorpdfstring{$2^L$}{2^L}-进全息地址、有限 Grothendieck 秩障碍与 Lee--Yang 二次分辨率的接口压缩}
在单向量 \(2^L\)-进全息像集的精确自相似、地址动力系统的维数刚性、可消去幺半群的 Grothendieck 完成判据以及 Lee--Yang 三次数域链的共同三次基域都已确定之后，还可把这三条接口进一步压缩为下列五条直接后果。

\begin{theorem}[全息地址像的严格不交自相似分解]\label{thm:xi-time-part62ca-holographic-address-strict-selfsimilarity}
沿用结论 \ref{con:xi-time-part62b-rankone-hologram-selfsimilar-measure-dimension} 与结论 \ref{con:xi-time-part62c-godel-tate-address-fullshift-zeta-entropy} 的记号。设 \(E\) 为有限事件字母表，\(\mathrm{code}:E\hookrightarrow\{0,\dots,M\}\subset\ZZ\) 为单射，\(B=2^L>M\)，并记
\[
D:=\mathrm{code}(E),
\qquad
X(\mathbf e):=\sum_{t\ge 1}\mathrm{code}(e_t)\,B^{t-1}v
\qquad
(\mathbf e=(e_t)_{t\ge 1}\in E^{\NN}),
\]
\[
\mathcal X_E:=X(E^{\NN})\subset T_v.
\]
则 \(\mathcal X_E\) 满足严格的不交自相似分解
\[
\mathcal X_E
=
\bigsqcup_{a\in D}\bigl(av+B\mathcal X_E\bigr).
\]
更精确地，每一块 \(av+B\mathcal X_E\) 都是 \(\mathcal X_E\) 与某个 \(BT_v\)-陪集的交，因而在 \(\mathcal X_E\) 的相对拓扑下是 clopen。
\end{theorem}
\begin{proof}
由结论 \ref{con:xi-time-part62b-rankone-hologram-selfsimilar-measure-dimension} 的第 \(1\) 条，对 \(D=\mathrm{code}(E)\) 已有
\[
\mathcal X_D=\bigsqcup_{d\in D}\bigl(dv+B\mathcal X_D\bigr),
\]
且每一块都可写成 \(\mathcal X_D\) 与某个 \(BT_v\)-陪集的交。由于 \(\mathrm{code}\) 为单射，\(|D|=|E|\)，并且按定义 \(\mathcal X_E=\mathcal X_D\)。代换记号即得所述分解。证毕。
\end{proof}

\begin{theorem}[全息地址像的 Hausdorff 与 Minkowski 维数完全一致]\label{thm:xi-time-part62ca-holographic-address-exact-dimension}
沿用定理 \ref{thm:xi-time-part62ca-holographic-address-strict-selfsimilarity} 的记号，并在 \(\mathcal X_E\) 上赋予 \(B\)-进超度量
\[
d_B(x,y):=
\begin{cases}
0,& x=y,\\[2pt]
B^{-\nu_B(x-y)},& x\neq y,
\end{cases}
\qquad
\nu_B(z):=\max\{n\ge 0:\ z\in B^nT_v\}.
\]
则
\[
\dim_{\mathrm H}(\mathcal X_E)
=
\underline{\dim}_{\mathrm M}(\mathcal X_E)
=
\overline{\dim}_{\mathrm M}(\mathcal X_E)
=
\frac{\log |E|}{\log B}.
\]
\end{theorem}
\begin{proof}
由定理 \ref{thm:xi-time-part62ca-holographic-address-strict-selfsimilarity}，长度 \(n\) 的前缀给出 \(|E|^n\) 个两两不交的 cylinder，而每个 cylinder 都恰是某个 \(B^nT_v\)-陪集与 \(\mathcal X_E\) 的交，因此其直径正好为 \(B^{-n}\)。故尺度 \(B^{-n}\) 下的最少覆盖数满足
\[
N(B^{-n})=|E|^n.
\]
于是
\[
\underline{\dim}_{\mathrm M}(\mathcal X_E)
=
\overline{\dim}_{\mathrm M}(\mathcal X_E)
=
\lim_{n\to\infty}\frac{\log |E|^n}{-\log B^{-n}}
=
\frac{\log |E|}{\log B}.
\]

另一方面，结论 \ref{con:xi-time-part62c-godel-tate-address-ahlfors-regularity} 已给出 \(\mathcal X_E\) 上的等权 Bernoulli 推前测度是
\[
s:=\frac{\log |E|}{\log B}
\]
维 Ahlfors 正则测度，故由质量分布原理
\[
\dim_{\mathrm H}(\mathcal X_E)=s.
\]
与上式合并即得结论。证毕。
\end{proof}

\begin{corollary}[任意有限值连续观测都因子化于有限前缀]\label{cor:xi-time-part62ca-finite-valued-continuous-observables-prefix-factorization}
\leanverified{paper\_xi\_time\_part62ca\_finite\_valued\_continuous\_observables\_prefix\_factorization}
沿用定理 \ref{thm:xi-time-part62ca-holographic-address-strict-selfsimilarity} 的记号。设 \(A\) 为有限离散集，且
\[
F:\mathcal X_E\to A
\]
连续。则存在某个整数 \(n\ge 1\) 与映射
\[
\overline F:E^n\to A
\]
使得对任意 \(\mathbf e=(e_t)_{t\ge 1}\in E^{\NN}\) 都有
\[
F\!\bigl(X(\mathbf e)\bigr)=\overline F(e_1,\dots,e_n).
\]
等价地，存在映射
\[
\widetilde F:T_v/B^nT_v\to A
\]
使
\[
F=\widetilde F\circ \pi_n|_{\mathcal X_E},
\]
其中 \(\pi_n:T_v\to T_v/B^nT_v\) 为自然商映射。
\end{corollary}
\begin{proof}
由结论 \ref{con:xi-time-part62c-godel-tate-address-fullshift-zeta-entropy}，地址映射
\[
X:E^{\NN}\longrightarrow \mathcal X_E
\]
是同胚。故
\[
G:=F\circ X:E^{\NN}\to A
\]
连续。由于 \(A\) 离散，\(G^{-1}(a)\) 对每个 \(a\in A\) 都是 clopen。再因 \(E^{\NN}\) 的拓扑基由有限前缀 cylinder 构成，且 \(E^{\NN}\) 紧致，存在某个统一层数 \(n\) 使得每个长度 \(n\) 的 cylinder 都完全包含在某一条纤维 \(G^{-1}(a)\) 中。于是 \(G\) 在每个长度 \(n\) cylinder 上常值，可定义
\[
\overline F(e_1,\dots,e_n):=G(\mathbf e)
\]
其中 \(\mathbf e\) 为任意以 \((e_1,\dots,e_n)\) 为前缀的无限延拓；该定义与延拓无关。由 \(G=F\circ X\) 即得
\[
F\!\bigl(X(\mathbf e)\bigr)=\overline F(e_1,\dots,e_n).
\]

再由结论 \ref{con:xi-time-part62c-godel-tate-prefix-oracle-budget} 的第 \(1\) 条，长度 \(n\) 的前缀与 \(T_v/B^nT_v\) 中落在 \(\mathcal X_E\) 上的相应剩余类一一对应。故可先定义
\[
\widetilde F_0:\pi_n(\mathcal X_E)\to A,
\]
再把 \(\widetilde F_0\) 向整个有限商 \(T_v/B^nT_v\) 任意延拓为 \(\widetilde F\)，于是满足
\[
F=\widetilde F\circ \pi_n|_{\mathcal X_E}.
\]
证毕。
\end{proof}

\begin{theorem}[有限 Grothendieck 有理秩目标中不存在乘法对象的同态嵌入]\label{thm:xi-time-part62ca-no-multiplicative-embedding-into-finite-grothendieck-rank}
\leanverified{paper\_xi\_time\_part62ca\_no\_multiplicative\_embedding\_into\_finite\_grothendieck\_rank}
设 \(M\) 为交换可消去幺半群，并满足
\[
\operatorname{rank}_{\QQ}\!\bigl(G(M)\otimes_{\ZZ}\QQ\bigr)<\infty.
\]
则不存在幺半群单射
\[
\NN_{\times}\hookrightarrow M.
\]
\end{theorem}
\begin{proof}
反设存在单射幺半群同态
\[
f:\NN_{\times}\hookrightarrow M.
\]
由命题 \ref{prop:killo-grothendieck-completion-preserves-injection}，Grothendieck 完成后得到群单射
\[
G(f):G(\NN_{\times})\hookrightarrow G(M).
\]
而由定理 \ref{thm:killo-prime-freedom-non-finitizability} 的第 \(1\) 条，
\[
G(\NN_{\times})\cong \bigoplus_{p\in\mathbb P}\ZZ.
\]
再与 \(\QQ\) 张量。由于 \(\QQ\) 作为 \(\ZZ\)-模是平坦的，单射 \(G(f)\) 仍保持为单射，从而得到
\[
G(\NN_{\times})\otimes_{\ZZ}\QQ
\cong
\bigoplus_{p\in\mathbb P}\QQ,
\]
其 \(\QQ\)-维数可数无穷。于是 \(G(f)\otimes_{\ZZ}\QQ\) 给出一个从可数无穷维 \(\QQ\)-向量空间到
\[
G(M)\otimes_{\ZZ}\QQ
\]
的单射，这与后者具有有限 \(\QQ\)-维数矛盾。故所述单射不存在。证毕。
\end{proof}

\begin{theorem}[Lee--Yang 三次分支不变量共享同一唯一二次分辨子域]\label{thm:xi-time-part62ca-leyang-branch-invariants-common-quadratic-resolvent}
沿用定理 \ref{thm:xi-terminal-zm-stokes-leyang-shared-artin-representation}、推论 \ref{cor:xi-terminal-zm-stokes-leyang-common-quadratic-resolvent} 与定理 \ref{thm:xi-time-part62d-leyang-signed-amplitude-exact-degree6} 的记号，记
\[
u:=\kappa_*^2,
\qquad
b:=B^2,
\qquad
K:=\QQ(u)=\QQ(b)=\QQ(\lambda_*).
\]
设 \(L_K\) 为三次数域 \(K/\QQ\) 的伽罗瓦闭包。则 \(u\)、\(b\) 与 \(\lambda_*\) 的三次最小多项式的分裂域全都等于 \(L_K\)，并且共有同一个唯一非平凡二次子域
\[
\QQ(\sqrt{-111}).
\]
等价地，Lee--Yang 三次核 \(P_{\mathrm{LY}}(y)=256y^3+411y^2+165y+32\) 所定义分裂域的唯一二次分辨子域，恰与三元分支不变量 \((u,b,\lambda_*)\) 的共同二次分辨子域一致。
\end{theorem}
\begin{proof}
由定理 \ref{thm:xi-terminal-zm-stokes-leyang-shared-artin-representation}，
\[
\QQ(u)=\QQ(b)=\QQ(\lambda_*)=K.
\]
因此 \(u\)、\(b\) 与 \(\lambda_*\) 的三次最小多项式都生成同一非伽罗瓦三次数域 \(K\)，其分裂域必同为 \(L_K\)。

再由推论 \ref{cor:xi-terminal-zm-stokes-leyang-common-quadratic-resolvent}，\(L_K/\QQ\) 的唯一二次子域为
\[
\QQ(\sqrt{-111}).
\]
另一方面，定理 \ref{thm:xi-time-part9g-leyang-cubic-discriminant} 已证明 Lee--Yang 三次核 \(P_{\mathrm{LY}}(y)\) 的分裂域同样具有唯一二次子域
\[
\QQ(\sqrt{-111}).
\]
故 Lee--Yang 三次核与分支不变量 \((u,b,\lambda_*)\) 在二次分辨层上严格汇合到同一个域
\[
\QQ(\sqrt{-111}).
\]
证毕。
\end{proof}
\paragraph{有理玫瑰读出、universal solenoid 极限与虚部可见层的核消隐}
在有限素数局部化群的 prime-register solenoid、无穷 prime-register 的有限可见因子化，以及 RH 的 Carath\'eodory--Toeplitz 视界接口均已建立之后，还可把有理玫瑰族、Lissajous 族与分母扫描的几何读出压缩为下列一条连续链。其核心是：经典平面曲线可视为字符读出在圆商与奇异环之间的阴影，而相位补全后的扫描坐标则把分母层级严格识别为 universal solenoid 的有限层因子。

\begin{theorem}[有限 prime-register solenoid 上的 Kronecker 读出与有理玫瑰/Lissajous 的统一提升]\label{thm:xi-time-part62d-solenoid-kronecker-rose-lissajous-lift}
沿用定理 \ref{thm:killo-localization-circle-dimension-prime-ledger-splitting} 的记号。设 \(S\subset\mathbb P\) 为有限素数集，并记
\[
G_S:=\ZZ[S^{-1}],
\qquad
\Sigma_S:=\widehat{G_S},
\qquad
0\longrightarrow K_S:=\prod_{p\in S}\ZZ_p
\longrightarrow
\Sigma_S
\xrightarrow{\ \pi\ }
\TT
\longrightarrow 0.
\]
定义 Kronecker 读出轨道
\[
\kappa_S:\RR\longrightarrow \Sigma_S,
\qquad
\bigl(\kappa_S(t)\bigr)(q):=e^{2\pi iqt}
\qquad (q\in G_S).
\]
并对每个 \(q\in G_S\) 记字符
\[
\chi_q:\Sigma_S\to\TT,
\qquad
\chi_q(x):=x(q).
\]
则成立：
\begin{enumerate}
  \item \(\kappa_S\) 是连续单射群同态，且 \(\overline{\kappa_S(\RR)}=\Sigma_S\)。
  \item 对任意 \(q\in G_S\)，有理玫瑰轨道
  \[
  \gamma_q(\theta):=\cos(q\theta)e^{i\theta}
  \qquad (\theta\in\RR)
  \]
  满足严格恒等式
  \[
  \gamma_q(\theta)
  =
  \frac12\Bigl(
  \chi_{1+q}\bigl(\kappa_S(\theta/2\pi)\bigr)
  +
  \chi_{1-q}\bigl(\kappa_S(\theta/2\pi)\bigr)
  \Bigr).
  \]
  \item 对任意整数 \(a,b\in\ZZ\) 与相位 \(\delta\in\RR\)，Lissajous 轨道
  \[
  L_{a,b,\delta}(t):=(\cos(at+\delta),\cos(bt))
  \]
  满足
  \[
  L_{a,b,\delta}(t)
  =
  \Bigl(
  \Re\bigl(e^{i\delta}\chi_a(\kappa_S(t/2\pi))\bigr),\
  \Re\bigl(\chi_b(\kappa_S(t/2\pi))\bigr)
  \Bigr).
  \]
\end{enumerate}
\end{theorem}
\begin{proof}
把 \(\Sigma_S\) 视为 \(\TT^{G_S}\) 的闭子群，则 \(\kappa_S\) 的每个坐标函数
\[
t\longmapsto e^{2\pi iqt}
\qquad (q\in G_S)
\]
都连续，故 \(\kappa_S\) 连续；又显然
\[
\kappa_S(t+t')=\kappa_S(t)\kappa_S(t'),
\]
故其为群同态。

若 \(\kappa_S(t)=\kappa_S(t')\)，则对 \(q=1\) 有
\[
e^{2\pi i(t-t')}=1,
\]
从而 \(t-t'\in\ZZ\)。再对任意
\[
N=\prod_{p\in S}p^{k_p}
\]
取 \(q=1/N\in G_S\)，得到
\[
e^{2\pi i(t-t')/N}=1,
\]
故 \(t-t'\in N\ZZ\)。由于此结论对所有这样的 \(N\) 同时成立，只能有 \(t-t'=0\)。因此 \(\kappa_S\) 单射。

对稠密性，只需计算其湮没子。若 \(q\in G_S\) 满足
\[
\chi_q(\kappa_S(t))=1
\qquad
(\forall\, t\in\RR),
\]
则
\[
e^{2\pi iqt}=1
\qquad
(\forall\, t\in\RR),
\]
于是必有 \(q=0\)。故 \(\kappa_S(\RR)\) 的湮没子平凡。由 Pontryagin 对偶中“闭子群由其湮没子确定”的标准结论，\(\overline{\kappa_S(\RR)}=\Sigma_S\)。

对第 \(2\) 条，直接代入
\[
\chi_{1\pm q}\bigl(\kappa_S(\theta/2\pi)\bigr)=e^{i(1\pm q)\theta}
\]
即得
\[
\frac12\Bigl(e^{i(1+q)\theta}+e^{i(1-q)\theta}\Bigr)
=
e^{i\theta}\frac{e^{iq\theta}+e^{-iq\theta}}{2}
=
\cos(q\theta)e^{i\theta}.
\]

对第 \(3\) 条，同理有
\[
\chi_a(\kappa_S(t/2\pi))=e^{iat},
\qquad
\chi_b(\kappa_S(t/2\pi))=e^{ibt}.
\]
取实部便得到
\[
\Re\bigl(e^{i\delta}\chi_a(\kappa_S(t/2\pi))\bigr)=\cos(at+\delta),
\qquad
\Re\bigl(\chi_b(\kappa_S(t/2\pi))\bigr)=\cos(bt).
\]
证毕。
\end{proof}

\begin{proposition}[universal solenoid 的相位补全分母扫描与有限层因子化]\label{prop:xi-time-part62d-universal-solenoid-phase-completed-scan}
记
\[
\Sigma_{\mathbb P}:=\widehat{\QQ}.
\]
则有短正合列
\[
0\longrightarrow \prod_{p\in\mathbb P}\ZZ_p
\longrightarrow \Sigma_{\mathbb P}
\xrightarrow{\ \pi\ }
\TT
\longrightarrow 0,
\]
并且
\[
\Sigma_{\mathbb P}
\cong
\varprojlim_{N\ge 1}\widehat{(1/N)\ZZ}
\cong
\varprojlim_{N\ge 1}\TT,
\]
其中当 \(N\mid M\) 时连接映射为 \(z\mapsto z^{M/N}\)。

对任意 \(q\in\QQ\)，定义相位补全后的扫描坐标
\[
\Phi_q^{\cos}(x):=\frac12\bigl(\chi_{1+q}(x)+\chi_{1-q}(x)\bigr),
\qquad
\Phi_q^{\sin}(x):=\frac{1}{2i}\bigl(\chi_{1+q}(x)-\chi_{1-q}(x)\bigr).
\]
则成立：
\begin{enumerate}
  \item 对每个 \(q\in\QQ\)，有
  \[
  \Phi_q^{\cos}(x)+i\Phi_q^{\sin}(x)=\chi_{1+q}(x).
  \]
  \item 映射
  \[
  \Phi^{\mathrm{ph}}:\Sigma_{\mathbb P}\longrightarrow \prod_{q\in\QQ}\CC^2,
  \qquad
  x\longmapsto \bigl((\Phi_q^{\cos}(x),\Phi_q^{\sin}(x))\bigr)_{q\in\QQ}
  \]
  是拓扑嵌入。
  \item 若 \(E\subset\QQ\) 为有限集，且对每个 \(q\in E\) 记 \(d(q)\) 为其既约正分母，并令
  \[
  N_E:=\operatorname{lcm}\{d(q):\ q\in E\},
  \qquad
  H_{N_E}:=(1/N_E)\ZZ,
  \]
  则有限投影
  \[
  \Phi_E^{\mathrm{ph}}:\Sigma_{\mathbb P}\longrightarrow \prod_{q\in E}\CC^2
  \]
  仅依赖于自然商映射
  \[
  \pi_{N_E}:\Sigma_{\mathbb P}\twoheadrightarrow \widehat{H_{N_E}}\cong \TT.
  \]
\end{enumerate}
因此，任意有限分母扫描表都是 universal solenoid 的有限层因子观测，而全部相容有限层的逆极限恰恢复 \(\Sigma_{\mathbb P}\)。
\end{proposition}
\begin{proof}
由
\[
0\longrightarrow \ZZ\longrightarrow \QQ\longrightarrow \QQ/\ZZ\longrightarrow 0,
\qquad
\QQ/\ZZ\cong \bigoplus_{p\in\mathbb P}C_{p^\infty},
\]
对偶化即得所示短正合列。

再由
\[
\QQ=\varinjlim_{N\ge 1}(1/N)\ZZ
\]
（偏序由整除关系给出），Pontryagin 对偶把该直极限送为逆极限，从而
\[
\Sigma_{\mathbb P}
\cong
\varprojlim_{N\ge 1}\widehat{(1/N)\ZZ}.
\]
把 \(\widehat{(1/N)\ZZ}\) 识别为 \(\TT\) 后，若 \(N\mid M\)，则包含
\[
(1/N)\ZZ\hookrightarrow (1/M)\ZZ
\]
的对偶映射满足
\[
z\longmapsto z^{M/N}.
\]

第 \(1\) 条由定义立即得到：
\[
\Phi_q^{\cos}(x)+i\Phi_q^{\sin}(x)
=
\frac12\bigl(\chi_{1+q}(x)+\chi_{1-q}(x)\bigr)
+\frac12\bigl(\chi_{1+q}(x)-\chi_{1-q}(x)\bigr)
=
\chi_{1+q}(x).
\]

对第 \(2\) 条，由上一式可知 \(\Phi^{\mathrm{ph}}(x)\) 决定全部字符
\[
\chi_r(x)
\qquad (r\in\QQ),
\]
因为可取 \(r=1+q\)。而 \(\QQ\) 正是 \(\Sigma_{\mathbb P}\) 的对偶群，其字符分离点，故 \(\Phi^{\mathrm{ph}}\) 单射。连续性显然，而 \(\Sigma_{\mathbb P}\) 紧、\(\prod_{q\in\QQ}\CC^2\) Hausdorff，因此连续单射自动是到像上的同胚，即拓扑嵌入。

对第 \(3\) 条，若 \(q\in E\)，则 \(q\in H_{N_E}\)，从而
\[
1\pm q\in H_{N_E}.
\]
故 \(\chi_{1\pm q}\) 都只依赖于 \(\Sigma_{\mathbb P}\) 在 \(\widehat{H_{N_E}}\) 上的投影，也即只依赖于 \(\pi_{N_E}\)。于是 \(\Phi_q^{\cos}\) 与 \(\Phi_q^{\sin}\) 都经由 \(\pi_{N_E}\) 因子化，有限投影 \(\Phi_E^{\mathrm{ph}}\) 亦然。

最后，由 \(\QQ=\bigcup_{N\ge 1}(1/N)\ZZ\)，任一有理分母观测都落在某一有限层 \(H_N\) 中；而全部这些有限层通过整除偏序相容拼接，便恢复整个对偶群 \(\QQ\)，从而恢复 \(\Sigma_{\mathbb P}\)。证毕。
\end{proof}

\begin{corollary}[有理玫瑰的一阶字符承载对象由分母素数支撑唯一锁定]\label{cor:xi-time-part62d-rational-rose-minimal-prime-carrier}
\leanverified{paper\_xi\_time\_part62d\_rational\_rose\_minimal\_prime\_carrier}
设 \(q=n/d\in\QQ\) 为既约有理数，\(d\ge 1\)，并记 \(S(q)\subset\mathbb P\) 为 \(d\) 的素因子集合。则在定理 \ref{thm:xi-time-part62d-solenoid-kronecker-rose-lissajous-lift} 的一阶字符读出口径下，玫瑰轨道
\[
\gamma_q(\theta)=\cos(q\theta)e^{i\theta}
\]
可由 \(\Sigma_{S(q)}\) 实现，且不能由任何严格更小的 \(\Sigma_{S'}\)（\(S'\subsetneq S(q)\)）实现。并且
\[
p\in S(q)
\iff
\prod_{r\in S(q)}\ZZ_r\ \text{的}\ p\text{-主因子非平凡}.
\]
因此，最小承载对象的 prime-register 核完全恢复分母的素数支撑。
\end{corollary}
\begin{proof}
由定理 \ref{thm:xi-time-part62d-solenoid-kronecker-rose-lissajous-lift}，\(\gamma_q\) 的读出公式只使用字符 \(\chi_{1+q}\) 与 \(\chi_{1-q}\)。这两者在 \(\Sigma_S=\widehat{\ZZ[S^{-1}]}\) 上有定义，当且仅当
\[
1+q,\ 1-q\in \ZZ[S^{-1}],
\]
等价地，当且仅当 \(q\in \ZZ[S^{-1}]\)。而对既约有理数 \(q=n/d\)，这又等价于 \(d\) 的全部素因子都属于 \(S\)。故最小可行素数集恰为 \(S(q)\)。

若 \(S'\subsetneq S(q)\)，则 \(q\notin \ZZ[S'^{-1}]\)，从而上述一阶字符公式在 \(\Sigma_{S'}\) 上无法成立，因此 \(\Sigma_{S(q)}\) 在该读出口径下确为最小承载对象。

最后，由定理 \ref{thm:killo-localization-circle-dimension-prime-ledger-splitting}，\(\Sigma_{S(q)}\) 的核
\[
K_{S(q)}=\prod_{p\in S(q)}\ZZ_p
\]
满足
\[
p\in S(q)
\iff
K_{S(q)}\ \text{的}\ p\text{-主因子非平凡}.
\]
故分母的素数支撑可由 prime-register 核完全恢复。证毕。
\end{proof}

\begin{theorem}[有理玫瑰轨道的最小周期与花瓣数分类]\label{thm:xi-time-part62d-rational-rose-period-petal-classification}
\leanverified{paper\_xi\_time\_part62d\_rational\_rose\_period\_petal\_classification}
设 \(q=n/d\) 为正既约有理数，其中 \(n,d\in\NN\)。定义
\[
\gamma_q(\theta):=\cos\!\Bigl(\frac{n}{d}\theta\Bigr)e^{i\theta}.
\]
并把
\[
\gamma_q(\RR)\setminus\{0\}
\]
的连通分支闭包称为该轨道的花瓣。则其最小周期 \(T(q)\) 与花瓣数 \(P(q)\) 满足
\[
T(q)=
\begin{cases}
\pi d, & n\equiv d\equiv 1\pmod 2,\\[1mm]
2\pi d, & \text{否则},
\end{cases}
\qquad
P(q)=
\begin{cases}
n, & n\equiv d\equiv 1\pmod 2,\\[1mm]
2n, & \text{否则}.
\end{cases}
\]
\end{theorem}
\begin{proof}
把 \(\gamma_q\) 改写为
\[
\gamma_q(\theta)
=
\frac12\Bigl(
e^{i(1+n/d)\theta}+e^{i(1-n/d)\theta}
\Bigr).
\]
令
\[
\alpha:=d+n,
\qquad
\beta:=d-n.
\]
若 \(T>0\) 是其周期，则必须同时满足
\[
e^{i\alpha T/d}=1,
\qquad
e^{i\beta T/d}=1.
\]
故最小正周期为
\[
T(q)
=
\frac{2\pi d}{\gcd(\alpha,\beta)}
=
\frac{2\pi d}{\gcd(d+n,d-n)}
=
\frac{2\pi d}{\gcd(d+n,2n)}.
\]
由于 \((n,d)=1\)，有
\[
\gcd(d+n,2n)=\gcd(d+n,2).
\]
因此当且仅当 \(d+n\) 为偶数时分母为 \(2\)。又 \((n,d)=1\) 下，\(d+n\) 为偶数等价于 \(n,d\) 同为奇数。于是得到所示 \(T(q)\)。

现计数花瓣。原点恰在
\[
\theta_k:=\frac{d(2k+1)\pi}{2n}
\qquad (k\in\ZZ)
\]
处出现，因为此时 \(\cos(n\theta/d)=0\)。因此每个开区间
\[
I_k:=(\theta_k,\theta_{k+1})
\]
的像都落在 \(\gamma_q(\RR)\setminus\{0\}\) 的某个连通分支中，而全部花瓣恰由这些区间像的闭包给出。

再注意到步长 \(\pi d/n\) 的平移满足
\[
\gamma_q\!\Bigl(\theta+\frac{\pi d}{n}\Bigr)
=
\cos\!\Bigl(\frac{n}{d}\theta+\pi\Bigr)e^{i\theta}e^{i\pi d/n}
=
\omega\,\gamma_q(\theta),
\]
其中
\[
\omega:=-e^{i\pi d/n}=e^{i\pi(d+n)/n}.
\]
故相邻零区间上的像互为绕原点的旋转；全部花瓣正是一朵基准花瓣在循环群 \(\langle \omega\rangle\) 作用下的轨道。于是花瓣数等于 \(\omega\) 的阶。

而 \(\omega^m=1\) 当且仅当
\[
\frac{m(d+n)}{n}\in 2\ZZ.
\]
又因 \((n,d)=1\)，有 \(\gcd(d+n,n)=1\)。因此最小这样的 \(m\) 为
\[
m=
\begin{cases}
n, & d+n\ \text{为偶数},\\
2n, & d+n\ \text{为奇数}.
\end{cases}
\]
结合前述奇偶判别，即得
\[
P(q)=
\begin{cases}
n, & n\equiv d\equiv 1\pmod 2,\\
2n, & \text{否则}.
\end{cases}
\]
证毕。
\end{proof}

\begin{proposition}[Lissajous 自交点作为 Klein 四元反演像的轨道交会]\label{prop:xi-time-part62d-lissajous-klein-orbit-intersection}
设整数对 \((a,b)\neq (0,0)\)，相位 \(\delta\in\RR\)，并记
\[
L_{a,b,\delta}(t):=(\cos(at+\delta),\cos(bt)).
\]
定义其环面提升
\[
\widetilde L_{a,b,\delta}(t):=(e^{i(at+\delta)},e^{ibt})\in\TT^2,
\]
投影
\[
\Pi:\TT^2\to[-1,1]^2,
\qquad
\Pi(z_1,z_2):=(\Re z_1,\Re z_2),
\]
以及 Klein 四元群
\[
V_4:=\{(\varepsilon_1,\varepsilon_2):\ \varepsilon_j\in\{\pm 1\}\}
\]
在 \(\TT^2\) 上的作用
\[
(\varepsilon_1,\varepsilon_2)\cdot(z_1,z_2):=(z_1^{\varepsilon_1},z_2^{\varepsilon_2}).
\]
记
\[
g:=\gcd(|a|,|b|),
\qquad
T_{a,b}:=\frac{2\pi}{g}.
\]
则 \(L_{a,b,\delta}=\Pi\circ \widetilde L_{a,b,\delta}\)，且对任意
\[
t,t'\in[0,T_{a,b}),
\qquad
t\neq t',
\]
有严格等价
\[
L_{a,b,\delta}(t)=L_{a,b,\delta}(t')
\iff
\exists\,(\varepsilon_1,\varepsilon_2)\in V_4\setminus\{(1,1)\}
\ \text{使}\ 
\widetilde L_{a,b,\delta}(t)
=
(\varepsilon_1,\varepsilon_2)\cdot
\widetilde L_{a,b,\delta}(t').
\]
因此，Lissajous 的全部自交点都来源于提升轨道与其三条非平凡反演像之间的交会。
\end{proposition}
\begin{proof}
恒等式 \(L_{a,b,\delta}=\Pi\circ \widetilde L_{a,b,\delta}\) 由 \(\Re(e^{i\theta})=\cos\theta\) 立即得到。

现取 \(t,t'\in[0,T_{a,b})\)。若
\[
L_{a,b,\delta}(t)=L_{a,b,\delta}(t'),
\]
则
\[
\Re\bigl(e^{i(at+\delta)}\bigr)=\Re\bigl(e^{i(at'+\delta)}\bigr),
\qquad
\Re\bigl(e^{ibt}\bigr)=\Re\bigl(e^{ibt'}\bigr).
\]
在单位圆上，两个点实部相同当且仅当二者相等或互为逆元，因此存在
\[
\varepsilon_1,\varepsilon_2\in\{\pm1\}
\]
使
\[
e^{i(at+\delta)}=\bigl(e^{i(at'+\delta)}\bigr)^{\varepsilon_1},
\qquad
e^{ibt}=\bigl(e^{ibt'}\bigr)^{\varepsilon_2}.
\]
这正是
\[
\widetilde L_{a,b,\delta}(t)
=
(\varepsilon_1,\varepsilon_2)\cdot
\widetilde L_{a,b,\delta}(t').
\]

若 \((\varepsilon_1,\varepsilon_2)=(1,1)\)，则
\[
 a(t-t')\in 2\pi\ZZ,
\qquad
 b(t-t')\in 2\pi\ZZ.
\]
取 Bézout 系数 \(u,v\in\ZZ\) 使
\[
ua+vb=g.
\]
则
\[
g(t-t')=u\,a(t-t')+v\,b(t-t')\in 2\pi\ZZ,
\]
从而
\[
t-t'\in (2\pi/g)\ZZ=T_{a,b}\ZZ.
\]
由于 \(t,t'\in[0,T_{a,b})\)，只能有 \(t=t'\)，与假设矛盾。故此时 \((\varepsilon_1,\varepsilon_2)\neq (1,1)\)。

反向推出同样直接：若存在非平凡 \((\varepsilon_1,\varepsilon_2)\) 使
\[
\widetilde L_{a,b,\delta}(t)
=
(\varepsilon_1,\varepsilon_2)\cdot
\widetilde L_{a,b,\delta}(t'),
\]
则两坐标分别具有相同实部，因而
\[
\Pi\bigl(\widetilde L_{a,b,\delta}(t)\bigr)
=
\Pi\bigl(\widetilde L_{a,b,\delta}(t')\bigr),
\]
即 \(L_{a,b,\delta}(t)=L_{a,b,\delta}(t')\)。证毕。
\end{proof}

\begin{theorem}[Fibonacci 挠点并集饱和单位圆的全部有限挠点]\label{thm:xi-time-part62da-fibonacci-torsion-saturates-circle}
\leanverified{paper\_xi\_time\_part62da\_fibonacci\_torsion\_saturates\_circle}
定义 Fibonacci 挠点并集
\[
\mathrm{Tor}_F(\TT):=\bigcup_{m\ge 1}\mu_{F_{m+2}},
\]
其中 \(\mu_n\subset\TT\) 表示 \(n\) 次单位根群。则有严格恒等式
\[
\boxed{\;
\mathrm{Tor}_F(\TT)=\mathrm{Tor}(\TT):=\bigcup_{N\ge 1}\mu_N.\;}
\]
\end{theorem}
\begin{proof}
显然
\[
\mathrm{Tor}_F(\TT)\subseteq \mathrm{Tor}(\TT).
\]
反向包含方面，任取 \(\zeta\in\mu_N\)。由定理 \ref{thm:xi-visible-arithmetic-fibonacci-adic-profinite-coincidence} 的证明，对该 \(N\) 存在某个 \(m\ge 1\) 使
\[
N\mid F_{m+2}.
\]
于是
\[
\zeta^{F_{m+2}}=1,
\]
故 \(\zeta\in \mu_{F_{m+2}}\subseteq \mathrm{Tor}_F(\TT)\)。从而 \(\mathrm{Tor}(\TT)\subseteq \mathrm{Tor}_F(\TT)\)，两侧相等。证毕。
\end{proof}

\begin{theorem}[Fibonacci 分母子系统恢复 universal solenoid]\label{thm:xi-time-part62da-fibonacci-solenoid-cofinal-universal}
\leanverified{paper\_xi\_time\_part62da\_fibonacci\_solenoid\_cofinal\_universal}
记
\[
S_F:=\{F_{m+2}:\ m\ge 1\}\subset \NN.
\]
把 \(S_F\) 视为整除偏序 \((\NN,\mid)\) 的共终子系统，并定义对应逆极限
\[
\Sigma_F:=\varprojlim_{N\in S_F}\TT,
\]
其中当 \(N\mid N'\) 时连接映射取为 \(z\mapsto z^{N'/N}\)。则有规范同构
\[
\boxed{\;
\Sigma_F\cong \Sigma_{\mathbb P}=\widehat{\QQ}\cong \varprojlim_{N\ge 1}\TT.\;}
\]
\end{theorem}
\begin{proof}
由定理 \ref{thm:xi-visible-arithmetic-fibonacci-adic-profinite-coincidence} 的证明，对任意 \(N\ge 1\) 都存在某个 \(m\ge 1\) 使
\[
N\mid F_{m+2}.
\]
故 \(S_F\) 在整除偏序中对 \(\NN\) 共终。

另一方面，由命题 \ref{prop:xi-time-part62d-universal-solenoid-phase-completed-scan}，
\[
\Sigma_{\mathbb P}\cong \varprojlim_{N\ge 1}\TT
\]
且连接映射正是整除关系诱导的幂映射。逆极限对共终限制不变，故
\[
\varprojlim_{N\ge 1}\TT\cong \varprojlim_{N\in S_F}\TT=\Sigma_F.
\]
合并即得所示规范同构。证毕。
\end{proof}

\begin{theorem}[Fibonacci 分母子系统的单位群逆极限恢复 profinite 单位群]\label{thm:xi-time-part62da-fibonacci-unit-limit-zhat-units}
\leanverified{paper\_xi\_time\_part62da\_fibonacci\_unit\_limit\_zhat\_units}
定义
\[
U_F:=\varprojlim_{m}(\ZZ/F_{m+2}\ZZ)^\times,
\]
其中连接映射由整除关系下的标准降模同态给出。则存在规范拓扑群同构
\[
\boxed{\;
U_F\cong \widehat{\ZZ}^{\,\times}.\;}
\]
\end{theorem}
\begin{proof}
由定理 \ref{thm:xi-time-part62da-fibonacci-solenoid-cofinal-universal}，集合 \(S_F=\{F_{m+2}\}\) 在整除偏序中共终。故对单位群系统同样可作共终限制：
\[
\varprojlim_{N\ge 1}(\ZZ/N\ZZ)^\times
\cong
\varprojlim_{N\in S_F}(\ZZ/N\ZZ)^\times
=
U_F.
\]
而左端正是 \(\widehat{\ZZ}^{\,\times}\) 的经典表示。证毕。
\end{proof}

\begin{corollary}[Dirichlet 扭曲在 Fibonacci 分辨率上内禀实现]\label{cor:xi-time-part62da-dirichlet-twists-intrinsic-on-fibonacci-levels}
\leanverified{paper\_xi\_time\_part62da\_dirichlet\_twists\_intrinsic\_on\_fibonacci\_levels}
任取正整数 \(N\) 与 Dirichlet 特征
\[
\chi:(\ZZ/N\ZZ)^\times\to \CC^\times.
\]
则存在某个 \(m\ge 1\) 与群同态
\[
\widetilde\chi:(\ZZ/F_{m+2}\ZZ)^\times\to \CC^\times
\]
使得
\[
\widetilde\chi=\chi\circ \rho_{m,N},
\]
其中 \(\rho_{m,N}\) 为标准降模映射
\[
(\ZZ/F_{m+2}\ZZ)^\times\twoheadrightarrow (\ZZ/N\ZZ)^\times.
\]
因此，任意有限模 Dirichlet 扭曲都可在某一级 Fibonacci 分辨率上内禀实现；命题 \ref{prop:real-input-40-output-dirichlet-nonrh} 中的单位根扭曲只是其中最基本的特例。
\end{corollary}
\begin{proof}
由定理 \ref{thm:xi-visible-arithmetic-fibonacci-adic-profinite-coincidence} 的证明，可取 \(m\ge 1\) 使 \(N\mid F_{m+2}\)。于是标准降模映射
\[
\rho_{m,N}:(\ZZ/F_{m+2}\ZZ)^\times\twoheadrightarrow (\ZZ/N\ZZ)^\times
\]
为满射。令
\[
\widetilde\chi:=\chi\circ \rho_{m,N}
\]
便得到在 Fibonacci 模数层上的提升。再由定理 \ref{thm:fold-congruence-mod-semirings}，
\[
\Omega_m/\!\equiv_m\cong \ZZ/(F_{m+2}\ZZ),
\]
故该扭曲首先在可见同余层的单位扇区上给出规范乘法权重；再经由商映射
\[
\Omega_m\twoheadrightarrow \Omega_m/\!\equiv_m
\]
即可拉回到对应的微观态单位扇区。若需把它扩展为定义在全部微观态上的权重，只需再对非单位扇区作任意固定延拓。证毕。
\end{proof}

\begin{proposition}[有理玫瑰在原点的唯一聚合奇异性与切锥方向数]\label{prop:xi-time-part62da-rational-rose-origin-tangent-cone}
\leanverified{paper\_xi\_time\_part62da\_rational\_rose\_origin\_tangent\_cone}
设 \(n,d\in\NN\) 互素，且 \(d>n\)。定义复平面参数曲线
\[
z_{n,d}(t):=e^{idt}\cos(nt)
\qquad (t\in\RR).
\]
则：
\begin{enumerate}
  \item 参数曲线 \(z_{n,d}\) 处处浸入，但其像在原点 \(0\) 处存在自交聚合；并且原点是唯一由 \(\cos(nt)=0\) 所产生的聚合像点。
  \item 原点处不同切线方向的条数恰为
  \[
  \boxed{\;
  \#\mathrm{Tan}_0(z_{n,d})
  =
  n.\;}
  \]
  全部方向在 \([0,\pi)\) 内等角分布。
\end{enumerate}
\end{proposition}
\begin{proof}
首先，\(z_{n,d}(t)=0\) 当且仅当 \(\cos(nt)=0\)，即
\[
t=t_k:=\frac{(2k+1)\pi}{2n}
\qquad (k\in\ZZ).
\]
因此原点恰由这一离散时刻族聚合得到。另一方面，若 \(z_{n,d}(t)=0\)，则必满足 \(\cos(nt)=0\)，故在该参数表示下原点是唯一零像聚合点。

再求导得
\[
z'_{n,d}(t)
=
e^{idt}\bigl(-n\sin(nt)+id\cos(nt)\bigr).
\]
在 \(t=t_k\) 处有 \(\cos(nt_k)=0\) 且 \(\sin(nt_k)=\pm 1\)，从而
\[
z'_{n,d}(t_k)=\mp n\,e^{idt_k}\neq 0.
\]
故曲线处处浸入，原点处的奇异性来自多分支聚合，而非尖点。

原点处分支的切线方向由
\[
\arg z'_{n,d}(t_k)\equiv dt_k\pmod\pi
\]
决定。于是不同方向数等于集合
\[
\left\{\frac{d(2k+1)\pi}{2n}\bmod \pi:\ 0\le k<2n\right\}
\]
的基数。两项给出同一方向当且仅当
\[
d(2k+1)\equiv d(2k'+1)\pmod{2n},
\]
即
\[
n\mid d(k-k').
\]
由于 \(\gcd(n,d)=1\)，这等价于
\[
n\mid (k-k').
\]
因此在 \(0\le k<2n\) 的范围内，每个方向类恰由 \(k\) 与 \(k+n\) 这两个参数点组成，故不同切线方向总数为 \(n\)。

再看其分布。把方向写成
\[
\frac{d(2k+1)\pi}{2n}
\equiv
\frac{d\pi}{2n}+\frac{dk\pi}{n}
\pmod\pi.
\]
由于 \(\gcd(d,n)=1\)，乘以 \(d\) 在 \(\ZZ/n\ZZ\) 上只是一个置换，因此这些方向与
\[
\frac{\pi}{2n}+\frac{k\pi}{n}
\qquad (0\le k<n)
\]
给出的 \(n\) 个方向完全相同，从而在 \([0,\pi)\) 内等角分布。证毕。
\end{proof}

\leanverified{paper\_xi\_time\_part62da\_lissajous\_two\_nonimmersed\_points}
\begin{proposition}[Lissajous 的非沉浸点判别与精确二点数目]\label{prop:xi-time-part62da-lissajous-two-nonimmersed-points}
设 \(a,b\in\NN\) 互素，\(\delta\in\RR\)，并定义
\[
\gamma_{a,b,\delta}(t):=(\cos(at+\delta),\cos(bt))
\qquad (t\in\RR/2\pi\ZZ).
\]
则 \(\gamma_{a,b,\delta}\) 存在非沉浸点当且仅当
\[
\boxed{\;
\delta\in \frac{\pi}{b}\ZZ.\;}
\]
在此情形下，非沉浸点恰有 \(2\) 个。更精确地，若 \(k\) 是模 \(b\) 下满足
\[
ak+\frac{b\delta}{\pi}\equiv 0\pmod b
\]
的唯一剩余类，则对应参数值为
\[
\boxed{\;
t=\frac{k\pi}{b},
\qquad
t=\frac{(k+b)\pi}{b}\pmod{2\pi}.\;}
\]
\end{proposition}
\begin{proof}
由定理 \ref{thm:xi-lissajous-phase-singular-ring-immersion}，
\[
\gamma'_{a,b,\delta}(t)=0
\]
当且仅当
\[
\sin(bt)=0,
\qquad
\sin(at+\delta)=0.
\]
前者给出
\[
t=\frac{k\pi}{b}
\qquad(k\in\ZZ).
\]
代入第二式，得到同余条件
\[
ak+\frac{b\delta}{\pi}\equiv 0\pmod b.
\]
由于 \(\gcd(a,b)=1\)，\(a\) 在模 \(b\) 下可逆，因此该同余有解当且仅当
\[
\frac{b\delta}{\pi}\in \ZZ,
\]
即 \(\delta\in \frac{\pi}{b}\ZZ\)；并且解 \(k\) 在模 \(b\) 下唯一。

在 \([0,2\pi)\) 中，与这一模 \(b\) 解对应的参数恰为
\[
\frac{k\pi}{b}
\quad\text{与}\quad
\frac{(k+b)\pi}{b},
\]
二者相差 \(\pi\)，给出两个不同点。除此之外不再有别的 \(k\) 模 \(b\) 解，故非沉浸点恰有两处。证毕。
\end{proof}


在上述命题与命题 \ref{prop:xi-time-part62d-universal-solenoid-phase-completed-scan} 之间，实部通道给出经典玫瑰/Lissajous 的几何阴影，而与之互补的虚部通道则正是相位补全所必需的附加数据。由此可把 RH 一侧的圆盘完成化接口组织为如下猜想性口径。

\begin{conjecture}[虚部读出下的 prime-register 核消隐与 RH 的 Toeplitz 因子化]\label{conj:xi-time-part62d-imaginary-channel-prime-kernel-rh}
沿用命题 \ref{prop:xi-time-part62d-universal-solenoid-phase-completed-scan} 的相位补全坐标与定理 \ref{thm:xi-disk-zero-pole-index-reduction} 的圆盘完成化接口。设
\[
\Sigma_{\mathbb P}=\widehat{\QQ}
\twoheadrightarrow
\TT
\]
为 universal solenoid 的圆商。则 RH 应等价于如下核消隐原则：在与 \(\mathscr C_\zeta\) 对应的 Weyl--Carath\'eodory/Toeplitz 可见层上，全部边界数据都经由该圆商因子化；等价地，凡对 prime-register 核
\[
\prod_{p\in\mathbb P}\ZZ_p
\]
真正敏感的虚部通道分量都必须消失。若某一有限层仍残留此类核敏感分量，则相应 Toeplitz 截断不再维持 Carath\'eodory--PSD 正性，并应通过定理 \ref{thm:xi-disk-zero-pole-index-reduction} 诱发盘内零点，从而否定 RH。
\end{conjecture}

\paragraph{单位圆输入、奇异环障碍与结式阈值的审计闭环}
在有理玫瑰/Lissajous 的 prime-register 提升、单位圆触环与相位退化环已经建立之后，还可把二尺度二项式输运、Chebyshev 消元、Sylvester 结式与 RH 切片规范进一步压缩为下列一条闭合链。其核心是：单位圆只承载可见轨道，而真正控制分歧障碍与阈值分离的是位于复内层的一条临界圆支撑及其一元消元对象。

\begin{theorem}[二尺度二项式输运的单位圆层与临界圆支撑]\label{thm:xi-time-part62d-two-scale-binomial-transport-singular-circle}
设 \(m>n\ge 0\) 为整数，\(\alpha,\beta\in\CC^\times\)。定义
\[
F(z):=\alpha z^m+\beta z^n
\qquad (z\in\CC^\times).
\]
记其临界点集合为
\[
\mathcal S(F):=\{z\in\CC^\times:\ F'(z)=0\}.
\]
则：
\begin{enumerate}
  \item 若 \(n=0\)，则
  \[
  \mathcal S(F)=\varnothing.
  \]
  \item 若 \(1\le n<m\)，则
  \[
  \boxed{\;
  \mathcal S(F)=
  \left\{
  z\in\CC^\times:\ 
  z^{m-n}=-\frac{\beta n}{\alpha m}
  \right\}.\;}
  \]
  特别地，全部临界点都位于同一 Euclid 圆
  \[
  |z|=r_\ast,
  \qquad
  r_\ast:=
  \left|\frac{\beta n}{\alpha m}\right|^{1/(m-n)},
  \]
  并且恰由 \(m-n\) 个等角分布点组成。
\end{enumerate}
进一步，在 \(\CC^\times\setminus\mathcal S(F)\) 上有 \(F'(z)\neq 0\)，故 \(F\) 在该域上局部共形；因而局部逆分支的全部分歧都只可能发生在临界值集合 \(F(\mathcal S(F))\) 上。
\end{theorem}
\begin{proof}
直接求导得到
\[
F'(z)=\alpha m z^{m-1}+\beta n z^{n-1}
=
z^{n-1}\bigl(\alpha m z^{m-n}+\beta n\bigr).
\]
若 \(n=0\)，则
\[
F'(z)=\alpha m z^{m-1},
\]
而 \(z\in\CC^\times\) 上从不为零，故 \(\mathcal S(F)=\varnothing\)。

若 \(1\le n<m\)，则 \(z^{n-1}\neq 0\)，故
\[
F'(z)=0
\iff
\alpha m z^{m-n}+\beta n=0
\iff
z^{m-n}=-\frac{\beta n}{\alpha m}.
\]
于是 \(\mathcal S(F)\) 由 \(m-n\) 个 \((m-n)\)-次根构成，它们具有相同模长
\[
\left|-\frac{\beta n}{\alpha m}\right|^{1/(m-n)}
=
\left|\frac{\beta n}{\alpha m}\right|^{1/(m-n)},
\]
并按单位根等角分布。最后，\(F'(z)\neq 0\) 时由复分析中的逆函数定理知 \(F\) 在该点局部共形，局部逆分支的分歧只可能出现于临界值处。证毕。
\end{proof}

\begin{corollary}[有理玫瑰轨道的二项式输运实现与隐藏临界圆半径]\label{cor:xi-time-part62d-rose-binomial-transport-hidden-radius}
\leanverified{paper\_xi\_time\_part62d\_rose\_binomial\_transport\_hidden\_radius}
设 \(d>n\ge 1\) 互素，并定义
\[
F_{n,d}(z):=\frac12\bigl(z^{d+n}+z^{d-n}\bigr).
\]
则对一切 \(\theta\in\RR\)，若取
\[
z=e^{i\theta/d}\in\TT,
\]
则有严格恒等式
\[
\boxed{\;
F_{n,d}(e^{i\theta/d})
=
\cos\!\Bigl(\frac{n}{d}\theta\Bigr)e^{i\theta}.\;}
\]
因此定理 \ref{thm:xi-time-part62d-solenoid-kronecker-rose-lissajous-lift} 与命题 \ref{prop:xi-rose-unit-circle-touch-ring-parity-bifurcation} 中的有理玫瑰轨道正是 \(F_{n,d}(\TT)\) 的像。

其临界圆半径为
\[
\boxed{\;
r_\ast(n,d)=
\left(\frac{d-n}{d+n}\right)^{1/(2n)}.\;}
\]
并且满足完全闭式
\[
\boxed{\;
\log\frac{1}{r_\ast(n,d)}
=
\frac{1}{n}\operatorname{arctanh}\!\Bigl(\frac{n}{d}\Bigr).\;}
\]
\end{corollary}
\begin{proof}
由 Euler 公式，
\[
\cos\!\Bigl(\frac{n}{d}\theta\Bigr)e^{i\theta}
=
\frac12
\left(
e^{i\frac{n}{d}\theta}+e^{-i\frac{n}{d}\theta}
\right)e^{i\theta}
=
\frac12
\left(
e^{i\frac{d+n}{d}\theta}+e^{i\frac{d-n}{d}\theta}
\right),
\]
而 \(z=e^{i\theta/d}\)，故
\[
\frac12
\left(
e^{i\frac{d+n}{d}\theta}+e^{i\frac{d-n}{d}\theta}
\right)
=
\frac12\bigl(z^{d+n}+z^{d-n}\bigr)
=
F_{n,d}(z).
\]

再把定理 \ref{thm:xi-time-part62d-two-scale-binomial-transport-singular-circle} 应用于
\[
\alpha=\beta=\frac12,
\qquad
m=d+n,
\qquad
n_0=d-n,
\]
则
\[
m-n_0=2n,
\qquad
\frac{\beta n_0}{\alpha m}=\frac{d-n}{d+n}.
\]
因此
\[
r_\ast(n,d)
=
\left(\frac{d-n}{d+n}\right)^{1/(2n)}.
\]
最后，
\[
\operatorname{arctanh}(x)=\frac12\log\frac{1+x}{1-x}
\qquad(|x|<1),
\]
代入 \(x=n/d\in(0,1)\) 后得到
\[
\frac{1}{n}\operatorname{arctanh}\!\Bigl(\frac{n}{d}\Bigr)
=
\frac{1}{2n}\log\frac{d+n}{d-n}
=
\log\frac{1}{r_\ast(n,d)}.
\]
证毕。
\end{proof}

\begin{theorem}[有理玫瑰族的单位圆反推与奇异环消元公式]\label{thm:xi-time-part62d-rational-rose-singular-ring-elimination}
\leanverified{paper\_xi\_time\_part62d\_rational\_rose\_singular\_ring\_elimination}
设 \(n,d\in\NN\) 互素，并考虑复平面参数曲线
\[
\gamma_{n,d}:\qquad
z(\theta):=\cos\!\Bigl(\frac{n}{d}\theta\Bigr)e^{i\theta}.
\]
在复化平面上引入坐标 \((z,w)\in\CC^2\)，其中实切片对应 \(w=\bar z\)。记第一类 Chebyshev 多项式为 \(T_d\)。则 \(\gamma_{n,d}(\RR)\) 的 Zariski 闭包满足单一消元关系
\[
\boxed{\;
z^{2n}+w^{2n}-2z^n w^n\,T_d(2zw-1)=0.\;}
\]
因此在实平面切片 \(z=x+iy\)、\(w=x-iy\) 上有等价方程
\[
\boxed{\;
\Re\!\bigl((x+iy)^{2n}\bigr)
=(x^2+y^2)^n\,T_d\!\bigl(2(x^2+y^2)-1\bigr).\;}
\]
特别地，单位圆参数所反推出的玫瑰奇异环在 \(\RR[x,y]\) 中被压缩为单一显式关系。

进一步，若记
\[
F_{n,d}(z,w):=z^{2n}+w^{2n}-2z^n w^n\,T_d(2zw-1),
\]
则原点的切锥由其最低次齐次项
\[
F_{n,d}^{[2n]}(z,w)
=
z^{2n}+w^{2n}-2(-1)^d z^n w^n
=
\bigl(z^n-(-1)^d w^n\bigr)^2
\]
给出。于是原点的重数恰为 \(2n\)，且切锥分裂为
\[
z=\zeta\,w,
\qquad
\zeta^n=(-1)^d,
\]
这 \(n\) 条互异方向，各自以重数 \(2\) 出现。等价地，在实切片上其齐次主项可写成
\[
\Re(z^{2n})-(-1)^d(z\bar z)^n.
\]
\end{theorem}
\begin{proof}
取
\[
u:=e^{i\theta/d}.
\]
则有
\[
z=\frac12\bigl(u^{d+n}+u^{d-n}\bigr),
\qquad
w=\frac12\bigl(u^{-(d+n)}+u^{-(d-n)}\bigr).
\]
于是立刻得到
\[
\frac{z}{w}=u^{2d},
\qquad
4zw=(u^n+u^{-n})^2.
\]
若记
\[
U:=u^{2n},
\]
则
\[
U+U^{-1}=4zw-2.
\]
由 Chebyshev 恒等式
\[
U^d+U^{-d}
=
2T_d\!\Bigl(\frac{U+U^{-1}}{2}\Bigr)
\]
可知
\[
\Bigl(\frac{z}{w}\Bigr)^n+\Bigl(\frac{w}{z}\Bigr)^n
=
U^d+U^{-d}
=
2T_d(2zw-1).
\]
两边同乘 \(z^n w^n\)，便得到
\[
z^{2n}+w^{2n}-2z^n w^n\,T_d(2zw-1)=0.
\]
这就给出所示消元关系。

在实切片上代入
\[
w=\bar z=x-iy,
\qquad
zw=x^2+y^2,
\qquad
z^{2n}+w^{2n}=2\Re(z^{2n}),
\]
即得
\[
\Re\!\bigl((x+iy)^{2n}\bigr)
=(x^2+y^2)^n\,T_d\!\bigl(2(x^2+y^2)-1\bigr).
\]

最后，由
\[
T_d(-1)=(-1)^d
\]
且 \(2zw\) 具有次数 \(2\)，可知 \(F_{n,d}\) 在原点的最低次齐次项正是
\[
z^{2n}+w^{2n}-2(-1)^d z^n w^n
=
\bigl(z^n-(-1)^d w^n\bigr)^2.
\]
故原点重数为 \(2n\)，而切锥由满足 \(\zeta^n=(-1)^d\) 的 \(n\) 条互异方向 \(z=\zeta w\) 组成，并且每条方向都以重数 \(2\) 出现。证毕。
\end{proof}

\begin{theorem}[Lissajous 图形的单位圆干涉闭包与相位参数的代数化]\label{thm:xi-time-part62d-lissajous-chebyshev-implicit-equation}
\leanverified{paper\_xi\_time\_part62d\_lissajous\_chebyshev\_implicit\_equation}
设 \(a,b\in\NN\)、\(\delta\in\RR\)，并定义相移 Lissajous 轨道
\[
x(t)=\cos(at+\delta),
\qquad
y(t)=\cos(bt).
\]
记第一类 Chebyshev 多项式为 \(T_k\)。则该轨道满足单一代数约束
\[
\boxed{\;
T_b(x)^2
-2\cos(b\delta)\,T_b(x)\,T_a(y)
+T_a(y)^2
=
\sin^2(b\delta).\;}
\]
特别地，当
\[
\cos(b\delta)=0
\]
时，上式退化为
\[
\boxed{\;
T_b(x)^2+T_a(y)^2=1.\;}
\]
因此，相位参数 \(\delta\) 在坐标环中只通过 \(\cos(b\delta)\)（等价地 \(\sin^2(b\delta)\)）进入。
\end{theorem}
\begin{proof}
令
\[
\alpha:=abt,
\qquad
c:=\cos(b\delta),
\qquad
s:=\sin(b\delta).
\]
则
\[
T_a(y)=\cos(abt)=\cos\alpha.
\]
另一方面，由 \(x(t)=\cos(at+\delta)\)，有
\[
T_b(x)=T_b(\cos(at+\delta))
=
\cos(\alpha+b\delta)
=
c\cos\alpha-s\sin\alpha
=
c\,T_a(y)-s\sin\alpha.
\]
因此
\[
\bigl(T_b(x)-c\,T_a(y)\bigr)^2=s^2\sin^2\alpha.
\]
再由
\[
\sin^2\alpha=1-\cos^2\alpha=1-T_a(y)^2,
\]
得到
\[
\bigl(T_b(x)-c\,T_a(y)\bigr)^2
=
s^2\bigl(1-T_a(y)^2\bigr).
\]
展开并整理即得
\[
T_b(x)^2-2c\,T_b(x)\,T_a(y)+T_a(y)^2=s^2,
\]
也即
\[
T_b(x)^2
-2\cos(b\delta)\,T_b(x)\,T_a(y)
+T_a(y)^2
=
\sin^2(b\delta).
\]
当 \(\cos(b\delta)=0\) 时，\(s^2=1\)，故立得
\[
T_b(x)^2+T_a(y)^2=1.
\]
证毕。
\end{proof}

\begin{proposition}[相移 Lissajous 复延拓的临界值由 Sylvester 结式完全刻画]\label{prop:xi-time-part62d-lissajous-branch-spectrum-resultant}
\leanverified{paper\_xi\_time\_part62d\_lissajous\_branch\_spectrum\_resultant}
设 \(a,b\in\NN\)、\(\delta\in\RR\)，并定义
\[
P_{a,b,\delta}(z):=
\frac12\Bigl(
e^{i\delta}z^a+e^{-i\delta}z^{-a}
+iz^b+iz^{-b}
\Bigr)
\qquad(z\in\CC^\times).
\]
记其临界值集合为
\[
\Sigma_{a,b,\delta}:=
\Bigl\{
w\in\CC:\ \exists z\in\CC^\times,\
P_{a,b,\delta}(z)=w,\
zP_{a,b,\delta}'(z)=0
\Bigr\}.
\]
再令 \(M:=\max\{a,b\}\)，并定义关于 \(z\) 的两个多项式
\[
\widetilde P_{a,b,\delta}(z,w):=
2z^M\bigl(P_{a,b,\delta}(z)-w\bigr),
\]
\[
\widetilde Q_{a,b,\delta}(z):=
2z^M\bigl(zP_{a,b,\delta}'(z)\bigr).
\]
则一元多项式
\[
\boxed{\;
R_{a,b,\delta}(w):=
\operatorname{Res}_{z}
\bigl(
\widetilde P_{a,b,\delta}(z,w),\
\widetilde Q_{a,b,\delta}(z)
\bigr)
\in\CC[w]\;}
\]
满足
\[
\boxed{\;
\Sigma_{a,b,\delta}
=
\{w\in\CC:\ R_{a,b,\delta}(w)=0\}.\;}
\]
因此相移 Lissajous 的全部分歧值构成一个有限、可由 Sylvester 结式直接审计的一元代数对象。
\end{proposition}
\begin{proof}
由定义，
\[
P_{a,b,\delta}(z)=w,\qquad zP_{a,b,\delta}'(z)=0
\]
在 \(z\in\CC^\times\) 上同时成立，当且仅当
\[
\widetilde P_{a,b,\delta}(z,w)=0,
\qquad
\widetilde Q_{a,b,\delta}(z)=0
\]
同时成立。这里乘以 \(2z^M\) 只是在 \(\CC^\times\) 上清除 Laurent 分母，并不改变共同零点。

因此对固定 \(w\)，存在某个 \(z\in\CC^\times\) 使二式同时成立，当且仅当这两个关于 \(z\) 的多项式有公共根；按结式的定义，这又等价于
\[
R_{a,b,\delta}(w)=0.
\]
由此得到
\[
\Sigma_{a,b,\delta}
=
\{w\in\CC:\ R_{a,b,\delta}(w)=0\}.
\]
证毕。
\end{proof}

\begin{remark}[复嵌入下的轴交换不是结构对称]\label{rem:xi-time-part62d-axis-exchange-not-structural-symmetry}
命题 \ref{prop:xi-time-part62d-lissajous-branch-spectrum-resultant} 给出的 \(P_{a,b,\delta}\) 一般并不满足自互反或共轭自反的 Laurent 系数约束，因此其临界值集合 \(\Sigma_{a,b,\delta}\) 通常并不享有普适的实轴共轭封闭性。换言之，在复读出
\[
w=x+iy
\]
中交换实轴与虚轴并不是一个内部对称操作，而只是切片规范的改变；这正是下一个判别原理的几何背景。
\end{remark}

\begin{theorem}[实轴与虚轴的交换只对应切片规范而不产生新的 RH 信息]\label{thm:xi-time-part62d-axis-rotation-is-slice-gauge}
\leanverified{paper\_xi\_time\_part62d\_axis\_rotation\_is\_slice\_gauge}
设 \(I\subset\RR\) 为连通区间，\(f:I\to\CC\) 连续且在 \(I\) 上无零点。则存在连续相位函数
\[
\vartheta:I\to\RR
\]
使得
\[
e^{-i\vartheta(t)}f(t)\in\RR_{>0}
\qquad(\forall t\in I).
\]
进一步，若存在常数相位 \(\phi\in\RR\) 使
\[
e^{-i\phi}f(t)\in\RR
\qquad(\forall t\in I),
\]
则 \(\arg f(t)\) 在 \(I\) 上必为常数（模 \(\pi\)）。

因此，任何固定轴交换
\[
f\longmapsto e^{-i\phi}f
\]
都不能把本质二维的消失条件 \(f(t)=0\) 简化为新的单轴抵消条件；除非原函数在该区间内本就具有常相位。就 RH 的 Carath\'eodory--Toeplitz 切片语言而言，这说明换轴只改变读出规范，而不引入新的结构信息。
\end{theorem}
\begin{proof}
由于 \(f\) 在 \(I\) 上无零点，\(\arg f\) 可在连通区间 \(I\) 上取连续支。令
\[
\vartheta(t):=\arg f(t),
\]
则
\[
e^{-i\vartheta(t)}f(t)=|f(t)|>0,
\]
从而第一条成立。

若存在常数 \(\phi\) 使 \(e^{-i\phi}f(t)\in\RR\) 对全部 \(t\in I\) 成立，则
\[
\arg f(t)\in \phi+\pi\ZZ
\qquad(\forall t\in I).
\]
而 \(\arg f(t)\) 是连通集上的连续函数，只能落在离散集 \(\phi+\pi\ZZ\) 的单一点上，故 \(\arg f(t)\) 在 \(I\) 上必为常数（模 \(\pi\)）。这就说明固定相位旋转只能在常相位情形下把复值读出压扁到单实轴上。证毕。
\end{proof}

\begin{conjecture}[单位圆可见层的临界圆边界分离与 RH 审计闭环]\label{conj:xi-time-part62d-singular-circle-boundary-separation-rh}
沿用定理 \ref{thm:xi-time-part62d-two-scale-binomial-transport-singular-circle}、命题 \ref{prop:xi-time-part62d-lissajous-branch-spectrum-resultant}、猜想 \ref{conj:xi-time-part62d-imaginary-channel-prime-kernel-rh}、定理 \ref{thm:xi-disk-zero-pole-index-reduction} 与定义 \ref{def:kernel-newman-threshold} 的语言。预期对每一层级 \(m\)，都存在一个回文自对偶的完成化输运核 \(P_m\)，满足：
\begin{enumerate}
  \item \(P_m|_{\TT}\) 给出该层的可见边界数据；
  \item \(zP_m'(z)=0\) 的解全部支撑在若干条临界圆上；
  \item 相应分歧值集合 \(\Sigma_m\) 由一个有限 Sylvester 结式 \(R_m\) 完全决定；
  \item RH 等价于 \(\Sigma_m\) 在完成化坐标中对边界 \(\TT\) 的严格分离；
  \item 首次失去分离的参数恰给出该层的 Newman 型阈值，并且这些阈值沿层级推进单调收敛。
\end{enumerate}
换言之，RH 应可被重写为一个“单位圆可见层与临界圆障碍永不接触”的审计原则；其失败见证则由某一层的结式根直接输出，从而形成有限消元、阈值触边与离线验证器之间的闭环。
\end{conjecture}

\paragraph{无限秩可消去幺半群、局部化直和账本与态射矩阵的统一分类}
在可消去幺半群的 Grothendieck 完成刚性、有限素数局部化数系的圆维--prime-ledger 分裂、以及单对象 \(\mathrm{Hom}\)-分类与自覆盖分类已经确立之后，还可把无限秩发散、有限直和的二轴账本与态射零图统一压缩为下列六条结构结论。

\begin{theorem}[无限 Grothendieck 秩可消去幺半群的结构层半圆维发散与实现层极小性]\label{thm:xi-infinite-grothendieck-rank-halfdimension-divergence}
\leanverified{paper\_xi\_infinite\_grothendieck\_rank\_halfdimension\_divergence}
设 \(M\) 为交换可消去幺半群，并满足定义 \ref{def:half-circle-dimension} 的嵌入/完备性口径。若 \(M\) 的底层集合可数，且其 Grothendieck 群完成满足
\[
\operatorname{rk}_{\ZZ}G(M)=\infty,
\]
则
\[
\cdim^{\mathrm{hom}}_{+}(M)=+\infty,
\qquad
\cdim^{\mathrm{impl}}(M)=\frac12.
\]
\end{theorem}
\begin{proof}
反设 \(\cdim^{\mathrm{hom}}_{+}(M)<\infty\)。则按定义存在某个有限 \(k\) 与单射幺半群同态
\[
M\hookrightarrow (\NN^k,+).
\]
由命题 \ref{prop:killo-grothendieck-completion-preserves-injection}，Grothendieck 完成后得到群单射
\[
G(M)\hookrightarrow G(\NN^k,+)\cong \ZZ^k.
\]
然而 \(\ZZ^k\) 的任意子群都至多具有有限秩 \(k\)，这与 \(\operatorname{rk}_{\ZZ}G(M)=\infty\) 矛盾。故
\[
\cdim^{\mathrm{hom}}_{+}(M)=+\infty.
\]

另一方面，由 \(M\) 可数可得存在集合单射
\[
M\hookrightarrow \NN,
\]
故
\[
\cdim^{\mathrm{impl}}(M)\le \frac12.
\]
又因 \(\operatorname{rk}_{\ZZ}G(M)=\infty\)，群 \(G(M)\) 非平凡，从而 \(M\) 不可能是单点集，故不可能单射进入 \(\NN^0\)。因此 \(\cdim^{\mathrm{impl}}(M)\neq 0\)，只可能取最小非零值
\[
\cdim^{\mathrm{impl}}(M)=\frac12.
\]
证毕。
\end{proof}

\begin{theorem}[有限局部化直和的二轴账本分解与圆维公式]\label{thm:xi-localized-directsum-twoaxis-ledger}
\leanverified{paper\_xi\_localized\_directsum\_twoaxis\_ledger}
设
\[
\mathbf S=(S_1,\dots,S_r),
\qquad
S_i\subset\mathbb P\ \text{有限},
\]
并定义
\[
G_{\mathbf S}:=\bigoplus_{i=1}^{r}G_{S_i},
\qquad
m_p(\mathbf S):=\#\{\,i:\ p\in S_i\,\}.
\]
则存在规范短正合列
\[
0\longrightarrow \ZZ^r\longrightarrow G_{\mathbf S}
\longrightarrow \bigoplus_{p\in\mathbb P} C_{p^\infty}^{\,m_p(\mathbf S)}
\longrightarrow 0,
\]
其 Pontryagin 对偶为
\[
0\longrightarrow \prod_{p\in\mathbb P}\ZZ_p^{\,m_p(\mathbf S)}
\longrightarrow \widehat{G_{\mathbf S}}
\longrightarrow \TT^r
\longrightarrow 0.
\]
此外，
\[
\cdim_{\mathrm{fr}}(G_{\mathbf S})=r.
\]
若记上述对偶短正合列的核为
\[
K_{\mathbf S}:=\ker\!\bigl(\widehat{G_{\mathbf S}}\twoheadrightarrow \TT^r\bigr),
\]
则
\[
K_{\mathbf S}\cong \prod_{p\in\mathbb P}\ZZ_p^{\,m_p(\mathbf S)}
\cong
\widehat{G_{\mathbf S}}/(\widehat{G_{\mathbf S}})^0.
\]
等价地，对每个素数 \(p\) 都有
\[
(K_{\mathbf S})_p\cong \ZZ_p^{\,m_p(\mathbf S)},
\]
故多重度函数 \(p\mapsto m_p(\mathbf S)\) 可由 \(\widehat{G_{\mathbf S}}\) 的全不连通商完全恢复。
\end{theorem}
\begin{proof}
对每个 \(i\)，定理 \ref{thm:killo-localization-circle-dimension-prime-ledger-splitting} 给出短正合列
\[
0\longrightarrow \ZZ\longrightarrow G_{S_i}
\longrightarrow \bigoplus_{p\in S_i}C_{p^\infty}
\longrightarrow 0.
\]
对 \(i=1,\dots,r\) 取有限直和，得到
\[
0\longrightarrow \ZZ^r\longrightarrow \bigoplus_{i=1}^{r}G_{S_i}
\longrightarrow \bigoplus_{i=1}^{r}\bigoplus_{p\in S_i}C_{p^\infty}
\longrightarrow 0.
\]
把右端按素数重新分组，便得到
\[
\bigoplus_{i=1}^{r}\bigoplus_{p\in S_i}C_{p^\infty}
\cong
\bigoplus_{p\in\mathbb P} C_{p^\infty}^{\,m_p(\mathbf S)},
\]
其中仅有限多个素数给出非平凡直和项。这就得到第一条短正合列。

对其作 Pontryagin 对偶，并使用
\[
\widehat{C_{p^\infty}}\cong \ZZ_p,
\qquad
\widehat{\ZZ}\cong \TT,
\]
即得
\[
0\longrightarrow \prod_{p\in\mathbb P}\ZZ_p^{\,m_p(\mathbf S)}
\longrightarrow \widehat{G_{\mathbf S}}
\longrightarrow \TT^r
\longrightarrow 0.
\]

再看圆维。每个 \(G_{S_i}\subset \QQ\) 都是秩 \(1\) 的无挠群，因此
\[
G_{\mathbf S}=\bigoplus_{i=1}^{r}G_{S_i}
\]
亦为有限秩无挠群，且
\[
\operatorname{rank}_{\QQ}(G_{\mathbf S})=r.
\]
由定理 \ref{thm:conclusion-cdim-finite-rank-extension}，
\[
\cdim_{\mathrm{fr}}(G_{\mathbf S})=r
\qquad\text{且}\qquad
(\widehat{G_{\mathbf S}})^0\cong \TT^r.
\]

最后比较全不连通核与连通分量。设
\[
\pi:\widehat{G_{\mathbf S}}\twoheadrightarrow \TT^r
\]
为上式中的商映射。因为 \(\prod_p \ZZ_p^{\,m_p(\mathbf S)}\) 为紧全不连通群，故其与 \((\widehat{G_{\mathbf S}})^0\) 的交只能平凡。另一方面，连续满同态把连通分量映到连通分量，因此
\[
\pi\bigl((\widehat{G_{\mathbf S}})^0\bigr)=\TT^r.
\]
于是 \(\pi|_{(\widehat{G_{\mathbf S}})^0}\) 为同构，从而
\[
\widehat{G_{\mathbf S}}/(\widehat{G_{\mathbf S}})^0
\cong
\ker(\pi)=K_{\mathbf S}.
\]
结合
\[
K_{\mathbf S}\cong \prod_{p\in\mathbb P}\ZZ_p^{\,m_p(\mathbf S)}
\]
即可得到每个 \(p\)-主因子
\[
(K_{\mathbf S})_p\cong \ZZ_p^{\,m_p(\mathbf S)}.
\]
故全部多重度 \(m_p(\mathbf S)\) 都由全不连通商唯一恢复。证毕。
\end{proof}

\leanverified{paper\_xi\_localized\_directsum\_twoaxis\_ledger\_cancellation}
\begin{corollary}[二轴账本在局部化直和族上的严格消去律]\label{cor:xi-localized-directsum-twoaxis-ledger-cancellation}
设 \(\mathbf S,\mathbf T,\mathbf U\) 为有限素数集的有限族。若
\[
G_{\mathbf S}\oplus G_{\mathbf U}\cong G_{\mathbf T}\oplus G_{\mathbf U}
\]
作为离散交换群同构，则
\[
\#\mathbf S=\#\mathbf T,
\qquad
m_p(\mathbf S)=m_p(\mathbf T)\quad(\forall p\in\mathbb P).
\]
\end{corollary}
\begin{proof}
由定理 \ref{thm:xi-localized-directsum-twoaxis-ledger}，
\[
\cdim_{\mathrm{fr}}(G_{\mathbf S}\oplus G_{\mathbf U})
=\#\mathbf S+\#\mathbf U,
\]
同理
\[
\cdim_{\mathrm{fr}}(G_{\mathbf T}\oplus G_{\mathbf U})
=\#\mathbf T+\#\mathbf U.
\]
同构保持 \(\cdim_{\mathrm{fr}}\)，故
\[
\#\mathbf S+\#\mathbf U=\#\mathbf T+\#\mathbf U,
\]
从而
\[
\#\mathbf S=\#\mathbf T.
\]

再对给定同构作 Pontryagin 对偶，并分别取全不连通商。由定理 \ref{thm:xi-localized-directsum-twoaxis-ledger}，
\[
\widehat{G_{\mathbf S}\oplus G_{\mathbf U}}/
\bigl(\widehat{G_{\mathbf S}\oplus G_{\mathbf U}}\bigr)^0
\cong
\prod_{p\in\mathbb P}\ZZ_p^{\,m_p(\mathbf S)+m_p(\mathbf U)},
\]
\[
\widehat{G_{\mathbf T}\oplus G_{\mathbf U}}/
\bigl(\widehat{G_{\mathbf T}\oplus G_{\mathbf U}}\bigr)^0
\cong
\prod_{p\in\mathbb P}\ZZ_p^{\,m_p(\mathbf T)+m_p(\mathbf U)}.
\]
同构保持各个 \(p\)-主因子的 \(\ZZ_p\)-幂次，因此对每个素数 \(p\) 都有
\[
m_p(\mathbf S)+m_p(\mathbf U)
=
m_p(\mathbf T)+m_p(\mathbf U).
\]
消去共同项即得
\[
m_p(\mathbf S)=m_p(\mathbf T).
\]
证毕。
\end{proof}

\begin{theorem}[有限素数截断的局部化账本实现需要且仅需要 \texorpdfstring{$B$}{B} 条通道]\label{thm:xi-localized-directsum-prime-truncation-minimal-channels}
\leanverified{paper\_xi\_localized\_directsum\_prime\_truncation\_minimal\_channels}
记 \(p_1<p_2<\cdots\) 为递增素数序列。对 \(B\ge 1\)，定义有限素数截断幺半群
\[
M_B:=\langle p_1,\dots,p_B\rangle_{\times}\subset \NN^\times.
\]
则对任意有限素数集 \(S_1,\dots,S_r\)，成立精确等价
\[
\boxed{\;
\exists\ \text{幺半群单射 }\ 
M_B\hookrightarrow G_{S_1}\oplus\cdots\oplus G_{S_r}
\iff
r\ge B.\;}
\]
特别地，使前 \(B\) 个 prime-register 在有限局部化直和账本中得到严格同态实现所需的最小通道数恰为 \(B\)。
\end{theorem}
\begin{proof}
指数向量映射
\[
p_1^{a_1}\cdots p_B^{a_B}\longmapsto (a_1,\dots,a_B)
\]
给出幺半群同构
\[
M_B\cong (\NN^B,+).
\]
因此只需讨论 \((\NN^B,+)\) 是否可单射进入给定的加法幺半群
\(
G_{S_1}\oplus\cdots\oplus G_{S_r}
\)。

先证充分性。若 \(r\ge B\)，则每个 \(G_{S_i}\) 都含有整数子群 \(\ZZ\subset G_{S_i}\)。故当 \(r=B\) 时可定义
\[
\Phi:\NN^B\longrightarrow G_{S_1}\oplus\cdots\oplus G_{S_B},
\qquad
(a_1,\dots,a_B)\longmapsto (a_1,\dots,a_B),
\]
这是显然的幺半群单射。若 \(r>B\)，再把右端嵌入到前 \(B\) 个坐标并在其余坐标填零即可。

再证必要性。设存在幺半群单射
\[
\Psi:\NN^B\hookrightarrow G_{S_1}\oplus\cdots\oplus G_{S_r}.
\]
记 \(e_1,\dots,e_B\) 为 \(\NN^B\) 的标准生成元，并令
\[
v_i:=\Psi(e_i)\in G_{S_1}\oplus\cdots\oplus G_{S_r}\subset \QQ^r.
\]
若 \(r<B\)，则 \(v_1,\dots,v_B\) 在 \(\QQ^r\) 中必线性相关，故存在不全为零的整数 \(z_1,\dots,z_B\) 使
\[
\sum_{i=1}^B z_i v_i=0.
\]
把 \(z_i\) 写成正负部分之差
\[
z_i=m_i-n_i,
\qquad
m_i,n_i\in \NN,
\]
则
\[
\sum_{i=1}^B m_i v_i=\sum_{i=1}^B n_i v_i.
\]
于是由幺半群同态性得到
\[
\Psi(m_1,\dots,m_B)=\Psi(n_1,\dots,n_B).
\]
又因 \((z_i)\) 不全为零，故 \((m_i)\neq (n_i)\)，这与 \(\Psi\) 的单射性矛盾。故必有 \(r\ge B\)。证毕。
\end{proof}

\begin{theorem}[局部化直和对象族的 \texorpdfstring{$\mathrm{Hom}$}{Hom} 群稀疏矩阵分类]\label{thm:xi-localized-directsum-hom-sparse-matrix}
\leanverified{paper\_xi\_localized\_directsum\_hom\_sparse\_matrix}
设
\[
\mathbf S=(S_1,\dots,S_r),
\qquad
\mathbf T=(T_1,\dots,T_s),
\]
并记
\[
G_{\mathbf S}:=\bigoplus_{i=1}^{r}G_{S_i},
\qquad
G_{\mathbf T}:=\bigoplus_{j=1}^{s}G_{T_j}.
\]
则存在自然同构
\[
\operatorname{Hom}(G_{\mathbf S},G_{\mathbf T})
\cong
\bigoplus_{j=1}^{s}\ \bigoplus_{\substack{1\le i\le r\\ S_i\subseteq T_j}} G_{T_j}.
\]
等价地，每个群同态都唯一对应于一个 \(s\times r\) 稀疏块矩阵 \((x_{ji})\)，其中
\[
x_{ji}\in
\begin{cases}
G_{T_j},& S_i\subseteq T_j,\\
0,& S_i\not\subseteq T_j,
\end{cases}
\]
而该矩阵对元素 \((y_i)_{i=1}^{r}\in G_{\mathbf S}\) 的作用为
\[
(y_i)_{i=1}^{r}
\longmapsto
\left(
\sum_{\substack{1\le i\le r\\ S_i\subseteq T_j}}x_{ji}y_i
\right)_{j=1}^{s}.
\]
\end{theorem}
\begin{proof}
有限直和的泛性质给出
\[
\operatorname{Hom}\!\Bigl(\bigoplus_{i=1}^{r}G_{S_i},\bigoplus_{j=1}^{s}G_{T_j}\Bigr)
\cong
\bigoplus_{j=1}^{s}\bigoplus_{i=1}^{r}\operatorname{Hom}(G_{S_i},G_{T_j}).
\]
对每一对 \((i,j)\)，定理 \ref{thm:xi-cdim-localization-hom-category-classification} 给出
\[
\operatorname{Hom}(G_{S_i},G_{T_j})
\cong
\begin{cases}
G_{T_j},& S_i\subseteq T_j,\\[1mm]
0,& S_i\not\subseteq T_j.
\end{cases}
\]
逐项代入即得所示直和分解。

在 \(S_i\subseteq T_j\) 时，该同构把 \(x_{ji}\in G_{T_j}\) 送到乘法映射
\[
m_{x_{ji}}:G_{S_i}\to G_{T_j},
\qquad
y\mapsto x_{ji}y.
\]
把全部块拼接起来，便得到所述稀疏矩阵表示。唯一性来自上述直和分解与单块分类的唯一性。证毕。
\end{proof}

\begin{corollary}[局部化直和对象族的端同态环由素数支撑偏序支配]\label{cor:xi-localized-directsum-end-ring-poset}
\leanverified{paper\_xi\_localized\_directsum\_end\_ring\_poset}
沿用定理 \ref{thm:xi-localized-directsum-hom-sparse-matrix} 的记号，并取 \(\mathbf T=\mathbf S\)。则
\[
\operatorname{End}(G_{\mathbf S})
\cong
\bigoplus_{j=1}^{r}\ \bigoplus_{\substack{1\le i\le r\\ S_i\subseteq S_j}} G_{S_j}
\]
作为加法群同构。

若先把相同的素数支撑合并为同一对角块，再按严格包含关系的某个线性扩张重排这些块，则 \(\operatorname{End}(G_{\mathbf S})\) 可视为一个稀疏块上三角环：仅当对应支撑满足包含关系时，块位 \((j,i)\) 才允许非零；若某个支撑类 \(U\) 重复出现 \(n\) 次，则相应对角块自然识别为 \(M_n(G_U)\)。

特别地，若各 \(S_i\) 两两不可比，则
\[
\operatorname{End}(G_{\mathbf S})\cong \prod_{i=1}^{r}G_{S_i};
\]
若
\[
S_1\subsetneq S_2\subsetneq \cdots \subsetneq S_r,
\]
则每个上三角块都非零。
\end{corollary}
\begin{proof}
把定理 \ref{thm:xi-localized-directsum-hom-sparse-matrix} 应用于 \(\mathbf T=\mathbf S\) 即得加法群分解。环乘法由矩阵复合给出，而包含关系的传递性
\[
S_i\subseteq S_j,\quad S_j\subseteq S_k
\quad\Longrightarrow\quad
S_i\subseteq S_k
\]
保证零块模式在复合下封闭，因此在兼容包含偏序的重排下得到块上三角结构。

若各 \(S_i\) 两两不可比，则对角外不可能有 \(S_i\subseteq S_j\)（\(i\neq j\)），故仅剩对角块，从而
\[
\operatorname{End}(G_{\mathbf S})\cong \prod_{i=1}^{r}G_{S_i}.
\]
若 \(S_1\subsetneq \cdots \subsetneq S_r\) 构成严格链，则对每个 \(i\le j\) 都有 \(S_i\subseteq S_j\)，故所有上三角块均允许非零。证毕。
\end{proof}

\noindent
上述六条结果把附录中的局部命题收束为一条统一结构链：定理 \ref{thm:killo-prime-freedom-non-finitizability} 与定理 \ref{thm:xi-integer-ellipse-half-circle-dimension-gap} 可分别视为其中在 \(\NN_{\times}\) 与整数轴对齐椭圆半群上的两个显式原型。一方面，只要 Grothendieck 完成的自由秩发散，结构层同态半圆维便必然不可有限化，而可数实现仍停留在最小非零值 \(\frac12\)；另一方面，有限局部化直和对象则分裂为连通圆维轴 \(\#\mathbf S\) 与全不连通 prime-ledger 轴 \((m_p(\mathbf S))_{p\in\mathbb P}\)，并且前 \(B\) 个 prime-register 的有限截断恰需要 \(B\) 条局部化通道才能严格实现。这两条轴同时控制对象的可加不变量、严格消去律以及 \(\mathrm{Hom}\)/\(\mathrm{End}\) 的零--非零矩阵图，从而把有限素数局部化对象族组织成一个由素数支撑偏序驱动的加性结构体系。

\paragraph{黄金均衡地址动力学与 Smith 周期窗口的双重有限化}
在 Gödel--Zeckendorf 地址极限的严格自仿射闭合已经建立、模核 Dirichlet 级数的 Euler 因子与 Smith 完备性已经确定之后，还可把符号动力学侧与算术周期侧继续压缩为两类彼此独立而又同样刚性的有限证书。

\begin{conclusion}[Gödel--Zeckendorf 地址极限与黄金均衡子移位的严格共轭]\label{con:xi-time-part62d-zg-address-golden-shift-conjugacy}
沿用定理 \ref{thm:xi-zg-address-binary-self-affine} 的记号，记
\[
\Xi_B:\Sigma_{\mathrm{ZG}}\longrightarrow K_{B,v},
\qquad
\Xi_B\bigl((z_t)_{t\ge 1}\bigr):=\sum_{t\ge 1}z_tB^{t-1}v,
\]
并令 \(\sigma:\Sigma_{\mathrm{ZG}}\to\Sigma_{\mathrm{ZG}}\) 为单侧移位
\[
\sigma(z_1,z_2,z_3,\dots)=(z_2,z_3,\dots).
\]
由定理 \ref{thm:xi-zg-address-binary-self-affine} 的严格不交并分解
\[
K_{B,v}=BK_{B,v}\sqcup\bigl(v+B^2K_{B,v}\bigr),
\]
定义分段映射
\[
\mathcal S_B:K_{B,v}\to K_{B,v},
\qquad
\mathcal S_B(x):=
\begin{cases}
x/B,& x\in BK_{B,v},\\[2pt]
(x-v)/B,& x\in v+B^2K_{B,v}.
\end{cases}
\]
则 \(\mathcal S_B\) 良定义，且
\[
\Xi_B\circ \sigma=\mathcal S_B\circ \Xi_B.
\]
因此 \((K_{B,v},\mathcal S_B)\) 与黄金均衡一侧子移位 \((\Sigma_{\mathrm{ZG}},\sigma)\) 拓扑共轭。
\end{conclusion}
\begin{proof}
由定理 \ref{thm:xi-zg-address-binary-self-affine}，地址映射 \(\Xi_B\) 在 \(\Sigma_{\mathrm{ZG}}\) 上单射，且其像正是 \(K_{B,v}\)。又 \(\Sigma_{\mathrm{ZG}}\) 是紧空间 \(\{0,1\}^{\NN}\) 的闭子集，而 \(K_{B,v}\subset L_v\) 位于 Hausdorff 的 \(2\)-进地址线中。由定理 \ref{thm:xi-time-part9g-holographic-prefix-isometry-on-line}，\(\Xi_B\) 是连续的，故它是从 \(\Sigma_{\mathrm{ZG}}\) 到 \(K_{B,v}\) 的同胚。

定理 \ref{thm:xi-zg-address-binary-self-affine} 已给出严格不交并分解
\[
K_{B,v}=BK_{B,v}\sqcup\bigl(v+B^2K_{B,v}\bigr),
\]
故 \(\mathcal S_B\) 的两条定义分支互不冲突，从而良定义。

现验证共轭关系。取任意 \(z=(z_t)_{t\ge 1}\in\Sigma_{\mathrm{ZG}}\)。若 \(z_1=0\)，则
\[
\Xi_B(z)
=
\sum_{t\ge 2}z_tB^{t-1}v
=
B\sum_{u\ge 1}z_{u+1}B^{u-1}v
=
B\,\Xi_B(\sigma z),
\]
故
\[
\mathcal S_B(\Xi_B(z))=\Xi_B(\sigma z).
\]

若 \(z_1=1\)，则 admissibility 强制 \(z_2=0\)，于是
\[
\Xi_B(z)
=
v+\sum_{t\ge 3}z_tB^{t-1}v
=
v+B^2\sum_{u\ge 1}z_{u+2}B^{u-1}v.
\]
另一方面，
\[
\Xi_B(\sigma z)
=
\Xi_B(0,z_3,z_4,\dots)
=
B\sum_{u\ge 1}z_{u+2}B^{u-1}v,
\]
从而
\[
\mathcal S_B(\Xi_B(z))
=
\frac{\Xi_B(z)-v}{B}
=
\Xi_B(\sigma z).
\]
两种情形合并即得
\[
\Xi_B\circ \sigma=\mathcal S_B\circ \Xi_B.
\]
因此 \((K_{B,v},\mathcal S_B)\) 与 \((\Sigma_{\mathrm{ZG}},\sigma)\) 拓扑共轭。证毕。
\end{proof}

\begin{conclusion}[地址动力系统的 Lucas 周期谱、拓扑熵与 Artin--Mazur \texorpdfstring{$\zeta$}{zeta} 闭式]\label{con:xi-time-part62d-zg-address-lucas-entropy-zeta}
沿用结论 \ref{con:xi-time-part62d-zg-address-golden-shift-conjugacy} 的记号，并令黄金均衡子移位的邻接矩阵为
\[
M_{\mathrm G}:=
\begin{pmatrix}
1&1\\
1&0
\end{pmatrix}.
\]
则对任意 \(n\ge 1\)，\(\mathcal S_B\) 的 \(n\)-周期点数满足
\[
\#\operatorname{Fix}(\mathcal S_B^n)=\tr(M_{\mathrm G}^n)=L_n,
\]
其中 \(L_n\) 为第 \(n\) 个 Lucas 数。因而：
\begin{enumerate}
  \item \(\mathcal S_B\) 的拓扑熵为
  \[
  h_{\mathrm{top}}(\mathcal S_B)=\log\varphi;
  \]
  \item 其 Artin--Mazur \(\zeta\) 函数满足
  \[
  \zeta_{\mathcal S_B}(z)
  :=
  \exp\!\left(
  \sum_{n\ge 1}\frac{\#\operatorname{Fix}(\mathcal S_B^n)}{n}z^n
  \right)
  =
  \frac{1}{1-z-z^2};
  \]
  \item 若记 primitive \(n\)-周期轨道数为 \(\mathcal O_n\)，则
  \[
  \mathcal O_n
  =
  \frac{1}{n}\sum_{d\mid n}\mu(d)L_{n/d}.
  \]
\end{enumerate}
\end{conclusion}
\begin{proof}
由结论 \ref{con:xi-time-part62d-zg-address-golden-shift-conjugacy}，\((K_{B,v},\mathcal S_B)\) 与 \((\Sigma_{\mathrm{ZG}},\sigma)\) 共轭。后者正是禁词 \(11\) 的一侧子移位，其邻接矩阵即 \(M_{\mathrm G}\)。因此对每个 \(n\ge 1\)，
\[
\#\operatorname{Fix}(\mathcal S_B^n)
=
\#\operatorname{Fix}(\sigma^n)
=
\tr(M_{\mathrm G}^n).
\]

又 Fibonacci 矩阵恒等式给出
\[
M_{\mathrm G}^n=
\begin{pmatrix}
F_{n+1}&F_n\\
F_n&F_{n-1}
\end{pmatrix},
\]
故
\[
\tr(M_{\mathrm G}^n)=F_{n+1}+F_{n-1}=L_n.
\]
这就得到周期点公式。

拓扑熵等于邻接矩阵谱半径的对数，故
\[
h_{\mathrm{top}}(\mathcal S_B)=\log\rho(M_{\mathrm G})=\log\varphi.
\]

再由有限型子移位的标准行列式公式，
\[
\sum_{n\ge 1}\frac{\tr(M_{\mathrm G}^n)}{n}z^n
=
-\log\det(I-zM_{\mathrm G}),
\]
于是
\[
\zeta_{\mathcal S_B}(z)
=
\det(I-zM_{\mathrm G})^{-1}
=
\frac{1}{1-z-z^2}.
\]

最后，周期点数与 primitive 周期轨道数满足
\[
\#\operatorname{Fix}(\mathcal S_B^n)=\sum_{d\mid n}d\,\mathcal O_d.
\]
对该恒等式施行 Möbius 反演，即得
\[
\mathcal O_n
=
\frac{1}{n}\sum_{d\mid n}\mu(d)\#\operatorname{Fix}(\mathcal S_B^{n/d})
=
\frac{1}{n}\sum_{d\mid n}\mu(d)L_{n/d}.
\]
证毕。
\end{proof}

\begin{conclusion}[Smith 模核归一化计数的严格周期性、普通生成函数与 Ces\`aro 平均]\label{con:xi-time-part62d-smith-kappa-periodic-ogf-cesaro}
设 \(A\in M_{m\times n}(\ZZ)\) 满足 \(\operatorname{rank}_{\QQ}(A)=m\)，其 Smith 不变量为
\[
d_1\mid d_2\mid\cdots\mid d_m.
\]
沿用推论 \ref{cor:killo-kernel-size-general-modulus} 的记号，对任意 \(b\ge 1\) 定义归一化模核计数
\[
\kappa_A(b)
:=
\frac{\#\ker(A_b)}{b^{\,n-m}}
=
\prod_{i=1}^{m}\gcd(d_i,b).
\]
则成立：
\begin{enumerate}
  \item \(\kappa_A\) 以 \(d_m\) 为周期，即
  \[
  \kappa_A(b+d_m)=\kappa_A(b)\qquad(\forall\,b\ge 1);
  \]
  \item 其普通生成函数
  \[
  \mathscr K_A(z):=\sum_{b\ge 1}\kappa_A(b)z^b
  \]
  是有理函数，并且精确为
  \[
  \mathscr K_A(z)
  =
  \frac{\sum_{r=1}^{d_m}\kappa_A(r)z^r}{1-z^{d_m}};
  \]
  \item 存在精确 Ces\`aro 平均
  \[
  \mu_A
  :=
  \lim_{X\to\infty}\frac{1}{X}\sum_{b\le X}\kappa_A(b)
  =
  \frac{1}{d_m}\sum_{r=1}^{d_m}\kappa_A(r),
  \]
  并且
  \[
  \sum_{b\le X}\kappa_A(b)=\mu_A X+O_A(1).
  \]
\end{enumerate}
\end{conclusion}
\begin{proof}
由推论 \ref{cor:killo-kernel-size-general-modulus}，
\[
\kappa_A(b)=\prod_{i=1}^{m}\gcd(d_i,b).
\]
由于每个 \(d_i\mid d_m\)，故对任意 \(b\ge 1\) 与每个 \(i\)，都有
\[
\gcd(d_i,b+d_m)=\gcd(d_i,b).
\]
逐项相乘即得
\[
\kappa_A(b+d_m)=\kappa_A(b),
\]
从而 \(\kappa_A\) 以 \(d_m\) 为周期。

于是
\[
\mathscr K_A(z)
=
\sum_{q\ge 0}\sum_{r=1}^{d_m}\kappa_A(qd_m+r)z^{qd_m+r}
=
\sum_{q\ge 0}\sum_{r=1}^{d_m}\kappa_A(r)z^{qd_m+r},
\]
故
\[
\mathscr K_A(z)
=
\left(\sum_{r=1}^{d_m}\kappa_A(r)z^r\right)\sum_{q\ge 0}z^{qd_m}
=
\frac{\sum_{r=1}^{d_m}\kappa_A(r)z^r}{1-z^{d_m}}.
\]

再写
\[
X=qd_m+s,\qquad 0\le s<d_m.
\]
利用周期性可得
\[
\sum_{b\le X}\kappa_A(b)
=
q\sum_{r=1}^{d_m}\kappa_A(r)+O_A(1)
=
\frac{X}{d_m}\sum_{r=1}^{d_m}\kappa_A(r)+O_A(1).
\]
除以 \(X\) 并令 \(X\to\infty\)，即得
\[
\mu_A=\frac{1}{d_m}\sum_{r=1}^{d_m}\kappa_A(r),
\]
且
\[
\sum_{b\le X}\kappa_A(b)=\mu_A X+O_A(1).
\]
证毕。
\end{proof}

\begin{conclusion}[归一化模核 Dirichlet 级数的 Hurwitz 分解与主极点留数]\label{con:xi-time-part62d-smith-kappa-hurwitz-residue}
沿用结论 \ref{con:xi-time-part62d-smith-kappa-periodic-ogf-cesaro} 的记号，并定义
\[
Z_A(s):=\sum_{b\ge 1}\frac{\kappa_A(b)}{b^s},
\qquad
\Re s>1.
\]
则该级数正是定理 \ref{thm:xi-kernel-loss-divisibility-valuation-rational-euler} 中的归一化 Dirichlet 级数 \(\Xi_A(s)\)，并满足 Hurwitz--zeta 分解
\[
Z_A(s)
=
d_m^{-s}\sum_{r=1}^{d_m}\kappa_A(r)\,
\zeta\!\left(s,\frac{r}{d_m}\right).
\]
因此 \(Z_A(s)\) 在 \(s=1\) 处具有唯一简单极点，且
\[
\Res_{s=1}Z_A(s)
=
\mu_A
=
\frac{1}{d_m}\sum_{r=1}^{d_m}\kappa_A(r).
\]
\end{conclusion}
\begin{proof}
由 \(\#\ker(A_b)=b^{\,n-m}\kappa_A(b)\)，可知
\[
Z_A(s)=\sum_{b\ge 1}\frac{\#\ker(A_b)}{b^{\,n-m+s}},
\]
故它与定理 \ref{thm:xi-kernel-loss-divisibility-valuation-rational-euler} 中的 \(\Xi_A(s)\) 一致。

再由结论 \ref{con:xi-time-part62d-smith-kappa-periodic-ogf-cesaro} 的周期性，
\[
\kappa_A(qd_m+r)=\kappa_A(r)\qquad(q\ge 0,\ 1\le r\le d_m).
\]
于是
\[
Z_A(s)
=
\sum_{r=1}^{d_m}\kappa_A(r)\sum_{q\ge 0}(qd_m+r)^{-s}
=
d_m^{-s}\sum_{r=1}^{d_m}\kappa_A(r)\sum_{q\ge 0}\left(q+\frac{r}{d_m}\right)^{-s}.
\]
右端内层级数正是 Hurwitz zeta，故得到所示分解式。

已知对任意 \(a\in(0,1]\)，\(\zeta(s,a)\) 在 \(s=1\) 处的留数都等于 \(1\)。因此
\[
\Res_{s=1}Z_A(s)
=
d_m^{-1}\sum_{r=1}^{d_m}\kappa_A(r)
=
\mu_A.
\]
证毕。
\end{proof}

\begin{conclusion}[模核总数求和主项的 Smith 周期平均律]\label{con:xi-time-part62d-smith-kernel-summatory-periodic-average}
沿用结论 \ref{con:xi-time-part62d-smith-kappa-periodic-ogf-cesaro} 的记号，并记
\[
u:=n-m\ge 0.
\]
则模核总数满足精确渐近
\[
\sum_{b\le X}\#\ker(A_b)
=
\frac{\mu_A}{u+1}X^{u+1}+O_A(X^u),
\]
其中
\[
\mu_A=\frac{1}{d_m}\sum_{r=1}^{d_m}\prod_{i=1}^{m}\gcd(d_i,r).
\]
特别地，当 \(A\) 为方阵满秩时，即 \(u=0\)，有
\[
\sum_{b\le X}\#\ker(A_b)=\mu_A X+O_A(1).
\]
\end{conclusion}
\begin{proof}
由推论 \ref{cor:killo-kernel-size-general-modulus}，
\[
\#\ker(A_b)=b^u\kappa_A(b).
\]
再按模 \(d_m\) 分组：
\[
\sum_{b\le X}\#\ker(A_b)
=
\sum_{r=1}^{d_m}\kappa_A(r)
\sum_{\substack{q\ge 0\\ qd_m+r\le X}}(qd_m+r)^u.
\]
对固定的 \(r\)，记
\[
Q_r:=\left\lfloor \frac{X-r}{d_m}\right\rfloor.
\]
则内层和为
\[
\sum_{q=0}^{Q_r}(qd_m+r)^u.
\]
这是一关于 \(Q_r\) 的次数 \(u+1\) 多项式，其首项系数为 \(d_m^u/(u+1)\)。又 \(Q_r=X/d_m+O_A(1)\)，故
\[
\sum_{q=0}^{Q_r}(qd_m+r)^u
=
\frac{d_m^u}{u+1}Q_r^{u+1}+O_A(Q_r^u)
=
\frac{X^{u+1}}{d_m(u+1)}+O_A(X^u),
\]
其中误差在有限个 \(r\) 上一致。

将此代回并对 \(r=1,\dots,d_m\) 求和，得到
\[
\sum_{b\le X}\#\ker(A_b)
=
\left(\frac{1}{d_m}\sum_{r=1}^{d_m}\kappa_A(r)\right)\frac{X^{u+1}}{u+1}
+O_A(X^u)
=
\frac{\mu_A}{u+1}X^{u+1}+O_A(X^u).
\]
当 \(u=0\) 时即化为
\[
\sum_{b\le X}\#\ker(A_b)=\mu_A X+O_A(1).
\]
证毕。
\end{proof}

\begin{conclusion}[地址 \texorpdfstring{$\zeta$}{zeta} 与 Smith 留数的双重有限化分界]\label{con:xi-time-part62d-dynamical-arithmetic-finitization-separation}
在同一“地址--账本--投影”框架内，存在两种本质不同的有限化机制：
\begin{enumerate}
  \item Gödel--Zeckendorf 地址动力学 \((K_{B,v},\mathcal S_B)\) 的 Artin--Mazur \(\zeta\) 函数满足
  \[
  \zeta_{\mathcal S_B}(z)=\frac{1}{1-z-z^2},
  \]
  其极点结构由 admissible 语法的 Fibonacci 递推矩阵 \(M_{\mathrm G}\) 控制；
  \item 整数线性投影 \(A\) 的归一化模核 Dirichlet 级数
  \[
  Z_A(s)=\sum_{b\ge 1}\frac{\kappa_A(b)}{b^s}
  \]
  在 \(s=1\) 处的主极点留数满足
  \[
  \Res_{s=1}Z_A(s)=\mu_A,
  \]
  而 \(\mu_A\) 完全由长度 \(d_m\) 的 Smith 周期窗口决定。
\end{enumerate}
因此，前者属于由递推图压缩无限时间结构的符号有限化，后者属于由周期窗口压缩全体模数层行为的算术有限化；二者都把无穷对象压缩到有限数据，但其有限证书彼此不可替代。
\end{conclusion}
\begin{proof}
第一条正是结论 \ref{con:xi-time-part62d-zg-address-lucas-entropy-zeta} 的内容：\(\zeta_{\mathcal S_B}\) 完全由禁词 \(11\) 所定义的 Fibonacci 邻接矩阵控制。第二条则由结论 \ref{con:xi-time-part62d-smith-kappa-hurwitz-residue} 给出：\(\Res_{s=1}Z_A(s)\) 完全等于一个长度 \(d_m\) 的有限周期平均。

两者相比，地址动力学侧的有限数据是递推矩阵及其闭路计数，算术模核侧的有限数据则是最长 Smith 不变量所确定的周期窗口。前者压缩的是无限时间符号链，后者压缩的是全部模数层的算术起伏，故两类有限化机制在结构上彼此独立而不可互代。证毕。
\end{proof}

\noindent
上述六条结论把本文中原本分离的两条闭合链并置到同一有限证书框架内：Gödel--Zeckendorf 地址侧由禁词语法与 Fibonacci 递推矩阵给出动力学闭式，而 Smith 模核侧则由最长 Smith 周期窗口给出平均、生成函数、留数与求和主项的完全有限表达。前者压缩的是无限时间符号结构，后者压缩的是全体模数层算术结构；二者共同说明，本文的可审计压缩不仅存在，而且能够在动力学与算术两个方向上同时落实为显式可计算的有限数据。

\paragraph{自由单子地址化、Tauberian 留数与 dyadic torsor 的统一闭合}
前述规范 Gödel 半直积、单向量 \(2^L\)-进全息像、Zeckendorf--Gödel Dirichlet 链与 Lee--Yang dyadic 分支逆极限已经分别得到结构描述。若重新回到 H.160--H.164 的原始对象记号，则这四条主线还可压缩为下列一组彼此直接衔接的结论。

\begin{theorem}[H.161 记号下全息像集的同胚与严格 cylinder 分解]\label{thm:xi-time-part62d-hologram-image-cantor-homeomorphism}
\leanverified{paper\_xi\_time\_part62d\_hologram\_image\_cantor\_homeomorphism}
设 \(E\) 为有限事件集，\(\mathrm{code}:E\hookrightarrow \NN\) 为单射，并记
\[
M:=\max_{e\in E}\mathrm{code}(e),
\qquad
B:=2^L>M,
\qquad
T:=T_2(E_{\min}).
\]
取非零向量 \(v\in T\)，并对任意无限事件流
\[
\mathbf e=(e_1,e_2,\dots)\in E^{\NN}
\]
定义
\[
X(\mathbf e):=\sum_{t\ge 1}\mathrm{code}(e_t)\,B^{t-1}v\in T,
\qquad
\Lambda_{E,v}:=X(E^{\NN}).
\]
则：
\begin{enumerate}
  \item \(X:E^{\NN}\to \Lambda_{E,v}\) 是拓扑同胚；
  \item 若记 \(D:=\mathrm{code}(E)\)，则
  \[
  \Lambda_{E,v}
  =
  \bigsqcup_{a\in D}\bigl(av+B\Lambda_{E,v}\bigr);
  \]
  \item 对任意前缀
  \[
  \mathbf a=(e_1,\dots,e_n)\in E^n
  \]
  与相应 cylinder
  \[
  [\mathbf a]:=\{\mathbf f\in E^{\NN}:\ \mathrm{pref}_n(\mathbf f)=\mathbf a\},
  \]
  若定义
  \[
  X_n(\mathbf a):=\sum_{t=1}^{n}\mathrm{code}(e_t)\,B^{t-1}v,
  \]
  则有
  \[
  X([\mathbf a])=X_n(\mathbf a)+B^n\Lambda_{E,v};
  \]
  \item \(\Lambda_{E,v}\) 是紧、全不连通的；若 \(|E|\ge 2\)，则 \(\Lambda_{E,v}\) 无孤立点。
\end{enumerate}
\end{theorem}
\begin{proof}
由于 \(|B|_2<1\) 且 \(\mathrm{code}(e_t)\in\{0,\dots,M\}\) 有界，定义 \(X(\mathbf e)\) 的级数在完备 \(\ZZ_2\)-模 \(T_v:=\ZZ_2v\subset T\) 中收敛。

现证单射。若 \(X(\mathbf e)=X(\mathbf e')\)，则
\[
\sum_{t\ge 1}\bigl(\mathrm{code}(e_t)-\mathrm{code}(e_t')\bigr)B^{t-1}v=0.
\]
由于 \(T\) 对 \(\ZZ_2\) 无挠且 \(v\neq 0\)，这等价于
\[
\sum_{t\ge 1}c_tB^{t-1}=0
\qquad \text{于 } \ZZ_2,
\]
其中
\(
c_t:=\mathrm{code}(e_t)-\mathrm{code}(e_t')\in\{-(B-1),\dots,B-1\}
\)。
若存在最小 \(n\) 使 \(c_n\neq 0\)，则
\[
\sum_{t\ge 1}c_tB^{t-1}=B^{n-1}(c_n+Bu)
\]
对某个 \(u\in \ZZ_2\) 成立。因 \(0<|c_n|<B\)，有 \(c_n\notin B\ZZ_2\)，从而 \(c_n+Bu\neq 0\)，矛盾。故所有 \(c_t=0\)，于是 \(\mathbf e=\mathbf e'\)。

若两条事件流前 \(n\) 位相同，则其像之差属于 \(B^nT_v\)，故 \(X\) 连续。又 \(E^{\NN}\) 在乘积拓扑下紧，而 \(T\) 为 Hausdorff，故连续单射 \(X\) 自动是到其像 \(\Lambda_{E,v}\) 的拓扑同胚。这证明了第 \(1\) 条。

把事件流写成首位与尾部
\[
\mathbf e=(e_1,\sigma\mathbf e),
\]
则
\[
X(\mathbf e)=\mathrm{code}(e_1)v+B\,X(\sigma\mathbf e),
\]
从而
\[
\Lambda_{E,v}\subseteq \bigcup_{a\in D}(av+B\Lambda_{E,v}).
\]
反向包含显然成立，故并式成立。若 \(av+B\Lambda_{E,v}\) 与 \(a'v+B\Lambda_{E,v}\) 相交，则存在首位分别为 \(a\) 与 \(a'\) 的两条事件流有相同像，与单射性矛盾，故并为不交并。这证明了第 \(2\) 条。

对第 \(3\) 条，只需把前 \(n\) 位与尾部拆开：
\[
X(\mathbf f)=\sum_{t=1}^{n}\mathrm{code}(e_t)\,B^{t-1}v
+B^n\sum_{u\ge 1}\mathrm{code}(f_{n+u})\,B^{u-1}v
=
X_n(\mathbf a)+B^nX(\sigma^n\mathbf f).
\]
当 \(\mathbf f\) 遍历 \([\mathbf a]\) 时，\(\sigma^n\mathbf f\) 恰遍历整个 \(E^{\NN}\)，故
\[
X([\mathbf a])=X_n(\mathbf a)+B^n\Lambda_{E,v}.
\]

最后，\(E^{\NN}\) 为紧且全不连通，故同胚像 \(\Lambda_{E,v}\) 也紧且全不连通。若 \(|E|\ge 2\)，则任意 cylinder 都可再按下一位分解为至少两个更小 cylinder，因此 \(E^{\NN}\) 无孤立点；同胚转移给出 \(\Lambda_{E,v}\) 无孤立点。证毕。
\end{proof}

\begin{theorem}[全息像集在 Tate 模块中的 Haar 零测与维数公式]\label{thm:xi-time-part62d-hologram-image-zero-haar-dimension}
沿用定理 \ref{thm:xi-time-part62d-hologram-image-cantor-homeomorphism} 的记号。赋予
\(
T\cong \ZZ_2^2
\)
标准 \(2\)-进上确界度量与归一化 Haar 测度 \(\mu_T\)。则
\[
\mu_T(\Lambda_{E,v})=0.
\]
并且 Hausdorff 维数、上 Minkowski 维数与下 Minkowski 维数一致，且满足
\[
\dim_{\mathrm H}(\Lambda_{E,v})
=
\overline{\dim}_{\mathrm M}(\Lambda_{E,v})
=
\underline{\dim}_{\mathrm M}(\Lambda_{E,v})
=
\frac{\log |E|}{L\log 2}.
\]
\end{theorem}
\begin{proof}
取 primitive 向量 \(u\in T\) 与整数 \(m\ge 0\) 使 \(v=2^m u\)，再补出一向量 \(w\in T\) 使
\[
T=\ZZ_2u\oplus \ZZ_2w.
\]
记
\[
\pi:T\longrightarrow \ZZ_2,
\qquad
\pi(au+bw):=b.
\]
则 \(T_v=\ZZ_2v\subseteq \ZZ_2u=\ker \pi\)。对每个 \(n\ge 0\)，\(\pi^{-1}(2^n\ZZ_2)\) 是 \(T\) 的开子群，且因 \(\pi_*\mu_T\) 为 \(\ZZ_2\) 上的 Haar 概率测度，有
\[
\mu_T\bigl(\pi^{-1}(2^n\ZZ_2)\bigr)=2^{-n}.
\]
由于
\[
T_v\subseteq \bigcap_{n\ge 0}\pi^{-1}(2^n\ZZ_2),
\]
故
\[
\mu_T(T_v)\le 2^{-n}\qquad (n\ge 0).
\]
令 \(n\to\infty\) 得 \(\mu_T(T_v)=0\)，从而
\[
\mu_T(\Lambda_{E,v})=0
\]
因为 \(\Lambda_{E,v}\subseteq T_v\)。

另一方面，映射
\[
\iota_v:\ZZ_2\longrightarrow T_v,
\qquad
a\longmapsto av
\]
在所选分解下满足
\[
|\,\iota_v(a)-\iota_v(b)\,|_2
=
|2^m(a-b)u|_2
=
2^{-m}|a-b|_2.
\]
故 \(\iota_v\) 是从 \(\ZZ_2\) 到 \(T_v\) 的相似映射，特别是双 Lipschitz 同构；而 \(T_v\hookrightarrow T\) 只是该相似空间的等距嵌入。于是 \(\Lambda_{E,v}\) 在 \(T\) 中的 Hausdorff 与 Minkowski 维数，等于其在 rank-one 地址线 \(T_v\) 中的相应维数。

再由结论 \ref{con:xi-time-part62b-rankone-hologram-selfsimilar-measure-dimension}，地址线口径下有
\[
\dim_{\mathrm H}(\Lambda_{E,v})
=
\dim_{\mathrm M}(\Lambda_{E,v})
=
\frac{\log |D|}{\log B}.
\]
由于 \(|D|=|E|\) 且 \(B=2^L\)，即得
\[
\frac{\log |D|}{\log B}=\frac{\log |E|}{L\log 2}.
\]
从而三种维数全部同值。证毕。
\end{proof}

\begin{theorem}[Zeckendorf--Gödel Dirichlet 级数的 Tauberian 简单极点与临界 abscissa]\label{thm:xi-time-part62d-zg-dirichlet-tauberian-pole-abscissa}
记
\[
D_{\mathrm{ZG}}(s):=\sum_{n\in \mathrm{ZG}}n^{-s}
\qquad (\Re(s)>1).
\]
则有
\[
\lim_{\sigma\downarrow 1}(\sigma-1)D_{\mathrm{ZG}}(\sigma)=\delta_{\mathrm{ZG}}=\delta(\mathrm{ZG}),
\]
因此 \(D_{\mathrm{ZG}}(s)\) 在 \(s=1\) 处具有留数 \(\delta(\mathrm{ZG})\) 的简单极点。进一步，
\[
\sigma_c(D_{\mathrm{ZG}})=1.
\]
\end{theorem}
\begin{proof}
这正是定理 \ref{thm:xi-zg-abel-residue-log-density} 的第一条与第三条的直接转写。第一条给出 Abel 型边界律
\[
\lim_{\sigma\downarrow 1}(\sigma-1)D_{\mathrm{ZG}}(\sigma)=\delta_{\mathrm{ZG}},
\]
而 \(\delta_{\mathrm{ZG}}>0\) 表明 \(s=1\) 处确为简单极点。其第三条同时给出收敛 abscissa
\[
\sigma_c(D_{\mathrm{ZG}})=1.
\]
证毕。
\end{proof}

\begin{theorem}[Lee--Yang 分支逆极限恰分裂为五个 \(T_2(E_{\min})\)-torsor]\label{thm:xi-time-part62d-leyang-branch-inverse-limit-five-torsors}
沿用 H.164 的 branch 记号。记 \(R(y_n)\) 为第 \(n\) 层分歧除子，并把其中五个平行有限 branch 族写成
\[
\Pi_n^{\mathrm{fin}}:=\bigsqcup_{i=0}^{4}\Pi_{i,n},
\qquad
\Pi_{i,n}=P_i+E_{\min}[2^n].
\]
定义分支逆极限
\[
\mathcal R_\infty:=\varprojlim_n \Pi_n^{\mathrm{fin}}
=
\bigsqcup_{i=0}^{4}\mathcal R_\infty^{(i)},
\qquad
\mathcal R_\infty^{(i)}:=\varprojlim_n \Pi_{i,n},
\]
其中过渡映射为 \([2]\) 的限制。则对每个 \(i\)，\(\mathcal R_\infty^{(i)}\) 都是
\[
T_2(E_{\min})=\varprojlim_n E_{\min}[2^n]
\]
上的主齐性空间；亦即每个 \(\mathcal R_\infty^{(i)}\) 都是一个 \(T_2(E_{\min})\)-torsor。特别地，一旦在第 \(i\) 分支上选定一条相容提升 \(\mathbf P_i\in \mathcal R_\infty^{(i)}\)，便有
\[
\mathcal R_\infty^{(i)}\cong \mathbf P_i+T_2(E_{\min}),
\]
从而非规范地得到
\[
\mathcal R_\infty\cong \bigsqcup_{i=0}^{4}\ZZ_2^2.
\]
\end{theorem}
\begin{proof}
由推论 \ref{cor:xi-terminal-zm-leyang-finite-branch-regular-4ary-address}，每一层第 \(i\) 个分支都是
\[
\Pi_{i,n}=P_i+E_{\min}[2^n]
\]
的平移余类，且映射
\[
[2]:\Pi_{i,n+1}\longrightarrow \Pi_{i,n}
\]
与 \(E_{\min}[2^{n+1}]\to E_{\min}[2^n]\) 的平移结构相容。于是逆极限
\[
\mathcal R_\infty^{(i)}=\varprojlim_n \Pi_{i,n}
\]
继承了
\[
\varprojlim_n E_{\min}[2^n]=T_2(E_{\min})
\]
的自由且传递的平移作用，因此 \(\mathcal R_\infty^{(i)}\) 是一个 \(T_2(E_{\min})\)-torsor。

一旦取定相容提升 \(\mathbf P_i\in \mathcal R_\infty^{(i)}\)，平移映射给出同构
\[
T_2(E_{\min})\stackrel{\sim}{\longrightarrow}\mathcal R_\infty^{(i)},
\qquad
\mathbf Q\longmapsto \mathbf P_i+\mathbf Q.
\]
再由 \(T_2(E_{\min})\cong \ZZ_2^2\)，便得到各分支的非规范 \(\ZZ_2^2\) 识别；对五个分支作不交并即得结论。证毕。
\end{proof}

\begin{theorem}[事件流地址塔与 Lee--Yang 分支塔的共同 pro-\(2\) 截断律]\label{thm:xi-time-part62d-common-pro2-truncation-law}
沿用定理 \ref{thm:xi-time-part62d-hologram-image-cantor-homeomorphism}、定理 \ref{thm:xi-time-part9g-holographic-prefix-concatenation-ultrametric} 与定理 \ref{thm:xi-time-part62d-leyang-branch-inverse-limit-five-torsors} 的记号。对任意事件流 \(\mathbf e\in E^{\NN}\) 与每个 \(n\ge 1\)，记
\[
\mathcal A_n(\mathbf e):=X(\mathbf e)\bmod B^nT\in T/B^nT.
\]
则 \((\mathcal A_n(\mathbf e))_{n\ge 1}\) 在自然商映射下构成相容逆系统点列，并且
\[
\mathcal A_n(\mathbf e)=X_n\!\bigl(\mathrm{pref}_n(\mathbf e)\bigr)+B^nT.
\]
因而映射
\[
\mathbf e\longmapsto \bigl(\mathcal A_n(\mathbf e)\bigr)_{n\ge 1}
\]
是单射。

另一方面，对每个 \(n\ge 0\) 与 \(i\in\{0,\dots,4\}\)，记
\[
\mathcal R_n^{(i)}:=P_i+E_{\min}[2^n].
\]
则 \((\mathcal R_n^{(i)},[2])_{n\ge 0}\) 构成五条并行的 pro-\(2\) 逆系统，并满足
\[
R(y_n)=\bigsqcup_{i=0}^{4}\mathcal R_n^{(i)},
\qquad
\deg R(y_n)=5\cdot 4^n.
\]
因此，事件流地址塔与 Lee--Yang 分支塔共享同一二进截断法则：前者按共尾子序列 \(B^nT=2^{Ln}T\) 记录单分量地址的精细化，后者按乘二投影记录五分量分支族的同步精细化。
\end{theorem}
\begin{proof}
由定理 \ref{thm:xi-time-part62d-hologram-image-cantor-homeomorphism} 的 cylinder 分解，对每个 \(n\ge 1\) 都有
\[
X(\mathbf e)\in X_n\!\bigl(\mathrm{pref}_n(\mathbf e)\bigr)+B^n\Lambda_{E,v}\subseteq X_n\!\bigl(\mathrm{pref}_n(\mathbf e)\bigr)+B^nT,
\]
故
\[
\mathcal A_n(\mathbf e)=X_n\!\bigl(\mathrm{pref}_n(\mathbf e)\bigr)+B^nT.
\]
若 \(m\ge n\)，再模去 \(B^nT\) 只会删除更深层尾项，故 \((\mathcal A_n(\mathbf e))_{n\ge 1}\) 在自然商映射下相容。

若两条事件流 \(\mathbf e,\mathbf e'\) 满足
\[
\mathcal A_n(\mathbf e)=\mathcal A_n(\mathbf e')
\qquad (n\ge 1),
\]
则
\[
X(\mathbf e)-X(\mathbf e')\in \bigcap_{n\ge 1}B^nT.
\]
由于 \(B=2^L\) 且 \(T\cong \ZZ_2^2\)，有 \(\bigcap_{n\ge 1}B^nT=\{0\}\)，故 \(X(\mathbf e)=X(\mathbf e')\)。再由定理 \ref{thm:xi-time-part62d-hologram-image-cantor-homeomorphism} 的单射性，得到 \(\mathbf e=\mathbf e'\)。

第二部分中，\(\mathcal R_n^{(i)}\) 正是定理 \ref{thm:xi-time-part62d-leyang-branch-inverse-limit-five-torsors} 所用的 \(\Pi_{i,n}\)。由推论 \ref{cor:xi-terminal-zm-leyang-finite-branch-regular-4ary-address}，映射
\[
[2]:\mathcal R_{n+1}^{(i)}\longrightarrow \mathcal R_n^{(i)}
\]
对每个 \(i\) 都与 torsion 塔 \(E_{\min}[2^{n+1}]\to E_{\min}[2^n]\) 相容，因此给出五条并行的 pro-\(2\) 逆系统。

又因为
\[
\mathcal R_0^{(i)}=P_i+E_{\min}[1]=\{P_i\},
\]
并且每一步 \([2]\) 都是四重覆盖，遂由归纳得到
\[
|\mathcal R_n^{(i)}|=4^n.
\]
对五个分支求和便有
\[
\deg R(y_n)=\sum_{i=0}^{4}|\mathcal R_n^{(i)}|=5\cdot 4^n.
\]
故结论成立。证毕。
\end{proof}

\noindent
定理 \ref{thm:xi-time-part9g-canonical-submonoid-free-monoid} 与本段五条结论共同表明，H.160--H.164 实际上闭合为一条严格的结构链：规范编码把半直积账本压缩为有限字母自由单子；faithful \(2^L\)-进仿射实现把无限事件流送入 Tate 模块中的 Cantor 型地址像；prime-index hard-core 约束再经 Dirichlet 生成函数外显为简单极点与临界 abscissa 的 Tauberian 定律；事件流地址塔与 Lee--Yang 分支塔进一步被统一到同一类 pro-\(2\) 截断框架中；而 Lee--Yang dyadic 分支塔的逆极限最终分裂为五个 \(T_2(E_{\min})\)-torsor。于是，符号事件流、\(2\)-进地址几何、Euler 型解析约束与 branch 逆极限几何便在同一框架中被统一为可逆且可审计的闭合对象。

\paragraph{全息地址像的自仿射动力学与 Bernoulli 精确维度测度}
在 H.161 记号下的全息像集同胚、严格 cylinder 分解以及 \(B\)-进全移位共轭已经确立之后，同一地址模型还可进一步给出分支逆映射的严格自仿射动力学，以及任意 Bernoulli 事件源在地址像上的精确维度闭式。

\begin{theorem}[全息地址像的严格自仿射 Cantor 动力学与分支逆映射]\label{thm:xi-time-part62dc-hologram-selfaffine-cantor-dynamics}
\leanverified{paper\_xi\_time\_part62dc\_hologram\_selfaffine\_cantor\_dynamics}
沿用定理 \ref{thm:xi-time-part62d-hologram-image-cantor-homeomorphism} 的记号，并记
\[
K_{E,v}:=\Lambda_{E,v}=X(E^{\NN}).
\]
对每个 \(a\in E\)，定义分支片
\[
K_{E,v}^{(a)}:=\mathrm{code}(a)v+B K_{E,v}\subseteq K_{E,v},
\]
以及其上的分支逆映射
\[
\psi_a:K_{E,v}^{(a)}\to K_{E,v},
\qquad
\psi_a(x):=B^{-1}\!\bigl(x-\mathrm{code}(a)v\bigr).
\]
则成立：
\begin{enumerate}
  \item 地址像具有严格不交并分解
  \[
  K_{E,v}=\bigsqcup_{a\in E}K_{E,v}^{(a)};
  \]
  \item 逐分支定义的映射
  \[
  S_{E,v}:K_{E,v}\to K_{E,v},
  \qquad
  S_{E,v}\big|_{K_{E,v}^{(a)}}:=\psi_a,
  \]
  良定义且连续；
  \item 若 \(\sigma:E^{\NN}\to E^{\NN}\) 为左移，则
  \[
  X\circ \sigma=S_{E,v}\circ X.
  \]
  更具体地，对任意
  \[
  \mathbf e=(e_1,e_2,\dots)\in E^{\NN}
  \]
  都有
  \[
  X(\sigma \mathbf e)=\psi_{e_1}\!\bigl(X(\mathbf e)\bigr).
  \]
\end{enumerate}
因此，\(K_{E,v}\) 并非仅是一个可注入的 \(2^L\)-进地址簇，而是由 \(|E|\) 个仿射分支所组成、并与单侧全移位严格共轭的自仿射 Cantor 动力系统。
\end{theorem}
\begin{proof}
定理 \ref{thm:xi-time-part62d-hologram-image-cantor-homeomorphism} 已证明 \(X:E^{\NN}\to K_{E,v}\) 为拓扑同胚，并且
\[
K_{E,v}
=
\bigsqcup_{a\in E}\bigl(\mathrm{code}(a)v+B K_{E,v}\bigr).
\]
这正是第 \(1\) 条。

若 \(x\in K_{E,v}^{(a)}\)，则由上述不交并分解可唯一写成
\[
x=\mathrm{code}(a)v+B y
\qquad (y\in K_{E,v}),
\]
从而
\[
\psi_a(x)=y
\]
良定义。由于 \(x\mapsto B^{-1}(x-\mathrm{code}(a)v)\) 在 \(T_2(E_{\min})\) 上是连续仿射映射，而每个分支片 \(K_{E,v}^{(a)}\) 在 \(K_{E,v}\) 的相对拓扑下为 clopen，故逐分支拼接得到的 \(S_{E,v}\) 也是良定义且连续的。

对任意 \(\mathbf e=(e_1,e_2,\dots)\in E^{\NN}\)，由地址级数的首项分裂式，
\[
X(\mathbf e)
=
\mathrm{code}(e_1)v+B X(\sigma \mathbf e).
\]
移项后即得
\[
X(\sigma \mathbf e)
=
B^{-1}\!\bigl(X(\mathbf e)-\mathrm{code}(e_1)v\bigr)
=
\psi_{e_1}\!\bigl(X(\mathbf e)\bigr).
\]
于是
\[
X\circ \sigma=S_{E,v}\circ X.
\]
故 \(S_{E,v}=X\circ \sigma\circ X^{-1}\)，从而 \((K_{E,v},S_{E,v})\) 与单侧全移位 \((E^{\NN},\sigma)\) 拓扑共轭。证毕。
\end{proof}

\begin{theorem}[任意 Bernoulli 事件源在全息地址像上诱导精确维度测度]\label{thm:xi-time-part62dc-hologram-bernoulli-exact-dimensional}
沿用定理 \ref{thm:xi-time-part62dc-hologram-selfaffine-cantor-dynamics} 的记号，并在 \(K_{E,v}\) 上赋予定理 \ref{thm:xi-time-part9g-holographic-prefix-isometry-on-line} 的 \(B\)-进超度量。设
\[
\mathbf p=(p_a)_{a\in E}
\]
为 \(E\) 上的概率向量，
\[
\nu_{\mathbf p}:=
\bigotimes_{n\ge 1}\left(\sum_{a\in E}p_a\delta_a\right)
\]
为相应的 Bernoulli 乘积测度，并记其地址推前
\[
\mu_{\mathbf p}:=X_*\nu_{\mathbf p}.
\]
再记 Shannon 熵
\[
H(\mathbf p):=-\sum_{a\in E}p_a\log p_a.
\]
则 \(\mu_{\mathbf p}\) 为精确维度测度，并且对 \(\mu_{\mathbf p}\)-几乎处处的 \(x\in K_{E,v}\) 都有
\[
\dim_{\mathrm{loc}}(\mu_{\mathbf p},x)
:=
\lim_{n\to\infty}
\frac{\log \mu_{\mathbf p}\!\bigl(\mathrm B_{K_{E,v}}(x,B^{-n})\bigr)}{\log B^{-n}}
=
\frac{H(\mathbf p)}{\log B}.
\]
因而
\[
\dim_{\mathrm H}(\mu_{\mathbf p})=\frac{H(\mathbf p)}{\log B}.
\]

特别地，当 \(\mathbf p\) 为均匀分布时，对任意 \(x\in K_{E,v}\) 与任意 \(n\ge 1\)，都有
\[
\mu_{\mathrm{unif}}\!\bigl(\mathrm B_{K_{E,v}}(x,B^{-n})\bigr)
=
|E|^{-n}
=
B^{-n\log_B|E|}.
\]
因此 \(\mu_{\mathrm{unif}}\) 是维数
\[
\frac{\log |E|}{\log B}
\]
的 Ahlfors 正则自然测度；尤其
\[
\dim_{\mathrm H}(K_{E,v})=\frac{\log |E|}{\log B}.
\]
\end{theorem}
\begin{proof}
取
\[
x=X(\mathbf e)\in K_{E,v},
\qquad
\mathbf e=(e_1,e_2,\dots)\in E^{\NN}.
\]
对每个 \(n\ge 1\)，记前缀
\[
a_n:=(e_1,\dots,e_n)\in E^n.
\]
由定理 \ref{thm:xi-time-part9g-holographic-prefix-isometry-on-line} 的柱--球对应，
\[
X([a_n])
=
\bigl(X_n(a_n)+B^nT_v\bigr)\cap K_{E,v}
=
\mathrm B_{K_{E,v}}(x,B^{-n}),
\]
其中 \([a_n]\subset E^{\NN}\) 为以 \(a_n\) 为前缀的 cylinder。于是
\[
\mu_{\mathbf p}\!\bigl(\mathrm B_{K_{E,v}}(x,B^{-n})\bigr)
=
\nu_{\mathbf p}([a_n])
=
p_{e_1}\cdots p_{e_n}.
\]
因此
\[
-\frac{1}{n}\log
\mu_{\mathbf p}\!\bigl(\mathrm B_{K_{E,v}}(x,B^{-n})\bigr)
=
\frac1n\sum_{j=1}^{n}\bigl(-\log p_{e_j}\bigr).
\]
对 \(\nu_{\mathbf p}\)-几乎处处的 \(\mathbf e\)，右端由独立同分布序列的大数定律收敛到
\[
\mathbb E_{\nu_{\mathbf p}}[-\log p_{e_1}]
=
H(\mathbf p).
\]
而
\[
\log B^{-n}=-n\log B,
\]
故对 \(\mu_{\mathbf p}\)-几乎处处的 \(x=X(\mathbf e)\)，局部维数满足
\[
\lim_{n\to\infty}
\frac{\log \mu_{\mathbf p}\!\bigl(\mathrm B_{K_{E,v}}(x,B^{-n})\bigr)}{\log B^{-n}}
=
\frac{H(\mathbf p)}{\log B}.
\]
这就证明了 \(\mu_{\mathbf p}\) 的 exact-dimensional 性与
\[
\dim_{\mathrm H}(\mu_{\mathbf p})=\frac{H(\mathbf p)}{\log B}.
\]

若 \(\mathbf p\) 为均匀分布，则每个长度 \(n\) 的 cylinder 的 \(\nu_{\mathrm{unif}}\)-测度都等于 \(|E|^{-n}\)，故上式立刻给出
\[
\mu_{\mathrm{unif}}\!\bigl(\mathrm B_{K_{E,v}}(x,B^{-n})\bigr)=|E|^{-n}.
\]
于是 \(\mu_{\mathrm{unif}}\) 在自然离散尺度 \(B^{-n}\) 上恰满足幂律
\[
|E|^{-n}=B^{-n\log_B|E|},
\]
从而恢复结论 \ref{con:xi-time-part62c-godel-tate-address-ahlfors-regularity} 的 Ahlfors 正则性；再结合其中的维数公式，即得
\[
\dim_{\mathrm H}(K_{E,v})=\frac{\log |E|}{\log B}.
\]
证毕。
\end{proof}

\noindent
由此，H.160--H.161 的 \(2^L\)-进全息地址链在动力学与测度两侧都达到严格闭合：地址像不仅与单侧全移位逐分支共轭，而且任意 Bernoulli 事件源都在其上诱导出熵除以 \(\log B\) 的精确维度测度；均匀源则恢复既有的 Ahlfors 正则几何常数。于是，符号流、非阿基米德自仿射动力学与信息维数三者在同一地址模型中被完全统一。

\paragraph{全息 Bernoulli 测度的不动点、自相似质量与精确维数}
在全息地址像的自仿射动力学、圆柱--球严格同构以及 Bernoulli 推前测度的 exact-dimensional 性都已建立之后，还可把这一地址系统进一步压缩为以下三条可直接调用的测度闭式。

\begin{proposition}[\texorpdfstring{$2^L$}{2^L}-进全息点的仿射 IFS 固定点方程]\label{prop:xi-time-part62dca-hologram-affine-ifs-fixedpoint-equation}
\leanverified{paper\_xi\_time\_part62dca\_hologram\_affine\_ifs\_fixedpoint\_equation}
沿用定理 \ref{thm:xi-time-part62dc-hologram-selfaffine-cantor-dynamics} 的记号。设 \(E\) 为有限事件集，
\[
\mathrm{code}:E\hookrightarrow \NN,
\qquad
B=2^L>\max_{a\in E}\mathrm{code}(a),
\qquad
T:=T_2(E_{\min}),
\]
并记
\[
X(\mathbf e)=\sum_{t\ge 1}\mathrm{code}(e_t)B^{t-1}v
\qquad
\bigl(\mathbf e=(e_t)_{t\ge 1}\in E^{\NN}\bigr).
\]
对每个 \(a\in E\)，定义仿射映射
\[
F_a:T\longrightarrow T,
\qquad
F_a(x):=\mathrm{code}(a)\,v+B x.
\]
再取任意严格正概率向量
\[
\pi=(\pi_a)_{a\in E},
\]
并记 \(\mathbf P_\pi\) 为 \(E^{\NN}\) 上的 Bernoulli 乘积测度，
\[
\mu_\pi:=X_*\mathbf P_\pi.
\]
则成立：
\begin{enumerate}
  \item \(\mu_\pi\) 满足精确不动点方程
        \[
        \mu_\pi=\sum_{a\in E}\pi_a\,(F_a)_*\mu_\pi;
        \]
  \item 对任意有限词
        \[
        \mathbf a=(a_1,\dots,a_n)\in E^n,
        \]
        若记
        \[
        U_{\mathbf a}:=F_{a_1}\circ\cdots\circ F_{a_n}(T),
        \]
        则
        \[
        \mu_\pi(U_{\mathbf a})=\pi_{a_1}\cdots \pi_{a_n}.
        \]
\end{enumerate}
\end{proposition}
\begin{proof}
对任意 \(a\in E\) 与任意尾流 \(\mathbf e'=(e_2,e_3,\dots)\in E^{\NN}\)，由地址级数的首项分裂式直接得到
\[
X(a,\mathbf e')
=
\mathrm{code}(a)\,v+B X(\mathbf e')
=
F_a(X(\mathbf e')).
\]
按第一位事件对 \(\mathbf P_\pi\) 作条件分解，便有
\[
\mu_\pi
=
X_*\mathbf P_\pi
=
\sum_{a\in E}\pi_a\,(F_a)_*X_*\mathbf P_\pi
=
\sum_{a\in E}\pi_a\,(F_a)_*\mu_\pi,
\]
从而第 \(1\) 条成立。

再对 \(\mathbf a=(a_1,\dots,a_n)\in E^n\)，记 cylinder 集
\[
[\mathbf a]
:=
\{\mathbf e\in E^{\NN}:\ e_1=a_1,\dots,e_n=a_n\}.
\]
由上式迭代可得
\[
X([\mathbf a])=U_{\mathbf a}\cap X(E^{\NN}).
\]
另一方面，\(\mathbf P_\pi([\mathbf a])=\pi_{a_1}\cdots \pi_{a_n}\)。因此
\[
\mu_\pi(U_{\mathbf a})
=
\mu_\pi\!\bigl(U_{\mathbf a}\cap X(E^{\NN})\bigr)
=
\mathbf P_\pi([\mathbf a])
=
\pi_{a_1}\cdots \pi_{a_n}.
\]
第 \(2\) 条成立。证毕。
\end{proof}

\leanverified{paper\_xi\_time\_part62dca\_hologram\_bernoulli\_exact\_dimension\_formula}
\begin{theorem}[全息 Bernoulli 测度的精确维数公式]\label{thm:xi-time-part62dca-hologram-bernoulli-exact-dimension-formula}
沿用命题 \ref{prop:xi-time-part62dca-hologram-affine-ifs-fixedpoint-equation} 的记号，并记
\[
H(\pi):=-\sum_{a\in E}\pi_a\log \pi_a.
\]
则 \(\mu_\pi\) 为 exact-dimensional 测度，并且对 \(\mu_\pi\)-几乎处处的
\[
x\in X(E^{\NN})
\]
都有
\[
\lim_{r\downarrow 0}\frac{\log \mu_\pi(B_T(x,r))}{\log r}
=
\frac{H(\pi)}{\log B},
\]
其中 \(B_T(x,r)\) 表示 \(T\) 上的 \(B\)-进超度量球。特别地，
\[
\dim_{\mathrm H}(\mu_\pi)=\frac{H(\pi)}{\log B}.
\]
\end{theorem}
\begin{proof}
取 \(x=X(\mathbf e)\)，其中
\[
\mathbf e=(e_1,e_2,\dots)\in E^{\NN}.
\]
对每个 \(n\ge 1\)，记前缀
\[
\mathbf a_n:=(e_1,\dots,e_n).
\]
由结论 \ref{con:xi-time-part62dgb-hologram-cylinder-ball-isometry}，半径 \(B^{-n}\) 的球恰对应于长度 \(n\) 的 cylinder 像：
\[
B_T(x,B^{-n})
=
U_{\mathbf a_n},
\qquad
U_{\mathbf a_n}\cap X(E^{\NN})=X([\mathbf a_n]).
\]
于是由命题 \ref{prop:xi-time-part62dca-hologram-affine-ifs-fixedpoint-equation}，
\[
\mu_\pi\!\bigl(B_T(x,B^{-n})\bigr)
=
\mu_\pi(U_{\mathbf a_n})
=
\prod_{t=1}^{n}\pi_{e_t}.
\]
从而
\[
\frac{\log \mu_\pi(B_T(x,B^{-n}))}{\log B^{-n}}
=
\frac{-\frac1n\sum_{t=1}^{n}\log \pi_{e_t}}{\log B}.
\]
在 \(\mathbf P_\pi\) 下，\((e_t)\) 为独立同分布序列，因此由强大数定律，
\[
\frac1n\sum_{t=1}^{n}(-\log \pi_{e_t})\longrightarrow H(\pi)
\qquad \mathbf P_\pi\text{-几乎处处}.
\]
经 \(X\) 推前后即得
\[
\lim_{n\to\infty}
\frac{\log \mu_\pi(B_T(x,B^{-n}))}{\log B^{-n}}
=
\frac{H(\pi)}{\log B}
\qquad \mu_\pi\text{-几乎处处}.
\]
又由于 \(T\) 上的该超度量具有离散尺度结构，对每个充分小的半径 \(r\) 都存在唯一的 \(n\) 使相应球与某个半径 \(B^{-n}\) 的球一致，故上式等价于
\[
\lim_{r\downarrow 0}\frac{\log \mu_\pi(B_T(x,r))}{\log r}
=
\frac{H(\pi)}{\log B}.
\]
因此 \(\mu_\pi\) exact-dimensional，且其 Hausdorff 维数即为该常数。证毕。
\end{proof}

\leanverified{paper\_xi\_time\_part62dca\_hologram\_bernoulli\_lq\_spectrum}
\begin{theorem}[全息 Bernoulli 推前测度的精确 \texorpdfstring{$L^q$}{Lq}-谱]\label{thm:xi-time-part62dca-hologram-bernoulli-lq-spectrum}
沿用定理 \ref{thm:xi-time-part62dca-hologram-bernoulli-exact-dimension-formula} 的记号，并额外假设
\[
\pi_a>0
\qquad
(a\in E).
\]
对任意 \(q\in\RR\)，定义第 \(n\) 层 cylinder 质量和
\[
Z_n(q):=\sum_{\mathbf a\in E^n}\mu_\pi(U_{\mathbf a})^q.
\]
则极限
\[
\tau_{\mu_\pi}(q)
:=
\lim_{n\to\infty}\frac{\log Z_n(q)}{\log B^{-n}}
\]
存在，并满足显式公式
\[
\boxed{\;
\tau_{\mu_\pi}(q)
=
-\frac{\log\sum_{a\in E}\pi_a^q}{\log B}
=
-\frac{\log\sum_{a\in E}\pi_a^q}{L\log 2}.\;}
\]
特别地，\(\mu_\pi\) 的 exact-dimensional 性与维数公式
\[
\dim_{\mathrm H}(\mu_\pi)=\frac{H(\pi)}{\log B}
\]
正是这一地址模型在 \(q=1\) 处的熵极限读出。
\end{theorem}
\begin{proof}
由命题 \ref{prop:xi-time-part62dca-hologram-affine-ifs-fixedpoint-equation}，对任意
\[
\mathbf a=(a_1,\dots,a_n)\in E^n
\]
都有
\[
\mu_\pi(U_{\mathbf a})=\pi_{a_1}\cdots \pi_{a_n}.
\]
因此
\[
Z_n(q)
=
\sum_{(a_1,\dots,a_n)\in E^n}(\pi_{a_1}\cdots \pi_{a_n})^q
=
\left(\sum_{a\in E}\pi_a^q\right)^n.
\]
于是
\[
\frac{\log Z_n(q)}{\log B^{-n}}
=
\frac{n\log\sum_{a\in E}\pi_a^q}{-n\log B}
=
-\frac{\log\sum_{a\in E}\pi_a^q}{\log B}.
\]
右端与 \(n\) 无关，故极限存在并等于所示常数。再由
\[
\log B=L\log 2
\]
得到第二种写法。

最后，定理 \ref{thm:xi-time-part62dca-hologram-bernoulli-exact-dimension-formula} 已经给出
\[
\dim_{\mathrm H}(\mu_\pi)=\frac{H(\pi)}{\log B},
\]
故这里的 \(\tau\)-公式把此前的 exact-dimensional 结论与全体 \(q\)-矩谱统一到同一个闭式中。证毕。
\end{proof}

\begin{corollary}[均匀全息像的 Ahlfors 正则性与精确 Hausdorff 量级]\label{cor:xi-time-part62dca-hologram-uniform-ahlfors-regularity}
\leanverified{paper\_xi\_time\_part62dca\_hologram\_uniform\_ahlfors\_regularity}
在定理 \ref{thm:xi-time-part62dca-hologram-bernoulli-exact-dimension-formula} 的设定下，取均匀分布
\[
\pi_a\equiv |E|^{-1},
\qquad
s:=\frac{\log |E|}{\log B},
\qquad
\mathcal X:=X(E^{\NN})\subset T.
\]
则成立：
\begin{enumerate}
  \item 对任意 \(\mathbf a\in E^n\)，有
        \[
        \mu_{\mathrm{unif}}(U_{\mathbf a})
        =
        |E|^{-n}
        =
        B^{-ns};
        \]
  \item 对任意 \(x\in\mathcal X\) 与任意 \(n\ge 0\)，有
        \[
        \mu_{\mathrm{unif}}(B_T(x,B^{-n}))
        =
        B^{-ns};
        \]
  \item \(\mathcal X\) 是 \(s\)-Ahlfors 正则的，并且
        \[
        0<\mathcal H^s(\mathcal X)<\infty.
        \]
\end{enumerate}
\end{corollary}
\begin{proof}
第 \(1\) 条由命题 \ref{prop:xi-time-part62dca-hologram-affine-ifs-fixedpoint-equation} 直接得到，因为均匀分布下
\[
\pi_{a_1}\cdots \pi_{a_n}=|E|^{-n}=B^{-ns}.
\]

再由结论 \ref{con:xi-time-part62dgb-hologram-cylinder-ball-isometry}，\(\mathcal X\) 上半径 \(B^{-n}\) 的球恰等于某个 \(U_{\mathbf a}\)，其与 \(\mathcal X\) 的交正对应于长度 \(n\) 的 cylinder 像。故第 \(2\) 条立即由第 \(1\) 条推出。

第 \(2\) 条表明，在全部离散尺度 \(r=B^{-n}\) 上都存在精确体积律
\[
\mu_{\mathrm{unif}}(B_T(x,r))=r^s
\qquad (x\in\mathcal X).
\]
这正是超度量口径下的 \(s\)-Ahlfors 正则性。于是由标准质量分布原理与覆盖估计，得到
\[
0<\mathcal H^s(\mathcal X)<\infty.
\]
证毕。
\end{proof}

\begin{corollary}[均匀全息像上的 Hausdorff 测度识别]\label{cor:xi-time-part62dca-hologram-hausdorff-measure-identification}
沿用推论 \ref{cor:xi-time-part62dca-hologram-uniform-ahlfors-regularity} 的记号，记
\[
\mathcal X:=X(E^{\NN}),
\qquad
\nu_s:=\mathcal H^s\!\restriction_{\mathcal X}.
\]
则
\[
\mathcal H^s(\mathcal X)=1,
\qquad
\nu_s=\mu_{\mathrm{unif}}.
\]
特别地，
\[
\mu_{\mathrm{unif}}
=
\frac{\mathcal H^s\!\restriction_{\mathcal X}}{\mathcal H^s(\mathcal X)}.
\]
\end{corollary}
\begin{proof}
若 \(|E|=1\)，则 \(\mathcal X\) 为单点集、\(s=0\)，而 \(\mu_{\mathrm{unif}}\) 与 \(\mathcal H^0\!\restriction_{\mathcal X}\) 都是该点上的 Dirac 概率测度，结论显然成立。
以下设 \(|E|\ge 2\)。

由结论 \ref{con:xi-time-part62dgb-hologram-cylinder-ball-isometry}，\(\mathcal X\) 中半径 \(B^{-n}\) 的闭球恰对应于长度 \(n\) 的 cylinder 像 \(U_{\mathbf a}\)。
再由推论 \ref{cor:xi-time-part62dca-hologram-uniform-ahlfors-regularity}，
\[
\mu_{\mathrm{unif}}(U_{\mathbf a})=B^{-ns}.
\]
由于 \(|E|\ge 2\)，每个长度 \(n\) cylinder 都含有两条在第 \(n+1\) 位分歧的延拓，故
\[
\operatorname{diam}(U_{\mathbf a})=B^{-n},
\qquad
\mu_{\mathrm{unif}}(U_{\mathbf a})=\operatorname{diam}(U_{\mathbf a})^s.
\]

固定任意 cylinder \(U_{\mathbf a}\)。
在当前离散超度量口径下，\(\mathcal H^s\) 可由球覆盖计算。
若 \(U_{\mathbf a}\subset \bigcup_i V_i\) 是任意由 \(\mathcal X\) 中闭球组成的可数覆盖，则每个 \(V_i\) 也是某个 cylinder，故
\[
\mu_{\mathrm{unif}}(U_{\mathbf a})
\le
\sum_i \mu_{\mathrm{unif}}(V_i)
=
\sum_i \operatorname{diam}(V_i)^s.
\]
对所有直径受控的球覆盖取下确界，得到
\[
\mathcal H^s(U_{\mathbf a})\ge \mu_{\mathrm{unif}}(U_{\mathbf a}).
\]

反过来，对任意 \(m\ge |\mathbf a|\)，\(U_{\mathbf a}\) 可由其所有长度 \(m\) 的后继 cylinder 两两不交覆盖；这类后继共有 \(|E|^{m-|\mathbf a|}\) 个，每个直径都是 \(B^{-m}\)。于是
\[
\mathcal H^s(U_{\mathbf a})
\le
|E|^{m-|\mathbf a|}B^{-ms}
=
B^{-|\mathbf a|s}
=
\mu_{\mathrm{unif}}(U_{\mathbf a}).
\]
故
\[
\mathcal H^s(U_{\mathbf a})=\mu_{\mathrm{unif}}(U_{\mathbf a})
\]
对每个 cylinder 都成立。

取空前缀对应的 \(U_{\emptyset}=\mathcal X\)，便得
\[
\mathcal H^s(\mathcal X)=\mu_{\mathrm{unif}}(\mathcal X)=1.
\]
最后，cylinder 族构成 \(\mathcal X\) 的可数 clopen \(\pi\)-系统，并生成其 Borel \(\sigma\)-代数。
有限 Borel 测度 \(\nu_s\) 与 \(\mu_{\mathrm{unif}}\) 在全部 cylinder 上一致，故由测度唯一性，
\[
\nu_s=\mu_{\mathrm{unif}}.
\]
归一化公式随即成立。证毕。
\end{proof}

\begin{corollary}[等概率全息地址测度的 Rényi 维数与多重分形谱完全塌缩]\label{cor:xi-time-part62dca-hologram-uniform-monofractal-collapse}
\leanverified{paper\_xi\_time\_part62dca\_hologram\_uniform\_monofractal\_collapse}
沿用推论 \ref{cor:xi-time-part62dca-hologram-hausdorff-measure-identification} 的记号，并设
\[
s:=\frac{\log |E|}{\log B}
=
\frac{\log |E|}{L\log 2}.
\]
则均匀测度 \(\mu_{\mathrm{unif}}\) 满足
\[
\tau_{\mu_{\mathrm{unif}}}(q)=(q-1)s
\qquad
(q\in\RR).
\]
若对 \(q\neq 1\) 定义 Rényi 维数
\[
D_q(\mu_{\mathrm{unif}}):=\frac{\tau_{\mu_{\mathrm{unif}}}(q)}{q-1},
\]
并在 \(q=1\) 处取连续延拓，则
\[
\boxed{\;
D_q(\mu_{\mathrm{unif}})=s
\qquad
(q\in\RR).\;}
\]
进一步，\(\mu_{\mathrm{unif}}\) 的多重分形谱塌缩为单点：
\[
\boxed{\;
f(\alpha)=
\begin{cases}
s,& \alpha=s,\\
-\infty,& \alpha\neq s.
\end{cases}\;}
\]
\end{corollary}
\begin{proof}
把均匀概率
\[
\pi_a=|E|^{-1}
\]
代入定理 \ref{thm:xi-time-part62dca-hologram-bernoulli-lq-spectrum}，得到
\[
\sum_{a\in E}\pi_a^q
=
|E|\cdot |E|^{-q}
=
|E|^{1-q}.
\]
故
\[
\tau_{\mu_{\mathrm{unif}}}(q)
=
-\frac{\log |E|^{1-q}}{\log B}
=
\frac{(q-1)\log|E|}{\log B}
=(q-1)s.
\]
于是对 \(q\neq 1\)，
\[
D_q(\mu_{\mathrm{unif}})
=
\frac{\tau_{\mu_{\mathrm{unif}}}(q)}{q-1}
=
s.
\]
而 \(q\to 1\) 的连续延拓同样给出 \(D_1=s\)。

另一方面，由推论 \ref{cor:xi-time-part62dca-hologram-uniform-ahlfors-regularity}，
\[
\mu_{\mathrm{unif}}(B_T(x,B^{-n}))=B^{-ns}
\qquad
(x\in \mathcal X,\ n\ge 0),
\]
故每一点的局部维数都等于 \(s\)。因此
\[
\{x:\dim_{\mathrm{loc}}(\mu_{\mathrm{unif}},x)=s\}=\mathcal X,
\]
而对 \(\alpha\neq s\)，相应层集为空。再由推论
\ref{cor:xi-time-part62dca-hologram-hausdorff-measure-identification}
中的
\[
\mathcal H^s(\mathcal X)=1,
\]
得到该唯一非空层集的 Hausdorff 维数正是 \(s\)。这就给出所示的单点多重分形谱。证毕。
\end{proof}

\noindent
因此，\(2^L\)-进全息地址并不仅把 Bernoulli 时间流嵌入一个 Cantor 型像集，而且把源分布的 Shannon 熵、逐前缀 cylinder 质量以及均匀源下的 Hausdorff 量级同时压缩为完全显式的自相似闭式。

\paragraph{环境全息地址的超度量闭式与 Zeckendorf--Gödel 链的 Riccati 纯量化}
在 H.160--H.163 的 faithful \(2^L\)-进地址实现、环境 Tate 模块中的 cylinder--ball 严格识别、以及 Zeckendorf--Gödel 矩阵递推已经固定之后，还可把两条支线进一步压缩为四条彼此衔接的闭式：地址像侧获得环境超度量下的精确等距、自相似分划与分段仿射动力学；算术侧则可把矩阵传递链降到沿素数指标耦合的标量 Riccati 系统与无穷乘积。这里的纯量化依然保留相邻指标耦合，因而并不把 \(D_{\mathrm{ZG}}\) 退化为独立素因子的标量 Euler 积。

\begin{theorem}[环境 Tate 模块中全息像集的精确自相似、超度量等距与维数闭式]\label{thm:xi-time-part62dcb-hologram-ambient-isometry-dimension}
\leanverified{paper\_xi\_time\_part62dcb\_hologram\_ambient\_isometry\_dimension}
沿用定理 \ref{thm:xi-time-part62d-hologram-image-cantor-homeomorphism} 的记号，并进一步假设
\[
v\in T\setminus 2T.
\]
记
\[
K_B:=\Lambda_{E,v}=X(E^{\NN})\subset T:=T_2(E_{\min})\cong \ZZ_2^2.
\]
对任意 \(x,y\in K_B\)，定义
\[
N_B(x,y):=\sup\{n\ge 0:\ x\equiv y\pmod{B^nT}\},
\qquad
d_B(x,y):=B^{-N_B(x,y)},
\]
并约定 \(B^{-\infty}=0\)。对任意 \(\mathbf e,\mathbf f\in E^{\NN}\)，定义最长公共前缀长度
\[
m(\mathbf e,\mathbf f):=\sup\{n\ge 0:\ \mathrm{pref}_n(\mathbf e)=\mathrm{pref}_n(\mathbf f)\},
\]
以及前缀超度量
\[
\rho_B(\mathbf e,\mathbf f):=B^{-m(\mathbf e,\mathbf f)}.
\]
则成立：
\begin{enumerate}
  \item \(K_B\) 具有严格不交自相似分解
  \[
  K_B=\bigsqcup_{a\in E}\bigl(\mathrm{code}(a)v+B K_B\bigr).
  \]
  \item 编码映射
  \[
  \Xi:E^{\NN}\to K_B,
  \qquad
  \Xi(\mathbf e):=X(\mathbf e)
  \]
  是从 \((E^{\NN},\rho_B)\) 到 \((K_B,d_B)\) 的等距同构。
  \item Hausdorff 维数与上下 Minkowski 维数一致，且
  \[
  \dim_{\mathrm H}(K_B;d_B)
  =
  \overline{\dim}_{\mathrm M}(K_B;d_B)
  =
  \underline{\dim}_{\mathrm M}(K_B;d_B)
  =
  \frac{\log |E|}{\log B}.
  \]
  此外，\(K_B\) 在 \(T\cong \ZZ_2^2\) 的归一化 Haar 测度下为零测集。
\end{enumerate}
\end{theorem}
\begin{proof}
定理 \ref{thm:xi-time-part62d-hologram-image-cantor-homeomorphism} 已给出
\[
K_B
=
\bigsqcup_{a\in E}\bigl(\mathrm{code}(a)v+B K_B\bigr),
\]
故第 \(1\) 条成立。

现证第 \(2\) 条。由于 \(v\in T\setminus 2T\)，故 \(v\) 为 primitive 向量；由定理
\ref{thm:xi-time-part62d-hologram-image-zero-haar-dimension} 证明中的分解，可取某个
\(w\in T\) 使
\[
T=\ZZ_2v\oplus \ZZ_2w.
\]
对任意 \(n\ge 0\)，有
\[
T_v\cap B^nT
=
\ZZ_2v
\cap
\bigl(B^n\ZZ_2v\oplus B^n\ZZ_2w\bigr)
=
B^n\ZZ_2v
=
B^nT_v.
\]
因此，若 \(x,y\in K_B\subset T_v\)，则
\[
x-y\in B^nT
\Longleftrightarrow
x-y\in B^nT_v.
\]
换言之，环境模 \(T\) 上的赋值与地址线 \(T_v\) 上的赋值在 \(K_B\) 上完全一致。

再取 \(\mathbf e,\mathbf f\in E^{\NN}\)，记 \(m:=m(\mathbf e,\mathbf f)\)。由定理
\ref{thm:xi-time-part9g-holographic-prefix-isometry-on-line}，
\[
\Xi(\mathbf e)-\Xi(\mathbf f)\in B^mT_v\setminus B^{m+1}T_v.
\]
结合上式即得
\[
\Xi(\mathbf e)-\Xi(\mathbf f)\in B^mT\setminus B^{m+1}T,
\]
从而
\[
N_B(\Xi(\mathbf e),\Xi(\mathbf f))=m,
\qquad
d_B(\Xi(\mathbf e),\Xi(\mathbf f))=B^{-m}=\rho_B(\mathbf e,\mathbf f).
\]
故 \(\Xi\) 为等距单射。又由定理
\ref{thm:xi-time-part62d-hologram-image-cantor-homeomorphism}，\(\Xi\) 已是到 \(K_B\) 的双射，故它是等距同构。

现证第 \(3\) 条。由第 \(2\) 条，\(K_B\) 中半径 \(B^{-n}\) 的球恰对应于长度 \(n\) 的
cylinder 像，而这类 cylinder 恰有 \(|E|^n\) 个，且彼此不交并覆盖 \(K_B\)。因此最少覆盖数满足
\[
N(K_B,B^{-n})=|E|^n.
\]
于是
\[
\frac{\log N(K_B,B^{-n})}{-\log B^{-n}}
=
\frac{n\log |E|}{n\log B}
=
\frac{\log |E|}{\log B},
\]
从而上下 Minkowski 维数同值，且等于 \(\log|E|/\log B\)。

对 Hausdorff 维数，若 \(s>\log|E|/\log B\)，则上述深度 \(n\) 覆盖给出
\[
\mathcal H^s_{B^{-n}}(K_B)
\le
|E|^nB^{-ns}
=
\bigl(|E|B^{-s}\bigr)^n\to 0,
\]
故 \(\dim_{\mathrm H}(K_B)\le \log|E|/\log B\)。反之，任意直径不超过 \(B^{-n}\) 的集合至多落入一个深度
\(n\) cylinder 中，否则将与第 \(2\) 条的等距性矛盾；因此任意此类覆盖至少需要 \(|E|^n\) 个集合，从而当
\(s<\log|E|/\log B\) 时，
\[
\mathcal H^s_{B^{-n}}(K_B)\ge |E|^nB^{-ns}\to \infty.
\]
故 \(\dim_{\mathrm H}(K_B)=\log|E|/\log B\)。

最后，\(T\cong \ZZ_2^2\) 且 \(B=2^L\)，故商 \(T/B^nT\) 的基数为 \(B^{2n}\)，每个陪集的 Haar
测度为 \(B^{-2n}\)。由第 \(1\) 条，\(K_B\) 可被 \(|E|^n\) 个 \(B^nT\)-陪集覆盖，因此
\[
\mu_T(K_B)\le |E|^nB^{-2n}.
\]
又 \(|E|\le B\)，故 \(|E|/B^2\le 1/B<1\)，令 \(n\to\infty\) 得 \(\mu_T(K_B)=0\)。证毕。
\end{proof}

\begin{theorem}[全息地址 Cantor 集上的分段仿射动力学与 Artin--Mazur \texorpdfstring{$\zeta$}{zeta} 闭式]\label{thm:xi-time-part62dcb-hologram-piecewise-affine-zeta}
\leanverified{paper\_xi\_time\_part62dcb\_hologram\_piecewise\_affine\_zeta}
沿用定理 \ref{thm:xi-time-part62dcb-hologram-ambient-isometry-dimension} 的记号。对每个
\(a\in E\)，定义
\[
K_a:=\mathrm{code}(a)v+B K_B.
\]
再定义分段仿射映射
\[
R_B:K_B\to K_B,
\qquad
R_B(x):=\frac{x-\mathrm{code}(a)v}{B}
\quad (x\in K_a).
\]
则 \(\{K_a\}_{a\in E}\) 构成 \(K_B\) 的两两不交 clopen 分划，且
\[
R_B\circ \Xi=\Xi\circ \sigma,
\]
其中 \(\sigma:E^{\NN}\to E^{\NN}\) 为单侧满移位。特别地，
\[
h_{\mathrm{top}}(R_B)=\log |E|,
\qquad
\#\operatorname{Fix}(R_B^n)=|E|^n
\quad (n\ge 1),
\]
并且其 Artin--Mazur \(\zeta\) 函数满足
\[
\zeta_{R_B}(z)
:=
\exp\!\left(\sum_{n\ge 1}\frac{\#\operatorname{Fix}(R_B^n)}{n}z^n\right)
=
\frac{1}{1-|E|z}.
\]
\end{theorem}
\begin{proof}
由定理 \ref{thm:xi-time-part62dcb-hologram-ambient-isometry-dimension} 的第 \(1\) 条，
\[
K_B=\bigsqcup_{a\in E}K_a.
\]
又每个 \(K_a\) 都可写成 \(K_B\) 与某个 \(BT\)-陪集的交，故在 \(K_B\) 的相对拓扑下为 clopen。
于是 \(R_B\) 在各片上由同一仿射公式给出，因而良定义且连续。

对任意 \(\mathbf e=(e_1,e_2,\dots)\in E^{\NN}\)，由地址级数首项分裂式，
\[
\Xi(\mathbf e)=\mathrm{code}(e_1)v+B\Xi(\sigma\mathbf e).
\]
于是 \(\Xi(\mathbf e)\in K_{e_1}\)，并且
\[
R_B(\Xi(\mathbf e))
=
\frac{\Xi(\mathbf e)-\mathrm{code}(e_1)v}{B}
=
\Xi(\sigma\mathbf e).
\]
故
\[
R_B\circ \Xi=\Xi\circ \sigma.
\]
因此 \((K_B,R_B)\) 与 \((E^{\NN},\sigma)\) 拓扑共轭。

单侧满移位的拓扑熵为 \(\log|E|\)，故
\[
h_{\mathrm{top}}(R_B)=\log|E|.
\]
又 \(\sigma^n\) 的不动点恰为长度 \(n\) 的周期词无限重复，故
\[
\#\operatorname{Fix}(\sigma^n)=|E|^n.
\]
共轭保持不动点计数，从而
\[
\#\operatorname{Fix}(R_B^n)=|E|^n.
\]
代入 Artin--Mazur \(\zeta\) 函数定义，
\[
\zeta_{R_B}(z)
=
\exp\!\left(\sum_{n\ge 1}\frac{|E|^n}{n}z^n\right)
=
\exp\!\bigl(-\log(1-|E|z)\bigr)
=
\frac{1}{1-|E|z}.
\]
证毕。
\end{proof}

\begin{theorem}[Zeckendorf--Gödel 自然密度的 Riccati--Euler 标量分解]\label{thm:xi-time-part62dcb-zg-density-riccati-euler}
\leanverified{paper\_xi\_time\_part62dcb\_zg\_density\_riccati\_euler}
沿用命题 \ref{prop:xi-zg-prefix-recursion-monotone-tail-certificate} 的记号。设
\[
(a_0,b_0)=(1,0),
\]
并对 \(i\ge 1\) 递推地定义
\[
a_i=\frac{p_i}{p_i+1}(a_{i-1}+b_{i-1}),
\qquad
b_i=\frac{1}{p_i+1}a_{i-1}.
\]
再令
\[
r_i:=\frac{b_i}{a_i}\quad (i\ge 1),
\qquad
u_i:=a_i+b_i\quad (i\ge 0),
\qquad
r_0:=0.
\]
则成立：
\begin{enumerate}
  \item 比率 \(r_i\) 满足 Riccati 递推
  \[
  r_i=\frac{1}{p_i(1+r_{i-1})}
  \qquad (i\ge 1).
  \]
  \item 合法前缀质量满足 Euler 型乘积递推
  \[
  u_i
  =
  u_{i-1}
  \left(
  1-\frac{r_{i-1}}{(p_i+1)(1+r_{i-1})}
  \right),
  \]
  从而
  \[
  u_n
  =
  \prod_{i=1}^{n}
  \left(
  1-\frac{r_{i-1}}{(p_i+1)(1+r_{i-1})}
  \right).
  \]
  \item 自然密度具有纯量闭式
  \[
  \delta_{\mathrm{ZG}}
  =
  \frac{1}{\zeta(2)}
  \prod_{i=1}^{\infty}
  \left(
  1-\frac{r_{i-1}}{(p_i+1)(1+r_{i-1})}
  \right),
  \]
  且该无穷乘积绝对收敛并严格为正。
\end{enumerate}
\end{theorem}
\begin{proof}
由递推式可知一切 \(a_i,b_i\) 都严格正，故 \(r_i\) 良定义。直接计算得
\[
r_i
=
\frac{b_i}{a_i}
=
\frac{a_{i-1}/(p_i+1)}{p_i(a_{i-1}+b_{i-1})/(p_i+1)}
=
\frac{1}{p_i(1+r_{i-1})},
\]
从而第 \(1\) 条成立。

再由 \(u_i=a_i+b_i\) 以及
\[
a_{i-1}=\frac{u_{i-1}}{1+r_{i-1}},
\]
得到
\[
u_i
=
\frac{p_i}{p_i+1}u_{i-1}
+\frac{1}{p_i+1}\frac{u_{i-1}}{1+r_{i-1}}
=
u_{i-1}
\frac{p_i(1+r_{i-1})+1}{(p_i+1)(1+r_{i-1})}.
\]
整理即得
\[
u_i
=
u_{i-1}
\left(
1-\frac{r_{i-1}}{(p_i+1)(1+r_{i-1})}
\right),
\]
迭代后给出第 \(2\) 条的有限乘积表示。

由命题 \ref{prop:xi-zg-prefix-recursion-monotone-tail-certificate}，
\[
\delta_{\mathrm{ZG}}
=
\frac{1}{\zeta(2)}\lim_{n\to\infty}u_n.
\]
因此只需证明有限乘积收敛到严格正极限。由第 \(1\) 条与 \(r_0=0\) 归纳可得
\[
0<r_i\le \frac{1}{p_i}
\qquad (i\ge 1).
\]
于是对 \(i\ge 2\)，有
\[
0\le
\frac{r_{i-1}}{(p_i+1)(1+r_{i-1})}
\le
\frac{1}{p_{i-1}(p_i+1)}.
\]
而 \(p_i\ge i+1\)，故
\[
\sum_{i\ge 2}\frac{1}{p_{i-1}(p_i+1)}<\infty.
\]
由 \(\sum |\log(1-t_i)|<\infty\) 的标准判据可知，无穷乘积
\[
\prod_{i=1}^{\infty}
\left(
1-\frac{r_{i-1}}{(p_i+1)(1+r_{i-1})}
\right)
\]
绝对收敛。又每个因子都落在 \((0,1]\) 内，故其极限严格正。第 \(3\) 条成立。证毕。
\end{proof}

\begin{theorem}[Zeckendorf--Gödel 自然密度常数的行比值 Riccati 压缩与 Euler 乘积表示]\label{thm:xi-time-part62dcb-zg-density-row-ratio-riccati-product}
\leanverified{paper\_xi\_time\_part62dcb\_zg\_density\_row\_ratio\_riccati\_product}
沿用命题 \ref{prop:xi-zg-prefix-recursion-monotone-tail-certificate} 的记号。对每个 \(n\ge 1\)，定义
\[
M_n:=
\begin{pmatrix}
\dfrac{p_n}{p_n+1} & \dfrac{p_n}{p_n+1}\\[1mm]
\dfrac{1}{p_n+1} & 0
\end{pmatrix},
\]
\[
(x_n,y_n):=(1,1)\,M_1M_2\cdots M_n,
\qquad
r_n:=\frac{y_n}{x_n},
\qquad
r_0:=1.
\]
则有递推关系
\[
r_n=\frac{p_n}{p_n+r_{n-1}},
\qquad
x_n=x_{n-1}\frac{p_n+r_{n-1}}{p_n+1}.
\]
并且 \(ZG\) 的自然密度常数满足
\[
\delta_{\mathrm{ZG}}
=
\frac{1}{\zeta(2)}\lim_{n\to\infty}x_n
=
\frac{1}{\zeta(2)}
\prod_{n=1}^{\infty}\frac{p_n+r_{n-1}}{p_n+1}.
\]
换言之，归一化前缀质量的 \(2\times 2\) 素数传递链可沿行比值坐标完全压缩为一条一维 Riccati--Euler 递推。
\end{theorem}
\begin{proof}
由命题 \ref{prop:xi-zg-prefix-recursion-monotone-tail-certificate}，
\[
(a_0,b_0)=(1,0),
\qquad
u_n:=a_n+b_n,
\]
并且对 \(n\ge 1\) 有
\[
a_n=\frac{p_n}{p_n+1}(a_{n-1}+b_{n-1}),
\qquad
b_n=\frac{a_{n-1}}{p_n+1}.
\]
因此
\[
(u_n,a_n)
=(u_{n-1},a_{n-1})
\begin{pmatrix}
\dfrac{p_n}{p_n+1} & \dfrac{p_n}{p_n+1}\\[1mm]
\dfrac{1}{p_n+1} & 0
\end{pmatrix}
=(u_{n-1},a_{n-1})M_n.
\]
再由 \((u_0,a_0)=(1,1)\) 归纳可得
\[
(x_n,y_n)=(u_n,a_n)
\qquad (n\ge 0).
\]
特别地，
\[
\delta_{\mathrm{ZG}}
=
\frac{1}{\zeta(2)}\lim_{n\to\infty}u_n
=
\frac{1}{\zeta(2)}\lim_{n\to\infty}x_n.
\]

另一方面，由矩阵乘法，
\[
x_n=\frac{p_nx_{n-1}+y_{n-1}}{p_n+1},
\qquad
y_n=\frac{p_nx_{n-1}}{p_n+1}.
\]
于是
\[
r_n
=
\frac{y_n}{x_n}
=
\frac{p_nx_{n-1}}{p_nx_{n-1}+y_{n-1}}
=
\frac{p_n}{p_n+r_{n-1}},
\]
\[
x_n
=
x_{n-1}\frac{p_n+r_{n-1}}{p_n+1}.
\]
对后一式迭代，得到
\[
x_n=\prod_{j=1}^{n}\frac{p_j+r_{j-1}}{p_j+1}.
\]
最后令 \(n\to\infty\)，即得所示 Euler 乘积表示。证毕。
\end{proof}

\begin{corollary}[Zeckendorf--Gödel 密度乘积的二阶素数误差可和性与严格正性]\label{cor:xi-time-part62dcb-zg-density-second-order-summability}
\leanverified{paper\_xi\_time\_part62dcb\_zg\_density\_second\_order\_summability}
沿用定理 \ref{thm:xi-time-part62dcb-zg-density-row-ratio-riccati-product} 的记号。则对一切 \(n\ge 1\)，有
\[
0<1-r_n=\frac{r_{n-1}}{p_n+r_{n-1}}\le \frac{1}{p_n},
\]
并且对一切 \(n\ge 2\)，有
\[
0<
1-\frac{p_n+r_{n-1}}{p_n+1}
=
\frac{1-r_{n-1}}{p_n+1}
\le
\frac{1}{p_{n-1}(p_n+1)}.
\]
因此
\[
\sum_{n\ge 2}\left|1-\frac{p_n+r_{n-1}}{p_n+1}\right|<\infty,
\]
故
\[
\prod_{n=1}^{\infty}\frac{p_n+r_{n-1}}{p_n+1}
\]
绝对收敛到一个严格正数；特别地，
\[
\delta_{\mathrm{ZG}}>0.
\]
\end{corollary}
\begin{proof}
由定理 \ref{thm:xi-time-part62dcb-zg-density-row-ratio-riccati-product} 的递推
\[
r_n=\frac{p_n}{p_n+r_{n-1}}
\]
直接得到
\[
1-r_n=\frac{r_{n-1}}{p_n+r_{n-1}},
\]
并由 \(0<r_{n-1}\le 1\) 推出
\[
0<1-r_n\le \frac{1}{p_n}.
\]
再由
\[
1-\frac{p_n+r_{n-1}}{p_n+1}
=
\frac{1-r_{n-1}}{p_n+1}
\]
以及上一步对 \(r_{n-1}\) 的估计
\[
1-r_{n-1}\le \frac{1}{p_{n-1}},
\]
便得第二组不等式。

由于 \(p_n\ge n+1\)，故
\[
\sum_{n\ge 2}\frac{1}{p_{n-1}(p_n+1)}<\infty.
\]
于是无穷乘积的偏离 \(1\) 的量绝对可和，从而该乘积绝对收敛。又每个因子都位于 \((0,1]\) 内，故其极限严格为正。最后乘上 \(1/\zeta(2)>0\)，即得
\[
\delta_{\mathrm{ZG}}>0.
\]
证毕。
\end{proof}

\begin{theorem}[Zeckendorf--Gödel Dirichlet 链的 Riccati 纯量化与正留数简单极点]\label{thm:xi-time-part62dcb-zg-dirichlet-riccati-residue}
\leanverified{paper\_xi\_time\_part62dcb\_zg\_dirichlet\_riccati\_residue}
沿用定理 \ref{thm:killo-zg-dirichlet-matrix-euler} 的递推记号。对正实参数 \(\sigma>1\)，记
\[
D_{\mathrm{ZG}}(\sigma):=\sum_{n\in \mathrm{ZG}}n^{-\sigma},
\]
并定义
\[
R_n(\sigma):=\frac{B_n(\sigma)}{A_n(\sigma)},
\qquad
R_0(\sigma):=0.
\]
则成立：
\begin{enumerate}
  \item 比率满足标量 Riccati 递推
  \[
  R_n(\sigma)=\frac{p_n^{-\sigma}}{1+R_{n-1}(\sigma)}
  \qquad (n\ge 1).
  \]
  \item \(D_{\mathrm{ZG}}(\sigma)\) 具有沿素数指标链耦合的标量无穷乘积表示
  \[
  D_{\mathrm{ZG}}(\sigma)
  =
  \prod_{n=1}^{\infty}\bigl(1+R_n(\sigma)\bigr),
  \qquad
  \sigma>1.
  \]
  \item 当 \(\sigma\to 1^+\) 时，
  \[
  (\sigma-1)D_{\mathrm{ZG}}(\sigma)\longrightarrow \delta_{\mathrm{ZG}}>0.
  \]
  因而 \(D_{\mathrm{ZG}}\) 在 \(s=1\) 处具有正留数简单极点，且该留数恰等于自然密度 \(\delta_{\mathrm{ZG}}\)。
\end{enumerate}
这里第 \(2\) 条中的无穷乘积是对矩阵传递链的标量化重写，而不是独立素因子的 Euler 因子分解；因此与定理 \ref{thm:xi-time-part62dgc-zg-matrix-euler-no-scalar-product} 并不矛盾。
\end{theorem}
\begin{proof}
由定理 \ref{thm:killo-zg-dirichlet-matrix-euler} 的递推，
\[
A_n(\sigma)=A_{n-1}(\sigma)+B_{n-1}(\sigma),
\qquad
B_n(\sigma)=p_n^{-\sigma}A_{n-1}(\sigma),
\qquad
(A_0,B_0)=(1,0).
\]
由于 \(\sigma>1\) 时一切量皆严格正，\(R_n(\sigma)\) 良定义，并且
\[
R_n(\sigma)
=
\frac{B_n(\sigma)}{A_n(\sigma)}
=
\frac{p_n^{-\sigma}A_{n-1}(\sigma)}{A_{n-1}(\sigma)+B_{n-1}(\sigma)}
=
\frac{p_n^{-\sigma}}{1+R_{n-1}(\sigma)}.
\]
第 \(1\) 条成立。

再记
\[
Z_n(\sigma):=A_n(\sigma)+B_n(\sigma).
\]
由上式可知 \(A_n(\sigma)=Z_{n-1}(\sigma)\)，从而
\[
Z_n(\sigma)
=
A_n(\sigma)\bigl(1+R_n(\sigma)\bigr)
=
Z_{n-1}(\sigma)\bigl(1+R_n(\sigma)\bigr).
\]
结合 \(Z_0(\sigma)=1\)，归纳得到
\[
Z_n(\sigma)=\prod_{j=1}^{n}\bigl(1+R_j(\sigma)\bigr).
\]
另一方面，定理 \ref{thm:killo-zg-dirichlet-matrix-euler} 已证明
\[
D_{\mathrm{ZG}}(\sigma)=\lim_{n\to\infty}Z_n(\sigma).
\]
又由第 \(1\) 条与 \(R_0(\sigma)=0\) 归纳可得
\[
0\le R_n(\sigma)\le p_n^{-\sigma},
\]
故
\[
\sum_{n\ge 1}R_n(\sigma)\le \sum_{n\ge 1}p_n^{-\sigma}<\infty.
\]
于是无穷乘积 \(\prod_{n\ge 1}(1+R_n(\sigma))\) 绝对收敛，并收敛到
\(D_{\mathrm{ZG}}(\sigma)\)。第 \(2\) 条得证。

最后，第 \(3\) 条正是定理 \ref{thm:xi-zg-abel-residue-log-density} 的 Abel 边界律：
\[
\lim_{\sigma\downarrow 1}(\sigma-1)D_{\mathrm{ZG}}(\sigma)=\delta_{\mathrm{ZG}}.
\]
由于 \(\delta_{\mathrm{ZG}}>0\)，故 \(s=1\) 确为正留数简单极点，且留数恰等于自然密度。证毕。
\end{proof}

\noindent
由此，H.160--H.161 的 \(2^L\)-进全息地址链与 H.162--H.163 的 Zeckendorf--Gödel 算术链在同一层面进一步闭合：地址像在环境 Tate 模块中不仅保持严格等距的 Cantor 自相似结构，而且其自然动力学能够用显式分段仿射公式与 Artin--Mazur \(\zeta\) 函数完全刻画；与此同时，prime-index hard-core 约束虽不能分解为独立素因子的标量 Euler 乘积，却仍可沿素数指标链被纯量化为 Riccati 递推与绝对收敛无穷乘积，并在 \(s=1\) 处以正留数简单极点与自然密度严格闭合。

\paragraph{可寻址 Gödel 外置的源依赖最优复杂度与 majorization 序结构}
在全息地址像的精确维度测度已经取得闭式、而逐坐标可寻址 Gödel 整数化的纯估值极小化器也已完成分类之后，同一 prime-address 框架还可进一步给出任意源分布下的期望复杂度闭式，以及由 majorization 完全控制的复杂度序结构。

\begin{theorem}[可寻址 Gödel 整数化在任意源分布下的最小期望复杂度闭式]\label{thm:xi-time-part62dd-addressable-godel-expected-complexity-source-optimality}
\leanverified{paper\_xi\_time\_part62dd\_addressable\_godel\_expected\_complexity\_source\_optimality}
沿用定理 \ref{thm:xi-addressable-godel-product-lower-bound-classification} 的设定。设 \(\mathbb P\) 为 \(\Sigma^T\) 上任意概率分布，记随机字
\[
S=(S_1,\dots,S_T)\sim \mathbb P,
\]
并对每个 \(1\le t\le T\) 记其边缘分布
\[
\pi_t(a):=\mathbb P(S_t=a)
\qquad (a\in\Sigma).
\]
把 \(\pi_t\) 的坐标按非增序排列为
\[
\pi_t^\downarrow(0)\ge \pi_t^\downarrow(1)\ge \cdots \ge \pi_t^\downarrow(q-1).
\]
对任意逐坐标可寻址编码
\[
E_T:\Sigma^T\hookrightarrow \NN_{\ge 1},
\]
定义期望对数复杂度
\[
\mathcal L(E_T;\mathbb P):=\mathbb E_{\mathbb P}\bigl[\log E_T(S)\bigr].
\]
则有精确闭式
\[
\inf_{E_T}\mathcal L(E_T;\mathbb P)
=
\sum_{t=1}^{T}
\left(
\sum_{j=0}^{q-1}j\,\pi_t^\downarrow(j)
\right)\log p_t.
\]

该下界总可由纯估值编码达到：若对每个 \(t\) 选取双射
\[
\iota_t:\Sigma\xrightarrow{\sim}\{0,1,\dots,q-1\}
\]
满足
\[
\pi_t(a)>\pi_t(b)\Longrightarrow \iota_t(a)<\iota_t(b),
\]
则
\[
E_T^{\min}(s):=\prod_{t=1}^{T}p_t^{\,\iota_t(s_t)}
\qquad (s\in\Sigma^T)
\]
满足
\[
\mathcal L(E_T^{\min};\mathbb P)
=
\sum_{t=1}^{T}
\left(
\sum_{j=0}^{q-1}j\,\pi_t^\downarrow(j)
\right)\log p_t.
\]

更一般地，任一极小化器都必须在 \(\mathbb P\)-几乎处处与某个上述纯估值编码一致。若进一步假设 \(\mathbb P\) 在 \(\Sigma^T\) 上满支撑，则全部极小化器恰为这类按边缘概率逆序排列的纯估值编码。
\end{theorem}
\begin{proof}
对每个 \(t\) 与 \(a\in\Sigma\)，定义估值支持
\[
\mathcal V_{t,a}:=
\{v_{p_t}(E_T(s)):\ s_t=a\}\subset \ZZ_{\ge 0},
\qquad
m_{t,a}:=\min \mathcal V_{t,a}.
\]
由于 \(s_t\) 由 \(v_{p_t}(E_T(s))\) 经 \(\delta_t\) 单独解码，不同字母 \(a\neq b\) 的估值支持必两两不交。故
\[
\{m_{t,a}:a\in\Sigma\}
\]
由 \(q\) 个互异的非负整数组成。

于是对任意 \(s\in\Sigma^T\)，有
\[
v_{p_t}(E_T(s))\ge m_{t,s_t}
\qquad (1\le t\le T),
\]
从而
\[
\begin{aligned}
\mathcal L(E_T;\mathbb P)
&=
\sum_{s\in\Sigma^T}\mathbb P(s)\log E_T(s)\\
&=
\sum_{s\in\Sigma^T}\mathbb P(s)\sum_{p}v_p(E_T(s))\log p\\
&\ge
\sum_{t=1}^{T}\log p_t
\sum_{s\in\Sigma^T}\mathbb P(s)\,v_{p_t}(E_T(s))\\
&\ge
\sum_{t=1}^{T}\log p_t
\sum_{a\in\Sigma}\pi_t(a)\,m_{t,a}.
\end{aligned}
\]

现固定 \(t\)。把 \(\{m_{t,a}:a\in\Sigma\}\) 递增排列为
\[
n_{t,0}<n_{t,1}<\cdots <n_{t,q-1}.
\]
由于它们是 \(q\) 个互异非负整数，必有
\[
n_{t,j}\ge j
\qquad (0\le j\le q-1).
\]
对给定的递增序列 \((n_{t,j})\)，重排不等式表明
\[
\sum_{a\in\Sigma}\pi_t(a)\,m_{t,a}
\ge
\sum_{j=0}^{q-1}\pi_t^\downarrow(j)\,n_{t,j}
\ge
\sum_{j=0}^{q-1}j\,\pi_t^\downarrow(j).
\]
代回即可得到
\[
\mathcal L(E_T;\mathbb P)
\ge
\sum_{t=1}^{T}
\left(
\sum_{j=0}^{q-1}j\,\pi_t^\downarrow(j)
\right)\log p_t.
\]

反过来，若取满足概率逆序条件的双射 \(\iota_t\)，并定义
\[
E_T^{\min}(s)=\prod_{t=1}^{T}p_t^{\,\iota_t(s_t)},
\]
则显然仍满足逐坐标可寻址性，只需取
\[
\delta_t=\iota_t^{-1}.
\]
并且
\[
\log E_T^{\min}(s)=\sum_{t=1}^{T}\iota_t(s_t)\log p_t.
\]
对 \(\mathbb P\) 取期望，得到
\[
\mathcal L(E_T^{\min};\mathbb P)
=
\sum_{t=1}^{T}
\left(
\sum_{a\in\Sigma}\pi_t(a)\iota_t(a)
\right)\log p_t
=
\sum_{t=1}^{T}
\left(
\sum_{j=0}^{q-1}j\,\pi_t^\downarrow(j)
\right)\log p_t.
\]
故上述下界可达。

若某个编码达到极小值，则上面链中的全部不等式都必须取等。于是首先，所有非指定素数的估值贡献在 \(\mathbb P\)-几乎处处都必须消失；其次，对每个 \(t\)，随机变量
\[
v_{p_t}(E_T(S))-m_{t,S_t}
\]
在 \(\mathbb P\)-几乎处处都必须为零。故任一极小化器都在 \(\mathbb P\)-几乎处处与某个纯估值编码一致。若 \(\mathbb P\) 在 \(\Sigma^T\) 上满支撑，则上述等式事实上对每个 \(s\in\Sigma^T\) 都成立；于是编码在每个坐标上只能依赖于单个字母，并且其指数集合必为 \(\{0,\dots,q-1\}\) 的某个按边缘概率逆序排列，故全部极小化器恰为所述纯估值编码。证毕。
\end{proof}

\begin{corollary}[最小期望复杂度关于 majorization 为 Schur--凹]\label{cor:xi-time-part62dd-addressable-godel-majorization-schur-concavity}
\leanverified{paper\_xi\_time\_part62dd\_addressable\_godel\_majorization\_schur\_concavity}
在定理 \ref{thm:xi-time-part62dd-addressable-godel-expected-complexity-source-optimality} 的记号下，定义单坐标代价泛函
\[
\mathfrak c_q(\pi):=\sum_{j=0}^{q-1}j\,\pi^\downarrow(j),
\qquad
\pi\in\Delta(\Sigma),\ |\Sigma|=q.
\]
则 \(\mathfrak c_q\) 关于 majorization 是 Schur--凹的：若
\[
\pi\succ \sigma,
\]
则
\[
\mathfrak c_q(\pi)\le \mathfrak c_q(\sigma).
\]
因此
\[
\mathcal L_{\min}(\mathbb P)
:=
\inf_{E_T}\mathcal L(E_T;\mathbb P)
=
\sum_{t=1}^{T}\mathfrak c_q(\pi_t)\log p_t
\]
在每个边缘变量上都是 Schur--凹函数。

特别地：
\begin{enumerate}
  \item 若每个 \(\pi_t\) 都是点质量，则
  \[
  \mathcal L_{\min}(\mathbb P)=0.
  \]
  \item 若每个 \(\pi_t\) 都是均匀分布 \(u_q\)，则
  \[
  \mathcal L_{\min}(\mathbb P)
  =
  \frac{q-1}{2}\sum_{t=1}^{T}\log p_t,
  \]
  并且这在同字母表大小的全部源中取得最大值。
\end{enumerate}
\end{corollary}
\begin{proof}
把 \(\mathfrak c_q\) 改写为
\[
\mathfrak c_q(\pi)
=
\sum_{j=0}^{q-1}j\,\pi^\downarrow(j)
=
\sum_{k=1}^{q-1}\sum_{j=k}^{q-1}\pi^\downarrow(j)
=
\sum_{k=1}^{q-1}\left(1-\sum_{j=0}^{k-1}\pi^\downarrow(j)\right).
\]
若 \(\pi\succ \sigma\)，则按 majorization 的定义，
\[
\sum_{j=0}^{k-1}\pi^\downarrow(j)\ge \sum_{j=0}^{k-1}\sigma^\downarrow(j)
\qquad (1\le k\le q-1).
\]
逐项相减并求和即得
\[
\mathfrak c_q(\pi)\le \mathfrak c_q(\sigma).
\]
故 \(\mathfrak c_q\) 为 Schur--凹函数。

由定理 \ref{thm:xi-time-part62dd-addressable-godel-expected-complexity-source-optimality} 的闭式，
\[
\mathcal L_{\min}(\mathbb P)=\sum_{t=1}^{T}\mathfrak c_q(\pi_t)\log p_t,
\]
且 \(\log p_t>0\)，故 \(\mathcal L_{\min}\) 在每个边缘分布上同样是 Schur--凹的。

若 \(\pi\) 为点质量，则 \(\pi^\downarrow=(1,0,\dots,0)\)，故
\[
\mathfrak c_q(\pi)=0.
\]
若 \(\pi=u_q\) 为均匀分布，则
\[
\mathfrak c_q(u_q)=\frac1q\sum_{j=0}^{q-1}j=\frac{q-1}{2}.
\]
又点质量 majorize 任意分布，而任意分布都 majorize 均匀分布，故 Schur--凹性说明点质量给出最小值，均匀分布给出最大值。将其代回上式即得全部结论。证毕。
\end{proof}

\noindent
至此，可寻址 Gödel 外置在源依赖口径下也达到终端闭式：最优期望复杂度不再只在均匀源上可算，而是对任意 \(\mathbb P\) 都完全分解为各坐标边缘的 Schur--凹代价之和。由此，prime-register 地址化的最优复杂度理论从均匀平均律提升为一条完整的 source-dependent 序结构。

\paragraph{Lee--Yang 分支地址几何的精确反演阈值、wreath 仿射作用与高斯谱极限}
在定理 \ref{thm:xi-time-part62d-leyang-branch-inverse-limit-five-torsors} 已将 dyadic 分支逆极限识别为五个 \(T_2(E_{\min})\)-torsor，而定理 \ref{thm:xi-time-part9g-canonical-faithful-minimal-radix-threshold} 已把 \(H_{\mathrm{can}}\) 在 Tate 模块上的规范 \(B\)-进仿射实现锁定为忠实之后，还可从有限层四叉抬升与超立方体地址模型中继续抽取出下列三条后果。

\leanverified{paper\_xi\_time\_part62d\_leyang\_lstep\_oracle\_budget}
\begin{theorem}[Lee--Yang 分支 \texorpdfstring{$\ell$}{l}-步抬升的最小二进 oracle 预算]\label{thm:xi-time-part62d-leyang-lstep-oracle-budget}
沿用推论 \ref{cor:xi-terminal-zm-leyang-finite-branch-regular-4ary-address} 的记号。对任意 \(n,\ell\ge 1\)，在有限 branch 集上定义分层投影
\[
q_{n+\ell,n}:\Pi_{n+\ell}^{\mathrm{fin}}\longrightarrow \Pi_n^{\mathrm{fin}},
\qquad
q_{n+\ell,n}\big|_{\Pi_{i,n+\ell}}:=[2^\ell].
\]
若存在编码/解码对
\[
c:\Pi_{n+\ell}^{\mathrm{fin}}\to \{0,1\}^{m},
\qquad
D:\Pi_n^{\mathrm{fin}}\times \{0,1\}^{m}\to \Pi_{n+\ell}^{\mathrm{fin}}
\]
满足
\[
D\bigl(q_{n+\ell,n}(y),c(y)\bigr)=y
\qquad
\bigl(y\in \Pi_{n+\ell}^{\mathrm{fin}}\bigr),
\]
则成立：
\begin{enumerate}
  \item 对每个 \(x\in \Pi_n^{\mathrm{fin}}\)，有精确纤维大小
  \[
  \bigl|q_{n+\ell,n}^{-1}(x)\bigr|=4^\ell;
  \]
  \item 必有
  \[
  m\ge 2\ell;
  \]
  \item 上述下界可达：存在满足同一恢复条件的编码/解码对，且恰有
  \[
  m=2\ell.
  \]
\end{enumerate}
\end{theorem}
\begin{proof}
由推论 \ref{cor:xi-terminal-zm-leyang-finite-branch-regular-4ary-address} 的第二条，对每个分支 \(i\) 与每个 \(x\in \Pi_{i,n}\)，映射
\[
[2^\ell]:\Pi_{i,n+\ell}\to \Pi_{i,n}
\]
是正则 \(4^\ell\)-重覆盖，故每个纤维都恰有 \(4^\ell\) 个点。这证明了第一条。

现证第二条。固定任意 \(x\in \Pi_n^{\mathrm{fin}}\)。若 \(y,y'\in q_{n+\ell,n}^{-1}(x)\) 且 \(c(y)=c(y')\)，则
\[
D(x,c(y))=D(x,c(y')).
\]
由恢复条件得到 \(y=y'\)。因此 \(c\) 在每个纤维上的限制都必须是单射，从而
\[
2^m\ge \bigl|q_{n+\ell,n}^{-1}(x)\bigr|=4^\ell.
\]
取二进对数即得
\[
m\ge 2\ell.
\]

最后证可达性。由于 \(\Pi_n^{\mathrm{fin}}\) 为有限集，可取一截面
\[
s:\Pi_n^{\mathrm{fin}}\to \Pi_{n+\ell}^{\mathrm{fin}},
\qquad
q_{n+\ell,n}\circ s=\mathrm{id}.
\]
对任意 \(x\in \Pi_n^{\mathrm{fin}}\)，由同一条正则覆盖结论，其纤维可写成
\[
q_{n+\ell,n}^{-1}(x)=s(x)+E_{\min}[2^\ell].
\]
再固定一组 \(\ZZ/2^\ell\ZZ\)-基，从而得到群同构
\[
E_{\min}[2^\ell]\cong (\ZZ/2^\ell\ZZ)^2,
\]
并进一步固定一个双射
\[
\beta_\ell:E_{\min}[2^\ell]\stackrel{\sim}{\longrightarrow}\{0,1\}^{2\ell}.
\]
于是对任意 \(y\in q_{n+\ell,n}^{-1}(x)\)，存在唯一
\[
Q(y)\in E_{\min}[2^\ell]
\]
使
\[
y=s(x)+Q(y).
\]
定义
\[
c(y):=\beta_\ell\bigl(Q(y)\bigr)\in\{0,1\}^{2\ell},
\]
并令
\[
D\bigl(x,\beta_\ell(Q)\bigr):=s(x)+Q.
\]
则对一切 \(y\in \Pi_{n+\ell}^{\mathrm{fin}}\) 都有
\[
D\bigl(q_{n+\ell,n}(y),c(y)\bigr)=y.
\]
故 \(2\ell\) bit 确已足够。证毕。
\end{proof}

\begin{theorem}[Lee--Yang 分支逆极限上的 wreath 型忠实仿射作用]\label{thm:xi-time-part62d-leyang-branch-wreath-affine}
沿用定理 \ref{thm:xi-time-part62d-leyang-branch-inverse-limit-five-torsors} 的记号，并对每个 \(i\in\{0,\dots,4\}\) 选取一条相容提升
\[
\mathbf P_i\in \mathcal R_\infty^{(i)}.
\]
记
\[
\tau_i:T_2(E_{\min})\stackrel{\sim}{\longrightarrow}\mathcal R_\infty^{(i)},
\qquad
x\longmapsto \mathbf P_i+x.
\]
则 \(H_{\mathrm{can}}^{5}\) 通过各分量上的 transported affine action 作用于
\[
\mathcal R_\infty=\bigsqcup_{i=0}^{4}\mathcal R_\infty^{(i)},
\]
而 \(S_5\) 通过置换五个分量指标作用于同一不交并；该置换作用规范化前述分量作用，因而诱导出一个忠实作用
\[
H_{\mathrm{can}}^{5}\rtimes S_5\ \curvearrowright\ \mathcal R_\infty.
\]
特别地，在再选定 \(T_2(E_{\min})\) 的一组 \(\ZZ_2\)-基之后，\(\mathcal R_\infty\) 非规范地识别为五片 \(\ZZ_2^2\)，而上式遂成为五片 \(2\)-进仿射平面上的 wreath 型忠实仿射作用。
\end{theorem}
\begin{proof}
由定理 \ref{thm:xi-time-part62d-leyang-branch-inverse-limit-five-torsors}，每个分量 \(\mathcal R_\infty^{(i)}\) 都是 \(T_2(E_{\min})\)-torsor；一旦选定 \(\mathbf P_i\)，映射 \(\tau_i\) 便给出与 \(T_2(E_{\min})\) 的识别。于是可把
\[
\Psi|_{H_{\mathrm{can}}}:H_{\mathrm{can}}\to \operatorname{Aff}\bigl(T_2(E_{\min})\bigr)
\]
经 \(\tau_i\) 传送到第 \(i\) 个分量，从而得到直积 \(H_{\mathrm{can}}^{5}\) 的分量仿射作用。

另一方面，\(S_5\) 在指标集合 \(\{0,\dots,4\}\) 上的自然置换显然给出 \(\mathcal R_\infty\) 的分量置换作用；并且它把第 \(i\) 个分量上的 affine action 共轭到第 \(\sigma(i)\) 个分量上的同型 action，故规范化 \(H_{\mathrm{can}}^{5}\) 的分量作用。于是得到半直积
\[
H_{\mathrm{can}}^{5}\rtimes S_5
\]
在 \(\mathcal R_\infty\) 上的作用。该分量标签置换与定理 \ref{thm:xi-terminal-zm-leyang-linear-twist-parameter-map-s5} 给出的五分支 \(S_5\)-单值化口径相容。

现证忠实性。若 \(\sigma\in S_5\) 非平凡，则它必移动至少一个分量 \(\mathcal R_\infty^{(i)}\)，故整体作用非平凡。若 \(\sigma=1\) 而
\[
(h_0,\dots,h_4)\in H_{\mathrm{can}}^{5}
\]
非平凡，则至少有一个 \(h_i\neq 1\)。在本文既定设定中 \(B=2^L>M\)，故由定理 \ref{thm:xi-time-part9g-canonical-faithful-minimal-radix-threshold}，\(\Psi|_{H_{\mathrm{can}}}\) 为忠实表示；于是 \(h_i\) 在 \(T_2(E_{\min})\) 上作用非平凡，传回 \(\mathcal R_\infty^{(i)}\) 后仍非平凡，从而整体作用非平凡。故该半直积作用忠实。证毕。
\end{proof}

\begin{theorem}[有限层 Lee--Yang 分支图的显式谱与高斯极限]\label{thm:xi-time-part62d-leyang-branch-graph-gaussian-spectrum}
对每个 \(n\ge 1\)，在每个分量
\[
\Pi_{i,n}=P_i+E_{\min}[2^n]
\]
上取定理 \ref{thm:xi-coherence-affine-subcube} 的 \(4\)-进地址标识
\[
\Gamma_n:\Pi_{i,n}\stackrel{\sim}{\longrightarrow}\{0,1,2,3\}^{n},
\]
再把每个 \(4\)-进 digit 拆成两位二进 digit，从而把 \(\Pi_{i,n}\) 识别为 \(\{0,1\}^{2n}\) 并传送标准 \(Q_{2n}\) 图结构。令
\[
G_n:=\bigsqcup_{i=0}^{4}\Pi_{i,n}
\]
为五分量 Lee--Yang 分支图。记其邻接算子为 \(A_n\)，组合 Laplacian 为
\[
L_n:=2nI-A_n.
\]
则成立：
\begin{enumerate}
  \item 邻接谱满足
  \[
  \operatorname{Spec}(A_n)=\{\,2n-2k:\ 0\le k\le 2n\,\},
  \]
  且相应重数为
  \[
  \operatorname{mult}_{A_n}(2n-2k)=5\binom{2n}{k}.
  \]
  \item 经验谱测度
  \[
  \mu_n:=
  \frac{1}{5\cdot 4^n}
  \sum_{k=0}^{2n}5\binom{2n}{k}\,
  \delta_{\frac{2n-2k}{\sqrt{2n}}}
  \]
  弱收敛到标准高斯分布 \(N(0,1)\)。
  \item 对任意 \(t\ge 0\)，热迹具有闭式
  \[
  \Tr\bigl(e^{-tL_n}\bigr)=5\bigl(1+e^{-2t}\bigr)^{2n}.
  \]
\end{enumerate}
\end{theorem}
\begin{proof}
由定理 \ref{thm:xi-coherence-affine-subcube} 的地址模型，对每个 \(i\) 与 \(n\)，分量
\[
\Pi_{i,n}
\]
可与 \(\{0,1,2,3\}^{n}\) 双射；把每一位 \(4\)-进数字展开成两位二进数字，便得到与 \(\{0,1\}^{2n}\) 的双射。将 Hamming 距离 \(1\) 的标准边关系传送回 \(\Pi_{i,n}\)，每个分量遂在图同构意义下就是一个 \(2n\)-维超立方体 \(Q_{2n}\)。故
\[
G_n\cong \bigsqcup_{i=0}^{4}Q_{2n}.
\]

对单个 \(Q_{2n}\)，由定理 \ref{thm:xi-hypercube-weighted-laplacian-heat-trace-factorization} 取权重
\[
w_1=\cdots=w_{2n}=1
\]
可得其 Laplacian 特征值为
\[
2k\qquad(0\le k\le 2n),
\]
重数为 \(\binom{2n}{k}\)。由于
\[
L_n=2nI-A_n,
\]
对应的邻接特征值便是
\[
2n-2k,
\]
重数仍为 \(\binom{2n}{k}\)。五个连通分量彼此不交，因此整体图 \(G_n\) 的重数恰乘以 \(5\)。这证明了第一条。

对第二条，若 \(K_n\sim \mathrm{Bin}(2n,\tfrac12)\)，则 \(\mu_n\) 正是随机变量
\[
\frac{2n-2K_n}{\sqrt{2n}}
\]
的分布。由 de Moivre--Laplace 中心极限定理，
\[
\frac{2K_n-2n}{\sqrt{2n}}\Longrightarrow N(0,1),
\]
取负号不改变极限分布，故 \(\mu_n\Longrightarrow N(0,1)\)。

最后，由同一超立方体 Laplacian 谱分解，单个 \(Q_{2n}\) 的热迹为
\[
\sum_{k=0}^{2n}\binom{2n}{k}e^{-2tk}
=
\bigl(1+e^{-2t}\bigr)^{2n}.
\]
再乘以五个连通分量的贡献，即得
\[
\Tr\bigl(e^{-tL_n}\bigr)=5\bigl(1+e^{-2t}\bigr)^{2n}.
\]
证毕。
\end{proof}

\noindent
定理 \ref{thm:xi-time-part62d-leyang-branch-inverse-limit-five-torsors} 连同上述三条已经把 H.159--H.164 的 Lee--Yang 主线压缩到同一条可审计的地址几何链：有限层 \(4^\ell\)-重抬升把精确反演的最小外置信息锁定在 \(2\ell\) bit；逆极限上的五分量 torsor 分解把 \(H_{\mathrm{can}}\) 的忠实仿射表示提升为 wreath 型全局对称；同一 dyadic 地址结构在有限层又生成五个 \(Q_{2n}\) 超立方体分支图，从而使邻接谱、热迹与 Gaussian 极限全部显式可算。基于同一超立方体模型，邻接特征多项式、自同构群、Kirchhoff 森林复杂度、规范地址移位的极大熵结构以及 Ihara \(\zeta\) 还可继续写成完全闭式。

\paragraph{Lee--Yang 分支图的图论闭式与地址移位的极大熵结构}
在上一段已经把有限层 Lee--Yang branch 图确定为五个 \(Q_{2n}\) 的不交并之后，同一超立方体模型还可继续给出邻接特征多项式、全自同构群、Kirchhoff 森林复杂度、规范地址移位的极大熵测度以及 Ihara \(\zeta\) 的完全闭式。

\begin{theorem}[有限层 Lee--Yang 分支图的邻接特征多项式与邻接热迹闭式]\label{thm:xi-time-part62d-leyang-branch-adjacency-charpoly-heat}
沿用定理 \ref{thm:xi-time-part62d-leyang-branch-graph-gaussian-spectrum} 的记号。则
\[
\chi_{G_n}(x):=\det(xI-A_n)
=
\prod_{k=0}^{2n}\bigl(x-(2n-2k)\bigr)^{5\binom{2n}{k}}.
\]
并且对任意 \(t\in\CC\) 都有
\[
\Tr\!\bigl(e^{tA_n}\bigr)
=
5\sum_{k=0}^{2n}\binom{2n}{k}e^{t(2n-2k)}
=
5(2\cosh t)^{2n}.
\]
\end{theorem}
\begin{proof}
定理 \ref{thm:xi-time-part62d-leyang-branch-graph-gaussian-spectrum} 已给出 \(A_n\) 的全部特征值与重数，即
\[
2n-2k
\qquad(0\le k\le 2n)
\]
的重数为 \(5\binom{2n}{k}\)。特征多项式即为按重数取乘积，得到第一式。

同一谱分解同时给出
\[
\Tr\!\bigl(e^{tA_n}\bigr)
=
5\sum_{k=0}^{2n}\binom{2n}{k}e^{t(2n-2k)}.
\]
把 \(e^{2nt}\) 提出并作二项式展开，便得
\[
\Tr\!\bigl(e^{tA_n}\bigr)
=
5e^{2nt}\sum_{k=0}^{2n}\binom{2n}{k}e^{-2kt}
=
5(e^t+e^{-t})^{2n}
=
5(2\cosh t)^{2n}.
\]
证毕。
\end{proof}

\begin{theorem}[有限层 Lee--Yang 分支图的全自同构群为超八面体花环积]\label{thm:xi-time-part62d-leyang-branch-graph-automorphism-wreath}
沿用定理 \ref{thm:xi-time-part62d-leyang-branch-graph-gaussian-spectrum} 的记号。则有自然群同构
\[
\operatorname{Aut}(G_n)
\cong
\operatorname{Aut}(Q_{2n})\wr S_5
\cong
\bigl((\ZZ/2\ZZ)^{2n}\rtimes S_{2n}\bigr)\wr S_5.
\]
\end{theorem}
\begin{proof}
由定理 \ref{thm:xi-time-part62d-leyang-branch-graph-gaussian-spectrum}，
\[
G_n\cong \bigsqcup_{i=0}^{4}Q_{2n}.
\]
任一图自同构必须将连通分量映到连通分量，因此先给出一个 \(S_5\) 的分量置换；在每个分量内部，则独立给出一个 \(Q_{2n}\) 的图自同构。于是
\[
\operatorname{Aut}(G_n)\cong \operatorname{Aut}(Q_{2n})^5\rtimes S_5
=
\operatorname{Aut}(Q_{2n})\wr S_5.
\]

另一方面，对单个超立方体 \(Q_d\)，任选一个顶点作为原点像，再任选该顶点处 \(d\) 条边方向的排列，便唯一确定一个图自同构；因此
\[
\operatorname{Aut}(Q_d)\cong (\ZZ/2\ZZ)^d\rtimes S_d.
\]
取 \(d=2n\) 即得所示公式。证毕。
\end{proof}

\begin{theorem}[有限层 Lee--Yang 分支图的 Kirchhoff 森林数闭式]\label{thm:xi-time-part62d-leyang-branch-kirchhoff-forest}
记 \(\tau^\sharp(G_n)\) 为 \(G_n\) 的组分保持生成森林数，即在每个连通分量内各取一棵生成树而得到的生成森林总数。则
\[
\tau^\sharp(G_n)
=
\Biggl(
2^{2^{2n}-2n-1}
\prod_{k=1}^{2n}k^{\binom{2n}{k}}
\Biggr)^5.
\]
\end{theorem}
\begin{proof}
由定理 \ref{thm:xi-time-part62d-leyang-branch-graph-gaussian-spectrum}，\(G_n\) 是五个 \(Q_{2n}\) 的不交并。对单个 \(Q_{2n}\)，定理 \ref{thm:xi-hypercube-weighted-spanning-tree-spectrum-product} 在 \(k=2n\) 且 \(w_1=\cdots=w_{2n}=1\) 时给出
\[
\tau(Q_{2n})
=
2^{2^{2n}-2n-1}
\prod_{\varnothing\neq S\subseteq[2n]}|S|.
\]
按子集基数分组即得
\[
\tau(Q_{2n})
=
2^{2^{2n}-2n-1}
\prod_{k=1}^{2n}k^{\binom{2n}{k}}.
\]
而组分保持生成森林恰是在五个互不连通且彼此同构的分量中独立各选一棵生成树，故
\[
\tau^\sharp(G_n)=\tau(Q_{2n})^5.
\]
代入上式即得结论。证毕。
\end{proof}

\begin{theorem}[Lee--Yang branch lamination 的五个遍历极大熵测度]\label{thm:xi-time-part62d-leyang-branch-maximal-entropy-measures}
沿用定理 \ref{thm:xi-time-part62d-leyang-branch-inverse-limit-five-torsors} 的记号。对每个 \(i\in\{0,\dots,4\}\)，记 \(m_i\) 为通过任一基点识别
\[
\mathcal R_\infty^{(i)}\cong T_2(E_{\min})\cong \ZZ_2^2
\]
传送到第 \(i\) 个分支分量上的 Haar 概率测度。为避免与推论 \ref{cor:xi-time-part62a-leyang-branch-artin-mazur-zeta} 中的纤维内乘二端同态混淆，下文记
\[
\sigma_{\mathrm{LY}}:\mathcal R_\infty\to \mathcal R_\infty
\]
为由正则四叉地址截断诱导的规范地址移位。则成立：
\begin{enumerate}
  \item 对每个 \(i\)，\(\sigma_{\mathrm{LY}}\big|_{\mathcal R_\infty^{(i)}}\) 与四符号单侧满移位拓扑共轭；
  \item 其拓扑熵满足
  \[
  h_{\mathrm{top}}(\sigma_{\mathrm{LY}})=\log 4=2\log 2;
  \]
  \item \(\sigma_{\mathrm{LY}}\) 的遍历极大熵测度恰为 \(m_0,\dots,m_4\)，而全部极大熵测度恰为
  \[
  \Biggl\{
  \sum_{i=0}^{4}a_i m_i:\ a_i\ge 0,\ \sum_{i=0}^{4}a_i=1
  \Biggr\}.
  \]
\end{enumerate}
\end{theorem}
\begin{proof}
由定理 \ref{thm:xi-time-part62d-leyang-branch-graph-gaussian-spectrum} 中的地址识别与推论 \ref{cor:xi-terminal-zm-leyang-finite-branch-regular-4ary-address} 的层间四重覆盖结构，对每个 \(i\) 都可把逆极限分量 \(\mathcal R_\infty^{(i)}\) 识别为
\[
\{0,1,2,3\}^{\NN},
\]
并且 \([2]\) 在该地址模型中恰对应于删去首位 \(4\)-进 digit 的左移。因此
\[
\sigma_{\mathrm{LY}}\big|_{\mathcal R_\infty^{(i)}}
\]
与四符号单侧满移位拓扑共轭，第一条成立。

四符号单侧满移位的拓扑熵为 \(\log 4\)；有限个不交并的拓扑熵等于各分量熵的最大值，故
\[
h_{\mathrm{top}}(\sigma_{\mathrm{LY}})=\log 4=2\log 2.
\]

再由标准符号动力学，四符号单侧满移位存在唯一遍历极大熵测度，即等权 Bernoulli 测度。把 \(\{0,1,2,3\}\) 重新写成二比特对后，长度 \(r\) 的 cylinder 恰对应于 \(\ZZ_2^2\) 中某个 \(2^r\ZZ_2^2\)-陪集，因此该等权 Bernoulli 测度在所选分支识别下正对应于 Haar 概率测度 \(m_i\)。于是每个分支分量恰贡献一个遍历极大熵测度。

最后，五个分量 \(\mathcal R_\infty^{(i)}\) 都是开闭且 \(\sigma_{\mathrm{LY}}\)-不变的，因此任一遍历不变概率测度都必须支撑在其中某一个分量上；故遍历极大熵测度只能是 \(m_0,\dots,m_4\)。一般极大熵测度再由遍历分解定理给出为它们的凸组合。证毕。
\end{proof}

\begin{theorem}[有限层 Lee--Yang 分支图的 Ihara \texorpdfstring{$\zeta$}{zeta} 闭式]\label{thm:xi-time-part62d-leyang-branch-ihara-zeta}
记 \(Z_{G_n}^{\mathrm{Ih}}(u)\) 为 \(G_n\) 的 Ihara \(\zeta\) 函数。则
\[
Z_{G_n}^{\mathrm{Ih}}(u)
=
\Biggl[
(1-u^2)^{-(2n-2)2^{2n-1}}
\prod_{k=0}^{2n}
\bigl(1-(2n-2k)u+(2n-1)u^2\bigr)^{-\binom{2n}{k}}
\Biggr]^5.
\]
\end{theorem}
\begin{proof}
由定理 \ref{thm:xi-time-part62d-leyang-branch-graph-gaussian-spectrum}，\(G_n\) 是五个 \(Q_{2n}\) 的不交并。故只需先计算单个 \(Q_{2n}\) 的 Ihara \(\zeta\) 函数，再取五次方。

对 \(d\)-正则图 \(Q_d\)，标准 Ihara--Bass 行列式公式给出
\[
\bigl(Z_{Q_d}^{\mathrm{Ih}}(u)\bigr)^{-1}
=
(1-u^2)^{|E(Q_d)|-|V(Q_d)|}
\det\!\bigl(I-uA_{Q_d}+(d-1)u^2I\bigr).
\]
这里
\[
|V(Q_d)|=2^d,
\qquad
|E(Q_d)|=d\,2^{d-1},
\]
并且 \(A_{Q_d}\) 的谱为
\[
d-2k
\qquad(0\le k\le d)
\]
其重数为 \(\binom{d}{k}\)。因此
\[
Z_{Q_d}^{\mathrm{Ih}}(u)
=
(1-u^2)^{-(d-2)2^{d-1}}
\prod_{k=0}^{d}
\bigl(1-(d-2k)u+(d-1)u^2\bigr)^{-\binom{d}{k}}.
\]
取 \(d=2n\) 即得单个 \(Q_{2n}\) 的闭式。再由 Ihara \(\zeta\) 对不交并取乘法，
\[
Z_{G_n}^{\mathrm{Ih}}(u)=\bigl(Z_{Q_{2n}}^{\mathrm{Ih}}(u)\bigr)^5,
\]
从而得到所示公式。证毕。
\end{proof}

\noindent
至此，有限层 branch 图的邻接特征多项式、闭路热迹、Kirchhoff 复杂度、全自同构群与无回溯闭路 \(\zeta\) 函数，以及逆极限 branch lamination 的极大熵动力学，便在同一条 Lee--Yang 四叉地址几何链上同时获得了完全显式的闭式描述。

\paragraph{Lee--Yang 分支图的 wreath 谱几何、dyadic affine Tate sheet 与 ZG 主极点残数}
在 Lee--Yang 有限分支图的超立方体地址、层间正则四叉结构与 Zeckendorf--Gödel 素数码的 Dirichlet 密度已经确定之后，还可把 H.158--H.164 的终端几何进一步压缩为四项可直接调用的后果。它们分别固定有限层 branch 图的全自同构、谱与 Ihara--Bass 数据，branch 塔逆极限的 affine Tate sheet 结构，以及 \(D_{\mathrm{ZG}}\) 在 \(s=1\) 处的主极点常数。

\begin{theorem}[Lee--Yang 第 \texorpdfstring{$n$}{n} 层分支图的超八面体花环对称]\label{thm:xi-time-part62dh-leyang-branch-hyperoctahedral-wreath}
沿用定理 \ref{thm:xi-time-part62d-leyang-branch-graph-gaussian-spectrum} 的记号。则有自然群同构
\[
\operatorname{Aut}(G_n)
\cong
\bigl((\ZZ/2\ZZ)^{2n}\rtimes S_{2n}\bigr)\wr S_5.
\]
\end{theorem}
\begin{proof}
这正是定理 \ref{thm:xi-time-part62d-leyang-branch-graph-automorphism-wreath} 的内容。证毕。
\end{proof}
\leanverified{paper\_xi\_time\_part62dh\_leyang\_branch\_hyperoctahedral\_wreath}

\begin{corollary}[Lee--Yang 第 \texorpdfstring{$n$}{n} 层分支图的谱、Kirchhoff 森林与 Ihara--Bass 闭式]\label{cor:xi-time-part62dh-leyang-branch-spectrum-forest-ihara}
\leanverified{paper\_xi\_time\_part62dh\_leyang\_branch\_spectrum\_forest\_ihara}
沿用定理 \ref{thm:xi-time-part62d-leyang-branch-graph-gaussian-spectrum} 的记号，记组合 Laplacian 为
\[
L_n:=2nI-A_n.
\]
则成立：
\begin{enumerate}
  \item 对每个 \(0\le k\le 2n\)，有
  \[
  \operatorname{mult}_{A_n}(2n-2k)=5\binom{2n}{k},
  \qquad
  \operatorname{mult}_{L_n}(2k)=5\binom{2n}{k}.
  \]
  \item 若记 \(\tau^\sharp(G_n)\) 为在每个连通分量内各取一棵生成树所得的组分保持生成森林数，则
  \[
  \tau^\sharp(G_n)
  =
  \Biggl(
  2^{2^{2n}-2n-1}
  \prod_{k=1}^{2n}k^{\binom{2n}{k}}
  \Biggr)^5.
  \]
  \item Ihara--Bass \(\zeta\) 函数满足
  \[
  \bigl(Z_{G_n}^{\mathrm{Ih}}(u)\bigr)^{-1}
  =
  \Biggl[
  (1-u^2)^{(2n-2)2^{2n-1}}
  \prod_{k=0}^{2n}
  \bigl(1-(2n-2k)u+(2n-1)u^2\bigr)^{\binom{2n}{k}}
  \Biggr]^5.
  \]
\end{enumerate}
\end{corollary}
\begin{proof}
第一条的邻接谱重数正是定理 \ref{thm:xi-time-part62d-leyang-branch-graph-gaussian-spectrum} 的第一条；再由
\[
L_n=2nI-A_n
\]
立即得到 Laplacian 谱重数公式。第二条即定理 \ref{thm:xi-time-part62d-leyang-branch-kirchhoff-forest}。第三条把定理 \ref{thm:xi-time-part62d-leyang-branch-ihara-zeta} 改写为逆函数形式便得。证毕。
\end{proof}

\begin{theorem}[Lee--Yang 分支塔逆极限的 affine Tate sheet 分解与层间正则四重覆盖]\label{thm:xi-time-part62dh-leyang-branch-affine-tate-regular-cover}
\leanverified{paper\_xi\_time\_part62dh\_leyang\_branch\_affine\_tate\_regular\_cover}
沿用推论 \ref{cor:xi-terminal-zm-leyang-finite-branch-regular-4ary-address} 的记号，并记
\[
\mathcal R_\infty^{(i)}:=\varprojlim_n\Pi_{i,n},
\qquad
\mathcal R_\infty:=\bigsqcup_{i=0}^{4}\mathcal R_\infty^{(i)}.
\]
则对每个 \(i\in\{0,\dots,4\}\) 与每个 \(n\ge 0\) 都有：
\begin{enumerate}
  \item 映射
  \[
  [2]:\Pi_{i,n+1}\longrightarrow \Pi_{i,n}
  \]
  是正则四重覆盖，并且
  \[
  \operatorname{Deck}\bigl([2]:\Pi_{i,n+1}\to \Pi_{i,n}\bigr)
  \cong
  E_{\min}[2]
  \cong
  (\ZZ/2\ZZ)^2.
  \]
  \item 一旦在第 \(i\) 个分支上选定相容提升 \(\mathbf P_i\in\mathcal R_\infty^{(i)}\)，便有
  \[
  \mathcal R_\infty^{(i)}\cong \mathbf P_i+T_2(E_{\min}),
  \qquad
  \mathcal R_\infty
  \cong
  \bigsqcup_{i=0}^{4}\bigl(\mathbf P_i+T_2(E_{\min})\bigr).
  \]
\end{enumerate}
特别地，Lee--Yang dyadic 分支塔的逆极限恰由五片彼此不交的 affine \(T_2(E_{\min})\)-sheet 组成。
\end{theorem}
\begin{proof}
推论 \ref{cor:xi-terminal-zm-leyang-finite-branch-regular-4ary-address} 的第一条已经给出
\[
[2]:\Pi_{i,n+1}\to\Pi_{i,n}
\]
为正则四重覆盖，且每个纤维在平移作用下都是 \(E_{\min}[2]\) 的主齐次空间。相应的 deck 变换正由核 \(E_{\min}[2]\) 的平移实现，因此 deck 群自然同构于 \(E_{\min}[2]\cong(\ZZ/2\ZZ)^2\)。这证明了第一条。

第二条则由定理 \ref{thm:xi-time-part62d-leyang-branch-inverse-limit-five-torsors} 与推论 \ref{cor:xi-terminal-zm-leyang-finite-branch-regular-4ary-address} 的第三条直接给出。证毕。
\end{proof}

\begin{corollary}[Zeckendorf--Gödel Dirichlet 级数主极点的自然密度残数]\label{cor:xi-time-part62dh-zg-dirichlet-density-residue}
\leanverified{paper\_xi\_time\_part62dh\_zg\_dirichlet\_density\_residue}
记
\[
D_{\mathrm{ZG}}(s):=\sum_{n\in\mathrm{ZG}}n^{-s}
\qquad (\Re(s)>1).
\]
则有
\[
\lim_{\sigma\downarrow 1}(\sigma-1)D_{\mathrm{ZG}}(\sigma)
=
\delta_{\mathrm{ZG}}=\delta(\mathrm{ZG}),
\qquad
\lim_{\sigma\downarrow 1}\frac{D_{\mathrm{ZG}}(\sigma)}{\zeta(\sigma)}
=
\delta_{\mathrm{ZG}}.
\]
因而 \(D_{\mathrm{ZG}}(s)\) 在 \(s=1\) 处具有留数 \(\delta(\mathrm{ZG})\) 的简单极点。
\end{corollary}
\begin{proof}
第一条正是定理 \ref{thm:xi-time-part62a-zg-simple-pole-density-residue}，第二条正是推论 \ref{cor:xi-time-part62a-zg-zeta-normalized-density-limit}。二者合并即得所述主极点残数公式。证毕。
\end{proof}

\noindent
由此，H.158--H.164 的 Lee--Yang 主线在结论层被再次压缩为一组彼此闭合的终端对象：有限层 branch 图的超八面体 wreath 对称、显式谱/森林/Ihara 数据与 dyadic 逆极限的五片 affine Tate sheet 彼此兼容，而 Zeckendorf--Gödel 像集的自然密度则同时外显为其 Dirichlet 级数在 \(s=1\) 的主极点残数。于是，超立方体地址、\(2\)-进仿射塔与 Euler 型解析常数便在同一审计框架内合流。

\paragraph{Lee--Yang 分支抬升、三域 Frobenius 张量独立与纤维同伦奇偶锁定}
在 Lee--Yang 分支三次数域的振幅平方刚性、三域 Chebotarev 直积结构以及纤维独立集复形的同伦分类已经确立之后，还可进一步抽取出下列十条彼此衔接的结构后果。一方面，有符号振幅相对于无符号分支数据给出一个不能在三次基域内消去的真正二次抬升；中间数条把三域 Frobenius 观测从张量独立进一步压缩为二次符号层与 \(L_4\) 共轭类层上的 Lee--Yang 分裂型刚性及其联合密度闭式；最后几条则把纤维多重度的奇偶、最大非可缩层的相对损失及其模 \(6\) 的 Fibonacci 拓扑税分别锁定为 \(\FF_2\)-同伦理指纹、四常数谱与显式缺口公式。

\begin{theorem}[Lee--Yang 有符号振幅在三次分支域上的非平方性与恰好六次代数度]\label{thm:xi-time-part62d-leyang-signed-amplitude-exact-degree6}
沿用定理 \ref{thm:xi-terminal-zm-kappa-square-cubic-field-s3}、\ref{thm:xi-terminal-zm-stokes-amplitude-square-cubic-rigidity} 与 \ref{thm:xi-terminal-zm-b2-kappa2-rational-interconstruction} 的记号，记
\[
K:=\QQ(u)=\QQ(\lambda_*)=\QQ(b),
\qquad
b:=B^2.
\]
则 \(K/\QQ\) 为三次扩张，并且
\[
b\notin K^{\times 2},
\qquad
[\QQ(B):K]=2,
\qquad
[\QQ(B):\QQ]=6.
\]
\end{theorem}
\begin{proof}
由定理 \ref{thm:xi-terminal-zm-b2-kappa2-rational-interconstruction}，
\[
K=\QQ(b)=\QQ(u)=\QQ(\lambda_*).
\]
又由定理 \ref{thm:xi-terminal-zm-stokes-amplitude-square-cubic-rigidity}，\(b\) 满足不可约三次方程
\[
162b^3+1593b^2+1998b+1184=0,
\]
故
\[
[K:\QQ]=[\QQ(b):\QQ]=3.
\]

把该多项式首一化为
\[
x^3+\frac{1593}{162}x^2+\frac{1998}{162}x+\frac{1184}{162},
\]
于是
\[
N_{K/\QQ}(b)=(-1)^3\frac{1184}{162}=-\frac{592}{81}<0.
\]
若 \(b=c^2\) 属于 \(K^{\times 2}\)，则
\[
N_{K/\QQ}(b)=N_{K/\QQ}(c)^2
\]
必须是 \(\QQ\) 中的平方，特别应为非负有理数，与上式矛盾。故
\[
b\notin K^{\times 2}.
\]

因此 \(X^2-b\in K[X]\) 不可约，从而
\[
[K(B):K]=2.
\]
又因 \(b=B^2\in \QQ(B)\)，故 \(K=\QQ(b)\subseteq \QQ(B)\)，于是
\[
K(B)=\QQ(B).
\]
由塔式公式得到
\[
[\QQ(B):\QQ]=[\QQ(B):K]\,[K:\QQ]=2\cdot 3=6.
\]
证毕。
\end{proof}

\begin{corollary}[Lee--Yang 分支数据的无符号/有符号代数寄存器度断裂]\label{cor:xi-time-part62d-leyang-unsigned-signed-register-gap}
在定理 \ref{thm:xi-time-part62d-leyang-signed-amplitude-exact-degree6} 的设定下，令 \(t\) 为任一代数元。则：
\begin{enumerate}
  \item 若
  \[
  u,\ \lambda_*,\ b\in \QQ(t),
  \]
  则
  \[
  [\QQ(t):\QQ]\ge 3;
  \]
  \item 若进一步
  \[
  B\in \QQ(t),
  \]
  则
  \[
  [\QQ(t):\QQ]\ge 6.
  \]
\end{enumerate}
因此，Lee--Yang 分支点的无符号数据 \((u,\lambda_*,b)\) 与有符号振幅 \(B\) 之间存在不可压缩的二倍代数复杂度断裂：前者的最小代数寄存器度为 \(3\)，后者的最小代数寄存器度为 \(6\)。
\end{corollary}
\begin{proof}
由定理 \ref{thm:xi-time-part62d-leyang-signed-amplitude-exact-degree6}，
\[
K=\QQ(u)=\QQ(\lambda_*)=\QQ(b),
\qquad
[K:\QQ]=3,
\qquad
[\QQ(B):\QQ]=6.
\]
若 \(u,\lambda_*,b\in \QQ(t)\)，则 \(K\subseteq \QQ(t)\)，故
\[
[\QQ(t):\QQ]\ge [K:\QQ]=3.
\]
若再有 \(B\in \QQ(t)\)，则 \(\QQ(B)\subseteq \QQ(t)\)，于是
\[
[\QQ(t):\QQ]\ge [\QQ(B):\QQ]=6.
\]
证毕。
\end{proof}

\begin{theorem}[三域 Frobenius 指纹的全部混合矩严格张量分解]\label{thm:xi-time-part62d-triple-frobenius-classfunction-tensor-independence}
沿用定理 \ref{thm:killo-leyang-lattes-triple-product-galois} 的记号，记
\[
L:=L_{10}L_{\mathrm{LY}}L_4,
\qquad
\Gal(L/\QQ)\cong S_{10}\times S_3\times (S_4\times C_2).
\]
对每个在 \(L/\QQ\) 中不分歧的素数 \(p\)，记
\[
\sigma_{10}(p):=\mathrm{Frob}_p|_{L_{10}},
\qquad
\sigma_3(p):=\mathrm{Frob}_p|_{L_{\mathrm{LY}}},
\qquad
\bigl(\sigma_4(p),\varepsilon_4(p)\bigr):=\mathrm{Frob}_p|_{L_4}\in S_4\times C_2.
\]
任取类函数
\[
f:S_{10}\to\CC,
\qquad
g:S_3\to\CC,
\qquad
h:S_4\times C_2\to\CC,
\]
并定义群平均
\[
\langle f\rangle_{S_{10}}:=\frac{1}{10!}\sum_{\tau\in S_{10}}f(\tau),
\qquad
\langle g\rangle_{S_3}:=\frac{1}{6}\sum_{\upsilon\in S_3}g(\upsilon),
\]
\[
\langle h\rangle_{S_4\times C_2}:=
\frac{1}{48}\sum_{\eta\in S_4\times C_2}h(\eta).
\]
若记
\[
\sum_{p\le x}^{\prime}
\]
为仅对在 \(L/\QQ\) 中不分歧的素数求和，则有
\[
\lim_{x\to\infty}\frac{1}{\pi(x)}
\sum_{p\le x}^{\prime}
f\!\bigl(\sigma_{10}(p)\bigr)\,
g\!\bigl(\sigma_3(p)\bigr)\,
h\!\bigl(\sigma_4(p),\varepsilon_4(p)\bigr)
=
\langle f\rangle_{S_{10}}\,
\langle g\rangle_{S_3}\,
\langle h\rangle_{S_4\times C_2}.
\]
特别地，若定义中心化类函数
\[
\widetilde f:=f-\langle f\rangle_{S_{10}},
\qquad
\widetilde g:=g-\langle g\rangle_{S_3},
\qquad
\widetilde h:=h-\langle h\rangle_{S_4\times C_2},
\]
则
\[
\lim_{x\to\infty}\frac{1}{\pi(x)}
\sum_{p\le x}^{\prime}
\widetilde f\!\bigl(\sigma_{10}(p)\bigr)\,
\widetilde g\!\bigl(\sigma_3(p)\bigr)\,
\widetilde h\!\bigl(\sigma_4(p),\varepsilon_4(p)\bigr)
=0.
\]
因而一切中心化混合矩均消失。更一般地，在相应的素数密度极限分布下，一切涉及至少两个不同因子的混合 cumulant 全部消失。
\end{theorem}
\begin{proof}
定理 \ref{thm:killo-leyang-lattes-triple-class-function-factorization} 已经给出同一极限公式，只是归一化分母写成
\[
\pi_{\mathrm{tr}}^*(x):=
\#\{p\le x:\ p\ \text{在}\ L/\QQ\ \text{中不分歧}\}.
\]
由于分歧素数只有有限多个，故
\[
\frac{\pi_{\mathrm{tr}}^*(x)}{\pi(x)}\longrightarrow 1
\qquad(x\to\infty),
\]
从而把分母 \(\pi_{\mathrm{tr}}^*(x)\) 替换为 \(\pi(x)\) 不改变极限值，第一条公式成立。

对中心化函数应用同一公式便得
\[
\langle \widetilde f\rangle_{S_{10}}=
\langle \widetilde g\rangle_{S_3}=
\langle \widetilde h\rangle_{S_4\times C_2}=0,
\]
故中心化三重乘积的极限平均为零。

进一步，上述极限恒等式表明三元 Frobenius 观测的极限分布正是
\[
\mu_{10}\otimes \mu_3\otimes \mu_4,
\]
其中 \(\mu_{10},\mu_3,\mu_4\) 分别是三个有限群上的均匀类分布。乘积分布的标准 cumulant 性质说明：只要某个 cumulant 同时涉及至少两个不同因子的观测量，它便必为零。证毕。
\end{proof}

\begin{theorem}[并入 \texorpdfstring{$L_4$}{L4} 后 Lee--Yang 三次分裂型在任意二次符号层上仍保持刚性独立]\label{thm:xi-time-part62d-chi10-l4-conditioned-leyang-splitting-rigidity}
沿用定理 \ref{thm:xi-time-part62d-triple-frobenius-classfunction-tensor-independence} 的记号，并记
\[
g(y):=256y^3+411y^2+165y+32.
\]
对不分歧素数 \(p\)，定义
\[
\chi_{10}(p):=\mathrm{sgn}\!\bigl(\sigma_{10}(p)\bigr)\in\{\pm1\}.
\]
设 \(C\subset S_4\times C_2\) 为任意共轭类，\(\varepsilon\in\{\pm1\}\)。则在条件事件
\[
\chi_{10}(p)=\varepsilon,
\qquad
\bigl(\sigma_4(p),\varepsilon_4(p)\bigr)\in C
\]
之下，\(g\bmod p\) 的三种分裂型密度仍保持不变：
\[
\boxed{\;
\delta\!\bigl((1)(2)\,\bigm|\,\chi_{10}(p)=\varepsilon,\ (\sigma_4(p),\varepsilon_4(p))\in C\bigr)=\frac12,\;}
\]
\[
\boxed{\;
\delta\!\bigl((1)(1)(1)\,\bigm|\,\chi_{10}(p)=\varepsilon,\ (\sigma_4(p),\varepsilon_4(p))\in C\bigr)=\frac16,\;}
\]
\[
\boxed{\;
\delta\!\bigl((3)\,\bigm|\,\chi_{10}(p)=\varepsilon,\ (\sigma_4(p),\varepsilon_4(p))\in C\bigr)=\frac13.\;}
\]
亦即：即便额外固定来自 \(L_{10}\) 的二次符号与来自 \(L_4\) 的任意共轭类信息，Lee--Yang 三次因子的模 \(p\) 分裂统计仍保持严格不偏置。
\end{theorem}
\begin{proof}
由定理 \ref{thm:xi-time-part62d-triple-frobenius-classfunction-tensor-independence}，
\[
\Gal(L_{10}L_{\mathrm{LY}}L_4/\QQ)\cong S_{10}\times S_3\times (S_4\times C_2)
\]
上的 Frobenius 极限分布是三个因子的乘积分布。因而 \(\chi_{10}(p)\) 只依赖 \(S_{10}\) 因子，事件 \((\sigma_4(p),\varepsilon_4(p))\in C\) 只依赖 \(S_4\times C_2\) 因子，而 \(g\bmod p\) 的分裂型只依赖 \(S_3\) 因子。

于是，在固定前两个条件之后，\(\sigma_3(p)\) 仍按 \(S_3\) 中共轭类大小给出的原始 Chebotarev 分布出现。再由结论 \ref{con:xi-terminal-zm-leyang-cubic-modp-splitting-density}，三种分裂型分别对应于 \(S_3\) 中的换位类、恒等类与三循环类，其类大小依次为 \(3,1,2\)。故条件密度仍为
\[
\frac{3}{6}=\frac12,\qquad \frac{1}{6},\qquad \frac{2}{6}=\frac13.
\]
证毕。
\end{proof}

\begin{corollary}[三域联合条件下的 Lee--Yang 分裂型具有完全显式的联合密度表]\label{cor:xi-time-part62d-chi10-l4-leyang-joint-density-table}
沿用定理 \ref{thm:xi-time-part62d-chi10-l4-conditioned-leyang-splitting-rigidity} 的记号。对任意 \(\varepsilon\in\{\pm1\}\) 与任意共轭类 \(C\subset S_4\times C_2\)，有
\[
\boxed{\;
\delta\!\bigl(\chi_{10}(p)=\varepsilon,\ (\sigma_4(p),\varepsilon_4(p))\in C,\ g\bmod p\text{ 型为 }(1)(2)\bigr)
=\frac{|C|}{192},\;}
\]
\[
\boxed{\;
\delta\!\bigl(\chi_{10}(p)=\varepsilon,\ (\sigma_4(p),\varepsilon_4(p))\in C,\ g\bmod p\text{ 型为 }(1)(1)(1)\bigr)
=\frac{|C|}{576},\;}
\]
\[
\boxed{\;
\delta\!\bigl(\chi_{10}(p)=\varepsilon,\ (\sigma_4(p),\varepsilon_4(p))\in C,\ g\bmod p\text{ 型为 }(3)\bigr)
=\frac{|C|}{288}.\;}
\]
特别地，
\[
\delta\!\bigl(\chi_{10}(p)=\varepsilon,\ (\sigma_4(p),\varepsilon_4(p))\in C\bigr)=\frac{|C|}{96}.
\]
\end{corollary}
\begin{proof}
由附录定理 \ref{thm:killo-leyang-quadratic-character-decoupling}，
\[
\delta\!\bigl(\chi_{10}(p)=\varepsilon\bigr)=\frac12.
\]
由定理 \ref{thm:xi-time-part62d-triple-frobenius-classfunction-tensor-independence}，
\[
\delta\!\bigl((\sigma_4(p),\varepsilon_4(p))\in C\bigr)=\frac{|C|}{48}.
\]
再由定理 \ref{thm:xi-time-part62d-chi10-l4-conditioned-leyang-splitting-rigidity}，在给定这两个条件后，三种 Lee--Yang 分裂型的条件密度分别为 \(1/2,1/6,1/3\)。三者相乘即得
\[
\frac12\cdot\frac{|C|}{48}\cdot\frac12=\frac{|C|}{192},
\qquad
\frac12\cdot\frac{|C|}{48}\cdot\frac16=\frac{|C|}{576},
\qquad
\frac12\cdot\frac{|C|}{48}\cdot\frac13=\frac{|C|}{288}.
\]
把三项相加即得最后一式。证毕。
\end{proof}

\begin{proposition}[纤维多重度奇偶的 \texorpdfstring{$\FF_2$}{F2}-同伦指纹锁定]\label{prop:xi-time-part62d-fiber-parity-f2-homotopy-fingerprint}
设 \(m\ge 1\)，\(x\in X_m\) 为稳定类型。若其候选位点诱导图分解为路径的不交并
\[
\Gamma_x\cong \bigsqcup_{i=1}^{r}P_{\ell_i},
\qquad
d_m(x)=\prod_{i=1}^{r}F_{\ell_i+2},
\]
并记
\[
K_x:=\Ind(\Gamma_x),
\qquad
\tau(x):=\sum_{i=1}^{r}\left\lfloor\frac{\ell_i+1}{3}\right\rfloor,
\]
则有同余
\[
d_m(x)\equiv
\sum_{j\ge -1}\dim_{\FF_2}\widetilde H_j(K_x;\FF_2)
\equiv
\widetilde\chi(K_x)
\pmod 2.
\]
特别地，若 \(K_x\simeq K_y\) 同伦等价，则
\[
d_m(x)\equiv d_m(y)\pmod 2.
\]
\end{proposition}
\begin{proof}
由定理 \ref{thm:killo-fold-fiber-homotopy-parity-equivalence}，\(|K_x|\) 可缩当且仅当 \(d_m(x)\) 为偶数，而 \(|K_x|\simeq S^{\tau(x)-1}\) 当且仅当 \(d_m(x)\) 为奇数。

若 \(d_m(x)\) 为偶数，则 \(|K_x|\) 可缩，故所有约化同调群均为零，从而
\[
\sum_{j\ge -1}\dim_{\FF_2}\widetilde H_j(K_x;\FF_2)=0
\equiv d_m(x)\pmod 2.
\]
同时 \(\widetilde\chi(K_x)=0\)，故第二个同余亦成立。

若 \(d_m(x)\) 为奇数，则 \(|K_x|\simeq S^{\tau(x)-1}\)。由推论 \ref{cor:pom-fiber-independence-complex-homology}，其约化同调仅在唯一维度上非零且秩为 \(1\)。因此
\[
\sum_{j\ge -1}\dim_{\FF_2}\widetilde H_j(K_x;\FF_2)=1
\equiv d_m(x)\pmod 2.
\]
又
\[
\widetilde\chi(K_x)=\widetilde\chi\!\bigl(S^{\tau(x)-1}\bigr)=(-1)^{\tau(x)-1}\equiv 1\pmod 2,
\]
故第二个同余也成立。

最后，\(\FF_2\)-约化 Betti 数总和与约化 Euler 特征都是同伦不变量，因此若 \(K_x\simeq K_y\)，则二者对应相等，进而
\[
d_m(x)\equiv d_m(y)\pmod 2.
\]
证毕。
\end{proof}

\begin{theorem}[最大非可缩纤维相对缺口的模 \texorpdfstring{$6$}{6} 四常数谱]\label{thm:xi-time-part62d-max-noncontractible-fiber-gap-four-constants}
定义
\[
D_m:=\max_{x\in X_m}d_m(x),
\qquad
\widetilde D_m:=\max\{d_m(x):x\in X_m,\ |K_x|\not\simeq *\},
\qquad
\delta_m:=1-\frac{\widetilde D_m}{D_m}.
\]
则沿模 \(6\) 的剩余类有极限
\[
\lim_{\substack{m\to\infty\\ m\equiv 0,4\!\!\!\!\pmod 6}}\delta_m=0,
\qquad
\lim_{\substack{m\to\infty\\ m\equiv 2\!\!\!\!\pmod 6}}\delta_m=\varphi^{-5},
\]
\[
\lim_{\substack{m\to\infty\\ m\equiv 1,5\!\!\!\!\pmod 6}}\delta_m=\frac{1}{2}\varphi^{-5},
\qquad
\lim_{\substack{m\to\infty\\ m\equiv 3\!\!\!\!\pmod 6}}\delta_m=\frac{1}{2}\bigl(\varphi^{-5}+\varphi^{-9}\bigr).
\]
因而
\[
\liminf_{m\to\infty}\delta_m=0,
\qquad
\limsup_{m\to\infty}\delta_m=\varphi^{-5}.
\]
\end{theorem}
\begin{proof}
由定理 \ref{thm:pom-max-noncontractible-fiber-mod6-phase}，
\[
\widetilde D_m=
\begin{cases}
D_m, & m\equiv 0,4\pmod 6,\\[2pt]
D_m^{(2)}, & m\equiv 1,5\pmod 6,\\[2pt]
D_m^{(3)}, & m\equiv 2,3\pmod 6.
\end{cases}
\]
再由定理 \ref{thm:pom-max-fiber}、定理 \ref{thm:pom-second-max-fiber-closed-form}、定理 \ref{thm:pom-third-max-fiber-even-closed-form} 与注 \ref{rem:pom-third-max-fiber-odd-closed-form}，
\[
D_{2k}=F_{k+2},
\qquad
D_{2k+1}=2F_{k+1},
\]
\[
D_{2k+1}^{(2)}=2F_{k+1}-F_{k-4},
\qquad
D_{2k}^{(3)}=F_{k+2}-F_{k-3},
\]
\[
D_{2k+1}^{(3)}=2F_{k+1}-F_{k-4}-F_{k-8}.
\]

于是分四类计算。

若 \(m\equiv 0,4\pmod 6\)，则 \(\widetilde D_m=D_m\)，故 \(\delta_m=0\)。

若 \(m=2k\equiv 2\pmod 6\)，则
\[
\delta_{2k}=1-\frac{F_{k+2}-F_{k-3}}{F_{k+2}}
=\frac{F_{k-3}}{F_{k+2}}
\longrightarrow \varphi^{-5}.
\]

若 \(m=2k+1\equiv 1,5\pmod 6\)，则
\[
\delta_{2k+1}
=1-\frac{2F_{k+1}-F_{k-4}}{2F_{k+1}}
=\frac{F_{k-4}}{2F_{k+1}}
\longrightarrow \frac{1}{2}\varphi^{-5}.
\]

若 \(m=2k+1\equiv 3\pmod 6\)，则
\[
\delta_{2k+1}
=1-\frac{2F_{k+1}-F_{k-4}-F_{k-8}}{2F_{k+1}}
=\frac{F_{k-4}+F_{k-8}}{2F_{k+1}}
\longrightarrow \frac{1}{2}\bigl(\varphi^{-5}+\varphi^{-9}\bigr).
\]
这里用到了 Fibonacci 比值极限
\[
\frac{F_{n-r}}{F_n}\longrightarrow \varphi^{-r}.
\]
四个剩余类的极限于是全部得到，故下极限与上极限分别为
\[
0
\qquad\text{和}\qquad
\varphi^{-5}.
\]
证毕。
\end{proof}

\begin{corollary}[最大非可缩纤维的拓扑税具有严格的模 \texorpdfstring{$6$}{6} 相图与极限常数]\label{cor:xi-time-part62d-max-noncontractible-fiber-tax-mod6}
沿用定理 \ref{thm:xi-time-part62d-max-noncontractible-fiber-gap-four-constants} 的记号，并对 \(m\ge 17\) 定义
\[
\Delta_m:=D_m-\widetilde D_m.
\]
则有精确分段公式
\[
\boxed{
\Delta_m=
\begin{cases}
0,& m\equiv 0,4\pmod6,\\[1mm]
F_{(m-9)/2},& m\equiv 1,5\pmod6,\\[1mm]
F_{m/2-3},& m\equiv 2\pmod6,\\[1mm]
F_{(m-9)/2}+F_{(m-17)/2},& m\equiv 3\pmod6.
\end{cases}}
\]
并且相应的非可缩最大纤维相对税率极限满足
\[
\lim_{\substack{m\to\infty\\ m\equiv 0,4\!\!\!\!\pmod6}}\frac{\widetilde D_m}{D_m}=1,
\]
\[
\lim_{\substack{m\to\infty\\ m\equiv 1,5\!\!\!\!\pmod6}}\frac{\widetilde D_m}{D_m}
=1-\frac12\varphi^{-5},
\]
\[
\lim_{\substack{m\to\infty\\ m\equiv 2\!\!\!\!\pmod6}}\frac{\widetilde D_m}{D_m}
=1-\varphi^{-5},
\]
\[
\lim_{\substack{m\to\infty\\ m\equiv 3\!\!\!\!\pmod6}}\frac{\widetilde D_m}{D_m}
=1-\frac12\bigl(\varphi^{-5}+\varphi^{-9}\bigr).
\]
因此，强制最大纤维从可缩层跃迁到非可缩层所支付的代价并非抽象误差，而是一条由模 \(6\) 剩余类完全控制的 Fibonacci 拓扑税相图。
\end{corollary}
\begin{proof}
由定理 \ref{thm:pom-max-noncontractible-fiber-mod6-phase}，
\[
\widetilde D_m=
\begin{cases}
D_m,& m\equiv 0,4\pmod 6,\\[2pt]
D_m^{(2)},& m\equiv 1,5\pmod 6,\\[2pt]
D_m^{(3)},& m\equiv 2,3\pmod 6.
\end{cases}
\]
再由定理 \ref{thm:pom-max-fiber}、定理 \ref{thm:pom-second-max-fiber-closed-form}、定理 \ref{thm:pom-third-max-fiber-even-closed-form} 与注 \ref{rem:pom-third-max-fiber-odd-closed-form}，
\[
D_{2k}=F_{k+2},\qquad D_{2k+1}=2F_{k+1},
\]
\[
D_{2k+1}^{(2)}=2F_{k+1}-F_{k-4},\qquad
D_{2k}^{(3)}=F_{k+2}-F_{k-3},
\]
\[
D_{2k+1}^{(3)}=2F_{k+1}-F_{k-4}-F_{k-8}.
\]

若 \(m\equiv 0,4\pmod 6\)，则 \(\widetilde D_m=D_m\)，故 \(\Delta_m=0\)。

若 \(m=2k+1\equiv 1,5\pmod 6\)，则
\[
\Delta_m=D_{2k+1}-D_{2k+1}^{(2)}=F_{k-4}=F_{(m-9)/2}.
\]

若 \(m=2k\equiv 2\pmod 6\)，则
\[
\Delta_m=D_{2k}-D_{2k}^{(3)}=F_{k-3}=F_{m/2-3}.
\]

若 \(m=2k+1\equiv 3\pmod 6\)，则
\[
\Delta_m=D_{2k+1}-D_{2k+1}^{(3)}=F_{k-4}+F_{k-8}
=F_{(m-9)/2}+F_{(m-17)/2}.
\]
这就得到精确拓扑税公式。

对比值极限，只需把
\[
\frac{\widetilde D_m}{D_m}=1-\frac{\Delta_m}{D_m}
\]
与 Fibonacci 比值极限
\[
\frac{F_{n-r}}{F_n}\longrightarrow \varphi^{-r}
\]
结合即可。例如当 \(m=2k+1\equiv 1,5\pmod 6\) 时，
\[
\frac{\widetilde D_m}{D_m}
=
1-\frac{F_{k-4}}{2F_{k+1}}
\longrightarrow
1-\frac12\varphi^{-5}.
\]
其余三类完全同理。证毕。
\end{proof}

\begin{proposition}[三次分支域上的平方根抬升构成非平凡 \texorpdfstring{$\mu_2$}{mu2}-torsor]\label{prop:xi-time-part62d-leyang-square-root-nontrivial-torsor}
沿用定理 \ref{thm:xi-time-part62d-leyang-signed-amplitude-exact-degree6} 的记号。则有限 \'etale \(K\)-概形
\[
Y_b:=\operatorname{Spec}K[X]/(X^2-b)
\longrightarrow
\operatorname{Spec}K
\]
是一个非平凡 \(\mu_2\)-torsor。等价地，在 \(\overline K\) 上把其几何纤维识别为根集
\[
\{+B,-B\},
\]
则该二点集在符号作用
\[
B\longmapsto -B
\]
下构成一个真正的 \(\ZZ/2\ZZ\)-torsor，并且不存在 \(K\)-有理截面把 \(B\) 从 \(b=B^2\) 中单值恢复出来。
\end{proposition}
\begin{proof}
由定理 \ref{thm:xi-time-part62d-leyang-signed-amplitude-exact-degree6}，
\[
b\notin K^{\times 2}.
\]
故多项式
\[
X^2-b
\]
在 \(K[X]\) 中不可约，从而 \(Y_b(K)=\varnothing\)。这说明上述二次覆盖在 \(K\) 上没有截面，因而不是平凡 torsor。

另一方面，在 \(\overline K\) 上有分解
\[
X^2-b=(X-B)(X+B),
\]
故 \(Y_b(\overline K)=\{+B,-B\}\)。群 \(\mu_2=\{\pm 1\}\) 通过乘法
\[
(\pm 1)\cdot B:=\pm B
\]
在该二点集上自由且传递地作用，因此它确为一个 \(\mu_2\)-torsor；在忘却群结构而仅保留交换元的意义下，也即一个非平凡 \(\ZZ/2\ZZ\)-torsor。不存在 \(K\)-有理截面正是 \(Y_b(K)=\varnothing\) 的等价表述。证毕。
\end{proof}

\paragraph{纤维独立集复形的约化 Euler--Witten 指数与奇偶指纹}

\begin{theorem}[纤维独立集复形的约化 Euler--Witten 指数完全等价于纤维奇偶]\label{thm:xi-time-part62dg-fiber-reduced-euler-witten-parity}
\leanverified{paper\_xi\_time\_part62dg\_fiber\_reduced\_euler\_witten\_parity}
设 \(m\ge 1\) 且 \(x\in X_m\)。沿用定理 \ref{thm:killo-fold-fiber-homotopy-parity-equivalence} 的记号，记
\[
\Gamma_x\cong\bigsqcup_{i=1}^{r}P_{\ell_i},
\qquad
d_m(x)=\prod_{i=1}^{r}F_{\ell_i+2},
\]
\[
K_x:=\operatorname{Ind}^{\Delta}(\Gamma_x),
\qquad
\tau(x):=\sum_{i=1}^{r}\left\lfloor\frac{\ell_i+1}{3}\right\rfloor.
\]
定义约化 Euler--Witten 指数
\[
\mathfrak W(x):=\widetilde\chi(K_x).
\]
则有严格公式
\[
\mathfrak W(x)=
\begin{cases}
0,& d_m(x)\ \text{为偶},\\[1mm]
(-1)^{\tau(x)-1},& d_m(x)\ \text{为奇}.
\end{cases}
\]
特别地，
\[
|\mathfrak W(x)|=\mathbf 1_{\{d_m(x)\ \mathrm{odd}\}},
\qquad
\mathfrak W(x)\neq 0\iff |K_x|\not\simeq *.
\]
\end{theorem}
\begin{proof}
由定理 \ref{thm:killo-fold-fiber-homotopy-parity-equivalence}，
\[
|K_x|\ \text{可缩}
\quad\Longleftrightarrow\quad
d_m(x)\ \text{为偶},
\]
并且
\[
|K_x|\simeq S^{\tau(x)-1}
\quad\Longleftrightarrow\quad
d_m(x)\ \text{为奇}.
\]

若 \(d_m(x)\) 为偶，则 \(K_x\) 可缩，从而
\[
\widetilde\chi(K_x)=0.
\]
若 \(d_m(x)\) 为奇，则
\[
|K_x|\simeq S^{\tau(x)-1},
\]
故
\[
\widetilde\chi(K_x)=\widetilde\chi\!\bigl(S^{\tau(x)-1}\bigr)=(-1)^{\tau(x)-1}.
\]
这就得到所示分段公式。

进一步，第一条特别推出
\[
|\mathfrak W(x)|=
\begin{cases}
0,& d_m(x)\ \text{为偶},\\
1,& d_m(x)\ \text{为奇},
\end{cases}
\]
即
\[
|\mathfrak W(x)|=\mathbf 1_{\{d_m(x)\ \mathrm{odd}\}}.
\]
而 \(\mathfrak W(x)\neq 0\) 与 \(d_m(x)\) 为奇等价，后者又与 \(|K_x|\not\simeq *\) 等价，故最后一条亦成立。证毕。
\end{proof}

\noindent
上述十条结果把 Lee--Yang 分支三次数域链、三域 Frobenius 直积统计与 Fold 纤维同伦相图压缩为同一条刚性主线：无符号分支数据与其有符号抬升之间存在严格的三次/六次代数度断裂，并在三次基域上表现为不可平凡化的二值 torsor；与此同时，三域 Frobenius 观测不仅在素数密度意义下完全张量分解，而且在固定 \(L_{10}\) 二次符号层与 \(L_4\) 共轭类层之后，Lee--Yang 三次分裂型仍保持严格不偏置，并给出完整的联合密度表；最终，纤维多重度的奇偶、最大非可缩顶层损失及其模 \(6\) 的 Fibonacci 拓扑税则分别由 \(\FF_2\)-同伦指纹、四个黄金常数与显式缺口公式完全锁定。于是，代数分支抬升、Chebotarev 统计独立性与纤维同伦相图便在同一二值刚性层上汇合。

\paragraph{三域类函数正交与最大非可缩纤维的六步重整化}
在三域 Frobenius 类函数统计的完全张量分解、\(L_{10}\) 与 \(L_4\) 条件下的 Lee--Yang 分裂型刚性，以及最大非可缩纤维的模 \(6\) 相图都已经固定之后，还可把这两条链进一步压缩为下列三条高层结论。

\begin{theorem}[三域直积下的类函数 Chebotarev 平均完全因子化]\label{thm:xi-time-part62di-triple-classfunction-chebotarev-factorization}
\leanverified{paper\_xi\_time\_part62di\_triple\_classfunction\_chebotarev\_factorization}
沿用定理 \ref{thm:xi-time-part62d-triple-frobenius-classfunction-tensor-independence} 的记号，记
\[
L:=L_{10}L_{\mathrm{LY}}L_4,
\qquad
\Gal(L/\QQ)\cong S_{10}\times S_3\times (S_4\times C_2).
\]
对每个在 \(L/\QQ\) 中不分歧的素数 \(p\)，记
\[
\sigma_{10}(p)\in S_{10},
\qquad
\sigma_3(p)\in S_3,
\qquad
\tau(p):=\bigl(\sigma_4(p),\varepsilon_4(p)\bigr)\in S_4\times C_2.
\]
对任意类函数
\[
f\in \mathrm{Cl}(S_{10}),
\qquad
g\in \mathrm{Cl}(S_3),
\qquad
h\in \mathrm{Cl}(S_4\times C_2),
\]
定义群平均
\[
\langle f\rangle_{S_{10}}:=\frac{1}{10!}\sum_{\gamma\in S_{10}}f(\gamma),
\qquad
\langle g\rangle_{S_3}:=\frac{1}{6}\sum_{\delta\in S_3}g(\delta),
\]
\[
\langle h\rangle_{S_4\times C_2}:=
\frac{1}{48}\sum_{\eta\in S_4\times C_2}h(\eta).
\]
则有极限恒等式
\[
\lim_{x\to\infty}\frac{1}{\pi(x)}
\sum_{p\le x}^{\prime}
f\!\bigl(\sigma_{10}(p)\bigr)\,
g\!\bigl(\sigma_3(p)\bigr)\,
h\!\bigl(\tau(p)\bigr)
=
\langle f\rangle_{S_{10}}\,
\langle g\rangle_{S_3}\,
\langle h\rangle_{S_4\times C_2},
\]
其中 \(\sum_{p\le x}^{\prime}\) 表示仅对在 \(L/\QQ\) 中不分歧的素数求和。
\end{theorem}
\begin{proof}
这正是定理 \ref{thm:xi-time-part62d-triple-frobenius-classfunction-tensor-independence} 的第一条公式，只是把
\[
\bigl(\sigma_4(p),\varepsilon_4(p)\bigr)
\]
重新记作 \(\tau(p)\)。由于 \(L/\QQ\) 的分歧素数只有有限多个，以 \(\pi(x)\) 归一化与以不分歧素数计数归一化给出同一极限值，故所示公式成立。证毕。
\end{proof}

\begin{corollary}[\texorpdfstring{$S_3$}{S3} 分裂型与外侧中心化观测的严格正交]\label{cor:xi-time-part62di-s3-external-centered-orthogonality}
\leanverified{paper\_xi\_time\_part62di\_s3\_external\_centered\_orthogonality}
在定理 \ref{thm:xi-time-part62di-triple-classfunction-chebotarev-factorization} 的设定下，记
\[
G_{\mathrm{ext}}:=S_{10}\times (S_4\times C_2).
\]
设
\[
g\in \mathrm{Cl}(S_3),
\qquad
h\in \mathrm{Cl}(G_{\mathrm{ext}})
\]
满足零均值条件
\[
\langle g\rangle_{S_3}=0,
\qquad
\langle h\rangle_{G_{\mathrm{ext}}}=0.
\]
则
\[
\lim_{x\to\infty}\frac{1}{\pi(x)}
\sum_{p\le x}^{\prime}
g\!\bigl(\sigma_3(p)\bigr)\,
h\!\bigl(\sigma_{10}(p),\tau(p)\bigr)
=0.
\]
特别地，对 Lee--Yang 三次多项式
\[
p_{\mathrm{LY}}(y):=256y^3+411y^2+165y+32
\]
在模 \(p\) 下的任意中心化分裂型指示函数，都与来自 \(L_{10}\) 与 \(L_4\) 两侧的任意中心化类函数观测严格正交。
\end{corollary}
\begin{proof}
定义三域直积群上的类函数
\[
\Phi:S_{10}\times S_3\times (S_4\times C_2)\longrightarrow \CC,
\qquad
\Phi(\gamma,\delta,\eta):=g(\delta)\,h(\gamma,\eta).
\]
由 \(G_{\mathrm{ext}}\times S_3\) 的直积结构与 Chebotarev 定理，
\[
\lim_{x\to\infty}\frac{1}{\pi(x)}
\sum_{p\le x}^{\prime}
g\!\bigl(\sigma_3(p)\bigr)\,
h\!\bigl(\sigma_{10}(p),\tau(p)\bigr)
=
\frac{1}{|G_{\mathrm{ext}}|\cdot 6}
\sum_{\gamma\in S_{10}}
\sum_{\delta\in S_3}
\sum_{\eta\in S_4\times C_2}
\Phi(\gamma,\delta,\eta).
\]
而右端按直积求和分解为
\[
\left(\frac{1}{6}\sum_{\delta\in S_3}g(\delta)\right)
\left(
\frac{1}{|G_{\mathrm{ext}}|}
\sum_{(\gamma,\eta)\in G_{\mathrm{ext}}}
h(\gamma,\eta)
\right)
=
\langle g\rangle_{S_3}\,\langle h\rangle_{G_{\mathrm{ext}}},
\]
在零均值假设下即为 \(0\)。最后一句只需把 \(g\) 取为三种 Lee--Yang 分裂型的中心化指示函数即可。证毕。
\end{proof}

\begin{theorem}[最大非可缩纤维沿每个模 \texorpdfstring{$6$}{6} 相位都满足同一六步重整化递推]\label{thm:xi-time-part62di-max-noncontractible-fiber-sixstep-renormalization}
\leanverified{paper\_xi\_time\_part62di\_max\_noncontractible\_fiber\_sixstep\_renormalization}
对每个同余类 \(\ell\in\{0,1,2,3,4,5\}\)，在 \(6r+\ell\ge 17\) 的范围内定义
\[
a_r^{(\ell)}:=\widetilde D_{6r+\ell},
\]
其中 \(\widetilde D_m\) 为最大非可缩纤维。则对每个 \(\ell\) 都存在 \(r_0(\ell)\)，使得对一切 \(r\ge r_0(\ell)\) 成立统一递推
\[
a_{r+2}^{(\ell)}=4a_{r+1}^{(\ell)}+a_r^{(\ell)}.
\]
因此：
\begin{enumerate}
  \item 若定义
        \[
        A_\ell(z):=\sum_{r\ge r_0(\ell)}a_r^{(\ell)}z^r,
        \]
        则每个 \(A_\ell(z)\) 都是有理函数，且分母统一为
        \[
        1-4z-z^2.
        \]
  \item 对每个相位 \(\ell\)，六步增长率都存在且相同：
        \[
        \lim_{r\to\infty}\frac{a_{r+1}^{(\ell)}}{a_r^{(\ell)}}=\varphi^3.
        \]
\end{enumerate}
\end{theorem}
\begin{proof}
由定理 \ref{thm:pom-max-noncontractible-fiber-mod6-phase}，\(\widetilde D_m\) 在六个剩余类上分别由顶层谱的前三层之一给出。再结合定理 \ref{thm:pom-max-fiber}、定理 \ref{thm:pom-second-max-fiber-closed-form}、定理 \ref{thm:pom-third-max-fiber-even-closed-form} 与注 \ref{rem:pom-third-max-fiber-odd-closed-form}，对充分大的 \(r\) 有显式公式
\[
a_r^{(0)}=F_{3r+2},
\qquad
a_r^{(4)}=F_{3r+4},
\]
\[
a_r^{(1)}=2F_{3r+1}-F_{3r-4},
\qquad
a_r^{(5)}=2F_{3r+3}-F_{3r-2},
\]
\[
a_r^{(2)}=F_{3r+3}-F_{3r-2},
\qquad
a_r^{(3)}=2F_{3r+2}-F_{3r-3}-F_{3r-7}.
\]
因此每条相位序列都是若干 Fibonacci 三步子序列 \(F_{3r+c}\) 的有理线性组合。

现固定整数 \(c\)，并记
\[
u_r^{(c)}:=F_{3r+c}.
\]
由 Binet 公式
\[
F_n=\frac{\varphi^n-\psi^n}{\sqrt5},
\qquad
\psi=-\varphi^{-1},
\]
得到
\[
u_r^{(c)}
=
\frac{\varphi^{3r+c}-\psi^{3r+c}}{\sqrt5}
=
\alpha_c(\varphi^3)^r+\beta_c(-\varphi^{-3})^r
\]
对某些常数 \(\alpha_c,\beta_c\in\QQ(\sqrt5)\) 成立。而 \(\varphi^3\) 与 \(-\varphi^{-3}\) 正是二次方程
\[
\lambda^2-4\lambda-1=0
\]
的两根，故每个 \(u_r^{(c)}\) 都满足递推
\[
u_{r+2}^{(c)}=4u_{r+1}^{(c)}+u_r^{(c)}.
\]
线性组合保持同一递推，于是每个 \(a_r^{(\ell)}\) 也满足
\[
a_{r+2}^{(\ell)}=4a_{r+1}^{(\ell)}+a_r^{(\ell)}
\]
对充分大的 \(r\) 成立。

这立即推出生成函数 \(A_\ell(z)\) 满足
\[
(1-4z-z^2)A_\ell(z)
\]
为某个多项式，从而 \(A_\ell(z)\) 是分母统一为 \(1-4z-z^2\) 的有理函数。

最后，由上述显式公式，每个 \(a_r^{(\ell)}\) 都具有形如
\[
a_r^{(\ell)}=c_\ell\,\varphi^{3r}+O(\varphi^{-3r})
\qquad
(c_\ell>0)
\]
的展开，因此
\[
\frac{a_{r+1}^{(\ell)}}{a_r^{(\ell)}}\longrightarrow \varphi^3.
\]
证毕。
\end{proof}

\noindent
因此，三域直积结构在结论层不仅迫使任意类函数素数平均完全乘积化，而且进一步把 \(S_3\) 分裂型与一切外侧中心化观测严格解耦；与此同时，最大非可缩纤维虽然在模 \(6\) 上呈现分段相图，但一旦按步长 \(6\) 采样，六个相位又统一塌缩到同一个二阶重整化动力学
\[
\lambda^2-4\lambda-1=0.
\]

\paragraph{Chebotarev 直积的符号立方体、单位控制与全分裂稀释}
在三域 Frobenius 类函数观测的完全张量独立与 Lee--Yang 分裂型的联合密度表已经建立之后，还可把二次特征层、\(S_3\) 分解型的信息流向以及 Latt\`es 六次因子的全分裂作用进一步压缩为下列三条结构后果。

\begin{theorem}[三重直积域的最大初等 \texorpdfstring{$2$}{2}-阿贝尔商恰为 \texorpdfstring{$C_2^4$}{C2^4}，从而恰有十五个非平凡二次子域]\label{thm:xi-time-part62e-triple-product-abelianization-c2-four}
\leanverified{paper\_xi\_time\_part62e\_triple\_product\_abelianization\_c2\_four}
沿用定理 \ref{thm:killo-leyang-lattes-triple-product-galois} 的记号，记
\[
G:=\Gal(L_{10}L_{\mathrm{LY}}L_4/\QQ)
\cong
S_{10}\times S_3\times (S_4\times C_2).
\]
则其阿贝尔化满足
\[
\boxed{\;
G^{\mathrm{ab}}
\cong
S_{10}^{\mathrm{ab}}\times S_3^{\mathrm{ab}}\times (S_4\times C_2)^{\mathrm{ab}}
\cong C_2^4.\;}
\]
特别地，\(L_{10}L_{\mathrm{LY}}L_4\) 中恰有
\[
\boxed{\;2^4-1=15\;}
\]
个非平凡二次子域。

更具体地，固定一个同构 \(\iota:C_2\xrightarrow{\sim}\{\pm1\}\) 后，可取四个基本二次特征为
\[
\chi_{10}:=\operatorname{sgn}\circ \mathrm{pr}_{S_{10}},
\qquad
\chi_{\mathrm{LY}}:=\operatorname{sgn}\circ \mathrm{pr}_{S_3},
\]
\[
\chi_{4,\mathrm{sgn}}:=\operatorname{sgn}\circ \mathrm{pr}_{S_4},
\qquad
\chi_{4,\mathrm{aux}}:=\iota\circ \mathrm{pr}_{C_2},
\]
它们生成全部二次特征群
\[
\operatorname{Hom}(G,\{\pm1\})\cong C_2^4.
\]
\end{theorem}
\begin{proof}
由定理 \ref{thm:killo-leyang-lattes-triple-product-galois}，
\[
G\cong S_{10}\times S_3\times (S_4\times C_2).
\]
而对一切 \(n\ge 2\) 都有
\[
S_n^{\mathrm{ab}}\cong C_2,
\]
并且
\[
(S_4\times C_2)^{\mathrm{ab}}
\cong
S_4^{\mathrm{ab}}\times C_2
\cong
C_2\times C_2.
\]
因此
\[
G^{\mathrm{ab}}
\cong
C_2\times C_2\times (C_2\times C_2)
\cong
C_2^4.
\]

由附录定理 \ref{thm:killo-leyang-quadratic-character-decoupling}，前两条特征 \(\chi_{10}\) 与 \(\chi_{\mathrm{LY}}\) 分别对应 \(L_{10}\) 与 \(L_{\mathrm{LY}}\) 的唯一非平凡二次正规子域；后两条特征 \(\chi_{4,\mathrm{sgn}},\chi_{4,\mathrm{aux}}\) 则给出第三因子 \((S_4\times C_2)^{\mathrm{ab}}\cong C_2^2\) 的两条基方向。故四者生成全部
\[
\operatorname{Hom}(G,\{\pm1\})\cong C_2^4.
\]

最后，非平凡二次子域与非零二次特征等价，也与指数 \(2\) 子群等价。由于 \(C_2^4\) 上非零线性泛函恰有 \(2^4-1=15\) 个，故相应的非平凡二次子域也恰有 \(15\) 个。证毕。
\end{proof}

\begin{theorem}[在完整的 \texorpdfstring{$C_2^4$}{C2^4} 符号立方体中，\texorpdfstring{$\chi_{\mathrm{LY}}$}{chi_LY} 是 \texorpdfstring{$g\bmod p$}{g mod p} 分解型的唯一有效位]\label{thm:xi-time-part62e-chi-ly-sole-controller}
沿用定理 \ref{thm:xi-time-part62d-triple-frobenius-classfunction-tensor-independence} 与
\ref{thm:xi-time-part62e-triple-product-abelianization-c2-four} 的记号。对不分歧素数 \(p\)，定义
\[
\chi_{10}(p):=\operatorname{sgn}\!\bigl(\sigma_{10}(p)\bigr),
\qquad
\chi_{\mathrm{LY}}(p):=\operatorname{sgn}\!\bigl(\sigma_3(p)\bigr),
\]
\[
\chi_{4,\mathrm{sgn}}(p):=\operatorname{sgn}\!\bigl(\sigma_4(p)\bigr),
\qquad
\chi_{4,\mathrm{aux}}(p):=\iota\!\bigl(\varepsilon_4(p)\bigr)\in\{\pm1\}.
\]
则成立：
\begin{enumerate}
  \item
  \[
  \chi_{\mathrm{LY}}(p)=-1
  \iff
  g\bmod p\ \text{的分解型为 }(1)(2).
  \]
  \item
  \[
  \chi_{\mathrm{LY}}(p)=+1
  \iff
  g\bmod p\ \text{的分解型属于 }\{(1)(1)(1),(3)\},
  \]
  并且
  \[
  \delta\!\bigl((1)(1)(1)\mid \chi_{\mathrm{LY}}=+1\bigr)=\frac13,
  \qquad
  \delta\!\bigl((3)\mid \chi_{\mathrm{LY}}=+1\bigr)=\frac23.
  \]
  \item 对任意固定的
  \[
  (\epsilon_{10},\epsilon_{4,\mathrm{sgn}},\epsilon_{4,\mathrm{aux}})\in\{\pm1\}^3,
  \]
  上述关于 \(g\bmod p\) 的条件分布保持不变。等价地，
  \[
  \delta\!\bigl((1)(2)\mid \chi_{\mathrm{LY}}=-1,\ \chi_{10}=\epsilon_{10},\ \chi_{4,\mathrm{sgn}}=\epsilon_{4,\mathrm{sgn}},\ \chi_{4,\mathrm{aux}}=\epsilon_{4,\mathrm{aux}}\bigr)=1,
  \]
  \[
  \delta\!\bigl((1)(1)(1)\mid \chi_{\mathrm{LY}}=+1,\ \chi_{10}=\epsilon_{10},\ \chi_{4,\mathrm{sgn}}=\epsilon_{4,\mathrm{sgn}},\ \chi_{4,\mathrm{aux}}=\epsilon_{4,\mathrm{aux}}\bigr)=\frac13,
  \]
  \[
  \delta\!\bigl((3)\mid \chi_{\mathrm{LY}}=+1,\ \chi_{10}=\epsilon_{10},\ \chi_{4,\mathrm{sgn}}=\epsilon_{4,\mathrm{sgn}},\ \chi_{4,\mathrm{aux}}=\epsilon_{4,\mathrm{aux}}\bigr)=\frac23.
  \]
\end{enumerate}
因此，全部 Frobenius 数据在 \(g\)-分解型意义上形成一个严格的 \(16\times 3\) 微胞分解，其中只有 \(\chi_{\mathrm{LY}}\) 真正携带 \(S_3\) 分解型信息，而其余三位只实施均匀细分。
\end{theorem}
\begin{proof}
由结论 \ref{con:xi-terminal-zm-leyang-cubic-modp-splitting-density}，\(g\bmod p\) 的三种分解型与 \(S_3\) 的三类共轭类一一对应：恒等类对应 \((1)(1)(1)\)，换位类对应 \((1)(2)\)，三循环类对应 \((3)\)。

另一方面，\(\chi_{\mathrm{LY}}=\operatorname{sgn}\circ\mathrm{pr}_{S_3}\) 在 \(S_3\) 上恰对换位类取 \(-1\)，对恒等类与三循环类取 \(+1\)。因此
\[
\chi_{\mathrm{LY}}(p)=-1
\iff
\sigma_3(p)\ \text{为换位类}
\iff
g\bmod p\ \text{型为 }(1)(2),
\]
这给出第 \(1\) 条。

若 \(\chi_{\mathrm{LY}}(p)=+1\)，则 \(\sigma_3(p)\) 只能落在恒等类或三循环类。其类大小分别为 \(1\) 与 \(2\)，故条件分布按 \(1:2\) 分配，于是
\[
\delta\!\bigl((1)(1)(1)\mid \chi_{\mathrm{LY}}=+1\bigr)=\frac{1}{1+2}=\frac13,
\qquad
\delta\!\bigl((3)\mid \chi_{\mathrm{LY}}=+1\bigr)=\frac{2}{1+2}=\frac23.
\]
这给出第 \(2\) 条。

对第 \(3\) 条，由定理 \ref{thm:xi-time-part62d-triple-frobenius-classfunction-tensor-independence}，\(\chi_{10}\) 仅依赖 \(S_{10}\) 因子，而 \(\chi_{4,\mathrm{sgn}},\chi_{4,\mathrm{aux}}\) 仅依赖 \(S_4\times C_2\) 因子；它们与只依赖 \(S_3\) 因子的 \(g\bmod p\) 分解型完全独立。故在固定后三者取值时，\(S_3\) 方向的条件分布不变。

进一步，由附录定理 \ref{thm:killo-leyang-quadratic-character-decoupling}，
\[
\delta\!\bigl(\chi_{10}=\epsilon_{10}\bigr)=\frac12.
\]
而在 \(S_4\times C_2\) 上，\(\chi_{4,\mathrm{sgn}}\) 与 \(\chi_{4,\mathrm{aux}}\) 分别是两个独立的 \(C_2\) 坐标，故
\[
\delta\!\bigl(\chi_{4,\mathrm{sgn}}=\epsilon_{4,\mathrm{sgn}},\ \chi_{4,\mathrm{aux}}=\epsilon_{4,\mathrm{aux}}\bigr)=\frac14.
\]
于是每个固定三元组 \((\epsilon_{10},\epsilon_{4,\mathrm{sgn}},\epsilon_{4,\mathrm{aux}})\) 在任一 \(\chi_{\mathrm{LY}}\)-扇区内都占据密度 \(1/8\)。这正表明其余三位只把每个 \(\chi_{\mathrm{LY}}\)-扇区均匀细分为 \(2^3\) 份。证毕。
\end{proof}

\begin{theorem}[Latt\`es 六次域的全分裂对任意 \texorpdfstring{$S_{10}\times S_3$}{S10 x S3} 事件实施精确 \texorpdfstring{$1/48$}{1/48} 稀释]\label{thm:xi-time-part62e-l4-fullsplit-universal-dilution}
沿用定理 \ref{thm:xi-time-part62d-triple-frobenius-classfunction-tensor-independence} 的记号，并令
\[
B_4^{\mathrm{split}}
:=
\{\text{\(C\bmod p\) 全分裂}\}.
\]
设
\[
E\subset S_{10}\times S_3
\]
为任意共轭类并；等价地，\(E\) 为任意只依赖 \(\bigl(\sigma_{10}(p),\sigma_3(p)\bigr)\) 的可审计事件。则有统一密度公式
\[
\boxed{\;
\delta\!\bigl(E\cap B_4^{\mathrm{split}}\bigr)=\frac{1}{48}\,\delta(E).\;}
\]

特别地，令
\[
A:=\{\text{\(P_{10}\bmod p\) 含一次因子}\},
\qquad
B_{\mathrm{split}}:=\{\text{\(g\bmod p\) 全分裂}\},
\]
则除已知恒等式
\[
\delta\!\bigl(A\cap B_{\mathrm{split}}\cap B_4^{\mathrm{split}}\bigr)
=
\frac{28319}{12902400}
\]
之外，还可得到
\[
\boxed{\;
\delta\!\bigl(A\cap\{g\bmod p\text{ 型为 }(1)(2)\}\cap B_4^{\mathrm{split}}\bigr)
=
\frac{28319}{4300800}.\;}
\]
\[
\boxed{\;
\delta\!\bigl(A\cap\{g\bmod p\text{ 型为 }(3)\}\cap B_4^{\mathrm{split}}\bigr)
=
\frac{28319}{6451200}.\;}
\]
\end{theorem}
\begin{proof}
在第三因子 \(S_4\times C_2\) 中，“\(C\bmod p\) 全分裂”等价于 Frobenius 落在恒等元共轭类 \(\{(1,1)\}\)。而事件 \(E\) 作为 \(S_{10}\times S_3\) 的共轭类并，对应于直积群
\[
S_{10}\times S_3\times (S_4\times C_2)
\]
中的共轭类并
\[
E\times\{(1,1)\}.
\]
由定理 \ref{thm:killo-leyang-lattes-triple-product-galois} 的直积结构与 Chebotarev 密度定理，
\[
\delta\!\bigl(E\cap B_4^{\mathrm{split}}\bigr)
=
\frac{|E|}{|S_{10}\times S_3|}\cdot \frac{1}{|S_4\times C_2|}
=
\delta(E)\cdot \frac{1}{48}.
\]
这就得到了统一稀释律。

对特例，由推论 \ref{cor:fold-gauge-anomaly-P10-modp-splitting-densities}，
\[
\delta(A)=\frac{28319}{44800}.
\]
再由结论 \ref{con:xi-terminal-zm-leyang-cubic-modp-splitting-density}，
\[
\delta\!\bigl(g\bmod p\text{ 型为 }(1)(2)\bigr)=\frac12,
\qquad
\delta(B_{\mathrm{split}})=\frac16,
\qquad
\delta\!\bigl(g\bmod p\text{ 型为 }(3)\bigr)=\frac13.
\]
由于 \(A\) 与 \(g\bmod p\) 的分解型分别只依赖 \(S_{10}\) 与 \(S_3\) 因子，它们彼此独立，故
\[
\delta\!\bigl(A\cap\{g\bmod p\text{ 型为 }(1)(2)\}\bigr)
=
\frac{28319}{44800}\cdot\frac12
=
\frac{28319}{89600},
\]
\[
\delta\!\bigl(A\cap B_{\mathrm{split}}\bigr)
=
\frac{28319}{44800}\cdot\frac16
=
\frac{28319}{268800},
\]
\[
\delta\!\bigl(A\cap\{g\bmod p\text{ 型为 }(3)\}\bigr)
=
\frac{28319}{44800}\cdot\frac13
=
\frac{28319}{134400}.
\]
把统一稀释律分别乘以 \(1/48\)，即得
\[
\frac{28319}{89600}\cdot \frac{1}{48}=\frac{28319}{4300800},
\qquad
\frac{28319}{268800}\cdot \frac{1}{48}=\frac{28319}{12902400},
\qquad
\frac{28319}{134400}\cdot \frac{1}{48}=\frac{28319}{6451200}.
\]
证毕。
\end{proof}

\paragraph{双指纹审计的张量闭合、最小二值统计与强乘法律}
在 \(P_{10}\times L_{\mathrm{LY}}\) 的双域 Chebotarev 乘积律、三域类函数张量分解、\(\chi_{\mathrm{LY}}\) 的唯一有效位以及 Latt\`es 六次域的统一稀释律已经确立之后，还可把双指纹审计进一步压缩为下列五条闭合后果。它们依次对应标准素数平均下的类函数张量因子化、一次因子事件与 Lee--Yang 三次的三种分裂型的完整联合密度、由“是否有根”与判别式符号构成的最小二值统计量、仅保留单一符号位时的精确信息亏损，以及任意有限布尔组合上的强乘法律。

\begin{theorem}[标准素数平均下的双指纹类函数张量因子化]\label{thm:xi-time-part62ea-dualfingerprint-classfunction-prime-average-factorization}
\leanverified{paper\_xi\_time\_part62ea\_dualfingerprint\_classfunction\_prime\_average\_factorization}
沿用命题 \ref{prop:fold-gauge-anomaly-P10-leyang-linear-disjointness} 的记号。对任意类函数
\[
f:S_{10}\to\CC,
\qquad
h:S_3\to\CC,
\]
对每个共同不分歧素数 \(p\) 记
\[
\sigma_{10}(p):=\mathrm{Frob}_p(L_{10}/\QQ)\in S_{10},
\qquad
\sigma_3(p):=\mathrm{Frob}_p(L_{\mathrm{LY}}/\QQ)\in S_3.
\]
则有极限恒等式
\[
\lim_{x\to\infty}
\frac{1}{\pi(x)}
\sum_{\substack{p\le x\\ p\nmid \Disc(P_{10})\Disc(P_{\mathrm{LY}})}}
f\!\bigl(\sigma_{10}(p)\bigr)\,
h\!\bigl(\sigma_3(p)\bigr)
=
\left(\frac{1}{10!}\sum_{\tau\in S_{10}}f(\tau)\right)
\left(\frac{1}{6}\sum_{\upsilon\in S_3}h(\upsilon)\right).
\]
特别地，若定义中心化类函数
\[
\widetilde f:=
f-\frac{1}{10!}\sum_{\tau\in S_{10}}f(\tau),
\qquad
\widetilde h:=
h-\frac{1}{6}\sum_{\upsilon\in S_3}h(\upsilon),
\]
则
\[
\lim_{x\to\infty}
\frac{1}{\pi(x)}
\sum_{\substack{p\le x\\ p\nmid \Disc(P_{10})\Disc(P_{\mathrm{LY}})}}
\widetilde f\!\bigl(\sigma_{10}(p)\bigr)\,
\widetilde h\!\bigl(\sigma_3(p)\bigr)
=0.
\]
\end{theorem}
\begin{proof}
由定理 \ref{thm:fold-gauge-anomaly-P10-leyang-class-function-tensor-independence}，
\[
\lim_{x\to\infty}
\frac{1}{\pi_{10,\mathrm{LY}}^\ast(x)}
\sum_{\substack{p\le x\\ p\nmid \Disc(P_{10})\Disc(P_{\mathrm{LY}})}}
f\!\bigl(\sigma_{10}(p)\bigr)\,
h\!\bigl(\sigma_3(p)\bigr)
=
\left(\frac{1}{10!}\sum_{\tau\in S_{10}}f(\tau)\right)
\left(\frac{1}{6}\sum_{\upsilon\in S_3}h(\upsilon)\right),
\]
其中
\[
\pi_{10,\mathrm{LY}}^\ast(x):=
\#\{p\le x:\ p\nmid \Disc(P_{10})\Disc(P_{\mathrm{LY}})\}.
\]
由于坏素数只有有限多个，
\[
\frac{\pi_{10,\mathrm{LY}}^\ast(x)}{\pi(x)}\longrightarrow 1
\qquad(x\to\infty),
\]
故把分母由 \(\pi_{10,\mathrm{LY}}^\ast(x)\) 改为 \(\pi(x)\) 不改变极限值，从而得到第一式。

第二式只需把 \(f,h\) 分别替换为 \(\widetilde f,\widetilde h\)。此时两者的群平均均为 \(0\)，故右端为 \(0\)。证毕。
\end{proof}

\begin{corollary}[一次因子事件与 Lee--Yang 三种分裂型的完整联合密度]\label{cor:xi-time-part62ea-A-leyang-threeway-density}
\leanverified{paper\_xi\_time\_part62ea\_a\_leyang\_threeway\_density}
沿用定理 \ref{thm:xi-time-part62ea-dualfingerprint-classfunction-prime-average-factorization} 的记号，并令
\[
g(y):=P_{\mathrm{LY}}(y)=256y^3+411y^2+165y+32.
\]
对共同不分歧素数 \(p\)，定义事件
\[
A:=\{\text{\(P_{10}\bmod p\) 含一次因子}\},
\]
\[
B_{\mathrm{split}}:=\{\text{\(g\bmod p\) 全分裂}\},
\qquad
B_{\mathrm{root}}:=\{\text{\(g\bmod p\) 至少有一根}\},
\]
\[
B_{\mathrm{single}}:=B_{\mathrm{root}}\setminus B_{\mathrm{split}},
\qquad
B_{\mathrm{noroot}}:=B_{\mathrm{root}}^{c}.
\]
则有
\[
\delta(A\cap B_{\mathrm{split}})=\frac{28319}{268800},
\]
\[
\delta(A\cap B_{\mathrm{single}})=\frac{28319}{89600},
\]
\[
\delta(A\cap B_{\mathrm{noroot}})=\frac{28319}{134400}.
\]
特别地，
\[
\delta(A\mid B_{\mathrm{split}})
=
\delta(A\mid B_{\mathrm{single}})
=
\delta(A\mid B_{\mathrm{noroot}})
=
\frac{28319}{44800}.
\]
\end{corollary}
\begin{proof}
由推论 \ref{cor:fold-gauge-anomaly-P10-modp-splitting-densities}，
\[
\delta(A)=\frac{28319}{44800}.
\]
由结论 \ref{con:xi-terminal-zm-leyang-cubic-modp-splitting-density}，
\[
\delta(B_{\mathrm{split}})=\frac16,
\qquad
\delta(B_{\mathrm{single}})=\frac12,
\qquad
\delta(B_{\mathrm{noroot}})=\frac13.
\]
再由推论 \ref{cor:fold-gauge-anomaly-P10-leyang-chebotarev-product-law}，\(A\) 与每个只依赖 \(S_3\) 分量的事件都严格独立，故
\[
\delta(A\cap E)=\delta(A)\delta(E)
\]
对 \(E\in\{B_{\mathrm{split}},B_{\mathrm{single}},B_{\mathrm{noroot}}\}\) 成立。代入上述边际密度即得三条联合密度公式。最后对各式分别除以 \(\delta(E)\) 便得到条件密度恒等式。证毕。
\end{proof}

\begin{theorem}[Lee--Yang 三分裂型的最小二值充分统计量]\label{thm:xi-time-part62ea-leyang-minimal-binary-sufficient-statistic}
沿用推论 \ref{cor:xi-time-part62ea-A-leyang-threeway-density} 的记号，并记
\[
\Sigma_{\mathrm{LY}}(p)\in\{(1)(1)(1),(1)(2),(3)\}
\]
为 \(g\bmod p\) 的分解型随机变量。定义二值观测对
\[
\mathcal O(p):=
\bigl(\mathbf 1_{B_{\mathrm{root}}}(p),\chi_{\mathrm{LY}}(p)\bigr)
\in\{0,1\}\times\{\pm1\}.
\]
则 \(\mathcal O\) 的像恰为
\[
\{(1,+1),(1,-1),(0,+1)\},
\]
并且存在双射
\[
(1)(1)(1)\longleftrightarrow (1,+1),
\qquad
(1)(2)\longleftrightarrow (1,-1),
\qquad
(3)\longleftrightarrow (0,+1).
\]
因而 \((0,-1)\) 从不出现，且 \(\mathcal O\) 与 \(\Sigma_{\mathrm{LY}}\) 等价地编码同一三值信息。特别地，在一切由二值坐标组成并能逐素完全恢复 \(\Sigma_{\mathrm{LY}}\) 的观测中，\(\mathcal O\) 具有最小坐标数；因此它构成 Lee--Yang 三分裂型的最小充分二值统计量。
\end{theorem}
\begin{proof}
由结论 \ref{con:xi-terminal-zm-leyang-cubic-modp-splitting-density}，
\[
\Sigma_{\mathrm{LY}}=(1)(1)(1)
\iff
g\bmod p\ \text{全分裂},
\]
\[
\Sigma_{\mathrm{LY}}=(1)(2)
\iff
g\bmod p\ \text{恰有一个一次因子},
\]
\[
\Sigma_{\mathrm{LY}}=(3)
\iff
g\bmod p\ \text{在 }\FF_p\text{ 上无根}.
\]
另一方面，由定理 \ref{thm:xi-time-part62e-chi-ly-sole-controller}，
\[
\chi_{\mathrm{LY}}(p)=-1
\iff
\Sigma_{\mathrm{LY}}(p)=(1)(2),
\]
\[
\chi_{\mathrm{LY}}(p)=+1
\iff
\Sigma_{\mathrm{LY}}(p)\in\{(1)(1)(1),(3)\}.
\]
再由 \(\mathbf 1_{B_{\mathrm{root}}}(p)=1\) 当且仅当 \(\Sigma_{\mathrm{LY}}(p)\in\{(1)(1)(1),(1)(2)\}\)，立刻得到上述对应表，故 \(\mathcal O\) 与 \(\Sigma_{\mathrm{LY}}\) 双射，因而充分性成立。

若只允许一个二值坐标，则观测值至多只有 \(2\) 种，而 \(\Sigma_{\mathrm{LY}}\) 有 \(3\) 种可能状态，因此不可能由单一二值观测完全恢复。故任何完备恢复 \(\Sigma_{\mathrm{LY}}\) 的二值编码都至少需要两个坐标；而 \(\mathcal O\) 恰以两个二值坐标实现该恢复，故其为最小。证毕。
\end{proof}

\begin{corollary}[单一符号位压缩的精确信息亏损]\label{cor:xi-time-part62ea-chi-ly-exact-information-loss}
\leanverified{paper\_xi\_time\_part62ea\_chi\_ly\_exact\_information\_loss}
沿用定理 \ref{thm:xi-time-part62ea-leyang-minimal-binary-sufficient-statistic} 的记号，并以 \(\log_2\) 为底定义 Shannon 熵。则有
\[
H(\Sigma_{\mathrm{LY}}\mid \chi_{\mathrm{LY}})
=
\frac12\,h_2\!\left(\frac13\right),
\qquad
h_2(x):=-x\log_2 x-(1-x)\log_2(1-x).
\]
等价地，
\[
H(\Sigma_{\mathrm{LY}})
=
1+\frac12\,h_2\!\left(\frac13\right),
\qquad
I(\Sigma_{\mathrm{LY}};\chi_{\mathrm{LY}})=1.
\]
数值上，
\[
H(\Sigma_{\mathrm{LY}}\mid \chi_{\mathrm{LY}})
\approx 0.4591479170\ \text{bits}.
\]
\end{corollary}
\begin{proof}
由定理 \ref{thm:killo-leyang-quadratic-character-decoupling}，
\[
\delta(\chi_{\mathrm{LY}}=+1)=\delta(\chi_{\mathrm{LY}}=-1)=\frac12.
\]
再由定理 \ref{thm:xi-time-part62ea-leyang-minimal-binary-sufficient-statistic}，
\[
\chi_{\mathrm{LY}}=-1
\Longrightarrow
\Sigma_{\mathrm{LY}}=(1)(2),
\]
故
\[
H(\Sigma_{\mathrm{LY}}\mid \chi_{\mathrm{LY}}=-1)=0.
\]
而在 \(\chi_{\mathrm{LY}}=+1\) 的条件下，只剩两种可能：
\[
\Sigma_{\mathrm{LY}}=(1)(1)(1)\quad\text{或}\quad \Sigma_{\mathrm{LY}}=(3).
\]
其条件概率由 \(S_3\) 中恒等类与三循环类的类大小比 \(1:2\) 给出，即
\[
\mathbb P\!\bigl(\Sigma_{\mathrm{LY}}=(1)(1)(1)\mid \chi_{\mathrm{LY}}=+1\bigr)=\frac13,
\]
\[
\mathbb P\!\bigl(\Sigma_{\mathrm{LY}}=(3)\mid \chi_{\mathrm{LY}}=+1\bigr)=\frac23.
\]
因此
\[
H(\Sigma_{\mathrm{LY}}\mid \chi_{\mathrm{LY}}=+1)=h_2\!\left(\frac13\right).
\]
对 \(\chi_{\mathrm{LY}}=\pm1\) 按各自概率 \(1/2\) 加权，得到
\[
H(\Sigma_{\mathrm{LY}}\mid \chi_{\mathrm{LY}})
=
\frac12\,h_2\!\left(\frac13\right).
\]

又由于 \(\chi_{\mathrm{LY}}\) 是 \(\Sigma_{\mathrm{LY}}\) 的确定函数且取值等分布，故
\[
H(\chi_{\mathrm{LY}})=1,
\qquad
I(\Sigma_{\mathrm{LY}};\chi_{\mathrm{LY}})=H(\chi_{\mathrm{LY}})=1.
\]
将链式公式
\[
H(\Sigma_{\mathrm{LY}})
=
I(\Sigma_{\mathrm{LY}};\chi_{\mathrm{LY}})
+H(\Sigma_{\mathrm{LY}}\mid \chi_{\mathrm{LY}})
\]
代入即可得到第二式。证毕。
\end{proof}

\begin{theorem}[双指纹有限布尔组合的强乘法律]\label{thm:xi-time-part62ea-dualfingerprint-boolean-strong-product-law}
\leanverified{paper\_xi\_time\_part62ea\_dualfingerprint\_boolean\_strong\_product\_law}
仍在共同不分歧素数的密度概率空间上，令 \(\mathcal A_{10}\) 为由全部事件
\[
\{\sigma_{10}(p)\in C\}
\qquad
\bigl(C\subset S_{10}\ \text{为共轭类}\bigr)
\]
生成的有限布尔代数，令 \(\mathcal A_3\) 为由全部事件
\[
\{\sigma_3(p)\in D\}
\qquad
\bigl(D\subset S_3\ \text{为共轭类}\bigr)
\]
生成的有限布尔代数，并记
\[
\chi_{10}(p):=\mathrm{sgn}\!\bigl(\sigma_{10}(p)\bigr).
\]
则对任意
\[
E\in\mathcal A_{10},
\qquad
F\in\mathcal A_3,
\]
都有精确乘法密度律
\[
\delta(E\cap F)=\delta(E)\,\delta(F).
\]
特别地，凡是仅由 \(A\)、\(\{\chi_{10}=+1\}\)、\(\{\chi_{10}=-1\}\) 及其他 \(S_{10}\)-循环型事件经有限次并、交、补得到的事件，均与仅由 \(B_{\mathrm{split}}\)、\(B_{\mathrm{single}}\)、\(B_{\mathrm{noroot}}\)、\(B_{\mathrm{root}}\)、\(\{\chi_{\mathrm{LY}}=+1\}\)、\(\{\chi_{\mathrm{LY}}=-1\}\) 及其他 \(S_3\)-分裂型事件经有限次并、交、补得到的事件严格独立。
\end{theorem}
\begin{proof}
由有限群上的布尔代数结构，每个 \(E\in\mathcal A_{10}\) 都可唯一写成若干 \(S_{10}\) 共轭类的有限不交并：
\[
E=\{\sigma_{10}(p)\in \widetilde E\},
\qquad
\widetilde E\subset S_{10}\ \text{为共轭类并},
\]
同理每个 \(F\in\mathcal A_3\) 都可写成
\[
F=\{\sigma_3(p)\in \widetilde F\},
\qquad
\widetilde F\subset S_3\ \text{为共轭类并}.
\]
因此群上的指示函数
\[
f:=\mathbf 1_{\widetilde E},
\qquad
h:=\mathbf 1_{\widetilde F}
\]
都是类函数。把它们代入定理
\ref{thm:xi-time-part62ea-dualfingerprint-classfunction-prime-average-factorization}，得到
\[
\delta(E\cap F)
=
\left(\frac{|\widetilde E|}{10!}\right)
\left(\frac{|\widetilde F|}{6}\right)
=
\delta(E)\,\delta(F).
\]
这就证明了强乘法律。

最后，\(A\) 是 \(S_{10}\) 中“存在不动点”所对应的共轭类并；\(\{\chi_{10}=\pm1\}\) 是偶/奇置换所对应的共轭类并；而 \(B_{\mathrm{split}},B_{\mathrm{single}},B_{\mathrm{noroot}},B_{\mathrm{root}},\{\chi_{\mathrm{LY}}=\pm1\}\) 全部都是 \(S_3\) 共轭类事件或其并。故列举出的有限布尔组合都落在相应布尔代数中，从而自动满足上式。证毕。
\end{proof}

\paragraph{辅助符号链、外侧审计免疫与固定窗口的张量障碍}
在定理 \ref{thm:xi-time-part62d-triple-frobenius-classfunction-tensor-independence} 已将三域 Frobenius 类函数平均压缩为严格张量乘积，而定理 \ref{thm:xi-time-part62e-triple-product-abelianization-c2-four} 与推论 \ref{cor:xi-time-part62ea-A-leyang-threeway-density} 已分别给出 \(C_2^4\) 的二次特征层与一次因子事件的双域联合密度之后，还可继续抽取出下列四条后果与一个结构性猜想。

\begin{corollary}[三条二次符号链 \texorpdfstring{$(\chi_{10},\chi_{\mathrm{LY}},\chi_4)$}{(chi10, chiLY, chi4)} 在 \texorpdfstring{$\{\pm1\}^3$}{pm1^3} 上严格等分布]\label{cor:xi-time-part62eb-three-sign-equidistribution}
沿用定理 \ref{thm:xi-time-part62e-triple-product-abelianization-c2-four} 的记号，并记
\[
\chi_4:=\chi_{4,\mathrm{aux}}=\iota\circ \mathrm{pr}_{C_2}:S_4\times C_2\to\{\pm1\}.
\]
则对任意
\[
(\varepsilon_{10},\varepsilon_{\mathrm{LY}},\varepsilon_4)\in\{\pm1\}^3
\]
都有
\[
\delta\!\bigl(
\chi_{10}(p)=\varepsilon_{10},
\chi_{\mathrm{LY}}(p)=\varepsilon_{\mathrm{LY}},
\chi_4(p)=\varepsilon_4
\bigr)
=
\frac18.
\]
从而任意两条符号被固定以后，第三条仍保持均匀分布。例如
\[
\delta\!\bigl(
\chi_4(p)=+1\mid
\chi_{10}(p)=\varepsilon_{10},
\chi_{\mathrm{LY}}(p)=\varepsilon_{\mathrm{LY}}
\bigr)=\frac12.
\]
\[
\delta\!\bigl(
\chi_4(p)=-1\mid
\chi_{10}(p)=\varepsilon_{10},
\chi_{\mathrm{LY}}(p)=\varepsilon_{\mathrm{LY}}
\bigr)=\frac12.
\]
\end{corollary}
\begin{proof}
取类函数
\[
f(\tau):=\mathbf 1_{\{\operatorname{sgn}(\tau)=\varepsilon_{10}\}},
\qquad
g(\upsilon):=\mathbf 1_{\{\operatorname{sgn}(\upsilon)=\varepsilon_{\mathrm{LY}}\}},
\qquad
h(\eta):=\mathbf 1_{\{\chi_4(\eta)=\varepsilon_4\}}.
\]
由 \(A_{10}\subset S_{10}\) 与 \(A_3\subset S_3\) 都是指数 \(2\) 子群，得
\[
\langle f\rangle_{S_{10}}=\frac12,
\qquad
\langle g\rangle_{S_3}=\frac12.
\]
而在 \(S_4\times C_2\) 中，\(\chi_4\) 只读取 \(C_2\) 坐标，故
\[
\#\{\eta\in S_4\times C_2:\ \chi_4(\eta)=\varepsilon_4\}=24,
\]
于是
\[
\langle h\rangle_{S_4\times C_2}=\frac{24}{48}=\frac12.
\]
代入定理 \ref{thm:xi-time-part62d-triple-frobenius-classfunction-tensor-independence} 即得联合密度为
\[
\frac12\cdot \frac12\cdot \frac12=\frac18.
\]
条件概率公式再由 Bayes 公式立得。证毕。
\end{proof}

\leanverified{paper\_xi\_time\_part62eb\_leyang\_external\_audit\_immunity}
\begin{proposition}[Lee--Yang 分裂型对一切外侧审计事件严格免疫]\label{prop:xi-time-part62eb-leyang-external-audit-immunity}
沿用定理 \ref{thm:xi-time-part62d-triple-frobenius-classfunction-tensor-independence} 的记号，并记
\[
\sigma_{\mathrm{LY}}(p):=\sigma_3(p)\in S_3,
\qquad
\sigma_{4,\mathrm{tot}}(p):=\bigl(\sigma_4(p),\varepsilon_4(p)\bigr)\in S_4\times C_2.
\]
设
\[
\mathcal E\subset S_{10}\times (S_4\times C_2)
\]
为任意非空共轭类并，并记对应的 prime 事件为
\[
E_{\mathcal E}:=
\bigl\{
p:\ (\sigma_{10}(p),\sigma_{4,\mathrm{tot}}(p))\in \mathcal E
\bigr\}.
\]
则对 \(S_3\) 的任意共轭类 \(D\subset S_3\)，都有
\[
\delta\!\bigl(
\sigma_{\mathrm{LY}}(p)\in D\mid E_{\mathcal E}
\bigr)
=
\frac{|D|}{6}.
\]
特别地，\(g\bmod p\) 的三种分裂型
\[
(1)(1)(1),\qquad (1)(2),\qquad (3)
\]
在任意外侧条件 \(E_{\mathcal E}\) 之下仍分别具有条件密度
\[
\frac16,\qquad \frac12,\qquad \frac13.
\]
\end{proposition}
\begin{proof}
记
\[
G_{\mathrm{ext}}:=S_{10}\times (S_4\times C_2).
\]
由定理 \ref{thm:xi-time-part62d-triple-frobenius-classfunction-tensor-independence} 所依赖的三重直积结构，
\[
\Gal(L_{10}L_{\mathrm{LY}}L_4/\QQ)\cong G_{\mathrm{ext}}\times S_3.
\]
因此对任意非空共轭类并 \(\mathcal E\subset G_{\mathrm{ext}}\) 与任意共轭类 \(D\subset S_3\)，Chebotarev 密度定理给出
\[
\delta\!\bigl(E_{\mathcal E}\cap\{\sigma_{\mathrm{LY}}(p)\in D\}\bigr)
=
\frac{|\mathcal E|}{|G_{\mathrm{ext}}|}\cdot \frac{|D|}{6}
=
\delta(E_{\mathcal E})\cdot \frac{|D|}{6}.
\]
对 \(\delta(E_{\mathcal E})>0\) 归一化便得
\[
\delta\!\bigl(
\sigma_{\mathrm{LY}}(p)\in D\mid E_{\mathcal E}
\bigr)
=
\frac{|D|}{6}.
\]
再把 \(S_3\) 的恒等类、换位类与三循环类大小 \(1,3,2\) 代入，即得三种分裂型的条件密度
\[
\frac16,\qquad \frac36=\frac12,\qquad \frac26=\frac13.
\]
证毕。
\end{proof}

\begin{theorem}[十次线性因子事件、辅助二次符号与 Lee--Yang 分裂型的六个显式联合密度]\label{thm:xi-time-part62eb-A-chi4-leyang-six-densities}
沿用推论 \ref{cor:xi-time-part62ea-A-leyang-threeway-density} 的记号，尤其
\[
A:=\{\text{\(P_{10}\bmod p\) 含一次因子}\},
\]
以及
\[
B_{\mathrm{split}},\qquad
B_{\mathrm{single}},\qquad
B_{\mathrm{noroot}},\qquad
B_{\mathrm{root}}=B_{\mathrm{split}}\sqcup B_{\mathrm{single}},
\]
其中 \(B_{\mathrm{split}},B_{\mathrm{single}},B_{\mathrm{noroot}}\) 分别对应 \(g\bmod p\) 的分裂型
\[
(1)(1)(1),\qquad (1)(2),\qquad (3).
\]
再令 \(\chi_4=\chi_{4,\mathrm{aux}}\) 如推论 \ref{cor:xi-time-part62eb-three-sign-equidistribution}。则对每个 \(\varepsilon\in\{\pm1\}\)，有
\[
\delta\!\bigl(A\cap\{\chi_4=\varepsilon\}\cap B_{\mathrm{split}}\bigr)
=
\frac{28319}{537600},
\]
\[
\delta\!\bigl(A\cap\{\chi_4=\varepsilon\}\cap B_{\mathrm{single}}\bigr)
=
\frac{28319}{179200},
\]
\[
\delta\!\bigl(A\cap\{\chi_4=\varepsilon\}\cap B_{\mathrm{noroot}}\bigr)
=
\frac{28319}{268800}.
\]
从而
\[
\delta\!\bigl(A\cap\{\chi_4=\varepsilon\}\cap B_{\mathrm{root}}\bigr)
=
\frac{28319}{134400}.
\]
\end{theorem}
\begin{proof}
由推论 \ref{cor:xi-time-part62ea-A-leyang-threeway-density}，
\[
\delta(A)=\frac{28319}{44800}.
\]
由推论 \ref{cor:xi-time-part62eb-three-sign-equidistribution}，
\[
\delta(\chi_4=\varepsilon)=\frac12.
\]
因此外侧事件
\[
E_\varepsilon:=A\cap\{\chi_4=\varepsilon\}
\]
的密度为
\[
\delta(E_\varepsilon)=\frac{28319}{44800}\cdot \frac12=\frac{28319}{89600}.
\]
而 \(E_\varepsilon\) 只依赖 \(S_{10}\times (S_4\times C_2)\) 因子，故由命题 \ref{prop:xi-time-part62eb-leyang-external-audit-immunity}，在该条件下 Lee--Yang 三种分裂型的条件密度仍分别为
\[
\frac16,\qquad \frac12,\qquad \frac13.
\]
于是
\[
\delta(E_\varepsilon\cap B_{\mathrm{split}})
=
\frac{28319}{89600}\cdot \frac16
=
\frac{28319}{537600},
\]
\[
\delta(E_\varepsilon\cap B_{\mathrm{single}})
=
\frac{28319}{89600}\cdot \frac12
=
\frac{28319}{179200},
\]
\[
\delta(E_\varepsilon\cap B_{\mathrm{noroot}})
=
\frac{28319}{89600}\cdot \frac13
=
\frac{28319}{268800}.
\]
最后，由
\[
B_{\mathrm{root}}=B_{\mathrm{split}}\sqcup B_{\mathrm{single}}
\]
且两者不交，故
\[
\delta(E_\varepsilon\cap B_{\mathrm{root}})
=
\frac{28319}{537600}+\frac{28319}{179200}
=
\frac{28319}{134400}.
\]
证毕。
\end{proof}

\begin{conjecture}[固定有限审计域窗口的张量刚性障碍]\label{conj:xi-time-part62eb-fixed-audit-window-tensor-rigidity}
设
\[
\{L_i/\QQ\}_{i=1}^{r}
\]
为一族固定的 Galois 扩张，满足各 \(L_i/\QQ\) 的伽罗瓦群 \(G_i\) 固定有限，且两两线性互不相交。记
\[
\mathrm{CF}(G_i)
\]
为 \(G_i\) 上的类函数代数。则任何仅由固定多个 Frobenius 因子
\[
\bigl(\mathrm{Frob}_p(L_1/\QQ),\dots,\mathrm{Frob}_p(L_r/\QQ)\bigr)
\]
上的固定类函数、固定二次特征与固定有限布尔组合所生成的 prime-fingerprint 审计体系，在自然密度极限下都只通过有限维交换张量代数
\[
\bigotimes_{i=1}^{r}\mathrm{CF}(G_i)
\]
及其由共轭类并所生成的有限布尔代数起作用。因而，只要域族、群族与线性互不相交模式固定而不随输入规模增长，此类 Chebotarev 指纹审计不足以单独承载 NP-完全型的组合非对称性。若要在 Frobenius 指纹层面逼近 \(P\)–\(NP\) 分离，必须允许审计域链、Galois 群规模或交域失效机制随输入长度增长。
\end{conjecture}


\noindent
上述结果表明：三域 Frobenius 直积结构在 \(C_2^4\) 符号层、Lee--Yang 三分裂型与 Latt\`es 六次全分裂条件上不仅给出唯一有效位与 \(1/48\) 稀释律，而且在忽略第三因子后继续诱导出 \(S_{10}\times S_3\) 双指纹的标准素数平均张量分解、一次因子事件与 Lee--Yang 三次的三种分裂型的完整联合密度、由 \((\mathbf 1_{B_{\mathrm{root}}},\chi_{\mathrm{LY}})\) 承载的最小二值充分统计，以及任意有限布尔组合上的强乘法律。进一步，辅助 \(C_2\) 符号 \(\chi_4\) 与 \((\chi_{10},\chi_{\mathrm{LY}})\) 形成严格的三位等分布，任意外侧审计事件条件化都不改变 Lee--Yang 的 \(S_3\) 分裂分布，而一次因子事件与辅助符号的三重联合密度也被迫完全有理化。于是，固定有限审计域窗口所能生成的全部指纹统计便进一步暴露为一个有限维张量对象。

\paragraph{Gödel--Zeckendorf 地址临界测度、负谱单支障碍与原子激活正规形}
在 Gödel--Zeckendorf 地址极限的严格自仿射闭合、纯碰撞位势的 kernel RH 临界律、primitive 原子剥离以及激活分支的局部解析展开已经建立之后，还可进一步抽取出下列七条彼此衔接的结构后果。第一条把黄金地址极限的临界几何常数完全闭式化；中间四条把 RH 阈值、解析缓冲层、两步原子扇区及其 Abel 有限部校正分别压缩为单支障碍、显式长度、一维 Witt 校正与闭式常数位移；最后两条则把激活分支在 \(u=1\) 附近的局部坐标先提升为逆正规形，再进一步压缩为无二次项的二阶沉默。

\begin{theorem}[Gödel--Zeckendorf 地址极限的临界 Hausdorff 测度闭式]\label{thm:xi-time-part62d-zg-critical-hausdorff-measure}
\leanverified{paper\_xi\_time\_part62d\_zg\_critical\_hausdorff\_measure}
沿用定理 \ref{thm:xi-zg-address-binary-self-affine}、\ref{thm:xi-time-part9g-holographic-prefix-isometry-on-line} 与
\ref{thm:xi-zg-address-minkowski-dimension-content} 的记号。记
\[
s_B:=\frac{\log\varphi}{\log B}.
\]
则 \(K_{B,v}\) 的 \(s_B\)-维 Hausdorff 测度满足
\[
\mathcal H^{s_B}(K_{B,v})=\frac{\varphi^2}{\sqrt5}.
\]
\end{theorem}
\begin{proof}
由定理 \ref{thm:xi-zg-address-minkowski-dimension-content}，对每个 \(n\ge 0\)，覆盖 \(K_{B,v}\) 所需的最少 \(B^{-n}\)-球个数恰为
\[
N_{B,v}(n)=F_{n+2}.
\]
因此用这组球直接覆盖可得
\[
\mathcal H_{B^{-n}}^{s_B}(K_{B,v})
\le
F_{n+2}B^{-ns_B}.
\]

另一方面，由定理 \ref{thm:xi-time-part9g-holographic-prefix-isometry-on-line} 的柱--球对应，长度 \(n\) 的不同 admissible 前缀落在彼此不同的 \(B^nL_v\)-余类中。故任一直径不超过 \(B^{-n}\) 的集合都至多与其中一个余类相交，从而任意这类覆盖都至少需要 \(F_{n+2}\) 个集合。于是
\[
\mathcal H_{B^{-n}}^{s_B}(K_{B,v})
\ge
F_{n+2}B^{-ns_B}.
\]
故对每个 \(n\) 都有
\[
\mathcal H_{B^{-n}}^{s_B}(K_{B,v})
=
F_{n+2}B^{-ns_B}.
\]

再由 \(B^{s_B}=\varphi\) 与 Binet 公式，
\[
F_{n+2}B^{-ns_B}=F_{n+2}\varphi^{-n}\longrightarrow \frac{\varphi^2}{\sqrt5}.
\]
当 \(n\to\infty\) 时 \(B^{-n}\downarrow 0\)，故由 Hausdorff 外测度的单调极限得到
\[
\mathcal H^{s_B}(K_{B,v})
=
\lim_{n\to\infty}\mathcal H_{B^{-n}}^{s_B}(K_{B,v})
=
\frac{\varphi^2}{\sqrt5}.
\]
证毕。
\end{proof}

\begin{theorem}[核版 RH/Ramanujan 障碍由唯一负谱支完全控制]\label{thm:xi-time-part62d-collision-rh-negative-branch-control}
沿用命题 \ref{prop:real-input-40-collision-all-real} 与
\ref{prop:real-input-40-collision-rhk-window} 的记号。设
\[
f(\lambda;u)=\lambda^3-2\lambda^2-(u+1)\lambda+u
\]
的三条实根按
\[
\lambda_+(u)>1,\qquad 0<\lambda_0(u)<1,\qquad \lambda_-(u)<0
\]
排序，并记
\[
\Lambda(u):=\max\{|\lambda_0(u)|,|\lambda_-(u)|,1\}.
\]
则 kernel RH 条件等价于
\[
\textsf{RH}(u)\iff
\begin{cases}
\text{恒真}, & 0<u\le 1,\\[1mm]
(-\lambda_-(u))^2\le \lambda_+(u), & u>1.
\end{cases}
\]
特别地，\(\lambda_0(u)\) 从不决定 RH 窗口边界；唯一的障碍分支始终是负根 \(\lambda_-(u)\)。
\end{theorem}
\begin{proof}
由命题 \ref{prop:real-input-40-collision-all-real}，
\[
0<\lambda_0(u)<1
\qquad(\forall\,u>0),
\]
故 \(|\lambda_0(u)|<1\) 对全部正参数恒成立。

当 \(0<u\le 1\) 时，同一命题给出
\[
\lambda_-(u)\in[-1,0),
\]
于是
\[
|\lambda_-(u)|\le 1,\qquad |\lambda_0(u)|<1.
\]
从而
\[
\Lambda(u)=1.
\]
又因 \(\lambda_+(u)>1\)，故
\[
1\le \lambda_+(u)^{1/2},
\]
于是 \(\textsf{RH}(u)\) 自动成立。

当 \(u>1\) 时，命题 \ref{prop:real-input-40-collision-all-real} 给出
\[
\lambda_-(u)<-1.
\]
因此
\[
|\lambda_-(u)|>1>|\lambda_0(u)|,
\]
从而
\[
\Lambda(u)=|\lambda_-(u)|=-\lambda_-(u).
\]
代回命题 \ref{prop:real-input-40-collision-rhk-window} 的判据
\[
\textsf{RH}(u)\iff \Lambda(u)\le \lambda_+(u)^{1/2}
\]
便得
\[
\textsf{RH}(u)\iff -\lambda_-(u)\le \lambda_+(u)^{1/2}
\iff
(-\lambda_-(u))^2\le \lambda_+(u).
\]
证毕。
\end{proof}

\begin{proposition}[RH 临界参数的结果式消元范式]\label{prop:xi-time-part62d-collision-rh-resultant-template}
设
\[
f(\lambda;u)\in \RR[\lambda,u]
\]
是关于 \(\lambda\) 的首一三次多项式，并且在某个区间 \(I\subset(0,\infty)\) 上具有三条实根分支
\[
\lambda_+(u)>\lambda_0(u)>\lambda_-(u),
\qquad
\lambda_+(u)>0>\lambda_-(u).
\]
若在某个 \(u\in I\) 处发生 Ramanujan 型达界事件
\[
|\lambda_-(u)|=\sqrt{\lambda_+(u)},
\]
则必有
\[
\operatorname{Res}_{\lambda}\bigl(f(\lambda;u),f(\lambda^2;u)\bigr)=0.
\]
对纯碰撞三次
\[
f(\lambda;u)=\lambda^3-2\lambda^2-(u+1)\lambda+u
\]
应用该范式，可得精确因式分解
\[
\operatorname{Res}_{\lambda}\bigl(f(\lambda;u),f(\lambda^2;u)\bigr)
=
2u\bigl(u^5+4u^4+3u^3-96u^2-42u-14\bigr).
\]
因此正参数域中的 Ramanujan 临界候选恰由五次方程
\[
u^5+4u^4+3u^3-96u^2-42u-14=0
\]
给出。
\end{proposition}
\begin{proof}
若 \(|\lambda_-(u)|=\sqrt{\lambda_+(u)}\)，由于 \(\lambda_-(u)<0\)，便有
\[
\lambda_-(u)^2=\lambda_+(u).
\]
取 \(\lambda=\lambda_-(u)\)，则
\[
f(\lambda;u)=0,
\qquad
f(\lambda^2;u)=f(\lambda_+(u);u)=0.
\]
故 \(f(\lambda;u)\) 与 \(f(\lambda^2;u)\) 在变量 \(\lambda\) 上有公共根，于是其结果式必为零。

对给定三次
\[
f(\lambda;u)=\lambda^3-2\lambda^2-(u+1)\lambda+u,
\]
直接计算 Sylvester 结果式即得
\[
\operatorname{Res}_{\lambda}\bigl(f(\lambda;u),f(\lambda^2;u)\bigr)
=
2u\bigl(u^5+4u^4+3u^3-96u^2-42u-14\bigr).
\]
由此可知所有正参数临界候选都必须落在该五次方程的正根集合中。再结合命题 \ref{prop:real-input-40-collision-rhk-window} 的唯一性结论，便得到真正的临界参数 \(u_R\)。证毕。
\end{proof}

\begin{corollary}[解析但非 RH 实轴相的显式长度与负谱触发机制]\label{cor:xi-time-part62d-analytic-nonrh-window-length}
沿用定理 \ref{thm:xi-rh-critical-before-nearest-complex-branch} 与
\ref{thm:xi-time-part62d-collision-rh-negative-branch-control} 的记号。记
\[
\delta_{\mathrm{an}}:=R_\theta-\theta_R.
\]
则
\[
\delta_{\mathrm{an}}
=
R_\theta-\theta_R
\approx
1.48342289346087>0.
\]
并且对任意
\[
\theta\in(\theta_R,R_\theta)
\]
都有
\[
(-\lambda_-(e^\theta))^2>\lambda_+(e^\theta).
\]
换言之，在这一非空实区间内，Perron 压力分支仍保持解析，而 kernel RH 的失效已经由负谱支的平方越界单独触发。
\end{corollary}
\begin{proof}
由定理 \ref{thm:xi-rh-critical-before-nearest-complex-branch}，
\[
\theta_R<R_\theta,
\qquad
R_\theta\approx 2.76230289346087,
\qquad
\theta_R\approx 1.27888.
\]
故
\[
\delta_{\mathrm{an}}=R_\theta-\theta_R\approx 1.48342289346087.
\]

若 \(\theta\in(\theta_R,R_\theta)\)，则首先 \(|\theta|<R_\theta\)，因而 Perron 分支 \(P_2(\theta)\) 仍位于解析收敛盘内部。另一方面 \(u=e^\theta>\!u_R>1\)，由定理 \ref{thm:xi-time-part62d-collision-rh-negative-branch-control}，
\[
\textsf{RH}(u)\iff (-\lambda_-(u))^2\le \lambda_+(u).
\]
而此时 RH 已失效，故必有
\[
(-\lambda_-(u))^2>\lambda_+(u).
\]
把 \(u=e^\theta\) 代回即得所示不等式。证毕。
\end{proof}

\begin{theorem}[两步原子扇区在 primitive 层的一维 Witt 塌缩]\label{thm:xi-time-part62d-atomic-two-step-witt-rankone-collapse}
沿用命题 \ref{prop:atomic-2-orbit-split-logM} 与定理
\ref{thm:xi-atomic-witt-cycle-primitive-exact-peeling} 的记号。记
\[
P(z,u)=\sum_{n\ge 1}p_n(u)z^n,
\qquad
P_{\mathrm{core}}(z,u)=\sum_{n\ge 1}p_n^{\mathrm{core}}(u)z^n.
\]
则：
\begin{enumerate}
  \item 对一切 \(n\neq 2\)，有
  \[
  p_n(u)=p_n^{\mathrm{core}}(u);
  \]
  \item 在 \(n=2\) 处唯一出现原子修正
  \[
  p_2(u)=u+p_2^{\mathrm{core}}(u).
  \]
\end{enumerate}
此外，命题 \ref{prop:atomic-2-orbit-split-logM} 所嵌审计输出表明，branch 层满足
\[
(|\tau|,E)=(2,1)
\]
的两步回路恰有 \(7\) 条。因而这 \(7\) 条几何上彼此不同的两步 branch 回路在 primitive--Witt 层的净效应严格压缩为单一 monomial
\[
u z^2.
\]
等价地，
\[
\dim_{\QQ(u)}
\operatorname{span}\{\,p_n(u)-p_n^{\mathrm{core}}(u):n\ge 1\,\}=1.
\]
\end{theorem}
\begin{proof}
定理 \ref{thm:xi-atomic-witt-cycle-primitive-exact-peeling} 已给出精确剥离
\[
P(z,u)=u z^2+P_{\mathrm{core}}(z,u).
\]
逐项比较系数，立即得到
\[
p_n(u)=p_n^{\mathrm{core}}(u)\quad(n\neq 2),
\qquad
p_2(u)=u+p_2^{\mathrm{core}}(u).
\]
故 full/core 的 primitive 差分仅支撑在长度 \(2\) 的单一坐标上，其 \(\QQ(u)\)-线性张成空间因此恰为一维。

另一方面，命题 \ref{prop:atomic-2-orbit-split-logM} 同页嵌入的审计输出已经逐项列出 branch 层全部满足
\[
(|\tau|,E)=(2,1)
\]
的两步回路，并给出其总数恰为 \(7\)。于是这 \(7\) 条 branch 两步回路在 primitive--Witt 层的总校正只能通过唯一项
\[
u z^2
\]
出现。证毕。
\end{proof}

\begin{theorem}[七条两步 branch 回路在 primitive--Witt 与 Abel 有限部中只留下单原子校正]\label{thm:xi-time-part62d-seven-two-step-branch-loops-single-atom-finitepart-shift}
沿用命题 \ref{prop:atomic-2-orbit-split-logM} 与定理
\ref{thm:xi-time-part62d-atomic-two-step-witt-rankone-collapse} 的记号。令
\[
P(z)=\sum_{n\ge 1}p_nz^n,
\qquad
P_{\mathrm{core}}(z)=\sum_{n\ge 1}p_n^{\mathrm{core}}z^n
\]
表示 \(u=1\) 时的 primitive 系列，并记
\[
P(z)=-\log(1-\lambda z)+\log\mathfrak M+o(1),
\]
\[
P_{\mathrm{core}}(z)=-\log(1-\lambda z)+\log\mathfrak M_{\mathrm{core}}+o(1)
\qquad
(z\to z_\star),
\]
其中
\[
\lambda=\varphi^2,
\qquad
z_\star=\lambda^{-1}=\varphi^{-2}.
\]
则
\[
p_2=1+p_2^{\mathrm{core}},
\qquad
p_n=p_n^{\mathrm{core}}
\quad(n\neq 2).
\]
特别地，branch 层中恰为 \(7\) 条的两步回路在 primitive--Witt 层只留下唯一原子
\[
z^2.
\]
并且 Abel 有限部常数满足精确位移
\[
\log\mathfrak M-\log\mathfrak M_{\mathrm{core}}
=
z_\star^2
=
\varphi^{-4}
=
\frac{7-3\sqrt5}{2}.
\]
等价地，
\[
\log\mathfrak M_{\mathrm{core}}
=
\log\mathfrak M-\varphi^{-4}.
\]
\end{theorem}
\begin{proof}
将定理 \ref{thm:xi-time-part62d-atomic-two-step-witt-rankone-collapse} 在 \(u=1\) 处专门化，即得
\[
P(z)=z^2+P_{\mathrm{core}}(z).
\]
逐项比较系数便得到
\[
p_2=1+p_2^{\mathrm{core}},
\qquad
p_n=p_n^{\mathrm{core}}
\quad(n\neq 2).
\]
同一定理已经说明，branch 层满足
\[
(|\tau|,E)=(2,1)
\]
的两步回路恰有 \(7\) 条，故它们在 primitive--Witt 层的净效应正是唯一原子 \(z^2\)。

另一方面，命题 \ref{prop:atomic-2-orbit-split-logM} 给出精确常数分解
\[
\log\mathfrak M
=
\log\mathfrak M_{\mathrm{core}}+z_\star^2.
\]
再代入
\[
z_\star=\varphi^{-2}
\]
即得
\[
\log\mathfrak M-\log\mathfrak M_{\mathrm{core}}
=
\varphi^{-4}
=
\frac{7-3\sqrt5}{2}.
\]
证毕。
\end{proof}

\begin{theorem}[激活分支在 \texorpdfstring{$u=1$}{u=1} 的逆正规形与无二次重整化]\label{thm:xi-time-part62d-activated-branch-inverse-normalform}
沿用注 \ref{rem:real-input-40-activated-branch-series} 与定理
\ref{thm:xi-output-potential-activated-branch-rigidity} 的记号。令
\[
h:=u-1,
\qquad
\lambda:=\lambda_{\mathrm{act}}(u).
\]
则正向展开
\[
\lambda
=
h+3h^3-h^4+18h^5-22h^6+O(h^7)
\]
在局部可逆，并且其反函数满足显式逆正规形
\[
h
=
\lambda-3\lambda^3+\lambda^4+9\lambda^5+\lambda^6+O(\lambda^7).
\]
特别地：
\begin{enumerate}
  \item 正向展开与逆正规形中都不存在二次项；
  \item 第一非线性修正项严格为三次项；
  \item 在逆正规形一侧，该三次系数恰为 \(-3\)。
\end{enumerate}
\end{theorem}
\begin{proof}
由注 \ref{rem:real-input-40-activated-branch-series}，
\[
\lambda
=
h+3h^3-h^4+18h^5-22h^6+O(h^7).
\]
设其局部反函数写成
\[
h
=
\lambda+a_2\lambda^2+a_3\lambda^3+a_4\lambda^4+a_5\lambda^5+a_6\lambda^6+O(\lambda^7).
\]
把此式代回上面的正向展开，并按 \(\lambda\) 的幂次逐项比较。

\(\lambda^2\) 项给出
\[
a_2=0.
\]
\(\lambda^3\) 项给出
\[
a_3+3=0,
\]
故
\[
a_3=-3.
\]
\(\lambda^4\) 项给出
\[
a_4-1=0,
\]
故
\[
a_4=1.
\]
\(\lambda^5\) 项给出
\[
a_5+9a_3+18=0,
\]
从而
\[
a_5=9.
\]
\(\lambda^6\) 项给出
\[
a_6+9a_4-4a_3-22=0,
\]
故
\[
a_6=1.
\]
于是
\[
h
=
\lambda-3\lambda^3+\lambda^4+9\lambda^5+\lambda^6+O(\lambda^7).
\]
所述三条推论随即成立。证毕。
\end{proof}

\begin{theorem}[激活支的二阶沉默]\label{thm:xi-time-part62d-activated-branch-second-order-silence}
\leanverified{paper\_xi\_time\_part62d\_activated\_branch\_second\_order\_silence}
沿用定理 \ref{thm:xi-time-part62d-activated-branch-inverse-normalform} 的记号，记
\[
h:=u-1.
\]
则激活分支满足显式展开
\[
\lambda_{\mathrm{act}}(u)
=
h+3h^3-h^4+18h^5-22h^6+O(h^7),
\]
从而
\[
\lambda_{\mathrm{act}}(u)-(u-1)=O\!\bigl((u-1)^3\bigr).
\]
特别地，激活支在 \(u=1\) 邻域内不存在二次弯曲项，首个非线性修正严格出现在三次项。
\end{theorem}
\begin{proof}
定理 \ref{thm:xi-time-part62d-activated-branch-inverse-normalform} 的正向展开已给出
\[
\lambda_{\mathrm{act}}(u)
=
h+3h^3-h^4+18h^5-22h^6+O(h^7).
\]
从该式直接读出二次项系数为 \(0\)，并把 \(h=u-1\) 代回，即得
\[
\lambda_{\mathrm{act}}(u)-(u-1)=O\!\bigl((u-1)^3\bigr).
\]
证毕。
\end{proof}

\noindent
上述七条结果把 H.160--H.162 的 \(B\)-进地址链与 D.509--D.532 的纯谱/Witt/激活链压缩为同一层二值刚性：一方面，Gödel--Zeckendorf 地址极限在临界维数处给出完全显式的黄金常数；另一方面，kernel RH 的破裂不再依赖多支比较，而被负谱单支完全控制，并在解析域内部留下显式长度的缓冲层；与此同时，两步原子扇区在 primitive 层只留下一个一维方向，而激活分支的局部坐标直到二阶都保持严格线性，首个非线性修正只能在三阶显影。于是，地址几何、谱阈值、Witt 剥离与局部正规形便被统一压缩为同一条可审计的刚性主线。

\paragraph{黄金柱测度、Smith 最小周期与局部化商序的有限刚性}
在黄金地址极限的临界测度闭式、Smith 周期平均律以及局部化数系的同态/商分类已经建立之后，还可把柱级质量、最小周期与对偶商序继续压缩为下列五条闭合结论。其共同特征是：无限地址链、全部模数层与全部连续商结构，分别由有限前缀、最大 Smith 不变量与素数集包含偏序这三类极小证书完全控制。

\begin{conclusion}[Gödel--Zeckendorf 地址极限的标准 Hausdorff 概率测度在柱集上具有严格黄金闭式]\label{con:xi-time-part62da-zg-cylinder-hausdorff-probability}
沿用定理 \ref{thm:xi-zg-address-binary-self-affine} 与定理
\ref{thm:xi-time-part62d-zg-critical-hausdorff-measure} 的记号，记
\[
s_B:=\frac{\log\varphi}{\log B},
\qquad
\mu_{B,v}(A):=
\frac{\sqrt5}{\varphi^2}\,\mathcal H^{s_B}(A\cap K_{B,v})
\qquad(A\subseteq L_v\ \text{Borel}).
\]
对任意长度 \(n\) 的 admissible 前缀
\[
a=(a_1,\dots,a_n)\in\{0,1\}^n,
\qquad
a_ia_{i+1}=0\quad(1\le i<n),
\]
定义柱集
\[
K[a]:=
\Bigl\{
\Xi_B(z):\ z\in \Sigma_{\mathrm{ZG}},\ z_1=a_1,\dots,z_n=a_n
\Bigr\}
\subset K_{B,v},
\]
并记
\[
\ell(a):=a_n\in\{0,1\}.
\]
则有精确公式
\[
\mu_{B,v}(K[a])=\varphi^{-(n+\ell(a))}.
\]
特别地，
\[
\mu_{B,v}(K[0])=\varphi^{-1},
\qquad
\mu_{B,v}(K[1])=\varphi^{-2}.
\]
\end{conclusion}
\begin{proof}
由定理 \ref{thm:xi-time-part62d-zg-critical-hausdorff-measure}，
\[
\mathcal H^{s_B}(K_{B,v})=\frac{\varphi^2}{\sqrt5}.
\]
再记
\[
\tau_a:=\sum_{t=1}^{n}a_tB^{t-1}v.
\]

若 \(a_n=0\)，则尾部仍可接任意 admissible 序列，故
\[
K[a]=\tau_a+B^nK_{B,v}.
\]
若 \(a_n=1\)，则 admissibility 强制下一位必为 \(0\)，从而
\[
K[a]=K[a0]=\tau_a+B^{n+1}K_{B,v}.
\]

Hausdorff 测度在地址线 \(L_v\) 上平移不变，而映射 \(x\mapsto B^r x\) 把距离缩放为 \(B^{-r}\)。因此对任意 Borel 集 \(E\subseteq L_v\) 都有
\[
\mathcal H^{s_B}(\tau+B^rE)
=
B^{-rs_B}\mathcal H^{s_B}(E)
=
\varphi^{-r}\mathcal H^{s_B}(E).
\]
于是当 \(a_n=0\) 时，
\[
\mathcal H^{s_B}(K[a])
=
\varphi^{-n}\mathcal H^{s_B}(K_{B,v})
=
\varphi^{-n}\cdot \frac{\varphi^2}{\sqrt5},
\]
当 \(a_n=1\) 时，
\[
\mathcal H^{s_B}(K[a])
=
\varphi^{-(n+1)}\mathcal H^{s_B}(K_{B,v})
=
\varphi^{-(n+1)}\cdot \frac{\varphi^2}{\sqrt5}.
\]
两种情形统一写成
\[
\mathcal H^{s_B}(K[a])
=
\varphi^{-(n+\ell(a))}\cdot \frac{\varphi^2}{\sqrt5}.
\]
再乘上归一化因子 \(\sqrt5/\varphi^2\)，即得
\[
\mu_{B,v}(K[a])=\varphi^{-(n+\ell(a))}.
\]
证毕。
\end{proof}

\begin{conclusion}[Gödel--Zeckendorf 地址极限上的标准 Hausdorff 概率测度是唯一的黄金 conformal 测度]\label{con:xi-time-part62da-zg-unique-golden-conformal-measure}
沿用结论 \ref{con:xi-time-part62da-zg-cylinder-hausdorff-probability} 的记号，并定义两条收缩分支
\[
\psi_0(x):=Bx,
\qquad
\psi_1(x):=v+B^2x.
\]
则 \(K_{B,v}\) 满足严格不交并分解
\[
K_{B,v}=\psi_0(K_{B,v})\sqcup \psi_1(K_{B,v}),
\]
并且 \(\mu_{B,v}\) 是唯一满足
\[
\mu
=
\varphi^{-1}(\psi_0)_\#\mu+\varphi^{-2}(\psi_1)_\#\mu
\]
的 \(K_{B,v}\) 上 Borel 概率测度。
\end{conclusion}
\begin{proof}
定理 \ref{thm:xi-zg-address-binary-self-affine} 已给出严格不交并分解
\[
K_{B,v}=BK_{B,v}\sqcup\bigl(v+B^2K_{B,v}\bigr)
=
\psi_0(K_{B,v})\sqcup \psi_1(K_{B,v}).
\]
又由结论 \ref{con:xi-time-part62da-zg-cylinder-hausdorff-probability}，
\[
\mu_{B,v}(K[0])=\varphi^{-1},
\qquad
\mu_{B,v}(K[1])=\varphi^{-2},
\]
其中
\[
K[0]=\psi_0(K_{B,v}),
\qquad
K[1]=\psi_1(K_{B,v}).
\]

对任意 Borel 集 \(A\subseteq K_{B,v}\)，由上面两块互不相交并结合 Hausdorff 测度的缩放关系，可得
\[
\mu_{B,v}(\psi_0(A))=\varphi^{-1}\mu_{B,v}(A),
\qquad
\mu_{B,v}(\psi_1(A))=\varphi^{-2}\mu_{B,v}(A).
\]
于是对任意 Borel 集 \(E\subseteq K_{B,v}\)，有
\[
\mu_{B,v}(E)
=
\mu_{B,v}(E\cap K[0])+\mu_{B,v}(E\cap K[1])
=
\varphi^{-1}\mu_{B,v}(\psi_0^{-1}(E))
+\varphi^{-2}\mu_{B,v}(\psi_1^{-1}(E)),
\]
即
\[
\mu_{B,v}
=
\varphi^{-1}(\psi_0)_\#\mu_{B,v}
+\varphi^{-2}(\psi_1)_\#\mu_{B,v}.
\]

反过来，设 \(\nu\) 为任意满足该自相似方程的 Borel 概率测度。由于分支像互不相交，对任意 Borel 集 \(A\subseteq K_{B,v}\) 立即推出
\[
\nu(\psi_0(A))=\varphi^{-1}\nu(A),
\qquad
\nu(\psi_1(A))=\varphi^{-2}\nu(A).
\]
特别地，
\[
\nu(K[0])=\varphi^{-1},
\qquad
\nu(K[1])=\varphi^{-2}.
\]

现取任意 admissible 前缀 \(a\)。若 \(a_n=1\)，则 \(K[a]=K[a0]\)；故记
\[
\widetilde a:=
\begin{cases}
a,& a_n=0,\\
a0,& a_n=1.
\end{cases}
\]
则 \(K[a]=K[\widetilde a]\)，且 \(\widetilde a\) 唯一分解为若干块 \(0\) 与 \(10\) 的串接：
\[
\widetilde a=\beta_1\cdots \beta_m,
\qquad
\beta_j\in\{0,10\}.
\]
若记
\[
\psi_{0}:=\psi_0,
\qquad
\psi_{10}:=\psi_1,
\]
则
\[
K[a]
=
\psi_{\beta_1}\circ\cdots\circ\psi_{\beta_m}(K_{B,v}).
\]
反复应用上面的缩放关系，得到
\[
\nu(K[a])
=
\prod_{j=1}^{m}w(\beta_j),
\qquad
w(0)=\varphi^{-1},
\quad
w(10)=\varphi^{-2}.
\]
而 \(\widetilde a\) 的总长度正是 \(n+\ell(a)\)，故
\[
\nu(K[a])=\varphi^{-(n+\ell(a))}.
\]
这与结论 \ref{con:xi-time-part62da-zg-cylinder-hausdorff-probability} 对 \(\mu_{B,v}\) 给出的柱质量完全一致。由于柱集生成 \(K_{B,v}\) 的 Borel \(\sigma\)-代数，故
\[
\nu=\mu_{B,v}.
\]
证毕。
\end{proof}

\begin{conclusion}[Smith 归一化模核计数的最小正周期恰等于最大 Smith 不变量]\label{con:xi-time-part62da-smith-kappa-minimal-period}
设 \(A\in M_{m\times n}(\ZZ)\) 满足
\[
\operatorname{rank}_{\QQ}(A)=m,
\]
其 Smith 不变量为
\[
d_1\mid d_2\mid\cdots\mid d_m.
\]
定义归一化模核计数
\[
\kappa_A(b):=
\frac{\#\ker(A_b)}{b^{\,n-m}}
=
\prod_{i=1}^{m}\gcd(d_i,b)
\qquad(b\ge 1).
\]
则 \(\kappa_A\) 的最小正周期恰为 \(d_m\)。并且其普通生成函数与 Ces\`aro 平均仍满足精确公式
\[
\mathscr K_A(z):=\sum_{b\ge 1}\kappa_A(b)z^b
=
\frac{\sum_{r=1}^{d_m}\kappa_A(r)z^r}{1-z^{d_m}},
\]
\[
\mu_A
:=
\lim_{X\to\infty}\frac{1}{X}\sum_{b\le X}\kappa_A(b)
=
\frac{1}{d_m}\sum_{r=1}^{d_m}\kappa_A(r).
\]
\end{conclusion}
\begin{proof}
结论 \ref{con:xi-time-part62d-smith-kappa-periodic-ogf-cesaro} 已证明 \(d_m\) 是 \(\kappa_A\) 的一个周期，并给出了上面两条周期分解公式。故只需证明不存在更小的正周期。

设 \(p\) 为 \(\kappa_A\) 的最小正周期。由于 \(d_m\) 也是周期，最小周期整除任意周期，故
\[
p\mid d_m.
\]
若 \(p<d_m\)，则
\[
\kappa_A(d_m)=\prod_{i=1}^{m}\gcd(d_i,d_m)=\prod_{i=1}^{m}d_i.
\]
另一方面，由 \(p\mid d_m\) 可得
\[
\gcd(d_m,d_m-p)=\gcd(d_m,p)=p<d_m,
\]
而对每个 \(i<m\) 都有
\[
\gcd(d_i,d_m-p)\le d_i.
\]
于是
\[
\kappa_A(d_m-p)
=
\prod_{i=1}^{m}\gcd(d_i,d_m-p)
<
\prod_{i=1}^{m}d_i
=
\kappa_A(d_m),
\]
这与 \(p\) 为周期所要求的
\[
\kappa_A(d_m-p)=\kappa_A(d_m)
\]
矛盾。故只能有
\[
p=d_m.
\]
因此 \(\kappa_A\) 的最小正周期恰为 \(d_m\)，而所示普通生成函数与 Ces\`aro 平均公式即由该周期分解立得。证毕。
\end{proof}

\begin{conclusion}[局部化数系之间的加法同态具有严格乘法刚性与完全分类]\label{con:xi-time-part62da-localization-hom-rigidity}
设 \(S,T\subset\mathbb P\) 为有限素数集，并记
\[
G_S:=\ZZ[S^{-1}],
\qquad
G_T:=\ZZ[T^{-1}],
\]
视为加法群。则有自然同构
\[
\operatorname{Hom}(G_S,G_T)\cong
\begin{cases}
G_T,& S\subseteq T,\\[1mm]
0,& S\not\subseteq T.
\end{cases}
\]
更具体地：
\begin{enumerate}
  \item 任意群同态 \(\phi:G_S\to G_T\) 都唯一写成
  \[
  \phi(x)=qx
  \qquad(x\in G_S)
  \]
  对某个 \(q\in\QQ\)；
  \item 非零同态存在当且仅当 \(S\subseteq T\)；
  \item 自同构群满足
  \[
  \operatorname{Aut}(G_S)
  \cong
  \Bigl\{\pm\prod_{p\in S}p^{n_p}: n_p\in\ZZ\Bigr\}
  \cong C_2\times \ZZ^{|S|}.
  \]
\end{enumerate}
\end{conclusion}
\begin{proof}
任取群同态 \(\phi:G_S\to G_T\)。由于 \(G_S\subset\QQ\) 是秩 \(1\) 的无挠加法群，\(\phi\) 被 \(\phi(1)\) 唯一决定。记
\[
q:=\phi(1)\in G_T\subset\QQ.
\]
对任意
\[
x=\frac{a}{s}\in G_S,
\qquad
s\in\ZZ_{>0},
\]
且 \(s\) 的全部素因子属于 \(S\)，都有 \(sx=a\)。施行 \(\phi\) 得
\[
s\,\phi(x)=a\,q.
\]
由于 \(G_T\subset\QQ\) 无挠，该方程在 \(G_T\) 中至多有一个解，而显然 \(qx\) 就是其在 \(\QQ\) 中的解，故
\[
\phi(x)=qx.
\]
这就证明了第一条。

若 \(S\subseteq T\)，则任取 \(q\in G_T\)，乘法映射
\[
x\longmapsto qx
\]
都把 \(G_S\) 送入 \(G_T\)，故得到 \(\operatorname{Hom}(G_S,G_T)\cong G_T\)。

若 \(S\not\subseteq T\)，取 \(p\in S\setminus T\)。若 \(\phi\neq 0\)，则上式中的 \(q=\phi(1)\neq 0\)。但对每个 \(n\ge 1\)，
\[
\phi(p^{-n})=q\,p^{-n}\in G_T.
\]
由于 \(p\notin T\)，群 \(G_T\) 中元素的分母不允许出现任意高次的 \(p\)，故对充分大的 \(n\)，\(q\,p^{-n}\notin G_T\)，矛盾。于是只能有 \(q=0\)，从而
\[
\operatorname{Hom}(G_S,G_T)=0.
\]
第二条得证。

最后，当 \(T=S\) 时，\(\phi(x)=qx\) 为自同构当且仅当 \(q,q^{-1}\in G_S\)。这恰等价于
\[
q=\pm\prod_{p\in S}p^{n_p},
\qquad
n_p\in\ZZ.
\]
于是得到所示自同构群公式。证毕。
\end{proof}

\begin{corollary}[Pontryagin 对偶上的连续满射序由素数集包含关系完全控制]\label{cor:xi-time-part62da-localization-pontryagin-surjection-poset}
沿用结论 \ref{con:xi-time-part62da-localization-hom-rigidity} 的记号。则以下命题等价：
\begin{enumerate}
  \item 存在连续满射
  \[
  \Pi_{T\to S}:\widehat{G_T}\twoheadrightarrow \widehat{G_S};
  \]
  \item \(S\subseteq T\)。
\end{enumerate}
并且在 \(S\subseteq T\) 时，由包含映射 \(G_S\hookrightarrow G_T\) 对偶化得到的规范满射 \(\Pi_{T\to S}\) 满足
\[
\ker(\Pi_{T\to S})\cong \prod_{p\in T\setminus S}\ZZ_p.
\]
\end{corollary}
\begin{proof}
若 \(S\subseteq T\)，则定理 \ref{thm:xi-cdim-localization-quotient-solenoid-surjections} 已给出规范短正合列
\[
0\longrightarrow \prod_{p\in T\setminus S}\ZZ_p
\longrightarrow \widehat{G_T}
\xrightarrow{\ \Pi_{T\to S}\ }
\widehat{G_S}
\longrightarrow 0,
\]
故连续满射存在，且其核正是
\[
\prod_{p\in T\setminus S}\ZZ_p.
\]

反之，若存在连续满射
\[
f:\widehat{G_T}\twoheadrightarrow \widehat{G_S},
\]
则由 Pontryagin 对偶的反变等价，\(f\) 对偶化为一个单射
\[
f^\vee:G_S\hookrightarrow G_T.
\]
该单射非零，故由结论 \ref{con:xi-time-part62da-localization-hom-rigidity}，必有
\[
S\subseteq T.
\]
证毕。
\end{proof}

\noindent
上述五条结果把地址链、Smith 周期链与局部化商链进一步压缩到三类有限证书：黄金柱质量只取决于前缀长度及末位，模核计数的全部周期结构只取决于最大 Smith 不变量，而 Pontryagin 对偶上的全部连续商序只取决于素数集的包含偏序。由此，分形测度、模层算术与紧 Abel 对偶三侧的无限结构便同时获得了严格可审计的有限闭式。

\paragraph{纤维同伦相图、局部化圆周覆盖与 bin-fold 饱和异常集}
在最大非可缩纤维的 Fibonacci 拓扑税、局部化数系对偶 solenoid 的 \(S\)-自由覆盖律以及 bin-fold 饱和阈值的 Baker 型有限性已经分别建立之后，还可把这三条链进一步压缩为下列五条直接后果。

\begin{theorem}[最大纤维的同伦相图具有精确的模 \texorpdfstring{$6$}{6} 分岔]\label{thm:xi-time-part62db-maxfiber-homotopy-phase-mod6}
沿用定理 \ref{thm:xi-time-part62d-max-noncontractible-fiber-gap-four-constants} 与推论 \ref{cor:xi-time-part62d-max-noncontractible-fiber-tax-mod6} 的记号。对 \(m\ge 17\)，记
\[
D_m:=\max_{x\in X_m}d_m(x),
\qquad
\widetilde D_m:=\max\{d_m(x):x\in X_m,\ |K_x|\not\simeq *\}.
\]
则最大纤维本身非可缩，当且仅当
\[
D_m=\widetilde D_m
\iff
m\equiv 0,4\pmod 6.
\]
等价地，
\[
|K_x|\not\simeq *
\quad\text{对某个满足 }d_m(x)=D_m\text{ 的极大纤维成立}
\iff
m\equiv 0,4\pmod 6.
\]
而在其余四个剩余类中，最大纤维必可缩，且其与最大非可缩纤维之间的精确缺口满足
\[
D_m-\widetilde D_m=
\begin{cases}
F_{(m-9)/2},& m\equiv 1,5\pmod 6,\\[1mm]
F_{m/2-3},& m\equiv 2\pmod 6,\\[1mm]
F_{(m-9)/2}+F_{(m-17)/2},& m\equiv 3\pmod 6.
\end{cases}
\]
\end{theorem}
\begin{proof}
由推论 \ref{cor:xi-time-part62d-max-noncontractible-fiber-tax-mod6}，
\[
\Delta_m:=D_m-\widetilde D_m
\]
满足
\[
\Delta_m=
\begin{cases}
0,& m\equiv 0,4\pmod 6,\\[1mm]
F_{(m-9)/2},& m\equiv 1,5\pmod 6,\\[1mm]
F_{m/2-3},& m\equiv 2\pmod 6,\\[1mm]
F_{(m-9)/2}+F_{(m-17)/2},& m\equiv 3\pmod 6.
\end{cases}
\]
因此
\[
D_m=\widetilde D_m
\iff
\Delta_m=0
\iff
m\equiv 0,4\pmod 6.
\]
而 \(\widetilde D_m\) 的定义正是全部非可缩纤维大小中的最大值，故
\[
D_m=\widetilde D_m
\]
当且仅当某个极大纤维已经非可缩；若 \(\widetilde D_m<D_m\)，则所有大小达到 \(D_m\) 的纤维都只能是可缩的。于是得到所述精确分岔与显式缺口公式。证毕。
\end{proof}

\begin{corollary}[最大非可缩纤维的拓扑保真比具有唯一最坏相]\label{cor:xi-time-part62db-maxfiber-fidelity-worst-phase}
沿用定理 \ref{thm:xi-time-part62db-maxfiber-homotopy-phase-mod6} 的记号，并定义
\[
R_m:=\frac{\widetilde D_m}{D_m}.
\]
则沿模 \(6\) 子序列分别有极限
\[
\lim_{\substack{m\to\infty\\ m\equiv 0,4\!\!\!\!\pmod 6}}R_m=1,
\]
\[
\lim_{\substack{m\to\infty\\ m\equiv 1,5\!\!\!\!\pmod 6}}R_m
=
1-\frac{1}{2\varphi^5},
\]
\[
\lim_{\substack{m\to\infty\\ m\equiv 2\!\!\!\!\pmod 6}}R_m
=
1-\varphi^{-5},
\]
\[
\lim_{\substack{m\to\infty\\ m\equiv 3\!\!\!\!\pmod 6}}R_m
=
1-\frac{\varphi^{-5}+\varphi^{-9}}{2}.
\]
并且这四个极限常数满足严格排序
\[
1-\varphi^{-5}
<
1-\frac{\varphi^{-5}+\varphi^{-9}}{2}
<
1-\frac{1}{2\varphi^5}
<
1.
\]
因而模 \(6\) 的唯一最坏拓扑相恰出现在
\[
m\equiv 2\pmod 6.
\]
\end{corollary}
\begin{proof}
推论 \ref{cor:xi-time-part62d-max-noncontractible-fiber-tax-mod6} 已给出四条极限公式。只需比较它们即可。

由于
\[
0<\varphi^{-9}<\varphi^{-5},
\]
故
\[
\varphi^{-5}
>
\frac{\varphi^{-5}+\varphi^{-9}}{2}
>
\frac12\varphi^{-5}
>
0.
\]
两侧同时取 \(1-\) 即得
\[
1-\varphi^{-5}
<
1-\frac{\varphi^{-5}+\varphi^{-9}}{2}
<
1-\frac{1}{2\varphi^5}
<
1.
\]
因此四个极限值中最小者唯一是
\[
1-\varphi^{-5},
\]
它对应的正是 \(m\equiv 2\pmod 6\) 子序列。证毕。
\end{proof}

\begin{theorem}[局部化数系对偶 \texorpdfstring{$\widehat{G_S}$}{G-hat-S} 上保圆周商的自同态由整数参数唯一分类]\label{thm:xi-time-part62db-localized-solenoid-circle-quotient-lifts}
沿用定理 \ref{thm:killo-localization-circle-dimension-prime-ledger-splitting} 与定理 \ref{thm:xi-localized-basic-extension-strictly-nonsplit} 的记号。设 \(S\subset\mathbb P\) 为有限素数集，记
\[
G_S:=\ZZ[S^{-1}],
\qquad
\Sigma_S:=\widehat{G_S},
\]
并写 Pontryagin 对偶短正合列为
\[
0\longrightarrow \prod_{p\in S}\ZZ_p
\longrightarrow \Sigma_S
\xrightarrow{\ \pi_S\ }
\TT
\longrightarrow 0.
\]
对任意整数 \(n\neq 0\)，把 \(|n|\) 分解为
\[
|n|=n_S\,n_{S^\perp},
\qquad
n_S:=\prod_{p\in S}p^{v_p(n)},
\qquad
n_{S^\perp}:=\prod_{p\notin S}p^{v_p(n)}.
\]
则存在唯一连续群同态
\[
\Phi_{S,n}:\Sigma_S\to \Sigma_S
\]
满足
\[
\pi_S\circ \Phi_{S,n}=[n]\circ \pi_S,
\]
其中 \([n]:\TT\to\TT\) 为 \(n\) 倍覆盖。并且
\[
\ker(\Phi_{S,n})\cong \ZZ/n_{S^\perp}\ZZ,
\qquad
\deg(\Phi_{S,n})=n_{S^\perp}.
\]
换言之，\(S\)-内素因子在 \(\Sigma_S\) 上全部表现为可逆方向，只有 \(S\)-外素因子真正贡献有限覆盖核。
\end{theorem}
\begin{proof}
由 Pontryagin 对偶的反变等价，条件
\[
\pi_S\circ \Phi_{S,n}=[n]\circ \pi_S
\]
对偶化后等价于存在加法群同态
\[
\phi_{S,n}:G_S\to G_S
\]
使得它在圆周商对偶 \(\ZZ\hookrightarrow G_S\) 上满足
\[
\phi_{S,n}(1)=n.
\]
而由结论 \ref{con:xi-time-part62da-localization-hom-rigidity}，任意加法群同态 \(G_S\to G_S\) 都唯一是“乘以某个有理数”的映射；既然它在 \(1\) 上取值为 \(n\)，便只能是
\[
\phi_{S,n}(x)=nx.
\]
因此对应的连续 lift \(\Phi_{S,n}\) 存在且唯一。

再由定理 \ref{thm:xi-localized-solenoid-endomorphism-kernel-cokernel-sfree-numerator} 对 \(x=n\) 的特例，
\[
\ker(\Phi_{S,n})\cong \ZZ/|n|_{S^\perp}\ZZ=\ZZ/n_{S^\perp}\ZZ.
\]
同一定理同时表明 \(\Phi_{S,n}\) 为满同态。故它是有限核的连续满射，而其覆盖次数正等于核的基数，于是
\[
\deg(\Phi_{S,n})=\#\ker(\Phi_{S,n})=n_{S^\perp}.
\]
证毕。
\end{proof}

\begin{corollary}[保圆周商 lift 的刚性与可逆素判别]\label{cor:xi-time-part62db-localized-solenoid-circle-quotient-rigidity}
沿用定理 \ref{thm:xi-time-part62db-localized-solenoid-circle-quotient-lifts} 的记号。则成立：
\begin{enumerate}
  \item 若连续群同态 \(\Phi:\Sigma_S\to\Sigma_S\) 满足
  \[
  \pi_S\circ \Phi=\pi_S,
  \]
  则
  \[
  \Phi=\mathrm{id}_{\Sigma_S}.
  \]
  \item \(\Phi_{S,n}\) 是拓扑自同构，当且仅当
  \[
  n_{S^\perp}=1,
  \]
  等价地，\(n\) 的全部素因子都落在 \(S\) 中。
  \item 因而，全部保圆周商的连续自同构恰对应于整数半群
  \[
  \left\{
  \pm\prod_{p\in S}p^{a_p}: a_p\in\ZZ_{\ge 0}
  \right\},
  \]
  而对每个 \(p\notin S\)，\(\Phi_{S,p}\) 都是严格的 \(p\)-重覆盖而非自同构。
\end{enumerate}
\end{corollary}
\begin{proof}
若 \(\pi_S\circ \Phi=\pi_S\)，则它正是定理 \ref{thm:xi-time-part62db-localized-solenoid-circle-quotient-lifts} 中 \(n=1\) 的情形。由唯一性，必有
\[
\Phi=\Phi_{S,1}=\mathrm{id}_{\Sigma_S}.
\]

第二条由同一定理中的核公式立即得到：
\[
\Phi_{S,n}\ \text{为自同构}
\iff
\ker(\Phi_{S,n})=0
\iff
n_{S^\perp}=1.
\]
而
\[
n_{S^\perp}=1
\]
恰等价于 \(n\) 的素因子全部属于 \(S\)。

第三条只是把这一条件改写为显式的整数半群形式。若 \(p\notin S\)，则
\[
p_{S^\perp}=p,
\]
故
\[
\ker(\Phi_{S,p})\cong \ZZ/p\ZZ,
\qquad
\deg(\Phi_{S,p})=p,
\]
于是它是严格的 \(p\)-重覆盖而不是自同构。证毕。
\end{proof}

\begin{proposition}[bin-fold 组合上界饱和的有限性与已审计完全列表]\label{prop:xi-time-part62db-binfold-saturation-audited-finite-list}
沿用定义 \ref{def:K-of-m} 与推论 \ref{cor:xi-foldbin-maxfiber-saturation-finite-occurrence} 的记号。对每个 \(m\ge 1\)，令 \(K(m)\) 满足
\[
F_{K(m)+1}\le 2^m-1<F_{K(m)+2},
\]
并记
\[
L(m):=K(m)-m,
\qquad
d_{\max}^{\mathrm{bin}}(m):=\max_{w\in X_m}d_m^{\mathrm{bin}}(w).
\]
则最大纤维达到组合上界
\[
d_{\max}^{\mathrm{bin}}(m)=F_{L(m)+2}
\]
只能发生有限次。并且在已审计窗口 \(m\le 400\) 内，恰有
\[
d_{\max}^{\mathrm{bin}}(m)=F_{L(m)+2}
\iff
m\in\{1,2,3,4,5,7,12\}.
\]
因此在整个区间
\[
13\le m\le 400
\]
上，严格不等式
\[
d_{\max}^{\mathrm{bin}}(m)<F_{L(m)+2}
\]
一律成立。
\end{proposition}
\begin{proof}
推论 \ref{cor:xi-foldbin-maxfiber-saturation-finite-occurrence} 已经给出两件事：其一，结合注 \ref{rem:fold-bin-finite} 的 Baker--W\"ustholz/Matveev 型线性对数形式下界，上界饱和
\[
d_{\max}^{\mathrm{bin}}(m)=F_{L(m)+2}
\]
只能对有限多个 \(m\) 发生；其二，在可复核扫描窗口 \(m\le 400\) 内，全部饱和点恰由表 \ref{tab:fold_binary_saturation_points} 审计为
\[
\{1,2,3,4,5,7,12\}.
\]
故在该窗口内，上述等式成立当且仅当 \(m\) 属于该集合。于是其补集 \(13\le m\le 400\) 上自动满足严格不等式
\[
d_{\max}^{\mathrm{bin}}(m)<F_{L(m)+2}.
\]
证毕。
\end{proof}

\noindent
综上，最大非可缩纤维的相图、局部化数系对偶 \(\Sigma_S=\widehat{G_S}\) 的保圆周商覆盖分类，以及 bin-fold 组合上界饱和的有限异常集，都被压缩为有限可核验的刚性证书：前者只由模 \(6\) 剩余类决定，中者只由 \(S\)-外素因子决定，后者则在 Baker 型有限性与有限窗审计的结合下收缩为一个显式有限集合。

\paragraph{全息地址动力学的闭式、Smith 剖面的全局实现与 window-\texorpdfstring{$6$}{6} 的余根对偶}
在 H.160--H.161 的 \(2^L\)-进全息地址系统、H.154--H.156 的 Smith 剖面与局部 Euler 因子链，以及附录 F.1 的 window-\(6\) 循环核根字典都已建立之后，原先分散的三条线路还可进一步压缩为若干条终端结论：地址像侧获得自然动力学与 Bernoulli 推前测度的闭式，Smith 剖面侧获得有限素数支撑下的全局实现与局部 Euler 因子的完全判据，而 window-\(6\) 的 cyclic 核则在逐点层面严格实现 \(B_3\) 与 \(C_3\) 的 coroot 对偶。

\begin{theorem}[\texorpdfstring{$(2^L)$}{(2^L)}-进全息像集上的自然动力学与一侧满移位严格共轭]\label{thm:xi-time-part62de-hologram-natural-dynamics-fullshift}
沿用定理 \ref{thm:xi-time-part62dc-hologram-selfaffine-cantor-dynamics} 的记号，并记
\[
\Lambda_{E,v}:=X(E^{\NN}),
\qquad
\mathcal F:=S_{E,v}:\Lambda_{E,v}\to \Lambda_{E,v}.
\]
则成立：
\begin{enumerate}
  \item \(\mathcal F\) 良定义，且与一侧满移位 \(\sigma:E^{\NN}\to E^{\NN}\) 满足严格共轭关系
  \[
  \mathcal F\circ X=X\circ \sigma;
  \]
  \item \(\Lambda_{E,v}\) 具有严格不交并分解
  \[
  \Lambda_{E,v}
  =
  \bigsqcup_{a\in \mathrm{code}(E)}\bigl(av+B\Lambda_{E,v}\bigr);
  \]
  \item 其拓扑熵、周期点计数与 Artin--Mazur \(\zeta\) 函数满足
  \[
  h_{\mathrm{top}}(\mathcal F)=\log |E|,
  \qquad
  \#\operatorname{Fix}(\mathcal F^n)=|E|^n
  \quad (n\ge 1),
  \]
  \[
  \zeta_{\mathcal F}(z)
  :=
  \exp\!\left(\sum_{n\ge 1}\frac{\#\operatorname{Fix}(\mathcal F^n)}{n}z^n\right)
  =
  \frac{1}{1-|E|z}.
  \]
\end{enumerate}
\end{theorem}
\begin{proof}
定理 \ref{thm:xi-time-part62dc-hologram-selfaffine-cantor-dynamics} 已给出 \(\mathcal F=S_{E,v}\) 的良定义性以及
\[
X\circ \sigma=\mathcal F\circ X.
\]
这就得到第 \(1\) 条。又同一定理证明了
\[
\Lambda_{E,v}
=
\bigsqcup_{e\in E}\bigl(\mathrm{code}(e)v+B\Lambda_{E,v}\bigr),
\]
而 \(\mathrm{code}:E\hookrightarrow \{0,\dots,M\}\) 为单射，故也可重写为按 \(\mathrm{code}(E)\) 标号的不交并分解，第 \(2\) 条成立。

由 \(X:\,E^{\NN}\to \Lambda_{E,v}\) 为拓扑同胚且 \(X\circ \sigma=\mathcal F\circ X\)，可知 \((\Lambda_{E,v},\mathcal F)\) 与一侧满移位 \((E^{\NN},\sigma)\) 拓扑共轭。于是
\[
h_{\mathrm{top}}(\mathcal F)=h_{\mathrm{top}}(\sigma)=\log |E|.
\]
另一方面，\(\sigma^n\) 的不动点恰为所有 \(n\)-周期重复词
\[
(e_1,\dots,e_n,e_1,\dots,e_n,\dots),
\]
故
\[
\#\operatorname{Fix}(\sigma^n)=|E|^n.
\]
共轭保持不动点计数，遂有
\[
\#\operatorname{Fix}(\mathcal F^n)=|E|^n.
\]
将其代入 Artin--Mazur \(\zeta\) 函数定义，得到
\[
\zeta_{\mathcal F}(z)
=
\exp\!\left(\sum_{n\ge 1}\frac{|E|^n}{n}z^n\right)
=
\exp\!\bigl(-\log(1-|E|z)\bigr)
=
\frac{1}{1-|E|z}.
\]
证毕。
\end{proof}

\begin{theorem}[\texorpdfstring{$(2^L)$}{(2^L)}-进全息像上的 Bernoulli 推前测度在环境模空间中仍为精确维度]\label{thm:xi-time-part62de-hologram-ambient-bernoulli-exact-dimensional}
沿用定理 \ref{thm:xi-time-part62de-hologram-natural-dynamics-fullshift} 的记号。设
\[
\mathbf p=(p_e)_{e\in E},
\qquad
p_e>0,
\qquad
\sum_{e\in E}p_e=1,
\]
\[
\mathbb P_{\mathbf p}
:=
\bigotimes_{n\ge 1}\left(\sum_{e\in E}p_e\delta_e\right),
\qquad
\mu_{\mathbf p}:=X_*\mathbb P_{\mathbf p}.
\]
在环境模空间 \(T=T_2(E_{\min})\) 上定义
\[
\nu_B(z):=\max\{n\ge 0:\ z\in B^nT\}
\qquad (z\neq 0),
\]
并约定 \(\nu_B(0)=+\infty\)；再定义 \(B\)-进超度量
\[
d_B(x,y):=2^{-L\nu_B(x-y)}.
\]
则 \(\mu_{\mathbf p}\) 为 exact-dimensional 测度，并且对 \(\mu_{\mathbf p}\)-几乎处处的 \(x\in \Lambda_{E,v}\) 都有
\[
\lim_{r\downarrow 0}
\frac{\log \mu_{\mathbf p}\!\bigl(\overline{\mathrm B}_{d_B}(x,r)\bigr)}{\log r}
=
\frac{H(\mathbf p)}{L\log 2},
\qquad
H(\mathbf p):=-\sum_{e\in E}p_e\log p_e.
\]
特别地，当 \(\mathbf p\) 为均匀分布时，
\[
\dim_{\mathrm H}(\mu_{\mathrm{unif}})
=
\frac{\log |E|}{L\log 2}.
\]
\end{theorem}
\begin{proof}
记
\[
\lambda_v:=\max\{m\ge 0:\ v\in B^mT\},
\]
则 \(\lambda_v<\infty\)，并可唯一写成
\[
v=B^{\lambda_v}u,
\qquad
u\in T\setminus BT.
\]
对 \(\mathbf e=(e_1,e_2,\dots)\in E^{\NN}\) 定义
\[
X_u(\mathbf e):=\sum_{t\ge 1}\mathrm{code}(e_t)B^{t-1}u.
\]
则
\[
X(\mathbf e)=B^{\lambda_v}X_u(\mathbf e).
\]
由定理 \ref{thm:xi-time-part9g-holographic-prefix-concatenation-ultrametric} 作用于向量 \(u\)，若两条序列 \(\mathbf e,\mathbf e'\) 的最长公共前缀长度恰为 \(n\)，则
\[
X_u(\mathbf e)-X_u(\mathbf e')
\in
B^nT\setminus B^{n+1}T.
\]
乘上 \(B^{\lambda_v}\) 后得到
\[
X(\mathbf e)-X(\mathbf e')
\in
B^{n+\lambda_v}T\setminus B^{n+1+\lambda_v}T.
\]
因此，对任意
\[
x=X(\mathbf e)\in \Lambda_{E,v}
\]
及其长度为 \(n\) 的前缀 cylinder
\[
[e_1\cdots e_n]
:=
\{\mathbf f\in E^{\NN}:\ f_1=e_1,\dots,f_n=e_n\},
\]
都有
\[
X([e_1\cdots e_n])
=
\overline{\mathrm B}_{d_B}\!\bigl(x,2^{-L(n+\lambda_v)}\bigr)\cap \Lambda_{E,v}.
\]
由于 \(\mu_{\mathbf p}\) 支持在 \(\Lambda_{E,v}\) 上，故
\[
\mu_{\mathbf p}\!\bigl(\overline{\mathrm B}_{d_B}(x,2^{-L(n+\lambda_v)})\bigr)
=
\mathbb P_{\mathbf p}([e_1\cdots e_n])
=
p_{e_1}\cdots p_{e_n}.
\]
于是
\[
\frac{
\log \mu_{\mathbf p}\!\bigl(\overline{\mathrm B}_{d_B}(x,2^{-L(n+\lambda_v)})\bigr)
}{
\log 2^{-L(n+\lambda_v)}
}
=
\frac{
\sum_{j=1}^{n}\log p_{e_j}
}{
-(n+\lambda_v)L\log 2
}.
\]
对 \(\mathbb P_{\mathbf p}\)-几乎处处的 \(\mathbf e\)，强大数定律给出
\[
\frac1n\sum_{j=1}^{n}(-\log p_{e_j})\to H(\mathbf p),
\]
而
\[
\frac{n}{n+\lambda_v}\to 1,
\]
故对 \(\mu_{\mathbf p}\)-几乎处处的 \(x=X(\mathbf e)\) 有
\[
\lim_{n\to\infty}
\frac{
\log \mu_{\mathbf p}\!\bigl(\overline{\mathrm B}_{d_B}(x,2^{-L(n+\lambda_v)})\bigr)
}{
\log 2^{-L(n+\lambda_v)}
}
=
\frac{H(\mathbf p)}{L\log 2}.
\]

又由于 \(d_B\) 的非零取值只可能是 \(\{2^{-Lm}:m\ge 0\}\)，故任意 \(r>0\) 足够小时都落在某个离散尺度区间
\[
2^{-L(m+1)}<r\le 2^{-Lm},
\]
并满足
\[
\overline{\mathrm B}_{d_B}(x,r)=\overline{\mathrm B}_{d_B}(x,2^{-Lm}).
\]
因此 \(r\downarrow 0\) 的极限与沿离散尺度 \(2^{-Lm}\) 的极限相同，遂得 exact-dimensional 性与所示局部维数公式。

若 \(\mathbf p\) 为均匀分布，则
\[
H(\mathbf p)=\log |E|,
\]
代回即可得到
\[
\dim_{\mathrm H}(\mu_{\mathrm{unif}})
=
\frac{\log |E|}{L\log 2}.
\]
证毕。
\end{proof}

\begin{theorem}[Smith 局部剖面族的全局可实现性]\label{thm:xi-time-part62de-smith-profile-global-realization}
\leanverified{paper\_xi\_time\_part62de\_smith\_profile\_global\_realization}
对每个素数 \(p\)，设
\[
f_p:\ZZ_{\ge 0}\to \ZZ_{\ge 0}
\]
满足 \(f_p(0)=0\)，并定义一阶差分
\[
\Delta_p(k):=f_p(k)-f_p(k-1)
\qquad (k\ge 1).
\]
若对每个 \(p\) 都有 \((\Delta_p(k))_{k\ge 1}\) 是非负、非增且最终为零的整数列，并且只有有限多个素数 \(p\) 满足 \(f_p\not\equiv 0\)，则存在某个满秩整数矩阵 \(A\)。更具体地，记
\[
m_0:=\max_{p}\Delta_p(1),
\qquad
m:=\max\{1,m_0\},
\]
则可取某个最小尺寸的方阵对角矩阵
\[
A=\operatorname{diag}(d_1,\dots,d_m)\in M_m(\ZZ),
\qquad
d_1\mid \cdots\mid d_m,
\]
使得对所有素数 \(p\) 与所有 \(k\ge 1\) 都有
\[
f_p(k)
=
\sum_{i=1}^{m}\min\bigl(k,v_p(d_i)\bigr)
=
\log_p K_p(A;k),
\]
其中
\[
K_p(A;k):=\#\ker(A\bmod p^k).
\]
此外，在这一最小尺寸口径下，族 \(\{f_p\}_p\) 唯一决定 \(A\) 的 Smith 不变量 \(d_1,\dots,d_m\)。
\end{theorem}
\begin{proof}
若对所有素数 \(p\) 都有 \(f_p\equiv 0\)，则取
\[
A=(1)\in M_1(\ZZ)
\]
即可；此时 \(m_0=0\)、\(m=1\)，且唯一的最小 Smith 不变量为 \(d_1=1\)。以下设至少存在一个素数使 \(f_p\not\equiv 0\)，于是 \(m=m_0\ge 1\)。

对每个固定的 \(p\)，由于
\[
\Delta_p(k)\in\{0,1,\dots,m\},
\]
且其为非增并最终为零，定理 \ref{thm:xi-smith-loss-spectrum-conjugate-partition-realization} 作用于行秩 \(m\) 的情形可得一个唯一的非增整数列
\[
\widetilde e_1(p)\ge \widetilde e_2(p)\ge \cdots \ge \widetilde e_m(p)\ge 0
\]
使得
\[
f_p(k)=\sum_{i=1}^{m}\min\bigl(k,\widetilde e_i(p)\bigr).
\]
把它反向重排为
\[
0\le e_1(p)\le e_2(p)\le \cdots \le e_m(p),
\]
则同一公式仍成立：
\[
f_p(k)=\sum_{i=1}^{m}\min\bigl(k,e_i(p)\bigr).
\]
又因只有有限多个 \(p\) 非平凡，故对每个 \(1\le i\le m\) 可以定义有限乘积
\[
d_i:=\prod_{p}p^{e_i(p)}\in \ZZ_{>0}.
\]
由 \(e_1(p)\le \cdots \le e_m(p)\) 对每个 \(p\) 都成立，立即得到
\[
d_1\mid d_2\mid \cdots \mid d_m.
\]
取
\[
A:=\operatorname{diag}(d_1,\dots,d_m).
\]
则 \(A\) 满秩，且其 Smith 不变量正是 \(d_1,\dots,d_m\)。于是由 Smith 模核公式，
\[
K_p(A;k)
=
\prod_{i=1}^{m}p^{\min(k,v_p(d_i))}
=
p^{\sum_{i=1}^{m}\min(k,e_i(p))}
=
p^{f_p(k)}.
\]
取 \(\log_p\) 便得所求公式。

唯一性方面，对每个素数 \(p\)，定理 \ref{thm:xi-smith-loss-spectrum-conjugate-partition-realization} 已说明 \(f_p\) 在固定长度 \(m\) 的口径下唯一决定整列 \((e_i(p))_{i=1}^{m}\)。把这些逐素数赋值重新拼合，便唯一确定每个 \(d_i\)。由于此处 \(m\) 已取最小可能值，故最小实现的 Smith 不变量唯一。证毕。
\end{proof}

\begin{corollary}[局部 Euler 因子的完全可实现性判据]\label{cor:xi-time-part62de-local-euler-factor-realization-criterion}
对每个素数 \(p\)，设
\[
\Xi_p(T)=\sum_{k\ge 0}a_{p,k}T^k
\in \ZZ[[T]].
\]
则族 \(\{\Xi_p\}_p\) 来自某个满秩整数矩阵 \(A\) 的局部因子
\[
\Xi_{A,p}(T):=\sum_{k\ge 0}\#\ker(A\bmod p^k)\,p^{-k(n-m)}T^k
\]
当且仅当满足下列条件：
\begin{enumerate}
  \item 对每个 \(p\) 都有 \(a_{p,0}=1\)；
  \item 对每个 \(p\) 与每个 \(k\ge 0\)，\(a_{p,k}\) 都是 \(p\) 的非负整数幂；
  \item 若定义
  \[
  f_p(k):=\log_p a_{p,k},
  \]
  则 \(f_p(0)=0\)，且
  \[
  \Delta_p(k):=f_p(k)-f_p(k-1)\qquad (k\ge 1)
  \]
  为非负、非增且最终为零的整数列；
  \item 只有有限多个素数 \(p\) 满足 \(f_p\not\equiv 0\)，等价地只有有限多个 \(p\) 满足
  \[
  \Xi_p(T)\neq \frac{1}{1-T}.
  \]
\end{enumerate}
在这些条件成立时，实现它们的最小 Smith 不变量唯一。
\end{corollary}
\begin{proof}
必要性来自局部模核公式。若 \(\Xi_p=\Xi_{A,p}\) 来自某个满秩整数矩阵 \(A\)，则由定理 \ref{thm:xi-smith-local-euler-complete-invariant} 及其先前公式，
\[
\Xi_{A,p}(T)=\sum_{k\ge 0}p^{f_{A,p}(k)}T^k,
\]
其中 \(f_{A,p}(0)=0\)，且其差分序列
\[
f_{A,p}(k)-f_{A,p}(k-1)=\#\{i:\ e_i(p)\ge k\}
\]
自动非负、非增并最终为零。又只有有限多个素数整除 Smith 不变量的乘积，因此只有有限多个 \(p\) 使 \(f_{A,p}\not\equiv 0\)；对其余素数有 \(a_{p,k}=1\)，故 \(\Xi_p(T)=(1-T)^{-1}\)。

反过来，若上述四条成立，则由定理 \ref{thm:xi-time-part62de-smith-profile-global-realization}，族 \(\{f_p\}_p\) 可由某个满秩整数矩阵 \(A\) 的 Smith 赋值谱全局实现。该构造给出的 \(A\) 为方阵，故 \(n=m\)。于是对每个 \(p\) 与 \(k\ge 0\)，有
\[
\#\ker(A\bmod p^k)\,p^{-k(n-m)}=p^{f_p(k)}=a_{p,k},
\]
从而
\[
\Xi_{A,p}(T)=\sum_{k\ge 0}a_{p,k}T^k=\Xi_p(T).
\]
唯一性则由同一定理中最小 Smith 不变量的唯一决定性立即推出。证毕。
\end{proof}

\begin{theorem}[window-\texorpdfstring{$6$}{6} 循环核同时实现 \texorpdfstring{$B_3$}{B3} 根系与 \texorpdfstring{$C_3$}{C3} 余根系]\label{thm:xi-time-part62de-window6-b3c3-coroot-duality}
沿用表 \ref{tab:fold6_b3c3_root_dictionary_b3} 与表 \ref{tab:fold6_b3c3_root_dictionary_c3} 的记号。定义两张字典
\[
\alpha_B:X_6^{\mathrm{cyc}}\to \Phi(B_3),
\qquad
\alpha_C:X_6^{\mathrm{cyc}}\to \Phi(C_3),
\]
其中
\[
\Phi(B_3)=\{\pm e_i,\ \pm e_i\pm e_j:\ 1\le i<j\le 3\},
\]
\[
\Phi(C_3)=\{\pm 2e_i,\ \pm e_i\pm e_j:\ 1\le i<j\le 3\},
\]
并取标准 Euclidean 内积 \((\cdot,\cdot)\)。
则对每个 \(w\in X_6^{\mathrm{cyc}}\) 都有严格公式
\[
\alpha_C(w)
=
\alpha_B(w)^\vee
:=
\frac{2\,\alpha_B(w)}{(\alpha_B(w),\alpha_B(w))}.
\]
因此：
\begin{enumerate}
  \item \(X_6^{\mathrm{cyc}}\) 同时实现了 \(B_3\) 根系与其 Langlands 对偶 \(C_3=B_3^\vee\) 的余根系；
  \item 六个 weight-\(1\) 元在 \(B_3\) 规范下恰对应六条短根 \(\{\pm e_1,\pm e_2,\pm e_3\}\)，在 coroot 变换下正好变为 \(C_3\) 的六条长根 \(\{\pm 2e_1,\pm 2e_2,\pm 2e_3\}\)；
  \item 其余十二个元素在 \(B_3\) 规范下恰对应长根 \(\{\pm e_i\pm e_j\}\)，其平方范数为 \(2\)，故 coroot 变换保持它们不变，并恰给出 \(C_3\) 的十二条短根。
\end{enumerate}
\end{theorem}
\begin{proof}
由表 \ref{tab:fold6_b3c3_root_dictionary_b3} 可直接读出，六个 weight-\(1\) 元恰被送到
\[
\{\pm e_1,\pm e_2,\pm e_3\},
\]
而由表 \ref{tab:fold6_b3c3_root_dictionary_c3}，同一批元素恰被送到
\[
\{\pm 2e_1,\pm 2e_2,\pm 2e_3\}.
\]
对这六个向量，\((\alpha_B(w),\alpha_B(w))=1\)，故
\[
\frac{2\,\alpha_B(w)}{(\alpha_B(w),\alpha_B(w))}=2\alpha_B(w)=\alpha_C(w).
\]

再看其余十二个元素。两张表都把它们送到同一组向量
\[
\{\pm e_i\pm e_j:\ 1\le i<j\le 3\}.
\]
对每个此类向量都有
\[
(\alpha_B(w),\alpha_B(w))=2,
\]
从而
\[
\frac{2\,\alpha_B(w)}{(\alpha_B(w),\alpha_B(w))}
=
\alpha_B(w)
=
\alpha_C(w).
\]
故 coroot 公式对所有 \(w\in X_6^{\mathrm{cyc}}\) 逐点成立。

而 \(B_3\) 的余根系按定义正是
\[
\Phi(B_3^\vee)=\Phi(C_3),
\]
所列三条后果遂全部成立。证毕。
\end{proof}

\noindent
由此，H.160--H.161 的全息地址链、H.154--H.156 的 Smith--Euler 链与附录 F.1 的 window-\(6\) 根字典便在同一结论层面完成闭合：自由符号流可被 \(2^L\)-进几何严格承载，局部 valuation profile 的全部有限支撑数据可被全局 Smith 正规形完全吸收，而 \(X_6^{\mathrm{cyc}}\) 则在同一组离散状态上同时实现根与余根的 Langlands 对偶。

\paragraph{全息地址系统的严格仿射满移位模型与 primitive 周期闭式}
在 H.160--H.161 的 faithful \(2^L\)-进地址实现、像集同胚、严格 cylinder 分解与自然动力学闭式都已确立之后，还可把该地址系统进一步压缩为一条更直接的动力学表述：地址级数满足精确的首位递推，其模 \(B^n\) 余类严格记录有限前缀，而地址像本身就是一侧满移位的严格仿射模型；相应 entropy、周期点、primitive 周期轨与 Artin--Mazur \(\zeta\) 函数全部由字母表基数单独决定。

\begin{theorem}[H.160--H.161 原始记号下全息地址级数的首位递推、自仿射分解与一侧满移位共轭]\label{thm:xi-time-part62deb-hologram-affine-fullshift-conjugacy}
\leanverified{paper\_xi\_time\_part62deb\_hologram\_affine\_fullshift\_conjugacy}
沿用定理 \ref{thm:xi-time-part62d-hologram-image-cantor-homeomorphism} 的记号。设 \(E\) 为有限事件集，
\[
\mathrm{code}:E\hookrightarrow \{0,\dots,M\}\subset \ZZ
\]
为单射，
\[
B:=2^L>M,
\qquad
\mathcal X_E:=X(E^{\NN})\subset T_2(E_{\min}).
\]
记 \(\sigma:E^{\NN}\to E^{\NN}\) 为左移。则对任意
\[
\mathbf e=(e_1,e_2,\dots)\in E^{\NN}
\]
都有首位递推
\[
X(\mathbf e)=\mathrm{code}(e_1)v+B\,X(\sigma\mathbf e).
\]
从而
\[
\mathcal X_E
=
\bigsqcup_{a\in E}\bigl(\mathrm{code}(a)v+B\mathcal X_E\bigr).
\]
对每个 \(a\in E\)，在分支片
\[
\mathrm{code}(a)v+B\mathcal X_E
\]
上定义仿射映射
\[
\mathcal T(x):=\frac{x-\mathrm{code}(a)v}{B}.
\]
则 \(\mathcal T:\mathcal X_E\to \mathcal X_E\) 良定义，并满足严格共轭关系
\[
\mathcal T\circ X=X\circ \sigma.
\]
特别地，\((\mathcal X_E,\mathcal T)\) 与 \((E^{\NN},\sigma)\) 拓扑共轭。
\end{theorem}
\begin{proof}
若 \(\mathbf e=(e_t)_{t\ge 1}\in E^{\NN}\)，则
\[
X(\mathbf e)
=
\mathrm{code}(e_1)v
+B\sum_{t\ge 1}\mathrm{code}(e_{t+1})B^{t-1}v
=
\mathrm{code}(e_1)v+B\,X(\sigma\mathbf e).
\]
这就证明了首位递推，且立即得到
\[
\mathcal X_E\subseteq \bigcup_{a\in E}\bigl(\mathrm{code}(a)v+B\mathcal X_E\bigr).
\]
反向包含显然成立，故并式成立。

若存在 \(a\neq a'\in E\) 使
\[
\bigl(\mathrm{code}(a)v+B\mathcal X_E\bigr)\cap
\bigl(\mathrm{code}(a')v+B\mathcal X_E\bigr)\neq\varnothing,
\]
则可取 \(\mathbf e,\mathbf f\in E^{\NN}\) 满足 \(e_1=a\)、\(f_1=a'\) 且 \(X(\mathbf e)=X(\mathbf f)\)，这与定理 \ref{thm:xi-time-part62d-hologram-image-cantor-homeomorphism} 给出的 \(X\) 的单射性矛盾。故该并为不交并。

因此对每个 \(x\in \mathcal X_E\) 恰有唯一 \(a\in E\) 使
\[
x\in \mathrm{code}(a)v+B\mathcal X_E,
\]
从而 \(\mathcal T\) 良定义，且其像仍落在 \(\mathcal X_E\)。

对任意 \(\mathbf e\in E^{\NN}\)，由上式
\[
\mathcal T(X(\mathbf e))
=
\frac{X(\mathbf e)-\mathrm{code}(e_1)v}{B}
=
\sum_{t\ge 1}\mathrm{code}(e_{t+1})B^{t-1}v
=
X(\sigma\mathbf e).
\]
故
\[
\mathcal T\circ X=X\circ \sigma.
\]
再由定理 \ref{thm:xi-time-part62d-hologram-image-cantor-homeomorphism}，\(X\) 为到 \(\mathcal X_E\) 的拓扑同胚，故上述关系即给出严格共轭。证毕。
\end{proof}

\begin{corollary}[全息地址的模 \texorpdfstring{$B^n$}{Bn} 余类与长度 \texorpdfstring{$n$}{n} 前缀的严格对应]\label{cor:xi-time-part62deb-hologram-prefix-congruence-bijection}
\leanverified{paper\_xi\_time\_part62deb\_hologram\_prefix\_congruence\_bijection}
在定理 \ref{thm:xi-time-part62deb-hologram-affine-fullshift-conjugacy} 的设定下，记
\[
T:=T_2(E_{\min}).
\]
则对任意 \(\mathbf e,\mathbf f\in E^{\NN}\) 与任意 \(n\ge 1\)，下列两条命题等价：
\[
X(\mathbf e)\equiv X(\mathbf f)\pmod{B^nT},
\]
\[
(e_1,\dots,e_n)=(f_1,\dots,f_n).
\]
特别地，长度 \(n\) 的 cylinder 族与 \(\mathcal X_E\) 在商空间 \(T/B^nT\) 中所占据的深度 \(n\) 余类严格一一对应；因此，有限前缀在模 \(B^n\) 意义下是严格可逆的。
\end{corollary}
\begin{proof}
若前 \(n\) 位相同，则由定理 \ref{thm:xi-time-part62deb-hologram-affine-fullshift-conjugacy} 的首位递推迭代 \(n\) 次，得到
\[
X(\mathbf e)
=
\sum_{t=1}^{n}\mathrm{code}(e_t)B^{t-1}v
+B^nX(\sigma^n\mathbf e),
\]
\[
X(\mathbf f)
=
\sum_{t=1}^{n}\mathrm{code}(f_t)B^{t-1}v
+B^nX(\sigma^n\mathbf f).
\]
由于两者前 \(n\) 位相同，故
\[
X(\mathbf e)-X(\mathbf f)\in B^nT,
\]
即
\[
X(\mathbf e)\equiv X(\mathbf f)\pmod{B^nT}.
\]

反之，若
\[
X(\mathbf e)\equiv X(\mathbf f)\pmod{B^nT}
\]
而前 \(n\) 位并不相同，令 \(m<n\) 为首次分歧位置减一，则定理 \ref{thm:killo-2adic-holographic-exact-cylinder-separation} 给出
\[
X(\mathbf e)-X(\mathbf f)\in B^mT\setminus B^{m+1}T.
\]
但由 \(m+1\le n\) 可知 \(B^nT\subseteq B^{m+1}T\)，这与
\[
X(\mathbf e)-X(\mathbf f)\in B^nT
\]
矛盾。故前 \(n\) 位必须相同。

最后，对任意长度 \(n\) 前缀 \(a\in E^n\)，其 cylinder 恰由上述等价关系对应于 \(\mathcal X_E\) 在某一唯一 \(B^nT\)-余类中的截面。证毕。
\end{proof}

\begin{corollary}[全息地址系统的拓扑熵、周期点、primitive 轨道与 Artin--Mazur \texorpdfstring{$\zeta$}{zeta} 全闭式]\label{cor:xi-time-part62deb-hologram-fullshift-entropy-primitive-zeta}
\leanverified{paper\_xi\_time\_part62deb\_hologram\_fullshift\_entropy\_primitive\_zeta}
在定理 \ref{thm:xi-time-part62deb-hologram-affine-fullshift-conjugacy} 的设定下，对任意 \(n\ge 1\)，映射 \(\mathcal T\) 满足
\[
h_{\mathrm{top}}(\mathcal T)=\log |E|,
\qquad
\#\operatorname{Fix}(\mathcal T^n)=|E|^n.
\]
若记 \(\operatorname{Prim}(n)\) 为 \(\mathcal T\) 的最小周期恰为 \(n\) 的 primitive 周期轨道数，则
\[
\operatorname{Prim}(n)
=
\frac1n\sum_{d\mid n}\mu(d)\,|E|^{n/d},
\]
其中 \(\mu\) 为经典 Möbius 函数。相应地，Artin--Mazur 动力学 \(\zeta\) 函数满足
\[
\zeta_{\mathcal T}(z)
:=
\exp\!\left(
\sum_{n\ge 1}\frac{\#\operatorname{Fix}(\mathcal T^n)}{n}z^n
\right)
=
\frac{1}{1-|E|z}.
\]
\end{corollary}
\begin{proof}
由定理 \ref{thm:xi-time-part62deb-hologram-affine-fullshift-conjugacy}，\((\mathcal X_E,\mathcal T)\) 与一侧满移位 \((E^{\NN},\sigma)\) 拓扑共轭，故其全部动力学不变量均与 \(\sigma\) 相同。因而
\[
h_{\mathrm{top}}(\mathcal T)=h_{\mathrm{top}}(\sigma)=\log |E|,
\qquad
\#\operatorname{Fix}(\mathcal T^n)=\#\operatorname{Fix}(\sigma^n)=|E|^n.
\]
再由周期点与 primitive 轨道之间的标准关系
\[
\#\operatorname{Fix}(\mathcal T^n)=\sum_{d\mid n} d\,\operatorname{Prim}(d),
\]
对该卷积施行 Möbius 反演，即得
\[
\operatorname{Prim}(n)
=
\frac1n\sum_{d\mid n}\mu(d)\,|E|^{n/d}.
\]
最后，将
\[
\#\operatorname{Fix}(\mathcal T^n)=|E|^n
\]
代入 Artin--Mazur \(\zeta\) 函数定义，得到
\[
\zeta_{\mathcal T}(z)
=
\exp\!\left(\sum_{n\ge 1}\frac{|E|^n}{n}z^n\right)
=
\exp\!\bigl(-\log(1-|E|z)\bigr)
=
\frac{1}{1-|E|z}.
\]
证毕。
\end{proof}

\paragraph{有限局部化 rank-one 轴的全局障碍与 Pontryagin 对偶账本}
在有限 Grothendieck 有理秩目标中的乘法嵌入障碍、有限局部化直和账本的二轴分解以及有限素数截断对象的最小通道数都已建立之后，H.130--H.133 的离散/紧对偶链还可继续压缩为如下四条结论：任何有限多个 rank-one 局部化轴只足以容纳有限多个独立素数方向，而完整的乘法幺半群则严格需要可数无穷多轴；该障碍在 Pontryagin 对偶后等价地表现为有限个圆因子与有限个局部化 solenoid 因子所组成的紧交换审计平台无法忠实承载全部乘法结构。

\begin{theorem}[任意有限个自由整数轴与局部化 rank-one 轴的直和都无法忠实承载 \texorpdfstring{$\NN_{\times}$}{N×}]\label{thm:xi-time-part62debb-finite-localization-axes-no-full-multiplicative-host}
设 \(a,b\in \ZZ_{\ge 0}\)，\(\mathbf S=(S_1,\dots,S_b)\) 为有限素数集族，并定义
\[
\mathcal L_{a,\mathbf S}
:=
\ZZ^a\oplus G_{S_1}\oplus\cdots\oplus G_{S_b}.
\]
则不存在幺半群单射
\[
\NN_{\times}\hookrightarrow \mathcal L_{a,\mathbf S}.
\]
\leanverified{paper\_xi\_time\_part62debb\_finite\_localization\_axes\_no\_full\_multiplicative\_host}
\end{theorem}
\begin{proof}
由定理 \ref{thm:killo-localization-circle-dimension-prime-ledger-splitting}，每个 \(G_{S_j}\) 都满足
\[
G_{S_j}\otimes_{\ZZ}\QQ\cong \QQ.
\]
因此
\[
\mathcal L_{a,\mathbf S}\otimes_{\ZZ}\QQ
\cong
\QQ^{a+b},
\]
从而
\[
\operatorname{rank}_{\QQ}
\bigl(\mathcal L_{a,\mathbf S}\otimes_{\ZZ}\QQ\bigr)=a+b<\infty.
\]
又 \(\mathcal L_{a,\mathbf S}\) 本身就是交换可消去幺半群，且其 Grothendieck 完成仍为 \(\mathcal L_{a,\mathbf S}\) 本身，故定理
\ref{thm:xi-time-part62ca-no-multiplicative-embedding-into-finite-grothendieck-rank}
可直接应用，得到所述单射不存在。证毕。
\end{proof}

\begin{theorem}[有限素数支撑自由乘法子幺半群的最小局部化轴阈值恰等于生成素数个数]\label{thm:xi-time-part62debb-finite-prime-support-localized-axis-threshold}
设 \(p_1,\dots,p_T\) 为两两不同的素数，并记
\[
M_T:=\langle p_1,\dots,p_T\rangle_{\times}\subset \NN_{\times}.
\]
对任意 \(a,b\in \ZZ_{\ge 0}\) 与任意有限素数集族 \(\mathbf S=(S_1,\dots,S_b)\)，有精确等价
\[
\exists\ \text{幺半群单射 }\ 
M_T\hookrightarrow \mathcal L_{a,\mathbf S}
\iff
a+b\ge T.
\]
换言之，有限局部化轴族能够忠实承载 \(T\) 个独立素数方向，当且仅当其总有理秩至少为 \(T\)；该阈值严格且最优。
\leanverified{paper\_xi\_time\_part62debb\_finite\_prime\_support\_localized\_axis\_threshold}
\end{theorem}
\begin{proof}
记
\[
T_1=\cdots=T_a=\varnothing,
\qquad
T_{a+j}=S_j\quad (1\le j\le b).
\]
由于
\[
G_{\varnothing}=\ZZ,
\]
可将目标账本重写为
\[
\mathcal L_{a,\mathbf S}
\cong
\bigoplus_{i=1}^{a+b}G_{T_i}.
\]
另一方面，无论所选 \(p_1,\dots,p_T\) 是哪 \(T\) 个两两不同的素数，指数向量映射都给出规范同构
\[
M_T\cong (\NN^T,+),
\]
故 \(M_T\) 与定理
\ref{thm:xi-localized-directsum-prime-truncation-minimal-channels}
中的标准 \(T\)-素数截断幺半群仅差一个生成元重标号。
于是定理 \ref{thm:xi-localized-directsum-prime-truncation-minimal-channels}
恰可直接应用，并给出
\[
M_T\hookrightarrow \mathcal L_{a,\mathbf S}
\iff
a+b\ge T.
\]
证毕。
\end{proof}

\begin{corollary}[完整乘法幺半群的局部化轴复杂度恰为 \texorpdfstring{$\aleph_0$}{aleph0}]\label{cor:xi-time-part62debb-full-multiplicative-localized-axis-complexity}
约定 \(G_{\varnothing}:=\ZZ\)。对任意交换幺半群 \(M\)，定义其局部化轴复杂度为
\[
\ell_{\mathrm{loc}}(M)
:=
\min\left\{
|I|:\ \exists\ (S_i)_{i\in I},\ S_i\subset \mathbb P\ \text{有限},
M\hookrightarrow \bigoplus_{i\in I}G_{S_i}
\right\},
\]
只要右端集合非空。则
\[
\ell_{\mathrm{loc}}(\NN_{\times})=\aleph_0.
\]
\end{corollary}
\begin{proof}
由定理 \ref{thm:xi-time-part62debb-finite-localization-axes-no-full-multiplicative-host}，任何有限个局部化轴的直和都无法承载 \(\NN_{\times}\)，故
\[
\ell_{\mathrm{loc}}(\NN_{\times})\ge \aleph_0.
\]
另一方面，按素数估值给出的映射
\[
v:\NN_{\times}\longrightarrow \bigoplus_{p\in\mathbb P}\ZZ
\subset
\bigoplus_{p\in\mathbb P}G_{\{p\}},
\qquad
n\longmapsto \bigl(v_p(n)\bigr)_{p\in\mathbb P},
\]
是标准幺半群单射。由于素数集 \(\mathbb P\) 可数，右端只使用了可数无穷多个局部化轴，故
\[
\ell_{\mathrm{loc}}(\NN_{\times})\le \aleph_0.
\]
两式合并即得
\[
\ell_{\mathrm{loc}}(\NN_{\times})=\aleph_0.
\]
证毕。
\end{proof}

\leanverified{paper\_xi\_time\_part62debb\_full\_multiplicative\_localized\_axis\_complexity}

\begin{theorem}[有限个圆因子与局部化 solenoid 因子的 Pontryagin 对偶账本无法忠实编码全部乘法]\label{thm:xi-time-part62debb-finite-circle-solenoid-dual-ledger-obstruction}
设
\[
K_{a,\mathbf S}
:=
\TT^a\times \widehat{G_{S_1}}\times\cdots\times \widehat{G_{S_b}},
\]
其中 \(a,b\in\ZZ_{\ge 0}\)，且 \(S_1,\dots,S_b\subset \mathbb P\) 皆有限。则不存在幺半群单射
\[
\NN_{\times}\hookrightarrow \widehat{K_{a,\mathbf S}}.
\]
等价地，任何只由有限多个圆相位因子与有限多个局部化 prime-ledger solenoid 因子组成的紧交换审计平台，都不可能在其 Pontryagin 对偶账本上忠实编码完整乘法结构。
\end{theorem}
\begin{proof}
由 Pontryagin 对偶与双对偶识别，
\[
\widehat{K_{a,\mathbf S}}
\cong
\widehat{\TT^a}\oplus
\widehat{\widehat{G_{S_1}}}\oplus\cdots\oplus \widehat{\widehat{G_{S_b}}}
\cong
\ZZ^a\oplus G_{S_1}\oplus\cdots\oplus G_{S_b}
=
\mathcal L_{a,\mathbf S}.
\]
若存在幺半群单射
\[
\NN_{\times}\hookrightarrow \widehat{K_{a,\mathbf S}},
\]
则经由上述同构便得到
\[
\NN_{\times}\hookrightarrow \mathcal L_{a,\mathbf S},
\]
这与定理
\ref{thm:xi-time-part62debb-finite-localization-axes-no-full-multiplicative-host}
矛盾。故所述单射不存在。证毕。
\end{proof}

\paragraph{局部化数系的素数邻接自然扩张与 window-\texorpdfstring{$6$}{6} 根字典的重量径向几何}
在局部化数系对偶 solenoid 的圆周商结构、自然扩张的分支寄存器实现以及 window-\(6\) 的 \(B_3/C_3\) 根字典均已确立之后，还可把素数邻接与循环核几何进一步压缩为下列五条结构命题。前两条刻画 \(S\)-外素数在 Pontryagin 对偶端诱导的自然扩张；后三条说明 \(X_6^{\mathrm{cyc}}\) 的 Weyl 轨道、根多胞体组合型与双字典之间的变换全部由 Hamming 重量这一单一离散参数控制。

\begin{theorem}[向 \texorpdfstring{$\widehat{G_S}$}{G-hat-S} 邻接一个 \texorpdfstring{$S$}{S}-外素数的自然扩张恰为 \texorpdfstring{$\widehat{G_{S\cup\{p\}}}$}{G-hat-(S union p)}]\label{thm:xi-time-part62dea-localized-prime-adjunction-natural-extension}
设 \(S\subset \mathbb P\) 为有限素数集，记
\[
G_S:=\ZZ[S^{-1}],
\qquad
\Sigma_S:=\widehat{G_S},
\]
以及 Pontryagin 对偶短正合列
\[
0\longrightarrow \prod_{q\in S}\ZZ_q
\longrightarrow \Sigma_S
\xrightarrow{\ \pi_S\ }
\TT
\longrightarrow 0.
\]
取素数 \(p\notin S\)，定义
\[
m_p:G_S\to G_S,
\qquad
x\longmapsto px,
\qquad
\Phi_{S,p}:=\widehat{m_p}:\Sigma_S\to\Sigma_S.
\]
则有：
\begin{enumerate}
  \item \(\Phi_{S,p}\) 的每条纤维恰有 \(p\) 个点；
  \item 其自然扩张
  \[
  (\Sigma_S)^{\leftarrow}_{\Phi_{S,p}}
  :=
  \varprojlim\bigl(
  \Sigma_S\xleftarrow{\ \Phi_{S,p}\ }
  \Sigma_S\xleftarrow{\ \Phi_{S,p}\ }\cdots
  \bigr)
  \]
  存在规范同构
  \[
  (\Sigma_S)^{\leftarrow}_{\Phi_{S,p}}
  \cong
  \widehat{G_{S\cup\{p\}}};
  \]
  \item 因而相应单寄存器分支指标实现的最小可行字母表基数恰为 \(p\)。
\end{enumerate}
\end{theorem}
\begin{proof}
由于 \(m_p\) 为单射，\(\Phi_{S,p}\) 是连续满同态，并且
\[
\ker(\Phi_{S,p})^\vee
\cong
\operatorname{coker}(m_p)
=
G_S/pG_S.
\]
当 \(p\notin S\) 时，\(G_S\) 中每个元素都可写成
\[
x=\frac{a}{\prod_{q\in S}q^{k_q}},
\]
且分母在 \(\ZZ/p\ZZ\) 中可逆，故模 \(p\) 约化诱导同构
\[
G_S/pG_S\cong \ZZ/p\ZZ.
\]
再取 Pontryagin 对偶，得到
\[
\ker(\Phi_{S,p})\cong \ZZ/p\ZZ.
\]
因此每条纤维都是 \(\ker(\Phi_{S,p})\) 的平移陪集，恰有 \(p\) 个点。这证明了第 \(1\) 条。

再看自然扩张。Pontryagin 对偶把紧 Abel 群的逆极限送为离散 Abel 群的直极限，因此
\[
\bigl((\Sigma_S)^{\leftarrow}_{\Phi_{S,p}}\bigr)^\vee
\cong
\varinjlim\bigl(
G_S\xrightarrow{\ m_p\ }G_S\xrightarrow{\ m_p\ }\cdots
\bigr).
\]
该直极限可由
\[
[x,k]\longmapsto p^{-k}x
\]
规范识别为
\[
\bigcup_{k\ge 0}p^{-k}G_S.
\]
而
\[
\bigcup_{k\ge 0}p^{-k}G_S
=
\ZZ[(S\cup\{p\})^{-1}]
=
G_{S\cup\{p\}},
\]
因为右端恰由分母素因子只允许落在 \(S\cup\{p\}\) 中的有理数组成。故
\[
\bigl((\Sigma_S)^{\leftarrow}_{\Phi_{S,p}}\bigr)^\vee
\cong
G_{S\cup\{p\}},
\]
再取 Pontryagin 对偶即得
\[
(\Sigma_S)^{\leftarrow}_{\Phi_{S,p}}
\cong
\widehat{G_{S\cup\{p\}}}.
\]
这证明了第 \(2\) 条。

最后，由第 \(1\) 条知
\[
\sup_{y\in\Sigma_S}\#\Phi_{S,p}^{-1}(y)=p.
\]
于是定理 \ref{thm:xi-natural-extension-minimal-d-alphabet-realization} 立即给出：该自然扩张的单寄存器分支指标实现所需最小字母表基数恰为 \(p\)。证毕。
\end{proof}

\begin{corollary}[有限素数集的邻接顺序无关，且附加素方向的核恰为 \texorpdfstring{$\prod_{p\in T}\ZZ_p$}{product Zp}]\label{cor:xi-time-part62dea-prime-adjunction-order-independent-kernel-product}
沿用定理 \ref{thm:xi-time-part62dea-localized-prime-adjunction-natural-extension} 的记号。取有限集合
\[
T\subset \mathbb P\setminus S.
\]
则有：
\begin{enumerate}
  \item 无论按何种顺序逐个把 \(T\) 中素数作为自然扩张方向邻接进去，最终得到的逆极限对象都规范同构于
  \[
  \widehat{G_{S\cup T}};
  \]
  \item 并存在规范短正合列
  \[
  0\longrightarrow \prod_{p\in T}\ZZ_p
  \longrightarrow \widehat{G_{S\cup T}}
  \longrightarrow \widehat{G_S}
  \longrightarrow 0.
  \]
\end{enumerate}
\end{corollary}
\begin{proof}
设 \(T=\{p_1,\dots,p_r\}\)。对任意排列 \(\sigma\in S_r\)，反复应用定理 \ref{thm:xi-time-part62dea-localized-prime-adjunction-natural-extension}，得到一串规范同构
\[
\widehat{G_S}
\leadsto
\widehat{G_{S\cup\{p_{\sigma(1)}\}}}
\leadsto
\widehat{G_{S\cup\{p_{\sigma(1)},p_{\sigma(2)}\}}}
\leadsto \cdots \leadsto
\widehat{G_{S\cup T}}.
\]
由于每一步只是在当前局部化集合中并入一个额外素数，而集合并运算可交换，故最终对象与邻接顺序无关。这证明了第 \(1\) 条。

对第 \(2\) 条，由定理 \ref{thm:killo-localization-circle-dimension-prime-ledger-splitting} 分别作用于 \(S\) 与 \(S\cup T\)，有短正合列
\[
0\longrightarrow \prod_{q\in S}\ZZ_q
\longrightarrow \widehat{G_S}
\xrightarrow{\ \pi_S\ }
\TT
\longrightarrow 0,
\]
\[
0\longrightarrow \prod_{q\in S\cup T}\ZZ_q
\longrightarrow \widehat{G_{S\cup T}}
\xrightarrow{\ \pi_{S\cup T}\ }
\TT
\longrightarrow 0.
\]
另一方面，推论 \ref{cor:xi-time-part62da-localization-pontryagin-surjection-poset} 在 \(S\subseteq S\cup T\) 时给出规范满射
\[
\widehat{G_{S\cup T}}\twoheadrightarrow \widehat{G_S},
\]
其核为
\[
\prod_{p\in (S\cup T)\setminus S}\ZZ_p.
\]
而
\[
(S\cup T)\setminus S=T,
\qquad
\prod_{q\in S\cup T}\ZZ_q
\cong
\Bigl(\prod_{q\in S}\ZZ_q\Bigr)\times
\Bigl(\prod_{p\in T}\ZZ_p\Bigr),
\]
故附加核恰为
\[
\prod_{p\in T}\ZZ_p.
\]
证毕。
\end{proof}

\begin{theorem}[window-\texorpdfstring{$6$}{6} 上 Weyl 轨道分裂完全由 Hamming 重量决定]\label{thm:xi-time-part62dea-window6-weight-determines-weyl-orbits}
\leanverified{paper\_xi\_time\_part62dea\_window6\_weight\_determines\_weyl\_orbits}
沿用定理 \ref{thm:window6-cyclic-weight-threshold-root-length} 与表 \ref{tab:fold6_b3c3_root_dictionary_b3}、表 \ref{tab:fold6_b3c3_root_dictionary_c3} 的记号。window-\(6\) 的刚性循环核满足
\[
|X_6|=21,
\qquad
|X_6^{\mathrm{cyc}}|=18,
\qquad
|X_6^{\mathrm{bdry}}|=3,
\]
以及
\[
\mathrm{wt}\text{-hist}(X_6^{\mathrm{cyc}})
=
\{0:1,\ 1:6,\ 2:9,\ 3:2\}.
\]
定义两张字典
\[
r_B:X_6^{\mathrm{cyc}}\to \Phi(B_3),
\qquad
r_C:X_6^{\mathrm{cyc}}\to \Phi(C_3).
\]
则有：
\begin{enumerate}
  \item 在 \(B_3\) 字典中，
  \[
  \mathrm{wt}(w)=1
  \iff
  r_B(w)\in \{\pm e_1,\pm e_2,\pm e_3\},
  \]
  即恰对应 \(B_3\) 的六条短根；
  \item 在 \(C_3\) 字典中，
  \[
  \mathrm{wt}(w)=1
  \iff
  r_C(w)\in \{\pm 2e_1,\pm 2e_2,\pm 2e_3\},
  \]
  即恰对应 \(C_3\) 的六条长根；
  \item 因而两张字典下的 Weyl 轨道分裂都完全由 Hamming 重量决定；
  \item 对应六点轨道的指示函数满足
  \[
  \mathbf 1_{\{\mathrm{wt}=1\}}(w)
  =
  \frac12\,\mathrm{wt}(w)\bigl(\mathrm{wt}(w)-2\bigr)\bigl(\mathrm{wt}(w)-3\bigr).
  \]
\end{enumerate}
\end{theorem}
\begin{proof}
由上述重量直方图可知
\[
\#\{w\in X_6^{\mathrm{cyc}}:\ \mathrm{wt}(w)=1\}=6.
\]
再由表 \ref{tab:fold6_b3c3_root_dictionary_b3} 逐项读出，这六个 weight-\(1\) 元恰对应
\[
(\pm1,0,0),\ (0,\pm1,0),\ (0,0,\pm1),
\]
即 \(B_3\) 的全部六条短根；其余十二个元素对应
\[
(\pm1,\pm1,0),\ (\pm1,0,\pm1),\ (0,\pm1,\pm1),
\]
即 \(B_3\) 的全部十二条长根。故第 \(1\) 条成立。

同样，由表 \ref{tab:fold6_b3c3_root_dictionary_c3} 可见，同一批 weight-\(1\) 元恰对应
\[
(\pm2,0,0),\ (0,\pm2,0),\ (0,0,\pm2),
\]
即 \(C_3\) 的六条长根；其余十二个元素仍对应
\[
(\pm1,\pm1,0),\ (\pm1,0,\pm1),\ (0,\pm1,\pm1),
\]
即 \(C_3\) 的十二条短根。故第 \(2\) 条成立。

于是两张字典下的 Weyl 轨道都被同一条件
\[
\mathrm{wt}(w)=1
\]
与其补集完全分离。这证明了第 \(3\) 条。

最后，在有限集合 \(\{0,1,2,3\}\) 上考虑唯一满足
\[
P(1)=1,
\qquad
P(0)=P(2)=P(3)=0
\]
的次数至多 \(3\) 的插值多项式。Lagrange 插值给出
\[
P(s)=\frac{(s-0)(s-2)(s-3)}{(1-0)(1-2)(1-3)}
=
\frac12\,s(s-2)(s-3).
\]
令 \(s=\mathrm{wt}(w)\) 即得第 \(4\) 条。证毕。
\end{proof}

\begin{theorem}[两张字典诱导两种不同的根多胞体：一者为 cuboctahedron，一者为尺度 \texorpdfstring{$2$}{2} 的 octahedron]\label{thm:xi-time-part62dea-window6-b3c3-root-polytopes-not-similar}
沿用定理 \ref{thm:xi-time-part62dea-window6-weight-determines-weyl-orbits} 的记号，定义
\[
P_B:=\operatorname{conv}\{r_B(w):\ w\in X_6^{\mathrm{cyc}}\},
\qquad
P_C:=\operatorname{conv}\{r_C(w):\ w\in X_6^{\mathrm{cyc}}\}.
\]
则有：
\begin{enumerate}
  \item
  \[
  P_B
  =
  \{x\in\RR^3:\ |x_i|\le 1,\ |x_1|+|x_2|+|x_3|\le 2\},
  \]
  即标准 cuboctahedron，且
  \[
  \operatorname{Vol}(P_B)=\frac{20}{3}.
  \]
  其中十二条长根是顶点，而六条短根位于六个正方形面的中心；
  \item
  \[
  P_C
  =
  \{x\in\RR^3:\ |x_1|+|x_2|+|x_3|\le 2\},
  \]
  即尺度 \(2\) 的正交叉多胞体，且
  \[
  \operatorname{Vol}(P_C)=\frac{32}{3}.
  \]
  其中六条长根是顶点，而十二条短根位于各棱的中点；
  \item 因而 \(P_B\) 与 \(P_C\) 具有不同的组合类型；同一 \(18\) 个循环态在两张字典下分别张成 cuboctahedron 与尺度 \(2\) 的 octahedron。
\end{enumerate}
\end{theorem}
\begin{proof}
由定理 \ref{thm:xi-time-part62dea-window6-weight-determines-weyl-orbits}，
\[
r_B(X_6^{\mathrm{cyc}})
=
\{\pm e_i,\ \pm e_i\pm e_j:\ 1\le i<j\le 3\}.
\]
对每个 \(i\)，短根 \(\pm e_i\) 都可写成相应四个长根的平均，例如
\[
e_1=
\frac14\bigl((1,1,0)+(1,-1,0)+(1,0,1)+(1,0,-1)\bigr),
\]
故
\[
P_B=\operatorname{conv}\{\pm e_i\pm e_j:\ 1\le i<j\le 3\}.
\]
而这正是标准 cuboctahedron 的顶点集，其半空间表示为
\[
\{x\in\RR^3:\ |x_i|\le 1,\ |x_1|+|x_2|+|x_3|\le 2\}.
\]
在第一卦限中，该区域等于单位立方体
\[
[0,1]^3
\]
去掉角四面体
\[
\{x_i\ge 0,\ x_1+x_2+x_3>2\},
\]
故第一卦限体积为
\[
1-\frac16=\frac56.
\]
乘上八个卦限，得到
\[
\operatorname{Vol}(P_B)=8\cdot \frac56=\frac{20}{3}.
\]
并且上式表明 \(e_1\) 位于面
\[
\{x_1=1,\ |x_2|+|x_3|\le 1\}
\]
的中心，其余 \(\pm e_2,\pm e_3\) 情形同理。

另一方面，
\[
r_C(X_6^{\mathrm{cyc}})
=
\{\pm 2e_i,\ \pm e_i\pm e_j:\ 1\le i<j\le 3\}.
\]
每条短根都可以写成两条长根的中点，例如
\[
(1,1,0)=\frac12\bigl((2,0,0)+(0,2,0)\bigr),
\]
故
\[
P_C=\operatorname{conv}\{\pm 2e_1,\pm 2e_2,\pm 2e_3\}.
\]
这正是尺度 \(2\) 的 \(\ell^1\)-球
\[
\{x\in\RR^3:\ |x_1|+|x_2|+|x_3|\le 2\},
\]
亦即三维正交叉多胞体。由三维交叉多胞体体积公式，
\[
\operatorname{Vol}(P_C)=\frac{2^3\cdot 2^3}{3!}=\frac{32}{3}.
\]
并且每条短根都是一条棱的中点。

最后，\(P_B\) 的顶点数为 \(12\)，而 \(P_C\) 的顶点数为 \(6\)。因此二者不可能具有相同的组合类型，更不可能由同一 Euclid 相似变换相互得到。证毕。
\end{proof}

\begin{theorem}[从 \texorpdfstring{$B_3$}{B3} 到 \texorpdfstring{$C_3$}{C3} 的字典变换由唯一的重量径向伸缩给出，且该变换不是线性的]\label{thm:xi-time-part62dea-window6-b3c3-unique-radial-weight-scaling}
沿用定理 \ref{thm:xi-time-part62dea-window6-weight-determines-weyl-orbits} 的记号。则有：
\begin{enumerate}
  \item 不存在任何线性变换 \(A\in GL_3(\RR)\) 使得
  \[
  r_C(w)=A\,r_B(w)
  \qquad(\forall\,w\in X_6^{\mathrm{cyc}});
  \]
  \item 然而存在唯一仅依赖 Hamming 重量的函数
  \[
  \lambda:\{0,1,2,3\}\to \RR_{>0}
  \]
  使得
  \[
  r_C(w)=\lambda(\mathrm{wt}(w))\,r_B(w)
  \qquad(\forall\,w\in X_6^{\mathrm{cyc}});
  \]
  \item 该函数显式为
  \[
  \lambda(1)=2,
  \qquad
  \lambda(0)=\lambda(2)=\lambda(3)=1,
  \]
  因而其在 \(\{0,1,2,3\}\) 上的唯一次数至多 \(3\) 插值多项式为
  \[
  \lambda(s)=1+\frac12\,s(s-2)(s-3).
  \]
\end{enumerate}
\end{theorem}
\begin{proof}
先证第 \(1\) 条。由两张表可直接读出
\[
r_B(\texttt{100000})=e_1,\qquad r_C(\texttt{100000})=2e_1,
\]
\[
r_B(\texttt{010000})=e_2,\qquad r_C(\texttt{010000})=2e_2,
\]
\[
r_B(\texttt{001000})=e_3,\qquad r_C(\texttt{001000})=2e_3.
\]
若存在 \(A\in GL_3(\RR)\) 满足
\[
r_C(w)=A\,r_B(w)\qquad(\forall\,w),
\]
则必有
\[
Ae_i=2e_i\qquad(i=1,2,3),
\]
从而
\[
A=2I_3.
\]
但对 \(\texttt{000000}\)（其 Hamming 重量为 \(0\)），表 \ref{tab:fold6_b3c3_root_dictionary_b3} 与表 \ref{tab:fold6_b3c3_root_dictionary_c3} 同时给出
\[
r_B(\texttt{000000})=r_C(\texttt{000000})=(1,1,0).
\]
于是
\[
Ar_B(\texttt{000000})
=
2(1,1,0)
\neq
(1,1,0)
=
r_C(\texttt{000000}),
\]
矛盾。故不存在这样的线性变换。

再证第 \(2\) 与第 \(3\) 条。由两张表逐项比较可见：
\[
\mathrm{wt}(w)=1\Longrightarrow r_C(w)=2\,r_B(w),
\]
而当
\[
\mathrm{wt}(w)\in\{0,2,3\}
\]
时都有
\[
r_C(w)=r_B(w).
\]
因此唯一可能的重量函数只能取值为
\[
\lambda(1)=2,
\qquad
\lambda(0)=\lambda(2)=\lambda(3)=1,
\]
并且它确实逐点满足
\[
r_C(w)=\lambda(\mathrm{wt}(w))\,r_B(w).
\]
唯一性显然，因为 \(\mathrm{wt}(w)\) 在 \(X_6^{\mathrm{cyc}}\) 上恰取遍 \(\{0,1,2,3\}\)。

最后，由定理 \ref{thm:xi-time-part62dea-window6-weight-determines-weyl-orbits} 已知
\[
\mathbf 1_{\{\mathrm{wt}=1\}}(w)
=
\frac12\,\mathrm{wt}(w)\bigl(\mathrm{wt}(w)-2\bigr)\bigl(\mathrm{wt}(w)-3\bigr),
\]
故
\[
\lambda(\mathrm{wt}(w))
=
1+\mathbf 1_{\{\mathrm{wt}=1\}}(w)
=
1+\frac12\,\mathrm{wt}(w)\bigl(\mathrm{wt}(w)-2\bigr)\bigl(\mathrm{wt}(w)-3\bigr).
\]
把变量记为 \(s\)，便得到所示唯一插值多项式
\[
\lambda(s)=1+\frac12\,s(s-2)(s-3).
\]
证毕。
\end{proof}

\noindent
由此，局部化数系的素数邻接塔与 window-\(6\) 的重量壳层几何在同一组结论中闭合：向 \(\widehat{G_S}\) 邻接一个 \(S\)-外素数，恰对应一次 \(p\)-分支自然扩张，并在逆极限上精确恢复 \(\widehat{G_{S\cup\{p\}}}\)；与此同时，\(X_6^{\mathrm{cyc}}\) 的 Weyl 轨道、根多胞体与 \(B_3/C_3\) 双字典之间的唯一非线性对应，又全部被 Hamming 重量所统摄。于是，prime-register 邻接与重量径向几何便在同一有限可审计框架内获得统一描述。

\paragraph{纤维相图的比例极限、容量曲线的跳跃反演、prime-shadow 衰减积与 Lee--Yang 三次同步}
在最大非可缩纤维的模 \(6\) 相图、连续容量曲线与尾谱/矩谱的互逆公式、奇素数单位根扭曲的四阶展开，以及 Lee--Yang 分支参数与振幅平方所张成的共同三次数域都已建立之后，这几条彼此分离的线路还可继续压缩为五个终端闭式：前两条把纤维尺度与容量层析写成完全可恢复的有限数据，第三条把最大非可缩/绝对最大纤维之比的极限点谱完全显式化，第四条把奇素数扭曲的总衰减锁定为严格正的 Euler 常数，最后一条则把 \(u\)-三次与 \(b\)-三次的“同数域性”提升为模 \(p\) 分解型的逐点同步。

\begin{theorem}[最大非可缩纤维的缺口公式、六相极限律与指数级别保持]\label{thm:xi-time-part62df-maxfiber-gap-ratio-limit}
沿用定理 \ref{thm:xi-time-part62db-maxfiber-homotopy-phase-mod6} 与推论 \ref{cor:xi-time-part62d-max-noncontractible-fiber-tax-mod6} 的记号。对 \(m\ge 17\)，记
\[
D_m:=\max_{x\in X_m}d_m(x),
\qquad
\widetilde D_m:=\max\{d_m(x):x\in X_m,\ |K_x|\not\simeq *\},
\]
\[
\Delta_m:=D_m-\widetilde D_m.
\]
则有精确闭式
\[
\Delta_m=
\begin{cases}
0,& m\equiv 0,4\pmod 6,\\[2mm]
F_{(m-9)/2},& m\equiv 1,5\pmod 6,\\[2mm]
F_{m/2-3},& m\equiv 2\pmod 6,\\[2mm]
F_{(m-9)/2}+F_{(m-17)/2},& m\equiv 3\pmod 6,
\end{cases}
\]
并且沿模 \(6\) 子序列的比例极限满足
\[
\lim_{\substack{m\to\infty\\ m\equiv 0,4\!\!\!\!\pmod 6}}
\frac{\widetilde D_m}{D_m}=1,
\]
\[
\lim_{\substack{m\to\infty\\ m\equiv 1,5\!\!\!\!\pmod 6}}
\frac{\widetilde D_m}{D_m}
=
1-\frac{1}{2}\varphi^{-5},
\]
\[
\lim_{\substack{m\to\infty\\ m\equiv 2\!\!\!\!\pmod 6}}
\frac{\widetilde D_m}{D_m}
=
1-\varphi^{-5},
\]
\[
\lim_{\substack{m\to\infty\\ m\equiv 3\!\!\!\!\pmod 6}}
\frac{\widetilde D_m}{D_m}
=
1-\frac{1}{2}\bigl(\varphi^{-5}+\varphi^{-9}\bigr).
\]
特别地，非可缩性约束并不改变最大纤维的 Fibonacci 指数级别；它只在模 \(6\) 的三类残余相位 \(\{1,5\}\)、\(\{2\}\)、\(\{3\}\) 上削去显式常数比例，而在 \(0,4\) 两类上完全不削减主极值。
\end{theorem}
\begin{proof}
推论 \ref{cor:xi-time-part62d-max-noncontractible-fiber-tax-mod6} 已给出 \(\Delta_m\) 的精确分段公式：
\[
\Delta_m=
\begin{cases}
0,& m\equiv 0,4\pmod 6,\\[1mm]
F_{(m-9)/2},& m\equiv 1,5\pmod 6,\\[1mm]
F_{m/2-3},& m\equiv 2\pmod 6,\\[1mm]
F_{(m-9)/2}+F_{(m-17)/2},& m\equiv 3\pmod 6.
\end{cases}
\]
同一推论又给出了四条比值极限公式。逐条汇总即得所述结论。
\end{proof}

\begin{corollary}[最大非可缩纤维与绝对最大纤维之比不存在极限，且其极限点由模 \texorpdfstring{$6$}{6} 相位完全给出]\label{cor:xi-time-part62df-maxfiber-ratio-nonconvergence-mod6}
沿用定理 \ref{thm:xi-time-part62df-maxfiber-gap-ratio-limit} 的记号。则序列
\[
\frac{\widetilde D_m}{D_m}
\]
不收敛；其全部子序列极限恰为
\[
1,\qquad
1-\frac{1}{2}\varphi^{-5},\qquad
1-\varphi^{-5},\qquad
1-\frac{1}{2}\bigl(\varphi^{-5}+\varphi^{-9}\bigr).
\]
因而其极限点集合恰含四个彼此不同的值。
\end{corollary}
\begin{proof}
定理 \ref{thm:xi-time-part62df-maxfiber-gap-ratio-limit} 已给出沿
\[
m\equiv 0,4\pmod 6,\qquad
m\equiv 1,5\pmod 6,\qquad
m\equiv 2\pmod 6,\qquad
m\equiv 3\pmod 6
\]
这四类无限子序列的极限分别为
\[
1,\qquad
1-\frac{1}{2}\varphi^{-5},\qquad
1-\varphi^{-5},\qquad
1-\frac{1}{2}\bigl(\varphi^{-5}+\varphi^{-9}\bigr).
\]
因此这四个数全部都是子序列极限。

反过来，任取任意子序列 \((m_j)\)。模 \(6\) 只有有限个剩余类，故可再抽取一条子序列，使其落在某个固定剩余类上。对该再抽取子序列应用定理 \ref{thm:xi-time-part62df-maxfiber-gap-ratio-limit}，便知其极限必为上面四个常数之一。故不存在其他极限点。

最后，这四个常数显然互不相同，故原序列不收敛，且极限点集合恰由上述四个值组成。证毕。
\end{proof}

\begin{theorem}[连续容量曲线对纤维多重度谱的完全可逆性]\label{thm:xi-time-part62df-capacity-curve-complete-inversion}
\leanverified{paper\_xi\_time\_part62df\_capacity\_curve\_complete\_inversion}
固定分辨率 \(m\)。定义尾计数函数、连续容量曲线与整数直方图
\[
N_m(t):=\#\{x\in X_m:\ d_m(x)\ge t\},
\qquad
\mathcal C_m^{\mathrm{cont}}(T):=\sum_{x\in X_m}\min\bigl(d_m(x),T\bigr),
\]
\[
h_m(n):=\#\{x\in X_m:\ d_m(x)=n\}
\qquad (n\ge 1).
\]
则：
\begin{enumerate}
  \item 对每个整数 \(n\ge 1\)，都有精确跳跃反演公式
  \[
  h_m(n)=N_m(n)-N_m(n+1)
  =\partial_T^{-}\mathcal C_m^{\mathrm{cont}}(n)-\partial_T^{+}\mathcal C_m^{\mathrm{cont}}(n).
  \]
  因而 \(\mathcal C_m^{\mathrm{cont}}\) 唯一决定全部纤维多重度多重集 \(\{d_m(x):x\in X_m\}\)。
  \item 对任意 \(q>1\)，有精确积分公式
  \[
  S_q(m):=\sum_{x\in X_m}d_m(x)^q
  =
  q(q-1)\int_0^\infty
  T^{q-2}\bigl(2^m-\mathcal C_m^{\mathrm{cont}}(T)\bigr)\,\dd T.
  \]
  \item 对任意函数 \(F:\ZZ_{\ge 1}\to \RR\)，都有
  \[
  \sum_{x\in X_m}F\bigl(d_m(x)\bigr)
  =
  \sum_{n\ge 1}h_m(n)F(n),
  \]
  故任意仅依赖于多重度多重集的对称泛函都由 \(\mathcal C_m^{\mathrm{cont}}\) 唯一恢复。
\end{enumerate}
\end{theorem}
\begin{proof}
由定理 \ref{thm:xi-capacity-tail-histogram-moment-complete-statistic}，
\[
h_m(n)=N_m(n)-N_m(n+1)
\qquad(n\ge 1),
\]
且
\[
\mathcal C_m^{\mathrm{cont}}(T)=\int_0^T N_m(t)\,\dd t.
\]
由于所有 \(d_m(x)\) 都是正整数，\(N_m\) 是右连续阶梯函数，而 \(\mathcal C_m^{\mathrm{cont}}\) 是折点只出现在正整数处的连续分段线性函数。于是对每个 \(n\ge 1\)，其左导数与右导数分别为
\[
\partial_T^{-}\mathcal C_m^{\mathrm{cont}}(n)=N_m(n),
\qquad
\partial_T^{+}\mathcal C_m^{\mathrm{cont}}(n)=N_m(n+1).
\]
相减即得
\[
h_m(n)
=
N_m(n)-N_m(n+1)
=
\partial_T^{-}\mathcal C_m^{\mathrm{cont}}(n)-\partial_T^{+}\mathcal C_m^{\mathrm{cont}}(n),
\]
从而第一条成立。

再记
\[
R_m(T):=2^m-\mathcal C_m^{\mathrm{cont}}(T).
\]
由于
\[
\sum_{x\in X_m}d_m(x)=2^m,
\]
可重写为
\[
R_m(T)=\sum_{x\in X_m}\bigl(d_m(x)-T\bigr)_+.
\]
于是对几乎处处的 \(T>0\)，有
\[
-R_m'(T)=N_m(T).
\]
由定理 \ref{thm:xi-capacity-tail-histogram-moment-complete-statistic} 的 Mellin--Stieltjes 公式，
\[
S_q(m)=q\int_0^\infty T^{q-1}N_m(T)\,\dd T
=
-q\int_0^\infty T^{q-1}R_m'(T)\,\dd T.
\]
当 \(q>1\) 时，边界项满足
\[
\lim_{T\downarrow 0}T^{q-1}R_m(T)=0,
\qquad
\lim_{T\to\infty}T^{q-1}R_m(T)=0,
\]
因为 \(R_m(0)=2^m\) 而 \(R_m(T)=0\) 对一切充分大的 \(T\) 都成立。故分部积分给出
\[
S_q(m)=q(q-1)\int_0^\infty T^{q-2}R_m(T)\,\dd T
=
q(q-1)\int_0^\infty T^{q-2}\bigl(2^m-\mathcal C_m^{\mathrm{cont}}(T)\bigr)\,\dd T.
\]
这就证明了第二条。

最后，由第一条可从 \(\mathcal C_m^{\mathrm{cont}}\) 唯一恢复全部 \(h_m(n)\)。于是对任意 \(F\) 都有
\[
\sum_{x\in X_m}F(d_m(x))
=
\sum_{n\ge 1}h_m(n)F(n),
\]
右端完全由 \(\mathcal C_m^{\mathrm{cont}}\) 决定，第三条成立。证毕。
\end{proof}

\begin{theorem}[奇素数单位根扭曲的 Euler 衰减积收敛到正值]\label{thm:xi-time-part62df-odd-prime-twist-euler-positive-limit}
沿用定理 \ref{thm:xi-time-part9p-odd-prime-shadow-product} 的记号，并记
\[
\lambda:=\rho(1)=\varphi^2,
\qquad
r_p:=\left|\frac{\rho(e^{2\pi i/p})}{\lambda}\right|
\qquad (p\ge 3\ \text{为奇素数}).
\]
则有：
\begin{enumerate}
  \item 对数级数绝对收敛，
  \[
  \sum_{p\ge 3}\bigl|\log r_p\bigr|<\infty;
  \]
  \item Euler 乘积
  \[
  \prod_{p\ge 3}r_p
  \]
  收敛到某个常数
  \[
  \mathcal C_{\mathrm{prime}}\in(0,1);
  \]
  \item 若记
  \[
  \kappa_\infty:=\frac{6\sqrt5-5}{250},
  \]
  则当 \(X\to\infty\) 时有
  \[
  \sum_{\substack{p\le X\\ p\ge 3}}\log r_p
  =
  -\kappa_\infty(2\pi)^2
  \sum_{\substack{p\le X\\ p\ge 3}}\frac{1}{p^2}
  +O(1),
  \]
  从而
  \[
  \prod_{\substack{p\le X\\ p\ge 3}}r_p
  =
  \mathcal C_{\mathrm{prime}}+o(1).
  \]
\end{enumerate}
\end{theorem}
\begin{proof}
定理 \ref{thm:xi-time-part9p-odd-prime-shadow-product} 的证明已经给出
\[
\log r_p
=
-\frac{6\sqrt5-5}{250}\left(\frac{2\pi}{p}\right)^2
+\frac{7+24\sqrt5}{75000}\left(\frac{2\pi}{p}\right)^4
+O(p^{-6}).
\]
由于素数级数
\[
\sum_{p\ge 3}p^{-2},
\qquad
\sum_{p\ge 3}p^{-4},
\qquad
\sum_{p\ge 3}p^{-6}
\]
全部收敛，故
\[
\sum_{p\ge 3}|\log r_p|<\infty.
\]
于是
\[
\sum_{p\ge 3}\log r_p
\]
绝对收敛，从而 Euler 乘积
\[
\prod_{p\ge 3}r_p
=
\exp\!\left(\sum_{p\ge 3}\log r_p\right)
\]
收敛到某个严格正的实数 \(\mathcal C_{\mathrm{prime}}>0\)。

再由定理 \ref{thm:xi-time-part9p-odd-prime-shadow-product}，对每个奇素数 \(p\) 都有
\[
0<r_p<1,
\]
故全部部分乘积都严格小于 \(1\)，极限常数满足
\[
0<\mathcal C_{\mathrm{prime}}<1.
\]

对部分和，把上式从 \(p=3\) 到 \(p\le X\) 求和，得到
\[
\sum_{\substack{p\le X\\ p\ge 3}}\log r_p
=
-\kappa_\infty(2\pi)^2
\sum_{\substack{p\le X\\ p\ge 3}}p^{-2}
+\frac{7+24\sqrt5}{75000}(2\pi)^4
\sum_{\substack{p\le X\\ p\ge 3}}p^{-4}
+O\!\left(\sum_{\substack{p\le X\\ p\ge 3}}p^{-6}\right).
\]
后两项在 \(X\to\infty\) 时有界，故
\[
\sum_{\substack{p\le X\\ p\ge 3}}\log r_p
=
-\kappa_\infty(2\pi)^2
\sum_{\substack{p\le X\\ p\ge 3}}\frac{1}{p^2}
+O(1).
\]
最后，由
\[
\sum_{\substack{p\le X\\ p\ge 3}}\log r_p
\longrightarrow
\sum_{p\ge 3}\log r_p
=
\log \mathcal C_{\mathrm{prime}},
\]
指数化即得
\[
\prod_{\substack{p\le X\\ p\ge 3}}r_p
=
\mathcal C_{\mathrm{prime}}+o(1).
\]
证毕。
\end{proof}

\begin{theorem}[Lee--Yang 分支参数与振幅平方的模 \texorpdfstring{$p$}{p} 分解型同步]\label{thm:xi-time-part62df-leyang-u-b-splitting-synchrony}
定义三次多项式
\[
f_u(U):=324U^3-540U^2+333U-74,
\]
\[
f_b(T):=162T^3+1593T^2+1998T+1184.
\]
设 \(u\) 与 \(b\) 分别为其根。沿用定理 \ref{thm:xi-terminal-zm-stokes-leyang-shared-artin-representation} 的记号，则
\[
\QQ(u)=\QQ(b)=\QQ(\lambda_*)=:K,
\]
且其共同伽罗瓦闭包记为 \(L\)，满足
\[
\Gal(L/\QQ)\cong S_3.
\]
再由既有判别式公式，
\[
\Disc(f_u)=-2^6\cdot 3^9\cdot 37,
\qquad
\Disc(f_b)=-2^2\cdot 3^9\cdot 11^2\cdot 37\cdot 109^2.
\]
于是对每个素数
\[
p\nmid 2\cdot 3\cdot 11\cdot 37\cdot 109,
\]
多项式 \(f_u\bmod p\) 与 \(f_b\bmod p\) 在 \(\FF_p\) 上具有完全相同的分解型。因而三种同步分解事件的自然密度分别为
\[
\delta\bigl(f_u,f_b\ \text{同时全分裂}\bigr)=\frac16,
\]
\[
\delta\bigl(f_u,f_b\ \text{同时为 }(1)(2)\text{ 型}\bigr)=\frac12,
\]
\[
\delta\bigl(f_u,f_b\ \text{同时不可约}\bigr)=\frac13.
\]
\end{theorem}
\begin{proof}
由定理 \ref{thm:xi-terminal-zm-stokes-leyang-shared-artin-representation}，
\[
\QQ(u)=\QQ(b)=:K,
\]
并且 \(K/\QQ\) 的伽罗瓦闭包就是同一个 \(S_3\)-扩张 \(L/\QQ\)。因此
\[
\Gal(L/K)
\]
同时是 \(u\) 与 \(b\) 这两组三共轭根中任一根的稳定子。换言之，\(\Gal(L/\QQ)\) 在 \(f_u\) 的三根集与在 \(f_b\) 的三根集上的置换表示都同构于其在左陪集空间
\[
\Gal(L/\QQ)\big/\Gal(L/K)
\]
上的自然三点传递作用。故对任意未分歧素数 \(p\)，Frobenius 共轭类
\[
\Frob_p(L/\QQ)
\]
在这两个三点集合上诱导出完全相同的循环型。

而对三次多项式，模 \(p\) 的分解型正由该循环型决定：恒等类对应全分裂 \((1)(1)(1)\)，换位类对应 \((1)(2)\)，三循环类对应不可约 \((3)\)。于是对一切
\[
p\nmid \Disc(f_u)\Disc(f_b)
\]
都得到
\[
f_u\bmod p
\quad\text{与}\quad
f_b\bmod p
\]
具有相同分解型。由上面的判别式公式，这正等价于
\[
p\nmid 2\cdot 3\cdot 11\cdot 37\cdot 109.
\]

最后，对共同分裂域 \(L/\QQ\) 应用 Chebotarev 密度定理。因为
\[
|S_3|=6,
\]
其三类共轭类的大小分别为
\[
1,\qquad 3,\qquad 2,
\]
故三种同步分解事件的自然密度分别为
\[
\frac{1}{6},\qquad \frac{3}{6}=\frac12,\qquad \frac{2}{6}=\frac13.
\]
证毕。
\end{proof}

\noindent
由此，纤维同伦理相图、连续容量层析、奇素数 prime-shadow 衰减与 Lee--Yang 三次分裂统计便进一步收束为同一类有限可核验对象：最大非可缩层的损失由模 \(6\) 与 Fibonacci 差额完全控制，容量曲线本身已经是纤维多重度谱的完整坐标，而所有奇素数单位根扭曲的总谱损失则只留下一个严格正的 Euler 常数；与此同时，\(u\)-三次与 \(b\)-三次在全部好素数上的分解型同步表明，这两种观测并非仅共享同一三次数域，而是在算术局部层面也保持逐点同谱。

\paragraph{Lee--Yang 分支圆柱的 \texorpdfstring{$2$}{2}-进重整化与相位/激活芽的阶数分离}
在五分支 \(T_2(E_{\min})\)-torsor 分解、规范 \(B\)-进仿射作用、单位根扭曲的四阶局部展开以及激活谱支的唯一小根性都已固定之后，branch tower 的 dyadic 圆柱几何与输出位势的局部解析芽还可继续压缩为四条闭式：分支塔在逆极限中形成对规范 Gödel 仿射半群闭合的 clopen cylinder 体系，而 phase twist 与 activation 在原点附近则分属不同的零点阶层级；由此，模 \(2\) 奇偶扭曲与高模奇素数扭曲的强度分离，以及 \(u=1\) 邻域简单根的解析诞生，都可被写成严格的局部结论。

\begin{theorem}[Lee--Yang 分支圆柱体系的 \texorpdfstring{$2$}{2}-进仿射重整化]\label{thm:xi-time-part62dg-leyang-cylinder-affine-renormalization}
沿用定理 \ref{thm:xi-time-part62d-leyang-branch-inverse-limit-five-torsors}、定理 \ref{thm:xi-time-part62d-leyang-branch-wreath-affine} 与定理 \ref{thm:xi-time-part9g-canonical-faithful-minimal-radix-threshold} 的记号。记
\[
T:=T_2(E_{\min})\cong \ZZ_2^2.
\]
在五个分支分量上分别选取一条相容提升
\[
\mathbf P_i\in \mathcal R_\infty^{(i)}
\qquad (0\le i\le 4),
\]
从而把
\[
\mathcal R_\infty
\cong
\bigsqcup_{i=0}^{4}\bigl(\mathbf P_i+T\bigr)
\]
视为五片 \(2\)-进仿射平面。对每个 \(n\ge 0\)，定义深度 \(n\) 的分支圆柱体系
\[
\mathcal R_n^{\mathrm{cyl}}
:=
\bigsqcup_{i=0}^{4}\bigl(\mathbf P_i+2^nT\bigr).
\]
则：
\begin{enumerate}
  \item 对每个 \(n\ge 0\)，第 \(i\) 个圆柱 \(\mathbf P_i+2^nT\) 在第 \(n\) 层投影下精确对应于有限 branch 余类
  \[
  P_i+E_{\min}[2^n].
  \]
  因而 \(\mathcal R_n^{\mathrm{cyl}}\) 恰是第 \(n\) 层 Lee--Yang 分支集的五分量 clopen cylinder 表示。
  \item 若 \(B=2^L\)，并仍记
  \[
  \Psi(r,k)(x):=B^k x+\operatorname{ev}_B(r)
  \qquad (x\in T)
  \]
  为各分量在仿射模型 \(\mathbf P_i+T\) 上搬运后的作用公式，则对任意 \(n\ge 0\)、任意 \(0\le i\le 4\) 都有
  \[
  \Psi(r,k)\bigl(\mathbf P_i+2^nT\bigr)
  =
  B^k\mathbf P_i+\operatorname{ev}_B(r)+2^{n+Lk}T.
  \]
  因而
  \[
  \Psi(r,k)\bigl(\mathcal R_n^{\mathrm{cyl}}\bigr)
  =
  \bigcup_{i=0}^{4}
  \Bigl(
  B^k\mathbf P_i+\operatorname{ev}_B(r)+2^{n+Lk}T
  \Bigr).
  \]
\end{enumerate}
换言之，Lee--Yang 分支塔在逆极限中并非仅给出有限层离散点列，而是形成一个对规范 Gödel 仿射半群闭合的 \(2\)-进 clopen 圆柱体系。
\end{theorem}
\begin{proof}
由定理 \ref{thm:xi-time-part62d-leyang-branch-inverse-limit-five-torsors}，每个分量 \(\mathcal R_\infty^{(i)}\) 都是 \(T\) 上的主齐性空间；一旦选定 \(\mathbf P_i\)，便可把它写成仿射余类 \(\mathbf P_i+T\)。

对第一条，只需考察逆极限的第 \(n\) 层坐标投影。由
\[
T=\varprojlim_n E_{\min}[2^n]
\]
的定义，在识别 \(T\cong \ZZ_2^2\) 之后，第 \(n\) 层投影的核正是 \(2^nT\)。因此在第 \(i\) 个分量上，具有相同第 \(n\) 层坐标的全部点恰构成同一个圆柱 \(\mathbf P_i+2^nT\)，而该固定坐标正是有限层 branch 余类 \(P_i+E_{\min}[2^n]\)。于是 \(\mathcal R_n^{\mathrm{cyl}}\) 正是第 \(n\) 层分支集的五分量 clopen cylinder 表示。

对第二条，由定理 \ref{thm:xi-time-part62d-leyang-branch-wreath-affine}，在所选仿射坐标下，各分量上搬运后的仿射作用仍由
\[
x\longmapsto B^k x+\operatorname{ev}_B(r)
\]
给出。故
\[
\Psi(r,k)\bigl(\mathbf P_i+2^nT\bigr)
=
B^k\mathbf P_i+\operatorname{ev}_B(r)+B^k(2^nT).
\]
又因 \(B=2^L\)，故
\[
B^k(2^nT)=2^{n+Lk}T,
\]
从而得到所示公式。把五个分量并起来便得到 \(\Psi(r,k)(\mathcal R_n^{\mathrm{cyl}})\) 的表达式。证毕。
\end{proof}

\begin{theorem}[相位扭曲芽与激活芽的零点阶严格分离]\label{thm:xi-time-part62dg-phase-activation-germ-order-separation}
令
\[
g(t):=-\log\left|\frac{\rho(e^{it})}{\lambda}\right|,
\qquad
h(s):=\lambda_{\mathrm{act}}(1+s),
\qquad
\lambda:=\rho(1)=\varphi^2.
\]
则不存在局部实解析双射 \(\phi\)（满足 \(\phi(0)=0\) 且 \(\phi'(0)\neq 0\)）与非零常数 \(c\)，使得
\[
h(s)=c\,g(\phi(s))
\]
在 \(s=0\) 的邻域内成立。若再记
\[
q(t):=1-\left|\frac{\rho(e^{it})}{\lambda}\right|,
\]
则同样不存在局部实解析双射把 \(h\) 与 \(q\) 的芽等同。因而相位读出与权重激活在原点附近属于两种零点阶严格不同的解析谱类。
\end{theorem}
\begin{proof}
由推论 \ref{cor:real-input-40-collision-twist-fourth}，
\[
\log\left|\frac{\rho(e^{it})}{\lambda}\right|
=
-\frac{6\sqrt5-5}{250}\,t^2+O(t^4),
\]
故
\[
g(t)=\frac{6\sqrt5-5}{250}\,t^2+O(t^4),
\]
从而 \(g\) 在原点的零点阶为 \(2\)。

另一方面，由注 \ref{rem:real-input-40-activated-branch-series}，
\[
h(s)
=
s+3s^3-s^4+18s^5-22s^6+O(s^7),
\]
故 \(h\) 在原点的零点阶为 \(1\)。

而局部实解析双射复合以及乘以非零常数都不改变零点阶，因此不可能存在
\[
h(s)=c\,g(\phi(s))
\]
这样的局部表示。

再看第二部分。由
\[
\exp\bigl(-g(t)\bigr)
=
\left|\frac{\rho(e^{it})}{\lambda}\right|
\]
与 \(g(t)=O(t^2)\)，可展开得
\[
q(t)=1-e^{-g(t)}=g(t)+O(g(t)^2)
=
\frac{6\sqrt5-5}{250}\,t^2+O(t^4),
\]
故 \(q\) 在原点的零点阶同样为 \(2\)。于是 \(h\) 与 \(q\) 的芽也不可能经局部实解析双射而等同。证毕。
\end{proof}

\begin{corollary}[模 \texorpdfstring{$2$}{2} 奇偶扭曲与高模奇素数扭曲的强度分离]\label{cor:xi-time-part62dg-parity-vs-large-prime-twist-separation}
设
\[
\rho_-:=\rho\!\bigl(W(-1)\bigr),
\qquad
r_p:=\left|\rho\!\left(e^{2\pi i/p}\right)\right|
\qquad (p\ge 3\ \text{为奇素数}),
\]
并记
\[
\lambda:=\rho(1)=\varphi^2.
\]
则
\[
\frac{\rho_-}{\lambda}\approx 0.8424525579<1,
\]
并且存在 \(p_0\)，使得对一切奇素数 \(p\ge p_0\) 都有
\[
\rho_-<r_p<\lambda.
\]
进一步，
\[
\log\frac{r_p}{\lambda}
=
-\frac{6\sqrt5-5}{250}\left(\frac{2\pi}{p}\right)^2
+\frac{7+24\sqrt5}{75000}\left(\frac{2\pi}{p}\right)^4
+O(p^{-6}).
\]
因而
\[
\frac{r_p}{\lambda}\to 1
\qquad (p\to\infty,\ p\ \text{为奇素数}),
\]
而模 \(2\) 奇偶扭曲保留一个固定的 \(O(1)\) 谱壁；二者不处于同一强度尺度。
\end{corollary}
\begin{proof}
由定理 \ref{thm:xi-real-input-40-collision-mod2-trace-three-spectrum}，\(\rho_-\) 是三次方程
\[
x^3-2x^2-1=0
\]
的最大实根；数值上有
\[
\rho_-\approx 2.2055694304,
\qquad
\frac{\rho_-}{\lambda}\approx 0.8424525579.
\]

另一方面，推论 \ref{cor:real-input-40-collision-twist-fourth} 在 \(t_p=2\pi/p\) 处给出
\[
\log\frac{r_p}{\lambda}
=
-\frac{6\sqrt5-5}{250}\left(\frac{2\pi}{p}\right)^2
+\frac{7+24\sqrt5}{75000}\left(\frac{2\pi}{p}\right)^4
+O(p^{-6}),
\]
因此 \(r_p/\lambda\to 1\)。

又由定理 \ref{thm:xi-time-part9p-odd-prime-shadow-product} 的严格上界部分，对每个奇素数 \(p\) 都有
\[
0<r_p<\lambda.
\]
由于 \(\rho_-/\lambda<1\)，从 \(r_p/\lambda\to 1\) 立刻推出：存在 \(p_0\) 使得对全部奇素数 \(p\ge p_0\)，都有
\[
\frac{r_p}{\lambda}>\frac{\rho_-}{\lambda},
\]
即 \(r_p>\rho_-\)。综上得到
\[
\rho_-<r_p<\lambda
\]
对一切 \(p\ge p_0\) 成立。证毕。
\end{proof}

\begin{theorem}[激活谱支的 Weierstrass 分离与简单根诞生]\label{thm:xi-time-part62dg-activated-branch-weierstrass-separation}
令
\[
G(\lambda,u)=\lambda^8F(\lambda^{-1},u),
\]
其中
\[
\Delta(z,u)=(1-uz^2)F(z,u).
\]
则在 \((\lambda,u)=(0,1)\) 的某邻域内，存在唯一解析函数 \(\lambda_{\mathrm{act}}(u)\) 与唯一解析函数 \(\widetilde G(\lambda,u)\)，使得
\[
G(\lambda,u)
=
\bigl(\lambda-\lambda_{\mathrm{act}}(u)\bigr)\widetilde G(\lambda,u),
\qquad
\widetilde G(0,1)=1,
\]
并且
\[
\lambda_{\mathrm{act}}(u)
=
(u-1)+3(u-1)^3-(u-1)^4+18(u-1)^5-22(u-1)^6
+O\!\bigl((u-1)^7\bigr).
\]
特别地，存在 \(\varepsilon,\delta>0\)，使得当 \(|u-1|<\delta\) 时，圆盘 \(|\lambda|<\varepsilon\) 内恰有且仅有一条零点分支，即
\[
\lambda=\lambda_{\mathrm{act}}(u);
\]
其余全部谱支与该圆盘保持正距离。因而 \(u=1\) 附近发生的不是分支碰撞，而是一条简单特征值从 \(\lambda=0\) 的解析诞生。
\end{theorem}
\begin{proof}
定理 \ref{thm:xi-output-potential-activated-branch-rigidity} 已给出 \(\lambda_{\mathrm{act}}(u)\) 的唯一小根性与 Weierstrass 因式分解；注 \ref{rem:real-input-40-activated-branch-series} 则给出了所需的高阶展开。故只需说明
\[
\widetilde G(0,1)=1.
\]

为此，由命题 \ref{prop:real-input-40-activated-branch-linear}，
\[
G(0,1)=0,
\qquad
\partial_\lambda G(0,1)=1\neq 0.
\]
而把 \(\lambda=\lambda_{\mathrm{act}}(u)\) 代入分解式并对 \(\lambda\) 求导，可得
\[
\partial_\lambda G\bigl(\lambda_{\mathrm{act}}(u),u\bigr)
=
\widetilde G\bigl(\lambda_{\mathrm{act}}(u),u\bigr).
\]
在 \((\lambda,u)=(0,1)\) 处代入便得到
\[
\widetilde G(0,1)=\partial_\lambda G(0,1)=1.
\]

由于 \(\widetilde G\) 在 \((0,1)\) 邻域解析且在该点非零，存在 \(\varepsilon,\delta>0\) 使得当 \(|u-1|<\delta\)、\(|\lambda|<\varepsilon\) 时，
\[
\widetilde G(\lambda,u)\neq 0.
\]
于是 \(G(\lambda,u)=0\) 在这一小圆盘内只能由
\[
\lambda=\lambda_{\mathrm{act}}(u)
\]
给出；其余零点分支全部被排除在外。故 \(u=1\) 邻域的局部谱结构并非多支碰撞，而是一条简单根的解析诞生。证毕。
\end{proof}

\noindent
由此，Lee--Yang 分支逆极限的五分量 torsor 几何与输出位势在原点附近的解析芽结构便被进一步压缩到同一条可审计链上：一方面，branch tower 在 \(2\)-进 Tate 模块中形成对仿射半群稳定的 clopen cylinder 体系；另一方面，相位扭曲、模 \(2\) 奇偶壁与激活谱支则通过零点阶、固定谱壁及 Weierstrass 分离呈现出彼此不可混同的局部层级。

\paragraph{激活谱支的无穷远逃逸、剥离谱隙与碰撞配分函数闭式}
在激活分支的逆正规形、Weierstrass 分离与碰撞势的三次压力方程已经建立之后，还可把 \(u=1\) 邻域的零尺度根、无穷远逃逸根、剥离后的解析稳定性以及四态有限体积配分函数进一步压缩为下列五条结论。前两条把局部小根重写为无穷远 Laurent 分支；第三条给出剥离后的统一谱隙；后两条则把同一碰撞势在有限维压力骨架上的配分函数完全闭式化。

\begin{conclusion}[激活分支与 \texorpdfstring{$\lambda=u-1$}{lambda=u-1} 的三阶接触及其逆正规形]\label{con:xi-time-part62dga-activated-third-order-contact-reversion}
沿用命题 \ref{prop:real-input-40-activated-branch-linear} 的记号，记
\[
h:=u-1,
\qquad
\lambda:=\lambda_{\mathrm{act}}(u).
\]
则有显式展开
\[
\lambda_{\mathrm{act}}(u)
=
h+3h^3-h^4+18h^5-22h^6+O(h^7),
\]
从而
\[
\lambda_{\mathrm{act}}(u)-h
=
3h^3-h^4+18h^5-22h^6+O(h^7).
\]
因此曲线 \(\lambda=\lambda_{\mathrm{act}}(u)\) 与直线 \(\lambda=u-1\) 在 \((u,\lambda)=(1,0)\) 处恰有三阶接触。并且其局部逆关系满足
\[
h
=
\lambda-3\lambda^3+\lambda^4+9\lambda^5+\lambda^6+O(\lambda^7).
\]
\end{conclusion}
\begin{proof}
正向展开由注 \ref{rem:real-input-40-activated-branch-series} 直接给出，而逆正规形正是定理 \ref{thm:xi-time-part62d-activated-branch-inverse-normalform} 的结论。由于二次项严格缺失而三次项系数为 \(3\neq 0\)，故 \(\lambda_{\mathrm{act}}(u)-(u-1)\) 的消失阶恰为 \(3\)，亦即两条曲线在 \((1,0)\) 处具有恰好三阶接触。证毕。
\end{proof}

\begin{conclusion}[逃向无穷远的激活根分支与 \texorpdfstring{$z_{\mathrm{act}}$}{z_act} 的 Laurent 展开]\label{con:xi-time-part62dga-activated-escape-root-laurent}
沿用分解
\[
\Delta(z,u)=(1-uz^2)F(z,u),
\qquad
G(\lambda,u)=\lambda^8F(\lambda^{-1},u),
\]
并定义
\[
z_{\mathrm{act}}(u):=\lambda_{\mathrm{act}}(u)^{-1}.
\]
则存在 \(\varepsilon>0\)，使得在穿孔邻域
\[
0<|u-1|<\varepsilon
\]
内，方程 \(\Delta(z,u)=0\) 恰有一条满足 \(|z(u)|\to\infty\)（当 \(u\to 1\)）的亚纯根分支，即 \(z_{\mathrm{act}}(u)\)。其 Laurent 展开为
\[
z_{\mathrm{act}}(u)
=
\frac{1}{u-1}
-3(u-1)
+(u-1)^2
-9(u-1)^3
+16(u-1)^4
-32(u-1)^5
+O((u-1)^6).
\]
特别地，沿实轴有
\[
u\downarrow 1 \Longrightarrow z_{\mathrm{act}}(u)\to +\infty,
\qquad
u\uparrow 1 \Longrightarrow z_{\mathrm{act}}(u)\to -\infty.
\]
\end{conclusion}
\begin{proof}
由命题 \ref{prop:real-input-40-activated-branch-linear} 与定理 \ref{thm:xi-time-part62dg-activated-branch-weierstrass-separation}，\(\lambda_{\mathrm{act}}(u)\) 是 \(G(\lambda,u)=0\) 在 \((0,1)\) 邻域内唯一的小根分支，且当 \(u\neq 1\) 足够接近 \(1\) 时 \(\lambda_{\mathrm{act}}(u)\neq 0\)。因此
\[
F(z_{\mathrm{act}}(u),u)=0,
\qquad
z_{\mathrm{act}}(u)=\lambda_{\mathrm{act}}(u)^{-1},
\]
并且 \(|z_{\mathrm{act}}(u)|\to\infty\) 当 \(u\to 1\)。因子 \(1-uz^2\) 仅贡献两条有界分支 \(z=\pm u^{-1/2}\)，故 \(\Delta(z,u)=0\) 中逃向无穷远的根分支只能有这一条。
\par
为求 Laurent 展开，只需把命题 \ref{prop:real-input-40-pressure-freezing} 中的显式多项式 \(G(\lambda,u)\) 再向高一阶做系数匹配，即得
\[
\lambda_{\mathrm{act}}(u)
=
h+3h^3-h^4+18h^5-22h^6+114h^7+O(h^8),
\qquad
h:=u-1.
\]
于是
\[
z_{\mathrm{act}}(u)
=
\frac{1}{h}\,
\frac{1}{1+3h^2-h^3+18h^4-22h^5+114h^6+O(h^7)},
\]
对右端作 Laurent 展开便得到
\[
z_{\mathrm{act}}(u)
=
h^{-1}-3h+h^2-9h^3+16h^4-32h^5+O(h^6).
\]
最后一条由主项 \(h^{-1}\) 直接推出。证毕。
\end{proof}

\begin{conclusion}[剥离激活分支后的统一谱隙与解析稳定性]\label{con:xi-time-part62dga-deflated-uniform-gap}
沿用 Weierstrass 分解
\[
G(\lambda,u)=\bigl(\lambda-\lambda_{\mathrm{act}}(u)\bigr)\widetilde G(\lambda,u).
\]
则存在常数 \(\eta>0\) 与 \(\delta_0>0\)，使得对一切 \(|u-1|<\delta_0\) 都有：
\begin{enumerate}
  \item \(\widetilde G(\lambda,u)\) 在 \(|\lambda|<\eta\) 内解析且无零点；
  \item 除激活分支 \(\lambda_{\mathrm{act}}(u)\) 外，\(G(\lambda,u)=0\) 的其余七个根全部满足
  \[
  |\lambda|\ge \eta.
  \]
\end{enumerate}
因而在 \(u=1\) 邻域内，先剥离激活分支再处理 \(\widetilde G\) 会把零尺度根与 \(O(1)\) 其余谱支严格分离。
\end{conclusion}
\begin{proof}
定理 \ref{thm:xi-time-part62dg-activated-branch-weierstrass-separation} 已给出解析分解
\[
G(\lambda,u)
=
\bigl(\lambda-\lambda_{\mathrm{act}}(u)\bigr)\widetilde G(\lambda,u),
\qquad
\widetilde G(0,1)=1.
\]
另一方面，\(\lambda=0\) 是 \(G(\lambda,1)\) 的简单根，故 \(G(\lambda,1)=0\) 的其余七个根都与原点保持正距离。取
\[
\eta<
\frac12\min\{|\zeta|:\ \zeta\neq 0,\ G(\zeta,1)=0\}.
\]
由多项式根对参数的连续依赖性，缩小到充分小的 \(\delta_0\) 后，当 \(|u-1|<\delta_0\) 时，除 \(\lambda_{\mathrm{act}}(u)\) 外的其余七个根仍全部落在 \(|\lambda|\ge \eta\) 上。由于 \(\widetilde G(\cdot,u)\) 的零点恰是这七个根，故 \(\widetilde G(\lambda,u)\) 在 \(|\lambda|<\eta\) 内无零点。证毕。
\end{proof}

\begin{conclusion}[碰撞势四态配分函数的四阶递推与有理生成函数]\label{con:xi-time-part62dga-collision-partition-recurrence-ogf}
沿用命题 \ref{prop:real-input-40-collision-pressure} 的四态权矩阵表示。记
\[
K:=
\begin{pmatrix}
1&1\\
1&0
\end{pmatrix},
\qquad
A:=K\otimes K,
\qquad
D(u):=\operatorname{diag}(1,1,1,u),
\qquad
W(u):=D(u)A,
\]
并令
\[
\mathbf 1:=(1,1,1,1)^{\mathsf T}.
\]
定义长度 \(n\) 的配分函数
\[
Z_n(u):=\mathbf 1^{\mathsf T}W(u)^n\mathbf 1
\qquad(n\ge 0).
\]
则对一切 \(n\ge 0\) 都有精确递推
\[
Z_{n+4}(u)
=
Z_{n+3}(u)+(u+3)Z_{n+2}(u)+Z_{n+1}(u)-uZ_n(u).
\]
其初值为
\[
Z_0(u)=4,\qquad
Z_1(u)=u+8,\qquad
Z_2(u)=5u+20,\qquad
Z_3(u)=u^2+15u+48.
\]
因而普通生成函数
\[
\mathcal Z_u(t):=\sum_{n\ge 0}Z_n(u)t^n
\]
满足完全闭式
\[
\mathcal Z_u(t)
=
\frac{4+(u+4)t-ut^3}{1-t-(u+3)t^2-t^3+ut^4}.
\]
\end{conclusion}
\begin{proof}
由命题 \ref{prop:real-input-40-collision-pressure}，
\[
\det(\lambda I-W(u))
=(\lambda+1)\bigl(\lambda^3-2\lambda^2-(u+1)\lambda+u\bigr)
=\lambda^4-\lambda^3-(u+3)\lambda^2-\lambda+u.
\]
因此 Cayley--Hamilton 定理给出
\[
W(u)^4-W(u)^3-(u+3)W(u)^2-W(u)+uI=0.
\]
左乘 \(\mathbf 1^{\mathsf T}W(u)^n\)、右乘 \(\mathbf 1\)，便得四阶递推。
\par
又由
\[
W(u)=
\begin{pmatrix}
1&1&1&1\\
1&0&1&0\\
1&1&0&0\\
u&0&0&0
\end{pmatrix},
\]
直接计算得
\[
W(u)\mathbf 1=(4,2,2,u)^{\mathsf T},
\qquad
W(u)^2\mathbf 1=(8+u,6,6,4u)^{\mathsf T},
\]
\[
W(u)^3\mathbf 1=(20+5u,14+u,14+u,8u+u^2)^{\mathsf T}.
\]
逐项求和即得
\[
Z_0(u)=4,\qquad
Z_1(u)=u+8,\qquad
Z_2(u)=5u+20,\qquad
Z_3(u)=u^2+15u+48.
\]
最后把递推乘以 \(t^{n+4}\) 并对 \(n\ge 0\) 求和，代入上述初值得
\[
\bigl(1-t-(u+3)t^2-t^3+ut^4\bigr)\mathcal Z_u(t)
=
4+(u+4)t-ut^3,
\]
从而得到所示闭式。证毕。
\end{proof}

\begin{conclusion}[无权截面上的 \texorpdfstring{$F_{n+3}^2$}{F_{n+3}^2} 平方公式与三阶降阶递推]\label{con:xi-time-part62dga-collision-partition-fibonacci-square}
在结论 \ref{con:xi-time-part62dga-collision-partition-recurrence-ogf} 的记号下，取 \(u=1\)。若 \(F_0=0,F_1=1\) 为 Fibonacci 数列，则对一切 \(n\ge 0\) 都有
\[
Z_n(1)=F_{n+3}^2.
\]
相应地，序列 \(\{Z_n(1)\}_{n\ge 0}\) 满足三阶降阶递推
\[
Z_{n+3}(1)=2Z_{n+2}(1)+2Z_{n+1}(1)-Z_n(1)
\qquad(n\ge 0).
\]
该三阶递推严格低于结论 \ref{con:xi-time-part62dga-collision-partition-recurrence-ogf} 在 \(u=1\) 处给出的四阶递推，反映了无权截面上的额外谱约化。
\end{conclusion}
\begin{proof}
记
\[
\mathbf 1_2:=(1,1)^{\mathsf T},
\]
则 \(\mathbf 1=\mathbf 1_2\otimes \mathbf 1_2\)。当 \(u=1\) 时，
\[
W(1)=K\otimes K,
\]
于是
\[
Z_n(1)
=
(\mathbf 1_2^{\mathsf T}\otimes \mathbf 1_2^{\mathsf T})
(K^n\otimes K^n)
(\mathbf 1_2\otimes \mathbf 1_2)
=
\bigl(\mathbf 1_2^{\mathsf T}K^n\mathbf 1_2\bigr)^2.
\]
再记
\[
s_n:=\mathbf 1_2^{\mathsf T}K^n\mathbf 1_2.
\]
由于 \(K^2=K+I\)，序列 \(s_n\) 满足 Fibonacci 递推
\[
s_{n+2}=s_{n+1}+s_n,
\qquad
s_0=2,\quad s_1=3.
\]
故 \(s_n=F_{n+3}\)，从而
\[
Z_n(1)=s_n^2=F_{n+3}^2.
\]
最后，由 Binet 公式可知 \(\{F_{n+3}^2\}\) 的三条特征根为
\[
\varphi^2,\qquad -1,\qquad \varphi^{-2},
\]
故其特征多项式为
\[
(x+1)(x^2-3x+1)=x^3-2x^2-2x+1,
\]
对应递推正是
\[
Z_{n+3}(1)=2Z_{n+2}(1)+2Z_{n+1}(1)-Z_n(1).
\]
证毕。
\end{proof}

\noindent
由此，\(u=1\) 邻域的激活分支不再只是“最高次项消失”所留下的一条小根：它在 \(z\)-平面上对应唯一的无穷远逃逸根，在剥离之后与其余七根保持统一谱隙；与此同时，同一碰撞势的四态压力骨架又给出完全显式的有限体积配分函数，而无权截面则进一步塌缩为 Fibonacci 数列的平方。于是局部激活几何、解析稳定性与全局配分闭式便被压缩为同一条有限维、可审计的闭合链条。

\paragraph{全息圆柱--球同构与 Smith 亏损的单极点压缩}
在 H.160--H.161 的 \(\,2^L\)-进全息地址链、以及 H.154--H.156 的 Smith 模核与 Euler 因子链已经固定之后，还可把前缀几何与模数亏损进一步压缩为下列四条结论。前两条把事件流前缀精确识别为 Tate 模块中的余类球，并由此恢复环境维数与零测性；后两条则说明归一化模核函数完全塌缩到 \(D\) 的除子格上，而相应 Dirichlet 级数只保留一个 Riemann \(\zeta\) 极点与有限多个坏素数修正。

\begin{conclusion}[\texorpdfstring{$2^L$}{2^L}-进全息地址的圆柱--球严格同构]\label{con:xi-time-part62dgb-hologram-cylinder-ball-isometry}
沿用定理 \ref{thm:xi-time-part62d-hologram-image-cantor-homeomorphism} 与定理 \ref{thm:killo-2adic-holographic-exact-cylinder-separation} 的记号。设 \(E\) 为有限事件集，
\[
\mathrm{code}:E\hookrightarrow \NN,\qquad
M:=\max_{e\in E}\mathrm{code}(e),\qquad
B:=2^L>M,
\]
并记
\[
T:=T_2(E_{\min})\cong \ZZ_2^2,\qquad
v\in T\setminus 2T.
\]
对任意 \(\mathbf e=(e_t)_{t\ge 1}\in E^{\NN}\)，定义
\[
X(\mathbf e):=\sum_{t\ge 1}\mathrm{code}(e_t)\,B^{t-1}v,
\qquad
\mathcal X_{E,v}:=X(E^{\NN})=\Lambda_{E,v}\subset T.
\]
再对 \(\mathbf e,\mathbf e'\in E^{\NN}\) 定义公共前缀长度
\[
\mathrm{pref}_n(\mathbf e):=(e_1,\dots,e_n),
\qquad
\ell(\mathbf e,\mathbf e')
:=
\sup\bigl\{n\ge 0:\ \mathrm{pref}_n(\mathbf e)=\mathrm{pref}_n(\mathbf e')\bigr\}
\in \NN\cup\{\infty\},
\]
并定义事件流上的柱超距
\[
d_E(\mathbf e,\mathbf e')
:=
\begin{cases}
0,& \mathbf e=\mathbf e',\\[2pt]
B^{-\ell(\mathbf e,\mathbf e')},& \mathbf e\neq \mathbf e'.
\end{cases}
\]
在 \(T\) 上，对 \(z\neq 0\) 记
\[
\nu_B(z):=\max\{n\ge 0:\ z\in B^nT\},
\]
并定义环境 \(B\)-进超距
\[
d_T(x,y)
:=
\begin{cases}
0,& x=y,\\[2pt]
B^{-\nu_B(x-y)},& x\neq y.
\end{cases}
\]
则 \(X:(E^{\NN},d_E)\to (\mathcal X_{E,v},d_T)\) 是等距嵌入。更精确地，对任意长度 \(n\) 的前缀
\[
a=(a_1,\dots,a_n)\in E^n,
\qquad
[a]:=\{\mathbf e\in E^{\NN}:\ \mathrm{pref}_n(\mathbf e)=a\},
\]
若记
\[
x_a:=\sum_{t=1}^{n}\mathrm{code}(a_t)\,B^{t-1}v,
\]
则有精确圆柱--球公式
\[
X([a])=\bigl(x_a+B^nT\bigr)\cap \mathcal X_{E,v}.
\]
\end{conclusion}
\begin{proof}
若 \(\mathbf e=\mathbf e'\)，则两侧距离都为零。现设 \(\mathbf e\neq \mathbf e'\)，并记
\[
\ell(\mathbf e,\mathbf e')=n.
\]
则首次分歧出现在第 \(n+1\) 位。由于 \(B=2^L\) 且 \(v\notin 2T\)，必有 \(v\notin BT\)，故可应用定理 \ref{thm:killo-2adic-holographic-exact-cylinder-separation}，得到
\[
X(\mathbf e)-X(\mathbf e')
\in
B^nT\setminus B^{n+1}T.
\]
因此
\[
\nu_B\bigl(X(\mathbf e)-X(\mathbf e')\bigr)=n,
\]
从而
\[
d_T\bigl(X(\mathbf e),X(\mathbf e')\bigr)=B^{-n}=d_E(\mathbf e,\mathbf e').
\]
这就证明了等距性。
\par
再证圆柱--球公式。若 \(\mathbf e\in[a]\)，则
\[
X(\mathbf e)-x_a
=
\sum_{t>n}\mathrm{code}(e_t)\,B^{t-1}v\in B^nT,
\]
故
\[
X([a])\subseteq \bigl(x_a+B^nT\bigr)\cap\mathcal X_{E,v}.
\]
反之，若 \(X(\mathbf e)\in x_a+B^nT\)，则 \(X(\mathbf e)\equiv x_a\pmod{B^nT}\)。任取 \(\mathbf f\in[a]\)，同样有 \(X(\mathbf f)\equiv x_a\pmod{B^nT}\)，故
\[
X(\mathbf e)-X(\mathbf f)\in B^nT.
\]
由定理 \ref{thm:killo-2adic-holographic-exact-cylinder-separation} 的同余--前缀等价式，得到
\[
\mathrm{pref}_n(\mathbf e)=\mathrm{pref}_n(\mathbf f)=a,
\]
即 \(\mathbf e\in[a]\)。故
\[
X([a])=\bigl(x_a+B^nT\bigr)\cap \mathcal X_{E,v}.
\]
证毕。
\end{proof}

\begin{conclusion}[事件流全息像的精确 Hausdorff/Minkowski 维数与 Haar 零测]\label{con:xi-time-part62dgb-hologram-exact-dimension-zero-haar}
沿用结论 \ref{con:xi-time-part62dgb-hologram-cylinder-ball-isometry} 的记号。对每个 \(n\ge 0\)，记
\[
N_{E,v}(B^{-n})
\]
为用环境空间 \(T\) 中半径 \(B^{-n}\) 的闭球覆盖 \(\mathcal X_{E,v}\) 所需的最小球数。则有精确覆盖公式
\[
N_{E,v}(B^{-n})=|E|^n.
\]
因而
\[
\dim_H(\mathcal X_{E,v})
=
\underline{\dim}_M(\mathcal X_{E,v})
=
\overline{\dim}_M(\mathcal X_{E,v})
=
\frac{\log |E|}{\log B}.
\]
此外，若 \(\mu_T\) 为 \(T\cong \ZZ_2^2\) 的归一化 Haar 概率测度，则
\[
\mu_T(\mathcal X_{E,v})=0.
\]
\end{conclusion}
\begin{proof}
由结论 \ref{con:xi-time-part62dgb-hologram-cylinder-ball-isometry}，对每个 \(n\ge 0\)，每个长度 \(n\) 的前缀 \(a\in E^n\) 都对应一个非空集合
\[
X([a])=\bigl(x_a+B^nT\bigr)\cap \mathcal X_{E,v}.
\]
不同前缀对应的这类集合彼此不交：若 \(a\neq b\) 而两者落在同一个 \(B^nT\)-陪集上，则任取 \(\mathbf e\in[a]\)、\(\mathbf f\in[b]\)，便有
\[
X(\mathbf e)-X(\mathbf f)\in B^nT,
\]
这与结论 \ref{con:xi-time-part62dgb-hologram-cylinder-ball-isometry} 的前缀等距性矛盾。由于长度 \(n\) 的前缀共有 \(|E|^n\) 个，故任意半径 \(B^{-n}\) 的环境球至多覆盖其中一个 cylinder；另一方面，这 \(|E|^n\) 个球的并恰好覆盖整个 \(\mathcal X_{E,v}\)。因此
\[
N_{E,v}(B^{-n})=|E|^n.
\]
于是上下 Minkowski 维数都满足
\[
\underline{\dim}_M(\mathcal X_{E,v})
=
\overline{\dim}_M(\mathcal X_{E,v})
=
\lim_{n\to\infty}\frac{\log N_{E,v}(B^{-n})}{-\log(B^{-n})}
=
\frac{\log|E|}{\log B}.
\]
再取 \(E\) 上的均匀 Bernoulli 测度，并记其地址推前为 \(\mu_{\mathrm{unif}}\)。由定理 \ref{thm:xi-time-part62dc-hologram-bernoulli-exact-dimensional}，
\[
\dim_H(\mu_{\mathrm{unif}})=\frac{\log|E|}{\log B},
\]
且 \(\mu_{\mathrm{unif}}\) 支撑于 \(\mathcal X_{E,v}\)。故
\[
\dim_H(\mathcal X_{E,v})\ge \frac{\log|E|}{\log B}.
\]
另一方面，Hausdorff 维数不超过上 Minkowski 维数，因而
\[
\dim_H(\mathcal X_{E,v})\le \frac{\log|E|}{\log B}.
\]
两式合并即得 Hausdorff 维数公式。
\par
最后，由 \(T/B^nT\cong (\ZZ/2^{Ln}\ZZ)^2\)，每个半径 \(B^{-n}\) 的环境球的 Haar 质量都等于
\[
|T/B^nT|^{-1}=2^{-2Ln}=B^{-2n}.
\]
因此
\[
\mu_T(\mathcal X_{E,v})
\le
N_{E,v}(B^{-n})\,B^{-2n}
=
\left(\frac{|E|}{B^2}\right)^n.
\]
又因 \(\mathrm{code}\) 为单射且 \(B>M\)，故
\[
|E|\le M+1\le B,
\]
从而
\[
\frac{|E|}{B^2}\le \frac{1}{B}<1.
\]
令 \(n\to\infty\) 即得 \(\mu_T(\mathcal X_{E,v})=0\)。证毕。
\end{proof}

\begin{conclusion}[模数信息亏损函数的 \texorpdfstring{$D$}{D}-饱和因子化与除子格严格单调嵌入]\label{con:xi-time-part62dgb-smith-d-saturation-divisor-embedding}
设 \(A\in M_{m\times n}(\ZZ)\) 满足 \(\operatorname{rank}_{\QQ}(A)=m\)，其 Smith 不变量为
\[
d_1\mid d_2\mid\cdots\mid d_m.
\]
定义
\[
D:=\operatorname{lcm}(d_1,\dots,d_m),
\qquad
\Lambda_A(b):=\frac{\#\ker(A_b)}{b^{\,n-m}}
\qquad(b\ge 1),
\]
其中 \(A_b:(\ZZ/b\ZZ)^n\to(\ZZ/b\ZZ)^m\) 为模 \(b\) 诱导映射。则有：
\begin{enumerate}
  \item \textbf{\(D\)-饱和因子化：}
  \[
  \Lambda_A(b)=\Lambda_A\bigl(\gcd(b,D)\bigr)
  \qquad(b\ge 1).
  \]
  \item \textbf{互素乘法性：}若 \(\gcd(b_1,b_2)=1\)，则
  \[
  \Lambda_A(b_1b_2)=\Lambda_A(b_1)\Lambda_A(b_2).
  \]
  \item \textbf{除子格严格单调性：}若 \(g,h\mid D\) 且 \(g\mid h\) 但 \(g\neq h\)，则
  \[
  \Lambda_A(g)<\Lambda_A(h).
  \]
  \item \textbf{值域完全压缩到 \(D\) 的除子格：}
  \[
  \{\Lambda_A(b):\ b\ge 1\}
  =
  \{\Lambda_A(g):\ g\mid D\},
  \qquad
  \#\{\Lambda_A(b):\ b\ge 1\}=\tau(D),
  \]
  其中 \(\tau(D)\) 为 \(D\) 的除子个数。并且
  \[
  \max_{b\ge 1}\Lambda_A(b)=\prod_{i=1}^{m}d_i,
  \]
  且当且仅当 \(D\mid b\) 时取到该最大值。
\end{enumerate}
\end{conclusion}
\begin{proof}
由推论 \ref{cor:killo-kernel-size-general-modulus}，
\[
\Lambda_A(b)=\prod_{i=1}^{m}\gcd(d_i,b).
\]
对每个素数 \(p\)，记
\[
e_i(p):=v_p(d_i),\qquad e_p:=v_p(D)=\max_i e_i(p),
\qquad
f_p(k):=\sum_{i=1}^{m}\min(k,e_i(p)).
\]
则有素数分解
\[
\Lambda_A(b)=\prod_{p\mid D}p^{\,f_p(v_p(b))}.
\]
由于 \(\min(v_p(b),e_i(p))=\min(v_p(\gcd(b,D)),e_i(p))\)，立即得到
\[
\Lambda_A(b)=\Lambda_A(\gcd(b,D)),
\]
这证明了第一条。
\par
若 \(\gcd(b_1,b_2)=1\)，则对每个 \(i\) 有
\[
\gcd(d_i,b_1b_2)=\gcd(d_i,b_1)\gcd(d_i,b_2),
\]
逐项相乘便得
\[
\Lambda_A(b_1b_2)=\Lambda_A(b_1)\Lambda_A(b_2),
\]
故第二条成立。
\par
再证第三条。若 \(g,h\mid D\) 且 \(g\mid h\) 但 \(g\neq h\)，则存在素数 \(p\mid D\) 使
\[
v_p(g)=k<\ell=v_p(h)\le e_p.
\]
而对 \(0\le t<e_p\)，有
\[
f_p(t+1)-f_p(t)=\#\{i:\ e_i(p)\ge t+1\}>0,
\]
因为 \(t+1\le e_p\) 时至少存在某个 \(i\) 使 \(e_i(p)=e_p\ge t+1\)。因此 \(f_p\) 在 \(\{0,1,\dots,e_p\}\) 上严格递增，故
\[
v_p(\Lambda_A(h))=f_p(\ell)>f_p(k)=v_p(\Lambda_A(g)).
\]
其余素数处指数均不减，从而
\[
\Lambda_A(g)<\Lambda_A(h).
\]
\par
第四条中，值域包含式由第一条立即推出。现证单射：若 \(g,h\mid D\) 且
\[
\Lambda_A(g)=\Lambda_A(h),
\]
则对每个 \(p\mid D\) 有
\[
f_p(v_p(g))=v_p(\Lambda_A(g))=v_p(\Lambda_A(h))=f_p(v_p(h)).
\]
由 \(f_p\) 的严格递增性可知 \(v_p(g)=v_p(h)\) 对所有 \(p\mid D\) 都成立，从而 \(g=h\)。故
\[
g\longmapsto \Lambda_A(g)\qquad(g\mid D)
\]
是单射，因而
\[
\#\{\Lambda_A(b):\ b\ge 1\}=\#\{g:\ g\mid D\}=\tau(D).
\]
最后，由
\[
\Lambda_A(g)=\prod_{i=1}^{m}\gcd(d_i,g)\le \prod_{i=1}^{m}d_i
\qquad(g\mid D)
\]
知最大值至多为 \(\prod_i d_i\)。当 \(D\mid b\) 时，\(\gcd(d_i,b)=d_i\) 对每个 \(i\) 都成立，故
\[
\Lambda_A(b)=\prod_{i=1}^{m}d_i.
\]
反之，若 \(\Lambda_A(b)=\prod_i d_i\)，则由第一条可把 \(b\) 替换为 \(g:=\gcd(b,D)\mid D\)，并得到
\[
\Lambda_A(g)=\prod_{i=1}^{m}d_i=\Lambda_A(D).
\]
由第三条的严格单调性只能有 \(g=D\)，即 \(D\mid b\)。证毕。
\end{proof}

\begin{conclusion}[模数核增长 Dirichlet 级数的单一 \texorpdfstring{$\zeta$}{zeta} 极点与有限 Euler 修正]\label{con:xi-time-part62dgb-kernel-dirichlet-single-zeta-pole}
沿用结论 \ref{con:xi-time-part62dgb-smith-d-saturation-divisor-embedding} 的记号，并定义
\[
\mathcal Z_A(s):=\sum_{b\ge 1}\frac{\#\ker(A_b)}{b^s},
\qquad
u:=n-m.
\]
对每个素数 \(p\mid D\)，记
\[
e_p:=v_p(D),\qquad
\nu_p:=\sum_{i=1}^{m}e_i(p),\qquad
f_p(k):=\sum_{i=1}^{m}\min(k,e_i(p)).
\]
则有显式分解
\[
\mathcal Z_A(s)
=
\zeta(s-u)\prod_{p\mid D}P_{A,p}(p^{-s}),
\]
其中每个局部修正因子都是多项式
\[
P_{A,p}(Y)
:=
\bigl(1-p^uY\bigr)\sum_{k=0}^{e_p-1}p^{\,f_p(k)+uk}Y^k
+p^{\,\nu_p+u e_p}Y^{e_p}.
\]
因此 \(\mathcal Z_A(s)\) 延拓为整个复平面上的亚纯函数，并且其唯一可能的极点来自单一因子 \(\zeta(s-u)\)；特别地，\(\mathcal Z_A(s)\) 至多在
\[
s=u+1=n-m+1
\]
处具有一个简单极点。
\end{conclusion}
\begin{proof}
由推论 \ref{cor:killo-kernel-size-general-modulus}，
\[
\#\ker(A_b)=b^u\Lambda_A(b)
=
b^u\prod_{i=1}^{m}\gcd(d_i,b).
\]
右端关于 \(b\) 是积性函数，故 \(\mathcal Z_A\) 具有 Euler 乘积
\[
\mathcal Z_A(s)=\prod_p \mathcal Z_{A,p}(s),
\qquad
\mathcal Z_{A,p}(s):=\sum_{k\ge 0}\frac{\#\ker(A_{p^k})}{p^{ks}}.
\]
对 \(k\ge 0\)，有
\[
\#\ker(A_{p^k})
=
p^{uk+f_p(k)},
\]
故
\[
\mathcal Z_{A,p}(s)=\sum_{k\ge 0}p^{\,f_p(k)+uk}p^{-ks}.
\]
若 \(p\nmid D\)，则 \(e_i(p)=0\) 对一切 \(i\) 都成立，故 \(f_p(k)=0\)，于是
\[
\mathcal Z_{A,p}(s)=\sum_{k\ge 0}p^{k(u-s)}=\frac{1}{1-p^{u-s}}.
\]
若 \(p\mid D\)，则当 \(k\ge e_p\) 时，
\[
f_p(k)=\sum_{i=1}^{m}e_i(p)=\nu_p.
\]
因此
\[
\mathcal Z_{A,p}(s)
=
\sum_{k=0}^{e_p-1}p^{\,f_p(k)+uk}p^{-ks}
+p^{\nu_p}\sum_{k\ge e_p}p^{k(u-s)}.
\]
对尾项作几何求和得
\[
\mathcal Z_{A,p}(s)
=
\sum_{k=0}^{e_p-1}p^{\,f_p(k)+uk}p^{-ks}
+\frac{p^{\,\nu_p+u e_p}p^{-e_p s}}{1-p^{u-s}}.
\]
令 \(Y:=p^{-s}\)，并把上式乘以 \((1-p^uY)\)，便得到
\[
\bigl(1-p^uY\bigr)\mathcal Z_{A,p}(s)=P_{A,p}(Y),
\]
即
\[
\mathcal Z_{A,p}(s)=\frac{P_{A,p}(p^{-s})}{1-p^{u-s}}.
\]
于是
\[
\mathcal Z_A(s)
=
\prod_{p\nmid D}\frac{1}{1-p^{u-s}}
\prod_{p\mid D}\frac{P_{A,p}(p^{-s})}{1-p^{u-s}}
=
\zeta(s-u)\prod_{p\mid D}P_{A,p}(p^{-s}).
\]
由于 \(p^{-s}\) 是整函数而 \(P_{A,p}\) 为多项式，故有限乘积
\[
\prod_{p\mid D}P_{A,p}(p^{-s})
\]
是整函数。因而 \(\mathcal Z_A(s)\) 的全部亚纯结构都来自单一因子 \(\zeta(s-u)\)，从而唯一可能的极点位于 \(s=u+1\)，且其阶数至多为一。证毕。
\end{proof}

\noindent
由此，H.160--H.161 的全息地址链与 H.154--H.156 的 Smith 模核链便在同一层结论中完成了进一步压缩：事件流前缀与环境 Tate 模块中的 \(B^nT\)-余类球严格对应，故时间复杂度能够被精确转写为空间覆盖指数；而模数核增长则完全塌缩到 \(D\) 的除子格与有限个坏素数局部因子之上，其全球 Dirichlet 解析结构只保留一个 Riemann \(\zeta\) 主极点。于是前缀几何、模层亏损与 Euler 乘积三者被统一压缩为同一组有限证书。

\paragraph{模核总数主项的有限 Euler 修正常数闭式}
在模数核增长 Dirichlet 级数的单一 \(\zeta\) 极点、Smith 周期平均律以及局部 \(p\)-窗指数列都已经固定之后，原始核计数部分和的主项常数还可进一步完全显式化为有限 Euler 乘积。

\begin{theorem}[整数线性投影核计数的 Tauberian 主项与有限 Euler 修正常数]\label{thm:xi-time-part62dgba-smith-kernel-tauberian-finite-euler-leading-constant}
沿用结论 \ref{con:xi-time-part62dgb-kernel-dirichlet-single-zeta-pole} 的记号。设
\[
A\in M_{m\times n}(\ZZ),
\qquad
\operatorname{rank}_{\QQ}(A)=m,
\]
其 Smith 不变量为
\[
d_1\mid d_2\mid \cdots \mid d_m.
\]
对每个 \(b\ge 1\)，记
\[
A_b:(\ZZ/b\ZZ)^n\longrightarrow (\ZZ/b\ZZ)^m
\]
为 \(A\) 在模 \(b\) 下的诱导映射，并定义
\[
a_b:=\#\ker(A_b),
\qquad
D:=\operatorname{lcm}(d_1,\dots,d_m),
\qquad
u:=n-m.
\]
对每个素数 \(p\mid D\)，再记
\[
e_p:=v_p(D),
\qquad
\nu_p:=\sum_{i=1}^{m}e_i(p),
\qquad
f_p(k):=\sum_{i=1}^{m}\min(k,e_i(p)).
\]
则有精确渐近
\[
\sum_{b\le X}a_b
\sim
\frac{G_A}{u+1}X^{u+1}
=
\frac{G_A}{n-m+1}X^{n-m+1}
\qquad (X\to\infty),
\]
其中有限 Euler 常数为
\[
G_A
:=
\prod_{p\mid D}
\left(
(1-p^{-1})\sum_{k=0}^{e_p-1}p^{f_p(k)-k}
+p^{\nu_p-e_p}
\right).
\]
\end{theorem}
\begin{proof}
由结论 \ref{con:xi-time-part62dgb-kernel-dirichlet-single-zeta-pole}，
\[
\mathcal Z_A(s):=\sum_{b\ge 1}\frac{a_b}{b^s}
=
\zeta(s-u)\prod_{p\mid D}P_{A,p}(p^{-s}),
\]
其中局部因子满足
\[
P_{A,p}(Y)
:=
\bigl(1-p^uY\bigr)\sum_{k=0}^{e_p-1}p^{f_p(k)+uk}Y^k
+p^{\nu_p+u e_p}Y^{e_p}.
\]
在极点
\[
s_0:=u+1=n-m+1
\]
处，代入 \(Y=p^{-s_0}=p^{-u-1}\) 即得
\[
p^uY=p^{-1},
\qquad
p^{f_p(k)+uk}Y^k=p^{f_p(k)-k},
\qquad
p^{\nu_p+u e_p}Y^{e_p}=p^{\nu_p-e_p}.
\]
因此
\[
P_{A,p}(p^{-s_0})
=
(1-p^{-1})\sum_{k=0}^{e_p-1}p^{f_p(k)-k}
+p^{\nu_p-e_p}.
\]
从而
\[
\Res_{s=s_0}\mathcal Z_A(s)
=
\prod_{p\mid D}P_{A,p}(p^{-s_0})
=
G_A.
\]

另一方面，结论 \ref{con:xi-time-part62d-smith-kernel-summatory-periodic-average} 已证明
\[
\sum_{b\le X}a_b
\sim
\frac{\Res_{s=s_0}\mathcal Z_A(s)}{s_0}\,X^{s_0}.
\]
将上式中的留数与 \(s_0=u+1=n-m+1\) 代入，便得到
\[
\sum_{b\le X}a_b
\sim
\frac{G_A}{u+1}X^{u+1}
=
\frac{G_A}{n-m+1}X^{n-m+1}.
\]
证毕。
\end{proof}

\noindent
因此，模数核增长的主项只在幂指数上记住自由度差 \(n-m\)，而全部 Smith 复杂性则被压缩进一个仅由坏素数与有限 \(p\)-窗指数列构成的 Euler 常数 \(G_A\)。

\paragraph{ZG 的矩阵 Euler 刚性与 Lee--Yang 五分支圆柱的芯核塌缩}
在 Zeckendorf--Gödel 素数码的矩阵 Euler 传递链、自然密度主极点，以及 Lee--Yang 分支塔的 affine Tate sheet 与 clopen cylinder 体系都已固定之后，还可把算术侧与 dyadic 几何侧进一步压缩为下列两条定理。前一条把 \(ZG\) 的解析结构明确识别为矩阵 Euler 而非标量 Euler；后一条则说明，围绕各分支所选相容 lift 的 dyadic 圆柱过滤，其芯核极限恰收缩为五个刚性基点。

\begin{theorem}[\texorpdfstring{$ZG$}{ZG} 的 Dirichlet 级数具有矩阵 Euler 积而不存在标量 Euler 积]\label{thm:xi-time-part62dgc-zg-matrix-euler-no-scalar-product}
\leanverified{paper\_xi\_time\_part62dgc\_zg\_matrix\_euler\_no\_scalar\_product}
记
\[
D_{\mathrm{ZG}}(s):=\sum_{n\in \mathrm{ZG}}n^{-s},
\qquad
\mathbf 1:=(1,1)^{\mathsf T},
\qquad
e_1:=(1,0)^{\mathsf T}.
\]
则成立：
\begin{enumerate}
  \item 对 \(\Re(s)>1\)，有严格的 \(2\times 2\) 矩阵 Euler 积表示
  \[
  D_{\mathrm{ZG}}(s)
  =
  \lim_{N\to\infty}
  \mathbf 1^{\mathsf T}
  \Biggl(
  \prod_{i=1}^{N}
  \begin{pmatrix}
  1&1\\[2pt]
  p_i^{-s}&0
  \end{pmatrix}
  \Biggr)e_1.
  \]
  \item 集合 \(\mathrm{ZG}\subset\NN\) 具有严格正的自然密度
  \[
  \delta_{\mathrm{ZG}}>0.
  \]
  \item 不存在任何标量 Euler 因子分解
  \[
  D_{\mathrm{ZG}}(s)
  =
  \prod_{p}
  \bigl(1+a_p p^{-s}+a_{p,2}p^{-2s}+\cdots\bigr)
  \]
  能在某个右半平面上表示同一 Dirichlet 级数并保持其系数仍等于集合指示函数 \(1_{\mathrm{ZG}}(n)\)。等价地，\(1_{\mathrm{ZG}}\) 不是乘法函数。
\end{enumerate}
因此，\(\mathrm{ZG}\) 的 Euler 结构本质上来自素数指标链上的一维 hard-core 约束，而不能压缩为独立标量局部因子的乘积。
\end{theorem}
\begin{proof}
第一条正是定理 \ref{thm:killo-zg-dirichlet-matrix-euler} 的内容。第二条由定理 \ref{con:xi-zg-density-closed-form-tail-bound} 与定理 \ref{thm:xi-zg-abel-residue-log-density} 直接给出。

现证第三条。若存在所示标量 Euler 因子分解，并且其 Dirichlet 系数仍为 \(1_{\mathrm{ZG}}(n)\)，则对应系数函数必为乘法函数。另一方面，\(\mathrm{ZG}\) 的元素由素数指标链上的独立集编码：对每个 \(i\ge 1\)，单点集合 \(\{i\}\) 与 \(\{i+1\}\) 都是合法独立集，故
\[
p_i\in \mathrm{ZG},
\qquad
p_{i+1}\in \mathrm{ZG}.
\]
但 \(\{i,i+1\}\) 含相邻指标，不再是合法独立集，因此
\[
p_ip_{i+1}\notin \mathrm{ZG}.
\]
于是
\[
1_{\mathrm{ZG}}(p_i)\,1_{\mathrm{ZG}}(p_{i+1})=1,
\qquad
1_{\mathrm{ZG}}(p_ip_{i+1})=0,
\]
这与乘法性矛盾。故不存在这样的标量 Euler 因子分解。证毕。
\end{proof}

\begin{theorem}[Lee--Yang 五分支 dyadic 圆柱过滤的芯核塌缩]\label{thm:xi-time-part62dgc-leyang-fivecore-dyadic-core-collapse}
沿用定理 \ref{thm:xi-time-part62dg-leyang-cylinder-affine-renormalization} 的记号，记
\[
T:=T_2(E_{\min})\cong \ZZ_2^2,
\]
并在五个分支分量上分别选取相容提升
\[
\mathbf P_i\in \mathcal R_\infty^{(i)}
\qquad (0\le i\le 4).
\]
定义深度 \(n\) 的分支圆柱并
\[
\mathcal C_n
:=
\bigsqcup_{i=0}^{4}\bigl(\mathbf P_i+2^nT\bigr).
\]
则有：
\begin{enumerate}
  \item \(\mathcal C_{n+1}\subset \mathcal C_n\) 对一切 \(n\ge 0\) 成立；
  \item 若在每一片 affine sheet \(\mathbf P_i+T\) 上赋予由 \(T\) 的归一化 Haar 概率测度搬运得到的测度，并把五片测度求和记为 \(\mu_{\mathrm{Haar}}\)，则
  \[
  \mu_{\mathrm{Haar}}(\mathcal C_n)=5\cdot 4^{-n};
  \]
  \item 其芯核交集恰为
  \[
  \bigcap_{n\ge 0}\mathcal C_n=\{\mathbf P_0,\mathbf P_1,\mathbf P_2,\mathbf P_3,\mathbf P_4\}.
  \]
\end{enumerate}
特别地，由有限层正则四叉地址结构，
\[
\deg R(y_n)=5\cdot 4^n,
\]
而这一级数增长相对于上述圆柱过滤只记录五个所选 lift 周围的 dyadic 地址熵，并不再产生额外的芯核厚度。
\end{theorem}
\begin{proof}
第一条直接由
\[
2^{n+1}T\subset 2^nT
\]
推出。

对第二条，由定理 \ref{thm:xi-time-part62dg-leyang-cylinder-affine-renormalization}，每个 \(\mathbf P_i+2^nT\) 都是第 \(i\) 个 affine sheet 中的 clopen cylinder。又因
\[
T/2^nT\cong (\ZZ/2^n\ZZ)^2,
\]
其指数为 \(4^n\)，故每个半径 \(2^{-n}\) 的 cylinder 的归一化 Haar 质量都等于 \(4^{-n}\)。五片 sheet 彼此不交，因此
\[
\mu_{\mathrm{Haar}}(\mathcal C_n)
=
\sum_{i=0}^{4}\mu_{\mathrm{Haar}}(\mathbf P_i+2^nT)
=
5\cdot 4^{-n}.
\]

再证第三条。对每个固定 \(i\)，有
\[
\bigcap_{n\ge 0}\bigl(\mathbf P_i+2^nT\bigr)
=
\mathbf P_i+\bigcap_{n\ge 0}2^nT.
\]
而 \(T\cong\ZZ_2^2\) 是 Hausdorff 完备的 \(2\)-进模，故
\[
\bigcap_{n\ge 0}2^nT=\{0\}.
\]
于是
\[
\bigcap_{n\ge 0}\bigl(\mathbf P_i+2^nT\bigr)=\{\mathbf P_i\}.
\]
由于五片 sheet 彼此不交，便得到
\[
\bigcap_{n\ge 0}\mathcal C_n=\{\mathbf P_0,\mathbf P_1,\mathbf P_2,\mathbf P_3,\mathbf P_4\}.
\]

最后，由推论 \ref{cor:xi-terminal-zm-leyang-finite-branch-regular-4ary-address}，每个有限层分支分量 \(\Pi_{i,n}\) 都与一个 \(4^n\) 点地址块对应，故
\[
|\Pi_n^{\mathrm{fin}}|
=
\sum_{i=0}^{4}|\Pi_{i,n}|
=
5\cdot 4^n.
\]
而按定理 \ref{thm:xi-time-part62d-leyang-branch-inverse-limit-five-torsors} 的记号，\(R(y_n)\) 的有限 branch 次数正由 \(\Pi_n^{\mathrm{fin}}\) 计数，故
\[
\deg R(y_n)=|\Pi_n^{\mathrm{fin}}|=5\cdot 4^n.
\]
证毕。
\end{proof}

\noindent
由此，prime-index hard-core 约束与 Lee--Yang dyadic branch 过滤在同一层面完成了进一步压缩：\(\mathrm{ZG}\) 的 Dirichlet 解析结构只能以矩阵 Euler 的方式忠实呈现，而五分支圆柱体系在选定相容 lift 之后又收缩为五点芯核。于是，算术耦合的不可标量化与 dyadic 地址的五核芯化在同一组有限证书中闭合。

\paragraph{ZG 的素数 J-分式、黄金 Cantor 地址与 Archimedean/\texorpdfstring{$2^L$}{2^L}-进规模分离}
在 \(ZG\) 的矩阵 Euler 递推、自然密度的前缀单调极限以及 H.161 的 \(2^L\)-进地址像同胚都已固定之后，还可把算术密度与非阿基米德地址几何进一步压缩为下列三条结果：前缀密度具有显式的 prime \(J\)-分式与无穷乘积表示；相应地址像恰是黄金均值子移位的 Cantor 型实现，并具有 \(\log\varphi/\log B\) 的精确维数；从而 \(ZG\) 在 Archimedean 计数侧的正密度与其在 \(2^L\)-进地址侧的严格亚一维几何之间形成一条明确的尺度分离。

\begin{theorem}[ZG 自然密度的素数 \texorpdfstring{$J$}{J}-分式、无穷乘积与定量正性]\label{thm:xi-time-part62dgd-zg-density-prime-jfraction-positive}
记
\[
\mathbf 1:=(1,1)^{\mathsf T},
\qquad
e_1:=(1,0)^{\mathsf T},
\]
并定义
\[
\binom{A_n}{B_n}
:=
\left(
\prod_{i=1}^{n}
\begin{pmatrix}
\dfrac{p_i}{p_i+1} & \dfrac{p_i}{p_i+1}\\[1ex]
\dfrac{1}{p_i+1} & 0
\end{pmatrix}
\right)e_1,
\qquad
t_n:=A_n+B_n,
\qquad
r_n:=\frac{B_n}{A_n},
\]
且约定 \(r_0:=0\)。则 \(\delta_{\mathrm{ZG}}=\mathrm{dens}(ZG)\) 满足：
\begin{enumerate}
  \item 对每个 \(n\ge 1\)，有递推
  \[
  r_n=\frac{1}{p_n(1+r_{n-1})},
  \]
  从而得到 \(J\)-分式展开
  \[
  r_n=
  \cfrac{1}{p_n\!\left(1+\cfrac{1}{p_{n-1}\!\left(1+\cfrac{1}{\ddots\left(1+\frac1{p_1}\right)}\right)}\right)}.
  \]
  \item 对每个 \(n\ge 1\)，有
  \[
  t_n
  =
  t_{n-1}\left(1-\frac{r_{n-1}}{(p_n+1)(1+r_{n-1})}\right).
  \]
  因而
  \[
  \delta_{\mathrm{ZG}}
  =
  \frac1{\zeta(2)}\lim_{n\to\infty}t_n
  =
  \frac1{\zeta(2)}
  \prod_{n\ge 2}
  \left(1-\frac{r_{n-1}}{(p_n+1)(1+r_{n-1})}\right).
  \]
  \item 若记
  \[
  \delta_n:=\frac{t_n}{\zeta(2)},
  \]
  则 \(\delta_n\downarrow \delta_{\mathrm{ZG}}\in(0,1)\)，并且对每个 \(n\ge 1\) 有显式尾估计
  \[
  0<\delta_n-\delta_{\mathrm{ZG}}
  \le
  \delta_n\sum_{j\ge n+1}\frac{1}{p_{j-1}(p_j+1)}.
  \]
\end{enumerate}
\end{theorem}
\begin{proof}
命题 \ref{prop:xi-zg-prefix-recursion-monotone-tail-certificate} 中的归一化前缀量 \((a_n,b_n)\) 满足
\[
(a_0,b_0)=(1,0),
\]
\[
a_n=(a_{n-1}+b_{n-1})\frac{p_n}{p_n+1},
\qquad
b_n=\frac{a_{n-1}}{p_n+1}.
\]
而这里按矩阵定义的 \((A_n,B_n)\) 由乘法递推也满足同一组初值与同一递推式。故
\[
A_n=a_n,
\qquad
B_n=b_n,
\qquad
t_n=A_n+B_n=u_n.
\]
于是
\[
r_n=\frac{B_n}{A_n}
=
\frac{1}{p_n}\frac{A_{n-1}}{A_{n-1}+B_{n-1}}
=
\frac{1}{p_n(1+r_{n-1})},
\]
第一条得证；递归展开即给出所述 \(J\)-分式。

再由
\[
t_n=A_n+B_n
=
\frac{p_n}{p_n+1}(A_{n-1}+B_{n-1})
+\frac{1}{p_n+1}A_{n-1},
\]
以及
\[
A_{n-1}=\frac{t_{n-1}}{1+r_{n-1}},
\]
可得
\[
t_n
=
t_{n-1}\frac{p_n+\frac{1}{1+r_{n-1}}}{p_n+1}
=
t_{n-1}\left(1-\frac{r_{n-1}}{(p_n+1)(1+r_{n-1})}\right).
\]
这就得到第二条的递推式。又由 \(t_1=1\) 与命题
\ref{prop:xi-zg-prefix-recursion-monotone-tail-certificate}
中的
\[
\delta_{\mathrm{ZG}}=\frac{1}{\zeta(2)}\lim_{n\to\infty}t_n
\]
可知对 \(n\ge 2\) 迭代上式便得到所示无穷乘积表示。

现证第三条。由 \(r_0=0\) 与第一条递推归纳得到
\[
0<r_n\le \frac1{p_n}\qquad (n\ge 1).
\]
记
\[
c_j:=\frac{r_{j-1}}{(p_j+1)(1+r_{j-1})}\qquad (j\ge 2).
\]
则
\[
0<c_j\le \frac{1}{p_{j-1}(p_j+1)}.
\]
由第二条可得
\[
t_\infty
=
t_n\prod_{j\ge n+1}(1-c_j),
\]
从而
\[
0<t_n-t_\infty
=
t_n\left(1-\prod_{j\ge n+1}(1-c_j)\right).
\]
对正数列 \((c_j)\) 应用不等式
\[
1-\prod_{j\ge n+1}(1-c_j)\le \sum_{j\ge n+1}c_j
\]
便得
\[
0<t_n-t_\infty
\le
t_n\sum_{j\ge n+1}\frac{1}{p_{j-1}(p_j+1)}.
\]
两边同时乘以 \(\zeta(2)^{-1}\)，即得
\[
0<\delta_n-\delta_{\mathrm{ZG}}
\le
\delta_n\sum_{j\ge n+1}\frac{1}{p_{j-1}(p_j+1)}.
\]
又命题 \ref{prop:xi-zg-prefix-recursion-monotone-tail-certificate} 已证明 \(t_n\) 单调下降并收敛到 \(\zeta(2)\delta_{\mathrm{ZG}}\)，故 \(\delta_n\downarrow \delta_{\mathrm{ZG}}\in(0,1)\)。证毕。
\end{proof}

\begin{theorem}[ZG 的 \texorpdfstring{$2^L$}{2^L}-进地址像是黄金 Cantor 集，且维数为 \texorpdfstring{$\log\varphi/\log B$}{log phi / log B}]\label{thm:xi-time-part62dgd-zg-address-golden-cantor-dimension}
设
\[
\Sigma_\varphi
:=
\{\,z=(z_k)_{k\ge 1}\in \{0,1\}^{\NN}: z_kz_{k+1}=0\ \forall k\,\},
\]
固定 \(T:=T_2(E_{\min})\) 中一个非零向量 \(v\)，并令 \(B:=2^L\ge 2\)。定义地址映射
\[
X_B:\Sigma_\varphi\longrightarrow T,
\qquad
X_B(z):=\sum_{k\ge 1}z_kB^{k-1}v,
\]
以及像集
\[
\mathcal X_B:=X_B(\Sigma_\varphi).
\]
则成立：
\begin{enumerate}
  \item \(X_B\) 是 \(\Sigma_\varphi\) 到 \(\mathcal X_B\) 的拓扑同胚。若定义
  \[
  \mathcal T:=X_B\circ \sigma\circ X_B^{-1},
  \]
  其中 \(\sigma\) 为黄金均值子移位上的左移，则 \((\mathcal X_B,\mathcal T)\) 与 \((\Sigma_\varphi,\sigma)\) 拓扑共轭。
  \item 在由最长公共前缀等价地由模 \(B^mT\) 同余诱导的 \(B\)-进超度量下，
  \[
  \dim_{\mathrm H}(\mathcal X_B)
  =
  \underline{\dim}_{\mathrm M}(\mathcal X_B)
  =
  \overline{\dim}_{\mathrm M}(\mathcal X_B)
  =
  \frac{\log\varphi}{\log B}.
  \]
  \item 若记 \(N_m(\mathcal X_B)\) 为 \(\mathcal X_B\) 在尺度 \(B^{-m}\) 上的最少球覆盖数，则
  \[
  N_m(\mathcal X_B)=F_{m+2},
  \]
  并且对
  \[
  s:=\frac{\log\varphi}{\log B}
  \]
  有
  \[
  \lim_{m\to\infty}N_m(\mathcal X_B)\,B^{-ms}
  =
  \frac{\varphi^2}{\sqrt5}.
  \]
\end{enumerate}
\end{theorem}
\begin{proof}
取 \(E=\{0,1\}\) 且 \(\mathrm{code}(0)=0\)、\(\mathrm{code}(1)=1\)。由于 \(B=2^L>1\)，定理
\ref{thm:xi-time-part62d-hologram-image-cantor-homeomorphism}
适用于该二元字母表，并给出地址映射
\[
X:E^{\NN}\to \Lambda_{E,v}
\]
是到其像的拓扑同胚。由于 \(\Sigma_\varphi\subset E^{\NN}\) 是闭的移位不变子集，故其限制
\[
X_B:=X|_{\Sigma_\varphi}
:
\Sigma_\varphi\to \mathcal X_B
\]
仍是拓扑同胚。由此定义
\[
\mathcal T:=X_B\circ \sigma\circ X_B^{-1}
\]
便立即得到第一条。

对第二、三条，记 \([\mathbf a]_\varphi\subset \Sigma_\varphi\) 为长度 \(m\) 的合法前缀 \(\mathbf a=(a_1,\dots,a_m)\) 所对应的圆柱集。由定理
\ref{thm:xi-time-part62d-hologram-image-cantor-homeomorphism}
的前缀同余性质，所有 \(z\in[\mathbf a]_\varphi\) 的像都满足
\[
X_B(z)\equiv \sum_{k=1}^{m}a_kB^{k-1}v\pmod{B^mT}.
\]
不同合法前缀对应不同的模 \(B^mT\) 余类，因此在 \(\mathcal X_B\) 上的长度 \(m\) 圆柱集恰给出一族两两不交的半径 \(B^{-m}\) 的闭开球。于是尺度 \(B^{-m}\) 的最少覆盖数恰等于长度 \(m\) 的合法前缀个数。该个数即为标准 Fibonacci 计数
\[
N_m(\mathcal X_B)=F_{m+2}.
\]

由此立得
\[
\underline{\dim}_{\mathrm M}(\mathcal X_B)
=
\overline{\dim}_{\mathrm M}(\mathcal X_B)
=
\lim_{m\to\infty}\frac{\log F_{m+2}}{m\log B}
=
\frac{\log\varphi}{\log B}.
\]
对 Hausdorff 维数，若 \(s>\log\varphi/\log B\)，则长度 \(m\) 圆柱集构成一个直径 \(B^{-m}\) 的覆盖，并满足
\[
F_{m+2}B^{-ms}\longrightarrow 0,
\]
故 \(\mathcal H^s(\mathcal X_B)=0\)。反之，若 \(s<\log\varphi/\log B\)，则任何直径不超过 \(B^{-m}\) 的覆盖都至少需要 \(F_{m+2}\) 个集合，因而其 \(s\)-维代价至少为
\[
F_{m+2}B^{-ms}\longrightarrow \infty.
\]
于是 \(\mathcal H^s(\mathcal X_B)=\infty\)。因此
\[
\dim_{\mathrm H}(\mathcal X_B)=\frac{\log\varphi}{\log B},
\]
第二条得证。

最后，由 Binet 公式
\[
F_{m+2}=\frac{\varphi^{m+2}-(-\varphi^{-1})^{m+2}}{\sqrt5}
\]
可得
\[
F_{m+2}\varphi^{-m}\longrightarrow \frac{\varphi^2}{\sqrt5}.
\]
而
\[
B^{-ms}=\varphi^{-m},
\]
故
\[
N_m(\mathcal X_B)\,B^{-ms}
=
F_{m+2}\varphi^{-m}
\longrightarrow
\frac{\varphi^2}{\sqrt5}.
\]
第三条成立。证毕。
\end{proof}

\begin{corollary}[ZG 的 Archimedean 正密度与 \texorpdfstring{$2^L$}{2^L}-进亚一维地址几何严格分离]\label{cor:xi-time-part62dgd-zg-archimedean-nonarchimedean-scale-separation}
沿用定理 \ref{thm:xi-time-part62dgd-zg-density-prime-jfraction-positive} 与定理 \ref{thm:xi-time-part62dgd-zg-address-golden-cantor-dimension} 的记号。则 \(ZG\) 同时满足
\[
0<\delta_{\mathrm{ZG}}<1,
\]
并且其 \(2^L\)-进地址像 \(\mathcal X_B\) 满足
\[
\dim_{\mathrm H}(\mathcal X_B)
=
\frac{\log\varphi}{\log B}
<1.
\]
因此，\(ZG\) 在整数轴上的自然计数中是正密度对象，而在 \(2^L\)-进地址侧却形成一个严格亚一维的 Cantor 集。
\end{corollary}
\begin{proof}
第一式由定理 \ref{thm:xi-zg-abel-residue-log-density} 与定理
\ref{thm:xi-time-part62dgd-zg-density-prime-jfraction-positive}
给出。第二式是定理
\ref{thm:xi-time-part62dgd-zg-address-golden-cantor-dimension}
的直接结论。又因 \(B=2^L\ge 2>\varphi\)，故
\[
\frac{\log\varphi}{\log B}<1.
\]
证毕。
\end{proof}

\paragraph{有限 ghost 窗的整值重建、方差强迫与钩型湮灭}
Schur 断层的 Frobenius 整性、\(p\)-典型 ghost 可审计性与特征标筛热力学判据，还可进一步闭合为一条有限可恢复且可定量强迫的结论链：局部模 \(p^N\) 断层能够经 Chinese remainder 升格为整值通道重建，而循环型方差既给出非平凡通道的最小振幅与活跃支持下界，也把钩型湮灭压缩为关于低阶压力函数的显式单变量优势条件。

\leanverified{paper\_xi\_time\_part63\_schur\_crt\_exact\_recovery}
\begin{theorem}[有限素数窗的 Chinese remainder 精确恢复：由有限 \texorpdfstring{$p$}{p}-幂 ghost 窗整值重建全部 Schur 通道]\label{thm:xi-time-part63-schur-crt-exact-recovery}
固定分辨率 \(m\) 与整数 \(q\ge 2\)。对每个分拆 \(\lambda\vdash q\)，记
\[
s_\lambda(m):=\frac{T_{q,\lambda}(m)}{f^\lambda}\in\ZZ,
\]
其中 \(T_{q,\lambda}(m)\) 为 Schur 通道迹、\(f^\lambda=\dim S^\lambda\)。
取互异素数幂模
\[
Q:=\prod_{j=1}^{J}p_j^{N_j}.
\]
假设对每个 \(j\)，都已由命题
\ref{prop:pom-schur-channels-mod-pN-from-ptypical-ghost}
所述有限 \(p_j\)-幂 ghost 窗确定全部同余类
\[
s_\lambda(m)\pmod{p_j^{N_j}}
\qquad(\lambda\vdash q).
\]
若
\[
Q>2q!\,2^{mq},
\]
则全部整数
\[
s_\lambda(m)\qquad(\lambda\vdash q)
\]
都由这些有限 ghost 窗唯一精确恢复。因而全部 \(T_{q,\lambda}(m)\) 亦可精确恢复；再由推论
\ref{cor:pom-schur-traces-moments-bijection}，所有低阶碰撞矩
\[
S_1(m),S_2(m),\dots,S_q(m)
\]
亦随之唯一确定。
\end{theorem}
\begin{proof}
由推论 \ref{cor:pom-schur-trace-frobenius-character-identification}，
\[
s_\lambda(m)=\frac{T_{q,\lambda}(m)}{f^\lambda}\in\ZZ.
\]
另一方面，命题
\ref{prop:pom-schur-channels-mod-pN-from-ptypical-ghost}
表明：对每个 \(j\)，有限 \(p_j\)-幂 ghost 窗足以确定
\(
s_\lambda(m)\bmod p_j^{N_j}
\)
对一切 \(\lambda\vdash q\) 的同余类。
因此由 Chinese remainder theorem，可唯一得到
\[
s_\lambda(m)\pmod Q.
\]

剩余只需建立一个绝对值上界。由定理
\ref{thm:pom-schur-cycle-index-tomography}，
\[
s_\lambda(m)
=
\frac{1}{q!}\sum_{\sigma\in S_q}\chi^\lambda(\sigma)\,\mathcal C_\sigma(m),
\]
其中
\[
\mathcal C_\sigma(m)=\prod_{r\ge 1}S_r(m)^{c_r(\sigma)}.
\]
若 \(\sigma\) 的循环型为 \(1^{c_1}2^{c_2}\cdots\)，则
\(
\sum_{r\ge 1}r\,c_r(\sigma)=q
\)。
又由
\[
S_r(m)=\sum_{x\in X_m}d_m(x)^r\le
\left(\sum_{x\in X_m}d_m(x)\right)^r
=2^{mr},
\]
得到
\[
|\mathcal C_\sigma(m)|\le 2^{mq}.
\]
同时对任意 \(\sigma\in S_q\) 都有
\(
|\chi^\lambda(\sigma)|\le f^\lambda\le q!
\)。
因此
\[
|s_\lambda(m)|
\le
\frac{1}{q!}\sum_{\sigma\in S_q}
|\chi^\lambda(\sigma)|\,|\mathcal C_\sigma(m)|
\le
q!\,2^{mq}.
\]
当 \(Q>2q!\,2^{mq}\) 时，模 \(Q\) 的剩余类唯一决定区间
\(
(-Q/2,Q/2)
\)
中的整数；故 \(s_\lambda(m)\) 由其模 \(Q\) 同余类唯一恢复。
再乘回 \(f^\lambda\) 即得 \(T_{q,\lambda}(m)=f^\lambda s_\lambda(m)\) 的精确恢复。
最后，由推论 \ref{cor:pom-schur-traces-moments-bijection}，固定 \(q\) 的全通道迹表与低阶碰撞矩表逐 \(m\) 互译，从而 \(S_1(m),\dots,S_q(m)\) 亦被唯一确定。证毕。
\end{proof}

\begin{theorem}[非平凡 Schur 通道的方差强迫与活跃支持下界]\label{thm:xi-time-part63-schur-variance-forcing-support-lower-bound}
\leanverified{paper\_xi\_time\_part63\_schur\_variance\_forcing\_support\_lower\_bound}
固定 \(q\ge 2\) 与分辨率 \(m\)。定义循环型碰撞函数的平均与方差
\[
\overline{\mathcal C}(m):=\frac{1}{q!}\sum_{\sigma\in S_q}\mathcal C_\sigma(m),
\qquad
V_q(m):=\frac{1}{q!}\sum_{\sigma\in S_q}
\bigl(\mathcal C_\sigma(m)-\overline{\mathcal C}(m)\bigr)^2.
\]
再定义活跃非平凡 Schur 通道集合
\[
\mathcal A_q(m):=
\{\lambda\vdash q:\ \lambda\neq(q),\ T_{q,\lambda}(m)\neq 0\}.
\]
则
\[
V_q(m)=0
\qquad\Longleftrightarrow\qquad
\mathcal A_q(m)=\varnothing.
\]
若 \(V_q(m)>0\)，则必有 \(\mathcal A_q(m)\neq\varnothing\)，并且
\[
\boxed{\;
\max_{\lambda\in\mathcal A_q(m)}
\frac{|T_{q,\lambda}(m)|}{f^\lambda}
\ge
\sqrt{\frac{V_q(m)}{p(q)-1}}\;}
\]
以及
\[
\boxed{\;
|\mathcal A_q(m)|
\ge
\frac{V_q(m)}
{\displaystyle
\max_{\lambda\vdash q,\ \lambda\neq(q)}
\Bigl(\frac{|T_{q,\lambda}(m)|}{f^\lambda}\Bigr)^2}\; } ,
\]
其中 \(p(q)\) 为 \(q\) 的分拆数。
\end{theorem}
\begin{proof}
由定理 \ref{thm:pom-schur-variance-decomposition}，
\[
V_q(m)
=
\sum_{\lambda\vdash q,\ \lambda\neq(q)}
\left(\frac{T_{q,\lambda}(m)}{f^\lambda}\right)^2.
\]
又由于 \(T_{q,\lambda}(m)\in\ZZ\)，上式右端为非负实数之和，故
\[
V_q(m)=0
\qquad\Longleftrightarrow\qquad
T_{q,\lambda}(m)=0\ \text{对所有 }\lambda\neq(q),
\]
这与 \(\mathcal A_q(m)=\varnothing\) 等价；亦即推论
\ref{cor:pom-schur-nontrivial-channel-activation}
的定量版。

现设 \(V_q(m)>0\)，则 \(\mathcal A_q(m)\neq\varnothing\)。由上式立即有
\[
V_q(m)
\le
\bigl(p(q)-1\bigr)
\max_{\lambda\vdash q,\ \lambda\neq(q)}
\left(\frac{|T_{q,\lambda}(m)|}{f^\lambda}\right)^2.
\]
而零通道不影响最大值，故可把右端最大值限制到 \(\mathcal A_q(m)\)，从而得到第一条不等式。

另一方面，
\[
V_q(m)
=
\sum_{\lambda\in\mathcal A_q(m)}
\left(\frac{|T_{q,\lambda}(m)|}{f^\lambda}\right)^2
\le
|\mathcal A_q(m)|
\max_{\lambda\vdash q,\ \lambda\neq(q)}
\left(\frac{|T_{q,\lambda}(m)|}{f^\lambda}\right)^2,
\]
移项即得第二条下界。证毕。
\end{proof}

\begin{theorem}[单变量压力优势强制全部非平凡钩型通道湮灭]\label{thm:xi-time-part63-elementary-pressure-gap-kills-hooks}
\leanverified{paper\_xi\_time\_part63\_elementary\_pressure\_gap\_kills\_hooks}
在定理 \ref{thm:pom-partition-gap-positive-kills-hooks} 的设定下，对 \(q\ge 2\) 定义
\[
\Delta_{\mathrm{elem}}(q):=
P(q)-
\max_{1\le r\le q-1}\bigl(P(r)+(q-r)\log 2\bigr).
\]
若
\[
\Delta_{\mathrm{elem}}(q)>0,
\]
则对所有非平凡钩型分拆
\[
\lambda=(q-k,1^k),\qquad 1\le k\le q-1,
\]
都有
\[
\boxed{\;
B_{q,\lambda}=0,
\qquad
T_{q,\lambda}(m)=0\ \ \forall m\ge 1.\;}
\]
\end{theorem}
\begin{proof}
沿用 \S\ref{subsubsec:pom-schur-character-sieve-thermodynamic} 的分拆压力记号
\[
P(\nu):=\sum_{r\ge 1}c_r(\nu)\,P(r)\qquad(\nu\vdash q).
\]
又由
\[
S_r(m)\le S_1(m)^r=2^{mr}
\]
得
\[
P(r)\le r\log 2\qquad(1\le r\le q).
\]

现任取非平凡分拆 \(\nu\vdash q\)、\(\nu\neq(q)\)。
则存在某个 \(r_0\le q-1\) 使 \(c_{r_0}(\nu)\ge 1\)。
于是
\[
\begin{aligned}
P(\nu)
&=
\sum_{r\ge 1}c_r(\nu)\,P(r)\\
&\le
P(r_0)+
\sum_{r\neq r_0}c_r(\nu)\,r\log 2\\
&=
P(r_0)+(q-r_0)\log 2\\
&\le
\max_{1\le r\le q-1}\bigl(P(r)+(q-r)\log 2\bigr).
\end{aligned}
\]
因此
\[
\Delta_{\mathrm{part}}(q)
:=
P(q)-\max_{\substack{\nu\vdash q\\ \nu\neq(q)}}P(\nu)
\ge
\Delta_{\mathrm{elem}}(q)>0.
\]
由定理 \ref{thm:pom-partition-gap-positive-kills-hooks}，所有非平凡钩型
\(\lambda=(q-k,1^k)\) 必满足 \(B_{q,\lambda}=0\)，等价地
\(
T_{q,\lambda}(m)=0
\)
对一切 \(m\ge 1\) 成立。证毕。
\end{proof}

\noindent
这三条结果共同表明，Schur 断层的有限 \(p\)-进可审计性并不止于局部模数层：一旦收集到足够大的有限素数窗，整值通道与低阶碰撞矩便可被完全重建；而在指数率层面，循环型方差既是非平凡通道激活的精确能量，也把钩型通道的消失归约为关于 \(P(1),\dots,P(q)\) 的显式优势判据。

\paragraph{交替 Schur 特征多项式、Newton 刚性、hook 差分曲率与 Ces\`aro--Fourier 留数恢复}
在 Schur 通道的 Frobenius 标识、Jacobi--Trudi/Euler 展开以及 torsion-\(q\) 准周期主项的 residue-class Fourier 反演已经建立之后，还可把交替通道、第一枚非平凡 hook 通道与循环留数向量进一步压缩为下列四条直接后果。它们分别表明：交替 Schur 通道本身已经给出纤维谱的特征多项式；其 Newton 刚性可作为严格的非均匀谱判据；\((q-1,1)\)-hook 层只是平凡层的一步差分曲率；而准周期主项中的 Fourier 留数又可由单一系数序列的 Ces\`aro--Fourier 投影直接恢复。

\begin{theorem}[交替 Schur 通道给出纤维谱的精确特征多项式]\label{thm:xi-time-part63a-alternating-schur-characteristic-polynomial}
固定分辨率 \(m\)，记
\[
N:=|X_m|,
\qquad
D_m:=\operatorname{diag}\bigl(d_m(x)\bigr)_{x\in X_m}.
\]
对 \(q\ge 0\) 定义交替 Schur 通道
\[
A_q^{(m)}:=
\begin{cases}
1, & q=0,\\[2pt]
T_{q,(1^q)}(m), & 1\le q\le N,\\[2pt]
0, & q>N.
\end{cases}
\]
再定义生成多项式
\[
E_m(z):=\sum_{q=0}^{N}A_q^{(m)}z^q.
\]
则有严格恒等式
\[
\boxed{\;
E_m(z)=\prod_{x\in X_m}\bigl(1+d_m(x)z\bigr).\;}
\]
特别地：
\begin{enumerate}
  \item \(A_q^{(m)}=0\) 当且仅当 \(q>N\)；
  \item \(E_m(z)\) 的全部零点均为负实数，且恰为
  \[
  z=-\frac{1}{d_m(x)}
  \qquad (x\in X_m),
  \]
  并按多重度计入；
  \item 纤维多重度多重集 \(\{d_m(x)\}_{x\in X_m}\) 由单一序列
  \[
  \{A_q^{(m)}\}_{0\le q\le N}
  \]
  唯一恢复。
\end{enumerate}
\end{theorem}
\begin{proof}
枚举
\[
X_m=\{x_1,\dots,x_N\},
\qquad
d_i:=d_m(x_i).
\]
由推论 \ref{cor:pom-schur-trace-frobenius-character-identification}，
\[
T_{q,(1^q)}(m)=s_{(1^q)}(D_m)=e_q(D_m)=e_q(d_1,\dots,d_N),
\]
因为分拆 \((1^q)\) 的 Specht 维数为 \(1\)，且
\[
s_{(1^q)}=e_q.
\]
于是对 \(0\le q\le N\) 有
\[
A_q^{(m)}=e_q(d_1,\dots,d_N).
\]
基本对称函数的标准生成恒等式给出
\[
\sum_{q=0}^{N}e_q(d_1,\dots,d_N)z^q
=
\prod_{i=1}^{N}(1+d_i z),
\]
从而得到所示公式。
又因为 \(d_i>0\) 对一切 \(i\) 成立，故
\[
e_q(d_1,\dots,d_N)>0
\qquad(0\le q\le N),
\]
而对 \(q>N\) 则按定义 \(A_q^{(m)}=0\)。这就证明了第 \(1\) 条。

由于所有 \(d_i>0\)，右端因子分解立即表明全部零点都位于负实轴，并且恰为 \(-1/d_i\)。反之，\(E_m\) 是首项为 \(\prod_i d_i\) 的次数 \(N\) 多项式，而其零点多重集唯一决定 \(\{d_i\}\)；等价地，系数列 \(\{A_q^{(m)}\}_{0\le q\le N}\) 唯一决定该多项式，故也唯一决定纤维多重度多重集。证毕。
\end{proof}

\begin{theorem}[交替 Schur 通道的 Newton 刚性与非均匀谱严格判别]\label{thm:xi-time-part63a-alternating-schur-newton-rigidity}
沿用定理 \ref{thm:xi-time-part63a-alternating-schur-characteristic-polynomial} 的记号，并记
\[
N:=|X_m|.
\]
则对每个 \(1\le q\le N-1\)，成立 Newton 严格不等式
\[
\boxed{\;
\bigl(A_q^{(m)}\bigr)^2
\ge
\frac{q+1}{q}\cdot \frac{N-q+1}{N-q}\,
A_{q-1}^{(m)}A_{q+1}^{(m)}.\;}
\]
并且等号对某个 \(q\) 成立，当且仅当全部纤维多重度完全相等：
\[
d_m(x)\equiv d
\qquad (x\in X_m).
\]
等价地，归一化序列
\[
\frac{A_q^{(m)}}{\binom{N}{q}}
\qquad(0\le q\le N)
\]
严格对数凹；其失去严格性恰当刻画平坦纤维谱情形。
\end{theorem}
\begin{proof}
由定理 \ref{thm:xi-time-part63a-alternating-schur-characteristic-polynomial}，
\[
A_q^{(m)}=e_q(d_1,\dots,d_N)
\qquad(0\le q\le N),
\]
其中 \(d_i>0\)。Newton 不等式断言
\[
\left(\frac{e_q}{\binom{N}{q}}\right)^2
\ge
\frac{e_{q-1}}{\binom{N}{q-1}}
\frac{e_{q+1}}{\binom{N}{q+1}}
\qquad(1\le q\le N-1).
\]
整理组合系数即得
\[
e_q^2
\ge
\frac{\binom{N}{q}^2}{\binom{N}{q-1}\binom{N}{q+1}}\,e_{q-1}e_{q+1}
=
\frac{q+1}{q}\cdot\frac{N-q+1}{N-q}\,e_{q-1}e_{q+1},
\]
这正是所示不等式。

关于等号条件，对正变量组 \((d_1,\dots,d_N)\)，Newton 不等式的经典严格性判定表明：若某一级出现等号，则全部变量相等；反之若 \(d_1=\cdots=d_N=d\)，则
\[
A_q^{(m)}=\binom{N}{q}d^q,
\]
从而归一化序列 \(\binom{N}{q}^{-1}A_q^{(m)}=d^q\) 为几何列，故在每一级都取等号。证毕。
\end{proof}

\begin{theorem}[hook 通道的闭式差分律：第一非平凡 Schur 层由平凡层决定]\label{thm:xi-time-part63a-hook-channel-closed-difference-law}
\leanverified{paper\_xi\_time\_part63a\_hook\_channel\_closed\_difference\_law}
对 \(q\ge 2\) 定义平凡 Schur 通道与第一枚非平凡 hook 通道
\[
H_q(m):=T_{q,(q)}(m)=h_q(D_m),
\qquad
K_q(m):=T_{q,(q-1,1)}(m).
\]
再记
\[
S_1(m):=\sum_{x\in X_m}d_m(x)=|\Omega_m|.
\]
则成立完全显式恒等式
\[
\boxed{\;
K_q(m)=(q-1)\bigl(S_1(m)H_{q-1}(m)-H_q(m)\bigr).\;}
\]
因此对一切 \(q\ge 2\)，都有
\[
\boxed{\;
H_q(m)\le S_1(m)\,H_{q-1}(m).\;}
\]
并且等号成立当且仅当至多存在一条非零纤维；在当前 \(X_m\) 作为非空纤维支撑的记号下，这等价于 \(|X_m|=1\)。
\end{theorem}
\begin{proof}
由推论 \ref{cor:pom-schur-trace-frobenius-character-identification}，
\[
K_q(m)=f^{(q-1,1)}\,s_{(q-1,1)}(D_m),
\qquad
H_q(m)=h_q(D_m).
\]
而 hook 分拆 \((q-1,1)\) 的 Specht 维数为
\[
f^{(q-1,1)}=q-1.
\]
再由 Jacobi--Trudi 恒等式，
\[
s_{(q-1,1)}=h_{q-1}h_1-h_q.
\]
代入并注意
\[
h_1(D_m)=\sum_{x\in X_m}d_m(x)=S_1(m),
\]
便得
\[
K_q(m)=(q-1)\bigl(S_1(m)H_{q-1}(m)-H_q(m)\bigr).
\]

为证明不等式与等号条件，仍写 \(X_m=\{x_1,\dots,x_N\}\) 且 \(d_i=d_m(x_i)\)。由完全对称函数的定义，
\[
h_{q-1}(d)\,h_1(d)
=
\sum_{|\alpha|=q-1}d^\alpha\sum_{i=1}^{N}d_i
=
\sum_{|\beta|=q}\#\operatorname{supp}(\beta)\,d^\beta,
\]
其中
\[
d^\beta:=\prod_{i=1}^{N}d_i^{\beta_i},
\qquad
\operatorname{supp}(\beta):=\{i:\beta_i>0\}.
\]
另一方面，
\[
h_q(d)=\sum_{|\beta|=q}d^\beta.
\]
故
\[
h_{q-1}(d)h_1(d)-h_q(d)
=
\sum_{|\beta|=q}\bigl(\#\operatorname{supp}(\beta)-1\bigr)d^\beta\ge 0.
\]
这就得到
\[
H_q(m)\le S_1(m)H_{q-1}(m).
\]
并且等号成立当且仅当上式右端所有项都为零，即所有出现的单项式都满足 \(\#\operatorname{supp}(\beta)=1\)。这等价于变量组 \((d_1,\dots,d_N)\) 中至多有一个非零项。由于当前记号下 \(x\in X_m\) 当且仅当 \(d_m(x)\ge 1\)，故这又等价于 \(|X_m|=1\)。证毕。
\end{proof}

\begin{theorem}[torsion-\texorpdfstring{$q$}{q} 准周期主项的 Ces\`aro--Fourier 投影公式]\label{thm:xi-time-part63a-torsion-q-cesaro-fourier-projection}\leanverified{paper\_xi\_time\_part63a\_torsion\_q\_cesaro\_fourier\_projection}
沿用定理 \ref{thm:pom-torsion-q-artin-factorization-residue-fourier} 的设定。设
\[
a_n:=[z^n]\zeta(z)
=
\rho^n\sum_{j=0}^{q-1}c_j\omega_q^{jn}
+
O(\rho^n\theta^n),
\qquad
\omega_q:=e^{2\pi i/q},
\qquad
0<\theta<1.
\]
则对每个 \(j=0,1,\dots,q-1\)，都有严格恢复公式
\[
\boxed{\;
c_j=
\lim_{N\to\infty}
\frac{1}{N}\sum_{n=0}^{N-1}\rho^{-n}a_n\,\omega_q^{-jn}.\;}
\]
更强地，对任意 \(r\in\{0,\dots,q-1\}\)，residue-class 主项极限
\[
A_r:=
\lim_{N\to\infty}\rho^{-(qN+r)}a_{qN+r}
\]
存在，并与 \((c_j)_{0\le j\le q-1}\) 互为离散 Fourier 变换：
\[
\boxed{\;
A_r=\sum_{j=0}^{q-1}c_j\,\omega_q^{jr},
\qquad
c_j=\frac{1}{q}\sum_{r=0}^{q-1}A_r\,\omega_q^{-jr}.\;}
\]
\end{theorem}
\begin{proof}
对固定 \(j\)，由假设展开得
\[
\rho^{-n}a_n\,\omega_q^{-jn}
=
\sum_{\ell=0}^{q-1}c_\ell\,\omega_q^{(\ell-j)n}
+
O(\theta^n).
\]
于是
\[
\frac{1}{N}\sum_{n=0}^{N-1}\rho^{-n}a_n\,\omega_q^{-jn}
=
\sum_{\ell=0}^{q-1}c_\ell
\left(
\frac{1}{N}\sum_{n=0}^{N-1}\omega_q^{(\ell-j)n}
\right)
+
\frac{1}{N}\sum_{n=0}^{N-1}O(\theta^n).
\]
当 \(\ell=j\) 时，括号中平均值恒等于 \(1\)。当 \(\ell\neq j\) 时，
\[
\sum_{n=0}^{N-1}\omega_q^{(\ell-j)n}
=
\frac{1-\omega_q^{(\ell-j)N}}{1-\omega_q^{\ell-j}},
\]
故其 Ces\`aro 平均为 \(O(N^{-1})\to 0\)。最后一项亦满足
\[
\frac{1}{N}\sum_{n=0}^{N-1}O(\theta^n)=O(N^{-1})\to 0,
\]
因为 \(\sum_{n\ge 0}\theta^n<\infty\)。因此极限存在且恰等于 \(c_j\)，得第一式。

再令 \(n=qN+r\)。由同一准周期展开，
\[
\rho^{-(qN+r)}a_{qN+r}
=
\sum_{j=0}^{q-1}c_j\,\omega_q^{j(qN+r)}
+
O(\theta^{qN+r})
=
\sum_{j=0}^{q-1}c_j\,\omega_q^{jr}
+
O(\theta^{qN+r}),
\]
因为 \(\omega_q^{jqN}=1\)。令 \(N\to\infty\)，得到
\[
A_r=\sum_{j=0}^{q-1}c_j\,\omega_q^{jr}.
\]
最后，对有限循环群 \(C_q\) 应用标准 Fourier 反演，即得
\[
c_j=\frac{1}{q}\sum_{r=0}^{q-1}A_r\,\omega_q^{-jr}.
\]
证毕。
\end{proof}

\noindent
这四条结果把 Schur 等变分解进一步压缩为一条完全显式的代数--解析主线：交替通道直接暴露纤维谱的特征多项式，Newton 刚性把非均匀谱变成严格可判别的代数条件，第一枚 hook 通道只测量平凡层相对于一步上界的差分曲率，而 torsion-\(q\) 准周期主项中的 Fourier 留数则可由原始系数序列经 Ces\`aro--Fourier 投影直接恢复。于是，Schur 断层的谱读出、偏置刚性与循环留数重建便在同一条有限可审计链上闭合。

\paragraph{固定阶 Schur 包的严格逆变换与 primitive 同余层析的精确能量账本}
在 Schur 断层成像已经按分拆压缩、而阿贝尔 primitive 同余分层的角色 Fourier 展开与商单调性也已建立之后，还可把这两条链进一步闭合为一层完全可逆且可继续压缩的层析结论：固定阶 \(q\) 的 Schur 包本身已经构成全部循环型碰撞多项式的完备坐标，其中单循环型又进一步只落在 hook 家族；与此同时，primitive 同余分层不仅在空间侧与角色侧之间形成无损 Fourier 对偶，而且还能沿子群格按 kernel 精确分层并作 Möbius 完全反演；相应地，群标签细化所新增的全部可见偏差仍可逐不可下降角色通道精确记账。

\leanverified{paper\_xi\_time\_part63b\_fixedq\_schur\_packet\_exact\_inversion}
\begin{theorem}[固定 \(q\) 的 Schur 包对循环型碰撞多项式的严格逆变换]\label{thm:xi-time-part63b-fixedq-schur-packet-exact-inversion}
固定整数 \(q\ge 2\)。对任意分拆
\[
\mu=1^{c_1(\mu)}2^{c_2(\mu)}\cdots \vdash q,
\]
定义对应的循环型碰撞多项式
\[
\mathcal C_\mu(m):=\prod_{r\ge 1}S_r(m)^{\,c_r(\mu)}.
\]
则对每个 \(\mu\vdash q\) 都有严格恒等式
\begin{equation}\label{eq:xi-time-part63b-fixedq-schur-packet-exact-inversion}
\boxed{\;
\mathcal C_\mu(m)
=
\sum_{\lambda\vdash q}\frac{\chi^\lambda(\mu)}{f^\lambda}\,T_{q,\lambda}(m)\; }.
\end{equation}
因此，对固定 \(q\) 与固定分辨率 \(m\)，Schur 包
\[
\{T_{q,\lambda}(m)\}_{\lambda\vdash q}
\]
与全部循环型碰撞多项式
\[
\{\mathcal C_\mu(m)\}_{\mu\vdash q}
\]
之间构成双向线性同构。
\end{theorem}
\begin{proof}
由命题 \ref{prop:pom-schur-tomography-by-partitions}，
\[
\frac{T_{q,\lambda}(m)}{f^\lambda}
=
\sum_{\mu\vdash q}\frac{\chi^\lambda(\mu)}{z_\mu}\,\mathcal C_\mu(m),
\]
其中
\[
z_\mu:=\prod_{r\ge 1}r^{c_r(\mu)}c_r(\mu)!,
\]
并且 \(|\mathrm{Cl}(\mu)|=q!/z_\mu\)。

现固定 \(\nu\vdash q\)。将上式两边乘以 \(\chi^\lambda(\nu)\) 并对 \(\lambda\vdash q\) 求和，得到
\[
\sum_{\lambda\vdash q}\frac{\chi^\lambda(\nu)}{f^\lambda}\,T_{q,\lambda}(m)
=
\sum_{\mu\vdash q}\frac{\mathcal C_\mu(m)}{z_\mu}
\sum_{\lambda\vdash q}\chi^\lambda(\mu)\chi^\lambda(\nu).
\]
由对称群特征标的列正交关系
\[
\sum_{\lambda\vdash q}\chi^\lambda(\mu)\chi^\lambda(\nu)
=
z_\mu\,\mathbf 1_{\{\mu=\nu\}},
\]
右端即化为 \(\mathcal C_\nu(m)\)。由于 \(\nu\) 任意，式
\eqref{eq:xi-time-part63b-fixedq-schur-packet-exact-inversion} 对一切 \(\mu\vdash q\) 成立。双向可逆性随之得到。证毕。
\end{proof}

\begin{corollary}[单个固定阶 \(q\) 的 Schur 包恢复全部 \(r\le q\) 阶碰撞矩与 primitive 计数]\label{cor:xi-time-part63b-fixedq-schur-packet-recovers-low-moments-primitives}
\leanverified{paper\_xi\_time\_part63b\_fixedq\_schur\_packet\_recovers\_low\_moments\_primitives}
对每个 \(1\le r\le q\)，记
\[
\mu_r:=(r,1^{q-r})\vdash q.
\]
则有闭式恢复公式
\begin{equation}\label{eq:xi-time-part63b-fixedq-schur-packet-recovers-sr}
\boxed{\;
S_r(m)
=
S_1(m)^{-(q-r)}
\sum_{\lambda\vdash q}\frac{\chi^\lambda(\mu_r)}{f^\lambda}\,T_{q,\lambda}(m)\; }.
\end{equation}
在 bin-fold 情形 \(S_1(m)=2^m\)，故
\begin{equation}\label{eq:xi-time-part63b-fixedq-schur-packet-recovers-sr-binfold}
\boxed{\;
S_r(m)
=
2^{-m(q-r)}
\sum_{\lambda\vdash q}\frac{\chi^\lambda(\mu_r)}{f^\lambda}\,T_{q,\lambda}(m)\; }.
\end{equation}
进一步，若原始纤维坐标定义为
\[
b_n(m):=\frac{1}{n}\sum_{d\mid n}\mu(d)\,S_{n/d}(m)
\qquad (n\ge 1),
\]
则对一切 \(n\le q\)，\(b_n(m)\) 也由同一个固定阶 \(q\) 的 Schur 包唯一恢复。
\end{corollary}
\begin{proof}
对分拆 \(\mu_r=(r,1^{q-r})\)，恰有一个长度为 \(r\) 的循环与 \(q-r\) 个不动点，故
\[
\mathcal C_{\mu_r}(m)=S_r(m)\,S_1(m)^{q-r}.
\]
代入定理 \ref{thm:xi-time-part63b-fixedq-schur-packet-exact-inversion} 即得
\eqref{eq:xi-time-part63b-fixedq-schur-packet-recovers-sr}；再用 bin-fold 语境中的
\(
S_1(m)=2^m
\)
便得到 \eqref{eq:xi-time-part63b-fixedq-schur-packet-recovers-sr-binfold}。

一旦 \(\{S_r(m)\}_{1\le r\le q}\) 已知，命题
\ref{prop:pom-schur-trace-jacobi-trudi-euler-primitive-fiber} 中的 Möbius 反演公式便逐项给出全部
\(
b_n(m)
\)
\((n\le q)\) 的唯一恢复。证毕。
\end{proof}

\begin{theorem}[单循环型的 hook--Schur 精确投影公式]\label{thm:xi-time-part63b-qcycle-hook-schur-projection}
\leanverified{paper\_xi\_time\_part63b\_qcycle\_hook\_schur\_projection}
对任意整数 \(q\ge 2\)，第 \(q\) 阶碰撞矩
\[
S_q(m):=\sum_{x\in X_m}d_m(x)^q
\]
满足完全闭式
\begin{equation}\label{eq:xi-time-part63b-qcycle-hook-schur-projection}
\boxed{\;
S_q(m)
=
\sum_{j=0}^{q-1}
\frac{(-1)^j}{\binom{q-1}{j}}\,
T_{q,(q-j,1^j)}(m)\; }.
\end{equation}
换言之，在全部长度为 \(q\) 的 Schur 通道中，单循环型 \((q)\) 只投影到 hook 家族。
\end{theorem}
\begin{proof}
在定理 \ref{thm:xi-time-part63b-fixedq-schur-packet-exact-inversion} 中取
\[
\mu=(q),
\]
得到
\[
\mathcal C_{(q)}(m)
=
\sum_{\lambda\vdash q}\frac{\chi^\lambda((q))}{f^\lambda}\,T_{q,\lambda}(m).
\]
而单个 \(q\)-循环对应的循环型碰撞多项式正是
\[
\mathcal C_{(q)}(m)=S_q(m).
\]
由 Murnaghan--Nakayama 定理，对 \(q\)-循环有
\[
\chi^\lambda((q))
=
\begin{cases}
(-1)^j,& \lambda=(q-j,1^j),\\[2pt]
0,& \text{其余 }\lambda.
\end{cases}
\]
另一方面，对 hook 分拆 \((q-j,1^j)\)，hook-length 公式给出
\[
f^{(q-j,1^j)}=\binom{q-1}{j}.
\]
将这两条显式公式代回上式，即得
\eqref{eq:xi-time-part63b-qcycle-hook-schur-projection}。证毕。
\end{proof}

\begin{corollary}[primitive 计数的纯 hook--Schur 闭式]\label{cor:xi-time-part63b-primitive-hook-schur-closed-form}
\leanverified{paper\_xi\_time\_part63b\_primitive\_hook\_schur\_closed\_form}
对任意 \(n\ge 1\)，命题 \ref{prop:pom-schur-trace-jacobi-trudi-euler-primitive-fiber} 中的 primitive 计数满足
\begin{equation}\label{eq:xi-time-part63b-primitive-hook-schur-closed-form}
\boxed{\;
b_n(m)
=
\frac{1}{n}\sum_{d\mid n}\mu(d)
\sum_{j=0}^{n/d-1}
\frac{(-1)^j}{\binom{n/d-1}{j}}\,
T_{n/d,(n/d-j,1^j)}(m)\; }.
\end{equation}
因此，全部 primitive 纤维计数都可仅由 hook 通道塔严格恢复。
\end{corollary}
\begin{proof}
由命题 \ref{prop:pom-schur-trace-jacobi-trudi-euler-primitive-fiber} 的 Möbius 反演公式，
\[
b_n(m)=\frac{1}{n}\sum_{d\mid n}\mu(d)\,S_{n/d}(m).
\]
对其中每个 \(r=n/d\)，若 \(r=1\)，则
\[
S_1(m)=T_{1,(1)}(m);
\]
若 \(r\ge 2\)，则由定理 \ref{thm:xi-time-part63b-qcycle-hook-schur-projection} 有
\[
S_r(m)
=
\sum_{j=0}^{r-1}
\frac{(-1)^j}{\binom{r-1}{j}}\,
T_{r,(r-j,1^j)}(m).
\]
逐项代回即可得到
\eqref{eq:xi-time-part63b-primitive-hook-schur-closed-form}。证毕。
\end{proof}

\begin{theorem}[primitive 同余分层的严格 Fourier 反演与均匀化判据]\label{thm:xi-time-part63b-primitive-congruence-fourier-inversion-uniformity}
\leanverified{paper\_xi\_time\_part63b\_primitive\_congruence\_fourier\_inversion\_uniformity}
沿用命题 \ref{prop:kernel-chebotarev-fourier} 与命题 \ref{prop:kernel-chebotarev-mobius-adams} 的记号。设 \(G\) 为有限阿贝尔群，\(u\in(0,1]\)。则对每个 \(n\ge 1\) 与每个角色 \(\chi\in\widehat G\)，都有严格反演公式
\begin{equation}\label{eq:xi-time-part63b-primitive-congruence-fourier-inversion}
\boxed{\;
B_n(\chi;u)
=
\sum_{c\in G}\chi(c)\,B_{n,c}(u)\; }.
\end{equation}
并且下列两条严格等价：
\begin{equation}\label{eq:xi-time-part63b-primitive-congruence-uniformity-criterion}
\boxed{\;
B_{n,c}(u)\ \text{与}\ c\ \text{无关}
\iff
B_n(\chi;u)=0\quad(\forall\,\chi\neq \mathbf 1)\; }.
\end{equation}
换言之，primitive 同余分层在 \(G\) 上完全均匀，当且仅当全部非平凡角色通道严格熄灭。
\end{theorem}
\begin{proof}
由命题 \ref{prop:kernel-chebotarev-fourier}，
\[
B_{n,c}(u)
=
\frac{1}{|G|}
\sum_{\eta\in\widehat G}\overline{\eta(c)}\,B_n(\eta;u).
\]
现对固定 \(\psi\in\widehat G\)，将上式乘以 \(\psi(c)\) 并对 \(c\in G\) 求和，得到
\[
\sum_{c\in G}\psi(c)\,B_{n,c}(u)
=
\frac{1}{|G|}
\sum_{\eta\in\widehat G}B_n(\eta;u)
\sum_{c\in G}\psi(c)\overline{\eta(c)}.
\]
由有限阿贝尔群角色正交关系，
\[
\sum_{c\in G}\psi(c)\overline{\eta(c)}
=
|G|\,\mathbf 1_{\{\psi=\eta\}},
\]
右端化为 \(B_n(\psi;u)\)，从而得到
\eqref{eq:xi-time-part63b-primitive-congruence-fourier-inversion}。

若 \(B_{n,c}(u)\) 与 \(c\) 无关，记其共同值为 \(\beta_n(u)\)。则对任意非平凡 \(\chi\neq \mathbf 1\)，有
\[
B_n(\chi;u)
=
\beta_n(u)\sum_{c\in G}\chi(c)=0.
\]
反之，若全部非平凡 \(B_n(\chi;u)\) 都为零，则由命题
\ref{prop:kernel-chebotarev-fourier}，
\[
B_{n,c}(u)
=
\frac{1}{|G|}B_n(\mathbf 1;u),
\]
与 \(c\) 无关。故 \eqref{eq:xi-time-part63b-primitive-congruence-uniformity-criterion} 成立。证毕。
\end{proof}

\begin{theorem}[primitive 角色谱的 subgroup--Möbius 完全反演]\label{thm:xi-time-part63b-primitive-kernel-layer-subgroup-mobius}
沿用命题 \ref{prop:kernel-chebotarev-fourier} 的记号。设 \(G\) 为有限阿贝尔群，固定 \(n\ge 1\) 与 \(u\in(0,1]\)。对每个子群 \(H\le G\)，定义子群求和
\[
\mathcal B_H(n,u):=\sum_{h\in H}B_{n,h}(u),
\]
以及 kernel 精确层
\[
\mathcal K_H(n,u)
:=
\sum_{\substack{\chi\in \widehat G\\ \ker\chi=H}} B_n(\chi;u).
\]
则有两条严格恒等式
\begin{equation}\label{eq:xi-time-part63b-primitive-kernel-layer-forward}
\boxed{\;
\frac{|G|}{|H|}\,\mathcal B_H(n,u)
=
\sum_{L\ge H}\mathcal K_L(n,u)\; }.
\end{equation}
\begin{equation}\label{eq:xi-time-part63b-primitive-kernel-layer-mobius}
\boxed{\;
\mathcal K_H(n,u)
=
\sum_{L\ge H}\mu_{\mathcal L(G)}(H,L)\,
\frac{|G|}{|L|}\,\mathcal B_L(n,u)\; }.
\end{equation}
其中 \(\mathcal L(G)\) 为 \(G\) 的子群格，\(\mu_{\mathcal L(G)}\) 为其 Möbius 函数。因而，primitive 同余层析不仅可恢复活动角色支撑，而且能够沿角色的 kernel 精确分层并作完全反演。
\end{theorem}
\begin{proof}
由命题 \ref{prop:kernel-chebotarev-fourier}，
\[
B_{n,h}(u)
=
\frac{1}{|G|}
\sum_{\chi\in\widehat G}\overline{\chi(h)}\,B_n(\chi;u).
\]
对 \(h\in H\) 求和，得到
\[
\mathcal B_H(n,u)
=
\frac{1}{|G|}
\sum_{\chi\in\widehat G}
B_n(\chi;u)\sum_{h\in H}\overline{\chi(h)}.
\]
有限阿贝尔群角色在子群上的正交关系给出
\[
\sum_{h\in H}\overline{\chi(h)}
=
\begin{cases}
|H|,& H\subseteq \ker\chi,\\[2pt]
0,& \text{否则}.
\end{cases}
\]
故
\[
\frac{|G|}{|H|}\,\mathcal B_H(n,u)
=
\sum_{\chi:\,H\subseteq \ker\chi}B_n(\chi;u).
\]
再按 kernel 的精确层分组，右端可写为
\[
\sum_{L\ge H}\ \sum_{\substack{\chi\in\widehat G\\ \ker\chi=L}}B_n(\chi;u)
=
\sum_{L\ge H}\mathcal K_L(n,u),
\]
从而得到
\eqref{eq:xi-time-part63b-primitive-kernel-layer-forward}。

现在在子群格 \((\mathcal L(G),\subseteq)\) 上记
\[
F(H):=\frac{|G|}{|H|}\,\mathcal B_H(n,u),
\qquad
G(H):=\mathcal K_H(n,u).
\]
则上式恰为
\[
F(H)=\sum_{L\ge H}G(L).
\]
对该偏序集应用 Möbius 反演，即得
\eqref{eq:xi-time-part63b-primitive-kernel-layer-mobius}。证毕。
\end{proof}

\begin{corollary}[循环群情形的约数格闭式：按 conductor 与精确阶分层]\label{cor:xi-time-part63b-cyclic-kernel-layer-divisor-mobius}
令 \(G=C_N\) 为阶 \(N\) 的循环群。对每个约数 \(d\mid N\)，记 \(H_d\le C_N\) 为唯一阶为 \(d\) 的子群，并定义
\[
\mathcal K_d(n,u)
:=
\sum_{\substack{\chi\in \widehat{C_N}\\ \ker\chi=H_d}} B_n(\chi;u),
\qquad
\mathcal B_d(n,u)
:=
\sum_{h\in H_d}B_{n,h}(u).
\]
则有完全显式的约数格公式
\begin{equation}\label{eq:xi-time-part63b-cyclic-kernel-layer-forward}
\boxed{\;
\frac{N}{d}\,\mathcal B_d(n,u)
=
\sum_{\substack{e\mid N\\ d\mid e}}\mathcal K_e(n,u)\; }.
\end{equation}
\begin{equation}\label{eq:xi-time-part63b-cyclic-kernel-layer-mobius}
\boxed{\;
\mathcal K_d(n,u)
=
\sum_{\substack{e\mid N\\ d\mid e}}
\mu\!\left(\frac{e}{d}\right)\frac{N}{e}\,\mathcal B_e(n,u)\; }.
\end{equation}
若把 \(\widehat{C_N}\) 参数化为
\[
\chi_a(x):=e^{2\pi i ax/N}
\qquad (0\le a<N,\ x\in \ZZ/N\ZZ),
\]
则
\[
\ker\chi_a=H_{\gcd(a,N)},
\]
因而
\begin{equation}\label{eq:xi-time-part63b-cyclic-kernel-layer-gcd}
\boxed{\;
\mathcal K_d(n,u)
=
\sum_{\substack{0\le a<N\\ \gcd(a,N)=d}} B_n(\chi_a;u)\; }.
\end{equation}
等价地，\(\mathcal K_d(n,u)\) 正是精确 conductor \(N/d\) 的角色层总和。
\end{corollary}
\begin{proof}
对循环群 \(C_N\)，子群格与约数格
\[
\{d:\ d\mid N\}
\]
在关系
\[
d\mid e \iff H_d\le H_e
\]
下同构。因此定理 \ref{thm:xi-time-part63b-primitive-kernel-layer-subgroup-mobius} 的正向分解立即化为
\[
\frac{N}{d}\,\mathcal B_d(n,u)
=
\sum_{\substack{e\mid N\\ d\mid e}}\mathcal K_e(n,u),
\]
即
\eqref{eq:xi-time-part63b-cyclic-kernel-layer-forward}。

同理，区间 \([H_d,H_e]\) 上的子群格 Möbius 函数正是整数 Möbius 函数
\[
\mu_{\mathcal L(C_N)}(H_d,H_e)=\mu(e/d),
\]
从而
\eqref{eq:xi-time-part63b-cyclic-kernel-layer-mobius} 成立。

最后，对角色
\[
\chi_a(x)=e^{2\pi i ax/N},
\]
其核大小等于 \(\gcd(a,N)\)，故
\[
\ker\chi_a=H_{\gcd(a,N)}.
\]
于是
\eqref{eq:xi-time-part63b-cyclic-kernel-layer-gcd} 随即得到。证毕。
\end{proof}

\begin{theorem}[primitive 同余细化的精确能量损失公式]\label{thm:xi-time-part63b-primitive-quotient-energy-loss-exact}
\leanverified{paper\_xi\_time\_part63b\_primitive\_quotient\_energy\_loss\_exact}
沿用推论 \ref{cor:kernel-chebotarev-plancherel-energy} 与推论 \ref{cor:kernel-chebotarev-L2-monotone} 的记号。设
\[
\pi:G_2\twoheadrightarrow G_1
\]
为有限阿贝尔群满同态，并记
\[
\pi^\ast:\widehat{G_1}\hookrightarrow \widehat{G_2},
\qquad
\chi_1\longmapsto \chi_1\circ\pi.
\]
设细层 primitive 同余分层 \(B_{n,c_2}^{(2)}(u)\) 在粗层的推前由
\[
B_{n,c_1}^{(1)}(u):=\sum_{c_2\in\pi^{-1}(c_1)}B_{n,c_2}^{(2)}(u)
\qquad(c_1\in G_1)
\]
给出。定义两层的平均值与零均值偏差
\[
\overline B_n^{(i)}(u)
:=
\frac{1}{|G_i|}\sum_{c_i\in G_i}B_{n,c_i}^{(i)}(u)
=
\frac{1}{|G_i|}B_n^{(i)}(\mathbf 1;u),
\]
\[
\varepsilon_n^{(i)}(c_i;u)
:=
B_{n,c_i}^{(i)}(u)-\overline B_n^{(i)}(u)
\qquad(i=1,2).
\]
则有严格恒等式
\begin{equation}\label{eq:xi-time-part63b-primitive-quotient-energy-loss-exact}
\boxed{\;
\sum_{c_2\in G_2}\bigl|\varepsilon_n^{(2)}(c_2;u)\bigr|^2
-
|\ker\pi|
\sum_{c_1\in G_1}\bigl|\varepsilon_n^{(1)}(c_1;u)\bigr|^2
=
\frac{1}{|G_2|}
\sum_{\chi\in \widehat G_2\setminus \pi^\ast(\widehat G_1)}
\bigl|B_n^{(2)}(\chi;u)\bigr|^2\; }.
\end{equation}
因此，群标签细化所新增的全部 primitive 偏差能量，恰等于那些不能下降到粗商的非平凡角色通道能量总和。
\end{theorem}
\begin{proof}
对两层分别应用推论 \ref{cor:kernel-chebotarev-plancherel-energy}，得到
\[
\sum_{c_i\in G_i}\bigl|\varepsilon_n^{(i)}(c_i;u)\bigr|^2
=
\frac{1}{|G_i|}
\sum_{\chi_i\in\widehat G_i\setminus\{\mathbf 1\}}
\bigl|B_n^{(i)}(\chi_i;u)\bigr|^2
\qquad(i=1,2).
\]

另一方面，由粗层分层是细层推前可知，对任意 \(\chi_1\in\widehat G_1\) 有
\[
\begin{aligned}
B_n^{(1)}(\chi_1;u)
&=
\sum_{c_1\in G_1}\chi_1(c_1)\,B_{n,c_1}^{(1)}(u)\\
&=
\sum_{c_1\in G_1}\chi_1(c_1)
\sum_{c_2\in\pi^{-1}(c_1)}B_{n,c_2}^{(2)}(u)\\
&=
\sum_{c_2\in G_2}(\chi_1\circ\pi)(c_2)\,B_{n,c_2}^{(2)}(u)\\
&=
B_n^{(2)}(\pi^\ast\chi_1;u).
\end{aligned}
\]
再用 \(|G_2|=|G_1|\,|\ker\pi|\)，便得到
\[
|\ker\pi|
\sum_{c_1\in G_1}\bigl|\varepsilon_n^{(1)}(c_1;u)\bigr|^2
=
\frac{1}{|G_2|}
\sum_{\chi\in \pi^\ast(\widehat G_1)\setminus\{\mathbf 1\}}
\bigl|B_n^{(2)}(\chi;u)\bigr|^2.
\]
将其从细层恒等式
\[
\sum_{c_2\in G_2}\bigl|\varepsilon_n^{(2)}(c_2;u)\bigr|^2
=
\frac{1}{|G_2|}
\sum_{\chi\in \widehat G_2\setminus\{\mathbf 1\}}
\bigl|B_n^{(2)}(\chi;u)\bigr|^2
\]
中相减，即得
\eqref{eq:xi-time-part63b-primitive-quotient-energy-loss-exact}。证毕。
\end{proof}

\noindent
由此，固定阶 Schur 包与阿贝尔 primitive 同余层各自都获得了更细的逆向层析：在 Schur 侧，全部循环型碰撞多项式由固定阶包完备编码，而单循环型 \(S_q(m)\) 乃至 primitive 计数 \(b_n(m)\) 又可进一步压缩到 hook 通道塔；在算术侧，空间同余类与角色通道除了构成无损 Fourier 对偶外，还可沿子群格按 kernel 精确分层并作 Möbius 完全反演，在循环群上更退化为约数格上的整数 Möbius 闭式；最后，商细化所增加的全部可见偏差仍可精确分解为不能下降到粗商的新增角色能量。于是，Schur 断层的 hook 压缩、Chebotarev primitive 的 kernel 层析与商细化能量账本便在同一层结论中完成闭合。

\paragraph{Schur 通道化的 Frobenius 逆变换、Cauchy 主核与完全充分统计闭合}
在固定阶 Schur 包的严格逆变换、交替 Schur 通道的特征多项式识别以及高阶单包完备性都已建立之后，8.1959--8.1965 的 Schur 断层理论还可进一步压缩为如下六条闭合命题：循环型碰撞矩与 Schur 通道迹之间存在严格可逆的 Frobenius 坐标变换；全部 Schur 断层由单一 Cauchy 主核一次性生成；平凡与反对称通道互为符号倒数；反对称包自动携带全负实根与 \(\mathrm{PF}_\infty\) 刚性；而平凡包与反对称包各自单独都足以无损恢复全部 Schur 断层。

\begin{theorem}[Schur 通道迹与循环型碰撞矩之间的严格可逆 Frobenius 变换]\label{thm:xi-time-part63c-frobenius-exact-inversion-cycle-moments}
\leanverified{paper\_xi\_time\_part63c\_frobenius\_exact\_inversion\_cycle\_moments}
固定整数 \(q\ge 2\)。对任意分拆
\[
\mu=1^{c_1(\mu)}2^{c_2(\mu)}\cdots\vdash q,
\]
定义循环型碰撞矩
\[
C_\mu(m):=\prod_{r\ge 1}S_r(m)^{\,c_r(\mu)}.
\]
则对每个 \(\mu\vdash q\) 都有精确反演公式
\[
\boxed{\;
C_\mu(m)
=
\sum_{\lambda\vdash q}\chi^\lambda(\mu)\,\frac{T_{q,\lambda}(m)}{f^\lambda}.\;}
\]
特别地，对单圈型 \(\mu=(q)\) 有
\[
\boxed{\;
S_q(m)
=
\sum_{\lambda\vdash q}\chi^\lambda((q))\,\frac{T_{q,\lambda}(m)}{f^\lambda}.\;}
\]
因此，固定 \(q\) 的全部 Schur 通道迹族
\[
\left\{\frac{T_{q,\lambda}(m)}{f^\lambda}\right\}_{\lambda\vdash q}
\]
与同阶全部循环型碰撞矩族
\[
\{C_\mu(m)\}_{\mu\vdash q}
\]
彼此线性等价；二者只是同一共轭类函数在 Schur 基与幂和基中的两组坐标。
\end{theorem}
\begin{proof}
这正是定理 \ref{thm:xi-time-part63b-fixedq-schur-packet-exact-inversion} 的改写，其中该定理的记号
\[
\mathcal C_\mu(m):=\prod_{r\ge 1}S_r(m)^{c_r(\mu)}
\]
与此处的 \(C_\mu(m)\) 完全一致。对 \(\mu=(q)\)，由
\[
C_{(q)}(m)=S_q(m)
\]
立得第二式。证毕。
\end{proof}

\begin{theorem}[全部 Schur 断层由单一 Cauchy 主核统一生成]\label{thm:xi-time-part63c-schur-cauchy-master-kernel}
设 \(\mathbf y=(y_1,y_2,\dots)\) 为任意有限支撑变量族，并定义 Schur 总生成函数
\[
\mathcal K_m(\mathbf y)
:=
\sum_{\lambda}
\frac{T_{|\lambda|,\lambda}(m)}{f^\lambda}\,
s_\lambda(\mathbf y),
\]
其中求和遍历全部分拆并把空分拆的 \(1\) 项计入。则有严格恒等式
\[
\boxed{\;
\mathcal K_m(\mathbf y)
=
\prod_{x\in X_m}\prod_{j\ge 1}\frac{1}{1-d_m(x)y_j}.\;}
\]
等价地，
\[
\boxed{\;
\mathcal K_m(\mathbf y)
=
\prod_{j\ge 1}H_m(y_j),\qquad
H_m(z):=1+\sum_{q\ge 1}T_{q,(q)}(m)\,z^q.\;}
\]
因此，整个 Schur 通道塔并非若干彼此孤立的固定 \(q\) 层，而是由纤维退化谱
\[
\{d_m(x)\}_{x\in X_m}
\]
唯一决定的单一 Cauchy 主核。
\leanverified{paper\_xi\_time\_part63c\_schur\_cauchy\_master\_kernel}
\end{theorem}
\begin{proof}
由推论 \ref{cor:pom-schur-trace-frobenius-character-identification}，
\[
\frac{T_{|\lambda|,\lambda}(m)}{f^\lambda}=s_\lambda(D_m),
\]
其中 \(D_m=\operatorname{diag}(d_m(x))_{x\in X_m}\)。故
\[
\mathcal K_m(\mathbf y)
=
\sum_{\lambda}s_\lambda(D_m)\,s_\lambda(\mathbf y).
\]
对 \(D_m\) 的对角元多重集与变量族 \(\mathbf y\) 应用经典 Cauchy 恒等式，即得
\[
\mathcal K_m(\mathbf y)
=
\prod_{x\in X_m}\prod_{j\ge 1}(1-d_m(x)y_j)^{-1}.
\]
另一方面，由命题 \ref{prop:pom-schur-trace-jacobi-trudi-euler-primitive-fiber}，
\[
H_m(z)=\prod_{x\in X_m}(1-d_m(x)z)^{-1}.
\]
对每个 \(j\) 取 \(z=y_j\) 并相乘，即得第二个公式。证毕。
\end{proof}

\leanverified{paper\_xi\_time\_part63c\_trivial\_alternating\_sign\_reciprocal\_duality}
\begin{theorem}[平凡通道与反对称通道之间的严格符号倒数对偶]\label{thm:xi-time-part63c-trivial-alternating-sign-reciprocal-duality}
记 \(n:=|X_m|\)，并按空分拆约定
\[
T_{0,\varnothing}(m):=1.
\]
定义平凡通道与反对称通道生成函数
\[
H_m(z):=1+\sum_{q\ge 1}T_{q,(q)}(m)\,z^q
=
\sum_{q\ge 0}h_q(D_m)\,z^q,
\]
\[
E_m(z):=T_{0,\varnothing}(m)+\sum_{q=1}^{n}T_{q,(1^q)}(m)\,z^q
=
\sum_{q=0}^{n}e_q(D_m)\,z^q.
\]
则有精确恒等式
\[
\boxed{\;
E_m(z)=H_m(-z)^{-1}.\;}
\]
等价地，
\[
\boxed{\;
E_m(z)=\prod_{x\in X_m}(1+d_m(x)z),\qquad
Q_m(z):=\prod_{x\in X_m}(1-d_m(x)z)=E_m(-z).\;}
\]
因此，反对称通道包
\[
\{T_{0,\varnothing}(m)\}\cup\{T_{q,(1^q)}(m)\}_{1\le q\le n}
\]
本身就是纤维退化谱多重集 \(\{d_m(x)\}_{x\in X_m}\) 的特征多项式系数列；其零点恰为
\[
-d_m(x)^{-1}
\]
（按重数计）。换言之，仅凭反对称通道一包即可完整恢复全部纤维多重度谱。
\end{theorem}
\begin{proof}
这正是定理 \ref{thm:xi-time-part63a-alternating-schur-characteristic-polynomial} 与命题
\ref{prop:pom-schur-trace-jacobi-trudi-euler-primitive-fiber}
中
\[
H_m(z)=\prod_{x\in X_m}(1-d_m(x)z)^{-1}
\]
的直接合并。由
\[
E_m(z)=\prod_{x\in X_m}(1+d_m(x)z)
\]
与
\[
H_m(-z)=\prod_{x\in X_m}(1+d_m(x)z)^{-1}
\]
相乘即得
\[
E_m(z)H_m(-z)=1.
\]
再由
\[
Q_m(z)=\prod_{x\in X_m}(1-d_m(x)z)=E_m(-z)
\]
便知反对称通道包的系数列恰是 \(Q_m\) 的系数列，而 \(Q_m\) 的根多重集正为 \(\{d_m(x)^{-1}\}\)，故可唯一恢复全部纤维多重度谱。证毕。
\end{proof}

\begin{corollary}[反对称 Schur 包是 \texorpdfstring{$\mathrm{PF}_\infty$}{PFinfty} 序列并满足严格 Newton 链]\label{cor:xi-time-part63c-alternating-pf-infty-newton-rigidity}
\leanverified{paper\_xi\_time\_part63c\_alternating\_pf\_infty\_newton\_rigidity}
沿用定理 \ref{thm:xi-time-part63c-trivial-alternating-sign-reciprocal-duality} 的记号，记
\[
E_q(m):=
\begin{cases}
T_{0,\varnothing}(m),& q=0,\\[2pt]
T_{q,(1^q)}(m),& 1\le q\le n.
\end{cases}
\]
则生成多项式
\[
E_m(z)=\sum_{q=0}^{n}E_q(m)z^q
\]
的全部零点都位于负实轴，因此系数列
\[
\bigl(E_q(m)\bigr)_{q=0}^{n}
\]
是 \(\mathrm{PF}_\infty\) 序列。特别地，对每个 \(1\le q\le n-1\) 都有 Newton 不等式
\[
\boxed{\;
\frac{E_q(m)^2}{\binom{n}{q}^2}
\ge
\frac{E_{q-1}(m)}{\binom{n}{q-1}}
\frac{E_{q+1}(m)}{\binom{n}{q+1}}.\;}
\]
若纤维退化谱不全相等，则上述不等式对一切 \(1\le q\le n-1\) 严格成立。
\end{corollary}
\begin{proof}
由定理 \ref{thm:xi-time-part63c-trivial-alternating-sign-reciprocal-duality}，
\[
E_m(z)=\prod_{x\in X_m}(1+d_m(x)z),
\]
故其全部零点均为负实数 \(-d_m(x)^{-1}\)。因此带非负系数的多项式 \(E_m(z)\) 具有全负实根，系数列遂为 \(\mathrm{PF}_\infty\) 序列。Newton 不等式是这一事实的标准推论；而严格性则由定理
\ref{thm:xi-time-part63a-alternating-schur-newton-rigidity}
给出的等号判据立即得到：只有当全部 \(d_m(x)\) 完全相等时，归一化 Newton 链才会失去严格性。证毕。
\end{proof}

\begin{theorem}[平凡通道包、反对称通道包与全部 Schur 断层之间的完全充分统计关系]\label{thm:xi-time-part63c-trivial-alternating-schur-complete-sufficiency}
记 \(n:=|X_m|\)。定义三组数据
\[
\mathfrak T_m^{\mathrm{triv}}
:=
\{T_{0,\varnothing}(m)\}\cup\{T_{q,(q)}(m)\}_{q\ge 1},
\qquad
\mathfrak T_m^{\mathrm{alt}}
:=
\{T_{0,\varnothing}(m)\}\cup\{T_{q,(1^q)}(m)\}_{1\le q\le n},
\]
\[
\mathfrak T_m^{\mathrm{Schur}}
:=
\{T_{q,\lambda}(m)\}_{q\ge 1,\ \lambda\vdash q}.
\]
则这三组数据彼此无损等价。更精确地：
\begin{enumerate}
  \item 由 \(\mathfrak T_m^{\mathrm{triv}}\) 的前 \(n+1\) 项，可唯一恢复
  \[
  Q_m(z)=\prod_{x\in X_m}(1-d_m(x)z),
  \]
  故唯一恢复全部 \(\{d_m(x)\}_{x\in X_m}\)；
  \item 由 \(\mathfrak T_m^{\mathrm{alt}}\) 可直接写出
  \[
  \boxed{\;
  Q_m(z)=T_{0,\varnothing}(m)+\sum_{q=1}^{n}(-1)^q\,T_{q,(1^q)}(m)\,z^q,\;}
  \]
  因而同样唯一恢复全部 \(\{d_m(x)\}_{x\in X_m}\)；
  \item 一旦 \(\{d_m(x)\}_{x\in X_m}\) 已知，则全部 Schur 断层由
  \[
  T_{q,\lambda}(m)=f^\lambda s_\lambda(D_m)
  \]
  唯一给出。
\end{enumerate}
因此，平凡通道包与反对称通道包各自单独都可视为全部 Schur 断层的完全充分统计量。
\end{theorem}
\begin{proof}
第 \(1\) 条正是命题 \ref{prop:pom-trivial-channel-finite-decidability} 的内容：平凡通道生成函数
\[
H_m(z)=1+\sum_{q\ge 1}T_{q,(q)}(m)z^q
\]
满足
\[
Q_m(z)H_m(z)=1,
\]
故由前 \(n+1\) 项即可递归恢复 \(Q_m\) 的系数，从而恢复全部纤维多重度谱。

第 \(2\) 条由定理 \ref{thm:xi-time-part63c-trivial-alternating-sign-reciprocal-duality} 立即得到，因为
\[
Q_m(z)=E_m(-z)
=
T_{0,\varnothing}(m)+\sum_{q=1}^{n}(-1)^qT_{q,(1^q)}(m)z^q.
\]

对第 \(3\) 条，一旦 \(\{d_m(x)\}_{x\in X_m}\) 已知，便知 \(D_m\) 的对角元多重集。于是由推论
\ref{cor:pom-schur-trace-frobenius-character-identification}，
\[
T_{q,\lambda}(m)=f^\lambda s_\lambda(D_m)
\]
对一切 \(q\ge 1\) 与 \(\lambda\vdash q\) 都被唯一确定。

反向恢复则是显然的，因为 \(\mathfrak T_m^{\mathrm{Schur}}\) 本身包含所有平凡分拆 \((q)\) 与符号分拆 \((1^q)\) 的通道值。故三组数据彼此无损等价。证毕。
\end{proof}

\begin{corollary}[window-\texorpdfstring{$6$}{6} 的全部 Schur 断层由一个 \texorpdfstring{$21$}{21} 次特征多项式完全锁定]\label{cor:xi-time-part63c-window6-schur-packet-locked-by-charpoly}
\leanverified{paper\_xi\_time\_part63c\_window6\_schur\_packet\_locked\_by\_charpoly}
在 \(m=6\) 时，由直方图
\[
N_{6,2}=8,\qquad N_{6,3}=4,\qquad N_{6,4}=9,
\qquad |X_6|=21,
\]
有
\[
\boxed{\;
Q_6(z)=\prod_{x\in X_6}(1-d_6(x)z)
=(1-2z)^8(1-3z)^4(1-4z)^9.\;}
\]
相应地，
\[
\boxed{\;
E_6(z)=Q_6(-z)=(1+2z)^8(1+3z)^4(1+4z)^9,\qquad
H_6(z)=Q_6(z)^{-1}.\;}
\]
因此，window-\(6\) 的全部 Schur 断层、全部碰撞矩 \(\{S_r(6)\}_{r\ge 1}\) 以及全部原始纤维坐标 \(\{b_n(6)\}_{n\ge 1}\)，都由这一条 \(21\) 次特征多项式无损决定。
\end{corollary}
\begin{proof}
由推论 \ref{cor:fold-bin-m6-fiber-hist}，window-\(6\) 的纤维多重度只取 \(2,3,4\) 三个值，并且对应重数分别为 \(8,4,9\)。将其代入定理
\ref{thm:xi-time-part63c-trivial-alternating-sign-reciprocal-duality}
中的一般公式
\[
Q_m(z)=\prod_{x\in X_m}(1-d_m(x)z)
\]
便得所示 \(Q_6(z)\)、\(E_6(z)=Q_6(-z)\) 与 \(H_6(z)=Q_6(z)^{-1}\)。
最后，由定理
\ref{thm:xi-time-part63c-trivial-alternating-schur-complete-sufficiency}
与命题 \ref{prop:pom-schur-trace-jacobi-trudi-euler-primitive-fiber}，一旦 \(Q_6\) 已知，全部纤维多重度谱、全部 Schur 断层、全部碰撞矩及全部原始纤维坐标便同时唯一确定。证毕。
\end{proof}

\paragraph{平凡通道窗口、Schur--Plancherel 极限与 generic pressure fan 比值律}
在固定阶 Schur 包的严格逆变换、primitive 坐标的 Euler 分解以及分拆压力相图都已闭合之后，还可把最新一段 Schur 断层后段材料进一步压缩为如下六条结果：前两条刻画平凡通道窗口与 primitive 坐标之间的单位上三角整自同构；其后两条给出恒等扇区主导下的 Schur--Plancherel 极限及非平凡能量比例；最后两条把 generic pressure chamber 中的通道比较与 Perron-purity 下的 hook 湮灭压缩为有限 partition fan 上的显式角色比值与线性压力差。

\begin{theorem}[平凡通道与 primitive 坐标之间的单位上三角整自同构]\label{thm:xi-time-part63c-trivial-channel-unitriangular-primitive-automorphism}
\leanverified{paper\_xi\_time\_part63c\_trivial\_channel\_unitriangular\_primitive\_automorphism}
固定分辨率 \(m\)，并记
\[
T_q(m):=T_{q,(q)}(m)
\qquad(q\ge 1),
\]
其中 \(T_{q,(q)}(m)=h_q(D_m)\) 为平凡 Schur 通道。沿用命题
\ref{prop:pom-schur-trace-jacobi-trudi-euler-primitive-fiber}
中的 primitive 坐标
\[
b_n(m):=\frac{1}{n}\sum_{d\mid n}\mu(d)\,S_{n/d}(m),
\]
以及 Euler 分解
\[
H_m(z)=\prod_{n\ge 1}(1-z^n)^{-b_n(m)}=\sum_{q\ge 0}T_q(m)z^q.
\]
则对每个 \(q\ge 1\)，存在整系数多项式
\[
P_q\in \ZZ[b_1,\dots,b_{q-1}]
\]
使得
\[
\boxed{\;
T_q(m)=b_q(m)+P_q\bigl(b_1(m),\dots,b_{q-1}(m)\bigr).\;}
\]
因此，对任意 \(Q\ge 1\)，坐标变换
\[
(b_1,\dots,b_Q)\longmapsto (T_1,\dots,T_Q)
\]
是 \(\ZZ^Q\) 上的单位上三角多项式自同构。其逆变换可递归写为
\[
\boxed{\;
b_q(m)=T_q(m)-P_q\bigl(b_1(m),\dots,b_{q-1}(m)\bigr)
\qquad(1\le q\le Q).\;}
\]
\end{theorem}
\begin{proof}
由命题 \ref{prop:pom-schur-trace-trivial-channel-primitive-fiber-expansion}，
\[
T_q(m)=
\sum_{\substack{(k_n)_{n\ge 1}\subset\NN\\ \sum_{n\ge 1}n k_n=q}}
\prod_{n\ge 1}\binom{b_n(m)+k_n-1}{k_n}.
\]
在全部满足
\[
\sum_{n\ge 1}n k_n=q
\]
的多重指标中，唯一含有 \(b_q(m)\) 的可行项只能是
\[
k_q=1,
\qquad
k_n=0\quad(n\neq q),
\]
其贡献恰为
\[
\binom{b_q(m)}{1}=b_q(m).
\]
确实，若 \(k_q\ge 2\)，则 \(qk_q\ge 2q>q\) 不可行；若 \(k_q=1\) 且某个 \(k_n>0\)（\(n<q\)），则
\[
\sum_{r\ge 1}r k_r\ge q+n>q,
\]
同样不可能发生。故其余全部可行项都只涉及 \(\{b_1(m),\dots,b_{q-1}(m)\}\)。把这些项之和记为
\[
P_q\bigl(b_1(m),\dots,b_{q-1}(m)\bigr),
\]
便得到所示上三角展开。

由于对每个 \(q\) 都有
\[
T_q=b_q+\text{(仅依赖于 }b_1,\dots,b_{q-1}\text{ 的整系数多项式)},
\]
故该坐标变换在每一级上都具有对角元为 \(1\) 的上三角形式，可逐级回代唯一反解。于是它给出 \(\ZZ^Q\) 上的单位上三角多项式自同构。证毕。
\end{proof}

\begin{corollary}[有限平凡通道窗口对全部低阶 Schur 断层的完备充分性]\label{cor:xi-time-part63c-finite-trivial-window-complete-schur-statistic}
固定 \(Q\ge 1\)。则对任意分拆 \(\lambda\vdash q\le Q\)，存在整系数多项式
\[
\mathcal P_{q,\lambda}^{(Q)}\in \ZZ[T_1,\dots,T_Q]
\]
使得
\[
\boxed{\;
\frac{T_{q,\lambda}(m)}{f^\lambda}
=
\mathcal P_{q,\lambda}^{(Q)}\bigl(T_1(m),\dots,T_Q(m)\bigr).\;}
\]
反之，对每个 \(q\le Q\)，存在整系数多项式
\[
\mathcal R_q^{(Q)}\in \ZZ[T_1,\dots,T_q]
\]
使得
\[
\boxed{\;
b_q(m)=\mathcal R_q^{(Q)}\bigl(T_1(m),\dots,T_q(m)\bigr).\;}
\]
因此，对固定 \(Q\)，下列两组数据逐 \(m\) 严格等价：
\[
\{T_1(m),\dots,T_Q(m)\}
\quad\Longleftrightarrow\quad
\left\{\frac{T_{q,\lambda}(m)}{f^\lambda}:1\le q\le Q,\ \lambda\vdash q\right\}.
\]
\end{corollary}
\begin{proof}
由推论 \ref{cor:pom-schur-trace-all-channels-primitive-fiber-coordinate}，
\[
\frac{T_{q,\lambda}(m)}{f^\lambda}\in \ZZ[b_1(m),\dots,b_q(m)]
\qquad(\lambda\vdash q).
\]
再由定理 \ref{thm:xi-time-part63c-trivial-channel-unitriangular-primitive-automorphism}，每个 \(b_j(m)\) 都可递归写成 \(\ZZ[T_1,\dots,T_j]\) 中的多项式。将这些表达代入上式，即得
\[
\frac{T_{q,\lambda}(m)}{f^\lambda}
=
\mathcal P_{q,\lambda}^{(Q)}\bigl(T_1(m),\dots,T_Q(m)\bigr)
\]
的整系数多项式表达。

反向公式正是定理 \ref{thm:xi-time-part63c-trivial-channel-unitriangular-primitive-automorphism} 的回代构造；把其中右端的 \(b_1,\dots,b_{q-1}\) 继续用已恢复的 \(T_1,\dots,T_{q-1}\) 代换，即得 \(\mathcal R_q^{(Q)}\)。

于是，固定窗口 \(\{T_1,\dots,T_Q\}\) 与全部低阶 Schur 断层
\[
\left\{\frac{T_{q,\lambda}}{f^\lambda}:q\le Q\right\}
\]
互为整多项式重编码，故两组数据逐 \(m\) 严格等价。证毕。
\end{proof}

\begin{theorem}[恒等扇区严格主导时的 Schur--Plancherel 极限]\label{thm:xi-time-part63c-identity-sector-schur-plancherel-limit}
固定 \(q\ge 2\)，并假设恒等扇区严格主导：
\[
\frac{P(1)}{1}>\max_{2\le r\le q}\frac{P(r)}{r}.
\]
则存在 \(\varepsilon>0\)，使对任意 \(\lambda\vdash q\) 都有
\[
T_{q,\lambda}(m)
=
\frac{(f^\lambda)^2}{q!}\,\kappa_1^{\,q}\,e^{mqP(1)}
+O\!\left(e^{m(qP(1)-\varepsilon)}\right).
\]
因此，若定义归一化通道权
\[
\Pi_{m,q}(\lambda):=
\frac{T_{q,\lambda}(m)}{\sum_{\mu\vdash q}T_{q,\mu}(m)},
\]
则有极限
\[
\boxed{\;
\Pi_{m,q}(\lambda)\longrightarrow \frac{(f^\lambda)^2}{q!}
\qquad(m\to\infty).\;}
\]
\end{theorem}
\begin{proof}
由推论 \ref{cor:pom-identity-sector-collapse-all-schur-channels}，对每个 \(\lambda\vdash q\)，都有
\[
T_{q,\lambda}(m)
=
\frac{(f^\lambda)^2}{q!}\,\kappa_1^{\,q}\,e^{mqP(1)}
+O\!\left(e^{m(qP(1)-\varepsilon)}\right)
\]
对某个统一的 \(\varepsilon>0\) 成立。

对所有 \(\mu\vdash q\) 求和，并利用经典恒等式
\[
\sum_{\mu\vdash q}(f^\mu)^2=q!,
\]
得到
\[
\sum_{\mu\vdash q}T_{q,\mu}(m)
=
\kappa_1^{\,q}e^{mqP(1)}
+O\!\left(e^{m(qP(1)-\varepsilon)}\right).
\]
故
\[
\Pi_{m,q}(\lambda)
=
\frac{\frac{(f^\lambda)^2}{q!}\,\kappa_1^{\,q}e^{mqP(1)}+O(e^{m(qP(1)-\varepsilon)})}
{\kappa_1^{\,q}e^{mqP(1)}+O(e^{m(qP(1)-\varepsilon)})}.
\]
分子分母同除以 \(\kappa_1^{\,q}e^{mqP(1)}\) 并令 \(m\to\infty\)，即得
\[
\Pi_{m,q}(\lambda)\to \frac{(f^\lambda)^2}{q!}.
\]
证毕。
\end{proof}

\begin{corollary}[Schur--Plancherel 能量与非平凡方差的普适极限比例]\label{cor:xi-time-part63c-schur-plancherel-energy-variance-ratio}
沿用定理 \ref{thm:xi-time-part63c-identity-sector-schur-plancherel-limit} 的假设。定义
\[
E_q(m):=
\frac{1}{q!}\sum_{\sigma\in S_q}\mathcal C_\sigma(m)^2
=
\sum_{\lambda\vdash q}\left(\frac{T_{q,\lambda}(m)}{f^\lambda}\right)^2,
\]
\[
V_q(m):=
\frac{1}{q!}\sum_{\sigma\in S_q}\bigl(\mathcal C_\sigma(m)-\overline{\mathcal C}(m)\bigr)^2
=
\sum_{\lambda\vdash q,\ \lambda\neq(q)}\left(\frac{T_{q,\lambda}(m)}{f^\lambda}\right)^2,
\]
其中
\[
\overline{\mathcal C}(m):=\frac{1}{q!}\sum_{\sigma\in S_q}\mathcal C_\sigma(m)=T_{q,(q)}(m).
\]
则有精确渐近
\[
\boxed{\;
E_q(m)\sim \frac{\kappa_1^{2q}}{q!}\,e^{2mqP(1)},
\qquad
V_q(m)\sim \frac{q!-1}{q!^2}\,\kappa_1^{2q}e^{2mqP(1)}.\;}
\]
特别地，
\[
\boxed{\;
\frac{V_q(m)}{E_q(m)}\longrightarrow 1-\frac{1}{q!}.\;}
\]
\end{corollary}
\begin{proof}
由定理 \ref{thm:pom-schur-plancherel-energy-identity} 与定理 \ref{thm:pom-schur-variance-decomposition}，\(E_q(m)\) 与 \(V_q(m)\) 分别由全部 Schur 通道与全部非平凡 Schur 通道的平方和给出。

再由定理 \ref{thm:xi-time-part63c-identity-sector-schur-plancherel-limit}，
\[
\frac{T_{q,\lambda}(m)}{f^\lambda}
=
\frac{f^\lambda}{q!}\,\kappa_1^{\,q}e^{mqP(1)}
+O\!\left(e^{m(qP(1)-\varepsilon)}\right).
\]
平方后得到
\[
\left(\frac{T_{q,\lambda}(m)}{f^\lambda}\right)^2
=
\frac{(f^\lambda)^2}{q!^2}\,\kappa_1^{2q}e^{2mqP(1)}
+O\!\left(e^{m(2qP(1)-\varepsilon)}\right).
\]
对 \(\lambda\vdash q\) 求和，并利用
\[
\sum_{\lambda\vdash q}(f^\lambda)^2=q!,
\]
得
\[
E_q(m)\sim \frac{\kappa_1^{2q}}{q!}\,e^{2mqP(1)}.
\]

对 \(V_q(m)\) 只需在求和中去掉平凡分拆 \((q)\)。由于
\[
f^{(q)}=1,
\qquad
\sum_{\lambda\vdash q,\ \lambda\neq(q)}(f^\lambda)^2=q!-1,
\]
故
\[
V_q(m)\sim \frac{q!-1}{q!^2}\,\kappa_1^{2q}e^{2mqP(1)}.
\]
最后相除即得
\[
\frac{V_q(m)}{E_q(m)}\to 1-\frac{1}{q!}.
\]
证毕。
\end{proof}

\begin{theorem}[generic pressure chamber 内任意两条 Schur 通道的严格指数比值律]\label{thm:xi-time-part63c-generic-pressure-fan-ratio-law}
固定 \(q\ge 2\)，并把压力向量写为
\[
p=(P(1),\dots,P(q))\in\RR^q.
\]
假设
\[
p\notin \bigcup\mathcal H_q,
\]
其中 \(\mathcal H_q\) 为定理 \ref{thm:pom-schur-pressure-fan-diagram} 的共振超平面族。则对每个 \(\lambda\vdash q\)，存在唯一主导分拆
\[
\nu_\lambda(p)\in \mathrm{Supp}(\lambda)
\]
与某个 \(\varepsilon_\lambda(p)>0\)，使
\[
T_{q,\lambda}(m)
=
f^\lambda\frac{\chi^\lambda(\nu_\lambda(p))}{z_{\nu_\lambda(p)}}K(\nu_\lambda(p))\,e^{m\mathcal P(\nu_\lambda(p))}
+O\!\left(e^{m(\mathcal P(\nu_\lambda(p))-\varepsilon_\lambda(p))}\right).
\]
由此得到：
\begin{enumerate}
  \item 若两条通道 \(\lambda,\mu\vdash q\) 共享同一主导分拆
  \[
  \nu_\lambda(p)=\nu_\mu(p)=:\nu,
  \]
  则
  \[
  \boxed{\;
  \frac{T_{q,\lambda}(m)}{T_{q,\mu}(m)}
  \longrightarrow
  \frac{f^\lambda\chi^\lambda(\nu)}{f^\mu\chi^\mu(\nu)}.\;}
  \]
  \item 若
  \[
  \mathcal P(\nu_\lambda(p))>\mathcal P(\nu_\mu(p)),
  \]
  则
  \[
  \boxed{\;
  \frac{1}{m}\log\left|\frac{T_{q,\lambda}(m)}{T_{q,\mu}(m)}\right|
  \longrightarrow
  \mathcal P(\nu_\lambda(p))-\mathcal P(\nu_\mu(p)).\;}
  \]
\end{enumerate}
\end{theorem}
\begin{proof}
第一段渐近正是定理 \ref{thm:pom-schur-pressure-fan-diagram} 的内容。

若 \(\nu_\lambda(p)=\nu_\mu(p)=\nu\)，则两式主指数相同，且共有因子
\[
\frac{K(\nu)}{z_\nu}e^{m\mathcal P(\nu)}
\]
在比值中完全约去，从而
\[
\frac{T_{q,\lambda}(m)}{T_{q,\mu}(m)}
=
\frac{f^\lambda\chi^\lambda(\nu)+O(e^{-m\varepsilon})}
{f^\mu\chi^\mu(\nu)+O(e^{-m\varepsilon})}
\]
对某个 \(\varepsilon>0\) 成立。由于 \(\nu\in\mathrm{Supp}(\lambda)\cap\mathrm{Supp}(\mu)\)，主系数均非零，令 \(m\to\infty\) 即得第一条。

若
\[
\mathcal P(\nu_\lambda(p))>\mathcal P(\nu_\mu(p)),
\]
则
\[
T_{q,\lambda}(m)=C_\lambda e^{m\mathcal P(\nu_\lambda(p))}\bigl(1+O(e^{-m\varepsilon_\lambda})\bigr),
\]
\[
T_{q,\mu}(m)=C_\mu e^{m\mathcal P(\nu_\mu(p))}\bigl(1+O(e^{-m\varepsilon_\mu})\bigr),
\]
其中 \(C_\lambda,C_\mu\neq 0\)。因此
\[
\log\left|\frac{T_{q,\lambda}(m)}{T_{q,\mu}(m)}\right|
=
m\bigl(\mathcal P(\nu_\lambda(p))-\mathcal P(\nu_\mu(p))\bigr)+O(1),
\]
两边除以 \(m\) 并令 \(m\to\infty\) 即得第二条。证毕。
\end{proof}

\begin{theorem}[generic Perron-purity 强制唯一最大分拆为 \texorpdfstring{$(q)$}{(q)} 并导致全部非平凡 hook 通道湮灭]\label{thm:xi-time-part63c-generic-perron-purity-kills-all-hooks}
固定 \(q\ge 2\)，并作如下假设：
\begin{enumerate}
  \item \(A_q\) 原始；
  \item 其 Perron 根
  \[
  \rho_q=e^{P(q)}
  \]
  简单，且仅出现在平凡通道 \(\lambda=(q)\)，即对所有 \(\lambda\neq(q)\) 都有
  \[
  \rho(B_{q,\lambda})<\rho_q;
  \]
  \item 压力向量 \(p=(P(1),\dots,P(q))\) 满足
  \[
  p\notin \bigcup\mathcal H_q.
  \]
\end{enumerate}
则成立：
\begin{enumerate}
  \item 唯一全局最大化的分拆必为
  \[
  (q),
  \]
  等价地
  \[
  \boxed{\;
  P(q)>\mathcal P(\nu)\qquad(\forall\,\nu\vdash q,\ \nu\neq(q)).\;}
  \]
  \item 分拆共振间隙严格为正：
  \[
  \boxed{\;
  \Delta_{\mathrm{part}}(q)>0.\;}
  \]
  \item 对所有非平凡 hook 分拆
  \[
  \lambda=(q-k,1^k),
  \qquad
  1\le k\le q-1,
  \]
  都有
  \[
  \boxed{\;
  B_{q,\lambda}=0,
  \qquad
  T_{q,\lambda}(m)=0\ \ \forall m\ge 1.\;}
  \]
\end{enumerate}
\end{theorem}
\begin{proof}
设 \(\nu_*\) 为全局最大化分拆。若 \(\nu_*\neq(q)\)，则由 generic 假设
\[
p\notin\bigcup\mathcal H_q
\]
与定理 \ref{thm:pom-schur-pressure-fan-diagram} 可知，全体分拆压力 \(\{\mathcal P(\nu)\}_{\nu\vdash q}\) 两两不同，故
\[
\mathcal R_{\max}=\{\nu_*\}
\]
为单点。

现考虑共振层权函数
\[
w(\nu):=K(\nu)\mathbf 1_{\{\nu=\nu_*\}}.
\]
由于 \(\nu_*\neq(q)\)，该函数不是常值类函数。于是由命题 \ref{prop:pom-schur-resonant-layer-total-cancellation}，必存在某个非平凡分拆 \(\lambda\neq(q)\) 使对应 Fourier 系数
\[
b_\lambda=
\frac{\chi^\lambda(\nu_*)}{z_{\nu_*}}K(\nu_*)
\]
非零。特别地，\(\nu_*\in\mathrm{Supp}(\lambda)\)。

再由定理 \ref{thm:pom-schur-pressure-fan-diagram}，该 \(\lambda\) 通道满足
\[
T_{q,\lambda}(m)
=
f^\lambda\frac{\chi^\lambda(\nu_*)}{z_{\nu_*}}K(\nu_*)\,e^{m\mathcal P(\nu_*)}
+O\!\left(e^{m(\mathcal P(\nu_*)-\varepsilon_\lambda)}\right),
\]
故其主指数率为 \(\mathcal P(\nu_*)\)。由定理
\ref{thm:pom-schur-channel-radius-character-sieve}，
\[
\rho(B_{q,\lambda})=e^{\mathcal P(\nu_*)}.
\]
而 \(\nu_*\neq(q)\) 且全局最大，结合 generic 假设可知
\[
\mathcal P(\nu_*)>\mathcal P((q))=P(q),
\]
从而
\[
\rho(B_{q,\lambda})=e^{\mathcal P(\nu_*)}>e^{P(q)}=\rho_q.
\]
这与 Perron-purity 假设
\[
\rho(B_{q,\lambda})<\rho_q\qquad(\lambda\neq(q))
\]
矛盾。故唯一全局最大化分拆只能是 \((q)\)。这证明了第 \(1\) 条。

由于 generic 假设排除了与任何其他分拆的平手，故 \((q)\) 一旦成为唯一最大值，就自动得到
\[
\Delta_{\mathrm{part}}(q)=P(q)-\max_{\nu\neq(q)}\mathcal P(\nu)>0.
\]
第 \(2\) 条成立。

最后，第 \(3\) 条由定理 \ref{thm:pom-partition-gap-positive-kills-hooks} 立即推出：严格正的分拆间隙强制所有非平凡 hook 通道满足
\[
B_{q,\lambda}=0,
\]
等价地
\[
T_{q,\lambda}(m)=\Tr(B_{q,\lambda}^m)=0
\qquad(\forall\,m\ge 1).
\]
证毕。
\end{proof}

\noindent
综上，平凡通道窗口与 primitive 坐标之间并非仅有可逆关系，而是形成单位上三角的整系数多项式自同构；因此固定窗口内的全部 Schur 断层都不再携带超出平凡通道的新代数信息。进一步，在恒等扇区严格主导时，Schur 通道归一化后刚性收敛到对称群的 Plancherel 测度，而其非平凡总能量比例趋于纯表示论常数 \(1-\frac{1}{q!}\)。与此同时，在 generic pressure chamber 中，任意两条通道的全部渐近比较都压缩为主导分拆上的角色比值与线性压力差；一旦再叠加 Perron-purity，则唯一最大分拆被强制锁定为 \((q)\)，并由此导致全部非平凡 hook 通道整体湮灭。


\paragraph{有限 jet 决定性、三模碰撞缺口、无限分支账本与 Smith/prime-axis 刚性}
在 fold-gauge 常数的 Bernoulli 递阶剥离、Fibonacci 共振双频的碰撞地板、自然扩张的分层 prime-register 外置、Smith 局部亏损谱的完全分类以及整数轴对齐椭圆的自由原子分解都已建立之后，还可进一步抽取出下列跨链结论。它们分别把有限奇指数 jet 的决定性、零模重整化的失效、无限深度历史的有限层账本不可能性、模 \(p^k\) 核增长谱对 Smith 正规形的完备判别，以及整数轴对齐椭圆幺半群上同态的 prime-axis 刚性统一为一组有限决定性命题。

\begin{theorem}[折叠规范常数的有限 \texorpdfstring{$\varphi$}{phi}-adic jet 完全决定有限偶 \texorpdfstring{$\zeta$}{zeta} 塔]\label{thm:xi-time-part63b-logcm-phiadic-jet-determines-even-zeta}
\leanverified{paper\_xi\_time\_part63b\_logcm\_phiadic\_jet\_determines\_even\_zeta}
记
\[
q:=\frac{\varphi}{2}\in(0,1),
\qquad
L_m:=\log C_m-\log(2\pi).
\]
则对任意固定 \(R\ge 1\)，\(L_m\) 在精度
\[
O\!\bigl(q^{(2R+1)m}\bigr)
\]
下的截断渐近类唯一决定 \(B_2,\dots,B_{2R}\)，从而唯一决定
\[
\zeta(2),\zeta(4),\dots,\zeta(2R).
\]
等价地，fold-gauge 常数的有限 \(\varphi\)-adic jet 已经完整编码了前 \(R\) 个偶 Bernoulli 数与前 \(R\) 个偶值 \(\zeta\)-塔。
\end{theorem}
\begin{proof}
由定理 \ref{thm:fold-bin-gauge-constant-stirling-bernoulli-hierarchy}，对任意固定 \(R\ge 1\) 都有展开
\[
L_m
=
\sum_{r=1}^{R}
\frac{B_{2r}}{r(2r-1)}
\bigl(\varphi^{-1}+\varphi^{2r-3}\bigr)q^{(2r-1)m}
+O\!\bigl(q^{(2R+1)m}\bigr).
\]
因此该截断 jet 的自由系数恰为
\[
a_r:=
\frac{B_{2r}}{r(2r-1)}
\bigl(\varphi^{-1}+\varphi^{2r-3}\bigr)
\qquad(1\le r\le R).
\]
另一方面，定理 \ref{thm:xi-time-part60ab-gauge-constant-recursive-bernoulli-stripping} 已经给出逐阶极限消元反演：这些系数可由 \(L_m\) 递归无损恢复。又因
\[
\varphi^{-1}+\varphi^{2r-3}>0
\qquad(r\ge 1),
\]
故 \(a_r\) 与 \(B_{2r}\) 一一对应，从而前 \(R\) 个偶 Bernoulli 数由该 jet 唯一确定。

最后，由标准恒等式
\[
B_{2r}
=
(-1)^{r-1}\frac{2(2r)!}{(2\pi)^{2r}}\zeta(2r)
\qquad(r\ge 1),
\]
便可逐项恢复 \(\zeta(2),\dots,\zeta(2R)\)。证毕。
\end{proof}

\begin{corollary}[\texorpdfstring{$(\log C_m)$}{(log Cm)} 的完整渐近 germ 对偶值 \texorpdfstring{$\zeta$}{zeta} 塔是忠实不变量]\label{cor:xi-time-part63b-logcm-asymptotic-germ-faithful-even-zeta}
\leanverified{paper\_xi\_time\_part63b\_logcm\_asymptotic\_germ\_faithful\_even\_zeta}
沿用定理 \ref{thm:xi-time-part63b-logcm-phiadic-jet-determines-even-zeta} 的记号。设 \((C_m)_{m\ge 1}\) 与 \((C_m')_{m\ge 1}\) 为两组折叠规范常数序列，并记
\[
L_m:=\log C_m-\log(2\pi),
\qquad
L_m':=\log C_m'-\log(2\pi).
\]
对每个 \(r\ge 1\)，记 \(B_{2r}(C)\)、\(\zeta_C(2r)\) 分别为由序列 \((C_m)\) 依定理 \ref{thm:xi-time-part63b-logcm-phiadic-jet-determines-even-zeta} 所恢复出的第 \(2r\) 个 Bernoulli 数与偶值 \(\zeta\)-数据；对 \((C_m')\) 亦作同样约定。
若对每个 \(R\ge 1\) 都有
\[
L_m-L_m'
=
O\!\bigl(q^{(2R+1)m}\bigr)
\qquad(m\to\infty),
\]
则
\[
B_{2r}(C)=B_{2r}(C'),
\qquad
\zeta_C(2r)=\zeta_{C'}(2r)
\qquad(\forall r\ge 1).
\]
换言之，\((\log C_m)\) 在全部奇指数层上的渐近 germ，对偶值 \(\zeta\)-塔给出一个忠实不变量。
\end{corollary}
\begin{proof}
任取 \(R\ge 1\)。假设 \(L_m\) 与 \(L_m'\) 在 \(O(q^{(2R+1)m})\) 精度上一致，则它们的前 \(R\) 阶 \(\varphi\)-adic jet 相同。由定理 \ref{thm:xi-time-part63b-logcm-phiadic-jet-determines-even-zeta}，两者的
\[
B_2,\dots,B_{2R}
\]
以及
\[
\zeta(2),\dots,\zeta(2R)
\]
必逐项一致。由于 \(R\) 任意，遂得对每个 \(r\ge 1\) 都有
\[
B_{2r}(C)=B_{2r}(C'),
\qquad
\zeta_C(2r)=\zeta_{C'}(2r).
\]
证毕。
\end{proof}

\begin{theorem}[三模不可去缺口与零模不可重整化]\label{thm:xi-time-part63b-three-mode-gap-zero-mode-nonrenormalizable}
\leanverified{paper\_xi\_time\_part63b\_three\_mode\_gap\_zero\_mode\_nonrenormalizable}
设 \(\mathsf{Col}_m\) 为 bin-fold 的碰撞概率。则相对于均匀基线 \(F_{m+2}^{-1}\)，有
\[
\liminf_{m\to\infty}
F_{m+2}\left(\mathsf{Col}_m-\frac{1}{F_{m+2}}\right)
\ge
2C_\varphi^2
>
0.
\]
并且这一正下界已经由三模
\[
k=0,\qquad k=F_m,\qquad k=F_{m+1}
\]
单独强制。因而任何只依赖零模归一化的均匀化重整，都不可能消除此常数量级缺口。
\end{theorem}
\begin{proof}
由定理 \ref{thm:xi-time-part9p-chi2-ranktwo-resonant-floor}，
\[
\chi^2\!\bigl(p_m^{\mathrm{bin}}\Vert u_m\bigr)
=
F_{m+2}\mathsf{Col}_m-1,
\]
并且若删去两条 Fibonacci 共振频率 \(F_m,F_{m+1}\)，则仍有
\[
\chi^2\!\bigl(p_m^{\mathrm{bin}}\Vert u_m\bigr)-\chi^2_{\mathrm{cut}}(m)
=
a_m(F_m)^2+a_m(F_{m+1})^2
=
2C_\varphi^2+o(1).
\]
于是
\[
F_{m+2}\left(\mathsf{Col}_m-\frac{1}{F_{m+2}}\right)
=
\chi^2\!\bigl(p_m^{\mathrm{bin}}\Vert u_m\bigr)
\ge
2C_\varphi^2+o(1),
\]
取下极限即得所示不等式。

这里 \(k=0\) 模只贡献均匀基线
\(
F_{m+2}^{-1}
\),
而严格正的超额全部来自两条非零频率 \(F_m,F_{m+1}\)。因此，只在零模上作归一化而不消去这对 bulk 共振的任何方案，都不可能把碰撞概率重新压回
\(
F_{m+2}^{-1}+o(F_{m+2}^{-1})
\)
的尺度。证毕。
\end{proof}

\begin{theorem}[无限深度有限纤维塔的有限层 prime-register 不可能性]\label{thm:xi-time-part63b-infinite-depth-finite-prime-register-impossibility}
\leanverified{paper\_xi\_time\_part63b\_infinite\_depth\_finite\_prime\_register\_impossibility}
设
\[
X_0\xrightarrow{\pi_0}X_1\xrightarrow{\pi_1}X_2\xrightarrow{\pi_2}\cdots
\]
为可数集合上的有限纤维满射链，并假设对每个 \(j\ge 0\) 都存在 \(y_j\in X_{j+1}\) 使
\[
\#\pi_j^{-1}(y_j)\ge 2.
\]
记其逆极限为
\[
\widehat X
:=
\varprojlim(X_j,\pi_j)
=
\Bigl\{(x_0,x_1,x_2,\dots)\in \prod_{j\ge 0}X_j:\ \pi_j(x_j)=x_{j+1}\ \forall j\ge 0\Bigr\}.
\]
则在定理 \ref{thm:xi-layered-prime-register-sharp-natural-extension} 的分层 prime-register 外置语义下，不存在任何仅使用有限多个预先固定 prime slices 的精确单射实现，能够在 \(\widehat X\) 上同时保持全历史可恢复性。

更精确地，对任意 \(k\ge 1\)，深度 \(k\) 前缀的最坏情形精确恢复至少需要 \(k\) 个彼此独立的 prime slices；而在常系数情形 \(X_j=X\)、\(\pi_j=T\) 下，文稿中的自然扩张分支指标编码恰以一侧无穷序列实现这一极限最小方案。
\end{theorem}
\begin{proof}
固定 \(k\ge 1\)。把定理 \ref{thm:xi-layered-prime-register-sharp-natural-extension} 应用于截断链
\[
X_0\xrightarrow{\pi_0}X_1\xrightarrow{\pi_1}\cdots\xrightarrow{\pi_{k-1}}X_k.
\]
由于每一层都存在至少一个真实分支点，最坏情形下深度 \(k\) 的精确分层外置至少需要 \(k\) 个有效 prime slices；并且该定理同时给出恰在每层写入一枚有效素数的单射构造，说明这个下界是可达的。

反设整个逆极限 \(\widehat X\) 存在某个只使用有限多个 prime slices 的精确单射外置。把该编码限制到前 \(k\) 层坐标，便得到同一有限 prime-slice 预算下的深度 \(k\) 精确恢复方案。取 \(k\) 大于可用 prime slices 的总数，即与上面的必要性矛盾。故有限层 prime-register 不可能精确实现整个无穷历史。

在常系数情形 \(X_j=X\)、\(\pi_j=T\) 下，定理 \ref{thm:killo-natural-extension-branch-register} 已把自然扩张共轭为右移插入系统
\[
(x_0,x_{-1},x_{-2},\dots)\longmapsto
\bigl(x_0,(\beta(x_{-1}),\beta(x_{-2}),\dots)\bigr),
\]
每一步只写入一个新的分支指标；因此它正是有限深度最优外置在逆极限上的极限闭合。证毕。
\end{proof}

\begin{theorem}[Smith 等价的核增长谱判别准则]\label{thm:xi-time-part63b-smith-kernel-growth-spectrum-criterion}
\leanverified{paper\_xi\_time\_part63b\_smith\_kernel\_growth\_spectrum\_criterion}
设
\[
A,B\in M_{m\times n}(\ZZ),
\qquad
\operatorname{rank}_{\QQ}(A)=\operatorname{rank}_{\QQ}(B)=m.
\]
则下列条件等价：
\[
\#\ker(A\bmod p^k)=\#\ker(B\bmod p^k)
\qquad
(\forall p\ \text{素数},\ \forall k\ge 1);
\]
\[
A,\ B\ \text{具有完全相同的 Smith 不变量}.
\]
换言之，整数线性投影在全部模 \(p^k\) 上的核增长谱，恰是其 Smith 正规形的完备判别式。
\end{theorem}
\begin{proof}
设 \(A\) 的 Smith 不变量为
\[
d_1^{A}\mid \cdots \mid d_m^{A},
\]
并记
\[
e_i^{A}(p):=v_p(d_i^{A}).
\]
对每个素数 \(p\) 与整数 \(k\ge 1\)，由定理 \ref{thm:xi-kernel-loss-divisibility-valuation-rational-euler}，
\[
\#\ker(A\bmod p^k)
=
p^{\,k(n-m)+\sum_{i=1}^{m}\min(k,e_i^{A}(p))}.
\]
因此若定义
\[
f_{A,p}(k):=
\log_p\#\ker(A\bmod p^k)-k(n-m),
\]
则
\[
f_{A,p}(k)=\sum_{i=1}^{m}\min(k,e_i^{A}(p)).
\]
于是其一阶差分满足
\[
\Delta_{A,p}(k):=
f_{A,p}(k)-f_{A,p}(k-1)
=
\#\{i:\ e_i^{A}(p)\ge k\}.
\]
故序列 \(\bigl(\#\ker(A\bmod p^k)\bigr)_{k\ge 1}\) 唯一决定多重集
\[
\{e_i^{A}(p)\}_{i=1}^{m}
\]
；逐素数拼合便唯一决定 \(A\) 的全部 Smith 不变量。对 \(B\) 亦然。

因此，若 \(A\) 与 \(B\) 的核增长谱逐点相同，则对每个 \(p\) 都有
\[
f_{A,p}(k)=f_{B,p}(k)
\qquad(\forall k\ge 1),
\]
从而
\[
\{e_i^{A}(p)\}_{i=1}^{m}
=
\{e_i^{B}(p)\}_{i=1}^{m}.
\]
逐素数合并便得两者 Smith 正规形相同。反向 implication 显然：相同的 Smith 不变量给出相同的模 \(p^k\) 核大小公式。证毕。
\end{proof}

\begin{theorem}[信息亏损 \texorpdfstring{$\zeta$}{zeta} 函数的 Euler 因子化与 Smith 刚性]\label{thm:xi-time-part63b-loss-zeta-euler-factorization-smith-rigidity}
\leanverified{paper\_xi\_time\_part63b\_loss\_zeta\_euler\_factorization\_smith\_rigidity}
设
\[
A\in M_{m\times n}(\ZZ),
\qquad
\operatorname{rank}_{\QQ}(A)=m,
\]
其 Smith 不变量为
\[
d_1\mid\cdots\mid d_m.
\]
对每个素数 \(p\) 记
\[
e_i(p):=v_p(d_i).
\]
定义信息亏损 \(\zeta\)-函数
\[
\mathcal Z_A(s)
:=
\sum_{b\ge 1}\frac{\#\ker(A_b)}{b^{\,s+n-m}}
\qquad(\Re s\gg 1).
\]
则有 Euler 因子化
\[
\mathcal Z_A(s)
=
\prod_p
\left(
\sum_{k\ge 0}
p^{-ks+\sum_{i=1}^{m}\min(k,e_i(p))}
\right).
\]
特别地，对同型尺寸的满秩整数矩阵 \(A,B\)，有
\[
\mathcal Z_A=\mathcal Z_B
\quad\Longleftrightarrow\quad
A,\ B\ \text{具有相同的 Smith 正规形}.
\]
\end{theorem}
\begin{proof}
由推论 \ref{cor:killo-kernel-size-general-modulus}，
\[
\#\ker(A_b)
=
b^{\,n-m}\prod_{i=1}^{m}\gcd(d_i,b).
\]
因此归一化系数
\[
a_A(b):=
\frac{\#\ker(A_b)}{b^{\,n-m}}
=
\prod_{i=1}^{m}\gcd(d_i,b)
\]
是乘法函数，从而 \(\mathcal Z_A(s)=\sum_{b\ge 1}a_A(b)b^{-s}\) 具有 Euler 乘积。对 \(b=p^k\) 进一步有
\[
a_A(p^k)
=
\prod_{i=1}^{m}p^{\min(k,e_i(p))}
=
p^{\sum_{i=1}^{m}\min(k,e_i(p))},
\]
故逐素数展开正是
\[
\mathcal Z_A(s)
=
\prod_p
\left(
\sum_{k\ge 0}
p^{-ks+\sum_{i=1}^{m}\min(k,e_i(p))}
\right).
\]

若 \(\mathcal Z_A=\mathcal Z_B\)，则两者的 Dirichlet 系数逐项相同，特别地对每个素数幂 \(p^k\) 都有
\[
\#\ker(A\bmod p^k)\,p^{-k(n-m)}
=
\#\ker(B\bmod p^k)\,p^{-k(n-m)}.
\]
于是
\[
\#\ker(A\bmod p^k)=\#\ker(B\bmod p^k)
\qquad(\forall p,\ \forall k\ge 1).
\]
由定理 \ref{thm:xi-time-part63b-smith-kernel-growth-spectrum-criterion}，\(A\) 与 \(B\) 具有相同的 Smith 正规形。反向 implication 显然：相同的 Smith 不变量给出相同的 Euler 因子。证毕。
\end{proof}

\begin{proposition}[整数轴对齐椭圆单子的 prime-axis 刚性]\label{prop:xi-time-part63b-integer-axis-ellipse-prime-axis-rigidity}
沿用定理 \ref{thm:xi-integer-ellipse-free-commutative-monoid} 的记号，令
\[
\mathsf E_{\NN}
:=
\{E_{a,b}:x^2/a^2+y^2/b^2=1,\ a,b\in\NN_{>0}\},
\qquad
E_{a,b}\odot E_{c,d}:=E_{ac,bd}.
\]
则
\[
(\mathsf E_{\NN},\odot)
\cong
\bigl(\NN^{(\mathcal P)}\times \NN^{(\mathcal P)},+\bigr).
\]
因而对任意交换幺半群 \((M,+)\)，每个幺半群同态
\[
J:(\mathsf E_{\NN},\odot)\to (M,+)
\]
都唯一由素数轴生成元
\[
E_{p,1},\qquad E_{1,p}
\qquad(p\in\mathcal P)
\]
的像所决定。特别地，当 \(M=\RR\) 时，存在唯一系数族 \((\alpha_p)_{p\in\mathcal P}\)、\((\beta_p)_{p\in\mathcal P}\) 使得对任意 \(a,b\in\NN_{>0}\) 都有
\[
J(E_{a,b})
=
\sum_{p}\alpha_p\,v_p(a)
+
\sum_{p}\beta_p\,v_p(b).
\]
\end{proposition}
\begin{proof}
定理 \ref{thm:xi-integer-ellipse-free-commutative-monoid} 已给出幺半群同构
\[
(\NN^{(\mathcal P)}\times \NN^{(\mathcal P)},+)
\cong
(\mathsf E_{\NN},\odot),
\qquad
(r_1,r_2)\longmapsto E_{N(r_1),N(r_2)}.
\]
其中 \(\NN^{(\mathcal P)}\) 是以素数集 \(\mathcal P\) 为基的自由交换幺半群。故 \((\mathsf E_{\NN},\odot)\) 本身就是以
\[
\{E_{p,1},E_{1,p}:p\in\mathcal P\}
\]
为原子集的自由交换幺半群。

自由对象上的幺半群同态由生成元像唯一决定，于是任意
\[
J:(\mathsf E_{\NN},\odot)\to (M,+)
\]
都完全由 \(J(E_{p,1})\) 与 \(J(E_{1,p})\) 唯一决定。

当 \(M=\RR\) 时，记
\[
\alpha_p:=J(E_{p,1}),
\qquad
\beta_p:=J(E_{1,p}).
\]
又由唯一原子分解
\[
E_{a,b}
=
\left(\bigodot_{p}E_{p,1}^{\odot v_p(a)}\right)
\odot
\left(\bigodot_{p}E_{1,p}^{\odot v_p(b)}\right),
\]
应用幺半群同态性即可得到
\[
J(E_{a,b})
=
\sum_{p}\alpha_p\,v_p(a)
+
\sum_{p}\beta_p\,v_p(b).
\]
由于原子分解唯一，系数族亦唯一。证毕。
\end{proof}
\leanverified{paper\_xi\_time\_part63b\_integer\_axis\_ellipse\_prime\_axis\_rigidity}

\paragraph{Peter--Weyl 因子化与覆盖 \texorpdfstring{$\zeta$}{zeta} 极点秩序的维数加权律}
在正规覆盖的 Peter--Weyl 块对角化下，覆盖行列式按不可约通道作 Artin 因子化。由此，任一给定特征值对覆盖 \(\zeta\) 主点的贡献都可逐通道累加，并被严格记录为“表示维数 \(\times\) 通道重数”的和。

\begin{theorem}[正则覆盖 \texorpdfstring{$\zeta$}{zeta} 的极点阶数等于不可约通道维数加权的特征值重数和]\label{thm:xi-time-part64-cover-zeta-pole-order-weighted-channel-multiplicity}
沿用命题 \ref{prop:kernel-artin-factorization} 的记号。
设 \(G\) 为有限群，\(M_{K,G}(u)\) 为群扩张邻接算子，并记
\[
\det(I-zM_{K,G}(u))
=
\prod_{\rho\in\widehat G}
\det(I-zM_{K,\rho}(u))^{d_\rho},
\qquad
d_\rho:=\dim\rho.
\]
取任意复数 \(\Lambda\neq 0\)，并令 \(m_\rho(\Lambda)\) 表示 \(\Lambda\) 作为 \(M_{K,\rho}(u)\) 特征值的代数重数。
则覆盖 \(\zeta\) 函数
\[
\zeta_{K,G}(z,u):=\det(I-zM_{K,G}(u))^{-1}
\]
在 \(z=\Lambda^{-1}\) 处的极点阶数恰为
\[
\boxed{\;
\operatorname{ord}_{z=\Lambda^{-1}}\zeta_{K,G}(z,u)
=
\sum_{\rho\in\widehat G} d_\rho\,m_\rho(\Lambda)\; }.
\]
\end{theorem}
\begin{proof}
对每个 \(\rho\in\widehat G\)，在 \(z=\Lambda^{-1}\) 的邻域内可写
\[
\det(I-zM_{K,\rho}(u))
=(1-z\Lambda)^{m_\rho(\Lambda)}h_\rho(z),
\qquad
h_\rho(\Lambda^{-1})\neq 0.
\]
将其代入命题 \ref{prop:kernel-artin-factorization} 的 Artin 因子化，得到
\[
\det(I-zM_{K,G}(u))
=
\prod_{\rho\in\widehat G}
\bigl((1-z\Lambda)^{m_\rho(\Lambda)}h_\rho(z)\bigr)^{d_\rho}
\]
\[
=(1-z\Lambda)^{\sum_{\rho}d_\rho m_\rho(\Lambda)}\,H(z),
\qquad
H(\Lambda^{-1})\neq 0.
\]
取倒数即得
\[
\zeta_{K,G}(z,u)
=(1-z\Lambda)^{-\sum_{\rho}d_\rho m_\rho(\Lambda)}H(z)^{-1},
\]
故其在 \(z=\Lambda^{-1}\) 处的极点阶数恰为
\(
\sum_{\rho}d_\rho m_\rho(\Lambda)
\)。
证毕。
\end{proof}

\begin{corollary}[主极点的覆盖稳定性判据]\label{cor:xi-time-part64-cover-zeta-leading-pole-stability}
沿用定理 \ref{thm:xi-time-part64-cover-zeta-pole-order-weighted-channel-multiplicity} 的记号，并令
\[
\lambda_1(u)
\]
表示平凡通道 \(M_{K,\mathbf 1}(u)\) 的 Perron 根。
若满足：
\begin{enumerate}
  \item \(\lambda_1(u)\) 在平凡通道中是简单特征值；
  \item 对所有非平凡 \(\rho\neq \mathbf 1\)，都有
  \[
  \operatorname{spr}\!\bigl(M_{K,\rho}(u)\bigr)<\lambda_1(u),
  \]
  其中 \(\operatorname{spr}\) 表示谱半径，
\end{enumerate}
则覆盖 \(\zeta\) 在主点 \(z=\lambda_1(u)^{-1}\) 处满足
\[
\boxed{\;
\operatorname{ord}_{z=\lambda_1(u)^{-1}}\zeta_{K,G}(z,u)=1\; }.
\]
亦即有限正则覆盖在主谱层面不改变极点阶数，只可能修改次主通道。

更一般地，若恰有一族非平凡通道
\[
\mathcal R\subseteq \widehat G\setminus\{\mathbf 1\}
\]
满足
\(
\lambda_1(u)\in\Spec(M_{K,\rho}(u))
\)
对每个 \(\rho\in\mathcal R\) 成立，而对其余非平凡通道 \(\eta\notin\mathcal R\) 有
\(
\lambda_1(u)\notin\Spec(M_{K,\eta}(u))
\)，
则
\[
\operatorname{ord}_{z=\lambda_1(u)^{-1}}\zeta_{K,G}(z,u)
=
1+\sum_{\rho\in\mathcal R}d_\rho\,m_\rho(\lambda_1(u)).
\]
\end{corollary}
\begin{proof}
在平凡通道 \(\mathbf 1\) 上，\(\lambda_1(u)\) 的代数重数为 \(1\)。
若对所有非平凡 \(\rho\neq\mathbf 1\) 都有
\(
\operatorname{spr}(M_{K,\rho}(u))<\lambda_1(u)
\)，
则这些通道根本不可能含有特征值 \(\lambda_1(u)\)，故
\[
m_\rho(\lambda_1(u))=0
\qquad(\rho\neq\mathbf 1).
\]
由定理 \ref{thm:xi-time-part64-cover-zeta-pole-order-weighted-channel-multiplicity}，
\[
\operatorname{ord}_{z=\lambda_1(u)^{-1}}\zeta_{K,G}(z,u)
=
d_{\mathbf 1}\,m_{\mathbf 1}(\lambda_1(u))
=1.
\]

一般情形下，仍由同一定理，
\[
\operatorname{ord}_{z=\lambda_1(u)^{-1}}\zeta_{K,G}(z,u)
=
d_{\mathbf 1}\,m_{\mathbf 1}(\lambda_1(u))
+\sum_{\rho\in\mathcal R}d_\rho\,m_\rho(\lambda_1(u)),
\]
而 \(d_{\mathbf 1}=1\)、\(m_{\mathbf 1}(\lambda_1(u))=1\)，遂得所示公式。证毕。
\end{proof}

\noindent
于是 Peter--Weyl/Artin 因子化不仅把覆盖 \(\zeta\) 分解为有限多个不可约通道；它还把主极点的升阶机制完全刚性化：覆盖只能按表示维数与通道共振重数的乘积累加极点阶数，而在非平凡通道不触及主谱时，主极点保持简单。

\paragraph{非阿贝尔通道能量账本与覆盖 Perron 主相的平凡通道刚性}
在非阿贝尔 Chebotarev--Plancherel 能量恒等式、群扩张邻接算子的 Peter--Weyl 块对角化以及覆盖 \(\zeta\) 的 Artin 因子化已经建立之后，还可进一步抽取出下列两条严格后果：商映射在类函数层面损失的 \(L^2\) 能量可逐不可约通道精确记账，而正规覆盖算子的 Perron 根则始终由平凡通道完全控制。

\begin{theorem}[非阿贝尔商映射的精确能量损失公式]\label{thm:xi-time-part64-nonabelian-quotient-energy-loss-exact}
\leanverified{paper\_xi\_time\_part64\_nonabelian\_quotient\_energy\_loss\_exact}
沿用命题 \ref{prop:kernel-chebotarev-plancherel-nonabelian} 的记号。设
\[
\varphi:G\twoheadrightarrow H
\]
为有限群满同态，并记
\[
K:=\ker\varphi.
\]
令 \(\varepsilon\in \mathrm{Cl}(G)\) 为零均值类函数，即
\[
\sum_{g\in G}\varepsilon(g)=0.
\]
对每个 \(\rho\in\widehat G\) 定义不可约通道幅度
\[
B_G(\rho):=\sum_{g\in G}\varepsilon(g)\,\overline{\chi_\rho(g)},
\]
并定义纤维平均推前
\[
(\varphi_\ast\varepsilon)(h):=\frac1{|K|}\sum_{g\in\varphi^{-1}(h)}\varepsilon(g)
\qquad(h\in H).
\]
再记
\[
\mathrm{Inf}_\varphi(\widehat H)
:=
\{\widetilde \sigma:\ \sigma\in \widehat H,\ \widetilde \sigma=\sigma\circ\varphi\}
\subseteq \widehat G.
\]
则在归一化均匀测度对应的 \(L^2\) 范数下有严格恒等式
\[
\boxed{\;
\|\varepsilon\|_{L^2(G)}^2-\|\varphi_\ast\varepsilon\|_{L^2(H)}^2
=
\frac1{|G|^2}
\sum_{\rho\in\widehat G\setminus \mathrm{Inf}_\varphi(\widehat H)}
\bigl|B_G(\rho)\bigr|^2.\;}
\]
\end{theorem}
\begin{proof}
先由
\[
\sum_{h\in H}(\varphi_\ast\varepsilon)(h)
=
\frac1{|K|}\sum_{g\in G}\varepsilon(g)=0
\]
可知 \(\varphi_\ast\varepsilon\) 也是 \(H\) 上的零均值类函数。

对 \(G\) 应用命题 \ref{prop:kernel-chebotarev-plancherel-nonabelian}，得到
\[
\|\varepsilon\|_{L^2(G)}^2
=
\frac1{|G|^2}
\sum_{\rho\in\widehat G\setminus\{\mathbf 1_G\}}
\bigl|B_G(\rho)\bigr|^2.
\]

对 \(H\) 上的 \(\varphi_\ast\varepsilon\) 同样应用该命题。若 \(\sigma\in \widehat H\)，定义
\[
B_H(\sigma):=\sum_{h\in H}(\varphi_\ast\varepsilon)(h)\,\overline{\chi_\sigma(h)}.
\]
则由纤维平均定义，
\[
\begin{aligned}
B_H(\sigma)
&=
\sum_{h\in H}
\frac1{|K|}
\sum_{g\in\varphi^{-1}(h)}
\varepsilon(g)\,\overline{\chi_\sigma(h)}\\
&=
\frac1{|K|}
\sum_{g\in G}
\varepsilon(g)\,\overline{\chi_\sigma(\varphi(g))}\\
&=
\frac1{|K|}
\sum_{g\in G}
\varepsilon(g)\,\overline{\chi_{\widetilde \sigma}(g)}\\
&=
\frac1{|K|}B_G(\widetilde \sigma).
\end{aligned}
\]
由 \(|G|=|H||K|\) 遂得
\[
\frac1{|H|^2}\bigl|B_H(\sigma)\bigr|^2
=
\frac1{|G|^2}\bigl|B_G(\widetilde \sigma)\bigr|^2.
\]
因此
\[
\|\varphi_\ast\varepsilon\|_{L^2(H)}^2
=
\frac1{|G|^2}
\sum_{\rho\in \mathrm{Inf}_\varphi(\widehat H)\setminus\{\mathbf 1_G\}}
\bigl|B_G(\rho)\bigr|^2.
\]
从 \(\|\varepsilon\|_{L^2(G)}^2\) 的通道展开中减去上式，即得所示恒等式。证毕。
\end{proof}

\begin{theorem}[正规覆盖算子的 Perron 根完全由平凡通道控制]\label{thm:xi-time-part64-cover-perron-root-trivial-channel}
沿用命题 \ref{prop:kernel-peter-weyl-block-diagonalization}、命题 \ref{prop:kernel-artin-factorization} 以及定义 \ref{def:kernel-rh} 的记号，并假设 \(M_K(u)\) 原始。记
\[
\lambda_1(u):=\rho\!\bigl(M_K(u)\bigr)=\rho\!\bigl(M_{K,\mathbf 1}(u)\bigr).
\]
则对每个不可约表示 \(\pi\in\widehat G\)，都有
\[
\boxed{\;
\rho\!\bigl(M_{K,\pi}(u)\bigr)\le \lambda_1(u).\;}
\]
并且整个覆盖算子的谱半径满足
\[
\boxed{\;
\rho\!\bigl(M_{K,G}(u)\bigr)=\lambda_1(u).\;}
\]
进一步，若记
\[
\eta(u):=\max_{\pi\neq \mathbf 1}\rho\!\bigl(M_{K,\pi}(u)\bigr),
\qquad
\Lambda(u):=\max\{|\lambda_2(u)|,\eta(u)\},
\]
则下列两条等价：
\begin{enumerate}
  \item 平凡通道的次特征值与全部非平凡通道谱半径共同满足平方根界，即
  \[
  \Lambda(u)\le \lambda_1(u)^{1/2};
  \]
  \item \(\textsf{RH}_K(u)\) 与 \(\textsf{GRH}_{K,G}(u)\) 同时成立。
\end{enumerate}
\end{theorem}
\begin{proof}
由命题 \ref{prop:kernel-peter-weyl-block-diagonalization}，
\[
M_{K,G}(u)
\ \cong\
\bigoplus_{\pi\in\widehat G}
\bigl(M_{K,\pi}(u)\otimes I_{\dim\pi}\bigr),
\]
故
\[
\rho\!\bigl(M_{K,G}(u)\bigr)
=
\max_{\pi\in\widehat G}\rho\!\bigl(M_{K,\pi}(u)\bigr).
\]
因此只需证明对每个 \(\pi\) 都有
\(
\rho(M_{K,\pi}(u))\le \lambda_1(u)
\)。

固定 \(\pi\in\widehat G\)。由于 \(G\) 有限，\(\pi\) 可酉化；取 \(V_\pi\) 上一个 \(G\)-不变 Hermitian 范数，使 \(\pi(g)\) 全为酉算子。设
\[
x=(x_i)_{i\in V}\in V_\pi^{\,V}
\]
是 \(M_{K,\pi}(u)\) 的特征向量，对应特征值 \(\mu\)。按定义
\[
(M_{K,\pi}(u)x)_i
=
\sum_{e:i\to j}u^{\kappa(e)}\,\pi(g(e))\,x_j.
\]
逐坐标取范数，并记
\[
h_i:=\|x_i\|\ge 0,
\qquad
h:=(h_i)_{i\in V}\neq 0,
\]
则由三角不等式与酉不变性，
\[
\begin{aligned}
|\mu|\,h_i
&=
\bigl\|(M_{K,\pi}(u)x)_i\bigr\|\\
&\le
\sum_{e:i\to j}u^{\kappa(e)}\,\|x_j\|\\
&=
(M_K(u)h)_i.
\end{aligned}
\]
故
\[
|\mu|\,h\le M_K(u)h
\]
在非负锥上逐坐标成立。由 Perron--Frobenius 的 Collatz--Wielandt 公式，
\[
|\mu|\le \rho(M_K(u))=\lambda_1(u).
\]
对 \(M_{K,\pi}(u)\) 的所有特征值取上确界，便得
\[
\rho\!\bigl(M_{K,\pi}(u)\bigr)\le \lambda_1(u).
\]

另一方面，平凡通道 \(\pi=\mathbf 1\) 上恰有
\[
\rho\!\bigl(M_{K,\mathbf 1}(u)\bigr)=\lambda_1(u),
\]
因此
\[
\rho\!\bigl(M_{K,G}(u)\bigr)=\lambda_1(u).
\]

最后，由定义
\[
\Lambda(u)=\max\{|\lambda_2(u)|,\eta(u)\},
\]
第 \(1\) 条本身即是
\[
\Lambda(u)\le \lambda_1(u)^{1/2},
\]
而这正等价于 \(\textsf{RH}_K(u)\) 与 \(\textsf{GRH}_{K,G}(u)\) 同时成立。证毕。
\end{proof}

\paragraph{Chebotarev 商等号、最小承载商与极值见证集}
在阿贝尔与非阿贝尔两侧的 Chebotarev--Plancherel 能量恒等式、商单调性以及二主项见证定理都已建立之后，还可把“偏差在商下何时无损”“全部可见偏差究竟由哪个最小商承载”以及“最优见证集如何从极大通道中被唯一选出”进一步压缩为下列四条接口。

\begin{theorem}[阿贝尔 Chebotarev 偏差在商映射下的等号判据]\label{thm:xi-time-part64aa-abelian-quotient-equality-criterion}
\leanverified{paper\_xi\_time\_part64aa\_abelian\_quotient\_equality\_criterion}
沿用推论 \ref{cor:kernel-chebotarev-L2-monotone} 的记号。设
\[
\pi:G_2\twoheadrightarrow G_1
\]
为有限阿贝尔群满同态，并记
\[
\pi^\ast:\widehat{G_1}\hookrightarrow \widehat{G_2},
\qquad
\chi\longmapsto \chi\circ\pi.
\]
对固定 \(u\in(0,1]\) 与 \(n\ge 1\)，设细层与粗层的相对 Haar 密度分别为
\[
r_n^{(2)}(\,\cdot\,;u)\in \CC^{G_2},
\qquad
r_n^{(1)}(\,\cdot\,;u)\in \CC^{G_1},
\]
并满足
\[
r_n^{(1)}=\mathbb E_\pi(r_n^{(2)}),
\]
其中 \(\mathbb E_\pi\) 表示沿 \(\pi\) 的纤维平均。则下列三条等价：
\begin{enumerate}
  \item
  \[
  \boxed{\;
  \|r_n^{(1)}(\cdot;u)-1\|_{L^2(G_1)}
  =
  \|r_n^{(2)}(\cdot;u)-1\|_{L^2(G_2)}.\;}
  \]
  \item \(r_n^{(2)}(\cdot;u)\) 在 \(\pi\)-纤维上常值，即
  \[
  \boxed{\;
  r_n^{(2)}(c;u)=r_n^{(2)}(c';u)
  \qquad
  \text{只要 }\pi(c)=\pi(c').\;}
  \]
  \item 若把细层扭曲 primitive 谱写作
  \[
  B_n^{(2)}(\chi;u)
  \qquad
  (\chi\in \widehat{G_2}),
  \]
  则一切不来自 \(\pi^\ast(\widehat{G_1})\) 的非平凡角色通道都严格消失：
  \[
  \boxed{\;
  B_n^{(2)}(\chi;u)=0
  \qquad
  \forall\,\chi\in \widehat{G_2}\setminus \pi^\ast(\widehat{G_1}).\;}
  \]
\end{enumerate}
换言之，商映射下 \(L^2\) 偏差不发生任何损失，当且仅当全部可见 Chebotarev 能量已经完全落在该商所允许的角色通道内。
\end{theorem}
\begin{proof}
由推论 \ref{cor:kernel-chebotarev-L2-monotone} 的证明，
\[
r_n^{(1)}-1=\mathbb E_\pi(r_n^{(2)}-1).
\]
通过把 \(L^2(G_1)\) 中的函数 \(h\) 拉回为 \(h\circ\pi\)，可把 \(L^2(G_1)\) 识别为 \(G_2\) 上的纤维常值函数子空间；在此识别下，\(\mathbb E_\pi\) 正是 \(L^2(G_2)\) 到该子空间的正交投影。因此对任意 \(f\in L^2(G_2)\) 都有
\[
\|\mathbb E_\pi f\|_2=\|f\|_2
\iff
f\in \operatorname{Im}(\mathbb E_\pi).
\]
取 \(f=r_n^{(2)}-1\)，便得第 \(1\) 条与第 \(2\) 条等价。

再证第 \(2\) 条与第 \(3\) 条等价。有限阿贝尔群上，沿 \(\pi\) 纤维常值的函数恰是那些在 \(\ker\pi\) 上平凡的角色所张成的子空间；等价地，
\[
\operatorname{Im}(\mathbb E_\pi)
=
\operatorname{span}\bigl(\pi^\ast(\widehat{G_1})\bigr).
\]
另一方面，由命题 \ref{prop:kernel-chebotarev-fourier} 与
\[
r_n^{(2)}(c;u)
=
\frac{|G_2|\,B_{n,c}^{(2)}(u)}{B_n^{(2)}(\mathbf{1};u)}
\]
可得
\[
r_n^{(2)}(c;u)-1
=
\frac{1}{B_n^{(2)}(\mathbf{1};u)}
\sum_{\chi\in \widehat{G_2}\setminus\{\mathbf 1\}}
\overline{\chi(c)}\,B_n^{(2)}(\chi;u).
\]
因此 \(r_n^{(2)}-1\) 属于 \(\operatorname{span}(\pi^\ast(\widehat{G_1}))\)，当且仅当所有不属于 \(\pi^\ast(\widehat{G_1})\) 的角色系数全部为零。这正是第 \(3\) 条。证毕。
\end{proof}

\begin{theorem}[有限阿贝尔标签系统存在唯一最小偏差承载商]\label{thm:xi-time-part64aa-abelian-minimal-deviation-carrier-quotient}
\leanverified{paper\_xi\_time\_part64aa\_abelian\_minimal\_deviation\_carrier\_quotient}
在定理 \ref{thm:xi-time-part64aa-abelian-quotient-equality-criterion} 的阿贝尔设定下，固定 \(u\in(0,1]\) 与 \(n\ge 1\)。定义活跃角色集
\[
\mathcal A_{n,u}
:=
\{\chi\in \widehat G\setminus\{\mathbf 1\}:\ B_n(\chi;u)\neq 0\},
\]
并定义其核交
\[
K_{n,u}
:=
\bigcap_{\chi\in \mathcal A_{n,u}}\ker\chi,
\]
其中在 \(\mathcal A_{n,u}=\varnothing\) 时约定 \(K_{n,u}=G\)。则有：
\begin{enumerate}
  \item 存在唯一函数
  \[
  \widetilde r_{n,u}:G/K_{n,u}\to \CC
  \]
  使得
  \[
  \boxed{\;
  r_n(c;u)=\widetilde r_{n,u}(cK_{n,u})
  \qquad(\forall\,c\in G).\;}
  \]
  \item 对任意有限商
  \[
  \pi:G\twoheadrightarrow H,
  \qquad
  N:=\ker\pi,
  \]
  下列三条等价：
  \[
  \boxed{\;
  r_n(\cdot;u)\ \text{通过 }\pi\text{ 因子化}\;}
  \iff
  \boxed{\;
  N\subseteq K_{n,u}\;}
  \iff
  \boxed{\;
  \|r_n^{H}-1\|_{L^2(H)}
  =
  \|r_n^{G}-1\|_{L^2(G)}.\;}
  \]
  其中 \(r_n^{G}:=r_n\)，而 \(r_n^{H}\) 表示由 \(\pi\) 推前得到的商层相对 Haar 密度。
  \item 因而 \(G/K_{n,u}\) 是保持全部 \(n\)-步 Chebotarev 偏差的唯一最小商：任何其他保持全部偏差的商都唯一地再向下满射到它。
\end{enumerate}
\end{theorem}
\begin{proof}
先证第 \(1\) 条。若 \(k\in K_{n,u}\)，则对每个活跃角色 \(\chi\in\mathcal A_{n,u}\) 都有 \(\chi(k)=1\)。由定理 \ref{thm:xi-time-part64aa-abelian-quotient-equality-criterion} 证明中同一 Fourier 展开，
\[
r_n(ck;u)-1
=
\frac{1}{B_n(\mathbf 1;u)}
\sum_{\chi\in\mathcal A_{n,u}}
\overline{\chi(c)}\,\overline{\chi(k)}\,B_n(\chi;u)
=
r_n(c;u)-1.
\]
故 \(r_n(\cdot;u)\) 在 \(K_{n,u}\)-陪集上常值，从而通过商 \(G/K_{n,u}\) 因子化。唯一性由商映射满射立得。

再证第 \(2\) 条。若 \(r_n(\cdot;u)\) 通过 \(\pi\) 因子化，则它在每条 \(N\)-陪集上常值。由定理 \ref{thm:xi-time-part64aa-abelian-quotient-equality-criterion}，这等价于所有活跃角色都在 \(N\) 上平凡，即
\[
N\subseteq \ker\chi
\qquad
(\forall\,\chi\in\mathcal A_{n,u}),
\]
从而 \(N\subseteq K_{n,u}\)。反之若 \(N\subseteq K_{n,u}\)，则每个活跃角色都在 \(N\) 上平凡，故同一 Fourier 展开只含来自 \(\widehat{G/N}\) 的角色，于是 \(r_n(\cdot;u)\) 必通过 \(\pi\) 因子化。最后，把定理 \ref{thm:xi-time-part64aa-abelian-quotient-equality-criterion} 应用于 \(G\twoheadrightarrow H\)，便得到与 \(L^2\) 等号的等价性。

第 \(3\) 条是第 \(2\) 条的直接推论：若某商 \(G/N\) 保持全部偏差，则必有 \(N\subseteq K_{n,u}\)，因此存在唯一满同态
\[
G/N\twoheadrightarrow G/K_{n,u}.
\]
故 \(G/K_{n,u}\) 在所有保持偏差的商中是唯一最小者。证毕。
\end{proof}

\begin{theorem}[非阿贝尔 Chebotarev--Plancherel 商单调性的等号判据与最小类函数商]\label{thm:xi-time-part64aa-nonabelian-equality-criterion-minimal-class-quotient}
\leanverified{paper\_xi\_time\_part64aa\_nonabelian\_equality\_criterion\_minimal\_class\_quotient}
沿用命题 \ref{prop:kernel-chebotarev-plancherel-nonabelian} 与定理 \ref{thm:xi-time-part64-nonabelian-quotient-energy-loss-exact} 的记号。设 \(G\) 为有限群，\(\varepsilon\in\mathrm{Cl}(G)\) 为零均值类函数，并令
\[
\varphi:G\twoheadrightarrow H
\]
为满同态。把 \(\varepsilon\) 的不可约特征标展开写为
\[
\varepsilon=\sum_{\rho\in\widehat G\setminus\{\mathbf 1\}} a_\rho\,\chi_\rho.
\]
则下列三条等价：
\begin{enumerate}
  \item
  \[
  \boxed{\;
  \|\varphi_\ast\varepsilon\|_{L^2(H)}
  =
  \|\varepsilon\|_{L^2(G)}.\;}
  \]
  \item \(\varepsilon\) 在 \(\varphi\)-纤维上常值，即
  \[
  \boxed{\;
  \varepsilon(g)=\varepsilon(g')
  \qquad
  \text{只要 }\varphi(g)=\varphi(g').\;}
  \]
  \item 一切不通过 \(H\) 的不可约通道都不出现，即
  \[
  \boxed{\;
  a_\rho=0
  \qquad
  \forall\,\rho\in\widehat G
  \text{ 满足 }
  \ker\varphi\nsubseteq \ker\rho.\;}
  \]
\end{enumerate}

进一步，定义
\[
K(\varepsilon)
:=
\bigcap_{a_\rho\neq 0}\ker\rho,
\]
并在 \(\varepsilon=0\) 时约定 \(K(\varepsilon)=G\)。则 \(\varepsilon\) 唯一通过商 \(G/K(\varepsilon)\) 因子化，并且该商是保持全部类函数能量的唯一最小商。
\end{theorem}
\begin{proof}
先证第 \(1\) 条与第 \(2\) 条等价。命题 \ref{prop:kernel-chebotarev-plancherel-nonabelian} 中的 \(\varphi_\ast\) 正是类函数空间上的纤维平均条件期望；而把 \(\mathrm{Cl}(H)\) 中的类函数 \(f\) 拉回为 \(f\circ\varphi\) 后，\(\mathrm{Cl}(H)\) 恰可识别为 \(G\) 上纤维常值类函数子空间。在此识别下，
\[
\|\varphi_\ast\varepsilon\|_2=\|\varepsilon\|_2
\iff
\varepsilon\in \operatorname{Im}(\varphi_\ast),
\]
而右侧正是纤维常值性。

再证第 \(1\) 条与第 \(3\) 条等价。由定理 \ref{thm:xi-time-part64-nonabelian-quotient-energy-loss-exact}，
\[
\|\varepsilon\|_{L^2(G)}^2-\|\varphi_\ast\varepsilon\|_{L^2(H)}^2
=
\frac1{|G|^2}
\sum_{\rho\notin \mathrm{Inf}_\varphi(\widehat H)}
|B_G(\rho)|^2.
\]
又由
\[
a_\rho=\langle \varepsilon,\chi_\rho\rangle
=
\frac{1}{|G|}B_G(\rho),
\]
可知上式右端为零，当且仅当一切不来自 \(\mathrm{Inf}_\varphi(\widehat H)\) 的 \(\rho\) 都满足 \(a_\rho=0\)。而 \(\rho\) 属于 \(\mathrm{Inf}_\varphi(\widehat H)\) 的充要条件正是 \(\ker\varphi\subseteq \ker\rho\)。故第 \(1\) 条与第 \(3\) 条等价。

最后证明最小商断言。若 \(k\in K(\varepsilon)\)，则对每个出现的 \(\rho\) 都有 \(\rho(k)=I\)，故
\[
\chi_\rho(gk)=\chi_\rho(g)
\qquad(\forall\,g\in G).
\]
对全部出现通道求和便得 \(\varepsilon(gk)=\varepsilon(g)\)，从而 \(\varepsilon\) 通过 \(G/K(\varepsilon)\) 因子化。反之，若 \(\varepsilon\) 通过某个商 \(G/N\) 因子化，则由上面三条等价性可知每个出现的 \(\rho\) 都满足 \(N\subseteq \ker\rho\)，于是 \(N\subseteq K(\varepsilon)\)。因此存在唯一满同态
\[
G/N\twoheadrightarrow G/K(\varepsilon),
\]
故 \(G/K(\varepsilon)\) 是唯一最小类函数商。证毕。
\end{proof}

\begin{theorem}[二主项最优见证集由极大通道特征函数的符号层给出]\label{thm:xi-time-part64aa-optimal-chebotarev-witness-sign-slice}
沿用定理 \ref{thm:kernel-chebotarev-second-main-term-witness} 的设定。固定 \(u\in(0,1]\)，并取达到
\[
\Lambda(u)=\max_{\rho\neq \mathbf 1}\operatorname{spr}(M_{K,\rho}(u))
\]
的一个非平凡通道 \(\rho_\ast\) 及其对应本征值 \(\mu_\ast\)，并假设该通道在 \(\mu_\ast\) 处具有谱隙。对每个共轭不变集合 \(A\subseteq G\)，记
\[
U_A:=\pi^{-1}(A),
\]
并把定理 \ref{thm:kernel-chebotarev-second-main-term-witness} 中的二主项系数记作 \(c_A:=c_{U_A}\)。则存在唯一零均值类函数
\[
\psi_{\ast,u}\in \mathrm{Cl}(G)
\]
使得对所有共轭不变集合 \(A\subseteq G\) 都有
\[
\boxed{\;
c_A=\langle 1_A,\psi_{\ast,u}\rangle_{\mathrm{Cl}(G)}.\;}
\]
对每个 \(\theta\in\RR\)，再定义共轭类并
\[
A_\theta
:=
\bigcup_{\substack{C\in \mathrm{Conj}(G)\\
\Re(e^{-i\theta}\psi_{\ast,u}(C))>0}} C.
\]
则最优见证振幅满足精确变分公式
\[
\boxed{\;
\sup_{A\subseteq G\ \mathrm{class\ union}}|c_A|
=
\frac12
\sup_{\theta\in\RR}
\sum_{C\in \mathrm{Conj}(G)}
\frac{|C|}{|G|}
\left|\Re\!\bigl(e^{-i\theta}\psi_{\ast,u}(C)\bigr)\right|.\;}
\]
并且每个达到上确界的共轭不变集合都可取为某个 \(A_\theta\)，在满足
\[
\Re(e^{-i\theta}\psi_{\ast,u}(C))=0
\]
的零层共轭类上可任意取舍。
\end{theorem}
\begin{proof}
定理 \ref{thm:kernel-chebotarev-second-main-term-witness} 的证明已经表明：\(c_A\) 来自 \(1_A\) 在固定通道 \(\rho_\ast\) 上的 Fourier 分量与固定左右本征向量（等价地，谱投影）的线性配对。因此映射
\[
1_A\longmapsto c_A
\]
在类函数空间 \(\mathrm{Cl}(G)\) 上线性延拓为一个连续复线性泛函 \(L_{\ast,u}\)。由于 \(\mathrm{Cl}(G)\) 是有限维 Hilbert 空间，存在唯一 \(\psi_{\ast,u}\in \mathrm{Cl}(G)\) 使
\[
L_{\ast,u}(f)=\langle f,\psi_{\ast,u}\rangle_{\mathrm{Cl}(G)}
\qquad(\forall\,f\in \mathrm{Cl}(G)).
\]
特别地，对 \(f=1_A\) 即得所示表示公式。

又因当 \(A=G\) 时，\(1_A\) 只有平凡通道分量，而 \(c_G\) 是非平凡极大通道的二主项系数，故 \(c_G=0\)。于是
\[
\langle 1,\psi_{\ast,u}\rangle_{\mathrm{Cl}(G)}=0,
\]
即 \(\psi_{\ast,u}\) 为零均值类函数。

对任意固定 \(\theta\in\RR\)，定义实值零均值类函数
\[
f_\theta:=\Re(e^{-i\theta}\psi_{\ast,u}).
\]
则
\[
|c_A|
=
\sup_{\theta\in\RR}\Re(e^{-i\theta}c_A)
=
\sup_{\theta\in\RR}\langle 1_A,f_\theta\rangle_{\mathrm{Cl}(G)}.
\]
当 \(\theta\) 固定时，\(1_A\) 在每个共轭类 \(C\) 上只能取 \(0\) 或 \(1\)。因此最大化
\[
\langle 1_A,f_\theta\rangle_{\mathrm{Cl}(G)}
=
\sum_{C\in \mathrm{Conj}(G)}\frac{|C|}{|G|}\,1_A(C)\,f_\theta(C)
\]
的最优策略，正是选取所有满足 \(f_\theta(C)>0\) 的共轭类并；这恰给出 \(A_\theta\)。对应最大值为
\[
\sum_{f_\theta(C)>0}\frac{|C|}{|G|}f_\theta(C).
\]
又因 \(f_\theta\) 零均值，
\[
\sum_{C}\frac{|C|}{|G|}f_\theta(C)=0,
\]
故正部分与负部分的总质量相等，从而
\[
\sum_{f_\theta(C)>0}\frac{|C|}{|G|}f_\theta(C)
=
\frac12
\sum_{C\in \mathrm{Conj}(G)}\frac{|C|}{|G|}\,|f_\theta(C)|.
\]
对 \(\theta\) 再取上确界，即得所示变分公式。最后，在满足 \(f_\theta(C)=0\) 的共轭类上，无论取入或取出都不改变目标值，因此一切极值集都可表示为某个 \(A_\theta\) 加上若干零层共轭类的任意选择。证毕。
\end{proof}

\paragraph{primitive Chebotarev 类分布的 Haar 收敛率、Plancherel 偏差指数与 Artin cofactor 半径}
在非阿贝尔 Chebotarev--Plancherel 能量恒等式、正规覆盖算子的 Perron 主相刚性、二主项见证集合以及最小类函数商都已建立之后，还可把 primitive Chebotarev / Peter--Weyl 通道化链进一步压缩为下列四条直接后果。它们表明：一旦把类分布偏差写成 centered class function，平凡通道便会在归一化后完全消失；于是真正控制统计收敛率、平方能量与去主因子 cofactor 最近零点半径的，恰是非平凡通道主谱半径 \(\eta(u)\)。

\begin{theorem}[primitive Chebotarev 类分布到 Haar 分布的精确指数收敛率]\label{thm:xi-time-part64ab-primitive-chebotarev-haar-tv-exponent}
沿用定理 \ref{thm:xi-time-part64-cover-perron-root-trivial-channel}、定理
\ref{thm:kernel-weighted-prime-orbit}、定理
\ref{thm:kernel-chebotarev-second-main-term-witness} 与定理
\ref{thm:xi-time-part64aa-optimal-chebotarev-witness-sign-slice} 的记号。设
\[
\lambda_1(u):=\rho(M_{K,\mathbf 1}(u)),
\qquad
\eta(u):=\max_{\rho\neq \mathbf 1}\rho(M_{K,\rho}(u)).
\]
对每个共轭类 \(C\in \mathrm{Conj}(G)\)，记 primitive Chebotarev 计数为
\[
B_{n,C}(u),
\]
并记总计数为
\[
B_n(u):=\sum_{C\in \mathrm{Conj}(G)}B_{n,C}(u).
\]
定义共轭类分布
\[
p_{n,u}(C):=\frac{B_{n,C}(u)}{B_n(u)},
\qquad
\mu_{\mathrm{Haar}}(C):=\frac{|C|}{|G|},
\]
以及总变差距离
\[
\|p_{n,u}-\mu_{\mathrm{Haar}}\|_{\mathrm{TV}}
:=
\sup_{A\subseteq G\ \mathrm{class\ union}}
\bigl|p_{n,u}(A)-\mu_{\mathrm{Haar}}(A)\bigr|.
\]
则有
\[
\boxed{\;
\|p_{n,u}-\mu_{\mathrm{Haar}}\|_{\mathrm{TV}}
=
O\!\left(\Bigl(\frac{\eta(u)}{\lambda_1(u)}\Bigr)^n\right).\;}
\]

若进一步存在唯一非平凡通道 \(\rho_\ast\neq \mathbf 1\) 与本征值 \(\mu_\ast\in\Spec(M_{K,\rho_\ast}(u))\) 满足
\[
|\mu_\ast|=\eta(u),
\]
且该通道在 \(\mu_\ast\) 处具有谱隙，则
\[
\boxed{\;
\limsup_{n\to\infty}
\|p_{n,u}-\mu_{\mathrm{Haar}}\|_{\mathrm{TV}}^{1/n}
=
\frac{\eta(u)}{\lambda_1(u)}.\;}
\]
\end{theorem}
\begin{proof}
对任意共轭类并 \(A\subseteq G\)，记
\[
B_{n,A}(u):=\sum_{C\subseteq A}B_{n,C}(u),
\qquad
\varepsilon_n(A;u):=
B_{n,A}(u)-\mu_{\mathrm{Haar}}(A)\,B_n(u).
\]
如定理 \ref{thm:kernel-chebotarev-second-main-term-witness} 的证明所示，\(B_{n,A}(u)\) 可按 Peter--Weyl 通道分解为平凡通道与全部非平凡通道的线性组合；其中平凡通道部分恰为
\[
\mu_{\mathrm{Haar}}(A)\,B_n(u).
\]
因此 \(\varepsilon_n(A;u)\) 只由非平凡通道贡献组成。对每个非平凡 \(\rho\neq \mathbf 1\)，标准谱半径估计给出相应通道 primitive 系数为
\[
O\!\left(\frac{\eta(u)^n}{n}\right),
\]
故
\[
\varepsilon_n(A;u)=O\!\left(\frac{\eta(u)^n}{n}\right).
\]
由于共轭类并的总数有限，上述估计可在 \(A\) 上统一。

另一方面，由定理 \ref{thm:kernel-weighted-prime-orbit}，
\[
B_n(u)=\frac{\lambda_1(u)^n}{n}+O\!\left(\frac{|\lambda_2(u)|^n}{n}\right)
\sim \frac{\lambda_1(u)^n}{n}.
\]
于是对任意共轭类并 \(A\) 都有
\[
\bigl|p_{n,u}(A)-\mu_{\mathrm{Haar}}(A)\bigr|
=
\frac{|\varepsilon_n(A;u)|}{B_n(u)}
=
O\!\left(\Bigl(\frac{\eta(u)}{\lambda_1(u)}\Bigr)^n\right).
\]
对 \(A\) 取上确界，即得总变差上界。

若唯一主导通道 \(\rho_\ast\) 存在，则由定理
\ref{thm:kernel-chebotarev-second-main-term-witness} 与定理
\ref{thm:xi-time-part64aa-optimal-chebotarev-witness-sign-slice}，可取某个共轭类并
\[
A_\ast\subseteq G
\]
使得对应二主项系数 \(c_{A_\ast}\neq 0\)。于是沿某个无穷子序列有
\[
|\varepsilon_n(A_\ast;u)|
\ge
c\,\frac{\eta(u)^n}{n}
\]
成立，其中 \(c>0\)。再除以
\[
B_n(u)\sim \frac{\lambda_1(u)^n}{n}
\]
得到
\[
\bigl|p_{n,u}(A_\ast)-\mu_{\mathrm{Haar}}(A_\ast)\bigr|
\ge
c'\Bigl(\frac{\eta(u)}{\lambda_1(u)}\Bigr)^n
\]
在无穷多个 \(n\) 上成立。由总变差的定义，
\[
\|p_{n,u}-\mu_{\mathrm{Haar}}\|_{\mathrm{TV}}
\ge
\bigl|p_{n,u}(A_\ast)-\mu_{\mathrm{Haar}}(A_\ast)\bigr|,
\]
故得到相反方向的 \(\limsup\) 下界。与前面的上界合并，即得结论。证毕。
\end{proof}

\begin{theorem}[Chebotarev--Plancherel 偏差的平方能量具有精确指数率]\label{thm:xi-time-part64ab-chebotarev-plancherel-l2-exponent}
沿用定理 \ref{thm:xi-time-part64ab-primitive-chebotarev-haar-tv-exponent} 的记号。定义 centered class-count 偏差
\[
\varepsilon_n(C;u):=
B_{n,C}(u)-\mu_{\mathrm{Haar}}(C)\,B_n(u),
\]
并定义其按 Haar 权归一化的平方能量
\[
\mathcal E_n(u)
:=
\sum_{C\in\mathrm{Conj}(G)}
\frac{|\varepsilon_n(C;u)|^2}{\mu_{\mathrm{Haar}}(C)}.
\]
则有
\[
\boxed{\;
\mathcal E_n(u)
=
O\!\left(\frac{\eta(u)^{2n}}{n^2}\right).\;}
\]
若再处于定理 \ref{thm:xi-time-part64ab-primitive-chebotarev-haar-tv-exponent} 的唯一主导通道设定下，则
\[
\boxed{\;
\limsup_{n\to\infty}\mathcal E_n(u)^{1/n}
=
\eta(u)^2.\;}
\]
\end{theorem}
\begin{proof}
由定义，
\[
\varepsilon_n(C;u)=B_n(u)\bigl(p_{n,u}(C)-\mu_{\mathrm{Haar}}(C)\bigr),
\]
故
\[
\mathcal E_n(u)
=
B_n(u)^2
\sum_{C\in\mathrm{Conj}(G)}
\frac{|p_{n,u}(C)-\mu_{\mathrm{Haar}}(C)|^2}{\mu_{\mathrm{Haar}}(C)}.
\]
由于共轭类个数有限，而定理
\ref{thm:xi-time-part64ab-primitive-chebotarev-haar-tv-exponent} 已给出逐类一致估计
\[
p_{n,u}(C)-\mu_{\mathrm{Haar}}(C)
=
O\!\left(\Bigl(\frac{\eta(u)}{\lambda_1(u)}\Bigr)^n\right),
\]
遂得
\[
\sum_{C\in\mathrm{Conj}(G)}
\frac{|p_{n,u}(C)-\mu_{\mathrm{Haar}}(C)|^2}{\mu_{\mathrm{Haar}}(C)}
=
O\!\left(\Bigl(\frac{\eta(u)}{\lambda_1(u)}\Bigr)^{2n}\right).
\]
再由
\[
B_n(u)\sim \frac{\lambda_1(u)^n}{n},
\]
即得
\[
\mathcal E_n(u)
=
O\!\left(\frac{\eta(u)^{2n}}{n^2}\right).
\]

在唯一主导通道设定下，利用加权 Cauchy--Schwarz 不等式
\[
\left(\sum_{C}|x_C|\right)^2
\le
\left(\sum_C \frac{|x_C|^2}{\mu_{\mathrm{Haar}}(C)}\right)
\left(\sum_C \mu_{\mathrm{Haar}}(C)\right),
\]
并取
\[
x_C:=p_{n,u}(C)-\mu_{\mathrm{Haar}}(C),
\]
便得到
\[
4\|p_{n,u}-\mu_{\mathrm{Haar}}\|_{\mathrm{TV}}^2
\le
\sum_{C}
\frac{|p_{n,u}(C)-\mu_{\mathrm{Haar}}(C)|^2}{\mu_{\mathrm{Haar}}(C)}.
\]
因此
\[
\mathcal E_n(u)
\ge
4B_n(u)^2\,\|p_{n,u}-\mu_{\mathrm{Haar}}\|_{\mathrm{TV}}^2.
\]
再代入
\[
B_n(u)\sim \frac{\lambda_1(u)^n}{n}
\]
与定理 \ref{thm:xi-time-part64ab-primitive-chebotarev-haar-tv-exponent} 的精确 \(\limsup\) 公式，便得
\[
\limsup_{n\to\infty}\mathcal E_n(u)^{1/n}\ge \eta(u)^2.
\]
与前面上界给出的
\[
\limsup_{n\to\infty}\mathcal E_n(u)^{1/n}\le \eta(u)^2
\]
合并，即得所示等式。证毕。
\end{proof}

\begin{theorem}[Artin cofactor 行列式的最近零点半径精确等于非平凡通道主谱倒数]\label{thm:xi-time-part64ab-artin-cofactor-nearest-zero-radius}
沿用命题 \ref{prop:kernel-artin-factorization} 与定理 \ref{thm:xi-time-part64-cover-perron-root-trivial-channel} 的记号。定义去主因子 cofactor 行列式
\[
\Delta_G^{\mathrm{cof}}(z,u)
:=
\frac{\det(I-zM_{K,G}(u))}{\det(I-zM_{K,\mathbf 1}(u))}
=
\prod_{\rho\neq \mathbf 1}\det(I-zM_{K,\rho}(u))^{d_\rho}.
\]
记其离原点最近的零点半径为
\[
R_{\mathrm{cof}}(u)
:=
\min\bigl\{|z|:\ \Delta_G^{\mathrm{cof}}(z,u)=0\bigr\}.
\]
则有
\[
\boxed{\;
R_{\mathrm{cof}}(u)=\frac{1}{\eta(u)}.\;}
\]
特别地，若满足核版 GRH 型平方根界
\[
\eta(u)\le \lambda_1(u)^{1/2},
\]
则
\[
\boxed{\;
R_{\mathrm{cof}}(u)\ge \lambda_1(u)^{-1/2}.\;}
\]
\end{theorem}
\begin{proof}
由命题 \ref{prop:kernel-artin-factorization}，
\[
\Delta_G^{\mathrm{cof}}(z,u)
=
\prod_{\rho\neq \mathbf 1}\det(I-zM_{K,\rho}(u))^{d_\rho}.
\]
因此它的全部零点恰由所有非平凡通道矩阵 \(M_{K,\rho}(u)\) 的特征值倒数构成。故最近零点半径满足
\[
R_{\mathrm{cof}}(u)
=
\min_{\rho\neq \mathbf 1}\ \min_{\lambda\in\Spec(M_{K,\rho}(u))}
\frac{1}{|\lambda|}
=
\frac{1}{\max_{\rho\neq \mathbf 1}\rho(M_{K,\rho}(u))}
=
\frac{1}{\eta(u)}.
\]
这就证明了第一式。

若再有
\[
\eta(u)\le \lambda_1(u)^{1/2},
\]
则取倒数即得
\[
R_{\mathrm{cof}}(u)\ge \lambda_1(u)^{-1/2}.
\]
证毕。
\end{proof}

\begin{theorem}[群商不会恶化 Chebotarev 误差指数率]\label{thm:xi-time-part64ab-quotient-monotonicity-error-exponent}
设
\[
\varphi:G\twoheadrightarrow H
\]
为有限群满同态。把核 \(K\) 的边标签经 \(\varphi\) 推送后得到 \(H\)-标号商核，记其非平凡通道主谱半径为
\[
\eta_G(u),\qquad \eta_H(u),
\]
而平凡通道 Perron 根统一记为
\[
\lambda_1(u).
\]
则有
\[
\boxed{\;
\eta_H(u)\le \eta_G(u).\;}
\]
因而同时成立
\[
\boxed{\;
\frac{\eta_H(u)}{\lambda_1(u)}
\le
\frac{\eta_G(u)}{\lambda_1(u)},
\qquad
R_{\mathrm{cof},H}(u)\ge R_{\mathrm{cof},G}(u).\;}
\]
换言之，群商不会恶化 primitive Chebotarev 类分布向 Haar 分布的指数收敛率，也不会缩小去主因子 cofactor 行列式的最近零点半径。
\end{theorem}
\begin{proof}
任取 \(H\) 的一个非平凡不可约表示
\[
\sigma\in\widehat H\setminus\{\mathbf 1\}.
\]
其 inflation
\[
\widetilde\sigma:=\sigma\circ\varphi
\]
是 \(G\) 的非平凡不可约表示。由通道矩阵的定义，
\[
M_{K_H,\sigma}(u)=M_{K,\widetilde\sigma}(u).
\]
因此
\[
\rho(M_{K_H,\sigma}(u))
=
\rho(M_{K,\widetilde\sigma}(u))
\le
\eta_G(u).
\]
对所有非平凡 \(\sigma\) 取上确界，便得
\[
\eta_H(u)\le \eta_G(u).
\]

另一方面，平凡表示在推送下保持平凡，故
\[
M_{K_H,\mathbf 1}(u)=M_{K,\mathbf 1}(u),
\]
从而平凡通道 Perron 根不变，即两侧都等于 \(\lambda_1(u)\)。于是立得
\[
\frac{\eta_H(u)}{\lambda_1(u)}
\le
\frac{\eta_G(u)}{\lambda_1(u)}.
\]
最后，由定理 \ref{thm:xi-time-part64ab-artin-cofactor-nearest-zero-radius}
\[
R_{\mathrm{cof}}(u)=\frac{1}{\eta(u)},
\]
可知
\[
R_{\mathrm{cof},H}(u)=\frac{1}{\eta_H(u)}
\ge
\frac{1}{\eta_G(u)}
=
R_{\mathrm{cof},G}(u).
\]
证毕。
\end{proof}

\noindent
由此，Chebotarev--Plancherel / Peter--Weyl 通道化链在 primitive 层进一步闭合：归一化共轭类分布向 Haar 分布的精确指数率不再由平凡通道内部的次谱所控制，而是完全由非平凡通道主谱半径 \(\eta(u)\) 决定；相应的平方偏差能量则以 \(\eta(u)^2\) 为精确指数率衰减；在解析侧，Artin 因子化中的去主因子 cofactor 最近零点半径也正好等于 \(\eta(u)^{-1}\)；而在函子侧，群商只会压低 \(\eta(u)\)，因而只能改进而不会恶化这些统计与解析指数。于是，primitive Chebotarev 收敛、Plancherel 偏差、Artin cofactor 半径与 quotient monotonicity 便被压缩为同一个非平凡通道主谱参数。

\paragraph{Chebotarev 商熵、可见角色能量与平方根统计障碍}
在 primitive 同余层的严格 Fourier 反演、subgroup--Möbius 完全反演、商映射下的精确能量损失、非阿贝尔通道能量账本、最小偏差承载商以及主导通道见证集合都已建立之后，还可把群商粗化、Plancherel 通道分裂、次主奇点与 RH/GRH 型统计后果进一步压缩为下列四条结构后果。它们分别把商映射的 \(L^2\) 收缩提升为相对熵链式律，把可见偏差写成可见/隐藏角色能量的精确正交分解，把主导扭曲通道的次主解析障碍严格识别为表示维数加权极点，并把平方根谱隙条件转译为 TV/\(\chi^2\)/KL 的统一指数律及其见证性反证。

\begin{theorem}[群商下 Chebotarev 分布相对均匀基线的 KL 链式律]\label{thm:xi-time-part64ac-abelian-quotient-kl-chain}
沿用推论 \ref{cor:kernel-chebotarev-L2-monotone} 的记号。设
\[
\pi:G_2\twoheadrightarrow G_1
\]
为有限阿贝尔群满同态，并记 \(K:=\ker\pi\)。对固定 \(u\in(0,1]\) 与 \(n\ge 1\)，定义两层 primitive Chebotarev 概率分布
\[
\mu_{n,u}^{(i)}(c_i):=\frac{B_{n,c_i}^{(i)}(u)}{B_n(u)},
\qquad
B_n(u):=B_n^{(1)}(\mathbf{1};u)=B_n^{(2)}(\mathbf{1};u),
\qquad
i\in\{1,2\},
\]
并记 \(u_{G_i}\) 为 \(G_i\) 上的均匀分布。对每个满足 \(\mu_{n,u}^{(1)}(c_1)>0\) 的 \(c_1\in G_1\)，定义纤维条件分布
\[
\mu_{n,u}^{(c_1),\mathrm{fib}}(c_2)
:=
\frac{\mu_{n,u}^{(2)}(c_2)}{\mu_{n,u}^{(1)}(c_1)},
\qquad
c_2\in \pi^{-1}(c_1),
\]
并记 \(u_{\pi^{-1}(c_1)}\) 为该纤维上的均匀分布。则有严格链式恒等式
\[
\boxed{\;
D_{\mathrm{KL}}\!\bigl(\mu_{n,u}^{(2)}\Vert u_{G_2}\bigr)
=
D_{\mathrm{KL}}\!\bigl(\mu_{n,u}^{(1)}\Vert u_{G_1}\bigr)
+
\sum_{c_1\in G_1}\mu_{n,u}^{(1)}(c_1)\,
D_{\mathrm{KL}}\!\bigl(\mu_{n,u}^{(c_1),\mathrm{fib}}\Vert u_{\pi^{-1}(c_1)}\bigr)\; }.
\]
其中对 \(\mu_{n,u}^{(1)}(c_1)=0\) 的纤维约定对应项为 \(0\)。特别地，商映射所丢失的全部相对熵恰等于纤维条件非均匀性的平均 KL 成本。
\leanverified{paper\_xi\_time\_part64ac\_abelian\_quotient\_kl\_chain}
\end{theorem}
\begin{proof}
由于 \(\mu_{n,u}^{(1)}=\pi_\ast\mu_{n,u}^{(2)}\) 且 \(\pi\) 为群满同态，每条纤维大小都等于 \(|K|\)，故
\[
u_{G_2}(c_2)
=
\frac{1}{|G_2|}
=
\frac{1}{|G_1|\,|K|}
=
u_{G_1}(\pi(c_2))\,u_{\pi^{-1}(\pi(c_2))}(c_2).
\]
于是对任意满足 \(\mu_{n,u}^{(2)}(c_2)>0\) 的 \(c_2\in G_2\)，有
\[
\log\frac{\mu_{n,u}^{(2)}(c_2)}{u_{G_2}(c_2)}
=
\log\frac{\mu_{n,u}^{(1)}(\pi(c_2))}{u_{G_1}(\pi(c_2))}
+
\log\frac{\mu_{n,u}^{(\pi(c_2)),\mathrm{fib}}(c_2)}{u_{\pi^{-1}(\pi(c_2))}(c_2)}.
\]
对 \(\mu_{n,u}^{(2)}\) 取期望并按纤维分组，即得所示分解；这也是引理 \ref{lem:pom-fiberwise-kl} 在等势群商情形下的直接专门化。证毕。
\end{proof}

\begin{theorem}[可见/隐藏角色能量的精确正交分解]\label{thm:xi-time-part64ac-visible-hidden-character-energy-splitting}
沿用定理 \ref{thm:xi-time-part64ac-abelian-quotient-kl-chain} 的设定，并记
\[
r_n^{(i)}(c_i;u):=|G_i|\,\mu_{n,u}^{(i)}(c_i)
=
\frac{|G_i|\,B_{n,c_i}^{(i)}(u)}{B_n(u)},
\qquad i\in\{1,2\}.
\]
把 \(\widehat{G_1}\) 通过 \(\pi^\ast\) 识别为
\[
K^\perp:=\{\chi\in \widehat{G_2}:\ \chi|_K\equiv 1\}\subseteq \widehat{G_2},
\]
并记细层角色幅度
\[
B_n^{(2)}(\chi;u):=\sum_{c_2\in G_2}\chi(c_2)\,B_{n,c_2}^{(2)}(u).
\]
则有两组精确分解：
\[
\boxed{\;
\bigl\|r_n^{(2)}(\cdot;u)-1\bigr\|_{L^2(G_2)}^2
=
\bigl\|r_n^{(1)}(\cdot;u)-1\bigr\|_{L^2(G_1)}^2
+
\frac{1}{|G_2|}
\sum_{c_1\in G_1}\sum_{c_2\in\pi^{-1}(c_1)}
\bigl|r_n^{(2)}(c_2;u)-r_n^{(1)}(c_1;u)\bigr|^2\; }.
\]
\[
\boxed{\;
\bigl\|r_n^{(1)}(\cdot;u)-1\bigr\|_{L^2(G_1)}^2
=
\frac{1}{B_n(u)^2}
\sum_{\chi\in K^\perp\setminus\{\mathbf{1}\}}
\bigl|B_n^{(2)}(\chi;u)\bigr|^2\; }.
\]
\[
\boxed{\;
\bigl\|r_n^{(2)}(\cdot;u)-1\bigr\|_{L^2(G_2)}^2
-
\bigl\|r_n^{(1)}(\cdot;u)-1\bigr\|_{L^2(G_1)}^2
=
\frac{1}{B_n(u)^2}
\sum_{\chi\in \widehat{G_2}\setminus K^\perp}
\bigl|B_n^{(2)}(\chi;u)\bigr|^2\; }.
\]
因而，群商仅保留在 \(K\) 上平凡的可见角色通道，并精确删除全部隐藏通道能量。
\end{theorem}
\begin{proof}
由推论 \ref{cor:kernel-chebotarev-L2-monotone} 的证明，
\[
r_n^{(1)}(\cdot;u)\circ\pi
=
\mathbb{E}_\pi\!\bigl[r_n^{(2)}(\cdot;u)\bigr]
\]
正是 \(r_n^{(2)}(\cdot;u)\) 沿纤维常值子空间的条件期望。故在 \(L^2(G_2)\) 中可应用正交投影的勾股恒等式：
\[
\bigl\|r_n^{(2)}-1\bigr\|_2^2
=
\bigl\|r_n^{(1)}\circ\pi-1\bigr\|_2^2
+
\bigl\|r_n^{(2)}-r_n^{(1)}\circ\pi\bigr\|_2^2.
\]
按纤维展开第二项，即得第一式。

另一方面，推论 \ref{cor:kernel-chebotarev-plancherel-energy} 应用于细层给出
\[
\bigl\|r_n^{(2)}(\cdot;u)-1\bigr\|_{L^2(G_2)}^2
=
\frac{1}{B_n(u)^2}
\sum_{\chi\in \widehat{G_2}\setminus\{\mathbf{1}\}}
\bigl|B_n^{(2)}(\chi;u)\bigr|^2.
\]
而粗层角色正好对应于 \(K^\perp=\pi^\ast(\widehat{G_1})\)，故同理
\[
\bigl\|r_n^{(1)}(\cdot;u)-1\bigr\|_{L^2(G_1)}^2
=
\frac{1}{B_n(u)^2}
\sum_{\chi\in K^\perp\setminus\{\mathbf{1}\}}
\bigl|B_n^{(2)}(\chi;u)\bigr|^2,
\]
得到第二式。将两式相减即得第三式；这也是定理
\ref{thm:xi-time-part63b-primitive-quotient-energy-loss-exact}
在相对 Haar 密度口径下的直接重写。证毕。
\end{proof}

\begin{theorem}[主导非平凡通道的极点阶与见证振荡]\label{thm:xi-time-part64ac-dominant-channel-pole-order-witness-oscillation}
沿用定理 \ref{thm:xi-time-part64-cover-zeta-pole-order-weighted-channel-multiplicity} 与定理 \ref{thm:kernel-chebotarev-second-main-term-witness} 的记号。固定 \(u\in(0,1]\)，并假设存在唯一非平凡通道 \(\rho_\star\neq \mathbf{1}\) 与其简单本征值 \(\mu_\star\in\Spec(M_{K,\rho_\star}(u))\) 满足
\[
|\mu_\star|
=
\max_{\rho\neq \mathbf{1}}\rho\!\bigl(M_{K,\rho}(u)\bigr)
=:\eta(u),
\]
且该通道在 \(\mu_\star\) 处具有谱隙。则有：
\begin{enumerate}
  \item 去主因子覆盖 \(\zeta\) 在 \(z=\mu_\star^{-1}\) 处具有精确极点阶
  \[
  \boxed{\;
  \operatorname{ord}_{z=\mu_\star^{-1}}
  \frac{\zeta_{K,G}(z,u)}{L_{K,\mathbf{1}}(z,u)}
  =
  \dim\rho_\star\; }.
  \]
  \item 存在某个常数 \(\eta_-(u)<\eta(u)\)，使得对任意共轭不变集合 \(A\subseteq G\)，若其指示函数 \(1_A\) 在 \(\rho_\star\)-通道上的 Fourier 分量非零，则记
  \[
  B_{n,A}(u):=\sum_{C\subseteq A}B_{n,C}(u),
  \qquad
  \mu_{\mathrm{Haar}}(A):=\frac{|A|}{|G|},
  \]
  并存在 \(c_A\neq 0\) 使
  \[
  \boxed{\;
  B_{n,A}(u)
  =
  \mu_{\mathrm{Haar}}(A)\frac{\lambda_1(u)^n}{n}
  +
  \Re\!\bigl(c_A\mu_\star^n\bigr)\frac1n
  +
  O\!\left(\frac{\eta_-(u)^n}{n}\right)\; }.
  \]
  特别地，存在常数 \(c(A)>0\) 使
  \[
  \boxed{\;
  \left|B_{n,A}(u)-\mu_{\mathrm{Haar}}(A)\frac{\lambda_1(u)^n}{n}\right|
  \ge
  c(A)\frac{\eta(u)^n}{n}\; }
  \]
  在无穷多个 \(n\) 上成立。
\end{enumerate}
\end{theorem}
\begin{proof}
由 Artin 因子化
\[
\frac{\zeta_{K,G}(z,u)}{L_{K,\mathbf{1}}(z,u)}
=
\prod_{\rho\neq \mathbf{1}}L_{K,\rho}(z,u)^{\dim\rho}
\]
可知，该商在 \(z=\mu_\star^{-1}\) 处的极点阶只由非平凡通道决定。再由定理 \ref{thm:xi-time-part64-cover-zeta-pole-order-weighted-channel-multiplicity}，
\[
\operatorname{ord}_{z=\mu_\star^{-1}}
\frac{\zeta_{K,G}(z,u)}{L_{K,\mathbf{1}}(z,u)}
=
\sum_{\rho\in\widehat G\setminus\{\mathbf{1}\}}(\dim\rho)\,m_\rho(\mu_\star),
\]
其中 \(m_\rho(\mu_\star)\) 为 \(\mu_\star\) 在 \(M_{K,\rho}(u)\) 中的代数重数。由唯一主导通道假设与 \(\mu_\star\) 的简单性可知，\(\rho_\star\) 贡献 \(m_{\rho_\star}(\mu_\star)=1\)，而其余非平凡通道都不含该特征值，故极点阶等于 \(\dim\rho_\star\)。

第二部分正是定理 \ref{thm:kernel-chebotarev-second-main-term-witness} 的直接应用：唯一主导通道给出单一的次主振荡模态，谱隙使所有其余通道并入某个 \(\eta_-(u)<\eta(u)\) 的统一误差项；若 \(1_A\) 在 \(\rho_\star\)-通道上的分量非零，则对应系数 \(c_A\neq 0\)，从而沿无穷多个 \(n\) 得到 \(\Omega(\eta(u)^n/n)\) 的见证振荡下界。证毕。
\end{proof}

\begin{theorem}[GRH 型平方根谱隙的 TV/\texorpdfstring{$\chi^2$}{chi2}/KL 统计后果]\label{thm:xi-time-part64ac-grh-tv-chi2-kl-barrier}
沿用定理 \ref{thm:xi-time-part64ab-primitive-chebotarev-haar-tv-exponent} 的设定，并进一步假设 \(G\) 为有限阿贝尔群。定义
\[
\mu_{n,u}(c):=\frac{B_{n,c}(u)}{B_n(u)},
\qquad
u_G(c):=\frac{1}{|G|},
\qquad
r_n(c;u):=|G|\,\mu_{n,u}(c),
\]
以及
\[
\eta(u):=\max_{\chi\neq \mathbf{1}}\rho\!\bigl(M_{K,\chi}(u)\bigr).
\]
则有统一指数界
\[
\boxed{\;
\bigl\|\mu_{n,u}-u_G\bigr\|_{\mathrm{TV}}
=
O\!\left(\Bigl(\frac{\eta(u)}{\lambda_1(u)}\Bigr)^n\right),\;}
\]
\[
\boxed{\;
\chi^2\!\bigl(\mu_{n,u}\Vert u_G\bigr)
=
O\!\left(\Bigl(\frac{\eta(u)}{\lambda_1(u)}\Bigr)^{2n}\right),\qquad
D_{\mathrm{KL}}\!\bigl(\mu_{n,u}\Vert u_G\bigr)
=
O\!\left(\Bigl(\frac{\eta(u)}{\lambda_1(u)}\Bigr)^{2n}\right).\;}
\]
特别地，若 \(\textsf{GRH}_{K,G}(u)\) 成立，即
\[
\eta(u)\le \lambda_1(u)^{1/2},
\]
则
\[
\boxed{\;
\bigl\|\mu_{n,u}-u_G\bigr\|_{\mathrm{TV}}=O\!\bigl(\lambda_1(u)^{-n/2}\bigr),\qquad
\chi^2\!\bigl(\mu_{n,u}\Vert u_G\bigr),\
D_{\mathrm{KL}}\!\bigl(\mu_{n,u}\Vert u_G\bigr)
=
O\!\bigl(\lambda_1(u)^{-n}\bigr).\;}
\]
反之，若定理 \ref{thm:xi-time-part64ac-dominant-channel-pole-order-witness-oscillation} 的唯一主导通道假设成立，且
\[
\eta(u)>\lambda_1(u)^{1/2},
\]
则存在某个集合 \(A\subseteq G\) 与常数 \(c(A)>0\)，使
\[
\boxed{\;
\bigl|\mu_{n,u}(A)-u_G(A)\bigr|
\ge
c(A)\Bigl(\frac{\eta(u)}{\lambda_1(u)}\Bigr)^n\; }
\]
在无穷多个 \(n\) 上成立。因而此时不可能存在统一的平方根级 TV 收敛律。
\end{theorem}
\begin{proof}
阿贝尔情形下，共轭类分布即元素分布，故定理 \ref{thm:xi-time-part64ab-primitive-chebotarev-haar-tv-exponent} 直接给出 TV 上界。另一方面，
\[
\chi^2\!\bigl(\mu_{n,u}\Vert u_G\bigr)
=
\frac{1}{|G|}\sum_{c\in G}\bigl(r_n(c;u)-1\bigr)^2.
\]
由推论 \ref{cor:kernel-chebotarev-plancherel-energy} 的相对 Haar 密度版本，
\[
\chi^2\!\bigl(\mu_{n,u}\Vert u_G\bigr)
=
\frac{1}{B_n(u)^2}
\sum_{\chi\in\widehat G\setminus\{\mathbf{1}\}}
\bigl|B_n(\chi;u)\bigr|^2.
\]
对每个非平凡角色 \(\chi\neq \mathbf{1}\)，由命题 \ref{prop:kernel-chebotarev-mobius-adams} 与谱半径估计有
\[
B_n(\chi;u)=O\!\left(\frac{\eta(u)^n}{n}\right),
\]
而 \(B_n(u)\sim \lambda_1(u)^n/n\) 由定理 \ref{thm:kernel-weighted-prime-orbit} 给出。由于非平凡角色数有限，遂得
\[
\chi^2\!\bigl(\mu_{n,u}\Vert u_G\bigr)
=
O\!\left(\Bigl(\frac{\eta(u)}{\lambda_1(u)}\Bigr)^{2n}\right).
\]
再由标准不等式
\[
D_{\mathrm{KL}}(P\Vert Q)\le \chi^2(P\Vert Q),
\]
立即得到 KL 上界。若 \(\eta(u)\le \lambda_1(u)^{1/2}\)，则把该平方根界代入上述三式即得所示结论。

反向方向中，若唯一主导通道存在且 \(\eta(u)>\lambda_1(u)^{1/2}\)，则由定理
\ref{thm:xi-time-part64ac-dominant-channel-pole-order-witness-oscillation}
可取某个 \(A\subseteq G\) 使
\[
\left|B_{n,A}(u)-u_G(A)\frac{\lambda_1(u)^n}{n}\right|
\ge
c(A)\frac{\eta(u)^n}{n}
\]
在无穷多个 \(n\) 上成立。再除以
\[
B_n(u)\sim \frac{\lambda_1(u)^n}{n}
\]
即得
\[
\bigl|\mu_{n,u}(A)-u_G(A)\bigr|
\ge
c'(A)\Bigl(\frac{\eta(u)}{\lambda_1(u)}\Bigr)^n
\]
沿无穷多个 \(n\) 成立。由于总变差是事件偏差的上确界，若再有统一的 \(O(\lambda_1(u)^{-n/2})\) 上界，则会与 \(\eta(u)>\lambda_1(u)^{1/2}\) 矛盾。证毕。
\end{proof}

\noindent
因此，primitive Chebotarev 的 quotient pushforward 不仅满足 \(L^2\) 收缩，而且在相对熵层面具有严格链式记账；相应的可见偏差可精确拆分为可见/隐藏角色通道能量；一旦某个非平凡通道主导全部次主谱，它便同时决定去主因子覆盖 \(\zeta\) 的最近奇点阶数与 Chebotarev 见证集合的振荡量级；而 RH/GRH 型平方根谱隙条件则恰把这些通道参数进一步压缩为 TV/\(\chi^2\)/KL 的统一统计指数与尖锐障碍。于是，商粗化、角色 Fourier、Artin 因子化与平方根误差预算便在同一条有限维通道链上闭合。

\paragraph{三域符号盒、最小退化逐点判据与 fold-gauge 的奇次重整化}
在三域 Frobenius 直积的完全张量分解、审计偶数窗最小退化 Fibonacci 律以及 fold-gauge 缺陷序列的 Bernoulli 剥离已经建立之后，还可把条件独立、子覆盖猜想与非线性正规形继续压缩为下列六条后果。前半部分把 \(L_4\) 因子、四条二次特征与全分裂事件收束为严格的符号盒均分与条件独立；后半部分则把最小退化扇区的几何猜想收束为逐点多重度判据，并把 fold-gauge 缺陷的重整化结构提升到五次与纯奇次两个层面。

\begin{theorem}[\texorpdfstring{$L_4$}{L4} 事件对 \texorpdfstring{$(\chi_{10},\chi_{\mathrm{LY}})$}{(chi10, chiLY)} 的完全条件独立]\label{thm:xi-time-part64a-l4-event-complete-conditional-independence}
沿用定理 \ref{thm:xi-time-part62d-triple-frobenius-classfunction-tensor-independence} 的记号。对不分歧素数 \(p\)，记
\[
\sigma_{10}(p)\in S_{10},
\qquad
\sigma_{3}(p)\in S_3,
\qquad
\sigma_{4,\mathrm{tot}}(p):=\bigl(\sigma_4(p),\varepsilon_4(p)\bigr)\in S_4\times C_2,
\]
并定义二次特征
\[
\chi_{10}(p):=\operatorname{sgn}\!\bigl(\sigma_{10}(p)\bigr),
\qquad
\chi_{\mathrm{LY}}(p):=\operatorname{sgn}\!\bigl(\sigma_3(p)\bigr).
\]
任取 \(S_4\times C_2\) 的共轭类并集 \(\mathcal E\)，定义事件
\[
E_{\mathcal E}:=\{\,p:\ \sigma_{4,\mathrm{tot}}(p)\in\mathcal E\,\}.
\]
则对任意 \((\varepsilon_1,\varepsilon_2)\in\{\pm1\}^2\)，有
\[
\delta\!\bigl(
E_{\mathcal E}\cap\{\chi_{10}=\varepsilon_1,\ \chi_{\mathrm{LY}}=\varepsilon_2\}
\bigr)
=
\frac{1}{4}\cdot \frac{|\mathcal E|}{48}.
\]
从而
\[
\delta\!\bigl(
E_{\mathcal E}\mid \chi_{10}=\varepsilon_1,\ \chi_{\mathrm{LY}}=\varepsilon_2
\bigr)
=
\frac{|\mathcal E|}{48}.
\]
特别地，若记
\[
E_{\mathrm{split}}:=\{\text{\(C(X)\bmod p\) 全分裂}\},
\]
则对任意 \((\varepsilon_1,\varepsilon_2)\in\{\pm1\}^2\)，有
\[
\delta\!\bigl(E_{\mathrm{split}}\mid \chi_{10}=\varepsilon_1,\ \chi_{\mathrm{LY}}=\varepsilon_2\bigr)=\frac{1}{48}.
\]
\end{theorem}
\begin{proof}
由定理 \ref{thm:xi-time-part62d-triple-frobenius-classfunction-tensor-independence}，
\[
\Gal(L_{10}L_{\mathrm{LY}}L_4/\QQ)
\cong
S_{10}\times S_3\times (S_4\times C_2),
\]
并且三个因子的 Frobenius 极限分布严格乘积化。

取类函数
\[
f(\tau):=\mathbf 1_{\{\operatorname{sgn}(\tau)=\varepsilon_1\}},
\qquad
g(\upsilon):=\mathbf 1_{\{\operatorname{sgn}(\upsilon)=\varepsilon_2\}},
\qquad
h(\eta):=\mathbf 1_{\mathcal E}(\eta).
\]
由于 \(A_{10}\subset S_{10}\) 与 \(A_3\subset S_3\) 均为指数 \(2\) 子群，故
\[
\langle f\rangle_{S_{10}}=\frac12,
\qquad
\langle g\rangle_{S_3}=\frac12.
\]
而
\[
\langle h\rangle_{S_4\times C_2}=\frac{|\mathcal E|}{48}.
\]
将这三式代入定理 \ref{thm:xi-time-part62d-triple-frobenius-classfunction-tensor-independence} 的乘积公式，便得
\[
\delta\!\bigl(
E_{\mathcal E}\cap\{\chi_{10}=\varepsilon_1,\ \chi_{\mathrm{LY}}=\varepsilon_2\}
\bigr)
=
\frac12\cdot\frac12\cdot\frac{|\mathcal E|}{48}
=
\frac14\cdot\frac{|\mathcal E|}{48}.
\]
再除以
\[
\delta(\chi_{10}=\varepsilon_1,\ \chi_{\mathrm{LY}}=\varepsilon_2)=\frac14
\]
即得条件密度公式。

最后，由定理 \ref{thm:xi-time-part62e-l4-fullsplit-universal-dilution} 的证明，\(C(X)\bmod p\) 全分裂恰对应第三因子的恒等元共轭类，故
\[
|\mathcal E|=1,
\]
于是
\[
\delta\!\bigl(E_{\mathrm{split}}\mid \chi_{10}=\varepsilon_1,\ \chi_{\mathrm{LY}}=\varepsilon_2\bigr)=\frac{1}{48}.
\]
证毕。
\end{proof}

\begin{theorem}[四重二次特征的 \texorpdfstring{$\{\pm1\}^4$}{pm1^4} 等分布]\label{thm:xi-time-part64a-four-quadratic-character-equidistribution}
沿用定理 \ref{thm:xi-time-part62e-triple-product-abelianization-c2-four} 的记号，定义
\[
\chi_{4,\mathrm{sgn}}(p):=\operatorname{sgn}\!\bigl(\sigma_4(p)\bigr),
\qquad
\chi_{4,\mathrm{aux}}(p):=\iota\!\bigl(\varepsilon_4(p)\bigr)\in\{\pm1\},
\]
其中 \(\iota:C_2\xrightarrow{\sim}\{\pm1\}\) 为定理
\ref{thm:xi-time-part62e-triple-product-abelianization-c2-four} 中固定的同构。则四元组
\[
\bigl(
\chi_{10}(p),\chi_{\mathrm{LY}}(p),\chi_{4,\mathrm{sgn}}(p),\chi_{4,\mathrm{aux}}(p)
\bigr)
\]
在 \(\{\pm1\}^4\) 上等分布；亦即对任意
\[
(\varepsilon_1,\varepsilon_2,\varepsilon_3,\varepsilon_4)\in\{\pm1\}^4,
\]
都有
\[
\delta\!\Bigl(
\chi_{10}=\varepsilon_1,\ \chi_{\mathrm{LY}}=\varepsilon_2,\ 
\chi_{4,\mathrm{sgn}}=\varepsilon_3,\ \chi_{4,\mathrm{aux}}=\varepsilon_4
\Bigr)
=
\frac{1}{16}.
\]
\end{theorem}
\begin{proof}
仍由定理 \ref{thm:xi-time-part62d-triple-frobenius-classfunction-tensor-independence}，三个因子的 Frobenius 统计严格乘积化。前两个因子的符号条件各占一半，故
\[
\delta(\chi_{10}=\varepsilon_1,\ \chi_{\mathrm{LY}}=\varepsilon_2)=\frac14.
\]

再看第三因子。定义映射
\[
\Theta:S_4\times C_2\longrightarrow \{\pm1\}\times\{\pm1\},
\qquad
\Theta(\sigma,\epsilon):=\bigl(\operatorname{sgn}(\sigma),\iota(\epsilon)\bigr).
\]
该映射显然满射，且其核为
\[
\ker\Theta=A_4\times\{1\},
\]
大小为 \(12\)。故每个符号盒
\[
\Theta^{-1}(\varepsilon_3,\varepsilon_4)
\]
都恰有 \(12\) 个元素，从而
\[
\delta(\chi_{4,\mathrm{sgn}}=\varepsilon_3,\ \chi_{4,\mathrm{aux}}=\varepsilon_4)
=
\frac{12}{48}
=
\frac14.
\]
于是对全部四重符号盒，
\[
\delta\!\Bigl(
\chi_{10}=\varepsilon_1,\ \chi_{\mathrm{LY}}=\varepsilon_2,\ 
\chi_{4,\mathrm{sgn}}=\varepsilon_3,\ \chi_{4,\mathrm{aux}}=\varepsilon_4
\Bigr)
=
\frac14\cdot\frac14
=
\frac{1}{16}.
\]
证毕。
\end{proof}

\begin{corollary}[三域混合事件的显式密度闭式]\label{cor:xi-time-part64a-triple-mixed-event-explicit-density}
沿用推论 \ref{cor:killo-leyang-joint-density-explicit} 的记号，定义
\[
A:=\{\text{\(P_{10}\bmod p\) 含一次因子}\},
\qquad
B_{\mathrm{root}}:=\{\text{\(g\bmod p\) 至少有一根}\},
\]
并令
\[
C_{\mathrm{split}}:=\{\text{\(C(X)\bmod p\) 全分裂}\}.
\]
则有显式密度
\[
\delta(A\cap B_{\mathrm{root}}\cap C_{\mathrm{split}})
=
\frac{28319}{3225600}.
\]
\end{corollary}
\begin{proof}
由推论 \ref{cor:killo-leyang-joint-density-explicit}，
\[
\delta(A\cap B_{\mathrm{root}})=\frac{28319}{67200}.
\]
另一方面，事件 \(A\cap B_{\mathrm{root}}\) 只依赖 \(S_{10}\times S_3\) 两个因子，而
\(C_{\mathrm{split}}\) 只依赖第三因子。故由定理
\ref{thm:xi-time-part62e-l4-fullsplit-universal-dilution}，
\[
\delta(A\cap B_{\mathrm{root}}\cap C_{\mathrm{split}})
=
\frac{1}{48}\,\delta(A\cap B_{\mathrm{root}}).
\]
代入上式即得
\[
\delta(A\cap B_{\mathrm{root}}\cap C_{\mathrm{split}})
=
\frac{1}{48}\cdot \frac{28319}{67200}
=
\frac{28319}{3225600}.
\]
证毕。
\end{proof}

\begin{proposition}[最小退化扇区猜想的逐点等价判据]\label{prop:xi-time-part64a-minsector-pointwise-equivalence}
沿用定义 \ref{def:fold-bin-degeneracy-sectors} 与猜想
\ref{conj:fold-bin-min-sector-regular-subcover} 的记号。设 \(m\ge 6\) 为偶数，并假设
\[
d_{\min}(m)=F_{m/2},
\qquad
\bigl|\mathcal S_{m,d_{\min}(m)}\bigr|=F_m=|X_{m-2}|.
\]
定义候选嵌入
\[
\iota_{m-2}:X_{m-2}\hookrightarrow X_m,
\qquad
\iota_{m-2}(v):=(v,0,1).
\]
则猜想
\[
\mathcal S_{m,d_{\min}(m)}=\iota_{m-2}(X_{m-2})
\]
与下述逐点陈述完全等价：
\[
d_m^{\mathrm{bin}}(v,0,1)=F_{m/2}
\qquad
(\forall\,v\in X_{m-2}).
\]
\end{proposition}
\begin{proof}
若
\[
\mathcal S_{m,d_{\min}(m)}=\iota_{m-2}(X_{m-2}),
\]
则对每个 \(v\in X_{m-2}\)，像点 \(\iota_{m-2}(v)=(v,0,1)\) 属于最小退化扇区，故
\[
d_m^{\mathrm{bin}}(v,0,1)=d_{\min}(m)=F_{m/2}.
\]

反之，若对所有 \(v\in X_{m-2}\) 都有
\[
d_m^{\mathrm{bin}}(v,0,1)=F_{m/2}=d_{\min}(m),
\]
则
\[
\iota_{m-2}(X_{m-2})\subseteq \mathcal S_{m,d_{\min}(m)}.
\]
又 \(\iota_{m-2}\) 显然为单射，因此
\[
\bigl|\iota_{m-2}(X_{m-2})\bigr|
=
|X_{m-2}|
=
F_m
=
\bigl|\mathcal S_{m,d_{\min}(m)}\bigr|.
\]
包含关系与基数相同遂推出
\[
\mathcal S_{m,d_{\min}(m)}=\iota_{m-2}(X_{m-2}).
\]
证毕。
\end{proof}

\begin{theorem}[fold-gauge 缺陷的五次重整化正规形]\label{thm:xi-time-part64a-fold-gauge-quintic-renormalization-normal-form}
沿用定理 \ref{thm:xi-time-part60ab-gauge-constant-recursive-bernoulli-stripping} 与
\ref{thm:xi-fold-gauge-defect-cubic-renormalization-normal-form} 的记号，记
\[
E_m:=\log C_m-\log(2\pi),
\qquad
q:=\frac{\varphi}{2}.
\]
则当 \(m\to\infty\) 时有精确提升
\[
E_{m+1}
=
qE_m
+\frac{3(2+\sqrt5)}{32}\,E_m^3
-\frac{27(655+294\sqrt5)}{4480}\,E_m^5
+O(E_m^7).
\]
特别地，
\[
\lim_{m\to\infty}
\frac{
E_{m+1}
-qE_m
-\frac{3(2+\sqrt5)}{32}\,E_m^3
}{
E_m^5
}
=
-\frac{27(655+294\sqrt5)}{4480}.
\]
\end{theorem}
\begin{proof}
由定理 \ref{thm:xi-time-part60ab-gauge-constant-recursive-bernoulli-stripping} 在 \(r=1,2,3\) 处的系数公式，
\[
E_m
=
aq^m+bq^{3m}+cq^{5m}+O(q^{7m}),
\]
其中
\[
a=\frac{1}{3\varphi},
\qquad
b=
\frac{B_4}{2\cdot 3}\,(\varphi^{-1}+\varphi)
=
-\frac{\sqrt5}{180},
\]
\[
c=
\frac{B_6}{3\cdot 5}\,(\varphi^{-1}+\varphi^3)
=
\frac{1}{42\cdot 15}\cdot 3\varphi
=
\frac{\varphi}{210}.
\]
现取 \(\alpha,\beta\) 使
\[
E_{m+1}-qE_m-\alpha E_m^3-\beta E_m^5=O(q^{7m}),
\]
并比较 \(q^{3m}\) 与 \(q^{5m}\) 的系数。

首先，
\[
E_{m+1}-qE_m
=
b(q^3-q)q^{3m}+c(q^5-q)q^{5m}+O(q^{7m}).
\]
另一方面，
\[
E_m^3=a^3q^{3m}+3a^2b\,q^{5m}+O(q^{7m}),
\qquad
E_m^5=a^5q^{5m}+O(q^{7m}).
\]
比较 \(q^{3m}\) 项可得
\[
\alpha=\frac{b(q^3-q)}{a^3}
=
\frac{3(2+\sqrt5)}{32}.
\]
再比较 \(q^{5m}\) 项，得到
\[
\beta
=
\frac{c(q^5-q)-3\alpha a^2b}{a^5}
=
-\frac{27(655+294\sqrt5)}{4480}.
\]
于是
\[
E_{m+1}
=
qE_m
+\frac{3(2+\sqrt5)}{32}\,E_m^3
-\frac{27(655+294\sqrt5)}{4480}\,E_m^5
+O(q^{7m}).
\]
又因 \(a>0\)，有
\[
E_m\sim aq^m,
\]
故 \(q^{7m}=O(E_m^7)\)。这就得到所示五次正规形。

最后，将等式两边减去前三项并除以 \(E_m^5\)，再利用 \(E_m^7= o(E_m^5)\)，便得到极限公式。证毕。
\end{proof}

\begin{theorem}[fold-gauge 重整化的纯奇次性]\label{thm:xi-time-part64a-fold-gauge-purely-odd-renormalization}
沿用定理 \ref{thm:xi-time-part60ab-gauge-constant-recursive-bernoulli-stripping} 的记号。存在唯一形式幂级数
\[
\mathcal R(z)=\frac{\varphi}{2}z+\sum_{j\ge 1}c_{2j+1}z^{2j+1}
\in z\mathbb R[[z^2]]
\]
使得对任意 \(N\ge 1\)，若记其截断式
\[
\mathcal R_N(z):=\frac{\varphi}{2}z+\sum_{j=1}^{N}c_{2j+1}z^{2j+1},
\]
则当 \(m\to\infty\) 时有
\[
E_{m+1}=\mathcal R_N(E_m)+O(E_m^{2N+3}).
\]
换言之，fold-gauge 缺陷的全部有限阶重整化正规形只含奇次幂，而一切偶次非线性项在所有阶上全部消失。
\end{theorem}
\begin{proof}
由定理 \ref{thm:xi-time-part60ab-gauge-constant-recursive-bernoulli-stripping}，
\[
E_m=\sum_{r\ge 1}a_r q^{(2r-1)m},
\qquad
q:=\frac{\varphi}{2},
\]
其中
\[
a_1=\frac{1}{3\varphi}\neq 0.
\]
令
\[
F(y):=\sum_{r\ge 1}a_r y^{2r-1}\in y\mathbb R[[y^2]].
\]
若记 \(y_m:=q^m\)，则
\[
E_m=F(y_m).
\]

由于 \(F\) 的线性项系数 \(a_1\) 非零，\(F\) 在形式幂级数意义下存在唯一复合逆
\[
F^{-1}(z)\in z\mathbb R[[z]].
\]
又 \(F\) 为奇级数，故 \(F^{-1}\) 亦为奇级数，从而
\[
F^{-1}(z)\in z\mathbb R[[z^2]].
\]
定义
\[
\mathcal R(z):=F\!\bigl(q\,F^{-1}(z)\bigr).
\]
因为 \(F\) 与 \(F^{-1}\) 都是奇级数，故
\[
\mathcal R(z)\in z\mathbb R[[z^2]].
\]

现验证该级数给出所求重整化。由 \(y_{m+1}=qy_m\)，有
\[
E_{m+1}
=
F(y_{m+1})
=
F(qy_m)
=
F\!\bigl(q\,F^{-1}(E_m)\bigr)
=
\mathcal R(E_m).
\]
这说明 \(\mathcal R\) 就是所求的形式重整化映射。

再证唯一性。若另有 \(\widetilde{\mathcal R}(z)\in z\mathbb R[[z]]\) 满足
\[
F(qy)=\widetilde{\mathcal R}(F(y))
\]
作为形式幂级数恒等式成立，则以 \(y=F^{-1}(z)\) 代入得到
\[
\widetilde{\mathcal R}(z)=F\!\bigl(qF^{-1}(z)\bigr)=\mathcal R(z),
\]
故 \(\mathcal R\) 唯一。

最后，对任意 \(N\ge 1\)，令 \(\mathcal R_N\) 为 \(\mathcal R\) 的 \(2N+1\) 阶截断。则
\[
\mathcal R(z)-\mathcal R_N(z)=O(z^{2N+3})
\qquad(z\to 0).
\]
由于 \(E_m\to 0\)，代入 \(z=E_m\) 即得
\[
E_{m+1}
=
\mathcal R(E_m)
=
\mathcal R_N(E_m)+O(E_m^{2N+3}).
\]
证毕。
\end{proof}

\noindent
上述六条结果把三条先前分立的闭合链进一步压缩到两个统一结构层：在算术侧，三域 Frobenius 数据被四维 \(C_2\)-符号盒完全组织，并使 \(L_4\) 因子仅以无偏条件化与 \(1/48\) 稀释的方式介入；在 bin-fold/gauge 侧，最小退化子覆盖猜想被还原为尾部固定 \(01\) 的逐点多重度判据，而 fold-gauge 缺陷的全部高阶重整化则被唯一的纯奇形式自映射严格控制。

\paragraph{部分可逆尾量、二原子信息塌缩、乘法外置与条件正交的横向闭合}
在预言机容量闭式、bin-fold 相对均匀基线的两点 \(f\)-散度极限、稳定地址半环的 PTIME 双向可解释性、一般 Ostrowski 数字层的规范化函子性、三轴 CRT 最小窗以及十次分歧/Lee--Yang 三次的 Chebotarev 乘法律已经建立之后，还可把部分可逆误差、信息几何塌缩、数字层局域性、三轴最小窗及跨指纹条件独立进一步压缩为下列八条结论。前四条把最优部分读出的全部损失先后还原为纤维超额尾量、Neyman--Pearson 对偶与矩谱驱动的非渐近失败界；中间三条分别把渐近 \(f\)-几何压缩到唯一二原子似然比测度、把 Zeckendorf/Ostrowski 层的局域可加性与乘法账本不可有限化严格分离，并指出三轴最小窗在最小层面即不可纯化；最后一条则表明十次线性因子指纹对 Lee--Yang 的符号/根/全分裂粗读出保持完全条件不变。

\begin{theorem}[最优部分可逆率的溢出恒等式]\label{thm:xi-time-part64b-optimal-partial-invertibility-overflow}
\leanverified{paper\_xi\_time\_part64b\_optimal\_partial\_invertibility\_overflow}
沿用定理 \ref{thm:fold-oracle-capacity-closed-form} 的记号。记
\[
d_m(x):=\bigl|\Fold_m^{-1}(x)\bigr|,
\qquad
T:=2^B,
\]
\[
C_m(B):=\mathcal C_m(B)=\sum_{x\in X_m}\min\bigl(d_m(x),T\bigr),
\qquad
\mathrm{Succ}_m(B):=\frac{C_m(B)}{2^m}.
\]
则有严格恒等式
\[
\boxed{\;
1-\mathrm{Succ}_m(B)
=
2^{-m}\sum_{x\in X_m}\bigl(d_m(x)-T\bigr)_+.\;}
\]
换言之，全部不可逆损失恰等于超过阈值 \(T\) 的纤维超额体积之归一化总和。
\end{theorem}
\begin{proof}
对任意整数 \(d\ge 0\)，都有逐点恒等式
\[
\min(d,T)=d-(d-T)_+.
\]
将其代入容量闭式得
\[
C_m(B)
=
\sum_{x\in X_m}\Bigl[d_m(x)-\bigl(d_m(x)-T\bigr)_+\Bigr].
\]
又
\[
\sum_{x\in X_m}d_m(x)=|\Omega_m|=2^m,
\]
故
\[
C_m(B)=2^m-\sum_{x\in X_m}\bigl(d_m(x)-T\bigr)_+.
\]
两边同除以 \(2^m\)，便得到
\[
1-\mathrm{Succ}_m(B)
=
2^{-m}\sum_{x\in X_m}\bigl(d_m(x)-T\bigr)_+.
\]
证毕。
\end{proof}

\begin{corollary}[近完美读出的尖锐必要条件]\label{cor:xi-time-part64b-near-perfect-readout-sharp-necessary-condition}
沿用定理 \ref{thm:xi-time-part64b-optimal-partial-invertibility-overflow} 的记号，并记
\[
M_m:=\max_{x\in X_m}d_m(x).
\]
若对某个 \(\varepsilon\ge 0\) 有
\[
\mathrm{Succ}_m(B)\ge 1-\varepsilon,
\]
则必有
\[
\boxed{\;
2^B\ge M_m-\varepsilon\,2^m.\;}
\]
特别地：
\begin{enumerate}
  \item 当 \(\varepsilon=0\) 时，
  \[
  2^B\ge M_m
  \]
  是零误差全可逆的必要条件；
  \item 若 \(2^B<M_m\)，则必有严格损失下界
  \[
  \boxed{\;
  1-\mathrm{Succ}_m(B)\ge \frac{M_m-2^B}{2^m}.\;}
  \]
\end{enumerate}
\end{corollary}
\begin{proof}
由定理 \ref{thm:xi-time-part64b-optimal-partial-invertibility-overflow}，
\[
1-\mathrm{Succ}_m(B)
=
2^{-m}\sum_{x\in X_m}\bigl(d_m(x)-2^B\bigr)_+.
\]
由于对某个达到最大值的 \(x_\ast\in X_m\) 有 \(d_m(x_\ast)=M_m\)，故
\[
1-\mathrm{Succ}_m(B)
\ge
2^{-m}\bigl(M_m-2^B\bigr)_+.
\]
若 \(\mathrm{Succ}_m(B)\ge 1-\varepsilon\)，则
\[
\varepsilon\ge \frac{\bigl(M_m-2^B\bigr)_+}{2^m}.
\]
若 \(M_m\le 2^B\)，所示结论显然成立；若 \(M_m>2^B\)，则上式化为
\[
M_m-2^B\le \varepsilon\,2^m,
\]
即
\[
2^B\ge M_m-\varepsilon\,2^m.
\]
其余两条断言分别对应 \(\varepsilon=0\) 与 \(M_m>2^B\) 的特例。证毕。
\end{proof}

\begin{theorem}[预言机容量的 Neyman--Pearson 对偶：推前分布与稳定层均匀基线的精确差异表达]\label{thm:xi-time-part64b-oracle-capacity-neyman-pearson-duality}
沿用定理 \ref{thm:fold-oracle-capacity-closed-form} 的记号。定义推前分布与稳定层均匀基线
\[
\mu_m(x):=\frac{d_m(x)}{2^m},
\qquad
u_m(x):=\frac{1}{|X_m|}
\qquad (x\in X_m),
\]
并令阈值 \(T:=2^B\)。再定义差异泛函
\[
\mathfrak D_m(B):=
\sup_{A\subseteq X_m}
\Bigl(
\mu_m(A)-2^{B-m}|X_m|\,u_m(A)
\Bigr).
\]
则有严格恒等式
\[
\mathrm{Succ}_m(B)=1-\mathfrak D_m(B),
\]
并且上确界恰由阈值集
\[
A_B^\star:=\{x\in X_m:\ d_m(x)>2^B\}
\]
取得。
\end{theorem}
\begin{proof}
由于
\[
|X_m|\,u_m(A)=|A|,
\]
故对任意 \(A\subseteq X_m\)，有
\[
\mu_m(A)-2^{B-m}|X_m|\,u_m(A)
=
\sum_{x\in A}\left(\frac{d_m(x)}{2^m}-\frac{2^B}{2^m}\right)
=
\frac{1}{2^m}\sum_{x\in A}\bigl(d_m(x)-2^B\bigr).
\]
上式在且仅在取遍所有正项时达到最大，因此
\[
\mathfrak D_m(B)
=
\frac{1}{2^m}
\sum_{x:\,d_m(x)>2^B}\bigl(d_m(x)-2^B\bigr),
\]
且最大集正是 \(A_B^\star\)。

另一方面，由定理 \ref{thm:xi-time-part64b-optimal-partial-invertibility-overflow} 的逐点恒等式
\[
\min(d,2^B)=d-(d-2^B)_+
\]
可得
\[
C_m(B)
=
\sum_{x\in X_m}\min\bigl(d_m(x),2^B\bigr)
=
2^m-\sum_{x\in X_m}\bigl(d_m(x)-2^B\bigr)_+.
\]
因此
\[
\mathfrak D_m(B)
=
\frac{2^m-C_m(B)}{2^m}
=
1-\mathrm{Succ}_m(B).
\]
证毕。
\end{proof}

\begin{theorem}[幂和诱导的非渐近失败率指数界与可计算预算下界]\label{thm:xi-time-part64b-power-sum-failure-bound-budget-lower}
对任意 \(q>1\)，定义幂和
\[
S_q(m):=\sum_{x\in X_m}d_m(x)^q.
\]
则对任意实预算 \(B\in\mathbb R\) 都有失败率上界
\[
1-\mathrm{Succ}_m(B)
\le
\frac{S_q(m)}{2^m\,2^{(q-1)B}}.
\]
因此，若希望
\[
\mathrm{Succ}_m(B)\ge 1-\varepsilon
\qquad (\varepsilon\in(0,1)),
\]
则充分条件为
\[
B\ge
\frac{\log_2 S_q(m)-m-\log_2\varepsilon}{q-1}.
\]
对右端再对 \(q>1\) 取下确界，便得到完全由矩谱驱动的最强可计算预算下界。
\end{theorem}
\begin{proof}
记
\[
T:=2^B.
\]
由定理 \ref{thm:xi-time-part64b-optimal-partial-invertibility-overflow}，
\[
2^m-C_m(B)=\sum_{x\in X_m}(d_m(x)-T)_+.
\]
对任意 \(d\ge 0\) 与 \(q>1\)，成立初等不等式
\[
(d-T)_+\le \frac{d^q}{T^{q-1}}.
\]
确实，若 \(d\le T\)，左端为 \(0\)；若 \(d>T\)，则
\[
d-T\le d\le d\left(\frac{d}{T}\right)^{q-1}=\frac{d^q}{T^{q-1}}.
\]
逐项求和即得
\[
2^m-C_m(B)
\le
\frac{1}{T^{q-1}}\sum_{x\in X_m}d_m(x)^q
=
\frac{S_q(m)}{2^{(q-1)B}}.
\]
两边同除以 \(2^m\)，便有
\[
1-\mathrm{Succ}_m(B)\le \frac{S_q(m)}{2^m\,2^{(q-1)B}}.
\]
若右端不超过 \(\varepsilon\)，则 \(\mathrm{Succ}_m(B)\ge 1-\varepsilon\)。对上式整理即得
\[
B\ge
\frac{\log_2 S_q(m)-m-\log_2\varepsilon}{q-1}.
\]
证毕。
\end{proof}

\begin{theorem}[折叠信息几何的二点塌缩原理]\label{thm:xi-time-part64b-fold-information-geometry-two-point-collapse}
沿用命题 \ref{prop:fold-bin-f-divergence-collapse} 的记号。定义
\[
a:=\frac{\varphi^2}{\sqrt5},
\qquad
b:=\frac{\varphi}{\sqrt5},
\qquad
\nu_\varphi:=\frac1\varphi\,\delta_a+\frac1{\varphi^2}\,\delta_b.
\]
则对任意凸函数 \(f:(0,\infty)\to\RR\) 满足 \(f(1)=0\)，其对应的 Csisz\'ar \(f\)-散度满足
\[
\boxed{\;
\lim_{m\to\infty}D_f(\mu_m\Vert u_m)
=
\int_{(0,\infty)}f(t)\,\dd\nu_\varphi(t)
=
\frac1\varphi\,f\!\left(\frac{\varphi^2}{\sqrt5}\right)
+\frac1{\varphi^2}\,f\!\left(\frac{\varphi}{\sqrt5}\right).\;}
\]
因此，\((\mu_m,u_m)\) 的全部渐近 \(f\)-信息几何并不依赖于完整纤维谱，而完全塌缩到两个似然比原子 \(a,b\) 上；换言之，折叠多重度在信息论极限中仅保留二相可辨结构。
\end{theorem}
\begin{proof}
命题 \ref{prop:fold-bin-f-divergence-collapse} 已直接给出
\[
D_f(\mu_m\Vert u_m)
=
\frac1\varphi\,f\!\left(\frac{\varphi^2}{\sqrt5}\right)
+\frac1{\varphi^2}\,f\!\left(\frac{\varphi}{\sqrt5}\right)
+O_f\bigl((\varphi/2)^m\bigr),
\]
故所示极限公式立即成立。

另一方面，同一命题的证明还给出一致展开
\[
\frac{\mu_m(w)}{u_m(w)}
=
\frac{\varphi^{2-w_m}}{\sqrt5}+O\bigl((\varphi/2)^m\bigr)
\qquad (w\in X_m),
\]
因而极限似然比仅可能落在
\[
\frac{\varphi^2}{\sqrt5}=a,
\qquad
\frac{\varphi}{\sqrt5}=b
\]
两点上。再由
\[
u_m(w_m=0)=\frac{F_{m+1}}{F_{m+2}}\longrightarrow \frac1\varphi,
\qquad
u_m(w_m=1)=\frac{F_m}{F_{m+2}}\longrightarrow \frac1{\varphi^2},
\]
可知极限似然比测度恰为
\[
\nu_\varphi=\frac1\varphi\,\delta_a+\frac1{\varphi^2}\,\delta_b.
\]
于是任意由似然比驱动并在该一致收敛下稳定的 \(f\)-统计，都只能因子化通过 \(\nu_\varphi\)。证毕。
\end{proof}

\begin{theorem}[局域可加性与非局域可乘性的严格偏振]\label{thm:xi-time-part64b-local-additivity-nonlocal-multiplicativity-polarization}
沿用定义 \ref{def:stable-add}、\ref{def:stable-mul}，定理 \ref{thm:normalize-then-truncate-functorial}，注 \ref{rem:zeckendorf-normalize-fst} 以及命题 \ref{prop:killo-implementation-vs-structural-half-circle-dimension} 的记号。则在本文构造的 Zeckendorf/Ostrowski 数字层中，规范化与加法具有纯数字串实现：可在有限状态换能器或局部重写框架内完成，并在分辨率口径上线性、在比特口径上为多项式时间。

然而，若坚持“乘法 \(\mapsto\) 有限寄存器加法”的严格同态语义，则不存在任何有限秩压缩：既不存在单射幺半群同态
\[
\NN_{\times}\hookrightarrow (\NN^k,+)
\qquad(k<\infty),
\]
亦不存在有限生成交换群 \(L\) 使得
\[
\Phi(ab)=\Phi(a)+\Phi(b),
\qquad
\Phi:\NN_{\times}\to \ZZ\oplus L
\]
为单射。

并且
\[
\boxed{\;
\cdim^{\mathrm{hom}}_{+}(\NN_{\times})=+\infty,
\qquad
\cdim^{\mathrm{impl}}(\NN_{\times})=\frac12.\;}
\]
因此，该体系中的精确算术必然分裂为两层：一层是局域、有限状态、纯数字的可加层；另一层是任何有限加法账本都无法吸收的外置乘法层。
\end{theorem}
\begin{proof}
先看数字层的局域实现。定理 \ref{thm:normalize-then-truncate-functorial} 已把一般 Ostrowski 数字层的“先规范形再截断”固定为定义性的局部机制，而注 \ref{rem:zeckendorf-normalize-fst} 进一步指出，在黄金分支上 Zeckendorf/Ostrowski 规范化与在线加法可由纯数字串层的有限状态换能器或局部重写实现，从而在分辨率口径上具有线性复杂度。另一方面，定理 \ref{thm:xi-stable-semiring-ptime-biinterpretability} 又给出值层回拉实现的多项式时间界：稳定地址层与自然数半环在 PTIME 内双向可解释，故在比特口径上，加法、规范化以及经由值层回拉的乘法都具有多项式时间实现。

再看“乘法 \(\mapsto\) 有限寄存器加法”的严格同态化是否可能。推论 \ref{cor:killo-no-finite-additive-register-linearization} 断言：对任意有限 \(k\)，不存在单射幺半群同态
\[
\NN_{\times}\hookrightarrow (\NN^k,+).
\]
定理 \ref{thm:killo-godel-compression-not-finite-rank-homologizable} 又断言：若 \(L\) 有限生成，则不存在单射
\[
\Phi:\NN_{\times}\to \ZZ\oplus L
\]
满足
\[
\Phi(ab)=\Phi(a)+\Phi(b).
\]
因此，一切有限秩加法账本上的严格乘法同态化都被排除。

最后，由命题 \ref{prop:killo-implementation-vs-structural-half-circle-dimension}，
\[
\cdim^{\mathrm{hom}}_{+}(\NN_{\times})=+\infty,
\qquad
\cdim^{\mathrm{impl}}(\NN_{\times})=\frac12.
\]
这正表明：作为集合，\(\NN_{\times}\) 仍可嵌入单寄存器实现层；但作为乘法幺半群，其结构层同态半圆维不可有限化。故局域数字可加性与乘法外置性形成严格偏振。证毕。
\end{proof}

\begin{theorem}[三轴最小窗的非纯化定理]\label{thm:xi-time-part64b-three-axis-minimum-window-non-purification}
设 \(m_\star\) 为命题 \ref{prop:crt-235-min-depth} 给出的三轴窗最小分辨率。则
\[
\boxed{\;
m_\star=58,
\qquad
m_\star+2=60.\;}
\]
其对应底层 Fibonacci 模数满足
\[
\boxed{\;
F_{60}
=
2^4\cdot 3^2\cdot 5\cdot 11\cdot 31\cdot 41\cdot 61\cdot 2521.\;}
\]
因此，第一次三轴共现并非仅引入目标三条算术方向 \(2,3,5\)；它被严格迫使同时携带五条额外素轴 \(11,31,41,61,2521\)。换言之，三轴统一在最小可见层上即不可纯化。
\end{theorem}
\begin{proof}
命题 \ref{prop:crt-235-min-depth} 已给出最小三轴窗的精确位置
\[
m_\star=58,
\qquad
m_\star+2=60.
\]
另一方面，注 \ref{rem:crt-235-m58-factorization} 已给出对应 Fibonacci 模数的精确分解
\[
F_{60}
=
1548008755920
=
2^4\cdot 3^2\cdot 5\cdot 11\cdot 31\cdot 41\cdot 61\cdot 2521.
\]
故在最小三轴层中，除目标三轴 \(2,3,5\) 外，还不可避免地并入五条额外素轴
\[
11,\ 31,\ 41,\ 61,\ 2521.
\]
这就说明三轴最小窗并不能被“纯三轴”算术所单独承载，而天然带有额外的剩余素数方向。证毕。
\end{proof}

\begin{theorem}[十次指纹对 Lee--Yang 粗读出的条件不变性]\label{thm:xi-time-part64b-p10-fingerprint-leyang-coarse-readout-invariance}
记
\[
A:=\{\text{\(P_{10}\bmod p\) 含一次因子}\},
\]
并定义 Lee--Yang 侧三个粗事件
\[
E_1:=\{\chi_{\mathrm{LY}}(p)=+1\},
\qquad
E_2:=\{\chi_{\mathrm{LY}}(p)=-1\},
\qquad
E_3:=B_{\mathrm{root}}:=\{\text{\(g\bmod p\) 至少有一根}\}.
\]
再记
\[
B_{\mathrm{split}}:=\{\text{\(g\bmod p\) 全分裂}\}.
\]
则自然密度满足
\[
\boxed{\;
\delta(A\mid E_1)
=
\delta(A\mid E_2)
=
\delta(A\mid E_3)
=
\delta(A)
=
\frac{28319}{44800}.\;}
\]
并且
\[
\boxed{\;
\delta(A\mid B_{\mathrm{split}})
=
\frac{28319}{44800}.\;}
\]
因此，十次模分解指纹与 Lee--Yang 的符号读出、存在根读出及全分裂读出在统计上完全正交：任何此类粗 Lee--Yang 观测都不改变十次线性因子事件的密度。
\end{theorem}
\begin{proof}
由命题 \ref{prop:fold-gauge-anomaly-P10-leyang-linear-disjointness} 与推论
\ref{cor:fold-gauge-anomaly-P10-leyang-chebotarev-product-law}，任意只依赖 \(S_{10}\) 因子的共轭类事件与任意只依赖 \(S_3\) 因子的共轭类事件，其联合自然密度都等于边际密度之积。

事件 \(A\) 只依赖 \(P_{10}\) 分裂域的 \(S_{10}\) 因子；而
\[
E_1,\ E_2,\ B_{\mathrm{root}},\ B_{\mathrm{split}}
\]
都只依赖 Lee--Yang 三次分裂域的 \(S_3\) 因子。故对任意
\[
E\in\{E_1,E_2,B_{\mathrm{root}},B_{\mathrm{split}}\}
\]
都成立
\[
\delta(A\cap E)=\delta(A)\,\delta(E).
\]

再由推论 \ref{cor:fold-gauge-anomaly-P10-modp-splitting-densities}，
\[
\delta(A)=\frac{28319}{44800}.
\]
另一方面，在 \(S_3\) 中，偶置换与奇置换各占 \(3\) 个元素，故
\[
\delta(E_1)=\delta(E_2)=\frac12.
\]
又由结论 \ref{con:xi-terminal-zm-leyang-cubic-modp-splitting-density}，\(g\bmod p\) 的分裂型与 \(S_3\) 的循环型一一对应：恒等类对应全分裂，换位类对应恰有一根，三循环类对应不可约。三类大小分别为 \(1,3,2\)，故
\[
\delta(B_{\mathrm{split}})=\frac16,
\qquad
\delta(B_{\mathrm{root}})=\frac{1+3}{6}=\frac23.
\]

由于上述四个条件事件的密度均为正，分别相除即可得到
\[
\delta(A\mid E_1)=\delta(A),
\qquad
\delta(A\mid E_2)=\delta(A),
\]
\[
\delta(A\mid B_{\mathrm{root}})=\delta(A),
\qquad
\delta(A\mid B_{\mathrm{split}})=\delta(A).
\]
将 \(\delta(A)=28319/44800\) 代回，即得全部公式。证毕。
\end{proof}

\noindent
上述六条结果把容量亏损、极限信息几何、数字层实现、三轴最小窗与 Chebotarev 条件化五条原先分立的链进一步压缩为同一组有限证书：部分可逆误差由纤维超额尾量完全记录，渐近似然比只剩两原子，乘法的外置性由同态半圆维无穷障碍强制，最小三轴窗天然携带额外素轴，而十次指纹对 Lee--Yang 粗读出保持严格条件不变。于是，恢复论、信息论、实现论、模窗算术与分裂型统计之间的横向耦合便获得了统一的有限闭合。

\paragraph{stop-loss 容量曲线、幂散度单标量刚性、平衡剖面与最小扇区 Witten 电荷}
在临界预言机容量双折线、二原子极限似然比、\(\chi^2\)-型极限常数对幂散度族的坐标化、折叠多重度的 Robin--Hood 调平机制以及审计偶窗最小退化扇区的同伦相位都已建立之后，还可把这些链条进一步压缩为下列七条彼此衔接的对象层后果。

\begin{theorem}[临界 oracle 曲线恰为极限似然比的 stop-loss 变换]\label{thm:xi-time-part64ba-critical-oracle-stoploss-transform}
沿用定理 \ref{thm:xi-foldbin-dyadic-capacity-critical-window-limit} 与定理 \ref{thm:xi-time-part64b-fold-information-geometry-two-point-collapse} 的记号。定义
\[
r_0:=\frac{\varphi^2}{\sqrt5},
\qquad
r_1:=\frac{\varphi}{\sqrt5},
\]
并令 \(R_\infty\) 为二原子随机变量
\[
R_\infty=
\begin{cases}
r_0,& \text{概率 } \varphi^{-1},\\[1mm]
r_1,& \text{概率 } \varphi^{-2}.
\end{cases}
\]
再记临界尺度下的最优成功率极限曲线为
\[
S(\tau):=\mathcal S_{\mathrm{bin}}(\tau)
\qquad (\tau\ge 0),
\]
即
\[
S(\tau)=
\begin{cases}
\dfrac{\varphi^2}{\sqrt5}\,\tau,& 0\le \tau\le \varphi^{-1},\\[3mm]
\dfrac{1}{\varphi\sqrt5}+\dfrac{\varphi}{\sqrt5}\,\tau,& \varphi^{-1}\le \tau\le 1,\\[3mm]
1,& \tau\ge 1.
\end{cases}
\]
则对一切 \(\tau\ge 0\)，成立精确恒等式
\[
\boxed{\;
S(\tau)=\EE\!\left[\min\{R_\infty,\tau r_0\}\right]
=
\int_{0}^{\tau r_0}\PP(R_\infty>t)\,\dd t.\;}
\]
特别地，当 \(\tau\notin\{\varphi^{-1},1\}\) 时，
\[
\boxed{\;
S'(\tau)=r_0\,\PP(R_\infty>\tau r_0).\;}
\]
因此，临界 oracle 曲线的全部折点恰由极限似然比尾分布的 stop-loss 变换完全产生。
\end{theorem}
\begin{proof}
记
\[
\lambda:=\tau r_0.
\]
分三段讨论。

若 \(0\le \tau\le \varphi^{-1}\)，则
\[
\lambda\le r_1,
\]
故
\[
\EE[\min\{R_\infty,\lambda\}]=\lambda=\frac{\varphi^2}{\sqrt5}\tau.
\]

若 \(\varphi^{-1}\le \tau\le 1\)，则
\[
r_1\le \lambda\le r_0,
\]
于是
\[
\EE[\min\{R_\infty,\lambda\}]
=
\varphi^{-2}r_1+\varphi^{-1}\lambda
=
\frac{1}{\varphi\sqrt5}+\frac{\varphi}{\sqrt5}\tau.
\]

若 \(\tau\ge 1\)，则
\[
\lambda\ge r_0,
\]
故
\[
\EE[\min\{R_\infty,\lambda\}]
=
\EE[R_\infty]
=
\varphi^{-1}r_0+\varphi^{-2}r_1
=1.
\]
这与 \(S(\tau)\) 的三段闭式完全一致。

第二个等式是 stop-loss 变换的标准层蛋糕恒等式
\[
\EE[\min\{X,\lambda\}]=\int_0^\lambda \PP(X>t)\,\dd t
\qquad (X\ge 0),
\]
应用于 \(X=R_\infty\) 即得。对 \(\tau\) 在非折点处求导，链式法则给出
\[
S'(\tau)=r_0\,\PP(R_\infty>\tau r_0).
\]
证毕。
\end{proof}

\begin{theorem}[单个 \texorpdfstring{$\chi^2$}{chi2} 型极限常数决定整个 power-divergence 谱，并满足二阶常系数微分方程]\label{thm:xi-time-part64ba-chi2-rigid-powerdivergence-ode}
记
\[
\chi:=D_{f_2}^\infty.
\]
则对每个 \(\alpha>1\)，全部极限 power-divergence 常数都可由 \(\chi\) 显式恢复：
\[
\boxed{\;
D_{f_\alpha,\infty}
=
5^{-\alpha/2}
\left[
\left(\frac{5\chi+3}{2}\right)^{2\alpha-1}
+
\left(\frac{5\chi+3}{2}\right)^{\alpha-2}
\right]-1.\;}
\]
若再令
\[
E(\alpha):=D_{f_\alpha,\infty}+1,
\]
\[
c_\chi:=\log\frac{(5\chi+3)^2}{4\sqrt5},
\qquad
d_\chi:=\log\frac{5\chi+3}{2\sqrt5},
\]
则 \(E\) 满足精确线性常微分方程
\[
\boxed{\;
E''(\alpha)-(c_\chi+d_\chi)E'(\alpha)+c_\chi d_\chi\,E(\alpha)=0.\;}
\]
因此，整个极限 power-divergence 族并非自由函数，而是由单个二阶统计常数 \(\chi\) 刚性锁定的二维指数簇。
\end{theorem}
\begin{proof}
由结论 \ref{con:xi-time-part9l-chi2-generates-power-divergence-and-cumulants}，
\[
D_{f_2}^\infty=\frac{2\varphi-3}{5},
\qquad
\varphi=\frac{5D_{f_2}^\infty+3}{2}.
\]
代入 \(\chi=D_{f_2}^\infty\) 后得到
\[
\varphi=\frac{5\chi+3}{2}.
\]
同一结论又给出一般公式
\[
D_{f_\alpha}^\infty
=
\frac{1}{5^{\alpha/2}}\bigl(\varphi^{2\alpha-1}+\varphi^{\alpha-2}\bigr)-1.
\]
把 \(\varphi=(5\chi+3)/2\) 代入，即得第一条闭式。

现看 \(E(\alpha)=D_{f_\alpha,\infty}+1\)。由上式，
\[
E(\alpha)
=
5^{-\alpha/2}\bigl(\varphi^{2\alpha-1}+\varphi^{\alpha-2}\bigr)
=
\varphi^{-1}\left(\frac{\varphi^2}{\sqrt5}\right)^\alpha
+
\varphi^{-2}\left(\frac{\varphi}{\sqrt5}\right)^\alpha.
\]
再用
\[
\varphi=\frac{5\chi+3}{2}
\]
重写两条指数率，便得
\[
\frac{\varphi^2}{\sqrt5}
=
\frac{(5\chi+3)^2}{4\sqrt5}
=
e^{c_\chi},
\qquad
\frac{\varphi}{\sqrt5}
=
\frac{5\chi+3}{2\sqrt5}
=
e^{d_\chi}.
\]
故
\[
E(\alpha)=A e^{c_\chi\alpha}+B e^{d_\chi\alpha},
\qquad
A=\varphi^{-1},
\qquad
B=\varphi^{-2}.
\]
任意两指数和都满足
\[
(\partial_\alpha-c_\chi)(\partial_\alpha-d_\chi)E=0,
\]
亦即
\[
E''-(c_\chi+d_\chi)E'+c_\chi d_\chi E=0.
\]
证毕。
\end{proof}

\begin{corollary}[黄金几何审计常数由同一个 \texorpdfstring{$\chi^2$}{chi2} 型极限标量完全恢复]\label{cor:xi-time-part64ba-chi2-recovers-golden-geometric-audits}
仍记 \(\chi=D_{f_2}^\infty\)。则黄金分支上既有 Wasserstein/星差异审计常数可全部写成 \(\chi\) 的有理分式：
\[
\boxed{\;
\lim_{\substack{n\to\infty\\ n\ \mathrm{even}}}F_n D_n^\ast
=
1+\frac{1}{5\chi+2}.\;}
\]
\[
\boxed{\;
\lim_{\substack{n\to\infty\\ n\ \mathrm{even}}}F_n T_n
=
\frac12+\frac{1}{2(5\chi+2)}.\;}
\]
\[
\boxed{\;
\lim_{\substack{n\to\infty\\ n\ \mathrm{odd}}}F_n T_n
=
\frac{17}{30}-\frac{1}{2(5\chi+2)}.\;}
\]
故同一个 \(\chi^2\)-型极限常数不仅决定信息散度族，而且刚性控制黄金分支的几何审计常数。
\end{corollary}
\begin{proof}
文稿在黄金分支上已给出
\[
\lim_{\substack{n\to\infty\\ n\ \mathrm{even}}}F_nT_n
=
\frac12+\frac{1}{2\sqrt5},
\]
\[
\lim_{\substack{n\to\infty\\ n\ \mathrm{odd}}}F_nT_n
=
\frac12-\frac{1}{2\sqrt5}+\frac{1}{15}
=
\frac{17}{30}-\frac{1}{2\sqrt5},
\]
并且偶数支上满足
\[
T_n=\frac12 D_n^\ast.
\]
另一方面，由结论 \ref{con:xi-time-part9l-chi2-generates-power-divergence-and-cumulants}，
\[
\sqrt5=5\chi+2.
\]
把该恒等式代入上面三条公式，便得到所示全部闭式。证毕。
\end{proof}

\begin{theorem}[每一级折叠多重度剖面都可由有限 Robin--Hood 链严格平衡化；平衡剖面位于其双随机轨道中]\label{thm:xi-time-part64ba-fold-multiplicity-majorization-balancing}
\leanverified{paper\_xi\_time\_part64ba\_fold\_multiplicity\_majorization\_balancing}
记
\[
F:=F_{m+2},
\qquad
M:=2^m,
\qquad
M=aF+\rho,
\qquad
0\le \rho<F.
\]
将折叠多重度剖面写为
\[
d^{(m)}:=(d_m(u))_{u\in \ZZ/(F\ZZ)}\in \ZZ_{\ge 0}^{F},
\]
并定义平衡剖面
\[
b^{(m)}:=
(\underbrace{a+1,\dots,a+1}_{\rho\text{ 个}},
\underbrace{a,\dots,a}_{F-\rho\text{ 个}}).
\]
则在按非增序重排后，有主序关系
\[
\boxed{\;
d^{(m),\downarrow}\succ b^{(m)}.\;}
\]
等价地，存在一个 \(F\times F\) 双随机矩阵 \(\Pi_m\) 使
\[
\boxed{\;
b^{(m)}=\Pi_m\,d^{(m)}.\;}
\]
更强地，可由有限次平衡化操作
\[
(v_{u_1},v_{u_2})\longmapsto (v_{u_1}-1,v_{u_2}+1)
\qquad (v_{u_1}\ge v_{u_2}+2)
\]
把 \(d^{(m)}\) 明确送到 \(b^{(m)}\)。
\end{theorem}
\begin{proof}
结论 \ref{con:xi-time-part9l-fold-fourier-energy-arithmetic-lower-bound} 的证明已经给出如下严格调平机制：若某个整数向量 \(v\in\ZZ_{\ge 0}^{F}\) 满足
\[
\sum_u v_u=M
\]
且存在指标 \(u_1,u_2\) 使
\[
v_{u_1}\ge v_{u_2}+2,
\]
则执行一步 Robin--Hood 操作
\[
(v_{u_1},v_{u_2})\mapsto (v_{u_1}-1,v_{u_2}+1)
\]
会严格降低平方偏差和，因此有限步之后必然终止于某个任意两坐标相差至多 \(1\) 的向量。

由于总和固定为
\[
M=aF+\rho,
\]
满足“任意两坐标相差至多 \(1\)” 的整数向量只能是恰有 \(\rho\) 个分量等于 \(a+1\)、其余 \(F-\rho\) 个分量等于 \(a\) 的排列，也就是 \(b^{(m)}\) 的一个置换。把这些分量按非增序排列，唯一得到
\[
b^{(m)}.
\]
因此 \(d^{(m)}\) 经过有限次 Robin--Hood 调平后确实到达 \(b^{(m)}\)。

而 Hardy--Littlewood--Pólya 判别表明：一个向量可经有限次 Robin--Hood 操作被送到另一个向量，当且仅当前者 majorize 后者。故
\[
d^{(m),\downarrow}\succ b^{(m)}.
\]
再由 Hardy--Littlewood--Pólya 定理与 Birkhoff--von Neumann 定理，主序关系等价于存在双随机矩阵 \(\Pi_m\) 使
\[
b^{(m)}=\Pi_m d^{(m)}.
\]
证毕。
\end{proof}

\begin{corollary}[平衡剖面给出全阶 Rényi 熵的统一上包络]\label{cor:xi-time-part64ba-balanced-profile-universal-renyi-envelope}
在定理 \ref{thm:xi-time-part64ba-fold-multiplicity-majorization-balancing} 的记号下，定义
\[
p_m(u):=\frac{d_m(u)}{2^m},
\qquad
p_m^{\mathrm{bal}}:=\frac{b^{(m)}}{2^m}.
\]
则对任意 \(\alpha>0\)，Rényi 熵满足
\[
\boxed{\;
H_\alpha(p_m)\le H_\alpha(p_m^{\mathrm{bal}}).\;}
\]
其中当 \(\alpha\neq 1\) 时，
\[
H_\alpha(p):=\frac{1}{1-\alpha}\log\sum_u p(u)^\alpha,
\]
而上包络显式为
\[
\boxed{\;
H_\alpha(p_m)
\le
\frac{1}{1-\alpha}
\log\!\left(
\frac{(F-\rho)a^\alpha+\rho(a+1)^\alpha}{(aF+\rho)^\alpha}
\right).\;}
\]
在 Shannon 极限 \(\alpha\to 1\) 下，
\[
\boxed{\;
H(p_m)\le
-\frac{(F-\rho)a}{M}\log\frac{a}{M}
-\frac{\rho(a+1)}{M}\log\frac{a+1}{M}.\;}
\]
并且上述各式取等当且仅当多重度剖面本身已经平衡。
\end{corollary}
\begin{proof}
由定理 \ref{thm:xi-time-part64ba-fold-multiplicity-majorization-balancing}，
\[
d^{(m),\downarrow}\succ b^{(m)}.
\]
归一化后仍有
\[
p_m\succ p_m^{\mathrm{bal}}.
\]
而 Rényi 熵对一切 \(\alpha>0\) 都是 Schur--凹函数，故
\[
H_\alpha(p_m)\le H_\alpha(p_m^{\mathrm{bal}}).
\]
把平衡剖面
\[
p_m^{\mathrm{bal}}=
\left(
\underbrace{\frac{a+1}{M},\dots,\frac{a+1}{M}}_{\rho\text{ 个}},
\underbrace{\frac{a}{M},\dots,\frac{a}{M}}_{F-\rho\text{ 个}}
\right)
\]
直接代入 Rényi 熵定义，便得到
\[
H_\alpha(p_m^{\mathrm{bal}})
=
\frac{1}{1-\alpha}
\log\!\left(
\frac{(F-\rho)a^\alpha+\rho(a+1)^\alpha}{M^\alpha}
\right).
\]
其中 \(M=aF+\rho\)，故即为题述上界。令 \(\alpha\to 1\) 即得 Shannon 形式。

最后，由定理 \ref{thm:xi-time-part64ba-fold-multiplicity-majorization-balancing} 中 Robin--Hood 降平过程的严格性可知，若 \(d^{(m)}\) 不是平衡剖面，则主序关系严格；而 Rényi 熵在此情形下严格 Schur--凹，故等号仅当 \(d^{(m)}\) 本身已平衡时成立。证毕。
\end{proof}

\begin{theorem}[已审计偶窗的最小退化扇区携带一条显式可求的总 Witten 电荷]\label{thm:xi-time-part64ba-audited-even-minsector-total-witten-charge}
对审计偶窗
\[
m\in\{6,8,10,12\},
\]
定义最小退化扇区
\[
S_{m,\min}:=\{x\in X_m:\ d_m(x)=d_{\min}(m)\},
\]
以及交错 Witten 指数
\[
W_x:=\sum_{S\in \mathsf{Ind}(\Gamma_x)}(-1)^{|S|}.
\]
再定义总 Witten 电荷
\[
\mathcal W_{\min}(m):=\sum_{x\in S_{m,\min}}W_x.
\]
则有精确公式
\[
\boxed{\;
\mathcal W_{\min}(m)=
\begin{cases}
0,& m=6,12,\\[1mm]
-F_m,& m=8,10.
\end{cases}\;}
\]
特别地，
\[
\boxed{\;
\frac{\mathcal W_{\min}(m)}{|S_{m,\min}|}=
\begin{cases}
0,& m=6,12,\\[1mm]
-1,& m=8,10.
\end{cases}\;}
\]
因此，最小退化扇区在审计偶窗上并不只是经历“可缩/不可缩”二分，而且携带一条从 \(0\) 到 \(-F_m\) 的整扇区总电荷跃迁。
\end{theorem}
\begin{proof}
由定理 \ref{thm:fold-bin-min-degeneracy-fib-audited}，
\[
d_{\min}(6)=2,\qquad d_{\min}(8)=3,\qquad d_{\min}(10)=5,\qquad d_{\min}(12)=8,
\]
且
\[
|S_{m,\min}|=F_m.
\]

另一方面，附录定理 \ref{thm:derived-fiber-indcomplex-alternating-witten-parity} 给出
\[
W_x=
\begin{cases}
0,& d_m(x)\ \text{为偶},\\[1mm]
(-1)^{\tau(x)},& d_m(x)\ \text{为奇}.
\end{cases}
\]
故当 \(m=6,12\) 时，\(d_{\min}(m)=2,8\) 为偶，于是对每个
\[
x\in S_{m,\min}
\]
都有
\[
W_x=0.
\]
因此
\[
\mathcal W_{\min}(m)=0.
\]

当 \(m=8,10\) 时，\(d_{\min}(m)=3,5\) 为奇。又由定理 \ref{thm:xi-audited-even-minsector-homotopy-mod6-phase} 的证明，在这两种情形中都必有
\[
\tau(x)=1
\qquad (x\in S_{m,\min}),
\]
因为
\[
3=F_4,\qquad 5=F_5
\]
分别只对应单一非平凡 Fibonacci 因子。于是
\[
W_x=(-1)^{\tau(x)}=-1
\qquad (x\in S_{m,\min}).
\]
故
\[
\mathcal W_{\min}(m)= -|S_{m,\min}|=-F_m.
\]
最后再除以 \(|S_{m,\min}|=F_m\)，便得到平均电荷公式。证毕。
\end{proof}

\begin{corollary}[最小退化扇区的非平凡拓扑质量密度在审计偶窗上仅取两值 \(0\) 与 \(1\)]\label{cor:xi-time-part64ba-audited-even-minsector-topological-mass-density}
定义
\[
\mathfrak m_{\min}(m)
:=
\frac{1}{|S_{m,\min}|}\sum_{x\in S_{m,\min}}|W_x|.
\]
则对审计偶窗 \(m\in\{6,8,10,12\}\) 有
\[
\boxed{\;
\mathfrak m_{\min}(m)=
\begin{cases}
0,& m=6,12,\\[1mm]
1,& m=8,10.
\end{cases}\;}
\]
换言之，在这些审计偶窗上，最小退化扇区要么整体拓扑平凡，要么每条纤维都携带单位绝对 Witten 电荷；不存在混合相。
\end{corollary}
\begin{proof}
由附录定理 \ref{thm:derived-fiber-indcomplex-alternating-witten-parity}，
\[
|W_x|=
\begin{cases}
0,& d_m(x)\ \text{为偶},\\[1mm]
1,& d_m(x)\ \text{为奇}.
\end{cases}
\]
再结合定理 \ref{thm:xi-time-part64ba-audited-even-minsector-total-witten-charge} 中对四个审计偶窗的逐窗判定，便立即得到所示两值结论。证毕。
\end{proof}

\noindent
由此，临界 oracle 曲线、幂散度谱、折叠多重度剖面与最小退化扇区的拓扑相位不再是彼此平行的局部现象：一方面，临界容量双折线正是极限似然比的 stop-loss 影像，而整个 power-divergence 族又由单个 \(\chi^2\)-型极限常数刚性锁定；另一方面，Robin--Hood 调平把任意多重度剖面推向唯一平衡端点，并使这一平衡端点同时成为全阶 Rényi 熵的统一上包络；最后，审计偶窗上的最小退化扇区不只在同伦型上区分为可缩相与球面相，而且在交错 Witten 电荷层面进一步裂解为零电荷相与满密度单位电荷相。于是，统计 stop-loss、二阶信息坐标、双随机平衡化与拓扑电荷跃迁最终被压缩进同一条可审计主线。

\paragraph{双分歧场、同余三次专化与 Lee--Yang/\texorpdfstring{$19$}{19} 次层的乘法化 Frobenius 统计}
前述链条已经分别建立了十次分歧域 \(L_{10}\)、九次分歧域 \(L_9\)、同余三次专化域 \(M_t\)、Lee--Yang 三次域与十九次指数域 \(L_H\) 的满对称伽罗瓦性及其两两独立性。
下面把这些分散的算术接口进一步压缩为两条互不重叠而可交叉编译的统计主轴：一条是
\[
S_{10}\times S_9\times S_3
\]
上的三重分解型与符号盒规律，另一条是
\[
S_{10}\times S_3\times S_{19}
\]
上的 Lee--Yang 根出现事件与十九次固定点统计的完全张量化。

\begin{theorem}[双分歧场与同余三次专化场的三重直积伽罗瓦群]\label{thm:xi-time-part64bb-p10-p9-trigonal-threefold-galois-independence}
沿用定理 \ref{thm:fold-gauge-anomaly-P10-P9-linear-disjointness} 与定理
\ref{thm:fold-gauge-anomaly-second-trigonal-ht-specialization-s3-disjoint} 的记号。
设 \(L_{10}/\QQ\) 与 \(L_9/\QQ\) 分别为十次分歧多项式 \(P_{10}\) 与九次分歧多项式 \(P_9\) 的分裂域。
对整数
\[
t\ge 5,
\qquad
t\not\equiv 1\pmod 3,
\]
令 \(h_t(u)\in\QQ[u]\) 为三次专化多项式，其分裂域记为 \(M_t\)。
则
\[
M_t\cap L_{10}L_9=\QQ,
\qquad
\Gal(L_{10}L_9M_t/\QQ)\cong S_{10}\times S_9\times S_3.
\]
特别地，
\[
[L_{10}L_9M_t:\QQ]=10!\cdot 9!\cdot 6.
\]
\end{theorem}
\begin{proof}
由定理 \ref{thm:fold-gauge-anomaly-P10-P9-linear-disjointness}，
\[
L_{10}\cap L_9=\QQ,
\qquad
\Gal(L_{10}L_9/\QQ)\cong S_{10}\times S_9.
\]
而由定理 \ref{thm:fold-gauge-anomaly-second-trigonal-ht-specialization-s3-disjoint}，
\[
\Gal(M_t/\QQ)\cong S_3,
\qquad
M_t\cap L_{10}=\QQ,
\qquad
M_t\cap L_9=\QQ.
\]
故只须排除 \(M_t\) 与合成域 \(L_{10}L_9\) 之间出现新的公共正规子域。

记
\[
E:=M_t\cap L_{10}L_9.
\]
由于两边都是 \(\QQ\) 上伽罗瓦扩张，故 \(E/\QQ\) 亦为伽罗瓦扩张。
又因为 \(\Gal(M_t/\QQ)\cong S_3\)，\(M_t\) 的唯一非平凡真正规子域只能是其判别式二次域
\[
K_t:=\QQ\!\left(\sqrt{\Disc_u(h_t)}\right).
\]
同一定理已给出
\[
\Disc_u(h_t)=-(t-1)P_9(t)<0,
\]
故 \(K_t\) 为虚二次域。

另一方面，
\[
\Gal(L_{10}L_9/\QQ)\cong S_{10}\times S_9
\]
的交换化为
\[
(S_{10}\times S_9)^{\mathrm{ab}}\cong C_2\times C_2.
\]
因此 \(L_{10}L_9\) 的二次正规子域恰对应三条非平凡符号特征，分别为
\[
\QQ(\sqrt{66269989}),
\qquad
\QQ(\sqrt{33605193}),
\qquad
\QQ(\sqrt{66269989\cdot 33605193}),
\]
它们全都是实二次域。
于是 \(K_t\) 不可能嵌入 \(L_{10}L_9\)，从而 \(E\neq K_t\)。
故只能 \(E=\QQ\)。

交域平凡且两边皆为伽罗瓦扩张，遂得到线性互素与直积型伽罗瓦群同构
\[
\Gal(L_{10}L_9M_t/\QQ)\cong \Gal(L_{10}L_9/\QQ)\times \Gal(M_t/\QQ)
\cong S_{10}\times S_9\times S_3.
\]
次数公式随之立即成立。证毕。
\end{proof}

\begin{corollary}[三重分解型的完全乘法密度与 \texorpdfstring{$1/270$}{1/270} 定律]\label{cor:xi-time-part64bb-p10-p9-trigonal-chebotarev-1over270}
在定理 \ref{thm:xi-time-part64bb-p10-p9-trigonal-threefold-galois-independence} 的设定下，
对任意分拆
\[
\lambda\vdash 10,
\qquad
\mu\vdash 9,
\qquad
\nu\vdash 3,
\]
记 \(C_\lambda\subset S_{10}\)、\(C_\mu\subset S_9\)、\(C_\nu\subset S_3\) 为对应循环型共轭类。
则除去有限个分歧素数后，满足
\[
\Frob_p\in C_\lambda\times C_\mu\times C_\nu
\]
的素数集合具有自然密度
\[
\delta(\lambda,\mu,\nu)
=
\frac{|C_\lambda|}{10!}\cdot
\frac{|C_\mu|}{9!}\cdot
\frac{|C_\nu|}{6}.
\]
特别地，同时满足
\[
P_{10}\bmod p\ \text{不可约},
\qquad
P_9\bmod p\ \text{不可约},
\qquad
h_t\bmod p\ \text{不可约}
\]
的素数集合密度为
\[
\frac{1}{10}\cdot\frac{1}{9}\cdot\frac{1}{3}
=
\frac{1}{270},
\]
而三者同时全分裂的密度为
\[
\frac{1}{10!}\cdot \frac{1}{9!}\cdot \frac{1}{6}.
\]
\end{corollary}
\begin{proof}
由定理 \ref{thm:xi-time-part64bb-p10-p9-trigonal-threefold-galois-independence}，
\[
\Gal(L_{10}L_9M_t/\QQ)\cong S_{10}\times S_9\times S_3.
\]
对直积群应用 Chebotarev 密度定理，便得到一般共轭类的乘法密度公式。

对特例，只需注意：\(n\) 次多项式模 \(p\) 不可约对应 \(S_n\) 中的 \(n\)-循环，其比例为 \(1/n\)；全分裂对应恒等元类，其比例为 \(1/n!\)。
而在 \(S_3\) 中，不可约对应 \(3\)-循环类，大小为 \(2\)，故比例为 \(2/6=1/3\)。
代入即可得到 \(1/270\) 与全分裂密度公式。证毕。
\end{proof}

\begin{theorem}[三重判别特征立方体的等分布与三次分裂律的条件不变性]\label{thm:xi-time-part64bb-p10-p9-trigonal-signcube-equidistribution}
在定理 \ref{thm:xi-time-part64bb-p10-p9-trigonal-threefold-galois-independence} 的设定下，对不分歧素数 \(p\) 定义
\[
\chi_{10}(p):=\sgn(\sigma_{10}(p))\in\{\pm 1\},
\qquad
\chi_9(p):=\sgn(\sigma_9(p))\in\{\pm 1\},
\qquad
\chi_t(p):=\sgn(\sigma_t(p))\in\{\pm 1\},
\]
其中
\[
\sigma_{10}(p)\in S_{10},
\qquad
\sigma_9(p)\in S_9,
\qquad
\sigma_t(p)\in S_3
\]
分别表示相应分裂域上的 Frobenius 元。
则有：
\begin{enumerate}
  \item 三元符号向量
  \[
  \bigl(\chi_{10}(p),\chi_9(p),\chi_t(p)\bigr)
  \]
  在 \(\{\pm1\}^3\) 上完全等分布；每一符号模式的自然密度均为
  \[
  \frac{1}{8}.
  \]
  \item 对任意给定的双符号条件
  \[
  \chi_{10}(p)=\varepsilon_1,
  \qquad
  \chi_9(p)=\varepsilon_2
  \qquad
  (\varepsilon_1,\varepsilon_2\in\{\pm1\}),
  \]
  三次专化 \(h_t\bmod p\) 的三种分解型条件密度保持不变：
  \[
  \delta\bigl((1)(1)(1)\mid \chi_{10}=\varepsilon_1,\chi_9=\varepsilon_2\bigr)=\frac16,
  \]
  \[
  \delta\bigl((1)(2)\mid \chi_{10}=\varepsilon_1,\chi_9=\varepsilon_2\bigr)=\frac12,
  \]
  \[
  \delta\bigl((3)\mid \chi_{10}=\varepsilon_1,\chi_9=\varepsilon_2\bigr)=\frac13.
  \]
\end{enumerate}
换言之，十次分歧通道与九次分歧通道的全部二次判别偏置，对三次专化层的局部分裂统计完全透明。
\end{theorem}
\begin{proof}
由定理 \ref{thm:xi-time-part64bb-p10-p9-trigonal-threefold-galois-independence}，
\[
\Gal(L_{10}L_9M_t/\QQ)\cong S_{10}\times S_9\times S_3.
\]
对三个因子分别取符号特征，得到满射
\[
S_{10}\times S_9\times S_3\twoheadrightarrow C_2\times C_2\times C_2.
\]
Chebotarev 定理在该阿贝尔商上给出八个符号盒的等分布，故每种符号模式的密度都是 \(1/8\)。

再看条件分布。
给定 \((\chi_{10},\chi_9)\) 只会约束 \(S_{10}\times S_9\) 坐标，而 \(h_t\bmod p\) 的分解型仅由 \(S_3\) 坐标决定。
由于直积群上的 Frobenius 分布完全乘法化，条件化后 \(S_3\) 的共轭类比例不变。
而 \(S_3\) 的三个共轭类大小分别为 \(1,3,2\)，对应分解型
\[
(1)(1)(1),\qquad (1)(2),\qquad (3),
\]
故条件密度分别为
\[
\frac16,\qquad \frac12,\qquad \frac13.
\]
证毕。
\end{proof}

\begin{theorem}[十九次固定点统计与 \texorpdfstring{$(10,3)$}{(10,3)} 分裂指纹的三向独立]\label{thm:xi-time-part64bb-p10-leyang-h19-fixedpoint-tensorization}
沿用定理 \ref{thm:fold-gauge-anomaly-P10-leyang-H-triple-product} 的记号。
令 \(L_H/\QQ\) 为十九次指数域对应的 \(S_{19}\)-扩张，并对不分歧素数 \(p\) 记事件
\[
E_r(p):=\text{Frobenius 在 \(19\) 个根上恰有 \(r\) 个不动点}
\qquad
(0\le r\le 19),
\]
等价地，十九次多项式 \(H\bmod p\) 恰有 \(r\) 个一次因子。
再定义
\[
I_{10}(p):=\{\text{\(P_{10}\bmod p\) 不可约}\},
\]
\[
R(p):=\{\text{\(g(y)=256y^3+411y^2+165y+32\bmod p\) 至少有一根}\},
\]
\[
S(p):=\{\text{\(g\bmod p\) 全分裂}\}.
\]
则有精确密度公式
\[
\delta\bigl(I_{10}\cap R\cap E_r\bigr)
=
\frac1{10}\cdot \frac23\cdot \frac{\binom{19}{r}\,!(19-r)}{19!},
\]
以及
\[
\delta\bigl(I_{10}\cap S\cap E_r\bigr)
=
\frac1{10}\cdot \frac16\cdot \frac{\binom{19}{r}\,!(19-r)}{19!}.
\]
这表明：十次分歧的不可约性、Lee--Yang 三次层的根出现事件、以及十九次指数层的一次因子数，在自然密度意义下完全张量化。
\end{theorem}
\begin{proof}
由定理 \ref{thm:fold-gauge-anomaly-P10-leyang-H-triple-product}，
\[
\Gal(L_{10}L_{\mathrm{LY}}L_H/\QQ)\cong S_{10}\times S_3\times S_{19},
\]
而推论 \ref{cor:fold-gauge-anomaly-P10-leyang-H-chebotarev-triple-product} 给出三因子的 Frobenius 类分布完全乘法化。

另一方面，由命题 \ref{prop:fold-gauge-anomaly-H-rencontres-fixed-point-density}，
\[
\delta(E_r)=\frac{\binom{19}{r}\,!(19-r)}{19!}.
\]
在 \(S_{10}\) 中，\(I_{10}\) 对应 \(10\)-循环类，故
\[
\delta(I_{10})=\frac1{10}.
\]
在 \(S_3\) 中，“至少有一根”对应恒等类与换位类之并，故
\[
\delta(R)=\frac16+\frac12=\frac23,
\]
而“全分裂”对应恒等类，故
\[
\delta(S)=\frac16.
\]
将这三个边际密度相乘，即得
\[
\delta\bigl(I_{10}\cap R\cap E_r\bigr)
=
\delta(I_{10})\,\delta(R)\,\delta(E_r)
=
\frac1{10}\cdot \frac23\cdot \frac{\binom{19}{r}\,!(19-r)}{19!},
\]
以及
\[
\delta\bigl(I_{10}\cap S\cap E_r\bigr)
=
\delta(I_{10})\,\delta(S)\,\delta(E_r)
=
\frac1{10}\cdot \frac16\cdot \frac{\binom{19}{r}\,!(19-r)}{19!}.
\]
证毕。
\end{proof}


\paragraph{bin-fold 满秩化、原子有限部位移与碰撞势内域大偏差}
在 bin-fold 两态一致渐近、规范群 Stirling 账本、真实输入 \(40\) 状态核的原子 Euler 剥离以及碰撞位势的解析/大偏差链已经建立之后，还可进一步抽取出下列十三条直接后果。它们分别刻画 fold-gauge 直积分解的中心与阿贝尔化精确结构、阿贝尔化的最终满秩化、中心的最终消失与纤维奇偶荷最小完备秩的严格稳定、规范复杂度相对异常维数的指数超载、体积密度对平均纤维信息项的线性响应斜率、纯碰撞误差指数缺口的 \(\ell^p\) 临界幂、原子因子的 primitive 与 Abel 有限部刚性位移、平方根归一化慢模的三歧相变与 \(u=1\) 邻域的统一亚临界性、RH-sharp 临界点上一阶滤子的精确消模、误差指数 \(\beta=\tfrac12\) 的余维一尖点接触，以及 RH 失效在 Taylor 解析域内部的可见性与均值点速率函数的显式二次胚。

\begin{theorem}[fold-gauge 直积分解的中心与阿贝尔化精确结构]\label{thm:xi-time-part65-binfold-gauge-center-abelianization-exact}
沿用命题 \ref{prop:fold-bin-gauge-decomposition} 的记号，记
\[
n_2(m):=\#\{w\in X_m:\ d_m^{\mathrm{bin}}(w)=2\},
\qquad
n_{\ge 2}(m):=\#\{w\in X_m:\ d_m^{\mathrm{bin}}(w)\ge 2\}.
\]
则有同构
\[
Z\!\bigl(G_m^{\mathrm{bin}}\bigr)\cong (\ZZ/2\ZZ)^{n_2(m)},
\qquad
\bigl(G_m^{\mathrm{bin}}\bigr)^{\mathrm{ab}}\cong (\ZZ/2\ZZ)^{n_{\ge 2}(m)}.
\]
特别地，每一条二点纤维都贡献一个中心 sheet-flip，而所有 \(d_m^{\mathrm{bin}}(w)\ge 3\) 的对称群因子都具有平凡中心。
\end{theorem}
\begin{proof}
由命题 \ref{prop:fold-bin-gauge-decomposition}，
\[
G_m^{\mathrm{bin}}
\cong
\prod_{w\in X_m}\mathfrak S_{d_m^{\mathrm{bin}}(w)}.
\]
对称群满足
\[
Z(\mathfrak S_1)=1,
\qquad
Z(\mathfrak S_2)=\mathfrak S_2\cong \ZZ/2\ZZ,
\qquad
Z(\mathfrak S_d)=1\quad(d\ge 3),
\]
以及
\[
\mathfrak S_1^{\mathrm{ab}}=1,
\qquad
\mathfrak S_d^{\mathrm{ab}}\cong \ZZ/2\ZZ\quad(d\ge 2).
\]
再由直积群的中心与阿贝尔化分别按因子直积分解，立得
\[
Z\!\bigl(G_m^{\mathrm{bin}}\bigr)\cong (\ZZ/2\ZZ)^{n_2(m)},
\qquad
\bigl(G_m^{\mathrm{bin}}\bigr)^{\mathrm{ab}}\cong (\ZZ/2\ZZ)^{n_{\ge 2}(m)}.
\]
证毕。
\end{proof}
\leanverified{paper\_xi\_time\_part65\_binfold\_gauge\_center\_abelianization\_exact}

\begin{theorem}[bin-fold 规范群阿贝尔化的最终满秩化]\label{thm:xi-time-part65-binfold-abel-eventual-full-rank}
沿用定理 \ref{thm:xi-time-part65-binfold-gauge-center-abelianization-exact} 与
\ref{thm:fold-bin-two-state-asymptotic} 的记号，并记
\[
N_m^{(2)}:=\#\{w\in X_m:\ d_m^{\mathrm{bin}}(w)\ge 2\}.
\]
则存在 \(m_0\)，使得对所有 \(m\ge m_0\) 都有
\[
\boxed{\;
\bigl(G_m^{\mathrm{bin}}\bigr)^{\mathrm{ab}}
\cong
(\ZZ/2\ZZ)^{F_{m+2}}.\;}
\]
等价地，bin-fold 规范群的全部稳定类型在充分高分辨率下都产生彼此独立的 \((\ZZ/2\ZZ)\)-奇偶通道，而阿贝尔化秩最终不再遗漏任何稳定类型。
\end{theorem}
\begin{proof}
由定理 \ref{thm:fold-bin-two-state-asymptotic} 的一致对数展开，
\[
\log d_m^{\mathrm{bin}}(w)
=
m\log\!\Bigl(\frac{2}{\varphi}\Bigr)-w_m\log\varphi
+O\!\Bigl(\Bigl(\frac{\varphi}{2}\Bigr)^m\Bigr)
\qquad (w\in X_m).
\]
由于 \(w_m\in\{0,1\}\)，故对全部 \(w\in X_m\) 一致有
\[
\log d_m^{\mathrm{bin}}(w)
\ge
m\log\!\Bigl(\frac{2}{\varphi}\Bigr)-\log\varphi
+O\!\Bigl(\Bigl(\frac{\varphi}{2}\Bigr)^m\Bigr).
\]
而 \(\log(2/\varphi)>0\)，故右端随 \(m\to\infty\) 一致趋于 \(+\infty\)。于是存在 \(m_0\) 使当 \(m\ge m_0\) 时，对所有 \(w\in X_m\) 都有 \(d_m^{\mathrm{bin}}(w)>1\)。由于 \(d_m^{\mathrm{bin}}(w)\) 为正整数，遂得
\[
d_m^{\mathrm{bin}}(w)\ge 2
\qquad(\forall\,w\in X_m,\ m\ge m_0),
\]
即
\[
N_m^{(2)}=|X_m|=F_{m+2}.
\]
由定理 \ref{thm:xi-time-part65-binfold-gauge-center-abelianization-exact}，
\[
\bigl(G_m^{\mathrm{bin}}\bigr)^{\mathrm{ab}}
\cong
(\ZZ/2\ZZ)^{n_{\ge 2}(m)}.
\]
再代入
\[
n_{\ge 2}(m)=N_m^{(2)}=F_{m+2},
\]
即得结论。证毕。
\end{proof}

\begin{theorem}[bin-fold 规范群中心的最终消失与纤维奇偶荷最小完备秩的严格稳定]\label{thm:xi-time-part65-binfold-center-vanishing-signature-stability}
沿用定理 \ref{thm:xi-time-part65-binfold-gauge-center-abelianization-exact} 与
\ref{thm:xi-time-part65-binfold-abel-eventual-full-rank} 的记号。则存在 \(m_0\)，使得对所有 \(m\ge m_0\) 都有
\[
\boxed{\;
Z\!\bigl(G_m^{\mathrm{bin}}\bigr)=1,
\qquad
\bigl(G_m^{\mathrm{bin}}\bigr)^{\mathrm{ab}}
\cong
(\ZZ/2\ZZ)^{F_{m+2}}.\;}
\]
从而
\[
\boxed{\;
\operatorname{rank}_{\FF_2}\!\bigl((G_m^{\mathrm{bin}})^{\mathrm{ab}}\bigr)=F_{m+2}.\;}
\]
并且定理 \ref{thm:xi-foldbin-parity-charge-split-exact-minrank} 中的纤维奇偶荷最小完备秩在同一范围内亦严格等于 \(F_{m+2}\)。
\end{theorem}
\begin{proof}
由定理 \ref{thm:fold-bin-two-state-asymptotic} 的一致两态展开，存在常数 \(C>0\) 使对充分大的 \(m\) 与任意 \(w\in X_m\) 都有
\[
d_m^{\mathrm{bin}}(w)
\ge
\left(\frac{2}{\varphi}\right)^m
\left(\varphi^{-1}-C\left(\frac{\varphi}{2}\right)^m\right).
\]
由于
\[
\left(\frac{2}{\varphi}\right)^m
\left(\varphi^{-1}-C\left(\frac{\varphi}{2}\right)^m\right)
=
\varphi^{-1}\left(\frac{2}{\varphi}\right)^m-C
\longrightarrow +\infty,
\]
故存在 \(m_0\)，使得当 \(m\ge m_0\) 时，对所有 \(w\in X_m\) 都有
\[
d_m^{\mathrm{bin}}(w)\ge 3.
\]

于是
\[
n_2(m)=0,
\qquad
n_{\ge 2}(m)=|X_m|=F_{m+2}.
\]
由定理 \ref{thm:xi-time-part65-binfold-gauge-center-abelianization-exact}，
\[
Z\!\bigl(G_m^{\mathrm{bin}}\bigr)\cong (\ZZ/2\ZZ)^{n_2(m)}=1,
\]
\[
\bigl(G_m^{\mathrm{bin}}\bigr)^{\mathrm{ab}}
\cong
(\ZZ/2\ZZ)^{n_{\ge 2}(m)}
\cong
(\ZZ/2\ZZ)^{F_{m+2}}.
\]
这就给出了中心消失与交换化维数公式。

最后，由定理 \ref{thm:xi-foldbin-parity-charge-split-exact-minrank}，纤维奇偶荷的最小完备秩恰等于
\[
N_m:=\#\{w\in X_m:\ d_m^{\mathrm{bin}}(w)\ge 2\}.
\]
而此处对所有 \(w\in X_m\) 已有 \(d_m^{\mathrm{bin}}(w)\ge 3\)，故
\[
N_m=|X_m|=F_{m+2}.
\]
于是最小完备秩在同一范围内严格稳定为 \(F_{m+2}\)。证毕。
\end{proof}

\begin{theorem}[规范群复杂度相对异常维数的指数超载律]\label{thm:xi-time-part65-gauge-complexity-anomaly-overload}
沿用定理 \ref{thm:xi-time-part65-binfold-abel-eventual-full-rank} 与定理 \ref{thm:xi-foldbin-gauge-density-exact-thermodynamic-constant} 的记号。则当 \(m\to\infty\) 时，
\[
\boxed{\;
\frac{\log_2|G_m^{\mathrm{bin}}|}
{\dim_{\FF_2}\!\bigl((G_m^{\mathrm{bin}})^{\mathrm{ab}}\bigr)}
=
\frac{\sqrt5}{\varphi^2\log 2}
\left(\frac{2}{\varphi}\right)^m
\left(
m\log\!\Bigl(\frac{2}{\varphi}\Bigr)
-1-\frac{\log\varphi}{1+\varphi^2}
+o(1)
\right).\;}
\]
特别地，
\[
\boxed{\;
\frac{\log_2|G_m^{\mathrm{bin}}|}
{\dim_{\FF_2}\!\bigl((G_m^{\mathrm{bin}})^{\mathrm{ab}}\bigr)}
\longrightarrow +\infty,\;}
\]
并且其主导增长尺度恰为 \((2/\varphi)^m\)。
\end{theorem}
\begin{proof}
由定理 \ref{thm:xi-foldbin-gauge-density-exact-thermodynamic-constant}，
\[
g_m
=
m\log\!\Bigl(\frac{2}{\varphi}\Bigr)
-1-\frac{\log\varphi}{1+\varphi^2}
+O\!\Bigl(m\Bigl(\frac{\varphi}{2}\Bigr)^m\Bigr).
\]
由于 \(g_m=2^{-m}\log|G_m^{\mathrm{bin}}|\)，故
\[
\log_2|G_m^{\mathrm{bin}}|
=
\frac{2^m}{\log 2}
\left(
m\log\!\Bigl(\frac{2}{\varphi}\Bigr)
-1-\frac{\log\varphi}{1+\varphi^2}
+o(1)
\right).
\]

另一方面，由定理 \ref{thm:xi-time-part65-binfold-abel-eventual-full-rank}，对充分大的 \(m\) 有
\[
\dim_{\FF_2}\!\bigl((G_m^{\mathrm{bin}})^{\mathrm{ab}}\bigr)=F_{m+2}.
\]
再用 Binet 公式
\[
F_{m+2}=\frac{\varphi^{m+2}}{\sqrt5}+o(\varphi^m),
\]
便得到
\[
\frac{\log_2|G_m^{\mathrm{bin}}|}
{\dim_{\FF_2}\!\bigl((G_m^{\mathrm{bin}})^{\mathrm{ab}}\bigr)}
=
\frac{\sqrt5}{\varphi^2\log 2}
\left(\frac{2}{\varphi}\right)^m
\left(
m\log\!\Bigl(\frac{2}{\varphi}\Bigr)
-1-\frac{\log\varphi}{1+\varphi^2}
+o(1)
\right).
\]
由于 \(\log(2/\varphi)>0\) 且 \(2/\varphi>1\)，右端趋于 \(+\infty\)，并且主导增长尺度正是 \((2/\varphi)^m\)。证毕。
\end{proof}

\begin{theorem}[规范体积密度与纤维信息项之间的线性响应斜率]\label{thm:xi-time-part65-binfold-gauge-linear-response-slope}
沿用命题 \ref{prop:fold-bin-gauge-density-two-scale} 的记号。则
\[
\boxed{\;
\lim_{m\to\infty}
\frac{g_m-\kappa_m+1}{m(\varphi/2)^m}
=
\frac{\varphi^2}{2\sqrt5}\log\!\Bigl(\frac{2}{\varphi}\Bigr).\;}
\]
并且利用 \(|X_m|=F_{m+2}\sim \varphi^{m+2}/\sqrt5\) 可进一步改写为
\[
\boxed{\;
\frac{2^m}{|X_m|}\,\bigl(g_m-\kappa_m+1\bigr)
=
\frac12\,m\log\!\Bigl(\frac{2}{\varphi}\Bigr)+O(1).\;}
\]
特别地，对充分大的 \(m\) 有
\[
g_m>\kappa_m-1.
\]
\end{theorem}
\begin{proof}
由命题 \ref{prop:fold-bin-gauge-density-two-scale}，
\[
g_m-\kappa_m+1
=
\frac{(\varphi/2)^m}{2\sqrt5}
\Bigl(
\varphi^2 m\log\!\Bigl(\frac{2}{\varphi}\Bigr)
+\varphi^2\log(2\pi)-\log\varphi
\Bigr)
+O\!\Bigl(\Bigl(\frac{\varphi}{2}\Bigr)^m\Bigr).
\]
两边同除以 \(m(\varphi/2)^m\)，得到
\[
\frac{g_m-\kappa_m+1}{m(\varphi/2)^m}
=
\frac{1}{2\sqrt5}
\left(
\varphi^2\log\!\Bigl(\frac{2}{\varphi}\Bigr)
+\frac{\varphi^2\log(2\pi)-\log\varphi}{m}
\right)
+O\!\left(\frac1m\right),
\]
从而首个极限成立。

另一方面，由 Binet 公式，
\[
\frac{|X_m|}{2^m}
=
\frac{F_{m+2}}{2^m}
=
\frac{\varphi^2}{\sqrt5}\Bigl(\frac{\varphi}{2}\Bigr)^m\bigl(1+o(1)\bigr),
\]
故
\[
\frac{2^m}{|X_m|}
=
\frac{\sqrt5}{\varphi^2}\Bigl(\frac{2}{\varphi}\Bigr)^m\bigl(1+o(1)\bigr).
\]
将其与上式相乘，便得
\[
\frac{2^m}{|X_m|}\,\bigl(g_m-\kappa_m+1\bigr)
=
\frac12\,m\log\!\Bigl(\frac{2}{\varphi}\Bigr)+O(1).
\]
由于 \(\log(2/\varphi)>0\)，右端对充分大 \(m\) 为正，因此 \(g_m-\kappa_m+1>0\)，即
\(
g_m>\kappa_m-1
\)。
证毕。
\end{proof}

\begin{theorem}[纯碰撞误差指数缺口的 \texorpdfstring{$\ell^p$}{lp} 临界幂为 \texorpdfstring{$p=2$}{p=2}]\label{thm:xi-time-part65-collision-error-gap-lp-threshold}
沿用推论 \ref{cor:real-input-40-collision-error-phase} 的记号，并把误差指数 \(\beta(u)\) 限制在整数参数 \(u=n\ge 3\) 上。则对任意 \(p>0\)，级数
\[
\sum_{n\ge 3}\bigl(1-\beta(n)\bigr)^p
\]
满足精确判别
\[
\boxed{\;
\sum_{n\ge 3}\bigl(1-\beta(n)\bigr)^p
\begin{cases}
<\infty,& p\ge 2,\\[2pt]
=\infty,& 0<p<2.
\end{cases}}
\]
换言之，临界幂指数恰为 \(p=2\)，并且边界点的额外对数增益已经足以恢复可和性。
\end{theorem}
\begin{proof}
由推论 \ref{cor:real-input-40-collision-error-phase}，
\[
\beta(u)
=
1-\frac{2}{\sqrt u\,\log u}
+O\!\left(\frac{1}{u\log u}\right)
\qquad (u\to\infty).
\]
因此存在常数 \(A,B>0\) 及充分大的 \(n_0\)，使得对所有 \(n\ge n_0\)，
\[
\frac{A}{\sqrt n\,\log n}
\le
1-\beta(n)
\le
\frac{B}{\sqrt n\,\log n}.
\]
从而
\[
\bigl(1-\beta(n)\bigr)^p
\asymp
\frac{1}{n^{p/2}(\log n)^p}.
\]
于是由积分判别或 Cauchy 凝聚判别，
\[
\sum_{n\ge n_0}\frac{1}{n^{p/2}(\log n)^p}
\]
在 \(p>2\) 时显然收敛，在 \(0<p<2\) 时显然发散；而在边界点 \(p=2\) 时，它等价于
\[
\sum_{n\ge n_0}\frac{1}{n(\log n)^2},
\]
该级数收敛。故所示判别成立。证毕。
\end{proof}

\begin{theorem}[原子 Euler 因子的 primitive 单点刚性与 Abel 有限部的精确乘法位移]\label{thm:xi-time-part65-atomic-primitive-abel-shift}
沿用命题 \ref{prop:atomic-2-orbit-split-logM} 与定理
\ref{thm:xi-atomic-witt-cycle-primitive-exact-peeling} 的记号。写
\[
P(z,u)=\sum_{n\ge 1}p_n(u)z^n,
\qquad
P_{\mathrm{core}}(z,u)=\sum_{n\ge 1}p_n^{\mathrm{core}}(u)z^n.
\]
则对一切 \(n\ge 1\) 都有
\[
\boxed{\;
p_n(u)=p_n^{\mathrm{core}}(u)+u\,\mathbf 1_{\{n=2\}}.\;}
\]
特别地，在未加权情形 \(u=1\) 下，
\[
\boxed{\;
p_n=p_n^{\mathrm{core}}+\mathbf 1_{\{n=2\}}.\;}
\]
进一步，若 \(\mathfrak M\) 与 \(\mathfrak M_{\mathrm{core}}\) 分别表示 \(u=1\) 时 full/core 的 Abel 有限部常数，则
\[
\boxed{\;
\log\mathfrak M
=
\log\mathfrak M_{\mathrm{core}}+\varphi^{-4},\qquad
\mathfrak M=e^{\varphi^{-4}}\mathfrak M_{\mathrm{core}}.\;}
\]
因此，原子因子 \((1-uz^2)^{-1}\) 在 primitive 层只作用于长度 \(2\) 的单一坐标，而在 Abel 有限部上只留下一个完全闭式的乘法位移。
\end{theorem}
\begin{proof}
由定理 \ref{thm:xi-atomic-witt-cycle-primitive-exact-peeling}，
\[
P(z,u)=u z^2+P_{\mathrm{core}}(z,u).
\]
逐项比较 \(z^n\) 的系数，便得
\[
p_n(u)=p_n^{\mathrm{core}}(u)+u\,\mathbf 1_{\{n=2\}},
\]
从而在 \(u=1\) 时有
\(
p_n=p_n^{\mathrm{core}}+\mathbf 1_{\{n=2\}}
\)。

另一方面，命题 \ref{prop:atomic-2-orbit-split-logM} 已给出
\[
\log\mathfrak M=\log\mathfrak M_{\mathrm{core}}+z_\star^2.
\]
在本核的未加权情形下，\(z_\star=\varphi^{-2}\)，故
\[
z_\star^2=\varphi^{-4}.
\]
将其代回并指数化，便得
\[
\mathfrak M=e^{\varphi^{-4}}\mathfrak M_{\mathrm{core}}.
\]
证毕。
\end{proof}

\begin{theorem}[碰撞位势的 RH 失效在原点 Taylor 解析域内可直接见证]\label{thm:xi-time-part65-collision-taylor-visible-nonrh-window}
沿用定理 \ref{thm:xi-rh-critical-before-nearest-complex-branch} 的记号。则有严格包含
\[
\boxed{\;
(\theta_R,R_\theta)
\subset
\Bigl\{
\theta\in\RR_{>0}:\ 
P_2(\theta)=\sum_{n\ge 0}\frac{P_2^{(n)}(0)}{n!}\theta^n,
\ \textsf{RH}(e^\theta)\ \text{失效}
\Bigr\}.\;}
\]
更一般地，对任意紧区间 \(K\subset(\theta_R,R_\theta)\)，上述 Taylor 级数都在 \(K\) 上绝对且一致收敛，而 kernel RH 条件在 \(K\) 上处处失效。
\end{theorem}
\begin{proof}
由定理 \ref{thm:xi-rh-critical-before-nearest-complex-branch}，
\[
\theta_R<R_\theta
\]
且对每个 \(\theta\in(\theta_R,R_\theta)\)，一方面 \(|\theta|<R_\theta\)，另一方面 \(u=e^\theta>u_R\)，故 \(\textsf{RH}(e^\theta)\) 已失效。

由同一定理，\(P_2\) 在 \(\theta=0\) 的 Taylor 收敛半径正是 \(R_\theta\)。因此对任意 \(|\theta|<R_\theta\)，都有
\[
P_2(\theta)=\sum_{n\ge 0}\frac{P_2^{(n)}(0)}{n!}\theta^n.
\]
特别地，这对全部 \(\theta\in(\theta_R,R_\theta)\) 成立，从而得到所示包含关系。

若 \(K\subset(\theta_R,R_\theta)\) 为紧区间，则可取 \(r\) 使
\[
\sup_{\theta\in K}|\theta|<r<R_\theta.
\]
于是幂级数在闭区间 \(K\) 上由半径 \(r\) 的 Cauchy 估计控制，从而绝对且一致收敛。与此同时，由 \(K\subset(\theta_R,R_\theta)\) 仍知 \(\textsf{RH}(e^\theta)\) 在 \(K\) 上处处失效。证毕。
\end{proof}

\begin{theorem}[RH 窗口内 Taylor 截断深度的 Cauchy 指数界]\label{thm:xi-time-part65-collision-taylor-depth-exp-bound}
沿用定理 \ref{thm:xi-rh-critical-before-nearest-complex-branch} 的记号。对任意
\[
\theta_R<r<R_\theta,
\]
记
\[
M(r):=\sup_{|\theta|=r}|P_2(\theta)|.
\]
若
\[
T_N(\theta):=\sum_{n=0}^N\frac{P_2^{(n)}(0)}{n!}\theta^n
\]
为 \(P_2\) 在 \(\theta=0\) 处的 \(N\) 阶 Taylor 多项式，则对一切 \(N\ge 0\) 有一致余项界
\[
\sup_{|\theta|\le \theta_R}|P_2(\theta)-T_N(\theta)|
\le
\frac{M(r)}{1-\theta_R/r}\left(\frac{\theta_R}{r}\right)^{N+1}.
\]
因此，为使 RH 窗口内的截断误差不超过任意给定的 \(\varepsilon\in(0,1)\)，只需取
\[
N\ge
\frac{\log(\varepsilon^{-1})}{\log(r/\theta_R)}+O_r(1).
\]
在形式极限 \(r\uparrow R_\theta\) 下，上述深度需求呈现严格指数级缩放
\[
N=
\Theta\!\left(
\frac{\log(\varepsilon^{-1})}{\log(R_\theta/\theta_R)}
\right).
\]
\end{theorem}
\begin{proof}
由定理 \ref{thm:xi-rh-critical-before-nearest-complex-branch}，\(P_2\) 在圆盘 \(|\theta|<R_\theta\) 内解析。故当 \(r<R_\theta\) 时，Cauchy 估计给出 Taylor 系数满足
\[
\left|\frac{P_2^{(n)}(0)}{n!}\right|
\le
\frac{M(r)}{r^n}.
\]
于是对任意 \(|\theta|\le \theta_R<r\)，有
\[
|P_2(\theta)-T_N(\theta)|
\le
\sum_{n\ge N+1}\frac{M(r)}{r^n}|\theta|^n
\le
\sum_{n\ge N+1}M(r)\left(\frac{\theta_R}{r}\right)^n
=
\frac{M(r)}{1-\theta_R/r}\left(\frac{\theta_R}{r}\right)^{N+1}.
\]
这就得到一致余项界。

若要求右端不超过 \(\varepsilon\)，对数化并整理即可得
\[
N\ge
\frac{\log(\varepsilon^{-1})}{\log(r/\theta_R)}+O_r(1).
\]
当 \(r\) 固定时，这是关于 \(\log(\varepsilon^{-1})\) 的线性增长；而当 \(r\uparrow R_\theta\) 时，主导比例常数逼近 \(\log(R_\theta/\theta_R)^{-1}\)，故得到所述 \(\Theta\)-级深度律。证毕。
\end{proof}

\begin{theorem}[碰撞位势大偏差率函数在均值点的显式二次胚]\label{thm:xi-time-part65-collision-rate-quadratic-germ}
沿用定理 \ref{thm:xi-real-input-40-collision-rate-cramer-edgeworth} 的记号，记
\[
\mu:=P_2'(0)=\frac{3-\sqrt5}{10},
\]
并把归一化 Legendre--Fenchel 对偶写为
\[
I_2(x):=\sup_{\theta\in\RR}\Bigl(x\theta-\bigl(P_2(\theta)-P_2(0)\bigr)\Bigr).
\]
则在 \(h\to 0\) 时，
\[
\boxed{\;
I_2(\mu+h)
=
\frac{125}{12\sqrt5-10}\,h^2+O(h^3).\;}
\]
等价地，若采用未归一化对偶
\[
I(x):=\sup_{\theta\in\RR}\bigl(x\theta-P_2(\theta)\bigr),
\]
则
\[
\boxed{\;
I(\mu+h)
=
-2\log\varphi
+\frac{125}{12\sqrt5-10}\,h^2+O(h^3).\;}
\]
\end{theorem}
\begin{proof}
由定理 \ref{thm:xi-real-input-40-collision-rate-cramer-edgeworth}，
\[
I_2(\mu+h)
=
\frac{h^2}{2\kappa_2}+O(h^3),
\qquad
\kappa_2=P_2''(0)=\frac{6\sqrt5-5}{125}.
\]
因此
\[
\frac{1}{2\kappa_2}
=
\frac{125}{12\sqrt5-10},
\]
从而得到第一条展开。

另一方面，
\[
I(x)=I_2(x)-P_2(0).
\]
而在 \(u=1\) 时，\(\rho(1)=\varphi^2\)，故
\[
P_2(0)=\log \rho(1)=\log \varphi^2=2\log\varphi.
\]
将其代回，便得未归一化版本
\[
I(\mu+h)
=
-2\log\varphi+\frac{125}{12\sqrt5-10}\,h^2+O(h^3).
\]
证毕。
\end{proof}

\paragraph{平方根归一化的三歧相变与 \texorpdfstring{$u=1$}{u=1} 审计窗口的统一亚临界性}
在纯碰撞三次核的 RH/Ramanujan 窗口、负谱支控制以及实参数审计窗口已经建立之后，还可把平方根归一化下的慢模行为与 \(u=1\) 邻域的统一亚临界性压缩为下列两条结果。

\begin{theorem}[平方根归一化的三歧相变]\label{thm:xi-collision-square-root-normalization-trichotomy}
\leanverified{paper\_xi\_collision\_square\_root\_normalization\_trichotomy}
沿用定理 \ref{thm:xi-time-part62d-collision-rh-negative-branch-control} 与命题 \ref{prop:real-input-40-collision-rhk-window} 的记号。记 Perron 根为 \(\lambda_+(u)\)，次谱半径为
\[
\Lambda(u):=\max\{|\lambda_0(u)|,|\lambda_-(u)|,1\},
\]
并定义平方根归一化的慢模比
\[
\mathfrak S_n(u):=
\frac{\Lambda(u)^n}{\lambda_+(u)^{n/2}}
\qquad
(n\ge 1).
\]
则存在唯一临界参数 \(u_R>1\)，满足
\[
\Lambda(u)\le \lambda_+(u)^{1/2}
\iff
0<u\le u_R,
\qquad
u_R\approx 3.59262181344.
\]
从而：
\begin{enumerate}
  \item 若 \(0<u<u_R\)，则 \(\mathfrak S_n(u)\) 以指数速率趋于 \(0\)；
  \item 若 \(u=u_R\)，则 \(\mathfrak S_n(u)\equiv 1\) 对一切 \(n\ge 1\) 成立；
  \item 若 \(u>u_R\)，则 \(\mathfrak S_n(u)\) 以指数速率发散到 \(+\infty\)。
\end{enumerate}
\end{theorem}
\begin{proof}
定义
\[
r(u):=\frac{\Lambda(u)}{\lambda_+(u)^{1/2}}.
\]
命题 \ref{prop:real-input-40-collision-rhk-window} 已给出
\[
r(u)\le 1
\iff
0<u\le u_R,
\]
且临界参数 \(u_R\) 在正实轴上唯一。于是对 \(u<u_R\)、\(u=u_R\) 与 \(u>u_R\) 分别有
\[
r(u)<1,
\qquad
r(u)=1,
\qquad
r(u)>1.
\]
另一方面，
\[
\mathfrak S_n(u)=r(u)^n.
\]
三种情形遂立即成立。证毕。
\end{proof}

\begin{corollary}[\texorpdfstring{$u=1$}{u=1} 审计窗口中的统一亚临界性]\label{cor:xi-collision-u1-audit-window-uniform-subcriticality}
\leanverified{paper\_xi\_collision\_u1\_audit\_window\_uniform\_subcriticality}
注 \ref{rem:real-input-40-output-potential-degeneracy-phase} 在实参数方向上给出无退化阈值
\[
\theta_0=0.0533477827\ldots,
\qquad
|\theta|<\theta_0
\Longrightarrow
u=e^\theta\in(0.932837\ldots,1.054796\ldots).
\]
记
\[
I_0:=[e^{-\theta_0},e^{\theta_0}].
\]
又因
\[
u_R\approx 3.59262181344,
\]
故
\[
I_0\Subset (0,u_R).
\]
于是存在常数 \(q_0\in(0,1)\)，使得对一切 \(u\in I_0\) 与一切 \(n\ge 1\)，都有
\[
\Lambda(u)^n\le q_0^n\,\lambda_+(u)^{n/2}.
\]
\end{corollary}
\begin{proof}
由定理 \ref{thm:xi-collision-square-root-normalization-trichotomy}，函数
\[
r(u):=\frac{\Lambda(u)}{\lambda_+(u)^{1/2}}
\]
在区间 \((0,u_R)\) 上严格小于 \(1\)。另一方面，\(\Lambda(u)\) 与 \(\lambda_+(u)\) 都由有限维矩阵的特征值模给出，因此 \(r(u)\) 在正实轴上连续。

由 \(\theta_0=0.0533477827\ldots\) 可得
\[
e^{\theta_0}=1.054796\ldots<u_R,
\]
且显然 \(e^{-\theta_0}>0\)，故紧区间 \(I_0=[e^{-\theta_0},e^{\theta_0}]\) 满足
\[
I_0\Subset (0,u_R).
\]
于是可定义
\[
q_0:=\sup_{u\in I_0}r(u).
\]
由连续函数在紧集上取到上确界且 \(r(u)<1\) 于 \(I_0\) 上处处成立，遂得
\[
0<q_0<1.
\]
因此对任意 \(u\in I_0\) 与任意 \(n\ge 1\)，
\[
\Lambda(u)^n
=
r(u)^n\lambda_+(u)^{n/2}
\le
q_0^n\,\lambda_+(u)^{n/2}.
\]
证毕。
\end{proof}

在平方根归一化的三歧相变之外，RH-sharp 临界点本身还允许进一步压缩出一个单步代数滤子与一个余维一的误差指数尖点。

\begin{theorem}[临界平方根模的一阶湮灭]\label{thm:xi-time-part65-critical-sqrt-mode-first-order-annihilation}
沿用命题 \ref{prop:real-input-40-collision-rhk-window} 的记号，设
\[
u_\ast:=u_R,
\qquad
s:=-\lambda_-(u_\ast)>0.
\]
则
\[
\lambda_+(u_\ast)=s^2,
\qquad
\lambda_-(u_\ast)=-s.
\]
进一步，对任意形如
\[
a_n=A\,\lambda_+(u_\ast)^n+B\,\lambda_-(u_\ast)^n+r_n
\qquad (n\ge 0)
\]
的标量序列，只要存在 \(\eta<s\) 使
\[
r_n=O(\eta^n),
\]
则一阶滤子
\[
\mathcal L_s a_n:=a_{n+1}+s\,a_n
\]
满足
\[
\mathcal L_s a_n
=(s^2+s)A\,s^{2n}+\widetilde r_n,
\qquad
\widetilde r_n=O(\eta^n)=o(s^n).
\]
特别地，负谱模 \(\lambda_-(u_\ast)^n\) 被 \(\mathcal L_s\) 单步精确湮灭，而保留下来的主导项只来自临界 Perron 模 \(s^{2n}\)。
\end{theorem}
\begin{proof}
注 \ref{rem:collision-cubic-rhsharp-quintic} 已给出
\[
\lambda_+(u_\ast)=s^2,
\qquad
\lambda_-(u_\ast)=-s.
\]
因此
\[
\mathcal L_s\bigl(\lambda_-(u_\ast)^n\bigr)
=(-s)^{n+1}+s(-s)^n
=0,
\]
而
\[
\mathcal L_s\bigl(\lambda_+(u_\ast)^n\bigr)
=s^{2n+2}+s^{2n+1}
=(s^2+s)s^{2n}.
\]
若 \(r_n=O(\eta^n)\) 且 \(\eta<s\)，则
\[
\mathcal L_s r_n=r_{n+1}+s\,r_n=O(\eta^n)=o(s^n).
\]
将上述三式按线性叠加，便得到所示展开。证毕。
\end{proof}

\begin{corollary}[平方根律的尖点接触与非厚相性]\label{cor:xi-time-part65-sqrt-law-cusp-nonthick-phase}
沿用推论 \ref{cor:real-input-40-collision-error-phase} 的记号，定义
\[
\beta(u):=\frac{\log \Lambda(u)}{\log \lambda_+(u)}.
\]
则
\[
\beta(u)=0
\qquad (0<u\le 1),
\]
\[
\beta(u)\le \frac12
\qquad (1<u\le u_R),
\]
且等号仅在 \(u=u_R\) 发生；此外，
\[
\beta(u)>\frac12
\qquad (u>u_R),
\qquad
\beta(u)\to 1
\qquad (u\to\infty).
\]
因此，\(\beta=\frac12\) 的平方根级误差并不形成参数上的开相区，而只在唯一临界参数 \(u_R\) 处与相图发生余维一接触。
\end{corollary}
\begin{proof}
推论 \ref{cor:real-input-40-collision-error-phase} 已逐段给出 \(\beta(u)\) 在
\((0,1]\)、\((1,u_R]\) 与 \((u_R,\infty)\) 上的精确刻画，并证明
\(\beta(u)\to 1\) 当 \(u\to\infty\)。将其直接合并，便得到全部结论。证毕。
\end{proof}


\noindent
这十二条结果把 bin-fold 末位两态律、规范群 Stirling 账本、原子 Euler 手术与碰撞位势的解析几何再度压缩为同一条闭合主线：一方面，阿贝尔化在高分辨率极限中达到全部稳定类型的最终满秩，而中心在同一极限中彻底消失，从而纤维奇偶荷的最小完备秩被严格锁定为 \(F_{m+2}\)；进一步，完整规范群复杂度相对于可见异常维数并不同阶，而是以 \((2/\varphi)^m\) 呈现严格的指数超载；另一方面，其与平均纤维信息项之间仍保留一条由黄金比锁定的线性响应斜率，纯碰撞误差指数的缺口则只以临界幂 \(p=2\) 的精确边界可和方式衰减；再者，原子因子在 primitive 与 Abel 有限部两层都表现为完全局域的单点位移，而平方根归一化的慢模比不仅在唯一临界参数 \(u_R\) 处发生严格三歧相变，而且该临界点本身已退化为一个可由一阶滤子单步消去的负谱节点，\(\beta=\tfrac12\) 也只以余维一尖点方式触及相图；最后，RH 失效不仅先于最近复分歧发生，而且整个失效区间都仍由 \(\theta=0\) 的 Taylor 胚直接控制，而均值点的大偏差率函数曲率常数同样被 \(\sqrt5\) 完全闭式锁定。由此，bin-fold 的群论刚性、原子谱的局域手术与碰撞位势的局部热力学几何便被统一收束到同一组可审计的终端后果之中。

\paragraph{Lee--Yang 三次闭包、Latt\`es 六次正交、\texorpdfstring{$\chi^2$}{chi2} 缺口与 Jacobian 半单刚性}
在 Lee--Yang 分支不变量的共同三次数域、bin-fold 碰撞缺口的 Fibonacci 共振强制，以及 \(S_4\)--Hurwitz Jacobian 的可见等型块分解已经确立之后，还可进一步抽取出下列一组彼此兼容的刚性后果。

\begin{theorem}[Lee--Yang 分支几何的单参数有理统一化]\label{thm:xi-time-part65b-leyang-branch-b-birational-uniformization}
沿用定理 \ref{thm:xi-terminal-zm-b2-kappa2-rational-interconstruction}、定理 \ref{thm:xi-terminal-zm-stokes-t-elliptic-branch-coordinate} 与定理 \ref{thm:xi-terminal-zm-leyang-elliptic-structure} 的记号。令 \((\lambda_*,y_*)\) 为非平凡 Lee--Yang 分歧点，并记
\[
u:=\kappa_*^2,
\qquad
b:=B^2.
\]
则 \(u\)、\(\lambda_*\) 与 \(y_*\) 都可写成 \(b\) 的有理函数；具体地，
\[
u=\frac{2294-8463b-918b^2}{14388},
\]
\[
\lambda_*
=
\frac{918b^2+8463b+4900}{2(918b^2+8463b+7298)},
\]
\[
y_*
=
-\frac{3597\,b\,(918b^2+8463b+4900)^2}
{(918b^2+8463b-2294)(918b^2+8463b+7298)^2}.
\]
因此
\[
\QQ(b)=\QQ(u)=\QQ(\lambda_*)=\QQ(y_*).
\]
换言之，Lee--Yang 分支点的局部几何不变量 \(u,\lambda_*,y_*,b\) 并非只是在同一三次数域中同居，而可由单一参数 \(b\) 作显式有理统一化。
\end{theorem}
\begin{proof}
由定理 \ref{thm:xi-terminal-zm-b2-kappa2-rational-interconstruction}，
\[
b
=
-2\cdot
\frac{24708u^2-28628u+7733}{51084u^2-44412u+8735}.
\]
同一定理并已给出 \(\QQ(b)=\QQ(u)\)，故 \(u\in \QQ(b)\) 必可唯一写成
\[
u=x_0+x_1b+x_2b^2
\qquad (x_0,x_1,x_2\in \QQ).
\]
把这一待定表达代回上式，并结合 \(b\) 的三次关系
\[
162b^3+1593b^2+1998b+1184=0
\]
在 \(\QQ(b)\) 中约化，即唯一解得
\[
u=\frac{2294-8463b-918b^2}{14388}.
\]

再由定理 \ref{thm:xi-terminal-zm-kappa-square-cubic-field-s3}，
\[
\lambda_*=\frac{3(2u-1)}{4(3u-2)}.
\]
把前式代入并化简，便得到所示 \(\lambda_*(b)\) 的表达。

另一方面，定理 \ref{thm:xi-terminal-zm-stokes-t-elliptic-branch-coordinate} 给出
\[
4\lambda_*y_*=4\lambda_*^3-3\lambda_*^2-2\lambda_*+1.
\]
再把 \(\lambda_*(b)\) 代入并整理，即得所示 \(y_*(b)\) 的表达。

由定理 \ref{thm:xi-terminal-zm-b2-kappa2-rational-interconstruction} 已知
\[
\QQ(b)=\QQ(u)=\QQ(\lambda_*).
\]
又由于 \(y_*\) 是定理 \ref{thm:xi-time-part9g-leyang-cubic-discriminant} 中不可约三次
\[
P_{\mathrm{LY}}(y)=256y^3+411y^2+165y+32
\]
的根，故 \([\QQ(y_*):\QQ]=3\)。而上式已证明 \(y_*\in \QQ(b)\)，于是只能有
\[
\QQ(y_*)=\QQ(b).
\]
遂得全部域相等式。证毕。
\end{proof}

\begin{theorem}[Lee--Yang 分支不变量链的共同二次分辨子域与统一 \texorpdfstring{$S_3$}{S3} 闭包]\label{thm:xi-time-part65b-leyang-common-s3-closure}
沿用定理 \ref{thm:xi-time-part65b-leyang-branch-b-birational-uniformization} 的记号，记
\[
u:=\kappa_*^2,
\qquad
b:=B^2,
\qquad
K:=\QQ(u)=\QQ(b)=\QQ(\lambda_*)=\QQ(y_*).
\]
设 \(u\) 与 \(b\) 的三次最小多项式分别为
\[
324u^3-540u^2+333u-74,
\qquad
162b^3+1593b^2+1998b+1184,
\]
其判别式分别记为
\[
\Delta_u=-2^6\cdot 3^9\cdot 37,
\qquad
\Delta_b=-2^2\cdot 3^9\cdot 11^2\cdot 37\cdot 109^2.
\]
若 \(L_{\mathrm{LY}}\) 记为三次数域 \(K/\QQ\) 的伽罗瓦闭包，则有
\[
\mathrm{sf}(\Delta_u)=\mathrm{sf}(\Delta_b)=-111,
\]
并且
\[
\Gal(L_{\mathrm{LY}}/\QQ)\cong S_3,
\qquad
L_{\mathrm{LY}}^{A_3}=\QQ(\sqrt{-111}).
\]
等价地，\(u\)、\(b\)、\(\lambda_*\) 与 \(y_*\) 都落在同一非伽罗瓦三次数域 \(K\) 中，其共同伽罗瓦闭包唯一的非平凡二次正规子域恰为 \(\QQ(\sqrt{-111})\)。
\end{theorem}
\begin{proof}
由定理 \ref{thm:xi-time-part65b-leyang-branch-b-birational-uniformization} 与定理 \ref{thm:xi-terminal-zm-stokes-leyang-shared-artin-representation}，
\[
\QQ(u)=\QQ(b)=\QQ(\lambda_*)=\QQ(y_*)=:K,
\qquad
\Gal(L_{\mathrm{LY}}/\QQ)\cong S_3.
\]
因此 \(u\)、\(b\)、\(\lambda_*\) 与 \(y_*\) 都生成同一非伽罗瓦三次数域 \(K\)，其分裂域必同为 \(L_{\mathrm{LY}}\)。

另一方面，
\[
\Delta_u
=
-2^6\cdot 3^9\cdot 37
=
(2^3\cdot 3^4)^2\cdot (-111),
\]
\[
\Delta_b
=
-2^2\cdot 3^9\cdot 11^2\cdot 37\cdot 109^2
=
(2\cdot 3^4\cdot 11\cdot 109)^2\cdot (-111),
\]
故
\[
\mathrm{sf}(\Delta_u)=\mathrm{sf}(\Delta_b)=-111.
\]
对不可约三次域，其伽罗瓦闭包的唯一非平凡二次正规子域由判别式的平方自由核决定，因此
\[
L_{\mathrm{LY}}^{A_3}=\QQ(\sqrt{-111}).
\]
证毕。
\end{proof}

\begin{proposition}[Lee--Yang 三次数域的两组显式整生成元]\label{prop:xi-time-part65b-leyang-u-b-integral-generators}
沿用定理 \ref{thm:xi-time-part65b-leyang-common-s3-closure} 的记号，定义
\[
x:=18u,
\qquad
y:=18b.
\]
则 \(x\) 与 \(y\) 都是代数整数，并分别满足首一整系数三次方程
\[
x^3-30x^2+333x-1332=0,
\]
\[
y^3+177y^2+3996y+42624=0.
\]
特别地，
\[
\QQ(x)=\QQ(u)=K,
\qquad
\QQ(y)=\QQ(b)=K.
\]
\end{proposition}
\begin{proof}
由定理 \ref{thm:xi-time-part65b-leyang-common-s3-closure}，
\[
324u^3-540u^2+333u-74=0.
\]
代入 \(u=x/18\) 并清分母，得到
\[
x^3-30x^2+333x-1332=0.
\]
同理，由
\[
162b^3+1593b^2+1998b+1184=0
\]
代入 \(b=y/18\) 并清分母，得到
\[
y^3+177y^2+3996y+42624=0.
\]
两式皆为首一整系数多项式，故 \(x\) 与 \(y\) 都是代数整数。
又因 \(18\in \QQ^\times\)，缩放不改变所生成的 \(\QQ\)-子域，因此
\[
\QQ(x)=\QQ(u)=K,
\qquad
\QQ(y)=\QQ(b)=K.
\]
证毕。
\end{proof}

\begin{corollary}[共同 \texorpdfstring{$S_3$}{S3} 闭包与 Latt\`es 六次域的线性互素性]\label{cor:xi-time-part65b-leyang-lattes-linear-disjointness}
设
\[
C(X)=X^6-8X^5+24X^4-36X^3+34X^2-20X+4
\]
的分裂域为 \(L_4\)，并满足
\[
\Disc(C)=-2^{19}\cdot 37^3,
\qquad
\Gal(L_4/\QQ)\cong S_4\times C_2.
\]
则定理 \ref{thm:xi-time-part65b-leyang-common-s3-closure} 中的共同闭包 \(L_{\mathrm{LY}}\) 与 \(L_4\) 线性互素：
\[
L_{\mathrm{LY}}\cap L_4=\QQ.
\]
因而
\[
\Gal(L_{\mathrm{LY}}L_4/\QQ)
\cong
S_3\times (S_4\times C_2).
\]
\end{corollary}
\begin{proof}
由定理 \ref{thm:xi-time-part65b-leyang-common-s3-closure}，
\[
L_{\mathrm{LY}}^{A_3}=\QQ(\sqrt{-111}),
\]
这是 \(L_{\mathrm{LY}}/\QQ\) 的唯一非平凡二次正规子域。于是交域
\[
F:=L_{\mathrm{LY}}\cap L_4
\]
作为 \(L_{\mathrm{LY}}/\QQ\) 的伽罗瓦子扩张，只能是 \(\QQ\) 或 \(\QQ(\sqrt{-111})\)。

若 \(F=\QQ(\sqrt{-111})\)，则素数 \(3\) 在 \(F/\QQ\) 中分歧；然而由
\[
\Disc(C)=-2^{19}\cdot 37^3
\]
可知 \(L_4/\QQ\) 只在 \(\{2,37\}\) 处分歧，故 \(3\) 在 \(L_4\) 的任一子扩张中都不分歧，矛盾。故只能有
\[
L_{\mathrm{LY}}\cap L_4=\QQ.
\]
最后，两边皆为 \(\QQ\) 上的伽罗瓦扩张，故复合域的伽罗瓦群为直积：
\[
\Gal(L_{\mathrm{LY}}L_4/\QQ)
\cong
\Gal(L_{\mathrm{LY}}/\QQ)\times \Gal(L_4/\QQ)
\cong
S_3\times (S_4\times C_2).
\]
证毕。
\end{proof}

\begin{theorem}[临界共振常数强制 \texorpdfstring{$\chi^2$}{chi2} 碰撞缺口具有正的极限下界]\label{thm:xi-time-part65b-cphi-chi2-gap-floor}
沿用定义 \ref{def:fold-normalized-amplitude}、命题 \ref{prop:fold-bin-chi2-col} 与定理 \ref{thm:fold-critical-resonance-constant} 的记号。记
\[
a_m(k):=\frac{|\widehat d_m(k)|}{2^m},
\qquad
C_\varphi:=\lim_{m\to\infty}a_m(F_m)=\lim_{m\to\infty}a_m(F_{m+1})>0.
\]
若 \(p_m^{\mathrm{bin}}\) 为 bin-fold 推前分布，\(u_m\) 为 \(X_m\) 上的均匀分布，则
\[
\liminf_{m\to\infty}\chi^2\!\bigl(p_m^{\mathrm{bin}}\|u_m\bigr)\ge 2C_\varphi^2.
\]
等价地，
\[
\liminf_{m\to\infty}
F_{m+2}\Bigl(\mathsf{Col}_m-\frac{1}{F_{m+2}}\Bigr)
\ge
2C_\varphi^2.
\]
\end{theorem}
\begin{proof}
由命题 \ref{prop:fold-bin-chi2-col}，
\[
\chi^2\!\bigl(p_m^{\mathrm{bin}}\|u_m\bigr)
=
|X_m|\,\mathsf{Col}_m-1
=
F_{m+2}\Bigl(\mathsf{Col}_m-\frac{1}{F_{m+2}}\Bigr).
\]
另一方面，结论 \ref{con:xi-time-part9m-cphi-forces-collision-gap-and-peak-bias} 已给出
\[
\liminf_{m\to\infty}
F_{m+2}\Bigl(\mathsf{Col}_m-\frac{1}{F_{m+2}}\Bigr)
\ge
2C_\varphi^2.
\]
将二者合并即得
\[
\liminf_{m\to\infty}\chi^2\!\bigl(p_m^{\mathrm{bin}}\|u_m\bigr)\ge 2C_\varphi^2.
\]
第二式只是上述恒等式的等价改写。证毕。
\end{proof}

\begin{theorem}[可见等型块代数的显式半单嵌入与 N\'eron--Severi 秩下界]\label{thm:xi-time-part65b-jacobian-visible-semisimple-ns-floor}
沿用定理 \ref{thm:killo-s4-genus49-jacobian-computable-isogeny} 的 \(\QQ\)-同源分解
\[
J(X)\sim_{\QQ}J(H)\times J(Z)^2\times J(C_F)^3\times P^3.
\]
则 \(\operatorname{End}_{\QQ}^0(J(X))\) 含有显式半单 \(\QQ\)-子代数
\[
\mathcal A_X:=\QQ\times M_2(\QQ)\times M_3(\QQ)\times M_3(\QQ)
\hookrightarrow
\operatorname{End}_{\QQ}^0(J(X)),
\]
从而
\[
\dim_{\QQ}\operatorname{End}_{\QQ}^0(J(X))\ge 1+4+9+9=23.
\]
进一步，在右端乘积上取任意乘积极化，并把相应 Rosati 对合记为 \(\dagger_{\mathrm{vis}}\)。则
\[
\dim_{\QQ}(\mathcal A_X)^{\dagger_{\mathrm{vis}}}
=
1+\frac{2\cdot 3}{2}+\frac{3\cdot 4}{2}+\frac{3\cdot 4}{2}
=
16.
\]
特别地，
\[
\operatorname{rank}\operatorname{NS}(J(X)_{\overline{\QQ}})\ge 16.
\]
\end{theorem}
\begin{proof}
定理 \ref{thm:xi-time-part9n1b-jacobian-projector-idempotent-algebra} 已经给出显式嵌入
\[
\mathcal A_X
\hookrightarrow
\operatorname{End}_{\QQ}^0(J(X))
\]
以及
\[
\dim_{\QQ}\operatorname{End}_{\QQ}^0(J(X))\ge 23.
\]
故只需证明 Rosati 不动维数与 N\'eron--Severi 秩下界。

取一个实现上述 \(\QQ\)-同源分解的同源映射，并在乘积
\[
A\times B^2\times C^3\times D^3
\qquad
\bigl(A:=J(H),\ B:=J(Z),\ C:=J(C_F),\ D:=P\bigr)
\]
上选取乘积极化 \(\lambda_{\mathrm{prod}}\)。在该模型下，\(\mathcal A_X\) 恰由四个可见等型块的重数空间上的标量矩阵作用给出，因此 \(\lambda_{\mathrm{prod}}\) 的 Rosati 对合在 \(\QQ\) 因子上恒等，在 \(M_r(\QQ)\) 因子上与转置等价。于是其不动子空间维数分别为
\[
1,\qquad \frac{2\cdot 3}{2}=3,\qquad \frac{3\cdot 4}{2}=6,\qquad \frac{3\cdot 4}{2}=6,
\]
从而
\[
\dim_{\QQ}(\mathcal A_X)^{\dagger_{\mathrm{vis}}}=1+3+6+6=16.
\]

最后，同源映射的拉回在 N\'eron--Severi 群上于有理化后给出同构，而对任意极化阿贝尔簇 \(A\) 均有标准识别
\[
\operatorname{NS}(A_{\overline{\QQ}})\otimes_{\ZZ}\QQ
\cong
\operatorname{End}^0_{\overline{\QQ}}(A)^{\dagger=1}.
\]
故上述 \(16\) 维 Rosati 不动子空间传回 \(J(X)\) 后仍给出
\[
\operatorname{rank}\operatorname{NS}(J(X)_{\overline{\QQ}})\ge 16.
\]
证毕。
\end{proof}

\noindent
以上五条结果把 Lee--Yang 分支几何的单参数有理统一化、共同 \(S_3\)-闭包、Latt\`es 六次闭包、黄金共振二阶缺口与 \(S_4\)--Hurwitz Jacobian 的可见块分解进一步压缩到同一层级：一方面，\(b\) 不仅与 \(u\)、\(\lambda_*\) 及 \(y_*\) 生成同一三次数域，而且把全部分支点局部几何显式统一化；该共同三次数域的伽罗瓦闭包唯一的非平凡二次正规子域仍被 \(\QQ(\sqrt{-111})\) 完全锁定，并与 \(L_4\) 严格正交；另一方面，临界共振常数把 bin-fold 与均匀基线之间的 \(\chi^2\) 偏离锁定为严格正的极限底噪，而 Jacobian 的可见等型块则单独强制出一个 \(23\) 维有理半单子代数与至少 \(16\) 维的 N\'eron--Severi 秩下界。


\paragraph{escort 冻结图、归一化 residue law 与输出位势的一阶非高斯非对称}
在冻结相压力谱的线性支正、escort 态向最大纤维层的指数集中、真实输入 \(40\) 状态核的单原子 primitive/residue 剥离，以及输出位势的准幂极限定律与局部 Legendre 正规形已经建立之后，还可把这些链进一步压缩为一组互相衔接的结构后果：首先，escort--Rényi 熵率不仅在 \((a,s)\)-平面上给出双参数冻结边界，而且在更强阈值上提升为由 \(\log K_m\) 控制的有限尺度主项，并进一步诱导出最大纤维扇区在预言机容量中的显式指数贡献与侧信息临界斜率；随后，单原子的秩一局域手术被提升为归一化 periodic residue law 的精确指数、正则单纯形首项轨道与 centered residue 各向同性；最后，中心极限层的一阶 Edgeworth 偏斜与局部大偏差层的三次不对称项仍由同一输出位势常数 \(\kappa_3\) 严格锁定。

\leanverified{paper\_xi\_time\_part65c\_escort\_renyi\_two\_parameter\_freezing\_line}
\begin{theorem}[escort--Rényi 熵的双参数冻结线]\label{thm:xi-time-part65c-escort-renyi-two-parameter-freezing-line}
沿用定理 \ref{thm:xi-fixed-freezing-escort-renyi-spectrum-collapse} 与定理 \ref{thm:fold-pressure-freezing-threshold} 的记号。对整数 \(a\ge 1\) 定义 escort 分布
\[
\pi_{m,a}(x):=\frac{d_m(x)^a}{S_a(m)},
\qquad
S_a(m):=\sum_{x\in X_m}d_m(x)^a.
\]
对整数 \(s\ge 2\) 定义其 Rényi 熵
\[
H_s(\pi_{m,a})
:=
\frac{1}{1-s}\log\sum_{x\in X_m}\pi_{m,a}(x)^s.
\]
若整数压力极限
\[
P_q:=\lim_{m\to\infty}\frac{1}{m}\log S_q(m)
\qquad(q\ge 1)
\]
存在，则 escort--Rényi 熵率
\[
\mathfrak h_s(a):=
\lim_{m\to\infty}\frac{1}{m}H_s(\pi_{m,a})
\]
存在且满足严格闭式
\[
\boxed{\;
\mathfrak h_s(a)=\frac{sP_a-P_{as}}{s-1}.\;}
\]
特别地：
\begin{enumerate}
  \item 若 \(a>a_{\mathrm c}\)，则对每个整数 \(s\ge 2\) 都有
  \[
  \boxed{\;\mathfrak h_s(a)=g_*.\;}
  \]
  \item 更一般地，若 \(a\le a_{\mathrm c}\) 但 \(as>a_{\mathrm c}\)，则
  \[
  \boxed{\;
  \mathfrak h_s(a)
  =
  \frac{sP_a-(as\,\alpha_*+g_*)}{s-1}.\;}
  \]
\end{enumerate}
因此，原始压力谱的冻结阈值位于 \(a=a_{\mathrm c}\)，而 escort--Rényi 熵的相应双参数冻结边界则位于
\[
\boxed{\;as=a_{\mathrm c}.\;}
\]
\end{theorem}
\begin{proof}
由定义，
\[
\sum_{x\in X_m}\pi_{m,a}(x)^s
=
\sum_{x\in X_m}\frac{d_m(x)^{as}}{S_a(m)^s}
=
\frac{S_{as}(m)}{S_a(m)^s}.
\]
故
\[
H_s(\pi_{m,a})
=
\frac{1}{1-s}\bigl(\log S_{as}(m)-s\log S_a(m)\bigr).
\]
两边除以 \(m\) 并令 \(m\to\infty\)，即得
\[
\mathfrak h_s(a)=\frac{sP_a-P_{as}}{s-1}.
\]

若 \(a>a_{\mathrm c}\)，则由于 \(s\ge 2\)，自动有 \(as>a>a_{\mathrm c}\)。由定理 \ref{thm:fold-pressure-freezing-threshold}，
\[
P_a=a\alpha_*+g_*,
\qquad
P_{as}=as\,\alpha_*+g_*.
\]
代回上式得到
\[
\mathfrak h_s(a)
=
\frac{s(a\alpha_*+g_*)-(as\,\alpha_*+g_*)}{s-1}
=
g_*.
\]

若 \(a\le a_{\mathrm c}\) 但 \(as>a_{\mathrm c}\)，则仍有
\[
P_{as}=as\,\alpha_*+g_*,
\]
而 \(P_a\) 保持一般形式，代回即得第二条。最后，由 \(P_{as}\) 是否进入冻结支正只取决于 \(as\) 是否超过 \(a_{\mathrm c}\)，故双参数冻结边界正是直线 \(as=a_{\mathrm c}\)。证毕。
\end{proof}

\begin{corollary}[冻结相内全部整数阶 Rényi 熵率塌缩到同一常数]\label{cor:xi-time-part65c-escort-all-integer-renyi-collapse}
沿用定理 \ref{thm:xi-time-part65c-escort-renyi-two-parameter-freezing-line} 的记号。若 \(a>a_{\mathrm c}\)，则对每个整数 \(s\ge 2\) 都有
\[
\boxed{\;
\lim_{m\to\infty}\frac{1}{m}H_s(\pi_{m,a})=g_*.\;}
\]
再与定理 \ref{thm:xi-fixed-freezing-escort-renyi-spectrum-collapse} 在 \(q=\infty\) 的结论合并，得到：在冻结区 \(a>a_{\mathrm c}\) 内，escort 分布 \(\pi_{m,a}\) 的全部整数阶 Rényi 熵率，包括 \(s=\infty\) 的 \(\min\)-entropy 率，都在极限中塌缩为同一常数 \(g_*\)。
\end{corollary}
\begin{proof}
对任意整数 \(s\ge 2\)，由 \(a>a_{\mathrm c}\) 可知 \(as>a_{\mathrm c}\)。于是定理 \ref{thm:xi-time-part65c-escort-renyi-two-parameter-freezing-line} 的第一条直接给出
\[
\lim_{m\to\infty}\frac{1}{m}H_s(\pi_{m,a})=g_*.
\]
另一方面，定理 \ref{thm:xi-fixed-freezing-escort-renyi-spectrum-collapse} 的第三条已给出
\[
\lim_{m\to\infty}\frac{1}{m}H_\infty(\pi_{m,a})=g_*.
\]
二者合并即得结论。证毕。
\end{proof}

\begin{theorem}[冻结相 escort--Rényi 熵的有限尺度主项由基态简并度精确控制]\label{thm:xi-time-part65c-escort-renyi-finite-size-mainterm}
\leanverified{paper\_xi\_time\_part65c\_escort\_renyi\_finite\_size\_mainterm}
沿用结论 \ref{con:xi-time-part57ba-freezing-escort-groundstate-fdiv-closed}、推论 \ref{cor:xi-fixed-freezing-escort-groundstate-tv-identity} 与定理 \ref{thm:xi-fixed-freezing-power-sum-ratio-locking} 的记号。对任意实数 \(s>1\)，定义更强阈值
\[
a_{\mathrm c,s}:=
\frac{\log\varphi-\frac{1}{s}g_*}{\alpha_*-\alpha_*^{(2)}}.
\]
则当 \(a>a_{\mathrm c,s}\) 时，随着 \(m\to\infty\)，有
\[
\boxed{\;
\sum_{x\in X_m}\pi_{m,a}(x)^s
=
K_m^{1-s}\bigl(1+o(1)\bigr),\;}
\]
\[
\boxed{\;
H_s(\pi_{m,a})=\log K_m+o(1).\;}
\]
换言之，在半直线 \(a>a_{\mathrm c,s}\) 上，冻结 escort--Rényi 熵不仅在指数率上塌缩到 \(g_*\)，其完整有限尺度主项也被基态简并度 \(\log K_m\) 精确控制。
\end{theorem}
\begin{proof}
由结论 \ref{con:xi-time-part57ba-freezing-escort-groundstate-fdiv-closed}，可写
\[
\pi_{m,a}=p_{m,a}\sigma_m+(1-p_{m,a})\eta_{m,a},
\]
其中 \(\sigma_m=\mathrm{Unif}(A_m^{\max})\)，且 \(\sigma_m\) 与 \(\eta_{m,a}\) 的支撑不交。故
\[
\sum_{x\in X_m}\pi_{m,a}(x)^s
=
p_{m,a}^s\sum_x \sigma_m(x)^s
+(1-p_{m,a})^s\sum_x \eta_{m,a}(x)^s.
\]
由于 \(\sigma_m\) 在 \(A_m^{\max}\) 上均匀，故
\[
\sum_x \sigma_m(x)^s=K_m^{1-s}.
\]
又因 \(s>1\)，任意概率分布都满足 \(\sum_x \eta_{m,a}(x)^s\le 1\)。于是
\[
p_{m,a}^sK_m^{1-s}
\le
\sum_{x\in X_m}\pi_{m,a}(x)^s
\le
p_{m,a}^sK_m^{1-s}+(1-p_{m,a})^s.
\]
将上式除以 \(K_m^{1-s}\) 得
\[
p_{m,a}^s
\le
K_m^{s-1}\sum_{x\in X_m}\pi_{m,a}(x)^s
\le
p_{m,a}^s+K_m^{s-1}(1-p_{m,a})^s.
\]
因此只需验证
\[
K_m^{s-1}(1-p_{m,a})^s\to 0.
\]
由推论 \ref{cor:xi-fixed-freezing-escort-groundstate-tv-identity}，
\[
1-p_{m,a}\le \exp\!\bigl(-m\Delta(a)+o(m)\bigr),
\qquad
\Delta(a):=a(\alpha_*-\alpha_*^{(2)})-(\log\varphi-g_*)>0.
\]
另一方面，定理 \ref{thm:xi-fixed-freezing-power-sum-ratio-locking} 给出
\[
\log K_m=g_*m+o(m).
\]
故
\[
\log\!\Bigl(K_m^{s-1}(1-p_{m,a})^s\Bigr)
\le
m\bigl((s-1)g_*-s\Delta(a)\bigr)+o(m).
\]
当且仅当
\[
s\Delta(a)>(s-1)g_*
\]
时，右端趋于 \(-\infty\)。将 \(\Delta(a)\) 的定义代回并整理，恰得
\[
a>\frac{\log\varphi-\frac{1}{s}g_*}{\alpha_*-\alpha_*^{(2)}}=a_{\mathrm c,s}.
\]
于是
\[
K_m^{s-1}\sum_{x\in X_m}\pi_{m,a}(x)^s\to 1,
\]
即
\[
\sum_{x\in X_m}\pi_{m,a}(x)^s
=
K_m^{1-s}\bigl(1+o(1)\bigr).
\]
两边取对数并乘以 \((1-s)^{-1}\)，便得到
\[
H_s(\pi_{m,a})=\log K_m+o(1).
\]
证毕。
\end{proof}

\begin{theorem}[冻结相最大纤维扇区给出的预言机容量指数下界]\label{thm:xi-time-part65ca-frozen-max-sector-oracle-capacity-lower-bound}
\leanverified{paper\_xi\_time\_part65ca\_frozen\_max\_sector\_oracle\_capacity\_lower\_bound}
沿用结论 \ref{con:xi-time-part57ba-freezing-escort-groundstate-fdiv-closed}、
推论 \ref{cor:xi-fixed-freezing-escort-groundstate-tv-identity} 与
定理 \ref{thm:xi-fixed-freezing-power-sum-ratio-locking} 的记号，并固定
\[
a>a_{\mathrm c}.
\]
记
\[
M_m:=\max_{x\in X_m}d_m(x),
\qquad
A_m^{\max}:=\{x\in X_m:\ d_m(x)=M_m\},
\qquad
K_m:=|A_m^{\max}|.
\]
若整数预算序列 \((B_m)_{m\ge 1}\) 满足
\[
B_m=\beta m+o(m)
\qquad(\beta\ge 0),
\]
则离散预言机容量
\[
C_m(B_m):=\sum_{x\in X_m}\min\bigl(d_m(x),2^{B_m}\bigr)
\]
满足严格指数下界
\[
\boxed{\;
\liminf_{m\to\infty}\frac{1}{m}\log C_m(B_m)
\ge
g_*+\min\{\alpha_*,\beta\log 2\}.\;}
\]
\end{theorem}
\begin{proof}
由定义，对每个 \(m\) 都有
\[
C_m(B_m)
\ge
\sum_{x\in A_m^{\max}}\min\bigl(M_m,2^{B_m}\bigr)
=
K_m\min\bigl(M_m,2^{B_m}\bigr).
\]
另一方面，定理 \ref{thm:xi-fixed-freezing-power-sum-ratio-locking} 已给出
\[
K_m=e^{g_*m+o(m)},
\qquad
M_m=e^{\alpha_*m+o(m)}.
\]
又由
\[
B_m=\beta m+o(m)
\]
可得
\[
2^{B_m}=e^{(\beta\log 2)m+o(m)}.
\]
故
\[
\min\bigl(M_m,2^{B_m}\bigr)
=
\exp\!\Bigl(
\min\{\alpha_*,\beta\log 2\}\,m+o(m)
\Bigr).
\]
将上述估计代回容量下界并取对数、除以 \(m\)、再取 \(\liminf\)，即得结论。
\end{proof}

\begin{corollary}[最大纤维扇区的侧信息临界斜率]\label{cor:xi-time-part65ca-frozen-max-sector-sideinfo-critical-slope}
\leanverified{paper\_xi\_time\_part65ca\_frozen\_max\_sector\_sideinfo\_critical\_slope}
沿用定理 \ref{thm:xi-time-part65ca-frozen-max-sector-oracle-capacity-lower-bound} 的记号，定义最优成功率
\[
\mathrm{Succ}_m(B):=\frac{C_m(B)}{2^m},
\]
以及最大纤维扇区单独贡献的成功率下界
\[
\mathrm{Succ}_m^{\max}(B)
:=
2^{-m}K_m\min\bigl(M_m,2^B\bigr)
\le
\mathrm{Succ}_m(B).
\]
则对满足 \(B_m=\beta m+o(m)\) 的任意整数预算序列，都有
\[
\boxed{\;
\liminf_{m\to\infty}\frac{1}{m}\log \mathrm{Succ}_m^{\max}(B_m)
\ge
-\bigl[\log 2-g_*-\min\{\alpha_*,\beta\log 2\}\bigr]_+,\;}
\]
其中 \([t]_+:=\max\{t,0\}\)。

特别地，若记
\[
\beta_c^{\max}
:=
\begin{cases}
\dfrac{\log 2-g_*}{\log 2},
& \log 2-g_*\le \alpha_*,\\[8pt]
+\infty,
& \log 2-g_*>\alpha_*,
\end{cases}
\]
则：
\begin{enumerate}
  \item 若 \(\beta>\beta_c^{\max}\)，则最大纤维扇区单独给出的成功率下界不再呈指数衰减；
  \item 若 \(\beta<\beta_c^{\max}\)，则该扇区的单独贡献仍保持指数级稀薄；
  \item 若 \(\log 2-g_*>\alpha_*\)，则即便侧信息按线性速率增长，最大纤维扇区本身也不足以闭合微态熵缺口。
\end{enumerate}
\end{corollary}
\begin{proof}
由定理 \ref{thm:xi-time-part65ca-frozen-max-sector-oracle-capacity-lower-bound}，
\[
\liminf_{m\to\infty}\frac{1}{m}\log \mathrm{Succ}_m^{\max}(B_m)
\ge
g_*+\min\{\alpha_*,\beta\log 2\}-\log 2.
\]
另一方面，由
\[
K_mM_m\le \sum_{x\in X_m}d_m(x)=2^m
\]
可得
\[
g_*+\alpha_*\le \log 2,
\]
故上式右端恒不为正，于是等价地可改写为所示的
\([\,\cdot\,]_+\) 形式。

再看临界斜率。若
\[
\log 2-g_*\le \alpha_*,
\]
则使
\[
g_*+\min\{\alpha_*,\beta\log 2\}-\log 2
\]
首次达到 \(0\) 的唯一有限值为
\[
\beta=\frac{\log 2-g_*}{\log 2}.
\]
若
\[
\log 2-g_*>\alpha_*,
\]
则对一切有限 \(\beta\) 都有
\[
g_*+\min\{\alpha_*,\beta\log 2\}-\log 2
\le
g_*+\alpha_*-\log 2<0,
\]
因而不存在有限临界斜率，只能记为 \(+\infty\)。三条解释正是这两种情形的直接改写。
\end{proof}


\begin{theorem}[单原子 centered residue 扰动向量构成正则单纯形]\label{thm:xi-time-part65c-atomic-centered-residue-simplex}
沿用定理 \ref{thm:xi-atomic-witt-energy-residue-l2-law} 与结论 \ref{con:xi-time-part57ba-core-full-residue-rankone-perturbation} 的记号。固定模数 \(m\ge 2\)，并对每个 \(\ell\in \ZZ/m\ZZ\) 定义 centered residue 扰动向量
\[
\Delta^{(\ell)}\in\RR^{\ZZ/m\ZZ},
\qquad
\Delta^{(\ell)}_r:=2\,\mathbf 1_{\{r=\ell\}}-\frac{2}{m}
\qquad (r\in\ZZ/m\ZZ).
\]
等价地，若 even primitive/residue 注入发生在剩余类 \(r\equiv \ell\pmod m\)，则 \(\Delta^{(\ell)}\) 正是相应未中心化原子 \(2e_\ell\) 去掉全局平均 \(2m^{-1}\mathbf 1\) 后得到的 centered residue 差。
若记 \(\omega_m:=e^{2\pi i/m}\)，并采用离散 Fourier 变换
\[
\widehat x(j):=\sum_{r\in\ZZ/m\ZZ}x_r\,\omega_m^{jr},
\qquad j\in\ZZ/m\ZZ,
\]
则
\[
\widehat{\Delta^{(\ell)}}(0)=0,
\qquad
\widehat{\Delta^{(\ell)}}(j)=2\,\omega_m^{j\ell}
\quad (j\neq 0).
\]
换言之，centered residue 向量在非平凡 Fourier 模上恰给出等模相位轨道。

记零均值超平面
\[
\mathcal H_m:=\left\{x\in\RR^{\ZZ/m\ZZ}:\ \sum_{r\in\ZZ/m\ZZ}x_r=0\right\}.
\]
则成立：
\begin{enumerate}
  \item 对每个 \(\ell\) 都有 \(\Delta^{(\ell)}\in\mathcal H_m\)；
  \item 内积满足精确公式
  \[
  \boxed{\;
  \bigl\langle \Delta^{(\ell)},\Delta^{(\ell')}\bigr\rangle
  =
  4\left(\delta_{\ell,\ell'}-\frac{1}{m}\right).\;}
  \]
  \item 因而族 \(\{\Delta^{(\ell)}\}_{\ell\in\ZZ/m\ZZ}\) 在 \(\mathcal H_m\) 中构成一个边长恒定的正则单纯形顶点集；
  \item 它们张成整个 \((m-1)\)-维超平面 \(\mathcal H_m\)。
\end{enumerate}
\end{theorem}
\begin{proof}
由定义，
\[
\sum_{r\in\ZZ/m\ZZ}\Delta_r^{(\ell)}
=
2-\frac{2}{m}\cdot m
=
0,
\]
故 \(\Delta^{(\ell)}\in\mathcal H_m\)。

再计算内积。若 \(\ell=\ell'\)，则
\[
\bigl\|\Delta^{(\ell)}\bigr\|_2^2
=
\left(2-\frac{2}{m}\right)^2
+(m-1)\left(\frac{2}{m}\right)^2
=
4\left(1-\frac{1}{m}\right).
\]
若 \(\ell\neq \ell'\)，则恰有两个坐标分别取值
\[
2-\frac{2}{m}
\quad\text{与}\quad
-\frac{2}{m},
\]
其余 \(m-2\) 个坐标都等于 \(-2/m\)。因此
\[
\bigl\langle \Delta^{(\ell)},\Delta^{(\ell')}\bigr\rangle
=
2\left(2-\frac{2}{m}\right)\left(-\frac{2}{m}\right)
+(m-2)\left(\frac{2}{m}\right)^2
=
-\frac{4}{m}.
\]
合并两种情形，得到统一公式
\[
\bigl\langle \Delta^{(\ell)},\Delta^{(\ell')}\bigr\rangle
=
4\left(\delta_{\ell,\ell'}-\frac{1}{m}\right).
\]

于是全部向量具有相同范数、两两内积恒定，故在其仿射张成空间中构成正则单纯形。又
\[
\sum_{\ell\in\ZZ/m\ZZ}\Delta^{(\ell)}=0,
\]
故它们至多张成 \((m-1)\) 维空间。其 Gram 矩阵正是
\[
4\left(I-\frac{1}{m}J\right),
\]
而该矩阵的秩为 \(m-1\)。因此其线性张成空间维数恰为 \(m-1\)，并且只能等于同维的零均值超平面 \(\mathcal H_m\)。证毕。
\end{proof}

\begin{corollary}[单原子 centered residue 扰动在零均值子空间上的各向同性白噪声律]\label{cor:xi-time-part65c-atomic-centered-residue-isotropic}
\leanverified{paper\_xi\_time\_part65c\_atomic\_centered\_residue\_isotropic}
沿用定理 \ref{thm:xi-time-part65c-atomic-centered-residue-simplex} 的记号，定义均匀平均协方差算子
\[
\Sigma_m:=
\frac{1}{m}\sum_{\ell\in\ZZ/m\ZZ}
\Delta^{(\ell)}\bigl(\Delta^{(\ell)}\bigr)^{\mathsf T}.
\]
则有严格闭式
\[
\boxed{\;
\Sigma_m=\frac{4}{m}\left(I-\frac{1}{m}J\right),\;}
\]
其中 \(J\) 为全 \(1\) 矩阵。于是：
\begin{enumerate}
  \item \(\Sigma_m\) 在零均值超平面 \(\mathcal H_m\) 上恰等于
  \[
  \frac{4}{m}I_{\mathcal H_m},
  \]
  因而完全各向同性；
  \item 对任意 \(v\in\mathcal H_m\)，都有精确能量恒等式
  \[
  \boxed{\;
  \frac{1}{m}\sum_{\ell\in\ZZ/m\ZZ}
  \bigl|\langle v,\Delta^{(\ell)}\rangle\bigr|^2
  =
  \frac{4}{m}\|v\|_2^2.\;}
  \]
\end{enumerate}
换言之，single atom 在 centered residue 空间中的周期轨道平均不偏向任何方向；其二次统计贡献在 \(\mathcal H_m\) 上表现为一个完全各向同性的白噪声核。
\end{corollary}
\begin{proof}
由
\[
\Delta^{(\ell)}=2e_\ell-\frac{2}{m}\mathbf 1,
\]
可得
\[
\Sigma_m
=
\frac{1}{m}\sum_{\ell}
\left(2e_\ell-\frac{2}{m}\mathbf 1\right)
\left(2e_\ell-\frac{2}{m}\mathbf 1\right)^{\mathsf T}.
\]
逐项展开：
\[
\frac{1}{m}\sum_{\ell}4e_\ell e_\ell^{\mathsf T}=\frac{4}{m}I,
\]
\[
\frac{1}{m}\sum_{\ell}\left(
-\frac{4}{m}e_\ell\mathbf 1^{\mathsf T}
-\frac{4}{m}\mathbf 1 e_\ell^{\mathsf T}
\right)
=
-\frac{8}{m^2}J,
\]
\[
\frac{1}{m}\sum_{\ell}\frac{4}{m^2}\mathbf 1\mathbf 1^{\mathsf T}
=
\frac{4}{m^2}J.
\]
合并得
\[
\Sigma_m=\frac{4}{m}I-\frac{4}{m^2}J
=
\frac{4}{m}\left(I-\frac{1}{m}J\right).
\]
这便得到第一式。

若 \(v\in\mathcal H_m\)，则 \(Jv=0\)。因此
\[
v^{\mathsf T}\Sigma_m v
=
\frac{4}{m}v^{\mathsf T}v
=
\frac{4}{m}\|v\|_2^2.
\]
另一方面，
\[
v^{\mathsf T}\Sigma_m v
=
\frac{1}{m}\sum_{\ell}
\bigl|\langle v,\Delta^{(\ell)}\rangle\bigr|^2.
\]
两式合并即得所示能量恒等式。证毕。
\end{proof}

\begin{theorem}[归一化 periodic residue law 偏离均匀分布的精确指数]\label{thm:xi-time-part65c-normalized-residue-law-exact-exponent}
沿用结论 \ref{con:xi-time-part57ba-core-full-residue-rankone-perturbation} 与推论 \ref{cor:xi-single-atom-centered-residue-l2-exponent-truncation} 的记号。固定模数 \(m\ge 2\)，定义归一化 periodic residue law
\[
\mu_{m,r}^{\mathrm{full}}(n)
:=
\frac{A_{m,r}^{\mathrm{full}}(n)}{a_n^{\mathrm{full}}(1)},
\qquad
\mu_{m,\bullet}^{\mathrm{full}}(n)
:=
\bigl(\mu_{m,r}^{\mathrm{full}}(n)\bigr)_{r\in\ZZ/m\ZZ},
\qquad
u_m:=\frac{1}{m}\mathbf 1_m.
\]
\leanverified{paper\_xi\_time\_part65c\_normalized\_residue\_law\_exact\_exponent}
再记
\[
\lambda:=\rho(M(1))=\varphi^2,
\qquad
\rho_m^{\mathrm{core}}
:=
\max_{1\le j\le m-1}\rho\!\bigl(M_{\mathrm{core}}(\omega_m^j)\bigr).
\]
则对 \(\CC^{\ZZ/m\ZZ}\) 上任意固定向量范数 \(\|\cdot\|\)，都有
\[
\boxed{\;
\limsup_{n\to\infty}
\bigl\|
\mu_{m,\bullet}^{\mathrm{full}}(n)-u_m
\bigr\|^{1/n}
=
\frac{\max\{1,\rho_m^{\mathrm{core}}\}}{\lambda}.\;}
\]
\end{theorem}
\begin{proof}
记 centered residue 向量
\[
\widetilde A_{m,\bullet}^{\mathrm{full}}(n)
:=
\Bigl(
A_{m,r}^{\mathrm{full}}(n)-\frac{a_n^{\mathrm{full}}(1)}{m}
\Bigr)_{r\in\ZZ/m\ZZ}.
\]
则
\[
\mu_{m,\bullet}^{\mathrm{full}}(n)-u_m
=
\frac{\widetilde A_{m,\bullet}^{\mathrm{full}}(n)}{a_n^{\mathrm{full}}(1)}.
\]
由推论 \ref{cor:xi-single-atom-centered-residue-l2-exponent-truncation}，
\[
\limsup_{n\to\infty}
\bigl\|\widetilde A_{m,\bullet}^{\mathrm{full}}(n)\bigr\|_2^{1/n}
=
\max\{1,\rho_m^{\mathrm{core}}\}.
\]
由于 \(\CC^{\ZZ/m\ZZ}\) 有限维，任意两个固定范数等价，故上式中的指数不依赖于范数选择，从而
\[
\limsup_{n\to\infty}
\bigl\|\widetilde A_{m,\bullet}^{\mathrm{full}}(n)\bigr\|^{1/n}
=
\max\{1,\rho_m^{\mathrm{core}}\}.
\]
另一方面，无扭曲 periodic 总数满足
\[
\lim_{n\to\infty}\bigl(a_n^{\mathrm{full}}(1)\bigr)^{1/n}
=
\rho(M(1))
=
\lambda,
\]
这是无扭曲 Perron 根的直接结果。将上述两式相除即可得到题述结论。证毕。
\end{proof}

\begin{theorem}[原子主导区中归一化 residue law 的单纯形首项轨道]\label{thm:xi-time-part65c-atom-dominated-normalized-residue-profile}
沿用定理 \ref{thm:xi-time-part65c-atomic-centered-residue-simplex} 与定理 \ref{thm:xi-time-part65c-normalized-residue-law-exact-exponent} 的记号。若
\[
\rho_m^{\mathrm{core}}<1,
\]
则对偶数 \(n=2q\)，记 \(r_q\in\ZZ/m\ZZ\) 为 \(q\) 的剩余类，有
\[
\boxed{\;
\frac{a_{2q}^{\mathrm{full}}(1)}{2}
\bigl(\mu_{m,\bullet}^{\mathrm{full}}(2q)-u_m\bigr)
=
e_{r_q}-u_m+o(1)
=
\frac12\Delta^{(r_q)}+o(1),\;}
\]
而对奇数 \(n\) 有
\[
\boxed{\;
a_n^{\mathrm{full}}(1)\bigl(\mu_{m,\bullet}^{\mathrm{full}}(n)-u_m\bigr)=o(1).\;}
\]
因此在 \(\rho_m^{\mathrm{core}}<1\) 的原子主导区中，归一化 residue law 的首项偏差并非任意零均值向量，而是沿定理 \ref{thm:xi-time-part65c-atomic-centered-residue-simplex} 的正则单纯形顶点轨道
\[
\{e_r-u_m:\ r\in\ZZ/m\ZZ\}
\]
周期移动。
\end{theorem}
\begin{proof}
由结论 \ref{con:xi-time-part57ba-core-full-residue-rankone-perturbation}，对偶数 \(n=2q\) 有
\[
A_{m,r}^{\mathrm{full}}(2q)
=
A_{m,r}^{\mathrm{core}}(2q)+2\,\mathbf 1_{\{r\equiv q\;(\mathrm{mod}\,m)\}},
\qquad
a_{2q}^{\mathrm{full}}(1)=a_{2q}^{\mathrm{core}}(1)+2.
\]
因此
\[
\mu_{m,\bullet}^{\mathrm{full}}(2q)-u_m
=
\frac{1}{a_{2q}^{\mathrm{full}}(1)}
\Bigl(
\widetilde A_{m,\bullet}^{\mathrm{core}}(2q)+2(e_{r_q}-u_m)
\Bigr).
\]
另一方面，由 \(\rho_m^{\mathrm{core}}<1\) 与推论 \ref{cor:xi-single-atom-centered-residue-l2-exponent-truncation} 的第一条，
\[
\bigl\|\widetilde A_{m,\bullet}^{\mathrm{core}}(n)\bigr\|_2=o(1),
\]
故由范数等价亦有
\[
\bigl\|\widetilde A_{m,\bullet}^{\mathrm{core}}(n)\bigr\|=o(1)
\]
在任意固定范数下成立。把上式乘以 \(a_{2q}^{\mathrm{full}}(1)/2\) 即得偶数情形。

当 \(n\) 为奇数时，单原子 residue 修正消失，故
\[
\mu_{m,\bullet}^{\mathrm{full}}(n)-u_m
=
\frac{\widetilde A_{m,\bullet}^{\mathrm{core}}(n)}{a_n^{\mathrm{full}}(1)}.
\]
再乘以 \(a_n^{\mathrm{full}}(1)\) 并利用同一 \(o(1)\) 估计，即得奇数情形。最后，由定理 \ref{thm:xi-time-part65c-atomic-centered-residue-simplex} 中
\(
\Delta^{(r)}=2(e_r-u_m)
\)
的定义，首项轨道正是同一定理给出的正则单纯形顶点轨道。证毕。
\end{proof}

\begin{theorem}[输出计数中心窗口的一阶 Edgeworth 修正与严格左偏]\label{thm:xi-time-part65c-output-edgeworth-first-correction}
沿用命题 \ref{prop:real-input-40-output-cumulants-closed} 与定理 \ref{thm:xi-real-input-40-output-edgeworth-quasipowers} 的记号。记
\[
\alpha_0:=P'(0)=0.367376207875073606\ldots,
\qquad
\sigma^2:=P''(0)=0.0487412456183376392\ldots,
\]
\[
\kappa_3:=P^{(3)}(0)=-0.0158881374258564278\ldots<0,
\qquad
\kappa_4:=P^{(4)}(0)=0.00568478997567274949\ldots>0.
\]
若 \(X_n\) 为定理 \ref{thm:xi-real-input-40-output-edgeworth-quasipowers} 中的输出 \(1\) 总数，并定义标准化变量
\[
Y_n:=\frac{X_n-n\alpha_0}{\sigma\sqrt n},
\]
则在中心窗口内一致有
\[
\boxed{\;
\PP(Y_n\le x)
=
\Phi(x)
+\frac{\kappa_3}{6\sigma^3\sqrt n}(1-x^2)\phi(x)
+O(n^{-1}).\;}
\]
特别地，标准化偏度常数
\[
\boxed{\;
\gamma_1:=\frac{\kappa_3}{\sigma^3}<0.\;}
\]
因此输出计数在一阶 Edgeworth 层面严格左偏。
\end{theorem}
\begin{proof}
定理 \ref{thm:xi-real-input-40-output-edgeworth-quasipowers} 已给出
\[
\PP\!\left(
\frac{X_n-n\mu}{\sqrt{n\kappa_2}}
\le x
\right)
=
\Phi(x)
+\frac{\kappa_3}{6\kappa_2^{3/2}\sqrt n}(1-x^2)\phi(x)
+O(n^{-1})
\]
在中心窗口内一致成立。将其中 \(\mu,\kappa_2\) 分别改写为 \(\alpha_0,\sigma^2\) 即得第一式。

再由定理 \ref{thm:xi-real-input-40-output-edgeworth-nongaussian} 与表 \ref{tab:real_input_40_output_potential_cumulants_closed}，
\[
\kappa_3<0,
\]
从而 \(\gamma_1=\kappa_3/\sigma^3<0\)。证毕。
\end{proof}

\begin{corollary}[Edgeworth 首项与局部大偏差三次项的同源符号锁定]\label{cor:xi-time-part65c-edgeworth-local-ldp-sign-lock}
沿用定理 \ref{thm:xi-time-part65c-output-edgeworth-first-correction} 与定理 \ref{thm:xi-output-potential-local-legendre-birkhoff-normal-form} 的记号。则
\[
\operatorname{sign}\!\left(\frac{\kappa_3}{6\sigma^3}\right)
=
\operatorname{sign}\!\left(\frac{\kappa_3}{6\sigma^6}\right)
=
\operatorname{sign}(\kappa_3).
\]
等价地，中心极限层的一阶 Edgeworth 偏斜与局部 Legendre 正规形中的三次偏斜常数具有同一符号源。对当前核，由 \(\kappa_3<0\) 可得
\[
I_{\mathrm{out}}(\alpha_0+\delta)-I_{\mathrm{out}}(\alpha_0-\delta)
=
-\frac{\kappa_3}{3\sigma^6}\delta^3+O(\delta^5)>0
\qquad (\delta>0\ \text{充分小}),
\]
故中心窗口严格左偏，而局部大偏差层的右偏离严格更昂贵。
\end{corollary}
\begin{proof}
由定理 \ref{thm:xi-time-part65c-output-edgeworth-first-correction}，一阶 Edgeworth 系数为 \(\kappa_3/(6\sigma^3)\)。由定理 \ref{thm:xi-output-potential-local-legendre-birkhoff-normal-form}，
\[
I_{\mathrm{out}}(\alpha_0+\delta)
=
\frac{\delta^2}{2\sigma^2}
-
\left(\frac{\kappa_3}{6\sigma^6}\right)\delta^3
+O(\delta^4).
\]
因此控制两种首阶奇次不对称的常数分别为 \(\kappa_3/(6\sigma^3)\) 与 \(\kappa_3/(6\sigma^6)\)，它们显然具有相同符号。再对 \(\delta\) 与 \(-\delta\) 相减，即得
\[
I_{\mathrm{out}}(\alpha_0+\delta)-I_{\mathrm{out}}(\alpha_0-\delta)
=
-\frac{\kappa_3}{3\sigma^6}\delta^3+O(\delta^5).
\]
由 \(\kappa_3<0\) 立得题述不等式。证毕。
\end{proof}

\noindent
以上诸结果把冻结相 pressure/escort 链、单原子 residue 手术与输出位势的首阶非高斯修正进一步压缩到同一层：一方面，escort--Rényi 熵率不仅在 \((a,s)\)-平面上沿 \(as=a_{\mathrm c}\) 发生严格冻结，而且在更强半直线 \(a>a_{\mathrm c,s}\) 上把有限尺度主项提升为 \(\log K_m\)，并把最大纤维扇区对预言机容量的贡献压缩为指数下界 \(g_*+\min\{\alpha_*,\beta\log 2\}\) 及其相应侧信息临界斜率；另一方面，单原子的一维局域手术不仅在 centered residue 空间中生成正则单纯形，而且把归一化 periodic residue law 的指数精确锁定为 \(\max\{1,\rho_m^{\mathrm{core}}\}/\lambda\)，并在 \(\rho_m^{\mathrm{core}}<1\) 时强迫首项沿该单纯形顶点轨道移动；最后，输出计数在中心窗口中的左偏 Edgeworth 修正与局部大偏差的右侧代价增强又由同一三阶累积量 \(\kappa_3\) 严格统一。由此，冻结热力学、原子局域手术与首阶非高斯非对称便在“冻结边界 + 有限尺度主项 + 侧信息斜率 + 单纯形轨道 + 偏斜常数”的共同语言下闭合为同一条可审计主线。


\paragraph{宏观共振模、周期化绝对连续障碍与熵缺口尺度失配}
在 Fibonacci/Lucas 共振子列的正极限、Pisot--Bernoulli 周期化障碍、最大纤维偏置、谱零点稀疏律以及碰撞/KL 恒等式已经建立之后，还可进一步压缩出下列六条彼此衔接的直接后果。它们分别刻画宏观共振模的对数丰度、周期化测度的非 Rajchman 与绝对连续性排斥、最大纤维相对均值的无界放大、相对均匀基线的 \(\sqrt{M}\)-尺度偏离刚性、零点相对密度对碰撞尺度的根本失效，以及 Shannon 熵缺口在自然 \(1/M\) 标度上的发散。

\begin{theorem}[宏观共振模的对数丰度]\label{thm:xi-time-part66-macro-resonance-log-abundance}
沿用定理 \ref{thm:fold-critical-resonance-constant-integer-ladder} 与推论
\ref{cor:fold-resonance-ladder-fib-luc-limits} 的记号，记
\[
c_{\mathrm{Fib}}:=\lim_{n\to\infty}C_\varphi(F_n)>0.
\]
对任意常数
\[
0<\eta<c_{\mathrm{Fib}},
\]
存在 \(N_\eta\) 与 \(c_\eta>0\)，使得对所有 \(N\ge N_\eta\) 都有
\[
\boxed{\;
\#\{1\le u\le N:\ C_\varphi(u)\ge \eta\}
\ge
c_\eta\log N.\;}
\]
因此，非零共振峰并非孤立稀见事件，而是在对数尺度上持续出现的宏观族。
\end{theorem}
\leanverified{paper\_xi\_time\_part66\_macro\_resonance\_log\_abundance}
\begin{proof}
由推论 \ref{cor:fold-resonance-ladder-fib-luc-limits}，
\[
C_\varphi(F_n)\longrightarrow c_{\mathrm{Fib}}>0.
\]
故对任意 \(0<\eta<c_{\mathrm{Fib}}\)，存在 \(n_0\) 使得对所有 \(n\ge n_0\) 都有
\[
C_\varphi(F_n)\ge \eta.
\]
于是当 \(N\) 充分大时，
\[
\#\{1\le u\le N:\ C_\varphi(u)\ge \eta\}
\ge
\#\{n\ge n_0:\ F_n\le N\}.
\]
再由 Binet 公式，
\[
\#\{n:\ F_n\le N\}
=
\frac{\log N}{\log\varphi}+O(1).
\]
因此存在常数 \(c_\eta>0\) 与 \(N_\eta\)，使得对所有 \(N\ge N_\eta\) 都有
\[
\#\{1\le u\le N:\ C_\varphi(u)\ge \eta\}\ge c_\eta\log N.
\]
证毕。
\end{proof}

\begin{corollary}[\texorpdfstring{$\mu_{\varphi^{-1}}^{(2)}$}{mu} 的模 \texorpdfstring{$2$}{2} 周期化既非 Rajchman，亦无任何 \texorpdfstring{$L^p$}{Lp} 密度]\label{cor:xi-time-part66-periodized-bernoulli-nonrajchman-no-lp}
令
\[
q:=\varphi^{-1},
\]
并令 \(\mu_q^{(2)}\) 如定理 \ref{thm:fold-pisot-bernoulli-convolution-representation}。记
\[
\nu:=\bigl(\RR\to \mathbb T_2=\RR/(2\ZZ)\bigr)_\#\mu_q^{(2)}
\]
为其模 \(2\) 周期化测度。则：
\begin{enumerate}
  \item \(\nu\) 的 Fourier 系数不趋于 \(0\)，故 \(\nu\) 不是 Rajchman 测度；
  \item \(\nu\) 不可能具有任何 \(L^p(\mathbb T_2)\) 密度，其中 \(1\le p\le \infty\)。
\end{enumerate}
\end{corollary}
\begin{proof}
由推论 \ref{cor:fold-resonance-ladder-fib-luc-limits}，记
\[
c_{\mathrm{Fib}}:=\lim_{n\to\infty}C_\varphi(F_n)>0.
\]
再由定理 \ref{thm:xi-pisot-bernoulli-fibonacci-direction-nonrajchman}，
\[
\bigl|\widehat\nu(F_n)\bigr|
=
\bigl|\widehat{\mu_q^{(2)}}(\pi F_n)\bigr|
=
C_\varphi(F_n)
\longrightarrow
c_{\mathrm{Fib}}>0,
\]
故 \(\widehat\nu(u)\not\to 0\)，从而 \(\nu\) 不是 Rajchman 测度。

若 \(\nu\) 具有某个 \(L^1(\mathbb T_2)\) 密度，则由 Riemann--Lebesgue 引理，其 Fourier 系数必趋于 \(0\)，与上式矛盾。另一方面，\(\mathbb T_2\) 的 Haar 测度有限，故对任意 \(1\le p\le\infty\) 都有
\[
L^p(\mathbb T_2)\subseteq L^1(\mathbb T_2).
\]
因此 \(\nu\) 不可能具有任何 \(L^p\) 密度。证毕。
\end{proof}
\leanverified{paper\_xi\_time\_part66\_periodized\_bernoulli\_nonrajchman\_no\_lp}

\begin{theorem}[保 Fibonacci 共振的 Fourier 乘子不可能把周期化测度送入 \texorpdfstring{$L^2$}{L2}]\label{thm:xi-time-part66-fibonacci-visible-multiplier-no-l2}
\leanverified{paper\_xi\_time\_part66\_fibonacci\_visible\_multiplier\_no\_l2}
沿用推论 \ref{cor:xi-time-part66-periodized-bernoulli-nonrajchman-no-lp} 的记号。设
\[
\sigma:\ZZ\to \CC
\]
为任意符号，并把 \(T_\sigma\nu\) 视为 Fourier 系数满足
\[
\widehat{T_\sigma\nu}(u)=\sigma(u)\widehat\nu(u)
\qquad (u\in\ZZ)
\]
的周期分布。若存在常数 \(c>0\) 与 \(n_0\)，使得
\[
|\sigma(F_n)|\ge c
\qquad (n\ge n_0),
\]
则 \(T_\sigma\nu\) 不可能由任何 \(L^2(\mathbb T_2)\) 密度表示。
\end{theorem}
\begin{proof}
由推论 \ref{cor:xi-time-part66-periodized-bernoulli-nonrajchman-no-lp} 的证明，
\[
|\widehat\nu(F_n)|=C_\varphi(F_n)\longrightarrow c_{\mathrm{Fib}}>0.
\]
故存在 \(n_1\) 使得对一切 \(n\ge n_1\) 都有
\[
|\widehat\nu(F_n)|\ge \frac{c_{\mathrm{Fib}}}{2}.
\]
于是对任意 \(N\ge \max\{n_0,n_1\}\)，都有
\[
\sum_{u\in\ZZ}\bigl|\widehat{T_\sigma\nu}(u)\bigr|^2
\ge
\sum_{n\ge N}\bigl|\sigma(F_n)\widehat\nu(F_n)\bigr|^2
\ge
\frac{c^2c_{\mathrm{Fib}}^2}{4}\sum_{n\ge N}1
=+\infty.
\]
若 \(T_\sigma\nu\) 由某个 \(f\in L^2(\mathbb T_2)\) 表示，则 Parseval 恒等式给出
\[
\sum_{u\in\ZZ}\bigl|\widehat{T_\sigma\nu}(u)\bigr|^2
=
\sum_{u\in\ZZ}|\widehat f(u)|^2
<
\infty,
\]
矛盾。故 \(T_\sigma\nu\) 不可能具有 \(L^2(\mathbb T_2)\) 密度。证毕。
\end{proof}

\begin{theorem}[最大纤维相对均值的放大因子无界]\label{thm:xi-time-part66-maxfiber-over-mean-diverges}
\leanverified{paper\_xi\_time\_part66\_maxfiber\_over\_mean\_diverges}
沿用定理 \ref{thm:xi-time-part9m-maxfiber-uniformity-gap-chain} 与推论
\ref{cor:fold-collision-scale-diverges} 的记号。记
\[
M:=F_{m+2},
\qquad
p_m(r):=\frac{d_m(r)}{2^m},
\qquad
M_m:=\max_r d_m(r),
\qquad
\overline d_m:=\frac{2^m}{M}.
\]
则对每个 \(m\) 都有
\[
\boxed{\;
\frac{M_m}{\overline d_m}\ge M\,\mathsf{Col}_m.\;}
\]
从而
\[
\boxed{\;
\frac{M_m}{\overline d_m}\longrightarrow +\infty.\;}
\]
换言之，最大纤维相对均值的偏置本身在尺度极限下发散。
\end{theorem}
\begin{proof}
设
\[
p_{\max}:=\max_r p_m(r)=\frac{M_m}{2^m}.
\]
由 \(\sum_r p_m(r)=1\) 可知
\[
\mathsf{Col}_m=\sum_r p_m(r)^2\le p_{\max}\sum_r p_m(r)=p_{\max}.
\]
两边乘以 \(M\)，得到
\[
M\,\mathsf{Col}_m\le M p_{\max}=M\frac{M_m}{2^m}=\frac{M_m}{\overline d_m}.
\]
这就证明了第一条不等式。

再由推论 \ref{cor:fold-collision-scale-diverges}，
\[
M\,\mathsf{Col}_m=F_{m+2}\mathsf{Col}_m\longrightarrow +\infty.
\]
结合上述下界，立即得到
\[
\frac{M_m}{\overline d_m}\longrightarrow +\infty.
\]
证毕。
\end{proof}

\begin{theorem}[零点相对密度永远不是碰撞尺度的充分统计量]\label{thm:xi-time-part66-zero-density-never-sufficient}
\leanverified{paper\_xi\_time\_part66\_zero\_density\_never\_sufficient}
记
\[
M:=F_{m+2},
\qquad
\eta_m:=\frac{|\mathcal Z_m|}{M},
\]
其中 \(\mathcal Z_m\subset \ZZ/(M\ZZ)\) 为谱零点集合。设
\[
\Phi:[0,1]\to [0,\infty)
\]
在 \(0\) 的某邻域内有界。则存在 \(m_0\)，使得对所有 \(m\ge m_0\) 都有
\[
\boxed{\;
M\,\mathsf{Col}_m>\Phi(\eta_m).\;}
\]
因此，任何仅依赖零点相对密度并在 \(0\) 邻域有界的零点几何指标，都不可能控制归一化碰撞尺度。
\end{theorem}
\begin{proof}
先证 \(\eta_m\to 0\)。若 \(3\nmid (m+2)\)，则由定理
\ref{thm:killo-fold-zero-spectrum-affine-congruence}，
\[
\mathcal Z_m=\varnothing,
\qquad
\eta_m=0.
\]
若 \(3\mid (m+2)\)，则 \(M\) 为偶数，定理 \ref{thm:fold-zero-density-sparse} 给出
\[
\eta_m=\frac{|\mathcal Z_m|}{M}
\le
\tau(m+2)\,\frac{F_{\lfloor (m+2)/2\rfloor}}{F_{m+2}}
=
\tau(m+2)\,O(\varphi^{-m/2})\to 0.
\]
故在全体 \(m\) 上都有 \(\eta_m\to 0\)。

由于 \(\Phi\) 在 \(0\) 的某邻域内有界，存在 \(\delta>0\) 与 \(K<\infty\)，使得
\[
0\le t\le \delta\quad\Longrightarrow\quad \Phi(t)\le K.
\]
由 \(\eta_m\to 0\)，存在 \(m_1\) 使得对所有 \(m\ge m_1\) 都有 \(\eta_m\le \delta\)，故
\[
\Phi(\eta_m)\le K.
\]
另一方面，由推论 \ref{cor:fold-collision-scale-diverges}，
\[
M\,\mathsf{Col}_m\longrightarrow +\infty.
\]
故存在 \(m_2\) 使得对所有 \(m\ge m_2\) 都有
\[
M\,\mathsf{Col}_m>K.
\]
取 \(m_0:=\max\{m_1,m_2\}\)，便得
\[
M\,\mathsf{Col}_m>\Phi(\eta_m)
\qquad (m\ge m_0).
\]
证毕。
\end{proof}

\begin{theorem}[相对均匀基线的 \texorpdfstring{$\sqrt{M}$}{sqrtM} 尺度偏离不可消除]\label{thm:xi-time-part66-uniform-baseline-sqrtM-rigidity}
\leanverified{paper\_xi\_time\_part66\_uniform\_baseline\_sqrtm\_rigidity}
沿用结论 \ref{con:xi-fold-collision-spectral-trace-parseval} 与结论
\ref{con:xi-fold-chi2-constant-residual-energy} 的记号。记
\[
M:=F_{m+2},
\qquad
p_m(r):=\frac{d_m(r)}{2^m},
\qquad
u(r):=\frac{1}{M}
\qquad
\bigl(r\in \ZZ/(M\ZZ)\bigr).
\]
则有
\[
\boxed{\;
\liminf_{m\to\infty}\sqrt{M}\,\|p_m-u\|_{\ell^2}
\ge
\sqrt{2}\,C_\varphi.\;}
\]
以及
\[
\boxed{\;
\liminf_{m\to\infty}\sqrt{M}\,\|p_m-u\|_{\mathrm{TV}}
\ge
\frac{C_\varphi}{\sqrt{2}}.\;}
\]
因此，折叠推前分布相对均匀基线的偏离并非 \(o(M^{-1/2})\) 级扰动，而在中心极限定标下保留严格正的黄金常数下界。
\end{theorem}
\begin{proof}
由结论 \ref{con:xi-fold-collision-spectral-trace-parseval}，
\[
\|p_m-u\|_{\ell^2}^2
:=
\sum_{r}\bigl(p_m(r)-u(r)\bigr)^2
=
\mathsf{Col}_m-\frac1M.
\]
另一方面，结论 \ref{con:xi-fold-chi2-constant-residual-energy} 给出
\[
\liminf_{m\to\infty}M\Bigl(\mathsf{Col}_m-\frac1M\Bigr)\ge 2C_\varphi^2.
\]
两式合并即得
\[
\liminf_{m\to\infty}M\,\|p_m-u\|_{\ell^2}^2\ge 2C_\varphi^2,
\]
从而
\[
\liminf_{m\to\infty}\sqrt{M}\,\|p_m-u\|_{\ell^2}\ge \sqrt2\,C_\varphi.
\]
再由
\[
\|p_m-u\|_{\mathrm{TV}}
=
\frac12\|p_m-u\|_{\ell^1}
\ge
\frac12\|p_m-u\|_{\ell^2},
\]
便得到
\[
\liminf_{m\to\infty}\sqrt{M}\,\|p_m-u\|_{\mathrm{TV}}
\ge
\frac{1}{2}\cdot \sqrt2\,C_\varphi
=
\frac{C_\varphi}{\sqrt2}.
\]
证毕。
\end{proof}

\begin{theorem}[Shannon 熵缺口在 \texorpdfstring{$1/M$}{1/M} 标度上发散]\label{thm:xi-time-part66-shannon-gap-times-M-diverges}
\leanverified{paper\_xi\_time\_part66\_shannon\_gap\_times\_m\_diverges}
沿用结论 \ref{con:xi-fold-collision-spectral-trace-parseval} 与推论
\ref{cor:fold-collision-scale-diverges} 的记号。记
\[
M:=F_{m+2},
\qquad
p_m(r):=\frac{d_m(r)}{2^m},
\qquad
u(r):=\frac{1}{M}.
\]
则有下界
\[
\boxed{\;
D_{\mathrm{KL}}(p_m\|u)\ge \frac12\Bigl(\mathsf{Col}_m-\frac1M\Bigr).\;}
\]
特别地，
\[
\boxed{\;
M\bigl(\log M-H(p_m)\bigr)
\ge
\frac12\bigl(M\,\mathsf{Col}_m-1\bigr)
\longrightarrow +\infty.\;}
\]
因此，相对均匀分布的 Shannon 熵缺口乘以自然尺度 \(M\) 后并不趋于有限常数，而是发散到无穷。
\end{theorem}
\begin{proof}
由结论 \ref{con:xi-fold-collision-spectral-trace-parseval}，
\[
\sum_r\bigl(p_m(r)-u(r)\bigr)^2
=
\mathsf{Col}_m-\frac1M.
\]
另一方面，Pinsker 不等式给出
\[
D_{\mathrm{KL}}(p_m\|u)\ge 2\,\|p_m-u\|_{\mathrm{TV}}^2.
\]
而
\[
\|p_m-u\|_{\mathrm{TV}}
=
\frac12\|p_m-u\|_{\ell^1}
\ge
\frac12\|p_m-u\|_{\ell^2},
\]
故
\[
D_{\mathrm{KL}}(p_m\|u)
\ge
2\cdot \frac14\|p_m-u\|_{\ell^2}^2
=
\frac12\sum_r\bigl(p_m(r)-u(r)\bigr)^2
=
\frac12\Bigl(\mathsf{Col}_m-\frac1M\Bigr).
\]
再由
\[
D_{\mathrm{KL}}(p_m\|u)=\log M-H(p_m),
\]
两边乘以 \(M\)，得到
\[
M\bigl(\log M-H(p_m)\bigr)
\ge
\frac12\bigl(M\,\mathsf{Col}_m-1\bigr).
\]
最后由推论 \ref{cor:fold-collision-scale-diverges}，\(M\,\mathsf{Col}_m\to+\infty\)，故右端趋于 \(+\infty\)。证毕。
\end{proof}

\noindent
这六条结果把共振闭包常数、周期化 Fourier 采样、谱零点集合、碰撞尺度与 Shannon 熵缺口几条主线真正闭合起来：一方面，Fibonacci 共振并非只提供个别正极限，而在阈值意义下形成对数丰度的宏观模族，并迫使周期化测度同时失去 Rajchman 性与全部绝对连续可能；另一方面，碰撞尺度的发散不仅会被最大纤维相对均值直接放大到无界量级，还表明相对均匀基线的 \(\ell^2\) 与总变差偏离在 \(\sqrt{M}\) 定标下保留严格正下界，并且任何局部有界的零点密度指标都不足以控制归一化碰撞强度；最终，这一非均匀性又在信息论层面外显为 \(M(\log M-H(p_m))\to+\infty\) 的熵缺口发散。由此，有限窗口折叠的非均匀性便不再能被解释为少数零点位置的几何副产物，而必须被视为一族对数丰度的宏观共振模在振幅、峰值与熵三重层面上的共同驱动。

\paragraph{角色投影、平滑塌缩与外置算术核的函子分裂}
在模剩余 Fourier 反演、window-\(6\) 群商障碍、末位充分化、输出位势的解析腔室刚性、二步原子的 primitive/finite-part 位移、零温四态地基层分裂以及局部化数系的 Pontryagin 对偶链已经建立之后，还可把离散波纹与连续平滑之间的关系进一步压缩为下列七条彼此闭合的结论。它们分别刻画非平凡角色扇区、非自由群商障碍、末位二点终端塌缩、单解析相内的波滑分裂、两步原子的 Abel 有限部位移、零温平滑包络的有限状态化，以及连续相位层之后仍不可消去的 profinite 素数核。

\begin{theorem}[角色投影的扭曲扇区分裂]\label{thm:xi-time-part67-character-projection-twisted-sectors}
\leanverified{paper\_xi\_time\_part67\_character\_projection\_twisted\_sectors}
固定整数 \(q\ge 2\)。设 \(B_{q,a}(n)\) 为定义在 \(a\in \ZZ/(q\ZZ)\) 上的计数函数，并令
\[
\widehat B_{q,j}(n):=\sum_{a\in \ZZ/(q\ZZ)}B_{q,a}(n)\,\omega_q^{ja},
\qquad
\omega_q:=e^{2\pi i/q}.
\]
则有 Fourier 反演
\[
B_{q,a}(n)=\frac1q\sum_{j=0}^{q-1}\widehat B_{q,j}(n)\,\omega_q^{-ja},
\qquad
\frac1q\sum_{b\in \ZZ/(q\ZZ)}B_{q,b}(n)=\frac1q\,\widehat B_{q,0}(n),
\]
因而中心化波纹恰为
\[
B_{q,a}(n)-\frac1q\,\widehat B_{q,0}(n)
=
\frac1q\sum_{j=1}^{q-1}\widehat B_{q,j}(n)\,\omega_q^{-ja}.
\]
对真实输入 \(40\) 状态核取 \(B_{q,a}(n)=A_{q,a}(n)\) 或 \(B_{q,a}(n)=P_{q,a}(n)\)，则
\[
\widehat A_{q,j}(n)=a_n(\omega_q^j),
\qquad
\widehat P_{q,j}(n)=p_n(\omega_q^j),
\]
从而平滑均值恰是平凡角色扇区，而全部离散波纹严格由非平凡单位根扭曲 \(M(\omega_q^j)\) 承载。进一步，若记
\[
\rho_q:=\max_{1\le j\le q-1}\rho\!\bigl(M(\omega_q^j)\bigr),
\qquad
\theta_q:=\frac{\log\rho_q}{\log\lambda},
\qquad
\lambda=\varphi^2,
\]
则这些非平凡扇区在 periodic 与 primitive 两层均以 \(O(\rho_q^n)\) 衰减；并且存在明确模数 \(q\) 使 \(\theta_q>\frac12\)，且该性质沿倍数塔向上遗传。
\end{theorem}
\begin{proof}
有限循环群 \(\ZZ/(q\ZZ)\) 上的离散 Fourier 反演直接给出前两式；中心化公式只需把 \(j=0\) 项单独提出。

对真实输入 \(40\) 状态核，定理 \ref{thm:xi-dirichlet-residue-fourier-parseval-exact} 已在 periodic 与 primitive 两层分别给出
\[
a_n(\omega_q^j)=\sum_{a\in \ZZ/(q\ZZ)}A_{q,a}(n)\,\omega_q^{ja},
\qquad
p_n(\omega_q^j)=\sum_{a\in \ZZ/(q\ZZ)}P_{q,a}(n)\,\omega_q^{ja},
\]
故扭曲 trace 正是 residue counts 的非平凡角色分量。命题 \ref{prop:real-input-40-output-dirichlet-nonrh} 进一步给出：这些扭曲分量在 periodic/primitive 两层都以 \(O(\rho_q^n)\) 衰减；并且存在明确模数使 \(\theta_q>\frac12\)，其中 \(q=2\) 已足够。最后，由结论 \ref{con:xi-time-part9l-dirichlet-twist-divisibility-monotonicity}，一旦某个模数满足 \(\theta_q>\frac12\)，则其全部倍数也都满足同一严格不等式。证毕。
\end{proof}

\begin{theorem}[window-\texorpdfstring{$6$}{6} 平滑包络的非自由群商障碍]\label{thm:xi-time-part67-window6-nonfree-group-quotient-obstruction}
设
\[
\Omega_6:=\{0,1\}^6,
\qquad
F_6:=\Fold_6^{\mathrm{bin}}:\Omega_6\to X_6.
\]
若试图把 \(F_6\) 的纤维压缩解释为某个固定有限群的轨道商，或解释为某个阿贝尔线性陪集划分，则该方案必失败。更具体地，\(F_6\) 的纤维直方图
\[
2:8,\qquad 3:4,\qquad 4:9
\]
迫使任何与之同轨道的群作用都必须带三层稳定子分层；而在保持 fold-gauge 纤维独立置换语义时，完备异常签名仍需
\[
\Aut(\Fold_6^{\mathrm{bin}})^{\mathrm{ab}}\cong (\ZZ/2\ZZ)^{21}.
\]
因此，审计窗中的平滑包络所消去的并非单一统一对称冗余，而只能是随纤维变化的稳定子残差；其自然承载对象是带变稳定子的群胚型商，而非单一群商。
\end{theorem}
\begin{proof}
定理 \ref{thm:xi-window6-group-quotient-isotropy-strata} 已证明：任何把 \(F_6\) 纤维分解实现为群轨道分解的方案，都不可能是自由作用，也不可能退化为等势轨道；并且轨道大小只能取 \(2,3,4\)，从而稳定子恰分成三层。于是，固定自由群商与阿贝尔线性陪集解释同时被排除。

另一方面，定理 \ref{thm:xi-window6-oracle-anomaly-coordinate-bifurcation} 证明：若继续保持 fold-gauge 的纤维独立置换语义，则完备异常信息的最小承载对象仍为
\(
(\ZZ/2\ZZ)^{21}
\)；
任何把异常坐标压到 \(21\) 以下的方案都必须引入跨纤维关系，从而破坏单一纤维直积语义。故这里被“平滑”掉的并不是某个固定群的统一冗余，而是随纤维变化的稳定子残差；用群胚而非单一群来描述，正是这一变稳定子机制的自然形式化。证毕。
\end{proof}

\begin{theorem}[末位投影的渐近充分性与终端二点塌缩]\label{thm:xi-time-part67-lastbit-terminal-csiszar-collapse}
设
\[
\ell_m(w):=w_m,
\qquad
\mu_m(w):=\frac{d_m^{\mathrm{bin}}(w)}{2^m},
\qquad
u_m(w):=\frac{1}{|X_m|}
\qquad (w\in X_m).
\]
则 \(\ell_m\) 对分布对 \((\mu_m,u_m)\) 是渐近充分统计量：对任意凸函数 \(f:(0,\infty)\to\RR\) 满足 \(f(1)=0\)，有
\[
D_f(\mu_m\Vert u_m)
=
D_f\!\bigl((\ell_m)_\#\mu_m\Vert (\ell_m)_\#u_m\bigr)
+O_f\!\Bigl(\Bigl(\frac{\varphi}{2}\Bigr)^m\Bigr).
\]
并且若记
\[
p_*:=\frac{1}{\varphi\sqrt5},
\qquad
q_*:=\varphi^{-2},
\]
则
\[
\lim_{m\to\infty}D_f(\mu_m\Vert u_m)
=
D_f\!\bigl(\mathrm{Bernoulli}(p_*)\Vert \mathrm{Bernoulli}(q_*)\bigr)
=
\frac1\varphi\,f\!\left(\frac{\varphi^2}{\sqrt5}\right)
+
\frac1{\varphi^2}\,f\!\left(\frac{\varphi}{\sqrt5}\right).
\]
换言之，相对于稳定层均匀基线的全部渐近 Csisz\'ar 几何，都在末位二点实验上终端塌缩。
\end{theorem}
\begin{proof}
定理 \ref{thm:xi-fold-lastbit-asymptotic-sufficiency} 已在全变差与 Le Cam 意义下给出末位投影 \(\ell_m\) 的渐近充分性；下面只需把这一充分化提升到 \(f\)-散度层面。

由定理 \ref{thm:xi-time-part64b-fold-information-geometry-two-point-collapse}，
\[
\frac{\mu_m(w)}{u_m(w)}
=
\frac{\varphi^{2-\ell_m(w)}}{\sqrt5}
+O\!\Bigl(\Bigl(\frac{\varphi}{2}\Bigr)^m\Bigr)
\qquad (w\in X_m)
\]
在 \(w\) 上一致成立。于是存在紧区间 \(K\subset(0,\infty)\)，使得对充分大的 \(m\)，上式右端与左端都落在 \(K\) 内。由于凸函数在紧子区间上局部 Lipschitz，故对每个固定 \(f\) 有
\[
D_f(\mu_m\Vert u_m)
=
\sum_{b\in\{0,1\}}u_m(\ell_m=b)\,
f\!\left(\frac{\varphi^{2-b}}{\sqrt5}\right)
+O_f\!\Bigl(\Bigl(\frac{\varphi}{2}\Bigr)^m\Bigr).
\]

另一方面，设
\[
p_m:=((\ell_m)_\#\mu_m)(\{1\}),
\qquad
q_m:=((\ell_m)_\#u_m)(\{1\})=\frac{F_m}{F_{m+2}}.
\]
对上式在每个末位块上取 \(u_m\)-平均，得到
\[
\frac{p_m}{q_m}
=
\frac{\varphi}{\sqrt5}
+O\!\Bigl(\Bigl(\frac{\varphi}{2}\Bigr)^m\Bigr),
\qquad
\frac{1-p_m}{1-q_m}
=
\frac{\varphi^2}{\sqrt5}
+O\!\Bigl(\Bigl(\frac{\varphi}{2}\Bigr)^m\Bigr),
\]
从而
\[
D_f\!\bigl((\ell_m)_\#\mu_m\Vert (\ell_m)_\#u_m\bigr)
=
q_m\,f\!\left(\frac{p_m}{q_m}\right)
+(1-q_m)\,f\!\left(\frac{1-p_m}{1-q_m}\right)
\]
与前一式只差 \(O_f((\varphi/2)^m)\)。这便证明了渐近充分性。

最后，由
\[
q_m=\frac{F_m}{F_{m+2}}\longrightarrow \frac1{\varphi^2},
\qquad
1-q_m=\frac{F_{m+1}}{F_{m+2}}\longrightarrow \frac1{\varphi},
\]
以及上面的两个比值极限，立即得到
\[
\lim_{m\to\infty}D_f\!\bigl((\ell_m)_\#\mu_m\Vert (\ell_m)_\#u_m\bigr)
=
\frac1\varphi\,f\!\left(\frac{\varphi^2}{\sqrt5}\right)
+\frac1{\varphi^2}\,f\!\left(\frac{\varphi}{\sqrt5}\right).
\]
再与刚证得的误差估计合并，即得所示极限公式。证毕。
\end{proof}

\begin{theorem}[单一解析腔室内的波滑分裂]\label{thm:xi-time-part67-single-analytic-chamber-wave-smooth-splitting}
在真实输入 \(40\) 状态核的输出位势参数化
\[
u=e^\theta
\]
下，审计窗口
\[
|\theta|<\theta_0,
\qquad
\theta_0=0.0533477827\ldots,
\]
满足
\[
u\in (e^{-\theta_0},e^{\theta_0})
\Subset
(0,u_R),
\qquad
u_R\approx 3.59262181344.
\]
因而该窗口既不跨越正实退化点，也完全停留在
\[
\Lambda(u)<\lambda_+(u)^{1/2}
\]
的统一亚临界区间内。对于任何可归约为解析 \(\arg\max\) 读出的折叠谱观测量，若其在此窗口内出现离散波纹与连续平滑的分裂，则该分裂必发生于同一解析相内部，而不能归因于谱分支碰撞或正实相变。
\end{theorem}
\begin{proof}
推论 \ref{cor:xi-collision-u1-audit-window-uniform-subcriticality} 已证明：
\[
|\theta|<\theta_0
\Longrightarrow
u=e^\theta\in(0.932837\ldots,1.054796\ldots)\Subset(0,u_R),
\]
从而整个 \(u=1\) 邻域审计窗都处在统一亚临界腔室内，不会接触平方根门槛 \(\Lambda(u)=\lambda_+(u)^{1/2}\) 的临界端点。

另一方面，命题 \ref{prop:xi-time-part9m3-analytic-chamber-rigidity-sampling-readout} 说明：在任何不穿越退化点的实解析区间上，凡由解析 \(\arg\max\) 读出定义的纤维计数、容量与矩谱，都只能在离散交叉集之外分段常值。于是，在本窗口内若仍观察到“波纹--平滑”二分，则其来源只能是同一解析腔室内部不同读出函子的分裂，而不能被解释为谱支交换或正实相变。证毕。
\end{proof}

\begin{theorem}[两步原子剥离与 Abel 有限部位移]\label{thm:xi-time-part67-two-step-atomic-peeling-finitepart-shift}
对真实输入 \(40\) 状态核，有精确分解
\[
\Delta(z,u)=(1-uz^2)\,F(z,u),
\qquad
P(z,u)=uz^2+P_{\mathrm{core}}(z,u).
\]
在未加权情形 \(u=1\) 下，若 \(z_\star=\varphi^{-2}\)，并记 \(\mathfrak M,\mathfrak M_{\mathrm{core}}\) 分别为 full/core 的 Abel 有限部常数，则
\[
\log \mathfrak M
=
\log \mathfrak M_{\mathrm{core}}+\varphi^{-4},
\qquad
\varphi^{-4}=\frac{7-3\sqrt5}{2}.
\]
因此，由离散波纹统计转入平滑统计并非中性平均，而是一次精确的两步 primitive Euler 原子剥离。
\end{theorem}
\begin{proof}
定理 \ref{thm:xi-time-part65-atomic-primitive-abel-shift} 已直接给出
\[
P(z,u)=uz^2+P_{\mathrm{core}}(z,u),
\qquad
\log\mathfrak M=\log\mathfrak M_{\mathrm{core}}+\varphi^{-4}.
\]
其中 \(\varphi^{-4}=(7-3\sqrt5)/2\) 是同一定理中的闭式常数。于是 full/core 之差在 primitive 层只作用于长度 \(2\) 的单一原子，而在 Abel 有限部层只留下一个完全闭式的代数位移。证毕。
\end{proof}

\begin{theorem}[零温平滑包络的有限状态化]\label{thm:xi-time-part67-zero-temp-envelope-finite-stateization}
对真实输入 \(40\) 状态核的输出位势 \(P(\theta)=\log\lambda(e^\theta)\)，有
\[
P'(\theta)\longrightarrow \frac12
\qquad (\theta\to+\infty),
\]
并且满足平方根端点律
\[
\frac12-P'(\theta)=O(e^{-\theta/2}).
\]
另一方面，在零温缩放
\[
z=\sqrt{w/u}
\]
下，
\[
\lim_{u\to+\infty}\Delta\!\left(\sqrt{w/u},u\right)
=
(1-w)\det(I-wB_\infty),
\]
其中
\[
B_\infty=
\begin{pmatrix}
0&1&1&0\\
0&0&1&0\\
0&0&3&1\\
1&0&0&3
\end{pmatrix}.
\]
因此，零温极限中被选出的平滑包络并非整个 \(40\) 状态核的粗粒化投影，而是由显式四状态地基 SFT 所控制的有限状态对象；原子层仅留下因子 \((1-w)\)，而主导极点由地基层而非原子层决定。
\end{theorem}
\begin{proof}
定理 \ref{thm:xi-output-potential-zero-temp-square-root-critical-law} 给出
\[
\frac12-P'(\theta)=\frac{d}{2c}e^{-\theta/2}+O(e^{-\theta}),
\]
故零温端点的饱和斜率确为 \(1/2\)。

再由定理 \ref{thm:xi-zero-temp-atomic-ground-exact-splitting}，
\[
\lim_{u\to+\infty}\Delta\!\left(\sqrt{w/u},u\right)
=(1-w)\det(I-wB_\infty),
\]
这说明零温解析对象严格分裂为一个 primitive 原子因子与一个四状态地基 SFT 的动力学行列式。最后，推论 \ref{cor:xi-zero-temp-leading-pole-ground-dominance} 进一步表明正实最内层零点来自 \(\det(I-wB_\infty)\) 而非 \((1-w)\)，故主导平滑尺度由四状态地基层决定。证毕。
\end{proof}

\begin{theorem}[连续相位层与外置算术核]\label{thm:xi-time-part67-continuous-phase-external-arithmetic-kernel}
设 \(\varnothing\neq S\subset\mathbb P\) 为有限素数集，并记
\[
G_S:=\ZZ[S^{-1}].
\]
则其 Pontryagin 对偶满足短正合列
\[
0\longrightarrow \prod_{p\in S}\ZZ_p
\longrightarrow \widehat{G_S}
\longrightarrow \TT
\longrightarrow 0,
\]
并且该 profinite 核 \(\prod_{p\in S}\ZZ_p\) 可由相位采样函数 \(\Sigma^{\mathrm{ph}}_{G_S}\) 唯一恢复。因而圆相位因子 \(\TT\) 只是连通可见层，真正的算术内容则以外置 prime-register 核的形式保留在 profinite 方向上。

更进一步，不存在任何有限秩加法账本能够把正整数乘法幺半群忠实严格同态化；因此，上述 profinite 核不能被有限多个加法寄存器吸收。
\end{theorem}
\begin{proof}
定理 \ref{thm:xi-localized-phase-sampling-closed-form} 已证明：\(\Sigma^{\mathrm{ph}}_{G_S}\) 唯一恢复素数集 \(S\)，并因而唯一恢复对偶短正合列中的 profinite 核
\[
\prod_{p\in S}\ZZ_p.
\]
定理 \ref{thm:xi-localized-basic-extension-strictly-nonsplit} 进一步说明该短正合列在 \(S\neq\varnothing\) 时严格不分裂，因此 \(\TT\) 不能替代整个对偶对象；它只给出连通商，而非全部算术信息。

最后，定理 \ref{thm:xi-time-part64b-local-additivity-nonlocal-multiplicativity-polarization} 已排除一切把
\(
(\NN_{\times},\times)
\)
严格忠实同态化到有限秩加法账本的方案。故这里的 prime-register profinite 核确实是不可由有限加法寄存器消化的外置算术部分。证毕。
\end{proof}

\noindent
上述七条结果表明，离散间隔统计中的波纹并不应被理解为连续对象自身的粗糙化。在本文框架内，它们恰是非平凡角色扇区、window-\(6\) 的变稳定子残差、两步 primitive 原子与 prime-register profinite 核在离散读出层的共同显影；相应地，平滑极限则由平凡角色投影、末位二点塌缩、Abel 有限部反项剥离以及零温地基层选择共同构成。因而，“波”与“滑”并非同一对象的两种视觉描画，而是同一算术动力学对象在不同函子层上的两种范畴像。

\paragraph{谱临界、输出慢模、端点有限部与熵缺损常数的统一闭合}
在碰撞位势的解析半径、输出位势的四阶累积结构、地基层端点能垒、两步 primitive 原子、极端偏置有限部以及 bin-fold 的信息位率阈值都已建立之后，还可把这些分散结论进一步压缩为下列七条彼此闭合的终端后果。它们分别刻画碰撞 RH 临界与解析分歧的严格分离、输出模扭曲慢模的四阶刚性、端点代价与地基熵的精确配平、输出大偏差相图的双坐标代数化、七条分支二回路在 primitive 层的单原子塌缩、极端偏置有限部对单一四次代数参数的依赖，以及熵缺损常数在三种信息尺度上的同一化。

\begin{theorem}[碰撞 RH 临界先于解析分歧出现]\label{thm:xi-time-part68-collision-rh-critical-before-analytic-branch}
设
\[
P_2(\theta):=\log \rho(e^\theta),
\]
并记 \(R_\theta\) 为其在 \(\theta=0\) 处 Taylor 级数的收敛半径。又设 \(u_R>1\) 为核版 RH/Ramanujan 条件
\[
\Lambda(u)\le \lambda_1(u)^{1/2}
\]
失效的唯一正临界点，\(\theta_R:=\log u_R\)。则
\[
0<\theta_R<R_\theta.
\]
更具体地，
\[
u_R\approx 3.59262181344,\qquad
\theta_R\approx 1.27888,\qquad
R_\theta\approx 2.76230289346087.
\]
因此，碰撞核的 RH/Ramanujan 失效并非由解析延拓的首个分歧点触发，而是在解析域内部由实谱不等式
\[
\Lambda(u)\le \lambda_1(u)^{1/2}
\]
的翻转所触发；该临界属于纯谱型临界，而非分支型临界。
\end{theorem}
\begin{proof}
定理 \ref{thm:xi-rh-critical-before-nearest-complex-branch} 已证明：\(P_2\) 的全部代数分歧点恰由判别式方程
\[
4u^3+25u^2+88u+8=0
\]
的三根经
\[
\theta=\log u_j+2\pi i k
\qquad (k\in\ZZ)
\]
产生，且最近分歧点的模长正是
\[
R_\theta\approx 2.76230289346087.
\]
同一定理又给出：\(u_R\) 是五次方程
\[
u^5+4u^4+3u^3-96u^2-42u-14=0
\]
的唯一正根，\(\theta_R=\log u_R\approx 1.27888\)，并且
\[
\textsf{RH}(u)\Longleftrightarrow 0<u\le u_R.
\]
比较数值便得
\[
0<\theta_R<R_\theta.
\]
于是对任意 \(\theta\in(\theta_R,R_\theta)\)，一方面 \(|\theta|<R_\theta\)，故 \(P_2(\theta)\) 仍位于 \(\theta=0\) 的解析收敛盘内；另一方面 \(u=e^\theta>u_R\)，故 RH/Ramanujan 条件已经失效。再由同一定理中的
\[
\Lambda(u_R)=\lambda_1(u_R)^{1/2},
\]
可知临界边界确由谱不等式的等号翻转刻画，而非由解析分支奇性触发。证毕。
\end{proof}

\begin{theorem}[输出模扭曲慢模的四阶累积刚性]\label{thm:xi-time-part68-output-twist-slowmode-fourth-cumulant-rigidity}
记
\[
\lambda(u):=\rho(M(u)),
\qquad
\lambda:=\lambda(1)=\varphi^2,
\qquad
P(\theta):=\log \lambda(e^\theta),
\qquad
\kappa_2:=P''(0),
\qquad
\kappa_4:=P^{(4)}(0).
\]
则
\[
\kappa_2,\kappa_4\in\QQ(\sqrt5),
\qquad
\kappa_2>0.
\]
对主单位根扭曲通道设
\[
t_m:=\frac{2\pi}{m},
\qquad
\rho_{m,1}:=\rho\!\bigl(M(e^{it_m})\bigr).
\]
则有严格渐近展开
\[
\log \frac{\rho_{m,1}}{\lambda}
=
-\frac{\kappa_2}{2}\Bigl(\frac{2\pi}{m}\Bigr)^2
+\frac{\kappa_4}{24}\Bigl(\frac{2\pi}{m}\Bigr)^4
+O(m^{-6}).
\]
并且在已审计模数 \(m=5,10,20\) 上，四阶预测的绝对误差分别约为
\[
1.08\times 10^{-6},\qquad
2.78\times 10^{-8},\qquad
6.21\times 10^{-10},
\]
而仅保留二阶项时的相应误差分别约为
\[
5.67\times 10^{-4},\qquad
3.66\times 10^{-5},\qquad
2.30\times 10^{-6}.
\]
因此，输出模扭曲的近平凡慢模在四阶以内已被前四阶累积量刚性锁定；四阶项并非可忽略修饰，而是恢复慢模几何所必需的主导校正。
\end{theorem}
\begin{proof}
命题 \ref{prop:real-input-40-output-cumulants-closed} 已给出
\[
P^{(k)}(0)\in\QQ(\sqrt5),
\]
特别是 \(\kappa_2,\kappa_4\in\QQ(\sqrt5)\) 且 \(\kappa_2>0\)。再由定理 \ref{thm:xi-output-potential-dirichlet-fixed-j-fourth-order} 在 \(j=1\) 时的结论，
\[
\log \frac{\rho_{m,1}}{\lambda}
=
-\frac{\kappa_2}{2}\Bigl(\frac{2\pi}{m}\Bigr)^2
+\frac{\kappa_4}{24}\Bigl(\frac{2\pi}{m}\Bigr)^4
+O(m^{-6}),
\]
这就给出所示四阶展开。

另一方面，附录注 \ref{rem:real-input-40-dirichlet-cumulant-error-order} 已逐项记录 \(m=5,10,20\) 上的四阶预测误差与纯二阶高斯近似误差，恰为题述数值。再由定理 \ref{thm:xi-output-potential-dirichlet-fourth-order-gain}，
\[
\frac{\widetilde\varepsilon_m^{(2)}}{\widetilde\varepsilon_m^{(4)}}=\Omega(m^2),
\]
故四阶修正不仅改善常数，而且在阶数量级上严格压缩误差。于是四阶项确为慢模重建所必需的主导校正。证毕。
\end{proof}

\begin{proposition}[端点代价缺口恰等于地基熵]\label{prop:xi-time-part68-endpoint-gap-ground-entropy-identity}
\leanverified{paper\_xi\_time\_part68\_endpoint\_gap\_ground\_entropy\_identity}
设 \(I(\alpha)\) 为输出位势
\[
P(\theta)=\log\lambda(e^\theta)
\]
的归一化大偏差速率函数，并记
\[
\alpha_{\max}:=\lim_{\theta\to+\infty}P'(\theta)=\frac12,
\qquad
h_{\mathrm{ground}}:=\log c.
\]
则有精确恒等式
\[
I(0)-I(\alpha_{\max})=h_{\mathrm{ground}}.
\]
更具体地，
\[
I(0)=\log(\varphi^2),
\qquad
I\!\left(\frac12\right)=\log\!\left(\frac{\varphi^2}{c}\right),
\qquad
h_{\mathrm{ground}}=\log c.
\]
因此，右端点的大偏差代价相对于左端点的全部减免，恰由地基层的组合熵指数精确补偿。
\end{proposition}
\begin{proof}
定理 \ref{thm:xi-output-potential-endpoint-gap-ground-state} 已直接给出
\[
I(0)-I\!\left(\frac12\right)=h_{\mathrm{ground}}.
\]
而推论 \ref{cor:real-input-40-alpha-max} 说明
\[
\alpha_{\max}=\frac12,
\]
推论 \ref{cor:real-input-40-ground-entropy} 则给出
\[
h_{\mathrm{ground}}=\log c.
\]
再结合推论 \ref{cor:real-input-40-rate-endpoints} 中的端点闭式
\[
I(0)=\log(\varphi^2),
\qquad
I(1/2)=\log(\varphi^2/c),
\]
即可得到题述恒等式与具体数值表达。证毕。
\end{proof}

\begin{theorem}[输出大偏差相图的代数一致化]\label{thm:xi-time-part68-output-ldp-phase-diagram-algebraization}
设 \(u=e^\theta\)，\(\lambda(u)\) 为方程
\[
G(\lambda,u)=0
\]
的 Perron 正根，并定义平衡态下输出 \(1\) 的典型密度
\[
\alpha(u):=\frac{\mathrm d}{\mathrm d\theta}\log\lambda(e^\theta).
\]
则存在非零多项式
\[
F(\alpha,u)\in\QQ[\alpha,u],
\qquad
H(\alpha,\lambda)\in\QQ[\alpha,\lambda]
\]
使得沿 Perron 分支恒有
\[
F(\alpha(u),u)=0,
\qquad
H(\alpha(u),\lambda(u))=0.
\]
并且归一化大偏差速率函数满足参数公式
\[
I(\alpha(u))
=
\alpha(u)\log u-\log\frac{\lambda(u)}{\lambda(1)}.
\]
因此，输出大偏差相图在 \((\alpha,u)\) 与 \((\alpha,\lambda)\) 两组坐标中都可代数化；真正的超越性只保留在 Legendre 对偶中的单一对数项。
\end{theorem}
\begin{proof}
命题 \ref{prop:real-input-40-pressure-freezing} 已给出输出位势的特征多项式
\[
G(\lambda,u)\in\QQ[\lambda,u],
\]
并说明 \(\lambda(u)\) 是其 Perron 正根。另一方面，附录 \ref{app:real-input-40-zeta-u} 中的输出位势大偏差代数参数化已把 \(\alpha(u)\) 写成
\[
\alpha(u)=\left.\frac{N(\lambda,u)}{D(\lambda,u)}\right|_{\lambda=\lambda(u)},
\qquad
N,D\in\QQ[\lambda,u].
\]
因此沿 Perron 分支有
\[
\alpha\,D(\lambda,u)-N(\lambda,u)=0,
\qquad
G(\lambda,u)=0.
\]
对 \(\lambda\) 取结果式并约去共同因子，便得到某个非零多项式
\[
F(\alpha,u)\in\QQ[\alpha,u]
\]
满足 \(F(\alpha(u),u)=0\)。同理，对 \(u\) 取结果式并约去共同因子，可得某个非零多项式
\[
H(\alpha,\lambda)\in\QQ[\alpha,\lambda]
\]
满足 \(H(\alpha(u),\lambda(u))=0\)。

最后，由同一附录中的参数公式，或者直接由严格凸压力的驻点 Legendre 公式，
\[
I(\alpha(u))
=
\alpha(u)\log u-\log\frac{\lambda(u)}{\lambda(1)}.
\]
于是 \((\alpha,u)\) 与 \((\alpha,\lambda)\) 两种坐标下的相图都由代数关系控制，而超越部分只出现在该单一对数项中。证毕。
\end{proof}

\begin{theorem}[七条分支二回路在 primitive 层塌缩为唯一原子]\label{thm:xi-time-part68-seven-branch-two-cycles-collapse-single-atom}
设输出位势的 primitive 生成函数为
\[
P(z,u)=\sum_{n\ge 1}p_n(u)z^n.
\]
则在真实输入 \(40\) 状态核上有精确分解
\[
P(z,u)=uz^2+P_{\mathrm{core}}(z,u).
\]
另一方面，branch 层审计表明，满足
\[
(|\tau|,E)=(2,1)
\]
的两步回路恰有 \(7\) 条。因而，这 \(7\) 条几何上可区分的 branch 二回路在 primitive--Witt 化之后并不留下 \(7\) 个独立原子，而是完全塌缩为唯一长度 \(2\) 原子
\[
uz^2.
\]
\end{theorem}
\begin{proof}
定理 \ref{thm:xi-time-part62d-atomic-two-step-witt-rankone-collapse} 已直接给出
\[
P(z,u)=uz^2+P_{\mathrm{core}}(z,u),
\]
并同时指出 branch 层中满足
\[
(|\tau|,E)=(2,1)
\]
的两步回路总数恰为 \(7\)。逐项比较系数可知，full/core 的 primitive 差分只支撑在长度 \(2\) 的单一坐标上；等价地，
\[
\dim_{\QQ(u)}
\operatorname{span}\{\,p_n(u)-p_n^{\mathrm{core}}(u):n\ge 1\,\}=1.
\]
因此，这七条 branch 两步回路在 primitive--Witt 层只能留下唯一原子 \(uz^2\)。证毕。
\end{proof}

\begin{theorem}[极端偏置有限部由单一四次代数参数完全控制]\label{thm:xi-time-part68-extreme-bias-finitepart-single-quartic-control}
记
\[
b:=b_\infty:=\rho(B_\infty)^{-1}\in(0,1),
\]
等价地，\(b\) 是四次方程
\[
1-6b+9b^2-b^3-b^4=0
\]
在区间 \((0,1)\) 内的最小正根。则极端偏置 \(u\to+\infty\) 下，主极点常数 \(C_\infty\) 与 Abel 有限部常数 \(\mathfrak M_\infty\) 满足
\[
C_\infty=
\frac{1}{2b(1-b)(6-18b+3b^2+4b^3)},
\]
\[
\log\mathfrak M_\infty
=
\log C_\infty
-\sum_{m\ge 2}\frac{\mu(m)}{m}
\log\Bigl[(1-b^m)\bigl(1-6b^m+9b^{2m}-b^{3m}-b^{4m}\bigr)\Bigr].
\]
数值上
\[
b\approx 0.27807480214113,\qquad
C_\infty\approx 1.89745149363,\qquad
\log\mathfrak M_\infty\approx 0.282290255262.
\]
因此，极端偏置下的全部有限部数据都由单一四次代数参数 \(b\) 完整统辖。
\end{theorem}
\begin{proof}
命题 \ref{prop:real-input-40-logM-infty} 已直接证明：令
\[
b=\lim_{u\to\infty}u\,z_\star(u)^2=\frac{1}{c^2},
\]
则 \(b\) 满足题述四次方程，并且 \(C_\infty\) 与 \(\log\mathfrak M_\infty\) 的闭式恰为上式。另一方面，推论 \ref{cor:xi-zero-temp-leading-pole-ground-dominance} 说明
\[
b=b_\infty=\rho(B_\infty)^{-1}
\]
正是零温缩放极限中控制主导正实零点的位置参数。故无论从主极点常数还是从 Abel 有限部常数看，极端偏置下的主导有限部数据都确实由同一个四次代数参数 \(b\) 完整决定。证毕。
\end{proof}

\begin{theorem}[熵缺损常数的三重同一化]\label{thm:xi-time-part68-entropy-deficit-triple-identification}
记
\[
\mu_m(w):=\frac{d_m^{\mathrm{bin}}(w)}{2^m},
\qquad
W_m\sim\mu_m,
\qquad
\gamma_m:=\frac{1}{m}\log d_m^{\mathrm{bin}}(W_m).
\]
则对任意 \(q>0\)（其中 \(q=1\) 按 Shannon 熵解释）都有
\[
\lim_{m\to\infty}\frac{1}{m}\bigl(m\log 2-H_q(\mu_m)\bigr)
=
\log\frac{2}{\varphi},
\qquad
\gamma_m\longrightarrow \log\frac{2}{\varphi}.
\]
若一列二进制预算 \(B_m\in\NN\) 满足
\[
\mathrm{Succ}_m^{\mathrm{bin}}(B_m)\ge 1-\varepsilon
\qquad
(0<\varepsilon<1\ \text{固定}),
\]
则必有
\[
\liminf_{m\to\infty}\frac{B_m}{m}
\ge
\log_2\frac{2}{\varphi}.
\]
因此，同一个常数 \(\log(2/\varphi)\) 同时控制全部正阶 Rényi 熵的广延缺损、典型微态的纤维多重度指数以及保持固定成功率所需的最小外置比特率。
\end{theorem}
\begin{proof}
由定理 \ref{thm:fold-bin-renyi-rate-collapse}，
\[
\frac{1}{m}H_q(\mu_m)\longrightarrow \log\varphi
\qquad (q>0),
\]
并且
\[
\gamma_m\longrightarrow \log\frac{2}{\varphi}.
\]
于是
\[
\frac{1}{m}\bigl(m\log 2-H_q(\mu_m)\bigr)
=
\log 2-\frac{1}{m}H_q(\mu_m)
\longrightarrow
\log 2-\log\varphi
=
\log\frac{2}{\varphi},
\]
这就得到前两条极限。

另一方面，若用 \(B_m\) 个二进制外置信息位实现恢复，则其状态数为 \(M_m=2^{B_m}\)。推论 \ref{cor:xi-reconstruction-cardinality-strong-converse} 对 \(\delta=1-\varepsilon\) 给出
\[
B_m
\ge
m-\log_2F_{m+2}+\log_2(1-\varepsilon).
\]
两边同除以 \(m\)，再用 Binet 渐近
\[
\log_2F_{m+2}=m\log_2\varphi+O(1),
\]
便得
\[
\liminf_{m\to\infty}\frac{B_m}{m}\ge \log_2\frac{2}{\varphi}.
\]
故题述三种量确由同一个常数统一控制。证毕。
\end{proof}

\noindent
上述七条结果表明，碰撞核的 RH 破裂、输出单位根慢模、地基层端点能垒、两步 primitive 原子、零温有限部常数与 bin-fold 的熵缺损阈值并不是彼此分离的局部现象。它们共同揭示：一方面，碰撞侧的临界首先发生在解析盘内部，因而属于纯谱型翻转；另一方面，输出侧的慢模、端点代价与大偏差相图都受同一条代数压力曲线及其四阶累积结构支配；最终，这种谱--代数刚性又在 primitive 原子剥离、零温有限部与信息预算下界中外显为同一组可审计的不变量。因而，从碰撞 RH 窗口到输出大偏差，再到部分可逆预算与极端偏置有限部，全文中的主导常数与临界门槛实际上共享同一条谱--组合--信息几何闭合链。

\paragraph{低预算线性相、二点实验塌缩、局部化周期谱与形式线性化的闭合后果}
在审计偶窗的最小退化 Fibonacci 律、末位两层塌缩、局部化 solenoid 的自覆盖分类、素数寄存器的幂等估值算术化以及 fold-gauge 缺陷的纯奇次重整化已经建立之后，还可把这些链条继续压缩为下列七条结果。它们分别刻画审计偶窗中的低预算严格线性段及其首个偏离阈值、bin-fold 输出实验的 Le Cam 二点塌缩、一般局部化有理乘法自同态的固定点与 Artin--Mazur \(\zeta\) 函数、相应周期点几何的指数率、素数寄存器幂等元的精确闭式计数，以及 fold-gauge 缺陷的形式 Koenigs 线性化。

\begin{theorem}[bin-fold 审计偶数窗上的低预算严格线性律]\label{thm:xi-time-part69-audited-even-pure-pigeonhole-linear-phase}
对审计偶数窗
\[
m\in\{6,8,10,12\},
\]
定义
\[
d_{\min}(m):=\min_{w\in X_m}d_m^{\mathrm{bin}}(w),
\qquad
C_m^{\mathrm{bin}}(B):=\sum_{w\in X_m}\min\bigl(d_m^{\mathrm{bin}}(w),2^B\bigr),
\]
\[
\mathrm{Succ}_m^{\mathrm{bin}}(B):=\frac{C_m^{\mathrm{bin}}(B)}{2^m}.
\]
若整数预算 \(B\ge 0\) 满足
\[
2^B\le d_{\min}(m)=F_{m/2},
\]
则有精确公式
\[
C_m^{\mathrm{bin}}(B)=2^B\,F_{m+2},
\qquad
\mathrm{Succ}_m^{\mathrm{bin}}(B)=2^{B-m}F_{m+2}.
\]
因此，在最小退化阈值之前，bin-fold 的每增加一比特辅助信息都会使可逆微态总数严格翻倍，且不出现任何次级修正。
\end{theorem}
\begin{proof}
定理 \ref{thm:fold-bin-min-degeneracy-fib-audited} 已证明：当
\[
m\in\{6,8,10,12\}
\]
时，
\[
d_{\min}(m)=F_{m/2}.
\]
若 \(2^B\le F_{m/2}\)，则对所有 \(w\in X_m\) 都有
\[
2^B\le d_{\min}(m)\le d_m^{\mathrm{bin}}(w),
\]
从而
\[
\min\bigl(d_m^{\mathrm{bin}}(w),2^B\bigr)=2^B.
\]
因此
\[
C_m^{\mathrm{bin}}(B)
=
\sum_{w\in X_m}2^B
=
2^B|X_m|.
\]
另一方面，由稳定词计数律
\[
|X_m|=F_{m+2},
\]
即得
\[
C_m^{\mathrm{bin}}(B)=2^BF_{m+2}.
\]
再除以 \(2^m\)，得到
\[
\mathrm{Succ}_m^{\mathrm{bin}}(B)=2^{B-m}F_{m+2}.
\]
证毕。
\end{proof}

\begin{corollary}[bin-fold 容量曲线的首个折点恰位于 \texorpdfstring{$F_{m/2}$}{F(m/2)}]\label{cor:xi-time-part69-binfold-first-kink}
\leanverified{paper\_xi\_time\_part69\_binfold\_first\_kink}
在定理 \ref{thm:xi-time-part69-audited-even-pure-pigeonhole-linear-phase} 的设定下，容量曲线
\[
B\longmapsto C_m^{\mathrm{bin}}(B)
\]
的首个偏离线性区间发生在最小整数
\[
B_{\mathrm{kink}}(m):=\min\{B\in\NN:\ 2^B>F_{m/2}\}
=
\bigl\lfloor \log_2 F_{m/2}\bigr\rfloor+1.
\]
更准确地，
\[
C_m^{\mathrm{bin}}(B)=2^BF_{m+2}
\quad\Longleftrightarrow\quad
2^B\le F_{m/2},
\]
而一旦 \(2^B>F_{m/2}\)，便有严格不等式
\[
C_m^{\mathrm{bin}}(B)<2^BF_{m+2}.
\]
\end{corollary}
\begin{proof}
“若”方向即定理 \ref{thm:xi-time-part69-audited-even-pure-pigeonhole-linear-phase}。反之，由 \(X_m\) 的有限性与
\[
d_{\min}(m)=F_{m/2}
\]
可取 \(w_0\in X_m\) 使
\[
d_m^{\mathrm{bin}}(w_0)=F_{m/2}.
\]
若 \(2^B>F_{m/2}\)，则
\[
\min\bigl(d_m^{\mathrm{bin}}(w_0),2^B\bigr)=F_{m/2}<2^B,
\]
而其余各项至多为 \(2^B\)。于是
\[
C_m^{\mathrm{bin}}(B)<2^B|X_m|=2^BF_{m+2}.
\]
故首个折点恰由阈值 \(F_{m/2}\) 决定。证毕。
\end{proof}

\begin{theorem}[bin-fold 统计实验的 Le Cam 二点塌缩]\label{thm:xi-time-part69-binfold-lecam-two-point-collapse}
\leanverified{paper\_xi\_time\_part69\_binfold\_lecam\_two\_point\_collapse}
设
\[
\mu_m(w):=\frac{d_m^{\mathrm{bin}}(w)}{2^m},
\qquad
u_m(w):=\frac{1}{|X_m|},
\qquad
\ell_m(w):=w_m\in\{0,1\},
\]
并定义两块
\[
A_0:=\{w\in X_m:\ w_m=0\},
\qquad
A_1:=\{w\in X_m:\ w_m=1\}.
\]
令
\[
T_m:=\ell_m:X_m\to\{0,1\},
\qquad
K_m(i,\cdot):=\mathrm{Unif}(A_i)\qquad(i=0,1).
\]
再定义极限二点实验
\[
Q_\infty:=\left(\frac1\varphi,\frac1{\varphi^2}\right),
\qquad
P_\infty:=\left(\frac{\varphi}{\sqrt5},\frac{1}{\varphi\sqrt5}\right),
\]
其中坐标顺序对应 \((0,1)\)。则原始两点统计实验
\[
\mathcal E_m:=\{u_m,\mu_m\}\quad\text{定义在 }X_m
\]
与二元实验
\[
\mathcal E_\infty:=\{Q_\infty,P_\infty\}\quad\text{定义在 }\{0,1\}
\]
在 Le Cam 意义下渐近等价。更具体地，
\[
D_{\mathrm{TV}}\!\bigl(T_{m\#}u_m,Q_\infty\bigr)=O(\varphi^{-2m}),
\qquad
D_{\mathrm{TV}}\!\bigl(T_{m\#}\mu_m,P_\infty\bigr)=O\!\Bigl(\Bigl(\frac{\varphi}{2}\Bigr)^m\Bigr),
\]
并且
\[
D_{\mathrm{TV}}\!\bigl(u_m,K_{m\#}Q_\infty\bigr)=O(\varphi^{-2m}),
\qquad
D_{\mathrm{TV}}\!\bigl(\mu_m,K_{m\#}P_\infty\bigr)=O\!\Bigl(\Bigl(\frac{\varphi}{2}\Bigr)^m\Bigr).
\]
因此
\[
\Delta(\mathcal E_m,\mathcal E_\infty)\xrightarrow[m\to\infty]{}0.
\]
\end{theorem}
\begin{proof}
先看均匀基线 \(u_m\)。由
\[
|A_0|=F_{m+1},
\qquad
|A_1|=F_m,
\qquad
|X_m|=F_{m+2},
\]
可得
\[
T_{m\#}u_m=
\left(\frac{F_{m+1}}{F_{m+2}},\frac{F_m}{F_{m+2}}\right)
\]
且由 Binet 公式比较相邻 Fibonacci 比值可知，其两坐标分别与 \(Q_\infty\) 的对应坐标相差 \(O(\varphi^{-2m})\)。特别地，
\[
D_{\mathrm{TV}}\!\bigl(T_{m\#}u_m,Q_\infty\bigr)=O(\varphi^{-2m}).
\]
另一方面，\(u_m\) 在每个块 \(A_i\) 上本来就是严格均匀分布，因此
\[
u_m=K_{m\#}(T_{m\#}u_m).
\]
由 Markov 核的收缩性，
\[
D_{\mathrm{TV}}\!\bigl(u_m,K_{m\#}Q_\infty\bigr)
\le
D_{\mathrm{TV}}\!\bigl(T_{m\#}u_m,Q_\infty\bigr)
=
O(\varphi^{-2m}).
\]

再看 \(\mu_m\)。定理 \ref{thm:xi-fold-lastbit-asymptotic-sufficiency} 已给出
\[
D_{\mathrm{TV}}\!\bigl(\mu_m,K_{m\#}(T_{m\#}\mu_m)\bigr)
=
O\!\Bigl(\Bigl(\frac{\varphi}{2}\Bigr)^m\Bigr).
\]
另一方面，定理 \ref{thm:xi-fold-lastbit-two-layer-collapse} 证明 \(\mu_m\) 与末位两层模型 \(\nu_m(w)=\varphi^{-m-w_m}\) 在全变差意义下相差
\[
O\!\Bigl(\Bigl(\frac{\varphi}{2}\Bigr)^m\Bigr).
\]
而由
\[
\nu_m(A_0)=\varphi^{-m}F_{m+1}=\frac{\varphi}{\sqrt5}+O(\varphi^{-2m}),
\qquad
\nu_m(A_1)=\varphi^{-m-1}F_m=\frac{1}{\varphi\sqrt5}+O(\varphi^{-2m}),
\]
再与
\[
D_{\mathrm{TV}}(\mu_m,\nu_m)=O\!\Bigl(\Bigl(\frac{\varphi}{2}\Bigr)^m\Bigr)
\]
合并，可知
\[
D_{\mathrm{TV}}\!\bigl(T_{m\#}\mu_m,P_\infty\bigr)
=
O\!\Bigl(\Bigl(\frac{\varphi}{2}\Bigr)^m\Bigr).
\]
于是再用 Markov 核收缩性，
\[
\begin{aligned}
D_{\mathrm{TV}}\!\bigl(\mu_m,K_{m\#}P_\infty\bigr)
&\le
D_{\mathrm{TV}}\!\bigl(\mu_m,K_{m\#}(T_{m\#}\mu_m)\bigr)
+
D_{\mathrm{TV}}\!\bigl(K_{m\#}(T_{m\#}\mu_m),K_{m\#}P_\infty\bigr)\\
&\le
O\!\Bigl(\Bigl(\frac{\varphi}{2}\Bigr)^m\Bigr)
+
D_{\mathrm{TV}}\!\bigl(T_{m\#}\mu_m,P_\infty\bigr)\\
&=
O\!\Bigl(\Bigl(\frac{\varphi}{2}\Bigr)^m\Bigr).
\end{aligned}
\]

最后，对参数集 \(\{0,1\}\) 上的两点实验，前向缺陷由核 \(T_m\) 控制，后向缺陷由核 \(K_m\) 控制，因此 Le Cam 距离满足
\[
\Delta(\mathcal E_m,\mathcal E_\infty)
\le
\max\Bigl\{
D_{\mathrm{TV}}(T_{m\#}u_m,Q_\infty),
D_{\mathrm{TV}}(T_{m\#}\mu_m,P_\infty),
D_{\mathrm{TV}}(u_m,K_{m\#}Q_\infty),
D_{\mathrm{TV}}(\mu_m,K_{m\#}P_\infty)
\Bigr\},
\]
右端趋于 \(0\)，故结论成立。证毕。
\end{proof}

\begin{theorem}[局部化 solenoid 上有理乘法自同态的固定点与 Artin--Mazur \texorpdfstring{$\zeta$}{zeta} 函数]\label{thm:xi-time-part69-localized-rational-endomorphism-fixedpoints-zeta}
\leanverified{paper\_xi\_time\_part69\_localized\_rational\_endomorphism\_fixedpoints\_zeta}
设 \(S\subset\mathbb P\) 为有限素数集，并记
\[
G_S:=\ZZ[S^{-1}],
\qquad
\Sigma_S:=\widehat{G_S}.
\]
取有理数
\[
r=\frac{a}{b}\in G_S,
\qquad
\gcd(a,b)=1,
\]
其中 \(b\) 的全部素因子都属于 \(S\)，并设 \(f_r:\Sigma_S\to\Sigma_S\) 为乘法映射
\[
x\longmapsto rx\qquad (x\in G_S)
\]
的 Pontryagin 对偶连续自同态。若 \(|a|\neq |b|\)，则对每个 \(n\ge 1\)，都有
\[
\#\mathrm{Fix}(f_r^n)=|a^n-b^n|_{S^\perp},
\]
其中对任意非零整数 \(N\)，记
\[
|N|_{S^\perp}:=(|N|)_{S^\perp}
\]
为其 \(S\)-自由部分。并且 Artin--Mazur 动力学 \(\zeta\) 函数满足
\[
\zeta_{f_r}(z)
:=
\exp\!\left(
\sum_{n\ge 1}\frac{\#\mathrm{Fix}(f_r^n)}{n}z^n
\right)
=
\exp\!\left(
\sum_{n\ge 1}\frac{|a^n-b^n|_{S^\perp}}{n}z^n
\right).
\]
换言之，局部化 solenoid 上的周期点几何只感知 \(a^n-b^n\) 的 \(S\)-外部。
\end{theorem}
\begin{proof}
由定理 \ref{thm:xi-localized-solenoid-self-cover-degree-classification}，任意连续自同态都唯一对应于某个 \(r\in G_S\)，其对偶离散同态为
\[
x\longmapsto rx
\qquad (x\in G_S).
\]
因此对每个 \(n\ge 1\)，
\[
f_r^n
\]
在对偶离散侧对应于乘法
\[
x\longmapsto r^n x.
\]
按定理 \ref{thm:xi-localized-solenoid-fixedpoints-artin-mazur-series} 的同样论证，
\[
\mathrm{Fix}(f_r^n)\cong \widehat{G_S/(r^n-1)G_S}.
\]

现把
\[
r^n-1=\frac{a^n-b^n}{b^n}
\]
代入。由于 \(b\) 的全部素因子都属于 \(S\)，故 \(b^n\in G_S^\times\) 是单位，从而
\[
(r^n-1)G_S=(a^n-b^n)G_S.
\]
又因 \(|a|\neq |b|\)，故 \(a^n-b^n\neq 0\)。于是由定理 \ref{thm:xi-localized-quotient-exact-sfree-externalization}，
\[
G_S/(a^n-b^n)G_S
\cong
\ZZ/|a^n-b^n|_{S^\perp}\ZZ.
\]
取 Pontryagin 对偶便得到
\[
\#\mathrm{Fix}(f_r^n)=|a^n-b^n|_{S^\perp}.
\]
再将该固定点计数代入 Artin--Mazur \(\zeta\) 函数定义，即得所示级数公式。证毕。
\end{proof}

\begin{theorem}[局部化有理自覆盖的周期点指数率只由分子分母的最大模控制]\label{thm:xi-time-part69-localized-rational-endomorphism-growth-rate}
\leanverified{paper\_xi\_time\_part69\_localized\_rational\_endomorphism\_growth\_rate}
沿用定理 \ref{thm:xi-time-part69-localized-rational-endomorphism-fixedpoints-zeta} 的记号，并进一步假设
\[
|a|\neq |b|.
\]
则有极限
\[
\lim_{n\to\infty}\frac{1}{n}\log \#\mathrm{Fix}(f_r^n)
=
\log\max\{|a|,|b|\}.
\]
换言之，局部化 solenoid 上有理自覆盖的周期点指数率不受 \(S\)-内素因子的影响，而仅由分子与分母模长中的较大者决定。
\end{theorem}
\begin{proof}
由定理 \ref{thm:xi-time-part69-localized-rational-endomorphism-fixedpoints-zeta} 证明中的同一对偶化论证，对每个 \(n\ge 1\) 都有
\[
\mathrm{Fix}(f_r^n)\cong \widehat{G_S/(r^n-1)G_S}.
\]
又因 \(|a|\neq |b|\)，故 \(a^n-b^n\neq 0\)。并且
\[
r^n-1=\frac{a^n-b^n}{b^n},
\]
其中 \(b^n\in G_S^\times\) 为单位，于是
\[
(r^n-1)G_S=(a^n-b^n)G_S.
\]
再由定理 \ref{thm:xi-localized-quotient-exact-sfree-externalization}，
\[
\#\mathrm{Fix}(f_r^n)=|a^n-b^n|_{S^\perp}.
\]
因此只需比较 \(a^n-b^n\) 的 \(S\)-内与 \(S\)-外部分的增长。

先固定 \(p\in S\)。若 \(p\mid ab\)，则由 \(\gcd(a,b)=1\) 可知 \(p\) 不能同时整除 \(a\) 与 \(b\)，从而对任意 \(n\ge 1\) 都有
\[
p\nmid (a^n-b^n).
\]
若 \(p\nmid ab\)，则 \(a/b\) 在 \((\ZZ/p\ZZ)^\times\) 中有定义。记其模 \(p\) 的乘法阶为 \(t_p\)。只有在 \(t_p\mid n\) 时，\(p\) 才可能整除 \(a^n-b^n\)；并且由标准 lifting-the-exponent 估计，存在仅依赖于 \(p,a,b\) 的常数 \(C_p\)，使得
\[
v_p(a^n-b^n)\le C_p+v_p(n)
\qquad(n\ge 1).
\]
由于 \(S\) 有限，遂得
\[
\sum_{p\in S}v_p(a^n-b^n)\log p=O(\log n),
\]
因而
\[
\frac{1}{n}\sum_{p\in S}v_p(a^n-b^n)\log p\longrightarrow 0.
\]

另一方面，若 \(|a|>|b|\)，则
\[
|a^n-b^n|
=
|a|^n\Bigl|1-(b/a)^n\Bigr|
=
|a|^n(1+o(1)),
\]
故
\[
\frac{1}{n}\log |a^n-b^n|\longrightarrow \log|a|.
\]
若 \(|b|>|a|\)，同理得到极限为 \(\log|b|\)。

最后，由
\[
\log |a^n-b^n|_{S^\perp}
=
\log |a^n-b^n|
-\sum_{p\in S}v_p(a^n-b^n)\log p
\]
即知
\[
\frac{1}{n}\log |a^n-b^n|_{S^\perp}
\longrightarrow
\log\max\{|a|,|b|\}.
\]
再代回
\(\#\mathrm{Fix}(f_r^n)=|a^n-b^n|_{S^\perp}\)，即得所示极限。证毕。
\end{proof}

\begin{theorem}[素数寄存器半群幂等元的精确计数律]\label{thm:xi-time-part69-prime-register-idempotent-count-law}
\leanverified{paper\_xi\_time\_part69\_prime\_register\_idempotent\_count\_law}
沿用定理 \ref{thm:killo-finite-semigroup-prime-valuation-arithmetization} 的记号，设
\[
S_n:=\left\{\prod_{i=0}^{n-1}q_i^{a_i}: a_i\in\{0,\dots,n-1\}\right\},
\]
并记
\[
E_n:=\{N\in S_n:\ N\star N=N\},
\qquad
E_{n,k}:=\{N\in E_n:\ |\operatorname{Im}(f_N)|=k\}
\qquad (1\le k\le n).
\]
则有精确计数公式
\[
|E_{n,k}|=\binom{n}{k}k^{n-k}
\qquad (1\le k\le n),
\]
从而
\[
|E_n|=\sum_{k=1}^{n}\binom{n}{k}k^{n-k}.
\]
换言之，素数寄存器中的幂等投影按像集大小完全分层，而其总数与 \([n]\) 上幂等自映射的经典计数严格一致。
\end{theorem}
\begin{proof}
由定理 \ref{thm:killo-finite-semigroup-prime-valuation-arithmetization} 与推论 \ref{cor:killo-projection-idempotent-prime-register}，\(N\in S_n\) 为幂等元当且仅当对应映射
\[
f_N:[n]\to[n]
\]
为幂等变换。

固定 \(k\in\{1,\dots,n\}\)，并设 \(|\operatorname{Im}(f_N)|=k\)。记
\[
A:=\operatorname{Im}(f_N)\subseteq [n],
\qquad |A|=k.
\]
幂等条件
\[
f_N\circ f_N=f_N
\]
等价于：每个 \(a\in A\) 都必须是定点，即
\[
f_N(a)=a
\qquad (a\in A),
\]
而对每个 \(x\in[n]\setminus A\)，其像只需任意落在 \(A\) 中。因而，一旦像集 \(A\) 固定，幂等变换数恰为
\[
k^{n-k}.
\]
再对全部 \(\binom{n}{k}\) 个 \(k\)-元像集求和，即得
\[
|E_{n,k}|=\binom{n}{k}k^{n-k}.
\]
最后对 \(k=1,\dots,n\) 求和，得到
\[
|E_n|=\sum_{k=1}^{n}\binom{n}{k}k^{n-k}.
\]
证毕。
\end{proof}

\begin{theorem}[fold-gauge 缺陷的形式 Koenigs 线性化]\label{thm:xi-time-part69-fold-gauge-formal-koenigs-linearization}
记
\[
E_m:=\log C_m-\log(2\pi),
\qquad
x:=\frac{\varphi}{2}\in(0,1).
\]
由定理 \ref{thm:xi-time-part60ab-gauge-constant-recursive-bernoulli-stripping}，存在奇形式级数
\[
F(y):=\sum_{r\ge 1}a_r y^{2r-1}\in y\RR[[y^2]],
\qquad
a_1=\frac{1}{3\varphi}\neq 0,
\]
使
\[
E_m=F(x^m).
\]
则存在唯一奇形式幂级数
\[
\mathcal K(z)=3\varphi\,z+\kappa_3 z^3+\kappa_5 z^5+\cdots \in z\RR[[z^2]]
\]
使得对任意 \(N\ge 1\)，若记 \(\mathcal K_N\) 为其截断到 \(z^{2N+1}\) 的多项式，则
\[
\mathcal K_N(E_m)=x^m+O\!\bigl(x^{(2N+1)m}\bigr)
\qquad (m\to\infty).
\]
等价地，若 \(\mathcal R(z)\in z\RR[[z^2]]\) 为定理 \ref{thm:xi-time-part64a-fold-gauge-purely-odd-renormalization} 中的唯一形式重整化映射，则
\[
\mathcal K(\mathcal R(z))=x\,\mathcal K(z)
\]
在形式幂级数意义下成立。其首个非平凡系数为
\[
\kappa_3=\frac{135+63\sqrt5}{40}.
\]
\end{theorem}
\begin{proof}
由于 \(a_1\neq 0\)，级数 \(F\) 在形式幂级数环中存在唯一复合逆
\[
F^{-1}(z)\in z\RR[[z]].
\]
又因为 \(F\) 为奇级数，其逆也必为奇级数，于是
\[
\mathcal K(z):=F^{-1}(z)\in z\RR[[z^2]].
\]
由构造立得
\[
\mathcal K(E_m)=\mathcal K(F(x^m))=x^m
\]
作为形式恒等式成立。若 \(\mathcal K_N\) 为截断到 \(z^{2N+1}\) 的多项式，则
\[
\mathcal K(z)-\mathcal K_N(z)=O(z^{2N+3}),
\]
而 \(E_m\sim a_1x^m\)，故
\[
\mathcal K_N(E_m)=x^m+O(x^{(2N+3)m}),
\]
特别地也有较弱的
\[
\mathcal K_N(E_m)=x^m+O(x^{(2N+1)m}).
\]

另一方面，定理 \ref{thm:xi-time-part64a-fold-gauge-purely-odd-renormalization} 已给出唯一形式重整化映射
\[
\mathcal R(z)=F\!\bigl(xF^{-1}(z)\bigr).
\]
于是
\[
\mathcal K(\mathcal R(z))
=
F^{-1}\!\bigl(F(xF^{-1}(z))\bigr)
=
xF^{-1}(z)
=
x\,\mathcal K(z),
\]
这正是 Koenigs 线性化方程。

最后计算首项系数。由 \(r=1,2\) 的系数公式，
\[
a_1=\frac{1}{3\varphi},
\qquad
a_2=\frac{B_4}{2\cdot 3}\,(\varphi^{-1}+\varphi)
=
-\frac{\sqrt5}{180},
\]
其中用到 \(B_4=-1/30\) 与 \(\varphi+\varphi^{-1}=\sqrt5\)。设
\[
\mathcal K(z)=k_1 z+k_3 z^3+O(z^5).
\]
把它代入
\[
\mathcal K(F(y))=y
\]
并比较 \(y\) 与 \(y^3\) 的系数，得到
\[
k_1a_1=1,
\qquad
k_1a_2+k_3a_1^3=0.
\]
第一式给出
\[
k_1=\frac{1}{a_1}=3\varphi.
\]
第二式给出
\[
k_3=-\frac{a_2}{a_1^4}
=
\frac{135+63\sqrt5}{40}.
\]
故 \(\kappa_3=k_3\)。证毕。
\end{proof}

\noindent
上述七条结果表明，审计偶窗的低预算容量、其首个非线性折点、bin-fold 的完整统计实验、局部化 solenoid 的周期点几何、素数寄存器的有限幂等结构以及 fold-gauge 缺陷的形式重整化之间存在同一类低维闭合机制：在有限分辨率层面，它表现为严格线性段、精确折点与精确的幂等计数；在统计层面，它表现为末位二点实验对全模型的 Le Cam 塌缩；在局部化动力学层面，它既表现为 \(S\)-外部因子对周期点计数的完全控制，也表现为指数率仅由 \(\max\{|a|,|b|\}\) 锁定；而在形式动力学层面，它最终收束为 Koenigs 坐标下的严格乘法线性化。于是，有限预算盲区、容量折点、二点统计极限、局部化周期谱与重整化正规形不再是彼此分离的技术现象，而被统一压缩为同一条可审计的降维主线。

\paragraph{单位自同构、周期指数与 prime-register 盲性}
在一般局部化有理自同态的固定点与指数率已经获得闭式之后，单位参数区间还可进一步收束为一条更强的动力学链：连续自同构群完全由单位格给出，周期点公式退化为 \(S\)-自由差分的精确读数，而遍历、mixing 与同构分类之间则出现了严格分离。

\begin{theorem}[局部化 solenoid 的连续自同构群完全分类]\label{thm:xi-time-part69-localized-unit-automorphism-group-classification}
\leanverified{paper\_xi\_time\_part69\_localized\_unit\_automorphism\_group\_classification}
设 \(S\subset\mathbb P\) 为有限素数集，并记
\[
G_S:=\ZZ[S^{-1}],
\qquad
\Sigma_S:=\widehat{G_S}.
\]
则连续群自同构群满足显式同构
\[
\boxed{
\operatorname{Aut}_{\mathrm{cts}}(\Sigma_S)
\cong
G_S^\times
=
\left\{\pm\prod_{p\in S}p^{n_p}:n_p\in\ZZ\right\}
\cong
C_2\times \ZZ^{|S|}.
}
\]
\end{theorem}
\begin{proof}
定理 \ref{thm:xi-localized-solenoid-self-cover-degree-classification} 已给出
\[
\operatorname{Aut}_{\mathrm{cts}}(\Sigma_S)\cong G_S^\times
=
\left\{\pm\prod_{p\in S}p^{n_p}:n_p\in\ZZ\right\}.
\]
由于 \(S\) 有限，映射
\[
\varepsilon\prod_{p\in S}p^{n_p}
\longmapsto
\bigl(\varepsilon,(n_p)_{p\in S}\bigr),
\qquad
\varepsilon\in\{\pm1\},
\]
给出该单位群与 \(C_2\times \ZZ^{|S|}\) 的群同构，故结论成立。证毕。
\end{proof}

\begin{theorem}[单位比值自同构的周期点公式]\label{thm:xi-time-part69-unit-automorphism-periodic-point-formula}
\leanverified{paper\_xi\_time\_part69\_unit\_automorphism\_periodic\_point\_formula}
沿用定理 \ref{thm:xi-time-part69-localized-unit-automorphism-group-classification} 的记号。取
\[
q\in G_S^\times\setminus\{\pm1\},
\]
并将其写成既约形式
\[
q=\frac{a}{b},
\qquad
a\in\ZZ\setminus\{0\},
\qquad
b\in\NN,
\qquad
\gcd(a,b)=1,
\qquad
\operatorname{Supp}(|ab|)\subseteq S.
\]
令
\[
\alpha_q:\Sigma_S\to\Sigma_S
\]
为与离散端乘法 \(x\mapsto qx\) 对偶的连续群自同构。则对任意 \(m\ge 1\)，有
\[
\boxed{
\operatorname{Fix}(\alpha_q^m)
\cong
\ZZ/|a^m-b^m|_{S^\perp}\ZZ,
\qquad
\#\operatorname{Fix}(\alpha_q^m)=|a^m-b^m|_{S^\perp}.
}
\]
\end{theorem}
\begin{proof}
由定理 \ref{thm:xi-time-part69-localized-unit-automorphism-group-classification}，\(\alpha_q\) 确为 \(\Sigma_S\) 的连续群自同构。再把定理 \ref{thm:xi-time-part69-localized-rational-endomorphism-fixedpoints-zeta} 应用于 \(r=q=a/b\)，便得到
\[
\#\operatorname{Fix}(\alpha_q^m)=|a^m-b^m|_{S^\perp},
\]
以及对应的有限循环群结构。证毕。
\end{proof}

\begin{theorem}[单位自同构的周期指数与 prime-register 的多项式级盲性]\label{thm:xi-time-part69-unit-automorphism-polynomial-primeblindness}
\leanpartial{paper\_xi\_time\_part69\_unit\_automorphism\_polynomial\_primeblindness}{当前 Lean 目标按给定签名为假：取 $a=b=1$ 时假设成立而结论要求正的下界，但左侧恒为 $0$。}
沿用定理 \ref{thm:xi-time-part69-unit-automorphism-periodic-point-formula} 的记号，并设
\[
q\neq \pm1.
\]
则有极限
\[
\boxed{
\lim_{m\to\infty}\frac1m\log \#\operatorname{Fix}(\alpha_q^m)
=
\log\max\{|a|,b\}.
}
\]
更精确地，存在常数 \(C=C(q,S)>0\)，使得对充分大的 \(m\) 都有
\[
\boxed{
C^{-1}\,\frac{\max\{|a|,b\}^{\,m}}{m^{|S|}}
\le
\#\operatorname{Fix}(\alpha_q^m)
\le
C\,\max\{|a|,b\}^{\,m}.
}
\]
因此，prime-register \(S\) 只能改变周期点计数的多项式级修正，而不能改变其指数主率。
\end{theorem}
\begin{proof}
由定理 \ref{thm:xi-time-part69-unit-automorphism-periodic-point-formula}，
\[
\#\operatorname{Fix}(\alpha_q^m)
=
|a^m-b^m|_{S^\perp}
=
\frac{|a^m-b^m|}{\prod_{p\in S}p^{v_p(a^m-b^m)}}.
\]
记
\[
M:=\max\{|a|,b\},
\qquad
\lambda:=\frac{\min\{|a|,b\}}{M}\in(0,1).
\]
由于 \(q\neq \pm1\)，故 \(|a|\neq b\)。若 \(M=|a|\)，则
\[
|a^m-b^m|
=
M^m\left|\operatorname{sgn}(a)^m-\lambda^m\right|;
\]
若 \(M=b\)，则
\[
|a^m-b^m|
=
M^m\left|(\operatorname{sgn}(a)\lambda)^m-1\right|.
\]
两种情形都表明：存在常数 \(c_0,C_0>0\)，使得对充分大的 \(m\)，有
\[
c_0 M^m\le |a^m-b^m|\le C_0 M^m.
\]

现估计 \(S\)-内素因子。对固定 \(p\in S\)：

若 \(p\mid ab\)，则由 \(\gcd(a,b)=1\) 可知 \(p\) 只能整除 \(a,b\) 之一，从而对任意 \(m\ge 1\) 都有
\[
v_p(a^m-b^m)=0.
\]
在这种情形下取 \(C_p:=0\)。
若 \(p\nmid ab\)，则 \(a/b\) 在 \((\ZZ/p\ZZ)^\times\) 中有定义。记其模 \(p\) 的乘法阶为 \(d_p\)。当 \(d_p\nmid m\) 时，显然
\[
v_p(a^m-b^m)=0.
\]
当 \(d_p\mid m\) 时，由标准的 lifting-the-exponent 估计（在 \(p=2\) 时使用其 dyadic 变体）可知，存在仅依赖于 \(p,a,b\) 的常数 \(C_p\)，使得
\[
v_p(a^m-b^m)\le C_p+v_p(m).
\]
因此对所有 \(p\in S\) 与所有 \(m\ge1\)，统一有
\[
v_p(a^m-b^m)\le C_p+v_p(m).
\]
于是
\[
\prod_{p\in S}p^{v_p(a^m-b^m)}
\le
\left(\prod_{p\in S}p^{C_p}\right)\prod_{p\in S}p^{v_p(m)}
\le
C_1\,m^{|S|},
\]
其中 \(C_1:=\prod_{p\in S}p^{C_p}\)。

把上述估计与 \(c_0 M^m\le |a^m-b^m|\le C_0 M^m\) 合并，即得
\[
\#\operatorname{Fix}(\alpha_q^m)
=
\frac{|a^m-b^m|}{\prod_{p\in S}p^{v_p(a^m-b^m)}}
\ge
\frac{c_0}{C_1}\,\frac{M^m}{m^{|S|}}
\]
以及
\[
\#\operatorname{Fix}(\alpha_q^m)\le C_0 M^m.
\]
重新吸收常数便得到所示双边界。两边取对数、除以 \(m\) 并令 \(m\to\infty\)，立得指数极限。证毕。
\end{proof}

\begin{theorem}[rank-\texorpdfstring{$1$}{1} 局部化世界中遍历与 mixing 完全塌缩为 \texorpdfstring{$q\neq\pm1$}{q!=±1}]\label{thm:xi-time-part69-unit-automorphism-ergodic-mixing-collapse}
\leanverified{paper\_xi\_time\_part69\_unit\_automorphism\_ergodic\_mixing\_collapse}
沿用定理 \ref{thm:xi-time-part69-localized-unit-automorphism-group-classification} 的记号。对任意 \(q\in G_S^\times\)，记
\[
\alpha_q:\Sigma_S\to\Sigma_S
\]
为对偶于乘法 \(x\mapsto qx\) 的连续群自同构。则在 Haar 概率测度意义下，
\[
\boxed{
\alpha_q\ \text{遍历}
\iff
\alpha_q\ \text{mixing}
\iff
q\neq \pm1.
}
\]
\end{theorem}
\begin{proof}
记 \(\mu_S\) 为 \(\Sigma_S\) 上的 Haar 概率测度。对每个 \(x\in G_S\)，记对应字符为
\[
\chi_x(\omega):=\omega(x)
\qquad(\omega\in \Sigma_S).
\]
则 Pontryagin 对偶关系给出
\[
\chi_x\circ \alpha_q^n=\chi_{q^n x}
\qquad(n\ge 0).
\]
因而对任意 \(x,y\in G_S\)，利用字符正交性有
\[
\int_{\Sigma_S}(\chi_x\circ \alpha_q^n)\,\overline{\chi_y}\,d\mu_S
=
\int_{\Sigma_S}\chi_{q^n x-y}\,d\mu_S
=
\begin{cases}
1,& q^n x=y,\\
0,& q^n x\neq y.
\end{cases}
\]

现设 \(q\neq \pm1\)。若存在 \(x,y\neq 0\) 与两个不同整数 \(m>n\) 使
\[
q^m x=y=q^n x,
\]
则
\[
(q^{m-n}-1)x=0.
\]
由于 \(G_S=\ZZ[S^{-1}]\subset\QQ\) 无挠，必有 \(q^{m-n}=1\)，与 \(q\neq \pm1\) 矛盾。故对任意非零 \(x,y\in G_S\)，方程 \(q^n x=y\) 至多成立一次；于是上式对所有充分大的 \(n\) 恒为 \(0\)。

由三角多项式在 \(L^2(\Sigma_S,\mu_S)\) 中稠密，便得到：对任意零均值函数 \(f,g\in L^2(\Sigma_S,\mu_S)\)，有
\[
\langle f\circ \alpha_q^n,g\rangle\longrightarrow 0.
\]
故 \(\alpha_q\) mixing，从而也遍历。

反之，若 \(q=1\)，则 \(\alpha_q=\mathrm{id}\)，显然不遍历。若 \(q=-1\)，取任意非零 \(x\in G_S\)。则
\[
\chi_x\circ \alpha_{-1}^{2n}=\chi_x
\qquad(n\ge 0),
\]
故 \(\alpha_{-1}\) 不是 mixing。并且
\[
\Re\chi_x\circ \alpha_{-1}
=
\Re\chi_{-x}
=
\Re\chi_x,
\]
而 \(\Re\chi_x\) 非常值，故 \(\alpha_{-1}\) 亦不遍历。综上得到所示等价关系。证毕。
\end{proof}

\begin{corollary}[同一指数主率下存在无穷多 pairwise non-isomorphic 的 mixing 局部化系统]\label{cor:xi-time-part69-same-exponential-rate-infinite-mixing-family}
固定任意既约有理数
\[
q=\frac{a}{b}>1,
\qquad
a,b\in\NN,
\qquad
\gcd(a,b)=1.
\]
对任意有限素数集
\[
S\supseteq \operatorname{Supp}(ab),
\]
记
\[
\Sigma_S:=\widehat{\ZZ[S^{-1}]},
\qquad
\alpha_q:\Sigma_S\to\Sigma_S
\]
为与离散端乘法 \(x\mapsto qx\) 对偶的连续群自同构。则：
\begin{enumerate}
  \item 对所有这类 \(S\)，都有
  \[
  \lim_{m\to\infty}\frac1m\log \#\operatorname{Fix}(\alpha_q^m)=\log a;
  \]
  \item 若 \(S\neq T\)，则
  \[
  \Sigma_S\not\cong \Sigma_T;
  \]
  \item 因而对每个固定的 \(q>1\)，存在无穷多个 pairwise non-isomorphic 的 compact connected mixing 系统
  \[
  (\Sigma_S,\alpha_q),
  \]
  它们具有完全相同的周期指数主率 \(\log a\)。
\end{enumerate}
\end{corollary}
\begin{proof}
第 \(1\) 条是定理 \ref{thm:xi-time-part69-unit-automorphism-polynomial-primeblindness} 的直接特例，因为此时
\[
\max\{|a|,b\}=a.
\]
第 \(2\) 条由定理 \ref{thm:xi-time-part56-localized-pontryagin-isomorphism-rigidity} 立即得到。

对第 \(3\) 条，只需注意：包含 \(\operatorname{Supp}(ab)\) 的有限素数集 \(S\) 有无穷多个；由定理 \ref{thm:xi-time-part69-unit-automorphism-ergodic-mixing-collapse}，对每个这样的 \(S\)，因 \(q>1\) 故 \(\alpha_q\) mixing；再结合前两条，便得到一列两两不同构而指数主率完全相同的局部化动力系统。证毕。
\end{proof}

\paragraph{有限角色审计的单圆塌缩与 prime-register 尾部失明}
在局部化 solenoid 的连续自同构、周期点指数与 prime-register 的多项式级盲性已经闭合之后，有限角色族的审计能力还可进一步压缩为一条双重塌缩链：在连通方向上，任意有限角色族都退化为单一 primitive 圆相位；在全不连通方向上，它们对 prime-register 核只能读取有限深度截断，因此有限 trigonometric 审计代数同时失去高维相位自由度与无限 \(p\)-进尾部。

\begin{theorem}[局部化 solenoid 上任意有限角色族都塌缩为单一 primitive 相位]\label{thm:xi-time-part69-finite-character-family-single-primitive-phase}
\leanverified{paper\_xi\_time\_part69\_finite\_character\_family\_single\_primitive\_phase}
设 \(S\subset\mathbb P\) 为有限素数集，并记
\[
G_S:=\ZZ[S^{-1}],
\qquad
\Sigma_S:=\widehat{G_S}.
\]
任取有限族连续角色
\[
\chi_1,\dots,\chi_r\in \widehat{\Sigma_S}\cong G_S.
\]
则存在一个角色 \(\chi_\ast\in\widehat{\Sigma_S}\) 与整数
\[
n_1,\dots,n_r\in\ZZ
\]
使得
\[
\chi_j=\chi_\ast^{\,n_j}
\qquad(1\le j\le r).
\]
若 \(\chi_1,\dots,\chi_r\) 不全为平凡角色，则 \(\chi_\ast\) 可取为 primitive 角色，并且在此规范下唯一到反演；相应地，整数族 \((n_1,\dots,n_r)\) 自动满足
\[
\gcd(n_1,\dots,n_r)=1.
\]
\end{theorem}
\begin{proof}
在 Pontryagin 对偶同构
\(
\widehat{\Sigma_S}\cong G_S
\)
下，令 \(\chi_j\) 对应于 \(q_j\in G_S\)。记
\[
H:=\langle q_1,\dots,q_r\rangle\subseteq G_S\subseteq \QQ.
\]
取正整数 \(D\) 使得
\[
Dq_j\in\ZZ
\qquad(1\le j\le r).
\]
则 \(DH\) 是 \(\ZZ\) 的有限生成子群，故必为循环群。于是存在唯一整数 \(d\ge 0\) 使
\[
DH=d\ZZ.
\]

若 \(d=0\)，则 \(H=\{0\}\)，亦即全部 \(q_j=0\)，从而所有 \(\chi_j\) 都平凡。此时取 \(\chi_\ast=1\) 与 \(n_j=0\) 即可。

若 \(d>0\)，则
\[
H=\frac{d}{D}\ZZ.
\]
记
\[
q_\ast:=\frac{d}{D}\in G_S.
\]
则 \(H=\ZZ q_\ast\)，故存在整数 \(n_j\) 使
\[
q_j=n_j q_\ast
\qquad(1\le j\le r).
\]
令 \(\chi_\ast\) 为与 \(q_\ast\) 对应的角色，则
\[
\chi_j=\chi_{q_j}=\chi_{q_\ast}^{\,n_j}=\chi_\ast^{\,n_j}.
\]
若 \(q_\ast'\) 也是 \(H\) 的生成元，则 \(q_\ast'=\pm q_\ast\)，故对应角色只差反演 \(\chi_\ast^{-1}\)。最后，由
\[
H=\langle q_1,\dots,q_r\rangle=\langle n_1 q_\ast,\dots,n_r q_\ast\rangle
\]
可知 \(\langle n_1,\dots,n_r\rangle=\ZZ\)，从而 \(\gcd(n_1,\dots,n_r)=1\)。证毕。
\end{proof}

\begin{corollary}[有限相位审计的全部 Borel 结构都等价于单圆相位审计]\label{cor:xi-time-part69-finite-phase-audit-borel-collapse-single-circle}
沿用定理 \ref{thm:xi-time-part69-finite-character-family-single-primitive-phase} 的记号，并定义联合相位映射
\[
\boldsymbol{\chi}:\Sigma_S\to \TT^r,
\qquad
x\longmapsto \bigl(\chi_1(x),\dots,\chi_r(x)\bigr).
\]
则存在一个角色 \(\chi_\ast:\Sigma_S\to\TT\) 与整数指数单项式映射
\[
M:\TT\to \TT^r,
\qquad
M(z):=(z^{n_1},\dots,z^{n_r}),
\]
使得
\[
\boldsymbol{\chi}=M\circ \chi_\ast.
\]
从而对任意 Borel 集 \(B\subseteq \TT^r\)，都存在 Borel 集 \(B_\ast\subseteq \TT\) 使
\[
\boldsymbol{\chi}^{-1}(B)=\chi_\ast^{-1}(B_\ast).
\]
并且有限角色族所生成的 \(\sigma\)-代数恰满足
\[
\sigma(\chi_1,\dots,\chi_r)=\sigma(\chi_\ast).
\]
\end{corollary}
\begin{proof}
由定理 \ref{thm:xi-time-part69-finite-character-family-single-primitive-phase}，存在 \(\chi_\ast\) 与整数 \(n_1,\dots,n_r\) 使
\[
\chi_j=\chi_\ast^{\,n_j}.
\]
故对任意 \(x\in \Sigma_S\)，有
\[
\boldsymbol{\chi}(x)
=
\bigl(\chi_\ast(x)^{n_1},\dots,\chi_\ast(x)^{n_r}\bigr)
=
M(\chi_\ast(x)),
\]
即 \(\boldsymbol{\chi}=M\circ \chi_\ast\)。

任取 Borel 集 \(B\subseteq \TT^r\)，令
\[
B_\ast:=M^{-1}(B)\subseteq \TT.
\]
由于 \(M\) 连续，\(B_\ast\) 仍为 Borel 集，并且
\[
\boldsymbol{\chi}^{-1}(B)
=
\chi_\ast^{-1}(M^{-1}(B))
=
\chi_\ast^{-1}(B_\ast).
\]
这给出
\[
\sigma(\chi_1,\dots,\chi_r)\subseteq \sigma(\chi_\ast).
\]

若全部 \(\chi_j\) 平凡，则两侧 \(\sigma\)-代数都只含空集与全集，结论已得。下设并非全部平凡。由定理 \ref{thm:xi-time-part69-finite-character-family-single-primitive-phase}，
\[
\gcd(n_1,\dots,n_r)=1.
\]
故可取整数 \(a_1,\dots,a_r\) 使
\[
\sum_{j=1}^{r}a_j n_j=1.
\]
定义连续映射
\[
N:\TT^r\to \TT,
\qquad
N(z_1,\dots,z_r):=\prod_{j=1}^{r}z_j^{a_j}.
\]
则对任意 \(x\in \Sigma_S\)，有
\[
N(\boldsymbol{\chi}(x))
=
\prod_{j=1}^{r}\chi_j(x)^{a_j}
=
\chi_\ast(x)^{\sum_j a_j n_j}
=
\chi_\ast(x).
\]
于是
\[
\chi_\ast=N\circ \boldsymbol{\chi}.
\]
对任意 Borel 集 \(C\subseteq \TT\)，有
\[
\chi_\ast^{-1}(C)=\boldsymbol{\chi}^{-1}(N^{-1}(C)),
\]
从而
\[
\sigma(\chi_\ast)\subseteq \sigma(\chi_1,\dots,\chi_r).
\]
综上得到两侧 \(\sigma\)-代数相等。证毕。
\end{proof}

\begin{theorem}[有限角色族对 prime-register 核只能看到有限层截断]\label{thm:xi-time-part69-finite-character-prime-ledger-finite-depth}
沿用定理 \ref{thm:cdim-star-z1s-dual-extension} 的短正合列
\[
0\longrightarrow K_S:=\prod_{p\in S}\ZZ_p
\longrightarrow \Sigma_S
\longrightarrow \TT
\longrightarrow 0.
\]
给定任意有限角色族 \(\chi_1,\dots,\chi_r\in \widehat{\Sigma_S}\cong G_S\)，令 \(q_j\in G_S\) 为其对应元素，并对每个 \(p\in S\) 定义
\[
m_{j,p}:=\max\{0,-v_p(q_j)\}.
\]
当 \(q_j\neq 0\) 时，等价地，\(q_j\) 可唯一写成
\[
q_j=a_j\prod_{p\in S}p^{-m_{j,p}},
\qquad
a_j\in\ZZ,
\qquad
v_p(a_j)=0\ \ (p\in S).
\]
若 \(q_j=0\)，则取全部 \(m_{j,p}=0\)。
再令
\[
M_p:=\max_{1\le j\le r}m_{j,p},
\qquad
U(\chi_1,\dots,\chi_r):=\prod_{p\in S}p^{M_p}\ZZ_p\subseteq K_S.
\]
则：
\begin{enumerate}
  \item 每个限制角色 \(\chi_j|_{K_S}\) 都在 \(U(\chi_1,\dots,\chi_r)\) 的陪集上恒定；
  \item 若记
  \[
  T_{\boldsymbol{\chi}}:=\{p\in S:\ M_p\ge 1\},
  \]
  则联合映射
  \[
  K_S\longrightarrow \TT^r,
  \qquad
  x\longmapsto \bigl(\chi_1(x),\dots,\chi_r(x)\bigr)
  \]
  必经由有限商
  \[
  K_S/U(\chi_1,\dots,\chi_r)
  \cong
  \prod_{p\in T_{\boldsymbol{\chi}}}\ZZ/p^{M_p}\ZZ
  \]
  因子化；
  \item 因而任何有限角色审计都严格失明于 \(K_S\) 的无限 \(p\)-进尾部。
\end{enumerate}
\end{theorem}
\begin{proof}
由定理 \ref{thm:cdim-star-z1s-dual-extension}，有互为 Pontryagin 对偶的短正合列
\[
0\longrightarrow \ZZ
\longrightarrow G_S
\longrightarrow \bigoplus_{p\in S}C_{p^\infty}
\longrightarrow 0
\]
与
\[
0\longrightarrow K_S
\longrightarrow \Sigma_S
\longrightarrow \TT
\longrightarrow 0.
\]
因此，\(\chi_j\) 在 \(K_S\) 上的限制正对应于 \(q_j\) 在商群
\[
G_S/\ZZ\cong \bigoplus_{p\in S}C_{p^\infty}
\]
中的像。

固定 \(j\) 与 \(p\in S\)。由 \(m_{j,p}=\max\{0,-v_p(q_j)\}\) 的定义，\(q_j\) 在 \(G_S/\ZZ\) 中的 \(p\)-主分量落在 \(C_{p^\infty}\) 的唯一 \(p^{m_{j,p}}\)-阶循环子群内。对偶化以后，这意味着 \(\chi_j|_{K_S}\) 的第 \(p\)-坐标分量只能通过有限商
\[
\ZZ_p\twoheadrightarrow \ZZ_p/p^{m_{j,p}}\ZZ_p\cong \ZZ/p^{m_{j,p}}\ZZ
\]
被读取；等价地，
\[
p^{m_{j,p}}\ZZ_p
\subseteq
\ker\bigl(\chi_j|_{\ZZ_p}\bigr).
\]
于是定义
\[
U_j:=\prod_{p\in S}p^{m_{j,p}}\ZZ_p\subseteq K_S,
\]
便有
\[
U_j\subseteq \ker(\chi_j|_{K_S}).
\]

对全部 \(j\) 取交，得到
\[
\bigcap_{j=1}^{r}U_j
=
\prod_{p\in S}p^{M_p}\ZZ_p
=
U(\chi_1,\dots,\chi_r).
\]
故每个 \(\chi_j|_{K_S}\) 都在 \(U(\chi_1,\dots,\chi_r)\) 的陪集上恒定，第 \(1\) 条成立。

由于 \(U(\chi_1,\dots,\chi_r)\) 是 \(K_S\) 的开子群，联合映射必经由商群 \(K_S/U(\chi_1,\dots,\chi_r)\) 因子化。又由各坐标逐素数分解，
\[
K_S/U(\chi_1,\dots,\chi_r)
\cong
\prod_{p\in T_{\boldsymbol{\chi}}}\ZZ/p^{M_p}\ZZ.
\]
这给出第 \(2\) 条。最后，所有更深层的 \(p\)-进尾部都落在 \(U(\chi_1,\dots,\chi_r)\) 内，因而不会改变任何有限角色的取值，第 \(3\) 条随即成立。证毕。
\end{proof}

\begin{theorem}[有限角色生成的 trigonometric 审计代数在一致闭包后恰为单圆代数]\label{thm:xi-time-part69-finite-trigonometric-audit-algebra-single-circle-finite-tail}
设
\[
\mathcal A(\chi_1,\dots,\chi_r)\subset C(\Sigma_S)
\]
为由有限角色族 \(\chi_1,\dots,\chi_r\) 及其复共轭生成的含幺 \(^\ast\)-代数。令 \(\chi_\ast\) 为定理 \ref{thm:xi-time-part69-finite-character-family-single-primitive-phase} 给出的生成角色，并令
\[
U:=U(\chi_1,\dots,\chi_r)\subseteq K_S
\]
为定理 \ref{thm:xi-time-part69-finite-character-prime-ledger-finite-depth} 的开子群。则成立：
\begin{enumerate}
  \item 代数恒等式
  \[
  \mathcal A(\chi_1,\dots,\chi_r)
  =
  \chi_\ast^{*}\bigl(\CC[z,z^{-1}]\bigr);
  \]
  \item 若 \(\chi_1,\dots,\chi_r\) 不全平凡，则
  \[
  \overline{\mathcal A(\chi_1,\dots,\chi_r)}^{\|\cdot\|_\infty}
  =
  \chi_\ast^{*}C(\TT)
  \cong
  C(\TT);
  \]
  若全部平凡，则
  \[
  \overline{\mathcal A(\chi_1,\dots,\chi_r)}^{\|\cdot\|_\infty}
  =
  \CC\mathbf 1;
  \]
  \item 由 \(\mathcal A(\chi_1,\dots,\chi_r)\) 生成的任意有限 Borel 分割的每个非空原子都包含某个 \(U\)-陪集；特别地，它不可能把 \(K_S\) 的点分离到有限精度以下。
\end{enumerate}
\end{theorem}
\begin{proof}
由定理 \ref{thm:xi-time-part69-finite-character-family-single-primitive-phase}，存在整数 \(n_1,\dots,n_r\) 使
\[
\chi_j=\chi_\ast^{\,n_j}.
\]
因此 \(\mathcal A(\chi_1,\dots,\chi_r)\) 的每个元素都是 \(\chi_\ast\) 的 Laurent 多项式，故有包含
\[
\mathcal A(\chi_1,\dots,\chi_r)
\subseteq
\chi_\ast^{*}\bigl(\CC[z,z^{-1}]\bigr).
\]

若全部 \(\chi_j\) 平凡，则两边都等于常数代数，结论立即成立。下设并非全部平凡。由定理 \ref{thm:xi-time-part69-finite-character-family-single-primitive-phase}，可取 \(\chi_\ast\) primitive，且
\[
\gcd(n_1,\dots,n_r)=1.
\]
故可取整数 \(a_1,\dots,a_r\) 使
\[
\sum_{j=1}^{r}a_j n_j=1.
\]
于是
\[
\chi_\ast=\prod_{j=1}^{r}\chi_j^{a_j}.
\]
由于负指数可由复共轭实现，故 \(\chi_\ast\in \mathcal A(\chi_1,\dots,\chi_r)\)。从而
\[
\chi_\ast^{*}\bigl(\CC[z,z^{-1}]\bigr)
\subseteq
\mathcal A(\chi_1,\dots,\chi_r),
\]
第 \(1\) 条得证。

对第 \(2\) 条，只需证明非平凡情形下 \(\chi_\ast\) 满射到 \(\TT\)。由于 \(G_S\) 无挠，\(\Sigma_S=\widehat{G_S}\) 是紧连通交换群；故 \(\chi_\ast(\Sigma_S)\) 是 \(\TT\) 的闭连通子群。又 \(\chi_\ast\) 非平凡，故其像不能为 \(\{1\}\)，从而
\[
\chi_\ast(\Sigma_S)=\TT.
\]
于是 \(\chi_\ast^{*}:C(\TT)\to C(\Sigma_S)\) 是等距单射。由 Stone--Weierstrass 定理，
\[
\overline{\CC[z,z^{-1}]}^{\|\cdot\|_\infty}=C(\TT),
\]
拉回后即得
\[
\overline{\mathcal A(\chi_1,\dots,\chi_r)}^{\|\cdot\|_\infty}
=
\chi_\ast^{*}C(\TT)
\cong
C(\TT).
\]

最后，由定理 \ref{thm:xi-time-part69-finite-character-prime-ledger-finite-depth}，每个生成元 \(\chi_j\) 都在 \(U\)-陪集上恒定。故 \(\mathcal A(\chi_1,\dots,\chi_r)\) 的每个元素，以及其一致闭包中的每个元素，也都在 \(U\)-陪集上恒定。于是由这些函数所生成的全部 Borel 集都是 \(U\)-饱和集，即都是 \(U\)-陪集的并。因而任意由该代数决定的有限 Borel 分割，其每个非空原子也必为 \(U\)-陪集的并；特别地，取其中任一点 \(x\)，整个陪集 \(xU\) 都包含在该原子中。故该分割不可能把 \(K_S\) 的点分离到比 \(U\) 更细的精度。证毕。
\end{proof}

\paragraph{有限连续可见性、有限指数/有限商谱与连续自同态环的闭合压缩}
在局部化数系的对偶短正合列、有限商分类、连续自同态与有限角色审计已经分别获得闭式之后，还可把连续有限可见性、有限指数格、有限商轮廓与连续动力学几条线路进一步压缩为下列若干条直接后果。

\begin{theorem}[局部化数系对偶在有限连续读出下的完全塌缩]\label{thm:xi-time-part69c-localized-dual-finite-visible-collapse}
\leanverified{paper\_xi\_time\_part69c\_localized\_dual\_finite\_visible\_collapse}
设 \(S\subset\mathbb P\) 为有限素数集，并记
\[
G_S:=\ZZ[S^{-1}],
\qquad
\Sigma_S:=\widehat{G_S}.
\]
则 \(\Sigma_S\) 是连通紧交换群。因而对任意有限离散集 \(A\) 与任意连续映射
\[
F:\Sigma_S\longrightarrow A,
\]
映射 \(F\) 必为常值。
\end{theorem}
\begin{proof}
由定理 \ref{thm:cdim-star-z1s-dual-extension}，\(\Sigma_S\) 为紧交换群，并落在短正合列
\[
0\longrightarrow \prod_{p\in S}\ZZ_p
\longrightarrow \Sigma_S
\longrightarrow \TT
\longrightarrow 0.
\]
另一方面，
\[
G_S=\ZZ[S^{-1}]\subset \QQ
\]
作为离散加法群无挠。按 Pontryagin 对偶的标准连通性判据，\(\Sigma_S\) 必连通。

现令 \(A\) 为有限离散集，\(F:\Sigma_S\to A\) 连续。则 \(F(\Sigma_S)\) 是 \(A\) 的连通子空间；而有限离散空间的连通子空间只能是单点，故 \(F(\Sigma_S)\) 仅含一个元素，从而 \(F\) 为常值。证毕。
\end{proof}

\begin{theorem}[局部化数系有限商阶谱的完全闭式]\label{thm:xi-time-part69c-localized-finite-quotient-order-spectrum}
\leanverified{paper\_xi\_time\_part69c\_localized\_finite\_quotient\_order\_spectrum}
沿用定理 \ref{thm:xi-time-part69c-localized-dual-finite-visible-collapse} 的记号，并定义
\[
\mathcal Q(G_S)
:=
\left\{
n\ge 1:\ \exists\ \text{满同态 }G_S\twoheadrightarrow \ZZ/n\ZZ
\right\}.
\]
则有严格等式
\[
\mathcal Q(G_S)
=
\left\{
n\ge 1:\ \gcd\!\left(n,\prod_{p\in S}p\right)=1
\right\}.
\]
等价地，\(G_S\) 的全部有限商恰为循环群
\[
\ZZ/n\ZZ
\qquad
\bigl(\gcd(n,\prod_{p\in S}p)=1\bigr).
\]
\end{theorem}
\begin{proof}
由定理 \ref{thm:xi-localized-finite-quotient-phase-layer-classification} 第 \(1\) 条，对任意 \(n\ge 1\) 都有
\[
G_S/nG_S\cong \ZZ/n_{S^\perp}\ZZ.
\]
若
\[
\gcd\!\left(n,\prod_{p\in S}p\right)=1,
\]
则 \(n_{S^\perp}=n\)，故自然商映射给出满同态
\[
G_S\twoheadrightarrow G_S/nG_S\cong \ZZ/n\ZZ,
\]
于是 \(n\in \mathcal Q(G_S)\)。

反之，若 \(n\in \mathcal Q(G_S)\)，则存在满同态
\[
G_S\twoheadrightarrow \ZZ/n\ZZ.
\]
这说明 \(\ZZ/n\ZZ\) 是 \(G_S\) 的有限商。再由定理
\ref{thm:xi-localized-finite-quotient-phase-layer-classification}
第 \(2\) 条，\(G_S\) 的每个有限商都必为阶数与 \(S\) 中所有素数互素的循环群。故 \(n\) 的全部素因子都不属于 \(S\)，亦即
\[
\gcd\!\left(n,\prod_{p\in S}p\right)=1.
\]
于是所示集合恒等式成立；最后一句只是同一定理的重述。证毕。
\end{proof}

\begin{theorem}[局部化数系的有限指数子群完全分类]\label{thm:xi-time-part69c-localized-finite-index-subgroup-classification}
沿用定理 \ref{thm:xi-time-part69c-localized-finite-quotient-order-spectrum} 的记号。则 \(G_S\) 的任意有限指数子群 \(H\le G_S\) 都唯一形如
\[
H=nG_S
\]
其中
\[
n\ge 1,
\qquad
\gcd\!\left(n,\prod_{p\in S}p\right)=1.
\]
并且其指数满足精确公式
\[
[G_S:nG_S]=n.
\]
\leanverified{paper\_xi\_time\_part69c\_localized\_finite\_index\_subgroup\_classification}
\end{theorem}
\begin{proof}
先设
\[
\gcd\!\left(n,\prod_{p\in S}p\right)=1.
\]
定义模 \(n\) 约化映射
\[
\rho_n:G_S\to \ZZ/n\ZZ,
\qquad
\rho_n\!\left(\frac{a}{\prod_{p\in S}p^{k_p}}\right)
:=
a\cdot\left(\prod_{p\in S}p^{k_p}\right)^{-1}\pmod n.
\]
由于每个 \(p\in S\) 在 \(\ZZ/n\ZZ\) 中可逆，\(\rho_n\) 良定义。它显然满射，因为整数子群 \(\ZZ\subset G_S\) 已满射到 \(\ZZ/n\ZZ\)。

现证
\[
\ker(\rho_n)=nG_S.
\]
包含 \(nG_S\subseteq \ker(\rho_n)\) 直接成立。反向地，若
\[
x=\frac{a}{b}\in G_S,
\qquad
b=\prod_{p\in S}p^{k_p},
\]
且 \(\rho_n(x)=0\)，则
\[
a\,b^{-1}\equiv 0\pmod n.
\]
由于 \(\gcd(b,n)=1\)，\(b^{-1}\) 在 \(\ZZ/n\ZZ\) 中可逆，故 \(a\equiv 0\pmod n\)。写 \(a=nc\)，便有
\[
x=\frac{nc}{b}=n\cdot \frac{c}{b}\in nG_S.
\]
因此
\[
G_S/nG_S\cong \ZZ/n\ZZ,
\]
从而
\[
[G_S:nG_S]=n.
\]

反之，设 \(H\le G_S\) 为有限指数子群，并记
\[
A:=G_S/H.
\]
则 \(A\) 为有限商。由定理
\ref{thm:xi-time-part69c-localized-finite-quotient-order-spectrum}，存在唯一满足
\[
\gcd\!\left(n,\prod_{p\in S}p\right)=1
\]
的正整数 \(n\)，使得
\[
A\cong \ZZ/n\ZZ.
\]
因而 \([G_S:H]=n\)，并且 \(nA=0\)。这说明对每个 \(x\in G_S\) 都有
\[
nx\in H,
\]
故
\[
nG_S\subseteq H.
\]
而第一部分已证明
\[
[G_S:nG_S]=n=[G_S:H].
\]
既然 \(nG_S\subseteq H\) 且两者指数相同，遂得
\[
H=nG_S.
\]
唯一性则由指数唯一确定。证毕。
\end{proof}

\begin{corollary}[有限指数子群格与 \(S\)-自由整除格的反同构及其子群 \texorpdfstring{$\zeta$}{zeta} 函数]\label{cor:xi-time-part69c-localized-finite-index-lattice-zeta}
沿用定理 \ref{thm:xi-time-part69c-localized-finite-index-subgroup-classification} 的记号。则映射
\[
n\longmapsto nG_S
\qquad
\left(
\gcd\!\left(n,\prod_{p\in S}p\right)=1
\right)
\]
把 \(S\)-自由正整数的整除偏序反同构到 \(G_S\) 的有限指数子群格。更具体地，对任意与 \(S\) 互素的正整数 \(m,n\)，有
\[
nG_S\subseteq mG_S
\quad\Longleftrightarrow\quad
m\mid n,
\]
并且
\[
nG_S\cap mG_S=\operatorname{lcm}(n,m)\,G_S,
\]
\[
nG_S+mG_S=\gcd(n,m)\,G_S.
\]
因此 \(G_S\) 的子群 \(\zeta\) 函数满足完全闭式
\[
\zeta_{G_S}^{\le}(s)
:=
\sum_{H\le_f G_S}[G_S:H]^{-s}
=
\sum_{\gcd(n,\prod_{p\in S}p)=1}n^{-s}
=
\zeta(s)\prod_{p\in S}(1-p^{-s}).
\]
\end{corollary}
\begin{proof}
先证包含判据。若 \(m\mid n\)，写 \(n=mk\)，则
\[
nG_S=m(kG_S)\subseteq mG_S,
\]
故 \(nG_S\subseteq mG_S\)。

反之，若 \(nG_S\subseteq mG_S\)，则特别有 \(n\in mG_S\)，故存在 \(q\in G_S\) 使
\[
n=mq.
\]
于是
\[
q=\frac{n}{m}\in G_S.
\]
但 \(n,m\) 都与 \(S\) 互素，故既约分数 \(n/m\) 的分母不含任何 \(S\)-中素数。要使 \(n/m\in G_S=\ZZ[S^{-1}]\)，其分母只能为 \(1\)。因此 \(n/m\in \ZZ\)，即 \(m\mid n\)。

现设
\[
K:=nG_S\cap mG_S.
\]
由定理 \ref{thm:xi-time-part69c-localized-finite-index-subgroup-classification}，存在唯一 \(S\)-自由正整数 \(k\) 使
\[
K=kG_S.
\]
由 \(K\subseteq nG_S\) 与 \(K\subseteq mG_S\)，按上面的包含判据得
\[
n\mid k,
\qquad
m\mid k,
\]
故 \(\operatorname{lcm}(n,m)\mid k\)。另一方面，
\[
\operatorname{lcm}(n,m)\,G_S\subseteq nG_S\cap mG_S=K,
\]
再次应用包含判据得
\[
k\mid \operatorname{lcm}(n,m).
\]
故 \(k=\operatorname{lcm}(n,m)\)，从而
\[
nG_S\cap mG_S=\operatorname{lcm}(n,m)\,G_S.
\]

再设
\[
L:=nG_S+mG_S.
\]
同理，\(L=\ell G_S\) 对某个唯一 \(S\)-自由正整数 \(\ell\) 成立。由
\[
nG_S\subseteq L,
\qquad
mG_S\subseteq L,
\]
可知
\[
\ell\mid n,
\qquad
\ell\mid m,
\]
故 \(\ell\mid \gcd(n,m)\)。另一方面，取整数 \(u,v\) 满足
\[
un+vm=\gcd(n,m).
\]
则对任意 \(x\in G_S\)，有
\[
\gcd(n,m)\,x=u(nx)+v(mx)\in nG_S+mG_S=L.
\]
因此
\[
\gcd(n,m)\,G_S\subseteq L=\ell G_S,
\]
故由包含判据得
\[
\ell\mid \gcd(n,m).
\]
另一方面，由 \(n,m\) 都是 \(\gcd(n,m)\) 的整数倍，显然
\[
nG_S\subseteq \gcd(n,m)\,G_S,
\qquad
mG_S\subseteq \gcd(n,m)\,G_S.
\]
从而
\[
L=nG_S+mG_S\subseteq \gcd(n,m)\,G_S,
\]
再次应用包含判据得到
\[
\gcd(n,m)\mid \ell.
\]
于是 \(\ell=\gcd(n,m)\)，从而
\[
nG_S+mG_S=\gcd(n,m)\,G_S.
\]

最后，由定理 \ref{thm:xi-time-part69c-localized-finite-index-subgroup-classification}，全部有限指数子群恰由
\[
H=nG_S
\qquad
\left(
\gcd\!\left(n,\prod_{p\in S}p\right)=1
\right)
\]
枚举，且其指数正为 \(n\)。故
\[
\zeta_{G_S}^{\le}(s)
=
\sum_{\gcd(n,\prod_{p\in S}p)=1}n^{-s}.
\]
按 Euler 乘积展开即得
\[
\sum_{\gcd(n,\prod_{p\in S}p)=1}n^{-s}
=
\prod_{p\notin S}(1-p^{-s})^{-1}
=
\zeta(s)\prod_{p\in S}(1-p^{-s}).
\]
证毕。
\end{proof}

\begin{corollary}[有限商轮廓对素数账本的完全恢复]\label{cor:xi-time-part69c-localized-finite-quotient-profile-recovers-primes}
\leanverified{paper\_xi\_time\_part69c\_localized\_finite\_quotient\_profile\_recovers\_primes}
对任意有限素数集 \(S,T\subset\mathbb P\)，有
\[
\mathcal Q(G_S)=\mathcal Q(G_T)
\quad\Longleftrightarrow\quad
S=T.
\]
等价地，对任意素数 \(p\) 都有
\[
p\in S
\quad\Longleftrightarrow\quad
\forall n\in \mathcal Q(G_S),\ p\nmid n.
\]
\end{corollary}
\begin{proof}
由定理 \ref{thm:xi-time-part69c-localized-finite-quotient-order-spectrum}，
\[
\mathcal Q(G_S)
=
\left\{
n\ge 1:\ \forall p\in S,\ p\nmid n
\right\}.
\]
因此 \(\mathcal Q(G_S)\) 恰记录了哪些素数被完全排除在全部有限商阶谱之外。若
\[
\mathcal Q(G_S)=\mathcal Q(G_T),
\]
则被排除的素数集合一致，故 \(S=T\)。反向蕴含显然。

第二条只是同一刻画按单个素数逐点展开：若 \(p\in S\)，则由上式 \(p\) 不整除任意 \(n\in\mathcal Q(G_S)\)；反之若 \(p\notin S\)，则 \(n=p\) 已属于 \(\mathcal Q(G_S)\)。证毕。
\end{proof}

\begin{theorem}[局部化数系的连续自同态环完全由局部化环穷尽]\label{thm:xi-time-part69c-localized-dual-endomorphism-ring-exhaustion}
沿用定理 \ref{thm:xi-time-part69c-localized-dual-finite-visible-collapse} 的记号。则加法群 \(G_S\) 的端同态环满足
\[
\operatorname{End}(G_S)\cong \ZZ[S^{-1}],
\]
其中同构由
\[
q\longmapsto m_q,
\qquad
m_q(x):=qx
\]
给出。因而其对偶紧群的连续端同态环满足
\[
\operatorname{End}_{\mathrm{cts}}(\Sigma_S)\cong \ZZ[S^{-1}].
\]
进一步，连续自同构群满足
\[
\operatorname{Aut}_{\mathrm{cts}}(\Sigma_S)
\cong
\operatorname{Aut}(G_S)
\cong
\left\{
\pm\prod_{p\in S}p^{n_p}: n_p\in\ZZ
\right\}.
\]
\end{theorem}
\begin{proof}
定理 \ref{thm:xi-localized-endomorphism-ring-solenoid-indecomposable} 的前两条已经给出
\[
\operatorname{End}(G_S)\cong G_S=\ZZ[S^{-1}],
\]
并且该同构正由乘法映射 \(q\mapsto m_q\) 实现。Pontryagin 对偶把紧交换群上的连续端同态环识别为离散端端同态环的反环，故
\[
\operatorname{End}_{\mathrm{cts}}(\Sigma_S)\cong \operatorname{End}(G_S)^{\mathrm{op}}.
\]
由于 \(\ZZ[S^{-1}]\) 为交换环，其反环与自身一致，于是
\[
\operatorname{End}_{\mathrm{cts}}(\Sigma_S)\cong \ZZ[S^{-1}].
\]

最后，由定理 \ref{thm:xi-time-part69-localized-unit-automorphism-group-classification}，
\[
\operatorname{Aut}_{\mathrm{cts}}(\Sigma_S)
\cong
G_S^\times
=
\left\{
\pm\prod_{p\in S}p^{n_p}: n_p\in\ZZ
\right\}.
\]
而离散端自同构正对应于 \(\ZZ[S^{-1}]\) 的可逆乘法元，故同一公式也描述 \(\operatorname{Aut}(G_S)\)。证毕。
\end{proof}
\leanverified{paper\_xi\_time\_part69c\_localized\_dual\_endomorphism\_ring\_exhaustion}

\noindent
由此，局部化数系的对偶对象在连续有限值读出层面完全塌缩为单点；而在有限指数与有限商层面，其全部可见结构又分别被精确压缩为 \(S\)-自由整数的整除半格、对应的子群 \(\zeta\) 函数 Euler 因子以及有限商阶谱。进一步，在连续自同态动力学层面，全部端同态与自同构又完全坍缩为局部化环 \(\ZZ[S^{-1}]\) 及其单位群本身。于是，素数账本、有限指数格、有限商轮廓与连续动力学四者被同一局部化数据闭合地锁定。

\paragraph{局部化态射矩阵、同构刚性与 Gödel--Tate 仿射动力学的终端分类}
在局部化数系之间的 \(\operatorname{Hom}\)-分类、有限局部化直和的稀疏矩阵态射结构，以及 Gödel--Tate 仿射作用的唯一不动点、闭不变核与轨道闭包都已建立之后，还可把有限局部化账本与 \(2\)-进仿射表示两条主线进一步压缩为下列一组分类性结论。

\begin{theorem}[局部化数系 \texorpdfstring{$G_S=\ZZ[S^{-1}]$}{G\_S=Z[S^{-1}]} 之间的 \texorpdfstring{$\operatorname{Hom}$}{Hom}/\texorpdfstring{$\operatorname{Aut}$}{Aut} 完全分类]\label{thm:xi-time-part69d-localized-hom-aut-complete-classification}
\leanverified{paper\_xi\_time\_part69d\_localized\_hom\_aut\_complete\_classification}
设 \(S,T\subset\mathbb P\) 为有限素数集，并记
\[
G_S:=\ZZ[S^{-1}],
\qquad
G_T:=\ZZ[T^{-1}]
\]
为加法群。则任意群同态 \(\phi:G_S\to G_T\) 都唯一写成
\[
\phi(x)=qx
\qquad (x\in G_S)
\]
的形式，其中 \(q\in \QQ\)。进一步，
\[
\operatorname{Hom}(G_S,G_T)\cong
\begin{cases}
G_T,& S\subseteq T,\\[1mm]
0,& S\not\subseteq T,
\end{cases}
\]
其中同构由 \(q\mapsto (x\mapsto qx)\) 给出。特别地，
\[
\operatorname{End}(G_S)\cong G_S
\]
作为环同构成立，而自同构群满足
\[
\operatorname{Aut}(G_S)\cong G_S^\times
=
\left\{
\pm \prod_{p\in S}p^{n_p}:\ n_p\in\ZZ,\ \text{仅有限个 }n_p\neq 0
\right\}.
\]
\end{theorem}
\begin{proof}
定理 \ref{thm:xi-localized-integers-cross-hom-classification} 已给出
\[
\operatorname{Hom}(G_S,G_T)\cong
\begin{cases}
G_T,& S\subseteq T,\\[1mm]
0,& S\not\subseteq T.
\end{cases}
\]
其证明机制同时说明每个同态都由 \(\phi(1)\) 唯一决定。事实上，设 \(\phi:G_S\to G_T\) 为任意群同态，并记 \(q:=\phi(1)\in G_T\subset \QQ\)。对任意
\[
x=\frac{a}{b}\in G_S,
\qquad
a\in\ZZ,
\qquad
b=\prod_{p\in S}p^{m_p},
\]
在加法群 \(G_S\) 中有 \(bx=a\)。施加 \(\phi\) 得
\[
b\,\phi(x)=\phi(a)=aq.
\]
由于 \(G_T\subset \QQ\) 为无挠群，上式在 \(\QQ\) 中唯一解为
\[
\phi(x)=\frac{aq}{b}=qx.
\]
故全部同态均唯一来自有理数乘法。

当 \(S\subseteq T\) 时，\(q\in G_T\) 必满足 \(qG_S\subseteq G_T\)，故得到全部同态；当 \(S\not\subseteq T\) 时，定理 \ref{thm:xi-localized-integers-cross-hom-classification} 说明只能有零同态。

再由定理 \ref{thm:xi-localized-integers-endomorphism-automorphism-explicit}，
\[
\operatorname{End}(G_S)\cong G_S
\]
且
\[
\operatorname{Aut}(G_S)\cong G_S^\times.
\]
而局部化环 \(\ZZ[S^{-1}]\) 的单位恰为
\[
\left\{
\pm \prod_{p\in S}p^{n_p}:\ n_p\in\ZZ,\ \text{仅有限个 }n_p\neq 0
\right\},
\]
于是结论成立。
\end{proof}

\begin{theorem}[局部化数系紧对偶的连续自同态环完全显式]\label{thm:xi-time-part69d-localized-solenoid-endomorphism-ring}
设 \(S\subset\mathbb P\) 为有限素数集，并记
\[
G_S:=\ZZ[S^{-1}],
\qquad
\Sigma_S:=\widehat{G_S}.
\]
则存在自然环同构
\[
\operatorname{End}_{\mathrm{cts}}(\Sigma_S)\cong \ZZ[S^{-1}].
\]
\end{theorem}
\begin{proof}
Pontryagin 对偶把连续自同态反变地送到离散侧自同态，因此给出自然环反同构
\[
\operatorname{End}_{\mathrm{cts}}(\Sigma_S)
\cong
\operatorname{End}_{\ZZ}(G_S)^{\mathrm{op}}.
\]
而定理 \ref{thm:xi-time-part69d-localized-hom-aut-complete-classification} 与定理 \ref{thm:xi-localized-integers-endomorphism-automorphism-explicit} 已给出
\[
\operatorname{End}_{\ZZ}(G_S)\cong \ZZ[S^{-1}]
\]
作为环同构成立。由于 \(\ZZ[S^{-1}]\) 为交换环，其对立环与自身相同，故结论成立。证毕。
\end{proof}
\leanverified{paper\_xi\_time\_part69d\_localized\_solenoid\_endomorphism\_ring}

\begin{corollary}[局部化数系紧对偶的连续自同构群正是 \texorpdfstring{$S$}{S}-unit 群]\label{cor:xi-time-part69d-localized-solenoid-automorphism-sunits}
沿用定理 \ref{thm:xi-time-part69d-localized-solenoid-endomorphism-ring} 的记号，则
\[
\operatorname{Aut}_{\mathrm{cts}}(\Sigma_S)
\cong
\ZZ[S^{-1}]^\times
=
\left\{
\pm\prod_{p\in S}p^{n_p}: n_p\in\ZZ,\ \text{仅有限个 }n_p\neq 0
\right\}.
\]
\end{corollary}
\begin{proof}
Pontryagin 对偶在自同构层面给出自然同构
\[
\operatorname{Aut}_{\mathrm{cts}}(\Sigma_S)\cong \operatorname{Aut}_{\ZZ}(G_S).
\]
再由定理 \ref{thm:xi-time-part69d-localized-hom-aut-complete-classification} 与定理 \ref{thm:xi-localized-integers-endomorphism-automorphism-explicit}，
\[
\operatorname{Aut}_{\ZZ}(G_S)\cong \ZZ[S^{-1}]^\times.
\]
这正是所述带符号的 \(S\)-unit 群。证毕。
\end{proof}
\leanverified{paper\_xi\_time\_part69d\_localized\_solenoid\_automorphism\_sunits}

\begin{theorem}[有限局部化直和的矩阵 \texorpdfstring{$\operatorname{Hom}$}{Hom} 分类与同构判别]\label{thm:xi-time-part69d-localized-directsum-matrix-isomorphism-criterion}
\leanverified{paper\_xi\_time\_part69d\_localized\_directsum\_matrix\_isomorphism\_criterion}
设
\[
\mathbf S=(S_1,\dots,S_r),
\qquad
\mathbf T=(T_1,\dots,T_s)
\]
为有限素数集族，并记
\[
G_{\mathbf S}:=\bigoplus_{i=1}^{r}G_{S_i},
\qquad
G_{\mathbf T}:=\bigoplus_{j=1}^{s}G_{T_j}.
\]
则任意群同态 \(\Phi:G_{\mathbf S}\to G_{\mathbf T}\) 唯一对应于一个 \(s\times r\) 矩阵
\[
Q=(q_{ji})
\]
满足
\[
q_{ji}\in
\begin{cases}
G_{T_j},& S_i\subseteq T_j,\\
0,& S_i\not\subseteq T_j,
\end{cases}
\]
且作用为
\[
\Phi(x_1,\dots,x_r)=
\left(
\sum_{i=1}^{r}q_{1i}x_i,\ \dots,\ \sum_{i=1}^{r}q_{si}x_i
\right).
\]
因而
\[
\operatorname{Hom}(G_{\mathbf S},G_{\mathbf T})
\cong
\bigoplus_{j=1}^{s}\bigoplus_{i=1}^{r}\operatorname{Hom}(G_{S_i},G_{T_j}).
\]
进一步，
\[
G_{\mathbf S}\cong G_{\mathbf T}
\iff
r=s\ \text{且在重排后 }S_i=T_i\ \ (1\le i\le r).
\]
\end{theorem}
\begin{proof}
前两条正是定理 \ref{thm:xi-localized-directsum-hom-sparse-matrix} 的内容。下面证明同构判别。

若在重排后 \(S_i=T_i\)，则逐分量恒等映射的直和显然给出
\(
G_{\mathbf S}\cong G_{\mathbf T}
\)。

反过来，设
\[
\Phi:G_{\mathbf S}\xrightarrow{\sim} G_{\mathbf T}
\]
为一同构，\(\Psi=\Phi^{-1}\)。对任意有限素数集 \(U\)，记
\[
A_U:=\bigoplus_{\,i:\,S_i=U}G_{S_i}\subset G_{\mathbf S},
\qquad
B_U:=\bigoplus_{\,j:\,T_j=U}G_{T_j}\subset G_{\mathbf T}.
\]
在有限集合族
\(
\{S_1,\dots,S_r,T_1,\dots,T_s\}
\)
中取一个关于包含关系极大的素数集 \(U\)。由单对象 \(\operatorname{Hom}\)-分类，若 \(S_i=U\)，则 \(G_{S_i}\) 只能非零映入那些满足 \(U\subseteq T_j\) 的目标分量；而 \(U\) 极大迫使这等价于 \(T_j=U\)。因此
\[
\Phi(A_U)\subseteq B_U.
\]
同理，
\[
\Psi(B_U)\subseteq A_U.
\]
于是 \(\Phi|_{A_U}:A_U\to B_U\) 为同构。另一方面，\(A_U\) 与 \(B_U\) 都是若干个 \(G_U\) 的有限直和，而每个 \(G_U\) 的 \(\QQ\)-秩均为 \(1\)。故
\[
\operatorname{rank}_{\QQ}(A_U)=\#\{i:\ S_i=U\},
\qquad
\operatorname{rank}_{\QQ}(B_U)=\#\{j:\ T_j=U\}.
\]
同构保持 \(\QQ\)-秩，从而
\[
\#\{i:\ S_i=U\}=\#\{j:\ T_j=U\}.
\]
删去两侧全部等于 \(U\) 的分量后，对剩余有限集合族重复同样论证。由于待比较的素数集总数有限，归纳终止并得到两族完全相同的多重集。故
\[
r=s
\]
且经重排后
\[
S_i=T_i\qquad(1\le i\le r).
\]
结论成立。
\end{proof}

\leanverified{paper\_xi\_time\_part69d\_godel\_tate\_affine\_centralizer\_complete}
\begin{theorem}[Gödel--Tate 仿射表示的中心化子完全描述]\label{thm:xi-time-part69d-godel-tate-affine-centralizer-complete}
沿用定理 \ref{thm:killo-godel-semidir-affine-2adic} 与定理 \ref{thm:xi-godel-tate-affine-closed-subgroup-classification} 的记号。记
\[
T:=T_2(E_{\min})\cong \ZZ_2^2,
\qquad
H:=\mathcal R\rtimes_\Sigma \NN,
\]
\[
\Psi(r,k)(x):=B^k x+\operatorname{ev}_B(r),
\qquad
\operatorname{ev}_B(r)=\sum_{t\ge 1}r(p_t)B^{t-1}v\in L_v:=\ZZ_2v.
\]
则 \(\Psi(H)\) 在连续仿射自映射半群
\(
\mathrm{Aff}_{\mathrm{cts}}(T)
\)
中的中心化子恰为
\[
\operatorname{Cent}_{\mathrm{Aff}_{\mathrm{cts}}(T)}(\Psi(H))
=
\left\{
x\mapsto Ax:\ A\in \operatorname{End}_{\ZZ_2}(T),\ A|_{L_v}=\mathrm{id}_{L_v}
\right\}.
\]
若固定基底 \((v,w)\) 使
\(
T=\ZZ_2v\oplus \ZZ_2w
\)，则上述中心化子恰对应于矩阵集合
\[
\left\{
\begin{pmatrix}
1 & b\\
0 & d
\end{pmatrix}
:\ b,d\in\ZZ_2
\right\}
\subset M_2(\ZZ_2),
\]
而其可逆部分，即
\(
\mathrm{Aff}_{\mathrm{cts}}^\times(T)
\)
中的中心化子，为
\[
\left\{
\begin{pmatrix}
1 & b\\
0 & d
\end{pmatrix}
:\ b\in\ZZ_2,\ d\in\ZZ_2^\times
\right\}.
\]
\end{theorem}
\begin{proof}
任取连续仿射映射
\[
\Phi(x)=Ax+t,
\qquad
A\in \operatorname{End}_{\ZZ_2}(T),
\qquad
t\in T,
\]
并要求它与一切 \(\Psi(r,k)\) 交换。

先令 \(k=0\)。此时
\[
\Psi(r,0)(x)=x+\operatorname{ev}_B(r).
\]
交换关系给出
\[
\Phi\bigl(x+\operatorname{ev}_B(r)\bigr)
=
Ax+A\operatorname{ev}_B(r)+t,
\]
\[
\Psi(r,0)(\Phi(x))
=
Ax+t+\operatorname{ev}_B(r).
\]
故对一切 \(r\in\mathcal R\) 都有
\[
A\operatorname{ev}_B(r)=\operatorname{ev}_B(r).
\]
而定理 \ref{thm:xi-godel-tate-affine-closed-subgroup-classification} 的证明已经给出
\[
\overline{\operatorname{ev}_B(\mathcal R)}=L_v.
\]
由 \(A\) 的连续性，推出
\[
A|_{L_v}=\mathrm{id}_{L_v}.
\]

再令 \(r=0\)。对任意 \(k\ge 1\)，交换关系给出
\[
\Phi(B^k x)=AB^k x+t,
\qquad
\Psi(0,k)(\Phi(x))=B^kAx+B^k t.
\]
由于 \(B^k=2^{Lk}\) 是标量，故 \(AB^k=B^kA\)。于是必须有
\[
t=B^k t
\qquad (\forall k\ge 1).
\]
取 \(k=1\)，即
\[
(1-B)t=0.
\]
而 \(B=2^L\)，故
\[
1-B=1-2^L\in \ZZ_2^\times
\]
为单位，从而 \(t=0\)。

因此任何中心化元都必为线性映射 \(x\mapsto Ax\)，且满足 \(A|_{L_v}=\mathrm{id}_{L_v}\)。

反过来，若 \(A\in \operatorname{End}_{\ZZ_2}(T)\) 且 \(A|_{L_v}=\mathrm{id}_{L_v}\)，则对任意 \((r,k)\in H\) 与 \(x\in T\)，有
\[
A(\Psi(r,k)(x))
=
A(B^k x+\operatorname{ev}_B(r))
=
B^kAx+\operatorname{ev}_B(r)
=
\Psi(r,k)(Ax),
\]
故 \(x\mapsto Ax\) 确实中心化 \(\Psi(H)\)。

最后，在基底 \((v,w)\) 下，条件 \(A(v)=v\) 恰意味着矩阵第一列为 \((1,0)^\top\)，故
\[
A=
\begin{pmatrix}
1 & b\\
0 & d
\end{pmatrix},
\qquad
b,d\in\ZZ_2.
\]
而该矩阵可逆当且仅当 \(d\in\ZZ_2^\times\)。结论成立。
\end{proof}

\begin{theorem}[Gödel--Tate 仿射动力系统中非平凡周期轨道完全不存在]\label{thm:xi-time-part69d-godel-tate-affine-no-nontrivial-periodic-orbits}
\leanverified{paper\_xi\_time\_part69d\_godel\_tate\_affine\_no\_nontrivial\_periodic\_orbits}
沿用定理 \ref{thm:xi-time-part69d-godel-tate-affine-centralizer-complete} 的记号。对任意 \((r,k)\in H\)，考虑仿射映射
\[
\Psi(r,k):T\to T,
\qquad
x\mapsto B^k x+\operatorname{ev}_B(r).
\]
则其周期点结构满足：
\begin{enumerate}
  \item 若 \(k\ge 1\)，则 \(\Psi(r,k)\) 恰有唯一不动点
  \[
  x_{r,k}=(1-B^k)^{-1}\operatorname{ev}_B(r),
  \]
  且它也是 \(\Psi(r,k)\) 的唯一周期点；特别地不存在任何真周期 \(n>1\) 轨道。
  \item 若 \(k=0\) 且 \((r,0)\neq e\)，则 \(\Psi(r,0)\) 没有任何周期点。
\end{enumerate}
因此，\(H\) 的任一非单位元素在该 \(2\)-进仿射作用下都不产生非平凡周期轨道。
\end{theorem}
\begin{proof}
若 \(k\ge 1\)，则定理 \ref{thm:xi-time-part9g-affine-semigroup-unique-fixed-point} 已给出唯一不动点
\[
x_{r,k}=(1-B^k)^{-1}\operatorname{ev}_B(r).
\]
记 \(\Phi:=\Psi(r,k)\)。由不动点方程
\[
x_{r,k}=B^k x_{r,k}+\operatorname{ev}_B(r)
\]
可得对任意 \(x\in T\)
\[
\Phi(x)-x_{r,k}=B^k(x-x_{r,k}).
\]
归纳即得
\[
\Phi^n(x)-x_{r,k}=B^{kn}(x-x_{r,k})
\qquad (n\ge 1).
\]
若 \(x\) 为周期点，则存在 \(n\ge 1\) 使 \(\Phi^n(x)=x\)，故
\[
(1-B^{kn})(x-x_{r,k})=0.
\]
而
\[
1-B^{kn}=1-2^{Lkn}\in \ZZ_2^\times
\]
仍为单位，因此必有 \(x=x_{r,k}\)。故唯一周期点就是唯一不动点。

若 \(k=0\)，则
\[
\Psi(r,0)(x)=x+\operatorname{ev}_B(r).
\]
若 \((r,0)\neq e\)，则由表示 \(\Psi\) 的忠实性可知
\[
\operatorname{ev}_B(r)\neq 0.
\]
若存在周期点 \(x\)，则对某个 \(n\ge 1\) 有
\[
x+n\operatorname{ev}_B(r)=x,
\]
即
\[
n\operatorname{ev}_B(r)=0.
\]
但 \(T\cong \ZZ_2^2\) 是无挠 \(\ZZ_2\)-模，故这不可能发生。于是 \(\Psi(r,0)\) 没有任何周期点。
\end{proof}

\noindent
由此，有限局部化账本与 \(2\)-进 Gödel--Tate 仿射表示两条主线都达到分类级终点：在代数侧，单对象与有限直和对象的全部态射已被局部化乘法与稀疏矩阵完全穷尽；在动力学侧，仿射表示的外部对称只剩下对地址线 \(L_v\) 恒等的线性中心化子，而一切非单位仿射元都不产生任何非平凡周期轨道。于是，素数账本的线性可见性与 Gödel 半直积的 \(2\)-进动力学刚性在同一框架中同时闭合为严格的终端分类。


\paragraph{六模拓扑相图、精确暗块、boundary 奇偶盲性与纯鸽巢截断}
在最大非可缩纤维的模 \(6\) 比例极限、审计偶窗最小退化扇区的奇偶同伦相型、window-\(6\) 的 boundary sheet-parity 中心分裂，以及黄金分支在 Fibonacci 分母上的 \(\mathsf T\)/\(D^\ast\) 严格奇偶双相都已固定之后，还可进一步抽取出下列四条跨链后果。它们分别刻画精确暗块与顶端非可缩纤维的非耦合共存、根字典对 boundary 奇偶三格的严格盲性、审计偶窗最弱退化扇区在阿贝尔化中的精确 Fibonacci 份额，以及最小退化 Fibonacci 律一旦试图脱离审计窗口便会遭遇的首个纯计数截断。

\begin{corollary}[精确暗块与拓扑主纤维的非耦合共存]\label{cor:xi-time-part69c-darkblock-topfiber-decoupling}
沿用推论 \ref{cor:xi-binfold-zero-block-collision-gap-decoupling} 与定理 \ref{thm:xi-time-part62db-maxfiber-homotopy-phase-mod6} 的记号，记
\[
D_m:=\max_{x\in X_m}d_m(x),
\qquad
\widetilde D_m:=\max\{d_m(x):x\in X_m,\ |K_x|\not\simeq *\},
\qquad
\bar d_m:=\frac{2^m}{F_{m+2}}.
\]
若
\[
m\equiv 4\pmod 6,
\]
则同时成立
\[
|\mathcal Z_m|\ge F_{(m+2)/2},
\qquad
\widetilde D_m=D_m,
\]
并且
\[
\liminf_{\substack{m\to\infty\\ m\equiv 4\!\!\!\!\pmod 6}}
\frac{D_m}{\bar d_m}
\ge
1+C_\varphi.
\]
因此，半阶指数规模的精确暗块与不发生拓扑折损的顶端非可缩纤维，可以在同一分辨率相上严格共存。
\end{corollary}
\begin{proof}
当 \(m\equiv 4\pmod 6\) 时，\(m+2\equiv 0\pmod 6\)。故由推论 \ref{cor:xi-binfold-zero-block-collision-gap-decoupling}，
\[
|\mathcal Z_m|\ge F_{(m+2)/2}.
\]
同一同余类上，定理 \ref{thm:xi-time-part62db-maxfiber-homotopy-phase-mod6} 直接给出
\[
\widetilde D_m=D_m.
\]
另一方面，结论 \ref{con:xi-time-part9m-cphi-forces-collision-gap-and-peak-bias} 已证明
\[
\liminf_{m\to\infty}\frac{M_m}{\overline d_m}\ge 1+C_\varphi.
\]
将其中记号 \(M_m,\overline d_m\) 分别改写为 \(D_m,\bar d_m\)，再限制到子序列 \(m\equiv 4\pmod 6\)，便得到第三式。证毕。
\end{proof}

\begin{theorem}[window-\texorpdfstring{$6$}{6} 根字典的 boundary 盲性余维恰为 \texorpdfstring{$3$}{3}]\label{thm:xi-time-part69c-window6-root-boundary-blind-codim3}
\leanverified{paper\_xi\_time\_part69c\_window6\_root\_boundary\_blind\_codim3}
沿用定理 \ref{thm:xi-window6-cyclic-kernel-bc3-root-datum} 与定理 \ref{thm:xi-window6-center-exact-splitting-bdry-cyclic} 的记号。通过定理 \ref{thm:xi-foldbin-gauge-representation-distribution-law-m6-spectrum} 的直积分解，识别
\[
\bigl(G_6^{\mathrm{bin}}\bigr)^{\mathrm{ab}}
\cong
(\ZZ/2\ZZ)^{X_6}.
\]
令
\[
\rho:
\bigl(G_6^{\mathrm{bin}}\bigr)^{\mathrm{ab}}
\cong
(\ZZ/2\ZZ)^{X_6}
\longrightarrow
(\ZZ/2\ZZ)^{X_6^{\mathrm{cyc}}}
\]
为忘去 \(X_6^{\mathrm{bdry}}\) 三个坐标的自然投影，并记
\[
\iota_{\mathrm{ab}}:
Z\!\bigl(G_6^{\mathrm{bin}}\bigr)
\hookrightarrow
\bigl(G_6^{\mathrm{bin}}\bigr)^{\mathrm{ab}}
\]
为自然交换化嵌入。则
\[
\ker\rho
\cong
(\ZZ/2\ZZ)^3.
\]
更精确地，
\[
\ker\rho
=
\iota_{\mathrm{ab}}(Z_{\partial}),
\qquad
Z_{\partial}
=
\langle \tau_w:\ w\in X_6^{\mathrm{bdry}}\rangle
\subset Z\!\bigl(G_6^{\mathrm{bin}}\bigr).
\]
因此，任何仅经由 \(X_6^{\mathrm{cyc}}\) 上的 \(B_3/C_3\) 根字典读取的可观测量，都必然对这一规范三维 boundary 奇偶子格失明。
\end{theorem}
\begin{proof}
由定理 \ref{thm:xi-foldbin-gauge-representation-distribution-law-m6-spectrum}，
\[
G_6^{\mathrm{bin}}
\cong
\prod_{w\in X_6}\mathfrak S_{d_6(w)}.
\]
由于 \(m=6\) 时全部纤维大小只可能为 \(2,3,4\)，而
\[
\mathfrak S_2^{\mathrm{ab}}\cong
\mathfrak S_3^{\mathrm{ab}}\cong
\mathfrak S_4^{\mathrm{ab}}
\cong
\ZZ/2\ZZ,
\]
故交换化给出规范识别
\[
\bigl(G_6^{\mathrm{bin}}\bigr)^{\mathrm{ab}}
\cong
\prod_{w\in X_6}\mathfrak S_{d_6(w)}^{\mathrm{ab}}
\cong
(\ZZ/2\ZZ)^{X_6}.
\]
再由定理 \ref{thm:xi-window6-cyclic-kernel-bc3-root-datum}，
\[
X_6=X_6^{\mathrm{cyc}}\sqcup X_6^{\mathrm{bdry}},
\qquad
|X_6^{\mathrm{cyc}}|=18,
\qquad
|X_6^{\mathrm{bdry}}|=3.
\]
因此忘去 boundary 三坐标的投影 \(\rho\) 的核恰为
\[
(\ZZ/2\ZZ)^{X_6^{\mathrm{bdry}}}
\cong
(\ZZ/2\ZZ)^3.
\]

另一方面，定理 \ref{thm:xi-window6-center-exact-splitting-bdry-cyclic} 已经把
\[
Z_{\partial}
=
\langle \tau_w:\ w\in X_6^{\mathrm{bdry}}\rangle
\cong
(\ZZ/2\ZZ)^{X_6^{\mathrm{bdry}}}
\]
识别为 \(Z(G_6^{\mathrm{bin}})\) 的规范直和因子。故
\[
\ker\rho=\iota_{\mathrm{ab}}(Z_{\partial})\cong (\ZZ/2\ZZ)^3.
\]
凡仅通过 \(X_6^{\mathrm{cyc}}\) 的根字典读取的可观测量，都必经由 \(\rho\) 因子化，因而在 \(\ker\rho=\iota_{\mathrm{ab}}(Z_{\partial})\) 上必为常值。证毕。
\end{proof}

\begin{corollary}[审计偶窗最弱退化扇区的奇偶份额恰为 \texorpdfstring{$F_m/F_{m+2}$}{F_m/F_{m+2}}]\label{cor:xi-time-part69c-audited-even-minsector-abel-share}
\leanverified{paper\_xi\_time\_part69c\_audited\_even\_minsector\_abel\_share}
对审计偶窗
\[
m\in\{6,8,10,12\},
\]
记
\[
\mathcal S_{m,*}:=\mathcal S_{m,d_{\min}(m)},
\qquad
\Lambda_m^{\min}:=(\ZZ/2\ZZ)^{\mathcal S_{m,*}}
\subset
\bigl(G_m^{\mathrm{bin}}\bigr)^{\mathrm{ab}},
\]
其中 \(\Lambda_m^{\min}\) 由定理 \ref{thm:xi-audited-even-abel-center-fibonacci-splitting} 的规范直和分裂给出。则
\[
\operatorname{rank}_{\FF_2}\Lambda_m^{\min}=F_m,
\qquad
\operatorname{rank}_{\FF_2}\!\bigl(G_m^{\mathrm{bin}}\bigr)^{\mathrm{ab}}=F_{m+2},
\]
从而
\[
\frac{\operatorname{rank}_{\FF_2}\Lambda_m^{\min}}
{\operatorname{rank}_{\FF_2}\!\bigl(G_m^{\mathrm{bin}}\bigr)^{\mathrm{ab}}}
=
\frac{F_m}{F_{m+2}}.
\]
具体地，
\[
\frac{8}{21},\qquad
\frac{21}{55},\qquad
\frac{55}{144},\qquad
\frac{144}{377}
\]
分别对应 \(m=6,8,10,12\)。
\end{corollary}
\begin{proof}
由定理 \ref{thm:xi-audited-even-abel-center-fibonacci-splitting}，
\[
\bigl(G_m^{\mathrm{bin}}\bigr)^{\mathrm{ab}}
\cong
(\ZZ/2\ZZ)^{\mathcal S_{m,*}}
\oplus
(\ZZ/2\ZZ)^{X_m\setminus\mathcal S_{m,*}},
\]
并且
\[
|\mathcal S_{m,*}|=F_m.
\]
故
\[
\operatorname{rank}_{\FF_2}\Lambda_m^{\min}=F_m.
\]
另一方面，审计偶窗上由定理 \ref{thm:xi-foldbin-audited-even-center-collapse-abel-full-rank}，
\[
\bigl(G_m^{\mathrm{bin}}\bigr)^{\mathrm{ab}}
\cong
(\ZZ/2\ZZ)^{|X_m|}.
\]
再由稳定类型计数律 \(|X_m|=F_{m+2}\)，得到
\[
\operatorname{rank}_{\FF_2}\!\bigl(G_m^{\mathrm{bin}}\bigr)^{\mathrm{ab}}=F_{m+2}.
\]
于是所示比值公式成立，四个具体分数只需代入 \(m=6,8,10,12\) 即得。证毕。
\end{proof}

\begin{theorem}[审计偶窗最小退化 Fibonacci 外推的首个纯计数截断]\label{thm:xi-time-part69c-minsector-fib-extension-pigeonhole-cutoff}
\leanverified{paper\_xi\_time\_part69c\_minsector\_fib\_extension\_pigeonhole\_cutoff}
把审计恒等式
\[
d_{\min}(m)=F_{m/2}
\qquad
\bigl(m\in\{6,8,10,12\}\bigr)
\]
视为全部偶窗上的延拓候选。则纯计数必要条件已经强制
\[
2^m\ge F_{m+2}F_{m/2}.
\]
该条件在
\[
m=22
\]
时仍成立，但在
\[
m=24
\]
时首次失效；更精确地，
\[
F_{24}F_{11}=46368\cdot 89=4126752<2^{22}=4194304,
\]
而
\[
F_{26}F_{12}=121393\cdot 144=17480592>2^{24}=16777216.
\]
因此，若试图把审计偶窗上的最小退化 Fibonacci 律无条件外推到全部偶窗，则其首个纯鸽巢截断最迟发生在 \(m=24\)。
\end{theorem}
\begin{proof}
若对某个偶数 \(m\) 仍有
\[
d_{\min}(m)=F_{m/2},
\]
则每条纤维都至少含有 \(F_{m/2}\) 个微态。由稳定类型计数律 \(|X_m|=F_{m+2}\) 与恒等式
\[
2^m=\sum_{x\in X_m}d_m(x),
\]
立即得到必要条件
\[
2^m\ge |X_m|\,d_{\min}(m)=F_{m+2}F_{m/2}.
\]
对 \(m=22\) 与 \(m=24\) 直接代入，便得到
\[
F_{24}F_{11}=46368\cdot 89=4126752<2^{22}=4194304,
\]
\[
F_{26}F_{12}=121393\cdot 144=17480592>2^{24}=16777216.
\]
因此这一纯计数必要条件在 \(m=22\) 仍未破裂，而在 \(m=24\) 已经失效；于是任何无条件外推都不可能延续到 \(m=24\) 之后。证毕。
\end{proof}

\noindent
由此，最大非可缩纤维的模 \(6\) 相图、半阶精确暗块、window-\(6\) 根字典对 boundary 奇偶的三维盲核，以及黄金分支在 Fibonacci 分母上的双相输运常数，便不再构成彼此平行的孤立事实：前两者说明同一同余相上可以同时出现大规模谱消失与无折损拓扑主纤维，中者表明可见的根字典天然遗漏一条规范的中心奇偶方向，而纯鸽巢截断又说明最小退化 Fibonacci 律不能脱离审计窗口作无条件延展。因此，该折叠系统的可见图样同时受同余类拓扑、精确暗块与侧向采样异常三类彼此不可约机制共同支配。

\paragraph{二原子经验谱、临界容量反演、Bernoulli 残差比与奇素数缺陷尾谱}
在 bin-fold 的两态一致渐近、critical-scale 容量极限曲线、fold-gauge 常数的 Bernoulli 递归剥离以及输出位势的奇素数单位根扭曲四阶展开已经建立之后，还可继续抽取出下列七条后果。它们分别刻画稳定层均匀经验谱的严格两点弱极限、连续容量极限曲线的分布二阶导数、目标成功率下的最优辅助比特预算、由 \((\log C_m)\) 序列恢复 Bernoulli 塔时的 Richardson 残差比、奇素数模 RH 缺陷的可和性与其 \(\ell^\alpha\) 临界指数，以及最小非平凡 prime-twist 谱半径对 \(\lambda\) 的严格二次接近律与最终单调逼近。

\begin{theorem}[bin-fold 纤维谱经自然缩放后的稳定层两点弱极限]\label{thm:xi-time-part70-binfold-uniform-empirical-two-point-law}
\leanverified{paper\_xi\_time\_part70\_binfold\_uniform\_empirical\_two\_point\_law}
定义尺度化多重度
\[
Z_m^{\mathrm{bin}}(w):=\left(\frac{\varphi}{2}\right)^m d_m^{\mathrm{bin}}(w)
\qquad (w\in X_m),
\]
以及稳定层上的均匀经验测度
\[
\nu_m:=\frac{1}{|X_m|}\sum_{w\in X_m}\delta_{Z_m^{\mathrm{bin}}(w)}.
\]
再记
\[
A_0^{(m)}:=\{w\in X_m:\ w_m=0\},
\qquad
A_1^{(m)}:=\{w\in X_m:\ w_m=1\}.
\]
则有弱收敛
\[
\nu_m\xrightarrow{\mathrm w}\nu_\infty
:=
\frac{1}{\varphi}\,\delta_1+\frac{1}{\varphi^2}\,\delta_{\varphi^{-1}}.
\]
等价地，对任意有界连续函数 \(f\)，都有
\[
\lim_{m\to\infty}\frac{1}{|X_m|}\sum_{w\in X_m}f\!\bigl(Z_m^{\mathrm{bin}}(w)\bigr)
=
\frac{1}{\varphi}f(1)+\frac{1}{\varphi^2}f(\varphi^{-1}).
\]
\end{theorem}
\begin{proof}
由定理 \ref{thm:fold-bin-two-state-asymptotic}，存在 \(\varepsilon_m\to 0\) 使得对全部 \(w\in X_m\) 一致有
\[
\bigl|Z_m^{\mathrm{bin}}(w)-\varphi^{-w_m}\bigr|\le \varepsilon_m.
\]
因此对充分大的 \(m\)，所有 \(Z_m^{\mathrm{bin}}(w)\) 都落在某个固定紧区间内。任取有界连续函数 \(f\)，其在该紧区间上一致连续。于是若定义
\[
\eta_{m,0}:=\sup_{w\in A_0^{(m)}}\bigl|f(Z_m^{\mathrm{bin}}(w))-f(1)\bigr|,
\qquad
\eta_{m,1}:=\sup_{w\in A_1^{(m)}}\bigl|f(Z_m^{\mathrm{bin}}(w))-f(\varphi^{-1})\bigr|,
\]
则 \(\eta_{m,0},\eta_{m,1}\to 0\)。从而
\[
\frac{1}{|X_m|}\sum_{w\in X_m}f\!\bigl(Z_m^{\mathrm{bin}}(w)\bigr)
=
\frac{|A_0^{(m)}|}{|X_m|}\bigl(f(1)+o(1)\bigr)
+
\frac{|A_1^{(m)}|}{|X_m|}\bigl(f(\varphi^{-1})+o(1)\bigr).
\]
再由稳定词计数
\[
|A_0^{(m)}|=F_{m+1},
\qquad
|A_1^{(m)}|=F_m,
\qquad
|X_m|=F_{m+2},
\]
以及 Fibonacci 比值极限
\[
\frac{F_{m+1}}{F_{m+2}}\to \frac{1}{\varphi},
\qquad
\frac{F_m}{F_{m+2}}\to \frac{1}{\varphi^2},
\]
便得到所示极限。证毕。
\end{proof}

\begin{theorem}[连续容量极限曲线的分布二阶导数恰为两原子测度]\label{thm:xi-time-part70-continuous-capacity-two-atom-second-derivative}
\leanverified{paper\_xi\_time\_part70\_continuous\_capacity\_two\_atom\_second\_derivative}
定义尺度化连续容量
\[
\widehat C_m(s):=
2^{-m}\,\mathcal C_m^{\mathrm{bin,cont}}\!\left(
s\left(\frac{2}{\varphi}\right)^m
\right)
\qquad (s\ge 0),
\]
并记其极限函数
\[
L(s):=\frac{\varphi}{\sqrt5}\min\{1,s\}+\frac{1}{\sqrt5}\min\{\varphi^{-1},s\}.
\]
则在分布意义下有精确恒等式
\[
-D^2L
=
\frac{1}{\sqrt5}\,\delta_{\varphi^{-1}}
+
\frac{\varphi}{\sqrt5}\,\delta_1.
\]
换言之，极限容量曲线只有两个不可消去的折点，其位置分别位于 \(s=\varphi^{-1}\) 与 \(s=1\)，对应原子质量分别为 \(1/\sqrt5\) 与 \(\varphi/\sqrt5\)。
\end{theorem}
\begin{proof}
这正是定理 \ref{thm:xi-foldbin-underresolution-two-atomic-curvature-measure} 中曲率测度公式的直接重写。若将
\[
L(s)=
\begin{cases}
\dfrac{\varphi^2}{\sqrt5}\,s, & 0\le s\le \varphi^{-1},\\[6pt]
\dfrac{\varphi}{\sqrt5}\,s+\dfrac{1}{\varphi\sqrt5}, & \varphi^{-1}\le s\le 1,\\[8pt]
1, & s\ge 1
\end{cases}
\]
逐段求导，则 \(L'\) 在 \(s=\varphi^{-1}\) 处的跃降为
\[
\frac{\varphi^2}{\sqrt5}-\frac{\varphi}{\sqrt5}=\frac{1}{\sqrt5},
\]
在 \(s=1\) 处的跃降为
\[
\frac{\varphi}{\sqrt5}-0=\frac{\varphi}{\sqrt5}.
\]
因此分布导数满足
\[
-D^2L
=
\frac{1}{\sqrt5}\,\delta_{\varphi^{-1}}
+
\frac{\varphi}{\sqrt5}\,\delta_1.
\]
证毕。
\end{proof}

\begin{theorem}[目标成功率下最优辅助比特预算的严格分段阈值]\label{thm:xi-time-part70-target-success-optimal-budget-threshold}
\leanverified{paper\_xi\_time\_part70\_target\_success\_optimal\_budget\_threshold}
沿用定理 \ref{thm:xi-foldbin-dyadic-capacity-critical-window-limit} 的记号。对任意目标成功率 \(\sigma\in(0,1)\)，定义最小预算
\[
B_m^\star(\sigma):=
\min\bigl\{B\in\mathbb N:\ \mathrm{Succ}_m^{\mathrm{bin}}(B)\ge \sigma\bigr\}.
\]
再定义连续阈值常数
\[
s_\sigma:=
\begin{cases}
\dfrac{\sqrt5}{\varphi^2}\,\sigma,
& 0<\sigma\le \dfrac{\varphi}{\sqrt5},\\[8pt]
\dfrac{\sqrt5\,\sigma-\varphi^{-1}}{\varphi},
& \dfrac{\varphi}{\sqrt5}\le \sigma<1.
\end{cases}
\]
则有
\[
B_m^\star(\sigma)
=
m\log_2\!\left(\frac{2}{\varphi}\right)+\log_2 s_\sigma+O(1).
\]
更一般地，若整数序列 \((B_m)_{m\ge 1}\) 满足
\[
2^{B_m}\left(\frac{\varphi}{2}\right)^m\longrightarrow s\in(0,\infty),
\]
则
\[
\mathrm{Succ}_m^{\mathrm{bin}}(B_m)\longrightarrow L(s),
\]
其中 \(L\) 即定理 \ref{thm:xi-time-part70-continuous-capacity-two-atom-second-derivative} 的极限曲线；因此 \(s_\sigma\) 正是方程 \(L(s)=\sigma\) 的唯一解。
\end{theorem}
\begin{proof}
第二个结论正是定理 \ref{thm:xi-foldbin-dyadic-capacity-critical-window-limit} 与推论 \ref{cor:xi-foldbin-critical-scale-success-curve-inversion} 的内容，其中极限成功曲线即
\[
L(s)=
\begin{cases}
\dfrac{\varphi^2}{\sqrt5}\,s, & 0\le s\le \varphi^{-1},\\[6pt]
\dfrac{\varphi}{\sqrt5}\,s+\dfrac{1}{\varphi\sqrt5}, & \varphi^{-1}\le s\le 1,\\[8pt]
1, & s\ge 1.
\end{cases}
\]
因此 \(s_\sigma=L^{-1}(\sigma)\) 的显式公式恰为上式两段线性方程的解。

现证明 \(B_m^\star(\sigma)\) 的渐近公式。由于 \(\sigma\in(0,1)\)，故 \(s_\sigma\in(0,1)\)。取
\[
0<\varepsilon<\min\{s_\sigma,1-s_\sigma\}.
\]
由 \(L\) 在 \([0,1]\) 上严格递增可知
\[
L(s_\sigma-\varepsilon)<\sigma<L(s_\sigma+\varepsilon).
\]
再由 critical-scale 极限，分别取
\[
\widetilde B_m^-:=\left\lfloor \log_2\!\Bigl((s_\sigma-\varepsilon)\Bigl(\frac{2}{\varphi}\Bigr)^m\Bigr)\right\rfloor,
\qquad
\widetilde B_m^+:=\left\lceil \log_2\!\Bigl((s_\sigma+\varepsilon)\Bigl(\frac{2}{\varphi}\Bigr)^m\Bigr)\right\rceil,
\]
则当 \(m\) 充分大时有
\[
\mathrm{Succ}_m^{\mathrm{bin}}(\widetilde B_m^-)<\sigma,
\qquad
\mathrm{Succ}_m^{\mathrm{bin}}(\widetilde B_m^+)\ge \sigma.
\]
由于 \(\mathrm{Succ}_m^{\mathrm{bin}}(B)\) 关于 \(B\) 单调不减，而 \(B_m^\star(\sigma)\) 按定义是达到成功率阈值的最小整数预算，遂得
\[
\widetilde B_m^-+1\le B_m^\star(\sigma)\le \widetilde B_m^+.
\]
把上式展开后便得到
\[
\log_2\!\Bigl((s_\sigma-\varepsilon)\Bigl(\frac{2}{\varphi}\Bigr)^m\Bigr)
\le
B_m^\star(\sigma)
\le
\log_2\!\Bigl((s_\sigma+\varepsilon)\Bigl(\frac{2}{\varphi}\Bigr)^m\Bigr)+1.
\]
由于 \(\varepsilon\) 固定，端点与
\[
m\log_2\!\left(\frac{2}{\varphi}\right)+\log_2 s_\sigma
\]
之差都是有界量，故结论成立。证毕。
\end{proof}

\begin{theorem}[由 \texorpdfstring{$(\log C_m)$}{(log Cm)} 序列恢复 Bernoulli 塔时的 Richardson 残差比极限]\label{thm:xi-time-part70-logcm-richardson-residual-ratio}
\leanverified{paper\_xi\_time\_part70\_logcm\_richardson\_residual\_ratio}
记
\[
q:=\frac{\varphi}{2}\in(0,1),
\qquad
E_m:=\log C_m-\log(2\pi),
\]
\[
a_r:=
\frac{B_{2r}}{r(2r-1)}
\bigl(\varphi^{-1}+\varphi^{2r-3}\bigr)
\qquad (r\ge 1).
\]
对每个固定整数 \(R\ge 0\)，定义 \(R\) 阶 Richardson 残差
\[
\mathcal R_m^{(R)}
:=
E_m-\sum_{r=1}^{R}a_r q^{(2r-1)m}.
\]
则有严格比值极限
\[
\frac{\mathcal R_{m+1}^{(R)}}{\mathcal R_m^{(R)}}
\longrightarrow
q^{2R+1}
=
\left(\frac{\varphi}{2}\right)^{2R+1}.
\]
换言之，一旦减去前 \(R\) 层 Bernoulli 模态，剩余误差的相邻比值必然锁定到奇数次黄金幂 \((\varphi/2)^{2R+1}\)。
\end{theorem}
\begin{proof}
对任意固定 \(R\ge 0\)，将定理 \ref{thm:fold-bin-gauge-constant-stirling-bernoulli-hierarchy} 应用于 \(R+1\) 阶，得到
\[
E_m
=
\sum_{r=1}^{R+1}a_r q^{(2r-1)m}
+O\!\bigl(q^{(2R+3)m}\bigr).
\]
因此
\[
\mathcal R_m^{(R)}
=
a_{R+1}q^{(2R+1)m}
+O\!\bigl(q^{(2R+3)m}\bigr).
\]
又因为偶 Bernoulli 数满足 \(B_{2r}\neq 0\)（\(r\ge 1\)），而 \(\varphi^{-1}+\varphi^{2r-3}>0\)，故
\[
a_{R+1}\neq 0.
\]
于是
\[
\mathcal R_m^{(R)}
=
a_{R+1}q^{(2R+1)m}\bigl(1+O(q^{2m})\bigr),
\]
从而
\[
\frac{\mathcal R_{m+1}^{(R)}}{\mathcal R_m^{(R)}}
=
\frac{a_{R+1}q^{(2R+1)(m+1)}\bigl(1+O(q^{2m})\bigr)}
{a_{R+1}q^{(2R+1)m}\bigl(1+O(q^{2m})\bigr)}
\longrightarrow
q^{2R+1}.
\]
证毕。
\end{proof}

\begin{theorem}[奇素数模 RH 缺陷总质量的可和性]\label{thm:xi-time-part70-odd-prime-rh-defect-summable-mass}
\leanverified{paper\_xi\_time\_part70\_odd\_prime\_rh\_defect\_summable\_mass}
记
\[
\lambda:=\varphi^2,
\qquad
\rho_{p,1}:=\rho\!\bigl(M(e^{2\pi i/p})\bigr)
\qquad
(p\ge 3\ \text{为奇素数}),
\]
并定义扭曲指数与其缺陷
\[
\theta_{p,1}:=\frac{\log \rho_{p,1}}{\log\lambda},
\qquad
\delta_p:=1-\theta_{p,1}.
\]
则有显式展开
\[
\delta_p
=
\frac{4\pi^2(6\sqrt5-5)}{250\log(\varphi^2)}\,\frac{1}{p^2}
+O\!\left(\frac{1}{p^4}\right).
\]
从而
\[
\sum_{\substack{p\ge 3\\ p\ \text{为奇素数}}}\delta_p<\infty,
\qquad
\sum_{\substack{p>P\\ p\ \text{为奇素数}}}\delta_p=O(P^{-1}).
\]
因此，尽管奇素数方向上的平方根门槛失效在频率上呈余有限出现，其缺陷幅度却形成可和尾谱。
\end{theorem}
\begin{proof}
由定理 \ref{thm:xi-output-potential-odd-prime-rh-obstruction-index-to-one}，
\[
\theta_{p,1}
=
1-\frac{6\sqrt5-5}{250\log\lambda}\Bigl(\frac{2\pi}{p}\Bigr)^2
+\frac{7+24\sqrt5}{75000\log\lambda}\Bigl(\frac{2\pi}{p}\Bigr)^4
+O(p^{-6}).
\]
将 \(\lambda=\varphi^2\) 代入并移项，便得到
\[
\delta_p
=
\frac{4\pi^2(6\sqrt5-5)}{250\log(\varphi^2)}\,\frac{1}{p^2}
+O(p^{-4}).
\]
特别地，存在常数 \(c_1,c_2>0\) 与奇素数阈值 \(p_0\)，使得对所有奇素数 \(p\ge p_0\) 都有
\[
\frac{c_1}{p^2}\le \delta_p\le \frac{c_2}{p^2}.
\]
于是
\[
\sum_{\substack{p\ge p_0\\ p\ \text{为奇素数}}}\delta_p
\ll
\sum_{n\ge p_0}\frac{1}{n^2}
<\infty,
\]
从而总和收敛。对尾和，同样有
\[
\sum_{\substack{p>P\\ p\ \text{为奇素数}}}\delta_p
\ll
\sum_{n>P}\frac{1}{n^2}
=
O(P^{-1}).
\]
证毕。
\end{proof}

\begin{corollary}[奇素数 RH 缺陷尾谱的 \texorpdfstring{$\ell^\alpha$}{ell alpha} 临界指数恰为 \texorpdfstring{$1/2$}{1/2}]\label{cor:xi-time-part70-odd-prime-rh-defect-ell-alpha-threshold}
\leanverified{paper\_xi\_time\_part70\_odd\_prime\_rh\_defect\_ell\_alpha\_threshold}
沿用定理 \ref{thm:xi-time-part70-odd-prime-rh-defect-summable-mass} 的记号。则对任意 \(\alpha>0\)，有
\[
(\delta_p)_{p\ \text{为奇素数}}\in \ell^\alpha
\qquad\Longleftrightarrow\qquad
\alpha>\frac12.
\]
\end{corollary}
\begin{proof}
由定理 \ref{thm:xi-time-part70-odd-prime-rh-defect-summable-mass}，存在常数 \(c_1,c_2>0\) 与 \(p_0\) 使得对全部奇素数 \(p\ge p_0\) 都有
\[
c_1p^{-2}\le \delta_p\le c_2p^{-2}.
\]
于是
\[
\delta_p^\alpha\asymp p^{-2\alpha}
\qquad (p\to\infty,\ p\ \text{为奇素数}).
\]
而素数级数满足
\[
\sum_{p}p^{-s}<\infty
\qquad\Longleftrightarrow\qquad
s>1.
\]
因此
\[
\sum_{\substack{p\ge 3\\ p\ \text{为奇素数}}}\delta_p^\alpha<\infty
\qquad\Longleftrightarrow\qquad
\sum_{\substack{p\ge 3\\ p\ \text{为奇素数}}}p^{-2\alpha}<\infty
\qquad\Longleftrightarrow\qquad
2\alpha>1,
\]
即所求结论。证毕。
\end{proof}

\begin{theorem}[奇素数最小非平凡单位根扭曲的谱半径二次接近律与最终严格单调逼近]\label{thm:xi-time-part70-prime-twist-spectral-radius-quadratic-approach}
\leanverified{paper\_xi\_time\_part70\_prime\_twist\_spectral\_radius\_quadratic\_approach}
沿用推论 \ref{cor:real-input-40-collision-twist-fourth} 的记号，并记
\[
\lambda:=\rho(1)=\varphi^2,
\qquad
\rho_p:=\rho\!\left(e^{2\pi i/p}\right)
\qquad
(p\ge 3\ \text{为奇素数}).
\]
则成立：
\begin{enumerate}
  \item 有精确二次接近律
  \[
  \boxed{\;
  p^2\left(1-\frac{\rho_p}{\lambda}\right)
  \longrightarrow
  \frac{6\sqrt5-5}{250}\,(2\pi)^2.\;}
  \]
  \item 存在奇素数阈值 \(p_0\)，使得对任意奇素数 \(p<q\) 且 \(p,q\ge p_0\)，都有
  \[
  \boxed{\;
  \rho_p<\rho_q<\lambda.\;}
  \]
\end{enumerate}
换言之，最小非平凡 prime twist 的谱半径不仅以严格 \(p^{-2}\) 主项从下方逼近 \(\lambda\)，而且在充分大素数区间上按模数严格单调恢复。
\end{theorem}
\begin{proof}
由推论 \ref{cor:real-input-40-collision-twist-fourth}，当 \(t\to 0\) 时有
\[
\log\frac{\rho(t)}{\lambda}
=
-A t^2+B t^4+O(t^6),
\]
其中
\[
A:=\frac{6\sqrt5-5}{250}>0,
\qquad
B:=\frac{7+24\sqrt5}{75000}.
\]
将其指数化，得到
\[
\frac{\rho(t)}{\lambda}
=
\exp\!\bigl(-A t^2+B t^4+O(t^6)\bigr)
=
1-A t^2+O(t^4).
\]
对 \(t=t_p:=2\pi/p\) 代入，即得
\[
1-\frac{\rho_p}{\lambda}
=
A\left(\frac{2\pi}{p}\right)^2+O(p^{-4}),
\]
从而
\[
p^2\left(1-\frac{\rho_p}{\lambda}\right)
\longrightarrow
A(2\pi)^2
=
\frac{6\sqrt5-5}{250}\,(2\pi)^2,
\]
这就证明了第一条。

再令
\[
R(t):=\frac{\rho(t)}{\lambda}.
\]
由上式可得
\[
\log R(t)=-A t^2+B t^4+O(t^6).
\]
对其求导，得到
\[
\frac{R'(t)}{R(t)}
=
-2At+4Bt^3+O(t^5)
=
-2At\bigl(1+O(t^2)\bigr).
\]
因此存在 \(t_0>0\)，使得当 \(0<t\le t_0\) 时，
\[
\frac{R'(t)}{R(t)}<0.
\]
又因 \(R(t)>0\)，可知 \(R'(t)<0\)；故 \(R\) 在 \((0,t_0]\) 上严格递减。另一方面，由
\[
R(t)=1-A t^2+O(t^4)
\]
可知在同一小邻域内 \(R(t)<1\)。

现取奇素数阈值 \(p_0\) 使得
\[
\frac{2\pi}{p_0}\le t_0.
\]
若 \(p<q\) 且 \(p,q\ge p_0\)，则
\[
t_p=\frac{2\pi}{p}>\frac{2\pi}{q}=t_q,
\]
而 \(R\) 在 \((0,t_0]\) 上严格递减，遂得
\[
R(t_p)<R(t_q)<1.
\]
两边同乘以 \(\lambda\)，即
\[
\rho_p<\rho_q<\lambda.
\]
第二条得证。证毕。
\end{proof}

\noindent
上述七条结果表明，二原子塌缩并不只停留在 bin-fold 的输出实验或连续容量曲线上，而会同时外显为稳定层均匀经验谱的严格弱极限、目标成功率反演中的唯一分段阈值、\((\log C_m)\) 的奇次黄金幂残差比，以及奇素数单位根扭曲缺陷的可和尾谱与谱半径恢复律。因而，\(\{\varphi^{-1},1\}\) 这两个黄金比例尺度既支配欠分辨容量几何，也支配 Bernoulli 递归剥离后的剩余模态；而在奇素数单位根扭曲方向上，坏模数的频率与缺陷的振幅不仅被严格分离为“余有限出现”与“总质量可和”两种不同层级，而且最小非平凡扭曲的谱半径还会按严格 \(p^{-2}\) 主项最终单调地从下方逼近 \(\lambda\)。于是，稳定层二点律、资源阈值反演、Richardson 诊断与奇素数 RH 缺陷最终被压缩进同一条黄金尺度主线。

\paragraph{稳定层二态统计极限与规范体积修正的一阶账本}
在稳定层均匀经验谱的两点弱极限、统一层矩函数以及 fold-gauge 常数恢复链已经建立之后，还可把二态一致渐近对均匀层平均量与规范体积修正的支配压缩为下列四条直接后果。它们分别给出缩放退化度的量化两点律、全部缩放矩与幂和首项常数、均匀平均对数退化度的二态展开，以及 \(g_m-\kappa_m+1\) 的精确首项。

\begin{theorem}[均匀稳定类型下缩放退化度的量化两点极限律]\label{thm:xi-time-part70a-uniform-scaled-degeneracy-quantitative}
\leanverified{paper\_xi\_time\_part70a\_uniform\_scaled\_degeneracy\_quantitative}
令
\[
u_m(w):=\frac{1}{|X_m|}\qquad (w\in X_m),
\]
并取 \(W_m\sim u_m\)。定义缩放退化度随机变量
\[
Y_m:=\left(\frac{\varphi}{2}\right)^m d_m^{\mathrm{bin}}(W_m).
\]
则对任意有界 Lipschitz 函数 \(f:\RR\to\RR\)，都有
\[
\mathbb E[f(Y_m)]
=
\frac{F_{m+1}}{F_{m+2}}\,f(1)
+
\frac{F_m}{F_{m+2}}\,f(\varphi^{-1})
+O\!\left(\left(\frac{\varphi}{2}\right)^m\right),
\]
其中隐含常数仅依赖于 \(f\) 的 Lipschitz 常数。特别地，\(Y_m\) 依分布收敛到两点随机变量
\[
Y_\infty=
\begin{cases}
1,&\text{概率 } \varphi^{-1},\\[3pt]
\varphi^{-1},&\text{概率 } \varphi^{-2}.
\end{cases}
\]
\end{theorem}
\begin{proof}
由定理 \ref{thm:fold-bin-two-state-asymptotic}，存在常数 \(C>0\)，使得对充分大的 \(m\) 与全部 \(w\in X_m\) 都一致有
\[
\left|
\left(\frac{\varphi}{2}\right)^m d_m^{\mathrm{bin}}(w)-\varphi^{-w_m}
\right|
\le C\left(\frac{\varphi}{2}\right)^m.
\]
因此
\[
\bigl|Y_m-\varphi^{-W_m(m)}\bigr|
\le C\left(\frac{\varphi}{2}\right)^m
\]
几乎处处成立。若记 \(L_f\) 为 \(f\) 的 Lipschitz 常数，则
\[
\bigl|f(Y_m)-f(\varphi^{-W_m(m)})\bigr|
\le
L_f C\left(\frac{\varphi}{2}\right)^m.
\]
对均匀分布取期望，得到
\[
\mathbb E[f(Y_m)]
=
\mathbb E[f(\varphi^{-W_m(m)})]
+O\!\left(\left(\frac{\varphi}{2}\right)^m\right).
\]
另一方面，
\[
\mathbb E[f(\varphi^{-W_m(m)})]
=
\frac{F_{m+1}}{F_{m+2}}\,f(1)
+
\frac{F_m}{F_{m+2}}\,f(\varphi^{-1}),
\]
因为
\[
\mathbb P(W_m(m)=0)=\frac{F_{m+1}}{F_{m+2}},
\qquad
\mathbb P(W_m(m)=1)=\frac{F_m}{F_{m+2}}.
\]
这就得到所示定量公式。再用
\[
\frac{F_{m+1}}{F_{m+2}}\to\varphi^{-1},
\qquad
\frac{F_m}{F_{m+2}}\to\varphi^{-2},
\]
即可得到所述两点极限分布。证毕。
\end{proof}

\begin{corollary}[统一缩放矩极限与 \texorpdfstring{$S_q^{\mathrm{bin}}(m)$}{Sqbin(m)} 的首项常数]\label{cor:xi-time-part70a-scaled-moments-leading-constant}
对任意实数 \(q\ge 0\)，都有
\[
\lim_{m\to\infty}
\frac{1}{|X_m|}
\sum_{w\in X_m}
\left(
\left(\frac{\varphi}{2}\right)^m d_m^{\mathrm{bin}}(w)
\right)^q
=
\frac{1}{\varphi}+\varphi^{-q-2}.
\]
等价地，若记
\[
S_q^{\mathrm{bin}}(m):=\sum_{w\in X_m}\bigl(d_m^{\mathrm{bin}}(w)\bigr)^q,
\]
则有精确首项渐近
\[
\boxed{\;
S_q^{\mathrm{bin}}(m)
\sim
\frac{\varphi+\varphi^{-q}}{\sqrt5}
\left(\frac{2^q}{\varphi^{q-1}}\right)^m.\;}
\]
\end{corollary}
\begin{proof}
对固定 \(q\ge 0\)，由上一条定理作用于函数 \(f_q(y):=y^q\) 即可。事实上，由定理 \ref{thm:xi-time-part70a-uniform-scaled-degeneracy-quantitative} 的证明，对充分大的 \(m\) 与几乎处处的样本点都有
\[
Y_m\in\left[\frac{\varphi^{-1}}{2},2\right],
\]
故 \(f_q\) 在该固定紧区间上 Lipschitz。于是
\[
\frac{1}{|X_m|}
\sum_{w\in X_m}
\left(
\left(\frac{\varphi}{2}\right)^m d_m^{\mathrm{bin}}(w)
\right)^q
\to
\frac{1}{\varphi}+\varphi^{-q-2}.
\]
将上式改写为
\[
S_q^{\mathrm{bin}}(m)
=
|X_m|
\left(\frac{2}{\varphi}\right)^{mq}
\left(\frac{1}{\varphi}+\varphi^{-q-2}+o(1)\right),
\]
再用 Binet 公式
\[
|X_m|=F_{m+2}\sim \frac{\varphi^{m+2}}{\sqrt5},
\]
便得到
\[
S_q^{\mathrm{bin}}(m)
\sim
\frac{\varphi^{m+2}}{\sqrt5}
\left(\frac{2}{\varphi}\right)^{mq}
\left(\varphi^{-1}+\varphi^{-q-2}\right)
=
\frac{\varphi+\varphi^{-q}}{\sqrt5}
\left(\frac{2^q}{\varphi^{q-1}}\right)^m.
\]
证毕。
\end{proof}

\begin{theorem}[均匀稳定层平均对数退化度的二态展开]\label{thm:xi-time-part70a-uniform-average-logdeg-two-state}
定义均匀稳定层平均对数退化度
\[
\overline{\ell}_m
:=
\frac{1}{|X_m|}\sum_{w\in X_m}\log d_m^{\mathrm{bin}}(w).
\]
则有
\[
\boxed{\;
\overline{\ell}_m
=
m\log\!\Bigl(\frac{2}{\varphi}\Bigr)
-\frac{F_m}{F_{m+2}}\log\varphi
+O\!\left(\left(\frac{\varphi}{2}\right)^m\right).\;}
\]
因而
\[
\boxed{\;
\overline{\ell}_m
=
m\log\!\Bigl(\frac{2}{\varphi}\Bigr)
-\varphi^{-2}\log\varphi
+O\!\left(\left(\frac{\varphi}{2}\right)^m\right).\;}
\]
\end{theorem}
\begin{proof}
由定理 \ref{thm:fold-bin-two-state-asymptotic} 的一致对数展开，
\[
\log d_m^{\mathrm{bin}}(w)
=
m\log\!\Bigl(\frac{2}{\varphi}\Bigr)
-w_m\log\varphi
+O\!\left(\left(\frac{\varphi}{2}\right)^m\right)
\qquad (w\in X_m),
\]
且误差在 \(w\) 上一致。对全部 \(w\in X_m\) 求平均，得到
\[
\overline{\ell}_m
=
m\log\!\Bigl(\frac{2}{\varphi}\Bigr)
-\frac{1}{|X_m|}\sum_{w\in X_m}w_m\,\log\varphi
+O\!\left(\left(\frac{\varphi}{2}\right)^m\right).
\]
而
\[
\sum_{w\in X_m}w_m=|A_1|=F_m,
\qquad
|X_m|=F_{m+2},
\]
故
\[
\frac{1}{|X_m|}\sum_{w\in X_m}w_m=\frac{F_m}{F_{m+2}}.
\]
代回即得第一式。再用 Binet 公式给出的比值极限
\[
\frac{F_m}{F_{m+2}}=\varphi^{-2}+O(\varphi^{-2m}),
\]
便得到第二式。证毕。
\end{proof}
\leanverified{paper\_xi\_time\_part70a\_uniform\_average\_logdeg\_two\_state}

\leanverified{paper\_xi\_time\_part70a\_gauge\_fiber\_gap\_first\_order}
\begin{theorem}[规范体积与纤维信息差的一阶首项]\label{thm:xi-time-part70a-gauge-fiber-gap-first-order}
定义
\[
\mu_m(w):=\frac{d_m^{\mathrm{bin}}(w)}{2^m},
\qquad
\kappa_m:=\sum_{w\in X_m}\mu_m(w)\log d_m^{\mathrm{bin}}(w),
\qquad
g_m:=2^{-m}\log|G_m^{\mathrm{bin}}|,
\]
并沿用推论 \ref{cor:fold-bin-recover-2pi} 的规范常数
\[
C_m
:=
\exp\!\left(
\frac{2^{m+1}}{|X_m|}\bigl(g_m-\kappa_m+1\bigr)-\overline{\ell}_m
\right).
\]
则有
\[
\boxed{\;
\frac{2^{m+1}}{|X_m|}\bigl(g_m-\kappa_m+1\bigr)
=
m\log\!\Bigl(\frac{2}{\varphi}\Bigr)
-\frac{F_m}{F_{m+2}}\log\varphi
+\log(2\pi)
+O\!\left(\left(\frac{\varphi}{2}\right)^m\right).\;}
\]
因此
\[
\boxed{\;
g_m-\kappa_m+1
=
\frac{\varphi^2}{2\sqrt5}
\Bigl(
m\log\!\Bigl(\frac{2}{\varphi}\Bigr)
-\varphi^{-2}\log\varphi
+\log(2\pi)
\Bigr)
\left(\frac{\varphi}{2}\right)^m
+O\!\left(\left(\frac{\varphi}{2}\right)^{2m}\right).\;}
\]
\end{theorem}
\begin{proof}
由 \(C_m\) 的定义可得精确恒等式
\[
\frac{2^{m+1}}{|X_m|}\bigl(g_m-\kappa_m+1\bigr)=\log C_m+\overline{\ell}_m.
\]
另一方面，推论 \ref{cor:fold-bin-recover-2pi} 给出
\[
\log C_m=\log(2\pi)+O\!\left(\left(\frac{\varphi}{2}\right)^m\right),
\]
而定理 \ref{thm:xi-time-part70a-uniform-average-logdeg-two-state} 已证明
\[
\overline{\ell}_m
=
m\log\!\Bigl(\frac{2}{\varphi}\Bigr)
-\frac{F_m}{F_{m+2}}\log\varphi
+O\!\left(\left(\frac{\varphi}{2}\right)^m\right).
\]
相加便得到第一条公式。

再由 Binet 公式，
\[
\frac{|X_m|}{2^{m+1}}
=
\frac{F_{m+2}}{2^{m+1}}
=
\frac{\varphi^2}{2\sqrt5}\left(\frac{\varphi}{2}\right)^m
+O\!\bigl((2\varphi)^{-m}\bigr),
\]
且
\[
\frac{F_m}{F_{m+2}}=\varphi^{-2}+O(\varphi^{-2m}).
\]
用这两式乘第一条公式，得到
\[
g_m-\kappa_m+1
=
\frac{\varphi^2}{2\sqrt5}
\Bigl(
m\log\!\Bigl(\frac{2}{\varphi}\Bigr)
-\varphi^{-2}\log\varphi
+\log(2\pi)
\Bigr)
\left(\frac{\varphi}{2}\right)^m
+O\!\bigl(m(2\varphi)^{-m}\bigr)
+O\!\left(\left(\frac{\varphi}{2}\right)^{2m}\right).
\]
由于
\[
\frac{1}{2\varphi}<\left(\frac{\varphi}{2}\right)^2,
\]
故
\[
m(2\varphi)^{-m}=O\!\left(\left(\frac{\varphi}{2}\right)^{2m}\right).
\]
于是第二条公式成立。证毕。
\end{proof}

\noindent
由此，bin-fold 的两态一致渐近不但控制稳定层均匀经验谱与全部缩放矩，而且把均匀平均的对数退化度与规范体积修正的首项常数同时固定在同一末位二态账本之上。于是，均匀层二点律、幂和首项、平均对数势与规范体积差额最终被收束为同一条可显式求值的黄金比例主线。


\paragraph{bin-fold 退化度的全实矩闭式与 Bernoulli 塔的逆矩驱动}
在稳定层均匀经验谱的两点极限、全部非负缩放矩的首项常数以及 fold-gauge 常数的 Stirling--Bernoulli 层级已经建立之后，还可把 bin-fold 退化度的实参数矩谱与 Bernoulli 模态系数进一步压缩为下列两条直接接口。

\begin{theorem}[bin-fold 退化度的全实矩闭式]\label{thm:xi-time-part70aa-binfold-all-real-moments-closed}
对每个 \(m\ge 1\)，记
\[
A_0^{(m)}:=\{w\in X_m:\ w_m=0\},
\qquad
A_1^{(m)}:=\{w\in X_m:\ w_m=1\},
\]
并定义稳定层均匀实矩
\[
M_s^{\mathrm{bin}}(m)
:=
\frac{1}{|X_m|}
\sum_{w\in X_m}\bigl(d_m^{\mathrm{bin}}(w)\bigr)^s
\qquad (s\in\RR).
\]
则对任意固定实数 \(s\)，都有渐近公式
\[
\boxed{\;
M_s^{\mathrm{bin}}(m)
=
\left(\frac{2}{\varphi}\right)^{sm}
\left(
\varphi^{-1}+\varphi^{-s-2}
\right)
+
O_s\!\left(
\left(\frac{2}{\varphi}\right)^{sm}
\left(\frac{\varphi}{2}\right)^m
\right).\;}
\]
等价地，
\[
\boxed{\;
\lim_{m\to\infty}
\left(\frac{\varphi}{2}\right)^{sm}
M_s^{\mathrm{bin}}(m)
=
\varphi^{-1}+\varphi^{-s-2}.\;}
\]
\end{theorem}
\begin{proof}
记
\[
\varepsilon_m:=\left(\frac{\varphi}{2}\right)^m.
\]
由定理 \ref{thm:fold-bin-two-state-asymptotic}，存在常数 \(C>0\)，使对充分大的 \(m\) 与全部 \(w\in X_m\) 都有
\[
\left(\frac{\varphi}{2}\right)^m d_m^{\mathrm{bin}}(w)
=
\begin{cases}
1+O(\varepsilon_m), & w\in A_0^{(m)},\\[4pt]
\varphi^{-1}+O(\varepsilon_m), & w\in A_1^{(m)},
\end{cases}
\]
且误差在 \(w\) 上一致。由于两块主值 \(1\) 与 \(\varphi^{-1}\) 都严格远离 \(0\)，对任意固定 \(s\in\RR\)，幂函数 \(x\mapsto x^s\) 在这两个主值附近都可作一致一阶展开，从而
\[
\bigl(d_m^{\mathrm{bin}}(w)\bigr)^s
=
\left(\frac{2}{\varphi}\right)^{sm}
\begin{cases}
1+O_s(\varepsilon_m), & w\in A_0^{(m)},\\[4pt]
\varphi^{-s}+O_s(\varepsilon_m), & w\in A_1^{(m)},
\end{cases}
\]
其中隐含常数仅依赖于 \(s\)。

对两块分别求和并除以 \(|X_m|=F_{m+2}\)，得到
\[
M_s^{\mathrm{bin}}(m)
=
\left(\frac{2}{\varphi}\right)^{sm}
\left(
\frac{F_{m+1}+\varphi^{-s}F_m}{F_{m+2}}
+
O_s(\varepsilon_m)
\right),
\]
因为
\[
|A_0^{(m)}|=F_{m+1},
\qquad
|A_1^{(m)}|=F_m.
\]
再用 Binet 公式给出的比值渐近
\[
\frac{F_{m+1}}{F_{m+2}}=\varphi^{-1}+O\!\left(\varphi^{-2m}\right),
\qquad
\frac{F_m}{F_{m+2}}=\varphi^{-2}+O\!\left(\varphi^{-2m}\right),
\]
并注意 \(\varphi^{-2m}=O(\varepsilon_m)\)，便得到
\[
M_s^{\mathrm{bin}}(m)
=
\left(\frac{2}{\varphi}\right)^{sm}
\left(
\varphi^{-1}+\varphi^{-s-2}
+
O_s(\varepsilon_m)
\right).
\]
这就给出所示渐近式；第二条极限式只是将主项重新整理后的等价写法。证毕。
\end{proof}

\begin{corollary}[Bernoulli 级数系数的逆矩解释]\label{cor:xi-time-part70aa-bernoulli-coefficient-inverse-moment-interpretation}
对每个整数 \(r\ge 1\)，定义奇数阶逆矩
\[
I_{2r-1}(m)
:=
\frac{1}{|X_m|}
\sum_{w\in X_m}\bigl(d_m^{\mathrm{bin}}(w)\bigr)^{-(2r-1)}.
\]
则有
\[
\boxed{\;
I_{2r-1}(m)
=
\left(\frac{\varphi}{2}\right)^{(2r-1)m}
\left(
\varphi^{-1}+\varphi^{2r-3}
\right)
+
O_r\!\left(
\left(\frac{\varphi}{2}\right)^{2rm}
\right).\;}
\]
从而
\[
\boxed{\;
\lim_{m\to\infty}
\left(\frac{2}{\varphi}\right)^{(2r-1)m}
I_{2r-1}(m)
=
\varphi^{-1}+\varphi^{2r-3}.\;}
\]
另一方面，对任意固定 \(R\ge 1\)，定理 \ref{thm:fold-bin-gauge-constant-stirling-bernoulli-hierarchy} 给出
\[
\log C_m
=
\log(2\pi)
+
\sum_{r=1}^{R}
\frac{B_{2r}}{r(2r-1)}
\bigl(\varphi^{-1}+\varphi^{2r-3}\bigr)
\left(\frac{\varphi}{2}\right)^{(2r-1)m}
+
O\!\left(\left(\frac{\varphi}{2}\right)^{(2R+1)m}\right).
\]
因此，第 \(r\) 级 Bernoulli 模态系数恰可写为
\[
\frac{B_{2r}}{r(2r-1)}
\cdot
\lim_{m\to\infty}
\left(\frac{2}{\varphi}\right)^{(2r-1)m}
I_{2r-1}(m).
\]
\end{corollary}
\begin{proof}
在定理 \ref{thm:xi-time-part70aa-binfold-all-real-moments-closed} 中取
\[
s=-(2r-1)
\]
即可得到
\[
I_{2r-1}(m)=M_{-(2r-1)}^{\mathrm{bin}}(m),
\]
并由
\[
\varphi^{-(-2r+1)-2}=\varphi^{2r-3}
\]
得到所示主项常数。误差项则满足
\[
\left(\frac{2}{\varphi}\right)^{-(2r-1)m}
\left(\frac{\varphi}{2}\right)^m
=
\left(\frac{\varphi}{2}\right)^{2rm}.
\]
最后，把该极限常数与定理 \ref{thm:fold-bin-gauge-constant-stirling-bernoulli-hierarchy} 中对应的第 \(r\) 级系数逐项比较，即得最后一式。证毕。
\end{proof}

\noindent
这表明，fold-gauge 常数序列中的 Bernoulli 塔并不是脱离 bin-fold 退化度谱而附加的解析修正；相反，它正由稳定层退化度的奇数阶逆矩常数逐级驱动。


\paragraph{末位二态分层、逆阈值几何与 Bernoulli 塔的 Mellin 编码}
在稳定层均匀经验谱的两点极限、escort 温度族的末位 Bernoulli 塌缩、中心临界容量骨架的双折点闭式以及实参数矩谱与 Bernoulli 模态系数之间的桥接已经建立之后，还可把阈值/采样读出、路径/相位重加权与 Stirling--Bernoulli 层级进一步压缩为同一条末位二态主线。下列结果分别刻画最终严格分层、分位二原子塌缩、归一化逆阈值几何、escort 熵谱的完整常数项，以及容量曲线对 Bernoulli 塔几何权重的 Mellin 编码。

\begin{theorem}[末位相位的最终严格分层]\label{thm:xi-time-part70ab-terminal-bit-eventual-strict-stratification}
沿用记号
\[
d_m(w):=d_m^{\mathrm{bin}}(w),\qquad
A_i^{(m)}:=\{w\in X_m:\ w_m=i\}\quad(i=0,1),
\qquad
A_m:=\left(\frac{2}{\varphi}\right)^m.
\]
则存在整数 \(m_0\)，使得对一切 \(m\ge m_0\)，都有严格分离
\[
\min_{w\in A_0^{(m)}} d_m(w)
>
\max_{v\in A_1^{(m)}} d_m(v).
\]
因此在充分高分辨率下，纤维多重度的全序被末位 \(w_m\) 强制划分为两层。
\end{theorem}
\begin{proof}
由定理 \ref{thm:fold-bin-two-state-asymptotic}，存在绝对常数 \(C>0\)，使对充分大的 \(m\) 与全部 \(x\in X_m\) 都有
\[
\bigl|d_m(x)-\varphi^{-x_m}A_m\bigr|\le C.
\]
若 \(w\in A_0^{(m)}\) 且 \(v\in A_1^{(m)}\)，则
\[
d_m(w)-d_m(v)
\ge
A_m-C-\bigl(\varphi^{-1}A_m+C\bigr)
=
\varphi^{-2}A_m-2C.
\]
由于 \(A_m\to\infty\)，右端对充分大的 \(m\) 为正，结论成立。
\end{proof}

\begin{corollary}[极值纤维的末位定位与最小尺度失配率]\label{cor:xi-time-part70ab-extremal-fibers-terminal-bit-localization}
定义
\[
d_{\min}(m):=\min_{w\in X_m}d_m(w),
\qquad
M_m:=\max_{w\in X_m}d_m(w).
\]
则对充分大的 \(m\)，有
\[
\arg\min_{w\in X_m} d_m(w)\subseteq A_1^{(m)},
\qquad
\arg\max_{w\in X_m} d_m(w)\subseteq A_0^{(m)}.
\]
并且
\[
d_{\min}(m)=\varphi^{-1}A_m+O(1),
\qquad
M_m=A_m+O(1),
\qquad
\frac{M_m}{d_{\min}(m)}\longrightarrow \varphi.
\]
此外，沿偶数分辨率 \(m\to\infty\) 有
\[
\frac{d_{\min}(m)}{F_{m/2}}\longrightarrow 0.
\]
\end{corollary}
\begin{proof}
由定理 \ref{thm:xi-time-part70ab-terminal-bit-eventual-strict-stratification}，对充分大的 \(m\)，最小值只能取自 \(A_1^{(m)}\)，最大值只能取自 \(A_0^{(m)}\)。另一方面，由定理 \ref{thm:fold-bin-two-state-asymptotic} 的一致渐近，
\[
d_m(w)=\varphi^{-1}A_m+O(1)\quad(w\in A_1^{(m)}),
\qquad
d_m(w)=A_m+O(1)\quad(w\in A_0^{(m)}),
\]
且误差在各块上统一，故
\[
d_{\min}(m)=\varphi^{-1}A_m+O(1),
\qquad
M_m=A_m+O(1).
\]
于是 \(M_m/d_{\min}(m)\to \varphi\)。

再沿偶数分辨率比较 \(F_{m/2}\)。由 Binet 公式，
\[
F_{m/2}\sim \frac{\varphi^{m/2}}{\sqrt5},
\qquad
d_{\min}(m)\sim \varphi^{-1}\left(\frac{2}{\varphi}\right)^m,
\]
从而
\[
\frac{d_{\min}(m)}{F_{m/2}}
\sim
\sqrt5\,\varphi^{-1}
\left(\frac{2}{\varphi^{3/2}}\right)^m.
\]
由于 \(\varphi^3=2+\sqrt5>4\)，故 \(\varphi^{3/2}>2\)，于是上式趋于 \(0\)。
\end{proof}

\begin{theorem}[中心投影的分位二原子塌缩]\label{thm:xi-time-part70ab-uniform-quantile-two-atom-collapse}
设 \(U_m\) 服从 \(X_m\) 上的均匀分布，并定义尺度化随机变量
\[
Y_m:=\left(\frac{\varphi}{2}\right)^m d_m(U_m).
\]
再定义左连续分位函数
\[
Q_m(\alpha):=\inf\{t\ge 0:\ \mathbf P(Y_m\le t)\ge \alpha\},
\qquad
0<\alpha<1.
\]
则对每个
\[
\alpha\in(0,1)\setminus\{\varphi^{-2}\},
\]
都有极限
\[
Q_m(\alpha)\longrightarrow Q(\alpha),
\qquad
Q(\alpha)=
\begin{cases}
\varphi^{-1}, & 0<\alpha<\varphi^{-2},\\[2mm]
1, & \varphi^{-2}<\alpha<1.
\end{cases}
\]
并且在临界分位 \(\alpha=\varphi^{-2}\) 处出现奇偶振荡：
\[
Q_{2n}(\varphi^{-2})\longrightarrow 1,
\qquad
Q_{2n+1}(\varphi^{-2})\longrightarrow \varphi^{-1}.
\]
\end{theorem}
\begin{proof}
取
\[
0<\varepsilon<\frac{1-\varphi^{-1}}{3}.
\]
由定理 \ref{thm:fold-bin-two-state-asymptotic}，对充分大的 \(m\) 与全部 \(w\in X_m\) 都有
\[
\left|\left(\frac{\varphi}{2}\right)^m d_m(w)-\varphi^{-w_m}\right|
\le
C\left(\frac{\varphi}{2}\right)^m<\varepsilon.
\]
因此当 \(m\) 充分大时，
\[
Y_m\in[\varphi^{-1}-\varepsilon,\varphi^{-1}+\varepsilon]
\quad\text{在 }A_1^{(m)}\text{ 上},
\]
\[
Y_m\in[1-\varepsilon,1+\varepsilon]
\quad\text{在 }A_0^{(m)}\text{ 上}.
\]
故分位值只由块质量
\[
\mathbf P\bigl(Y_m\le \varphi^{-1}+\varepsilon\bigr)
=
\frac{|A_1^{(m)}|}{|X_m|}
=
\frac{F_m}{F_{m+2}}
\]
决定。若 \(\alpha<\varphi^{-2}\) 或 \(\alpha>\varphi^{-2}\)，由
\[
\frac{F_m}{F_{m+2}}\longrightarrow \varphi^{-2}
\]
知对充分大的 \(m\)，比较关系 \(\alpha\lessgtr F_m/F_{m+2}\) 不再改变，于是 \(Q_m(\alpha)\) 分别落在左原子簇与右原子簇中，令 \(\varepsilon\to 0\) 即得所示极限。

最后，由 Binet 公式可直接计算
\[
\frac{F_m}{F_{m+2}}-\varphi^{-2}
=
\frac{(-1)^{m+1}(1-\varphi^{-4})\varphi^{-m}}{\sqrt5\,F_{m+2}},
\]
其符号恰为 \((-1)^{m+1}\)。因此当 \(m\) 为奇数时，
\[
\frac{F_m}{F_{m+2}}>\varphi^{-2},
\]
故 \(Q_m(\varphi^{-2})\) 落在左原子簇；当 \(m\) 为偶数时，
\[
\frac{F_m}{F_{m+2}}<\varphi^{-2},
\]
故 \(Q_m(\varphi^{-2})\) 落在右原子簇。令 \(\varepsilon\to 0\) 即得两条子序列极限。
\end{proof}

\begin{theorem}[escort 温度族的分位相图]\label{thm:xi-time-part70ab-escort-quantile-phase-diagram}
对任意 \(q\ge 0\)，令
\[
W_m^{(q)}\sim \pi_{m,q},
\qquad
Y_m^{(q)}:=\left(\frac{\varphi}{2}\right)^m d_m\bigl(W_m^{(q)}\bigr),
\qquad
\theta(q):=\frac{1}{1+\varphi^{q+1}}.
\]
并定义左连续分位函数
\[
Q_{m,q}(\alpha):=\inf\{t\ge 0:\ \mathbf P(Y_m^{(q)}\le t)\ge \alpha\},
\qquad
0<\alpha<1.
\]
则对每个固定 \(q\ge 0\) 与每个
\[
\alpha\in(0,1)\setminus\{\theta(q)\},
\]
都有极限
\[
Q_{m,q}(\alpha)\longrightarrow Q_q(\alpha),
\qquad
Q_q(\alpha)=
\begin{cases}
\varphi^{-1}, & 0<\alpha<\theta(q),\\[2mm]
1, & \theta(q)<\alpha<1.
\end{cases}
\]
\end{theorem}
\begin{proof}
同样取
\[
0<\varepsilon<\frac{1-\varphi^{-1}}{3}.
\]
由定理 \ref{thm:fold-bin-two-state-asymptotic}，对充分大的 \(m\) 与全部 \(w\in X_m\) 都有
\[
\left|\left(\frac{\varphi}{2}\right)^m d_m(w)-\varphi^{-w_m}\right|<\varepsilon.
\]
因此当 \(m\) 充分大时，
\[
Y_m^{(q)}\in[\varphi^{-1}-\varepsilon,\varphi^{-1}+\varepsilon]
\quad\text{在事件 }W_m^{(q)}(m)=1\text{ 上},
\]
\[
Y_m^{(q)}\in[1-\varepsilon,1+\varepsilon]
\quad\text{在事件 }W_m^{(q)}(m)=0\text{ 上}.
\]
另一方面，由命题 \ref{prop:fold-bin-escort-lastbit}，
\[
\mathbf P\bigl(W_m^{(q)}(m)=1\bigr)=\theta(q)+O\!\left(\left(\frac{\varphi}{2}\right)^m\right).
\]
于是当 \(\alpha<\theta(q)\) 时，充分大的 \(m\) 上有 \(\alpha<\mathbf P(W_m^{(q)}(m)=1)\)，故 \(Q_{m,q}(\alpha)\) 落在左原子簇；当 \(\alpha>\theta(q)\) 时，则 \(Q_{m,q}(\alpha)\) 落在右原子簇。令 \(\varepsilon\to 0\) 即得结论。
\end{proof}

\begin{theorem}[中心容量逆问题的闭式解]\label{thm:xi-time-part70ab-center-capacity-inverse-threshold-closed}
定义归一化连续容量
\[
\Psi_m(\tau):=
2^{-m}\,\mathcal C_m^{\mathrm{bin,cont}}(\tau A_m),
\qquad
\tau\ge 0,
\]
并定义达到水平 \(\rho\) 的最小归一化阈值
\[
\tau_\ast(\rho):=
\inf\left\{\tau\ge 0:\ \liminf_{m\to\infty}\Psi_m(\tau)\ge \rho\right\},
\qquad
0\le \rho\le 1.
\]
则有严格闭式
\[
\tau_\ast(\rho)=
\begin{cases}
\displaystyle \frac{\sqrt5}{\varphi^2}\,\rho,
& 0\le \rho\le \dfrac{\varphi}{\sqrt5},\\[3mm]
\displaystyle \frac{\sqrt5\,\rho-\varphi^{-1}}{\varphi},
& \dfrac{\varphi}{\sqrt5}\le \rho\le 1.
\end{cases}
\]
\end{theorem}
\begin{proof}
由定理 \ref{thm:xi-time-part70b-center-capacity-two-kink-skeleton}，
\[
\Psi_m(\tau)\longrightarrow \Psi(\tau):=
\frac{\varphi}{\sqrt5}\min\{1,\tau\}
+
\frac{1}{\sqrt5}\min\{\varphi^{-1},\tau\}.
\]
因此
\[
\tau_\ast(\rho)=\inf\{\tau\ge 0:\ \Psi(\tau)\ge \rho\}.
\]
而
\[
\Psi(\tau)=
\begin{cases}
\dfrac{\varphi^2}{\sqrt5}\tau, & 0\le \tau\le \varphi^{-1},\\[6pt]
\dfrac{\varphi}{\sqrt5}\tau+\dfrac{1}{\varphi\sqrt5}, & \varphi^{-1}\le \tau\le 1,\\[8pt]
1, & \tau\ge 1,
\end{cases}
\]
故广义逆由两段线性方程直接求得，即为所示公式。
\end{proof}

\begin{theorem}[escort 响应的逆温度阈值公式]\label{thm:xi-time-part70ab-escort-response-inverse-threshold-closed}
对固定 \(q\ge 0\)，定义 escort 归一化响应
\[
\Psi_{m,q}^{\mathrm{esc}}(\tau):=
\mathbb E_{\pi_{m,q}}
\left[
\min\!\left(1,\frac{\tau A_m}{d_m(W_m^{(q)})}\right)
\right],
\qquad
\tau\ge 0.
\]
再定义达到水平 \(\rho\) 的最小归一化阈值
\[
\tau_{q,\ast}(\rho):=
\inf\left\{\tau\ge 0:\ \liminf_{m\to\infty}\Psi_{m,q}^{\mathrm{esc}}(\tau)\ge \rho\right\},
\qquad
0\le \rho\le 1.
\]
记
\[
\rho_c(q):=\theta(q)+\varphi^{-1}\bigl(1-\theta(q)\bigr).
\]
则有闭式
\[
\tau_{q,\ast}(\rho)=
\begin{cases}
\displaystyle \frac{\rho}{1+(\varphi-1)\theta(q)},
& 0\le \rho\le \rho_c(q),\\[3mm]
\displaystyle \frac{\rho-\theta(q)}{1-\theta(q)},
& \rho_c(q)\le \rho\le 1.
\end{cases}
\]
\end{theorem}
\begin{proof}
由定理 \ref{thm:xi-time-part70b-escort-two-kink-response}，
\[
\Psi_{m,q}^{\mathrm{esc}}(\tau)\longrightarrow
\Psi_q^{\mathrm{esc}}(\tau):=
\bigl(1-\theta(q)\bigr)\min\{1,\tau\}
+
\theta(q)\min\{1,\varphi\tau\}.
\]
因此
\[
\tau_{q,\ast}(\rho)=\inf\{\tau\ge 0:\ \Psi_q^{\mathrm{esc}}(\tau)\ge \rho\}.
\]
把 \(\Psi_q^{\mathrm{esc}}\) 在 \([0,\varphi^{-1}]\) 与 \([\varphi^{-1},1]\) 上分段展开，得到
\[
\Psi_q^{\mathrm{esc}}(\tau)=
\begin{cases}
\bigl(1+(\varphi-1)\theta(q)\bigr)\tau,
& 0\le \tau\le \varphi^{-1},\\[6pt]
\theta(q)+\bigl(1-\theta(q)\bigr)\tau,
& \varphi^{-1}\le \tau\le 1,\\[6pt]
1, & \tau\ge 1.
\end{cases}
\]
其连接点函数值为
\[
\Psi_q^{\mathrm{esc}}(\varphi^{-1})
=
\theta(q)+\varphi^{-1}\bigl(1-\theta(q)\bigr)
=
\rho_c(q).
\]
故广义逆由上述两段线性函数分别求反即可。
\end{proof}

\begin{theorem}[escort 熵谱的完整常数项]\label{thm:xi-time-part70ab-escort-renyi-entropy-full-constant}
对固定 \(q\ge 0\)、\(a>0\) 且 \(a\neq 1\)，定义 escort Rényi 熵
\[
H_a(\pi_{m,q})
:=
\frac{1}{1-a}\log\sum_{w\in X_m}\pi_{m,q}(w)^a.
\]
则有渐近展开
\[
H_a(\pi_{m,q})
=
m\log\varphi
-\frac12\log 5
+
\frac{1}{1-a}
\log\!\Bigl(
\varphi^{1-a}\bigl(1-\theta(q)\bigr)^a+\theta(q)^a
\Bigr)
+o(1).
\]
当 \(a\to 1\) 时，得到 Shannon 熵常数项
\[
H(\pi_{m,q})
=
m\log\varphi
-\frac12\log 5
+
\bigl(1-\theta(q)\bigr)\log\varphi
+
h\bigl(\theta(q)\bigr)
+o(1),
\]
其中
\[
h(t):=-t\log t-(1-t)\log(1-t).
\]
\end{theorem}
\begin{proof}
记
\[
S_t(m):=\sum_{w\in X_m} d_m(w)^t.
\]
由定理 \ref{thm:xi-time-part70aa-binfold-all-real-moments-closed} 与 Binet 公式，对每个固定实数 \(t\) 都有
\[
S_t(m)
=
\frac{\varphi+\varphi^{-t}}{\sqrt5}
\bigl(2^t\varphi^{1-t}\bigr)^m
\bigl(1+o(1)\bigr).
\]
因此
\[
\sum_{w\in X_m}\pi_{m,q}(w)^a
=
\frac{S_{aq}(m)}{S_q(m)^a}
=
5^{(a-1)/2}\,
\varphi^{(1-a)m}\,
\frac{\varphi+\varphi^{-aq}}{(\varphi+\varphi^{-q})^a}
\bigl(1+o(1)\bigr).
\]
取对数并除以 \(1-a\)，得到
\[
H_a(\pi_{m,q})
=
m\log\varphi
-\frac12\log 5
+
\frac{1}{1-a}
\log\!\frac{\varphi+\varphi^{-aq}}{(\varphi+\varphi^{-q})^a}
+o(1).
\]
再由
\[
\theta(q)=\frac{\varphi^{-q}}{\varphi+\varphi^{-q}},
\qquad
1-\theta(q)=\frac{\varphi}{\varphi+\varphi^{-q}},
\]
可将最后一项重写为
\[
\frac{\varphi+\varphi^{-aq}}{(\varphi+\varphi^{-q})^a}
=
\varphi^{1-a}\bigl(1-\theta(q)\bigr)^a+\theta(q)^a.
\]
这就给出第一条公式。

当 \(a\to 1\) 时，对
\[
F_q(a):=
\log\!\Bigl(
\varphi^{1-a}\bigl(1-\theta(q)\bigr)^a+\theta(q)^a
\Bigr)
\]
应用 l'Hospital 法则，得到
\[
\lim_{a\to 1}\frac{F_q(a)}{1-a}
=
-F_q'(1)
=
\bigl(1-\theta(q)\bigr)\log\varphi+h\bigl(\theta(q)\bigr).
\]
代回即得 Shannon 常数项。
\end{proof}

\begin{theorem}[容量极限曲线的负矩 Mellin 反演]\label{thm:xi-time-part70ab-capacity-limit-negative-moment-mellin}
定义中心容量极限
\[
m_\ast(\tau):=
\lim_{m\to\infty}
\frac{\sqrt5}{\varphi^2}\,
2^{-m}\,\mathcal C_m^{\mathrm{bin,cont}}(\tau A_m),
\qquad
\tau\ge 0.
\]
则对任意实数 \(s>0\)，成立严格恒等式
\[
s(s+1)\int_0^\infty \tau^{-s-2}\bigl(\tau-m_\ast(\tau)\bigr)\,\dd\tau
=
\varphi^{-1}+\varphi^{s-2}.
\]
\end{theorem}
\begin{proof}
由定理 \ref{thm:xi-time-part70b-center-capacity-two-kink-skeleton}，
\[
m_\ast(\tau)=
\varphi^{-1}\min\{1,\tau\}
+
\varphi^{-2}\min\{\varphi^{-1},\tau\}.
\]
因而若定义双原子概率测度
\[
\mu_\ast:=\varphi^{-2}\delta_{\varphi^{-1}}+\varphi^{-1}\delta_1,
\]
则有
\[
m_\ast(\tau)=\int \min\{y,\tau\}\,\dd\mu_\ast(y).
\]
对任意固定 \(y>0\) 与 \(s>0\)，直接计算得
\[
\int_0^\infty \tau^{-s-2}\bigl(\tau-\min\{y,\tau\}\bigr)\,\dd\tau
=
\int_y^\infty \tau^{-s-2}(\tau-y)\,\dd\tau
=
\frac{y^{-s}}{s(s+1)}.
\]
两边乘以 \(s(s+1)\) 并对 \(\mu_\ast\) 积分，即得
\[
s(s+1)\int_0^\infty \tau^{-s-2}\bigl(\tau-m_\ast(\tau)\bigr)\,\dd\tau
=
\int y^{-s}\,\dd\mu_\ast(y)
=
\varphi^{-1}+\varphi^{-2}\varphi^s
=
\varphi^{-1}+\varphi^{s-2}.
\]
证毕。
\end{proof}

\begin{corollary}[Stirling--Bernoulli 几何系数的容量解释]\label{cor:xi-time-part70ab-bernoulli-hierarchy-capacity-interpretation}
对任意整数 \(r\ge 1\)，都有
\[
\varphi^{-1}+\varphi^{2r-3}
=
(2r-1)(2r)\int_0^\infty \tau^{-2r-1}\bigl(\tau-m_\ast(\tau)\bigr)\,\dd\tau.
\]
从而 Stirling--Bernoulli 层级中第 \(r\) 级系数可写为
\[
\frac{B_{2r}}{r(2r-1)}
\bigl(\varphi^{-1}+\varphi^{2r-3}\bigr)
=
\frac{2B_{2r}}{r}
\int_0^\infty \tau^{-2r-1}\bigl(\tau-m_\ast(\tau)\bigr)\,\dd\tau.
\]
特别地，对任意固定 \(R\ge 1\)，定理 \ref{thm:fold-bin-gauge-constant-stirling-bernoulli-hierarchy} 的展开可重写为
\[
\log C_m
=
\log(2\pi)
+
\sum_{r=1}^{R}
\frac{2B_{2r}}{r}
\left(
\int_0^\infty \tau^{-2r-1}\bigl(\tau-m_\ast(\tau)\bigr)\,\dd\tau
\right)
\left(\frac{\varphi}{2}\right)^{(2r-1)m}
+
O\!\left(\left(\frac{\varphi}{2}\right)^{(2R+1)m}\right).
\]
\end{corollary}
\begin{proof}
在定理 \ref{thm:xi-time-part70ab-capacity-limit-negative-moment-mellin} 中取 \(s=2r-1\) 即得第一式。再把该恒等式代入定理 \ref{thm:fold-bin-gauge-constant-stirling-bernoulli-hierarchy} 的系数表达，便得到后两式。
\end{proof}

\noindent
由此，末位 \(0/1\) 的二态骨架不只控制纤维多重度的最终排序与均匀/escort 采样下的分位相图，还把中心容量曲线的广义逆、escort 温度响应的广义逆以及 \(\log C_m\) 的 Stirling--Bernoulli 几何权重全部还原为同一组双原子数据。于是，阈值/采样轴、路径/相位轴与 Bernoulli 塔的解析层级便在同一条末位二态主线上闭合。

\paragraph{全息地址前缀商的精确熵律与最小判别比特}
在 H.160--H.161 的 faithful \(2^L\)-进地址表示、精确柱分离与圆柱--球识别已经固定之后，有限深度地址商的可见信息量还可进一步压缩为一条完全闭式的计数律。

\begin{theorem}[\texorpdfstring{$2^L$}{2^L}-进全息地址前缀商的精确计数与最小判别比特]\label{thm:xi-time-part70ac-hologram-prefix-quotient-entropy-bits}
设 \(E\) 为有限事件集，
\[
\mathrm{code}:E\hookrightarrow \NN,
\qquad
M:=\max_{e\in E}\mathrm{code}(e),
\qquad
B:=2^L>M,
\]
并记
\[
T:=T_2(E_{\min}),
\qquad
v\in T\setminus 2T,
\]
其中
\[
X(\mathbf e):=\sum_{t\ge 1}\mathrm{code}(e_t)\,B^{t-1}v
\qquad
\bigl(\mathbf e=(e_t)_{t\ge 1}\in E^{\NN}\bigr).
\]
并定义
\[
\mathcal X_{E,v}:=X(E^{\NN}).
\]
对每个 \(n\ge 1\)，定义第 \(n\) 级地址商集
\[
\mathcal A_n
:=
\left\{
X(\mathbf e)\bmod B^nT:\ \mathbf e\in E^{\NN}
\right\}
\subseteq T/B^nT.
\]
再对每个长度 \(n\) 前缀
\[
\mathbf a=(a_1,\dots,a_n)\in E^n
\]
记
\[
[\mathbf a]
:=
\left\{
\mathbf e\in E^{\NN}:\ \mathrm{pref}_n(\mathbf e)=\mathbf a
\right\},
\qquad
x_{\mathbf a}
:=
\sum_{t=1}^{n}\mathrm{code}(a_t)\,B^{t-1}v.
\]
则成立：
\begin{enumerate}
  \item 地址商集满足精确计数公式
  \[
  |\mathcal A_n|=|E|^n.
  \]
  \item 若定义深度 \(n\) 的 clopen 地址柱代数
  \[
  \mathcal B_n
  :=
  \left\{
  \bigcup_{\mathbf a\in U}X([\mathbf a]):\ U\subseteq E^n
  \right\},
  \]
  则
  \[
  |\mathcal B_n|=2^{|E|^n}.
  \]
  \item 若 \(A\) 为有限集，且读出映射
  \[
  R:\mathcal X_{E,v}\longrightarrow A
  \]
  满足对任意不同前缀 \(\mathbf a\neq \mathbf b\in E^n\) 及任意
  \[
  x\in X([\mathbf a]),
  \qquad
  y\in X([\mathbf b]),
  \]
  都有 \(R(x)\neq R(y)\)，则
  \[
  |A|\ge |E|^n.
  \]
  因而区分全部长度 \(n\) 前缀所需的最小判别比特数恰为
  \[
  \left\lceil n\log_2|E|\right\rceil.
  \]
\end{enumerate}
\end{theorem}
\begin{proof}
固定 \(n\ge 1\)。定义映射
\[
\Phi_n:E^n\longrightarrow \mathcal A_n,
\qquad
\mathbf a\longmapsto x_{\mathbf a}+B^nT.
\]
先证满射。任取 \(\mathbf e=(e_t)_{t\ge 1}\in E^{\NN}\)，并记
\[
\mathbf a:=\mathrm{pref}_n(\mathbf e)=(e_1,\dots,e_n).
\]
则
\[
X(\mathbf e)-x_{\mathbf a}
=
\sum_{t>n}\mathrm{code}(e_t)\,B^{t-1}v
\in B^nT,
\]
故
\[
X(\mathbf e)\bmod B^nT=\Phi_n(\mathbf a).
\]
因此 \(\Phi_n\) 满射。

再证单射。若
\[
\Phi_n(\mathbf a)=\Phi_n(\mathbf b),
\]
则
\[
x_{\mathbf a}-x_{\mathbf b}\in B^nT.
\]
任取
\[
\mathbf e\in[\mathbf a],
\qquad
\mathbf f\in[\mathbf b].
\]
由结论 \ref{con:xi-time-part62dgb-hologram-cylinder-ball-isometry}，
\[
X(\mathbf e)\equiv x_{\mathbf a}\pmod{B^nT},
\qquad
X(\mathbf f)\equiv x_{\mathbf b}\pmod{B^nT},
\]
从而
\[
X(\mathbf e)\equiv X(\mathbf f)\pmod{B^nT}.
\]
再由同一结论中的同余--前缀等价式，只能有
\[
\mathrm{pref}_n(\mathbf e)=\mathrm{pref}_n(\mathbf f),
\]
亦即 \(\mathbf a=\mathbf b\)。故 \(\Phi_n\) 为双射，于是
\[
|\mathcal A_n|=|E^n|=|E|^n,
\]
第 \(1\) 条得证。

对第 \(2\) 条，由定理 \ref{thm:xi-time-part62d-hologram-image-cantor-homeomorphism}，映射
\[
X:E^{\NN}\longrightarrow \mathcal X_{E,v}
\]
是拓扑同胚。长度 \(n\) 的 cylinder 集族
\[
\{[\mathbf a]:\ \mathbf a\in E^n\}
\]
在 \(E^{\NN}\) 中构成一个由 \(|E|^n\) 个原子组成的 clopen 分割，故其像
\[
\{X([\mathbf a]):\ \mathbf a\in E^n\}
\]
在 \(\mathcal X_{E,v}\) 中同样构成一个由 \(|E|^n\) 个原子组成的 clopen 分割。于是由这些原子生成的 Boolean 代数正是全部原子并的集合族，故
\[
|\mathcal B_n|=2^{|E|^n}.
\]

再证第 \(3\) 条。对每个 \(\mathbf a\in E^n\)，取一点
\[
x_{\mathbf a}^{\ast}\in X([\mathbf a]),
\]
这是可能的，因为每个 cylinder 都非空。若 \(\mathbf a\neq \mathbf b\)，则按假设
\[
R(x_{\mathbf a}^{\ast})\neq R(x_{\mathbf b}^{\ast}).
\]
因此映射
\[
E^n\longrightarrow A,
\qquad
\mathbf a\longmapsto R(x_{\mathbf a}^{\ast})
\]
是单射，故
\[
|A|\ge |E|^n.
\]
这给出输出值个数的下界。

反向上，由定理 \ref{thm:xi-time-part62d-hologram-image-cantor-homeomorphism}，\(X\) 为双射，故前缀读出
\[
R_n:\mathcal X_{E,v}\longrightarrow E^n,
\qquad
R_n(x):=\mathrm{pref}_n(X^{-1}(x))
\]
良定义，并且恰取 \(|E|^n\) 个值。于是最小所需输出值个数正好等于 \(|E|^n\)。将有限输出值编码为二进制字，最小判别比特数遂为
\[
\left\lceil \log_2|E|^n\right\rceil
=
\left\lceil n\log_2|E|\right\rceil.
\]
证毕。
\end{proof}

\noindent
因此，深度 \(n\) 的地址可见层并非只在数量级上与 \(|E|^n\) 同阶，而是恰由 \(|E|^n\) 个前缀原子严格组成；相应有限值读出的最小复杂度也被精确锁定为 \(\left\lceil n\log_2|E|\right\rceil\) 比特。于是，时间推进所带来的地址精度提升被压缩为一条完全刚性的有限计数律。

\paragraph{bin-fold 输出分布的全阶信息几何常数层闭式}
在输出分布的 Rényi 熵率塌缩、隐藏对数多重度的二原子 Laplace 极限、规范体积密度的两尺度展开以及碰撞矩的两态主项都已建立之后，还可把 bin-fold 推前分布 \(\mu_m\) 的常数层信息几何进一步压缩为下列四条闭式结论。

\begin{theorem}[输出分布的全阶 Rényi 熵具有精确常数漂移]\label{thm:xi-time-part70ad-binfold-renyi-entropy-constant-drift}
\leanverified{paper\_xi\_time\_part70ad\_binfold\_renyi\_entropy\_constant\_drift}
记
\[
\mu_m(w):=\frac{d_m^{\mathrm{bin}}(w)}{2^m}
\qquad (w\in X_m),
\]
并对 \(q>0\)、\(q\neq 1\) 定义
\[
H_q(\mu_m):=\frac{1}{1-q}\log\sum_{w\in X_m}\mu_m(w)^q.
\]
则对每个固定 \(q>0\)、\(q\neq 1\)，都有精确渐近
\[
\boxed{\;
H_q(\mu_m)
=
m\log\varphi
+
\frac{1}{1-q}\log\!\Bigl(\frac{\varphi+\varphi^{-q}}{\sqrt5}\Bigr)
+
O_q\!\Bigl(\Bigl(\frac{\varphi}{2}\Bigr)^m\Bigr).\;}
\]
特别地，
\[
\boxed{\;
\lim_{m\to\infty}\frac{1}{m}H_q(\mu_m)=\log\varphi
\qquad (q>0,\ q\neq 1).\;}
\]
\end{theorem}
\begin{proof}
记
\[
A_0^{(m)}:=\{w\in X_m:\ w_m=0\},
\qquad
A_1^{(m)}:=\{w\in X_m:\ w_m=1\}.
\]
由定理 \ref{thm:fold-bin-two-state-asymptotic} 的一致主项，对每个固定 \(q>0\) 都存在常数 \(C_q>0\)，使得对充分大的 \(m\) 与全部 \(w\in X_m\) 都有
\[
\bigl(d_m^{\mathrm{bin}}(w)\bigr)^q
=
\left(\frac{2}{\varphi}\right)^{qm}
\begin{cases}
1+O_q\!\bigl((\varphi/2)^m\bigr), & w\in A_0^{(m)},\\[4pt]
\varphi^{-q}+O_q\!\bigl((\varphi/2)^m\bigr), & w\in A_1^{(m)}.
\end{cases}
\]
故
\[
S_q^{\mathrm{bin}}(m):=\sum_{w\in X_m}\bigl(d_m^{\mathrm{bin}}(w)\bigr)^q
=
\left(\frac{2}{\varphi}\right)^{qm}
\Bigl(
F_{m+1}+\varphi^{-q}F_m
\Bigr)
\Bigl(1+O_q\!\bigl((\varphi/2)^m\bigr)\Bigr),
\]
因为 \(|A_0^{(m)}|=F_{m+1}\)、\(|A_1^{(m)}|=F_m\)。于是
\[
\sum_{w\in X_m}\mu_m(w)^q
=
2^{-qm}S_q^{\mathrm{bin}}(m)
=
\varphi^{-qm}
\Bigl(
F_{m+1}+\varphi^{-q}F_m
\Bigr)
\Bigl(1+O_q\!\bigl((\varphi/2)^m\bigr)\Bigr).
\]
再由 Binet 公式，
\[
F_{m+1}+\varphi^{-q}F_m
=
\frac{\varphi^m}{\sqrt5}\bigl(\varphi+\varphi^{-q}\bigr)+O(\varphi^{-m}),
\]
而 \(O(\varphi^{-m})=O((\varphi/2)^m\varphi^m)\)，故
\[
\sum_{w\in X_m}\mu_m(w)^q
=
\varphi^{(1-q)m}
\frac{\varphi+\varphi^{-q}}{\sqrt5}
\Bigl(1+O_q\!\bigl((\varphi/2)^m\bigr)\Bigr).
\]
右端主因子严格为正，取对数并除以 \(1-q\) 即得
\[
H_q(\mu_m)
=
m\log\varphi
+
\frac{1}{1-q}\log\!\Bigl(\frac{\varphi+\varphi^{-q}}{\sqrt5}\Bigr)
+
O_q\!\Bigl(\Bigl(\frac{\varphi}{2}\Bigr)^m\Bigr).
\]
再除以 \(m\) 并令 \(m\to\infty\)，便得到熵率极限。证毕。
\end{proof}

\leanverified{paper\_xi\_time\_part70ad\_hidden\_logmultiplicity\_exponential\_moments}
\begin{theorem}[中心化对数多重度的全指数矩极限与全部累积量闭式]\label{thm:xi-time-part70ad-hidden-logmultiplicity-exponential-moments}
定义
\[
Y_m(w):=\log d_m^{\mathrm{bin}}(w)-m\log\!\Bigl(\frac{2}{\varphi}\Bigr),
\qquad
W_m\sim \mu_m.
\]
则对任意固定复数 \(t\in\CC\)，都有
\[
\boxed{\;
\EE_{\mu_m}\!\bigl[e^{tY_m(W_m)}\bigr]
=
\frac{\varphi+\varphi^{-(t+1)}}{\sqrt5}
+
O_t\!\Bigl(\Bigl(\frac{\varphi}{2}\Bigr)^m\Bigr).\;}
\]
因此 \(Y_m(W_m)\) 在分布上并且在全部指数矩意义下收敛到两点随机变量
\[
Y_\infty=
\begin{cases}
0,& \text{概率 } \dfrac{\varphi}{\sqrt5},\\[6pt]
-\log\varphi,& \text{概率 } \dfrac{1}{\varphi\sqrt5}.
\end{cases}
\]
若记
\[
p:=\frac{1}{\varphi\sqrt5},
\qquad
\ell:=\log\varphi,
\]
则其极限累积量生成函数为
\[
K_\infty(t)=\log\bigl((1-p)+pe^{-t\ell}\bigr),
\]
并且前四阶极限累积量满足
\[
\boxed{\;
\kappa_1^\infty=-p\ell,\qquad
\kappa_2^\infty=p(1-p)\ell^2,\qquad
\kappa_3^\infty=-p(1-p)(1-2p)\ell^3,\qquad
\kappa_4^\infty=p(1-p)(1-6p+6p^2)\ell^4.\;}
\]
\end{theorem}
\begin{proof}
记 \(B_m:=W_m(m)\in\{0,1\}\)。由定理 \ref{thm:fold-bin-two-state-asymptotic} 的一致对数展开，
\[
Y_m(w)=-w_m\log\varphi+O\!\Bigl(\Bigl(\frac{\varphi}{2}\Bigr)^m\Bigr)
\qquad (w\in X_m),
\]
于是对任意固定 \(t\in\CC\)，都有一致展开
\[
e^{tY_m(w)}
=
\varphi^{-t w_m}
\Bigl(1+O_t\!\Bigl(\Bigl(\frac{\varphi}{2}\Bigr)^m\Bigr)\Bigr).
\]
对 \(\mu_m\) 取期望，并利用命题 \ref{prop:fold-bin-escort-lastbit} 在 \(q=1\) 处的末位质量公式
\[
\mu_m(B_m=0)=\frac{\varphi}{\sqrt5}+O\!\Bigl(\Bigl(\frac{\varphi}{2}\Bigr)^m\Bigr),
\qquad
\mu_m(B_m=1)=\frac{1}{\varphi\sqrt5}+O\!\Bigl(\Bigl(\frac{\varphi}{2}\Bigr)^m\Bigr),
\]
便得到
\[
\EE_{\mu_m}\!\bigl[e^{tY_m(W_m)}\bigr]
=
\frac{\varphi}{\sqrt5}
+
\frac{1}{\varphi\sqrt5}\varphi^{-t}
+
O_t\!\Bigl(\Bigl(\frac{\varphi}{2}\Bigr)^m\Bigr)
=
\frac{\varphi+\varphi^{-(t+1)}}{\sqrt5}
+
O_t\!\Bigl(\Bigl(\frac{\varphi}{2}\Bigr)^m\Bigr).
\]
这就得到指数矩极限。

由于 \(Y_m\) 一致有界，上式同时给出所有指数矩收敛，并推出 \(Y_m(W_m)\Rightarrow Y_\infty\)，其中
\[
\PP(Y_\infty=0)=1-p=\frac{\varphi}{\sqrt5},
\qquad
\PP(Y_\infty=-\ell)=p=\frac{1}{\varphi\sqrt5}.
\]
故极限累积量生成函数即为
\[
K_\infty(t)=\log\EE[e^{tY_\infty}]
=
\log\bigl((1-p)+pe^{-t\ell}\bigr).
\]
对 \(t=0\) 逐阶求导，便得到
\[
\kappa_1^\infty=-p\ell,
\qquad
\kappa_2^\infty=p(1-p)\ell^2,
\]
\[
\kappa_3^\infty=-p(1-p)(1-2p)\ell^3,
\qquad
\kappa_4^\infty=p(1-p)(1-6p+6p^2)\ell^4.
\]
证毕。
\end{proof}

\begin{theorem}[规范群熵的精确首项与常数漂移]\label{thm:xi-time-part70ad-gauge-entropy-leading-drift}
\leanverified{paper\_xi\_time\_part70ad\_gauge\_entropy\_leading\_drift}
记
\[
G_m^{\mathrm{bin}}
\cong
\prod_{w\in X_m}S_{d_m^{\mathrm{bin}}(w)},
\qquad
H_m^{\mathrm{gauge}}:=\log_2|G_m^{\mathrm{bin}}|.
\]
则其自然对数版本满足
\[
\boxed{\;
\frac{1}{2^m}\log |G_m^{\mathrm{bin}}|
=
m\log\!\Bigl(\frac{2}{\varphi}\Bigr)
-1
-\frac{\log\varphi}{\varphi\sqrt5}
+o(1).\;}
\]
等价地，以 bits 计，
\[
\boxed{\;
H_m^{\mathrm{gauge}}
=
2^m\left[
m\log_2\!\Bigl(\frac{2}{\varphi}\Bigr)
-\frac{1}{\log 2}
\Bigl(
1+\frac{\log\varphi}{\varphi\sqrt5}
\Bigr)
+o(1)
\right].\;}
\]
\end{theorem}
\begin{proof}
由定理 \ref{thm:xi-fold-hidden-logmultiplicity-kappa-gauge-two-scale-explicit}，若记
\[
g_m:=2^{-m}\log|G_m^{\mathrm{bin}}|,
\]
则有
\[
g_m
=
m\log\!\Bigl(\frac{2}{\varphi}\Bigr)
-\Bigl(1+\frac{\log\varphi}{1+\varphi^2}\Bigr)
+O\!\Bigl(\Bigl(\frac{\varphi}{2}\Bigr)^m\Bigr).
\]
再用恒等式
\[
\frac{1}{1+\varphi^2}=\frac{1}{\varphi\sqrt5},
\]
即得
\[
\frac{1}{2^m}\log |G_m^{\mathrm{bin}}|
=
m\log\!\Bigl(\frac{2}{\varphi}\Bigr)
-1
-\frac{\log\varphi}{\varphi\sqrt5}
+o(1).
\]
最后两边乘以 \(2^m/\log 2\)，便得到 bits 形式。证毕。
\end{proof}

\begin{theorem}[bin-fold 输出分布的 Tsallis 熵同样具有闭式常数层]\label{thm:xi-time-part70ad-binfold-tsallis-entropy-asymptotics}
对 \(q>0\)、\(q\neq 1\)，定义 Tsallis 熵
\[
T_q(\mu_m):=\frac{1-\sum_{w\in X_m}\mu_m(w)^q}{q-1}.
\]
则有如下两段渐近：
\begin{enumerate}
  \item 若 \(q>1\)，则
  \[
  \boxed{\;
  T_q(\mu_m)
  =
  \frac{1}{q-1}
  -
  \frac{1}{q-1}\,
  \varphi^{(1-q)m}
  \frac{\varphi+\varphi^{-q}}{\sqrt5}
  +
  O_q\!\Bigl(\varphi^{(1-q)m}\Bigl(\frac{\varphi}{2}\Bigr)^m\Bigr).\;}
  \]
  特别地，
  \[
  \boxed{\;
  T_q(\mu_m)\longrightarrow \frac{1}{q-1}.\;}
  \]
  \item 若 \(0<q<1\)，则
  \[
  \boxed{\;
  T_q(\mu_m)
  =
  \frac{1}{1-q}\,
  \varphi^{(1-q)m}
  \frac{\varphi+\varphi^{-q}}{\sqrt5}
  \bigl(1+o(1)\bigr).\;}
  \]
  因而
  \[
  \boxed{\;
  \lim_{m\to\infty}\frac{1}{m}\log T_q(\mu_m)
  =
  (1-q)\log\varphi.\;}
  \]
\end{enumerate}
\end{theorem}
\begin{proof}
由定理 \ref{thm:xi-time-part70ad-binfold-renyi-entropy-constant-drift} 的证明中间式，
\[
\sum_{w\in X_m}\mu_m(w)^q
=
\varphi^{(1-q)m}
\frac{\varphi+\varphi^{-q}}{\sqrt5}
\Bigl(1+O_q\!\Bigl(\Bigl(\frac{\varphi}{2}\Bigr)^m\Bigr)\Bigr).
\]
若 \(q>1\)，则 \(1-q<0\)，上式趋于 \(0\)。代入定义
\[
T_q(\mu_m)
=
\frac{1}{q-1}\left(1-\sum_{w\in X_m}\mu_m(w)^q\right),
\]
便立即得到第一段公式与极限。

若 \(0<q<1\)，则 \(1-q>0\)，故上式指数发散。于是
\[
T_q(\mu_m)
=
\frac{\sum_w\mu_m(w)^q-1}{1-q}
=
\frac{1}{1-q}\,
\varphi^{(1-q)m}
\frac{\varphi+\varphi^{-q}}{\sqrt5}
\bigl(1+o(1)\bigr),
\]
因为被减去的常数 \(1\) 相对于主项可忽略。最后取对数并除以 \(m\)，即可得到指数率
\[
\lim_{m\to\infty}\frac{1}{m}\log T_q(\mu_m)
=(1-q)\log\varphi.
\]
证毕。
\end{proof}

\noindent
由此，bin-fold 输出分布的全阶信息几何在常数层也达到闭合：Rényi 熵不只具有统一熵率，还具有显式 \(O(1)\) 漂移；中心化对数多重度不只具有一阶均值极限，而是整个指数矩与全部有限阶累积量都收敛到同一个二原子模型；规范群熵的指数体量则带有可审计的首项斜率与常数偏移；而 Tsallis 熵在 \(q>1\) 与 \(0<q<1\) 两侧的饱和值与增长率也由同一碰撞矩常数统一给出。于是，两态末位机制、碰撞矩闭式与规范体积展开最终在同一信息几何层面完全合流。


\paragraph{中心双折点容量骨架、escort 温度重加权与 dark band 半阶尺度}
在 bin-fold 的中心临界容量极限、escort 温度族末位塌缩、尾部单跳阶梯以及六倍窗零块下界都已建立之后，还可把阈值/采样轴、相位/路径信息轴与 Fourier 侧 dark band 进一步压缩为同一组闭合接口。其共同特征是：中心层容量骨架始终只含两个不可消去折点，而温度参数只改变同一双折点骨架上的权重，不改变骨架本身。

\begin{theorem}[中心临界容量骨架的双折点闭式]\label{thm:xi-time-part70b-center-capacity-two-kink-skeleton}
定义
\[
A_m:=\left(\frac{2}{\varphi}\right)^m,
\qquad
\Psi_m(\tau):=
2^{-m}\,\mathcal C_m^{\mathrm{bin,cont}}(\tau A_m)
\qquad (\tau\ge 0).
\]
则对每个固定 \(\tau\ge 0\)，都有
\[
\Psi_m(\tau)\longrightarrow \Psi(\tau)
:=
\frac{\varphi}{\sqrt5}\min\{1,\tau\}
+
\frac{1}{\sqrt5}\min\{\varphi^{-1},\tau\}.
\]
等价地，
\[
\Psi(\tau)=
\begin{cases}
\dfrac{\varphi^2}{\sqrt5}\,\tau, & 0\le \tau\le \varphi^{-1},\\[6pt]
\dfrac{\varphi}{\sqrt5}\,\tau+\dfrac{1}{\varphi\sqrt5}, & \varphi^{-1}\le \tau\le 1,\\[8pt]
1, & \tau\ge 1.
\end{cases}
\]
特别地，中心临界容量骨架只有两个不可消去的折点，位置分别位于
\[
\tau=\varphi^{-1},
\qquad
\tau=1.
\]
\end{theorem}
\begin{proof}
这是定理 \ref{thm:xi-foldbin-dyadic-capacity-critical-window-limit} 与定理 \ref{thm:xi-time-part70-continuous-capacity-two-atom-second-derivative} 的直接改写。分段表达式由极限函数逐段展开得到，而折点位置则由分布二阶导数
\[
-D^2\Psi
=
\frac{1}{\sqrt5}\,\delta_{\varphi^{-1}}
+
\frac{\varphi}{\sqrt5}\,\delta_1
\]
唯一确定。证毕。
\end{proof}

\begin{theorem}[escort 温度族对中心双折点骨架的重加权响应]\label{thm:xi-time-part70b-escort-two-kink-response}
对任意固定 \(q\ge 0\)，令 \(W_m^{(q)}\sim \pi_{m,q}\)，并定义
\[
\Psi_{m,q}^{\mathrm{esc}}(\tau)
:=
\mathbb E_{\pi_{m,q}}
\left[
\min\!\left(
1,\frac{\tau A_m}{d_m^{\mathrm{bin}}(W_m^{(q)})}
\right)
\right]
\qquad (\tau\ge 0).
\]
则对每个固定 \(\tau\ge 0\)，极限
\[
\Psi_q^{\mathrm{esc}}(\tau):=
\lim_{m\to\infty}\Psi_{m,q}^{\mathrm{esc}}(\tau)
\]
存在，并满足闭式
\[
\Psi_q^{\mathrm{esc}}(\tau)
=
\bigl(1-\theta(q)\bigr)\min\{1,\tau\}
+
\theta(q)\min\{1,\varphi\tau\},
\qquad
\theta(q):=\frac{1}{1+\varphi^{q+1}}.
\]
因此，escort 温度参数 \(q\) 并不改变中心容量骨架的两个折点位置，而只在同一双折点框架上重新分配两条线性支的权重。
\end{theorem}
\begin{proof}
由定理 \ref{thm:fold-bin-two-state-asymptotic}，存在常数 \(C>0\)，使对充分大的 \(m\) 与全部 \(w\in X_m\) 都有
\[
\left|
\frac{d_m^{\mathrm{bin}}(w)}{A_m}-\varphi^{-w_m}
\right|
\le C\left(\frac{\varphi}{2}\right)^m.
\]
故对充分大的 \(m\)，上述比值一律落在紧区间
\[
\left[\frac{\varphi^{-1}}{2},\,2\right]
\]
内。对固定 \(\tau\ge 0\)，函数
\[
x\longmapsto \min\!\left(1,\frac{\tau}{x}\right)
\]
在该紧区间上 Lipschitz，因而存在常数 \(C_\tau>0\) 使
\[
\min\!\left(1,\frac{\tau A_m}{d_m^{\mathrm{bin}}(w)}\right)
=
\min\!\bigl(1,\tau\varphi^{w_m}\bigr)
+O\!\left(\left(\frac{\varphi}{2}\right)^m\right)
\]
在 \(w\in X_m\) 上一致成立。对 \(\pi_{m,q}\) 取期望，得到
\[
\Psi_{m,q}^{\mathrm{esc}}(\tau)
=
\mathbb E_{\pi_{m,q}}\!\bigl[\min(1,\tau\varphi^{W_m^{(q)}(m)})\bigr]
+O\!\left(\left(\frac{\varphi}{2}\right)^m\right).
\]
另一方面，由命题 \ref{prop:fold-bin-escort-lastbit}，
\[
\mathbb P\bigl(W_m^{(q)}(m)=1\bigr)
=
\theta(q)+O\!\left(\left(\frac{\varphi}{2}\right)^m\right),
\]
\[
\mathbb P\bigl(W_m^{(q)}(m)=0\bigr)
=
1-\theta(q)+O\!\left(\left(\frac{\varphi}{2}\right)^m\right).
\]
于是
\[
\Psi_{m,q}^{\mathrm{esc}}(\tau)
=
\bigl(1-\theta(q)\bigr)\min\{1,\tau\}
+
\theta(q)\min\{1,\varphi\tau\}
+O\!\left(\left(\frac{\varphi}{2}\right)^m\right).
\]
令 \(m\to\infty\)，即得所示极限公式。证毕。
\end{proof}

\begin{corollary}[中心双折点骨架上的温度斜率可辨识性]\label{cor:xi-time-part70b-escort-slope-identifiability}
沿用定理 \ref{thm:xi-time-part70b-escort-two-kink-response} 的记号。则在中间区间
\[
\tau\in (\varphi^{-1},1)
\]
上有
\[
\Psi_q^{\mathrm{esc}}(\tau)
=
\bigl(1-\theta(q)\bigr)\tau+\theta(q),
\]
从而导数恒为
\[
\frac{\dd}{\dd\tau}\Psi_q^{\mathrm{esc}}(\tau)
=
\sigma_q
:=
1-\theta(q)\in(0,1).
\]
并且仅由该斜率 \(\sigma_q\) 即可唯一恢复温度参数 \(q\)：
\[
q=\log_{\varphi}\!\left(\frac{\sigma_q}{1-\sigma_q}\right)-1.
\]
\end{corollary}
\begin{proof}
当 \(\tau\in(\varphi^{-1},1)\) 时，
\[
\min\{1,\tau\}=\tau,
\qquad
\min\{1,\varphi\tau\}=1,
\]
故定理 \ref{thm:xi-time-part70b-escort-two-kink-response} 中的极限曲线退化为
\[
\Psi_q^{\mathrm{esc}}(\tau)
=
\bigl(1-\theta(q)\bigr)\tau+\theta(q).
\]
于是导数恒为 \(1-\theta(q)\)。再由
\[
\theta(q)=\frac{1}{1+\varphi^{q+1}}
\]
可得
\[
\frac{1-\theta(q)}{\theta(q)}=\varphi^{q+1}.
\]
代入 \(\sigma_q=1-\theta(q)\) 并整理，即得所示反演公式。证毕。
\end{proof}

\begin{corollary}[固定末位层的阈值层纹由单跳尾和谱完全恢复]\label{cor:xi-time-part70b-lastbit-tail-jump-recovery}
沿用定理 \ref{thm:xi-foldbin-lastbit-tail-spectrum-single-jump-staircase} 的记号，并把其中的末位指标 \(i\) 记为 \(b\in\{0,1\}\)。定义
\[
\delta_{m,b}(V):=
D_{m,b}(V)
=
\#\bigl\{t\in \mathcal T_{m,b}: t\le 2^m-1-V\bigr\}.
\]
则对一切相容的 \(V\) 都有精确跳跃公式
\[
\delta_{m,b}(V)-\delta_{m,b}(V+1)
=
\mathbf 1_{\{\,2^m-1-V\in\mathcal T_{m,b}\,\}}.
\]
因而合法尾和集可由跳点完全恢复：
\[
\mathcal T_{m,b}
=
\bigl\{2^m-1-V:\ \delta_{m,b}(V)-\delta_{m,b}(V+1)=1\bigr\}.
\]
特别地，在固定末位相位内，每一次向下跳变都恰对应一个合法尾和被阈值切除。
\end{corollary}
\begin{proof}
这只是定理 \ref{thm:xi-foldbin-lastbit-tail-spectrum-single-jump-staircase} 的盒装恒等式与其逆向重写。证毕。
\end{proof}

\begin{theorem}[六倍窗 exact dark band 的半阶指数尺度]\label{thm:xi-time-part70b-dark-band-half-order-scale}
\leanverified{paper\_xi\_time\_part70b\_dark\_band\_half\_order\_scale}
沿用定理 \ref{thm:fold-zero-density-sparse} 与推论 \ref{cor:fold-zero-half-index-multiple6} 的记号，记
\[
\mathcal Z_m:=\{k\in \ZZ/(F_{m+2}\ZZ):\ \widehat d_m(k)=0\},
\]
并令 \(\tau_{\mathrm{div}}(n)\) 表示 \(n\) 的约数个数函数。若
\[
m+2\equiv 0\pmod 6,
\]
则有双边界
\[
F_{(m+2)/2}
\le
|\mathcal Z_m|
\le
\tau_{\mathrm{div}}(m+2)\,F_{(m+2)/2}.
\]
从而沿该六倍窗子序列，
\[
\log|\mathcal Z_m|
=
\frac{m}{2}\log\varphi+O(\log m),
\qquad
\frac{1}{m}\log|\mathcal Z_m|
\longrightarrow
\frac12\log\varphi.
\]
换言之，exact dark band 在六倍窗上具有 \(\varphi^{m/2}\) 级的半阶指数尺度。
\end{theorem}
\begin{proof}
下界正是推论 \ref{cor:fold-zero-half-index-multiple6} 的内容。上界则由定理 \ref{thm:fold-zero-density-sparse} 给出：
\[
|\mathcal Z_m|
\le
\tau_{\mathrm{div}}(m+2)\,F_{\lfloor (m+2)/2\rfloor}
=
\tau_{\mathrm{div}}(m+2)\,F_{(m+2)/2}.
\]
另一方面，Binet 公式给出
\[
\log F_n=n\log\varphi+O(1),
\]
而初等约数界
\[
\tau_{\mathrm{div}}(n)\le 2\sqrt n
\qquad (n\ge 1)
\]
蕴含
\[
\log\tau_{\mathrm{div}}(m+2)=O(\log m).
\]
把这些估计代入上下界，便得到
\[
\log|\mathcal Z_m|
=
\frac{m}{2}\log\varphi+O(\log m),
\]
并进一步推出极限
\[
\frac{1}{m}\log|\mathcal Z_m|\to \frac12\log\varphi.
\]
证毕。
\end{proof}

\noindent
上述结果把中心读出、escort 重加权、阈值层纹与 Fourier 侧零块压缩到同一条双折点--跳点主线上：在阈值/采样轴上，中心容量骨架只允许两个折点 \(\varphi^{-1}\) 与 \(1\)；在相位/路径信息轴上，温度参数 \(q\) 仅通过 \(\theta(q)\) 重新分配两支权重；在固定末位相位内，层纹由合法尾和集的单跳删除精确生成；而六倍窗上的 exact dark band 则在 Fourier 侧取得 \(\varphi^{m/2}\) 级的半阶指数尺度。因而，相应可视化中的连续着色纹理与局部波动若不改变这些离散骨架，便不能改变中心容量、尾和跳点结构与 dark-band 的主尺度。


\paragraph{谱零点的高阶 \texorpdfstring{$\ZZ_2$}{Z2} 轨道抵消}
在谱零点的单坐标判别、余类并结构与六倍窗 half-order dark band 已经建立之后，还可进一步把单个按位翻转对合的相消推广为一个多坐标 \(\ZZ_2\)-轨道抵消机制。其结果是：一旦某个频率 \(k\) 触发若干独立坐标上的 \(-1\) 相位翻转，则整个 \(\Omega_m\) 会按自由的 \((\ZZ/2\ZZ)^r\)-轨道分块，而每个轨道上的相位和都逐块归零。

\begin{theorem}[谱零点的高阶 \texorpdfstring{$\ZZ_2$}{Z2} 轨道抵消]\label{thm:xi-time-part70c-spectrum-zero-z2-orbit-cancellation}
沿用定理 \ref{thm:fold-spectrum-zero-criterion} 与推论 \ref{cor:fold-spectrum-zero-involution} 的记号，记
\[
M:=F_{m+2},
\qquad
\widehat d_m(k)=\sum_{\omega\in\Omega_m}e^{-2\pi i kN(\omega)/M}.
\]
对任意 \(k\in \ZZ/(M\ZZ)\)，定义
\[
J(k):=
\left\{
j\in\{1,\dots,m\}:\ 2kF_{j+1}\equiv M\pmod{2M}
\right\},
\]
并令
\[
\Phi_k(\omega):=e^{-2\pi i kN(\omega)/M}
\qquad
(\omega\in \Omega_m).
\]
对每个 \(j\in J(k)\)，记
\[
\iota_j:\Omega_m\to \Omega_m,
\qquad
\iota_j(\omega):=\omega\oplus e_j.
\]
再定义
\[
G_k:=\langle \iota_j:\ j\in J(k)\rangle.
\]
则成立：
\begin{enumerate}
  \item \(G_k\cong (\ZZ/2\ZZ)^{|J(k)|}\)，且其在 \(\Omega_m\) 上的作用自由；
  \item 对任意
  \[
  g=\prod_{j\in S}\iota_j\in G_k,
  \]
  都有
  \[
  \Phi_k(g\omega)=(-1)^{|S|}\Phi_k(\omega)
  \qquad
  (\omega\in\Omega_m);
  \]
  \item 若 \(J(k)\neq \varnothing\)，则对任意 \(G_k\)-轨道 \(\mathcal O\subset \Omega_m\)，都有
  \[
  \sum_{\omega\in\mathcal O}\Phi_k(\omega)=0;
  \]
  特别地，
  \[
  \widehat d_m(k)=0
  \qquad\Longleftrightarrow\qquad
  J(k)\neq \varnothing.
  \]
\end{enumerate}
\end{theorem}
\begin{proof}
由定理 \ref{thm:fold-spectrum-zero-criterion}，
\[
\widehat d_m(k)=0
\qquad\Longleftrightarrow\qquad
\exists\,j\in\{1,\dots,m\}\ \text{使}\ 2kF_{j+1}\equiv M\pmod{2M},
\]
这正等价于 \(J(k)\neq \varnothing\)。

对每个 \(j\in J(k)\)，推论 \ref{cor:fold-spectrum-zero-involution} 已给出：\(\iota_j\) 是 \(\Omega_m\) 上无固定点的对合，且
\[
\Phi_k(\iota_j\omega)=-\Phi_k(\omega).
\]
若 \(j,j'\in J(k)\) 且 \(j\neq j'\)，则 \(\iota_j,\iota_{j'}\) 只翻转不同坐标，因此彼此交换，且各自平方为恒等。于是
\[
G_k\cong (\ZZ/2\ZZ)^{|J(k)|}.
\]
任一非平凡元素 \(g=\prod_{j\in S}\iota_j\) 都会翻转至少一个坐标，故不可能固定任何 \(\omega\in\Omega_m\)；从而 \(G_k\) 在 \(\Omega_m\) 上自由作用。

再对
\[
g=\prod_{j\in S}\iota_j
\]
重复使用 \(\Phi_k(\iota_j\omega)=-\Phi_k(\omega)\)，便得
\[
\Phi_k(g\omega)=(-1)^{|S|}\Phi_k(\omega).
\]
因此当 \(J(k)\neq \varnothing\) 时，对任意轨道
\[
\mathcal O=G_k\cdot \omega
\]
都有
\[
\sum_{\eta\in\mathcal O}\Phi_k(\eta)
=
\Phi_k(\omega)\sum_{S\subseteq J(k)}(-1)^{|S|}
=
\Phi_k(\omega)(1-1)^{|J(k)|}
=0.
\]
把 \(\Omega_m\) 分解为 \(G_k\)-轨道的并，即得
\[
\widehat d_m(k)=\sum_{\omega\in\Omega_m}\Phi_k(\omega)=0.
\]
与第一步的零点判别合并，即得最后的充要条件。证毕。
\end{proof}


\paragraph{六倍窗零点密度的精确指数率与 KL--最大纤维耦合}
在谱零点的高阶 \texorpdfstring{$(\ZZ_2)^r$}{(Z2)r} 轨道抵消已经给出逐点见证之后，还可继续沿频域--信息主线抽取出下列两条结果。它们分别把六倍窗子序列上的归一化零点密度锁定在严格的半维指数率 \(\varphi^{-m/2+o(m)}\)，并把输出分布相对于均匀分布的 KL 缺口与最大纤维超额压缩为同一条确定性控制不等式。

\begin{theorem}[六倍窗上的零点密度具有精确指数率]\label{thm:xi-time-part70c-zero-density-exact-ratio-exponent-multiple6}
\leanverified{paper\_xi\_time\_part70c\_zero\_density\_exact\_ratio\_exponent\_multiple6}
若
\[
m+2\equiv 0\pmod 6,
\qquad
d:=\frac{m+2}{2},
\qquad
M:=F_{m+2},
\qquad
\mathcal Z_m:=\{k\in \ZZ/(M\ZZ):\ \widehat d_m(k)=0\},
\]
则有极限
\[
\lim_{\substack{m\to\infty\\ m+2\equiv 0\ (6)}}
\frac1m\log\frac{|\mathcal Z_m|}{F_{m+2}}
=
-\frac12\log\varphi.
\]
等价地，在六倍窗子序列上，
\[
\frac{|\mathcal Z_m|}{F_{m+2}}
=
\varphi^{-m/2+o(m)}.
\]
\end{theorem}
\begin{proof}
下界方面，由推论 \ref{cor:fold-zero-half-index-multiple6}，
\[
|\mathcal Z_m|
\ge
F_d.
\]
因此
\[
\frac{|\mathcal Z_m|}{F_{m+2}}
\ge
\frac{F_d}{F_{m+2}}.
\]
由 Binet 公式 \(F_n=\varphi^{n+o(n)}\) 得
\[
\liminf_{\substack{m\to\infty\\ m+2\equiv 0\ (6)}}
\frac1m\log\frac{|\mathcal Z_m|}{F_{m+2}}
\ge
\lim_{m\to\infty}
\frac1m\log\frac{F_d}{F_{m+2}}
=
-\frac12\log\varphi.
\]

上界方面，由定理 \ref{thm:fold-zero-density-sparse}，
\[
\frac{|\mathcal Z_m|}{F_{m+2}}
\le
\tau(m+2)\frac{F_{\lfloor (m+2)/2\rfloor}}{F_{m+2}}.
\]
又由于
\[
\tau(n)=e^{o(n)},
\qquad
F_{\lfloor (m+2)/2\rfloor}= \varphi^{m/2+o(m)},
\qquad
F_{m+2}=\varphi^{m+o(m)},
\]
故
\[
\limsup_{\substack{m\to\infty\\ m+2\equiv 0\ (6)}}
\frac1m\log\frac{|\mathcal Z_m|}{F_{m+2}}
\le
-\frac12\log\varphi.
\]
上下界相合，即得结论。
\end{proof}

\begin{theorem}[最大纤维超额由 KL 缺口确定性控制]\label{thm:xi-time-part70c-kl-gap-controls-maxfiber-excess}
设
\[
p_m(r):=\frac{d_m(r)}{2^m},
\qquad
M:=F_{m+2},
\qquad
u(r):=\frac1M,
\]
并记
\[
M_m:=\max_r d_m(r),
\qquad
d_m^{\mathrm{avg}}:=\frac{2^m}{M}.
\]
则有确定性耦合不等式
\[
\boxed{\;
M_m-d_m^{\mathrm{avg}}
\le
2^m\,\|p_m-u\|_{\mathrm{TV}}
\le
2^m\sqrt{\frac12\,D_{\mathrm{KL}}(p_m\|u)}.\;}
\]
特别地，一旦 KL 缺口被控制，最大纤维相对于平均纤维的超额便被同步控制。
\end{theorem}
\begin{proof}
首先由
\[
M_m=2^m\max_r p_m(r),
\qquad
d_m^{\mathrm{avg}}=\frac{2^m}{M}=2^m u(r),
\]
可得
\[
M_m-d_m^{\mathrm{avg}}
=
2^m\left(\max_r p_m(r)-\frac1M\right)
\le
2^m\max_r|p_m(r)-u(r)|.
\]
另一方面，对任意概率分布 \(p,u\) 都有
\[
\max_r|p(r)-u(r)|
\le
\frac12\sum_r|p(r)-u(r)|
=
\|p-u\|_{\mathrm{TV}},
\]
故
\[
M_m-d_m^{\mathrm{avg}}
\le
2^m\,\|p_m-u\|_{\mathrm{TV}}.
\]
最后应用 Pinsker 不等式
\[
\|p_m-u\|_{\mathrm{TV}}
\le
\sqrt{\frac12\,D_{\mathrm{KL}}(p_m\|u)},
\]
便得到所示结论。证毕。
\end{proof}

\noindent
由此，谱零点结构在三个层级上完成闭合：在单频层级，高阶 \texorpdfstring{$(\ZZ_2)^r$}{(Z2)r} 轨道抵消给出逐点见证；在密度层级，六倍窗上的归一化零点比例被锁定在 \(\varphi^{-m/2+o(m)}\)；在信息层级，KL 缺口与最大纤维超额被耦合为同一条确定性不等式。于是，轨道抵消、半维稀疏率与信息偏置便被压缩进同一条完全可审计的频域主线。


\paragraph{末位终板刚性、临界预言机振荡与 Riesz 能量常数}
在末位相位的最终严格分层、bin-fold 规范群的最终满秩化、critical-scale 预言机容量骨架、escort 温度族的双折点响应以及六倍窗 exact dark band 的半阶指数尺度已经确立之后，还可把终板几何、商结构障碍、临界预算振荡与二阶能量常数进一步压缩为同一条后续链。在本段中统一记
\[
d_m(w):=d_m^{\mathrm{bin}}(w),
\qquad
\lambda:=\frac{2}{\varphi},
\qquad
\Omega_m:=\{0,1\}^m.
\]
下列结果分别刻画最小扇区的末位实现、大窗口上的全局非商化刚性、异常签名维数的 Fibonacci 饱和、临界比特率下的不可压平簇集带、escort 温度族的临界响应带、中心纤维场的精确二阶常数、碰撞概率与非平凡 Fourier 能量的精确主项，以及最大纤维超额的显式主系数。

\begin{theorem}[最小扇区的末位实现与终板逼近]\label{thm:xi-time-part70d-minsector-terminal-sheet-realization}
定义
\[
d_{\min}(m):=\min_{v\in X_m}d_m(v),
\qquad
A_1^{(m)}:=\{w\in X_m:\ w_m=1\},
\qquad
S_{m,\min}:=\{w\in X_m:\ d_m(w)=d_{\min}(m)\}.
\]
若“子覆盖同构”恒等式
\[
|S_{m,\min}|=|X_{m-2}|
\]
沿某个趋于无穷的偶数序列成立，则沿同一序列有
\[
\frac{|S_{m,\min}\triangle A_1^{(m)}|}{|X_m|}\longrightarrow 0.
\]
\end{theorem}
\leanverified{paper\_xi\_time\_part70d\_minsector\_terminal\_sheet\_realization}
\begin{proof}
由定理 \ref{thm:xi-time-part70ab-terminal-bit-eventual-strict-stratification}，对充分大的 \(m\) 有
\[
S_{m,\min}\subseteq A_1^{(m)}.
\]
另一方面，在假设下
\[
\frac{|S_{m,\min}|}{|X_m|}
=
\frac{|X_{m-2}|}{|X_m|}
=
\frac{F_m}{F_{m+2}}
\longrightarrow
\varphi^{-2},
\]
这也正是推论 \ref{cor:fold-bin-minsector-density-phi-minus2} 的内容。又由稳定词计数，
\[
|A_1^{(m)}|=F_m,
\qquad
|X_m|=F_{m+2},
\]
故
\[
\frac{|A_1^{(m)}|}{|X_m|}=\frac{F_m}{F_{m+2}}\longrightarrow \varphi^{-2}.
\]
由于 \(S_{m,\min}\subseteq A_1^{(m)}\) 且两者相对密度极限相同，遂有
\[
\frac{|A_1^{(m)}\setminus S_{m,\min}|}{|X_m|}\longrightarrow 0.
\]
而此时
\[
S_{m,\min}\triangle A_1^{(m)}=A_1^{(m)}\setminus S_{m,\min},
\]
故结论成立。
\end{proof}

\begin{theorem}[大窗口上的全局非商化刚性]\label{thm:xi-time-part70d-global-nonquotient-rigidity}
存在整数 \(m_0\)，使得对一切 \(m\ge m_0\)，等价关系
\[
\omega\sim_m\eta
\qquad\Longleftrightarrow\qquad
\Fold_m^{\mathrm{bin}}(\omega)=\Fold_m^{\mathrm{bin}}(\eta)
\]
既不能由任何有限群在 \(\Omega_m\) 上的自由作用之轨道划分实现，也不能来自任何有限阿贝尔群上的线性陪集划分。更准确地，不存在有限阿贝尔群 \(A_m\)、子群 \(H_m\le A_m\) 及双射
\[
\iota_m:\Omega_m\longrightarrow A_m
\]
使得
\[
\omega\sim_m\eta
\qquad\Longleftrightarrow\qquad
\iota_m(\omega)-\iota_m(\eta)\in H_m.
\]
\end{theorem}
\leanverified{paper\_xi\_time\_part70d\_global\_nonquotient\_rigidity}
\begin{proof}
记
\[
A_0^{(m)}:=\{w\in X_m:\ w_m=0\},
\qquad
A_1^{(m)}:=\{w\in X_m:\ w_m=1\}.
\]
由定理 \ref{thm:xi-time-part70ab-terminal-bit-eventual-strict-stratification}，存在 \(m_0\) 使得当 \(m\ge m_0\) 时，
\[
\min_{w\in A_0^{(m)}}d_m(w)>\max_{v\in A_1^{(m)}}d_m(v).
\]
由于 \(A_0^{(m)}\) 与 \(A_1^{(m)}\) 都非空，故对每个 \(m\ge m_0\)，纤维大小在同一窗口内至少取两个不同整数值。

若 \(\sim_m\) 来自某个有限群 \(G_m\) 的自由作用，则每个轨道大小都恒等于 \(|G_m|\)，因而所有纤维必同势；矛盾。

若 \(\sim_m\) 来自上式所述的线性陪集划分，则每个等价类都是某个陪集的原像，其大小恒等于 \(|H_m|\)，故所有纤维同样必同势；再度矛盾。于是两类实现均不可能存在。
\end{proof}

\begin{theorem}[异常签名维数的渐近饱和]\label{thm:xi-time-part70d-anomaly-signature-saturation}
记
\[
G_m^{\mathrm{bin}}:=\prod_{w\in X_m}\mathfrak S_{d_m(w)}.
\]
则存在整数 \(m_0\)，使得对一切 \(m\ge m_0\)，都有
\[
\bigl(G_m^{\mathrm{bin}}\bigr)^{\mathrm{ab}}
\cong
(\ZZ/2\ZZ)^{|X_m|}
\cong
(\ZZ/2\ZZ)^{F_{m+2}}.
\]
因而任何保持“纤维独立置换”语义、并在 \(\FF_2\) 上完备分离 \(\bigl(G_m^{\mathrm{bin}}\bigr)^{\mathrm{ab}}\) 全部坐标的异常见证系统，其最小比特数恰为
\[
F_{m+2}.
\]
\end{theorem}
\begin{proof}
这正是定理 \ref{thm:xi-time-part65-binfold-center-vanishing-signature-stability} 的重述。该定理已证明：对充分大的 \(m\)，有
\[
\bigl(G_m^{\mathrm{bin}}\bigr)^{\mathrm{ab}}
\cong
(\ZZ/2\ZZ)^{F_{m+2}},
\qquad
\operatorname{rank}_{\FF_2}\!\bigl((G_m^{\mathrm{bin}})^{\mathrm{ab}}\bigr)=F_{m+2},
\]
并且纤维奇偶荷的最小完备秩在同一范围内严格等于 \(F_{m+2}\)。这恰好给出所述异常签名维数结论。
\end{proof}

\begin{theorem}[均匀微态输入下的临界比特率相变与簇集区间]\label{thm:xi-time-part70d-critical-bitrate-phase-cluster}
\leanverified{paper\_xi\_time\_part70d\_critical\_bitrate\_phase\_cluster}
记
\[
\alpha:=\log_2\!\left(\frac{2}{\varphi}\right),
\qquad
\mathrm{Succ}_m(B):=\mathrm{Succ}_m^{\mathrm{bin}}(B).
\]
则成立：
\begin{enumerate}
  \item 若 \(\rho<\alpha\)，则
  \[
  \mathrm{Succ}_m(\lfloor \rho m\rfloor)\longrightarrow 0.
  \]
  \item 若 \(\rho>\alpha\)，则
  \[
  \mathrm{Succ}_m(\lfloor \rho m\rfloor)\longrightarrow 1.
  \]
  \item 对任意固定 \(\beta\in\RR\)，令
  \[
  B_m:=\lfloor \alpha m+\beta\rfloor.
  \]
  则序列 \(\bigl(\mathrm{Succ}_m(B_m)\bigr)_{m\ge 1}\) 不收敛，其全部簇集点恰组成闭区间
  \[
  \left[\frac{\varphi^2}{2\sqrt5},\,1\right].
  \]
\end{enumerate}
\end{theorem}
\begin{proof}
记 \(\lambda:=2/\varphi\)，则 \(\alpha=\log_2\lambda\)。

若 \(\rho<\alpha\)，则由平凡上界
\[
C_m^{\mathrm{bin}}(B)\le 2^B|X_m|
\]
得到
\[
\mathrm{Succ}_m(\lfloor \rho m\rfloor)
\le
\frac{2^{\lfloor \rho m\rfloor}F_{m+2}}{2^m}
=
\frac{2^{\lfloor \rho m\rfloor}}{\lambda^m}\cdot \frac{F_{m+2}}{\varphi^m}
\longrightarrow 0,
\]
因为 \(2^{\lfloor \rho m\rfloor}/\lambda^m\to 0\) 且 \(F_{m+2}/\varphi^m\to \varphi^2/\sqrt5\)。

若 \(\rho>\alpha\)，则 \(2^{\lfloor \rho m\rfloor}/\lambda^m\to+\infty\)。由推论 \ref{cor:xi-time-part70ab-extremal-fibers-terminal-bit-localization}，
\[
M_m:=\max_{w\in X_m}d_m(w)=\lambda^m+O(1).
\]
故对充分大的 \(m\) 有 \(2^{\lfloor \rho m\rfloor}>M_m\)，从而所有纤维都完全通过，即
\[
\mathrm{Succ}_m(\lfloor \rho m\rfloor)=1.
\]

现取临界序列 \(B_m=\lfloor \alpha m+\beta\rfloor\)。定义
\[
\theta_m:=\alpha m+\beta-\lfloor \alpha m+\beta\rfloor\in[0,1),
\qquad
\tau_m:=\frac{2^{B_m}}{\lambda^m}=2^{-\theta_m}\in[1/2,1].
\]
若 \(\log_2\varphi\in\QQ\)，则存在整数 \(p,q\ge 1\) 使 \(\varphi^q=2^p\in\QQ\)，这与 \(\varphi^q\notin\QQ\) 矛盾。故 \(\alpha=1-\log_2\varphi\) 为无理数，从而 \((\theta_m)\) 在 \([0,1]\) 上稠密，\((\tau_m)\) 的簇集点恰为整个区间 \([1/2,1]\)。

任取 \(\tau\in[1/2,1]\) 并取子列 \(m_j\) 使 \(\tau_{m_j}\to\tau\)。由定理 \ref{thm:xi-foldbin-dyadic-capacity-critical-window-limit}，
\[
\mathrm{Succ}_{m_j}(B_{m_j})
\longrightarrow
\Psi(\tau),
\]
其中
\[
\Psi(\tau)=
\begin{cases}
\dfrac{\varphi^2}{\sqrt5}\,\tau, & 0\le \tau\le \varphi^{-1},\\[6pt]
\dfrac{\varphi}{\sqrt5}\,\tau+\dfrac{1}{\varphi\sqrt5}, & \varphi^{-1}\le \tau\le 1,\\[8pt]
1, & \tau\ge 1.
\end{cases}
\]
因此 \(\bigl(\mathrm{Succ}_m(B_m)\bigr)\) 的全部簇集点恰为连续单调函数 \(\Psi\) 在 \([1/2,1]\) 上的像，即
\[
\Psi([1/2,1])
=
\left[\Psi(1/2),\,\Psi(1)\right]
=
\left[\frac{\varphi^2}{2\sqrt5},\,1\right].
\]
该区间非单点，故序列不收敛。
\end{proof}

\begin{theorem}[escort 温度族下的临界响应带]\label{thm:xi-time-part70d-escort-critical-response-band}
\leanverified{paper\_xi\_time\_part70d\_escort\_critical\_response\_band}
对任意 \(q\ge 0\)，定义 escort 加权的阈值响应
\[
\mathrm{Succ}_{m}^{(q)}(B)
:=
\sum_{w\in X_m}\pi_{m,q}(w)\min\!\left(1,\frac{2^B}{d_m(w)}\right),
\qquad
\theta(q):=\frac{1}{1+\varphi^{q+1}}.
\]
则同一临界比特率
\[
\alpha=\log_2\!\left(\frac{2}{\varphi}\right)
\]
给出 sharp transition：
\begin{enumerate}
  \item 若 \(\rho<\alpha\)，则
  \[
  \mathrm{Succ}_{m}^{(q)}(\lfloor \rho m\rfloor)\longrightarrow 0.
  \]
  \item 若 \(\rho>\alpha\)，则
  \[
  \mathrm{Succ}_{m}^{(q)}(\lfloor \rho m\rfloor)\longrightarrow 1.
  \]
  \item 若对固定 \(\beta\in\RR\) 令
  \[
  B_m=\lfloor \alpha m+\beta\rfloor,
  \]
  则序列 \(\bigl(\mathrm{Succ}_{m}^{(q)}(B_m)\bigr)_{m\ge 1}\) 不收敛，其全部簇集点恰为
  \[
  \left[\frac{1+(\varphi-1)\theta(q)}{2},\,1\right].
  \]
\end{enumerate}
\end{theorem}
\begin{proof}
若 \(\rho<\alpha\)，则 \(2^{\lfloor \rho m\rfloor}/\lambda^m\to 0\)。由定理 \ref{thm:fold-bin-two-state-asymptotic}，存在常数 \(C>0\) 使对充分大的 \(m\) 与全部 \(w\in X_m\) 都有
\[
d_m(w)\ge \lambda^m\left(\varphi^{-1}-C\Bigl(\frac{\varphi}{2}\Bigr)^m\right).
\]
于是
\[
\mathrm{Succ}_{m}^{(q)}(\lfloor \rho m\rfloor)
\le
\frac{2^{\lfloor \rho m\rfloor}}{\lambda^m\left(\varphi^{-1}-C(\varphi/2)^m\right)}
\longrightarrow 0.
\]

若 \(\rho>\alpha\)，则记
\[
M_m:=\max_{w\in X_m}d_m(w).
\]
由推论 \ref{cor:xi-time-part70ab-extremal-fibers-terminal-bit-localization}，\(M_m=\lambda^m+O(1)\)。故对充分大的 \(m\) 有 \(2^{\lfloor \rho m\rfloor}>M_m\)，于是每个纤维都完全通过，从而
\[
\mathrm{Succ}_{m}^{(q)}(\lfloor \rho m\rfloor)=1.
\]

现令 \(B_m=\lfloor \alpha m+\beta\rfloor\)，并仍记
\[
\tau_m:=\frac{2^{B_m}}{\lambda^m}\in[1/2,1].
\]
由上一条定理的证明，\((\tau_m)\) 的全部簇集点恰为 \([1/2,1]\)。

记
\[
I:=[1/2,1],
\qquad
J:=\left[\frac{\varphi^{-1}}{2},\,2\right].
\]
函数
\[
f(\tau,x):=\min\!\left(1,\frac{\tau}{x}\right)
\]
在紧集 \(I\times J\) 上联合 Lipschitz。再由定理 \ref{thm:fold-bin-two-state-asymptotic}，
\[
\left|\frac{d_m(w)}{\lambda^m}-\varphi^{-w_m}\right|
\le
C\Bigl(\frac{\varphi}{2}\Bigr)^m
\qquad (w\in X_m),
\]
故对充分大的 \(m\) 与全部 \(w\in X_m\) 都有
\[
\min\!\left(1,\frac{2^{B_m}}{d_m(w)}\right)
=
f\!\left(\tau_m,\frac{d_m(w)}{\lambda^m}\right)
=
f\!\left(\tau_m,\varphi^{-w_m}\right)
+O\!\left(\left(\frac{\varphi}{2}\right)^m\right),
\]
且误差对 \(w\) 与 \(\tau_m\in I\) 一致。对 \(\pi_{m,q}\) 取期望，并用命题 \ref{prop:fold-bin-escort-lastbit}
\[
\mathbf P_{\pi_{m,q}}(w_m=1)=\theta(q)+O\!\left(\left(\frac{\varphi}{2}\right)^m\right),
\qquad
\mathbf P_{\pi_{m,q}}(w_m=0)=1-\theta(q)+O\!\left(\left(\frac{\varphi}{2}\right)^m\right),
\]
得到
\[
\mathrm{Succ}_{m}^{(q)}(B_m)
=
\bigl(1-\theta(q)\bigr)\min\{1,\tau_m\}
+
\theta(q)\min\{1,\varphi\tau_m\}
+o(1).
\]
任取 \(\tau\in[1/2,1]\) 及子列 \(m_j\) 使 \(\tau_{m_j}\to\tau\)，便有
\[
\mathrm{Succ}_{m_j}^{(q)}(B_{m_j})
\longrightarrow
\Psi_q^{\mathrm{esc}}(\tau),
\qquad
\Psi_q^{\mathrm{esc}}(\tau)
:=
\bigl(1-\theta(q)\bigr)\min\{1,\tau\}
+
\theta(q)\min\{1,\varphi\tau\}.
\]
故全部簇集点恰为 \(\Psi_q^{\mathrm{esc}}([1/2,1])\)。由于 \(\Psi_q^{\mathrm{esc}}\) 在 \([1/2,1]\) 上连续单调，故该像为区间
\[
\left[\Psi_q^{\mathrm{esc}}(1/2),\,\Psi_q^{\mathrm{esc}}(1)\right].
\]
又因 \(1/2<\varphi^{-1}\)，故
\[
\Psi_q^{\mathrm{esc}}(1/2)
=
\frac{1-\theta(q)}{2}+\frac{\varphi\theta(q)}{2}
=
\frac{1+(\varphi-1)\theta(q)}{2},
\qquad
\Psi_q^{\mathrm{esc}}(1)=1.
\]
于是所求簇集区间即为
\[
\left[\frac{1+(\varphi-1)\theta(q)}{2},\,1\right].
\]
其长度严格为正，故序列不收敛。
\end{proof}

\begin{theorem}[中心纤维场的精确二阶常数]\label{thm:xi-time-part70d-central-fiber-second-moment}
\leanverified{paper\_xi\_time\_part70d\_central\_fiber\_second\_moment}
令 \(U_m\sim \mathrm{Unif}(X_m)\)，并定义尺度化纤维场
\[
Y_m:=\left(\frac{\varphi}{2}\right)^m d_m(U_m).
\]
则有极限
\[
\mathbb E(Y_m^2)\longrightarrow 2\varphi^{-2},
\qquad
\operatorname{Var}(Y_m)\longrightarrow \varphi^{-7}.
\]
\end{theorem}
\begin{proof}
由推论 \ref{cor:xi-time-part70a-scaled-moments-leading-constant}，当 \(q=1,2\) 时分别有
\[
\mathbb E(Y_m)
\longrightarrow
\frac{1}{\varphi}+\varphi^{-3},
\qquad
\mathbb E(Y_m^2)
\longrightarrow
\frac{1}{\varphi}+\varphi^{-4}.
\]
利用恒等式 \(\varphi^2=\varphi+1\) 与 \(\sqrt5=2\varphi-1\)，可将上式化简为
\[
\frac{1}{\varphi}+\varphi^{-4}=2\varphi^{-2},
\qquad
\frac{1}{\varphi}+\varphi^{-3}=\frac{\sqrt5}{\varphi^2}.
\]
因此
\[
\operatorname{Var}(Y_m)
=
\mathbb E(Y_m^2)-\bigl(\mathbb E(Y_m)\bigr)^2
\longrightarrow
2\varphi^{-2}-\frac{5}{\varphi^4}.
\]
再用 \(\varphi^2=\varphi+1\) 化简，得到
\[
2\varphi^{-2}-\frac{5}{\varphi^4}=\varphi^{-7}.
\]
证毕。
\end{proof}

\begin{theorem}[碰撞概率与非平凡 Fourier 能量的精确主项]\label{thm:xi-time-part70d-collision-riesz-main-terms}
\leanverified{paper\_xi\_time\_part70d\_collision\_riesz\_main\_terms}
沿用定理 \ref{thm:xi-time-part70d-central-fiber-second-moment} 的记号，并沿用定理 \ref{thm:xi-time-part9p-chi2-ranktwo-resonant-floor} 的 Fourier 记号。记
\[
\mathsf{Col}_m:=
\frac{1}{4^m}\sum_{w\in X_m}d_m(w)^2.
\]
则有
\[
\mathsf{Col}_m\sim \frac{2}{\sqrt5}\,\varphi^{-m},
\qquad
F_{m+2}\,\mathsf{Col}_m\longrightarrow \frac{2\varphi^2}{5}.
\]
进一步，非平凡 Fourier 能量满足
\[
4^{-m}\sum_{k\ne 0}\bigl|\widehat d_m(k)\bigr|^2
\longrightarrow
\frac{1}{5\varphi^3}.
\]
\end{theorem}
\begin{proof}
由 \(U_m\sim\mathrm{Unif}(X_m)\) 的定义，
\[
\mathsf{Col}_m
=
\frac{|X_m|}{4^m}\,\mathbb E\!\bigl(d_m(U_m)^2\bigr).
\]
又因为
\[
d_m(U_m)=\lambda^m Y_m,
\qquad
\lambda=\frac{2}{\varphi},
\]
故
\[
\mathsf{Col}_m
=
|X_m|\,\varphi^{-2m}\,\mathbb E(Y_m^2).
\]
由 \(|X_m|=F_{m+2}\) 与定理 \ref{thm:xi-time-part70d-central-fiber-second-moment}，
\[
\mathsf{Col}_m
=
F_{m+2}\varphi^{-2m}\bigl(2\varphi^{-2}+o(1)\bigr).
\]
再用 Binet 公式
\[
F_{m+2}\sim \frac{\varphi^{m+2}}{\sqrt5},
\]
便得
\[
\mathsf{Col}_m
\sim
\frac{\varphi^{m+2}}{\sqrt5}\cdot \varphi^{-2m}\cdot 2\varphi^{-2}
=
\frac{2}{\sqrt5}\,\varphi^{-m}.
\]
两边再乘以 \(F_{m+2}\) 即得
\[
F_{m+2}\,\mathsf{Col}_m\longrightarrow \frac{2\varphi^2}{5}.
\]

另一方面，由定理 \ref{thm:xi-time-part9p-chi2-ranktwo-resonant-floor} 的第一条精确恒等式，
\[
F_{m+2}\,\mathsf{Col}_m-1
=
4^{-m}\sum_{k\ne 0}\bigl|\widehat d_m(k)\bigr|^2.
\]
取极限并用
\[
\frac{2\varphi^2}{5}-1=\frac{1}{5\varphi^3},
\]
便得到
\[
4^{-m}\sum_{k\ne 0}\bigl|\widehat d_m(k)\bigr|^2
\longrightarrow
\frac{1}{5\varphi^3}.
\]
证毕。
\end{proof}

\begin{corollary}[最大纤维超额的精确主项]\label{cor:xi-time-part70d-maxfiber-excess-exact-main-term}
\leanverified{paper\_xi\_time\_part70d\_maxfiber\_excess\_exact\_main\_term}
记
\[
\bar d_m:=\frac{2^m}{F_{m+2}},
\qquad
M_m:=\max_{w\in X_m}d_m(w).
\]
则有
\[
M_m-\bar d_m
=
\varphi^{-4}\left(\frac{2}{\varphi}\right)^m
+o\!\left(\left(\frac{2}{\varphi}\right)^m\right).
\]
\end{corollary}
\begin{proof}
由推论 \ref{cor:xi-time-part70ab-extremal-fibers-terminal-bit-localization}，
\[
M_m=\left(\frac{2}{\varphi}\right)^m+O(1).
\]
另一方面，由 Binet 公式，
\[
\bar d_m
=
\frac{2^m}{F_{m+2}}
=
\frac{\sqrt5}{\varphi^2}\left(\frac{2}{\varphi}\right)^m
+o\!\left(\left(\frac{2}{\varphi}\right)^m\right).
\]
由于 \(O(1)=o((2/\varphi)^m)\)，相减即得
\[
M_m-\bar d_m
=
\left(1-\frac{\sqrt5}{\varphi^2}\right)\left(\frac{2}{\varphi}\right)^m
+o\!\left(\left(\frac{2}{\varphi}\right)^m\right).
\]
再用 \(\sqrt5=2\varphi-1\) 与 \(\varphi-1=\varphi^{-1}\) 化简：
\[
1-\frac{\sqrt5}{\varphi^2}
=
1-\frac{2\varphi-1}{\varphi^2}
=
\frac{(\varphi-1)^2}{\varphi^2}
=
\varphi^{-4}.
\]
故结论成立。
\end{proof}

\noindent
由此，末位 \(1\) 终板不再只是最小纤维的容纳层，而在“子覆盖同构”外推下渐近地与整个最小扇区重合；与此同时，充分大窗口上的折叠关系也被彻底排除出自由商与线性陪集两类等势模型。再向临界资源侧推进，dyadic 预算与 escort 温度都在同一比特率 \(\log_2(2/\varphi)\) 处出现不可压平的簇集带，而其下边界由黄金比例与温度参数共同给出显式常数。最后，中心纤维场的二阶矩、碰撞概率、非平凡 Fourier 能量与最大纤维超额都锁定到闭式黄金常数，从而终板刚性、资源相变与频域能量三条先前分立的链被压缩为同一条可审计的 Fibonacci 后续主线。

\paragraph{window-\texorpdfstring{$6$}{6} 抽象轨道型的最小循环实现与几何障碍的范畴分离}
在 window-\(6\) 的纤维多值性、自由作用障碍与线性陪集障碍都已确立之后，还可进一步把“抽象轨道大小多重集”与“真实立方体上的几何实现”彻底分开：就纯 \(G\)-集合而言，直方图 \(2^8,3^4,4^9\) 实际上已经能由最小的 \(12\) 阶循环群实现；因此 window-\(6\) 的真正障碍并不位于 Burnside 型轨道计数层，而位于把该轨道型忠实落到 \(\Omega_6\) 上时所必须满足的自由性与线性几何兼容条件。

\begin{theorem}[window-\texorpdfstring{$6$}{6} 的抽象轨道型可由最小阶循环群 \texorpdfstring{$C_{12}$}{C12} 实现]\label{thm:xi-time-part70da-window6-abstract-orbittype-minimal-c12}
设
\[
\mathfrak m:=2^8\,3^4\,4^9
\]
表示轨道大小多重集，其中大小为 \(2,3,4\) 的轨道个数分别为 \(8,4,9\)。则：
\begin{enumerate}
  \item 若某有限群 \(G\) 的某个有限 \(G\)-集合具有轨道型 \(\mathfrak m\)，则必有
  \[
  12\mid |G|.
  \]
  \item 该下界是锐的。更精确地，循环群 \(C_{12}\) 已足够：若对每个 \(d\mid 12\) 记 \(C_d\le C_{12}\) 为唯一的阶为 \(d\) 的子群，则
  \[
  Y:=8(C_{12}/C_6)\ \sqcup\ 4(C_{12}/C_4)\ \sqcup\ 9(C_{12}/C_3)
  \]
  是一个 \(64\) 点 \(C_{12}\)-集合，其轨道大小多重集恰为 \(\mathfrak m\)。
\end{enumerate}
因此，就抽象轨道算术而言，实现 window-\(6\) 直方图所需的最小群阶恰为
\[
12.
\]
\end{theorem}
\begin{proof}
若有限群 \(G\) 的某个轨道分解中出现大小为 \(2,3,4\) 的轨道，则每个轨道大小都必须整除 \(|G|\)。因此
\[
2\mid |G|,
\qquad
3\mid |G|,
\qquad
4\mid |G|,
\]
从而
\[
\operatorname{lcm}(2,3,4)=12\mid |G|.
\]
这就给出第一条与最小群阶的下界。

现取 \(G=C_{12}\)。由于循环群对每个除数都只有唯一子群，故存在唯一的阶 \(6,4,3\) 子群 \(C_6,C_4,C_3\)。相应陪集空间的基数分别为
\[
|C_{12}/C_6|=2,
\qquad
|C_{12}/C_4|=3,
\qquad
|C_{12}/C_3|=4.
\]
因此上式定义的离散并
\[
Y:=8(C_{12}/C_6)\ \sqcup\ 4(C_{12}/C_4)\ \sqcup\ 9(C_{12}/C_3)
\]
确为一个 \(C_{12}\)-集合，并且其轨道大小多重集恰为
\[
2^8\,3^4\,4^9.
\]
总点数为
\[
8\cdot 2+4\cdot 3+9\cdot 4=16+12+36=64.
\]
故 \(C_{12}\) 已实现该抽象轨道型。结合上界即知最小群阶恰为 \(12\)。证毕。
\end{proof}

\begin{corollary}[window-\texorpdfstring{$6$}{6} 的压缩障碍是几何性的，而非纯轨道计数性的]\label{cor:xi-time-part70da-window6-obstruction-geometric-not-burnside}
对 \(\Omega_6=\{0,1\}^6\) 上的 bin-fold 纤维划分，成立如下分层：
\begin{enumerate}
  \item 在抽象有限 \(G\)-集合范畴中，轨道型
  \[
  2^8\,3^4\,4^9
  \]
  已由定理 \ref{thm:xi-time-part70da-window6-abstract-orbittype-minimal-c12} 中的 \(C_{12}\)-集合实现。
  \item 在真实立方体 \(\Omega_6\) 上，该纤维划分不可能来自任何有限群的自由作用轨道分解。
  \item 同样地，该纤维划分也不可能来自任何有限阿贝尔群上的线性陪集划分。
\end{enumerate}
因此，window-\(6\) 的压缩障碍并不是 Burnside 型轨道大小不可实现性，而是一个严格更强的几何--线性兼容性障碍。
\end{corollary}
\begin{proof}
第一条正是定理 \ref{thm:xi-time-part70da-window6-abstract-orbittype-minimal-c12}。第二、三条则由推论 \ref{cor:foldbin6-degeneracy-multivalued-nonfree} 直接给出：window-\(6\) 的纤维直方图
\[
2:8,\qquad 3:4,\qquad 4:9
\]
不能来自任何自由有限群作用，也不能来自任何有限阿贝尔群上的线性陪集划分。将三条结论并列，便得到所述范畴分层。证毕。
\end{proof}

\paragraph{整振幅、仿射压力与 Stokes 归一化的单核闭合}
在稳定层两态一致渐近、bin-fold 退化度的全实矩闭式、相对稳定层均匀基线的 Rényi 散度极限，以及 fold-gauge 常数的 Stirling--Bernoulli 奇层展开均已建立之后，还可把这些此前分散出现的接口压缩为一条复参数解析骨架：同一个整函数振幅同时控制 bin-fold 的复矩主项、输出分布的全轴压力、Rényi 极限曲线的规范延拓以及 \(\log C_m\) 尾部的 Bernoulli 采样。由此，一切归一化无关的有限窗对数观测都退化为该整振幅在负奇整数点上的读出，而剩余自由度严格只保留一个乘法常数。

\begin{theorem}[bin-fold 全复矩谱的整振幅律]\label{thm:xi-time-part70e-binfold-entire-amplitude-law}
\leanverified{paper\_xi\_time\_part70e\_binfold\_entire\_amplitude\_law}
记
\[
\varrho:=\frac{\varphi}{2}\in(0,1),
\qquad
d_m(w):=d_m^{\mathrm{bin}}(w),
\qquad
S_m(z):=\sum_{w\in X_m} d_m(w)^z
\qquad (z\in\CC),
\]
其中 \(d_m(w)^z:=\exp(z\log d_m(w))\) 取实正数上的主支。则对任意紧集 \(K\subset\CC\)，都有一致渐近
\[
S_m(z)
=
\Bigl(\frac{2}{\varphi}\Bigr)^{zm}
\Bigl(F_{m+1}+\varphi^{-z}F_m\Bigr)
\Bigl(1+O_K(\varrho^m)\Bigr)
\qquad (z\in K).
\]
从而
\[
\mathcal A_m(z):=
\varphi^{-m}\Bigl(\frac{\varphi}{2}\Bigr)^{zm}S_m(z)
\]
在 \(\CC\) 的每个紧集上一致收敛到
\[
\mathcal A(z):=\frac{\varphi+\varphi^{-z}}{\sqrt5},
\]
因此 \(\mathcal A\) 为整函数。
\end{theorem}
\begin{proof}
由定理 \ref{thm:fold-bin-two-state-asymptotic}，存在 \(\varepsilon_m(w)\) 满足
\[
d_m(w)=\varphi^{-w_m}\Bigl(\frac{2}{\varphi}\Bigr)^m\bigl(1+\varepsilon_m(w)\bigr),
\qquad
\sup_{w\in X_m}|\varepsilon_m(w)|=O(\varrho^m).
\]
固定紧集 \(K\subset\CC\)。对充分大的 \(m\) 有 \(\sup_{w}|\varepsilon_m(w)|<1/2\)，因而
\[
\log\bigl(1+\varepsilon_m(w)\bigr)=O(\varrho^m)
\]
在 \(w\in X_m\) 上一致成立。由于 \(K\) 有界，遂有
\[
(1+\varepsilon_m(w))^z
=
\exp\!\bigl(z\log(1+\varepsilon_m(w))\bigr)
=
1+O_K(\varrho^m)
\]
在 \(z\in K\) 与 \(w\in X_m\) 上一致成立。于是
\[
d_m(w)^z
=
\Bigl(\frac{2}{\varphi}\Bigr)^{zm}\varphi^{-zw_m}
\Bigl(1+O_K(\varrho^m)\Bigr).
\]
把 \(X_m\) 按末位分拆为
\[
A_0^{(m)}:=\{w\in X_m:w_m=0\},
\qquad
A_1^{(m)}:=\{w\in X_m:w_m=1\},
\]
并用
\[
|A_0^{(m)}|=F_{m+1},
\qquad
|A_1^{(m)}|=F_m,
\]
便得到
\[
S_m(z)
=
\Bigl(\frac{2}{\varphi}\Bigr)^{zm}
\Bigl(F_{m+1}+\varphi^{-z}F_m\Bigr)
\Bigl(1+O_K(\varrho^m)\Bigr).
\]
再乘以 \(\varphi^{-m}(\varphi/2)^{zm}\)，并用 Binet 渐近
\[
\varphi^{-m}F_{m+1}\to \frac{\varphi}{\sqrt5},
\qquad
\varphi^{-m}F_m\to \frac{1}{\sqrt5},
\]
即可得到
\[
\mathcal A_m(z)\to \frac{\varphi+\varphi^{-z}}{\sqrt5}
\]
且收敛在每个紧集上一致。极限函数由显式公式显然为整函数。证毕。
\end{proof}

\begin{corollary}[全实轴仿射压力与单谱塌缩]\label{cor:xi-time-part70e-affine-pressure-monofractal}
\leanverified{paper\_xi\_time\_part70e\_affine\_pressure\_monofractal}
记
\[
\mu_m(w):=\frac{d_m^{\mathrm{bin}}(w)}{2^m},
\qquad
\Xi_m(t):=\sum_{w\in X_m}\mu_m(w)^t
\qquad (t\in\RR).
\]
则极限
\[
\tau_\mu(t):=\lim_{m\to\infty}\frac1m\log \Xi_m(t)
\]
对全部 \(t\in\RR\) 存在，并且满足闭式
\[
\tau_\mu(t)=(1-t)\log\varphi.
\]
因而 \(\tau_\mu\) 在全实轴上仿射，其 Legendre--Fenchel 变换退化为单点谱。等价地，局部信息密度满足一致塌缩
\[
-\frac1m\log \mu_m(w)=\log\varphi+O(m^{-1})
\qquad (w\in X_m),
\]
其中误差在 \(w\) 上一致；并且相对于
\[
\log|X_m|=(m+2)\log\varphi-\frac12\log 5+o(1)
\]
的归一化局部维数亦有
\[
-\frac{\log\mu_m(w)}{\log|X_m|}=1+O(m^{-1})
\qquad (w\in X_m),
\]
其中误差同样在 \(w\) 上一致。
\end{corollary}
\begin{proof}
由定理 \ref{thm:xi-time-part70e-binfold-entire-amplitude-law} 在实轴上的结论，
\[
\Xi_m(t)=2^{-tm}S_m(t)
=
\varphi^{-tm}\Bigl(F_{m+1}+\varphi^{-t}F_m\Bigr)\Bigl(1+O_t(\varrho^m)\Bigr).
\]
又 \(F_{m+1}+\varphi^{-t}F_m=\Theta_t(\varphi^m)\)，故
\[
\frac1m\log\Xi_m(t)\to (1-t)\log\varphi.
\]
这就得到 \(\tau_\mu(t)\) 的显式公式。由于 \(\tau_\mu\) 在全轴上仿射，其 Legendre--Fenchel 变换只能在唯一斜率处非平凡，故相应多重分形谱退化为单点。

另一方面，由定理 \ref{thm:fold-bin-two-state-asymptotic}，
\[
\mu_m(w)=\frac{d_m^{\mathrm{bin}}(w)}{2^m}
=
\varphi^{-m-w_m}\bigl(1+O(\varrho^m)\bigr)
\]
在 \(w\in X_m\) 上一致成立。因此
\[
-\log\mu_m(w)=m\log\varphi+w_m\log\varphi+O(\varrho^m)
\]
一致成立，两边除以 \(m\) 即得
\[
-\frac1m\log\mu_m(w)=\log\varphi+O(m^{-1})
\]
在 \(w\) 上一致。再用 \(|X_m|=F_{m+2}\) 的 Binet 渐近
\[
\log|X_m|=(m+2)\log\varphi-\frac12\log5+o(1),
\]
即可得到最后一式。证毕。
\end{proof}

\begin{theorem}[规范化 Rényi 极限曲线与负奇点 Bernoulli 采样]\label{thm:xi-time-part70e-renyi-curve-negative-odd-sampling}
\leanverified{paper\_xi\_time\_part70e\_renyi\_curve\_negative\_odd\_sampling}
对 \(\alpha>0\)，记 \(D_\alpha^\infty\) 为定理 \ref{thm:fold-bin-renyi-divergence-limit} 中的极限 Rényi 散度，并定义规范化曲线
\[
\mathcal R(\alpha):=
5^{(\alpha-1)/2}\varphi^{2-2\alpha}
\exp\!\bigl((\alpha-1)D_\alpha^\infty\bigr).
\]
则对全部 \(\alpha>0\) 都有严格恒等式
\[
\mathcal R(\alpha)=\frac{\varphi+\varphi^{-\alpha}}{\sqrt5}=\mathcal A(\alpha),
\]
其中 \(\mathcal A\) 为定理 \ref{thm:xi-time-part70e-binfold-entire-amplitude-law} 的整振幅。因而 \(\mathcal R\) 由正实轴数据唯一延拓为整函数，并与 \(\mathcal A\) 完全一致。进一步，对任意固定整数 \(R\ge 1\)，有
\[
\log C_m-\log(2\pi)
=
\frac{\sqrt5}{\varphi^2}
\sum_{r=1}^{R}
\frac{B_{2r}}{r(2r-1)}
\mathcal R\bigl(-(2r-1)\bigr)\,
\varrho^{(2r-1)m}
+
O\!\bigl(\varrho^{(2R+1)m}\bigr).
\]
换言之，fold-gauge 常数的尾部模态恰由规范化 Rényi 极限曲线在负奇整数点的取值经 Bernoulli 变换生成。
\end{theorem}
\begin{proof}
由定理 \ref{thm:fold-bin-renyi-divergence-limit}，当 \(\alpha\neq 1\) 时有
\[
D_\alpha^\infty
=
\frac{(2\alpha-2)\log\varphi+\log(\varphi+\varphi^{-\alpha})-\alpha\log\sqrt5}{\alpha-1}.
\]
故
\[
\exp\!\bigl((\alpha-1)D_\alpha^\infty\bigr)
=
\frac{\varphi^{2\alpha-2}}{5^{\alpha/2}}
\bigl(\varphi+\varphi^{-\alpha}\bigr)
=
\frac{\varphi^{2\alpha-2}}{5^{(\alpha-1)/2}}
\cdot
\frac{\varphi+\varphi^{-\alpha}}{\sqrt5}.
\]
两边再乘以 \(5^{(\alpha-1)/2}\varphi^{2-2\alpha}\)，便得到
\[
\mathcal R(\alpha)=\frac{\varphi+\varphi^{-\alpha}}{\sqrt5}.
\]
而当 \(\alpha=1\) 时，左端按定义等于 \(1\)，右端则满足
\[
\frac{\varphi+\varphi^{-1}}{\sqrt5}=1,
\]
故同样成立。于是 \(\mathcal R\) 在正实轴上与显式函数 \((\varphi+\varphi^{-z})/\sqrt5\) 一致，因此其唯一整延拓正是定理 \ref{thm:xi-time-part70e-binfold-entire-amplitude-law} 中的 \(\mathcal A\)。

另一方面，由定理 \ref{thm:fold-bin-gauge-constant-stirling-bernoulli-hierarchy}，
\[
\log C_m-\log(2\pi)
=
\sum_{r=1}^{R}
\frac{B_{2r}}{r(2r-1)}
\bigl(\varphi^{-1}+\varphi^{2r-3}\bigr)\varrho^{(2r-1)m}
+O\!\bigl(\varrho^{(2R+1)m}\bigr).
\]
而由刚刚得到的闭式，对每个 \(r\ge 1\) 都有
\[
\frac{\sqrt5}{\varphi^2}\mathcal R\bigl(-(2r-1)\bigr)
=
\frac{\sqrt5}{\varphi^2}\cdot
\frac{\varphi+\varphi^{2r-1}}{\sqrt5}
=
\varphi^{-1}+\varphi^{2r-3}.
\]
代回即可得到所示采样公式。证毕。
\end{proof}

\begin{theorem}[零和有限窗对数观测的热力学全决定性]\label{thm:xi-time-part70e-zero-sum-window-thermo-determinacy}
\leanverified{paper\_xi\_time\_part70e\_zero\_sum\_window\_thermo\_determinacy}
令移位算子 \(T\) 作用于序列 \(f=(f_m)_{m\ge 1}\) 为
\[
(Tf)_m:=f_{m+1}.
\]
取任意多项式
\[
Q(\lambda)=\sum_{j=0}^{J}c_j\lambda^j
\]
满足零和条件
\[
Q(1)=\sum_{j=0}^{J}c_j=0.
\]
则对任意固定整数 \(R\ge 1\)，有
\[
Q(T)\log C_m
=
\frac{\sqrt5}{\varphi^2}
\sum_{r=1}^{R}
\frac{B_{2r}}{r(2r-1)}
\mathcal R\bigl(-(2r-1)\bigr)\,
Q\!\bigl(\varrho^{2r-1}\bigr)\,
\varrho^{(2r-1)m}
+
O\!\bigl(\varrho^{(2R+1)m}\bigr).
\]
因此，一切有限窗、零和的对数观测量，例如
\[
\log\frac{C_{m+1}}{C_m},
\qquad
\log C_{m+2}-2\log C_{m+1}+\log C_m,
\]
在渐近层面都完全由 \(\mathcal R\) 决定；其中规范常数 \(\log(2\pi)\) 严格消失。
\end{theorem}
\begin{proof}
由定理 \ref{thm:xi-time-part70e-renyi-curve-negative-odd-sampling}，
\[
\log C_m
=
\log(2\pi)
+
\frac{\sqrt5}{\varphi^2}
\sum_{r=1}^{R}
\frac{B_{2r}}{r(2r-1)}
\mathcal R\bigl(-(2r-1)\bigr)\,
\varrho^{(2r-1)m}
+
O\!\bigl(\varrho^{(2R+1)m}\bigr).
\]
对上式施加算子 \(Q(T)\)。首先
\[
Q(T)\log(2\pi)=Q(1)\log(2\pi)=0.
\]
其次，对任意 \(r\ge 1\) 都有
\[
T\bigl(\varrho^{(2r-1)m}\bigr)=\varrho^{2r-1}\varrho^{(2r-1)m},
\]
故
\[
Q(T)\bigl(\varrho^{(2r-1)m}\bigr)
=
Q\!\bigl(\varrho^{2r-1}\bigr)\varrho^{(2r-1)m}.
\]
逐项代回即可得到所示公式。证毕。
\end{proof}

\begin{theorem}[热力学 germ 的一维 Stokes 归一化刚性]\label{thm:xi-time-part70e-stokes-normalization-rigidity}
设 \((G_m)_{m\ge 1}\) 为一列正数，并假定存在常数 \(c\in\RR\)，使得对每个固定整数 \(R\ge 1\) 都有
\[
\log G_m
=
c+
\frac{\sqrt5}{\varphi^2}
\sum_{r=1}^{R}
\frac{B_{2r}}{r(2r-1)}
\mathcal R\bigl(-(2r-1)\bigr)\,
\varrho^{(2r-1)m}
+
O\!\bigl(\varrho^{(2R+1)m}\bigr).
\]
则
\[
\log G_m-\log C_m
=
c-\log(2\pi)+O\!\bigl(\varrho^{(2R+1)m}\bigr)
\qquad (m\to\infty),
\]
从而
\[
\frac{G_m}{C_m}\longrightarrow e^{\,c-\log(2\pi)}.
\]
亦即：一旦整函数 \(\mathcal R=\mathcal A\) 被固定，整个 fold-gauge germ 在全部归一化无关方向上都已刚性化；剩余自由度严格只保留一个乘法常数，而本文中的 \(2\pi\) 正是该唯一自由度的规范取值。
\leanverified{paper\_xi\_time\_part70e\_stokes\_normalization\_rigidity}
\end{theorem}
\begin{proof}
由定理 \ref{thm:xi-time-part70e-renyi-curve-negative-odd-sampling}，
\[
\log C_m
=
\log(2\pi)+
\frac{\sqrt5}{\varphi^2}
\sum_{r=1}^{R}
\frac{B_{2r}}{r(2r-1)}
\mathcal R\bigl(-(2r-1)\bigr)\,
\varrho^{(2r-1)m}
+
O\!\bigl(\varrho^{(2R+1)m}\bigr).
\]
将其与 \(\log G_m\) 的假设展开相减，级数主项逐项抵消，即得
\[
\log G_m-\log C_m
=
c-\log(2\pi)+O\!\bigl(\varrho^{(2R+1)m}\bigr).
\]
取 \(R=1\) 并令 \(m\to\infty\)，得到
\[
\log G_m-\log C_m\to c-\log(2\pi).
\]
指数化即得所示比值极限。证毕。
\end{proof}

\begin{conjecture}[连续射影谱的精确闭式与整振幅插值]\label{conj:xi-time-part70e-projective-pressure-entire-amplitude}
沿用附录 \ref{app:pom-projective-pressure-operator} 的射影转移算子 \(\mathcal L_t\)（\(t\ge 1\)）。则其主谱半径满足精确闭式
\[
\rho(\mathcal L_t)=2^t\varphi^{1-t}
\qquad (t\ge 1),
\]
等价地，连续归一化压力
\[
\log\rho(\mathcal L_t)-t\log 2=(1-t)\log\varphi
\]
在整个半轴上恰为一条仿射直线。进一步，存在一种自然归一化，使得主特征数据在 \(t\)-轴上的解析延拓恰给出定理 \ref{thm:xi-time-part70e-binfold-entire-amplitude-law} 的整振幅
\[
\mathcal A(t)=\frac{\varphi+\varphi^{-t}}{\sqrt5}.
\]
若该猜想成立，则定理 \ref{thm:xi-time-part70e-renyi-curve-negative-odd-sampling} 中的尾部采样公式便不再只是 bin-fold moment asymptotics 的结果，而可被提升为连续射影主谱在负奇整数点的 Bernoulli 采样律。
\end{conjecture}
\begin{remark}
该猜想与现有证据链完全同向。首先，定理 \ref{thm:pom-real-q-pressure-analytic-interpolation} 已把整数 moment-kernel 塔嵌入一条实解析压力曲线；定理 \ref{thm:pom-projective-operator-degenerates-to-moment-kernel} 与定理 \ref{thm:xi-projective-moment-tower-anchors-normalized-pressure} 则说明，整数节点上的主谱数据正是离散碰撞谱 \(r_q\)。其次，命题 \ref{prop:fold-bin-power-sum-asymptotic} 已经给出
\[
r_q=2^q\varphi^{1-q}
\qquad (q\in\ZZ_{\ge 2}),
\]
而定理 \ref{thm:xi-time-part70e-binfold-entire-amplitude-law} 与定理 \ref{thm:xi-time-part70e-renyi-curve-negative-odd-sampling} 又表明，同一条离散谱线在复参数方向上伴随一个刚性的整振幅 \(\mathcal A\)，并且该振幅的负奇采样恰重建 \(\log C_m\) 的 Bernoulli 奇层。因而，从离散 moment-kernel、Rényi 极限曲线与 fold-gauge 尾部采样三侧同时观察，全部可见证据都指向同一连续谱公式。
\end{remark}

\noindent
由此得到的图景是：本文框架中的 \(\pi\)-机制并不承载于压力函数的曲率，而承载于一条全轴仿射的热力学骨架之上的整函数振幅及其负奇采样。换言之，指数率本身已经被压平为单一直线；真正保留归一化信息的，是该直线上方的整振幅 \(\mathcal A=\mathcal R\)，以及由 Bernoulli 变换读出的唯一 Stokes 型常数 \(2\pi\)。于是，bin-fold 的复矩谱、Rényi 极限曲线、有限窗消模器与 fold-gauge 常数 germ 不再是并列现象，而被收束为同一条解析主线。


\paragraph{零余类有限交截、nerve 的 flag 刚性、加权 clique--Euler 并集公式与共振整函数系数刚性}
在零点余类刚性、约数驱动的谱零点余类并结构、零点削减碰撞上界，以及共振梯整函数的 Laguerre--P\'olya 闭包、实零点结构与截断误差已经建立之后，还可继续沿同一条算术--解析主线抽取出下列七条结果。它们分别刻画零余类族任意有限交截的广义中国剩余判据、零点余类覆盖的 Helly--2 刚性与 nerve 的 flag 闭式、谱零点集合的精确加权 clique--Euler 公式、碰撞上界的约数兼容图闭式、共振整函数 \(\mathcal F_\varphi\) 的对数 Taylor 级数与倒数零矩、其偶系数的显式 Newton 递推与严格对数凹性，以及截断深度对相对误差预算的可审计控制。

\begin{theorem}[零余类族的任意有限交截由两两兼容性完全控制]\label{thm:xi-time-part71-zero-coset-finite-intersection-pairwise-compatibility}
设 \(M\) 为偶整数，且
\[
g_1,\dots,g_r\mid \frac{M}{2}.
\]
对每个 \(g\mid M/2\)，定义
\[
S_g(M):=
\left\{
k\in \ZZ/(M\ZZ):
\ k\equiv \frac{M}{2g}\pmod{M/g}
\right\}.
\]
则对任意有限族 \(\{g_1,\dots,g_r\}\)，都有
\[
\bigcap_{j=1}^{r}S_{g_j}(M)\neq\varnothing
\iff
2\gcd(g_i,g_j)\mid (g_j-g_i)
\qquad
(1\le i<j\le r).
\]
并且一旦交截非空，就有精确基数公式
\[
\left|
\bigcap_{j=1}^{r}S_{g_j}(M)
\right|
=
\gcd(g_1,\dots,g_r).
\]
\end{theorem}
\begin{proof}
由定理 \ref{thm:fold-zero-coset-rigidity}，
\[
S_{g_j}(M)
=
\left\{
k\in \ZZ/(M\ZZ):
\ k\equiv \frac{M}{2g_j}\pmod{M/g_j}
\right\}
\qquad (1\le j\le r).
\]
因此
\[
\bigcap_{j=1}^{r}S_{g_j}(M)\neq\varnothing
\]
当且仅当同余组
\[
k\equiv \frac{M}{2g_j}\pmod{M/g_j}
\qquad (j=1,\dots,r)
\]
可解。广义中国剩余定理断言：这一线性同余组可解，当且仅当任意两式相容。另一方面，引理 \ref{lem:fold-zero-coset-intersection} 已给出两式相容的充要条件恰为
\[
2\gcd(g_i,g_j)\mid (g_j-g_i).
\]
于是得到第一条等价关系。

现设该系统可解。由于解集模数为
\[
\operatorname{lcm}\!\left(\frac{M}{g_1},\dots,\frac{M}{g_r}\right)
=
\frac{M}{\gcd(g_1,\dots,g_r)},
\]
故在模 \(M\) 的意义下，解恰形成一个单一同余类的提升，其基数为
\[
\frac{M}{M/\gcd(g_1,\dots,g_r)}
=
\gcd(g_1,\dots,g_r).
\]
证毕。
\end{proof}

\begin{corollary}[零点余类覆盖的 Helly--2 刚性与 nerve 的 flag 闭式]\label{cor:xi-time-part71-zero-coset-nerve-flag}
沿用定理 \ref{thm:xi-time-part71-zero-coset-finite-intersection-pairwise-compatibility} 的记号。任取有限集合
\[
\mathcal G\subseteq
\left\{
g:\ g\mid \frac{M}{2}
\right\},
\]
并考虑覆盖
\[
\mathcal U_{\mathcal G}:=\{S_g(M):\ g\in\mathcal G\}.
\]
定义其 nerve 为
\[
\mathcal N(\mathcal U_{\mathcal G})
:=
\left\{
\sigma\subseteq \mathcal G:\ \bigcap_{g\in \sigma}S_g(M)\neq\varnothing
\right\}.
\]
再在顶点集 \(\mathcal G\) 上定义算术兼容图 \(\Gamma_{\mathcal G}\)：对不同的 \(g,h\in\mathcal G\)，规定
\[
g\sim h
\iff
2\gcd(g,h)\mid (h-g).
\]
则 \(\mathcal N(\mathcal U_{\mathcal G})\) 恰为 \(\Gamma_{\mathcal G}\) 的 clique 复形。特别地，\(\mathcal N(\mathcal U_{\mathcal G})\) 是一个 flag 复形：任意顶点集若两两相邻，则自动张成一个单纯形。
\end{corollary}
\begin{proof}
由定理 \ref{thm:xi-time-part71-zero-coset-finite-intersection-pairwise-compatibility}，对任意有限子集
\[
\sigma\subseteq \mathcal G
\]
都有
\[
\bigcap_{g\in \sigma}S_g(M)\neq\varnothing
\iff
2\gcd(g,h)\mid (h-g)
\qquad
(\forall\, g\neq h\in \sigma).
\]
右侧正表明 \(\sigma\) 在图 \(\Gamma_{\mathcal G}\) 中形成一个 clique。故
\[
\sigma\in \mathcal N(\mathcal U_{\mathcal G})
\iff
\sigma\in \mathrm{Cliq}(\Gamma_{\mathcal G}),
\]
即 nerve 恰为该兼容图的 clique 复形。于是它自动是 flag 复形。证毕。
\end{proof}


\begin{theorem}[谱零点集合的精确加权 clique--Euler 公式]\label{thm:xi-time-part71-zero-spectrum-weighted-clique-euler}
\leanverified{paper\_xi\_time\_part71\_zero\_spectrum\_weighted\_clique\_euler}
设
\[
M:=F_{m+2},
\qquad
\mathcal Z_m:=\{k\in \ZZ/(M\ZZ):\ \widehat d_m(k)=0\},
\]
并假设 \(M\) 为偶数。定义
\[
\mathcal D_m:=
\left\{
d:\ d\mid (m+2),\ d<m+2,\ F_d\mid \frac{M}{2}
\right\}.
\]
在顶点集 \(\mathcal D_m\) 上定义兼容图 \(\Gamma_m\)：对不同顶点 \(d,e\in\mathcal D_m\)，规定
\[
d\sim e
\iff
2F_{\gcd(d,e)}\mid (F_e-F_d).
\]
若记 \(\mathrm{Cliq}(\Gamma_m)\) 为 \(\Gamma_m\) 的全部非空 clique，且对每个 clique \(C\) 记
\[
\gcd(C):=\gcd\{d:\ d\in C\},
\]
则谱零点集合满足严格公式
\[
|\mathcal Z_m|
=
\sum_{\varnothing\neq C\in \mathrm{Cliq}(\Gamma_m)}
(-1)^{|C|+1}F_{\gcd(C)}.
\]
\end{theorem}
\begin{proof}
由推论 \ref{cor:fold-zero-coset-union}，当 \(M\) 为偶数时，
\[
\mathcal Z_m
=
\bigcup_{\substack{1\le j\le m\\ g_j\mid M/2}}S_{g_j}(M),
\qquad
g_j=\gcd(F_{j+1},M)=F_{\gcd(j+1,m+2)}.
\]
去重之后，恰可写成
\[
\mathcal Z_m=\bigcup_{d\in\mathcal D_m}S_{F_d}(M).
\]

现取任意非空子集 \(I=\{d_1,\dots,d_r\}\subseteq \mathcal D_m\)。由定理 \ref{thm:xi-time-part71-zero-coset-finite-intersection-pairwise-compatibility}，
\[
\bigcap_{d\in I}S_{F_d}(M)\neq\varnothing
\]
当且仅当对所有 \(1\le i<j\le r\) 都有
\[
2\gcd(F_{d_i},F_{d_j})\mid (F_{d_j}-F_{d_i}).
\]
再由 Fibonacci 最大公因子恒等式
\[
\gcd(F_a,F_b)=F_{\gcd(a,b)},
\]
上式等价于
\[
2F_{\gcd(d_i,d_j)}\mid (F_{d_j}-F_{d_i}),
\]
也即 \(I\) 在 \(\Gamma_m\) 中形成一个 clique。

若交截非空，则再由定理 \ref{thm:xi-time-part71-zero-coset-finite-intersection-pairwise-compatibility}，
\[
\left|
\bigcap_{d\in I}S_{F_d}(M)
\right|
=
\gcd(F_{d_1},\dots,F_{d_r})
=
F_{\gcd(d_1,\dots,d_r)}
=
F_{\gcd(I)}.
\]
因此对并集
\[
\mathcal Z_m=\bigcup_{d\in\mathcal D_m}S_{F_d}(M)
\]
施行包含--排斥，便得到
\[
|\mathcal Z_m|
=
\sum_{\varnothing\neq C\in \mathrm{Cliq}(\Gamma_m)}
(-1)^{|C|+1}F_{\gcd(C)}.
\]
证毕。
\end{proof}

\begin{corollary}[碰撞上界可压缩为约数兼容图的加权 Euler 特征]\label{cor:xi-time-part71-collision-upper-bound-weighted-clique-euler}
\leanverified{paper\_xi\_time\_part71\_collision\_upper\_bound\_weighted\_clique\_euler}
沿用定理 \ref{thm:xi-time-part71-zero-spectrum-weighted-clique-euler} 的记号。则碰撞概率满足严格上界
\[
\mathsf{Col}_m
\le
1-\frac{1}{F_{m+2}}
\sum_{\varnothing\neq C\in \mathrm{Cliq}(\Gamma_m)}
(-1)^{|C|+1}F_{\gcd(C)}.
\]
特别地，若 \(\Gamma_m\) 不含三角形，亦即其非空 clique 仅有单点与边，则上界简化为
\[
\mathsf{Col}_m
\le
1-\frac{1}{F_{m+2}}
\left(
\sum_{d\in\mathcal D_m}F_d
-
\sum_{\{d,e\}\in E(\Gamma_m)}F_{\gcd(d,e)}
\right).
\]
\end{corollary}
\begin{proof}
由定理 \ref{thm:fold-collision-zero-reduction}，
\[
\mathsf{Col}_m\le 1-\frac{|\mathcal Z_m|}{F_{m+2}}.
\]
将定理 \ref{thm:xi-time-part71-zero-spectrum-weighted-clique-euler} 的精确公式代入，即得第一式。

若 \(\Gamma_m\) 不含三角形，则全部非空 clique 仅由单点与边组成，包含--排斥求和自动截断为两层：
\[
|\mathcal Z_m|
=
\sum_{d\in\mathcal D_m}F_d
-
\sum_{\{d,e\}\in E(\Gamma_m)}F_{\gcd(d,e)}.
\]
再代回碰撞上界即可。证毕。
\end{proof}

\begin{theorem}[共振整函数的对数 Taylor 级数与倒数零矩闭式]\label{thm:xi-time-part71-resonance-entire-log-taylor-zero-moments}
\leanverified{paper\_xi\_time\_part71\_resonance\_entire\_log\_taylor\_zero\_moments}
沿用定理 \ref{thm:fold-resonance-entire-lp} 的记号，记
\[
\mathcal F_\varphi(z):=\prod_{n=2}^{\infty}\cos(\pi z\varphi^{-n}).
\]
则在最近正零点半径
\[
|z|<\frac{\varphi^2}{2}
\]
内，有严格对数展开
\[
\log \mathcal F_\varphi(z)
=
-\sum_{r=1}^{\infty}\frac{\sigma_r}{r}\,z^{2r},
\qquad
\sigma_r:=
(2^{2r}-1)\zeta(2r)\frac{\varphi^{-2r}}{\varphi^{2r}-1}.
\]
等价地，正零点倒数偶矩满足
\[
\sum_{\substack{\rho\in\mathcal Z(\mathcal F_\varphi)\\ \rho>0}}\rho^{-2r}
=
(2^{2r}-1)\zeta(2r)\frac{\varphi^{-2r}}{\varphi^{2r}-1},
\]
而全零点求和式为
\[
\sum_{\rho\in\mathcal Z(\mathcal F_\varphi)}\rho^{-2r}
=
2(2^{2r}-1)\zeta(2r)\frac{\varphi^{-2r}}{\varphi^{2r}-1}.
\]
\end{theorem}
\begin{proof}
由定理 \ref{thm:fold-resonance-entire-lp}，\(\mathcal F_\varphi\) 为整函数，且
\[
\mathcal Z(\mathcal F_\varphi)
=
\left\{
\left(k+\tfrac12\right)\varphi^n:\ k\in\ZZ,\ n\ge 2
\right\}.
\]
当 \(|z|<\varphi^2/2\) 时，对所有 \(n\ge 2\) 都有
\[
|\pi z\varphi^{-n}|<\frac{\pi}{2},
\]
故可使用经典展开
\[
\log\cos x
=
-\sum_{r=1}^{\infty}\frac{(2^{2r}-1)\zeta(2r)}{r\pi^{2r}}\,x^{2r}
\qquad
(|x|<\pi/2).
\]
逐项代入 \(x=\pi z\varphi^{-n}\) 并交换求和，得到
\[
\log \mathcal F_\varphi(z)
=
\sum_{n=2}^{\infty}\log\cos(\pi z\varphi^{-n})
=
-\sum_{r=1}^{\infty}\frac{(2^{2r}-1)\zeta(2r)}{r}
\left(
\sum_{n=2}^{\infty}\varphi^{-2rn}
\right)z^{2r}.
\]
而
\[
\sum_{n=2}^{\infty}\varphi^{-2rn}
=
\frac{\varphi^{-4r}}{1-\varphi^{-2r}}
=
\frac{\varphi^{-2r}}{\varphi^{2r}-1},
\]
遂得
\[
\log \mathcal F_\varphi(z)
=
-\sum_{r=1}^{\infty}\frac{\sigma_r}{r}\,z^{2r}.
\]

另一方面，由定理 \ref{thm:fold-resonance-entire-lp} 的零点描述，
\[
\mathcal Z(\mathcal F_\varphi)\cap (0,\infty)
=
\left\{
\left(k+\tfrac12\right)\varphi^n:\ k\ge 0,\ n\ge 2
\right\}.
\]
故
\[
\sum_{\substack{\rho\in\mathcal Z(\mathcal F_\varphi)\\ \rho>0}}\rho^{-2r}
=
\sum_{n=2}^{\infty}\varphi^{-2rn}\sum_{k\ge 0}\left(k+\tfrac12\right)^{-2r}.
\]
又有
\[
\sum_{k\ge 0}\left(k+\tfrac12\right)^{-2r}
=
2^{2r}\sum_{k\ge 0}(2k+1)^{-2r}
=
(2^{2r}-1)\zeta(2r),
\]
故正零点公式成立。全零点关于原点成对对称，因此再乘以 \(2\) 即得最后一式。证毕。
\end{proof}

\begin{theorem}[共振整函数偶系数的显式 Newton 递推与严格对数凹性]\label{thm:xi-time-part71-resonance-entire-even-coefficient-newton-logconcavity}
\leanverified{paper\_xi\_time\_part71\_resonance\_entire\_even\_coefficient\_newton\_logconcavity}
将 \(\mathcal F_\varphi\) 写成偶幂级数
\[
\mathcal F_\varphi(z)=\sum_{n\ge 0}a_{2n}z^{2n},
\qquad
a_{2n+1}=0,
\]
并定义
\[
b_n:=(-1)^n a_{2n}
\qquad (n\ge 0).
\]
则对全部 \(n\ge 0\) 都有 \(b_n>0\)，并满足精确 Newton 递推
\[
n\,b_n
=
\sum_{r=1}^{n}(-1)^{r-1}\sigma_r\,b_{n-r}
\qquad (n\ge 1),
\]
其中 \(\sigma_r\) 如定理 \ref{thm:xi-time-part71-resonance-entire-log-taylor-zero-moments} 所示。进一步，序列 \((b_n)_{n\ge 0}\) 严格对数凹：
\[
b_n^2>b_{n-1}b_{n+1}
\qquad (n\ge 1).
\]
因此偶系数满足严格交替符号律
\[
a_0>0,\qquad a_2<0,\qquad a_4>0,\qquad \dots,
\]
且相邻比值 \((b_{n+1}/b_n)_{n\ge 0}\) 严格递减。
\end{theorem}
\begin{proof}
先由定理 \ref{thm:xi-time-part71-resonance-entire-log-taylor-zero-moments} 定义形式幂级数
\[
G(t):=\mathcal F_\varphi(\sqrt t)
=
\sum_{n\ge 0}a_{2n}t^n
=
\sum_{n\ge 0}(-1)^n b_n t^n.
\]
则
\[
\log G(t)
=
-\sum_{r=1}^{\infty}\frac{\sigma_r}{r}t^r.
\]
对 \(t\) 求导得
\[
\frac{G'(t)}{G(t)}
=
-\sum_{r=1}^{\infty}\sigma_r t^{r-1}.
\]
把
\[
G(t)=\sum_{n\ge 0}(-1)^n b_n t^n
\]
代入，并比较 \(t^{n-1}\) 的系数，即得
\[
n\,(-1)^n b_n
=
-\sum_{r=1}^{n}\sigma_r(-1)^{n-r}b_{n-r},
\]
等价于
\[
n\,b_n
=
\sum_{r=1}^{n}(-1)^{r-1}\sigma_r\,b_{n-r}.
\]

接着证明 \(b_n>0\) 与严格对数凹性。由定理 \ref{thm:fold-resonance-entire-lp}，
\(\mathcal F_\varphi\) 的正零点可记为
\[
0<\rho_1<\rho_2<\cdots,
\]
且由定理 \ref{thm:xi-time-part71-resonance-entire-log-taylor-zero-moments} 在 \(r=1\) 时的公式，
\[
\sum_{j\ge 1}\rho_j^{-2}<\infty.
\]
因此
\[
G(t)=\prod_{j=1}^{\infty}\left(1-\frac{t}{\rho_j^2}\right)
\]
在 \(t=0\) 邻域内成立，并且
\[
b_n=e_n(\rho_1^{-2},\rho_2^{-2},\dots)>0
\]
是正数列 \((\rho_j^{-2})_{j\ge 1}\) 的第 \(n\) 个初等对称函数。

对每个 \(N>n\)，定义截断多项式
\[
G_N(t):=\prod_{j=1}^{N}\left(1-\frac{t}{\rho_j^2}\right)
=
\sum_{s=0}^{N}(-1)^s b_s^{(N)}t^s.
\]
则 \(b_s^{(N)}\to b_s\)（\(s\) 固定）且 \(b_s^{(N)}>0\)。Newton 不等式给出
\[
\bigl(b_n^{(N)}\bigr)^2
\ge
\frac{n+1}{n}\frac{N-n+1}{N-n}\,
b_{n-1}^{(N)}b_{n+1}^{(N)}.
\]
令 \(N\to\infty\)，得到
\[
b_n^2
\ge
\frac{n+1}{n}\,b_{n-1}b_{n+1}
>
b_{n-1}b_{n+1}.
\]
这就证明了严格对数凹性。由于 \(b_n>0\)，于是
\[
a_{2n}=(-1)^n b_n
\]
严格交替。再将
\[
b_n^2>b_{n-1}b_{n+1}
\]
两边除以 \(b_{n-1}b_n>0\)，即可得到
\[
\frac{b_{n+1}}{b_n}<\frac{b_n}{b_{n-1}},
\]
故相邻比值严格递减。证毕。
\end{proof}

\begin{corollary}[共振梯截断深度的相对误差预算]\label{cor:xi-time-part71-resonance-truncation-relative-error-budget}
\leanverified{paper\_xi\_time\_part71\_resonance\_truncation\_relative\_error\_budget}
沿用命题 \ref{prop:fold-resonance-truncation-error} 的记号。对任意 \(\varepsilon>0\)，当 \(u=0\) 时有
\[
\frac{C_{\varphi,N}(0)}{C_\varphi(0)}=1
\qquad
(\forall N\ge 2).
\]
若 \(u\neq 0\)、\(N\ge 2\) 且
\[
N\ \ge\
-1+\max\!\left\{
\log_\varphi(2\pi|u|),\
\frac12\log_\varphi\!\left(
\frac{\pi^2u^2}{(1-\varphi^{-2})\log(1+\varepsilon)}
\right)
\right\},
\]
则有相对误差界
\[
1\le \frac{C_{\varphi,N}(u)}{C_\varphi(u)}\le 1+\varepsilon.
\]
\end{corollary}
\begin{proof}
当 \(u=0\) 时，
\[
C_{\varphi,N}(0)=C_\varphi(0)=1,
\]
结论平凡。以下设 \(u\neq 0\)。

由命题 \ref{prop:fold-resonance-truncation-error}，只要满足
\[
\pi|u|\varphi^{-(N+1)}\le \frac12,
\]
就有
\[
C_{\varphi,N}(u)\exp\!\left(
-\frac{\pi^2u^2\varphi^{-2(N+1)}}{1-\varphi^{-2}}
\right)
\le
C_\varphi(u)
\le
C_{\varphi,N}(u).
\]
因此
\[
1
\le
\frac{C_{\varphi,N}(u)}{C_\varphi(u)}
\le
\exp\!\left(
\frac{\pi^2u^2\varphi^{-2(N+1)}}{1-\varphi^{-2}}
\right).
\]
若进一步要求
\[
\frac{\pi^2u^2\varphi^{-2(N+1)}}{1-\varphi^{-2}}
\le
\log(1+\varepsilon),
\]
则立得
\[
\frac{C_{\varphi,N}(u)}{C_\varphi(u)}\le 1+\varepsilon.
\]
最后，把上述两条约束分别对 \(N\) 求解，便得到
\[
N\ \ge\
-1+\max\!\left\{
\log_\varphi(2\pi|u|),\
\frac12\log_\varphi\!\left(
\frac{\pi^2u^2}{(1-\varphi^{-2})\log(1+\varepsilon)}
\right)
\right\}.
\]
证毕。
\end{proof}

\noindent
上述七条结果表明，Fibonacci 模零谱与共振整函数 \(\mathcal F_\varphi\) 已经形成一条严格可逆的算术--解析双轨：一方面，谱零点集合 \(\mathcal Z_m\) 不再只是若干候选余类的松散并，而可被零余类覆盖的 flag nerve、约数兼容图的 clique 结构及其加权 Euler 特征精确决定，并进一步把碰撞上界压缩为有限图论计算；另一方面，\(\mathcal F_\varphi\) 的局部解析则不再停留于 Laguerre--P\'olya 性与 Jensen 多项式全实根性这样的抽象实根性，而会在对数 Taylor 系数、倒数零矩、偶系数 Newton 递推与截断深度预算中外显为完全可审计的 \(\zeta(2r)\)--\(\varphi\) 谱矩闭式。于是，零余类几何、有限交 nerve、碰撞上界与共振常数计算深度便被统一收束到同一条由 Fibonacci 约数格与黄金比例谱矩共同支配的闭合链条之中。

\paragraph{奇倍嵌套链、约数格指标容斥、线性零点增量与六倍窗密度指数率}
在零余类族的 \(\gcd\)-刚性交结构与谱零点集合的加权 clique--Euler 公式已经建立之后，还可继续沿约数格方向抽取出下列四条直接后果。它们分别刻画整除链上的奇倍嵌套--偶倍排斥二分律、以相容指标书写的约数格精确容斥公式、pairwise-incompatible 约数族所导致的线性零点增量下界，以及六倍窗上谱零点相对密度的精确指数率。

\begin{theorem}[零余类沿整除链的奇倍嵌套与偶倍排斥]\label{thm:xi-time-part72-zero-coset-divisibility-nested-disjoint-dichotomy}
\leanverified{paper\_xi\_time\_part72\_zero\_coset\_divisibility\_nested\_disjoint\_dichotomy}
设 \(M\) 为偶整数，并令
\[
g\mid h\mid \frac{M}{2}.
\]
则有完全二分律
\[
S_g(M)\subseteq S_h(M)
\iff
\frac{h}{g}\ \text{为奇数},
\]
\[
S_g(M)\cap S_h(M)=\varnothing
\iff
\frac{h}{g}\ \text{为偶数}.
\]
特别地，对任意整除链
\[
g_1\mid g_2\mid \cdots \mid g_r\mid \frac{M}{2},
\]
若所有相邻商
\[
\frac{g_{i+1}}{g_i}\qquad (1\le i<r)
\]
都为奇数，则
\[
S_{g_1}(M)\subseteq S_{g_2}(M)\subseteq \cdots \subseteq S_{g_r}(M),
\qquad
\left|\bigcup_{i=1}^{r}S_{g_i}(M)\right|=g_r.
\]
反之，只要存在某个 \(i\) 使 \(g_{i+1}/g_i\) 为偶数，便有
\[
S_{g_i}(M)\cap S_{g_{i+1}}(M)=\varnothing.
\]
\end{theorem}
\begin{proof}
由定理 \ref{thm:xi-fold-zero-coset-v2-intersection-rigidity}，
\[
S_g(M)\cap S_h(M)=
\begin{cases}
S_{\gcd(g,h)}(M),& v_2(g)=v_2(h),\\[2mm]
\varnothing,& v_2(g)\neq v_2(h).
\end{cases}
\]
在 \(g\mid h\) 的情形，
\[
\gcd(g,h)=g.
\]
又
\[
v_2(g)=v_2(h)
\iff
v_2(h/g)=0
\iff
\frac{h}{g}\ \text{为奇数}.
\]
因此当 \(h/g\) 为奇数时，
\[
S_g(M)\cap S_h(M)=S_g(M),
\]
从而 \(S_g(M)\subseteq S_h(M)\)；当 \(h/g\) 为偶数时，
\[
v_2(g)\neq v_2(h),
\]
故交集为空。

对整除链逐步应用上述二分律，便得到链式包含关系与相邻偶商导致的互斥结论。若全部相邻商皆为奇数，则整条链单调嵌套，故其并集等于最大项 \(S_{g_r}(M)\)，于是
\[
\left|\bigcup_{i=1}^{r}S_{g_i}(M)\right|=|S_{g_r}(M)|=g_r.
\]
证毕。
\end{proof}

\begin{theorem}[谱零点集合的约数格指标容斥公式]\label{thm:xi-time-part72-zero-spectrum-divisor-lattice-indicator-inclusion-exclusion}
\leanverified{paper\_xi\_time\_part72\_zero\_spectrum\_divisor\_lattice\_indicator\_inclusion\_exclusion}
沿用定理 \ref{thm:xi-time-part71-zero-spectrum-weighted-clique-euler} 的记号，令
\[
\mathcal D_m:=
\left\{
d:\ d\mid (m+2),\ d<m+2,\ F_d\mid \frac{F_{m+2}}{2}
\right\}.
\]
对每个非空子集 \(I\subseteq \mathcal D_m\)，定义
\[
\gcd(I):=\gcd\{d:\ d\in I\},
\]
并置相容指标
\[
\mathcal C(I)
\iff
2F_{\gcd(d,e)}\mid (F_e-F_d)
\qquad
(\forall\, d,e\in I,\ d\neq e).
\]
则谱零点集合满足完全显式的约数格容斥展开
\[
|\mathcal Z_m|
=
\sum_{\varnothing\neq I\subseteq \mathcal D_m}
(-1)^{|I|+1}\mathbf 1_{\mathcal C(I)}\,F_{\gcd(I)}.
\]
换言之，\(\mathcal Z_m\) 的基数可由 \((m+2)\) 的约数格与 Fibonacci 差值的整除判别完全决定，而无需再对频率 \(k\) 作逐点扫描。
\end{theorem}
\begin{proof}
由定理 \ref{thm:xi-time-part71-zero-spectrum-weighted-clique-euler}，
\[
|\mathcal Z_m|
=
\sum_{\varnothing\neq C\in \mathrm{Cliq}(\Gamma_m)}
(-1)^{|C|+1}F_{\gcd(C)},
\]
其中 \(\Gamma_m\) 的边关系定义为
\[
d\sim e
\iff
2F_{\gcd(d,e)}\mid (F_e-F_d).
\]
因此，对任意非空子集 \(I\subseteq \mathcal D_m\)，
\[
I\in \mathrm{Cliq}(\Gamma_m)
\iff
\mathcal C(I).
\]
于是把 clique 求和改写为对全部非空子集的求和，并以指标函数 \(\mathbf 1_{\mathcal C(I)}\) 过滤不相容子集，便得到
\[
|\mathcal Z_m|
=
\sum_{\varnothing\neq I\subseteq \mathcal D_m}
(-1)^{|I|+1}\mathbf 1_{\mathcal C(I)}\,F_{\gcd(I)}.
\]
证毕。
\end{proof}

\begin{corollary}[pairwise-incompatible 约数族给出线性零点增量下界]\label{cor:xi-time-part72-zero-spectrum-pairwise-incompatible-linear-lower-bound}
\leanverified{paper\_xi\_time\_part72\_zero\_spectrum\_pairwise\_incompatible\_linear\_lower\_bound}
沿用定理 \ref{thm:xi-time-part72-zero-spectrum-divisor-lattice-indicator-inclusion-exclusion} 的记号。若子集
\[
\mathcal A\subseteq \mathcal D_m
\]
满足对任意不同的 \(d,e\in\mathcal A\) 都有
\[
2F_{\gcd(d,e)}\nmid (F_e-F_d),
\]
则相应余类族两两不交，从而
\[
|\mathcal Z_m|
\ge
\sum_{d\in\mathcal A}F_d.
\]
特别地，当 \(m+2\equiv 0\pmod 6\) 且
\[
d:=\frac{m+2}{2},
\]
由推论 \ref{cor:fold-zero-half-index-multiple6} 已知
\[
F_d\mid \frac{F_{m+2}}{2},
\qquad
|\mathcal Z_m|\ge F_d.
\]
若再存在 \(e\in\mathcal D_m\) 满足
\[
2F_{\gcd(d,e)}\nmid (F_e-F_d),
\]
则立刻提升为
\[
|\mathcal Z_m|\ge F_d+F_e.
\]
\end{corollary}
\begin{proof}
由定理 \ref{thm:xi-time-part71-zero-coset-finite-intersection-pairwise-compatibility}，对不同的 \(d,e\in \mathcal A\)，条件
\[
2F_{\gcd(d,e)}\nmid (F_e-F_d)
\]
等价于
\[
S_{F_d}(F_{m+2})\cap S_{F_e}(F_{m+2})=\varnothing.
\]
故族
\[
\{S_{F_d}(F_{m+2}):\ d\in\mathcal A\}
\]
两两不交。于是
\[
\left|
\bigcup_{d\in\mathcal A}S_{F_d}(F_{m+2})
\right|
=
\sum_{d\in\mathcal A}|S_{F_d}(F_{m+2})|
=
\sum_{d\in\mathcal A}F_d,
\]
其中最后一步使用定理 \ref{thm:fold-zero-coset-rigidity} 中
\[
|S_g(M)|=g.
\]
再由推论 \ref{cor:fold-zero-coset-union}，
\[
\bigcup_{d\in\mathcal A}S_{F_d}(F_{m+2})
\subseteq
\mathcal Z_m,
\]
遂得所示下界。

最后一段只需把 \(d=(m+2)/2\) 的六倍窗见证余类与一个与之 pairwise-incompatible 的 \(e\) 一并代入上述结论即可。证毕。
\end{proof}

\begin{theorem}[六倍窗上谱零点相对密度的精确指数率]\label{thm:xi-time-part72-zero-density-exact-exponential-rate}
沿用定理 \ref{thm:xi-time-part70b-dark-band-half-order-scale} 的记号。若
\[
m+2\equiv 0\pmod 6,
\qquad
M:=F_{m+2},
\]
则有极限公式
\[
\lim_{\substack{m\to\infty\\ m+2\equiv 0\!\!\!\pmod 6}}
\frac{1}{m}\log\frac{|\mathcal Z_m|}{F_{m+2}}
=
-\frac12\log\varphi.
\]
等价地，
\[
\frac{|\mathcal Z_m|}{F_{m+2}}
=
\exp\!\left(
-\frac{m}{2}\log\varphi+o(m)
\right)
\]
沿六倍窗子序列成立。
\end{theorem}
\begin{proof}
由定理 \ref{thm:xi-time-part70b-dark-band-half-order-scale}，
\[
\frac{1}{m}\log |\mathcal Z_m|
\longrightarrow
\frac12\log\varphi
\]
沿子序列 \(m+2\equiv 0\pmod 6\) 成立。另一方面，Binet 公式给出
\[
\log F_{m+2}=(m+2)\log\varphi+O(1),
\]
故
\[
\frac{1}{m}\log F_{m+2}\longrightarrow \log\varphi.
\]
两式相减即可得到
\[
\frac{1}{m}\log\frac{|\mathcal Z_m|}{F_{m+2}}
\longrightarrow
\frac12\log\varphi-\log\varphi
=
-\frac12\log\varphi.
\]
等价表述只是把极限式改写为指数渐近。证毕。
\end{proof}

\noindent
上述四条结果表明，零余类几何在有限交截判据与加权 clique--Euler 公式之外，还具有一条更细的约数格层级结构：沿整除链，余类并不会出现任意复杂的部分重叠，而只能在“奇商嵌套”与“偶商互斥”两种情形之间二分；一旦把这一二分律与相容指标容斥结合，任何 pairwise-incompatible 约数簇都会立刻转化为 \(\mathcal Z_m\) 的线性可加增量；而在六倍窗上，这些局部增量又与 half-order dark band 的双边界闭合到同一个精确指数率 \(\varphi^{-m/2}\)。于是，谱零点集合的算术几何便不仅可被精确计数，而且可在约数格上作系统的局部增量构造，并在特定同余窗上获得完全锁定的密度衰减主尺度。

\paragraph{奇商偏序下的零余类压缩生成、唯一极小覆盖与精确计数}
在定理 \ref{thm:xi-time-part71-zero-coset-finite-intersection-pairwise-compatibility} 已给出零余类族任意有限交截的完全分类、定理 \ref{thm:xi-time-part72-zero-coset-divisibility-nested-disjoint-dichotomy} 已把整除链上的交叠压缩为奇倍嵌套与偶倍排斥二分律之后，这条主线还可继续向“表示压缩”推进一步：零点集合 \(\mathcal Z_m\) 并不需要由全部候选余类来生成，而可唯一压缩到奇商偏序下的极大元族；相应的精确计数也可从全约数格容斥进一步压缩为该极大生成族上的包含--排斥闭式。

\begin{theorem}[零余类之间的全局包含关系由奇商整除完全刻画]\label{thm:xi-time-part72a-zero-coset-global-inclusion-odd-divisibility}
设 \(M\) 为偶整数，且 \(g,h\mid M/2\)。则
\[
S_g(M)\subseteq S_h(M)
\iff
g\mid h\ \text{且}\ \frac{h}{g}\ \text{为奇数}.
\]
\end{theorem}
\begin{proof}
若 \(g\mid h\) 且 \(h/g\) 为奇数，则结论已由定理 \ref{thm:xi-time-part72-zero-coset-divisibility-nested-disjoint-dichotomy} 的第一条给出。

反之，设 \(S_g(M)\subseteq S_h(M)\)。则
\[
S_g(M)\cap S_h(M)=S_g(M),
\]
因而交集非空。由定理 \ref{thm:xi-time-part71-zero-coset-finite-intersection-pairwise-compatibility} 在 \(r=2\) 的情形，
\[
\left|S_g(M)\cap S_h(M)\right|=\gcd(g,h).
\]
又由定理 \ref{thm:fold-zero-coset-rigidity}，
\[
|S_g(M)|=g.
\]
于是
\[
\gcd(g,h)=|S_g(M)\cap S_h(M)|=|S_g(M)|=g,
\]
从而 \(g\mid h\)。再由交集非空判据，
\[
2\gcd(g,h)\mid (h-g),
\]
即
\[
2g\mid g\left(\frac{h}{g}-1\right).
\]
故 \(2\mid (h/g-1)\)，也就是 \(h/g\) 为奇数。证毕。
\end{proof}

\begin{theorem}[谱零点集合在奇商偏序下的唯一极小生成族]\label{thm:xi-time-part72a-zero-spectrum-maximal-odd-divisibility-compression}
\leanverified{paper\_xi\_time\_part72a\_zero\_spectrum\_maximal\_odd\_divisibility\_compression}
沿用推论 \ref{cor:fold-zero-coset-union} 与推论 \ref{cor:xi-fold-zero-set-v2-semilattice-decomposition} 的记号。记
\[
M:=F_{m+2},
\qquad
\mathcal D_m:=
\left\{
d:\ d\mid (m+2),\ d<m+2,\ F_d\mid \frac{M}{2}
\right\},
\]
\[
\mathcal G_m:=\{F_d:\ d\in \mathcal D_m\}.
\]
则
\[
\mathcal Z_m=\bigcup_{g\in \mathcal G_m}S_g(M).
\]
在 \(\mathcal G_m\) 上定义偏序
\[
g\preceq h
\iff
S_g(M)\subseteq S_h(M),
\]
等价地，由定理 \ref{thm:xi-time-part72a-zero-coset-global-inclusion-odd-divisibility}，
\[
g\preceq h
\iff
g\mid h\ \text{且}\ \frac{h}{g}\ \text{为奇数}.
\]
记 \(\mathcal M_m\subseteq \mathcal G_m\) 为全部 \(\preceq\)-极大元组成的集合。则成立：
\begin{enumerate}
  \item 谱零点集合具有压缩表示
  \[
  \mathcal Z_m=\bigcup_{g\in \mathcal M_m}S_g(M).
  \]
  \item \(\mathcal M_m\) 是满足上式的唯一极小生成族：若 \(\mathcal A\subseteq \mathcal G_m\) 也满足
  \[
  \mathcal Z_m=\bigcup_{g\in \mathcal A}S_g(M),
  \]
  则必有
  \[
  \mathcal M_m\subseteq \mathcal A.
  \]
\end{enumerate}
\end{theorem}
\begin{proof}
第一条只需删除全部非极大元。事实上，若 \(g\in \mathcal G_m\setminus \mathcal M_m\)，则存在 \(h\in \mathcal G_m\) 使 \(g\prec h\)。由定理 \ref{thm:xi-time-part72a-zero-coset-global-inclusion-odd-divisibility}，
\[
S_g(M)\subseteq S_h(M),
\]
故从并集中去掉 \(S_g(M)\) 不改变并集。对所有非极大元依次执行此操作，即得
\[
\mathcal Z_m=\bigcup_{g\in \mathcal M_m}S_g(M).
\]

现证唯一极小性。任取 \(g\in \mathcal M_m\)。若 \(h\in \mathcal G_m\setminus\{g\}\) 且
\[
S_g(M)\cap S_h(M)\neq \varnothing,
\]
则由定理 \ref{thm:xi-time-part71-zero-coset-finite-intersection-pairwise-compatibility}，
\[
S_g(M)\cap S_h(M)=S_{\gcd(g,h)}(M)
\]
且
\[
\gcd(g,h)<g,
\]
否则若 \(\gcd(g,h)=g\)，则 \(g\mid h\)，再由交非空判据可得 \(h/g\) 为奇数，从而 \(g\prec h\)，与 \(g\) 的极大性矛盾。

取 \(S_g(M)\) 的标准代表点
\[
k_g:=\frac{M}{2g}\in S_g(M).
\]
若 \(k_g\in S_h(M)\) 对某个 \(h\neq g\) 成立，则由上式知
\[
k_g\in S_g(M)\cap S_h(M)=S_{\gcd(g,h)}(M)
\]
且 \(d:=\gcd(g,h)<g\)。由于交集非空，\(g/d\) 必为奇数。于是 \(S_d(M)\subseteq S_g(M)\)，并且按定理 \ref{thm:fold-zero-coset-rigidity} 的显式描述，\(S_d(M)\) 在 \(S_g(M)\) 内对应于模 \(g/d\) 的单一余类
\[
t\equiv \frac{g/d-1}{2}\pmod{g/d}.
\]
而点 \(k_g\) 对应于参数 \(t=0\)。因为 \(g/d>1\) 且为奇数，
\[
\frac{g/d-1}{2}\not\equiv 0\pmod{g/d},
\]
故 \(k_g\notin S_d(M)\)，矛盾。

因此
\[
k_g\notin \bigcup_{h\in \mathcal G_m\setminus\{g\}}S_h(M),
\]
从而
\[
S_g(M)\nsubseteq \bigcup_{h\in \mathcal G_m\setminus\{g\}}S_h(M).
\]
于是任何生成整个 \(\mathcal Z_m\) 的子族都必须包含 \(g\)。由于 \(g\in \mathcal M_m\) 任意，遂得
\[
\mathcal M_m\subseteq \mathcal A.
\]
这就证明了唯一极小性。证毕。
\end{proof}

\begin{theorem}[极大奇商生成族上的谱零点集合精确容斥公式]\label{thm:xi-time-part72a-zero-spectrum-maximal-family-inclusion-exclusion}
\leanverified{paper\_xi\_time\_part72a\_zero\_spectrum\_maximal\_family\_inclusion\_exclusion}
沿用定理 \ref{thm:xi-time-part72a-zero-spectrum-maximal-odd-divisibility-compression} 的记号，写
\[
\mathcal M_m=\{g_1,\dots,g_r\}.
\]
则谱零点集合的基数满足精确公式
\[
|\mathcal Z_m|
=
\sum_{\varnothing\neq I\subseteq \{1,\dots,r\}}
(-1)^{|I|+1}
\mathbf 1_{\mathrm{comp}(I)}
\gcd(g_i:\ i\in I),
\]
其中相容指标 \(\mathbf 1_{\mathrm{comp}(I)}\) 定义为：当且仅当对全部 \(i,j\in I\) 都有
\[
2\gcd(g_i,g_j)\mid (g_j-g_i)
\]
时，\(\mathbf 1_{\mathrm{comp}(I)}=1\)；否则 \(\mathbf 1_{\mathrm{comp}(I)}=0\)。
\end{theorem}
\begin{proof}
由定理 \ref{thm:xi-time-part72a-zero-spectrum-maximal-odd-divisibility-compression}，
\[
\mathcal Z_m=\bigcup_{\nu=1}^{r}S_{g_\nu}(M).
\]
对该有限并应用包含--排斥，
\[
|\mathcal Z_m|
=
\sum_{\varnothing\neq I\subseteq \{1,\dots,r\}}
(-1)^{|I|+1}
\left|
\bigcap_{i\in I}S_{g_i}(M)
\right|.
\]
再由定理 \ref{thm:xi-time-part71-zero-coset-finite-intersection-pairwise-compatibility}，交截非空当且仅当上述两两相容条件成立，而一旦非空就有
\[
\left|
\bigcap_{i\in I}S_{g_i}(M)
\right|
=
\gcd(g_i:\ i\in I).
\]
代回即得所示公式。证毕。
\end{proof}

\begin{corollary}[极大生成族 pairwise-incompatible 时的精确直和公式]\label{cor:xi-time-part72a-zero-spectrum-maximal-family-pairwise-incompatible-exact-sum}
\leanverified{paper\_xi\_time\_part72a\_zero\_spectrum\_maximal\_family\_pairwise\_incompatible\_exact\_sum}
沿用定理 \ref{thm:xi-time-part72a-zero-spectrum-maximal-odd-divisibility-compression} 的记号。若极大生成族 \(\mathcal M_m=\{g_1,\dots,g_r\}\) 满足任意两元都不相容，即
\[
2\gcd(g_i,g_j)\nmid (g_j-g_i)
\qquad (i\neq j),
\]
则对应余类两两不交，从而
\[
|\mathcal Z_m|=\sum_{\nu=1}^{r}g_\nu.
\]
\end{corollary}
\begin{proof}
由定理 \ref{thm:xi-time-part71-zero-coset-finite-intersection-pairwise-compatibility}，上述条件等价于
\[
S_{g_i}(M)\cap S_{g_j}(M)=\varnothing
\qquad (i\neq j).
\]
故
\[
\mathcal Z_m=\bigsqcup_{\nu=1}^{r}S_{g_\nu}(M).
\]
再由定理 \ref{thm:fold-zero-coset-rigidity} 的 \(|S_g(M)|=g\)，可得
\[
|\mathcal Z_m|
=
\sum_{\nu=1}^{r}|S_{g_\nu}(M)|
=
\sum_{\nu=1}^{r}g_\nu.
\]
证毕。
\end{proof}

\noindent
上述结果表明，零余类族的算术几何还具有一层此前未被完全显化的“压缩表示”刚性：有限交截虽由全体候选余类的 pairwise 相容性控制，但真正生成 \(\mathcal Z_m\) 的并不需要整张约数图，而只需要奇商偏序下的极大元族；这一压缩不仅是存在性的，而且是唯一极小的。相应地，谱零点计数也可从全约数格的指标容斥，进一步压缩为极大生成族上的精确包含--排斥；而在 pairwise-incompatible 的情形下，全部交叠复杂性彻底消失，精确计数退化为一条纯粹的 Fibonacci 权重直和公式。于是，零余类交叠、包含、压缩表示与 exact counting 四者就在同一条 odd-divisibility 主线上获得了闭合。
\paragraph{pairwise-incompatible 零块族的同步信息削减}
在推论 \ref{cor:xi-time-part72-zero-spectrum-pairwise-incompatible-linear-lower-bound} 已把 pairwise-incompatible 约数族转化为 \(|\mathcal Z_m|\) 的线性可加增量，而定理 \ref{thm:xi-fold-zero-dyadic-tower-synchronous-information-gap-improvement} 已在 dyadic 塔特例下把这一增量同步传递到碰撞、KL、峰值与最大纤维偏离之后，可以看到：这一同步机制本身并不依赖 dyadic 结构，而只依赖 pairwise incompatibility 所保证的零块互不相交。

\begin{theorem}[pairwise-incompatible Fibonacci 约数族触发的同步信息削减]\label{thm:xi-time-part72b-pairwise-incompatible-synchronous-information-gap}
沿用定理 \ref{thm:xi-time-part72-zero-spectrum-divisor-lattice-indicator-inclusion-exclusion} 的记号，并沿用定理 \ref{thm:fold-zero-uncertainty} 的记号。设
\[
\mathcal A\subseteq \mathcal D_m
\]
满足对任意不同的 \(d,e\in\mathcal A\)，都有
\[
2F_{\gcd(d,e)}\nmid (F_e-F_d).
\]
记
\[
L_{\mathcal A}:=\sum_{d\in\mathcal A}F_d.
\]
则
\[
|\mathcal Z_m|\ge L_{\mathcal A}.
\]
因而同时成立
\[
\mathsf{Col}_m
\le
1-\frac{L_{\mathcal A}}{F_{m+2}},
\qquad
D_{\mathrm{KL}}(p_m\|u_m)
\le
\log\!\bigl(F_{m+2}-L_{\mathcal A}\bigr),
\]
\[
S_m\le 2^m\sqrt{F_{m+2}-1-L_{\mathcal A}},
\]
\[
M_m-\bar d_m
\le
2^m\sqrt{\frac{F_{m+2}-1-L_{\mathcal A}}{F_{m+2}}}.
\]
\end{theorem}
\begin{proof}
由推论 \ref{cor:xi-time-part72-zero-spectrum-pairwise-incompatible-linear-lower-bound}，
\[
|\mathcal Z_m|\ge \sum_{d\in\mathcal A}F_d=L_{\mathcal A}.
\]
将该下界分别代入定理 \ref{thm:fold-collision-zero-reduction} 与定理
\ref{thm:fold-zero-uncertainty} 的第 \(2\)、\(3\)、\(4\) 条，便得到所示碰撞上界、KL 上界、Fourier 峰值上界以及最大纤维偏离上界。证毕。
\end{proof}

\noindent
因此，定理 \ref{thm:xi-fold-zero-dyadic-tower-synchronous-information-gap-improvement} 只是这一一般机制在 dyadic 塔族上的特例。真正驱动同步削减的并非 \(v_2\)-塔本身，而是 Fibonacci 零块之间的 pairwise incompatibility：只要若干约数对应的零块能够同时并入 \(\mathcal Z_m\) 而不发生交叠，它们所贡献的零点增量便会无损地传递到碰撞、相对熵、峰值振幅与最大纤维偏离这四个量上。


\paragraph{二维原子层、非刚性迹系统、激活零尺度根与零温四态核心的三尺度分裂}
在真实输入 \(40\) 状态核的二维加权 \(\zeta\) 分母闭式、一参数输出位势的激活分支，以及零温四态地基层已经建立之后，还可把原子层、非刚性残余谱、局部激活层与零温核心层进一步压缩为下列若干条彼此衔接的结构后果。前两条仍对应对原子因子的中心投影式相位盲读出；紧随其后的两条把去除 rigid branches 后的全部非刚性迹压缩为单一八次 core 因子的 Newton--Cayley 闭系统，并进一步抽出 \(v=1\) 超曲面与 \((u,v)=(0,1)\) 尖点处的谱秩塌缩；随后的几条则转入相位敏感的纤维内读出，分别刻画 \(u=1\) 处的极点阶、Abel 有限部平移、激活分支的横截性、去支后的局部无碰撞窗口以及正实轴上的局部谱壁垒；最后再把局部奇性独占与零温 Perron 幅度统一到同一条谱分裂链中。

\begin{theorem}[二维加权 \texorpdfstring{$\zeta$}{zeta} 的原子因子在 ghost 层对碰撞权完全失明]\label{thm:xi-time-part73-twoparam-atomic-ghost-u-blindness}
沿用定理 \ref{thm:real-input-40-det-uv-2plus8} 与推论 \ref{cor:real-input-40-zeta-uv-sqrtv-eigs} 的记号，并记
\[
\zeta(z;u,v)
=
\frac{1}{\Delta(z;u,v)}
=
\frac{1}{1-vz^2}\,\zeta_{\mathrm{core}}(z;u,v),
\qquad
\zeta_{\mathrm{core}}(z;u,v):=-P_8(z;u,v)^{-1}.
\]
定义原子 ghost 系数 \(a_n^{\mathrm{atom}}(u,v)\) 为
\[
\log\frac{1}{1-vz^2}
=
\sum_{n\ge 1}\frac{a_n^{\mathrm{atom}}(u,v)}{n}z^n.
\]
则对任意 \(n\ge 1\) 都有精确公式
\[
\boxed{\;
a_n^{\mathrm{atom}}(u,v)
=
\begin{cases}
2v^{n/2},& 2\mid n,\\[4pt]
0,& 2\nmid n,
\end{cases}\;}
\]
因而该原子层对碰撞参数 \(u\) 完全独立。若进一步取 \(v>0\)，则同一公式可重写为
\[
\boxed{\;
a_n^{\mathrm{atom}}(u,v)
=
(\sqrt v)^n+(-\sqrt v)^n.\;}
\]
特别地，所有奇阶 ghost/trace 审计都对因子 \((1-vz^2)^{-1}\) 完全失明：对任意 \(m\ge 0\)，其对
\[
\operatorname{Tr}(M_{u,v}^{\,2m+1})
\]
的原子贡献恒为 \(0\)。
\end{theorem}
\begin{proof}
由
\[
-\log(1-vz^2)=\sum_{k\ge 1}\frac{v^k}{k}z^{2k}
\]
并与 ghost 记号
\[
\log\frac{1}{1-vz^2}
=
\sum_{n\ge 1}\frac{a_n^{\mathrm{atom}}(u,v)}{n}z^n
\]
逐项比较系数，立得
\[
a_{2k}^{\mathrm{atom}}(u,v)=2v^k,
\qquad
a_{2k+1}^{\mathrm{atom}}(u,v)=0.
\]
这已经给出对 \(u\) 的完全独立性。

当 \(v>0\) 时，推论 \ref{cor:real-input-40-zeta-uv-sqrtv-eigs} 表明因子 \(1-vz^2\) 在谱上对应两条 rigid branches
\[
\lambda_\pm=\pm\sqrt v,
\]
故其对 \(n\)-步迹的贡献正是
\[
(\sqrt v)^n+(-\sqrt v)^n,
\]
而这与上式完全一致。于是奇数阶贡献必为零，从而任何只读取奇阶 ghost/trace 的审计函数都无法检测该原子因子。证毕。
\end{proof}

\begin{theorem}[去除 rigid branches 后的非刚性迹由八次 core 因子完全编码]\label{thm:xi-time-part73-nonrigid-trace-newton-cayley-closure}
沿用定理 \ref{thm:real-input-40-det-uv-2plus8} 与推论 \ref{cor:real-input-40-zeta-uv-sqrtv-eigs} 的记号。定义正规化 core 多项式
\[
\mathcal Q_8(z;u,v):=-P_8(z;u,v),
\qquad
\mathcal Q_8(0;u,v)=1.
\]
并写成
\[
\mathcal Q_8(z;u,v)=1+c_1z+\cdots+c_8z^8,
\]
其中若 \(\deg_z\mathcal Q_8<8\)，则约定高阶系数 \(c_j\) 自动取 \(0\)。当 \(v>0\) 时，再定义非刚性迹
\[
R_n(u,v):=\operatorname{Tr}(M_{u,v}^n)-(\sqrt v)^n-(-\sqrt v)^n
\qquad(n\ge 1).
\]
则成立：
\begin{enumerate}
  \item 生成函数恒等式
  \[
  \boxed{\;
  -\frac{z\,\partial_z\mathcal Q_8(z;u,v)}{\mathcal Q_8(z;u,v)}
  =
  \sum_{n\ge 1}R_n(u,v)\,z^n.\;}
  \]
  \item 当 \(1\le n\le 8\) 时，
  \[
  \boxed{\;
  R_n+c_1R_{n-1}+\cdots+c_{n-1}R_1+n\,c_n=0.\;}
  \]
  \item 当 \(n\ge 9\) 时，
  \[
  \boxed{\;
  R_n+c_1R_{n-1}+c_2R_{n-2}+\cdots+c_8R_{n-8}=0.\;}
  \]
\end{enumerate}
因此，二维加权核除去刚性对 \(\pm\sqrt v\) 之后的全部非刚性迹谱，完全由单一八次多项式 \(\mathcal Q_8\) 编码，并形成阶数至多为 \(8\) 的闭线性递推系统。
\end{theorem}
\begin{proof}
由定理 \ref{thm:real-input-40-det-uv-2plus8}，
\[
\Delta(z;u,v)=(1-vz^2)\mathcal Q_8(z;u,v).
\]
再由推论 \ref{cor:real-input-40-zeta-uv-sqrtv-eigs}，当 \(v>0\) 时矩阵 \(M_{u,v}\) 的两条 rigid branches 恰为
\[
\lambda_\pm=\pm\sqrt v.
\]
因此存在至多 \(8\) 个非刚性非零特征值
\[
\eta_1,\dots,\eta_d
\qquad(d\le 8),
\]
使得
\[
\mathcal Q_8(z;u,v)=\prod_{j=1}^{d}(1-z\eta_j),
\qquad
R_n(u,v)=\sum_{j=1}^{d}\eta_j^n.
\]
于是
\[
-\frac{z\,\partial_z\mathcal Q_8(z;u,v)}{\mathcal Q_8(z;u,v)}
=
\sum_{j=1}^{d}\frac{z\eta_j}{1-z\eta_j}
=
\sum_{j=1}^{d}\sum_{n\ge 1}\eta_j^n z^n
=
\sum_{n\ge 1}R_n(u,v)\,z^n,
\]
得到第一条。

其后两条只是标准 Newton--Girard 恒等式在
\[
\mathcal Q_8(z;u,v)=1+c_1z+\cdots+c_8z^8
\]
与幂和 \(R_n=\sum_j\eta_j^n\) 上的专门化：对 \(1\le n\le 8\) 得到非齐次关系
\[
R_n+c_1R_{n-1}+\cdots+c_{n-1}R_1+n\,c_n=0,
\]
而对 \(n\ge 9\) 则转为齐次关系
\[
R_n+c_1R_{n-1}+\cdots+c_8R_{n-8}=0.
\]
若 \(d<8\)，则高阶系数 \(c_{d+1},\dots,c_8\) 本来即为零，故同一公式仍然成立。证毕。
\end{proof}

\begin{theorem}[二维加权核在 \texorpdfstring{$v=1$}{v=1} 上发生共维一谱秩塌缩]\label{thm:xi-time-part73-v1-rank-collapse-codim2-tip}
沿用定理 \ref{thm:xi-time-part73-nonrigid-trace-newton-cayley-closure} 的记号。则有：
\begin{enumerate}
  \item 对任意 \(u\)，在超曲面 \(v=1\) 上都成立
  \[
  \deg_z \Delta(z;u,1)\le 9,
  \qquad
  \dim\ker M_{u,1}\ge 31.
  \]
  \item 在共维二点 \((u,v)=(0,1)\) 上，显式有
  \[
  P_8(z;0,1)=-2z^5-6z^4-z^3+5z^2+z-1,
  \]
  从而
  \[
  \deg_z \Delta(z;0,1)\le 7,
  \qquad
  \dim\ker M_{0,1}\ge 33.
  \]
  \item 相应地，上一条定理中的非刚性迹递推在 \(v=1\) 上自动降为阶数至多 \(7\) 的系统，而在 \((u,v)=(0,1)\) 处进一步降为阶数至多 \(5\) 的系统。
\end{enumerate}
\end{theorem}
\begin{proof}
由定理 \ref{thm:real-input-40-det-uv-2plus8}，
\[
P_8(z;u,v)
=(u^2v^4-u^2v^3)z^8+\cdots
=u^2v^3(v-1)z^8+\cdots.
\]
因此当 \(v=1\) 时，\(z^8\) 项恒消失，故
\[
\deg_z P_8(z;u,1)\le 7.
\]
再由
\[
\Delta(z;u,1)=(z^2-1)P_8(z;u,1),
\]
便得到
\[
\deg_z\Delta(z;u,1)\le 2+7=9.
\]
另一方面，\(\Delta(z;u,v)=\det(I-zM_{u,v})\) 的 \(z\)-次数恰等于 \(M_{u,v}\) 的全部非零特征值总数（按代数重数计）。由于 \(M_{u,1}\) 为 \(40\times 40\) 矩阵，这意味着
\[
\dim\ker M_{u,1}\ge 40-9=31.
\]

若进一步取 \((u,v)=(0,1)\)，则把 \(u=0,v=1\) 直接代入定理 \ref{thm:real-input-40-det-uv-2plus8} 的显式八次式可得
\[
P_8(z;0,1)=-2z^5-6z^4-z^3+5z^2+z-1,
\]
故
\[
\deg_z P_8(z;0,1)\le 5,
\qquad
\deg_z \Delta(z;0,1)\le 2+5=7.
\]
于是
\[
\dim\ker M_{0,1}\ge 40-7=33.
\]

最后，由定理 \ref{thm:xi-time-part73-nonrigid-trace-newton-cayley-closure}，非刚性迹递推的阶数恰由多项式 \(\mathcal Q_8=-P_8\) 的次数控制。故在 \(v=1\) 上其阶数至多为 \(7\)，而在 \((u,v)=(0,1)\) 处至多为 \(5\)。证毕。
\end{proof}


\begin{theorem}[二维加权 \texorpdfstring{$\zeta$}{zeta} 在碰撞零温缩放下坍缩为单一原子模型]\label{thm:xi-time-part73-twoparam-collision-freezing-single-atom-collapse}
沿用定理 \ref{thm:real-input-40-det-uv-2plus8} 的记号。固定 \(w>0\) 与 \(v>0\)，并令
\[
z=\sqrt{\frac{w}{u}},
\qquad
u\to+\infty.
\]
则有极限
\[
\boxed{\;
\lim_{u\to+\infty}\Delta\!\left(\sqrt{\frac{w}{u}};u,v\right)
=
1-vw.\;}
\]
更精确地，
\[
\boxed{\;
P_8\!\left(\sqrt{\frac{w}{u}};u,v\right)
=
-1+vw+O(u^{-1/2}),\;}
\]
从而
\[
\boxed{\;
\Delta\!\left(\sqrt{\frac{w}{u}};u,v\right)
=
1-vw+O(u^{-1/2}).\;}
\]
因此只要 \(vw\neq 1\)，便等价地有
\[
\boxed{\;
\zeta\!\left(\sqrt{\frac{w}{u}};u,v\right)
\longrightarrow
\frac{1}{1-vw}.\;}
\]
换言之，在碰撞零温缩放下，二维 \((u,v)\)-联合势的核心部分完全退化为长度 \(2\)、权重 \(v\) 的单原子 Euler 模型。
\end{theorem}
\begin{proof}
将定理 \ref{thm:real-input-40-det-uv-2plus8} 中 \(P_8(z;u,v)\) 的显式八次式代入
\[
z=\sqrt{\frac{w}{u}}.
\]
逐项估计可见：常数项给出 \(-1\)，而 \(z^2\) 项中的主贡献来自
\[
(uv+5v)z^2=vw+5vw\,u^{-1}.
\]
其余各项都至少带有 \(u^{-1/2}\) 的衰减；更具体地，\(z,z^3,z^5,z^7\) 项为 \(O(u^{-1/2})\)，而 \(z^4,z^6,z^8\) 项为 \(O(u^{-1})\)。故
\[
P_8\!\left(\sqrt{\frac{w}{u}};u,v\right)
=
-1+vw+O(u^{-1/2}).
\]

另一方面，
\[
vz^2-1
=
\frac{vw}{u}-1
=
-1+O(u^{-1}).
\]
于是
\[
\Delta\!\left(\sqrt{\frac{w}{u}};u,v\right)
=
\Bigl(-1+O(u^{-1})\Bigr)
\Bigl(-1+vw+O(u^{-1/2})\Bigr)
=
1-vw+O(u^{-1/2}),
\]
从而得到第一条极限。若 \(vw\neq 1\)，则极限不为零，故取倒数即得 \(\zeta\) 的极限公式。证毕。
\end{proof}

\begin{theorem}[\texorpdfstring{$u=1$}{u=1} 处 \texorpdfstring{$\zeta$}{zeta} 的双--三极点律]\label{thm:xi-time-part73-u1-zeta-double-triple-pole}
沿用命题 \ref{prop:real-input-40-output-spectral-collisions} 的记号，并记
\[
\zeta(z,1):=\Delta(z,1)^{-1},
\qquad
\zeta_{\mathrm{core}}(z,1):=Q(z,1)^{-1}.
\]
则有精确分解
\[
\Delta(z,1)=-(z-1)^2(z+1)^3(z^2-3z+1)(z^2-z-1).
\]
因此 \(\zeta(z,1)\) 在 \(z=1\) 具有二阶极点，在 \(z=-1\) 具有三阶极点。去除原子 Witt 因子 \((1-z^2)^{-1}\) 之后，核心 \(\zeta_{\mathrm{core}}(z,1)\) 在 \(z=1\)、\(z=-1\) 分别仅余一阶与二阶极点。
\end{theorem}
\begin{proof}
由命题 \ref{prop:real-input-40-output-spectral-collisions}，
\[
Q(z,1)=(z-1)(z+1)^2(z^2-3z+1)(z^2-z-1).
\]
另一方面，
\[
\Delta(z,1)=(1-z^2)Q(z,1)=-(z-1)(z+1)Q(z,1),
\]
将两式相乘即得所示因式分解。于是 \(\zeta(z,1)=\Delta(z,1)^{-1}\) 在 \(z=1\)、\(z=-1\) 的极点阶分别为 \(2\) 与 \(3\)。而
\[
\zeta_{\mathrm{core}}(z,1)=Q(z,1)^{-1}
\]
只除去原子因子 \((1-z^2)^{-1}\)，故两处极点阶分别降低 \(1\)。证毕。
\end{proof}

\begin{theorem}[原子二周期对 Abel 有限部的精确平移]\label{thm:xi-time-part73-two-cycle-abel-finitepart-shift}
沿用定理 \ref{thm:xi-time-part65-atomic-primitive-abel-shift} 的记号。在 \(u=1\) 的同一核上，设 \(z_\star=\varphi^{-2}\)。则 primitive--core 分解满足
\[
P_{\mathrm{core}}(z)=P(z)-z^2,
\qquad
\log\mathfrak M=\log\mathfrak M_{\mathrm{core}}+z_\star^2.
\]
因此
\[
\log\mathfrak M-\log\mathfrak M_{\mathrm{core}}
=
\varphi^{-4}
=
\frac{7-3\sqrt5}{2},
\]
亦即
\[
\mathfrak M_{\mathrm{core}}=e^{-\varphi^{-4}}\mathfrak M.
\]
\end{theorem}
\begin{proof}
由定理 \ref{thm:xi-time-part65-atomic-primitive-abel-shift}，
\[
P(z,u)=uz^2+P_{\mathrm{core}}(z,u).
\]
在 \(u=1\) 处专门化即得
\[
P_{\mathrm{core}}(z)=P(z)-z^2.
\]
同一定理并给出
\[
\log\mathfrak M=\log\mathfrak M_{\mathrm{core}}+z_\star^2.
\]
又在本核上 \(z_\star=\varphi^{-2}\)，故
\[
z_\star^2=\varphi^{-4}=\frac{7-3\sqrt5}{2}.
\]
代回并指数化，便得到所示 Abel 有限部平移公式。证毕。
\end{proof}

\begin{theorem}[激活分支的横截性与二次项消失]\label{thm:xi-time-part73-activated-branch-unique-zero-scale-singularity}
沿用命题 \ref{prop:real-input-40-activated-branch-linear}、注 \ref{rem:real-input-40-activated-branch-series} 与定理 \ref{thm:xi-output-potential-activated-branch-rigidity} 的记号，并记
\[
G(\lambda,u):=\lambda^8F(\lambda^{-1},u).
\]
则在 \(u=1\) 的某邻域内存在唯一解析分支 \(\lambda_{\mathrm{act}}(u)\) 满足
\[
G(\lambda_{\mathrm{act}}(u),u)=0,
\qquad
\lambda_{\mathrm{act}}(1)=0,
\]
并且其 Taylor 展开满足
\[
\boxed{\;
\lambda_{\mathrm{act}}(u)
=(u-1)+3(u-1)^3-(u-1)^4+18(u-1)^5-22(u-1)^6
+O((u-1)^7).\;}
\]
特别地，二次项严格消失，且
\[
\lambda'_{\mathrm{act}}(1)=1,
\]
从而该分支穿过 \(\lambda=0\) 的方式是横截而非抛物切触。进一步，存在常数 \(\rho>0\) 与 \(\varepsilon>0\)，使得对一切满足 \(|u-1|<\varepsilon\) 的实参数，都有
\[
\boxed{\;
|\lambda_j(u)|\ge \rho
\qquad
(2\le j\le 8).\;}
\]
\end{theorem}
\begin{proof}
由命题 \ref{prop:real-input-40-activated-branch-linear}，
\[
\partial_\lambda G(0,1)=1\neq 0,
\qquad
\lambda_{\mathrm{act}}(1)=0.
\]
故由隐函数定理，\(G(\lambda,u)=0\) 在 \((0,1)\) 邻域内确定唯一解析分支 \(\lambda_{\mathrm{act}}(u)\)，且
\[
\lambda'_{\mathrm{act}}(1)=1.
\]
另一方面，注 \ref{rem:real-input-40-activated-branch-series} 给出
\[
\lambda_{\mathrm{act}}(u)
=(u-1)+3(u-1)^3-(u-1)^4+18(u-1)^5-22(u-1)^6
+O((u-1)^7),
\]
故二次项系数严格为 \(0\)，而线性项系数为 \(1\)，从而该分支以横截方式穿过 \(\lambda=0\)。

最后，定理 \ref{thm:xi-output-potential-activated-branch-rigidity} 已证明：在 \((\lambda,u)=(0,1)\) 的某邻域内，唯一的小根就是 \(\lambda_{\mathrm{act}}(u)\)，其余七条根分支全部远离 \(0\)。因此存在 \(r>0\) 与 \(\varepsilon_0>0\)，使得当 \(|u-1|<\varepsilon_0\) 时，除 \(\lambda_{\mathrm{act}}(u)\) 外的所有根都落在 \(|\lambda|\ge r\) 上。取 \(\rho:=r\)、\(\varepsilon:=\varepsilon_0\) 即得最后一条。证毕。
\end{proof}

\begin{theorem}[激活分支的三次归一化刚性与去支实谱的无碰撞窗口]\label{thm:xi-time-part73-activated-branch-deflation-window}
沿用定理 \ref{thm:xi-time-part73-activated-branch-unique-zero-scale-singularity} 的记号，并记
\[
I_\ast:=
\bigl(0.93283726248596492776\ldots,\ 1.0547964213679192275\ldots\bigr).
\]
定义三次归一化 defect germ
\[
\Xi(u):=\frac{\lambda_{\mathrm{act}}(u)-(u-1)}{(u-1)^3}.
\]
又定义去支因子
\[
\widetilde G(\lambda,u):=\frac{G(\lambda,u)}{\lambda-\lambda_{\mathrm{act}}(u)}.
\]
则成立：
\begin{enumerate}
  \item \(\Xi(u)\) 在 \(u=1\) 处可解析延拓，并满足
  \[
  \Xi(1)=3.
  \]
  \item \(\widetilde G\) 在 \((\lambda,u)=(0,1)\) 的某邻域内解析；
  \item 对每个实参数 \(u\in I_\ast\setminus\{1\}\)，多项式 \(\widetilde G(\cdot,u)\) 的全部实根都是单根。换言之，去除激活分支之后，审计窗口 \(I_\ast\) 内的剩余实谱不再发生双根碰撞。
\end{enumerate}
\end{theorem}
\begin{proof}
由注 \ref{rem:real-input-40-activated-branch-series}，
\[
\lambda_{\mathrm{act}}(u)-(u-1)
=
3(u-1)^3-(u-1)^4+18(u-1)^5-22(u-1)^6+O((u-1)^7).
\]
故
\[
\Xi(u)
=
3-(u-1)+18(u-1)^2-22(u-1)^3+O((u-1)^4),
\]
从而 \(\Xi\) 在 \(u=1\) 处具有可去奇性，并且
\[
\Xi(1)=3.
\]

第二条正是定理 \ref{thm:xi-output-potential-activated-branch-rigidity} 的 Weierstrass 因除结论。

最后，设某个 \(u_0\in I_\ast\setminus\{1\}\) 使 \(\widetilde G(\cdot,u_0)\) 存在实双根 \(\lambda_0\)。由于
\[
G(0,u)=u^3(1-u),
\]
且 \(u_0\neq 1\)，故 \(\lambda_0\neq 0\)。记
\[
z_0:=\lambda_0^{-1}\in \RR.
\]
在 \(\lambda\neq 0\) 上，变换 \(z=\lambda^{-1}\) 为解析双射，而
\[
G(\lambda,u)=\lambda^8F(\lambda^{-1},u)
\]
只差一个处处非零因子 \(\lambda^8\)，故 \(\lambda_0\) 为 \(\widetilde G(\cdot,u_0)\) 的实双根会强迫 \(z_0\) 成为 core 因子 \(F(z,u_0)\) 的实双根；沿一参数输出位势的记号，这等价于 \(Q(z,u_0)\) 在实轴上出现重根。可这与注 \ref{rem:real-input-40-output-potential-degeneracy-phase} 矛盾：该注指出在正实轴上除 \(u=1\) 外，全部退化点恰为
\[
0.008460\ldots,\ 0.745534\ldots,\ 0.932837\ldots,\ 1.054796\ldots,\ 2.090392\ldots,\ 15.522820\ldots,
\]
因而在开区间 \(I_\ast\setminus\{1\}\) 内不存在任何重根参数。故 \(\widetilde G(\cdot,u)\) 的实根在该窗口内全为单根。证毕。
\end{proof}

\begin{corollary}[围绕 \texorpdfstring{$u=1$}{u=1} 的局部谱病态完全由唯一激活小根承担]\label{cor:xi-time-part73-single-activated-branch-all-local-illconditioning}
沿用定理 \ref{thm:xi-time-part73-activated-branch-unique-zero-scale-singularity} 与定理 \ref{thm:xi-time-part73-activated-branch-deflation-window} 的记号。则存在 \(u=1\) 的邻域 \(\mathcal U\) 与在 \((\lambda,u)=(0,1)\) 某邻域内解析的函数 \(\widetilde G(\lambda,u)\)，使得
\[
G(\lambda,u)=\bigl(\lambda-\lambda_{\mathrm{act}}(u)\bigr)\widetilde G(\lambda,u)
\qquad(u\in\mathcal U).
\]
并且存在常数 \(c>0\)，使得对每个 \(u\in\mathcal U\)，\(\widetilde G(\cdot,u)\) 的任一零点 \(\mu(u)\) 都满足
\[
|\mu(u)|\ge c.
\]
因此，\(u=1\) 邻域内全部由小根诱导的局部谱病态都由唯一分支 \(\lambda_{\mathrm{act}}(u)\) 承担；去除该分支之后，其余七支统一恢复到 \(O(1)\) 尺度。
\end{corollary}
\begin{proof}
定理 \ref{thm:xi-time-part73-activated-branch-deflation-window} 的第二条给出
\[
\widetilde G(\lambda,u):=\frac{G(\lambda,u)}{\lambda-\lambda_{\mathrm{act}}(u)}
\]
在 \((\lambda,u)=(0,1)\) 的某邻域内解析，故所示 Weierstrass 因除成立。

另一方面，定理 \ref{thm:xi-time-part73-activated-branch-unique-zero-scale-singularity} 已证明：在 \(u=1\) 的某邻域内，除 \(\lambda_{\mathrm{act}}(u)\) 外，其余全部根都与 \(0\) 保持一致正距离。于是可取某个常数 \(c>0\)，使 \(\widetilde G(\cdot,u)\) 的所有零点统一满足
\[
|\mu(u)|\ge c.
\]
这恰表明全部局部小尺度病态只由 \(\lambda_{\mathrm{act}}(u)\) 一支承担，而去支后的剩余谱支都停留在 \(O(1)\) 距离之外。证毕。
\end{proof}

\begin{theorem}[正实轴上的谱“护城河”]\label{thm:xi-time-part73-positive-real-spectral-moat}
仍在真实输入 \(40\) 状态核的输出位势加权下，取
\[
\Delta(z,u)=(1-uz^2)Q(z,u),
\qquad
u>0.
\]
则正实参数区间
\[
1<u<4
\]
构成一个严格的谱护城河：在该区间内同时不存在显式原子分支 \(z=\pm u^{-1/2}\) 与 \(Q\)-谱的碰撞，也不存在 \(Q\)-谱在整数点 \(z=\pm1\) 的共振。更精确地，
\[
Q(\pm u^{-1/2},u)=0
\iff
u=1,
\]
且
\[
Q(1,u)=-u(u-4)(u-1)(u+1),
\qquad
Q(-1,u)=-(u-1)^2(u^2+6u-2).
\]
故在 \(u>1\) 的正实轴上，首个整数级共振恰发生于
\[
u=4,
\qquad
\lambda=1.
\]
\end{theorem}
\begin{proof}
命题 \ref{prop:real-input-40-output-spectral-collisions} 已给出
\[
Q\!\left(\frac{1}{\sqrt u},u\right)=\frac{2(\sqrt u-1)(\sqrt u+1)^2}{u^{3/2}},
\qquad
Q\!\left(-\frac{1}{\sqrt u},u\right)=\frac{2(\sqrt u-1)^2(\sqrt u+1)}{u^{3/2}},
\]
因而在 \(u>0\) 上，
\[
Q(\pm u^{-1/2},u)=0 \iff u=1.
\]
这说明在 \(1<u<4\) 内不存在原子分支与 \(Q\)-谱的碰撞。

同一命题又给出
\[
Q(1,u)=-u(u-4)(u-1)(u+1),
\qquad
Q(-1,u)=-(u-1)^2(u^2+6u-2).
\]
当 \(1<u<4\) 时，前式显然非零；而
\[
u^2+6u-2>0
\qquad (u>1),
\]
故后式亦非零。因此在 \(1<u<4\) 内不存在 \(Q\)-谱于 \(z=\pm1\) 的整数级共振。最后，\(u=4\) 时 \(Q(1,4)=0\)，从而给出首个正实整数级共振 \(\lambda=1\)。证毕。
\end{proof}

\begin{theorem}[主 Perron 根相对于原子谱支存在严格 \texorpdfstring{$\sqrt u$}{sqrt(u)} 级分离，并从上方逼近零温四态 Perron 常数]\label{thm:xi-time-part73-perron-vs-atomic-sqrtu-separation}
沿用推论 \ref{cor:real-input-40-factor-sqrtu-eigs}、命题 \ref{prop:real-input-40-pressure-freezing} 与命题 \ref{prop:real-input-40-output-mirror-branch} 的记号。记 \(\lambda(u)\) 为输出势的 Perron 正根，并令
\[
P(\theta):=\log\lambda(e^\theta),
\qquad
c^2=\rho(B_\infty),
\]
其中 \(c\approx 1.89635299351377\)、\(d\approx 0.807943317323508\)，且 \(c^2\) 是四次方程
\[
y^4-6y^3+9y^2-y-1=0
\]
的最大正根。则当 \(u\to+\infty\) 时有：
\[
\boxed{\;
\lambda(u)-\sqrt u
=
(c-1)\sqrt u+d+O(u^{-1/2}),\;}
\]
因而
\[
\boxed{\;
\lambda(u)-\sqrt u\longrightarrow +\infty.\;}
\]
进一步，
\[
\boxed{\;
\frac{\lambda(u)^2}{u}
=
\rho(B_\infty)+\frac{2cd}{\sqrt u}+O(u^{-1}),\;}
\]
特别地，由于 \(d>0\)，该比值从上方逼近 \(\rho(B_\infty)\)。
最后，若取 \(u=e^\theta\)，则
\[
\boxed{\;
P(\theta)-\frac{\theta}{2}
=
\log c+\frac{d}{c}e^{-\theta/2}+O(e^{-\theta})
\longrightarrow
\log c>0.\;}
\]
因此显式原子谱支 \(\pm\sqrt u\) 只决定零温端点的斜率 \(1/2\)，而主导自由能的幅度常数则完全由零温四态核心 \(B_\infty\) 决定。
\end{theorem}
\begin{proof}
由命题 \ref{prop:real-input-40-output-mirror-branch} 的 Puiseux 展开，
\[
\lambda(u)=c\sqrt u+d+\frac{e}{\sqrt u}+\frac{f}{u}+O(u^{-3/2}).
\]
与显式原子谱支 \(\sqrt u\) 相减，便得到
\[
\lambda(u)-\sqrt u
=
(c-1)\sqrt u+d+O(u^{-1/2}).
\]
由于 \(c>1\)，主项 \((c-1)\sqrt u\) 严格为正并趋于无穷，从而得到第一条。

再将同一展开平方：
\[
\lambda(u)^2
=
c^2u+2cd\sqrt u+O(1).
\]
两边同除以 \(u\)，并用 \(c^2=\rho(B_\infty)\)，即得
\[
\frac{\lambda(u)^2}{u}
=
\rho(B_\infty)+\frac{2cd}{\sqrt u}+O(u^{-1}).
\]
由于 \(d>0\)，该展开表明收敛从上方发生。

最后令 \(u=e^\theta\)。由
\[
\lambda(e^\theta)
=
c\,e^{\theta/2}
\left(
1+\frac{d}{c}e^{-\theta/2}+O(e^{-\theta})
\right)
\]
可得
\[
P(\theta)
=
\log\lambda(e^\theta)
=
\frac{\theta}{2}+\log c+\frac{d}{c}e^{-\theta/2}+O(e^{-\theta}),
\]
这就给出最后一条公式。证毕。
\end{proof}

\paragraph{二维切片塌缩与原子 Witt 校正的变量分裂}
在二维加权 \(\zeta(z;u,v)\) 的 \((2+8)\) 次显式分解、两条 rigid branches \(\pm\sqrt v\) 与一参数输出位势中的原子 Witt 剥离都已建立之后，还可把输出完全抑制切片与二维/一维原子分量的变量分工进一步压缩为下列两条接口。

\begin{theorem}[二维加权 \texorpdfstring{$\zeta(z;u,v)$}{zeta(z;u,v)} 在 \texorpdfstring{$v=0$}{v=0} 切片上的完全谱塌缩]\label{thm:xi-time-part73a-twoparam-vzero-total-spectral-collapse}
\leanverified{paper\_xi\_time\_part73a\_twoparam\_vzero\_total\_spectral\_collapse}
沿用定理 \ref{thm:real-input-40-det-uv-2plus8} 的记号。则对任意碰撞参数 \(u\) 都有严格恒等式
\[
\boxed{\;
\Delta(z;u,0)=1-z.\;}
\]
因此：
\begin{enumerate}
  \item \(M_{u,0}\) 的非零谱恰为单点 \(\{1\}\)；
  \item 由于 \(M_{u,0}\in M_{40}(\CC)\)，其零特征值的代数重数恰为 \(39\)；
  \item 相应加权 Artin--Mazur \(\zeta\) 函数满足
  \[
  \boxed{\;
  \zeta(z;u,0)=\frac{1}{1-z}.\;}
  \]
  \item 若定义 ghost 坐标与 primitive 生成函数
  \[
  a_n(u,0):=\Tr(M_{u,0}^n),
  \qquad
  P(z;u,0):=\sum_{n\ge 1}p_n(u,0)z^n,
  \]
  则有
  \[
  \boxed{\;
  a_n(u,0)=1\qquad(n\ge 1),\;}
  \]
  以及
  \[
  \boxed{\;
  P(z;u,0)=z,\qquad
  p_1(u,0)=1,\quad p_n(u,0)=0\ (n\ge 2).\;}
  \]
\end{enumerate}
换言之，一旦完全抑制输出位 \(e=1\)，二维模型中的碰撞权 \(u\) 在谱层面完全失明。
\end{theorem}
\begin{proof}
由定理 \ref{thm:real-input-40-det-uv-2plus8}，
\[
\Delta(z;u,v)=(v z^2-1)P_8(z;u,v),
\]
其中
\[
P_8(z;u,v)
=
(u^2 v^4-u^2 v^3)z^8+\cdots +(uv+5v)z^2+z-1.
\]
将 \(v=0\) 代入，则除末两项外其余各项全都消失，从而
\[
P_8(z;u,0)=z-1.
\]
于是
\[
\Delta(z;u,0)=(-1)(z-1)=1-z.
\]
这就给出第一条恒等式。

又由于该加权矩阵作用在真实输入 \(40\) 状态核上，故 \(M_{u,0}\) 为 \(40\times 40\) 矩阵。特征多项式的非平凡部分既然仅为 \(1-z\)，则唯一非零特征值只能是 \(1\)，其余 \(39\) 个特征值全为 \(0\)。于是
\[
\zeta(z;u,0)=\Delta(z;u,0)^{-1}=(1-z)^{-1}.
\]
从而
\[
\log\zeta(z;u,0)=\sum_{n\ge 1}\frac{z^n}{n},
\]
故 ghost 系数满足 \(a_n(u,0)=1\) 对所有 \(n\ge 1\) 成立；同一式也表明 primitive 生成函数恰为
\[
P(z;u,0)=z.
\]
于是 \(p_1(u,0)=1\) 且 \(p_n(u,0)=0\)（\(n\ge 2\)）。证毕。
\end{proof}

\begin{corollary}[二维/一维原子 Witt 校正的变量分裂与单层激活]\label{cor:xi-time-part73a-atomic-witt-bivariate-univariate-variable-rigidity}
分别考虑二维加权模型 \(\zeta(z;u,v)\) 与一参数输出位势模型 \(\zeta(z,u)\)。
\begin{enumerate}
  \item 在二维模型中，若定义 primitive 生成函数
  \[
  P(z;u,v):=\sum_{n\ge 1}p_n(u,v)z^n,
  \qquad
  P_{\mathrm{core}}(z;u,v):=\sum_{n\ge 1}p_n^{\mathrm{core}}(u,v)z^n,
  \]
  以及 ghost 坐标
  \[
  \log\zeta(z;u,v)=\sum_{n\ge 1}\frac{a_n(u,v)}{n}z^n,
  \qquad
  \log\zeta_{\mathrm{core}}(z;u,v)=\sum_{n\ge 1}\frac{a_n^{\mathrm{core}}(u,v)}{n}z^n,
  \]
  则有
  \[
  \boxed{\;
  P(z;u,v)=v z^2+P_{\mathrm{core}}(z;u,v),\;}
  \]
  \[
  \boxed{\;
  a_n(u,v)=
  \begin{cases}
  2v^{n/2}+a_n^{\mathrm{core}}(u,v), & 2\mid n,\\[4pt]
  a_n^{\mathrm{core}}(u,v), & 2\nmid n.
  \end{cases}\;}
  \]
  因而对任意 \(r\ge 1\) 与任意 \(n\ge 1\)，都有
  \[
  \partial_u^r\!\bigl(p_n(u,v)-p_n^{\mathrm{core}}(u,v)\bigr)\equiv 0,
  \qquad
  \partial_u^r\!\bigl(a_n(u,v)-a_n^{\mathrm{core}}(u,v)\bigr)\equiv 0.
  \]
  \item 在一参数输出位势模型中，若记主极点 \(z_\star(u)=\lambda(u)^{-1}\)，则原子 Witt 校正恰在 primitive 的 \(n=2\) 层激活，并满足
  \[
  \boxed{\;
  p_n(u)=p_n^{\mathrm{core}}(u)+u\,\mathbf 1_{\{n=2\}},\;}
  \]
  \[
  \boxed{\;
  a_n(u)=a_n^{\mathrm{core}}(u)+2u^{n/2}\mathbf 1_{\{2\mid n\}},\;}
  \]
  \[
  \boxed{\;
  \log\mathfrak M(u)-\log\mathfrak M_{\mathrm{core}}(u)=u\,z_\star(u)^2.\;}
  \]
  特别地，在无权点 \(u=1\) 上，
  \[
  \log\mathfrak M(1)-\log\mathfrak M_{\mathrm{core}}(1)
  =
  z_\star(1)^2
  =
  \varphi^{-4}.
  \]
\end{enumerate}
\end{corollary}
\begin{proof}
对二维模型，固定任意 \(u\) 并把 \(v\) 视为加权变量。由定理 \ref{thm:real-input-40-det-uv-2plus8}，
\[
\zeta(z;u,v)=\frac{1}{1-vz^2}\,\zeta_{\mathrm{core}}(z;u,v).
\]
于是定理 \ref{thm:xi-finite-euler-primitive-surgery-additivity} 在单原子情形 \((m,E,\ell)=(1,1,2)\) 下直接给出
\[
P(z;u,v)=v z^2+P_{\mathrm{core}}(z;u,v).
\]
另一方面，定理 \ref{thm:xi-time-part73-twoparam-atomic-ghost-u-blindness} 已证明
\[
a_n(u,v)-a_n^{\mathrm{core}}(u,v)=2v^{n/2}\mathbf 1_{\{2\mid n\}}.
\]
由于上述二维原子修正只含 \(v\) 而不含 \(u\)，故关于 \(u\) 的任意高阶导数恒为零。

对一参数模型，全部公式都已由定理 \ref{thm:xi-real-input40-evenlength-abel-shift-rigidity} 给出；最后在 \(u=1\) 处再用该定理中的专门化
\[
z_\star(1)=\varphi^{-2}
\]
便得到
\[
z_\star(1)^2=\varphi^{-4}.
\]
证毕。
\end{proof}


\paragraph{固定参数 Möbius 反演下的偶长 necklace 校正}
在上文采用参数提升型 primitive 生成函数之外，若固定 \(u,v\) 并仅对长度指标施行 Möbius 反演，则二维原子因子的偶长校正还可进一步压缩为一条显式 necklace 闭式。

\begin{theorem}[二维加权 periodic 谱中的普适偶长原子校正]\label{thm:xi-time-part73c-periodic-evenlength-atomic-correction}
沿用定理 \ref{thm:real-input-40-det-uv-2plus8} 与定理 \ref{thm:xi-time-part73-nonrigid-trace-newton-cayley-closure} 的记号。记
\[
a_n(u,v):=\operatorname{Tr}(M_{u,v}^n),
\qquad
\mathcal Q_8(z;u,v):=-P_8(z;u,v),
\]
并由
\[
-\log \mathcal Q_8(z;u,v)=\sum_{n\ge 1}\frac{\widetilde a_n(u,v)}{n}z^n
\]
定义 residual periodic 系数 \(\widetilde a_n(u,v)\)。则对任意 \(n\ge 1\)，有
\[
a_n(u,v)=\widetilde a_n(u,v)+2v^{n/2}\mathbf 1_{\{2\mid n\}}.
\]
等价地，对每个 \(m\ge 0\) 有
\[
a_{2m+1}(u,v)=\widetilde a_{2m+1}(u,v),
\]
而对每个 \(m\ge 1\) 有
\[
a_{2m}(u,v)=\widetilde a_{2m}(u,v)+2v^m.
\]
特别地，该偶长校正与碰撞参数 \(u\) 完全无关。
\leanverified{evenLengthCorrection}
\leanverified{evenLengthCorrection\_odd}
\leanverified{evenLengthCorrection\_even}
\leanverified{evenLengthCorrection\_cases}
\end{theorem}
\begin{proof}
由定理 \ref{thm:real-input-40-det-uv-2plus8}，
\[
\Delta(z;u,v)=(vz^2-1)P_8(z;u,v)=(1-vz^2)\mathcal Q_8(z;u,v).
\]
因此
\[
\zeta(z;u,v)=\frac{1}{\Delta(z;u,v)}
=
\frac{1}{1-vz^2}\,\mathcal Q_8(z;u,v)^{-1}.
\]
又由 \(\mathcal Q_8(0;u,v)=1\)，取形式对数得到
\[
\log\zeta(z;u,v)
=
-\log(1-vz^2)-\log \mathcal Q_8(z;u,v).
\]
另一方面，
\[
\log\zeta(z;u,v)=\sum_{n\ge 1}\frac{a_n(u,v)}{n}z^n,
\]
且
\[
-\log(1-vz^2)=\sum_{m\ge 1}\frac{v^m}{m}z^{2m}.
\]
把这三式逐项比较，便得
\[
\frac{a_n(u,v)-\widetilde a_n(u,v)}{n}
=
\begin{cases}
\dfrac{v^{n/2}}{n/2},& 2\mid n,\\[6pt]
0,& 2\nmid n,
\end{cases}
\]
亦即
\[
a_n(u,v)-\widetilde a_n(u,v)=2v^{n/2}\mathbf 1_{\{2\mid n\}}.
\]
最后一条只是同一公式的直接重述。证毕。
\end{proof}

\begin{corollary}[固定参数 Möbius 本原系数的偶长 necklace 校正]\label{cor:xi-time-part73c-fixed-parameter-necklace-correction}
固定参数 \(u,v\)。定义固定参数 Möbius 本原系数 \(p_n^{\mathrm{fix}}(u,v)\)、\(\widetilde p_n^{\mathrm{fix}}(u,v)\) 为
\[
a_n(u,v)=\sum_{d\mid n} d\,p_d^{\mathrm{fix}}(u,v),
\qquad
\widetilde a_n(u,v)=\sum_{d\mid n} d\,\widetilde p_d^{\mathrm{fix}}(u,v).
\]
则对任意 \(n\ge 1\)，有
\[
p_n^{\mathrm{fix}}(u,v)-\widetilde p_n^{\mathrm{fix}}(u,v)
=
\frac{2}{n}
\sum_{\substack{d\mid n\\ 2\mid n/d}}
\mu(d)\,v^{n/(2d)}.
\]
特别地：
\[
n\ \text{为奇数}
\quad\Longrightarrow\quad
p_n^{\mathrm{fix}}(u,v)=\widetilde p_n^{\mathrm{fix}}(u,v),
\]
\[
p_{2m}^{\mathrm{fix}}(u,v)-\widetilde p_{2m}^{\mathrm{fix}}(u,v)
=
\frac{1}{m}\sum_{d\mid m}\mu(d)\,v^{m/d}
\qquad (m\ge 1).
\]
若记
\[
M_m(v):=\frac{1}{m}\sum_{d\mid m}\mu(d)\,v^{m/d},
\]
则 \(M_m(v)\) 正是第 \(m\) 个 necklace 多项式。因而固定参数的 primitive 校正只支撑于偶长层，并在长度 \(2m\) 处恰等于 \(M_m(v)\)。
\leanverified{necklaceCorrectionKernel}
\leanverified{necklaceCorrectionKernel\_odd}
\leanverified{necklaceCorrectionKernel\_even}
\leanverified{necklaceCorrectionKernel\_odd\_eq\_zero}
\leanverified{paper\_xi\_time\_part73c\_fixed\_parameter\_necklace\_correction}
\end{corollary}
\begin{proof}
由经典 Möbius 反演，
\[
n\,p_n^{\mathrm{fix}}(u,v)=\sum_{d\mid n}\mu(d)\,a_{n/d}(u,v),
\]
\[
n\,\widetilde p_n^{\mathrm{fix}}(u,v)=\sum_{d\mid n}\mu(d)\,\widetilde a_{n/d}(u,v).
\]
两式相减，并代入定理 \ref{thm:xi-time-part73c-periodic-evenlength-atomic-correction}，得到
\[
n\Bigl(
p_n^{\mathrm{fix}}(u,v)-\widetilde p_n^{\mathrm{fix}}(u,v)
\Bigr)
=
2\sum_{\substack{d\mid n\\ 2\mid n/d}}
\mu(d)\,v^{n/(2d)}.
\]
除以 \(n\) 即得第一式。

若 \(n\) 为奇数，则条件 \(2\mid n/d\) 不可能成立，故校正为零。若 \(n=2m\)，则
\[
d\mid 2m,\ 2\mid \frac{2m}{d}
\quad\Longleftrightarrow\quad
d\mid m,
\]
于是上式化为
\[
p_{2m}^{\mathrm{fix}}(u,v)-\widetilde p_{2m}^{\mathrm{fix}}(u,v)
=
\frac{1}{m}\sum_{d\mid m}\mu(d)\,v^{m/d}.
\]
最后一条只是把右端识别为标准 necklace 多项式。证毕。
\end{proof}


\noindent
以上这些结果表明，真实输入 \(40\) 状态核在二维原子层、非刚性残余谱、\(u=1\) 的局部激活层与 \(u\to+\infty\) 的零温核心层之间存在严格的多尺度分裂：一方面，二维加权 \(\zeta\) 的单原子因子在 ghost 层只向偶长度注入 \(2v^{n/2}\)，并对碰撞权 \(u\) 完全失明；去除两条 rigid branches \(\pm\sqrt v\) 之后，其余全部非刚性迹又自动坍缩为由单一八次 core 因子控制的 Newton--Cayley 递推系统，并在 \(v=1\) 超曲面与 \((u,v)=(0,1)\) 尖点上继续发生可量化的谱秩塌缩；当进一步压到输出完全抑制的切片 \(v=0\) 时，整个二维谱又严格塌缩为唯一一个长度 \(1\) 的 primitive 轨道。另一方面，在碰撞零温缩放下，整个二维核心因子进一步退化为同一原子的单参数模型 \(1-vw\)；在一参数输出位势上，同一原子 Witt 机制只在 primitive 的 \(n=2\) 层激活，并把 Abel 有限部刚性平移为 \(u z_\star(u)^2\)；再者，激活分支以无二次项的横截方式穿过 \(\lambda=0\)，其三次归一化 defect 在 \(u=1\) 处锁定为常数 \(3\)，去支后的剩余实谱则在最近退化点所围成的审计窗口内恢复为无碰撞解析族；而正实轴上的区间 \(1<u<4\) 又把原子碰撞与整数级共振共同隔离在外。最后，真正主导大 \(u\) 幅度常数的并不是显式原子谱支 \(\pm\sqrt u\)，而是零温四态核心 \(B_\infty\) 的 Perron 常数。因而在读出口径上，前述几条对应于中心投影式的相位盲读出，而其后的非刚性迹系统、局部极点、激活分支与谱壁垒则共同构成相位敏感的纤维内读出。

\paragraph{Lee--Yang 振幅域、有限支持账本障碍与 Smith--椭圆寄存器极值}
在 Stokes--Lee--Yang 三次数域的共同域识别、自然扩张的 prime-register 外置、Smith 信息亏损谱的离散曲率恢复，以及周期化 Bernoulli 共振与整数轴对齐椭圆半群的素数寄存器结构已经建立之后，还可进一步压缩出下列九条直接后果。它们分别刻画 Lee--Yang 振幅常数 \(B\) 的六次全复扩张、无限前史的有限支持 prime-slice 账本不可能性、固定总估值下 Smith 信息亏损的 sharp 极值原理，以及周期化 Fourier 共振与整数轴对齐椭圆在结构层和实现层上的刚性裂解。

\begin{theorem}[Lee--Yang 振幅常数 \texorpdfstring{$B$}{B} 的六次不可约性]\label{thm:xi-time-part74-leyang-amplitude-sixth-irreducibility}
沿用定理 \ref{thm:xi-terminal-zm-stokes-amplitude-square-cubic-rigidity} 与定理 \ref{thm:xi-terminal-zm-stokes-leyang-shared-artin-representation} 的记号，记
\[
b:=B^2,
\qquad
K:=\QQ(b)=\QQ(u)=\QQ(\lambda_*).
\]
则
\[
[\QQ(B):\QQ]=6,
\]
并且 \(B\) 在 \(\QQ\) 上满足的 primitive 不可约整式恰为
\[
\boxed{\;
162X^6+1593X^4+1998X^2+1184.\;}
\]
\end{theorem}
\begin{proof}
由定理 \ref{thm:xi-terminal-zm-stokes-amplitude-square-cubic-rigidity}，
\[
162b^3+1593b^2+1998b+1184=0
\]
是 \(b\) 在 \(\QQ\) 上的不可约三次式。再由定理
\ref{thm:xi-terminal-zm-stokes-leyang-shared-artin-representation}，
\[
K=\QQ(b)=\QQ(u)=\QQ(\lambda_*),
\qquad
[K:\QQ]=3.
\]

若 \(B\in K\)，则 \(b=B^2\) 为 \(K\) 中平方元，于是其范数必为 \(\QQ\) 中平方：
\[
N_{K/\QQ}(b)=N_{K/\QQ}(B)^2\in \QQ_{\ge 0}.
\]
另一方面，由三次极小多项式的首末系数可得
\[
N_{K/\QQ}(b)=(-1)^3\frac{1184}{162}=-\frac{592}{81},
\]
这在 \(\QQ\) 中不可能是平方，矛盾。故 \(B\notin K\)，从而 \(X^2-b\) 在 \(K[X]\) 中不可约，并且
\[
[\QQ(B):K]=2.
\]
于是
\[
[\QQ(B):\QQ]=2[K:\QQ]=6.
\]

显然 \(B\) 满足
\[
162B^6+1593B^4+1998B^2+1184=0.
\]
由于该整式次数为 \(6=[\QQ(B):\QQ]\)，故它就是 \(B\) 的 primitive 不可约整式。证毕。
\end{proof}

\begin{corollary}[Lee--Yang 振幅域的签名与二次扩张结构]\label{cor:xi-time-part74-leyang-amplitude-field-signature}
令
\[
L:=\QQ(B),
\qquad
K:=\QQ(B^2)=\QQ(b).
\]
则 \(L/K\) 为二次扩张；域 \(K\) 的签名为 \((1,1)\)，而 \(L\) 的签名为 \((0,3)\)。特别地，\(L\) 是全复六次域，而其立方子域 \(K\) 恰有一个实嵌入。
\end{corollary}
\begin{proof}
由定理 \ref{thm:xi-time-part74-leyang-amplitude-sixth-irreducibility}，\([L:K]=2\)。再由定理 \ref{thm:xi-terminal-zm-stokes-amplitude-square-cubic-rigidity}，\(K\) 由不可约三次式
\[
f(t)=162t^3+1593t^2+1998t+1184
\]
生成，而推论 \ref{cor:xi-terminal-zm-stokes-b2-quadratic-resolvent} 给出的判别式平方自由部分为 \(-111\)，故 \(\Disc(f)<0\)。因此 \(K\) 的签名为 \((1,1)\)。

又因为 \(f\) 的系数全为正，所以 \(f(t)\) 在正实轴上无根；而三次实系数多项式必有实根，故其唯一实根只能是负数。于是对 \(K\) 的唯一实嵌入 \(\sigma:K\hookrightarrow \RR\)，有
\[
\sigma(b)<0.
\]
因此方程 \(x^2=\sigma(b)\) 在 \(\RR\) 中无解，故 \(\sigma\) 的任何延拓都不给出 \(L\) 的实嵌入。其余两条嵌入本来就是复嵌入，延拓后仍成共轭复对。由于 \([L:\QQ]=6\)，遂得 \(L\) 的签名为 \((0,3)\)。证毕。
\end{proof}

\begin{theorem}[自然扩张的有限支持 prime-slice 外置不可能性]\label{thm:xi-time-part74-natural-extension-finite-support-primeslice-impossibility}
\leanverified{paper\_xi\_time\_part74\_natural\_extension\_finite\_support\_primeslice\_impossibility}
设 \(T:X\twoheadrightarrow X\) 为可数集合上的有限纤维满射，并记其自然扩张为
\[
X^{\leftarrow}
:=
\left\{
(x_0,x_{-1},x_{-2},\dots)\in X^{\NN}:\ T(x_{-j})=x_{-j+1}\ \ (j\ge 1)
\right\}.
\]
对每个 \(N\ge 1\)，再记 \(N\)-步前史空间
\[
X^{\leftarrow,N}
:=
\left\{
(x_0,x_{-1},\dots,x_{-N})\in X^{N+1}:\ T(x_{-j})=x_{-j+1}\ \ (1\le j\le N)
\right\}.
\]
假设对每个 \(N\ge 1\) 以及每一层 \(0\le j\le N-1\)，都存在某个纤维点使该层纤维大小至少为 \(2\)。在定理 \ref{thm:killo-layered-prime-slices-finite-depth} 的分层 prime-slice 语义下，不存在任何单射
\[
\Psi:X^{\leftarrow}\longrightarrow X\times \ZZ_{>0},
\qquad
\Psi(\mathbf x)=\bigl(x_0,H(\mathbf x)\bigr),
\]
其编码分解为
\[
H(\mathbf x)=\prod_{j\ge 0}h_j(\mathbf x),
\]
并满足：
\begin{enumerate}
  \item 每个 \(h_j(\mathbf x)\) 的全部素因子都限制在预先固定的 prime slice \(\mathcal P_j\)；
  \item 存在统一常数 \(M\)，使得对所有 \(\mathbf x\in X^{\leftarrow}\) 与所有 \(j\ge M\) 都有
  \[
  h_j(\mathbf x)=1.
  \]
\end{enumerate}
\end{theorem}
\begin{proof}
由定理 \ref{thm:killo-natural-extension-branch-register}，自然扩张 \((X^{\leftarrow},\sigma)\) 与右移插入的分支寄存器系统共轭。另一方面，定理 \ref{thm:xi-time-part63b-infinite-depth-finite-prime-register-impossibility} 已经证明：只要每一层都存在真实分支，则任何保持全历史可恢复性的精确分层 prime-register 外置都不可能只使用有限多个 prime slices。

若存在所述编码 \(\Psi\) 与统一常数 \(M\)，则对每个 \(\mathbf x\in X^{\leftarrow}\)，整数 \(H(\mathbf x)\) 的全部非平凡素因子都落在前 \(M\) 个 prime slices 中；也就是说，这一外置实际上只使用了有限多个 prime slices。于是它正与定理 \ref{thm:xi-time-part63b-infinite-depth-finite-prime-register-impossibility} 的必要性结论矛盾。故这种有限支持外置不存在。证毕。
\end{proof}

\begin{corollary}[\texorpdfstring{$N$}{N} 步前史的最优活动 prime-slice 复杂度恰等于深度]\label{cor:xi-time-part74-finite-past-primeslice-complexity-depth}
\leanverified{paper\_xi\_time\_part74\_finite\_past\_primeslice\_complexity\_depth}
在上述设定下，假设对每个 \(N\ge 1\) 的前史空间 \(X^{\leftarrow,N}\)，其每一层都存在真实分支。记
\[
\mathfrak c_N
\]
为所有单射分层 prime-slice 外置
\[
\Phi_N:X^{\leftarrow,N}\hookrightarrow X\times \ZZ_{>0},
\qquad
\Phi_N(x_0,\dots,x_{-N})=(x_0,H_N(x_0,\dots,x_{-N}))
\]
在最坏情形下所需的最少非平凡 prime slices 数目。则
\[
\boxed{\;
\mathfrak c_N=N.\;}
\]
\end{corollary}
\begin{proof}
在 \(X^{\leftarrow,N}\) 上，每一层 \(0,\dots,N-1\) 都是活跃分支层。故定理 \ref{thm:xi-layered-primeslice-complexity-active-depth} 直接给出下界
\[
\mathfrak c_N\ge N.
\]
另一方面，定理 \ref{thm:killo-layered-prime-slices-finite-depth} 构造了恰在每一层写入一枚有效素数的单射外置，从而
\[
\mathfrak c_N\le N.
\]
两式合并即得 \(\mathfrak c_N=N\)。证毕。
\end{proof}

\begin{theorem}[固定总估值下 Smith 信息亏损的 sharp 极值原理]\label{thm:xi-time-part74-smith-loss-fixed-total-valuation-extremal}
\leanverified{paper\_xi\_time\_part74\_smith\_loss\_fixed\_total\_valuation\_extremal}
设 \(A\in M_{m\times n}(\ZZ)\) 满足 \(\operatorname{rank}_{\QQ}(A)=m\)，其 Smith 不变量为
\[
d_1\mid d_2\mid \cdots \mid d_m.
\]
固定素数 \(p\) 与整数 \(k\ge 1\)，记
\[
e_i:=v_p(d_i),
\qquad
E:=\sum_{i=1}^{m}e_i,
\qquad
f_{p,k}(A):=\sum_{i=1}^{m}\min(k,e_i).
\]
把 \(E\) 作 Euclidean 分解
\[
E=mq+r,
\qquad
0\le r<m.
\]
若只固定总估值 \(E\) 而允许 \((e_1,\dots,e_m)\) 在全部 Smith 估值谱中变化，则
\[
\boxed{\;
\min f_{p,k}(A)=\min(k,E),\;}
\]
\[
\boxed{\;
\max f_{p,k}(A)=(m-r)\min(k,q)+r\min(k,q+1).\;}
\]
极小值由集中谱
\[
(0,\dots,0,E)
\]
达到；极大值由最均衡谱
\[
(\underbrace{q,\dots,q}_{m-r},\underbrace{q+1,\dots,q+1}_{r})
\]
达到。
\end{theorem}
\begin{proof}
记
\[
\psi_k(t):=\min(k,t)
\qquad (t\in \ZZ_{\ge 0}).
\]
则
\[
f_{p,k}(A)=\sum_{i=1}^{m}\psi_k(e_i).
\]
函数 \(\psi_k\) 的离散导数
\[
\delta_k(t):=\psi_k(t)-\psi_k(t-1)
\qquad (t\ge 1)
\]
满足 \(\delta_k(t)=1\) 当 \(t\le k\)，而 \(\delta_k(t)=0\) 当 \(t>k\)。因此 \(\delta_k\) 单调不增，\(\psi_k\) 是离散凹函数。

现取任意两个坐标 \(a\le b-2\)。把它们替换为 \(a+1,b-1\) 时，总和保持不变，而函数值变化为
\[
\bigl(\psi_k(a+1)-\psi_k(a)\bigr)-\bigl(\psi_k(b)-\psi_k(b-1)\bigr)
=
\delta_k(a+1)-\delta_k(b)\ge 0.
\]
因此，每一次“均衡化”操作都不会减小 \(\sum_i\psi_k(e_i)\)。反复执行直到任意两坐标之差至多为 \(1\)，便必然到达唯一的最均衡谱
\[
(\underbrace{q,\dots,q}_{m-r},\underbrace{q+1,\dots,q+1}_{r}),
\]
从而得到极大值
\[
(m-r)\psi_k(q)+r\psi_k(q+1)
=(m-r)\min(k,q)+r\min(k,q+1).
\]

反向地，每一次把质量从较小坐标转移到较大坐标的“集中化”操作都不会增大 \(\sum_i\psi_k(e_i)\)。不断集中后，最终得到极端谱
\[
(0,\dots,0,E),
\]
其对应值为
\[
\psi_k(E)=\min(k,E).
\]
于是这给出极小值。

最后，二者都可由对角 Smith 形直接实现：取矩阵
\[
\operatorname{diag}(p^{e_1},\dots,p^{e_m})
\]
即可得到相应估值谱；若需固定列数 \(n\ge m\)，只需在右侧附加 \(n-m\) 个零列。故上下界都是 sharp。证毕。
\end{proof}

\begin{corollary}[一般模数下 Smith 信息亏损的逐素数 sharp 极值]\label{cor:xi-time-part74-general-modulus-kernel-loss-extremal}
\leanverified{paper\_xi\_time\_part74\_general\_modulus\_kernel\_loss\_extremal}
沿用定理 \ref{thm:xi-time-part74-smith-loss-fixed-total-valuation-extremal} 的记号。对任意模数
\[
b=\prod_p p^{k_p},
\]
记
\[
E_p:=v_p(d_1\cdots d_m)=mq_p+r_p,
\qquad
0\le r_p<m.
\]
若只固定每个素数方向的总估值 \(E_p\) 而允许 Smith 谱在各素数方向重新分配，则诱导映射
\[
A_b:(\ZZ/b\ZZ)^n\to(\ZZ/b\ZZ)^m
\]
满足 sharp bounds
\[
\boxed{\;
(n-m)\log b+\sum_p \min(k_p,E_p)\log p
\le
\log \#\ker(A_b),\;}
\]
\[
\boxed{\;
\log \#\ker(A_b)
\le
(n-m)\log b+\sum_p\Bigl[(m-r_p)\min(k_p,q_p)+r_p\min(k_p,q_p+1)\Bigr]\log p.\;}
\]
并且上下界都可由适当的 Smith 谱精确达到。
\end{corollary}
\begin{proof}
由推论 \ref{cor:killo-kernel-size-general-modulus}，
\[
\log \#\ker(A_b)
=
(n-m)\log b+\sum_p\left(\sum_{i=1}^{m}\min(k_p,e_i(p))\right)\log p,
\]
其中 \(e_i(p):=v_p(d_i)\)。对每个固定素数 \(p\)，在总估值
\[
\sum_{i=1}^{m}e_i(p)=E_p
\]
已定的条件下，定理 \ref{thm:xi-time-part74-smith-loss-fixed-total-valuation-extremal} 给出
\[
\min \sum_i\min(k_p,e_i(p))=\min(k_p,E_p),
\]
以及
\[
\max \sum_i\min(k_p,e_i(p))
=(m-r_p)\min(k_p,q_p)+r_p\min(k_p,q_p+1).
\]
逐素数求和便得到所示上下界。

Sharpness 也逐素数独立实现：对每个 \(p\)，任选达到极小或极大的非降估值谱 \((e_1(p),\dots,e_m(p))\)。再定义
\[
d_i:=\prod_p p^{e_i(p)}.
\]
由于每个素数方向上都保持非降序，故仍有
\[
d_1\mid d_2\mid \cdots \mid d_m.
\]
于是这些 \(d_i\) 确实给出某个 Smith 链，并且在各素数方向同时达到相应极值，从而整体界是 sharp。证毕。
\end{proof}

\begin{theorem}[临界周期化测度在 Fibonacci 双频上的正极限与非 Rajchman 性]\label{thm:xi-time-part74-periodized-bernoulli-fibonacci-limit-nonrajchman}
令
\[
q:=\varphi^{-1},
\qquad
\nu:=\bigl(\RR\to \mathbb T_2=\RR/(2\ZZ)\bigr)_\#\mu_q^{(2)}.
\]
再记
\[
c_{\mathrm{Fib}}:=\lim_{n\to\infty}C_\varphi(F_n)>0
\]
为推论 \ref{cor:fold-resonance-ladder-fib-luc-limits} 给出的 Fibonacci 方向极限常数。则
\[
\bigl|\widehat \nu(F_m)\bigr|\longrightarrow c_{\mathrm{Fib}},
\qquad
\bigl|\widehat \nu(F_{m+1})\bigr|\longrightarrow c_{\mathrm{Fib}}.
\]
特别地，\(\nu\) 不是 Rajchman 测度；因而 \(\nu\) 不可能对 \(\mathbb T_2\) 上的 Haar 测度绝对连续。
\end{theorem}
\begin{proof}
由定理 \ref{thm:fold-pisot-bernoulli-convolution-representation}，
\[
\widehat \nu(u)=\widehat{\mu_q^{(2)}}(\pi u)
\qquad (u\in \ZZ),
\]
并且
\[
\bigl|\widehat{\mu_q^{(2)}}(\pi u)\bigr|=C_\varphi(u).
\]
故对任意 \(m\ge 1\)，有
\[
\bigl|\widehat \nu(F_m)\bigr|=C_\varphi(F_m),
\qquad
\bigl|\widehat \nu(F_{m+1})\bigr|=C_\varphi(F_{m+1}).
\]
再由推论 \ref{cor:fold-resonance-ladder-fib-luc-limits}，
\[
C_\varphi(F_m)\to c_{\mathrm{Fib}},
\qquad
C_\varphi(F_{m+1})\to c_{\mathrm{Fib}},
\]
其中 \(c_{\mathrm{Fib}}>0\)。于是 \(\widehat \nu(u)\) 沿 Fibonacci 双频子列不趋于 \(0\)，从而 \(\nu\) 不是 Rajchman 测度。

若 \(\nu\) 对 Haar 测度绝对连续，则在有限测度空间 \(\mathbb T_2\) 上它必具有某个 \(L^1\) 密度；Riemann--Lebesgue 引理遂推出 \(\widehat \nu(u)\to 0\)，与上式矛盾。故 \(\nu\) 不可能绝对连续。证毕。
\end{proof}

\begin{corollary}[周期化临界测度全部卷积幂的共振刚性]\label{cor:xi-time-part74-periodized-bernoulli-convolution-powers-nonrajchman}
在定理 \ref{thm:xi-time-part74-periodized-bernoulli-fibonacci-limit-nonrajchman} 的设定下，对任意整数 \(r\ge 1\)，卷积幂 \(\nu^{*r}\) 满足
\[
\bigl|\widehat{\nu^{*r}}(F_m)\bigr|
\longrightarrow
c_{\mathrm{Fib}}^r>0.
\]
特别地，\(\nu^{*r}\) 对所有 \(r\ge 1\) 都不是 Rajchman 测度，也都不可能对 Haar 测度绝对连续。
\end{corollary}
\begin{proof}
Fourier 变换把卷积化为乘法：
\[
\widehat{\nu^{*r}}(u)=\widehat \nu(u)^r.
\]
由定理 \ref{thm:xi-time-part74-periodized-bernoulli-fibonacci-limit-nonrajchman}，
\[
\bigl|\widehat \nu(F_m)\bigr|\to c_{\mathrm{Fib}}>0.
\]
于是
\[
\bigl|\widehat{\nu^{*r}}(F_m)\bigr|
=
\bigl|\widehat \nu(F_m)\bigr|^r
\longrightarrow
c_{\mathrm{Fib}}^r>0.
\]
因此 \(\nu^{*r}\) 同样不是 Rajchman 测度。若它绝对连续，则其 Fourier 系数应当趋于 \(0\)，再由 Riemann--Lebesgue 引理得到矛盾。证毕。
\end{proof}

\begin{theorem}[整数轴对齐椭圆半群在结构层与实现层的半圆维裂解]\label{thm:xi-time-part74-integer-ellipse-structure-implementation-split}
令
\[
\mathsf E_{\NN}
:=
\left\{
E_{a,b}:\ \frac{x^2}{a^2}+\frac{y^2}{b^2}=1,\ a,b\in\NN_{>0}
\right\},
\qquad
E_{a,b}\odot E_{c,d}:=E_{ac,bd}.
\]
则：
\begin{enumerate}
  \item 作为交换幺半群，
  \[
  (\mathsf E_{\NN},\odot)\cong (\NN^{(\mathcal P)}\times \NN^{(\mathcal P)},+)
  \cong (\NN_{>0}\times \NN_{>0},\times);
  \]
  \item 其实现层半圆维满足
  \[
  \cdim^{\mathrm{impl}}(\mathsf E_{\NN})=\frac12;
  \]
  \item 其结构层同态半圆维满足
  \[
  \cdim^{\mathrm{hom}}_{+}(\mathsf E_{\NN},\odot)=+\infty.
  \]
\end{enumerate}
亦即，整数轴对齐椭圆半群在可实现性意义下仍只是可数对象，而在严格同态线性化意义下却需要可数无穷个彼此独立的素数轴。
\end{theorem}
\begin{proof}
第一条正是定理 \ref{thm:killo-ellipse-diagonal-prime-register-equivalence} 的内容。因而 \(\mathsf E_{\NN}\) 与 \(\NN_{>0}\times \NN_{>0}\) 等势，特别是一个可数无限集。由命题 \ref{prop:killo-implementation-vs-structural-half-circle-dimension} 的定义口径，可数无限集的实现层半圆维恰为 \(\frac12\)，从而得到第二条。

再看结构层。子幺半群
\[
\mathsf E_x:=\{E_{a,1}:a\in\NN_{>0}\}\subset \mathsf E_{\NN}
\]
满足
\[
(\mathsf E_x,\odot)\cong (\NN_{\times},\times).
\]
若存在某个有限 \(k\) 与单射幺半群同态
\[
\Phi:(\mathsf E_{\NN},\odot)\hookrightarrow (\NN^k,+),
\]
则其限制将给出
\[
\NN_{\times}\hookrightarrow (\NN^k,+)
\]
的单射幺半群同态，这与定理 \ref{thm:killo-prime-freedom-non-finitizability} 矛盾。故不存在任何有限秩加法寄存器对 \(\mathsf E_{\NN}\) 给出忠实同态实现，也即
\[
\cdim^{\mathrm{hom}}_{+}(\mathsf E_{\NN},\odot)=+\infty.
\]
证毕。
\end{proof}

\noindent
这九条结果把 Lee--Yang 分支振幅、自然扩张账本、Smith 信息亏损与椭圆 prime-register 几条后段链条进一步压缩成统一结构：一方面，Stokes 幅值平方虽与 Lee--Yang 三次共享同一立方子域，但其平方根振幅 \(B\) 必然逃离该三次数域，转而落入一个全复六次扩张；另一方面，有限深度前史虽然能够以逐层单素数的最优方式压缩，但这种最优性一旦进入无限前史便不再可有限支持化，从而把 prime-slice 账本刚性提升为真正的逆极限障碍；再者，Smith 信息亏损不仅由离散曲率完全恢复，而且在固定总估值约束下满足精确的集中--均衡极值律，并逐素数推送到一般模数的 sharp 核增长界；最后，Fibonacci 方向的周期化共振既在全部卷积幂上保持非 Rajchman 刚性，又与整数轴对齐椭圆半群的可数无穷素数轴结构一道，展示出实现层可数而结构层不可有限线性化的根本裂解。

\paragraph{前缀圆柱的 dyadic 商几何、分歧阴影屏蔽与共振尾的双层机制}
在全息地址的 \(2^L\)-进圆柱--球识别、规范地址单子的仿射忠实实现、十次分歧域与 Lee--Yang/Latt\`es 三域 Frobenius 张量分解，以及共振闭包常数的对数发散已经建立之后，还可进一步抽取出下列六条彼此闭合的后果。它们分别刻画事件前缀在 Tate 模块商空间中的精确余类几何、有限维忠实实现必然停留在仿射地址层、十次分歧阴影对外部扩张的普适分歧屏蔽、三分歧体系上由共轭类并生成的布尔独立律，以及碰撞异常由刚性有限核与不可截断无限尾共同驱动的双层机制。

\begin{theorem}[事件前缀圆柱与 \texorpdfstring{$B^nT$}{B^nT}-陪集的严格等价]\label{thm:xi-time-part75-hologram-prefix-cylinder-coset-equivalence}
\leanverified{paper\_xi\_time\_part75\_hologram\_prefix\_cylinder\_coset\_equivalence}
沿用结论 \ref{con:xi-time-part62dgb-hologram-cylinder-ball-isometry} 的记号。设
\[
T:=T_2(E_{\min})\cong \ZZ_2^2,
\qquad
\mathrm{code}:E\hookrightarrow \NN,
\qquad
B:=2^L>M:=\max_{e\in E}\mathrm{code}(e),
\]
并固定 \(v\in T\setminus 2T\)。对任意事件流 \(\mathbf e=(e_t)_{t\ge 1}\in E^{\NN}\)，记
\[
X(\mathbf e):=\sum_{t\ge 1}\mathrm{code}(e_t)\,B^{t-1}v,
\qquad
\mathcal X:=X(E^{\NN})\subset T.
\]
对任意长度 \(n\) 前缀
\[
a=(a_1,\dots,a_n)\in E^n
\]
记
\[
[a]:=\{\mathbf e\in E^{\NN}:\ e_1=a_1,\dots,e_n=a_n\},
\qquad
x(a):=\sum_{t=1}^{n}\mathrm{code}(a_t)\,B^{t-1}v.
\]
则有精确恒等式
\[
\boxed{\;
X([a])=\mathcal X\cap \bigl(x(a)+B^nT\bigr).\;}
\]
因此，长度 \(n\) 前缀与 \(\mathcal X\) 在商空间 \(T/B^nT\) 中所占据的深度 \(n\) 陪集之间存在严格一一对应。
\end{theorem}
\begin{proof}
前一条盒装恒等式正是结论 \ref{con:xi-time-part62dgb-hologram-cylinder-ball-isometry} 的圆柱--球公式在当前记号下的重写。

现证后一条对应。由上式可知，每个前缀 \(a\in E^n\) 都确定一个非空占用陪集
\[
x(a)+B^nT\in T/B^nT.
\]
若两个前缀 \(a,b\in E^n\) 给出同一占用陪集，则
\[
\mathcal X\cap \bigl(x(a)+B^nT\bigr)
=
\mathcal X\cap \bigl(x(b)+B^nT\bigr),
\]
亦即
\[
X([a])=X([b]).
\]
由于 \(X:E^{\NN}\to \mathcal X\) 为单射，遂有 \([a]=[b]\)，从而 \(a=b\)。故该对应既满且单，因而为双射。证毕。
\end{proof}

\begin{corollary}[全息地址像集的商占用计数与精确盒维数]\label{cor:xi-time-part75-hologram-image-quotient-box-dimension}
\leanverified{paper\_xi\_time\_part75\_hologram\_image\_quotient\_box\_dimension}
沿用定理 \ref{thm:xi-time-part75-hologram-prefix-cylinder-coset-equivalence} 的记号。则对每个 \(n\ge 1\)，像集 \(\mathcal X\) 在商空间 \(T/B^nT\) 中恰占据
\[
|E|^n
\]
个陪集，而全部深度 \(n\) 陪集总数为
\[
|T/B^nT|=B^{2n}.
\]
因而
\[
\underline{\dim}_{\mathrm B}\mathcal X
=
\overline{\dim}_{\mathrm B}\mathcal X
=
\dim_{\mathrm B}\mathcal X
=
\frac{\log |E|}{\log B}
=
\frac{\log_2|E|}{L}.
\]
\end{corollary}
\begin{proof}
由定理 \ref{thm:xi-time-part75-hologram-prefix-cylinder-coset-equivalence}，长度 \(n\) 前缀与 \(\mathcal X\) 在 \(T/B^nT\) 中的占用陪集一一对应，故占用数正好为
\[
|E^n|=|E|^n.
\]
另一方面，由
\[
T\cong \ZZ_2^2,
\qquad
B=2^L,
\]
可知
\[
T/B^nT\cong (\ZZ/B^n\ZZ)^2,
\]
故其基数为 \(B^{2n}\)。

因此在尺度 \(B^{-n}\) 下，覆盖 \(\mathcal X\) 所需的最小环境球数恰为 \(|E|^n\)。于是上下盒维数都由极限
\[
\lim_{n\to\infty}\frac{\log |E|^n}{-\log(B^{-n})}
=
\lim_{n\to\infty}\frac{n\log|E|}{n\log B}
=
\frac{\log|E|}{\log B}
\]
给出。最后一式仅由 \(B=2^L\) 改写为
\[
\frac{\log|E|}{\log B}=\frac{\log_2|E|}{L}.
\]
证毕。
\end{proof}

\begin{theorem}[有限维忠实实现的仿射必然性]\label{thm:xi-time-part75-faithful-finite-dimensional-realization-affine-necessity}
\leanverified{paper\_xi\_time\_part75\_faithful\_finite\_dimensional\_realization\_affine\_necessity}
沿用定理 \ref{thm:killo-godel-semidir-affine-2adic} 与定理 \ref{thm:xi-time-part9g-canonical-submonoid-free-monoid} 的记号。设
\[
T:=T_2(E_{\min})\cong \ZZ_2^2,
\qquad
H_{\mathrm{can}}\subset \mathcal R\rtimes_{\Sigma}\NN
\]
为规范地址单子，且 \(B=2^L>M\)。则存在忠实单子同态
\[
\Psi:H_{\mathrm{can}}\hookrightarrow \mathrm{Aff}(T).
\]
另一方面，不存在任何有限 \(k\) 与单射幺半群同态
\[
\NN_{\times}\hookrightarrow (\NN^k,+),
\]
更一般地，也不存在任何到有限生成交换可消去加法幺半群或有限生成交换群的忠实乘法账本线性化。特别地，
\[
\cdim^{\mathrm{hom}}_{+}(\NN_{\times})=+\infty,
\qquad
\cdim^{\mathrm{impl}}(\NN_{\times})=\frac12.
\]
因此，本项目中能够在有限秩对象上保持忠实性的层必然是仿射--非交换的地址动力学层，而不可能是有限秩--交换的严格加法账本层。
\end{theorem}
\begin{proof}
第一部分正是定理 \ref{thm:killo-godel-semidir-affine-2adic} 在规范子单子 \(H_{\mathrm{can}}\) 上的忠实性结论。由于
\[
T_2(E_{\min})\cong \ZZ_2^2,
\]
该忠实实现确实落在固定秩 \(2\) 的 \(2\)-进 Tate 模块的仿射自同构半群中。

第二部分由定理 \ref{thm:killo-prime-freedom-non-finitizability} 直接推出：\(\NN_{\times}\) 不能嵌入任何有限生成交换可消去幺半群。有限生成交换群在加法口径下当然也是有限生成交换可消去幺半群，故同样被排除。最后，命题 \ref{prop:killo-implementation-vs-structural-half-circle-dimension} 已给出
\[
\cdim^{\mathrm{hom}}_{+}(\NN_{\times})=+\infty,
\qquad
\cdim^{\mathrm{impl}}(\NN_{\times})=\frac12.
\]
这就得到所述严格分裂。证毕。
\end{proof}

\begin{theorem}[\texorpdfstring{$S_{10}$}{S10} 阴影的普适分歧屏蔽判据]\label{thm:xi-time-part75-p10-shadow-universal-ramification-shielding}
\leanverified{paper\_xi\_time\_part75\_p10\_shadow\_universal\_ramification\_shielding}
设 \(L_{10}/\QQ\) 为命题 \ref{prop:fold-gauge-anomaly-P10-galois-discriminant} 中十次分歧多项式 \(P_{10}\) 的分裂域，则
\[
\Gal(L_{10}/\QQ)\cong S_{10},
\]
并且其唯一非平凡真 Galois 子域为
\[
K_{10}:=L_{10}^{A_{10}}=\QQ(\sqrt{66269989}),
\qquad
66269989=1931\cdot 34319.
\]
令 \(M/\QQ\) 为任意 Galois 扩张。若
\[
\operatorname{Ram}(M/\QQ)\cap\{1931,34319\}=\varnothing,
\]
则
\[
L_{10}\cap M=\QQ.
\]
\end{theorem}
\begin{proof}
由命题 \ref{prop:fold-gauge-anomaly-P10-galois-discriminant} 与推论 \ref{cor:fold-gauge-anomaly-p10-unique-quadratic-subfield}，
\[
\Gal(L_{10}/\QQ)\cong S_{10},
\qquad
K_{10}=L_{10}^{A_{10}}=\QQ(\sqrt{66269989}),
\]
而 \(S_{10}\) 的唯一非平凡真正规子群正是 \(A_{10}\)。因此，\(L_{10}/\QQ\) 的唯一非平凡真 Galois 子域只能是 \(K_{10}\)。

现设
\[
F:=L_{10}\cap M.
\]
由于 \(L_{10}/\QQ\) 与 \(M/\QQ\) 都是 Galois 扩张，故 \(F/\QQ\) 亦为 Galois 扩张，并且 \(F\subseteq L_{10}\)。遂有
\[
F\in\{\QQ,K_{10}\}.
\]
若 \(F\neq \QQ\)，则必有 \(F=K_{10}\subseteq M\)。

另一方面，命题 \ref{prop:fold-gauge-anomaly-P10-galois-discriminant} 已说明 \(K_{10}/\QQ\) 的域判别式为
\[
66269989=1931\cdot 34319,
\]
故 \(K_{10}\) 恰在素数 \(1931,34319\) 处分歧。若 \(K_{10}\subseteq M\)，则这些素数必在 \(M/\QQ\) 中分歧，这与
\[
\operatorname{Ram}(M/\QQ)\cap\{1931,34319\}=\varnothing
\]
矛盾。故 \(F=\QQ\)。证毕。
\end{proof}

\begin{theorem}[三分歧体系上由共轭类并生成的布尔独立律]\label{thm:xi-time-part75-triple-frobenius-boolean-independence}
\leanverified{paper\_xi\_time\_part75\_triple\_frobenius\_boolean\_independence}
沿用定理 \ref{thm:xi-time-part62d-triple-frobenius-classfunction-tensor-independence} 的记号。记
\[
G:=S_{10}\times S_3\times (S_4\times C_2),
\]
并对共同不分歧素数 \(p\) 记相应 Frobenius 分量为
\[
\sigma_{10}(p)\in S_{10},
\qquad
\sigma_{\mathrm{LY}}(p)\in S_3,
\qquad
\sigma_4^{\mathrm{tot}}(p)\in S_4\times C_2.
\]
任取共轭类并
\[
\Omega_{10}\subset S_{10},
\qquad
\Omega_{\mathrm{LY}}\subset S_3,
\qquad
\Omega_4\subset S_4\times C_2,
\]
并定义事件
\[
\mathcal E_{10}:=\{\sigma_{10}(p)\in \Omega_{10}\},
\qquad
\mathcal E_{\mathrm{LY}}:=\{\sigma_{\mathrm{LY}}(p)\in \Omega_{\mathrm{LY}}\},
\qquad
\mathcal E_4:=\{\sigma_4^{\mathrm{tot}}(p)\in \Omega_4\}.
\]
则自然密度满足精确乘法律
\[
\boxed{\;
\delta(\mathcal E_{10}\cap \mathcal E_{\mathrm{LY}}\cap \mathcal E_4)
=
\frac{|\Omega_{10}|}{10!}\cdot
\frac{|\Omega_{\mathrm{LY}}|}{6}\cdot
\frac{|\Omega_4|}{48}.\;}
\]
特别地，由三个因子的共轭类并分别生成的三个有限 Boolean 代数在自然密度意义下相互独立；因而对任一因子的条件化，都不改变其余两个因子的共轭类分布。
\end{theorem}
\begin{proof}
对三个因子分别取指标类函数
\[
f:=\mathbf 1_{\Omega_{10}},
\qquad
g:=\mathbf 1_{\Omega_{\mathrm{LY}}},
\qquad
h:=\mathbf 1_{\Omega_4}.
\]
则
\[
\langle f\rangle_{S_{10}}=\frac{|\Omega_{10}|}{10!},
\qquad
\langle g\rangle_{S_3}=\frac{|\Omega_{\mathrm{LY}}|}{6},
\qquad
\langle h\rangle_{S_4\times C_2}=\frac{|\Omega_4|}{48}.
\]
定理 \ref{thm:xi-time-part62d-triple-frobenius-classfunction-tensor-independence} 对任意类函数给出严格张量分解，于是
\[
\delta(\mathcal E_{10}\cap \mathcal E_{\mathrm{LY}}\cap \mathcal E_4)
=
\langle f\rangle_{S_{10}}\,
\langle g\rangle_{S_3}\,
\langle h\rangle_{S_4\times C_2},
\]
即得所示公式。

又因为“共轭类并”在每个有限群上本身就构成一个有限 Boolean 代数，而上述公式对任意这样的并集都成立，所以三个因子对应的 Boolean 代数两两独立且三重独立。条件分布不变性只是将联合密度公式除以相应条件事件的正密度所得。证毕。
\end{proof}

\begin{theorem}[共振异常的刚性有限核与不可截断无限尾]\label{thm:xi-time-part75-resonance-anomaly-rigid-core-infinite-tail}
\leanverified{paper\_xi\_time\_part75\_resonance\_anomaly\_rigid\_core\_infinite\_tail}
记
\[
\Gamma_m:=F_{m+2}\,\mathsf{Col}_m.
\]
则该归一化碰撞尺度具有严格的两层分裂。
\begin{enumerate}
  \item 存在由两条 bulk 共振频率
  \[
  k=F_m,\qquad k=F_{m+1}
  \]
  锁定的刚性有限核，并且
  \[
  \Gamma_m\ge 1+a_m(F_m)^2+a_m(F_{m+1})^2,
  \]
  从而
  \[
  \liminf_{m\to\infty}\Gamma_m\ge 1+2C_\varphi^2>1.
  \]
  \item 对任意固定 \(U\in\NN\)，都有
  \[
  \liminf_{m\to\infty}\Gamma_m
  \ge
  1+2\sum_{u=1}^{U}C_\varphi(u)^2.
  \]
  由于
  \[
  \sum_{u=1}^{\infty}C_\varphi(u)^2=+\infty,
  \]
  故
  \[
  \Gamma_m\longrightarrow +\infty.
  \]
\end{enumerate}
因此，该异常既不是单频效应，也不是任何有限频率集合可以封闭解释的现象：两条 Fibonacci bulk 模只给出一个不可消去的刚性有限核，而整体无界增长则由不可截断的共振稀疏尾共同驱动。
\end{theorem}
\begin{proof}
由定义 \ref{def:fold-normalized-amplitude} 与 Parseval 口径，
\[
\Gamma_m=F_{m+2}\,\mathsf{Col}_m=\sum_{k\in \ZZ/(F_{m+2}\ZZ)}a_m(k)^2.
\]
其中零频项恒为
\[
a_m(0)^2=1.
\]
又由定理 \ref{thm:fold-critical-resonance-constant}，
\[
a_m(F_m)\longrightarrow C_\varphi,
\qquad
a_m(F_{m+1})\longrightarrow C_\varphi.
\]
故只取三项 \(k=0,F_m,F_{m+1}\) 即得
\[
\Gamma_m\ge 1+a_m(F_m)^2+a_m(F_{m+1})^2,
\]
从而
\[
\liminf_{m\to\infty}\Gamma_m\ge 1+2C_\varphi^2.
\]
这就给出第一层的刚性有限核。

另一方面，推论 \ref{cor:fold-collision-gap-resonance-ladder-U} 断言：对任意固定 \(U\in\NN\)，都有
\[
\liminf_{m\to\infty}\Gamma_m
=
\liminf_{m\to\infty}F_{m+2}\,\mathsf{Col}_m
\ge
1+2\sum_{u=1}^{U}C_\varphi(u)^2.
\]
而定理 \ref{thm:fold-sigma-phi-diverges} 给出
\[
1+2\sum_{u=1}^{\infty}C_\varphi(u)^2=+\infty.
\]
故对任意 \(M>0\)，总可取 \(U\) 使
\[
1+2\sum_{u=1}^{U}C_\varphi(u)^2\ge M.
\]
再代回上式，得到
\[
\liminf_{m\to\infty}\Gamma_m\ge M.
\]
由于 \(M\) 任意，遂有 \(\liminf_{m\to\infty}\Gamma_m=+\infty\)，亦即 \(\Gamma_m\to+\infty\)。证毕。
\end{proof}

\noindent
这六条结果把地址几何、实现层分裂、分歧阴影、三域 Frobenius 独立性与碰撞异常的解析尾部进一步压缩成一组可直接横向调用的接口：首先，事件时间流在 \(T_2(E_{\min})\cong \ZZ_2^2\) 中并非只给出近似的 \(2^L\)-进前缀结构，而是以陪集商的方式精确外显，并留下完全闭式的盒维数 \(\log_2|E|/L\)；其次，有限维忠实实现虽能在仿射地址动力学层上严格成立，但任何试图把乘法层压平到有限秩交换加法账本的方案都会被素数自由度障碍排除；再者，十次分歧域的全部非平凡 Galois 阴影都锁定在二次域 \(\QQ(\sqrt{66269989})\) 上，因此一切避开素数 \(1931,34319\) 的外部 Galois 扩张都必须与之算术透明；最后，三域 Frobenius 事件不仅在具体类事件上满足乘法律，而且在由各因子共轭类并生成的 Boolean 代数层面依然严格独立，而折叠碰撞异常则同时呈现出由两条 Fibonacci bulk 模固定下来的刚性有限核与由全部 \(C_\varphi(u)\) 稀疏尾推向无界的整体增长。

\paragraph{无限共振尾、自然扩张深度与冻结端面的双相收束}
在碰撞缺口的多共振下界、谱零点几何的指数稀疏、自然扩张的 prime-register 极小性、幂散度极限常数的黄金参数可辨识性，以及冻结相 escort 的共同基态塌缩都已建立之后，还可进一步把这些此前分散出现的链条压缩为一条双相机制：一方面，碰撞侧的主尺度异常既不能由任何有限共振预算封闭，也不能由零点计数单独控制；另一方面，冻结侧的高阶 escort 与压力射线却在同一基态微正则壳层上发生严格塌缩与刚性外推。二者之间的结构接口，则由自然扩张在有限深度与无限持续深度之间出现的计数型相变，以及单一幂散度标量对整条极限信息几何塔的决定性所给出。

\begin{theorem}[碰撞缺口的无限共振秩]\label{thm:xi-time-part75a-collision-gap-infinite-resonance-rank}
\leanverified{paper\_xi\_time\_part75a\_collision\_gap\_infinite\_resonance\_rank}
记缩放碰撞缺口
\[
\Gamma_m:=F_{m+2}\!\left(\mathsf{Col}_m-\frac{1}{F_{m+2}}\right),
\]
并对每个 \(U\in\NN\) 定义有限共振预算
\[
R_U:=2\sum_{u=1}^{U}C_\varphi(u)^2.
\]
则对任意固定 \(U\)，都有
\[
\liminf_{m\to\infty}\bigl(\Gamma_m-R_U\bigr)=+\infty.
\]
\end{theorem}
\begin{proof}
任取 \(U\in\NN\) 与 \(M>0\)。由定理 \ref{thm:fold-sigma-phi-diverges}，
\[
\sum_{u=1}^{\infty}C_\varphi(u)^2=+\infty,
\]
故可取 \(V>U\) 使
\[
2\sum_{u=U+1}^{V}C_\varphi(u)^2>M.
\]
另一方面，推论 \ref{cor:fold-collision-gap-resonance-ladder-U} 给出
\[
\liminf_{m\to\infty}F_{m+2}\,\mathsf{Col}_m
\ge
1+2\sum_{u=1}^{V}C_\varphi(u)^2.
\]
两边同减去 \(1\)，得到
\[
\liminf_{m\to\infty}\Gamma_m
\ge
2\sum_{u=1}^{V}C_\varphi(u)^2.
\]
再减去 \(R_U\) 即得
\[
\liminf_{m\to\infty}\bigl(\Gamma_m-R_U\bigr)
\ge
2\sum_{u=U+1}^{V}C_\varphi(u)^2
>
M.
\]
由于 \(M\) 任意，结论成立。证毕。
\end{proof}

\begin{theorem}[零点计数不足以控制真实均匀性缺口]\label{thm:xi-time-part75a-zero-counting-insufficient-for-uniformity-gap}
\leanverified{paper\_xi\_time\_part75a\_zero\_counting\_insufficient\_for\_uniformity\_gap}
不存在函数 \(\omega:[0,1]\to[0,\infty)\) 与数列 \((\varepsilon_m)_{m\ge 1}\) 满足
\[
\varepsilon_m\downarrow 0,
\qquad
\omega(t)\xrightarrow[t\downarrow 0]{}0,
\]
并且对所有充分大的 \(m\) 都有
\[
\Gamma_m
\le
\omega\!\left(\frac{|\mathcal Z_m|}{F_{m+2}}\right)+\varepsilon_m,
\]
其中 \(\Gamma_m\) 取自定理 \ref{thm:xi-time-part75a-collision-gap-infinite-resonance-rank}，而
\[
\mathcal Z_m:=\{k\in \ZZ/(F_{m+2}\ZZ):\ \widehat d_m(k)=0\}
\]
为谱零点集合。换言之，谱零点密度的衰减不能单独上界真实的碰撞非均匀性。
\end{theorem}
\begin{proof}
记
\[
\delta_m:=\frac{|\mathcal Z_m|}{F_{m+2}}.
\]
若 \(3\nmid(m+2)\)，则由定理 \ref{thm:killo-fold-zero-spectrum-affine-congruence}，
\[
\mathcal Z_m=\varnothing,
\qquad
\delta_m=0.
\]
若 \(3\mid(m+2)\)，则由定理 \ref{thm:fold-zero-density-sparse}，
\[
\delta_m
\le
\tau(m+2)\,O(\varphi^{-m/2})
\longrightarrow 0.
\]
因此沿全序列都有 \(\delta_m\to 0\)。

另一方面，由推论 \ref{cor:fold-critical-resonance-gap}，
\[
\mathsf{Col}_m
\ge
\frac{1+2C_\varphi^2+o(1)}{F_{m+2}},
\]
从而
\[
\Gamma_m
=
F_{m+2}\!\left(\mathsf{Col}_m-\frac{1}{F_{m+2}}\right)
\ge
2C_\varphi^2+o(1).
\]
由于 \(C_\varphi>0\)，故 \(\liminf_{m\to\infty}\Gamma_m>0\)。

若存在题设中的 \(\omega\) 与 \(\varepsilon_m\)，则由 \(\delta_m\to 0\)、\(\omega(t)\to 0\) 与 \(\varepsilon_m\to 0\) 可得
\[
\omega(\delta_m)+\varepsilon_m\longrightarrow 0.
\]
这将迫使 \(\Gamma_m\to 0\)，与上面的严格正下极限矛盾。故所述控制不可能成立。证毕。
\end{proof}

\begin{theorem}[持续歧义迫使自然扩张需要无限尾寄存器]\label{thm:xi-time-part75a-persistent-ambiguity-forces-infinite-tail}
\leanverified{paper\_xi\_time\_part75a\_persistent\_ambiguity\_forces\_infinite\_tail}
设 \(T:X\twoheadrightarrow X\) 为有限纤维满射，并记其自然扩张为
\[
X^{\leftarrow}
:=
\left\{
(x_0,x_{-1},x_{-2},\dots)\in X^{\NN}:\ T(x_{-j})=x_{-j+1}\ \ (j\ge 1)
\right\}.
\]
对每个 \(y\in X\) 固定一组纤维单射
\[
\iota_y:T^{-1}(y)\hookrightarrow \NN,
\]
并记
\[
\beta(x):=\iota_{T(x)}(x).
\]
于是自然扩张的分支指标编码可写为
\[
\mathcal C(x_0,x_{-1},x_{-2},\dots)
:=
\bigl(x_0,\beta(x_{-1}),\beta(x_{-2}),\dots\bigr).
\]
对 \(L\ge 1\) 定义有限尾截断
\[
\Pi_L(x_0,x_{-1},x_{-2},\dots)
:=
\bigl(x_0,\beta(x_{-1}),\dots,\beta(x_{-L})\bigr).
\]
若存在某条逆轨道
\[
\mathbf x=(x_0,x_{-1},x_{-2},\dots)\in X^{\leftarrow}
\]
满足
\[
\#T^{-1}(x_{-j})\ge 2
\qquad
\text{对无穷多个 }j,
\]
则对每个 \(L\ge 1\)，映射 \(\Pi_L\) 都不是单射。
\end{theorem}
\begin{proof}
固定 \(L\ge 1\)。由假设，可取 \(n>L\) 使
\[
\#T^{-1}(x_{-n+1})\ge 2.
\]
于是存在不同点
\[
y_n\neq y_n',
\qquad
T(y_n)=T(y_n')=x_{-n+1}.
\]
将原逆轨道的前 \(n-1\) 层固定为
\[
x_0,x_{-1},\dots,x_{-n+1},
\]
再利用 \(T\) 的满射性，分别从 \(y_n\) 与 \(y_n'\) 继续向后选取无限前史，得到两条逆轨道
\[
\mathbf y,\mathbf y'\in X^{\leftarrow}.
\]
按构造，\(\mathbf y\) 与 \(\mathbf y'\) 在层深 \(1,\dots,n-1\) 上完全一致，只在第 \(n\) 层首次分叉。由于 \(n>L\)，它们的前 \(L\) 个分支指标完全相同，从而
\[
\Pi_L(\mathbf y)=\Pi_L(\mathbf y').
\]
但 \(\mathbf y\neq \mathbf y'\)，故 \(\Pi_L\) 非单射。证毕。
\end{proof}

\begin{corollary}[分层外置资源的有限--无限深度相变]\label{cor:xi-time-part75a-layered-externalization-depth-phase-transition}
\leanverified{paper\_xi\_time\_part75a\_layered\_externalization\_depth\_phase\_transition}
在定理 \ref{thm:xi-layered-prime-register-sharp-natural-extension} 的分层 prime-register 语义下，若有限深度 \(k\) 的每一层都存在至少一个二歧义纤维，则最坏情形下的精确外置最小成本恰为 \(k\) 个有效 prime slices，且该下界可以达到。

相反，若满足定理 \ref{thm:xi-time-part75a-persistent-ambiguity-forces-infinite-tail} 的持续歧义条件，则不存在任何只使用有限多个活动 prime slices 的精确实现。等价地，
\[
\text{有限深度 }k
\quad\Longrightarrow\quad
\text{最小 prime-slice 成本 }=k,
\]
\[
\text{无限持续深度}
\quad\Longrightarrow\quad
\text{不存在有限成本精确实现}.
\]
\end{corollary}
\begin{proof}
有限深度部分正是定理 \ref{thm:xi-layered-prime-register-sharp-natural-extension} 的内容：在每一层都存在真实分支时，精确分层外置的最坏情形下界恰为层数 \(k\)，并且逐层写入一枚有效素数的构造达到该下界。

若持续歧义沿某条逆轨道无穷出现，则定理 \ref{thm:xi-time-part75a-persistent-ambiguity-forces-infinite-tail} 表明：任意有限前缀截断 \(\Pi_L\) 都失去单射性。对 prime-register 账本而言，这意味着任何只保留有限层活动 slices 的实现都不足以忠实恢复全部前史；其一般形式的不可能性正由定理 \ref{thm:xi-time-part63b-infinite-depth-finite-prime-register-impossibility} 给出。故无限持续深度不再允许任何有限成本的精确实现。证毕。
\end{proof}

\begin{theorem}[单一幂散度标量决定整个极限散度塔]\label{thm:xi-time-part75a-single-power-divergence-determines-limit-tower}
\leanverified{paper\_xi\_time\_part75a\_single\_power\_divergence\_determines\_limit\_tower}
对任意 \(\alpha>1\)，记
\[
f_\alpha(t):=t^\alpha-1-\alpha(t-1),
\qquad
\mathcal D_\alpha:=D_{f_\alpha}^\infty.
\]
则成立：
\begin{enumerate}
  \item
  \[
  \mathcal D_\alpha
  =
  \frac{1}{5^{\alpha/2}}
  \bigl(\varphi^{2\alpha-1}+\varphi^{\alpha-2}\bigr)-1;
  \]
  \item 对固定 \(\alpha>1\)，映射
  \[
  \xi\longmapsto
  \frac{1}{5^{\alpha/2}}
  \bigl(\xi^{2\alpha-1}+\xi^{\alpha-2}\bigr)-1
  \qquad (\xi>1)
  \]
  严格递增，因此单一标量 \(\mathcal D_\alpha\) 即可唯一反演 \(\varphi\)；
  \item 一旦某个固定 \(\alpha>1\) 的 \(\mathcal D_\alpha\) 已知，则对一切 \(\beta>1\)，全部极限散度
  \[
  \mathcal D_\beta
  \]
  都被唯一确定；并且似然比随机变量 \(\mu_m/u_m\) 的极限两点律
  \[
  \frac{\mu_m}{u_m}
  \Longrightarrow
  \frac1\varphi\,\delta_{\varphi^2/\sqrt5}
  +
  \frac1{\varphi^2}\,\delta_{\varphi/\sqrt5}
  \]
  也被同一标量唯一锁定。
\end{enumerate}
\end{theorem}
\begin{proof}
前两条正是命题 \ref{prop:fold-bin-phi-identifiable-power-divergence} 的闭式与严格单调反演。

若某个固定 \(\alpha>1\) 的 \(\mathcal D_\alpha\) 已知，则由第 \(2\) 条可唯一恢复 \(\varphi\)。将该 \(\varphi\) 代回第 \(1\) 条，便得到对每个 \(\beta>1\) 的唯一值
\[
\mathcal D_\beta
=
\frac{1}{5^{\beta/2}}
\bigl(\varphi^{2\beta-1}+\varphi^{\beta-2}\bigr)-1.
\]
这证明了整个幂散度极限塔的唯一决定性。

最后，命题 \ref{prop:fold-bin-f-divergence-collapse} 与命题 \ref{prop:xi-foldbin-f-divergence-two-point-experiment} 已把 bin-fold 的渐近 \(f\)-散度几何统一压缩为唯一二点实验，其似然比原子正为
\[
\frac{\varphi^2}{\sqrt5},
\qquad
\frac{\varphi}{\sqrt5},
\]
相应权重为
\[
\frac1\varphi,
\qquad
\frac1{\varphi^2}.
\]
因此该两点极限律亦由同一 \(\varphi\) 唯一决定。证毕。
\end{proof}

\begin{theorem}[冻结区 escort 测度的微正则塌缩]\label{thm:xi-time-part75a-frozen-escort-microcanonical-collapse}
设 \(a,b>a_{\mathrm c}\)，并沿用定理 \ref{thm:xi-fixed-freezing-escort-renyi-spectrum-collapse} 的记号。记
\[
A_m^{\max}:=\{x\in X_m:\ d_m(x)=M_m\},
\qquad
U_m^{\max}:=\mathrm{Unif}(A_m^{\max}).
\]
则对 \(t\in\{a,b\}\)，存在概率测度 \(\nu_{m,t}\) 使
\[
\pi_{m,t}
=
p_{m,t}\,U_m^{\max}
+
\bigl(1-p_{m,t}\bigr)\nu_{m,t},
\qquad
p_{m,t}:=\pi_{m,t}(A_m^{\max}),
\]
并且
\[
1-p_{m,t}
\le
\exp\!\bigl(-m\Delta(t)+o(m)\bigr).
\]
从而
\[
\|\pi_{m,a}-\pi_{m,b}\|_{\mathrm{TV}}
\le
\exp\!\bigl(-m\Delta(a)+o(m)\bigr)
+
\exp\!\bigl(-m\Delta(b)+o(m)\bigr).
\]
换言之，冻结区内全部 escort 测度都以指数速度塌缩到同一个随 \(m\) 变化的微正则核心 \(U_m^{\max}\)。
\end{theorem}
\begin{proof}
对任意 \(t>a_{\mathrm c}\)，由于 \(\pi_{m,t}(x)\propto d_m(x)^t\)，而 \(d_m(x)\) 在 \(A_m^{\max}\) 上恒等于 \(M_m\)，故 \(\pi_{m,t}\) 在 \(A_m^{\max}\) 上的条件分布必为均匀分布 \(U_m^{\max}\)。于是存在概率测度 \(\nu_{m,t}\) 使
\[
\pi_{m,t}
=
p_{m,t}\,U_m^{\max}
+
\bigl(1-p_{m,t}\bigr)\nu_{m,t},
\qquad
p_{m,t}=\pi_{m,t}(A_m^{\max}).
\]

再由定理 \ref{thm:xi-fixed-freezing-escort-renyi-spectrum-collapse} 第 \(4\) 条，
\[
1-p_{m,t}
=
1-\pi_{m,t}(A_m^{\max})
\le
\exp\!\bigl(-m\Delta(t)+o(m)\bigr).
\]
另一方面，
\[
\|\pi_{m,t}-U_m^{\max}\|_{\mathrm{TV}}
=
1-p_{m,t},
\]
故由三角不等式，
\[
\|\pi_{m,a}-\pi_{m,b}\|_{\mathrm{TV}}
\le
\|\pi_{m,a}-U_m^{\max}\|_{\mathrm{TV}}
+
\|U_m^{\max}-\pi_{m,b}\|_{\mathrm{TV}}
=
(1-p_{m,a})+(1-p_{m,b}).
\]
代入上述指数界即得结论。证毕。
\end{proof}

\begin{corollary}[冻结压力射线的两点确定律]\label{cor:xi-time-part75a-freezing-pressure-ray-two-point-determination}
\leanverified{paper\_xi\_time\_part75a\_freezing\_pressure\_ray\_two\_point\_determination}
任取 \(a_0\neq a_1>a_{\mathrm c}\)。若已知对应压力 \(P_{a_0},P_{a_1}\)，则对一切 \(a>a_{\mathrm c}\)，整条冻结压力射线
\[
P_a
\]
以及共同的基态简并指数
\[
h_\infty(a)
\]
都被唯一确定。更具体地，
\[
P_a
=
\frac{a_1-a}{a_1-a_0}\,P_{a_0}
+
\frac{a-a_0}{a_1-a_0}\,P_{a_1},
\]
且
\[
h_\infty(a)
=
P_{a_0}-a_0\frac{P_{a_1}-P_{a_0}}{a_1-a_0}
=
P_{a_1}-a_1\frac{P_{a_1}-P_{a_0}}{a_1-a_0}.
\]
\end{corollary}
\begin{proof}
由定理 \ref{thm:fold-pressure-freezing-threshold}，对每个 \(a>a_{\mathrm c}\) 都有
\[
P_a=a\alpha_*+g_*.
\]
因此两点 \((a_0,P_{a_0})\)、\((a_1,P_{a_1})\) 唯一确定仿射函数的斜率与截距，即
\[
\alpha_*=\frac{P_{a_1}-P_{a_0}}{a_1-a_0},
\qquad
g_*=P_{a_0}-a_0\frac{P_{a_1}-P_{a_0}}{a_1-a_0}.
\]
把它们代回即可得到所示插值公式。

另一方面，由定理 \ref{thm:xi-fixed-freezing-pressure-address-cost-residual-decomposition}，冻结相的 min-entropy 率满足
\[
h_\infty(a)=g_*.
\]
故 \(h_\infty(a)\) 亦由同一对数据唯一决定，并与 \(a\) 无关。证毕。
\end{proof}


\paragraph{Lee--Yang branch 图的超立方体封闭、Smith 满射筛、椭圆原子长度与共振临界点云}
在有限层 Lee--Yang branch 图的显式超立方体几何、Smith 核增长的全模数平台、整数轴对齐椭圆幺半群的自由原子分解，以及 Fibonacci 共振整函数的实零点结构已经建立之后，还可继续抽取出下列五条直接后果。它们分别把 branch 图的超立方体谱与 Ihara--Bass 数据整理为统一接口，把模数满射性压缩为最大 Smith 不变量上的素数筛条件，把整数轴对齐椭圆的素数原子长度与右可除偏序写成显式算术公式，并把共振整函数导函数的临界点云锁定到与原零点云相同的线性密度和缩放均匀律。

\begin{theorem}[Lee--Yang 第 \texorpdfstring{$n$}{n} 层分支图的超立方体谱与热迹封闭]\label{thm:xi-time-part76-leyang-branch-hypercube-spectrum-heat}
\leanverified{paper\_xi\_time\_part76\_leyang\_branch\_hypercube\_spectrum\_heat}
沿用定理 \ref{thm:xi-time-part62d-leyang-branch-graph-gaussian-spectrum} 的记号，记邻接算子为 \(A_n\)，组合 Laplacian 为
\[
L_n:=2nI-A_n.
\]
则成立：
\begin{enumerate}
  \item 有图同构
  \[
  G_n\cong \bigsqcup_{i=0}^{4}Q_{2n}.
  \]
  \item 邻接谱满足
  \[
  \operatorname{Spec}(A_n)=\{\,2n-2j:\ 0\le j\le 2n\,\},
  \qquad
  \operatorname{mult}_{A_n}(2n-2j)=5\binom{2n}{j}.
  \]
  \item 对任意 \(t\ge 0\)，热迹具有闭式
  \[
  \Tr\!\bigl(e^{-tL_n}\bigr)=5\sum_{j=0}^{2n}\binom{2n}{j}e^{-2jt}
  =
  5\bigl(1+e^{-2t}\bigr)^{2n}.
  \]
\end{enumerate}
\end{theorem}
\begin{proof}
定理 \ref{thm:xi-time-part62d-leyang-branch-graph-gaussian-spectrum} 的地址识别已经给出：每个分量 \(\Pi_{i,n}\) 都与 \(Q_{2n}\) 图同构，故
\[
G_n\cong \bigsqcup_{i=0}^{4}Q_{2n}.
\]
这证明了第一条。

对单个 \(Q_{2n}\)，Walsh 角色基对角化其邻接算子，得到特征值
\[
2n-2j\qquad(0\le j\le 2n)
\]
及重数 \(\binom{2n}{j}\)。五个分量彼此不交，因此整体图 \(G_n\) 的相应重数恰乘以 \(5\)，得到第二条。

最后，由
\[
L_n=2nI-A_n
\]
可知 \(L_n\) 的特征值为 \(2j\)（\(0\le j\le 2n\)），重数仍为 \(5\binom{2n}{j}\)。于是
\[
\Tr\!\bigl(e^{-tL_n}\bigr)
=
5\sum_{j=0}^{2n}\binom{2n}{j}e^{-2jt}.
\]
再对二项式展开，即得
\[
\Tr\!\bigl(e^{-tL_n}\bigr)=5\bigl(1+e^{-2t}\bigr)^{2n}.
\]
证毕。
\end{proof}

\begin{corollary}[Lee--Yang 第 \texorpdfstring{$n$}{n} 层分支图的 Ihara--Bass 闭式]\label{cor:xi-time-part76-leyang-branch-ihara-bass-closed}
沿用定理 \ref{thm:xi-time-part76-leyang-branch-hypercube-spectrum-heat} 的记号。则 \(G_n\) 的 Ihara \(\zeta\) 函数满足
\[
\bigl(Z_{G_n}^{\mathrm{Ih}}(u)\bigr)^{-1}
=
(1-u^2)^{5(2n-2)2^{2n-1}}
\prod_{j=0}^{2n}
\bigl(1-(2n-2j)u+(2n-1)u^2\bigr)^{5\binom{2n}{j}}.
\]
\end{corollary}
\begin{proof}
由定理 \ref{thm:xi-time-part76-leyang-branch-hypercube-spectrum-heat}，
\[
G_n\cong \bigsqcup_{i=0}^{4}Q_{2n}.
\]
Ihara \(\zeta\) 函数对不交并取乘法，因此只需计算单个 \(Q_{2n}\) 的情形。对 \(d\)-正则图 \(Q_d\)，Ihara--Bass 行列式公式给出
\[
\bigl(Z_{Q_d}^{\mathrm{Ih}}(u)\bigr)^{-1}
=(1-u^2)^{(d-2)2^{d-1}}
\prod_{j=0}^{d}
\bigl(1-(d-2j)u+(d-1)u^2\bigr)^{\binom{d}{j}},
\]
其中已代入 \(Q_d\) 的邻接谱。取 \(d=2n\) 后再取五次方，即得所示闭式。证毕。
\end{proof}

\begin{theorem}[Smith 模满射性的最大不变量筛律]\label{thm:xi-time-part76-smith-surjective-moduli-prime-sieve}
设 \(A\in M_{m\times n}(\ZZ)\) 满足 \(\operatorname{rank}_{\QQ}(A)=m\)，其 Smith 不变量为
\[
d_1\mid d_2\mid \cdots\mid d_m.
\]
对任意 \(b\ge 1\)，记模 \(b\) 诱导映射为
\[
A_b:(\ZZ/b\ZZ)^n\longrightarrow (\ZZ/b\ZZ)^m.
\]
则以下条件等价：
\begin{enumerate}
  \item \(A_b\) 为满射；
  \item \(\prod_{i=1}^{m}\gcd(d_i,b)=1\)；
  \item \(\gcd(d_m,b)=1\)。
\end{enumerate}
进一步，记
\[
D:=\operatorname{rad}(d_m).
\]
则满射模数组成的集合
\[
\mathcal S_A:=\{\,b\ge 1:\ A_b\ \text{为满射}\,\}
\]
具有自然密度
\[
\delta(\mathcal S_A)=\frac{\varphi(D)}{D},
\]
并且满足精确计数
\[
\#(\mathcal S_A\cap [1,X])
=
\frac{\varphi(D)}{D}X+O(1)
\qquad (X\to\infty).
\]
\end{theorem}
\begin{proof}
由推论 \ref{cor:killo-kernel-size-general-modulus}，
\[
\#\ker(A_b)=b^{n-m}\prod_{i=1}^{m}\gcd(d_i,b).
\]
于是
\[
\#\operatorname{im}(A_b)
=
\frac{b^n}{\#\ker(A_b)}
=
\frac{b^m}{\prod_{i=1}^{m}\gcd(d_i,b)}.
\]
因此 \(\#\operatorname{im}(A_b)=b^m\) 当且仅当
\[
\prod_{i=1}^{m}\gcd(d_i,b)=1,
\]
这就得到第 \(1\iff 2\) 条。

由于 \(d_1\mid \cdots\mid d_m\)，一旦 \(\gcd(d_m,b)=1\)，便自动有
\[
\gcd(d_i,b)=1\qquad(1\le i\le m),
\]
故第 \(2\iff 3\) 条亦成立。

于是
\[
\mathcal S_A=\{\,b\ge 1:\ \gcd(b,d_m)=1\,\}
=
\{\,b\ge 1:\ \gcd(b,D)=1\,\}.
\]
而指标函数 \(1_{\gcd(\,\cdot\,,D)=1}\) 以 \(D\) 为周期，且在每个完整模 \(D\) 周期中恰有 \(\varphi(D)\) 个取值为 \(1\)。因此
\[
\#(\mathcal S_A\cap [1,X])
=
\frac{\varphi(D)}{D}X+O(1).
\]
自然密度公式随之立即成立。证毕。
\end{proof}

\begin{theorem}[整数轴对齐椭圆幺半群的原子长度与右可除偏序]\label{thm:xi-time-part76-integer-ellipse-atomic-length-divisibility}
\leanverified{paper\_xi\_time\_part76\_integer\_ellipse\_atomic\_length\_divisibility}
沿用定理 \ref{thm:xi-integer-ellipse-free-commutative-monoid} 的记号。对每个素数 \(p\in\mathcal P\)，定义
\[
\mathsf e_p:=E_{p,1},
\qquad
\mathsf f_p:=E_{1,p},
\qquad
\mathcal A:=\{\,\mathsf e_p,\mathsf f_p:\ p\in\mathcal P\,\}.
\]
则对任意 \(E_{a,b}\in\mathsf E_{\NN}\)，都有唯一分解
\[
E_{a,b}
=
\left(\bigodot_{p\in\mathcal P}\mathsf e_p^{\odot v_p(a)}\right)
\odot
\left(\bigodot_{p\in\mathcal P}\mathsf f_p^{\odot v_p(b)}\right),
\]
其中仅有限多个因子不同于 \(E_{1,1}\)。此外，若 \(\ell_{\mathcal A}\) 表示相对于生成系 \(\mathcal A\) 的最短词长，则
\[
\ell_{\mathcal A}(E_{a,b})=\Omega(a)+\Omega(b),
\qquad
\Omega(n):=\sum_{p\in\mathcal P}v_p(n).
\]
最后，若定义右可除偏序
\[
E_{a,b}\preceq E_{c,d}
\iff
\exists\,E_{u,v}\in\mathsf E_{\NN}\ \text{使}\ E_{c,d}=E_{a,b}\odot E_{u,v},
\]
则
\[
E_{a,b}\preceq E_{c,d}
\iff
a\mid c\ \text{且}\ b\mid d.
\]
\end{theorem}
\begin{proof}
定理 \ref{thm:xi-integer-ellipse-free-commutative-monoid} 已给出幺半群同构
\[
(\NN^{(\mathcal P)}\times \NN^{(\mathcal P)},+)
\cong
(\mathsf E_{\NN},\odot),
\qquad
(r_1,r_2)\longmapsto E_{N(r_1),N(r_2)}.
\]
在该坐标下，\(\mathsf e_p\) 与 \(\mathsf f_p\) 分别对应基向量 \((\delta_p,0)\) 与 \((0,\delta_p)\)。因此
\[
\bigl((v_p(a))_{p},(v_p(b))_{p}\bigr)
=
\sum_{p}v_p(a)(\delta_p,0)+\sum_{p}v_p(b)(0,\delta_p)
\]
是自由交换幺半群中的唯一基分解；传回 \((\mathsf E_{\NN},\odot)\) 即得所示唯一原子分解。

同样在该自由坐标下，相对于基向量族 \(\{(\delta_p,0),(0,\delta_p)\}\) 的最短词长恰为 \(\ell^1\)-范数，故
\[
\ell_{\mathcal A}(E_{a,b})
=
\sum_{p}v_p(a)+\sum_{p}v_p(b)
=
\Omega(a)+\Omega(b).
\]

最后，
\[
E_{c,d}=E_{a,b}\odot E_{u,v}
\iff
(c,d)=(au,bv),
\]
而这等价于 \(a\mid c\) 且 \(b\mid d\)。证毕。
\end{proof}

\begin{theorem}[共振整函数导函数的实临界点交错与缩放均匀律]\label{thm:xi-time-part76-resonance-derivative-critical-cloud-uniform}
沿用定理 \ref{thm:fold-resonance-entire-lp} 的记号，记
\[
\mathcal F_\varphi(z):=\prod_{n=2}^{\infty}\cos(\pi z\varphi^{-n}),
\qquad
N_\varphi'(R):=
\#\bigl(\mathcal Z(\mathcal F_\varphi')\cap[-R,R]\bigr).
\]
则成立：
\begin{enumerate}
  \item \(\mathcal F_\varphi'\) 的全部零点都实且单，并与 \(\mathcal F_\varphi\) 的零点严格交错；
  \item
  \[
  N_\varphi'(R)=2R+O(\log R)
  \qquad(R\to\infty);
  \]
  \item 若定义缩放经验测度
  \[
  \nu_R'
  :=
  \frac{1}{N_\varphi'(R)}
  \sum_{\substack{\xi\in\mathcal Z(\mathcal F_\varphi')\\ |\xi|\le R}}
  \delta_{\xi/R},
  \]
  则
  \[
  \nu_R'\Longrightarrow \frac12\,\mathbf 1_{[-1,1]}(x)\,\dd x
  \qquad (R\to\infty).
  \]
\end{enumerate}
\end{theorem}
\begin{proof}
由定理 \ref{thm:fold-resonance-entire-lp}，\(\mathcal F_\varphi\) 属于 Laguerre--P\'olya 类，且其零点全部为实单零点。经典 Laguerre 分离定理遂给出：\(\mathcal F_\varphi'\) 的零点也全部为实单零点，并可按增序排列为
\[
\cdots<\xi_{-1}<\xi_0<\xi_1<\cdots
\]
使得若
\[
\cdots<\zeta_{-1}<\zeta_0<\zeta_1<\cdots
\]
是 \(\mathcal F_\varphi\) 的零点排列，则对每个 \(j\in\ZZ\) 都有
\[
\zeta_j<\xi_j<\zeta_{j+1}.
\]
这证明了第一条。

由上述严格交错，对任意有界闭区间 \(I\subset\RR\)，两类零点在 \(I\) 内的个数至多只会因两端边界效应相差常数。更准确地，
\[
\Bigl|
\#\bigl(\mathcal Z(\mathcal F_\varphi')\cap I\bigr)
-
\#\bigl(\mathcal Z(\mathcal F_\varphi)\cap I\bigr)
\Bigr|
\le 2.
\]
取 \(I=[-R,R]\)，再由定理 \ref{cor:xi-resonance-entire-zero-count-explicit-linear-density}，
\[
\#\bigl(\mathcal Z(\mathcal F_\varphi)\cap[-R,R]\bigr)=2R+O(\log R),
\]
便得
\[
N_\varphi'(R)=2R+O(\log R).
\]
这证明了第二条。

现取任意固定区间 \([a,b]\subset(-1,1)\)。由定理 \ref{thm:fold-resonance-entire-lp} 的显式零点集
\[
\mathcal Z(\mathcal F_\varphi)
=
\Bigl\{\bigl(k+\tfrac12\bigr)\varphi^n:\ k\in\ZZ,\ n\ge 2\Bigr\},
\]
对每个固定层 \(n\ge 2\)，落在 \([aR,bR]\) 中的零点数为
\[
(b-a)R\varphi^{-n}+O(1).
\]
只有 \(O(\log R)\) 个层可能贡献非零项，而
\[
\sum_{n=2}^{\infty}\varphi^{-n}=1,
\]
故求和得到
\[
\#\bigl(\mathcal Z(\mathcal F_\varphi)\cap[aR,bR]\bigr)
=(b-a)R+O(\log R).
\]
再由交错性带来的 \(O(1)\) 计数差，得到
\[
\#\bigl(\mathcal Z(\mathcal F_\varphi')\cap[aR,bR]\bigr)
=(b-a)R+O(\log R).
\]
与第二条合并可得
\[
\nu_R'([a,b])
=
\frac{(b-a)R+O(\log R)}{2R+O(\log R)}
\longrightarrow
\frac{b-a}{2}.
\]
由于 \(\frac12\mathbf 1_{[-1,1]}(x)\,\dd x\) 对区间端点不赋质量，Portmanteau 定理即给出
\[
\nu_R'\Longrightarrow \frac12\,\mathbf 1_{[-1,1]}(x)\,\dd x.
\]
证毕。
\end{proof}

\noindent
这五条结果把 Lee--Yang branch 图、Smith 模层、整数轴对齐椭圆与 Fibonacci 共振整函数再度压到同一条可直接调用的闭链：有限层 branch 图在图论层面不只保持超立方体式的邻接谱，而且连 Ihara--Bass 数据也完全封闭；模数满射性则不再需要逐个 Smith 因子检查，而只由最大不变量的素数支撑控制；整数轴对齐椭圆的素数原子坐标进一步把几何对象族转写成显式的长度与可除序公式；最后，共振整函数的导函数临界点云并未偏离原零点云的宏观几何，而是以同样的线性密度与区间均匀律继承了整个实轴上的缩放统计。

\paragraph{单原子 Dirichlet 壳、residue 有限部与 cyclotomic 几何}
在单原子因子的周期--primitive 精确剥离、periodic residue 的 Fourier 反演以及 full/core residue 差的 rank-one 局域化已经建立之后，还可把同一原子在 Dirichlet 层继续压缩为一组完全闭式的后果。它们分别把 twisted 通道上的 ghost 贡献识别为精确 polylog 壳，把 residue 类上的 Dirichlet 级数写成 Hurwitz--cyclotomic 壳并统一主极点残数，把 residue-wise 有限部常数的原子位移压缩为显式 digamma 项，并把 centered residue 扰动提升为零均值超平面上的等角紧框架及其精确投影能量律。

\begin{theorem}[single atom 的 twisted Dirichlet 壳是一个精确 polylog 壳]\label{thm:xi-time-part77-atomic-twisted-dirichlet-polylog-shell}
沿用定理 \ref{thm:xi-atomic-witt-cycle-primitive-exact-peeling} 的记号。对固定模数 \(m\ge 2\) 与扭曲通道 \(j\in \ZZ/m\ZZ\)，定义原子 periodic Dirichlet 级数
\[
\mathcal D^{\mathrm{cycle}}_{m,j}(s)
:=
\sum_{n\ge 1}a_n^{\mathrm{cycle}}(\omega_m^j)\,n^{-s}
\qquad (\Re s>1).
\]
则有完全闭式
\[
\mathcal D^{\mathrm{cycle}}_{m,j}(s)
:=
2^{1-s}\operatorname{Li}_s(\omega_m^j),
\]
其中 \(j=0\) 时按 \(\operatorname{Li}_s(1)=\zeta(s)\) 解释。特别地：
\begin{enumerate}
  \item 当 \(j=0\) 时，
  \[
  \mathcal D^{\mathrm{cycle}}_{m,0}(s)=2^{1-s}\zeta(s),
  \]
  因而在 \(s=1\) 处具有简单极点；
  \item 当 \(1\le j\le m-1\) 时，\(\mathcal D^{\mathrm{cycle}}_{m,j}(s)\) 解析延拓为整个函数。
\end{enumerate}
\end{theorem}
\begin{proof}
由定理 \ref{thm:xi-atomic-witt-cycle-primitive-exact-peeling}，
\[
a_n^{\mathrm{cycle}}(u)=
\begin{cases}
2u^{n/2},& 2\mid n,\\
0,& 2\nmid n.
\end{cases}
\]
代入 \(u=\omega_m^j\) 得
\[
a_{2\ell}^{\mathrm{cycle}}(\omega_m^j)=2\omega_m^{j\ell},
\qquad
a_{2\ell+1}^{\mathrm{cycle}}(\omega_m^j)=0.
\]
因此
\[
\mathcal D^{\mathrm{cycle}}_{m,j}(s)
:=
\sum_{\ell\ge 1}2\,\omega_m^{j\ell}(2\ell)^{-s}
:=
2^{1-s}\sum_{\ell\ge 1}\frac{\omega_m^{j\ell}}{\ell^s}
:=
2^{1-s}\operatorname{Li}_s(\omega_m^j),
\]
得到所示闭式。

当 \(j=0\) 时，上式退化为
\[
\mathcal D^{\mathrm{cycle}}_{m,0}(s)=2^{1-s}\zeta(s),
\]
故在 \(s=1\) 处有简单极点。

现设 \(j\neq 0\)。标准 root-of-unity 分解给出
\[
\operatorname{Li}_s(\omega_m^j)
:=
m^{-s}\sum_{r=1}^{m}\omega_m^{jr}\,\zeta\!\left(s,\frac{r}{m}\right).
\]
对每个 \(r\)，Hurwitz zeta \(\zeta(s,r/m)\) 的唯一极点都位于 \(s=1\)，且留数同为 \(1\)。但
\[
\sum_{r=1}^{m}\omega_m^{jr}=0
\qquad (j\neq 0),
\]
故 \(s=1\) 处的极点严格相消，从而 \(\operatorname{Li}_s(\omega_m^j)\) 解析延拓为整个函数。乘上整因子 \(2^{1-s}\) 后，\(\mathcal D^{\mathrm{cycle}}_{m,j}(s)\) 亦为整个函数。证毕。
\end{proof}

\begin{theorem}[single atom 的 periodic residue Dirichlet 级数是 Hurwitz--cyclotomic 壳，且全部 residue 主极点残数完全一致]\label{thm:xi-time-part77-atomic-residue-dirichlet-hurwitz-shell}
沿用定理 \ref{thm:xi-atomic-witt-energy-residue-l2-law} 与结论 \ref{con:xi-time-part57ba-atomic-residue-delta-flat-extremizer} 的记号。对固定模数 \(m\ge 2\) 与 residue 类 \(r\in \ZZ/m\ZZ\)，定义
\[
\mathcal A^{\mathrm{cycle}}_{m,r}(s)
:=
\sum_{n\ge 1}A^{\mathrm{cycle}}_{m,r}(n)\,n^{-s}
\qquad (\Re s>1).
\]
当 \(r\neq 0\) 时，下文把同一符号 \(r\) 也用于其唯一代表元 \(r\in\{1,\dots,m-1\}\)。
则成立：
\begin{enumerate}
  \item 对 \(r=0\)，有
  \[
  \mathcal A^{\mathrm{cycle}}_{m,0}(s)
  =
  2^{1-s}m^{-s}\zeta(s).
  \]
  \item 对 \(1\le r\le m-1\)，有
  \[
  \mathcal A^{\mathrm{cycle}}_{m,r}(s)
  =
  2^{1-s}m^{-s}\zeta\!\left(s,\frac{r}{m}\right).
  \]
\end{enumerate}
特别地，对每个 residue 类 \(r\in \ZZ/m\ZZ\)，\(\mathcal A^{\mathrm{cycle}}_{m,r}(s)\) 在 \(s=1\) 处都有简单极点，并且
\[
\operatorname*{Res}_{s=1}\mathcal A^{\mathrm{cycle}}_{m,r}(s)=\frac{1}{m}.
\]
\end{theorem}
\begin{proof}
由结论 \ref{con:xi-time-part57ba-atomic-residue-delta-flat-extremizer}，
\[
A_{m,r}^{\mathrm{cycle}}(n)
:=
2\,\mathbf 1_{\{2\mid n\}}\mathbf 1_{\{r\equiv n/2\;(\mathrm{mod}\,m)\}}.
\]
若 \(r=0\)，则非零项恰对应于 \(n=2mq\)（\(q\ge 1\)），故
\[
\mathcal A^{\mathrm{cycle}}_{m,0}(s)
:=
2\sum_{q\ge 1}(2mq)^{-s}
:=
2^{1-s}m^{-s}\zeta(s).
\]

若 \(1\le r\le m-1\)，则非零项恰对应于
\[
n=2(r+mq)\qquad (q\ge 0),
\]
故
\[
\mathcal A^{\mathrm{cycle}}_{m,r}(s)
:=
2\sum_{q\ge 0}(2(r+mq))^{-s}
:=
2^{1-s}m^{-s}\sum_{q\ge 0}\left(q+\frac{r}{m}\right)^{-s}
:=
2^{1-s}m^{-s}\zeta\!\left(s,\frac{r}{m}\right).
\]
这就得到前两条公式。

已知 \(\zeta(s)\) 与全部 Hurwitz zeta \(\zeta(s,a)\)（\(0<a\le 1\)）在 \(s=1\) 的留数都等于 \(1\)。于是
\[
\operatorname*{Res}_{s=1}\mathcal A^{\mathrm{cycle}}_{m,r}(s)
:=
\left.2^{1-s}m^{-s}\right|_{s=1}
:=
\frac{1}{m}.
\]
证毕。
\end{proof}

\begin{corollary}[single atom 对 residue-wise Mertens 常数只产生显式 digamma 位移]\label{cor:xi-time-part77-atomic-residue-finitepart-digamma-shift}
对固定模数 \(m\ge 2\) 与 residue 类 \(r\in \ZZ/m\ZZ\)，把 single atom 壳在 \(s=1\) 的 Laurent 展开记为
\[
\mathcal A^{\mathrm{cycle}}_{m,r}(s)
:=
\frac{1/m}{s-1}
+\mathcal M^{\mathrm{cycle}}_{m,r}
+O(s-1).
\]
则
\[
\mathcal M^{\mathrm{cycle}}_{m,0}
:=
\frac{\gamma-\log(2m)}{m},
\]
\[
\mathcal M^{\mathrm{cycle}}_{m,r}
:=
\frac{-\psi(r/m)-\log(2m)}{m}
\qquad (1\le r\le m-1),
\]
其中 \(\gamma\) 为 Euler 常数，\(\psi=\Gamma'/\Gamma\) 为 digamma 函数。

进一步，定义 full/core residue Dirichlet 级数
\[
\mathcal A^{\mathrm{full}}_{m,r}(s)
:=
\sum_{n\ge 1}A^{\mathrm{full}}_{m,r}(n)\,n^{-s},
\qquad
\mathcal A^{\mathrm{core}}_{m,r}(s)
:=
\sum_{n\ge 1}A^{\mathrm{core}}_{m,r}(n)\,n^{-s}.
\]
若把相应 Laurent 常数项记为
\[
\mathcal A^{\star}_{m,r}(s)
:=
\frac{\kappa^{\star}_{m,r}}{s-1}
+\log\mathfrak M^{\star}_{m,r}
+O(s-1),
\qquad
\star\in\{\mathrm{full},\mathrm{core}\},
\]
并在 \(r\neq 0\) 时仍以 \(r\in\{1,\dots,m-1\}\) 表示其标准代表元，则有
\[
\log\mathfrak M^{\mathrm{full}}_{m,0}
-\log\mathfrak M^{\mathrm{core}}_{m,0}
:=
\mathcal M^{\mathrm{cycle}}_{m,0}
=
\frac{\gamma-\log(2m)}{m},
\]
\[
\log\mathfrak M^{\mathrm{full}}_{m,r}
-\log\mathfrak M^{\mathrm{core}}_{m,r}
:=
\mathcal M^{\mathrm{cycle}}_{m,r}
=
\frac{-\psi(r/m)-\log(2m)}{m}
\qquad (1\le r\le m-1).
\]
因此，single atom 在 residue-wise finite part 层并不与 core 混合，而只留下由 cyclotomic/Hurwitz 数据完全决定的一族 digamma 位移。
\end{corollary}
\begin{proof}
由定理 \ref{thm:xi-time-part77-atomic-residue-dirichlet-hurwitz-shell} 的闭式表达式与标准展开
\[
\zeta(s)=\frac{1}{s-1}+\gamma+O(s-1),
\qquad
\zeta(s,a)=\frac{1}{s-1}-\psi(a)+O(s-1),
\]
以及
\[
2^{1-s}m^{-s}
=
\frac{1}{m}\Bigl(1-(s-1)\log(2m)+O((s-1)^2)\Bigr),
\]
逐项相乘，即得 \(\mathcal M^{\mathrm{cycle}}_{m,0}\) 与 \(\mathcal M^{\mathrm{cycle}}_{m,r}\) 的两条显式公式。

另一方面，由定义，
\[
A^{\mathrm{full}}_{m,r}(n)=A^{\mathrm{core}}_{m,r}(n)+A^{\mathrm{cycle}}_{m,r}(n),
\]
故在 \(\Re s>1\) 上
\[
\mathcal A^{\mathrm{full}}_{m,r}(s)
:=
\mathcal A^{\mathrm{core}}_{m,r}(s)+\mathcal A^{\mathrm{cycle}}_{m,r}(s).
\]
于是 \(\log\mathfrak M^{\mathrm{full}}_{m,r}-\log\mathfrak M^{\mathrm{core}}_{m,r}\) 恰等于 \(\mathcal A^{\mathrm{cycle}}_{m,r}(s)\) 在 \(s=1\) 处的 Laurent 常数项，即 \(\mathcal M^{\mathrm{cycle}}_{m,r}\)。证毕。
\end{proof}

\begin{theorem}[single atom 的 centered residue 扰动在零均值超平面上构成等角紧框架]\label{thm:xi-time-part77-atomic-centered-residue-tight-frame}
沿用定理 \ref{thm:xi-time-part65c-atomic-centered-residue-simplex} 的记号。对固定模数 \(m\ge 2\)，令
\[
\mathcal H_m
:=
\left\{x\in\RR^{\ZZ/m\ZZ}:\ \sum_{r\in\ZZ/m\ZZ}x_r=0\right\}.
\]
则族 \(\{\Delta^{(\ell)}\}_{\ell\in\ZZ/m\ZZ}\) 满足：
\begin{enumerate}
  \item 对任意 \(\ell,\ell'\in\ZZ/m\ZZ\)，有 \(\Delta^{(\ell)}\in\mathcal H_m\) 且
  \[
  \bigl\langle \Delta^{(\ell)},\Delta^{(\ell')}\bigr\rangle
  =
  4\left(\delta_{\ell,\ell'}-\frac{1}{m}\right).
  \]
  \item 因而 \(\{\Delta^{(\ell)}\}_{\ell\in\ZZ/m\ZZ}\) 在 \(\mathcal H_m\) 中构成一个等角紧框架；更精确地，
  \[
  \frac{1}{m}\sum_{\ell\in\ZZ/m\ZZ}
  \Delta^{(\ell)}\bigl(\Delta^{(\ell)}\bigr)^{\mathsf T}
  =
  \frac{4}{m}\left(I-\frac{1}{m}J\right),
  \]
  其中 \(J\) 为全 \(1\) 矩阵。
  \item 对任意子空间 \(V\subset\mathcal H_m\) 及其正交投影 \(Q_V\)，若 \(\dim V=d\)，则
  \[
  \frac{1}{m}\sum_{\ell\in\ZZ/m\ZZ}
  \bigl\|Q_V\Delta^{(\ell)}\bigr\|_2^2
  =
  \frac{4d}{m}.
  \]
\end{enumerate}
换言之，single atom 的 centered residue shell 在零均值超平面上对任意 \(d\)-维线性探针的平均能量注入都只由维数 \(d\) 决定，而与探针方向无关。
\end{theorem}
\begin{proof}
第 \(1\) 条正是定理 \ref{thm:xi-time-part65c-atomic-centered-residue-simplex} 的内容。又由推论 \ref{cor:xi-time-part65c-atomic-centered-residue-isotropic}，
\[
\frac{1}{m}\sum_{\ell\in\ZZ/m\ZZ}
\Delta^{(\ell)}\bigl(\Delta^{(\ell)}\bigr)^{\mathsf T}
=
\frac{4}{m}\left(I-\frac{1}{m}J\right).
\]
其中 \(I-\frac{1}{m}J\) 正是到零均值超平面 \(\mathcal H_m\) 的正交投影，故第 \(2\) 条给出 \(\{\Delta^{(\ell)}\}\) 在 \(\mathcal H_m\) 上的紧框架性质；而第 \(1\) 条中的两两内积公式又表明该框架是等角的。

最后，对任意 \(V\subset\mathcal H_m\) 的正交投影 \(Q_V\) 取迹可得
\[
\frac{1}{m}\sum_{\ell\in\ZZ/m\ZZ}
\bigl\|Q_V\Delta^{(\ell)}\bigr\|_2^2
=
\operatorname{tr}\!\left(
Q_V\cdot
\frac{1}{m}\sum_{\ell\in\ZZ/m\ZZ}
\Delta^{(\ell)}\bigl(\Delta^{(\ell)}\bigr)^{\mathsf T}
\right)
=
\frac{4}{m}\operatorname{tr}\!\left(Q_V\left(I-\frac{1}{m}J\right)\right).
\]
由于 \(V\subset\mathcal H_m=\mathbf 1^\perp\)，故 \(Q_V\mathbf 1=0\)，从而 \(Q_VJ=0\)。因此
\[
\operatorname{tr}\!\left(Q_V\left(I-\frac{1}{m}J\right)\right)
=
\operatorname{tr}(Q_V)
=
d,
\]
便得到
\[
\frac{1}{m}\sum_{\ell\in\ZZ/m\ZZ}
\bigl\|Q_V\Delta^{(\ell)}\bigr\|_2^2
=
\frac{4d}{m}.
\]
证毕。
\end{proof}

\noindent
这四条结果把单原子因子的 residue / twisted 接口进一步压缩到纯 Dirichlet 层的完全显式外壳：在 twisted 通道上，它恰等于一个根单位 polylog 壳；在 residue 通道上，它恰等于一族 Hurwitz--cyclotomic 壳，并把全部主极点残数统一锁定为 \(1/m\)；相应地，full/core residue-wise finite part 常数之间的差也不再携带任何隐藏动力学信息，而只是一组显式的 digamma 位移；最后，centered residue 扰动不仅给出正则单纯形几何，而且在零均值超平面上形成等角紧框架，使任意低维线性探针所接收到的平均能量严格服从维数比例律。

\paragraph{相位采样 \texorpdfstring{$\zeta$}{zeta} 的 Hurwitz 闭式、群刚性重建与异常分裂 torsor 势函数}
沿用定义 \ref{def:cdim-zeta-function}、命题 \ref{prop:cdim-residual-quotient-probe}、定理 \ref{thm:cdim-anomaly-quadratic-law}、命题 \ref{prop:cdim-discrete-anomaly-ledger} 与命题 \ref{prop:cdim-anomaly-splitting-torsor-obstruction} 的记号。对有限生成离散交换群
\[
G\cong \ZZ^r\oplus T
\]
记
\[
a_n(G):=\card{\operatorname{Hom}(G,\ZZ/n\ZZ)},
\qquad
\zeta_G^{\mathrm{ph}}(s):=\sum_{n\ge 1}\frac{a_n(G)}{n^s},
\]
并把 \(\zeta_G^{\mathrm{ph}}\) 视为前文 \(\zeta_G\) 的相位采样记号。由有限采样计数、异常群分解与分裂截面的 torsor 结构，可进一步压缩出下列五条结论。

\begin{theorem}[相位采样 \texorpdfstring{$\zeta$}{zeta} 的 Hurwitz 分解与唯一极点]\label{thm:xi-time-part78-phase-zeta-hurwitz-pole}
\leanverified{paper\_xi\_time\_part78\_phase\_zeta\_hurwitz\_pole}
设 \(G\cong \ZZ^r\oplus T\) 为有限生成离散交换群，其中 \(T\) 有限。记
\[
e:=\exp(T),
\qquad
c_T(n):=\card{T/nT}.
\]
则成立：
\begin{enumerate}
  \item \(c_T(n)\) 仅依赖于 \(\gcd(n,e)\)，因而是以 \(e\) 为周期的算术函数；
  \item 在 \(\Re(s)>r+1\) 上有 Hurwitz 分解
  \[
  \zeta_G^{\mathrm{ph}}(s)
  =e^{\,r-s}\sum_{a=1}^{e}c_T(a)\,\zeta\!\left(s-r,\frac{a}{e}\right);
  \]
  \item \(\zeta_G^{\mathrm{ph}}(s)\) 亚纯延拓到整个复平面，并在 \(s=r+1\) 处有唯一单极点，其留数为
  \[
  \operatorname*{Res}_{s=r+1}\zeta_G^{\mathrm{ph}}(s)
  =\frac{1}{e}\sum_{a=1}^{e}c_T(a).
  \]
\end{enumerate}
\end{theorem}
\begin{proof}
由命题 \ref{prop:cdim-residual-quotient-probe}，
\[
a_n(G)=n^r\card{T/nT}=n^r c_T(n).
\]
将 \(T\) 分解为 \(p\)-初等分量 \(T=\bigoplus_p T_{(p)}\)。对任意素数 \(p\) 与任意整数 \(n\)，乘以 \(n\) 在 \(T_{(p)}\) 上可写成
\[
n=p^{v_p(n)}u,\qquad p\nmid u,
\]
而乘以 \(u\) 为自同构，故
\[
\card{T_{(p)}/nT_{(p)}}=\card{T_{(p)}/p^{v_p(n)}T_{(p)}}.
\]
若 \(p^{\lambda_p}\Vert e\)，则 \(T_{(p)}\) 的指数为 \(p^{\lambda_p}\)，从而上式只依赖于 \(\min(v_p(n),\lambda_p)\)，亦即只依赖于 \(\gcd(n,p^{\lambda_p})\)。对所有 \(p\mid e\) 取乘积便得 \(c_T(n)\) 只依赖于 \(\gcd(n,e)\)，故以 \(e\) 为周期。

于是当 \(\Re(s)>r+1\) 时，
\[
\zeta_G^{\mathrm{ph}}(s)
=\sum_{n\ge 1}c_T(n)\,n^{r-s}
=\sum_{a=1}^{e}c_T(a)\sum_{m\ge 0}(me+a)^{r-s}.
\]
将 \(me+a=e\bigl(m+\frac{a}{e}\bigr)\) 提出，得到
\[
\zeta_G^{\mathrm{ph}}(s)
=e^{\,r-s}\sum_{a=1}^{e}c_T(a)\sum_{m\ge 0}\left(m+\frac{a}{e}\right)^{-(s-r)}
=e^{\,r-s}\sum_{a=1}^{e}c_T(a)\,\zeta\!\left(s-r,\frac{a}{e}\right).
\]
这便给出所述 Hurwitz 分解。

每个 Hurwitz zeta \(\zeta(s-r,a/e)\) 都只在 \(s-r=1\) 处有单极点，留数为 \(1\)。因此 \(\zeta_G^{\mathrm{ph}}\) 亚纯延拓到整个复平面，并且只能在 \(s=r+1\) 处出现极点；其留数为
\[
e^{\,r-(r+1)}\sum_{a=1}^{e}c_T(a)=\frac1e\sum_{a=1}^{e}c_T(a).
\]
由于 \(c_T(a)\ge 1\) 对所有 \(a\) 都成立，该留数严格为正，故该极点确为唯一单极点。证毕。
\end{proof}

\begin{theorem}[相位采样计数部分和的精确主项与 \texorpdfstring{$O(N^r)$}{O(N^r)} 余项]\label{thm:xi-time-part78-phase-zeta-partial-sum-sharp}
沿用定理 \ref{thm:xi-time-part78-phase-zeta-hurwitz-pole} 的记号，并记
\[
A_G(N):=\sum_{n\le N}a_n(G).
\]
则当 \(N\to\infty\) 时，
\[
A_G(N)=\frac{C_G}{r+1}N^{r+1}+O(N^r),
\qquad
C_G:=\frac1e\sum_{a=1}^{e}c_T(a).
\]
特别地，
\[
C_G=\operatorname*{Res}_{s=r+1}\zeta_G^{\mathrm{ph}}(s).
\]
\end{theorem}
\begin{proof}
由定理 \ref{thm:xi-time-part78-phase-zeta-hurwitz-pole}，函数 \(c_T(n)\) 以 \(e\) 为周期，故
\[
A_G(N)
=\sum_{a=1}^{e}c_T(a)\sum_{\substack{n\le N\\ n\equiv a\!\!\!\pmod e}}n^r.
\]
对每个 \(a\in\{1,\dots,e\}\) 记
\[
M_a:=\left\lfloor\frac{N-a}{e}\right\rfloor.
\]
则内层和可写为
\[
\sum_{\substack{n\le N\\ n\equiv a\!\!\!\pmod e}}n^r
=\sum_{m=0}^{M_a}(em+a)^r.
\]
由于 \((em+a)^r\) 是 \(m\) 的 \(r\) 次多项式，Faulhaber 求和公式给出
\[
\sum_{m=0}^{M_a}(em+a)^r
=\frac{e^r}{r+1}M_a^{r+1}+O(M_a^r),
\]
其中隐含常数对 \(a=1,\dots,e\) 一致，因为 \(e\) 固定且 \(a\) 只在有限集合中变动。又
\[
M_a=\frac{N}{e}+O(1)
\]
同样对 \(a\) 一致，于是
\[
\sum_{m=0}^{M_a}(em+a)^r
=\frac{1}{e(r+1)}N^{r+1}+O(N^r).
\]
代回并对 \(a\) 求和，得到
\[
A_G(N)
=\sum_{a=1}^{e}c_T(a)\left(\frac{1}{e(r+1)}N^{r+1}+O(N^r)\right)
=\frac{C_G}{r+1}N^{r+1}+O(N^r).
\]
最后，\(C_G=\operatorname*{Res}_{s=r+1}\zeta_G^{\mathrm{ph}}(s)\) 正是定理 \ref{thm:xi-time-part78-phase-zeta-hurwitz-pole} 的留数公式。证毕。
\end{proof}

\begin{theorem}[相位采样 \texorpdfstring{$\zeta$}{zeta} 的完全分类刚性]\label{thm:xi-time-part78-phase-zeta-complete-classification}
设 \(G,H\) 为有限生成离散交换群。若
\[
\zeta_G^{\mathrm{ph}}(s)\equiv \zeta_H^{\mathrm{ph}}(s)
\]
作为亚纯函数恒等成立，则
\[
G\cong H.
\]
等价地，映射
\[
G\longmapsto \zeta_G^{\mathrm{ph}}(s)
\]
是有限生成离散交换群同构类的完备不变量。
\end{theorem}
\begin{proof}
在某个足够右的半平面内，两端都由绝对收敛的 Dirichlet 级数
\[
\zeta_G^{\mathrm{ph}}(s)=\sum_{n\ge 1}\frac{a_n(G)}{n^s},
\qquad
\zeta_H^{\mathrm{ph}}(s)=\sum_{n\ge 1}\frac{a_n(H)}{n^s}
\]
给出。亚纯恒等因而在该半平面内也成立。由 Dirichlet 级数唯一性定理，
\[
a_n(G)=a_n(H)\qquad(\forall n\ge 1).
\]
再由推论 \ref{cor:xi-cdim-zeta-fingerprint-rigidity}（等价地，定理 \ref{thm:cdim-phase-spectrum-reconstruction}），可知这迫使 \(G\cong H\)。证毕。
\end{proof}

\begin{theorem}[异常分裂 torsor 的精确势函数]\label{thm:xi-time-part78-anomaly-splitting-torsor-cardinality}
\leanverified{paper\_xi\_time\_part78\_anomaly\_splitting\_torsor\_cardinality}
设
\[
G\cong \ZZ^r\oplus F
\]
为有限生成离散交换群，其中 \(F\) 有限。则分裂截面集合 \(\mathrm{Spl}(G)\) 为有限集，并满足
\[
\card{\mathrm{Spl}(G)}
=\bigl(\card{F}^{\,r}\,\card{\wedge^2 F}\bigr)^{\binom r2}.
\]
更具体地，若对每个素数 \(p\) 写
\[
F_{(p)}\cong \bigoplus_{j=1}^{t_p}\ZZ/p^{\lambda_{p,j}}\ZZ,
\qquad
\lambda_{p,1}\ge \cdots\ge \lambda_{p,t_p}\ge 1,
\]
则
\[
\card{\wedge^2 F}
=\prod_p p^{\sum_{j=2}^{t_p}(j-1)\lambda_{p,j}},
\]
从而
\[
\card{\mathrm{Spl}(G)}
=\prod_p
p^{\binom r2\left(
r\sum_{j=1}^{t_p}\lambda_{p,j}
+\sum_{j=2}^{t_p}(j-1)\lambda_{p,j}
\right)}.
\]
特别地，当 \(r\le 1\) 或 \(F=0\) 时，\(\mathrm{Spl}(G)\) 仅含一个元素。
\end{theorem}
\begin{proof}
由命题 \ref{prop:cdim-anomaly-splitting-torsor-obstruction}，
\(\mathrm{Spl}(G)\) 是
\[
\operatorname{Hom}\bigl(\mathcal D(G),\mathcal A(G)^0\bigr)
\]
的主齐次空间，因而
\[
\card{\mathrm{Spl}(G)}
=\card{\operatorname{Hom}\bigl(\mathcal D(G),\mathcal A(G)^0\bigr)}.
\]
再由定理 \ref{thm:cdim-anomaly-quadratic-law}，
\[
\mathcal A(G)^0\cong \TT^{\binom r2},
\]
故
\[
\operatorname{Hom}\bigl(\mathcal D(G),\mathcal A(G)^0\bigr)
\cong
\operatorname{Hom}\bigl(\mathcal D(G),\TT\bigr)^{\binom r2}.
\]
由于 \(\mathcal D(G)\) 为有限交换群，Pontryagin 对偶给出
\[
\card{\operatorname{Hom}\bigl(\mathcal D(G),\TT\bigr)}=\card{\mathcal D(G)},
\]
从而
\[
\card{\mathrm{Spl}(G)}=\card{\mathcal D(G)}^{\binom r2}.
\]

另一方面，命题 \ref{prop:cdim-discrete-anomaly-ledger} 给出
\[
\mathcal D(G)\cong
\operatorname{Hom}(\ZZ^r\otimes F,\TT)\times \operatorname{Hom}(\wedge^2F,\TT).
\]
因此
\[
\card{\mathcal D(G)}
=\card{\ZZ^r\otimes F}\,\card{\wedge^2F}
=\card{F}^{\,r}\,\card{\wedge^2F},
\]
因为 \(\ZZ^r\otimes F\cong F^r\)。于是得到
\[
\card{\mathrm{Spl}(G)}
=\bigl(\card{F}^{\,r}\,\card{\wedge^2F}\bigr)^{\binom r2}.
\]

为求 \(\card{\wedge^2F}\)，先按主分解写 \(F=\bigoplus_p F_{(p)}\)。由不同素数方向的张量积消失，
\[
\wedge^2F\cong \bigoplus_p \wedge^2F_{(p)}.
\]
再对
\[
F_{(p)}\cong \bigoplus_{j=1}^{t_p}\ZZ/p^{\lambda_{p,j}}\ZZ
\]
应用有限交换 \(p\)-群的外平方公式，得到
\[
\wedge^2F_{(p)}
\cong
\bigoplus_{1\le i<j\le t_p}\ZZ/p^{\min(\lambda_{p,i},\lambda_{p,j})}\ZZ.
\]
由于 \(\lambda_{p,1}\ge\cdots\ge\lambda_{p,t_p}\)，当 \(i<j\) 时有
\(\min(\lambda_{p,i},\lambda_{p,j})=\lambda_{p,j}\)，故
\[
\card{\wedge^2F_{(p)}}
=p^{\sum_{j=2}^{t_p}(j-1)\lambda_{p,j}}.
\]
对所有 \(p\) 取乘积便得
\[
\card{\wedge^2F}
=\prod_p p^{\sum_{j=2}^{t_p}(j-1)\lambda_{p,j}}.
\]
再与
\[
\card{F}=\prod_p p^{\sum_{j=1}^{t_p}\lambda_{p,j}}
\]
合并，即得最后的素数分解式。

若 \(r\le 1\)，则 \(\binom r2=0\)，故 \(\card{\mathrm{Spl}(G)}=1\)。若 \(F=0\)，则 \(\mathcal D(G)=0\)，同样得到 \(\mathrm{Spl}(G)\) 为单点。证毕。
\end{proof}

\begin{corollary}[异常分裂 torsor 的素数支撑与严格下界]\label{cor:xi-time-part78-anomaly-splitting-prime-support-lower-bound}
沿用定理 \ref{thm:xi-time-part78-anomaly-splitting-torsor-cardinality} 的记号，并设 \(r\ge 2\)。则对任意素数 \(p\) 都有
\[
p\mid \card{\mathrm{Spl}(G)}
\quad\Longleftrightarrow\quad
F_{(p)}\neq 0.
\]
等价地，
\[
\{p:\ p\mid \card{\mathrm{Spl}(G)}\}
=
\{p:\ F_{(p)}\neq 0\}.
\]
若进一步 \(F\neq 0\)，并记
\[
p_{\min}(F):=\min\{p:\ F_{(p)}\neq 0\},
\]
则
\[
\card{\mathrm{Spl}(G)}
\ge p_{\min}(F)^{\,r\binom r2}.
\]
\end{corollary}
\begin{proof}
由定理 \ref{thm:xi-time-part78-anomaly-splitting-torsor-cardinality} 的显式分解式，
\[
v_p\!\bigl(\card{\mathrm{Spl}(G)}\bigr)
=\binom r2\left(
r\sum_{j=1}^{t_p}\lambda_{p,j}
+\sum_{j=2}^{t_p}(j-1)\lambda_{p,j}
\right).
\]
若 \(F_{(p)}=0\)，则所有 \(\lambda_{p,j}\) 均为零，从而
\(
v_p(\card{\mathrm{Spl}(G)})=0
\)。
反之若 \(F_{(p)}\neq 0\)，则
\(
\sum_{j=1}^{t_p}\lambda_{p,j}\ge 1
\)；
又因 \(r\ge 2\)，故
\[
v_p\!\bigl(\card{\mathrm{Spl}(G)}\bigr)
\ge \binom r2\,r>0.
\]
于是
\[
p\mid \card{\mathrm{Spl}(G)}
\quad\Longleftrightarrow\quad
F_{(p)}\neq 0.
\]

若 \(F\neq 0\)，取其最小支撑素数 \(p_{\min}(F)\)。对该素数有
\[
v_{p_{\min}(F)}\!\bigl(\card{\mathrm{Spl}(G)}\bigr)\ge \binom r2\,r,
\]
其余素数因子贡献非负，故
\[
\card{\mathrm{Spl}(G)}
\ge p_{\min}(F)^{\,r\binom r2}.
\]
证毕。
\end{proof}

\noindent
这五条结果把有限生成离散交换群的 phase-sampling 侧统计与二阶异常分裂的离散扭量压缩到同一条可审计主线上：相位采样 \(\zeta\) 不仅具有有限 Hurwitz 闭式并把唯一主极点精确锁定在 \(s=\cdim(G)+1\)，而且其亚纯函数本身已经足以完整恢复群的同构类型；与此同时，异常分裂截面的全部非唯一性也不再只是存在性障碍，而被一个可显式计算的有限 torsor 势函数完全控制，其素数支撑与大小下界均由 torsion 主分解直接给出。

\paragraph{MOM 乘法分级、交叉异常阈值与 window-\texorpdfstring{$6$}{6} 扇区读出的信息压缩}
在相位采样 \(\zeta\) 的 Hurwitz 闭式、群刚性重建以及异常分裂 torsor 的势函数已经闭合之后，前文关于四生成元投影词重写、交叉异常坐标化、对称幂 moment-kernel 编译与 window-\(6\) 的取向扇区结构还可继续压缩为同一条更高层的审计链：MOM 前缀指数给出乘法分级，交叉异常族给出轴可分解性的最小完备审计坐标，而对称幂编译与四扇区读出则分别把“显影成本”与“信息损失”严格量化。

\begin{theorem}[四生成元投影词演算的乘法分级分裂]\label{thm:xi-time-part79-pom-multiplicative-graded-splitting}
在定义 \ref{def:pom-four-gen-rewrite-rules}、定义 \ref{def:pom-algorithmic-equivalence-nf-anom} 与命题 \ref{prop:pom-confluence-mod-anom-audited} 的有界审计切片内，令
\[
\mathcal W_4
\]
为由
\(
\mathrm{PROJ}_u,\ E_m,\ \mathrm{LIFT}_G,\ \mathrm{MOM}_q\ (q\ge 2)
\)
生成的自由单半群。记
\[
\Pi_{\mathrm{MOM}}:\mathcal W_4\to (\NN_{\ge 1},\times)
\]
为命题 \ref{prop:pom-mom-prefix-product-invariant} 的 MOM 前缀指数。称 \(w_1,w_2\in\mathcal W_4\) 结构等价，记为
\(
w_1\equiv_{\mathrm{str}}w_2
\)，若二者都可仅用 \((\mathrm{RMOM})\) 与 \((\mathrm{RMOM\text{-}BUBBLE})\) 规约为
\[
w_i\ \Rightarrow\ \mathrm{MOM}_{Q_i}\circ v_i,
\qquad
v_i\ \text{不含}\ \mathrm{MOM},
\]
并满足
\[
Q_1=Q_2
\qquad\text{且}\qquad
v_1\sim v_2,
\]
其中 \(\sim\) 为定义 \ref{def:pom-algorithmic-equivalence-nf-anom} 的三生成元算法等价。
则存在规范分裂
\[
\mathcal W_4/\equiv_{\mathrm{str}}
\ \cong\
\bigsqcup_{Q\ge 1}\mathfrak E_Q,
\qquad
\mathfrak E_Q:=
\{[w]_{\equiv_{\mathrm{str}}}:\ \Pi_{\mathrm{MOM}}(w)=Q\}.
\]
并且每个结构类都存在一个前缀正规代表
\[
\mathrm{MOM}_{Q}\circ N,
\]
其中 \(N\) 为命题 \ref{prop:pom-three-gen-interface-normal-form} 的三生成元接口正规形。更精确地，在固定 \(Q\) 的扇区 \(\mathfrak E_Q\) 内，结构类恰由不含 MOM 的残余片段的
\[
\bigl(\mathrm{NF},[\mathrm{Anom}]\bigr)
\]
唯一决定。
\leanverified{paper\_xi\_time\_part79\_pom\_multiplicative\_graded\_splitting}
\end{theorem}
\begin{proof}
命题 \ref{prop:pom-mom-prefix-product-invariant} 断言 \(\Pi_{\mathrm{MOM}}\) 在允许重写下严格不变，因此不同 \(Q\) 的词不可能在结构等价中相遇。
另一方面，命题 \ref{prop:pom-mom-interface-normal-form} 给出一条规范的归一化策略：先用 \((\mathrm{RMOM})\) 合并全部 moment 门，再用 \((\mathrm{RMOM\text{-}BUBBLE})\) 将其外提为唯一前缀 \(\mathrm{MOM}_Q\)，从而把任意 \(w\) 化为 \(\mathrm{MOM}_Q\circ v\) 的形式，其中 \(v\) 落回三生成元片段。
对该不含 MOM 的残余片段，命题 \ref{prop:pom-confluence-mod-anom-audited} 与定义 \ref{def:pom-algorithmic-equivalence-nf-anom} 表明，其路径不变量正是接口正规形与异常上同调类；因此在固定 \(Q\) 内，结构类恰由 \((\mathrm{NF}(v),[\mathrm{Anom}(v)])\) 分类。综上，整个商集按 MOM 前缀指数分裂为互不连通的乘法分级扇区。证毕。
\end{proof}

\begin{theorem}[高阶交叉异常审计的锐完备阈值]\label{thm:xi-time-part79-cross-anom-sharp-completeness-threshold}
\leanverified{paper\_xi\_time\_part79\_cross\_anom\_sharp\_completeness\_threshold}
设 \(k\ge 3\)，并令 \(G_1,\dots,G_k\) 为有限阿贝尔群。固定 \(\theta\) 后，把
\[
f(\chi):=\log\mathfrak{M}_{\chi}(\theta),
\qquad
\chi\in \widehat G_1\times\cdots\times\widehat G_k,
\]
视为常数层载荷，并记其 Möbius 交叉分量为 \(\Delta_S(\chi_S;\theta)\)。
则全体高阶交叉异常
\[
\{\Delta_S:\ |S|\ge 2\}
\]
构成轴可分解性的最小完备审计家族，具体地：
\begin{enumerate}
  \item \(f\) 轴可分解当且仅当对一切 \(|S|\ge 2\) 都有
  \[
  \Delta_S(\chi_S;\theta)\equiv 0.
  \]
  \item 对任意固定的 \(r\in\{3,\dots,k\}\)，只要前 \(r\) 条轴各含一个非平凡角色，则任何仅检查全部
  \[
  \{\Delta_S:\ 2\le |S|<r\}
  \]
  的审计协议都不是完备的。更强地，存在一个纯 \(r\)-体载荷 \(f_r\)，使得
  \[
  \Delta_S^{(r)}\equiv 0\qquad(2\le |S|<r),
  \qquad
  \Delta_{[r]}^{(r)}\not\equiv 0.
  \]
\end{enumerate}
因此，一旦省略任意阶数 \(r\ge 3\) 的检测，就会留下纯 \(r\)-体异常的结构性盲区。
\end{theorem}
\begin{proof}
第 (1) 条正是定理 \ref{thm:pom-axis-decomposability-iff} 的内容。

对第 (2) 条，取前 \(r\) 个轴上的非平凡角色 \(a_i\in\widehat G_i\setminus\{\chi_{i,0}\}\)，并定义
\[
f_r(\chi):=
\begin{cases}
1,& \chi=(a_1,\dots,a_r,\chi_{r+1,0},\dots,\chi_{k,0}),\\
0,& \text{其他情形}.
\end{cases}
\]
若 \(S\subseteq[k]\) 满足 \(2\le |S|<r\)，则在 \(\Delta_S\) 的 Möbius 差分中，每一项都至少有一个前 \(r\) 坐标被置为平凡角色，故所有取值均为 \(0\)，从而
\(
\Delta_S^{(r)}\equiv 0
\)。
但当 \(S=[r]\) 且 \(\chi_{[r]}=(a_1,\dots,a_r)\) 时，差分和中只有满集项保留，故
\[
\Delta_{[r]}^{(r)}(a_1,\dots,a_r;\theta)=1\neq 0.
\]
于是 \(f_r\) 不是轴可分解的，却逃过了所有低于 \(r\) 阶的交叉异常审计。证毕。
\end{proof}

\begin{theorem}[异常显影阈值的多项式编译成本]\label{thm:xi-time-part79-anom-threshold-polynomial-compilation-cost}
设 \(w\) 为任意投影词，并记
\[
a:=[\mathrm{Anom}(w)]
\]
为其异常上同调类。若 \(a\neq 0\)，并且底层在线核的内部状态集为有限集 \(\mathcal Q\)，则对任意阈值 \(\tau>0\) 都存在整数 \(q\ge 2\)，使得：
\begin{enumerate}
  \item 某个异常坐标 \(\chi\) 满足
  \[
  \bigl|[\mathrm{Anom}(w^{\boxtimes q})]_{\chi}\bigr|\ge \tau;
  \]
  \item \(q\)-重 moment-kernel 可在命题 \ref{prop:pom-moment-symmetric-power} 的对称幂商核上实现，其状态数至多为
  \[
  \binom{q+|\mathcal Q|-1}{|\mathcal Q|-1};
  \]
  \item 因而实现阈值 \(\tau\) 的最小状态成本满足上界
  \[
  \mathrm{Cost}_{\min}(\tau)=O\!\bigl(\tau^{|\mathcal Q|-1}\bigr).
  \]
\end{enumerate}
与此相对，若采用未压缩的 \(q\)-份同步直积实现，则状态数为 \(|\mathcal Q|^q\)，即指数规模。
\end{theorem}
\begin{proof}
由命题 \ref{prop:pom-anom-gap-amplification}，有
\[
[\mathrm{Anom}(w^{\boxtimes q})]=q\cdot a.
\]
取一坐标 \(\chi\) 使 \(a_\chi\neq 0\)，并令
\[
q:=\max\!\left\{2,\left\lceil \frac{\tau}{|a_\chi|}\right\rceil\right\}.
\]
则
\[
\bigl|[\mathrm{Anom}(w^{\boxtimes q})]_{\chi}\bigr|
=
q|a_\chi|
\ge \tau.
\]
另一方面，命题 \ref{prop:pom-moment-symmetric-power} 给出 \(q\)-重碰撞核的对称幂商态空间
\[
\mathcal N_q=\left\{(n_s)_{s\in\mathcal Q}\in\ZZ_{\ge 0}^{\mathcal Q}:\ \sum_{s\in\mathcal Q}n_s=q\right\},
\]
其基数正是
\[
|\mathcal N_q|=\binom{q+|\mathcal Q|-1}{|\mathcal Q|-1}.
\]
由于 \(q=O(\tau)\)，故最小状态成本至多按 \(\tau^{|\mathcal Q|-1}\) 多项式增长。若不做对称压缩，则直接同步直积的状态空间为 \(\mathcal Q^q\)，大小为 \(|\mathcal Q|^q\)，故为指数级。证毕。
\end{proof}

\begin{theorem}[window-\texorpdfstring{$6$}{6} 取向扇区读出的 \texorpdfstring{$19$}{19} 比特压缩下界]\label{thm:xi-time-part79-window6-sector-19bit-collapse}
沿用推论 \ref{cor:xi-window6-boundary-parity-unique-minimal-faithful-quotient} 的记号，记
\[
\mathcal A_6:=\bigl(G_6^{\mathrm{bin}}\bigr)^{\mathrm{ab}}\cong (\ZZ/2\ZZ)^{21}
\]
为 window-\(6\) 的完备二元异常签名空间。
另一方面，由命题 \ref{prop:fold-groupoid-z2x2-sector-by-orientation} 与定理 \ref{thm:op-algebra-fold-weak-hopf-z2x2-sector-closure}，window-\(6\) 的取向指纹只给出四个弱 Hopf 扇区
\[
\CC[\mathcal G_6]\cong \bigoplus_{(\alpha,\beta)\in\{\pm 1\}^2} H_6^{\alpha,\beta}.
\]
则任意只通过取向扇区标签
\[
\sigma:\mathcal A_6\to \{\pm 1\}^2
\]
进行的异常分类，都至少丢失 \(19\) 比特信息。更精确地，至少存在一个扇区，其纤维满足
\[
\#\sigma^{-1}(\alpha,\beta)\ge 2^{19}.
\]
\end{theorem}
\begin{proof}
由 \(\mathcal A_6\cong(\ZZ/2\ZZ)^{21}\)，可知
\[
|\mathcal A_6|=2^{21}.
\]
而取向扇区标签集 \(\{\pm1\}^2\) 的基数为 \(4\)。
因此任何 \(\sigma:\mathcal A_6\to\{\pm1\}^2\) 的像至多有 \(4\) 个点。由鸽巢原理，至少有一个标签的原像大小不小于
\[
\frac{2^{21}}{4}=2^{19}.
\]
换言之，取向扇区读出至多保留 \(2\) 个二元比特，而完整异常签名需要 \(21\) 个独立比特；两者之间的不可约信息缺口至少为 \(19\) 比特。证毕。
\end{proof}

\begin{conjecture}[取向扇区读出的线性失配律]\label{conj:xi-time-part79-orientation-sector-linear-mismatch}
设 \(m\) 满足所有 bin-fold 纤维都非平凡，并令 \(\mathcal A_m\) 为保持纤维独立置换语义的最小完备二元异常签名空间。则任意仅通过取向指纹的四扇区标签
\[
\sigma_m:\mathcal A_m\to \{\pm 1\}^2
\]
实现的异常读出，在所有与弱 Hopf 扇区分裂相容的实现中，都应满足不可改进的信息损失下界
\[
\mathrm{loss}(\sigma_m)\ge |X_m|-2.
\]
换言之，四扇区数保持常值，而完备异常维数应随分辨率按 \(|X_m|\) 线性增长；window-\(6\) 的 \(19\) 比特缺口只是这条普遍失配律在首个非平凡窗口上的精确实例。
\end{conjecture}

\noindent
其直接根据在于：命题 \ref{prop:fold-groupoid-z2x2-sector-by-orientation} 与定理 \ref{thm:op-algebra-fold-weak-hopf-z2x2-sector-closure} 已对一般 \(m\) 固定四扇区框架，而定理 \ref{thm:gut-fiber-parity-minimal-complete-bits} 在“全部纤维均非平凡”的情形下给出最小完备比特数恰为 \(|X_m|\)；因此随分辨率增长而线性放大的失配，至少在现有证书链上已呈现出稳定的全局形态。

\noindent
由此，MOM 前缀指数把四生成元投影词演算分裂为彼此隔离的乘法扇区；交叉异常族把轴可分解性的审计门槛精确提升为“全阶且不可降阶”的完备家族；对称幂商核则表明任意显影阈值仍可在多项式后端体积内实现；而当读出被限制为 window-\(6\) 的四个取向扇区时，完备异常所需的二十一比特信息被迫压缩为两比特粗标签。这条链把“可重写性”“可审计性”“可编译性”与“不可压缩性”锁进了同一组离散不变量。












\paragraph{二步 primitive 全系数、偶模坏模律与误差指数的反向成本}
在七条分支二回路的 primitive 塌缩、模 \(2\) 的最小坏模证人、正实退化点列表以及纯碰撞误差指数的相变已经建立之后，还可进一步抽取出下列五条完全闭式的后果。它们分别把最短 primitive 层的全系数、其在 ghost 层的偶次迹塔、偶数模上的坏模普遍性、\(\theta=0\) 邻域无退化窗的极大性，以及 \(\beta(u)\to 1\) 的前向渐近所对应的反向参数成本压缩为显式公式。

\begin{theorem}[真实输入输出位势的二步 primitive 系数恰为 \texorpdfstring{$8u$}{8u}]\label{thm:xi-time-part68a-output-two-step-primitive-eight-u}
\leanverified{paper\_xi\_time\_part68a\_output\_two\_step\_primitive\_eight\_u}
沿用定理 \ref{thm:xi-time-part62d-atomic-two-step-witt-rankone-collapse} 的记号，记
\[
P(z,u)=\sum_{n\ge 1}p_n(u)z^n.
\]
则二步 primitive 系数满足精确恒等式
\[
p_2(u)=8u.
\]
\end{theorem}
\begin{proof}
由定理 \ref{thm:xi-time-part62d-atomic-two-step-witt-rankone-collapse}，
\[
p_2(u)=u+p_2^{\mathrm{core}}(u).
\]
同一定理并已指出：branch 层中满足
\[
(|\tau|,E)=(2,1)
\]
的两步回路恰有 \(7\) 条。按 primitive 生成函数的定义，每一条长度 \(2\)、能量 \(1\) 的 primitive 回路都恰贡献一个 monomial \(u z^2\)；因此
\[
p_2^{\mathrm{core}}(u)=7u.
\]
代回即得
\[
p_2(u)=u+7u=8u.
\]
证毕。
\end{proof}

\begin{corollary}[长度 \texorpdfstring{$2$}{2} primitive 扇区在 ghost 层给出完全显式的偶次迹塔]\label{cor:xi-time-part68a-output-two-step-primitive-even-ghost-tower}
沿用定理 \ref{thm:xi-time-part68a-output-two-step-primitive-eight-u} 的记号。定义由长度 \(2\) primitive 扇区单独诱导的 ghost 级数为
\[
\sum_{n\ge 1}\frac{a_n^{\langle 2\rangle}(u)}{n}z^n
:=
\sum_{m\ge 1}\frac{p_2(u^m)}{m}z^{2m}.
\]
则对一切 \(m\ge 1\) 都有
\[
a_{2m}^{\langle 2\rangle}(u)=16u^m,
\qquad
a_{2m+1}^{\langle 2\rangle}(u)=0.
\]
\end{corollary}
\begin{proof}
由定理 \ref{thm:xi-time-part68a-output-two-step-primitive-eight-u}，
\[
p_2(u^m)=8u^m
\qquad (m\ge 1).
\]
故
\[
\sum_{n\ge 1}\frac{a_n^{\langle 2\rangle}(u)}{n}z^n
=
\sum_{m\ge 1}\frac{8u^m}{m}z^{2m}.
\]
将右端与
\[
\sum_{m\ge 1}\frac{a_{2m}^{\langle 2\rangle}(u)}{2m}z^{2m}
\]
逐项比较，立得
\[
\frac{a_{2m}^{\langle 2\rangle}(u)}{2m}=\frac{8u^m}{m},
\]
从而
\[
a_{2m}^{\langle 2\rangle}(u)=16u^m.
\]
而右端仅含偶次幂，故全部奇次系数必为 \(0\)。证毕。
\end{proof}

\begin{theorem}[每一个偶数模都是坏模数]\label{thm:xi-time-part68a-collision-all-even-moduli-bad}
沿用定理 \ref{thm:xi-real-input-40-collision-minimal-modulus-two-witness} 的记号。对每个模数 \(m\ge 2\) 与扭曲通道 \(1\le j\le m-1\)，记
\[
\omega_m:=e^{2\pi i/m},
\qquad
\rho_{m,j}:=\rho\!\bigl(W(\omega_m^j)\bigr),
\qquad
\lambda:=\rho(W(1))=\varphi^2.
\]
则对每一个偶数 \(m\) 都存在某个 \(j\) 使
\[
\rho_{m,j}>\sqrt{\lambda}=\varphi.
\]
更具体地，取 \(j=m/2\) 即有
\[
\rho_{m,m/2}=\rho(W(-1))>\varphi.
\]
因此全部偶数模都是平方根门槛失效的坏模数。
\end{theorem}
\begin{proof}
若 \(m\) 为偶数，则取
\[
j=\frac{m}{2}.
\]
此时
\[
\omega_m^{\,m/2}=e^{\pi i}=-1,
\]
故
\[
\rho_{m,m/2}=\rho\!\bigl(W(-1)\bigr).
\]
而定理 \ref{thm:xi-real-input-40-collision-minimal-modulus-two-witness} 已证明
\[
\rho\!\bigl(W(-1)\bigr)>\sqrt{\lambda}=\varphi.
\]
于是该偶数模 \(m\) 至少在通道 \(j=m/2\) 上突破平方根门槛。证毕。
\end{proof}
\leanverified{paper\_xi\_time\_part68a\_collision\_all\_even\_moduli\_bad}

\begin{theorem}[\texorpdfstring{$\theta=0$}{theta=0} 邻域正实参数的最大对称无退化窗]\label{thm:xi-time-part68a-output-positive-real-maximal-nondegenerate-window}
沿用注 \ref{rem:real-input-40-output-potential-degeneracy-phase} 的记号。记距离 \(u=1\) 最近的两个正实退化点为
\[
u_-=0.9328372624859649\ldots,
\qquad
u_+=1.0547964213679192\ldots,
\]
并定义
\[
\theta_0:=\min\{-\log u_-,\,\log u_+\}
=0.0533477827\ldots.
\]
则：
\begin{enumerate}
  \item 当 \(|\theta|<\theta_0\) 时，沿正实参数 \(u=e^\theta\) 不会跨越任何正实退化点；
  \item 若 \(\theta_1>\theta_0\)，则对称区间 \((-\theta_1,\theta_1)\) 必跨越至少一个正实退化点。
\end{enumerate}
因而 \((-\theta_0,\theta_0)\) 正是以 \(\theta=0\) 为中心的最大对称无退化实窗。
\end{theorem}
\begin{proof}
注 \ref{rem:real-input-40-output-potential-degeneracy-phase} 已给出全部正实退化点
\[
u\in\{0.008460\ldots,\ 0.745534\ldots,\ 0.932837\ldots,\ 1.054796\ldots,\ 2.090392\ldots,\ 15.522820\ldots\},
\]
并指出离 \(u=1\) 最近的两点正是 \(u_-\) 与 \(u_+\)。

若 \(|\theta|<\theta_0\)，则
\[
-\log u_-<\theta<\log u_+,
\]
从而
\[
u_-<e^\theta<u_+.
\]
因此 \(u=e^\theta\) 落在开区间 \((u_-,u_+)\) 内，不会跨越任何正实退化点。这证明了第 \(1\) 条。

反之，若 \(\theta_1>\theta_0\)，则由 \(\theta_0\) 的定义可知，必有
\[
\theta_1>-\log u_-
\quad\text{或}\quad
\theta_1>\log u_+.
\]
在前一种情形，区间 \((-\theta_1,\theta_1)\) 包含点 \(\log u_-<0\)；在后一种情形，它包含点 \(\log u_+>0\)。两者都意味着该对称区间至少跨越一个正实退化点。故第 \(2\) 条成立，极大性随之得到。证毕。
\end{proof}

\begin{theorem}[纯碰撞误差指数的 Lambert--\texorpdfstring{$\mathscr W$}{W} 反演律]\label{thm:xi-time-part68a-collision-beta-lambertw-inversion}
沿用推论 \ref{cor:real-input-40-collision-error-phase} 的记号，并令
\[
\varepsilon:=1-\beta(u)\downarrow 0
\qquad (u\to+\infty).
\]
记 \(\mathscr W\) 为 Lambert--\(\mathscr W\) 函数主支，即
\[
\mathscr W(x)e^{\mathscr W(x)}=x
\qquad (x>0).
\]
则有反演渐近
\[
\sqrt{u}
=
\frac{1}{\varepsilon\,\mathscr W(1/\varepsilon)}(1+o(1)),
\qquad
u
=
\frac{1}{\varepsilon^2\,\mathscr W(1/\varepsilon)^2}(1+o(1)).
\]
特别地，由
\[
\mathscr W(x)=\log x-\log\log x+o(1)
\qquad (x\to+\infty)
\]
可得粗尺度资源律
\[
u\asymp \frac{1}{\varepsilon^2\log^2(1/\varepsilon)}.
\]
\end{theorem}
\begin{proof}
由推论 \ref{cor:real-input-40-collision-error-phase}，
\[
\beta(u)=1-\frac{2}{\sqrt u\,\log u}+O\!\left(\frac{1}{u\log u}\right)
\qquad (u\to+\infty).
\]
令
\[
x:=\sqrt u.
\]
则 \(\log u=2\log x\)，故
\[
\varepsilon
=
\frac{1}{x\log x}+O\!\left(\frac{1}{x^2\log x}\right)
=
\frac{1}{x\log x}(1+o(1)).
\]
于是
\[
x\log x=\frac{1}{\varepsilon}(1+o(1)).
\]
再令
\[
y:=\log x,
\]
则 \(x=e^y\)，从而
\[
y e^y=\frac{1}{\varepsilon}(1+o(1)).
\]
按 Lambert--\(\mathscr W\) 函数定义，
\[
y=\mathscr W\!\left(\frac{1}{\varepsilon}(1+o(1))\right).
\]
又由于
\[
\mathscr W'(t)=\frac{\mathscr W(t)}{t(1+\mathscr W(t))}=O(t^{-1})
\qquad (t\to+\infty),
\]
故
\[
\mathscr W\!\left(\frac{1}{\varepsilon}(1+o(1))\right)
=
\mathscr W\!\left(\frac{1}{\varepsilon}\right)+o(1).
\]
于是
\[
x=e^y
=
\frac{(1/\varepsilon)(1+o(1))}{\mathscr W(1/\varepsilon)+o(1)}
=
\frac{1}{\varepsilon\,\mathscr W(1/\varepsilon)}(1+o(1)).
\]
这就得到
\[
\sqrt u=\frac{1}{\varepsilon\,\mathscr W(1/\varepsilon)}(1+o(1)).
\]
两边平方即得
\[
u=\frac{1}{\varepsilon^2\,\mathscr W(1/\varepsilon)^2}(1+o(1)).
\]
最后由
\[
\mathscr W(1/\varepsilon)\sim \log(1/\varepsilon)
\qquad (\varepsilon\downarrow 0)
\]
可知
\[
u\asymp \frac{1}{\varepsilon^2\log^2(1/\varepsilon)}.
\]
证毕。
\end{proof}

\noindent
由此，真实输入 \(40\) 状态核的最短 primitive 层、偶模单位根扭曲、正实退化窗与纯碰撞误差指数之间形成了一条更细的封闭链：二步 primitive 坐标已在 \(8u\) 处完全闭式化，并在 ghost 层外显为一整条偶次迹塔；奇偶扭曲迫使全部偶数模统一越过平方根门槛，而正实参数的中心对称无退化窗则恰由 \(\theta_0\) 取尽；最后，\(\beta(u)\to 1\) 的慢模塌缩不只给出前向渐近，而且可反解为带 Lambert--\(\mathscr W\) 校正的显式参数成本律。

\paragraph{围绕 \texorpdfstring{$u=1$}{u=1} 的正实谱室与代数 \texorpdfstring{$\pm1$}{±1} 锚点}
在正实退化点列表、中心对称无退化窗以及整数级共振校准点都已确定之后，还可把真实输入 \(40\) 状态核的正实谱几何进一步压缩为下列两条定理。前一条把围绕 \(u=1\) 的最大正实谱室与其对数半径完全闭式化；后一条则把 \(\lambda=\pm1\) 的两枚正实代数锚点提升为非退化谱证书。

\begin{theorem}[围绕 \texorpdfstring{$u=1$}{u=1} 的正实最大无额外碰撞谱室]\label{thm:xi-time-part68aa-output-positive-real-maximal-collisionfree-chamber}
\leanverified{paper\_xi\_time\_part68aa\_output\_positive\_real\_maximal\_collisionfree\_chamber}
沿用命题 \ref{prop:real-input-40-output-spectral-collisions} 与注 \ref{rem:real-input-40-output-potential-degeneracy-phase} 的记号，记
\[
\Delta(z,u)=(1-uz^2)Q(z,u).
\]
则在正实参数轴 \(u>0\) 上，除结构性点 \(u=1\) 外，\(Q(\cdot,u)\) 出现重根当且仅当
\[
u\in
\left\{
0.0084602654608787966132,\ 
0.74553449580298956225,\ 
0.93283726248596492776,\ 
1.0547964213679192275,\ 
2.0903923780490870311,\ 
15.522820704426852031
\right\}.
\]
因此，以 \(u=1\) 为中心且不再跨越其他正实碰撞点的最大开区间恰为
\[
\bigl(0.93283726248596492776,\ 1.0547964213679192275\bigr).
\]
若改用
\[
\theta:=\log u
\]
作为局部参数，则以 \(\theta=0\) 为中心的最大对称无额外碰撞区间为
\[
|\theta|<\theta_0,
\qquad
\theta_0
:=
\min\!\Bigl\{
\log 1.0547964213679192275,\ 
-\log 0.93283726248596492776
\Bigr\}
=
0.0533477827\ldots.
\]
\end{theorem}
\begin{proof}
注 \ref{rem:real-input-40-output-potential-degeneracy-phase} 已列出 \(u>0\) 上全部正实退化点，并指出离 \(u=1\) 最近的两点恰为
\[
u_-=0.93283726248596492776,\qquad
u_+=1.0547964213679192275.
\]
故在去掉这些碰撞点之后，包含 \(u=1\) 的最大连通开区间只能是
\[
(u_-,u_+)
=
\bigl(0.93283726248596492776,\ 1.0547964213679192275\bigr).
\]
这就给出了第一条结论。

再取 \(\theta=\log u\)。由于对数映射在 \((0,+\infty)\) 上严格单调，上一开区间对应于
\[
\bigl(\log u_-,\ \log u_+\bigr).
\]
若要求得到以 \(\theta=0\) 为中心的对称区间，则其最大半径只能取左右两端到 \(0\) 的较小距离，即
\[
\theta_0=\min\{-\log u_-,\,\log u_+\}.
\]
代入上述数值便得
\[
\theta_0=0.0533477827\ldots.
\]
证毕。
\end{proof}

\begin{theorem}[非退化 \texorpdfstring{$\pm1$}{±1} 谱锚的代数精确性]\label{thm:xi-time-part68aa-output-algebraic-plusminus-one-nondegenerate-anchors}
沿用
\[
\Delta(z,u)=\det(I-zM(u))=(1-uz^2)Q(z,u)
\]
的记号。则存在两枚完全代数的正实参数
\[
u_+=4,
\qquad
u_-=-3+\sqrt{11},
\]
使得矩阵 \(M(u)\) 分别具有特征值
\[
\lambda=1
\qquad\text{与}\qquad
\lambda=-1.
\]
更强地，这两个谱锚都落在正实碰撞集合之外，因而对应特征值均为非退化谱点。换言之，\(\lambda=\pm1\) 在这两处都给出完全代数且可直接审计的正实谱校准点。
\end{theorem}
\begin{proof}
由推论 \ref{cor:real-input-40-integer-resonance-calibration}，
\[
Q(1,4)=0,
\qquad
Q(-1,-3+\sqrt{11})=0.
\]
当 \(u\neq 1\) 时，
\[
\Delta(\pm1,u)=(1-u)Q(\pm1,u),
\]
故
\[
\Delta(1,4)=0,
\qquad
\Delta(-1,-3+\sqrt{11})=0.
\]
按
\[
\Delta(z,u)=\det(I-zM(u)),
\]
可知 \(z=\pm1\) 为零点等价于 \(M(u)\) 具有特征值 \(\lambda=1/z=\pm1\)。因此在 \(u_+=4\) 与 \(u_-=-3+\sqrt{11}\) 处分别得到 \(\lambda=1\) 与 \(\lambda=-1\)。

另一方面，注 \ref{rem:real-input-40-output-potential-degeneracy-phase} 已给出全部正实碰撞参数。显然
\[
4\notin
\left\{
0.008460\ldots,\ 0.745534\ldots,\ 0.932837\ldots,\ 1.054796\ldots,\ 2.090392\ldots,\ 15.522820\ldots
\right\},
\]
且
\[
-3+\sqrt{11}\approx 0.31662479
\]
同样不在该集合中，并且二者都不等于结构性点 \(u=1\)。因此对应零点都是简单根，亦即 \(\lambda=\pm1\) 在这两处都是非退化谱点。证毕。
\end{proof}

\noindent
由此，围绕平衡点 \(u=1\) 的正实谱室不再只由局部展开半径间接刻画，而是被一张完整的碰撞点表和一个最优对数半径精确锁定；与此同时，\(\lambda=\pm1\) 的两枚正实代数锚点又把整数级共振提升为真正的非退化谱校准点。于是，局部可展域与全局代数锚定在同一条正实谱证书链上闭合。




\paragraph{dyadic 屏幕 exactization、边界码退化、矩层析二歧与读出精度分离}
在 dyadic 体--界全息链的精确屏幕、边界像 LDPC 阴影、有限矩盒完备性与对角边界译码阈值都已固定之后，还可继续抽取出下列四条派生结果。它们分别把部分屏幕的 exactization 代价压缩为加权拟阵极小化，把 dyadic 边界像码的相对距离退化写成闭式，把固定阶边界矩与分辨率自适应矩盒之间的分界提升为严格二歧，并把单标量 Joukowsky--Gödel 读出与多通道边界译码之间的精度差距量化为指数分离。

\begin{theorem}[部分屏幕 exactization 的加权拟阵极小化与拓扑缺口恒等式]\label{thm:xi-time-part65a-weighted-exactization-matroid-top-gap}
\leanverified{paper\_xi\_time\_part65a\_weighted\_exactization\_matroid\_top\_gap}
沿用定理 \ref{thm:xi-time-part9fa-exact-screen-minimal-size-matroid-basis} 与定理 \ref{thm:spg-screen-rank-matroid-supermodularity} 的记号，并令
\[
M_{m,n}:=\bigl(\overrightarrow{\mathcal F_m^{n-1}},r\bigr)
\]
为由可见秩函数 \(r(\cdot)\) 定义的边界行拟阵，记其总秩为
\[
\rho_{m,n}:=2^{mn}.
\]
给定任意非负权
\[
c:\overrightarrow{\mathcal F_m^{n-1}}\to \RR_{\ge 0},
\]
并对任意部分屏幕 \(S\subseteq \overrightarrow{\mathcal F_m^{n-1}}\) 定义 exactization 的最小追加代价
\[
\mathcal C_c(S):=
\min\Bigl\{
\sum_{\vec f\in T\setminus S}c(\vec f):
\ T\supseteq S,\ \ker f_T=0
\Bigr\}.
\]
则有
\[
\mathcal C_c(S)
=
\min\Bigl\{
\sum_{\vec f\in B}c(\vec f):
\ B\subseteq \overrightarrow{\mathcal F_m^{n-1}}\setminus S,\
B\text{ 是 }M_{m,n}/S\text{ 的一组基}
\Bigr\}.
\]
特别地，当 \(c\equiv 1\) 时，
\[
\mathcal C_1(S)
=
\rho_{m,n}-r(S)
=
2^{mn}-r(S)
=
\delta(S)
=
\rank_{\ZZ}\widetilde H_{n-1}(A_{\neg S};\ZZ).
\]
\end{theorem}
\begin{proof}
由定理 \ref{thm:xi-time-part9fa-exact-screen-minimal-size-matroid-basis}，精确屏幕恰由边界行拟阵的基控制，且其总秩为
\(
\rho_{m,n}=2^{mn}
\)。
对任意 \(B\subseteq \overrightarrow{\mathcal F_m^{n-1}}\setminus S\)，收缩拟阵 \(M_{m,n}/S\) 的秩函数满足
\[
r_{M_{m,n}/S}(B)=r(S\cup B)-r(S).
\]
因此，\(B\) 为 \(M_{m,n}/S\) 的基当且仅当
\[
r(S\cup B)=\rho_{m,n},
\qquad
|B|=\rho_{m,n}-r(S),
\]
也即当且仅当 \(S\cup B\) 是一个达到满秩而不含冗余追加面的 exactization。

现取任意 \(T\supseteq S\) 满足 \(\ker f_T=0\)。由定理
\ref{thm:xi-time-part9fa-exact-screen-minimal-size-matroid-basis} 的首部分，\(\ker f_T=0\) 等价于 \(r(T)=\rho_{m,n}\)，故 \(T\setminus S\) 在收缩拟阵 \(M_{m,n}/S\) 中张成满秩子集。于是可取某组基
\[
B\subseteq T\setminus S.
\]
由于权函数 \(c\) 非负，便有
\[
\sum_{\vec f\in B}c(\vec f)\le \sum_{\vec f\in T\setminus S}c(\vec f).
\]
这说明任意 exactization 的追加代价都不小于某组收缩拟阵基的权和；反过来，任一基 \(B\) 都给出 \(T=S\cup B\) 的 exactization。故第一条公式成立。

当 \(c\equiv 1\) 时，收缩拟阵的每一组基都具有基数
\(
\rho_{m,n}-r(S)
\)，于是
\[
\mathcal C_1(S)=\rho_{m,n}-r(S)=2^{mn}-r(S).
\]
再由定理 \ref{thm:spg-screen-rank-matroid-supermodularity}，
\[
\delta(S)=\rho_{m,n}-r(S),
\]
而定理 \ref{thm:xi-time-part9fa-screen-kernel-top-homology} 给出
\[
\delta(S)=\rank_{\ZZ}\widetilde H_{n-1}(A_{\neg S};\ZZ).
\]
合并即得结论。证毕。
\end{proof}

\begin{corollary}[dyadic 边界像码的正码率--零相对距离退化与固定唯一纠错半径]\label{cor:xi-time-part65a-dyadic-boundary-code-not-asymptotically-good}
\leanverified{paper\_xi\_time\_part65a\_dyadic\_boundary\_code\_not\_asymptotically\_good}
固定 \(n\ge 1\)，并沿用定理 \ref{thm:xi-time-part9fa-dyadic-boundary-image-ldpc} 的记号。则 dyadic 边界像码族
\[
\bigl(\mathcal C_{m,n}\bigr)_{m\ge 1}
\]
的相对距离满足
\[
\Delta_{m,n}
:=
\frac{d_{\min}(\mathcal C_{m,n})}{N_{m,n}}
=
\frac{2}{(2^m+1)\,2^{m(n-1)}}
\xrightarrow[m\to\infty]{}
0.
\]
另一方面，其码率满足
\[
\operatorname{Rate}_{m,n}
:=
\frac{\dim_{\FF_2}\mathcal C_{m,n}}{N_{m,n}}
=
\frac{2^m}{n(2^m+1)}
\xrightarrow[m\to\infty]{}
\frac1n.
\]
因此，对任意固定 \(n\)，码族 \((\mathcal C_{m,n})_{m\ge 1}\) 不是渐近优码族。并且任何最优最大似然唯一译码的可保证纠错半径都满足
\[
t_{m,n}^{\operatorname{ML}}
\le
\left\lfloor \frac{d_{\min}(\mathcal C_{m,n})-1}{2}\right\rfloor
=
n-1.
\]
特别地，在最坏情形的对抗错误模型下，可保证唯一纠错的边界面数与分辨率 \(m\) 无关。
\end{corollary}
\begin{proof}
由定理 \ref{thm:xi-time-part9fa-dyadic-boundary-image-ldpc}，
\[
N_{m,n}=n(2^m+1)2^{m(n-1)},
\qquad
\dim_{\FF_2}\mathcal C_{m,n}=2^{mn},
\qquad
d_{\min}(\mathcal C_{m,n})=2n.
\]
直接代入即可得到
\[
\Delta_{m,n}
=
\frac{2n}{n(2^m+1)2^{m(n-1)}}
=
\frac{2}{(2^m+1)2^{m(n-1)}}
\to 0
\]
以及
\[
\operatorname{Rate}_{m,n}
=
\frac{2^{mn}}{n(2^m+1)2^{m(n-1)}}
=
\frac{2^m}{n(2^m+1)}
\to \frac1n.
\]
故该码族具有正码率极限而相对距离趋于零，从而不是渐近优码族。

对于任意二元线性码，统一唯一纠错半径至多为
\[
\left\lfloor \frac{d_{\min}-1}{2}\right\rfloor.
\]
将 \(d_{\min}(\mathcal C_{m,n})=2n\) 代入便得
\[
t_{m,n}^{\operatorname{ML}}\le n-1.
\]
证毕。
\end{proof}

\begin{theorem}[固定阶边界矩不可分与分辨率自适应矩盒完备的严格二歧]\label{thm:xi-time-part65a-fixed-order-boundary-moment-vs-adaptive-box-dichotomy}
固定 \(n\ge 2\)。对每个分辨率 \(m\ge 1\)，设
\[
\mathcal A_m\subseteq \NN^n
\]
为一族多重指标。若存在常数 \(R<\infty\) 使得
\[
\mathcal A_m\subseteq \{\alpha\in\NN^n:\ |\alpha|\le R\}
\qquad(\forall m\ge 1),
\]
则边界通量矩字典
\[
U\longmapsto \bigl(M_\alpha(U)\bigr)_{\alpha\in\mathcal A_m}
\]
不可能在全部 \(m\)-级 dyadic polycube 上同时单射。

相反，若取分辨率自适应矩盒
\[
\mathcal B_m:=\{0,\dots,2^m-1\}^n,
\]
则体积分矩字典
\[
U\longmapsto \bigl(\mu_\alpha(U)\bigr)_{\alpha\in\mathcal B_m}
\]
在全部 \(m\)-级 dyadic polycube 上为单射；并且边界 Gödel 整数在固定分辨率类内唯一决定 \(U\)。
\end{theorem}
\begin{proof}
若存在统一总次数上界 \(R\)，则对 \(m=R+1\) 应用定理
\ref{thm:spg-prouhet-thue-morse-obstruction-dyadic-polyclube-flux-moments}，可得两个互异的 \(m\)-级 dyadic polycube
\[
U\neq V
\]
满足
\[
M_\alpha(U)=M_\alpha(V)
\qquad
(\forall \alpha\in\NN^n,\ |\alpha|\le R).
\]
由于 \(\mathcal A_m\subseteq \{|\alpha|\le R\}\)，遂有
\[
\bigl(M_\alpha(U)\bigr)_{\alpha\in\mathcal A_m}
=
\bigl(M_\alpha(V)\bigr)_{\alpha\in\mathcal A_m},
\]
故该字典在 \(m=R+1\) 时已非单射，从而不可能在全部分辨率上同时单射。

另一方面，定理 \ref{thm:xi-time-part9fa-dyadic-finite-moment-completeness} 已证明：对每个固定 \(m\)，矩盒
\[
\mathcal B_m=\{0,\dots,2^m-1\}^n
\]
给出的体积分矩字典是单射。再由推论
\ref{cor:spg-boundary-godel-finite-moment-completeness}，边界 Gödel 整数在固定分辨率类内唯一决定 dyadic polycube。本定理遂得证。证毕。
\end{proof}

\begin{theorem}[Joukowsky--Gödel 单标量读出与多通道边界译码的指数精度分离]\label{thm:xi-time-part65a-scalar-vs-multichannel-exponential-precision-gap}
\leanverified{paper\_xi\_time\_part65a\_scalar\_vs\_multichannel\_exponential\_precision\_gap}
取互异素数 \(p_1,\dots,p_k\)，并沿用定理 \ref{thm:spg-single-parameter-log-readout-vs-diagonal-multichannel-bifurcation} 的单标量读出
\[
u(b):=\sum_{i=1}^k b_i\log p_i,
\qquad
b\in\{0,1\}^k.
\]
若要求仅由 \(u(b)\) 的单标量观测，在统一绝对误差阈值 \(\eta_k\) 下对全部 \(2^k\) 个状态保证无碰撞唯一译码，则必有
\[
\eta_k
\le
\frac{1}{2(2^k-1)}
\sum_{i=1}^k\log p_i.
\]

另一方面，多通道对角边界译码允许的统一相对误差阈值为
\[
\varepsilon_\ast
:=
\frac{p_{\min}-1}{p_{\min}+1},
\qquad
p_{\min}:=\min_{1\le i\le k}p_i,
\]
并且该阈值由定理 \ref{thm:spg-single-parameter-log-readout-vs-diagonal-multichannel-bifurcation} 给出充分性、由命题 \ref{prop:spg-relative-error-threshold-sharpness} 给出最坏情形下的尖锐必要性。

因此，任意试图用单一标量 \(u(b)\) 取代多通道边界译码的方案，都必然支付指数级绝对精度代价。若进一步取 \(p_i\) 为前 \(k\) 个素数，则
\[
\eta_k=O\!\left(\frac{k\log k}{2^k}\right),
\qquad
\varepsilon_\ast=\frac13.
\]
\end{theorem}
\begin{proof}
由定理 \ref{thm:spg-single-parameter-log-readout-vs-diagonal-multichannel-bifurcation}，存在 \(b\neq b'\) 使
\[
|u(b)-u(b')|
\le
\frac{1}{2^k-1}\sum_{i=1}^k\log p_i.
\]
若要求在绝对误差不超过 \(\eta_k\) 的条件下对全部状态无碰撞唯一译码，则以各点 \(u(b)\) 为中心、半径 \(\eta_k\) 的闭区间必须两两不交。于是必要地有
\[
2\eta_k
\le
\min_{b\neq b'}|u(b)-u(b')|
\le
\frac{1}{2^k-1}\sum_{i=1}^k\log p_i,
\]
从而得到
\[
\eta_k
\le
\frac{1}{2(2^k-1)}
\sum_{i=1}^k\log p_i.
\]

多通道情形的充分性由定理
\ref{thm:spg-single-parameter-log-readout-vs-diagonal-multichannel-bifurcation} 直接给出，而阈值
\(
\varepsilon_\ast=\frac{p_{\min}-1}{p_{\min}+1}
\)
在最坏情形下的必要性由命题
\ref{prop:spg-relative-error-threshold-sharpness} 给出，故该相对误差门槛是 sharp 的。

若 \(p_i\) 取前 \(k\) 个素数，则 \(p_{\min}=2\)，故
\[
\varepsilon_\ast=\frac{2-1}{2+1}=\frac13.
\]
同时由素数定理的标准推论，
\[
\sum_{i=1}^k\log p_i=\vartheta(p_k)\sim p_k\sim k\log k,
\]
代回上式即得
\[
\eta_k=O\!\left(\frac{k\log k}{2^k}\right).
\]
证毕。
\end{proof}

\noindent
这四条派生结果表明，dyadic 边界协议中的可恢复性门槛可分解为三类严格量化的结构参数：部分屏幕的秩亏由收缩拟阵基选择与顶维同调秩共同控制，边界像码的渐近限制由相对距离退化而非码率塌缩主导，矩层析的完备性必须随分辨率同步扩张，而边界读出的鲁棒性则依赖于多通道坐标而不能被无损压缩为单一实参数。

\paragraph{Gödel--Tate 地址像的严格自相似闭合与链内算子格的 Boolean 算术随机化}
在单向量 \(2^L\)-进全息像集的同胚、自相似、维数与 Haar 二歧律已经确定之后，Gödel--Tate 地址链还可进一步压缩为一条显式的 Cantor--Bernoulli 闭合；另一方面，链内算子格的平方自由 Gödel 算术化也可进一步收束为一组完全显式的 Boolean--概率闭式。下列四条结果分别统一这两条主线的拓扑、测度、格论与随机结构。

\begin{theorem}[Gödel--Tate 地址像的严格自相似、移位共轭与 clopen cylinder 分解]\label{thm:xi-time-part65d-godeltate-selfsimilar-shift-cylinder}
设 \(E\) 为有限原子事件集，\(\mathrm{code}:E\hookrightarrow \NN\) 为单射，并记
\[
M:=\max_{e\in E}\mathrm{code}(e),
\qquad
B:=2^L>M,
\qquad
T:=T_2(E_{\min}).
\]
取任意非零向量 \(v\in T\)，并记
\[
D:=\mathrm{code}(E)\subseteq \{0,\dots,B-1\},
\qquad
T_v:=\ZZ_2v\subseteq T.
\]
对任意无限事件流 \(\mathbf e=(e_t)_{t\ge 1}\in E^{\NN}\)，定义
\[
\Xi(\mathbf e):=X(\mathbf e):=\sum_{t\ge 1}\mathrm{code}(e_t)\,B^{t-1}v,
\qquad
\mathcal X_E:=\Xi(E^{\NN})\subseteq T_v.
\]
则成立：
\begin{enumerate}
  \item \(\mathcal X_E\) 是紧、全不连通集，并满足严格不交分解
  \[
  \mathcal X_E=\bigsqcup_{a\in D}\bigl(av+B\mathcal X_E\bigr).
  \]
  \item 记左移算子
  \[
  \sigma:E^{\NN}\to E^{\NN},
  \qquad
  \sigma(e_1,e_2,\dots)=(e_2,e_3,\dots).
  \]
  则 \(\Xi:E^{\NN}\to \mathcal X_E\) 为拓扑同胚；并且若定义
  \[
  S_B(av+Bx):=x
  \qquad
  (a\in D,\ x\in \mathcal X_E),
  \]
  则 \(S_B\) 良定义且满足共轭关系
  \[
  S_B\circ \Xi=\Xi\circ \sigma.
  \]
  等价地，对任意 \(\mathbf e\in E^{\NN}\)，有
  \[
  \Xi(\mathbf e)=\mathrm{code}(e_1)v+B\,\Xi(\sigma\mathbf e).
  \]
  \item 记 \(N:=|E|=|D|\)。对每个 \(n\ge 1\) 与每个 digit 前缀
  \[
  (a_1,\dots,a_n)\in D^n,
  \]
  定义
  \[
  \mathcal X_E[a_1,\dots,a_n]
  :=
  \left(\sum_{t=1}^{n}a_tB^{t-1}v\right)+B^n\mathcal X_E.
  \]
  则这 \(N^n\) 个集合两两不交，且都是 \(\mathcal X_E\) 中的 clopen cylinder，并满足
  \[
  \mathcal X_E=\bigsqcup_{(a_1,\dots,a_n)\in D^n}\mathcal X_E[a_1,\dots,a_n].
  \]
\end{enumerate}
\end{theorem}
\begin{proof}
由定理 \ref{thm:xi-time-part62d-hologram-image-cantor-homeomorphism}，地址映射
\[
\Xi:E^{\NN}\longrightarrow \mathcal X_E
\]
已知是拓扑同胚，并且满足严格不交分解
\[
\mathcal X_E=\bigsqcup_{a\in D}(av+B\mathcal X_E).
\]
因此第 \(1\) 条与第 \(2\) 条中的同胚性立即成立；而对任意
\(
\mathbf e=(e_1,e_2,\dots)
\)，按定义拆开首位即可得到
\[
\Xi(\mathbf e)=\mathrm{code}(e_1)v+B\,\Xi(\sigma\mathbf e),
\]
故 \(S_B\circ \Xi=\Xi\circ \sigma\)，并且 \(S_B\) 由上述严格不交分解决定而良定义。

对第 \(3\) 条，由同一定理，对任意事件前缀 \(\mathbf a=(e_1,\dots,e_n)\in E^n\) 及其 cylinder
\[
[\mathbf a]:=\{\mathbf f\in E^{\NN}:\ \mathrm{pref}_n(\mathbf f)=\mathbf a\},
\]
都有
\[
\Xi([\mathbf a])
=
\left(\sum_{t=1}^{n}\mathrm{code}(e_t)\,B^{t-1}v\right)+B^n\mathcal X_E.
\]
由于 \(\mathrm{code}\) 单射，\(D^n\) 中每个 digit 前缀与 \(E^n\) 中唯一事件前缀一一对应，故上式正是所定义的 \(\mathcal X_E[a_1,\dots,a_n]\)。而 \(\{[\mathbf a]:\mathbf a\in E^n\}\) 在 \(E^{\NN}\) 中本来就是由 \(N^n\) 个两两不交 clopen 集组成的分割，经同胚 \(\Xi\) 送到 \(\mathcal X_E\) 后，便得到所述 \(N^n\) 个两两不交 clopen cylinder 的不交并分解。证毕。
\end{proof}

\begin{theorem}[Gödel--Tate 地址像的精确 \texorpdfstring{$2$}{2}-进维数与 Haar 测度二分律]\label{thm:xi-time-part65d-godeltate-dimension-haar-dichotomy}
\leanverified{paper\_xi\_time\_part65d\_godeltate\_dimension\_haar\_dichotomy}
沿用定理 \ref{thm:xi-time-part65d-godeltate-selfsimilar-shift-cylinder} 的记号，并记
\[
N:=|E|=|D|.
\]
则 \(\mathcal X_E\subseteq T_v\cong \ZZ_2\) 的 Hausdorff 维数与上、下 Minkowski 维数存在且相等，并满足
\[
\dim_{\mathrm H}(\mathcal X_E)
=
\dim_{\mathrm M}(\mathcal X_E)
=
\frac{\log N}{\log B}
=
\frac{\log N}{L\log 2}.
\]
此外，\(\mathcal X_E\) 关于地址线 \(T_v\) 的归一化 Haar 测度满足二分律：
\begin{enumerate}
  \item 若 \(N<B\)，则
  \[
  \mu_{T_v}(\mathcal X_E)=0,
  \]
  并且 \(\mathcal X_E\) 在 \(T_v\) 中内部为空；特别地，当 \(2\le N<B\) 时，\(\mathcal X_E\) 为 nowhere dense 的 Cantor 型集；
  \item 若 \(N=B\)，则等价地 \(D=\{0,1,\dots,B-1\}\)，并且
  \[
  \mathcal X_E=T_v,
  \qquad
  \mu_{T_v}(\mathcal X_E)=1.
  \]
\end{enumerate}
\end{theorem}
\begin{proof}
由结论 \ref{con:xi-time-part62b-rankone-hologram-selfsimilar-measure-dimension}，对任意有限 digit 集 \(D\subseteq\{0,\dots,B-1\}\)，其单向量 \(2^L\)-进全息像集在地址线口径下都满足
\[
\dim_{\mathrm H}=\dim_{\mathrm M}=\frac{\log|D|}{\log B},
\]
并且 Haar 测度恰为
\[
\mu_{T_v}=
\begin{cases}
1,& |D|=B,\\
0,& |D|<B.
\end{cases}
\]
将 \(|D|=N\) 代入，立即得到维数公式及测度二分律。

若 \(N<B\)，同一结论还给出 \(\mathcal X_E\) 在 \(T_v\) 中内部为空。另一方面，由定理 \ref{thm:xi-time-part65d-godeltate-selfsimilar-shift-cylinder}，\(\mathcal X_E\) 为紧集，因而在紧 Hausdorff 群 \(T_v\) 中闭。闭且内部为空即意味着 nowhere dense。

若 \(N=B\)，则因 \(D\subseteq\{0,\dots,B-1\}\) 且 \(|D|=B\)，只能有
\[
D=\{0,1,\dots,B-1\}.
\]
此时同一结论给出 \(\mathcal X_E=T_v\)。证毕。
\end{proof}

\leanverified{paper\_xi\_time\_part65d\_chain\_interior\_boolean\_arithmeticization}
\begin{theorem}[链内算子格的 Boolean 算术化、Möbius 函数与特征多项式]\label{thm:xi-time-part65d-chain-interior-boolean-arithmeticization}
对链
\[
C_n=\{0<1<\cdots<n-1\},
\]
记
\[
\mathcal I_n:=\operatorname{Int}(C_n)
\]
为其全部内算子所成偏序格。沿用定理 \ref{thm:killo-chain-interior-godel-gcd-lcm} 的记号，对每个 \(I\in \mathcal I_n\) 记
\[
\kappa(I):=\prod_{i\in \mathrm{Fix}(I)}q_i,
\qquad
Q_n:=\prod_{i=0}^{n-1}q_i.
\]
则成立：
\begin{enumerate}
  \item \((\mathcal I_n,\preceq)\) 是秩为 \(n-1\) 的 Boolean 格；其秩函数可取为
  \[
  \operatorname{rk}(I):=|\mathrm{Fix}(I)|-1.
  \]
  \item 映射 \(\kappa\) 给出格同构
  \[
  \mathcal I_n
  \cong
  \{d\mid Q_n:\ q_0\mid d\},
  \]
  其中右侧按整除序排序；并且对任意 \(I,J\in \mathcal I_n\)，有
  \[
  \kappa(I\wedge J)=\gcd(\kappa(I),\kappa(J)),
  \qquad
  \kappa(I\vee J)=\mathrm{lcm}(\kappa(I),\kappa(J)).
  \]
  \item 对任意 \(I\preceq J\)，其 Möbius 函数满足
  \[
  \mu_{\mathcal I_n}(I,J)
  =
  (-1)^{\operatorname{rk}(J)-\operatorname{rk}(I)}.
  \]
  \item 其特征多项式满足
  \[
  \chi_{\mathcal I_n}(t)=(t-1)^{n-1}.
  \]
\end{enumerate}
\end{theorem}
\begin{proof}
由定理 \ref{thm:killo-chain-interior-godel-gcd-lcm}，映射
\[
I\longmapsto \mathrm{Fix}(I)
\]
把 \(\mathcal I_n\) 规范识别为全部满足 \(0\in F\subseteq\{0,\dots,n-1\}\) 的子集格，并把 meet/join 分别送到交与并。去掉强制出现的元素 \(0\) 后，便得到
\[
\mathcal I_n\cong \mathcal P(\{1,\dots,n-1\}),
\]
故 \(\mathcal I_n\) 是秩为 \(n-1\) 的 Boolean 格，且秩函数正是
\(
|\mathrm{Fix}(I)|-1
\)。

同一条定理还给出
\[
\kappa(I)=\prod_{i\in \mathrm{Fix}(I)}q_i,
\]
因此 \(\kappa(\mathcal I_n)\) 恰为 \(Q_n\) 的全部含 \(q_0\) 因子的平方自由约数；由于 \(Q_n\) 本身平方自由，这正是集合
\[
\{d\mid Q_n:\ q_0\mid d\}.
\]
同时 \(\kappa\) 将 meet/join 精确送到 \(\gcd/\mathrm{lcm}\)，故第 \(2\) 条成立。

对 Möbius 函数与特征多项式，定理
\ref{thm:xi-chain-interior-boolean-mobius-characteristic-maxchains} 与推论
\ref{cor:xi-chain-interior-dirichlet-mobius-characteristic-closure} 已给出
\[
\mu_{\mathcal I_n}(I,J)=(-1)^{|\mathrm{Fix}(J)\setminus \mathrm{Fix}(I)|}
\]
以及
\[
\chi_{\mathcal I_n}(t)=(t-1)^{n-1}.
\]
而由秩函数定义，
\[
|\mathrm{Fix}(J)\setminus \mathrm{Fix}(I)|
=
\bigl(|\mathrm{Fix}(J)|-1\bigr)-\bigl(|\mathrm{Fix}(I)|-1\bigr)
=
\operatorname{rk}(J)-\operatorname{rk}(I),
\]
遂得第 \(3\) 条。证毕。
\end{proof}

\leanverified{paper\_xi\_time\_part65d\_chain\_interior\_interval\_squarefree\_divisor\_isomorphism}
\begin{theorem}[链内算子区间与平方自由整除区间的精确算术同构]\label{thm:xi-time-part65d-chain-interior-interval-squarefree-divisor-isomorphism}
沿用定理 \ref{thm:xi-time-part65d-chain-interior-boolean-arithmeticization} 的记号。对任意
\[
I\preceq J\in \mathcal I_n,
\]
定义
\[
m_{I,J}:=\frac{\kappa(J)}{\kappa(I)}.
\]
则 \(m_{I,J}\) 是平方自由整数，并且区间
\[
[I,J]:=\{K\in \mathcal I_n:\ I\preceq K\preceq J\}
\]
与整除格
\[
\{d\in\NN:\ d\mid m_{I,J}\}
\]
之间存在规范格同构
\[
\Phi_{I,J}:[I,J]\longrightarrow \{d\in\NN:\ d\mid m_{I,J}\},
\qquad
\Phi_{I,J}(K):=\frac{\kappa(K)}{\kappa(I)}.
\]
并且对任意 \(K_1,K_2\in [I,J]\)，有
\[
\Phi_{I,J}(K_1\wedge K_2)
=
\gcd\!\bigl(\Phi_{I,J}(K_1),\Phi_{I,J}(K_2)\bigr),
\]
\[
\Phi_{I,J}(K_1\vee K_2)
=
\mathrm{lcm}\!\bigl(\Phi_{I,J}(K_1),\Phi_{I,J}(K_2)\bigr).
\]
\end{theorem}
\begin{proof}
由定理 \ref{thm:killo-chain-interior-godel-gcd-lcm}，
\[
I\longmapsto \mathrm{Fix}(I)
\]
把 \(\mathcal I_n\) 识别为满足 \(0\in F\subseteq\{0,\dots,n-1\}\) 的子集格，并把偏序送到包含关系。因而
\[
I\preceq K\preceq J
\quad\Longleftrightarrow\quad
\mathrm{Fix}(I)\subseteq \mathrm{Fix}(K)\subseteq \mathrm{Fix}(J).
\]
于是
\[
m_{I,J}
=
\frac{\kappa(J)}{\kappa(I)}
=
\prod_{i\in \mathrm{Fix}(J)\setminus \mathrm{Fix}(I)}q_i
\]
确为平方自由整数。

对任意 \(K\in[I,J]\)，同样有
\[
\Phi_{I,J}(K)
=
\frac{\kappa(K)}{\kappa(I)}
=
\prod_{i\in \mathrm{Fix}(K)\setminus \mathrm{Fix}(I)}q_i.
\]
当 \(K\) 在区间 \([I,J]\) 内变化时，\(\mathrm{Fix}(K)\setminus \mathrm{Fix}(I)\) 恰遍历
\(\mathrm{Fix}(J)\setminus \mathrm{Fix}(I)\) 的全部子集，故 \(\Phi_{I,J}(K)\) 恰遍历 \(m_{I,J}\) 的全部正因子。于是 \(\Phi_{I,J}\) 为双射。

再由定理 \ref{thm:killo-chain-interior-godel-gcd-lcm}，
\[
\kappa(K_1\wedge K_2)=\gcd(\kappa(K_1),\kappa(K_2)),
\qquad
\kappa(K_1\vee K_2)=\mathrm{lcm}(\kappa(K_1),\kappa(K_2)).
\]
而 \(\Phi_{I,J}(K_\nu)\) 的素因子都来自 \(\mathrm{Fix}(J)\setminus \mathrm{Fix}(I)\)，故与 \(\kappa(I)\) 互素，从而
\[
\gcd(\kappa(K_1),\kappa(K_2))
=
\kappa(I)\,
\gcd\!\bigl(\Phi_{I,J}(K_1),\Phi_{I,J}(K_2)\bigr),
\]
\[
\mathrm{lcm}(\kappa(K_1),\kappa(K_2))
=
\kappa(I)\,
\mathrm{lcm}\!\bigl(\Phi_{I,J}(K_1),\Phi_{I,J}(K_2)\bigr).
\]
两边同除以 \(\kappa(I)\) 即得所述运算相容性。证毕。
\end{proof}

\begin{corollary}[链内算子区间的 Möbius 函数与特征多项式的整数退化]\label{cor:xi-time-part65d-chain-interval-mobius-characteristic-arithmetic}
\leanverified{paper\_xi\_time\_part65d\_chain\_interval\_mobius\_characteristic\_arithmetic}
沿用定理 \ref{thm:xi-time-part65d-chain-interior-interval-squarefree-divisor-isomorphism} 的记号。则对任意 \(I\preceq J\)，有
\[
\mu_{\mathcal I_n}(I,J)=\mu(m_{I,J}),
\]
其中右端 \(\mu\) 表示经典整数 Möbius 函数。若 \(\omega(m)\) 表示整数 \(m\) 的不同素因子个数，则区间 \([I,J]\) 的特征多项式满足
\[
\chi_{[I,J]}(t)=(t-1)^{\omega(m_{I,J})}.
\]
\end{corollary}
\begin{proof}
把
\[
m_{I,J}=\prod_{r=1}^{s}p_r
\]
写成互异素数乘积，则 \(\{d\in\NN:\ d\mid m_{I,J}\}\) 恰是秩为 \(s=\omega(m_{I,J})\) 的 Boolean 格。由定理
\ref{thm:xi-time-part65d-chain-interior-interval-squarefree-divisor-isomorphism}，区间 \([I,J]\) 与该 Boolean 格同构，故其 Möbius 函数为
\[
(-1)^s=\mu(m_{I,J}),
\]
而特征多项式即
\[
(t-1)^s=(t-1)^{\omega(m_{I,J})}.
\]
证毕。
\end{proof}

\begin{theorem}[链内算子算术化的 Dirichlet 生成函数、矩母函数与 Hoeffding 型偏差界]\label{thm:xi-time-part65d-chain-interior-dirichlet-mgf-hoeffding}
\leanverified{paper\_xi\_time\_part65d\_chain\_interior\_dirichlet\_mgf\_hoeffding}
沿用定理 \ref{thm:xi-time-part65d-chain-interior-boolean-arithmeticization} 的记号，并在 \(\mathcal I_n\) 上取均匀分布。定义随机变量
\[
K_n(I):=\kappa(I)=\prod_{i\in \mathrm{Fix}(I)}q_i.
\]
则成立：
\begin{enumerate}
  \item 对任意 \(s\in \CC\)，其 Dirichlet 生成函数满足
  \[
  \sum_{I\in \mathcal I_n}K_n(I)^{-s}
  =
  q_0^{-s}\prod_{i=1}^{n-1}(1+q_i^{-s}).
  \]
  \item 对任意 \(t\in \RR\)，\(\log K_n\) 的矩母函数满足
  \[
  \EE\!\left[e^{t\log K_n}\right]
  =
  q_0^t\prod_{i=1}^{n-1}\frac{1+q_i^t}{2}.
  \]
  \item 其均值与方差为
  \[
  \EE[\log K_n]
  =
  \log q_0+\frac12\sum_{i=1}^{n-1}\log q_i,
  \qquad
  \mathrm{Var}(\log K_n)
  =
  \frac14\sum_{i=1}^{n-1}(\log q_i)^2.
  \]
  \item 对任意 \(u>0\)，有 Hoeffding 型偏差界
  \[
  \PP\!\left(\left|\log K_n-\EE[\log K_n]\right|\ge u\right)
  \le
  2\exp\!\left(
  -\frac{2u^2}{\sum_{i=1}^{n-1}(\log q_i)^2}
  \right).
  \]
\end{enumerate}
\end{theorem}
\begin{proof}
由定理 \ref{thm:xi-chain-interior-random-godel-euler-logvariance}，存在独立同分布随机变量
\[
\xi_1,\dots,\xi_{n-1}\stackrel{\mathrm{i.i.d.}}{\sim}\mathrm{Bernoulli}\!\left(\frac12\right)
\]
使得
\[
K_n=q_0\prod_{i=1}^{n-1}q_i^{\xi_i},
\qquad
\log K_n=\log q_0+\sum_{i=1}^{n-1}\xi_i\log q_i.
\]

第 \(1\) 条也可由定理 \ref{thm:xi-chain-interior-godel-euler-mobius-factorization} 直接得到；等价地，按 \(\xi_i\in\{0,1\}\) 求和可写为
\[
\sum_{I\in \mathcal I_n}K_n(I)^{-s}
=
q_0^{-s}\sum_{\xi_i\in\{0,1\}}
\prod_{i=1}^{n-1}q_i^{-s\xi_i}
=
q_0^{-s}\prod_{i=1}^{n-1}(1+q_i^{-s}).
\]

对第 \(2\) 条，由独立性，
\[
\EE\!\left[e^{t\log K_n}\right]
=
\EE\!\left[K_n^t\right]
=
q_0^t\prod_{i=1}^{n-1}\EE\!\left[q_i^{t\xi_i}\right]
=
q_0^t\prod_{i=1}^{n-1}\frac{1+q_i^t}{2}.
\]

第 \(3\) 条则由
\[
\EE[\xi_i]=\frac12,
\qquad
\mathrm{Var}(\xi_i)=\frac14
\]
与独立性直接得到：
\[
\EE[\log K_n]
=
\log q_0+\sum_{i=1}^{n-1}\EE[\xi_i]\log q_i
=
\log q_0+\frac12\sum_{i=1}^{n-1}\log q_i,
\]
\[
\mathrm{Var}(\log K_n)
=
\sum_{i=1}^{n-1}\mathrm{Var}(\xi_i)(\log q_i)^2
=
\frac14\sum_{i=1}^{n-1}(\log q_i)^2.
\]

最后定义中心化独立变量
\[
X_i:=\left(\xi_i-\frac12\right)\log q_i
\in
\left[-\frac12\log q_i,\frac12\log q_i\right].
\]
则
\[
\log K_n-\EE[\log K_n]=\sum_{i=1}^{n-1}X_i,
\]
且每个 \(X_i\) 的取值区间长度为 \(\log q_i\)。由 Hoeffding 不等式，
\[
\PP\!\left(\left|\sum_{i=1}^{n-1}X_i\right|\ge u\right)
\le
2\exp\!\left(
-\frac{2u^2}{\sum_{i=1}^{n-1}(\log q_i)^2}
\right),
\]
即得第 \(4\) 条。证毕。
\end{proof}

\noindent
由此，Gödel--Tate 地址像在 Tate 模块中闭合为一个严格的 \(B\)-进 Cantor--Bernoulli 模型，而链内算子格则在平方自由算术化下闭合为一个 Boolean 整除格与对数可加的 Bernoulli 随机模型。前者把地址几何与移位动力学统一到同一条 clopen 自相似链上，后者则把纯序结构直接转写为可完全积化、可求矩并可作非渐近偏差估计的算术对象。

\paragraph{Toeplitz 主子式的离散凹锥、Blaschke 原点模维数证书、Rouch\'e 边界裕度与 window-\(6\) parity 的几何压缩}
在 Toeplitz--Verblunsky 乘积分解、离切片 Blaschke 指标、Jensen 缺陷的视界纯度判据以及 window-\(6\) boundary uplift 的二层/三层分岔都已明确之后，还可继续抽取出下列五条后果。它们分别把 Toeplitz 主子式链在对数坐标下精确识别为离散凹锥，把单标量 \(|B_{\mathcal Z}(0)|\) 升格为残差维数的 sharp 证书，把 Hurwitz 型零自由逼近进一步有限化为一列圆周上的 Rouch\'e 裕度不等式，并把 window-\(6\) 的 \((\ZZ_2)^3\) boundary parity 先压缩为局域几何只保留的对角 \(\ZZ_2\) 方向，再识别为不能在“无额外协议门、无额外寄存器”的口径下跨 \(m=7,8\) 自然延拓的最小锚点现象。

\begin{theorem}[Toeplitz 主子式锥在对数坐标下恰为离散凹锥]\label{thm:xi-time-part65e-toeplitz-local-logconcavity-realizability}
\leanverified{paper\_xi\_time\_part65e\_toeplitz\_local\_logconcavity\_realizability}
设 \(N\ge 0\)，并给定正数列
\[
\Delta_{-1}:=1,
\qquad
\Delta_0>0,\ \Delta_1>0,\dots,\Delta_N>0.
\]
定义
\[
u_k:=\log \Delta_k
\qquad(-1\le k\le N).
\]
则以下两命题等价：
\begin{enumerate}
  \item 存在某个 Carath\'eodory 函数的截断 Toeplitz 主子式链
  \[
  \mathbf T_k=(c_{j-\ell})_{0\le j,\ell\le k}
  \qquad(0\le k\le N)
  \]
  使得
  \[
  \det \mathbf T_k=\Delta_k
  \qquad(0\le k\le N);
  \]
  \item 对所有 \(0\le k\le N-1\)，都有
  \[
  u_{k+1}-2u_k+u_{k-1}\le 0.
  \]
\end{enumerate}
并且在该情形下，对每个 \(0\le k\le N-1\) 都有严格恒等式
\[
\boxed{\;
u_{k+1}-2u_k+u_{k-1}
=
\log(1-|\alpha_k|^2)\le 0,
\;}
\]
从而相应 Verblunsky 系数的模长由主子式序列唯一恢复为
\[
|\alpha_k|^2
=
1-\exp\!\bigl(u_{k+1}-2u_k+u_{k-1}\bigr)
\qquad(0\le k\le N-1).
\]
换言之，Toeplitz--PSD 主子式链在忘却相位后恰等价于离散凹锥条件。
\end{theorem}
\begin{proof}
当 \(N=0\) 时结论平凡：取常值 Carath\'eodory 函数 \(\mathscr C\equiv \Delta_0\) 即可。以下设 \(N\ge 1\)。

若存在所述 Toeplitz 主子式链，则由定理 \ref{thm:xi-toeplitz-principal-minor-discrete-curvature-recovery}，对 \(k=0\) 有
\[
1-|\alpha_0|^2
=
\frac{\Delta_1}{\Delta_0^2}
=
\exp(u_1-2u_0+u_{-1}),
\]
而对 \(1\le k\le N-1\) 有
\[
1-|\alpha_k|^2
=
\frac{\Delta_{k+1}\Delta_{k-1}}{\Delta_k^2}
=
\exp(u_{k+1}-2u_k+u_{k-1}).
\]
两侧取对数即得
\[
u_{k+1}-2u_k+u_{k-1}
=
\log(1-|\alpha_k|^2)\le 0,
\]
并同时得到模长恢复公式。

反过来，若已知
\[
u_{k+1}-2u_k+u_{k-1}\le 0
\qquad(0\le k\le N-1),
\]
则可定义
\[
\rho_k:=\exp(u_{k+1}-2u_k+u_{k-1})\in(0,1],
\qquad
r_k:=\rho_k,
\qquad
|\alpha_k|:=\sqrt{1-\rho_k}\in[0,1)
\qquad(0\le k\le N-1),
\]
并任选相位；例如取全部 \(\alpha_k\in[0,1)\subset\RR\)。再令
\[
\alpha_j:=0
\qquad(j\ge N).
\]
由 Schur--Carath\'eodory 对应，存在某个 Carath\'eodory 函数 \(\mathscr C\) 具有这列 Verblunsky 系数。记其 Toeplitz 主子式为
\[
\widetilde\Delta_k:=\det \widetilde{\mathbf T}_k
\qquad(k\ge 0).
\]
由定理 \ref{thm:xi-toeplitz-principal-minor-discrete-curvature-recovery}，有
\[
\widetilde\Delta_0=\Delta_0,
\qquad
\widetilde\Delta_1
=
\Delta_0^2(1-|\alpha_0|^2)
=
\Delta_0^2\rho_0
=
\Delta_1.
\]
对每个 \(1\le k\le N-1\)，同一定理还给出递推
\[
\widetilde\Delta_{k+1}
=(1-|\alpha_k|^2)\frac{\widetilde\Delta_k^2}{\widetilde\Delta_{k-1}}
=r_k\frac{\widetilde\Delta_k^2}{\widetilde\Delta_{k-1}}.
\]
若已知 \(\widetilde\Delta_{k-1}=\Delta_{k-1}\) 与 \(\widetilde\Delta_k=\Delta_k\)，则
\[
\widetilde\Delta_{k+1}
=r_k\frac{\Delta_k^2}{\Delta_{k-1}}
=\Delta_{k+1}.
\]
归纳即得 \(\widetilde\Delta_k=\Delta_k\) 对所有 \(0\le k\le N\) 成立，故所需实现性得证。最后，模长唯一恢复公式已由上面的曲率恒等式直接给出。证毕。
\end{proof}

\begin{theorem}[Blaschke 原点模对离切片残差维数的 sharp 上界]\label{thm:xi-time-part65e-blaschke-origin-modulus-dimension-bound}
\leanverified{paper\_xi\_time\_part65e\_blaschke\_origin\_modulus\_dimension\_bound}
沿用定义 \ref{def:xi-offslice-ind-weight} 的记号，并设 \(\mathcal Z_{+}\neq\varnothing\) 为有限离切片零点多重集，\(B_{\mathcal Z}\) 为其对应的有限 Blaschke 乘积。定义
\[
\delta_{\min}(\mathcal Z)
:=
\min_{s\in\mathcal Z_{+}}
\Bigl(\Re(s)-\frac12\Bigr)
>0.
\]
则有 sharp 估计
\[
\mathrm{Ind}(\mathcal Z)
\le
\frac{\mathrm W(\mathcal Z)}{2L\,\delta_{\min}(\mathcal Z)}
=
\frac{-\log|B_{\mathcal Z}(0)|}{2L\,\delta_{\min}(\mathcal Z)}.
\]
等价地，
\[
\dim K_{B_{\mathcal Z}}
\le
\frac{-\log|B_{\mathcal Z}(0)|}{2L\,\delta_{\min}(\mathcal Z)}.
\]
并且等号成立当且仅当全部离切片零点具有相同的实部偏差，即
\[
\Re(s)-\frac12\equiv \delta_{\min}(\mathcal Z)
\qquad(s\in\mathcal Z_{+}).
\]
\end{theorem}
\begin{proof}
由定理 \ref{thm:xi-offslice-realpart-sum-law}，
\[
\mathrm W(\mathcal Z)
=
2L\sum_{s\in\mathcal Z_{+}}
\Bigl(\Re(s)-\frac12\Bigr).
\]
由于每一项都不小于 \(\delta_{\min}(\mathcal Z)\)，便有
\[
\mathrm W(\mathcal Z)\ge 2L\,\mathrm{Ind}(\mathcal Z)\,\delta_{\min}(\mathcal Z),
\]
即
\[
\mathrm{Ind}(\mathcal Z)
\le
\frac{\mathrm W(\mathcal Z)}{2L\,\delta_{\min}(\mathcal Z)}.
\]
再由定理 \ref{thm:xi-offslice-blaschke-origin-modulus}，
\[
\mathrm W(\mathcal Z)=-\log|B_{\mathcal Z}(0)|,
\]
从而得到第一条闭式。另一方面，命题 \ref{prop:xi-offslice-model-space-dimension} 给出
\[
\dim K_{B_{\mathcal Z}}=\mathrm{Ind}(\mathcal Z),
\]
遂得第二条闭式。

等号情形同样直接：上式中的和达到下界，当且仅当每个求和项都等于 \(\delta_{\min}(\mathcal Z)\)，亦即全部 \(\Re(s)-\tfrac12\) 完全相同。证毕。
\end{proof}

\begin{theorem}[Blaschke 原点模在 strip 约束下的残差维数下界]\label{thm:xi-time-part65e-strip-blaschke-origin-modulus-index-lower-bound}
沿用定义 \ref{def:xi-offslice-ind-weight} 的记号，并设 \(\mathcal Z_{+}\neq\varnothing\) 为有限离切片零点多重集，\(B_{\mathcal Z}\) 为其对应的有限 Blaschke 乘积。
若存在常数 \(\eta>0\) 使
\[
0<\Re(s)-\frac12\le \eta
\qquad(s\in\mathcal Z_{+}),
\]
则有 sharp 下界
\[
\boxed{\;
\mathrm{Ind}(\mathcal Z)
\ge
\frac{\mathrm W(\mathcal Z)}{2L\,\eta}
=
\frac{-\log|B_{\mathcal Z}(0)|}{2L\,\eta}.\;}
\]
等价地，
\[
\boxed{\;
\dim K_{B_{\mathcal Z}}
\ge
\frac{-\log|B_{\mathcal Z}(0)|}{2L\,\eta}.\;}
\]
并且等号成立当且仅当全部离切片零点具有相同的实部偏差
\[
\Re(s)-\frac12\equiv \eta
\qquad(s\in\mathcal Z_{+}).
\]
\end{theorem}
\begin{proof}
由定理 \ref{thm:xi-offslice-realpart-sum-law}，
\[
\mathrm W(\mathcal Z)
=
2L\sum_{s\in\mathcal Z_{+}}
\Bigl(\Re(s)-\frac12\Bigr).
\]
由于每一项都不超过 \(\eta\)，故
\[
\mathrm W(\mathcal Z)\le 2L\,\eta\,\mathrm{Ind}(\mathcal Z),
\]
从而
\[
\mathrm{Ind}(\mathcal Z)\ge \frac{\mathrm W(\mathcal Z)}{2L\,\eta}.
\]
再由定理 \ref{thm:xi-offslice-blaschke-origin-modulus}，
\(
\mathrm W(\mathcal Z)=-\log|B_{\mathcal Z}(0)|
\)，
即得第一条闭式。
另一方面，由命题 \ref{prop:xi-offslice-model-space-dimension}，
\[
\dim K_{B_{\mathcal Z}}=\mathrm{Ind}(\mathcal Z),
\]
于是第二条闭式随即成立。

等号成立当且仅当上式求和中的每一项都达到其上界 \(\eta\)，亦即全部 \(\Re(s)-\tfrac12\) 完全相同且等于 \(\eta\)。证毕。
\end{proof}

\begin{theorem}[Rouch\'e 边界裕度所给出的 RH 有限证书]\label{thm:xi-time-part65e-rouche-boundary-margin-rh-certificate}
沿用 \eqref{eq:xi-ramanujan-horizon-Fzeta} 的圆盘完成化函数 \(F_{\zeta}\)。设 \(\{F_N\}_{N\ge 1}\) 为一列解析函数 \(F_N:\mathbb{D}\to\CC\)，并满足定理 \ref{thm:xi-hurwitz-horizon-reduction} 的自对偶完成化一致性。若存在一列半径
\[
\rho_N\uparrow 1
\qquad(N\to\infty)
\]
使得对每个 \(N\) 都有
\[
F_N\text{ 在 }\overline{D(0,\rho_N)}\text{ 内无零点}
\]
且
\[
\sup_{|w|=\rho_N}|F_{\zeta}(w)-F_N(w)|
<
\inf_{|w|=\rho_N}|F_N(w)|,
\]
则 RH 成立。
\end{theorem}
\begin{proof}
固定 \(N\)。由于 \(F_N\) 在 \(\overline{D(0,\rho_N)}\) 内无零点，故
\[
\inf_{|w|=\rho_N}|F_N(w)|>0.
\]
于是由 Rouch\'e 定理，\(F_{\zeta}\) 与 \(F_N\) 在圆盘 \(D(0,\rho_N)\) 内具有相同的零点总数（按重数计）。而 \(F_N\) 在该圆盘内无零点，故 \(F_{\zeta}\) 在 \(D(0,\rho_N)\) 内亦无零点。

由于 \(\rho_N\uparrow 1\)，对任意 \(0<\rho<1\) 都可取 \(N\) 充分大使 \(\rho<\rho_N\)。于是 \(F_{\zeta}\) 在 \(D(0,\rho)\) 内无零点。因为这对每个 \(\rho<1\) 都成立，便知 \(F_{\zeta}\) 在整个 \(\mathbb D\) 内无零点。最后由定理 \ref{thm:app-rh-iff-disk-zero-free}，RH 成立。证毕。
\end{proof}

\begin{theorem}[window-\(6\) 的几何可实现 boundary parity 子群恰为对角 \texorpdfstring{$\ZZ_2$}{Z2}，且精确损失两比特]\label{thm:xi-time-part65e-window6-geometric-diagonal-z2-parity}
记
\[
P_6:=C_{\mathrm{bdry}}\cong (\ZZ/2\ZZ)^3
\]
为命题 \ref{prop:xi-time-part60ab3-window6-boundary-center-residual} 的 boundary parity 中心子群。若只允许保持稳定类型标签且由 \(\Aut(Q_6)\) 实现的强局域 geometric uplift，则其几何可实现的 boundary parity 翻转子群恰为
\[
P_6^{\mathrm{geo}}=\langle\zeta\rangle\cong \ZZ/2\ZZ,
\]
其中 \(\zeta\) 为定理 \ref{thm:xi-time-part53da-boundary-z6-torsor} 中同时翻转三条两点纤维的全局 sheet-flip。更精确地，在坐标识别
\[
P_6\cong (\ZZ/2\ZZ)^3
\]
下有
\[
P_6^{\mathrm{geo}}=\langle(1,1,1)\rangle.
\]
因此
\[
P_6/P_6^{\mathrm{geo}}\cong (\ZZ/2\ZZ)^2.
\]
换言之，局域几何可实现性仅保留三位 boundary parity 的同步对角方向，并精确剥离两个独立的二值自由度。
\end{theorem}
\begin{proof}
由命题 \ref{prop:xi-time-part60ab3-window6-boundary-center-residual}，三条 boundary 二点纤维各自产生一个彼此独立的中心 sheet-flip，故
\[
P_6\cong (\ZZ/2\ZZ)^3.
\]
另一方面，由定理 \ref{thm:xi-time-part53da-boundary-z6-torsor}，\(\widetilde X_6^{\mathrm{bdry}}\) 上存在一个全局对合
\[
\zeta:\widetilde X_6^{\mathrm{bdry}}\to \widetilde X_6^{\mathrm{bdry}},
\]
它在每条两点纤维上都施加唯一非平凡交换，因此在上述三维 parity 坐标中正对应于同步翻转 \((1,1,1)\)。

再由结论 \ref{con:xi-terminal-window6-local-uplift-so10-unique}，在保持稳定类型标签 \(F_6\) 不变且由 \(\Aut(Q_6)\) 实现的强局域 uplift 口径下，除平凡位移外仅允许 \(\delta=\pm 34\)。而在 window-\(6\) 的 boundary uplift 中，同一纤维内两点之差正恒等于 \(34\)（定理 \ref{thm:xi-time-part53da-boundary-z6-torsor} 的设定），故这唯一允许的几何非平凡分支恰好就是 \(\zeta\)。因此
\[
P_6^{\mathrm{geo}}=\langle \zeta\rangle=\langle(1,1,1)\rangle\cong \ZZ/2\ZZ.
\]
最后，\((\ZZ/2\ZZ)^3\) 除以其一维对角子群后得到
\[
(\ZZ/2\ZZ)^3/\langle(1,1,1)\rangle\cong (\ZZ/2\ZZ)^2,
\]
从而 cokernel 精确具有两位二进自由度。证毕。
\end{proof}

\begin{theorem}[window-\(6\) boundary parity 的无协议门跨尺度 \(2\)-挠中心延拓不存在]\label{thm:xi-time-part65e-window6-no-gateless-z2-central-extension}
不存在某个固定紧连通群 \(G\) 与一族仅依赖于分辨率的规则
\[
\mathfrak F_m:X_m^{\mathrm{bdry}}\longrightarrow Z(G)[2]
\qquad(m\in\{6,7,8\}),
\]
其中
\[
Z(G)[2]:=\{z\in Z(G):z^2=e\},
\]
使得同时满足下列三条条件：
\begin{enumerate}
  \item 在 \(m=6\) 时，\(\mathfrak F_6\) 恢复推论 \ref{cor:terminal-delta-parity-rule} 中三条 boundary 纤维所给出的独立 \((\ZZ_2)^3\)-sheet parity；
  \item 该规则对 dyadic uplift 的 boundary 纤维置换保持自然；
  \item 构造过程中不引入任何额外协议门或额外读出寄存器。
\end{enumerate}
\end{theorem}
\begin{proof}
反设存在这样一族规则。条件 (1) 表明：在 \(m=6\) 的 boundary 扇区上，每条二点纤维都被赋予一个非平凡的二值中心选择；这正是推论 \ref{cor:terminal-delta-parity-rule} 中由
\[
\Delta_6^{\mathrm{bdry}}=\{0,34\}
\]
所诱导的 sheet parity。若条件 (2) 与 (3) 也同时成立，则这一二值对象必须沿 dyadic uplift 在更高分辨率的 boundary 纤维上继续以“置换自然且无需外加侧信息”的方式存在。

然而，命题 \ref{prop:terminal-sign-covariant-label-boundary-bifurcation} 已经给出严格判据：这样的符号协变二值对象存在当且仅当纤维基数等于 \(2\)。同一命题并明确指出，在 boundary uplift 口径下有
\[
m=6:\ \#(\mathrm{Fold}^{\mathrm{bin}}_6)^{-1}(w)=2,
\]
\[
m=7,8:\ \#(\mathrm{Fold}^{\mathrm{bin}}_m)^{-1}(w)=3,
\qquad
\Delta_7^{\mathrm{bdry}}=\{0,55,89\},
\qquad
\Delta_8^{\mathrm{bdry}}=\{0,89,144\}.
\]
因此 \(m=7,8\) 的三点 boundary 纤维上根本不存在所需的置换自然二值延拓，这与条件 (2)--(3) 矛盾。故这样的统一规则不存在。

最后，备注 \ref{rem:window6-sheet-parity-nonpersistent} 已指出：若欲把 window-\(6\) 的 \(\ZZ_2\) sheet parity 提升为跨尺度残余选择律，就必须额外施加可审计协议门以重新锁定二层覆盖，或引入额外读出寄存器。并且定理 \ref{thm:xi-boundary-uplift-mixed-radix-register-lower-bound} 进一步量化了这一下限：若要求完整继承 window-\(6\) 的三位 sheet parity，则必有
\[
r_{\mathrm{ext}}(7)\ge 11,
\qquad
r_{\mathrm{ext}}(8)\ge 16.
\]
证毕。
\end{proof}

\noindent
由此，Toeplitz 主子式、离切片 Blaschke 数据、圆盘完成化边界比较与 window-\(6\) boundary uplift 这五条原本分属不同模块的链条，都被进一步压缩为有限局域证书：Toeplitz 可实现性不再需要搜索整个正定锥，而可在对数坐标下完全转写为离散凹锥；Blaschke 原点模不再只是总偏离负载，而在给定最小偏离阈值时直接控制残差维数；RH 的一个充分条件不再需要整列紧集一致逼近，而只需一列逼近边界的圆周裕度不等式；至于 window-\(6\) 的 \((\ZZ_2)^3\) boundary parity，则在局域几何 uplift 口径下仅保留同步对角 \(\ZZ_2\) 方向，而其跨尺度延拓仍必须显式支付新的协议成本。

\paragraph{部分屏幕观测的后验熵、唯一回放阈值、压缩树补全与单面边际律}
在部分边界观测的相对同调核、屏幕图连通分量公式、最小审计补全成本与可见秩函数的拟阵结构已经确立之后，这条主线还可继续压缩为下列五条直接后果。它们分别把 \(\FF_2\)-后验分布写成完全显式的公平比特模型，把唯一回放屏幕的极小阈值识别为全局边界多重图的生成树阈值，把任意部分屏幕的极小精确化补全面集完全刻画为压缩多重图的生成树，把带权精确化编译为最小生成树问题，并把单张观测面的边际信息增益提升为严格的 \(0\text{--}1\) 律。

\begin{theorem}[部分屏幕观测的后验分布完全等价于独立公平比特]\label{thm:xi-time-part65f-screen-posterior-fair-bits}
\leanverified{paper\_xi\_time\_part65f\_screen\_posterior\_fair\_bits}
固定 \(m,n\ge 1\)，并沿用定理 \ref{thm:spg-partial-boundary-screen-kernel-relative-homology}、定理 \ref{thm:spg-screen-kernel-connected-components} 与推论 \ref{cor:spg-partial-boundary-fiber-cardinality-f2} 的记号。将系数域取为 \(\FF_2\)，令
\[
f_S:C_n(\mathcal Q_m;\FF_2)\longrightarrow \FF_2^{S},
\qquad
Y:=f_S(X),
\]
其中 \(X\) 在 \(C_n(\mathcal Q_m;\FF_2)\) 上服从均匀分布。记
\[
c_S:=c(\Gamma_S).
\]
则对每个 \(y\in \operatorname{im}(f_S)\)，纤维 \(f_S^{-1}(y)\) 是维数 \(c_S-1\) 的仿射 \(\FF_2\)-子空间；特别地，\(f_S^{-1}(y)\) 与 \(\FF_2^{\,c_S-1}\) 仿射同构。

若以下 Shannon 熵与互信息均以 \(\log_2\) 为底，则有
\[
H(X\mid Y)=c_S-1,
\]
\[
H(Y)=2^{mn}-c_S+1,
\]
\[
I(X;Y)=2^{mn}-c_S+1.
\]
\end{theorem}
\begin{proof}
顶维链群 \(C_n(\mathcal Q_m;\FF_2)\) 的维数等于顶维胞腔数 \(2^{mn}\)，故
\[
H(X)=\log_2|C_n(\mathcal Q_m;\FF_2)|=2^{mn}.
\]
另一方面，定理 \ref{thm:spg-screen-kernel-connected-components} 的“分量常值”论证在系数域 \(\FF_2\) 上逐字适用，从而得到
\[
\ker f_S\cong \FF_2^{\,c_S-1}.
\]
由于 \(f_S\) 为 \(\FF_2\)-线性映射，每个非空纤维 \(f_S^{-1}(y)\) 都是某个 \(\ker f_S\) 的仿射平移，因此与 \(\FF_2^{\,c_S-1}\) 仿射同构。

又由推论 \ref{cor:spg-partial-boundary-fiber-cardinality-f2}，
\[
|f_S^{-1}(y)|=2^{\dim_{\FF_2}\ker f_S}=2^{c_S-1}
\qquad
(y\in\operatorname{im}(f_S)),
\]
右端与 \(y\) 无关。故在均匀先验下，条件分布 \(X\mid Y=y\) 在纤维上均匀，因而
\[
H(X\mid Y=y)=\log_2|f_S^{-1}(y)|=c_S-1
\]
对所有可达 \(y\) 都成立，从而
\[
H(X\mid Y)=c_S-1.
\]

由链式法则，
\[
H(Y)=H(X)-H(X\mid Y)=2^{mn}-c_S+1.
\]
最后，\(Y=f_S(X)\) 是 \(X\) 的确定函数，因此
\[
I(X;Y)=H(Y)-H(Y\mid X)=H(Y)=2^{mn}-c_S+1.
\]
证毕。
\end{proof}

\begin{theorem}[唯一回放屏幕的极小规模与生成树刻画]\label{thm:xi-time-part65f-exact-screen-spanning-tree-threshold}
沿用定理 \ref{thm:spg-screen-kernel-connected-components} 的边界多重图 \(\Gamma\) 记号：其顶点集为全部 \(n\)-胞腔外加外部顶点 \(\infty\)，并把每个可观测取向 \((n-1)\)-面视为 \(\Gamma\) 的一条边；内部面连接其两侧顶维胞腔，外边界面连接相应胞腔与 \(\infty\)。对任意屏幕 \(S\subset \overrightarrow{\mathcal F_m^{n-1}}\)，记其所诱导的生成子多重图为 \(\Gamma_S\)。

则成立：
\begin{enumerate}
  \item \(f_S\) 单射当且仅当 \(\Gamma_S\) 连通；
  \item 任意使 \(f_S\) 单射的屏幕都满足
        \[
        |S|\ge 2^{mn};
        \]
  \item 等号成立当且仅当 \(\Gamma_S\) 是 \(\Gamma\) 的一棵生成树；
  \item 因而极小唯一回放屏幕的总数恰为 \(\tau(\Gamma)\)，其中 \(\tau(\Gamma)\) 表示 \(\Gamma\) 的生成树数。
\end{enumerate}
\end{theorem}
\begin{proof}
第 \(1\) 条正是推论 \ref{cor:spg-screen-injectivity-and-audit-cost-components} 的内容。

由于 \(\mathcal Q_m\) 的顶维胞腔数为 \(2^{mn}\)，故
\[
|V(\Gamma)|=2^{mn}+1.
\]
任意使 \(f_S\) 单射的屏幕都对应一个连通 spanning sub-multigraph \(\Gamma_S\)。而任意连通图至少有 \(|V|-1\) 条边，因此
\[
|S|=|E(\Gamma_S)|\ge |V(\Gamma)|-1=2^{mn}.
\]
这就得到第 \(2\) 条。

若 \(|S|=2^{mn}\) 且 \(f_S\) 单射，则 \(\Gamma_S\) 连通并且
\[
|E(\Gamma_S)|=|V(\Gamma)|-1,
\]
故 \(\Gamma_S\) 必为生成树。反之，任一生成树都连通且边数恰为 \(|V(\Gamma)|-1=2^{mn}\)，于是对应屏幕必为极小唯一回放屏幕。于是第 \(3\) 条成立。

第 \(4\) 条随即成立：极小唯一回放屏幕与 \(\Gamma\) 的生成树一一对应，其总数即为 \(\tau(\Gamma)\)。证毕。
\end{proof}
\leanverified{paper\_xi\_time\_part65f\_exact\_screen\_spanning\_tree\_threshold}

\begin{theorem}[任意部分屏幕的极小精确化补全面集由压缩图生成树完全刻画]\label{thm:xi-time-part65f-screen-quotient-tree-completion}
沿用定理 \ref{thm:spg-screen-kernel-connected-components} 的记号。给定任意部分屏幕 \(S\subset \overrightarrow{\mathcal F_m^{n-1}}\)，设
\[
\Gamma_S=\bigsqcup_{i=1}^{c}\mathcal C_i,
\qquad
c:=c(\Gamma_S).
\]
将每个连通分量 \(\mathcal C_i\) 缩并为一个顶点，并把 \(\Gamma\setminus E(\Gamma_S)\) 中连接不同分量的边保留下来、删去缩并后产生的回路，得到压缩多重图
\[
\Gamma/\Gamma_S.
\]
以下把每条跨分量补全边与其在 \(\Gamma/\Gamma_S\) 中对应的边仍记为同一符号。
则成立：
\begin{enumerate}
  \item 同一个整数 \(c-1\) 同时等于最小圆维审计补全成本与最少追加观测面数，即
        \[
        \min_{(R,r_S)}\cdim(R)
        =
        c-1
        =
        \min\Bigl\{
        |T|:\ T\subseteq E(\Gamma)\setminus E(\Gamma_S),\ \Gamma_{S\cup T}\text{ 连通}
        \Bigr\};
        \]
  \item 对任意 \(T\subseteq E(\Gamma)\setminus E(\Gamma_S)\) 且 \(|T|=c-1\)，屏幕 \(S\cup T\) 达到唯一回放当且仅当 \(T\) 在压缩图 \(\Gamma/\Gamma_S\) 中的像是一棵生成树；
  \item 因而全部极小精确化补全面集的总数恰为
        \[
        \tau(\Gamma/\Gamma_S).
        \]
\end{enumerate}
\end{theorem}
\begin{proof}
第一条中的左端等式正是推论 \ref{cor:spg-screen-injectivity-and-audit-cost-components}。

现证右端等式。由于 \(\Gamma_S\) 恰有 \(c\) 个连通分量，若要通过继续添加观测面使 \(\Gamma_{S\cup T}\) 连通，则 \(T\) 在压缩图上必须把这 \(c\) 个缩并顶点连成一个连通图。任意 \(c\) 点连通图至少含 \(c-1\) 条边，因此
\[
|T|\ge c-1.
\]
另一方面，压缩图任取一棵生成树，其对应的原图边集 \(T\) 恰含 \(c-1\) 条边，并使 \(\Gamma_{S\cup T}\) 连通，从而上界可达。故右端最小值也等于 \(c-1\)。

再证第 \(2\) 条。若 \(|T|=c-1\) 且 \(\Gamma_{S\cup T}\) 连通，则 \(T\) 在压缩图上的像给出一个 \(c\) 点连通图，且边数恰为 \(c-1\)，故必为生成树。反之，若该像是一棵生成树，则对应的原图边集把 \(c\) 个缩并顶点连通，从而 \(\Gamma_{S\cup T}\) 连通，推论 \ref{cor:spg-screen-injectivity-and-audit-cost-components} 遂给出 \(f_{S\cup T}\) 单射。

第 \(3\) 条随之成立：压缩图的每棵生成树与一个极小精确化补全面集唯一对应，故总数恰为 \(\tau(\Gamma/\Gamma_S)\)。证毕。
\end{proof}
\leanverified{paper\_xi\_time\_part65f\_screen\_quotient\_tree\_completion}

\begin{theorem}[带权精确化可完全编译为压缩图的最小生成树问题]\label{thm:xi-time-part65f-weighted-screen-completion-mst}
在定理 \ref{thm:xi-time-part65f-screen-quotient-tree-completion} 的设定下，对每条补全边赋予非负代价
\[
\mathrm{cost}:E(\Gamma)\setminus E(\Gamma_S)\longrightarrow \RR_{\ge 0}.
\]
对任意补全面集 \(T\subseteq E(\Gamma)\setminus E(\Gamma_S)\)，定义总代价
\[
\mathrm{Cost}(T):=\sum_{e\in T}\mathrm{cost}(e).
\]
则
\[
\min\Bigl\{\mathrm{Cost}(T):\ \Gamma_{S\cup T}\text{ 连通}\Bigr\}
=
\min\Bigl\{\mathrm{Cost}(T):\ T\text{ 在 }\Gamma/\Gamma_S\text{ 中是一棵生成树}\Bigr\}.
\]
特别地：
\begin{enumerate}
  \item 任一压缩图最小生成树对应的补全面集都给出成本最优的精确化方案；
  \item 反之，任一按包含极小的成本最优补全面集，其压缩像必为 \(\Gamma/\Gamma_S\) 的一棵最小生成树。
\end{enumerate}
\end{theorem}
\begin{proof}
任取使 \(\Gamma_{S\cup T}\) 连通的补全面集 \(T\)。其在压缩图上的像是一个连通生成子多重图。取其任意生成树 \(T_{\mathrm{tr}}\subseteq T\)。由于代价非负，
\[
\mathrm{Cost}(T_{\mathrm{tr}})\le \mathrm{Cost}(T).
\]
因此，在全部可行补全面集中最小化总代价，等价于在压缩图的全部生成树中最小化总代价。第一个等式得证。

于是任一压缩图最小生成树都直接给出成本最优补全，这证明了第 \(1\) 条。

若 \(T\) 是按包含极小的成本最优补全面集，而其压缩像含有回路，则删去该回路上一条边后图仍连通；由于边权非负，总代价不会增加。若总代价严格减小，则与成本最优性矛盾；若总代价不变，则与按包含极小矛盾。故其压缩像无回路。既然它仍连通并覆盖全部缩并顶点，便是一棵生成树；再由成本最优性可知它必是最小生成树。第 \(2\) 条成立。证毕。
\end{proof}

\begin{theorem}[单张观测面的边际信息增益具有严格的 \(0\text{--}1\) 律]\label{thm:xi-time-part65f-single-face-zero-one-marginal-law}
\leanverified{paper\_xi\_time\_part65f\_single\_face\_zero\_one\_marginal\_law}
沿用定理 \ref{thm:spg-screen-rank-matroid-supermodularity} 的记号，令
\[
r(S):=\rank_{\QQ}\bigl(\operatorname{im}(f_S\otimes\QQ)\bigr),
\qquad
\delta(S):=\rank_{\QQ}\bigl(\ker(f_S\otimes\QQ)\bigr).
\]
再把任一未观测取向 \((n-1)\)-面与其在 \(\Gamma\) 中对应的边同记为 \(e\in E(\Gamma)\setminus E(\Gamma_S)\)。若 \(e\) 的两个端点分别落在 \(\Gamma_S\) 的连通分量 \(\mathcal C_i,\mathcal C_j\) 中，则有精确公式
\[
\delta(S)-\delta(S\cup\{e\})
=
\begin{cases}
1,& \mathcal C_i\neq \mathcal C_j,\\[4pt]
0,& \mathcal C_i=\mathcal C_j,
\end{cases}
\]
等价地，
\[
r(S\cup\{e\})-r(S)
=
\begin{cases}
1,& \mathcal C_i\neq \mathcal C_j,\\[4pt]
0,& \mathcal C_i=\mathcal C_j.
\end{cases}
\]
\end{theorem}
\begin{proof}
由定理 \ref{thm:spg-screen-kernel-connected-components} 与推论 \ref{cor:spg-screen-injectivity-and-audit-cost-components}，
\[
\delta(S)=c(\Gamma_S)-1.
\]
若 \(e\) 的两个端点位于同一连通分量中，则向 \(\Gamma_S\) 添加 \(e\) 不改变连通分量数，因此
\[
c(\Gamma_{S\cup\{e\}})=c(\Gamma_S),
\qquad
\delta(S\cup\{e\})=\delta(S).
\]

若 \(e\) 连接两个不同分量，则添加 \(e\) 恰好把这两个分量合并，从而
\[
c(\Gamma_{S\cup\{e\}})=c(\Gamma_S)-1,
\]
进而
\[
\delta(S\cup\{e\})=c(\Gamma_{S\cup\{e\}})-1=\delta(S)-1.
\]
这就得到第一条 \(0\text{--}1\) 公式。

再由定理 \ref{thm:spg-screen-rank-matroid-supermodularity} 的秩亏恒等式
\[
\delta(S)=\rank_{\QQ}C_n(\mathcal Q_m;\QQ)-r(S)
=
2^{mn}-r(S),
\]
立刻推出
\[
r(S\cup\{e\})-r(S)=\delta(S)-\delta(S\cup\{e\}),
\]
从而得到第二条公式。证毕。
\end{proof}

\noindent
综上，部分边界观测这条链已经在概率、图论、优化与拟阵四个层面同时闭合：纤维常数不再只是一个计数事实，而是完整的后验公平比特模型；唯一回放屏幕的极小阈值不再只是 \(2^{mn}\) 的抽象秩值，而是全局边界多重图的生成树阈值；任意部分屏幕的极小精确化不再只是“还差 \(c(\Gamma_S)-1\) 张面”，而是由压缩图的生成树与最小生成树完全组织；与此同时，单张新观测面所带来的信息秩增量也不再只是超模意义下的平均改善，而是具有严格可判定的 \(0\text{--}1\) 局域律。

\paragraph{局部化奇异环的 \(S\)-根单值化、逆极限相位账本、Fold 相位游走与 Dirichlet 轴交换阻断}
在有限 prime-register solenoid 的 Kronecker 读出、有理玫瑰/Lissajous 的几何提升、局部化相位采样的素数恢复与 Fold 谱的乘积分解已经确立之后，还可把 singular ring 侧的根提升、分母一致性、干涉能量的概率读出以及 Dirichlet 频率上的换轴障碍继续压缩为下列一组闭式后果。其核心是：\(\Sigma_S=\widehat{\ZZ[S^{-1}]}\) 不仅是圆周商上的 prime-register 扩张，而且恰是全部 \(S\)-光滑有理幂的普适单值化对象；一旦把这一普适对象与 Fold 的乘积谱和 Dirichlet 频率的线性无关性并置，分母账本、概率干涉与换轴阻断三条链便可统一落到同一相位审计框架之内。

\begin{theorem}[奇异环上的 \(S\)-根提升与单位圆有理幂的普适单值化]\label{thm:xi-time-part65g-solenoid-s-root-universal-single-valuedization}
沿用定理 \ref{thm:cdim-star-z1s-dual-extension} 的记号。设 \(S\subset\mathbb P\) 为有限素数集，并记
\[
G_S:=\ZZ[S^{-1}],
\qquad
\Sigma_S:=\widehat{G_S},
\qquad
\chi_q(x):=x(q)
\quad (q\in G_S,\ x\in \Sigma_S).
\]
令
\[
\pi:=\chi_1:\Sigma_S\to \TT.
\]
则 \(\pi\) 为满同态。对每个 \(S\)-光滑正整数 \(d\)，定义
\[
\rho_d:=\chi_{1/d}:\Sigma_S\to \TT.
\]
则对任意 \(x\in \Sigma_S\) 都有
\[
(\rho_d(x))^d=\pi(x),
\]
并且若 \(q=m/d\in G_S\) 为既约分式且 \(d\) 为 \(S\)-光滑，则
\[
\chi_q(x)=\rho_d(x)^m.
\]

更进一步，\(\Sigma_S\) 在如下意义下具有普适性：若 \(K\) 为紧交换群，存在满同态
\[
\pi_K:K\twoheadrightarrow \TT
\]
与一族连续同态
\[
\rho_d^K:K\to \TT
\]
（\(d\) 遍历全部 \(S\)-光滑正整数）满足
\[
(\rho_d^K)^d=\pi_K,
\qquad
(\rho_{d'}^K)^{d'/d}=\rho_d^K
\quad (d\mid d'),
\]
则存在唯一连续满同态
\[
F:K\twoheadrightarrow \Sigma_S
\]
使得
\[
\pi_K=\pi\circ F,
\qquad
\rho_d^K=\rho_d\circ F
\]
对全部 \(S\)-光滑 \(d\) 都成立。
\end{theorem}
\begin{proof}
由字符乘法法则，\(\chi_{q_1}\chi_{q_2}=\chi_{q_1+q_2}\)。因此对每个 \(S\)-光滑 \(d\) 都有
\[
(\rho_d)^d
=
\chi_{1/d}^d
=
\chi_{d\cdot (1/d)}
=
\chi_1
=
\pi,
\]
而对 \(q=m/d\in G_S\) 则有
\[
\chi_q
=
\chi_{m\cdot(1/d)}
=
\chi_{1/d}^{\,m}
=
\rho_d^{\,m}.
\]

由定理 \ref{thm:cdim-star-z1s-dual-extension}，\(\Sigma_S\) 连通。又 \(\chi_1\) 非平凡，故其像是 \(\TT\) 的非平凡闭连通子群，只能等于 \(\TT\)。于是 \(\pi\) 满射。

对普适性，令 \(\gamma_d\in K^\wedge\) 为 \(\rho_d^K\) 的对偶元素。则
\[
d\gamma_d=\gamma_1,
\qquad
\frac{d'}{d}\gamma_{d'}=\gamma_d
\quad (d\mid d').
\]
定义映射
\[
\iota:G_S\longrightarrow K^\wedge,
\qquad
\iota\!\left(\frac{m}{d}\right):=m\gamma_d.
\]
需证其良定义。若 \(m/d=m'/d'\)，取 \(S\)-光滑公倍数 \(\ell\) 使 \(d\mid \ell\)、\(d'\mid \ell\)。则
\[
m\gamma_d
=
m\frac{\ell}{d}\gamma_\ell
=
\frac{m\ell}{d}\gamma_\ell
=
\frac{m'\ell}{d'}\gamma_\ell
=
m'\frac{\ell}{d'}\gamma_\ell
=
m'\gamma_{d'},
\]
故 \(\iota\) 良定义。并且
\[
\iota\!\left(\frac{m}{d}+\frac{m'}{d'}\right)
=
\iota\!\left(\frac{m\ell/d+m'\ell/d'}{\ell}\right)
=
\left(\frac{m\ell}{d}+\frac{m'\ell}{d'}\right)\gamma_\ell
=
m\gamma_d+m'\gamma_{d'}
\]
说明 \(\iota\) 为群同态。

再证 \(\iota\) 单射。若 \(\iota(m/d)=0\)，则 \(m\gamma_d=0\)，等价于
\[
(\rho_d^K)^m=1.
\]
而 \(\rho_d^K(K)\) 是 \(\TT\) 的闭子群，且其 \(d\)-次幂像等于 \(\pi_K(K)=\TT\)。因此 \(\rho_d^K(K)\) 不可能是有限循环群，只能等于 \(\TT\)。于是 \(z^m=1\) 对所有 \(z\in\TT\) 都成立，故 \(m=0\)，从而 \(m/d=0\)。所以 \(\iota\) 单射。

Pontryagin 对偶化后得到连续满同态
\[
F:=\widehat{\iota}:K\twoheadrightarrow \Sigma_S.
\]
对任意 \(S\)-光滑 \(d\)，\(\rho_d\circ F\) 的对偶元素正是
\[
\iota(1/d)=\gamma_d,
\]
故 \(\rho_d\circ F=\rho_d^K\)。同理 \(\pi\circ F=\pi_K\)。唯一性则来自
\[
\widehat{\Sigma_S}\cong G_S
\]
由全部元素 \(m/d\) 生成，而 \(\chi_{m/d}=\rho_d^m\) 已完全由上述相容关系决定。证毕。
\end{proof}

\begin{corollary}[有理玫瑰与 Lissajous 在奇异环上的同态提升]\label{cor:xi-time-part65g-rose-lissajous-singular-ring-homomorphic-lift}
沿用定理 \ref{thm:xi-time-part65g-solenoid-s-root-universal-single-valuedization} 与定理 \ref{thm:xi-time-part62d-solenoid-kronecker-rose-lissajous-lift} 的记号，记
\[
\iota_S:=\kappa_S:\RR\longrightarrow \Sigma_S,
\qquad
\iota_S(t)(q):=e^{2\pi iqt}
\quad (q\in G_S).
\]
取 \(S\)-光滑正整数 \(d\) 与整数 \(n\) 满足 \(\gcd(n,d)=1\)，定义
\[
\mathcal R_{n,d}(x):=\frac12\Bigl(\rho_d(x)^{d+n}+\rho_d(x)^{d-n}\Bigr).
\]
则对任意 \(\theta\in\RR\) 都有
\[
\mathcal R_{n,d}\bigl(\iota_S(\theta/2\pi)\bigr)
=
e^{i\theta}\cos\!\Bigl(\frac{n}{d}\theta\Bigr).
\]
因此有理玫瑰轨道
\[
\theta\longmapsto e^{i\theta}\cos\!\Bigl(\frac{n}{d}\theta\Bigr)
\]
恰是奇异环字符读出的单值提升。

若 \(a=a_0/d\)、\(b=b_0/d\in G_S\) 且 \(a_0,b_0\in\ZZ\)，定义
\[
\mathcal L_{a,b,\delta}(x):=
\Bigl(
\Re\bigl(e^{i\delta}\rho_d(x)^{a_0}\bigr),\
\Re\bigl(\rho_d(x)^{b_0}\bigr)
\Bigr).
\]
则对任意 \(t\in\RR\) 都有
\[
\mathcal L_{a,b,\delta}\bigl(\iota_S(t/2\pi)\bigr)
=
\bigl(\cos(at+\delta),\cos(bt)\bigr).
\]
故一切分母仅由 \(S\) 中素数支撑的有理玫瑰与有理 Lissajous 轨道，都可在 \(\Sigma_S\) 上以单值同态方式实现。
\end{corollary}
\begin{proof}
由定理 \ref{thm:xi-time-part62d-solenoid-kronecker-rose-lissajous-lift}，对任意 \(q\in G_S\) 与 \(u\in\RR\) 都有
\[
\chi_q\bigl(\iota_S(u/2\pi)\bigr)=e^{iqu}.
\]
特别地，
\[
\rho_d\bigl(\iota_S(u/2\pi)\bigr)=\chi_{1/d}\bigl(\iota_S(u/2\pi)\bigr)=e^{iu/d}.
\]
于是取 \(u=\theta\) 得
\[
\mathcal R_{n,d}\bigl(\iota_S(\theta/2\pi)\bigr)
=
\frac12\Bigl(
e^{i(d+n)\theta/d}+e^{i(d-n)\theta/d}
\Bigr)
=
e^{i\theta}\cos\!\Bigl(\frac{n}{d}\theta\Bigr).
\]
这证明了玫瑰公式。

同理，取 \(u=t\) 得
\[
\rho_d\bigl(\iota_S(t/2\pi)\bigr)^{a_0}=e^{ia_0t/d}=e^{iat},
\qquad
\rho_d\bigl(\iota_S(t/2\pi)\bigr)^{b_0}=e^{ib_0t/d}=e^{ibt},
\]
故
\[
\mathcal L_{a,b,\delta}\bigl(\iota_S(t/2\pi)\bigr)
=
\bigl(\Re(e^{i\delta}e^{iat}),\Re(e^{ibt})\bigr)
=
\bigl(\cos(at+\delta),\cos(bt)\bigr).
\]
证毕。
\end{proof}

\begin{proposition}[局部化奇异环的逆极限相位账本刻画]\label{prop:xi-time-part65g-solenoid-inverse-limit-phase-ledger}
沿用定理 \ref{thm:xi-time-part65g-solenoid-s-root-universal-single-valuedization} 的记号。记
\[
\mathcal D_S:=\left\{
d\in\NN_{\ge 1}:\ \text{\(d\) 的全部素因子都属于 }S
\right\},
\]
并按整除关系赋予其有向结构。对 \(d\mid d'\) 定义连接映射
\[
\varphi_{d',d}:\TT\to\TT,
\qquad
\varphi_{d',d}(z):=z^{d'/d}.
\]
则映射
\[
\Sigma_S\longrightarrow \varprojlim_{d\in\mathcal D_S}\bigl(\TT,\varphi_{d',d}\bigr),
\qquad
x\longmapsto \bigl(\rho_d(x)\bigr)_{d\in\mathcal D_S}
\]
是拓扑同构。

等价地，一个相位账本 \((z_d)_{d\in\mathcal D_S}\in \TT^{\mathcal D_S}\) 来自某个 \(x\in\Sigma_S\) 当且仅当它满足一致性
\[
z_d=z_{d'}^{\,d'/d}
\qquad (d\mid d').
\]
在该情形下，相应字符满足
\[
x\!\left(\frac{m}{d}\right)=z_d^{\,m}
\qquad
\left(\frac{m}{d}\in G_S\right).
\]
\end{proposition}
\begin{proof}
有直极限分解
\[
G_S
=
\bigcup_{d\in\mathcal D_S}\frac{1}{d}\ZZ
=
\varinjlim_{d\in\mathcal D_S}\frac{1}{d}\ZZ,
\]
其中当 \(d\mid d'\) 时，\((1/d)\ZZ\hookrightarrow (1/d')\ZZ\) 为自然包含。Pontryagin 对偶把离散群直极限送为紧群逆极限，因此
\[
\Sigma_S
\cong
\varprojlim_{d\in\mathcal D_S}\widehat{(1/d)\ZZ}.
\]
把每个 \(\widehat{(1/d)\ZZ}\) 识别为 \(\TT\) 后，自然包含的对偶恰为
\[
z\longmapsto z^{d'/d},
\]
故得到所示逆系统。

显式地，若 \(x\in\Sigma_S\)，记
\[
z_d:=x(1/d)=\rho_d(x).
\]
当 \(d\mid d'\) 时，
\[
z_{d'}^{\,d'/d}
=
x(1/d')^{d'/d}
=
x\!\left(\frac{d'}{d}\cdot \frac{1}{d'}\right)
=
x(1/d)
=
z_d,
\]
故 \((z_d)\) 一致。

反之，给定一致族 \((z_d)_{d\in\mathcal D_S}\)，对任意 \(q=m/d\in G_S\) 定义
\[
x(q):=z_d^{\,m}.
\]
若 \(m/d=m'/d'\)，取 \(\ell\in\mathcal D_S\) 使 \(d\mid \ell\)、\(d'\mid \ell\)。则由一致性
\[
z_d^{\,m}
=
\left(z_\ell^{\,\ell/d}\right)^m
=
z_\ell^{\,m\ell/d}
=
z_\ell^{\,m'\ell/d'}
=
\left(z_\ell^{\,\ell/d'}\right)^{m'}
=
z_{d'}^{\,m'},
\]
故定义与表示无关。并且
\[
x(q+q')=x(q)x(q')
\]
由同分母计算立即可得，所以 \(x\in\operatorname{Hom}(G_S,\TT)=\Sigma_S\)。显然这两个构造互为逆映射。紧 Hausdorff 范畴中的连续双射即为同胚，故结论成立。证毕。
\end{proof}

\begin{theorem}[有限频率集的根闭包最小载体与 prime-register 共同核的素数恢复]\label{thm:xi-time-part65g-frequency-prime-support-kernel-recovery}
\leanverified{paper\_xi\_time\_part65g\_frequency\_prime\_support\_kernel\_recovery}
设 \(Q=\{q_1,\dots,q_r\}\subset\QQ^\times\) 为有限频率族，并记 \(S(Q)\subset\mathbb P\) 为其既约分母中出现的全部素数。再记
\[
\Sigma_{S(Q)}:=\widehat{\ZZ[S(Q)^{-1}]},
\qquad
K_{S(Q)}:=\prod_{p\in S(Q)}\ZZ_p.
\]
则成立：
\begin{enumerate}
  \item 若一个紧交换群 \(K\) 配备满同态 \(\pi_K:K\twoheadrightarrow \TT\) 及其对全部 \(S(Q)\)-光滑正整数的兼容根提升 \(\{\rho_d^K\}\)，并且对每个 \(q_j=m_j/d_j\)（既约）都以
  \[
  (\rho_{d_j}^K)^{m_j}
  \]
  作为对应频率 \(q_j\) 的读出，
  则存在唯一连续满同态
  \[
  F:K\twoheadrightarrow \Sigma_{S(Q)}
  \]
  使得全部 \(\rho_d^K\) 以及每个频率读出 \((\rho_{d_j}^K)^{m_j}\) 都由 \(\Sigma_{S(Q)}\) 上的标准字符经 \(F\) 拉回。换言之，一旦把 \(Q\) 的实现口径要求为对全部 \(S(Q)\)-光滑分母闭合的根一致读出，则 \(\Sigma_{S(Q)}\) 是普适最小载体。
  \item 对 \(\Sigma_{S(Q)}\) 上的角色族 \(\chi_{q_1},\dots,\chi_{q_r}\)，定义
  \[
  M_p:=\max_{1\le j\le r}\max\{0,-v_p(q_j)\}
  \qquad (p\in S(Q)).
  \]
  则有开子群包含
  \[
  \prod_{p\in S(Q)}p^{M_p}\ZZ_p
  \subseteq
  \bigcap_{j=1}^{r}\ker\bigl(\chi_{q_j}|_{K_{S(Q)}}\bigr),
  \]
  并且对每个素数 \(p\) 都有
  \[
  p\in S(Q)
  \iff
  M_p\ge 1
  \iff
  \bigcap_{j=1}^{r}\ker\bigl(\chi_{q_j}|_{K_{S(Q)}}\bigr)
  \text{ 含有非平凡开 }p\text{-进子群}.
  \]
  因而分母素数支撑可由共同 prime-register 核完全恢复。
\end{enumerate}
\end{theorem}
\begin{proof}
第 \(1\) 条是定理 \ref{thm:xi-time-part65g-solenoid-s-root-universal-single-valuedization} 在 \(S=S(Q)\) 时的直接应用；对每个 \(q_j=m_j/d_j\)，由该定理中的恒等式
\[
\chi_{q_j}=\rho_{d_j}^{\,m_j}
\]
可知 \(Q\) 的全部频率读出都被同一普适对象 \(\Sigma_{S(Q)}\) 吸收。

对第 \(2\) 条，把定理 \ref{thm:xi-time-part69-finite-character-prime-ledger-finite-depth} 应用于角色族
\[
\chi_{q_1},\dots,\chi_{q_r}\in \widehat{\Sigma_{S(Q)}}\cong \ZZ[S(Q)^{-1}].
\]
该定理给出开子群
\[
U(\chi_{q_1},\dots,\chi_{q_r})
=
\prod_{p\in S(Q)}p^{M_p}\ZZ_p
\]
并满足
\[
U(\chi_{q_1},\dots,\chi_{q_r})
\subseteq
\bigcap_{j=1}^{r}\ker\bigl(\chi_{q_j}|_{K_{S(Q)}}\bigr).
\]
另一方面，按 \(M_p\) 的定义，
\[
M_p\ge 1
\iff
\exists\,j\ \text{使}\ v_p(q_j)<0
\iff
p\ \text{出现在某个 }q_j\text{ 的既约分母中}
\iff
p\in S(Q).
\]
若 \(p\in S(Q)\)，则 \(p^{M_p}\ZZ_p\) 是 \(\ZZ_p\) 的非平凡开子群，因而共同核中含有非平凡开 \(p\)-进子群；反之若 \(p\notin S(Q)\)，则 \(K_{S(Q)}\) 本身并无 \(p\)-主方向，更不存在此类 \(p\)-进开子群。故所述等价成立。证毕。
\end{proof}

\begin{theorem}[Fold 多重度 Fourier 能量的三值相位游走表示]\label{thm:xi-time-part65g-fold-energy-threevalued-phase-walk}
沿用定理 \ref{thm:fold-spectrum-factorization} 的记号。对固定 \(m\ge 1\)，记
\[
M:=F_{m+2},
\qquad
\widehat d_m(k)=\prod_{j=1}^{m}\left(1+e^{-2\pi i kF_{j+1}/M}\right).
\]
令 \(\tau_1,\dots,\tau_m\) 独立同分布，并满足
\[
\PP(\tau_j=0)=\frac12,
\qquad
\PP(\tau_j=\pm 1)=\frac14.
\]
定义随机相位游走
\[
X_m:=
\sum_{j=1}^{m}\tau_j\frac{F_{j+1}}{F_{m+2}}
\in
\RR/\ZZ.
\]
则对任意整数 \(k\) 都有
\[
\EE\bigl[e^{2\pi i kX_m}\bigr]
=
\prod_{j=1}^{m}
\cos^2\!\Bigl(\pi k\frac{F_{j+1}}{F_{m+2}}\Bigr)
=
\frac{|\widehat d_m(k)|^2}{2^{2m}}.
\]
特别地，归一化 Fourier 能量剖面既是一个有限 Riesz 乘积，也是一个显式三值相位游走的特征函数。
\end{theorem}
\begin{proof}
由独立性，
\[
\EE\bigl[e^{2\pi i kX_m}\bigr]
=
\prod_{j=1}^{m}
\EE\!\left[
e^{2\pi i k\tau_jF_{j+1}/F_{m+2}}
\right].
\]
对任意实数 \(\theta\)，单步特征函数满足
\[
\EE[e^{i\theta\tau_j}]
=
\frac12+\frac14e^{i\theta}+\frac14e^{-i\theta}
=
\frac12(1+\cos\theta)
=
\cos^2(\theta/2).
\]
取
\[
\theta=2\pi k\frac{F_{j+1}}{F_{m+2}}
\]
便得到
\[
\EE\!\left[
e^{2\pi i k\tau_jF_{j+1}/F_{m+2}}
\right]
=
\cos^2\!\Bigl(\pi k\frac{F_{j+1}}{F_{m+2}}\Bigr).
\]
于是
\[
\EE\bigl[e^{2\pi i kX_m}\bigr]
=
\prod_{j=1}^{m}
\cos^2\!\Bigl(\pi k\frac{F_{j+1}}{F_{m+2}}\Bigr).
\]

另一方面，由定理 \ref{thm:fold-spectrum-factorization}，
\[
\frac{|\widehat d_m(k)|^2}{2^{2m}}
=
\prod_{j=1}^{m}
\frac{|1+e^{-2\pi i kF_{j+1}/F_{m+2}}|^2}{4}
=
\prod_{j=1}^{m}
\cos^2\!\Bigl(\pi k\frac{F_{j+1}}{F_{m+2}}\Bigr),
\]
因为
\[
|1+e^{-i\theta}|^2=4\cos^2(\theta/2).
\]
两式合并即得结论。证毕。
\end{proof}

\begin{theorem}[Dirichlet 频率上单一时间平移的实虚轴交换阻断]\label{thm:xi-time-part65g-dirichlet-axis-swap-obstruction}
设 \(A\subset\NN\) 为有限集，\(a_n\in\RR\) 为实系数，并定义 Dirichlet 多项式
\[
D(t):=\sum_{n\in A}a_n n^{-it}
=
\sum_{n\in A}a_n e^{-it\log n}.
\]
若存在 \(t_0\in\RR\) 使得对一切 \(t\in\RR\) 都有
\[
\bigl(\Re D(t+t_0),\Im D(t+t_0)\bigr)
=
\bigl(\Im D(t),\Re D(t)\bigr),
\]
则必有
\[
D\equiv 0.
\]
特别地，任何非零实系数 Dirichlet 多项式都不可能通过单一时间平移实现实轴与虚轴的全局交换。
\end{theorem}
\begin{proof}
记
\[
D(t)=x(t)+iy(t).
\]
题设等价于
\[
x(t+t_0)+iy(t+t_0)=y(t)+ix(t)=i\overline{D(t)}.
\]
由于所有 \(a_n\) 都是实数，
\[
\overline{D(t)}=D(-t),
\]
故得到函数恒等式
\[
D(t+t_0)=iD(-t)
\qquad (\forall t\in\RR).
\]
展开两边可得
\[
\sum_{n\in A}a_n e^{-it_0\log n}e^{-it\log n}
=
i\sum_{n\in A}a_n e^{it\log n}.
\]

指数函数族 \(\{e^{it\lambda}\}_{\lambda\in\Lambda}\) 在 \(\lambda\) 互异时于 \(\RR\) 上线性无关。对每个 \(n>1\)，左侧只出现频率 \(-\log n<0\)，右侧只出现频率 \(+\log n>0\)；因此这些频率的系数必须全部为零，从而
\[
a_n=0
\qquad (n>1).
\]
剩下唯一可能的频率是 \(\log 1=0\)。比较常数项得
\[
a_1=ia_1,
\]
故 \(a_1=0\)。于是全部系数皆为零，即 \(D\equiv 0\)。证毕。
\end{proof}

由此，局部化奇异环一侧的 \(S\)-光滑根提升、逆极限相位账本、prime-register 共同核恢复与 Fold 能量的概率化读出便形成了一条内部闭合的单值化链；与此同时，Dirichlet 频率侧试图以单一时间平移完成的全局换轴则被指数频率的线性无关性严格排除。于是“根单值化可行、相位账本可逆、干涉能量可概率化、单轴换位不可行”四个层面被压缩为同一条可审计的 singular-ring 主线。
\paragraph{压力缺口的乘法射线势、escort 单比特实验与 Hellinger--Chernoff 母核}
在整数压力轴的冻结斜率恢复、escort--Rényi 熵率的压力差公式、末位统计的渐近充分化、escort 熵常数的闭式化以及 Fisher--Rao 几何已经确立之后，还可继续压缩出下列六条彼此闭合的后果。它们分别把乘法细化中的压力损失写成一条精确射线势，把可除偏序上的斜率单调与整倍数冻结传播压缩为纯序结构命题，把有界温度 escort 实验族在 Blackwell 意义下指数塌缩为单比特 Bernoulli 家族，并把全阶 Rényi 常数层、Shannon 温度耗散以及二温度判别几何统一到同一个 logistic 母核之下。

\begin{theorem}[压力缺口是乘法熵亏损的精确射线势]\label{thm:xi-time-part65h-pressure-gap-ray-potential}
沿用定理 \ref{thm:xi-time-part65c-escort-renyi-two-parameter-freezing-line} 与定理 \ref{thm:xi-time-part57b-pressure-spectrum-slope-intercept-recovery} 的记号。设整数压力极限
\[
P_a:=\lim_{m\to\infty}\frac1m\log S_a(m)
\qquad (a\ge 1)
\]
都存在，并记
\[
\alpha_*:=\lim_{q\to\infty}\frac{P_q}{q}.
\]
对每个整数 \(a\ge 1\) 记 escort 分布
\[
\pi_{m,a}(x):=\frac{d_m(x)^a}{S_a(m)}.
\]
再对整数 \(a\ge 1\)、\(b\ge 2\) 定义
\[
\delta(a):=\frac{P_a}{a}-\alpha_*,
\qquad
H_{\mathrm{loss}}(a,b):=bP_a-P_{ab}.
\]
再记 escort--Rényi 熵率
\[
h_b^{\mathrm{esc}}(a):=
\lim_{m\to\infty}\frac1m H_b(\pi_{m,a}).
\]
其存在性由定理 \ref{thm:xi-time-part65c-escort-renyi-two-parameter-freezing-line} 保证。则对任意整数 \(a\ge 1\)、\(b\ge 2\) 都有严格恒等式
\[
\delta(a)-\delta(ab)=\frac{H_{\mathrm{loss}}(a,b)}{ab}.
\]
并且
\[
\delta(a)-\delta(ab)=\frac{b-1}{ab}\,h_b^{\mathrm{esc}}(a)\ge 0.
\]
更进一步，沿任意乘法射线都有绝对收敛展开
\[
\delta(a)=\sum_{k=0}^{\infty}\frac{H_{\mathrm{loss}}(ab^k,b)}{ab^{k+1}}.
\]
\leanverified{paper\_xi\_time\_part65h\_pressure\_gap\_ray\_potential}
\end{theorem}
\begin{proof}
第一式只是代数恒等式：
\[
\delta(a)-\delta(ab)
=
\left(\frac{P_a}{a}-\alpha_*\right)
-\left(\frac{P_{ab}}{ab}-\alpha_*\right)
=
\frac{bP_a-P_{ab}}{ab}
=
\frac{H_{\mathrm{loss}}(a,b)}{ab}.
\]

另一方面，由定理 \ref{thm:xi-time-part65c-escort-renyi-two-parameter-freezing-line}，
\[
h_b^{\mathrm{esc}}(a)
=
\frac{bP_a-P_{ab}}{b-1}
=
\frac{H_{\mathrm{loss}}(a,b)}{b-1}.
\]
代回即得
\[
\delta(a)-\delta(ab)=\frac{b-1}{ab}\,h_b^{\mathrm{esc}}(a)\ge 0,
\]
其中非负性来自 Rényi 熵率本身非负。

最后，对任意 \(N\ge 1\) 作有限望远镜求和：
\[
\delta(a)-\delta(ab^N)
=
\sum_{k=0}^{N-1}\bigl(\delta(ab^k)-\delta(ab^{k+1})\bigr)
=
\sum_{k=0}^{N-1}\frac{H_{\mathrm{loss}}(ab^k,b)}{ab^{k+1}}.
\]
由定理 \ref{thm:xi-time-part57b-pressure-spectrum-slope-intercept-recovery}，
\[
\frac{P_q}{q}\longrightarrow \alpha_*
\qquad (q\to\infty),
\]
故当 \(N\to\infty\) 时，
\[
\delta(ab^N)=\frac{P_{ab^N}}{ab^N}-\alpha_*\longrightarrow 0.
\]
于是得到
\[
\delta(a)=\sum_{k=0}^{\infty}\frac{H_{\mathrm{loss}}(ab^k,b)}{ab^{k+1}}.
\]
又由于各项都可写成
\[
\frac{H_{\mathrm{loss}}(ab^k,b)}{ab^{k+1}}
=
\frac{b-1}{ab^{k+1}}\,h_b^{\mathrm{esc}}(ab^k)\ge 0,
\]
该级数因而绝对收敛。证毕。
\end{proof}

\begin{corollary}[可除偏序上的斜率单调、几何逼近与整倍数冻结传播]\label{cor:xi-time-part65h-divisibility-slope-freezing-propagation}
\leanverified{paper\_xi\_time\_part65h\_divisibility\_slope\_freezing\_propagation}
沿用定理 \ref{thm:xi-time-part65h-pressure-gap-ray-potential} 的记号。则成立：
\begin{enumerate}
  \item 若 \(a\mid c\)，则
  \[
  \frac{P_c}{c}\le \frac{P_a}{a}.
  \]
  \item 对任意整数 \(a\ge 1\)、\(b\ge 2\)、\(N\ge 0\)，都有
  \[
  0\le \frac{P_{ab^N}}{ab^N}-\alpha_*
  \le \frac{\log\varphi}{ab^N}.
  \]
  \item 若对某个整数 \(a_0\ge 1\) 恰有
  \[
  P_{a_0}=a_0\alpha_*,
  \]
  则对每个 \(n\ge 1\) 都有
  \[
  P_{a_0n}=a_0n\,\alpha_*.
  \]
\end{enumerate}
\end{corollary}
\begin{proof}
若 \(c=a\)，结论平凡。若 \(c>a\) 且 \(a\mid c\)，写 \(c=ab\) 且 \(b\ge 2\)。由定理 \ref{thm:xi-time-part65h-pressure-gap-ray-potential}，
\[
\frac{P_a}{a}-\frac{P_c}{c}
=
\delta(a)-\delta(ab)
=
\frac{b-1}{ab}\,h_b^{\mathrm{esc}}(a)\ge 0,
\]
故第 \(1\) 条成立。

对第 \(2\) 条，令 \(n:=ab^N\)。由定义，
\[
M_m^n\le S_n(m)\le |X_m|\,M_m^n.
\]
两边取对数后除以 \(mn\)，再令 \(m\to\infty\)，得到
\[
\alpha_*\le \frac{P_n}{n}\le \alpha_*+\frac{1}{n}\lim_{m\to\infty}\frac1m\log|X_m|.
\]
而 \(|X_m|=F_{m+2}\)，故
\[
\lim_{m\to\infty}\frac1m\log|X_m|=\log\varphi.
\]
于是
\[
0\le \frac{P_n}{n}-\alpha_*\le \frac{\log\varphi}{n}.
\]
代回 \(n=ab^N\) 即得所示估计。

最后，若 \(P_{a_0}=a_0\alpha_*\)，则对任意 \(n\ge 1\)，由第 \(1\) 条与第 \(2\) 条分别得到
\[
\alpha_*\le \frac{P_{a_0n}}{a_0n}\le \frac{P_{a_0}}{a_0}=\alpha_*.
\]
故必有
\[
\frac{P_{a_0n}}{a_0n}=\alpha_*,
\]
即
\[
P_{a_0n}=a_0n\,\alpha_*.
\]
证毕。
\end{proof}

\begin{theorem}[有界温度 escort 实验族与单比特 Bernoulli 实验的 Blackwell 指数等价]\label{thm:xi-time-part65h-escort-blackwell-onebit-equivalence}
沿用定理 \ref{thm:killo-fold-bin-escort-onebit-fisher} 的记号。固定 \(Q<\infty\)，定义统计实验族
\[
\mathcal E_m^{(Q)}:=\{\pi_{m,q}:0\le q\le Q\},
\]
以及极限单比特实验族
\[
\mathcal B^{(Q)}:=\{\beta_q:0\le q\le Q\},
\qquad
\beta_q:=\mathrm{Bernoulli}(\theta(q)),
\qquad
\theta(q):=\frac{1}{1+\varphi^{q+1}}.
\]
再定义
\[
\Delta_Q(m):=
\inf_{K,L}\sup_{0\le q\le Q}
\max\Bigl\{
D_{\mathrm{TV}}(K_{\#}\pi_{m,q},\beta_q),\
D_{\mathrm{TV}}(L_{\#}\beta_q,\pi_{m,q})
\Bigr\},
\]
其中 \(K:X_m\rightsquigarrow\{0,1\}\) 与 \(L:\{0,1\}\rightsquigarrow X_m\) 取遍全部 Markov 核。则存在常数 \(C_Q>0\) 使得
\[
\Delta_Q(m)\le C_Q\Bigl(\frac{\varphi}{2}\Bigr)^m.
\]
因此，对整个有界温度区间而言，escort 统计实验在 Blackwell 意义下以指数速度塌缩为由末位比特承载的一维 logistic Bernoulli 家族。
\end{theorem}
\begin{proof}
记
\[
A_i:=\{w\in X_m:w_m=i\},
\qquad
U_{m,i}:=\mathrm{Unif}(A_i)
\qquad (i=0,1).
\]
定义前向核
\[
K_m(w):=w_m\in\{0,1\},
\]
以及反向核
\[
L_m(i,\cdot):=U_{m,i}.
\]
由定理 \ref{thm:killo-fold-bin-escort-onebit-fisher}，存在两态 Gibbs 近似
\[
\nu_{m,q}(w):=\frac{u_m(w)\varphi^{-qw_m}}{Z_{m,q}},
\qquad
Z_{m,q}:=\frac{F_{m+1}+\varphi^{-q}F_m}{F_{m+2}},
\]
并且
\[
\sup_{0\le q\le Q}\|\pi_{m,q}-\nu_{m,q}\|_{\mathrm{TV}}
=
O_Q\!\Bigl(\Bigl(\frac{\varphi}{2}\Bigr)^m\Bigr).
\]
由于 \(\nu_{m,q}\) 在每个块 \(A_i\) 上严格均匀，可把它写成
\[
\nu_{m,q}=(1-\theta_m(q))U_{m,0}+\theta_m(q)U_{m,1},
\]
其中
\[
\theta_m(q):=\nu_{m,q}(A_1)
=
\frac{\varphi^{-q}F_m}{F_{m+1}+\varphi^{-q}F_m}.
\]

先看前向方向。由 \(K_m\) 的定义，
\[
K_{m\#}\pi_{m,q}=\mathrm{Bernoulli}(\theta_m(q)).
\]
而
\[
\theta_m(q)=\frac{\varphi^{-q}}{F_{m+1}/F_m+\varphi^{-q}},
\qquad
\theta(q)=\frac{\varphi^{-q}}{\varphi+\varphi^{-q}}.
\]
又因为
\[
\frac{F_{m+1}}{F_m}=\varphi+O(\varphi^{-2m}),
\]
并且 \(q\in[0,Q]\) 落在紧区间上，所以
\[
\sup_{0\le q\le Q}|\theta_m(q)-\theta(q)|
=
O_Q(\varphi^{-2m})
=
O_Q\!\Bigl(\Bigl(\frac{\varphi}{2}\Bigr)^m\Bigr).
\]
因此
\[
\sup_{0\le q\le Q}
D_{\mathrm{TV}}(K_{m\#}\pi_{m,q},\beta_q)
=
O_Q\!\Bigl(\Bigl(\frac{\varphi}{2}\Bigr)^m\Bigr).
\]

再看反向方向。由构造，
\[
L_{m\#}\beta_q=(1-\theta(q))U_{m,0}+\theta(q)U_{m,1}.
\]
而与 \(\nu_{m,q}\) 相比，两者仅块权重不同，故
\[
\|L_{m\#}\beta_q-\nu_{m,q}\|_{\mathrm{TV}}
=
|\theta(q)-\theta_m(q)|
=
O_Q\!\Bigl(\Bigl(\frac{\varphi}{2}\Bigr)^m\Bigr).
\]
再与 \(\pi_{m,q}\) 到 \(\nu_{m,q}\) 的估计合并，得到
\[
\sup_{0\le q\le Q}
D_{\mathrm{TV}}(L_{m\#}\beta_q,\pi_{m,q})
=
O_Q\!\Bigl(\Bigl(\frac{\varphi}{2}\Bigr)^m\Bigr).
\]
对 \(\Delta_Q(m)\) 取上述这对核 \((K_m,L_m)\) 作为候选，即得所示上界。证毕。
\end{proof}

\begin{theorem}[bin-fold escort 的全阶 Rényi 熵在常数层具有显式温度剖面]\label{thm:xi-time-part65h-escort-renyi-o1-temperature-profile}
沿用定理 \ref{thm:killo-fold-bin-escort-onebit-fisher} 的记号。对固定 \(q\ge 0\) 与 \(\alpha>0\)，定义 escort Rényi 熵
\[
H_{\alpha}(\pi_{m,q})
:=
\begin{cases}
\dfrac{1}{1-\alpha}\log\sum_{w\in X_m}\pi_{m,q}(w)^\alpha,
& \alpha\neq 1,\\[2mm]
-\sum_{w\in X_m}\pi_{m,q}(w)\log\pi_{m,q}(w),
& \alpha=1.
\end{cases}
\]
则当 \(m\to\infty\) 时，若 \(\alpha\neq 1\)，有
\[
H_{\alpha}(\pi_{m,q})
=
m\log\varphi
-\frac12\log 5
+
\frac{\alpha\log(\varphi+\varphi^{-q})-\log(\varphi+\varphi^{-\alpha q})}{\alpha-1}
+O\!\Bigl(\Bigl(\frac{\varphi}{2}\Bigr)^m\Bigr),
\]
而当 \(\alpha=1\) 时，有
\[
H(\pi_{m,q})
=
m\log\varphi
-\frac12\log 5
+\log(\varphi+\varphi^{-q})
+\frac{q\log\varphi}{1+\varphi^{q+1}}
+O\!\Bigl(\Bigl(\frac{\varphi}{2}\Bigr)^m\Bigr).
\]
因此，escort 温度族不只在熵率层面塌缩到 \(\log\varphi\)；其全部 \(O(1)\) 剩余结构也由 \(q\) 的显式闭式函数完全控制。
\end{theorem}
\begin{proof}
先设 \(\alpha\neq 1\)。由定义，
\[
\sum_{w\in X_m}\pi_{m,q}(w)^\alpha
=
\frac{S_{\alpha q}^{\mathrm{bin}}(m)}{(S_q^{\mathrm{bin}}(m))^\alpha}.
\]
而命题 \ref{prop:fold-bin-power-sum-asymptotic} 给出，对每个固定 \(t\ge 0\)，
\[
S_t^{\mathrm{bin}}(m)
=
\Bigl(\frac{2}{\varphi}\Bigr)^{tm}
\bigl(F_{m+1}+\varphi^{-t}F_m\bigr)
\Bigl(1+O\!\Bigl(\Bigl(\frac{\varphi}{2}\Bigr)^m\Bigr)\Bigr).
\]
将 \(t=q\) 与 \(t=\alpha q\) 代入，即得
\[
\log\sum_{w\in X_m}\pi_{m,q}(w)^\alpha
=
\log\bigl(F_{m+1}+\varphi^{-\alpha q}F_m\bigr)
-\alpha\log\bigl(F_{m+1}+\varphi^{-q}F_m\bigr)
+O\!\Bigl(\Bigl(\frac{\varphi}{2}\Bigr)^m\Bigr).
\]
再由 Binet 公式，
\[
F_{m+1}+\varphi^{-u}F_m
=
\frac{\varphi^m}{\sqrt5}\bigl(\varphi+\varphi^{-u}\bigr)+O(\varphi^{-m}),
\]
故
\[
\log\bigl(F_{m+1}+\varphi^{-u}F_m\bigr)
=
m\log\varphi-\frac12\log5+\log(\varphi+\varphi^{-u})
+O(\varphi^{-2m}).
\]
代回并除以 \(1-\alpha\)，便得到所示 Rényi 公式。

再看 Shannon 情形。由定理 \ref{thm:xi-fold-escort-renyi-shannon-constants} 的证明所用恒等式，
\[
H(\pi_{m,q})
=
\log S_q^{\mathrm{bin}}(m)-q\,\kappa_m^{\mathrm{bin}}(q),
\]
其中
\[
\kappa_m^{\mathrm{bin}}(q):=
\sum_{w\in X_m}\pi_{m,q}(w)\log d_m^{\mathrm{bin}}(w).
\]
由推论 \ref{cor:xi-time-part9o-escort-fiberinfo-constantterm-inversion}，
\[
\kappa_m^{\mathrm{bin}}(q)
=
m\log\!\Bigl(\frac{2}{\varphi}\Bigr)
-\frac{\log\varphi}{1+\varphi^{q+1}}
+O\!\Bigl(\Bigl(\frac{\varphi}{2}\Bigr)^m\Bigr).
\]
另一方面，把 \(t=q\) 代入上述幂和展开并取对数，得到
\[
\log S_q^{\mathrm{bin}}(m)
=
qm\log\!\Bigl(\frac{2}{\varphi}\Bigr)
+m\log\varphi
-\frac12\log5
+\log(\varphi+\varphi^{-q})
+O\!\Bigl(\Bigl(\frac{\varphi}{2}\Bigr)^m\Bigr).
\]
两式合并即得
\[
H(\pi_{m,q})
=
m\log\varphi
-\frac12\log5
+\log(\varphi+\varphi^{-q})
+\frac{q\log\varphi}{1+\varphi^{q+1}}
+O\!\Bigl(\Bigl(\frac{\varphi}{2}\Bigr)^m\Bigr).
\]
证毕。
\end{proof}

\begin{theorem}[escort Shannon 常数满足精确 de Bruijn--Fisher 身份]\label{thm:xi-time-part65h-escort-shannon-debruijn-fisher}
沿用定理 \ref{thm:xi-time-part65h-escort-renyi-o1-temperature-profile} 的记号。定义 Shannon 常数层
\[
\mathfrak h(q):=
\lim_{m\to\infty}\bigl(H(\pi_{m,q})-m\log\varphi\bigr).
\]
则有闭式
\[
\mathfrak h(q)
=
-\frac12\log5+\log(\varphi+\varphi^{-q})
+\frac{q\log\varphi}{1+\varphi^{q+1}}.
\]
进一步，记
\[
\theta(q):=\frac{1}{1+\varphi^{q+1}},
\qquad
\mathcal I(q):=(\log\varphi)^2\theta(q)\bigl(1-\theta(q)\bigr).
\]
则成立精确微分恒等式
\[
\mathfrak h'(q)=-q\,\mathcal I(q),
\]
以及积分表示
\[
\mathfrak h(q)=\mathfrak h(0)-\int_0^q t\,\mathcal I(t)\,\dd t.
\]
其中 \(\mathcal I(q)\) 正是结论 \ref{con:xi-time-part9o-escort-fisher-rao-closed} 给出的极限 Fisher 信息密度。
\end{theorem}
\begin{proof}
第一条公式正是定理 \ref{thm:xi-time-part65h-escort-renyi-o1-temperature-profile} 在 \(\alpha=1\) 时的常数项。

令
\[
A(q):=\varphi+\varphi^{-q}.
\]
则
\[
\theta(q)=\frac{\varphi^{-q}}{A(q)}=\frac{1}{1+\varphi^{q+1}},
\]
并且
\[
\mathfrak h(q)
=
-\frac12\log5+\log A(q)+q\,\theta(q)\log\varphi.
\]
直接求导得
\[
\frac{A'(q)}{A(q)}
=
-\frac{\varphi^{-q}\log\varphi}{A(q)}
=
-\theta(q)\log\varphi,
\]
而
\[
\theta'(q)
=
-(\log\varphi)\theta(q)\bigl(1-\theta(q)\bigr).
\]
因此
\[
\mathfrak h'(q)
=
-\theta(q)\log\varphi+\theta(q)\log\varphi+q\,\theta'(q)\log\varphi
=
-q(\log\varphi)^2\theta(q)\bigl(1-\theta(q)\bigr)
=
-q\,\mathcal I(q).
\]
最后对 \(q\) 从 \(0\) 到任意给定值积分，即得
\[
\mathfrak h(q)=\mathfrak h(0)-\int_0^q t\,\mathcal I(t)\,\dd t.
\]
证毕。
\end{proof}

\begin{theorem}[escort 温度间的 Hellinger--Chernoff 几何由闭式母核完全支配]\label{thm:xi-time-part65h-escort-hellinger-chernoff-kernel}
沿用定理 \ref{thm:killo-fold-bin-escort-onebit-fisher} 的记号。对任意 \(p,q\ge 0\) 与 \(s\in[0,1]\)，定义混合亲和量
\[
\mathcal M_s^{(m)}(p,q):=
\sum_{w\in X_m}\pi_{m,p}(w)^s\pi_{m,q}(w)^{1-s}.
\]
则极限
\[
\mathcal M_s(p,q):=\lim_{m\to\infty}\mathcal M_s^{(m)}(p,q)
\]
存在，并满足精确闭式
\[
\mathcal M_s(p,q)
=
\frac{\varphi+\varphi^{-(sp+(1-s)q)}}
{(\varphi+\varphi^{-p})^s(\varphi+\varphi^{-q})^{1-s}}.
\]
特别地，Hellinger 亲和量
\[
\mathcal H(p,q):=\mathcal M_{1/2}(p,q)
\]
满足
\[
\mathcal H(p,q)
=
\frac{\varphi+\varphi^{-(p+q)/2}}
{\sqrt{(\varphi+\varphi^{-p})(\varphi+\varphi^{-q})}},
\]
而 Chernoff 母核可写为
\[
\mathcal C(p,q):=\inf_{0\le s\le 1}\mathcal M_s(p,q).
\]
因此，两温度 escort 之间的全部 Hellinger--Chernoff 判别几何都由单参数显式核
\[
s\longmapsto \mathcal M_s(p,q)
\]
完全控制。
\end{theorem}
\begin{proof}
对任意 \(s\in[0,1]\)，由 escort 分布的定义，
\[
\pi_{m,p}(w)^s\pi_{m,q}(w)^{1-s}
=
\frac{d_m^{\mathrm{bin}}(w)^{sp+(1-s)q}}
{\bigl(S_p^{\mathrm{bin}}(m)\bigr)^s\bigl(S_q^{\mathrm{bin}}(m)\bigr)^{1-s}}.
\]
故有严格恒等式
\[
\mathcal M_s^{(m)}(p,q)
=
\frac{S_{sp+(1-s)q}^{\mathrm{bin}}(m)}
{\bigl(S_p^{\mathrm{bin}}(m)\bigr)^s\bigl(S_q^{\mathrm{bin}}(m)\bigr)^{1-s}}.
\]
对分子与分母分别应用命题 \ref{prop:fold-bin-power-sum-asymptotic}，
\[
S_t^{\mathrm{bin}}(m)
=
\Bigl(\frac{2}{\varphi}\Bigr)^{tm}
\bigl(F_{m+1}+\varphi^{-t}F_m\bigr)
\Bigl(1+O\!\Bigl(\Bigl(\frac{\varphi}{2}\Bigr)^m\Bigr)\Bigr)
\qquad (t\ge 0).
\]
其中指数因子完全抵消。
于是
\[
\mathcal M_s^{(m)}(p,q)
=
\frac{F_{m+1}+\varphi^{-(sp+(1-s)q)}F_m}
{\bigl(F_{m+1}+\varphi^{-p}F_m\bigr)^s
\bigl(F_{m+1}+\varphi^{-q}F_m\bigr)^{1-s}}
\Bigl(1+O\!\Bigl(\Bigl(\frac{\varphi}{2}\Bigr)^m\Bigr)\Bigr).
\]
再由 Binet 公式
\[
F_{m+1}+\varphi^{-u}F_m
=
\frac{\varphi^m}{\sqrt5}\bigl(\varphi+\varphi^{-u}\bigr)+O(\varphi^{-m}),
\]
便得到
\[
\mathcal M_s^{(m)}(p,q)
=
\frac{\varphi+\varphi^{-(sp+(1-s)q)}}
{(\varphi+\varphi^{-p})^s(\varphi+\varphi^{-q})^{1-s}}
+O\!\Bigl(\Bigl(\frac{\varphi}{2}\Bigr)^m\Bigr).
\]
令 \(m\to\infty\) 即得所示极限核。

取 \(s=\tfrac12\) 立即得到 Hellinger 亲和量公式；而 Chernoff 母核不过是同一闭式核在区间 \([0,1]\) 上的下确界。证毕。
\end{proof}

由此，escort 主线中的整数压力差、末位单比特充分化、Rényi 常数层、Fisher 温度耗散与 Hellinger--Chernoff 判别核便被进一步压缩到同一条闭合链上：乘法偏序一侧的全部亏损由射线势 \(\delta\) 精确编码，而统计实验一侧的全部有限温度几何则完全退化为由 \(\theta(q)=1/(1+\varphi^{q+1})\) 所驱动的一维 logistic 母模型。


\paragraph{Hankel 缺口商环、除子层级与高阶异常的交互阶刚性}
在过拟合递推的 Hankel--null 缺口已被识别为自由阿贝尔群，而常数层高阶交叉异常的 Möbius 分解与 moment 放大律也已闭合之后，还可把这两条链进一步压缩为一组更强的结构定理：前者给出原始缺口模的规范商环模型及其中间层级的除子格分类，后者给出多轴异常的最小交互阶刻画及其在 moment 编译下的严格保持。

\begin{theorem}[过拟合 Hankel 缺口模的规范商环模型]\label{thm:xi-time-part62af-hankel-gap-quotient-ring-model}
沿用定理 \ref{thm:pom-hankel-syndrome-gap-rank-defect} 的记号。设
\[
P(\lambda)=P_{\min}(\lambda)\,R(\lambda),
\qquad
d_{\min}:=\deg P_{\min},
\qquad
\Delta:=\deg R,
\]
其中 \(P_{\min},R\in\ZZ[\lambda]\) 均首一，且 \(P_{\min}\) 为从偏移 \(m_0\) 起的首一最小湮灭多项式。
对任意 \(n\ge \deg P+1\)，记
\[
\mathsf Q_n(P;s)
:=
\ker_{\ZZ}\bigl((\mathsf H^{(m_0)}_n)^{\mathsf T}\bigr)\big/\mathcal M_n(P).
\]
则映射
\[
\Theta_R:\ZZ[\lambda]_{<\Delta}\longrightarrow \mathsf Q_n(P;s),\qquad
U(\lambda)\longmapsto
\Bigl[\operatorname{coeff}_{<n}\!\bigl(U(\lambda)P_{\min}(\lambda)\bigr)\Bigr]
\]
为同构。经由首一多项式 \(R\) 的唯一余式代表，仍记其诱导同构为
\[
\Theta_R:\ZZ[\lambda]/(R)\xrightarrow{\sim}\mathsf Q_n(P;s),
\]
并有规范识别
\[
\boxed{\;
\mathsf Q_n(P;s)\ \cong\ \ZZ[\lambda]/(R).\;}
\]
若记 \(\bar\sigma\) 为截断移位在 \(\mathsf Q_n(P;s)\) 上诱导的端同态，则在 \(\Theta_R\) 下，
\[
\bar\sigma
\quad\longleftrightarrow\quad
\mathsf T_R:\ZZ[\lambda]_{<\Delta}\to \ZZ[\lambda]_{<\Delta},
\qquad
\mathsf T_R(U):=\operatorname{rem}(\lambda U,R).
\]
因此 \(\bar\sigma\otimes_{\ZZ}\QQ\) 的特征多项式恰为 \(R(\lambda)\)。
\end{theorem}
\begin{proof}
由定理 \ref{thm:pom-hankel-syndrome-gap-rank-defect}，
\[
\mathsf Q_n(P;s)
\cong
\ZZ[\lambda]_{<n-d_{\min}}\Big/ R\cdot \ZZ[\lambda]_{<n-\deg P}.
\]
任取 \(A(\lambda)\in \ZZ[\lambda]_{<n-d_{\min}}\)。由于 \(R\) 首一，带余除法给出唯一分解
\[
A(\lambda)=R(\lambda)Q(\lambda)+U(\lambda),
\qquad
\deg U<\Delta.
\]
又因
\[
\deg A<n-d_{\min}=(n-\deg P)+\Delta,
\]
故必有 \(\deg Q<n-\deg P\)。因此上述商群的每个商类都由唯一余式 \(U\in \ZZ[\lambda]_{<\Delta}\) 代表，从而得到 \(\Theta_R\) 的双射性；这也正是
\(
\mathsf Q_n(P;s)\cong \ZZ[\lambda]/(R)
\)
的规范实现。

再看移位作用。把 \(\ZZ^n\) 与 \(\ZZ[\lambda]_{<n}\) 以系数向量同一，并令
\[
\sigma\bigl(\operatorname{coeff}_{<n}(A)\bigr):=
\operatorname{coeff}_{<n}(\lambda A).
\]
由于 \(\mathcal M_n(P_{\min})\) 与 \(\mathcal M_n(P)\) 都是截断倍乘子模，二者均对 \(\sigma\) 不变，故 \(\sigma\) 在商模 \(\mathsf Q_n(P;s)\) 上诱导出 \(\bar\sigma\)。对任意 \(U\in\ZZ[\lambda]_{<\Delta}\)，有
\[
\bar\sigma\bigl(\Theta_R(U)\bigr)
=
\Bigl[\operatorname{coeff}_{<n}\!\bigl(\lambda U P_{\min}\bigr)\Bigr].
\]
把 \(\lambda U\) 对 \(R\) 作带余除法，写成
\[
\lambda U=R Q+\operatorname{rem}(\lambda U,R).
\]
则
\[
\lambda U P_{\min}
=
Q P + \operatorname{rem}(\lambda U,R)\,P_{\min},
\]
其中 \(Q P\) 的截断系数向量落在 \(\mathcal M_n(P)\) 中，故在商模内消失，从而
\[
\bar\sigma\bigl(\Theta_R(U)\bigr)
=
\Theta_R\!\bigl(\operatorname{rem}(\lambda U,R)\bigr).
\]
这就证明了 \(\bar\sigma\) 与 \(\mathsf T_R\) 的共轭关系。最后，在基
\(
1,\lambda,\dots,\lambda^{\Delta-1}
\)
下，\(\mathsf T_R\otimes_{\ZZ}\QQ\) 的矩阵正是 \(R\) 的 Frobenius 伴随块，故其特征多项式恰为 \(R\)。证毕。
\end{proof}

\begin{theorem}[过拟合综合征层级的除子格分类]\label{thm:xi-time-part62af-hankel-gap-divisor-lattice}
\leanverified{paper\_xi\_time\_part62af\_hankel\_gap\_divisor\_lattice}
沿用定理 \ref{thm:xi-time-part62af-hankel-gap-quotient-ring-model} 的记号。对每个首一因子 \(D\mid R\)，定义中间综合征子模
\[
\mathcal N_n(D):=\mathcal M_n(P_{\min}D)
\subseteq
\ker_{\ZZ}\bigl((\mathsf H^{(m_0)}_n)^{\mathsf T}\bigr).
\]
则对任意首一因子 \(D_1,D_2\mid R\)，有
\[
\boxed{\;
D_1\mid D_2
\quad\Longleftrightarrow\quad
\mathcal N_n(D_2)\subseteq \mathcal N_n(D_1),\;}
\]
以及
\[
\boxed{\;
\mathcal N_n(D_1)\cap \mathcal N_n(D_2)
=
\mathcal N_n\bigl(\operatorname{lcm}(D_1,D_2)\bigr).\;}
\]
此外，在定理 \ref{thm:xi-time-part62af-hankel-gap-quotient-ring-model} 的同构
\[
\Theta_R:\ZZ[\lambda]/(R)\xrightarrow{\sim}\mathsf Q_n(P;s)
\]
下，有精确识别
\[
\boxed{\;
\Theta_R^{-1}\!\bigl(\mathcal N_n(D)/\mathcal M_n(P)\bigr)
=
(D)/(R)\subseteq \ZZ[\lambda]/(R).\;}
\]
因而
\[
\boxed{\;
\rank_{\ZZ}\bigl(\mathcal N_n(D)/\mathcal M_n(P)\bigr)
=
\Delta-\deg D,\;}
\qquad
\boxed{\;
\rank_{\ZZ}\Bigl(\mathsf Q_n(P;s)\big/\bigl(\mathcal N_n(D)/\mathcal M_n(P)\bigr)\Bigr)
=
\deg D.\;}
\]
\end{theorem}
\begin{proof}
由命题 \ref{prop:pom-hankel-syndrome-divisibility-reversal-lcm}，对
\[
P_1=P_{\min}D_1,
\qquad
P_2=P_{\min}D_2
\]
应用整除--包含反序对应，立得
\[
D_1\mid D_2
\iff
P_{\min}D_1\mid P_{\min}D_2
\iff
\mathcal N_n(D_2)\subseteq \mathcal N_n(D_1),
\]
且交公式同理由该命题给出。

再证商环识别。取 \(U\in \ZZ[\lambda]_{<\Delta}\)。由定理
\ref{thm:xi-time-part62af-hankel-gap-quotient-ring-model}，类
\[
\Theta_R(U)
=
\Bigl[\operatorname{coeff}_{<n}(U P_{\min})\Bigr]
\]
落入 \(\mathcal N_n(D)/\mathcal M_n(P)\) 当且仅当
\[
\operatorname{coeff}_{<n}(U P_{\min})\in \mathcal M_n(P_{\min}D).
\]
由于
\[
\deg(UP_{\min})<\Delta+d_{\min}=\deg P<n,
\]
此时截断无损；故上式等价于存在 \(Q\in\ZZ[\lambda]\) 使
\[
U P_{\min}=Q\,P_{\min}D.
\]
在整环 \(\ZZ[\lambda]\) 中约去首一多项式 \(P_{\min}\)，便得到
\[
U=Q D.
\]
也就是说，\(\Theta_R(U)\) 恰当且仅当 \(U\) 落在理想 \((D)\) 中。于是
\[
\Theta_R^{-1}\!\bigl(\mathcal N_n(D)/\mathcal M_n(P)\bigr)=(D)/(R).
\]

若写 \(R=D\widetilde R\)，则每个 \((D)/(R)\) 中的类都唯一写成
\[
D V\pmod R,
\qquad
V\in\ZZ[\lambda]_{<\deg\widetilde R},
\]
故其为秩 \(\deg\widetilde R=\Delta-\deg D\) 的自由阿贝尔群。后一个秩公式则由
\[
\rank_{\ZZ}\mathsf Q_n(P;s)=\Delta
\]
与前式相减即得。证毕。
\end{proof}

\begin{theorem}[高阶交叉异常的最小交互阶刻画]\label{thm:xi-time-part62af-cross-anom-interaction-degree}
沿用定义 \ref{def:pom-higher-cross-anomaly} 与定理 \ref{thm:pom-mobius-complete-decomposition} 的记号，记
\[
f_\theta(\chi):=\log\mathfrak{M}_\chi(\theta),
\qquad
\chi\in\widehat G=\prod_{i=1}^{k}\widehat{G_i}.
\]
对 \(1\le r\le k\)，称 \(f_\theta\) 的交互阶不超过 \(r\)，若存在函数族
\[
L_T:\prod_{i\in T}\widehat{G_i}\to\RR
\qquad
\bigl(\varnothing\neq T\subseteq [k],\ |T|\le r\bigr)
\]
使得对一切 \(\chi\in\widehat G\) 都有
\[
f_\theta(\chi)
=
\sum_{\substack{\varnothing\neq T\subseteq [k]\\ |T|\le r}}
L_T(\chi_T).
\]
记满足该条件的最小 \(r\) 为 \(\deg_{\mathrm{int}}f_\theta\)，并约定当
\(
f_\theta\equiv 0
\)
时
\(
\deg_{\mathrm{int}}f_\theta:=0
\)
（等价地，\(\max\varnothing:=0\)）。
则对任意 \(1\le r\le k\)，成立严格等价
\[
\boxed{\;
\deg_{\mathrm{int}}f_\theta\le r
\quad\Longleftrightarrow\quad
\Delta_S(\cdot;\theta)\equiv 0
\ \text{对所有}\ |S|>r.\;}
\]
特别地，
\[
\boxed{\;
\deg_{\mathrm{int}}f_\theta
=
\max\bigl\{|S|:\ \Delta_S(\cdot;\theta)\not\equiv 0\bigr\}.\;}
\]
\end{theorem}
\begin{proof}
若 \(\deg_{\mathrm{int}}f_\theta\le r\)，则
\[
f_\theta=\sum_{|T|\le r}L_T(\chi_T).
\]
任取 \(|S|>r\)。对任一求和项 \(L_T(\chi_T)\)，总可取某个 \(i\in S\setminus T\)，于是定义
\ref{def:pom-higher-cross-anomaly} 证明中出现的差分算子 \(\delta_i\) 满足
\[
\delta_i L_T=0.
\]
因此
\[
\delta_S L_T=0
\qquad
(\forall\,|T|\le r),
\]
从而
\[
\Delta_S(\cdot;\theta)=\delta_S f_\theta\equiv 0.
\]

反之，若所有 \(|S|>r\) 的 \(\Delta_S(\cdot;\theta)\) 都恒为零，则由定理
\ref{thm:pom-mobius-complete-decomposition}，
\[
f_\theta(\chi)
=
\sum_{\varnothing\neq S\subseteq [k]}\Delta_S(\chi_S;\theta)
=
\sum_{\substack{\varnothing\neq S\subseteq [k]\\ |S|\le r}}
\Delta_S(\chi_S;\theta).
\]
取 \(L_S:=\Delta_S(\cdot;\theta)\) 即得到一个交互阶不超过 \(r\) 的表示。

最后一个公式只是上述等价的最小化表达：\(\deg_{\mathrm{int}}f_\theta\) 恰等于最高非零 Möbius 分量的阶数。证毕。
\end{proof}

\begin{theorem}[moment 编译下交互阶的刚性保持与纯三体失明]\label{thm:xi-time-part62af-moment-interaction-degree-rigidity}
设 \(G_1,\dots,G_k\) 为有限阿贝尔群，记
\[
\widehat G:=\prod_{i=1}^{k}\widehat{G_i},
\qquad
\chi^0:=(\chi_{1,0},\dots,\chi_{k,0}).
\]
取任意投影词 \(w\)，并以常数层角色坐标定义其异常函数
\[
a_w:\widehat G\to \RR,
\qquad
a_w(\chi):=[\mathrm{Anom}(w)]_\chi,
\qquad
a_w(\chi^0):=0.
\]
对每个 \(q\ge 2\)，记
\[
w^{\boxtimes q}:=\mathrm{MOM}_q\circ w.
\]
由命题 \ref{prop:pom-anom-gap-amplification} 与推论
\ref{cor:pom-anom-gap-amplification-completeness} 的坐标化读法，常数层异常满足
\[
a_{w^{\boxtimes q}}=q\,a_w.
\]
再按定义 \ref{def:pom-higher-cross-anomaly} 的同一 Möbius 差分记
\[
\Delta_S(w):=\Delta_S(a_w)
\qquad
(\varnothing\neq S\subseteq [k]),
\]
并定义
\[
\deg_{\mathrm{int}}(w):=\deg_{\mathrm{int}}(a_w).
\]
则对任意 \(q\ge 2\) 与任意非空 \(S\subseteq [k]\)，有
\[
\boxed{\;
\Delta_S(w^{\boxtimes q})=q\,\Delta_S(w).\;}
\]
特别地，
\[
\boxed{\;
\deg_{\mathrm{int}}(w^{\boxtimes q})
=
\deg_{\mathrm{int}}(w).\;}
\]
进一步，若 \(w\) 满足
\[
\Delta_{\{i,j\}}(w)\equiv 0
\qquad
(\forall\, i\neq j),
\]
且对某个三元组 \(i<j<k\) 有
\[
\Delta_{\{i,j,k\}}(w)\not\equiv 0,
\]
则对一切 \(q\ge 2\) 仍成立
\[
\Delta_{\{i,j\}}(w^{\boxtimes q})\equiv 0
\qquad
(\forall\, i\neq j),
\]
以及
\[
\Delta_{\{i,j,k\}}(w^{\boxtimes q})
=
q\,\Delta_{\{i,j,k\}}(w)
\not\equiv 0.
\]
因此，任何仅依赖两体交叉异常的全矩阶审计都无法恢复这类纯三体不可分性。
\end{theorem}
\begin{proof}
由
\(
a_{w^{\boxtimes q}}=q\,a_w
\)
及 Möbius 差分的线性性，立得
\[
\Delta_S(w^{\boxtimes q})
=
\Delta_S(q\,a_w)
=
q\,\Delta_S(w)
\qquad
(\varnothing\neq S\subseteq [k]).
\]
于是
\[
\Delta_S(w)\equiv 0
\iff
\Delta_S(w^{\boxtimes q})\equiv 0,
\]
故非零支撑集合在 moment 编译下保持不变。由定理
\ref{thm:xi-time-part62af-cross-anom-interaction-degree} 的最小交互阶公式，即得
\[
\deg_{\mathrm{int}}(w^{\boxtimes q})
=
\deg_{\mathrm{int}}(w).
\]

若 \(w\) 为纯三体型，则所有两体分量原本恒为零，而某个三体分量非零；把上式分别应用于 \(S=\{i,j\}\) 与 \(S=\{i,j,k\}\) 即得所示结论。命题
\ref{prop:pom-pairwise-cross-anom-insufficient} 进一步保证这种制度确实存在。证毕。
\end{proof}

\noindent
因此，过拟合递推所留下的全部原始 Hankel 缺口并非任意自由方向，而是由单一首一因子 \(R\) 的商环与除子格完全控制；与此同时，多轴异常的真实复杂度也并不由两体审计决定，而由最高非零 Möbius 分量的阶数决定，并在 moment 放大下保持严格不变。

\paragraph{PRIMEGAME 的估值格点共轭、有限支撑分解与局部化 solenoid 宿主}
在自然扩张的外置预算、prime-register 账本语义与局部化 solenoid 终对象已经确立之后，固定 PRIMEGAME 微码的算术动力学还可进一步压缩为一条有限活跃支撑链：整数状态、first-fit 控制与局部化紧化在同一估值坐标中严格闭合。
沿用定义 \ref{def:pom-fractran}、命题 \ref{prop:pom-fractran-prime-translation} 与定理 \ref{thm:xi-time-part56-solenoid-terminal-prime-register-recovery} 的记号。固定一条十四指令 FRACTRAN 程序
\[
F=\left(\frac{a_1}{b_1},\dots,\frac{a_{14}}{b_{14}}\right),
\]
并称之为 PRIMEGAME。设其活跃素数集为
\[
S_*:=\{2,3,5,7,11,13,17,19,23,29\},
\]
且 PRIMEGAME 的全部分子与分母都只含 \(S_*\) 中的素因子。记
\[
\NN_{S_*}:=\left\{\prod_{p\in S_*}p^{x_p}:x_p\in \ZZ_{\ge 0}\right\},
\qquad
P_*:=\prod_{p\in S_*}p,
\qquad
G_{S_*}:=\ZZ[S_*^{-1}],
\]
并约定 \((\ZZ_{\ge 0})^{S_*}\) 与 \(\ZZ^{S_*}\) 上的偏序 \(x\ge y\) 总按逐坐标解释。

\begin{theorem}[PRIMEGAME 的精确 prime-register 共轭]\label{thm:xi-time-part62ae-primegame-exact-prime-register-conjugacy}
定义估值坐标
\[
\Psi_{S_*}:\NN_{S_*}\longrightarrow (\ZZ_{\ge 0})^{S_*},
\qquad
\Psi_{S_*}(n):=(v_p(n))_{p\in S_*}.
\]
对每个 \(1\le i\le 14\)，记
\[
\beta_i^{(*)}:=(v_p(b_i))_{p\in S_*}\in (\ZZ_{\ge 0})^{S_*},
\qquad
\delta_i^{(*)}:=(v_p(a_i)-v_p(b_i))_{p\in S_*}\in \ZZ^{S_*}.
\]
则 PRIMEGAME 的一步演化 \(T_F\) 在 \(\NN_{S_*}\) 上与估值格点上的分段平移
\[
\widetilde T_F(x):=x+\delta_{i(x)}^{(*)},
\qquad
i(x):=\min\{\,i:\ x\ge \beta_i^{(*)}\,\}
\]
严格共轭，即
\[
\Psi_{S_*}\circ T_F=\widetilde T_F\circ \Psi_{S_*}
\]
在 \(T_F\) 的定义域上恒成立。
\end{theorem}
\begin{proof}
命题 \ref{prop:pom-fractran-prime-translation} 已在全素数寄存器纤维上给出同一结论；此处只需把坐标限制到活跃支撑 \(S_*\)。对任意 \(n\in \NN_{S_*}\) 与任意指令 \(a_i/b_i\)，有
\[
\frac{a_i}{b_i}\ \text{可用于}\ n
\iff
b_i\mid n
\iff
v_p(n)\ge v_p(b_i)\quad(\forall\, p\in S_*),
\]
因为 PRIMEGAME 的全部分母都只含 \(S_*\)-素因子。故在 \(\Psi_{S_*}\)-坐标中，可行性恰由
\[
\Psi_{S_*}(n)\ge \beta_i^{(*)}
\]
刻画。若第 \(i\) 条指令被 first-fit 规则选中，则
\[
T_F(n)=n\frac{a_i}{b_i},
\]
从而对每个 \(p\in S_*\) 都有
\[
v_p(T_F(n))=v_p(n)+v_p(a_i)-v_p(b_i).
\]
于是
\[
\Psi_{S_*}(T_F(n))
=
\Psi_{S_*}(n)+\delta_i^{(*)}.
\]
再由 first-fit 选择定义，所取索引正是
\[
i(\Psi_{S_*}(n))
=
\min\{\,i:\ \Psi_{S_*}(n)\ge \beta_i^{(*)}\,\}.
\]
故共轭关系成立。证毕。
\end{proof}

\begin{theorem}[PRIMEGAME 的活跃支撑分解]\label{thm:xi-time-part62ae-primegame-active-support-splitting}
任取 \(n\in \ZZ_{>0}\)，则存在唯一分解
\[
n=u\,m,
\qquad
u\in \NN_{S_*},
\qquad
\gcd(m,P_*)=1.
\]
并且 PRIMEGAME 的一步演化满足：
\begin{enumerate}
  \item first-fit 选中的指令索引只由 \(u\) 决定，与 \(m\) 无关；
  \item 若 \(T_F(u)\) 有定义，则
  \[
  T_F(n)=m\,T_F(u);
  \]
  若 \(T_F(u)\) 无定义，则 \(T_F(n)\) 亦无定义。
\end{enumerate}
特别地，从初值 \(n_0=2\) 出发的 PRIMEGAME 轨道完全落在 \(\NN_{S_*}\) 内。
\end{theorem}
\begin{proof}
唯一分解来自整数的素因子分离：把 \(n\) 的 \(S_*\)-主部分并入 \(u\)，其余素因子并入 \(m\) 即得，而且此分解显然唯一。对任意指令 \(a_i/b_i\)，由于 \(\gcd(m,b_i)=1\) 且 \(b_i\) 的全部素因子都落在 \(S_*\)，有
\[
b_i\mid n
\iff
b_i\mid u.
\]
因此可行指令集合以及由此决定的 first-fit 索引只依赖于 \(u\)。若第 \(i\) 条指令适用于 \(u\)，则
\[
T_F(n)
=
n\frac{a_i}{b_i}
=
u\,m\frac{a_i}{b_i}
=
m\left(u\frac{a_i}{b_i}\right)
=
m\,T_F(u),
\]
其中最后一步利用了 PRIMEGAME 的全部分子也只含 \(S_*\)-素因子，故外部因子 \(m\) 完全惰性。若 \(u\) 上无可用指令，则 \(n\) 上也无可用指令。最后，由 \(2\in S_*\) 可知初值 \(n_0=2\) 对应 \(m=1\)；再由上式归纳即得整条轨道始终满足 \(m=1\)，从而完全落在 \(\NN_{S_*}\) 内。证毕。
\end{proof}

\begin{corollary}[PRIMEGAME 的有限阶段 solenoidal 宿主]\label{cor:xi-time-part62ae-primegame-finite-solenoidal-host}
PRIMEGAME 的一步控制律以及一切仅依赖 \(S_*\)-估值坐标的轨道观测，都先天地因子化到 \(G_{S_*}=\ZZ[S_*^{-1}]\) 的局部化阶段。其自然紧化宿主为有限阶段 arithmetic solenoid
\[
\Sigma_{S_*}:=\widehat{G_{S_*}},
\]
并满足 Pontryagin 对偶短正合列
\[
0\longrightarrow \prod_{p\in S_*}\ZZ_p
\longrightarrow \Sigma_{S_*}
\longrightarrow \TT
\longrightarrow 0.
\]
因此，对固定的 PRIMEGAME 微码而言，universal solenoid \(\widehat{\QQ}\) 并非操作层所必需；只有在把 PRIMEGAME 置入程序族、Mellin 读出或 \(\zeta\)-型完成时，才需要引入通用阶段。
\end{corollary}
\begin{proof}
定理 \ref{thm:xi-time-part62ae-primegame-active-support-splitting} 表明，PRIMEGAME 的一步选择与更新完全由 \(S_*\)-支撑部分 \(u\in \NN_{S_*}\) 决定，而与外部因子 \(m\) 解耦。因此任何只读取 \(S_*\)-估值坐标的轨道观测，都可经由 \(\NN_{S_*}\subset G_{S_*}\) 因子化。
另一方面，定理 \ref{thm:xi-time-part56-solenoid-terminal-prime-register-recovery} 与定理 \ref{thm:xi-localized-phase-sampling-closed-form} 已给出 \(G_{S_*}\) 的 Pontryagin 对偶紧化正是
\[
\Sigma_{S_*}=\widehat{G_{S_*}},
\]
并带有所示短正合列。故 PRIMEGAME 的自然紧化宿主已在有限局部化阶段闭合。证毕。
\end{proof}

\begin{corollary}[固定 FRACTRAN 微码的有限素数支撑刚性]\label{cor:xi-time-part62ae-fixed-fractran-finite-prime-support-rigidity}
任意固定长度的 FRACTRAN 程序
\[
F=\left(\frac{a_1}{b_1},\dots,\frac{a_k}{b_k}\right)
\]
都只激活有限个 prime-register 轴，即
\[
S(F):=\bigcup_{i=1}^k\{\,p\in \mathbb P:\ v_p(a_i b_i)>0\,\}
\]
为有限集，并且其一步 first-fit 控制与更新完全由该有限活跃支撑决定。因而单个固定程序表不能仅凭微码自身实现真正的无限素数控制；若要出现 universal-solenoidal 现象，必须把无穷性置于程序族、观测族或 Mellin/\(\zeta\) 完成过程中。
\end{corollary}
\begin{proof}
定理 \ref{thm:xi-time-part62ae-primegame-active-support-splitting} 的分解论证只用到了“程序表有限”与“全部分子分母的素因子支撑有限”这两点，因此对任意固定长度的 FRACTRAN 程序可逐字适用。把其中的 \(S_*\) 换成 \(S(F)\) 即得结论。证毕。
\end{proof}

因此，PRIMEGAME 在本文口径下可被精确识别为：哥德尔素数估值坐标给出状态空间，FRACTRAN 的 first-fit 规则给出微码执行，而有限局部化 solenoid 给出其紧化宿主。若继续推进谱侧分析，则自然对象将是 \(\Sigma_{S_*}\) 上的 transfer kernel、轨道 \(\zeta\) 与 Mellin 型读出。

\paragraph{单一 first-fit 微码的 classical 阻断、共尾素数耗尽与惰性谱扇区}
在 PRIMEGAME 的有限活跃支撑分解、局部化 solenoidal 宿主与有限素数账本刚性已经确定之后，若进一步要求单一 first-fit 微码直接抵达 classical \(\zeta\)/\(\Xi\) 的 Euler/Fredholm 层，则 prime-support、账本秩与有限阶段谱读出之间还会出现如下结构障碍。

\begin{theorem}[单一 first-fit 微码的经典 Euler 乘积阻断]\label{thm:xi-time-part62ae-single-firstfit-classical-euler-obstruction}
设
\[
F=\left(\frac{a_1}{b_1},\dots,\frac{a_k}{b_k}\right)
\]
为任意有限 FRACTRAN-compatible first-fit 微码，并记其活跃素数集为
\[
S(F):=\{\,p\in\mathbb P:\ v_p(a_i b_i)>0\ \text{for some }1\le i\le k\,\}.
\]
若某条证明路线坚持：其全部算术语义都由 \(F\) 的 prime-register 动力学，经结构保持的局部化 solenoidal 完成读出；并且该读出所诱导的 Euler/Fredholm 行列式在文稿既定的通用层语义下应与 classical \(\zeta\) 或完成函数 \(\Xi\) 一致，则该路线不可能由单一 \(F\) 实现。
\end{theorem}
\begin{proof}
记
\[
\NN_{S(F)}
:=
\left\{
\prod_{p\in S(F)}p^{x_p}:x_p\in\ZZ_{\ge 0}
\right\}.
\]
由推论 \ref{cor:xi-time-part62ae-fixed-fractran-finite-prime-support-rigidity}，\(F\) 的一步控制与更新完全由有限支撑 \(S(F)\) 决定。更具体地，若
\[
n=u\,m,
\qquad
u\in \NN_{S(F)},
\qquad
\gcd\!\Bigl(m,\prod_{p\in S(F)}p\Bigr)=1,
\]
则 first-fit 选中的指令索引只由 \(u\) 决定，且一旦 \(T_F(u)\) 有定义便有
\[
T_F(n)=m\,T_F(u).
\]
因此 \(F\) 的 prime-register 语义严格因子化到局部化阶段
\[
G_{S(F)}=\ZZ[S(F)^{-1}]
\]
之上。

再由定理 \ref{thm:killo-localization-circle-dimension-prime-ledger-splitting}，
\[
0\longrightarrow \prod_{p\in S(F)}\ZZ_p
\longrightarrow \widehat{G_{S(F)}}
\longrightarrow \TT
\longrightarrow 0,
\]
且 \(S(F)\) 可由核 \(\prod_{p\in S(F)}\ZZ_p\) 完全恢复。故任何仅从 \(\widehat{G_{S(F)}}\) 及其 Pontryagin 对偶账本 functorially 构造的对象，在结构层都不可能生成 \(S(F)\) 之外的新素数方向。

然而，文稿在 \S\ref{subsec:profinite-operator-l2} 中把“统一全部模数”的寄存器提升算子安置于 \(L^2(\widehat{\ZZ})\) 上，并在定理 \ref{thm:mellin-unitary-slice-unique}、定理 \ref{thm:app-horizon-toeplitz-lmi} 与定理 \ref{thm:xi-slice-horizon-carath-toeplitz} 中把 classical \(\Xi\)-切片的 Mellin--Plancherel 幺正读出及其 Toeplitz--Carath\'eodory 正性放在这一通用层上刻画。若单一有限微码已足以恢复 classical \(\zeta\) 或 \(\Xi\)，则该通用层中的全部素数尺度必须可由某个有限 \(S(F)\)-支撑对象结构性再现，矛盾。故结论成立。
\end{proof}

\begin{theorem}[共尾素数耗尽]\label{thm:xi-time-part62ae-cofinal-prime-exhaustion}
设 \(\mathscr S\) 为某一 FRACTRAN-compatible 路线中实际出现的全部有限 prime-support 之族。若该路线的目标是 classical \(\zeta\)/\(\Xi\) 而非固定有限核上的 RH 证书，则必存在一列有限素数集
\[
S_1\subset S_2\subset \cdots \subset \mathbb P,
\qquad
\bigcup_{n\ge 1}S_n=\mathbb P,
\]
并且每个 \(S_n\) 都可写成 \(\mathscr S\) 中有限多个支撑的并。特别地，通用 solenoidal 宿主并非附加修饰，而是由 prime-support 逻辑强制出的共尾极限对象。
\leanverified{paper\_xi\_time\_part62ae\_cofinal\_prime\_exhaustion}
\end{theorem}
\begin{proof}
若
\[
\bigcup_{S\in\mathscr S}S\neq \mathbb P,
\]
则可取某个素数 \(p_0\notin \bigcup_{S\in\mathscr S}S\)。于是对每个 \(S\in\mathscr S\)，定理 \ref{thm:killo-localization-circle-dimension-prime-ledger-splitting} 都给出
\[
\ker\!\bigl(\widehat{G_S}\to\TT\bigr)\cong \prod_{p\in S}\ZZ_p,
\]
其 \(p_0\)-主因子必为平凡。另一方面，定理 \ref{thm:xi-time-part62ac-localized-prime-adjoining-natural-extension} 表明：只有显式邻接素数 \(p_0\) 才会在自然扩张中新增一条 \(\ZZ_{p_0}\)-方向；推论 \ref{cor:xi-time-part62ac-prime-adjoining-order-independent-kernel-product} 又说明反复邻接后的最终对象只取决于实际并入的素数并。既然 \(p_0\) 从未出现，则任何由这些有限阶段及其共尾组合得到的对象都不可能在结构层看见 \(p_0\)。这与 classical Euler 乘积必须包含 \(p_0\)-尺度矛盾，故并集只能是 \(\mathbb P\)。

现取素数的固定枚举
\[
\mathbb P=\{p_1,p_2,\dots\}.
\]
对每个 \(n\ge 1\)，由上式并集等于 \(\mathbb P\)，可选 \(S^{(n)}\in\mathscr S\) 使 \(p_n\in S^{(n)}\)。定义
\[
S_n:=\bigcup_{j=1}^{n}S^{(j)}.
\]
则每个 \(S_n\) 皆有限，且
\[
S_1\subset S_2\subset\cdots,
\qquad
\bigcup_{n\ge 1}S_n=\mathbb P.
\]
再由推论 \ref{cor:xi-time-part62ac-prime-adjoining-order-independent-kernel-product}，每个有限阶段都对应规范短正合列
\[
0\longrightarrow \prod_{p\in S_n}\ZZ_p
\longrightarrow \Sigma_{S_n}
\longrightarrow \TT
\longrightarrow 0,
\qquad
\Sigma_{S_n}:=\widehat{G_{S_n}},
\]
并且各阶段之间的过渡正是逐素数邻接所诱导的共尾系统。故通用 solenoidal 宿主只能作为这些有限阶段的共尾极限出现，而不可能绕过 prime-support 账本被直接获得。证毕。
\end{proof}

\begin{theorem}[异常项与无穷秩账本的二择一]\label{thm:xi-time-part62ae-anomaly-infinite-rank-dichotomy}
设某条 Gödel--FRACTRAN 路线试图从 prime-register 微码抵达 classical \(\zeta\)/\(\Xi\)。若其 prime-support 最终覆盖全部素数，则至少发生下列两类机制之一：
\begin{enumerate}
  \item 承载乘法语义的有限阶段账本秩无界增长，并仅在共尾极限后进入通用层；
  \item 乘法到账本的传递在某一层失去严格同态性，而必须由异常签名、counterterm strictification 或普适可交换化群商承担补偿。
\end{enumerate}
特别地，不存在同时满足“prime-support 共尾覆盖全部素数、账本秩一致有界、乘法语义在各层皆严格同态”的证明链。
\end{theorem}
\begin{proof}
由定理 \ref{thm:xi-time-part62ae-cofinal-prime-exhaustion}，可取一列有限支撑
\[
S_1\subset S_2\subset\cdots\subset\mathbb P,
\qquad
\bigcup_{n\ge 1}S_n=\mathbb P.
\]
若在某一阶段中，\(S_n\) 内已有素数方向在账本中发生不可逆合并，从而不再可区分，则该路线上“乘法到账本”的传递已不再保持严格忠实语义，因而已经落入第二分支。故以下只需排除“各阶段都严格忠实同态化且账本秩一致有界”的可能性。

反设存在一致常数 \(R\)，使每一有限阶段的账本都可在秩至多 \(R\) 的加法/局部化对象上严格同态实现，并且在结构层忠实保留 \(S_n\) 中全部素数方向。由于 \(|S_n|\to\infty\)，可取 \(n\) 使
\[
|S_n|>R.
\]
记
\[
M_{S_n}:=\langle p:\ p\in S_n\rangle_{\times}\subset \NN^\times.
\]
则 \(M_{S_n}\cong (\NN^{|S_n|},+)\) 为由 \(|S_n|\) 个独立素数方向生成的自由交换幺半群。若其可在秩至多 \(R\) 的局部化账本中严格忠实实现，便与定理 \ref{thm:xi-time-part62debb-finite-prime-support-localized-axis-threshold} 及定理 \ref{thm:xi-localized-directsum-prime-truncation-minimal-channels} 矛盾。故只要严格同态性在各层保持，账本秩就不可能一致有界；这给出第一分支。

若反而坚持账本秩一致有界，则上面的矛盾表明：在某一有限阶段，乘法到账本的传递必不能继续保持严格同态。文稿内部对此类失配的可见补偿机制正是异常签名与 strictification 链：定义 \ref{def:pom-anom-signature} 给出异常对象，命题 \ref{prop:pom-counterterm-anom-cancel} 给出常数层对消律，命题 \ref{prop:pom-bc-quotient-strictify} 与命题 \ref{prop:pom-bc-quotient-universal} 给出普适可交换化群商，而定理 \ref{thm:dephys-offline-hair-curvature-reduction} 把这一失配压缩为“正规形—异常类—strictification”的统一归约。故任何绕开有限秩矛盾的路径都必须落入第二分支。特别地，三项条件不可能同时成立。证毕。
\end{proof}

\begin{corollary}[Ramanujan 窗口不可单独升级为 classical RH]\label{cor:xi-time-part62ae-ramanujan-window-no-classical-upgrade}
任意固定有限核家族上的 RH-sharp/Ramanujan 窗口结论，即便其边界由唯一临界参数与显式代数方程精确给出，也不能单独升级为 classical RH。特别地，命题 \ref{prop:real-input-40-collision-rhk-window} 给出的 \(4\times 4\) 碰撞位势窗口
\[
\Lambda(u)\le \lambda_1(u)^{1/2}
\iff
0<u\le u_R
\]
只构成一个有限核 RH 证书，而不构成 classical \(\Xi\) 的全局闭合。
\end{corollary}
\begin{proof}
命题 \ref{prop:real-input-40-collision-rhk-window} 已给出一个边界完全显式的有限维窗口，其参数 \(u_R\) 由唯一正实五次根确定。然 \S\ref{subsubsec:xi-horizon-ramanujan-feasible-reduction} 已明确指出：从 Ramanujan 机制出发逼近真实 RH，必须超出“固定有限维 \(+\) 单尺度 Dirichlet 化”的范式，进入增长维数的逼近序列或无穷维 Fredholm 行列式口径。

另一方面，定理 \ref{thm:xi-time-part62ae-single-firstfit-classical-euler-obstruction}、定理 \ref{thm:xi-time-part62ae-cofinal-prime-exhaustion} 与定理 \ref{thm:xi-time-part62ae-anomaly-infinite-rank-dichotomy} 表明：若目标是 classical \(\zeta\)/\(\Xi\)，则必须同时面对全素数共尾支撑、无界账本秩或异常补偿分支。单个固定有限核家族并不携带这些结构条件，故其 Ramanujan 窗口不能单独升级为 classical RH。证毕。
\end{proof}

\begin{proposition}[真实输入 \texorpdfstring{$40$}{40} 状态核的刚性惰性谱对]\label{prop:xi-time-part62ae-real-input40-rigid-inert-spectral-pair}
沿用定理 \ref{thm:real-input-40-det-uv-2plus8} 与推论 \ref{cor:real-input-40-zeta-uv-sqrtv-eigs} 的记号。则二维加权行列式
\[
\Delta(z;u,v)=\det(I-zM_{u,v})=(vz^2-1)P_8(z;u,v)
\]
对任意 \(u\) 恒带有一对参数刚性谱支
\[
\lambda_\pm=\pm\sqrt v.
\]
因此：
\begin{enumerate}
  \item 对每个整数 \(r\ge 1\)，都有
  \[
  \partial_u^{\,r}\log(vz^2-1)=0;
  \]
  \item 任何仅由 \(\partial_u^{\,r}\log\Delta(z;u,v)\) 所提取的谱统计，都只能读取 core 因子 \(P_8(z;u,v)\) 的变化，而对刚性谱对 \(\{\pm\sqrt v\}\) 完全失明；
  \item 因而任何仅以碰撞参数 \(u\) 的变形来读取“素数影子”的方案，都不可能在全谱层面保持忠实；若要与通用素数生成元层或 Mellin 幺正切片比对，必须先对该刚性扇区作商化或重整化。
\end{enumerate}
\end{proposition}
\begin{proof}
定理 \ref{thm:real-input-40-det-uv-2plus8} 与推论 \ref{cor:real-input-40-zeta-uv-sqrtv-eigs} 直接给出
\[
\Delta(z;u,v)=(vz^2-1)P_8(z;u,v),
\qquad
\lambda_\pm=\pm\sqrt v,
\]
其中 \(\lambda_\pm\) 与 \(u\) 无关。故对任意 \(r\ge 1\)，
\[
\partial_u^{\,r}\log(vz^2-1)=0,
\]
第一条成立。

进一步，
\[
\log\Delta(z;u,v)=\log(vz^2-1)+\log P_8(z;u,v),
\]
故一切 \(u\)-微分都自动杀掉刚性因子，只留下 \(P_8\) 的贡献。等价地，由定理 \ref{thm:xi-time-part73-twoparam-atomic-ghost-u-blindness} 与定理 \ref{thm:xi-time-part73-nonrigid-trace-newton-cayley-closure}，任何仅依赖 \(u\)-导数的 ghost/trace 审计都只能看到去除 \(\pm\sqrt v\) 后的非刚性迹谱 \(R_n(u,v)\)。

于是，若某个谱读出函子只使用 \(u\)-变形，则它会把“整条 determinant family”与“剔除刚性因子后的 core family”识别为同一对象；这在全谱层面不是忠实读出。故若要把该二维核接到通用素数生成元层或 Mellin 幺正切片的比较问题上，必须先将刚性惰性扇区 \((vz^2-1)\) 商去或重整化。证毕。
\end{proof}

\begin{conjecture}[Cofinal PRIMEGAME--\texorpdfstring{$\Xi$}{Xi} 极限定理]\label{conj:xi-time-part62ae-cofinal-primegame-xi-limit}
存在一族按有限素数集 \(S\subset\mathbb P\) 过滤的 FRACTRAN-compatible first-fit 微码 \(F_S\)，以及一族作用于 \(L^2(\widehat{\ZZ})\) 的迹类算子 \(\mathcal L_S\)，使得若记
\[
\widetilde\Delta_S(z):=\det(I-z\mathcal L_S),
\]
则满足：
\begin{enumerate}
  \item \(F_S\) 的 prime-register 语义因子化到
  \[
  G_S=\ZZ[S^{-1}];
  \]
  \item 若 \(S'\subset S\)，则 \(\widetilde\Delta_{S'}\) 与 \(\widetilde\Delta_S\) 在限制映射下相容；
  \item 当 \(S\nearrow \mathbb P\) 时，\(\mathcal L_S\) 在迹范数中收敛到某个迹类极限 \(\mathcal L_\infty\)，从而
  \[
  \widetilde\Delta_S(z)\longrightarrow
  \widetilde\Delta_\infty(z):=\det(I-z\mathcal L_\infty)
  \]
  于紧子集上一致成立；
  \item \(\widetilde\Delta_\infty\) 的 Mellin--Plancherel 幺正切片与圆盘完成化的 \(\Xi\)-截面一致。
\end{enumerate}
在此设定下，classical RH 等价于极限视界 Carath\'eodory 函数的 Toeplitz--PSD 层级成立。
\end{conjecture}

\paragraph{自然扩张层深、黄金共振 Schatten 排斥与单参数热力学塔的 RH 不可识别性}
在共尾 prime-support、有限阶段 solenoidal 宿主与 \(40\) 状态核的刚性谱扇区已经确定之后，还可把这三条链进一步压缩为一组彼此咬合的后果：自然扩张层深在局部化 prime-ledger 上具有严格线性下界；黄金共振振幅在算子侧排斥一切有限反项的 Schatten 正则化；而单参数热力学导数塔则不足以恢复圆盘完成化所需的 Toeplitz--Carath\'eodory 证书。

\begin{theorem}[自然扩张层深的 prime-slice 线性下界]\label{thm:xi-time-part62ae-natural-extension-prime-slice-linear-lower-bound}
沿用定理 \ref{thm:xi-layered-prime-register-sharp-natural-extension} 的设定。设
\[
\Psi_k(x_0):=(x_k,H_k(x_0)),
\qquad
H_k(x_0):=\prod_{j=0}^{k-1}h_j(x_0),
\]
且对每个 \(j<k\) 都存在两两不交的素数集 \(P_j\subset\mathbb P\)，使 \(h_j(x_0)\) 的全部素因子都落在 \(P_j\) 中。定义
\[
\operatorname{supp}_{\mathbb P}(h_j)
:=
\left\{
p\in\mathbb P:\ \exists\,x_0\in X_0,\ v_p\bigl(h_j(x_0)\bigr)>0
\right\}.
\]
若对每个 \(j<k\) 都存在 \(y_j\in X_{j+1}\) 使
\[
\#\pi_j^{-1}(y_j)\ge 2,
\]
且 \(\Psi_k\) 为单射，则有
\[
\#\Bigl(\bigcup_{j=0}^{k-1}\operatorname{supp}_{\mathbb P}(h_j)\Bigr)\ge k.
\]
更强地，对每个 \(j<k\) 都存在某个素数 \(p_j\in P_j\) 与某个微态 \(x_0^{(j)}\in X_0\) 使
\[
v_{p_j}\bigl(h_j(x_0^{(j)})\bigr)>0.
\]
\end{theorem}
\begin{proof}
由定理 \ref{thm:xi-layered-prime-register-sharp-natural-extension}，对每个 \(j<k\) 都存在某个微态 \(x_0^{(j)}\in X_0\) 使
\[
h_j\bigl(x_0^{(j)}\bigr)\neq 1.
\]
因此 \(h_j(x_0^{(j)})\) 至少含有一个素因子 \(p_j\in P_j\)，并满足
\[
v_{p_j}\bigl(h_j(x_0^{(j)})\bigr)>0.
\]
由于各 \(P_j\) 两两不交，故这些 \(p_j\) 两两不同，从而
\[
\#\Bigl(\bigcup_{j=0}^{k-1}\operatorname{supp}_{\mathbb P}(h_j)\Bigr)\ge k.
\]
证毕。
\end{proof}

\begin{corollary}[有限阶段 arithmetic solenoid 的分支历史容量]\label{cor:xi-time-part62ae-finite-solenoid-branch-history-capacity}
设 \(S\subset\mathbb P\) 为有限集，并记
\[
G_S=\ZZ[S^{-1}],
\qquad
0\longrightarrow \prod_{p\in S}\ZZ_p
\longrightarrow \widehat{G_S}
\longrightarrow \TT
\longrightarrow 0
\]
为定理 \ref{thm:xi-time-part56-solenoid-terminal-prime-register-recovery} 所给出的 arithmetic solenoid 短正合列。则任何 prime-register 方向受限于 \(S\) 的自然扩张忠实编码，至多只能承载 \(|S|\) 层非平凡分支历史。特别地，若某个有限纤维自映射 \(T:X\twoheadrightarrow X\) 的自然扩张在无穷多层上都存在真实分支，则不存在任何固定有限 \(S\) 使 \(\widehat{G_S}\) 忠实承载其全部前史；只能取一列
\[
S_1\subset S_2\subset\cdots,
\qquad
\bigcup_{n\ge 1}S_n=\mathbb P,
\]
并进入相应的共尾逆极限。
\end{corollary}
\begin{proof}
定理 \ref{thm:xi-time-part56-solenoid-terminal-prime-register-recovery} 与定理 \ref{thm:xi-localized-phase-sampling-closed-form} 表明：\(\widehat{G_S}\) 的全不连通 prime-ledger 核恰为
\[
\prod_{p\in S}\ZZ_p,
\]
因而其可用素数方向完全由 \(S\) 决定。若某个深度为 \(k\) 的自然扩张前史能在该有限阶段上被忠实编码，则它必然给出一个所有素因子都落在 \(S\) 中的单射分层 prime-register 外置。由定理 \ref{thm:xi-time-part62ae-natural-extension-prime-slice-linear-lower-bound}，此时至少需要 \(k\) 个互异素数方向，故 \(k\le |S|\)。

若 \(T\) 的自然扩张在无穷多层上均存在真实分支，则对任意 \(k\) 都可取长度 \(k\) 的前史截断，并由上段推出任何忠实编码都需要至少 \(k\) 个互异素数方向。于是任意固定有限 \(S\) 最终都会失效，只能取 \(S_n\nearrow\mathbb P\) 的共尾系统。证毕。
\end{proof}

\begin{theorem}[黄金共振振幅的 Hilbert--Schmidt 排斥]\label{thm:xi-time-part62ae-golden-resonance-hilbert-schmidt-obstruction}
设 \((e_u)_{u\ge 1}\) 为 Hilbert 空间 \(\mathcal H\) 中一组标准正交系，并令
\[
R
:=
\sum_{u\ge 1}C_\varphi(u)\,\langle\cdot,e_u\rangle e_u
\]
为其在自然极大定义域上的对角算子。则
\[
R\notin \mathcal S_2(\mathcal H).
\]
更强地，对任意 \(N\ge 0\)，尾算子
\[
R_{>N}
:=
\sum_{u>N}C_\varphi(u)\,\langle\cdot,e_u\rangle e_u
\]
亦不属于 \(\mathcal S_2(\mathcal H)\)。特别地，
\[
R\notin \mathcal S_1(\mathcal H).
\]
\end{theorem}
\begin{proof}
由定理 \ref{thm:fold-sigma-phi-diverges}，
\[
\Sigma_\varphi
:=
1+2\sum_{u=1}^{\infty}C_\varphi(u)^2
=+\infty,
\]
从而
\[
\sum_{u\ge 1}|C_\varphi(u)|^2=+\infty.
\]
这也正与推论 \ref{cor:fold-bernoulli-fourier-sampling-not-l2} 中 Fourier 采样序列的非平方可和性一致。对角算子的 Hilbert--Schmidt 范数满足
\[
\|R\|_{\mathcal S_2}^2
=
\sum_{u\ge 1}|C_\varphi(u)|^2,
\]
故 \(\|R\|_{\mathcal S_2}=+\infty\)，于是 \(R\notin\mathcal S_2(\mathcal H)\)。

同理，对任意固定 \(N\ge 0\)，有
\[
\|R_{>N}\|_{\mathcal S_2}^2
=
\sum_{u>N}|C_\varphi(u)|^2
=+\infty,
\]
故 \(R_{>N}\notin \mathcal S_2(\mathcal H)\)。最后，由 \(\mathcal S_1(\mathcal H)\subset \mathcal S_2(\mathcal H)\) 即得 \(R\notin \mathcal S_1(\mathcal H)\)。证毕。
\end{proof}

\begin{corollary}[有限反项 Fredholm 封口的不可能性]\label{cor:xi-time-part62ae-no-finite-counterterm-fredholm-closure}
在定理 \ref{thm:xi-time-part62ae-golden-resonance-hilbert-schmidt-obstruction} 的设定下，不存在任何有限秩算子 \(F\) 使
\[
R-F\in \mathcal S_1(\mathcal H).
\]
等价地，任何试图把原始黄金共振振幅直接作为 Mellin/Fourier 正交通道特征值的 Fredholm--Hilbert--P\'olya 封口，都不可能通过有限个 counterterm 获得迹类正则化。
\end{corollary}
\begin{proof}
若存在有限秩算子 \(F\) 使 \(R-F\in\mathcal S_1(\mathcal H)\)，则 \(R-F\in\mathcal S_2(\mathcal H)\)。另一方面，有限秩算子必属于 \(\mathcal S_2(\mathcal H)\)，故
\[
R=(R-F)+F\in\mathcal S_2(\mathcal H),
\]
与定理 \ref{thm:xi-time-part62ae-golden-resonance-hilbert-schmidt-obstruction} 矛盾。证毕。
\end{proof}

\begin{theorem}[单参数热力学塔对 RH 证书的不可识别性]\label{thm:xi-time-part62ae-oneparam-thermodynamic-tower-rh-indistinguishability}
固定 \(v>0\)。沿用定理 \ref{thm:real-input-40-det-uv-2plus8} 与推论 \ref{cor:real-input-40-zeta-uv-sqrtv-eigs} 的记号，定义
\[
\Delta_{\mathrm{full}}(z;u)
:=
\Delta(z;u,v)
=(vz^2-1)P_8(z;u,v),
\qquad
\Delta_{\mathrm{red}}(z;u):=P_8(z;u,v).
\]
对每个 \(r\ge 1\)，定义
\[
\mathcal C_r(z;u):=\partial_u^r\log\Delta(z;u,v).
\]
则对一切 \(r\ge 1\) 都有
\[
\partial_u^r\log \Delta_{\mathrm{full}}(z;u)
=
\partial_u^r\log \Delta_{\mathrm{red}}(z;u).
\]
因此，任何仅由 \(u\)-压力、\(u\)-矩、\(u\)-大偏差或其有限代数组合所生成、并最终只依赖这条 \(u\)-导数塔的热力学观察子，都无法区分 \(\Delta_{\mathrm{full}}\) 与 \(\Delta_{\mathrm{red}}\)。

然而，一旦把这两族对象送入文稿既定的圆盘完成化--Toeplitz--PSD 接口，则相应的 RH 正性证书并不相同。更一般地，若 \(B(z)\) 为任一与 \(u\) 无关的非常数完成因子，则 \(B(z)\Delta_{\mathrm{red}}(z;u)\) 与 \(\Delta_{\mathrm{red}}(z;u)\) 具有完全相同的全部 \(u\)-导数塔，但其 Toeplitz--Carath\'eodory 边界数据因
\[
\partial_z\log\!\bigl(B(z)\Delta_{\mathrm{red}}(z;u)\bigr)
=
\partial_z\log \Delta_{\mathrm{red}}(z;u)+\partial_z\log B(z)
\]
而被改写。特别地，取 \(B(z)=1-vz^2\) 时，\(\Delta_{\mathrm{full}}\) 与 \(\Delta_{\mathrm{red}}\) 在热力学塔上不可区分，却在完成化证书端保持可区分。
\end{theorem}
\begin{proof}
由
\[
\log\Delta_{\mathrm{full}}(z;u)
=
\log(vz^2-1)+\log P_8(z;u,v),
\]
并注意到 \(\log(vz^2-1)\) 与 \(u\) 无关，便对每个 \(r\ge 1\) 得到
\[
\partial_u^r\log\Delta_{\mathrm{full}}(z;u)
=
\partial_u^r\log P_8(z;u,v)
=
\partial_u^r\log\Delta_{\mathrm{red}}(z;u).
\]
因此任何仅由这条 \(u\)-导数塔构成的热力学观察子都无法区分二者。

另一方面，定理 \ref{thm:mellin-unitary-slice-unique} 与定理 \ref{thm:xi-slice-horizon-carath-toeplitz} 表明：classical RH 端所使用的是完成对象的 Mellin--Plancherel 幺正切片及其诱导的 Toeplitz--Carath\'eodory 边界数据，而非单一参数的热力学导数塔。若将完成对象乘上一个与 \(u\) 无关的非常数因子 \(B(z)\)，则全部 \(u\)-导数塔保持不变，但 \(\partial_z\log B(z)\) 会向边界系数额外注入一个非零序列，从而改变 Toeplitz 数据。取 \(B(z)=1-vz^2\) 即得所述具体例子。故从单参数热力学塔到 RH 证书的映射不是单射。证毕。
\end{proof}

\begin{corollary}[\texorpdfstring{$40$}{40} 状态核通向 classical RH 的最小额外输入]\label{cor:xi-time-part62ae-realinput40-minimal-extra-rh-input}
任何仅依赖 \(u\)-单参数热力学数据的 \(40\) 状态核--RH 方案，在逻辑上都不完备。若要抵达文稿中的 classical RH 终端证书，则至少必须补入下列三类信息之一：
\begin{enumerate}
  \item 一个与 \(u\) 独立的第二形变参数；
  \item 由分支历史或惯性谱扇区所决定的规范反因子；
  \item 直接给定的 Toeplitz--Carath\'eodory 边界数据。
\end{enumerate}
\end{corollary}
\begin{proof}
若单参数 \(u\)-热力学塔已足以恢复完成化 RH 证书，则定理 \ref{thm:xi-time-part62ae-oneparam-thermodynamic-tower-rh-indistinguishability} 中的 \(\Delta_{\mathrm{full}}\) 与 \(\Delta_{\mathrm{red}}\) 应在 RH 端不可区分；但该定理恰表明它们在全部 \(u\)-导数塔上不可区分，而在 Toeplitz--Carath\'eodory 证书端可区分，矛盾。故必须补入至少一类与单参数热力学塔独立的额外信息。证毕。
\end{proof}

\begin{conjecture}[分支惯性反因子与 reduced determinant 的共尾 \texorpdfstring{$\Xi$}{Xi} 极限]\label{conj:xi-time-part62ae-branch-inertia-counterfactor-cofinal-xi-limit}
猜想 \ref{conj:xi-time-part62ae-cofinal-primegame-xi-limit} 还应具有如下规范分解：存在一族按有限素数集 \(S\subset\mathbb P\) 过滤的 FRACTRAN-compatible / Fibonacci-compatible 内核 \(\Delta_S\)，以及与之伴随的一族规范反因子 \(B_S(z)\)，使得若记
\[
\Delta_S^{\mathrm{red}}(z):=\frac{\Delta_S(z)}{B_S(z)},
\]
则满足：
\begin{enumerate}
  \item \(B_S\) 恰吸收一切对所选热力学形变不可见、但对完成化证书可见的惯性谱支；
  \item \(\Delta_S^{\mathrm{red}}\) 对自然扩张分支历史保持忠实；
  \item 若 \(S'\subset S\)，则 \(B_{S'}\) 与 \(B_S\) 在限制映射下相容；
  \item 当 \(S\nearrow\mathbb P\) 时，\(\Delta_S^{\mathrm{red}}\) 在适当的 Fredholm--Mellin 口径下收敛到某个极限 \(\Delta_\infty\)，且 \(\Delta_\infty\) 的 Mellin--Plancherel 幺正切片与文稿圆盘完成化中的 \(\Xi\)-截面一致。
\end{enumerate}
在此设定下，classical RH 等价于 \(\Delta_\infty\) 所诱导的 Toeplitz--Carath\'eodory 数据全阶正性。
\end{conjecture}

综上，PRIMEGAME/FRACTRAN--solenoid--Mellin 路线在本文口径下呈现出一条明确的三重结构障碍：
\[
\boxed{\;
\text{分支历史要求共尾 prime-ledger，}\quad
\text{黄金共振排斥有限反项 Fredholm 封口，}\quad
\text{单参数热力学塔不能恢复 RH 的 Toeplitz 证书。}\;}
\]
因此，真正尚待建立的并非更多有限核实验，而是一套定义在共尾 prime-register 家族上的规范反因子理论，并要求其同时与自然扩张分支历史、有限阶段惯性谱扇区的剥离，以及 Mellin--Plancherel / Toeplitz--Carath\'eodory 终端接口严格相容。
\paragraph{有限分支前缀、Smith 曲率谱与几何对象的素数寄存器化}
在自然扩张的分支指标寄存器、有限深度 prime-slice 外置、Smith 信息亏损谱与整数轴对齐椭圆半群已经分别建立之后，还可把这四条线索压缩成一条更高阶的结构链：有限分支前缀给出 prime-slice 外置的规范前缀模型，Smith 曲率谱把全部模核损失压缩为离散边界函数，而几何对象的素数寄存器化则把这种刚性直接推送到结构层圆维。

\begin{theorem}[有限深度素数片外置是分支寄存器前缀的乘法封装]\label{thm:xi-finite-depth-primeslice-branch-prefix-compression}
设
\[
X_0\xrightarrow{\pi_0}X_1\xrightarrow{\pi_1}\cdots\xrightarrow{\pi_{k-1}}X_k
\]
为可数集合之间的有限纤维满射链。对每个 \(j=0,\dots,k-1\) 与每个 \(y\in X_{j+1}\)，固定单射
\[
\iota_{j,y}:\pi_j^{-1}(y)\hookrightarrow \NN.
\]
对 \(x_0\in X_0\) 递推定义
\[
x_{j+1}:=\pi_j(x_j),
\qquad
\beta_j(x_0):=\iota_{j,x_{j+1}}(x_j)\in\NN,
\]
并记有限深度分支寄存器
\[
\Psi_k(x_0):=(x_k,\beta_0(x_0),\dots,\beta_{k-1}(x_0))
\in X_k\times \NN^k.
\]
则 \(\Psi_k\) 为单射。

进一步，沿用定理 \ref{thm:killo-layered-prime-slices-finite-depth} 的两两不交 prime slices \((\mathcal P_j)_{0\le j\le k-1}\)，并为每层固定单射
\[
\alpha_j:X_{j+1}\times\NN\hookrightarrow \mathcal P_j.
\]
则存在一个定义在 \(\Psi_k(X_0)\) 上的规范映射
\[
\mathfrak M_k:\Psi_k(X_0)\longrightarrow X_k\times \ZZ_{>0}
\]
使得有限深度 prime-slice 外置
\[
\Phi_k(x_0):=
\left(
x_k,\ \prod_{j=0}^{k-1}\alpha_j(x_{j+1},\beta_j(x_0))
\right)
\]
满足
\[
\boxed{\;
\Phi_k=\mathfrak M_k\circ \Psi_k.\;}
\]
换言之，有限深度的 prime-slice 外置正是分支寄存器前缀的乘法封装；当该链来自某个有限纤维自映射的 \(k\)-步前史截断时，它恰是定理 \ref{thm:killo-natural-extension-branch-register} 与定理 \ref{thm:xi-prime-register-history-inverse-limit} 中自然扩张分支寄存器的长度 \(k\) 前缀压缩。
\end{theorem}
\begin{proof}
先证单射。若给定
\[
(x_k,a_0,\dots,a_{k-1})\in \Psi_k(X_0),
\]
则可自后向前唯一恢复：
\[
x_{k-1}=\iota_{k-1,x_k}^{-1}(a_{k-1}),
\qquad
x_{k-2}=\iota_{k-2,x_{k-1}}^{-1}(a_{k-2}),
\quad \dots,\quad
x_0=\iota_{0,x_1}^{-1}(a_0).
\]
因此 \(\Psi_k\) 在其像上可逆，从而为单射。

现定义 \(\mathfrak M_k\)。对任意 \((x_k,a_0,\dots,a_{k-1})\in \Psi_k(X_0)\)，先按上一步递推恢复唯一的中间态 \(x_{k-1},\dots,x_0\)，再置
\[
\mathfrak M_k(x_k,a_0,\dots,a_{k-1})
:=
\left(
x_k,\ \prod_{j=0}^{k-1}\alpha_j(x_{j+1},a_j)
\right).
\]
该定义良好，因为所有 \(x_j\) 已由 \((x_k,a_0,\dots,a_{k-1})\) 唯一确定。对真实输入 \(x_0\in X_0\)，代入 \(a_j=\beta_j(x_0)\) 即得
\[
(\mathfrak M_k\circ \Psi_k)(x_0)
=
\left(
x_k,\ \prod_{j=0}^{k-1}\alpha_j(x_{j+1},\beta_j(x_0))
\right)
=
\Phi_k(x_0).
\]
故 \(\Phi_k=\mathfrak M_k\circ \Psi_k\)。最后一条只是在 \(T\)-前史链情形下把 \(\Psi_k\) 识别为定理 \ref{thm:xi-prime-register-history-inverse-limit} 中截断分支码 \(\mathcal C_k\) 的一般链版本。证毕。
\end{proof}
\leanverified{paper\_xi\_finite\_depth\_primeslice\_branch\_prefix\_compression}

\begin{theorem}[分层素数片复杂度的精确最优值等于活跃分支深度]\label{thm:xi-layered-primeslice-complexity-active-depth}
在上述有限纤维链设定下，定义活跃分支层集合
\[
J:=
\left\{
\,0\le j\le k-1:\ \exists\,y\in X_{j+1}\ \text{使}\ \#\pi_j^{-1}(y)\ge 2
\right\}.
\]
在分层 prime-slice 外置语义下，定义其复杂度
\[
\mathfrak p(\pi_\bullet)
\]
为所有单射编码
\[
\Xi(x_0)=\bigl(x_k,H(x_0)\bigr),
\qquad
H(x_0)=\prod_{j=0}^{k-1}h_j(x_0),
\]
中非平凡 prime slices 的最小可能数目，其中每个 \(h_j(x_0)\) 的全部素因子均限制在某个预先指定的层片 \(\mathcal P_j\)。则有精确公式
\[
\boxed{\;
\mathfrak p(\pi_\bullet)=|J|.\;}
\]
特别地，若每一层都存在真正分支，即 \(J=\{0,\dots,k-1\}\)，则
\[
\boxed{\;
\mathfrak p(\pi_\bullet)=k.\;}
\]
\end{theorem}
\leanverified{paper\_xi\_layered\_primeslice\_complexity\_active\_depth}
\begin{proof}
这正是定理 \ref{thm:xi-prime-slice-nontrivial-layer-exact-minimality} 的内容。该定理证明：任一单射分层 prime-slice 编码的非平凡片数都不小于 \(|J|\)，而且存在构造恰好在每个活跃层使用一枚有效 prime slice、在非活跃层取平凡因子 \(1\) 的单射实现。故最小值正是 \(|J|\)。证毕。
\end{proof}

\begin{theorem}[Smith 信息亏损谱的离散曲率完全刻画]\label{thm:xi-smith-loss-spectrum-discrete-curvature-complete-characterization}
设 \(f:\ZZ_{\ge 0}\to \ZZ_{\ge 0}\) 满足 \(f(0)=0\)，并定义一阶差分
\[
\Delta(k):=f(k)-f(k-1)\qquad(k\ge 1).
\]
则下列两条等价：
\begin{enumerate}
\item 存在某个有限多重集 \(E\subset \ZZ_{\ge 1}\) 使
\[
f(k)=\sum_{e\in E}\min(k,e)
\qquad(k\ge 0);
\]
\item 序列 \((\Delta(k))_{k\ge 1}\) 非负、单调不增且最终为零。
\end{enumerate}
在这些条件成立时，\(E\) 由离散曲率唯一恢复：对每个 \(t\ge 1\)，其重数满足
\[
\boxed{\;
\#\{e\in E:\ e=t\}
=
\Delta(t)-\Delta(t+1)
=
2f(t)-f(t-1)-f(t+1).\;}
\]
因此，任何 \(\,p\)-进 Smith 信息亏损谱都可被视为一个有限 Young 边界函数，而确切的初等因子多重集正是该边界的离散曲率。
\leanverified{paper\_xi\_smith\_loss\_spectrum\_discrete\_curvature\_complete\_characterization}
\end{theorem}
\begin{proof}
定理 \ref{thm:xi-smith-padic-loss-spectrum-classification} 已证明：函数 \(f\) 可以写成
\[
f(k)=\sum_{e\ge 1}m_e\min(k,e)
\]
当且仅当其前向差分 \(\Delta(k)\) 非负、非增且有限支撑；并且此时重数 \(m_e\) 唯一确定。推论 \ref{cor:xi-smith-second-discrete-curvature-recovery} 进一步给出
\[
m_t=\Delta(t)-\Delta(t+1)=2f(t)-f(t-1)-f(t+1).
\]
将 \(m_t\) 解释为 \(E\) 中取值 \(t\) 的重数，即得全部结论。证毕。
\end{proof}

\begin{corollary}[Smith 核谱函数的离散凹性、平台截断与二阶差分反演]\label{cor:xi-smith-kernel-spectrum-discrete-concavity-plateau-inversion}
\leanverified{paper\_xi\_smith\_kernel\_spectrum\_discrete\_concavity\_plateau\_inversion}
设 \(A\in M_{m\times n}(\ZZ)\) 满足 \(\operatorname{rank}_{\QQ}(A)=m\)，其 Smith 不变量为
\[
d_1\mid d_2\mid\cdots\mid d_m.
\]
固定素数 \(p\)，并记
\[
K_p(k):=\#\ker(A_{p^k})
\qquad(k\ge 1),
\]
\[
f_p(k):=\log_p K_p(k)-k(n-m)
\qquad(k\ge 1),
\qquad
f_p(0):=0.
\]
再记
\[
e_i(p):=v_p(d_i),
\qquad
E_p:=\max_{1\le i\le m}e_i(p),
\qquad
m_{p,\ell}:=\#\{i:\ e_i(p)=\ell\}
\quad(\ell\ge 1).
\]
则成立：
\begin{enumerate}
  \item \(f_p\) 是单调不减的离散凹函数，即
  \[
  f_p(k+1)-f_p(k)\ge f_p(k+2)-f_p(k+1)
  \qquad(k\ge 0);
  \]
  \item 当 \(k\ge E_p\) 时，\(f_p\) 进入平台：
  \[
  f_p(k)=\sum_{i=1}^{m}e_i(p);
  \]
  \item 对每个 \(\ell\ge 1\)，精确重数由二阶差分恢复：
  \[
  m_{p,\ell}
  =
  2f_p(\ell)-f_p(\ell-1)-f_p(\ell+1).
  \]
\end{enumerate}
\end{corollary}
\begin{proof}
由定理 \ref{thm:killo-smith-loss-spectrum}，
\[
f_p(k)=\sum_{i=1}^{m}\min\bigl(k,e_i(p)\bigr)
\qquad(k\ge 1),
\]
并且
\[
\Delta_p(k):=f_p(k)-f_p(k-1)=\#\{i:\ e_i(p)\ge k\}
\qquad(k\ge 1).
\]
因此 \(\Delta_p(k)\) 关于 \(k\) 非负且单调不增，故 \(f_p\) 单调不减且离散凹，这证明了第 \(1\) 条。

若 \(k\ge E_p\)，则 \(\min(k,e_i(p))=e_i(p)\) 对所有 \(i\) 都成立，故
\[
f_p(k)=\sum_{i=1}^{m}e_i(p),
\]
得到第 \(2\) 条。

最后，由上式直接有
\[
m_{p,\ell}
=
\#\{i:\ e_i(p)\ge \ell\}-\#\{i:\ e_i(p)\ge \ell+1\}
=
\Delta_p(\ell)-\Delta_p(\ell+1),
\]
从而
\[
m_{p,\ell}
=
\bigl(f_p(\ell)-f_p(\ell-1)\bigr)-\bigl(f_p(\ell+1)-f_p(\ell)\bigr)
=
2f_p(\ell)-f_p(\ell-1)-f_p(\ell+1).
\]
证毕。
\end{proof}

\begin{theorem}[模数信息亏损是 divisibility 格上的精确 valuation，且具有有理 Euler 因子]\label{thm:xi-kernel-loss-divisibility-valuation-rational-euler}
\leanverified{paper\_xi\_kernel\_loss\_divisibility\_valuation\_rational\_euler}
设 \(A\in M_{m\times n}(\ZZ)\) 满足 \(\operatorname{rank}_{\QQ}(A)=m\)，其 Smith 不变量为
\[
d_1\mid d_2\mid\cdots\mid d_m.
\]
对每个 \(b\ge 1\)，记
\[
A_b:(\ZZ/b\ZZ)^n\to(\ZZ/b\ZZ)^m
\]
为 \(A\) 的模 \(b\) 诱导映射。对任意正整数 \(b\)，定义归一化信息亏损
\[
L_A(b):=
\log \#\ker(A_b)-(n-m)\log b
=
\sum_{i=1}^{m}\log \gcd(d_i,b).
\]
则成立：
\begin{enumerate}
\item 对任意 \(b,c\ge 1\)，有精确 valuation 恒等式
\[
\boxed{\;
L_A(\operatorname{lcm}(b,c))+L_A(\gcd(b,c))
=
L_A(b)+L_A(c).\;}
\]
\item 定义 Dirichlet 生成函数
\[
\Xi_A(s):=\sum_{b\ge 1}\frac{\#\ker(A_b)}{b^{\,n-m+s}}
\qquad(\Re s>1).
\]
若记 \(e_i(p):=v_p(d_i)\)、\(E_p:=\max_i e_i(p)\)、\(\sigma_p:=\sum_i e_i(p)\)，则
\[
\boxed{\;
\Xi_A(s)=\prod_{p}\Xi_{A,p}(s),\qquad
\Xi_{A,p}(s)=\sum_{k=0}^{\infty}p^{-ks+\sum_{i=1}^{m}\min(k,e_i(p))}.\;}
\]
并且每个局部因子都是 \(p^{-s}\) 的有理函数，更精确地，
\[
\boxed{\;
\Xi_{A,p}(s)
=
\sum_{k=0}^{E_p-1}p^{-ks+\sum_i\min(k,e_i(p))}
\;+\;
\frac{p^{-E_p s+\sigma_p}}{1-p^{-s}}.\;}
\]
\end{enumerate}
\end{theorem}
\begin{proof}
由推论 \ref{cor:killo-kernel-size-general-modulus}，
\[
L_A(b)=\sum_{p}\left(\sum_{i=1}^{m}\min(v_p(b),e_i(p))\right)\log p.
\]
故要证第一式，只需逐素数验证
\[
\sum_i \min(\max\{u,v\},e_i)+\sum_i \min(\min\{u,v\},e_i)
=
\sum_i \min(u,e_i)+\sum_i \min(v,e_i),
\]
其中 \(u=v_p(b)\)、\(v=v_p(c)\)。而对每个固定 \(e\) 都有
\[
\min(\max\{u,v\},e)+\min(\min\{u,v\},e)=\min(u,e)+\min(v,e),
\]
逐项求和即得 valuation 恒等式。

第二式同样由推论 \ref{cor:killo-kernel-size-general-modulus} 与定理 \ref{thm:xi-kernel-dirichlet-zeta-finite-euler-correction} 直接推出。对 \(b=p^k\)，有
\[
\frac{\#\ker(A_{p^k})}{p^{k(n-m)}}
=
p^{\sum_i\min(k,e_i(p))},
\]
故局部因子为
\[
\Xi_{A,p}(s)=\sum_{k\ge 0}p^{-ks+\sum_i\min(k,e_i(p))}.
\]
当 \(k\ge E_p\) 时，\(\min(k,e_i(p))=e_i(p)\) 对所有 \(i\) 恒定，尾部即成为几何级数
\[
\sum_{k\ge E_p}p^{-ks+\sigma_p}
=
\frac{p^{-E_p s+\sigma_p}}{1-p^{-s}},
\]
于是得到所示有理表达。证毕。
\end{proof}

\begin{corollary}[模核大小在整除格上满足精确乘法恒等式]\label{cor:xi-kernel-size-gcd-lcm-exact-multiplicativity}
\leanverified{paper\_xi\_kernel\_size\_gcd\_lcm\_exact\_multiplicativity}
沿用定理 \ref{thm:xi-kernel-loss-divisibility-valuation-rational-euler} 的记号。对任意 \(b,c\ge 1\)，都有
\[
\#\ker(A_{\gcd(b,c)})\cdot \#\ker(A_{\operatorname{lcm}(b,c)})
=
\#\ker(A_b)\cdot \#\ker(A_c).
\]
等价地，模核大小函数在除去自由因子 \(b^{\,n-m}\) 之后，正是正整数整除格上的一个严格 valuation。
\end{corollary}
\begin{proof}
由定理 \ref{thm:xi-kernel-loss-divisibility-valuation-rational-euler}，
\[
L_A(\operatorname{lcm}(b,c))+L_A(\gcd(b,c))=L_A(b)+L_A(c).
\]
再由定义
\[
L_A(t)=\log\#\ker(A_t)-(n-m)\log t
\qquad(t\ge 1),
\]
以及恒等式
\[
\log\gcd(b,c)+\log\operatorname{lcm}(b,c)=\log b+\log c,
\]
可得
\[
\log\#\ker(A_{\gcd(b,c)})+\log\#\ker(A_{\operatorname{lcm}(b,c)})
=
\log\#\ker(A_b)+\log\#\ker(A_c).
\]
两边取指数即得所示精确乘法关系。证毕。
\end{proof}

\begin{theorem}[模核 Dirichlet 级数的 Euler--rational 形态与 Smith 完备性]\label{thm:xi-kernel-dirichlet-euler-rational-smith-completeness}
设 \(A\in M_{m\times n}(\ZZ)\) 满足 \(\operatorname{rank}_{\QQ}(A)=m\)，其 Smith 不变量为
\[
d_1\mid d_2\mid\cdots\mid d_m,\qquad d_i>0.
\]
对每个正整数 \(b\)，记
\[
A_b:(\ZZ/b\ZZ)^n\to(\ZZ/b\ZZ)^m
\]
为模 \(b\) 的诱导映射，并定义 Dirichlet 级数
\[
\mathcal K_A(s):=\sum_{b\ge 1}\frac{\#\ker(A_b)}{b^s}.
\]
则成立：
\[
\#\ker(A_b)=b^{\,n-m}\prod_{i=1}^{m}\gcd(d_i,b),
\]
从而 \(b\mapsto \#\ker(A_b)\) 为积性函数，故
\[
\mathcal K_A(s)=\prod_p \mathcal K_{A,p}(s),
\]
其中局部因子为
\[
\mathcal K_{A,p}(s)
=
\sum_{k\ge 0}p^{-k(s-(n-m))+f_p(k)},
\qquad
f_p(k):=\sum_{i=1}^{m}\min\bigl(e_i(p),k\bigr),
\]
并记
\[
e_i(p):=v_p(d_i),\qquad
E_p:=\max_i e_i(p),\qquad
\sigma_p:=\sum_{i=1}^{m}e_i(p).
\]
则更精确地，
\[
\mathcal K_{A,p}(s)
=
\sum_{k=0}^{E_p-1}p^{-k(s-(n-m))+f_p(k)}
\;+\;
\frac{p^{-E_p(s-(n-m))+\sigma_p}}{1-p^{-(s-(n-m))}},
\]
因此每个局部因子都是 \(p^{-s}\) 的有理函数。

若再记
\[
\Delta_A:=\prod_{i=1}^{m}d_i,
\]
则
\[
\mathcal K_A(s)
=
\zeta\!\bigl(s-(n-m)\bigr)
\prod_{p\mid \Delta_A}
\Bigl(1-p^{-(s-(n-m))}\Bigr)\mathcal K_{A,p}(s).
\]
因而 \(\mathcal K_A(s)\) 的收敛横坐标恰为
\[
\sigma_c=n-m+1,
\]
并在该点具有唯一简单极点。更进一步，在固定行秩 \(m\) 的口径下，\(\mathcal K_A(s)\) 唯一决定全部 Smith 不变量 \((d_1,\dots,d_m)\)。
\end{theorem}
\begin{proof}
由推论 \ref{cor:killo-kernel-size-general-modulus}，
\[
\#\ker(A_b)=b^{\,n-m}\prod_{i=1}^{m}\gcd(d_i,b).
\]
右端关于 \(b\) 为积性函数，故 \(\mathcal K_A\) 具有 Euler 乘积分解。对 \(b=p^k\)，再由同一推论与定理 \ref{thm:killo-smith-loss-spectrum}，
\[
\#\ker(A_{p^k})
=
p^{k(n-m)}\prod_{i=1}^{m}p^{\min(e_i(p),k)}
=
p^{k(n-m)+f_p(k)}.
\]
于是
\[
\mathcal K_{A,p}(s)
=
\sum_{k\ge 0}\frac{\#\ker(A_{p^k})}{p^{ks}}
=
\sum_{k\ge 0}p^{-k(s-(n-m))+f_p(k)}.
\]
当 \(k\ge E_p\) 时，\(\min(e_i(p),k)=e_i(p)\) 对所有 \(i\) 都成立，故 \(f_p(k)=\sigma_p\)，尾部因而化为几何级数
\[
\sum_{k\ge E_p}p^{-k(s-(n-m))+\sigma_p}
=
\frac{p^{-E_p(s-(n-m))+\sigma_p}}{1-p^{-(s-(n-m))}}.
\]
这就得到局部因子的有理表达式。

若 \(p\nmid \Delta_A\)，则全部 \(e_i(p)=0\)，从而
\[
\mathcal K_{A,p}(s)=\sum_{k\ge 0}p^{-k(s-(n-m))}
=
\frac{1}{1-p^{-(s-(n-m))}}.
\]
把这些好素数因子提出，即得
\[
\mathcal K_A(s)
=
\zeta\!\bigl(s-(n-m)\bigr)
\prod_{p\mid \Delta_A}
\Bigl(1-p^{-(s-(n-m))}\Bigr)\mathcal K_{A,p}(s).
\]
右侧修正项是有限 Euler 乘积，因此在 \(s=n-m+1\) 邻域内解析且非零；于是 \(\mathcal K_A\) 的收敛横坐标与 \(\zeta(s-(n-m))\) 相同，恰为 \(n-m+1\)，并在该点具有唯一简单极点。

最后证明 Smith 完备性。Dirichlet 级数 \(\mathcal K_A(s)\) 唯一决定全部系数
\[
a_b:=\#\ker(A_b).
\]
极点位置给出 \(n-m=\sigma_c-1\)。因此对每个素数 \(p\) 与每个 \(k\ge 1\)，可由
\[
a_{p^k}=p^{k(n-m)+f_p(k)}
\]
恢复
\[
f_p(k)=v_p(a_{p^k})-k(n-m).
\]
再由推论 \ref{cor:xi-smith-second-discrete-curvature-recovery}，
\[
\#\{i:\ e_i(p)=k\}
=
2f_p(k)-f_p(k-1)-f_p(k+1)
\qquad(k\ge 1),
\]
故每个素数 \(p\) 上的赋值多重集 \(\{e_i(p)\}_{i=1}^{m}\) 都被唯一确定。由于行秩 \(m\) 已固定，零赋值的个数亦随之确定。记
\[
0\le e_1^{\uparrow}(p)\le \cdots \le e_m^{\uparrow}(p)
\]
为该多重集的非降序排列，则逐项合并便唯一恢复整个 Smith 链
\[
d_i=\prod_p p^{e_i^{\uparrow}(p)},
\qquad(1\le i\le m).
\]
证毕。
\end{proof}

\begin{theorem}[整数轴对齐椭圆半群的结构层同态半圆维发散，而实现层仅为 \texorpdfstring{$1/2$}{1/2}]\label{thm:xi-integer-ellipse-half-circle-dimension-gap}
沿用定理 \ref{thm:killo-ellipse-diagonal-prime-register-equivalence} 的记号，令
\[
\mathsf E_{\NN}
:=
\left\{
E_{a,b}:\ \frac{x^2}{a^2}+\frac{y^2}{b^2}=1,\ a,b\in\NN_{>0}
\right\},
\qquad
E_{a,b}\odot E_{c,d}:=E_{ac,bd}.
\]
则
\[
\boxed{\;
\cdim^{\mathrm{hom}}_{+}(\mathsf E_{\NN},\odot)=+\infty,
\qquad
\cdim^{\mathrm{impl}}(\mathsf E_{\NN})=\frac12.\;}
\]
等价地，不存在任何有限加法寄存器系统 \((\NN^k,+)\) 能对整个整数轴对齐椭圆半群给出严格同态的完整编码；但作为纯集合，\(\mathsf E_{\NN}\) 仍只具有可数无限集的实现层复杂度。
\end{theorem}
\begin{proof}
对任意有限素数集 \(S\subset\mathbb P\)，定义子幺半群
\[
\mathsf E_S:=\{E_{a,1}:a\in \NN_S^\times\}\subset \mathsf E_{\NN}.
\]
映射
\[
\NN_S^\times\longrightarrow \mathsf E_S,\qquad a\longmapsto E_{a,1}
\]
是幺半群同构，故 \(\mathsf E_{\NN}\) 对任意 \(|S|=N\) 都包含一份 \(\NN_S^\times\) 的同态拷贝。于是由推论 \ref{cor:xi-unbounded-prime-directions-force-infinite-half-circle-dimension}，
\[
\cdim^{\mathrm{hom}}_{+}(\mathsf E_{\NN},\odot)=+\infty.
\]

另一方面，\(\mathsf E_{\NN}\) 与 \(\NN_{>0}\times \NN_{>0}\) 等势，因而是可数无限集。故存在集合单射
\[
\mathsf E_{\NN}\hookrightarrow \NN,
\]
从而
\[
\cdim^{\mathrm{impl}}(\mathsf E_{\NN})\le \frac12.
\]
又实现层半圆维为 \(0\) 当且仅当集合有限（见命题 \ref{prop:killo-implementation-vs-structural-half-circle-dimension} 的定义口径），故
\[
\cdim^{\mathrm{impl}}(\mathsf E_{\NN})=\frac12.
\]
证毕。
\end{proof}

\begin{corollary}[整数轴对齐椭圆半群不存在任何有限加法账本的忠实同态实现]\label{cor:xi-integer-ellipse-no-faithful-finite-additive-ledger}
沿用定理 \ref{thm:xi-integer-ellipse-half-circle-dimension-gap} 的记号，并记
\[
\mathsf E_x:=\{E_{a,1}:a\in\NN_{>0}\}\subset \mathsf E_{\NN}.
\]
则不存在任何有限整数 \(k\) 与单射幺半群同态
\[
\Phi:(\mathsf E_{\NN},\odot)\hookrightarrow (\NN^k,+).
\]
另一方面，\(\mathsf E_{\NN}\) 作为可数集合仍可单射进入 \(\NN\)。因此，对整数轴对齐椭圆半群而言，结构层的忠实同态实现与实现层的可寻址集合压缩发生严格分离。
\end{corollary}
\begin{proof}
映射
\[
\NN_{>0}\longrightarrow \mathsf E_x,\qquad
a\longmapsto E_{a,1}
\]
满足
\[
E_{a,1}\odot E_{c,1}=E_{ac,1},
\]
故给出幺半群同构
\[
(\NN_{>0},\times)\cong (\mathsf E_x,\odot).
\]
若存在某个有限 \(k\) 与单射幺半群同态
\[
\Phi:(\mathsf E_{\NN},\odot)\hookrightarrow (\NN^k,+),
\]
则其限制
\[
\Phi|_{\mathsf E_x}:(\NN_{>0},\times)\hookrightarrow (\NN^k,+)
\]
仍为单射幺半群同态。这与定理 \ref{thm:killo-prime-freedom-non-finitizability} 矛盾，因为 \((\NN^k,+)\) 是有限生成、交换且可消去的幺半群。故这样的 \(\Phi\) 不存在。

另一方面，\(\mathsf E_{\NN}\) 与 \(\NN_{>0}\times \NN_{>0}\) 等势，故是可数无限集，从而存在集合单射
\[
\mathsf E_{\NN}\hookrightarrow \NN.
\]
这就给出了所述的严格分离。证毕。
\end{proof}

\noindent
上述这些结果把附录 \(H\) 的若干局部构造压缩为一条统一刚性链：有限深度 branch register 是 prime-slice 外置的前缀模型，prime-slice 的最小复杂度由活跃分支层精确锁定；进入模数投影后，全部信息亏损不仅被 Smith 曲率谱完全分类，而且在整除格上形成精确 valuation 与原始模核乘法恒等式；最终，这条素数寄存器刚性可以直接投射到整数椭圆半群这样的显式几何对象，迫使其在结构层上表现出无限的加法寄存器秩，而实现层仍停留在可数可寻址的最低非零复杂度。

\paragraph{自然扩张寄存器的逆极限规范性、Smith 局部亏损谱的自由化与 Euler 完备不变量}
在有限深度 prime-slice 前缀压缩、Smith 离散曲率恢复与整数椭圆原子分解之外，附录 \(H\) 中的 branch-register 语义还可继续压缩为逆极限、有限步过去的极小完备前缀、局部亏损谱的自由原子结构以及全模数核计数的 Smith 饱和阈值层的若干闭式结论。

\begin{theorem}[自然扩张分支寄存器系统是有限深度截断的严格逆极限]\label{thm:xi-branch-register-strict-inverse-limit}
\leanverified{paper\_xi\_branch\_register\_strict\_inverse\_limit}
沿用定理 \ref{thm:killo-natural-extension-branch-register} 的记号。对每个 \(k\ge 1\)，记
\[
X^{\leftarrow,k}
:=
\{(x_0,x_{-1},\dots,x_{-k})\in X^{k+1}:T(x_{-j})=x_{-j+1}\ \ (1\le j\le k)\},
\]
\[
\mathcal C_k(x_0,x_{-1},\dots,x_{-k})
:=
\bigl(x_0,(\beta(x_{-1}),\dots,\beta(x_{-k}))\bigr),
\]
\[
\widehat X_k^{\mathrm{br}}
:=
\mathcal C_k(X^{\leftarrow,k})
\subset X\times \NN^k.
\]
定义截断投影
\[
\pi_{k+1,k}^{\mathrm{br}}:\widehat X_{k+1}^{\mathrm{br}}\to \widehat X_k^{\mathrm{br}},
\qquad
\pi_{k+1,k}^{\mathrm{br}}(x,a_1,\dots,a_{k+1})=(x,a_1,\dots,a_k),
\]
以及有限深度更新
\[
\widehat T_k^{\mathrm{br}}(x,a_1,\dots,a_k)
:=
\begin{cases}
(T(x),\beta(x)),& k=1,\\
(T(x),\beta(x),a_1,\dots,a_{k-1}),& k\ge 2.
\end{cases}
\]
则 \((\widehat X_k^{\mathrm{br}},\pi_{k+1,k}^{\mathrm{br}})_{k\ge 1}\) 构成投影系统，并且存在规范同构
\[
\widehat X
\cong
\varprojlim_{k}\widehat X_k^{\mathrm{br}}
\]
把定理 \ref{thm:killo-natural-extension-branch-register} 中的 \(\widehat T\) 识别为
\[
\widehat T
=
\varprojlim_{k}\widehat T_k^{\mathrm{br}}.
\]
\end{theorem}
\begin{proof}
对每个 \(k\ge 1\)，定义截断映射
\[
\rho_k:\widehat X\to \widehat X_k^{\mathrm{br}},
\qquad
\rho_k(x,(a_1,a_2,\dots))=(x,a_1,\dots,a_k).
\]
由定理 \ref{thm:killo-natural-extension-branch-register}，\(\widehat X=\mathcal C(X^{\leftarrow})\)。若
\[
\mathcal C(x_0,x_{-1},x_{-2},\dots)=(x_0,(a_1,a_2,\dots)),
\]
则
\[
\rho_k(x_0,(a_1,a_2,\dots))
=
\mathcal C_k(x_0,x_{-1},\dots,x_{-k})\in \widehat X_k^{\mathrm{br}},
\]
故 \(\rho_k\) 良定义，且显然满足
\[
\pi_{k+1,k}^{\mathrm{br}}\circ \rho_{k+1}=\rho_k.
\]
于是得到自然映射
\[
\Theta:\widehat X\to \varprojlim_k \widehat X_k^{\mathrm{br}},
\qquad
\Theta(z):=(\rho_k(z))_{k\ge 1}.
\]

\(\Theta\) 的单射性来自前缀决定性：若 \(\Theta(z)=\Theta(z')\)，则二者全部有限前缀一致，从而其 \(\NN^\NN\) 坐标完全一致，故 \(z=z'\)。

再证满射。任取相容族
\[
(\xi_k)_{k\ge 1}\in \varprojlim_k \widehat X_k^{\mathrm{br}},
\qquad
\xi_k=(x^{(k)},a_1^{(k)},\dots,a_k^{(k)}).
\]
由相容性
\[
\pi_{k+1,k}^{\mathrm{br}}(\xi_{k+1})=\xi_k
\]
立得 \(x^{(k)}\) 对 \(k\) 稳定，且对每个固定 \(j\ge 1\)，当 \(k\ge j\) 时 \(a_j^{(k)}\) 亦稳定。记其稳定值为
\[
x_0,\qquad a_j\in \NN\ \ (j\ge 1).
\]
于是得到一个候选点
\[
z:=(x_0,(a_1,a_2,\dots))\in X\times \NN^\NN.
\]
另一方面，对每个 \(k\)，由 \(\xi_k\in \widehat X_k^{\mathrm{br}}\) 可知存在唯一有限历史
\[
(x_0,x_{-1}^{(k)},\dots,x_{-k}^{(k)})\in X^{\leftarrow,k}
\]
满足
\[
\mathcal C_k(x_0,x_{-1}^{(k)},\dots,x_{-k}^{(k)})
=
(x_0,a_1,\dots,a_k).
\]
由 \(\mathcal C_k\) 的递推反演，长度 \(k+1\) 与长度 \(k\) 的历史在前 \(k\) 层必然一致；因此这些有限历史彼此兼容，定义出唯一无限前史
\[
(x_0,x_{-1},x_{-2},\dots)\in X^{\leftarrow}.
\]
按构造有
\[
\mathcal C(x_0,x_{-1},x_{-2},\dots)=z,
\]
故 \(z\in \widehat X\) 且 \(\Theta(z)=(\xi_k)_k\)。因此 \(\Theta\) 为双射。

最后，对任意 \((x,(a_1,a_2,\dots))\in \widehat X\)，由定理 \ref{thm:killo-natural-extension-branch-register} 的右移插入公式，
\[
\widehat T(x,(a_1,a_2,\dots))=(T(x),(\beta(x),a_1,a_2,\dots)).
\]
截断到前 \(k\) 项便得到
\[
\rho_k\bigl(\widehat T(x,(a_1,a_2,\dots))\bigr)
=
\widehat T_k^{\mathrm{br}}(x,a_1,\dots,a_k).
\]
故 \(\Theta\) 将 \(\widehat T\) 识别为逆极限动力学 \(\varprojlim_k \widehat T_k^{\mathrm{br}}\)。证毕。
\end{proof}

\begin{theorem}[分支指标寄存器编码的规范共轭不变性]\label{thm:xi-branch-register-coding-conjugacy-invariance}
仍沿用定理 \ref{thm:killo-natural-extension-branch-register} 的设定。若另取一组纤维单射
\[
\iota_y':T^{-1}(y)\hookrightarrow \NN,
\qquad
\beta'(x):=\iota'_{T(x)}(x),
\]
并由此得到编码
\[
\mathcal C':X^{\leftarrow}\to X\times\NN^\NN,
\qquad
\widehat X':=\mathcal C'(X^{\leftarrow}),
\qquad
\widehat T':=\mathcal C'\circ \sigma\circ (\mathcal C')^{-1},
\]
则存在规范双射
\[
\Gamma:\widehat X\to \widehat X'
\]
满足
\[
\Gamma\circ \widehat T=\widehat T'\circ \Gamma.
\]
更具体地，若
\[
\mathcal C^{-1}(x_0,(a_1,a_2,\dots))=(x_0,x_{-1},x_{-2},\dots),
\]
则
\[
\Gamma(x_0,(a_1,a_2,\dots))
=
\bigl(x_0,(\beta'(x_{-1}),\beta'(x_{-2}),\dots)\bigr).
\]
等价地，改变纤维枚举 \((\iota_y)_y\) 只在各条纤维内部实施局部字母重命名，而不改变自然扩张的动力学共轭类。
\end{theorem}
\begin{proof}
定义
\[
\Gamma:=\mathcal C'\circ \mathcal C^{-1}:\widehat X\to \widehat X'.
\]
由于 \(\mathcal C,\mathcal C'\) 都是各自像空间上的双射，\(\Gamma\) 亦为双射。并且
\[
\Gamma\circ \widehat T
=
\mathcal C'\circ \mathcal C^{-1}\circ \mathcal C\circ \sigma\circ \mathcal C^{-1}
=
\mathcal C'\circ \sigma\circ \mathcal C^{-1}
=
\widehat T'\circ \mathcal C'\circ \mathcal C^{-1}
=
\widehat T'\circ \Gamma.
\]
这就得到所需共轭关系。

显式公式只需把 \((x_0,(a_1,a_2,\dots))\) 先经 \(\mathcal C^{-1}\) 恢复为唯一前史
\[
(x_0,x_{-1},x_{-2},\dots)\in X^{\leftarrow},
\]
再送入 \(\mathcal C'\)。因此
\[
\Gamma(x_0,(a_1,a_2,\dots))
=
\mathcal C'(x_0,x_{-1},x_{-2},\dots)
=
\bigl(x_0,(\beta'(x_{-1}),\beta'(x_{-2}),\dots)\bigr).
\]
上述映射完全由两套编码 \(\mathcal C,\mathcal C'\) 决定，因而给出了所需的规范共轭。证毕。
\end{proof}

\leanverified{paper\_xi\_branch\_register\_finite\_past\_minimal\_complete\_prefix}
\begin{theorem}[有限纤维满射的 \texorpdfstring{$n$}{n}-步过去由首端点与前 \texorpdfstring{$n$}{n} 个分支指标完全且唯一决定]\label{thm:xi-branch-register-finite-past-minimal-complete-prefix}
仍沿用定理 \ref{thm:killo-natural-extension-branch-register} 的设定。对每个 \(n\ge 1\)，记
\[
X^{\leftarrow,n}
:=
\{(x_0,x_{-1},\dots,x_{-n})\in X^{n+1}:T(x_{-j})=x_{-j+1}\ \ (1\le j\le n)\},
\]
\[
\mathcal C_n(x_0,x_{-1},\dots,x_{-n})
:=
\bigl(x_0,\beta(x_{-1}),\dots,\beta(x_{-n})\bigr),
\]
\[
\widehat X_n^{\mathrm{br}}
:=
\mathcal C_n(X^{\leftarrow,n})\subset X\times \NN^n.
\]
则 \(\mathcal C_n:X^{\leftarrow,n}\to \widehat X_n^{\mathrm{br}}\) 为双射。等价地，对任意 \(x\in X\)，映射
\[
T^{-n}(x)\longrightarrow \NN^n,
\qquad
y\longmapsto \bigl(\beta(T^{n-1}y),\beta(T^{n-2}y),\dots,\beta(y)\bigr)
\]
在其像上为双射。

特别地，\((x_0,\beta(x_{-1}),\dots,\beta(x_{-n}))\) 构成 \(n\)-步过去的极小完备前缀统计量：若
\[
F:X^{\leftarrow,n}\to Y
\]
满足存在映射 \(G:\widehat X_n^{\mathrm{br}}\to Y\) 使
\[
F=G\circ \mathcal C_n,
\]
并且 \(F\) 在 \(X^{\leftarrow,n}\) 上单射，则 \(G\) 必在 \(\widehat X_n^{\mathrm{br}}\) 上单射。因而不存在任何严格弱于 \(\mathcal C_n\) 而仍能在全部 \(n\)-步前史上保持单射的信息压缩。
\end{theorem}
\begin{proof}
由定理 \ref{thm:killo-natural-extension-branch-register}，无限分支寄存器编码
\[
\mathcal C(x_0,x_{-1},x_{-2},\dots)=\bigl(x_0,(\beta(x_{-1}),\beta(x_{-2}),\dots)\bigr)
\]
可按递推公式唯一反演：
\[
x_{-1}=\iota_{x_0}^{-1}\!\bigl(\beta(x_{-1})\bigr),\qquad
x_{-2}=\iota_{x_{-1}}^{-1}\!\bigl(\beta(x_{-2})\bigr),\qquad \dots
\]
故一旦给定
\[
\bigl(x_0,a_1,\dots,a_n\bigr)\in \widehat X_n^{\mathrm{br}},
\]
就可逐步唯一恢复
\[
x_{-1}=\iota_{x_0}^{-1}(a_1),\qquad
x_{-2}=\iota_{x_{-1}}^{-1}(a_2),\qquad \dots,\qquad
x_{-n}=\iota_{x_{-(n-1)}}^{-1}(a_n).
\]
因此 \(\mathcal C_n\) 在其像上可逆，从而为双射。

对固定 \(x\in X\)，映射
\[
T^{-n}(x)\longrightarrow X^{\leftarrow,n},
\qquad
y\longmapsto \bigl(x,T^{n-1}y,T^{n-2}y,\dots,y\bigr)
\]
显然为双射；再与 \(\mathcal C_n\) 复合，便得到所示从 \(T^{-n}(x)\) 到 \(\NN^n\) 的像上双射。

最后证明极小性。若 \(F=G\circ \mathcal C_n\) 且 \(F\) 单射，取任意
\[
u,v\in \widehat X_n^{\mathrm{br}}
\]
满足 \(G(u)=G(v)\)。由于 \(\mathcal C_n\) 对其像为双射，可取唯一
\[
z,z'\in X^{\leftarrow,n}
\]
使 \(\mathcal C_n(z)=u\)、\(\mathcal C_n(z')=v\)。于是
\[
F(z)=G(\mathcal C_n(z))=G(u)=G(v)=G(\mathcal C_n(z'))=F(z').
\]
由 \(F\) 的单射性得 \(z=z'\)，再由 \(\mathcal C_n\) 的单射性得 \(u=v\)。故 \(G\) 在 \(\widehat X_n^{\mathrm{br}}\) 上亦为单射，这正是所述极小完备性。证毕。
\end{proof}

\begin{corollary}[最坏情形下 \texorpdfstring{$n$}{n}-步过去的 prime-slice 可寻址维数恰等于深度]\label{cor:xi-branch-register-depth-primeslice-addressability}
设 \(T:X\twoheadrightarrow X\) 为非单射有限纤维满射，并固定 \(n\ge 1\)。考虑所有单射分层乘法账本
\[
\Psi_n(y)=\bigl(T^n(y),H_n(y)\bigr),
\qquad
H_n(y)=\prod_{j=0}^{n-1}h_j(y),
\]
其中每个 \(h_j(y)\) 的全部素因子均限制在预先指定且两两不交的 prime slice \(\mathcal P_j\)。定义
\[
\operatorname{psupp}(\Psi_n):=\#\{\,0\le j\le n-1:\ h_j\not\equiv 1\,\}.
\]
则
\[
\inf_{\Psi_n\ \mathrm{injective}}\operatorname{psupp}(\Psi_n)=n.
\]
因此，自然扩张的深度不仅是分支寄存器前缀的最小完备长度，也正是最坏情形下不可压缩的 prime-slice 账本维数。
\end{corollary}
\begin{proof}
对长度 \(n\) 的前史塔取
\[
X_0=X_1=\cdots=X_n=X,
\qquad
\pi_0=\pi_1=\cdots=\pi_{n-1}=T.
\]
由于 \(T\) 非单射，存在某个 \(y\in X\) 使 \(\#T^{-1}(y)\ge 2\)，故该塔的每一层都满足“存在真正分支点”的最坏情形假设。于是定理 \ref{thm:xi-layered-prime-register-sharp-natural-extension} 立即给出：任意单射分层 prime-register 外置都必须在全部 \(n\) 层上保留非平凡的 prime slice，而且存在恰好每层写入一枚有效素数的单射实现。故上式成立。

再由定理 \ref{thm:xi-branch-register-finite-past-minimal-complete-prefix}，\(\mathcal C_n\) 已经给出 \(n\)-步过去的极小完备前缀统计量。因此同一深度 \(n\) 同时刻画了加性分支寄存器的最小完备长度与乘法 prime-slice 账本的最小可寻址维数。证毕。
\end{proof}

\leanverified{paper\_xi\_smith\_local\_loss\_profiles\_free\_commutative\_monoid}
\begin{theorem}[\texorpdfstring{$p$}{p}-局部信息亏损剖面的自由交换幺半群结构]\label{thm:xi-smith-local-loss-profiles-free-commutative-monoid}
固定素数 \(p\)。令 \(\mathfrak F_p\) 为所有由有限多重集 \(E\subset \ZZ_{>0}\) 诱导的函数
\[
f_E:\ZZ_{\ge 0}\to \ZZ_{\ge 0},
\qquad
f_E(k):=\sum_{e\in E}\min(k,e),
\]
组成的集合，其中 \(f_E(0)=0\)。赋予 \(\mathfrak F_p\) 逐点加法。对每个 \(t\ge 1\)，定义
\[
u_t(k):=\min(k,t)\qquad(k\ge 0).
\]
则 \((\mathfrak F_p,+)\) 是由 \(\{u_t\}_{t\ge 1}\) 自由生成的交换幺半群。更具体地，每个 \(f\in\mathfrak F_p\) 都唯一表示为
\[
f=\sum_{t\ge 1}m_tu_t,
\qquad
m_t\in \ZZ_{\ge 0},
\]
且仅有限多个 \(m_t\) 非零；其系数可由离散曲率唯一恢复：
\[
m_t=\Delta_f(t)-\Delta_f(t+1),
\qquad
\Delta_f(k):=f(k)-f(k-1).
\]
特别地，任意满行秩整数矩阵的 \(\,p\)-局部 Smith 信息亏损谱都属于 \(\mathfrak F_p\)。
\end{theorem}
\begin{proof}
若 \(E,F\) 为有限多重集，则
\[
f_E+f_F=f_{E\sqcup F},
\]
其中 \(E\sqcup F\) 表示按重数计的并。因此 \(\mathfrak F_p\) 在逐点加法下确为交换幺半群。又由定义
\[
f_E=\sum_{t\ge 1}m_tu_t,
\qquad
m_t:=\#\{e\in E:\ e=t\},
\]
可知 \(\{u_t\}_{t\ge 1}\) 生成 \(\mathfrak F_p\)。

唯一性可由前述 Smith 谱分类直接读出。事实上，定理 \ref{thm:xi-smith-padic-loss-spectrum-classification} 已说明，这类函数恰由非负、非增且最终为零的一阶差分所刻画；而推论 \ref{cor:xi-smith-second-discrete-curvature-recovery} 与定理 \ref{thm:xi-smith-loss-spectrum-discrete-curvature-complete-characterization} 给出
\[
\Delta_f(t)-\Delta_f(t+1)
=
2f(t)-f(t-1)-f(t+1),
\]
其值正是指数 \(t\) 的重数。故一旦
\[
f=\sum_{t\ge 1}m_tu_t=\sum_{t\ge 1}m_t'u_t,
\]
就必有
\[
m_t=\Delta_f(t)-\Delta_f(t+1)=m_t'
\qquad(\forall\,t\ge 1).
\]
这就证明了自由性。最后，若 \(f=f_p\) 为某个满行秩整数矩阵的 \(\,p\)-局部信息亏损谱，则由定理 \ref{thm:killo-smith-loss-spectrum}
\[
f_p(k)=\sum_{i=1}^{m}\min(k,e_i(p)),
\]
故 \(f_p\in \mathfrak F_p\)。证毕。
\end{proof}

\begin{theorem}[局部 Euler 因子对 Smith 正规形的完备性]\label{thm:xi-smith-local-euler-complete-invariant}
沿用定理 \ref{thm:xi-kernel-loss-divisibility-valuation-rational-euler} 的记号。对每个素数 \(p\)，把局部因子视为形式幂级数
\[
\Xi_{A,p}(T)
:=
\sum_{k\ge 0}p^{f_p(k)}T^k.
\]
则 \(\Xi_{A,p}(T)\) 唯一决定整个序列 \((f_p(k))_{k\ge 1}\)，从而唯一决定多重集 \(\{e_i(p)\}_{i=1}^{m}\)。因而全局 Dirichlet 级数
\[
\Xi_A(s)
:=
\sum_{b\ge 1}\frac{\#\ker(A_b)}{b^{\,n-m+s}}
\]
唯一决定 \(A\) 的全部 Smith 不变量。更一般地，若另一个满行秩整数矩阵
\[
B\in M_{m_B\times n_B}(\ZZ)
\]
满足
\[
\Xi_A=\Xi_B,
\]
则 \(A\) 与 \(B\) 具有相同的 Smith 正规形。
\end{theorem}
\begin{proof}
由定理 \ref{thm:xi-kernel-loss-divisibility-valuation-rational-euler}，
\[
\Xi_{A,p}(T)=\sum_{k\ge 0}p^{f_p(k)}T^k.
\]
因此对每个 \(k\ge 1\)，其 \(T^k\) 项系数恰为 \(p^{f_p(k)}\)，从而
\[
f_p(k)=\log_p\!\bigl([T^k]\Xi_{A,p}(T)\bigr).
\]
局部亏损谱一旦给定，定理 \ref{thm:killo-smith-loss-spectrum} 与定理 \ref{thm:xi-smith-loss-spectrum-discrete-curvature-complete-characterization} 即可无损恢复多重集 \(\{e_i(p)\}_{i=1}^{m}\)。

再看全局唯一性。若 \(\Xi_A=\Xi_B\) 作为 Dirichlet 级数相等，则对每个正整数 \(b\) 都有
\[
\frac{\#\ker(A_b)}{b^{\,n-m}}
=
\frac{\#\ker(B_b)}{b^{\,n_B-m_B}},
\]
亦即两者的归一化系数逐项相同。特别地，对任意素数幂 \(b=p^k\)，得到
\[
p^{f_{A,p}(k)}
=
\frac{\#\ker(A\bmod p^k)}{p^{k(n-m)}}
=
\frac{\#\ker(B\bmod p^k)}{p^{k(n_B-m_B)}}
=
p^{f_{B,p}(k)}.
\]
故
\[
f_{A,p}(k)=f_{B,p}(k)\qquad(\forall\,p,\ \forall\,k\ge 1).
\]
于是对每个素数 \(p\)，两矩阵的 \(\{e_i(p)\}\) 完全一致；逐素数合并便得全部 Smith 不变量一致，亦即 Smith 正规形相同。证毕。
\end{proof}

\begin{corollary}[自由原子结构在整数轴对齐椭圆半群中的几何实现]\label{cor:xi-free-atomic-structure-ellipse-realization}
定理 \ref{thm:xi-smith-local-loss-profiles-free-commutative-monoid} 与定理 \ref{thm:xi-integer-ellipse-free-commutative-monoid} 表明，\(\,p\)-局部信息亏损谱幺半群与整数轴对齐椭圆半群都属于以可数原子集为基的自由交换幺半群：前者的原子为 \(\{u_t\}_{t\ge 1}\)，后者的原子为
\[
\mathcal A=\{E_{p,1},E_{1,p}:p\in\mathcal P\}.
\]
因此
\[
\operatorname{Aut}_{\mathrm{Mon}}(\mathsf E_{\NN},\odot)\cong \mathrm{Sym}(\mathcal A)
\]
可视为附录 \(H\) 中自由原子机制在显式几何对象类上的完全对称实现。
\end{corollary}
\begin{proof}
定理 \ref{thm:xi-smith-local-loss-profiles-free-commutative-monoid} 给出 \(\mathfrak F_p\) 的自由原子分解，定理 \ref{thm:xi-integer-ellipse-free-commutative-monoid} 的第一条给出 \((\mathsf E_{\NN},\odot)\) 的自由原子分解，而其第三条正给出所述自同构群公式。证毕。
\end{proof}

\begin{corollary}[Smith 亏损谱的二阶差分、平台起点与总亏损统一恢复]\label{cor:xi-smith-loss-spectrum-plateau-start-total-loss}
设 \(A\in M_{m\times n}(\ZZ)\) 满足 \(\operatorname{rank}_{\QQ}(A)=m\)。固定素数 \(p\)，并沿用定理 \ref{thm:killo-smith-loss-spectrum} 的记号：
\[
f_p(0):=0,\qquad
\Delta_p(k):=f_p(k)-f_p(k-1)\quad(k\ge 1),
\]
\[
e_i(p):=v_p(d_i),\qquad
m_{p,j}:=\#\{i:\ e_i(p)=j\}\quad(j\ge 1).
\]
则对每个 \(j\ge 1\) 都有
\[
m_{p,j}
=
\Delta_p(j)-\Delta_p(j+1)
=
2f_p(j)-f_p(j-1)-f_p(j+1).
\]
再记
\[
E_p:=\max_i e_i(p).
\]
则
\[
E_p
=
\max\bigl(\{k\ge 1:\Delta_p(k)>0\}\cup\{0\}\bigr)
=
\min\{k\ge 0:\ f_p(k)=f_p(k+1)\},
\]
并且
\[
\sum_{i=1}^{m}e_i(p)=\lim_{k\to\infty}f_p(k).
\]
换言之，\(f_p\) 是一条最终常值的非减离散凹函数，而其二阶差分、平台起点与平台值已经分别精确编码了各个 \(\,p\)-进高度的重数、最大高度与总亏损。
\end{corollary}
\begin{proof}
二阶差分公式与极限公式分别正是推论 \ref{cor:xi-smith-second-discrete-curvature-recovery} 与推论 \ref{cor:xi-smith-volume-recovered-from-padic-platform} 的内容。

再由推论 \ref{cor:xi-smith-kernel-spectrum-discrete-concavity-plateau-inversion}，\(f_p\) 为非减离散凹函数，并在 \(k\ge E_p\) 后进入平台。另一方面，
\[
\Delta_p(k)=\#\{i:\ e_i(p)\ge k\}
\]
表明 \(\Delta_p(k)>0\) 当且仅当存在某个 \(i\) 满足 \(e_i(p)\ge k\)，亦即 \(k\le E_p\)。故
\[
E_p=\max\bigl(\{k\ge 1:\Delta_p(k)>0\}\cup\{0\}\bigr).
\]
又因为
\[
f_p(k+1)-f_p(k)=\Delta_p(k+1),
\]
故 \(f_p(k)=f_p(k+1)\) 当且仅当 \(\Delta_p(k+1)=0\)，其最早发生恰在 \(k=E_p\)。这就得到第二个恢复公式。证毕。
\end{proof}

\begin{theorem}[全模数核计数的严格 Smith 饱和平台]\label{thm:xi-smith-global-kernel-dmax-saturation-platform}
设 \(A\in M_{m\times n}(\ZZ)\) 满足 \(\operatorname{rank}_{\QQ}(A)=m\)，其 Smith 不变量为
\[
d_1\mid d_2\mid\cdots\mid d_m.
\]
对每个 \(N\ge 1\)，记
\[
A_N:(\ZZ/N\ZZ)^n\to(\ZZ/N\ZZ)^m,
\qquad
K_A(N):=\#\ker(A_N).
\]
若记 \(e_i(p):=v_p(d_i)\)、\(f_p(k):=\sum_{i=1}^{m}\min(k,e_i(p))\)，并设
\[
N=\prod_p p^{\nu_p(N)},
\qquad
\Delta_A:=\prod_{i=1}^{m}d_i,
\qquad
d_{\max}:=d_m,
\]
则
\[
K_A(N)
=
N^{\,n-m}\prod_{i=1}^{m}\gcd(d_i,N)
=
N^{\,n-m}\prod_{p}p^{\,f_p(\nu_p(N))}.
\]
特别地，
\[
K_A(N)\le N^{\,n-m}\Delta_A,
\]
并且等号成立当且仅当
\[
d_{\max}\mid N.
\]
因而全模数核计数在整除格上存在一条严格的 Smith 饱和平台：
\[
K_A(N)=N^{\,n-m}\Delta_A
\qquad(d_{\max}\mid N),
\]
并且 \(d_{\max}\) 正是使全部局部亏损同时饱和的最小阈值。
\end{theorem}
\begin{proof}
由推论 \ref{cor:killo-kernel-size-general-modulus}，
\[
K_A(N)=N^{\,n-m}\prod_{i=1}^{m}\gcd(d_i,N).
\]
再把每个 \(\gcd(d_i,N)\) 按素数分解，便得
\[
\prod_{i=1}^{m}\gcd(d_i,N)
=
\prod_{p}p^{\sum_{i=1}^{m}\min(\nu_p(N),e_i(p))}
=
\prod_{p}p^{f_p(\nu_p(N))}.
\]
这证明了第一式。

对每个 \(i\) 都有 \(\gcd(d_i,N)\le d_i\)，故
\[
K_A(N)\le N^{\,n-m}\prod_{i=1}^{m}d_i
=
N^{\,n-m}\Delta_A.
\]
并且取等当且仅当
\[
\gcd(d_i,N)=d_i
\qquad(1\le i\le m),
\]
亦即当且仅当每个 \(d_i\mid N\)。由于
\[
d_1\mid d_2\mid\cdots\mid d_m,
\]
上述条件又等价于 \(d_m\mid N\)，即 \(d_{\max}\mid N\)。

因此达到上界的模数恰为 \(d_{\max}\) 的倍数，故一旦模数吸收最大 Smith 不变量，全部局部亏损便同时饱和；反之若 \(d_{\max}\nmid N\)，则至少有一个因子 \(\gcd(d_i,N)\) 严格小于 \(d_i\)，从而上界不能取到。故 \(d_{\max}\) 正是所述最小饱和阈值。证毕。
\end{proof}

\noindent
由此，附录 \(H\) 中的结构链可进一步收束为
\[
\text{有限深度分支寄存器}
\Longrightarrow
\text{\(n\)-步过去的极小完备前缀}
\Longrightarrow
\text{自然扩张的严格逆极限},
\]
\[
\text{自然扩张深度}
\Longleftrightarrow
\text{最坏 prime-slice 账本维数},
\]
\[
\text{Smith 局部亏损谱}
\Longrightarrow
\text{二阶曲率与平台起点恢复}
\Longrightarrow
\text{自由原子分解}
\Longrightarrow
\text{Euler 完备不变量}
\Longrightarrow
\text{全模数饱和阈值},
\]
而推论 \ref{cor:xi-free-atomic-structure-ellipse-realization} 则说明同一自由原子机制在整数轴对齐椭圆半群上获得了显式几何实现与全对称自同构群；定理 \ref{thm:xi-smith-global-kernel-dmax-saturation-platform} 则进一步表明，这条局部到全局的链最终由单一最大 Smith 不变量锁定其整除阈值。

\paragraph{分支指标像空间的坐标逆极限与全深度 prime-register 的无限秩障碍}
在定理 \ref{thm:xi-branch-register-strict-inverse-limit} 已把自然扩张的分支寄存器模型识别为有限深度截断的严格逆极限之后，还可把这一事实重写为像空间坐标本身的逆极限完备性；进一步，当每一层都存在真实分支时，忠实的全深度 prime-register 账本便不再是有限层 prime-slice 最优化的简单延长，而会在结构层上强制出现可数无穷个彼此独立的素数方向。

\begin{theorem}[分支指标像空间本身就是有限截断坐标系的规范逆极限]\label{thm:xi-branch-register-image-coordinate-inverse-limit}
\leanverified{paper\_xi\_branch\_register\_image\_coordinate\_inverse\_limit}
沿用定理 \ref{thm:killo-natural-extension-branch-register} 的记号。对每个 \(n\ge 1\)，定义前缀投影
\[
\mathrm{pr}_{\le n}:\widehat X\to X\times \NN^n,
\qquad
\mathrm{pr}_{\le n}\bigl(x,(a_1,a_2,\dots)\bigr):=(x,a_1,\dots,a_n),
\]
并记
\[
Y_n:=\mathrm{pr}_{\le n}(\widehat X)
=
\Bigl\{(x_0,a_1,\dots,a_n)\in X\times\NN^n:\ \exists\,\widehat x\in\widehat X,
\mathrm{pr}_{\le n}(\widehat x)=(x_0,a_1,\dots,a_n)\Bigr\}.
\]
定义截断映射
\[
\rho_{n+1,n}:Y_{n+1}\to Y_n,
\qquad
\rho_{n+1,n}(x_0,a_1,\dots,a_{n+1})=(x_0,a_1,\dots,a_n).
\]
则规范映射
\[
\Theta:\widehat X\to \varprojlim(Y_n,\rho_{n+1,n}),
\qquad
\Theta(\widehat x):=\bigl(\mathrm{pr}_{\le n}(\widehat x)\bigr)_{n\ge 1},
\]
为双射。

进一步，对每个 \(n\ge 1\) 定义有限深度右移插入系统
\[
\widehat T_n:Y_n\to Y_n,
\qquad
\widehat T_n(x_0,a_1,\dots,a_n):=
\begin{cases}
(T(x_0),\beta(x_0)),& n=1,\\
(T(x_0),\beta(x_0),a_1,\dots,a_{n-1}),& n\ge 2.
\end{cases}
\]
则在上述识别下，
\[
\widehat T=\varprojlim_n \widehat T_n.
\]
换言之，分支指标模型 \(\widehat X\) 本身就是有限深度右移插入系统的逆极限对象。
\end{theorem}
\begin{proof}
由定理 \ref{thm:xi-branch-register-finite-past-minimal-complete-prefix}，\(n\)-步过去的分支寄存器前缀
\[
(x_0,a_1,\dots,a_n)
\]
与唯一有限前史
\[
(x_0,x_{-1},\dots,x_{-n})\in X^{\leftarrow,n}
\]
一一对应。反过来，由 \(T\) 为满射，任意有限前史都可继续向后逐层选取原像，因而总能延拓为某个无限前史。故 \(Y_n\) 与定理 \ref{thm:xi-branch-register-strict-inverse-limit} 中的
\[
\widehat X_n^{\mathrm{br}}=\mathcal C_n(X^{\leftarrow,n})
\]
是同一对象，只是记号改写而已。

\(\Theta\) 的单射性是直接的：若两点在全部有限前缀上相同，则其 \(\widehat X\subset X\times\NN^\NN\) 中的全部坐标都相同，因而两点一致。

再证满射。任取一致族
\[
y^{(n)}=(x_0^{(n)},a_1^{(n)},\dots,a_n^{(n)})\in Y_n,
\qquad
\rho_{n+1,n}(y^{(n+1)})=y^{(n)}.
\]
由一致性可知 \(x_0^{(n)}\) 稳定，且对每个固定 \(j\ge 1\)，当 \(n\ge j\) 时 \(a_j^{(n)}\) 亦稳定。记稳定值为
\[
x_0,\qquad a_j\in\NN\quad (j\ge 1).
\]
由于每个 \(y^{(n)}\in Y_n\)，由定理 \ref{thm:xi-branch-register-finite-past-minimal-complete-prefix} 可恢复唯一有限前史
\[
(x_0,x_{-1}^{(n)},\dots,x_{-n}^{(n)})\in X^{\leftarrow,n}.
\]
而 \(\rho_{n+1,n}(y^{(n+1)})=y^{(n)}\) 又迫使这些有限前史在重叠部分完全一致，故它们定义出唯一无限前史
\[
(x_0,x_{-1},x_{-2},\dots)\in X^{\leftarrow}.
\]
将其送入定理 \ref{thm:killo-natural-extension-branch-register} 的编码映射，得到
\[
\widehat x:=\mathcal C(x_0,x_{-1},x_{-2},\dots)\in \widehat X,
\]
并满足
\[
\mathrm{pr}_{\le n}(\widehat x)=y^{(n)}\qquad(\forall n\ge 1).
\]
故 \(\Theta\) 为满射。

最后，由定理 \ref{thm:killo-natural-extension-branch-register} 的右移插入公式
\[
\widehat T(x,(a_1,a_2,\dots))=(T(x),(\beta(x),a_1,a_2,\dots)),
\]
对每个 \(n\) 取前缀便得到
\[
\mathrm{pr}_{\le n}\!\bigl(\widehat T(\widehat x)\bigr)
=
\widehat T_n\!\bigl(\mathrm{pr}_{\le n}(\widehat x)\bigr).
\]
因此 \(\Theta\) 将 \(\widehat T\) 精确识别为有限深度右移插入系统的逆极限动力学。证毕。
\end{proof}

\begin{theorem}[忠实的全深度 prime-register 实现必然具有无限素数秩，因而排斥任何有限生成加法账本]\label{thm:xi-infinite-depth-prime-ledger-infinite-rank-obstruction}
\leanverified{paper\_xi\_infinite\_depth\_prime\_ledger\_infinite\_rank\_obstruction}
设 \(T:X\twoheadrightarrow X\) 为有限纤维满射，并假设其逆向分支深度无界：对每个 \(N\ge 1\)，都存在长度 \(N\) 的前史片段
\[
x_{-N}\xrightarrow{\,T\,}x_{-N+1}\xrightarrow{\,T\,}\cdots\xrightarrow{\,T\,}x_{-1}\xrightarrow{\,T\,}x_0
\]
满足
\[
\#T^{-1}(x_{-j+1})\ge 2
\qquad(1\le j\le N).
\]
称一个全深度 prime-register 账本实现为忠实，若它对每个有限深度 \(N\) 的前缀截断都与定理 \ref{thm:xi-layered-prime-register-sharp-natural-extension} 的精确分层外置语义相容。

则任意忠实的全深度实现都必须使用无穷多个彼此独立的 prime slices。更具体地，若从每一层的有效 prime slice 中选取一枚代表素数
\[
q_j\in \mathcal P_j\qquad (j\ge 0),
\]
并定义自由交换幺半群
\[
M_\infty
:=
\left\{
\prod_{j\ge 0}q_j^{a_j}:\ a_j\in\ZZ_{\ge 0},\ a_j=0\ \text{仅有限多个例外}
\right\},
\]
则 \(M_\infty\) 在唯一分解意义下同构于 \(\NN_{\times}\)。

特别地，不存在任何有限生成交换群 \(L\) 使忠实的全深度 prime-register 账本能够严格同态化到
\[
(\ZZ\oplus L,+)
\]
这样的有限秩加法账本上。
\end{theorem}
\begin{proof}
第一部分正是定理 \ref{thm:xi-natural-extension-infinite-prime-register-necessity} 的内容：在逆向分支深度无界时，任何精确实现整个自然扩张的分层 prime-register 外置都不可能只使用有限多个 prime slices。结合上一定理，\(\widehat X\) 已被识别为全部有限深度前缀系统的逆极限，因此忠实的全深度账本若与每个有限前缀的精确恢复相容，就必须在每一层保留一个独立的有效 prime direction。

由 \((q_j)_{j\ge 0}\) 两两不同且都是素数，唯一分解给出幺半群同构
\[
\vartheta:\NN_{\times}\xrightarrow{\ \sim\ } M_\infty,
\]
其在素数生成元上由
\[
\vartheta(p_{j+1})=q_j
\qquad(j\ge 0)
\]
确定。

反设存在有限生成交换群 \(L\)，使忠实的全深度账本可严格同态化到 \((\ZZ\oplus L,+)\) 上。则其 restriction 在上述独立素数方向上给出一个幺半群单射
\[
M_\infty\hookrightarrow (\ZZ\oplus L,+).
\]
与 \(\vartheta\) 复合便得到
\[
\NN_{\times}\hookrightarrow (\ZZ\oplus L,+)
\]
的单射幺半群同态。
但 \((\ZZ\oplus L,+)\) 是有限生成交换可消去幺半群，这与定理 \ref{thm:killo-prime-freedom-non-finitizability} 矛盾；等价地，这也与定理 \ref{thm:killo-godel-compression-not-finite-rank-homologizable} 所排除的有限秩同态化账本机制相冲突。故此类有限生成加法账本不存在。证毕。
\end{proof}

\noindent
配合定理 \ref{thm:xi-terminal-fixed-freezing-tv-collapse}、定理 \ref{thm:xi-fixed-freezing-escort-bounded-observable-collapse} 以及推论 \ref{cor:xi-time-part57bb-freezing-groundstate-subset-equidistribution} 可见：自然扩张首先在分支指标像空间层面构成严格逆极限对象；冻结 escort 态随后在统计层面指数塌缩到最大纤维上的微正则分布；而一旦要求把这种全深度可逆账本忠实地外置为 prime-register，其结构层便不可避免地转入可数无穷素数秩，从而排斥一切有限秩加法线性化。

\paragraph{自然扩张前缀的 Gödel 算术化与 \texorpdfstring{$p$}{p}-局部核增长的单尾极点}
在自然扩张分支寄存器的有限前缀、prime-slice 历史压缩与 Smith 局部亏损谱的 Euler 因子化都已建立之后，还可把有限前史地址与 \(p\)-局部核增长的解析尾部进一步压缩为下列两条接口。

\begin{corollary}[分层素数外置恰是自然扩张前缀地址的 Gödel 化]\label{cor:xi-time-part62ba-primeslice-natural-extension-prefix-godelization}
\leanverified{paper\_xi\_time\_part62ba\_primeslice\_natural\_extension\_prefix\_godelization}
设 \(X\) 为可数集合，\(T:X\twoheadrightarrow X\) 为有限纤维满射，并沿用定理 \ref{thm:killo-natural-extension-branch-register}、定理 \ref{thm:xi-prime-register-history-inverse-limit} 与定理 \ref{thm:xi-branch-register-finite-past-minimal-complete-prefix} 的记号。对任意 \(k\ge 1\)，在自映射前史塔
\[
X_0=X_1=\cdots=X_k=X,
\qquad
\pi_0=\pi_1=\cdots=\pi_{k-1}=T
\]
上，定义长度 \(k\) 的截断地址
\[
\mathcal C_k(x_0,x_{-1},\dots,x_{-k})
:=
\bigl(x_0,\beta(x_{-1}),\dots,\beta(x_{-k})\bigr).
\]
再取定定理 \ref{thm:xi-prime-register-history-inverse-limit} 中的分层单射
\[
\alpha_\ell:X^{\leftarrow,\ell-1}\times\NN\hookrightarrow \mathcal P_{\ell-1}
\qquad (1\le \ell\le k),
\]
并令
\[
G_k(x_0,\dots,x_{-k})
:=
\prod_{\ell=1}^{k}
\alpha_\ell\!\bigl((x_0,\dots,x_{-\ell+1}),\beta(x_{-\ell})\bigr).
\]
则对任意固定首端点 \(x_0\in X\)，长度 \(k\) 的逆轨道前缀
\[
(x_{-1},\dots,x_{-k})
\]
与整数 \(G_k(x_0,\dots,x_{-k})\) 之间构成双射。等价地，定理 \ref{thm:killo-layered-prime-slices-finite-depth} 的分层素数外置，在自映射前史塔上恰是长度 \(k\) 的自然扩张地址 \(\mathcal C_k\) 的 Gödel 编码。
\end{corollary}
\begin{proof}
定理 \ref{thm:xi-branch-register-finite-past-minimal-complete-prefix} 说明：对固定 \(x_0\) 与 \(k\)，截断地址 \(\mathcal C_k\) 与长度 \(k\) 的前史前缀彼此唯一决定。另一方面，定理 \ref{thm:xi-prime-register-history-inverse-limit} 已证明：由 \((x_0,G_k)\) 可确定性恢复唯一历史 \((x_0,\dots,x_{-k})\)，等价地可恢复唯一的 \(\mathcal C_k(x_0,\dots,x_{-k})\)。因此在固定首端点 \(x_0\) 后，整数 \(G_k\) 与长度 \(k\) 的逆轨道前缀双向确定，结论成立。证毕。
\end{proof}

\begin{theorem}[\texorpdfstring{$p$}{p}-局部核增长序列的局部 \texorpdfstring{$\zeta$}{zeta} 函数只有唯一尾极点]\label{thm:xi-time-part62ba-padic-kernel-growth-local-zeta-single-tail-pole}
\leanverified{paper\_xi\_time\_part62ba\_padic\_kernel\_growth\_local\_zeta\_single\_tail\_pole}
设 \(A\in M_{m\times n}(\ZZ)\) 满足 \(\operatorname{rank}_{\QQ}(A)=m\)，其 Smith 不变量为
\[
d_1\mid d_2\mid\cdots\mid d_m.
\]
固定素数 \(p\)，并定义
\[
K_p(k):=\#\ker(A\bmod p^k)
\qquad (k\ge 1),
\]
\[
K_p(0):=1,
\]
\[
e_i(p):=v_p(d_i)
\qquad (1\le i\le m),
\qquad
E_p:=\max_{1\le i\le m}e_i(p).
\]
再定义 \(p\)-局部核增长的 ordinary generating function
\[
\mathcal Z_{A,p}^{\mathrm{loc}}(z)
:=
\sum_{k\ge 0}K_p(k)z^k.
\]
则 \(\mathcal Z_{A,p}^{\mathrm{loc}}\) 为有理函数，并且具有显式分解
\[
\boxed{\;
\mathcal Z_{A,p}^{\mathrm{loc}}(z)
=
\sum_{k=0}^{E_p-1}K_p(k)z^k
+
\frac{K_p(E_p)\,z^{E_p}}{1-p^{\,n-m}z}.\;}
\]
换言之，\(p\)-局部核增长在有限窗口 \(\{0,1,\dots,E_p\}\) 之后进入严格几何尾部；其尾部唯一极点位于
\[
\boxed{\;
z=p^{-(n-m)}.\;}
\]
\end{theorem}
\begin{proof}
由定理 \ref{thm:killo-smith-loss-spectrum}，对每个 \(k\ge 1\) 都有
\[
K_p(k)
=
p^{\,k(n-m)+\sum_{i=1}^{m}\min(k,e_i(p))}.
\]
若 \(E_p=0\)，则一切 \(e_i(p)=0\)，从而
\[
K_p(k)=p^{k(n-m)}
\qquad (k\ge 0),
\]
于是
\[
\mathcal Z_{A,p}^{\mathrm{loc}}(z)=\frac{1}{1-p^{\,n-m}z},
\]
结论显然成立。以下设 \(E_p\ge 1\)。

当 \(k\ge E_p\) 时，对一切 \(i\) 都有 \(\min(k,e_i(p))=e_i(p)\)，故
\[
K_p(k)
=
p^{\,k(n-m)+\sum_{i=1}^{m}e_i(p)}
=
K_p(E_p)\,p^{(k-E_p)(n-m)}.
\]
于是尾部级数成为公比 \(p^{n-m}z\) 的几何级数：
\[
\sum_{k\ge E_p}K_p(k)z^k
=
K_p(E_p)z^{E_p}\sum_{r\ge 0}\bigl(p^{n-m}z\bigr)^r
=
\frac{K_p(E_p)\,z^{E_p}}{1-p^{\,n-m}z}.
\]
再加上前 \(E_p\) 项，即得所示“有限前缀 + 几何尾部”分解。由于分母只有 \(1-p^{\,n-m}z\) 这一项，故尾部唯一极点恰位于 \(z=p^{-(n-m)}\)。证毕。
\end{proof}

\noindent
因此，prime-slice 外置与自然扩张地址之间的关系并不是两套并列装置，而是同一有限前史证书的加性与乘法两种坐标；与之平行，Smith 的 \(p\)-局部核增长也把全部非平凡复杂性压缩到有限前缀，而把尾部解析结构严格收束为唯一一个几何极点。

\paragraph{Fibonacci 递推 Fourier 采样、零点缩放均匀律与局部化 solenoid 的自覆盖分类}
Pisot--Bernoulli Fourier 采样、Fibonacci 共振整函数的零点缩放统计，以及有限素数局部化 solenoid 的内部自同态结构，还可进一步压缩为下列四条可直接调用的结论。

\begin{theorem}[Pisot--Bernoulli Fourier 采样在任意非零 Fibonacci 递推方向上不发生 Rajchman 衰减]\label{thm:xi-pisot-bernoulli-fibonacci-direction-nonrajchman}
\leanverified{paper\_xi\_pisot\_bernoulli\_fibonacci\_direction\_nonrajchman}
令 \(q:=\varphi^{-1}\in(0,1)\)，并令 \(\mu_q^{(2)}\) 为定理 \ref{thm:fold-pisot-bernoulli-convolution-representation} 中的 Pisot 参数 Bernoulli 卷积移位测度。取任意非零整数初值对
\[
(u_0,u_1)\in \ZZ^2\setminus\{(0,0)\},
\]
并令 \((u_n)_{n\ge 0}\) 满足 Fibonacci 递推
\[
u_{n+2}=u_{n+1}+u_n.
\]
再记
\[
\psi:=-\varphi^{-1},\qquad
\alpha:=\frac{u_1-u_0\psi}{\sqrt5},\qquad
\beta:=\frac{u_0\varphi-u_1}{\sqrt5}.
\]
则极限
\[
\lim_{n\to\infty}\Bigl|\widehat{\mu_q^{(2)}}(\pi u_n)\Bigr|
\]
存在且严格为正，并满足显式公式
\[
\boxed{\;
\lim_{n\to\infty}\Bigl|\widehat{\mu_q^{(2)}}(\pi u_n)\Bigr|
=
c(u_0,u_1)
:=
\Bigl|\cos(\pi\alpha)\Bigr|
\prod_{k=1}^{\infty}\Bigl|\cos(\pi\alpha\varphi^{-k})\Bigr|
\prod_{k=1}^{\infty}
\Bigl|\cos(\pi\beta\varphi^{-k})\Bigr|
>0.
\;}
\]
此外 \(|u_n|\to\infty\)。若 \(\nu\) 记 \(\mu_q^{(2)}\) 在 \(\mathbb T_2:=\RR/(2\ZZ)\) 上的模 \(2\) 周期化测度，则
\[
\Bigl|\widehat{\nu}(u_n)\Bigr|
=
\Bigl|\widehat{\mu_q^{(2)}}(\pi u_n)\Bigr|
\longrightarrow c(u_0,u_1)>0,
\]
从而 \(\nu\) 沿任意非平凡 Fibonacci 递推逃逸序列都不发生 Rajchman 衰减。
\end{theorem}
\begin{proof}
由定理 \ref{thm:fold-pisot-bernoulli-convolution-representation}，
\[
\Bigl|\widehat{\mu_q^{(2)}}(\pi u)\Bigr|=C_\varphi(u)
\qquad(u\in\ZZ).
\]
另一方面，定理 \ref{thm:fold-resonance-ladder-fibonacci-directional-limit} 对任意非零初值对给出
\[
\lim_{n\to\infty}C_\varphi(u_n)
=
\Bigl|\cos(\pi\alpha)\Bigr|
\prod_{k=1}^{\infty}\Bigl|\cos(\pi\alpha\varphi^{-k})\Bigr|
\prod_{k=1}^{\infty}
\Bigl|\cos(\pi\beta\varphi^{-k})\Bigr|
>0.
\]
两式合并即得所示 Fourier 采样极限公式。

再证 \(|u_n|\to\infty\)。由 Binet 公式，
\[
u_n=\alpha\varphi^n+\beta\psi^n.
\]
若 \(\alpha=0\)，则
\[
u_1=u_0\psi.
\]
由于 \(\psi=-\varphi^{-1}\notin \QQ\)，这与 \(u_0,u_1\in\ZZ\) 同时成立只能推出 \(u_0=u_1=0\)，与假设矛盾。故 \(\alpha\neq 0\)。于是
\[
u_n=\alpha\varphi^n\Bigl(1+\frac{\beta}{\alpha}(\psi/\varphi)^n\Bigr)
=
\alpha\varphi^n(1+o(1)),
\]
从而
\[
|u_n|=|\alpha|\varphi^n(1+o(1))\to\infty.
\]

最后，由推论 \ref{cor:fold-bernoulli-fourier-sampling-not-l2} 的证明，\(\nu\) 的 Fourier 系数满足
\[
\widehat{\nu}(u)=\widehat{\mu_q^{(2)}}(\pi u)
\qquad(u\in\ZZ).
\]
既然 \(|u_n|\to\infty\) 且 \(|\widehat{\mu_q^{(2)}}(\pi u_n)|\to c(u_0,u_1)>0\)，便得到
\[
\Bigl|\widehat{\nu}(u_n)\Bigr|\to c(u_0,u_1)>0,
\]
从而 \(\nu\) 沿该逃逸序列不发生 Rajchman 衰减。证毕。
\end{proof}

\begin{corollary}[任意非零 Fibonacci 递推方向上的对数级平方能量律]\label{cor:xi-pisot-bernoulli-fibonacci-direction-log-energy}
沿用定理 \ref{thm:xi-pisot-bernoulli-fibonacci-direction-nonrajchman} 的记号，并记
\[
c(u_0,u_1)
:=
\lim_{n\to\infty}\Bigl|\widehat{\mu_q^{(2)}}(\pi u_n)\Bigr|
>0.
\]
则当 \(N\to\infty\) 时，
\[
\boxed{\;
\sum_{\substack{n\ge 0\\ |u_n|\le N}}
\Bigl|\widehat{\mu_q^{(2)}}(\pi u_n)\Bigr|^2
=
\bigl(c(u_0,u_1)^2+o(1)\bigr)\log_\varphi N.
\;}
\]
特别地，
\[
\sum_{n\ge 0}\Bigl|\widehat{\mu_q^{(2)}}(\pi u_n)\Bigr|^2=+\infty.
\]
\end{corollary}
\begin{proof}
由定理 \ref{thm:xi-pisot-bernoulli-fibonacci-direction-nonrajchman}，
\[
u_n=\alpha\varphi^n+\beta\psi^n,
\qquad
\alpha\neq 0,
\]
因而
\[
|u_n|=|\alpha|\varphi^n(1+o(1)).
\]
设
\[
M(N):=\#\{n\ge 0:\ |u_n|\le N\}.
\]
则上式等价于
\[
M(N)=\log_\varphi N-\log_\varphi|\alpha|+O(1)
=
\log_\varphi N+O(1).
\]

另一方面，由
\[
\Bigl|\widehat{\mu_q^{(2)}}(\pi u_n)\Bigr|^2\to c(u_0,u_1)^2
\]
可知其部分和满足 Ces\`aro 渐近
\[
\sum_{n=0}^{M}\Bigl|\widehat{\mu_q^{(2)}}(\pi u_n)\Bigr|^2
=
c(u_0,u_1)^2\,M+o(M)
\qquad(M\to\infty).
\]
取 \(M=M(N)\) 即得
\[
\sum_{\substack{n\ge 0\\ |u_n|\le N}}
\Bigl|\widehat{\mu_q^{(2)}}(\pi u_n)\Bigr|^2
=
c(u_0,u_1)^2\,M(N)+o(M(N))
=
\bigl(c(u_0,u_1)^2+o(1)\bigr)\log_\varphi N.
\]
由于右端随 \(N\to\infty\) 对数发散，第二条随即成立。证毕。
\end{proof}

\begin{theorem}[Fibonacci 共振整函数的缩放零点计数测度趋于区间均匀分布]\label{thm:xi-fold-resonance-scaled-zero-measure-uniform}
沿用定理 \ref{thm:fold-resonance-entire-lp} 与定理 \ref{thm:xi-fold-resonance-zero-count-selfsimilar-linear} 的记号，记
\[
\mathcal F_\varphi(z):=\prod_{n=2}^{\infty}\cos(\pi z\varphi^{-n}),
\]
\[
N_\varphi(R):=\#\bigl(\mathcal Z(\mathcal F_\varphi)\cap[-R,R]\bigr),
\qquad
\nu_R:=\frac{1}{N_\varphi(R)}\sum_{\substack{\zeta\in\mathcal Z(\mathcal F_\varphi)\\ |\zeta|\le R}}\delta_{\zeta/R}.
\]
则当 \(R\to\infty\) 时，
\[
\boxed{\;
\nu_R \Longrightarrow \frac12\,\mathbf 1_{[-1,1]}(x)\,\dd x.
\;}
\]
特别地，
\[
\boxed{\;
N_\varphi(R)=2R+O(\log R).\;}
\]
\end{theorem}
\begin{proof}
由定理 \ref{thm:fold-resonance-entire-lp}，
\[
\mathcal Z(\mathcal F_\varphi)
=
\Bigl\{(k+\tfrac12)\varphi^n:\ k\in\ZZ,\ n\ge 2\Bigr\}.
\]
固定闭区间 \([a,b]\subset(-1,1)\)。对每个固定层 \(n\ge 2\)，落在区间 \([aR,bR]\) 中的零点个数等于满足
\[
aR\le (k+\tfrac12)\varphi^n\le bR
\]
的整数 \(k\) 的个数，因此为
\[
(b-a)R\varphi^{-n}+O(1),
\]
其中误差仅来自端点取整。只有当 \(\varphi^n\lesssim R\) 时这一层才可能贡献非零点，故有效层数为 \(O(\log R)\)。于是求和得到
\[
\#\bigl(\mathcal Z(\mathcal F_\varphi)\cap[aR,bR]\bigr)
=
(b-a)R\sum_{n=2}^{\infty}\varphi^{-n}+O(\log R).
\]
由于
\[
\sum_{n=2}^{\infty}\varphi^{-n}=1,
\]
遂得
\[
\#\bigl(\mathcal Z(\mathcal F_\varphi)\cap[aR,bR]\bigr)
=
(b-a)R+O(\log R).
\]
取 \([a,b]=[-1,1]\) 即得
\[
N_\varphi(R)=2R+O(\log R),
\]
这与定理 \ref{thm:xi-fold-resonance-zero-count-selfsimilar-linear} 的线性主项一致。

于是对任意 \([a,b]\subset(-1,1)\)，有
\[
\nu_R([a,b])
=
\frac{(b-a)R+O(\log R)}{2R+O(\log R)}
\longrightarrow
\frac{b-a}{2}.
\]
而 \(\frac12\mathbf 1_{[-1,1]}(x)\,\dd x\) 对区间端点不赋质量，故由 Portmanteau 定理可知
\[
\nu_R \Longrightarrow \frac12\,\mathbf 1_{[-1,1]}(x)\,\dd x.
\]
证毕。
\end{proof}

\begin{theorem}[局部化 solenoid 的连续自覆盖分类与精确度数公式]\label{thm:xi-localized-solenoid-self-cover-degree-classification}
设 \(S\subset\mathbb P\) 为有限素数集，并记
\[
G_S:=\ZZ[S^{-1}]
\]
及其 Pontryagin 对偶 \(\widehat{G_S}\)。则每个连续群自同态
\[
f:\widehat{G_S}\to \widehat{G_S}
\]
都唯一对应于某个 \(r\in G_S\)，其对偶离散同态为
\[
f^\vee(x)=rx
\qquad(x\in G_S).
\]
其中有限度连续自覆盖恰对应于 \(r\neq 0\) 的情形。对任意非零 \(r\in G_S\)，可唯一写成
\[
r=u\,n,
\qquad
u\in G_S^\times=\ZZ[S^{-1}]^\times,
\qquad
n\in\NN,
\qquad
\gcd\!\Bigl(n,\prod_{p\in S}p\Bigr)=1.
\]
若 \(f=f_r\) 为对应的连续自覆盖，则
\[
\boxed{\;
\ker(f)\cong \ZZ/n\ZZ,
\qquad
\deg(f)=n.
\;}
\]
特别地，
\[
\boxed{\;
\{\deg(f): f\ \text{为}\ \widehat{G_S}\ \text{的有限度连续自覆盖}\}
=
\{n\in\NN:\ p\mid n \Rightarrow p\notin S\},
\;}
\]
并且
\[
\boxed{\;
\operatorname{Aut}_{\mathrm{cts}}(\widehat{G_S})
\cong
\ZZ[S^{-1}]^\times
=
\left\{\pm\prod_{p\in S}p^{m_p}:m_p\in\ZZ\right\}.
\;}
\]
等价地，\(f\) 为同构当且仅当 \(n=1\)。
\end{theorem}
\begin{proof}
由定理 \ref{thm:xi-localized-endomorphism-ring-solenoid-indecomposable} 与推论 \ref{cor:xi-localized-integers-dual-automorphism-units}，
\[
\operatorname{End}_{\mathrm{cts}}(\widehat{G_S})
\cong
\operatorname{End}(G_S)^{\mathrm{op}}
\cong
G_S.
\]
因此每个连续自同态都唯一对应于某个 \(r\in G_S\)，并在对偶离散群上表现为乘法
\[
x\mapsto rx.
\]
若 \(r=0\)，则 \(f\) 为常值映射，显然不是自覆盖。以下设 \(r\neq 0\)。

由于 \(G_S=\ZZ[S^{-1}]\subset \QQ\)，\(r\) 的每个素数赋值都良定义；把属于 \(S\) 的全部素数赋值并入单位项，便得到唯一分解
\[
r=u\,n,
\]
其中 \(u\in G_S^\times\)，而 \(n\) 为正的 \(S\)-自由整数。由于乘以非零 \(r\) 在无挠群 \(G_S\subset\QQ\) 上单射，故 \(f^\vee\) 为单射。Pontryagin 对偶把单射送为满射，于是 \(f\) 为满同态。

又因 \(u\) 为单位，
\[
rG_S=u(nG_S)=nG_S.
\]
因此由定理 \ref{thm:xi-localized-finite-quotient-phase-layer-classification}，
\[
G_S/rG_S
\cong
G_S/nG_S
\cong
\ZZ/n\ZZ.
\]
对离散短正合列
\[
0\longrightarrow G_S\xrightarrow{\ \times r\ }G_S\longrightarrow G_S/rG_S\longrightarrow 0
\]
作 Pontryagin 对偶，得到
\[
0\longrightarrow (G_S/rG_S)^\wedge\longrightarrow \widehat{G_S}\xrightarrow{\ f\ }\widehat{G_S}\longrightarrow 0,
\]
从而
\[
\ker(f)\cong (G_S/rG_S)^\wedge\cong \ZZ/n\ZZ.
\]
特别地，\(\ker(f)\) 有限且
\[
\#\ker(f)=n.
\]
而紧交换群之间的连续满同态若核有限，则为有限度覆盖，其度数等于核的基数，故
\[
\deg(f)=\#\ker(f)=n.
\]

反之，对任意与 \(S\) 互素的正整数 \(n\)，乘法映射
\[
\times n:G_S\to G_S
\]
是单射，其对偶给出一个核同构于 \(\ZZ/n\ZZ\) 的连续自覆盖，因此上述度数集合恰好被全部 \(S\)-自由正整数实现。

最后，\(f\) 为同构当且仅当 \(\ker(f)\) 平凡，当且仅当 \(n=1\)。这等价于 \(r\in G_S^\times\)，故
\[
\operatorname{Aut}_{\mathrm{cts}}(\widehat{G_S})\cong G_S^\times
=
\left\{\pm\prod_{p\in S}p^{m_p}:m_p\in\ZZ\right\},
\]
与推论 \ref{cor:xi-localized-integers-dual-automorphism-units} 一致。证毕。
\end{proof}

\noindent
这些结果说明：Pisot--Bernoulli 卷积在任意非零 Fibonacci 递推方向上的 Fourier 采样都保留严格正的共振极限，因此非 Rajchman 障碍并非少数特殊子列现象；与之并行，Fibonacci 共振整函数虽由几何放大的半整数格层叠构成，其缩放零点计数测度却在宏观上趋于 \([-1,1]\) 上的均匀律；而有限素数局部化 solenoid 的全部有限度自覆盖，则完全由 \(S\)-自由整数度数控制，\(S\)-内素数只保留为零代价的单位自同构。




\paragraph{共振闭包级数、Pisot--Bernoulli 卷积与实零梯的若干解析后果}
在推论 \ref{cor:fold-resonance-ladder-fib-luc-limits} 的 Fibonacci/Lucas 方向极限、
定理 \ref{thm:fold-pisot-bernoulli-convolution-representation} 的 Fourier 采样表象以及
定理 \ref{thm:fold-resonance-entire-lp} 的 Laguerre--P\'olya 闭包已经建立之后，
还可进一步抽出下列若干条彼此兼容的结果。

\begin{theorem}[共振闭包级数的显式对数下界与 Fibonacci--Lucas 双子列强化]\label{thm:xi-resonance-closure-log-lower-bound-fib-luc}
\leanverified{paper\_xi\_resonance\_closure\_log\_lower\_bound\_fib\_luc}
沿用定义 \ref{def:fold-critical-resonance-ladder} 的记号，并令
\[
c_{\mathrm{Fib}}:=\lim_{n\to\infty}C_\varphi(F_n),
\qquad
c_{\mathrm{Luc}}:=\lim_{n\to\infty}C_\varphi(L_n),
\]
其中 \((F_n)_{n\ge 0}\) 与 \((L_n)_{n\ge 0}\) 分别为 Fibonacci 与 Lucas 数列。再定义部分共振闭包级数
\[
\Sigma_\varphi(N):=1+2\sum_{u=1}^{N}C_\varphi(u)^2.
\]
则有显式下界
\[
\boxed{\;
\liminf_{N\to\infty}
\frac{1}{\log N}\sum_{u=1}^{N}C_\varphi(u)^2
\ge
\frac{c_{\mathrm{Fib}}^2+c_{\mathrm{Luc}}^2}{\log\varphi}.\;}
\]
等价地，
\[
\boxed{\;
\Sigma_\varphi(N)\ge
\frac{2(c_{\mathrm{Fib}}^2+c_{\mathrm{Luc}}^2)}{\log\varphi}\,\log N+O(1).\;}
\]
因此，定理 \ref{thm:fold-sigma-phi-diverges} 中的抽象对数发散下界可由
Fibonacci 与 Lucas 两条递推方向的极限常数显式强化。
\end{theorem}

% AUTO-GENERATED by scripts/exp_fold_bulk_resonance_fibonacci_directional_limits.py
\[
\begin{aligned}
C_{\varphi,\mathrm{Fib}}^{(\le 200)}:=\Bigl|\cos\!\Bigl(\frac{\pi}{\sqrt5}\Bigr)\Bigr|\prod_{k=1}^{200}\Bigl|\cos\!\Bigl(\frac{\pi}{\sqrt5\,\varphi^{k}}\Bigr)\Bigr|^{2}&\approx 0.0424974234054296,\\
C_{\varphi,\mathrm{Luc}}^{(\le 200)}:=\prod_{k=1}^{200}\Bigl|\cos\bigl(\pi\varphi^{-k}\bigr)\Bigr|^{2}&\approx 0.0066134930353441.
\end{aligned}
\]
% AUTO-GENERATED by scripts/exp_fold_bulk_resonance_fibonacci_directional_limits.py
\[
\begin{aligned}
\frac{c_{\mathrm{Fib}}^2+c_{\mathrm{Luc}}^2}{\log\varphi}&\approx 0.003843981361014525,\\
\frac{2(c_{\mathrm{Fib}}^2+c_{\mathrm{Luc}}^2)}{\log\varphi}&\approx 0.007687962722029051.
\end{aligned}
\]

\begin{proof}
任取 \(\varepsilon>0\)。由推论 \ref{cor:fold-resonance-ladder-fib-luc-limits}，
存在 \(n_0(\varepsilon)\) 使得对一切 \(n\ge n_0(\varepsilon)\) 都有
\[
C_\varphi(F_n)^2\ge c_{\mathrm{Fib}}^2-\varepsilon,
\qquad
C_\varphi(L_n)^2\ge c_{\mathrm{Luc}}^2-\varepsilon.
\]
另一方面，对 \(n\ge 3\) 由恒等式
\[
L_n=F_{n-1}+F_{n+1}
\]
以及 \(F_{n-1}<F_n\) 得
\[
F_{n+1}<L_n<F_{n+1}+F_n=F_{n+2},
\]
故 \(L_n\) 在 \(n\ge 3\) 时不可能为 Fibonacci 数。于是两集合
\[
\{F_n:\ n\ge n_0(\varepsilon)\},
\qquad
\{L_n:\ n\ge \max\{3,n_0(\varepsilon)\}\}
\]
最终互不相交。

从而对充分大的 \(N\) 有
\[
\sum_{u=1}^{N}C_\varphi(u)^2
\ge
\sum_{\substack{n\ge n_0(\varepsilon)\\ F_n\le N}}C_\varphi(F_n)^2
+
\sum_{\substack{n\ge n_0(\varepsilon)\\ L_n\le N}}C_\varphi(L_n)^2,
\]
故
\[
\sum_{u=1}^{N}C_\varphi(u)^2
\ge
(c_{\mathrm{Fib}}^2-\varepsilon)\,\#\{n:\ F_n\le N\}
+
(c_{\mathrm{Luc}}^2-\varepsilon)\,\#\{n:\ L_n\le N\}
+O(1).
\]
由 Binet 公式
\[
F_n=\frac{\varphi^n-(-\varphi^{-1})^n}{\sqrt5},
\qquad
L_n=\varphi^n+(-\varphi^{-1})^n,
\]
可得
\[
\#\{n:\ F_n\le N\}=\frac{\log N}{\log\varphi}+O(1),
\qquad
\#\{n:\ L_n\le N\}=\frac{\log N}{\log\varphi}+O(1).
\]
代回并先令 \(N\to\infty\)，再令 \(\varepsilon\downarrow 0\)，便得到第一条盒装式。
第二条仅由 \(\Sigma_\varphi(N)=1+2\sum_{u\le N}C_\varphi(u)^2\) 立即推出。证毕。
\end{proof}

\begin{theorem}[共振闭包常数序列在任意有限 \texorpdfstring{$\ell^p$}{lp} 上均不可和]\label{thm:xi-resonance-closure-no-finite-lp-summability}
\leanverified{paper\_xi\_resonance\_closure\_no\_finite\_lp\_summability}
沿用定理 \ref{thm:xi-resonance-closure-log-lower-bound-fib-luc} 的记号。对任意
\[
0<p<\infty
\]
与任意
\[
0<\eta<c_{\mathrm{Fib}},
\]
定义阈值计数函数
\[
R_\eta(N):=\#\{1\le u\le N:\ C_\varphi(u)\ge \eta\}.
\]
则存在整数 \(n_0(\eta)\) 使得
\[
R_\eta(N)\ge \#\{n\ge n_0(\eta):\ F_n\le N\}
=
\frac{\log N}{\log\varphi}+O(1).
\]
因此
\[
\sum_{u=1}^{N}C_\varphi(u)^p
\ge
\eta^p\frac{\log N}{\log\varphi}+O(1),
\]
从而
\[
\bigl(C_\varphi(u)\bigr)_{u\in\ZZ}\notin \ell^p(\ZZ)
\qquad
(\forall\,0<p<\infty).
\]
\end{theorem}
\begin{proof}
由推论 \ref{cor:fold-resonance-ladder-fib-luc-limits}，
\[
C_\varphi(F_n)\longrightarrow c_{\mathrm{Fib}}>0.
\]
故对任意 \(0<\eta<c_{\mathrm{Fib}}\)，存在 \(n_0(\eta)\) 使得当 \(n\ge n_0(\eta)\) 时恒有
\[
C_\varphi(F_n)\ge \eta.
\]
于是对每个 \(N\) 都有
\[
R_\eta(N)\ge \#\{n\ge n_0(\eta):\ F_n\le N\}.
\]
再由 Binet 公式可知 \(F_n\asymp \varphi^n\)，从而
\[
\#\{n:\ F_n\le N\}
=
\frac{\log N}{\log\varphi}+O(1).
\]
这就得到阈值计数下界。

随后由定义直接有
\[
\sum_{u=1}^{N}C_\varphi(u)^p
\ge
\sum_{\substack{1\le u\le N\\ C_\varphi(u)\ge \eta}}C_\varphi(u)^p
\ge
\eta^p R_\eta(N),
\]
代入上式即得
\[
\sum_{u=1}^{N}C_\varphi(u)^p
\ge
\eta^p\frac{\log N}{\log\varphi}+O(1).
\]
令 \(N\to\infty\)，可见正半轴上的 \(p\)-次幂和已经发散，因此双侧序列
\[
\bigl(C_\varphi(u)\bigr)_{u\in\ZZ}
\]
不属于任何有限 \(\ell^p(\ZZ)\)。证毕。
\end{proof}

\begin{theorem}[共振闭包常数不存在最终单调的一阶包络]\label{thm:xi-resonance-closure-no-eventually-monotone-envelope}
沿用推论 \ref{cor:fold-resonance-ladder-no-decay-no-limit} 的记号。若
\[
r:\NN\to\RR
\]
是最终单调序列，则不可能成立
\[
C_\varphi(u)-r(u)\longrightarrow 0
\qquad (u\to\infty).
\]
更强地，若 \(r(u)>0\) 对充分大的 \(u\) 成立，则也不可能有
\[
\frac{C_\varphi(u)}{r(u)}\longrightarrow 1
\qquad (u\to\infty).
\]
\end{theorem}
\begin{proof}
先证第一条。若存在最终单调的 \(r(u)\) 使
\[
C_\varphi(u)-r(u)\longrightarrow 0,
\]
则由于 \(0\le C_\varphi(u)\le 1\)，序列 \(r(u)\) 亦最终有界。最终单调且有界的实序列必收敛，设
\[
r(u)\longrightarrow \ell.
\]
于是
\[
C_\varphi(u)=r(u)+o(1)\longrightarrow \ell,
\]
这与推论 \ref{cor:fold-resonance-ladder-no-decay-no-limit} 所给出的“整数方向不存在单一极限”矛盾。故第一条不可能。

再证第二条。若 \(r(u)>0\) 最终成立且
\[
\frac{C_\varphi(u)}{r(u)}\longrightarrow 1,
\]
则对充分大的 \(u\) 有
\[
\frac12\le \frac{C_\varphi(u)}{r(u)}\le \frac32.
\]
由于 \(0\le C_\varphi(u)\le 1\)，可得最终
\[
0<r(u)\le 2.
\]
因此 \(r(u)\) 仍是最终单调且最终有界的序列，故必收敛。另一方面，
\[
C_\varphi(u)-r(u)
=
r(u)\left(\frac{C_\varphi(u)}{r(u)}-1\right)\longrightarrow 0.
\]
这就回到了第一条已排除的情形，矛盾。证毕。
\end{proof}

\begin{theorem}[Pisot--Bernoulli 卷积的偶阶累积量闭式]\label{thm:xi-pisot-bernoulli-even-cumulants-closed}
令 \(q=\varphi^{-1}\in(0,1)\)，并定义随机级数
\[
X=\sum_{n=2}^{\infty}\varepsilon_n q^n,
\qquad
\mathbb P(\varepsilon_n=\pm 1)=\frac12,
\]
其分布记为 \(\mu_q^{(2)}\)。则对每个整数 \(r\ge 1\)，
\(X\) 的第 \(2r\) 阶累积量满足显式公式
\[
\boxed{\;
\kappa_{2r}(X)
=
(-1)^{r-1}
\frac{2^{2r}(2^{2r}-1)|B_{2r}|}{2r}
\cdot
\frac{\varphi^{-4r}}{1-\varphi^{-2r}}.\;}
\]
特别地，
\[
\kappa_2(X)=\varphi^{-3},
\qquad
\kappa_4(X)=-2\,\frac{\varphi^{-8}}{1-\varphi^{-4}},
\qquad
\kappa_6(X)=16\,\frac{\varphi^{-12}}{1-\varphi^{-6}}.
\]
\end{theorem}
\begin{proof}
由定理 \ref{thm:fold-pisot-bernoulli-convolution-representation}，
\[
\mathcal F_\varphi(z)
=
\prod_{n=2}^{\infty}\cos(\pi z\varphi^{-n})
=
\widehat{\mu_q^{(2)}}(\pi z).
\]
于是
\[
\log \mathcal F_\varphi(z)
=
\sum_{r\ge 1}\kappa_{2r}(X)\frac{(i\pi z)^{2r}}{(2r)!}
=
\sum_{r\ge 1}(-1)^r\kappa_{2r}(X)\frac{\pi^{2r}z^{2r}}{(2r)!},
\]
其中奇阶项由对称性自动消失。

当 \(|z|<\varphi^2/2\) 时，对所有 \(n\ge 2\) 都有
\(
|\pi z\varphi^{-n}|<\pi/2
\)，
故可用经典展开
\[
\log\cos x
=
-\sum_{r=1}^{\infty}
\frac{2^{2r}(2^{2r}-1)|B_{2r}|}{2r(2r)!}\,x^{2r}
\qquad (|x|<\pi/2)
\]
逐项求和：
\[
\log \mathcal F_\varphi(z)
=
-\sum_{n=2}^{\infty}\sum_{r=1}^{\infty}
\frac{2^{2r}(2^{2r}-1)|B_{2r}|}{2r(2r)!}
(\pi z\varphi^{-n})^{2r}.
\]
由于
\[
\sum_{n=2}^{\infty}\sum_{r=1}^{\infty}
\left|
\frac{2^{2r}(2^{2r}-1)|B_{2r}|}{2r(2r)!}
(\pi z\varphi^{-n})^{2r}
\right|
<\infty,
\]
可以交换求和顺序，得到
\[
\log \mathcal F_\varphi(z)
=
-\sum_{r=1}^{\infty}
\frac{2^{2r}(2^{2r}-1)|B_{2r}|}{2r(2r)!}
\pi^{2r}z^{2r}
\sum_{n=2}^{\infty}\varphi^{-2rn}.
\]
又
\[
\sum_{n=2}^{\infty}\varphi^{-2rn}
=
\frac{\varphi^{-4r}}{1-\varphi^{-2r}}.
\]
比较 \(\log \mathcal F_\varphi(z)\) 的 \(z^{2r}\) 系数，即得
\[
(-1)^r\kappa_{2r}(X)\frac{\pi^{2r}}{(2r)!}
=
-
\frac{2^{2r}(2^{2r}-1)|B_{2r}|}{2r(2r)!}
\pi^{2r}
\cdot
\frac{\varphi^{-4r}}{1-\varphi^{-2r}},
\]
整理便得到盒装公式。将 \(r=1,2,3\) 代入即得三条特例。证毕。
\end{proof}

\begin{theorem}[Pisot--Bernoulli 卷积偶矩序列的 Tur\'an--Newton 约束]\label{thm:xi-pisot-bernoulli-even-moment-turan-newton}
沿用定理 \ref{thm:xi-pisot-bernoulli-even-cumulants-closed} 的记号，并记
\[
m_{2r}:=\int_{\RR}x^{2r}\,d\mu_q^{(2)}(x)
\qquad (r\ge 0).
\]
则对每个整数 \(r\ge 1\) 都有严格不等式
\[
\boxed{\;
m_{2r}^2
>
\frac{(2r)(2r-1)}{(2r+1)(2r+2)}\,
m_{2r-2}m_{2r+2}.\;}
\]
等价地，序列
\[
\gamma_r:=\frac{\pi^{2r}}{(2r)!}\,m_{2r}
\qquad (r\ge 0)
\]
满足严格对数凹性
\[
\gamma_r^2>\gamma_{r-1}\gamma_{r+1}
\qquad (r\ge 1).
\]
\end{theorem}
\begin{proof}
由定理 \ref{thm:fold-resonance-entire-lp}，\(\mathcal F_\varphi\) 的全部零点均为实单零点，
且关于原点成对对称。将正零点按升序记为 \(\{\xi_j\}_{j\ge 1}\subset (0,\infty)\)，则
\[
\mathcal F_\varphi(z)=\prod_{j=1}^{\infty}\left(1-\frac{z^2}{\xi_j^2}\right)
\]
在紧集上一致收敛。另一方面，由
\[
\mathcal F_\varphi(z)=\widehat{\mu_q^{(2)}}(\pi z)
=
\sum_{r=0}^{\infty}(-1)^r\frac{\pi^{2r}m_{2r}}{(2r)!}\,z^{2r},
\]
可知
\[
\gamma_r=\frac{\pi^{2r}}{(2r)!}\,m_{2r}
\]
正是该偶幂级数的系数。

对每个 \(N>r\)，定义有限截断
\[
P_N(z):=\prod_{j=1}^{N}\left(1-\frac{z^2}{\xi_j^2}\right)
=
\sum_{s=0}^{N}(-1)^s\gamma_s^{(N)}z^{2s}.
\]
则 \(\gamma_s^{(N)}\) 是正数列 \((\xi_1^{-2},\dots,\xi_N^{-2})\) 的第 \(s\) 个初等对称多项式。
Newton 不等式给出
\[
\bigl(\gamma_r^{(N)}\bigr)^2
\ge
\frac{r+1}{r}\frac{N-r+1}{N-r}\,
\gamma_{r-1}^{(N)}\gamma_{r+1}^{(N)}
>
\gamma_{r-1}^{(N)}\gamma_{r+1}^{(N)}.
\]
又因 \(P_N\to \mathcal F_\varphi\) 在零点外的紧集上一致，从而幂级数系数逐项收敛：
\[
\gamma_s^{(N)}\xrightarrow[N\to\infty]{}\gamma_s
\qquad (s\ \text{固定}).
\]
令 \(N\to\infty\)，得到
\[
\gamma_r^2\ge \frac{r+1}{r}\,\gamma_{r-1}\gamma_{r+1}
>
\gamma_{r-1}\gamma_{r+1},
\]
这就证明了 \((\gamma_r)\) 的严格对数凹性。

最后把
\[
\gamma_r=\frac{\pi^{2r}}{(2r)!}\,m_{2r}
\]
代回，可得
\[
\left(\frac{\pi^{2r}}{(2r)!}m_{2r}\right)^2
\ge
\frac{\pi^{2r-2}}{(2r-2)!}m_{2r-2}\cdot
\frac{\pi^{2r+2}}{(2r+2)!}m_{2r+2}.
\]
约去 \(\pi^{4r}\) 并整理阶乘比
\[
\frac{(2r)!^2}{(2r-2)!(2r+2)!}
=
\frac{(2r)(2r-1)}{(2r+1)(2r+2)},
\]
即得所示矩不等式。证毕。
\end{proof}

\begin{theorem}[共振整函数的精确零点计数与 Weyl 型主项]\label{cor:xi-resonance-entire-zero-count-explicit-linear-density}
\leanverified{paper\_xi\_resonance\_entire\_zero\_count\_explicit\_linear\_density}
沿用定理 \ref{thm:fold-resonance-entire-lp} 的记号，并定义
\[
N_\varphi(R):=\#\{z\in \mathcal Z(\mathcal F_\varphi):\ |z|\le R\}.
\]
则对每个 \(R>0\) 都有精确公式
\[
\boxed{\;
N_\varphi(R)=
2\sum_{n\ge 2}\left\lfloor \frac{R}{\varphi^n}+\frac12\right\rfloor.\;}
\]
特别地，
\[
\boxed{\;
N_\varphi(R)=2R+O(\log R)
\qquad (R\to\infty).\;}
\]
因此 \(\mathcal F_\varphi\) 的零点收敛指数恰为 \(1\)。
\end{theorem}
\begin{proof}
由定理 \ref{thm:fold-resonance-entire-lp}，
\[
\mathcal Z(\mathcal F_\varphi)
=
\Bigl\{\bigl(k+\tfrac12\bigr)\varphi^n:\ k\in\ZZ,\ n\ge 2\Bigr\}.
\]
对固定 \(n\ge 2\)，第 \(n\) 层零点恰为
\[
\Bigl\{\bigl(k+\tfrac12\bigr)\varphi^n:\ k\in\ZZ\Bigr\}.
\]
其中满足
\[
\bigl|\bigl(k+\tfrac12\bigr)\varphi^n\bigr|\le R
\]
等价于
\[
\left|k+\frac12\right|\le R\varphi^{-n}.
\]
满足该不等式的整数 \(k\) 个数恰为
\[
2\left\lfloor \frac{R}{\varphi^n}+\frac12\right\rfloor.
\]
于是对全部 \(n\ge 2\) 求和便得到所示精确公式；其中当 \(R\varphi^{-n}<\frac12\) 时，对应项自动为 \(0\)，故该无穷和实际上只有有限个非零项。

再由
\[
2\left\lfloor \frac{R}{\varphi^n}+\frac12\right\rfloor
=
\frac{2R}{\varphi^n}+O(1)
\]
对满足 \(R\varphi^{-n}\ge \frac12\) 的指标求和。记
\[
L(R):=\left\lfloor \log_\varphi(2R)\right\rfloor,
\]
则对充分大的 \(R\) 有
\[
N_\varphi(R)
=
2R\sum_{n=2}^{L(R)}\varphi^{-n}
+O(\log R).
\]
而
\[
\sum_{n=2}^{\infty}\varphi^{-n}
=
\frac{\varphi^{-2}}{1-\varphi^{-1}}
=
1,
\]
故
\[
N_\varphi(R)=2R+O(\log R).
\]
零点收敛指数为 \(1\) 因而立即成立；这也与推论
\ref{cor:xi-fold-resonance-zero-density-exponent-one} 的结论一致。证毕。
\end{proof}

\paragraph{Cartwright 精确型、周期化不可平滑化与对数导数的极点骨架}
在共振整函数 \(\mathcal F_\varphi\) 的精确零点计数已经确定之后，Pisot--Bernoulli 周期化测度与其对数导数的解析结构还可进一步压缩为下列四条直接后果。它们分别刻画 \(\mathcal F_\varphi\) 的 Cartwright 精确型、模 \(2\) 周期化测度任意有限卷积幂的不可平滑化、周期化 Fourier 能量的显式对数下界，以及对数导数在极点层面的主值型 Mittag--Leffler 骨架。

\begin{corollary}[共振整函数属于 Cartwright 类，且指数型恰为 \texorpdfstring{$\pi$}{pi}]\label{cor:xi-time-part59ab-resonance-cartwright-exact-type}
沿用定理 \ref{thm:fold-resonance-entire-lp} 与定理 \ref{cor:xi-resonance-entire-zero-count-explicit-linear-density} 的记号，则
\[
\boxed{\;
\mathcal F_\varphi \text{ 属于 Cartwright 类，且其指数型 }
\tau(\mathcal F_\varphi)=\pi.\;}
\]
进一步，其指示函数满足
\[
\boxed{\;
h_{\mathcal F_\varphi}(\theta)=\pi|\sin\theta|
\qquad (\theta\in\RR).\;}
\]
\end{corollary}
\begin{proof}
由定理 \ref{thm:fold-resonance-entire-lp} 的全平面增长界
\[
|\mathcal F_\varphi(z)|\le \exp(\pi|\Im z|)
\]
与实轴上
\[
|\mathcal F_\varphi(x)|\le 1
\qquad (x\in\RR)
\]
可得
\[
\int_{\RR}\frac{\log^+|\mathcal F_\varphi(x)|}{1+x^2}\,\dd x=0,
\]
故 \(\mathcal F_\varphi\) 属于 Cartwright 类。另一方面，定理
\ref{cor:xi-resonance-entire-zero-count-explicit-linear-density} 已给出
\[
N_\varphi(R)=2R+O(\log R).
\]
对 Cartwright 类整函数应用 Levinson--Cartwright 的零点密度定理，若其指数型为 \(\tau\)，则
\[
N_\varphi(R)=\frac{2\tau}{\pi}R+o(R).
\]
比较主项系数即得 \(\tau=\pi\)。最后，由附录定理
\ref{thm:fold-resonance-cartwright-indicator}，
\[
h_{\mathcal F_\varphi}(\theta)=\pi|\sin\theta|.
\]
证毕。
\end{proof}

\begin{theorem}[Pisot--Bernoulli 周期化测度的任意有限卷积幂都不会平滑化]\label{thm:xi-time-part59ab-periodized-bernoulli-no-finite-smoothing}
\leanverified{paper\_xi\_time\_part59ab\_periodized\_bernoulli\_no\_finite\_smoothing}
令
\[
q:=\varphi^{-1},
\qquad
\nu:=\bigl(\RR\to \mathbb T_2=\RR/(2\ZZ)\bigr)_\#\mu_q^{(2)},
\]
并对每个整数 \(r\ge 1\) 记 \(\nu^{*r}\) 为 \(\nu\) 在 \(\mathbb T_2\) 上的 \(r\)-重卷积。再记
\[
c_{\mathrm{Fib}}:=\lim_{n\to\infty}C_\varphi(F_n)>0,
\qquad
c_{\mathrm{Luc}}:=\lim_{n\to\infty}C_\varphi(L_n)>0.
\]
则对任意整数 \(r\ge 1\) 都有
\[
\boxed{\;
\bigl|\widehat{\nu^{*r}}(F_n)\bigr|\longrightarrow c_{\mathrm{Fib}}^{\,r},
\qquad
\bigl|\widehat{\nu^{*r}}(L_n)\bigr|\longrightarrow c_{\mathrm{Luc}}^{\,r}.\;}
\]
从而：
\begin{enumerate}
  \item \(\nu^{*r}\) 对任意 \(r\ge 1\) 都不是 Rajchman 测度；
  \item \(\nu^{*r}\) 对任意 \(r\ge 1\) 都不可能具有任何 \(L^p(\mathbb T_2)\) 密度，其中 \(1\le p\le\infty\)。
\end{enumerate}
\end{theorem}
\begin{proof}
由定理 \ref{thm:xi-pisot-bernoulli-fibonacci-direction-nonrajchman} 的证明，
\[
\widehat\nu(u)=\widehat{\mu_q^{(2)}}(\pi u)
\qquad (u\in\ZZ).
\]
而卷积的 Fourier 乘法律给出
\[
\widehat{\nu^{*r}}(u)=\widehat\nu(u)^r
\qquad (u\in\ZZ).
\]
于是
\[
\bigl|\widehat{\nu^{*r}}(u)\bigr|
=
\bigl|\widehat\nu(u)\bigr|^r
=
\bigl|\widehat{\mu_q^{(2)}}(\pi u)\bigr|^r.
\]
再由推论 \ref{cor:fold-resonance-ladder-fib-luc-limits}，
\[
\bigl|\widehat\nu(F_n)\bigr|=C_\varphi(F_n)\longrightarrow c_{\mathrm{Fib}}>0,
\qquad
\bigl|\widehat\nu(L_n)\bigr|=C_\varphi(L_n)\longrightarrow c_{\mathrm{Luc}}>0,
\]
故
\[
\bigl|\widehat{\nu^{*r}}(F_n)\bigr|\longrightarrow c_{\mathrm{Fib}}^{\,r}>0,
\qquad
\bigl|\widehat{\nu^{*r}}(L_n)\bigr|\longrightarrow c_{\mathrm{Luc}}^{\,r}>0.
\]
因此 \(\widehat{\nu^{*r}}(u)\) 沿无穷子列不趋于 \(0\)，从而 \(\nu^{*r}\) 不是 Rajchman 测度。

若 \(\nu^{*r}\) 具有某个 \(L^1(\mathbb T_2)\) 密度，则由 Riemann--Lebesgue 引理，其 Fourier 系数必满足
\[
\widehat{\nu^{*r}}(u)\longrightarrow 0
\qquad (|u|\to\infty),
\]
与上式矛盾。另一方面，\(\mathbb T_2\) 的 Haar 测度有限，故对任意 \(1\le p\le\infty\) 都有
\[
L^p(\mathbb T_2)\subseteq L^1(\mathbb T_2).
\]
因此任何 \(L^p\) 密度也不可能存在。证毕。
\end{proof}

\begin{corollary}[周期化 Fourier 能量的对数发散具有显式 Fibonacci 下界]\label{cor:xi-time-part59ab-periodized-log-energy-fib}
沿用定理 \ref{thm:xi-time-part59ab-periodized-bernoulli-no-finite-smoothing} 的记号，并定义
\[
\mathcal E(N):=\sum_{|u|\le N}|\widehat\nu(u)|^2.
\]
则有定量下界
\[
\boxed{\;
\liminf_{N\to\infty}\frac{\mathcal E(N)}{\log N}
\ge
\frac{2c_{\mathrm{Fib}}^2}{\log\varphi}.\;}
\]
其中
\[
\boxed{\;
c_{\mathrm{Fib}}
=
\Bigl|\cos\!\Bigl(\frac{\pi}{\sqrt5}\Bigr)\Bigr|
\prod_{k=1}^{\infty}
\Bigl|\cos\!\Bigl(\frac{\pi}{\sqrt5\,\varphi^k}\Bigr)\Bigr|^2
>0.\;}
\]
特别地，
\[
\mathcal E(N)\neq o(\log N).
\]
\end{corollary}
\begin{proof}
由推论 \ref{cor:fold-resonance-ladder-fib-luc-limits}，
\[
\bigl|\widehat\nu(F_n)\bigr|
=
C_\varphi(F_n)\longrightarrow c_{\mathrm{Fib}}>0,
\]
且 \(c_{\mathrm{Fib}}\) 的显式乘积公式正是该推论中的第一式。故对任意
\(\varepsilon\in(0,c_{\mathrm{Fib}})\)，存在 \(n_0(\varepsilon)\) 使得当 \(n\ge n_0(\varepsilon)\) 时都有
\[
\bigl|\widehat\nu(F_n)\bigr|\ge c_{\mathrm{Fib}}-\varepsilon.
\]
再由 \(|\widehat\nu(-u)|=|\widehat\nu(u)|\)，对充分大的 \(N\) 有
\[
\mathcal E(N)
\ge
2\sum_{\substack{n\ge n_0(\varepsilon)\\ F_n\le N}}
\bigl|\widehat\nu(F_n)\bigr|^2
\ge
2(c_{\mathrm{Fib}}-\varepsilon)^2\,
\#\{n\ge n_0(\varepsilon):\ F_n\le N\}.
\]
由 Binet 公式
\[
F_n=\frac{\varphi^n-\psi^n}{\sqrt5},
\qquad
\psi:=-\varphi^{-1},
\]
可得
\[
\#\{n:\ F_n\le N\}=\frac{\log N}{\log\varphi}+O(1).
\]
于是
\[
\liminf_{N\to\infty}\frac{\mathcal E(N)}{\log N}
\ge
\frac{2(c_{\mathrm{Fib}}-\varepsilon)^2}{\log\varphi}.
\]
令 \(\varepsilon\downarrow 0\) 即得结论。证毕。
\end{proof}

\begin{theorem}[共振整函数对数导数的自相似 Riccati 递归与主值 Mittag--Leffler 骨架]\label{thm:xi-time-part59ab-resonance-logderivative-mittagleffler}
沿用定理 \ref{thm:fold-resonance-entire-lp} 的记号，并在
\(\CC\setminus\mathcal Z(\mathcal F_\varphi)\) 上定义
\[
L_\varphi(z):=\frac{\mathcal F_\varphi'(z)}{\mathcal F_\varphi(z)}.
\]
则成立严格递归
\[
\boxed{\;
L_\varphi(z)=
-\pi\varphi^{-2}\tan(\pi z\varphi^{-2})
+\varphi^{-1}L_\varphi(z/\varphi).\;}
\]
对任意整数 \(N\ge 2\) 都有迭代公式
\[
L_\varphi(z)=
-\pi\sum_{n=2}^{N}\varphi^{-n}\tan(\pi z\varphi^{-n})
+\varphi^{-(N-1)}L_\varphi(z/\varphi^{N-1}).
\]
并在离零点集的紧集上一致收敛到
\[
\boxed{\;
L_\varphi(z)=
-\pi\sum_{n=2}^{\infty}\varphi^{-n}\tan(\pi z\varphi^{-n}),
\qquad
z\in\CC\setminus\mathcal Z(\mathcal F_\varphi).\;}
\]
进一步，对每个固定 \(n\) 的内层求和按 \(k\in\ZZ\) 的对称主值理解时，有主值型 Mittag--Leffler 展开
\[
\boxed{\;
L_\varphi(z)=
-\sum_{n=2}^{\infty}
\operatorname{v.p.}\!\sum_{k\in\ZZ}
\frac{1}{z-(k+\tfrac12)\varphi^n},
\qquad
z\in\CC\setminus\mathcal Z(\mathcal F_\varphi).\;}
\]
\end{theorem}
\begin{proof}
前三式正是定理 \ref{thm:xi-fold-resonance-logderivative-selfsimilar} 的内容；这里只是把其中的 \(\Psi_\varphi\) 改记为 \(L_\varphi\)，并把迭代步数改写为 \(N-1\) 次。

对每个固定 \(n\ge 2\)，经典部分分式恒等式
\[
\pi\tan(\pi w)=\operatorname{v.p.}\!\sum_{k\in\ZZ}\frac{1}{w-(k+\tfrac12)}
\]
在 \(w=z\varphi^{-n}\) 处给出
\[
-\pi\varphi^{-n}\tan(\pi z\varphi^{-n})
=
-\operatorname{v.p.}\!\sum_{k\in\ZZ}
\frac{1}{z-(k+\tfrac12)\varphi^n}.
\]
再对 \(n\ge 2\) 求和，并利用上一步已经得到的 tangent 级数在
\(\CC\setminus\mathcal Z(\mathcal F_\varphi)\) 的紧集上局部一致收敛，便得到所示主值型 Mittag--Leffler 展开。证毕。
\end{proof}


\paragraph{全微态阈值的极端矩恢复、冻结压力分解与零温地基层分裂}
在全微态精确反演的有限层比特阈值、冻结相的双指数分离、原子 Euler 因子的 primitive 剥离以及零温四态地基层闭式已经建立之后，还可进一步抽取出下列六条跨模块后果。

\begin{corollary}[全可逆阈值的极端矩重建]\label{cor:xi-full-inversion-threshold-extreme-moment-recovery}
沿用定理 \ref{thm:xi-full-microstate-exact-inversion-bitrate-threshold} 的记号。对每个固定分辨率 \(m\) 与每个 \(q>0\)，记
\[
S_q(m):=\sum_{x\in X_m}d_m(x)^q.
\]
则有
\[
\boxed{\;
M_m=\lim_{q\to\infty}S_q(m)^{1/q}.\;}
\]
从而
\[
\boxed{\;
B_m^{\mathrm{full}}
=
\left\lceil
\log_2\!\left(\lim_{q\to\infty}S_q(m)^{1/q}\right)
\right\rceil.\;}
\]
因此，全微态精确反演阈值不仅由最大纤维规模 \(M_m\) 唯一确定，也由正阶幂和矩谱 \(\{S_q(m)\}_{q>0}\) 的端点极限唯一恢复。
\end{corollary}
\begin{proof}
对固定 \(m\)，向量
\[
\mathbf d^{(m)}:=\bigl(d_m(x)\bigr)_{x\in X_m}
\]
只有有限个坐标。有限维范数极限定理给出
\[
\lim_{q\to\infty}\left(\sum_{x\in X_m}d_m(x)^q\right)^{1/q}
=
\max_{x\in X_m}d_m(x)
=
M_m.
\]
这便得到第一式。再将其代入定理 \ref{thm:xi-full-microstate-exact-inversion-bitrate-threshold} 的恒等式
\[
B_m^{\mathrm{full}}=\left\lceil \log_2 M_m\right\rceil
\]
即可得到第二式。证毕。
\end{proof}

\begin{theorem}[冻结相压力的地址成本--残余简并分解]\label{thm:xi-fixed-freezing-pressure-address-cost-residual-decomposition}
沿用定理 \ref{thm:xi-full-microstate-exact-inversion-bitrate-threshold}、定理 \ref{thm:fold-pressure-freezing-threshold} 与定理 \ref{thm:xi-fixed-freezing-escort-renyi-spectrum-collapse} 的记号。对任意 \(a>a_{\mathrm c}\)，记
\[
\beta_{\mathrm{inv}}
:=
\lim_{m\to\infty}\frac{B_m^{\mathrm{full}}}{m}
=
\frac{\alpha_*}{\log 2},
\]
并记冻结相的 min-entropy 率为
\[
h_\infty(a)
:=
\lim_{m\to\infty}\frac{1}{m}
\Bigl(-\log\max_{x\in X_m}\pi_{m,a}(x)\Bigr)
=
g_*.
\]
则冻结相压力满足严格分解
\[
\boxed{\;
P_a
=
a(\log 2)\,\beta_{\mathrm{inv}}+h_\infty(a).\;}
\]
换言之，冻结相压力恰分解为“实现全微态精确反演所需的指数地址成本”与“冻结端面的残余简并熵率”两部分。
\end{theorem}
\begin{proof}
由定理 \ref{thm:fold-pressure-freezing-threshold}，对任意 \(a>a_{\mathrm c}\) 都有
\[
P_a=a\alpha_*+g_*.
\]
另一方面，定理 \ref{thm:xi-full-microstate-exact-inversion-bitrate-threshold} 给出
\[
\beta_{\mathrm{inv}}=\frac{\alpha_*}{\log 2},
\]
而定理 \ref{thm:xi-fixed-freezing-escort-renyi-spectrum-collapse} 在 \(q=\infty\) 时给出
\[
h_\infty(a)=g_*.
\]
将这两式代回即可得到
\[
P_a
=
a\alpha_*+g_*
=
a(\log 2)\,\beta_{\mathrm{inv}}+h_\infty(a).
\]
证毕。
\end{proof}

\begin{corollary}[冻结相的斜率--截距反演公式]\label{cor:xi-fixed-freezing-slope-intercept-inversion}
沿用定理 \ref{thm:xi-fixed-freezing-pressure-address-cost-residual-decomposition} 的记号。则有
\[
\boxed{\;
\lim_{q\to\infty}\frac{P_q}{q}=\alpha_*,
\qquad
\beta_{\mathrm{inv}}
=
\frac{1}{\log 2}\lim_{q\to\infty}\frac{P_q}{q}.\;}
\]
并且对任意 \(a>a_{\mathrm c}\)，成立
\[
\boxed{\;
h_\infty(a)
=
P_a-a\lim_{q\to\infty}\frac{P_q}{q}.\;}
\]
因此，冻结相的截距并非独立参数，而是由压力谱减去其无穷远斜率后精确恢复。
\end{corollary}
\begin{proof}
由定理 \ref{thm:xi-time-part57b-pressure-spectrum-slope-intercept-recovery}，
\[
\lim_{q\to\infty}\frac{P_q}{q}=\alpha_*.
\]
再与
\[
\beta_{\mathrm{inv}}=\frac{\alpha_*}{\log 2}
\]
合并，便得第二式。

另一方面，定理
\ref{thm:xi-fixed-freezing-pressure-address-cost-residual-decomposition}
给出
\[
P_a=a\alpha_*+h_\infty(a).
\]
将 \(\alpha_*\) 以
\[
\lim_{q\to\infty}\frac{P_q}{q}
\]
替换，即得
\[
h_\infty(a)
=
P_a-a\lim_{q\to\infty}\frac{P_q}{q}.
\]
证毕。
\end{proof}

\begin{theorem}[冻结相中高阶幂和比值的指数锁定]\label{thm:xi-fixed-freezing-power-sum-ratio-locking}
在严格基态间隙
\[
\alpha_*^{(2)}<\alpha_*
\]
下，沿用定理 \ref{thm:xi-fixed-freezing-escort-renyi-spectrum-collapse} 的记号。记
\[
M_m:=\max_{x\in X_m}d_m(x),
\qquad
K_m:=\#\{x\in X_m:\ d_m(x)=M_m\},
\]
\[
M_m^{(2)}:=\max\{d_m(x):\ d_m(x)<M_m\},
\qquad
S_t(m):=\sum_{x\in X_m}d_m(x)^t.
\]
若次大纤维不存在，则约定 \(M_m^{(2)}:=1\)。对任意固定整数 \(k\ge 1\) 与任意
\[
a>a_{\mathrm c}:=\frac{\log\varphi-g_*}{\alpha_*-\alpha_*^{(2)}},
\]
都有
\[
\boxed{\;
\frac{S_{a+k}(m)}{S_a(m)}
=
M_m^k\Bigl(1+O\bigl(e^{-m\Delta(a)+o(m)}\bigr)\Bigr),\;}
\]
其中
\[
\Delta(a):=a(\alpha_*-\alpha_*^{(2)})-(\log\varphi-g_*)>0.
\]
\end{theorem}
\begin{proof}
把
\[
R_t(m):=\sum_{x:\,d_m(x)<M_m}d_m(x)^t
\]
记作基态外余项，则对任意 \(t>0\) 都有精确分解
\[
S_t(m)=K_mM_m^t+R_t(m).
\]
并且
\[
0\le R_t(m)\le |X_m|\,(M_m^{(2)})^t.
\]
由
\[
M_m=e^{\alpha_* m+o(m)},
\qquad
K_m=e^{g_* m+o(m)},
\qquad
|X_m|=e^{(\log\varphi)m+o(m)},
\]
以及
\[
\limsup_{m\to\infty}\frac1m\log M_m^{(2)}=\alpha_*^{(2)},
\]
可得
\[
\frac{R_a(m)}{K_mM_m^a}
\le
\exp\!\bigl(-m\Delta(a)+o(m)\bigr).
\]
因此
\[
S_a(m)=K_mM_m^a\Bigl(1+O\bigl(e^{-m\Delta(a)+o(m)}\bigr)\Bigr).
\]
同理，
\[
\frac{R_{a+k}(m)}{K_mM_m^{a+k}}
\le
\exp\!\bigl(-m\Delta(a+k)+o(m)\bigr)
\le
\exp\!\bigl(-m\Delta(a)+o(m)\bigr),
\]
因为
\[
\Delta(a+k)=\Delta(a)+k(\alpha_*-\alpha_*^{(2)})\ge \Delta(a).
\]
故
\[
S_{a+k}(m)=K_mM_m^{a+k}\Bigl(1+O\bigl(e^{-m\Delta(a)+o(m)}\bigr)\Bigr).
\]
两式相除便得到
\[
\frac{S_{a+k}(m)}{S_a(m)}
=
M_m^k\,
\frac{1+O\bigl(e^{-m\Delta(a)+o(m)}\bigr)}
{1+O\bigl(e^{-m\Delta(a)+o(m)}\bigr)}
=
M_m^k\Bigl(1+O\bigl(e^{-m\Delta(a)+o(m)}\bigr)\Bigr).
\]
证毕。
\end{proof}

\begin{corollary}[冻结 escort 测度的正幂矩塌缩]\label{cor:xi-fixed-freezing-escort-positive-moment-collapse}
沿用定理 \ref{thm:xi-fixed-freezing-power-sum-ratio-locking} 的记号。对任意固定整数 \(k\ge 1\)，若
\[
X\sim \pi_{m,a},
\qquad
\pi_{m,a}(x):=\frac{d_m(x)^a}{S_a(m)},
\qquad
a>a_{\mathrm c},
\]
则有
\[
\boxed{\;
\EE_{\pi_{m,a}}\!\left[\left(\frac{d_m(X)}{M_m}\right)^k\right]
=
1+O\bigl(e^{-m\Delta(a)+o(m)}\bigr),\;}
\]
并且
\[
\boxed{\;
\EE_{\pi_{m,a}}\!\left[\left(1-\frac{d_m(X)}{M_m}\right)^k\right]
=
O\bigl(e^{-m\Delta(a)+o(m)}\bigr).\;}
\]
\end{corollary}
\begin{proof}
由定义直接得到
\[
\EE_{\pi_{m,a}}\!\left[\left(\frac{d_m(X)}{M_m}\right)^k\right]
=
\frac{S_{a+k}(m)}{M_m^kS_a(m)}.
\]
代入定理 \ref{thm:xi-fixed-freezing-power-sum-ratio-locking} 即得第一式。

对第二式，作二项展开
\[
\left(1-\frac{d_m(X)}{M_m}\right)^k
=
\sum_{j=0}^{k}(-1)^j\binom{k}{j}\left(\frac{d_m(X)}{M_m}\right)^j.
\]
取 \(\pi_{m,a}\)-期望后，\(j=0\) 项恒等于 \(1\)，而对每个 \(1\le j\le k\) 由第一式有
\[
\EE_{\pi_{m,a}}\!\left[\left(\frac{d_m(X)}{M_m}\right)^j\right]
=
1+O\bigl(e^{-m\Delta(a)+o(m)}\bigr).
\]
于是
\[
\EE_{\pi_{m,a}}\!\left[\left(1-\frac{d_m(X)}{M_m}\right)^k\right]
=
\sum_{j=0}^{k}(-1)^j\binom{k}{j}
+O\bigl(e^{-m\Delta(a)+o(m)}\bigr).
\]
再利用
\[
\sum_{j=0}^{k}(-1)^j\binom{k}{j}=0
\]
便得结论。证毕。
\end{proof}

\begin{theorem}[零温缩放极限的原子--地基精确分裂]\label{thm:xi-zero-temp-atomic-ground-exact-splitting}
沿用附录 \ref{app:real-input-40-zeta-u} 的记号。对真实输入 \(40\) 状态核，
\[
\Delta(z,u)=(1-uz^2)\,F(z,u).
\]
再令
\[
B_\infty=
\begin{pmatrix}
0&1&1&0\\
0&0&1&0\\
0&0&3&1\\
1&0&0&3
\end{pmatrix}.
\]
则对任意固定 \(w\in\CC\)，在零温缩放
\[
z=\sqrt{w/u}
\]
下有
\[
\boxed{\;
\lim_{u\to+\infty}\Delta\!\left(\sqrt{w/u},u\right)
=
(1-w)\det(I-wB_\infty).\;}
\]
等价地，零温窗口内的极限解析对象严格分裂为一个 primitive/Big-Witt 原子因子 \((1-w)\) 与一个四状态地基 SFT 的动力学行列式 \(\det(I-wB_\infty)\)。
\end{theorem}
\begin{proof}
由附录 \ref{app:real-input-40-zeta-u}，
\[
\begin{aligned}
F(z,u)
&=
1-z-6uz^2+2uz^3+(9u^2-u)z^4\\
&\qquad +(u-3u^2)z^5-(u^3+2u^2)z^6\\
&\qquad +(4u^3-3u^2)z^7+(u^3-u^4)z^8.
\end{aligned}
\]
把
\[
z=\sqrt{w/u}
\]
代入，得到
\[
\begin{aligned}
F\!\left(\sqrt{w/u},u\right)
&=
1-u^{-1/2}w^{1/2}-6w+2u^{-1/2}w^{3/2}\\
&\qquad +(9-u^{-1})w^2+(u^{-3/2}-3u^{-1/2})w^{5/2}\\
&\qquad -(1+2u^{-1})w^3+(4u^{-1/2}-3u^{-3/2})w^{7/2}\\
&\qquad +(u^{-1}-1)w^4.
\end{aligned}
\]
令 \(u\to+\infty\)，所有带负半整数次幂的项都消失，于是
\[
\lim_{u\to+\infty}F\!\left(\sqrt{w/u},u\right)
=
1-6w+9w^2-w^3-w^4.
\]
另一方面，由定理 \ref{thm:xi-zero-temperature-fourstate-artin-mazur-zeta-primitive-asymptotic}，
\[
1-6w+9w^2-w^3-w^4=\det(I-wB_\infty).
\]
再由
\[
1-u\left(\sqrt{w/u}\right)^2=1-w,
\]
便得到
\[
\lim_{u\to+\infty}\Delta\!\left(\sqrt{w/u},u\right)
=
(1-w)\det(I-wB_\infty).
\]
证毕。
\end{proof}

\begin{corollary}[零温主极点由地基层而非原子层支配]\label{cor:xi-zero-temp-leading-pole-ground-dominance}
沿用定理 \ref{thm:xi-zero-temp-atomic-ground-exact-splitting} 的记号。记
\[
y_\infty:=\rho(B_\infty)\approx 3.596154676009,
\qquad
b_\infty:=y_\infty^{-1}\approx 0.278074802141.
\]
则极限多项式
\[
(1-w)\det(I-wB_\infty)
\]
在正实轴上的最内层零点为
\[
\boxed{\;w=b_\infty,\;}
\]
而非原子零点 \(w=1\)。因此，零温缩放窗口中的主导自由能尺度完全由四状态地基层的 Perron 根 \(y_\infty\) 决定，而不由原子因子 \((1-w)\) 决定。
\end{corollary}
\begin{proof}
由定理 \ref{thm:xi-zero-temperature-fourstate-artin-mazur-zeta-primitive-asymptotic}，
\[
\det(I-wB_\infty)=1-6w+9w^2-w^3-w^4.
\]
其零点恰为 \(B_\infty\) 各特征值的倒数。由于 \(B_\infty\) 为非负不可约矩阵，其 Perron 根
\[
y_\infty=\rho(B_\infty)>1
\]
给出唯一最小正实倒数
\[
b_\infty=y_\infty^{-1}\in(0,1).
\]
另一方面，原子因子 \((1-w)\) 的正实零点为 \(w=1\)。由于
\[
b_\infty<1,
\]
在全部正实零点中更靠近原点的是 \(w=b_\infty\) 而非 \(w=1\)。故主导尺度来自地基层行列式 \(\det(I-wB_\infty)\)，而非原子因子。证毕。
\end{proof}

\begin{theorem}[零温地基四次数域与 Lee--Yang 三次数域的线性互不相交]\label{thm:xi-zero-temp-quartic-leyang-cubic-linear-disjoint}
记
\[
K_\infty:=\QQ(y_\infty),
\qquad
y_\infty=\rho(B_\infty),
\]
其中 \(y_\infty\) 满足
\[
y^4-6y^3+9y^2-y-1=0.
\]
再令
\[
K_{\mathrm{LY}}:=\QQ(u)=\QQ(b_{\mathrm{LY}}),
\]
其中 \(u\) 与 \(b_{\mathrm{LY}}\) 分别为文中 Lee--Yang 分支系数与振幅平方常数所生成的共同三次数域。则有
\[
\boxed{\;
[K_\infty:\QQ]=4,\qquad [K_{\mathrm{LY}}:\QQ]=3,\qquad
K_\infty\cap K_{\mathrm{LY}}=\QQ.\;}
\]
从而
\[
\boxed{\;
[K_\infty K_{\mathrm{LY}}:\QQ]=12,\;}
\]
并且 \(K_\infty\) 与 \(K_{\mathrm{LY}}\) 在 \(\QQ\) 上线性互不相交。
\end{theorem}
\begin{proof}
由定理 \ref{thm:xi-zero-temperature-fourstate-perron-parry-closed-form}，多项式
\[
y^4-6y^3+9y^2-y-1
\]
在 \(\QQ\) 上不可约，故
\[
[K_\infty:\QQ]=4.
\]
另一方面，由推论 \ref{cor:xi-terminal-zm-leyang-cubic-rational-observable-primitivity}，
\[
K_{\mathrm{LY}}=\QQ(u)=\QQ(b_{\mathrm{LY}})
\]
且
\[
[K_{\mathrm{LY}}:\QQ]=3.
\]

设
\[
E:=K_\infty\cap K_{\mathrm{LY}}.
\]
由塔式公式，
\[
[E:\QQ]\mid [K_\infty:\QQ]=4,
\qquad
[E:\QQ]\mid [K_{\mathrm{LY}}:\QQ]=3.
\]
因此 \([E:\QQ]\) 同时整除 \(4\) 与 \(3\)，只能等于 \(1\)。于是
\[
K_\infty\cap K_{\mathrm{LY}}=\QQ.
\]
再由塔式公式，
\[
[K_\infty K_{\mathrm{LY}}:\QQ]
=
\frac{[K_\infty:\QQ][K_{\mathrm{LY}}:\QQ]}{[K_\infty\cap K_{\mathrm{LY}}:\QQ]}
=
4\cdot 3
=
12.
\]
由于特征 \(0\) 下有限扩张皆可分，乘积次数达到乘积上界即等价于线性互不相交，故结论成立。证毕。
\end{proof}






\paragraph{Smith 信息亏损谱、有限历史 prime-register 压缩与 Euler 原子的严格手术律}
下列结果把定理 \ref{thm:killo-smith-loss-spectrum}、\ref{thm:killo-layered-prime-slices-finite-depth}、\ref{thm:killo-natural-extension-branch-register} 与命题 \ref{prop:atomic-2-orbit-split-logM}、\ref{prop:primitive-orbit-surgery} 接入同一条“线性投影--分支寄存器--prime-slice 历史--primitive 手术”链条。

\begin{theorem}[Smith 信息亏损谱与一步精确反演的最小分支寄存器]\label{thm:xi-smith-loss-minimal-branch-register}
\leanverified{paper\_xi\_smith\_loss\_minimal\_branch\_register}
设
\[
A\in M_{m\times n}(\ZZ),
\qquad
\operatorname{rank}_{\QQ}(A)=m,
\]
其 Smith 不变量记为
\[
d_1\mid d_2\mid\cdots\mid d_m.
\]
固定素数 \(p\) 与整数 \(k\ge 1\)，并将
\[
T_{A,p^k}:(\ZZ/p^k\ZZ)^n\longrightarrow (\ZZ/p^k\ZZ)^m,
\qquad
x\longmapsto Ax\bmod p^k
\]
视为到其像上的满射。记
\[
K_p(k):=\#\ker(A\bmod p^k),
\qquad
e_i(p):=v_p(d_i),
\qquad
f_p(k):=\sum_{i=1}^{m}\min\{k,e_i(p)\}.
\]
若存在映射
\[
R:(\ZZ/p^k\ZZ)^n\longrightarrow \mathcal R
\]
使得
\[
x\longmapsto \bigl(T_{A,p^k}(x),R(x)\bigr)
\]
为单射，则必有
\[
|\mathcal R|\ge K_p(k)=p^{\,k(n-m)+f_p(k)}.
\]
该下界可达。特别地，
\[
\log_p|\mathcal R|\ge k(n-m)+f_p(k),
\qquad
\log_2|\mathcal R|\ge \bigl(k(n-m)+f_p(k)\bigr)\log_2 p.
\]
\end{theorem}
\begin{proof}
由定理 \ref{thm:killo-smith-loss-spectrum}，
\[
K_p(k)=p^{\,k(n-m)+f_p(k)}.
\]
由于 \(T_{A,p^k}\) 是有限 Abel 群之间的同态，其任意非空纤维都是 \(\ker(A\bmod p^k)\) 的仿射陪集，故每条纤维的基数都等于 \(K_p(k)\)。

若 \(x\mapsto (T_{A,p^k}(x),R(x))\) 单射，则对每个 \(y\in\operatorname{Im}(T_{A,p^k})\)，限制映射
\[
R:T_{A,p^k}^{-1}(y)\to \mathcal R
\]
必须单射。于是 \(|\mathcal R|\) 至少不小于任一纤维的大小，即
\(
|\mathcal R|\ge K_p(k)
\)。

反之，取一个固定集合 \(\mathcal R_0\) 满足 \(|\mathcal R_0|=K_p(k)\)。由于 \((\ZZ/p^k\ZZ)^n\) 有限，对每条纤维 \(T_{A,p^k}^{-1}(y)\) 任选一个双射
\[
\rho_y:T_{A,p^k}^{-1}(y)\xrightarrow{\sim}\mathcal R_0.
\]
定义 \(R(x):=\rho_{T_{A,p^k}(x)}(x)\)，则同一宏态纤维内由 \(R\) 区分，不同宏态由 \(T_{A,p^k}\) 区分，从而
\[
x\longmapsto \bigl(T_{A,p^k}(x),R(x)\bigr)
\]
为单射，故下界可达。
\end{proof}

\begin{corollary}[Smith 指数的层寿命与逐层寄存器增量]\label{cor:xi-smith-index-layer-lifetime-register}
\leanverified{paper\_xi\_smith\_index\_layer\_lifetime\_register}
沿用定理 \ref{thm:xi-smith-loss-minimal-branch-register} 的记号，定义
\[
B_p(k):=\log_p K_p(k)=k(n-m)+f_p(k),
\qquad
f_p(0):=0,
\]
\[
\Delta_p(k):=f_p(k)-f_p(k-1)
\qquad(k\ge 1).
\]
则有
\[
\Delta_p(k)=\#\{i:\ e_i(p)\ge k\},
\]
从而
\[
B_p(k)-B_p(k-1)=(n-m)+\Delta_p(k).
\]
因此自由核秩 \(n-m\) 在每一层都贡献恒定的 \(p\)-元寄存器长度，而每个 Smith 指数 \(e_i(p)\) 则精确贡献一个寿命为 \(e_i(p)\) 的附加 digit 轴。若进一步 \(m=n\)，则
\[
\sum_{k\ge 1}\Delta_p(k)=\sum_{i=1}^{n}e_i(p)=v_p(\det A).
\]
\end{corollary}
\begin{proof}
由定义
\[
f_p(k)=\sum_{i=1}^{m}\min\{k,e_i(p)\}.
\]
故
\[
\Delta_p(k)
=\sum_{i=1}^{m}\bigl(\min\{k,e_i(p)\}-\min\{k-1,e_i(p)\}\bigr)
=\#\{i:\ e_i(p)\ge k\},
\]
因为对任意整数 \(e\ge 0\) 都有
\[
\min\{k,e\}-\min\{k-1,e\}=\mathbf 1_{\{e\ge k\}}.
\]
将该式代入
\(
B_p(k)=k(n-m)+f_p(k)
\)
即可得到
\[
B_p(k)-B_p(k-1)=(n-m)+\Delta_p(k).
\]
若 \(m=n\)，则 \(\Delta_p(k)\) 只有有限支撑，故
\[
\sum_{k\ge 1}\Delta_p(k)=f_p(\infty)=\sum_{i=1}^{n}e_i(p)=v_p(\det A).
\]
\end{proof}

\begin{theorem}[有限历史的 prime-register 压缩与自然扩张的逆极限实现]\label{thm:xi-prime-register-history-inverse-limit}
设 \(X\) 为可数集合，\(T:X\twoheadrightarrow X\) 为有限纤维满射，并固定定理 \ref{thm:killo-natural-extension-branch-register} 中的分支指标
\[
\beta(x):=\iota_{T(x)}(x).
\]
对每个 \(k\ge 1\)，定义 \(k\)-步历史空间
\[
X^{\leftarrow,k}
:=
\{(x_0,x_{-1},\dots,x_{-k})\in X^{k+1}:\ T(x_{-j})=x_{-j+1}\ \ (1\le j\le k)\},
\]
以及截断分支码
\[
\mathcal C_k(x_0,\dots,x_{-k})
:=
\bigl(x_0,\bigl(\beta(x_{-1}),\dots,\beta(x_{-k})\bigr)\bigr)\in X\times \NN^k.
\]
则对每个 \(k\ge 1\) 都存在单射
\[
\Phi_k:X^{\leftarrow,k}\hookrightarrow X\times \ZZ_{>0},
\qquad
\Phi_k(x_0,\dots,x_{-k})=\bigl(x_0,G_k(x_0,\dots,x_{-k})\bigr),
\]
满足：
\begin{enumerate}
  \item \(G_k\) 具有严格分层分解
  \[
  G_k=\prod_{\ell=1}^{k}q_\ell,
  \qquad
  q_\ell\in\mathcal P_{\ell-1},
  \]
  其中 \((\mathcal P_j)_{j\ge 0}\) 为定理 \ref{thm:killo-layered-prime-slices-finite-depth} 的不交素数片；
  \item 由 \((x_0,G_k)\) 可确定性恢复唯一历史 \((x_0,\dots,x_{-k})\)，等价地可恢复 \(\mathcal C_k(x_0,\dots,x_{-k})\)；
  \item 若记 \(\widehat X_k:=\Phi_k(X^{\leftarrow,k})\)，并对 \((x_0,\prod_{\ell=1}^{k+1}q_\ell)\in\widehat X_{k+1}\) 定义
  \[
  \tau_{k+1,k}\!\left(x_0,\prod_{\ell=1}^{k+1}q_\ell\right)
  :=
  \left(x_0,\prod_{\ell=1}^{k}q_\ell\right),
  \]
  则 \((\widehat X_k,\tau_{k+1,k})\) 构成投影系统，其逆极限与定理 \ref{thm:killo-natural-extension-branch-register} 中的 \(\widehat X\) 规范同构；因此，经该同构输运的动力学与自然扩张 \((X^{\leftarrow},\sigma)\) 共轭。
\end{enumerate}
\leanverified{paper\_xi\_prime\_register\_history\_inverse\_limit}
\end{theorem}
\begin{proof}
由定理 \ref{thm:killo-layered-prime-slices-finite-depth}，素数集合分解为两两不交的可数片
\[
\mathcal P=\bigsqcup_{j\ge 0}\mathcal P_j.
\]
对每个 \(\ell\ge 1\)，集合 \(X^{\leftarrow,\ell-1}\times\NN\) 可数，故可取单射
\[
\alpha_\ell:X^{\leftarrow,\ell-1}\times\NN\hookrightarrow \mathcal P_{\ell-1}.
\]
对任意 \((x_0,\dots,x_{-k})\in X^{\leftarrow,k}\)，定义
\[
q_\ell:=\alpha_\ell\!\bigl((x_0,\dots,x_{-\ell+1}),\beta(x_{-\ell})\bigr)
\qquad(1\le \ell\le k),
\]
以及
\[
G_k(x_0,\dots,x_{-k}):=\prod_{\ell=1}^{k}q_\ell.
\]
于是 \(q_\ell\in\mathcal P_{\ell-1}\)，从而得到第一条。

现由 \((x_0,G_k)\) 恢复历史。由于素数片两两不交，唯一分解定理可从 \(G_k\) 中逐片抽取出每个 \(q_\ell\)。已知 \(x_0\) 后，由 \(q_1\) 与 \(\alpha_1\) 的像上反演可恢复 \(\beta(x_{-1})\)；再由 \(\beta(x_{-1})=\iota_{x_0}(x_{-1})\) 与 \(\iota_{x_0}\) 的像上反演，可唯一恢复 \(x_{-1}\)。随后已知 \((x_0,x_{-1})\) 与 \(q_2\)，同理恢复 \(\beta(x_{-2})\) 及 \(x_{-2}\)。如此归纳，依次恢复全部 \(x_{-\ell}\)（\(1\le \ell\le k\)），从而得到唯一历史与唯一的 \(\mathcal C_k\)。第二条成立。

再看相容性。对任意 \((x_0,\dots,x_{-(k+1)})\in X^{\leftarrow,k+1}\)，有
\[
\Phi_{k+1}(x_0,\dots,x_{-(k+1)})
=
\left(x_0,\prod_{\ell=1}^{k+1}q_\ell\right),
\]
其中最后一层因子 \(q_{k+1}\in\mathcal P_k\) 只编码新增的最深一级前史 \(x_{-(k+1)}\)。因此
\[
\tau_{k+1,k}\circ \Phi_{k+1}
=
\Phi_k\circ \pi_{k+1,k},
\]
其中
\[
\pi_{k+1,k}(x_0,\dots,x_{-(k+1)})=(x_0,\dots,x_{-k}).
\]
故 \((\widehat X_k,\tau_{k+1,k})\) 确为投影系统。

最后说明逆极限识别。一个相容族 \(\bigl((x_0,G_k)\bigr)_{k\ge 1}\in\varprojlim \widehat X_k\) 通过逐层抽取素数片因子，唯一给出一列整数
\[
a_\ell:=\beta(x_{-\ell})\in\NN
\qquad(\ell\ge 1),
\]
从而唯一确定
\[
\bigl(x_0,(a_1,a_2,\dots)\bigr)\in \widehat X.
\]
反之，任取 \(\widehat X\) 中元素 \(\bigl(x_0,(\beta(x_{-1}),\beta(x_{-2}),\dots)\bigr)\)，按照上述构造截取前 \(k\) 项即可得到 \(\widehat X_k\) 中的兼容族。于是
\[
\varprojlim \widehat X_k\cong \widehat X.
\]
再由定理 \ref{thm:killo-natural-extension-branch-register}，
\((\widehat X,\widehat T)\) 与 \((X^{\leftarrow},\sigma)\) 共轭，故把 \(\widehat T\) 通过上述同构输运到 \(\varprojlim\widehat X_k\) 上，即得到与自然扩张共轭的逆极限动力学。
\end{proof}

\begin{theorem}[有限多 Euler 原子的 primitive 手术可加律]\label{thm:xi-finite-euler-primitive-surgery-additivity}
\leanverified{paper\_xi\_finite\_euler\_primitive\_surgery\_additivity}
设某个加权 \(\zeta\)-函数可写成
\[
\zeta(z,u)
=
\left(\prod_{\nu=1}^{s}(1-u^{E_\nu}z^{\ell_\nu})^{-m_\nu}\right)
\zeta_{\mathrm{core}}(z,u),
\]
其中
\[
m_\nu\in\ZZ_{>0},
\qquad
E_\nu\in\ZZ_{\ge 0},
\qquad
\ell_\nu\in\ZZ_{\ge 1},
\]
并假设对每个 \(\nu\) 都有
\[
1-u^{E_\nu}z_\star(u)^{\ell_\nu}\neq 0,
\]
其中 \(z_\star(u)=\lambda(u)^{-1}\) 为 \(\zeta_{\mathrm{core}}\) 的主极点。令 \(P(z,u)\)、\(P_{\mathrm{core}}(z,u)\) 分别为 \(\zeta(z,u)\)、\(\zeta_{\mathrm{core}}(z,u)\) 的 primitive 生成函数，\(\log\mathfrak M(u)\)、\(\log\mathfrak M_{\mathrm{core}}(u)\) 为对应的 Abel 有限部常数。则有严格加法公式
\[
P(z,u)=P_{\mathrm{core}}(z,u)+\sum_{\nu=1}^{s}m_\nu u^{E_\nu}z^{\ell_\nu},
\]
\[
\log\mathfrak M(u)
=
\log\mathfrak M_{\mathrm{core}}(u)
+\sum_{\nu=1}^{s}m_\nu u^{E_\nu}z_\star(u)^{\ell_\nu}.
\]
特别地，有限多个可剥离 Euler 原子之间不存在任何交叉修正项。
\end{theorem}
\begin{proof}
把分解式取对数，得到
\[
\log\zeta(z,u)
=
\sum_{\nu=1}^{s}-m_\nu\log(1-u^{E_\nu}z^{\ell_\nu})
+\log\zeta_{\mathrm{core}}(z,u).
\]
对两边施加 Witt/M\"obius 变换
\[
P(z,u)=\sum_{r\ge 1}\frac{\mu(r)}{r}\log\zeta(z^r,u^r),
\]
便有
\[
P(z,u)
=
P_{\mathrm{core}}(z,u)
+\sum_{\nu=1}^{s}m_\nu
\left(
-\sum_{r\ge 1}\frac{\mu(r)}{r}\log(1-u^{rE_\nu}z^{r\ell_\nu})
\right).
\]
对形式变量 \(X\) 使用标准恒等式
\[
-\sum_{r\ge 1}\frac{\mu(r)}{r}\log(1-X^r)=X,
\]
并取 \(X=u^{E_\nu}z^{\ell_\nu}\)，即得
\[
P(z,u)=P_{\mathrm{core}}(z,u)+\sum_{\nu=1}^{s}m_\nu u^{E_\nu}z^{\ell_\nu}.
\]

由于各原子因子在 \(z=z_\star(u)\) 处均解析且不为零，主极点仍由 \(\zeta_{\mathrm{core}}\) 控制。于是可逐项将上式代入 Abel 有限部常数的定义；等价地，也可对每个原子重复应用命题 \ref{prop:primitive-orbit-surgery}。由线性性立即得到
\[
\log\mathfrak M(u)
=
\log\mathfrak M_{\mathrm{core}}(u)
+\sum_{\nu=1}^{s}m_\nu u^{E_\nu}z_\star(u)^{\ell_\nu}.
\]
当 \(s=1\)、\((m_1,E_1,\ell_1)=(1,1,2)\) 时，即回到真实输入 \(40\) 状态核中命题 \ref{prop:atomic-2-orbit-split-logM} 的长度 \(2\) 原子剥离。
\end{proof}

\paragraph{有限模层析、ghost 闭式与 Abel 位移}

\begin{theorem}[原子 primitive 轨道手术的有限同余层析]\label{thm:xi-primitive-surgery-finite-congruence-tomography}
设某个加权有限矩阵族/有限型子移位满足
\[
\zeta(z,u)=\frac{1}{(1-u^Ez^\ell)^m}\,\zeta_{\mathrm{core}}(z,u),
\qquad
m\in\ZZ_{\ge 1},\ \ell\in\ZZ_{\ge 1},\ E\in\ZZ_{\ge 0},
\]
并令 primitive 生成函数满足
\[
P(z,u)=P_{\mathrm{core}}(z,u)+m\,u^Ez^\ell.
\]
对任意模数 \(q\ge 2\) 与剩余类 \(a\in\ZZ/q\ZZ\)，定义 root-of-unity 滤波的 primitive 生成函数
\[
P_{a,q}(z,u):=\frac1q\sum_{j=0}^{q-1}\omega_q^{-aj}P(\omega_q^j z,u),
\qquad
\omega_q:=e^{2\pi i/q},
\]
并对 \(P_{\mathrm{core}}\) 同样定义 \(P^{\mathrm{core}}_{a,q}(z,u)\)。则成立精确公式
\[
\boxed{\;
P_{a,q}(z,u)-P^{\mathrm{core}}_{a,q}(z,u)
=
m\,u^E z^\ell\,
\mathbf 1_{\{a\equiv \ell\;(\mathrm{mod}\,q)\}}.\;}
\]
因此，若取两两互素模数 \(q_1,\dots,q_r\) 满足
\[
Q:=q_1\cdots q_r>\ell,
\]
并已知对每个 \(q_i\) 都存在唯一剩余类 \(a_i\in\ZZ/q_i\ZZ\) 使
\[
P_{a_i,q_i}(z,u)-P^{\mathrm{core}}_{a_i,q_i}(z,u)\neq 0,
\]
则由中国剩余定理可唯一恢复 \(\ell\)。进一步，若在两个不同正参数 \(u_1\neq u_2\) 处测得相应非零系数
\[
c_i^{(j)}
:=
\left[z^\ell\right]
\bigl(
P_{a_i,q_i}(z,u_j)-P^{\mathrm{core}}_{a_i,q_i}(z,u_j)
\bigr)
=
m\,u_j^E
\qquad (j=1,2),
\]
则 \(E\) 与 \(m\) 亦由此唯一恢复：
\[
E=\frac{\log(c_i^{(2)}/c_i^{(1)})}{\log(u_2/u_1)},
\qquad
m=\frac{c_i^{(1)}}{u_1^E}.
\]
\end{theorem}
\begin{proof}
把
\[
P(z,u)-P_{\mathrm{core}}(z,u)=m\,u^Ez^\ell
\]
代入 root-of-unity 滤波的定义，得到
\[
\begin{aligned}
P_{a,q}(z,u)-P^{\mathrm{core}}_{a,q}(z,u)
&=
\frac1q\sum_{j=0}^{q-1}\omega_q^{-aj}
\bigl(P(\omega_q^j z,u)-P_{\mathrm{core}}(\omega_q^j z,u)\bigr)\\
&=
\frac1q\sum_{j=0}^{q-1}\omega_q^{-aj}\,m\,u^E(\omega_q^j z)^\ell\\
&=
m\,u^Ez^\ell\cdot
\frac1q\sum_{j=0}^{q-1}\omega_q^{j(\ell-a)}.
\end{aligned}
\]
有限 Fourier 求和给出
\[
\frac1q\sum_{j=0}^{q-1}\omega_q^{j(\ell-a)}
=
\mathbf 1_{\{a\equiv \ell\;(\mathrm{mod}\,q)\}},
\]
从而得到所示精确公式。

于是每个模数 \(q_i\) 都唯一暴露 \(\ell\bmod q_i\)。因 \(q_1,\dots,q_r\) 两两互素，中国剩余定理给出唯一的
\[
\ell \pmod Q.
\]
又因 \(0\le \ell<Q\)，故 \(\ell\) 本身唯一确定。已知 \(\ell\) 后，任何非零滤波系数都等于 \(m\,u^E\)。在两个不同正参数 \(u_1\neq u_2\) 处比较，便唯一恢复
\[
E=\frac{\log(c_i^{(2)}/c_i^{(1)})}{\log(u_2/u_1)},
\qquad
m=\frac{c_i^{(1)}}{u_1^E}.
\]
证毕。
\end{proof}

\begin{theorem}[原子手术的 ghost 级数闭式与 Abel 有限部的精确位移]\label{thm:xi-atomic-surgery-ghost-abel-shift-general}
沿用定理 \ref{thm:xi-primitive-surgery-finite-congruence-tomography} 的设定。定义 periodic/ghost 系数 \(a_n(u)\)、\(a_n^{\mathrm{core}}(u)\) 为
\[
\log \zeta(z,u)=\sum_{n\ge 1}\frac{a_n(u)}{n}z^n,
\qquad
\log \zeta_{\mathrm{core}}(z,u)=\sum_{n\ge 1}\frac{a_n^{\mathrm{core}}(u)}{n}z^n.
\]
则有精确 ghost 修正公式
\[
\boxed{\;
a_n(u)-a_n^{\mathrm{core}}(u)
=
m\ell\,u^{En/\ell}\,\mathbf 1_{\{\ell\mid n\}}.\;}
\]
进一步，若记主增长率为 \(\lambda(u)=z_\star(u)^{-1}\)，并定义 Abel 有限部常数
\[
\log\mathfrak M(u):=\sum_{n\ge 1}\left(\frac{p_n(u)}{\lambda(u)^n}-\frac1n\right),
\]
其中 \(P(z,u)=\sum_{n\ge 1}p_n(u)z^n\)，而 \(\log\mathfrak M_{\mathrm{core}}(u)\) 对 \(P_{\mathrm{core}}\) 同样定义，则有
\[
\boxed{\;
\log\mathfrak M(u)-\log\mathfrak M_{\mathrm{core}}(u)
=
m\,u^Ez_\star(u)^\ell.\;}
\]
\end{theorem}
\begin{proof}
由
\[
\zeta(z,u)=\frac{1}{(1-u^Ez^\ell)^m}\,\zeta_{\mathrm{core}}(z,u)
\]
直接取对数可得
\[
\log\zeta(z,u)-\log\zeta_{\mathrm{core}}(z,u)
=
-m\log(1-u^Ez^\ell)
=
m\sum_{r\ge 1}\frac{u^{Er}z^{\ell r}}{r}.
\]
另一方面，
\[
\log\zeta(z,u)-\log\zeta_{\mathrm{core}}(z,u)
=
\sum_{n\ge 1}\frac{a_n(u)-a_n^{\mathrm{core}}(u)}{n}z^n.
\]
比较两边系数，若 \(n=\ell r\) 则
\[
\frac{a_n(u)-a_n^{\mathrm{core}}(u)}{n}
=
\frac{m\,u^{Er}}{r},
\]
亦即
\[
a_n(u)-a_n^{\mathrm{core}}(u)
=
m\ell\,u^{En/\ell}.
\]
若 \(\ell\nmid n\)，则对应系数为 \(0\)。这就得到 ghost 修正公式。

至于 Abel 有限部的位移，定理 \ref{thm:xi-finite-euler-primitive-surgery-additivity} 在单原子情形 \(s=1\) 已经给出
\[
\log\mathfrak M(u)-\log\mathfrak M_{\mathrm{core}}(u)
=
m\,u^Ez_\star(u)^\ell.
\]
证毕。
\end{proof}

\paragraph{原子 Witt 手术的压力不变量与 Abel 有限部位移}
现有 primitive 剥离、ghost 闭式与有限部常数公式还可进一步压缩为一条统一的谱手术定理：可剥离的单原子 Euler 因子只改写有限长度的 primitive 坐标、其倍长 ghost 坐标与 Abel 有限部常数，而不改动主极点、压力函数及其 Legendre 几何。

\begin{theorem}[可剥离原子 Witt 因子的谱手术不变量]\label{thm:xi-atomic-witt-surgery-pressure-ldp-abel-invariance}
\leanverified{paper\_xi\_atomic\_witt\_surgery\_pressure\_ldp\_abel\_invariance}
设
\[
\zeta(z,u)=\frac{1}{(1-u^Ez^\ell)^m}\,\zeta_{\mathrm{core}}(z,u),
\qquad
m\in\ZZ_{\ge 1},\ \ell\in\ZZ_{\ge 1},\ E\in\ZZ_{\ge 0},
\]
并记
\[
\mathcal P(z,u):=\sum_{n\ge 1}p_n(u)z^n,
\qquad
\mathcal P_{\mathrm{core}}(z,u):=\sum_{n\ge 1}p_n^{\mathrm{core}}(u)z^n,
\]
\[
\log\zeta(z,u)=\sum_{n\ge 1}\frac{a_n(u)}{n}z^n,
\qquad
\log\zeta_{\mathrm{core}}(z,u)=\sum_{n\ge 1}\frac{a_n^{\mathrm{core}}(u)}{n}z^n.
\]
设 \(z_\star(u)\) 为 \(\zeta_{\mathrm{core}}\) 的主极点，并假设
\[
1-u^Ez_\star(u)^\ell\neq 0,
\]
即该原子因子在主极点处解析且不为零。并记 \(\lambda(u)\) 为 \(\zeta(z,u)\) 的主增长率，且定义
\[
\lambda_{\mathrm{core}}(u):=z_\star(u)^{-1},
\qquad
P(\theta):=\log\lambda(e^\theta),
\qquad
P_{\mathrm{core}}(\theta):=\log\lambda_{\mathrm{core}}(e^\theta).
\]
则有：
\begin{enumerate}
  \item \textbf{压力与全部导数不变：}
  \[
  \boxed{\;
  \lambda(u)=\lambda_{\mathrm{core}}(u),\qquad
  P(\theta)=P_{\mathrm{core}}(\theta),\qquad
  P^{(k)}(\theta)=P_{\mathrm{core}}^{(k)}(\theta)\ \ (k\ge 1).\;}
  \]
  \item \textbf{Legendre 大偏差率函数不变：}
  若
  \[
  I(\alpha):=\sup_{\theta\in\RR}\bigl(\theta\alpha-(P(\theta)-P(0))\bigr),
  \]
  \[
  I_{\mathrm{core}}(\alpha):=\sup_{\theta\in\RR}\bigl(\theta\alpha-(P_{\mathrm{core}}(\theta)-P_{\mathrm{core}}(0))\bigr),
  \]
  则
  \[
  \boxed{\;I(\alpha)=I_{\mathrm{core}}(\alpha).\;}
  \]
  \item \textbf{primitive 轨道只在长度 \(\ell\) 处发生原子增量：}
  \[
  \boxed{\;
  p_n(u)=p_n^{\mathrm{core}}(u)+m\,u^E\,\mathbf 1_{\{n=\ell\}}.\;}
  \]
  \item \textbf{periodic/ghost 计数只在 \(\ell\)-倍长度出现刚性增量：}
  \[
  \boxed{\;
  a_n(u)=a_n^{\mathrm{core}}(u)+m\ell\,u^{En/\ell}\,\mathbf 1_{\{\ell\mid n\}}.\;}
  \]
  \item \textbf{Abel 有限部常数发生精确平移：}
  若
  \[
  \log\mathfrak M(u):=\sum_{n\ge 1}\left(\frac{p_n(u)}{\lambda(u)^n}-\frac1n\right),
  \]
  且 \(\log\mathfrak M_{\mathrm{core}}(u)\) 对 \(\mathcal P_{\mathrm{core}}\) 同样定义，则
  \[
  \boxed{\;
  \log\mathfrak M(u)=\log\mathfrak M_{\mathrm{core}}(u)+m\,u^Ez_\star(u)^\ell.\;}
  \]
\end{enumerate}
\end{theorem}
\begin{proof}
由定理 \ref{thm:xi-finite-euler-primitive-surgery-additivity}，
\[
\mathcal P(z,u)=\mathcal P_{\mathrm{core}}(z,u)+m\,u^Ez^\ell,
\]
逐项比较系数即得第 \(3\) 条。

由定理 \ref{thm:xi-atomic-surgery-ghost-abel-shift-general}，直接得到第 \(4\) 条与第 \(5\) 条。

由于 \(1-u^Ez_\star(u)^\ell\neq 0\)，原子因子在 \(z=z_\star(u)\) 处既不产生极点也不消去极点，故总体 \(\zeta\) 与
\(\zeta_{\mathrm{core}}\) 的主极点一致，从而
\[
\lambda(u)=\lambda_{\mathrm{core}}(u).
\]
取对数便得
\[
P(\theta)=P_{\mathrm{core}}(\theta),
\]
于是全部导数一致，第 \(1\) 条成立。

第 \(2\) 条随即由 Legendre 变换仅依赖压力函数这一事实立即推出。证毕。
\end{proof}
\begin{theorem}[真实输入 \texorpdfstring{$40$}{40} 状态核的偶长度刚性项与零温 Abel 位移]\label{thm:xi-real-input40-evenlength-abel-shift-rigidity}
沿用附录 \ref{app:real-input-40-zeta-u} 的记号。对真实输入 \(40\) 状态核，
\[
\Delta(z,u)=(1-uz^2)\,F(z,u),
\qquad
\zeta(z,u)=\Delta(z,u)^{-1},
\qquad
\zeta_{\mathrm{core}}(z,u):=F(z,u)^{-1}.
\]
记
\[
\mathcal P(z,u)=\sum_{n\ge 1}p_n(u)z^n,
\qquad
\mathcal P_{\mathrm{core}}(z,u)=\sum_{n\ge 1}p_n^{\mathrm{core}}(u)z^n,
\]
\[
\log\zeta(z,u)=\sum_{n\ge 1}\frac{a_n(u)}{n}z^n,
\qquad
\log\zeta_{\mathrm{core}}(z,u)=\sum_{n\ge 1}\frac{a_n^{\mathrm{core}}(u)}{n}z^n,
\]
并令 \(z_\star(u)=\lambda(u)^{-1}\) 为由 core 因子确定的主极点。则有：
\begin{enumerate}
  \item primitive 轨道只在长度 \(2\) 处出现唯一原子修正：
  \[
  \boxed{\;
  p_n(u)=p_n^{\mathrm{core}}(u)+u\,\mathbf 1_{\{n=2\}}.\;}
  \]
  \item periodic/ghost 计数满足严格的偶长度刚性：
  \[
  \boxed{\;
  a_n(u)=a_n^{\mathrm{core}}(u)+2u^{n/2}\,\mathbf 1_{\{2\mid n\}}.\;}
  \]
  特别地，
  \[
  a_{2m+1}(u)=a_{2m+1}^{\mathrm{core}}(u),
  \qquad
  a_{2m}(u)-a_{2m}^{\mathrm{core}}(u)=2u^m.
  \]
  \item Abel 有限部常数的位移满足
  \[
  \boxed{\;
  \log\mathfrak M(u)-\log\mathfrak M_{\mathrm{core}}(u)=u\,z_\star(u)^2.\;}
  \]
  若记
  \[
  b:=\lim_{u\to\infty}u\,z_\star(u)^2,
  \]
  则该极限存在，且 \(b\) 是方程
  \[
  1-6b+9b^2-b^3-b^4=0
  \]
  的最大正根，并满足
  \[
  \boxed{\;
  \lim_{u\to\infty}
  \bigl(\log\mathfrak M(u)-\log\mathfrak M_{\mathrm{core}}(u)\bigr)=b
  \approx 0.27807480214113.\;}
  \]
  \item 在无权点 \(u=1\) 处，
  \[
  \boxed{\;
  \log\mathfrak M(1)-\log\mathfrak M_{\mathrm{core}}(1)=z_\star(1)^2
  =\varphi^{-4}
  =\frac{7-3\sqrt5}{2}.\;}
  \]
\end{enumerate}
因此，长度 \(2\) 的原子因子对压力函数及其大偏差几何完全透明，而其全部可见效应恰被压缩为偶长度 periodic 刚性项与 Abel 有限部的精确位移。
\end{theorem}
\begin{proof}
将定理 \ref{thm:xi-atomic-witt-surgery-pressure-ldp-abel-invariance} 应用于
\[
(m,\ell,E)=(1,2,1)
\]
即可得到第 \(1\) 条至第 \(3\) 条的有限 \(u\) 公式。

零温极限
\[
u\,z_\star(u)^2\longrightarrow b
\]
的存在性、四次方程表征与数值近似由命题 \ref{prop:real-input-40-logM-infty} 给出，故
第 \(3\) 条中的极限位移随之成立。

在 \(u=1\) 时，由附录 \ref{app:real-input-40-zeta-u} 的无权主极点
\[
z_\star(1)=\varphi^{-2}
\]
立得
\[
\log\mathfrak M(1)-\log\mathfrak M_{\mathrm{core}}(1)=z_\star(1)^2=\varphi^{-4}
=\frac{7-3\sqrt5}{2}.
\]
证毕。
\end{proof}
\begin{corollary}[无权两步 primitive 系数的精确识别]\label{cor:xi-real-input40-primitive-two-step-exact-count}
\leanverified{paper\_xi\_real\_input40\_primitive\_two\_step\_exact\_count}
沿用定理 \ref{thm:xi-real-input40-evenlength-abel-shift-rigidity} 的记号，并取 \(u=1\)。
若再结合命题 \ref{prop:atomic-2-orbit-split-logM} 所嵌 branch-level 审计给出的两步 primitive 回路总数
\[
(|\tau|,E)=(2,1)\quad\text{恰有 }7\text{ 条},
\]
则
\[
\boxed{\;
p_2=7,
\qquad
p_2^{\mathrm{core}}=6.\;}
\]
\end{corollary}
\begin{proof}
由定理 \ref{thm:xi-real-input40-evenlength-abel-shift-rigidity} 在 \(u=1\) 时的 primitive 剥离，
\[
p_2=1+p_2^{\mathrm{core}}.
\]
另一方面，branch-level 审计已经给出全部满足
\[
(|\tau|,E)=(2,1)
\]
的两步 primitive 回路总数恰为 \(7\)，故
\[
p_2=7.
\]
代回前式立即得到
\[
p_2^{\mathrm{core}}=6.
\]
证毕。
\end{proof}

\paragraph{自然扩张的周期守恒、prime-slice 精确极小性与 Smith 核增长的 Euler 定律}
沿用定理 \ref{thm:killo-layered-prime-slices-finite-depth}、推论
\ref{cor:killo-layered-necessity-minimality}、定理
\ref{thm:killo-natural-extension-branch-register}、定理
\ref{thm:killo-smith-loss-spectrum} 与推论
\ref{cor:killo-kernel-size-general-modulus} 的记号。下列结果把
“自然扩张--有限历史 prime-slice--Smith 信息亏损谱--全模数核增长”
四段对象压缩为同一条可直接调用的结论链。

\begin{theorem}[自然扩张寄存器实现保持全部周期点计数与 Artin--Mazur \texorpdfstring{$\zeta$}{zeta} 函数]\label{thm:xi-natural-extension-periodic-orbit-zeta-preservation}
设 \(T:X\twoheadrightarrow X\) 为有限纤维满射，\(X^{\leftarrow}\) 为定理
\ref{thm:killo-natural-extension-branch-register} 的自然扩张，\(\sigma\) 为其左移同构，
\[
\mathcal C:X^{\leftarrow}\xrightarrow{\sim}\widehat X,
\qquad
\widehat T:=\mathcal C\circ \sigma\circ \mathcal C^{-1}.
\]
则对每个整数 \(n\ge 1\)，映射
\[
\pi_0:\operatorname{Fix}(\sigma^n)\longrightarrow \operatorname{Fix}(T^n),
\qquad
(x_0,x_{-1},x_{-2},\dots)\longmapsto x_0
\]
为自然双射。因而也有自然双射
\[
\operatorname{Fix}(\widehat T^n)\cong \operatorname{Fix}(T^n).
\]
特别地，若这些不动点集合逐个有限，则 Artin--Mazur \(\zeta\)-函数满足严格恒等式
\[
\zeta_{\widehat T}(z)
:=
\exp\!\left(\sum_{n\ge 1}\frac{\#\operatorname{Fix}(\widehat T^n)}{n}z^n\right)
=
\exp\!\left(\sum_{n\ge 1}\frac{\#\operatorname{Fix}(T^n)}{n}z^n\right)
=:
\zeta_T(z).
\]
进一步，若两侧拓扑熵均由周期点增长率定义，则
\[
h_{\mathrm{top}}(\widehat T)=h_{\mathrm{top}}(T).
\]
\leanverified{paper\_xi\_natural\_extension\_periodic\_orbit\_zeta\_preservation}
\end{theorem}
\begin{proof}
对任意
\[
\mathbf x=(x_0,x_{-1},x_{-2},\dots)\in X^{\leftarrow},
\]
由自然扩张的定义可知
\[
\sigma^n(\mathbf x)
=
\bigl(T^n(x_0),T^{n-1}(x_0),\dots,T(x_0),x_0,x_{-1},x_{-2},\dots\bigr).
\]
故若 \(\sigma^n(\mathbf x)=\mathbf x\)，则必有
\[
T^n(x_0)=x_0,
\qquad
x_{-r}=T^{n-r}(x_0)\quad(1\le r\le n),
\]
并且由继续比较坐标得到
\[
x_{-(qn+r)}=T^{n-r}(x_0)
\qquad(q\ge 0,\ 1\le r\le n).
\]
因此 \(\mathbf x\) 由 \(x_0\in\operatorname{Fix}(T^n)\) 唯一决定。

反之，任取 \(x_0\in\operatorname{Fix}(T^n)\)，按上式定义
\[
x_{-(qn+r)}:=T^{n-r}(x_0)
\qquad(q\ge 0,\ 1\le r\le n),
\]
则由 \(T^n(x_0)=x_0\) 可直接验证
\[
T(x_{-j})=x_{-j+1}\qquad(j\ge 1),
\]
故所得 \(\mathbf x\in X^{\leftarrow}\)；并且再由构造立即有
\[
\sigma^n(\mathbf x)=\mathbf x.
\]
从而 \(\pi_0\) 给出 \(\operatorname{Fix}(\sigma^n)\) 与 \(\operatorname{Fix}(T^n)\) 的双射。

由定理 \ref{thm:killo-natural-extension-branch-register}，
\[
(\widehat X,\widehat T)\cong (X^{\leftarrow},\sigma),
\]
因此
\[
\operatorname{Fix}(\widehat T^n)=\mathcal C\bigl(\operatorname{Fix}(\sigma^n)\bigr)
\cong
\operatorname{Fix}(T^n).
\]
若各不动点集合有限，则两侧 Artin--Mazur \(\zeta\)-函数的指数级数逐项相同，故 \(\zeta_{\widehat T}=\zeta_T\)。若拓扑熵由
\[
h_{\mathrm{top}}(S)
=
\limsup_{n\to\infty}\frac1n\log \#\operatorname{Fix}(S^n)
\]
给出，则由于对应周期点计数对一切 \(n\) 完全一致，立得
\(
h_{\mathrm{top}}(\widehat T)=h_{\mathrm{top}}(T)
\)。
证毕。
\end{proof}

\begin{theorem}[有限层 prime-slice 编码的非平凡片数精确极小性]\label{thm:xi-prime-slice-nontrivial-layer-exact-minimality}
\leanverified{paper\_xi\_prime\_slice\_nontrivial\_layer\_exact\_minimality}
设
\[
X_0\xrightarrow{\pi_0}X_1\xrightarrow{\pi_1}\cdots\xrightarrow{\pi_{k-1}}X_k
\]
为有限纤维满射塔。定义非平凡分支层集合
\[
J:=
\left\{
\,0\le j\le k-1:\ \exists\,y\in X_{j+1}\ \text{使}\ \#\pi_j^{-1}(y)\ge 2
\right\}.
\]
考虑任意按层分解的 prime-slice 型映射
\[
\Psi(x_0)=\bigl(x_k,H(x_0)\bigr),
\qquad
H(x_0)=\prod_{j=0}^{k-1}h_j(x_0),
\]
其中每个 \(h_j(x_0)\) 的全部素因子均落在预先指定且两两不交的素数片
\(\mathcal P_j\)，并假设 \(\Psi\) 为单射。定义其非平凡片数
\[
N_{\mathrm{slice}}(\Psi):=\#\{\,j:\ h_j\not\equiv 1\,\}.
\]
则必有严格下界
\[
N_{\mathrm{slice}}(\Psi)\ge |J|.
\]
反之，存在一条这类单射 \(\Psi\) 使得
\[
N_{\mathrm{slice}}(\Psi)=|J|.
\]
故 prime-slice 编码所需的最小有效片数恰等于非平凡分支层数 \(|J|\)。
\end{theorem}
\begin{proof}
先证下界。若 \(j\in J\)，则存在 \(y_j\in X_{j+1}\) 使
\[
\#\pi_j^{-1}(y_j)\ge 2.
\]
取不同点
\[
u_j\neq v_j\in \pi_j^{-1}(y_j).
\]
再固定其余层的选择，并沿塔向下选取兼容前史、向上取公共像，得到两点
\[
x_0',x_0''\in X_0
\]
使它们在第 \(j\) 层的分支选择不同，而在其余各层的对应片语义下保持相同的选取，并且终端宏态同为某个 \(x_k\in X_k\)。
若 \(h_j\equiv 1\)，则第 \(j\) 层对乘积账本不贡献任何区分信息；由于其余层的选择已固定，分层分离便给出
\[
H(x_0')=H(x_0'').
\]
于是
\[
\Psi(x_0')=\Psi(x_0''),
\]
与单射性矛盾。故对每个 \(j\in J\) 都必须有
\[
h_j\not\equiv 1.
\]
因此每个 \(j\in J\) 都必须贡献一个非平凡片，故
\[
N_{\mathrm{slice}}(\Psi)\ge |J|.
\]

再证上界。把 \(J\) 写成严格递增序列
\[
J=\{j_1<j_2<\cdots<j_r\},
\qquad r=|J|.
\]
若 \(j\notin J\)，则 \(\pi_j\) 的每条纤维都是单点；由于 \(\pi_j\) 还是满射，故 \(\pi_j\) 实为双射，从而
\[
x_j=\pi_j^{-1}(x_{j+1})
\]
可由上一层唯一恢复，无需任何额外 prime-slice 信息。

对每个 \(j_\nu\in J\)，沿用定理 \ref{thm:killo-layered-prime-slices-finite-depth} 的构造，取单射
\[
\alpha_{j_\nu}:X_{j_\nu+1}\times\NN\hookrightarrow \mathcal P_{j_\nu}
\]
及纤维枚举
\[
\iota_{j_\nu,y}:\pi_{j_\nu}^{-1}(y)\hookrightarrow \NN.
\]
对 \(x_0\in X_0\) 记其各层像为 \(x_j\in X_j\)，并定义
\[
h_{j_\nu}(x_0):=
\alpha_{j_\nu}\!\bigl(x_{j_\nu+1},\iota_{j_\nu,x_{j_\nu+1}}(x_{j_\nu})\bigr),
\qquad
h_j(x_0):=1\quad(j\notin J).
\]
置
\[
H(x_0):=\prod_{j=0}^{k-1}h_j(x_0)
=
\prod_{\nu=1}^{r}h_{j_\nu}(x_0),
\qquad
\Psi(x_0):=(x_k,H(x_0)).
\]
显然 \(N_{\mathrm{slice}}(\Psi)=r=|J|\)。

现说明 \(\Psi\) 单射。给定 \((x_k,H)\) 属于其像。由素数片两两不交与唯一分解，可从 \(H\) 中逐片抽出全部非平凡因子
\[
h_{j_1},\dots,h_{j_r}.
\]
从顶层 \(x_k\) 开始向下恢复。若当前层 \(j\notin J\)，则
\[
x_j=\pi_j^{-1}(x_{j+1})
\]
唯一确定；若 \(j=j_\nu\in J\)，则由 \(x_{j_\nu+1}\) 与素数因子 \(h_{j_\nu}\) 经
\(\alpha_{j_\nu}\) 的像上反演恢复唯一索引，再经
\(\iota_{j_\nu,x_{j_\nu+1}}\) 的像上反演恢复唯一 \(x_{j_\nu}\in\pi_{j_\nu}^{-1}(x_{j_\nu+1})\)。
如此逐层下推，唯一恢复 \(x_0\)。故 \(\Psi\) 为单射。

综上，最小非平凡片数恰等于 \(|J|\)。证毕。
\end{proof}

\begin{theorem}[Smith 信息亏损谱的离散曲率原子精确等于 \texorpdfstring{$p$}{p}-初等因子重数]\label{thm:xi-smith-loss-discrete-curvature-atoms}
\leanverified{paper\_xi\_smith\_loss\_discrete\_curvature\_atoms}
设 \(A\in M_{m\times n}(\ZZ)\) 且 \(\operatorname{rank}_{\QQ}(A)=m\)。对固定素数 \(p\)，沿用定理
\ref{thm:killo-smith-loss-spectrum} 的记号，令
\[
e_i(p):=v_p(d_i),
\qquad
f_p(k):=\sum_{i=1}^{m}\min(k,e_i(p))
\qquad(k\ge 1),
\]
并约定 \(f_p(0):=0\)。则：
\begin{enumerate}
\item \(f_p:\ZZ_{\ge 0}\to\ZZ_{\ge 0}\) 为离散凹函数，即对每个 \(k\ge 1\) 都有
\[
f_p(k+1)-2f_p(k)+f_p(k-1)\le 0.
\]
\item 对每个 \(k\ge 1\)，其二阶差分的负值恰等于赋值 \(k\) 的初等因子重数：
\[
-\bigl(f_p(k+1)-2f_p(k)+f_p(k-1)\bigr)
=
\#\{i:\ e_i(p)=k\}.
\]
\end{enumerate}
因而 \(f_p\) 的整个图像可以直接视为 \(\{e_i(p)\}\) 的离散 Newton 图，而无需额外反演算法。
\end{theorem}
\begin{proof}
由定理 \ref{thm:killo-smith-loss-spectrum}，
\[
\Delta_p(k):=f_p(k)-f_p(k-1)=\#\{i:\ e_i(p)\ge k\}
\qquad(k\ge 1).
\]
由于右端随 \(k\) 单调不增，故
\[
\Delta_p(k+1)\le \Delta_p(k),
\]
亦即
\[
f_p(k+1)-f_p(k)\le f_p(k)-f_p(k-1).
\]
这等价于
\[
f_p(k+1)-2f_p(k)+f_p(k-1)\le 0,
\]
从而 \(f_p\) 离散凹。

再对 \(\Delta_p\) 作一次差分，得到
\[
\Delta_p(k)-\Delta_p(k+1)
=
\#\{i:\ e_i(p)\ge k\}-\#\{i:\ e_i(p)\ge k+1\}
=
\#\{i:\ e_i(p)=k\}.
\]
另一方面，
\[
\Delta_p(k)-\Delta_p(k+1)
=
\bigl(f_p(k)-f_p(k-1)\bigr)-\bigl(f_p(k+1)-f_p(k)\bigr)
\]
\[
=
-\bigl(f_p(k+1)-2f_p(k)+f_p(k-1)\bigr).
\]
比较两式即得
\[
-\bigl(f_p(k+1)-2f_p(k)+f_p(k-1)\bigr)
=
\#\{i:\ e_i(p)=k\}.
\]
证毕。
\end{proof}

\begin{theorem}[模数核大小 Dirichlet 生成函数的 Euler 因子分解与单极点定律]\label{thm:xi-smith-kernel-dirichlet-euler-single-pole}
\leanverified{paper\_xi\_smith\_kernel\_dirichlet\_euler\_single\_pole}
沿用推论 \ref{cor:killo-kernel-size-general-modulus} 的设定。定义
\[
\mathcal Z_A(s):=\sum_{b\ge 1}\frac{\#\ker(A_b)}{b^s},
\qquad
\Delta_A:=\prod_{i=1}^{m}d_i.
\]
对每个素数 \(p\)，记
\[
e_i(p):=v_p(d_i),
\qquad
Z_{A,p}(s):=
\sum_{k\ge 0}p^{\,k(n-m-s)+\sum_{i=1}^{m}\min(k,e_i(p))}.
\]
则成立：
\begin{enumerate}
\item \(\mathcal Z_A(s)\) 为积性 Dirichlet 级数，并具有 Euler 因子分解
\[
\mathcal Z_A(s)=\prod_{p}Z_{A,p}(s)
=
\zeta\!\bigl(s-(n-m)\bigr)
\prod_{p\mid \Delta_A}\bigl(1-p^{\,n-m-s}\bigr)Z_{A,p}(s).
\]
\item 对每个固定素数 \(p\)，局部因子 \(Z_{A,p}(s)\) 是 \(p^{-s}\) 的有理函数。若按非降序写
\[
0=:e_0(p)\le e_1(p)\le\cdots\le e_m(p),
\qquad
E_j(p):=\sum_{i=1}^{j}e_i(p),
\]
则有显式闭式
\[
Z_{A,p}(s)
=
\sum_{j=0}^{m-1}
p^{E_j(p)}
\frac{
p^{e_j(p)(n-j-s)}-p^{e_{j+1}(p)(n-j-s)}
}{
1-p^{\,n-j-s}
}
\]
\[
\qquad\qquad
+\;
p^{E_m(p)}
\frac{
p^{e_m(p)(n-m-s)}
}{
1-p^{\,n-m-s}
}.
\]
\item \(\mathcal Z_A(s)\) 在半平面 \(\Re(s)>n-m\) 上亚纯，并且在
\[
s_0:=n-m+1
\]
处具有唯一简单极点，其留数为
\[
\Res_{s=s_0}\mathcal Z_A(s)
=
\prod_{p\mid \Delta_A}\Bigl(1-\frac1p\Bigr)Z_{A,p}(s_0).
\]
\item 因而由 Wiener--Ikehara 定理，
\[
\sum_{b\le X}\#\ker(A_b)
\sim
\frac{\mathfrak R_A}{n-m+1}\,X^{n-m+1},
\qquad
\mathfrak R_A:=\Res_{s=s_0}\mathcal Z_A(s).
\]
\end{enumerate}
\end{theorem}
\begin{proof}
由推论 \ref{cor:killo-kernel-size-general-modulus}，
\[
\#\ker(A_b)=b^{\,n-m}\prod_{i=1}^{m}\gcd(d_i,b).
\]
右端关于 \(b\) 为积性函数，因此 Dirichlet 级数 \(\mathcal Z_A(s)\) 具有 Euler 乘积分解
\[
\mathcal Z_A(s)=\prod_p Z_{A,p}(s),
\]
其中
\[
Z_{A,p}(s)
=
\sum_{k\ge 0}\frac{\#\ker(A_{p^k})}{p^{ks}}
=
\sum_{k\ge 0}p^{\,k(n-m-s)+\sum_{i=1}^{m}\min(k,e_i(p))}.
\]
这给出第 \(1\) 条的第一式。

若 \(p\nmid \Delta_A\)，则所有 \(e_i(p)=0\)，故
\[
Z_{A,p}(s)=\sum_{k\ge 0}p^{k(n-m-s)}=\frac{1}{1-p^{\,n-m-s}}.
\]
把这些“好素数”因子提出，即得
\[
\mathcal Z_A(s)
=
\zeta\!\bigl(s-(n-m)\bigr)
\prod_{p\mid \Delta_A}\bigl(1-p^{\,n-m-s}\bigr)Z_{A,p}(s),
\]
从而得到第 \(1\) 条的第二式。

现证第 \(2\) 条。对固定素数 \(p\)，把 \(k\)-求和按区间
\[
e_j(p)\le k<e_{j+1}(p)\qquad(0\le j\le m-1)
\]
以及尾区间
\[
k\ge e_m(p)
\]
分块。在第 \(j\) 个有限区间上，有
\[
\sum_{i=1}^{m}\min(k,e_i(p))
=
E_j(p)+(m-j)k,
\]
故对应求和块为
\[
\sum_{k=e_j(p)}^{e_{j+1}(p)-1}p^{\,E_j(p)}p^{\,k(n-j-s)}
=
p^{E_j(p)}
\frac{
p^{e_j(p)(n-j-s)}-p^{e_{j+1}(p)(n-j-s)}
}{
1-p^{\,n-j-s}
}.
\]
在尾区间上则有
\[
\sum_{i=1}^{m}\min(k,e_i(p))=E_m(p),
\]
因而
\[
\sum_{k\ge e_m(p)}p^{\,k(n-m-s)+E_m(p)}
=
p^{E_m(p)}
\frac{
p^{e_m(p)(n-m-s)}
}{
1-p^{\,n-m-s}
}.
\]
把这些分块相加即得所示闭式。由于它是有限和与一个几何尾项的组合，故 \(Z_{A,p}(s)\) 的确是 \(p^{-s}\) 的有理函数。

再证第 \(3\) 条。把
\[
F_A(s):=\prod_{p\mid \Delta_A}\bigl(1-p^{\,n-m-s}\bigr)Z_{A,p}(s)
\]
记作坏素数修正因子。由上式可见，对每个 \(p\mid\Delta_A\)，
\[
\bigl(1-p^{\,n-m-s}\bigr)Z_{A,p}(s)
\]
在 \(s\) 上是有限个指数函数 \(p^{-rs}\) 的线性组合，故在整个复平面上全纯。于是 \(F_A(s)\) 为有限乘积，亦为全纯函数。故
\[
\mathcal Z_A(s)=\zeta\!\bigl(s-(n-m)\bigr)F_A(s)
\]
在 \(\Re(s)>n-m\) 上的唯一极点只能来自 \(\zeta(s-(n-m))\) 在 \(s_0=n-m+1\) 处的简单极点。

在 \(s=s_0\) 处，
\[
\Res_{s=s_0}\zeta\!\bigl(s-(n-m)\bigr)=1,
\]
故
\[
\Res_{s=s_0}\mathcal Z_A(s)=F_A(s_0)
=
\prod_{p\mid \Delta_A}\Bigl(1-\frac1p\Bigr)Z_{A,p}(s_0).
\]
第 \(3\) 条成立。

最后，由 \(\#\ker(A_b)\ge 0\) 对一切 \(b\) 成立，且 \(\mathcal Z_A(s)\) 在 \(\Re(s)\ge s_0\) 上除 \(s_0\) 之外无其他极点，Wiener--Ikehara 定理适用，从而得到
\[
\sum_{b\le X}\#\ker(A_b)
\sim
\frac{\mathfrak R_A}{s_0}\,X^{s_0}
=
\frac{\mathfrak R_A}{n-m+1}\,X^{n-m+1},
\]
其中 \(\mathfrak R_A=\Res_{s=s_0}\mathcal Z_A(s)\)。证毕。
\end{proof}

\begin{theorem}[模核 Dirichlet 级数的 Euler--Stokes 尾分解与唯一极点留数]\label{thm:xi-smith-kernel-dirichlet-euler-stokes-tail-residue}
\leanverified{paper\_xi\_smith\_kernel\_dirichlet\_euler\_stokes\_tail\_residue}
沿用定理 \ref{thm:xi-smith-kernel-dirichlet-euler-single-pole} 的记号。对每个素数 \(p\)，记
\[
\mathcal Z_{A,p}(s):=
\sum_{k\ge 0}p^{-k(s-(n-m))+\sum_{i=1}^{m}\min(k,e_i(p))},
\qquad
E_p:=\max_i e_i(p).
\]
则
\[
\boxed{\;
\mathcal Z_A(s)
=
\zeta\!\bigl(s-(n-m)\bigr)
\prod_{p\mid \Delta_A}
\Bigl(1-p^{-(s-(n-m))}\Bigr)\,
\mathcal Z_{A,p}(s).\;}
\]
并且每个坏素数局部因子都具有显式“有限前缀 + 几何尾部”分解
\[
\boxed{\;
\mathcal Z_{A,p}(s)
=
\sum_{k=0}^{E_p-1}
p^{-k(s-(n-m))+\sum_{i=1}^{m}\min(k,e_i(p))}
\;+\;
\frac{
p^{-E_p(s-(n-m))+\sum_{i=1}^{m}e_i(p)}
}{
1-p^{-(s-(n-m))}}
\;}
\]
从而 \(\mathcal Z_A(s)\) 在半平面 \(\Re(s)>n-m\) 上亚纯，并且在
\[
s_0:=n-m+1
\]
处具有唯一简单极点，其留数可写为
\[
\boxed{\;
\Res_{s=s_0}\mathcal Z_A(s)
=
\prod_{p\mid \Delta_A}
\Bigl(1-\frac1p\Bigr)
\sum_{k\ge 0}p^{-k+\sum_{i=1}^{m}\min(k,e_i(p))}.
\;}
\]
\end{theorem}
\begin{proof}
定理 \ref{thm:xi-smith-kernel-dirichlet-euler-single-pole} 已给出
\[
\mathcal Z_A(s)
=
\zeta\!\bigl(s-(n-m)\bigr)
\prod_{p\mid \Delta_A}
\Bigl(1-p^{-(s-(n-m))}\Bigr)\,
Z_{A,p}(s),
\]
其中
\[
Z_{A,p}(s)
=
\sum_{k\ge 0}p^{-k(s-(n-m))+\sum_{i=1}^{m}\min(k,e_i(p))}.
\]
故只需把 \(Z_{A,p}(s)\) 改记为 \(\mathcal Z_{A,p}(s)\)。

若 \(k\ge E_p\)，则对一切 \(i\) 都有 \(\min(k,e_i(p))=e_i(p)\)，于是
\[
\sum_{i=1}^{m}\min(k,e_i(p))=\sum_{i=1}^{m}e_i(p)
\]
稳定为常数。因此局部级数尾部成为公比
\(
p^{-(s-(n-m))}
\)
的几何级数：
\[
\sum_{k\ge E_p}
p^{-k(s-(n-m))+\sum_{i=1}^{m}e_i(p)}
=
\frac{
p^{-E_p(s-(n-m))+\sum_{i=1}^{m}e_i(p)}
}{
1-p^{-(s-(n-m))}}.
\]
加上前 \(E_p\) 项即得所示 Euler--Stokes 尾分解。

定理 \ref{thm:xi-smith-kernel-dirichlet-euler-single-pole} 同时已证明 \(\mathcal Z_A(s)\) 在 \(\Re(s)>n-m\) 上只有 \(s_0=n-m+1\) 这一唯一简单极点，并满足
\[
\Res_{s=s_0}\mathcal Z_A(s)
=
\prod_{p\mid \Delta_A}\Bigl(1-\frac1p\Bigr)\mathcal Z_{A,p}(s_0).
\]
把
\[
s_0-(n-m)=1
\]
代入 \(\mathcal Z_{A,p}(s)\) 即得
\[
\mathcal Z_{A,p}(s_0)
=
\sum_{k\ge 0}p^{-k+\sum_{i=1}^{m}\min(k,e_i(p))},
\]
从而得到所示留数公式。证毕。
\end{proof}

由此，附录中的四段对象在本链上闭合为
\[
\text{有限纤维自映射}
\Longrightarrow
\text{自然扩张寄存器实现}
\Longrightarrow
\text{周期点与 Artin--Mazur }\zeta\text{ 的严格守恒},
\]
\[
\text{Smith 核大小谱 }K_p(k)
\Longrightarrow
\text{离散曲率原子 }\#\{i:e_i(p)=k\}
\Longrightarrow
\text{Dirichlet 生成函数的 Euler 因子化}
\Longrightarrow
\text{全模数核增长的单极点定律}.
\]

\paragraph{bin-fold 稳定统计的一位塌缩与有限局部化的 \(p\)-进核刚性}
下列结果把定理 \ref{thm:fold-bin-two-state-asymptotic}、命题 \ref{prop:fold-bin-lastbit-sufficient-tv}、
定理 \ref{thm:fold-bin-renyi-divergence-limit}、推论 \ref{cor:xi-localized-integers-order-dual-quotient}、
定理 \ref{thm:cdim-star-z1s-dual-extension} 与定理
\ref{thm:conclusion-cdim-finite-rank-extension} 组织为同一条链：
可见稳定统计在指数精度上塌缩为末位两层，而可逆算术的分类信息则完全驻留在全不连通
\(p\)-进核中。

\begin{theorem}[稳定型分布的末位充分性与两层模型塌缩]\label{thm:xi-fold-lastbit-two-layer-collapse}
设
\[
X_m:=\{w\in\{0,1\}^m:\ w_iw_{i+1}=0\ \ (1\le i<m)\},
\qquad
\mu_m(w):=\frac{d_m^{\mathrm{bin}}(w)}{2^m}.
\]
定义参考分布
\[
\nu_m(w):=\varphi^{-m-w_m},
\qquad
w=(w_1,\dots,w_m)\in X_m.
\]
则 \(\nu_m\) 是 \(X_m\) 上的概率分布，并且对一切充分大的 \(m\) 与一切 \(w\in X_m\) 都有一致估计
\[
\mu_m(w)=\nu_m(w)\Bigl(1+O\!\Bigl(\Bigl(\frac{\varphi}{2}\Bigr)^m\Bigr)\Bigr).
\]
从而
\[
D_{\mathrm{TV}}(\mu_m,\nu_m)=O\!\Bigl(\Bigl(\frac{\varphi}{2}\Bigr)^m\Bigr).
\]
进一步，对每个 \(b\in\{0,1\}\) 记
\[
X_m^{(b)}:=\{w\in X_m:\ w_m=b\},
\]
并令 \(u_m^{(b)}\) 为 \(X_m^{(b)}\) 上的均匀分布。则有
\[
D_{\mathrm{TV}}\bigl(\mu_m(\,\cdot\,|X_m^{(b)}),u_m^{(b)}\bigr)
=O\!\Bigl(\Bigl(\frac{\varphi}{2}\Bigr)^m\Bigr).
\]
换言之，在稳定型层面，除末位 \(w_m\) 外的全部组合自由度都以指数速度趋于层内均匀化。
\end{theorem}
\begin{proof}
由定理 \ref{thm:fold-bin-two-state-asymptotic}，
\[
d_m^{\mathrm{bin}}(w)=\varphi^{-w_m}\Bigl(\frac{2}{\varphi}\Bigr)^m+O(1)
\]
对 \(w\in X_m\) 一致成立。两边除以 \(2^m\) 得
\[
\mu_m(w)=\varphi^{-m-w_m}+O(2^{-m}).
\]
由于
\[
\varphi^{-m-1}\le \nu_m(w)\le \varphi^{-m},
\]
故可把余项重写为
\[
\mu_m(w)=\nu_m(w)\Bigl(1+O\!\Bigl(\Bigl(\frac{\varphi}{2}\Bigr)^m\Bigr)\Bigr),
\]
且误差在 \(w\in X_m\) 上一致。

现证 \(\nu_m\) 归一。由
\[
|X_m^{(0)}|=F_{m+1},
\qquad
|X_m^{(1)}|=F_m
\]
与 Binet 公式可得恒等式
\[
F_{m+1}+\varphi^{-1}F_m=\varphi^m.
\]
因此
\[
\sum_{w\in X_m}\nu_m(w)
=\varphi^{-m}|X_m^{(0)}|+\varphi^{-m-1}|X_m^{(1)}|
=\varphi^{-m}\bigl(F_{m+1}+\varphi^{-1}F_m\bigr)=1.
\]

再由逐点相对误差估计求和，
\[
\sum_{w\in X_m}|\mu_m(w)-\nu_m(w)|
\le
C\Bigl(\frac{\varphi}{2}\Bigr)^m\sum_{w\in X_m}\nu_m(w)
=
C\Bigl(\frac{\varphi}{2}\Bigr)^m,
\]
故
\[
D_{\mathrm{TV}}(\mu_m,\nu_m)
=O\!\Bigl(\Bigl(\frac{\varphi}{2}\Bigr)^m\Bigr).
\]

最后固定 \(b\in\{0,1\}\)。由于 \(\nu_m\) 在 \(X_m^{(b)}\) 上为常数，故
\[
\nu_m(\,\cdot\,|X_m^{(b)})=u_m^{(b)}.
\]
将前述相对误差在 \(X_m^{(b)}\) 上求和得
\[
\mu_m(X_m^{(b)})=\nu_m(X_m^{(b)})\Bigl(1+O\!\Bigl(\Bigl(\frac{\varphi}{2}\Bigr)^m\Bigr)\Bigr).
\]
于是对任意 \(w\in X_m^{(b)}\)，
\[
\mu_m(w|X_m^{(b)})
=\frac{\mu_m(w)}{\mu_m(X_m^{(b)})}
=\frac{\nu_m(w)}{\nu_m(X_m^{(b)})}
\Bigl(1+O\!\Bigl(\Bigl(\frac{\varphi}{2}\Bigr)^m\Bigr)\Bigr)
=u_m^{(b)}(w)\Bigl(1+O\!\Bigl(\Bigl(\frac{\varphi}{2}\Bigr)^m\Bigr)\Bigr),
\]
且误差在 \(w\in X_m^{(b)}\) 上一致。再对 \(X_m^{(b)}\) 求和即得条件化后的总变差界。
\end{proof}
\leanverified{paper\_xi\_fold\_lastbit\_two\_layer\_collapse}

\begin{theorem}[Rényi 熵常数谱的闭式]\label{thm:xi-fold-renyi-constant-spectrum-closed-form}
对每个固定的 \(q>0\)、\(q\neq 1\)，定义
\[
H_q(\mu_m):=\frac{1}{1-q}\log\sum_{w\in X_m}\mu_m(w)^q.
\]
则有闭式展开
\[
H_q(\mu_m)
=
m\log\varphi+c(q)+O\!\Bigl(\Bigl(\frac{\varphi}{2}\Bigr)^m\Bigr),
\]
其中
\[
c(q):=\frac{1}{1-q}\log\!\Bigl(\frac{\varphi+\varphi^{-q}}{\sqrt5}\Bigr).
\]
并且 \(c(q)\) 在 \(q=1\) 处可连续延拓为
\[
c(1)=\frac{\log\varphi}{\varphi\sqrt5},
\]
相应地 Shannon 熵满足
\[
H(\mu_m)=m\log\varphi+\frac{\log\varphi}{\varphi\sqrt5}+o(1).
\]
因此 bin-fold 的全部 Rényi 熵族不仅具有统一的塌缩斜率 \(\log\varphi\)，而且其二阶常数项亦可闭式写出。
\end{theorem}
\begin{proof}
令 \(u_m(w):=1/|X_m|\) 为 \(X_m\) 上的均匀分布。
由定理 \ref{thm:fold-bin-renyi-divergence-limit}，对每个固定 \(q>0\)、\(q\neq 1\)，有
\[
D_q(\mu_m\Vert u_m)=D_q^\infty+O\!\Bigl(\Bigl(\frac{\varphi}{2}\Bigr)^m\Bigr),
\]
其中
\[
D_q^\infty
=
\frac{(2q-2)\log\varphi+\log(\varphi+\varphi^{-q})-q\log\sqrt5}{q-1}.
\]
由于 \(u_m\) 为均匀分布，Rényi 散度与 Rényi 熵满足恒等式
\[
D_q(\mu_m\Vert u_m)=\log|X_m|-H_q(\mu_m).
\]
又因 \(|X_m|=F_{m+2}\)，Binet 公式给出
\[
\log|X_m|=(m+2)\log\varphi-\frac12\log 5+O(\varphi^{-2m}).
\]
代入即得
\[
H_q(\mu_m)
=
m\log\varphi+
\Bigl(2\log\varphi-\frac12\log 5-D_q^\infty\Bigr)
+O\!\Bigl(\Bigl(\frac{\varphi}{2}\Bigr)^m\Bigr).
\]
把括号内常数化简，得到
\[
2\log\varphi-\frac12\log 5-D_q^\infty
=
\frac{1}{1-q}\log\!\Bigl(\frac{\varphi+\varphi^{-q}}{\sqrt5}\Bigr)
=c(q).
\]

当 \(q=1\) 时，命题 \ref{prop:fold-bin-kl-constant} 给出
\[
H(\mu_m)=\log|X_m|-D_{\mathrm{KL}}(\mu_m\Vert u_m)
=m\log\varphi+\frac{\log\varphi}{\varphi\sqrt5}+o(1).
\]
最后，对 \(c(q)\) 在 \(q=1\) 处应用 l'Hospital 法则，便得
\[
\lim_{q\to 1}c(q)=\frac{\log\varphi}{\varphi\sqrt5}.
\]
证毕。
\end{proof}

\paragraph{escort 温度轴上的末位 Gibbs 极限、两原子涨落与熵常数闭式}
前两条处理原始推前分布 \(\mu_m=\mu_{m,1}\)。
对整条 escort 温度轴，同一两态结构实际上同时控制末位 Gibbs 偏置、对数多重度的极限分布、一阶局部信息修正、Rényi--Shannon 常数以及局部指数谱。

\begin{theorem}[escort 族的末位 Gibbs 极限]\label{thm:xi-fold-escort-lastbit-gibbs-limit}
\leanverified{paper\_xi\_fold\_escort\_lastbit\_gibbs\_limit}
对任意固定 \(q\ge 0\)，定义 escort 测度
\[
\mu_{m,q}(w):=\frac{(d_m^{\mathrm{bin}}(w))^{q}}{S_q^{\mathrm{bin}}(m)},
\qquad
w\in X_m.
\]
则有极限
\[
\mu_{m,q}(w_m=1)\longrightarrow \frac{1}{1+\varphi^{q+1}},
\qquad
\mu_{m,q}(w_m=0)\longrightarrow \frac{\varphi^{q+1}}{1+\varphi^{q+1}}.
\]
特别地，\(\mu_{m,1}=\mu_m\)，并且
\[
\mu_m(w_m=1)\longrightarrow \frac{1}{1+\varphi^2}=\frac{1}{\varphi\sqrt5},
\qquad
\mu_m(w_m=0)\longrightarrow \frac{\varphi^2}{1+\varphi^2}=\frac{\varphi}{\sqrt5}.
\]
\end{theorem}
\begin{proof}
该极限正是命题 \ref{prop:fold-bin-escort-lastbit} 的第一部分。
对 \(q=1\)，只需使用恒等式 \(1+\varphi^2=\varphi\sqrt5\) 即可得到所示化简。
\end{proof}

\begin{theorem}[escort 对数多重度的两原子极限分布]\label{thm:xi-fold-escort-log-multiplicity-two-atom}
\leanverified{paper\_xi\_fold\_escort\_log\_multiplicity\_two\_atom}
沿用定理 \ref{thm:xi-fold-escort-lastbit-gibbs-limit} 的记号，并令 \(W_{m,q}\sim\mu_{m,q}\)。
定义中心化对数多重度
\[
X_m^{(q)}:=\log d_m^{\mathrm{bin}}(W_{m,q})-m\log\!\Bigl(\frac{2}{\varphi}\Bigr).
\]
则
\[
X_m^{(q)}\Longrightarrow \nu_q
:=
\frac{\varphi^{q+1}}{1+\varphi^{q+1}}\delta_{0}
+\frac{1}{1+\varphi^{q+1}}\delta_{-\log\varphi}.
\]
等价地，对任意固定 \(s\in\CC\)，有
\[
\lim_{m\to\infty}\mathbb E_{\mu_{m,q}}\!\bigl[e^{sX_m^{(q)}}\bigr]
=
\frac{\varphi^{q+1}+\varphi^{-s}}{1+\varphi^{q+1}}.
\]
从而
\[
\lim_{m\to\infty}\mathbb E_{\mu_{m,q}}[X_m^{(q)}]
=-\frac{\log\varphi}{1+\varphi^{q+1}},
\]
\[
\lim_{m\to\infty}\operatorname{Var}_{\mu_{m,q}}(X_m^{(q)})
=
\frac{\varphi^{q+1}}{(1+\varphi^{q+1})^2}(\log\varphi)^2.
\]
\end{theorem}
\begin{proof}
由定理 \ref{thm:fold-bin-two-state-asymptotic} 的一致对数展开，
\[
\log d_m^{\mathrm{bin}}(w)
=
m\log\!\Bigl(\frac{2}{\varphi}\Bigr)-w_m\log\varphi+O\!\Bigl(\Bigl(\frac{\varphi}{2}\Bigr)^m\Bigr)
\]
对 \(w\in X_m\) 一致成立，因此
\[
X_m^{(q)}=-W_{m,q}(m)\log\varphi+o_{\mathbb P}(1).
\]
再与定理 \ref{thm:xi-fold-escort-lastbit-gibbs-limit} 的末位极限分布合并，即得所示两原子极限。
变换公式、均值与方差都只是该两点分布的直接计算。
\end{proof}

\begin{corollary}[局部信息密度的一阶离散修正律]\label{cor:xi-fold-local-information-density-first-order-discrete-law}
令 \(W_m\sim\mu_m\)，并定义
\[
\gamma_m:=\frac{1}{m}\log d_m^{\mathrm{bin}}(W_m).
\]
再令 \(B\) 为参数
\[
\mathbb P(B=1)=\frac{1}{\varphi\sqrt5},
\qquad
\mathbb P(B=0)=\frac{\varphi}{\sqrt5}
\]
的 Bernoulli 变量。则有分布收敛
\[
m\left(\gamma_m-\log\!\frac{2}{\varphi}\right)
\Longrightarrow
-(\log\varphi)B.
\]
并且
\[
\lim_{m\to\infty}
m\left(\mathbb E[\gamma_m]-\log\!\frac{2}{\varphi}\right)
=
-\frac{\log\varphi}{\varphi\sqrt5},
\]
\[
\lim_{m\to\infty}
m^2\operatorname{Var}(\gamma_m)
=
(\log\varphi)^2
\frac{1}{\varphi\sqrt5}
\left(
1-\frac{1}{\varphi\sqrt5}
\right).
\]
因此 \(\gamma_m\) 的首个非平凡修正并不呈现 \(\sqrt m\)-Gaussian 涨落，而是在 \(m^{-1}\) 尺度上塌缩为末位几何所决定的两点律。
\end{corollary}
\begin{proof}
在定理 \ref{thm:xi-fold-escort-log-multiplicity-two-atom} 中取 \(q=1\)。
由于 \(\mu_{m,1}=\mu_m\)，且
\[
X_m^{(1)}
=
\log d_m^{\mathrm{bin}}(W_m)-m\log\!\Bigl(\frac{2}{\varphi}\Bigr)
=
m\left(\gamma_m-\log\!\frac{2}{\varphi}\right),
\]
故该定理立即给出
\[
m\left(\gamma_m-\log\!\frac{2}{\varphi}\right)
\Longrightarrow
\frac{\varphi}{\sqrt5}\delta_0+\frac{1}{\varphi\sqrt5}\delta_{-\log\varphi},
\]
这正是随机变量 \(-(\log\varphi)B\) 的分布。

又由定理 \ref{thm:xi-fold-escort-log-multiplicity-two-atom} 的证明，
\[
X_m^{(1)}=-W_m(m)\log\varphi+O\!\Bigl(\Bigl(\frac{\varphi}{2}\Bigr)^m\Bigr),
\]
且误差一致于样本点，从而 \((X_m^{(1)})_{m\ge 1}\) 一致有界。
因此其一、二阶矩都收敛到极限分布的相应矩。于是
\[
\lim_{m\to\infty}\mathbb E[X_m^{(1)}]
=
\mathbb E[-(\log\varphi)B]
=
-\frac{\log\varphi}{\varphi\sqrt5},
\]
\[
\lim_{m\to\infty}\operatorname{Var}(X_m^{(1)})
=
\operatorname{Var}((\log\varphi)B)
=
(\log\varphi)^2
\frac{1}{\varphi\sqrt5}
\left(
1-\frac{1}{\varphi\sqrt5}
\right).
\]
最后用
\[
X_m^{(1)}=m\left(\gamma_m-\log\!\frac{2}{\varphi}\right)
\]
即可得到所示均值与方差极限。证毕。
\end{proof}

\begin{theorem}[escort 温度族的 Rényi--Shannon 常数闭式]\label{thm:xi-fold-escort-renyi-shannon-constants}
固定 \(q\ge 0\)，并对任意 \(r>0\)、\(r\neq 1\) 定义 Rényi 熵
\[
H_r(\mu_{m,q}):=\frac{1}{1-r}\log\sum_{w\in X_m}\mu_{m,q}(w)^r.
\]
则有展开
\[
H_r(\mu_{m,q})=m\log\varphi+C_{q,r}+O\!\Bigl(\Bigl(\frac{\varphi}{2}\Bigr)^m\Bigr),
\]
其中
\[
C_{q,r}
:=
\frac{1}{1-r}\left[
\log\!\Bigl(\frac{\varphi+\varphi^{-qr}}{\sqrt5}\Bigr)
-r\log\!\Bigl(\frac{\varphi+\varphi^{-q}}{\sqrt5}\Bigr)
\right].
\]
相应的 Shannon 熵满足
\[
H(\mu_{m,q})=m\log\varphi+C_q+O\!\Bigl(\Bigl(\frac{\varphi}{2}\Bigr)^m\Bigr),
\]
其中
\[
C_q
:=
\log\!\Bigl(\frac{\varphi+\varphi^{-q}}{\sqrt5}\Bigr)
+\frac{q\log\varphi}{1+\varphi^{q+1}}.
\]
\end{theorem}
\begin{proof}
由定义
\[
\sum_{w\in X_m}\mu_{m,q}(w)^r=\frac{S_{qr}^{\mathrm{bin}}(m)}{(S_q^{\mathrm{bin}}(m))^r}.
\]
记
\[
A(a):=\log\!\Bigl(\frac{\varphi+\varphi^{-a}}{\sqrt5}\Bigr).
\]
由命题 \ref{prop:fold-bin-power-sum-asymptotic}，
\[
\log S_a^{\mathrm{bin}}(m)
=
am\log\!\Bigl(\frac{2}{\varphi}\Bigr)+m\log\varphi+A(a)
+O\!\Bigl(\Bigl(\frac{\varphi}{2}\Bigr)^m\Bigr).
\]
因此
\[
H_r(\mu_{m,q})
=
m\log\varphi+\frac{A(qr)-rA(q)}{1-r}
+O\!\Bigl(\Bigl(\frac{\varphi}{2}\Bigr)^m\Bigr),
\]
即得 \(C_{q,r}\)。

对 Shannon 熵，利用恒等式
\[
H(\mu_{m,q})=\log S_q^{\mathrm{bin}}(m)-q\,\mathbb E_{\mu_{m,q}}\!\bigl[\log d_m^{\mathrm{bin}}(W_{m,q})\bigr]
\]
并代入命题 \ref{prop:fold-bin-escort-lastbit} 给出的纤维信息展开，即得所示 \(C_q\)。
\end{proof}

\begin{theorem}[escort 全族的严格单点多重分形塌缩]\label{thm:xi-fold-escort-single-point-multifractal-collapse}
对每个固定 \(q\ge 0\)，定义局部指数
\[
\alpha_{m,q}(w):=-\frac{1}{m}\log\mu_{m,q}(w),
\qquad
w\in X_m.
\]
则有一致极限
\[
\sup_{w\in X_m}\bigl|\alpha_{m,q}(w)-\log\varphi\bigr|\longrightarrow 0.
\]
更精确地，
\[
\alpha_{m,q}(w)
=
\log\varphi
+\frac{q\,w_m\log\varphi+\log\!\bigl(\frac{\varphi+\varphi^{-q}}{\sqrt5}\bigr)}{m}
+o\!\Bigl(\frac{1}{m}\Bigr)
\]
在 \(w\in X_m\) 上一致成立。

若进一步定义离散谱函数
\[
f_q(\alpha)
:=
\lim_{\varepsilon\downarrow 0}\limsup_{m\to\infty}
\frac{1}{m}\log
\#\bigl\{w\in X_m:\ |\alpha_{m,q}(w)-\alpha|<\varepsilon\bigr\},
\]
则
\[
f_q(\alpha)=
\begin{cases}
\log\varphi,& \alpha=\log\varphi,\\[1mm]
-\infty,& \alpha\neq \log\varphi.
\end{cases}
\]
\end{theorem}
\begin{proof}
由定义
\[
\alpha_{m,q}(w)=-\frac{q}{m}\log d_m^{\mathrm{bin}}(w)+\frac{1}{m}\log S_q^{\mathrm{bin}}(m).
\]
再将定理 \ref{thm:fold-bin-two-state-asymptotic} 的一致对数展开与上一个证明中的
\(\log S_q^{\mathrm{bin}}(m)\) 展开代入，即得所示统一公式。
因此 \(\sup_{w}|\alpha_{m,q}(w)-\log\varphi|\to 0\)。

若 \(\alpha\neq\log\varphi\)，取 \(\varepsilon<|\alpha-\log\varphi|/2\)，则对充分大 \(m\) 上述集合为空，
故 \(f_q(\alpha)=-\infty\)。
若 \(\alpha=\log\varphi\)，则对任意固定 \(\varepsilon>0\)，充分大 \(m\) 时全部 \(w\in X_m\) 都落入该带，于是
\[
f_q(\log\varphi)=\lim_{m\to\infty}\frac{1}{m}\log|X_m|=\log\varphi.
\]
证毕。
\end{proof}

\begin{corollary}[escort--uniform 相对熵的显式温度律]\label{cor:xi-fold-escort-uniform-kl-temperature-law}
令 \(u_m\) 为 \(X_m\) 上的均匀分布。则对每个固定 \(q\ge 0\)，有
\[
D_{\mathrm{KL}}(\mu_{m,q}\Vert u_m)
=
D(q)+O\!\Bigl(\Bigl(\frac{\varphi}{2}\Bigr)^m\Bigr),
\]
其中
\[
D(q)
:=
2\log\varphi-\log(\varphi+\varphi^{-q})
-\frac{q\log\varphi}{1+\varphi^{q+1}}.
\]
特别地，当 \(q=1\) 时，
\[
D(1)=2\log\varphi-\frac12\log 5-\frac{\log\varphi}{\varphi\sqrt5},
\]
从而恢复命题 \ref{prop:fold-bin-kl-constant} 的极限常数。
\end{corollary}
\begin{proof}
由恒等式
\[
D_{\mathrm{KL}}(\mu_{m,q}\Vert u_m)=\log|X_m|-H(\mu_{m,q})
\]
以及定理 \ref{thm:xi-fold-escort-renyi-shannon-constants}，
\[
\log|X_m|=(m+2)\log\varphi-\frac12\log 5+O(\varphi^{-2m}),
\]
即可得到所示公式。
当 \(q=1\) 时，只需再用 \(\varphi+\varphi^{-1}=\sqrt5\) 与 \(1+\varphi^2=\varphi\sqrt5\) 化简。
\end{proof}

\begin{theorem}[有限素数局部化的 \(p\)-进核刚性分类]\label{thm:xi-localized-integers-padic-kernel-rigidity}
对有限素数集 \(S\subset\mathbb P\)，记
\[
G_S:=\ZZ[S^{-1}],
\qquad
K_S:=\prod_{p\in S}\ZZ_p.
\]
则有短正合列
\[
0\longrightarrow K_S\longrightarrow \widehat{G_S}\longrightarrow \TT\longrightarrow 0,
\qquad
(\widehat{G_S})^0\cong \TT,
\qquad
\cdim_{\mathrm{fr}}(G_S)=1.
\]
并且对任意有限素数集 \(S,T\subset\mathbb P\)，成立
\[
G_S\cong G_T
\iff
K_S\cong K_T
\iff
S=T.
\]
因此有限局部化群的连通相位部分完全相同，而其全部可分辨算术信息都由全不连通
\(p\)-进核 \(K_S\) 完全承载。
\end{theorem}
\begin{proof}
定理 \ref{thm:cdim-star-z1s-dual-extension} 给出短正合列
\[
0\longrightarrow K_S\longrightarrow \widehat{G_S}\longrightarrow \TT\longrightarrow 0.
\]
由于 \(K_S\) 全不连通，连通分量只能映到 \(\TT\) 的连通分量，故
\(
(\widehat{G_S})^0\cong \TT
\)。
另一方面，\(G_S\otimes_{\ZZ}\QQ\cong \QQ\) 为秩 \(1\) 的有限秩无挠群，故由定理
\ref{thm:conclusion-cdim-finite-rank-extension}，
\[
\cdim_{\mathrm{fr}}(G_S)=1.
\]

若 \(K_S\cong K_T\)，则对任意素数 \(p\) 考察商群 \(K_S/pK_S\)。
当 \(p\in S\) 时，
\[
\ZZ_p/p\ZZ_p\cong \FF_p\neq 0;
\]
而当 \(q\neq p\) 时，\(p\) 在 \(\ZZ_q\) 中可逆，故
\[
\ZZ_q/p\ZZ_q=0.
\]
于是
\[
K_S/pK_S\neq 0
\iff
p\in S.
\]
同理
\(
K_T/pK_T\neq 0\iff p\in T
\)。
由 \(K_S\cong K_T\) 可知上述商群同时非零或同时为零，故 \(S=T\)。

反向蕴含 \(S=T\Rightarrow G_S\cong G_T\) 与 \(S=T\Rightarrow K_S\cong K_T\) 显然。
而 \(G_S\cong G_T\Rightarrow S=T\) 亦可由推论
\ref{cor:xi-localized-integers-order-dual-quotient} 直接得到。
故三者等价。证毕。
\end{proof}
\leanverified{paper\_xi\_localized\_integers\_padic\_kernel\_rigidity}

\begin{corollary}[有限局部化塔的逐素数增量律]\label{cor:xi-localized-integers-prime-axis-increment-law}
设 \(S\subset T\subset\mathbb P\) 为有限素数集。则有自然短正合列
\[
0\longrightarrow G_S\longrightarrow G_T\longrightarrow
\bigoplus_{p\in T\setminus S}\ZZ(p^\infty)\longrightarrow 0,
\]
其 Pontryagin 对偶为
\[
0\longrightarrow \prod_{p\in T\setminus S}\ZZ_p
\longrightarrow \widehat{G_T}
\longrightarrow \widehat{G_S}
\longrightarrow 0.
\]
特别地，每加入一个新素数 \(p\)，对偶端恰增加一条新的 \(p\)-进轴，而连通商
\(\TT\) 与有限秩圆维 \(\cdim_{\mathrm{fr}}=1\) 保持不变。
\end{corollary}
\begin{proof}
由
\[
0\to \ZZ\to G_T\to \bigoplus_{p\in T}\ZZ(p^\infty)\to 0,
\qquad
0\to \ZZ\to G_S\to \bigoplus_{p\in S}\ZZ(p^\infty)\to 0
\]
及包含 \(G_S\subset G_T\) 可得
\[
G_T/G_S
\cong
(G_T/\ZZ)/(G_S/\ZZ)
\cong
\bigoplus_{p\in T\setminus S}\ZZ(p^\infty),
\]
从而得到第一条短正合列。对偶化并使用
\(
\widehat{\ZZ(p^\infty)}\cong \ZZ_p
\)
即得第二条。

最后，由定理 \ref{thm:xi-localized-integers-padic-kernel-rigidity}，
\(\widehat{G_S}\) 与 \(\widehat{G_T}\) 的连通分量均同构于 \(\TT\)，且
\(\cdim_{\mathrm{fr}}(G_S)=\cdim_{\mathrm{fr}}(G_T)=1\)。
故新增信息完全进入全不连通核。证毕。
\end{proof}
\leanverified{paper\_xi\_localized\_integers\_prime\_axis\_increment\_law}

\begin{theorem}[连通与有理不变量的统一不可区分性]\label{thm:xi-localized-integers-connected-rational-blindness}
设 \(\mathfrak I\) 为定义在有限素数局部化群族
\[
\{G_S=\ZZ[S^{-1}]:\ S\subset\mathbb P,\ |S|<\infty\}
\]
上的同构不变量，并满足下列任一条件：
\begin{enumerate}
  \item \(\mathfrak I(G)\) 仅依赖于 \(G\otimes_{\ZZ}\QQ\)；
  \item \(\mathfrak I(G)\) 仅依赖于 \((\widehat G)^0\)；
  \item \(\mathfrak I(G)\) 仅依赖于 \(\cdim_{\mathrm{fr}}(G)\)。
\end{enumerate}
则对任意有限 \(S,T\subset\mathbb P\) 都有
\[
\mathfrak I(G_S)=\mathfrak I(G_T).
\]
特别地，任何只读取有理部分、连通部分或有限秩圆维的一维不变量，都不可能对该族给出完备分类。
\end{theorem}
\begin{proof}
对任意有限 \(S\subset\mathbb P\)，都有
\[
G_S\otimes_{\ZZ}\QQ\cong \QQ.
\]
又由定理 \ref{thm:xi-localized-integers-padic-kernel-rigidity}，
\[
(\widehat{G_S})^0\cong \TT,
\qquad
\cdim_{\mathrm{fr}}(G_S)=1.
\]
因此若 \(\mathfrak I\) 只通过上述任一对象因子化，则它在所有 \(G_S\) 上都取同一数值。

另一方面，定理 \ref{thm:xi-localized-integers-padic-kernel-rigidity} 已证明
\[
G_S\cong G_T
\iff
S=T.
\]
故这类不变量都不能区分不同的素数支撑，更不可能给出完备分类。证毕。
\end{proof}
\leanverified{paper\_xi\_localized\_integers\_connected\_rational\_blindness}

\paragraph{纤维群胚代数的中心谱、Fuglede--Kadison 体积、容量截断与 window-\texorpdfstring{$6$}{6} 异常分层}
下列结果把命题 \ref{prop:fold-bin-groupoid-algebra-wedderburn}、定理
\ref{thm:op-algebra-fold-watatani-index-equals-multiplicity-field}、定理
\ref{thm:fold-oracle-capacity-closed-form}、命题
\ref{prop:fold-oracle-capacity-bimeasure-identity}、定理
\ref{thm:fold-bin-gauge-volume-stirling}、定理
\ref{thm:xi-foldbin-gauge-representation-distribution-law-m6-spectrum} 与定理
\ref{thm:xi-window6-center-exact-splitting-bdry-cyclic} 组织为同一条中心算子链：纤维多重度场既是群胚代数中心条件期望的最小中心值指数元，也是规范体积、离散容量与
window-\(6\) 异常签名分层的统一谱变量。

\begin{theorem}[纤维群胚代数的最小中心值指数元与谱顶端]\label{thm:xi-foldbin-minimal-center-valued-index-spectrum-top}
设
\[
A_m:=\CC[\mathcal R_m^{\mathrm{bin}}]
\cong
\bigoplus_{w\in X_m}M_{d_m^{\mathrm{bin}}(w)}(\CC),
\qquad
Z_m:=Z(A_m)=\bigoplus_{w\in X_m}\CC z_w,
\]
其中 \(z_w\) 为对应于第 \(w\)-块的最小中心幂等元。定义逐块归一化迹条件期望
\[
E_m\!\left(\bigoplus_{w\in X_m}a_w\right)
:=
\bigoplus_{w\in X_m}
\frac{\operatorname{Tr}(a_w)}{d_m^{\mathrm{bin}}(w)}I_{d_m^{\mathrm{bin}}(w)},
\]
并记
\[
d_w:=d_m^{\mathrm{bin}}(w),
\qquad
\Lambda_m:=\sum_{w\in X_m}d_w z_w\in Z_m.
\]
再定义
\[
M_m:=\max_{w\in X_m}d_w,
\qquad
K_m:=\#\{w\in X_m:\ d_w=M_m\}.
\]
则有：
\begin{enumerate}
  \item 对任意 \(x\in A_m^+\)，都有
  \[
  E_m(x)\ge \Lambda_m^{-1}x,
  \qquad
  \Lambda_m^{-1}=\sum_{w\in X_m}d_w^{-1}z_w.
  \]
  \item 若 \(\Gamma=\sum_{w\in X_m}\gamma_w z_w\in Z_m\) 满足 \(\gamma_w>0\) 且
  \[
  E_m(x)\ge \Gamma^{-1}x
  \qquad
  \forall x\in A_m^+,
  \]
  则 \(\Gamma\ge \Lambda_m\)。因而 \(\Lambda_m\) 是 \(E_m\) 的最小中心值指数元。
  \item \(\Lambda_m\) 的标量指数恰为
  \[
  \|\Lambda_m\|_\infty=M_m,
  \]
  且中心左乘算子
  \[
  L_{\Lambda_m}:Z_m\to Z_m,
  \qquad
  z\mapsto \Lambda_m z
  \]
  在特征值 \(M_m\) 处的几何重数恰为 \(K_m\)。
  \item 若极限
  \[
  \alpha_*:=\lim_{m\to\infty}\frac{1}{m}\log M_m,
  \qquad
  g_*:=\lim_{m\to\infty}\frac{1}{m}\log K_m
  \]
  存在，则
  \[
  \alpha_*=
  \lim_{m\to\infty}\frac{1}{m}\log\|\Lambda_m\|_\infty,
  \qquad
  g_*=
  \lim_{m\to\infty}\frac{1}{m}\log
  \dim\ker\!\bigl(M_m I_{Z_m}-L_{\Lambda_m}\bigr).
  \]
\end{enumerate}
\end{theorem}
\begin{proof}
对任意 \(x\in A_m^+\) 写成块对角形式
\[
x=\bigoplus_{w\in X_m}x_w,
\qquad
x_w\in M_{d_w}(\CC)^+.
\]
在单个矩阵块上，
\[
E_m(x)_w=\frac{\operatorname{Tr}(x_w)}{d_w}I_{d_w}.
\]
又对任意正矩阵 \(x_w\)，都有基本不等式
\[
x_w\le \operatorname{Tr}(x_w)\,I_{d_w}.
\]
因此
\[
E_m(x)_w\ge \frac{1}{d_w}x_w.
\]
逐块并合即得
\[
E_m(x)\ge \left(\sum_{w\in X_m}d_w^{-1}z_w\right)x=\Lambda_m^{-1}x,
\]
从而得到第 \(1\) 条。

现证最小性。设 \(\Gamma=\sum_w\gamma_w z_w\) 满足 \(\gamma_w>0\) 且
\[
E_m(x)\ge \Gamma^{-1}x
\qquad
\forall x\in A_m^+.
\]
固定 \(w\in X_m\)，取第 \(w\)-块中的秩一投影 \(p_w\)。则
\[
E_m(p_w)=\frac{1}{d_w}I_{d_w}.
\]
若 \(\gamma_w<d_w\)，则在 \(\operatorname{Ran}(p_w)\) 上有
\[
\left(E_m(p_w)-\Gamma^{-1}p_w\right)\big|_{\operatorname{Ran}(p_w)}
=
\left(\frac{1}{d_w}-\frac{1}{\gamma_w}\right)I_{\operatorname{Ran}(p_w)},
\]
该算子严格为负，矛盾。故 \(\gamma_w\ge d_w\) 对所有 \(w\) 都成立，即
\(\Gamma\ge \Lambda_m\)。第 \(2\) 条得证。

第 \(3\) 条是 \(\Lambda_m\) 在中心基 \(\{z_w\}_{w\in X_m}\) 上按坐标对角的直接结果：其特征值恰为 \(\{d_w\}_{w\in X_m}\)，故算子范数等于最大坐标 \(M_m\)，而顶端特征空间由满足 \(d_w=M_m\) 的坐标轴张成，维数正是 \(K_m\)。

第 \(4\) 条只是把第 \(3\) 条的两条精确恒等式写成指数增长率。证毕。
\end{proof}

\begin{theorem}[中心值指数元的 Fuglede--Kadison 对数体积与规范体积常数差]\label{thm:xi-foldbin-fkdet-logvolume-kappa}
沿用定理 \ref{thm:xi-foldbin-minimal-center-valued-index-spectrum-top} 的记号，并记
\[
\mathcal G_m^{\mathrm{bin}}:=\Aut(\Fold_m^{\mathrm{bin}}),
\qquad
g_m:=2^{-m}\log|\mathcal G_m^{\mathrm{bin}}|,
\]
\[
\kappa_m:=2^{-m}\sum_{w\in X_m}d_w\log d_w.
\]
在 \(A_m\) 上定义概率迹
\[
\tau_m\!\left(\bigoplus_{w\in X_m}a_w\right)
:=
2^{-m}\sum_{w\in X_m}\operatorname{Tr}(a_w).
\]
则
\[
\log\Delta_{\tau_m}(\Lambda_m):=\tau_m(\log\Lambda_m)=\kappa_m.
\]
因此
\[
|\mathcal G_m^{\mathrm{bin}}|^{1/2^m}
=
e^{-1+O\!\left(m\left(\frac{\varphi}{2}\right)^m\right)}
\Delta_{\tau_m}(\Lambda_m),
\]
并且
\[
\lim_{m\to\infty}
\frac{|\mathcal G_m^{\mathrm{bin}}|^{1/2^m}}{\Delta_{\tau_m}(\Lambda_m)}
=
e^{-1}.
\]
\end{theorem}
\begin{proof}
由 \(\sum_{w\in X_m}d_w=2^m\)，有
\[
\tau_m(1_{A_m})=2^{-m}\sum_{w\in X_m}d_w=1,
\]
故 \(\tau_m\) 确为概率迹。又因 \(\Lambda_m\) 在第 \(w\)-块上等于 \(d_w I_{d_w}\)，故其函数演算满足
\[
\log\Lambda_m=\sum_{w\in X_m}(\log d_w)\,z_w.
\]
于是
\[
\tau_m(\log\Lambda_m)
=
2^{-m}\sum_{w\in X_m}\operatorname{Tr}\!\bigl((\log d_w)I_{d_w}\bigr)
=
2^{-m}\sum_{w\in X_m}d_w\log d_w
=
\kappa_m,
\]
从而
\[
\log\Delta_{\tau_m}(\Lambda_m)=\kappa_m.
\]

另一方面，定理 \ref{thm:fold-bin-gauge-volume-stirling} 给出
\[
g_m=\kappa_m-1+O\!\left(m\left(\frac{\varphi}{2}\right)^m\right).
\]
对两边取指数，并用
\[
e^{g_m}=|\mathcal G_m^{\mathrm{bin}}|^{1/2^m},
\qquad
e^{\kappa_m}=\Delta_{\tau_m}(\Lambda_m),
\]
即得
\[
|\mathcal G_m^{\mathrm{bin}}|^{1/2^m}
=
e^{-1+O\!\left(m\left(\frac{\varphi}{2}\right)^m\right)}
\Delta_{\tau_m}(\Lambda_m).
\]
令 \(m\to\infty\) 即得极限比值 \(e^{-1}\)。证毕。
\end{proof}

\begin{theorem}[预言机容量是中心值指数元的谱截断迹]\label{thm:xi-foldbin-oracle-capacity-central-spectral-truncation}
沿用定理 \ref{thm:xi-foldbin-minimal-center-valued-index-spectrum-top} 的记号。对整数 \(B\ge 0\)，定义中心谱截断
\[
\Lambda_m\wedge 2^B
:=
\sum_{w\in X_m}\min(d_w,2^B)\,z_w\in Z_m,
\]
并记
\[
1_{\Lambda_m\le 2^B}:=\sum_{d_w\le 2^B}z_w,
\qquad
1_{\Lambda_m>2^B}:=\sum_{d_w>2^B}z_w.
\]
令 \(C_m(B)\) 为定理 \ref{thm:fold-oracle-capacity-closed-form} 的离散预言机容量。再定义中心上的三个迹
\[
\operatorname{Tr}_{Z_m}\!\left(\sum_{w\in X_m}c_w z_w\right):=\sum_{w\in X_m}c_w,
\qquad
u_m^Z:=|X_m|^{-1}\operatorname{Tr}_{Z_m},
\]
\[
\tau_m^Z\!\left(\sum_{w\in X_m}c_w z_w\right)
:=
2^{-m}\sum_{w\in X_m}d_w c_w.
\]
则有严格恒等式
\[
C_m(B)=\operatorname{Tr}_{Z_m}\!\bigl(\Lambda_m\wedge 2^B\bigr),
\]
以及
\[
\frac{C_m(B)}{2^m}
=
\tau_m^Z\!\bigl(1_{\Lambda_m\le 2^B}\bigr)
+
2^{B-m}|X_m|\,u_m^Z\!\bigl(1_{\Lambda_m>2^B}\bigr).
\]
特别地，
\[
C_m(B)\le 2^B\dim Z_m=2^B|X_m|.
\]
\end{theorem}
\begin{proof}
由于 \(\Lambda_m\) 在中心基 \(\{z_w\}\) 上对角，且其谱值恰为 \(\{d_w\}_{w\in X_m}\)，函数演算直接给出
\[
\Lambda_m\wedge 2^B
=
\sum_{w\in X_m}\min(d_w,2^B)\,z_w.
\]
对其取计数迹 \(\operatorname{Tr}_{Z_m}\) 得
\[
\operatorname{Tr}_{Z_m}\!\bigl(\Lambda_m\wedge 2^B\bigr)
=
\sum_{w\in X_m}\min(d_w,2^B).
\]
而右端正是定理 \ref{thm:fold-oracle-capacity-closed-form} 的容量闭式，故得到第一条恒等式。

再将上式按 \(d_w\le 2^B\) 与 \(d_w>2^B\) 分拆并除以 \(2^m\)，得
\[
\frac{C_m(B)}{2^m}
=
2^{-m}\sum_{d_w\le 2^B}d_w
+
2^{-m}\sum_{d_w>2^B}2^B.
\]
第一项恰为
\[
\tau_m^Z\!\bigl(1_{\Lambda_m\le 2^B}\bigr),
\]
第二项则等于
\[
2^{B-m}\#\{w\in X_m:\ d_w>2^B\}
=
2^{B-m}|X_m|\,u_m^Z\!\bigl(1_{\Lambda_m>2^B}\bigr),
\]
于是得到第二条公式。这正是命题 \ref{prop:fold-oracle-capacity-bimeasure-identity} 在中心谱上的算子化改写。

最后，由点态不等式
\[
\Lambda_m\wedge 2^B\le 2^B\,1_{Z_m}
\]
取 \(\operatorname{Tr}_{Z_m}\) 即得
\[
C_m(B)\le 2^B\operatorname{Tr}_{Z_m}(1_{Z_m})=2^B|X_m|.
\]
证毕。
\end{proof}

\begin{theorem}[window-\texorpdfstring{$6$}{6} 异常签名空间的规范三层分解]\label{thm:xi-window6-anomaly-signature-three-layer-splitting}
沿用定理 \ref{thm:xi-window6-center-exact-splitting-bdry-cyclic}、定理
\ref{thm:xi-foldbin-gauge-representation-distribution-law-m6-spectrum} 与定义
\ref{def:fold-bin-gauge-sign} 的记号。记
\[
C_{\partial}:=(\ZZ/2\ZZ)^{X_6^{\mathrm{bdry}}}\cong(\ZZ/2\ZZ)^3,
\qquad
C_{\mathrm{int}}:=(\ZZ/2\ZZ)^{X_{6,\mathrm{cyc},2}}\cong(\ZZ/2\ZZ)^5,
\]
并定义
\[
C_{\mathrm{bulk}}
:=
\bigl(S_3^{\,4}\times S_4^{\,9}\bigr)^{\mathrm{ab}}
\cong
(\ZZ/2\ZZ)^4\oplus(\ZZ/2\ZZ)^9
\cong
(\ZZ/2\ZZ)^{13}.
\]
则
\[
Z\!\bigl(\mathcal G_6^{\mathrm{bin}}\bigr)
\cong
C_{\partial}\oplus C_{\mathrm{int}}
\cong
(\ZZ/2\ZZ)^3\oplus(\ZZ/2\ZZ)^5
\cong
(\ZZ/2\ZZ)^8.
\]
并且在符号映射 \(\mathrm{sgn}_6\) 给出的 \(21\) 比特坐标下，
\[
\bigl(\mathcal G_6^{\mathrm{bin}}\bigr)^{\mathrm{ab}}
\cong
C_{\partial}\oplus C_{\mathrm{int}}\oplus C_{\mathrm{bulk}}
\cong
(\ZZ/2\ZZ)^3\oplus(\ZZ/2\ZZ)^5\oplus(\ZZ/2\ZZ)^{13}.
\]
因此 window-\(6\) 的异常签名空间具有规范三层分解：\(3\) 维 boundary 中心层、\(5\) 维非 boundary 的
\(d=2\) 中心层，以及 \(13\) 维非中心 bulk 符号层。
\end{theorem}
\begin{proof}
第一条分解正是定理 \ref{thm:xi-window6-center-exact-splitting-bdry-cyclic} 的结论：全部中心恰由八条
\(d=2\) 纤维给出，并按 boundary 三条与其余五条非 boundary 的 \(d=2\) 纤维裂解为
\[
Z\!\bigl(\mathcal G_6^{\mathrm{bin}}\bigr)
\cong
C_{\partial}\oplus C_{\mathrm{int}}
\cong
(\ZZ/2\ZZ)^3\oplus(\ZZ/2\ZZ)^5.
\]

另一方面，由定理 \ref{thm:xi-foldbin-gauge-representation-distribution-law-m6-spectrum}，
\[
\mathcal G_6^{\mathrm{bin}}
\cong
S_2^{\,8}\times S_3^{\,4}\times S_4^{\,9},
\qquad
\bigl(\mathcal G_6^{\mathrm{bin}}\bigr)^{\mathrm{ab}}
\cong
(\ZZ/2\ZZ)^8\oplus(\ZZ/2\ZZ)^4\oplus(\ZZ/2\ZZ)^9.
\]
其中第一个 \((\ZZ/2\ZZ)^8\) 正对应全部 \(d=2\) 纤维的符号坐标；由上一步，它又规范地分裂为
\[
(\ZZ/2\ZZ)^8\cong C_{\partial}\oplus C_{\mathrm{int}}.
\]
其余 \((\ZZ/2\ZZ)^4\oplus(\ZZ/2\ZZ)^9\) 则来自 \(S_3^{\,4}\times S_4^{\,9}\) 的符号荷，构成
\(C_{\mathrm{bulk}}\)。

最后，由于 \(S_3\) 与 \(S_4\) 的中心都平凡，这 \(13\) 个比特不可能来自中心层；它们只在交换子商与符号映射口径下可见。故得到所述三层分解。证毕。
\end{proof}

\begin{theorem}[window-\texorpdfstring{$6$}{6} boundary 奇偶子格在交换化中的规范直和实现]\label{thm:xi-window6-boundary-parity-canonical-direct-summand}
沿用定理 \ref{thm:xi-window6-anomaly-signature-three-layer-splitting} 与定理
\ref{thm:xi-window6-center-exact-splitting-bdry-cyclic} 的记号，并记
\[
G_6^{\mathrm{bin}}:=\mathcal G_6^{\mathrm{bin}},
\qquad
Z_{\partial}:=\langle \tau_w:\ w\in X_6^{\mathrm{bdry}}\rangle
\cong
(\ZZ/2\ZZ)^3
\subset Z(G_6^{\mathrm{bin}}).
\]
若
\[
\operatorname{ab}_6:G_6^{\mathrm{bin}}\longrightarrow \bigl(G_6^{\mathrm{bin}}\bigr)^{\mathrm{ab}}
\]
为交换化映射，则 \(\operatorname{ab}_6|_{Z_{\partial}}\) 为单同态，其像恰为
\[
C_{\partial}=(\ZZ/2\ZZ)^{X_6^{\mathrm{bdry}}}.
\]
从而有规范直和分解
\[
\boxed{\;
\bigl(G_6^{\mathrm{bin}}\bigr)^{\mathrm{ab}}
\cong
C_{\partial}\oplus (\ZZ/2\ZZ)^{X_6\setminus X_6^{\mathrm{bdry}}}
\cong
(\ZZ/2\ZZ)^3\oplus(\ZZ/2\ZZ)^{18}.\;}
\]
其中第一因子由中心子群 \(Z_{\partial}\) 在交换化中的像规范实现。
\end{theorem}
\begin{proof}
定理 \ref{thm:xi-window6-center-exact-splitting-bdry-cyclic} 已将 boundary 三条 \(d=2\) 纤维所生成的中心层识别为
\[
C_{\partial}\cong (\ZZ/2\ZZ)^{X_6^{\mathrm{bdry}}}.
\]
因此三条中心 involution \(\tau_w\) 在交换化中的像彼此独立，\(\operatorname{ab}_6|_{Z_{\partial}}\) 为单同态，且其像正是 \(C_{\partial}\)。

另一方面，定理 \ref{thm:xi-window6-anomaly-signature-three-layer-splitting} 给出
\[
\bigl(G_6^{\mathrm{bin}}\bigr)^{\mathrm{ab}}
\cong
C_{\partial}\oplus C_{\mathrm{int}}\oplus C_{\mathrm{bulk}}
\cong
(\ZZ/2\ZZ)^3\oplus(\ZZ/2\ZZ)^5\oplus(\ZZ/2\ZZ)^{13}.
\]
将后两项合并即得
\[
C_{\mathrm{int}}\oplus C_{\mathrm{bulk}}
\cong
(\ZZ/2\ZZ)^{18}
\cong
(\ZZ/2\ZZ)^{X_6\setminus X_6^{\mathrm{bdry}}}.
\]
于是得到所示规范直和分解，其中第一因子恰由 \(Z_{\partial}\) 的像实现。证毕。
\end{proof}

\begin{corollary}[window-\texorpdfstring{$6$}{6} 的 boundary parity 是 \texorpdfstring{$21$}{21} 维异常签名空间的唯一三维最小忠实商]\label{cor:xi-window6-boundary-parity-unique-minimal-faithful-quotient}
沿用定理 \ref{thm:xi-window6-boundary-parity-canonical-direct-summand} 的记号，并记
\[
\mathcal A_6:=\bigl(G_6^{\mathrm{bin}}\bigr)^{\mathrm{ab}}
\cong
(\ZZ/2\ZZ)^{21},
\qquad
P_6:=C_{\partial}\cong(\ZZ/2\ZZ)^3.
\]
设 \(B\) 为任意初等 \(2\)-群，\(\phi:\mathcal A_6\to B\) 为群同态，并假设 \(\phi|_{P_6}\) 为单同态。则有：
\begin{enumerate}
  \item \(\dim_{\FF_2}B\ge 3\)；
  \item 当且仅当 \(\dim_{\FF_2}B=3\) 时，\(\phi\) 与沿某个 \(18\) 维补空间的投影等价。
\end{enumerate}
等价地，在保持全部 boundary parity 信息的前提下，\(\mathcal A_6\) 的最小忠实商恰为
\[
\mathcal A_6/K\cong P_6\cong (\ZZ/2\ZZ)^3,
\]
其中 \(K\subset \mathcal A_6\) 为某个 \(18\) 维补空间；因而其核维数必且仅必为 \(18\)。
\end{corollary}
\begin{proof}
由于 \(P_6\cong(\ZZ/2\ZZ)^3\) 且 \(\phi|_{P_6}\) 单射，像 \(\phi(P_6)\) 是 \(B\) 的一个三维 \(\FF_2\)-子空间，故必有 \(\dim_{\FF_2}B\ge 3\)。

若 \(\dim_{\FF_2}B=3\)，则 \(\phi|_{P_6}:P_6\to B\) 为同构。此时 \(\phi\) 的像即为 \(B\)，故 \(\operatorname{rank}\phi=3\)，从而
\[
\dim_{\FF_2}\ker\phi=21-3=18.
\]
又因 \(\phi|_{P_6}\) 单射，故
\[
\ker\phi\cap P_6=\{0\}.
\]
因此 \(\ker\phi\) 给出 \(P_6\) 的一个 \(18\) 维补空间，并有
\[
\mathcal A_6=P_6\oplus\ker\phi.
\]
于是 \(\phi\) 与沿 \(\ker\phi\) 的投影
\[
\mathcal A_6\twoheadrightarrow \mathcal A_6/\ker\phi\cong P_6
\]
在识别 \(P_6\cong B\) 后等价。

反之，若 \(K\subset \mathcal A_6\) 为任一与 \(P_6\) 互补的 \(18\) 维子空间，则自然投影
\[
\mathcal A_6\twoheadrightarrow \mathcal A_6/K\cong P_6\cong(\ZZ/2\ZZ)^3
\]
在 \(P_6\) 上显然保持忠实，且目标维数恰为 \(3\)。故三维正是最小可达维数。证毕。
\end{proof}

由此，纤维多重度场在群胚代数层面实现如下统一链：
\[
\CC[\mathcal R_m^{\mathrm{bin}}]
\Longrightarrow
\Lambda_m=\sum_{w\in X_m}d_w z_w
\Longrightarrow
\begin{cases}
\text{最小中心值指数元},\\
\log\Delta_{\tau_m}(\Lambda_m)=\kappa_m,\\
C_m(B)=\operatorname{Tr}_{Z_m}(\Lambda_m\wedge 2^B).
\end{cases}
\]
而在 \(m=6\) 时，\(\Lambda_6\) 的 \(d=2\) 谱层进一步对应于
\[
Z\!\bigl(\mathcal G_6^{\mathrm{bin}}\bigr)
\cong
(\ZZ/2\ZZ)^3_{\partial}\oplus(\ZZ/2\ZZ)^5_{\mathrm{int}}
\subset
\bigl(\mathcal G_6^{\mathrm{bin}}\bigr)^{\mathrm{ab}}
\cong
(\ZZ/2\ZZ)^3\oplus(\ZZ/2\ZZ)^5\oplus(\ZZ/2\ZZ)^{13}.
\]

\paragraph{审计偶数窗上的最小退化扇区、群胚维数、规范群熵与黄金密度分裂}
在定理 \ref{thm:xi-foldbin-minimal-center-valued-index-spectrum-top} 已把纤维多重度场压缩为中心谱变量之后，审计偶数窗
\[
m\in\{6,8,10,12\}
\]
上由最小退化扇区驱动的碰撞、Abel 化、群胚中心与规范群熵还可进一步收束为一组完全对象层结论。沿用定义 \ref{def:fold-bin-degeneracy-sectors}、定理 \ref{thm:fold-bin-min-degeneracy-fib-audited}、推论 \ref{cor:fold-bin-minsector-complement-maxfiber-audited}、定理 \ref{thm:xi-foldbin-collision-probability-normalized-groupoid-dimension}、定理 \ref{thm:xi-foldbin-parity-charge-split-exact-minrank}、命题 \ref{prop:fold-bin-groupoid-algebra-wedderburn}、注记 \ref{rem:fold-bin-gauge-entropy-scale} 与推论 \ref{cor:fold-bin-minsector-density-phi-minus2} 的记号，记
\[
d_*:=d_{\min}(m)=F_{m/2},
\qquad
\mathcal S_{m,*}:=\mathcal S_{m,d_*},
\qquad
r_m:=|X_m\setminus \mathcal S_{m,*}|=D_{2m-2}=F_{m+1}.
\]
为避免与定理 \ref{thm:xi-foldbin-minimal-center-valued-index-spectrum-top} 中的谱顶端记号 \(M_m\) 混淆，记补扇区总质量
\[
T_m:=2^m-F_m d_*,
\qquad
T_m=q_m r_m+s_m,\qquad 0\le s_m<r_m.
\]
并置
\[
A_m:=\CC[\mathcal R_m^{\mathrm{bin}}],
\qquad
A_{m,\min}:=\bigoplus_{w\in\mathcal S_{m,*}}M_{d_*}(\CC),
\qquad
A_{m,\mathrm{out}}:=\bigoplus_{w\in X_m\setminus\mathcal S_{m,*}}
M_{d_m^{\mathrm{bin}}(w)}(\CC).
\]

\begin{theorem}[审计偶数窗连续容量曲线的首折点具有精确 Fibonacci 斜率跳跃]\label{thm:xi-audited-even-first-capacity-kink-fibonacci-jump}
沿用定理 \ref{thm:xi-foldbin-continuous-capacity-right-derivative-abelian-center} 的连续容量记号，定义
\[
\mathcal C_m^{\mathrm{cont}}(t):=\sum_{w\in X_m}\min\bigl(d_m^{\mathrm{bin}}(w),t\bigr),
\qquad t\ge 0.
\]
则对每个审计偶数窗 \(m\in\{6,8,10,12\}\)，首折点 \(t=d_*=F_{m/2}\) 附近成立精确分段公式
\[
\boxed{\;
\mathcal C_m^{\mathrm{cont}}(t)=F_{m+2}\,t
\qquad
(0\le t\le d_*),\;}
\]
\[
\boxed{\;
\mathcal C_m^{\mathrm{cont}}(t)=F_m\,d_*+F_{m+1}\,t
\qquad
(d_*\le t\le d_*+1).\;}
\]
从而在首折点处的左右导数满足
\[
\boxed{\;
\partial_t^{-}\mathcal C_m^{\mathrm{cont}}(d_*)=F_{m+2},
\qquad
\partial_t^{+}\mathcal C_m^{\mathrm{cont}}(d_*)=F_{m+1},\;}
\]
并且首个斜率跳跃精确等于
\[
\boxed{\;
\partial_t^{-}\mathcal C_m^{\mathrm{cont}}(d_*)-\partial_t^{+}\mathcal C_m^{\mathrm{cont}}(d_*)
=
F_m.\;}
\]
\end{theorem}
\begin{proof}
由定理 \ref{thm:fold-bin-min-degeneracy-fib-audited}，
\[
d_m^{\mathrm{bin}}(w)\ge d_*
\qquad (w\in X_m),
\]
且
\[
|\mathcal S_{m,*}|=F_m.
\]
因此当 \(0\le t\le d_*\) 时，每一条纤维都尚未饱和，从而
\[
\mathcal C_m^{\mathrm{cont}}(t)
=
\sum_{w\in X_m} t
=
|X_m|\,t
=
F_{m+2}\,t.
\]

再由推论 \ref{cor:fold-bin-minsector-complement-maxfiber-audited}，
\[
|X_m\setminus \mathcal S_{m,*}|=r_m=F_{m+1}.
\]
而对 \(w\in\mathcal S_{m,*}\) 有
\[
d_m^{\mathrm{bin}}(w)=d_*,
\]
对 \(w\in X_m\setminus \mathcal S_{m,*}\) 则由最小退化度定义知
\[
d_m^{\mathrm{bin}}(w)\ge d_*+1.
\]
故当 \(d_*\le t\le d_*+1\) 时，最小退化扇区中的纤维已经饱和，而补扇区中的纤维仍处于线性段，于是
\[
\mathcal C_m^{\mathrm{cont}}(t)
=
\sum_{w\in\mathcal S_{m,*}} d_*
+
\sum_{w\in X_m\setminus \mathcal S_{m,*}} t
=
F_m\,d_*+F_{m+1}\,t.
\]
这就给出了所示两段精确公式。

最后，两段都是线性函数，故首折点处的左右导数分别等于对应斜率
\[
F_{m+2},
\qquad
F_{m+1}.
\]
它们之差为
\[
F_{m+2}-F_{m+1}=F_m,
\]
即首个斜率跳跃恰等于最小退化扇区的 Fibonacci 基数。证毕。
\end{proof}

\begin{theorem}[群胚代数维数与碰撞矩的严格等同及最小扇区平衡下界]\label{thm:xi-audited-even-groupoid-dimension-collision-balanced-lower-bound}
\leanverified{paper\_xi\_audited\_even\_groupoid\_dimension\_collision\_balanced\_lower\_bound}
有严格恒等式
\[
\boxed{\;
\dim_{\CC}A_m
=
\sum_{w\in X_m}\bigl(d_m^{\mathrm{bin}}(w)\bigr)^2
=
4^m\,\mathsf{Col}^{\mathrm{bin}}_m.\;}
\]
进一步，
\[
\boxed{\;
\dim_{\CC}A_m
\ge
F_m d_*^2+(r_m-s_m)q_m^2+s_m(q_m+1)^2.\;}
\]
并且等号成立当且仅当补扇区中的全部纤维大小恰只取 \(q_m\) 与 \(q_m+1\) 两值，其中恰有 \(s_m\) 条取 \(q_m+1\)。
\end{theorem}
\begin{proof}
由定理 \ref{thm:xi-foldbin-collision-probability-normalized-groupoid-dimension}，
\[
\dim_{\CC}A_m
=
\sum_{w\in X_m}\bigl(d_m^{\mathrm{bin}}(w)\bigr)^2
=
4^m\,\mathsf{Col}^{\mathrm{bin}}_m,
\]
第一条结论成立。

再证下界。由定理 \ref{thm:fold-bin-min-degeneracy-fib-audited}，
\[
|\mathcal S_{m,*}|=F_m,
\qquad
d_m^{\mathrm{bin}}(w)=d_*
\quad (w\in\mathcal S_{m,*}),
\]
而由推论 \ref{cor:fold-bin-minsector-complement-maxfiber-audited}，
\[
|X_m\setminus \mathcal S_{m,*}|=r_m.
\]
将补扇区上的纤维大小重记为
\[
e_1,\dots,e_{r_m},
\]
则
\[
e_i\ge d_*+1,
\qquad
\sum_{i=1}^{r_m}e_i=T_m,
\]
并且
\[
\dim_{\CC}A_m
=
F_m d_*^2+\sum_{i=1}^{r_m}e_i^2.
\]
因此只需在固定长度 \(r_m\) 与固定总和 \(T_m\) 的约束下最小化
\(
\sum_i e_i^2
\)。

若某一对 \(a,b\) 满足 \(a\ge b+2\)，则
\[
a^2+b^2-(a-1)^2-(b+1)^2
=
2(a-b)-2
>0.
\]
故在保持总和不变的前提下，用 \((a-1,b+1)\) 代替 \((a,b)\) 会严格降低平方和。不断迭代可知，极小值只能在所有坐标彼此至多相差 \(1\) 时取得。把
\[
T_m=q_m r_m+s_m,\qquad 0\le s_m<r_m
\]
代入，唯一可能的极小配置即为
\[
\underbrace{q_m,\dots,q_m}_{r_m-s_m\text{ 次}},
\underbrace{q_m+1,\dots,q_m+1}_{s_m\text{ 次}}.
\]
于是
\[
\sum_{i=1}^{r_m}e_i^2
\ge
(r_m-s_m)q_m^2+s_m(q_m+1)^2,
\]
从而得到所示下界。

最后说明等号条件。对四个审计偶数窗直接代入可得
\[
q_m\in\{3,5,8,12\},
\qquad
d_*\in\{2,3,5,8\},
\]
故总有 \(q_m\ge d_*+1\)，从而上述平衡配置确为补扇区约束 \(e_i\ge d_*+1\) 下的可行点。于是等号成立当且仅当补扇区中的全部纤维大小恰为 \(q_m\) 与 \(q_m+1\) 两值，且其中有 \(s_m\) 条取 \(q_m+1\)。证毕。
\end{proof}

\begin{theorem}[审计偶数窗上的 Abel 化与群胚代数中心的 Fibonacci 双尺度分裂]\label{thm:xi-audited-even-abel-center-fibonacci-splitting}
\leanverified{paper\_xi\_audited\_even\_abel\_center\_fibonacci\_splitting}
有规范直和分裂
\[
\boxed{\;
\bigl(G_m^{\mathrm{bin}}\bigr)^{\mathrm{ab}}
\cong
(\ZZ/2\ZZ)^{\mathcal S_{m,*}}
\oplus
(\ZZ/2\ZZ)^{X_m\setminus\mathcal S_{m,*}}
\cong
(\ZZ/2\ZZ)^{F_m}\oplus(\ZZ/2\ZZ)^{D_{2m-2}}.\;}
\]
并且
\[
\boxed{\;
Z(A_m)
\cong
Z(A_{m,\min})\oplus Z(A_{m,\mathrm{out}})
\cong
\CC^{F_m}\oplus\CC^{D_{2m-2}}.\;}
\]
其中第一直和因子恰由最小退化扇区 \(\mathcal S_{m,*}\) 提供，第二直和因子恰由其补集提供。
\end{theorem}
\begin{proof}
由定理 \ref{thm:xi-foldbin-parity-charge-split-exact-minrank}，
\[
\bigl(G_m^{\mathrm{bin}}\bigr)^{\mathrm{ab}}
\cong
(\ZZ/2\ZZ)^{N_m},
\qquad
N_m:=\#\{w\in X_m:\ d_m^{\mathrm{bin}}(w)\ge 2\}.
\]
而在审计偶数窗上，定理 \ref{thm:fold-bin-min-degeneracy-fib-audited} 给出
\[
d_{\min}(m)=d_*\ge 2,
\]
故每条纤维都满足 \(d_m^{\mathrm{bin}}(w)\ge 2\)，从而
\[
N_m=|X_m|.
\]
于是
\[
\bigl(G_m^{\mathrm{bin}}\bigr)^{\mathrm{ab}}
\cong
(\ZZ/2\ZZ)^{X_m}.
\]
再由集合分拆
\[
X_m=\mathcal S_{m,*}\sqcup(X_m\setminus\mathcal S_{m,*}),
\]
得到
\[
(\ZZ/2\ZZ)^{X_m}
\cong
(\ZZ/2\ZZ)^{\mathcal S_{m,*}}
\oplus
(\ZZ/2\ZZ)^{X_m\setminus\mathcal S_{m,*}}.
\]
最后由
\[
|\mathcal S_{m,*}|=F_m,
\qquad
|X_m\setminus\mathcal S_{m,*}|=D_{2m-2},
\]
便得第一条分裂。

对中心，由命题 \ref{prop:fold-bin-groupoid-algebra-wedderburn}，
\[
A_m
\cong
\bigoplus_{w\in X_m}M_{d_m^{\mathrm{bin}}(w)}(\CC)
=
A_{m,\min}\oplus A_{m,\mathrm{out}}.
\]
有限维半单代数的中心对直和可加，因此
\[
Z(A_m)\cong Z(A_{m,\min})\oplus Z(A_{m,\mathrm{out}}).
\]
又每个矩阵块的中心均为 \(\CC\)，故
\[
Z(A_{m,\min})\cong \CC^{\mathcal S_{m,*}},
\qquad
Z(A_{m,\mathrm{out}})\cong \CC^{X_m\setminus\mathcal S_{m,*}}.
\]
取维数并代入上述两组基数，即得
\[
Z(A_m)\cong \CC^{F_m}\oplus \CC^{D_{2m-2}}.
\]
证毕。
\end{proof}

\begin{theorem}[规范群熵的最优平衡下界与补扇区熵缺口]\label{thm:xi-audited-even-gauge-entropy-balanced-lower-bound}
\leanverified{paper\_xi\_audited\_even\_gauge\_entropy\_balanced\_lower\_bound}
沿用注记 \ref{rem:fold-bin-gauge-entropy-scale} 的 bits 记号
\[
H_m^{\mathrm{gauge}}
:=
\log_2|G_m^{\mathrm{bin}}|
=
\sum_{w\in X_m}\log_2\bigl(d_m^{\mathrm{bin}}(w)!\bigr).
\]
则有最优下界
\[
\boxed{\;
H_m^{\mathrm{gauge}}
\ge
F_m\log_2(d_*!)
+
(r_m-s_m)\log_2(q_m!)
+
s_m\log_2((q_m+1)!).\;}
\]
并且等号成立当且仅当补扇区中的全部纤维大小恰只取 \(q_m\) 与 \(q_m+1\) 两值。

等价地，若定义补扇区熵缺口
\[
\Delta H_m:=H_m^{\mathrm{gauge}}-F_m\log_2(d_*!),
\]
则
\[
\boxed{\;
\Delta H_m
\ge
(r_m-s_m)\log_2(q_m!)
+
s_m\log_2((q_m+1)!).\;}
\]
\end{theorem}
\begin{proof}
由定义，最小扇区上的每条纤维大小都等于 \(d_*\)，其条数为 \(F_m\)。把补扇区上的纤维大小仍记为
\[
e_1,\dots,e_{r_m},
\qquad
\sum_{i=1}^{r_m}e_i=T_m.
\]
则
\[
H_m^{\mathrm{gauge}}
=
F_m\log_2(d_*!)
+
\sum_{i=1}^{r_m}\log_2(e_i!).
\]

定义
\[
f(n):=\log_2(n!)
\qquad (n\in\ZZ_{\ge 1}).
\]
则
\[
f(n+1)-f(n)=\log_2(n+1),
\]
故差分严格递增，亦即 \(f\) 在正整数上严格离散凸。若某一对 \(a,b\) 满足 \(a\ge b+2\)，则
\[
f(a)+f(b)-f(a-1)-f(b+1)
=
\log_2 a-\log_2(b+1)
>0.
\]
因此在固定总和 \(T_m\) 与固定项数 \(r_m\) 的约束下，\(\sum_i f(e_i)\) 的极小值同样由“尽可能均匀”的配置取得，也就是全部 \(e_i\) 只取
\[
q_m,\qquad q_m+1
\]
两值，其中恰有 \(s_m\) 个取 \(q_m+1\)。由定理 \ref{thm:xi-audited-even-groupoid-dimension-collision-balanced-lower-bound} 证明末尾的可行性说明，该配置在四个审计偶数窗上确为可行。故
\[
\sum_{i=1}^{r_m}\log_2(e_i!)
\ge
(r_m-s_m)\log_2(q_m!)
+
s_m\log_2((q_m+1)!),
\]
代回即得第一条下界。

若取等，则上述离散平衡过程中每一步都不能继续降低函数值，于是全部 \(e_i\) 必彼此至多相差 \(1\)，亦即只可能取 \(q_m\) 与 \(q_m+1\) 两值；反之若补扇区纤维正好取这两值，则显然达到下界。第二条关于 \(\Delta H_m\) 的公式只是把最小扇区项移至左端。证毕。
\end{proof}

\begin{theorem}[子覆盖密度延拓下的 Abel 秩密度与中心密度黄金分裂]\label{thm:xi-audited-even-minsector-algebraic-density-golden-splitting}
\leanverified{paper\_xi\_audited\_even\_minsector\_algebraic\_density\_golden\_splitting}
若“子覆盖同构”恒等式
\[
|\mathcal S_{m,*}|=|X_{m-2}|
\]
沿某个趋于无穷的偶数序列成立，则在同一序列上有
\[
\boxed{\;
\frac{
\operatorname{rank}_{\FF_2}\bigl((\ZZ/2\ZZ)^{\mathcal S_{m,*}}\bigr)
}{
\operatorname{rank}_{\FF_2}\!\bigl((G_m^{\mathrm{bin}})^{\mathrm{ab}}\bigr)
}
=
\frac{\dim_{\CC}Z(A_{m,\min})}{\dim_{\CC}Z(A_m)}
=
\frac{F_m}{F_{m+2}}
\longrightarrow
\varphi^{-2}.\;}
\]
并且补扇区的对应密度满足
\[
\boxed{\;
\frac{
\operatorname{rank}_{\FF_2}\bigl((\ZZ/2\ZZ)^{X_m\setminus\mathcal S_{m,*}}\bigr)
}{
\operatorname{rank}_{\FF_2}\!\bigl((G_m^{\mathrm{bin}})^{\mathrm{ab}}\bigr)
}
=
\frac{\dim_{\CC}Z(A_{m,\mathrm{out}})}{\dim_{\CC}Z(A_m)}
=
\frac{F_{m+1}}{F_{m+2}}
\longrightarrow
\varphi^{-1}.\;}
\]
\end{theorem}
\begin{proof}
由定理 \ref{thm:xi-audited-even-abel-center-fibonacci-splitting}，
\[
\operatorname{rank}_{\FF_2}\bigl((\ZZ/2\ZZ)^{\mathcal S_{m,*}}\bigr)
=
|\mathcal S_{m,*}|,
\qquad
\operatorname{rank}_{\FF_2}\!\bigl((G_m^{\mathrm{bin}})^{\mathrm{ab}}\bigr)
=
|X_m|,
\]
并且
\[
\dim_{\CC}Z(A_{m,\min})=|\mathcal S_{m,*}|,
\qquad
\dim_{\CC}Z(A_m)=|X_m|.
\]
因此最小扇区的 Abel 秩密度与中心密度完全一致，均等于
\[
\frac{|\mathcal S_{m,*}|}{|X_m|}.
\]
在假设下，推论 \ref{cor:fold-bin-minsector-density-phi-minus2} 已证明
\[
\frac{|\mathcal S_{m,*}|}{|X_m|}
\longrightarrow
\varphi^{-2},
\]
故第一条极限成立。

对补扇区，同理有
\[
\frac{
\operatorname{rank}_{\FF_2}\bigl((\ZZ/2\ZZ)^{X_m\setminus\mathcal S_{m,*}}\bigr)
}{
\operatorname{rank}_{\FF_2}\!\bigl((G_m^{\mathrm{bin}})^{\mathrm{ab}}\bigr)
}
=
\frac{|X_m\setminus\mathcal S_{m,*}|}{|X_m|}
=
\frac{F_{m+1}}{F_{m+2}},
\]
以及
\[
\frac{\dim_{\CC}Z(A_{m,\mathrm{out}})}{\dim_{\CC}Z(A_m)}
=
\frac{|X_m\setminus\mathcal S_{m,*}|}{|X_m|}
=
\frac{F_{m+1}}{F_{m+2}}.
\]
再由
\[
\frac{F_{m+1}}{F_{m+2}}
=
1-\frac{F_m}{F_{m+2}}
\longrightarrow
1-\varphi^{-2}
=
\varphi^{-1},
\]
即得第二条极限。证毕。
\end{proof}

\begin{theorem}[审计偶数窗最小退化扇区的模 \texorpdfstring{$6$}{6} 同伦相图]\label{thm:xi-audited-even-minsector-homotopy-mod6-phase}
\leanverified{paper\_xi\_audited\_even\_minsector\_homotopy\_mod6\_phase}
沿用上述记号。对任意 \(x\in \mathcal S_{m,*}\)，记
\[
K_x:=\mathrm{Ind}^{\Delta}(\Gamma_x).
\]
其中 \(\tau(x)\) 沿用定理 \ref{thm:pom-fiber-independence-complex-classification} 的记号。
则在审计偶数窗 \(m\in\{6,8,10,12\}\) 上成立：
\begin{enumerate}
  \item 若 \(m\equiv 0\pmod 6\)，则对所有 \(x\in \mathcal S_{m,*}\) 都有 \(|K_x|\) 可缩；
  \item 若 \(m\equiv 2,4\pmod 6\)，则对所有 \(x\in \mathcal S_{m,*}\) 都有
  \[
  |K_x|\simeq S^{\tau(x)-1}.
  \]
\end{enumerate}
特别地，
\[
m=6,12\ \Longrightarrow\ \mathcal S_{m,*}\ \text{上的纤维独立集复形全部可缩},
\]
\[
m=8,10\ \Longrightarrow\ \mathcal S_{m,*}\ \text{上的纤维独立集复形全部为非可缩球面}.
\]
\end{theorem}
\begin{proof}
对任意 \(x\in \mathcal S_{m,*}\)，在本段的 bin-fold 记号下有
\[
d_m(x)=d_m^{\mathrm{bin}}(x)=d_*=F_{m/2},
\]
其中最后一步由定理 \ref{thm:fold-bin-min-degeneracy-fib-audited} 给出。
另一方面，Fibonacci 奇偶律满足
\[
F_n\equiv 0\pmod 2
\quad\Longleftrightarrow\quad
3\mid n.
\]
因此
\[
d_m(x)\ \text{为偶}
\quad\Longleftrightarrow\quad
3\mid \frac{m}{2}
\quad\Longleftrightarrow\quad
m\equiv 0\pmod 6,
\]
并且
\[
d_m(x)\ \text{为奇}
\quad\Longleftrightarrow\quad
m\equiv 2,4\pmod 6.
\]
再由定理 \ref{thm:pom-fiber-parity-homotopy-equivalence}，
\[
|K_x|\ \text{可缩}
\quad\Longleftrightarrow\quad
d_m(x)\ \text{为偶},
\]
\[
|K_x|\simeq S^{\tau(x)-1}
\quad\Longleftrightarrow\quad
d_m(x)\ \text{为奇},
\]
结论遂得。证毕。
\end{proof}

由此，在审计偶数窗上，最小退化扇区不再只是一个局域计数对象，而是同时控制以下五层结构：
\[
\mathcal S_{m,*}
\Longrightarrow
\mathcal S_{m,*}\ \text{上的}\ |K_x|\ \text{按}\ m\bmod 6\ \text{切换于可缩相与球面相}
\Longrightarrow
\dim_{\CC}A_m=4^m\mathsf{Col}^{\mathrm{bin}}_m
\Longrightarrow
\bigl(G_m^{\mathrm{bin}}\bigr)^{\mathrm{ab}}\ \text{与}\ Z(A_m)\ \text{的 Fibonacci 双裂解}
\Longrightarrow
H_m^{\mathrm{gauge}}\ \text{的最优平衡下界}
\Longrightarrow
(\varphi^{-2},\varphi^{-1})\ \text{型代数密度分裂}.
\]

\paragraph{最小退化扇区的严格同伦相位切换}

\begin{corollary}[审计偶数窗最小退化扇区经历严格的模 \texorpdfstring{$6$}{6} 同伦相变]\label{cor:xi-time-part55ab-audited-even-minsector-strict-homotopy-transition}
沿用定义 \ref{def:fold-bin-degeneracy-sectors}、定理 \ref{thm:fold-bin-min-degeneracy-fib-audited} 与定理 \ref{thm:xi-audited-even-minsector-homotopy-mod6-phase} 的记号。对审计偶数窗
\[
m\in\{6,8,10,12\},
\]
记
\[
d_{\min}(m)=F_{m/2},
\qquad
\mathcal S_{m,d_{\min}(m)}:=\{x\in X_m:\ d_m(x)=d_{\min}(m)\}.
\]
并且
\[
\bigl|\mathcal S_{m,d_{\min}(m)}\bigr|=F_m=|X_{m-2}|.
\]
对每个 \(x\in \mathcal S_{m,d_{\min}(m)}\)，再记
\[
K_x:=\operatorname{Ind}^{\Delta}(\Gamma_x),
\qquad
\tau(x):=\sum_i\left\lfloor\frac{\ell_i+1}{3}\right\rfloor .
\]
则对每个 \(x\in \mathcal S_{m,d_{\min}(m)}\)，其纤维独立集复形满足
\[
|K_x|\simeq
\begin{cases}
*,& m\equiv 0\pmod 6,\\[1mm]
S^{\tau(x)-1},& m\equiv 2,4\pmod 6.
\end{cases}
\]
特别地，
\[
m=6,12\ \Longrightarrow\ \mathcal S_{m,d_{\min}(m)}\ \text{上的纤维独立集复形全部可缩},
\]
\[
m=8,10\ \Longrightarrow\ \mathcal S_{m,d_{\min}(m)}\ \text{上的纤维独立集复形全部非可缩}.
\]
\end{corollary}
\begin{proof}
由定理 \ref{thm:fold-bin-min-degeneracy-fib-audited}，在所述审计偶数窗上有
\[
d_{\min}(m)=F_{m/2},
\]
并且最小退化扇区正是
\[
\mathcal S_{m,*}=\mathcal S_{m,d_{\min}(m)}.
\]
同一定理还给出
\[
\bigl|\mathcal S_{m,d_{\min}(m)}\bigr|=\bigl|\mathcal S_{m,*}\bigr|=F_m=|X_{m-2}|.
\]
另一方面，定理 \ref{thm:xi-audited-even-minsector-homotopy-mod6-phase} 已证明：对任意
\[
x\in \mathcal S_{m,*},
\]
当 \(m\equiv 0\pmod 6\) 时 \(|K_x|\) 可缩，而当 \(m\equiv 2,4\pmod 6\) 时
\[
|K_x|\simeq S^{\tau(x)-1}.
\]
将 \(\mathcal S_{m,*}\) 以 \(\mathcal S_{m,d_{\min}(m)}\) 重写，即得结论。证毕。
\end{proof}

\paragraph{审计偶数窗首折点之前的 dyadic 线性段}
连续容量曲线在首个折点处的精确分段公式，进一步把 dyadic 预算下的预言机容量与成功率写成完全显式的前阈值线性段。

\begin{corollary}[审计偶数窗上 dyadic 预算的首段完全线性]\label{cor:xi-audited-even-dyadic-first-capacity-linear-regime}
\leanverified{paper\_xi\_audited\_even\_dyadic\_first\_capacity\_linear\_regime}
对每个审计偶数窗 \(m\in\{6,8,10,12\}\)，定义
\[
C_m(B):=\mathcal C_m^{\mathrm{cont}}(2^B),
\qquad
\mathrm{Succ}_m^{\mathrm{bin}}(B):=2^{-m}C_m(B).
\]
若
\[
0\le B\le \log_2 F_{m/2},
\qquad\text{等价地}\qquad
2^B\le d_{\min}(m)=F_{m/2},
\]
则
\[
C_m(B)=2^B\,F_{m+2},
\qquad
\mathrm{Succ}_m^{\mathrm{bin}}(B)=2^{B-m}F_{m+2}.
\]
因此在首个折点之前，每增加 \(1\) bit，预言机成功率恰好加倍；而一旦越过阈值 \(B=\log_2 F_{m/2}\)，连续容量的边际斜率便由定理 \ref{thm:xi-audited-even-first-capacity-kink-fibonacci-jump} 中的 \(F_{m+2}\) 跳降为 \(F_{m+1}\)。
\end{corollary}
\begin{proof}
由定理 \ref{thm:xi-audited-even-first-capacity-kink-fibonacci-jump}，当
\[
0\le t\le d_{\min}(m)=F_{m/2}
\]
时有
\[
\mathcal C_m^{\mathrm{cont}}(t)=F_{m+2}\,t.
\]
把 \(t=2^B\) 代入，便得到
\[
C_m(B)=\mathcal C_m^{\mathrm{cont}}(2^B)=2^B\,F_{m+2}
\]
对全部 \(0\le B\le \log_2 F_{m/2}\) 成立。两边再除以 \(2^m\)，即得
\[
\mathrm{Succ}_m^{\mathrm{bin}}(B)=2^{B-m}F_{m+2}.
\]
若 \(B+1\le \log_2 F_{m/2}\)，则
\[
\mathrm{Succ}_m^{\mathrm{bin}}(B+1)=2\,\mathrm{Succ}_m^{\mathrm{bin}}(B),
\]
故该段确为精确倍增律。最后，越过首折点后的斜率跳降正是定理 \ref{thm:xi-audited-even-first-capacity-kink-fibonacci-jump} 的右导数结论。证毕。
\end{proof}

\paragraph{window-\texorpdfstring{$6$}{6} 微态 Hilbert 空间的可见/隐藏分解与纤维谱几何}
在定理 \ref{thm:xi-foldbin-gauge-representation-distribution-law-m6-spectrum} 的 \(m=6\) 规范群直积分解、定理 \ref{thm:xi-window6-center-exact-splitting-bdry-cyclic} 的 boundary 中心直和因子，以及定理 \ref{thm:xi-window6-cyclic-kernel-bc3-root-datum} 与表 \ref{tab:fold6_b3c3_root_dictionary_b3}、表 \ref{tab:fold6_b3c3_root_dictionary_c3} 的 \(B_3/C_3\) 根字典已经固定之后，还可把 window-\(6\) 的微态层提升为一条严格的有限维表示论--谱几何链：所有根几何只出现于 gauge 固定点商，而隐藏部门完全由各纤维的零和模组成，并在同纤维团图上量子化为三条有限质量通道。

\begin{theorem}[window-\texorpdfstring{$6$}{6} 微态 Hilbert 空间的规范分解]\label{thm:xi-time-part55d-window6-microstate-hilbert-gauge-splitting}
\leanverified{paper\_xi\_time\_part55d\_window6\_microstate\_hilbert\_gauge\_splitting}
沿用定理 \ref{thm:xi-foldbin-gauge-representation-distribution-law-m6-spectrum} 的记号，记
\[
\Omega_6:=\{0,1,\dots,63\},
\qquad
F_w:=\Fold_6^{\mathrm{bin}}{}^{-1}(w),
\qquad
d(w):=|F_w|
\quad (w\in X_6),
\]
并定义
\[
\mathcal H_6:=\ell^2(\Omega_6)
=
\bigoplus_{w\in X_6}\ell^2(F_w).
\]
在
\[
G_6^{\mathrm{bin}}
=
\prod_{w\in X_6}\mathfrak S_{d(w)}
\]
的自然置换表示下，存在规范正交分解
\[
\mathcal H_6=\mathcal H_{\mathrm{vis}}\oplus \mathcal H_{\mathrm{hid}},
\qquad
\mathcal H_{\mathrm{vis}}
:=
\bigoplus_{w\in X_6}\CC \mathbf 1_{F_w}
=
\mathcal H_6^{G_6^{\mathrm{bin}}},
\]
并且
\[
\mathcal H_{\mathrm{hid}}
\cong
\bigoplus_{w:\,d(w)=4}\mathrm{Std}_4^{(w)}
\oplus
\bigoplus_{w:\,d(w)=3}\mathrm{Std}_3^{(w)}
\oplus
\bigoplus_{w:\,d(w)=2}\mathrm{sgn}_2^{(w)}.
\]
因此
\[
\dim \mathcal H_{\mathrm{vis}}=21,
\qquad
\dim \mathcal H_{\mathrm{hid}}=64-21=43,
\]
且
\[
\dim \mathcal H_{\mathrm{hid}}
=
9\cdot 3+4\cdot 2+8\cdot 1
=
27+8+8.
\]
\end{theorem}
\begin{proof}
由于 \(\{F_w:\ w\in X_6\}\) 构成 \(\Omega_6\) 的不交划分，故
\[
\mathcal H_6
=
\bigoplus_{w\in X_6}\ell^2(F_w)
\]
是规范正交直和。对每个 \(w\in X_6\)，都有
\[
\ell^2(F_w)=\CC \mathbf 1_{F_w}\oplus \mathbf 1_{F_w}^{\perp}.
\]
其中 \(\mathfrak S_{d(w)}\) 在 \(\CC \mathbf 1_{F_w}\) 上平凡作用，在 \(\mathbf 1_{F_w}^{\perp}\) 上给出零和超平面表示。于是当 \(d(w)=4\) 时得到三维标准表示 \(\mathrm{Std}_4\)；当 \(d(w)=3\) 时得到二维标准表示 \(\mathrm{Std}_3\)；当 \(d(w)=2\) 时，\(\mathbf 1_{F_w}^{\perp}\) 为一维符号表示 \(\mathrm{sgn}_2\)。

对整个直积群 \(G_6^{\mathrm{bin}}\) 而言，每个因子 \(\mathfrak S_{d(w)}\) 只作用在对应的 \(\ell^2(F_w)\) 上，因此把上述分解逐纤维相加，便得到所示
\[
\mathcal H_6=\mathcal H_{\mathrm{vis}}\oplus \mathcal H_{\mathrm{hid}}
\]
及隐藏部门的不可约分块。

最后，\(\mathcal H_6^{G_6^{\mathrm{bin}}}\) 恰由各纤维上的常数函数组成，故等于 \(\mathcal H_{\mathrm{vis}}\)。维数计算则由推论 \ref{cor:fold-bin-m6-fiber-hist} 的直方图
\[
4^{\times 9},\qquad 3^{\times 4},\qquad 2^{\times 8}
\]
直接给出。证毕。
\end{proof}

\begin{corollary}[window-\texorpdfstring{$6$}{6} 的 \texorpdfstring{$B_3/C_3$}{B3/C3} 根几何只存在于可见固定点商]\label{cor:xi-time-part55d-window6-root-geometry-visible-fixed-quotient}
沿用定理 \ref{thm:xi-time-part55d-window6-microstate-hilbert-gauge-splitting} 的记号，定义
\[
\mathcal H_{\mathrm{vis}}^{\mathrm{cyc}}
:=
\bigoplus_{w\in X_6^{\mathrm{cyc}}}\CC \mathbf 1_{F_w},
\qquad
\mathcal H_{\mathrm{vis}}^{\mathrm{bdry}}
:=
\bigoplus_{w\in X_6^{\mathrm{bdry}}}\CC \mathbf 1_{F_w}.
\]
则
\[
\mathcal H_{\mathrm{vis}}
=
\mathcal H_{\mathrm{vis}}^{\mathrm{cyc}}
\oplus
\mathcal H_{\mathrm{vis}}^{\mathrm{bdry}}
\cong
\CC^{18}\oplus \CC^3.
\]
并且表 \ref{tab:fold6_b3c3_root_dictionary_b3} 与表 \ref{tab:fold6_b3c3_root_dictionary_c3} 的 \(B_3/C_3\) 字典仅定义在 \(\mathcal H_{\mathrm{vis}}^{\mathrm{cyc}}\) 的自然基上。换言之，window-\(6\) 的根系几何并非微态层上的结构，而是 gauge 固定点商中的可见几何。
\end{corollary}
\begin{proof}
由定理 \ref{thm:xi-window6-cyclic-kernel-bc3-root-datum}，
\[
X_6=X_6^{\mathrm{cyc}}\sqcup X_6^{\mathrm{bdry}},
\qquad
|X_6^{\mathrm{cyc}}|=18,
\qquad
|X_6^{\mathrm{bdry}}|=3.
\]
因此 \(\mathcal H_{\mathrm{vis}}\) 按 cyclic/boundary 精确裂解为
\[
\mathcal H_{\mathrm{vis}}^{\mathrm{cyc}}\oplus \mathcal H_{\mathrm{vis}}^{\mathrm{bdry}}
\cong
\CC^{18}\oplus \CC^3.
\]
另一方面，表 \ref{tab:fold6_b3c3_root_dictionary_b3} 与表 \ref{tab:fold6_b3c3_root_dictionary_c3} 只对 \(X_6^{\mathrm{cyc}}\) 的 \(18\) 个稳定词给出确定性根字典，而不涉及 boundary 三词。故根几何完全落在 \(\mathcal H_{\mathrm{vis}}^{\mathrm{cyc}}\) 上。证毕。
\end{proof}

\begin{theorem}[规范作用像代数与交换子代数的双侧刚性]\label{thm:xi-time-part55d-window6-action-image-commutant-rigidity}
\leanverified{paper\_xi\_time\_part55d\_window6\_action\_image\_commutant\_rigidity}
设
\[
\rho:\CC[G_6^{\mathrm{bin}}]\longrightarrow \operatorname{End}(\mathcal H_6)
\]
为自然置换表示的线性延拓。则其像代数满足
\[
\rho\!\bigl(\CC[G_6^{\mathrm{bin}}]\bigr)
\cong
\CC
\oplus
M_3(\CC)^{\oplus 9}
\oplus
M_2(\CC)^{\oplus 4}
\oplus
\CC^{\oplus 8},
\]
而其交换子代数满足
\[
\rho\!\bigl(\CC[G_6^{\mathrm{bin}}]\bigr)'
\cong
M_{21}(\CC)
\oplus
\CC^{\oplus 9}
\oplus
\CC^{\oplus 4}
\oplus
\CC^{\oplus 8}
\cong
M_{21}(\CC)\oplus \CC^{\oplus 21}.
\]
特别地，\(T\in \operatorname{End}(\mathcal H_6)\) 为自伴且与 \(G_6^{\mathrm{bin}}\) 作用可交换，当且仅当
\[
T
=
A_{\mathrm{vis}}
\oplus
\bigoplus_{w\in X_6}\lambda_w I_{\mathbf 1_{F_w}^{\perp}},
\qquad
A_{\mathrm{vis}}=A_{\mathrm{vis}}^*\in \operatorname{End}(\mathcal H_{\mathrm{vis}}),
\quad
\lambda_w\in \RR.
\]
因此，一切与 gauge 作用等变的非平凡耦合都被压缩在 \(21\) 维可见商 \(\mathcal H_{\mathrm{vis}}\) 中；隐藏部门上既不存在跨纤维耦合，也不存在可见--隐藏混合耦合。
\end{theorem}
\begin{proof}
由定理 \ref{thm:xi-time-part55d-window6-microstate-hilbert-gauge-splitting}，
\[
\mathcal H_6
\cong
\mathbf 1^{\oplus 21}
\oplus
\bigoplus_{w:\,d(w)=4}\mathrm{Std}_4^{(w)}
\oplus
\bigoplus_{w:\,d(w)=3}\mathrm{Std}_3^{(w)}
\oplus
\bigoplus_{w:\,d(w)=2}\mathrm{sgn}_2^{(w)}.
\]
对每个 \(w\)，隐藏块 \(\mathbf 1_{F_w}^{\perp}\) 上只有第 \(w\) 个因子 \(\mathfrak S_{d(w)}\) 非平凡作用，其余因子都平凡。因此每个隐藏块都是 \(G_6^{\mathrm{bin}}\) 的不可约表示；若 \(w\neq w'\)，则把表示限制到第 \(w\) 个因子即可看出：\(\mathbf 1_{F_w}^{\perp}\) 上该因子非平凡，而 \(\mathbf 1_{F_{w'}}^{\perp}\) 上该因子平凡，故这些隐藏块两两不同构。

另一方面，Burnside 定理给出
\[
\CC[\mathfrak S_4]\big|_{\mathrm{Std}_4}\cong M_3(\CC),
\qquad
\CC[\mathfrak S_3]\big|_{\mathrm{Std}_3}\cong M_2(\CC),
\qquad
\CC[\mathfrak S_2]\big|_{\mathrm{sgn}_2}\cong \CC.
\]
而 \(21\) 条常数线共同构成平凡表示的 \(21\) 重等型分量，群代数在其上只通过标量作用。于是像代数正是
\[
\CC
\oplus
M_3(\CC)^{\oplus 9}
\oplus
M_2(\CC)^{\oplus 4}
\oplus
\CC^{\oplus 8}.
\]

再由 Schur 引理，交换子在平凡表示的 \(21\) 重等型块上给出全部矩阵代数
\[
\operatorname{End}(\mathcal H_{\mathrm{vis}})\cong M_{21}(\CC),
\]
而在每个 multiplicity-one 的隐藏不可约块上只能给出标量。故
\[
\rho\!\bigl(\CC[G_6^{\mathrm{bin}}]\bigr)'
\cong
M_{21}(\CC)\oplus \CC^{\oplus 21}.
\]
最后，自伴条件只要求 \(A_{\mathrm{vis}}\) 自伴且各标量 \(\lambda_w\) 取实数。由不同不可约分量之间不存在非零交织算子，隐藏部门上不可能出现跨纤维项，也不可能与可见部门混合。证毕。
\end{proof}

\begin{theorem}[粗粒化算子的三值奇异谱]\label{thm:xi-time-part55d-window6-coarse-graining-three-singular-values}
\leanverified{paper\_xi\_time\_part55d\_window6\_coarse\_graining\_three\_singular\_values}
在 \(\ell^2(X_6)\) 的标准正交基 \((e_w)_{w\in X_6}\) 下，定义粗粒化算子
\[
B:\mathcal H_6\longrightarrow \ell^2(X_6),
\qquad
(Bf)(w):=\sum_{n\in F_w}f(n),
\]
并记
\[
D:=\operatorname{diag}(d(w))_{w\in X_6}.
\]
则成立：
\begin{enumerate}
  \item
  \[
  B^*B=\bigoplus_{w\in X_6}J_{d(w)},
  \]
  其中 \(J_d\) 为 \(d\times d\) 全 \(1\) 矩阵；
  \item \(B\) 的非零奇异值恰为
  \[
  \{2^{(9)},\ \sqrt3^{(4)},\ \sqrt2^{(8)}\};
  \]
  \item
  \[
  \ker B=\mathcal H_{\mathrm{hid}},
  \qquad
  \operatorname{rank}B=21;
  \]
  \item 若定义
  \[
  \widetilde B:=D^{-1/2}B,
  \]
  则
  \[
  \widetilde B\widetilde B^*=I_{\ell^2(X_6)},
  \qquad
  \widetilde B^*\widetilde B=P_{\mathrm{vis}},
  \]
  其中 \(P_{\mathrm{vis}}\) 为到 \(\mathcal H_{\mathrm{vis}}\) 的正交投影。特别地，\(\ell^2(X_6)\) 与 \(\mathcal H_{\mathrm{vis}}\) 规范等距同构。
\end{enumerate}
\end{theorem}
\begin{proof}
对每个 \(w\in X_6\)，算子 \(B\) 在 \(\ell^2(F_w)\) 上的限制只是求和泛函：
\[
B|_{\ell^2(F_w)}(f)
=
\left(\sum_{n\in F_w}f(n)\right)e_w.
\]
故
\[
B^*e_w=\mathbf 1_{F_w},
\]
从而 \(B^*B\) 在第 \(w\) 个纤维块上正是 \(J_{d(w)}\)，这就得到第 \(1\) 条。

而 \(J_d\) 的特征值为 \(d\)（重数 \(1\)）与 \(0\)（重数 \(d-1\)），所以 \(B\) 的非零奇异值恰为 \(\sqrt{d(w)}\)。由推论 \ref{cor:fold-bin-m6-fiber-hist} 的直方图，\(d(w)=4,3,2\) 的纤维个数分别为 \(9,4,8\)，于是得到第 \(2\) 条。

同样，\(J_d\) 的零空间正是零和超平面 \(\mathbf 1_{F_w}^{\perp}\)，故
\[
\ker B
=
\bigoplus_{w\in X_6}\mathbf 1_{F_w}^{\perp}
=
\mathcal H_{\mathrm{hid}}.
\]
由
\[
\dim \mathcal H_6=64,
\qquad
\dim \ker B=43
\]
立得 \(\operatorname{rank}B=21\)，从而第 \(3\) 条成立。

最后，直接计算得
\[
BB^*=D,
\]
故
\[
\widetilde B\widetilde B^*
=
D^{-1/2}BB^*D^{-1/2}
=
I_{\ell^2(X_6)}.
\]
另一方面，\(\widetilde B^*\widetilde B\) 在第 \(w\) 个纤维块上等于
\[
d(w)^{-1}J_{d(w)},
\]
这正是到 \(\CC \mathbf 1_{F_w}\) 的正交投影。逐块相加便得到
\[
\widetilde B^*\widetilde B=P_{\mathrm{vis}}.
\]
因此 \(\widetilde B\) 把 \(\mathcal H_{\mathrm{vis}}\) 等距识别为 \(\ell^2(X_6)\)。证毕。
\end{proof}

\begin{theorem}[同纤维团图的质量量子化与热核三指数分解]\label{thm:xi-time-part55d-window6-fiber-clique-mass-quantization}
在 \(\Omega_6\) 上定义图 \(\Gamma_6\)：若且唯若 \(u\neq v\) 且
\[
\Fold_6^{\mathrm{bin}}(u)=\Fold_6^{\mathrm{bin}}(v),
\]
则在 \(u,v\) 之间连一条边。记 \(A_{\Gamma_6}\) 为其邻接矩阵，\(L_{\Gamma_6}\) 为其图拉普拉斯。则
\[
\Gamma_6\cong 9K_4\sqcup 4K_3\sqcup 8K_2.
\]
从而
\[
\operatorname{Spec}(A_{\Gamma_6})
=
\{3^{(9)},\ 2^{(4)},\ 1^{(8)},\ (-1)^{(43)}\},
\]
\[
\operatorname{Spec}(L_{\Gamma_6})
=
\{0^{(21)},\ 4^{(27)},\ 3^{(8)},\ 2^{(8)}\}.
\]
若 \(P_0,P_2,P_3,P_4\) 为 \(L_{\Gamma_6}\) 对应于特征值 \(0,2,3,4\) 的谱投影，则
\[
P_0=P_{\mathrm{vis}},
\qquad
\mathcal H_{\mathrm{hid}}
=
P_2\mathcal H_6\oplus P_3\mathcal H_6\oplus P_4\mathcal H_6,
\]
并且
\[
\dim P_2\mathcal H_6=8,
\qquad
\dim P_3\mathcal H_6=8,
\qquad
\dim P_4\mathcal H_6=27.
\]
进一步，对任意 \(t\ge 0\) 都有精确热核分解
\[
e^{-tL_{\Gamma_6}}
=
P_{\mathrm{vis}}
+
e^{-2t}P_2
+
e^{-3t}P_3
+
e^{-4t}P_4.
\]
因此隐藏部门不含零模，而严格量子化为三条有限质量通道 \(2,3,4\)。
\end{theorem}
\begin{proof}
按定义，\(\Gamma_6\) 的每个连通分支恰是某条纤维 \(F_w\) 上的完全图 \(K_{d(w)}\)。由推论 \ref{cor:fold-bin-m6-fiber-hist}，window-\(6\) 的纤维直方图为
\[
4^{\times 9},\qquad 3^{\times 4},\qquad 2^{\times 8},
\]
故
\[
\Gamma_6\cong 9K_4\sqcup 4K_3\sqcup 8K_2.
\]

对完全图 \(K_d\)，邻接谱为
\[
\{(d-1)^{(1)},\ (-1)^{(d-1)}\},
\]
图拉普拉斯谱为
\[
\{0^{(1)},\ d^{(d-1)}\}.
\]
逐分支相加即得所示总谱。

特别地，\(L_{\Gamma_6}\) 的零特征空间正是每个连通分支上的常数函数空间，因此恰为
\[
\bigoplus_{w\in X_6}\CC \mathbf 1_{F_w}
=
\mathcal H_{\mathrm{vis}}.
\]
相应的正特征空间则是各分支上的零和空间，因而等于 \(\mathcal H_{\mathrm{hid}}\)。其中 \(K_2\) 贡献 \(8\) 维特征值 \(2\) 空间，\(K_3\) 贡献 \(4\cdot 2=8\) 维特征值 \(3\) 空间，\(K_4\) 贡献 \(9\cdot 3=27\) 维特征值 \(4\) 空间。热核公式最后由谱定理直接给出。证毕。
\end{proof}

\begin{conjecture}[六顶点/GFF 极限中的可见--隐藏严格分离]\label{conj:xi-time-part55d-window6-sixvertex-gff-visible-hidden-separation}
沿用定理 \ref{thm:xi-time-part55d-window6-fiber-clique-mass-quantization} 的谱投影 \(P_{\mathrm{vis}},P_2,P_3,P_4\)。设某一临界六顶点尺度极限下，高度函数在适当归一化后收敛到标量高斯自由场。则任意随网格尺度 \(\delta\downarrow 0\) 变化的 window-\(6\) 局部观测量提升
\[
\mathcal O_\delta\in \mathcal H_6
\]
都应满足正交分解
\[
\mathcal O_\delta
=
P_{\mathrm{vis}}\mathcal O_\delta
+
P_2\mathcal O_\delta
+
P_3\mathcal O_\delta
+
P_4\mathcal O_\delta,
\]
并且：
\begin{enumerate}
  \item 只有 \(P_{\mathrm{vis}}\mathcal O_\delta\) 可能产生临界长程极限；
  \item \(P_2\mathcal O_\delta\)、\(P_3\mathcal O_\delta\)、\(P_4\mathcal O_\delta\) 全部只贡献局部且紧的有限质量修正；
  \item window-\(6\) 的隐藏部门不可能再生成第二个独立的无质量 Gaussian 场；
  \item 表 \ref{tab:fold6_b3c3_root_dictionary_b3} 与表 \ref{tab:fold6_b3c3_root_dictionary_c3} 所编码的 \(B_3/C_3\) 根系几何，仅是可见固定点商上的组合记账，而不是新的连续自由场来源。
\end{enumerate}
\end{conjecture}

\begin{remark}
该猜想的要点不在于再次排除额外高斯模，而在于把全部连续极限障碍精确压缩为隐藏部门的三条离散质量 \(2,3,4\)。若其成立，则 window-\(6\) 与六顶点/GFF 的接口将被压缩为一个单一命题：临界性只驻留于 \(21\) 维可见商，而全部隐藏模式只能以有限质量方式附着其上。
\end{remark}

\paragraph{同余类残差、审计容量硬界、可见算术计数与零温地基层}
前述 Dirichlet 扭曲谱、审计偶数窗上的 dyadic 容量、有限分辨率可见算术以及输出位势的零温极限，还可继续压缩为下列五条直接后果。周期层的同余类平方根障碍已由定理 \ref{thm:xi-dirichlet-mode-forces-congruence-residual} 给出；下条表明，经 Witt/M\"obius 投影之后，这一障碍在 primitive 层并不会消失；最后一条则把地基层熵、端点极点尺度与零温矩阵的谱半径压缩为同一条平方律。

\begin{theorem}[Witt/M\"obius 投影下的 primitive 同余类平方根障碍]\label{thm:xi-primitive-congruence-residual-nonrh}
沿用命题 \ref{prop:real-input-40-output-dirichlet-nonrh} 与定理 \ref{thm:xi-dirichlet-mode-forces-congruence-residual} 的记号。对每个模数 \(m\ge 2\) 与剩余类 \(a\in\ZZ/m\ZZ\)，定义 primitive 层的同余计数
\[
P_n(a;m):=\#\{\mathcal O:\ \mathcal O\text{ 为长度 }n\text{ 的 primitive 轨道},\ E(\mathcal O)\equiv a\pmod m\},
\qquad
P_n:=\sum_{a\in\ZZ/m\ZZ}P_n(a;m).
\]
若
\[
\theta_m:=\max_{1\le j\le m-1}\frac{\log\rho_{m,j}}{\log\lambda}>\frac12,
\qquad
\lambda=\rho(M(1))=\varphi^2,
\]
则存在某个剩余类 \(a_\ast\in\ZZ/m\ZZ\) 使
\[
\limsup_{n\to\infty}
\frac{1}{n\log\lambda}
\log\Bigl|P_n(a_\ast;m)-\frac{P_n}{m}\Bigr|
\ge
\theta_m
>
\frac12.
\]
因此，在该模数下，primitive 层的同余类分布同样不可能满足任何 RH 形的平方根指数误差界。
\end{theorem}
\begin{proof}
取 \(j_\ast\in\{1,\dots,m-1\}\) 使
\[
\rho_{m,j_\ast}=\lambda^{\theta_m}=:\rho>\lambda^{1/2}.
\]
定义 primitive Fourier 系数
\[
\widehat P_n(j):=\sum_{a\in\ZZ/m\ZZ}P_n(a;m)\omega_m^{ja}.
\]
由命题 \ref{prop:real-input-40-output-dirichlet-nonrh} 中的 Witt/M\"obius 反演，
\[
\widehat P_n(j_\ast)
=
p_n(\omega_m^{j_\ast})
=
\frac{1}{n}\sum_{d\mid n}\mu(d)\,a_{n/d}(\omega_m^{j_\ast d}).
\]
另一方面，定理 \ref{thm:xi-dirichlet-mode-forces-congruence-residual} 的证明已给出
\[
\limsup_{n\to\infty}|a_n(\omega_m^{j_\ast})|^{1/n}=\rho,
\]
故存在子序列 \(n_r\to\infty\) 满足
\[
|a_{n_r}(\omega_m^{j_\ast})|=\rho^{\,n_r+o(n_r)}.
\]

现估计 Möbius 反演中的其余约数项。对任意 \(d\ge 2\)，由逐项取绝对值与 Perron 单调性，
\[
\rho\!\bigl(M(\omega_m^{j_\ast d})\bigr)\le \rho(M(1))=\lambda,
\]
从而
\[
a_{n_r/d}(\omega_m^{j_\ast d})=O(\lambda^{n_r/d})=O(\lambda^{n_r/2})=o(\rho^{n_r}),
\]
因为 \(\rho>\lambda^{1/2}\)。又约数函数满足 \(\tau(n_r)=e^{o(n_r)}\)，故
\[
\sum_{\substack{d\mid n_r\\ d\ge 2}}\mu(d)\,a_{n_r/d}(\omega_m^{j_\ast d})=o(\rho^{n_r}).
\]
于是沿该子序列有
\[
\widehat P_{n_r}(j_\ast)
=
\frac{1}{n_r}a_{n_r}(\omega_m^{j_\ast})+o\!\left(\frac{\rho^{n_r}}{n_r}\right),
\]
进而
\[
\limsup_{n\to\infty}|\widehat P_n(j_\ast)|^{1/n}\ge \rho.
\]

记 primitive 层的同余残差为
\[
b_n^{\mathrm{prim}}(a;m):=P_n(a;m)-\frac{P_n}{m}.
\]
则对 \(j\neq 0\) 有
\[
\sum_{a\in\ZZ/m\ZZ} b_n^{\mathrm{prim}}(a;m)\,\omega_m^{ja}
=
\widehat P_n(j),
\]
因为 \(\sum_a\omega_m^{ja}=0\)。故
\[
|\widehat P_n(j_\ast)|
\le
\sum_{a\in\ZZ/m\ZZ}|b_n^{\mathrm{prim}}(a;m)|
\le
m\max_{a\in\ZZ/m\ZZ}|b_n^{\mathrm{prim}}(a;m)|.
\]
取 \(n\to\infty\) 的上极限并开 \(n\) 次方，得到
\[
\limsup_{n\to\infty}
\left(
\max_{a\in\ZZ/m\ZZ}\Bigl|P_n(a;m)-\frac{P_n}{m}\Bigr|
\right)^{1/n}
\ge
\rho=\lambda^{\theta_m}.
\]
由于剩余类只有有限个，必存在某个 \(a_\ast\in\ZZ/m\ZZ\) 使相应单类残差的上极限至少为 \(\lambda^{\theta_m}\)。两边取对数并除以 \(n\log\lambda\)，即得所示结论。证毕。
\end{proof}

\leanverified{paper\_xi\_foldbin\_audited\_window\_oracle\_hard\_cap}
\begin{theorem}[审计偶数窗上最小退化扇区强制预言机成功率硬上界]\label{thm:xi-foldbin-audited-window-oracle-hard-cap}
沿用定理 \ref{thm:xi-foldbin-oracle-capacity-central-spectral-truncation} 的记号，并定义
\[
\mathrm{Succ}_m^{\mathrm{bin}}(B):=\frac{C_m(B)}{2^m},
\qquad
(x)_+:=\max\{x,0\}.
\]
若 \(m\in\{6,8,10,12\}\)，则对任意整数 \(B\ge 0\) 都有
\[
\boxed{\;
\mathrm{Succ}_m^{\mathrm{bin}}(B)
\le
1-\frac{F_m}{2^m}\bigl(F_{m/2}-2^B\bigr)_+.\;}
\]
特别地，当 \(B<\log_2 F_{m/2}\) 时，
\[
1-\mathrm{Succ}_m^{\mathrm{bin}}(B)
\ge
\frac{F_m}{2^m}\bigl(F_{m/2}-2^B\bigr)
>
0.
\]
\end{theorem}
\begin{proof}
由定理 \ref{thm:fold-bin-min-degeneracy-fib-audited}，当 \(m\in\{6,8,10,12\}\) 时，恰有 \(F_m\) 条纤维达到最小大小
\[
d_{\min}(m)=F_{m/2}.
\]
在这 \(F_m\) 条纤维上，容量公式中的每一项至多贡献 \(2^B\) 个可恢复微态；因此它们总共至多贡献
\[
F_m\min(F_{m/2},2^B).
\]
其余纤维即使全部完全恢复，也至多贡献
\[
2^m-F_mF_{m/2}
\]
个微态。于是
\[
\begin{aligned}
C_m(B)
&\le
\bigl(2^m-F_mF_{m/2}\bigr)+F_m\min(F_{m/2},2^B)\\
&=
2^m-F_m\bigl(F_{m/2}-2^B\bigr)_+.
\end{aligned}
\]
两边除以 \(2^m\) 即得所示成功率上界。证毕。
\end{proof}

\begin{theorem}[有限分辨率可见算术的单位元与幂等元计数闭式]\label{thm:xi-visible-arithmetic-unit-idempotent-count}
\leanverified{paper\_xi\_visible\_arithmetic\_unit\_idempotent\_count}
沿用定理 \ref{thm:fold-congruence-mod-semirings} 的有限分辨率可见算术记号，记
\[
R_m^{\mathrm{vis}}:=\Omega_m/\!\equiv_m
\cong
\ZZ/(F_{m+2}\ZZ).
\]
则其乘法结构满足
\[
\boxed{\;
\#\bigl(R_m^{\mathrm{vis}}\bigr)^\times
=
\phi_{\mathrm{Euler}}(F_{m+2}),\;}
\]
\[
\boxed{\;
\#\{e\in R_m^{\mathrm{vis}}:\ e^2=e\}
=
2^{\omega(F_{m+2})},\;}
\]
其中 \(\phi_{\mathrm{Euler}}\) 为 Euler totient，\(\omega(N)\) 表示 \(N\) 的不同素因子个数。
\end{theorem}
\begin{proof}
由定理 \ref{thm:fold-congruence-mod-semirings}，
\[
R_m^{\mathrm{vis}}\cong \ZZ/(F_{m+2}\ZZ)
\]
作为半环同构成立。故其乘法单位元类与模 \(F_{m+2}\) 的单位元一一对应，数量正是
\[
\phi_{\mathrm{Euler}}(F_{m+2}).
\]

再把 Fibonacci 模数分解为
\[
F_{m+2}=\prod_{i=1}^{r}p_i^{\nu_i},
\qquad
r=\omega(F_{m+2}).
\]
中国剩余定理给出
\[
\ZZ/(F_{m+2}\ZZ)\cong \prod_{i=1}^{r}\ZZ/(p_i^{\nu_i}\ZZ).
\]
在每个局部因子 \(\ZZ/(p_i^{\nu_i}\ZZ)\) 中，幂等方程 \(e^2=e\) 只有 \(e=0,1\) 两个解；因此全局幂等元可在每个坐标独立选取 \(0\) 或 \(1\)，总数恰为 \(2^r\)。这与定理 \ref{thm:xi-fold-congruence-semiring-singularity-nilpotent-profile} 中给出的幂等元计数完全一致。证毕。
\end{proof}

\begin{theorem}[Fibonacci 同余拓扑与 profinite 拓扑完全一致]\label{thm:xi-visible-arithmetic-fibonacci-adic-profinite-coincidence}
\leanverified{paper\_xi\_visible\_arithmetic\_fibonacci\_adic\_profinite\_coincidence}
沿用定理 \ref{thm:fold-congruence-mod-semirings} 的记号。令
\[
\mathcal N_F:=\{F_{m+2}\ZZ:\ m\ge 1\}
\]
并令 \(\tau_F\) 为使 \(\mathcal N_F\) 中每个理想都成为 \(0\) 邻域的最弱线性拓扑。则 \(\tau_F\) 与 \(\ZZ\) 上的 profinite 拓扑 \(\tau_{\mathrm{pro}}\) 完全一致。因而 \(\tau_F\)-完备化满足规范拓扑环同构
\[
\boxed{\;
\widehat{\ZZ}_F\cong \widehat{\ZZ}.\;}
\]
\end{theorem}
\begin{proof}
先证 \(\tau_F\preceq\tau_{\mathrm{pro}}\)。对每个 \(m\ge 1\)，理想
\[
F_{m+2}\ZZ
\]
本身就是 profinite 邻域基 \(\{N\ZZ:\ N\ge 1\}\) 中的一项，故由 \(\mathcal N_F\) 生成的拓扑不强于 profinite 拓扑。

再证反向包含。任取 \(N\ge 1\)。考虑 Fibonacci 递推在有限集
\[
(\ZZ/N\ZZ)^2
\]
上诱导的状态演化
\[
(F_k,F_{k+1})\longmapsto (F_{k+1},F_{k+2}).
\]
由于状态空间有限，轨道必最终周期；而初态 \((F_0,F_1)=(0,1)\) 在周期中必再现。故存在某个 \(L=L(N)\ge 1\) 使
\[
(F_L,F_{L+1})\equiv (0,1)\pmod N.
\]
于是 \(N\mid F_L\)。若 \(N>1\)，则必有 \(L\ge 3\)，从而可写 \(L=m+2\)；对 \(N=1\) 则结论平凡。因而对任意 \(N\ge 1\) 都存在 \(m\ge 1\) 使
\[
F_{m+2}\ZZ\subseteq N\ZZ.
\]
这说明每个 profinite 邻域都包含某个 Fibonacci 邻域，故 \(\tau_{\mathrm{pro}}\preceq\tau_F\)。

综上 \(\tau_F=\tau_{\mathrm{pro}}\)。两种线性拓扑既然一致，其 Hausdorff 完备化作为同一连续有限商系统的逆极限必自然同构，即
\[
\widehat{\ZZ}_F\cong \widehat{\ZZ}.
\]
证毕。
\end{proof}

\begin{corollary}[可见同余半环族对全部有限模算术的共终性]\label{cor:xi-visible-arithmetic-fibonacci-cofinal-quotients}
\leanverified{paper\_xi\_visible\_arithmetic\_fibonacci\_cofinal\_quotients}
沿用定理 \ref{thm:xi-visible-arithmetic-fibonacci-adic-profinite-coincidence} 与定理 \ref{thm:fold-congruence-mod-semirings} 的记号。对任意正整数 \(N\)，存在某个 \(m\ge 1\) 与满射半环同态
\[
\Psi_{m,N}:R_m^{\mathrm{vis}}=\Omega_m/\!\equiv_m\longtwoheadrightarrow \ZZ/(N\ZZ).
\]
等价地，有限半环族
\[
\{\,\Omega_m/\!\equiv_m\,\}_{m\ge 1}
\]
在所有整数商环 \(\ZZ/(N)\) 的意义下是共终的：每个 \(\ZZ/(N)\) 都是某个 \(\Omega_m/\!\equiv_m\) 的商。
\end{corollary}
\begin{proof}
由定理 \ref{thm:xi-visible-arithmetic-fibonacci-adic-profinite-coincidence} 的证明，对任意 \(N\ge 1\) 可取 \(m\ge 1\) 使
\[
N\mid F_{m+2}.
\]
于是标准降模映射给出满射环同态
\[
\rho_{m,N}:\ZZ/(F_{m+2}\ZZ)\longtwoheadrightarrow \ZZ/(N\ZZ).
\]
再由定理 \ref{thm:fold-congruence-mod-semirings}，
\[
R_m^{\mathrm{vis}}=\Omega_m/\!\equiv_m\cong \ZZ/(F_{m+2}\ZZ),
\]
将 \(\rho_{m,N}\) 经此同构拉回，即得所需的满射半环同态 \(\Psi_{m,N}\)。证毕。
\end{proof}

\begin{theorem}[输出位势零温极限的正熵四态地基层与饱和密度上界]\label{thm:xi-real-input-40-output-freezing-positive-entropy-halffill}
\leanverified{paper\_xi\_real\_input\_40\_output\_freezing\_positive\_entropy\_halffill}
沿用命题 \ref{prop:real-input-40-pressure-freezing}、推论 \ref{cor:real-input-40-ground-entropy} 与定理 \ref{thm:killo-real-input-40-positive-entropy-freezing} 的记号。记
\[
P(\theta):=\log\lambda(e^\theta).
\]
则端点斜率满足
\[
\lim_{\theta\to-\infty}P'(\theta)=0,
\qquad
\lim_{\theta\to+\infty}P'(\theta)=\frac12.
\]
并且零温缩放极限对应于四态一步 SFT \(\Sigma_\infty\)，其转移矩阵为
\[
B_\infty=
\begin{pmatrix}
0&1&1&0\\
0&0&1&0\\
0&0&3&1\\
1&0&0&3
\end{pmatrix},
\]
其地基层熵满足
\[
h_{\mathrm{ground}}
=
\log c
=
\frac12\log\rho(B_\infty)
>
0.
\]
因此，极端正偏置下的冻结端面不可能塌缩为有限个周期轨道的并；同时，可见输出 \(1\) 的典型密度始终被压在区间 \([0,\tfrac12]\) 内。
\end{theorem}
\begin{proof}
命题 \ref{prop:real-input-40-pressure-freezing} 已给出
\[
\lim_{\theta\to-\infty}P'(\theta)=0,
\qquad
\lim_{\theta\to+\infty}P'(\theta)=\frac12,
\]
故极端偏置下的典型输出 \(1\) 密度只能落在 \([0,\tfrac12]\) 内。

另一方面，推论 \ref{cor:real-input-40-ground-entropy} 与定理 \ref{thm:killo-real-input-40-positive-entropy-freezing} 说明：饱和端点对应的零温地基并非有限个最大平均能量周期轨道，而是由上式 \(B_\infty\) 所定义的四态 SFT 支配，且其地基层熵满足
\[
h_{\mathrm{ground}}=\log c=\frac12\log\rho(B_\infty)>0.
\]
若冻结端面仅由有限个周期轨道支撑，则其拓扑熵必为 \(0\)，与 \(h_{\mathrm{ground}}>0\) 矛盾。故极端正偏置极限不可能塌缩为有限周期轨道并。证毕。
\end{proof}

\begin{corollary}[正负偏置零温极限分别冻结于 half-filling 与空相]\label{cor:xi-real-input-40-output-freezing-half-filling-empty-phase}
沿用定理 \ref{thm:xi-real-input-40-output-freezing-positive-entropy-halffill} 的记号。记 \(\mu_\theta\) 为输出位势 \(P(\theta)\) 所对应的平衡态族，并记
\[
\mu_\infty^{+}:=\lim_{\theta\to+\infty}\mu_\theta,
\qquad
\mu_\infty^{-}:=\lim_{\theta\to-\infty}\mu_\theta
\]
为两端零温极限测度。则对观测量 \(\mathbf 1_{\mathrm{out}=1}\) 有
\[
\boxed{\;
\int \mathbf 1_{\mathrm{out}=1}\,\dd\mu_\infty^{+}
=
\frac12,\qquad
\int \mathbf 1_{\mathrm{out}=1}\,\dd\mu_\infty^{-}
=
0.\;}
\]
换言之，正偏置零温极限并不饱和到全占据，而是冻结于严格的 half-filling 地基态；负偏置零温极限则冻结于空相。
\end{corollary}
\begin{proof}
对一参数平衡态族，压力导数给出相应观测量的平衡态期望，即
\[
P'(\theta)=\int \mathbf 1_{\mathrm{out}=1}\,\dd\mu_\theta.
\]
由定理 \ref{thm:xi-real-input-40-output-freezing-positive-entropy-halffill}，
\[
\lim_{\theta\to-\infty}P'(\theta)=0,
\qquad
\lim_{\theta\to+\infty}P'(\theta)=\frac12.
\]
令 \(\theta\to\pm\infty\) 即得
\[
\int \mathbf 1_{\mathrm{out}=1}\,\dd\mu_\infty^{-}=0,
\qquad
\int \mathbf 1_{\mathrm{out}=1}\,\dd\mu_\infty^{+}=\frac12.
\]
又同一定理已把 \(\mu_\infty^{+}\) 识别为由四态矩阵 \(B_\infty\) 的正熵地基层所支配的零温极限，故其饱和值严格停在 \(1/2\) 而非 \(1\)。证毕。
\end{proof}
\leanverified{paper\_xi\_real\_input\_40\_output\_freezing\_half\_filling\_empty\_phase}

\leanverified{paper\_xi\_real\_input\_40\_ground\_entropy\_endpoint\_square\_law}
\begin{theorem}[地基层熵、端点极点尺度与零温矩阵谱半径的平方律]\label{thm:xi-real-input-40-ground-entropy-endpoint-square-law}
沿用定理 \ref{thm:xi-real-input-40-output-freezing-positive-entropy-halffill} 的记号，并记
\[
z_\star(u):=\lambda(u)^{-1},
\qquad
b:=\rho(B_\infty)^{-1}.
\]
则有精确恒等式
\[
\boxed{\;
h_{\mathrm{ground}}
=
\log c
=
-\frac12\log b
=
\frac12\log\rho(B_\infty).\;}
\]
并且当 \(u\to+\infty\) 时，
\[
\boxed{\;
u\,z_\star(u)^2\to b=e^{-2h_{\mathrm{ground}}},
\qquad
z_\star(u)\sim e^{-h_{\mathrm{ground}}}u^{-1/2},
\qquad
\lambda(u)\sim e^{h_{\mathrm{ground}}}u^{1/2}.\;}
\]
\end{theorem}
\begin{proof}
由定理 \ref{thm:xi-real-input-40-output-freezing-positive-entropy-halffill}，
\[
h_{\mathrm{ground}}=\log c=\frac12\log\rho(B_\infty).
\]
故
\[
b=\rho(B_\infty)^{-1}=c^{-2}=e^{-2h_{\mathrm{ground}}},
\]
即得第一条盒装恒等式。

另一方面，命题 \ref{prop:real-input-40-pressure-freezing} 给出
\[
\lambda(u)=c\,u^{1/2}(1+o(1))
\qquad(u\to+\infty).
\]
取倒数便得
\[
z_\star(u)=\lambda(u)^{-1}=c^{-1}u^{-1/2}(1+o(1))
=
e^{-h_{\mathrm{ground}}}u^{-1/2}(1+o(1)).
\]
于是
\[
u\,z_\star(u)^2\to c^{-2}=b=e^{-2h_{\mathrm{ground}}},
\]
再将上式取倒数即得
\[
\lambda(u)\sim e^{h_{\mathrm{ground}}}u^{1/2}.
\]
证毕。
\end{proof}



\paragraph{escort 二态极限、相对熵与 Shannon 常数的闭式回收}
在 escort 温度轴的末位 Gibbs 极限、对数多重度两原子分布以及 Rényi--Shannon 常数谱都已固定之后，还可把 bin-fold 的实参数 escort 族进一步压缩为下列四条直接后果。它们把末位边际、缩放退化度、对数退化度的一二阶矩、相对于稳定层均匀基线的 KL 缺口以及 Shannon 熵常数统一归结为单一 Bernoulli 参数
\[
\theta(a):=\frac{1}{1+\varphi^{a+1}}
\qquad (a\ge 0).
\]

\begin{theorem}[escort 纤维测度的二态极限律]\label{thm:xi-time-part54ab-escort-two-state-limit-law}
\leanverified{paper\_xi\_time\_part54ab\_escort\_two\_state\_limit\_law}
对任意实数 \(a\ge 0\)，定义 escort 纤维测度
\[
\pi_{m,a}(w):=\frac{d_m^{\mathrm{bin}}(w)^a}{S_a^{\mathrm{bin}}(m)},
\qquad
S_a^{\mathrm{bin}}(m):=\sum_{u\in X_m}d_m^{\mathrm{bin}}(u)^a.
\]
令 \(W_m^{(a)}\sim \pi_{m,a}\)，并记
\[
p_1(a):=\theta(a)=\frac{1}{1+\varphi^{a+1}},
\qquad
p_0(a):=1-\theta(a)=\frac{\varphi^{a+1}}{1+\varphi^{a+1}}.
\]
则有
\[
\mathbb P\bigl(W_m^{(a)}(m)=1\bigr)\longrightarrow p_1(a),
\qquad
\mathbb P\bigl(W_m^{(a)}(m)=0\bigr)\longrightarrow p_0(a).
\]
进一步，若定义缩放退化度
\[
Y_m^{(a)}:=\Bigl(\frac{\varphi}{2}\Bigr)^m d_m^{\mathrm{bin}}(W_m^{(a)}),
\]
则
\[
Y_m^{(a)}\Longrightarrow Y_\infty^{(a)},
\qquad
Y_\infty^{(a)}=
\begin{cases}
1,& \text{概率 }p_0(a),\\[4pt]
\varphi^{-1},& \text{概率 }p_1(a).
\end{cases}
\]
\end{theorem}
\begin{proof}
末位边际极限正是定理 \ref{thm:xi-fold-escort-lastbit-gibbs-limit} 在参数 \(q=a\) 处的结论。

再记
\[
X_m^{(a)}:=\log d_m^{\mathrm{bin}}(W_m^{(a)})-m\log\!\Bigl(\frac{2}{\varphi}\Bigr).
\]
则
\[
X_m^{(a)}=\log Y_m^{(a)}.
\]
由定理 \ref{thm:xi-fold-escort-log-multiplicity-two-atom} 在 \(q=a\) 处的结论，
\[
X_m^{(a)}\Longrightarrow
p_0(a)\delta_0+p_1(a)\delta_{-\log\varphi}.
\]
对上式应用连续映射 \(x\mapsto e^x\)，便得到
\[
Y_m^{(a)}\Longrightarrow
p_0(a)\delta_1+p_1(a)\delta_{\varphi^{-1}},
\]
即所述二点极限律。证毕。
\end{proof}

\leanverified{paper\_xi\_time\_part54ab\_escort\_logdefect\_moments}
\begin{theorem}[escort 对数退化度的一阶矩与方差闭式]\label{thm:xi-time-part54ab-escort-logdefect-moments}
沿用定理 \ref{thm:xi-time-part54ab-escort-two-state-limit-law} 的记号。则当 \(m\to\infty\) 时有
\[
\mathbb E_{\pi_{m,a}}\!\bigl[\log d_m^{\mathrm{bin}}(W_m^{(a)})\bigr]
=
m\log\!\Bigl(\frac{2}{\varphi}\Bigr)
-\frac{\log\varphi}{1+\varphi^{a+1}}
+O\!\Bigl(\Bigl(\frac{\varphi}{2}\Bigr)^m\Bigr),
\]
以及
\[
\operatorname{Var}_{\pi_{m,a}}\!\bigl(\log d_m^{\mathrm{bin}}(W_m^{(a)})\bigr)
=
\frac{\varphi^{a+1}}{(1+\varphi^{a+1})^2}(\log\varphi)^2
+O\!\Bigl(\Bigl(\frac{\varphi}{2}\Bigr)^m\Bigr).
\]
\end{theorem}
\begin{proof}
第一式正是命题 \ref{prop:fold-bin-escort-lastbit} 在 \(q=a\) 处的一阶矩展开。

再记
\[
\varepsilon_m:=\Bigl(\frac{\varphi}{2}\Bigr)^m.
\]
由定理 \ref{thm:fold-bin-two-state-asymptotic} 的一致对数展开，
\[
\log d_m^{\mathrm{bin}}(w)
=
m\log\!\Bigl(\frac{2}{\varphi}\Bigr)-w_m\log\varphi+\rho_m(w),
\qquad
\sup_{w\in X_m}|\rho_m(w)|=O(\varepsilon_m).
\]
因此
\[
\log d_m^{\mathrm{bin}}(W_m^{(a)})
=
m\log\!\Bigl(\frac{2}{\varphi}\Bigr)
-(\log\varphi)\,W_m^{(a)}(m)
+O_{L^\infty}(\varepsilon_m).
\]
常数平移不影响方差，故
\[
\operatorname{Var}_{\pi_{m,a}}\!\bigl(\log d_m^{\mathrm{bin}}(W_m^{(a)})\bigr)
=
(\log\varphi)^2\operatorname{Var}_{\pi_{m,a}}\!\bigl(W_m^{(a)}(m)\bigr)
+O(\varepsilon_m).
\]
另一方面，命题 \ref{prop:fold-bin-escort-lastbit} 给出
\[
\mathbb P\bigl(W_m^{(a)}(m)=1\bigr)
=
\theta(a)+O(\varepsilon_m).
\]
由于 \(W_m^{(a)}(m)\) 是 Bernoulli 变量，便有
\[
\operatorname{Var}_{\pi_{m,a}}\!\bigl(W_m^{(a)}(m)\bigr)
=
\theta(a)\bigl(1-\theta(a)\bigr)+O(\varepsilon_m).
\]
代入并用
\[
\theta(a)\bigl(1-\theta(a)\bigr)
=
\frac{\varphi^{a+1}}{(1+\varphi^{a+1})^2},
\]
即得所示方差公式。证毕。
\end{proof}

\begin{theorem}[escort 相对于稳定层均匀基线的 KL 极限与 Bernoulli 约化]\label{thm:xi-time-part54ab-escort-uniform-kl-bernoulli-reduction}
\leanverified{paper\_xi\_time\_part54ab\_escort\_uniform\_kl\_bernoulli\_reduction}
设 \(u_m\) 为 \(X_m\) 上的均匀分布，
\[
u_m(w)=\frac{1}{|X_m|}=\frac{1}{F_{m+2}}.
\]
则对任意 \(a\ge 0\)，有
\[
D_{\mathrm{KL}}(\pi_{m,a}\Vert u_m)
=
D_{\mathrm{KL}}\!\Bigl(\mathrm{Bernoulli}(p_1(a))\Vert \mathrm{Bernoulli}(\varphi^{-2})\Bigr)
+O\!\Bigl(\Bigl(\frac{\varphi}{2}\Bigr)^m\Bigr).
\]
等价地，
\[
\lim_{m\to\infty}D_{\mathrm{KL}}(\pi_{m,a}\Vert u_m)
=
p_0(a)\log\!\bigl(p_0(a)\varphi\bigr)
+
p_1(a)\log\!\bigl(p_1(a)\varphi^2\bigr).
\]
\end{theorem}
\begin{proof}
由推论 \ref{cor:xi-fold-escort-uniform-kl-temperature-law} 在 \(q=a\) 处的公式，
\[
D_{\mathrm{KL}}(\pi_{m,a}\Vert u_m)
=
D(a)+O\!\Bigl(\Bigl(\frac{\varphi}{2}\Bigr)^m\Bigr),
\]
其中
\[
D(a)=
2\log\varphi-\log(\varphi+\varphi^{-a})
-\frac{a\log\varphi}{1+\varphi^{a+1}}.
\]
再用
\[
p_1(a)=\frac{1}{1+\varphi^{a+1}},
\qquad
p_0(a)=\frac{\varphi^{a+1}}{1+\varphi^{a+1}},
\qquad
1-\varphi^{-2}=\varphi^{-1},
\]
直接整理可得
\[
D(a)
=
p_0(a)\log\!\frac{p_0(a)}{\varphi^{-1}}
+
p_1(a)\log\!\frac{p_1(a)}{\varphi^{-2}},
\]
亦即
\[
D(a)=
D_{\mathrm{KL}}\!\Bigl(\mathrm{Bernoulli}(p_1(a))\Vert \mathrm{Bernoulli}(\varphi^{-2})\Bigr).
\]
这就得到第一式。第二式只是将同一 Bernoulli KL 按显式坐标写开。证毕。
\end{proof}

\begin{theorem}[escort 测度在稳定层上的 Shannon 熵闭式]\label{thm:xi-time-part54ab-escort-shannon-closed-form}
沿用定理 \ref{thm:xi-time-part54ab-escort-uniform-kl-bernoulli-reduction} 的记号。则 escort 测度 \(\pi_{m,a}\) 的 Shannon 熵满足
\[
H(\pi_{m,a})
=
m\log\varphi-\frac12\log 5
+\frac{\varphi^{a+1}}{1+\varphi^{a+1}}\log\varphi
+h\!\left(\frac{1}{1+\varphi^{a+1}}\right)
+O\!\Bigl(\Bigl(\frac{\varphi}{2}\Bigr)^m\Bigr),
\]
其中
\[
h(p):=-p\log p-(1-p)\log(1-p).
\]
等价地，
\[
H(\pi_{m,a})
=
\log F_{m+2}
-D_{\mathrm{KL}}\!\Bigl(\mathrm{Bernoulli}(p_1(a))\Vert \mathrm{Bernoulli}(\varphi^{-2})\Bigr)
+O\!\Bigl(\Bigl(\frac{\varphi}{2}\Bigr)^m\Bigr).
\]
\end{theorem}
\begin{proof}
由均匀基线的恒等式，
\[
H(\pi_{m,a})=\log|X_m|-D_{\mathrm{KL}}(\pi_{m,a}\Vert u_m)=\log F_{m+2}-D_{\mathrm{KL}}(\pi_{m,a}\Vert u_m).
\]
结合定理 \ref{thm:xi-time-part54ab-escort-uniform-kl-bernoulli-reduction}，立即得到第二个公式。

再由 Binet 渐近，
\[
\log F_{m+2}
=(m+2)\log\varphi-\frac12\log 5+O(\varphi^{-2m}).
\]
将
\[
D_{\mathrm{KL}}\!\Bigl(\mathrm{Bernoulli}(p_1(a))\Vert \mathrm{Bernoulli}(\varphi^{-2})\Bigr)
=
p_0(a)\log\!\bigl(p_0(a)\varphi\bigr)
+
p_1(a)\log\!\bigl(p_1(a)\varphi^2\bigr)
\]
代入，并用 \(p_0(a)+p_1(a)=1\) 整理，得到
\[
\begin{aligned}
H(\pi_{m,a})
&=
(m+2)\log\varphi-\frac12\log 5
-p_0(a)\log\!\bigl(p_0(a)\varphi\bigr)
-p_1(a)\log\!\bigl(p_1(a)\varphi^2\bigr)
+O\!\Bigl(\Bigl(\frac{\varphi}{2}\Bigr)^m\Bigr)\\
&=
m\log\varphi-\frac12\log 5
+p_0(a)\log\varphi
-p_0(a)\log p_0(a)-p_1(a)\log p_1(a)
+O\!\Bigl(\Bigl(\frac{\varphi}{2}\Bigr)^m\Bigr).
\end{aligned}
\]
最后注意
\[
-p_0(a)\log p_0(a)-p_1(a)\log p_1(a)=h\bigl(p_1(a)\bigr),
\]
便得第一式。证毕。
\end{proof}

\paragraph{自然扩张字母表极小性、整数椭圆模核共谱与分层 Shannon 预算}
沿用定理 \ref{thm:killo-natural-extension-branch-register}、定理
\ref{thm:killo-smith-loss-spectrum}、推论
\ref{cor:killo-kernel-size-general-modulus}、定理
\ref{thm:killo-ellipse-diagonal-prime-register-equivalence} 与定理
\ref{thm:xi-prime-slice-nontrivial-layer-exact-minimality} 的记号。前述链条已经分别给出自然扩张的分支指标寄存器实现、Smith 核谱的可逆读出、模数核大小的全局闭式，以及整数轴对齐椭圆的 prime-register 语义。由此还可进一步抽出如下四条彼此咬合的后果：H.153 型寄存器实现的最小有限字母表定律、整数椭圆的 primewise 无序模核共谱分类、对应 Dirichlet 级数的 \(\zeta\)-主因子，以及分层外置的平均 Shannon 预算律。

\begin{theorem}[自然扩张的 H.153 型分支寄存器实现具有精确最小有限字母表]\label{thm:xi-natural-extension-h153-minimal-finite-alphabet}
\leanverified{paper\_xi\_natural\_extension\_h153\_minimal\_finite\_alphabet}
设 \(X\) 为可数集合，\(T:X\twoheadrightarrow X\) 为有限纤维满射，并记
\[
M(T):=\sup_{y\in X}\#T^{-1}(y)\in \NN\cup\{\infty\}.
\]
对任意整数 \(M\ge 1\)，下列两条等价：
\begin{enumerate}
  \item 存在大小为 \(M\) 的有限字母表 \(A\)，以及一族单射
  \[
  \iota_y:T^{-1}(y)\hookrightarrow A\qquad(y\in X),
  \]
  使得由
  \[
  \beta_A(x):=\iota_{T(x)}(x)
  \]
  诱导的 H.153 型编码
  \[
  \mathcal C_A:X^{\leftarrow}\longrightarrow X\times A^{\NN},
  \qquad
  \mathcal C_A((x_0,x_{-1},x_{-2},\dots))
  :=
  \bigl(x_0,(\beta_A(x_{-1}),\beta_A(x_{-2}),\dots)\bigr)
  \]
  为单射，并把自然扩张左移共轭到右移插入系统
  \[
  \widehat T_A\bigl(x,(a_1,a_2,\dots)\bigr)
  :=
  \bigl(T(x),(\beta_A(x),a_1,a_2,\dots)\bigr).
  \]
  \item \(M\ge M(T)\)。
\end{enumerate}
特别地，若 \(M(T)<\infty\)，则 H.153 型自然扩张分支寄存器实现所需的最小有限字母表大小恰为 \(M(T)\)；若 \(M(T)=\infty\)，则不存在任何有限字母表实现。
\end{theorem}
\begin{proof}
先证必要性。若存在第 \(1\) 条所述实现，则对每个 \(y\in X\) 都已有单射
\[
\iota_y:T^{-1}(y)\hookrightarrow A.
\]
故
\[
\#T^{-1}(y)\le |A|=M
\qquad(\forall\,y\in X),
\]
从而
\[
M(T)=\sup_{y\in X}\#T^{-1}(y)\le M.
\]

再证充分性。若 \(M\ge M(T)\)，则对每个 \(y\in X\) 都有
\[
\#T^{-1}(y)\le M=|A|.
\]
因此可以逐点选取单射
\[
\iota_y:T^{-1}(y)\hookrightarrow A.
\]
由此定义 \(\beta_A\) 与 \(\mathcal C_A\)。定理 \ref{thm:killo-natural-extension-branch-register} 的递推反演证明逐字适用：给定
\[
\bigl(x_0,(a_1,a_2,\dots)\bigr)\in \mathcal C_A(X^{\leftarrow}),
\]
可沿像上反演逐步恢复唯一前史
\[
x_{-1}=\iota_{x_0}^{-1}(a_1),\qquad
x_{-2}=\iota_{x_{-1}}^{-1}(a_2),\ \dots,
\]
从而 \(\mathcal C_A\) 为单射；并且
\[
\mathcal C_A\circ \sigma=\widehat T_A\circ \mathcal C_A.
\]
故第 \(1\) 条成立。证毕。
\end{proof}

\begin{theorem}[整数轴对齐椭圆的模核谱只看见 primewise 无序轴指数]\label{thm:xi-integer-ellipse-mod-kernel-spectrum-primewise-unordered}
\leanverified{paper\_xi\_integer\_ellipse\_mod\_kernel\_spectrum\_primewise\_unordered}
沿用定理 \ref{thm:killo-ellipse-diagonal-prime-register-equivalence} 的椭圆记号。对
\[
E_{a,b},E_{c,d}\in \mathsf E_{\NN},
\qquad
A_{a,b}:=\operatorname{diag}(a,b),
\qquad
A_{c,d}:=\operatorname{diag}(c,d),
\]
并对每个素数 \(p\) 与整数 \(k\ge 1\) 定义
\[
K_p^{a,b}(k):=\#\ker(A_{a,b}\bmod p^k).
\]
则成立等价
\[
K_p^{a,b}(k)=K_p^{c,d}(k)\quad(\forall\,p,\ \forall\,k\ge 1)
\iff
\{v_p(a),v_p(b)\}=\{v_p(c),v_p(d)\}\quad(\forall\,p).
\]
亦即，两个有序整数轴对齐椭圆具有完全相同的模核谱，当且仅当每个素数上的两根轴指数只发生逐素数的坐标交换。

进一步，定义
\[
\omega_{\neq}(a,b):=\#\{\,p:\ v_p(a)\neq v_p(b)\,\}.
\]
则与 \(E_{a,b}\) 共谱的有序椭圆恰有
\[
2^{\omega_{\neq}(a,b)}
\]
个。
\end{theorem}
\begin{proof}
对对角矩阵
\[
A_{a,b}=\operatorname{diag}(a,b)
\]
作 Smith 归约时，其 Smith 不变量即为
\[
d_1=\gcd(a,b),
\qquad
d_2=\operatorname{lcm}(a,b).
\]
故对每个素数 \(p\)，若记
\[
\alpha_p:=v_p(a),
\qquad
\beta_p:=v_p(b),
\]
则
\[
v_p(d_1)=\min(\alpha_p,\beta_p),
\qquad
v_p(d_2)=\max(\alpha_p,\beta_p).
\]
由定理 \ref{thm:killo-smith-loss-spectrum}，模核谱
\[
\bigl(K_p^{a,b}(k)\bigr)_{p,k}
\]
可逆恢复每个 \(p\) 上的 Smith 赋值对
\[
\bigl(v_p(d_1),v_p(d_2)\bigr)
=
\bigl(\min(\alpha_p,\beta_p),\max(\alpha_p,\beta_p)\bigr).
\]
因此它所恢复的恰是无序二元组
\[
\{\alpha_p,\beta_p\}=\{v_p(a),v_p(b)\}.
\]
这就证明了前一条等价关系。

现计算共谱有序椭圆的数目。对每个满足 \(\alpha_p\neq \beta_p\) 的素数 \(p\)，保持无序集合
\[
\{\alpha_p,\beta_p\}
\]
不变时，只有两种可能的有序分配：把较小指数放在第一轴、较大指数放在第二轴，或反之交换。若 \(\alpha_p=\beta_p\)，则交换不产生新对象。由于不同素数上的选择彼此独立，故总共有
\[
2^{\omega_{\neq}(a,b)}
\]
个有序椭圆与 \(E_{a,b}\) 共谱。证毕。
\end{proof}

\begin{corollary}[整数轴对齐椭圆的模核 Dirichlet 级数是 \texorpdfstring{$\zeta$}{zeta} 的有限 Euler 修正]\label{cor:xi-integer-ellipse-mod-kernel-dirichlet-zeta}
沿用定理 \ref{thm:xi-integer-ellipse-mod-kernel-spectrum-primewise-unordered} 的记号。对每个 \(N\ge 1\)，定义模数核计数函数
\[
K_{a,b}(N):=\#\ker(A_{a,b}\bmod N).
\]
则有精确闭式
\[
K_{a,b}(N)=\gcd(a,N)\gcd(b,N).
\]
因此 Dirichlet 级数
\[
\mathcal Z_{a,b}(s):=\sum_{N\ge 1}\frac{K_{a,b}(N)}{N^s}
\]
在 \(\Re(s)>1\) 上绝对收敛，并具有 Euler 乘积分解
\[
\mathcal Z_{a,b}(s)
=
\prod_p
\left(
\sum_{k\ge 0}p^{-ks+\min(k,\alpha_p)+\min(k,\beta_p)}
\right),
\qquad
\alpha_p:=v_p(a),\ \beta_p:=v_p(b).
\]
更进一步，
\[
\mathcal Z_{a,b}(s)=\zeta(s)\,R_{a,b}(s),
\]
其中
\[
R_{a,b}(s):=
\prod_{p\mid ab}
\Bigl(1-p^{-s}\Bigr)
\sum_{k\ge 0}p^{-ks+\min(k,\alpha_p)+\min(k,\beta_p)}
\]
是仅在坏素数 \(p\mid ab\) 上出现的有限 Euler 修正。因而 \(\mathcal Z_{a,b}(s)\) 在 \(s=1\) 处具有唯一简单极点。

此外，
\[
\mathcal Z_{a,b}(s)=\mathcal Z_{c,d}(s)
\iff
\bigl(K_p^{a,b}(k)\bigr)_{p,k}=\bigl(K_p^{c,d}(k)\bigr)_{p,k},
\]
从而也等价于 \(E_{a,b}\) 与 \(E_{c,d}\) 在定理 \ref{thm:xi-integer-ellipse-mod-kernel-spectrum-primewise-unordered} 的意义下共谱。
\end{corollary}
\begin{proof}
由推论 \ref{cor:killo-kernel-size-general-modulus}，对任意 \(A\in M_{m\times n}(\ZZ)\) 都有
\[
\#\ker(A_N)=N^{\,n-m}\prod_{i=1}^{m}\gcd(d_i,N),
\]
其中 \(d_i\) 为 Smith 不变量。对
\[
A_{a,b}=\operatorname{diag}(a,b)
\]
有 \(m=n=2\)，且 Smith 不变量为
\[
d_1=\gcd(a,b),
\qquad
d_2=\operatorname{lcm}(a,b).
\]
故
\[
K_{a,b}(N)=\gcd(d_1,N)\gcd(d_2,N).
\]
逐素数比较即可验证
\[
\gcd(d_1,N)\gcd(d_2,N)=\gcd(a,N)\gcd(b,N),
\]
于是得到所示精确闭式。

由右端关于 \(N\) 的积性，Dirichlet 级数分解为 Euler 乘积
\[
\mathcal Z_{a,b}(s)
=
\prod_p \sum_{k\ge 0}\frac{K_{a,b}(p^k)}{p^{ks}}.
\]
而
\[
K_{a,b}(p^k)=p^{\min(k,\alpha_p)+\min(k,\beta_p)},
\]
故得到所示局部因子公式。

若 \(p\nmid ab\)，则 \(\alpha_p=\beta_p=0\)，从而
\[
\sum_{k\ge 0}p^{-ks+\min(k,\alpha_p)+\min(k,\beta_p)}
=
\sum_{k\ge 0}p^{-ks}
=
\frac{1}{1-p^{-s}}.
\]
把全部好素数因子提出，便得到
\[
\mathcal Z_{a,b}(s)
=
\zeta(s)
\prod_{p\mid ab}
\Bigl(1-p^{-s}\Bigr)
\sum_{k\ge 0}p^{-ks+\min(k,\alpha_p)+\min(k,\beta_p)}.
\]
由于坏素数只有有限多个，\(R_{a,b}(s)\) 为有限 Euler 乘积；并且每个局部修正因子都是 \(p^{-s}\) 的有理函数，故 \(R_{a,b}\) 在 \(s=1\) 的邻域全纯。又因各局部因子在 \(s=1\) 处都取正实值，故 \(R_{a,b}(1)\neq 0\)。因此 \(\mathcal Z_{a,b}(s)\) 只在 \(s=1\) 继承 \(\zeta(s)\) 的唯一简单极点。

最后，若 \(\mathcal Z_{a,b}(s)=\mathcal Z_{c,d}(s)\)，则两边 Dirichlet 系数逐项相同，特别地对每个素数幂 \(N=p^k\) 有
\[
K_p^{a,b}(k)=K_p^{c,d}(k).
\]
反向蕴含显然。再结合定理 \ref{thm:xi-integer-ellipse-mod-kernel-spectrum-primewise-unordered}，便得到与共谱条件的等价性。证毕。
\end{proof}

\begin{theorem}[分层外置的平均 Shannon 极小预算由层间条件熵完全控制]\label{thm:xi-layered-externalization-average-shannon-budget}
\leanverified{paper\_xi\_layered\_externalization\_average\_shannon\_budget}
设
\[
X_0\xrightarrow{\pi_0}X_1\xrightarrow{\pi_1}\cdots\xrightarrow{\pi_{k-1}}X_k
\]
为有限纤维满射塔，并在 \(X_0\) 上赋予任意概率分布 \(\mu\)。记
\[
Y_0\sim\mu,
\qquad
Y_{j+1}:=\pi_j(Y_j)\qquad(0\le j\le k-1),
\]
并对每个 \(j\) 与 \(y\in X_{j+1}\) 记纤维
\[
F_j(y):=\pi_j^{-1}(y).
\]

对每个 \(j\) 与 \(y\in X_{j+1}\)，取一个二进制前缀码
\[
c_{j,y}:F_j(y)\longrightarrow \{0,1\}^{\ast}.
\]
把分层侧信息定义为
\[
\Psi_c(x_0)
:=
\Bigl(
x_k,\,
c_{k-1,x_k}(x_{k-1})\cdots c_{1,x_2}(x_1)c_{0,x_1}(x_0)
\Bigr),
\]
其中 \(x_{j+1}=\pi_j(x_j)\)。记平均侧信息长度为
\[
L(c):=
\EE_\mu\!\left[
\sum_{j=0}^{k-1}
\bigl|c_{j,Y_{j+1}}(Y_j)\bigr|
\right].
\]
则有下界
\[
L(c)\ge
\frac{1}{\log 2}\sum_{j=0}^{k-1}H_\mu(Y_j\mid Y_{j+1}).
\]
并且存在一族二进制前缀码 \(\{c_{j,y}\}\) 使
\[
L(c)\le
\frac{1}{\log 2}\sum_{j=0}^{k-1}H_\mu(Y_j\mid Y_{j+1})+k.
\]
若上述条件熵总和有限，则由链式法则右端也可写为
\[
\frac{1}{\log 2}H_\mu(Y_0\mid Y_k)
\quad\text{与}\quad
\frac{1}{\log 2}H_\mu(Y_0\mid Y_k)+k.
\]
\end{theorem}
\begin{proof}
先证下界。固定一层 \(j\) 与一点 \(y\in X_{j+1}\)。由于
\[
c_{j,y}:F_j(y)\to\{0,1\}^{\ast}
\]
是二进制前缀码，Shannon 下界给出其在条件分布
\[
\mu(\cdot\mid Y_{j+1}=y)
\]
下的平均码长满足
\[
\sum_{x\in F_j(y)}
\mu(x\mid y)\,|c_{j,y}(x)|
\ge
\frac{1}{\log 2}
\left(
-\sum_{x\in F_j(y)}\mu(x\mid y)\log\mu(x\mid y)
\right).
\]
右端恰为
\[
\frac{1}{\log 2}H_\mu(Y_j\mid Y_{j+1}=y).
\]
再对 \(y\) 以 \(Y_{j+1}\) 的分布加权，即得
\[
\EE_\mu\bigl[|c_{j,Y_{j+1}}(Y_j)|\bigr]
\ge
\frac{1}{\log 2}H_\mu(Y_j\mid Y_{j+1}).
\]
对 \(j=0,\dots,k-1\) 求和，便得到
\[
L(c)\ge
\frac{1}{\log 2}\sum_{j=0}^{k-1}H_\mu(Y_j\mid Y_{j+1}).
\]

再证上界。对每个固定的 \((j,y)\)，在有限集 \(F_j(y)\) 上按条件分布
\[
\mu(\cdot\mid Y_{j+1}=y)
\]
构造 Shannon--Fano 码；若某些点具有零条件概率，则在有限纤维上任意补全为前缀码即可，而不会影响条件平均长度。于是可取 \(c_{j,y}\) 使
\[
\sum_{x\in F_j(y)}
\mu(x\mid y)\,|c_{j,y}(x)|
\le
\frac{1}{\log 2}H_\mu(Y_j\mid Y_{j+1}=y)+1.
\]
再对 \(y\) 加权，得到
\[
\EE_\mu\bigl[|c_{j,Y_{j+1}}(Y_j)|\bigr]
\le
\frac{1}{\log 2}H_\mu(Y_j\mid Y_{j+1})+1.
\]
对全部层求和，便有
\[
L(c)\le
\frac{1}{\log 2}\sum_{j=0}^{k-1}H_\mu(Y_j\mid Y_{j+1})+k.
\]

最后，若上述熵量皆有限，则由于 \(Y_{j+1}\) 是 \(Y_j\) 的确定函数，
\[
H_\mu(Y_j,Y_{j+1})=H_\mu(Y_j),
\]
从而
\[
H_\mu(Y_j\mid Y_{j+1})=H_\mu(Y_j)-H_\mu(Y_{j+1}).
\]
望远镜求和即得
\[
\sum_{j=0}^{k-1}H_\mu(Y_j\mid Y_{j+1})
=
H_\mu(Y_0)-H_\mu(Y_k)
=
H_\mu(Y_0\mid Y_k).
\]
证毕。
\end{proof}

\noindent
这些结果表明：H.153 型自然扩张分支寄存器实现的有限字母表消耗并非可任意压缩，而是恰由最大纤维基数锁定；H.154--H.156 在整数轴对齐椭圆语义下并不恢复逐素数“哪根轴承载哪一指数”，而只恢复每个素数上的无序二元组；相应模核 Dirichlet 级数则正是 \(\zeta(s)\) 的有限 Euler 修正。另一方面，H.151--H.152 所给出的“每个活跃层至少一枚有效 prime-slice”最坏情形定律，在平均意义下精确退化为由层间条件熵总和控制的 Shannon 极小预算。这便把自然扩张、Smith 模核谱、整数椭圆与分层外置的硬/软复杂度进一步压缩为同一条可审计结构链。


\paragraph{分层 prime-slice 的局部字母表极小性与 Smith \texorpdfstring{$p$}{p}-核谱的最短充分前缀}
在定理 \ref{thm:xi-prime-slice-nontrivial-layer-exact-minimality} 已把“哪些层必须激活 prime-slice”刻画为活跃分支层精确极小性，而定理 \ref{thm:xi-smith-loss-discrete-curvature-atoms} 已把 Smith 信息亏损谱的离散曲率原子写成 \(\{e_i(p)\}\) 的显式重数之后，还可继续把局部资源律与最短前缀律进一步压缩为下列三条结论。

\begin{theorem}[分层 prime-slice 外置的局部最小标签字母表恰等于最大纤维数]\label{thm:xi-layered-primeslice-local-alphabet-fibermax}
设
\[
X_0\xrightarrow{\pi_0}X_1\xrightarrow{\pi_1}\cdots\xrightarrow{\pi_{k-1}}X_k
\]
为可数集合上的有限纤维满射塔。沿用定理 \ref{thm:killo-layered-prime-slices-finite-depth} 的 prime-slice 记号
\[
\mathcal P_j:=\{p_{2^j(2n-1)}:n\ge 1\}
\qquad (0\le j\le k-1),
\]
并对每层定义最大纤维数
\[
B_j:=\sup_{y\in X_{j+1}}\#\pi_j^{-1}(y)<\infty.
\]
固定某个层号 \(j\)。考虑任意映射
\[
q_j:X_j\to \Lambda_j,
\qquad
\Lambda_j\subset \mathcal P_j,
\]
并要求
\[
(x_{j+1},q_j(x_j))
\]
能够在全部 \(x_j\in X_j\) 上唯一恢复 \(x_j\)。则所需的最小标签字母表大小恰为
\[
|\Lambda_j|_{\min}=B_j.
\]
更具体地：
\begin{enumerate}
  \item 存在只使用 \(\mathcal P_j\) 中恰好 \(B_j\) 个素数值的第 \(j\) 层标签，能够在该层全部纤维上保持可逆；
  \item 若 \(|\Lambda_j|<B_j\)，则不存在任何满足上述恢复要求的第 \(j\) 层标签。
\end{enumerate}
\end{theorem}
\leanverified{paper\_xi\_layered\_primeslice\_local\_alphabet\_fibermax}
\begin{proof}
先证必要性。取 \(y\in X_{j+1}\) 使
\[
\#\pi_j^{-1}(y)=B_j.
\]
若 \(|\Lambda_j|<B_j\)，则由抽屉原理，在同一纤维 \(\pi_j^{-1}(y)\) 中存在不同点
\[
u\neq v
\]
满足
\[
q_j(u)=q_j(v).
\]
但解码时 \(x_{j+1}=y\) 已知，而第 \(j\) 层额外只看到共同标签 \(q_j(u)=q_j(v)\)，故无法区分 \(u\) 与 \(v\)。这与“由 \((x_{j+1},q_j(x_j))\) 唯一恢复 \(x_j\)”矛盾。故必有
\[
|\Lambda_j|\ge B_j.
\]

再证充分性。对每个 \(y\in X_{j+1}\)，任选一个单射
\[
\iota_y:\pi_j^{-1}(y)\hookrightarrow [B_j]:=\{1,\dots,B_j\}.
\]
取 \(\mathcal P_j\) 中最小的 \(B_j\) 个素数，记为
\[
\Lambda_j^{\min}:=\{p_{2^j(2n-1)}:1\le n\le B_j\}\subset \mathcal P_j.
\]
固定一个双射
\[
\eta_j:[B_j]\xrightarrow{\sim}\Lambda_j^{\min},
\]
并定义
\[
q_j(x_j):=\eta_j\!\bigl(\iota_{\pi_j(x_j)}(x_j)\bigr)\in \Lambda_j^{\min}.
\]
由于 \(\iota_{\pi_j(x_j)}\) 在每条纤维上单射，故已知 \((x_{j+1},q_j(x_j))\) 时，可先由 \(\eta_j^{-1}\) 恢复索引，再由 \(\iota_{x_{j+1}}^{-1}\) 在该纤维上恢复唯一的 \(x_j\)。因此该层确可用恰好 \(B_j\) 个 prime labels 实现精确可逆。证毕。
\end{proof}

\begin{corollary}[分层 prime-slice 基础设施库存的显式 prime-index 公式]\label{cor:xi-layered-primeslice-inventory-prime-index-formula}
\leanverified{paper\_xi\_layered\_primeslice\_inventory\_prime\_index\_formula}
沿用定理 \ref{thm:xi-layered-primeslice-local-alphabet-fibermax} 的记号。若对每一层 \(j\) 都要求预先配置一个显式 prime-label 字母表 \(\Lambda_j\subset\mathcal P_j\)，并且该字母表在该层达到最小可逆规模 \(|\Lambda_j|=B_j\)，则最小可行标签集唯一地取为
\[
\Lambda_j^{\min}
=
\{p_{2^j},p_{3\cdot 2^j},p_{5\cdot 2^j},\dots,p_{(2B_j-1)2^j}\}.
\]
因而第 \(j\) 层所需的最小最大素数恰为
\[
p_{(2B_j-1)2^j}.
\]
此外，这种“逐层显式标签基础设施”所需的 distinct prime 总数恰为
\[
\sum_{j=0}^{k-1}B_j.
\]
若进一步存在统一上界
\[
B_j\le B\qquad(0\le j\le k-1),
\]
则其最小最大素数满足
\[
\max_{0\le j\le k-1}\max \Lambda_j^{\min}
=
p_{(2B-1)2^{k-1}},
\]
并且
\[
\log p_{(2B-1)2^{k-1}}
=
(k-1)\log 2+\log(2B-1)+\log k+O(1).
\]
\end{corollary}
\begin{proof}
由定理 \ref{thm:xi-layered-primeslice-local-alphabet-fibermax}，第 \(j\) 层达到可逆所需的最小标签数恰为 \(B_j\)。若要在 \(\mathcal P_j\) 内最小化所用素数值，则必然取该 prime slice 中最小的 \(B_j\) 个元素。由于
\[
\mathcal P_j=\{p_{2^j(2n-1)}:n\ge 1\},
\]
故最小标签集正是
\[
\Lambda_j^{\min}
=
\{p_{2^j},p_{3\cdot 2^j},\dots,p_{(2B_j-1)2^j}\}.
\]
其最大元素即为
\[
p_{(2B_j-1)2^j}.
\]

不同层的 prime slices \(\mathcal P_j\) 两两不交，故若每层都显式配置其全部最小标签字母表，则全塔所需 distinct primes 数正好是各层字母表大小之和：
\[
\sum_{j=0}^{k-1}|\Lambda_j^{\min}|=\sum_{j=0}^{k-1}B_j.
\]
这给出的是“逐层显式 prime-label 基础设施”的库存规模。若允许某些满足 \(B_j=1\) 的层直接退化为平凡因子 \(1\)，则全局最小非平凡 prime-slice 数仍由定理 \ref{thm:xi-prime-slice-nontrivial-layer-exact-minimality} 精确控制为活跃分支层数。

若 \(B_j\le B\)，则对每个 \(j\) 都有
\[
\max \Lambda_j^{\min}\le p_{(2B-1)2^j},
\]
而右端随 \(j\) 单调递增，故最大值在 \(j=k-1\) 处取得。于是
\[
\max_{0\le j\le k-1}\max \Lambda_j^{\min}=p_{(2B-1)2^{k-1}}.
\]
最后设
\[
N:=(2B-1)2^{k-1}.
\]
由素数定理的标准形式
\[
p_N=N(\log N+\log\log N-1+o(1))
\]
可得
\[
\log p_N=\log N+\log\log N+O(1).
\]
而
\[
\log N=(k-1)\log 2+\log(2B-1),
\]
并且
\[
\log\log N=\log k+O(1),
\]
从而得到所示对数渐近。证毕。
\end{proof}

\begin{theorem}[Smith \texorpdfstring{$p$}{p}-核谱的最短充分前缀与 plateau 检测门槛]\label{thm:xi-smith-p-kernel-shortest-sufficient-prefix}
\leanverified{paper\_xi\_smith\_p\_kernel\_shortest\_sufficient\_prefix}
设 \(A\in M_{m\times n}(\ZZ)\) 且 \(\operatorname{rank}_{\QQ}(A)=m\)。对固定素数 \(p\)，沿用定理 \ref{thm:killo-smith-loss-spectrum} 与定理 \ref{thm:xi-smith-loss-discrete-curvature-atoms} 的记号，记
\[
e_1(p),\dots,e_m(p)\in \ZZ_{\ge 0},
\qquad
f_p(k):=\sum_{i=1}^{m}\min(k,e_i(p))
\qquad(k\ge 1),
\]
并约定
\[
f_p(0):=0,
\qquad
E_p:=\max_{1\le i\le m}e_i(p).
\]
等价地，若记
\[
K_p(k):=\#\ker(A\bmod p^k),
\]
则
\[
f_p(k)=\log_p K_p(k)-k(n-m).
\]
则成立：
\begin{enumerate}
  \item 若 \(E_p\) 先验已知，则有限前缀
  \[
  \bigl(f_p(1),\dots,f_p(E_p)\bigr)
  \]
  已经充分且完全决定多重集 \(\{e_i(p)\}_{i=1}^{m}\)。
  \item 若 \(E_p\) 未知，则对所有矩阵的统一恢复而言，最短充分前缀恰为
  \[
  \bigl(f_p(1),\dots,f_p(E_p+1)\bigr).
  \]
  换言之，额外的 plateau 检测值 \(f_p(E_p+1)\) 既必要又充分。
  \item 设
  \[
  m_{p,\ell}:=\#\{i:e_i(p)=\ell\}
  \qquad(\ell\ge 1).
  \]
  则在 \(E_p\ge 1\) 时，有显式公式
  \[
  m_{p,\ell}
  =
  2f_p(\ell)-f_p(\ell-1)-f_p(\ell+1)
  \qquad(1\le \ell<E_p),
  \]
  \[
  m_{p,E_p}=f_p(E_p)-f_p(E_p-1).
  \]
  因而 \(p\)-主部分的全部正指数重数都由上述最短前缀显式给出；零指数重数则为
  \[
  m_{p,0}=m-\sum_{\ell=1}^{E_p}m_{p,\ell}.
  \]
\end{enumerate}
\end{theorem}
\begin{proof}
若 \(E_p=0\)，则全部 \(e_i(p)=0\)，此时 \(f_p(k)\equiv 0\)，结论平凡。以下设 \(E_p\ge 1\)。

由定理 \ref{thm:killo-smith-loss-spectrum}，
\[
\Delta_p(k):=f_p(k)-f_p(k-1)=\#\{i:e_i(p)\ge k\}
\qquad(k\ge 1).
\]
因此当 \(E_p\) 已知时，对 \(1\le \ell<E_p\) 有
\[
m_{p,\ell}
=
\Delta_p(\ell)-\Delta_p(\ell+1)
\]
\[
=
\bigl(f_p(\ell)-f_p(\ell-1)\bigr)-\bigl(f_p(\ell+1)-f_p(\ell)\bigr)
\]
\[
=
2f_p(\ell)-f_p(\ell-1)-f_p(\ell+1).
\]
而在顶层 \(\ell=E_p\) 处，由于不存在更大指数，故
\[
\Delta_p(E_p+1)=0,
\]
从而
\[
m_{p,E_p}=\Delta_p(E_p)=f_p(E_p)-f_p(E_p-1).
\]
于是所有 \(m_{p,\ell}\) 都由前缀
\[
\bigl(f_p(1),\dots,f_p(E_p)\bigr)
\]
显式恢复，进而整个多重集 \(\{e_i(p)\}\) 也被完全确定。这证明了第 \(1\) 条与第 \(3\) 条。

现证第 \(2\) 条的充分性。若只知道
\[
\bigl(f_p(1),\dots,f_p(E_p+1)\bigr),
\]
则可由其中最后一个平台位置恢复 \(E_p\) 本身：
\[
E_p=\max\{k\ge 1:\ f_p(k)>f_p(k-1)\}.
\]
事实上，对每个 \(1\le k\le E_p\)，至少存在一个 \(i\) 满足 \(e_i(p)\ge k\)，故
\[
\Delta_p(k)=\#\{i:e_i(p)\ge k\}\ge 1,
\]
即 \(f_p(k)>f_p(k-1)\)；而对 \(k=E_p+1\) 则
\[
\Delta_p(E_p+1)=0,
\qquad
f_p(E_p+1)=f_p(E_p).
\]
因此由长度 \(E_p+1\) 的前缀即可先识别 \(E_p\)，再由上一步的显式公式恢复全部 multiplicities，故其确为充分。

最后证必要性。若只给出到 \(E_p\) 为止的前缀，则无法统一地区分两个 \(p\)-指数多重集
\[
\{E_p\}
\qquad\text{与}\qquad
\{E_p+1\}.
\]
因为对所有 \(1\le k\le E_p\) 都有
\[
\min(k,E_p)=\min(k,E_p+1)=k,
\]
故二者给出完全相同的前缀
\[
\bigl(f_p(1),\dots,f_p(E_p)\bigr)=(1,2,\dots,E_p).
\]
只有再看一步
\[
f_p(E_p+1)
\]
才能区分平台是否已出现：第一种情形满足
\[
f_p(E_p+1)=f_p(E_p),
\]
第二种情形则有
\[
f_p(E_p+1)=f_p(E_p)+1.
\]
故在 \(E_p\) 未知的统一恢复问题中，长度 \(E_p+1\) 的前缀既必要又充分。证毕。
\end{proof}

\noindent
结合定理 \ref{thm:xi-natural-extension-h153-minimal-finite-alphabet} 可见，prime-slice 外置与自然扩张分支寄存器在资源层形成了两条完全一致的门槛律：有限深度塔的局部标签字母表大小精确等于逐层最大纤维数，而单寄存器自然扩张的全局最小字母表大小精确等于全局最大分支数。另一方面，定理 \ref{thm:xi-smith-p-kernel-shortest-sufficient-prefix} 又把 H.154 的整体可逆恢复进一步压缩为最短前缀律：若顶层指数未知，则额外的一步 plateau 检测正是唯一不可省去的信息。于是，局部可逆化资源与 \(p\)-核谱恢复复杂度在同一链上同时获得了精确的最小化刻画。

\paragraph{prime-slice 逆极限完成、Smith 线性投影可解性与整数椭圆的 profinite 可达性}
在定理 \ref{thm:killo-layered-prime-slices-finite-depth}、定理
\ref{thm:killo-natural-extension-branch-register}、定理
\ref{thm:killo-smith-loss-spectrum}、推论
\ref{cor:killo-kernel-size-general-modulus} 与定理
\ref{thm:killo-ellipse-diagonal-prime-register-equivalence} 已分别把有限深度 prime-slice 外置、自然扩张寄存器实现、Smith 核谱恢复、全模数核大小公式以及整数轴对齐椭圆的对角寄存器语义固定之后，这条链还可继续压缩为下列三条结果。它们分别给出有限深度前史的乘法型逆极限完成、一般整数线性投影的完整模 \(b\) 可解性与常纤维律，以及整数轴对齐椭圆在有限模与 profinite 极限中的精确可达性指数。

\begin{theorem}[有限深度前史的 prime-slice 乘法编码与严格逆极限完成]\label{thm:xi-time-part53ac-primeslice-truncated-history-inverse-limit}
\leanverified{paper\_xi\_time\_part53ac\_primeslice\_truncated\_history\_inverse\_limit}
设 \(X\) 为可数集合，\(T:X\twoheadrightarrow X\) 为有限纤维满射。对每个 \(k\ge 1\)，记
\[
X^{\leftarrow,k}
:=
\left\{
(x_0,x_{-1},\dots,x_{-k})\in X^{k+1}:
T(x_{-j})=x_{-j+1}\ \ (1\le j\le k)
\right\}.
\]
固定一列两两不交的 prime slices \((\mathcal P_j)_{j\ge 1}\)。则存在单射
\[
\Phi_k:X^{\leftarrow,k}\longrightarrow X\times \ZZ_{>0}
\]
满足
\[
\Phi_k(x_0,x_{-1},\dots,x_{-k})
=
\left(
x_0,\,
\prod_{j=1}^{k}q_j(x_0,x_{-1},\dots,x_{-k})
\right),
\qquad
q_j(\cdot)\in \mathcal P_j.
\]
若记
\[
\pi_{k+1,k}^{\leftarrow}(x_0,x_{-1},\dots,x_{-(k+1)})
:=
(x_0,x_{-1},\dots,x_{-k}),
\]
并定义
\[
\rho_{k+1,k}(n)
:=
\text{把 \(n\) 中属于 \(\mathcal P_{k+1}\) 的全部素因子剥去所得整数},
\]
则
\[
(\mathrm{id}_X\times \rho_{k+1,k})\circ \Phi_{k+1}
=
\Phi_k\circ \pi_{k+1,k}^{\leftarrow}.
\]
因此，若记
\[
\widehat X_k^{\mathrm{ps}}
:=
\Phi_k(X^{\leftarrow,k})\subset X\times \ZZ_{>0},
\]
则有规范同构
\[
\varprojlim_k \widehat X_k^{\mathrm{ps}}
\cong
\mathcal C(X^{\leftarrow})\subset X\times \NN^\NN,
\]
其中 \(\mathcal C\) 为定理 \ref{thm:killo-natural-extension-branch-register} 的自然扩张分支寄存器编码。

特别地，若对每个 \(1\le j\le k\) 都存在某个 \(y_j\in X\) 满足 \(\#T^{-1}(y_j)\ge 2\)，则深度 \(k\) 的精确 prime-slice 外置所需的最少有效 prime slices 数恰为 \(k\)。
\end{theorem}
\begin{proof}
对每个 \(y\in X\)，固定定理 \ref{thm:killo-natural-extension-branch-register} 中的纤维单射
\[
\iota_y:T^{-1}(y)\hookrightarrow \NN,
\qquad
\beta(x):=\iota_{T(x)}(x).
\]
由于 \(X\times \NN\) 可数而每个 \(\mathcal P_j\) 都是无限集，可取单射
\[
\alpha_j:X\times \NN\hookrightarrow \mathcal P_j
\qquad (j\ge 1).
\]
对任意
\[
\underline x=(x_0,x_{-1},\dots,x_{-k})\in X^{\leftarrow,k},
\]
定义
\[
q_j(\underline x)
:=
\alpha_j\bigl(x_{-(j-1)},\beta(x_{-j})\bigr)
\qquad (1\le j\le k),
\]
并令 \(\Phi_k\) 如定理中所示。

由于 prime slices 两两不交，唯一分解定理保证 \(\Phi_k(\underline x)\) 的第二坐标能够逐层恢复每个因子 \(q_j(\underline x)\)。先由
\[
\alpha_1^{-1}\bigl(q_1(\underline x)\bigr)
=
\bigl(x_0,\beta(x_{-1})\bigr)
\]
恢复 \(\beta(x_{-1})\)，再由 \(\iota_{x_0}^{-1}\) 恢复 \(x_{-1}\)；继而
\[
\alpha_2^{-1}\bigl(q_2(\underline x)\bigr)
=
\bigl(x_{-1},\beta(x_{-2})\bigr)
\]
恢复 \(x_{-2}\)，如此递推便可唯一恢复全部
\[
(x_{-1},\dots,x_{-k}).
\]
故 \(\Phi_k\) 为单射。

截断相容性直接来自定义：\(\pi_{k+1,k}^{\leftarrow}\) 删除最深层 \(x_{-(k+1)}\)，而 \(\rho_{k+1,k}\) 只删除属于 \(\mathcal P_{k+1}\) 的层因子 \(q_{k+1}\)，故
\[
(\mathrm{id}_X\times \rho_{k+1,k})\circ \Phi_{k+1}
=
\Phi_k\circ \pi_{k+1,k}^{\leftarrow}.
\]

于是得到自然映射
\[
\Theta:X^{\leftarrow}\longrightarrow \varprojlim_k \widehat X_k^{\mathrm{ps}},
\qquad
\Theta((x_0,x_{-1},x_{-2},\dots))
:=
\bigl(\Phi_k(x_0,x_{-1},\dots,x_{-k})\bigr)_{k\ge 1}.
\]
反过来，任取一个相容族
\[
\bigl((x^{(k)},N_k)\bigr)_{k\ge 1}\in \varprojlim_k \widehat X_k^{\mathrm{ps}}.
\]
由相容性可知 \(x^{(k)}\) 对 \(k\) 稳定，记其稳定值为 \(x_0\)。对每个 \(k\ge 1\)，由于
\[
(x_0,N_k)\in \widehat X_k^{\mathrm{ps}},
\]
存在唯一截断前史
\[
\underline x^{(k)}
:=
(x_0,x_{-1}^{(k)},\dots,x_{-k}^{(k)})\in X^{\leftarrow,k}
\]
满足
\[
\Phi_k(\underline x^{(k)})=(x_0,N_k).
\]
又因为
\[
(\mathrm{id}_X\times \rho_{k+1,k})(x_0,N_{k+1})=(x_0,N_k),
\]
结合前面已证的截断相容性可得
\[
\Phi_k\bigl(\pi_{k+1,k}^{\leftarrow}(\underline x^{(k+1)})\bigr)
=
(\mathrm{id}_X\times \rho_{k+1,k})\bigl(\Phi_{k+1}(\underline x^{(k+1)})\bigr)
=
(x_0,N_k)
=
\Phi_k(\underline x^{(k)}).
\]
由 \(\Phi_k\) 的单射性，
\[
\pi_{k+1,k}^{\leftarrow}(\underline x^{(k+1)})=\underline x^{(k)}.
\]
因此这些截断前史彼此兼容，定义出唯一无限前史
\[
(x_0,x_{-1},x_{-2},\dots)\in X^{\leftarrow}.
\]

对每个固定的 \(j\ge 1\)，整数 \(N_k\) 中来自 slice \(\mathcal P_j\) 的唯一素因子在 \(k\ge j\) 时保持不变，记作 \(q_j\in \mathcal P_j\)。再写
\[
\alpha_j^{-1}(q_j)=(u_{j-1},b_j)\in X\times \NN.
\]
按构造，对每个 \(j\ge 1\) 都有
\[
\alpha_j^{-1}(q_j)=\bigl(x_{-(j-1)},\beta(x_{-j})\bigr),
\]
故
\[
\mathcal C(x_0,x_{-1},x_{-2},\dots)
=
\bigl(x_0,(b_1,b_2,\dots)\bigr),
\]
而 \((b_j)_{j\ge 1}\) 与相容族一一对应，故这给出了
\[
\varprojlim_k \widehat X_k^{\mathrm{ps}}
\cong
\mathcal C(X^{\leftarrow})
\]
的规范双射。

最后，把定理 \ref{thm:xi-prime-slice-nontrivial-layer-exact-minimality} 应用于截断塔
\[
X^{\leftarrow,k}\to X^{\leftarrow,k-1}\to \cdots \to X^{\leftarrow,1}\to X,
\]
便知若每一层都存在真实分支，则任意精确 prime-slice 外置至少需要 \(k\) 个有效 slices。上面的构造恰在每一层使用一枚有效素数，故最小值正是 \(k\)。证毕。
\end{proof}

\begin{theorem}[Smith 正规形给出的模 \(b\) 线性投影可解性与常纤维律]\label{thm:xi-time-part53ac-smith-linear-solvability-constant-fibers}
\leanverified{paper\_xi\_time\_part53ac\_smith\_linear\_solvability\_constant\_fibers}
设 \(A\in M_{m\times n}(\ZZ)\) 满足 \(\operatorname{rank}_{\QQ}(A)=m\)，并取其 Smith 正规形
\[
UAV=\operatorname{diag}(d_1,\dots,d_m,0,\dots,0),
\qquad
d_1\mid d_2\mid \cdots \mid d_m,
\]
其中 \(U\in GL_m(\ZZ)\)、\(V\in GL_n(\ZZ)\)。对任意 \(b\ge 1\)，记
\[
A_b:(\ZZ/b\ZZ)^n\longrightarrow (\ZZ/b\ZZ)^m
\]
为模 \(b\) 诱导映射。若 \(y\in(\ZZ/b\ZZ)^m\) 且 \(y':=Uy\)，则同余方程
\[
Ax\equiv y\pmod b
\]
可解，当且仅当
\[
y_i'\in \gcd(d_i,b)\cdot (\ZZ/b\ZZ)
\qquad (1\le i\le m).
\]
等价地，取 \(y_i'\) 的任一整数代表后，有
\[
\gcd(d_i,b)\mid y_i'
\qquad (1\le i\le m).
\]
一旦可解，其解的个数与 \(y\) 无关，并且恰为
\[
\#\{x\in(\ZZ/b\ZZ)^n:\ Ax\equiv y\pmod b\}
=
b^{n-m}\prod_{i=1}^{m}\gcd(d_i,b).
\]
因此像集大小满足
\[
\#\operatorname{im}(A_b)
=
\frac{b^m}{\prod_{i=1}^{m}\gcd(d_i,b)}.
\]
\end{theorem}
\begin{proof}
由于 \(U,V\) 在模 \(b\) 后仍可逆，同余方程
\[
Ax\equiv y\pmod b
\]
等价于
\[
\operatorname{diag}(d_1,\dots,d_m,0,\dots,0)\,z\equiv y'\pmod b,
\qquad
z:=V^{-1}x.
\]
于是它分裂为
\[
d_i z_i\equiv y_i'\pmod b
\qquad (1\le i\le m),
\]
以及 \(n-m\) 个自由坐标 \(z_{m+1},\dots,z_n\)。

对单个一维同余
\[
d_i z_i\equiv y_i'\pmod b,
\]
标准初等数论给出：它可解，当且仅当
\[
y_i'\in \gcd(d_i,b)\cdot (\ZZ/b\ZZ),
\]
并且一旦可解，模 \(b\) 恰有 \(\gcd(d_i,b)\) 个解。故整个系统可解的充要条件正是逐坐标条件同时成立。若可解，则前 \(m\) 个坐标的解数为
\[
\prod_{i=1}^{m}\gcd(d_i,b),
\]
而后 \(n-m\) 个自由坐标各有 \(b\) 个取值，因此总解数恰为
\[
b^{n-m}\prod_{i=1}^{m}\gcd(d_i,b).
\]
这证明了解数公式。

最后，由于每个可达点的纤维大小都等于
\[
b^{n-m}\prod_{i=1}^{m}\gcd(d_i,b),
\]
故总点数分解为
\[
b^n
=
\#\operatorname{im}(A_b)\cdot
b^{n-m}\prod_{i=1}^{m}\gcd(d_i,b).
\]
整理即得
\[
\#\operatorname{im}(A_b)
=
\frac{b^m}{\prod_{i=1}^{m}\gcd(d_i,b)}.
\]
证毕。
\end{proof}

\begin{corollary}[整数轴对齐椭圆的模可达性判据与 profinite 指数]\label{cor:xi-time-part53ac-integer-ellipse-local-global-profinite-index}
\leanverified{paper\_xi\_time\_part53ac\_integer\_ellipse\_local\_global\_profinite\_index}
沿用定理 \ref{thm:killo-ellipse-diagonal-prime-register-equivalence} 的记号，并令
\[
A_{a,b}:=\operatorname{diag}(a,b).
\]
对任意 \(N\ge 1\) 与
\[
(u,v)\in (\ZZ/N\ZZ)^2,
\]
同余系统
\[
ax\equiv u\pmod N,
\qquad
by\equiv v\pmod N
\]
可解，当且仅当
\[
u\in \gcd(a,N)\cdot (\ZZ/N\ZZ),
\qquad
v\in \gcd(b,N)\cdot (\ZZ/N\ZZ).
\]
等价地，取整数代表后，
\[
\gcd(a,N)\mid u,
\qquad
\gcd(b,N)\mid v.
\]
一旦可解，解的个数恰为
\[
\gcd(a,N)\gcd(b,N).
\]
因此可达 residue 对的密度满足
\[
\delta_N(E_{a,b})
:=
\frac{\#\operatorname{im}(A_{a,b}\bmod N)}{N^2}
=
\frac{1}{\gcd(a,N)\gcd(b,N)}.
\]

进一步，在 profinite 完成上，
\[
A_{a,b}:\widehat{\ZZ}^2\longrightarrow \widehat{\ZZ}^2
\]
的像为
\[
\operatorname{im}(A_{a,b})
=
a\widehat{\ZZ}\times b\widehat{\ZZ},
\]
并且
\[
[\widehat{\ZZ}^2:\operatorname{im}(A_{a,b})]=ab.
\]
\end{corollary}
\begin{proof}
对角矩阵 \(A_{a,b}\) 使方程完全分裂为两个一维同余
\[
ax\equiv u\pmod N,
\qquad
by\equiv v\pmod N.
\]
每个方程分别可解，当且仅当右端属于对应的 \(\gcd\)-子群；一旦可解，其解数分别为 \(\gcd(a,N)\) 与 \(\gcd(b,N)\)。两式相乘便得到可解性判据与总解数公式。

于是可达 residue 对的总数为
\[
\frac{N}{\gcd(a,N)}\cdot \frac{N}{\gcd(b,N)},
\]
除以 \(N^2\) 即得密度公式。

对 profinite 版本，利用标准分解
\[
\widehat{\ZZ}\cong \prod_{p}\ZZ_p,
\]
乘以 \(a\) 的像正是
\[
a\widehat{\ZZ}=\prod_{p}p^{v_p(a)}\ZZ_p,
\]
其指数为
\[
\prod_{p}p^{v_p(a)}=a.
\]
同理，乘以 \(b\) 的像为 \(b\widehat{\ZZ}\)，指数为 \(b\)。因此
\[
\operatorname{im}(A_{a,b})=a\widehat{\ZZ}\times b\widehat{\ZZ},
\]
且
\[
[\widehat{\ZZ}^2:\operatorname{im}(A_{a,b})]
=
[\widehat{\ZZ}:a\widehat{\ZZ}]\,[\widehat{\ZZ}:b\widehat{\ZZ}]
=
ab.
\]
证毕。
\end{proof}

\noindent
由此，H.151--H.156 这一链条在乘法寄存器层面进一步闭合：有限深度前史不再只是分别编码为有限元组，而可被组织为单一整数寄存器中的严格逆极限系统；Smith 不变量也不再只控制核大小增长，而会完整决定模 \(b\) 线性投影的可达右端与常纤维大小；对整数轴对齐椭圆而言，这一结构又精确外显为有限模上的坐标 \(\gcd\)-可达性剖面与 profinite 极限中的闭开子群指数 \(ab\)。于是，自然扩张、prime-slice 外置、Smith 模核与整数椭圆的算术可达性便被压缩为同一条可审计的乘法编码链。


\paragraph{原子 Euler 因子的 ghost 支撑半群}
在单原子与有限多原子的 primitive 手术、ghost 级数闭式以及 Abel 有限部位移都已建立之后，还可把一般原子 Euler 因子在 periodic 层上的可见效应进一步压缩为一条纯支撑定律。

\begin{corollary}[一般原子 Euler 因子的 ghost 校正恰支撑在等差子半群]\label{cor:xi-time-part53aa-atomic-ghost-subsemigroup-support}
\leanverified{paper\_xi\_time\_part53aa\_atomic\_ghost\_subsemigroup\_support}
沿用定理 \ref{thm:xi-atomic-surgery-ghost-abel-shift-general} 的设定。设
\[
\zeta(z,u)=\frac{1}{(1-u^Ez^\ell)^m}\,\zeta_{\mathrm{core}}(z,u),
\qquad
m\in\ZZ_{\ge 1},\ \ell\ge 1,\ E\ge 0,
\]
并定义 periodic/ghost 系数
\[
\log \zeta(z,u)=\sum_{n\ge 1}\frac{a_n(u)}{n}z^n,
\qquad
\log \zeta_{\mathrm{core}}(z,u)=\sum_{n\ge 1}\frac{a_n^{\mathrm{core}}(u)}{n}z^n.
\]
则对每个 \(n\ge 1\) 都有显式分段公式
\[
\boxed{\;
a_n(u)
=
\begin{cases}
a_n^{\mathrm{core}}(u)+m\ell\,u^{En/\ell}, & \ell\mid n,\\[4pt]
a_n^{\mathrm{core}}(u), & \ell\nmid n.
\end{cases}\;}
\]
特别地，原子 Euler 因子对 ghost 层的全部修正恰好支撑在等差子半群
\[
\ell\NN=\{\ell,2\ell,3\ell,\dots\},
\]
而在其补集上，full/core 的 periodic 系数完全一致。
\end{corollary}
\begin{proof}
定理 \ref{thm:xi-atomic-surgery-ghost-abel-shift-general} 已给出精确公式
\[
a_n(u)-a_n^{\mathrm{core}}(u)
=
m\ell\,u^{En/\ell}\,\mathbf 1_{\{\ell\mid n\}}.
\]
将该式按 \(\ell\mid n\) 与 \(\ell\nmid n\) 两种情形拆开，即得到所示分段表达。因而一切 ghost 修正都严格集中在 \(\ell\NN\) 上，而其余长度不受原子因子影响。证毕。
\end{proof}

\noindent
因此，一般原子因子并不会对整个周期谱产生弥散型微扰；它只沿一个由长度 \(\ell\) 生成的等差子半群实施精确注入。


\paragraph{Rényi--2 熵亏损常数、\texorpdfstring{$BC_3$}{BC3} 根数据与 Kronecker 近共振分叉}
在 bin-fold 的共振碰撞缺口、window-\(6\) 的 \(B_3/C_3\) 根字典以及 Kronecker 分母刻度的 \(\mathsf{T}_q\) 闭式已经分别建立之后，还可进一步抽取出若干条横跨统计缺口、离散根数据与最优传输几何的解析后果。

\begin{theorem}[bin-fold 的不可消除 Rényi--2 熵亏损常数]\label{thm:xi-binfold-renyi2-entropy-defect-constant}
设 \(\mu_m\) 为 bin-fold 输出分布，
\[
\mathsf{Col}_m:=\sum_{w\in X_m}\mu_m(w)^2,
\qquad
H_2(\mu_m):=-\log \mathsf{Col}_m,
\qquad
H_2^{(2)}(\mu_m):=-\log_2 \mathsf{Col}_m.
\]
记 \(u_m\) 为 \(X_m\) 上的均匀分布，并令 \(C_\varphi\) 为定理 \ref{thm:fold-critical-resonance-constant} 的黄金共振常数。则
\[
\log F_{m+2}-H_2(\mu_m)
=
D_2(\mu_m\Vert u_m)
\ge
\log\!\bigl(1+2C_\varphi^2\bigr)+o(1).
\]
特别地，
\[
\log F_{m+2}-H_2(\mu_m)\ge 0.0959706185\ldots+o(1)\ \text{\rm(nat)},
\]
以及
\[
\log_2 F_{m+2}-H_2^{(2)}(\mu_m)\ge 0.1384563354\ldots+o(1)\ \text{\rm(bit)}.
\]
因此 bin-fold 的二阶信息损失不是随尺度消失的边界效应，而是一个不可消去的常数级熵亏损。
\end{theorem}
\begin{proof}
由定理 \ref{thm:xi-time-part9p-renyi2-obstruction}，
\[
D_2(\mu_m\Vert u_m)=\log\!\bigl(F_{m+2}\mathsf{Col}_m\bigr)
\ge
\log\!\bigl(1+2C_\varphi^2\bigr)+o(1).
\]
另一方面，
\[
\log\!\bigl(F_{m+2}\mathsf{Col}_m\bigr)
=
\log F_{m+2}-H_2(\mu_m).
\]
于是得到第一式。再除以 \(\log 2\) 即得以 bit 计的版本。数值常数则由
\(\log(1+2C_\varphi^2)\) 与 \(\log_2(1+2C_\varphi^2)\) 直接给出。证毕。
\end{proof}

\begin{corollary}[谱零点绝对规模与碰撞缺口主尺度的解耦]\label{cor:xi-binfold-zero-block-collision-gap-decoupling}
若 \(m+2\equiv 0\pmod 6\)，则
\[
|\mathcal Z_m|\ge F_{(m+2)/2},
\qquad
\frac{|\mathcal Z_m|}{F_{m+2}}
\le
\tau(m+2)\,\frac{F_{\lfloor (m+2)/2\rfloor}}{F_{m+2}}
=
\tau(m+2)\,O(\varphi^{-m/2}).
\]
在最小三轴窗 \(m=58\) 上更有精确值
\[
|\mathcal Z_{58}|=839415>F_{30}=832040.
\]
然而定理 \ref{thm:xi-binfold-renyi2-entropy-defect-constant} 同时给出
\[
\mathsf{Col}_m-\frac{1}{F_{m+2}}
\ge
\frac{2C_\varphi^2+o(1)}{F_{m+2}}.
\]
因此，即使谱零点集合在绝对规模上已经达到 \(\asymp \varphi^{m/2}\) 的指数级，其本身仍不足以决定碰撞缺口的首项常数；后者的主尺度依旧是 \(F_{m+2}^{-1}\)。
\end{corollary}
\begin{proof}
六倍窗下的半阶零块下界由结论 \ref{con:xi-fold-window6-kernel-half-index-lower-bound} 给出，而相对密度上界由定理 \ref{thm:fold-zero-density-sparse} 给出。\(m=58\) 的精确计数则见注 \ref{rem:fold-zero-lb-m58}。
另一方面，由定理 \ref{thm:xi-binfold-renyi2-entropy-defect-constant} 的等价改写，
\[
\mathsf{Col}_m-\frac{1}{F_{m+2}}
=
\frac{e^{D_2(\mu_m\Vert u_m)}-1}{F_{m+2}}
\ge
\frac{2C_\varphi^2+o(1)}{F_{m+2}}.
\]
于是零点集合可以在绝对规模上指数增大，而碰撞缺口的主项仍由 \(F_{m+2}^{-1}\) 这一更细尺度控制。证毕。
\end{proof}

\begin{theorem}[Fold\(_6\) 的刚性循环核给出规范的 \texorpdfstring{$BC_3$}{BC3} 根数据]\label{thm:xi-window6-cyclic-kernel-bc3-root-datum}
沿用附录 \ref{app:fold6-b3c3-seed-dictionary} 的刚性分裂
\[
X_6=X_6^{\mathrm{cyc}}\sqcup X_6^{\mathrm{bdry}},
\qquad
|X_6|=21,\quad |X_6^{\mathrm{cyc}}|=18,\quad |X_6^{\mathrm{bdry}}|=3.
\]
记
\[
X_{6,\mathrm{ax}}
:=
\{w\in X_6^{\mathrm{cyc}}:\ \mathrm{wt}(w)=1\},
\qquad
X_{6,\mathrm{off}}
:=
X_6^{\mathrm{cyc}}\setminus X_{6,\mathrm{ax}}.
\]
则 \(|X_{6,\mathrm{ax}}|=6\)、\(|X_{6,\mathrm{off}}|=12\)，并且：
\begin{enumerate}
  \item 在表 \ref{tab:fold6_b3c3_root_dictionary_b3} 的 \(B_3\) 规范下，\(X_{6,\mathrm{ax}}\) 恰对应
  \(\{\pm e_1,\pm e_2,\pm e_3\}\)，而 \(X_{6,\mathrm{off}}\) 恰对应
  \(\{\pm e_i\pm e_j:\ 1\le i<j\le 3\}\)；
  \item 在表 \ref{tab:fold6_b3c3_root_dictionary_c3} 的 \(C_3\) 规范下，\(X_{6,\mathrm{off}}\) 的对应向量完全不变，而 \(X_{6,\mathrm{ax}}\) 统一径向放大为
  \(\{\pm 2e_1,\pm 2e_2,\pm 2e_3\}\)。
\end{enumerate}
因此 \(X_6^{\mathrm{cyc}}\) 所携带的本征对象不是预先固定的 \(B_3\) 或 \(C_3\) 归一化，而是一组规范的 \(BC_3\) 根数据；\(B_3\) 与 \(C_3\) 只是对同一 \(6+12\) 组合壳层施加的两种径向归一化。
\end{theorem}
\begin{proof}
附录 \ref{app:fold6-b3c3-seed-dictionary} 给出
\[
|X_6^{\mathrm{bdry}}|=3,
\qquad
\mathrm{wt}\text{-hist}(X_6^{\mathrm{cyc}})=\{0:1,\ 1:6,\ 2:9,\ 3:2\},
\]
故 \(\mathrm{wt}(w)=1\) 的循环点恰有 \(6\) 个，其余循环点恰有 \(12\) 个。再由表
\ref{tab:fold6_b3c3_root_dictionary_b3} 逐项可见：这 \(6\) 个点正好对应 \(B_3\) 的六个轴根
\(\{\pm e_i\}\)，而其余 \(12\) 个点正好对应 \(B_3\) 的十二个非轴根
\(\{\pm e_i\pm e_j\}\)。

同样比较表 \ref{tab:fold6_b3c3_root_dictionary_c3} 可见：十二个非轴根的向量完全保持不变，而六个轴向点统一由 \(\pm e_i\) 改写为 \(\pm 2e_i\)。因此在两张字典之间真正不变的，是“六条轴向方向 \(+\) 十二条非轴向方向”的组合壳层，而非某一固定的根长归一化。该不变数据正是 \(BC_3\) 口径下的规范根数据；\(B_3\) 与 \(C_3\) 只是对同一方向数据所作的两种径向标定。证毕。
\end{proof}

\begin{theorem}[Fold\(_6\) 的 \texorpdfstring{$C_3$}{C3} 字典中 Hamming 权恰实现长短根分裂]\label{thm:xi-window6-c3-hamming-root-length-split}
沿用定理 \ref{thm:xi-window6-cyclic-kernel-bc3-root-datum} 的记号，并令
\[
\iota_C:X_6^{\mathrm{cyc}}\longrightarrow \Phi(C_3)
\]
表示表 \ref{tab:fold6_b3c3_root_dictionary_c3} 给出的 \(C_3\) 根字典。则有精确识别
\[
\iota_C\bigl(\{w\in X_6^{\mathrm{cyc}}:\ \mathrm{wt}(w)=1\}\bigr)
=
\{(\pm 2,0,0),(0,\pm 2,0),(0,0,\pm 2)\},
\]
\[
\iota_C\bigl(\{w\in X_6^{\mathrm{cyc}}:\ \mathrm{wt}(w)\neq 1\}\bigr)
=
\{(\pm 1,\pm 1,0),(\pm 1,0,\pm 1),(0,\pm 1,\pm 1)\}.
\]
特别地，在该 \(C_3\) 规范下，
\[
\mathrm{wt}(w)=1
\iff
\iota_C(w)\ \text{是长根},
\]
\[
\mathrm{wt}(w)\neq 1
\iff
\iota_C(w)\ \text{是短根}.
\]
因此 Hamming 权在该字典中并非附属标签，而恰好是区分两条根长 Weyl 轨道的纯组合判别式。
\leanverified{paper\_xi\_window6\_c3\_hamming\_root\_length\_split}
\end{theorem}
\begin{proof}
由定理 \ref{thm:xi-window6-cyclic-kernel-bc3-root-datum}，
\[
X_{6,\mathrm{ax}}
:=
\{w\in X_6^{\mathrm{cyc}}:\ \mathrm{wt}(w)=1\}
\]
恰有 \(6\) 个元素，并满足
\[
\iota_C(X_{6,\mathrm{ax}})
=
\{\pm 2e_1,\pm 2e_2,\pm 2e_3\}.
\]
同一定理还给出其补集
\[
X_{6,\mathrm{off}}:=X_6^{\mathrm{cyc}}\setminus X_{6,\mathrm{ax}}
\]
恰有 \(12\) 个元素，并满足
\[
\iota_C(X_{6,\mathrm{off}})
=
\{\pm e_i\pm e_j:\ 1\le i<j\le 3\}.
\]
将这些向量按坐标展开，便得到题设中的两组显式列表。最后，标准 \(C_3\) 根系中 \(\{\pm 2e_i\}\) 正是长根轨道，而 \(\{\pm e_i\pm e_j\}\) 正是短根轨道，故上述两条等价关系立即成立。证毕。
\end{proof}

\begin{theorem}[Kronecker 近共振窗中的 \texorpdfstring{$W_1$}{W1} 唯一极小点严格位于有理逼近的右侧]\label{thm:xi-kronecker-w1-right-shifted-minimizer}
设 \(p/q\) 为既约有理数，\(q\ge 2\)，并写
\[
\alpha=\frac{p}{q}+\delta,
\qquad
|\delta|<\frac{1}{q(q-1)}.
\]
令
\[
P_q(\alpha):=\{\{t\alpha\}:0\le t\le q-1\},
\qquad
\mu_q:=\frac{1}{q}\sum_{x\in P_q(\alpha)}\delta_x,
\qquad
\mathsf{T}_q(\alpha):=W_1(\mu_q,\lambda).
\]
则
\[
\mathsf{T}_q(\alpha)=\frac{1}{2q}+\frac{q-1}{2}\Bigl(\frac{p}{q}-\alpha\Bigr)
\qquad(\delta<0),
\]
而
\[
\mathsf{T}_q(\alpha)=\frac{1}{2q}-\frac{q-1}{2}\delta+\frac{q(q-1)(2q-1)}{6}\delta^2
\qquad(\delta>0).
\]
特别地，在右侧窗口 \(\delta>0\) 内，\(\mathsf{T}_q\) 存在唯一极小点
\[
\delta_*=\frac{3}{2q(2q-1)},
\qquad
\alpha_*=\frac{p}{q}+\frac{3}{2q(2q-1)},
\]
并且
\[
\mathsf{T}_q(\alpha_*)=\frac{5q-1}{8q(2q-1)}.
\]
因此最近收敛分数 \(p/q\) 本身并不是 Wasserstein 几何的极小点；真正的最优点严格位于其右侧。
\leanverified{paper\_xi\_kronecker\_w1\_right\_shifted\_minimizer}
\end{theorem}
\begin{proof}
条件 \(|\delta|<1/(q(q-1))\) 等价于
\[
(q-1)|\delta|<\frac{1}{q},
\]
这正保证每个点 \(\{t\alpha\}\) 在长度 \(1/q\) 的分箱中至多发生箱内位移而不越阶。因此定理 \ref{thm:kronecker-w1-denominator-closed-form} 的同一分箱求和论证在该更大窗口内仍然适用，得到上述左右两侧的闭式。

左侧 \(\delta<0\) 时，\(\mathsf{T}_q\) 对 \(\delta\) 为仿射函数，故不存在内部极小点。右侧 \(\delta>0\) 时，
\[
\mathsf{T}_q'(\delta)
=
-\frac{q-1}{2}+\frac{q(q-1)(2q-1)}{3}\delta,
\]
且二次项系数 \(q(q-1)(2q-1)/6\) 严格为正，故 \(\mathsf{T}_q\) 为严格凸二次式。令导数为零得到唯一驻点
\[
\delta_*=\frac{3}{2q(2q-1)}.
\]
又
\[
\delta_*-\frac{1}{q(q-1)}
=
-\frac{q+1}{2q(q-1)(2q-1)}<0,
\]
故该驻点确实落在允许的右侧窗口之内，因此它就是唯一极小点。把 \(\delta_*\) 代回右侧闭式即可得到
\[
\mathsf{T}_q(\alpha_*)=\frac{5q-1}{8q(2q-1)}.
\]
证毕。
\end{proof}

\begin{corollary}[星差异在右侧近共振窗中失明而 \texorpdfstring{$W_1$}{W1} 保持严格可分辨性]\label{cor:xi-kronecker-star-discrepancy-w1-branching}
\leanverified{paper\_xi\_kronecker\_star\_discrepancy\_w1\_branching}
沿用定理 \ref{thm:xi-kronecker-w1-right-shifted-minimizer} 的记号。则
\[
D_q^\ast(P_q(\alpha))
=
\frac{1}{q}+(q-1)\Bigl(\frac{p}{q}-\alpha\Bigr),
\qquad
\mathsf{T}_q(\alpha)=\frac12\,D_q^\ast(P_q(\alpha))
\qquad(\delta<0),
\]
而
\[
D_q^\ast(P_q(\alpha))=\frac{1}{q}
\qquad(\delta>0),
\]
同时 \(\mathsf{T}_q(\alpha)\) 仍按定理 \ref{thm:xi-kronecker-w1-right-shifted-minimizer} 的严格凸二次式发生变化。特别地，在最优点 \(\alpha_*\) 上，
\[
D_q^\ast(P_q(\alpha_*))=\frac{1}{q},
\qquad
\frac{2\mathsf{T}_q(\alpha_*)}{D_q^\ast(P_q(\alpha_*))}
=
\frac{5q-1}{4(2q-1)}
\longrightarrow \frac{5}{8}.
\]
因此在右侧最近收敛窗中，星差异对 \(\alpha\) 完全失明，而 \(\mathsf{T}_q\) 仍保留严格曲率与唯一极小点；两种审计量在该窗内发生结构性分叉。
\end{corollary}
\begin{proof}
当 \(\delta<0\) 时，仍在同一不越阶分箱条件下，首箱出现双占、末箱出现空缺；与定理 \ref{thm:kronecker-w1-denominator-closed-form} 的计数相同，立得
\[
D_q^\ast(P_q(\alpha))
=
\frac{1}{q}+(q-1)\Bigl(\frac{p}{q}-\alpha\Bigr),
\]
并满足 \(\mathsf{T}_q(\alpha)=\tfrac12 D_q^\ast(P_q(\alpha))\)。
当 \(\delta>0\) 时，每个箱 \([j/q,(j+1)/q)\) 恰含一点，故
\(
D_q^\ast(P_q(\alpha))=1/q
\)
为常数，而 \(\mathsf{T}_q(\alpha)\) 依定理 \ref{thm:xi-kronecker-w1-right-shifted-minimizer} 仍随 \(\delta\) 呈严格凸二次变化。

最后把定理 \ref{thm:xi-kronecker-w1-right-shifted-minimizer} 给出的
\[
\mathsf{T}_q(\alpha_*)=\frac{5q-1}{8q(2q-1)}
\]
代入即可得到比值公式。证毕。
\end{proof}

\begin{corollary}[好侧单箱室中的严格凸传输几何与星差异非完备性]\label{cor:xi-kronecker-good-side-convexity-star-incompleteness}
\leanverified{paper\_xi\_kronecker\_good\_side\_convexity\_star\_incompleteness}
沿用定理 \ref{thm:xi-kronecker-w1-right-shifted-minimizer} 与推论
\ref{cor:xi-kronecker-star-discrepancy-w1-branching} 的记号，并限制到好侧单箱室
\[
0<\delta<\frac{1}{q(q-1)}.
\]
则在该室内有精确闭式
\[
\mathsf{T}_q(\alpha)
=
\frac{1}{2q}
-\frac{q-1}{2}\,\delta
+\frac{q(q-1)(2q-1)}{6}\,\delta^2,
\]
同时
\[
D_q^\ast(P_q(\alpha))=\frac{1}{q}.
\]
因此：
\begin{enumerate}
  \item \(\mathsf{T}_q(\alpha)\) 在该室内是严格凸二次函数；
  \item 其唯一极小点为
  \[
  \delta_\ast=\frac{3}{2q(2q-1)};
  \]
  \item 极小值可写为
  \[
  \boxed{\;
  \mathsf{T}_q^{\min}
  =
  \frac{5q-1}{8q(2q-1)};\;}
  \]
  \item \(\mathsf{T}_q\) 在该室内至少扫描非退化区间
  \[
  \left[\frac{5q-1}{8q(2q-1)},\frac{1}{2q}\right),
  \]
  而 \(D_q^\ast\) 保持常值 \(1/q\)。因而在该局部室中，\(\mathsf{T}_q\) 不可能是 \(D_q^\ast\) 的任何函数。
\end{enumerate}
\end{corollary}
\begin{proof}
第一条闭式由定理 \ref{thm:xi-kronecker-w1-right-shifted-minimizer} 在 \(\delta>0\) 的分支直接给出，而
\[
D_q^\ast(P_q(\alpha))=\frac{1}{q}
\]
则由推论 \ref{cor:xi-kronecker-star-discrepancy-w1-branching} 给出。

对 \(\delta\) 求导得到
\[
\frac{d}{d\delta}\mathsf{T}_q(\alpha)
=
-\frac{q-1}{2}
+\frac{q(q-1)(2q-1)}{3}\delta,
\]
\[
\frac{d^2}{d\delta^2}\mathsf{T}_q(\alpha)
=
\frac{q(q-1)(2q-1)}{3}>0,
\]
故 \(\mathsf{T}_q\) 在好侧单箱室内严格凸，且唯一驻点即唯一极小点为
\[
\delta_\ast=\frac{3}{2q(2q-1)}.
\]
将 \(\delta_\ast\) 代回二次式，得到
\[
\mathsf{T}_q^{\min}
=
\frac{1}{2q}
-\frac{3(q-1)}{8q(2q-1)}
=
\frac{5q-1}{8q(2q-1)}.
\]

最后，由
\[
\lim_{\delta\downarrow 0}\mathsf{T}_q(\alpha)=\frac{1}{2q}
\]
以及 \(\delta=\delta_\ast\) 处达到唯一极小值，可知 \(\mathsf{T}_q\) 在该室内至少扫描所示非退化区间；而 \(D_q^\ast\) 在同一室内恒等于 \(1/q\)。若存在函数 \(\Phi\) 使 \(\mathsf{T}_q=\Phi(D_q^\ast)\)，则 \(\mathsf{T}_q\) 亦应为常数，矛盾。证毕。
\end{proof}

\paragraph{统计缺口、离散根数据与近共振传输的同层分叉}
上述结果表明：在 bin-fold 模块中，谱零点的绝对增殖与二阶统计缺口并不处于同一层级；在 window-\(6\) 模块中，循环核先天承载的是 \(BC_3\) 型组合根数据，而在 \(C_3\) 规范下 \(\mathrm{wt}(w)=1\) 又恰纯组合地区分出长根轨道，\(B_3/C_3\) 只是在此基础上施加的两种径向规范；在 Kronecker 近共振模块中，右侧窗口则出现“星差异平坦而 \(\mathsf{T}_q\) 仍具曲率”的严格分叉，并进一步表现为同一常值星差异之下的严格凸二次传输几何。三者合并之后，共振、零点、根字典与最优传输读出被压缩为一组可彼此对照的刚性接口。
\paragraph{多窗导出交叉、倍窗 Lucas 伴随扩张与 window-\texorpdfstring{$6$}{6} 根字典的外在性}
在折叠同余商的单模 Fibonacci 环表示、跨分辨率约化兼容性以及 window-\(6\) 的 \(BC_3\) 根数据都已固定之后，还可进一步抽取出下列跨窗口结论。它们分别把多窗分辨率的 gcd 型导出交叉、倍窗层的 Lucas 伴随扩张、window-\(14\) 对 window-\(6\) 核心的直因子承载，以及 \(B_3/C_3\) 根字典对商环算术的不可内化性统一到同一条有限算术几何主线上。

\begin{theorem}[多窗折叠同余商的导出交叉公式]\label{thm:xi-time-part51aa-derived-cross-fib-gcd-formula}
沿用附录定理 \ref{thm:fold-congruence-mod-semirings} 的记号，并记
\[
R_m:=\Omega_m/\!\equiv_m\cong \ZZ/(F_{m+2}\ZZ).
\]
给定 \(m_1,\dots,m_r\ge 1\)，令
\[
N_i:=F_{m_i+2},
\qquad
g:=\gcd(m_1+2,\dots,m_r+2).
\]
则在有界导出范畴 \(D^b(\ZZ)\) 中有
\[
H_q\!\left(
R_{m_1}\overset{\mathbf L}{\otimes}_{\ZZ}\cdots
\overset{\mathbf L}{\otimes}_{\ZZ}R_{m_r}
\right)
\cong
\bigl(\ZZ/(F_g\ZZ)\bigr)^{\oplus \binom{r-1}{q}}
\qquad(0\le q\le r-1),
\]
并且其余同调次数上的同调群均为零。特别地，当 \(g=1\) 时，整个导出张量积是零对象。
\end{theorem}
\begin{proof}
对每个 \(1\le i\le r\)，取两项自由分辨率
\[
K(N_i):=[\ZZ\xrightarrow{N_i}\ZZ],
\]
其中非零项放在同调次数 \(1\) 与 \(0\)。于是
\[
R_{m_1}\overset{\mathbf L}{\otimes}_{\ZZ}\cdots
\overset{\mathbf L}{\otimes}_{\ZZ}R_{m_r}
\]
可由 Koszul 复形
\[
K(N_1,\dots,N_r):=
K(N_1)\otimes_{\ZZ}\cdots\otimes_{\ZZ}K(N_r)
\]
表示。该复形可视为自由模 \(\Lambda^\bullet(\ZZ^r)\) 上由行向量
\[
v:=(N_1,\dots,N_r)
\]
定义的收缩微分。

对这条 \(1\times r\) 整数行向量作 Smith 约化，存在
\[
U\in \mathrm{GL}_r(\ZZ)
\]
使
\[
vU=(d,0,\dots,0),
\qquad
d:=\gcd(N_1,\dots,N_r).
\]
基变换 \(U\) 在 \(\Lambda^\bullet(\ZZ^r)\) 上诱导复形同构
\[
K(N_1,\dots,N_r)\cong K(d,0,\dots,0).
\]
另一方面，由 Fibonacci 最大公因子恒等式逐次迭代得到
\[
d
=
\gcd(F_{m_1+2},\dots,F_{m_r+2})
=
F_g.
\]

最后，
\[
K(d,0,\dots,0)=K(d)\otimes_{\ZZ}K(0)^{\otimes(r-1)}.
\]
其中 \(K(d)=[\ZZ\xrightarrow{d}\ZZ]\) 的唯一非零同调为
\[
H_0(K(d))\cong \ZZ/(d\ZZ),
\]
而
\[
K(0)=[\ZZ\xrightarrow{0}\ZZ]
\]
在次数 \(0\) 与 \(1\) 各有一份 \(\ZZ\)。因此
\[
K(0)^{\otimes(r-1)}
\]
在同调次数 \(q\) 上恰有
\[
\binom{r-1}{q}
\]
份 \(\ZZ\)。张量上 \(K(d)\) 后便得到
\[
H_q\!\bigl(K(N_1,\dots,N_r)\bigr)
\cong
\bigl(\ZZ/(d\ZZ)\bigr)^{\oplus \binom{r-1}{q}}
\qquad(0\le q\le r-1),
\]
且其余次数为零。代入 \(d=F_g\) 即得结论。若 \(g=1\)，则 \(d=F_1=1\)，此时 \(K(1)\) 为恰当复形，故整个导出张量积为零对象。证毕。
\end{proof}

\begin{corollary}[二窗折叠同余商的导出正交性具有正密度]\label{cor:xi-time-part51aa-derived-orthogonality-positive-density}
\leanverified{paper\_xi\_time\_part51aa\_derived\_orthogonality\_positive\_density}
沿用定理 \ref{thm:xi-time-part51aa-derived-cross-fib-gcd-formula} 的记号。对任意 \(m,n\ge 1\)，有
\[
R_m\overset{\mathbf L}{\otimes}_{\ZZ}R_n\simeq 0
\quad\Longleftrightarrow\quad
\gcd(m+2,n+2)=1.
\]
等价地，
\[
R_m\otimes_{\ZZ}R_n=0,
\qquad
\operatorname{Tor}_1^{\ZZ}(R_m,R_n)=0
\]
当且仅当 \(\gcd(m+2,n+2)=1\)。进一步，
\[
\frac{1}{X^2}
\#\Bigl\{
1\le m,n\le X:\ 
R_m\overset{\mathbf L}{\otimes}_{\ZZ}R_n\simeq 0
\Bigr\}
\longrightarrow
\frac{6}{\pi^2}
\qquad(X\to\infty).
\]
因此，按 \(\max\{m,n\}\le X\) 排序的全部窗口对中，导出正交对具有自然密度 \(6/\pi^2\)。
\end{corollary}
\begin{proof}
取定理 \ref{thm:xi-time-part51aa-derived-cross-fib-gcd-formula} 中 \(r=2\)。此时
\[
H_0\!\left(R_m\overset{\mathbf L}{\otimes}_{\ZZ}R_n\right)
\cong
H_1\!\left(R_m\overset{\mathbf L}{\otimes}_{\ZZ}R_n\right)
\cong
\ZZ/(F_g\ZZ),
\qquad
g:=\gcd(m+2,n+2).
\]
故导出张量积为零对象当且仅当 \(F_g=1\)，也即 \(g=1\)。而
\[
H_0\cong R_m\otimes_{\ZZ}R_n,
\qquad
H_1\cong \operatorname{Tor}_1^{\ZZ}(R_m,R_n),
\]
遂得等价表述。

最后，平移 \((m,n)\mapsto (m+2,n+2)\) 不改变互素对的自然密度，而经典计数给出
\[
\frac{1}{X^2}\#\{1\le a,b\le X:\ \gcd(a,b)=1\}\longrightarrow \frac{6}{\pi^2}.
\]
于是所示极限成立。证毕。
\end{proof}

\begin{theorem}[倍窗折叠同余商的 Lucas 伴随扩张]\label{thm:xi-time-part51aa-doubling-lucas-extension}
\leanverified{paper\_xi\_time\_part51aa\_doubling\_lucas\_extension}
沿用
\[
R_m:=\Omega_m/\!\equiv_m\cong \ZZ/(F_{m+2}\ZZ)
\]
的记号。固定 \(m\ge 1\)，并记
\[
\nu:=m+2,
\qquad
S_m:=R_{2m+2}\cong \ZZ/(F_{2\nu}\ZZ).
\]
再定义理想
\[
I_\nu:=F_\nu S_m\subset S_m.
\]
则存在短正合列
\[
0\longrightarrow I_\nu\longrightarrow S_m\longrightarrow R_m\longrightarrow 0,
\]
并且成立：
\begin{enumerate}
  \item \(I_\nu\) 作为加法群满足
  \[
  I_\nu\cong \ZZ/(L_\nu\ZZ),
  \qquad
  |I_\nu|=L_\nu;
  \]
  \item 其余切缺陷模满足
  \[
  I_\nu/I_\nu^2\cong \ZZ/\gcd(F_\nu,L_\nu)\ZZ;
  \]
  \item 更进一步，
  \[
  \gcd(F_\nu,L_\nu)=\gcd(F_\nu,2)=
  \begin{cases}
  1,& 3\nmid \nu,\\
  2,& 3\mid \nu.
  \end{cases}
  \]
  因而该扩张当且仅当 \(3\nmid \nu\) 时分裂；在此情形下有
  \[
  R_{2m+2}\cong R_m\times \ZZ/(L_\nu\ZZ).
  \]
  而当 \(3\mid \nu\) 时，扩张非分裂，且全部分裂障碍恰压缩为一维 \(\FF_2\)-级的余切缺陷
  \[
  I_\nu/I_\nu^2\cong \FF_2.
  \]
\end{enumerate}
\end{theorem}
\begin{proof}
由定理 \ref{thm:fold-congruence-mod-semirings}，
\[
S_m\cong \ZZ/(F_{2\nu}\ZZ).
\]
再由 Fibonacci doubling 恒等式
\[
F_{2\nu}=F_\nu(F_{\nu-1}+F_{\nu+1})=:F_\nu L_\nu,
\]
可得
\[
S_m\cong \ZZ/(F_\nu L_\nu\ZZ).
\]
对模 \(F_\nu\) 取商，得到满同态
\[
S_m\twoheadrightarrow \ZZ/(F_\nu\ZZ)\cong R_m,
\]
其核正是
\[
I_\nu=F_\nu S_m.
\]
于是短正合列成立。

定义映射
\[
\ZZ/(L_\nu\ZZ)\longrightarrow I_\nu,
\qquad
\overline t\longmapsto \overline{F_\nu t}.
\]
它显然是加法同态，且由于
\[
F_\nu t\equiv 0\pmod{F_\nu L_\nu}
\iff
L_\nu\mid t,
\]
故其亦为双射，从而
\[
I_\nu\cong \ZZ/(L_\nu\ZZ)
\]
作为加法群成立，并给出 \(|I_\nu|=L_\nu\)。

再看平方理想。由 \(I_\nu=F_\nu S_m\) 得
\[
I_\nu^2=F_\nu I_\nu.
\]
在上面的加法同构下，这对应于 \(\ZZ/(L_\nu\ZZ)\) 上乘以 \(F_\nu\) 的像，因此
\[
I_\nu/I_\nu^2
\cong
\bigl(\ZZ/(L_\nu\ZZ)\bigr)\big/F_\nu\bigl(\ZZ/(L_\nu\ZZ)\bigr)
\cong
\ZZ/\gcd(F_\nu,L_\nu)\ZZ.
\]

又由于
\[
L_\nu=F_{\nu-1}+F_{\nu+1}=2F_{\nu-1}+F_\nu,
\]
可得
\[
\gcd(F_\nu,L_\nu)
=
\gcd(F_\nu,2F_{\nu-1})
=
\gcd(F_\nu,2),
\]
因为 \(\gcd(F_\nu,F_{\nu-1})=1\)。而 Fibonacci 数列模 \(2\) 的周期为
\[
0,1,1,
\]
故 \(F_\nu\) 为偶数当且仅当 \(3\mid \nu\)。于是得到所示的二值判别。

若 \(3\nmid \nu\)，则 \(\gcd(F_\nu,L_\nu)=1\)，由中国剩余定理即得
\[
\ZZ/(F_\nu L_\nu\ZZ)\cong \ZZ/(F_\nu\ZZ)\times \ZZ/(L_\nu\ZZ),
\]
亦即扩张分裂。

反之，若 \(3\mid \nu\)，则
\[
I_\nu/I_\nu^2\cong \ZZ/(2\ZZ)=\FF_2\neq 0.
\]
若扩张分裂，则存在某个环 \(B\) 使
\[
S_m\cong R_m\times B,
\qquad
I_\nu\cong 0\times B.
\]
但此时
\[
I_\nu^2\cong 0\times B^2=0\times B\cong I_\nu,
\]
从而 \(I_\nu/I_\nu^2=0\)，矛盾。故 \(3\mid \nu\) 时扩张非分裂，且其全部障碍正由上式中的 \(\FF_2\) 缺陷刻画。证毕。
\end{proof}

\begin{corollary}[window-\texorpdfstring{$14$}{14} 对 window-\texorpdfstring{$6$}{6} 的严格直因子提升]\label{cor:xi-time-part51aa-window14-strict-direct-factor-lift}
取 \(m=6\)。则
\[
\nu=8,
\qquad
F_8=21,
\qquad
L_8=47,
\qquad
3\nmid 8.
\]
因此
\[
R_{14}\cong R_6\times \FF_{47}.
\]
并且
\[
H_0\!\left(R_6\overset{\mathbf L}{\otimes}_{\ZZ}R_{14}\right)
\cong
H_1\!\left(R_6\overset{\mathbf L}{\otimes}_{\ZZ}R_{14}\right)
\cong
R_6.
\]
因此，window-\(14\) 不是对 window-\(6\) 的单纯模数增厚，而是在保留 \(21\)-点算术核心的同时附加一个新的 \(47\)-元域因子；在导出意义下，这一 \(21\)-点核心被完整地保留在其倍窗后代中。
\end{corollary}
\begin{proof}
将 \(m=6\) 代入定理 \ref{thm:xi-time-part51aa-doubling-lucas-extension} 即得
\[
R_{14}\cong R_6\times \ZZ/(47\ZZ)=R_6\times \FF_{47}.
\]
另一方面，在定理 \ref{thm:xi-time-part51aa-derived-cross-fib-gcd-formula} 中取 \((m_1,m_2)=(6,14)\)，则
\[
g=\gcd(8,16)=8,
\qquad
F_g=F_8=21.
\]
于是
\[
H_0\!\left(R_6\overset{\mathbf L}{\otimes}_{\ZZ}R_{14}\right)
\cong
H_1\!\left(R_6\overset{\mathbf L}{\otimes}_{\ZZ}R_{14}\right)
\cong
\ZZ/(21\ZZ)
\cong
R_6.
\]
\leanverified{paper\_xi\_time\_part51aa\_window14\_strict\_direct\_factor\_lift}
证毕。
\end{proof}

\begin{proposition}[倍窗分裂的自然密度]\label{prop:xi-time-part51aa-doubling-split-density}
\leanverified{paper\_xi\_time\_part51aa\_doubling\_split\_density}
满足
\[
R_{2m+2}\cong R_m\times \ZZ/(L_{m+2}\ZZ)
\]
的分辨率 \(m\ge 1\) 构成一组自然密度为 \(2/3\) 的整数集合；其余密度 \(1/3\) 的分辨率恰对应唯一的 \(2\)-主非分裂支。
\end{proposition}
\begin{proof}
由定理 \ref{thm:xi-time-part51aa-doubling-lucas-extension}，分裂当且仅当
\[
3\nmid (m+2).
\]
而整数中不被 \(3\) 整除者的自然密度为 \(2/3\)。其补集自然密度为 \(1/3\)，并且正对应
\[
I_{m+2}/I_{m+2}^2\cong \FF_2
\]
的 \(2\)-主障碍支。证毕。
\end{proof}

\begin{theorem}[window-\texorpdfstring{$6$}{6} 的 \texorpdfstring{$B_3/C_3$}{B3/C3} 根字典不能内化为商环算术]\label{thm:xi-time-part51aa-window6-root-dictionary-externality}
\leanverified{paper\_xi\_time\_part51aa\_window6\_root\_dictionary\_externality}
沿用定理 \ref{thm:xi-window6-cyclic-kernel-bc3-root-datum} 的记号，并固定任意双射
\[
\beta:X_6\longrightarrow R_6\cong \ZZ/(21\ZZ).
\]
记
\[
A_{\mathrm{cyc}}:=\beta(X_6^{\mathrm{cyc}}),
\qquad
A_{\mathrm{ax}}:=\beta(X_{6,\mathrm{ax}}),
\qquad
A_{\mathrm{off}}:=\beta(X_{6,\mathrm{off}}).
\]
则这三个子集都不可能是：
\begin{enumerate}
  \item \(R_6\) 的理想；
  \item \(R_6\) 的加法子群之陪集；
  \item 任意环同态 \(R_6\to S\) 的像，尤其不可能是任何环端同态或自同态的像；
  \item 任意环同态 \(R_6\to S\) 的纤维。
\end{enumerate}
因此，window-\(6\) 的 \(B_3/C_3\) 根字典不可能由 \(R_6\) 的内部商环算术自行产生，而必然依赖于其外部的组合入射数据。
\end{theorem}
\begin{proof}
由定理 \ref{thm:xi-window6-cyclic-kernel-bc3-root-datum}，
\[
|X_6^{\mathrm{cyc}}|=18,
\qquad
|X_{6,\mathrm{ax}}|=6,
\qquad
|X_{6,\mathrm{off}}|=12.
\]
故
\[
|A_{\mathrm{cyc}}|=18,
\qquad
|A_{\mathrm{ax}}|=6,
\qquad
|A_{\mathrm{off}}|=12.
\]

另一方面，\(R_6\cong \ZZ/(21\ZZ)\) 的加法群是阶为 \(21\) 的循环群，因此其全部加法子群的基数只能是
\[
1,\ 3,\ 7,\ 21,
\]
陪集的基数亦同。理想作为特殊的加法子群，其基数当然也只能属于同一集合。

再设 \(\phi:R_6\to S\) 为任意环同态。则由第一同构定理，
\[
\operatorname{im}(\phi)\cong R_6/\ker(\phi),
\]
故
\[
|\operatorname{im}(\phi)|
=
\frac{|R_6|}{|\ker(\phi)|}
\]
必为 \(21\) 的约数；另一方面，每个纤维都是 \(\ker(\phi)\) 的加法陪集，因此其基数恰为 \(|\ker(\phi)|\)，也只能是
\[
1,\ 3,\ 7,\ 21.
\]

由于
\[
18,\ 6,\ 12\notin \{1,3,7,21\},
\]
上述三个集合不可能属于任一类。故 \(B_3/C_3\) 根字典不能被内化为 \(R_6\) 的商环算术结构。证毕。
\end{proof}

\begin{conjecture}[迭代倍窗塔的 \texorpdfstring{$2$}{2}-局部唯一障碍]\label{conj:xi-time-part51aa-doubling-tower-two-local-obstruction}
固定整数 \(\nu\ge 3\)，并定义
\[
S_k:=R_{2^k\nu-2}
\qquad(k\ge 0).
\]
则在整个倍窗塔中，每一步倍窗所新增的奇素数 Lucas 扇区在 pro-层面终将分裂为独立直因子；唯一能够跨层持续阻止完全分裂者，只可能是定理 \ref{thm:xi-time-part51aa-doubling-lucas-extension} 中出现的 \(2\)-主余切缺陷。等价地，倍窗机制的全部遗传性非半单障碍应当是纯 \(2\)-局部的。
\end{conjecture}

由此，最小统一分辨率的 lcm 型组织之外，还出现了与之对偶的 gcd 型导出交叉律；倍窗层不再表现为对旧模数的单纯放大，而是把 Lucas 伴随因子附加到既有 Fibonacci 核心之上；而 window-\(6\) 的 \(B_3/C_3\) 根字典则被严格排除在 \(R_6\) 的内部商环算术之外。因此，分辨率算术几何在固定模环分类之外，还自然携带跨窗口的导出几何与外部组合入射层。

\paragraph{Smith 局部--全局因子、稳定层中心相位与 window-\texorpdfstring{$6$}{6} 矩塔的共同有限化}
Smith 模核链、bin-fold 稳定层链与 window-\(6\) 有限尺度谱学在正文中已分别获得闭式终结。具体地，定理 \ref{thm:xi-smith-kernel-dirichlet-euler-stokes-tail-residue} 与定理 \ref{thm:xi-smith-normalized-kernel-positive-finite-euler-correction} 已把局部 Smith 模核级数压缩为有限前缀加几何尾部，并把归一化全局 Dirichlet 级数写成 \(\zeta(s)\) 乘有限 Euler 修正；定理 \ref{thm:xi-time-part60ab2-stable-spectrum-selfreciprocal-center} 与定理 \ref{thm:xi-time-part60ab2-stable-spectrum-fourier-phase-locking} 已给出 bin-fold 多重度谱的精确中心自反与 Fourier 相位锁定；命题 \ref{prop:xi-window6-complete-finite-scale-spectrum-closed-form} 连同定理 \ref{thm:xi-capacity-tail-histogram-moment-complete-statistic} 则表明 window-\(6\) 的全矩塔是一个三指数模型，其普通生成函数、最小三阶递推与直方图恢复都由有限数据完全决定。因而在这一闭链中，尚待补足的解析侧证书只剩 RH 边界在 Taylor 喷流层面的有限阶可审计性。

\begin{theorem}[核版 RH 边界的有限 Taylor 喷流证书]\label{thm:xi-time-part50be-rh-boundary-finite-taylor-jet-certificate}
\leanverified{paper\_xi\_time\_part50be\_rh\_boundary\_finite\_taylor\_jet\_certificate}
沿用推论 \ref{cor:real-input-40-collision-branch-radius}、命题 \ref{prop:real-input-40-collision-rhk-window} 与定理 \ref{thm:xi-time-part9pb-rh-transition-analytic-not-branch} 的记号。记
\[
P_2(\theta)=\log\rho(e^\theta),
\qquad
T_N(\theta):=\sum_{n=0}^{N}\frac{P_2^{(n)}(0)}{n!}\theta^n.
\]
取任意
\[
\theta_R<R<R_\theta,
\]
并定义
\[
M_R:=\max_{|z|=R}|P_2(z)|,
\qquad
q_R:=\frac{\theta_R}{R}\in(0,1).
\]
则对一切 \(N\ge 0\) 都有显式几何误差界
\[
\left|P_2(\theta_R)-T_N(\theta_R)\right|
\le
\frac{M_R}{1-q_R}q_R^{N+1},
\]
以及导数误差界
\[
\left|P_2'(\theta_R)-T_N'(\theta_R)\right|
\le
\frac{M_R}{R}\sum_{n=N+1}^{\infty}nq_R^{n-1}
=
\frac{M_R}{R}\,
\frac{(N+1)-Nq_R}{(1-q_R)^2}\,
q_R^{N}
\le
\frac{M_R}{R}\,
\frac{N+1}{(1-q_R)^2}\,
q_R^{N}.
\]
特别地，若
\[
\left|T_N'(\theta_R)\right|
>
\frac{M_R}{R}\,
\frac{N+1}{(1-q_R)^2}\,
q_R^{N},
\]
则 \(P_2'(\theta_R)\) 与 \(T_N'(\theta_R)\) 同号。因而核版 RH 临界边界 \(\theta_R\) 的函数值逼近与导数符号判定都可由有限阶 Taylor 喷流在几何速度下完成认证。
\end{theorem}
\begin{proof}
由定理 \ref{thm:xi-time-part9pb-rh-transition-analytic-not-branch}，
\[
0<\theta_R<R_\theta.
\]
故对任意 \(R\in(\theta_R,R_\theta)\)，函数 \(P_2\) 在闭圆盘 \(\{|z|\le R\}\) 上全纯且有界。由 Cauchy 估计，对所有 \(n\ge 0\) 有
\[
\left|\frac{P_2^{(n)}(0)}{n!}\right|
\le
\frac{M_R}{R^n}.
\]
于是
\[
\begin{aligned}
\left|P_2(\theta_R)-T_N(\theta_R)\right|
&\le
\sum_{n=N+1}^{\infty}
\left|\frac{P_2^{(n)}(0)}{n!}\right|
\theta_R^n\\
&\le
M_R\sum_{n=N+1}^{\infty}\left(\frac{\theta_R}{R}\right)^n
=
\frac{M_R}{1-q_R}q_R^{N+1},
\end{aligned}
\]
得到第一条。

对导数余项，同样有
\[
P_2'(\theta_R)-T_N'(\theta_R)
=
\sum_{n=N+1}^{\infty}
\frac{P_2^{(n)}(0)}{(n-1)!}\theta_R^{n-1},
\]
故
\[
\begin{aligned}
\left|P_2'(\theta_R)-T_N'(\theta_R)\right|
&\le
\sum_{n=N+1}^{\infty}
\frac{n!M_RR^{-n}}{(n-1)!}\theta_R^{n-1}\\
&=
\frac{M_R}{R}\sum_{n=N+1}^{\infty}nq_R^{n-1}.
\end{aligned}
\]
而
\[
\sum_{n=N+1}^{\infty}nq^{n-1}
=
\frac{(N+1)-Nq}{(1-q)^2}q^N
\qquad(0<q<1),
\]
代入 \(q=q_R\) 即得精确公式与上界。最后一条由
\[
\left|P_2'(\theta_R)-T_N'(\theta_R)\right|
<
\left|T_N'(\theta_R)\right|
\]
时符号不变的三角不等式立即推出。证毕。
\end{proof}

\noindent
由此，碰撞位势的 RH 临界边界不再仅以“解析域内部先于分歧发生”的存在性方式出现，而被进一步压缩为一条有限阶喷流证书；与前述 Smith 局部--全局 Euler 闭合、稳定层中心相位锁定以及 window-\(6\) 矩塔有限生成并置时，本文在算术、谱与热力学三侧的主要无限结构便都归结为显式有限数据上的可审计恢复问题。
在局部逆正规形之外，激活分支与其余根之间还存在一个对 \(u\approx 1\) 统一有效的正分离半径。

\begin{theorem}[激活分支在 \(u=1\) 邻域内的唯一小根性与严格尺度分离]\label{thm:xi-output-potential-activated-branch-uniform-scale-separation}
\leanverified{paper\_xi\_output\_potential\_activated\_branch\_uniform\_scale\_separation}
沿用定理 \ref{thm:xi-output-potential-activated-branch-rigidity} 的记号。则存在常数 \(\varepsilon>0\)、\(c>0\)，使得当
\[
|u-1|<\varepsilon
\]
时，方程
\[
G(\lambda,u)=0
\]
在圆盘
\[
|\lambda|<c
\]
内恰有一个根，且该根正是激活分支 \(\lambda_{\mathrm{act}}(u)\)。等价地，除 \(\lambda_{\mathrm{act}}(u)\) 之外的其余根全部满足
\[
|\lambda|\ge c.
\]
\end{theorem}
\begin{proof}
由命题 \ref{prop:real-input-40-activated-branch-linear}，
\[
G(0,1)=0,
\qquad
\partial_\lambda G(0,1)\neq 0.
\]
因此 \(\lambda=0\) 是 \(G(\lambda,1)\) 的简单根。由隐函数定理，存在 \(\varepsilon_0>0\) 与唯一解析分支 \(\lambda_{\mathrm{act}}(u)\) 定义在 \(|u-1|<\varepsilon_0\) 上，并满足
\[
\lambda_{\mathrm{act}}(1)=0,
\qquad
G(\lambda_{\mathrm{act}}(u),u)=0.
\]

记 \(G(\lambda,1)\) 的其余七个根组成的集合为 \(\Sigma_\ast\)。由于它们都非零，数量又有限，故
\[
2c:=\min_{\zeta\in\Sigma_\ast}|\zeta|>0.
\]
再由多项式根对系数的连续依赖性，缩小 \(\varepsilon_0\) 到某个 \(\varepsilon\in(0,\varepsilon_0)\) 后，可保证当 \(|u-1|<\varepsilon\) 时：
\begin{enumerate}
  \item 激活分支根 \(\lambda_{\mathrm{act}}(u)\) 仍满足 \(|\lambda_{\mathrm{act}}(u)|<c\)；
  \item 其余七个根都留在 \(|\lambda|\ge c\) 的区域内。
\end{enumerate}
于是 \(G(\lambda,u)=0\) 在 \(|\lambda|<c\) 内恰有一个根，且只能是 \(\lambda_{\mathrm{act}}(u)\)。证毕。
\end{proof}
